\documentclass[a4paper,11pt]{amsart}

% =========================================================================
%  Electromagnetismo I — Tema 2: Desarrollo multipolar
%  Apuntes formalizados. Compilar con: pdflatex + biber + pdflatex ×2
%  (o latexmk -pdf main.tex).  Todo el preámbulo está en preambulo.tex.
%  Los datos numéricos de las figuras «tipo Mathematica» se regeneran con
%  python3 scripts/generar_datos.py (ya se incluyen en figuras/datos/).
% =========================================================================
% =========================================================================
%  Electromagnetismo I — Tema 2: Desarrollo multipolar
%  Estilo: artículo científico minimalista, idéntico al de los temas 0 y 1
%  (amsart, Palatino, versalitas, enunciados numerados por sección), con la
%  paleta cambiada a magentas y rojizos claros sobre grises sobrios. La
%  notación vectorial es la de «Cálculo integral multivariable» (L. Navas) y
%  la de los temas anteriores; las expresiones son las del formulario.
% =========================================================================

% --- IDIOMA Y CODIFICACIÓN ------------------------------------------------
\usepackage[utf8]{inputenc}
\usepackage[T1]{fontenc}
\usepackage[spanish,es-tabla,es-nodecimaldot,es-noshorthands]{babel}

% --- TIPOGRAFÍA (Palatino, como en los temas 0 y 1) -------------------------
\usepackage{mathpazo}
\linespread{1.05}
\emergencystretch=2em
\usepackage[final]{microtype}

% --- MATEMÁTICAS -------------------------------------------------------------
\usepackage{amsmath, amssymb, mathtools}
\usepackage{bm}
\usepackage{mathrsfs}
\usepackage{esint}

% --- PÁGINA, COLORES Y GRÁFICOS --------------------------------------------
\usepackage[a4paper,margin=2.6cm]{geometry}
\usepackage{graphicx}
\usepackage{xcolor}
% Paleta base del Tema 2: magentas sobrios y rojizos claros
\definecolor{magprof}{HTML}{5B1A45}     % magenta muy oscuro (títulos, enunciados)
\definecolor{magacero}{HTML}{A0306E}    % magenta medio (enlaces, banda de portada)
\definecolor{magclaro}{HTML}{E2A6C4}    % magenta claro (marcos)
\definecolor{magfondo}{HTML}{FDF3F8}    % fondo tenue de definiciones
\definecolor{rosaclaro}{HTML}{F6D5DE}   % rosa pálido (rellenos)
\definecolor{rojoclaro}{HTML}{D9645A}   % rojizo claro (acento)
\definecolor{grispizarra}{HTML}{6B7280} % gris
\definecolor{grisclaro}{HTML}{D4D8DE}   % gris claro
\colorlet{acento}{rojoclaro}
% Paleta ampliada para campos escalares y vectoriales (magnitud creciente:
% índigo -> púrpura -> carmín -> naranja -> ámbar), la misma de los temas 0 y 1
\definecolor{cm0}{RGB}{38,30,112}
\definecolor{cm1}{RGB}{112,36,140}
\definecolor{cm2}{RGB}{190,52,118}
\definecolor{cm3}{RGB}{236,104,66}
\definecolor{cm4}{RGB}{252,184,42}
% Colores de apoyo de los diagramas geométricos (tonos sobrios)
\definecolor{rojoteo}{HTML}{C2415B}     % vectores destacados, resultados (carmín suave)
\definecolor{verdeteo}{HTML}{2F6F68}    % versores y normales (verde azulado, contraste)
\definecolor{pizarra}{HTML}{3F4C6B}     % cargas negativas, ejes secundarios
\colorlet{ocreteo}{rosaclaro}           % superficies y volúmenes (relleno)
% Colores de las figuras clonadas de Mathematica (se conservan los originales)
\definecolor{mmafondo}{HTML}{E6EEF8}    % fondo azul pálido de StreamPlot
\definecolor{mmamarco}{HTML}{5B7DB8}    % marco de la región
\definecolor{mmalinea}{HTML}{1E1E9A}    % líneas de campo
\definecolor{mmaequi}{HTML}{4A4A4A}     % equipotenciales superpuestas
\definecolor{mmarojo}{HTML}{C00000}     % vector dipolar y cargas
\definecolor{mmaazul}{HTML}{3E7CB1}     % serie 1 (Mathematica 97)
\definecolor{mmanaranja}{HTML}{E1A027}  % serie 2
\definecolor{mmaverde}{HTML}{8FA93C}    % serie 3
\definecolor{mmarojizo}{HTML}{E0644A}   % serie 4
\definecolor{mmavioleta}{HTML}{8A7DB8}  % serie 5

\usepackage{tikz}
\usepackage{tikz-3dplot}
\usetikzlibrary{babel, arrows.meta, calc, positioning, decorations.pathmorphing,
                decorations.markings, patterns, 3d, intersections, angles, quotes,
                fadings, shapes.geometric, spy}
\usepackage{pgfplots}
\pgfplotsset{compat=1.17}
\usepgfplotslibrary{colormaps, patchplots, fillbetween}
\pgfplotsset{colormap={campo}{
  rgb255=(38,30,112) rgb255=(112,36,140) rgb255=(190,52,118)
  rgb255=(236,104,66) rgb255=(252,184,42) rgb255=(255,226,110)}}
\pgfplotsset{colormap={magenta}{
  rgb255=(253,243,248) rgb255=(226,166,196) rgb255=(160,48,110) rgb255=(91,26,69)}}
\usepackage{booktabs}
\usepackage{needspace}
\usepackage{array}
\usepackage{multirow}
\usepackage{enumitem}
\setlist[itemize]{leftmargin=1.6em, itemsep=2pt, topsep=4pt}
\setlist[enumerate]{leftmargin=1.8em, itemsep=2pt, topsep=4pt}

% --- ENLACES Y REFERENCIAS CRUZADAS ----------------------------------------
\usepackage[colorlinks=true, linkcolor=magacero, citecolor=magacero,
            urlcolor=magacero, plainpages=false, pdfusetitle=false,
            pdftitle={Tema 2. Desarrollo multipolar},
            pdfauthor={Daniel Martín Guerrero}]{hyperref}
\usepackage[spanish,nameinlink]{cleveref}

% --- BIBLIOGRAFÍA -----------------------------------------------------------
\usepackage{csquotes}
\usepackage[backend=biber, style=numeric, sorting=none, giveninits=true,
            maxbibnames=6, minbibnames=3, doi=true, url=true,
            isbn=false]{biblatex}
\addbibresource{referencias.bib}
\DeclareFieldFormat[article,techreport,unpublished,book,inproceedings]{title}{#1\isdot}
\DeclareFieldFormat{url}{\url{#1}}
\renewcommand*{\bibfont}{\small}
\setlength{\bibitemsep}{3pt}

% --- CABECERAS ---------------------------------------------------------------
\pagestyle{headings}
\makeatletter
\def\ps@headings{%
  \let\@mkboth\markboth
  \def\@evenhead{\small\scshape\normalfont\thepage\hfil Electromagnetismo I}%
  \def\@oddhead{\small\scshape\normalfont Tema 2. Desarrollo multipolar\hfil\thepage}%
  \let\@evenfoot\@empty \let\@oddfoot\@empty}
\makeatother
\numberwithin{equation}{section}
% Subsecciones como títulos en bloque (no «run-in»), con aire antes y después
\usepackage{titlesec}
\titleformat{\subsection}[block]{\normalfont\bfseries\color{magprof}}{\thesubsection.}{0.6em}{}
\titlespacing*{\subsection}{0pt}{3ex plus 1ex minus .5ex}{1.6ex plus .4ex}
\setcounter{tocdepth}{2}
\makeatletter
\def\@@and{y}
\renewcommand{\keywordsname}{Palabras clave}
% Índice con sangrías fijas: el número de sección nunca se monta sobre el título
\def\l@section{\@tocline{1}{5pt plus 1pt}{0pt}{1.9em}{\bfseries\color{magprof}}}
\def\l@subsection{\@tocline{2}{0pt}{1.9em}{2.7em}{}}
\renewcommand{\tocsection}[3]{\indentlabel{\@ifnotempty{#2}{\makebox[1.9em][l]{#2.}}}#3}
\renewcommand{\tocsubsection}[3]{\indentlabel{\@ifnotempty{#2}{\makebox[2.7em][l]{#2.}}}#3}
\makeatother

% =========================================================================
%               NOTACIÓN Y MACROS DE USO FRECUENTE
% =========================================================================
% --- Diferencial, derivadas parciales y números -----------------------------
\newcommand{\dif}{\mathrm{d}}
\newcommand{\pd}[2]{\frac{\partial #1}{\partial #2}}
\newcommand{\pdd}[2]{\frac{\partial^{2} #1}{\partial #2^{2}}}
\newcommand{\pdm}[3]{\frac{\partial^{2} #1}{\partial #2\,\partial #3}}
\newcommand{\R}{\mathbb{R}}
\newcommand{\N}{\mathbb{N}}
\newcommand{\Z}{\mathbb{Z}}
\newcommand{\eul}{\mathrm{e}}

% --- Vectores, versores y operador nabla ---------------------------------
\newcommand{\vc}[1]{\vec{#1}}                       % vector con flecha
\newcommand{\uh}[1]{\hat{u}_{#1}}                   % versor û_#1
\newcommand{\ux}{\uh{x}}\newcommand{\uy}{\uh{y}}\newcommand{\uz}{\uh{z}}
\newcommand{\ur}{\uh{r}}\newcommand{\uth}{\uh{\theta}}\newcommand{\uph}{\uh{\varphi}}
\newcommand{\urho}{\uh{\rho}}\newcommand{\uR}{\uh{R}}\newcommand{\un}{\hat{n}}
\newcommand{\nab}{\vec{\nabla}}                     % ∇⃗
\newcommand{\nabp}{\vec{\nabla}'}                   % ∇⃗′ (coordenadas fuente)
\newcommand{\Lap}{\nabla^{2}}                       % laplaciano
\newcommand{\Lapp}{\nabla'^{2}}                     % laplaciano en coordenadas fuente
\newcommand{\norm}[1]{\lVert #1\rVert}
\newcommand{\abs}[1]{\lvert #1\rvert}
\newcommand{\prodesc}[2]{\langle #1,#2\rangle}
\newcommand{\rp}{\vc{r}\,'}                         % r⃗′ punto fuente
\newcommand{\dl}{\dif\vc{\ell}}                     % elemento de línea dirigido
\newcommand{\dS}{\dif\vc{S}}                        % elemento de área dirigido
\newcommand{\dsc}{\dif S}                           % elemento de área escalar
\newcommand{\dV}{\dif V}                            % elemento de volumen
\newcommand{\defeq}{\mathrel{\overset{\mathrm{def}}{=}}}
\newcommand{\bordeS}{\partial S}
\newcommand{\bordeV}{\partial V}
\DeclareMathOperator{\grad}{grad}
\DeclareMathOperator{\dive}{div}
\DeclareMathOperator{\rot}{rot}
\DeclareMathOperator{\sen}{sen}
\DeclareMathOperator{\senh}{senh}
\DeclareMathOperator{\sgn}{sgn}
\DeclareMathOperator{\Tr}{Tr}
\renewcommand{\sin}{\sen}                            % «sen» en todo el texto
% Tensor de Levi-Civita y delta de Kronecker
\newcommand{\eps}{\varepsilon}
\newcommand{\dpart}[1]{\partial_{#1}}

% --- Notación de electrostática (Tema 1) -------------------------------------
\newcommand{\eo}{\varepsilon_{0}}                    % permitividad del vacío
\newcommand{\qp}{q\,'}                              % carga fuente q'
\newcommand{\pot}{\phi}                              % potencial electrostático
\newcommand{\Qenc}{Q_{\mathrm{enc}}}                 % carga encerrada
\newcommand{\Ke}{\frac{1}{4\pi\eo}}                  % constante de Coulomb 1/(4 pi eps0)
\newcommand{\vE}{\vc{E}}
\newcommand{\vF}{\vc{F}}
\newcommand{\dlp}{\dif \ell\,'}                     % elemento de longitud fuente
\newcommand{\dSp}{\dif S\,'}                        % elemento de superficie fuente
\newcommand{\dVp}{\dif V\,'}                        % elemento de volumen fuente
\newcommand{\ent}[1]{\bigl\lvert #1\bigr\rvert}    % módulo grande

% --- Notación específica del Tema 2 (desarrollo multipolar) ------------------
\newcommand{\vp}{\vc{p}}                             % momento dipolar
\newcommand{\vl}{\vc{\ell}}                          % separación del dipolo (de -q a +q)
\newcommand{\vP}{\vc{P}}                             % polarización
\newcommand{\rhoP}{\rho_{P}}                         % densidad volúmica de polarización
\newcommand{\sigP}{\sigma_{P}}                       % densidad superficial de polarización
\newcommand{\Pl}[1]{P_{#1}}                          % polinomio de Legendre P_ℓ
\newcommand{\cosp}{\cos\theta'}                      % coseno del ángulo entre r⃗ y r⃗′
\newcommand{\Qt}{\mathsf{Q}}                         % tensor cuadrupolar (matriz)
\newcommand{\Qjk}{Q_{jk}}                            % componentes del tensor cuadrupolar
\newcommand{\cd}[1]{\ell_{#1}}                       % cosenos directores ℓ_x, ℓ_y, ℓ_z
\newcommand{\phimul}[1]{\phi_{#1}}                   % término ℓ-ésimo del desarrollo
\newcommand{\OOp}{\overrightarrow{OO'}}              % traslación del origen
\newcommand{\delt}[2]{\delta_{#1#2}}                 % delta de Kronecker
\newcommand{\ord}{\mathcal{O}}                       % símbolo de orden de Landau

% Estilos de cargas puntuales para los diagramas
\tikzset{
  cargapos/.style={circle, fill=rojoteo, draw=magprof, line width=0.4pt, inner sep=0pt,
                   minimum size=9pt, text=white, font=\scriptsize\bfseries},
  carganeg/.style={circle, fill=pizarra, draw=pizarra!70!black, line width=0.4pt, inner sep=0pt,
                   minimum size=9pt, text=white, font=\scriptsize\bfseries},
  cargapru/.style={circle, fill=white, draw=magprof, line width=0.6pt, inner sep=0pt,
                   minimum size=9pt, text=magprof, font=\scriptsize\bfseries},
}

% =========================================================================
%        ENTORNOS FORMALES (enunciados numerados, estilo artículo)
% =========================================================================
\usepackage{amsthm}
\usepackage[most]{tcolorbox}

% Enunciados en cursiva (resultados generales)
\newtheoremstyle{cienplano}{10pt}{10pt}{\itshape}{}%
  {\bfseries\color{magprof}}{.}{0.6em}%
  {\thmname{#1}\thmnumber{ #2}\thmnote{\normalfont\upshape\ (#3)}}
% Definiciones, cálculos y ejemplos (cuerpo en redonda)
\newtheoremstyle{ciendef}{10pt}{10pt}{\normalfont}{}%
  {\bfseries\color{magprof}}{.}{0.6em}%
  {\thmname{#1}\thmnumber{ #2}\thmnote{\normalfont\upshape\ (#3)}}
% Observaciones y notas (encabezado en cursiva)
\newtheoremstyle{cienobs}{8pt}{8pt}{\normalfont}{}%
  {\itshape\color{magprof}}{.}{0.6em}%
  {\thmname{#1}\thmnumber{ #2}\thmnote{\normalfont\upshape\ (#3)}}

\theoremstyle{cienplano}
\newtheorem{proposicion}{Proposición}[section]
\newtheorem{teorema}[proposicion]{Teorema}
\newtheorem{lema}[proposicion]{Lema}
\newtheorem{corolario}[proposicion]{Corolario}
\theoremstyle{ciendef}
\newtheorem{notacion}[proposicion]{Notación}
\newtheorem{ejemplo}[proposicion]{Ejemplo}
\newtheorem{ejercicio}[proposicion]{Ejercicio}
% Definiciones: único entorno enmarcado (como en los temas 0 y 1): caja fina con
% fondo tenue y título sobre el borde. Comparte contador con los enunciados.
\newcounter{definicion}
\renewcommand{\thedefinicion}{\theproposicion}
\newtcolorbox{definicionbox}[1]{%
  enhanced, breakable, sharp corners,
  colback=magfondo, colframe=magclaro, boxrule=0.4pt,
  left=9pt, right=9pt, top=5pt, bottom=5pt,
  before skip=12pt, after skip=12pt,
  attach boxed title to top left={yshift*=-\tcboxedtitleheight/2, xshift=8pt},
  boxed title style={colback=white, colframe=white, sharp corners, boxrule=0pt,
                     left=3pt, right=3pt, top=1pt, bottom=1pt},
  fonttitle=\bfseries\color{magprof}, coltitle=magprof,
  title={Definición~\thedefinicion\ifx\relax#1\relax\else\ (#1)\fi}}
\newenvironment{definicion}[1][]
  {\par\stepcounter{proposicion}\refstepcounter{definicion}\begin{definicionbox}{#1}}
  {\end{definicionbox}}

\theoremstyle{cienobs}
\newtheorem{observacion}[proposicion]{Observación}

% Resolución de ejercicios: misma maquinaria que «proof», con otro rótulo
\renewcommand{\qedsymbol}{\ensuremath{\blacksquare}}
\newenvironment{resolucion}[1][Resolución]
  {\begin{proof}[\itshape\color{magprof}#1]}{\end{proof}}

\crefname{proposicion}{proposición}{proposiciones}
\Crefname{proposicion}{Proposición}{Proposiciones}
\crefname{teorema}{teorema}{teoremas}
\Crefname{teorema}{Teorema}{Teoremas}
\crefname{lema}{lema}{lemas}
\Crefname{lema}{Lema}{Lemas}
\crefname{corolario}{corolario}{corolarios}
\Crefname{corolario}{Corolario}{Corolarios}
\crefname{definicion}{definición}{definiciones}
\Crefname{definicion}{Definición}{Definiciones}
\crefname{notacion}{notación}{notaciones}
\Crefname{notacion}{Notación}{Notaciones}
\crefname{ejemplo}{ejemplo}{ejemplos}
\Crefname{ejemplo}{Ejemplo}{Ejemplos}
\crefname{ejercicio}{ejercicio}{ejercicios}
\Crefname{ejercicio}{Ejercicio}{Ejercicios}
\crefname{observacion}{observación}{observaciones}
\Crefname{observacion}{Observación}{Observaciones}
\crefname{section}{sección}{secciones}
\Crefname{section}{Sección}{Secciones}
\crefname{subsection}{sección}{secciones}
\Crefname{subsection}{Sección}{Secciones}
\crefname{table}{tabla}{tablas}
\Crefname{table}{Tabla}{Tablas}
\crefname{figure}{figura}{figuras}
\Crefname{figure}{Figura}{Figuras}
\crefname{equation}{ecuación}{ecuaciones}
\Crefname{equation}{Ecuación}{Ecuaciones}
\crefformat{equation}{ecuación~(#2#1#3)}
\Crefformat{equation}{Ecuación~(#2#1#3)}
\crefrangeformat{equation}{ecuaciones~(#3#1#4)--(#5#2#6)}
\crefmultiformat{equation}{ecuaciones~(#2#1#3)}{ y~(#2#1#3)}{, (#2#1#3)}{ y~(#2#1#3)}

% Referencia a un panel de figura: «figura 5\,(b)», enlazada en bloque
\newcommand{\figpanel}[2]{\hyperref[#1]{figura~\ref*{#1}\,\textup{(#2)}}}
\newcommand{\Figpanel}[2]{\hyperref[#1]{Figura~\ref*{#1}\,\textup{(#2)}}}

% Pies de figura y tabla, en la línea de los temas anteriores
\usepackage[font=small, labelfont={sc,color=magprof}, labelsep=period,
            margin=0.6cm, justification=justified]{caption}

% =========================================================================
%        ESTILOS GRÁFICOS COMUNES (TikZ / pgfplots)
% =========================================================================
% =========================================================================
%  Estilos y macros gráficas comunes del Tema 2
%  - \mapcolor{t}: escala cromática «campo» (t en [0,1]); define el color mapc
%  - Estilos TikZ heredados de los temas 0 y 1 (ejes, vectores, versores...)
%  - \panelMMA: marco de las figuras clonadas de Mathematica (StreamPlot,
%    ContourPlot y VectorPlot) que dibuja los datos de figuras/datos/*.tex
% =========================================================================

% --- Escala cromática (5 tramos entre cm0 y cm4) ------------------------
\newcommand{\mapcolor}[1]{%
  \pgfmathsetmacro{\mt}{min(max(#1,0),0.9999)*4}%
  \pgfmathtruncatemacro{\mi}{floor(\mt)}%
  \pgfmathtruncatemacro{\mj}{\mi+1}%
  \pgfmathsetmacro{\mf}{100*(\mt-\mi)}%
  \colorlet{mapc}{cm\mj!\mf!cm\mi}}

% --- Estilos TikZ generales ------------------------------------------------
\tikzset{
  >=Stealth,
  eje/.style={->, line width=0.6pt, grispizarra},
  ejeprof/.style={->, line width=0.7pt, magprof},
  vect/.style={->, line width=1pt, magprof},
  vectrojo/.style={->, line width=1pt, rojoteo},
  vectverde/.style={->, line width=1pt, verdeteo},
  fino/.style={line width=0.4pt, grispizarra},
  punteado/.style={densely dotted, line width=0.5pt, grispizarra},
  etiq/.style={font=\small, text=black!85},
  cuadro/.style={draw=grispizarra!60, line width=0.4pt, fill=white, inner sep=2pt},
  trayecto/.style={line width=1.1pt, rojoteo},
  orienta/.style={decoration={markings, mark=at position 0.55 with {\arrow{Stealth[length=5pt]}}}, postaction={decorate}},
  % líneas de campo y equipotenciales en la paleta del documento
  lcampo/.style={magacero, line width=0.55pt, postaction={decorate},
                 decoration={markings, mark=at position 0.5 with {\arrow{Stealth[length=3.6pt]}}}},
  lcampoinv/.style={magacero, line width=0.55pt, postaction={decorate},
                 decoration={markings, mark=at position 0.5 with {\arrowreversed{Stealth[length=3.6pt]}}}},
  lequi/.style={pizarra!70, line width=0.45pt, densely dashed},
}

% --- Estilo base de ejes pgfplots (como en los temas anteriores) -----------
\pgfplotsset{
  tema/.style={
    tick label style={font=\scriptsize}, label style={font=\small},
    axis lines=left, axis line style={grispizarra},
    grid=major, grid style={grisclaro, line width=0.4pt},
    every tick/.style={grispizarra},
    legend style={font=\scriptsize, draw=none, fill=white, fill opacity=0.85,
                  text opacity=1},
    /pgf/number format/use comma,
  },
  % Ejes con marco completo, como los «Plot» de Mathematica de las figuras del tema
  marcoMMA/.style={
    tick label style={font=\scriptsize}, label style={font=\small},
    axis line style={grispizarra!80, line width=0.5pt},
    every tick/.style={grispizarra}, tick align=inside,
    grid=major, grid style={grisclaro!70, line width=0.3pt},
    /pgf/number format/use comma,
    legend style={font=\scriptsize, draw=none, fill=white, fill opacity=0.85,
                  text opacity=1},
  },
  mapa/.style={
    tick label style={font=\scriptsize}, label style={font=\small},
    axis line style={grispizarra}, every tick/.style={grispizarra},
    /pgf/number format/use comma,
    colormap name=campo,
  },
}

% =========================================================================
%  Paneles «tipo Mathematica» (ventana [-L,L]^2 en unidades de \ell)
%  Uso dentro de un tikzpicture con x=y=<escala>:
%    \panelMMA{L}{fondo}{radio del disco central}{contenido}
%  El contenido se recorta a la ventana; luego se dibujan el disco blanco
%  (región excluida, donde están las fuentes) y el marco con marcas.
% =========================================================================
% Estilos de los objetos que escribe scripts/generar_datos.py
\tikzset{
  mmaLC/.style={mmalinea, line width=0.42pt},
  mmaEQ/.style={mmaequi, line width=0.28pt},
  mmaEQsolo/.style={black, line width=0.4pt},
  mmaFL/.style={-{Stealth[length=3.3pt, width=2.7pt]}, mmalinea, line width=0.42pt},
}
% Macros por defecto (cada figura puede redefinirlas antes de \input)
\newcommand{\EQ}[1]{\draw[mmaEQ] plot coordinates {#1};}
\newcommand{\LC}[1]{\draw[mmaLC] plot coordinates {#1};}
\newcommand{\FL}[3]{\draw[mmaFL] (#1,#2) -- ++(#3:0.02);}
\newcommand{\VF}[5]{%
  \mapcolor{#5}%
  \pgfmathsetmacro{\hl}{2.0+2.4*#5}\pgfmathsetmacro{\lw}{0.32+0.45*#5}%
  \draw[-{Stealth[length=\hl pt,width=0.85*\hl pt]}, line width=\lw pt, mapc]
     ({#1-0.5*#3},{#2-0.5*#4}) -- ++(#3,#4);}

% Marco con marcas en -L, -L/2, 0, L/2, L y rótulos de los ejes
% \marcoMMA{L}{paso de marcas}{rótulo x}{rótulo y}
\colorlet{marcoc}{mmamarco}
\newcommand{\marcoMMA}[4]{%
  \draw[marcoc, line width=0.55pt] (-#1,-#1) rectangle (#1,#1);
  \foreach \t in {-#1,-#2,0,#2,#1}{%
    \draw[grispizarra, line width=0.35pt] (\t,-#1) -- ++(0,0.025*#1*2);
    \draw[grispizarra, line width=0.35pt] (-#1,\t) -- ++(0.025*#1*2,0);
    \draw[grispizarra, line width=0.35pt] (\t,#1) -- ++(0,-0.025*#1*2);
    \draw[grispizarra, line width=0.35pt] (#1,\t) -- ++(-0.025*#1*2,0);
    \node[below=1pt, font=\tiny, inner sep=1pt] at (\t,-#1) {$\pgfmathprintnumber{\t}$};
    \node[left=1pt, font=\tiny, inner sep=1pt] at (-#1,\t) {$\pgfmathprintnumber{\t}$};}
  \node[font=\small, anchor=north] at (0,-#1) [yshift=-9pt] {#3};
  \node[font=\small, anchor=south, rotate=90] at (-#1,0) [xshift=0pt, yshift=12pt] {#4};}

% \panelMMA{L}{color de fondo}{r0}{contenido}
\newcommand{\panelMMA}[4]{%
  \fill[#2] (-#1,-#1) rectangle (#1,#1);
  \begin{scope}\clip (-#1,-#1) rectangle (#1,#1); #4 \end{scope}
  \ifdim #3pt>0pt
    \fill[white] (0,0) circle (#3);
    \draw[mmamarco!80, line width=0.4pt] (0,0) circle (#3);
  \fi}

% Vector dipolar rojo en el centro de un panel (altura ±h)
\newcommand{\dipoloRojo}[1]{%
  \draw[-{Stealth[length=4pt,width=3.6pt]}, mmarojo, line width=1.3pt] (0,-#1) -- (0,#1);}

% =========================================================================
%  Primitivas 3D para tdplot (vista fija: tdplotsetmaincoords{70}{125}),
%  heredadas del Tema 1.
% =========================================================================
\def\VA{125}
% \tdCirc[opciones]{radio}{z}   (mitad frontal continua, trasera a trazos)
\newcommand{\tdCirc}[3][]{%
  \draw[#1, domain=\VA-180:\VA, samples=36, variable=\tdt]
    plot ({#2*cos(\tdt)},{#2*sin(\tdt)},{#3});
  \draw[#1, densely dashed, domain=\VA:\VA+180, samples=36, variable=\tdt]
    plot ({#2*cos(\tdt)},{#2*sin(\tdt)},{#3});}
% \tdEsfera{radio}{color}  (esfera sombreada con ecuador)
\newcommand{\tdEsfera}[2]{%
  \begin{scope}[tdplot_screen_coords]
    \shade[ball color=rosaclaro!80, opacity=0.45] (0,0) circle (#1);
    \draw[#2, line width=0.7pt] (0,0) circle (#1);
  \end{scope}
  \tdCirc[#2!60, line width=0.4pt]{#1}{0}}

% =========================================================================
%  Región de carga V' para los diagramas geométricos (secciones 3 y 4):
%  contorno suave, separado de los ejes, con sombreado radial tenue.
%  \regionV{desplazamiento}  (el contorno está centrado cerca de (2.25,2.7))
% =========================================================================
\def\contornoV{plot[smooth cycle, tension=0.75] coordinates
  {(0.95,2.5) (1.35,3.5) (2.45,3.95) (3.55,3.45) (3.6,2.25) (2.6,1.55) (1.45,1.65)}}
\newcommand{\regionV}{%
  \shade[inner color=rosaclaro!25, outer color=rosaclaro!90] \contornoV;
  \draw[magacero, line width=0.75pt] \contornoV;}
% Ejes cartesianos en perspectiva caballera con el origen marcado
% \ejesO{long z}{long y}{long x}
\newcommand{\ejesO}[3]{%
  \draw[eje] (0,0) -- (0,#1) node[above, etiq] {$z$};
  \draw[eje] (0,0) -- (#2,0) node[right, etiq] {$y$};
  \draw[eje] (0,0) -- ({-0.76*#3},{-0.65*#3}) node[below left=-2pt, etiq] {$x$};}
% Elemento de volumen (cubito en perspectiva) centrado en #1
\newcommand{\cubito}[1]{%
  \draw[magprof, line width=0.5pt, fill=white] ($(#1)+(-0.12,-0.12)$) rectangle ($(#1)+(0.12,0.12)$);
  \draw[magprof, line width=0.5pt, fill=rosaclaro!40] ($(#1)+(-0.12,0.12)$) -- ++(0.07,0.07) -- ++(0.24,0) -- ++(-0.07,-0.07) -- cycle;
  \draw[magprof, line width=0.5pt, fill=rosaclaro!70] ($(#1)+(0.12,0.12)$) -- ++(0.07,0.07) -- ++(0,-0.24) -- ++(-0.07,-0.07) -- cycle;}



% =========================================================================
%                     VARIABLES PARA LA PORTADA
% =========================================================================
\newcommand{\Universidad}{Universidad de Salamanca}
\newcommand{\Facultad}{Facultad de Ciencias}
\newcommand{\Titulacion}{Grado en Física}
\newcommand{\Asignatura}{Electromagnetismo I}
\newcommand{\NumTema}{Tema 2}
\newcommand{\TituloTema}{Desarrollo\\multipolar}
\newcommand{\tipotrabajo}{Apuntes formalizados con seminarios prácticos}
\newcommand{\miNombre}{Daniel Martín Guerrero}
\newcommand{\miProfesor}{Luis López Díaz}
\newcommand{\SubtituloTema}{Dipolo eléctrico, polarización, momentos multipolares\\y estimación del error}
\newcommand{\miUbicacion}{Salamanca, curso 2026/2027}

% =========================================================================
%                     DATOS DEL ARTÍCULO (cabecera)
% =========================================================================
\title{Desarrollo multipolar}
\author{Daniel Martín Guerrero}
\address{Grado en Física, Facultad de Ciencias, Universidad de Salamanca,
         Salamanca, España}
\date{}

% =========================================================================
%                                DOCUMENTO
% =========================================================================
\begin{document}
% --- PORTADA ----------------------------------------------------------------
\pagenumbering{alph}
%%%%%%%%%%%%%%%%%%%%%%%%%%%%%%%%%%%%%%%%%%%%%%%%%%%%%%%%%%%%%%%%%%%%%%%%
% Portada formal (estilo tesis), idéntica en estructura a la de los temas
% 0 y 1: logotipos, identificación institucional, título y autoría, con la
% banda vertical a la derecha. Novedad del Tema 2: paleta magenta y un
% emblema minimalista con las líneas de campo de un dipolo, r = k sen^2(θ).
% Requiere en main.tex las macros \Universidad, \Facultad, \Titulacion,
% \Asignatura, \NumTema, \TituloTema, \SubtituloTema, \tipotrabajo,
% \miNombre, \miProfesor, \miUbicacion.
%%%%%%%%%%%%%%%%%%%%%%%%%%%%%%%%%%%%%%%%%%%%%%%%%%%%%%%%%%%%%%%%%%%%%%%%

\begin{titlepage}
\thispagestyle{empty}
\pagecolor{white}

\tikzexternaldisable\begin{tikzpicture}[remember picture, overlay]
  \coordinate (O) at (current page.south west);

  % Velo muy tenue en la mitad superior (cabecera institucional)
  \shade[top color=magfondo, bottom color=white]
    ($(O)+(0,18.5cm)$) rectangle ($(O)+(21cm,29.7cm)$);

  % Banda vertical magenta, a la derecha del bloque de texto
  \fill[magacero] ($(O)+(19.0cm,2.6cm)$) rectangle ($(O)+(19.14cm,27.4cm)$);

  % Logotipos
  \node[anchor=north west, inner sep=0] at ($(O)+(2.4cm,27.7cm)$)
    {\includegraphics[height=1.55cm]{logo_usal.png}};
  \node[anchor=north east, inner sep=0] at ($(O)+(17.4cm,27.7cm)$)
    {\includegraphics[height=1.55cm]{logo_facultad.png}};
  \draw[grispizarra!60, line width=0.4pt] ($(O)+(2.4cm,25.7cm)$) -- ($(O)+(17.4cm,25.7cm)$);

  % Identificación institucional
  \node[anchor=north, align=center, inner sep=0, text width=15cm] at ($(O)+(9.9cm,25.2cm)$)
    {{\color{magprof}\scshape\Large\textls[120]{\Universidad}}\\[5pt]
     {\color{grispizarra}\scshape\large\textls[80]{\Facultad}}\\[3pt]
     {\color{grispizarra}\scshape\large\textls[80]{\Titulacion}}};

  % Asignatura y título
  \node[anchor=north, align=center, inner sep=0, text width=15cm] at ($(O)+(9.9cm,19.6cm)$)
    {{\color{magacero}\scshape\Large\textls[160]{\Asignatura}}\\[4pt]
     {\color{grispizarra}\scshape\large\textls[100]{\NumTema}}\\[22pt]
     {\color{magprof}\fontsize{34}{40}\selectfont\bfseries\TituloTema}\\[14pt]
     {\color{magacero}\rule{3.2cm}{0.8pt}}\\[14pt]
     {\color{black!80}\itshape\large\SubtituloTema}};

  % Emblema: líneas de campo de un dipolo puntual, r = k sen^2(θ)
  \begin{scope}[shift={($(O)+(9.9cm,11.05cm)$)}, x=0.95cm, y=0.95cm]
    \draw[magclaro!70, line width=0.35pt] (0,-1.15) -- (0,1.15);
    \foreach \k in {0.55,0.95,1.5,2.3,3.4}{%
      \foreach \s in {1,-1}{%
        \draw[magclaro, line width=0.5pt, domain=4:176, samples=90, variable=\t]
          plot ({\s*\k*sin(\t)^3},{\k*sin(\t)^2*cos(\t)});}}
    \fill[white] (0,0) circle (0.24);
    \draw[-{Stealth[length=4pt,width=3.4pt]}, magacero, line width=1.1pt] (0,-0.2) -- (0,0.24);
  \end{scope}

  % Epígrafe: las expresiones centrales del tema
  \node[anchor=north, align=center, inner sep=0, text width=15cm] at ($(O)+(9.9cm,9.35cm)$)
    {\color{grispizarra}\small
     $\displaystyle\phi(\vc{r})=\frac{1}{4\pi\eo}\sum_{\ell=0}^{\infty}\frac{1}{r^{\ell+1}}
        \int_{V'}r'^{\,\ell}\rho(\rp)\,\Pl{\ell}(\cosp)\,\dVp\qquad
      \vp=\int_{V'}\rp\rho(\rp)\,\dVp\qquad
      \frac{\phi_{n}}{\phi_{0}}<\Bigl(\frac{a}{r}\Bigr)^{n}$};

  % Autoría y dirección
  \draw[grispizarra!60, line width=0.4pt] ($(O)+(2.4cm,7.6cm)$) -- ($(O)+(17.4cm,7.6cm)$);
  \node[anchor=north west, inner sep=0, text width=7cm, align=left] at ($(O)+(2.4cm,7.2cm)$)
    {{\scshape\small\color{magacero}\textls[80]{Autor}}\\[4pt]
     {\large\color{black!85}\miNombre}\\[2pt]
     {\small\color{grispizarra}Grado en Física}};
  \node[anchor=north west, inner sep=0, text width=7cm, align=left] at ($(O)+(10.4cm,7.2cm)$)
    {{\scshape\small\color{magacero}\textls[80]{Profesor}}\\[4pt]
     {\large\color{black!85}\miProfesor}\\[2pt]
     {\small\color{grispizarra}\Asignatura}};

  % Tipo de documento, lugar y curso
  \node[anchor=south, align=center, inner sep=0, text width=15cm] at ($(O)+(9.9cm,2.9cm)$)
    {{\color{magprof}\scshape\large\textls[100]{\tipotrabajo}}\\[4pt]
     {\color{grispizarra}\scshape\small\textls[80]{\miUbicacion}}};
\end{tikzpicture}\tikzexternalenable

\end{titlepage}

\clearpage
\pagenumbering{arabic}
\setcounter{page}{1}

\begin{abstract}
El potencial de una distribución de carga arbitraria es, en principio, una integral que rara vez
sabemos resolver de forma cerrada. En este tema construimos su aproximación sistemática a grandes
distancias. Empezamos por el sistema más sencillo con carga neta nula, el dipolo eléctrico, cuyos
potencial y campo obtenemos por desarrollo de Taylor y expresamos en forma intrínseca; de él
nacen la polarización y las densidades de carga ligada. Después desarrollamos $1/R$ en serie de
polinomios de Legendre y obtenemos el desarrollo multipolar del potencial, cuyos tres primeros
términos (monopolar, dipolar y cuadrupolar) definen la carga neta, el momento dipolar y el tensor
cuadrupolar. Estudiamos cómo dependen del origen de coordenadas, las simetrías que reducen el número
de componentes independientes y, por último, cómo acotar el error cometido al truncar la serie.
Cada resultado se enuncia con rigor, se demuestra y se acompaña de gráficas; los seminarios finales
resuelven, con todo detalle, los problemas de la lista del tema.
\end{abstract}

\maketitle
\tableofcontents

% =========================================================================
\section{Introducción: el problema de calcular el potencial}\label{sec:intro}
% =========================================================================
En el tema anterior~\cite{TEMA1} obtuvimos las leyes que permiten calcular el campo y el potencial creados por
\emph{cualquier} distribución de carga en \emph{cualquier} punto del espacio. Conviene tenerlas a la vista,
porque este tema no añade ninguna ley nueva: se ocupa de cómo \emph{evaluar} en la práctica una de ellas.

% -------------------------------------------------------------------------
\subsection{Las tres descripciones de la electrostática}\label{sec:triangulo}
% -------------------------------------------------------------------------
La densidad de carga $\rho$, el potencial $\phi$ y el campo $\vE$ son tres descripciones equivalentes de un mismo problema
electrostático: conocida una de ellas, las otras dos quedan determinadas (\cref{fig:triangulo}).

\begin{proposicion}[Relaciones fundamentales de la electrostática~\cite{TEMA1}]\label{prop:triangulo}
Sea $\rho(\rp)$ una distribución de carga localizada en un volumen $V'$, y tomemos el origen de potencial en el infinito. Entonces,
para todo punto campo $\vc{r}$,
\begin{align}
  \phi(\vc{r})&=\frac{1}{4\pi\eo}\int_{V'}\frac{\rho(\rp)}{R}\,\dVp, &
  \vE(\vc{r})&=\frac{1}{4\pi\eo}\int_{V'}\frac{\rho(\rp)}{R^{3}}\,\vc{R}\,\dVp,\label{eq:integrales}\\[2pt]
  \vE(\vc{r})&=-\nab\phi(\vc{r}), &
  \phi(\vc{r})-\phi(\vc{r}_0)&=-\int_{\vc{r}_0}^{\vc{r}}\vE\cdot\dl,\label{eq:E-phi}\\[2pt]
  \nab\cdot\vE&=\frac{\rho}{\eo}, &
  \Lap\phi&=-\frac{\rho}{\eo},\label{eq:poisson}
\end{align}
donde $\vc{R}=\vc{r}-\rp$ y $R=\norm{\vc{R}}$.
\end{proposicion}

% =========================================================================
% Figura: el «triángulo» de la electrostática (apuntes manuscritos, pág. 1).
% Vértices rho, phi, E; en cada lado, la relación que permite pasar de una
% descripción a otra (integración directa, derivación o ecuación diferencial).
% Las flechas de ida van por fuera del triángulo y las de vuelta por dentro,
% con los rótulos en el lado correspondiente para que no se superpongan.
% =========================================================================
\begin{figure}[ht]
\centering
\begin{tikzpicture}[x=1cm, y=1cm,
    vert/.style={draw=magacero, line width=0.6pt, fill=magfondo, rounded corners=2pt,
                 inner sep=5pt, minimum width=1.45cm, minimum height=0.75cm, font=\normalsize},
    rel/.style={font=\footnotesize, text=black!85, inner sep=1.5pt},
    ida/.style={-{Stealth[length=5pt,width=4pt]}, line width=0.7pt, magacero, shorten >=17pt, shorten <=17pt},
    vuelta/.style={-{Stealth[length=5pt,width=4pt]}, line width=0.7pt, grispizarra, shorten >=17pt, shorten <=17pt}]
  \coordinate (A) at (0,4.2);   % rho
  \coordinate (B) at (-3.8,0);  % phi
  \coordinate (C) at (3.8,0);   % E
  % normales unitarias exteriores de los lados izquierdo y derecho
  \coordinate (nI) at (-0.742,0.671);
  \coordinate (nD) at (0.742,0.671);
  % lado izquierdo: rho -> phi (por fuera), phi -> rho (por dentro)
  \draw[ida]    ($(A)+0.13*(nI)$) -- ($(B)+0.13*(nI)$);
  \draw[vuelta] ($(B)-0.13*(nI)$) -- ($(A)-0.13*(nI)$);
  \node[rel, anchor=south east] at ($(A)!0.5!(B)+0.25*(nI)$)
       {$\displaystyle\phi=\frac{1}{4\pi\eo}\int_{V'}\frac{\rho}{R}\,\dVp$};
  \node[rel, anchor=north west] at ($(A)!0.57!(B)-0.2*(nI)$) {$\Lap\phi=-\dfrac{\rho}{\eo}$};
  % lado derecho: rho -> E (por fuera), E -> rho (por dentro)
  \draw[ida]    ($(A)+0.13*(nD)$) -- ($(C)+0.13*(nD)$);
  \draw[vuelta] ($(C)-0.13*(nD)$) -- ($(A)-0.13*(nD)$);
  \node[rel, anchor=south west] at ($(A)!0.5!(C)+0.25*(nD)$)
       {$\displaystyle\vE=\frac{1}{4\pi\eo}\int_{V'}\frac{\rho\,\vc{R}}{R^{3}}\,\dVp$};
  \node[rel, anchor=north east] at ($(A)!0.57!(C)-0.2*(nD)$) {$\nab\cdot\vE=\dfrac{\rho}{\eo}$};
  % base: phi -> E (derivación) y E -> phi (integral de línea)
  \draw[ida]    ($(B)+(0,0.13)$) -- ($(C)+(0,0.13)$);
  \draw[vuelta] ($(C)-(0,0.13)$) -- ($(B)-(0,0.13)$);
  \node[rel, anchor=south] at (0,0.22) {$\vE=-\nab\phi$};
  \node[rel, anchor=north] at (0,-0.22) {$\displaystyle\phi(\vc{r})-\phi(\vc{r}_0)=-\int_{\vc{r}_0}^{\vc{r}}\vE\cdot\dl$};
  % vértices
  \node[vert] at (A) {$\rho(\rp)$};
  \node[vert] at (B) {$\phi(\vc{r})$};
  \node[vert] at (C) {$\vE(\vc{r})$};
\end{tikzpicture}
\caption{Las tres descripciones equivalentes de un problema electrostático (\cref{prop:triangulo}). Las flechas magenta (exteriores) se recorren
\emph{integrando} sobre las fuentes o \emph{derivando}; las grises (interiores), resolviendo una ecuación diferencial o integrando a lo largo de un
camino. El desarrollo multipolar aproxima la flecha $\rho\to\phi$, que es la más difícil de recorrer en la práctica.}
\label{fig:triangulo}
\end{figure}


\begin{observacion}[Por qué preferimos el potencial]\label{obs:potencial-escalar}
De las dos integrales de \cref{eq:integrales}, la del potencial es siempre la más sencilla: su integrando es escalar (un número por
punto en lugar de tres) y decae como $1/R$ en lugar de como $1/R^{2}$. Por eso, en todo el tema calcularemos primero el potencial y
obtendremos después el campo derivando, $\vE=-\nab\phi$.
\end{observacion}

% -------------------------------------------------------------------------
\subsection{Planteamiento: una aproximación a grandes distancias}\label{sec:planteamiento}
% -------------------------------------------------------------------------
Aun así, la integral
\begin{equation}\label{eq:pot-integral}
  \phi(\vc{r})=\frac{1}{4\pi\eo}\int_{V'}\frac{\rho(\rp)}{\norm{\vc{r}-\rp}}\,\dVp
\end{equation}
es, en general, \emph{difícil de calcular en la práctica}: solo para distribuciones muy simétricas (esferas, cilindros, planos) admite una
expresión cerrada, y en los casos más generales ni siquiera conocemos con detalle $\rho$. Lo que sí sabemos es que si nos alejamos
suficientemente de la distribución, sus detalles dejan de importar y toda ella se ve como una \emph{carga puntual}. Entre ambos extremos
(conocer $\rho$ punto a punto o verla como un punto) hay un camino intermedio: caracterizar la distribución mediante unas pocas
cantidades, que dependen solo de \emph{algunos} detalles de la misma, y que bastan para obtener el potencial con la precisión deseada.

\begin{definicion}[Momentos multipolares]\label{def:multipolos}
Llamamos \emph{momentos multipolares} (o \emph{multipolos eléctricos}) de una distribución de carga a las cantidades
\[
  \int_{V'}\rho(\rp)\,\dVp,\qquad \int_{V'}\rho(\rp)\,\rp\,\dVp,\qquad \int_{V'}\rho(\rp)\bigl(3x'_jx'_k-\delt{j}{k}r'^{2}\bigr)\dVp,\ \dots
\]
(carga neta, momento dipolar, tensor cuadrupolar, \dots), que solo dependen de la distribución y del origen elegido, pero \emph{no} del
punto campo. El \emph{desarrollo multipolar} es la expresión del potencial como una serie cuyos términos son productos de un momento
multipolar por una función sencilla del punto campo, y que decaen cada vez más deprisa con la distancia.
\end{definicion}

El desarrollo permite así sustituir una integral distinta para cada punto campo por unas pocas integrales, calculadas una sola vez, y
\emph{sacar} las coordenadas del punto campo fuera de la integral. Construiremos el desarrollo en dos etapas. Primero estudiaremos el sistema más
sencillo de carga neta nula, el \emph{dipolo eléctrico} (\cref{sec:dipolo}), y su generalización a distribuciones continuas de dipolos, la
\emph{polarización} (\cref{sec:polarizacion}). Después desarrollaremos $1/R$ en serie de potencias de $r'/r$ y obtendremos el desarrollo
completo (\cref{sec:multipolar}), cuyo error de truncamiento acotaremos en la \cref{sec:error}. Los seminarios (\cref{sec:seminarios}) aplican
todo ello a los problemas de la lista del tema~\cite{PROBLEMAS-T2}. El desarrollo sigue los apuntes de teoría de la asignatura~\cite{APUNTES-T2}; como
bibliografía de apoyo, el dipolo eléctrico se trata en~\cite[cap.~3]{GRIFFITHS}, \cite[cap.~8]{WANGSNESS} y~\cite[cap.~4]{VICTORIANO}, y las expresiones multipolares en las
mismas referencias y, con más profundidad, en~\cite{JACKSON}.

% -------------------------------------------------------------------------
\subsection{Notación y convenios del tema}\label{sec:notacion}
% -------------------------------------------------------------------------
\begin{notacion}[Puntos campo y fuente]\label{not:puntos}
Seguimos la notación de los temas anteriores~\cite{TEMA0,TEMA1}, que a su vez adopta la de~\cite{NAVAS}:
\begin{itemize}
  \item $\vc{r}=(x,y,z)$ es el \emph{punto campo}, donde se evalúan $\phi$ y $\vE$, con coordenadas esféricas $(r,\theta,\varphi)$ y
        versores $\{\ur,\uth,\uph\}$; las magnitudes primadas, $\rp=(x',y',z')$, $r'=\norm{\rp}$, $\dVp$, $\nabp$, se refieren
        al \emph{punto fuente}.
  \item $\vc{R}=\vc{r}-\rp$ es el vector que va de la fuente al punto campo, $R=\norm{\vc{R}}$ y $\uR=\vc{R}/R$.
  \item $\theta'$ es el ángulo que forman $\vc{r}$ y $\rp$, de modo que $\vc{r}\cdot\rp=rr'\cosp$.
  \item $\cd{x},\cd{y},\cd{z}$ son los \emph{cosenos directores} del punto campo, es decir, las componentes cartesianas de $\ur$:
        \[
          \ur=\cd{x}\,\ux+\cd{y}\,\uy+\cd{z}\,\uz,\qquad
          \cd{x}=\frac{x}{r}=\sen\theta\cos\varphi,\quad \cd{y}=\frac{y}{r}=\sen\theta\sen\varphi,\quad \cd{z}=\frac{z}{r}=\cos\theta,
        \]
        que verifican $\cd{x}^{2}+\cd{y}^{2}+\cd{z}^{2}=1$.
  \item Los índices latinos $j,k\in\{x,y,z\}$ recorren las componentes cartesianas; $x'_j$ denota $x',y',z'$, y $\delt{j}{k}$ es la delta de
        Kronecker. Las sumas se escriben siempre de forma explícita, $\sum_{j,k}$.
\end{itemize}
\end{notacion}

\begin{observacion}[Tres usos de la letra $\ell$]\label{obs:ell}
Por fidelidad al formulario del curso~\cite{FORMULARIO}, la letra $\ell$ aparece en tres papeles que el contexto distingue sin ambigüedad:
como \emph{índice} del desarrollo ($\ell=0,1,2,\dots$, orden del multipolo), como \emph{vector} $\vl$ de separación de las cargas de un
dipolo (siempre con flecha) y como \emph{coseno director} $\cd{j}$ (siempre con subíndice cartesiano). En los seminarios, además, $\ell$ es
la longitud de una barra.
\end{observacion}

Como en el Tema~1, la constante de Coulomb es $1/(4\pi\eo)\simeq8{,}988\cdot10^{9}\,\mathrm{N\,m^{2}\,C^{-2}}$~\cite{CODATA} y todas las
distribuciones se suponen \emph{localizadas} (contenidas en una esfera de radio finito), de modo que el potencial se anula en el infinito.

% =========================================================================
\section{El dipolo eléctrico}\label{sec:dipolo}
% =========================================================================
Empezamos por la distribución más sencilla cuyo potencial a grandes distancias \emph{no} es el de una carga puntual: dos cargas
iguales y de signo contrario. Su carga neta es nula, de modo que el término $1/r$ desaparece, y lo que queda es el primer término
que codifica cómo está \emph{repartida} la carga.

% -------------------------------------------------------------------------
\subsection{Definición: dipolo físico y dipolo puntual}\label{sec:dipolo-def}
% -------------------------------------------------------------------------
\begin{definicion}[Dipolo eléctrico y momento dipolar]\label{def:dipolo}
Un \emph{dipolo eléctrico} (o \emph{dipolo físico}) es un sistema formado por dos cargas puntuales de igual módulo y signo contrario,
$+q$ y $-q$ ($q>0$), separadas una distancia $\ell$. Si $\vl$ es el vector que va de la carga negativa a la positiva, se llama
\emph{momento dipolar} del sistema al vector
\begin{equation}\label{eq:momento-dipolar}
  \boxed{\;\vp=q\,\vl\;}
\end{equation}
cuya unidad en el Sistema Internacional es el $\mathrm{C\,m}$.
\end{definicion}

\begin{definicion}[Dipolo puntual o elemental]\label{def:dipolo-puntual}
Se llama \emph{dipolo puntual} (o \emph{elemental}) al límite de un dipolo físico en el que la separación se anula y la carga crece
sin límite, de forma que su producto permanece constante:
\begin{equation}\label{eq:dipolo-puntual}
  \vp=\lim_{\substack{q\to\infty\\ \ell\to0}}q\,\vl .
\end{equation}
\end{definicion}

\begin{observacion}[Órdenes de magnitud]\label{obs:debye}
Los dipolos aparecen de forma natural en la materia: en muchas moléculas el centro de las cargas negativas (los electrones) no coincide con
el de las positivas (los núcleos). Las distancias involucradas son atómicas, del orden de $10^{-10}\,\mathrm{m}$, y por eso los momentos
moleculares se miden en \emph{debyes}, $1\,\mathrm{D}\simeq3{,}336\cdot10^{-30}\,\mathrm{C\,m}$, que es el momento de dos cargas elementales
separadas $0{,}21\,\text{\AA}$. La molécula de agua, por ejemplo, tiene $p\simeq1{,}85\,\mathrm{D}\simeq6{,}2\cdot10^{-30}\,\mathrm{C\,m}$. Vista
desde distancias macroscópicas, $r\gg\ell$, una molécula polar es, a todos los efectos, un dipolo puntual.
\end{observacion}

% -------------------------------------------------------------------------
\subsection{Potencial del dipolo}\label{sec:dipolo-potencial}
% -------------------------------------------------------------------------
Situamos la carga $-q$ en el origen y la carga $+q$ en el extremo de $\vl$ (\figpanel{fig:dipolo-geom}{a}), y llamamos $\alpha$ al ángulo
que forman $\vl$ y $\vc{r}$, de modo que $\vl\cdot\vc{r}=\ell r\cos\alpha$. Por el principio de superposición,
\begin{equation}\label{eq:pot-dipolo-exacto}
  \phi_{d}(\vc{r})=\phi_{+q}(\vc{r})+\phi_{-q}(\vc{r})=\frac{q}{4\pi\eo}\Bigl(\frac{1}{\norm{\vc{r}-\vl}}-\frac{1}{r}\Bigr).
\end{equation}
Esta expresión es exacta, pero nos interesa su comportamiento a distancias grandes comparadas con la separación, $\ell/r\ll1$.

% =========================================================================
% Figura: geometría del dipolo (apuntes manuscritos, págs. 2-3).
% (a) -q en el origen, +q en el extremo de l; punto campo P, ángulo alpha.
% (b) dipolo p = p u_z y base esférica en P: u_z = cos(t) u_r - sen(t) u_t.
% =========================================================================
\begin{figure}[ht]
\centering
\begin{minipage}[b]{0.48\textwidth}\centering
\begin{tikzpicture}[x=1cm, y=1cm, >=Stealth]
  % ejes (perspectiva caballera)
  \draw[eje] (0,0) -- (0,3.3) node[above, etiq] {$z$};
  \draw[eje] (0,0) -- (3.6,0) node[right, etiq] {$y$};
  \draw[eje] (0,0) -- (-1.5,-1.25) node[below left=-2pt, etiq] {$x$};
  % vectores l, r y R = r - l
  \coordinate (O) at (0,0);
  \coordinate (L) at (0.9,1.75);
  \coordinate (P) at (3.7,2.55);
  \draw[vect, rojoteo] (O) -- (L) node[midway, left=1pt, etiq, text=rojoteo] {$\vl$};
  \draw[vect] (O) -- (P) node[midway, below right=-1pt, etiq] {$\vc{r}$};
  \draw[->, line width=0.6pt, grispizarra, densely dashed] (L) -- (P) node[midway, above=1pt, etiq] {$\vc{r}-\vl$};
  % ángulo alpha entre l y r
  \pic[draw=magacero, line width=0.5pt, angle radius=0.75cm, "$\alpha$", angle eccentricity=1.35, font=\small]
      {angle=P--O--L};
  \node[carganeg] at (O) {$-$};
  \node[cargapos] at (L) {$+$};
  \node[left=4pt, etiq] at (O) {$-q$};
  \node[above left=1pt, etiq] at (L) {$+q$};
  \fill[magprof] (P) circle (1.5pt) node[right=2pt, etiq] {$P$};
\end{tikzpicture}\\[2pt]{\small\scshape (a)}
\end{minipage}\hfill
\begin{minipage}[b]{0.48\textwidth}\centering
\begin{tikzpicture}[x=1cm, y=1cm, >=Stealth]
  \def\t{38}\def\rr{3.0}
  \draw[eje] (0,-0.6) -- (0,3.6) node[above, etiq] {$z$};
  \draw[eje] (-0.6,0) -- (3.4,0) node[right, etiq] {$x$};
  \coordinate (O) at (0,0);
  \coordinate (Zp) at (0,1);
  \coordinate (P) at ({\rr*sin(\t)},{\rr*cos(\t)});
  \draw[vect] (O) -- (P) node[midway, below right=-1pt, etiq] {$\vc{r}$};
  \draw[->, line width=1.4pt, rojoteo] (0,-0.3) -- (0,0.75) node[left=2pt, etiq, text=rojoteo] {$\vp$};
  \pic[draw=magacero, line width=0.5pt, angle radius=0.95cm, "$\theta$", angle eccentricity=1.3, font=\small]
      {angle=P--O--Zp};
  % base esférica en P
  \draw[vectverde] (P) -- ++({sin(\t)},{cos(\t)}) node[above right=-2pt, etiq, text=verdeteo] {$\ur$};
  \draw[vectverde] (P) -- ++({cos(\t)},{-sin(\t)}) node[right=1pt, etiq, text=verdeteo] {$\uth$};
  % u_z en P y su descomposición
  \draw[->, line width=0.9pt, magacero] (P) -- ++(0,1) node[above=1pt, etiq, text=magacero] {$\uz$};
  % componentes de u_z: cos(t) u_r (a lo largo de u_r) y -sen(t) u_t
  \draw[punteado] ($(P)+({sin(\t)*cos(\t)},{cos(\t)*cos(\t)})$) -- ($(P)+(0,1)$);
  \fill[magprof] (P) circle (1.5pt) node[left=3pt, etiq] {$P$};
  \node[etiq, align=left, anchor=west, font=\footnotesize] at (1.5,0.55)
     {$\uz=\cos\theta\,\ur-\sen\theta\,\uth$};
\end{tikzpicture}\\[2pt]{\small\scshape (b)}
\end{minipage}
\caption{(a)~Dipolo físico con la carga $-q$ en el origen y $+q$ en el extremo de $\vl$; $\alpha$ es el ángulo entre $\vl$ y el vector de
posición $\vc{r}$ del punto campo $P$, cuya distancia a $+q$ es $\norm{\vc{r}-\vl}$. (b)~Con $\vp=p\,\uz$, el ángulo $\alpha$ es la colatitud
$\theta$ y el versor $\uz$ se descompone en la base esférica local de $P$ (la línea de puntos marca la proyección $\cos\theta$ de $\uz$ sobre $\ur$).}
\label{fig:dipolo-geom}
\end{figure}


\begin{lema}[Aproximación de primer orden]\label{lem:taylor}
Sea $a\in[-1,1]$ y $f(x)=(1+x^{2}-2ax)^{-1/2}$. Entonces, cuando $x\to0$,
\begin{equation}\label{eq:lema-taylor}
  f(x)=1+a\,x+\ord(x^{2}).
\end{equation}
\end{lema}
\begin{proof}
La función es de clase $C^{\infty}$ en un entorno de $x=0$, porque allí $1+x^{2}-2ax>0$. Su derivada es
\[
  f'(x)=-\frac{1}{2}\,(1+x^{2}-2ax)^{-3/2}\,(2x-2a),
\]
de modo que $f(0)=1$ y $f'(0)=-\frac12\,(1)^{-3/2}(-2a)=a$. El polinomio de Taylor de primer orden alrededor de $x=0$ (de Maclaurin) es
$f(0)+f'(0)\,x=1+ax$, y el resto es $\ord(x^{2})$ por el teorema de Taylor.
\end{proof}

\begin{proposicion}[Potencial de un dipolo]\label{prop:pot-dipolo}
A distancias $r\gg\ell$, el potencial del dipolo es
\begin{equation}\label{eq:pot-dipolo}
  \boxed{\;\phi_{d}(\vc{r})=\frac{\vp\cdot\ur}{4\pi\eo\,r^{2}}=\frac{\vp\cdot\vc{r}}{4\pi\eo\,r^{3}}\;}
\end{equation}
salvo términos de orden $q\ell^{2}/r^{3}=p\,\ell/r^{3}$. En particular, \cref{eq:pot-dipolo} es \emph{exacta} para un dipolo puntual.
\end{proposicion}
\begin{proof}
Desarrollamos el cuadrado del módulo y sacamos factor común $r^{2}$:
\[
  \frac{1}{\norm{\vc{r}-\vl}}=\bigl[(\vc{r}-\vl)^{2}\bigr]^{-1/2}=\bigl(r^{2}+\ell^{2}-2\,\vl\cdot\vc{r}\bigr)^{-1/2}
  =\frac{1}{r}\Bigl[1+\Bigl(\frac{\ell}{r}\Bigr)^{2}-2\cos\alpha\,\frac{\ell}{r}\Bigr]^{-1/2}.
\]
Es la función del \cref{lem:taylor} con $x=\ell/r$ y $a=\cos\alpha$, luego
\[
  \frac{1}{\norm{\vc{r}-\vl}}=\frac{1}{r}\Bigl(1+\frac{\ell}{r}\cos\alpha\Bigr)+\ord\Bigl(\frac{\ell^{2}}{r^{3}}\Bigr)
  =\frac{1}{r}\Bigl(1+\frac{\vl\cdot\ur}{r}\Bigr)+\ord\Bigl(\frac{\ell^{2}}{r^{3}}\Bigr).
\]
Sustituyendo en \cref{eq:pot-dipolo-exacto}, el término $1/r$ se cancela con el de la carga negativa:
\[
  \phi_{d}(\vc{r})=\frac{q}{4\pi\eo}\Bigl[\frac{1}{r}\Bigl(1+\frac{\vl\cdot\ur}{r}\Bigr)-\frac{1}{r}\Bigr]+\ord\Bigl(\frac{q\ell^{2}}{r^{3}}\Bigr)
  =\frac{q\,\vl\cdot\ur}{4\pi\eo\,r^{2}}+\ord\Bigl(\frac{p\,\ell}{r^{3}}\Bigr)=\frac{\vp\cdot\ur}{4\pi\eo\,r^{2}}+\ord\Bigl(\frac{p\,\ell}{r^{3}}\Bigr).
\]
En el límite del dipolo puntual, $p$ permanece fijo y $\ell\to0$, de modo que el resto se anula. La segunda forma de \cref{eq:pot-dipolo} se
sigue de $\ur=\vc{r}/r$.
\end{proof}

\begin{corolario}[Simetría de revolución]\label{cor:pot-dipolo-z}
Si el momento dipolar apunta según el eje $z$, $\vp=p\,\uz$, entonces $\vp\cdot\ur=p\cos\theta$ y
\begin{equation}\label{eq:pot-dipolo-z}
  \phi_{d}(r,\theta)=\frac{p\cos\theta}{4\pi\eo\,r^{2}} .
\end{equation}
El potencial no depende del ángulo azimutal $\varphi$: tiene simetría de revolución alrededor del eje del dipolo.
\end{corolario}
\begin{proof}
En la base esférica (\figpanel{fig:dipolo-geom}{b}), $\uz=\cos\theta\,\ur-\sen\theta\,\uth$, luego
$\uz\cdot\ur=(\cos\theta\,\ur-\sen\theta\,\uth)\cdot\ur=\cos\theta$, pues $\{\ur,\uth,\uph\}$ es ortonormal.
\end{proof}

\begin{observacion}[Decae más deprisa que el de una carga]\label{obs:decaimiento}
El potencial dipolar decae como $1/r^{2}$, más deprisa que el de una carga puntual, $q/(4\pi\eo r)$, que tiene simetría esférica. Es la
huella de la cancelación entre las dos cargas: vistas desde muy lejos casi se superponen, y solo sobrevive el efecto de su pequeña
separación. Además, $\phi_{d}$ ya no es isótropo: es positivo en el hemisferio hacia el que apunta $\vp$ ($\cos\theta>0$), negativo en el
opuesto y nulo en todo el plano ecuatorial $\theta=\pi/2$.
\end{observacion}

% -------------------------------------------------------------------------
\subsection{Campo eléctrico del dipolo}\label{sec:dipolo-campo}
% -------------------------------------------------------------------------
\begin{teorema}[Campo de un dipolo]\label{teo:campo-dipolo}
El campo eléctrico creado por un dipolo puntual $\vp$ situado en el origen es
\begin{equation}\label{eq:campo-dipolo}
  \boxed{\;\vE_{d}(\vc{r})=\frac{1}{4\pi\eo\,r^{3}}\bigl[3(\vp\cdot\ur)\,\ur-\vp\bigr]\;}
\end{equation}
y, si $\vp=p\,\uz$, en coordenadas esféricas,
\begin{equation}\label{eq:campo-dipolo-esf}
  \vE_{d}(\vc{r})=\frac{p}{4\pi\eo\,r^{3}}\bigl(2\cos\theta\,\ur+\sen\theta\,\uth\bigr).
\end{equation}
La forma \cref{eq:campo-dipolo} es intrínseca: no depende de la orientación del dipolo ni del sistema de coordenadas.
\end{teorema}
\begin{proof}[Demostración (en esféricas)]
Partimos de \cref{eq:pot-dipolo-z} y del gradiente en esféricas~\cite{TEMA0},
$\nab f=\pd{f}{r}\,\ur+\frac{1}{r}\pd{f}{\theta}\,\uth+\frac{1}{r\sen\theta}\pd{f}{\varphi}\,\uph$. Como $\phi_{d}$ no depende de $\varphi$,
\begin{align*}
  \vE_{d}&=-\nab\phi_{d}=-\pd{\phi_{d}}{r}\,\ur-\frac{1}{r}\pd{\phi_{d}}{\theta}\,\uth
  =-\frac{p\cos\theta}{4\pi\eo}\Bigl(-\frac{2}{r^{3}}\Bigr)\ur-\frac{1}{r}\,\frac{p}{4\pi\eo r^{2}}\,(-\sen\theta)\,\uth\\
  &=\frac{p}{4\pi\eo r^{3}}\bigl(2\cos\theta\,\ur+\sen\theta\,\uth\bigr).
\end{align*}
Para llegar a la forma intrínseca escribimos $2\cos\theta\,\ur=3\cos\theta\,\ur-\cos\theta\,\ur$ y usamos $\uz=\cos\theta\,\ur-\sen\theta\,\uth$:
\[
  2\cos\theta\,\ur+\sen\theta\,\uth=3\cos\theta\,\ur-\underbrace{\bigl(\cos\theta\,\ur-\sen\theta\,\uth\bigr)}_{\uz}=3\cos\theta\,\ur-\uz .
\]
Multiplicando por $p$ y recordando que $p\cos\theta=\vp\cdot\ur$ y $p\,\uz=\vp$ se obtiene \cref{eq:campo-dipolo}. Esta última ya no hace referencia
al eje $z$; como ambos miembros son vectores, vale para cualquier orientación de $\vp$.
\end{proof}

La forma intrínseca puede obtenerse también directamente, sin elegir ejes, derivando $\phi_{d}=\vp\cdot\vc{r}/(4\pi\eo r^{3})$. Es un buen
ejercicio de cálculo vectorial~\cite{TEMA0}.

\begin{proof}[Demostración (vectorial)]
Por la regla del producto $\nab(fg)=f\nab g+g\nab f$,
\[
  \vE_{d}=-\frac{1}{4\pi\eo}\,\nab\Bigl[(\vp\cdot\vc{r})\,\frac{1}{r^{3}}\Bigr]
  =-\frac{1}{4\pi\eo}\Bigl[\frac{1}{r^{3}}\,\nab(\vp\cdot\vc{r})+(\vp\cdot\vc{r})\,\nab\Bigl(\frac{1}{r^{3}}\Bigr)\Bigr].
\]
\emph{Primer gradiente.} Como $\vp$ es constante, en componentes, $\bigl[\nab(\vp\cdot\vc{r})\bigr]_{j}=\pd{}{x_j}\sum_{i}p_i x_i=\sum_i p_i\,\pd{x_i}{x_j}
=\sum_i p_i\,\delt{i}{j}=p_j$, es decir, $\nab(\vp\cdot\vc{r})=\vp$. (Equivalentemente: $\nab(\vp\cdot\vc{r})=(\vp\cdot\nab)\vc{r}=\vp$, porque
$\nab\times\vc{r}=\vc{0}$ y las derivadas de $\vp$ son nulas.)

\emph{Segundo gradiente.} Para una función radial, $\nab f(r)=f'(r)\,\ur$, luego $\nab\bigl(r^{-3}\bigr)=-3r^{-4}\,\ur$.

Sustituyendo, y escribiendo $\vp\cdot\vc{r}=r\,(\vp\cdot\ur)$,
\[
  \vE_{d}=-\frac{1}{4\pi\eo}\Bigl[\frac{\vp}{r^{3}}-\frac{3\,r\,(\vp\cdot\ur)}{r^{4}}\,\ur\Bigr]=\frac{1}{4\pi\eo r^{3}}\bigl[3(\vp\cdot\ur)\,\ur-\vp\bigr].\qedhere
\]
\end{proof}

\begin{observacion}[Estructura del campo dipolar]\label{obs:estructura}
\begin{enumerate}[label=(\alph*)]
  \item \emph{Dos contribuciones.} El campo es la suma de una componente paralela a $\ur$ (la dirección del punto campo) y otra
        \emph{antiparalela} a $\vp$.
  \item \emph{Ejes y ecuador.} $E_{r}=0$ en el plano ecuatorial, $\theta=\pi/2$, donde $\vE_{d}=-\vp/(4\pi\eo r^{3})$ es antiparalelo al momento;
        $E_{\theta}=0$ sobre el eje, $\theta=0,\pi$, donde $\vE_{d}=2\vp/(4\pi\eo r^{3})$ es paralelo al momento y \emph{el doble} de intenso que en el
        ecuador a la misma distancia.
  \item \emph{Decaimiento.} El campo dipolar cae como $1/r^{3}$, mucho más deprisa que el de una carga puntual, que lo hace como $1/r^{2}$
        (\cref{fig:potencias}). Si tuviésemos una carga puntual y un dipolo elemental en el mismo punto, a distancias inferiores a una dada
        prevalecería la contribución dipolar, mientras que a grandes distancias dominaría la de la carga.
\end{enumerate}
\end{observacion}

% =========================================================================
% Figura (apuntes de teoría, pág. II-3): decaimiento 1/r^2 frente a 1/r^3.
% Se cortan en r = 1: más cerca domina el término más rápido; más lejos,
% el más lento.
% =========================================================================
\begin{figure}[ht]
\centering
\begin{tikzpicture}
\begin{axis}[tema, width=8.6cm, height=5.2cm, xmin=0.4, xmax=4, ymin=0, ymax=4.2,
             xlabel={$r$ (u.\,arb.)}, ylabel={campo (u.\,arb.)}, grid=major,
             xtick={0.5,1,...,4}, ytick={0,1,...,4}, clip=true,
             legend pos=north east, legend cell align=left]
  \addplot[mmaazul, line width=1.1pt, domain=0.45:4, samples=150] {1/x^2};
  \addlegendentry{$1/r^{2}$ (carga)}
  \addplot[mmanaranja!90!black, line width=1.1pt, domain=0.6:4, samples=150] {1/x^3};
  \addlegendentry{$1/r^{3}$ (dipolo)}
  \draw[punteado] (axis cs:1,0) -- (axis cs:1,1);
  \node[font=\scriptsize, text=grispizarra, anchor=south west] at (axis cs:1.05,1.05) {$r=1$};
\end{axis}
\end{tikzpicture}
\caption{Comparación del decaimiento de una contribución monopolar ($1/r^{2}$, campo de una carga) y otra dipolar ($1/r^{3}$, campo de
un dipolo) con la misma intensidad en $r=1$. Para $r<1$ domina el término dipolar; para $r>1$, el de la carga, y la diferencia se agranda
rápidamente con la distancia.}
\label{fig:potencias}
\end{figure}


Las expresiones anteriores suponen el dipolo en el origen. Si está en un punto arbitrario $\rp$, basta sustituir $\vc{r}$ por la posición
relativa $\vc{R}=\vc{r}-\rp$; también puede comprobarse directamente sumando los campos de las dos cargas.

\begin{proposicion}[Dipolo en una posición arbitraria]\label{prop:dipolo-Rp}
El potencial y el campo de un dipolo puntual $\vp$ situado en $\rp$ son
\begin{equation}\label{eq:dipolo-R}
  \phi(\vc{r})=\frac{1}{4\pi\eo}\,\frac{\vp\cdot\vc{R}}{R^{3}},\qquad
  \vE(\vc{r})=\frac{1}{4\pi\eo}\Bigl[\frac{3(\vp\cdot\vc{R})\,\vc{R}}{R^{5}}-\frac{\vp}{R^{3}}\Bigr],\qquad \vc{R}=\vc{r}-\rp .
\end{equation}
\end{proposicion}
\begin{proof}
Sea $-q$ en $\rp$ y $+q$ en $\rp+\vc{d}$, con $\vp=q\vc{d}$. Por superposición de dos campos coulombianos,
\[
  \vE=\frac{q}{4\pi\eo}\Biggl[\frac{\vc{r}-(\rp+\vc{d})}{\norm{\vc{r}-(\rp+\vc{d})}^{3}}-\frac{\vc{r}-\rp}{\norm{\vc{r}-\rp}^{3}}\Biggr]
  =\frac{q}{4\pi\eo}\Biggl[\frac{\vc{R}-\vc{d}}{\norm{\vc{R}-\vc{d}}^{3}}-\frac{\vc{R}}{R^{3}}\Biggr].
\]
Si $d\ll R$, el término cuadrático $d^{2}/R^{2}$ es despreciable y, por el desarrollo binomial $(1+x)^{n}=1+nx+\ord(x^{2})$ con $n=-\frac32$,
\[
  \norm{\vc{R}-\vc{d}}^{-3}=\bigl[R^{2}+d^{2}-2\vc{R}\cdot\vc{d}\bigr]^{-3/2}=R^{-3}\Bigl[1-2\,\frac{\vc{R}\cdot\vc{d}}{R^{2}}+\frac{d^{2}}{R^{2}}\Bigr]^{-3/2}
  \simeq R^{-3}\Bigl[1+3\,\frac{\vc{R}\cdot\vc{d}}{R^{2}}\Bigr].
\]
Sustituyendo y despreciando de nuevo los términos de orden $d^{2}$,
\[
  \vE\simeq\frac{q}{4\pi\eo}\Bigl[(\vc{R}-\vc{d})\Bigl(\frac{1}{R^{3}}+\frac{3\,\vc{R}\cdot\vc{d}}{R^{5}}\Bigr)-\frac{\vc{R}}{R^{3}}\Bigr]
  \simeq\frac{q}{4\pi\eo}\Bigl[\frac{3(\vc{R}\cdot\vc{d})\,\vc{R}}{R^{5}}-\frac{\vc{d}}{R^{3}}\Bigr],
\]
y el límite $q\to\infty$, $d\to0$ con $q\vc{d}=\vp$ da el campo de \cref{eq:dipolo-R}. Para el potencial se procede igual con
$\norm{\vc{R}-\vc{d}}^{-1}\simeq R^{-1}\bigl[1+\vc{R}\cdot\vc{d}/R^{2}\bigr]$, lo que da
$\phi\simeq\frac{q}{4\pi\eo}\,\frac{\vc{R}\cdot\vc{d}}{R^{3}}\to\frac{1}{4\pi\eo}\frac{\vp\cdot\vc{R}}{R^{3}}$. Puede comprobarse que $\vE=-\nab\phi$.
\end{proof}

% -------------------------------------------------------------------------
\subsection{Líneas de campo y superficies equipotenciales}\label{sec:dipolo-lineas}
% -------------------------------------------------------------------------
\begin{proposicion}[Geometría del campo dipolar]\label{prop:lineas-dipolo}
Para un dipolo puntual $\vp=p\,\uz$, en cualquier plano meridiano:
\begin{enumerate}[label=(\roman*)]
  \item las equipotenciales $\phi_{d}=\phi_0\neq0$ son las curvas $r^{2}=K\cos\theta$, con $K=p/(4\pi\eo\phi_0)$, y la equipotencial $\phi_{d}=0$ es el plano
        ecuatorial;
  \item las líneas de campo son las curvas $r=C\sen^{2}\theta$, con $C>0$, más el propio eje $z$.
\end{enumerate}
Ambas familias son ortogonales en todo punto.
\end{proposicion}
\begin{proof}
(i) Es \cref{eq:pot-dipolo-z} igualada a $\phi_0$. (ii) Una línea de campo es tangente a $\vE$ en cada punto; en el plano meridiano, el desplazamiento
$\dif\vc{r}=\dif r\,\ur+r\,\dif\theta\,\uth$ ha de ser paralelo a $\vE$, es decir, $\dif r/E_r=r\,\dif\theta/E_\theta$. Con \cref{eq:campo-dipolo-esf},
\[
  \frac{\dif r}{r}=\frac{E_r}{E_\theta}\,\dif\theta=\frac{2\cos\theta}{\sen\theta}\,\dif\theta
  \;\Longrightarrow\;\ln r=2\ln(\sen\theta)+\text{cte}\;\Longrightarrow\; r=C\sen^{2}\theta .
\]
La ortogonalidad es general: $\vE=-\nab\phi$ es perpendicular a las superficies de nivel de $\phi$~\cite{TEMA1}.
\end{proof}

La \cref{fig:dipolo-mma} reproduce las representaciones del dipolo de las figuras del tema~\cite{FIGURAS-T2}: las equipotenciales,
las líneas de campo superpuestas a ellas y el campo vectorial, cuya intensidad decrece rápidamente con la distancia.

% =========================================================================
% Figura (clon de «Tema 2. Figuras», diapositiva «Dipolo eléctrico»):
% (a) equipotenciales (ContourPlot), (b) equipotenciales y líneas de campo
% (StreamPlot) y (c) campo eléctrico (VectorPlot) de un dipolo puntual
% p = q l u_z. Datos: figuras/datos/dip_puntual_*.tex. Se conservan los
% colores originales de Mathematica.
% =========================================================================
\begin{figure}[t]
\centering
\begin{minipage}[t]{0.325\textwidth}\centering
{\footnotesize Equipotenciales}\\[3pt]
\begin{tikzpicture}[x=0.198cm, y=0.198cm]
  \renewcommand{\EQ}[1]{\draw[mmaEQsolo] plot coordinates {#1};}%
  \panelMMA{10}{white}{0}{%
    % Dipolo puntual - equipotenciales
% Generado por scripts/generar_datos.py (no editar a mano).
\EQ{(1.224,-0.476) (1.250,-0.522) (1.275,-0.574) (1.296,-0.625) (1.313,-0.675) (1.326,-0.725) (1.337,-0.775) (1.345,-0.825) (1.350,-0.875) (1.353,-0.925) (1.354,-0.975) (1.352,-1.025) (1.347,-1.075) (1.340,-1.125) (1.331,-1.175) (1.319,-1.225) (1.305,-1.275) (1.288,-1.325) (1.269,-1.375) (1.246,-1.425) (1.221,-1.475) (1.194,-1.525) (1.162,-1.575) (1.128,-1.625) (1.089,-1.675) (1.050,-1.721) (1.000,-1.774) (0.950,-1.822) (0.900,-1.865) (0.856,-1.900) (0.800,-1.940) (0.750,-1.972) (0.702,-2.000) (0.656,-2.025) (0.604,-2.050) (0.550,-2.074) (0.500,-2.093) (0.450,-2.111) (0.400,-2.126) (0.350,-2.139) (0.300,-2.151) (0.250,-2.161) (0.200,-2.168) (0.150,-2.174) (0.100,-2.179) (0.050,-2.181) (0.000,-2.182) (-0.050,-2.181) (-0.100,-2.179) (-0.150,-2.174) (-0.200,-2.168) (-0.250,-2.161) (-0.300,-2.151) (-0.350,-2.139) (-0.400,-2.126) (-0.450,-2.111) (-0.500,-2.093) (-0.547,-2.075) (-0.600,-2.052) (-0.650,-2.028) (-0.700,-2.001) (-0.745,-1.975) (-0.800,-1.940) (-0.850,-1.904) (-0.900,-1.865) (-0.947,-1.825) (-0.999,-1.775) (-1.047,-1.725) (-1.089,-1.675) (-1.125,-1.629) (-1.162,-1.575) (-1.194,-1.525) (-1.221,-1.475) (-1.246,-1.425) (-1.269,-1.375) (-1.288,-1.325) (-1.305,-1.275) (-1.319,-1.225) (-1.331,-1.175) (-1.340,-1.125) (-1.347,-1.075) (-1.352,-1.025) (-1.354,-0.975) (-1.353,-0.925) (-1.350,-0.875) (-1.345,-0.825) (-1.337,-0.775) (-1.326,-0.725) (-1.313,-0.675) (-1.296,-0.625) (-1.275,-0.575) (-1.252,-0.525) (-1.225,-0.477)}
\EQ{(1.229,-0.446) (1.263,-0.500) (1.290,-0.550) (1.313,-0.600) (1.333,-0.650) (1.349,-0.700) (1.362,-0.750) (1.373,-0.800) (1.381,-0.850) (1.386,-0.900) (1.389,-0.950) (1.389,-1.000) (1.387,-1.050) (1.383,-1.100) (1.376,-1.150) (1.367,-1.200) (1.355,-1.250) (1.341,-1.300) (1.325,-1.349) (1.306,-1.400) (1.284,-1.450) (1.260,-1.500) (1.233,-1.550) (1.203,-1.600) (1.175,-1.642) (1.132,-1.700) (1.100,-1.740) (1.050,-1.796) (1.000,-1.846) (0.950,-1.891) (0.909,-1.925) (0.850,-1.970) (0.806,-2.000) (0.750,-2.035) (0.700,-2.063) (0.650,-2.089) (0.600,-2.112) (0.550,-2.133) (0.500,-2.152) (0.450,-2.169) (0.400,-2.184) (0.350,-2.197) (0.300,-2.208) (0.250,-2.218) (0.200,-2.225) (0.150,-2.231) (0.100,-2.236) (0.050,-2.238) (0.000,-2.239) (-0.050,-2.238) (-0.100,-2.236) (-0.150,-2.231) (-0.200,-2.225) (-0.250,-2.218) (-0.300,-2.208) (-0.350,-2.197) (-0.400,-2.184) (-0.450,-2.169) (-0.500,-2.152) (-0.550,-2.133) (-0.600,-2.112) (-0.650,-2.089) (-0.700,-2.063) (-0.750,-2.035) (-0.800,-2.004) (-0.843,-1.975) (-0.900,-1.932) (-0.940,-1.900) (-0.996,-1.850) (-1.046,-1.800) (-1.091,-1.750) (-1.125,-1.709) (-1.169,-1.650) (-1.200,-1.604) (-1.233,-1.550) (-1.260,-1.500) (-1.284,-1.450) (-1.306,-1.400) (-1.325,-1.350) (-1.341,-1.300) (-1.355,-1.250) (-1.367,-1.200) (-1.376,-1.150) (-1.383,-1.100) (-1.387,-1.050) (-1.389,-1.000) (-1.389,-0.950) (-1.386,-0.900) (-1.381,-0.850) (-1.373,-0.800) (-1.362,-0.750) (-1.349,-0.700) (-1.333,-0.650) (-1.313,-0.600) (-1.290,-0.550) (-1.263,-0.500) (-1.232,-0.450) (-1.229,-0.446)}
\EQ{(1.244,-0.431) (1.274,-0.475) (1.300,-0.518) (1.325,-0.565) (1.350,-0.619) (1.372,-0.675) (1.388,-0.725) (1.400,-0.775) (1.411,-0.825) (1.419,-0.875) (1.424,-0.925) (1.427,-0.975) (1.427,-1.025) (1.425,-1.075) (1.421,-1.125) (1.414,-1.175) (1.406,-1.225) (1.395,-1.275) (1.381,-1.325) (1.365,-1.375) (1.347,-1.425) (1.326,-1.475) (1.303,-1.525) (1.277,-1.575) (1.250,-1.621) (1.215,-1.675) (1.180,-1.725) (1.141,-1.775) (1.100,-1.822) (1.050,-1.874) (1.000,-1.922) (0.950,-1.965) (0.906,-2.000) (0.850,-2.040) (0.800,-2.073) (0.755,-2.100) (0.709,-2.125) (0.660,-2.150) (0.605,-2.175) (0.550,-2.198) (0.500,-2.216) (0.450,-2.233) (0.400,-2.247) (0.350,-2.260) (0.300,-2.271) (0.250,-2.280) (0.200,-2.287) (0.150,-2.293) (0.100,-2.297) (0.050,-2.299) (0.000,-2.300) (-0.050,-2.299) (-0.100,-2.297) (-0.150,-2.293) (-0.200,-2.287) (-0.250,-2.280) (-0.300,-2.271) (-0.350,-2.260) (-0.400,-2.247) (-0.450,-2.233) (-0.500,-2.216) (-0.550,-2.198) (-0.600,-2.177) (-0.650,-2.155) (-0.700,-2.130) (-0.750,-2.103) (-0.796,-2.075) (-0.850,-2.040) (-0.900,-2.004) (-0.950,-1.965) (-0.997,-1.925) (-1.049,-1.875) (-1.097,-1.825) (-1.141,-1.775) (-1.175,-1.731) (-1.215,-1.675) (-1.248,-1.625) (-1.275,-1.578) (-1.300,-1.531) (-1.325,-1.478) (-1.347,-1.425) (-1.365,-1.375) (-1.381,-1.325) (-1.395,-1.275) (-1.406,-1.225) (-1.414,-1.175) (-1.421,-1.125) (-1.425,-1.075) (-1.427,-1.025) (-1.427,-0.975) (-1.424,-0.925) (-1.419,-0.875) (-1.411,-0.825) (-1.400,-0.775) (-1.388,-0.725) (-1.372,-0.675) (-1.352,-0.625) (-1.330,-0.575) (-1.304,-0.525) (-1.275,-0.477) (-1.244,-0.431)}
\EQ{(1.250,-0.405) (1.285,-0.450) (1.318,-0.500) (1.347,-0.550) (1.373,-0.600) (1.395,-0.650) (1.413,-0.700) (1.429,-0.750) (1.442,-0.800) (1.452,-0.850) (1.460,-0.900) (1.465,-0.950) (1.468,-1.000) (1.468,-1.050) (1.467,-1.100) (1.463,-1.150) (1.457,-1.200) (1.448,-1.250) (1.438,-1.300) (1.425,-1.350) (1.410,-1.400) (1.392,-1.450) (1.372,-1.500) (1.350,-1.550) (1.325,-1.600) (1.300,-1.645) (1.266,-1.700) (1.232,-1.750) (1.200,-1.793) (1.154,-1.850) (1.109,-1.900) (1.059,-1.950) (1.004,-2.000) (0.950,-2.044) (0.909,-2.075) (0.850,-2.116) (0.800,-2.147) (0.751,-2.175) (0.704,-2.200) (0.652,-2.225) (0.600,-2.248) (0.550,-2.268) (0.500,-2.285) (0.450,-2.301) (0.400,-2.315) (0.350,-2.328) (0.300,-2.338) (0.250,-2.347) (0.200,-2.354) (0.150,-2.360) (0.100,-2.364) (0.050,-2.366) (0.000,-2.367) (-0.050,-2.366) (-0.100,-2.364) (-0.150,-2.360) (-0.200,-2.354) (-0.250,-2.347) (-0.300,-2.338) (-0.350,-2.328) (-0.400,-2.315) (-0.450,-2.301) (-0.500,-2.285) (-0.550,-2.268) (-0.600,-2.248) (-0.650,-2.226) (-0.700,-2.202) (-0.750,-2.176) (-0.795,-2.150) (-0.850,-2.116) (-0.900,-2.081) (-0.942,-2.050) (-1.000,-2.003) (-1.050,-1.958) (-1.100,-1.909) (-1.150,-1.854) (-1.195,-1.800) (-1.225,-1.760) (-1.266,-1.700) (-1.297,-1.650) (-1.325,-1.600) (-1.350,-1.550) (-1.372,-1.500) (-1.392,-1.450) (-1.410,-1.400) (-1.425,-1.350) (-1.438,-1.300) (-1.448,-1.250) (-1.457,-1.200) (-1.463,-1.150) (-1.467,-1.100) (-1.468,-1.050) (-1.468,-1.000) (-1.465,-0.950) (-1.460,-0.900) (-1.452,-0.850) (-1.442,-0.800) (-1.429,-0.750) (-1.413,-0.700) (-1.395,-0.650) (-1.373,-0.600) (-1.350,-0.555) (-1.325,-0.512) (-1.285,-0.450) (-1.250,-0.405)}
\EQ{(1.251,-0.374) (1.295,-0.425) (1.325,-0.465) (1.365,-0.525) (1.393,-0.575) (1.418,-0.625) (1.440,-0.675) (1.458,-0.725) (1.473,-0.775) (1.486,-0.825) (1.496,-0.875) (1.504,-0.925) (1.510,-0.975) (1.513,-1.025) (1.514,-1.075) (1.512,-1.125) (1.509,-1.175) (1.503,-1.225) (1.495,-1.275) (1.486,-1.325) (1.473,-1.375) (1.459,-1.425) (1.442,-1.475) (1.423,-1.525) (1.402,-1.575) (1.378,-1.625) (1.352,-1.675) (1.325,-1.721) (1.290,-1.775) (1.255,-1.825) (1.216,-1.875) (1.175,-1.923) (1.127,-1.975) (1.075,-2.025) (1.025,-2.070) (0.975,-2.110) (0.925,-2.147) (0.875,-2.182) (0.825,-2.213) (0.775,-2.242) (0.725,-2.268) (0.675,-2.292) (0.625,-2.314) (0.575,-2.334) (0.525,-2.352) (0.475,-2.369) (0.425,-2.383) (0.375,-2.396) (0.325,-2.407) (0.275,-2.416) (0.225,-2.424) (0.175,-2.430) (0.125,-2.435) (0.075,-2.438) (0.025,-2.440) (-0.050,-2.439) (-0.100,-2.437) (-0.150,-2.433) (-0.200,-2.427) (-0.250,-2.420) (-0.300,-2.412) (-0.350,-2.402) (-0.400,-2.390) (-0.450,-2.376) (-0.500,-2.361) (-0.550,-2.344) (-0.599,-2.325) (-0.650,-2.304) (-0.700,-2.280) (-0.750,-2.255) (-0.800,-2.228) (-0.846,-2.200) (-0.900,-2.165) (-0.950,-2.129) (-1.000,-2.090) (-1.048,-2.050) (-1.100,-2.002) (-1.150,-1.951) (-1.195,-1.900) (-1.236,-1.850) (-1.273,-1.800) (-1.307,-1.750) (-1.337,-1.700) (-1.365,-1.650) (-1.390,-1.600) (-1.413,-1.550) (-1.433,-1.500) (-1.450,-1.453) (-1.466,-1.400) (-1.480,-1.350) (-1.491,-1.300) (-1.500,-1.250) (-1.506,-1.200) (-1.511,-1.150) (-1.513,-1.100) (-1.514,-1.050) (-1.512,-1.000) (-1.507,-0.950) (-1.501,-0.900) (-1.492,-0.850) (-1.480,-0.800) (-1.466,-0.750) (-1.449,-0.700) (-1.429,-0.650) (-1.406,-0.600) (-1.379,-0.550) (-1.350,-0.502) (-1.314,-0.450) (-1.275,-0.401) (-1.251,-0.374)}
\EQ{(1.266,-0.359) (1.305,-0.400) (1.346,-0.450) (1.382,-0.500) (1.414,-0.550) (1.442,-0.600) (1.467,-0.650) (1.488,-0.700) (1.506,-0.750) (1.522,-0.800) (1.534,-0.850) (1.545,-0.900) (1.553,-0.950) (1.559,-1.000) (1.562,-1.050) (1.563,-1.100) (1.563,-1.150) (1.560,-1.200) (1.555,-1.250) (1.547,-1.300) (1.538,-1.350) (1.527,-1.400) (1.513,-1.450) (1.498,-1.500) (1.480,-1.550) (1.460,-1.600) (1.437,-1.650) (1.412,-1.700) (1.385,-1.750) (1.354,-1.800) (1.325,-1.844) (1.285,-1.900) (1.250,-1.944) (1.201,-2.000) (1.154,-2.050) (1.101,-2.100) (1.050,-2.145) (1.000,-2.185) (0.950,-2.222) (0.908,-2.250) (0.850,-2.287) (0.800,-2.315) (0.750,-2.342) (0.700,-2.366) (0.650,-2.388) (0.600,-2.409) (0.550,-2.427) (0.500,-2.443) (0.450,-2.458) (0.400,-2.471) (0.350,-2.483) (0.300,-2.493) (0.250,-2.501) (0.200,-2.508) (0.150,-2.513) (0.100,-2.517) (0.050,-2.519) (0.000,-2.520) (-0.050,-2.519) (-0.100,-2.517) (-0.150,-2.513) (-0.200,-2.508) (-0.250,-2.501) (-0.300,-2.493) (-0.350,-2.483) (-0.400,-2.471) (-0.450,-2.458) (-0.500,-2.443) (-0.550,-2.427) (-0.600,-2.409) (-0.650,-2.388) (-0.700,-2.366) (-0.750,-2.342) (-0.800,-2.315) (-0.850,-2.287) (-0.900,-2.256) (-0.945,-2.225) (-1.000,-2.185) (-1.044,-2.150) (-1.100,-2.101) (-1.150,-2.054) (-1.200,-2.001) (-1.245,-1.950) (-1.285,-1.900) (-1.321,-1.850) (-1.350,-1.807) (-1.385,-1.750) (-1.412,-1.700) (-1.437,-1.650) (-1.460,-1.600) (-1.480,-1.550) (-1.498,-1.500) (-1.513,-1.450) (-1.527,-1.400) (-1.538,-1.350) (-1.547,-1.300) (-1.555,-1.250) (-1.560,-1.200) (-1.563,-1.150) (-1.563,-1.100) (-1.562,-1.050) (-1.559,-1.000) (-1.553,-0.950) (-1.545,-0.900) (-1.534,-0.850) (-1.522,-0.800) (-1.506,-0.750) (-1.488,-0.700) (-1.467,-0.650) (-1.442,-0.600) (-1.414,-0.550) (-1.382,-0.500) (-1.350,-0.456) (-1.305,-0.400) (-1.266,-0.359)}
\EQ{(1.275,-0.337) (1.314,-0.375) (1.350,-0.414) (1.400,-0.475) (1.436,-0.525) (1.467,-0.575) (1.495,-0.625) (1.519,-0.675) (1.540,-0.725) (1.559,-0.775) (1.574,-0.825) (1.587,-0.875) (1.598,-0.925) (1.606,-0.975) (1.612,-1.025) (1.616,-1.075) (1.618,-1.125) (1.618,-1.175) (1.616,-1.225) (1.611,-1.275) (1.605,-1.325) (1.597,-1.375) (1.586,-1.425) (1.574,-1.475) (1.559,-1.525) (1.543,-1.575) (1.524,-1.625) (1.503,-1.675) (1.479,-1.725) (1.454,-1.775) (1.425,-1.825) (1.394,-1.875) (1.360,-1.925) (1.325,-1.973) (1.283,-2.025) (1.250,-2.063) (1.216,-2.100) (1.175,-2.141) (1.125,-2.187) (1.081,-2.225) (1.025,-2.269) (0.983,-2.300) (0.925,-2.339) (0.875,-2.370) (0.825,-2.398) (0.775,-2.425) (0.725,-2.449) (0.675,-2.471) (0.625,-2.492) (0.575,-2.510) (0.525,-2.527) (0.475,-2.542) (0.425,-2.555) (0.375,-2.567) (0.325,-2.578) (0.275,-2.586) (0.225,-2.594) (0.175,-2.599) (0.125,-2.604) (0.075,-2.607) (0.025,-2.608) (-0.050,-2.608) (-0.100,-2.605) (-0.150,-2.602) (-0.200,-2.597) (-0.250,-2.590) (-0.300,-2.582) (-0.350,-2.573) (-0.400,-2.562) (-0.450,-2.549) (-0.500,-2.535) (-0.550,-2.519) (-0.600,-2.501) (-0.650,-2.482) (-0.700,-2.460) (-0.750,-2.437) (-0.800,-2.412) (-0.850,-2.384) (-0.900,-2.355) (-0.946,-2.325) (-1.000,-2.288) (-1.050,-2.250) (-1.100,-2.209) (-1.139,-2.175) (-1.191,-2.125) (-1.239,-2.075) (-1.275,-2.034) (-1.323,-1.975) (-1.360,-1.925) (-1.394,-1.875) (-1.425,-1.825) (-1.450,-1.781) (-1.475,-1.734) (-1.500,-1.681) (-1.524,-1.625) (-1.543,-1.575) (-1.559,-1.525) (-1.574,-1.475) (-1.586,-1.425) (-1.597,-1.375) (-1.605,-1.325) (-1.611,-1.275) (-1.616,-1.225) (-1.618,-1.175) (-1.618,-1.125) (-1.616,-1.075) (-1.612,-1.025) (-1.606,-0.975) (-1.598,-0.925) (-1.587,-0.875) (-1.574,-0.825) (-1.559,-0.775) (-1.540,-0.725) (-1.519,-0.675) (-1.495,-0.625) (-1.467,-0.575) (-1.436,-0.525) (-1.400,-0.475) (-1.360,-0.425) (-1.325,-0.387) (-1.289,-0.350) (-1.275,-0.337)}
\EQ{(1.275,-0.308) (1.323,-0.350) (1.374,-0.400) (1.419,-0.450) (1.450,-0.489) (1.493,-0.550) (1.524,-0.600) (1.550,-0.647) (1.575,-0.698) (1.597,-0.750) (1.616,-0.800) (1.632,-0.850) (1.645,-0.900) (1.656,-0.950) (1.665,-1.000) (1.672,-1.050) (1.676,-1.100) (1.679,-1.150) (1.679,-1.200) (1.678,-1.250) (1.674,-1.300) (1.669,-1.350) (1.662,-1.400) (1.652,-1.450) (1.641,-1.500) (1.628,-1.550) (1.613,-1.600) (1.595,-1.650) (1.576,-1.700) (1.554,-1.750) (1.530,-1.800) (1.504,-1.850) (1.475,-1.900) (1.444,-1.950) (1.410,-2.000) (1.375,-2.047) (1.333,-2.100) (1.289,-2.150) (1.250,-2.191) (1.200,-2.240) (1.162,-2.275) (1.102,-2.325) (1.050,-2.365) (1.001,-2.400) (0.950,-2.434) (0.900,-2.464) (0.850,-2.492) (0.800,-2.518) (0.750,-2.543) (0.700,-2.565) (0.650,-2.585) (0.600,-2.604) (0.550,-2.621) (0.500,-2.636) (0.450,-2.650) (0.400,-2.662) (0.350,-2.672) (0.300,-2.682) (0.250,-2.689) (0.200,-2.696) (0.150,-2.700) (0.100,-2.704) (0.050,-2.706) (0.000,-2.707) (-0.050,-2.706) (-0.100,-2.704) (-0.150,-2.700) (-0.200,-2.696) (-0.250,-2.689) (-0.300,-2.682) (-0.350,-2.672) (-0.400,-2.662) (-0.450,-2.650) (-0.500,-2.636) (-0.550,-2.621) (-0.600,-2.604) (-0.650,-2.585) (-0.700,-2.565) (-0.750,-2.543) (-0.800,-2.518) (-0.850,-2.492) (-0.900,-2.464) (-0.950,-2.434) (-1.000,-2.401) (-1.050,-2.365) (-1.100,-2.327) (-1.150,-2.285) (-1.189,-2.250) (-1.241,-2.200) (-1.289,-2.150) (-1.325,-2.109) (-1.373,-2.050) (-1.410,-2.000) (-1.444,-1.950) (-1.475,-1.900) (-1.504,-1.850) (-1.530,-1.800) (-1.554,-1.750) (-1.575,-1.702) (-1.595,-1.650) (-1.613,-1.600) (-1.628,-1.550) (-1.641,-1.500) (-1.652,-1.450) (-1.662,-1.400) (-1.669,-1.350) (-1.674,-1.300) (-1.678,-1.250) (-1.679,-1.200) (-1.679,-1.150) (-1.676,-1.100) (-1.672,-1.050) (-1.665,-1.000) (-1.656,-0.950) (-1.645,-0.900) (-1.632,-0.850) (-1.616,-0.800) (-1.597,-0.750) (-1.576,-0.700) (-1.552,-0.650) (-1.525,-0.601) (-1.493,-0.550) (-1.458,-0.500) (-1.425,-0.458) (-1.375,-0.401) (-1.325,-0.352) (-1.275,-0.308)}
\EQ{(1.275,-0.280) (1.325,-0.319) (1.375,-0.362) (1.414,-0.400) (1.450,-0.438) (1.500,-0.497) (1.539,-0.550) (1.572,-0.600) (1.600,-0.648) (1.625,-0.696) (1.650,-0.750) (1.670,-0.800) (1.688,-0.850) (1.703,-0.900) (1.716,-0.950) (1.726,-1.000) (1.734,-1.050) (1.741,-1.100) (1.745,-1.150) (1.747,-1.200) (1.748,-1.250) (1.746,-1.300) (1.743,-1.350) (1.738,-1.400) (1.731,-1.450) (1.722,-1.500) (1.711,-1.550) (1.698,-1.600) (1.683,-1.650) (1.667,-1.700) (1.648,-1.750) (1.627,-1.800) (1.604,-1.850) (1.579,-1.900) (1.552,-1.950) (1.525,-1.995) (1.489,-2.050) (1.454,-2.100) (1.416,-2.150) (1.375,-2.200) (1.330,-2.250) (1.282,-2.300) (1.229,-2.350) (1.175,-2.397) (1.125,-2.437) (1.075,-2.475) (1.025,-2.509) (0.975,-2.542) (0.925,-2.571) (0.875,-2.599) (0.825,-2.625) (0.775,-2.649) (0.725,-2.671) (0.675,-2.691) (0.625,-2.710) (0.575,-2.727) (0.525,-2.742) (0.475,-2.756) (0.425,-2.768) (0.375,-2.779) (0.325,-2.789) (0.275,-2.797) (0.225,-2.804) (0.175,-2.809) (0.125,-2.813) (0.075,-2.816) (0.025,-2.817) (-0.050,-2.817) (-0.100,-2.815) (-0.150,-2.811) (-0.200,-2.807) (-0.250,-2.800) (-0.300,-2.793) (-0.350,-2.784) (-0.400,-2.774) (-0.450,-2.762) (-0.500,-2.749) (-0.550,-2.735) (-0.600,-2.719) (-0.650,-2.701) (-0.700,-2.681) (-0.750,-2.660) (-0.800,-2.637) (-0.850,-2.612) (-0.900,-2.586) (-0.950,-2.557) (-1.000,-2.526) (-1.050,-2.492) (-1.100,-2.456) (-1.141,-2.425) (-1.200,-2.376) (-1.250,-2.331) (-1.300,-2.282) (-1.350,-2.228) (-1.396,-2.175) (-1.436,-2.125) (-1.472,-2.075) (-1.506,-2.025) (-1.537,-1.975) (-1.566,-1.925) (-1.592,-1.875) (-1.616,-1.825) (-1.638,-1.775) (-1.658,-1.725) (-1.675,-1.675) (-1.691,-1.625) (-1.705,-1.575) (-1.716,-1.525) (-1.726,-1.475) (-1.734,-1.425) (-1.741,-1.375) (-1.745,-1.325) (-1.747,-1.275) (-1.748,-1.225) (-1.746,-1.175) (-1.743,-1.125) (-1.738,-1.075) (-1.730,-1.025) (-1.721,-0.975) (-1.709,-0.925) (-1.695,-0.875) (-1.679,-0.825) (-1.660,-0.775) (-1.639,-0.725) (-1.614,-0.675) (-1.587,-0.625) (-1.556,-0.575) (-1.525,-0.530) (-1.482,-0.475) (-1.438,-0.425) (-1.400,-0.386) (-1.362,-0.350) (-1.301,-0.300) (-1.275,-0.280)}
\EQ{(1.288,-0.262) (1.342,-0.300) (1.400,-0.347) (1.450,-0.391) (1.500,-0.442) (1.550,-0.499) (1.590,-0.550) (1.625,-0.600) (1.656,-0.650) (1.684,-0.700) (1.708,-0.750) (1.730,-0.800) (1.750,-0.850) (1.767,-0.900) (1.781,-0.950) (1.794,-1.000) (1.804,-1.050) (1.812,-1.100) (1.818,-1.150) (1.823,-1.200) (1.825,-1.250) (1.825,-1.300) (1.824,-1.350) (1.821,-1.400) (1.816,-1.450) (1.810,-1.500) (1.801,-1.550) (1.791,-1.600) (1.779,-1.650) (1.765,-1.700) (1.750,-1.748) (1.732,-1.800) (1.712,-1.850) (1.691,-1.900) (1.667,-1.950) (1.641,-2.000) (1.613,-2.050) (1.583,-2.100) (1.550,-2.149) (1.514,-2.200) (1.476,-2.250) (1.434,-2.300) (1.400,-2.339) (1.366,-2.375) (1.325,-2.416) (1.289,-2.450) (1.250,-2.484) (1.201,-2.525) (1.150,-2.565) (1.101,-2.600) (1.050,-2.634) (1.000,-2.666) (0.950,-2.695) (0.900,-2.722) (0.850,-2.748) (0.800,-2.771) (0.750,-2.793) (0.700,-2.813) (0.650,-2.831) (0.600,-2.848) (0.550,-2.864) (0.500,-2.878) (0.450,-2.890) (0.400,-2.901) (0.350,-2.911) (0.300,-2.919) (0.250,-2.926) (0.200,-2.932) (0.150,-2.937) (0.100,-2.940) (0.050,-2.942) (0.000,-2.943) (-0.050,-2.942) (-0.100,-2.940) (-0.150,-2.937) (-0.200,-2.932) (-0.250,-2.926) (-0.300,-2.919) (-0.350,-2.911) (-0.400,-2.901) (-0.450,-2.890) (-0.500,-2.878) (-0.550,-2.864) (-0.600,-2.848) (-0.650,-2.831) (-0.700,-2.813) (-0.750,-2.793) (-0.800,-2.771) (-0.850,-2.748) (-0.900,-2.722) (-0.950,-2.695) (-1.000,-2.666) (-1.050,-2.634) (-1.100,-2.601) (-1.150,-2.565) (-1.200,-2.526) (-1.250,-2.484) (-1.289,-2.450) (-1.325,-2.416) (-1.366,-2.375) (-1.400,-2.339) (-1.434,-2.300) (-1.475,-2.251) (-1.514,-2.200) (-1.550,-2.150) (-1.583,-2.100) (-1.613,-2.050) (-1.641,-2.000) (-1.667,-1.950) (-1.691,-1.900) (-1.712,-1.850) (-1.732,-1.800) (-1.749,-1.750) (-1.765,-1.700) (-1.779,-1.650) (-1.791,-1.600) (-1.801,-1.550) (-1.810,-1.500) (-1.816,-1.450) (-1.821,-1.400) (-1.824,-1.350) (-1.825,-1.300) (-1.825,-1.250) (-1.823,-1.200) (-1.818,-1.150) (-1.812,-1.100) (-1.804,-1.050) (-1.794,-1.000) (-1.781,-0.950) (-1.767,-0.900) (-1.750,-0.851) (-1.730,-0.800) (-1.708,-0.750) (-1.684,-0.700) (-1.656,-0.650) (-1.625,-0.601) (-1.590,-0.550) (-1.551,-0.500) (-1.508,-0.450) (-1.459,-0.400) (-1.404,-0.350) (-1.350,-0.307) (-1.307,-0.275) (-1.288,-0.262)}
\EQ{(1.300,-0.243) (1.350,-0.275) (1.400,-0.310) (1.450,-0.348) (1.500,-0.391) (1.550,-0.439) (1.600,-0.493) (1.647,-0.550) (1.684,-0.600) (1.717,-0.650) (1.747,-0.700) (1.774,-0.750) (1.798,-0.800) (1.819,-0.850) (1.838,-0.900) (1.855,-0.950) (1.869,-1.000) (1.881,-1.050) (1.892,-1.100) (1.900,-1.150) (1.906,-1.200) (1.911,-1.250) (1.914,-1.300) (1.915,-1.350) (1.914,-1.400) (1.911,-1.450) (1.907,-1.500) (1.901,-1.550) (1.894,-1.600) (1.885,-1.650) (1.873,-1.700) (1.861,-1.750) (1.846,-1.800) (1.830,-1.850) (1.812,-1.900) (1.792,-1.950) (1.770,-2.000) (1.746,-2.050) (1.720,-2.100) (1.692,-2.150) (1.661,-2.200) (1.629,-2.250) (1.593,-2.300) (1.555,-2.350) (1.515,-2.400) (1.475,-2.445) (1.425,-2.498) (1.375,-2.548) (1.325,-2.593) (1.275,-2.636) (1.226,-2.675) (1.175,-2.713) (1.125,-2.747) (1.083,-2.775) (1.025,-2.810) (0.975,-2.838) (0.925,-2.865) (0.875,-2.890) (0.825,-2.913) (0.775,-2.934) (0.725,-2.954) (0.675,-2.972) (0.625,-2.989) (0.575,-3.004) (0.525,-3.018) (0.475,-3.030) (0.425,-3.042) (0.375,-3.052) (0.325,-3.060) (0.275,-3.068) (0.225,-3.074) (0.175,-3.079) (0.125,-3.082) (0.075,-3.085) (0.025,-3.086) (-0.050,-3.086) (-0.100,-3.084) (-0.150,-3.081) (-0.200,-3.076) (-0.250,-3.071) (-0.300,-3.064) (-0.350,-3.056) (-0.400,-3.047) (-0.450,-3.036) (-0.500,-3.024) (-0.550,-3.011) (-0.600,-2.996) (-0.650,-2.980) (-0.700,-2.963) (-0.750,-2.944) (-0.800,-2.923) (-0.850,-2.901) (-0.900,-2.877) (-0.950,-2.852) (-0.999,-2.825) (-1.050,-2.795) (-1.100,-2.764) (-1.150,-2.730) (-1.192,-2.700) (-1.250,-2.656) (-1.300,-2.615) (-1.345,-2.575) (-1.398,-2.525) (-1.448,-2.475) (-1.493,-2.425) (-1.535,-2.375) (-1.575,-2.325) (-1.611,-2.275) (-1.645,-2.225) (-1.675,-2.178) (-1.706,-2.125) (-1.733,-2.075) (-1.758,-2.025) (-1.781,-1.975) (-1.802,-1.925) (-1.821,-1.875) (-1.838,-1.825) (-1.854,-1.775) (-1.867,-1.725) (-1.879,-1.675) (-1.889,-1.625) (-1.898,-1.575) (-1.905,-1.525) (-1.910,-1.475) (-1.913,-1.425) (-1.915,-1.375) (-1.914,-1.325) (-1.913,-1.275) (-1.909,-1.225) (-1.903,-1.175) (-1.896,-1.125) (-1.887,-1.075) (-1.875,-1.025) (-1.862,-0.975) (-1.847,-0.925) (-1.829,-0.875) (-1.809,-0.825) (-1.786,-0.775) (-1.761,-0.725) (-1.733,-0.675) (-1.701,-0.625) (-1.666,-0.575) (-1.627,-0.525) (-1.584,-0.475) (-1.536,-0.425) (-1.482,-0.375) (-1.425,-0.328) (-1.375,-0.292) (-1.325,-0.258) (-1.300,-0.243)}
\EQ{(1.300,-0.216) (1.350,-0.244) (1.400,-0.274) (1.450,-0.308) (1.500,-0.345) (1.550,-0.385) (1.595,-0.425) (1.646,-0.475) (1.692,-0.525) (1.725,-0.565) (1.770,-0.625) (1.800,-0.669) (1.834,-0.725) (1.862,-0.775) (1.887,-0.825) (1.909,-0.875) (1.929,-0.925) (1.947,-0.975) (1.962,-1.025) (1.975,-1.075) (1.987,-1.125) (1.997,-1.175) (2.004,-1.225) (2.011,-1.275) (2.015,-1.325) (2.017,-1.375) (2.018,-1.425) (2.017,-1.475) (2.015,-1.525) (2.011,-1.575) (2.005,-1.625) (1.998,-1.675) (1.989,-1.725) (1.979,-1.775) (1.967,-1.825) (1.953,-1.875) (1.937,-1.925) (1.920,-1.975) (1.901,-2.025) (1.880,-2.075) (1.858,-2.125) (1.833,-2.175) (1.807,-2.225) (1.778,-2.275) (1.750,-2.321) (1.714,-2.375) (1.679,-2.425) (1.641,-2.475) (1.601,-2.525) (1.557,-2.575) (1.510,-2.625) (1.475,-2.661) (1.434,-2.700) (1.378,-2.750) (1.325,-2.794) (1.285,-2.825) (1.225,-2.869) (1.180,-2.900) (1.125,-2.935) (1.075,-2.965) (1.025,-2.993) (0.975,-3.020) (0.925,-3.045) (0.875,-3.068) (0.825,-3.089) (0.775,-3.109) (0.725,-3.128) (0.675,-3.145) (0.625,-3.161) (0.576,-3.175) (0.525,-3.188) (0.476,-3.200) (0.425,-3.211) (0.375,-3.220) (0.325,-3.229) (0.275,-3.236) (0.225,-3.241) (0.175,-3.246) (0.125,-3.249) (0.075,-3.252) (0.025,-3.253) (-0.050,-3.252) (-0.100,-3.251) (-0.150,-3.248) (-0.200,-3.244) (-0.250,-3.239) (-0.300,-3.232) (-0.350,-3.225) (-0.400,-3.216) (-0.450,-3.206) (-0.500,-3.195) (-0.550,-3.182) (-0.600,-3.168) (-0.650,-3.153) (-0.700,-3.137) (-0.750,-3.119) (-0.798,-3.100) (-0.850,-3.079) (-0.900,-3.056) (-0.950,-3.032) (-1.000,-3.007) (-1.050,-2.980) (-1.100,-2.950) (-1.150,-2.919) (-1.200,-2.886) (-1.250,-2.851) (-1.300,-2.814) (-1.348,-2.775) (-1.400,-2.731) (-1.450,-2.685) (-1.486,-2.650) (-1.525,-2.610) (-1.575,-2.555) (-1.621,-2.500) (-1.661,-2.450) (-1.697,-2.400) (-1.731,-2.350) (-1.763,-2.300) (-1.793,-2.250) (-1.820,-2.200) (-1.846,-2.150) (-1.869,-2.100) (-1.891,-2.050) (-1.911,-2.000) (-1.929,-1.950) (-1.945,-1.900) (-1.960,-1.850) (-1.973,-1.800) (-1.984,-1.750) (-1.994,-1.700) (-2.002,-1.650) (-2.008,-1.600) (-2.013,-1.550) (-2.017,-1.500) (-2.018,-1.450) (-2.018,-1.400) (-2.016,-1.350) (-2.013,-1.300) (-2.008,-1.250) (-2.001,-1.200) (-1.992,-1.150) (-1.982,-1.100) (-1.969,-1.050) (-1.955,-1.000) (-1.938,-0.950) (-1.919,-0.900) (-1.898,-0.850) (-1.875,-0.801) (-1.850,-0.753) (-1.825,-0.710) (-1.787,-0.650) (-1.752,-0.600) (-1.713,-0.550) (-1.675,-0.506) (-1.625,-0.454) (-1.575,-0.407) (-1.525,-0.364) (-1.475,-0.326) (-1.425,-0.291) (-1.375,-0.259) (-1.325,-0.230) (-1.300,-0.216)}
\EQ{(1.300,-0.191) (1.350,-0.215) (1.400,-0.241) (1.450,-0.270) (1.499,-0.300) (1.550,-0.335) (1.600,-0.372) (1.650,-0.414) (1.691,-0.450) (1.742,-0.500) (1.788,-0.550) (1.825,-0.594) (1.868,-0.650) (1.900,-0.696) (1.934,-0.750) (1.963,-0.800) (1.989,-0.850) (2.013,-0.900) (2.034,-0.950) (2.053,-1.000) (2.070,-1.050) (2.085,-1.100) (2.098,-1.150) (2.109,-1.200) (2.118,-1.250) (2.126,-1.300) (2.132,-1.350) (2.137,-1.400) (2.139,-1.450) (2.141,-1.500) (2.140,-1.550) (2.138,-1.600) (2.135,-1.650) (2.130,-1.700) (2.123,-1.750) (2.115,-1.800) (2.106,-1.850) (2.095,-1.900) (2.082,-1.950) (2.068,-2.000) (2.052,-2.050) (2.035,-2.100) (2.016,-2.150) (1.995,-2.200) (1.972,-2.250) (1.950,-2.296) (1.925,-2.344) (1.900,-2.389) (1.864,-2.450) (1.831,-2.500) (1.800,-2.546) (1.760,-2.600) (1.725,-2.645) (1.679,-2.700) (1.634,-2.750) (1.600,-2.786) (1.562,-2.825) (1.525,-2.860) (1.481,-2.900) (1.425,-2.947) (1.375,-2.987) (1.325,-3.025) (1.275,-3.060) (1.225,-3.093) (1.175,-3.124) (1.132,-3.150) (1.087,-3.175) (1.040,-3.200) (0.990,-3.225) (0.936,-3.250) (0.879,-3.275) (0.825,-3.297) (0.775,-3.315) (0.725,-3.333) (0.675,-3.349) (0.625,-3.364) (0.575,-3.377) (0.525,-3.390) (0.475,-3.401) (0.425,-3.411) (0.375,-3.420) (0.325,-3.427) (0.275,-3.434) (0.225,-3.439) (0.175,-3.444) (0.125,-3.447) (0.075,-3.449) (0.025,-3.450) (-0.037,-3.450) (-0.100,-3.448) (-0.150,-3.445) (-0.200,-3.442) (-0.250,-3.437) (-0.300,-3.431) (-0.350,-3.424) (-0.400,-3.415) (-0.450,-3.406) (-0.500,-3.395) (-0.550,-3.384) (-0.600,-3.371) (-0.650,-3.356) (-0.700,-3.341) (-0.748,-3.325) (-0.800,-3.306) (-0.850,-3.287) (-0.900,-3.266) (-0.950,-3.244) (-1.000,-3.220) (-1.050,-3.195) (-1.100,-3.168) (-1.150,-3.139) (-1.200,-3.109) (-1.250,-3.077) (-1.300,-3.043) (-1.350,-3.006) (-1.391,-2.975) (-1.450,-2.927) (-1.500,-2.883) (-1.550,-2.836) (-1.587,-2.800) (-1.625,-2.760) (-1.675,-2.705) (-1.721,-2.650) (-1.760,-2.600) (-1.797,-2.550) (-1.831,-2.500) (-1.864,-2.450) (-1.894,-2.400) (-1.922,-2.350) (-1.948,-2.300) (-1.972,-2.250) (-1.995,-2.200) (-2.016,-2.150) (-2.035,-2.100) (-2.052,-2.050) (-2.068,-2.000) (-2.082,-1.950) (-2.095,-1.900) (-2.106,-1.850) (-2.115,-1.800) (-2.123,-1.750) (-2.130,-1.700) (-2.135,-1.650) (-2.138,-1.600) (-2.140,-1.550) (-2.141,-1.500) (-2.139,-1.450) (-2.137,-1.400) (-2.132,-1.350) (-2.126,-1.300) (-2.118,-1.250) (-2.109,-1.200) (-2.098,-1.150) (-2.085,-1.100) (-2.070,-1.050) (-2.053,-1.000) (-2.034,-0.950) (-2.013,-0.900) (-1.989,-0.850) (-1.963,-0.800) (-1.934,-0.750) (-1.903,-0.700) (-1.868,-0.650) (-1.830,-0.600) (-1.788,-0.550) (-1.750,-0.508) (-1.700,-0.459) (-1.650,-0.414) (-1.603,-0.375) (-1.550,-0.335) (-1.500,-0.301) (-1.450,-0.270) (-1.400,-0.241) (-1.350,-0.215) (-1.300,-0.191)}
\EQ{(1.300,-0.165) (1.350,-0.186) (1.400,-0.208) (1.450,-0.233) (1.500,-0.259) (1.550,-0.288) (1.600,-0.319) (1.646,-0.350) (1.700,-0.390) (1.744,-0.425) (1.800,-0.474) (1.850,-0.522) (1.900,-0.575) (1.943,-0.625) (1.975,-0.665) (2.018,-0.725) (2.050,-0.773) (2.082,-0.825) (2.109,-0.875) (2.134,-0.925) (2.157,-0.975) (2.178,-1.025) (2.197,-1.075) (2.213,-1.125) (2.228,-1.175) (2.242,-1.225) (2.253,-1.275) (2.263,-1.325) (2.271,-1.375) (2.277,-1.425) (2.282,-1.475) (2.286,-1.525) (2.288,-1.575) (2.288,-1.625) (2.288,-1.675) (2.285,-1.725) (2.281,-1.775) (2.276,-1.825) (2.269,-1.875) (2.261,-1.925) (2.252,-1.975) (2.241,-2.025) (2.228,-2.075) (2.215,-2.125) (2.199,-2.175) (2.182,-2.225) (2.164,-2.275) (2.144,-2.325) (2.122,-2.375) (2.100,-2.423) (2.075,-2.473) (2.050,-2.520) (2.025,-2.565) (1.989,-2.625) (1.956,-2.675) (1.925,-2.720) (1.885,-2.775) (1.850,-2.820) (1.805,-2.875) (1.761,-2.925) (1.725,-2.964) (1.690,-3.000) (1.650,-3.039) (1.611,-3.075) (1.554,-3.125) (1.500,-3.170) (1.450,-3.209) (1.400,-3.245) (1.358,-3.275) (1.300,-3.313) (1.250,-3.344) (1.200,-3.373) (1.152,-3.400) (1.104,-3.425) (1.053,-3.450) (1.000,-3.475) (0.950,-3.497) (0.900,-3.517) (0.850,-3.536) (0.800,-3.554) (0.750,-3.571) (0.700,-3.587) (0.650,-3.601) (0.600,-3.614) (0.550,-3.626) (0.500,-3.637) (0.450,-3.647) (0.400,-3.656) (0.350,-3.664) (0.300,-3.670) (0.250,-3.676) (0.200,-3.680) (0.150,-3.684) (0.100,-3.687) (0.050,-3.688) (0.000,-3.689) (-0.050,-3.688) (-0.100,-3.687) (-0.150,-3.684) (-0.200,-3.680) (-0.250,-3.676) (-0.300,-3.670) (-0.350,-3.664) (-0.400,-3.656) (-0.450,-3.647) (-0.500,-3.637) (-0.550,-3.626) (-0.600,-3.614) (-0.650,-3.601) (-0.700,-3.587) (-0.750,-3.571) (-0.800,-3.554) (-0.850,-3.536) (-0.900,-3.517) (-0.950,-3.497) (-1.000,-3.475) (-1.050,-3.452) (-1.100,-3.427) (-1.150,-3.401) (-1.197,-3.375) (-1.250,-3.344) (-1.300,-3.313) (-1.350,-3.280) (-1.394,-3.250) (-1.450,-3.209) (-1.494,-3.175) (-1.550,-3.129) (-1.600,-3.085) (-1.638,-3.050) (-1.675,-3.015) (-1.714,-2.975) (-1.750,-2.937) (-1.800,-2.881) (-1.846,-2.825) (-1.885,-2.775) (-1.922,-2.725) (-1.956,-2.675) (-1.989,-2.625) (-2.019,-2.575) (-2.047,-2.525) (-2.074,-2.475) (-2.099,-2.425) (-2.122,-2.375) (-2.144,-2.325) (-2.164,-2.275) (-2.182,-2.225) (-2.199,-2.175) (-2.215,-2.125) (-2.228,-2.075) (-2.241,-2.025) (-2.252,-1.975) (-2.261,-1.925) (-2.269,-1.875) (-2.276,-1.825) (-2.281,-1.775) (-2.285,-1.725) (-2.288,-1.675) (-2.288,-1.625) (-2.288,-1.575) (-2.286,-1.525) (-2.282,-1.475) (-2.277,-1.425) (-2.271,-1.375) (-2.263,-1.325) (-2.253,-1.275) (-2.242,-1.225) (-2.228,-1.175) (-2.213,-1.125) (-2.197,-1.075) (-2.178,-1.025) (-2.157,-0.975) (-2.134,-0.925) (-2.109,-0.875) (-2.082,-0.825) (-2.051,-0.775) (-2.018,-0.725) (-1.982,-0.675) (-1.950,-0.634) (-1.900,-0.575) (-1.853,-0.525) (-1.801,-0.475) (-1.750,-0.430) (-1.700,-0.390) (-1.650,-0.353) (-1.600,-0.319) (-1.550,-0.288) (-1.500,-0.259) (-1.450,-0.233) (-1.400,-0.208) (-1.350,-0.186) (-1.300,-0.165)}
\EQ{(1.300,-0.141) (1.350,-0.158) (1.400,-0.177) (1.450,-0.197) (1.500,-0.220) (1.550,-0.243) (1.600,-0.269) (1.650,-0.297) (1.697,-0.325) (1.750,-0.359) (1.800,-0.394) (1.841,-0.425) (1.900,-0.473) (1.950,-0.517) (2.000,-0.566) (2.050,-0.619) (2.099,-0.675) (2.138,-0.725) (2.174,-0.775) (2.208,-0.825) (2.239,-0.875) (2.267,-0.925) (2.293,-0.975) (2.317,-1.025) (2.339,-1.075) (2.359,-1.125) (2.377,-1.175) (2.394,-1.225) (2.409,-1.275) (2.422,-1.325) (2.433,-1.375) (2.443,-1.425) (2.451,-1.475) (2.458,-1.525) (2.464,-1.575) (2.468,-1.625) (2.470,-1.675) (2.472,-1.725) (2.472,-1.775) (2.470,-1.825) (2.467,-1.875) (2.463,-1.925) (2.458,-1.975) (2.451,-2.025) (2.443,-2.075) (2.434,-2.125) (2.423,-2.175) (2.411,-2.225) (2.398,-2.275) (2.383,-2.325) (2.367,-2.375) (2.350,-2.423) (2.330,-2.475) (2.310,-2.525) (2.288,-2.575) (2.265,-2.625) (2.240,-2.675) (2.213,-2.725) (2.184,-2.775) (2.154,-2.825) (2.125,-2.871) (2.089,-2.925) (2.053,-2.975) (2.015,-3.025) (1.975,-3.075) (1.933,-3.125) (1.888,-3.175) (1.850,-3.215) (1.800,-3.265) (1.763,-3.300) (1.725,-3.335) (1.679,-3.375) (1.625,-3.420) (1.575,-3.459) (1.525,-3.496) (1.475,-3.531) (1.425,-3.565) (1.375,-3.597) (1.328,-3.625) (1.285,-3.650) (1.239,-3.675) (1.191,-3.700) (1.141,-3.725) (1.087,-3.750) (1.031,-3.775) (0.975,-3.798) (0.925,-3.817) (0.875,-3.835) (0.825,-3.852) (0.775,-3.868) (0.725,-3.883) (0.675,-3.897) (0.625,-3.909) (0.575,-3.921) (0.525,-3.932) (0.475,-3.941) (0.425,-3.950) (0.375,-3.958) (0.325,-3.964) (0.275,-3.970) (0.225,-3.975) (0.175,-3.978) (0.125,-3.981) (0.075,-3.983) (0.025,-3.984) (-0.050,-3.984) (-0.100,-3.982) (-0.150,-3.980) (-0.200,-3.977) (-0.250,-3.972) (-0.300,-3.967) (-0.350,-3.961) (-0.400,-3.954) (-0.450,-3.946) (-0.500,-3.937) (-0.550,-3.926) (-0.600,-3.915) (-0.650,-3.903) (-0.700,-3.890) (-0.750,-3.876) (-0.800,-3.860) (-0.850,-3.844) (-0.900,-3.826) (-0.950,-3.808) (-1.000,-3.788) (-1.050,-3.767) (-1.100,-3.744) (-1.150,-3.721) (-1.200,-3.696) (-1.250,-3.669) (-1.300,-3.641) (-1.350,-3.612) (-1.400,-3.581) (-1.448,-3.550) (-1.500,-3.514) (-1.550,-3.478) (-1.600,-3.440) (-1.650,-3.400) (-1.700,-3.357) (-1.750,-3.312) (-1.790,-3.275) (-1.840,-3.225) (-1.875,-3.189) (-1.910,-3.150) (-1.950,-3.105) (-1.995,-3.050) (-2.034,-3.000) (-2.071,-2.950) (-2.100,-2.908) (-2.139,-2.850) (-2.170,-2.800) (-2.199,-2.750) (-2.225,-2.703) (-2.250,-2.654) (-2.275,-2.603) (-2.299,-2.550) (-2.320,-2.500) (-2.340,-2.450) (-2.358,-2.400) (-2.375,-2.351) (-2.391,-2.300) (-2.405,-2.250) (-2.417,-2.200) (-2.429,-2.150) (-2.439,-2.100) (-2.447,-2.050) (-2.455,-2.000) (-2.461,-1.950) (-2.466,-1.900) (-2.469,-1.850) (-2.471,-1.800) (-2.472,-1.750) (-2.471,-1.700) (-2.469,-1.650) (-2.466,-1.600) (-2.461,-1.550) (-2.455,-1.500) (-2.447,-1.450) (-2.438,-1.400) (-2.427,-1.350) (-2.415,-1.300) (-2.401,-1.250) (-2.386,-1.200) (-2.369,-1.150) (-2.350,-1.101) (-2.329,-1.050) (-2.306,-1.000) (-2.281,-0.950) (-2.253,-0.900) (-2.225,-0.852) (-2.192,-0.800) (-2.157,-0.750) (-2.125,-0.708) (-2.078,-0.650) (-2.033,-0.600) (-1.984,-0.550) (-1.931,-0.500) (-1.875,-0.452) (-1.825,-0.413) (-1.775,-0.376) (-1.725,-0.343) (-1.675,-0.312) (-1.625,-0.283) (-1.575,-0.256) (-1.525,-0.231) (-1.475,-0.208) (-1.425,-0.187) (-1.375,-0.167) (-1.327,-0.150) (-1.300,-0.141)}
\EQ{(1.300,-0.117) (1.350,-0.131) (1.400,-0.146) (1.450,-0.163) (1.500,-0.181) (1.549,-0.200) (1.600,-0.221) (1.650,-0.244) (1.700,-0.268) (1.750,-0.293) (1.800,-0.321) (1.849,-0.350) (1.900,-0.382) (1.950,-0.416) (1.997,-0.450) (2.050,-0.492) (2.090,-0.525) (2.145,-0.575) (2.196,-0.625) (2.243,-0.675) (2.287,-0.725) (2.325,-0.772) (2.365,-0.825) (2.400,-0.875) (2.432,-0.925) (2.462,-0.975) (2.490,-1.025) (2.516,-1.075) (2.540,-1.125) (2.562,-1.175) (2.582,-1.225) (2.600,-1.275) (2.617,-1.325) (2.632,-1.375) (2.646,-1.425) (2.658,-1.475) (2.669,-1.525) (2.679,-1.575) (2.687,-1.625) (2.693,-1.675) (2.699,-1.725) (2.703,-1.775) (2.706,-1.825) (2.707,-1.875) (2.708,-1.925) (2.707,-1.975) (2.705,-2.025) (2.701,-2.075) (2.697,-2.125) (2.691,-2.175) (2.684,-2.225) (2.676,-2.275) (2.667,-2.325) (2.656,-2.375) (2.644,-2.425) (2.631,-2.475) (2.617,-2.525) (2.602,-2.575) (2.585,-2.625) (2.567,-2.675) (2.547,-2.725) (2.527,-2.775) (2.504,-2.825) (2.481,-2.875) (2.456,-2.925) (2.430,-2.975) (2.402,-3.025) (2.375,-3.070) (2.341,-3.125) (2.308,-3.175) (2.275,-3.223) (2.237,-3.275) (2.200,-3.323) (2.158,-3.375) (2.125,-3.414) (2.075,-3.470) (2.025,-3.523) (1.975,-3.574) (1.925,-3.621) (1.875,-3.667) (1.836,-3.700) (1.776,-3.750) (1.725,-3.790) (1.677,-3.825) (1.625,-3.862) (1.575,-3.896) (1.530,-3.925) (1.475,-3.959) (1.425,-3.988) (1.375,-4.016) (1.325,-4.043) (1.275,-4.068) (1.225,-4.092) (1.175,-4.115) (1.125,-4.137) (1.075,-4.158) (1.025,-4.177) (0.975,-4.196) (0.925,-4.213) (0.875,-4.229) (0.825,-4.245) (0.775,-4.259) (0.725,-4.272) (0.675,-4.285) (0.625,-4.296) (0.575,-4.307) (0.525,-4.317) (0.475,-4.325) (0.425,-4.333) (0.375,-4.340) (0.325,-4.346) (0.275,-4.351) (0.225,-4.356) (0.175,-4.359) (0.125,-4.362) (0.075,-4.363) (0.025,-4.364) (-0.050,-4.364) (-0.100,-4.363) (-0.150,-4.361) (-0.200,-4.358) (-0.250,-4.354) (-0.300,-4.349) (-0.350,-4.343) (-0.400,-4.337) (-0.450,-4.329) (-0.500,-4.321) (-0.550,-4.312) (-0.600,-4.302) (-0.650,-4.291) (-0.700,-4.279) (-0.750,-4.266) (-0.800,-4.252) (-0.850,-4.237) (-0.900,-4.221) (-0.950,-4.204) (-1.000,-4.187) (-1.050,-4.168) (-1.100,-4.147) (-1.150,-4.126) (-1.200,-4.104) (-1.250,-4.080) (-1.300,-4.056) (-1.350,-4.030) (-1.400,-4.002) (-1.448,-3.975) (-1.500,-3.944) (-1.550,-3.912) (-1.600,-3.879) (-1.642,-3.850) (-1.700,-3.808) (-1.744,-3.775) (-1.800,-3.730) (-1.850,-3.688) (-1.894,-3.650) (-1.948,-3.600) (-1.999,-3.550) (-2.047,-3.500) (-2.093,-3.450) (-2.137,-3.400) (-2.175,-3.354) (-2.218,-3.300) (-2.250,-3.257) (-2.291,-3.200) (-2.325,-3.150) (-2.357,-3.100) (-2.387,-3.050) (-2.416,-3.000) (-2.443,-2.950) (-2.469,-2.900) (-2.493,-2.850) (-2.516,-2.800) (-2.537,-2.750) (-2.557,-2.700) (-2.575,-2.653) (-2.593,-2.600) (-2.610,-2.550) (-2.624,-2.500) (-2.638,-2.450) (-2.650,-2.400) (-2.662,-2.350) (-2.672,-2.300) (-2.680,-2.250) (-2.688,-2.200) (-2.694,-2.150) (-2.699,-2.100) (-2.703,-2.050) (-2.706,-2.000) (-2.707,-1.950) (-2.708,-1.900) (-2.707,-1.850) (-2.704,-1.800) (-2.701,-1.750) (-2.696,-1.700) (-2.690,-1.650) (-2.683,-1.600) (-2.674,-1.550) (-2.664,-1.500) (-2.652,-1.450) (-2.639,-1.400) (-2.625,-1.351) (-2.609,-1.300) (-2.591,-1.250) (-2.572,-1.200) (-2.551,-1.150) (-2.528,-1.100) (-2.503,-1.050) (-2.476,-1.000) (-2.450,-0.954) (-2.416,-0.900) (-2.383,-0.850) (-2.350,-0.805) (-2.308,-0.750) (-2.275,-0.711) (-2.225,-0.655) (-2.175,-0.604) (-2.125,-0.557) (-2.075,-0.513) (-2.029,-0.475) (-1.975,-0.434) (-1.927,-0.400) (-1.875,-0.366) (-1.825,-0.335) (-1.775,-0.307) (-1.725,-0.280) (-1.675,-0.255) (-1.625,-0.232) (-1.575,-0.211) (-1.525,-0.191) (-1.475,-0.172) (-1.425,-0.155) (-1.375,-0.139) (-1.325,-0.124) (-1.300,-0.117)}
\EQ{(1.300,-0.093) (1.350,-0.104) (1.400,-0.116) (1.450,-0.130) (1.500,-0.144) (1.550,-0.159) (1.600,-0.175) (1.650,-0.193) (1.700,-0.211) (1.750,-0.231) (1.800,-0.252) (1.850,-0.275) (1.900,-0.299) (1.950,-0.324) (1.997,-0.350) (2.041,-0.375) (2.100,-0.412) (2.150,-0.444) (2.194,-0.475) (2.250,-0.517) (2.292,-0.550) (2.350,-0.599) (2.400,-0.645) (2.450,-0.693) (2.500,-0.746) (2.548,-0.800) (2.589,-0.850) (2.625,-0.896) (2.664,-0.950) (2.697,-1.000) (2.725,-1.044) (2.759,-1.100) (2.786,-1.150) (2.812,-1.200) (2.836,-1.250) (2.858,-1.300) (2.879,-1.350) (2.898,-1.400) (2.916,-1.450) (2.932,-1.500) (2.947,-1.550) (2.960,-1.600) (2.972,-1.650) (2.983,-1.700) (2.993,-1.750) (3.001,-1.800) (3.008,-1.850) (3.014,-1.900) (3.019,-1.950) (3.023,-2.000) (3.025,-2.050) (3.027,-2.100) (3.027,-2.150) (3.026,-2.200) (3.025,-2.250) (3.022,-2.300) (3.018,-2.350) (3.013,-2.400) (3.006,-2.450) (2.999,-2.500) (2.991,-2.550) (2.981,-2.600) (2.971,-2.650) (2.959,-2.700) (2.947,-2.750) (2.933,-2.800) (2.918,-2.850) (2.902,-2.900) (2.885,-2.950) (2.866,-3.000) (2.847,-3.050) (2.826,-3.100) (2.804,-3.150) (2.781,-3.200) (2.756,-3.250) (2.730,-3.300) (2.703,-3.350) (2.675,-3.399) (2.645,-3.450) (2.613,-3.500) (2.580,-3.550) (2.550,-3.594) (2.509,-3.650) (2.475,-3.696) (2.432,-3.750) (2.400,-3.789) (2.350,-3.846) (2.301,-3.900) (2.253,-3.950) (2.203,-4.000) (2.151,-4.050) (2.100,-4.096) (2.050,-4.139) (2.007,-4.175) (1.950,-4.220) (1.900,-4.258) (1.850,-4.295) (1.807,-4.325) (1.750,-4.363) (1.700,-4.395) (1.651,-4.425) (1.600,-4.455) (1.550,-4.483) (1.500,-4.510) (1.450,-4.536) (1.400,-4.561) (1.350,-4.585) (1.300,-4.607) (1.250,-4.629) (1.200,-4.649) (1.150,-4.669) (1.100,-4.687) (1.050,-4.705) (1.000,-4.722) (0.950,-4.737) (0.900,-4.752) (0.850,-4.766) (0.800,-4.779) (0.750,-4.792) (0.700,-4.803) (0.650,-4.814) (0.600,-4.824) (0.550,-4.833) (0.500,-4.841) (0.450,-4.848) (0.400,-4.855) (0.350,-4.861) (0.300,-4.866) (0.250,-4.870) (0.200,-4.873) (0.150,-4.876) (0.100,-4.878) (0.050,-4.879) (0.000,-4.880) (-0.050,-4.879) (-0.100,-4.878) (-0.150,-4.876) (-0.200,-4.873) (-0.250,-4.870) (-0.300,-4.866) (-0.350,-4.861) (-0.400,-4.855) (-0.450,-4.848) (-0.500,-4.841) (-0.550,-4.833) (-0.600,-4.824) (-0.650,-4.814) (-0.700,-4.803) (-0.750,-4.792) (-0.800,-4.779) (-0.850,-4.766) (-0.900,-4.752) (-0.950,-4.737) (-1.000,-4.722) (-1.050,-4.705) (-1.100,-4.687) (-1.150,-4.669) (-1.198,-4.650) (-1.250,-4.629) (-1.300,-4.607) (-1.350,-4.585) (-1.400,-4.561) (-1.450,-4.536) (-1.500,-4.510) (-1.550,-4.483) (-1.600,-4.455) (-1.650,-4.426) (-1.700,-4.395) (-1.750,-4.363) (-1.800,-4.330) (-1.843,-4.300) (-1.900,-4.258) (-1.944,-4.225) (-2.000,-4.181) (-2.050,-4.139) (-2.095,-4.100) (-2.150,-4.051) (-2.200,-4.003) (-2.250,-3.954) (-2.300,-3.901) (-2.347,-3.850) (-2.390,-3.800) (-2.425,-3.758) (-2.472,-3.700) (-2.509,-3.650) (-2.546,-3.600) (-2.575,-3.558) (-2.613,-3.500) (-2.645,-3.450) (-2.675,-3.400) (-2.700,-3.356) (-2.725,-3.310) (-2.750,-3.262) (-2.775,-3.212) (-2.800,-3.159) (-2.825,-3.102) (-2.847,-3.050) (-2.866,-3.000) (-2.885,-2.950) (-2.902,-2.900) (-2.918,-2.850) (-2.933,-2.800) (-2.947,-2.750) (-2.959,-2.700) (-2.971,-2.650) (-2.981,-2.600) (-2.991,-2.550) (-2.999,-2.500) (-3.006,-2.450) (-3.013,-2.400) (-3.018,-2.350) (-3.022,-2.300) (-3.025,-2.250) (-3.026,-2.200) (-3.027,-2.150) (-3.027,-2.100) (-3.025,-2.050) (-3.023,-2.000) (-3.019,-1.950) (-3.014,-1.900) (-3.008,-1.850) (-3.001,-1.800) (-2.993,-1.750) (-2.983,-1.700) (-2.972,-1.650) (-2.960,-1.600) (-2.947,-1.550) (-2.932,-1.500) (-2.916,-1.450) (-2.898,-1.400) (-2.879,-1.350) (-2.858,-1.300) (-2.836,-1.250) (-2.812,-1.200) (-2.786,-1.150) (-2.759,-1.100) (-2.729,-1.050) (-2.700,-1.004) (-2.664,-0.950) (-2.628,-0.900) (-2.589,-0.850) (-2.550,-0.802) (-2.504,-0.750) (-2.457,-0.700) (-2.406,-0.650) (-2.351,-0.600) (-2.300,-0.557) (-2.261,-0.525) (-2.200,-0.479) (-2.158,-0.450) (-2.100,-0.412) (-2.050,-0.381) (-2.000,-0.352) (-1.951,-0.325) (-1.902,-0.300) (-1.850,-0.275) (-1.800,-0.252) (-1.750,-0.231) (-1.700,-0.211) (-1.650,-0.193) (-1.600,-0.175) (-1.550,-0.159) (-1.500,-0.144) (-1.450,-0.130) (-1.400,-0.116) (-1.350,-0.104) (-1.300,-0.093)}
\EQ{(1.300,-0.070) (1.350,-0.078) (1.400,-0.087) (1.450,-0.097) (1.500,-0.107) (1.550,-0.118) (1.600,-0.130) (1.650,-0.143) (1.700,-0.157) (1.750,-0.171) (1.800,-0.187) (1.850,-0.203) (1.900,-0.220) (1.950,-0.239) (2.000,-0.258) (2.050,-0.279) (2.098,-0.300) (2.150,-0.324) (2.200,-0.348) (2.250,-0.374) (2.298,-0.400) (2.342,-0.425) (2.400,-0.460) (2.450,-0.492) (2.500,-0.525) (2.550,-0.561) (2.600,-0.598) (2.650,-0.638) (2.694,-0.675) (2.750,-0.724) (2.800,-0.772) (2.850,-0.822) (2.899,-0.875) (2.943,-0.925) (2.985,-0.975) (3.024,-1.025) (3.060,-1.075) (3.095,-1.125) (3.125,-1.171) (3.159,-1.225) (3.188,-1.275) (3.216,-1.325) (3.241,-1.375) (3.266,-1.425) (3.289,-1.475) (3.310,-1.525) (3.330,-1.575) (3.349,-1.625) (3.367,-1.675) (3.383,-1.725) (3.398,-1.775) (3.412,-1.825) (3.425,-1.875) (3.437,-1.925) (3.447,-1.975) (3.456,-2.025) (3.465,-2.075) (3.472,-2.125) (3.479,-2.175) (3.484,-2.225) (3.488,-2.275) (3.491,-2.325) (3.494,-2.375) (3.495,-2.425) (3.496,-2.475) (3.495,-2.525) (3.494,-2.575) (3.491,-2.625) (3.488,-2.675) (3.483,-2.725) (3.478,-2.775) (3.472,-2.825) (3.465,-2.875) (3.457,-2.925) (3.448,-2.975) (3.438,-3.025) (3.427,-3.075) (3.415,-3.125) (3.403,-3.175) (3.389,-3.225) (3.374,-3.275) (3.359,-3.325) (3.342,-3.375) (3.325,-3.423) (3.306,-3.475) (3.286,-3.525) (3.265,-3.575) (3.243,-3.625) (3.220,-3.675) (3.196,-3.725) (3.171,-3.775) (3.145,-3.825) (3.118,-3.875) (3.089,-3.925) (3.059,-3.975) (3.028,-4.025) (3.000,-4.068) (2.962,-4.125) (2.927,-4.175) (2.890,-4.225) (2.852,-4.275) (2.812,-4.325) (2.775,-4.370) (2.728,-4.425) (2.683,-4.475) (2.637,-4.525) (2.600,-4.563) (2.563,-4.600) (2.525,-4.637) (2.485,-4.675) (2.430,-4.725) (2.375,-4.773) (2.325,-4.815) (2.281,-4.850) (2.225,-4.894) (2.183,-4.925) (2.125,-4.967) (2.078,-5.000) (2.025,-5.035) (1.975,-5.067) (1.925,-5.098) (1.881,-5.125) (1.825,-5.157) (1.775,-5.185) (1.725,-5.212) (1.675,-5.238) (1.625,-5.262) (1.575,-5.286) (1.525,-5.309) (1.475,-5.331) (1.425,-5.352) (1.375,-5.373) (1.325,-5.392) (1.275,-5.411) (1.225,-5.428) (1.175,-5.445) (1.125,-5.462) (1.075,-5.477) (1.025,-5.492) (0.975,-5.505) (0.925,-5.519) (0.875,-5.531) (0.825,-5.543) (0.775,-5.553) (0.725,-5.564) (0.675,-5.573) (0.625,-5.582) (0.575,-5.590) (0.525,-5.597) (0.475,-5.604) (0.425,-5.610) (0.375,-5.616) (0.325,-5.620) (0.275,-5.624) (0.225,-5.628) (0.175,-5.630) (0.125,-5.632) (0.075,-5.634) (0.025,-5.634) (-0.050,-5.634) (-0.100,-5.633) (-0.150,-5.631) (-0.200,-5.629) (-0.250,-5.626) (-0.300,-5.622) (-0.350,-5.618) (-0.400,-5.613) (-0.450,-5.607) (-0.500,-5.601) (-0.550,-5.594) (-0.600,-5.586) (-0.650,-5.578) (-0.700,-5.569) (-0.750,-5.559) (-0.800,-5.548) (-0.850,-5.537) (-0.899,-5.525) (-0.950,-5.512) (-1.000,-5.499) (-1.050,-5.484) (-1.100,-5.469) (-1.150,-5.453) (-1.200,-5.437) (-1.250,-5.420) (-1.300,-5.401) (-1.350,-5.382) (-1.400,-5.363) (-1.450,-5.342) (-1.500,-5.320) (-1.550,-5.298) (-1.599,-5.275) (-1.650,-5.250) (-1.700,-5.225) (-1.747,-5.200) (-1.793,-5.175) (-1.838,-5.150) (-1.900,-5.114) (-1.950,-5.083) (-2.000,-5.051) (-2.050,-5.019) (-2.100,-4.985) (-2.149,-4.950) (-2.200,-4.913) (-2.249,-4.875) (-2.300,-4.835) (-2.343,-4.800) (-2.400,-4.751) (-2.450,-4.707) (-2.500,-4.661) (-2.538,-4.625) (-2.575,-4.588) (-2.613,-4.550) (-2.650,-4.511) (-2.700,-4.457) (-2.750,-4.400) (-2.792,-4.350) (-2.825,-4.309) (-2.871,-4.250) (-2.908,-4.200) (-2.944,-4.150) (-2.975,-4.106) (-3.012,-4.050) (-3.044,-4.000) (-3.074,-3.950) (-3.100,-3.906) (-3.125,-3.862) (-3.150,-3.816) (-3.175,-3.768) (-3.200,-3.718) (-3.225,-3.665) (-3.250,-3.610) (-3.275,-3.551) (-3.296,-3.500) (-3.315,-3.450) (-3.333,-3.400) (-3.350,-3.351) (-3.366,-3.300) (-3.382,-3.250) (-3.396,-3.200) (-3.409,-3.150) (-3.421,-3.100) (-3.433,-3.050) (-3.443,-3.000) (-3.453,-2.950) (-3.461,-2.900) (-3.469,-2.850) (-3.475,-2.800) (-3.481,-2.750) (-3.486,-2.700) (-3.490,-2.650) (-3.492,-2.600) (-3.494,-2.550) (-3.495,-2.500) (-3.496,-2.450) (-3.495,-2.400) (-3.493,-2.350) (-3.490,-2.300) (-3.486,-2.250) (-3.481,-2.200) (-3.475,-2.150) (-3.469,-2.100) (-3.461,-2.050) (-3.452,-2.000) (-3.442,-1.950) (-3.431,-1.900) (-3.419,-1.850) (-3.405,-1.800) (-3.391,-1.750) (-3.375,-1.700) (-3.358,-1.650) (-3.340,-1.600) (-3.321,-1.550) (-3.300,-1.501) (-3.277,-1.450) (-3.254,-1.400) (-3.229,-1.350) (-3.202,-1.300) (-3.175,-1.252) (-3.144,-1.200) (-3.112,-1.150) (-3.078,-1.100) (-3.042,-1.050) (-3.004,-1.000) (-2.964,-0.950) (-2.925,-0.904) (-2.876,-0.850) (-2.828,-0.800) (-2.777,-0.750) (-2.725,-0.702) (-2.675,-0.659) (-2.634,-0.625) (-2.575,-0.579) (-2.525,-0.543) (-2.475,-0.508) (-2.425,-0.475) (-2.375,-0.444) (-2.325,-0.415) (-2.275,-0.387) (-2.225,-0.361) (-2.175,-0.336) (-2.125,-0.312) (-2.075,-0.290) (-2.025,-0.269) (-1.975,-0.248) (-1.925,-0.230) (-1.875,-0.212) (-1.825,-0.195) (-1.775,-0.179) (-1.725,-0.164) (-1.676,-0.150) (-1.625,-0.137) (-1.575,-0.124) (-1.525,-0.113) (-1.475,-0.102) (-1.425,-0.092) (-1.375,-0.082) (-1.325,-0.074) (-1.300,-0.070)}
\EQ{(1.300,-0.046) (1.350,-0.052) (1.400,-0.058) (1.450,-0.064) (1.500,-0.071) (1.550,-0.079) (1.600,-0.086) (1.650,-0.095) (1.700,-0.104) (1.750,-0.113) (1.800,-0.123) (1.850,-0.134) (1.900,-0.145) (1.950,-0.157) (2.000,-0.170) (2.050,-0.183) (2.100,-0.197) (2.150,-0.212) (2.200,-0.227) (2.250,-0.243) (2.300,-0.260) (2.350,-0.278) (2.400,-0.297) (2.450,-0.317) (2.500,-0.337) (2.550,-0.359) (2.600,-0.381) (2.650,-0.405) (2.700,-0.429) (2.750,-0.455) (2.800,-0.482) (2.850,-0.510) (2.900,-0.539) (2.950,-0.570) (2.998,-0.600) (3.050,-0.635) (3.100,-0.670) (3.141,-0.700) (3.200,-0.745) (3.250,-0.785) (3.298,-0.825) (3.350,-0.871) (3.400,-0.917) (3.450,-0.966) (3.500,-1.017) (3.550,-1.071) (3.598,-1.125) (3.640,-1.175) (3.675,-1.219) (3.718,-1.275) (3.750,-1.319) (3.789,-1.375) (3.822,-1.425) (3.850,-1.470) (3.883,-1.525) (3.912,-1.575) (3.939,-1.625) (3.965,-1.675) (3.990,-1.725) (4.014,-1.775) (4.036,-1.825) (4.057,-1.875) (4.077,-1.925) (4.096,-1.975) (4.114,-2.025) (4.131,-2.075) (4.147,-2.125) (4.162,-2.175) (4.176,-2.225) (4.189,-2.275) (4.201,-2.325) (4.213,-2.375) (4.223,-2.425) (4.232,-2.475) (4.241,-2.525) (4.249,-2.575) (4.256,-2.625) (4.262,-2.675) (4.267,-2.725) (4.271,-2.775) (4.275,-2.825) (4.278,-2.875) (4.280,-2.925) (4.281,-2.975) (4.281,-3.025) (4.281,-3.075) (4.280,-3.125) (4.278,-3.175) (4.275,-3.225) (4.272,-3.275) (4.268,-3.325) (4.263,-3.375) (4.257,-3.425) (4.250,-3.475) (4.243,-3.525) (4.235,-3.575) (4.226,-3.625) (4.217,-3.675) (4.206,-3.725) (4.195,-3.775) (4.184,-3.825) (4.171,-3.875) (4.157,-3.925) (4.143,-3.975) (4.128,-4.025) (4.112,-4.075) (4.096,-4.125) (4.078,-4.175) (4.060,-4.225) (4.041,-4.275) (4.021,-4.325) (4.000,-4.375) (3.979,-4.425) (3.956,-4.475) (3.933,-4.525) (3.908,-4.575) (3.883,-4.625) (3.857,-4.675) (3.830,-4.725) (3.802,-4.775) (3.775,-4.821) (3.743,-4.875) (3.711,-4.925) (3.679,-4.975) (3.650,-5.019) (3.611,-5.075) (3.576,-5.125) (3.539,-5.175) (3.501,-5.225) (3.462,-5.275) (3.425,-5.321) (3.380,-5.375) (3.336,-5.425) (3.300,-5.466) (3.250,-5.520) (3.200,-5.573) (3.150,-5.623) (3.100,-5.672) (3.050,-5.720) (3.000,-5.766) (2.961,-5.800) (2.904,-5.850) (2.850,-5.895) (2.800,-5.935) (2.750,-5.974) (2.700,-6.013) (2.650,-6.049) (2.600,-6.085) (2.550,-6.120) (2.505,-6.150) (2.450,-6.186) (2.400,-6.218) (2.350,-6.249) (2.306,-6.275) (2.250,-6.308) (2.200,-6.336) (2.150,-6.363) (2.100,-6.390) (2.050,-6.415) (2.000,-6.440) (1.950,-6.464) (1.900,-6.488) (1.850,-6.510) (1.800,-6.532) (1.750,-6.553) (1.700,-6.574) (1.650,-6.593) (1.600,-6.612) (1.550,-6.631) (1.500,-6.648) (1.450,-6.665) (1.400,-6.682) (1.350,-6.698) (1.300,-6.713) (1.250,-6.727) (1.200,-6.741) (1.150,-6.754) (1.100,-6.767) (1.050,-6.779) (1.000,-6.791) (0.950,-6.801) (0.900,-6.812) (0.850,-6.821) (0.800,-6.831) (0.750,-6.839) (0.700,-6.847) (0.650,-6.854) (0.600,-6.861) (0.550,-6.868) (0.500,-6.873) (0.450,-6.879) (0.400,-6.883) (0.350,-6.887) (0.300,-6.891) (0.250,-6.894) (0.200,-6.896) (0.150,-6.898) (0.100,-6.900) (0.050,-6.900) (0.000,-6.901) (-0.050,-6.900) (-0.100,-6.900) (-0.150,-6.898) (-0.200,-6.896) (-0.250,-6.894) (-0.300,-6.891) (-0.350,-6.887) (-0.400,-6.883) (-0.450,-6.879) (-0.500,-6.873) (-0.550,-6.868) (-0.600,-6.861) (-0.650,-6.854) (-0.700,-6.847) (-0.750,-6.839) (-0.800,-6.831) (-0.850,-6.821) (-0.900,-6.812) (-0.950,-6.801) (-1.000,-6.791) (-1.050,-6.779) (-1.100,-6.767) (-1.150,-6.754) (-1.200,-6.741) (-1.250,-6.727) (-1.300,-6.713) (-1.350,-6.698) (-1.400,-6.682) (-1.450,-6.665) (-1.500,-6.648) (-1.550,-6.631) (-1.600,-6.612) (-1.650,-6.593) (-1.697,-6.575) (-1.750,-6.553) (-1.800,-6.532) (-1.850,-6.510) (-1.900,-6.488) (-1.950,-6.464) (-2.000,-6.440) (-2.050,-6.415) (-2.100,-6.390) (-2.150,-6.363) (-2.200,-6.336) (-2.250,-6.308) (-2.300,-6.279) (-2.348,-6.250) (-2.400,-6.218) (-2.450,-6.186) (-2.500,-6.154) (-2.543,-6.125) (-2.600,-6.085) (-2.649,-6.050) (-2.700,-6.013) (-2.749,-5.975) (-2.800,-5.935) (-2.844,-5.900) (-2.900,-5.853) (-2.950,-5.810) (-2.990,-5.775) (-3.044,-5.725) (-3.097,-5.675) (-3.148,-5.625) (-3.198,-5.575) (-3.246,-5.525) (-3.292,-5.475) (-3.336,-5.425) (-3.375,-5.380) (-3.421,-5.325) (-3.462,-5.275) (-3.500,-5.227) (-3.539,-5.175) (-3.575,-5.126) (-3.611,-5.075) (-3.646,-5.025) (-3.675,-4.981) (-3.711,-4.925) (-3.743,-4.875) (-3.773,-4.825) (-3.800,-4.778) (-3.825,-4.734) (-3.850,-4.688) (-3.875,-4.641) (-3.900,-4.592) (-3.925,-4.541) (-3.950,-4.488) (-3.975,-4.433) (-4.000,-4.376) (-4.021,-4.325) (-4.041,-4.275) (-4.060,-4.225) (-4.078,-4.175) (-4.096,-4.125) (-4.112,-4.075) (-4.128,-4.025) (-4.143,-3.975) (-4.157,-3.925) (-4.171,-3.875) (-4.184,-3.825) (-4.195,-3.775) (-4.206,-3.725) (-4.217,-3.675) (-4.226,-3.625) (-4.235,-3.575) (-4.243,-3.525) (-4.250,-3.475) (-4.257,-3.425) (-4.263,-3.375) (-4.268,-3.325) (-4.272,-3.275) (-4.275,-3.225) (-4.278,-3.175) (-4.280,-3.125) (-4.281,-3.075) (-4.281,-3.025) (-4.281,-2.975) (-4.280,-2.925) (-4.278,-2.875) (-4.275,-2.825) (-4.271,-2.775) (-4.267,-2.725) (-4.262,-2.675) (-4.256,-2.625) (-4.249,-2.575) (-4.241,-2.525) (-4.232,-2.475) (-4.223,-2.425) (-4.213,-2.375) (-4.201,-2.325) (-4.189,-2.275) (-4.176,-2.225) (-4.162,-2.175) (-4.147,-2.125) (-4.131,-2.075) (-4.114,-2.025) (-4.096,-1.975) (-4.077,-1.925) (-4.057,-1.875) (-4.036,-1.825) (-4.014,-1.775) (-3.990,-1.725) (-3.965,-1.675) (-3.939,-1.625) (-3.912,-1.575) (-3.883,-1.525) (-3.853,-1.475) (-3.825,-1.430) (-3.789,-1.375) (-3.754,-1.325) (-3.718,-1.275) (-3.680,-1.225) (-3.640,-1.175) (-3.600,-1.127) (-3.554,-1.075) (-3.508,-1.025) (-3.459,-0.975) (-3.408,-0.925) (-3.355,-0.875) (-3.300,-0.827) (-3.250,-0.785) (-3.207,-0.750) (-3.150,-0.707) (-3.107,-0.675) (-3.050,-0.635) (-3.000,-0.602) (-2.950,-0.570) (-2.900,-0.539) (-2.850,-0.510) (-2.800,-0.482) (-2.750,-0.455) (-2.700,-0.429) (-2.650,-0.405) (-2.600,-0.381) (-2.550,-0.359) (-2.500,-0.337) (-2.450,-0.317) (-2.400,-0.297) (-2.350,-0.278) (-2.300,-0.260) (-2.250,-0.243) (-2.200,-0.227) (-2.150,-0.212) (-2.100,-0.197) (-2.050,-0.183) (-2.000,-0.170) (-1.950,-0.157) (-1.900,-0.145) (-1.850,-0.134) (-1.800,-0.123) (-1.750,-0.113) (-1.700,-0.104) (-1.650,-0.095) (-1.600,-0.086) (-1.550,-0.079) (-1.500,-0.071) (-1.450,-0.064) (-1.400,-0.058) (-1.350,-0.052) (-1.300,-0.046)}
\EQ{(1.300,-0.023) (1.350,-0.026) (1.400,-0.029) (1.450,-0.032) (1.500,-0.035) (1.550,-0.039) (1.600,-0.043) (1.650,-0.047) (1.700,-0.052) (1.750,-0.056) (1.800,-0.061) (1.850,-0.067) (1.900,-0.072) (1.950,-0.078) (2.000,-0.084) (2.050,-0.091) (2.100,-0.098) (2.150,-0.105) (2.200,-0.112) (2.250,-0.120) (2.300,-0.128) (2.350,-0.137) (2.400,-0.146) (2.450,-0.155) (2.500,-0.165) (2.550,-0.175) (2.600,-0.186) (2.650,-0.197) (2.700,-0.209) (2.750,-0.220) (2.800,-0.233) (2.850,-0.246) (2.900,-0.259) (2.950,-0.273) (3.000,-0.287) (3.050,-0.302) (3.100,-0.318) (3.150,-0.334) (3.199,-0.350) (3.250,-0.367) (3.300,-0.385) (3.350,-0.403) (3.400,-0.422) (3.450,-0.442) (3.500,-0.462) (3.550,-0.483) (3.600,-0.504) (3.646,-0.525) (3.700,-0.550) (3.750,-0.573) (3.800,-0.598) (3.850,-0.623) (3.900,-0.649) (3.949,-0.675) (3.994,-0.700) (4.038,-0.725) (4.100,-0.761) (4.150,-0.792) (4.200,-0.823) (4.250,-0.856) (4.300,-0.889) (4.350,-0.923) (4.400,-0.959) (4.450,-0.996) (4.500,-1.034) (4.550,-1.073) (4.600,-1.113) (4.644,-1.150) (4.700,-1.198) (4.750,-1.243) (4.800,-1.289) (4.838,-1.325) (4.875,-1.362) (4.913,-1.400) (4.950,-1.438) (4.985,-1.475) (5.025,-1.519) (5.074,-1.575) (5.117,-1.625) (5.150,-1.665) (5.198,-1.725) (5.237,-1.775) (5.274,-1.825) (5.310,-1.875) (5.345,-1.925) (5.375,-1.970) (5.411,-2.025) (5.443,-2.075) (5.473,-2.125) (5.500,-2.171) (5.525,-2.214) (5.550,-2.259) (5.575,-2.306) (5.598,-2.350) (5.623,-2.400) (5.648,-2.450) (5.671,-2.500) (5.694,-2.550) (5.716,-2.600) (5.737,-2.650) (5.757,-2.700) (5.777,-2.750) (5.796,-2.800) (5.814,-2.850) (5.831,-2.900) (5.848,-2.950) (5.863,-3.000) (5.879,-3.050) (5.893,-3.100) (5.907,-3.150) (5.920,-3.200) (5.933,-3.250) (5.944,-3.300) (5.956,-3.350) (5.966,-3.400) (5.976,-3.450) (5.985,-3.500) (5.994,-3.550) (6.002,-3.600) (6.010,-3.650) (6.017,-3.700) (6.023,-3.750) (6.029,-3.800) (6.034,-3.850) (6.038,-3.900) (6.042,-3.950) (6.046,-4.000) (6.049,-4.050) (6.051,-4.100) (6.053,-4.150) (6.054,-4.200) (6.054,-4.250) (6.054,-4.300) (6.054,-4.350) (6.053,-4.400) (6.051,-4.450) (6.049,-4.500) (6.047,-4.550) (6.043,-4.600) (6.040,-4.650) (6.035,-4.700) (6.031,-4.750) (6.025,-4.800) (6.019,-4.850) (6.013,-4.900) (6.006,-4.950) (5.998,-5.000) (5.990,-5.050) (5.982,-5.100) (5.973,-5.150) (5.963,-5.200) (5.953,-5.250) (5.942,-5.300) (5.931,-5.350) (5.919,-5.400) (5.906,-5.450) (5.893,-5.500) (5.880,-5.550) (5.866,-5.600) (5.851,-5.650) (5.836,-5.700) (5.820,-5.750) (5.804,-5.800) (5.787,-5.850) (5.769,-5.900) (5.751,-5.950) (5.732,-6.000) (5.713,-6.050) (5.693,-6.100) (5.673,-6.150) (5.652,-6.200) (5.630,-6.250) (5.608,-6.300) (5.585,-6.350) (5.561,-6.400) (5.537,-6.450) (5.512,-6.500) (5.487,-6.550) (5.461,-6.600) (5.434,-6.650) (5.406,-6.700) (5.378,-6.750) (5.350,-6.799) (5.320,-6.850) (5.289,-6.900) (5.258,-6.950) (5.226,-7.000) (5.194,-7.050) (5.160,-7.100) (5.126,-7.150) (5.091,-7.200) (5.056,-7.250) (5.025,-7.292) (4.981,-7.350) (4.950,-7.391) (4.904,-7.450) (4.864,-7.500) (4.825,-7.547) (4.781,-7.600) (4.738,-7.650) (4.700,-7.693) (4.650,-7.749) (4.603,-7.800) (4.555,-7.850) (4.507,-7.900) (4.457,-7.950) (4.407,-8.000) (4.355,-8.050) (4.302,-8.100) (4.250,-8.147) (4.200,-8.192) (4.162,-8.225) (4.104,-8.275) (4.050,-8.320) (4.000,-8.361) (3.952,-8.400) (3.900,-8.441) (3.855,-8.475) (3.800,-8.517) (3.755,-8.550) (3.700,-8.589) (3.650,-8.625) (3.600,-8.659) (3.550,-8.693) (3.501,-8.725) (3.450,-8.758) (3.400,-8.790) (3.350,-8.821) (3.302,-8.850) (3.250,-8.881) (3.200,-8.910) (3.150,-8.939) (3.100,-8.967) (3.050,-8.994) (3.000,-9.021) (2.950,-9.047) (2.900,-9.072) (2.850,-9.097) (2.800,-9.122) (2.750,-9.146) (2.700,-9.169) (2.650,-9.192) (2.600,-9.214) (2.550,-9.236) (2.500,-9.257) (2.450,-9.278) (2.400,-9.298) (2.350,-9.318) (2.300,-9.337) (2.250,-9.356) (2.200,-9.374) (2.150,-9.392) (2.100,-9.410) (2.050,-9.427) (2.000,-9.443) (1.950,-9.459) (1.900,-9.475) (1.850,-9.490) (1.800,-9.504) (1.750,-9.519) (1.700,-9.533) (1.650,-9.546) (1.600,-9.559) (1.550,-9.571) (1.500,-9.583) (1.450,-9.595) (1.400,-9.606) (1.350,-9.617) (1.300,-9.628) (1.250,-9.638) (1.200,-9.647) (1.150,-9.656) (1.100,-9.665) (1.050,-9.674) (1.000,-9.682) (0.950,-9.689) (0.900,-9.696) (0.850,-9.703) (0.800,-9.710) (0.750,-9.716) (0.700,-9.721) (0.650,-9.726) (0.600,-9.731) (0.550,-9.736) (0.500,-9.740) (0.450,-9.743) (0.400,-9.747) (0.350,-9.750) (0.300,-9.752) (0.250,-9.754) (0.200,-9.756) (0.150,-9.757) (0.100,-9.758) (0.050,-9.759) (0.000,-9.759) (-0.050,-9.759) (-0.100,-9.758) (-0.150,-9.757) (-0.200,-9.756) (-0.250,-9.754) (-0.300,-9.752) (-0.350,-9.750) (-0.400,-9.747) (-0.450,-9.743) (-0.500,-9.740) (-0.550,-9.736) (-0.600,-9.731) (-0.650,-9.726) (-0.700,-9.721) (-0.750,-9.716) (-0.800,-9.710) (-0.850,-9.703) (-0.900,-9.696) (-0.950,-9.689) (-1.000,-9.682) (-1.050,-9.674) (-1.100,-9.665) (-1.150,-9.656) (-1.200,-9.647) (-1.250,-9.638) (-1.300,-9.628) (-1.350,-9.617) (-1.400,-9.606) (-1.450,-9.595) (-1.500,-9.583) (-1.550,-9.571) (-1.600,-9.559) (-1.650,-9.546) (-1.700,-9.533) (-1.750,-9.519) (-1.800,-9.504) (-1.850,-9.490) (-1.899,-9.475) (-1.950,-9.459) (-2.000,-9.443) (-2.050,-9.427) (-2.100,-9.410) (-2.150,-9.392) (-2.198,-9.375) (-2.250,-9.356) (-2.300,-9.337) (-2.350,-9.318) (-2.400,-9.298) (-2.450,-9.278) (-2.500,-9.257) (-2.550,-9.236) (-2.600,-9.214) (-2.650,-9.192) (-2.700,-9.169) (-2.750,-9.146) (-2.800,-9.122) (-2.850,-9.097) (-2.900,-9.072) (-2.950,-9.047) (-3.000,-9.021) (-3.050,-8.994) (-3.100,-8.967) (-3.150,-8.939) (-3.200,-8.910) (-3.250,-8.881) (-3.300,-8.851) (-3.344,-8.825) (-3.400,-8.790) (-3.450,-8.758) (-3.500,-8.726) (-3.550,-8.693) (-3.600,-8.659) (-3.650,-8.625) (-3.700,-8.589) (-3.750,-8.553) (-3.800,-8.517) (-3.850,-8.479) (-3.900,-8.441) (-3.950,-8.401) (-4.000,-8.361) (-4.045,-8.325) (-4.100,-8.279) (-4.150,-8.236) (-4.191,-8.200) (-4.247,-8.150) (-4.300,-8.101) (-4.350,-8.055) (-4.400,-8.007) (-4.450,-7.957) (-4.500,-7.907) (-4.550,-7.856) (-4.600,-7.803) (-4.649,-7.750) (-4.694,-7.700) (-4.738,-7.650) (-4.775,-7.607) (-4.823,-7.550) (-4.864,-7.500) (-4.900,-7.455) (-4.943,-7.400) (-4.975,-7.359) (-5.019,-7.300) (-5.050,-7.258) (-5.091,-7.200) (-5.125,-7.152) (-5.160,-7.100) (-5.194,-7.050) (-5.225,-7.002) (-5.258,-6.950) (-5.289,-6.900) (-5.320,-6.850) (-5.349,-6.800) (-5.375,-6.755) (-5.400,-6.711) (-5.425,-6.666) (-5.450,-6.620) (-5.474,-6.575) (-5.500,-6.525) (-5.525,-6.475) (-5.549,-6.425) (-5.573,-6.375) (-5.596,-6.325) (-5.619,-6.275) (-5.641,-6.225) (-5.662,-6.175) (-5.683,-6.125) (-5.703,-6.075) (-5.723,-6.025) (-5.742,-5.975) (-5.760,-5.925) (-5.778,-5.875) (-5.795,-5.825) (-5.812,-5.775) (-5.828,-5.725) (-5.843,-5.675) (-5.858,-5.625) (-5.873,-5.575) (-5.887,-5.525) (-5.900,-5.475) (-5.913,-5.425) (-5.925,-5.375) (-5.936,-5.325) (-5.947,-5.275) (-5.958,-5.225) (-5.968,-5.175) (-5.977,-5.125) (-5.986,-5.075) (-5.994,-5.025) (-6.002,-4.975) (-6.009,-4.925) (-6.016,-4.875) (-6.022,-4.825) (-6.028,-4.775) (-6.033,-4.725) (-6.038,-4.675) (-6.042,-4.625) (-6.045,-4.575) (-6.048,-4.525) (-6.050,-4.475) (-6.052,-4.425) (-6.054,-4.375) (-6.054,-4.325) (-6.055,-4.275) (-6.054,-4.225) (-6.053,-4.175) (-6.052,-4.125) (-6.050,-4.075) (-6.047,-4.025) (-6.044,-3.975) (-6.040,-3.925) (-6.036,-3.875) (-6.031,-3.825) (-6.026,-3.775) (-6.020,-3.725) (-6.013,-3.675) (-6.006,-3.625) (-5.998,-3.575) (-5.990,-3.525) (-5.981,-3.475) (-5.971,-3.425) (-5.961,-3.375) (-5.950,-3.325) (-5.939,-3.275) (-5.926,-3.225) (-5.914,-3.175) (-5.900,-3.125) (-5.886,-3.075) (-5.871,-3.025) (-5.856,-2.975) (-5.839,-2.925) (-5.822,-2.875) (-5.805,-2.825) (-5.786,-2.775) (-5.767,-2.725) (-5.747,-2.675) (-5.727,-2.625) (-5.705,-2.575) (-5.683,-2.525) (-5.660,-2.475) (-5.636,-2.425) (-5.611,-2.375) (-5.585,-2.325) (-5.559,-2.275) (-5.531,-2.225) (-5.503,-2.175) (-5.475,-2.128) (-5.443,-2.075) (-5.411,-2.025) (-5.378,-1.975) (-5.350,-1.933) (-5.310,-1.875) (-5.275,-1.827) (-5.237,-1.775) (-5.200,-1.727) (-5.158,-1.675) (-5.125,-1.634) (-5.075,-1.576) (-5.030,-1.525) (-4.985,-1.475) (-4.950,-1.438) (-4.913,-1.400) (-4.875,-1.362) (-4.838,-1.325) (-4.800,-1.289) (-4.758,-1.250) (-4.702,-1.200) (-4.650,-1.155) (-4.600,-1.113) (-4.553,-1.075) (-4.500,-1.034) (-4.456,-1.000) (-4.400,-0.959) (-4.352,-0.925) (-4.300,-0.889) (-4.250,-0.856) (-4.203,-0.825) (-4.150,-0.792) (-4.100,-0.761) (-4.050,-0.732) (-4.000,-0.703) (-3.950,-0.676) (-3.902,-0.650) (-3.854,-0.625) (-3.805,-0.600) (-3.754,-0.575) (-3.701,-0.550) (-3.650,-0.527) (-3.600,-0.504) (-3.550,-0.483) (-3.500,-0.462) (-3.450,-0.442) (-3.400,-0.422) (-3.350,-0.403) (-3.300,-0.385) (-3.250,-0.367) (-3.200,-0.350) (-3.150,-0.334) (-3.100,-0.318) (-3.050,-0.302) (-3.000,-0.287) (-2.950,-0.273) (-2.900,-0.259) (-2.850,-0.246) (-2.800,-0.233) (-2.750,-0.220) (-2.700,-0.209) (-2.650,-0.197) (-2.600,-0.186) (-2.550,-0.175) (-2.500,-0.165) (-2.450,-0.155) (-2.400,-0.146) (-2.350,-0.137) (-2.300,-0.128) (-2.250,-0.120) (-2.200,-0.112) (-2.150,-0.105) (-2.100,-0.098) (-2.050,-0.091) (-2.000,-0.084) (-1.950,-0.078) (-1.900,-0.072) (-1.850,-0.067) (-1.800,-0.061) (-1.750,-0.056) (-1.700,-0.052) (-1.650,-0.047) (-1.600,-0.043) (-1.550,-0.039) (-1.500,-0.035) (-1.450,-0.032) (-1.400,-0.029) (-1.350,-0.026) (-1.302,-0.023)}
\EQ{(1.300,0.023) (1.350,0.026) (1.400,0.029) (1.450,0.032) (1.500,0.035) (1.550,0.039) (1.600,0.043) (1.650,0.047) (1.700,0.052) (1.750,0.056) (1.800,0.061) (1.850,0.067) (1.900,0.072) (1.950,0.078) (2.000,0.084) (2.050,0.091) (2.100,0.098) (2.150,0.105) (2.200,0.112) (2.250,0.120) (2.300,0.128) (2.350,0.137) (2.400,0.146) (2.450,0.155) (2.500,0.165) (2.550,0.175) (2.600,0.186) (2.650,0.197) (2.700,0.209) (2.750,0.220) (2.800,0.233) (2.850,0.246) (2.900,0.259) (2.950,0.273) (3.000,0.287) (3.050,0.302) (3.100,0.318) (3.150,0.334) (3.199,0.350) (3.250,0.367) (3.300,0.385) (3.350,0.403) (3.400,0.422) (3.450,0.442) (3.500,0.462) (3.550,0.483) (3.600,0.504) (3.646,0.525) (3.700,0.550) (3.750,0.573) (3.800,0.598) (3.850,0.623) (3.900,0.649) (3.949,0.675) (3.994,0.700) (4.038,0.725) (4.100,0.761) (4.150,0.792) (4.200,0.823) (4.250,0.856) (4.300,0.889) (4.350,0.923) (4.400,0.959) (4.450,0.996) (4.500,1.034) (4.550,1.073) (4.600,1.113) (4.644,1.150) (4.700,1.198) (4.750,1.243) (4.800,1.289) (4.838,1.325) (4.875,1.362) (4.913,1.400) (4.950,1.438) (4.985,1.475) (5.025,1.519) (5.074,1.575) (5.117,1.625) (5.150,1.665) (5.198,1.725) (5.237,1.775) (5.274,1.825) (5.310,1.875) (5.345,1.925) (5.375,1.970) (5.411,2.025) (5.443,2.075) (5.473,2.125) (5.500,2.171) (5.525,2.214) (5.550,2.259) (5.575,2.306) (5.598,2.350) (5.623,2.400) (5.648,2.450) (5.671,2.500) (5.694,2.550) (5.716,2.600) (5.737,2.650) (5.757,2.700) (5.777,2.750) (5.796,2.800) (5.814,2.850) (5.831,2.900) (5.848,2.950) (5.863,3.000) (5.879,3.050) (5.893,3.100) (5.907,3.150) (5.920,3.200) (5.933,3.250) (5.944,3.300) (5.956,3.350) (5.966,3.400) (5.976,3.450) (5.985,3.500) (5.994,3.550) (6.002,3.600) (6.010,3.650) (6.017,3.700) (6.023,3.750) (6.029,3.800) (6.034,3.850) (6.038,3.900) (6.042,3.950) (6.046,4.000) (6.049,4.050) (6.051,4.100) (6.053,4.150) (6.054,4.200) (6.054,4.250) (6.054,4.300) (6.054,4.350) (6.053,4.400) (6.051,4.450) (6.049,4.500) (6.047,4.550) (6.043,4.600) (6.040,4.650) (6.035,4.700) (6.031,4.750) (6.025,4.800) (6.019,4.850) (6.013,4.900) (6.006,4.950) (5.998,5.000) (5.990,5.050) (5.982,5.100) (5.973,5.150) (5.963,5.200) (5.953,5.250) (5.942,5.300) (5.931,5.350) (5.919,5.400) (5.906,5.450) (5.893,5.500) (5.880,5.550) (5.866,5.600) (5.851,5.650) (5.836,5.700) (5.820,5.750) (5.804,5.800) (5.787,5.850) (5.769,5.900) (5.751,5.950) (5.732,6.000) (5.713,6.050) (5.693,6.100) (5.673,6.150) (5.652,6.200) (5.630,6.250) (5.608,6.300) (5.585,6.350) (5.561,6.400) (5.537,6.450) (5.512,6.500) (5.487,6.550) (5.461,6.600) (5.434,6.650) (5.406,6.700) (5.378,6.750) (5.350,6.799) (5.320,6.850) (5.289,6.900) (5.258,6.950) (5.226,7.000) (5.194,7.050) (5.160,7.100) (5.126,7.150) (5.091,7.200) (5.056,7.250) (5.025,7.292) (4.981,7.350) (4.950,7.391) (4.904,7.450) (4.864,7.500) (4.825,7.547) (4.781,7.600) (4.738,7.650) (4.700,7.693) (4.650,7.749) (4.603,7.800) (4.555,7.850) (4.507,7.900) (4.457,7.950) (4.407,8.000) (4.355,8.050) (4.302,8.100) (4.250,8.147) (4.200,8.192) (4.162,8.225) (4.104,8.275) (4.050,8.320) (4.000,8.361) (3.952,8.400) (3.900,8.441) (3.855,8.475) (3.800,8.517) (3.755,8.550) (3.700,8.589) (3.650,8.625) (3.600,8.659) (3.550,8.693) (3.501,8.725) (3.450,8.758) (3.400,8.790) (3.350,8.821) (3.302,8.850) (3.250,8.881) (3.200,8.910) (3.150,8.939) (3.100,8.967) (3.050,8.994) (3.000,9.021) (2.950,9.047) (2.900,9.072) (2.850,9.097) (2.800,9.122) (2.750,9.146) (2.700,9.169) (2.650,9.192) (2.600,9.214) (2.550,9.236) (2.500,9.257) (2.450,9.278) (2.400,9.298) (2.350,9.318) (2.300,9.337) (2.250,9.356) (2.200,9.374) (2.150,9.392) (2.100,9.410) (2.050,9.427) (2.000,9.443) (1.950,9.459) (1.900,9.475) (1.850,9.490) (1.800,9.504) (1.750,9.519) (1.700,9.533) (1.650,9.546) (1.600,9.559) (1.550,9.571) (1.500,9.583) (1.450,9.595) (1.400,9.606) (1.350,9.617) (1.300,9.628) (1.250,9.638) (1.200,9.647) (1.150,9.656) (1.100,9.665) (1.050,9.674) (1.000,9.682) (0.950,9.689) (0.900,9.696) (0.850,9.703) (0.800,9.710) (0.750,9.716) (0.700,9.721) (0.650,9.726) (0.600,9.731) (0.550,9.736) (0.500,9.740) (0.450,9.743) (0.400,9.747) (0.350,9.750) (0.300,9.752) (0.250,9.754) (0.200,9.756) (0.150,9.757) (0.100,9.758) (0.050,9.759) (0.000,9.759) (-0.050,9.759) (-0.100,9.758) (-0.150,9.757) (-0.200,9.756) (-0.250,9.754) (-0.300,9.752) (-0.350,9.750) (-0.400,9.747) (-0.450,9.743) (-0.500,9.740) (-0.550,9.736) (-0.600,9.731) (-0.650,9.726) (-0.700,9.721) (-0.750,9.716) (-0.800,9.710) (-0.850,9.703) (-0.900,9.696) (-0.950,9.689) (-1.000,9.682) (-1.050,9.674) (-1.100,9.665) (-1.150,9.656) (-1.200,9.647) (-1.250,9.638) (-1.300,9.628) (-1.350,9.617) (-1.400,9.606) (-1.450,9.595) (-1.500,9.583) (-1.550,9.571) (-1.600,9.559) (-1.650,9.546) (-1.700,9.533) (-1.750,9.519) (-1.800,9.504) (-1.850,9.490) (-1.899,9.475) (-1.950,9.459) (-2.000,9.443) (-2.050,9.427) (-2.100,9.410) (-2.150,9.392) (-2.198,9.375) (-2.250,9.356) (-2.300,9.337) (-2.350,9.318) (-2.400,9.298) (-2.450,9.278) (-2.500,9.257) (-2.550,9.236) (-2.600,9.214) (-2.650,9.192) (-2.700,9.169) (-2.750,9.146) (-2.800,9.122) (-2.850,9.097) (-2.900,9.072) (-2.950,9.047) (-3.000,9.021) (-3.050,8.994) (-3.100,8.967) (-3.150,8.939) (-3.200,8.910) (-3.250,8.881) (-3.300,8.851) (-3.344,8.825) (-3.400,8.790) (-3.450,8.758) (-3.500,8.726) (-3.550,8.693) (-3.600,8.659) (-3.650,8.625) (-3.700,8.589) (-3.750,8.553) (-3.800,8.517) (-3.850,8.479) (-3.900,8.441) (-3.950,8.401) (-4.000,8.361) (-4.045,8.325) (-4.100,8.279) (-4.150,8.236) (-4.191,8.200) (-4.247,8.150) (-4.300,8.101) (-4.350,8.055) (-4.400,8.007) (-4.450,7.957) (-4.500,7.907) (-4.550,7.856) (-4.600,7.803) (-4.649,7.750) (-4.694,7.700) (-4.738,7.650) (-4.775,7.607) (-4.823,7.550) (-4.864,7.500) (-4.900,7.455) (-4.943,7.400) (-4.975,7.359) (-5.019,7.300) (-5.050,7.258) (-5.091,7.200) (-5.125,7.152) (-5.160,7.100) (-5.194,7.050) (-5.225,7.002) (-5.258,6.950) (-5.289,6.900) (-5.320,6.850) (-5.349,6.800) (-5.375,6.755) (-5.400,6.711) (-5.425,6.666) (-5.450,6.620) (-5.474,6.575) (-5.500,6.525) (-5.525,6.475) (-5.549,6.425) (-5.573,6.375) (-5.596,6.325) (-5.619,6.275) (-5.641,6.225) (-5.662,6.175) (-5.683,6.125) (-5.703,6.075) (-5.723,6.025) (-5.742,5.975) (-5.760,5.925) (-5.778,5.875) (-5.795,5.825) (-5.812,5.775) (-5.828,5.725) (-5.843,5.675) (-5.858,5.625) (-5.873,5.575) (-5.887,5.525) (-5.900,5.475) (-5.913,5.425) (-5.925,5.375) (-5.936,5.325) (-5.947,5.275) (-5.958,5.225) (-5.968,5.175) (-5.977,5.125) (-5.986,5.075) (-5.994,5.025) (-6.002,4.975) (-6.009,4.925) (-6.016,4.875) (-6.022,4.825) (-6.028,4.775) (-6.033,4.725) (-6.038,4.675) (-6.042,4.625) (-6.045,4.575) (-6.048,4.525) (-6.050,4.475) (-6.052,4.425) (-6.054,4.375) (-6.054,4.325) (-6.055,4.275) (-6.054,4.225) (-6.053,4.175) (-6.052,4.125) (-6.050,4.075) (-6.047,4.025) (-6.044,3.975) (-6.040,3.925) (-6.036,3.875) (-6.031,3.825) (-6.026,3.775) (-6.020,3.725) (-6.013,3.675) (-6.006,3.625) (-5.998,3.575) (-5.990,3.525) (-5.981,3.475) (-5.971,3.425) (-5.961,3.375) (-5.950,3.325) (-5.939,3.275) (-5.926,3.225) (-5.914,3.175) (-5.900,3.125) (-5.886,3.075) (-5.871,3.025) (-5.856,2.975) (-5.839,2.925) (-5.822,2.875) (-5.805,2.825) (-5.786,2.775) (-5.767,2.725) (-5.747,2.675) (-5.727,2.625) (-5.705,2.575) (-5.683,2.525) (-5.660,2.475) (-5.636,2.425) (-5.611,2.375) (-5.585,2.325) (-5.559,2.275) (-5.531,2.225) (-5.503,2.175) (-5.475,2.128) (-5.443,2.075) (-5.411,2.025) (-5.378,1.975) (-5.350,1.933) (-5.310,1.875) (-5.275,1.827) (-5.237,1.775) (-5.200,1.727) (-5.158,1.675) (-5.125,1.634) (-5.075,1.576) (-5.030,1.525) (-4.985,1.475) (-4.950,1.438) (-4.913,1.400) (-4.875,1.362) (-4.838,1.325) (-4.800,1.289) (-4.758,1.250) (-4.702,1.200) (-4.650,1.155) (-4.600,1.113) (-4.553,1.075) (-4.500,1.034) (-4.456,1.000) (-4.400,0.959) (-4.352,0.925) (-4.300,0.889) (-4.250,0.856) (-4.203,0.825) (-4.150,0.792) (-4.100,0.761) (-4.050,0.732) (-4.000,0.703) (-3.950,0.676) (-3.902,0.650) (-3.854,0.625) (-3.805,0.600) (-3.754,0.575) (-3.701,0.550) (-3.650,0.527) (-3.600,0.504) (-3.550,0.483) (-3.500,0.462) (-3.450,0.442) (-3.400,0.422) (-3.350,0.403) (-3.300,0.385) (-3.250,0.367) (-3.200,0.350) (-3.150,0.334) (-3.100,0.318) (-3.050,0.302) (-3.000,0.287) (-2.950,0.273) (-2.900,0.259) (-2.850,0.246) (-2.800,0.233) (-2.750,0.220) (-2.700,0.209) (-2.650,0.197) (-2.600,0.186) (-2.550,0.175) (-2.500,0.165) (-2.450,0.155) (-2.400,0.146) (-2.350,0.137) (-2.300,0.128) (-2.250,0.120) (-2.200,0.112) (-2.150,0.105) (-2.100,0.098) (-2.050,0.091) (-2.000,0.084) (-1.950,0.078) (-1.900,0.072) (-1.850,0.067) (-1.800,0.061) (-1.750,0.056) (-1.700,0.052) (-1.650,0.047) (-1.600,0.043) (-1.550,0.039) (-1.500,0.035) (-1.450,0.032) (-1.400,0.029) (-1.350,0.026) (-1.302,0.023)}
\EQ{(1.300,0.046) (1.350,0.052) (1.400,0.058) (1.450,0.064) (1.500,0.071) (1.550,0.079) (1.600,0.086) (1.650,0.095) (1.700,0.104) (1.750,0.113) (1.800,0.123) (1.850,0.134) (1.900,0.145) (1.950,0.157) (2.000,0.170) (2.050,0.183) (2.100,0.197) (2.150,0.212) (2.200,0.227) (2.250,0.243) (2.300,0.260) (2.350,0.278) (2.400,0.297) (2.450,0.317) (2.500,0.337) (2.550,0.359) (2.600,0.381) (2.650,0.405) (2.700,0.429) (2.750,0.455) (2.800,0.482) (2.850,0.510) (2.900,0.539) (2.950,0.570) (2.998,0.600) (3.050,0.635) (3.100,0.670) (3.141,0.700) (3.200,0.745) (3.250,0.785) (3.298,0.825) (3.350,0.871) (3.400,0.917) (3.450,0.966) (3.500,1.017) (3.550,1.071) (3.598,1.125) (3.640,1.175) (3.675,1.219) (3.718,1.275) (3.750,1.319) (3.789,1.375) (3.822,1.425) (3.850,1.470) (3.883,1.525) (3.912,1.575) (3.939,1.625) (3.965,1.675) (3.990,1.725) (4.014,1.775) (4.036,1.825) (4.057,1.875) (4.077,1.925) (4.096,1.975) (4.114,2.025) (4.131,2.075) (4.147,2.125) (4.162,2.175) (4.176,2.225) (4.189,2.275) (4.201,2.325) (4.213,2.375) (4.223,2.425) (4.232,2.475) (4.241,2.525) (4.249,2.575) (4.256,2.625) (4.262,2.675) (4.267,2.725) (4.271,2.775) (4.275,2.825) (4.278,2.875) (4.280,2.925) (4.281,2.975) (4.281,3.025) (4.281,3.075) (4.280,3.125) (4.278,3.175) (4.275,3.225) (4.272,3.275) (4.268,3.325) (4.263,3.375) (4.257,3.425) (4.250,3.475) (4.243,3.525) (4.235,3.575) (4.226,3.625) (4.217,3.675) (4.206,3.725) (4.195,3.775) (4.184,3.825) (4.171,3.875) (4.157,3.925) (4.143,3.975) (4.128,4.025) (4.112,4.075) (4.096,4.125) (4.078,4.175) (4.060,4.225) (4.041,4.275) (4.021,4.325) (4.000,4.375) (3.979,4.425) (3.956,4.475) (3.933,4.525) (3.908,4.575) (3.883,4.625) (3.857,4.675) (3.830,4.725) (3.802,4.775) (3.775,4.821) (3.743,4.875) (3.711,4.925) (3.679,4.975) (3.650,5.019) (3.611,5.075) (3.576,5.125) (3.539,5.175) (3.501,5.225) (3.462,5.275) (3.425,5.321) (3.380,5.375) (3.336,5.425) (3.300,5.466) (3.250,5.520) (3.200,5.573) (3.150,5.623) (3.100,5.672) (3.050,5.720) (3.000,5.766) (2.961,5.800) (2.904,5.850) (2.850,5.895) (2.800,5.935) (2.750,5.974) (2.700,6.013) (2.650,6.049) (2.600,6.085) (2.550,6.120) (2.505,6.150) (2.450,6.186) (2.400,6.218) (2.350,6.249) (2.306,6.275) (2.250,6.308) (2.200,6.336) (2.150,6.363) (2.100,6.390) (2.050,6.415) (2.000,6.440) (1.950,6.464) (1.900,6.488) (1.850,6.510) (1.800,6.532) (1.750,6.553) (1.700,6.574) (1.650,6.593) (1.600,6.612) (1.550,6.631) (1.500,6.648) (1.450,6.665) (1.400,6.682) (1.350,6.698) (1.300,6.713) (1.250,6.727) (1.200,6.741) (1.150,6.754) (1.100,6.767) (1.050,6.779) (1.000,6.791) (0.950,6.801) (0.900,6.812) (0.850,6.821) (0.800,6.831) (0.750,6.839) (0.700,6.847) (0.650,6.854) (0.600,6.861) (0.550,6.868) (0.500,6.873) (0.450,6.879) (0.400,6.883) (0.350,6.887) (0.300,6.891) (0.250,6.894) (0.200,6.896) (0.150,6.898) (0.100,6.900) (0.050,6.900) (0.000,6.901) (-0.050,6.900) (-0.100,6.900) (-0.150,6.898) (-0.200,6.896) (-0.250,6.894) (-0.300,6.891) (-0.350,6.887) (-0.400,6.883) (-0.450,6.879) (-0.500,6.873) (-0.550,6.868) (-0.600,6.861) (-0.650,6.854) (-0.700,6.847) (-0.750,6.839) (-0.800,6.831) (-0.850,6.821) (-0.900,6.812) (-0.950,6.801) (-1.000,6.791) (-1.050,6.779) (-1.100,6.767) (-1.150,6.754) (-1.200,6.741) (-1.250,6.727) (-1.300,6.713) (-1.350,6.698) (-1.400,6.682) (-1.450,6.665) (-1.500,6.648) (-1.550,6.631) (-1.600,6.612) (-1.650,6.593) (-1.697,6.575) (-1.750,6.553) (-1.800,6.532) (-1.850,6.510) (-1.900,6.488) (-1.950,6.464) (-2.000,6.440) (-2.050,6.415) (-2.100,6.390) (-2.150,6.363) (-2.200,6.336) (-2.250,6.308) (-2.300,6.279) (-2.348,6.250) (-2.400,6.218) (-2.450,6.186) (-2.500,6.154) (-2.543,6.125) (-2.600,6.085) (-2.649,6.050) (-2.700,6.013) (-2.749,5.975) (-2.800,5.935) (-2.844,5.900) (-2.900,5.853) (-2.950,5.810) (-2.990,5.775) (-3.044,5.725) (-3.097,5.675) (-3.148,5.625) (-3.198,5.575) (-3.246,5.525) (-3.292,5.475) (-3.336,5.425) (-3.375,5.380) (-3.421,5.325) (-3.462,5.275) (-3.500,5.227) (-3.539,5.175) (-3.575,5.126) (-3.611,5.075) (-3.646,5.025) (-3.675,4.981) (-3.711,4.925) (-3.743,4.875) (-3.773,4.825) (-3.800,4.778) (-3.825,4.734) (-3.850,4.688) (-3.875,4.641) (-3.900,4.592) (-3.925,4.541) (-3.950,4.488) (-3.975,4.433) (-4.000,4.376) (-4.021,4.325) (-4.041,4.275) (-4.060,4.225) (-4.078,4.175) (-4.096,4.125) (-4.112,4.075) (-4.128,4.025) (-4.143,3.975) (-4.157,3.925) (-4.171,3.875) (-4.184,3.825) (-4.195,3.775) (-4.206,3.725) (-4.217,3.675) (-4.226,3.625) (-4.235,3.575) (-4.243,3.525) (-4.250,3.475) (-4.257,3.425) (-4.263,3.375) (-4.268,3.325) (-4.272,3.275) (-4.275,3.225) (-4.278,3.175) (-4.280,3.125) (-4.281,3.075) (-4.281,3.025) (-4.281,2.975) (-4.280,2.925) (-4.278,2.875) (-4.275,2.825) (-4.271,2.775) (-4.267,2.725) (-4.262,2.675) (-4.256,2.625) (-4.249,2.575) (-4.241,2.525) (-4.232,2.475) (-4.223,2.425) (-4.213,2.375) (-4.201,2.325) (-4.189,2.275) (-4.176,2.225) (-4.162,2.175) (-4.147,2.125) (-4.131,2.075) (-4.114,2.025) (-4.096,1.975) (-4.077,1.925) (-4.057,1.875) (-4.036,1.825) (-4.014,1.775) (-3.990,1.725) (-3.965,1.675) (-3.939,1.625) (-3.912,1.575) (-3.883,1.525) (-3.853,1.475) (-3.825,1.430) (-3.789,1.375) (-3.754,1.325) (-3.718,1.275) (-3.680,1.225) (-3.640,1.175) (-3.600,1.127) (-3.554,1.075) (-3.508,1.025) (-3.459,0.975) (-3.408,0.925) (-3.355,0.875) (-3.300,0.827) (-3.250,0.785) (-3.207,0.750) (-3.150,0.707) (-3.107,0.675) (-3.050,0.635) (-3.000,0.602) (-2.950,0.570) (-2.900,0.539) (-2.850,0.510) (-2.800,0.482) (-2.750,0.455) (-2.700,0.429) (-2.650,0.405) (-2.600,0.381) (-2.550,0.359) (-2.500,0.337) (-2.450,0.317) (-2.400,0.297) (-2.350,0.278) (-2.300,0.260) (-2.250,0.243) (-2.200,0.227) (-2.150,0.212) (-2.100,0.197) (-2.050,0.183) (-2.000,0.170) (-1.950,0.157) (-1.900,0.145) (-1.850,0.134) (-1.800,0.123) (-1.750,0.113) (-1.700,0.104) (-1.650,0.095) (-1.600,0.086) (-1.550,0.079) (-1.500,0.071) (-1.450,0.064) (-1.400,0.058) (-1.350,0.052) (-1.300,0.046)}
\EQ{(1.300,0.070) (1.350,0.078) (1.400,0.087) (1.450,0.097) (1.500,0.107) (1.550,0.118) (1.600,0.130) (1.650,0.143) (1.700,0.157) (1.750,0.171) (1.800,0.187) (1.850,0.203) (1.900,0.220) (1.950,0.239) (2.000,0.258) (2.050,0.279) (2.098,0.300) (2.150,0.324) (2.200,0.348) (2.250,0.374) (2.298,0.400) (2.342,0.425) (2.400,0.460) (2.450,0.492) (2.500,0.525) (2.550,0.561) (2.600,0.598) (2.650,0.638) (2.694,0.675) (2.750,0.724) (2.800,0.772) (2.850,0.822) (2.899,0.875) (2.943,0.925) (2.985,0.975) (3.024,1.025) (3.060,1.075) (3.095,1.125) (3.125,1.171) (3.159,1.225) (3.188,1.275) (3.216,1.325) (3.241,1.375) (3.266,1.425) (3.289,1.475) (3.310,1.525) (3.330,1.575) (3.349,1.625) (3.367,1.675) (3.383,1.725) (3.398,1.775) (3.412,1.825) (3.425,1.875) (3.437,1.925) (3.447,1.975) (3.456,2.025) (3.465,2.075) (3.472,2.125) (3.479,2.175) (3.484,2.225) (3.488,2.275) (3.491,2.325) (3.494,2.375) (3.495,2.425) (3.496,2.475) (3.495,2.525) (3.494,2.575) (3.491,2.625) (3.488,2.675) (3.483,2.725) (3.478,2.775) (3.472,2.825) (3.465,2.875) (3.457,2.925) (3.448,2.975) (3.438,3.025) (3.427,3.075) (3.415,3.125) (3.403,3.175) (3.389,3.225) (3.374,3.275) (3.359,3.325) (3.342,3.375) (3.325,3.423) (3.306,3.475) (3.286,3.525) (3.265,3.575) (3.243,3.625) (3.220,3.675) (3.196,3.725) (3.171,3.775) (3.145,3.825) (3.118,3.875) (3.089,3.925) (3.059,3.975) (3.028,4.025) (3.000,4.068) (2.962,4.125) (2.927,4.175) (2.890,4.225) (2.852,4.275) (2.812,4.325) (2.775,4.370) (2.728,4.425) (2.683,4.475) (2.637,4.525) (2.600,4.563) (2.563,4.600) (2.525,4.637) (2.485,4.675) (2.430,4.725) (2.375,4.773) (2.325,4.815) (2.281,4.850) (2.225,4.894) (2.183,4.925) (2.125,4.967) (2.078,5.000) (2.025,5.035) (1.975,5.067) (1.925,5.098) (1.881,5.125) (1.825,5.157) (1.775,5.185) (1.725,5.212) (1.675,5.238) (1.625,5.262) (1.575,5.286) (1.525,5.309) (1.475,5.331) (1.425,5.352) (1.375,5.373) (1.325,5.392) (1.275,5.411) (1.225,5.428) (1.175,5.445) (1.125,5.462) (1.075,5.477) (1.025,5.492) (0.975,5.505) (0.925,5.519) (0.875,5.531) (0.825,5.543) (0.775,5.553) (0.725,5.564) (0.675,5.573) (0.625,5.582) (0.575,5.590) (0.525,5.597) (0.475,5.604) (0.425,5.610) (0.375,5.616) (0.325,5.620) (0.275,5.624) (0.225,5.628) (0.175,5.630) (0.125,5.632) (0.075,5.634) (0.025,5.634) (-0.050,5.634) (-0.100,5.633) (-0.150,5.631) (-0.200,5.629) (-0.250,5.626) (-0.300,5.622) (-0.350,5.618) (-0.400,5.613) (-0.450,5.607) (-0.500,5.601) (-0.550,5.594) (-0.600,5.586) (-0.650,5.578) (-0.700,5.569) (-0.750,5.559) (-0.800,5.548) (-0.850,5.537) (-0.899,5.525) (-0.950,5.512) (-1.000,5.499) (-1.050,5.484) (-1.100,5.469) (-1.150,5.453) (-1.200,5.437) (-1.250,5.420) (-1.300,5.401) (-1.350,5.382) (-1.400,5.363) (-1.450,5.342) (-1.500,5.320) (-1.550,5.298) (-1.599,5.275) (-1.650,5.250) (-1.700,5.225) (-1.747,5.200) (-1.793,5.175) (-1.838,5.150) (-1.900,5.114) (-1.950,5.083) (-2.000,5.051) (-2.050,5.019) (-2.100,4.985) (-2.149,4.950) (-2.200,4.913) (-2.249,4.875) (-2.300,4.835) (-2.343,4.800) (-2.400,4.751) (-2.450,4.707) (-2.500,4.661) (-2.538,4.625) (-2.575,4.588) (-2.613,4.550) (-2.650,4.511) (-2.700,4.457) (-2.750,4.400) (-2.792,4.350) (-2.825,4.309) (-2.871,4.250) (-2.908,4.200) (-2.944,4.150) (-2.975,4.106) (-3.012,4.050) (-3.044,4.000) (-3.074,3.950) (-3.100,3.906) (-3.125,3.862) (-3.150,3.816) (-3.175,3.768) (-3.200,3.718) (-3.225,3.665) (-3.250,3.610) (-3.275,3.551) (-3.296,3.500) (-3.315,3.450) (-3.333,3.400) (-3.350,3.351) (-3.366,3.300) (-3.382,3.250) (-3.396,3.200) (-3.409,3.150) (-3.421,3.100) (-3.433,3.050) (-3.443,3.000) (-3.453,2.950) (-3.461,2.900) (-3.469,2.850) (-3.475,2.800) (-3.481,2.750) (-3.486,2.700) (-3.490,2.650) (-3.492,2.600) (-3.494,2.550) (-3.495,2.500) (-3.496,2.450) (-3.495,2.400) (-3.493,2.350) (-3.490,2.300) (-3.486,2.250) (-3.481,2.200) (-3.475,2.150) (-3.469,2.100) (-3.461,2.050) (-3.452,2.000) (-3.442,1.950) (-3.431,1.900) (-3.419,1.850) (-3.405,1.800) (-3.391,1.750) (-3.375,1.700) (-3.358,1.650) (-3.340,1.600) (-3.321,1.550) (-3.300,1.501) (-3.277,1.450) (-3.254,1.400) (-3.229,1.350) (-3.202,1.300) (-3.175,1.252) (-3.144,1.200) (-3.112,1.150) (-3.078,1.100) (-3.042,1.050) (-3.004,1.000) (-2.964,0.950) (-2.925,0.904) (-2.876,0.850) (-2.828,0.800) (-2.777,0.750) (-2.725,0.702) (-2.675,0.659) (-2.634,0.625) (-2.575,0.579) (-2.525,0.543) (-2.475,0.508) (-2.425,0.475) (-2.375,0.444) (-2.325,0.415) (-2.275,0.387) (-2.225,0.361) (-2.175,0.336) (-2.125,0.312) (-2.075,0.290) (-2.025,0.269) (-1.975,0.248) (-1.925,0.230) (-1.875,0.212) (-1.825,0.195) (-1.775,0.179) (-1.725,0.164) (-1.676,0.150) (-1.625,0.137) (-1.575,0.124) (-1.525,0.113) (-1.475,0.102) (-1.425,0.092) (-1.375,0.082) (-1.325,0.074) (-1.300,0.070)}
\EQ{(1.300,0.093) (1.350,0.104) (1.400,0.116) (1.450,0.130) (1.500,0.144) (1.550,0.159) (1.600,0.175) (1.650,0.193) (1.700,0.211) (1.750,0.231) (1.800,0.252) (1.850,0.275) (1.900,0.299) (1.950,0.324) (1.997,0.350) (2.041,0.375) (2.100,0.412) (2.150,0.444) (2.194,0.475) (2.250,0.517) (2.292,0.550) (2.350,0.599) (2.400,0.645) (2.450,0.693) (2.500,0.746) (2.548,0.800) (2.589,0.850) (2.625,0.896) (2.664,0.950) (2.697,1.000) (2.725,1.044) (2.759,1.100) (2.786,1.150) (2.812,1.200) (2.836,1.250) (2.858,1.300) (2.879,1.350) (2.898,1.400) (2.916,1.450) (2.932,1.500) (2.947,1.550) (2.960,1.600) (2.972,1.650) (2.983,1.700) (2.993,1.750) (3.001,1.800) (3.008,1.850) (3.014,1.900) (3.019,1.950) (3.023,2.000) (3.025,2.050) (3.027,2.100) (3.027,2.150) (3.026,2.200) (3.025,2.250) (3.022,2.300) (3.018,2.350) (3.013,2.400) (3.006,2.450) (2.999,2.500) (2.991,2.550) (2.981,2.600) (2.971,2.650) (2.959,2.700) (2.947,2.750) (2.933,2.800) (2.918,2.850) (2.902,2.900) (2.885,2.950) (2.866,3.000) (2.847,3.050) (2.826,3.100) (2.804,3.150) (2.781,3.200) (2.756,3.250) (2.730,3.300) (2.703,3.350) (2.675,3.399) (2.645,3.450) (2.613,3.500) (2.580,3.550) (2.550,3.594) (2.509,3.650) (2.475,3.696) (2.432,3.750) (2.400,3.789) (2.350,3.846) (2.301,3.900) (2.253,3.950) (2.203,4.000) (2.151,4.050) (2.100,4.096) (2.050,4.139) (2.007,4.175) (1.950,4.220) (1.900,4.258) (1.850,4.295) (1.807,4.325) (1.750,4.363) (1.700,4.395) (1.651,4.425) (1.600,4.455) (1.550,4.483) (1.500,4.510) (1.450,4.536) (1.400,4.561) (1.350,4.585) (1.300,4.607) (1.250,4.629) (1.200,4.649) (1.150,4.669) (1.100,4.687) (1.050,4.705) (1.000,4.722) (0.950,4.737) (0.900,4.752) (0.850,4.766) (0.800,4.779) (0.750,4.792) (0.700,4.803) (0.650,4.814) (0.600,4.824) (0.550,4.833) (0.500,4.841) (0.450,4.848) (0.400,4.855) (0.350,4.861) (0.300,4.866) (0.250,4.870) (0.200,4.873) (0.150,4.876) (0.100,4.878) (0.050,4.879) (0.000,4.880) (-0.050,4.879) (-0.100,4.878) (-0.150,4.876) (-0.200,4.873) (-0.250,4.870) (-0.300,4.866) (-0.350,4.861) (-0.400,4.855) (-0.450,4.848) (-0.500,4.841) (-0.550,4.833) (-0.600,4.824) (-0.650,4.814) (-0.700,4.803) (-0.750,4.792) (-0.800,4.779) (-0.850,4.766) (-0.900,4.752) (-0.950,4.737) (-1.000,4.722) (-1.050,4.705) (-1.100,4.687) (-1.150,4.669) (-1.198,4.650) (-1.250,4.629) (-1.300,4.607) (-1.350,4.585) (-1.400,4.561) (-1.450,4.536) (-1.500,4.510) (-1.550,4.483) (-1.600,4.455) (-1.650,4.426) (-1.700,4.395) (-1.750,4.363) (-1.800,4.330) (-1.843,4.300) (-1.900,4.258) (-1.944,4.225) (-2.000,4.181) (-2.050,4.139) (-2.095,4.100) (-2.150,4.051) (-2.200,4.003) (-2.250,3.954) (-2.300,3.901) (-2.347,3.850) (-2.390,3.800) (-2.425,3.758) (-2.472,3.700) (-2.509,3.650) (-2.546,3.600) (-2.575,3.558) (-2.613,3.500) (-2.645,3.450) (-2.675,3.400) (-2.700,3.356) (-2.725,3.310) (-2.750,3.262) (-2.775,3.212) (-2.800,3.159) (-2.825,3.102) (-2.847,3.050) (-2.866,3.000) (-2.885,2.950) (-2.902,2.900) (-2.918,2.850) (-2.933,2.800) (-2.947,2.750) (-2.959,2.700) (-2.971,2.650) (-2.981,2.600) (-2.991,2.550) (-2.999,2.500) (-3.006,2.450) (-3.013,2.400) (-3.018,2.350) (-3.022,2.300) (-3.025,2.250) (-3.026,2.200) (-3.027,2.150) (-3.027,2.100) (-3.025,2.050) (-3.023,2.000) (-3.019,1.950) (-3.014,1.900) (-3.008,1.850) (-3.001,1.800) (-2.993,1.750) (-2.983,1.700) (-2.972,1.650) (-2.960,1.600) (-2.947,1.550) (-2.932,1.500) (-2.916,1.450) (-2.898,1.400) (-2.879,1.350) (-2.858,1.300) (-2.836,1.250) (-2.812,1.200) (-2.786,1.150) (-2.759,1.100) (-2.729,1.050) (-2.700,1.004) (-2.664,0.950) (-2.628,0.900) (-2.589,0.850) (-2.550,0.802) (-2.504,0.750) (-2.457,0.700) (-2.406,0.650) (-2.351,0.600) (-2.300,0.557) (-2.261,0.525) (-2.200,0.479) (-2.158,0.450) (-2.100,0.412) (-2.050,0.381) (-2.000,0.352) (-1.951,0.325) (-1.902,0.300) (-1.850,0.275) (-1.800,0.252) (-1.750,0.231) (-1.700,0.211) (-1.650,0.193) (-1.600,0.175) (-1.550,0.159) (-1.500,0.144) (-1.450,0.130) (-1.400,0.116) (-1.350,0.104) (-1.300,0.093)}
\EQ{(1.300,0.117) (1.350,0.131) (1.400,0.146) (1.450,0.163) (1.500,0.181) (1.549,0.200) (1.600,0.221) (1.650,0.244) (1.700,0.268) (1.750,0.293) (1.800,0.321) (1.849,0.350) (1.900,0.382) (1.950,0.416) (1.997,0.450) (2.050,0.492) (2.090,0.525) (2.145,0.575) (2.196,0.625) (2.243,0.675) (2.287,0.725) (2.325,0.772) (2.365,0.825) (2.400,0.875) (2.432,0.925) (2.462,0.975) (2.490,1.025) (2.516,1.075) (2.540,1.125) (2.562,1.175) (2.582,1.225) (2.600,1.275) (2.617,1.325) (2.632,1.375) (2.646,1.425) (2.658,1.475) (2.669,1.525) (2.679,1.575) (2.687,1.625) (2.693,1.675) (2.699,1.725) (2.703,1.775) (2.706,1.825) (2.707,1.875) (2.708,1.925) (2.707,1.975) (2.705,2.025) (2.701,2.075) (2.697,2.125) (2.691,2.175) (2.684,2.225) (2.676,2.275) (2.667,2.325) (2.656,2.375) (2.644,2.425) (2.631,2.475) (2.617,2.525) (2.602,2.575) (2.585,2.625) (2.567,2.675) (2.547,2.725) (2.527,2.775) (2.504,2.825) (2.481,2.875) (2.456,2.925) (2.430,2.975) (2.402,3.025) (2.375,3.070) (2.341,3.125) (2.308,3.175) (2.275,3.223) (2.237,3.275) (2.200,3.323) (2.158,3.375) (2.125,3.414) (2.075,3.470) (2.025,3.523) (1.975,3.574) (1.925,3.621) (1.875,3.667) (1.836,3.700) (1.776,3.750) (1.725,3.790) (1.677,3.825) (1.625,3.862) (1.575,3.896) (1.530,3.925) (1.475,3.959) (1.425,3.988) (1.375,4.016) (1.325,4.043) (1.275,4.068) (1.225,4.092) (1.175,4.115) (1.125,4.137) (1.075,4.158) (1.025,4.177) (0.975,4.196) (0.925,4.213) (0.875,4.229) (0.825,4.245) (0.775,4.259) (0.725,4.272) (0.675,4.285) (0.625,4.296) (0.575,4.307) (0.525,4.317) (0.475,4.325) (0.425,4.333) (0.375,4.340) (0.325,4.346) (0.275,4.351) (0.225,4.356) (0.175,4.359) (0.125,4.362) (0.075,4.363) (0.025,4.364) (-0.050,4.364) (-0.100,4.363) (-0.150,4.361) (-0.200,4.358) (-0.250,4.354) (-0.300,4.349) (-0.350,4.343) (-0.400,4.337) (-0.450,4.329) (-0.500,4.321) (-0.550,4.312) (-0.600,4.302) (-0.650,4.291) (-0.700,4.279) (-0.750,4.266) (-0.800,4.252) (-0.850,4.237) (-0.900,4.221) (-0.950,4.204) (-1.000,4.187) (-1.050,4.168) (-1.100,4.147) (-1.150,4.126) (-1.200,4.104) (-1.250,4.080) (-1.300,4.056) (-1.350,4.030) (-1.400,4.002) (-1.448,3.975) (-1.500,3.944) (-1.550,3.912) (-1.600,3.879) (-1.642,3.850) (-1.700,3.808) (-1.744,3.775) (-1.800,3.730) (-1.850,3.688) (-1.894,3.650) (-1.948,3.600) (-1.999,3.550) (-2.047,3.500) (-2.093,3.450) (-2.137,3.400) (-2.175,3.354) (-2.218,3.300) (-2.250,3.257) (-2.291,3.200) (-2.325,3.150) (-2.357,3.100) (-2.387,3.050) (-2.416,3.000) (-2.443,2.950) (-2.469,2.900) (-2.493,2.850) (-2.516,2.800) (-2.537,2.750) (-2.557,2.700) (-2.575,2.653) (-2.593,2.600) (-2.610,2.550) (-2.624,2.500) (-2.638,2.450) (-2.650,2.400) (-2.662,2.350) (-2.672,2.300) (-2.680,2.250) (-2.688,2.200) (-2.694,2.150) (-2.699,2.100) (-2.703,2.050) (-2.706,2.000) (-2.707,1.950) (-2.708,1.900) (-2.707,1.850) (-2.704,1.800) (-2.701,1.750) (-2.696,1.700) (-2.690,1.650) (-2.683,1.600) (-2.674,1.550) (-2.664,1.500) (-2.652,1.450) (-2.639,1.400) (-2.625,1.351) (-2.609,1.300) (-2.591,1.250) (-2.572,1.200) (-2.551,1.150) (-2.528,1.100) (-2.503,1.050) (-2.476,1.000) (-2.450,0.954) (-2.416,0.900) (-2.383,0.850) (-2.350,0.805) (-2.308,0.750) (-2.275,0.711) (-2.225,0.655) (-2.175,0.604) (-2.125,0.557) (-2.075,0.513) (-2.029,0.475) (-1.975,0.434) (-1.927,0.400) (-1.875,0.366) (-1.825,0.335) (-1.775,0.307) (-1.725,0.280) (-1.675,0.255) (-1.625,0.232) (-1.575,0.211) (-1.525,0.191) (-1.475,0.172) (-1.425,0.155) (-1.375,0.139) (-1.325,0.124) (-1.300,0.117)}
\EQ{(1.300,0.141) (1.350,0.158) (1.400,0.177) (1.450,0.197) (1.500,0.220) (1.550,0.243) (1.600,0.269) (1.650,0.297) (1.697,0.325) (1.750,0.359) (1.800,0.394) (1.841,0.425) (1.900,0.473) (1.950,0.517) (2.000,0.566) (2.050,0.619) (2.099,0.675) (2.138,0.725) (2.174,0.775) (2.208,0.825) (2.239,0.875) (2.267,0.925) (2.293,0.975) (2.317,1.025) (2.339,1.075) (2.359,1.125) (2.377,1.175) (2.394,1.225) (2.409,1.275) (2.422,1.325) (2.433,1.375) (2.443,1.425) (2.451,1.475) (2.458,1.525) (2.464,1.575) (2.468,1.625) (2.470,1.675) (2.472,1.725) (2.472,1.775) (2.470,1.825) (2.467,1.875) (2.463,1.925) (2.458,1.975) (2.451,2.025) (2.443,2.075) (2.434,2.125) (2.423,2.175) (2.411,2.225) (2.398,2.275) (2.383,2.325) (2.367,2.375) (2.350,2.423) (2.330,2.475) (2.310,2.525) (2.288,2.575) (2.265,2.625) (2.240,2.675) (2.213,2.725) (2.184,2.775) (2.154,2.825) (2.125,2.871) (2.089,2.925) (2.053,2.975) (2.015,3.025) (1.975,3.075) (1.933,3.125) (1.888,3.175) (1.850,3.215) (1.800,3.265) (1.763,3.300) (1.725,3.335) (1.679,3.375) (1.625,3.420) (1.575,3.459) (1.525,3.496) (1.475,3.531) (1.425,3.565) (1.375,3.597) (1.328,3.625) (1.285,3.650) (1.239,3.675) (1.191,3.700) (1.141,3.725) (1.087,3.750) (1.031,3.775) (0.975,3.798) (0.925,3.817) (0.875,3.835) (0.825,3.852) (0.775,3.868) (0.725,3.883) (0.675,3.897) (0.625,3.909) (0.575,3.921) (0.525,3.932) (0.475,3.941) (0.425,3.950) (0.375,3.958) (0.325,3.964) (0.275,3.970) (0.225,3.975) (0.175,3.978) (0.125,3.981) (0.075,3.983) (0.025,3.984) (-0.050,3.984) (-0.100,3.982) (-0.150,3.980) (-0.200,3.977) (-0.250,3.972) (-0.300,3.967) (-0.350,3.961) (-0.400,3.954) (-0.450,3.946) (-0.500,3.937) (-0.550,3.926) (-0.600,3.915) (-0.650,3.903) (-0.700,3.890) (-0.750,3.876) (-0.800,3.860) (-0.850,3.844) (-0.900,3.826) (-0.950,3.808) (-1.000,3.788) (-1.050,3.767) (-1.100,3.744) (-1.150,3.721) (-1.200,3.696) (-1.250,3.669) (-1.300,3.641) (-1.350,3.612) (-1.400,3.581) (-1.448,3.550) (-1.500,3.514) (-1.550,3.478) (-1.600,3.440) (-1.650,3.400) (-1.700,3.357) (-1.750,3.312) (-1.790,3.275) (-1.840,3.225) (-1.875,3.189) (-1.910,3.150) (-1.950,3.105) (-1.995,3.050) (-2.034,3.000) (-2.071,2.950) (-2.100,2.908) (-2.139,2.850) (-2.170,2.800) (-2.199,2.750) (-2.225,2.703) (-2.250,2.654) (-2.275,2.603) (-2.299,2.550) (-2.320,2.500) (-2.340,2.450) (-2.358,2.400) (-2.375,2.351) (-2.391,2.300) (-2.405,2.250) (-2.417,2.200) (-2.429,2.150) (-2.439,2.100) (-2.447,2.050) (-2.455,2.000) (-2.461,1.950) (-2.466,1.900) (-2.469,1.850) (-2.471,1.800) (-2.472,1.750) (-2.471,1.700) (-2.469,1.650) (-2.466,1.600) (-2.461,1.550) (-2.455,1.500) (-2.447,1.450) (-2.438,1.400) (-2.427,1.350) (-2.415,1.300) (-2.401,1.250) (-2.386,1.200) (-2.369,1.150) (-2.350,1.101) (-2.329,1.050) (-2.306,1.000) (-2.281,0.950) (-2.253,0.900) (-2.225,0.852) (-2.192,0.800) (-2.157,0.750) (-2.125,0.708) (-2.078,0.650) (-2.033,0.600) (-1.984,0.550) (-1.931,0.500) (-1.875,0.452) (-1.825,0.413) (-1.775,0.376) (-1.725,0.343) (-1.675,0.312) (-1.625,0.283) (-1.575,0.256) (-1.525,0.231) (-1.475,0.208) (-1.425,0.187) (-1.375,0.167) (-1.327,0.150) (-1.300,0.141)}
\EQ{(1.300,0.165) (1.350,0.186) (1.400,0.208) (1.450,0.233) (1.500,0.259) (1.550,0.288) (1.600,0.319) (1.646,0.350) (1.700,0.390) (1.744,0.425) (1.800,0.474) (1.850,0.522) (1.900,0.575) (1.943,0.625) (1.975,0.665) (2.018,0.725) (2.050,0.773) (2.082,0.825) (2.109,0.875) (2.134,0.925) (2.157,0.975) (2.178,1.025) (2.197,1.075) (2.213,1.125) (2.228,1.175) (2.242,1.225) (2.253,1.275) (2.263,1.325) (2.271,1.375) (2.277,1.425) (2.282,1.475) (2.286,1.525) (2.288,1.575) (2.288,1.625) (2.288,1.675) (2.285,1.725) (2.281,1.775) (2.276,1.825) (2.269,1.875) (2.261,1.925) (2.252,1.975) (2.241,2.025) (2.228,2.075) (2.215,2.125) (2.199,2.175) (2.182,2.225) (2.164,2.275) (2.144,2.325) (2.122,2.375) (2.100,2.423) (2.075,2.473) (2.050,2.520) (2.025,2.565) (1.989,2.625) (1.956,2.675) (1.925,2.720) (1.885,2.775) (1.850,2.820) (1.805,2.875) (1.761,2.925) (1.725,2.964) (1.690,3.000) (1.650,3.039) (1.611,3.075) (1.554,3.125) (1.500,3.170) (1.450,3.209) (1.400,3.245) (1.358,3.275) (1.300,3.313) (1.250,3.344) (1.200,3.373) (1.152,3.400) (1.104,3.425) (1.053,3.450) (1.000,3.475) (0.950,3.497) (0.900,3.517) (0.850,3.536) (0.800,3.554) (0.750,3.571) (0.700,3.587) (0.650,3.601) (0.600,3.614) (0.550,3.626) (0.500,3.637) (0.450,3.647) (0.400,3.656) (0.350,3.664) (0.300,3.670) (0.250,3.676) (0.200,3.680) (0.150,3.684) (0.100,3.687) (0.050,3.688) (0.000,3.689) (-0.050,3.688) (-0.100,3.687) (-0.150,3.684) (-0.200,3.680) (-0.250,3.676) (-0.300,3.670) (-0.350,3.664) (-0.400,3.656) (-0.450,3.647) (-0.500,3.637) (-0.550,3.626) (-0.600,3.614) (-0.650,3.601) (-0.700,3.587) (-0.750,3.571) (-0.800,3.554) (-0.850,3.536) (-0.900,3.517) (-0.950,3.497) (-1.000,3.475) (-1.050,3.452) (-1.100,3.427) (-1.150,3.401) (-1.197,3.375) (-1.250,3.344) (-1.300,3.313) (-1.350,3.280) (-1.394,3.250) (-1.450,3.209) (-1.494,3.175) (-1.550,3.129) (-1.600,3.085) (-1.638,3.050) (-1.675,3.015) (-1.714,2.975) (-1.750,2.937) (-1.800,2.881) (-1.846,2.825) (-1.885,2.775) (-1.922,2.725) (-1.956,2.675) (-1.989,2.625) (-2.019,2.575) (-2.047,2.525) (-2.074,2.475) (-2.099,2.425) (-2.122,2.375) (-2.144,2.325) (-2.164,2.275) (-2.182,2.225) (-2.199,2.175) (-2.215,2.125) (-2.228,2.075) (-2.241,2.025) (-2.252,1.975) (-2.261,1.925) (-2.269,1.875) (-2.276,1.825) (-2.281,1.775) (-2.285,1.725) (-2.288,1.675) (-2.288,1.625) (-2.288,1.575) (-2.286,1.525) (-2.282,1.475) (-2.277,1.425) (-2.271,1.375) (-2.263,1.325) (-2.253,1.275) (-2.242,1.225) (-2.228,1.175) (-2.213,1.125) (-2.197,1.075) (-2.178,1.025) (-2.157,0.975) (-2.134,0.925) (-2.109,0.875) (-2.082,0.825) (-2.051,0.775) (-2.018,0.725) (-1.982,0.675) (-1.950,0.634) (-1.900,0.575) (-1.853,0.525) (-1.801,0.475) (-1.750,0.430) (-1.700,0.390) (-1.650,0.353) (-1.600,0.319) (-1.550,0.288) (-1.500,0.259) (-1.450,0.233) (-1.400,0.208) (-1.350,0.186) (-1.300,0.165)}
\EQ{(1.300,0.191) (1.350,0.215) (1.400,0.241) (1.450,0.270) (1.499,0.300) (1.550,0.335) (1.600,0.372) (1.650,0.414) (1.691,0.450) (1.742,0.500) (1.788,0.550) (1.825,0.594) (1.868,0.650) (1.900,0.696) (1.934,0.750) (1.963,0.800) (1.989,0.850) (2.013,0.900) (2.034,0.950) (2.053,1.000) (2.070,1.050) (2.085,1.100) (2.098,1.150) (2.109,1.200) (2.118,1.250) (2.126,1.300) (2.132,1.350) (2.137,1.400) (2.139,1.450) (2.141,1.500) (2.140,1.550) (2.138,1.600) (2.135,1.650) (2.130,1.700) (2.123,1.750) (2.115,1.800) (2.106,1.850) (2.095,1.900) (2.082,1.950) (2.068,2.000) (2.052,2.050) (2.035,2.100) (2.016,2.150) (1.995,2.200) (1.972,2.250) (1.950,2.296) (1.925,2.344) (1.900,2.389) (1.864,2.450) (1.831,2.500) (1.800,2.546) (1.760,2.600) (1.725,2.645) (1.679,2.700) (1.634,2.750) (1.600,2.786) (1.562,2.825) (1.525,2.860) (1.481,2.900) (1.425,2.947) (1.375,2.987) (1.325,3.025) (1.275,3.060) (1.225,3.093) (1.175,3.124) (1.132,3.150) (1.087,3.175) (1.040,3.200) (0.990,3.225) (0.936,3.250) (0.879,3.275) (0.825,3.297) (0.775,3.315) (0.725,3.333) (0.675,3.349) (0.625,3.364) (0.575,3.377) (0.525,3.390) (0.475,3.401) (0.425,3.411) (0.375,3.420) (0.325,3.427) (0.275,3.434) (0.225,3.439) (0.175,3.444) (0.125,3.447) (0.075,3.449) (0.025,3.450) (-0.037,3.450) (-0.100,3.448) (-0.150,3.445) (-0.200,3.442) (-0.250,3.437) (-0.300,3.431) (-0.350,3.424) (-0.400,3.415) (-0.450,3.406) (-0.500,3.395) (-0.550,3.384) (-0.600,3.371) (-0.650,3.356) (-0.700,3.341) (-0.748,3.325) (-0.800,3.306) (-0.850,3.287) (-0.900,3.266) (-0.950,3.244) (-1.000,3.220) (-1.050,3.195) (-1.100,3.168) (-1.150,3.139) (-1.200,3.109) (-1.250,3.077) (-1.300,3.043) (-1.350,3.006) (-1.391,2.975) (-1.450,2.927) (-1.500,2.883) (-1.550,2.836) (-1.587,2.800) (-1.625,2.760) (-1.675,2.705) (-1.721,2.650) (-1.760,2.600) (-1.797,2.550) (-1.831,2.500) (-1.864,2.450) (-1.894,2.400) (-1.922,2.350) (-1.948,2.300) (-1.972,2.250) (-1.995,2.200) (-2.016,2.150) (-2.035,2.100) (-2.052,2.050) (-2.068,2.000) (-2.082,1.950) (-2.095,1.900) (-2.106,1.850) (-2.115,1.800) (-2.123,1.750) (-2.130,1.700) (-2.135,1.650) (-2.138,1.600) (-2.140,1.550) (-2.141,1.500) (-2.139,1.450) (-2.137,1.400) (-2.132,1.350) (-2.126,1.300) (-2.118,1.250) (-2.109,1.200) (-2.098,1.150) (-2.085,1.100) (-2.070,1.050) (-2.053,1.000) (-2.034,0.950) (-2.013,0.900) (-1.989,0.850) (-1.963,0.800) (-1.934,0.750) (-1.903,0.700) (-1.868,0.650) (-1.830,0.600) (-1.788,0.550) (-1.750,0.508) (-1.700,0.459) (-1.650,0.414) (-1.603,0.375) (-1.550,0.335) (-1.500,0.301) (-1.450,0.270) (-1.400,0.241) (-1.350,0.215) (-1.300,0.191)}
\EQ{(1.300,0.216) (1.350,0.244) (1.400,0.274) (1.450,0.308) (1.500,0.345) (1.550,0.385) (1.595,0.425) (1.646,0.475) (1.692,0.525) (1.725,0.565) (1.770,0.625) (1.800,0.669) (1.834,0.725) (1.862,0.775) (1.887,0.825) (1.909,0.875) (1.929,0.925) (1.947,0.975) (1.962,1.025) (1.975,1.075) (1.987,1.125) (1.997,1.175) (2.004,1.225) (2.011,1.275) (2.015,1.325) (2.017,1.375) (2.018,1.425) (2.017,1.475) (2.015,1.525) (2.011,1.575) (2.005,1.625) (1.998,1.675) (1.989,1.725) (1.979,1.775) (1.967,1.825) (1.953,1.875) (1.937,1.925) (1.920,1.975) (1.901,2.025) (1.880,2.075) (1.858,2.125) (1.833,2.175) (1.807,2.225) (1.778,2.275) (1.750,2.321) (1.714,2.375) (1.679,2.425) (1.641,2.475) (1.601,2.525) (1.557,2.575) (1.510,2.625) (1.475,2.661) (1.434,2.700) (1.378,2.750) (1.325,2.794) (1.285,2.825) (1.225,2.869) (1.180,2.900) (1.125,2.935) (1.075,2.965) (1.025,2.993) (0.975,3.020) (0.925,3.045) (0.875,3.068) (0.825,3.089) (0.775,3.109) (0.725,3.128) (0.675,3.145) (0.625,3.161) (0.576,3.175) (0.525,3.188) (0.476,3.200) (0.425,3.211) (0.375,3.220) (0.325,3.229) (0.275,3.236) (0.225,3.241) (0.175,3.246) (0.125,3.249) (0.075,3.252) (0.025,3.253) (-0.050,3.252) (-0.100,3.251) (-0.150,3.248) (-0.200,3.244) (-0.250,3.239) (-0.300,3.232) (-0.350,3.225) (-0.400,3.216) (-0.450,3.206) (-0.500,3.195) (-0.550,3.182) (-0.600,3.168) (-0.650,3.153) (-0.700,3.137) (-0.750,3.119) (-0.798,3.100) (-0.850,3.079) (-0.900,3.056) (-0.950,3.032) (-1.000,3.007) (-1.050,2.980) (-1.100,2.950) (-1.150,2.919) (-1.200,2.886) (-1.250,2.851) (-1.300,2.814) (-1.348,2.775) (-1.400,2.731) (-1.450,2.685) (-1.486,2.650) (-1.525,2.610) (-1.575,2.555) (-1.621,2.500) (-1.661,2.450) (-1.697,2.400) (-1.731,2.350) (-1.763,2.300) (-1.793,2.250) (-1.820,2.200) (-1.846,2.150) (-1.869,2.100) (-1.891,2.050) (-1.911,2.000) (-1.929,1.950) (-1.945,1.900) (-1.960,1.850) (-1.973,1.800) (-1.984,1.750) (-1.994,1.700) (-2.002,1.650) (-2.008,1.600) (-2.013,1.550) (-2.017,1.500) (-2.018,1.450) (-2.018,1.400) (-2.016,1.350) (-2.013,1.300) (-2.008,1.250) (-2.001,1.200) (-1.992,1.150) (-1.982,1.100) (-1.969,1.050) (-1.955,1.000) (-1.938,0.950) (-1.919,0.900) (-1.898,0.850) (-1.875,0.801) (-1.850,0.753) (-1.825,0.710) (-1.787,0.650) (-1.752,0.600) (-1.713,0.550) (-1.675,0.506) (-1.625,0.454) (-1.575,0.407) (-1.525,0.364) (-1.475,0.326) (-1.425,0.291) (-1.375,0.259) (-1.325,0.230) (-1.300,0.216)}
\EQ{(1.300,0.243) (1.350,0.275) (1.400,0.310) (1.450,0.348) (1.500,0.391) (1.550,0.439) (1.600,0.493) (1.647,0.550) (1.684,0.600) (1.717,0.650) (1.747,0.700) (1.774,0.750) (1.798,0.800) (1.819,0.850) (1.838,0.900) (1.855,0.950) (1.869,1.000) (1.881,1.050) (1.892,1.100) (1.900,1.150) (1.906,1.200) (1.911,1.250) (1.914,1.300) (1.915,1.350) (1.914,1.400) (1.911,1.450) (1.907,1.500) (1.901,1.550) (1.894,1.600) (1.885,1.650) (1.873,1.700) (1.861,1.750) (1.846,1.800) (1.830,1.850) (1.812,1.900) (1.792,1.950) (1.770,2.000) (1.746,2.050) (1.720,2.100) (1.692,2.150) (1.661,2.200) (1.629,2.250) (1.593,2.300) (1.555,2.350) (1.515,2.400) (1.475,2.445) (1.425,2.498) (1.375,2.548) (1.325,2.593) (1.275,2.636) (1.226,2.675) (1.175,2.713) (1.125,2.747) (1.083,2.775) (1.025,2.810) (0.975,2.838) (0.925,2.865) (0.875,2.890) (0.825,2.913) (0.775,2.934) (0.725,2.954) (0.675,2.972) (0.625,2.989) (0.575,3.004) (0.525,3.018) (0.475,3.030) (0.425,3.042) (0.375,3.052) (0.325,3.060) (0.275,3.068) (0.225,3.074) (0.175,3.079) (0.125,3.082) (0.075,3.085) (0.025,3.086) (-0.050,3.086) (-0.100,3.084) (-0.150,3.081) (-0.200,3.076) (-0.250,3.071) (-0.300,3.064) (-0.350,3.056) (-0.400,3.047) (-0.450,3.036) (-0.500,3.024) (-0.550,3.011) (-0.600,2.996) (-0.650,2.980) (-0.700,2.963) (-0.750,2.944) (-0.800,2.923) (-0.850,2.901) (-0.900,2.877) (-0.950,2.852) (-0.999,2.825) (-1.050,2.795) (-1.100,2.764) (-1.150,2.730) (-1.192,2.700) (-1.250,2.656) (-1.300,2.615) (-1.345,2.575) (-1.398,2.525) (-1.448,2.475) (-1.493,2.425) (-1.535,2.375) (-1.575,2.325) (-1.611,2.275) (-1.645,2.225) (-1.675,2.178) (-1.706,2.125) (-1.733,2.075) (-1.758,2.025) (-1.781,1.975) (-1.802,1.925) (-1.821,1.875) (-1.838,1.825) (-1.854,1.775) (-1.867,1.725) (-1.879,1.675) (-1.889,1.625) (-1.898,1.575) (-1.905,1.525) (-1.910,1.475) (-1.913,1.425) (-1.915,1.375) (-1.914,1.325) (-1.913,1.275) (-1.909,1.225) (-1.903,1.175) (-1.896,1.125) (-1.887,1.075) (-1.875,1.025) (-1.862,0.975) (-1.847,0.925) (-1.829,0.875) (-1.809,0.825) (-1.786,0.775) (-1.761,0.725) (-1.733,0.675) (-1.701,0.625) (-1.666,0.575) (-1.627,0.525) (-1.584,0.475) (-1.536,0.425) (-1.482,0.375) (-1.425,0.328) (-1.375,0.292) (-1.325,0.258) (-1.300,0.243)}
\EQ{(1.288,0.262) (1.342,0.300) (1.400,0.347) (1.450,0.391) (1.500,0.442) (1.550,0.499) (1.590,0.550) (1.625,0.600) (1.656,0.650) (1.684,0.700) (1.708,0.750) (1.730,0.800) (1.750,0.850) (1.767,0.900) (1.781,0.950) (1.794,1.000) (1.804,1.050) (1.812,1.100) (1.818,1.150) (1.823,1.200) (1.825,1.250) (1.825,1.300) (1.824,1.350) (1.821,1.400) (1.816,1.450) (1.810,1.500) (1.801,1.550) (1.791,1.600) (1.779,1.650) (1.765,1.700) (1.750,1.748) (1.732,1.800) (1.712,1.850) (1.691,1.900) (1.667,1.950) (1.641,2.000) (1.613,2.050) (1.583,2.100) (1.550,2.149) (1.514,2.200) (1.476,2.250) (1.434,2.300) (1.400,2.339) (1.366,2.375) (1.325,2.416) (1.289,2.450) (1.250,2.484) (1.201,2.525) (1.150,2.565) (1.101,2.600) (1.050,2.634) (1.000,2.666) (0.950,2.695) (0.900,2.722) (0.850,2.748) (0.800,2.771) (0.750,2.793) (0.700,2.813) (0.650,2.831) (0.600,2.848) (0.550,2.864) (0.500,2.878) (0.450,2.890) (0.400,2.901) (0.350,2.911) (0.300,2.919) (0.250,2.926) (0.200,2.932) (0.150,2.937) (0.100,2.940) (0.050,2.942) (0.000,2.943) (-0.050,2.942) (-0.100,2.940) (-0.150,2.937) (-0.200,2.932) (-0.250,2.926) (-0.300,2.919) (-0.350,2.911) (-0.400,2.901) (-0.450,2.890) (-0.500,2.878) (-0.550,2.864) (-0.600,2.848) (-0.650,2.831) (-0.700,2.813) (-0.750,2.793) (-0.800,2.771) (-0.850,2.748) (-0.900,2.722) (-0.950,2.695) (-1.000,2.666) (-1.050,2.634) (-1.100,2.601) (-1.150,2.565) (-1.200,2.526) (-1.250,2.484) (-1.289,2.450) (-1.325,2.416) (-1.366,2.375) (-1.400,2.339) (-1.434,2.300) (-1.475,2.251) (-1.514,2.200) (-1.550,2.150) (-1.583,2.100) (-1.613,2.050) (-1.641,2.000) (-1.667,1.950) (-1.691,1.900) (-1.712,1.850) (-1.732,1.800) (-1.749,1.750) (-1.765,1.700) (-1.779,1.650) (-1.791,1.600) (-1.801,1.550) (-1.810,1.500) (-1.816,1.450) (-1.821,1.400) (-1.824,1.350) (-1.825,1.300) (-1.825,1.250) (-1.823,1.200) (-1.818,1.150) (-1.812,1.100) (-1.804,1.050) (-1.794,1.000) (-1.781,0.950) (-1.767,0.900) (-1.750,0.851) (-1.730,0.800) (-1.708,0.750) (-1.684,0.700) (-1.656,0.650) (-1.625,0.601) (-1.590,0.550) (-1.551,0.500) (-1.508,0.450) (-1.459,0.400) (-1.404,0.350) (-1.350,0.307) (-1.307,0.275) (-1.288,0.262)}
\EQ{(1.275,0.280) (1.325,0.319) (1.375,0.362) (1.414,0.400) (1.450,0.438) (1.500,0.497) (1.539,0.550) (1.572,0.600) (1.600,0.648) (1.625,0.696) (1.650,0.750) (1.670,0.800) (1.688,0.850) (1.703,0.900) (1.716,0.950) (1.726,1.000) (1.734,1.050) (1.741,1.100) (1.745,1.150) (1.747,1.200) (1.748,1.250) (1.746,1.300) (1.743,1.350) (1.738,1.400) (1.731,1.450) (1.722,1.500) (1.711,1.550) (1.698,1.600) (1.683,1.650) (1.667,1.700) (1.648,1.750) (1.627,1.800) (1.604,1.850) (1.579,1.900) (1.552,1.950) (1.525,1.995) (1.489,2.050) (1.454,2.100) (1.416,2.150) (1.375,2.200) (1.330,2.250) (1.282,2.300) (1.229,2.350) (1.175,2.397) (1.125,2.437) (1.075,2.475) (1.025,2.509) (0.975,2.542) (0.925,2.571) (0.875,2.599) (0.825,2.625) (0.775,2.649) (0.725,2.671) (0.675,2.691) (0.625,2.710) (0.575,2.727) (0.525,2.742) (0.475,2.756) (0.425,2.768) (0.375,2.779) (0.325,2.789) (0.275,2.797) (0.225,2.804) (0.175,2.809) (0.125,2.813) (0.075,2.816) (0.025,2.817) (-0.050,2.817) (-0.100,2.815) (-0.150,2.811) (-0.200,2.807) (-0.250,2.800) (-0.300,2.793) (-0.350,2.784) (-0.400,2.774) (-0.450,2.762) (-0.500,2.749) (-0.550,2.735) (-0.600,2.719) (-0.650,2.701) (-0.700,2.681) (-0.750,2.660) (-0.800,2.637) (-0.850,2.612) (-0.900,2.586) (-0.950,2.557) (-1.000,2.526) (-1.050,2.492) (-1.100,2.456) (-1.141,2.425) (-1.200,2.376) (-1.250,2.331) (-1.300,2.282) (-1.350,2.228) (-1.396,2.175) (-1.436,2.125) (-1.472,2.075) (-1.506,2.025) (-1.537,1.975) (-1.566,1.925) (-1.592,1.875) (-1.616,1.825) (-1.638,1.775) (-1.658,1.725) (-1.675,1.675) (-1.691,1.625) (-1.705,1.575) (-1.716,1.525) (-1.726,1.475) (-1.734,1.425) (-1.741,1.375) (-1.745,1.325) (-1.747,1.275) (-1.748,1.225) (-1.746,1.175) (-1.743,1.125) (-1.738,1.075) (-1.730,1.025) (-1.721,0.975) (-1.709,0.925) (-1.695,0.875) (-1.679,0.825) (-1.660,0.775) (-1.639,0.725) (-1.614,0.675) (-1.587,0.625) (-1.556,0.575) (-1.525,0.530) (-1.482,0.475) (-1.438,0.425) (-1.400,0.386) (-1.362,0.350) (-1.301,0.300) (-1.275,0.280)}
\EQ{(1.275,0.308) (1.323,0.350) (1.374,0.400) (1.419,0.450) (1.450,0.489) (1.493,0.550) (1.524,0.600) (1.550,0.647) (1.575,0.698) (1.597,0.750) (1.616,0.800) (1.632,0.850) (1.645,0.900) (1.656,0.950) (1.665,1.000) (1.672,1.050) (1.676,1.100) (1.679,1.150) (1.679,1.200) (1.678,1.250) (1.674,1.300) (1.669,1.350) (1.662,1.400) (1.652,1.450) (1.641,1.500) (1.628,1.550) (1.613,1.600) (1.595,1.650) (1.576,1.700) (1.554,1.750) (1.530,1.800) (1.504,1.850) (1.475,1.900) (1.444,1.950) (1.410,2.000) (1.375,2.047) (1.333,2.100) (1.289,2.150) (1.250,2.191) (1.200,2.240) (1.162,2.275) (1.102,2.325) (1.050,2.365) (1.001,2.400) (0.950,2.434) (0.900,2.464) (0.850,2.492) (0.800,2.518) (0.750,2.543) (0.700,2.565) (0.650,2.585) (0.600,2.604) (0.550,2.621) (0.500,2.636) (0.450,2.650) (0.400,2.662) (0.350,2.672) (0.300,2.682) (0.250,2.689) (0.200,2.696) (0.150,2.700) (0.100,2.704) (0.050,2.706) (0.000,2.707) (-0.050,2.706) (-0.100,2.704) (-0.150,2.700) (-0.200,2.696) (-0.250,2.689) (-0.300,2.682) (-0.350,2.672) (-0.400,2.662) (-0.450,2.650) (-0.500,2.636) (-0.550,2.621) (-0.600,2.604) (-0.650,2.585) (-0.700,2.565) (-0.750,2.543) (-0.800,2.518) (-0.850,2.492) (-0.900,2.464) (-0.950,2.434) (-1.000,2.401) (-1.050,2.365) (-1.100,2.327) (-1.150,2.285) (-1.189,2.250) (-1.241,2.200) (-1.289,2.150) (-1.325,2.109) (-1.373,2.050) (-1.410,2.000) (-1.444,1.950) (-1.475,1.900) (-1.504,1.850) (-1.530,1.800) (-1.554,1.750) (-1.575,1.702) (-1.595,1.650) (-1.613,1.600) (-1.628,1.550) (-1.641,1.500) (-1.652,1.450) (-1.662,1.400) (-1.669,1.350) (-1.674,1.300) (-1.678,1.250) (-1.679,1.200) (-1.679,1.150) (-1.676,1.100) (-1.672,1.050) (-1.665,1.000) (-1.656,0.950) (-1.645,0.900) (-1.632,0.850) (-1.616,0.800) (-1.597,0.750) (-1.576,0.700) (-1.552,0.650) (-1.525,0.601) (-1.493,0.550) (-1.458,0.500) (-1.425,0.458) (-1.375,0.401) (-1.325,0.352) (-1.275,0.308)}
\EQ{(1.275,0.337) (1.314,0.375) (1.350,0.414) (1.400,0.475) (1.436,0.525) (1.467,0.575) (1.495,0.625) (1.519,0.675) (1.540,0.725) (1.559,0.775) (1.574,0.825) (1.587,0.875) (1.598,0.925) (1.606,0.975) (1.612,1.025) (1.616,1.075) (1.618,1.125) (1.618,1.175) (1.616,1.225) (1.611,1.275) (1.605,1.325) (1.597,1.375) (1.586,1.425) (1.574,1.475) (1.559,1.525) (1.543,1.575) (1.524,1.625) (1.503,1.675) (1.479,1.725) (1.454,1.775) (1.425,1.825) (1.394,1.875) (1.360,1.925) (1.325,1.973) (1.283,2.025) (1.250,2.063) (1.216,2.100) (1.175,2.141) (1.125,2.187) (1.081,2.225) (1.025,2.269) (0.983,2.300) (0.925,2.339) (0.875,2.370) (0.825,2.398) (0.775,2.425) (0.725,2.449) (0.675,2.471) (0.625,2.492) (0.575,2.510) (0.525,2.527) (0.475,2.542) (0.425,2.555) (0.375,2.567) (0.325,2.578) (0.275,2.586) (0.225,2.594) (0.175,2.599) (0.125,2.604) (0.075,2.607) (0.025,2.608) (-0.050,2.608) (-0.100,2.605) (-0.150,2.602) (-0.200,2.597) (-0.250,2.590) (-0.300,2.582) (-0.350,2.573) (-0.400,2.562) (-0.450,2.549) (-0.500,2.535) (-0.550,2.519) (-0.600,2.501) (-0.650,2.482) (-0.700,2.460) (-0.750,2.437) (-0.800,2.412) (-0.850,2.384) (-0.900,2.355) (-0.946,2.325) (-1.000,2.288) (-1.050,2.250) (-1.100,2.209) (-1.139,2.175) (-1.191,2.125) (-1.239,2.075) (-1.275,2.034) (-1.323,1.975) (-1.360,1.925) (-1.394,1.875) (-1.425,1.825) (-1.450,1.781) (-1.475,1.734) (-1.500,1.681) (-1.524,1.625) (-1.543,1.575) (-1.559,1.525) (-1.574,1.475) (-1.586,1.425) (-1.597,1.375) (-1.605,1.325) (-1.611,1.275) (-1.616,1.225) (-1.618,1.175) (-1.618,1.125) (-1.616,1.075) (-1.612,1.025) (-1.606,0.975) (-1.598,0.925) (-1.587,0.875) (-1.574,0.825) (-1.559,0.775) (-1.540,0.725) (-1.519,0.675) (-1.495,0.625) (-1.467,0.575) (-1.436,0.525) (-1.400,0.475) (-1.360,0.425) (-1.325,0.387) (-1.289,0.350) (-1.275,0.337)}
\EQ{(1.266,0.359) (1.305,0.400) (1.346,0.450) (1.382,0.500) (1.414,0.550) (1.442,0.600) (1.467,0.650) (1.488,0.700) (1.506,0.750) (1.522,0.800) (1.534,0.850) (1.545,0.900) (1.553,0.950) (1.559,1.000) (1.562,1.050) (1.563,1.100) (1.563,1.150) (1.560,1.200) (1.555,1.250) (1.547,1.300) (1.538,1.350) (1.527,1.400) (1.513,1.450) (1.498,1.500) (1.480,1.550) (1.460,1.600) (1.437,1.650) (1.412,1.700) (1.385,1.750) (1.354,1.800) (1.325,1.844) (1.285,1.900) (1.250,1.944) (1.201,2.000) (1.154,2.050) (1.101,2.100) (1.050,2.145) (1.000,2.185) (0.950,2.222) (0.908,2.250) (0.850,2.287) (0.800,2.315) (0.750,2.342) (0.700,2.366) (0.650,2.388) (0.600,2.409) (0.550,2.427) (0.500,2.443) (0.450,2.458) (0.400,2.471) (0.350,2.483) (0.300,2.493) (0.250,2.501) (0.200,2.508) (0.150,2.513) (0.100,2.517) (0.050,2.519) (0.000,2.520) (-0.050,2.519) (-0.100,2.517) (-0.150,2.513) (-0.200,2.508) (-0.250,2.501) (-0.300,2.493) (-0.350,2.483) (-0.400,2.471) (-0.450,2.458) (-0.500,2.443) (-0.550,2.427) (-0.600,2.409) (-0.650,2.388) (-0.700,2.366) (-0.750,2.342) (-0.800,2.315) (-0.850,2.287) (-0.900,2.256) (-0.945,2.225) (-1.000,2.185) (-1.044,2.150) (-1.100,2.101) (-1.150,2.054) (-1.200,2.001) (-1.245,1.950) (-1.285,1.900) (-1.321,1.850) (-1.350,1.807) (-1.385,1.750) (-1.412,1.700) (-1.437,1.650) (-1.460,1.600) (-1.480,1.550) (-1.498,1.500) (-1.513,1.450) (-1.527,1.400) (-1.538,1.350) (-1.547,1.300) (-1.555,1.250) (-1.560,1.200) (-1.563,1.150) (-1.563,1.100) (-1.562,1.050) (-1.559,1.000) (-1.553,0.950) (-1.545,0.900) (-1.534,0.850) (-1.522,0.800) (-1.506,0.750) (-1.488,0.700) (-1.467,0.650) (-1.442,0.600) (-1.414,0.550) (-1.382,0.500) (-1.350,0.456) (-1.305,0.400) (-1.266,0.359)}
\EQ{(1.251,0.374) (1.295,0.425) (1.325,0.465) (1.365,0.525) (1.393,0.575) (1.418,0.625) (1.440,0.675) (1.458,0.725) (1.473,0.775) (1.486,0.825) (1.496,0.875) (1.504,0.925) (1.510,0.975) (1.513,1.025) (1.514,1.075) (1.512,1.125) (1.509,1.175) (1.503,1.225) (1.495,1.275) (1.486,1.325) (1.473,1.375) (1.459,1.425) (1.442,1.475) (1.423,1.525) (1.402,1.575) (1.378,1.625) (1.352,1.675) (1.325,1.721) (1.290,1.775) (1.255,1.825) (1.216,1.875) (1.175,1.923) (1.127,1.975) (1.075,2.025) (1.025,2.070) (0.975,2.110) (0.925,2.147) (0.875,2.182) (0.825,2.213) (0.775,2.242) (0.725,2.268) (0.675,2.292) (0.625,2.314) (0.575,2.334) (0.525,2.352) (0.475,2.369) (0.425,2.383) (0.375,2.396) (0.325,2.407) (0.275,2.416) (0.225,2.424) (0.175,2.430) (0.125,2.435) (0.075,2.438) (0.025,2.440) (-0.050,2.439) (-0.100,2.437) (-0.150,2.433) (-0.200,2.427) (-0.250,2.420) (-0.300,2.412) (-0.350,2.402) (-0.400,2.390) (-0.450,2.376) (-0.500,2.361) (-0.550,2.344) (-0.599,2.325) (-0.650,2.304) (-0.700,2.280) (-0.750,2.255) (-0.800,2.228) (-0.846,2.200) (-0.900,2.165) (-0.950,2.129) (-1.000,2.090) (-1.048,2.050) (-1.100,2.002) (-1.150,1.951) (-1.195,1.900) (-1.236,1.850) (-1.273,1.800) (-1.307,1.750) (-1.337,1.700) (-1.365,1.650) (-1.390,1.600) (-1.413,1.550) (-1.433,1.500) (-1.450,1.453) (-1.466,1.400) (-1.480,1.350) (-1.491,1.300) (-1.500,1.250) (-1.506,1.200) (-1.511,1.150) (-1.513,1.100) (-1.514,1.050) (-1.512,1.000) (-1.507,0.950) (-1.501,0.900) (-1.492,0.850) (-1.480,0.800) (-1.466,0.750) (-1.449,0.700) (-1.429,0.650) (-1.406,0.600) (-1.379,0.550) (-1.350,0.502) (-1.314,0.450) (-1.275,0.401) (-1.251,0.374)}
\EQ{(1.250,0.405) (1.285,0.450) (1.318,0.500) (1.347,0.550) (1.373,0.600) (1.395,0.650) (1.413,0.700) (1.429,0.750) (1.442,0.800) (1.452,0.850) (1.460,0.900) (1.465,0.950) (1.468,1.000) (1.468,1.050) (1.467,1.100) (1.463,1.150) (1.457,1.200) (1.448,1.250) (1.438,1.300) (1.425,1.350) (1.410,1.400) (1.392,1.450) (1.372,1.500) (1.350,1.550) (1.325,1.600) (1.300,1.645) (1.266,1.700) (1.232,1.750) (1.200,1.793) (1.154,1.850) (1.109,1.900) (1.059,1.950) (1.004,2.000) (0.950,2.044) (0.909,2.075) (0.850,2.116) (0.800,2.147) (0.751,2.175) (0.704,2.200) (0.652,2.225) (0.600,2.248) (0.550,2.268) (0.500,2.285) (0.450,2.301) (0.400,2.315) (0.350,2.328) (0.300,2.338) (0.250,2.347) (0.200,2.354) (0.150,2.360) (0.100,2.364) (0.050,2.366) (0.000,2.367) (-0.050,2.366) (-0.100,2.364) (-0.150,2.360) (-0.200,2.354) (-0.250,2.347) (-0.300,2.338) (-0.350,2.328) (-0.400,2.315) (-0.450,2.301) (-0.500,2.285) (-0.550,2.268) (-0.600,2.248) (-0.650,2.226) (-0.700,2.202) (-0.750,2.176) (-0.795,2.150) (-0.850,2.116) (-0.900,2.081) (-0.942,2.050) (-1.000,2.003) (-1.050,1.958) (-1.100,1.909) (-1.150,1.854) (-1.195,1.800) (-1.225,1.760) (-1.266,1.700) (-1.297,1.650) (-1.325,1.600) (-1.350,1.550) (-1.372,1.500) (-1.392,1.450) (-1.410,1.400) (-1.425,1.350) (-1.438,1.300) (-1.448,1.250) (-1.457,1.200) (-1.463,1.150) (-1.467,1.100) (-1.468,1.050) (-1.468,1.000) (-1.465,0.950) (-1.460,0.900) (-1.452,0.850) (-1.442,0.800) (-1.429,0.750) (-1.413,0.700) (-1.395,0.650) (-1.373,0.600) (-1.350,0.555) (-1.325,0.512) (-1.285,0.450) (-1.250,0.405)}
\EQ{(1.244,0.431) (1.274,0.475) (1.300,0.518) (1.325,0.565) (1.350,0.619) (1.372,0.675) (1.388,0.725) (1.400,0.775) (1.411,0.825) (1.419,0.875) (1.424,0.925) (1.427,0.975) (1.427,1.025) (1.425,1.075) (1.421,1.125) (1.414,1.175) (1.406,1.225) (1.395,1.275) (1.381,1.325) (1.365,1.375) (1.347,1.425) (1.326,1.475) (1.303,1.525) (1.277,1.575) (1.250,1.621) (1.215,1.675) (1.180,1.725) (1.141,1.775) (1.100,1.822) (1.050,1.874) (1.000,1.922) (0.950,1.965) (0.906,2.000) (0.850,2.040) (0.800,2.073) (0.755,2.100) (0.709,2.125) (0.660,2.150) (0.605,2.175) (0.550,2.198) (0.500,2.216) (0.450,2.233) (0.400,2.247) (0.350,2.260) (0.300,2.271) (0.250,2.280) (0.200,2.287) (0.150,2.293) (0.100,2.297) (0.050,2.299) (0.000,2.300) (-0.050,2.299) (-0.100,2.297) (-0.150,2.293) (-0.200,2.287) (-0.250,2.280) (-0.300,2.271) (-0.350,2.260) (-0.400,2.247) (-0.450,2.233) (-0.500,2.216) (-0.550,2.198) (-0.600,2.177) (-0.650,2.155) (-0.700,2.130) (-0.750,2.103) (-0.796,2.075) (-0.850,2.040) (-0.900,2.004) (-0.950,1.965) (-0.997,1.925) (-1.049,1.875) (-1.097,1.825) (-1.141,1.775) (-1.175,1.731) (-1.215,1.675) (-1.248,1.625) (-1.275,1.578) (-1.300,1.531) (-1.325,1.478) (-1.347,1.425) (-1.365,1.375) (-1.381,1.325) (-1.395,1.275) (-1.406,1.225) (-1.414,1.175) (-1.421,1.125) (-1.425,1.075) (-1.427,1.025) (-1.427,0.975) (-1.424,0.925) (-1.419,0.875) (-1.411,0.825) (-1.400,0.775) (-1.388,0.725) (-1.372,0.675) (-1.352,0.625) (-1.330,0.575) (-1.304,0.525) (-1.275,0.477) (-1.244,0.431)}
\EQ{(1.229,0.446) (1.263,0.500) (1.290,0.550) (1.313,0.600) (1.333,0.650) (1.349,0.700) (1.362,0.750) (1.373,0.800) (1.381,0.850) (1.386,0.900) (1.389,0.950) (1.389,1.000) (1.387,1.050) (1.383,1.100) (1.376,1.150) (1.367,1.200) (1.355,1.250) (1.341,1.300) (1.325,1.349) (1.306,1.400) (1.284,1.450) (1.260,1.500) (1.233,1.550) (1.203,1.600) (1.175,1.642) (1.132,1.700) (1.100,1.740) (1.050,1.796) (1.000,1.846) (0.950,1.891) (0.909,1.925) (0.850,1.970) (0.806,2.000) (0.750,2.035) (0.700,2.063) (0.650,2.089) (0.600,2.112) (0.550,2.133) (0.500,2.152) (0.450,2.169) (0.400,2.184) (0.350,2.197) (0.300,2.208) (0.250,2.218) (0.200,2.225) (0.150,2.231) (0.100,2.236) (0.050,2.238) (0.000,2.239) (-0.050,2.238) (-0.100,2.236) (-0.150,2.231) (-0.200,2.225) (-0.250,2.218) (-0.300,2.208) (-0.350,2.197) (-0.400,2.184) (-0.450,2.169) (-0.500,2.152) (-0.550,2.133) (-0.600,2.112) (-0.650,2.089) (-0.700,2.063) (-0.750,2.035) (-0.800,2.004) (-0.843,1.975) (-0.900,1.932) (-0.940,1.900) (-0.996,1.850) (-1.046,1.800) (-1.091,1.750) (-1.125,1.709) (-1.169,1.650) (-1.200,1.604) (-1.233,1.550) (-1.260,1.500) (-1.284,1.450) (-1.306,1.400) (-1.325,1.350) (-1.341,1.300) (-1.355,1.250) (-1.367,1.200) (-1.376,1.150) (-1.383,1.100) (-1.387,1.050) (-1.389,1.000) (-1.389,0.950) (-1.386,0.900) (-1.381,0.850) (-1.373,0.800) (-1.362,0.750) (-1.349,0.700) (-1.333,0.650) (-1.313,0.600) (-1.290,0.550) (-1.263,0.500) (-1.232,0.450) (-1.229,0.446)}
\EQ{(1.224,0.476) (1.250,0.522) (1.275,0.574) (1.296,0.625) (1.313,0.675) (1.326,0.725) (1.337,0.775) (1.345,0.825) (1.350,0.875) (1.353,0.925) (1.354,0.975) (1.352,1.025) (1.347,1.075) (1.340,1.125) (1.331,1.175) (1.319,1.225) (1.305,1.275) (1.288,1.325) (1.269,1.375) (1.246,1.425) (1.221,1.475) (1.194,1.525) (1.162,1.575) (1.128,1.625) (1.089,1.675) (1.050,1.721) (1.000,1.774) (0.950,1.822) (0.900,1.865) (0.856,1.900) (0.800,1.940) (0.750,1.972) (0.702,2.000) (0.656,2.025) (0.604,2.050) (0.550,2.074) (0.500,2.093) (0.450,2.111) (0.400,2.126) (0.350,2.139) (0.300,2.151) (0.250,2.161) (0.200,2.168) (0.150,2.174) (0.100,2.179) (0.050,2.181) (0.000,2.182) (-0.050,2.181) (-0.100,2.179) (-0.150,2.174) (-0.200,2.168) (-0.250,2.161) (-0.300,2.151) (-0.350,2.139) (-0.400,2.126) (-0.450,2.111) (-0.500,2.093) (-0.547,2.075) (-0.600,2.052) (-0.650,2.028) (-0.700,2.001) (-0.745,1.975) (-0.800,1.940) (-0.850,1.904) (-0.900,1.865) (-0.947,1.825) (-0.999,1.775) (-1.047,1.725) (-1.089,1.675) (-1.125,1.629) (-1.162,1.575) (-1.194,1.525) (-1.221,1.475) (-1.246,1.425) (-1.269,1.375) (-1.288,1.325) (-1.305,1.275) (-1.319,1.225) (-1.331,1.175) (-1.340,1.125) (-1.347,1.075) (-1.352,1.025) (-1.354,0.975) (-1.353,0.925) (-1.350,0.875) (-1.345,0.825) (-1.337,0.775) (-1.326,0.725) (-1.313,0.675) (-1.296,0.625) (-1.275,0.575) (-1.252,0.525) (-1.225,0.477)}

    \draw[mmaEQsolo] (-10,0) -- (10,0);}
  \fill[white] (0,0) circle (2.0);
  \draw[-{Stealth[length=3.4pt,width=3pt]}, mmarojo, line width=1pt] (0,-1.15) -- (0,1.2);
  \fill (0,1.45) circle (0.3) node[right=1pt, font=\tiny, inner sep=0.5pt] {$q$};
  \fill (0,-1.45) circle (0.3) node[right=1pt, font=\tiny, inner sep=0.5pt] {$-q$};
  \colorlet{marcoc}{grispizarra!60}
  \marcoMMA{10}{5}{$x/\ell$}{$z/\ell$}
\end{tikzpicture}\\[1pt]{\small\scshape (a)}
\end{minipage}\hfill
\begin{minipage}[t]{0.325\textwidth}\centering
{\footnotesize Equipotenciales y líneas de campo}\\[3pt]
\begin{tikzpicture}[x=0.198cm, y=0.198cm]
  \panelMMA{10}{mmafondo}{1.2}{%
    % Dipolo puntual - equipotenciales
% Generado por scripts/generar_datos.py (no editar a mano).
\EQ{(1.224,-0.476) (1.250,-0.522) (1.275,-0.574) (1.296,-0.625) (1.313,-0.675) (1.326,-0.725) (1.337,-0.775) (1.345,-0.825) (1.350,-0.875) (1.353,-0.925) (1.354,-0.975) (1.352,-1.025) (1.347,-1.075) (1.340,-1.125) (1.331,-1.175) (1.319,-1.225) (1.305,-1.275) (1.288,-1.325) (1.269,-1.375) (1.246,-1.425) (1.221,-1.475) (1.194,-1.525) (1.162,-1.575) (1.128,-1.625) (1.089,-1.675) (1.050,-1.721) (1.000,-1.774) (0.950,-1.822) (0.900,-1.865) (0.856,-1.900) (0.800,-1.940) (0.750,-1.972) (0.702,-2.000) (0.656,-2.025) (0.604,-2.050) (0.550,-2.074) (0.500,-2.093) (0.450,-2.111) (0.400,-2.126) (0.350,-2.139) (0.300,-2.151) (0.250,-2.161) (0.200,-2.168) (0.150,-2.174) (0.100,-2.179) (0.050,-2.181) (0.000,-2.182) (-0.050,-2.181) (-0.100,-2.179) (-0.150,-2.174) (-0.200,-2.168) (-0.250,-2.161) (-0.300,-2.151) (-0.350,-2.139) (-0.400,-2.126) (-0.450,-2.111) (-0.500,-2.093) (-0.547,-2.075) (-0.600,-2.052) (-0.650,-2.028) (-0.700,-2.001) (-0.745,-1.975) (-0.800,-1.940) (-0.850,-1.904) (-0.900,-1.865) (-0.947,-1.825) (-0.999,-1.775) (-1.047,-1.725) (-1.089,-1.675) (-1.125,-1.629) (-1.162,-1.575) (-1.194,-1.525) (-1.221,-1.475) (-1.246,-1.425) (-1.269,-1.375) (-1.288,-1.325) (-1.305,-1.275) (-1.319,-1.225) (-1.331,-1.175) (-1.340,-1.125) (-1.347,-1.075) (-1.352,-1.025) (-1.354,-0.975) (-1.353,-0.925) (-1.350,-0.875) (-1.345,-0.825) (-1.337,-0.775) (-1.326,-0.725) (-1.313,-0.675) (-1.296,-0.625) (-1.275,-0.575) (-1.252,-0.525) (-1.225,-0.477)}
\EQ{(1.229,-0.446) (1.263,-0.500) (1.290,-0.550) (1.313,-0.600) (1.333,-0.650) (1.349,-0.700) (1.362,-0.750) (1.373,-0.800) (1.381,-0.850) (1.386,-0.900) (1.389,-0.950) (1.389,-1.000) (1.387,-1.050) (1.383,-1.100) (1.376,-1.150) (1.367,-1.200) (1.355,-1.250) (1.341,-1.300) (1.325,-1.349) (1.306,-1.400) (1.284,-1.450) (1.260,-1.500) (1.233,-1.550) (1.203,-1.600) (1.175,-1.642) (1.132,-1.700) (1.100,-1.740) (1.050,-1.796) (1.000,-1.846) (0.950,-1.891) (0.909,-1.925) (0.850,-1.970) (0.806,-2.000) (0.750,-2.035) (0.700,-2.063) (0.650,-2.089) (0.600,-2.112) (0.550,-2.133) (0.500,-2.152) (0.450,-2.169) (0.400,-2.184) (0.350,-2.197) (0.300,-2.208) (0.250,-2.218) (0.200,-2.225) (0.150,-2.231) (0.100,-2.236) (0.050,-2.238) (0.000,-2.239) (-0.050,-2.238) (-0.100,-2.236) (-0.150,-2.231) (-0.200,-2.225) (-0.250,-2.218) (-0.300,-2.208) (-0.350,-2.197) (-0.400,-2.184) (-0.450,-2.169) (-0.500,-2.152) (-0.550,-2.133) (-0.600,-2.112) (-0.650,-2.089) (-0.700,-2.063) (-0.750,-2.035) (-0.800,-2.004) (-0.843,-1.975) (-0.900,-1.932) (-0.940,-1.900) (-0.996,-1.850) (-1.046,-1.800) (-1.091,-1.750) (-1.125,-1.709) (-1.169,-1.650) (-1.200,-1.604) (-1.233,-1.550) (-1.260,-1.500) (-1.284,-1.450) (-1.306,-1.400) (-1.325,-1.350) (-1.341,-1.300) (-1.355,-1.250) (-1.367,-1.200) (-1.376,-1.150) (-1.383,-1.100) (-1.387,-1.050) (-1.389,-1.000) (-1.389,-0.950) (-1.386,-0.900) (-1.381,-0.850) (-1.373,-0.800) (-1.362,-0.750) (-1.349,-0.700) (-1.333,-0.650) (-1.313,-0.600) (-1.290,-0.550) (-1.263,-0.500) (-1.232,-0.450) (-1.229,-0.446)}
\EQ{(1.244,-0.431) (1.274,-0.475) (1.300,-0.518) (1.325,-0.565) (1.350,-0.619) (1.372,-0.675) (1.388,-0.725) (1.400,-0.775) (1.411,-0.825) (1.419,-0.875) (1.424,-0.925) (1.427,-0.975) (1.427,-1.025) (1.425,-1.075) (1.421,-1.125) (1.414,-1.175) (1.406,-1.225) (1.395,-1.275) (1.381,-1.325) (1.365,-1.375) (1.347,-1.425) (1.326,-1.475) (1.303,-1.525) (1.277,-1.575) (1.250,-1.621) (1.215,-1.675) (1.180,-1.725) (1.141,-1.775) (1.100,-1.822) (1.050,-1.874) (1.000,-1.922) (0.950,-1.965) (0.906,-2.000) (0.850,-2.040) (0.800,-2.073) (0.755,-2.100) (0.709,-2.125) (0.660,-2.150) (0.605,-2.175) (0.550,-2.198) (0.500,-2.216) (0.450,-2.233) (0.400,-2.247) (0.350,-2.260) (0.300,-2.271) (0.250,-2.280) (0.200,-2.287) (0.150,-2.293) (0.100,-2.297) (0.050,-2.299) (0.000,-2.300) (-0.050,-2.299) (-0.100,-2.297) (-0.150,-2.293) (-0.200,-2.287) (-0.250,-2.280) (-0.300,-2.271) (-0.350,-2.260) (-0.400,-2.247) (-0.450,-2.233) (-0.500,-2.216) (-0.550,-2.198) (-0.600,-2.177) (-0.650,-2.155) (-0.700,-2.130) (-0.750,-2.103) (-0.796,-2.075) (-0.850,-2.040) (-0.900,-2.004) (-0.950,-1.965) (-0.997,-1.925) (-1.049,-1.875) (-1.097,-1.825) (-1.141,-1.775) (-1.175,-1.731) (-1.215,-1.675) (-1.248,-1.625) (-1.275,-1.578) (-1.300,-1.531) (-1.325,-1.478) (-1.347,-1.425) (-1.365,-1.375) (-1.381,-1.325) (-1.395,-1.275) (-1.406,-1.225) (-1.414,-1.175) (-1.421,-1.125) (-1.425,-1.075) (-1.427,-1.025) (-1.427,-0.975) (-1.424,-0.925) (-1.419,-0.875) (-1.411,-0.825) (-1.400,-0.775) (-1.388,-0.725) (-1.372,-0.675) (-1.352,-0.625) (-1.330,-0.575) (-1.304,-0.525) (-1.275,-0.477) (-1.244,-0.431)}
\EQ{(1.250,-0.405) (1.285,-0.450) (1.318,-0.500) (1.347,-0.550) (1.373,-0.600) (1.395,-0.650) (1.413,-0.700) (1.429,-0.750) (1.442,-0.800) (1.452,-0.850) (1.460,-0.900) (1.465,-0.950) (1.468,-1.000) (1.468,-1.050) (1.467,-1.100) (1.463,-1.150) (1.457,-1.200) (1.448,-1.250) (1.438,-1.300) (1.425,-1.350) (1.410,-1.400) (1.392,-1.450) (1.372,-1.500) (1.350,-1.550) (1.325,-1.600) (1.300,-1.645) (1.266,-1.700) (1.232,-1.750) (1.200,-1.793) (1.154,-1.850) (1.109,-1.900) (1.059,-1.950) (1.004,-2.000) (0.950,-2.044) (0.909,-2.075) (0.850,-2.116) (0.800,-2.147) (0.751,-2.175) (0.704,-2.200) (0.652,-2.225) (0.600,-2.248) (0.550,-2.268) (0.500,-2.285) (0.450,-2.301) (0.400,-2.315) (0.350,-2.328) (0.300,-2.338) (0.250,-2.347) (0.200,-2.354) (0.150,-2.360) (0.100,-2.364) (0.050,-2.366) (0.000,-2.367) (-0.050,-2.366) (-0.100,-2.364) (-0.150,-2.360) (-0.200,-2.354) (-0.250,-2.347) (-0.300,-2.338) (-0.350,-2.328) (-0.400,-2.315) (-0.450,-2.301) (-0.500,-2.285) (-0.550,-2.268) (-0.600,-2.248) (-0.650,-2.226) (-0.700,-2.202) (-0.750,-2.176) (-0.795,-2.150) (-0.850,-2.116) (-0.900,-2.081) (-0.942,-2.050) (-1.000,-2.003) (-1.050,-1.958) (-1.100,-1.909) (-1.150,-1.854) (-1.195,-1.800) (-1.225,-1.760) (-1.266,-1.700) (-1.297,-1.650) (-1.325,-1.600) (-1.350,-1.550) (-1.372,-1.500) (-1.392,-1.450) (-1.410,-1.400) (-1.425,-1.350) (-1.438,-1.300) (-1.448,-1.250) (-1.457,-1.200) (-1.463,-1.150) (-1.467,-1.100) (-1.468,-1.050) (-1.468,-1.000) (-1.465,-0.950) (-1.460,-0.900) (-1.452,-0.850) (-1.442,-0.800) (-1.429,-0.750) (-1.413,-0.700) (-1.395,-0.650) (-1.373,-0.600) (-1.350,-0.555) (-1.325,-0.512) (-1.285,-0.450) (-1.250,-0.405)}
\EQ{(1.251,-0.374) (1.295,-0.425) (1.325,-0.465) (1.365,-0.525) (1.393,-0.575) (1.418,-0.625) (1.440,-0.675) (1.458,-0.725) (1.473,-0.775) (1.486,-0.825) (1.496,-0.875) (1.504,-0.925) (1.510,-0.975) (1.513,-1.025) (1.514,-1.075) (1.512,-1.125) (1.509,-1.175) (1.503,-1.225) (1.495,-1.275) (1.486,-1.325) (1.473,-1.375) (1.459,-1.425) (1.442,-1.475) (1.423,-1.525) (1.402,-1.575) (1.378,-1.625) (1.352,-1.675) (1.325,-1.721) (1.290,-1.775) (1.255,-1.825) (1.216,-1.875) (1.175,-1.923) (1.127,-1.975) (1.075,-2.025) (1.025,-2.070) (0.975,-2.110) (0.925,-2.147) (0.875,-2.182) (0.825,-2.213) (0.775,-2.242) (0.725,-2.268) (0.675,-2.292) (0.625,-2.314) (0.575,-2.334) (0.525,-2.352) (0.475,-2.369) (0.425,-2.383) (0.375,-2.396) (0.325,-2.407) (0.275,-2.416) (0.225,-2.424) (0.175,-2.430) (0.125,-2.435) (0.075,-2.438) (0.025,-2.440) (-0.050,-2.439) (-0.100,-2.437) (-0.150,-2.433) (-0.200,-2.427) (-0.250,-2.420) (-0.300,-2.412) (-0.350,-2.402) (-0.400,-2.390) (-0.450,-2.376) (-0.500,-2.361) (-0.550,-2.344) (-0.599,-2.325) (-0.650,-2.304) (-0.700,-2.280) (-0.750,-2.255) (-0.800,-2.228) (-0.846,-2.200) (-0.900,-2.165) (-0.950,-2.129) (-1.000,-2.090) (-1.048,-2.050) (-1.100,-2.002) (-1.150,-1.951) (-1.195,-1.900) (-1.236,-1.850) (-1.273,-1.800) (-1.307,-1.750) (-1.337,-1.700) (-1.365,-1.650) (-1.390,-1.600) (-1.413,-1.550) (-1.433,-1.500) (-1.450,-1.453) (-1.466,-1.400) (-1.480,-1.350) (-1.491,-1.300) (-1.500,-1.250) (-1.506,-1.200) (-1.511,-1.150) (-1.513,-1.100) (-1.514,-1.050) (-1.512,-1.000) (-1.507,-0.950) (-1.501,-0.900) (-1.492,-0.850) (-1.480,-0.800) (-1.466,-0.750) (-1.449,-0.700) (-1.429,-0.650) (-1.406,-0.600) (-1.379,-0.550) (-1.350,-0.502) (-1.314,-0.450) (-1.275,-0.401) (-1.251,-0.374)}
\EQ{(1.266,-0.359) (1.305,-0.400) (1.346,-0.450) (1.382,-0.500) (1.414,-0.550) (1.442,-0.600) (1.467,-0.650) (1.488,-0.700) (1.506,-0.750) (1.522,-0.800) (1.534,-0.850) (1.545,-0.900) (1.553,-0.950) (1.559,-1.000) (1.562,-1.050) (1.563,-1.100) (1.563,-1.150) (1.560,-1.200) (1.555,-1.250) (1.547,-1.300) (1.538,-1.350) (1.527,-1.400) (1.513,-1.450) (1.498,-1.500) (1.480,-1.550) (1.460,-1.600) (1.437,-1.650) (1.412,-1.700) (1.385,-1.750) (1.354,-1.800) (1.325,-1.844) (1.285,-1.900) (1.250,-1.944) (1.201,-2.000) (1.154,-2.050) (1.101,-2.100) (1.050,-2.145) (1.000,-2.185) (0.950,-2.222) (0.908,-2.250) (0.850,-2.287) (0.800,-2.315) (0.750,-2.342) (0.700,-2.366) (0.650,-2.388) (0.600,-2.409) (0.550,-2.427) (0.500,-2.443) (0.450,-2.458) (0.400,-2.471) (0.350,-2.483) (0.300,-2.493) (0.250,-2.501) (0.200,-2.508) (0.150,-2.513) (0.100,-2.517) (0.050,-2.519) (0.000,-2.520) (-0.050,-2.519) (-0.100,-2.517) (-0.150,-2.513) (-0.200,-2.508) (-0.250,-2.501) (-0.300,-2.493) (-0.350,-2.483) (-0.400,-2.471) (-0.450,-2.458) (-0.500,-2.443) (-0.550,-2.427) (-0.600,-2.409) (-0.650,-2.388) (-0.700,-2.366) (-0.750,-2.342) (-0.800,-2.315) (-0.850,-2.287) (-0.900,-2.256) (-0.945,-2.225) (-1.000,-2.185) (-1.044,-2.150) (-1.100,-2.101) (-1.150,-2.054) (-1.200,-2.001) (-1.245,-1.950) (-1.285,-1.900) (-1.321,-1.850) (-1.350,-1.807) (-1.385,-1.750) (-1.412,-1.700) (-1.437,-1.650) (-1.460,-1.600) (-1.480,-1.550) (-1.498,-1.500) (-1.513,-1.450) (-1.527,-1.400) (-1.538,-1.350) (-1.547,-1.300) (-1.555,-1.250) (-1.560,-1.200) (-1.563,-1.150) (-1.563,-1.100) (-1.562,-1.050) (-1.559,-1.000) (-1.553,-0.950) (-1.545,-0.900) (-1.534,-0.850) (-1.522,-0.800) (-1.506,-0.750) (-1.488,-0.700) (-1.467,-0.650) (-1.442,-0.600) (-1.414,-0.550) (-1.382,-0.500) (-1.350,-0.456) (-1.305,-0.400) (-1.266,-0.359)}
\EQ{(1.275,-0.337) (1.314,-0.375) (1.350,-0.414) (1.400,-0.475) (1.436,-0.525) (1.467,-0.575) (1.495,-0.625) (1.519,-0.675) (1.540,-0.725) (1.559,-0.775) (1.574,-0.825) (1.587,-0.875) (1.598,-0.925) (1.606,-0.975) (1.612,-1.025) (1.616,-1.075) (1.618,-1.125) (1.618,-1.175) (1.616,-1.225) (1.611,-1.275) (1.605,-1.325) (1.597,-1.375) (1.586,-1.425) (1.574,-1.475) (1.559,-1.525) (1.543,-1.575) (1.524,-1.625) (1.503,-1.675) (1.479,-1.725) (1.454,-1.775) (1.425,-1.825) (1.394,-1.875) (1.360,-1.925) (1.325,-1.973) (1.283,-2.025) (1.250,-2.063) (1.216,-2.100) (1.175,-2.141) (1.125,-2.187) (1.081,-2.225) (1.025,-2.269) (0.983,-2.300) (0.925,-2.339) (0.875,-2.370) (0.825,-2.398) (0.775,-2.425) (0.725,-2.449) (0.675,-2.471) (0.625,-2.492) (0.575,-2.510) (0.525,-2.527) (0.475,-2.542) (0.425,-2.555) (0.375,-2.567) (0.325,-2.578) (0.275,-2.586) (0.225,-2.594) (0.175,-2.599) (0.125,-2.604) (0.075,-2.607) (0.025,-2.608) (-0.050,-2.608) (-0.100,-2.605) (-0.150,-2.602) (-0.200,-2.597) (-0.250,-2.590) (-0.300,-2.582) (-0.350,-2.573) (-0.400,-2.562) (-0.450,-2.549) (-0.500,-2.535) (-0.550,-2.519) (-0.600,-2.501) (-0.650,-2.482) (-0.700,-2.460) (-0.750,-2.437) (-0.800,-2.412) (-0.850,-2.384) (-0.900,-2.355) (-0.946,-2.325) (-1.000,-2.288) (-1.050,-2.250) (-1.100,-2.209) (-1.139,-2.175) (-1.191,-2.125) (-1.239,-2.075) (-1.275,-2.034) (-1.323,-1.975) (-1.360,-1.925) (-1.394,-1.875) (-1.425,-1.825) (-1.450,-1.781) (-1.475,-1.734) (-1.500,-1.681) (-1.524,-1.625) (-1.543,-1.575) (-1.559,-1.525) (-1.574,-1.475) (-1.586,-1.425) (-1.597,-1.375) (-1.605,-1.325) (-1.611,-1.275) (-1.616,-1.225) (-1.618,-1.175) (-1.618,-1.125) (-1.616,-1.075) (-1.612,-1.025) (-1.606,-0.975) (-1.598,-0.925) (-1.587,-0.875) (-1.574,-0.825) (-1.559,-0.775) (-1.540,-0.725) (-1.519,-0.675) (-1.495,-0.625) (-1.467,-0.575) (-1.436,-0.525) (-1.400,-0.475) (-1.360,-0.425) (-1.325,-0.387) (-1.289,-0.350) (-1.275,-0.337)}
\EQ{(1.275,-0.308) (1.323,-0.350) (1.374,-0.400) (1.419,-0.450) (1.450,-0.489) (1.493,-0.550) (1.524,-0.600) (1.550,-0.647) (1.575,-0.698) (1.597,-0.750) (1.616,-0.800) (1.632,-0.850) (1.645,-0.900) (1.656,-0.950) (1.665,-1.000) (1.672,-1.050) (1.676,-1.100) (1.679,-1.150) (1.679,-1.200) (1.678,-1.250) (1.674,-1.300) (1.669,-1.350) (1.662,-1.400) (1.652,-1.450) (1.641,-1.500) (1.628,-1.550) (1.613,-1.600) (1.595,-1.650) (1.576,-1.700) (1.554,-1.750) (1.530,-1.800) (1.504,-1.850) (1.475,-1.900) (1.444,-1.950) (1.410,-2.000) (1.375,-2.047) (1.333,-2.100) (1.289,-2.150) (1.250,-2.191) (1.200,-2.240) (1.162,-2.275) (1.102,-2.325) (1.050,-2.365) (1.001,-2.400) (0.950,-2.434) (0.900,-2.464) (0.850,-2.492) (0.800,-2.518) (0.750,-2.543) (0.700,-2.565) (0.650,-2.585) (0.600,-2.604) (0.550,-2.621) (0.500,-2.636) (0.450,-2.650) (0.400,-2.662) (0.350,-2.672) (0.300,-2.682) (0.250,-2.689) (0.200,-2.696) (0.150,-2.700) (0.100,-2.704) (0.050,-2.706) (0.000,-2.707) (-0.050,-2.706) (-0.100,-2.704) (-0.150,-2.700) (-0.200,-2.696) (-0.250,-2.689) (-0.300,-2.682) (-0.350,-2.672) (-0.400,-2.662) (-0.450,-2.650) (-0.500,-2.636) (-0.550,-2.621) (-0.600,-2.604) (-0.650,-2.585) (-0.700,-2.565) (-0.750,-2.543) (-0.800,-2.518) (-0.850,-2.492) (-0.900,-2.464) (-0.950,-2.434) (-1.000,-2.401) (-1.050,-2.365) (-1.100,-2.327) (-1.150,-2.285) (-1.189,-2.250) (-1.241,-2.200) (-1.289,-2.150) (-1.325,-2.109) (-1.373,-2.050) (-1.410,-2.000) (-1.444,-1.950) (-1.475,-1.900) (-1.504,-1.850) (-1.530,-1.800) (-1.554,-1.750) (-1.575,-1.702) (-1.595,-1.650) (-1.613,-1.600) (-1.628,-1.550) (-1.641,-1.500) (-1.652,-1.450) (-1.662,-1.400) (-1.669,-1.350) (-1.674,-1.300) (-1.678,-1.250) (-1.679,-1.200) (-1.679,-1.150) (-1.676,-1.100) (-1.672,-1.050) (-1.665,-1.000) (-1.656,-0.950) (-1.645,-0.900) (-1.632,-0.850) (-1.616,-0.800) (-1.597,-0.750) (-1.576,-0.700) (-1.552,-0.650) (-1.525,-0.601) (-1.493,-0.550) (-1.458,-0.500) (-1.425,-0.458) (-1.375,-0.401) (-1.325,-0.352) (-1.275,-0.308)}
\EQ{(1.275,-0.280) (1.325,-0.319) (1.375,-0.362) (1.414,-0.400) (1.450,-0.438) (1.500,-0.497) (1.539,-0.550) (1.572,-0.600) (1.600,-0.648) (1.625,-0.696) (1.650,-0.750) (1.670,-0.800) (1.688,-0.850) (1.703,-0.900) (1.716,-0.950) (1.726,-1.000) (1.734,-1.050) (1.741,-1.100) (1.745,-1.150) (1.747,-1.200) (1.748,-1.250) (1.746,-1.300) (1.743,-1.350) (1.738,-1.400) (1.731,-1.450) (1.722,-1.500) (1.711,-1.550) (1.698,-1.600) (1.683,-1.650) (1.667,-1.700) (1.648,-1.750) (1.627,-1.800) (1.604,-1.850) (1.579,-1.900) (1.552,-1.950) (1.525,-1.995) (1.489,-2.050) (1.454,-2.100) (1.416,-2.150) (1.375,-2.200) (1.330,-2.250) (1.282,-2.300) (1.229,-2.350) (1.175,-2.397) (1.125,-2.437) (1.075,-2.475) (1.025,-2.509) (0.975,-2.542) (0.925,-2.571) (0.875,-2.599) (0.825,-2.625) (0.775,-2.649) (0.725,-2.671) (0.675,-2.691) (0.625,-2.710) (0.575,-2.727) (0.525,-2.742) (0.475,-2.756) (0.425,-2.768) (0.375,-2.779) (0.325,-2.789) (0.275,-2.797) (0.225,-2.804) (0.175,-2.809) (0.125,-2.813) (0.075,-2.816) (0.025,-2.817) (-0.050,-2.817) (-0.100,-2.815) (-0.150,-2.811) (-0.200,-2.807) (-0.250,-2.800) (-0.300,-2.793) (-0.350,-2.784) (-0.400,-2.774) (-0.450,-2.762) (-0.500,-2.749) (-0.550,-2.735) (-0.600,-2.719) (-0.650,-2.701) (-0.700,-2.681) (-0.750,-2.660) (-0.800,-2.637) (-0.850,-2.612) (-0.900,-2.586) (-0.950,-2.557) (-1.000,-2.526) (-1.050,-2.492) (-1.100,-2.456) (-1.141,-2.425) (-1.200,-2.376) (-1.250,-2.331) (-1.300,-2.282) (-1.350,-2.228) (-1.396,-2.175) (-1.436,-2.125) (-1.472,-2.075) (-1.506,-2.025) (-1.537,-1.975) (-1.566,-1.925) (-1.592,-1.875) (-1.616,-1.825) (-1.638,-1.775) (-1.658,-1.725) (-1.675,-1.675) (-1.691,-1.625) (-1.705,-1.575) (-1.716,-1.525) (-1.726,-1.475) (-1.734,-1.425) (-1.741,-1.375) (-1.745,-1.325) (-1.747,-1.275) (-1.748,-1.225) (-1.746,-1.175) (-1.743,-1.125) (-1.738,-1.075) (-1.730,-1.025) (-1.721,-0.975) (-1.709,-0.925) (-1.695,-0.875) (-1.679,-0.825) (-1.660,-0.775) (-1.639,-0.725) (-1.614,-0.675) (-1.587,-0.625) (-1.556,-0.575) (-1.525,-0.530) (-1.482,-0.475) (-1.438,-0.425) (-1.400,-0.386) (-1.362,-0.350) (-1.301,-0.300) (-1.275,-0.280)}
\EQ{(1.288,-0.262) (1.342,-0.300) (1.400,-0.347) (1.450,-0.391) (1.500,-0.442) (1.550,-0.499) (1.590,-0.550) (1.625,-0.600) (1.656,-0.650) (1.684,-0.700) (1.708,-0.750) (1.730,-0.800) (1.750,-0.850) (1.767,-0.900) (1.781,-0.950) (1.794,-1.000) (1.804,-1.050) (1.812,-1.100) (1.818,-1.150) (1.823,-1.200) (1.825,-1.250) (1.825,-1.300) (1.824,-1.350) (1.821,-1.400) (1.816,-1.450) (1.810,-1.500) (1.801,-1.550) (1.791,-1.600) (1.779,-1.650) (1.765,-1.700) (1.750,-1.748) (1.732,-1.800) (1.712,-1.850) (1.691,-1.900) (1.667,-1.950) (1.641,-2.000) (1.613,-2.050) (1.583,-2.100) (1.550,-2.149) (1.514,-2.200) (1.476,-2.250) (1.434,-2.300) (1.400,-2.339) (1.366,-2.375) (1.325,-2.416) (1.289,-2.450) (1.250,-2.484) (1.201,-2.525) (1.150,-2.565) (1.101,-2.600) (1.050,-2.634) (1.000,-2.666) (0.950,-2.695) (0.900,-2.722) (0.850,-2.748) (0.800,-2.771) (0.750,-2.793) (0.700,-2.813) (0.650,-2.831) (0.600,-2.848) (0.550,-2.864) (0.500,-2.878) (0.450,-2.890) (0.400,-2.901) (0.350,-2.911) (0.300,-2.919) (0.250,-2.926) (0.200,-2.932) (0.150,-2.937) (0.100,-2.940) (0.050,-2.942) (0.000,-2.943) (-0.050,-2.942) (-0.100,-2.940) (-0.150,-2.937) (-0.200,-2.932) (-0.250,-2.926) (-0.300,-2.919) (-0.350,-2.911) (-0.400,-2.901) (-0.450,-2.890) (-0.500,-2.878) (-0.550,-2.864) (-0.600,-2.848) (-0.650,-2.831) (-0.700,-2.813) (-0.750,-2.793) (-0.800,-2.771) (-0.850,-2.748) (-0.900,-2.722) (-0.950,-2.695) (-1.000,-2.666) (-1.050,-2.634) (-1.100,-2.601) (-1.150,-2.565) (-1.200,-2.526) (-1.250,-2.484) (-1.289,-2.450) (-1.325,-2.416) (-1.366,-2.375) (-1.400,-2.339) (-1.434,-2.300) (-1.475,-2.251) (-1.514,-2.200) (-1.550,-2.150) (-1.583,-2.100) (-1.613,-2.050) (-1.641,-2.000) (-1.667,-1.950) (-1.691,-1.900) (-1.712,-1.850) (-1.732,-1.800) (-1.749,-1.750) (-1.765,-1.700) (-1.779,-1.650) (-1.791,-1.600) (-1.801,-1.550) (-1.810,-1.500) (-1.816,-1.450) (-1.821,-1.400) (-1.824,-1.350) (-1.825,-1.300) (-1.825,-1.250) (-1.823,-1.200) (-1.818,-1.150) (-1.812,-1.100) (-1.804,-1.050) (-1.794,-1.000) (-1.781,-0.950) (-1.767,-0.900) (-1.750,-0.851) (-1.730,-0.800) (-1.708,-0.750) (-1.684,-0.700) (-1.656,-0.650) (-1.625,-0.601) (-1.590,-0.550) (-1.551,-0.500) (-1.508,-0.450) (-1.459,-0.400) (-1.404,-0.350) (-1.350,-0.307) (-1.307,-0.275) (-1.288,-0.262)}
\EQ{(1.300,-0.243) (1.350,-0.275) (1.400,-0.310) (1.450,-0.348) (1.500,-0.391) (1.550,-0.439) (1.600,-0.493) (1.647,-0.550) (1.684,-0.600) (1.717,-0.650) (1.747,-0.700) (1.774,-0.750) (1.798,-0.800) (1.819,-0.850) (1.838,-0.900) (1.855,-0.950) (1.869,-1.000) (1.881,-1.050) (1.892,-1.100) (1.900,-1.150) (1.906,-1.200) (1.911,-1.250) (1.914,-1.300) (1.915,-1.350) (1.914,-1.400) (1.911,-1.450) (1.907,-1.500) (1.901,-1.550) (1.894,-1.600) (1.885,-1.650) (1.873,-1.700) (1.861,-1.750) (1.846,-1.800) (1.830,-1.850) (1.812,-1.900) (1.792,-1.950) (1.770,-2.000) (1.746,-2.050) (1.720,-2.100) (1.692,-2.150) (1.661,-2.200) (1.629,-2.250) (1.593,-2.300) (1.555,-2.350) (1.515,-2.400) (1.475,-2.445) (1.425,-2.498) (1.375,-2.548) (1.325,-2.593) (1.275,-2.636) (1.226,-2.675) (1.175,-2.713) (1.125,-2.747) (1.083,-2.775) (1.025,-2.810) (0.975,-2.838) (0.925,-2.865) (0.875,-2.890) (0.825,-2.913) (0.775,-2.934) (0.725,-2.954) (0.675,-2.972) (0.625,-2.989) (0.575,-3.004) (0.525,-3.018) (0.475,-3.030) (0.425,-3.042) (0.375,-3.052) (0.325,-3.060) (0.275,-3.068) (0.225,-3.074) (0.175,-3.079) (0.125,-3.082) (0.075,-3.085) (0.025,-3.086) (-0.050,-3.086) (-0.100,-3.084) (-0.150,-3.081) (-0.200,-3.076) (-0.250,-3.071) (-0.300,-3.064) (-0.350,-3.056) (-0.400,-3.047) (-0.450,-3.036) (-0.500,-3.024) (-0.550,-3.011) (-0.600,-2.996) (-0.650,-2.980) (-0.700,-2.963) (-0.750,-2.944) (-0.800,-2.923) (-0.850,-2.901) (-0.900,-2.877) (-0.950,-2.852) (-0.999,-2.825) (-1.050,-2.795) (-1.100,-2.764) (-1.150,-2.730) (-1.192,-2.700) (-1.250,-2.656) (-1.300,-2.615) (-1.345,-2.575) (-1.398,-2.525) (-1.448,-2.475) (-1.493,-2.425) (-1.535,-2.375) (-1.575,-2.325) (-1.611,-2.275) (-1.645,-2.225) (-1.675,-2.178) (-1.706,-2.125) (-1.733,-2.075) (-1.758,-2.025) (-1.781,-1.975) (-1.802,-1.925) (-1.821,-1.875) (-1.838,-1.825) (-1.854,-1.775) (-1.867,-1.725) (-1.879,-1.675) (-1.889,-1.625) (-1.898,-1.575) (-1.905,-1.525) (-1.910,-1.475) (-1.913,-1.425) (-1.915,-1.375) (-1.914,-1.325) (-1.913,-1.275) (-1.909,-1.225) (-1.903,-1.175) (-1.896,-1.125) (-1.887,-1.075) (-1.875,-1.025) (-1.862,-0.975) (-1.847,-0.925) (-1.829,-0.875) (-1.809,-0.825) (-1.786,-0.775) (-1.761,-0.725) (-1.733,-0.675) (-1.701,-0.625) (-1.666,-0.575) (-1.627,-0.525) (-1.584,-0.475) (-1.536,-0.425) (-1.482,-0.375) (-1.425,-0.328) (-1.375,-0.292) (-1.325,-0.258) (-1.300,-0.243)}
\EQ{(1.300,-0.216) (1.350,-0.244) (1.400,-0.274) (1.450,-0.308) (1.500,-0.345) (1.550,-0.385) (1.595,-0.425) (1.646,-0.475) (1.692,-0.525) (1.725,-0.565) (1.770,-0.625) (1.800,-0.669) (1.834,-0.725) (1.862,-0.775) (1.887,-0.825) (1.909,-0.875) (1.929,-0.925) (1.947,-0.975) (1.962,-1.025) (1.975,-1.075) (1.987,-1.125) (1.997,-1.175) (2.004,-1.225) (2.011,-1.275) (2.015,-1.325) (2.017,-1.375) (2.018,-1.425) (2.017,-1.475) (2.015,-1.525) (2.011,-1.575) (2.005,-1.625) (1.998,-1.675) (1.989,-1.725) (1.979,-1.775) (1.967,-1.825) (1.953,-1.875) (1.937,-1.925) (1.920,-1.975) (1.901,-2.025) (1.880,-2.075) (1.858,-2.125) (1.833,-2.175) (1.807,-2.225) (1.778,-2.275) (1.750,-2.321) (1.714,-2.375) (1.679,-2.425) (1.641,-2.475) (1.601,-2.525) (1.557,-2.575) (1.510,-2.625) (1.475,-2.661) (1.434,-2.700) (1.378,-2.750) (1.325,-2.794) (1.285,-2.825) (1.225,-2.869) (1.180,-2.900) (1.125,-2.935) (1.075,-2.965) (1.025,-2.993) (0.975,-3.020) (0.925,-3.045) (0.875,-3.068) (0.825,-3.089) (0.775,-3.109) (0.725,-3.128) (0.675,-3.145) (0.625,-3.161) (0.576,-3.175) (0.525,-3.188) (0.476,-3.200) (0.425,-3.211) (0.375,-3.220) (0.325,-3.229) (0.275,-3.236) (0.225,-3.241) (0.175,-3.246) (0.125,-3.249) (0.075,-3.252) (0.025,-3.253) (-0.050,-3.252) (-0.100,-3.251) (-0.150,-3.248) (-0.200,-3.244) (-0.250,-3.239) (-0.300,-3.232) (-0.350,-3.225) (-0.400,-3.216) (-0.450,-3.206) (-0.500,-3.195) (-0.550,-3.182) (-0.600,-3.168) (-0.650,-3.153) (-0.700,-3.137) (-0.750,-3.119) (-0.798,-3.100) (-0.850,-3.079) (-0.900,-3.056) (-0.950,-3.032) (-1.000,-3.007) (-1.050,-2.980) (-1.100,-2.950) (-1.150,-2.919) (-1.200,-2.886) (-1.250,-2.851) (-1.300,-2.814) (-1.348,-2.775) (-1.400,-2.731) (-1.450,-2.685) (-1.486,-2.650) (-1.525,-2.610) (-1.575,-2.555) (-1.621,-2.500) (-1.661,-2.450) (-1.697,-2.400) (-1.731,-2.350) (-1.763,-2.300) (-1.793,-2.250) (-1.820,-2.200) (-1.846,-2.150) (-1.869,-2.100) (-1.891,-2.050) (-1.911,-2.000) (-1.929,-1.950) (-1.945,-1.900) (-1.960,-1.850) (-1.973,-1.800) (-1.984,-1.750) (-1.994,-1.700) (-2.002,-1.650) (-2.008,-1.600) (-2.013,-1.550) (-2.017,-1.500) (-2.018,-1.450) (-2.018,-1.400) (-2.016,-1.350) (-2.013,-1.300) (-2.008,-1.250) (-2.001,-1.200) (-1.992,-1.150) (-1.982,-1.100) (-1.969,-1.050) (-1.955,-1.000) (-1.938,-0.950) (-1.919,-0.900) (-1.898,-0.850) (-1.875,-0.801) (-1.850,-0.753) (-1.825,-0.710) (-1.787,-0.650) (-1.752,-0.600) (-1.713,-0.550) (-1.675,-0.506) (-1.625,-0.454) (-1.575,-0.407) (-1.525,-0.364) (-1.475,-0.326) (-1.425,-0.291) (-1.375,-0.259) (-1.325,-0.230) (-1.300,-0.216)}
\EQ{(1.300,-0.191) (1.350,-0.215) (1.400,-0.241) (1.450,-0.270) (1.499,-0.300) (1.550,-0.335) (1.600,-0.372) (1.650,-0.414) (1.691,-0.450) (1.742,-0.500) (1.788,-0.550) (1.825,-0.594) (1.868,-0.650) (1.900,-0.696) (1.934,-0.750) (1.963,-0.800) (1.989,-0.850) (2.013,-0.900) (2.034,-0.950) (2.053,-1.000) (2.070,-1.050) (2.085,-1.100) (2.098,-1.150) (2.109,-1.200) (2.118,-1.250) (2.126,-1.300) (2.132,-1.350) (2.137,-1.400) (2.139,-1.450) (2.141,-1.500) (2.140,-1.550) (2.138,-1.600) (2.135,-1.650) (2.130,-1.700) (2.123,-1.750) (2.115,-1.800) (2.106,-1.850) (2.095,-1.900) (2.082,-1.950) (2.068,-2.000) (2.052,-2.050) (2.035,-2.100) (2.016,-2.150) (1.995,-2.200) (1.972,-2.250) (1.950,-2.296) (1.925,-2.344) (1.900,-2.389) (1.864,-2.450) (1.831,-2.500) (1.800,-2.546) (1.760,-2.600) (1.725,-2.645) (1.679,-2.700) (1.634,-2.750) (1.600,-2.786) (1.562,-2.825) (1.525,-2.860) (1.481,-2.900) (1.425,-2.947) (1.375,-2.987) (1.325,-3.025) (1.275,-3.060) (1.225,-3.093) (1.175,-3.124) (1.132,-3.150) (1.087,-3.175) (1.040,-3.200) (0.990,-3.225) (0.936,-3.250) (0.879,-3.275) (0.825,-3.297) (0.775,-3.315) (0.725,-3.333) (0.675,-3.349) (0.625,-3.364) (0.575,-3.377) (0.525,-3.390) (0.475,-3.401) (0.425,-3.411) (0.375,-3.420) (0.325,-3.427) (0.275,-3.434) (0.225,-3.439) (0.175,-3.444) (0.125,-3.447) (0.075,-3.449) (0.025,-3.450) (-0.037,-3.450) (-0.100,-3.448) (-0.150,-3.445) (-0.200,-3.442) (-0.250,-3.437) (-0.300,-3.431) (-0.350,-3.424) (-0.400,-3.415) (-0.450,-3.406) (-0.500,-3.395) (-0.550,-3.384) (-0.600,-3.371) (-0.650,-3.356) (-0.700,-3.341) (-0.748,-3.325) (-0.800,-3.306) (-0.850,-3.287) (-0.900,-3.266) (-0.950,-3.244) (-1.000,-3.220) (-1.050,-3.195) (-1.100,-3.168) (-1.150,-3.139) (-1.200,-3.109) (-1.250,-3.077) (-1.300,-3.043) (-1.350,-3.006) (-1.391,-2.975) (-1.450,-2.927) (-1.500,-2.883) (-1.550,-2.836) (-1.587,-2.800) (-1.625,-2.760) (-1.675,-2.705) (-1.721,-2.650) (-1.760,-2.600) (-1.797,-2.550) (-1.831,-2.500) (-1.864,-2.450) (-1.894,-2.400) (-1.922,-2.350) (-1.948,-2.300) (-1.972,-2.250) (-1.995,-2.200) (-2.016,-2.150) (-2.035,-2.100) (-2.052,-2.050) (-2.068,-2.000) (-2.082,-1.950) (-2.095,-1.900) (-2.106,-1.850) (-2.115,-1.800) (-2.123,-1.750) (-2.130,-1.700) (-2.135,-1.650) (-2.138,-1.600) (-2.140,-1.550) (-2.141,-1.500) (-2.139,-1.450) (-2.137,-1.400) (-2.132,-1.350) (-2.126,-1.300) (-2.118,-1.250) (-2.109,-1.200) (-2.098,-1.150) (-2.085,-1.100) (-2.070,-1.050) (-2.053,-1.000) (-2.034,-0.950) (-2.013,-0.900) (-1.989,-0.850) (-1.963,-0.800) (-1.934,-0.750) (-1.903,-0.700) (-1.868,-0.650) (-1.830,-0.600) (-1.788,-0.550) (-1.750,-0.508) (-1.700,-0.459) (-1.650,-0.414) (-1.603,-0.375) (-1.550,-0.335) (-1.500,-0.301) (-1.450,-0.270) (-1.400,-0.241) (-1.350,-0.215) (-1.300,-0.191)}
\EQ{(1.300,-0.165) (1.350,-0.186) (1.400,-0.208) (1.450,-0.233) (1.500,-0.259) (1.550,-0.288) (1.600,-0.319) (1.646,-0.350) (1.700,-0.390) (1.744,-0.425) (1.800,-0.474) (1.850,-0.522) (1.900,-0.575) (1.943,-0.625) (1.975,-0.665) (2.018,-0.725) (2.050,-0.773) (2.082,-0.825) (2.109,-0.875) (2.134,-0.925) (2.157,-0.975) (2.178,-1.025) (2.197,-1.075) (2.213,-1.125) (2.228,-1.175) (2.242,-1.225) (2.253,-1.275) (2.263,-1.325) (2.271,-1.375) (2.277,-1.425) (2.282,-1.475) (2.286,-1.525) (2.288,-1.575) (2.288,-1.625) (2.288,-1.675) (2.285,-1.725) (2.281,-1.775) (2.276,-1.825) (2.269,-1.875) (2.261,-1.925) (2.252,-1.975) (2.241,-2.025) (2.228,-2.075) (2.215,-2.125) (2.199,-2.175) (2.182,-2.225) (2.164,-2.275) (2.144,-2.325) (2.122,-2.375) (2.100,-2.423) (2.075,-2.473) (2.050,-2.520) (2.025,-2.565) (1.989,-2.625) (1.956,-2.675) (1.925,-2.720) (1.885,-2.775) (1.850,-2.820) (1.805,-2.875) (1.761,-2.925) (1.725,-2.964) (1.690,-3.000) (1.650,-3.039) (1.611,-3.075) (1.554,-3.125) (1.500,-3.170) (1.450,-3.209) (1.400,-3.245) (1.358,-3.275) (1.300,-3.313) (1.250,-3.344) (1.200,-3.373) (1.152,-3.400) (1.104,-3.425) (1.053,-3.450) (1.000,-3.475) (0.950,-3.497) (0.900,-3.517) (0.850,-3.536) (0.800,-3.554) (0.750,-3.571) (0.700,-3.587) (0.650,-3.601) (0.600,-3.614) (0.550,-3.626) (0.500,-3.637) (0.450,-3.647) (0.400,-3.656) (0.350,-3.664) (0.300,-3.670) (0.250,-3.676) (0.200,-3.680) (0.150,-3.684) (0.100,-3.687) (0.050,-3.688) (0.000,-3.689) (-0.050,-3.688) (-0.100,-3.687) (-0.150,-3.684) (-0.200,-3.680) (-0.250,-3.676) (-0.300,-3.670) (-0.350,-3.664) (-0.400,-3.656) (-0.450,-3.647) (-0.500,-3.637) (-0.550,-3.626) (-0.600,-3.614) (-0.650,-3.601) (-0.700,-3.587) (-0.750,-3.571) (-0.800,-3.554) (-0.850,-3.536) (-0.900,-3.517) (-0.950,-3.497) (-1.000,-3.475) (-1.050,-3.452) (-1.100,-3.427) (-1.150,-3.401) (-1.197,-3.375) (-1.250,-3.344) (-1.300,-3.313) (-1.350,-3.280) (-1.394,-3.250) (-1.450,-3.209) (-1.494,-3.175) (-1.550,-3.129) (-1.600,-3.085) (-1.638,-3.050) (-1.675,-3.015) (-1.714,-2.975) (-1.750,-2.937) (-1.800,-2.881) (-1.846,-2.825) (-1.885,-2.775) (-1.922,-2.725) (-1.956,-2.675) (-1.989,-2.625) (-2.019,-2.575) (-2.047,-2.525) (-2.074,-2.475) (-2.099,-2.425) (-2.122,-2.375) (-2.144,-2.325) (-2.164,-2.275) (-2.182,-2.225) (-2.199,-2.175) (-2.215,-2.125) (-2.228,-2.075) (-2.241,-2.025) (-2.252,-1.975) (-2.261,-1.925) (-2.269,-1.875) (-2.276,-1.825) (-2.281,-1.775) (-2.285,-1.725) (-2.288,-1.675) (-2.288,-1.625) (-2.288,-1.575) (-2.286,-1.525) (-2.282,-1.475) (-2.277,-1.425) (-2.271,-1.375) (-2.263,-1.325) (-2.253,-1.275) (-2.242,-1.225) (-2.228,-1.175) (-2.213,-1.125) (-2.197,-1.075) (-2.178,-1.025) (-2.157,-0.975) (-2.134,-0.925) (-2.109,-0.875) (-2.082,-0.825) (-2.051,-0.775) (-2.018,-0.725) (-1.982,-0.675) (-1.950,-0.634) (-1.900,-0.575) (-1.853,-0.525) (-1.801,-0.475) (-1.750,-0.430) (-1.700,-0.390) (-1.650,-0.353) (-1.600,-0.319) (-1.550,-0.288) (-1.500,-0.259) (-1.450,-0.233) (-1.400,-0.208) (-1.350,-0.186) (-1.300,-0.165)}
\EQ{(1.300,-0.141) (1.350,-0.158) (1.400,-0.177) (1.450,-0.197) (1.500,-0.220) (1.550,-0.243) (1.600,-0.269) (1.650,-0.297) (1.697,-0.325) (1.750,-0.359) (1.800,-0.394) (1.841,-0.425) (1.900,-0.473) (1.950,-0.517) (2.000,-0.566) (2.050,-0.619) (2.099,-0.675) (2.138,-0.725) (2.174,-0.775) (2.208,-0.825) (2.239,-0.875) (2.267,-0.925) (2.293,-0.975) (2.317,-1.025) (2.339,-1.075) (2.359,-1.125) (2.377,-1.175) (2.394,-1.225) (2.409,-1.275) (2.422,-1.325) (2.433,-1.375) (2.443,-1.425) (2.451,-1.475) (2.458,-1.525) (2.464,-1.575) (2.468,-1.625) (2.470,-1.675) (2.472,-1.725) (2.472,-1.775) (2.470,-1.825) (2.467,-1.875) (2.463,-1.925) (2.458,-1.975) (2.451,-2.025) (2.443,-2.075) (2.434,-2.125) (2.423,-2.175) (2.411,-2.225) (2.398,-2.275) (2.383,-2.325) (2.367,-2.375) (2.350,-2.423) (2.330,-2.475) (2.310,-2.525) (2.288,-2.575) (2.265,-2.625) (2.240,-2.675) (2.213,-2.725) (2.184,-2.775) (2.154,-2.825) (2.125,-2.871) (2.089,-2.925) (2.053,-2.975) (2.015,-3.025) (1.975,-3.075) (1.933,-3.125) (1.888,-3.175) (1.850,-3.215) (1.800,-3.265) (1.763,-3.300) (1.725,-3.335) (1.679,-3.375) (1.625,-3.420) (1.575,-3.459) (1.525,-3.496) (1.475,-3.531) (1.425,-3.565) (1.375,-3.597) (1.328,-3.625) (1.285,-3.650) (1.239,-3.675) (1.191,-3.700) (1.141,-3.725) (1.087,-3.750) (1.031,-3.775) (0.975,-3.798) (0.925,-3.817) (0.875,-3.835) (0.825,-3.852) (0.775,-3.868) (0.725,-3.883) (0.675,-3.897) (0.625,-3.909) (0.575,-3.921) (0.525,-3.932) (0.475,-3.941) (0.425,-3.950) (0.375,-3.958) (0.325,-3.964) (0.275,-3.970) (0.225,-3.975) (0.175,-3.978) (0.125,-3.981) (0.075,-3.983) (0.025,-3.984) (-0.050,-3.984) (-0.100,-3.982) (-0.150,-3.980) (-0.200,-3.977) (-0.250,-3.972) (-0.300,-3.967) (-0.350,-3.961) (-0.400,-3.954) (-0.450,-3.946) (-0.500,-3.937) (-0.550,-3.926) (-0.600,-3.915) (-0.650,-3.903) (-0.700,-3.890) (-0.750,-3.876) (-0.800,-3.860) (-0.850,-3.844) (-0.900,-3.826) (-0.950,-3.808) (-1.000,-3.788) (-1.050,-3.767) (-1.100,-3.744) (-1.150,-3.721) (-1.200,-3.696) (-1.250,-3.669) (-1.300,-3.641) (-1.350,-3.612) (-1.400,-3.581) (-1.448,-3.550) (-1.500,-3.514) (-1.550,-3.478) (-1.600,-3.440) (-1.650,-3.400) (-1.700,-3.357) (-1.750,-3.312) (-1.790,-3.275) (-1.840,-3.225) (-1.875,-3.189) (-1.910,-3.150) (-1.950,-3.105) (-1.995,-3.050) (-2.034,-3.000) (-2.071,-2.950) (-2.100,-2.908) (-2.139,-2.850) (-2.170,-2.800) (-2.199,-2.750) (-2.225,-2.703) (-2.250,-2.654) (-2.275,-2.603) (-2.299,-2.550) (-2.320,-2.500) (-2.340,-2.450) (-2.358,-2.400) (-2.375,-2.351) (-2.391,-2.300) (-2.405,-2.250) (-2.417,-2.200) (-2.429,-2.150) (-2.439,-2.100) (-2.447,-2.050) (-2.455,-2.000) (-2.461,-1.950) (-2.466,-1.900) (-2.469,-1.850) (-2.471,-1.800) (-2.472,-1.750) (-2.471,-1.700) (-2.469,-1.650) (-2.466,-1.600) (-2.461,-1.550) (-2.455,-1.500) (-2.447,-1.450) (-2.438,-1.400) (-2.427,-1.350) (-2.415,-1.300) (-2.401,-1.250) (-2.386,-1.200) (-2.369,-1.150) (-2.350,-1.101) (-2.329,-1.050) (-2.306,-1.000) (-2.281,-0.950) (-2.253,-0.900) (-2.225,-0.852) (-2.192,-0.800) (-2.157,-0.750) (-2.125,-0.708) (-2.078,-0.650) (-2.033,-0.600) (-1.984,-0.550) (-1.931,-0.500) (-1.875,-0.452) (-1.825,-0.413) (-1.775,-0.376) (-1.725,-0.343) (-1.675,-0.312) (-1.625,-0.283) (-1.575,-0.256) (-1.525,-0.231) (-1.475,-0.208) (-1.425,-0.187) (-1.375,-0.167) (-1.327,-0.150) (-1.300,-0.141)}
\EQ{(1.300,-0.117) (1.350,-0.131) (1.400,-0.146) (1.450,-0.163) (1.500,-0.181) (1.549,-0.200) (1.600,-0.221) (1.650,-0.244) (1.700,-0.268) (1.750,-0.293) (1.800,-0.321) (1.849,-0.350) (1.900,-0.382) (1.950,-0.416) (1.997,-0.450) (2.050,-0.492) (2.090,-0.525) (2.145,-0.575) (2.196,-0.625) (2.243,-0.675) (2.287,-0.725) (2.325,-0.772) (2.365,-0.825) (2.400,-0.875) (2.432,-0.925) (2.462,-0.975) (2.490,-1.025) (2.516,-1.075) (2.540,-1.125) (2.562,-1.175) (2.582,-1.225) (2.600,-1.275) (2.617,-1.325) (2.632,-1.375) (2.646,-1.425) (2.658,-1.475) (2.669,-1.525) (2.679,-1.575) (2.687,-1.625) (2.693,-1.675) (2.699,-1.725) (2.703,-1.775) (2.706,-1.825) (2.707,-1.875) (2.708,-1.925) (2.707,-1.975) (2.705,-2.025) (2.701,-2.075) (2.697,-2.125) (2.691,-2.175) (2.684,-2.225) (2.676,-2.275) (2.667,-2.325) (2.656,-2.375) (2.644,-2.425) (2.631,-2.475) (2.617,-2.525) (2.602,-2.575) (2.585,-2.625) (2.567,-2.675) (2.547,-2.725) (2.527,-2.775) (2.504,-2.825) (2.481,-2.875) (2.456,-2.925) (2.430,-2.975) (2.402,-3.025) (2.375,-3.070) (2.341,-3.125) (2.308,-3.175) (2.275,-3.223) (2.237,-3.275) (2.200,-3.323) (2.158,-3.375) (2.125,-3.414) (2.075,-3.470) (2.025,-3.523) (1.975,-3.574) (1.925,-3.621) (1.875,-3.667) (1.836,-3.700) (1.776,-3.750) (1.725,-3.790) (1.677,-3.825) (1.625,-3.862) (1.575,-3.896) (1.530,-3.925) (1.475,-3.959) (1.425,-3.988) (1.375,-4.016) (1.325,-4.043) (1.275,-4.068) (1.225,-4.092) (1.175,-4.115) (1.125,-4.137) (1.075,-4.158) (1.025,-4.177) (0.975,-4.196) (0.925,-4.213) (0.875,-4.229) (0.825,-4.245) (0.775,-4.259) (0.725,-4.272) (0.675,-4.285) (0.625,-4.296) (0.575,-4.307) (0.525,-4.317) (0.475,-4.325) (0.425,-4.333) (0.375,-4.340) (0.325,-4.346) (0.275,-4.351) (0.225,-4.356) (0.175,-4.359) (0.125,-4.362) (0.075,-4.363) (0.025,-4.364) (-0.050,-4.364) (-0.100,-4.363) (-0.150,-4.361) (-0.200,-4.358) (-0.250,-4.354) (-0.300,-4.349) (-0.350,-4.343) (-0.400,-4.337) (-0.450,-4.329) (-0.500,-4.321) (-0.550,-4.312) (-0.600,-4.302) (-0.650,-4.291) (-0.700,-4.279) (-0.750,-4.266) (-0.800,-4.252) (-0.850,-4.237) (-0.900,-4.221) (-0.950,-4.204) (-1.000,-4.187) (-1.050,-4.168) (-1.100,-4.147) (-1.150,-4.126) (-1.200,-4.104) (-1.250,-4.080) (-1.300,-4.056) (-1.350,-4.030) (-1.400,-4.002) (-1.448,-3.975) (-1.500,-3.944) (-1.550,-3.912) (-1.600,-3.879) (-1.642,-3.850) (-1.700,-3.808) (-1.744,-3.775) (-1.800,-3.730) (-1.850,-3.688) (-1.894,-3.650) (-1.948,-3.600) (-1.999,-3.550) (-2.047,-3.500) (-2.093,-3.450) (-2.137,-3.400) (-2.175,-3.354) (-2.218,-3.300) (-2.250,-3.257) (-2.291,-3.200) (-2.325,-3.150) (-2.357,-3.100) (-2.387,-3.050) (-2.416,-3.000) (-2.443,-2.950) (-2.469,-2.900) (-2.493,-2.850) (-2.516,-2.800) (-2.537,-2.750) (-2.557,-2.700) (-2.575,-2.653) (-2.593,-2.600) (-2.610,-2.550) (-2.624,-2.500) (-2.638,-2.450) (-2.650,-2.400) (-2.662,-2.350) (-2.672,-2.300) (-2.680,-2.250) (-2.688,-2.200) (-2.694,-2.150) (-2.699,-2.100) (-2.703,-2.050) (-2.706,-2.000) (-2.707,-1.950) (-2.708,-1.900) (-2.707,-1.850) (-2.704,-1.800) (-2.701,-1.750) (-2.696,-1.700) (-2.690,-1.650) (-2.683,-1.600) (-2.674,-1.550) (-2.664,-1.500) (-2.652,-1.450) (-2.639,-1.400) (-2.625,-1.351) (-2.609,-1.300) (-2.591,-1.250) (-2.572,-1.200) (-2.551,-1.150) (-2.528,-1.100) (-2.503,-1.050) (-2.476,-1.000) (-2.450,-0.954) (-2.416,-0.900) (-2.383,-0.850) (-2.350,-0.805) (-2.308,-0.750) (-2.275,-0.711) (-2.225,-0.655) (-2.175,-0.604) (-2.125,-0.557) (-2.075,-0.513) (-2.029,-0.475) (-1.975,-0.434) (-1.927,-0.400) (-1.875,-0.366) (-1.825,-0.335) (-1.775,-0.307) (-1.725,-0.280) (-1.675,-0.255) (-1.625,-0.232) (-1.575,-0.211) (-1.525,-0.191) (-1.475,-0.172) (-1.425,-0.155) (-1.375,-0.139) (-1.325,-0.124) (-1.300,-0.117)}
\EQ{(1.300,-0.093) (1.350,-0.104) (1.400,-0.116) (1.450,-0.130) (1.500,-0.144) (1.550,-0.159) (1.600,-0.175) (1.650,-0.193) (1.700,-0.211) (1.750,-0.231) (1.800,-0.252) (1.850,-0.275) (1.900,-0.299) (1.950,-0.324) (1.997,-0.350) (2.041,-0.375) (2.100,-0.412) (2.150,-0.444) (2.194,-0.475) (2.250,-0.517) (2.292,-0.550) (2.350,-0.599) (2.400,-0.645) (2.450,-0.693) (2.500,-0.746) (2.548,-0.800) (2.589,-0.850) (2.625,-0.896) (2.664,-0.950) (2.697,-1.000) (2.725,-1.044) (2.759,-1.100) (2.786,-1.150) (2.812,-1.200) (2.836,-1.250) (2.858,-1.300) (2.879,-1.350) (2.898,-1.400) (2.916,-1.450) (2.932,-1.500) (2.947,-1.550) (2.960,-1.600) (2.972,-1.650) (2.983,-1.700) (2.993,-1.750) (3.001,-1.800) (3.008,-1.850) (3.014,-1.900) (3.019,-1.950) (3.023,-2.000) (3.025,-2.050) (3.027,-2.100) (3.027,-2.150) (3.026,-2.200) (3.025,-2.250) (3.022,-2.300) (3.018,-2.350) (3.013,-2.400) (3.006,-2.450) (2.999,-2.500) (2.991,-2.550) (2.981,-2.600) (2.971,-2.650) (2.959,-2.700) (2.947,-2.750) (2.933,-2.800) (2.918,-2.850) (2.902,-2.900) (2.885,-2.950) (2.866,-3.000) (2.847,-3.050) (2.826,-3.100) (2.804,-3.150) (2.781,-3.200) (2.756,-3.250) (2.730,-3.300) (2.703,-3.350) (2.675,-3.399) (2.645,-3.450) (2.613,-3.500) (2.580,-3.550) (2.550,-3.594) (2.509,-3.650) (2.475,-3.696) (2.432,-3.750) (2.400,-3.789) (2.350,-3.846) (2.301,-3.900) (2.253,-3.950) (2.203,-4.000) (2.151,-4.050) (2.100,-4.096) (2.050,-4.139) (2.007,-4.175) (1.950,-4.220) (1.900,-4.258) (1.850,-4.295) (1.807,-4.325) (1.750,-4.363) (1.700,-4.395) (1.651,-4.425) (1.600,-4.455) (1.550,-4.483) (1.500,-4.510) (1.450,-4.536) (1.400,-4.561) (1.350,-4.585) (1.300,-4.607) (1.250,-4.629) (1.200,-4.649) (1.150,-4.669) (1.100,-4.687) (1.050,-4.705) (1.000,-4.722) (0.950,-4.737) (0.900,-4.752) (0.850,-4.766) (0.800,-4.779) (0.750,-4.792) (0.700,-4.803) (0.650,-4.814) (0.600,-4.824) (0.550,-4.833) (0.500,-4.841) (0.450,-4.848) (0.400,-4.855) (0.350,-4.861) (0.300,-4.866) (0.250,-4.870) (0.200,-4.873) (0.150,-4.876) (0.100,-4.878) (0.050,-4.879) (0.000,-4.880) (-0.050,-4.879) (-0.100,-4.878) (-0.150,-4.876) (-0.200,-4.873) (-0.250,-4.870) (-0.300,-4.866) (-0.350,-4.861) (-0.400,-4.855) (-0.450,-4.848) (-0.500,-4.841) (-0.550,-4.833) (-0.600,-4.824) (-0.650,-4.814) (-0.700,-4.803) (-0.750,-4.792) (-0.800,-4.779) (-0.850,-4.766) (-0.900,-4.752) (-0.950,-4.737) (-1.000,-4.722) (-1.050,-4.705) (-1.100,-4.687) (-1.150,-4.669) (-1.198,-4.650) (-1.250,-4.629) (-1.300,-4.607) (-1.350,-4.585) (-1.400,-4.561) (-1.450,-4.536) (-1.500,-4.510) (-1.550,-4.483) (-1.600,-4.455) (-1.650,-4.426) (-1.700,-4.395) (-1.750,-4.363) (-1.800,-4.330) (-1.843,-4.300) (-1.900,-4.258) (-1.944,-4.225) (-2.000,-4.181) (-2.050,-4.139) (-2.095,-4.100) (-2.150,-4.051) (-2.200,-4.003) (-2.250,-3.954) (-2.300,-3.901) (-2.347,-3.850) (-2.390,-3.800) (-2.425,-3.758) (-2.472,-3.700) (-2.509,-3.650) (-2.546,-3.600) (-2.575,-3.558) (-2.613,-3.500) (-2.645,-3.450) (-2.675,-3.400) (-2.700,-3.356) (-2.725,-3.310) (-2.750,-3.262) (-2.775,-3.212) (-2.800,-3.159) (-2.825,-3.102) (-2.847,-3.050) (-2.866,-3.000) (-2.885,-2.950) (-2.902,-2.900) (-2.918,-2.850) (-2.933,-2.800) (-2.947,-2.750) (-2.959,-2.700) (-2.971,-2.650) (-2.981,-2.600) (-2.991,-2.550) (-2.999,-2.500) (-3.006,-2.450) (-3.013,-2.400) (-3.018,-2.350) (-3.022,-2.300) (-3.025,-2.250) (-3.026,-2.200) (-3.027,-2.150) (-3.027,-2.100) (-3.025,-2.050) (-3.023,-2.000) (-3.019,-1.950) (-3.014,-1.900) (-3.008,-1.850) (-3.001,-1.800) (-2.993,-1.750) (-2.983,-1.700) (-2.972,-1.650) (-2.960,-1.600) (-2.947,-1.550) (-2.932,-1.500) (-2.916,-1.450) (-2.898,-1.400) (-2.879,-1.350) (-2.858,-1.300) (-2.836,-1.250) (-2.812,-1.200) (-2.786,-1.150) (-2.759,-1.100) (-2.729,-1.050) (-2.700,-1.004) (-2.664,-0.950) (-2.628,-0.900) (-2.589,-0.850) (-2.550,-0.802) (-2.504,-0.750) (-2.457,-0.700) (-2.406,-0.650) (-2.351,-0.600) (-2.300,-0.557) (-2.261,-0.525) (-2.200,-0.479) (-2.158,-0.450) (-2.100,-0.412) (-2.050,-0.381) (-2.000,-0.352) (-1.951,-0.325) (-1.902,-0.300) (-1.850,-0.275) (-1.800,-0.252) (-1.750,-0.231) (-1.700,-0.211) (-1.650,-0.193) (-1.600,-0.175) (-1.550,-0.159) (-1.500,-0.144) (-1.450,-0.130) (-1.400,-0.116) (-1.350,-0.104) (-1.300,-0.093)}
\EQ{(1.300,-0.070) (1.350,-0.078) (1.400,-0.087) (1.450,-0.097) (1.500,-0.107) (1.550,-0.118) (1.600,-0.130) (1.650,-0.143) (1.700,-0.157) (1.750,-0.171) (1.800,-0.187) (1.850,-0.203) (1.900,-0.220) (1.950,-0.239) (2.000,-0.258) (2.050,-0.279) (2.098,-0.300) (2.150,-0.324) (2.200,-0.348) (2.250,-0.374) (2.298,-0.400) (2.342,-0.425) (2.400,-0.460) (2.450,-0.492) (2.500,-0.525) (2.550,-0.561) (2.600,-0.598) (2.650,-0.638) (2.694,-0.675) (2.750,-0.724) (2.800,-0.772) (2.850,-0.822) (2.899,-0.875) (2.943,-0.925) (2.985,-0.975) (3.024,-1.025) (3.060,-1.075) (3.095,-1.125) (3.125,-1.171) (3.159,-1.225) (3.188,-1.275) (3.216,-1.325) (3.241,-1.375) (3.266,-1.425) (3.289,-1.475) (3.310,-1.525) (3.330,-1.575) (3.349,-1.625) (3.367,-1.675) (3.383,-1.725) (3.398,-1.775) (3.412,-1.825) (3.425,-1.875) (3.437,-1.925) (3.447,-1.975) (3.456,-2.025) (3.465,-2.075) (3.472,-2.125) (3.479,-2.175) (3.484,-2.225) (3.488,-2.275) (3.491,-2.325) (3.494,-2.375) (3.495,-2.425) (3.496,-2.475) (3.495,-2.525) (3.494,-2.575) (3.491,-2.625) (3.488,-2.675) (3.483,-2.725) (3.478,-2.775) (3.472,-2.825) (3.465,-2.875) (3.457,-2.925) (3.448,-2.975) (3.438,-3.025) (3.427,-3.075) (3.415,-3.125) (3.403,-3.175) (3.389,-3.225) (3.374,-3.275) (3.359,-3.325) (3.342,-3.375) (3.325,-3.423) (3.306,-3.475) (3.286,-3.525) (3.265,-3.575) (3.243,-3.625) (3.220,-3.675) (3.196,-3.725) (3.171,-3.775) (3.145,-3.825) (3.118,-3.875) (3.089,-3.925) (3.059,-3.975) (3.028,-4.025) (3.000,-4.068) (2.962,-4.125) (2.927,-4.175) (2.890,-4.225) (2.852,-4.275) (2.812,-4.325) (2.775,-4.370) (2.728,-4.425) (2.683,-4.475) (2.637,-4.525) (2.600,-4.563) (2.563,-4.600) (2.525,-4.637) (2.485,-4.675) (2.430,-4.725) (2.375,-4.773) (2.325,-4.815) (2.281,-4.850) (2.225,-4.894) (2.183,-4.925) (2.125,-4.967) (2.078,-5.000) (2.025,-5.035) (1.975,-5.067) (1.925,-5.098) (1.881,-5.125) (1.825,-5.157) (1.775,-5.185) (1.725,-5.212) (1.675,-5.238) (1.625,-5.262) (1.575,-5.286) (1.525,-5.309) (1.475,-5.331) (1.425,-5.352) (1.375,-5.373) (1.325,-5.392) (1.275,-5.411) (1.225,-5.428) (1.175,-5.445) (1.125,-5.462) (1.075,-5.477) (1.025,-5.492) (0.975,-5.505) (0.925,-5.519) (0.875,-5.531) (0.825,-5.543) (0.775,-5.553) (0.725,-5.564) (0.675,-5.573) (0.625,-5.582) (0.575,-5.590) (0.525,-5.597) (0.475,-5.604) (0.425,-5.610) (0.375,-5.616) (0.325,-5.620) (0.275,-5.624) (0.225,-5.628) (0.175,-5.630) (0.125,-5.632) (0.075,-5.634) (0.025,-5.634) (-0.050,-5.634) (-0.100,-5.633) (-0.150,-5.631) (-0.200,-5.629) (-0.250,-5.626) (-0.300,-5.622) (-0.350,-5.618) (-0.400,-5.613) (-0.450,-5.607) (-0.500,-5.601) (-0.550,-5.594) (-0.600,-5.586) (-0.650,-5.578) (-0.700,-5.569) (-0.750,-5.559) (-0.800,-5.548) (-0.850,-5.537) (-0.899,-5.525) (-0.950,-5.512) (-1.000,-5.499) (-1.050,-5.484) (-1.100,-5.469) (-1.150,-5.453) (-1.200,-5.437) (-1.250,-5.420) (-1.300,-5.401) (-1.350,-5.382) (-1.400,-5.363) (-1.450,-5.342) (-1.500,-5.320) (-1.550,-5.298) (-1.599,-5.275) (-1.650,-5.250) (-1.700,-5.225) (-1.747,-5.200) (-1.793,-5.175) (-1.838,-5.150) (-1.900,-5.114) (-1.950,-5.083) (-2.000,-5.051) (-2.050,-5.019) (-2.100,-4.985) (-2.149,-4.950) (-2.200,-4.913) (-2.249,-4.875) (-2.300,-4.835) (-2.343,-4.800) (-2.400,-4.751) (-2.450,-4.707) (-2.500,-4.661) (-2.538,-4.625) (-2.575,-4.588) (-2.613,-4.550) (-2.650,-4.511) (-2.700,-4.457) (-2.750,-4.400) (-2.792,-4.350) (-2.825,-4.309) (-2.871,-4.250) (-2.908,-4.200) (-2.944,-4.150) (-2.975,-4.106) (-3.012,-4.050) (-3.044,-4.000) (-3.074,-3.950) (-3.100,-3.906) (-3.125,-3.862) (-3.150,-3.816) (-3.175,-3.768) (-3.200,-3.718) (-3.225,-3.665) (-3.250,-3.610) (-3.275,-3.551) (-3.296,-3.500) (-3.315,-3.450) (-3.333,-3.400) (-3.350,-3.351) (-3.366,-3.300) (-3.382,-3.250) (-3.396,-3.200) (-3.409,-3.150) (-3.421,-3.100) (-3.433,-3.050) (-3.443,-3.000) (-3.453,-2.950) (-3.461,-2.900) (-3.469,-2.850) (-3.475,-2.800) (-3.481,-2.750) (-3.486,-2.700) (-3.490,-2.650) (-3.492,-2.600) (-3.494,-2.550) (-3.495,-2.500) (-3.496,-2.450) (-3.495,-2.400) (-3.493,-2.350) (-3.490,-2.300) (-3.486,-2.250) (-3.481,-2.200) (-3.475,-2.150) (-3.469,-2.100) (-3.461,-2.050) (-3.452,-2.000) (-3.442,-1.950) (-3.431,-1.900) (-3.419,-1.850) (-3.405,-1.800) (-3.391,-1.750) (-3.375,-1.700) (-3.358,-1.650) (-3.340,-1.600) (-3.321,-1.550) (-3.300,-1.501) (-3.277,-1.450) (-3.254,-1.400) (-3.229,-1.350) (-3.202,-1.300) (-3.175,-1.252) (-3.144,-1.200) (-3.112,-1.150) (-3.078,-1.100) (-3.042,-1.050) (-3.004,-1.000) (-2.964,-0.950) (-2.925,-0.904) (-2.876,-0.850) (-2.828,-0.800) (-2.777,-0.750) (-2.725,-0.702) (-2.675,-0.659) (-2.634,-0.625) (-2.575,-0.579) (-2.525,-0.543) (-2.475,-0.508) (-2.425,-0.475) (-2.375,-0.444) (-2.325,-0.415) (-2.275,-0.387) (-2.225,-0.361) (-2.175,-0.336) (-2.125,-0.312) (-2.075,-0.290) (-2.025,-0.269) (-1.975,-0.248) (-1.925,-0.230) (-1.875,-0.212) (-1.825,-0.195) (-1.775,-0.179) (-1.725,-0.164) (-1.676,-0.150) (-1.625,-0.137) (-1.575,-0.124) (-1.525,-0.113) (-1.475,-0.102) (-1.425,-0.092) (-1.375,-0.082) (-1.325,-0.074) (-1.300,-0.070)}
\EQ{(1.300,-0.046) (1.350,-0.052) (1.400,-0.058) (1.450,-0.064) (1.500,-0.071) (1.550,-0.079) (1.600,-0.086) (1.650,-0.095) (1.700,-0.104) (1.750,-0.113) (1.800,-0.123) (1.850,-0.134) (1.900,-0.145) (1.950,-0.157) (2.000,-0.170) (2.050,-0.183) (2.100,-0.197) (2.150,-0.212) (2.200,-0.227) (2.250,-0.243) (2.300,-0.260) (2.350,-0.278) (2.400,-0.297) (2.450,-0.317) (2.500,-0.337) (2.550,-0.359) (2.600,-0.381) (2.650,-0.405) (2.700,-0.429) (2.750,-0.455) (2.800,-0.482) (2.850,-0.510) (2.900,-0.539) (2.950,-0.570) (2.998,-0.600) (3.050,-0.635) (3.100,-0.670) (3.141,-0.700) (3.200,-0.745) (3.250,-0.785) (3.298,-0.825) (3.350,-0.871) (3.400,-0.917) (3.450,-0.966) (3.500,-1.017) (3.550,-1.071) (3.598,-1.125) (3.640,-1.175) (3.675,-1.219) (3.718,-1.275) (3.750,-1.319) (3.789,-1.375) (3.822,-1.425) (3.850,-1.470) (3.883,-1.525) (3.912,-1.575) (3.939,-1.625) (3.965,-1.675) (3.990,-1.725) (4.014,-1.775) (4.036,-1.825) (4.057,-1.875) (4.077,-1.925) (4.096,-1.975) (4.114,-2.025) (4.131,-2.075) (4.147,-2.125) (4.162,-2.175) (4.176,-2.225) (4.189,-2.275) (4.201,-2.325) (4.213,-2.375) (4.223,-2.425) (4.232,-2.475) (4.241,-2.525) (4.249,-2.575) (4.256,-2.625) (4.262,-2.675) (4.267,-2.725) (4.271,-2.775) (4.275,-2.825) (4.278,-2.875) (4.280,-2.925) (4.281,-2.975) (4.281,-3.025) (4.281,-3.075) (4.280,-3.125) (4.278,-3.175) (4.275,-3.225) (4.272,-3.275) (4.268,-3.325) (4.263,-3.375) (4.257,-3.425) (4.250,-3.475) (4.243,-3.525) (4.235,-3.575) (4.226,-3.625) (4.217,-3.675) (4.206,-3.725) (4.195,-3.775) (4.184,-3.825) (4.171,-3.875) (4.157,-3.925) (4.143,-3.975) (4.128,-4.025) (4.112,-4.075) (4.096,-4.125) (4.078,-4.175) (4.060,-4.225) (4.041,-4.275) (4.021,-4.325) (4.000,-4.375) (3.979,-4.425) (3.956,-4.475) (3.933,-4.525) (3.908,-4.575) (3.883,-4.625) (3.857,-4.675) (3.830,-4.725) (3.802,-4.775) (3.775,-4.821) (3.743,-4.875) (3.711,-4.925) (3.679,-4.975) (3.650,-5.019) (3.611,-5.075) (3.576,-5.125) (3.539,-5.175) (3.501,-5.225) (3.462,-5.275) (3.425,-5.321) (3.380,-5.375) (3.336,-5.425) (3.300,-5.466) (3.250,-5.520) (3.200,-5.573) (3.150,-5.623) (3.100,-5.672) (3.050,-5.720) (3.000,-5.766) (2.961,-5.800) (2.904,-5.850) (2.850,-5.895) (2.800,-5.935) (2.750,-5.974) (2.700,-6.013) (2.650,-6.049) (2.600,-6.085) (2.550,-6.120) (2.505,-6.150) (2.450,-6.186) (2.400,-6.218) (2.350,-6.249) (2.306,-6.275) (2.250,-6.308) (2.200,-6.336) (2.150,-6.363) (2.100,-6.390) (2.050,-6.415) (2.000,-6.440) (1.950,-6.464) (1.900,-6.488) (1.850,-6.510) (1.800,-6.532) (1.750,-6.553) (1.700,-6.574) (1.650,-6.593) (1.600,-6.612) (1.550,-6.631) (1.500,-6.648) (1.450,-6.665) (1.400,-6.682) (1.350,-6.698) (1.300,-6.713) (1.250,-6.727) (1.200,-6.741) (1.150,-6.754) (1.100,-6.767) (1.050,-6.779) (1.000,-6.791) (0.950,-6.801) (0.900,-6.812) (0.850,-6.821) (0.800,-6.831) (0.750,-6.839) (0.700,-6.847) (0.650,-6.854) (0.600,-6.861) (0.550,-6.868) (0.500,-6.873) (0.450,-6.879) (0.400,-6.883) (0.350,-6.887) (0.300,-6.891) (0.250,-6.894) (0.200,-6.896) (0.150,-6.898) (0.100,-6.900) (0.050,-6.900) (0.000,-6.901) (-0.050,-6.900) (-0.100,-6.900) (-0.150,-6.898) (-0.200,-6.896) (-0.250,-6.894) (-0.300,-6.891) (-0.350,-6.887) (-0.400,-6.883) (-0.450,-6.879) (-0.500,-6.873) (-0.550,-6.868) (-0.600,-6.861) (-0.650,-6.854) (-0.700,-6.847) (-0.750,-6.839) (-0.800,-6.831) (-0.850,-6.821) (-0.900,-6.812) (-0.950,-6.801) (-1.000,-6.791) (-1.050,-6.779) (-1.100,-6.767) (-1.150,-6.754) (-1.200,-6.741) (-1.250,-6.727) (-1.300,-6.713) (-1.350,-6.698) (-1.400,-6.682) (-1.450,-6.665) (-1.500,-6.648) (-1.550,-6.631) (-1.600,-6.612) (-1.650,-6.593) (-1.697,-6.575) (-1.750,-6.553) (-1.800,-6.532) (-1.850,-6.510) (-1.900,-6.488) (-1.950,-6.464) (-2.000,-6.440) (-2.050,-6.415) (-2.100,-6.390) (-2.150,-6.363) (-2.200,-6.336) (-2.250,-6.308) (-2.300,-6.279) (-2.348,-6.250) (-2.400,-6.218) (-2.450,-6.186) (-2.500,-6.154) (-2.543,-6.125) (-2.600,-6.085) (-2.649,-6.050) (-2.700,-6.013) (-2.749,-5.975) (-2.800,-5.935) (-2.844,-5.900) (-2.900,-5.853) (-2.950,-5.810) (-2.990,-5.775) (-3.044,-5.725) (-3.097,-5.675) (-3.148,-5.625) (-3.198,-5.575) (-3.246,-5.525) (-3.292,-5.475) (-3.336,-5.425) (-3.375,-5.380) (-3.421,-5.325) (-3.462,-5.275) (-3.500,-5.227) (-3.539,-5.175) (-3.575,-5.126) (-3.611,-5.075) (-3.646,-5.025) (-3.675,-4.981) (-3.711,-4.925) (-3.743,-4.875) (-3.773,-4.825) (-3.800,-4.778) (-3.825,-4.734) (-3.850,-4.688) (-3.875,-4.641) (-3.900,-4.592) (-3.925,-4.541) (-3.950,-4.488) (-3.975,-4.433) (-4.000,-4.376) (-4.021,-4.325) (-4.041,-4.275) (-4.060,-4.225) (-4.078,-4.175) (-4.096,-4.125) (-4.112,-4.075) (-4.128,-4.025) (-4.143,-3.975) (-4.157,-3.925) (-4.171,-3.875) (-4.184,-3.825) (-4.195,-3.775) (-4.206,-3.725) (-4.217,-3.675) (-4.226,-3.625) (-4.235,-3.575) (-4.243,-3.525) (-4.250,-3.475) (-4.257,-3.425) (-4.263,-3.375) (-4.268,-3.325) (-4.272,-3.275) (-4.275,-3.225) (-4.278,-3.175) (-4.280,-3.125) (-4.281,-3.075) (-4.281,-3.025) (-4.281,-2.975) (-4.280,-2.925) (-4.278,-2.875) (-4.275,-2.825) (-4.271,-2.775) (-4.267,-2.725) (-4.262,-2.675) (-4.256,-2.625) (-4.249,-2.575) (-4.241,-2.525) (-4.232,-2.475) (-4.223,-2.425) (-4.213,-2.375) (-4.201,-2.325) (-4.189,-2.275) (-4.176,-2.225) (-4.162,-2.175) (-4.147,-2.125) (-4.131,-2.075) (-4.114,-2.025) (-4.096,-1.975) (-4.077,-1.925) (-4.057,-1.875) (-4.036,-1.825) (-4.014,-1.775) (-3.990,-1.725) (-3.965,-1.675) (-3.939,-1.625) (-3.912,-1.575) (-3.883,-1.525) (-3.853,-1.475) (-3.825,-1.430) (-3.789,-1.375) (-3.754,-1.325) (-3.718,-1.275) (-3.680,-1.225) (-3.640,-1.175) (-3.600,-1.127) (-3.554,-1.075) (-3.508,-1.025) (-3.459,-0.975) (-3.408,-0.925) (-3.355,-0.875) (-3.300,-0.827) (-3.250,-0.785) (-3.207,-0.750) (-3.150,-0.707) (-3.107,-0.675) (-3.050,-0.635) (-3.000,-0.602) (-2.950,-0.570) (-2.900,-0.539) (-2.850,-0.510) (-2.800,-0.482) (-2.750,-0.455) (-2.700,-0.429) (-2.650,-0.405) (-2.600,-0.381) (-2.550,-0.359) (-2.500,-0.337) (-2.450,-0.317) (-2.400,-0.297) (-2.350,-0.278) (-2.300,-0.260) (-2.250,-0.243) (-2.200,-0.227) (-2.150,-0.212) (-2.100,-0.197) (-2.050,-0.183) (-2.000,-0.170) (-1.950,-0.157) (-1.900,-0.145) (-1.850,-0.134) (-1.800,-0.123) (-1.750,-0.113) (-1.700,-0.104) (-1.650,-0.095) (-1.600,-0.086) (-1.550,-0.079) (-1.500,-0.071) (-1.450,-0.064) (-1.400,-0.058) (-1.350,-0.052) (-1.300,-0.046)}
\EQ{(1.300,-0.023) (1.350,-0.026) (1.400,-0.029) (1.450,-0.032) (1.500,-0.035) (1.550,-0.039) (1.600,-0.043) (1.650,-0.047) (1.700,-0.052) (1.750,-0.056) (1.800,-0.061) (1.850,-0.067) (1.900,-0.072) (1.950,-0.078) (2.000,-0.084) (2.050,-0.091) (2.100,-0.098) (2.150,-0.105) (2.200,-0.112) (2.250,-0.120) (2.300,-0.128) (2.350,-0.137) (2.400,-0.146) (2.450,-0.155) (2.500,-0.165) (2.550,-0.175) (2.600,-0.186) (2.650,-0.197) (2.700,-0.209) (2.750,-0.220) (2.800,-0.233) (2.850,-0.246) (2.900,-0.259) (2.950,-0.273) (3.000,-0.287) (3.050,-0.302) (3.100,-0.318) (3.150,-0.334) (3.199,-0.350) (3.250,-0.367) (3.300,-0.385) (3.350,-0.403) (3.400,-0.422) (3.450,-0.442) (3.500,-0.462) (3.550,-0.483) (3.600,-0.504) (3.646,-0.525) (3.700,-0.550) (3.750,-0.573) (3.800,-0.598) (3.850,-0.623) (3.900,-0.649) (3.949,-0.675) (3.994,-0.700) (4.038,-0.725) (4.100,-0.761) (4.150,-0.792) (4.200,-0.823) (4.250,-0.856) (4.300,-0.889) (4.350,-0.923) (4.400,-0.959) (4.450,-0.996) (4.500,-1.034) (4.550,-1.073) (4.600,-1.113) (4.644,-1.150) (4.700,-1.198) (4.750,-1.243) (4.800,-1.289) (4.838,-1.325) (4.875,-1.362) (4.913,-1.400) (4.950,-1.438) (4.985,-1.475) (5.025,-1.519) (5.074,-1.575) (5.117,-1.625) (5.150,-1.665) (5.198,-1.725) (5.237,-1.775) (5.274,-1.825) (5.310,-1.875) (5.345,-1.925) (5.375,-1.970) (5.411,-2.025) (5.443,-2.075) (5.473,-2.125) (5.500,-2.171) (5.525,-2.214) (5.550,-2.259) (5.575,-2.306) (5.598,-2.350) (5.623,-2.400) (5.648,-2.450) (5.671,-2.500) (5.694,-2.550) (5.716,-2.600) (5.737,-2.650) (5.757,-2.700) (5.777,-2.750) (5.796,-2.800) (5.814,-2.850) (5.831,-2.900) (5.848,-2.950) (5.863,-3.000) (5.879,-3.050) (5.893,-3.100) (5.907,-3.150) (5.920,-3.200) (5.933,-3.250) (5.944,-3.300) (5.956,-3.350) (5.966,-3.400) (5.976,-3.450) (5.985,-3.500) (5.994,-3.550) (6.002,-3.600) (6.010,-3.650) (6.017,-3.700) (6.023,-3.750) (6.029,-3.800) (6.034,-3.850) (6.038,-3.900) (6.042,-3.950) (6.046,-4.000) (6.049,-4.050) (6.051,-4.100) (6.053,-4.150) (6.054,-4.200) (6.054,-4.250) (6.054,-4.300) (6.054,-4.350) (6.053,-4.400) (6.051,-4.450) (6.049,-4.500) (6.047,-4.550) (6.043,-4.600) (6.040,-4.650) (6.035,-4.700) (6.031,-4.750) (6.025,-4.800) (6.019,-4.850) (6.013,-4.900) (6.006,-4.950) (5.998,-5.000) (5.990,-5.050) (5.982,-5.100) (5.973,-5.150) (5.963,-5.200) (5.953,-5.250) (5.942,-5.300) (5.931,-5.350) (5.919,-5.400) (5.906,-5.450) (5.893,-5.500) (5.880,-5.550) (5.866,-5.600) (5.851,-5.650) (5.836,-5.700) (5.820,-5.750) (5.804,-5.800) (5.787,-5.850) (5.769,-5.900) (5.751,-5.950) (5.732,-6.000) (5.713,-6.050) (5.693,-6.100) (5.673,-6.150) (5.652,-6.200) (5.630,-6.250) (5.608,-6.300) (5.585,-6.350) (5.561,-6.400) (5.537,-6.450) (5.512,-6.500) (5.487,-6.550) (5.461,-6.600) (5.434,-6.650) (5.406,-6.700) (5.378,-6.750) (5.350,-6.799) (5.320,-6.850) (5.289,-6.900) (5.258,-6.950) (5.226,-7.000) (5.194,-7.050) (5.160,-7.100) (5.126,-7.150) (5.091,-7.200) (5.056,-7.250) (5.025,-7.292) (4.981,-7.350) (4.950,-7.391) (4.904,-7.450) (4.864,-7.500) (4.825,-7.547) (4.781,-7.600) (4.738,-7.650) (4.700,-7.693) (4.650,-7.749) (4.603,-7.800) (4.555,-7.850) (4.507,-7.900) (4.457,-7.950) (4.407,-8.000) (4.355,-8.050) (4.302,-8.100) (4.250,-8.147) (4.200,-8.192) (4.162,-8.225) (4.104,-8.275) (4.050,-8.320) (4.000,-8.361) (3.952,-8.400) (3.900,-8.441) (3.855,-8.475) (3.800,-8.517) (3.755,-8.550) (3.700,-8.589) (3.650,-8.625) (3.600,-8.659) (3.550,-8.693) (3.501,-8.725) (3.450,-8.758) (3.400,-8.790) (3.350,-8.821) (3.302,-8.850) (3.250,-8.881) (3.200,-8.910) (3.150,-8.939) (3.100,-8.967) (3.050,-8.994) (3.000,-9.021) (2.950,-9.047) (2.900,-9.072) (2.850,-9.097) (2.800,-9.122) (2.750,-9.146) (2.700,-9.169) (2.650,-9.192) (2.600,-9.214) (2.550,-9.236) (2.500,-9.257) (2.450,-9.278) (2.400,-9.298) (2.350,-9.318) (2.300,-9.337) (2.250,-9.356) (2.200,-9.374) (2.150,-9.392) (2.100,-9.410) (2.050,-9.427) (2.000,-9.443) (1.950,-9.459) (1.900,-9.475) (1.850,-9.490) (1.800,-9.504) (1.750,-9.519) (1.700,-9.533) (1.650,-9.546) (1.600,-9.559) (1.550,-9.571) (1.500,-9.583) (1.450,-9.595) (1.400,-9.606) (1.350,-9.617) (1.300,-9.628) (1.250,-9.638) (1.200,-9.647) (1.150,-9.656) (1.100,-9.665) (1.050,-9.674) (1.000,-9.682) (0.950,-9.689) (0.900,-9.696) (0.850,-9.703) (0.800,-9.710) (0.750,-9.716) (0.700,-9.721) (0.650,-9.726) (0.600,-9.731) (0.550,-9.736) (0.500,-9.740) (0.450,-9.743) (0.400,-9.747) (0.350,-9.750) (0.300,-9.752) (0.250,-9.754) (0.200,-9.756) (0.150,-9.757) (0.100,-9.758) (0.050,-9.759) (0.000,-9.759) (-0.050,-9.759) (-0.100,-9.758) (-0.150,-9.757) (-0.200,-9.756) (-0.250,-9.754) (-0.300,-9.752) (-0.350,-9.750) (-0.400,-9.747) (-0.450,-9.743) (-0.500,-9.740) (-0.550,-9.736) (-0.600,-9.731) (-0.650,-9.726) (-0.700,-9.721) (-0.750,-9.716) (-0.800,-9.710) (-0.850,-9.703) (-0.900,-9.696) (-0.950,-9.689) (-1.000,-9.682) (-1.050,-9.674) (-1.100,-9.665) (-1.150,-9.656) (-1.200,-9.647) (-1.250,-9.638) (-1.300,-9.628) (-1.350,-9.617) (-1.400,-9.606) (-1.450,-9.595) (-1.500,-9.583) (-1.550,-9.571) (-1.600,-9.559) (-1.650,-9.546) (-1.700,-9.533) (-1.750,-9.519) (-1.800,-9.504) (-1.850,-9.490) (-1.899,-9.475) (-1.950,-9.459) (-2.000,-9.443) (-2.050,-9.427) (-2.100,-9.410) (-2.150,-9.392) (-2.198,-9.375) (-2.250,-9.356) (-2.300,-9.337) (-2.350,-9.318) (-2.400,-9.298) (-2.450,-9.278) (-2.500,-9.257) (-2.550,-9.236) (-2.600,-9.214) (-2.650,-9.192) (-2.700,-9.169) (-2.750,-9.146) (-2.800,-9.122) (-2.850,-9.097) (-2.900,-9.072) (-2.950,-9.047) (-3.000,-9.021) (-3.050,-8.994) (-3.100,-8.967) (-3.150,-8.939) (-3.200,-8.910) (-3.250,-8.881) (-3.300,-8.851) (-3.344,-8.825) (-3.400,-8.790) (-3.450,-8.758) (-3.500,-8.726) (-3.550,-8.693) (-3.600,-8.659) (-3.650,-8.625) (-3.700,-8.589) (-3.750,-8.553) (-3.800,-8.517) (-3.850,-8.479) (-3.900,-8.441) (-3.950,-8.401) (-4.000,-8.361) (-4.045,-8.325) (-4.100,-8.279) (-4.150,-8.236) (-4.191,-8.200) (-4.247,-8.150) (-4.300,-8.101) (-4.350,-8.055) (-4.400,-8.007) (-4.450,-7.957) (-4.500,-7.907) (-4.550,-7.856) (-4.600,-7.803) (-4.649,-7.750) (-4.694,-7.700) (-4.738,-7.650) (-4.775,-7.607) (-4.823,-7.550) (-4.864,-7.500) (-4.900,-7.455) (-4.943,-7.400) (-4.975,-7.359) (-5.019,-7.300) (-5.050,-7.258) (-5.091,-7.200) (-5.125,-7.152) (-5.160,-7.100) (-5.194,-7.050) (-5.225,-7.002) (-5.258,-6.950) (-5.289,-6.900) (-5.320,-6.850) (-5.349,-6.800) (-5.375,-6.755) (-5.400,-6.711) (-5.425,-6.666) (-5.450,-6.620) (-5.474,-6.575) (-5.500,-6.525) (-5.525,-6.475) (-5.549,-6.425) (-5.573,-6.375) (-5.596,-6.325) (-5.619,-6.275) (-5.641,-6.225) (-5.662,-6.175) (-5.683,-6.125) (-5.703,-6.075) (-5.723,-6.025) (-5.742,-5.975) (-5.760,-5.925) (-5.778,-5.875) (-5.795,-5.825) (-5.812,-5.775) (-5.828,-5.725) (-5.843,-5.675) (-5.858,-5.625) (-5.873,-5.575) (-5.887,-5.525) (-5.900,-5.475) (-5.913,-5.425) (-5.925,-5.375) (-5.936,-5.325) (-5.947,-5.275) (-5.958,-5.225) (-5.968,-5.175) (-5.977,-5.125) (-5.986,-5.075) (-5.994,-5.025) (-6.002,-4.975) (-6.009,-4.925) (-6.016,-4.875) (-6.022,-4.825) (-6.028,-4.775) (-6.033,-4.725) (-6.038,-4.675) (-6.042,-4.625) (-6.045,-4.575) (-6.048,-4.525) (-6.050,-4.475) (-6.052,-4.425) (-6.054,-4.375) (-6.054,-4.325) (-6.055,-4.275) (-6.054,-4.225) (-6.053,-4.175) (-6.052,-4.125) (-6.050,-4.075) (-6.047,-4.025) (-6.044,-3.975) (-6.040,-3.925) (-6.036,-3.875) (-6.031,-3.825) (-6.026,-3.775) (-6.020,-3.725) (-6.013,-3.675) (-6.006,-3.625) (-5.998,-3.575) (-5.990,-3.525) (-5.981,-3.475) (-5.971,-3.425) (-5.961,-3.375) (-5.950,-3.325) (-5.939,-3.275) (-5.926,-3.225) (-5.914,-3.175) (-5.900,-3.125) (-5.886,-3.075) (-5.871,-3.025) (-5.856,-2.975) (-5.839,-2.925) (-5.822,-2.875) (-5.805,-2.825) (-5.786,-2.775) (-5.767,-2.725) (-5.747,-2.675) (-5.727,-2.625) (-5.705,-2.575) (-5.683,-2.525) (-5.660,-2.475) (-5.636,-2.425) (-5.611,-2.375) (-5.585,-2.325) (-5.559,-2.275) (-5.531,-2.225) (-5.503,-2.175) (-5.475,-2.128) (-5.443,-2.075) (-5.411,-2.025) (-5.378,-1.975) (-5.350,-1.933) (-5.310,-1.875) (-5.275,-1.827) (-5.237,-1.775) (-5.200,-1.727) (-5.158,-1.675) (-5.125,-1.634) (-5.075,-1.576) (-5.030,-1.525) (-4.985,-1.475) (-4.950,-1.438) (-4.913,-1.400) (-4.875,-1.362) (-4.838,-1.325) (-4.800,-1.289) (-4.758,-1.250) (-4.702,-1.200) (-4.650,-1.155) (-4.600,-1.113) (-4.553,-1.075) (-4.500,-1.034) (-4.456,-1.000) (-4.400,-0.959) (-4.352,-0.925) (-4.300,-0.889) (-4.250,-0.856) (-4.203,-0.825) (-4.150,-0.792) (-4.100,-0.761) (-4.050,-0.732) (-4.000,-0.703) (-3.950,-0.676) (-3.902,-0.650) (-3.854,-0.625) (-3.805,-0.600) (-3.754,-0.575) (-3.701,-0.550) (-3.650,-0.527) (-3.600,-0.504) (-3.550,-0.483) (-3.500,-0.462) (-3.450,-0.442) (-3.400,-0.422) (-3.350,-0.403) (-3.300,-0.385) (-3.250,-0.367) (-3.200,-0.350) (-3.150,-0.334) (-3.100,-0.318) (-3.050,-0.302) (-3.000,-0.287) (-2.950,-0.273) (-2.900,-0.259) (-2.850,-0.246) (-2.800,-0.233) (-2.750,-0.220) (-2.700,-0.209) (-2.650,-0.197) (-2.600,-0.186) (-2.550,-0.175) (-2.500,-0.165) (-2.450,-0.155) (-2.400,-0.146) (-2.350,-0.137) (-2.300,-0.128) (-2.250,-0.120) (-2.200,-0.112) (-2.150,-0.105) (-2.100,-0.098) (-2.050,-0.091) (-2.000,-0.084) (-1.950,-0.078) (-1.900,-0.072) (-1.850,-0.067) (-1.800,-0.061) (-1.750,-0.056) (-1.700,-0.052) (-1.650,-0.047) (-1.600,-0.043) (-1.550,-0.039) (-1.500,-0.035) (-1.450,-0.032) (-1.400,-0.029) (-1.350,-0.026) (-1.302,-0.023)}
\EQ{(1.300,0.023) (1.350,0.026) (1.400,0.029) (1.450,0.032) (1.500,0.035) (1.550,0.039) (1.600,0.043) (1.650,0.047) (1.700,0.052) (1.750,0.056) (1.800,0.061) (1.850,0.067) (1.900,0.072) (1.950,0.078) (2.000,0.084) (2.050,0.091) (2.100,0.098) (2.150,0.105) (2.200,0.112) (2.250,0.120) (2.300,0.128) (2.350,0.137) (2.400,0.146) (2.450,0.155) (2.500,0.165) (2.550,0.175) (2.600,0.186) (2.650,0.197) (2.700,0.209) (2.750,0.220) (2.800,0.233) (2.850,0.246) (2.900,0.259) (2.950,0.273) (3.000,0.287) (3.050,0.302) (3.100,0.318) (3.150,0.334) (3.199,0.350) (3.250,0.367) (3.300,0.385) (3.350,0.403) (3.400,0.422) (3.450,0.442) (3.500,0.462) (3.550,0.483) (3.600,0.504) (3.646,0.525) (3.700,0.550) (3.750,0.573) (3.800,0.598) (3.850,0.623) (3.900,0.649) (3.949,0.675) (3.994,0.700) (4.038,0.725) (4.100,0.761) (4.150,0.792) (4.200,0.823) (4.250,0.856) (4.300,0.889) (4.350,0.923) (4.400,0.959) (4.450,0.996) (4.500,1.034) (4.550,1.073) (4.600,1.113) (4.644,1.150) (4.700,1.198) (4.750,1.243) (4.800,1.289) (4.838,1.325) (4.875,1.362) (4.913,1.400) (4.950,1.438) (4.985,1.475) (5.025,1.519) (5.074,1.575) (5.117,1.625) (5.150,1.665) (5.198,1.725) (5.237,1.775) (5.274,1.825) (5.310,1.875) (5.345,1.925) (5.375,1.970) (5.411,2.025) (5.443,2.075) (5.473,2.125) (5.500,2.171) (5.525,2.214) (5.550,2.259) (5.575,2.306) (5.598,2.350) (5.623,2.400) (5.648,2.450) (5.671,2.500) (5.694,2.550) (5.716,2.600) (5.737,2.650) (5.757,2.700) (5.777,2.750) (5.796,2.800) (5.814,2.850) (5.831,2.900) (5.848,2.950) (5.863,3.000) (5.879,3.050) (5.893,3.100) (5.907,3.150) (5.920,3.200) (5.933,3.250) (5.944,3.300) (5.956,3.350) (5.966,3.400) (5.976,3.450) (5.985,3.500) (5.994,3.550) (6.002,3.600) (6.010,3.650) (6.017,3.700) (6.023,3.750) (6.029,3.800) (6.034,3.850) (6.038,3.900) (6.042,3.950) (6.046,4.000) (6.049,4.050) (6.051,4.100) (6.053,4.150) (6.054,4.200) (6.054,4.250) (6.054,4.300) (6.054,4.350) (6.053,4.400) (6.051,4.450) (6.049,4.500) (6.047,4.550) (6.043,4.600) (6.040,4.650) (6.035,4.700) (6.031,4.750) (6.025,4.800) (6.019,4.850) (6.013,4.900) (6.006,4.950) (5.998,5.000) (5.990,5.050) (5.982,5.100) (5.973,5.150) (5.963,5.200) (5.953,5.250) (5.942,5.300) (5.931,5.350) (5.919,5.400) (5.906,5.450) (5.893,5.500) (5.880,5.550) (5.866,5.600) (5.851,5.650) (5.836,5.700) (5.820,5.750) (5.804,5.800) (5.787,5.850) (5.769,5.900) (5.751,5.950) (5.732,6.000) (5.713,6.050) (5.693,6.100) (5.673,6.150) (5.652,6.200) (5.630,6.250) (5.608,6.300) (5.585,6.350) (5.561,6.400) (5.537,6.450) (5.512,6.500) (5.487,6.550) (5.461,6.600) (5.434,6.650) (5.406,6.700) (5.378,6.750) (5.350,6.799) (5.320,6.850) (5.289,6.900) (5.258,6.950) (5.226,7.000) (5.194,7.050) (5.160,7.100) (5.126,7.150) (5.091,7.200) (5.056,7.250) (5.025,7.292) (4.981,7.350) (4.950,7.391) (4.904,7.450) (4.864,7.500) (4.825,7.547) (4.781,7.600) (4.738,7.650) (4.700,7.693) (4.650,7.749) (4.603,7.800) (4.555,7.850) (4.507,7.900) (4.457,7.950) (4.407,8.000) (4.355,8.050) (4.302,8.100) (4.250,8.147) (4.200,8.192) (4.162,8.225) (4.104,8.275) (4.050,8.320) (4.000,8.361) (3.952,8.400) (3.900,8.441) (3.855,8.475) (3.800,8.517) (3.755,8.550) (3.700,8.589) (3.650,8.625) (3.600,8.659) (3.550,8.693) (3.501,8.725) (3.450,8.758) (3.400,8.790) (3.350,8.821) (3.302,8.850) (3.250,8.881) (3.200,8.910) (3.150,8.939) (3.100,8.967) (3.050,8.994) (3.000,9.021) (2.950,9.047) (2.900,9.072) (2.850,9.097) (2.800,9.122) (2.750,9.146) (2.700,9.169) (2.650,9.192) (2.600,9.214) (2.550,9.236) (2.500,9.257) (2.450,9.278) (2.400,9.298) (2.350,9.318) (2.300,9.337) (2.250,9.356) (2.200,9.374) (2.150,9.392) (2.100,9.410) (2.050,9.427) (2.000,9.443) (1.950,9.459) (1.900,9.475) (1.850,9.490) (1.800,9.504) (1.750,9.519) (1.700,9.533) (1.650,9.546) (1.600,9.559) (1.550,9.571) (1.500,9.583) (1.450,9.595) (1.400,9.606) (1.350,9.617) (1.300,9.628) (1.250,9.638) (1.200,9.647) (1.150,9.656) (1.100,9.665) (1.050,9.674) (1.000,9.682) (0.950,9.689) (0.900,9.696) (0.850,9.703) (0.800,9.710) (0.750,9.716) (0.700,9.721) (0.650,9.726) (0.600,9.731) (0.550,9.736) (0.500,9.740) (0.450,9.743) (0.400,9.747) (0.350,9.750) (0.300,9.752) (0.250,9.754) (0.200,9.756) (0.150,9.757) (0.100,9.758) (0.050,9.759) (0.000,9.759) (-0.050,9.759) (-0.100,9.758) (-0.150,9.757) (-0.200,9.756) (-0.250,9.754) (-0.300,9.752) (-0.350,9.750) (-0.400,9.747) (-0.450,9.743) (-0.500,9.740) (-0.550,9.736) (-0.600,9.731) (-0.650,9.726) (-0.700,9.721) (-0.750,9.716) (-0.800,9.710) (-0.850,9.703) (-0.900,9.696) (-0.950,9.689) (-1.000,9.682) (-1.050,9.674) (-1.100,9.665) (-1.150,9.656) (-1.200,9.647) (-1.250,9.638) (-1.300,9.628) (-1.350,9.617) (-1.400,9.606) (-1.450,9.595) (-1.500,9.583) (-1.550,9.571) (-1.600,9.559) (-1.650,9.546) (-1.700,9.533) (-1.750,9.519) (-1.800,9.504) (-1.850,9.490) (-1.899,9.475) (-1.950,9.459) (-2.000,9.443) (-2.050,9.427) (-2.100,9.410) (-2.150,9.392) (-2.198,9.375) (-2.250,9.356) (-2.300,9.337) (-2.350,9.318) (-2.400,9.298) (-2.450,9.278) (-2.500,9.257) (-2.550,9.236) (-2.600,9.214) (-2.650,9.192) (-2.700,9.169) (-2.750,9.146) (-2.800,9.122) (-2.850,9.097) (-2.900,9.072) (-2.950,9.047) (-3.000,9.021) (-3.050,8.994) (-3.100,8.967) (-3.150,8.939) (-3.200,8.910) (-3.250,8.881) (-3.300,8.851) (-3.344,8.825) (-3.400,8.790) (-3.450,8.758) (-3.500,8.726) (-3.550,8.693) (-3.600,8.659) (-3.650,8.625) (-3.700,8.589) (-3.750,8.553) (-3.800,8.517) (-3.850,8.479) (-3.900,8.441) (-3.950,8.401) (-4.000,8.361) (-4.045,8.325) (-4.100,8.279) (-4.150,8.236) (-4.191,8.200) (-4.247,8.150) (-4.300,8.101) (-4.350,8.055) (-4.400,8.007) (-4.450,7.957) (-4.500,7.907) (-4.550,7.856) (-4.600,7.803) (-4.649,7.750) (-4.694,7.700) (-4.738,7.650) (-4.775,7.607) (-4.823,7.550) (-4.864,7.500) (-4.900,7.455) (-4.943,7.400) (-4.975,7.359) (-5.019,7.300) (-5.050,7.258) (-5.091,7.200) (-5.125,7.152) (-5.160,7.100) (-5.194,7.050) (-5.225,7.002) (-5.258,6.950) (-5.289,6.900) (-5.320,6.850) (-5.349,6.800) (-5.375,6.755) (-5.400,6.711) (-5.425,6.666) (-5.450,6.620) (-5.474,6.575) (-5.500,6.525) (-5.525,6.475) (-5.549,6.425) (-5.573,6.375) (-5.596,6.325) (-5.619,6.275) (-5.641,6.225) (-5.662,6.175) (-5.683,6.125) (-5.703,6.075) (-5.723,6.025) (-5.742,5.975) (-5.760,5.925) (-5.778,5.875) (-5.795,5.825) (-5.812,5.775) (-5.828,5.725) (-5.843,5.675) (-5.858,5.625) (-5.873,5.575) (-5.887,5.525) (-5.900,5.475) (-5.913,5.425) (-5.925,5.375) (-5.936,5.325) (-5.947,5.275) (-5.958,5.225) (-5.968,5.175) (-5.977,5.125) (-5.986,5.075) (-5.994,5.025) (-6.002,4.975) (-6.009,4.925) (-6.016,4.875) (-6.022,4.825) (-6.028,4.775) (-6.033,4.725) (-6.038,4.675) (-6.042,4.625) (-6.045,4.575) (-6.048,4.525) (-6.050,4.475) (-6.052,4.425) (-6.054,4.375) (-6.054,4.325) (-6.055,4.275) (-6.054,4.225) (-6.053,4.175) (-6.052,4.125) (-6.050,4.075) (-6.047,4.025) (-6.044,3.975) (-6.040,3.925) (-6.036,3.875) (-6.031,3.825) (-6.026,3.775) (-6.020,3.725) (-6.013,3.675) (-6.006,3.625) (-5.998,3.575) (-5.990,3.525) (-5.981,3.475) (-5.971,3.425) (-5.961,3.375) (-5.950,3.325) (-5.939,3.275) (-5.926,3.225) (-5.914,3.175) (-5.900,3.125) (-5.886,3.075) (-5.871,3.025) (-5.856,2.975) (-5.839,2.925) (-5.822,2.875) (-5.805,2.825) (-5.786,2.775) (-5.767,2.725) (-5.747,2.675) (-5.727,2.625) (-5.705,2.575) (-5.683,2.525) (-5.660,2.475) (-5.636,2.425) (-5.611,2.375) (-5.585,2.325) (-5.559,2.275) (-5.531,2.225) (-5.503,2.175) (-5.475,2.128) (-5.443,2.075) (-5.411,2.025) (-5.378,1.975) (-5.350,1.933) (-5.310,1.875) (-5.275,1.827) (-5.237,1.775) (-5.200,1.727) (-5.158,1.675) (-5.125,1.634) (-5.075,1.576) (-5.030,1.525) (-4.985,1.475) (-4.950,1.438) (-4.913,1.400) (-4.875,1.362) (-4.838,1.325) (-4.800,1.289) (-4.758,1.250) (-4.702,1.200) (-4.650,1.155) (-4.600,1.113) (-4.553,1.075) (-4.500,1.034) (-4.456,1.000) (-4.400,0.959) (-4.352,0.925) (-4.300,0.889) (-4.250,0.856) (-4.203,0.825) (-4.150,0.792) (-4.100,0.761) (-4.050,0.732) (-4.000,0.703) (-3.950,0.676) (-3.902,0.650) (-3.854,0.625) (-3.805,0.600) (-3.754,0.575) (-3.701,0.550) (-3.650,0.527) (-3.600,0.504) (-3.550,0.483) (-3.500,0.462) (-3.450,0.442) (-3.400,0.422) (-3.350,0.403) (-3.300,0.385) (-3.250,0.367) (-3.200,0.350) (-3.150,0.334) (-3.100,0.318) (-3.050,0.302) (-3.000,0.287) (-2.950,0.273) (-2.900,0.259) (-2.850,0.246) (-2.800,0.233) (-2.750,0.220) (-2.700,0.209) (-2.650,0.197) (-2.600,0.186) (-2.550,0.175) (-2.500,0.165) (-2.450,0.155) (-2.400,0.146) (-2.350,0.137) (-2.300,0.128) (-2.250,0.120) (-2.200,0.112) (-2.150,0.105) (-2.100,0.098) (-2.050,0.091) (-2.000,0.084) (-1.950,0.078) (-1.900,0.072) (-1.850,0.067) (-1.800,0.061) (-1.750,0.056) (-1.700,0.052) (-1.650,0.047) (-1.600,0.043) (-1.550,0.039) (-1.500,0.035) (-1.450,0.032) (-1.400,0.029) (-1.350,0.026) (-1.302,0.023)}
\EQ{(1.300,0.046) (1.350,0.052) (1.400,0.058) (1.450,0.064) (1.500,0.071) (1.550,0.079) (1.600,0.086) (1.650,0.095) (1.700,0.104) (1.750,0.113) (1.800,0.123) (1.850,0.134) (1.900,0.145) (1.950,0.157) (2.000,0.170) (2.050,0.183) (2.100,0.197) (2.150,0.212) (2.200,0.227) (2.250,0.243) (2.300,0.260) (2.350,0.278) (2.400,0.297) (2.450,0.317) (2.500,0.337) (2.550,0.359) (2.600,0.381) (2.650,0.405) (2.700,0.429) (2.750,0.455) (2.800,0.482) (2.850,0.510) (2.900,0.539) (2.950,0.570) (2.998,0.600) (3.050,0.635) (3.100,0.670) (3.141,0.700) (3.200,0.745) (3.250,0.785) (3.298,0.825) (3.350,0.871) (3.400,0.917) (3.450,0.966) (3.500,1.017) (3.550,1.071) (3.598,1.125) (3.640,1.175) (3.675,1.219) (3.718,1.275) (3.750,1.319) (3.789,1.375) (3.822,1.425) (3.850,1.470) (3.883,1.525) (3.912,1.575) (3.939,1.625) (3.965,1.675) (3.990,1.725) (4.014,1.775) (4.036,1.825) (4.057,1.875) (4.077,1.925) (4.096,1.975) (4.114,2.025) (4.131,2.075) (4.147,2.125) (4.162,2.175) (4.176,2.225) (4.189,2.275) (4.201,2.325) (4.213,2.375) (4.223,2.425) (4.232,2.475) (4.241,2.525) (4.249,2.575) (4.256,2.625) (4.262,2.675) (4.267,2.725) (4.271,2.775) (4.275,2.825) (4.278,2.875) (4.280,2.925) (4.281,2.975) (4.281,3.025) (4.281,3.075) (4.280,3.125) (4.278,3.175) (4.275,3.225) (4.272,3.275) (4.268,3.325) (4.263,3.375) (4.257,3.425) (4.250,3.475) (4.243,3.525) (4.235,3.575) (4.226,3.625) (4.217,3.675) (4.206,3.725) (4.195,3.775) (4.184,3.825) (4.171,3.875) (4.157,3.925) (4.143,3.975) (4.128,4.025) (4.112,4.075) (4.096,4.125) (4.078,4.175) (4.060,4.225) (4.041,4.275) (4.021,4.325) (4.000,4.375) (3.979,4.425) (3.956,4.475) (3.933,4.525) (3.908,4.575) (3.883,4.625) (3.857,4.675) (3.830,4.725) (3.802,4.775) (3.775,4.821) (3.743,4.875) (3.711,4.925) (3.679,4.975) (3.650,5.019) (3.611,5.075) (3.576,5.125) (3.539,5.175) (3.501,5.225) (3.462,5.275) (3.425,5.321) (3.380,5.375) (3.336,5.425) (3.300,5.466) (3.250,5.520) (3.200,5.573) (3.150,5.623) (3.100,5.672) (3.050,5.720) (3.000,5.766) (2.961,5.800) (2.904,5.850) (2.850,5.895) (2.800,5.935) (2.750,5.974) (2.700,6.013) (2.650,6.049) (2.600,6.085) (2.550,6.120) (2.505,6.150) (2.450,6.186) (2.400,6.218) (2.350,6.249) (2.306,6.275) (2.250,6.308) (2.200,6.336) (2.150,6.363) (2.100,6.390) (2.050,6.415) (2.000,6.440) (1.950,6.464) (1.900,6.488) (1.850,6.510) (1.800,6.532) (1.750,6.553) (1.700,6.574) (1.650,6.593) (1.600,6.612) (1.550,6.631) (1.500,6.648) (1.450,6.665) (1.400,6.682) (1.350,6.698) (1.300,6.713) (1.250,6.727) (1.200,6.741) (1.150,6.754) (1.100,6.767) (1.050,6.779) (1.000,6.791) (0.950,6.801) (0.900,6.812) (0.850,6.821) (0.800,6.831) (0.750,6.839) (0.700,6.847) (0.650,6.854) (0.600,6.861) (0.550,6.868) (0.500,6.873) (0.450,6.879) (0.400,6.883) (0.350,6.887) (0.300,6.891) (0.250,6.894) (0.200,6.896) (0.150,6.898) (0.100,6.900) (0.050,6.900) (0.000,6.901) (-0.050,6.900) (-0.100,6.900) (-0.150,6.898) (-0.200,6.896) (-0.250,6.894) (-0.300,6.891) (-0.350,6.887) (-0.400,6.883) (-0.450,6.879) (-0.500,6.873) (-0.550,6.868) (-0.600,6.861) (-0.650,6.854) (-0.700,6.847) (-0.750,6.839) (-0.800,6.831) (-0.850,6.821) (-0.900,6.812) (-0.950,6.801) (-1.000,6.791) (-1.050,6.779) (-1.100,6.767) (-1.150,6.754) (-1.200,6.741) (-1.250,6.727) (-1.300,6.713) (-1.350,6.698) (-1.400,6.682) (-1.450,6.665) (-1.500,6.648) (-1.550,6.631) (-1.600,6.612) (-1.650,6.593) (-1.697,6.575) (-1.750,6.553) (-1.800,6.532) (-1.850,6.510) (-1.900,6.488) (-1.950,6.464) (-2.000,6.440) (-2.050,6.415) (-2.100,6.390) (-2.150,6.363) (-2.200,6.336) (-2.250,6.308) (-2.300,6.279) (-2.348,6.250) (-2.400,6.218) (-2.450,6.186) (-2.500,6.154) (-2.543,6.125) (-2.600,6.085) (-2.649,6.050) (-2.700,6.013) (-2.749,5.975) (-2.800,5.935) (-2.844,5.900) (-2.900,5.853) (-2.950,5.810) (-2.990,5.775) (-3.044,5.725) (-3.097,5.675) (-3.148,5.625) (-3.198,5.575) (-3.246,5.525) (-3.292,5.475) (-3.336,5.425) (-3.375,5.380) (-3.421,5.325) (-3.462,5.275) (-3.500,5.227) (-3.539,5.175) (-3.575,5.126) (-3.611,5.075) (-3.646,5.025) (-3.675,4.981) (-3.711,4.925) (-3.743,4.875) (-3.773,4.825) (-3.800,4.778) (-3.825,4.734) (-3.850,4.688) (-3.875,4.641) (-3.900,4.592) (-3.925,4.541) (-3.950,4.488) (-3.975,4.433) (-4.000,4.376) (-4.021,4.325) (-4.041,4.275) (-4.060,4.225) (-4.078,4.175) (-4.096,4.125) (-4.112,4.075) (-4.128,4.025) (-4.143,3.975) (-4.157,3.925) (-4.171,3.875) (-4.184,3.825) (-4.195,3.775) (-4.206,3.725) (-4.217,3.675) (-4.226,3.625) (-4.235,3.575) (-4.243,3.525) (-4.250,3.475) (-4.257,3.425) (-4.263,3.375) (-4.268,3.325) (-4.272,3.275) (-4.275,3.225) (-4.278,3.175) (-4.280,3.125) (-4.281,3.075) (-4.281,3.025) (-4.281,2.975) (-4.280,2.925) (-4.278,2.875) (-4.275,2.825) (-4.271,2.775) (-4.267,2.725) (-4.262,2.675) (-4.256,2.625) (-4.249,2.575) (-4.241,2.525) (-4.232,2.475) (-4.223,2.425) (-4.213,2.375) (-4.201,2.325) (-4.189,2.275) (-4.176,2.225) (-4.162,2.175) (-4.147,2.125) (-4.131,2.075) (-4.114,2.025) (-4.096,1.975) (-4.077,1.925) (-4.057,1.875) (-4.036,1.825) (-4.014,1.775) (-3.990,1.725) (-3.965,1.675) (-3.939,1.625) (-3.912,1.575) (-3.883,1.525) (-3.853,1.475) (-3.825,1.430) (-3.789,1.375) (-3.754,1.325) (-3.718,1.275) (-3.680,1.225) (-3.640,1.175) (-3.600,1.127) (-3.554,1.075) (-3.508,1.025) (-3.459,0.975) (-3.408,0.925) (-3.355,0.875) (-3.300,0.827) (-3.250,0.785) (-3.207,0.750) (-3.150,0.707) (-3.107,0.675) (-3.050,0.635) (-3.000,0.602) (-2.950,0.570) (-2.900,0.539) (-2.850,0.510) (-2.800,0.482) (-2.750,0.455) (-2.700,0.429) (-2.650,0.405) (-2.600,0.381) (-2.550,0.359) (-2.500,0.337) (-2.450,0.317) (-2.400,0.297) (-2.350,0.278) (-2.300,0.260) (-2.250,0.243) (-2.200,0.227) (-2.150,0.212) (-2.100,0.197) (-2.050,0.183) (-2.000,0.170) (-1.950,0.157) (-1.900,0.145) (-1.850,0.134) (-1.800,0.123) (-1.750,0.113) (-1.700,0.104) (-1.650,0.095) (-1.600,0.086) (-1.550,0.079) (-1.500,0.071) (-1.450,0.064) (-1.400,0.058) (-1.350,0.052) (-1.300,0.046)}
\EQ{(1.300,0.070) (1.350,0.078) (1.400,0.087) (1.450,0.097) (1.500,0.107) (1.550,0.118) (1.600,0.130) (1.650,0.143) (1.700,0.157) (1.750,0.171) (1.800,0.187) (1.850,0.203) (1.900,0.220) (1.950,0.239) (2.000,0.258) (2.050,0.279) (2.098,0.300) (2.150,0.324) (2.200,0.348) (2.250,0.374) (2.298,0.400) (2.342,0.425) (2.400,0.460) (2.450,0.492) (2.500,0.525) (2.550,0.561) (2.600,0.598) (2.650,0.638) (2.694,0.675) (2.750,0.724) (2.800,0.772) (2.850,0.822) (2.899,0.875) (2.943,0.925) (2.985,0.975) (3.024,1.025) (3.060,1.075) (3.095,1.125) (3.125,1.171) (3.159,1.225) (3.188,1.275) (3.216,1.325) (3.241,1.375) (3.266,1.425) (3.289,1.475) (3.310,1.525) (3.330,1.575) (3.349,1.625) (3.367,1.675) (3.383,1.725) (3.398,1.775) (3.412,1.825) (3.425,1.875) (3.437,1.925) (3.447,1.975) (3.456,2.025) (3.465,2.075) (3.472,2.125) (3.479,2.175) (3.484,2.225) (3.488,2.275) (3.491,2.325) (3.494,2.375) (3.495,2.425) (3.496,2.475) (3.495,2.525) (3.494,2.575) (3.491,2.625) (3.488,2.675) (3.483,2.725) (3.478,2.775) (3.472,2.825) (3.465,2.875) (3.457,2.925) (3.448,2.975) (3.438,3.025) (3.427,3.075) (3.415,3.125) (3.403,3.175) (3.389,3.225) (3.374,3.275) (3.359,3.325) (3.342,3.375) (3.325,3.423) (3.306,3.475) (3.286,3.525) (3.265,3.575) (3.243,3.625) (3.220,3.675) (3.196,3.725) (3.171,3.775) (3.145,3.825) (3.118,3.875) (3.089,3.925) (3.059,3.975) (3.028,4.025) (3.000,4.068) (2.962,4.125) (2.927,4.175) (2.890,4.225) (2.852,4.275) (2.812,4.325) (2.775,4.370) (2.728,4.425) (2.683,4.475) (2.637,4.525) (2.600,4.563) (2.563,4.600) (2.525,4.637) (2.485,4.675) (2.430,4.725) (2.375,4.773) (2.325,4.815) (2.281,4.850) (2.225,4.894) (2.183,4.925) (2.125,4.967) (2.078,5.000) (2.025,5.035) (1.975,5.067) (1.925,5.098) (1.881,5.125) (1.825,5.157) (1.775,5.185) (1.725,5.212) (1.675,5.238) (1.625,5.262) (1.575,5.286) (1.525,5.309) (1.475,5.331) (1.425,5.352) (1.375,5.373) (1.325,5.392) (1.275,5.411) (1.225,5.428) (1.175,5.445) (1.125,5.462) (1.075,5.477) (1.025,5.492) (0.975,5.505) (0.925,5.519) (0.875,5.531) (0.825,5.543) (0.775,5.553) (0.725,5.564) (0.675,5.573) (0.625,5.582) (0.575,5.590) (0.525,5.597) (0.475,5.604) (0.425,5.610) (0.375,5.616) (0.325,5.620) (0.275,5.624) (0.225,5.628) (0.175,5.630) (0.125,5.632) (0.075,5.634) (0.025,5.634) (-0.050,5.634) (-0.100,5.633) (-0.150,5.631) (-0.200,5.629) (-0.250,5.626) (-0.300,5.622) (-0.350,5.618) (-0.400,5.613) (-0.450,5.607) (-0.500,5.601) (-0.550,5.594) (-0.600,5.586) (-0.650,5.578) (-0.700,5.569) (-0.750,5.559) (-0.800,5.548) (-0.850,5.537) (-0.899,5.525) (-0.950,5.512) (-1.000,5.499) (-1.050,5.484) (-1.100,5.469) (-1.150,5.453) (-1.200,5.437) (-1.250,5.420) (-1.300,5.401) (-1.350,5.382) (-1.400,5.363) (-1.450,5.342) (-1.500,5.320) (-1.550,5.298) (-1.599,5.275) (-1.650,5.250) (-1.700,5.225) (-1.747,5.200) (-1.793,5.175) (-1.838,5.150) (-1.900,5.114) (-1.950,5.083) (-2.000,5.051) (-2.050,5.019) (-2.100,4.985) (-2.149,4.950) (-2.200,4.913) (-2.249,4.875) (-2.300,4.835) (-2.343,4.800) (-2.400,4.751) (-2.450,4.707) (-2.500,4.661) (-2.538,4.625) (-2.575,4.588) (-2.613,4.550) (-2.650,4.511) (-2.700,4.457) (-2.750,4.400) (-2.792,4.350) (-2.825,4.309) (-2.871,4.250) (-2.908,4.200) (-2.944,4.150) (-2.975,4.106) (-3.012,4.050) (-3.044,4.000) (-3.074,3.950) (-3.100,3.906) (-3.125,3.862) (-3.150,3.816) (-3.175,3.768) (-3.200,3.718) (-3.225,3.665) (-3.250,3.610) (-3.275,3.551) (-3.296,3.500) (-3.315,3.450) (-3.333,3.400) (-3.350,3.351) (-3.366,3.300) (-3.382,3.250) (-3.396,3.200) (-3.409,3.150) (-3.421,3.100) (-3.433,3.050) (-3.443,3.000) (-3.453,2.950) (-3.461,2.900) (-3.469,2.850) (-3.475,2.800) (-3.481,2.750) (-3.486,2.700) (-3.490,2.650) (-3.492,2.600) (-3.494,2.550) (-3.495,2.500) (-3.496,2.450) (-3.495,2.400) (-3.493,2.350) (-3.490,2.300) (-3.486,2.250) (-3.481,2.200) (-3.475,2.150) (-3.469,2.100) (-3.461,2.050) (-3.452,2.000) (-3.442,1.950) (-3.431,1.900) (-3.419,1.850) (-3.405,1.800) (-3.391,1.750) (-3.375,1.700) (-3.358,1.650) (-3.340,1.600) (-3.321,1.550) (-3.300,1.501) (-3.277,1.450) (-3.254,1.400) (-3.229,1.350) (-3.202,1.300) (-3.175,1.252) (-3.144,1.200) (-3.112,1.150) (-3.078,1.100) (-3.042,1.050) (-3.004,1.000) (-2.964,0.950) (-2.925,0.904) (-2.876,0.850) (-2.828,0.800) (-2.777,0.750) (-2.725,0.702) (-2.675,0.659) (-2.634,0.625) (-2.575,0.579) (-2.525,0.543) (-2.475,0.508) (-2.425,0.475) (-2.375,0.444) (-2.325,0.415) (-2.275,0.387) (-2.225,0.361) (-2.175,0.336) (-2.125,0.312) (-2.075,0.290) (-2.025,0.269) (-1.975,0.248) (-1.925,0.230) (-1.875,0.212) (-1.825,0.195) (-1.775,0.179) (-1.725,0.164) (-1.676,0.150) (-1.625,0.137) (-1.575,0.124) (-1.525,0.113) (-1.475,0.102) (-1.425,0.092) (-1.375,0.082) (-1.325,0.074) (-1.300,0.070)}
\EQ{(1.300,0.093) (1.350,0.104) (1.400,0.116) (1.450,0.130) (1.500,0.144) (1.550,0.159) (1.600,0.175) (1.650,0.193) (1.700,0.211) (1.750,0.231) (1.800,0.252) (1.850,0.275) (1.900,0.299) (1.950,0.324) (1.997,0.350) (2.041,0.375) (2.100,0.412) (2.150,0.444) (2.194,0.475) (2.250,0.517) (2.292,0.550) (2.350,0.599) (2.400,0.645) (2.450,0.693) (2.500,0.746) (2.548,0.800) (2.589,0.850) (2.625,0.896) (2.664,0.950) (2.697,1.000) (2.725,1.044) (2.759,1.100) (2.786,1.150) (2.812,1.200) (2.836,1.250) (2.858,1.300) (2.879,1.350) (2.898,1.400) (2.916,1.450) (2.932,1.500) (2.947,1.550) (2.960,1.600) (2.972,1.650) (2.983,1.700) (2.993,1.750) (3.001,1.800) (3.008,1.850) (3.014,1.900) (3.019,1.950) (3.023,2.000) (3.025,2.050) (3.027,2.100) (3.027,2.150) (3.026,2.200) (3.025,2.250) (3.022,2.300) (3.018,2.350) (3.013,2.400) (3.006,2.450) (2.999,2.500) (2.991,2.550) (2.981,2.600) (2.971,2.650) (2.959,2.700) (2.947,2.750) (2.933,2.800) (2.918,2.850) (2.902,2.900) (2.885,2.950) (2.866,3.000) (2.847,3.050) (2.826,3.100) (2.804,3.150) (2.781,3.200) (2.756,3.250) (2.730,3.300) (2.703,3.350) (2.675,3.399) (2.645,3.450) (2.613,3.500) (2.580,3.550) (2.550,3.594) (2.509,3.650) (2.475,3.696) (2.432,3.750) (2.400,3.789) (2.350,3.846) (2.301,3.900) (2.253,3.950) (2.203,4.000) (2.151,4.050) (2.100,4.096) (2.050,4.139) (2.007,4.175) (1.950,4.220) (1.900,4.258) (1.850,4.295) (1.807,4.325) (1.750,4.363) (1.700,4.395) (1.651,4.425) (1.600,4.455) (1.550,4.483) (1.500,4.510) (1.450,4.536) (1.400,4.561) (1.350,4.585) (1.300,4.607) (1.250,4.629) (1.200,4.649) (1.150,4.669) (1.100,4.687) (1.050,4.705) (1.000,4.722) (0.950,4.737) (0.900,4.752) (0.850,4.766) (0.800,4.779) (0.750,4.792) (0.700,4.803) (0.650,4.814) (0.600,4.824) (0.550,4.833) (0.500,4.841) (0.450,4.848) (0.400,4.855) (0.350,4.861) (0.300,4.866) (0.250,4.870) (0.200,4.873) (0.150,4.876) (0.100,4.878) (0.050,4.879) (0.000,4.880) (-0.050,4.879) (-0.100,4.878) (-0.150,4.876) (-0.200,4.873) (-0.250,4.870) (-0.300,4.866) (-0.350,4.861) (-0.400,4.855) (-0.450,4.848) (-0.500,4.841) (-0.550,4.833) (-0.600,4.824) (-0.650,4.814) (-0.700,4.803) (-0.750,4.792) (-0.800,4.779) (-0.850,4.766) (-0.900,4.752) (-0.950,4.737) (-1.000,4.722) (-1.050,4.705) (-1.100,4.687) (-1.150,4.669) (-1.198,4.650) (-1.250,4.629) (-1.300,4.607) (-1.350,4.585) (-1.400,4.561) (-1.450,4.536) (-1.500,4.510) (-1.550,4.483) (-1.600,4.455) (-1.650,4.426) (-1.700,4.395) (-1.750,4.363) (-1.800,4.330) (-1.843,4.300) (-1.900,4.258) (-1.944,4.225) (-2.000,4.181) (-2.050,4.139) (-2.095,4.100) (-2.150,4.051) (-2.200,4.003) (-2.250,3.954) (-2.300,3.901) (-2.347,3.850) (-2.390,3.800) (-2.425,3.758) (-2.472,3.700) (-2.509,3.650) (-2.546,3.600) (-2.575,3.558) (-2.613,3.500) (-2.645,3.450) (-2.675,3.400) (-2.700,3.356) (-2.725,3.310) (-2.750,3.262) (-2.775,3.212) (-2.800,3.159) (-2.825,3.102) (-2.847,3.050) (-2.866,3.000) (-2.885,2.950) (-2.902,2.900) (-2.918,2.850) (-2.933,2.800) (-2.947,2.750) (-2.959,2.700) (-2.971,2.650) (-2.981,2.600) (-2.991,2.550) (-2.999,2.500) (-3.006,2.450) (-3.013,2.400) (-3.018,2.350) (-3.022,2.300) (-3.025,2.250) (-3.026,2.200) (-3.027,2.150) (-3.027,2.100) (-3.025,2.050) (-3.023,2.000) (-3.019,1.950) (-3.014,1.900) (-3.008,1.850) (-3.001,1.800) (-2.993,1.750) (-2.983,1.700) (-2.972,1.650) (-2.960,1.600) (-2.947,1.550) (-2.932,1.500) (-2.916,1.450) (-2.898,1.400) (-2.879,1.350) (-2.858,1.300) (-2.836,1.250) (-2.812,1.200) (-2.786,1.150) (-2.759,1.100) (-2.729,1.050) (-2.700,1.004) (-2.664,0.950) (-2.628,0.900) (-2.589,0.850) (-2.550,0.802) (-2.504,0.750) (-2.457,0.700) (-2.406,0.650) (-2.351,0.600) (-2.300,0.557) (-2.261,0.525) (-2.200,0.479) (-2.158,0.450) (-2.100,0.412) (-2.050,0.381) (-2.000,0.352) (-1.951,0.325) (-1.902,0.300) (-1.850,0.275) (-1.800,0.252) (-1.750,0.231) (-1.700,0.211) (-1.650,0.193) (-1.600,0.175) (-1.550,0.159) (-1.500,0.144) (-1.450,0.130) (-1.400,0.116) (-1.350,0.104) (-1.300,0.093)}
\EQ{(1.300,0.117) (1.350,0.131) (1.400,0.146) (1.450,0.163) (1.500,0.181) (1.549,0.200) (1.600,0.221) (1.650,0.244) (1.700,0.268) (1.750,0.293) (1.800,0.321) (1.849,0.350) (1.900,0.382) (1.950,0.416) (1.997,0.450) (2.050,0.492) (2.090,0.525) (2.145,0.575) (2.196,0.625) (2.243,0.675) (2.287,0.725) (2.325,0.772) (2.365,0.825) (2.400,0.875) (2.432,0.925) (2.462,0.975) (2.490,1.025) (2.516,1.075) (2.540,1.125) (2.562,1.175) (2.582,1.225) (2.600,1.275) (2.617,1.325) (2.632,1.375) (2.646,1.425) (2.658,1.475) (2.669,1.525) (2.679,1.575) (2.687,1.625) (2.693,1.675) (2.699,1.725) (2.703,1.775) (2.706,1.825) (2.707,1.875) (2.708,1.925) (2.707,1.975) (2.705,2.025) (2.701,2.075) (2.697,2.125) (2.691,2.175) (2.684,2.225) (2.676,2.275) (2.667,2.325) (2.656,2.375) (2.644,2.425) (2.631,2.475) (2.617,2.525) (2.602,2.575) (2.585,2.625) (2.567,2.675) (2.547,2.725) (2.527,2.775) (2.504,2.825) (2.481,2.875) (2.456,2.925) (2.430,2.975) (2.402,3.025) (2.375,3.070) (2.341,3.125) (2.308,3.175) (2.275,3.223) (2.237,3.275) (2.200,3.323) (2.158,3.375) (2.125,3.414) (2.075,3.470) (2.025,3.523) (1.975,3.574) (1.925,3.621) (1.875,3.667) (1.836,3.700) (1.776,3.750) (1.725,3.790) (1.677,3.825) (1.625,3.862) (1.575,3.896) (1.530,3.925) (1.475,3.959) (1.425,3.988) (1.375,4.016) (1.325,4.043) (1.275,4.068) (1.225,4.092) (1.175,4.115) (1.125,4.137) (1.075,4.158) (1.025,4.177) (0.975,4.196) (0.925,4.213) (0.875,4.229) (0.825,4.245) (0.775,4.259) (0.725,4.272) (0.675,4.285) (0.625,4.296) (0.575,4.307) (0.525,4.317) (0.475,4.325) (0.425,4.333) (0.375,4.340) (0.325,4.346) (0.275,4.351) (0.225,4.356) (0.175,4.359) (0.125,4.362) (0.075,4.363) (0.025,4.364) (-0.050,4.364) (-0.100,4.363) (-0.150,4.361) (-0.200,4.358) (-0.250,4.354) (-0.300,4.349) (-0.350,4.343) (-0.400,4.337) (-0.450,4.329) (-0.500,4.321) (-0.550,4.312) (-0.600,4.302) (-0.650,4.291) (-0.700,4.279) (-0.750,4.266) (-0.800,4.252) (-0.850,4.237) (-0.900,4.221) (-0.950,4.204) (-1.000,4.187) (-1.050,4.168) (-1.100,4.147) (-1.150,4.126) (-1.200,4.104) (-1.250,4.080) (-1.300,4.056) (-1.350,4.030) (-1.400,4.002) (-1.448,3.975) (-1.500,3.944) (-1.550,3.912) (-1.600,3.879) (-1.642,3.850) (-1.700,3.808) (-1.744,3.775) (-1.800,3.730) (-1.850,3.688) (-1.894,3.650) (-1.948,3.600) (-1.999,3.550) (-2.047,3.500) (-2.093,3.450) (-2.137,3.400) (-2.175,3.354) (-2.218,3.300) (-2.250,3.257) (-2.291,3.200) (-2.325,3.150) (-2.357,3.100) (-2.387,3.050) (-2.416,3.000) (-2.443,2.950) (-2.469,2.900) (-2.493,2.850) (-2.516,2.800) (-2.537,2.750) (-2.557,2.700) (-2.575,2.653) (-2.593,2.600) (-2.610,2.550) (-2.624,2.500) (-2.638,2.450) (-2.650,2.400) (-2.662,2.350) (-2.672,2.300) (-2.680,2.250) (-2.688,2.200) (-2.694,2.150) (-2.699,2.100) (-2.703,2.050) (-2.706,2.000) (-2.707,1.950) (-2.708,1.900) (-2.707,1.850) (-2.704,1.800) (-2.701,1.750) (-2.696,1.700) (-2.690,1.650) (-2.683,1.600) (-2.674,1.550) (-2.664,1.500) (-2.652,1.450) (-2.639,1.400) (-2.625,1.351) (-2.609,1.300) (-2.591,1.250) (-2.572,1.200) (-2.551,1.150) (-2.528,1.100) (-2.503,1.050) (-2.476,1.000) (-2.450,0.954) (-2.416,0.900) (-2.383,0.850) (-2.350,0.805) (-2.308,0.750) (-2.275,0.711) (-2.225,0.655) (-2.175,0.604) (-2.125,0.557) (-2.075,0.513) (-2.029,0.475) (-1.975,0.434) (-1.927,0.400) (-1.875,0.366) (-1.825,0.335) (-1.775,0.307) (-1.725,0.280) (-1.675,0.255) (-1.625,0.232) (-1.575,0.211) (-1.525,0.191) (-1.475,0.172) (-1.425,0.155) (-1.375,0.139) (-1.325,0.124) (-1.300,0.117)}
\EQ{(1.300,0.141) (1.350,0.158) (1.400,0.177) (1.450,0.197) (1.500,0.220) (1.550,0.243) (1.600,0.269) (1.650,0.297) (1.697,0.325) (1.750,0.359) (1.800,0.394) (1.841,0.425) (1.900,0.473) (1.950,0.517) (2.000,0.566) (2.050,0.619) (2.099,0.675) (2.138,0.725) (2.174,0.775) (2.208,0.825) (2.239,0.875) (2.267,0.925) (2.293,0.975) (2.317,1.025) (2.339,1.075) (2.359,1.125) (2.377,1.175) (2.394,1.225) (2.409,1.275) (2.422,1.325) (2.433,1.375) (2.443,1.425) (2.451,1.475) (2.458,1.525) (2.464,1.575) (2.468,1.625) (2.470,1.675) (2.472,1.725) (2.472,1.775) (2.470,1.825) (2.467,1.875) (2.463,1.925) (2.458,1.975) (2.451,2.025) (2.443,2.075) (2.434,2.125) (2.423,2.175) (2.411,2.225) (2.398,2.275) (2.383,2.325) (2.367,2.375) (2.350,2.423) (2.330,2.475) (2.310,2.525) (2.288,2.575) (2.265,2.625) (2.240,2.675) (2.213,2.725) (2.184,2.775) (2.154,2.825) (2.125,2.871) (2.089,2.925) (2.053,2.975) (2.015,3.025) (1.975,3.075) (1.933,3.125) (1.888,3.175) (1.850,3.215) (1.800,3.265) (1.763,3.300) (1.725,3.335) (1.679,3.375) (1.625,3.420) (1.575,3.459) (1.525,3.496) (1.475,3.531) (1.425,3.565) (1.375,3.597) (1.328,3.625) (1.285,3.650) (1.239,3.675) (1.191,3.700) (1.141,3.725) (1.087,3.750) (1.031,3.775) (0.975,3.798) (0.925,3.817) (0.875,3.835) (0.825,3.852) (0.775,3.868) (0.725,3.883) (0.675,3.897) (0.625,3.909) (0.575,3.921) (0.525,3.932) (0.475,3.941) (0.425,3.950) (0.375,3.958) (0.325,3.964) (0.275,3.970) (0.225,3.975) (0.175,3.978) (0.125,3.981) (0.075,3.983) (0.025,3.984) (-0.050,3.984) (-0.100,3.982) (-0.150,3.980) (-0.200,3.977) (-0.250,3.972) (-0.300,3.967) (-0.350,3.961) (-0.400,3.954) (-0.450,3.946) (-0.500,3.937) (-0.550,3.926) (-0.600,3.915) (-0.650,3.903) (-0.700,3.890) (-0.750,3.876) (-0.800,3.860) (-0.850,3.844) (-0.900,3.826) (-0.950,3.808) (-1.000,3.788) (-1.050,3.767) (-1.100,3.744) (-1.150,3.721) (-1.200,3.696) (-1.250,3.669) (-1.300,3.641) (-1.350,3.612) (-1.400,3.581) (-1.448,3.550) (-1.500,3.514) (-1.550,3.478) (-1.600,3.440) (-1.650,3.400) (-1.700,3.357) (-1.750,3.312) (-1.790,3.275) (-1.840,3.225) (-1.875,3.189) (-1.910,3.150) (-1.950,3.105) (-1.995,3.050) (-2.034,3.000) (-2.071,2.950) (-2.100,2.908) (-2.139,2.850) (-2.170,2.800) (-2.199,2.750) (-2.225,2.703) (-2.250,2.654) (-2.275,2.603) (-2.299,2.550) (-2.320,2.500) (-2.340,2.450) (-2.358,2.400) (-2.375,2.351) (-2.391,2.300) (-2.405,2.250) (-2.417,2.200) (-2.429,2.150) (-2.439,2.100) (-2.447,2.050) (-2.455,2.000) (-2.461,1.950) (-2.466,1.900) (-2.469,1.850) (-2.471,1.800) (-2.472,1.750) (-2.471,1.700) (-2.469,1.650) (-2.466,1.600) (-2.461,1.550) (-2.455,1.500) (-2.447,1.450) (-2.438,1.400) (-2.427,1.350) (-2.415,1.300) (-2.401,1.250) (-2.386,1.200) (-2.369,1.150) (-2.350,1.101) (-2.329,1.050) (-2.306,1.000) (-2.281,0.950) (-2.253,0.900) (-2.225,0.852) (-2.192,0.800) (-2.157,0.750) (-2.125,0.708) (-2.078,0.650) (-2.033,0.600) (-1.984,0.550) (-1.931,0.500) (-1.875,0.452) (-1.825,0.413) (-1.775,0.376) (-1.725,0.343) (-1.675,0.312) (-1.625,0.283) (-1.575,0.256) (-1.525,0.231) (-1.475,0.208) (-1.425,0.187) (-1.375,0.167) (-1.327,0.150) (-1.300,0.141)}
\EQ{(1.300,0.165) (1.350,0.186) (1.400,0.208) (1.450,0.233) (1.500,0.259) (1.550,0.288) (1.600,0.319) (1.646,0.350) (1.700,0.390) (1.744,0.425) (1.800,0.474) (1.850,0.522) (1.900,0.575) (1.943,0.625) (1.975,0.665) (2.018,0.725) (2.050,0.773) (2.082,0.825) (2.109,0.875) (2.134,0.925) (2.157,0.975) (2.178,1.025) (2.197,1.075) (2.213,1.125) (2.228,1.175) (2.242,1.225) (2.253,1.275) (2.263,1.325) (2.271,1.375) (2.277,1.425) (2.282,1.475) (2.286,1.525) (2.288,1.575) (2.288,1.625) (2.288,1.675) (2.285,1.725) (2.281,1.775) (2.276,1.825) (2.269,1.875) (2.261,1.925) (2.252,1.975) (2.241,2.025) (2.228,2.075) (2.215,2.125) (2.199,2.175) (2.182,2.225) (2.164,2.275) (2.144,2.325) (2.122,2.375) (2.100,2.423) (2.075,2.473) (2.050,2.520) (2.025,2.565) (1.989,2.625) (1.956,2.675) (1.925,2.720) (1.885,2.775) (1.850,2.820) (1.805,2.875) (1.761,2.925) (1.725,2.964) (1.690,3.000) (1.650,3.039) (1.611,3.075) (1.554,3.125) (1.500,3.170) (1.450,3.209) (1.400,3.245) (1.358,3.275) (1.300,3.313) (1.250,3.344) (1.200,3.373) (1.152,3.400) (1.104,3.425) (1.053,3.450) (1.000,3.475) (0.950,3.497) (0.900,3.517) (0.850,3.536) (0.800,3.554) (0.750,3.571) (0.700,3.587) (0.650,3.601) (0.600,3.614) (0.550,3.626) (0.500,3.637) (0.450,3.647) (0.400,3.656) (0.350,3.664) (0.300,3.670) (0.250,3.676) (0.200,3.680) (0.150,3.684) (0.100,3.687) (0.050,3.688) (0.000,3.689) (-0.050,3.688) (-0.100,3.687) (-0.150,3.684) (-0.200,3.680) (-0.250,3.676) (-0.300,3.670) (-0.350,3.664) (-0.400,3.656) (-0.450,3.647) (-0.500,3.637) (-0.550,3.626) (-0.600,3.614) (-0.650,3.601) (-0.700,3.587) (-0.750,3.571) (-0.800,3.554) (-0.850,3.536) (-0.900,3.517) (-0.950,3.497) (-1.000,3.475) (-1.050,3.452) (-1.100,3.427) (-1.150,3.401) (-1.197,3.375) (-1.250,3.344) (-1.300,3.313) (-1.350,3.280) (-1.394,3.250) (-1.450,3.209) (-1.494,3.175) (-1.550,3.129) (-1.600,3.085) (-1.638,3.050) (-1.675,3.015) (-1.714,2.975) (-1.750,2.937) (-1.800,2.881) (-1.846,2.825) (-1.885,2.775) (-1.922,2.725) (-1.956,2.675) (-1.989,2.625) (-2.019,2.575) (-2.047,2.525) (-2.074,2.475) (-2.099,2.425) (-2.122,2.375) (-2.144,2.325) (-2.164,2.275) (-2.182,2.225) (-2.199,2.175) (-2.215,2.125) (-2.228,2.075) (-2.241,2.025) (-2.252,1.975) (-2.261,1.925) (-2.269,1.875) (-2.276,1.825) (-2.281,1.775) (-2.285,1.725) (-2.288,1.675) (-2.288,1.625) (-2.288,1.575) (-2.286,1.525) (-2.282,1.475) (-2.277,1.425) (-2.271,1.375) (-2.263,1.325) (-2.253,1.275) (-2.242,1.225) (-2.228,1.175) (-2.213,1.125) (-2.197,1.075) (-2.178,1.025) (-2.157,0.975) (-2.134,0.925) (-2.109,0.875) (-2.082,0.825) (-2.051,0.775) (-2.018,0.725) (-1.982,0.675) (-1.950,0.634) (-1.900,0.575) (-1.853,0.525) (-1.801,0.475) (-1.750,0.430) (-1.700,0.390) (-1.650,0.353) (-1.600,0.319) (-1.550,0.288) (-1.500,0.259) (-1.450,0.233) (-1.400,0.208) (-1.350,0.186) (-1.300,0.165)}
\EQ{(1.300,0.191) (1.350,0.215) (1.400,0.241) (1.450,0.270) (1.499,0.300) (1.550,0.335) (1.600,0.372) (1.650,0.414) (1.691,0.450) (1.742,0.500) (1.788,0.550) (1.825,0.594) (1.868,0.650) (1.900,0.696) (1.934,0.750) (1.963,0.800) (1.989,0.850) (2.013,0.900) (2.034,0.950) (2.053,1.000) (2.070,1.050) (2.085,1.100) (2.098,1.150) (2.109,1.200) (2.118,1.250) (2.126,1.300) (2.132,1.350) (2.137,1.400) (2.139,1.450) (2.141,1.500) (2.140,1.550) (2.138,1.600) (2.135,1.650) (2.130,1.700) (2.123,1.750) (2.115,1.800) (2.106,1.850) (2.095,1.900) (2.082,1.950) (2.068,2.000) (2.052,2.050) (2.035,2.100) (2.016,2.150) (1.995,2.200) (1.972,2.250) (1.950,2.296) (1.925,2.344) (1.900,2.389) (1.864,2.450) (1.831,2.500) (1.800,2.546) (1.760,2.600) (1.725,2.645) (1.679,2.700) (1.634,2.750) (1.600,2.786) (1.562,2.825) (1.525,2.860) (1.481,2.900) (1.425,2.947) (1.375,2.987) (1.325,3.025) (1.275,3.060) (1.225,3.093) (1.175,3.124) (1.132,3.150) (1.087,3.175) (1.040,3.200) (0.990,3.225) (0.936,3.250) (0.879,3.275) (0.825,3.297) (0.775,3.315) (0.725,3.333) (0.675,3.349) (0.625,3.364) (0.575,3.377) (0.525,3.390) (0.475,3.401) (0.425,3.411) (0.375,3.420) (0.325,3.427) (0.275,3.434) (0.225,3.439) (0.175,3.444) (0.125,3.447) (0.075,3.449) (0.025,3.450) (-0.037,3.450) (-0.100,3.448) (-0.150,3.445) (-0.200,3.442) (-0.250,3.437) (-0.300,3.431) (-0.350,3.424) (-0.400,3.415) (-0.450,3.406) (-0.500,3.395) (-0.550,3.384) (-0.600,3.371) (-0.650,3.356) (-0.700,3.341) (-0.748,3.325) (-0.800,3.306) (-0.850,3.287) (-0.900,3.266) (-0.950,3.244) (-1.000,3.220) (-1.050,3.195) (-1.100,3.168) (-1.150,3.139) (-1.200,3.109) (-1.250,3.077) (-1.300,3.043) (-1.350,3.006) (-1.391,2.975) (-1.450,2.927) (-1.500,2.883) (-1.550,2.836) (-1.587,2.800) (-1.625,2.760) (-1.675,2.705) (-1.721,2.650) (-1.760,2.600) (-1.797,2.550) (-1.831,2.500) (-1.864,2.450) (-1.894,2.400) (-1.922,2.350) (-1.948,2.300) (-1.972,2.250) (-1.995,2.200) (-2.016,2.150) (-2.035,2.100) (-2.052,2.050) (-2.068,2.000) (-2.082,1.950) (-2.095,1.900) (-2.106,1.850) (-2.115,1.800) (-2.123,1.750) (-2.130,1.700) (-2.135,1.650) (-2.138,1.600) (-2.140,1.550) (-2.141,1.500) (-2.139,1.450) (-2.137,1.400) (-2.132,1.350) (-2.126,1.300) (-2.118,1.250) (-2.109,1.200) (-2.098,1.150) (-2.085,1.100) (-2.070,1.050) (-2.053,1.000) (-2.034,0.950) (-2.013,0.900) (-1.989,0.850) (-1.963,0.800) (-1.934,0.750) (-1.903,0.700) (-1.868,0.650) (-1.830,0.600) (-1.788,0.550) (-1.750,0.508) (-1.700,0.459) (-1.650,0.414) (-1.603,0.375) (-1.550,0.335) (-1.500,0.301) (-1.450,0.270) (-1.400,0.241) (-1.350,0.215) (-1.300,0.191)}
\EQ{(1.300,0.216) (1.350,0.244) (1.400,0.274) (1.450,0.308) (1.500,0.345) (1.550,0.385) (1.595,0.425) (1.646,0.475) (1.692,0.525) (1.725,0.565) (1.770,0.625) (1.800,0.669) (1.834,0.725) (1.862,0.775) (1.887,0.825) (1.909,0.875) (1.929,0.925) (1.947,0.975) (1.962,1.025) (1.975,1.075) (1.987,1.125) (1.997,1.175) (2.004,1.225) (2.011,1.275) (2.015,1.325) (2.017,1.375) (2.018,1.425) (2.017,1.475) (2.015,1.525) (2.011,1.575) (2.005,1.625) (1.998,1.675) (1.989,1.725) (1.979,1.775) (1.967,1.825) (1.953,1.875) (1.937,1.925) (1.920,1.975) (1.901,2.025) (1.880,2.075) (1.858,2.125) (1.833,2.175) (1.807,2.225) (1.778,2.275) (1.750,2.321) (1.714,2.375) (1.679,2.425) (1.641,2.475) (1.601,2.525) (1.557,2.575) (1.510,2.625) (1.475,2.661) (1.434,2.700) (1.378,2.750) (1.325,2.794) (1.285,2.825) (1.225,2.869) (1.180,2.900) (1.125,2.935) (1.075,2.965) (1.025,2.993) (0.975,3.020) (0.925,3.045) (0.875,3.068) (0.825,3.089) (0.775,3.109) (0.725,3.128) (0.675,3.145) (0.625,3.161) (0.576,3.175) (0.525,3.188) (0.476,3.200) (0.425,3.211) (0.375,3.220) (0.325,3.229) (0.275,3.236) (0.225,3.241) (0.175,3.246) (0.125,3.249) (0.075,3.252) (0.025,3.253) (-0.050,3.252) (-0.100,3.251) (-0.150,3.248) (-0.200,3.244) (-0.250,3.239) (-0.300,3.232) (-0.350,3.225) (-0.400,3.216) (-0.450,3.206) (-0.500,3.195) (-0.550,3.182) (-0.600,3.168) (-0.650,3.153) (-0.700,3.137) (-0.750,3.119) (-0.798,3.100) (-0.850,3.079) (-0.900,3.056) (-0.950,3.032) (-1.000,3.007) (-1.050,2.980) (-1.100,2.950) (-1.150,2.919) (-1.200,2.886) (-1.250,2.851) (-1.300,2.814) (-1.348,2.775) (-1.400,2.731) (-1.450,2.685) (-1.486,2.650) (-1.525,2.610) (-1.575,2.555) (-1.621,2.500) (-1.661,2.450) (-1.697,2.400) (-1.731,2.350) (-1.763,2.300) (-1.793,2.250) (-1.820,2.200) (-1.846,2.150) (-1.869,2.100) (-1.891,2.050) (-1.911,2.000) (-1.929,1.950) (-1.945,1.900) (-1.960,1.850) (-1.973,1.800) (-1.984,1.750) (-1.994,1.700) (-2.002,1.650) (-2.008,1.600) (-2.013,1.550) (-2.017,1.500) (-2.018,1.450) (-2.018,1.400) (-2.016,1.350) (-2.013,1.300) (-2.008,1.250) (-2.001,1.200) (-1.992,1.150) (-1.982,1.100) (-1.969,1.050) (-1.955,1.000) (-1.938,0.950) (-1.919,0.900) (-1.898,0.850) (-1.875,0.801) (-1.850,0.753) (-1.825,0.710) (-1.787,0.650) (-1.752,0.600) (-1.713,0.550) (-1.675,0.506) (-1.625,0.454) (-1.575,0.407) (-1.525,0.364) (-1.475,0.326) (-1.425,0.291) (-1.375,0.259) (-1.325,0.230) (-1.300,0.216)}
\EQ{(1.300,0.243) (1.350,0.275) (1.400,0.310) (1.450,0.348) (1.500,0.391) (1.550,0.439) (1.600,0.493) (1.647,0.550) (1.684,0.600) (1.717,0.650) (1.747,0.700) (1.774,0.750) (1.798,0.800) (1.819,0.850) (1.838,0.900) (1.855,0.950) (1.869,1.000) (1.881,1.050) (1.892,1.100) (1.900,1.150) (1.906,1.200) (1.911,1.250) (1.914,1.300) (1.915,1.350) (1.914,1.400) (1.911,1.450) (1.907,1.500) (1.901,1.550) (1.894,1.600) (1.885,1.650) (1.873,1.700) (1.861,1.750) (1.846,1.800) (1.830,1.850) (1.812,1.900) (1.792,1.950) (1.770,2.000) (1.746,2.050) (1.720,2.100) (1.692,2.150) (1.661,2.200) (1.629,2.250) (1.593,2.300) (1.555,2.350) (1.515,2.400) (1.475,2.445) (1.425,2.498) (1.375,2.548) (1.325,2.593) (1.275,2.636) (1.226,2.675) (1.175,2.713) (1.125,2.747) (1.083,2.775) (1.025,2.810) (0.975,2.838) (0.925,2.865) (0.875,2.890) (0.825,2.913) (0.775,2.934) (0.725,2.954) (0.675,2.972) (0.625,2.989) (0.575,3.004) (0.525,3.018) (0.475,3.030) (0.425,3.042) (0.375,3.052) (0.325,3.060) (0.275,3.068) (0.225,3.074) (0.175,3.079) (0.125,3.082) (0.075,3.085) (0.025,3.086) (-0.050,3.086) (-0.100,3.084) (-0.150,3.081) (-0.200,3.076) (-0.250,3.071) (-0.300,3.064) (-0.350,3.056) (-0.400,3.047) (-0.450,3.036) (-0.500,3.024) (-0.550,3.011) (-0.600,2.996) (-0.650,2.980) (-0.700,2.963) (-0.750,2.944) (-0.800,2.923) (-0.850,2.901) (-0.900,2.877) (-0.950,2.852) (-0.999,2.825) (-1.050,2.795) (-1.100,2.764) (-1.150,2.730) (-1.192,2.700) (-1.250,2.656) (-1.300,2.615) (-1.345,2.575) (-1.398,2.525) (-1.448,2.475) (-1.493,2.425) (-1.535,2.375) (-1.575,2.325) (-1.611,2.275) (-1.645,2.225) (-1.675,2.178) (-1.706,2.125) (-1.733,2.075) (-1.758,2.025) (-1.781,1.975) (-1.802,1.925) (-1.821,1.875) (-1.838,1.825) (-1.854,1.775) (-1.867,1.725) (-1.879,1.675) (-1.889,1.625) (-1.898,1.575) (-1.905,1.525) (-1.910,1.475) (-1.913,1.425) (-1.915,1.375) (-1.914,1.325) (-1.913,1.275) (-1.909,1.225) (-1.903,1.175) (-1.896,1.125) (-1.887,1.075) (-1.875,1.025) (-1.862,0.975) (-1.847,0.925) (-1.829,0.875) (-1.809,0.825) (-1.786,0.775) (-1.761,0.725) (-1.733,0.675) (-1.701,0.625) (-1.666,0.575) (-1.627,0.525) (-1.584,0.475) (-1.536,0.425) (-1.482,0.375) (-1.425,0.328) (-1.375,0.292) (-1.325,0.258) (-1.300,0.243)}
\EQ{(1.288,0.262) (1.342,0.300) (1.400,0.347) (1.450,0.391) (1.500,0.442) (1.550,0.499) (1.590,0.550) (1.625,0.600) (1.656,0.650) (1.684,0.700) (1.708,0.750) (1.730,0.800) (1.750,0.850) (1.767,0.900) (1.781,0.950) (1.794,1.000) (1.804,1.050) (1.812,1.100) (1.818,1.150) (1.823,1.200) (1.825,1.250) (1.825,1.300) (1.824,1.350) (1.821,1.400) (1.816,1.450) (1.810,1.500) (1.801,1.550) (1.791,1.600) (1.779,1.650) (1.765,1.700) (1.750,1.748) (1.732,1.800) (1.712,1.850) (1.691,1.900) (1.667,1.950) (1.641,2.000) (1.613,2.050) (1.583,2.100) (1.550,2.149) (1.514,2.200) (1.476,2.250) (1.434,2.300) (1.400,2.339) (1.366,2.375) (1.325,2.416) (1.289,2.450) (1.250,2.484) (1.201,2.525) (1.150,2.565) (1.101,2.600) (1.050,2.634) (1.000,2.666) (0.950,2.695) (0.900,2.722) (0.850,2.748) (0.800,2.771) (0.750,2.793) (0.700,2.813) (0.650,2.831) (0.600,2.848) (0.550,2.864) (0.500,2.878) (0.450,2.890) (0.400,2.901) (0.350,2.911) (0.300,2.919) (0.250,2.926) (0.200,2.932) (0.150,2.937) (0.100,2.940) (0.050,2.942) (0.000,2.943) (-0.050,2.942) (-0.100,2.940) (-0.150,2.937) (-0.200,2.932) (-0.250,2.926) (-0.300,2.919) (-0.350,2.911) (-0.400,2.901) (-0.450,2.890) (-0.500,2.878) (-0.550,2.864) (-0.600,2.848) (-0.650,2.831) (-0.700,2.813) (-0.750,2.793) (-0.800,2.771) (-0.850,2.748) (-0.900,2.722) (-0.950,2.695) (-1.000,2.666) (-1.050,2.634) (-1.100,2.601) (-1.150,2.565) (-1.200,2.526) (-1.250,2.484) (-1.289,2.450) (-1.325,2.416) (-1.366,2.375) (-1.400,2.339) (-1.434,2.300) (-1.475,2.251) (-1.514,2.200) (-1.550,2.150) (-1.583,2.100) (-1.613,2.050) (-1.641,2.000) (-1.667,1.950) (-1.691,1.900) (-1.712,1.850) (-1.732,1.800) (-1.749,1.750) (-1.765,1.700) (-1.779,1.650) (-1.791,1.600) (-1.801,1.550) (-1.810,1.500) (-1.816,1.450) (-1.821,1.400) (-1.824,1.350) (-1.825,1.300) (-1.825,1.250) (-1.823,1.200) (-1.818,1.150) (-1.812,1.100) (-1.804,1.050) (-1.794,1.000) (-1.781,0.950) (-1.767,0.900) (-1.750,0.851) (-1.730,0.800) (-1.708,0.750) (-1.684,0.700) (-1.656,0.650) (-1.625,0.601) (-1.590,0.550) (-1.551,0.500) (-1.508,0.450) (-1.459,0.400) (-1.404,0.350) (-1.350,0.307) (-1.307,0.275) (-1.288,0.262)}
\EQ{(1.275,0.280) (1.325,0.319) (1.375,0.362) (1.414,0.400) (1.450,0.438) (1.500,0.497) (1.539,0.550) (1.572,0.600) (1.600,0.648) (1.625,0.696) (1.650,0.750) (1.670,0.800) (1.688,0.850) (1.703,0.900) (1.716,0.950) (1.726,1.000) (1.734,1.050) (1.741,1.100) (1.745,1.150) (1.747,1.200) (1.748,1.250) (1.746,1.300) (1.743,1.350) (1.738,1.400) (1.731,1.450) (1.722,1.500) (1.711,1.550) (1.698,1.600) (1.683,1.650) (1.667,1.700) (1.648,1.750) (1.627,1.800) (1.604,1.850) (1.579,1.900) (1.552,1.950) (1.525,1.995) (1.489,2.050) (1.454,2.100) (1.416,2.150) (1.375,2.200) (1.330,2.250) (1.282,2.300) (1.229,2.350) (1.175,2.397) (1.125,2.437) (1.075,2.475) (1.025,2.509) (0.975,2.542) (0.925,2.571) (0.875,2.599) (0.825,2.625) (0.775,2.649) (0.725,2.671) (0.675,2.691) (0.625,2.710) (0.575,2.727) (0.525,2.742) (0.475,2.756) (0.425,2.768) (0.375,2.779) (0.325,2.789) (0.275,2.797) (0.225,2.804) (0.175,2.809) (0.125,2.813) (0.075,2.816) (0.025,2.817) (-0.050,2.817) (-0.100,2.815) (-0.150,2.811) (-0.200,2.807) (-0.250,2.800) (-0.300,2.793) (-0.350,2.784) (-0.400,2.774) (-0.450,2.762) (-0.500,2.749) (-0.550,2.735) (-0.600,2.719) (-0.650,2.701) (-0.700,2.681) (-0.750,2.660) (-0.800,2.637) (-0.850,2.612) (-0.900,2.586) (-0.950,2.557) (-1.000,2.526) (-1.050,2.492) (-1.100,2.456) (-1.141,2.425) (-1.200,2.376) (-1.250,2.331) (-1.300,2.282) (-1.350,2.228) (-1.396,2.175) (-1.436,2.125) (-1.472,2.075) (-1.506,2.025) (-1.537,1.975) (-1.566,1.925) (-1.592,1.875) (-1.616,1.825) (-1.638,1.775) (-1.658,1.725) (-1.675,1.675) (-1.691,1.625) (-1.705,1.575) (-1.716,1.525) (-1.726,1.475) (-1.734,1.425) (-1.741,1.375) (-1.745,1.325) (-1.747,1.275) (-1.748,1.225) (-1.746,1.175) (-1.743,1.125) (-1.738,1.075) (-1.730,1.025) (-1.721,0.975) (-1.709,0.925) (-1.695,0.875) (-1.679,0.825) (-1.660,0.775) (-1.639,0.725) (-1.614,0.675) (-1.587,0.625) (-1.556,0.575) (-1.525,0.530) (-1.482,0.475) (-1.438,0.425) (-1.400,0.386) (-1.362,0.350) (-1.301,0.300) (-1.275,0.280)}
\EQ{(1.275,0.308) (1.323,0.350) (1.374,0.400) (1.419,0.450) (1.450,0.489) (1.493,0.550) (1.524,0.600) (1.550,0.647) (1.575,0.698) (1.597,0.750) (1.616,0.800) (1.632,0.850) (1.645,0.900) (1.656,0.950) (1.665,1.000) (1.672,1.050) (1.676,1.100) (1.679,1.150) (1.679,1.200) (1.678,1.250) (1.674,1.300) (1.669,1.350) (1.662,1.400) (1.652,1.450) (1.641,1.500) (1.628,1.550) (1.613,1.600) (1.595,1.650) (1.576,1.700) (1.554,1.750) (1.530,1.800) (1.504,1.850) (1.475,1.900) (1.444,1.950) (1.410,2.000) (1.375,2.047) (1.333,2.100) (1.289,2.150) (1.250,2.191) (1.200,2.240) (1.162,2.275) (1.102,2.325) (1.050,2.365) (1.001,2.400) (0.950,2.434) (0.900,2.464) (0.850,2.492) (0.800,2.518) (0.750,2.543) (0.700,2.565) (0.650,2.585) (0.600,2.604) (0.550,2.621) (0.500,2.636) (0.450,2.650) (0.400,2.662) (0.350,2.672) (0.300,2.682) (0.250,2.689) (0.200,2.696) (0.150,2.700) (0.100,2.704) (0.050,2.706) (0.000,2.707) (-0.050,2.706) (-0.100,2.704) (-0.150,2.700) (-0.200,2.696) (-0.250,2.689) (-0.300,2.682) (-0.350,2.672) (-0.400,2.662) (-0.450,2.650) (-0.500,2.636) (-0.550,2.621) (-0.600,2.604) (-0.650,2.585) (-0.700,2.565) (-0.750,2.543) (-0.800,2.518) (-0.850,2.492) (-0.900,2.464) (-0.950,2.434) (-1.000,2.401) (-1.050,2.365) (-1.100,2.327) (-1.150,2.285) (-1.189,2.250) (-1.241,2.200) (-1.289,2.150) (-1.325,2.109) (-1.373,2.050) (-1.410,2.000) (-1.444,1.950) (-1.475,1.900) (-1.504,1.850) (-1.530,1.800) (-1.554,1.750) (-1.575,1.702) (-1.595,1.650) (-1.613,1.600) (-1.628,1.550) (-1.641,1.500) (-1.652,1.450) (-1.662,1.400) (-1.669,1.350) (-1.674,1.300) (-1.678,1.250) (-1.679,1.200) (-1.679,1.150) (-1.676,1.100) (-1.672,1.050) (-1.665,1.000) (-1.656,0.950) (-1.645,0.900) (-1.632,0.850) (-1.616,0.800) (-1.597,0.750) (-1.576,0.700) (-1.552,0.650) (-1.525,0.601) (-1.493,0.550) (-1.458,0.500) (-1.425,0.458) (-1.375,0.401) (-1.325,0.352) (-1.275,0.308)}
\EQ{(1.275,0.337) (1.314,0.375) (1.350,0.414) (1.400,0.475) (1.436,0.525) (1.467,0.575) (1.495,0.625) (1.519,0.675) (1.540,0.725) (1.559,0.775) (1.574,0.825) (1.587,0.875) (1.598,0.925) (1.606,0.975) (1.612,1.025) (1.616,1.075) (1.618,1.125) (1.618,1.175) (1.616,1.225) (1.611,1.275) (1.605,1.325) (1.597,1.375) (1.586,1.425) (1.574,1.475) (1.559,1.525) (1.543,1.575) (1.524,1.625) (1.503,1.675) (1.479,1.725) (1.454,1.775) (1.425,1.825) (1.394,1.875) (1.360,1.925) (1.325,1.973) (1.283,2.025) (1.250,2.063) (1.216,2.100) (1.175,2.141) (1.125,2.187) (1.081,2.225) (1.025,2.269) (0.983,2.300) (0.925,2.339) (0.875,2.370) (0.825,2.398) (0.775,2.425) (0.725,2.449) (0.675,2.471) (0.625,2.492) (0.575,2.510) (0.525,2.527) (0.475,2.542) (0.425,2.555) (0.375,2.567) (0.325,2.578) (0.275,2.586) (0.225,2.594) (0.175,2.599) (0.125,2.604) (0.075,2.607) (0.025,2.608) (-0.050,2.608) (-0.100,2.605) (-0.150,2.602) (-0.200,2.597) (-0.250,2.590) (-0.300,2.582) (-0.350,2.573) (-0.400,2.562) (-0.450,2.549) (-0.500,2.535) (-0.550,2.519) (-0.600,2.501) (-0.650,2.482) (-0.700,2.460) (-0.750,2.437) (-0.800,2.412) (-0.850,2.384) (-0.900,2.355) (-0.946,2.325) (-1.000,2.288) (-1.050,2.250) (-1.100,2.209) (-1.139,2.175) (-1.191,2.125) (-1.239,2.075) (-1.275,2.034) (-1.323,1.975) (-1.360,1.925) (-1.394,1.875) (-1.425,1.825) (-1.450,1.781) (-1.475,1.734) (-1.500,1.681) (-1.524,1.625) (-1.543,1.575) (-1.559,1.525) (-1.574,1.475) (-1.586,1.425) (-1.597,1.375) (-1.605,1.325) (-1.611,1.275) (-1.616,1.225) (-1.618,1.175) (-1.618,1.125) (-1.616,1.075) (-1.612,1.025) (-1.606,0.975) (-1.598,0.925) (-1.587,0.875) (-1.574,0.825) (-1.559,0.775) (-1.540,0.725) (-1.519,0.675) (-1.495,0.625) (-1.467,0.575) (-1.436,0.525) (-1.400,0.475) (-1.360,0.425) (-1.325,0.387) (-1.289,0.350) (-1.275,0.337)}
\EQ{(1.266,0.359) (1.305,0.400) (1.346,0.450) (1.382,0.500) (1.414,0.550) (1.442,0.600) (1.467,0.650) (1.488,0.700) (1.506,0.750) (1.522,0.800) (1.534,0.850) (1.545,0.900) (1.553,0.950) (1.559,1.000) (1.562,1.050) (1.563,1.100) (1.563,1.150) (1.560,1.200) (1.555,1.250) (1.547,1.300) (1.538,1.350) (1.527,1.400) (1.513,1.450) (1.498,1.500) (1.480,1.550) (1.460,1.600) (1.437,1.650) (1.412,1.700) (1.385,1.750) (1.354,1.800) (1.325,1.844) (1.285,1.900) (1.250,1.944) (1.201,2.000) (1.154,2.050) (1.101,2.100) (1.050,2.145) (1.000,2.185) (0.950,2.222) (0.908,2.250) (0.850,2.287) (0.800,2.315) (0.750,2.342) (0.700,2.366) (0.650,2.388) (0.600,2.409) (0.550,2.427) (0.500,2.443) (0.450,2.458) (0.400,2.471) (0.350,2.483) (0.300,2.493) (0.250,2.501) (0.200,2.508) (0.150,2.513) (0.100,2.517) (0.050,2.519) (0.000,2.520) (-0.050,2.519) (-0.100,2.517) (-0.150,2.513) (-0.200,2.508) (-0.250,2.501) (-0.300,2.493) (-0.350,2.483) (-0.400,2.471) (-0.450,2.458) (-0.500,2.443) (-0.550,2.427) (-0.600,2.409) (-0.650,2.388) (-0.700,2.366) (-0.750,2.342) (-0.800,2.315) (-0.850,2.287) (-0.900,2.256) (-0.945,2.225) (-1.000,2.185) (-1.044,2.150) (-1.100,2.101) (-1.150,2.054) (-1.200,2.001) (-1.245,1.950) (-1.285,1.900) (-1.321,1.850) (-1.350,1.807) (-1.385,1.750) (-1.412,1.700) (-1.437,1.650) (-1.460,1.600) (-1.480,1.550) (-1.498,1.500) (-1.513,1.450) (-1.527,1.400) (-1.538,1.350) (-1.547,1.300) (-1.555,1.250) (-1.560,1.200) (-1.563,1.150) (-1.563,1.100) (-1.562,1.050) (-1.559,1.000) (-1.553,0.950) (-1.545,0.900) (-1.534,0.850) (-1.522,0.800) (-1.506,0.750) (-1.488,0.700) (-1.467,0.650) (-1.442,0.600) (-1.414,0.550) (-1.382,0.500) (-1.350,0.456) (-1.305,0.400) (-1.266,0.359)}
\EQ{(1.251,0.374) (1.295,0.425) (1.325,0.465) (1.365,0.525) (1.393,0.575) (1.418,0.625) (1.440,0.675) (1.458,0.725) (1.473,0.775) (1.486,0.825) (1.496,0.875) (1.504,0.925) (1.510,0.975) (1.513,1.025) (1.514,1.075) (1.512,1.125) (1.509,1.175) (1.503,1.225) (1.495,1.275) (1.486,1.325) (1.473,1.375) (1.459,1.425) (1.442,1.475) (1.423,1.525) (1.402,1.575) (1.378,1.625) (1.352,1.675) (1.325,1.721) (1.290,1.775) (1.255,1.825) (1.216,1.875) (1.175,1.923) (1.127,1.975) (1.075,2.025) (1.025,2.070) (0.975,2.110) (0.925,2.147) (0.875,2.182) (0.825,2.213) (0.775,2.242) (0.725,2.268) (0.675,2.292) (0.625,2.314) (0.575,2.334) (0.525,2.352) (0.475,2.369) (0.425,2.383) (0.375,2.396) (0.325,2.407) (0.275,2.416) (0.225,2.424) (0.175,2.430) (0.125,2.435) (0.075,2.438) (0.025,2.440) (-0.050,2.439) (-0.100,2.437) (-0.150,2.433) (-0.200,2.427) (-0.250,2.420) (-0.300,2.412) (-0.350,2.402) (-0.400,2.390) (-0.450,2.376) (-0.500,2.361) (-0.550,2.344) (-0.599,2.325) (-0.650,2.304) (-0.700,2.280) (-0.750,2.255) (-0.800,2.228) (-0.846,2.200) (-0.900,2.165) (-0.950,2.129) (-1.000,2.090) (-1.048,2.050) (-1.100,2.002) (-1.150,1.951) (-1.195,1.900) (-1.236,1.850) (-1.273,1.800) (-1.307,1.750) (-1.337,1.700) (-1.365,1.650) (-1.390,1.600) (-1.413,1.550) (-1.433,1.500) (-1.450,1.453) (-1.466,1.400) (-1.480,1.350) (-1.491,1.300) (-1.500,1.250) (-1.506,1.200) (-1.511,1.150) (-1.513,1.100) (-1.514,1.050) (-1.512,1.000) (-1.507,0.950) (-1.501,0.900) (-1.492,0.850) (-1.480,0.800) (-1.466,0.750) (-1.449,0.700) (-1.429,0.650) (-1.406,0.600) (-1.379,0.550) (-1.350,0.502) (-1.314,0.450) (-1.275,0.401) (-1.251,0.374)}
\EQ{(1.250,0.405) (1.285,0.450) (1.318,0.500) (1.347,0.550) (1.373,0.600) (1.395,0.650) (1.413,0.700) (1.429,0.750) (1.442,0.800) (1.452,0.850) (1.460,0.900) (1.465,0.950) (1.468,1.000) (1.468,1.050) (1.467,1.100) (1.463,1.150) (1.457,1.200) (1.448,1.250) (1.438,1.300) (1.425,1.350) (1.410,1.400) (1.392,1.450) (1.372,1.500) (1.350,1.550) (1.325,1.600) (1.300,1.645) (1.266,1.700) (1.232,1.750) (1.200,1.793) (1.154,1.850) (1.109,1.900) (1.059,1.950) (1.004,2.000) (0.950,2.044) (0.909,2.075) (0.850,2.116) (0.800,2.147) (0.751,2.175) (0.704,2.200) (0.652,2.225) (0.600,2.248) (0.550,2.268) (0.500,2.285) (0.450,2.301) (0.400,2.315) (0.350,2.328) (0.300,2.338) (0.250,2.347) (0.200,2.354) (0.150,2.360) (0.100,2.364) (0.050,2.366) (0.000,2.367) (-0.050,2.366) (-0.100,2.364) (-0.150,2.360) (-0.200,2.354) (-0.250,2.347) (-0.300,2.338) (-0.350,2.328) (-0.400,2.315) (-0.450,2.301) (-0.500,2.285) (-0.550,2.268) (-0.600,2.248) (-0.650,2.226) (-0.700,2.202) (-0.750,2.176) (-0.795,2.150) (-0.850,2.116) (-0.900,2.081) (-0.942,2.050) (-1.000,2.003) (-1.050,1.958) (-1.100,1.909) (-1.150,1.854) (-1.195,1.800) (-1.225,1.760) (-1.266,1.700) (-1.297,1.650) (-1.325,1.600) (-1.350,1.550) (-1.372,1.500) (-1.392,1.450) (-1.410,1.400) (-1.425,1.350) (-1.438,1.300) (-1.448,1.250) (-1.457,1.200) (-1.463,1.150) (-1.467,1.100) (-1.468,1.050) (-1.468,1.000) (-1.465,0.950) (-1.460,0.900) (-1.452,0.850) (-1.442,0.800) (-1.429,0.750) (-1.413,0.700) (-1.395,0.650) (-1.373,0.600) (-1.350,0.555) (-1.325,0.512) (-1.285,0.450) (-1.250,0.405)}
\EQ{(1.244,0.431) (1.274,0.475) (1.300,0.518) (1.325,0.565) (1.350,0.619) (1.372,0.675) (1.388,0.725) (1.400,0.775) (1.411,0.825) (1.419,0.875) (1.424,0.925) (1.427,0.975) (1.427,1.025) (1.425,1.075) (1.421,1.125) (1.414,1.175) (1.406,1.225) (1.395,1.275) (1.381,1.325) (1.365,1.375) (1.347,1.425) (1.326,1.475) (1.303,1.525) (1.277,1.575) (1.250,1.621) (1.215,1.675) (1.180,1.725) (1.141,1.775) (1.100,1.822) (1.050,1.874) (1.000,1.922) (0.950,1.965) (0.906,2.000) (0.850,2.040) (0.800,2.073) (0.755,2.100) (0.709,2.125) (0.660,2.150) (0.605,2.175) (0.550,2.198) (0.500,2.216) (0.450,2.233) (0.400,2.247) (0.350,2.260) (0.300,2.271) (0.250,2.280) (0.200,2.287) (0.150,2.293) (0.100,2.297) (0.050,2.299) (0.000,2.300) (-0.050,2.299) (-0.100,2.297) (-0.150,2.293) (-0.200,2.287) (-0.250,2.280) (-0.300,2.271) (-0.350,2.260) (-0.400,2.247) (-0.450,2.233) (-0.500,2.216) (-0.550,2.198) (-0.600,2.177) (-0.650,2.155) (-0.700,2.130) (-0.750,2.103) (-0.796,2.075) (-0.850,2.040) (-0.900,2.004) (-0.950,1.965) (-0.997,1.925) (-1.049,1.875) (-1.097,1.825) (-1.141,1.775) (-1.175,1.731) (-1.215,1.675) (-1.248,1.625) (-1.275,1.578) (-1.300,1.531) (-1.325,1.478) (-1.347,1.425) (-1.365,1.375) (-1.381,1.325) (-1.395,1.275) (-1.406,1.225) (-1.414,1.175) (-1.421,1.125) (-1.425,1.075) (-1.427,1.025) (-1.427,0.975) (-1.424,0.925) (-1.419,0.875) (-1.411,0.825) (-1.400,0.775) (-1.388,0.725) (-1.372,0.675) (-1.352,0.625) (-1.330,0.575) (-1.304,0.525) (-1.275,0.477) (-1.244,0.431)}
\EQ{(1.229,0.446) (1.263,0.500) (1.290,0.550) (1.313,0.600) (1.333,0.650) (1.349,0.700) (1.362,0.750) (1.373,0.800) (1.381,0.850) (1.386,0.900) (1.389,0.950) (1.389,1.000) (1.387,1.050) (1.383,1.100) (1.376,1.150) (1.367,1.200) (1.355,1.250) (1.341,1.300) (1.325,1.349) (1.306,1.400) (1.284,1.450) (1.260,1.500) (1.233,1.550) (1.203,1.600) (1.175,1.642) (1.132,1.700) (1.100,1.740) (1.050,1.796) (1.000,1.846) (0.950,1.891) (0.909,1.925) (0.850,1.970) (0.806,2.000) (0.750,2.035) (0.700,2.063) (0.650,2.089) (0.600,2.112) (0.550,2.133) (0.500,2.152) (0.450,2.169) (0.400,2.184) (0.350,2.197) (0.300,2.208) (0.250,2.218) (0.200,2.225) (0.150,2.231) (0.100,2.236) (0.050,2.238) (0.000,2.239) (-0.050,2.238) (-0.100,2.236) (-0.150,2.231) (-0.200,2.225) (-0.250,2.218) (-0.300,2.208) (-0.350,2.197) (-0.400,2.184) (-0.450,2.169) (-0.500,2.152) (-0.550,2.133) (-0.600,2.112) (-0.650,2.089) (-0.700,2.063) (-0.750,2.035) (-0.800,2.004) (-0.843,1.975) (-0.900,1.932) (-0.940,1.900) (-0.996,1.850) (-1.046,1.800) (-1.091,1.750) (-1.125,1.709) (-1.169,1.650) (-1.200,1.604) (-1.233,1.550) (-1.260,1.500) (-1.284,1.450) (-1.306,1.400) (-1.325,1.350) (-1.341,1.300) (-1.355,1.250) (-1.367,1.200) (-1.376,1.150) (-1.383,1.100) (-1.387,1.050) (-1.389,1.000) (-1.389,0.950) (-1.386,0.900) (-1.381,0.850) (-1.373,0.800) (-1.362,0.750) (-1.349,0.700) (-1.333,0.650) (-1.313,0.600) (-1.290,0.550) (-1.263,0.500) (-1.232,0.450) (-1.229,0.446)}
\EQ{(1.224,0.476) (1.250,0.522) (1.275,0.574) (1.296,0.625) (1.313,0.675) (1.326,0.725) (1.337,0.775) (1.345,0.825) (1.350,0.875) (1.353,0.925) (1.354,0.975) (1.352,1.025) (1.347,1.075) (1.340,1.125) (1.331,1.175) (1.319,1.225) (1.305,1.275) (1.288,1.325) (1.269,1.375) (1.246,1.425) (1.221,1.475) (1.194,1.525) (1.162,1.575) (1.128,1.625) (1.089,1.675) (1.050,1.721) (1.000,1.774) (0.950,1.822) (0.900,1.865) (0.856,1.900) (0.800,1.940) (0.750,1.972) (0.702,2.000) (0.656,2.025) (0.604,2.050) (0.550,2.074) (0.500,2.093) (0.450,2.111) (0.400,2.126) (0.350,2.139) (0.300,2.151) (0.250,2.161) (0.200,2.168) (0.150,2.174) (0.100,2.179) (0.050,2.181) (0.000,2.182) (-0.050,2.181) (-0.100,2.179) (-0.150,2.174) (-0.200,2.168) (-0.250,2.161) (-0.300,2.151) (-0.350,2.139) (-0.400,2.126) (-0.450,2.111) (-0.500,2.093) (-0.547,2.075) (-0.600,2.052) (-0.650,2.028) (-0.700,2.001) (-0.745,1.975) (-0.800,1.940) (-0.850,1.904) (-0.900,1.865) (-0.947,1.825) (-0.999,1.775) (-1.047,1.725) (-1.089,1.675) (-1.125,1.629) (-1.162,1.575) (-1.194,1.525) (-1.221,1.475) (-1.246,1.425) (-1.269,1.375) (-1.288,1.325) (-1.305,1.275) (-1.319,1.225) (-1.331,1.175) (-1.340,1.125) (-1.347,1.075) (-1.352,1.025) (-1.354,0.975) (-1.353,0.925) (-1.350,0.875) (-1.345,0.825) (-1.337,0.775) (-1.326,0.725) (-1.313,0.675) (-1.296,0.625) (-1.275,0.575) (-1.252,0.525) (-1.225,0.477)}

    \draw[mmaEQ] (-10,0) -- (10,0);
    % Dipolo puntual - líneas de campo
% Generado por scripts/generar_datos.py (no editar a mano).
\LC{(1.082,0.489) (1.149,0.455) (1.212,0.414) (1.269,0.365) (1.319,0.309) (1.361,0.248) (1.394,0.180) (1.417,0.109) (1.428,0.035) (1.428,-0.040) (1.415,-0.114) (1.392,-0.185) (1.358,-0.252) (1.316,-0.314) (1.265,-0.369) (1.207,-0.417) (1.144,-0.458) (1.077,-0.491)}
\LC{(0.453,1.106) (0.495,1.168) (0.539,1.229) (0.583,1.290) (0.629,1.349) (0.676,1.408) (0.724,1.465) (0.773,1.522) (0.823,1.578) (0.873,1.633) (0.925,1.688) (0.978,1.741) (1.031,1.794) (1.085,1.846) (1.141,1.896) (1.197,1.946) (1.253,1.996) (1.311,2.044) (1.369,2.091) (1.428,2.137) (1.487,2.183) (1.548,2.228) (1.609,2.271) (1.670,2.314) (1.733,2.356) (1.795,2.397) (1.859,2.437) (1.923,2.476) (1.988,2.514) (2.053,2.551) (2.118,2.587) (2.185,2.622) (2.251,2.657) (2.319,2.690) (2.386,2.722) (2.454,2.753) (2.523,2.784) (2.592,2.813) (2.661,2.841) (2.731,2.869) (2.802,2.895) (2.872,2.920) (2.943,2.945) (3.014,2.968) (3.086,2.990) (3.158,3.011) (3.230,3.031) (3.303,3.050) (3.376,3.068) (3.449,3.085) (3.522,3.101) (3.596,3.116) (3.669,3.130) (3.743,3.142) (3.817,3.154) (3.892,3.164) (3.966,3.173) (4.041,3.181) (4.115,3.188) (4.190,3.194) (4.265,3.199) (4.340,3.203) (4.415,3.205) (4.490,3.206) (4.565,3.206) (4.640,3.205) (4.715,3.203) (4.790,3.200) (4.864,3.195) (4.939,3.189) (5.014,3.182) (5.088,3.174) (5.163,3.165) (5.237,3.154) (5.311,3.142) (5.385,3.129) (5.459,3.115) (5.532,3.099) (5.605,3.082) (5.678,3.064) (5.750,3.045) (5.822,3.024) (5.894,3.002) (5.965,2.979) (6.036,2.955) (6.107,2.929) (6.177,2.902) (6.246,2.874) (6.315,2.844) (6.384,2.813) (6.451,2.781) (6.518,2.748) (6.585,2.713) (6.651,2.677) (6.716,2.640) (6.780,2.601) (6.844,2.561) (6.906,2.520) (6.968,2.478) (7.029,2.434) (7.089,2.389) (7.148,2.343) (7.207,2.296) (7.264,2.247) (7.320,2.197) (7.374,2.146) (7.428,2.094) (7.481,2.040) (7.532,1.985) (7.582,1.930) (7.631,1.873) (7.679,1.815) (7.725,1.755) (7.769,1.695) (7.813,1.634) (7.854,1.572) (7.894,1.508) (7.933,1.444) (7.970,1.379) (8.006,1.313) (8.039,1.246) (8.071,1.178) (8.102,1.109) (8.130,1.040) (8.157,0.970) (8.182,0.899) (8.205,0.828) (8.226,0.756) (8.246,0.683) (8.263,0.610) (8.278,0.537) (8.292,0.463) (8.303,0.389) (8.313,0.315) (8.320,0.240) (8.326,0.165) (8.329,0.090) (8.331,0.015) (8.330,-0.060) (8.327,-0.134) (8.323,-0.209) (8.316,-0.284) (8.307,-0.359) (8.297,-0.433) (8.284,-0.507) (8.270,-0.580) (8.253,-0.653) (8.234,-0.726) (8.214,-0.798) (8.192,-0.870) (8.168,-0.941) (8.142,-1.011) (8.114,-1.081) (8.084,-1.150) (8.053,-1.218) (8.020,-1.285) (7.985,-1.352) (7.949,-1.417) (7.911,-1.482) (7.871,-1.546) (7.830,-1.608) (7.787,-1.670) (7.743,-1.731) (7.698,-1.790) (7.651,-1.849) (7.603,-1.906) (7.553,-1.963) (7.502,-2.018) (7.450,-2.072) (7.397,-2.124) (7.342,-2.176) (7.287,-2.226) (7.230,-2.276) (7.173,-2.324) (7.114,-2.370) (7.054,-2.416) (6.994,-2.460) (6.932,-2.503) (6.870,-2.544) (6.806,-2.585) (6.742,-2.624) (6.678,-2.662) (6.612,-2.698) (6.546,-2.733) (6.479,-2.767) (6.412,-2.800) (6.343,-2.831) (6.275,-2.862) (6.206,-2.890) (6.136,-2.918) (6.066,-2.944) (5.995,-2.969) (5.924,-2.993) (5.852,-3.015) (5.780,-3.036) (5.708,-3.056) (5.635,-3.075) (5.562,-3.092) (5.489,-3.108) (5.415,-3.123) (5.342,-3.137) (5.268,-3.149) (5.194,-3.160) (5.119,-3.170) (5.045,-3.179) (4.970,-3.187) (4.895,-3.193) (4.821,-3.198) (4.746,-3.202) (4.671,-3.205) (4.596,-3.206) (4.521,-3.206) (4.446,-3.206) (4.371,-3.204) (4.296,-3.201) (4.221,-3.196) (4.146,-3.191) (4.071,-3.184) (3.997,-3.177) (3.922,-3.168) (3.848,-3.158) (3.774,-3.147) (3.700,-3.135) (3.626,-3.122) (3.552,-3.107) (3.479,-3.092) (3.406,-3.075) (3.333,-3.058) (3.260,-3.039) (3.188,-3.020) (3.116,-2.999) (3.044,-2.977) (2.973,-2.954) (2.901,-2.930) (2.831,-2.906) (2.760,-2.880) (2.690,-2.853) (2.621,-2.825) (2.551,-2.796) (2.483,-2.766) (2.414,-2.735) (2.346,-2.703) (2.279,-2.670) (2.212,-2.637) (2.146,-2.602) (2.080,-2.566) (2.014,-2.529) (1.950,-2.492) (1.885,-2.453) (1.822,-2.413) (1.758,-2.373) (1.696,-2.331) (1.634,-2.289) (1.573,-2.246) (1.512,-2.201) (1.452,-2.156) (1.393,-2.110) (1.335,-2.063) (1.277,-2.016) (1.220,-1.967) (1.164,-1.917) (1.108,-1.867) (1.053,-1.815) (1.000,-1.763) (0.947,-1.710) (0.895,-1.656) (0.843,-1.601) (0.793,-1.545) (0.744,-1.489) (0.696,-1.432) (0.648,-1.373) (0.602,-1.314) (0.557,-1.254) (0.513,-1.194) (0.470,-1.132) (0.442,-1.091)}
\LC{(-0.453,1.106) (-0.495,1.168) (-0.539,1.229) (-0.583,1.290) (-0.629,1.349) (-0.676,1.408) (-0.724,1.465) (-0.773,1.522) (-0.823,1.578) (-0.873,1.633) (-0.925,1.688) (-0.978,1.741) (-1.031,1.794) (-1.085,1.846) (-1.141,1.896) (-1.197,1.946) (-1.253,1.996) (-1.311,2.044) (-1.369,2.091) (-1.428,2.137) (-1.487,2.183) (-1.548,2.228) (-1.609,2.271) (-1.670,2.314) (-1.733,2.356) (-1.795,2.397) (-1.859,2.437) (-1.923,2.476) (-1.988,2.514) (-2.053,2.551) (-2.118,2.587) (-2.185,2.622) (-2.251,2.657) (-2.319,2.690) (-2.386,2.722) (-2.454,2.753) (-2.523,2.784) (-2.592,2.813) (-2.661,2.841) (-2.731,2.869) (-2.802,2.895) (-2.872,2.920) (-2.943,2.945) (-3.014,2.968) (-3.086,2.990) (-3.158,3.011) (-3.230,3.031) (-3.303,3.050) (-3.376,3.068) (-3.449,3.085) (-3.522,3.101) (-3.596,3.116) (-3.669,3.130) (-3.743,3.142) (-3.817,3.154) (-3.892,3.164) (-3.966,3.173) (-4.041,3.181) (-4.115,3.188) (-4.190,3.194) (-4.265,3.199) (-4.340,3.203) (-4.415,3.205) (-4.490,3.206) (-4.565,3.206) (-4.640,3.205) (-4.715,3.203) (-4.790,3.200) (-4.864,3.195) (-4.939,3.189) (-5.014,3.182) (-5.088,3.174) (-5.163,3.165) (-5.237,3.154) (-5.311,3.142) (-5.385,3.129) (-5.459,3.115) (-5.532,3.099) (-5.605,3.082) (-5.678,3.064) (-5.750,3.045) (-5.822,3.024) (-5.894,3.002) (-5.965,2.979) (-6.036,2.955) (-6.107,2.929) (-6.177,2.902) (-6.246,2.874) (-6.315,2.844) (-6.384,2.813) (-6.451,2.781) (-6.518,2.748) (-6.585,2.713) (-6.651,2.677) (-6.716,2.640) (-6.780,2.601) (-6.844,2.561) (-6.906,2.520) (-6.968,2.478) (-7.029,2.434) (-7.089,2.389) (-7.148,2.343) (-7.207,2.296) (-7.264,2.247) (-7.320,2.197) (-7.374,2.146) (-7.428,2.094) (-7.481,2.040) (-7.532,1.985) (-7.582,1.930) (-7.631,1.873) (-7.679,1.815) (-7.725,1.755) (-7.769,1.695) (-7.813,1.634) (-7.854,1.572) (-7.894,1.508) (-7.933,1.444) (-7.970,1.379) (-8.006,1.313) (-8.039,1.246) (-8.071,1.178) (-8.102,1.109) (-8.130,1.040) (-8.157,0.970) (-8.182,0.899) (-8.205,0.828) (-8.226,0.756) (-8.246,0.683) (-8.263,0.610) (-8.278,0.537) (-8.292,0.463) (-8.303,0.389) (-8.313,0.315) (-8.320,0.240) (-8.326,0.165) (-8.329,0.090) (-8.331,0.015) (-8.330,-0.060) (-8.327,-0.134) (-8.323,-0.209) (-8.316,-0.284) (-8.307,-0.359) (-8.297,-0.433) (-8.284,-0.507) (-8.270,-0.580) (-8.253,-0.653) (-8.234,-0.726) (-8.214,-0.798) (-8.192,-0.870) (-8.168,-0.941) (-8.142,-1.011) (-8.114,-1.081) (-8.084,-1.150) (-8.053,-1.218) (-8.020,-1.285) (-7.985,-1.352) (-7.949,-1.417) (-7.911,-1.482) (-7.871,-1.546) (-7.830,-1.608) (-7.787,-1.670) (-7.743,-1.731) (-7.698,-1.790) (-7.651,-1.849) (-7.603,-1.906) (-7.553,-1.963) (-7.502,-2.018) (-7.450,-2.072) (-7.397,-2.124) (-7.342,-2.176) (-7.287,-2.226) (-7.230,-2.276) (-7.173,-2.324) (-7.114,-2.370) (-7.054,-2.416) (-6.994,-2.460) (-6.932,-2.503) (-6.870,-2.544) (-6.806,-2.585) (-6.742,-2.624) (-6.678,-2.662) (-6.612,-2.698) (-6.546,-2.733) (-6.479,-2.767) (-6.412,-2.800) (-6.343,-2.831) (-6.275,-2.862) (-6.206,-2.890) (-6.136,-2.918) (-6.066,-2.944) (-5.995,-2.969) (-5.924,-2.993) (-5.852,-3.015) (-5.780,-3.036) (-5.708,-3.056) (-5.635,-3.075) (-5.562,-3.092) (-5.489,-3.108) (-5.415,-3.123) (-5.342,-3.137) (-5.268,-3.149) (-5.194,-3.160) (-5.119,-3.170) (-5.045,-3.179) (-4.970,-3.187) (-4.895,-3.193) (-4.821,-3.198) (-4.746,-3.202) (-4.671,-3.205) (-4.596,-3.206) (-4.521,-3.206) (-4.446,-3.206) (-4.371,-3.204) (-4.296,-3.201) (-4.221,-3.196) (-4.146,-3.191) (-4.071,-3.184) (-3.997,-3.177) (-3.922,-3.168) (-3.848,-3.158) (-3.774,-3.147) (-3.700,-3.135) (-3.626,-3.122) (-3.552,-3.107) (-3.479,-3.092) (-3.406,-3.075) (-3.333,-3.058) (-3.260,-3.039) (-3.188,-3.020) (-3.116,-2.999) (-3.044,-2.977) (-2.973,-2.954) (-2.901,-2.930) (-2.831,-2.906) (-2.760,-2.880) (-2.690,-2.853) (-2.621,-2.825) (-2.551,-2.796) (-2.483,-2.766) (-2.414,-2.735) (-2.346,-2.703) (-2.279,-2.670) (-2.212,-2.637) (-2.146,-2.602) (-2.080,-2.566) (-2.014,-2.529) (-1.950,-2.492) (-1.885,-2.453) (-1.822,-2.413) (-1.758,-2.373) (-1.696,-2.331) (-1.634,-2.289) (-1.573,-2.246) (-1.512,-2.201) (-1.452,-2.156) (-1.393,-2.110) (-1.335,-2.063) (-1.277,-2.016) (-1.220,-1.967) (-1.164,-1.917) (-1.108,-1.867) (-1.053,-1.815) (-1.000,-1.763) (-0.947,-1.710) (-0.895,-1.656) (-0.843,-1.601) (-0.793,-1.545) (-0.744,-1.489) (-0.696,-1.432) (-0.648,-1.373) (-0.602,-1.314) (-0.557,-1.254) (-0.513,-1.194) (-0.470,-1.132) (-0.442,-1.091)}
\LC{(-1.082,0.489) (-1.149,0.455) (-1.212,0.414) (-1.269,0.365) (-1.319,0.309) (-1.361,0.248) (-1.394,0.180) (-1.417,0.109) (-1.428,0.035) (-1.428,-0.040) (-1.415,-0.114) (-1.392,-0.185) (-1.358,-0.252) (-1.316,-0.314) (-1.265,-0.369) (-1.207,-0.417) (-1.144,-0.458) (-1.077,-0.491)}
\LC{(4.893,-10.000) (4.845,-9.951) (4.797,-9.893) (4.749,-9.835) (4.702,-9.777) (4.655,-9.719) (4.607,-9.661) (4.560,-9.602) (4.514,-9.544) (4.467,-9.485) (4.420,-9.426) (4.374,-9.368) (4.328,-9.309) (4.281,-9.249) (4.235,-9.190) (4.190,-9.131) (4.144,-9.071) (4.099,-9.012) (4.053,-8.952) (4.008,-8.892) (3.963,-8.832) (3.918,-8.772) (3.873,-8.712) (3.829,-8.651) (3.785,-8.591) (3.740,-8.530) (3.696,-8.470) (3.652,-8.409) (3.609,-8.348) (3.565,-8.287) (3.522,-8.226) (3.479,-8.164) (3.435,-8.103) (3.393,-8.041) (3.350,-7.980) (3.307,-7.918) (3.265,-7.856) (3.223,-7.794) (3.181,-7.732) (3.139,-7.670) (3.097,-7.607) (3.056,-7.545) (3.015,-7.482) (2.973,-7.419) (2.932,-7.357) (2.892,-7.294) (2.851,-7.231) (2.811,-7.167) (2.771,-7.104) (2.730,-7.041) (2.691,-6.977) (2.651,-6.913) (2.612,-6.850) (2.572,-6.786) (2.533,-6.722) (2.494,-6.658) (2.456,-6.593) (2.417,-6.529) (2.379,-6.465) (2.341,-6.400) (2.303,-6.335) (2.265,-6.270) (2.228,-6.205) (2.190,-6.140) (2.153,-6.075) (2.116,-6.010) (2.080,-5.945) (2.043,-5.879) (2.007,-5.813) (1.971,-5.748) (1.935,-5.682) (1.899,-5.616) (1.864,-5.550) (1.829,-5.483) (1.794,-5.417) (1.759,-5.351) (1.724,-5.284) (1.690,-5.217) (1.656,-5.151) (1.622,-5.084) (1.588,-5.017) (1.555,-4.950) (1.522,-4.882) (1.489,-4.815) (1.456,-4.747) (1.423,-4.680) (1.391,-4.612) (1.359,-4.544) (1.327,-4.476) (1.296,-4.408) (1.265,-4.340) (1.234,-4.272) (1.203,-4.203) (1.172,-4.135) (1.142,-4.066) (1.112,-3.998) (1.082,-3.929) (1.053,-3.860) (1.024,-3.791) (0.995,-3.721) (0.966,-3.652) (0.938,-3.583) (0.910,-3.513) (0.882,-3.444) (0.854,-3.374) (0.827,-3.304) (0.800,-3.234) (0.773,-3.164) (0.747,-3.094) (0.721,-3.023) (0.695,-2.953) (0.670,-2.882) (0.645,-2.812) (0.620,-2.741) (0.595,-2.670) (0.571,-2.599) (0.547,-2.528) (0.524,-2.457) (0.501,-2.385) (0.478,-2.314) (0.455,-2.242) (0.433,-2.171) (0.412,-2.099) (0.390,-2.027) (0.369,-1.955) (0.349,-1.883) (0.328,-1.811) (0.309,-1.738) (0.289,-1.666) (0.283,-1.642)}
\LC{(-5.042,-10.000) (-4.993,-9.951) (-4.944,-9.894) (-4.895,-9.837) (-4.847,-9.780) (-4.798,-9.722) (-4.750,-9.665) (-4.702,-9.607) (-4.654,-9.550) (-4.606,-9.492) (-4.558,-9.434) (-4.511,-9.376) (-4.463,-9.318) (-4.416,-9.259) (-4.369,-9.201) (-4.322,-9.142) (-4.276,-9.084) (-4.229,-9.025) (-4.183,-8.966) (-4.136,-8.907) (-4.090,-8.848) (-4.044,-8.789) (-3.999,-8.729) (-3.953,-8.670) (-3.908,-8.610) (-3.862,-8.550) (-3.817,-8.490) (-3.772,-8.430) (-3.728,-8.370) (-3.683,-8.310) (-3.639,-8.249) (-3.594,-8.189) (-3.550,-8.128) (-3.506,-8.067) (-3.463,-8.006) (-3.419,-7.945) (-3.376,-7.884) (-3.333,-7.823) (-3.289,-7.761) (-3.247,-7.700) (-3.204,-7.638) (-3.161,-7.576) (-3.119,-7.514) (-3.077,-7.452) (-3.035,-7.390) (-2.993,-7.328) (-2.952,-7.265) (-2.910,-7.203) (-2.869,-7.140) (-2.828,-7.077) (-2.787,-7.015) (-2.747,-6.951) (-2.706,-6.888) (-2.666,-6.825) (-2.626,-6.762) (-2.586,-6.698) (-2.546,-6.634) (-2.507,-6.571) (-2.467,-6.507) (-2.428,-6.443) (-2.389,-6.379) (-2.351,-6.314) (-2.312,-6.250) (-2.274,-6.186) (-2.236,-6.121) (-2.198,-6.056) (-2.160,-5.991) (-2.123,-5.926) (-2.086,-5.861) (-2.049,-5.796) (-2.012,-5.731) (-1.975,-5.665) (-1.939,-5.600) (-1.903,-5.534) (-1.867,-5.468) (-1.831,-5.402) (-1.796,-5.336) (-1.760,-5.270) (-1.725,-5.204) (-1.690,-5.137) (-1.656,-5.071) (-1.621,-5.004) (-1.587,-4.937) (-1.553,-4.870) (-1.520,-4.803) (-1.486,-4.736) (-1.453,-4.669) (-1.420,-4.601) (-1.388,-4.534) (-1.355,-4.466) (-1.323,-4.399) (-1.291,-4.331) (-1.260,-4.263) (-1.228,-4.195) (-1.197,-4.126) (-1.166,-4.058) (-1.135,-3.990) (-1.105,-3.921) (-1.075,-3.852) (-1.045,-3.783) (-1.016,-3.714) (-0.986,-3.645) (-0.957,-3.576) (-0.929,-3.507) (-0.900,-3.438) (-0.872,-3.368) (-0.844,-3.298) (-0.817,-3.229) (-0.790,-3.159) (-0.763,-3.089) (-0.736,-3.019) (-0.710,-2.948) (-0.684,-2.878) (-0.658,-2.808) (-0.633,-2.737) (-0.608,-2.666) (-0.583,-2.595) (-0.559,-2.524) (-0.535,-2.453) (-0.511,-2.382) (-0.488,-2.311) (-0.465,-2.240) (-0.442,-2.168) (-0.420,-2.096) (-0.398,-2.025) (-0.377,-1.953) (-0.356,-1.881) (-0.335,-1.809) (-0.315,-1.736) (-0.295,-1.664)}
\LC{(10.000,-6.289) (9.931,-6.297) (9.856,-6.305) (9.782,-6.311) (9.707,-6.317) (9.632,-6.323) (9.557,-6.328) (9.482,-6.332) (9.407,-6.336) (9.332,-6.339) (9.257,-6.342) (9.182,-6.344) (9.107,-6.345) (9.032,-6.346) (8.957,-6.346) (8.882,-6.346) (8.808,-6.344) (8.733,-6.343) (8.658,-6.341) (8.583,-6.338) (8.508,-6.335) (8.433,-6.331) (8.358,-6.326) (8.283,-6.321) (8.208,-6.315) (8.134,-6.309) (8.059,-6.302) (7.984,-6.295) (7.910,-6.287) (7.835,-6.278) (7.761,-6.269) (7.686,-6.260) (7.612,-6.249) (7.538,-6.239) (7.464,-6.227) (7.390,-6.215) (7.316,-6.203) (7.242,-6.190) (7.168,-6.176) (7.094,-6.162) (7.021,-6.148) (6.947,-6.133) (6.874,-6.117) (6.801,-6.100) (6.728,-6.084) (6.655,-6.066) (6.582,-6.048) (6.509,-6.030) (6.437,-6.011) (6.364,-5.991) (6.292,-5.971) (6.220,-5.951) (6.148,-5.930) (6.076,-5.908) (6.004,-5.886) (5.933,-5.863) (5.862,-5.840) (5.791,-5.816) (5.720,-5.792) (5.649,-5.767) (5.578,-5.742) (5.508,-5.716) (5.438,-5.690) (5.367,-5.663) (5.298,-5.636) (5.228,-5.608) (5.159,-5.579) (5.089,-5.550) (5.020,-5.521) (4.952,-5.491) (4.883,-5.461) (4.815,-5.430) (4.747,-5.399) (4.679,-5.367) (4.611,-5.334) (4.544,-5.301) (4.476,-5.268) (4.410,-5.234) (4.343,-5.200) (4.276,-5.165) (4.210,-5.130) (4.144,-5.094) (4.079,-5.058) (4.013,-5.021) (3.948,-4.984) (3.883,-4.946) (3.819,-4.908) (3.755,-4.869) (3.691,-4.830) (3.627,-4.790) (3.564,-4.750) (3.501,-4.709) (3.438,-4.668) (3.375,-4.627) (3.313,-4.585) (3.251,-4.542) (3.190,-4.499) (3.129,-4.456) (3.068,-4.412) (3.008,-4.368) (2.947,-4.323) (2.888,-4.278) (2.828,-4.232) (2.769,-4.186) (2.710,-4.139) (2.652,-4.092) (2.594,-4.045) (2.536,-3.997) (2.479,-3.948) (2.422,-3.899) (2.366,-3.850) (2.309,-3.800) (2.254,-3.750) (2.198,-3.699) (2.144,-3.648) (2.089,-3.597) (2.035,-3.545) (1.981,-3.492) (1.928,-3.439) (1.875,-3.386) (1.823,-3.332) (1.771,-3.278) (1.720,-3.223) (1.669,-3.168) (1.618,-3.113) (1.568,-3.057) (1.519,-3.001) (1.470,-2.944) (1.421,-2.887) (1.373,-2.829) (1.326,-2.771) (1.279,-2.713) (1.233,-2.654) (1.187,-2.594) (1.141,-2.535) (1.097,-2.474) (1.052,-2.414)}
\LC{(-2.969,2.475) (-3.043,2.485) (-3.118,2.495) (-3.192,2.503) (-3.267,2.509) (-3.342,2.515) (-3.417,2.518) (-3.492,2.520) (-3.567,2.521) (-3.642,2.520) (-3.717,2.518) (-3.792,2.514) (-3.866,2.509) (-3.941,2.502) (-4.016,2.494) (-4.090,2.484) (-4.164,2.472) (-4.238,2.459) (-4.312,2.445) (-4.385,2.428) (-4.458,2.411) (-4.530,2.391) (-4.602,2.370) (-4.674,2.348) (-4.745,2.323) (-4.815,2.298) (-4.885,2.270) (-4.954,2.241) (-5.022,2.210) (-5.090,2.178) (-5.157,2.144) (-5.223,2.108) (-5.288,2.071) (-5.352,2.032) (-5.415,1.991) (-5.477,1.949) (-5.538,1.905) (-5.598,1.860) (-5.656,1.813) (-5.713,1.765) (-5.769,1.714) (-5.823,1.663) (-5.876,1.610) (-5.928,1.555) (-5.977,1.499) (-6.025,1.441) (-6.072,1.382) (-6.116,1.322) (-6.159,1.260) (-6.200,1.197) (-6.239,1.133) (-6.275,1.068) (-6.310,1.001) (-6.343,0.934) (-6.373,0.865) (-6.401,0.796) (-6.427,0.725) (-6.450,0.654) (-6.471,0.582) (-6.490,0.509) (-6.506,0.436) (-6.520,0.362) (-6.531,0.288) (-6.539,0.214) (-6.545,0.139) (-6.549,0.064) (-6.550,-0.011) (-6.548,-0.086) (-6.544,-0.161) (-6.537,-0.236) (-6.528,-0.310) (-6.516,-0.384) (-6.502,-0.458) (-6.485,-0.531) (-6.465,-0.603) (-6.444,-0.675) (-6.419,-0.746) (-6.393,-0.816) (-6.364,-0.885) (-6.333,-0.954) (-6.300,-1.021) (-6.265,-1.087) (-6.227,-1.152) (-6.188,-1.216) (-6.147,-1.279) (-6.103,-1.340) (-6.058,-1.400) (-6.012,-1.458) (-5.963,-1.515) (-5.913,-1.571) (-5.861,-1.625) (-5.808,-1.678) (-5.753,-1.729) (-5.697,-1.779) (-5.639,-1.827) (-5.580,-1.874) (-5.520,-1.918) (-5.459,-1.962) (-5.397,-2.004) (-5.333,-2.044) (-5.269,-2.082) (-5.203,-2.119) (-5.137,-2.154) (-5.070,-2.188) (-5.002,-2.219) (-4.934,-2.250) (-4.864,-2.278) (-4.794,-2.305) (-4.724,-2.331) (-4.653,-2.355) (-4.581,-2.377) (-4.509,-2.397) (-4.436,-2.416) (-4.363,-2.433) (-4.290,-2.449) (-4.216,-2.463) (-4.142,-2.476) (-4.068,-2.487) (-3.994,-2.496) (-3.919,-2.504) (-3.845,-2.511) (-3.770,-2.515) (-3.695,-2.519) (-3.620,-2.521) (-3.545,-2.521) (-3.470,-2.520) (-3.395,-2.517) (-3.320,-2.513) (-3.245,-2.508) (-3.170,-2.501) (-3.096,-2.492) (-3.022,-2.482) (-2.972,-2.475)}
\LC{(0.389,1.769) (0.413,1.840) (0.437,1.911) (0.463,1.982) (0.488,2.052) (0.514,2.123) (0.541,2.193) (0.568,2.263) (0.595,2.333) (0.623,2.402) (0.651,2.472) (0.680,2.541) (0.709,2.610) (0.738,2.679) (0.768,2.748) (0.798,2.817) (0.829,2.885) (0.860,2.954) (0.891,3.022) (0.923,3.090) (0.955,3.158) (0.987,3.225) (1.020,3.293) (1.053,3.360) (1.086,3.427) (1.120,3.494) (1.154,3.561) (1.189,3.628) (1.223,3.694) (1.258,3.760) (1.294,3.827) (1.329,3.893) (1.365,3.958) (1.402,4.024) (1.438,4.089) (1.475,4.155) (1.512,4.220) (1.550,4.285) (1.588,4.350) (1.626,4.414) (1.664,4.479) (1.703,4.543) (1.742,4.607) (1.781,4.671) (1.821,4.735) (1.860,4.798) (1.900,4.862) (1.941,4.925) (1.981,4.988) (2.022,5.051) (2.063,5.113) (2.105,5.176) (2.146,5.238) (2.188,5.301) (2.230,5.363) (2.273,5.424) (2.316,5.486) (2.358,5.548) (2.402,5.609) (2.445,5.670) (2.489,5.731) (2.533,5.792) (2.577,5.853) (2.621,5.913) (2.666,5.973) (2.711,6.033) (2.756,6.093) (2.801,6.153) (2.847,6.213) (2.892,6.272) (2.938,6.331) (2.985,6.390) (3.031,6.449) (3.078,6.508) (3.125,6.566) (3.172,6.625) (3.219,6.683) (3.267,6.741) (3.315,6.799) (3.363,6.856) (3.411,6.914) (3.459,6.971) (3.508,7.028) (3.557,7.085) (3.606,7.142) (3.655,7.198) (3.705,7.255) (3.754,7.311) (3.804,7.367) (3.854,7.423) (3.905,7.478) (3.955,7.534) (4.006,7.589) (4.057,7.644) (4.108,7.699) (4.159,7.754) (4.210,7.808) (4.262,7.863) (4.314,7.917) (4.366,7.971) (4.418,8.025) (4.471,8.078) (4.523,8.132) (4.576,8.185) (4.629,8.238) (4.682,8.291) (4.735,8.344) (4.789,8.396) (4.843,8.449) (4.896,8.501) (4.950,8.553) (5.005,8.605) (5.059,8.656) (5.114,8.708) (5.168,8.759) (5.223,8.810) (5.278,8.861) (5.334,8.912) (5.389,8.962) (5.445,9.013) (5.500,9.063) (5.556,9.113) (5.613,9.163) (5.669,9.212) (5.725,9.262) (5.782,9.311) (5.839,9.360) (5.895,9.409) (5.953,9.457) (6.010,9.506) (6.067,9.554) (6.125,9.602) (6.183,9.650) (6.240,9.698) (6.298,9.745) (6.357,9.793) (6.415,9.840) (6.473,9.887) (6.532,9.933) (6.591,9.980) (6.630,10.000)}
\LC{(-3.797,3.606) (-3.869,3.626) (-3.942,3.645) (-4.014,3.663) (-4.087,3.680) (-4.161,3.696) (-4.234,3.712) (-4.307,3.726) (-4.381,3.740) (-4.455,3.753) (-4.529,3.765) (-4.603,3.775) (-4.678,3.785) (-4.752,3.794) (-4.827,3.802) (-4.901,3.810) (-4.976,3.816) (-5.051,3.821) (-5.126,3.825) (-5.201,3.829) (-5.276,3.831) (-5.351,3.832) (-5.426,3.833) (-5.501,3.832) (-5.576,3.831) (-5.651,3.828) (-5.726,3.825) (-5.800,3.820) (-5.875,3.815) (-5.950,3.808) (-6.025,3.801) (-6.099,3.792) (-6.173,3.782) (-6.248,3.772) (-6.322,3.760) (-6.396,3.748) (-6.470,3.734) (-6.543,3.719) (-6.616,3.703) (-6.689,3.687) (-6.762,3.669) (-6.835,3.650) (-6.907,3.630) (-6.979,3.609) (-7.051,3.587) (-7.122,3.564) (-7.193,3.540) (-7.264,3.514) (-7.334,3.488) (-7.404,3.461) (-7.473,3.432) (-7.542,3.402) (-7.611,3.372) (-7.679,3.340) (-7.746,3.307) (-7.813,3.273) (-7.879,3.239) (-7.945,3.202) (-8.010,3.165) (-8.075,3.127) (-8.139,3.088) (-8.202,3.048) (-8.265,3.006) (-8.326,2.964) (-8.387,2.920) (-8.448,2.876) (-8.507,2.830) (-8.566,2.784) (-8.624,2.736) (-8.681,2.687) (-8.737,2.637) (-8.792,2.587) (-8.846,2.535) (-8.900,2.482) (-8.952,2.428) (-9.003,2.373) (-9.054,2.318) (-9.103,2.261) (-9.151,2.203) (-9.198,2.145) (-9.243,2.085) (-9.288,2.025) (-9.331,1.964) (-9.373,1.902) (-9.414,1.839) (-9.454,1.775) (-9.492,1.711) (-9.529,1.645) (-9.564,1.579) (-9.598,1.513) (-9.631,1.445) (-9.662,1.377) (-9.692,1.308) (-9.720,1.239) (-9.747,1.169) (-9.772,1.098) (-9.796,1.027) (-9.818,0.955) (-9.839,0.883) (-9.858,0.810) (-9.875,0.737) (-9.891,0.664) (-9.905,0.590) (-9.918,0.517) (-9.928,0.442) (-9.938,0.368) (-9.945,0.293) (-9.951,0.218) (-9.955,0.144) (-9.957,0.069) (-9.958,-0.006) (-9.957,-0.081) (-9.954,-0.156) (-9.950,-0.231) (-9.944,-0.306) (-9.936,-0.381) (-9.927,-0.455) (-9.916,-0.529) (-9.903,-0.603) (-9.888,-0.677) (-9.872,-0.750) (-9.855,-0.823) (-9.836,-0.895) (-9.815,-0.967) (-9.792,-1.039) (-9.768,-1.110) (-9.743,-1.180) (-9.716,-1.250) (-9.687,-1.320) (-9.657,-1.388) (-9.626,-1.457) (-9.593,-1.524) (-9.558,-1.591) (-9.522,-1.656) (-9.485,-1.722) (-9.447,-1.786) (-9.407,-1.850) (-9.366,-1.912) (-9.324,-1.974) (-9.280,-2.035) (-9.236,-2.096) (-9.190,-2.155) (-9.143,-2.213) (-9.094,-2.271) (-9.045,-2.327) (-8.995,-2.383) (-8.943,-2.437) (-8.891,-2.491) (-8.837,-2.544) (-8.783,-2.595) (-8.728,-2.646) (-8.671,-2.695) (-8.614,-2.744) (-8.556,-2.792) (-8.497,-2.838) (-8.438,-2.883) (-8.377,-2.928) (-8.316,-2.971) (-8.254,-3.013) (-8.191,-3.055) (-8.128,-3.095) (-8.064,-3.134) (-7.999,-3.172) (-7.934,-3.209) (-7.868,-3.245) (-7.802,-3.279) (-7.735,-3.313) (-7.667,-3.346) (-7.599,-3.377) (-7.530,-3.408) (-7.461,-3.437) (-7.392,-3.465) (-7.322,-3.492) (-7.252,-3.519) (-7.181,-3.544) (-7.110,-3.568) (-7.039,-3.591) (-6.967,-3.613) (-6.895,-3.633) (-6.823,-3.653) (-6.750,-3.672) (-6.677,-3.690) (-6.604,-3.706) (-6.531,-3.722) (-6.457,-3.736) (-6.383,-3.750) (-6.309,-3.762) (-6.235,-3.774) (-6.161,-3.784) (-6.086,-3.794) (-6.012,-3.802) (-5.937,-3.809) (-5.863,-3.816) (-5.788,-3.821) (-5.713,-3.825) (-5.638,-3.829) (-5.563,-3.831) (-5.488,-3.832) (-5.413,-3.833) (-5.338,-3.832) (-5.263,-3.831) (-5.188,-3.828) (-5.113,-3.825) (-5.038,-3.820) (-4.964,-3.815) (-4.889,-3.808) (-4.814,-3.801) (-4.740,-3.793) (-4.665,-3.784) (-4.591,-3.774) (-4.517,-3.763) (-4.443,-3.751) (-4.369,-3.738) (-4.295,-3.724) (-4.221,-3.709) (-4.148,-3.694) (-4.075,-3.677) (-4.002,-3.660) (-3.929,-3.641) (-3.857,-3.622) (-3.809,-3.609)}
\LC{(-0.048,-10.000) (-0.048,-9.947) (-0.047,-9.872) (-0.046,-9.797) (-0.046,-9.722) (-0.045,-9.647) (-0.045,-9.572) (-0.044,-9.497) (-0.044,-9.422) (-0.043,-9.347) (-0.043,-9.272) (-0.042,-9.197) (-0.042,-9.122) (-0.041,-9.047) (-0.041,-8.972) (-0.040,-8.897) (-0.040,-8.822) (-0.039,-8.747) (-0.039,-8.672) (-0.038,-8.597) (-0.038,-8.522) (-0.037,-8.447) (-0.037,-8.372) (-0.036,-8.297) (-0.036,-8.222) (-0.035,-8.147) (-0.035,-8.072) (-0.034,-7.997) (-0.034,-7.922) (-0.033,-7.847) (-0.033,-7.772) (-0.032,-7.697) (-0.032,-7.622) (-0.031,-7.547) (-0.031,-7.472) (-0.030,-7.397) (-0.030,-7.322) (-0.030,-7.247) (-0.029,-7.172) (-0.029,-7.097) (-0.028,-7.022) (-0.028,-6.947) (-0.027,-6.872) (-0.027,-6.797) (-0.026,-6.722) (-0.026,-6.647) (-0.026,-6.572) (-0.025,-6.497) (-0.025,-6.422) (-0.024,-6.347) (-0.024,-6.272) (-0.023,-6.197) (-0.023,-6.122) (-0.023,-6.047) (-0.022,-5.972) (-0.022,-5.897) (-0.021,-5.822) (-0.021,-5.747) (-0.020,-5.672) (-0.020,-5.597) (-0.020,-5.522) (-0.019,-5.447) (-0.019,-5.372) (-0.018,-5.297) (-0.018,-5.222) (-0.018,-5.147) (-0.017,-5.072) (-0.017,-4.997) (-0.017,-4.922) (-0.016,-4.847) (-0.016,-4.772) (-0.015,-4.697) (-0.015,-4.622) (-0.015,-4.547) (-0.014,-4.472) (-0.014,-4.397) (-0.014,-4.322) (-0.013,-4.247) (-0.013,-4.172) (-0.013,-4.097) (-0.012,-4.022) (-0.012,-3.947) (-0.012,-3.872) (-0.011,-3.797) (-0.011,-3.722) (-0.011,-3.647) (-0.010,-3.572) (-0.010,-3.497) (-0.010,-3.422) (-0.009,-3.347) (-0.009,-3.272) (-0.009,-3.197) (-0.008,-3.122) (-0.008,-3.047) (-0.008,-2.972) (-0.007,-2.897) (-0.007,-2.822) (-0.007,-2.747) (-0.007,-2.672) (-0.006,-2.597) (-0.006,-2.522) (-0.006,-2.447) (-0.006,-2.372) (-0.005,-2.297)}
\LC{(-0.488,1.887) (-0.516,1.956) (-0.545,2.025) (-0.574,2.094) (-0.604,2.163) (-0.634,2.232) (-0.665,2.300) (-0.696,2.369) (-0.727,2.437) (-0.759,2.504) (-0.792,2.572) (-0.825,2.639) (-0.858,2.707) (-0.892,2.774) (-0.926,2.840) (-0.961,2.907) (-0.996,2.973) (-1.031,3.039) (-1.067,3.105) (-1.103,3.171) (-1.140,3.236) (-1.177,3.301) (-1.214,3.366) (-1.252,3.431) (-1.290,3.496) (-1.329,3.560) (-1.368,3.624) (-1.407,3.688) (-1.446,3.752) (-1.486,3.816) (-1.526,3.879) (-1.567,3.942) (-1.608,4.005) (-1.649,4.067) (-1.691,4.130) (-1.733,4.192) (-1.775,4.254) (-1.818,4.316) (-1.860,4.377) (-1.904,4.439) (-1.947,4.500) (-1.991,4.560) (-2.035,4.621) (-2.080,4.682) (-2.124,4.742) (-2.169,4.802) (-2.215,4.861) (-2.260,4.921) (-2.306,4.980) (-2.352,5.039) (-2.399,5.098) (-2.446,5.157) (-2.493,5.215) (-2.540,5.273) (-2.588,5.331) (-2.635,5.389) (-2.684,5.447) (-2.732,5.504) (-2.781,5.561) (-2.830,5.618) (-2.879,5.674) (-2.928,5.731) (-2.978,5.787) (-3.028,5.843) (-3.078,5.899) (-3.129,5.954) (-3.179,6.009) (-3.230,6.064) (-3.282,6.119) (-3.333,6.174) (-3.385,6.228) (-3.437,6.282) (-3.489,6.336) (-3.542,6.389) (-3.594,6.443) (-3.647,6.496) (-3.700,6.549) (-3.754,6.602) (-3.807,6.654) (-3.861,6.706) (-3.915,6.758) (-3.969,6.810) (-4.024,6.862) (-4.079,6.913) (-4.134,6.964) (-4.189,7.015) (-4.244,7.065) (-4.300,7.116) (-4.355,7.166) (-4.411,7.216) (-4.468,7.265) (-4.524,7.315) (-4.581,7.364) (-4.638,7.413) (-4.695,7.461) (-4.752,7.510) (-4.809,7.558) (-4.867,7.606) (-4.925,7.654) (-4.983,7.701) (-5.041,7.749) (-5.100,7.796) (-5.158,7.842) (-5.217,7.889) (-5.276,7.935) (-5.335,7.981) (-5.395,8.027) (-5.454,8.073) (-5.514,8.118) (-5.574,8.163) (-5.634,8.208) (-5.694,8.252) (-5.755,8.297) (-5.816,8.341) (-5.876,8.385) (-5.937,8.428) (-5.999,8.472) (-6.060,8.515) (-6.122,8.558) (-6.183,8.600) (-6.245,8.643) (-6.307,8.685) (-6.369,8.727) (-6.432,8.768) (-6.494,8.810) (-6.557,8.851) (-6.620,8.892) (-6.683,8.932) (-6.746,8.973) (-6.810,9.013) (-6.873,9.053) (-6.937,9.092) (-7.001,9.132) (-7.065,9.171) (-7.129,9.210) (-7.193,9.248) (-7.257,9.287) (-7.322,9.325) (-7.387,9.363) (-7.452,9.400) (-7.517,9.437) (-7.582,9.475) (-7.647,9.511) (-7.713,9.548) (-7.778,9.584) (-7.844,9.620) (-7.910,9.656) (-7.976,9.692) (-8.042,9.727) (-8.109,9.762) (-8.175,9.797) (-8.242,9.831) (-8.308,9.866) (-8.375,9.900) (-8.442,9.933) (-8.509,9.967) (-8.577,10.000) (-8.599,10.000)}
\LC{(-1.957,3.712) (-2.008,3.767) (-2.059,3.821) (-2.111,3.876) (-2.163,3.930) (-2.215,3.983) (-2.268,4.036) (-2.322,4.089) (-2.375,4.142) (-2.429,4.194) (-2.484,4.245) (-2.538,4.297) (-2.593,4.348) (-2.649,4.398) (-2.704,4.448) (-2.761,4.498) (-2.817,4.547) (-2.874,4.596) (-2.931,4.645) (-2.988,4.693) (-3.046,4.741) (-3.104,4.789) (-3.162,4.836) (-3.221,4.883) (-3.280,4.929) (-3.339,4.975) (-3.399,5.020) (-3.459,5.065) (-3.519,5.110) (-3.579,5.155) (-3.640,5.199) (-3.701,5.242) (-3.763,5.285) (-3.824,5.328) (-3.886,5.370) (-3.948,5.412) (-4.011,5.454) (-4.073,5.495) (-4.136,5.536) (-4.199,5.576) (-4.263,5.616) (-4.327,5.656) (-4.391,5.695) (-4.455,5.734) (-4.519,5.772) (-4.584,5.810) (-4.649,5.848) (-4.714,5.885) (-4.780,5.921) (-4.845,5.958) (-4.911,5.994) (-4.977,6.029) (-5.043,6.064) (-5.110,6.099) (-5.177,6.133) (-5.244,6.167) (-5.311,6.200) (-5.378,6.233) (-5.446,6.265) (-5.514,6.298) (-5.582,6.329) (-5.650,6.360) (-5.718,6.391) (-5.787,6.422) (-5.855,6.452) (-5.924,6.481) (-5.994,6.510) (-6.063,6.539) (-6.132,6.567) (-6.202,6.595) (-6.272,6.622) (-6.342,6.649) (-6.412,6.676) (-6.482,6.702) (-6.553,6.727) (-6.623,6.753) (-6.694,6.777) (-6.765,6.802) (-6.836,6.825) (-6.908,6.849) (-6.979,6.872) (-7.051,6.894) (-7.122,6.916) (-7.194,6.938) (-7.266,6.959) (-7.338,6.980) (-7.410,7.000) (-7.483,7.020) (-7.555,7.039) (-7.628,7.058) (-7.700,7.076) (-7.773,7.094) (-7.846,7.112) (-7.919,7.129) (-7.992,7.145) (-8.066,7.161) (-8.139,7.177) (-8.212,7.192) (-8.286,7.207) (-8.360,7.221) (-8.433,7.235) (-8.507,7.248) (-8.581,7.261) (-8.655,7.273) (-8.729,7.285) (-8.803,7.297) (-8.877,7.308) (-8.952,7.318) (-9.026,7.328) (-9.100,7.337) (-9.175,7.346) (-9.249,7.355) (-9.324,7.363) (-9.399,7.371) (-9.473,7.378) (-9.548,7.384) (-9.623,7.390) (-9.697,7.396) (-9.772,7.401) (-9.847,7.406) (-9.922,7.410) (-9.997,7.413) (-10.000,7.414)}
\LC{(-10.000,-8.150) (-9.943,-8.140) (-9.868,-8.130) (-9.794,-8.119) (-9.720,-8.107) (-9.646,-8.095) (-9.572,-8.083) (-9.498,-8.070) (-9.424,-8.057) (-9.350,-8.044) (-9.277,-8.030) (-9.203,-8.016) (-9.130,-8.001) (-9.056,-7.986) (-8.983,-7.970) (-8.910,-7.954) (-8.836,-7.938) (-8.763,-7.921) (-8.690,-7.903) (-8.617,-7.886) (-8.545,-7.868) (-8.472,-7.849) (-8.399,-7.830) (-8.327,-7.811) (-8.255,-7.791) (-8.182,-7.771) (-8.110,-7.750) (-8.038,-7.729) (-7.966,-7.708) (-7.895,-7.686) (-7.823,-7.664) (-7.751,-7.641) (-7.680,-7.618) (-7.609,-7.595) (-7.538,-7.571) (-7.467,-7.547) (-7.396,-7.522) (-7.325,-7.497) (-7.255,-7.472) (-7.184,-7.446) (-7.114,-7.420) (-7.044,-7.393) (-6.974,-7.366) (-6.904,-7.339) (-6.834,-7.311) (-6.765,-7.283) (-6.695,-7.255) (-6.626,-7.226) (-6.557,-7.196) (-6.488,-7.167) (-6.420,-7.137) (-6.351,-7.106) (-6.283,-7.075) (-6.214,-7.044) (-6.146,-7.012) (-6.079,-6.980) (-6.011,-6.948) (-5.944,-6.915) (-5.876,-6.882) (-5.809,-6.849) (-5.742,-6.815) (-5.676,-6.780) (-5.609,-6.746) (-5.543,-6.711) (-5.477,-6.675) (-5.411,-6.640) (-5.345,-6.603) (-5.279,-6.567) (-5.214,-6.530) (-5.149,-6.493) (-5.084,-6.455) (-5.020,-6.417) (-4.955,-6.379) (-4.891,-6.340) (-4.827,-6.301) (-4.763,-6.261) (-4.700,-6.222) (-4.636,-6.181) (-4.573,-6.141) (-4.510,-6.100) (-4.448,-6.059) (-4.385,-6.017) (-4.323,-5.975) (-4.261,-5.933) (-4.200,-5.890) (-4.138,-5.847) (-4.077,-5.803) (-4.016,-5.760) (-3.956,-5.715) (-3.895,-5.671) (-3.835,-5.626) (-3.775,-5.581) (-3.716,-5.535) (-3.656,-5.489) (-3.597,-5.443) (-3.539,-5.397) (-3.480,-5.350) (-3.422,-5.302) (-3.364,-5.255) (-3.306,-5.207) (-3.249,-5.158) (-3.192,-5.110) (-3.135,-5.061) (-3.079,-5.011) (-3.023,-4.961) (-2.967,-4.911) (-2.911,-4.861) (-2.856,-4.810) (-2.801,-4.759) (-2.747,-4.708) (-2.692,-4.656) (-2.638,-4.604) (-2.585,-4.551) (-2.531,-4.499) (-2.479,-4.445) (-2.426,-4.392) (-2.374,-4.338) (-2.322,-4.284) (-2.270,-4.230) (-2.219,-4.175) (-2.168,-4.120) (-2.118,-4.064) (-2.068,-4.008) (-2.018,-3.952) (-1.968,-3.896) (-1.919,-3.839) (-1.871,-3.782) (-1.823,-3.724) (-1.775,-3.667) (-1.727,-3.609) (-1.680,-3.550) (-1.634,-3.491) (-1.587,-3.432) (-1.542,-3.373) (-1.496,-3.313) (-1.451,-3.253) (-1.407,-3.193) (-1.363,-3.132) (-1.319,-3.071) (-1.276,-3.010) (-1.233,-2.948) (-1.191,-2.886) (-1.149,-2.824)}
\LC{(3.767,2.687) (3.842,2.687) (3.917,2.685) (3.992,2.683) (4.066,2.678) (4.141,2.673) (4.216,2.665) (4.290,2.657) (4.365,2.647) (4.439,2.635) (4.513,2.622) (4.586,2.608) (4.660,2.592) (4.732,2.574) (4.805,2.555) (4.877,2.535) (4.949,2.513) (5.020,2.489) (5.091,2.464) (5.161,2.438) (5.231,2.410) (5.299,2.380) (5.368,2.349) (5.435,2.316) (5.502,2.282) (5.568,2.246) (5.633,2.209) (5.697,2.170) (5.760,2.130) (5.823,2.088) (5.884,2.045) (5.944,2.000) (6.003,1.954) (6.061,1.906) (6.117,1.856) (6.172,1.806) (6.226,1.753) (6.279,1.700) (6.330,1.645) (6.379,1.588) (6.427,1.531) (6.473,1.472) (6.518,1.411) (6.560,1.350) (6.601,1.287) (6.641,1.223) (6.678,1.158) (6.713,1.092) (6.747,1.025) (6.778,0.956) (6.807,0.887) (6.834,0.817) (6.859,0.747) (6.882,0.675) (6.902,0.603) (6.920,0.530) (6.936,0.457) (6.950,0.383) (6.961,0.309) (6.970,0.235) (6.976,0.160) (6.980,0.085) (6.981,0.010) (6.980,-0.065) (6.977,-0.140) (6.971,-0.215) (6.963,-0.289) (6.953,-0.364) (6.940,-0.437) (6.925,-0.511) (6.907,-0.584) (6.887,-0.656) (6.865,-0.728) (6.841,-0.799) (6.814,-0.869) (6.786,-0.938) (6.755,-1.007) (6.722,-1.074) (6.687,-1.141) (6.651,-1.206) (6.612,-1.270) (6.571,-1.333) (6.529,-1.395) (6.485,-1.456) (6.439,-1.515) (6.392,-1.573) (6.343,-1.630) (6.292,-1.685) (6.240,-1.739) (6.187,-1.792) (6.132,-1.843) (6.076,-1.893) (6.018,-1.941) (5.960,-1.988) (5.900,-2.033) (5.839,-2.077) (5.777,-2.119) (5.714,-2.160) (5.650,-2.199) (5.585,-2.237) (5.519,-2.273) (5.453,-2.308) (5.386,-2.341) (5.318,-2.372) (5.249,-2.402) (5.179,-2.431) (5.109,-2.458) (5.039,-2.483) (4.968,-2.507) (4.896,-2.529) (4.824,-2.550) (4.752,-2.569) (4.679,-2.587) (4.606,-2.604) (4.532,-2.618) (4.458,-2.632) (4.384,-2.644) (4.310,-2.654) (4.236,-2.663) (4.161,-2.671) (4.086,-2.677) (4.011,-2.682) (3.936,-2.685) (3.861,-2.687) (3.786,-2.687) (3.761,-2.687)}
\LC{(10.000,-7.569) (9.949,-7.564) (9.875,-7.558) (9.800,-7.553) (9.725,-7.546) (9.650,-7.539) (9.576,-7.532) (9.501,-7.524) (9.427,-7.516) (9.352,-7.508) (9.278,-7.498) (9.203,-7.489) (9.129,-7.479) (9.055,-7.468) (8.980,-7.457) (8.906,-7.445) (8.832,-7.434) (8.758,-7.421) (8.684,-7.408) (8.611,-7.395) (8.537,-7.381) (8.463,-7.367) (8.390,-7.352) (8.316,-7.337) (8.243,-7.321) (8.170,-7.305) (8.097,-7.288) (8.024,-7.271) (7.951,-7.254) (7.878,-7.236) (7.805,-7.218) (7.733,-7.199) (7.660,-7.179) (7.588,-7.160) (7.515,-7.140) (7.443,-7.119) (7.371,-7.098) (7.299,-7.076) (7.228,-7.054) (7.156,-7.032) (7.085,-7.009) (7.013,-6.986) (6.942,-6.962) (6.871,-6.938) (6.800,-6.914) (6.730,-6.889) (6.659,-6.863) (6.589,-6.837) (6.518,-6.811) (6.448,-6.784) (6.378,-6.757) (6.309,-6.730) (6.239,-6.702) (6.170,-6.673) (6.100,-6.645) (6.031,-6.615) (5.963,-6.586) (5.894,-6.555) (5.825,-6.525) (5.757,-6.494) (5.689,-6.463) (5.621,-6.431) (5.553,-6.399) (5.486,-6.366) (5.418,-6.333) (5.351,-6.300) (5.284,-6.266) (5.218,-6.232) (5.151,-6.197) (5.085,-6.162) (5.019,-6.126) (4.953,-6.091) (4.887,-6.054) (4.822,-6.018) (4.757,-5.980) (4.692,-5.943) (4.627,-5.905) (4.562,-5.867) (4.498,-5.828) (4.434,-5.789) (4.370,-5.749) (4.307,-5.710) (4.244,-5.669) (4.181,-5.629) (4.139,-5.601)}
\LC{(3.620,3.580) (3.691,3.603) (3.763,3.624) (3.835,3.645) (3.908,3.665) (3.980,3.684) (4.053,3.703) (4.126,3.720) (4.199,3.736) (4.272,3.752) (4.346,3.767) (4.420,3.781) (4.494,3.793) (4.568,3.805) (4.642,3.816) (4.716,3.827) (4.790,3.836) (4.865,3.844) (4.940,3.851) (5.014,3.858) (5.089,3.863) (5.164,3.868) (5.239,3.871) (5.314,3.874) (5.389,3.876) (5.464,3.876) (5.539,3.876) (5.614,3.875) (5.689,3.873) (5.764,3.870) (5.839,3.865) (5.914,3.860) (5.988,3.854) (6.063,3.847) (6.137,3.839) (6.212,3.830) (6.286,3.820) (6.360,3.808) (6.434,3.796) (6.508,3.783) (6.582,3.769) (6.655,3.754) (6.728,3.737) (6.801,3.720) (6.874,3.702) (6.947,3.682) (7.019,3.662) (7.091,3.640) (7.162,3.618) (7.233,3.594) (7.304,3.570) (7.375,3.544) (7.445,3.517) (7.514,3.489) (7.584,3.460) (7.652,3.431) (7.721,3.400) (7.788,3.367) (7.856,3.334) (7.922,3.300) (7.989,3.265) (8.054,3.229) (8.119,3.191) (8.184,3.153) (8.247,3.113) (8.311,3.073) (8.373,3.031) (8.435,2.988) (8.496,2.945) (8.556,2.900) (8.615,2.854) (8.674,2.807) (8.731,2.759) (8.788,2.710) (8.844,2.660) (8.899,2.610) (8.953,2.558) (9.007,2.505) (9.059,2.451) (9.110,2.396) (9.160,2.340) (9.209,2.284) (9.257,2.226) (9.304,2.167) (9.350,2.108) (9.394,2.048) (9.438,1.986) (9.480,1.924) (9.520,1.861) (9.560,1.798) (9.598,1.733) (9.635,1.668) (9.671,1.602) (9.705,1.535) (9.738,1.468) (9.769,1.400) (9.799,1.331) (9.828,1.261) (9.855,1.191) (9.880,1.121) (9.904,1.050) (9.927,0.978) (9.947,0.906) (9.967,0.834) (9.984,0.761) (10.000,0.688)}
\LC{(2.538,5.010) (2.587,5.067) (2.637,5.123) (2.687,5.179) (2.737,5.234) (2.788,5.290) (2.839,5.345) (2.890,5.400) (2.941,5.454) (2.993,5.508) (3.045,5.563) (3.097,5.616) (3.150,5.670) (3.203,5.723) (3.256,5.776) (3.309,5.829) (3.363,5.882) (3.416,5.934) (3.470,5.986) (3.525,6.037) (3.579,6.089) (3.634,6.140) (3.689,6.191) (3.745,6.242) (3.800,6.292) (3.856,6.342) (3.912,6.392) (3.968,6.442) (4.025,6.491) (4.081,6.540) (4.138,6.589) (4.196,6.637) (4.253,6.685) (4.311,6.733) (4.369,6.781) (4.427,6.829) (4.485,6.876) (4.544,6.923) (4.602,6.969) (4.661,7.016) (4.720,7.062) (4.780,7.107) (4.840,7.153) (4.899,7.198) (4.959,7.243) (5.020,7.288) (5.080,7.332) (5.141,7.376) (5.202,7.420) (5.263,7.464) (5.324,7.507) (5.385,7.550) (5.447,7.593) (5.509,7.635) (5.571,7.677) (5.633,7.719) (5.695,7.761) (5.758,7.802) (5.821,7.843) (5.884,7.884) (5.947,7.925) (6.010,7.965) (6.074,8.005) (6.137,8.044) (6.201,8.084) (6.265,8.123) (6.329,8.162) (6.394,8.200) (6.458,8.238) (6.523,8.276) (6.588,8.314) (6.653,8.351) (6.718,8.388) (6.783,8.425) (6.849,8.462) (6.915,8.498) (6.980,8.534) (7.046,8.569) (7.113,8.605) (7.179,8.640) (7.245,8.674) (7.312,8.709) (7.379,8.743) (7.446,8.777) (7.513,8.810) (7.580,8.844) (7.647,8.877) (7.715,8.909) (7.782,8.942) (7.850,8.974) (7.918,9.005) (7.986,9.037) (8.054,9.068) (8.123,9.099) (8.191,9.130) (8.260,9.160) (8.329,9.190) (8.398,9.220) (8.467,9.249) (8.536,9.278) (8.605,9.307) (8.674,9.335) (8.744,9.363) (8.813,9.391) (8.883,9.419) (8.953,9.446) (9.023,9.473) (9.093,9.500) (9.163,9.526) (9.234,9.552) (9.304,9.578) (9.375,9.603) (9.446,9.628) (9.516,9.653) (9.587,9.677) (9.658,9.702) (9.729,9.725) (9.801,9.749) (9.872,9.772) (9.943,9.795) (10.000,9.818)}
\LC{(1.312,1.428) (1.381,1.456) (1.451,1.482) (1.523,1.506) (1.594,1.528) (1.667,1.548) (1.739,1.566) (1.813,1.582) (1.886,1.595) (1.961,1.607) (2.035,1.616) (2.110,1.623) (2.184,1.628) (2.259,1.631) (2.334,1.631) (2.409,1.629) (2.484,1.625) (2.559,1.618) (2.633,1.609) (2.708,1.597) (2.781,1.583) (2.854,1.567) (2.927,1.548) (2.999,1.527) (3.070,1.503) (3.140,1.477) (3.209,1.448) (3.278,1.417) (3.345,1.383) (3.410,1.346) (3.474,1.307) (3.537,1.266) (3.597,1.222) (3.656,1.176) (3.713,1.127) (3.768,1.075) (3.820,1.022) (3.870,0.966) (3.918,0.908) (3.963,0.848) (4.004,0.785) (4.043,0.721) (4.079,0.655) (4.111,0.588) (4.140,0.518) (4.165,0.448) (4.187,0.376) (4.205,0.303) (4.219,0.230) (4.229,0.155) (4.235,0.080) (4.238,0.006) (4.236,-0.069) (4.230,-0.144) (4.221,-0.219) (4.207,-0.292) (4.190,-0.365) (4.169,-0.437) (4.144,-0.508) (4.116,-0.577) (4.084,-0.645) (4.049,-0.712) (4.010,-0.776) (3.969,-0.839) (3.925,-0.899) (3.878,-0.958) (3.828,-1.014) (3.776,-1.068) (3.721,-1.119) (3.665,-1.169) (3.606,-1.215) (3.546,-1.260) (3.483,-1.301) (3.420,-1.341) (3.354,-1.378) (3.288,-1.412) (3.220,-1.444) (3.151,-1.473) (3.080,-1.499) (3.009,-1.524) (2.938,-1.545) (2.865,-1.564) (2.792,-1.581) (2.718,-1.595) (2.644,-1.607) (2.570,-1.617) (2.495,-1.624) (2.420,-1.628) (2.345,-1.631) (2.270,-1.631) (2.195,-1.629) (2.121,-1.624) (2.046,-1.617) (1.971,-1.608) (1.897,-1.597) (1.824,-1.584) (1.750,-1.568) (1.677,-1.551) (1.605,-1.531) (1.533,-1.510) (1.462,-1.486) (1.391,-1.460) (1.322,-1.432) (1.299,-1.423)}
\LC{(2.736,-10.000) (2.707,-9.935) (2.678,-9.866) (2.648,-9.797) (2.619,-9.728) (2.591,-9.659) (2.562,-9.589) (2.533,-9.520) (2.504,-9.451) (2.476,-9.382) (2.447,-9.312) (2.419,-9.243) (2.391,-9.173) (2.363,-9.104) (2.335,-9.034) (2.307,-8.964) (2.279,-8.895) (2.252,-8.825) (2.224,-8.755) (2.197,-8.685) (2.169,-8.616) (2.142,-8.546) (2.115,-8.476) (2.088,-8.406) (2.061,-8.336) (2.034,-8.266) (2.008,-8.196) (1.981,-8.125) (1.955,-8.055) (1.929,-7.985) (1.902,-7.915) (1.876,-7.844) (1.851,-7.774) (1.825,-7.704) (1.799,-7.633) (1.773,-7.563) (1.748,-7.492) (1.723,-7.421) (1.697,-7.351) (1.672,-7.280) (1.647,-7.209) (1.623,-7.139) (1.598,-7.068) (1.573,-6.997) (1.549,-6.926) (1.524,-6.855) (1.500,-6.784) (1.476,-6.713) (1.452,-6.642) (1.428,-6.571) (1.405,-6.500) (1.381,-6.429) (1.358,-6.357) (1.334,-6.286) (1.311,-6.215) (1.288,-6.143) (1.265,-6.072) (1.243,-6.000) (1.220,-5.929) (1.197,-5.857) (1.175,-5.786) (1.153,-5.714) (1.131,-5.642) (1.109,-5.571) (1.087,-5.499) (1.066,-5.427) (1.044,-5.355) (1.023,-5.283) (1.001,-5.211) (0.980,-5.139) (0.959,-5.067) (0.939,-4.995) (0.918,-4.923) (0.898,-4.851) (0.877,-4.779) (0.857,-4.707) (0.837,-4.634) (0.817,-4.562) (0.798,-4.490) (0.778,-4.417) (0.759,-4.345) (0.740,-4.272)}
\LC{(-3.539,5.679) (-3.596,5.728) (-3.653,5.776) (-3.711,5.824) (-3.768,5.872) (-3.826,5.920) (-3.885,5.967) (-3.943,6.014) (-4.002,6.061) (-4.061,6.107) (-4.120,6.153) (-4.179,6.199) (-4.239,6.244) (-4.299,6.289) (-4.359,6.334) (-4.420,6.378) (-4.480,6.422) (-4.541,6.466) (-4.602,6.509) (-4.664,6.553) (-4.725,6.595) (-4.787,6.638) (-4.849,6.680) (-4.911,6.722) (-4.974,6.764) (-5.037,6.805) (-5.099,6.846) (-5.162,6.886) (-5.226,6.926) (-5.289,6.966) (-5.353,7.006) (-5.417,7.045) (-5.481,7.084) (-5.545,7.123) (-5.610,7.161) (-5.674,7.199) (-5.739,7.237) (-5.804,7.274) (-5.870,7.311) (-5.935,7.347) (-6.001,7.384) (-6.067,7.420) (-6.133,7.455) (-6.199,7.490) (-6.265,7.525) (-6.332,7.560) (-6.398,7.594) (-6.465,7.628) (-6.532,7.662) (-6.600,7.695) (-6.667,7.728) (-6.735,7.761) (-6.802,7.793) (-6.870,7.825) (-6.938,7.856) (-7.006,7.887) (-7.075,7.918) (-7.143,7.949) (-7.212,7.979) (-7.281,8.009) (-7.350,8.038) (-7.419,8.067) (-7.488,8.096) (-7.558,8.125) (-7.627,8.153) (-7.697,8.180) (-7.767,8.208) (-7.837,8.235) (-7.907,8.261) (-7.977,8.288) (-8.047,8.314) (-8.118,8.339) (-8.188,8.365) (-8.259,8.389) (-8.330,8.414) (-8.401,8.438) (-8.472,8.462) (-8.543,8.485) (-8.615,8.509) (-8.686,8.531) (-8.758,8.554) (-8.829,8.576) (-8.901,8.597) (-8.973,8.619) (-9.045,8.640) (-9.117,8.660) (-9.190,8.680) (-9.262,8.700) (-9.334,8.720) (-9.407,8.739) (-9.479,8.757) (-9.552,8.776) (-9.625,8.794) (-9.698,8.811) (-9.771,8.829) (-9.844,8.845) (-9.917,8.862) (-9.990,8.878) (-10.000,8.883)}
\LC{(8.623,2.249) (8.673,2.193) (8.722,2.136) (8.769,2.078) (8.815,2.019) (8.860,1.959) (8.904,1.898) (8.947,1.836) (8.988,1.773) (9.028,1.710) (9.066,1.646) (9.103,1.580) (9.139,1.514) (9.173,1.448) (9.206,1.380) (9.237,1.312) (9.267,1.243) (9.295,1.174) (9.321,1.103) (9.346,1.033) (9.370,0.961) (9.391,0.890) (9.411,0.817) (9.430,0.745) (9.446,0.671) (9.461,0.598) (9.474,0.524) (9.486,0.450) (9.496,0.376) (9.504,0.301) (9.510,0.226) (9.514,0.151) (9.517,0.076) (9.518,0.001) (9.517,-0.073) (9.514,-0.148) (9.510,-0.223) (9.504,-0.298) (9.496,-0.373) (9.486,-0.447) (9.475,-0.521) (9.462,-0.595) (9.447,-0.668) (9.430,-0.742) (9.412,-0.814) (9.392,-0.887) (9.370,-0.958) (9.347,-1.030) (9.322,-1.101) (9.296,-1.171) (9.268,-1.240) (9.238,-1.309) (9.207,-1.377) (9.174,-1.445) (9.140,-1.512) (9.105,-1.578) (9.068,-1.643) (9.029,-1.707) (8.990,-1.771) (8.948,-1.834) (8.906,-1.896) (8.862,-1.956) (8.817,-2.016) (8.771,-2.075) (8.724,-2.134) (8.675,-2.191) (8.625,-2.247) (8.574,-2.302) (8.522,-2.356) (8.469,-2.409) (8.415,-2.461) (8.360,-2.512) (8.304,-2.562) (8.248,-2.611) (8.190,-2.659) (8.131,-2.705) (8.071,-2.751) (8.011,-2.795) (7.950,-2.839) (7.888,-2.881) (7.825,-2.922) (7.762,-2.962) (7.698,-3.001) (7.633,-3.039) (7.567,-3.076) (7.501,-3.111) (7.435,-3.146) (7.368,-3.179) (7.300,-3.211) (7.232,-3.242) (7.163,-3.272) (7.094,-3.301) (7.024,-3.328) (6.954,-3.355) (6.883,-3.380) (6.812,-3.405) (6.741,-3.428) (6.669,-3.450) (6.597,-3.471) (6.525,-3.491) (6.452,-3.509) (6.379,-3.527) (6.306,-3.544) (6.233,-3.559) (6.159,-3.574) (6.085,-3.587) (6.011,-3.599) (5.937,-3.610) (5.863,-3.620) (5.788,-3.629) (5.714,-3.637) (5.639,-3.644) (5.564,-3.650) (5.490,-3.655) (5.415,-3.658) (5.340,-3.661) (5.265,-3.663) (5.190,-3.663) (5.115,-3.663) (5.040,-3.662) (4.965,-3.659) (4.890,-3.656) (4.815,-3.651)}
\LC{(9.552,-10.000) (9.482,-9.976) (9.412,-9.949) (9.342,-9.921) (9.273,-9.894) (9.203,-9.866) (9.134,-9.837) (9.064,-9.809) (8.995,-9.780) (8.926,-9.751) (8.857,-9.722) (8.788,-9.692) (8.719,-9.662) (8.651,-9.632) (8.582,-9.601) (8.514,-9.571) (8.445,-9.539) (8.377,-9.508) (8.309,-9.476) (8.241,-9.444) (8.174,-9.412) (8.106,-9.380) (8.039,-9.347) (7.971,-9.314) (7.904,-9.281) (7.837,-9.247) (7.770,-9.213) (7.703,-9.179) (7.637,-9.144) (7.570,-9.110) (7.504,-9.075) (7.438,-9.039) (7.372,-9.004) (7.306,-8.968) (7.240,-8.932) (7.175,-8.895) (7.109,-8.859) (7.044,-8.822) (6.979,-8.785) (6.914,-8.747) (6.849,-8.709) (6.784,-8.671) (6.720,-8.633) (6.656,-8.594) (6.592,-8.556) (6.528,-8.516) (6.464,-8.477) (6.400,-8.437) (6.337,-8.398) (6.273,-8.357) (6.210,-8.317) (6.147,-8.276) (6.084,-8.235) (6.022,-8.194) (5.959,-8.152) (5.897,-8.110) (5.835,-8.068) (5.773,-8.026) (5.711,-7.983) (5.650,-7.941) (5.588,-7.897) (5.527,-7.854) (5.466,-7.810) (5.405,-7.766) (5.345,-7.722) (5.284,-7.678) (5.224,-7.633) (5.164,-7.588) (5.104,-7.543) (5.045,-7.497) (4.985,-7.452) (4.926,-7.406) (4.867,-7.359) (4.808,-7.313) (4.750,-7.266) (4.691,-7.219) (4.633,-7.172) (4.575,-7.124) (4.517,-7.076) (4.460,-7.028) (4.402,-6.980) (4.345,-6.931) (4.288,-6.882) (4.232,-6.833) (4.175,-6.784) (4.119,-6.734) (4.063,-6.684) (4.007,-6.634) (3.952,-6.584) (3.896,-6.533) (3.841,-6.482) (3.786,-6.431) (3.732,-6.380) (3.677,-6.328) (3.623,-6.276) (3.569,-6.224) (3.515,-6.172) (3.462,-6.119) (3.409,-6.066) (3.356,-6.013) (3.303,-5.960) (3.251,-5.906) (3.199,-5.852) (3.147,-5.798) (3.095,-5.744) (3.044,-5.689) (2.993,-5.634) (2.942,-5.579) (2.891,-5.524) (2.841,-5.468) (2.791,-5.412) (2.741,-5.356) (2.691,-5.300) (2.642,-5.243) (2.593,-5.187) (2.545,-5.130) (2.496,-5.072) (2.448,-5.015) (2.400,-4.957) (2.353,-4.899) (2.306,-4.841) (2.259,-4.782) (2.212,-4.723) (2.166,-4.664) (2.120,-4.605) (2.074,-4.546) (2.029,-4.486) (1.984,-4.426) (1.939,-4.366) (1.894,-4.305) (1.850,-4.245) (1.807,-4.184) (1.763,-4.123) (1.720,-4.061) (1.677,-4.000) (1.635,-3.938) (1.593,-3.876) (1.551,-3.813) (1.510,-3.751) (1.482,-3.709)}
\LC{(-10.000,-5.600) (-9.942,-5.616) (-9.869,-5.631) (-9.795,-5.646) (-9.721,-5.659) (-9.648,-5.673) (-9.574,-5.685) (-9.500,-5.697) (-9.425,-5.708) (-9.351,-5.719) (-9.277,-5.729) (-9.202,-5.738) (-9.128,-5.747) (-9.053,-5.754) (-8.979,-5.762) (-8.904,-5.768) (-8.829,-5.774) (-8.754,-5.779) (-8.679,-5.784) (-8.605,-5.788) (-8.530,-5.791) (-8.455,-5.793) (-8.380,-5.795) (-8.305,-5.796) (-8.230,-5.797) (-8.155,-5.797) (-8.080,-5.796) (-8.005,-5.795) (-7.930,-5.793) (-7.855,-5.790) (-7.780,-5.787) (-7.705,-5.783) (-7.630,-5.779) (-7.555,-5.773) (-7.480,-5.768) (-7.406,-5.761) (-7.331,-5.754) (-7.256,-5.747) (-7.182,-5.738) (-7.107,-5.729) (-7.033,-5.720) (-6.959,-5.710) (-6.885,-5.699) (-6.810,-5.688) (-6.736,-5.676) (-6.662,-5.663) (-6.589,-5.650) (-6.515,-5.636) (-6.441,-5.622) (-6.368,-5.607) (-6.294,-5.591) (-6.221,-5.575) (-6.148,-5.558) (-6.075,-5.541) (-6.002,-5.523) (-5.930,-5.504) (-5.857,-5.485) (-5.785,-5.465) (-5.713,-5.445) (-5.641,-5.424) (-5.569,-5.403) (-5.497,-5.380) (-5.426,-5.358) (-5.354,-5.335) (-5.283,-5.311) (-5.212,-5.287) (-5.141,-5.262) (-5.071,-5.236) (-5.001,-5.210) (-4.930,-5.184) (-4.860,-5.156) (-4.791,-5.129) (-4.721,-5.100) (-4.652,-5.072) (-4.583,-5.042) (-4.514,-5.012) (-4.446,-4.982) (-4.377,-4.951) (-4.309,-4.919) (-4.242,-4.887) (-4.174,-4.854) (-4.107,-4.821) (-4.040,-4.787) (-3.973,-4.753) (-3.907,-4.718) (-3.841,-4.683) (-3.775,-4.647) (-3.709,-4.611) (-3.644,-4.574) (-3.579,-4.536) (-3.514,-4.498) (-3.450,-4.460) (-3.386,-4.421) (-3.322,-4.381) (-3.259,-4.341) (-3.196,-4.301) (-3.133,-4.259) (-3.071,-4.218) (-3.009,-4.176) (-2.947,-4.133) (-2.886,-4.090) (-2.825,-4.046) (-2.764,-4.002) (-2.704,-3.957) (-2.644,-3.912) (-2.584,-3.867) (-2.525,-3.820) (-2.467,-3.774)}
\LC{(10.000,-4.876) (9.939,-4.901) (9.868,-4.925) (9.797,-4.949) (9.725,-4.972) (9.653,-4.993) (9.581,-5.015) (9.509,-5.035) (9.437,-5.055) (9.364,-5.073) (9.292,-5.092) (9.219,-5.109) (9.145,-5.125) (9.072,-5.141) (8.999,-5.156) (8.925,-5.171) (8.851,-5.184) (8.777,-5.197) (8.703,-5.209) (8.629,-5.220) (8.555,-5.231) (8.481,-5.241) (8.406,-5.250) (8.332,-5.258) (8.257,-5.266) (8.182,-5.273) (8.108,-5.279) (8.033,-5.285) (7.958,-5.289) (7.883,-5.293) (7.808,-5.297) (7.733,-5.299) (7.658,-5.301) (7.583,-5.302) (7.508,-5.303) (7.433,-5.302) (7.358,-5.301) (7.283,-5.300) (7.208,-5.297) (7.133,-5.294) (7.058,-5.290) (6.984,-5.286) (6.909,-5.281) (6.834,-5.275) (6.759,-5.268) (6.685,-5.261) (6.610,-5.253) (6.535,-5.245) (6.461,-5.236) (6.387,-5.226) (6.312,-5.215) (6.238,-5.204) (6.164,-5.192) (6.090,-5.179) (6.017,-5.166) (5.943,-5.152) (5.869,-5.137) (5.796,-5.122) (5.723,-5.106) (5.649,-5.090) (5.576,-5.073) (5.504,-5.055) (5.431,-5.036) (5.358,-5.017) (5.286,-4.998) (5.214,-4.977) (5.142,-4.956) (5.070,-4.934) (4.998,-4.912) (4.927,-4.889) (4.856,-4.866) (4.785,-4.842) (4.714,-4.817) (4.643,-4.791) (4.573,-4.765) (4.503,-4.739) (4.433,-4.711) (4.363,-4.684) (4.294,-4.655) (4.248,-4.636)}
\LC{(-0.445,2.886) (-0.462,2.959) (-0.479,3.032) (-0.497,3.105) (-0.515,3.178) (-0.533,3.251) (-0.551,3.323) (-0.570,3.396) (-0.588,3.469) (-0.607,3.541) (-0.626,3.614) (-0.645,3.686) (-0.665,3.759) (-0.684,3.831) (-0.704,3.904) (-0.724,3.976) (-0.744,4.048) (-0.765,4.120) (-0.785,4.192) (-0.806,4.265) (-0.827,4.337) (-0.848,4.409) (-0.869,4.481) (-0.891,4.552) (-0.912,4.624) (-0.934,4.696) (-0.956,4.768) (-0.978,4.839) (-1.000,4.911) (-1.023,4.983) (-1.045,5.054) (-1.068,5.126) (-1.091,5.197) (-1.114,5.268) (-1.137,5.340) (-1.160,5.411) (-1.184,5.482) (-1.208,5.553) (-1.231,5.624) (-1.255,5.695) (-1.280,5.766) (-1.304,5.837) (-1.328,5.908) (-1.353,5.979) (-1.378,6.050) (-1.403,6.121) (-1.428,6.191) (-1.453,6.262) (-1.478,6.333) (-1.504,6.403) (-1.529,6.474) (-1.555,6.544) (-1.581,6.615) (-1.607,6.685) (-1.633,6.755) (-1.659,6.825) (-1.686,6.896) (-1.712,6.966) (-1.739,7.036) (-1.766,7.106) (-1.793,7.176) (-1.820,7.246) (-1.847,7.316) (-1.875,7.385) (-1.902,7.455) (-1.930,7.525) (-1.958,7.595) (-1.986,7.664) (-2.014,7.734) (-2.042,7.803) (-2.070,7.873) (-2.099,7.942) (-2.127,8.011) (-2.156,8.081) (-2.185,8.150) (-2.214,8.219) (-2.243,8.288) (-2.272,8.357) (-2.301,8.426) (-2.331,8.495) (-2.360,8.564) (-2.390,8.633) (-2.420,8.702) (-2.450,8.771) (-2.480,8.840) (-2.510,8.908) (-2.540,8.977) (-2.570,9.045) (-2.601,9.114) (-2.632,9.182) (-2.662,9.251) (-2.693,9.319) (-2.724,9.387) (-2.755,9.456) (-2.786,9.524) (-2.818,9.592) (-2.849,9.660) (-2.881,9.728) (-2.912,9.796) (-2.944,9.864) (-2.976,9.932) (-3.008,10.000) (-3.019,10.000)}
\LC{(-10.000,-9.569) (-9.950,-9.548) (-9.878,-9.527) (-9.807,-9.505) (-9.735,-9.484) (-9.663,-9.461) (-9.592,-9.439) (-9.520,-9.416) (-9.449,-9.393) (-9.378,-9.369) (-9.307,-9.346) (-9.236,-9.321) (-9.165,-9.297) (-9.094,-9.272) (-9.023,-9.247) (-8.953,-9.222) (-8.882,-9.196) (-8.812,-9.170) (-8.742,-9.143) (-8.672,-9.117) (-8.602,-9.090) (-8.532,-9.062) (-8.462,-9.035) (-8.393,-9.007) (-8.323,-8.978) (-8.254,-8.950) (-8.185,-8.921) (-8.116,-8.892) (-8.047,-8.862) (-7.978,-8.832) (-7.909,-8.802) (-7.841,-8.771) (-7.772,-8.741) (-7.704,-8.709) (-7.636,-8.678) (-7.568,-8.646) (-7.500,-8.614) (-7.432,-8.582) (-7.365,-8.549) (-7.298,-8.516) (-7.230,-8.483) (-7.163,-8.449) (-7.096,-8.415) (-7.030,-8.381) (-6.963,-8.347) (-6.897,-8.312) (-6.830,-8.277) (-6.764,-8.241) (-6.698,-8.206) (-6.633,-8.170) (-6.567,-8.133) (-6.501,-8.097) (-6.436,-8.060) (-6.371,-8.023) (-6.306,-7.985) (-6.241,-7.947) (-6.177,-7.909) (-6.112,-7.871) (-6.048,-7.832) (-5.984,-7.793) (-5.920,-7.754) (-5.857,-7.714) (-5.793,-7.674) (-5.730,-7.634) (-5.667,-7.593) (-5.604,-7.553) (-5.541,-7.512) (-5.478,-7.470) (-5.416,-7.429) (-5.354,-7.387) (-5.292,-7.344) (-5.230,-7.302) (-5.168,-7.259) (-5.107,-7.216) (-5.046,-7.173) (-4.985,-7.129) (-4.924,-7.085) (-4.864,-7.041) (-4.803,-6.996) (-4.743,-6.951) (-4.683,-6.906) (-4.624,-6.861) (-4.564,-6.815) (-4.505,-6.769) (-4.446,-6.723) (-4.387,-6.676) (-4.328,-6.629) (-4.270,-6.582) (-4.212,-6.535) (-4.154,-6.487) (-4.096,-6.439) (-4.039,-6.391) (-3.982,-6.342) (-3.925,-6.294) (-3.868,-6.245) (-3.812,-6.195) (-3.755,-6.146) (-3.699,-6.096) (-3.644,-6.045)}
\LC{(-2.919,-10.000) (-2.888,-9.932) (-2.857,-9.864) (-2.826,-9.795) (-2.795,-9.727) (-2.764,-9.658) (-2.734,-9.590) (-2.703,-9.521) (-2.673,-9.453) (-2.642,-9.384) (-2.612,-9.316) (-2.582,-9.247) (-2.552,-9.178) (-2.523,-9.109) (-2.493,-9.040) (-2.463,-8.971) (-2.434,-8.902) (-2.405,-8.833) (-2.375,-8.764) (-2.346,-8.695) (-2.317,-8.626) (-2.288,-8.557) (-2.260,-8.488) (-2.231,-8.418) (-2.203,-8.349) (-2.174,-8.279) (-2.146,-8.210) (-2.118,-8.140) (-2.090,-8.071) (-2.062,-8.001) (-2.034,-7.932) (-2.006,-7.862) (-1.979,-7.792) (-1.951,-7.722) (-1.924,-7.652) (-1.897,-7.583) (-1.870,-7.513) (-1.843,-7.443) (-1.816,-7.373) (-1.790,-7.302) (-1.763,-7.232) (-1.737,-7.162) (-1.710,-7.092) (-1.684,-7.021) (-1.658,-6.951) (-1.632,-6.881) (-1.607,-6.810) (-1.581,-6.740) (-1.556,-6.669) (-1.530,-6.599) (-1.505,-6.528) (-1.480,-6.457) (-1.455,-6.387) (-1.431,-6.316) (-1.406,-6.245) (-1.381,-6.174) (-1.357,-6.103) (-1.333,-6.032) (-1.309,-5.961) (-1.285,-5.890) (-1.261,-5.819) (-1.238,-5.748) (-1.214,-5.676) (-1.191,-5.605) (-1.168,-5.534) (-1.145,-5.462) (-1.122,-5.391) (-1.099,-5.320) (-1.076,-5.248) (-1.054,-5.176) (-1.032,-5.105) (-1.010,-5.033) (-0.988,-4.961) (-0.966,-4.890) (-0.944,-4.818) (-0.923,-4.746) (-0.902,-4.674) (-0.880,-4.602) (-0.859,-4.530) (-0.839,-4.458) (-0.818,-4.386)}
\LC{(2.151,2.078) (2.223,2.099) (2.296,2.118) (2.368,2.135) (2.442,2.152) (2.515,2.166) (2.589,2.179) (2.663,2.191) (2.738,2.200) (2.812,2.209) (2.887,2.215) (2.962,2.220) (3.037,2.223) (3.112,2.225) (3.187,2.225) (3.262,2.223) (3.336,2.220) (3.411,2.215) (3.486,2.208) (3.560,2.199) (3.635,2.189) (3.709,2.176) (3.782,2.163) (3.856,2.147) (3.929,2.129) (4.001,2.110) (4.073,2.089) (4.144,2.066) (4.215,2.041) (4.285,2.014) (4.355,1.986) (4.423,1.955) (4.491,1.923) (4.558,1.889) (4.624,1.853) (4.688,1.815) (4.752,1.776) (4.815,1.734) (4.876,1.691) (4.936,1.646) (4.994,1.599) (5.051,1.550) (5.107,1.500) (5.160,1.447) (5.213,1.393) (5.263,1.338) (5.311,1.281) (5.358,1.222) (5.402,1.161) (5.445,1.099) (5.485,1.036) (5.523,0.971) (5.558,0.905) (5.592,0.838) (5.622,0.770) (5.651,0.700) (5.676,0.630) (5.699,0.558) (5.719,0.486) (5.736,0.413) (5.751,0.340) (5.763,0.266) (5.772,0.191) (5.778,0.116) (5.781,0.041) (5.781,-0.034) (5.778,-0.109) (5.772,-0.183) (5.764,-0.258) (5.752,-0.332) (5.738,-0.406) (5.721,-0.479) (5.701,-0.551) (5.679,-0.622) (5.653,-0.693) (5.625,-0.763) (5.595,-0.831) (5.562,-0.899) (5.527,-0.965) (5.489,-1.030) (5.449,-1.093) (5.407,-1.155) (5.363,-1.216) (5.316,-1.275) (5.268,-1.332) (5.218,-1.388) (5.166,-1.442) (5.112,-1.494) (5.057,-1.545) (5.000,-1.594) (4.942,-1.641) (4.882,-1.686) (4.821,-1.730) (4.759,-1.771) (4.695,-1.811) (4.630,-1.849) (4.565,-1.885) (4.498,-1.920) (4.430,-1.952) (4.362,-1.983) (4.293,-2.011) (4.222,-2.038) (4.152,-2.063) (4.080,-2.086) (4.009,-2.108) (3.936,-2.127) (3.863,-2.145) (3.790,-2.161) (3.716,-2.175) (3.642,-2.187) (3.568,-2.198) (3.494,-2.207) (3.419,-2.214) (3.344,-2.219) (3.269,-2.223) (3.194,-2.225) (3.119,-2.225) (3.044,-2.224) (2.969,-2.220) (2.895,-2.216) (2.820,-2.209) (2.745,-2.201) (2.671,-2.192) (2.597,-2.180) (2.523,-2.167) (2.449,-2.153) (2.376,-2.137) (2.303,-2.120) (2.230,-2.101) (2.158,-2.080)}
\LC{(-0.002,2.978) (-0.002,3.053) (-0.002,3.128) (-0.002,3.203) (-0.003,3.278) (-0.003,3.353) (-0.003,3.428) (-0.003,3.503) (-0.003,3.578) (-0.003,3.653) (-0.003,3.728) (-0.003,3.803) (-0.003,3.878) (-0.003,3.953) (-0.003,4.028) (-0.004,4.103) (-0.004,4.178) (-0.004,4.253) (-0.004,4.328) (-0.004,4.403) (-0.004,4.478) (-0.004,4.553) (-0.004,4.628) (-0.004,4.703) (-0.004,4.778) (-0.005,4.853) (-0.005,4.928) (-0.005,5.003) (-0.005,5.078) (-0.005,5.153) (-0.005,5.228) (-0.005,5.303) (-0.005,5.378) (-0.005,5.453) (-0.006,5.528) (-0.006,5.603) (-0.006,5.678) (-0.006,5.753) (-0.006,5.828) (-0.006,5.903) (-0.006,5.978) (-0.006,6.053) (-0.006,6.128) (-0.007,6.203) (-0.007,6.278) (-0.007,6.353) (-0.007,6.428) (-0.007,6.503) (-0.007,6.578) (-0.007,6.653) (-0.007,6.728) (-0.008,6.803) (-0.008,6.878) (-0.008,6.953) (-0.008,7.028) (-0.008,7.103) (-0.008,7.178) (-0.008,7.253) (-0.008,7.328) (-0.009,7.403) (-0.009,7.478) (-0.009,7.553) (-0.009,7.628) (-0.009,7.703) (-0.009,7.778) (-0.009,7.853) (-0.010,7.928) (-0.010,8.003) (-0.010,8.078) (-0.010,8.153) (-0.010,8.228) (-0.010,8.303) (-0.010,8.378) (-0.011,8.453) (-0.011,8.528) (-0.011,8.603) (-0.011,8.678) (-0.011,8.753) (-0.011,8.828) (-0.011,8.903) (-0.012,8.978) (-0.012,9.053) (-0.012,9.128) (-0.012,9.203) (-0.012,9.278) (-0.012,9.353) (-0.012,9.428) (-0.013,9.503) (-0.013,9.578) (-0.013,9.653) (-0.013,9.728) (-0.013,9.803) (-0.013,9.878) (-0.013,9.953) (-0.014,10.000)}
\LC{(4.675,4.277) (4.748,4.294) (4.821,4.311) (4.894,4.327) (4.967,4.342) (5.041,4.357) (5.115,4.370) (5.189,4.383) (5.263,4.395) (5.337,4.406) (5.411,4.417) (5.485,4.427) (5.560,4.435) (5.634,4.444) (5.709,4.451) (5.784,4.457) (5.859,4.463) (5.933,4.468) (6.008,4.472) (6.083,4.475) (6.158,4.477) (6.233,4.479) (6.308,4.480) (6.383,4.480) (6.458,4.479) (6.533,4.477) (6.608,4.474) (6.683,4.471) (6.758,4.466) (6.833,4.461) (6.907,4.455) (6.982,4.448) (7.057,4.440) (7.131,4.432) (7.206,4.422) (7.280,4.412) (7.354,4.401) (7.428,4.388) (7.502,4.375) (7.576,4.361) (7.649,4.347) (7.722,4.331) (7.796,4.314) (7.869,4.297) (7.941,4.278) (8.014,4.259) (8.086,4.239) (8.158,4.218) (8.230,4.196) (8.301,4.173) (8.372,4.149) (8.443,4.124) (8.513,4.099) (8.584,4.072) (8.653,4.045) (8.723,4.016) (8.792,3.987) (8.860,3.957) (8.929,3.926) (8.996,3.893) (9.064,3.860) (9.131,3.827) (9.197,3.792) (9.263,3.756) (9.328,3.719) (9.393,3.682) (9.458,3.643) (9.521,3.604) (9.585,3.563) (9.647,3.522) (9.709,3.480) (9.771,3.437) (9.831,3.393) (9.891,3.348) (9.951,3.302) (10.000,3.255)}
\LC{(-1.771,2.911) (-1.827,2.961) (-1.884,3.010) (-1.941,3.058) (-1.999,3.106) (-2.057,3.154) (-2.116,3.200) (-2.175,3.246) (-2.234,3.292) (-2.295,3.337) (-2.355,3.381) (-2.416,3.425) (-2.477,3.468) (-2.539,3.511) (-2.601,3.552) (-2.664,3.594) (-2.727,3.634) (-2.790,3.674) (-2.854,3.714) (-2.918,3.753) (-2.983,3.791) (-3.048,3.828) (-3.113,3.865) (-3.179,3.902) (-3.245,3.937) (-3.311,3.972) (-3.378,4.007) (-3.445,4.040) (-3.512,4.074) (-3.580,4.106) (-3.648,4.138) (-3.716,4.169) (-3.784,4.199) (-3.853,4.229) (-3.922,4.258) (-3.992,4.287) (-4.061,4.315) (-4.131,4.342) (-4.201,4.369) (-4.272,4.394) (-4.342,4.420) (-4.413,4.444) (-4.484,4.468) (-4.556,4.491) (-4.627,4.514) (-4.699,4.535) (-4.771,4.556) (-4.843,4.577) (-4.915,4.596) (-4.988,4.615) (-5.061,4.634) (-5.134,4.651) (-5.207,4.668) (-5.280,4.684) (-5.353,4.700) (-5.427,4.715) (-5.501,4.729) (-5.574,4.742) (-5.648,4.754) (-5.722,4.766) (-5.797,4.777) (-5.871,4.788) (-5.945,4.798) (-6.020,4.806) (-6.094,4.815) (-6.169,4.822) (-6.244,4.829) (-6.318,4.835) (-6.393,4.840) (-6.468,4.844) (-6.543,4.848) (-6.618,4.851) (-6.693,4.853) (-6.768,4.855) (-6.843,4.855) (-6.918,4.855) (-6.993,4.854) (-7.068,4.852) (-7.143,4.850) (-7.218,4.847) (-7.293,4.843) (-7.367,4.838) (-7.442,4.832) (-7.517,4.826) (-7.592,4.819) (-7.666,4.811) (-7.741,4.802) (-7.815,4.792) (-7.889,4.782) (-7.963,4.770) (-8.037,4.758) (-8.111,4.746) (-8.185,4.732) (-8.259,4.717) (-8.332,4.702) (-8.405,4.686) (-8.478,4.669) (-8.551,4.651) (-8.624,4.633) (-8.696,4.613) (-8.769,4.593) (-8.840,4.572) (-8.912,4.550) (-8.984,4.527) (-9.055,4.503) (-9.126,4.479) (-9.196,4.454) (-9.267,4.427) (-9.337,4.400) (-9.406,4.372) (-9.475,4.344) (-9.544,4.314) (-9.613,4.284) (-9.681,4.253) (-9.749,4.220) (-9.816,4.187) (-9.883,4.154) (-9.950,4.119) (-10.000,4.084)}
\LC{(2.995,4.844) (3.051,4.893) (3.108,4.942) (3.166,4.990) (3.223,5.038) (3.281,5.086) (3.339,5.133) (3.398,5.180) (3.456,5.227) (3.516,5.273) (3.575,5.319) (3.635,5.364) (3.694,5.410) (3.755,5.454) (3.815,5.499) (3.876,5.543) (3.937,5.586) (3.998,5.630) (4.060,5.673) (4.121,5.715) (4.183,5.757) (4.246,5.799) (4.308,5.840) (4.371,5.881) (4.434,5.922) (4.497,5.962) (4.561,6.002) (4.624,6.042) (4.688,6.081) (4.753,6.120) (4.817,6.158) (4.882,6.196) (4.947,6.234) (5.012,6.271) (5.077,6.308) (5.143,6.344) (5.208,6.380) (5.274,6.416) (5.341,6.451) (5.407,6.486) (5.474,6.521) (5.540,6.555) (5.607,6.588) (5.675,6.622) (5.742,6.655) (5.809,6.687) (5.877,6.719) (5.945,6.751) (6.013,6.782) (6.082,6.813) (6.150,6.844) (6.219,6.874) (6.288,6.904) (6.357,6.933) (6.426,6.962) (6.495,6.991) (6.565,7.019) (6.634,7.046) (6.704,7.074) (6.774,7.101) (6.845,7.127) (6.915,7.153) (6.985,7.179) (7.056,7.204) (7.127,7.229) (7.198,7.253) (7.269,7.277) (7.340,7.301) (7.411,7.324) (7.483,7.347) (7.554,7.369) (7.626,7.391) (7.698,7.413) (7.770,7.434) (7.842,7.454) (7.914,7.474) (7.987,7.494) (8.059,7.514) (8.132,7.533) (8.204,7.551) (8.277,7.569) (8.350,7.587) (8.423,7.604) (8.496,7.621) (8.569,7.637) (8.642,7.653) (8.716,7.668) (8.789,7.683) (8.863,7.698) (8.937,7.712) (9.010,7.726) (9.084,7.739) (9.158,7.752) (9.232,7.764) (9.306,7.776) (9.380,7.788) (9.454,7.799) (9.529,7.810) (9.603,7.820) (9.677,7.829) (9.752,7.839) (9.826,7.847) (9.901,7.856) (9.975,7.864) (10.000,7.866)}
\LC{(-1.118,1.217) (-1.188,1.245) (-1.258,1.270) (-1.330,1.293) (-1.402,1.314) (-1.474,1.332) (-1.548,1.348) (-1.622,1.362) (-1.696,1.372) (-1.770,1.380) (-1.845,1.386) (-1.920,1.389) (-1.995,1.389) (-2.070,1.387) (-2.145,1.382) (-2.219,1.374) (-2.294,1.363) (-2.367,1.349) (-2.440,1.333) (-2.513,1.313) (-2.585,1.291) (-2.655,1.266) (-2.725,1.238) (-2.793,1.207) (-2.860,1.172) (-2.925,1.135) (-2.988,1.095) (-3.050,1.052) (-3.109,1.006) (-3.166,0.958) (-3.220,0.906) (-3.272,0.852) (-3.321,0.795) (-3.366,0.735) (-3.409,0.673) (-3.447,0.609) (-3.482,0.543) (-3.514,0.474) (-3.541,0.405) (-3.563,0.333) (-3.582,0.260) (-3.596,0.187) (-3.605,0.112) (-3.609,0.037) (-3.609,-0.038) (-3.605,-0.112) (-3.596,-0.187) (-3.582,-0.261) (-3.563,-0.333) (-3.540,-0.405) (-3.513,-0.475) (-3.482,-0.543) (-3.447,-0.609) (-3.409,-0.673) (-3.366,-0.735) (-3.321,-0.795) (-3.272,-0.852) (-3.220,-0.906) (-3.166,-0.958) (-3.109,-1.006) (-3.049,-1.052) (-2.988,-1.095) (-2.925,-1.135) (-2.859,-1.172) (-2.793,-1.207) (-2.724,-1.238) (-2.655,-1.266) (-2.584,-1.291) (-2.513,-1.313) (-2.440,-1.333) (-2.367,-1.349) (-2.293,-1.363) (-2.219,-1.374) (-2.145,-1.382) (-2.070,-1.387) (-1.995,-1.389) (-1.920,-1.389) (-1.845,-1.386) (-1.770,-1.380) (-1.696,-1.372) (-1.621,-1.361) (-1.548,-1.348) (-1.474,-1.332) (-1.401,-1.314) (-1.329,-1.293) (-1.258,-1.270) (-1.188,-1.245) (-1.118,-1.217)}
\LC{(-10.000,-2.969) (-9.946,-3.019) (-9.890,-3.069) (-9.833,-3.118) (-9.775,-3.165) (-9.716,-3.212) (-9.657,-3.258) (-9.597,-3.303) (-9.537,-3.347) (-9.475,-3.391) (-9.413,-3.433) (-9.351,-3.474) (-9.287,-3.515) (-9.224,-3.554) (-9.159,-3.593) (-9.094,-3.630) (-9.029,-3.667) (-8.963,-3.702) (-8.897,-3.737) (-8.830,-3.771) (-8.762,-3.804) (-8.694,-3.836) (-8.626,-3.867) (-8.557,-3.897) (-8.488,-3.926) (-8.419,-3.954) (-8.349,-3.981) (-8.278,-4.007) (-8.208,-4.033) (-8.137,-4.057) (-8.066,-4.080) (-7.994,-4.103) (-7.922,-4.124) (-7.850,-4.145) (-7.778,-4.165) (-7.705,-4.184) (-7.632,-4.201) (-7.559,-4.218) (-7.486,-4.234) (-7.413,-4.249) (-7.339,-4.264) (-7.265,-4.277) (-7.191,-4.289) (-7.117,-4.301) (-7.043,-4.311) (-6.968,-4.321) (-6.894,-4.330) (-6.819,-4.338) (-6.745,-4.345) (-6.670,-4.351) (-6.595,-4.356) (-6.520,-4.361) (-6.445,-4.364) (-6.370,-4.367) (-6.295,-4.368) (-6.220,-4.369) (-6.145,-4.369) (-6.070,-4.369) (-5.995,-4.367) (-5.920,-4.364) (-5.846,-4.361) (-5.771,-4.357) (-5.696,-4.352) (-5.621,-4.346) (-5.546,-4.339) (-5.472,-4.332) (-5.397,-4.323) (-5.323,-4.314) (-5.248,-4.304) (-5.174,-4.293) (-5.100,-4.282) (-5.051,-4.273)}
\LC{(-2.119,1.838) (-2.193,1.851) (-2.267,1.863) (-2.341,1.872) (-2.416,1.880) (-2.491,1.886) (-2.566,1.890) (-2.641,1.892) (-2.716,1.892) (-2.791,1.890) (-2.866,1.885) (-2.940,1.879) (-3.015,1.871) (-3.089,1.861) (-3.163,1.849) (-3.237,1.834) (-3.310,1.818) (-3.383,1.799) (-3.455,1.778) (-3.526,1.756) (-3.597,1.730) (-3.667,1.703) (-3.736,1.674) (-3.804,1.642) (-3.871,1.608) (-3.936,1.573) (-4.001,1.534) (-4.064,1.494) (-4.126,1.452) (-4.186,1.407) (-4.245,1.360) (-4.302,1.311) (-4.357,1.261) (-4.410,1.208) (-4.461,1.153) (-4.510,1.096) (-4.557,1.037) (-4.601,0.977) (-4.643,0.914) (-4.682,0.850) (-4.718,0.785) (-4.752,0.718) (-4.782,0.649) (-4.810,0.580) (-4.835,0.509) (-4.856,0.437) (-4.874,0.364) (-4.889,0.291) (-4.901,0.216) (-4.909,0.142) (-4.914,0.067) (-4.915,-0.008) (-4.913,-0.083) (-4.908,-0.158) (-4.899,-0.232) (-4.887,-0.306) (-4.871,-0.379) (-4.852,-0.452) (-4.830,-0.524) (-4.805,-0.594) (-4.776,-0.664) (-4.745,-0.732) (-4.711,-0.799) (-4.674,-0.864) (-4.634,-0.927) (-4.592,-0.989) (-4.547,-1.050) (-4.500,-1.108) (-4.451,-1.164) (-4.399,-1.219) (-4.346,-1.271) (-4.290,-1.322) (-4.233,-1.370) (-4.174,-1.417) (-4.113,-1.461) (-4.051,-1.503) (-3.987,-1.543) (-3.923,-1.580) (-3.857,-1.616) (-3.789,-1.649) (-3.721,-1.680) (-3.652,-1.709) (-3.582,-1.736) (-3.511,-1.761) (-3.440,-1.783) (-3.367,-1.803) (-3.295,-1.821) (-3.221,-1.838) (-3.148,-1.851) (-3.074,-1.863) (-2.999,-1.873) (-2.925,-1.881) (-2.850,-1.887) (-2.775,-1.890) (-2.700,-1.892) (-2.625,-1.891) (-2.550,-1.889) (-2.475,-1.885) (-2.400,-1.879) (-2.326,-1.871) (-2.251,-1.861) (-2.177,-1.849) (-2.128,-1.840)}
\LC{(-4.280,5.092) (-4.347,5.126) (-4.414,5.159) (-4.482,5.192) (-4.549,5.225) (-4.617,5.257) (-4.685,5.288) (-4.754,5.319) (-4.822,5.349) (-4.891,5.379) (-4.960,5.409) (-5.029,5.437) (-5.099,5.466) (-5.168,5.494) (-5.238,5.521) (-5.308,5.548) (-5.379,5.574) (-5.449,5.600) (-5.520,5.625) (-5.590,5.650) (-5.661,5.674) (-5.733,5.697) (-5.804,5.720) (-5.876,5.743) (-5.947,5.765) (-6.019,5.786) (-6.091,5.807) (-6.163,5.827) (-6.236,5.847) (-6.308,5.866) (-6.381,5.885) (-6.454,5.903) (-6.527,5.921) (-6.600,5.938) (-6.673,5.954) (-6.746,5.970) (-6.819,5.986) (-6.893,6.001) (-6.967,6.015) (-7.040,6.028) (-7.114,6.042) (-7.188,6.054) (-7.262,6.066) (-7.336,6.078) (-7.410,6.088) (-7.485,6.099) (-7.559,6.108) (-7.634,6.117) (-7.708,6.126) (-7.783,6.134) (-7.857,6.141) (-7.932,6.148) (-8.007,6.154) (-8.082,6.160) (-8.156,6.165) (-8.231,6.169) (-8.306,6.173) (-8.381,6.176) (-8.456,6.179) (-8.531,6.181) (-8.606,6.182) (-8.681,6.183) (-8.756,6.183) (-8.831,6.183) (-8.906,6.182) (-8.981,6.180) (-9.056,6.178) (-9.131,6.175) (-9.206,6.172) (-9.281,6.168) (-9.356,6.163) (-9.430,6.158) (-9.505,6.152) (-9.580,6.145) (-9.654,6.138) (-9.729,6.130) (-9.804,6.122) (-9.878,6.113) (-9.952,6.103) (-10.000,6.096)}
\LC{(1.461,2.722) (1.513,2.776) (1.565,2.830) (1.618,2.883) (1.671,2.936) (1.725,2.988) (1.779,3.040) (1.834,3.091) (1.890,3.142) (1.945,3.192) (2.002,3.242) (2.058,3.291) (2.115,3.339) (2.173,3.387) (2.231,3.435) (2.290,3.482) (2.349,3.528) (2.408,3.574) (2.468,3.619) (2.528,3.664) (2.588,3.708) (2.649,3.752) (2.711,3.795) (2.773,3.837) (2.835,3.879) (2.897,3.921) (2.960,3.962) (3.024,4.002) (3.087,4.042) (3.151,4.081) (3.215,4.119) (3.280,4.157) (3.345,4.195) (3.410,4.232) (3.476,4.268) (3.542,4.304) (3.608,4.339) (3.675,4.373) (3.742,4.407) (3.809,4.441) (3.876,4.473) (3.944,4.506) (4.012,4.537) (4.080,4.568) (4.149,4.599) (4.218,4.628) (4.287,4.658) (4.356,4.686) (4.426,4.714) (4.496,4.741) (4.566,4.768) (4.636,4.794) (4.706,4.820) (4.777,4.845) (4.848,4.869) (4.919,4.893) (4.991,4.916) (5.062,4.938) (5.134,4.960) (5.206,4.981) (5.278,5.002) (5.350,5.022) (5.423,5.041) (5.496,5.060) (5.568,5.078) (5.641,5.095) (5.714,5.112) (5.788,5.128) (5.861,5.143) (5.935,5.158) (6.008,5.172) (6.082,5.185) (6.156,5.198) (6.230,5.210) (6.304,5.222) (6.378,5.233) (6.453,5.243) (6.527,5.252) (6.602,5.261) (6.676,5.269) (6.751,5.276) (6.825,5.283) (6.900,5.289) (6.975,5.294) (7.050,5.299) (7.125,5.303) (7.200,5.306) (7.275,5.309) (7.350,5.311) (7.425,5.312) (7.500,5.312) (7.575,5.312) (7.650,5.311) (7.725,5.310) (7.800,5.307) (7.874,5.304) (7.949,5.300) (8.024,5.296) (8.099,5.291) (8.174,5.285) (8.249,5.278) (8.323,5.271) (8.398,5.262) (8.472,5.254) (8.547,5.244) (8.621,5.234) (8.695,5.223) (8.769,5.211) (8.843,5.198) (8.917,5.185) (8.990,5.171) (9.064,5.156) (9.137,5.140) (9.211,5.124) (9.284,5.107) (9.356,5.089) (9.429,5.071) (9.502,5.051) (9.574,5.031) (9.646,5.010) (9.718,4.989) (9.789,4.966) (9.861,4.943) (9.932,4.919) (10.000,4.894)}
\LC{(1.369,4.863) (1.399,4.932) (1.429,5.000) (1.460,5.069) (1.491,5.137) (1.522,5.205) (1.553,5.273) (1.585,5.341) (1.617,5.409) (1.649,5.477) (1.681,5.545) (1.713,5.613) (1.746,5.680) (1.779,5.748) (1.812,5.815) (1.845,5.882) (1.878,5.949) (1.912,6.017) (1.945,6.084) (1.979,6.150) (2.013,6.217) (2.048,6.284) (2.082,6.350) (2.117,6.417) (2.152,6.483) (2.187,6.550) (2.222,6.616) (2.258,6.682) (2.293,6.748) (2.329,6.814) (2.365,6.880) (2.401,6.945) (2.437,7.011) (2.474,7.077) (2.510,7.142) (2.547,7.207) (2.584,7.273) (2.621,7.338) (2.659,7.403) (2.696,7.468) (2.734,7.533) (2.772,7.597) (2.810,7.662) (2.848,7.727) (2.886,7.791) (2.925,7.855) (2.964,7.920) (3.002,7.984) (3.041,8.048) (3.081,8.112) (3.120,8.176) (3.159,8.239) (3.199,8.303) (3.239,8.367) (3.279,8.430) (3.319,8.494) (3.359,8.557) (3.399,8.620) (3.440,8.683) (3.481,8.746) (3.522,8.809) (3.563,8.872) (3.604,8.934) (3.645,8.997) (3.687,9.059) (3.728,9.122) (3.770,9.184) (3.812,9.246) (3.854,9.308) (3.896,9.370) (3.939,9.432) (3.981,9.494) (4.024,9.556) (4.067,9.617) (4.110,9.679) (4.153,9.740) (4.196,9.802) (4.239,9.863) (4.283,9.924) (4.326,9.985) (4.341,10.000)}
\LC{(-1.407,4.702) (-1.439,4.770) (-1.472,4.837) (-1.504,4.905) (-1.537,4.973) (-1.570,5.040) (-1.603,5.107) (-1.636,5.174) (-1.670,5.241) (-1.704,5.308) (-1.738,5.375) (-1.772,5.442) (-1.806,5.508) (-1.841,5.575) (-1.876,5.641) (-1.911,5.708) (-1.947,5.774) (-1.982,5.840) (-2.018,5.906) (-2.054,5.972) (-2.090,6.037) (-2.126,6.103) (-2.163,6.168) (-2.200,6.234) (-2.237,6.299) (-2.274,6.364) (-2.311,6.429) (-2.349,6.494) (-2.386,6.559) (-2.424,6.624) (-2.463,6.688) (-2.501,6.753) (-2.539,6.817) (-2.578,6.881) (-2.617,6.945) (-2.656,7.009) (-2.695,7.073) (-2.735,7.137) (-2.774,7.201) (-2.814,7.264) (-2.854,7.328) (-2.894,7.391) (-2.935,7.454) (-2.975,7.517) (-3.016,7.580) (-3.057,7.643) (-3.098,7.706) (-3.139,7.769) (-3.180,7.831) (-3.222,7.894) (-3.264,7.956) (-3.306,8.018) (-3.348,8.080) (-3.390,8.142) (-3.432,8.204) (-3.475,8.266) (-3.518,8.327) (-3.561,8.389) (-3.604,8.450) (-3.647,8.512) (-3.690,8.573) (-3.734,8.634) (-3.778,8.695) (-3.822,8.756) (-3.866,8.816) (-3.910,8.877) (-3.954,8.937) (-3.999,8.998) (-4.043,9.058) (-4.088,9.118) (-4.133,9.178) (-4.178,9.238) (-4.224,9.298) (-4.269,9.357) (-4.315,9.417) (-4.360,9.476) (-4.406,9.536) (-4.452,9.595) (-4.499,9.654) (-4.545,9.713) (-4.591,9.772) (-4.638,9.831) (-4.685,9.889) (-4.732,9.948) (-4.779,10.000)}
\LC{(0.449,3.419) (0.464,3.492) (0.479,3.566) (0.494,3.639) (0.509,3.713) (0.524,3.786) (0.540,3.860) (0.556,3.933) (0.571,4.006) (0.587,4.080) (0.603,4.153) (0.619,4.226) (0.636,4.299) (0.652,4.373) (0.669,4.446) (0.686,4.519) (0.703,4.592) (0.720,4.665) (0.737,4.738) (0.754,4.811) (0.772,4.884) (0.789,4.957) (0.807,5.030) (0.825,5.102) (0.843,5.175) (0.861,5.248) (0.879,5.321) (0.897,5.394) (0.916,5.466) (0.934,5.539) (0.953,5.612) (0.972,5.684) (0.991,5.757) (1.010,5.829) (1.029,5.902) (1.048,5.974) (1.068,6.047) (1.087,6.119) (1.107,6.191) (1.127,6.264) (1.147,6.336) (1.167,6.408) (1.187,6.481) (1.207,6.553) (1.227,6.625) (1.248,6.697) (1.268,6.769) (1.289,6.841) (1.310,6.913) (1.331,6.986) (1.352,7.057) (1.373,7.129) (1.394,7.201) (1.415,7.273) (1.437,7.345) (1.458,7.417) (1.480,7.489) (1.502,7.561) (1.524,7.632) (1.546,7.704) (1.568,7.776) (1.590,7.847) (1.612,7.919) (1.635,7.991) (1.657,8.062) (1.680,8.134) (1.702,8.205) (1.725,8.277) (1.748,8.348) (1.771,8.419) (1.794,8.491) (1.817,8.562) (1.841,8.633) (1.864,8.705) (1.887,8.776) (1.911,8.847) (1.935,8.918) (1.958,8.989) (1.982,9.060) (2.006,9.132) (2.030,9.203) (2.054,9.274) (2.079,9.345) (2.103,9.416) (2.127,9.486) (2.152,9.557) (2.176,9.628) (2.201,9.699) (2.226,9.770) (2.251,9.841) (2.276,9.911) (2.301,9.982) (2.309,10.000)}
\LC{(10.000,-9.013) (9.929,-8.996) (9.856,-8.978) (9.784,-8.960) (9.711,-8.942) (9.638,-8.924) (9.565,-8.905) (9.493,-8.886) (9.421,-8.866) (9.348,-8.846) (9.276,-8.826) (9.204,-8.805) (9.132,-8.784) (9.060,-8.763) (8.988,-8.741) (8.917,-8.719) (8.845,-8.696) (8.774,-8.674) (8.702,-8.650) (8.631,-8.627) (8.560,-8.603) (8.489,-8.579) (8.418,-8.554) (8.347,-8.529) (8.277,-8.504) (8.206,-8.478) (8.136,-8.453) (8.066,-8.426) (7.996,-8.400) (7.926,-8.373) (7.856,-8.345) (7.786,-8.318) (7.717,-8.289) (7.647,-8.261) (7.578,-8.232) (7.509,-8.203) (7.440,-8.174) (7.371,-8.144) (7.302,-8.114) (7.234,-8.084) (7.165,-8.053) (7.097,-8.022) (7.029,-7.990) (6.961,-7.959) (6.893,-7.927) (6.825,-7.894) (6.758,-7.861) (6.691,-7.828) (6.624,-7.795) (6.557,-7.761) (6.490,-7.727) (6.423,-7.692) (6.357,-7.658) (6.290,-7.623) (6.224,-7.587) (6.158,-7.551) (6.093,-7.515) (6.027,-7.479) (5.962,-7.442) (5.897,-7.405) (5.832,-7.368) (5.767,-7.330) (5.702,-7.292) (5.638,-7.254) (5.573,-7.215) (5.509,-7.176) (5.445,-7.137) (5.382,-7.097) (5.318,-7.057) (5.255,-7.017) (5.192,-6.976) (5.129,-6.935) (5.067,-6.894) (5.004,-6.852) (4.942,-6.810) (4.880,-6.768) (4.818,-6.726) (4.757,-6.683) (4.716,-6.654)}
\LC{(4.277,5.274) (4.342,5.310) (4.408,5.346) (4.475,5.381) (4.541,5.415) (4.608,5.449) (4.675,5.483) (4.742,5.516) (4.810,5.549) (4.878,5.581) (4.946,5.613) (5.014,5.644) (5.082,5.675) (5.151,5.705) (5.220,5.735) (5.289,5.764) (5.358,5.793) (5.427,5.822) (5.497,5.850) (5.567,5.877) (5.637,5.904) (5.707,5.930) (5.777,5.956) (5.848,5.982) (5.919,6.006) (5.989,6.031) (6.061,6.055) (6.132,6.078) (6.203,6.101) (6.275,6.123) (6.347,6.145) (6.418,6.167) (6.490,6.188) (6.563,6.208) (6.635,6.228) (6.707,6.247) (6.780,6.266) (6.853,6.284) (6.926,6.302) (6.999,6.319) (7.072,6.336) (7.145,6.352) (7.218,6.368) (7.292,6.383) (7.365,6.398) (7.439,6.412) (7.513,6.425) (7.587,6.439) (7.661,6.451) (7.735,6.463) (7.809,6.474) (7.883,6.485) (7.957,6.496) (8.032,6.506) (8.106,6.515) (8.180,6.524) (8.255,6.532) (8.330,6.539) (8.404,6.547) (8.479,6.553) (8.554,6.559) (8.629,6.565) (8.703,6.569) (8.778,6.574) (8.853,6.578) (8.928,6.581) (9.003,6.583) (9.078,6.586) (9.153,6.587) (9.228,6.588) (9.303,6.588) (9.378,6.588) (9.453,6.587) (9.528,6.586) (9.603,6.584) (9.678,6.582) (9.753,6.579) (9.828,6.575) (9.903,6.571) (9.977,6.566) (10.000,6.564)}
\LC{(0.444,6.194) (0.452,6.268) (0.460,6.343) (0.468,6.418) (0.477,6.492) (0.485,6.567) (0.493,6.641) (0.502,6.716) (0.510,6.790) (0.518,6.865) (0.527,6.939) (0.535,7.014) (0.544,7.088) (0.553,7.163) (0.561,7.237) (0.570,7.312) (0.579,7.386) (0.588,7.461) (0.596,7.535) (0.605,7.610) (0.614,7.684) (0.623,7.759) (0.632,7.833) (0.641,7.908) (0.650,7.982) (0.660,8.056) (0.669,8.131) (0.678,8.205) (0.687,8.280) (0.697,8.354) (0.706,8.429) (0.715,8.503) (0.725,8.577) (0.734,8.652) (0.744,8.726) (0.753,8.801) (0.763,8.875) (0.773,8.949) (0.782,9.024) (0.792,9.098) (0.802,9.172) (0.812,9.247) (0.822,9.321) (0.831,9.395) (0.841,9.470) (0.851,9.544) (0.861,9.618) (0.871,9.693) (0.881,9.767) (0.892,9.841) (0.902,9.916) (0.912,9.990) (0.915,10.000)}
\LC{(-0.458,5.643) (-0.467,5.717) (-0.476,5.792) (-0.485,5.866) (-0.495,5.941) (-0.504,6.015) (-0.513,6.090) (-0.523,6.164) (-0.532,6.238) (-0.542,6.313) (-0.552,6.387) (-0.561,6.461) (-0.571,6.536) (-0.581,6.610) (-0.591,6.685) (-0.601,6.759) (-0.611,6.833) (-0.621,6.908) (-0.631,6.982) (-0.641,7.056) (-0.651,7.130) (-0.661,7.205) (-0.672,7.279) (-0.682,7.353) (-0.692,7.428) (-0.703,7.502) (-0.713,7.576) (-0.724,7.650) (-0.734,7.725) (-0.745,7.799) (-0.756,7.873) (-0.767,7.947) (-0.777,8.022) (-0.788,8.096) (-0.799,8.170) (-0.810,8.244) (-0.821,8.318) (-0.832,8.392) (-0.843,8.467) (-0.854,8.541) (-0.866,8.615) (-0.877,8.689) (-0.888,8.763) (-0.900,8.837) (-0.911,8.912) (-0.922,8.986) (-0.934,9.060) (-0.945,9.134) (-0.957,9.208) (-0.969,9.282) (-0.980,9.356) (-0.992,9.430) (-1.004,9.504) (-1.016,9.578) (-1.028,9.652) (-1.039,9.726) (-1.051,9.801) (-1.063,9.875) (-1.076,9.949) (-1.088,10.000)}
\LC{(1.579,1.049) (1.654,1.043) (1.728,1.033) (1.802,1.020) (1.875,1.002) (1.947,0.981) (2.017,0.956) (2.086,0.927) (2.154,0.894) (2.219,0.857) (2.282,0.816) (2.342,0.771) (2.399,0.723) (2.453,0.671) (2.503,0.615) (2.549,0.556) (2.591,0.493) (2.628,0.428) (2.660,0.360) (2.686,0.290) (2.707,0.218) (2.722,0.144) (2.731,0.070) (2.733,-0.005) (2.730,-0.080) (2.720,-0.154) (2.704,-0.228) (2.683,-0.299) (2.656,-0.369) (2.623,-0.437) (2.586,-0.502) (2.544,-0.564) (2.497,-0.622) (2.446,-0.678) (2.392,-0.729) (2.334,-0.777) (2.274,-0.822) (2.210,-0.862) (2.145,-0.898) (2.077,-0.931) (2.008,-0.959) (1.937,-0.984) (1.865,-1.005) (1.792,-1.022) (1.718,-1.035) (1.644,-1.044) (1.594,-1.048)}
\LC{(1.474,-10.000) (1.457,-9.936) (1.441,-9.863) (1.425,-9.789) (1.409,-9.716) (1.393,-9.643) (1.377,-9.570) (1.361,-9.496) (1.345,-9.423) (1.329,-9.350) (1.313,-9.276) (1.298,-9.203) (1.282,-9.130) (1.266,-9.056) (1.251,-8.983) (1.235,-8.910) (1.220,-8.836) (1.205,-8.763) (1.189,-8.689) (1.174,-8.616) (1.159,-8.542) (1.144,-8.469) (1.129,-8.395) (1.114,-8.322) (1.099,-8.248) (1.084,-8.175) (1.070,-8.101) (1.055,-8.028) (1.040,-7.954) (1.026,-7.881) (1.012,-7.807) (0.997,-7.734) (0.983,-7.660) (0.969,-7.586) (0.954,-7.513) (0.940,-7.439) (0.926,-7.365) (0.912,-7.292) (0.898,-7.218) (0.884,-7.144) (0.871,-7.070) (0.857,-6.997) (0.843,-6.923) (0.830,-6.849) (0.816,-6.775) (0.803,-6.702) (0.790,-6.628) (0.776,-6.554) (0.763,-6.480) (0.750,-6.406) (0.737,-6.332) (0.724,-6.259) (0.711,-6.185) (0.699,-6.111) (0.686,-6.037) (0.673,-5.963) (0.661,-5.889) (0.648,-5.815) (0.636,-5.741) (0.623,-5.667) (0.611,-5.593) (0.599,-5.519) (0.587,-5.445) (0.575,-5.371) (0.571,-5.346)}
\LC{(-7.036,-10.000) (-6.975,-9.958) (-6.914,-9.914) (-6.853,-9.870) (-6.793,-9.826) (-6.732,-9.782) (-6.672,-9.737) (-6.612,-9.692) (-6.552,-9.647) (-6.492,-9.602) (-6.433,-9.556) (-6.373,-9.510) (-6.314,-9.464) (-6.255,-9.418) (-6.196,-9.372) (-6.137,-9.326) (-6.078,-9.279) (-6.020,-9.232) (-5.961,-9.185) (-5.903,-9.137) (-5.845,-9.090) (-5.787,-9.042) (-5.730,-8.994) (-5.672,-8.946) (-5.615,-8.898) (-5.558,-8.849) (-5.501,-8.800) (-5.444,-8.751) (-5.387,-8.702) (-5.331,-8.653) (-5.274,-8.603) (-5.218,-8.554) (-5.162,-8.504) (-5.107,-8.454) (-5.051,-8.403) (-4.996,-8.353) (-4.940,-8.302) (-4.885,-8.251) (-4.830,-8.200) (-4.776,-8.149) (-4.721,-8.097) (-4.667,-8.045) (-4.613,-7.993) (-4.559,-7.941) (-4.505,-7.889) (-4.452,-7.836) (-4.398,-7.784) (-4.345,-7.731) (-4.292,-7.678) (-4.239,-7.624) (-4.187,-7.571) (-4.134,-7.517) (-4.082,-7.463) (-4.030,-7.409) (-3.978,-7.355) (-3.927,-7.301) (-3.876,-7.246) (-3.824,-7.191) (-3.773,-7.136) (-3.723,-7.081) (-3.672,-7.025) (-3.622,-6.970) (-3.572,-6.914) (-3.522,-6.858) (-3.472,-6.802) (-3.423,-6.745) (-3.374,-6.689) (-3.325,-6.632) (-3.276,-6.575) (-3.227,-6.518) (-3.179,-6.461) (-3.131,-6.403) (-3.083,-6.345) (-3.035,-6.287) (-2.988,-6.229) (-2.941,-6.171) (-2.894,-6.113) (-2.847,-6.054) (-2.800,-5.995) (-2.754,-5.936) (-2.708,-5.877) (-2.662,-5.817) (-2.617,-5.758) (-2.572,-5.698) (-2.527,-5.638) (-2.482,-5.578) (-2.437,-5.517) (-2.393,-5.457) (-2.349,-5.396) (-2.320,-5.356)}
\LC{(10.000,-2.658) (9.952,-2.713) (9.900,-2.767) (9.847,-2.819) (9.793,-2.871) (9.738,-2.922) (9.682,-2.972) (9.625,-3.021) (9.568,-3.069) (9.509,-3.117) (9.450,-3.163) (9.391,-3.208) (9.330,-3.253) (9.269,-3.296) (9.207,-3.339) (9.145,-3.381) (9.082,-3.421) (9.018,-3.461) (8.954,-3.500) (8.889,-3.537) (8.824,-3.574) (8.758,-3.610) (8.691,-3.645) (8.625,-3.679) (8.557,-3.711) (8.489,-3.743) (8.421,-3.774) (8.352,-3.804) (8.283,-3.833) (8.214,-3.861) (8.144,-3.889) (8.073,-3.915) (8.003,-3.940) (7.932,-3.964) (7.860,-3.987) (7.789,-4.010) (7.717,-4.031) (7.645,-4.051) (7.572,-4.071) (7.500,-4.089) (7.427,-4.107) (7.354,-4.124) (7.280,-4.139) (7.207,-4.154) (7.133,-4.168) (7.059,-4.181) (6.985,-4.193) (6.911,-4.204) (6.837,-4.214) (6.762,-4.223) (6.688,-4.231) (6.613,-4.239) (6.538,-4.245) (6.463,-4.251) (6.389,-4.255) (6.314,-4.259) (6.239,-4.262) (6.164,-4.264) (6.089,-4.265) (6.014,-4.265) (5.939,-4.265) (5.864,-4.263) (5.789,-4.261) (5.714,-4.258) (5.639,-4.253) (5.564,-4.248) (5.489,-4.243) (5.415,-4.236) (5.340,-4.228) (5.266,-4.220) (5.191,-4.211) (5.117,-4.201) (5.043,-4.190) (4.969,-4.178) (4.895,-4.165) (4.821,-4.152) (4.747,-4.138) (4.674,-4.123) (4.600,-4.107) (4.527,-4.090) (4.454,-4.073) (4.382,-4.055) (4.309,-4.036) (4.237,-4.016) (4.165,-3.995) (4.093,-3.974) (4.021,-3.952) (3.950,-3.929) (3.879,-3.905) (3.808,-3.880) (3.737,-3.855) (3.667,-3.829) (3.597,-3.802) (3.527,-3.774) (3.458,-3.746) (3.388,-3.717) (3.320,-3.687) (3.251,-3.657) (3.183,-3.625) (3.115,-3.593) (3.048,-3.560) (2.981,-3.527) (2.914,-3.492) (2.848,-3.457) (2.782,-3.422) (2.716,-3.385) (2.651,-3.348) (2.586,-3.310) (2.522,-3.272)}
\LC{(-10.000,-4.265) (-9.942,-4.298) (-9.875,-4.330) (-9.806,-4.362) (-9.738,-4.392) (-9.669,-4.422) (-9.600,-4.451) (-9.530,-4.479) (-9.460,-4.506) (-9.390,-4.532) (-9.320,-4.558) (-9.249,-4.583) (-9.178,-4.606) (-9.106,-4.629) (-9.035,-4.652) (-8.963,-4.673) (-8.891,-4.694) (-8.818,-4.713) (-8.746,-4.732) (-8.673,-4.751) (-8.600,-4.768) (-8.527,-4.784) (-8.454,-4.800) (-8.380,-4.815) (-8.306,-4.829) (-8.233,-4.842) (-8.159,-4.855) (-8.085,-4.867) (-8.010,-4.878) (-7.936,-4.888) (-7.862,-4.897) (-7.787,-4.906) (-7.712,-4.913) (-7.638,-4.920) (-7.563,-4.926) (-7.488,-4.932) (-7.413,-4.936) (-7.339,-4.940) (-7.264,-4.943) (-7.189,-4.946) (-7.114,-4.947) (-7.039,-4.948) (-6.964,-4.948) (-6.889,-4.947) (-6.814,-4.946) (-6.739,-4.944) (-6.664,-4.941) (-6.589,-4.937) (-6.514,-4.933) (-6.439,-4.927) (-6.364,-4.921) (-6.290,-4.915) (-6.215,-4.907) (-6.140,-4.899) (-6.066,-4.890) (-5.992,-4.881) (-5.917,-4.870) (-5.843,-4.859) (-5.769,-4.848) (-5.695,-4.835) (-5.621,-4.822) (-5.548,-4.808) (-5.474,-4.794) (-5.401,-4.778)}
\LC{(-10.000,-7.079) (-9.926,-7.079) (-9.851,-7.078) (-9.776,-7.076) (-9.701,-7.074) (-9.626,-7.072) (-9.552,-7.069) (-9.477,-7.066) (-9.402,-7.062) (-9.327,-7.057) (-9.252,-7.052) (-9.177,-7.046) (-9.102,-7.040) (-9.028,-7.034) (-8.953,-7.027) (-8.878,-7.019) (-8.804,-7.011) (-8.729,-7.002) (-8.655,-6.993) (-8.581,-6.984) (-8.506,-6.973) (-8.432,-6.963) (-8.358,-6.952) (-8.284,-6.940) (-8.210,-6.928) (-8.136,-6.915) (-8.062,-6.902) (-7.988,-6.888) (-7.915,-6.874) (-7.841,-6.860) (-7.768,-6.845) (-7.694,-6.829) (-7.621,-6.813) (-7.548,-6.796) (-7.475,-6.779) (-7.402,-6.762) (-7.329,-6.744) (-7.256,-6.725) (-7.184,-6.706) (-7.111,-6.687) (-7.039,-6.667) (-6.967,-6.646) (-6.895,-6.625) (-6.823,-6.604) (-6.751,-6.582) (-6.680,-6.560) (-6.608,-6.537) (-6.537,-6.514) (-6.466,-6.490) (-6.395,-6.466) (-6.324,-6.441) (-6.253,-6.416) (-6.183,-6.391) (-6.112,-6.365) (-6.042,-6.338) (-5.972,-6.311) (-5.902,-6.284) (-5.833,-6.256) (-5.763,-6.227) (-5.694,-6.199) (-5.625,-6.169) (-5.556,-6.140) (-5.488,-6.110) (-5.419,-6.079) (-5.351,-6.048) (-5.283,-6.016) (-5.215,-5.985) (-5.147,-5.952) (-5.080,-5.919) (-5.013,-5.886)}
\LC{(6.534,-10.000) (6.476,-9.962) (6.418,-9.915) (6.360,-9.868) (6.302,-9.820) (6.244,-9.772) (6.186,-9.724) (6.129,-9.676) (6.072,-9.628) (6.014,-9.579) (5.957,-9.530) (5.901,-9.481) (5.844,-9.432) (5.787,-9.383) (5.731,-9.333) (5.675,-9.284) (5.619,-9.234) (5.563,-9.184) (5.507,-9.134) (5.452,-9.083) (5.396,-9.033) (5.341,-8.982) (5.286,-8.931) (5.231,-8.880) (5.177,-8.828) (5.122,-8.777) (5.068,-8.725) (5.013,-8.673) (4.959,-8.621) (4.906,-8.569) (4.852,-8.517) (4.799,-8.464) (4.745,-8.411) (4.692,-8.358) (4.639,-8.305) (4.586,-8.252) (4.534,-8.198) (4.482,-8.145) (4.429,-8.091) (4.377,-8.037) (4.325,-7.983) (4.274,-7.928) (4.222,-7.874) (4.171,-7.819) (4.120,-7.764) (4.069,-7.709) (4.019,-7.654) (3.968,-7.598) (3.918,-7.542) (3.868,-7.487) (3.818,-7.431) (3.768,-7.374) (3.719,-7.318) (3.669,-7.262) (3.620,-7.205) (3.571,-7.148) (3.523,-7.091) (3.474,-7.034) (3.426,-6.976) (3.378,-6.919) (3.330,-6.861) (3.282,-6.803) (3.235,-6.745) (3.188,-6.687) (3.141,-6.628) (3.094,-6.570) (3.047,-6.511) (3.001,-6.452) (2.955,-6.393) (2.909,-6.333) (2.863,-6.274) (2.818,-6.214) (2.773,-6.154) (2.728,-6.094) (2.683,-6.034) (2.639,-5.974) (2.594,-5.913) (2.550,-5.852) (2.521,-5.812)}
\LC{(-2.798,6.367) (-2.842,6.428) (-2.886,6.488) (-2.931,6.548) (-2.976,6.608) (-3.021,6.668) (-3.067,6.728) (-3.112,6.788) (-3.158,6.847) (-3.204,6.906) (-3.250,6.965) (-3.297,7.024) (-3.343,7.083) (-3.390,7.141) (-3.437,7.200) (-3.484,7.258) (-3.532,7.316) (-3.580,7.374) (-3.627,7.432) (-3.675,7.489) (-3.724,7.547) (-3.772,7.604) (-3.821,7.661) (-3.870,7.718) (-3.919,7.775) (-3.968,7.832) (-4.017,7.888) (-4.067,7.944) (-4.116,8.000) (-4.166,8.056) (-4.216,8.112) (-4.267,8.168) (-4.317,8.223) (-4.368,8.279) (-4.419,8.334) (-4.470,8.389) (-4.521,8.444) (-4.572,8.498) (-4.624,8.553) (-4.676,8.607) (-4.728,8.661) (-4.780,8.715) (-4.832,8.769) (-4.884,8.822) (-4.937,8.876) (-4.990,8.929) (-5.043,8.982) (-5.096,9.035) (-5.149,9.088) (-5.203,9.141) (-5.256,9.193) (-5.310,9.245) (-5.364,9.298) (-5.418,9.349) (-5.472,9.401) (-5.527,9.453) (-5.581,9.504) (-5.636,9.556) (-5.691,9.607) (-5.746,9.658) (-5.801,9.708) (-5.857,9.759) (-5.912,9.809) (-5.968,9.860) (-6.024,9.910) (-6.080,9.960) (-6.136,10.000)}
\LC{(-1.929,0.938) (-1.998,0.911) (-2.066,0.879) (-2.132,0.844) (-2.196,0.804) (-2.257,0.760) (-2.315,0.713) (-2.370,0.661) (-2.421,0.606) (-2.467,0.548) (-2.510,0.486) (-2.547,0.421) (-2.579,0.353) (-2.606,0.283) (-2.627,0.211) (-2.641,0.137) (-2.650,0.063) (-2.652,-0.012) (-2.648,-0.087) (-2.637,-0.161) (-2.621,-0.234) (-2.598,-0.306) (-2.570,-0.375) (-2.536,-0.442) (-2.497,-0.506) (-2.453,-0.567) (-2.405,-0.624) (-2.353,-0.678) (-2.297,-0.728) (-2.238,-0.775) (-2.176,-0.817) (-2.112,-0.855) (-2.045,-0.890) (-1.976,-0.920) (-1.930,-0.938)}
\LC{(-1.501,-10.000) (-1.484,-9.935) (-1.467,-9.862) (-1.451,-9.789) (-1.435,-9.716) (-1.418,-9.643) (-1.402,-9.569) (-1.386,-9.496) (-1.369,-9.423) (-1.353,-9.350) (-1.337,-9.276) (-1.321,-9.203) (-1.305,-9.130) (-1.290,-9.057) (-1.274,-8.983) (-1.258,-8.910) (-1.242,-8.837) (-1.227,-8.763) (-1.211,-8.690) (-1.196,-8.616) (-1.180,-8.543) (-1.165,-8.470) (-1.150,-8.396) (-1.135,-8.323) (-1.120,-8.249) (-1.104,-8.176) (-1.089,-8.102) (-1.075,-8.029) (-1.060,-7.955) (-1.045,-7.882) (-1.030,-7.808) (-1.016,-7.735) (-1.001,-7.661) (-0.986,-7.588) (-0.972,-7.514) (-0.958,-7.440) (-0.943,-7.367) (-0.929,-7.293) (-0.915,-7.219) (-0.901,-7.146) (-0.887,-7.072) (-0.873,-6.998) (-0.859,-6.925) (-0.845,-6.851) (-0.832,-6.777) (-0.818,-6.703) (-0.804,-6.630) (-0.791,-6.556) (-0.777,-6.482) (-0.764,-6.408) (-0.751,-6.334) (-0.738,-6.261) (-0.725,-6.187) (-0.712,-6.113) (-0.699,-6.039) (-0.686,-5.965) (-0.673,-5.891) (-0.660,-5.817) (-0.648,-5.743) (-0.635,-5.669) (-0.623,-5.596) (-0.610,-5.522) (-0.598,-5.448) (-0.586,-5.374) (-0.574,-5.300) (-0.562,-5.226) (-0.550,-5.152) (-0.538,-5.077) (-0.526,-5.003) (-0.522,-4.979)}
\LC{(-5.577,4.285) (-5.652,4.290) (-5.727,4.295) (-5.802,4.299) (-5.877,4.302) (-5.952,4.304) (-6.027,4.305) (-6.102,4.305) (-6.177,4.304) (-6.252,4.303) (-6.327,4.301) (-6.401,4.297) (-6.476,4.293) (-6.551,4.288) (-6.626,4.282) (-6.701,4.276) (-6.775,4.268) (-6.850,4.259) (-6.924,4.250) (-6.998,4.239) (-7.073,4.228) (-7.147,4.216) (-7.220,4.203) (-7.294,4.189) (-7.368,4.174) (-7.441,4.158) (-7.514,4.141) (-7.587,4.123) (-7.659,4.104) (-7.732,4.085) (-7.804,4.064) (-7.876,4.042) (-7.947,4.020) (-8.019,3.996) (-8.089,3.972) (-8.160,3.947) (-8.230,3.920) (-8.300,3.893) (-8.370,3.865) (-8.439,3.836) (-8.507,3.806) (-8.576,3.774) (-8.643,3.742) (-8.711,3.709) (-8.778,3.675) (-8.844,3.641) (-8.910,3.605) (-8.975,3.568) (-9.040,3.530) (-9.104,3.491) (-9.168,3.452) (-9.231,3.411) (-9.293,3.369) (-9.355,3.327) (-9.416,3.283) (-9.477,3.239) (-9.536,3.194) (-9.595,3.147) (-9.654,3.100) (-9.711,3.052) (-9.768,3.003) (-9.824,2.953) (-9.879,2.902) (-9.933,2.850) (-9.986,2.797) (-10.000,2.780)}

    % Dipolo puntual - puntas de flecha
% Generado por scripts/generar_datos.py (no editar a mano).
\FL{1.429}{-0.015}{-91.8}
\FL{1.467}{2.168}{37.2}
\FL{4.190}{3.194}{4.1}
\FL{6.989}{2.463}{-35.4}
\FL{8.330}{-0.035}{-90.7}
\FL{6.953}{-2.489}{-145.3}
\FL{4.146}{-3.191}{175.4}
\FL{1.433}{-2.141}{142.3}
\FL{-1.467}{2.168}{142.8}
\FL{-4.190}{3.194}{175.9}
\FL{-6.989}{2.463}{-144.6}
\FL{-8.330}{-0.035}{-89.3}
\FL{-6.953}{-2.489}{-34.7}
\FL{-4.146}{-3.191}{4.6}
\FL{-1.433}{-2.141}{37.7}
\FL{-1.429}{-0.015}{-88.2}
\FL{3.888}{-8.732}{126.6}
\FL{2.141}{-6.053}{119.5}
\FL{0.773}{-3.164}{110.7}
\FL{-4.014}{-8.749}{52.4}
\FL{-2.211}{-6.078}{59.7}
\FL{-0.799}{-3.182}{68.8}
\FL{8.283}{-6.321}{175.9}
\FL{4.997}{-5.511}{156.6}
\FL{2.162}{-3.665}{137.0}
\FL{-4.602}{2.370}{-163.1}
\FL{-6.550}{-0.011}{-89.7}
\FL{-4.581}{-2.377}{-16.5}
\FL{0.881}{2.999}{65.3}
\FL{2.174}{5.280}{56.0}
\FL{3.771}{7.329}{48.3}
\FL{5.631}{9.179}{41.4}
\FL{-5.151}{3.827}{177.3}
\FL{-7.813}{3.273}{-152.7}
\FL{-9.682}{1.331}{-113.2}
\FL{-9.677}{-1.343}{-66.6}
\FL{-7.802}{-3.279}{-27.1}
\FL{-5.138}{-3.826}{2.8}
\FL{-0.035}{-8.072}{89.6}
\FL{-0.013}{-4.222}{89.7}
\FL{-1.128}{3.214}{119.3}
\FL{-2.813}{5.599}{130.8}
\FL{-4.925}{7.654}{140.6}
\FL{-7.322}{9.325}{149.6}
\FL{-3.065}{4.757}{140.7}
\FL{-5.627}{6.350}{155.4}
\FL{-8.507}{7.248}{170.0}
\FL{-8.690}{-7.903}{13.5}
\FL{-6.169}{-7.023}{25.0}
\FL{-3.895}{-5.671}{36.6}
\FL{-1.952}{-3.877}{49.1}
\FL{5.299}{2.380}{-23.9}
\FL{6.981}{-0.015}{-90.4}
\FL{5.272}{-2.392}{-156.6}
\FL{8.439}{-7.362}{168.7}
\FL{5.486}{-6.366}{154.1}
\FL{4.915}{3.849}{5.5}
\FL{7.491}{3.499}{-21.9}
\FL{9.507}{1.882}{-57.2}
\FL{3.579}{6.089}{43.2}
\FL{5.947}{7.925}{32.5}
\FL{8.605}{9.307}{22.4}
\FL{3.187}{1.458}{-23.0}
\FL{3.151}{-1.473}{-158.1}
\FL{2.160}{-8.592}{111.3}
\FL{1.153}{-5.714}{107.2}
\FL{-4.995}{6.777}{146.6}
\FL{-8.259}{8.389}{160.8}
\FL{9.384}{0.914}{-73.4}
\FL{8.847}{-1.977}{-126.7}
\FL{6.306}{-3.544}{-167.7}
\FL{7.949}{-9.303}{153.6}
\FL{5.005}{-7.467}{142.4}
\FL{2.512}{-5.091}{130.1}
\FL{-8.679}{-5.784}{-3.2}
\FL{-6.027}{-5.529}{13.9}
\FL{-3.557}{-4.524}{30.4}
\FL{8.555}{-5.231}{-172.2}
\FL{5.625}{-5.084}{166.8}
\FL{-0.941}{4.720}{107.0}
\FL{-2.243}{8.288}{112.8}
\FL{-8.254}{-8.950}{22.6}
\FL{-5.066}{-7.187}{35.4}
\FL{-2.308}{-8.603}{67.3}
\FL{-1.253}{-5.795}{71.6}
\FL{3.758}{2.167}{-10.9}
\FL{5.781}{-0.009}{-90.3}
\FL{3.741}{-2.171}{-169.4}
\FL{-0.004}{4.753}{90.1}
\FL{-0.010}{8.253}{90.1}
\FL{6.083}{4.475}{2.1}
\FL{8.815}{3.977}{-23.7}
\FL{-2.961}{3.778}{149.5}
\FL{-5.722}{4.766}{171.2}
\FL{-8.648}{4.626}{-165.1}
\FL{4.561}{6.002}{32.0}
\FL{8.107}{7.526}{14.6}
\FL{-2.725}{1.238}{-156.7}
\FL{-2.701}{-1.247}{-22.5}
\FL{-8.897}{-3.737}{-27.2}
\FL{-6.345}{-4.367}{-1.5}
\FL{-3.980}{1.547}{-149.1}
\FL{-3.966}{-1.555}{-30.5}
\FL{-5.638}{5.666}{161.3}
\FL{-8.531}{6.181}{178.7}
\FL{2.629}{3.737}{35.5}
\FL{5.423}{5.041}{14.6}
\FL{8.497}{5.250}{-7.3}
\FL{2.013}{6.217}{62.9}
\FL{3.508}{8.788}{57.0}
\FL{-2.126}{6.103}{119.1}
\FL{-3.822}{8.756}{125.9}
\FL{0.825}{5.102}{76.2}
\FL{1.763}{8.396}{72.1}
\FL{8.584}{-8.611}{161.4}
\FL{5.940}{-7.430}{150.4}
\FL{5.637}{5.904}{20.8}
\FL{8.529}{6.557}{4.5}
\FL{0.666}{8.106}{83.0}
\FL{-0.752}{7.848}{98.2}
\FL{2.733}{-0.005}{-90.3}
\FL{0.983}{-7.660}{101.0}
\FL{-5.711}{-8.978}{39.9}
\FL{-3.341}{-6.651}{49.2}
\FL{8.932}{-3.512}{-149.6}
\FL{6.364}{-4.257}{-177.0}
\FL{3.714}{-3.846}{159.8}
\FL{-7.737}{-4.911}{-5.8}
\FL{-7.426}{-6.768}{13.6}
\FL{5.415}{-9.049}{137.6}
\FL{3.410}{-6.957}{129.9}
\FL{-4.368}{8.279}{132.6}
\FL{-2.652}{-0.012}{-89.2}
\FL{-0.967}{-7.489}{78.9}
\FL{-7.971}{4.012}{-161.9}
}
  \dipoloRojo{0.85}
  \marcoMMA{10}{5}{$x/\ell$}{$z/\ell$}
\end{tikzpicture}\\[1pt]{\small\scshape (b)}
\end{minipage}\hfill
\begin{minipage}[t]{0.325\textwidth}\centering
{\footnotesize Campo eléctrico}\\[3pt]
\begin{tikzpicture}[x=0.198cm, y=0.198cm]
  \panelMMA{10}{mmafondo}{1.2}{% Dipolo puntual - campo vectorial
% Generado por scripts/generar_datos.py (no editar a mano).
\VF{-10.000}{-10.000}{0.370}{0.123}{0.000}
\VF{-9.000}{-10.000}{0.357}{0.157}{0.000}
\VF{-8.000}{-10.000}{0.360}{0.204}{0.026}
\VF{-7.000}{-10.000}{0.361}{0.260}{0.061}
\VF{-6.000}{-10.000}{0.351}{0.320}{0.094}
\VF{-5.000}{-10.000}{0.327}{0.382}{0.124}
\VF{-4.000}{-10.000}{0.288}{0.442}{0.152}
\VF{-3.000}{-10.000}{0.234}{0.496}{0.174}
\VF{-2.000}{-10.000}{0.165}{0.540}{0.192}
\VF{-1.000}{-10.000}{0.086}{0.568}{0.202}
\VF{0.000}{-10.000}{-0.000}{0.578}{0.206}
\VF{1.000}{-10.000}{-0.086}{0.568}{0.202}
\VF{2.000}{-10.000}{-0.165}{0.540}{0.192}
\VF{3.000}{-10.000}{-0.234}{0.496}{0.174}
\VF{4.000}{-10.000}{-0.288}{0.442}{0.152}
\VF{5.000}{-10.000}{-0.327}{0.382}{0.124}
\VF{6.000}{-10.000}{-0.351}{0.320}{0.094}
\VF{7.000}{-10.000}{-0.361}{0.260}{0.061}
\VF{8.000}{-10.000}{-0.360}{0.204}{0.026}
\VF{9.000}{-10.000}{-0.357}{0.157}{0.000}
\VF{10.000}{-10.000}{-0.370}{0.123}{0.000}
\VF{-10.000}{-9.000}{0.380}{0.087}{0.000}
\VF{-9.000}{-9.000}{0.385}{0.128}{0.018}
\VF{-8.000}{-9.000}{0.402}{0.183}{0.057}
\VF{-7.000}{-9.000}{0.410}{0.245}{0.096}
\VF{-6.000}{-9.000}{0.404}{0.314}{0.134}
\VF{-5.000}{-9.000}{0.382}{0.388}{0.170}
\VF{-4.000}{-9.000}{0.341}{0.461}{0.202}
\VF{-3.000}{-9.000}{0.280}{0.529}{0.230}
\VF{-2.000}{-9.000}{0.200}{0.585}{0.251}
\VF{-1.000}{-9.000}{0.104}{0.621}{0.264}
\VF{0.000}{-9.000}{-0.000}{0.634}{0.268}
\VF{1.000}{-9.000}{-0.104}{0.621}{0.264}
\VF{2.000}{-9.000}{-0.200}{0.585}{0.251}
\VF{3.000}{-9.000}{-0.280}{0.529}{0.230}
\VF{4.000}{-9.000}{-0.341}{0.461}{0.202}
\VF{5.000}{-9.000}{-0.382}{0.388}{0.170}
\VF{6.000}{-9.000}{-0.404}{0.314}{0.134}
\VF{7.000}{-9.000}{-0.410}{0.245}{0.096}
\VF{8.000}{-9.000}{-0.402}{0.183}{0.057}
\VF{9.000}{-9.000}{-0.385}{0.128}{0.018}
\VF{10.000}{-9.000}{-0.380}{0.087}{0.000}
\VF{-10.000}{-8.000}{0.387}{0.045}{0.000}
\VF{-9.000}{-8.000}{0.419}{0.091}{0.043}
\VF{-8.000}{-8.000}{0.445}{0.148}{0.087}
\VF{-7.000}{-8.000}{0.461}{0.217}{0.131}
\VF{-6.000}{-8.000}{0.463}{0.296}{0.175}
\VF{-5.000}{-8.000}{0.446}{0.383}{0.217}
\VF{-4.000}{-8.000}{0.405}{0.473}{0.256}
\VF{-3.000}{-8.000}{0.338}{0.559}{0.289}
\VF{-2.000}{-8.000}{0.244}{0.631}{0.315}
\VF{-1.000}{-8.000}{0.129}{0.680}{0.332}
\VF{0.000}{-8.000}{-0.000}{0.697}{0.338}
\VF{1.000}{-8.000}{-0.129}{0.680}{0.332}
\VF{2.000}{-8.000}{-0.244}{0.631}{0.315}
\VF{3.000}{-8.000}{-0.338}{0.559}{0.289}
\VF{4.000}{-8.000}{-0.405}{0.473}{0.256}
\VF{5.000}{-8.000}{-0.446}{0.383}{0.217}
\VF{6.000}{-8.000}{-0.463}{0.296}{0.175}
\VF{7.000}{-8.000}{-0.461}{0.217}{0.131}
\VF{8.000}{-8.000}{-0.445}{0.148}{0.087}
\VF{9.000}{-8.000}{-0.419}{0.091}{0.043}
\VF{10.000}{-8.000}{-0.387}{0.045}{0.000}
\VF{-10.000}{-7.000}{0.408}{-0.004}{0.020}
\VF{-9.000}{-7.000}{0.449}{0.040}{0.067}
\VF{-8.000}{-7.000}{0.485}{0.098}{0.116}
\VF{-7.000}{-7.000}{0.513}{0.171}{0.166}
\VF{-6.000}{-7.000}{0.527}{0.259}{0.216}
\VF{-5.000}{-7.000}{0.519}{0.361}{0.266}
\VF{-4.000}{-7.000}{0.483}{0.471}{0.313}
\VF{-3.000}{-7.000}{0.412}{0.582}{0.355}
\VF{-2.000}{-7.000}{0.303}{0.678}{0.388}
\VF{-1.000}{-7.000}{0.161}{0.745}{0.409}
\VF{0.000}{-7.000}{-0.000}{0.769}{0.417}
\VF{1.000}{-7.000}{-0.161}{0.745}{0.409}
\VF{2.000}{-7.000}{-0.303}{0.678}{0.388}
\VF{3.000}{-7.000}{-0.412}{0.582}{0.355}
\VF{4.000}{-7.000}{-0.483}{0.471}{0.313}
\VF{5.000}{-7.000}{-0.519}{0.361}{0.266}
\VF{6.000}{-7.000}{-0.527}{0.259}{0.216}
\VF{7.000}{-7.000}{-0.513}{0.171}{0.166}
\VF{8.000}{-7.000}{-0.485}{0.098}{0.116}
\VF{9.000}{-7.000}{-0.449}{0.040}{0.067}
\VF{10.000}{-7.000}{-0.408}{-0.004}{0.020}
\VF{-10.000}{-6.000}{0.418}{-0.065}{0.037}
\VF{-9.000}{-6.000}{0.469}{-0.026}{0.088}
\VF{-8.000}{-6.000}{0.518}{0.029}{0.142}
\VF{-7.000}{-6.000}{0.561}{0.102}{0.198}
\VF{-6.000}{-6.000}{0.592}{0.197}{0.257}
\VF{-5.000}{-6.000}{0.600}{0.314}{0.316}
\VF{-4.000}{-6.000}{0.576}{0.448}{0.373}
\VF{-3.000}{-6.000}{0.506}{0.590}{0.426}
\VF{-2.000}{-6.000}{0.382}{0.722}{0.469}
\VF{-1.000}{-6.000}{0.207}{0.817}{0.497}
\VF{0.000}{-6.000}{-0.000}{0.852}{0.508}
\VF{1.000}{-6.000}{-0.207}{0.817}{0.497}
\VF{2.000}{-6.000}{-0.382}{0.722}{0.469}
\VF{3.000}{-6.000}{-0.506}{0.590}{0.426}
\VF{4.000}{-6.000}{-0.576}{0.448}{0.373}
\VF{5.000}{-6.000}{-0.600}{0.314}{0.316}
\VF{6.000}{-6.000}{-0.592}{0.197}{0.257}
\VF{7.000}{-6.000}{-0.561}{0.102}{0.198}
\VF{8.000}{-6.000}{-0.518}{0.029}{0.142}
\VF{9.000}{-6.000}{-0.469}{-0.026}{0.088}
\VF{10.000}{-6.000}{-0.418}{-0.065}{0.037}
\VF{-10.000}{-5.000}{0.413}{-0.138}{0.050}
\VF{-9.000}{-5.000}{0.473}{-0.109}{0.105}
\VF{-8.000}{-5.000}{0.536}{-0.063}{0.164}
\VF{-7.000}{-5.000}{0.597}{0.006}{0.227}
\VF{-6.000}{-5.000}{0.650}{0.101}{0.294}
\VF{-5.000}{-5.000}{0.685}{0.228}{0.364}
\VF{-4.000}{-5.000}{0.684}{0.388}{0.435}
\VF{-3.000}{-5.000}{0.626}{0.571}{0.503}
\VF{-2.000}{-5.000}{0.492}{0.754}{0.561}
\VF{-1.000}{-5.000}{0.274}{0.896}{0.601}
\VF{0.000}{-5.000}{-0.000}{0.950}{0.615}
\VF{1.000}{-5.000}{-0.274}{0.896}{0.601}
\VF{2.000}{-5.000}{-0.492}{0.754}{0.561}
\VF{3.000}{-5.000}{-0.626}{0.571}{0.503}
\VF{4.000}{-5.000}{-0.684}{0.388}{0.435}
\VF{5.000}{-5.000}{-0.685}{0.228}{0.364}
\VF{6.000}{-5.000}{-0.650}{0.101}{0.294}
\VF{7.000}{-5.000}{-0.597}{0.006}{0.227}
\VF{8.000}{-5.000}{-0.536}{-0.063}{0.164}
\VF{9.000}{-5.000}{-0.473}{-0.109}{0.105}
\VF{10.000}{-5.000}{-0.413}{-0.138}{0.050}
\VF{-10.000}{-4.000}{0.387}{-0.219}{0.060}
\VF{-9.000}{-4.000}{0.453}{-0.206}{0.118}
\VF{-8.000}{-4.000}{0.527}{-0.176}{0.182}
\VF{-7.000}{-4.000}{0.606}{-0.123}{0.251}
\VF{-6.000}{-4.000}{0.687}{-0.038}{0.327}
\VF{-5.000}{-4.000}{0.757}{0.088}{0.409}
\VF{-4.000}{-4.000}{0.798}{0.266}{0.496}
\VF{-3.000}{-4.000}{0.777}{0.496}{0.584}
\VF{-2.000}{-4.000}{0.648}{0.756}{0.665}
\VF{-1.000}{-4.000}{0.379}{0.979}{0.725}
\VF{0.000}{-4.000}{-0.000}{1.070}{0.747}
\VF{1.000}{-4.000}{-0.379}{0.979}{0.725}
\VF{2.000}{-4.000}{-0.648}{0.756}{0.665}
\VF{3.000}{-4.000}{-0.777}{0.496}{0.584}
\VF{4.000}{-4.000}{-0.798}{0.266}{0.496}
\VF{5.000}{-4.000}{-0.757}{0.088}{0.409}
\VF{6.000}{-4.000}{-0.687}{-0.038}{0.327}
\VF{7.000}{-4.000}{-0.606}{-0.123}{0.251}
\VF{8.000}{-4.000}{-0.527}{-0.176}{0.182}
\VF{9.000}{-4.000}{-0.453}{-0.206}{0.118}
\VF{10.000}{-4.000}{-0.387}{-0.219}{0.060}
\VF{-10.000}{-3.000}{0.333}{-0.303}{0.066}
\VF{-9.000}{-3.000}{0.399}{-0.310}{0.127}
\VF{-8.000}{-3.000}{0.477}{-0.305}{0.194}
\VF{-7.000}{-3.000}{0.569}{-0.280}{0.268}
\VF{-6.000}{-3.000}{0.674}{-0.225}{0.352}
\VF{-5.000}{-3.000}{0.786}{-0.122}{0.446}
\VF{-4.000}{-3.000}{0.890}{0.049}{0.551}
\VF{-3.000}{-3.000}{0.945}{0.315}{0.666}
\VF{-2.000}{-3.000}{0.870}{0.677}{0.782}
\VF{-1.000}{-3.000}{0.556}{1.051}{0.878}
\VF{0.000}{-3.000}{-0.000}{1.224}{0.917}
\VF{1.000}{-3.000}{-0.556}{1.051}{0.878}
\VF{2.000}{-3.000}{-0.870}{0.677}{0.782}
\VF{3.000}{-3.000}{-0.945}{0.315}{0.666}
\VF{4.000}{-3.000}{-0.890}{0.049}{0.551}
\VF{5.000}{-3.000}{-0.786}{-0.122}{0.446}
\VF{6.000}{-3.000}{-0.674}{-0.225}{0.352}
\VF{7.000}{-3.000}{-0.569}{-0.280}{0.268}
\VF{8.000}{-3.000}{-0.477}{-0.305}{0.194}
\VF{9.000}{-3.000}{-0.399}{-0.310}{0.127}
\VF{10.000}{-3.000}{-0.333}{-0.303}{0.066}
\VF{-10.000}{-2.000}{0.247}{-0.379}{0.069}
\VF{-9.000}{-2.000}{0.303}{-0.409}{0.131}
\VF{-8.000}{-2.000}{0.372}{-0.434}{0.200}
\VF{-7.000}{-2.000}{0.460}{-0.449}{0.277}
\VF{-6.000}{-2.000}{0.571}{-0.444}{0.366}
\VF{-5.000}{-2.000}{0.711}{-0.403}{0.469}
\VF{-4.000}{-2.000}{0.880}{-0.293}{0.591}
\VF{-3.000}{-2.000}{1.058}{-0.059}{0.736}
\VF{-2.000}{-2.000}{1.151}{0.384}{0.905}
\VF{-1.000}{-2.000}{0.846}{0.987}{1.000}
\VF{0.000}{-2.000}{-0.000}{1.300}{1.000}
\VF{1.000}{-2.000}{-0.846}{0.987}{1.000}
\VF{2.000}{-2.000}{-1.151}{0.384}{0.905}
\VF{3.000}{-2.000}{-1.058}{-0.059}{0.736}
\VF{4.000}{-2.000}{-0.880}{-0.293}{0.591}
\VF{5.000}{-2.000}{-0.711}{-0.403}{0.469}
\VF{6.000}{-2.000}{-0.571}{-0.444}{0.366}
\VF{7.000}{-2.000}{-0.460}{-0.449}{0.277}
\VF{8.000}{-2.000}{-0.372}{-0.434}{0.200}
\VF{9.000}{-2.000}{-0.303}{-0.409}{0.131}
\VF{10.000}{-2.000}{-0.247}{-0.379}{0.069}
\VF{-10.000}{-1.000}{0.133}{-0.434}{0.070}
\VF{-9.000}{-1.000}{0.165}{-0.483}{0.132}
\VF{-8.000}{-1.000}{0.207}{-0.535}{0.201}
\VF{-7.000}{-1.000}{0.263}{-0.589}{0.280}
\VF{-6.000}{-1.000}{0.340}{-0.643}{0.371}
\VF{-5.000}{-1.000}{0.451}{-0.691}{0.478}
\VF{-4.000}{-1.000}{0.614}{-0.717}{0.609}
\VF{-3.000}{-1.000}{0.865}{-0.672}{0.775}
\VF{-2.000}{-1.000}{1.233}{-0.411}{1.000}
\VF{-1.000}{-1.000}{1.233}{0.411}{1.000}
\VF{1.000}{-1.000}{-1.233}{0.411}{1.000}
\VF{2.000}{-1.000}{-1.233}{-0.411}{1.000}
\VF{3.000}{-1.000}{-0.865}{-0.672}{0.775}
\VF{4.000}{-1.000}{-0.614}{-0.717}{0.609}
\VF{5.000}{-1.000}{-0.451}{-0.691}{0.478}
\VF{6.000}{-1.000}{-0.340}{-0.643}{0.371}
\VF{7.000}{-1.000}{-0.263}{-0.589}{0.280}
\VF{8.000}{-1.000}{-0.207}{-0.535}{0.201}
\VF{9.000}{-1.000}{-0.165}{-0.483}{0.132}
\VF{10.000}{-1.000}{-0.133}{-0.434}{0.070}
\VF{-10.000}{0.000}{-0.000}{-0.453}{0.070}
\VF{-9.000}{0.000}{-0.000}{-0.510}{0.132}
\VF{-8.000}{0.000}{-0.000}{-0.573}{0.201}
\VF{-7.000}{0.000}{-0.000}{-0.645}{0.280}
\VF{-6.000}{0.000}{-0.000}{-0.728}{0.371}
\VF{-5.000}{0.000}{-0.000}{-0.826}{0.479}
\VF{-4.000}{0.000}{-0.000}{-0.946}{0.611}
\VF{-3.000}{0.000}{-0.000}{-1.100}{0.780}
\VF{-2.000}{0.000}{-0.000}{-1.300}{1.000}
\VF{2.000}{0.000}{0.000}{-1.300}{1.000}
\VF{3.000}{0.000}{0.000}{-1.100}{0.780}
\VF{4.000}{0.000}{0.000}{-0.946}{0.611}
\VF{5.000}{0.000}{0.000}{-0.826}{0.479}
\VF{6.000}{0.000}{0.000}{-0.728}{0.371}
\VF{7.000}{0.000}{0.000}{-0.645}{0.280}
\VF{8.000}{0.000}{0.000}{-0.573}{0.201}
\VF{9.000}{0.000}{0.000}{-0.510}{0.132}
\VF{10.000}{0.000}{0.000}{-0.453}{0.070}
\VF{-10.000}{1.000}{-0.133}{-0.434}{0.070}
\VF{-9.000}{1.000}{-0.165}{-0.483}{0.132}
\VF{-8.000}{1.000}{-0.207}{-0.535}{0.201}
\VF{-7.000}{1.000}{-0.263}{-0.589}{0.280}
\VF{-6.000}{1.000}{-0.340}{-0.643}{0.371}
\VF{-5.000}{1.000}{-0.451}{-0.691}{0.478}
\VF{-4.000}{1.000}{-0.614}{-0.717}{0.609}
\VF{-3.000}{1.000}{-0.865}{-0.672}{0.775}
\VF{-2.000}{1.000}{-1.233}{-0.411}{1.000}
\VF{-1.000}{1.000}{-1.233}{0.411}{1.000}
\VF{1.000}{1.000}{1.233}{0.411}{1.000}
\VF{2.000}{1.000}{1.233}{-0.411}{1.000}
\VF{3.000}{1.000}{0.865}{-0.672}{0.775}
\VF{4.000}{1.000}{0.614}{-0.717}{0.609}
\VF{5.000}{1.000}{0.451}{-0.691}{0.478}
\VF{6.000}{1.000}{0.340}{-0.643}{0.371}
\VF{7.000}{1.000}{0.263}{-0.589}{0.280}
\VF{8.000}{1.000}{0.207}{-0.535}{0.201}
\VF{9.000}{1.000}{0.165}{-0.483}{0.132}
\VF{10.000}{1.000}{0.133}{-0.434}{0.070}
\VF{-10.000}{2.000}{-0.247}{-0.379}{0.069}
\VF{-9.000}{2.000}{-0.303}{-0.409}{0.131}
\VF{-8.000}{2.000}{-0.372}{-0.434}{0.200}
\VF{-7.000}{2.000}{-0.460}{-0.449}{0.277}
\VF{-6.000}{2.000}{-0.571}{-0.444}{0.366}
\VF{-5.000}{2.000}{-0.711}{-0.403}{0.469}
\VF{-4.000}{2.000}{-0.880}{-0.293}{0.591}
\VF{-3.000}{2.000}{-1.058}{-0.059}{0.736}
\VF{-2.000}{2.000}{-1.151}{0.384}{0.905}
\VF{-1.000}{2.000}{-0.846}{0.987}{1.000}
\VF{0.000}{2.000}{0.000}{1.300}{1.000}
\VF{1.000}{2.000}{0.846}{0.987}{1.000}
\VF{2.000}{2.000}{1.151}{0.384}{0.905}
\VF{3.000}{2.000}{1.058}{-0.059}{0.736}
\VF{4.000}{2.000}{0.880}{-0.293}{0.591}
\VF{5.000}{2.000}{0.711}{-0.403}{0.469}
\VF{6.000}{2.000}{0.571}{-0.444}{0.366}
\VF{7.000}{2.000}{0.460}{-0.449}{0.277}
\VF{8.000}{2.000}{0.372}{-0.434}{0.200}
\VF{9.000}{2.000}{0.303}{-0.409}{0.131}
\VF{10.000}{2.000}{0.247}{-0.379}{0.069}
\VF{-10.000}{3.000}{-0.333}{-0.303}{0.066}
\VF{-9.000}{3.000}{-0.399}{-0.310}{0.127}
\VF{-8.000}{3.000}{-0.477}{-0.305}{0.194}
\VF{-7.000}{3.000}{-0.569}{-0.280}{0.268}
\VF{-6.000}{3.000}{-0.674}{-0.225}{0.352}
\VF{-5.000}{3.000}{-0.786}{-0.122}{0.446}
\VF{-4.000}{3.000}{-0.890}{0.049}{0.551}
\VF{-3.000}{3.000}{-0.945}{0.315}{0.666}
\VF{-2.000}{3.000}{-0.870}{0.677}{0.782}
\VF{-1.000}{3.000}{-0.556}{1.051}{0.878}
\VF{0.000}{3.000}{0.000}{1.224}{0.917}
\VF{1.000}{3.000}{0.556}{1.051}{0.878}
\VF{2.000}{3.000}{0.870}{0.677}{0.782}
\VF{3.000}{3.000}{0.945}{0.315}{0.666}
\VF{4.000}{3.000}{0.890}{0.049}{0.551}
\VF{5.000}{3.000}{0.786}{-0.122}{0.446}
\VF{6.000}{3.000}{0.674}{-0.225}{0.352}
\VF{7.000}{3.000}{0.569}{-0.280}{0.268}
\VF{8.000}{3.000}{0.477}{-0.305}{0.194}
\VF{9.000}{3.000}{0.399}{-0.310}{0.127}
\VF{10.000}{3.000}{0.333}{-0.303}{0.066}
\VF{-10.000}{4.000}{-0.387}{-0.219}{0.060}
\VF{-9.000}{4.000}{-0.453}{-0.206}{0.118}
\VF{-8.000}{4.000}{-0.527}{-0.176}{0.182}
\VF{-7.000}{4.000}{-0.606}{-0.123}{0.251}
\VF{-6.000}{4.000}{-0.687}{-0.038}{0.327}
\VF{-5.000}{4.000}{-0.757}{0.088}{0.409}
\VF{-4.000}{4.000}{-0.798}{0.266}{0.496}
\VF{-3.000}{4.000}{-0.777}{0.496}{0.584}
\VF{-2.000}{4.000}{-0.648}{0.756}{0.665}
\VF{-1.000}{4.000}{-0.379}{0.979}{0.725}
\VF{0.000}{4.000}{0.000}{1.070}{0.747}
\VF{1.000}{4.000}{0.379}{0.979}{0.725}
\VF{2.000}{4.000}{0.648}{0.756}{0.665}
\VF{3.000}{4.000}{0.777}{0.496}{0.584}
\VF{4.000}{4.000}{0.798}{0.266}{0.496}
\VF{5.000}{4.000}{0.757}{0.088}{0.409}
\VF{6.000}{4.000}{0.687}{-0.038}{0.327}
\VF{7.000}{4.000}{0.606}{-0.123}{0.251}
\VF{8.000}{4.000}{0.527}{-0.176}{0.182}
\VF{9.000}{4.000}{0.453}{-0.206}{0.118}
\VF{10.000}{4.000}{0.387}{-0.219}{0.060}
\VF{-10.000}{5.000}{-0.413}{-0.138}{0.050}
\VF{-9.000}{5.000}{-0.473}{-0.109}{0.105}
\VF{-8.000}{5.000}{-0.536}{-0.063}{0.164}
\VF{-7.000}{5.000}{-0.597}{0.006}{0.227}
\VF{-6.000}{5.000}{-0.650}{0.101}{0.294}
\VF{-5.000}{5.000}{-0.685}{0.228}{0.364}
\VF{-4.000}{5.000}{-0.684}{0.388}{0.435}
\VF{-3.000}{5.000}{-0.626}{0.571}{0.503}
\VF{-2.000}{5.000}{-0.492}{0.754}{0.561}
\VF{-1.000}{5.000}{-0.274}{0.896}{0.601}
\VF{0.000}{5.000}{0.000}{0.950}{0.615}
\VF{1.000}{5.000}{0.274}{0.896}{0.601}
\VF{2.000}{5.000}{0.492}{0.754}{0.561}
\VF{3.000}{5.000}{0.626}{0.571}{0.503}
\VF{4.000}{5.000}{0.684}{0.388}{0.435}
\VF{5.000}{5.000}{0.685}{0.228}{0.364}
\VF{6.000}{5.000}{0.650}{0.101}{0.294}
\VF{7.000}{5.000}{0.597}{0.006}{0.227}
\VF{8.000}{5.000}{0.536}{-0.063}{0.164}
\VF{9.000}{5.000}{0.473}{-0.109}{0.105}
\VF{10.000}{5.000}{0.413}{-0.138}{0.050}
\VF{-10.000}{6.000}{-0.418}{-0.065}{0.037}
\VF{-9.000}{6.000}{-0.469}{-0.026}{0.088}
\VF{-8.000}{6.000}{-0.518}{0.029}{0.142}
\VF{-7.000}{6.000}{-0.561}{0.102}{0.198}
\VF{-6.000}{6.000}{-0.592}{0.197}{0.257}
\VF{-5.000}{6.000}{-0.600}{0.314}{0.316}
\VF{-4.000}{6.000}{-0.576}{0.448}{0.373}
\VF{-3.000}{6.000}{-0.506}{0.590}{0.426}
\VF{-2.000}{6.000}{-0.382}{0.722}{0.469}
\VF{-1.000}{6.000}{-0.207}{0.817}{0.497}
\VF{0.000}{6.000}{0.000}{0.852}{0.508}
\VF{1.000}{6.000}{0.207}{0.817}{0.497}
\VF{2.000}{6.000}{0.382}{0.722}{0.469}
\VF{3.000}{6.000}{0.506}{0.590}{0.426}
\VF{4.000}{6.000}{0.576}{0.448}{0.373}
\VF{5.000}{6.000}{0.600}{0.314}{0.316}
\VF{6.000}{6.000}{0.592}{0.197}{0.257}
\VF{7.000}{6.000}{0.561}{0.102}{0.198}
\VF{8.000}{6.000}{0.518}{0.029}{0.142}
\VF{9.000}{6.000}{0.469}{-0.026}{0.088}
\VF{10.000}{6.000}{0.418}{-0.065}{0.037}
\VF{-10.000}{7.000}{-0.408}{-0.004}{0.020}
\VF{-9.000}{7.000}{-0.449}{0.040}{0.067}
\VF{-8.000}{7.000}{-0.485}{0.098}{0.116}
\VF{-7.000}{7.000}{-0.513}{0.171}{0.166}
\VF{-6.000}{7.000}{-0.527}{0.259}{0.216}
\VF{-5.000}{7.000}{-0.519}{0.361}{0.266}
\VF{-4.000}{7.000}{-0.483}{0.471}{0.313}
\VF{-3.000}{7.000}{-0.412}{0.582}{0.355}
\VF{-2.000}{7.000}{-0.303}{0.678}{0.388}
\VF{-1.000}{7.000}{-0.161}{0.745}{0.409}
\VF{0.000}{7.000}{0.000}{0.769}{0.417}
\VF{1.000}{7.000}{0.161}{0.745}{0.409}
\VF{2.000}{7.000}{0.303}{0.678}{0.388}
\VF{3.000}{7.000}{0.412}{0.582}{0.355}
\VF{4.000}{7.000}{0.483}{0.471}{0.313}
\VF{5.000}{7.000}{0.519}{0.361}{0.266}
\VF{6.000}{7.000}{0.527}{0.259}{0.216}
\VF{7.000}{7.000}{0.513}{0.171}{0.166}
\VF{8.000}{7.000}{0.485}{0.098}{0.116}
\VF{9.000}{7.000}{0.449}{0.040}{0.067}
\VF{10.000}{7.000}{0.408}{-0.004}{0.020}
\VF{-10.000}{8.000}{-0.387}{0.045}{0.000}
\VF{-9.000}{8.000}{-0.419}{0.091}{0.043}
\VF{-8.000}{8.000}{-0.445}{0.148}{0.087}
\VF{-7.000}{8.000}{-0.461}{0.217}{0.131}
\VF{-6.000}{8.000}{-0.463}{0.296}{0.175}
\VF{-5.000}{8.000}{-0.446}{0.383}{0.217}
\VF{-4.000}{8.000}{-0.405}{0.473}{0.256}
\VF{-3.000}{8.000}{-0.338}{0.559}{0.289}
\VF{-2.000}{8.000}{-0.244}{0.631}{0.315}
\VF{-1.000}{8.000}{-0.129}{0.680}{0.332}
\VF{0.000}{8.000}{0.000}{0.697}{0.338}
\VF{1.000}{8.000}{0.129}{0.680}{0.332}
\VF{2.000}{8.000}{0.244}{0.631}{0.315}
\VF{3.000}{8.000}{0.338}{0.559}{0.289}
\VF{4.000}{8.000}{0.405}{0.473}{0.256}
\VF{5.000}{8.000}{0.446}{0.383}{0.217}
\VF{6.000}{8.000}{0.463}{0.296}{0.175}
\VF{7.000}{8.000}{0.461}{0.217}{0.131}
\VF{8.000}{8.000}{0.445}{0.148}{0.087}
\VF{9.000}{8.000}{0.419}{0.091}{0.043}
\VF{10.000}{8.000}{0.387}{0.045}{0.000}
\VF{-10.000}{9.000}{-0.380}{0.087}{0.000}
\VF{-9.000}{9.000}{-0.385}{0.128}{0.018}
\VF{-8.000}{9.000}{-0.402}{0.183}{0.057}
\VF{-7.000}{9.000}{-0.410}{0.245}{0.096}
\VF{-6.000}{9.000}{-0.404}{0.314}{0.134}
\VF{-5.000}{9.000}{-0.382}{0.388}{0.170}
\VF{-4.000}{9.000}{-0.341}{0.461}{0.202}
\VF{-3.000}{9.000}{-0.280}{0.529}{0.230}
\VF{-2.000}{9.000}{-0.200}{0.585}{0.251}
\VF{-1.000}{9.000}{-0.104}{0.621}{0.264}
\VF{0.000}{9.000}{0.000}{0.634}{0.268}
\VF{1.000}{9.000}{0.104}{0.621}{0.264}
\VF{2.000}{9.000}{0.200}{0.585}{0.251}
\VF{3.000}{9.000}{0.280}{0.529}{0.230}
\VF{4.000}{9.000}{0.341}{0.461}{0.202}
\VF{5.000}{9.000}{0.382}{0.388}{0.170}
\VF{6.000}{9.000}{0.404}{0.314}{0.134}
\VF{7.000}{9.000}{0.410}{0.245}{0.096}
\VF{8.000}{9.000}{0.402}{0.183}{0.057}
\VF{9.000}{9.000}{0.385}{0.128}{0.018}
\VF{10.000}{9.000}{0.380}{0.087}{0.000}
\VF{-10.000}{10.000}{-0.370}{0.123}{0.000}
\VF{-9.000}{10.000}{-0.357}{0.157}{0.000}
\VF{-8.000}{10.000}{-0.360}{0.204}{0.026}
\VF{-7.000}{10.000}{-0.361}{0.260}{0.061}
\VF{-6.000}{10.000}{-0.351}{0.320}{0.094}
\VF{-5.000}{10.000}{-0.327}{0.382}{0.124}
\VF{-4.000}{10.000}{-0.288}{0.442}{0.152}
\VF{-3.000}{10.000}{-0.234}{0.496}{0.174}
\VF{-2.000}{10.000}{-0.165}{0.540}{0.192}
\VF{-1.000}{10.000}{-0.086}{0.568}{0.202}
\VF{0.000}{10.000}{0.000}{0.578}{0.206}
\VF{1.000}{10.000}{0.086}{0.568}{0.202}
\VF{2.000}{10.000}{0.165}{0.540}{0.192}
\VF{3.000}{10.000}{0.234}{0.496}{0.174}
\VF{4.000}{10.000}{0.288}{0.442}{0.152}
\VF{5.000}{10.000}{0.327}{0.382}{0.124}
\VF{6.000}{10.000}{0.351}{0.320}{0.094}
\VF{7.000}{10.000}{0.361}{0.260}{0.061}
\VF{8.000}{10.000}{0.360}{0.204}{0.026}
\VF{9.000}{10.000}{0.357}{0.157}{0.000}
\VF{10.000}{10.000}{0.370}{0.123}{0.000}
}
  \marcoMMA{10}{5}{$x/\ell$}{$z/\ell$}
\end{tikzpicture}\\[1pt]{\small\scshape (c)}
\end{minipage}
\caption{Dipolo eléctrico puntual $\vp=q\vl=p\,\uz$, en el plano meridiano $y=0$ (longitudes en unidades de $\ell$).
(a)~Equipotenciales $\phi=\text{cte}$, a intervalos iguales de potencial: se apiñan hacia el plano ecuatorial $z=0$, que es la
equipotencial $\phi=0$. (b)~Equipotenciales (gris) y líneas de campo (azul) con su sentido: salen del polo positivo, se curvan y
cortan ortogonalmente a las equipotenciales. (c)~Campo eléctrico: la longitud y el color de cada flecha crecen con
$\norm{\vE}$ (escala logarítmica, de azul a amarillo). El disco blanco central excluye la región $r\lesssim\ell$, donde la
aproximación de dipolo puntual pierde sentido. Recreación de~\cite{FIGURAS-T2}; las líneas de campo, con separación uniforme como en Mathematica, se han
calculado con el algoritmo de Jobard y Lefer~\cite{JOBARD}.}
\label{fig:dipolo-mma}
\end{figure}


% -------------------------------------------------------------------------
\subsection{Dipolo físico frente a dipolo puntual}\label{sec:fisico-puntual}
% -------------------------------------------------------------------------
Un dipolo real tiene una separación $\ell$ finita. ¿Hasta qué punto se parece su campo al del dipolo puntual? Fuera de una región de
tamaño comparable a $\ell$, las líneas de campo de ambos son prácticamente indistinguibles (\cref{fig:esquema-dipolo,fig:fisico-puntual}); solo
cerca de las cargas el dipolo físico revela su estructura. La siguiente proposición lo cuantifica.

% =========================================================================
% Figura (basada en «Tema 2. Figuras», diapositiva «Dipolo físico y dipolo
% puntual»): líneas de campo completas de un dipolo físico y de un dipolo
% puntual del mismo momento, como curvas de nivel de la función de corriente
% de Stokes (flujo constante entre dos líneas consecutivas). La circunferencia
% a trazos, de radio 1,5 l, marca la región fuera de la cual ambos campos son
% prácticamente indistinguibles. Datos: figuras/datos/esq_dip_*_lc.tex.
% =========================================================================
\begin{figure}[ht]
\centering
\newcommand{\esquema}[3]{% datos, eje (líneas Psi = 0), fuentes
  \begin{tikzpicture}[x=0.82cm, y=0.82cm]
    \renewcommand{\LC}[1]{\draw[magacero, line width=0.5pt] plot coordinates {##1};}%
    \fill[magfondo] (-3,-3) rectangle (3,3);
    \begin{scope}\clip (-3,-3) rectangle (3,3);
      \input{figuras/datos/#1}
      #2
    \end{scope}
    \draw[grispizarra, densely dashed, line width=0.6pt] (0,0) circle (1.5);
    \node[font=\scriptsize, text=grispizarra, fill=magfondo, fill opacity=0.85, text opacity=1, inner sep=1pt]
      at (-45:1.5) {$r=1{,}5\,\ell$};
    #3
    \draw[grispizarra!60, line width=0.8pt] (-3,-3) rectangle (3,3);
  \end{tikzpicture}}
\begin{minipage}[t]{0.47\textwidth}\centering
\esquema{esq_dip_fisico_lc.tex}%
  {\draw[magacero, line width=0.5pt] (0,0.5) -- (0,3) (0,-0.5) -- (0,-3);
   \draw[magacero, line width=0.8pt] (0,0.5) -- (0,-0.5);}%
  {\node[cargapos, minimum size=7pt, font=\tiny\bfseries] at (0,0.5) {$+$};
   \node[carganeg, minimum size=7pt, font=\tiny\bfseries] at (0,-0.5) {$-$};}
\\[2pt]{\small\scshape (a)\; dipolo físico}
\end{minipage}\hfill
\begin{minipage}[t]{0.47\textwidth}\centering
\esquema{esq_dip_puntual_lc.tex}%
  {\draw[magacero, line width=0.5pt] (0,0.05) -- (0,3) (0,-0.05) -- (0,-3);}%
  {\fill[white] (0,0) circle (0.36);
   \draw[magacero!60, line width=0.4pt] (0,0) circle (0.36);
   \draw[-{Stealth[length=3.4pt,width=3pt]}, rojoteo, line width=1.1pt] (0,-0.26) -- (0,0.28);
   \node[font=\scriptsize, text=rojoteo, inner sep=0.5pt] at (0.16,-0.08) [anchor=west] {$\vp$};}
\\[2pt]{\small\scshape (b)\; dipolo puntual}
\end{minipage}
\caption{Líneas de campo de (a)~un dipolo físico ($\pm q$ separadas $\ell$) y (b)~un dipolo puntual del mismo momento. Las líneas son las curvas
$\Psi=\text{cte}$ de la función de corriente, con $\Psi_{f}=\cos\alpha_{-}-\cos\alpha_{+}$ (ángulos medidos desde cada carga) y $\Psi_{p}=\ell\sen^{2}\theta/r$, y se trazan
con los mismos valores en ambos paneles. Dentro de la circunferencia de radio $1{,}5\,\ell$ el dipolo físico revela su estructura (las líneas nacen y mueren en las
cargas), mientras que las del dipolo puntual convergen todas en el origen (el pequeño disco blanco solo marca su posición); fuera de ella, ambas familias son prácticamente indistinguibles. Basada en~\cite{FIGURAS-T2}.}
\label{fig:esquema-dipolo}
\end{figure}


\begin{proposicion}[Discrepancia entre el dipolo físico y el puntual]\label{prop:discrepancia}
Sea un dipolo físico con $+q$ en $z=\ell/2$ y $-q$ en $z=-\ell/2$, y sea $\vE_{p}$ el campo del dipolo puntual equivalente, $\vp=q\ell\,\uz$, en el
origen. Definimos la discrepancia relativa
\[
  \varepsilon(r,\theta)=\frac{\norm{\vE_{p}-\vE_{f}}}{\norm{\vE_{f}}},
\]
donde $\vE_{f}$ es el campo exacto del dipolo físico. Sobre el eje ($\theta=0$) y en el plano ecuatorial ($\theta=\pi/2$), para $r>\ell/2$,
\begin{equation}\label{eq:discrepancia}
\begin{aligned}
  \varepsilon(r,0)&=1-\Bigl(1-\frac{\ell^{2}}{4r^{2}}\Bigr)^{2}=\frac{\ell^{2}}{2r^{2}}-\frac{\ell^{4}}{16r^{4}},\\
  \varepsilon\Bigl(r,\frac{\pi}{2}\Bigr)&=\Bigl(1+\frac{\ell^{2}}{4r^{2}}\Bigr)^{3/2}-1=\frac{3\ell^{2}}{8r^{2}}+\ord\Bigl(\frac{\ell^{4}}{r^{4}}\Bigr).
\end{aligned}
\end{equation}
En particular, la discrepancia es de orden $(\ell/r)^{2}$ y se anula cuando $r/\ell\to\infty$.
\end{proposicion}
\begin{proof}
\emph{Eje.} En $\vc{r}=r\,\uz$ los dos campos son paralelos a $\uz$. El del dipolo físico es
\[
  E_{f}=\frac{q}{4\pi\eo}\Bigl[\frac{1}{(r-\ell/2)^{2}}-\frac{1}{(r+\ell/2)^{2}}\Bigr]=\frac{q}{4\pi\eo}\,\frac{2r\ell}{(r^{2}-\ell^{2}/4)^{2}},
\]
y el del puntual, por \cref{obs:estructura}, $E_{p}=2q\ell/(4\pi\eo r^{3})$. Entonces $E_{p}/E_{f}=(r^{2}-\ell^{2}/4)^{2}/r^{4}=(1-\ell^{2}/4r^{2})^{2}<1$ y
$\varepsilon=1-E_{p}/E_{f}$.

\emph{Ecuador.} En $\vc{r}=r\,\ux$, las dos cargas están a la distancia $(r^{2}+\ell^{2}/4)^{1/2}$; las componentes $x$ de sus campos se cancelan y
las componentes $z$ se suman, de modo que $\vE_{f}=-\frac{q}{4\pi\eo}\,\frac{\ell}{(r^{2}+\ell^{2}/4)^{3/2}}\,\uz$, mientras que
$\vE_{p}=-\frac{q\ell}{4\pi\eo r^{3}}\,\uz$. Por tanto $\varepsilon=E_{p}/E_{f}-1=(1+\ell^{2}/4r^{2})^{3/2}-1$. Los desarrollos se obtienen con el binomio
de Newton.
\end{proof}

% =========================================================================
% Figura (clon de «Tema 2. Figuras»): equipotenciales y líneas de campo del
% dipolo puntual (izquierda) y del dipolo físico (derecha), como StreamPlot
% de Mathematica. Datos: figuras/datos/dip_{puntual,fisico}_*.tex.
% =========================================================================
\begin{figure}[ht]
\centering
\begin{minipage}[t]{0.49\textwidth}\centering
\begin{tikzpicture}[x=0.29cm, y=0.29cm]
  \panelMMA{10}{mmafondo}{1.2}{%
    % Dipolo puntual - equipotenciales
% Generado por scripts/generar_datos.py (no editar a mano).
\EQ{(1.224,-0.476) (1.250,-0.522) (1.275,-0.574) (1.296,-0.625) (1.313,-0.675) (1.326,-0.725) (1.337,-0.775) (1.345,-0.825) (1.350,-0.875) (1.353,-0.925) (1.354,-0.975) (1.352,-1.025) (1.347,-1.075) (1.340,-1.125) (1.331,-1.175) (1.319,-1.225) (1.305,-1.275) (1.288,-1.325) (1.269,-1.375) (1.246,-1.425) (1.221,-1.475) (1.194,-1.525) (1.162,-1.575) (1.128,-1.625) (1.089,-1.675) (1.050,-1.721) (1.000,-1.774) (0.950,-1.822) (0.900,-1.865) (0.856,-1.900) (0.800,-1.940) (0.750,-1.972) (0.702,-2.000) (0.656,-2.025) (0.604,-2.050) (0.550,-2.074) (0.500,-2.093) (0.450,-2.111) (0.400,-2.126) (0.350,-2.139) (0.300,-2.151) (0.250,-2.161) (0.200,-2.168) (0.150,-2.174) (0.100,-2.179) (0.050,-2.181) (0.000,-2.182) (-0.050,-2.181) (-0.100,-2.179) (-0.150,-2.174) (-0.200,-2.168) (-0.250,-2.161) (-0.300,-2.151) (-0.350,-2.139) (-0.400,-2.126) (-0.450,-2.111) (-0.500,-2.093) (-0.547,-2.075) (-0.600,-2.052) (-0.650,-2.028) (-0.700,-2.001) (-0.745,-1.975) (-0.800,-1.940) (-0.850,-1.904) (-0.900,-1.865) (-0.947,-1.825) (-0.999,-1.775) (-1.047,-1.725) (-1.089,-1.675) (-1.125,-1.629) (-1.162,-1.575) (-1.194,-1.525) (-1.221,-1.475) (-1.246,-1.425) (-1.269,-1.375) (-1.288,-1.325) (-1.305,-1.275) (-1.319,-1.225) (-1.331,-1.175) (-1.340,-1.125) (-1.347,-1.075) (-1.352,-1.025) (-1.354,-0.975) (-1.353,-0.925) (-1.350,-0.875) (-1.345,-0.825) (-1.337,-0.775) (-1.326,-0.725) (-1.313,-0.675) (-1.296,-0.625) (-1.275,-0.575) (-1.252,-0.525) (-1.225,-0.477)}
\EQ{(1.229,-0.446) (1.263,-0.500) (1.290,-0.550) (1.313,-0.600) (1.333,-0.650) (1.349,-0.700) (1.362,-0.750) (1.373,-0.800) (1.381,-0.850) (1.386,-0.900) (1.389,-0.950) (1.389,-1.000) (1.387,-1.050) (1.383,-1.100) (1.376,-1.150) (1.367,-1.200) (1.355,-1.250) (1.341,-1.300) (1.325,-1.349) (1.306,-1.400) (1.284,-1.450) (1.260,-1.500) (1.233,-1.550) (1.203,-1.600) (1.175,-1.642) (1.132,-1.700) (1.100,-1.740) (1.050,-1.796) (1.000,-1.846) (0.950,-1.891) (0.909,-1.925) (0.850,-1.970) (0.806,-2.000) (0.750,-2.035) (0.700,-2.063) (0.650,-2.089) (0.600,-2.112) (0.550,-2.133) (0.500,-2.152) (0.450,-2.169) (0.400,-2.184) (0.350,-2.197) (0.300,-2.208) (0.250,-2.218) (0.200,-2.225) (0.150,-2.231) (0.100,-2.236) (0.050,-2.238) (0.000,-2.239) (-0.050,-2.238) (-0.100,-2.236) (-0.150,-2.231) (-0.200,-2.225) (-0.250,-2.218) (-0.300,-2.208) (-0.350,-2.197) (-0.400,-2.184) (-0.450,-2.169) (-0.500,-2.152) (-0.550,-2.133) (-0.600,-2.112) (-0.650,-2.089) (-0.700,-2.063) (-0.750,-2.035) (-0.800,-2.004) (-0.843,-1.975) (-0.900,-1.932) (-0.940,-1.900) (-0.996,-1.850) (-1.046,-1.800) (-1.091,-1.750) (-1.125,-1.709) (-1.169,-1.650) (-1.200,-1.604) (-1.233,-1.550) (-1.260,-1.500) (-1.284,-1.450) (-1.306,-1.400) (-1.325,-1.350) (-1.341,-1.300) (-1.355,-1.250) (-1.367,-1.200) (-1.376,-1.150) (-1.383,-1.100) (-1.387,-1.050) (-1.389,-1.000) (-1.389,-0.950) (-1.386,-0.900) (-1.381,-0.850) (-1.373,-0.800) (-1.362,-0.750) (-1.349,-0.700) (-1.333,-0.650) (-1.313,-0.600) (-1.290,-0.550) (-1.263,-0.500) (-1.232,-0.450) (-1.229,-0.446)}
\EQ{(1.244,-0.431) (1.274,-0.475) (1.300,-0.518) (1.325,-0.565) (1.350,-0.619) (1.372,-0.675) (1.388,-0.725) (1.400,-0.775) (1.411,-0.825) (1.419,-0.875) (1.424,-0.925) (1.427,-0.975) (1.427,-1.025) (1.425,-1.075) (1.421,-1.125) (1.414,-1.175) (1.406,-1.225) (1.395,-1.275) (1.381,-1.325) (1.365,-1.375) (1.347,-1.425) (1.326,-1.475) (1.303,-1.525) (1.277,-1.575) (1.250,-1.621) (1.215,-1.675) (1.180,-1.725) (1.141,-1.775) (1.100,-1.822) (1.050,-1.874) (1.000,-1.922) (0.950,-1.965) (0.906,-2.000) (0.850,-2.040) (0.800,-2.073) (0.755,-2.100) (0.709,-2.125) (0.660,-2.150) (0.605,-2.175) (0.550,-2.198) (0.500,-2.216) (0.450,-2.233) (0.400,-2.247) (0.350,-2.260) (0.300,-2.271) (0.250,-2.280) (0.200,-2.287) (0.150,-2.293) (0.100,-2.297) (0.050,-2.299) (0.000,-2.300) (-0.050,-2.299) (-0.100,-2.297) (-0.150,-2.293) (-0.200,-2.287) (-0.250,-2.280) (-0.300,-2.271) (-0.350,-2.260) (-0.400,-2.247) (-0.450,-2.233) (-0.500,-2.216) (-0.550,-2.198) (-0.600,-2.177) (-0.650,-2.155) (-0.700,-2.130) (-0.750,-2.103) (-0.796,-2.075) (-0.850,-2.040) (-0.900,-2.004) (-0.950,-1.965) (-0.997,-1.925) (-1.049,-1.875) (-1.097,-1.825) (-1.141,-1.775) (-1.175,-1.731) (-1.215,-1.675) (-1.248,-1.625) (-1.275,-1.578) (-1.300,-1.531) (-1.325,-1.478) (-1.347,-1.425) (-1.365,-1.375) (-1.381,-1.325) (-1.395,-1.275) (-1.406,-1.225) (-1.414,-1.175) (-1.421,-1.125) (-1.425,-1.075) (-1.427,-1.025) (-1.427,-0.975) (-1.424,-0.925) (-1.419,-0.875) (-1.411,-0.825) (-1.400,-0.775) (-1.388,-0.725) (-1.372,-0.675) (-1.352,-0.625) (-1.330,-0.575) (-1.304,-0.525) (-1.275,-0.477) (-1.244,-0.431)}
\EQ{(1.250,-0.405) (1.285,-0.450) (1.318,-0.500) (1.347,-0.550) (1.373,-0.600) (1.395,-0.650) (1.413,-0.700) (1.429,-0.750) (1.442,-0.800) (1.452,-0.850) (1.460,-0.900) (1.465,-0.950) (1.468,-1.000) (1.468,-1.050) (1.467,-1.100) (1.463,-1.150) (1.457,-1.200) (1.448,-1.250) (1.438,-1.300) (1.425,-1.350) (1.410,-1.400) (1.392,-1.450) (1.372,-1.500) (1.350,-1.550) (1.325,-1.600) (1.300,-1.645) (1.266,-1.700) (1.232,-1.750) (1.200,-1.793) (1.154,-1.850) (1.109,-1.900) (1.059,-1.950) (1.004,-2.000) (0.950,-2.044) (0.909,-2.075) (0.850,-2.116) (0.800,-2.147) (0.751,-2.175) (0.704,-2.200) (0.652,-2.225) (0.600,-2.248) (0.550,-2.268) (0.500,-2.285) (0.450,-2.301) (0.400,-2.315) (0.350,-2.328) (0.300,-2.338) (0.250,-2.347) (0.200,-2.354) (0.150,-2.360) (0.100,-2.364) (0.050,-2.366) (0.000,-2.367) (-0.050,-2.366) (-0.100,-2.364) (-0.150,-2.360) (-0.200,-2.354) (-0.250,-2.347) (-0.300,-2.338) (-0.350,-2.328) (-0.400,-2.315) (-0.450,-2.301) (-0.500,-2.285) (-0.550,-2.268) (-0.600,-2.248) (-0.650,-2.226) (-0.700,-2.202) (-0.750,-2.176) (-0.795,-2.150) (-0.850,-2.116) (-0.900,-2.081) (-0.942,-2.050) (-1.000,-2.003) (-1.050,-1.958) (-1.100,-1.909) (-1.150,-1.854) (-1.195,-1.800) (-1.225,-1.760) (-1.266,-1.700) (-1.297,-1.650) (-1.325,-1.600) (-1.350,-1.550) (-1.372,-1.500) (-1.392,-1.450) (-1.410,-1.400) (-1.425,-1.350) (-1.438,-1.300) (-1.448,-1.250) (-1.457,-1.200) (-1.463,-1.150) (-1.467,-1.100) (-1.468,-1.050) (-1.468,-1.000) (-1.465,-0.950) (-1.460,-0.900) (-1.452,-0.850) (-1.442,-0.800) (-1.429,-0.750) (-1.413,-0.700) (-1.395,-0.650) (-1.373,-0.600) (-1.350,-0.555) (-1.325,-0.512) (-1.285,-0.450) (-1.250,-0.405)}
\EQ{(1.251,-0.374) (1.295,-0.425) (1.325,-0.465) (1.365,-0.525) (1.393,-0.575) (1.418,-0.625) (1.440,-0.675) (1.458,-0.725) (1.473,-0.775) (1.486,-0.825) (1.496,-0.875) (1.504,-0.925) (1.510,-0.975) (1.513,-1.025) (1.514,-1.075) (1.512,-1.125) (1.509,-1.175) (1.503,-1.225) (1.495,-1.275) (1.486,-1.325) (1.473,-1.375) (1.459,-1.425) (1.442,-1.475) (1.423,-1.525) (1.402,-1.575) (1.378,-1.625) (1.352,-1.675) (1.325,-1.721) (1.290,-1.775) (1.255,-1.825) (1.216,-1.875) (1.175,-1.923) (1.127,-1.975) (1.075,-2.025) (1.025,-2.070) (0.975,-2.110) (0.925,-2.147) (0.875,-2.182) (0.825,-2.213) (0.775,-2.242) (0.725,-2.268) (0.675,-2.292) (0.625,-2.314) (0.575,-2.334) (0.525,-2.352) (0.475,-2.369) (0.425,-2.383) (0.375,-2.396) (0.325,-2.407) (0.275,-2.416) (0.225,-2.424) (0.175,-2.430) (0.125,-2.435) (0.075,-2.438) (0.025,-2.440) (-0.050,-2.439) (-0.100,-2.437) (-0.150,-2.433) (-0.200,-2.427) (-0.250,-2.420) (-0.300,-2.412) (-0.350,-2.402) (-0.400,-2.390) (-0.450,-2.376) (-0.500,-2.361) (-0.550,-2.344) (-0.599,-2.325) (-0.650,-2.304) (-0.700,-2.280) (-0.750,-2.255) (-0.800,-2.228) (-0.846,-2.200) (-0.900,-2.165) (-0.950,-2.129) (-1.000,-2.090) (-1.048,-2.050) (-1.100,-2.002) (-1.150,-1.951) (-1.195,-1.900) (-1.236,-1.850) (-1.273,-1.800) (-1.307,-1.750) (-1.337,-1.700) (-1.365,-1.650) (-1.390,-1.600) (-1.413,-1.550) (-1.433,-1.500) (-1.450,-1.453) (-1.466,-1.400) (-1.480,-1.350) (-1.491,-1.300) (-1.500,-1.250) (-1.506,-1.200) (-1.511,-1.150) (-1.513,-1.100) (-1.514,-1.050) (-1.512,-1.000) (-1.507,-0.950) (-1.501,-0.900) (-1.492,-0.850) (-1.480,-0.800) (-1.466,-0.750) (-1.449,-0.700) (-1.429,-0.650) (-1.406,-0.600) (-1.379,-0.550) (-1.350,-0.502) (-1.314,-0.450) (-1.275,-0.401) (-1.251,-0.374)}
\EQ{(1.266,-0.359) (1.305,-0.400) (1.346,-0.450) (1.382,-0.500) (1.414,-0.550) (1.442,-0.600) (1.467,-0.650) (1.488,-0.700) (1.506,-0.750) (1.522,-0.800) (1.534,-0.850) (1.545,-0.900) (1.553,-0.950) (1.559,-1.000) (1.562,-1.050) (1.563,-1.100) (1.563,-1.150) (1.560,-1.200) (1.555,-1.250) (1.547,-1.300) (1.538,-1.350) (1.527,-1.400) (1.513,-1.450) (1.498,-1.500) (1.480,-1.550) (1.460,-1.600) (1.437,-1.650) (1.412,-1.700) (1.385,-1.750) (1.354,-1.800) (1.325,-1.844) (1.285,-1.900) (1.250,-1.944) (1.201,-2.000) (1.154,-2.050) (1.101,-2.100) (1.050,-2.145) (1.000,-2.185) (0.950,-2.222) (0.908,-2.250) (0.850,-2.287) (0.800,-2.315) (0.750,-2.342) (0.700,-2.366) (0.650,-2.388) (0.600,-2.409) (0.550,-2.427) (0.500,-2.443) (0.450,-2.458) (0.400,-2.471) (0.350,-2.483) (0.300,-2.493) (0.250,-2.501) (0.200,-2.508) (0.150,-2.513) (0.100,-2.517) (0.050,-2.519) (0.000,-2.520) (-0.050,-2.519) (-0.100,-2.517) (-0.150,-2.513) (-0.200,-2.508) (-0.250,-2.501) (-0.300,-2.493) (-0.350,-2.483) (-0.400,-2.471) (-0.450,-2.458) (-0.500,-2.443) (-0.550,-2.427) (-0.600,-2.409) (-0.650,-2.388) (-0.700,-2.366) (-0.750,-2.342) (-0.800,-2.315) (-0.850,-2.287) (-0.900,-2.256) (-0.945,-2.225) (-1.000,-2.185) (-1.044,-2.150) (-1.100,-2.101) (-1.150,-2.054) (-1.200,-2.001) (-1.245,-1.950) (-1.285,-1.900) (-1.321,-1.850) (-1.350,-1.807) (-1.385,-1.750) (-1.412,-1.700) (-1.437,-1.650) (-1.460,-1.600) (-1.480,-1.550) (-1.498,-1.500) (-1.513,-1.450) (-1.527,-1.400) (-1.538,-1.350) (-1.547,-1.300) (-1.555,-1.250) (-1.560,-1.200) (-1.563,-1.150) (-1.563,-1.100) (-1.562,-1.050) (-1.559,-1.000) (-1.553,-0.950) (-1.545,-0.900) (-1.534,-0.850) (-1.522,-0.800) (-1.506,-0.750) (-1.488,-0.700) (-1.467,-0.650) (-1.442,-0.600) (-1.414,-0.550) (-1.382,-0.500) (-1.350,-0.456) (-1.305,-0.400) (-1.266,-0.359)}
\EQ{(1.275,-0.337) (1.314,-0.375) (1.350,-0.414) (1.400,-0.475) (1.436,-0.525) (1.467,-0.575) (1.495,-0.625) (1.519,-0.675) (1.540,-0.725) (1.559,-0.775) (1.574,-0.825) (1.587,-0.875) (1.598,-0.925) (1.606,-0.975) (1.612,-1.025) (1.616,-1.075) (1.618,-1.125) (1.618,-1.175) (1.616,-1.225) (1.611,-1.275) (1.605,-1.325) (1.597,-1.375) (1.586,-1.425) (1.574,-1.475) (1.559,-1.525) (1.543,-1.575) (1.524,-1.625) (1.503,-1.675) (1.479,-1.725) (1.454,-1.775) (1.425,-1.825) (1.394,-1.875) (1.360,-1.925) (1.325,-1.973) (1.283,-2.025) (1.250,-2.063) (1.216,-2.100) (1.175,-2.141) (1.125,-2.187) (1.081,-2.225) (1.025,-2.269) (0.983,-2.300) (0.925,-2.339) (0.875,-2.370) (0.825,-2.398) (0.775,-2.425) (0.725,-2.449) (0.675,-2.471) (0.625,-2.492) (0.575,-2.510) (0.525,-2.527) (0.475,-2.542) (0.425,-2.555) (0.375,-2.567) (0.325,-2.578) (0.275,-2.586) (0.225,-2.594) (0.175,-2.599) (0.125,-2.604) (0.075,-2.607) (0.025,-2.608) (-0.050,-2.608) (-0.100,-2.605) (-0.150,-2.602) (-0.200,-2.597) (-0.250,-2.590) (-0.300,-2.582) (-0.350,-2.573) (-0.400,-2.562) (-0.450,-2.549) (-0.500,-2.535) (-0.550,-2.519) (-0.600,-2.501) (-0.650,-2.482) (-0.700,-2.460) (-0.750,-2.437) (-0.800,-2.412) (-0.850,-2.384) (-0.900,-2.355) (-0.946,-2.325) (-1.000,-2.288) (-1.050,-2.250) (-1.100,-2.209) (-1.139,-2.175) (-1.191,-2.125) (-1.239,-2.075) (-1.275,-2.034) (-1.323,-1.975) (-1.360,-1.925) (-1.394,-1.875) (-1.425,-1.825) (-1.450,-1.781) (-1.475,-1.734) (-1.500,-1.681) (-1.524,-1.625) (-1.543,-1.575) (-1.559,-1.525) (-1.574,-1.475) (-1.586,-1.425) (-1.597,-1.375) (-1.605,-1.325) (-1.611,-1.275) (-1.616,-1.225) (-1.618,-1.175) (-1.618,-1.125) (-1.616,-1.075) (-1.612,-1.025) (-1.606,-0.975) (-1.598,-0.925) (-1.587,-0.875) (-1.574,-0.825) (-1.559,-0.775) (-1.540,-0.725) (-1.519,-0.675) (-1.495,-0.625) (-1.467,-0.575) (-1.436,-0.525) (-1.400,-0.475) (-1.360,-0.425) (-1.325,-0.387) (-1.289,-0.350) (-1.275,-0.337)}
\EQ{(1.275,-0.308) (1.323,-0.350) (1.374,-0.400) (1.419,-0.450) (1.450,-0.489) (1.493,-0.550) (1.524,-0.600) (1.550,-0.647) (1.575,-0.698) (1.597,-0.750) (1.616,-0.800) (1.632,-0.850) (1.645,-0.900) (1.656,-0.950) (1.665,-1.000) (1.672,-1.050) (1.676,-1.100) (1.679,-1.150) (1.679,-1.200) (1.678,-1.250) (1.674,-1.300) (1.669,-1.350) (1.662,-1.400) (1.652,-1.450) (1.641,-1.500) (1.628,-1.550) (1.613,-1.600) (1.595,-1.650) (1.576,-1.700) (1.554,-1.750) (1.530,-1.800) (1.504,-1.850) (1.475,-1.900) (1.444,-1.950) (1.410,-2.000) (1.375,-2.047) (1.333,-2.100) (1.289,-2.150) (1.250,-2.191) (1.200,-2.240) (1.162,-2.275) (1.102,-2.325) (1.050,-2.365) (1.001,-2.400) (0.950,-2.434) (0.900,-2.464) (0.850,-2.492) (0.800,-2.518) (0.750,-2.543) (0.700,-2.565) (0.650,-2.585) (0.600,-2.604) (0.550,-2.621) (0.500,-2.636) (0.450,-2.650) (0.400,-2.662) (0.350,-2.672) (0.300,-2.682) (0.250,-2.689) (0.200,-2.696) (0.150,-2.700) (0.100,-2.704) (0.050,-2.706) (0.000,-2.707) (-0.050,-2.706) (-0.100,-2.704) (-0.150,-2.700) (-0.200,-2.696) (-0.250,-2.689) (-0.300,-2.682) (-0.350,-2.672) (-0.400,-2.662) (-0.450,-2.650) (-0.500,-2.636) (-0.550,-2.621) (-0.600,-2.604) (-0.650,-2.585) (-0.700,-2.565) (-0.750,-2.543) (-0.800,-2.518) (-0.850,-2.492) (-0.900,-2.464) (-0.950,-2.434) (-1.000,-2.401) (-1.050,-2.365) (-1.100,-2.327) (-1.150,-2.285) (-1.189,-2.250) (-1.241,-2.200) (-1.289,-2.150) (-1.325,-2.109) (-1.373,-2.050) (-1.410,-2.000) (-1.444,-1.950) (-1.475,-1.900) (-1.504,-1.850) (-1.530,-1.800) (-1.554,-1.750) (-1.575,-1.702) (-1.595,-1.650) (-1.613,-1.600) (-1.628,-1.550) (-1.641,-1.500) (-1.652,-1.450) (-1.662,-1.400) (-1.669,-1.350) (-1.674,-1.300) (-1.678,-1.250) (-1.679,-1.200) (-1.679,-1.150) (-1.676,-1.100) (-1.672,-1.050) (-1.665,-1.000) (-1.656,-0.950) (-1.645,-0.900) (-1.632,-0.850) (-1.616,-0.800) (-1.597,-0.750) (-1.576,-0.700) (-1.552,-0.650) (-1.525,-0.601) (-1.493,-0.550) (-1.458,-0.500) (-1.425,-0.458) (-1.375,-0.401) (-1.325,-0.352) (-1.275,-0.308)}
\EQ{(1.275,-0.280) (1.325,-0.319) (1.375,-0.362) (1.414,-0.400) (1.450,-0.438) (1.500,-0.497) (1.539,-0.550) (1.572,-0.600) (1.600,-0.648) (1.625,-0.696) (1.650,-0.750) (1.670,-0.800) (1.688,-0.850) (1.703,-0.900) (1.716,-0.950) (1.726,-1.000) (1.734,-1.050) (1.741,-1.100) (1.745,-1.150) (1.747,-1.200) (1.748,-1.250) (1.746,-1.300) (1.743,-1.350) (1.738,-1.400) (1.731,-1.450) (1.722,-1.500) (1.711,-1.550) (1.698,-1.600) (1.683,-1.650) (1.667,-1.700) (1.648,-1.750) (1.627,-1.800) (1.604,-1.850) (1.579,-1.900) (1.552,-1.950) (1.525,-1.995) (1.489,-2.050) (1.454,-2.100) (1.416,-2.150) (1.375,-2.200) (1.330,-2.250) (1.282,-2.300) (1.229,-2.350) (1.175,-2.397) (1.125,-2.437) (1.075,-2.475) (1.025,-2.509) (0.975,-2.542) (0.925,-2.571) (0.875,-2.599) (0.825,-2.625) (0.775,-2.649) (0.725,-2.671) (0.675,-2.691) (0.625,-2.710) (0.575,-2.727) (0.525,-2.742) (0.475,-2.756) (0.425,-2.768) (0.375,-2.779) (0.325,-2.789) (0.275,-2.797) (0.225,-2.804) (0.175,-2.809) (0.125,-2.813) (0.075,-2.816) (0.025,-2.817) (-0.050,-2.817) (-0.100,-2.815) (-0.150,-2.811) (-0.200,-2.807) (-0.250,-2.800) (-0.300,-2.793) (-0.350,-2.784) (-0.400,-2.774) (-0.450,-2.762) (-0.500,-2.749) (-0.550,-2.735) (-0.600,-2.719) (-0.650,-2.701) (-0.700,-2.681) (-0.750,-2.660) (-0.800,-2.637) (-0.850,-2.612) (-0.900,-2.586) (-0.950,-2.557) (-1.000,-2.526) (-1.050,-2.492) (-1.100,-2.456) (-1.141,-2.425) (-1.200,-2.376) (-1.250,-2.331) (-1.300,-2.282) (-1.350,-2.228) (-1.396,-2.175) (-1.436,-2.125) (-1.472,-2.075) (-1.506,-2.025) (-1.537,-1.975) (-1.566,-1.925) (-1.592,-1.875) (-1.616,-1.825) (-1.638,-1.775) (-1.658,-1.725) (-1.675,-1.675) (-1.691,-1.625) (-1.705,-1.575) (-1.716,-1.525) (-1.726,-1.475) (-1.734,-1.425) (-1.741,-1.375) (-1.745,-1.325) (-1.747,-1.275) (-1.748,-1.225) (-1.746,-1.175) (-1.743,-1.125) (-1.738,-1.075) (-1.730,-1.025) (-1.721,-0.975) (-1.709,-0.925) (-1.695,-0.875) (-1.679,-0.825) (-1.660,-0.775) (-1.639,-0.725) (-1.614,-0.675) (-1.587,-0.625) (-1.556,-0.575) (-1.525,-0.530) (-1.482,-0.475) (-1.438,-0.425) (-1.400,-0.386) (-1.362,-0.350) (-1.301,-0.300) (-1.275,-0.280)}
\EQ{(1.288,-0.262) (1.342,-0.300) (1.400,-0.347) (1.450,-0.391) (1.500,-0.442) (1.550,-0.499) (1.590,-0.550) (1.625,-0.600) (1.656,-0.650) (1.684,-0.700) (1.708,-0.750) (1.730,-0.800) (1.750,-0.850) (1.767,-0.900) (1.781,-0.950) (1.794,-1.000) (1.804,-1.050) (1.812,-1.100) (1.818,-1.150) (1.823,-1.200) (1.825,-1.250) (1.825,-1.300) (1.824,-1.350) (1.821,-1.400) (1.816,-1.450) (1.810,-1.500) (1.801,-1.550) (1.791,-1.600) (1.779,-1.650) (1.765,-1.700) (1.750,-1.748) (1.732,-1.800) (1.712,-1.850) (1.691,-1.900) (1.667,-1.950) (1.641,-2.000) (1.613,-2.050) (1.583,-2.100) (1.550,-2.149) (1.514,-2.200) (1.476,-2.250) (1.434,-2.300) (1.400,-2.339) (1.366,-2.375) (1.325,-2.416) (1.289,-2.450) (1.250,-2.484) (1.201,-2.525) (1.150,-2.565) (1.101,-2.600) (1.050,-2.634) (1.000,-2.666) (0.950,-2.695) (0.900,-2.722) (0.850,-2.748) (0.800,-2.771) (0.750,-2.793) (0.700,-2.813) (0.650,-2.831) (0.600,-2.848) (0.550,-2.864) (0.500,-2.878) (0.450,-2.890) (0.400,-2.901) (0.350,-2.911) (0.300,-2.919) (0.250,-2.926) (0.200,-2.932) (0.150,-2.937) (0.100,-2.940) (0.050,-2.942) (0.000,-2.943) (-0.050,-2.942) (-0.100,-2.940) (-0.150,-2.937) (-0.200,-2.932) (-0.250,-2.926) (-0.300,-2.919) (-0.350,-2.911) (-0.400,-2.901) (-0.450,-2.890) (-0.500,-2.878) (-0.550,-2.864) (-0.600,-2.848) (-0.650,-2.831) (-0.700,-2.813) (-0.750,-2.793) (-0.800,-2.771) (-0.850,-2.748) (-0.900,-2.722) (-0.950,-2.695) (-1.000,-2.666) (-1.050,-2.634) (-1.100,-2.601) (-1.150,-2.565) (-1.200,-2.526) (-1.250,-2.484) (-1.289,-2.450) (-1.325,-2.416) (-1.366,-2.375) (-1.400,-2.339) (-1.434,-2.300) (-1.475,-2.251) (-1.514,-2.200) (-1.550,-2.150) (-1.583,-2.100) (-1.613,-2.050) (-1.641,-2.000) (-1.667,-1.950) (-1.691,-1.900) (-1.712,-1.850) (-1.732,-1.800) (-1.749,-1.750) (-1.765,-1.700) (-1.779,-1.650) (-1.791,-1.600) (-1.801,-1.550) (-1.810,-1.500) (-1.816,-1.450) (-1.821,-1.400) (-1.824,-1.350) (-1.825,-1.300) (-1.825,-1.250) (-1.823,-1.200) (-1.818,-1.150) (-1.812,-1.100) (-1.804,-1.050) (-1.794,-1.000) (-1.781,-0.950) (-1.767,-0.900) (-1.750,-0.851) (-1.730,-0.800) (-1.708,-0.750) (-1.684,-0.700) (-1.656,-0.650) (-1.625,-0.601) (-1.590,-0.550) (-1.551,-0.500) (-1.508,-0.450) (-1.459,-0.400) (-1.404,-0.350) (-1.350,-0.307) (-1.307,-0.275) (-1.288,-0.262)}
\EQ{(1.300,-0.243) (1.350,-0.275) (1.400,-0.310) (1.450,-0.348) (1.500,-0.391) (1.550,-0.439) (1.600,-0.493) (1.647,-0.550) (1.684,-0.600) (1.717,-0.650) (1.747,-0.700) (1.774,-0.750) (1.798,-0.800) (1.819,-0.850) (1.838,-0.900) (1.855,-0.950) (1.869,-1.000) (1.881,-1.050) (1.892,-1.100) (1.900,-1.150) (1.906,-1.200) (1.911,-1.250) (1.914,-1.300) (1.915,-1.350) (1.914,-1.400) (1.911,-1.450) (1.907,-1.500) (1.901,-1.550) (1.894,-1.600) (1.885,-1.650) (1.873,-1.700) (1.861,-1.750) (1.846,-1.800) (1.830,-1.850) (1.812,-1.900) (1.792,-1.950) (1.770,-2.000) (1.746,-2.050) (1.720,-2.100) (1.692,-2.150) (1.661,-2.200) (1.629,-2.250) (1.593,-2.300) (1.555,-2.350) (1.515,-2.400) (1.475,-2.445) (1.425,-2.498) (1.375,-2.548) (1.325,-2.593) (1.275,-2.636) (1.226,-2.675) (1.175,-2.713) (1.125,-2.747) (1.083,-2.775) (1.025,-2.810) (0.975,-2.838) (0.925,-2.865) (0.875,-2.890) (0.825,-2.913) (0.775,-2.934) (0.725,-2.954) (0.675,-2.972) (0.625,-2.989) (0.575,-3.004) (0.525,-3.018) (0.475,-3.030) (0.425,-3.042) (0.375,-3.052) (0.325,-3.060) (0.275,-3.068) (0.225,-3.074) (0.175,-3.079) (0.125,-3.082) (0.075,-3.085) (0.025,-3.086) (-0.050,-3.086) (-0.100,-3.084) (-0.150,-3.081) (-0.200,-3.076) (-0.250,-3.071) (-0.300,-3.064) (-0.350,-3.056) (-0.400,-3.047) (-0.450,-3.036) (-0.500,-3.024) (-0.550,-3.011) (-0.600,-2.996) (-0.650,-2.980) (-0.700,-2.963) (-0.750,-2.944) (-0.800,-2.923) (-0.850,-2.901) (-0.900,-2.877) (-0.950,-2.852) (-0.999,-2.825) (-1.050,-2.795) (-1.100,-2.764) (-1.150,-2.730) (-1.192,-2.700) (-1.250,-2.656) (-1.300,-2.615) (-1.345,-2.575) (-1.398,-2.525) (-1.448,-2.475) (-1.493,-2.425) (-1.535,-2.375) (-1.575,-2.325) (-1.611,-2.275) (-1.645,-2.225) (-1.675,-2.178) (-1.706,-2.125) (-1.733,-2.075) (-1.758,-2.025) (-1.781,-1.975) (-1.802,-1.925) (-1.821,-1.875) (-1.838,-1.825) (-1.854,-1.775) (-1.867,-1.725) (-1.879,-1.675) (-1.889,-1.625) (-1.898,-1.575) (-1.905,-1.525) (-1.910,-1.475) (-1.913,-1.425) (-1.915,-1.375) (-1.914,-1.325) (-1.913,-1.275) (-1.909,-1.225) (-1.903,-1.175) (-1.896,-1.125) (-1.887,-1.075) (-1.875,-1.025) (-1.862,-0.975) (-1.847,-0.925) (-1.829,-0.875) (-1.809,-0.825) (-1.786,-0.775) (-1.761,-0.725) (-1.733,-0.675) (-1.701,-0.625) (-1.666,-0.575) (-1.627,-0.525) (-1.584,-0.475) (-1.536,-0.425) (-1.482,-0.375) (-1.425,-0.328) (-1.375,-0.292) (-1.325,-0.258) (-1.300,-0.243)}
\EQ{(1.300,-0.216) (1.350,-0.244) (1.400,-0.274) (1.450,-0.308) (1.500,-0.345) (1.550,-0.385) (1.595,-0.425) (1.646,-0.475) (1.692,-0.525) (1.725,-0.565) (1.770,-0.625) (1.800,-0.669) (1.834,-0.725) (1.862,-0.775) (1.887,-0.825) (1.909,-0.875) (1.929,-0.925) (1.947,-0.975) (1.962,-1.025) (1.975,-1.075) (1.987,-1.125) (1.997,-1.175) (2.004,-1.225) (2.011,-1.275) (2.015,-1.325) (2.017,-1.375) (2.018,-1.425) (2.017,-1.475) (2.015,-1.525) (2.011,-1.575) (2.005,-1.625) (1.998,-1.675) (1.989,-1.725) (1.979,-1.775) (1.967,-1.825) (1.953,-1.875) (1.937,-1.925) (1.920,-1.975) (1.901,-2.025) (1.880,-2.075) (1.858,-2.125) (1.833,-2.175) (1.807,-2.225) (1.778,-2.275) (1.750,-2.321) (1.714,-2.375) (1.679,-2.425) (1.641,-2.475) (1.601,-2.525) (1.557,-2.575) (1.510,-2.625) (1.475,-2.661) (1.434,-2.700) (1.378,-2.750) (1.325,-2.794) (1.285,-2.825) (1.225,-2.869) (1.180,-2.900) (1.125,-2.935) (1.075,-2.965) (1.025,-2.993) (0.975,-3.020) (0.925,-3.045) (0.875,-3.068) (0.825,-3.089) (0.775,-3.109) (0.725,-3.128) (0.675,-3.145) (0.625,-3.161) (0.576,-3.175) (0.525,-3.188) (0.476,-3.200) (0.425,-3.211) (0.375,-3.220) (0.325,-3.229) (0.275,-3.236) (0.225,-3.241) (0.175,-3.246) (0.125,-3.249) (0.075,-3.252) (0.025,-3.253) (-0.050,-3.252) (-0.100,-3.251) (-0.150,-3.248) (-0.200,-3.244) (-0.250,-3.239) (-0.300,-3.232) (-0.350,-3.225) (-0.400,-3.216) (-0.450,-3.206) (-0.500,-3.195) (-0.550,-3.182) (-0.600,-3.168) (-0.650,-3.153) (-0.700,-3.137) (-0.750,-3.119) (-0.798,-3.100) (-0.850,-3.079) (-0.900,-3.056) (-0.950,-3.032) (-1.000,-3.007) (-1.050,-2.980) (-1.100,-2.950) (-1.150,-2.919) (-1.200,-2.886) (-1.250,-2.851) (-1.300,-2.814) (-1.348,-2.775) (-1.400,-2.731) (-1.450,-2.685) (-1.486,-2.650) (-1.525,-2.610) (-1.575,-2.555) (-1.621,-2.500) (-1.661,-2.450) (-1.697,-2.400) (-1.731,-2.350) (-1.763,-2.300) (-1.793,-2.250) (-1.820,-2.200) (-1.846,-2.150) (-1.869,-2.100) (-1.891,-2.050) (-1.911,-2.000) (-1.929,-1.950) (-1.945,-1.900) (-1.960,-1.850) (-1.973,-1.800) (-1.984,-1.750) (-1.994,-1.700) (-2.002,-1.650) (-2.008,-1.600) (-2.013,-1.550) (-2.017,-1.500) (-2.018,-1.450) (-2.018,-1.400) (-2.016,-1.350) (-2.013,-1.300) (-2.008,-1.250) (-2.001,-1.200) (-1.992,-1.150) (-1.982,-1.100) (-1.969,-1.050) (-1.955,-1.000) (-1.938,-0.950) (-1.919,-0.900) (-1.898,-0.850) (-1.875,-0.801) (-1.850,-0.753) (-1.825,-0.710) (-1.787,-0.650) (-1.752,-0.600) (-1.713,-0.550) (-1.675,-0.506) (-1.625,-0.454) (-1.575,-0.407) (-1.525,-0.364) (-1.475,-0.326) (-1.425,-0.291) (-1.375,-0.259) (-1.325,-0.230) (-1.300,-0.216)}
\EQ{(1.300,-0.191) (1.350,-0.215) (1.400,-0.241) (1.450,-0.270) (1.499,-0.300) (1.550,-0.335) (1.600,-0.372) (1.650,-0.414) (1.691,-0.450) (1.742,-0.500) (1.788,-0.550) (1.825,-0.594) (1.868,-0.650) (1.900,-0.696) (1.934,-0.750) (1.963,-0.800) (1.989,-0.850) (2.013,-0.900) (2.034,-0.950) (2.053,-1.000) (2.070,-1.050) (2.085,-1.100) (2.098,-1.150) (2.109,-1.200) (2.118,-1.250) (2.126,-1.300) (2.132,-1.350) (2.137,-1.400) (2.139,-1.450) (2.141,-1.500) (2.140,-1.550) (2.138,-1.600) (2.135,-1.650) (2.130,-1.700) (2.123,-1.750) (2.115,-1.800) (2.106,-1.850) (2.095,-1.900) (2.082,-1.950) (2.068,-2.000) (2.052,-2.050) (2.035,-2.100) (2.016,-2.150) (1.995,-2.200) (1.972,-2.250) (1.950,-2.296) (1.925,-2.344) (1.900,-2.389) (1.864,-2.450) (1.831,-2.500) (1.800,-2.546) (1.760,-2.600) (1.725,-2.645) (1.679,-2.700) (1.634,-2.750) (1.600,-2.786) (1.562,-2.825) (1.525,-2.860) (1.481,-2.900) (1.425,-2.947) (1.375,-2.987) (1.325,-3.025) (1.275,-3.060) (1.225,-3.093) (1.175,-3.124) (1.132,-3.150) (1.087,-3.175) (1.040,-3.200) (0.990,-3.225) (0.936,-3.250) (0.879,-3.275) (0.825,-3.297) (0.775,-3.315) (0.725,-3.333) (0.675,-3.349) (0.625,-3.364) (0.575,-3.377) (0.525,-3.390) (0.475,-3.401) (0.425,-3.411) (0.375,-3.420) (0.325,-3.427) (0.275,-3.434) (0.225,-3.439) (0.175,-3.444) (0.125,-3.447) (0.075,-3.449) (0.025,-3.450) (-0.037,-3.450) (-0.100,-3.448) (-0.150,-3.445) (-0.200,-3.442) (-0.250,-3.437) (-0.300,-3.431) (-0.350,-3.424) (-0.400,-3.415) (-0.450,-3.406) (-0.500,-3.395) (-0.550,-3.384) (-0.600,-3.371) (-0.650,-3.356) (-0.700,-3.341) (-0.748,-3.325) (-0.800,-3.306) (-0.850,-3.287) (-0.900,-3.266) (-0.950,-3.244) (-1.000,-3.220) (-1.050,-3.195) (-1.100,-3.168) (-1.150,-3.139) (-1.200,-3.109) (-1.250,-3.077) (-1.300,-3.043) (-1.350,-3.006) (-1.391,-2.975) (-1.450,-2.927) (-1.500,-2.883) (-1.550,-2.836) (-1.587,-2.800) (-1.625,-2.760) (-1.675,-2.705) (-1.721,-2.650) (-1.760,-2.600) (-1.797,-2.550) (-1.831,-2.500) (-1.864,-2.450) (-1.894,-2.400) (-1.922,-2.350) (-1.948,-2.300) (-1.972,-2.250) (-1.995,-2.200) (-2.016,-2.150) (-2.035,-2.100) (-2.052,-2.050) (-2.068,-2.000) (-2.082,-1.950) (-2.095,-1.900) (-2.106,-1.850) (-2.115,-1.800) (-2.123,-1.750) (-2.130,-1.700) (-2.135,-1.650) (-2.138,-1.600) (-2.140,-1.550) (-2.141,-1.500) (-2.139,-1.450) (-2.137,-1.400) (-2.132,-1.350) (-2.126,-1.300) (-2.118,-1.250) (-2.109,-1.200) (-2.098,-1.150) (-2.085,-1.100) (-2.070,-1.050) (-2.053,-1.000) (-2.034,-0.950) (-2.013,-0.900) (-1.989,-0.850) (-1.963,-0.800) (-1.934,-0.750) (-1.903,-0.700) (-1.868,-0.650) (-1.830,-0.600) (-1.788,-0.550) (-1.750,-0.508) (-1.700,-0.459) (-1.650,-0.414) (-1.603,-0.375) (-1.550,-0.335) (-1.500,-0.301) (-1.450,-0.270) (-1.400,-0.241) (-1.350,-0.215) (-1.300,-0.191)}
\EQ{(1.300,-0.165) (1.350,-0.186) (1.400,-0.208) (1.450,-0.233) (1.500,-0.259) (1.550,-0.288) (1.600,-0.319) (1.646,-0.350) (1.700,-0.390) (1.744,-0.425) (1.800,-0.474) (1.850,-0.522) (1.900,-0.575) (1.943,-0.625) (1.975,-0.665) (2.018,-0.725) (2.050,-0.773) (2.082,-0.825) (2.109,-0.875) (2.134,-0.925) (2.157,-0.975) (2.178,-1.025) (2.197,-1.075) (2.213,-1.125) (2.228,-1.175) (2.242,-1.225) (2.253,-1.275) (2.263,-1.325) (2.271,-1.375) (2.277,-1.425) (2.282,-1.475) (2.286,-1.525) (2.288,-1.575) (2.288,-1.625) (2.288,-1.675) (2.285,-1.725) (2.281,-1.775) (2.276,-1.825) (2.269,-1.875) (2.261,-1.925) (2.252,-1.975) (2.241,-2.025) (2.228,-2.075) (2.215,-2.125) (2.199,-2.175) (2.182,-2.225) (2.164,-2.275) (2.144,-2.325) (2.122,-2.375) (2.100,-2.423) (2.075,-2.473) (2.050,-2.520) (2.025,-2.565) (1.989,-2.625) (1.956,-2.675) (1.925,-2.720) (1.885,-2.775) (1.850,-2.820) (1.805,-2.875) (1.761,-2.925) (1.725,-2.964) (1.690,-3.000) (1.650,-3.039) (1.611,-3.075) (1.554,-3.125) (1.500,-3.170) (1.450,-3.209) (1.400,-3.245) (1.358,-3.275) (1.300,-3.313) (1.250,-3.344) (1.200,-3.373) (1.152,-3.400) (1.104,-3.425) (1.053,-3.450) (1.000,-3.475) (0.950,-3.497) (0.900,-3.517) (0.850,-3.536) (0.800,-3.554) (0.750,-3.571) (0.700,-3.587) (0.650,-3.601) (0.600,-3.614) (0.550,-3.626) (0.500,-3.637) (0.450,-3.647) (0.400,-3.656) (0.350,-3.664) (0.300,-3.670) (0.250,-3.676) (0.200,-3.680) (0.150,-3.684) (0.100,-3.687) (0.050,-3.688) (0.000,-3.689) (-0.050,-3.688) (-0.100,-3.687) (-0.150,-3.684) (-0.200,-3.680) (-0.250,-3.676) (-0.300,-3.670) (-0.350,-3.664) (-0.400,-3.656) (-0.450,-3.647) (-0.500,-3.637) (-0.550,-3.626) (-0.600,-3.614) (-0.650,-3.601) (-0.700,-3.587) (-0.750,-3.571) (-0.800,-3.554) (-0.850,-3.536) (-0.900,-3.517) (-0.950,-3.497) (-1.000,-3.475) (-1.050,-3.452) (-1.100,-3.427) (-1.150,-3.401) (-1.197,-3.375) (-1.250,-3.344) (-1.300,-3.313) (-1.350,-3.280) (-1.394,-3.250) (-1.450,-3.209) (-1.494,-3.175) (-1.550,-3.129) (-1.600,-3.085) (-1.638,-3.050) (-1.675,-3.015) (-1.714,-2.975) (-1.750,-2.937) (-1.800,-2.881) (-1.846,-2.825) (-1.885,-2.775) (-1.922,-2.725) (-1.956,-2.675) (-1.989,-2.625) (-2.019,-2.575) (-2.047,-2.525) (-2.074,-2.475) (-2.099,-2.425) (-2.122,-2.375) (-2.144,-2.325) (-2.164,-2.275) (-2.182,-2.225) (-2.199,-2.175) (-2.215,-2.125) (-2.228,-2.075) (-2.241,-2.025) (-2.252,-1.975) (-2.261,-1.925) (-2.269,-1.875) (-2.276,-1.825) (-2.281,-1.775) (-2.285,-1.725) (-2.288,-1.675) (-2.288,-1.625) (-2.288,-1.575) (-2.286,-1.525) (-2.282,-1.475) (-2.277,-1.425) (-2.271,-1.375) (-2.263,-1.325) (-2.253,-1.275) (-2.242,-1.225) (-2.228,-1.175) (-2.213,-1.125) (-2.197,-1.075) (-2.178,-1.025) (-2.157,-0.975) (-2.134,-0.925) (-2.109,-0.875) (-2.082,-0.825) (-2.051,-0.775) (-2.018,-0.725) (-1.982,-0.675) (-1.950,-0.634) (-1.900,-0.575) (-1.853,-0.525) (-1.801,-0.475) (-1.750,-0.430) (-1.700,-0.390) (-1.650,-0.353) (-1.600,-0.319) (-1.550,-0.288) (-1.500,-0.259) (-1.450,-0.233) (-1.400,-0.208) (-1.350,-0.186) (-1.300,-0.165)}
\EQ{(1.300,-0.141) (1.350,-0.158) (1.400,-0.177) (1.450,-0.197) (1.500,-0.220) (1.550,-0.243) (1.600,-0.269) (1.650,-0.297) (1.697,-0.325) (1.750,-0.359) (1.800,-0.394) (1.841,-0.425) (1.900,-0.473) (1.950,-0.517) (2.000,-0.566) (2.050,-0.619) (2.099,-0.675) (2.138,-0.725) (2.174,-0.775) (2.208,-0.825) (2.239,-0.875) (2.267,-0.925) (2.293,-0.975) (2.317,-1.025) (2.339,-1.075) (2.359,-1.125) (2.377,-1.175) (2.394,-1.225) (2.409,-1.275) (2.422,-1.325) (2.433,-1.375) (2.443,-1.425) (2.451,-1.475) (2.458,-1.525) (2.464,-1.575) (2.468,-1.625) (2.470,-1.675) (2.472,-1.725) (2.472,-1.775) (2.470,-1.825) (2.467,-1.875) (2.463,-1.925) (2.458,-1.975) (2.451,-2.025) (2.443,-2.075) (2.434,-2.125) (2.423,-2.175) (2.411,-2.225) (2.398,-2.275) (2.383,-2.325) (2.367,-2.375) (2.350,-2.423) (2.330,-2.475) (2.310,-2.525) (2.288,-2.575) (2.265,-2.625) (2.240,-2.675) (2.213,-2.725) (2.184,-2.775) (2.154,-2.825) (2.125,-2.871) (2.089,-2.925) (2.053,-2.975) (2.015,-3.025) (1.975,-3.075) (1.933,-3.125) (1.888,-3.175) (1.850,-3.215) (1.800,-3.265) (1.763,-3.300) (1.725,-3.335) (1.679,-3.375) (1.625,-3.420) (1.575,-3.459) (1.525,-3.496) (1.475,-3.531) (1.425,-3.565) (1.375,-3.597) (1.328,-3.625) (1.285,-3.650) (1.239,-3.675) (1.191,-3.700) (1.141,-3.725) (1.087,-3.750) (1.031,-3.775) (0.975,-3.798) (0.925,-3.817) (0.875,-3.835) (0.825,-3.852) (0.775,-3.868) (0.725,-3.883) (0.675,-3.897) (0.625,-3.909) (0.575,-3.921) (0.525,-3.932) (0.475,-3.941) (0.425,-3.950) (0.375,-3.958) (0.325,-3.964) (0.275,-3.970) (0.225,-3.975) (0.175,-3.978) (0.125,-3.981) (0.075,-3.983) (0.025,-3.984) (-0.050,-3.984) (-0.100,-3.982) (-0.150,-3.980) (-0.200,-3.977) (-0.250,-3.972) (-0.300,-3.967) (-0.350,-3.961) (-0.400,-3.954) (-0.450,-3.946) (-0.500,-3.937) (-0.550,-3.926) (-0.600,-3.915) (-0.650,-3.903) (-0.700,-3.890) (-0.750,-3.876) (-0.800,-3.860) (-0.850,-3.844) (-0.900,-3.826) (-0.950,-3.808) (-1.000,-3.788) (-1.050,-3.767) (-1.100,-3.744) (-1.150,-3.721) (-1.200,-3.696) (-1.250,-3.669) (-1.300,-3.641) (-1.350,-3.612) (-1.400,-3.581) (-1.448,-3.550) (-1.500,-3.514) (-1.550,-3.478) (-1.600,-3.440) (-1.650,-3.400) (-1.700,-3.357) (-1.750,-3.312) (-1.790,-3.275) (-1.840,-3.225) (-1.875,-3.189) (-1.910,-3.150) (-1.950,-3.105) (-1.995,-3.050) (-2.034,-3.000) (-2.071,-2.950) (-2.100,-2.908) (-2.139,-2.850) (-2.170,-2.800) (-2.199,-2.750) (-2.225,-2.703) (-2.250,-2.654) (-2.275,-2.603) (-2.299,-2.550) (-2.320,-2.500) (-2.340,-2.450) (-2.358,-2.400) (-2.375,-2.351) (-2.391,-2.300) (-2.405,-2.250) (-2.417,-2.200) (-2.429,-2.150) (-2.439,-2.100) (-2.447,-2.050) (-2.455,-2.000) (-2.461,-1.950) (-2.466,-1.900) (-2.469,-1.850) (-2.471,-1.800) (-2.472,-1.750) (-2.471,-1.700) (-2.469,-1.650) (-2.466,-1.600) (-2.461,-1.550) (-2.455,-1.500) (-2.447,-1.450) (-2.438,-1.400) (-2.427,-1.350) (-2.415,-1.300) (-2.401,-1.250) (-2.386,-1.200) (-2.369,-1.150) (-2.350,-1.101) (-2.329,-1.050) (-2.306,-1.000) (-2.281,-0.950) (-2.253,-0.900) (-2.225,-0.852) (-2.192,-0.800) (-2.157,-0.750) (-2.125,-0.708) (-2.078,-0.650) (-2.033,-0.600) (-1.984,-0.550) (-1.931,-0.500) (-1.875,-0.452) (-1.825,-0.413) (-1.775,-0.376) (-1.725,-0.343) (-1.675,-0.312) (-1.625,-0.283) (-1.575,-0.256) (-1.525,-0.231) (-1.475,-0.208) (-1.425,-0.187) (-1.375,-0.167) (-1.327,-0.150) (-1.300,-0.141)}
\EQ{(1.300,-0.117) (1.350,-0.131) (1.400,-0.146) (1.450,-0.163) (1.500,-0.181) (1.549,-0.200) (1.600,-0.221) (1.650,-0.244) (1.700,-0.268) (1.750,-0.293) (1.800,-0.321) (1.849,-0.350) (1.900,-0.382) (1.950,-0.416) (1.997,-0.450) (2.050,-0.492) (2.090,-0.525) (2.145,-0.575) (2.196,-0.625) (2.243,-0.675) (2.287,-0.725) (2.325,-0.772) (2.365,-0.825) (2.400,-0.875) (2.432,-0.925) (2.462,-0.975) (2.490,-1.025) (2.516,-1.075) (2.540,-1.125) (2.562,-1.175) (2.582,-1.225) (2.600,-1.275) (2.617,-1.325) (2.632,-1.375) (2.646,-1.425) (2.658,-1.475) (2.669,-1.525) (2.679,-1.575) (2.687,-1.625) (2.693,-1.675) (2.699,-1.725) (2.703,-1.775) (2.706,-1.825) (2.707,-1.875) (2.708,-1.925) (2.707,-1.975) (2.705,-2.025) (2.701,-2.075) (2.697,-2.125) (2.691,-2.175) (2.684,-2.225) (2.676,-2.275) (2.667,-2.325) (2.656,-2.375) (2.644,-2.425) (2.631,-2.475) (2.617,-2.525) (2.602,-2.575) (2.585,-2.625) (2.567,-2.675) (2.547,-2.725) (2.527,-2.775) (2.504,-2.825) (2.481,-2.875) (2.456,-2.925) (2.430,-2.975) (2.402,-3.025) (2.375,-3.070) (2.341,-3.125) (2.308,-3.175) (2.275,-3.223) (2.237,-3.275) (2.200,-3.323) (2.158,-3.375) (2.125,-3.414) (2.075,-3.470) (2.025,-3.523) (1.975,-3.574) (1.925,-3.621) (1.875,-3.667) (1.836,-3.700) (1.776,-3.750) (1.725,-3.790) (1.677,-3.825) (1.625,-3.862) (1.575,-3.896) (1.530,-3.925) (1.475,-3.959) (1.425,-3.988) (1.375,-4.016) (1.325,-4.043) (1.275,-4.068) (1.225,-4.092) (1.175,-4.115) (1.125,-4.137) (1.075,-4.158) (1.025,-4.177) (0.975,-4.196) (0.925,-4.213) (0.875,-4.229) (0.825,-4.245) (0.775,-4.259) (0.725,-4.272) (0.675,-4.285) (0.625,-4.296) (0.575,-4.307) (0.525,-4.317) (0.475,-4.325) (0.425,-4.333) (0.375,-4.340) (0.325,-4.346) (0.275,-4.351) (0.225,-4.356) (0.175,-4.359) (0.125,-4.362) (0.075,-4.363) (0.025,-4.364) (-0.050,-4.364) (-0.100,-4.363) (-0.150,-4.361) (-0.200,-4.358) (-0.250,-4.354) (-0.300,-4.349) (-0.350,-4.343) (-0.400,-4.337) (-0.450,-4.329) (-0.500,-4.321) (-0.550,-4.312) (-0.600,-4.302) (-0.650,-4.291) (-0.700,-4.279) (-0.750,-4.266) (-0.800,-4.252) (-0.850,-4.237) (-0.900,-4.221) (-0.950,-4.204) (-1.000,-4.187) (-1.050,-4.168) (-1.100,-4.147) (-1.150,-4.126) (-1.200,-4.104) (-1.250,-4.080) (-1.300,-4.056) (-1.350,-4.030) (-1.400,-4.002) (-1.448,-3.975) (-1.500,-3.944) (-1.550,-3.912) (-1.600,-3.879) (-1.642,-3.850) (-1.700,-3.808) (-1.744,-3.775) (-1.800,-3.730) (-1.850,-3.688) (-1.894,-3.650) (-1.948,-3.600) (-1.999,-3.550) (-2.047,-3.500) (-2.093,-3.450) (-2.137,-3.400) (-2.175,-3.354) (-2.218,-3.300) (-2.250,-3.257) (-2.291,-3.200) (-2.325,-3.150) (-2.357,-3.100) (-2.387,-3.050) (-2.416,-3.000) (-2.443,-2.950) (-2.469,-2.900) (-2.493,-2.850) (-2.516,-2.800) (-2.537,-2.750) (-2.557,-2.700) (-2.575,-2.653) (-2.593,-2.600) (-2.610,-2.550) (-2.624,-2.500) (-2.638,-2.450) (-2.650,-2.400) (-2.662,-2.350) (-2.672,-2.300) (-2.680,-2.250) (-2.688,-2.200) (-2.694,-2.150) (-2.699,-2.100) (-2.703,-2.050) (-2.706,-2.000) (-2.707,-1.950) (-2.708,-1.900) (-2.707,-1.850) (-2.704,-1.800) (-2.701,-1.750) (-2.696,-1.700) (-2.690,-1.650) (-2.683,-1.600) (-2.674,-1.550) (-2.664,-1.500) (-2.652,-1.450) (-2.639,-1.400) (-2.625,-1.351) (-2.609,-1.300) (-2.591,-1.250) (-2.572,-1.200) (-2.551,-1.150) (-2.528,-1.100) (-2.503,-1.050) (-2.476,-1.000) (-2.450,-0.954) (-2.416,-0.900) (-2.383,-0.850) (-2.350,-0.805) (-2.308,-0.750) (-2.275,-0.711) (-2.225,-0.655) (-2.175,-0.604) (-2.125,-0.557) (-2.075,-0.513) (-2.029,-0.475) (-1.975,-0.434) (-1.927,-0.400) (-1.875,-0.366) (-1.825,-0.335) (-1.775,-0.307) (-1.725,-0.280) (-1.675,-0.255) (-1.625,-0.232) (-1.575,-0.211) (-1.525,-0.191) (-1.475,-0.172) (-1.425,-0.155) (-1.375,-0.139) (-1.325,-0.124) (-1.300,-0.117)}
\EQ{(1.300,-0.093) (1.350,-0.104) (1.400,-0.116) (1.450,-0.130) (1.500,-0.144) (1.550,-0.159) (1.600,-0.175) (1.650,-0.193) (1.700,-0.211) (1.750,-0.231) (1.800,-0.252) (1.850,-0.275) (1.900,-0.299) (1.950,-0.324) (1.997,-0.350) (2.041,-0.375) (2.100,-0.412) (2.150,-0.444) (2.194,-0.475) (2.250,-0.517) (2.292,-0.550) (2.350,-0.599) (2.400,-0.645) (2.450,-0.693) (2.500,-0.746) (2.548,-0.800) (2.589,-0.850) (2.625,-0.896) (2.664,-0.950) (2.697,-1.000) (2.725,-1.044) (2.759,-1.100) (2.786,-1.150) (2.812,-1.200) (2.836,-1.250) (2.858,-1.300) (2.879,-1.350) (2.898,-1.400) (2.916,-1.450) (2.932,-1.500) (2.947,-1.550) (2.960,-1.600) (2.972,-1.650) (2.983,-1.700) (2.993,-1.750) (3.001,-1.800) (3.008,-1.850) (3.014,-1.900) (3.019,-1.950) (3.023,-2.000) (3.025,-2.050) (3.027,-2.100) (3.027,-2.150) (3.026,-2.200) (3.025,-2.250) (3.022,-2.300) (3.018,-2.350) (3.013,-2.400) (3.006,-2.450) (2.999,-2.500) (2.991,-2.550) (2.981,-2.600) (2.971,-2.650) (2.959,-2.700) (2.947,-2.750) (2.933,-2.800) (2.918,-2.850) (2.902,-2.900) (2.885,-2.950) (2.866,-3.000) (2.847,-3.050) (2.826,-3.100) (2.804,-3.150) (2.781,-3.200) (2.756,-3.250) (2.730,-3.300) (2.703,-3.350) (2.675,-3.399) (2.645,-3.450) (2.613,-3.500) (2.580,-3.550) (2.550,-3.594) (2.509,-3.650) (2.475,-3.696) (2.432,-3.750) (2.400,-3.789) (2.350,-3.846) (2.301,-3.900) (2.253,-3.950) (2.203,-4.000) (2.151,-4.050) (2.100,-4.096) (2.050,-4.139) (2.007,-4.175) (1.950,-4.220) (1.900,-4.258) (1.850,-4.295) (1.807,-4.325) (1.750,-4.363) (1.700,-4.395) (1.651,-4.425) (1.600,-4.455) (1.550,-4.483) (1.500,-4.510) (1.450,-4.536) (1.400,-4.561) (1.350,-4.585) (1.300,-4.607) (1.250,-4.629) (1.200,-4.649) (1.150,-4.669) (1.100,-4.687) (1.050,-4.705) (1.000,-4.722) (0.950,-4.737) (0.900,-4.752) (0.850,-4.766) (0.800,-4.779) (0.750,-4.792) (0.700,-4.803) (0.650,-4.814) (0.600,-4.824) (0.550,-4.833) (0.500,-4.841) (0.450,-4.848) (0.400,-4.855) (0.350,-4.861) (0.300,-4.866) (0.250,-4.870) (0.200,-4.873) (0.150,-4.876) (0.100,-4.878) (0.050,-4.879) (0.000,-4.880) (-0.050,-4.879) (-0.100,-4.878) (-0.150,-4.876) (-0.200,-4.873) (-0.250,-4.870) (-0.300,-4.866) (-0.350,-4.861) (-0.400,-4.855) (-0.450,-4.848) (-0.500,-4.841) (-0.550,-4.833) (-0.600,-4.824) (-0.650,-4.814) (-0.700,-4.803) (-0.750,-4.792) (-0.800,-4.779) (-0.850,-4.766) (-0.900,-4.752) (-0.950,-4.737) (-1.000,-4.722) (-1.050,-4.705) (-1.100,-4.687) (-1.150,-4.669) (-1.198,-4.650) (-1.250,-4.629) (-1.300,-4.607) (-1.350,-4.585) (-1.400,-4.561) (-1.450,-4.536) (-1.500,-4.510) (-1.550,-4.483) (-1.600,-4.455) (-1.650,-4.426) (-1.700,-4.395) (-1.750,-4.363) (-1.800,-4.330) (-1.843,-4.300) (-1.900,-4.258) (-1.944,-4.225) (-2.000,-4.181) (-2.050,-4.139) (-2.095,-4.100) (-2.150,-4.051) (-2.200,-4.003) (-2.250,-3.954) (-2.300,-3.901) (-2.347,-3.850) (-2.390,-3.800) (-2.425,-3.758) (-2.472,-3.700) (-2.509,-3.650) (-2.546,-3.600) (-2.575,-3.558) (-2.613,-3.500) (-2.645,-3.450) (-2.675,-3.400) (-2.700,-3.356) (-2.725,-3.310) (-2.750,-3.262) (-2.775,-3.212) (-2.800,-3.159) (-2.825,-3.102) (-2.847,-3.050) (-2.866,-3.000) (-2.885,-2.950) (-2.902,-2.900) (-2.918,-2.850) (-2.933,-2.800) (-2.947,-2.750) (-2.959,-2.700) (-2.971,-2.650) (-2.981,-2.600) (-2.991,-2.550) (-2.999,-2.500) (-3.006,-2.450) (-3.013,-2.400) (-3.018,-2.350) (-3.022,-2.300) (-3.025,-2.250) (-3.026,-2.200) (-3.027,-2.150) (-3.027,-2.100) (-3.025,-2.050) (-3.023,-2.000) (-3.019,-1.950) (-3.014,-1.900) (-3.008,-1.850) (-3.001,-1.800) (-2.993,-1.750) (-2.983,-1.700) (-2.972,-1.650) (-2.960,-1.600) (-2.947,-1.550) (-2.932,-1.500) (-2.916,-1.450) (-2.898,-1.400) (-2.879,-1.350) (-2.858,-1.300) (-2.836,-1.250) (-2.812,-1.200) (-2.786,-1.150) (-2.759,-1.100) (-2.729,-1.050) (-2.700,-1.004) (-2.664,-0.950) (-2.628,-0.900) (-2.589,-0.850) (-2.550,-0.802) (-2.504,-0.750) (-2.457,-0.700) (-2.406,-0.650) (-2.351,-0.600) (-2.300,-0.557) (-2.261,-0.525) (-2.200,-0.479) (-2.158,-0.450) (-2.100,-0.412) (-2.050,-0.381) (-2.000,-0.352) (-1.951,-0.325) (-1.902,-0.300) (-1.850,-0.275) (-1.800,-0.252) (-1.750,-0.231) (-1.700,-0.211) (-1.650,-0.193) (-1.600,-0.175) (-1.550,-0.159) (-1.500,-0.144) (-1.450,-0.130) (-1.400,-0.116) (-1.350,-0.104) (-1.300,-0.093)}
\EQ{(1.300,-0.070) (1.350,-0.078) (1.400,-0.087) (1.450,-0.097) (1.500,-0.107) (1.550,-0.118) (1.600,-0.130) (1.650,-0.143) (1.700,-0.157) (1.750,-0.171) (1.800,-0.187) (1.850,-0.203) (1.900,-0.220) (1.950,-0.239) (2.000,-0.258) (2.050,-0.279) (2.098,-0.300) (2.150,-0.324) (2.200,-0.348) (2.250,-0.374) (2.298,-0.400) (2.342,-0.425) (2.400,-0.460) (2.450,-0.492) (2.500,-0.525) (2.550,-0.561) (2.600,-0.598) (2.650,-0.638) (2.694,-0.675) (2.750,-0.724) (2.800,-0.772) (2.850,-0.822) (2.899,-0.875) (2.943,-0.925) (2.985,-0.975) (3.024,-1.025) (3.060,-1.075) (3.095,-1.125) (3.125,-1.171) (3.159,-1.225) (3.188,-1.275) (3.216,-1.325) (3.241,-1.375) (3.266,-1.425) (3.289,-1.475) (3.310,-1.525) (3.330,-1.575) (3.349,-1.625) (3.367,-1.675) (3.383,-1.725) (3.398,-1.775) (3.412,-1.825) (3.425,-1.875) (3.437,-1.925) (3.447,-1.975) (3.456,-2.025) (3.465,-2.075) (3.472,-2.125) (3.479,-2.175) (3.484,-2.225) (3.488,-2.275) (3.491,-2.325) (3.494,-2.375) (3.495,-2.425) (3.496,-2.475) (3.495,-2.525) (3.494,-2.575) (3.491,-2.625) (3.488,-2.675) (3.483,-2.725) (3.478,-2.775) (3.472,-2.825) (3.465,-2.875) (3.457,-2.925) (3.448,-2.975) (3.438,-3.025) (3.427,-3.075) (3.415,-3.125) (3.403,-3.175) (3.389,-3.225) (3.374,-3.275) (3.359,-3.325) (3.342,-3.375) (3.325,-3.423) (3.306,-3.475) (3.286,-3.525) (3.265,-3.575) (3.243,-3.625) (3.220,-3.675) (3.196,-3.725) (3.171,-3.775) (3.145,-3.825) (3.118,-3.875) (3.089,-3.925) (3.059,-3.975) (3.028,-4.025) (3.000,-4.068) (2.962,-4.125) (2.927,-4.175) (2.890,-4.225) (2.852,-4.275) (2.812,-4.325) (2.775,-4.370) (2.728,-4.425) (2.683,-4.475) (2.637,-4.525) (2.600,-4.563) (2.563,-4.600) (2.525,-4.637) (2.485,-4.675) (2.430,-4.725) (2.375,-4.773) (2.325,-4.815) (2.281,-4.850) (2.225,-4.894) (2.183,-4.925) (2.125,-4.967) (2.078,-5.000) (2.025,-5.035) (1.975,-5.067) (1.925,-5.098) (1.881,-5.125) (1.825,-5.157) (1.775,-5.185) (1.725,-5.212) (1.675,-5.238) (1.625,-5.262) (1.575,-5.286) (1.525,-5.309) (1.475,-5.331) (1.425,-5.352) (1.375,-5.373) (1.325,-5.392) (1.275,-5.411) (1.225,-5.428) (1.175,-5.445) (1.125,-5.462) (1.075,-5.477) (1.025,-5.492) (0.975,-5.505) (0.925,-5.519) (0.875,-5.531) (0.825,-5.543) (0.775,-5.553) (0.725,-5.564) (0.675,-5.573) (0.625,-5.582) (0.575,-5.590) (0.525,-5.597) (0.475,-5.604) (0.425,-5.610) (0.375,-5.616) (0.325,-5.620) (0.275,-5.624) (0.225,-5.628) (0.175,-5.630) (0.125,-5.632) (0.075,-5.634) (0.025,-5.634) (-0.050,-5.634) (-0.100,-5.633) (-0.150,-5.631) (-0.200,-5.629) (-0.250,-5.626) (-0.300,-5.622) (-0.350,-5.618) (-0.400,-5.613) (-0.450,-5.607) (-0.500,-5.601) (-0.550,-5.594) (-0.600,-5.586) (-0.650,-5.578) (-0.700,-5.569) (-0.750,-5.559) (-0.800,-5.548) (-0.850,-5.537) (-0.899,-5.525) (-0.950,-5.512) (-1.000,-5.499) (-1.050,-5.484) (-1.100,-5.469) (-1.150,-5.453) (-1.200,-5.437) (-1.250,-5.420) (-1.300,-5.401) (-1.350,-5.382) (-1.400,-5.363) (-1.450,-5.342) (-1.500,-5.320) (-1.550,-5.298) (-1.599,-5.275) (-1.650,-5.250) (-1.700,-5.225) (-1.747,-5.200) (-1.793,-5.175) (-1.838,-5.150) (-1.900,-5.114) (-1.950,-5.083) (-2.000,-5.051) (-2.050,-5.019) (-2.100,-4.985) (-2.149,-4.950) (-2.200,-4.913) (-2.249,-4.875) (-2.300,-4.835) (-2.343,-4.800) (-2.400,-4.751) (-2.450,-4.707) (-2.500,-4.661) (-2.538,-4.625) (-2.575,-4.588) (-2.613,-4.550) (-2.650,-4.511) (-2.700,-4.457) (-2.750,-4.400) (-2.792,-4.350) (-2.825,-4.309) (-2.871,-4.250) (-2.908,-4.200) (-2.944,-4.150) (-2.975,-4.106) (-3.012,-4.050) (-3.044,-4.000) (-3.074,-3.950) (-3.100,-3.906) (-3.125,-3.862) (-3.150,-3.816) (-3.175,-3.768) (-3.200,-3.718) (-3.225,-3.665) (-3.250,-3.610) (-3.275,-3.551) (-3.296,-3.500) (-3.315,-3.450) (-3.333,-3.400) (-3.350,-3.351) (-3.366,-3.300) (-3.382,-3.250) (-3.396,-3.200) (-3.409,-3.150) (-3.421,-3.100) (-3.433,-3.050) (-3.443,-3.000) (-3.453,-2.950) (-3.461,-2.900) (-3.469,-2.850) (-3.475,-2.800) (-3.481,-2.750) (-3.486,-2.700) (-3.490,-2.650) (-3.492,-2.600) (-3.494,-2.550) (-3.495,-2.500) (-3.496,-2.450) (-3.495,-2.400) (-3.493,-2.350) (-3.490,-2.300) (-3.486,-2.250) (-3.481,-2.200) (-3.475,-2.150) (-3.469,-2.100) (-3.461,-2.050) (-3.452,-2.000) (-3.442,-1.950) (-3.431,-1.900) (-3.419,-1.850) (-3.405,-1.800) (-3.391,-1.750) (-3.375,-1.700) (-3.358,-1.650) (-3.340,-1.600) (-3.321,-1.550) (-3.300,-1.501) (-3.277,-1.450) (-3.254,-1.400) (-3.229,-1.350) (-3.202,-1.300) (-3.175,-1.252) (-3.144,-1.200) (-3.112,-1.150) (-3.078,-1.100) (-3.042,-1.050) (-3.004,-1.000) (-2.964,-0.950) (-2.925,-0.904) (-2.876,-0.850) (-2.828,-0.800) (-2.777,-0.750) (-2.725,-0.702) (-2.675,-0.659) (-2.634,-0.625) (-2.575,-0.579) (-2.525,-0.543) (-2.475,-0.508) (-2.425,-0.475) (-2.375,-0.444) (-2.325,-0.415) (-2.275,-0.387) (-2.225,-0.361) (-2.175,-0.336) (-2.125,-0.312) (-2.075,-0.290) (-2.025,-0.269) (-1.975,-0.248) (-1.925,-0.230) (-1.875,-0.212) (-1.825,-0.195) (-1.775,-0.179) (-1.725,-0.164) (-1.676,-0.150) (-1.625,-0.137) (-1.575,-0.124) (-1.525,-0.113) (-1.475,-0.102) (-1.425,-0.092) (-1.375,-0.082) (-1.325,-0.074) (-1.300,-0.070)}
\EQ{(1.300,-0.046) (1.350,-0.052) (1.400,-0.058) (1.450,-0.064) (1.500,-0.071) (1.550,-0.079) (1.600,-0.086) (1.650,-0.095) (1.700,-0.104) (1.750,-0.113) (1.800,-0.123) (1.850,-0.134) (1.900,-0.145) (1.950,-0.157) (2.000,-0.170) (2.050,-0.183) (2.100,-0.197) (2.150,-0.212) (2.200,-0.227) (2.250,-0.243) (2.300,-0.260) (2.350,-0.278) (2.400,-0.297) (2.450,-0.317) (2.500,-0.337) (2.550,-0.359) (2.600,-0.381) (2.650,-0.405) (2.700,-0.429) (2.750,-0.455) (2.800,-0.482) (2.850,-0.510) (2.900,-0.539) (2.950,-0.570) (2.998,-0.600) (3.050,-0.635) (3.100,-0.670) (3.141,-0.700) (3.200,-0.745) (3.250,-0.785) (3.298,-0.825) (3.350,-0.871) (3.400,-0.917) (3.450,-0.966) (3.500,-1.017) (3.550,-1.071) (3.598,-1.125) (3.640,-1.175) (3.675,-1.219) (3.718,-1.275) (3.750,-1.319) (3.789,-1.375) (3.822,-1.425) (3.850,-1.470) (3.883,-1.525) (3.912,-1.575) (3.939,-1.625) (3.965,-1.675) (3.990,-1.725) (4.014,-1.775) (4.036,-1.825) (4.057,-1.875) (4.077,-1.925) (4.096,-1.975) (4.114,-2.025) (4.131,-2.075) (4.147,-2.125) (4.162,-2.175) (4.176,-2.225) (4.189,-2.275) (4.201,-2.325) (4.213,-2.375) (4.223,-2.425) (4.232,-2.475) (4.241,-2.525) (4.249,-2.575) (4.256,-2.625) (4.262,-2.675) (4.267,-2.725) (4.271,-2.775) (4.275,-2.825) (4.278,-2.875) (4.280,-2.925) (4.281,-2.975) (4.281,-3.025) (4.281,-3.075) (4.280,-3.125) (4.278,-3.175) (4.275,-3.225) (4.272,-3.275) (4.268,-3.325) (4.263,-3.375) (4.257,-3.425) (4.250,-3.475) (4.243,-3.525) (4.235,-3.575) (4.226,-3.625) (4.217,-3.675) (4.206,-3.725) (4.195,-3.775) (4.184,-3.825) (4.171,-3.875) (4.157,-3.925) (4.143,-3.975) (4.128,-4.025) (4.112,-4.075) (4.096,-4.125) (4.078,-4.175) (4.060,-4.225) (4.041,-4.275) (4.021,-4.325) (4.000,-4.375) (3.979,-4.425) (3.956,-4.475) (3.933,-4.525) (3.908,-4.575) (3.883,-4.625) (3.857,-4.675) (3.830,-4.725) (3.802,-4.775) (3.775,-4.821) (3.743,-4.875) (3.711,-4.925) (3.679,-4.975) (3.650,-5.019) (3.611,-5.075) (3.576,-5.125) (3.539,-5.175) (3.501,-5.225) (3.462,-5.275) (3.425,-5.321) (3.380,-5.375) (3.336,-5.425) (3.300,-5.466) (3.250,-5.520) (3.200,-5.573) (3.150,-5.623) (3.100,-5.672) (3.050,-5.720) (3.000,-5.766) (2.961,-5.800) (2.904,-5.850) (2.850,-5.895) (2.800,-5.935) (2.750,-5.974) (2.700,-6.013) (2.650,-6.049) (2.600,-6.085) (2.550,-6.120) (2.505,-6.150) (2.450,-6.186) (2.400,-6.218) (2.350,-6.249) (2.306,-6.275) (2.250,-6.308) (2.200,-6.336) (2.150,-6.363) (2.100,-6.390) (2.050,-6.415) (2.000,-6.440) (1.950,-6.464) (1.900,-6.488) (1.850,-6.510) (1.800,-6.532) (1.750,-6.553) (1.700,-6.574) (1.650,-6.593) (1.600,-6.612) (1.550,-6.631) (1.500,-6.648) (1.450,-6.665) (1.400,-6.682) (1.350,-6.698) (1.300,-6.713) (1.250,-6.727) (1.200,-6.741) (1.150,-6.754) (1.100,-6.767) (1.050,-6.779) (1.000,-6.791) (0.950,-6.801) (0.900,-6.812) (0.850,-6.821) (0.800,-6.831) (0.750,-6.839) (0.700,-6.847) (0.650,-6.854) (0.600,-6.861) (0.550,-6.868) (0.500,-6.873) (0.450,-6.879) (0.400,-6.883) (0.350,-6.887) (0.300,-6.891) (0.250,-6.894) (0.200,-6.896) (0.150,-6.898) (0.100,-6.900) (0.050,-6.900) (0.000,-6.901) (-0.050,-6.900) (-0.100,-6.900) (-0.150,-6.898) (-0.200,-6.896) (-0.250,-6.894) (-0.300,-6.891) (-0.350,-6.887) (-0.400,-6.883) (-0.450,-6.879) (-0.500,-6.873) (-0.550,-6.868) (-0.600,-6.861) (-0.650,-6.854) (-0.700,-6.847) (-0.750,-6.839) (-0.800,-6.831) (-0.850,-6.821) (-0.900,-6.812) (-0.950,-6.801) (-1.000,-6.791) (-1.050,-6.779) (-1.100,-6.767) (-1.150,-6.754) (-1.200,-6.741) (-1.250,-6.727) (-1.300,-6.713) (-1.350,-6.698) (-1.400,-6.682) (-1.450,-6.665) (-1.500,-6.648) (-1.550,-6.631) (-1.600,-6.612) (-1.650,-6.593) (-1.697,-6.575) (-1.750,-6.553) (-1.800,-6.532) (-1.850,-6.510) (-1.900,-6.488) (-1.950,-6.464) (-2.000,-6.440) (-2.050,-6.415) (-2.100,-6.390) (-2.150,-6.363) (-2.200,-6.336) (-2.250,-6.308) (-2.300,-6.279) (-2.348,-6.250) (-2.400,-6.218) (-2.450,-6.186) (-2.500,-6.154) (-2.543,-6.125) (-2.600,-6.085) (-2.649,-6.050) (-2.700,-6.013) (-2.749,-5.975) (-2.800,-5.935) (-2.844,-5.900) (-2.900,-5.853) (-2.950,-5.810) (-2.990,-5.775) (-3.044,-5.725) (-3.097,-5.675) (-3.148,-5.625) (-3.198,-5.575) (-3.246,-5.525) (-3.292,-5.475) (-3.336,-5.425) (-3.375,-5.380) (-3.421,-5.325) (-3.462,-5.275) (-3.500,-5.227) (-3.539,-5.175) (-3.575,-5.126) (-3.611,-5.075) (-3.646,-5.025) (-3.675,-4.981) (-3.711,-4.925) (-3.743,-4.875) (-3.773,-4.825) (-3.800,-4.778) (-3.825,-4.734) (-3.850,-4.688) (-3.875,-4.641) (-3.900,-4.592) (-3.925,-4.541) (-3.950,-4.488) (-3.975,-4.433) (-4.000,-4.376) (-4.021,-4.325) (-4.041,-4.275) (-4.060,-4.225) (-4.078,-4.175) (-4.096,-4.125) (-4.112,-4.075) (-4.128,-4.025) (-4.143,-3.975) (-4.157,-3.925) (-4.171,-3.875) (-4.184,-3.825) (-4.195,-3.775) (-4.206,-3.725) (-4.217,-3.675) (-4.226,-3.625) (-4.235,-3.575) (-4.243,-3.525) (-4.250,-3.475) (-4.257,-3.425) (-4.263,-3.375) (-4.268,-3.325) (-4.272,-3.275) (-4.275,-3.225) (-4.278,-3.175) (-4.280,-3.125) (-4.281,-3.075) (-4.281,-3.025) (-4.281,-2.975) (-4.280,-2.925) (-4.278,-2.875) (-4.275,-2.825) (-4.271,-2.775) (-4.267,-2.725) (-4.262,-2.675) (-4.256,-2.625) (-4.249,-2.575) (-4.241,-2.525) (-4.232,-2.475) (-4.223,-2.425) (-4.213,-2.375) (-4.201,-2.325) (-4.189,-2.275) (-4.176,-2.225) (-4.162,-2.175) (-4.147,-2.125) (-4.131,-2.075) (-4.114,-2.025) (-4.096,-1.975) (-4.077,-1.925) (-4.057,-1.875) (-4.036,-1.825) (-4.014,-1.775) (-3.990,-1.725) (-3.965,-1.675) (-3.939,-1.625) (-3.912,-1.575) (-3.883,-1.525) (-3.853,-1.475) (-3.825,-1.430) (-3.789,-1.375) (-3.754,-1.325) (-3.718,-1.275) (-3.680,-1.225) (-3.640,-1.175) (-3.600,-1.127) (-3.554,-1.075) (-3.508,-1.025) (-3.459,-0.975) (-3.408,-0.925) (-3.355,-0.875) (-3.300,-0.827) (-3.250,-0.785) (-3.207,-0.750) (-3.150,-0.707) (-3.107,-0.675) (-3.050,-0.635) (-3.000,-0.602) (-2.950,-0.570) (-2.900,-0.539) (-2.850,-0.510) (-2.800,-0.482) (-2.750,-0.455) (-2.700,-0.429) (-2.650,-0.405) (-2.600,-0.381) (-2.550,-0.359) (-2.500,-0.337) (-2.450,-0.317) (-2.400,-0.297) (-2.350,-0.278) (-2.300,-0.260) (-2.250,-0.243) (-2.200,-0.227) (-2.150,-0.212) (-2.100,-0.197) (-2.050,-0.183) (-2.000,-0.170) (-1.950,-0.157) (-1.900,-0.145) (-1.850,-0.134) (-1.800,-0.123) (-1.750,-0.113) (-1.700,-0.104) (-1.650,-0.095) (-1.600,-0.086) (-1.550,-0.079) (-1.500,-0.071) (-1.450,-0.064) (-1.400,-0.058) (-1.350,-0.052) (-1.300,-0.046)}
\EQ{(1.300,-0.023) (1.350,-0.026) (1.400,-0.029) (1.450,-0.032) (1.500,-0.035) (1.550,-0.039) (1.600,-0.043) (1.650,-0.047) (1.700,-0.052) (1.750,-0.056) (1.800,-0.061) (1.850,-0.067) (1.900,-0.072) (1.950,-0.078) (2.000,-0.084) (2.050,-0.091) (2.100,-0.098) (2.150,-0.105) (2.200,-0.112) (2.250,-0.120) (2.300,-0.128) (2.350,-0.137) (2.400,-0.146) (2.450,-0.155) (2.500,-0.165) (2.550,-0.175) (2.600,-0.186) (2.650,-0.197) (2.700,-0.209) (2.750,-0.220) (2.800,-0.233) (2.850,-0.246) (2.900,-0.259) (2.950,-0.273) (3.000,-0.287) (3.050,-0.302) (3.100,-0.318) (3.150,-0.334) (3.199,-0.350) (3.250,-0.367) (3.300,-0.385) (3.350,-0.403) (3.400,-0.422) (3.450,-0.442) (3.500,-0.462) (3.550,-0.483) (3.600,-0.504) (3.646,-0.525) (3.700,-0.550) (3.750,-0.573) (3.800,-0.598) (3.850,-0.623) (3.900,-0.649) (3.949,-0.675) (3.994,-0.700) (4.038,-0.725) (4.100,-0.761) (4.150,-0.792) (4.200,-0.823) (4.250,-0.856) (4.300,-0.889) (4.350,-0.923) (4.400,-0.959) (4.450,-0.996) (4.500,-1.034) (4.550,-1.073) (4.600,-1.113) (4.644,-1.150) (4.700,-1.198) (4.750,-1.243) (4.800,-1.289) (4.838,-1.325) (4.875,-1.362) (4.913,-1.400) (4.950,-1.438) (4.985,-1.475) (5.025,-1.519) (5.074,-1.575) (5.117,-1.625) (5.150,-1.665) (5.198,-1.725) (5.237,-1.775) (5.274,-1.825) (5.310,-1.875) (5.345,-1.925) (5.375,-1.970) (5.411,-2.025) (5.443,-2.075) (5.473,-2.125) (5.500,-2.171) (5.525,-2.214) (5.550,-2.259) (5.575,-2.306) (5.598,-2.350) (5.623,-2.400) (5.648,-2.450) (5.671,-2.500) (5.694,-2.550) (5.716,-2.600) (5.737,-2.650) (5.757,-2.700) (5.777,-2.750) (5.796,-2.800) (5.814,-2.850) (5.831,-2.900) (5.848,-2.950) (5.863,-3.000) (5.879,-3.050) (5.893,-3.100) (5.907,-3.150) (5.920,-3.200) (5.933,-3.250) (5.944,-3.300) (5.956,-3.350) (5.966,-3.400) (5.976,-3.450) (5.985,-3.500) (5.994,-3.550) (6.002,-3.600) (6.010,-3.650) (6.017,-3.700) (6.023,-3.750) (6.029,-3.800) (6.034,-3.850) (6.038,-3.900) (6.042,-3.950) (6.046,-4.000) (6.049,-4.050) (6.051,-4.100) (6.053,-4.150) (6.054,-4.200) (6.054,-4.250) (6.054,-4.300) (6.054,-4.350) (6.053,-4.400) (6.051,-4.450) (6.049,-4.500) (6.047,-4.550) (6.043,-4.600) (6.040,-4.650) (6.035,-4.700) (6.031,-4.750) (6.025,-4.800) (6.019,-4.850) (6.013,-4.900) (6.006,-4.950) (5.998,-5.000) (5.990,-5.050) (5.982,-5.100) (5.973,-5.150) (5.963,-5.200) (5.953,-5.250) (5.942,-5.300) (5.931,-5.350) (5.919,-5.400) (5.906,-5.450) (5.893,-5.500) (5.880,-5.550) (5.866,-5.600) (5.851,-5.650) (5.836,-5.700) (5.820,-5.750) (5.804,-5.800) (5.787,-5.850) (5.769,-5.900) (5.751,-5.950) (5.732,-6.000) (5.713,-6.050) (5.693,-6.100) (5.673,-6.150) (5.652,-6.200) (5.630,-6.250) (5.608,-6.300) (5.585,-6.350) (5.561,-6.400) (5.537,-6.450) (5.512,-6.500) (5.487,-6.550) (5.461,-6.600) (5.434,-6.650) (5.406,-6.700) (5.378,-6.750) (5.350,-6.799) (5.320,-6.850) (5.289,-6.900) (5.258,-6.950) (5.226,-7.000) (5.194,-7.050) (5.160,-7.100) (5.126,-7.150) (5.091,-7.200) (5.056,-7.250) (5.025,-7.292) (4.981,-7.350) (4.950,-7.391) (4.904,-7.450) (4.864,-7.500) (4.825,-7.547) (4.781,-7.600) (4.738,-7.650) (4.700,-7.693) (4.650,-7.749) (4.603,-7.800) (4.555,-7.850) (4.507,-7.900) (4.457,-7.950) (4.407,-8.000) (4.355,-8.050) (4.302,-8.100) (4.250,-8.147) (4.200,-8.192) (4.162,-8.225) (4.104,-8.275) (4.050,-8.320) (4.000,-8.361) (3.952,-8.400) (3.900,-8.441) (3.855,-8.475) (3.800,-8.517) (3.755,-8.550) (3.700,-8.589) (3.650,-8.625) (3.600,-8.659) (3.550,-8.693) (3.501,-8.725) (3.450,-8.758) (3.400,-8.790) (3.350,-8.821) (3.302,-8.850) (3.250,-8.881) (3.200,-8.910) (3.150,-8.939) (3.100,-8.967) (3.050,-8.994) (3.000,-9.021) (2.950,-9.047) (2.900,-9.072) (2.850,-9.097) (2.800,-9.122) (2.750,-9.146) (2.700,-9.169) (2.650,-9.192) (2.600,-9.214) (2.550,-9.236) (2.500,-9.257) (2.450,-9.278) (2.400,-9.298) (2.350,-9.318) (2.300,-9.337) (2.250,-9.356) (2.200,-9.374) (2.150,-9.392) (2.100,-9.410) (2.050,-9.427) (2.000,-9.443) (1.950,-9.459) (1.900,-9.475) (1.850,-9.490) (1.800,-9.504) (1.750,-9.519) (1.700,-9.533) (1.650,-9.546) (1.600,-9.559) (1.550,-9.571) (1.500,-9.583) (1.450,-9.595) (1.400,-9.606) (1.350,-9.617) (1.300,-9.628) (1.250,-9.638) (1.200,-9.647) (1.150,-9.656) (1.100,-9.665) (1.050,-9.674) (1.000,-9.682) (0.950,-9.689) (0.900,-9.696) (0.850,-9.703) (0.800,-9.710) (0.750,-9.716) (0.700,-9.721) (0.650,-9.726) (0.600,-9.731) (0.550,-9.736) (0.500,-9.740) (0.450,-9.743) (0.400,-9.747) (0.350,-9.750) (0.300,-9.752) (0.250,-9.754) (0.200,-9.756) (0.150,-9.757) (0.100,-9.758) (0.050,-9.759) (0.000,-9.759) (-0.050,-9.759) (-0.100,-9.758) (-0.150,-9.757) (-0.200,-9.756) (-0.250,-9.754) (-0.300,-9.752) (-0.350,-9.750) (-0.400,-9.747) (-0.450,-9.743) (-0.500,-9.740) (-0.550,-9.736) (-0.600,-9.731) (-0.650,-9.726) (-0.700,-9.721) (-0.750,-9.716) (-0.800,-9.710) (-0.850,-9.703) (-0.900,-9.696) (-0.950,-9.689) (-1.000,-9.682) (-1.050,-9.674) (-1.100,-9.665) (-1.150,-9.656) (-1.200,-9.647) (-1.250,-9.638) (-1.300,-9.628) (-1.350,-9.617) (-1.400,-9.606) (-1.450,-9.595) (-1.500,-9.583) (-1.550,-9.571) (-1.600,-9.559) (-1.650,-9.546) (-1.700,-9.533) (-1.750,-9.519) (-1.800,-9.504) (-1.850,-9.490) (-1.899,-9.475) (-1.950,-9.459) (-2.000,-9.443) (-2.050,-9.427) (-2.100,-9.410) (-2.150,-9.392) (-2.198,-9.375) (-2.250,-9.356) (-2.300,-9.337) (-2.350,-9.318) (-2.400,-9.298) (-2.450,-9.278) (-2.500,-9.257) (-2.550,-9.236) (-2.600,-9.214) (-2.650,-9.192) (-2.700,-9.169) (-2.750,-9.146) (-2.800,-9.122) (-2.850,-9.097) (-2.900,-9.072) (-2.950,-9.047) (-3.000,-9.021) (-3.050,-8.994) (-3.100,-8.967) (-3.150,-8.939) (-3.200,-8.910) (-3.250,-8.881) (-3.300,-8.851) (-3.344,-8.825) (-3.400,-8.790) (-3.450,-8.758) (-3.500,-8.726) (-3.550,-8.693) (-3.600,-8.659) (-3.650,-8.625) (-3.700,-8.589) (-3.750,-8.553) (-3.800,-8.517) (-3.850,-8.479) (-3.900,-8.441) (-3.950,-8.401) (-4.000,-8.361) (-4.045,-8.325) (-4.100,-8.279) (-4.150,-8.236) (-4.191,-8.200) (-4.247,-8.150) (-4.300,-8.101) (-4.350,-8.055) (-4.400,-8.007) (-4.450,-7.957) (-4.500,-7.907) (-4.550,-7.856) (-4.600,-7.803) (-4.649,-7.750) (-4.694,-7.700) (-4.738,-7.650) (-4.775,-7.607) (-4.823,-7.550) (-4.864,-7.500) (-4.900,-7.455) (-4.943,-7.400) (-4.975,-7.359) (-5.019,-7.300) (-5.050,-7.258) (-5.091,-7.200) (-5.125,-7.152) (-5.160,-7.100) (-5.194,-7.050) (-5.225,-7.002) (-5.258,-6.950) (-5.289,-6.900) (-5.320,-6.850) (-5.349,-6.800) (-5.375,-6.755) (-5.400,-6.711) (-5.425,-6.666) (-5.450,-6.620) (-5.474,-6.575) (-5.500,-6.525) (-5.525,-6.475) (-5.549,-6.425) (-5.573,-6.375) (-5.596,-6.325) (-5.619,-6.275) (-5.641,-6.225) (-5.662,-6.175) (-5.683,-6.125) (-5.703,-6.075) (-5.723,-6.025) (-5.742,-5.975) (-5.760,-5.925) (-5.778,-5.875) (-5.795,-5.825) (-5.812,-5.775) (-5.828,-5.725) (-5.843,-5.675) (-5.858,-5.625) (-5.873,-5.575) (-5.887,-5.525) (-5.900,-5.475) (-5.913,-5.425) (-5.925,-5.375) (-5.936,-5.325) (-5.947,-5.275) (-5.958,-5.225) (-5.968,-5.175) (-5.977,-5.125) (-5.986,-5.075) (-5.994,-5.025) (-6.002,-4.975) (-6.009,-4.925) (-6.016,-4.875) (-6.022,-4.825) (-6.028,-4.775) (-6.033,-4.725) (-6.038,-4.675) (-6.042,-4.625) (-6.045,-4.575) (-6.048,-4.525) (-6.050,-4.475) (-6.052,-4.425) (-6.054,-4.375) (-6.054,-4.325) (-6.055,-4.275) (-6.054,-4.225) (-6.053,-4.175) (-6.052,-4.125) (-6.050,-4.075) (-6.047,-4.025) (-6.044,-3.975) (-6.040,-3.925) (-6.036,-3.875) (-6.031,-3.825) (-6.026,-3.775) (-6.020,-3.725) (-6.013,-3.675) (-6.006,-3.625) (-5.998,-3.575) (-5.990,-3.525) (-5.981,-3.475) (-5.971,-3.425) (-5.961,-3.375) (-5.950,-3.325) (-5.939,-3.275) (-5.926,-3.225) (-5.914,-3.175) (-5.900,-3.125) (-5.886,-3.075) (-5.871,-3.025) (-5.856,-2.975) (-5.839,-2.925) (-5.822,-2.875) (-5.805,-2.825) (-5.786,-2.775) (-5.767,-2.725) (-5.747,-2.675) (-5.727,-2.625) (-5.705,-2.575) (-5.683,-2.525) (-5.660,-2.475) (-5.636,-2.425) (-5.611,-2.375) (-5.585,-2.325) (-5.559,-2.275) (-5.531,-2.225) (-5.503,-2.175) (-5.475,-2.128) (-5.443,-2.075) (-5.411,-2.025) (-5.378,-1.975) (-5.350,-1.933) (-5.310,-1.875) (-5.275,-1.827) (-5.237,-1.775) (-5.200,-1.727) (-5.158,-1.675) (-5.125,-1.634) (-5.075,-1.576) (-5.030,-1.525) (-4.985,-1.475) (-4.950,-1.438) (-4.913,-1.400) (-4.875,-1.362) (-4.838,-1.325) (-4.800,-1.289) (-4.758,-1.250) (-4.702,-1.200) (-4.650,-1.155) (-4.600,-1.113) (-4.553,-1.075) (-4.500,-1.034) (-4.456,-1.000) (-4.400,-0.959) (-4.352,-0.925) (-4.300,-0.889) (-4.250,-0.856) (-4.203,-0.825) (-4.150,-0.792) (-4.100,-0.761) (-4.050,-0.732) (-4.000,-0.703) (-3.950,-0.676) (-3.902,-0.650) (-3.854,-0.625) (-3.805,-0.600) (-3.754,-0.575) (-3.701,-0.550) (-3.650,-0.527) (-3.600,-0.504) (-3.550,-0.483) (-3.500,-0.462) (-3.450,-0.442) (-3.400,-0.422) (-3.350,-0.403) (-3.300,-0.385) (-3.250,-0.367) (-3.200,-0.350) (-3.150,-0.334) (-3.100,-0.318) (-3.050,-0.302) (-3.000,-0.287) (-2.950,-0.273) (-2.900,-0.259) (-2.850,-0.246) (-2.800,-0.233) (-2.750,-0.220) (-2.700,-0.209) (-2.650,-0.197) (-2.600,-0.186) (-2.550,-0.175) (-2.500,-0.165) (-2.450,-0.155) (-2.400,-0.146) (-2.350,-0.137) (-2.300,-0.128) (-2.250,-0.120) (-2.200,-0.112) (-2.150,-0.105) (-2.100,-0.098) (-2.050,-0.091) (-2.000,-0.084) (-1.950,-0.078) (-1.900,-0.072) (-1.850,-0.067) (-1.800,-0.061) (-1.750,-0.056) (-1.700,-0.052) (-1.650,-0.047) (-1.600,-0.043) (-1.550,-0.039) (-1.500,-0.035) (-1.450,-0.032) (-1.400,-0.029) (-1.350,-0.026) (-1.302,-0.023)}
\EQ{(1.300,0.023) (1.350,0.026) (1.400,0.029) (1.450,0.032) (1.500,0.035) (1.550,0.039) (1.600,0.043) (1.650,0.047) (1.700,0.052) (1.750,0.056) (1.800,0.061) (1.850,0.067) (1.900,0.072) (1.950,0.078) (2.000,0.084) (2.050,0.091) (2.100,0.098) (2.150,0.105) (2.200,0.112) (2.250,0.120) (2.300,0.128) (2.350,0.137) (2.400,0.146) (2.450,0.155) (2.500,0.165) (2.550,0.175) (2.600,0.186) (2.650,0.197) (2.700,0.209) (2.750,0.220) (2.800,0.233) (2.850,0.246) (2.900,0.259) (2.950,0.273) (3.000,0.287) (3.050,0.302) (3.100,0.318) (3.150,0.334) (3.199,0.350) (3.250,0.367) (3.300,0.385) (3.350,0.403) (3.400,0.422) (3.450,0.442) (3.500,0.462) (3.550,0.483) (3.600,0.504) (3.646,0.525) (3.700,0.550) (3.750,0.573) (3.800,0.598) (3.850,0.623) (3.900,0.649) (3.949,0.675) (3.994,0.700) (4.038,0.725) (4.100,0.761) (4.150,0.792) (4.200,0.823) (4.250,0.856) (4.300,0.889) (4.350,0.923) (4.400,0.959) (4.450,0.996) (4.500,1.034) (4.550,1.073) (4.600,1.113) (4.644,1.150) (4.700,1.198) (4.750,1.243) (4.800,1.289) (4.838,1.325) (4.875,1.362) (4.913,1.400) (4.950,1.438) (4.985,1.475) (5.025,1.519) (5.074,1.575) (5.117,1.625) (5.150,1.665) (5.198,1.725) (5.237,1.775) (5.274,1.825) (5.310,1.875) (5.345,1.925) (5.375,1.970) (5.411,2.025) (5.443,2.075) (5.473,2.125) (5.500,2.171) (5.525,2.214) (5.550,2.259) (5.575,2.306) (5.598,2.350) (5.623,2.400) (5.648,2.450) (5.671,2.500) (5.694,2.550) (5.716,2.600) (5.737,2.650) (5.757,2.700) (5.777,2.750) (5.796,2.800) (5.814,2.850) (5.831,2.900) (5.848,2.950) (5.863,3.000) (5.879,3.050) (5.893,3.100) (5.907,3.150) (5.920,3.200) (5.933,3.250) (5.944,3.300) (5.956,3.350) (5.966,3.400) (5.976,3.450) (5.985,3.500) (5.994,3.550) (6.002,3.600) (6.010,3.650) (6.017,3.700) (6.023,3.750) (6.029,3.800) (6.034,3.850) (6.038,3.900) (6.042,3.950) (6.046,4.000) (6.049,4.050) (6.051,4.100) (6.053,4.150) (6.054,4.200) (6.054,4.250) (6.054,4.300) (6.054,4.350) (6.053,4.400) (6.051,4.450) (6.049,4.500) (6.047,4.550) (6.043,4.600) (6.040,4.650) (6.035,4.700) (6.031,4.750) (6.025,4.800) (6.019,4.850) (6.013,4.900) (6.006,4.950) (5.998,5.000) (5.990,5.050) (5.982,5.100) (5.973,5.150) (5.963,5.200) (5.953,5.250) (5.942,5.300) (5.931,5.350) (5.919,5.400) (5.906,5.450) (5.893,5.500) (5.880,5.550) (5.866,5.600) (5.851,5.650) (5.836,5.700) (5.820,5.750) (5.804,5.800) (5.787,5.850) (5.769,5.900) (5.751,5.950) (5.732,6.000) (5.713,6.050) (5.693,6.100) (5.673,6.150) (5.652,6.200) (5.630,6.250) (5.608,6.300) (5.585,6.350) (5.561,6.400) (5.537,6.450) (5.512,6.500) (5.487,6.550) (5.461,6.600) (5.434,6.650) (5.406,6.700) (5.378,6.750) (5.350,6.799) (5.320,6.850) (5.289,6.900) (5.258,6.950) (5.226,7.000) (5.194,7.050) (5.160,7.100) (5.126,7.150) (5.091,7.200) (5.056,7.250) (5.025,7.292) (4.981,7.350) (4.950,7.391) (4.904,7.450) (4.864,7.500) (4.825,7.547) (4.781,7.600) (4.738,7.650) (4.700,7.693) (4.650,7.749) (4.603,7.800) (4.555,7.850) (4.507,7.900) (4.457,7.950) (4.407,8.000) (4.355,8.050) (4.302,8.100) (4.250,8.147) (4.200,8.192) (4.162,8.225) (4.104,8.275) (4.050,8.320) (4.000,8.361) (3.952,8.400) (3.900,8.441) (3.855,8.475) (3.800,8.517) (3.755,8.550) (3.700,8.589) (3.650,8.625) (3.600,8.659) (3.550,8.693) (3.501,8.725) (3.450,8.758) (3.400,8.790) (3.350,8.821) (3.302,8.850) (3.250,8.881) (3.200,8.910) (3.150,8.939) (3.100,8.967) (3.050,8.994) (3.000,9.021) (2.950,9.047) (2.900,9.072) (2.850,9.097) (2.800,9.122) (2.750,9.146) (2.700,9.169) (2.650,9.192) (2.600,9.214) (2.550,9.236) (2.500,9.257) (2.450,9.278) (2.400,9.298) (2.350,9.318) (2.300,9.337) (2.250,9.356) (2.200,9.374) (2.150,9.392) (2.100,9.410) (2.050,9.427) (2.000,9.443) (1.950,9.459) (1.900,9.475) (1.850,9.490) (1.800,9.504) (1.750,9.519) (1.700,9.533) (1.650,9.546) (1.600,9.559) (1.550,9.571) (1.500,9.583) (1.450,9.595) (1.400,9.606) (1.350,9.617) (1.300,9.628) (1.250,9.638) (1.200,9.647) (1.150,9.656) (1.100,9.665) (1.050,9.674) (1.000,9.682) (0.950,9.689) (0.900,9.696) (0.850,9.703) (0.800,9.710) (0.750,9.716) (0.700,9.721) (0.650,9.726) (0.600,9.731) (0.550,9.736) (0.500,9.740) (0.450,9.743) (0.400,9.747) (0.350,9.750) (0.300,9.752) (0.250,9.754) (0.200,9.756) (0.150,9.757) (0.100,9.758) (0.050,9.759) (0.000,9.759) (-0.050,9.759) (-0.100,9.758) (-0.150,9.757) (-0.200,9.756) (-0.250,9.754) (-0.300,9.752) (-0.350,9.750) (-0.400,9.747) (-0.450,9.743) (-0.500,9.740) (-0.550,9.736) (-0.600,9.731) (-0.650,9.726) (-0.700,9.721) (-0.750,9.716) (-0.800,9.710) (-0.850,9.703) (-0.900,9.696) (-0.950,9.689) (-1.000,9.682) (-1.050,9.674) (-1.100,9.665) (-1.150,9.656) (-1.200,9.647) (-1.250,9.638) (-1.300,9.628) (-1.350,9.617) (-1.400,9.606) (-1.450,9.595) (-1.500,9.583) (-1.550,9.571) (-1.600,9.559) (-1.650,9.546) (-1.700,9.533) (-1.750,9.519) (-1.800,9.504) (-1.850,9.490) (-1.899,9.475) (-1.950,9.459) (-2.000,9.443) (-2.050,9.427) (-2.100,9.410) (-2.150,9.392) (-2.198,9.375) (-2.250,9.356) (-2.300,9.337) (-2.350,9.318) (-2.400,9.298) (-2.450,9.278) (-2.500,9.257) (-2.550,9.236) (-2.600,9.214) (-2.650,9.192) (-2.700,9.169) (-2.750,9.146) (-2.800,9.122) (-2.850,9.097) (-2.900,9.072) (-2.950,9.047) (-3.000,9.021) (-3.050,8.994) (-3.100,8.967) (-3.150,8.939) (-3.200,8.910) (-3.250,8.881) (-3.300,8.851) (-3.344,8.825) (-3.400,8.790) (-3.450,8.758) (-3.500,8.726) (-3.550,8.693) (-3.600,8.659) (-3.650,8.625) (-3.700,8.589) (-3.750,8.553) (-3.800,8.517) (-3.850,8.479) (-3.900,8.441) (-3.950,8.401) (-4.000,8.361) (-4.045,8.325) (-4.100,8.279) (-4.150,8.236) (-4.191,8.200) (-4.247,8.150) (-4.300,8.101) (-4.350,8.055) (-4.400,8.007) (-4.450,7.957) (-4.500,7.907) (-4.550,7.856) (-4.600,7.803) (-4.649,7.750) (-4.694,7.700) (-4.738,7.650) (-4.775,7.607) (-4.823,7.550) (-4.864,7.500) (-4.900,7.455) (-4.943,7.400) (-4.975,7.359) (-5.019,7.300) (-5.050,7.258) (-5.091,7.200) (-5.125,7.152) (-5.160,7.100) (-5.194,7.050) (-5.225,7.002) (-5.258,6.950) (-5.289,6.900) (-5.320,6.850) (-5.349,6.800) (-5.375,6.755) (-5.400,6.711) (-5.425,6.666) (-5.450,6.620) (-5.474,6.575) (-5.500,6.525) (-5.525,6.475) (-5.549,6.425) (-5.573,6.375) (-5.596,6.325) (-5.619,6.275) (-5.641,6.225) (-5.662,6.175) (-5.683,6.125) (-5.703,6.075) (-5.723,6.025) (-5.742,5.975) (-5.760,5.925) (-5.778,5.875) (-5.795,5.825) (-5.812,5.775) (-5.828,5.725) (-5.843,5.675) (-5.858,5.625) (-5.873,5.575) (-5.887,5.525) (-5.900,5.475) (-5.913,5.425) (-5.925,5.375) (-5.936,5.325) (-5.947,5.275) (-5.958,5.225) (-5.968,5.175) (-5.977,5.125) (-5.986,5.075) (-5.994,5.025) (-6.002,4.975) (-6.009,4.925) (-6.016,4.875) (-6.022,4.825) (-6.028,4.775) (-6.033,4.725) (-6.038,4.675) (-6.042,4.625) (-6.045,4.575) (-6.048,4.525) (-6.050,4.475) (-6.052,4.425) (-6.054,4.375) (-6.054,4.325) (-6.055,4.275) (-6.054,4.225) (-6.053,4.175) (-6.052,4.125) (-6.050,4.075) (-6.047,4.025) (-6.044,3.975) (-6.040,3.925) (-6.036,3.875) (-6.031,3.825) (-6.026,3.775) (-6.020,3.725) (-6.013,3.675) (-6.006,3.625) (-5.998,3.575) (-5.990,3.525) (-5.981,3.475) (-5.971,3.425) (-5.961,3.375) (-5.950,3.325) (-5.939,3.275) (-5.926,3.225) (-5.914,3.175) (-5.900,3.125) (-5.886,3.075) (-5.871,3.025) (-5.856,2.975) (-5.839,2.925) (-5.822,2.875) (-5.805,2.825) (-5.786,2.775) (-5.767,2.725) (-5.747,2.675) (-5.727,2.625) (-5.705,2.575) (-5.683,2.525) (-5.660,2.475) (-5.636,2.425) (-5.611,2.375) (-5.585,2.325) (-5.559,2.275) (-5.531,2.225) (-5.503,2.175) (-5.475,2.128) (-5.443,2.075) (-5.411,2.025) (-5.378,1.975) (-5.350,1.933) (-5.310,1.875) (-5.275,1.827) (-5.237,1.775) (-5.200,1.727) (-5.158,1.675) (-5.125,1.634) (-5.075,1.576) (-5.030,1.525) (-4.985,1.475) (-4.950,1.438) (-4.913,1.400) (-4.875,1.362) (-4.838,1.325) (-4.800,1.289) (-4.758,1.250) (-4.702,1.200) (-4.650,1.155) (-4.600,1.113) (-4.553,1.075) (-4.500,1.034) (-4.456,1.000) (-4.400,0.959) (-4.352,0.925) (-4.300,0.889) (-4.250,0.856) (-4.203,0.825) (-4.150,0.792) (-4.100,0.761) (-4.050,0.732) (-4.000,0.703) (-3.950,0.676) (-3.902,0.650) (-3.854,0.625) (-3.805,0.600) (-3.754,0.575) (-3.701,0.550) (-3.650,0.527) (-3.600,0.504) (-3.550,0.483) (-3.500,0.462) (-3.450,0.442) (-3.400,0.422) (-3.350,0.403) (-3.300,0.385) (-3.250,0.367) (-3.200,0.350) (-3.150,0.334) (-3.100,0.318) (-3.050,0.302) (-3.000,0.287) (-2.950,0.273) (-2.900,0.259) (-2.850,0.246) (-2.800,0.233) (-2.750,0.220) (-2.700,0.209) (-2.650,0.197) (-2.600,0.186) (-2.550,0.175) (-2.500,0.165) (-2.450,0.155) (-2.400,0.146) (-2.350,0.137) (-2.300,0.128) (-2.250,0.120) (-2.200,0.112) (-2.150,0.105) (-2.100,0.098) (-2.050,0.091) (-2.000,0.084) (-1.950,0.078) (-1.900,0.072) (-1.850,0.067) (-1.800,0.061) (-1.750,0.056) (-1.700,0.052) (-1.650,0.047) (-1.600,0.043) (-1.550,0.039) (-1.500,0.035) (-1.450,0.032) (-1.400,0.029) (-1.350,0.026) (-1.302,0.023)}
\EQ{(1.300,0.046) (1.350,0.052) (1.400,0.058) (1.450,0.064) (1.500,0.071) (1.550,0.079) (1.600,0.086) (1.650,0.095) (1.700,0.104) (1.750,0.113) (1.800,0.123) (1.850,0.134) (1.900,0.145) (1.950,0.157) (2.000,0.170) (2.050,0.183) (2.100,0.197) (2.150,0.212) (2.200,0.227) (2.250,0.243) (2.300,0.260) (2.350,0.278) (2.400,0.297) (2.450,0.317) (2.500,0.337) (2.550,0.359) (2.600,0.381) (2.650,0.405) (2.700,0.429) (2.750,0.455) (2.800,0.482) (2.850,0.510) (2.900,0.539) (2.950,0.570) (2.998,0.600) (3.050,0.635) (3.100,0.670) (3.141,0.700) (3.200,0.745) (3.250,0.785) (3.298,0.825) (3.350,0.871) (3.400,0.917) (3.450,0.966) (3.500,1.017) (3.550,1.071) (3.598,1.125) (3.640,1.175) (3.675,1.219) (3.718,1.275) (3.750,1.319) (3.789,1.375) (3.822,1.425) (3.850,1.470) (3.883,1.525) (3.912,1.575) (3.939,1.625) (3.965,1.675) (3.990,1.725) (4.014,1.775) (4.036,1.825) (4.057,1.875) (4.077,1.925) (4.096,1.975) (4.114,2.025) (4.131,2.075) (4.147,2.125) (4.162,2.175) (4.176,2.225) (4.189,2.275) (4.201,2.325) (4.213,2.375) (4.223,2.425) (4.232,2.475) (4.241,2.525) (4.249,2.575) (4.256,2.625) (4.262,2.675) (4.267,2.725) (4.271,2.775) (4.275,2.825) (4.278,2.875) (4.280,2.925) (4.281,2.975) (4.281,3.025) (4.281,3.075) (4.280,3.125) (4.278,3.175) (4.275,3.225) (4.272,3.275) (4.268,3.325) (4.263,3.375) (4.257,3.425) (4.250,3.475) (4.243,3.525) (4.235,3.575) (4.226,3.625) (4.217,3.675) (4.206,3.725) (4.195,3.775) (4.184,3.825) (4.171,3.875) (4.157,3.925) (4.143,3.975) (4.128,4.025) (4.112,4.075) (4.096,4.125) (4.078,4.175) (4.060,4.225) (4.041,4.275) (4.021,4.325) (4.000,4.375) (3.979,4.425) (3.956,4.475) (3.933,4.525) (3.908,4.575) (3.883,4.625) (3.857,4.675) (3.830,4.725) (3.802,4.775) (3.775,4.821) (3.743,4.875) (3.711,4.925) (3.679,4.975) (3.650,5.019) (3.611,5.075) (3.576,5.125) (3.539,5.175) (3.501,5.225) (3.462,5.275) (3.425,5.321) (3.380,5.375) (3.336,5.425) (3.300,5.466) (3.250,5.520) (3.200,5.573) (3.150,5.623) (3.100,5.672) (3.050,5.720) (3.000,5.766) (2.961,5.800) (2.904,5.850) (2.850,5.895) (2.800,5.935) (2.750,5.974) (2.700,6.013) (2.650,6.049) (2.600,6.085) (2.550,6.120) (2.505,6.150) (2.450,6.186) (2.400,6.218) (2.350,6.249) (2.306,6.275) (2.250,6.308) (2.200,6.336) (2.150,6.363) (2.100,6.390) (2.050,6.415) (2.000,6.440) (1.950,6.464) (1.900,6.488) (1.850,6.510) (1.800,6.532) (1.750,6.553) (1.700,6.574) (1.650,6.593) (1.600,6.612) (1.550,6.631) (1.500,6.648) (1.450,6.665) (1.400,6.682) (1.350,6.698) (1.300,6.713) (1.250,6.727) (1.200,6.741) (1.150,6.754) (1.100,6.767) (1.050,6.779) (1.000,6.791) (0.950,6.801) (0.900,6.812) (0.850,6.821) (0.800,6.831) (0.750,6.839) (0.700,6.847) (0.650,6.854) (0.600,6.861) (0.550,6.868) (0.500,6.873) (0.450,6.879) (0.400,6.883) (0.350,6.887) (0.300,6.891) (0.250,6.894) (0.200,6.896) (0.150,6.898) (0.100,6.900) (0.050,6.900) (0.000,6.901) (-0.050,6.900) (-0.100,6.900) (-0.150,6.898) (-0.200,6.896) (-0.250,6.894) (-0.300,6.891) (-0.350,6.887) (-0.400,6.883) (-0.450,6.879) (-0.500,6.873) (-0.550,6.868) (-0.600,6.861) (-0.650,6.854) (-0.700,6.847) (-0.750,6.839) (-0.800,6.831) (-0.850,6.821) (-0.900,6.812) (-0.950,6.801) (-1.000,6.791) (-1.050,6.779) (-1.100,6.767) (-1.150,6.754) (-1.200,6.741) (-1.250,6.727) (-1.300,6.713) (-1.350,6.698) (-1.400,6.682) (-1.450,6.665) (-1.500,6.648) (-1.550,6.631) (-1.600,6.612) (-1.650,6.593) (-1.697,6.575) (-1.750,6.553) (-1.800,6.532) (-1.850,6.510) (-1.900,6.488) (-1.950,6.464) (-2.000,6.440) (-2.050,6.415) (-2.100,6.390) (-2.150,6.363) (-2.200,6.336) (-2.250,6.308) (-2.300,6.279) (-2.348,6.250) (-2.400,6.218) (-2.450,6.186) (-2.500,6.154) (-2.543,6.125) (-2.600,6.085) (-2.649,6.050) (-2.700,6.013) (-2.749,5.975) (-2.800,5.935) (-2.844,5.900) (-2.900,5.853) (-2.950,5.810) (-2.990,5.775) (-3.044,5.725) (-3.097,5.675) (-3.148,5.625) (-3.198,5.575) (-3.246,5.525) (-3.292,5.475) (-3.336,5.425) (-3.375,5.380) (-3.421,5.325) (-3.462,5.275) (-3.500,5.227) (-3.539,5.175) (-3.575,5.126) (-3.611,5.075) (-3.646,5.025) (-3.675,4.981) (-3.711,4.925) (-3.743,4.875) (-3.773,4.825) (-3.800,4.778) (-3.825,4.734) (-3.850,4.688) (-3.875,4.641) (-3.900,4.592) (-3.925,4.541) (-3.950,4.488) (-3.975,4.433) (-4.000,4.376) (-4.021,4.325) (-4.041,4.275) (-4.060,4.225) (-4.078,4.175) (-4.096,4.125) (-4.112,4.075) (-4.128,4.025) (-4.143,3.975) (-4.157,3.925) (-4.171,3.875) (-4.184,3.825) (-4.195,3.775) (-4.206,3.725) (-4.217,3.675) (-4.226,3.625) (-4.235,3.575) (-4.243,3.525) (-4.250,3.475) (-4.257,3.425) (-4.263,3.375) (-4.268,3.325) (-4.272,3.275) (-4.275,3.225) (-4.278,3.175) (-4.280,3.125) (-4.281,3.075) (-4.281,3.025) (-4.281,2.975) (-4.280,2.925) (-4.278,2.875) (-4.275,2.825) (-4.271,2.775) (-4.267,2.725) (-4.262,2.675) (-4.256,2.625) (-4.249,2.575) (-4.241,2.525) (-4.232,2.475) (-4.223,2.425) (-4.213,2.375) (-4.201,2.325) (-4.189,2.275) (-4.176,2.225) (-4.162,2.175) (-4.147,2.125) (-4.131,2.075) (-4.114,2.025) (-4.096,1.975) (-4.077,1.925) (-4.057,1.875) (-4.036,1.825) (-4.014,1.775) (-3.990,1.725) (-3.965,1.675) (-3.939,1.625) (-3.912,1.575) (-3.883,1.525) (-3.853,1.475) (-3.825,1.430) (-3.789,1.375) (-3.754,1.325) (-3.718,1.275) (-3.680,1.225) (-3.640,1.175) (-3.600,1.127) (-3.554,1.075) (-3.508,1.025) (-3.459,0.975) (-3.408,0.925) (-3.355,0.875) (-3.300,0.827) (-3.250,0.785) (-3.207,0.750) (-3.150,0.707) (-3.107,0.675) (-3.050,0.635) (-3.000,0.602) (-2.950,0.570) (-2.900,0.539) (-2.850,0.510) (-2.800,0.482) (-2.750,0.455) (-2.700,0.429) (-2.650,0.405) (-2.600,0.381) (-2.550,0.359) (-2.500,0.337) (-2.450,0.317) (-2.400,0.297) (-2.350,0.278) (-2.300,0.260) (-2.250,0.243) (-2.200,0.227) (-2.150,0.212) (-2.100,0.197) (-2.050,0.183) (-2.000,0.170) (-1.950,0.157) (-1.900,0.145) (-1.850,0.134) (-1.800,0.123) (-1.750,0.113) (-1.700,0.104) (-1.650,0.095) (-1.600,0.086) (-1.550,0.079) (-1.500,0.071) (-1.450,0.064) (-1.400,0.058) (-1.350,0.052) (-1.300,0.046)}
\EQ{(1.300,0.070) (1.350,0.078) (1.400,0.087) (1.450,0.097) (1.500,0.107) (1.550,0.118) (1.600,0.130) (1.650,0.143) (1.700,0.157) (1.750,0.171) (1.800,0.187) (1.850,0.203) (1.900,0.220) (1.950,0.239) (2.000,0.258) (2.050,0.279) (2.098,0.300) (2.150,0.324) (2.200,0.348) (2.250,0.374) (2.298,0.400) (2.342,0.425) (2.400,0.460) (2.450,0.492) (2.500,0.525) (2.550,0.561) (2.600,0.598) (2.650,0.638) (2.694,0.675) (2.750,0.724) (2.800,0.772) (2.850,0.822) (2.899,0.875) (2.943,0.925) (2.985,0.975) (3.024,1.025) (3.060,1.075) (3.095,1.125) (3.125,1.171) (3.159,1.225) (3.188,1.275) (3.216,1.325) (3.241,1.375) (3.266,1.425) (3.289,1.475) (3.310,1.525) (3.330,1.575) (3.349,1.625) (3.367,1.675) (3.383,1.725) (3.398,1.775) (3.412,1.825) (3.425,1.875) (3.437,1.925) (3.447,1.975) (3.456,2.025) (3.465,2.075) (3.472,2.125) (3.479,2.175) (3.484,2.225) (3.488,2.275) (3.491,2.325) (3.494,2.375) (3.495,2.425) (3.496,2.475) (3.495,2.525) (3.494,2.575) (3.491,2.625) (3.488,2.675) (3.483,2.725) (3.478,2.775) (3.472,2.825) (3.465,2.875) (3.457,2.925) (3.448,2.975) (3.438,3.025) (3.427,3.075) (3.415,3.125) (3.403,3.175) (3.389,3.225) (3.374,3.275) (3.359,3.325) (3.342,3.375) (3.325,3.423) (3.306,3.475) (3.286,3.525) (3.265,3.575) (3.243,3.625) (3.220,3.675) (3.196,3.725) (3.171,3.775) (3.145,3.825) (3.118,3.875) (3.089,3.925) (3.059,3.975) (3.028,4.025) (3.000,4.068) (2.962,4.125) (2.927,4.175) (2.890,4.225) (2.852,4.275) (2.812,4.325) (2.775,4.370) (2.728,4.425) (2.683,4.475) (2.637,4.525) (2.600,4.563) (2.563,4.600) (2.525,4.637) (2.485,4.675) (2.430,4.725) (2.375,4.773) (2.325,4.815) (2.281,4.850) (2.225,4.894) (2.183,4.925) (2.125,4.967) (2.078,5.000) (2.025,5.035) (1.975,5.067) (1.925,5.098) (1.881,5.125) (1.825,5.157) (1.775,5.185) (1.725,5.212) (1.675,5.238) (1.625,5.262) (1.575,5.286) (1.525,5.309) (1.475,5.331) (1.425,5.352) (1.375,5.373) (1.325,5.392) (1.275,5.411) (1.225,5.428) (1.175,5.445) (1.125,5.462) (1.075,5.477) (1.025,5.492) (0.975,5.505) (0.925,5.519) (0.875,5.531) (0.825,5.543) (0.775,5.553) (0.725,5.564) (0.675,5.573) (0.625,5.582) (0.575,5.590) (0.525,5.597) (0.475,5.604) (0.425,5.610) (0.375,5.616) (0.325,5.620) (0.275,5.624) (0.225,5.628) (0.175,5.630) (0.125,5.632) (0.075,5.634) (0.025,5.634) (-0.050,5.634) (-0.100,5.633) (-0.150,5.631) (-0.200,5.629) (-0.250,5.626) (-0.300,5.622) (-0.350,5.618) (-0.400,5.613) (-0.450,5.607) (-0.500,5.601) (-0.550,5.594) (-0.600,5.586) (-0.650,5.578) (-0.700,5.569) (-0.750,5.559) (-0.800,5.548) (-0.850,5.537) (-0.899,5.525) (-0.950,5.512) (-1.000,5.499) (-1.050,5.484) (-1.100,5.469) (-1.150,5.453) (-1.200,5.437) (-1.250,5.420) (-1.300,5.401) (-1.350,5.382) (-1.400,5.363) (-1.450,5.342) (-1.500,5.320) (-1.550,5.298) (-1.599,5.275) (-1.650,5.250) (-1.700,5.225) (-1.747,5.200) (-1.793,5.175) (-1.838,5.150) (-1.900,5.114) (-1.950,5.083) (-2.000,5.051) (-2.050,5.019) (-2.100,4.985) (-2.149,4.950) (-2.200,4.913) (-2.249,4.875) (-2.300,4.835) (-2.343,4.800) (-2.400,4.751) (-2.450,4.707) (-2.500,4.661) (-2.538,4.625) (-2.575,4.588) (-2.613,4.550) (-2.650,4.511) (-2.700,4.457) (-2.750,4.400) (-2.792,4.350) (-2.825,4.309) (-2.871,4.250) (-2.908,4.200) (-2.944,4.150) (-2.975,4.106) (-3.012,4.050) (-3.044,4.000) (-3.074,3.950) (-3.100,3.906) (-3.125,3.862) (-3.150,3.816) (-3.175,3.768) (-3.200,3.718) (-3.225,3.665) (-3.250,3.610) (-3.275,3.551) (-3.296,3.500) (-3.315,3.450) (-3.333,3.400) (-3.350,3.351) (-3.366,3.300) (-3.382,3.250) (-3.396,3.200) (-3.409,3.150) (-3.421,3.100) (-3.433,3.050) (-3.443,3.000) (-3.453,2.950) (-3.461,2.900) (-3.469,2.850) (-3.475,2.800) (-3.481,2.750) (-3.486,2.700) (-3.490,2.650) (-3.492,2.600) (-3.494,2.550) (-3.495,2.500) (-3.496,2.450) (-3.495,2.400) (-3.493,2.350) (-3.490,2.300) (-3.486,2.250) (-3.481,2.200) (-3.475,2.150) (-3.469,2.100) (-3.461,2.050) (-3.452,2.000) (-3.442,1.950) (-3.431,1.900) (-3.419,1.850) (-3.405,1.800) (-3.391,1.750) (-3.375,1.700) (-3.358,1.650) (-3.340,1.600) (-3.321,1.550) (-3.300,1.501) (-3.277,1.450) (-3.254,1.400) (-3.229,1.350) (-3.202,1.300) (-3.175,1.252) (-3.144,1.200) (-3.112,1.150) (-3.078,1.100) (-3.042,1.050) (-3.004,1.000) (-2.964,0.950) (-2.925,0.904) (-2.876,0.850) (-2.828,0.800) (-2.777,0.750) (-2.725,0.702) (-2.675,0.659) (-2.634,0.625) (-2.575,0.579) (-2.525,0.543) (-2.475,0.508) (-2.425,0.475) (-2.375,0.444) (-2.325,0.415) (-2.275,0.387) (-2.225,0.361) (-2.175,0.336) (-2.125,0.312) (-2.075,0.290) (-2.025,0.269) (-1.975,0.248) (-1.925,0.230) (-1.875,0.212) (-1.825,0.195) (-1.775,0.179) (-1.725,0.164) (-1.676,0.150) (-1.625,0.137) (-1.575,0.124) (-1.525,0.113) (-1.475,0.102) (-1.425,0.092) (-1.375,0.082) (-1.325,0.074) (-1.300,0.070)}
\EQ{(1.300,0.093) (1.350,0.104) (1.400,0.116) (1.450,0.130) (1.500,0.144) (1.550,0.159) (1.600,0.175) (1.650,0.193) (1.700,0.211) (1.750,0.231) (1.800,0.252) (1.850,0.275) (1.900,0.299) (1.950,0.324) (1.997,0.350) (2.041,0.375) (2.100,0.412) (2.150,0.444) (2.194,0.475) (2.250,0.517) (2.292,0.550) (2.350,0.599) (2.400,0.645) (2.450,0.693) (2.500,0.746) (2.548,0.800) (2.589,0.850) (2.625,0.896) (2.664,0.950) (2.697,1.000) (2.725,1.044) (2.759,1.100) (2.786,1.150) (2.812,1.200) (2.836,1.250) (2.858,1.300) (2.879,1.350) (2.898,1.400) (2.916,1.450) (2.932,1.500) (2.947,1.550) (2.960,1.600) (2.972,1.650) (2.983,1.700) (2.993,1.750) (3.001,1.800) (3.008,1.850) (3.014,1.900) (3.019,1.950) (3.023,2.000) (3.025,2.050) (3.027,2.100) (3.027,2.150) (3.026,2.200) (3.025,2.250) (3.022,2.300) (3.018,2.350) (3.013,2.400) (3.006,2.450) (2.999,2.500) (2.991,2.550) (2.981,2.600) (2.971,2.650) (2.959,2.700) (2.947,2.750) (2.933,2.800) (2.918,2.850) (2.902,2.900) (2.885,2.950) (2.866,3.000) (2.847,3.050) (2.826,3.100) (2.804,3.150) (2.781,3.200) (2.756,3.250) (2.730,3.300) (2.703,3.350) (2.675,3.399) (2.645,3.450) (2.613,3.500) (2.580,3.550) (2.550,3.594) (2.509,3.650) (2.475,3.696) (2.432,3.750) (2.400,3.789) (2.350,3.846) (2.301,3.900) (2.253,3.950) (2.203,4.000) (2.151,4.050) (2.100,4.096) (2.050,4.139) (2.007,4.175) (1.950,4.220) (1.900,4.258) (1.850,4.295) (1.807,4.325) (1.750,4.363) (1.700,4.395) (1.651,4.425) (1.600,4.455) (1.550,4.483) (1.500,4.510) (1.450,4.536) (1.400,4.561) (1.350,4.585) (1.300,4.607) (1.250,4.629) (1.200,4.649) (1.150,4.669) (1.100,4.687) (1.050,4.705) (1.000,4.722) (0.950,4.737) (0.900,4.752) (0.850,4.766) (0.800,4.779) (0.750,4.792) (0.700,4.803) (0.650,4.814) (0.600,4.824) (0.550,4.833) (0.500,4.841) (0.450,4.848) (0.400,4.855) (0.350,4.861) (0.300,4.866) (0.250,4.870) (0.200,4.873) (0.150,4.876) (0.100,4.878) (0.050,4.879) (0.000,4.880) (-0.050,4.879) (-0.100,4.878) (-0.150,4.876) (-0.200,4.873) (-0.250,4.870) (-0.300,4.866) (-0.350,4.861) (-0.400,4.855) (-0.450,4.848) (-0.500,4.841) (-0.550,4.833) (-0.600,4.824) (-0.650,4.814) (-0.700,4.803) (-0.750,4.792) (-0.800,4.779) (-0.850,4.766) (-0.900,4.752) (-0.950,4.737) (-1.000,4.722) (-1.050,4.705) (-1.100,4.687) (-1.150,4.669) (-1.198,4.650) (-1.250,4.629) (-1.300,4.607) (-1.350,4.585) (-1.400,4.561) (-1.450,4.536) (-1.500,4.510) (-1.550,4.483) (-1.600,4.455) (-1.650,4.426) (-1.700,4.395) (-1.750,4.363) (-1.800,4.330) (-1.843,4.300) (-1.900,4.258) (-1.944,4.225) (-2.000,4.181) (-2.050,4.139) (-2.095,4.100) (-2.150,4.051) (-2.200,4.003) (-2.250,3.954) (-2.300,3.901) (-2.347,3.850) (-2.390,3.800) (-2.425,3.758) (-2.472,3.700) (-2.509,3.650) (-2.546,3.600) (-2.575,3.558) (-2.613,3.500) (-2.645,3.450) (-2.675,3.400) (-2.700,3.356) (-2.725,3.310) (-2.750,3.262) (-2.775,3.212) (-2.800,3.159) (-2.825,3.102) (-2.847,3.050) (-2.866,3.000) (-2.885,2.950) (-2.902,2.900) (-2.918,2.850) (-2.933,2.800) (-2.947,2.750) (-2.959,2.700) (-2.971,2.650) (-2.981,2.600) (-2.991,2.550) (-2.999,2.500) (-3.006,2.450) (-3.013,2.400) (-3.018,2.350) (-3.022,2.300) (-3.025,2.250) (-3.026,2.200) (-3.027,2.150) (-3.027,2.100) (-3.025,2.050) (-3.023,2.000) (-3.019,1.950) (-3.014,1.900) (-3.008,1.850) (-3.001,1.800) (-2.993,1.750) (-2.983,1.700) (-2.972,1.650) (-2.960,1.600) (-2.947,1.550) (-2.932,1.500) (-2.916,1.450) (-2.898,1.400) (-2.879,1.350) (-2.858,1.300) (-2.836,1.250) (-2.812,1.200) (-2.786,1.150) (-2.759,1.100) (-2.729,1.050) (-2.700,1.004) (-2.664,0.950) (-2.628,0.900) (-2.589,0.850) (-2.550,0.802) (-2.504,0.750) (-2.457,0.700) (-2.406,0.650) (-2.351,0.600) (-2.300,0.557) (-2.261,0.525) (-2.200,0.479) (-2.158,0.450) (-2.100,0.412) (-2.050,0.381) (-2.000,0.352) (-1.951,0.325) (-1.902,0.300) (-1.850,0.275) (-1.800,0.252) (-1.750,0.231) (-1.700,0.211) (-1.650,0.193) (-1.600,0.175) (-1.550,0.159) (-1.500,0.144) (-1.450,0.130) (-1.400,0.116) (-1.350,0.104) (-1.300,0.093)}
\EQ{(1.300,0.117) (1.350,0.131) (1.400,0.146) (1.450,0.163) (1.500,0.181) (1.549,0.200) (1.600,0.221) (1.650,0.244) (1.700,0.268) (1.750,0.293) (1.800,0.321) (1.849,0.350) (1.900,0.382) (1.950,0.416) (1.997,0.450) (2.050,0.492) (2.090,0.525) (2.145,0.575) (2.196,0.625) (2.243,0.675) (2.287,0.725) (2.325,0.772) (2.365,0.825) (2.400,0.875) (2.432,0.925) (2.462,0.975) (2.490,1.025) (2.516,1.075) (2.540,1.125) (2.562,1.175) (2.582,1.225) (2.600,1.275) (2.617,1.325) (2.632,1.375) (2.646,1.425) (2.658,1.475) (2.669,1.525) (2.679,1.575) (2.687,1.625) (2.693,1.675) (2.699,1.725) (2.703,1.775) (2.706,1.825) (2.707,1.875) (2.708,1.925) (2.707,1.975) (2.705,2.025) (2.701,2.075) (2.697,2.125) (2.691,2.175) (2.684,2.225) (2.676,2.275) (2.667,2.325) (2.656,2.375) (2.644,2.425) (2.631,2.475) (2.617,2.525) (2.602,2.575) (2.585,2.625) (2.567,2.675) (2.547,2.725) (2.527,2.775) (2.504,2.825) (2.481,2.875) (2.456,2.925) (2.430,2.975) (2.402,3.025) (2.375,3.070) (2.341,3.125) (2.308,3.175) (2.275,3.223) (2.237,3.275) (2.200,3.323) (2.158,3.375) (2.125,3.414) (2.075,3.470) (2.025,3.523) (1.975,3.574) (1.925,3.621) (1.875,3.667) (1.836,3.700) (1.776,3.750) (1.725,3.790) (1.677,3.825) (1.625,3.862) (1.575,3.896) (1.530,3.925) (1.475,3.959) (1.425,3.988) (1.375,4.016) (1.325,4.043) (1.275,4.068) (1.225,4.092) (1.175,4.115) (1.125,4.137) (1.075,4.158) (1.025,4.177) (0.975,4.196) (0.925,4.213) (0.875,4.229) (0.825,4.245) (0.775,4.259) (0.725,4.272) (0.675,4.285) (0.625,4.296) (0.575,4.307) (0.525,4.317) (0.475,4.325) (0.425,4.333) (0.375,4.340) (0.325,4.346) (0.275,4.351) (0.225,4.356) (0.175,4.359) (0.125,4.362) (0.075,4.363) (0.025,4.364) (-0.050,4.364) (-0.100,4.363) (-0.150,4.361) (-0.200,4.358) (-0.250,4.354) (-0.300,4.349) (-0.350,4.343) (-0.400,4.337) (-0.450,4.329) (-0.500,4.321) (-0.550,4.312) (-0.600,4.302) (-0.650,4.291) (-0.700,4.279) (-0.750,4.266) (-0.800,4.252) (-0.850,4.237) (-0.900,4.221) (-0.950,4.204) (-1.000,4.187) (-1.050,4.168) (-1.100,4.147) (-1.150,4.126) (-1.200,4.104) (-1.250,4.080) (-1.300,4.056) (-1.350,4.030) (-1.400,4.002) (-1.448,3.975) (-1.500,3.944) (-1.550,3.912) (-1.600,3.879) (-1.642,3.850) (-1.700,3.808) (-1.744,3.775) (-1.800,3.730) (-1.850,3.688) (-1.894,3.650) (-1.948,3.600) (-1.999,3.550) (-2.047,3.500) (-2.093,3.450) (-2.137,3.400) (-2.175,3.354) (-2.218,3.300) (-2.250,3.257) (-2.291,3.200) (-2.325,3.150) (-2.357,3.100) (-2.387,3.050) (-2.416,3.000) (-2.443,2.950) (-2.469,2.900) (-2.493,2.850) (-2.516,2.800) (-2.537,2.750) (-2.557,2.700) (-2.575,2.653) (-2.593,2.600) (-2.610,2.550) (-2.624,2.500) (-2.638,2.450) (-2.650,2.400) (-2.662,2.350) (-2.672,2.300) (-2.680,2.250) (-2.688,2.200) (-2.694,2.150) (-2.699,2.100) (-2.703,2.050) (-2.706,2.000) (-2.707,1.950) (-2.708,1.900) (-2.707,1.850) (-2.704,1.800) (-2.701,1.750) (-2.696,1.700) (-2.690,1.650) (-2.683,1.600) (-2.674,1.550) (-2.664,1.500) (-2.652,1.450) (-2.639,1.400) (-2.625,1.351) (-2.609,1.300) (-2.591,1.250) (-2.572,1.200) (-2.551,1.150) (-2.528,1.100) (-2.503,1.050) (-2.476,1.000) (-2.450,0.954) (-2.416,0.900) (-2.383,0.850) (-2.350,0.805) (-2.308,0.750) (-2.275,0.711) (-2.225,0.655) (-2.175,0.604) (-2.125,0.557) (-2.075,0.513) (-2.029,0.475) (-1.975,0.434) (-1.927,0.400) (-1.875,0.366) (-1.825,0.335) (-1.775,0.307) (-1.725,0.280) (-1.675,0.255) (-1.625,0.232) (-1.575,0.211) (-1.525,0.191) (-1.475,0.172) (-1.425,0.155) (-1.375,0.139) (-1.325,0.124) (-1.300,0.117)}
\EQ{(1.300,0.141) (1.350,0.158) (1.400,0.177) (1.450,0.197) (1.500,0.220) (1.550,0.243) (1.600,0.269) (1.650,0.297) (1.697,0.325) (1.750,0.359) (1.800,0.394) (1.841,0.425) (1.900,0.473) (1.950,0.517) (2.000,0.566) (2.050,0.619) (2.099,0.675) (2.138,0.725) (2.174,0.775) (2.208,0.825) (2.239,0.875) (2.267,0.925) (2.293,0.975) (2.317,1.025) (2.339,1.075) (2.359,1.125) (2.377,1.175) (2.394,1.225) (2.409,1.275) (2.422,1.325) (2.433,1.375) (2.443,1.425) (2.451,1.475) (2.458,1.525) (2.464,1.575) (2.468,1.625) (2.470,1.675) (2.472,1.725) (2.472,1.775) (2.470,1.825) (2.467,1.875) (2.463,1.925) (2.458,1.975) (2.451,2.025) (2.443,2.075) (2.434,2.125) (2.423,2.175) (2.411,2.225) (2.398,2.275) (2.383,2.325) (2.367,2.375) (2.350,2.423) (2.330,2.475) (2.310,2.525) (2.288,2.575) (2.265,2.625) (2.240,2.675) (2.213,2.725) (2.184,2.775) (2.154,2.825) (2.125,2.871) (2.089,2.925) (2.053,2.975) (2.015,3.025) (1.975,3.075) (1.933,3.125) (1.888,3.175) (1.850,3.215) (1.800,3.265) (1.763,3.300) (1.725,3.335) (1.679,3.375) (1.625,3.420) (1.575,3.459) (1.525,3.496) (1.475,3.531) (1.425,3.565) (1.375,3.597) (1.328,3.625) (1.285,3.650) (1.239,3.675) (1.191,3.700) (1.141,3.725) (1.087,3.750) (1.031,3.775) (0.975,3.798) (0.925,3.817) (0.875,3.835) (0.825,3.852) (0.775,3.868) (0.725,3.883) (0.675,3.897) (0.625,3.909) (0.575,3.921) (0.525,3.932) (0.475,3.941) (0.425,3.950) (0.375,3.958) (0.325,3.964) (0.275,3.970) (0.225,3.975) (0.175,3.978) (0.125,3.981) (0.075,3.983) (0.025,3.984) (-0.050,3.984) (-0.100,3.982) (-0.150,3.980) (-0.200,3.977) (-0.250,3.972) (-0.300,3.967) (-0.350,3.961) (-0.400,3.954) (-0.450,3.946) (-0.500,3.937) (-0.550,3.926) (-0.600,3.915) (-0.650,3.903) (-0.700,3.890) (-0.750,3.876) (-0.800,3.860) (-0.850,3.844) (-0.900,3.826) (-0.950,3.808) (-1.000,3.788) (-1.050,3.767) (-1.100,3.744) (-1.150,3.721) (-1.200,3.696) (-1.250,3.669) (-1.300,3.641) (-1.350,3.612) (-1.400,3.581) (-1.448,3.550) (-1.500,3.514) (-1.550,3.478) (-1.600,3.440) (-1.650,3.400) (-1.700,3.357) (-1.750,3.312) (-1.790,3.275) (-1.840,3.225) (-1.875,3.189) (-1.910,3.150) (-1.950,3.105) (-1.995,3.050) (-2.034,3.000) (-2.071,2.950) (-2.100,2.908) (-2.139,2.850) (-2.170,2.800) (-2.199,2.750) (-2.225,2.703) (-2.250,2.654) (-2.275,2.603) (-2.299,2.550) (-2.320,2.500) (-2.340,2.450) (-2.358,2.400) (-2.375,2.351) (-2.391,2.300) (-2.405,2.250) (-2.417,2.200) (-2.429,2.150) (-2.439,2.100) (-2.447,2.050) (-2.455,2.000) (-2.461,1.950) (-2.466,1.900) (-2.469,1.850) (-2.471,1.800) (-2.472,1.750) (-2.471,1.700) (-2.469,1.650) (-2.466,1.600) (-2.461,1.550) (-2.455,1.500) (-2.447,1.450) (-2.438,1.400) (-2.427,1.350) (-2.415,1.300) (-2.401,1.250) (-2.386,1.200) (-2.369,1.150) (-2.350,1.101) (-2.329,1.050) (-2.306,1.000) (-2.281,0.950) (-2.253,0.900) (-2.225,0.852) (-2.192,0.800) (-2.157,0.750) (-2.125,0.708) (-2.078,0.650) (-2.033,0.600) (-1.984,0.550) (-1.931,0.500) (-1.875,0.452) (-1.825,0.413) (-1.775,0.376) (-1.725,0.343) (-1.675,0.312) (-1.625,0.283) (-1.575,0.256) (-1.525,0.231) (-1.475,0.208) (-1.425,0.187) (-1.375,0.167) (-1.327,0.150) (-1.300,0.141)}
\EQ{(1.300,0.165) (1.350,0.186) (1.400,0.208) (1.450,0.233) (1.500,0.259) (1.550,0.288) (1.600,0.319) (1.646,0.350) (1.700,0.390) (1.744,0.425) (1.800,0.474) (1.850,0.522) (1.900,0.575) (1.943,0.625) (1.975,0.665) (2.018,0.725) (2.050,0.773) (2.082,0.825) (2.109,0.875) (2.134,0.925) (2.157,0.975) (2.178,1.025) (2.197,1.075) (2.213,1.125) (2.228,1.175) (2.242,1.225) (2.253,1.275) (2.263,1.325) (2.271,1.375) (2.277,1.425) (2.282,1.475) (2.286,1.525) (2.288,1.575) (2.288,1.625) (2.288,1.675) (2.285,1.725) (2.281,1.775) (2.276,1.825) (2.269,1.875) (2.261,1.925) (2.252,1.975) (2.241,2.025) (2.228,2.075) (2.215,2.125) (2.199,2.175) (2.182,2.225) (2.164,2.275) (2.144,2.325) (2.122,2.375) (2.100,2.423) (2.075,2.473) (2.050,2.520) (2.025,2.565) (1.989,2.625) (1.956,2.675) (1.925,2.720) (1.885,2.775) (1.850,2.820) (1.805,2.875) (1.761,2.925) (1.725,2.964) (1.690,3.000) (1.650,3.039) (1.611,3.075) (1.554,3.125) (1.500,3.170) (1.450,3.209) (1.400,3.245) (1.358,3.275) (1.300,3.313) (1.250,3.344) (1.200,3.373) (1.152,3.400) (1.104,3.425) (1.053,3.450) (1.000,3.475) (0.950,3.497) (0.900,3.517) (0.850,3.536) (0.800,3.554) (0.750,3.571) (0.700,3.587) (0.650,3.601) (0.600,3.614) (0.550,3.626) (0.500,3.637) (0.450,3.647) (0.400,3.656) (0.350,3.664) (0.300,3.670) (0.250,3.676) (0.200,3.680) (0.150,3.684) (0.100,3.687) (0.050,3.688) (0.000,3.689) (-0.050,3.688) (-0.100,3.687) (-0.150,3.684) (-0.200,3.680) (-0.250,3.676) (-0.300,3.670) (-0.350,3.664) (-0.400,3.656) (-0.450,3.647) (-0.500,3.637) (-0.550,3.626) (-0.600,3.614) (-0.650,3.601) (-0.700,3.587) (-0.750,3.571) (-0.800,3.554) (-0.850,3.536) (-0.900,3.517) (-0.950,3.497) (-1.000,3.475) (-1.050,3.452) (-1.100,3.427) (-1.150,3.401) (-1.197,3.375) (-1.250,3.344) (-1.300,3.313) (-1.350,3.280) (-1.394,3.250) (-1.450,3.209) (-1.494,3.175) (-1.550,3.129) (-1.600,3.085) (-1.638,3.050) (-1.675,3.015) (-1.714,2.975) (-1.750,2.937) (-1.800,2.881) (-1.846,2.825) (-1.885,2.775) (-1.922,2.725) (-1.956,2.675) (-1.989,2.625) (-2.019,2.575) (-2.047,2.525) (-2.074,2.475) (-2.099,2.425) (-2.122,2.375) (-2.144,2.325) (-2.164,2.275) (-2.182,2.225) (-2.199,2.175) (-2.215,2.125) (-2.228,2.075) (-2.241,2.025) (-2.252,1.975) (-2.261,1.925) (-2.269,1.875) (-2.276,1.825) (-2.281,1.775) (-2.285,1.725) (-2.288,1.675) (-2.288,1.625) (-2.288,1.575) (-2.286,1.525) (-2.282,1.475) (-2.277,1.425) (-2.271,1.375) (-2.263,1.325) (-2.253,1.275) (-2.242,1.225) (-2.228,1.175) (-2.213,1.125) (-2.197,1.075) (-2.178,1.025) (-2.157,0.975) (-2.134,0.925) (-2.109,0.875) (-2.082,0.825) (-2.051,0.775) (-2.018,0.725) (-1.982,0.675) (-1.950,0.634) (-1.900,0.575) (-1.853,0.525) (-1.801,0.475) (-1.750,0.430) (-1.700,0.390) (-1.650,0.353) (-1.600,0.319) (-1.550,0.288) (-1.500,0.259) (-1.450,0.233) (-1.400,0.208) (-1.350,0.186) (-1.300,0.165)}
\EQ{(1.300,0.191) (1.350,0.215) (1.400,0.241) (1.450,0.270) (1.499,0.300) (1.550,0.335) (1.600,0.372) (1.650,0.414) (1.691,0.450) (1.742,0.500) (1.788,0.550) (1.825,0.594) (1.868,0.650) (1.900,0.696) (1.934,0.750) (1.963,0.800) (1.989,0.850) (2.013,0.900) (2.034,0.950) (2.053,1.000) (2.070,1.050) (2.085,1.100) (2.098,1.150) (2.109,1.200) (2.118,1.250) (2.126,1.300) (2.132,1.350) (2.137,1.400) (2.139,1.450) (2.141,1.500) (2.140,1.550) (2.138,1.600) (2.135,1.650) (2.130,1.700) (2.123,1.750) (2.115,1.800) (2.106,1.850) (2.095,1.900) (2.082,1.950) (2.068,2.000) (2.052,2.050) (2.035,2.100) (2.016,2.150) (1.995,2.200) (1.972,2.250) (1.950,2.296) (1.925,2.344) (1.900,2.389) (1.864,2.450) (1.831,2.500) (1.800,2.546) (1.760,2.600) (1.725,2.645) (1.679,2.700) (1.634,2.750) (1.600,2.786) (1.562,2.825) (1.525,2.860) (1.481,2.900) (1.425,2.947) (1.375,2.987) (1.325,3.025) (1.275,3.060) (1.225,3.093) (1.175,3.124) (1.132,3.150) (1.087,3.175) (1.040,3.200) (0.990,3.225) (0.936,3.250) (0.879,3.275) (0.825,3.297) (0.775,3.315) (0.725,3.333) (0.675,3.349) (0.625,3.364) (0.575,3.377) (0.525,3.390) (0.475,3.401) (0.425,3.411) (0.375,3.420) (0.325,3.427) (0.275,3.434) (0.225,3.439) (0.175,3.444) (0.125,3.447) (0.075,3.449) (0.025,3.450) (-0.037,3.450) (-0.100,3.448) (-0.150,3.445) (-0.200,3.442) (-0.250,3.437) (-0.300,3.431) (-0.350,3.424) (-0.400,3.415) (-0.450,3.406) (-0.500,3.395) (-0.550,3.384) (-0.600,3.371) (-0.650,3.356) (-0.700,3.341) (-0.748,3.325) (-0.800,3.306) (-0.850,3.287) (-0.900,3.266) (-0.950,3.244) (-1.000,3.220) (-1.050,3.195) (-1.100,3.168) (-1.150,3.139) (-1.200,3.109) (-1.250,3.077) (-1.300,3.043) (-1.350,3.006) (-1.391,2.975) (-1.450,2.927) (-1.500,2.883) (-1.550,2.836) (-1.587,2.800) (-1.625,2.760) (-1.675,2.705) (-1.721,2.650) (-1.760,2.600) (-1.797,2.550) (-1.831,2.500) (-1.864,2.450) (-1.894,2.400) (-1.922,2.350) (-1.948,2.300) (-1.972,2.250) (-1.995,2.200) (-2.016,2.150) (-2.035,2.100) (-2.052,2.050) (-2.068,2.000) (-2.082,1.950) (-2.095,1.900) (-2.106,1.850) (-2.115,1.800) (-2.123,1.750) (-2.130,1.700) (-2.135,1.650) (-2.138,1.600) (-2.140,1.550) (-2.141,1.500) (-2.139,1.450) (-2.137,1.400) (-2.132,1.350) (-2.126,1.300) (-2.118,1.250) (-2.109,1.200) (-2.098,1.150) (-2.085,1.100) (-2.070,1.050) (-2.053,1.000) (-2.034,0.950) (-2.013,0.900) (-1.989,0.850) (-1.963,0.800) (-1.934,0.750) (-1.903,0.700) (-1.868,0.650) (-1.830,0.600) (-1.788,0.550) (-1.750,0.508) (-1.700,0.459) (-1.650,0.414) (-1.603,0.375) (-1.550,0.335) (-1.500,0.301) (-1.450,0.270) (-1.400,0.241) (-1.350,0.215) (-1.300,0.191)}
\EQ{(1.300,0.216) (1.350,0.244) (1.400,0.274) (1.450,0.308) (1.500,0.345) (1.550,0.385) (1.595,0.425) (1.646,0.475) (1.692,0.525) (1.725,0.565) (1.770,0.625) (1.800,0.669) (1.834,0.725) (1.862,0.775) (1.887,0.825) (1.909,0.875) (1.929,0.925) (1.947,0.975) (1.962,1.025) (1.975,1.075) (1.987,1.125) (1.997,1.175) (2.004,1.225) (2.011,1.275) (2.015,1.325) (2.017,1.375) (2.018,1.425) (2.017,1.475) (2.015,1.525) (2.011,1.575) (2.005,1.625) (1.998,1.675) (1.989,1.725) (1.979,1.775) (1.967,1.825) (1.953,1.875) (1.937,1.925) (1.920,1.975) (1.901,2.025) (1.880,2.075) (1.858,2.125) (1.833,2.175) (1.807,2.225) (1.778,2.275) (1.750,2.321) (1.714,2.375) (1.679,2.425) (1.641,2.475) (1.601,2.525) (1.557,2.575) (1.510,2.625) (1.475,2.661) (1.434,2.700) (1.378,2.750) (1.325,2.794) (1.285,2.825) (1.225,2.869) (1.180,2.900) (1.125,2.935) (1.075,2.965) (1.025,2.993) (0.975,3.020) (0.925,3.045) (0.875,3.068) (0.825,3.089) (0.775,3.109) (0.725,3.128) (0.675,3.145) (0.625,3.161) (0.576,3.175) (0.525,3.188) (0.476,3.200) (0.425,3.211) (0.375,3.220) (0.325,3.229) (0.275,3.236) (0.225,3.241) (0.175,3.246) (0.125,3.249) (0.075,3.252) (0.025,3.253) (-0.050,3.252) (-0.100,3.251) (-0.150,3.248) (-0.200,3.244) (-0.250,3.239) (-0.300,3.232) (-0.350,3.225) (-0.400,3.216) (-0.450,3.206) (-0.500,3.195) (-0.550,3.182) (-0.600,3.168) (-0.650,3.153) (-0.700,3.137) (-0.750,3.119) (-0.798,3.100) (-0.850,3.079) (-0.900,3.056) (-0.950,3.032) (-1.000,3.007) (-1.050,2.980) (-1.100,2.950) (-1.150,2.919) (-1.200,2.886) (-1.250,2.851) (-1.300,2.814) (-1.348,2.775) (-1.400,2.731) (-1.450,2.685) (-1.486,2.650) (-1.525,2.610) (-1.575,2.555) (-1.621,2.500) (-1.661,2.450) (-1.697,2.400) (-1.731,2.350) (-1.763,2.300) (-1.793,2.250) (-1.820,2.200) (-1.846,2.150) (-1.869,2.100) (-1.891,2.050) (-1.911,2.000) (-1.929,1.950) (-1.945,1.900) (-1.960,1.850) (-1.973,1.800) (-1.984,1.750) (-1.994,1.700) (-2.002,1.650) (-2.008,1.600) (-2.013,1.550) (-2.017,1.500) (-2.018,1.450) (-2.018,1.400) (-2.016,1.350) (-2.013,1.300) (-2.008,1.250) (-2.001,1.200) (-1.992,1.150) (-1.982,1.100) (-1.969,1.050) (-1.955,1.000) (-1.938,0.950) (-1.919,0.900) (-1.898,0.850) (-1.875,0.801) (-1.850,0.753) (-1.825,0.710) (-1.787,0.650) (-1.752,0.600) (-1.713,0.550) (-1.675,0.506) (-1.625,0.454) (-1.575,0.407) (-1.525,0.364) (-1.475,0.326) (-1.425,0.291) (-1.375,0.259) (-1.325,0.230) (-1.300,0.216)}
\EQ{(1.300,0.243) (1.350,0.275) (1.400,0.310) (1.450,0.348) (1.500,0.391) (1.550,0.439) (1.600,0.493) (1.647,0.550) (1.684,0.600) (1.717,0.650) (1.747,0.700) (1.774,0.750) (1.798,0.800) (1.819,0.850) (1.838,0.900) (1.855,0.950) (1.869,1.000) (1.881,1.050) (1.892,1.100) (1.900,1.150) (1.906,1.200) (1.911,1.250) (1.914,1.300) (1.915,1.350) (1.914,1.400) (1.911,1.450) (1.907,1.500) (1.901,1.550) (1.894,1.600) (1.885,1.650) (1.873,1.700) (1.861,1.750) (1.846,1.800) (1.830,1.850) (1.812,1.900) (1.792,1.950) (1.770,2.000) (1.746,2.050) (1.720,2.100) (1.692,2.150) (1.661,2.200) (1.629,2.250) (1.593,2.300) (1.555,2.350) (1.515,2.400) (1.475,2.445) (1.425,2.498) (1.375,2.548) (1.325,2.593) (1.275,2.636) (1.226,2.675) (1.175,2.713) (1.125,2.747) (1.083,2.775) (1.025,2.810) (0.975,2.838) (0.925,2.865) (0.875,2.890) (0.825,2.913) (0.775,2.934) (0.725,2.954) (0.675,2.972) (0.625,2.989) (0.575,3.004) (0.525,3.018) (0.475,3.030) (0.425,3.042) (0.375,3.052) (0.325,3.060) (0.275,3.068) (0.225,3.074) (0.175,3.079) (0.125,3.082) (0.075,3.085) (0.025,3.086) (-0.050,3.086) (-0.100,3.084) (-0.150,3.081) (-0.200,3.076) (-0.250,3.071) (-0.300,3.064) (-0.350,3.056) (-0.400,3.047) (-0.450,3.036) (-0.500,3.024) (-0.550,3.011) (-0.600,2.996) (-0.650,2.980) (-0.700,2.963) (-0.750,2.944) (-0.800,2.923) (-0.850,2.901) (-0.900,2.877) (-0.950,2.852) (-0.999,2.825) (-1.050,2.795) (-1.100,2.764) (-1.150,2.730) (-1.192,2.700) (-1.250,2.656) (-1.300,2.615) (-1.345,2.575) (-1.398,2.525) (-1.448,2.475) (-1.493,2.425) (-1.535,2.375) (-1.575,2.325) (-1.611,2.275) (-1.645,2.225) (-1.675,2.178) (-1.706,2.125) (-1.733,2.075) (-1.758,2.025) (-1.781,1.975) (-1.802,1.925) (-1.821,1.875) (-1.838,1.825) (-1.854,1.775) (-1.867,1.725) (-1.879,1.675) (-1.889,1.625) (-1.898,1.575) (-1.905,1.525) (-1.910,1.475) (-1.913,1.425) (-1.915,1.375) (-1.914,1.325) (-1.913,1.275) (-1.909,1.225) (-1.903,1.175) (-1.896,1.125) (-1.887,1.075) (-1.875,1.025) (-1.862,0.975) (-1.847,0.925) (-1.829,0.875) (-1.809,0.825) (-1.786,0.775) (-1.761,0.725) (-1.733,0.675) (-1.701,0.625) (-1.666,0.575) (-1.627,0.525) (-1.584,0.475) (-1.536,0.425) (-1.482,0.375) (-1.425,0.328) (-1.375,0.292) (-1.325,0.258) (-1.300,0.243)}
\EQ{(1.288,0.262) (1.342,0.300) (1.400,0.347) (1.450,0.391) (1.500,0.442) (1.550,0.499) (1.590,0.550) (1.625,0.600) (1.656,0.650) (1.684,0.700) (1.708,0.750) (1.730,0.800) (1.750,0.850) (1.767,0.900) (1.781,0.950) (1.794,1.000) (1.804,1.050) (1.812,1.100) (1.818,1.150) (1.823,1.200) (1.825,1.250) (1.825,1.300) (1.824,1.350) (1.821,1.400) (1.816,1.450) (1.810,1.500) (1.801,1.550) (1.791,1.600) (1.779,1.650) (1.765,1.700) (1.750,1.748) (1.732,1.800) (1.712,1.850) (1.691,1.900) (1.667,1.950) (1.641,2.000) (1.613,2.050) (1.583,2.100) (1.550,2.149) (1.514,2.200) (1.476,2.250) (1.434,2.300) (1.400,2.339) (1.366,2.375) (1.325,2.416) (1.289,2.450) (1.250,2.484) (1.201,2.525) (1.150,2.565) (1.101,2.600) (1.050,2.634) (1.000,2.666) (0.950,2.695) (0.900,2.722) (0.850,2.748) (0.800,2.771) (0.750,2.793) (0.700,2.813) (0.650,2.831) (0.600,2.848) (0.550,2.864) (0.500,2.878) (0.450,2.890) (0.400,2.901) (0.350,2.911) (0.300,2.919) (0.250,2.926) (0.200,2.932) (0.150,2.937) (0.100,2.940) (0.050,2.942) (0.000,2.943) (-0.050,2.942) (-0.100,2.940) (-0.150,2.937) (-0.200,2.932) (-0.250,2.926) (-0.300,2.919) (-0.350,2.911) (-0.400,2.901) (-0.450,2.890) (-0.500,2.878) (-0.550,2.864) (-0.600,2.848) (-0.650,2.831) (-0.700,2.813) (-0.750,2.793) (-0.800,2.771) (-0.850,2.748) (-0.900,2.722) (-0.950,2.695) (-1.000,2.666) (-1.050,2.634) (-1.100,2.601) (-1.150,2.565) (-1.200,2.526) (-1.250,2.484) (-1.289,2.450) (-1.325,2.416) (-1.366,2.375) (-1.400,2.339) (-1.434,2.300) (-1.475,2.251) (-1.514,2.200) (-1.550,2.150) (-1.583,2.100) (-1.613,2.050) (-1.641,2.000) (-1.667,1.950) (-1.691,1.900) (-1.712,1.850) (-1.732,1.800) (-1.749,1.750) (-1.765,1.700) (-1.779,1.650) (-1.791,1.600) (-1.801,1.550) (-1.810,1.500) (-1.816,1.450) (-1.821,1.400) (-1.824,1.350) (-1.825,1.300) (-1.825,1.250) (-1.823,1.200) (-1.818,1.150) (-1.812,1.100) (-1.804,1.050) (-1.794,1.000) (-1.781,0.950) (-1.767,0.900) (-1.750,0.851) (-1.730,0.800) (-1.708,0.750) (-1.684,0.700) (-1.656,0.650) (-1.625,0.601) (-1.590,0.550) (-1.551,0.500) (-1.508,0.450) (-1.459,0.400) (-1.404,0.350) (-1.350,0.307) (-1.307,0.275) (-1.288,0.262)}
\EQ{(1.275,0.280) (1.325,0.319) (1.375,0.362) (1.414,0.400) (1.450,0.438) (1.500,0.497) (1.539,0.550) (1.572,0.600) (1.600,0.648) (1.625,0.696) (1.650,0.750) (1.670,0.800) (1.688,0.850) (1.703,0.900) (1.716,0.950) (1.726,1.000) (1.734,1.050) (1.741,1.100) (1.745,1.150) (1.747,1.200) (1.748,1.250) (1.746,1.300) (1.743,1.350) (1.738,1.400) (1.731,1.450) (1.722,1.500) (1.711,1.550) (1.698,1.600) (1.683,1.650) (1.667,1.700) (1.648,1.750) (1.627,1.800) (1.604,1.850) (1.579,1.900) (1.552,1.950) (1.525,1.995) (1.489,2.050) (1.454,2.100) (1.416,2.150) (1.375,2.200) (1.330,2.250) (1.282,2.300) (1.229,2.350) (1.175,2.397) (1.125,2.437) (1.075,2.475) (1.025,2.509) (0.975,2.542) (0.925,2.571) (0.875,2.599) (0.825,2.625) (0.775,2.649) (0.725,2.671) (0.675,2.691) (0.625,2.710) (0.575,2.727) (0.525,2.742) (0.475,2.756) (0.425,2.768) (0.375,2.779) (0.325,2.789) (0.275,2.797) (0.225,2.804) (0.175,2.809) (0.125,2.813) (0.075,2.816) (0.025,2.817) (-0.050,2.817) (-0.100,2.815) (-0.150,2.811) (-0.200,2.807) (-0.250,2.800) (-0.300,2.793) (-0.350,2.784) (-0.400,2.774) (-0.450,2.762) (-0.500,2.749) (-0.550,2.735) (-0.600,2.719) (-0.650,2.701) (-0.700,2.681) (-0.750,2.660) (-0.800,2.637) (-0.850,2.612) (-0.900,2.586) (-0.950,2.557) (-1.000,2.526) (-1.050,2.492) (-1.100,2.456) (-1.141,2.425) (-1.200,2.376) (-1.250,2.331) (-1.300,2.282) (-1.350,2.228) (-1.396,2.175) (-1.436,2.125) (-1.472,2.075) (-1.506,2.025) (-1.537,1.975) (-1.566,1.925) (-1.592,1.875) (-1.616,1.825) (-1.638,1.775) (-1.658,1.725) (-1.675,1.675) (-1.691,1.625) (-1.705,1.575) (-1.716,1.525) (-1.726,1.475) (-1.734,1.425) (-1.741,1.375) (-1.745,1.325) (-1.747,1.275) (-1.748,1.225) (-1.746,1.175) (-1.743,1.125) (-1.738,1.075) (-1.730,1.025) (-1.721,0.975) (-1.709,0.925) (-1.695,0.875) (-1.679,0.825) (-1.660,0.775) (-1.639,0.725) (-1.614,0.675) (-1.587,0.625) (-1.556,0.575) (-1.525,0.530) (-1.482,0.475) (-1.438,0.425) (-1.400,0.386) (-1.362,0.350) (-1.301,0.300) (-1.275,0.280)}
\EQ{(1.275,0.308) (1.323,0.350) (1.374,0.400) (1.419,0.450) (1.450,0.489) (1.493,0.550) (1.524,0.600) (1.550,0.647) (1.575,0.698) (1.597,0.750) (1.616,0.800) (1.632,0.850) (1.645,0.900) (1.656,0.950) (1.665,1.000) (1.672,1.050) (1.676,1.100) (1.679,1.150) (1.679,1.200) (1.678,1.250) (1.674,1.300) (1.669,1.350) (1.662,1.400) (1.652,1.450) (1.641,1.500) (1.628,1.550) (1.613,1.600) (1.595,1.650) (1.576,1.700) (1.554,1.750) (1.530,1.800) (1.504,1.850) (1.475,1.900) (1.444,1.950) (1.410,2.000) (1.375,2.047) (1.333,2.100) (1.289,2.150) (1.250,2.191) (1.200,2.240) (1.162,2.275) (1.102,2.325) (1.050,2.365) (1.001,2.400) (0.950,2.434) (0.900,2.464) (0.850,2.492) (0.800,2.518) (0.750,2.543) (0.700,2.565) (0.650,2.585) (0.600,2.604) (0.550,2.621) (0.500,2.636) (0.450,2.650) (0.400,2.662) (0.350,2.672) (0.300,2.682) (0.250,2.689) (0.200,2.696) (0.150,2.700) (0.100,2.704) (0.050,2.706) (0.000,2.707) (-0.050,2.706) (-0.100,2.704) (-0.150,2.700) (-0.200,2.696) (-0.250,2.689) (-0.300,2.682) (-0.350,2.672) (-0.400,2.662) (-0.450,2.650) (-0.500,2.636) (-0.550,2.621) (-0.600,2.604) (-0.650,2.585) (-0.700,2.565) (-0.750,2.543) (-0.800,2.518) (-0.850,2.492) (-0.900,2.464) (-0.950,2.434) (-1.000,2.401) (-1.050,2.365) (-1.100,2.327) (-1.150,2.285) (-1.189,2.250) (-1.241,2.200) (-1.289,2.150) (-1.325,2.109) (-1.373,2.050) (-1.410,2.000) (-1.444,1.950) (-1.475,1.900) (-1.504,1.850) (-1.530,1.800) (-1.554,1.750) (-1.575,1.702) (-1.595,1.650) (-1.613,1.600) (-1.628,1.550) (-1.641,1.500) (-1.652,1.450) (-1.662,1.400) (-1.669,1.350) (-1.674,1.300) (-1.678,1.250) (-1.679,1.200) (-1.679,1.150) (-1.676,1.100) (-1.672,1.050) (-1.665,1.000) (-1.656,0.950) (-1.645,0.900) (-1.632,0.850) (-1.616,0.800) (-1.597,0.750) (-1.576,0.700) (-1.552,0.650) (-1.525,0.601) (-1.493,0.550) (-1.458,0.500) (-1.425,0.458) (-1.375,0.401) (-1.325,0.352) (-1.275,0.308)}
\EQ{(1.275,0.337) (1.314,0.375) (1.350,0.414) (1.400,0.475) (1.436,0.525) (1.467,0.575) (1.495,0.625) (1.519,0.675) (1.540,0.725) (1.559,0.775) (1.574,0.825) (1.587,0.875) (1.598,0.925) (1.606,0.975) (1.612,1.025) (1.616,1.075) (1.618,1.125) (1.618,1.175) (1.616,1.225) (1.611,1.275) (1.605,1.325) (1.597,1.375) (1.586,1.425) (1.574,1.475) (1.559,1.525) (1.543,1.575) (1.524,1.625) (1.503,1.675) (1.479,1.725) (1.454,1.775) (1.425,1.825) (1.394,1.875) (1.360,1.925) (1.325,1.973) (1.283,2.025) (1.250,2.063) (1.216,2.100) (1.175,2.141) (1.125,2.187) (1.081,2.225) (1.025,2.269) (0.983,2.300) (0.925,2.339) (0.875,2.370) (0.825,2.398) (0.775,2.425) (0.725,2.449) (0.675,2.471) (0.625,2.492) (0.575,2.510) (0.525,2.527) (0.475,2.542) (0.425,2.555) (0.375,2.567) (0.325,2.578) (0.275,2.586) (0.225,2.594) (0.175,2.599) (0.125,2.604) (0.075,2.607) (0.025,2.608) (-0.050,2.608) (-0.100,2.605) (-0.150,2.602) (-0.200,2.597) (-0.250,2.590) (-0.300,2.582) (-0.350,2.573) (-0.400,2.562) (-0.450,2.549) (-0.500,2.535) (-0.550,2.519) (-0.600,2.501) (-0.650,2.482) (-0.700,2.460) (-0.750,2.437) (-0.800,2.412) (-0.850,2.384) (-0.900,2.355) (-0.946,2.325) (-1.000,2.288) (-1.050,2.250) (-1.100,2.209) (-1.139,2.175) (-1.191,2.125) (-1.239,2.075) (-1.275,2.034) (-1.323,1.975) (-1.360,1.925) (-1.394,1.875) (-1.425,1.825) (-1.450,1.781) (-1.475,1.734) (-1.500,1.681) (-1.524,1.625) (-1.543,1.575) (-1.559,1.525) (-1.574,1.475) (-1.586,1.425) (-1.597,1.375) (-1.605,1.325) (-1.611,1.275) (-1.616,1.225) (-1.618,1.175) (-1.618,1.125) (-1.616,1.075) (-1.612,1.025) (-1.606,0.975) (-1.598,0.925) (-1.587,0.875) (-1.574,0.825) (-1.559,0.775) (-1.540,0.725) (-1.519,0.675) (-1.495,0.625) (-1.467,0.575) (-1.436,0.525) (-1.400,0.475) (-1.360,0.425) (-1.325,0.387) (-1.289,0.350) (-1.275,0.337)}
\EQ{(1.266,0.359) (1.305,0.400) (1.346,0.450) (1.382,0.500) (1.414,0.550) (1.442,0.600) (1.467,0.650) (1.488,0.700) (1.506,0.750) (1.522,0.800) (1.534,0.850) (1.545,0.900) (1.553,0.950) (1.559,1.000) (1.562,1.050) (1.563,1.100) (1.563,1.150) (1.560,1.200) (1.555,1.250) (1.547,1.300) (1.538,1.350) (1.527,1.400) (1.513,1.450) (1.498,1.500) (1.480,1.550) (1.460,1.600) (1.437,1.650) (1.412,1.700) (1.385,1.750) (1.354,1.800) (1.325,1.844) (1.285,1.900) (1.250,1.944) (1.201,2.000) (1.154,2.050) (1.101,2.100) (1.050,2.145) (1.000,2.185) (0.950,2.222) (0.908,2.250) (0.850,2.287) (0.800,2.315) (0.750,2.342) (0.700,2.366) (0.650,2.388) (0.600,2.409) (0.550,2.427) (0.500,2.443) (0.450,2.458) (0.400,2.471) (0.350,2.483) (0.300,2.493) (0.250,2.501) (0.200,2.508) (0.150,2.513) (0.100,2.517) (0.050,2.519) (0.000,2.520) (-0.050,2.519) (-0.100,2.517) (-0.150,2.513) (-0.200,2.508) (-0.250,2.501) (-0.300,2.493) (-0.350,2.483) (-0.400,2.471) (-0.450,2.458) (-0.500,2.443) (-0.550,2.427) (-0.600,2.409) (-0.650,2.388) (-0.700,2.366) (-0.750,2.342) (-0.800,2.315) (-0.850,2.287) (-0.900,2.256) (-0.945,2.225) (-1.000,2.185) (-1.044,2.150) (-1.100,2.101) (-1.150,2.054) (-1.200,2.001) (-1.245,1.950) (-1.285,1.900) (-1.321,1.850) (-1.350,1.807) (-1.385,1.750) (-1.412,1.700) (-1.437,1.650) (-1.460,1.600) (-1.480,1.550) (-1.498,1.500) (-1.513,1.450) (-1.527,1.400) (-1.538,1.350) (-1.547,1.300) (-1.555,1.250) (-1.560,1.200) (-1.563,1.150) (-1.563,1.100) (-1.562,1.050) (-1.559,1.000) (-1.553,0.950) (-1.545,0.900) (-1.534,0.850) (-1.522,0.800) (-1.506,0.750) (-1.488,0.700) (-1.467,0.650) (-1.442,0.600) (-1.414,0.550) (-1.382,0.500) (-1.350,0.456) (-1.305,0.400) (-1.266,0.359)}
\EQ{(1.251,0.374) (1.295,0.425) (1.325,0.465) (1.365,0.525) (1.393,0.575) (1.418,0.625) (1.440,0.675) (1.458,0.725) (1.473,0.775) (1.486,0.825) (1.496,0.875) (1.504,0.925) (1.510,0.975) (1.513,1.025) (1.514,1.075) (1.512,1.125) (1.509,1.175) (1.503,1.225) (1.495,1.275) (1.486,1.325) (1.473,1.375) (1.459,1.425) (1.442,1.475) (1.423,1.525) (1.402,1.575) (1.378,1.625) (1.352,1.675) (1.325,1.721) (1.290,1.775) (1.255,1.825) (1.216,1.875) (1.175,1.923) (1.127,1.975) (1.075,2.025) (1.025,2.070) (0.975,2.110) (0.925,2.147) (0.875,2.182) (0.825,2.213) (0.775,2.242) (0.725,2.268) (0.675,2.292) (0.625,2.314) (0.575,2.334) (0.525,2.352) (0.475,2.369) (0.425,2.383) (0.375,2.396) (0.325,2.407) (0.275,2.416) (0.225,2.424) (0.175,2.430) (0.125,2.435) (0.075,2.438) (0.025,2.440) (-0.050,2.439) (-0.100,2.437) (-0.150,2.433) (-0.200,2.427) (-0.250,2.420) (-0.300,2.412) (-0.350,2.402) (-0.400,2.390) (-0.450,2.376) (-0.500,2.361) (-0.550,2.344) (-0.599,2.325) (-0.650,2.304) (-0.700,2.280) (-0.750,2.255) (-0.800,2.228) (-0.846,2.200) (-0.900,2.165) (-0.950,2.129) (-1.000,2.090) (-1.048,2.050) (-1.100,2.002) (-1.150,1.951) (-1.195,1.900) (-1.236,1.850) (-1.273,1.800) (-1.307,1.750) (-1.337,1.700) (-1.365,1.650) (-1.390,1.600) (-1.413,1.550) (-1.433,1.500) (-1.450,1.453) (-1.466,1.400) (-1.480,1.350) (-1.491,1.300) (-1.500,1.250) (-1.506,1.200) (-1.511,1.150) (-1.513,1.100) (-1.514,1.050) (-1.512,1.000) (-1.507,0.950) (-1.501,0.900) (-1.492,0.850) (-1.480,0.800) (-1.466,0.750) (-1.449,0.700) (-1.429,0.650) (-1.406,0.600) (-1.379,0.550) (-1.350,0.502) (-1.314,0.450) (-1.275,0.401) (-1.251,0.374)}
\EQ{(1.250,0.405) (1.285,0.450) (1.318,0.500) (1.347,0.550) (1.373,0.600) (1.395,0.650) (1.413,0.700) (1.429,0.750) (1.442,0.800) (1.452,0.850) (1.460,0.900) (1.465,0.950) (1.468,1.000) (1.468,1.050) (1.467,1.100) (1.463,1.150) (1.457,1.200) (1.448,1.250) (1.438,1.300) (1.425,1.350) (1.410,1.400) (1.392,1.450) (1.372,1.500) (1.350,1.550) (1.325,1.600) (1.300,1.645) (1.266,1.700) (1.232,1.750) (1.200,1.793) (1.154,1.850) (1.109,1.900) (1.059,1.950) (1.004,2.000) (0.950,2.044) (0.909,2.075) (0.850,2.116) (0.800,2.147) (0.751,2.175) (0.704,2.200) (0.652,2.225) (0.600,2.248) (0.550,2.268) (0.500,2.285) (0.450,2.301) (0.400,2.315) (0.350,2.328) (0.300,2.338) (0.250,2.347) (0.200,2.354) (0.150,2.360) (0.100,2.364) (0.050,2.366) (0.000,2.367) (-0.050,2.366) (-0.100,2.364) (-0.150,2.360) (-0.200,2.354) (-0.250,2.347) (-0.300,2.338) (-0.350,2.328) (-0.400,2.315) (-0.450,2.301) (-0.500,2.285) (-0.550,2.268) (-0.600,2.248) (-0.650,2.226) (-0.700,2.202) (-0.750,2.176) (-0.795,2.150) (-0.850,2.116) (-0.900,2.081) (-0.942,2.050) (-1.000,2.003) (-1.050,1.958) (-1.100,1.909) (-1.150,1.854) (-1.195,1.800) (-1.225,1.760) (-1.266,1.700) (-1.297,1.650) (-1.325,1.600) (-1.350,1.550) (-1.372,1.500) (-1.392,1.450) (-1.410,1.400) (-1.425,1.350) (-1.438,1.300) (-1.448,1.250) (-1.457,1.200) (-1.463,1.150) (-1.467,1.100) (-1.468,1.050) (-1.468,1.000) (-1.465,0.950) (-1.460,0.900) (-1.452,0.850) (-1.442,0.800) (-1.429,0.750) (-1.413,0.700) (-1.395,0.650) (-1.373,0.600) (-1.350,0.555) (-1.325,0.512) (-1.285,0.450) (-1.250,0.405)}
\EQ{(1.244,0.431) (1.274,0.475) (1.300,0.518) (1.325,0.565) (1.350,0.619) (1.372,0.675) (1.388,0.725) (1.400,0.775) (1.411,0.825) (1.419,0.875) (1.424,0.925) (1.427,0.975) (1.427,1.025) (1.425,1.075) (1.421,1.125) (1.414,1.175) (1.406,1.225) (1.395,1.275) (1.381,1.325) (1.365,1.375) (1.347,1.425) (1.326,1.475) (1.303,1.525) (1.277,1.575) (1.250,1.621) (1.215,1.675) (1.180,1.725) (1.141,1.775) (1.100,1.822) (1.050,1.874) (1.000,1.922) (0.950,1.965) (0.906,2.000) (0.850,2.040) (0.800,2.073) (0.755,2.100) (0.709,2.125) (0.660,2.150) (0.605,2.175) (0.550,2.198) (0.500,2.216) (0.450,2.233) (0.400,2.247) (0.350,2.260) (0.300,2.271) (0.250,2.280) (0.200,2.287) (0.150,2.293) (0.100,2.297) (0.050,2.299) (0.000,2.300) (-0.050,2.299) (-0.100,2.297) (-0.150,2.293) (-0.200,2.287) (-0.250,2.280) (-0.300,2.271) (-0.350,2.260) (-0.400,2.247) (-0.450,2.233) (-0.500,2.216) (-0.550,2.198) (-0.600,2.177) (-0.650,2.155) (-0.700,2.130) (-0.750,2.103) (-0.796,2.075) (-0.850,2.040) (-0.900,2.004) (-0.950,1.965) (-0.997,1.925) (-1.049,1.875) (-1.097,1.825) (-1.141,1.775) (-1.175,1.731) (-1.215,1.675) (-1.248,1.625) (-1.275,1.578) (-1.300,1.531) (-1.325,1.478) (-1.347,1.425) (-1.365,1.375) (-1.381,1.325) (-1.395,1.275) (-1.406,1.225) (-1.414,1.175) (-1.421,1.125) (-1.425,1.075) (-1.427,1.025) (-1.427,0.975) (-1.424,0.925) (-1.419,0.875) (-1.411,0.825) (-1.400,0.775) (-1.388,0.725) (-1.372,0.675) (-1.352,0.625) (-1.330,0.575) (-1.304,0.525) (-1.275,0.477) (-1.244,0.431)}
\EQ{(1.229,0.446) (1.263,0.500) (1.290,0.550) (1.313,0.600) (1.333,0.650) (1.349,0.700) (1.362,0.750) (1.373,0.800) (1.381,0.850) (1.386,0.900) (1.389,0.950) (1.389,1.000) (1.387,1.050) (1.383,1.100) (1.376,1.150) (1.367,1.200) (1.355,1.250) (1.341,1.300) (1.325,1.349) (1.306,1.400) (1.284,1.450) (1.260,1.500) (1.233,1.550) (1.203,1.600) (1.175,1.642) (1.132,1.700) (1.100,1.740) (1.050,1.796) (1.000,1.846) (0.950,1.891) (0.909,1.925) (0.850,1.970) (0.806,2.000) (0.750,2.035) (0.700,2.063) (0.650,2.089) (0.600,2.112) (0.550,2.133) (0.500,2.152) (0.450,2.169) (0.400,2.184) (0.350,2.197) (0.300,2.208) (0.250,2.218) (0.200,2.225) (0.150,2.231) (0.100,2.236) (0.050,2.238) (0.000,2.239) (-0.050,2.238) (-0.100,2.236) (-0.150,2.231) (-0.200,2.225) (-0.250,2.218) (-0.300,2.208) (-0.350,2.197) (-0.400,2.184) (-0.450,2.169) (-0.500,2.152) (-0.550,2.133) (-0.600,2.112) (-0.650,2.089) (-0.700,2.063) (-0.750,2.035) (-0.800,2.004) (-0.843,1.975) (-0.900,1.932) (-0.940,1.900) (-0.996,1.850) (-1.046,1.800) (-1.091,1.750) (-1.125,1.709) (-1.169,1.650) (-1.200,1.604) (-1.233,1.550) (-1.260,1.500) (-1.284,1.450) (-1.306,1.400) (-1.325,1.350) (-1.341,1.300) (-1.355,1.250) (-1.367,1.200) (-1.376,1.150) (-1.383,1.100) (-1.387,1.050) (-1.389,1.000) (-1.389,0.950) (-1.386,0.900) (-1.381,0.850) (-1.373,0.800) (-1.362,0.750) (-1.349,0.700) (-1.333,0.650) (-1.313,0.600) (-1.290,0.550) (-1.263,0.500) (-1.232,0.450) (-1.229,0.446)}
\EQ{(1.224,0.476) (1.250,0.522) (1.275,0.574) (1.296,0.625) (1.313,0.675) (1.326,0.725) (1.337,0.775) (1.345,0.825) (1.350,0.875) (1.353,0.925) (1.354,0.975) (1.352,1.025) (1.347,1.075) (1.340,1.125) (1.331,1.175) (1.319,1.225) (1.305,1.275) (1.288,1.325) (1.269,1.375) (1.246,1.425) (1.221,1.475) (1.194,1.525) (1.162,1.575) (1.128,1.625) (1.089,1.675) (1.050,1.721) (1.000,1.774) (0.950,1.822) (0.900,1.865) (0.856,1.900) (0.800,1.940) (0.750,1.972) (0.702,2.000) (0.656,2.025) (0.604,2.050) (0.550,2.074) (0.500,2.093) (0.450,2.111) (0.400,2.126) (0.350,2.139) (0.300,2.151) (0.250,2.161) (0.200,2.168) (0.150,2.174) (0.100,2.179) (0.050,2.181) (0.000,2.182) (-0.050,2.181) (-0.100,2.179) (-0.150,2.174) (-0.200,2.168) (-0.250,2.161) (-0.300,2.151) (-0.350,2.139) (-0.400,2.126) (-0.450,2.111) (-0.500,2.093) (-0.547,2.075) (-0.600,2.052) (-0.650,2.028) (-0.700,2.001) (-0.745,1.975) (-0.800,1.940) (-0.850,1.904) (-0.900,1.865) (-0.947,1.825) (-0.999,1.775) (-1.047,1.725) (-1.089,1.675) (-1.125,1.629) (-1.162,1.575) (-1.194,1.525) (-1.221,1.475) (-1.246,1.425) (-1.269,1.375) (-1.288,1.325) (-1.305,1.275) (-1.319,1.225) (-1.331,1.175) (-1.340,1.125) (-1.347,1.075) (-1.352,1.025) (-1.354,0.975) (-1.353,0.925) (-1.350,0.875) (-1.345,0.825) (-1.337,0.775) (-1.326,0.725) (-1.313,0.675) (-1.296,0.625) (-1.275,0.575) (-1.252,0.525) (-1.225,0.477)}

    \draw[mmaEQ] (-10,0) -- (10,0);
    % Dipolo puntual - líneas de campo
% Generado por scripts/generar_datos.py (no editar a mano).
\LC{(1.082,0.489) (1.149,0.455) (1.212,0.414) (1.269,0.365) (1.319,0.309) (1.361,0.248) (1.394,0.180) (1.417,0.109) (1.428,0.035) (1.428,-0.040) (1.415,-0.114) (1.392,-0.185) (1.358,-0.252) (1.316,-0.314) (1.265,-0.369) (1.207,-0.417) (1.144,-0.458) (1.077,-0.491)}
\LC{(0.453,1.106) (0.495,1.168) (0.539,1.229) (0.583,1.290) (0.629,1.349) (0.676,1.408) (0.724,1.465) (0.773,1.522) (0.823,1.578) (0.873,1.633) (0.925,1.688) (0.978,1.741) (1.031,1.794) (1.085,1.846) (1.141,1.896) (1.197,1.946) (1.253,1.996) (1.311,2.044) (1.369,2.091) (1.428,2.137) (1.487,2.183) (1.548,2.228) (1.609,2.271) (1.670,2.314) (1.733,2.356) (1.795,2.397) (1.859,2.437) (1.923,2.476) (1.988,2.514) (2.053,2.551) (2.118,2.587) (2.185,2.622) (2.251,2.657) (2.319,2.690) (2.386,2.722) (2.454,2.753) (2.523,2.784) (2.592,2.813) (2.661,2.841) (2.731,2.869) (2.802,2.895) (2.872,2.920) (2.943,2.945) (3.014,2.968) (3.086,2.990) (3.158,3.011) (3.230,3.031) (3.303,3.050) (3.376,3.068) (3.449,3.085) (3.522,3.101) (3.596,3.116) (3.669,3.130) (3.743,3.142) (3.817,3.154) (3.892,3.164) (3.966,3.173) (4.041,3.181) (4.115,3.188) (4.190,3.194) (4.265,3.199) (4.340,3.203) (4.415,3.205) (4.490,3.206) (4.565,3.206) (4.640,3.205) (4.715,3.203) (4.790,3.200) (4.864,3.195) (4.939,3.189) (5.014,3.182) (5.088,3.174) (5.163,3.165) (5.237,3.154) (5.311,3.142) (5.385,3.129) (5.459,3.115) (5.532,3.099) (5.605,3.082) (5.678,3.064) (5.750,3.045) (5.822,3.024) (5.894,3.002) (5.965,2.979) (6.036,2.955) (6.107,2.929) (6.177,2.902) (6.246,2.874) (6.315,2.844) (6.384,2.813) (6.451,2.781) (6.518,2.748) (6.585,2.713) (6.651,2.677) (6.716,2.640) (6.780,2.601) (6.844,2.561) (6.906,2.520) (6.968,2.478) (7.029,2.434) (7.089,2.389) (7.148,2.343) (7.207,2.296) (7.264,2.247) (7.320,2.197) (7.374,2.146) (7.428,2.094) (7.481,2.040) (7.532,1.985) (7.582,1.930) (7.631,1.873) (7.679,1.815) (7.725,1.755) (7.769,1.695) (7.813,1.634) (7.854,1.572) (7.894,1.508) (7.933,1.444) (7.970,1.379) (8.006,1.313) (8.039,1.246) (8.071,1.178) (8.102,1.109) (8.130,1.040) (8.157,0.970) (8.182,0.899) (8.205,0.828) (8.226,0.756) (8.246,0.683) (8.263,0.610) (8.278,0.537) (8.292,0.463) (8.303,0.389) (8.313,0.315) (8.320,0.240) (8.326,0.165) (8.329,0.090) (8.331,0.015) (8.330,-0.060) (8.327,-0.134) (8.323,-0.209) (8.316,-0.284) (8.307,-0.359) (8.297,-0.433) (8.284,-0.507) (8.270,-0.580) (8.253,-0.653) (8.234,-0.726) (8.214,-0.798) (8.192,-0.870) (8.168,-0.941) (8.142,-1.011) (8.114,-1.081) (8.084,-1.150) (8.053,-1.218) (8.020,-1.285) (7.985,-1.352) (7.949,-1.417) (7.911,-1.482) (7.871,-1.546) (7.830,-1.608) (7.787,-1.670) (7.743,-1.731) (7.698,-1.790) (7.651,-1.849) (7.603,-1.906) (7.553,-1.963) (7.502,-2.018) (7.450,-2.072) (7.397,-2.124) (7.342,-2.176) (7.287,-2.226) (7.230,-2.276) (7.173,-2.324) (7.114,-2.370) (7.054,-2.416) (6.994,-2.460) (6.932,-2.503) (6.870,-2.544) (6.806,-2.585) (6.742,-2.624) (6.678,-2.662) (6.612,-2.698) (6.546,-2.733) (6.479,-2.767) (6.412,-2.800) (6.343,-2.831) (6.275,-2.862) (6.206,-2.890) (6.136,-2.918) (6.066,-2.944) (5.995,-2.969) (5.924,-2.993) (5.852,-3.015) (5.780,-3.036) (5.708,-3.056) (5.635,-3.075) (5.562,-3.092) (5.489,-3.108) (5.415,-3.123) (5.342,-3.137) (5.268,-3.149) (5.194,-3.160) (5.119,-3.170) (5.045,-3.179) (4.970,-3.187) (4.895,-3.193) (4.821,-3.198) (4.746,-3.202) (4.671,-3.205) (4.596,-3.206) (4.521,-3.206) (4.446,-3.206) (4.371,-3.204) (4.296,-3.201) (4.221,-3.196) (4.146,-3.191) (4.071,-3.184) (3.997,-3.177) (3.922,-3.168) (3.848,-3.158) (3.774,-3.147) (3.700,-3.135) (3.626,-3.122) (3.552,-3.107) (3.479,-3.092) (3.406,-3.075) (3.333,-3.058) (3.260,-3.039) (3.188,-3.020) (3.116,-2.999) (3.044,-2.977) (2.973,-2.954) (2.901,-2.930) (2.831,-2.906) (2.760,-2.880) (2.690,-2.853) (2.621,-2.825) (2.551,-2.796) (2.483,-2.766) (2.414,-2.735) (2.346,-2.703) (2.279,-2.670) (2.212,-2.637) (2.146,-2.602) (2.080,-2.566) (2.014,-2.529) (1.950,-2.492) (1.885,-2.453) (1.822,-2.413) (1.758,-2.373) (1.696,-2.331) (1.634,-2.289) (1.573,-2.246) (1.512,-2.201) (1.452,-2.156) (1.393,-2.110) (1.335,-2.063) (1.277,-2.016) (1.220,-1.967) (1.164,-1.917) (1.108,-1.867) (1.053,-1.815) (1.000,-1.763) (0.947,-1.710) (0.895,-1.656) (0.843,-1.601) (0.793,-1.545) (0.744,-1.489) (0.696,-1.432) (0.648,-1.373) (0.602,-1.314) (0.557,-1.254) (0.513,-1.194) (0.470,-1.132) (0.442,-1.091)}
\LC{(-0.453,1.106) (-0.495,1.168) (-0.539,1.229) (-0.583,1.290) (-0.629,1.349) (-0.676,1.408) (-0.724,1.465) (-0.773,1.522) (-0.823,1.578) (-0.873,1.633) (-0.925,1.688) (-0.978,1.741) (-1.031,1.794) (-1.085,1.846) (-1.141,1.896) (-1.197,1.946) (-1.253,1.996) (-1.311,2.044) (-1.369,2.091) (-1.428,2.137) (-1.487,2.183) (-1.548,2.228) (-1.609,2.271) (-1.670,2.314) (-1.733,2.356) (-1.795,2.397) (-1.859,2.437) (-1.923,2.476) (-1.988,2.514) (-2.053,2.551) (-2.118,2.587) (-2.185,2.622) (-2.251,2.657) (-2.319,2.690) (-2.386,2.722) (-2.454,2.753) (-2.523,2.784) (-2.592,2.813) (-2.661,2.841) (-2.731,2.869) (-2.802,2.895) (-2.872,2.920) (-2.943,2.945) (-3.014,2.968) (-3.086,2.990) (-3.158,3.011) (-3.230,3.031) (-3.303,3.050) (-3.376,3.068) (-3.449,3.085) (-3.522,3.101) (-3.596,3.116) (-3.669,3.130) (-3.743,3.142) (-3.817,3.154) (-3.892,3.164) (-3.966,3.173) (-4.041,3.181) (-4.115,3.188) (-4.190,3.194) (-4.265,3.199) (-4.340,3.203) (-4.415,3.205) (-4.490,3.206) (-4.565,3.206) (-4.640,3.205) (-4.715,3.203) (-4.790,3.200) (-4.864,3.195) (-4.939,3.189) (-5.014,3.182) (-5.088,3.174) (-5.163,3.165) (-5.237,3.154) (-5.311,3.142) (-5.385,3.129) (-5.459,3.115) (-5.532,3.099) (-5.605,3.082) (-5.678,3.064) (-5.750,3.045) (-5.822,3.024) (-5.894,3.002) (-5.965,2.979) (-6.036,2.955) (-6.107,2.929) (-6.177,2.902) (-6.246,2.874) (-6.315,2.844) (-6.384,2.813) (-6.451,2.781) (-6.518,2.748) (-6.585,2.713) (-6.651,2.677) (-6.716,2.640) (-6.780,2.601) (-6.844,2.561) (-6.906,2.520) (-6.968,2.478) (-7.029,2.434) (-7.089,2.389) (-7.148,2.343) (-7.207,2.296) (-7.264,2.247) (-7.320,2.197) (-7.374,2.146) (-7.428,2.094) (-7.481,2.040) (-7.532,1.985) (-7.582,1.930) (-7.631,1.873) (-7.679,1.815) (-7.725,1.755) (-7.769,1.695) (-7.813,1.634) (-7.854,1.572) (-7.894,1.508) (-7.933,1.444) (-7.970,1.379) (-8.006,1.313) (-8.039,1.246) (-8.071,1.178) (-8.102,1.109) (-8.130,1.040) (-8.157,0.970) (-8.182,0.899) (-8.205,0.828) (-8.226,0.756) (-8.246,0.683) (-8.263,0.610) (-8.278,0.537) (-8.292,0.463) (-8.303,0.389) (-8.313,0.315) (-8.320,0.240) (-8.326,0.165) (-8.329,0.090) (-8.331,0.015) (-8.330,-0.060) (-8.327,-0.134) (-8.323,-0.209) (-8.316,-0.284) (-8.307,-0.359) (-8.297,-0.433) (-8.284,-0.507) (-8.270,-0.580) (-8.253,-0.653) (-8.234,-0.726) (-8.214,-0.798) (-8.192,-0.870) (-8.168,-0.941) (-8.142,-1.011) (-8.114,-1.081) (-8.084,-1.150) (-8.053,-1.218) (-8.020,-1.285) (-7.985,-1.352) (-7.949,-1.417) (-7.911,-1.482) (-7.871,-1.546) (-7.830,-1.608) (-7.787,-1.670) (-7.743,-1.731) (-7.698,-1.790) (-7.651,-1.849) (-7.603,-1.906) (-7.553,-1.963) (-7.502,-2.018) (-7.450,-2.072) (-7.397,-2.124) (-7.342,-2.176) (-7.287,-2.226) (-7.230,-2.276) (-7.173,-2.324) (-7.114,-2.370) (-7.054,-2.416) (-6.994,-2.460) (-6.932,-2.503) (-6.870,-2.544) (-6.806,-2.585) (-6.742,-2.624) (-6.678,-2.662) (-6.612,-2.698) (-6.546,-2.733) (-6.479,-2.767) (-6.412,-2.800) (-6.343,-2.831) (-6.275,-2.862) (-6.206,-2.890) (-6.136,-2.918) (-6.066,-2.944) (-5.995,-2.969) (-5.924,-2.993) (-5.852,-3.015) (-5.780,-3.036) (-5.708,-3.056) (-5.635,-3.075) (-5.562,-3.092) (-5.489,-3.108) (-5.415,-3.123) (-5.342,-3.137) (-5.268,-3.149) (-5.194,-3.160) (-5.119,-3.170) (-5.045,-3.179) (-4.970,-3.187) (-4.895,-3.193) (-4.821,-3.198) (-4.746,-3.202) (-4.671,-3.205) (-4.596,-3.206) (-4.521,-3.206) (-4.446,-3.206) (-4.371,-3.204) (-4.296,-3.201) (-4.221,-3.196) (-4.146,-3.191) (-4.071,-3.184) (-3.997,-3.177) (-3.922,-3.168) (-3.848,-3.158) (-3.774,-3.147) (-3.700,-3.135) (-3.626,-3.122) (-3.552,-3.107) (-3.479,-3.092) (-3.406,-3.075) (-3.333,-3.058) (-3.260,-3.039) (-3.188,-3.020) (-3.116,-2.999) (-3.044,-2.977) (-2.973,-2.954) (-2.901,-2.930) (-2.831,-2.906) (-2.760,-2.880) (-2.690,-2.853) (-2.621,-2.825) (-2.551,-2.796) (-2.483,-2.766) (-2.414,-2.735) (-2.346,-2.703) (-2.279,-2.670) (-2.212,-2.637) (-2.146,-2.602) (-2.080,-2.566) (-2.014,-2.529) (-1.950,-2.492) (-1.885,-2.453) (-1.822,-2.413) (-1.758,-2.373) (-1.696,-2.331) (-1.634,-2.289) (-1.573,-2.246) (-1.512,-2.201) (-1.452,-2.156) (-1.393,-2.110) (-1.335,-2.063) (-1.277,-2.016) (-1.220,-1.967) (-1.164,-1.917) (-1.108,-1.867) (-1.053,-1.815) (-1.000,-1.763) (-0.947,-1.710) (-0.895,-1.656) (-0.843,-1.601) (-0.793,-1.545) (-0.744,-1.489) (-0.696,-1.432) (-0.648,-1.373) (-0.602,-1.314) (-0.557,-1.254) (-0.513,-1.194) (-0.470,-1.132) (-0.442,-1.091)}
\LC{(-1.082,0.489) (-1.149,0.455) (-1.212,0.414) (-1.269,0.365) (-1.319,0.309) (-1.361,0.248) (-1.394,0.180) (-1.417,0.109) (-1.428,0.035) (-1.428,-0.040) (-1.415,-0.114) (-1.392,-0.185) (-1.358,-0.252) (-1.316,-0.314) (-1.265,-0.369) (-1.207,-0.417) (-1.144,-0.458) (-1.077,-0.491)}
\LC{(4.893,-10.000) (4.845,-9.951) (4.797,-9.893) (4.749,-9.835) (4.702,-9.777) (4.655,-9.719) (4.607,-9.661) (4.560,-9.602) (4.514,-9.544) (4.467,-9.485) (4.420,-9.426) (4.374,-9.368) (4.328,-9.309) (4.281,-9.249) (4.235,-9.190) (4.190,-9.131) (4.144,-9.071) (4.099,-9.012) (4.053,-8.952) (4.008,-8.892) (3.963,-8.832) (3.918,-8.772) (3.873,-8.712) (3.829,-8.651) (3.785,-8.591) (3.740,-8.530) (3.696,-8.470) (3.652,-8.409) (3.609,-8.348) (3.565,-8.287) (3.522,-8.226) (3.479,-8.164) (3.435,-8.103) (3.393,-8.041) (3.350,-7.980) (3.307,-7.918) (3.265,-7.856) (3.223,-7.794) (3.181,-7.732) (3.139,-7.670) (3.097,-7.607) (3.056,-7.545) (3.015,-7.482) (2.973,-7.419) (2.932,-7.357) (2.892,-7.294) (2.851,-7.231) (2.811,-7.167) (2.771,-7.104) (2.730,-7.041) (2.691,-6.977) (2.651,-6.913) (2.612,-6.850) (2.572,-6.786) (2.533,-6.722) (2.494,-6.658) (2.456,-6.593) (2.417,-6.529) (2.379,-6.465) (2.341,-6.400) (2.303,-6.335) (2.265,-6.270) (2.228,-6.205) (2.190,-6.140) (2.153,-6.075) (2.116,-6.010) (2.080,-5.945) (2.043,-5.879) (2.007,-5.813) (1.971,-5.748) (1.935,-5.682) (1.899,-5.616) (1.864,-5.550) (1.829,-5.483) (1.794,-5.417) (1.759,-5.351) (1.724,-5.284) (1.690,-5.217) (1.656,-5.151) (1.622,-5.084) (1.588,-5.017) (1.555,-4.950) (1.522,-4.882) (1.489,-4.815) (1.456,-4.747) (1.423,-4.680) (1.391,-4.612) (1.359,-4.544) (1.327,-4.476) (1.296,-4.408) (1.265,-4.340) (1.234,-4.272) (1.203,-4.203) (1.172,-4.135) (1.142,-4.066) (1.112,-3.998) (1.082,-3.929) (1.053,-3.860) (1.024,-3.791) (0.995,-3.721) (0.966,-3.652) (0.938,-3.583) (0.910,-3.513) (0.882,-3.444) (0.854,-3.374) (0.827,-3.304) (0.800,-3.234) (0.773,-3.164) (0.747,-3.094) (0.721,-3.023) (0.695,-2.953) (0.670,-2.882) (0.645,-2.812) (0.620,-2.741) (0.595,-2.670) (0.571,-2.599) (0.547,-2.528) (0.524,-2.457) (0.501,-2.385) (0.478,-2.314) (0.455,-2.242) (0.433,-2.171) (0.412,-2.099) (0.390,-2.027) (0.369,-1.955) (0.349,-1.883) (0.328,-1.811) (0.309,-1.738) (0.289,-1.666) (0.283,-1.642)}
\LC{(-5.042,-10.000) (-4.993,-9.951) (-4.944,-9.894) (-4.895,-9.837) (-4.847,-9.780) (-4.798,-9.722) (-4.750,-9.665) (-4.702,-9.607) (-4.654,-9.550) (-4.606,-9.492) (-4.558,-9.434) (-4.511,-9.376) (-4.463,-9.318) (-4.416,-9.259) (-4.369,-9.201) (-4.322,-9.142) (-4.276,-9.084) (-4.229,-9.025) (-4.183,-8.966) (-4.136,-8.907) (-4.090,-8.848) (-4.044,-8.789) (-3.999,-8.729) (-3.953,-8.670) (-3.908,-8.610) (-3.862,-8.550) (-3.817,-8.490) (-3.772,-8.430) (-3.728,-8.370) (-3.683,-8.310) (-3.639,-8.249) (-3.594,-8.189) (-3.550,-8.128) (-3.506,-8.067) (-3.463,-8.006) (-3.419,-7.945) (-3.376,-7.884) (-3.333,-7.823) (-3.289,-7.761) (-3.247,-7.700) (-3.204,-7.638) (-3.161,-7.576) (-3.119,-7.514) (-3.077,-7.452) (-3.035,-7.390) (-2.993,-7.328) (-2.952,-7.265) (-2.910,-7.203) (-2.869,-7.140) (-2.828,-7.077) (-2.787,-7.015) (-2.747,-6.951) (-2.706,-6.888) (-2.666,-6.825) (-2.626,-6.762) (-2.586,-6.698) (-2.546,-6.634) (-2.507,-6.571) (-2.467,-6.507) (-2.428,-6.443) (-2.389,-6.379) (-2.351,-6.314) (-2.312,-6.250) (-2.274,-6.186) (-2.236,-6.121) (-2.198,-6.056) (-2.160,-5.991) (-2.123,-5.926) (-2.086,-5.861) (-2.049,-5.796) (-2.012,-5.731) (-1.975,-5.665) (-1.939,-5.600) (-1.903,-5.534) (-1.867,-5.468) (-1.831,-5.402) (-1.796,-5.336) (-1.760,-5.270) (-1.725,-5.204) (-1.690,-5.137) (-1.656,-5.071) (-1.621,-5.004) (-1.587,-4.937) (-1.553,-4.870) (-1.520,-4.803) (-1.486,-4.736) (-1.453,-4.669) (-1.420,-4.601) (-1.388,-4.534) (-1.355,-4.466) (-1.323,-4.399) (-1.291,-4.331) (-1.260,-4.263) (-1.228,-4.195) (-1.197,-4.126) (-1.166,-4.058) (-1.135,-3.990) (-1.105,-3.921) (-1.075,-3.852) (-1.045,-3.783) (-1.016,-3.714) (-0.986,-3.645) (-0.957,-3.576) (-0.929,-3.507) (-0.900,-3.438) (-0.872,-3.368) (-0.844,-3.298) (-0.817,-3.229) (-0.790,-3.159) (-0.763,-3.089) (-0.736,-3.019) (-0.710,-2.948) (-0.684,-2.878) (-0.658,-2.808) (-0.633,-2.737) (-0.608,-2.666) (-0.583,-2.595) (-0.559,-2.524) (-0.535,-2.453) (-0.511,-2.382) (-0.488,-2.311) (-0.465,-2.240) (-0.442,-2.168) (-0.420,-2.096) (-0.398,-2.025) (-0.377,-1.953) (-0.356,-1.881) (-0.335,-1.809) (-0.315,-1.736) (-0.295,-1.664)}
\LC{(10.000,-6.289) (9.931,-6.297) (9.856,-6.305) (9.782,-6.311) (9.707,-6.317) (9.632,-6.323) (9.557,-6.328) (9.482,-6.332) (9.407,-6.336) (9.332,-6.339) (9.257,-6.342) (9.182,-6.344) (9.107,-6.345) (9.032,-6.346) (8.957,-6.346) (8.882,-6.346) (8.808,-6.344) (8.733,-6.343) (8.658,-6.341) (8.583,-6.338) (8.508,-6.335) (8.433,-6.331) (8.358,-6.326) (8.283,-6.321) (8.208,-6.315) (8.134,-6.309) (8.059,-6.302) (7.984,-6.295) (7.910,-6.287) (7.835,-6.278) (7.761,-6.269) (7.686,-6.260) (7.612,-6.249) (7.538,-6.239) (7.464,-6.227) (7.390,-6.215) (7.316,-6.203) (7.242,-6.190) (7.168,-6.176) (7.094,-6.162) (7.021,-6.148) (6.947,-6.133) (6.874,-6.117) (6.801,-6.100) (6.728,-6.084) (6.655,-6.066) (6.582,-6.048) (6.509,-6.030) (6.437,-6.011) (6.364,-5.991) (6.292,-5.971) (6.220,-5.951) (6.148,-5.930) (6.076,-5.908) (6.004,-5.886) (5.933,-5.863) (5.862,-5.840) (5.791,-5.816) (5.720,-5.792) (5.649,-5.767) (5.578,-5.742) (5.508,-5.716) (5.438,-5.690) (5.367,-5.663) (5.298,-5.636) (5.228,-5.608) (5.159,-5.579) (5.089,-5.550) (5.020,-5.521) (4.952,-5.491) (4.883,-5.461) (4.815,-5.430) (4.747,-5.399) (4.679,-5.367) (4.611,-5.334) (4.544,-5.301) (4.476,-5.268) (4.410,-5.234) (4.343,-5.200) (4.276,-5.165) (4.210,-5.130) (4.144,-5.094) (4.079,-5.058) (4.013,-5.021) (3.948,-4.984) (3.883,-4.946) (3.819,-4.908) (3.755,-4.869) (3.691,-4.830) (3.627,-4.790) (3.564,-4.750) (3.501,-4.709) (3.438,-4.668) (3.375,-4.627) (3.313,-4.585) (3.251,-4.542) (3.190,-4.499) (3.129,-4.456) (3.068,-4.412) (3.008,-4.368) (2.947,-4.323) (2.888,-4.278) (2.828,-4.232) (2.769,-4.186) (2.710,-4.139) (2.652,-4.092) (2.594,-4.045) (2.536,-3.997) (2.479,-3.948) (2.422,-3.899) (2.366,-3.850) (2.309,-3.800) (2.254,-3.750) (2.198,-3.699) (2.144,-3.648) (2.089,-3.597) (2.035,-3.545) (1.981,-3.492) (1.928,-3.439) (1.875,-3.386) (1.823,-3.332) (1.771,-3.278) (1.720,-3.223) (1.669,-3.168) (1.618,-3.113) (1.568,-3.057) (1.519,-3.001) (1.470,-2.944) (1.421,-2.887) (1.373,-2.829) (1.326,-2.771) (1.279,-2.713) (1.233,-2.654) (1.187,-2.594) (1.141,-2.535) (1.097,-2.474) (1.052,-2.414)}
\LC{(-2.969,2.475) (-3.043,2.485) (-3.118,2.495) (-3.192,2.503) (-3.267,2.509) (-3.342,2.515) (-3.417,2.518) (-3.492,2.520) (-3.567,2.521) (-3.642,2.520) (-3.717,2.518) (-3.792,2.514) (-3.866,2.509) (-3.941,2.502) (-4.016,2.494) (-4.090,2.484) (-4.164,2.472) (-4.238,2.459) (-4.312,2.445) (-4.385,2.428) (-4.458,2.411) (-4.530,2.391) (-4.602,2.370) (-4.674,2.348) (-4.745,2.323) (-4.815,2.298) (-4.885,2.270) (-4.954,2.241) (-5.022,2.210) (-5.090,2.178) (-5.157,2.144) (-5.223,2.108) (-5.288,2.071) (-5.352,2.032) (-5.415,1.991) (-5.477,1.949) (-5.538,1.905) (-5.598,1.860) (-5.656,1.813) (-5.713,1.765) (-5.769,1.714) (-5.823,1.663) (-5.876,1.610) (-5.928,1.555) (-5.977,1.499) (-6.025,1.441) (-6.072,1.382) (-6.116,1.322) (-6.159,1.260) (-6.200,1.197) (-6.239,1.133) (-6.275,1.068) (-6.310,1.001) (-6.343,0.934) (-6.373,0.865) (-6.401,0.796) (-6.427,0.725) (-6.450,0.654) (-6.471,0.582) (-6.490,0.509) (-6.506,0.436) (-6.520,0.362) (-6.531,0.288) (-6.539,0.214) (-6.545,0.139) (-6.549,0.064) (-6.550,-0.011) (-6.548,-0.086) (-6.544,-0.161) (-6.537,-0.236) (-6.528,-0.310) (-6.516,-0.384) (-6.502,-0.458) (-6.485,-0.531) (-6.465,-0.603) (-6.444,-0.675) (-6.419,-0.746) (-6.393,-0.816) (-6.364,-0.885) (-6.333,-0.954) (-6.300,-1.021) (-6.265,-1.087) (-6.227,-1.152) (-6.188,-1.216) (-6.147,-1.279) (-6.103,-1.340) (-6.058,-1.400) (-6.012,-1.458) (-5.963,-1.515) (-5.913,-1.571) (-5.861,-1.625) (-5.808,-1.678) (-5.753,-1.729) (-5.697,-1.779) (-5.639,-1.827) (-5.580,-1.874) (-5.520,-1.918) (-5.459,-1.962) (-5.397,-2.004) (-5.333,-2.044) (-5.269,-2.082) (-5.203,-2.119) (-5.137,-2.154) (-5.070,-2.188) (-5.002,-2.219) (-4.934,-2.250) (-4.864,-2.278) (-4.794,-2.305) (-4.724,-2.331) (-4.653,-2.355) (-4.581,-2.377) (-4.509,-2.397) (-4.436,-2.416) (-4.363,-2.433) (-4.290,-2.449) (-4.216,-2.463) (-4.142,-2.476) (-4.068,-2.487) (-3.994,-2.496) (-3.919,-2.504) (-3.845,-2.511) (-3.770,-2.515) (-3.695,-2.519) (-3.620,-2.521) (-3.545,-2.521) (-3.470,-2.520) (-3.395,-2.517) (-3.320,-2.513) (-3.245,-2.508) (-3.170,-2.501) (-3.096,-2.492) (-3.022,-2.482) (-2.972,-2.475)}
\LC{(0.389,1.769) (0.413,1.840) (0.437,1.911) (0.463,1.982) (0.488,2.052) (0.514,2.123) (0.541,2.193) (0.568,2.263) (0.595,2.333) (0.623,2.402) (0.651,2.472) (0.680,2.541) (0.709,2.610) (0.738,2.679) (0.768,2.748) (0.798,2.817) (0.829,2.885) (0.860,2.954) (0.891,3.022) (0.923,3.090) (0.955,3.158) (0.987,3.225) (1.020,3.293) (1.053,3.360) (1.086,3.427) (1.120,3.494) (1.154,3.561) (1.189,3.628) (1.223,3.694) (1.258,3.760) (1.294,3.827) (1.329,3.893) (1.365,3.958) (1.402,4.024) (1.438,4.089) (1.475,4.155) (1.512,4.220) (1.550,4.285) (1.588,4.350) (1.626,4.414) (1.664,4.479) (1.703,4.543) (1.742,4.607) (1.781,4.671) (1.821,4.735) (1.860,4.798) (1.900,4.862) (1.941,4.925) (1.981,4.988) (2.022,5.051) (2.063,5.113) (2.105,5.176) (2.146,5.238) (2.188,5.301) (2.230,5.363) (2.273,5.424) (2.316,5.486) (2.358,5.548) (2.402,5.609) (2.445,5.670) (2.489,5.731) (2.533,5.792) (2.577,5.853) (2.621,5.913) (2.666,5.973) (2.711,6.033) (2.756,6.093) (2.801,6.153) (2.847,6.213) (2.892,6.272) (2.938,6.331) (2.985,6.390) (3.031,6.449) (3.078,6.508) (3.125,6.566) (3.172,6.625) (3.219,6.683) (3.267,6.741) (3.315,6.799) (3.363,6.856) (3.411,6.914) (3.459,6.971) (3.508,7.028) (3.557,7.085) (3.606,7.142) (3.655,7.198) (3.705,7.255) (3.754,7.311) (3.804,7.367) (3.854,7.423) (3.905,7.478) (3.955,7.534) (4.006,7.589) (4.057,7.644) (4.108,7.699) (4.159,7.754) (4.210,7.808) (4.262,7.863) (4.314,7.917) (4.366,7.971) (4.418,8.025) (4.471,8.078) (4.523,8.132) (4.576,8.185) (4.629,8.238) (4.682,8.291) (4.735,8.344) (4.789,8.396) (4.843,8.449) (4.896,8.501) (4.950,8.553) (5.005,8.605) (5.059,8.656) (5.114,8.708) (5.168,8.759) (5.223,8.810) (5.278,8.861) (5.334,8.912) (5.389,8.962) (5.445,9.013) (5.500,9.063) (5.556,9.113) (5.613,9.163) (5.669,9.212) (5.725,9.262) (5.782,9.311) (5.839,9.360) (5.895,9.409) (5.953,9.457) (6.010,9.506) (6.067,9.554) (6.125,9.602) (6.183,9.650) (6.240,9.698) (6.298,9.745) (6.357,9.793) (6.415,9.840) (6.473,9.887) (6.532,9.933) (6.591,9.980) (6.630,10.000)}
\LC{(-3.797,3.606) (-3.869,3.626) (-3.942,3.645) (-4.014,3.663) (-4.087,3.680) (-4.161,3.696) (-4.234,3.712) (-4.307,3.726) (-4.381,3.740) (-4.455,3.753) (-4.529,3.765) (-4.603,3.775) (-4.678,3.785) (-4.752,3.794) (-4.827,3.802) (-4.901,3.810) (-4.976,3.816) (-5.051,3.821) (-5.126,3.825) (-5.201,3.829) (-5.276,3.831) (-5.351,3.832) (-5.426,3.833) (-5.501,3.832) (-5.576,3.831) (-5.651,3.828) (-5.726,3.825) (-5.800,3.820) (-5.875,3.815) (-5.950,3.808) (-6.025,3.801) (-6.099,3.792) (-6.173,3.782) (-6.248,3.772) (-6.322,3.760) (-6.396,3.748) (-6.470,3.734) (-6.543,3.719) (-6.616,3.703) (-6.689,3.687) (-6.762,3.669) (-6.835,3.650) (-6.907,3.630) (-6.979,3.609) (-7.051,3.587) (-7.122,3.564) (-7.193,3.540) (-7.264,3.514) (-7.334,3.488) (-7.404,3.461) (-7.473,3.432) (-7.542,3.402) (-7.611,3.372) (-7.679,3.340) (-7.746,3.307) (-7.813,3.273) (-7.879,3.239) (-7.945,3.202) (-8.010,3.165) (-8.075,3.127) (-8.139,3.088) (-8.202,3.048) (-8.265,3.006) (-8.326,2.964) (-8.387,2.920) (-8.448,2.876) (-8.507,2.830) (-8.566,2.784) (-8.624,2.736) (-8.681,2.687) (-8.737,2.637) (-8.792,2.587) (-8.846,2.535) (-8.900,2.482) (-8.952,2.428) (-9.003,2.373) (-9.054,2.318) (-9.103,2.261) (-9.151,2.203) (-9.198,2.145) (-9.243,2.085) (-9.288,2.025) (-9.331,1.964) (-9.373,1.902) (-9.414,1.839) (-9.454,1.775) (-9.492,1.711) (-9.529,1.645) (-9.564,1.579) (-9.598,1.513) (-9.631,1.445) (-9.662,1.377) (-9.692,1.308) (-9.720,1.239) (-9.747,1.169) (-9.772,1.098) (-9.796,1.027) (-9.818,0.955) (-9.839,0.883) (-9.858,0.810) (-9.875,0.737) (-9.891,0.664) (-9.905,0.590) (-9.918,0.517) (-9.928,0.442) (-9.938,0.368) (-9.945,0.293) (-9.951,0.218) (-9.955,0.144) (-9.957,0.069) (-9.958,-0.006) (-9.957,-0.081) (-9.954,-0.156) (-9.950,-0.231) (-9.944,-0.306) (-9.936,-0.381) (-9.927,-0.455) (-9.916,-0.529) (-9.903,-0.603) (-9.888,-0.677) (-9.872,-0.750) (-9.855,-0.823) (-9.836,-0.895) (-9.815,-0.967) (-9.792,-1.039) (-9.768,-1.110) (-9.743,-1.180) (-9.716,-1.250) (-9.687,-1.320) (-9.657,-1.388) (-9.626,-1.457) (-9.593,-1.524) (-9.558,-1.591) (-9.522,-1.656) (-9.485,-1.722) (-9.447,-1.786) (-9.407,-1.850) (-9.366,-1.912) (-9.324,-1.974) (-9.280,-2.035) (-9.236,-2.096) (-9.190,-2.155) (-9.143,-2.213) (-9.094,-2.271) (-9.045,-2.327) (-8.995,-2.383) (-8.943,-2.437) (-8.891,-2.491) (-8.837,-2.544) (-8.783,-2.595) (-8.728,-2.646) (-8.671,-2.695) (-8.614,-2.744) (-8.556,-2.792) (-8.497,-2.838) (-8.438,-2.883) (-8.377,-2.928) (-8.316,-2.971) (-8.254,-3.013) (-8.191,-3.055) (-8.128,-3.095) (-8.064,-3.134) (-7.999,-3.172) (-7.934,-3.209) (-7.868,-3.245) (-7.802,-3.279) (-7.735,-3.313) (-7.667,-3.346) (-7.599,-3.377) (-7.530,-3.408) (-7.461,-3.437) (-7.392,-3.465) (-7.322,-3.492) (-7.252,-3.519) (-7.181,-3.544) (-7.110,-3.568) (-7.039,-3.591) (-6.967,-3.613) (-6.895,-3.633) (-6.823,-3.653) (-6.750,-3.672) (-6.677,-3.690) (-6.604,-3.706) (-6.531,-3.722) (-6.457,-3.736) (-6.383,-3.750) (-6.309,-3.762) (-6.235,-3.774) (-6.161,-3.784) (-6.086,-3.794) (-6.012,-3.802) (-5.937,-3.809) (-5.863,-3.816) (-5.788,-3.821) (-5.713,-3.825) (-5.638,-3.829) (-5.563,-3.831) (-5.488,-3.832) (-5.413,-3.833) (-5.338,-3.832) (-5.263,-3.831) (-5.188,-3.828) (-5.113,-3.825) (-5.038,-3.820) (-4.964,-3.815) (-4.889,-3.808) (-4.814,-3.801) (-4.740,-3.793) (-4.665,-3.784) (-4.591,-3.774) (-4.517,-3.763) (-4.443,-3.751) (-4.369,-3.738) (-4.295,-3.724) (-4.221,-3.709) (-4.148,-3.694) (-4.075,-3.677) (-4.002,-3.660) (-3.929,-3.641) (-3.857,-3.622) (-3.809,-3.609)}
\LC{(-0.048,-10.000) (-0.048,-9.947) (-0.047,-9.872) (-0.046,-9.797) (-0.046,-9.722) (-0.045,-9.647) (-0.045,-9.572) (-0.044,-9.497) (-0.044,-9.422) (-0.043,-9.347) (-0.043,-9.272) (-0.042,-9.197) (-0.042,-9.122) (-0.041,-9.047) (-0.041,-8.972) (-0.040,-8.897) (-0.040,-8.822) (-0.039,-8.747) (-0.039,-8.672) (-0.038,-8.597) (-0.038,-8.522) (-0.037,-8.447) (-0.037,-8.372) (-0.036,-8.297) (-0.036,-8.222) (-0.035,-8.147) (-0.035,-8.072) (-0.034,-7.997) (-0.034,-7.922) (-0.033,-7.847) (-0.033,-7.772) (-0.032,-7.697) (-0.032,-7.622) (-0.031,-7.547) (-0.031,-7.472) (-0.030,-7.397) (-0.030,-7.322) (-0.030,-7.247) (-0.029,-7.172) (-0.029,-7.097) (-0.028,-7.022) (-0.028,-6.947) (-0.027,-6.872) (-0.027,-6.797) (-0.026,-6.722) (-0.026,-6.647) (-0.026,-6.572) (-0.025,-6.497) (-0.025,-6.422) (-0.024,-6.347) (-0.024,-6.272) (-0.023,-6.197) (-0.023,-6.122) (-0.023,-6.047) (-0.022,-5.972) (-0.022,-5.897) (-0.021,-5.822) (-0.021,-5.747) (-0.020,-5.672) (-0.020,-5.597) (-0.020,-5.522) (-0.019,-5.447) (-0.019,-5.372) (-0.018,-5.297) (-0.018,-5.222) (-0.018,-5.147) (-0.017,-5.072) (-0.017,-4.997) (-0.017,-4.922) (-0.016,-4.847) (-0.016,-4.772) (-0.015,-4.697) (-0.015,-4.622) (-0.015,-4.547) (-0.014,-4.472) (-0.014,-4.397) (-0.014,-4.322) (-0.013,-4.247) (-0.013,-4.172) (-0.013,-4.097) (-0.012,-4.022) (-0.012,-3.947) (-0.012,-3.872) (-0.011,-3.797) (-0.011,-3.722) (-0.011,-3.647) (-0.010,-3.572) (-0.010,-3.497) (-0.010,-3.422) (-0.009,-3.347) (-0.009,-3.272) (-0.009,-3.197) (-0.008,-3.122) (-0.008,-3.047) (-0.008,-2.972) (-0.007,-2.897) (-0.007,-2.822) (-0.007,-2.747) (-0.007,-2.672) (-0.006,-2.597) (-0.006,-2.522) (-0.006,-2.447) (-0.006,-2.372) (-0.005,-2.297)}
\LC{(-0.488,1.887) (-0.516,1.956) (-0.545,2.025) (-0.574,2.094) (-0.604,2.163) (-0.634,2.232) (-0.665,2.300) (-0.696,2.369) (-0.727,2.437) (-0.759,2.504) (-0.792,2.572) (-0.825,2.639) (-0.858,2.707) (-0.892,2.774) (-0.926,2.840) (-0.961,2.907) (-0.996,2.973) (-1.031,3.039) (-1.067,3.105) (-1.103,3.171) (-1.140,3.236) (-1.177,3.301) (-1.214,3.366) (-1.252,3.431) (-1.290,3.496) (-1.329,3.560) (-1.368,3.624) (-1.407,3.688) (-1.446,3.752) (-1.486,3.816) (-1.526,3.879) (-1.567,3.942) (-1.608,4.005) (-1.649,4.067) (-1.691,4.130) (-1.733,4.192) (-1.775,4.254) (-1.818,4.316) (-1.860,4.377) (-1.904,4.439) (-1.947,4.500) (-1.991,4.560) (-2.035,4.621) (-2.080,4.682) (-2.124,4.742) (-2.169,4.802) (-2.215,4.861) (-2.260,4.921) (-2.306,4.980) (-2.352,5.039) (-2.399,5.098) (-2.446,5.157) (-2.493,5.215) (-2.540,5.273) (-2.588,5.331) (-2.635,5.389) (-2.684,5.447) (-2.732,5.504) (-2.781,5.561) (-2.830,5.618) (-2.879,5.674) (-2.928,5.731) (-2.978,5.787) (-3.028,5.843) (-3.078,5.899) (-3.129,5.954) (-3.179,6.009) (-3.230,6.064) (-3.282,6.119) (-3.333,6.174) (-3.385,6.228) (-3.437,6.282) (-3.489,6.336) (-3.542,6.389) (-3.594,6.443) (-3.647,6.496) (-3.700,6.549) (-3.754,6.602) (-3.807,6.654) (-3.861,6.706) (-3.915,6.758) (-3.969,6.810) (-4.024,6.862) (-4.079,6.913) (-4.134,6.964) (-4.189,7.015) (-4.244,7.065) (-4.300,7.116) (-4.355,7.166) (-4.411,7.216) (-4.468,7.265) (-4.524,7.315) (-4.581,7.364) (-4.638,7.413) (-4.695,7.461) (-4.752,7.510) (-4.809,7.558) (-4.867,7.606) (-4.925,7.654) (-4.983,7.701) (-5.041,7.749) (-5.100,7.796) (-5.158,7.842) (-5.217,7.889) (-5.276,7.935) (-5.335,7.981) (-5.395,8.027) (-5.454,8.073) (-5.514,8.118) (-5.574,8.163) (-5.634,8.208) (-5.694,8.252) (-5.755,8.297) (-5.816,8.341) (-5.876,8.385) (-5.937,8.428) (-5.999,8.472) (-6.060,8.515) (-6.122,8.558) (-6.183,8.600) (-6.245,8.643) (-6.307,8.685) (-6.369,8.727) (-6.432,8.768) (-6.494,8.810) (-6.557,8.851) (-6.620,8.892) (-6.683,8.932) (-6.746,8.973) (-6.810,9.013) (-6.873,9.053) (-6.937,9.092) (-7.001,9.132) (-7.065,9.171) (-7.129,9.210) (-7.193,9.248) (-7.257,9.287) (-7.322,9.325) (-7.387,9.363) (-7.452,9.400) (-7.517,9.437) (-7.582,9.475) (-7.647,9.511) (-7.713,9.548) (-7.778,9.584) (-7.844,9.620) (-7.910,9.656) (-7.976,9.692) (-8.042,9.727) (-8.109,9.762) (-8.175,9.797) (-8.242,9.831) (-8.308,9.866) (-8.375,9.900) (-8.442,9.933) (-8.509,9.967) (-8.577,10.000) (-8.599,10.000)}
\LC{(-1.957,3.712) (-2.008,3.767) (-2.059,3.821) (-2.111,3.876) (-2.163,3.930) (-2.215,3.983) (-2.268,4.036) (-2.322,4.089) (-2.375,4.142) (-2.429,4.194) (-2.484,4.245) (-2.538,4.297) (-2.593,4.348) (-2.649,4.398) (-2.704,4.448) (-2.761,4.498) (-2.817,4.547) (-2.874,4.596) (-2.931,4.645) (-2.988,4.693) (-3.046,4.741) (-3.104,4.789) (-3.162,4.836) (-3.221,4.883) (-3.280,4.929) (-3.339,4.975) (-3.399,5.020) (-3.459,5.065) (-3.519,5.110) (-3.579,5.155) (-3.640,5.199) (-3.701,5.242) (-3.763,5.285) (-3.824,5.328) (-3.886,5.370) (-3.948,5.412) (-4.011,5.454) (-4.073,5.495) (-4.136,5.536) (-4.199,5.576) (-4.263,5.616) (-4.327,5.656) (-4.391,5.695) (-4.455,5.734) (-4.519,5.772) (-4.584,5.810) (-4.649,5.848) (-4.714,5.885) (-4.780,5.921) (-4.845,5.958) (-4.911,5.994) (-4.977,6.029) (-5.043,6.064) (-5.110,6.099) (-5.177,6.133) (-5.244,6.167) (-5.311,6.200) (-5.378,6.233) (-5.446,6.265) (-5.514,6.298) (-5.582,6.329) (-5.650,6.360) (-5.718,6.391) (-5.787,6.422) (-5.855,6.452) (-5.924,6.481) (-5.994,6.510) (-6.063,6.539) (-6.132,6.567) (-6.202,6.595) (-6.272,6.622) (-6.342,6.649) (-6.412,6.676) (-6.482,6.702) (-6.553,6.727) (-6.623,6.753) (-6.694,6.777) (-6.765,6.802) (-6.836,6.825) (-6.908,6.849) (-6.979,6.872) (-7.051,6.894) (-7.122,6.916) (-7.194,6.938) (-7.266,6.959) (-7.338,6.980) (-7.410,7.000) (-7.483,7.020) (-7.555,7.039) (-7.628,7.058) (-7.700,7.076) (-7.773,7.094) (-7.846,7.112) (-7.919,7.129) (-7.992,7.145) (-8.066,7.161) (-8.139,7.177) (-8.212,7.192) (-8.286,7.207) (-8.360,7.221) (-8.433,7.235) (-8.507,7.248) (-8.581,7.261) (-8.655,7.273) (-8.729,7.285) (-8.803,7.297) (-8.877,7.308) (-8.952,7.318) (-9.026,7.328) (-9.100,7.337) (-9.175,7.346) (-9.249,7.355) (-9.324,7.363) (-9.399,7.371) (-9.473,7.378) (-9.548,7.384) (-9.623,7.390) (-9.697,7.396) (-9.772,7.401) (-9.847,7.406) (-9.922,7.410) (-9.997,7.413) (-10.000,7.414)}
\LC{(-10.000,-8.150) (-9.943,-8.140) (-9.868,-8.130) (-9.794,-8.119) (-9.720,-8.107) (-9.646,-8.095) (-9.572,-8.083) (-9.498,-8.070) (-9.424,-8.057) (-9.350,-8.044) (-9.277,-8.030) (-9.203,-8.016) (-9.130,-8.001) (-9.056,-7.986) (-8.983,-7.970) (-8.910,-7.954) (-8.836,-7.938) (-8.763,-7.921) (-8.690,-7.903) (-8.617,-7.886) (-8.545,-7.868) (-8.472,-7.849) (-8.399,-7.830) (-8.327,-7.811) (-8.255,-7.791) (-8.182,-7.771) (-8.110,-7.750) (-8.038,-7.729) (-7.966,-7.708) (-7.895,-7.686) (-7.823,-7.664) (-7.751,-7.641) (-7.680,-7.618) (-7.609,-7.595) (-7.538,-7.571) (-7.467,-7.547) (-7.396,-7.522) (-7.325,-7.497) (-7.255,-7.472) (-7.184,-7.446) (-7.114,-7.420) (-7.044,-7.393) (-6.974,-7.366) (-6.904,-7.339) (-6.834,-7.311) (-6.765,-7.283) (-6.695,-7.255) (-6.626,-7.226) (-6.557,-7.196) (-6.488,-7.167) (-6.420,-7.137) (-6.351,-7.106) (-6.283,-7.075) (-6.214,-7.044) (-6.146,-7.012) (-6.079,-6.980) (-6.011,-6.948) (-5.944,-6.915) (-5.876,-6.882) (-5.809,-6.849) (-5.742,-6.815) (-5.676,-6.780) (-5.609,-6.746) (-5.543,-6.711) (-5.477,-6.675) (-5.411,-6.640) (-5.345,-6.603) (-5.279,-6.567) (-5.214,-6.530) (-5.149,-6.493) (-5.084,-6.455) (-5.020,-6.417) (-4.955,-6.379) (-4.891,-6.340) (-4.827,-6.301) (-4.763,-6.261) (-4.700,-6.222) (-4.636,-6.181) (-4.573,-6.141) (-4.510,-6.100) (-4.448,-6.059) (-4.385,-6.017) (-4.323,-5.975) (-4.261,-5.933) (-4.200,-5.890) (-4.138,-5.847) (-4.077,-5.803) (-4.016,-5.760) (-3.956,-5.715) (-3.895,-5.671) (-3.835,-5.626) (-3.775,-5.581) (-3.716,-5.535) (-3.656,-5.489) (-3.597,-5.443) (-3.539,-5.397) (-3.480,-5.350) (-3.422,-5.302) (-3.364,-5.255) (-3.306,-5.207) (-3.249,-5.158) (-3.192,-5.110) (-3.135,-5.061) (-3.079,-5.011) (-3.023,-4.961) (-2.967,-4.911) (-2.911,-4.861) (-2.856,-4.810) (-2.801,-4.759) (-2.747,-4.708) (-2.692,-4.656) (-2.638,-4.604) (-2.585,-4.551) (-2.531,-4.499) (-2.479,-4.445) (-2.426,-4.392) (-2.374,-4.338) (-2.322,-4.284) (-2.270,-4.230) (-2.219,-4.175) (-2.168,-4.120) (-2.118,-4.064) (-2.068,-4.008) (-2.018,-3.952) (-1.968,-3.896) (-1.919,-3.839) (-1.871,-3.782) (-1.823,-3.724) (-1.775,-3.667) (-1.727,-3.609) (-1.680,-3.550) (-1.634,-3.491) (-1.587,-3.432) (-1.542,-3.373) (-1.496,-3.313) (-1.451,-3.253) (-1.407,-3.193) (-1.363,-3.132) (-1.319,-3.071) (-1.276,-3.010) (-1.233,-2.948) (-1.191,-2.886) (-1.149,-2.824)}
\LC{(3.767,2.687) (3.842,2.687) (3.917,2.685) (3.992,2.683) (4.066,2.678) (4.141,2.673) (4.216,2.665) (4.290,2.657) (4.365,2.647) (4.439,2.635) (4.513,2.622) (4.586,2.608) (4.660,2.592) (4.732,2.574) (4.805,2.555) (4.877,2.535) (4.949,2.513) (5.020,2.489) (5.091,2.464) (5.161,2.438) (5.231,2.410) (5.299,2.380) (5.368,2.349) (5.435,2.316) (5.502,2.282) (5.568,2.246) (5.633,2.209) (5.697,2.170) (5.760,2.130) (5.823,2.088) (5.884,2.045) (5.944,2.000) (6.003,1.954) (6.061,1.906) (6.117,1.856) (6.172,1.806) (6.226,1.753) (6.279,1.700) (6.330,1.645) (6.379,1.588) (6.427,1.531) (6.473,1.472) (6.518,1.411) (6.560,1.350) (6.601,1.287) (6.641,1.223) (6.678,1.158) (6.713,1.092) (6.747,1.025) (6.778,0.956) (6.807,0.887) (6.834,0.817) (6.859,0.747) (6.882,0.675) (6.902,0.603) (6.920,0.530) (6.936,0.457) (6.950,0.383) (6.961,0.309) (6.970,0.235) (6.976,0.160) (6.980,0.085) (6.981,0.010) (6.980,-0.065) (6.977,-0.140) (6.971,-0.215) (6.963,-0.289) (6.953,-0.364) (6.940,-0.437) (6.925,-0.511) (6.907,-0.584) (6.887,-0.656) (6.865,-0.728) (6.841,-0.799) (6.814,-0.869) (6.786,-0.938) (6.755,-1.007) (6.722,-1.074) (6.687,-1.141) (6.651,-1.206) (6.612,-1.270) (6.571,-1.333) (6.529,-1.395) (6.485,-1.456) (6.439,-1.515) (6.392,-1.573) (6.343,-1.630) (6.292,-1.685) (6.240,-1.739) (6.187,-1.792) (6.132,-1.843) (6.076,-1.893) (6.018,-1.941) (5.960,-1.988) (5.900,-2.033) (5.839,-2.077) (5.777,-2.119) (5.714,-2.160) (5.650,-2.199) (5.585,-2.237) (5.519,-2.273) (5.453,-2.308) (5.386,-2.341) (5.318,-2.372) (5.249,-2.402) (5.179,-2.431) (5.109,-2.458) (5.039,-2.483) (4.968,-2.507) (4.896,-2.529) (4.824,-2.550) (4.752,-2.569) (4.679,-2.587) (4.606,-2.604) (4.532,-2.618) (4.458,-2.632) (4.384,-2.644) (4.310,-2.654) (4.236,-2.663) (4.161,-2.671) (4.086,-2.677) (4.011,-2.682) (3.936,-2.685) (3.861,-2.687) (3.786,-2.687) (3.761,-2.687)}
\LC{(10.000,-7.569) (9.949,-7.564) (9.875,-7.558) (9.800,-7.553) (9.725,-7.546) (9.650,-7.539) (9.576,-7.532) (9.501,-7.524) (9.427,-7.516) (9.352,-7.508) (9.278,-7.498) (9.203,-7.489) (9.129,-7.479) (9.055,-7.468) (8.980,-7.457) (8.906,-7.445) (8.832,-7.434) (8.758,-7.421) (8.684,-7.408) (8.611,-7.395) (8.537,-7.381) (8.463,-7.367) (8.390,-7.352) (8.316,-7.337) (8.243,-7.321) (8.170,-7.305) (8.097,-7.288) (8.024,-7.271) (7.951,-7.254) (7.878,-7.236) (7.805,-7.218) (7.733,-7.199) (7.660,-7.179) (7.588,-7.160) (7.515,-7.140) (7.443,-7.119) (7.371,-7.098) (7.299,-7.076) (7.228,-7.054) (7.156,-7.032) (7.085,-7.009) (7.013,-6.986) (6.942,-6.962) (6.871,-6.938) (6.800,-6.914) (6.730,-6.889) (6.659,-6.863) (6.589,-6.837) (6.518,-6.811) (6.448,-6.784) (6.378,-6.757) (6.309,-6.730) (6.239,-6.702) (6.170,-6.673) (6.100,-6.645) (6.031,-6.615) (5.963,-6.586) (5.894,-6.555) (5.825,-6.525) (5.757,-6.494) (5.689,-6.463) (5.621,-6.431) (5.553,-6.399) (5.486,-6.366) (5.418,-6.333) (5.351,-6.300) (5.284,-6.266) (5.218,-6.232) (5.151,-6.197) (5.085,-6.162) (5.019,-6.126) (4.953,-6.091) (4.887,-6.054) (4.822,-6.018) (4.757,-5.980) (4.692,-5.943) (4.627,-5.905) (4.562,-5.867) (4.498,-5.828) (4.434,-5.789) (4.370,-5.749) (4.307,-5.710) (4.244,-5.669) (4.181,-5.629) (4.139,-5.601)}
\LC{(3.620,3.580) (3.691,3.603) (3.763,3.624) (3.835,3.645) (3.908,3.665) (3.980,3.684) (4.053,3.703) (4.126,3.720) (4.199,3.736) (4.272,3.752) (4.346,3.767) (4.420,3.781) (4.494,3.793) (4.568,3.805) (4.642,3.816) (4.716,3.827) (4.790,3.836) (4.865,3.844) (4.940,3.851) (5.014,3.858) (5.089,3.863) (5.164,3.868) (5.239,3.871) (5.314,3.874) (5.389,3.876) (5.464,3.876) (5.539,3.876) (5.614,3.875) (5.689,3.873) (5.764,3.870) (5.839,3.865) (5.914,3.860) (5.988,3.854) (6.063,3.847) (6.137,3.839) (6.212,3.830) (6.286,3.820) (6.360,3.808) (6.434,3.796) (6.508,3.783) (6.582,3.769) (6.655,3.754) (6.728,3.737) (6.801,3.720) (6.874,3.702) (6.947,3.682) (7.019,3.662) (7.091,3.640) (7.162,3.618) (7.233,3.594) (7.304,3.570) (7.375,3.544) (7.445,3.517) (7.514,3.489) (7.584,3.460) (7.652,3.431) (7.721,3.400) (7.788,3.367) (7.856,3.334) (7.922,3.300) (7.989,3.265) (8.054,3.229) (8.119,3.191) (8.184,3.153) (8.247,3.113) (8.311,3.073) (8.373,3.031) (8.435,2.988) (8.496,2.945) (8.556,2.900) (8.615,2.854) (8.674,2.807) (8.731,2.759) (8.788,2.710) (8.844,2.660) (8.899,2.610) (8.953,2.558) (9.007,2.505) (9.059,2.451) (9.110,2.396) (9.160,2.340) (9.209,2.284) (9.257,2.226) (9.304,2.167) (9.350,2.108) (9.394,2.048) (9.438,1.986) (9.480,1.924) (9.520,1.861) (9.560,1.798) (9.598,1.733) (9.635,1.668) (9.671,1.602) (9.705,1.535) (9.738,1.468) (9.769,1.400) (9.799,1.331) (9.828,1.261) (9.855,1.191) (9.880,1.121) (9.904,1.050) (9.927,0.978) (9.947,0.906) (9.967,0.834) (9.984,0.761) (10.000,0.688)}
\LC{(2.538,5.010) (2.587,5.067) (2.637,5.123) (2.687,5.179) (2.737,5.234) (2.788,5.290) (2.839,5.345) (2.890,5.400) (2.941,5.454) (2.993,5.508) (3.045,5.563) (3.097,5.616) (3.150,5.670) (3.203,5.723) (3.256,5.776) (3.309,5.829) (3.363,5.882) (3.416,5.934) (3.470,5.986) (3.525,6.037) (3.579,6.089) (3.634,6.140) (3.689,6.191) (3.745,6.242) (3.800,6.292) (3.856,6.342) (3.912,6.392) (3.968,6.442) (4.025,6.491) (4.081,6.540) (4.138,6.589) (4.196,6.637) (4.253,6.685) (4.311,6.733) (4.369,6.781) (4.427,6.829) (4.485,6.876) (4.544,6.923) (4.602,6.969) (4.661,7.016) (4.720,7.062) (4.780,7.107) (4.840,7.153) (4.899,7.198) (4.959,7.243) (5.020,7.288) (5.080,7.332) (5.141,7.376) (5.202,7.420) (5.263,7.464) (5.324,7.507) (5.385,7.550) (5.447,7.593) (5.509,7.635) (5.571,7.677) (5.633,7.719) (5.695,7.761) (5.758,7.802) (5.821,7.843) (5.884,7.884) (5.947,7.925) (6.010,7.965) (6.074,8.005) (6.137,8.044) (6.201,8.084) (6.265,8.123) (6.329,8.162) (6.394,8.200) (6.458,8.238) (6.523,8.276) (6.588,8.314) (6.653,8.351) (6.718,8.388) (6.783,8.425) (6.849,8.462) (6.915,8.498) (6.980,8.534) (7.046,8.569) (7.113,8.605) (7.179,8.640) (7.245,8.674) (7.312,8.709) (7.379,8.743) (7.446,8.777) (7.513,8.810) (7.580,8.844) (7.647,8.877) (7.715,8.909) (7.782,8.942) (7.850,8.974) (7.918,9.005) (7.986,9.037) (8.054,9.068) (8.123,9.099) (8.191,9.130) (8.260,9.160) (8.329,9.190) (8.398,9.220) (8.467,9.249) (8.536,9.278) (8.605,9.307) (8.674,9.335) (8.744,9.363) (8.813,9.391) (8.883,9.419) (8.953,9.446) (9.023,9.473) (9.093,9.500) (9.163,9.526) (9.234,9.552) (9.304,9.578) (9.375,9.603) (9.446,9.628) (9.516,9.653) (9.587,9.677) (9.658,9.702) (9.729,9.725) (9.801,9.749) (9.872,9.772) (9.943,9.795) (10.000,9.818)}
\LC{(1.312,1.428) (1.381,1.456) (1.451,1.482) (1.523,1.506) (1.594,1.528) (1.667,1.548) (1.739,1.566) (1.813,1.582) (1.886,1.595) (1.961,1.607) (2.035,1.616) (2.110,1.623) (2.184,1.628) (2.259,1.631) (2.334,1.631) (2.409,1.629) (2.484,1.625) (2.559,1.618) (2.633,1.609) (2.708,1.597) (2.781,1.583) (2.854,1.567) (2.927,1.548) (2.999,1.527) (3.070,1.503) (3.140,1.477) (3.209,1.448) (3.278,1.417) (3.345,1.383) (3.410,1.346) (3.474,1.307) (3.537,1.266) (3.597,1.222) (3.656,1.176) (3.713,1.127) (3.768,1.075) (3.820,1.022) (3.870,0.966) (3.918,0.908) (3.963,0.848) (4.004,0.785) (4.043,0.721) (4.079,0.655) (4.111,0.588) (4.140,0.518) (4.165,0.448) (4.187,0.376) (4.205,0.303) (4.219,0.230) (4.229,0.155) (4.235,0.080) (4.238,0.006) (4.236,-0.069) (4.230,-0.144) (4.221,-0.219) (4.207,-0.292) (4.190,-0.365) (4.169,-0.437) (4.144,-0.508) (4.116,-0.577) (4.084,-0.645) (4.049,-0.712) (4.010,-0.776) (3.969,-0.839) (3.925,-0.899) (3.878,-0.958) (3.828,-1.014) (3.776,-1.068) (3.721,-1.119) (3.665,-1.169) (3.606,-1.215) (3.546,-1.260) (3.483,-1.301) (3.420,-1.341) (3.354,-1.378) (3.288,-1.412) (3.220,-1.444) (3.151,-1.473) (3.080,-1.499) (3.009,-1.524) (2.938,-1.545) (2.865,-1.564) (2.792,-1.581) (2.718,-1.595) (2.644,-1.607) (2.570,-1.617) (2.495,-1.624) (2.420,-1.628) (2.345,-1.631) (2.270,-1.631) (2.195,-1.629) (2.121,-1.624) (2.046,-1.617) (1.971,-1.608) (1.897,-1.597) (1.824,-1.584) (1.750,-1.568) (1.677,-1.551) (1.605,-1.531) (1.533,-1.510) (1.462,-1.486) (1.391,-1.460) (1.322,-1.432) (1.299,-1.423)}
\LC{(2.736,-10.000) (2.707,-9.935) (2.678,-9.866) (2.648,-9.797) (2.619,-9.728) (2.591,-9.659) (2.562,-9.589) (2.533,-9.520) (2.504,-9.451) (2.476,-9.382) (2.447,-9.312) (2.419,-9.243) (2.391,-9.173) (2.363,-9.104) (2.335,-9.034) (2.307,-8.964) (2.279,-8.895) (2.252,-8.825) (2.224,-8.755) (2.197,-8.685) (2.169,-8.616) (2.142,-8.546) (2.115,-8.476) (2.088,-8.406) (2.061,-8.336) (2.034,-8.266) (2.008,-8.196) (1.981,-8.125) (1.955,-8.055) (1.929,-7.985) (1.902,-7.915) (1.876,-7.844) (1.851,-7.774) (1.825,-7.704) (1.799,-7.633) (1.773,-7.563) (1.748,-7.492) (1.723,-7.421) (1.697,-7.351) (1.672,-7.280) (1.647,-7.209) (1.623,-7.139) (1.598,-7.068) (1.573,-6.997) (1.549,-6.926) (1.524,-6.855) (1.500,-6.784) (1.476,-6.713) (1.452,-6.642) (1.428,-6.571) (1.405,-6.500) (1.381,-6.429) (1.358,-6.357) (1.334,-6.286) (1.311,-6.215) (1.288,-6.143) (1.265,-6.072) (1.243,-6.000) (1.220,-5.929) (1.197,-5.857) (1.175,-5.786) (1.153,-5.714) (1.131,-5.642) (1.109,-5.571) (1.087,-5.499) (1.066,-5.427) (1.044,-5.355) (1.023,-5.283) (1.001,-5.211) (0.980,-5.139) (0.959,-5.067) (0.939,-4.995) (0.918,-4.923) (0.898,-4.851) (0.877,-4.779) (0.857,-4.707) (0.837,-4.634) (0.817,-4.562) (0.798,-4.490) (0.778,-4.417) (0.759,-4.345) (0.740,-4.272)}
\LC{(-3.539,5.679) (-3.596,5.728) (-3.653,5.776) (-3.711,5.824) (-3.768,5.872) (-3.826,5.920) (-3.885,5.967) (-3.943,6.014) (-4.002,6.061) (-4.061,6.107) (-4.120,6.153) (-4.179,6.199) (-4.239,6.244) (-4.299,6.289) (-4.359,6.334) (-4.420,6.378) (-4.480,6.422) (-4.541,6.466) (-4.602,6.509) (-4.664,6.553) (-4.725,6.595) (-4.787,6.638) (-4.849,6.680) (-4.911,6.722) (-4.974,6.764) (-5.037,6.805) (-5.099,6.846) (-5.162,6.886) (-5.226,6.926) (-5.289,6.966) (-5.353,7.006) (-5.417,7.045) (-5.481,7.084) (-5.545,7.123) (-5.610,7.161) (-5.674,7.199) (-5.739,7.237) (-5.804,7.274) (-5.870,7.311) (-5.935,7.347) (-6.001,7.384) (-6.067,7.420) (-6.133,7.455) (-6.199,7.490) (-6.265,7.525) (-6.332,7.560) (-6.398,7.594) (-6.465,7.628) (-6.532,7.662) (-6.600,7.695) (-6.667,7.728) (-6.735,7.761) (-6.802,7.793) (-6.870,7.825) (-6.938,7.856) (-7.006,7.887) (-7.075,7.918) (-7.143,7.949) (-7.212,7.979) (-7.281,8.009) (-7.350,8.038) (-7.419,8.067) (-7.488,8.096) (-7.558,8.125) (-7.627,8.153) (-7.697,8.180) (-7.767,8.208) (-7.837,8.235) (-7.907,8.261) (-7.977,8.288) (-8.047,8.314) (-8.118,8.339) (-8.188,8.365) (-8.259,8.389) (-8.330,8.414) (-8.401,8.438) (-8.472,8.462) (-8.543,8.485) (-8.615,8.509) (-8.686,8.531) (-8.758,8.554) (-8.829,8.576) (-8.901,8.597) (-8.973,8.619) (-9.045,8.640) (-9.117,8.660) (-9.190,8.680) (-9.262,8.700) (-9.334,8.720) (-9.407,8.739) (-9.479,8.757) (-9.552,8.776) (-9.625,8.794) (-9.698,8.811) (-9.771,8.829) (-9.844,8.845) (-9.917,8.862) (-9.990,8.878) (-10.000,8.883)}
\LC{(8.623,2.249) (8.673,2.193) (8.722,2.136) (8.769,2.078) (8.815,2.019) (8.860,1.959) (8.904,1.898) (8.947,1.836) (8.988,1.773) (9.028,1.710) (9.066,1.646) (9.103,1.580) (9.139,1.514) (9.173,1.448) (9.206,1.380) (9.237,1.312) (9.267,1.243) (9.295,1.174) (9.321,1.103) (9.346,1.033) (9.370,0.961) (9.391,0.890) (9.411,0.817) (9.430,0.745) (9.446,0.671) (9.461,0.598) (9.474,0.524) (9.486,0.450) (9.496,0.376) (9.504,0.301) (9.510,0.226) (9.514,0.151) (9.517,0.076) (9.518,0.001) (9.517,-0.073) (9.514,-0.148) (9.510,-0.223) (9.504,-0.298) (9.496,-0.373) (9.486,-0.447) (9.475,-0.521) (9.462,-0.595) (9.447,-0.668) (9.430,-0.742) (9.412,-0.814) (9.392,-0.887) (9.370,-0.958) (9.347,-1.030) (9.322,-1.101) (9.296,-1.171) (9.268,-1.240) (9.238,-1.309) (9.207,-1.377) (9.174,-1.445) (9.140,-1.512) (9.105,-1.578) (9.068,-1.643) (9.029,-1.707) (8.990,-1.771) (8.948,-1.834) (8.906,-1.896) (8.862,-1.956) (8.817,-2.016) (8.771,-2.075) (8.724,-2.134) (8.675,-2.191) (8.625,-2.247) (8.574,-2.302) (8.522,-2.356) (8.469,-2.409) (8.415,-2.461) (8.360,-2.512) (8.304,-2.562) (8.248,-2.611) (8.190,-2.659) (8.131,-2.705) (8.071,-2.751) (8.011,-2.795) (7.950,-2.839) (7.888,-2.881) (7.825,-2.922) (7.762,-2.962) (7.698,-3.001) (7.633,-3.039) (7.567,-3.076) (7.501,-3.111) (7.435,-3.146) (7.368,-3.179) (7.300,-3.211) (7.232,-3.242) (7.163,-3.272) (7.094,-3.301) (7.024,-3.328) (6.954,-3.355) (6.883,-3.380) (6.812,-3.405) (6.741,-3.428) (6.669,-3.450) (6.597,-3.471) (6.525,-3.491) (6.452,-3.509) (6.379,-3.527) (6.306,-3.544) (6.233,-3.559) (6.159,-3.574) (6.085,-3.587) (6.011,-3.599) (5.937,-3.610) (5.863,-3.620) (5.788,-3.629) (5.714,-3.637) (5.639,-3.644) (5.564,-3.650) (5.490,-3.655) (5.415,-3.658) (5.340,-3.661) (5.265,-3.663) (5.190,-3.663) (5.115,-3.663) (5.040,-3.662) (4.965,-3.659) (4.890,-3.656) (4.815,-3.651)}
\LC{(9.552,-10.000) (9.482,-9.976) (9.412,-9.949) (9.342,-9.921) (9.273,-9.894) (9.203,-9.866) (9.134,-9.837) (9.064,-9.809) (8.995,-9.780) (8.926,-9.751) (8.857,-9.722) (8.788,-9.692) (8.719,-9.662) (8.651,-9.632) (8.582,-9.601) (8.514,-9.571) (8.445,-9.539) (8.377,-9.508) (8.309,-9.476) (8.241,-9.444) (8.174,-9.412) (8.106,-9.380) (8.039,-9.347) (7.971,-9.314) (7.904,-9.281) (7.837,-9.247) (7.770,-9.213) (7.703,-9.179) (7.637,-9.144) (7.570,-9.110) (7.504,-9.075) (7.438,-9.039) (7.372,-9.004) (7.306,-8.968) (7.240,-8.932) (7.175,-8.895) (7.109,-8.859) (7.044,-8.822) (6.979,-8.785) (6.914,-8.747) (6.849,-8.709) (6.784,-8.671) (6.720,-8.633) (6.656,-8.594) (6.592,-8.556) (6.528,-8.516) (6.464,-8.477) (6.400,-8.437) (6.337,-8.398) (6.273,-8.357) (6.210,-8.317) (6.147,-8.276) (6.084,-8.235) (6.022,-8.194) (5.959,-8.152) (5.897,-8.110) (5.835,-8.068) (5.773,-8.026) (5.711,-7.983) (5.650,-7.941) (5.588,-7.897) (5.527,-7.854) (5.466,-7.810) (5.405,-7.766) (5.345,-7.722) (5.284,-7.678) (5.224,-7.633) (5.164,-7.588) (5.104,-7.543) (5.045,-7.497) (4.985,-7.452) (4.926,-7.406) (4.867,-7.359) (4.808,-7.313) (4.750,-7.266) (4.691,-7.219) (4.633,-7.172) (4.575,-7.124) (4.517,-7.076) (4.460,-7.028) (4.402,-6.980) (4.345,-6.931) (4.288,-6.882) (4.232,-6.833) (4.175,-6.784) (4.119,-6.734) (4.063,-6.684) (4.007,-6.634) (3.952,-6.584) (3.896,-6.533) (3.841,-6.482) (3.786,-6.431) (3.732,-6.380) (3.677,-6.328) (3.623,-6.276) (3.569,-6.224) (3.515,-6.172) (3.462,-6.119) (3.409,-6.066) (3.356,-6.013) (3.303,-5.960) (3.251,-5.906) (3.199,-5.852) (3.147,-5.798) (3.095,-5.744) (3.044,-5.689) (2.993,-5.634) (2.942,-5.579) (2.891,-5.524) (2.841,-5.468) (2.791,-5.412) (2.741,-5.356) (2.691,-5.300) (2.642,-5.243) (2.593,-5.187) (2.545,-5.130) (2.496,-5.072) (2.448,-5.015) (2.400,-4.957) (2.353,-4.899) (2.306,-4.841) (2.259,-4.782) (2.212,-4.723) (2.166,-4.664) (2.120,-4.605) (2.074,-4.546) (2.029,-4.486) (1.984,-4.426) (1.939,-4.366) (1.894,-4.305) (1.850,-4.245) (1.807,-4.184) (1.763,-4.123) (1.720,-4.061) (1.677,-4.000) (1.635,-3.938) (1.593,-3.876) (1.551,-3.813) (1.510,-3.751) (1.482,-3.709)}
\LC{(-10.000,-5.600) (-9.942,-5.616) (-9.869,-5.631) (-9.795,-5.646) (-9.721,-5.659) (-9.648,-5.673) (-9.574,-5.685) (-9.500,-5.697) (-9.425,-5.708) (-9.351,-5.719) (-9.277,-5.729) (-9.202,-5.738) (-9.128,-5.747) (-9.053,-5.754) (-8.979,-5.762) (-8.904,-5.768) (-8.829,-5.774) (-8.754,-5.779) (-8.679,-5.784) (-8.605,-5.788) (-8.530,-5.791) (-8.455,-5.793) (-8.380,-5.795) (-8.305,-5.796) (-8.230,-5.797) (-8.155,-5.797) (-8.080,-5.796) (-8.005,-5.795) (-7.930,-5.793) (-7.855,-5.790) (-7.780,-5.787) (-7.705,-5.783) (-7.630,-5.779) (-7.555,-5.773) (-7.480,-5.768) (-7.406,-5.761) (-7.331,-5.754) (-7.256,-5.747) (-7.182,-5.738) (-7.107,-5.729) (-7.033,-5.720) (-6.959,-5.710) (-6.885,-5.699) (-6.810,-5.688) (-6.736,-5.676) (-6.662,-5.663) (-6.589,-5.650) (-6.515,-5.636) (-6.441,-5.622) (-6.368,-5.607) (-6.294,-5.591) (-6.221,-5.575) (-6.148,-5.558) (-6.075,-5.541) (-6.002,-5.523) (-5.930,-5.504) (-5.857,-5.485) (-5.785,-5.465) (-5.713,-5.445) (-5.641,-5.424) (-5.569,-5.403) (-5.497,-5.380) (-5.426,-5.358) (-5.354,-5.335) (-5.283,-5.311) (-5.212,-5.287) (-5.141,-5.262) (-5.071,-5.236) (-5.001,-5.210) (-4.930,-5.184) (-4.860,-5.156) (-4.791,-5.129) (-4.721,-5.100) (-4.652,-5.072) (-4.583,-5.042) (-4.514,-5.012) (-4.446,-4.982) (-4.377,-4.951) (-4.309,-4.919) (-4.242,-4.887) (-4.174,-4.854) (-4.107,-4.821) (-4.040,-4.787) (-3.973,-4.753) (-3.907,-4.718) (-3.841,-4.683) (-3.775,-4.647) (-3.709,-4.611) (-3.644,-4.574) (-3.579,-4.536) (-3.514,-4.498) (-3.450,-4.460) (-3.386,-4.421) (-3.322,-4.381) (-3.259,-4.341) (-3.196,-4.301) (-3.133,-4.259) (-3.071,-4.218) (-3.009,-4.176) (-2.947,-4.133) (-2.886,-4.090) (-2.825,-4.046) (-2.764,-4.002) (-2.704,-3.957) (-2.644,-3.912) (-2.584,-3.867) (-2.525,-3.820) (-2.467,-3.774)}
\LC{(10.000,-4.876) (9.939,-4.901) (9.868,-4.925) (9.797,-4.949) (9.725,-4.972) (9.653,-4.993) (9.581,-5.015) (9.509,-5.035) (9.437,-5.055) (9.364,-5.073) (9.292,-5.092) (9.219,-5.109) (9.145,-5.125) (9.072,-5.141) (8.999,-5.156) (8.925,-5.171) (8.851,-5.184) (8.777,-5.197) (8.703,-5.209) (8.629,-5.220) (8.555,-5.231) (8.481,-5.241) (8.406,-5.250) (8.332,-5.258) (8.257,-5.266) (8.182,-5.273) (8.108,-5.279) (8.033,-5.285) (7.958,-5.289) (7.883,-5.293) (7.808,-5.297) (7.733,-5.299) (7.658,-5.301) (7.583,-5.302) (7.508,-5.303) (7.433,-5.302) (7.358,-5.301) (7.283,-5.300) (7.208,-5.297) (7.133,-5.294) (7.058,-5.290) (6.984,-5.286) (6.909,-5.281) (6.834,-5.275) (6.759,-5.268) (6.685,-5.261) (6.610,-5.253) (6.535,-5.245) (6.461,-5.236) (6.387,-5.226) (6.312,-5.215) (6.238,-5.204) (6.164,-5.192) (6.090,-5.179) (6.017,-5.166) (5.943,-5.152) (5.869,-5.137) (5.796,-5.122) (5.723,-5.106) (5.649,-5.090) (5.576,-5.073) (5.504,-5.055) (5.431,-5.036) (5.358,-5.017) (5.286,-4.998) (5.214,-4.977) (5.142,-4.956) (5.070,-4.934) (4.998,-4.912) (4.927,-4.889) (4.856,-4.866) (4.785,-4.842) (4.714,-4.817) (4.643,-4.791) (4.573,-4.765) (4.503,-4.739) (4.433,-4.711) (4.363,-4.684) (4.294,-4.655) (4.248,-4.636)}
\LC{(-0.445,2.886) (-0.462,2.959) (-0.479,3.032) (-0.497,3.105) (-0.515,3.178) (-0.533,3.251) (-0.551,3.323) (-0.570,3.396) (-0.588,3.469) (-0.607,3.541) (-0.626,3.614) (-0.645,3.686) (-0.665,3.759) (-0.684,3.831) (-0.704,3.904) (-0.724,3.976) (-0.744,4.048) (-0.765,4.120) (-0.785,4.192) (-0.806,4.265) (-0.827,4.337) (-0.848,4.409) (-0.869,4.481) (-0.891,4.552) (-0.912,4.624) (-0.934,4.696) (-0.956,4.768) (-0.978,4.839) (-1.000,4.911) (-1.023,4.983) (-1.045,5.054) (-1.068,5.126) (-1.091,5.197) (-1.114,5.268) (-1.137,5.340) (-1.160,5.411) (-1.184,5.482) (-1.208,5.553) (-1.231,5.624) (-1.255,5.695) (-1.280,5.766) (-1.304,5.837) (-1.328,5.908) (-1.353,5.979) (-1.378,6.050) (-1.403,6.121) (-1.428,6.191) (-1.453,6.262) (-1.478,6.333) (-1.504,6.403) (-1.529,6.474) (-1.555,6.544) (-1.581,6.615) (-1.607,6.685) (-1.633,6.755) (-1.659,6.825) (-1.686,6.896) (-1.712,6.966) (-1.739,7.036) (-1.766,7.106) (-1.793,7.176) (-1.820,7.246) (-1.847,7.316) (-1.875,7.385) (-1.902,7.455) (-1.930,7.525) (-1.958,7.595) (-1.986,7.664) (-2.014,7.734) (-2.042,7.803) (-2.070,7.873) (-2.099,7.942) (-2.127,8.011) (-2.156,8.081) (-2.185,8.150) (-2.214,8.219) (-2.243,8.288) (-2.272,8.357) (-2.301,8.426) (-2.331,8.495) (-2.360,8.564) (-2.390,8.633) (-2.420,8.702) (-2.450,8.771) (-2.480,8.840) (-2.510,8.908) (-2.540,8.977) (-2.570,9.045) (-2.601,9.114) (-2.632,9.182) (-2.662,9.251) (-2.693,9.319) (-2.724,9.387) (-2.755,9.456) (-2.786,9.524) (-2.818,9.592) (-2.849,9.660) (-2.881,9.728) (-2.912,9.796) (-2.944,9.864) (-2.976,9.932) (-3.008,10.000) (-3.019,10.000)}
\LC{(-10.000,-9.569) (-9.950,-9.548) (-9.878,-9.527) (-9.807,-9.505) (-9.735,-9.484) (-9.663,-9.461) (-9.592,-9.439) (-9.520,-9.416) (-9.449,-9.393) (-9.378,-9.369) (-9.307,-9.346) (-9.236,-9.321) (-9.165,-9.297) (-9.094,-9.272) (-9.023,-9.247) (-8.953,-9.222) (-8.882,-9.196) (-8.812,-9.170) (-8.742,-9.143) (-8.672,-9.117) (-8.602,-9.090) (-8.532,-9.062) (-8.462,-9.035) (-8.393,-9.007) (-8.323,-8.978) (-8.254,-8.950) (-8.185,-8.921) (-8.116,-8.892) (-8.047,-8.862) (-7.978,-8.832) (-7.909,-8.802) (-7.841,-8.771) (-7.772,-8.741) (-7.704,-8.709) (-7.636,-8.678) (-7.568,-8.646) (-7.500,-8.614) (-7.432,-8.582) (-7.365,-8.549) (-7.298,-8.516) (-7.230,-8.483) (-7.163,-8.449) (-7.096,-8.415) (-7.030,-8.381) (-6.963,-8.347) (-6.897,-8.312) (-6.830,-8.277) (-6.764,-8.241) (-6.698,-8.206) (-6.633,-8.170) (-6.567,-8.133) (-6.501,-8.097) (-6.436,-8.060) (-6.371,-8.023) (-6.306,-7.985) (-6.241,-7.947) (-6.177,-7.909) (-6.112,-7.871) (-6.048,-7.832) (-5.984,-7.793) (-5.920,-7.754) (-5.857,-7.714) (-5.793,-7.674) (-5.730,-7.634) (-5.667,-7.593) (-5.604,-7.553) (-5.541,-7.512) (-5.478,-7.470) (-5.416,-7.429) (-5.354,-7.387) (-5.292,-7.344) (-5.230,-7.302) (-5.168,-7.259) (-5.107,-7.216) (-5.046,-7.173) (-4.985,-7.129) (-4.924,-7.085) (-4.864,-7.041) (-4.803,-6.996) (-4.743,-6.951) (-4.683,-6.906) (-4.624,-6.861) (-4.564,-6.815) (-4.505,-6.769) (-4.446,-6.723) (-4.387,-6.676) (-4.328,-6.629) (-4.270,-6.582) (-4.212,-6.535) (-4.154,-6.487) (-4.096,-6.439) (-4.039,-6.391) (-3.982,-6.342) (-3.925,-6.294) (-3.868,-6.245) (-3.812,-6.195) (-3.755,-6.146) (-3.699,-6.096) (-3.644,-6.045)}
\LC{(-2.919,-10.000) (-2.888,-9.932) (-2.857,-9.864) (-2.826,-9.795) (-2.795,-9.727) (-2.764,-9.658) (-2.734,-9.590) (-2.703,-9.521) (-2.673,-9.453) (-2.642,-9.384) (-2.612,-9.316) (-2.582,-9.247) (-2.552,-9.178) (-2.523,-9.109) (-2.493,-9.040) (-2.463,-8.971) (-2.434,-8.902) (-2.405,-8.833) (-2.375,-8.764) (-2.346,-8.695) (-2.317,-8.626) (-2.288,-8.557) (-2.260,-8.488) (-2.231,-8.418) (-2.203,-8.349) (-2.174,-8.279) (-2.146,-8.210) (-2.118,-8.140) (-2.090,-8.071) (-2.062,-8.001) (-2.034,-7.932) (-2.006,-7.862) (-1.979,-7.792) (-1.951,-7.722) (-1.924,-7.652) (-1.897,-7.583) (-1.870,-7.513) (-1.843,-7.443) (-1.816,-7.373) (-1.790,-7.302) (-1.763,-7.232) (-1.737,-7.162) (-1.710,-7.092) (-1.684,-7.021) (-1.658,-6.951) (-1.632,-6.881) (-1.607,-6.810) (-1.581,-6.740) (-1.556,-6.669) (-1.530,-6.599) (-1.505,-6.528) (-1.480,-6.457) (-1.455,-6.387) (-1.431,-6.316) (-1.406,-6.245) (-1.381,-6.174) (-1.357,-6.103) (-1.333,-6.032) (-1.309,-5.961) (-1.285,-5.890) (-1.261,-5.819) (-1.238,-5.748) (-1.214,-5.676) (-1.191,-5.605) (-1.168,-5.534) (-1.145,-5.462) (-1.122,-5.391) (-1.099,-5.320) (-1.076,-5.248) (-1.054,-5.176) (-1.032,-5.105) (-1.010,-5.033) (-0.988,-4.961) (-0.966,-4.890) (-0.944,-4.818) (-0.923,-4.746) (-0.902,-4.674) (-0.880,-4.602) (-0.859,-4.530) (-0.839,-4.458) (-0.818,-4.386)}
\LC{(2.151,2.078) (2.223,2.099) (2.296,2.118) (2.368,2.135) (2.442,2.152) (2.515,2.166) (2.589,2.179) (2.663,2.191) (2.738,2.200) (2.812,2.209) (2.887,2.215) (2.962,2.220) (3.037,2.223) (3.112,2.225) (3.187,2.225) (3.262,2.223) (3.336,2.220) (3.411,2.215) (3.486,2.208) (3.560,2.199) (3.635,2.189) (3.709,2.176) (3.782,2.163) (3.856,2.147) (3.929,2.129) (4.001,2.110) (4.073,2.089) (4.144,2.066) (4.215,2.041) (4.285,2.014) (4.355,1.986) (4.423,1.955) (4.491,1.923) (4.558,1.889) (4.624,1.853) (4.688,1.815) (4.752,1.776) (4.815,1.734) (4.876,1.691) (4.936,1.646) (4.994,1.599) (5.051,1.550) (5.107,1.500) (5.160,1.447) (5.213,1.393) (5.263,1.338) (5.311,1.281) (5.358,1.222) (5.402,1.161) (5.445,1.099) (5.485,1.036) (5.523,0.971) (5.558,0.905) (5.592,0.838) (5.622,0.770) (5.651,0.700) (5.676,0.630) (5.699,0.558) (5.719,0.486) (5.736,0.413) (5.751,0.340) (5.763,0.266) (5.772,0.191) (5.778,0.116) (5.781,0.041) (5.781,-0.034) (5.778,-0.109) (5.772,-0.183) (5.764,-0.258) (5.752,-0.332) (5.738,-0.406) (5.721,-0.479) (5.701,-0.551) (5.679,-0.622) (5.653,-0.693) (5.625,-0.763) (5.595,-0.831) (5.562,-0.899) (5.527,-0.965) (5.489,-1.030) (5.449,-1.093) (5.407,-1.155) (5.363,-1.216) (5.316,-1.275) (5.268,-1.332) (5.218,-1.388) (5.166,-1.442) (5.112,-1.494) (5.057,-1.545) (5.000,-1.594) (4.942,-1.641) (4.882,-1.686) (4.821,-1.730) (4.759,-1.771) (4.695,-1.811) (4.630,-1.849) (4.565,-1.885) (4.498,-1.920) (4.430,-1.952) (4.362,-1.983) (4.293,-2.011) (4.222,-2.038) (4.152,-2.063) (4.080,-2.086) (4.009,-2.108) (3.936,-2.127) (3.863,-2.145) (3.790,-2.161) (3.716,-2.175) (3.642,-2.187) (3.568,-2.198) (3.494,-2.207) (3.419,-2.214) (3.344,-2.219) (3.269,-2.223) (3.194,-2.225) (3.119,-2.225) (3.044,-2.224) (2.969,-2.220) (2.895,-2.216) (2.820,-2.209) (2.745,-2.201) (2.671,-2.192) (2.597,-2.180) (2.523,-2.167) (2.449,-2.153) (2.376,-2.137) (2.303,-2.120) (2.230,-2.101) (2.158,-2.080)}
\LC{(-0.002,2.978) (-0.002,3.053) (-0.002,3.128) (-0.002,3.203) (-0.003,3.278) (-0.003,3.353) (-0.003,3.428) (-0.003,3.503) (-0.003,3.578) (-0.003,3.653) (-0.003,3.728) (-0.003,3.803) (-0.003,3.878) (-0.003,3.953) (-0.003,4.028) (-0.004,4.103) (-0.004,4.178) (-0.004,4.253) (-0.004,4.328) (-0.004,4.403) (-0.004,4.478) (-0.004,4.553) (-0.004,4.628) (-0.004,4.703) (-0.004,4.778) (-0.005,4.853) (-0.005,4.928) (-0.005,5.003) (-0.005,5.078) (-0.005,5.153) (-0.005,5.228) (-0.005,5.303) (-0.005,5.378) (-0.005,5.453) (-0.006,5.528) (-0.006,5.603) (-0.006,5.678) (-0.006,5.753) (-0.006,5.828) (-0.006,5.903) (-0.006,5.978) (-0.006,6.053) (-0.006,6.128) (-0.007,6.203) (-0.007,6.278) (-0.007,6.353) (-0.007,6.428) (-0.007,6.503) (-0.007,6.578) (-0.007,6.653) (-0.007,6.728) (-0.008,6.803) (-0.008,6.878) (-0.008,6.953) (-0.008,7.028) (-0.008,7.103) (-0.008,7.178) (-0.008,7.253) (-0.008,7.328) (-0.009,7.403) (-0.009,7.478) (-0.009,7.553) (-0.009,7.628) (-0.009,7.703) (-0.009,7.778) (-0.009,7.853) (-0.010,7.928) (-0.010,8.003) (-0.010,8.078) (-0.010,8.153) (-0.010,8.228) (-0.010,8.303) (-0.010,8.378) (-0.011,8.453) (-0.011,8.528) (-0.011,8.603) (-0.011,8.678) (-0.011,8.753) (-0.011,8.828) (-0.011,8.903) (-0.012,8.978) (-0.012,9.053) (-0.012,9.128) (-0.012,9.203) (-0.012,9.278) (-0.012,9.353) (-0.012,9.428) (-0.013,9.503) (-0.013,9.578) (-0.013,9.653) (-0.013,9.728) (-0.013,9.803) (-0.013,9.878) (-0.013,9.953) (-0.014,10.000)}
\LC{(4.675,4.277) (4.748,4.294) (4.821,4.311) (4.894,4.327) (4.967,4.342) (5.041,4.357) (5.115,4.370) (5.189,4.383) (5.263,4.395) (5.337,4.406) (5.411,4.417) (5.485,4.427) (5.560,4.435) (5.634,4.444) (5.709,4.451) (5.784,4.457) (5.859,4.463) (5.933,4.468) (6.008,4.472) (6.083,4.475) (6.158,4.477) (6.233,4.479) (6.308,4.480) (6.383,4.480) (6.458,4.479) (6.533,4.477) (6.608,4.474) (6.683,4.471) (6.758,4.466) (6.833,4.461) (6.907,4.455) (6.982,4.448) (7.057,4.440) (7.131,4.432) (7.206,4.422) (7.280,4.412) (7.354,4.401) (7.428,4.388) (7.502,4.375) (7.576,4.361) (7.649,4.347) (7.722,4.331) (7.796,4.314) (7.869,4.297) (7.941,4.278) (8.014,4.259) (8.086,4.239) (8.158,4.218) (8.230,4.196) (8.301,4.173) (8.372,4.149) (8.443,4.124) (8.513,4.099) (8.584,4.072) (8.653,4.045) (8.723,4.016) (8.792,3.987) (8.860,3.957) (8.929,3.926) (8.996,3.893) (9.064,3.860) (9.131,3.827) (9.197,3.792) (9.263,3.756) (9.328,3.719) (9.393,3.682) (9.458,3.643) (9.521,3.604) (9.585,3.563) (9.647,3.522) (9.709,3.480) (9.771,3.437) (9.831,3.393) (9.891,3.348) (9.951,3.302) (10.000,3.255)}
\LC{(-1.771,2.911) (-1.827,2.961) (-1.884,3.010) (-1.941,3.058) (-1.999,3.106) (-2.057,3.154) (-2.116,3.200) (-2.175,3.246) (-2.234,3.292) (-2.295,3.337) (-2.355,3.381) (-2.416,3.425) (-2.477,3.468) (-2.539,3.511) (-2.601,3.552) (-2.664,3.594) (-2.727,3.634) (-2.790,3.674) (-2.854,3.714) (-2.918,3.753) (-2.983,3.791) (-3.048,3.828) (-3.113,3.865) (-3.179,3.902) (-3.245,3.937) (-3.311,3.972) (-3.378,4.007) (-3.445,4.040) (-3.512,4.074) (-3.580,4.106) (-3.648,4.138) (-3.716,4.169) (-3.784,4.199) (-3.853,4.229) (-3.922,4.258) (-3.992,4.287) (-4.061,4.315) (-4.131,4.342) (-4.201,4.369) (-4.272,4.394) (-4.342,4.420) (-4.413,4.444) (-4.484,4.468) (-4.556,4.491) (-4.627,4.514) (-4.699,4.535) (-4.771,4.556) (-4.843,4.577) (-4.915,4.596) (-4.988,4.615) (-5.061,4.634) (-5.134,4.651) (-5.207,4.668) (-5.280,4.684) (-5.353,4.700) (-5.427,4.715) (-5.501,4.729) (-5.574,4.742) (-5.648,4.754) (-5.722,4.766) (-5.797,4.777) (-5.871,4.788) (-5.945,4.798) (-6.020,4.806) (-6.094,4.815) (-6.169,4.822) (-6.244,4.829) (-6.318,4.835) (-6.393,4.840) (-6.468,4.844) (-6.543,4.848) (-6.618,4.851) (-6.693,4.853) (-6.768,4.855) (-6.843,4.855) (-6.918,4.855) (-6.993,4.854) (-7.068,4.852) (-7.143,4.850) (-7.218,4.847) (-7.293,4.843) (-7.367,4.838) (-7.442,4.832) (-7.517,4.826) (-7.592,4.819) (-7.666,4.811) (-7.741,4.802) (-7.815,4.792) (-7.889,4.782) (-7.963,4.770) (-8.037,4.758) (-8.111,4.746) (-8.185,4.732) (-8.259,4.717) (-8.332,4.702) (-8.405,4.686) (-8.478,4.669) (-8.551,4.651) (-8.624,4.633) (-8.696,4.613) (-8.769,4.593) (-8.840,4.572) (-8.912,4.550) (-8.984,4.527) (-9.055,4.503) (-9.126,4.479) (-9.196,4.454) (-9.267,4.427) (-9.337,4.400) (-9.406,4.372) (-9.475,4.344) (-9.544,4.314) (-9.613,4.284) (-9.681,4.253) (-9.749,4.220) (-9.816,4.187) (-9.883,4.154) (-9.950,4.119) (-10.000,4.084)}
\LC{(2.995,4.844) (3.051,4.893) (3.108,4.942) (3.166,4.990) (3.223,5.038) (3.281,5.086) (3.339,5.133) (3.398,5.180) (3.456,5.227) (3.516,5.273) (3.575,5.319) (3.635,5.364) (3.694,5.410) (3.755,5.454) (3.815,5.499) (3.876,5.543) (3.937,5.586) (3.998,5.630) (4.060,5.673) (4.121,5.715) (4.183,5.757) (4.246,5.799) (4.308,5.840) (4.371,5.881) (4.434,5.922) (4.497,5.962) (4.561,6.002) (4.624,6.042) (4.688,6.081) (4.753,6.120) (4.817,6.158) (4.882,6.196) (4.947,6.234) (5.012,6.271) (5.077,6.308) (5.143,6.344) (5.208,6.380) (5.274,6.416) (5.341,6.451) (5.407,6.486) (5.474,6.521) (5.540,6.555) (5.607,6.588) (5.675,6.622) (5.742,6.655) (5.809,6.687) (5.877,6.719) (5.945,6.751) (6.013,6.782) (6.082,6.813) (6.150,6.844) (6.219,6.874) (6.288,6.904) (6.357,6.933) (6.426,6.962) (6.495,6.991) (6.565,7.019) (6.634,7.046) (6.704,7.074) (6.774,7.101) (6.845,7.127) (6.915,7.153) (6.985,7.179) (7.056,7.204) (7.127,7.229) (7.198,7.253) (7.269,7.277) (7.340,7.301) (7.411,7.324) (7.483,7.347) (7.554,7.369) (7.626,7.391) (7.698,7.413) (7.770,7.434) (7.842,7.454) (7.914,7.474) (7.987,7.494) (8.059,7.514) (8.132,7.533) (8.204,7.551) (8.277,7.569) (8.350,7.587) (8.423,7.604) (8.496,7.621) (8.569,7.637) (8.642,7.653) (8.716,7.668) (8.789,7.683) (8.863,7.698) (8.937,7.712) (9.010,7.726) (9.084,7.739) (9.158,7.752) (9.232,7.764) (9.306,7.776) (9.380,7.788) (9.454,7.799) (9.529,7.810) (9.603,7.820) (9.677,7.829) (9.752,7.839) (9.826,7.847) (9.901,7.856) (9.975,7.864) (10.000,7.866)}
\LC{(-1.118,1.217) (-1.188,1.245) (-1.258,1.270) (-1.330,1.293) (-1.402,1.314) (-1.474,1.332) (-1.548,1.348) (-1.622,1.362) (-1.696,1.372) (-1.770,1.380) (-1.845,1.386) (-1.920,1.389) (-1.995,1.389) (-2.070,1.387) (-2.145,1.382) (-2.219,1.374) (-2.294,1.363) (-2.367,1.349) (-2.440,1.333) (-2.513,1.313) (-2.585,1.291) (-2.655,1.266) (-2.725,1.238) (-2.793,1.207) (-2.860,1.172) (-2.925,1.135) (-2.988,1.095) (-3.050,1.052) (-3.109,1.006) (-3.166,0.958) (-3.220,0.906) (-3.272,0.852) (-3.321,0.795) (-3.366,0.735) (-3.409,0.673) (-3.447,0.609) (-3.482,0.543) (-3.514,0.474) (-3.541,0.405) (-3.563,0.333) (-3.582,0.260) (-3.596,0.187) (-3.605,0.112) (-3.609,0.037) (-3.609,-0.038) (-3.605,-0.112) (-3.596,-0.187) (-3.582,-0.261) (-3.563,-0.333) (-3.540,-0.405) (-3.513,-0.475) (-3.482,-0.543) (-3.447,-0.609) (-3.409,-0.673) (-3.366,-0.735) (-3.321,-0.795) (-3.272,-0.852) (-3.220,-0.906) (-3.166,-0.958) (-3.109,-1.006) (-3.049,-1.052) (-2.988,-1.095) (-2.925,-1.135) (-2.859,-1.172) (-2.793,-1.207) (-2.724,-1.238) (-2.655,-1.266) (-2.584,-1.291) (-2.513,-1.313) (-2.440,-1.333) (-2.367,-1.349) (-2.293,-1.363) (-2.219,-1.374) (-2.145,-1.382) (-2.070,-1.387) (-1.995,-1.389) (-1.920,-1.389) (-1.845,-1.386) (-1.770,-1.380) (-1.696,-1.372) (-1.621,-1.361) (-1.548,-1.348) (-1.474,-1.332) (-1.401,-1.314) (-1.329,-1.293) (-1.258,-1.270) (-1.188,-1.245) (-1.118,-1.217)}
\LC{(-10.000,-2.969) (-9.946,-3.019) (-9.890,-3.069) (-9.833,-3.118) (-9.775,-3.165) (-9.716,-3.212) (-9.657,-3.258) (-9.597,-3.303) (-9.537,-3.347) (-9.475,-3.391) (-9.413,-3.433) (-9.351,-3.474) (-9.287,-3.515) (-9.224,-3.554) (-9.159,-3.593) (-9.094,-3.630) (-9.029,-3.667) (-8.963,-3.702) (-8.897,-3.737) (-8.830,-3.771) (-8.762,-3.804) (-8.694,-3.836) (-8.626,-3.867) (-8.557,-3.897) (-8.488,-3.926) (-8.419,-3.954) (-8.349,-3.981) (-8.278,-4.007) (-8.208,-4.033) (-8.137,-4.057) (-8.066,-4.080) (-7.994,-4.103) (-7.922,-4.124) (-7.850,-4.145) (-7.778,-4.165) (-7.705,-4.184) (-7.632,-4.201) (-7.559,-4.218) (-7.486,-4.234) (-7.413,-4.249) (-7.339,-4.264) (-7.265,-4.277) (-7.191,-4.289) (-7.117,-4.301) (-7.043,-4.311) (-6.968,-4.321) (-6.894,-4.330) (-6.819,-4.338) (-6.745,-4.345) (-6.670,-4.351) (-6.595,-4.356) (-6.520,-4.361) (-6.445,-4.364) (-6.370,-4.367) (-6.295,-4.368) (-6.220,-4.369) (-6.145,-4.369) (-6.070,-4.369) (-5.995,-4.367) (-5.920,-4.364) (-5.846,-4.361) (-5.771,-4.357) (-5.696,-4.352) (-5.621,-4.346) (-5.546,-4.339) (-5.472,-4.332) (-5.397,-4.323) (-5.323,-4.314) (-5.248,-4.304) (-5.174,-4.293) (-5.100,-4.282) (-5.051,-4.273)}
\LC{(-2.119,1.838) (-2.193,1.851) (-2.267,1.863) (-2.341,1.872) (-2.416,1.880) (-2.491,1.886) (-2.566,1.890) (-2.641,1.892) (-2.716,1.892) (-2.791,1.890) (-2.866,1.885) (-2.940,1.879) (-3.015,1.871) (-3.089,1.861) (-3.163,1.849) (-3.237,1.834) (-3.310,1.818) (-3.383,1.799) (-3.455,1.778) (-3.526,1.756) (-3.597,1.730) (-3.667,1.703) (-3.736,1.674) (-3.804,1.642) (-3.871,1.608) (-3.936,1.573) (-4.001,1.534) (-4.064,1.494) (-4.126,1.452) (-4.186,1.407) (-4.245,1.360) (-4.302,1.311) (-4.357,1.261) (-4.410,1.208) (-4.461,1.153) (-4.510,1.096) (-4.557,1.037) (-4.601,0.977) (-4.643,0.914) (-4.682,0.850) (-4.718,0.785) (-4.752,0.718) (-4.782,0.649) (-4.810,0.580) (-4.835,0.509) (-4.856,0.437) (-4.874,0.364) (-4.889,0.291) (-4.901,0.216) (-4.909,0.142) (-4.914,0.067) (-4.915,-0.008) (-4.913,-0.083) (-4.908,-0.158) (-4.899,-0.232) (-4.887,-0.306) (-4.871,-0.379) (-4.852,-0.452) (-4.830,-0.524) (-4.805,-0.594) (-4.776,-0.664) (-4.745,-0.732) (-4.711,-0.799) (-4.674,-0.864) (-4.634,-0.927) (-4.592,-0.989) (-4.547,-1.050) (-4.500,-1.108) (-4.451,-1.164) (-4.399,-1.219) (-4.346,-1.271) (-4.290,-1.322) (-4.233,-1.370) (-4.174,-1.417) (-4.113,-1.461) (-4.051,-1.503) (-3.987,-1.543) (-3.923,-1.580) (-3.857,-1.616) (-3.789,-1.649) (-3.721,-1.680) (-3.652,-1.709) (-3.582,-1.736) (-3.511,-1.761) (-3.440,-1.783) (-3.367,-1.803) (-3.295,-1.821) (-3.221,-1.838) (-3.148,-1.851) (-3.074,-1.863) (-2.999,-1.873) (-2.925,-1.881) (-2.850,-1.887) (-2.775,-1.890) (-2.700,-1.892) (-2.625,-1.891) (-2.550,-1.889) (-2.475,-1.885) (-2.400,-1.879) (-2.326,-1.871) (-2.251,-1.861) (-2.177,-1.849) (-2.128,-1.840)}
\LC{(-4.280,5.092) (-4.347,5.126) (-4.414,5.159) (-4.482,5.192) (-4.549,5.225) (-4.617,5.257) (-4.685,5.288) (-4.754,5.319) (-4.822,5.349) (-4.891,5.379) (-4.960,5.409) (-5.029,5.437) (-5.099,5.466) (-5.168,5.494) (-5.238,5.521) (-5.308,5.548) (-5.379,5.574) (-5.449,5.600) (-5.520,5.625) (-5.590,5.650) (-5.661,5.674) (-5.733,5.697) (-5.804,5.720) (-5.876,5.743) (-5.947,5.765) (-6.019,5.786) (-6.091,5.807) (-6.163,5.827) (-6.236,5.847) (-6.308,5.866) (-6.381,5.885) (-6.454,5.903) (-6.527,5.921) (-6.600,5.938) (-6.673,5.954) (-6.746,5.970) (-6.819,5.986) (-6.893,6.001) (-6.967,6.015) (-7.040,6.028) (-7.114,6.042) (-7.188,6.054) (-7.262,6.066) (-7.336,6.078) (-7.410,6.088) (-7.485,6.099) (-7.559,6.108) (-7.634,6.117) (-7.708,6.126) (-7.783,6.134) (-7.857,6.141) (-7.932,6.148) (-8.007,6.154) (-8.082,6.160) (-8.156,6.165) (-8.231,6.169) (-8.306,6.173) (-8.381,6.176) (-8.456,6.179) (-8.531,6.181) (-8.606,6.182) (-8.681,6.183) (-8.756,6.183) (-8.831,6.183) (-8.906,6.182) (-8.981,6.180) (-9.056,6.178) (-9.131,6.175) (-9.206,6.172) (-9.281,6.168) (-9.356,6.163) (-9.430,6.158) (-9.505,6.152) (-9.580,6.145) (-9.654,6.138) (-9.729,6.130) (-9.804,6.122) (-9.878,6.113) (-9.952,6.103) (-10.000,6.096)}
\LC{(1.461,2.722) (1.513,2.776) (1.565,2.830) (1.618,2.883) (1.671,2.936) (1.725,2.988) (1.779,3.040) (1.834,3.091) (1.890,3.142) (1.945,3.192) (2.002,3.242) (2.058,3.291) (2.115,3.339) (2.173,3.387) (2.231,3.435) (2.290,3.482) (2.349,3.528) (2.408,3.574) (2.468,3.619) (2.528,3.664) (2.588,3.708) (2.649,3.752) (2.711,3.795) (2.773,3.837) (2.835,3.879) (2.897,3.921) (2.960,3.962) (3.024,4.002) (3.087,4.042) (3.151,4.081) (3.215,4.119) (3.280,4.157) (3.345,4.195) (3.410,4.232) (3.476,4.268) (3.542,4.304) (3.608,4.339) (3.675,4.373) (3.742,4.407) (3.809,4.441) (3.876,4.473) (3.944,4.506) (4.012,4.537) (4.080,4.568) (4.149,4.599) (4.218,4.628) (4.287,4.658) (4.356,4.686) (4.426,4.714) (4.496,4.741) (4.566,4.768) (4.636,4.794) (4.706,4.820) (4.777,4.845) (4.848,4.869) (4.919,4.893) (4.991,4.916) (5.062,4.938) (5.134,4.960) (5.206,4.981) (5.278,5.002) (5.350,5.022) (5.423,5.041) (5.496,5.060) (5.568,5.078) (5.641,5.095) (5.714,5.112) (5.788,5.128) (5.861,5.143) (5.935,5.158) (6.008,5.172) (6.082,5.185) (6.156,5.198) (6.230,5.210) (6.304,5.222) (6.378,5.233) (6.453,5.243) (6.527,5.252) (6.602,5.261) (6.676,5.269) (6.751,5.276) (6.825,5.283) (6.900,5.289) (6.975,5.294) (7.050,5.299) (7.125,5.303) (7.200,5.306) (7.275,5.309) (7.350,5.311) (7.425,5.312) (7.500,5.312) (7.575,5.312) (7.650,5.311) (7.725,5.310) (7.800,5.307) (7.874,5.304) (7.949,5.300) (8.024,5.296) (8.099,5.291) (8.174,5.285) (8.249,5.278) (8.323,5.271) (8.398,5.262) (8.472,5.254) (8.547,5.244) (8.621,5.234) (8.695,5.223) (8.769,5.211) (8.843,5.198) (8.917,5.185) (8.990,5.171) (9.064,5.156) (9.137,5.140) (9.211,5.124) (9.284,5.107) (9.356,5.089) (9.429,5.071) (9.502,5.051) (9.574,5.031) (9.646,5.010) (9.718,4.989) (9.789,4.966) (9.861,4.943) (9.932,4.919) (10.000,4.894)}
\LC{(1.369,4.863) (1.399,4.932) (1.429,5.000) (1.460,5.069) (1.491,5.137) (1.522,5.205) (1.553,5.273) (1.585,5.341) (1.617,5.409) (1.649,5.477) (1.681,5.545) (1.713,5.613) (1.746,5.680) (1.779,5.748) (1.812,5.815) (1.845,5.882) (1.878,5.949) (1.912,6.017) (1.945,6.084) (1.979,6.150) (2.013,6.217) (2.048,6.284) (2.082,6.350) (2.117,6.417) (2.152,6.483) (2.187,6.550) (2.222,6.616) (2.258,6.682) (2.293,6.748) (2.329,6.814) (2.365,6.880) (2.401,6.945) (2.437,7.011) (2.474,7.077) (2.510,7.142) (2.547,7.207) (2.584,7.273) (2.621,7.338) (2.659,7.403) (2.696,7.468) (2.734,7.533) (2.772,7.597) (2.810,7.662) (2.848,7.727) (2.886,7.791) (2.925,7.855) (2.964,7.920) (3.002,7.984) (3.041,8.048) (3.081,8.112) (3.120,8.176) (3.159,8.239) (3.199,8.303) (3.239,8.367) (3.279,8.430) (3.319,8.494) (3.359,8.557) (3.399,8.620) (3.440,8.683) (3.481,8.746) (3.522,8.809) (3.563,8.872) (3.604,8.934) (3.645,8.997) (3.687,9.059) (3.728,9.122) (3.770,9.184) (3.812,9.246) (3.854,9.308) (3.896,9.370) (3.939,9.432) (3.981,9.494) (4.024,9.556) (4.067,9.617) (4.110,9.679) (4.153,9.740) (4.196,9.802) (4.239,9.863) (4.283,9.924) (4.326,9.985) (4.341,10.000)}
\LC{(-1.407,4.702) (-1.439,4.770) (-1.472,4.837) (-1.504,4.905) (-1.537,4.973) (-1.570,5.040) (-1.603,5.107) (-1.636,5.174) (-1.670,5.241) (-1.704,5.308) (-1.738,5.375) (-1.772,5.442) (-1.806,5.508) (-1.841,5.575) (-1.876,5.641) (-1.911,5.708) (-1.947,5.774) (-1.982,5.840) (-2.018,5.906) (-2.054,5.972) (-2.090,6.037) (-2.126,6.103) (-2.163,6.168) (-2.200,6.234) (-2.237,6.299) (-2.274,6.364) (-2.311,6.429) (-2.349,6.494) (-2.386,6.559) (-2.424,6.624) (-2.463,6.688) (-2.501,6.753) (-2.539,6.817) (-2.578,6.881) (-2.617,6.945) (-2.656,7.009) (-2.695,7.073) (-2.735,7.137) (-2.774,7.201) (-2.814,7.264) (-2.854,7.328) (-2.894,7.391) (-2.935,7.454) (-2.975,7.517) (-3.016,7.580) (-3.057,7.643) (-3.098,7.706) (-3.139,7.769) (-3.180,7.831) (-3.222,7.894) (-3.264,7.956) (-3.306,8.018) (-3.348,8.080) (-3.390,8.142) (-3.432,8.204) (-3.475,8.266) (-3.518,8.327) (-3.561,8.389) (-3.604,8.450) (-3.647,8.512) (-3.690,8.573) (-3.734,8.634) (-3.778,8.695) (-3.822,8.756) (-3.866,8.816) (-3.910,8.877) (-3.954,8.937) (-3.999,8.998) (-4.043,9.058) (-4.088,9.118) (-4.133,9.178) (-4.178,9.238) (-4.224,9.298) (-4.269,9.357) (-4.315,9.417) (-4.360,9.476) (-4.406,9.536) (-4.452,9.595) (-4.499,9.654) (-4.545,9.713) (-4.591,9.772) (-4.638,9.831) (-4.685,9.889) (-4.732,9.948) (-4.779,10.000)}
\LC{(0.449,3.419) (0.464,3.492) (0.479,3.566) (0.494,3.639) (0.509,3.713) (0.524,3.786) (0.540,3.860) (0.556,3.933) (0.571,4.006) (0.587,4.080) (0.603,4.153) (0.619,4.226) (0.636,4.299) (0.652,4.373) (0.669,4.446) (0.686,4.519) (0.703,4.592) (0.720,4.665) (0.737,4.738) (0.754,4.811) (0.772,4.884) (0.789,4.957) (0.807,5.030) (0.825,5.102) (0.843,5.175) (0.861,5.248) (0.879,5.321) (0.897,5.394) (0.916,5.466) (0.934,5.539) (0.953,5.612) (0.972,5.684) (0.991,5.757) (1.010,5.829) (1.029,5.902) (1.048,5.974) (1.068,6.047) (1.087,6.119) (1.107,6.191) (1.127,6.264) (1.147,6.336) (1.167,6.408) (1.187,6.481) (1.207,6.553) (1.227,6.625) (1.248,6.697) (1.268,6.769) (1.289,6.841) (1.310,6.913) (1.331,6.986) (1.352,7.057) (1.373,7.129) (1.394,7.201) (1.415,7.273) (1.437,7.345) (1.458,7.417) (1.480,7.489) (1.502,7.561) (1.524,7.632) (1.546,7.704) (1.568,7.776) (1.590,7.847) (1.612,7.919) (1.635,7.991) (1.657,8.062) (1.680,8.134) (1.702,8.205) (1.725,8.277) (1.748,8.348) (1.771,8.419) (1.794,8.491) (1.817,8.562) (1.841,8.633) (1.864,8.705) (1.887,8.776) (1.911,8.847) (1.935,8.918) (1.958,8.989) (1.982,9.060) (2.006,9.132) (2.030,9.203) (2.054,9.274) (2.079,9.345) (2.103,9.416) (2.127,9.486) (2.152,9.557) (2.176,9.628) (2.201,9.699) (2.226,9.770) (2.251,9.841) (2.276,9.911) (2.301,9.982) (2.309,10.000)}
\LC{(10.000,-9.013) (9.929,-8.996) (9.856,-8.978) (9.784,-8.960) (9.711,-8.942) (9.638,-8.924) (9.565,-8.905) (9.493,-8.886) (9.421,-8.866) (9.348,-8.846) (9.276,-8.826) (9.204,-8.805) (9.132,-8.784) (9.060,-8.763) (8.988,-8.741) (8.917,-8.719) (8.845,-8.696) (8.774,-8.674) (8.702,-8.650) (8.631,-8.627) (8.560,-8.603) (8.489,-8.579) (8.418,-8.554) (8.347,-8.529) (8.277,-8.504) (8.206,-8.478) (8.136,-8.453) (8.066,-8.426) (7.996,-8.400) (7.926,-8.373) (7.856,-8.345) (7.786,-8.318) (7.717,-8.289) (7.647,-8.261) (7.578,-8.232) (7.509,-8.203) (7.440,-8.174) (7.371,-8.144) (7.302,-8.114) (7.234,-8.084) (7.165,-8.053) (7.097,-8.022) (7.029,-7.990) (6.961,-7.959) (6.893,-7.927) (6.825,-7.894) (6.758,-7.861) (6.691,-7.828) (6.624,-7.795) (6.557,-7.761) (6.490,-7.727) (6.423,-7.692) (6.357,-7.658) (6.290,-7.623) (6.224,-7.587) (6.158,-7.551) (6.093,-7.515) (6.027,-7.479) (5.962,-7.442) (5.897,-7.405) (5.832,-7.368) (5.767,-7.330) (5.702,-7.292) (5.638,-7.254) (5.573,-7.215) (5.509,-7.176) (5.445,-7.137) (5.382,-7.097) (5.318,-7.057) (5.255,-7.017) (5.192,-6.976) (5.129,-6.935) (5.067,-6.894) (5.004,-6.852) (4.942,-6.810) (4.880,-6.768) (4.818,-6.726) (4.757,-6.683) (4.716,-6.654)}
\LC{(4.277,5.274) (4.342,5.310) (4.408,5.346) (4.475,5.381) (4.541,5.415) (4.608,5.449) (4.675,5.483) (4.742,5.516) (4.810,5.549) (4.878,5.581) (4.946,5.613) (5.014,5.644) (5.082,5.675) (5.151,5.705) (5.220,5.735) (5.289,5.764) (5.358,5.793) (5.427,5.822) (5.497,5.850) (5.567,5.877) (5.637,5.904) (5.707,5.930) (5.777,5.956) (5.848,5.982) (5.919,6.006) (5.989,6.031) (6.061,6.055) (6.132,6.078) (6.203,6.101) (6.275,6.123) (6.347,6.145) (6.418,6.167) (6.490,6.188) (6.563,6.208) (6.635,6.228) (6.707,6.247) (6.780,6.266) (6.853,6.284) (6.926,6.302) (6.999,6.319) (7.072,6.336) (7.145,6.352) (7.218,6.368) (7.292,6.383) (7.365,6.398) (7.439,6.412) (7.513,6.425) (7.587,6.439) (7.661,6.451) (7.735,6.463) (7.809,6.474) (7.883,6.485) (7.957,6.496) (8.032,6.506) (8.106,6.515) (8.180,6.524) (8.255,6.532) (8.330,6.539) (8.404,6.547) (8.479,6.553) (8.554,6.559) (8.629,6.565) (8.703,6.569) (8.778,6.574) (8.853,6.578) (8.928,6.581) (9.003,6.583) (9.078,6.586) (9.153,6.587) (9.228,6.588) (9.303,6.588) (9.378,6.588) (9.453,6.587) (9.528,6.586) (9.603,6.584) (9.678,6.582) (9.753,6.579) (9.828,6.575) (9.903,6.571) (9.977,6.566) (10.000,6.564)}
\LC{(0.444,6.194) (0.452,6.268) (0.460,6.343) (0.468,6.418) (0.477,6.492) (0.485,6.567) (0.493,6.641) (0.502,6.716) (0.510,6.790) (0.518,6.865) (0.527,6.939) (0.535,7.014) (0.544,7.088) (0.553,7.163) (0.561,7.237) (0.570,7.312) (0.579,7.386) (0.588,7.461) (0.596,7.535) (0.605,7.610) (0.614,7.684) (0.623,7.759) (0.632,7.833) (0.641,7.908) (0.650,7.982) (0.660,8.056) (0.669,8.131) (0.678,8.205) (0.687,8.280) (0.697,8.354) (0.706,8.429) (0.715,8.503) (0.725,8.577) (0.734,8.652) (0.744,8.726) (0.753,8.801) (0.763,8.875) (0.773,8.949) (0.782,9.024) (0.792,9.098) (0.802,9.172) (0.812,9.247) (0.822,9.321) (0.831,9.395) (0.841,9.470) (0.851,9.544) (0.861,9.618) (0.871,9.693) (0.881,9.767) (0.892,9.841) (0.902,9.916) (0.912,9.990) (0.915,10.000)}
\LC{(-0.458,5.643) (-0.467,5.717) (-0.476,5.792) (-0.485,5.866) (-0.495,5.941) (-0.504,6.015) (-0.513,6.090) (-0.523,6.164) (-0.532,6.238) (-0.542,6.313) (-0.552,6.387) (-0.561,6.461) (-0.571,6.536) (-0.581,6.610) (-0.591,6.685) (-0.601,6.759) (-0.611,6.833) (-0.621,6.908) (-0.631,6.982) (-0.641,7.056) (-0.651,7.130) (-0.661,7.205) (-0.672,7.279) (-0.682,7.353) (-0.692,7.428) (-0.703,7.502) (-0.713,7.576) (-0.724,7.650) (-0.734,7.725) (-0.745,7.799) (-0.756,7.873) (-0.767,7.947) (-0.777,8.022) (-0.788,8.096) (-0.799,8.170) (-0.810,8.244) (-0.821,8.318) (-0.832,8.392) (-0.843,8.467) (-0.854,8.541) (-0.866,8.615) (-0.877,8.689) (-0.888,8.763) (-0.900,8.837) (-0.911,8.912) (-0.922,8.986) (-0.934,9.060) (-0.945,9.134) (-0.957,9.208) (-0.969,9.282) (-0.980,9.356) (-0.992,9.430) (-1.004,9.504) (-1.016,9.578) (-1.028,9.652) (-1.039,9.726) (-1.051,9.801) (-1.063,9.875) (-1.076,9.949) (-1.088,10.000)}
\LC{(1.579,1.049) (1.654,1.043) (1.728,1.033) (1.802,1.020) (1.875,1.002) (1.947,0.981) (2.017,0.956) (2.086,0.927) (2.154,0.894) (2.219,0.857) (2.282,0.816) (2.342,0.771) (2.399,0.723) (2.453,0.671) (2.503,0.615) (2.549,0.556) (2.591,0.493) (2.628,0.428) (2.660,0.360) (2.686,0.290) (2.707,0.218) (2.722,0.144) (2.731,0.070) (2.733,-0.005) (2.730,-0.080) (2.720,-0.154) (2.704,-0.228) (2.683,-0.299) (2.656,-0.369) (2.623,-0.437) (2.586,-0.502) (2.544,-0.564) (2.497,-0.622) (2.446,-0.678) (2.392,-0.729) (2.334,-0.777) (2.274,-0.822) (2.210,-0.862) (2.145,-0.898) (2.077,-0.931) (2.008,-0.959) (1.937,-0.984) (1.865,-1.005) (1.792,-1.022) (1.718,-1.035) (1.644,-1.044) (1.594,-1.048)}
\LC{(1.474,-10.000) (1.457,-9.936) (1.441,-9.863) (1.425,-9.789) (1.409,-9.716) (1.393,-9.643) (1.377,-9.570) (1.361,-9.496) (1.345,-9.423) (1.329,-9.350) (1.313,-9.276) (1.298,-9.203) (1.282,-9.130) (1.266,-9.056) (1.251,-8.983) (1.235,-8.910) (1.220,-8.836) (1.205,-8.763) (1.189,-8.689) (1.174,-8.616) (1.159,-8.542) (1.144,-8.469) (1.129,-8.395) (1.114,-8.322) (1.099,-8.248) (1.084,-8.175) (1.070,-8.101) (1.055,-8.028) (1.040,-7.954) (1.026,-7.881) (1.012,-7.807) (0.997,-7.734) (0.983,-7.660) (0.969,-7.586) (0.954,-7.513) (0.940,-7.439) (0.926,-7.365) (0.912,-7.292) (0.898,-7.218) (0.884,-7.144) (0.871,-7.070) (0.857,-6.997) (0.843,-6.923) (0.830,-6.849) (0.816,-6.775) (0.803,-6.702) (0.790,-6.628) (0.776,-6.554) (0.763,-6.480) (0.750,-6.406) (0.737,-6.332) (0.724,-6.259) (0.711,-6.185) (0.699,-6.111) (0.686,-6.037) (0.673,-5.963) (0.661,-5.889) (0.648,-5.815) (0.636,-5.741) (0.623,-5.667) (0.611,-5.593) (0.599,-5.519) (0.587,-5.445) (0.575,-5.371) (0.571,-5.346)}
\LC{(-7.036,-10.000) (-6.975,-9.958) (-6.914,-9.914) (-6.853,-9.870) (-6.793,-9.826) (-6.732,-9.782) (-6.672,-9.737) (-6.612,-9.692) (-6.552,-9.647) (-6.492,-9.602) (-6.433,-9.556) (-6.373,-9.510) (-6.314,-9.464) (-6.255,-9.418) (-6.196,-9.372) (-6.137,-9.326) (-6.078,-9.279) (-6.020,-9.232) (-5.961,-9.185) (-5.903,-9.137) (-5.845,-9.090) (-5.787,-9.042) (-5.730,-8.994) (-5.672,-8.946) (-5.615,-8.898) (-5.558,-8.849) (-5.501,-8.800) (-5.444,-8.751) (-5.387,-8.702) (-5.331,-8.653) (-5.274,-8.603) (-5.218,-8.554) (-5.162,-8.504) (-5.107,-8.454) (-5.051,-8.403) (-4.996,-8.353) (-4.940,-8.302) (-4.885,-8.251) (-4.830,-8.200) (-4.776,-8.149) (-4.721,-8.097) (-4.667,-8.045) (-4.613,-7.993) (-4.559,-7.941) (-4.505,-7.889) (-4.452,-7.836) (-4.398,-7.784) (-4.345,-7.731) (-4.292,-7.678) (-4.239,-7.624) (-4.187,-7.571) (-4.134,-7.517) (-4.082,-7.463) (-4.030,-7.409) (-3.978,-7.355) (-3.927,-7.301) (-3.876,-7.246) (-3.824,-7.191) (-3.773,-7.136) (-3.723,-7.081) (-3.672,-7.025) (-3.622,-6.970) (-3.572,-6.914) (-3.522,-6.858) (-3.472,-6.802) (-3.423,-6.745) (-3.374,-6.689) (-3.325,-6.632) (-3.276,-6.575) (-3.227,-6.518) (-3.179,-6.461) (-3.131,-6.403) (-3.083,-6.345) (-3.035,-6.287) (-2.988,-6.229) (-2.941,-6.171) (-2.894,-6.113) (-2.847,-6.054) (-2.800,-5.995) (-2.754,-5.936) (-2.708,-5.877) (-2.662,-5.817) (-2.617,-5.758) (-2.572,-5.698) (-2.527,-5.638) (-2.482,-5.578) (-2.437,-5.517) (-2.393,-5.457) (-2.349,-5.396) (-2.320,-5.356)}
\LC{(10.000,-2.658) (9.952,-2.713) (9.900,-2.767) (9.847,-2.819) (9.793,-2.871) (9.738,-2.922) (9.682,-2.972) (9.625,-3.021) (9.568,-3.069) (9.509,-3.117) (9.450,-3.163) (9.391,-3.208) (9.330,-3.253) (9.269,-3.296) (9.207,-3.339) (9.145,-3.381) (9.082,-3.421) (9.018,-3.461) (8.954,-3.500) (8.889,-3.537) (8.824,-3.574) (8.758,-3.610) (8.691,-3.645) (8.625,-3.679) (8.557,-3.711) (8.489,-3.743) (8.421,-3.774) (8.352,-3.804) (8.283,-3.833) (8.214,-3.861) (8.144,-3.889) (8.073,-3.915) (8.003,-3.940) (7.932,-3.964) (7.860,-3.987) (7.789,-4.010) (7.717,-4.031) (7.645,-4.051) (7.572,-4.071) (7.500,-4.089) (7.427,-4.107) (7.354,-4.124) (7.280,-4.139) (7.207,-4.154) (7.133,-4.168) (7.059,-4.181) (6.985,-4.193) (6.911,-4.204) (6.837,-4.214) (6.762,-4.223) (6.688,-4.231) (6.613,-4.239) (6.538,-4.245) (6.463,-4.251) (6.389,-4.255) (6.314,-4.259) (6.239,-4.262) (6.164,-4.264) (6.089,-4.265) (6.014,-4.265) (5.939,-4.265) (5.864,-4.263) (5.789,-4.261) (5.714,-4.258) (5.639,-4.253) (5.564,-4.248) (5.489,-4.243) (5.415,-4.236) (5.340,-4.228) (5.266,-4.220) (5.191,-4.211) (5.117,-4.201) (5.043,-4.190) (4.969,-4.178) (4.895,-4.165) (4.821,-4.152) (4.747,-4.138) (4.674,-4.123) (4.600,-4.107) (4.527,-4.090) (4.454,-4.073) (4.382,-4.055) (4.309,-4.036) (4.237,-4.016) (4.165,-3.995) (4.093,-3.974) (4.021,-3.952) (3.950,-3.929) (3.879,-3.905) (3.808,-3.880) (3.737,-3.855) (3.667,-3.829) (3.597,-3.802) (3.527,-3.774) (3.458,-3.746) (3.388,-3.717) (3.320,-3.687) (3.251,-3.657) (3.183,-3.625) (3.115,-3.593) (3.048,-3.560) (2.981,-3.527) (2.914,-3.492) (2.848,-3.457) (2.782,-3.422) (2.716,-3.385) (2.651,-3.348) (2.586,-3.310) (2.522,-3.272)}
\LC{(-10.000,-4.265) (-9.942,-4.298) (-9.875,-4.330) (-9.806,-4.362) (-9.738,-4.392) (-9.669,-4.422) (-9.600,-4.451) (-9.530,-4.479) (-9.460,-4.506) (-9.390,-4.532) (-9.320,-4.558) (-9.249,-4.583) (-9.178,-4.606) (-9.106,-4.629) (-9.035,-4.652) (-8.963,-4.673) (-8.891,-4.694) (-8.818,-4.713) (-8.746,-4.732) (-8.673,-4.751) (-8.600,-4.768) (-8.527,-4.784) (-8.454,-4.800) (-8.380,-4.815) (-8.306,-4.829) (-8.233,-4.842) (-8.159,-4.855) (-8.085,-4.867) (-8.010,-4.878) (-7.936,-4.888) (-7.862,-4.897) (-7.787,-4.906) (-7.712,-4.913) (-7.638,-4.920) (-7.563,-4.926) (-7.488,-4.932) (-7.413,-4.936) (-7.339,-4.940) (-7.264,-4.943) (-7.189,-4.946) (-7.114,-4.947) (-7.039,-4.948) (-6.964,-4.948) (-6.889,-4.947) (-6.814,-4.946) (-6.739,-4.944) (-6.664,-4.941) (-6.589,-4.937) (-6.514,-4.933) (-6.439,-4.927) (-6.364,-4.921) (-6.290,-4.915) (-6.215,-4.907) (-6.140,-4.899) (-6.066,-4.890) (-5.992,-4.881) (-5.917,-4.870) (-5.843,-4.859) (-5.769,-4.848) (-5.695,-4.835) (-5.621,-4.822) (-5.548,-4.808) (-5.474,-4.794) (-5.401,-4.778)}
\LC{(-10.000,-7.079) (-9.926,-7.079) (-9.851,-7.078) (-9.776,-7.076) (-9.701,-7.074) (-9.626,-7.072) (-9.552,-7.069) (-9.477,-7.066) (-9.402,-7.062) (-9.327,-7.057) (-9.252,-7.052) (-9.177,-7.046) (-9.102,-7.040) (-9.028,-7.034) (-8.953,-7.027) (-8.878,-7.019) (-8.804,-7.011) (-8.729,-7.002) (-8.655,-6.993) (-8.581,-6.984) (-8.506,-6.973) (-8.432,-6.963) (-8.358,-6.952) (-8.284,-6.940) (-8.210,-6.928) (-8.136,-6.915) (-8.062,-6.902) (-7.988,-6.888) (-7.915,-6.874) (-7.841,-6.860) (-7.768,-6.845) (-7.694,-6.829) (-7.621,-6.813) (-7.548,-6.796) (-7.475,-6.779) (-7.402,-6.762) (-7.329,-6.744) (-7.256,-6.725) (-7.184,-6.706) (-7.111,-6.687) (-7.039,-6.667) (-6.967,-6.646) (-6.895,-6.625) (-6.823,-6.604) (-6.751,-6.582) (-6.680,-6.560) (-6.608,-6.537) (-6.537,-6.514) (-6.466,-6.490) (-6.395,-6.466) (-6.324,-6.441) (-6.253,-6.416) (-6.183,-6.391) (-6.112,-6.365) (-6.042,-6.338) (-5.972,-6.311) (-5.902,-6.284) (-5.833,-6.256) (-5.763,-6.227) (-5.694,-6.199) (-5.625,-6.169) (-5.556,-6.140) (-5.488,-6.110) (-5.419,-6.079) (-5.351,-6.048) (-5.283,-6.016) (-5.215,-5.985) (-5.147,-5.952) (-5.080,-5.919) (-5.013,-5.886)}
\LC{(6.534,-10.000) (6.476,-9.962) (6.418,-9.915) (6.360,-9.868) (6.302,-9.820) (6.244,-9.772) (6.186,-9.724) (6.129,-9.676) (6.072,-9.628) (6.014,-9.579) (5.957,-9.530) (5.901,-9.481) (5.844,-9.432) (5.787,-9.383) (5.731,-9.333) (5.675,-9.284) (5.619,-9.234) (5.563,-9.184) (5.507,-9.134) (5.452,-9.083) (5.396,-9.033) (5.341,-8.982) (5.286,-8.931) (5.231,-8.880) (5.177,-8.828) (5.122,-8.777) (5.068,-8.725) (5.013,-8.673) (4.959,-8.621) (4.906,-8.569) (4.852,-8.517) (4.799,-8.464) (4.745,-8.411) (4.692,-8.358) (4.639,-8.305) (4.586,-8.252) (4.534,-8.198) (4.482,-8.145) (4.429,-8.091) (4.377,-8.037) (4.325,-7.983) (4.274,-7.928) (4.222,-7.874) (4.171,-7.819) (4.120,-7.764) (4.069,-7.709) (4.019,-7.654) (3.968,-7.598) (3.918,-7.542) (3.868,-7.487) (3.818,-7.431) (3.768,-7.374) (3.719,-7.318) (3.669,-7.262) (3.620,-7.205) (3.571,-7.148) (3.523,-7.091) (3.474,-7.034) (3.426,-6.976) (3.378,-6.919) (3.330,-6.861) (3.282,-6.803) (3.235,-6.745) (3.188,-6.687) (3.141,-6.628) (3.094,-6.570) (3.047,-6.511) (3.001,-6.452) (2.955,-6.393) (2.909,-6.333) (2.863,-6.274) (2.818,-6.214) (2.773,-6.154) (2.728,-6.094) (2.683,-6.034) (2.639,-5.974) (2.594,-5.913) (2.550,-5.852) (2.521,-5.812)}
\LC{(-2.798,6.367) (-2.842,6.428) (-2.886,6.488) (-2.931,6.548) (-2.976,6.608) (-3.021,6.668) (-3.067,6.728) (-3.112,6.788) (-3.158,6.847) (-3.204,6.906) (-3.250,6.965) (-3.297,7.024) (-3.343,7.083) (-3.390,7.141) (-3.437,7.200) (-3.484,7.258) (-3.532,7.316) (-3.580,7.374) (-3.627,7.432) (-3.675,7.489) (-3.724,7.547) (-3.772,7.604) (-3.821,7.661) (-3.870,7.718) (-3.919,7.775) (-3.968,7.832) (-4.017,7.888) (-4.067,7.944) (-4.116,8.000) (-4.166,8.056) (-4.216,8.112) (-4.267,8.168) (-4.317,8.223) (-4.368,8.279) (-4.419,8.334) (-4.470,8.389) (-4.521,8.444) (-4.572,8.498) (-4.624,8.553) (-4.676,8.607) (-4.728,8.661) (-4.780,8.715) (-4.832,8.769) (-4.884,8.822) (-4.937,8.876) (-4.990,8.929) (-5.043,8.982) (-5.096,9.035) (-5.149,9.088) (-5.203,9.141) (-5.256,9.193) (-5.310,9.245) (-5.364,9.298) (-5.418,9.349) (-5.472,9.401) (-5.527,9.453) (-5.581,9.504) (-5.636,9.556) (-5.691,9.607) (-5.746,9.658) (-5.801,9.708) (-5.857,9.759) (-5.912,9.809) (-5.968,9.860) (-6.024,9.910) (-6.080,9.960) (-6.136,10.000)}
\LC{(-1.929,0.938) (-1.998,0.911) (-2.066,0.879) (-2.132,0.844) (-2.196,0.804) (-2.257,0.760) (-2.315,0.713) (-2.370,0.661) (-2.421,0.606) (-2.467,0.548) (-2.510,0.486) (-2.547,0.421) (-2.579,0.353) (-2.606,0.283) (-2.627,0.211) (-2.641,0.137) (-2.650,0.063) (-2.652,-0.012) (-2.648,-0.087) (-2.637,-0.161) (-2.621,-0.234) (-2.598,-0.306) (-2.570,-0.375) (-2.536,-0.442) (-2.497,-0.506) (-2.453,-0.567) (-2.405,-0.624) (-2.353,-0.678) (-2.297,-0.728) (-2.238,-0.775) (-2.176,-0.817) (-2.112,-0.855) (-2.045,-0.890) (-1.976,-0.920) (-1.930,-0.938)}
\LC{(-1.501,-10.000) (-1.484,-9.935) (-1.467,-9.862) (-1.451,-9.789) (-1.435,-9.716) (-1.418,-9.643) (-1.402,-9.569) (-1.386,-9.496) (-1.369,-9.423) (-1.353,-9.350) (-1.337,-9.276) (-1.321,-9.203) (-1.305,-9.130) (-1.290,-9.057) (-1.274,-8.983) (-1.258,-8.910) (-1.242,-8.837) (-1.227,-8.763) (-1.211,-8.690) (-1.196,-8.616) (-1.180,-8.543) (-1.165,-8.470) (-1.150,-8.396) (-1.135,-8.323) (-1.120,-8.249) (-1.104,-8.176) (-1.089,-8.102) (-1.075,-8.029) (-1.060,-7.955) (-1.045,-7.882) (-1.030,-7.808) (-1.016,-7.735) (-1.001,-7.661) (-0.986,-7.588) (-0.972,-7.514) (-0.958,-7.440) (-0.943,-7.367) (-0.929,-7.293) (-0.915,-7.219) (-0.901,-7.146) (-0.887,-7.072) (-0.873,-6.998) (-0.859,-6.925) (-0.845,-6.851) (-0.832,-6.777) (-0.818,-6.703) (-0.804,-6.630) (-0.791,-6.556) (-0.777,-6.482) (-0.764,-6.408) (-0.751,-6.334) (-0.738,-6.261) (-0.725,-6.187) (-0.712,-6.113) (-0.699,-6.039) (-0.686,-5.965) (-0.673,-5.891) (-0.660,-5.817) (-0.648,-5.743) (-0.635,-5.669) (-0.623,-5.596) (-0.610,-5.522) (-0.598,-5.448) (-0.586,-5.374) (-0.574,-5.300) (-0.562,-5.226) (-0.550,-5.152) (-0.538,-5.077) (-0.526,-5.003) (-0.522,-4.979)}
\LC{(-5.577,4.285) (-5.652,4.290) (-5.727,4.295) (-5.802,4.299) (-5.877,4.302) (-5.952,4.304) (-6.027,4.305) (-6.102,4.305) (-6.177,4.304) (-6.252,4.303) (-6.327,4.301) (-6.401,4.297) (-6.476,4.293) (-6.551,4.288) (-6.626,4.282) (-6.701,4.276) (-6.775,4.268) (-6.850,4.259) (-6.924,4.250) (-6.998,4.239) (-7.073,4.228) (-7.147,4.216) (-7.220,4.203) (-7.294,4.189) (-7.368,4.174) (-7.441,4.158) (-7.514,4.141) (-7.587,4.123) (-7.659,4.104) (-7.732,4.085) (-7.804,4.064) (-7.876,4.042) (-7.947,4.020) (-8.019,3.996) (-8.089,3.972) (-8.160,3.947) (-8.230,3.920) (-8.300,3.893) (-8.370,3.865) (-8.439,3.836) (-8.507,3.806) (-8.576,3.774) (-8.643,3.742) (-8.711,3.709) (-8.778,3.675) (-8.844,3.641) (-8.910,3.605) (-8.975,3.568) (-9.040,3.530) (-9.104,3.491) (-9.168,3.452) (-9.231,3.411) (-9.293,3.369) (-9.355,3.327) (-9.416,3.283) (-9.477,3.239) (-9.536,3.194) (-9.595,3.147) (-9.654,3.100) (-9.711,3.052) (-9.768,3.003) (-9.824,2.953) (-9.879,2.902) (-9.933,2.850) (-9.986,2.797) (-10.000,2.780)}

    % Dipolo puntual - puntas de flecha
% Generado por scripts/generar_datos.py (no editar a mano).
\FL{1.429}{-0.015}{-91.8}
\FL{1.467}{2.168}{37.2}
\FL{4.190}{3.194}{4.1}
\FL{6.989}{2.463}{-35.4}
\FL{8.330}{-0.035}{-90.7}
\FL{6.953}{-2.489}{-145.3}
\FL{4.146}{-3.191}{175.4}
\FL{1.433}{-2.141}{142.3}
\FL{-1.467}{2.168}{142.8}
\FL{-4.190}{3.194}{175.9}
\FL{-6.989}{2.463}{-144.6}
\FL{-8.330}{-0.035}{-89.3}
\FL{-6.953}{-2.489}{-34.7}
\FL{-4.146}{-3.191}{4.6}
\FL{-1.433}{-2.141}{37.7}
\FL{-1.429}{-0.015}{-88.2}
\FL{3.888}{-8.732}{126.6}
\FL{2.141}{-6.053}{119.5}
\FL{0.773}{-3.164}{110.7}
\FL{-4.014}{-8.749}{52.4}
\FL{-2.211}{-6.078}{59.7}
\FL{-0.799}{-3.182}{68.8}
\FL{8.283}{-6.321}{175.9}
\FL{4.997}{-5.511}{156.6}
\FL{2.162}{-3.665}{137.0}
\FL{-4.602}{2.370}{-163.1}
\FL{-6.550}{-0.011}{-89.7}
\FL{-4.581}{-2.377}{-16.5}
\FL{0.881}{2.999}{65.3}
\FL{2.174}{5.280}{56.0}
\FL{3.771}{7.329}{48.3}
\FL{5.631}{9.179}{41.4}
\FL{-5.151}{3.827}{177.3}
\FL{-7.813}{3.273}{-152.7}
\FL{-9.682}{1.331}{-113.2}
\FL{-9.677}{-1.343}{-66.6}
\FL{-7.802}{-3.279}{-27.1}
\FL{-5.138}{-3.826}{2.8}
\FL{-0.035}{-8.072}{89.6}
\FL{-0.013}{-4.222}{89.7}
\FL{-1.128}{3.214}{119.3}
\FL{-2.813}{5.599}{130.8}
\FL{-4.925}{7.654}{140.6}
\FL{-7.322}{9.325}{149.6}
\FL{-3.065}{4.757}{140.7}
\FL{-5.627}{6.350}{155.4}
\FL{-8.507}{7.248}{170.0}
\FL{-8.690}{-7.903}{13.5}
\FL{-6.169}{-7.023}{25.0}
\FL{-3.895}{-5.671}{36.6}
\FL{-1.952}{-3.877}{49.1}
\FL{5.299}{2.380}{-23.9}
\FL{6.981}{-0.015}{-90.4}
\FL{5.272}{-2.392}{-156.6}
\FL{8.439}{-7.362}{168.7}
\FL{5.486}{-6.366}{154.1}
\FL{4.915}{3.849}{5.5}
\FL{7.491}{3.499}{-21.9}
\FL{9.507}{1.882}{-57.2}
\FL{3.579}{6.089}{43.2}
\FL{5.947}{7.925}{32.5}
\FL{8.605}{9.307}{22.4}
\FL{3.187}{1.458}{-23.0}
\FL{3.151}{-1.473}{-158.1}
\FL{2.160}{-8.592}{111.3}
\FL{1.153}{-5.714}{107.2}
\FL{-4.995}{6.777}{146.6}
\FL{-8.259}{8.389}{160.8}
\FL{9.384}{0.914}{-73.4}
\FL{8.847}{-1.977}{-126.7}
\FL{6.306}{-3.544}{-167.7}
\FL{7.949}{-9.303}{153.6}
\FL{5.005}{-7.467}{142.4}
\FL{2.512}{-5.091}{130.1}
\FL{-8.679}{-5.784}{-3.2}
\FL{-6.027}{-5.529}{13.9}
\FL{-3.557}{-4.524}{30.4}
\FL{8.555}{-5.231}{-172.2}
\FL{5.625}{-5.084}{166.8}
\FL{-0.941}{4.720}{107.0}
\FL{-2.243}{8.288}{112.8}
\FL{-8.254}{-8.950}{22.6}
\FL{-5.066}{-7.187}{35.4}
\FL{-2.308}{-8.603}{67.3}
\FL{-1.253}{-5.795}{71.6}
\FL{3.758}{2.167}{-10.9}
\FL{5.781}{-0.009}{-90.3}
\FL{3.741}{-2.171}{-169.4}
\FL{-0.004}{4.753}{90.1}
\FL{-0.010}{8.253}{90.1}
\FL{6.083}{4.475}{2.1}
\FL{8.815}{3.977}{-23.7}
\FL{-2.961}{3.778}{149.5}
\FL{-5.722}{4.766}{171.2}
\FL{-8.648}{4.626}{-165.1}
\FL{4.561}{6.002}{32.0}
\FL{8.107}{7.526}{14.6}
\FL{-2.725}{1.238}{-156.7}
\FL{-2.701}{-1.247}{-22.5}
\FL{-8.897}{-3.737}{-27.2}
\FL{-6.345}{-4.367}{-1.5}
\FL{-3.980}{1.547}{-149.1}
\FL{-3.966}{-1.555}{-30.5}
\FL{-5.638}{5.666}{161.3}
\FL{-8.531}{6.181}{178.7}
\FL{2.629}{3.737}{35.5}
\FL{5.423}{5.041}{14.6}
\FL{8.497}{5.250}{-7.3}
\FL{2.013}{6.217}{62.9}
\FL{3.508}{8.788}{57.0}
\FL{-2.126}{6.103}{119.1}
\FL{-3.822}{8.756}{125.9}
\FL{0.825}{5.102}{76.2}
\FL{1.763}{8.396}{72.1}
\FL{8.584}{-8.611}{161.4}
\FL{5.940}{-7.430}{150.4}
\FL{5.637}{5.904}{20.8}
\FL{8.529}{6.557}{4.5}
\FL{0.666}{8.106}{83.0}
\FL{-0.752}{7.848}{98.2}
\FL{2.733}{-0.005}{-90.3}
\FL{0.983}{-7.660}{101.0}
\FL{-5.711}{-8.978}{39.9}
\FL{-3.341}{-6.651}{49.2}
\FL{8.932}{-3.512}{-149.6}
\FL{6.364}{-4.257}{-177.0}
\FL{3.714}{-3.846}{159.8}
\FL{-7.737}{-4.911}{-5.8}
\FL{-7.426}{-6.768}{13.6}
\FL{5.415}{-9.049}{137.6}
\FL{3.410}{-6.957}{129.9}
\FL{-4.368}{8.279}{132.6}
\FL{-2.652}{-0.012}{-89.2}
\FL{-0.967}{-7.489}{78.9}
\FL{-7.971}{4.012}{-161.9}
}
  \dipoloRojo{0.85}
  \node[font=\tiny, text=mmarojo] at (0.62,0.05) {$\vp$};
  \marcoMMA{10}{5}{$x/\ell$}{$z/\ell$}
\end{tikzpicture}\\[1pt]{\small\scshape (a)\; dipolo puntual}
\end{minipage}\hfill
\begin{minipage}[t]{0.49\textwidth}\centering
\begin{tikzpicture}[x=0.29cm, y=0.29cm]
  \panelMMA{10}{mmafondo}{1.2}{%
    % Dipolo físico - equipotenciales
% Generado por scripts/generar_datos.py (no editar a mano).
\EQ{(0.250,-0.010) (0.300,-0.012) (0.350,-0.013) (0.400,-0.015) (0.450,-0.018) (0.500,-0.021) (0.550,-0.024) (0.600,-0.028) (0.650,-0.032) (0.700,-0.037) (0.750,-0.043) (0.800,-0.049) (0.850,-0.056) (0.900,-0.064) (0.950,-0.073) (1.000,-0.082) (1.050,-0.093) (1.100,-0.104) (1.150,-0.116) (1.200,-0.130) (1.250,-0.145) (1.300,-0.161) (1.350,-0.178) (1.400,-0.196) (1.450,-0.216) (1.500,-0.238) (1.550,-0.261) (1.600,-0.286) (1.650,-0.313) (1.700,-0.341) (1.750,-0.372) (1.792,-0.400) (1.850,-0.441) (1.894,-0.475) (1.950,-0.521) (2.000,-0.566) (2.050,-0.615) (2.100,-0.668) (2.150,-0.725) (2.189,-0.775) (2.225,-0.824) (2.260,-0.875) (2.292,-0.925) (2.321,-0.975) (2.347,-1.025) (2.372,-1.075) (2.395,-1.125) (2.416,-1.175) (2.435,-1.225) (2.452,-1.275) (2.468,-1.325) (2.482,-1.375) (2.495,-1.425) (2.506,-1.475) (2.515,-1.525) (2.523,-1.575) (2.530,-1.625) (2.535,-1.675) (2.539,-1.725) (2.542,-1.775) (2.543,-1.825) (2.543,-1.875) (2.542,-1.925) (2.540,-1.975) (2.536,-2.025) (2.531,-2.075) (2.525,-2.125) (2.517,-2.175) (2.508,-2.225) (2.498,-2.275) (2.487,-2.325) (2.474,-2.375) (2.460,-2.425) (2.444,-2.475) (2.428,-2.525) (2.410,-2.575) (2.390,-2.625) (2.369,-2.675) (2.347,-2.725) (2.325,-2.771) (2.300,-2.821) (2.275,-2.867) (2.250,-2.912) (2.212,-2.975) (2.180,-3.025) (2.150,-3.069) (2.110,-3.125) (2.075,-3.172) (2.033,-3.225) (2.000,-3.264) (1.950,-3.321) (1.900,-3.374) (1.850,-3.424) (1.800,-3.472) (1.750,-3.517) (1.711,-3.550) (1.650,-3.600) (1.600,-3.638) (1.550,-3.675) (1.500,-3.709) (1.450,-3.742) (1.400,-3.774) (1.356,-3.800) (1.312,-3.825) (1.266,-3.850) (1.218,-3.875) (1.168,-3.900) (1.114,-3.925) (1.057,-3.950) (1.000,-3.974) (0.950,-3.993) (0.900,-4.011) (0.850,-4.028) (0.800,-4.044) (0.750,-4.059) (0.700,-4.073) (0.650,-4.086) (0.600,-4.098) (0.550,-4.108) (0.500,-4.118) (0.450,-4.127) (0.400,-4.135) (0.350,-4.142) (0.300,-4.148) (0.250,-4.153) (0.200,-4.157) (0.150,-4.161) (0.100,-4.163) (0.050,-4.164) (0.000,-4.165) (-0.050,-4.164) (-0.100,-4.163) (-0.150,-4.161) (-0.200,-4.157) (-0.250,-4.153) (-0.300,-4.148) (-0.350,-4.142) (-0.400,-4.135) (-0.450,-4.127) (-0.500,-4.118) (-0.550,-4.108) (-0.600,-4.098) (-0.650,-4.086) (-0.700,-4.073) (-0.750,-4.059) (-0.800,-4.044) (-0.850,-4.028) (-0.900,-4.011) (-0.950,-3.993) (-1.000,-3.974) (-1.050,-3.953) (-1.100,-3.931) (-1.150,-3.909) (-1.200,-3.884) (-1.250,-3.859) (-1.300,-3.832) (-1.350,-3.803) (-1.398,-3.775) (-1.450,-3.742) (-1.500,-3.709) (-1.549,-3.675) (-1.600,-3.638) (-1.649,-3.600) (-1.700,-3.559) (-1.740,-3.525) (-1.796,-3.475) (-1.849,-3.425) (-1.899,-3.375) (-1.946,-3.325) (-1.990,-3.275) (-2.025,-3.234) (-2.072,-3.175) (-2.110,-3.125) (-2.146,-3.075) (-2.175,-3.032) (-2.212,-2.975) (-2.242,-2.925) (-2.271,-2.875) (-2.298,-2.825) (-2.323,-2.775) (-2.347,-2.725) (-2.369,-2.675) (-2.390,-2.625) (-2.410,-2.575) (-2.428,-2.525) (-2.444,-2.475) (-2.460,-2.425) (-2.474,-2.375) (-2.487,-2.325) (-2.498,-2.275) (-2.508,-2.225) (-2.517,-2.175) (-2.525,-2.125) (-2.531,-2.075) (-2.536,-2.025) (-2.540,-1.975) (-2.542,-1.925) (-2.543,-1.875) (-2.543,-1.825) (-2.542,-1.775) (-2.539,-1.725) (-2.535,-1.675) (-2.530,-1.625) (-2.523,-1.575) (-2.515,-1.525) (-2.506,-1.475) (-2.495,-1.425) (-2.482,-1.375) (-2.468,-1.325) (-2.452,-1.275) (-2.435,-1.225) (-2.416,-1.175) (-2.395,-1.125) (-2.372,-1.075) (-2.347,-1.025) (-2.321,-0.975) (-2.292,-0.925) (-2.260,-0.875) (-2.226,-0.825) (-2.189,-0.775) (-2.150,-0.726) (-2.107,-0.675) (-2.060,-0.625) (-2.010,-0.575) (-1.954,-0.525) (-1.900,-0.480) (-1.850,-0.441) (-1.800,-0.406) (-1.754,-0.375) (-1.700,-0.341) (-1.650,-0.313) (-1.600,-0.286) (-1.550,-0.261) (-1.500,-0.238) (-1.450,-0.216) (-1.400,-0.196) (-1.350,-0.178) (-1.300,-0.161) (-1.250,-0.145) (-1.200,-0.130) (-1.150,-0.116) (-1.100,-0.104) (-1.050,-0.093) (-1.000,-0.082) (-0.950,-0.073) (-0.900,-0.064) (-0.850,-0.056) (-0.800,-0.049) (-0.750,-0.043) (-0.700,-0.037) (-0.650,-0.032) (-0.600,-0.028) (-0.550,-0.024) (-0.500,-0.021) (-0.450,-0.018) (-0.400,-0.015) (-0.350,-0.013) (-0.300,-0.012) (-0.250,-0.010)}
\EQ{(0.250,-0.009) (0.300,-0.010) (0.350,-0.012) (0.400,-0.014) (0.450,-0.016) (0.500,-0.019) (0.550,-0.022) (0.600,-0.025) (0.650,-0.029) (0.700,-0.034) (0.750,-0.039) (0.800,-0.044) (0.850,-0.051) (0.900,-0.058) (0.950,-0.065) (1.000,-0.074) (1.050,-0.083) (1.100,-0.094) (1.150,-0.105) (1.200,-0.117) (1.250,-0.130) (1.300,-0.144) (1.350,-0.159) (1.400,-0.176) (1.450,-0.194) (1.500,-0.213) (1.550,-0.233) (1.600,-0.255) (1.650,-0.279) (1.700,-0.304) (1.750,-0.331) (1.800,-0.360) (1.850,-0.391) (1.900,-0.425) (1.950,-0.460) (2.000,-0.498) (2.050,-0.540) (2.091,-0.575) (2.144,-0.625) (2.193,-0.675) (2.238,-0.725) (2.275,-0.769) (2.319,-0.825) (2.350,-0.868) (2.389,-0.925) (2.420,-0.975) (2.449,-1.025) (2.475,-1.074) (2.500,-1.124) (2.524,-1.175) (2.545,-1.225) (2.564,-1.275) (2.582,-1.325) (2.598,-1.375) (2.613,-1.425) (2.626,-1.475) (2.638,-1.525) (2.648,-1.575) (2.657,-1.625) (2.665,-1.675) (2.671,-1.725) (2.676,-1.775) (2.680,-1.825) (2.682,-1.875) (2.683,-1.925) (2.683,-1.975) (2.682,-2.025) (2.680,-2.075) (2.676,-2.125) (2.671,-2.175) (2.665,-2.225) (2.658,-2.275) (2.649,-2.325) (2.639,-2.375) (2.628,-2.425) (2.616,-2.475) (2.603,-2.525) (2.588,-2.575) (2.572,-2.625) (2.555,-2.675) (2.537,-2.725) (2.517,-2.775) (2.496,-2.825) (2.475,-2.871) (2.450,-2.923) (2.425,-2.972) (2.400,-3.019) (2.375,-3.063) (2.338,-3.125) (2.306,-3.175) (2.275,-3.222) (2.238,-3.275) (2.200,-3.325) (2.161,-3.375) (2.125,-3.419) (2.077,-3.475) (2.031,-3.525) (1.983,-3.575) (1.932,-3.625) (1.878,-3.675) (1.825,-3.722) (1.775,-3.763) (1.729,-3.800) (1.675,-3.841) (1.628,-3.875) (1.575,-3.912) (1.525,-3.944) (1.476,-3.975) (1.425,-4.005) (1.375,-4.034) (1.325,-4.061) (1.275,-4.087) (1.225,-4.111) (1.175,-4.134) (1.125,-4.156) (1.075,-4.177) (1.025,-4.197) (0.975,-4.216) (0.925,-4.233) (0.875,-4.250) (0.825,-4.266) (0.775,-4.280) (0.725,-4.294) (0.675,-4.306) (0.625,-4.318) (0.575,-4.329) (0.525,-4.338) (0.475,-4.347) (0.425,-4.355) (0.375,-4.362) (0.325,-4.368) (0.275,-4.374) (0.225,-4.378) (0.175,-4.381) (0.125,-4.384) (0.075,-4.386) (0.025,-4.387) (-0.050,-4.386) (-0.100,-4.385) (-0.150,-4.383) (-0.200,-4.380) (-0.250,-4.376) (-0.300,-4.371) (-0.350,-4.365) (-0.400,-4.359) (-0.450,-4.351) (-0.500,-4.343) (-0.550,-4.334) (-0.600,-4.323) (-0.650,-4.312) (-0.700,-4.300) (-0.750,-4.287) (-0.800,-4.273) (-0.850,-4.258) (-0.900,-4.242) (-0.949,-4.225) (-1.000,-4.207) (-1.050,-4.187) (-1.100,-4.167) (-1.150,-4.145) (-1.200,-4.123) (-1.248,-4.100) (-1.298,-4.075) (-1.345,-4.050) (-1.391,-4.025) (-1.434,-4.000) (-1.500,-3.960) (-1.550,-3.928) (-1.592,-3.900) (-1.650,-3.859) (-1.696,-3.825) (-1.750,-3.784) (-1.791,-3.750) (-1.850,-3.700) (-1.900,-3.655) (-1.950,-3.607) (-2.000,-3.557) (-2.050,-3.504) (-2.099,-3.450) (-2.141,-3.400) (-2.175,-3.358) (-2.219,-3.300) (-2.250,-3.258) (-2.290,-3.200) (-2.322,-3.150) (-2.350,-3.105) (-2.383,-3.050) (-2.410,-3.000) (-2.436,-2.950) (-2.461,-2.900) (-2.484,-2.850) (-2.506,-2.800) (-2.527,-2.750) (-2.546,-2.700) (-2.564,-2.650) (-2.580,-2.600) (-2.596,-2.550) (-2.610,-2.500) (-2.622,-2.450) (-2.634,-2.400) (-2.644,-2.350) (-2.653,-2.300) (-2.661,-2.250) (-2.668,-2.200) (-2.674,-2.150) (-2.678,-2.100) (-2.681,-2.050) (-2.683,-2.000) (-2.683,-1.950) (-2.683,-1.900) (-2.681,-1.850) (-2.678,-1.800) (-2.674,-1.750) (-2.668,-1.700) (-2.661,-1.650) (-2.653,-1.600) (-2.643,-1.550) (-2.632,-1.500) (-2.620,-1.450) (-2.606,-1.400) (-2.590,-1.350) (-2.573,-1.300) (-2.555,-1.250) (-2.534,-1.200) (-2.512,-1.150) (-2.488,-1.100) (-2.463,-1.050) (-2.435,-1.000) (-2.405,-0.950) (-2.375,-0.904) (-2.337,-0.850) (-2.300,-0.801) (-2.259,-0.750) (-2.225,-0.711) (-2.175,-0.657) (-2.125,-0.607) (-2.075,-0.561) (-2.033,-0.525) (-1.975,-0.479) (-1.925,-0.442) (-1.875,-0.408) (-1.825,-0.376) (-1.775,-0.346) (-1.725,-0.318) (-1.675,-0.292) (-1.625,-0.267) (-1.575,-0.244) (-1.525,-0.223) (-1.475,-0.203) (-1.425,-0.185) (-1.375,-0.168) (-1.325,-0.152) (-1.275,-0.137) (-1.225,-0.123) (-1.175,-0.111) (-1.125,-0.099) (-1.075,-0.088) (-1.025,-0.078) (-0.975,-0.070) (-0.925,-0.061) (-0.875,-0.054) (-0.825,-0.047) (-0.775,-0.041) (-0.725,-0.036) (-0.675,-0.031) (-0.625,-0.027) (-0.575,-0.023) (-0.525,-0.020) (-0.475,-0.017) (-0.425,-0.015) (-0.375,-0.013) (-0.325,-0.011) (-0.275,-0.010) (-0.250,-0.009)}
\EQ{(0.250,-0.008) (0.300,-0.009) (0.350,-0.011) (0.400,-0.012) (0.450,-0.014) (0.500,-0.017) (0.550,-0.019) (0.600,-0.022) (0.650,-0.026) (0.700,-0.030) (0.750,-0.034) (0.800,-0.039) (0.850,-0.045) (0.900,-0.051) (0.950,-0.058) (1.000,-0.066) (1.050,-0.074) (1.100,-0.083) (1.150,-0.093) (1.200,-0.104) (1.250,-0.115) (1.300,-0.128) (1.350,-0.141) (1.400,-0.156) (1.450,-0.172) (1.500,-0.188) (1.550,-0.206) (1.599,-0.225) (1.650,-0.246) (1.700,-0.268) (1.750,-0.292) (1.800,-0.317) (1.850,-0.344) (1.900,-0.372) (1.946,-0.400) (2.000,-0.435) (2.050,-0.470) (2.091,-0.500) (2.150,-0.547) (2.200,-0.589) (2.240,-0.625) (2.292,-0.675) (2.340,-0.725) (2.375,-0.764) (2.425,-0.824) (2.465,-0.875) (2.500,-0.924) (2.534,-0.975) (2.566,-1.025) (2.595,-1.075) (2.622,-1.125) (2.648,-1.175) (2.671,-1.225) (2.693,-1.275) (2.713,-1.325) (2.732,-1.375) (2.749,-1.425) (2.765,-1.475) (2.779,-1.525) (2.791,-1.575) (2.803,-1.625) (2.813,-1.675) (2.822,-1.725) (2.829,-1.775) (2.835,-1.825) (2.840,-1.875) (2.844,-1.925) (2.847,-1.975) (2.848,-2.025) (2.849,-2.075) (2.848,-2.125) (2.846,-2.175) (2.843,-2.225) (2.838,-2.275) (2.833,-2.325) (2.826,-2.375) (2.818,-2.425) (2.810,-2.475) (2.800,-2.525) (2.789,-2.575) (2.776,-2.625) (2.763,-2.675) (2.748,-2.725) (2.732,-2.775) (2.715,-2.825) (2.697,-2.875) (2.677,-2.925) (2.657,-2.975) (2.635,-3.025) (2.611,-3.075) (2.586,-3.125) (2.560,-3.175) (2.533,-3.225) (2.504,-3.275) (2.475,-3.322) (2.441,-3.375) (2.407,-3.425) (2.375,-3.470) (2.334,-3.525) (2.300,-3.569) (2.254,-3.625) (2.211,-3.675) (2.175,-3.715) (2.125,-3.767) (2.075,-3.818) (2.025,-3.865) (1.987,-3.900) (1.930,-3.950) (1.875,-3.996) (1.825,-4.035) (1.775,-4.073) (1.725,-4.109) (1.675,-4.143) (1.627,-4.175) (1.575,-4.208) (1.525,-4.238) (1.475,-4.266) (1.425,-4.294) (1.375,-4.320) (1.325,-4.345) (1.275,-4.369) (1.225,-4.392) (1.175,-4.413) (1.125,-4.434) (1.075,-4.453) (1.025,-4.472) (0.975,-4.489) (0.925,-4.506) (0.875,-4.521) (0.825,-4.536) (0.775,-4.549) (0.725,-4.562) (0.675,-4.574) (0.625,-4.585) (0.575,-4.595) (0.525,-4.604) (0.475,-4.612) (0.425,-4.620) (0.375,-4.626) (0.325,-4.632) (0.275,-4.637) (0.225,-4.641) (0.175,-4.644) (0.125,-4.647) (0.075,-4.649) (0.025,-4.649) (-0.050,-4.649) (-0.100,-4.648) (-0.150,-4.646) (-0.200,-4.643) (-0.250,-4.639) (-0.300,-4.635) (-0.350,-4.629) (-0.400,-4.623) (-0.450,-4.616) (-0.500,-4.608) (-0.550,-4.599) (-0.600,-4.590) (-0.650,-4.579) (-0.700,-4.568) (-0.750,-4.556) (-0.800,-4.543) (-0.850,-4.528) (-0.900,-4.513) (-0.950,-4.497) (-1.000,-4.480) (-1.050,-4.463) (-1.100,-4.444) (-1.150,-4.424) (-1.200,-4.402) (-1.250,-4.380) (-1.300,-4.357) (-1.350,-4.333) (-1.400,-4.307) (-1.450,-4.280) (-1.500,-4.252) (-1.546,-4.225) (-1.600,-4.192) (-1.650,-4.160) (-1.700,-4.126) (-1.750,-4.091) (-1.800,-4.054) (-1.850,-4.016) (-1.900,-3.975) (-1.950,-3.933) (-2.000,-3.888) (-2.041,-3.850) (-2.093,-3.800) (-2.142,-3.750) (-2.188,-3.700) (-2.225,-3.659) (-2.275,-3.600) (-2.315,-3.550) (-2.350,-3.504) (-2.389,-3.450) (-2.424,-3.400) (-2.457,-3.350) (-2.488,-3.300) (-2.518,-3.250) (-2.547,-3.200) (-2.573,-3.150) (-2.599,-3.100) (-2.623,-3.050) (-2.646,-3.000) (-2.667,-2.950) (-2.687,-2.900) (-2.706,-2.850) (-2.724,-2.800) (-2.740,-2.750) (-2.756,-2.700) (-2.770,-2.650) (-2.783,-2.600) (-2.794,-2.550) (-2.805,-2.500) (-2.814,-2.450) (-2.822,-2.400) (-2.830,-2.350) (-2.836,-2.300) (-2.841,-2.250) (-2.844,-2.200) (-2.847,-2.150) (-2.848,-2.100) (-2.849,-2.050) (-2.848,-2.000) (-2.846,-1.950) (-2.842,-1.900) (-2.838,-1.850) (-2.832,-1.800) (-2.826,-1.750) (-2.817,-1.700) (-2.808,-1.650) (-2.797,-1.600) (-2.785,-1.550) (-2.772,-1.500) (-2.757,-1.450) (-2.741,-1.400) (-2.723,-1.350) (-2.703,-1.300) (-2.683,-1.250) (-2.660,-1.200) (-2.635,-1.150) (-2.609,-1.100) (-2.581,-1.050) (-2.550,-1.000) (-2.518,-0.950) (-2.483,-0.900) (-2.450,-0.856) (-2.405,-0.800) (-2.362,-0.750) (-2.325,-0.709) (-2.275,-0.659) (-2.225,-0.611) (-2.184,-0.575) (-2.125,-0.526) (-2.075,-0.488) (-2.025,-0.452) (-1.975,-0.419) (-1.925,-0.387) (-1.875,-0.358) (-1.825,-0.330) (-1.775,-0.304) (-1.725,-0.280) (-1.675,-0.257) (-1.625,-0.236) (-1.575,-0.216) (-1.525,-0.197) (-1.475,-0.180) (-1.425,-0.164) (-1.375,-0.148) (-1.325,-0.134) (-1.275,-0.121) (-1.225,-0.109) (-1.175,-0.098) (-1.125,-0.088) (-1.075,-0.078) (-1.025,-0.070) (-0.975,-0.062) (-0.925,-0.055) (-0.875,-0.048) (-0.825,-0.042) (-0.775,-0.037) (-0.725,-0.032) (-0.675,-0.028) (-0.625,-0.024) (-0.575,-0.021) (-0.525,-0.018) (-0.475,-0.015) (-0.425,-0.013) (-0.375,-0.011) (-0.325,-0.010) (-0.275,-0.009) (-0.250,-0.008)}
\EQ{(0.250,-0.007) (0.300,-0.008) (0.350,-0.009) (0.400,-0.011) (0.450,-0.012) (0.500,-0.014) (0.550,-0.017) (0.600,-0.020) (0.650,-0.023) (0.700,-0.026) (0.750,-0.030) (0.800,-0.034) (0.850,-0.039) (0.900,-0.045) (0.950,-0.051) (1.000,-0.057) (1.050,-0.065) (1.100,-0.073) (1.150,-0.081) (1.200,-0.090) (1.250,-0.101) (1.300,-0.112) (1.350,-0.123) (1.400,-0.136) (1.450,-0.149) (1.500,-0.164) (1.550,-0.180) (1.600,-0.196) (1.650,-0.214) (1.700,-0.233) (1.750,-0.253) (1.800,-0.275) (1.850,-0.298) (1.900,-0.322) (1.950,-0.348) (2.000,-0.375) (2.050,-0.404) (2.100,-0.435) (2.150,-0.468) (2.196,-0.500) (2.250,-0.540) (2.294,-0.575) (2.350,-0.622) (2.400,-0.667) (2.450,-0.715) (2.500,-0.767) (2.550,-0.823) (2.594,-0.875) (2.625,-0.915) (2.669,-0.975) (2.700,-1.020) (2.735,-1.075) (2.765,-1.125) (2.793,-1.175) (2.820,-1.225) (2.844,-1.275) (2.867,-1.325) (2.888,-1.375) (2.908,-1.425) (2.926,-1.475) (2.943,-1.525) (2.959,-1.575) (2.973,-1.625) (2.986,-1.675) (2.997,-1.725) (3.008,-1.775) (3.017,-1.825) (3.024,-1.875) (3.031,-1.925) (3.037,-1.975) (3.041,-2.025) (3.044,-2.075) (3.047,-2.125) (3.048,-2.175) (3.048,-2.225) (3.047,-2.275) (3.044,-2.325) (3.041,-2.375) (3.037,-2.425) (3.032,-2.475) (3.025,-2.525) (3.018,-2.575) (3.009,-2.625) (3.000,-2.675) (2.989,-2.725) (2.977,-2.775) (2.964,-2.825) (2.950,-2.875) (2.935,-2.925) (2.919,-2.975) (2.902,-3.025) (2.883,-3.075) (2.864,-3.125) (2.843,-3.175) (2.821,-3.225) (2.798,-3.275) (2.775,-3.321) (2.750,-3.370) (2.725,-3.416) (2.700,-3.460) (2.662,-3.525) (2.630,-3.575) (2.600,-3.621) (2.563,-3.675) (2.527,-3.725) (2.489,-3.775) (2.450,-3.825) (2.408,-3.875) (2.375,-3.914) (2.325,-3.969) (2.275,-4.022) (2.225,-4.072) (2.175,-4.120) (2.125,-4.166) (2.087,-4.200) (2.028,-4.250) (1.975,-4.292) (1.933,-4.325) (1.875,-4.368) (1.829,-4.400) (1.775,-4.437) (1.725,-4.469) (1.676,-4.500) (1.625,-4.530) (1.575,-4.559) (1.525,-4.586) (1.475,-4.613) (1.425,-4.638) (1.375,-4.662) (1.325,-4.685) (1.275,-4.707) (1.225,-4.727) (1.175,-4.747) (1.125,-4.766) (1.075,-4.784) (1.025,-4.801) (0.975,-4.818) (0.925,-4.833) (0.875,-4.847) (0.825,-4.861) (0.775,-4.873) (0.725,-4.885) (0.675,-4.896) (0.625,-4.907) (0.575,-4.916) (0.525,-4.924) (0.475,-4.932) (0.425,-4.939) (0.375,-4.945) (0.325,-4.951) (0.275,-4.955) (0.225,-4.959) (0.175,-4.962) (0.125,-4.965) (0.075,-4.966) (0.025,-4.967) (-0.050,-4.967) (-0.100,-4.965) (-0.150,-4.963) (-0.200,-4.961) (-0.250,-4.957) (-0.300,-4.953) (-0.350,-4.948) (-0.400,-4.942) (-0.450,-4.936) (-0.500,-4.928) (-0.550,-4.920) (-0.600,-4.911) (-0.650,-4.902) (-0.700,-4.891) (-0.750,-4.880) (-0.800,-4.867) (-0.850,-4.854) (-0.900,-4.840) (-0.950,-4.825) (-1.000,-4.810) (-1.050,-4.793) (-1.100,-4.775) (-1.150,-4.757) (-1.200,-4.738) (-1.250,-4.717) (-1.300,-4.696) (-1.346,-4.675) (-1.400,-4.650) (-1.450,-4.625) (-1.499,-4.600) (-1.546,-4.575) (-1.591,-4.550) (-1.650,-4.516) (-1.700,-4.485) (-1.750,-4.453) (-1.793,-4.425) (-1.850,-4.385) (-1.899,-4.350) (-1.950,-4.312) (-1.997,-4.275) (-2.050,-4.231) (-2.100,-4.189) (-2.143,-4.150) (-2.197,-4.100) (-2.248,-4.050) (-2.296,-4.000) (-2.343,-3.950) (-2.387,-3.900) (-2.425,-3.855) (-2.470,-3.800) (-2.508,-3.750) (-2.545,-3.700) (-2.575,-3.658) (-2.614,-3.600) (-2.646,-3.550) (-2.675,-3.503) (-2.706,-3.450) (-2.734,-3.400) (-2.760,-3.350) (-2.786,-3.300) (-2.809,-3.250) (-2.832,-3.200) (-2.853,-3.150) (-2.874,-3.100) (-2.893,-3.050) (-2.911,-3.000) (-2.927,-2.950) (-2.943,-2.900) (-2.957,-2.850) (-2.971,-2.800) (-2.983,-2.750) (-2.994,-2.700) (-3.005,-2.650) (-3.014,-2.600) (-3.022,-2.550) (-3.029,-2.500) (-3.034,-2.450) (-3.039,-2.400) (-3.043,-2.350) (-3.046,-2.300) (-3.047,-2.250) (-3.048,-2.200) (-3.047,-2.150) (-3.046,-2.100) (-3.043,-2.050) (-3.039,-2.000) (-3.034,-1.950) (-3.028,-1.900) (-3.021,-1.850) (-3.012,-1.800) (-3.003,-1.750) (-2.992,-1.700) (-2.979,-1.650) (-2.966,-1.600) (-2.951,-1.550) (-2.935,-1.500) (-2.918,-1.450) (-2.898,-1.400) (-2.878,-1.350) (-2.856,-1.300) (-2.832,-1.250) (-2.807,-1.200) (-2.780,-1.150) (-2.751,-1.100) (-2.720,-1.050) (-2.686,-1.000) (-2.651,-0.950) (-2.613,-0.900) (-2.575,-0.852) (-2.530,-0.800) (-2.484,-0.750) (-2.435,-0.700) (-2.382,-0.650) (-2.325,-0.601) (-2.275,-0.560) (-2.230,-0.525) (-2.175,-0.485) (-2.125,-0.451) (-2.075,-0.419) (-2.025,-0.389) (-1.975,-0.361) (-1.925,-0.335) (-1.875,-0.310) (-1.825,-0.286) (-1.775,-0.264) (-1.725,-0.243) (-1.675,-0.223) (-1.625,-0.205) (-1.575,-0.188) (-1.525,-0.172) (-1.475,-0.157) (-1.425,-0.143) (-1.375,-0.129) (-1.325,-0.117) (-1.275,-0.106) (-1.225,-0.095) (-1.175,-0.086) (-1.125,-0.077) (-1.075,-0.069) (-1.025,-0.061) (-0.975,-0.054) (-0.925,-0.048) (-0.875,-0.042) (-0.825,-0.037) (-0.775,-0.032) (-0.725,-0.028) (-0.675,-0.024) (-0.625,-0.021) (-0.575,-0.018) (-0.525,-0.016) (-0.475,-0.013) (-0.425,-0.012) (-0.375,-0.010) (-0.325,-0.009) (-0.275,-0.008) (-0.250,-0.007)}
\EQ{(0.250,-0.006) (0.300,-0.007) (0.350,-0.008) (0.400,-0.009) (0.450,-0.011) (0.500,-0.012) (0.550,-0.014) (0.600,-0.017) (0.650,-0.019) (0.700,-0.022) (0.750,-0.026) (0.800,-0.030) (0.850,-0.034) (0.900,-0.038) (0.950,-0.044) (1.000,-0.049) (1.050,-0.055) (1.100,-0.062) (1.150,-0.069) (1.200,-0.077) (1.250,-0.086) (1.300,-0.095) (1.350,-0.105) (1.400,-0.116) (1.450,-0.128) (1.500,-0.140) (1.550,-0.153) (1.600,-0.167) (1.650,-0.183) (1.700,-0.199) (1.750,-0.216) (1.800,-0.234) (1.850,-0.253) (1.900,-0.273) (1.950,-0.295) (2.000,-0.318) (2.050,-0.342) (2.100,-0.367) (2.150,-0.394) (2.200,-0.423) (2.245,-0.450) (2.300,-0.485) (2.350,-0.519) (2.394,-0.550) (2.450,-0.593) (2.491,-0.625) (2.549,-0.675) (2.600,-0.721) (2.650,-0.770) (2.700,-0.821) (2.749,-0.875) (2.791,-0.925) (2.825,-0.967) (2.869,-1.025) (2.900,-1.069) (2.937,-1.125) (2.968,-1.175) (2.998,-1.225) (3.025,-1.274) (3.050,-1.322) (3.075,-1.373) (3.099,-1.425) (3.120,-1.475) (3.140,-1.525) (3.159,-1.575) (3.176,-1.625) (3.192,-1.675) (3.207,-1.725) (3.220,-1.775) (3.233,-1.825) (3.244,-1.875) (3.254,-1.925) (3.263,-1.975) (3.270,-2.025) (3.277,-2.075) (3.282,-2.125) (3.287,-2.175) (3.290,-2.225) (3.293,-2.275) (3.294,-2.325) (3.295,-2.375) (3.294,-2.425) (3.293,-2.475) (3.290,-2.525) (3.287,-2.575) (3.282,-2.625) (3.276,-2.675) (3.270,-2.725) (3.262,-2.775) (3.254,-2.825) (3.245,-2.875) (3.234,-2.925) (3.223,-2.975) (3.210,-3.025) (3.197,-3.075) (3.182,-3.125) (3.167,-3.175) (3.150,-3.225) (3.133,-3.275) (3.114,-3.325) (3.094,-3.375) (3.073,-3.425) (3.051,-3.475) (3.028,-3.525) (3.004,-3.575) (2.978,-3.625) (2.951,-3.675) (2.925,-3.722) (2.894,-3.775) (2.864,-3.825) (2.832,-3.875) (2.800,-3.923) (2.764,-3.975) (2.727,-4.025) (2.689,-4.075) (2.650,-4.125) (2.609,-4.175) (2.575,-4.214) (2.525,-4.270) (2.475,-4.324) (2.425,-4.375) (2.375,-4.424) (2.325,-4.471) (2.275,-4.516) (2.235,-4.550) (2.175,-4.600) (2.125,-4.640) (2.079,-4.675) (2.025,-4.715) (1.975,-4.750) (1.925,-4.784) (1.875,-4.817) (1.825,-4.848) (1.781,-4.875) (1.725,-4.907) (1.675,-4.935) (1.625,-4.962) (1.575,-4.988) (1.525,-5.013) (1.475,-5.037) (1.425,-5.059) (1.375,-5.081) (1.325,-5.102) (1.275,-5.122) (1.225,-5.141) (1.175,-5.159) (1.125,-5.177) (1.075,-5.193) (1.025,-5.209) (0.975,-5.223) (0.925,-5.238) (0.875,-5.251) (0.825,-5.263) (0.775,-5.275) (0.725,-5.286) (0.675,-5.296) (0.625,-5.305) (0.575,-5.314) (0.525,-5.322) (0.475,-5.329) (0.425,-5.335) (0.375,-5.341) (0.325,-5.346) (0.275,-5.350) (0.225,-5.354) (0.175,-5.357) (0.125,-5.359) (0.075,-5.360) (0.025,-5.361) (-0.050,-5.361) (-0.100,-5.360) (-0.150,-5.358) (-0.200,-5.355) (-0.250,-5.352) (-0.300,-5.348) (-0.350,-5.344) (-0.400,-5.338) (-0.450,-5.332) (-0.500,-5.325) (-0.550,-5.318) (-0.600,-5.310) (-0.650,-5.301) (-0.700,-5.291) (-0.750,-5.280) (-0.800,-5.269) (-0.850,-5.257) (-0.900,-5.244) (-0.950,-5.231) (-1.000,-5.216) (-1.050,-5.201) (-1.100,-5.185) (-1.150,-5.168) (-1.200,-5.150) (-1.250,-5.132) (-1.300,-5.112) (-1.350,-5.092) (-1.400,-5.070) (-1.450,-5.048) (-1.499,-5.025) (-1.550,-5.001) (-1.600,-4.975) (-1.648,-4.950) (-1.694,-4.925) (-1.738,-4.900) (-1.800,-4.863) (-1.850,-4.833) (-1.900,-4.801) (-1.950,-4.767) (-2.000,-4.733) (-2.045,-4.700) (-2.100,-4.659) (-2.144,-4.625) (-2.200,-4.580) (-2.250,-4.537) (-2.293,-4.500) (-2.347,-4.450) (-2.400,-4.400) (-2.450,-4.350) (-2.497,-4.300) (-2.543,-4.250) (-2.587,-4.200) (-2.625,-4.155) (-2.670,-4.100) (-2.708,-4.050) (-2.746,-4.000) (-2.775,-3.959) (-2.815,-3.900) (-2.848,-3.850) (-2.875,-3.807) (-2.909,-3.750) (-2.938,-3.700) (-2.965,-3.650) (-2.991,-3.600) (-3.016,-3.550) (-3.040,-3.500) (-3.062,-3.450) (-3.084,-3.400) (-3.104,-3.350) (-3.123,-3.300) (-3.141,-3.250) (-3.159,-3.200) (-3.175,-3.150) (-3.190,-3.100) (-3.204,-3.050) (-3.217,-3.000) (-3.228,-2.950) (-3.239,-2.900) (-3.249,-2.850) (-3.258,-2.800) (-3.266,-2.750) (-3.273,-2.700) (-3.279,-2.650) (-3.284,-2.600) (-3.288,-2.550) (-3.292,-2.500) (-3.294,-2.450) (-3.295,-2.400) (-3.295,-2.350) (-3.294,-2.300) (-3.292,-2.250) (-3.289,-2.200) (-3.285,-2.150) (-3.280,-2.100) (-3.274,-2.050) (-3.267,-2.000) (-3.258,-1.950) (-3.249,-1.900) (-3.238,-1.850) (-3.227,-1.800) (-3.214,-1.750) (-3.200,-1.701) (-3.184,-1.650) (-3.168,-1.600) (-3.150,-1.551) (-3.131,-1.500) (-3.110,-1.450) (-3.088,-1.400) (-3.064,-1.350) (-3.039,-1.300) (-3.012,-1.250) (-2.983,-1.200) (-2.953,-1.150) (-2.925,-1.107) (-2.887,-1.050) (-2.850,-1.000) (-2.811,-0.950) (-2.775,-0.906) (-2.727,-0.850) (-2.680,-0.800) (-2.630,-0.750) (-2.577,-0.700) (-2.525,-0.654) (-2.475,-0.613) (-2.427,-0.575) (-2.375,-0.537) (-2.325,-0.502) (-2.275,-0.469) (-2.225,-0.438) (-2.175,-0.408) (-2.125,-0.381) (-2.075,-0.354) (-2.025,-0.329) (-1.975,-0.306) (-1.925,-0.284) (-1.875,-0.263) (-1.825,-0.243) (-1.777,-0.225) (-1.725,-0.207) (-1.675,-0.190) (-1.625,-0.175) (-1.575,-0.160) (-1.525,-0.147) (-1.475,-0.134) (-1.425,-0.122) (-1.375,-0.111) (-1.325,-0.100) (-1.275,-0.091) (-1.225,-0.082) (-1.175,-0.073) (-1.125,-0.066) (-1.075,-0.059) (-1.025,-0.052) (-0.975,-0.046) (-0.925,-0.041) (-0.875,-0.036) (-0.825,-0.032) (-0.775,-0.028) (-0.725,-0.024) (-0.675,-0.021) (-0.625,-0.018) (-0.575,-0.016) (-0.525,-0.013) (-0.475,-0.012) (-0.425,-0.010) (-0.375,-0.009) (-0.325,-0.007) (-0.275,-0.007) (-0.250,-0.006)}
\EQ{(0.250,-0.005) (0.300,-0.006) (0.350,-0.007) (0.400,-0.008) (0.450,-0.009) (0.500,-0.010) (0.550,-0.012) (0.600,-0.014) (0.650,-0.016) (0.700,-0.019) (0.750,-0.021) (0.800,-0.025) (0.850,-0.028) (0.900,-0.032) (0.950,-0.036) (1.000,-0.041) (1.050,-0.046) (1.100,-0.052) (1.150,-0.058) (1.200,-0.064) (1.250,-0.072) (1.300,-0.079) (1.350,-0.088) (1.400,-0.097) (1.450,-0.106) (1.500,-0.116) (1.550,-0.127) (1.600,-0.139) (1.650,-0.151) (1.700,-0.165) (1.750,-0.179) (1.800,-0.193) (1.850,-0.209) (1.898,-0.225) (1.950,-0.243) (2.000,-0.262) (2.050,-0.282) (2.100,-0.302) (2.150,-0.324) (2.200,-0.347) (2.250,-0.371) (2.300,-0.397) (2.350,-0.423) (2.397,-0.450) (2.450,-0.481) (2.500,-0.513) (2.550,-0.546) (2.593,-0.575) (2.650,-0.617) (2.693,-0.650) (2.750,-0.696) (2.800,-0.739) (2.840,-0.775) (2.892,-0.825) (2.942,-0.875) (2.988,-0.925) (3.025,-0.967) (3.074,-1.025) (3.113,-1.075) (3.150,-1.125) (3.186,-1.175) (3.219,-1.225) (3.250,-1.274) (3.280,-1.325) (3.308,-1.375) (3.335,-1.425) (3.360,-1.475) (3.384,-1.525) (3.406,-1.575) (3.427,-1.625) (3.447,-1.675) (3.465,-1.725) (3.482,-1.775) (3.498,-1.825) (3.513,-1.875) (3.527,-1.925) (3.540,-1.975) (3.551,-2.025) (3.562,-2.075) (3.571,-2.125) (3.580,-2.175) (3.587,-2.225) (3.594,-2.275) (3.599,-2.325) (3.604,-2.375) (3.607,-2.425) (3.610,-2.475) (3.611,-2.525) (3.612,-2.575) (3.612,-2.625) (3.611,-2.675) (3.609,-2.725) (3.606,-2.775) (3.602,-2.825) (3.598,-2.875) (3.592,-2.925) (3.586,-2.975) (3.578,-3.025) (3.570,-3.075) (3.561,-3.125) (3.551,-3.175) (3.540,-3.225) (3.528,-3.275) (3.515,-3.325) (3.501,-3.375) (3.487,-3.425) (3.471,-3.475) (3.455,-3.525) (3.437,-3.575) (3.419,-3.625) (3.400,-3.673) (3.379,-3.725) (3.357,-3.775) (3.334,-3.825) (3.311,-3.875) (3.286,-3.925) (3.260,-3.975) (3.234,-4.025) (3.205,-4.075) (3.176,-4.125) (3.150,-4.168) (3.114,-4.225) (3.081,-4.275) (3.050,-4.321) (3.011,-4.375) (2.975,-4.424) (2.936,-4.475) (2.900,-4.520) (2.855,-4.575) (2.811,-4.625) (2.775,-4.666) (2.725,-4.720) (2.675,-4.771) (2.625,-4.821) (2.575,-4.869) (2.525,-4.915) (2.485,-4.950) (2.426,-5.000) (2.375,-5.042) (2.334,-5.075) (2.275,-5.120) (2.234,-5.150) (2.175,-5.192) (2.127,-5.225) (2.075,-5.260) (2.025,-5.291) (1.975,-5.322) (1.928,-5.350) (1.885,-5.375) (1.840,-5.400) (1.794,-5.425) (1.746,-5.450) (1.700,-5.473) (1.650,-5.497) (1.600,-5.521) (1.550,-5.543) (1.500,-5.565) (1.450,-5.586) (1.400,-5.606) (1.350,-5.625) (1.300,-5.643) (1.250,-5.661) (1.200,-5.677) (1.150,-5.693) (1.100,-5.709) (1.050,-5.723) (1.000,-5.737) (0.950,-5.750) (0.900,-5.762) (0.850,-5.774) (0.800,-5.785) (0.750,-5.795) (0.700,-5.805) (0.650,-5.813) (0.600,-5.822) (0.550,-5.829) (0.500,-5.836) (0.450,-5.842) (0.400,-5.848) (0.350,-5.853) (0.300,-5.857) (0.250,-5.860) (0.200,-5.863) (0.150,-5.866) (0.100,-5.867) (0.050,-5.868) (0.000,-5.868) (-0.050,-5.868) (-0.100,-5.867) (-0.150,-5.866) (-0.200,-5.863) (-0.250,-5.860) (-0.300,-5.857) (-0.350,-5.853) (-0.400,-5.848) (-0.450,-5.842) (-0.500,-5.836) (-0.550,-5.829) (-0.600,-5.822) (-0.650,-5.813) (-0.700,-5.805) (-0.750,-5.795) (-0.800,-5.785) (-0.850,-5.774) (-0.900,-5.762) (-0.950,-5.750) (-1.000,-5.737) (-1.050,-5.723) (-1.100,-5.709) (-1.150,-5.693) (-1.200,-5.677) (-1.250,-5.661) (-1.300,-5.643) (-1.349,-5.625) (-1.400,-5.606) (-1.450,-5.586) (-1.500,-5.565) (-1.550,-5.543) (-1.600,-5.521) (-1.650,-5.497) (-1.696,-5.475) (-1.746,-5.450) (-1.794,-5.425) (-1.840,-5.400) (-1.885,-5.375) (-1.928,-5.350) (-1.975,-5.322) (-2.025,-5.291) (-2.075,-5.260) (-2.125,-5.226) (-2.175,-5.192) (-2.225,-5.157) (-2.268,-5.125) (-2.325,-5.082) (-2.365,-5.050) (-2.425,-5.001) (-2.475,-4.959) (-2.513,-4.925) (-2.568,-4.875) (-2.621,-4.825) (-2.671,-4.775) (-2.720,-4.725) (-2.767,-4.675) (-2.811,-4.625) (-2.850,-4.580) (-2.896,-4.525) (-2.936,-4.475) (-2.974,-4.425) (-3.011,-4.375) (-3.047,-4.325) (-3.081,-4.275) (-3.114,-4.225) (-3.146,-4.175) (-3.175,-4.127) (-3.205,-4.075) (-3.234,-4.025) (-3.260,-3.975) (-3.286,-3.925) (-3.311,-3.875) (-3.334,-3.825) (-3.357,-3.775) (-3.379,-3.725) (-3.399,-3.675) (-3.419,-3.625) (-3.437,-3.575) (-3.455,-3.525) (-3.471,-3.475) (-3.487,-3.425) (-3.501,-3.375) (-3.515,-3.325) (-3.528,-3.275) (-3.540,-3.225) (-3.551,-3.175) (-3.561,-3.125) (-3.570,-3.075) (-3.578,-3.025) (-3.586,-2.975) (-3.592,-2.925) (-3.598,-2.875) (-3.602,-2.825) (-3.606,-2.775) (-3.609,-2.725) (-3.611,-2.675) (-3.612,-2.625) (-3.612,-2.575) (-3.611,-2.525) (-3.610,-2.475) (-3.607,-2.425) (-3.604,-2.375) (-3.599,-2.325) (-3.594,-2.275) (-3.587,-2.225) (-3.580,-2.175) (-3.571,-2.125) (-3.562,-2.075) (-3.551,-2.025) (-3.540,-1.975) (-3.527,-1.925) (-3.513,-1.875) (-3.498,-1.825) (-3.482,-1.775) (-3.465,-1.725) (-3.447,-1.675) (-3.427,-1.625) (-3.406,-1.575) (-3.384,-1.525) (-3.360,-1.475) (-3.335,-1.425) (-3.308,-1.375) (-3.280,-1.325) (-3.250,-1.275) (-3.219,-1.225) (-3.186,-1.175) (-3.150,-1.125) (-3.113,-1.075) (-3.075,-1.026) (-3.032,-0.975) (-2.988,-0.925) (-2.950,-0.884) (-2.900,-0.833) (-2.850,-0.784) (-2.813,-0.750) (-2.755,-0.700) (-2.700,-0.655) (-2.650,-0.617) (-2.600,-0.580) (-2.557,-0.550) (-2.500,-0.513) (-2.450,-0.481) (-2.400,-0.452) (-2.353,-0.425) (-2.307,-0.400) (-2.258,-0.375) (-2.207,-0.350) (-2.152,-0.325) (-2.100,-0.302) (-2.050,-0.282) (-2.000,-0.262) (-1.950,-0.243) (-1.900,-0.226) (-1.850,-0.209) (-1.800,-0.193) (-1.750,-0.179) (-1.700,-0.165) (-1.650,-0.151) (-1.600,-0.139) (-1.550,-0.127) (-1.500,-0.116) (-1.450,-0.106) (-1.400,-0.097) (-1.350,-0.088) (-1.300,-0.079) (-1.250,-0.072) (-1.200,-0.064) (-1.150,-0.058) (-1.100,-0.052) (-1.050,-0.046) (-1.000,-0.041) (-0.950,-0.036) (-0.900,-0.032) (-0.850,-0.028) (-0.800,-0.025) (-0.750,-0.021) (-0.700,-0.019) (-0.650,-0.016) (-0.600,-0.014) (-0.550,-0.012) (-0.500,-0.010) (-0.450,-0.009) (-0.400,-0.008) (-0.350,-0.007) (-0.300,-0.006) (-0.250,-0.005)}
\EQ{(0.250,-0.004) (0.300,-0.005) (0.350,-0.005) (0.400,-0.006) (0.450,-0.007) (0.500,-0.008) (0.550,-0.010) (0.600,-0.011) (0.650,-0.013) (0.700,-0.015) (0.750,-0.017) (0.800,-0.020) (0.850,-0.022) (0.900,-0.026) (0.950,-0.029) (1.000,-0.033) (1.050,-0.037) (1.100,-0.041) (1.150,-0.046) (1.200,-0.052) (1.250,-0.057) (1.300,-0.063) (1.350,-0.070) (1.400,-0.077) (1.450,-0.085) (1.500,-0.093) (1.550,-0.102) (1.600,-0.111) (1.650,-0.121) (1.700,-0.131) (1.750,-0.142) (1.800,-0.154) (1.850,-0.166) (1.900,-0.179) (1.950,-0.193) (2.000,-0.208) (2.050,-0.223) (2.100,-0.239) (2.150,-0.256) (2.200,-0.274) (2.250,-0.293) (2.300,-0.313) (2.350,-0.333) (2.400,-0.355) (2.450,-0.378) (2.497,-0.400) (2.548,-0.425) (2.596,-0.450) (2.642,-0.475) (2.686,-0.500) (2.750,-0.538) (2.800,-0.570) (2.847,-0.600) (2.900,-0.637) (2.950,-0.673) (3.000,-0.711) (3.050,-0.750) (3.100,-0.792) (3.138,-0.825) (3.193,-0.875) (3.245,-0.925) (3.294,-0.975) (3.341,-1.025) (3.375,-1.063) (3.425,-1.122) (3.467,-1.175) (3.500,-1.218) (3.542,-1.275) (3.575,-1.323) (3.609,-1.375) (3.640,-1.425) (3.670,-1.475) (3.698,-1.525) (3.725,-1.574) (3.750,-1.623) (3.775,-1.675) (3.798,-1.725) (3.820,-1.775) (3.841,-1.825) (3.860,-1.875) (3.878,-1.925) (3.896,-1.975) (3.912,-2.025) (3.927,-2.075) (3.941,-2.125) (3.954,-2.175) (3.966,-2.225) (3.977,-2.275) (3.988,-2.325) (3.997,-2.375) (4.005,-2.425) (4.013,-2.475) (4.020,-2.525) (4.025,-2.575) (4.030,-2.625) (4.034,-2.675) (4.037,-2.725) (4.040,-2.775) (4.041,-2.825) (4.042,-2.875) (4.042,-2.925) (4.041,-2.975) (4.039,-3.025) (4.037,-3.075) (4.033,-3.125) (4.029,-3.175) (4.024,-3.225) (4.018,-3.275) (4.012,-3.325) (4.004,-3.375) (3.996,-3.425) (3.987,-3.475) (3.977,-3.525) (3.967,-3.575) (3.955,-3.625) (3.943,-3.675) (3.930,-3.725) (3.916,-3.775) (3.901,-3.825) (3.886,-3.875) (3.869,-3.925) (3.852,-3.975) (3.834,-4.025) (3.815,-4.075) (3.795,-4.125) (3.775,-4.172) (3.752,-4.225) (3.729,-4.275) (3.706,-4.325) (3.681,-4.375) (3.655,-4.425) (3.629,-4.475) (3.601,-4.525) (3.575,-4.571) (3.543,-4.625) (3.512,-4.675) (3.480,-4.725) (3.450,-4.771) (3.413,-4.825) (3.378,-4.875) (3.341,-4.925) (3.303,-4.975) (3.264,-5.025) (3.225,-5.073) (3.181,-5.125) (3.138,-5.175) (3.100,-5.217) (3.050,-5.271) (3.000,-5.322) (2.950,-5.372) (2.900,-5.420) (2.850,-5.467) (2.800,-5.512) (2.756,-5.550) (2.700,-5.598) (2.650,-5.638) (2.603,-5.675) (2.550,-5.716) (2.504,-5.750) (2.450,-5.789) (2.400,-5.823) (2.350,-5.857) (2.300,-5.889) (2.250,-5.921) (2.202,-5.950) (2.150,-5.981) (2.100,-6.010) (2.050,-6.037) (2.000,-6.064) (1.950,-6.090) (1.900,-6.115) (1.850,-6.139) (1.800,-6.163) (1.750,-6.186) (1.700,-6.208) (1.650,-6.229) (1.600,-6.249) (1.550,-6.269) (1.500,-6.288) (1.450,-6.306) (1.400,-6.323) (1.350,-6.340) (1.300,-6.356) (1.250,-6.372) (1.200,-6.387) (1.150,-6.401) (1.100,-6.414) (1.050,-6.427) (1.000,-6.439) (0.950,-6.451) (0.900,-6.462) (0.850,-6.472) (0.800,-6.482) (0.750,-6.491) (0.700,-6.499) (0.650,-6.507) (0.600,-6.515) (0.550,-6.521) (0.500,-6.527) (0.450,-6.533) (0.400,-6.538) (0.350,-6.542) (0.300,-6.546) (0.250,-6.549) (0.200,-6.552) (0.150,-6.554) (0.100,-6.555) (0.050,-6.556) (0.000,-6.556) (-0.050,-6.556) (-0.100,-6.555) (-0.150,-6.554) (-0.200,-6.552) (-0.250,-6.549) (-0.300,-6.546) (-0.350,-6.542) (-0.400,-6.538) (-0.450,-6.533) (-0.500,-6.527) (-0.550,-6.521) (-0.600,-6.515) (-0.650,-6.507) (-0.700,-6.499) (-0.750,-6.491) (-0.800,-6.482) (-0.850,-6.472) (-0.900,-6.462) (-0.950,-6.451) (-1.000,-6.439) (-1.050,-6.427) (-1.100,-6.414) (-1.150,-6.401) (-1.200,-6.387) (-1.250,-6.372) (-1.300,-6.356) (-1.350,-6.340) (-1.400,-6.323) (-1.450,-6.306) (-1.500,-6.288) (-1.550,-6.269) (-1.598,-6.250) (-1.650,-6.229) (-1.700,-6.208) (-1.750,-6.186) (-1.800,-6.163) (-1.850,-6.139) (-1.900,-6.115) (-1.950,-6.090) (-2.000,-6.064) (-2.050,-6.037) (-2.100,-6.010) (-2.150,-5.981) (-2.200,-5.951) (-2.243,-5.925) (-2.300,-5.889) (-2.350,-5.857) (-2.397,-5.825) (-2.450,-5.789) (-2.500,-5.753) (-2.550,-5.716) (-2.600,-5.678) (-2.650,-5.638) (-2.697,-5.600) (-2.750,-5.556) (-2.800,-5.512) (-2.841,-5.475) (-2.895,-5.425) (-2.947,-5.375) (-2.997,-5.325) (-3.046,-5.275) (-3.092,-5.225) (-3.138,-5.175) (-3.175,-5.132) (-3.223,-5.075) (-3.264,-5.025) (-3.300,-4.979) (-3.341,-4.925) (-3.375,-4.879) (-3.413,-4.825) (-3.447,-4.775) (-3.475,-4.733) (-3.512,-4.675) (-3.543,-4.625) (-3.573,-4.575) (-3.600,-4.527) (-3.625,-4.482) (-3.650,-4.435) (-3.675,-4.387) (-3.700,-4.337) (-3.725,-4.284) (-3.750,-4.229) (-3.774,-4.175) (-3.795,-4.125) (-3.815,-4.075) (-3.834,-4.025) (-3.852,-3.975) (-3.869,-3.925) (-3.886,-3.875) (-3.901,-3.825) (-3.916,-3.775) (-3.930,-3.725) (-3.943,-3.675) (-3.955,-3.625) (-3.967,-3.575) (-3.977,-3.525) (-3.987,-3.475) (-3.996,-3.425) (-4.004,-3.375) (-4.012,-3.325) (-4.018,-3.275) (-4.024,-3.225) (-4.029,-3.175) (-4.033,-3.125) (-4.037,-3.075) (-4.039,-3.025) (-4.041,-2.975) (-4.042,-2.925) (-4.042,-2.875) (-4.041,-2.825) (-4.040,-2.775) (-4.037,-2.725) (-4.034,-2.675) (-4.030,-2.625) (-4.025,-2.575) (-4.020,-2.525) (-4.013,-2.475) (-4.005,-2.425) (-3.997,-2.375) (-3.988,-2.325) (-3.977,-2.275) (-3.966,-2.225) (-3.954,-2.175) (-3.941,-2.125) (-3.927,-2.075) (-3.912,-2.025) (-3.896,-1.975) (-3.878,-1.925) (-3.860,-1.875) (-3.841,-1.825) (-3.820,-1.775) (-3.798,-1.725) (-3.775,-1.675) (-3.751,-1.625) (-3.725,-1.575) (-3.700,-1.528) (-3.675,-1.483) (-3.640,-1.425) (-3.609,-1.375) (-3.576,-1.325) (-3.542,-1.275) (-3.505,-1.225) (-3.475,-1.185) (-3.427,-1.125) (-3.385,-1.075) (-3.350,-1.035) (-3.300,-0.981) (-3.250,-0.930) (-3.200,-0.882) (-3.150,-0.836) (-3.110,-0.800) (-3.050,-0.750) (-3.000,-0.711) (-2.953,-0.675) (-2.900,-0.637) (-2.850,-0.602) (-2.800,-0.570) (-2.750,-0.538) (-2.700,-0.508) (-2.650,-0.480) (-2.600,-0.452) (-2.550,-0.426) (-2.500,-0.401) (-2.450,-0.378) (-2.400,-0.355) (-2.350,-0.333) (-2.300,-0.313) (-2.250,-0.293) (-2.202,-0.275) (-2.150,-0.256) (-2.100,-0.239) (-2.050,-0.223) (-2.000,-0.208) (-1.950,-0.193) (-1.900,-0.179) (-1.850,-0.166) (-1.800,-0.154) (-1.750,-0.142) (-1.700,-0.131) (-1.650,-0.121) (-1.600,-0.111) (-1.550,-0.102) (-1.500,-0.093) (-1.450,-0.085) (-1.400,-0.077) (-1.350,-0.070) (-1.300,-0.063) (-1.250,-0.057) (-1.200,-0.052) (-1.150,-0.046) (-1.100,-0.041) (-1.050,-0.037) (-1.000,-0.033) (-0.950,-0.029) (-0.900,-0.026) (-0.850,-0.022) (-0.800,-0.020) (-0.750,-0.017) (-0.700,-0.015) (-0.650,-0.013) (-0.600,-0.011) (-0.550,-0.010) (-0.500,-0.008) (-0.450,-0.007) (-0.400,-0.006) (-0.350,-0.005) (-0.300,-0.005) (-0.250,-0.004)}
\EQ{(0.250,-0.003) (0.300,-0.003) (0.350,-0.004) (0.400,-0.005) (0.450,-0.005) (0.500,-0.006) (0.550,-0.007) (0.600,-0.008) (0.650,-0.010) (0.700,-0.011) (0.750,-0.013) (0.800,-0.015) (0.850,-0.017) (0.900,-0.019) (0.950,-0.022) (1.000,-0.025) (1.050,-0.028) (1.100,-0.031) (1.150,-0.035) (1.200,-0.039) (1.250,-0.043) (1.300,-0.047) (1.350,-0.052) (1.400,-0.058) (1.450,-0.063) (1.500,-0.070) (1.550,-0.076) (1.600,-0.083) (1.650,-0.090) (1.700,-0.098) (1.750,-0.106) (1.800,-0.115) (1.850,-0.124) (1.900,-0.134) (1.950,-0.144) (2.000,-0.155) (2.050,-0.166) (2.100,-0.178) (2.150,-0.191) (2.200,-0.204) (2.250,-0.218) (2.300,-0.232) (2.350,-0.247) (2.400,-0.263) (2.450,-0.279) (2.500,-0.296) (2.550,-0.314) (2.600,-0.333) (2.650,-0.353) (2.700,-0.373) (2.750,-0.394) (2.800,-0.416) (2.850,-0.439) (2.900,-0.463) (2.950,-0.488) (3.000,-0.514) (3.050,-0.541) (3.100,-0.569) (3.150,-0.599) (3.194,-0.625) (3.250,-0.661) (3.300,-0.694) (3.346,-0.725) (3.400,-0.764) (3.448,-0.800) (3.500,-0.840) (3.543,-0.875) (3.600,-0.923) (3.650,-0.967) (3.700,-1.013) (3.738,-1.050) (3.775,-1.087) (3.812,-1.125) (3.850,-1.165) (3.900,-1.221) (3.946,-1.275) (3.987,-1.325) (4.025,-1.373) (4.064,-1.425) (4.100,-1.475) (4.134,-1.525) (4.167,-1.575) (4.199,-1.625) (4.225,-1.668) (4.258,-1.725) (4.286,-1.775) (4.312,-1.825) (4.338,-1.875) (4.362,-1.925) (4.385,-1.975) (4.407,-2.025) (4.428,-2.075) (4.448,-2.125) (4.467,-2.175) (4.485,-2.225) (4.502,-2.275) (4.519,-2.325) (4.534,-2.375) (4.548,-2.425) (4.562,-2.475) (4.575,-2.525) (4.587,-2.575) (4.598,-2.625) (4.608,-2.675) (4.617,-2.725) (4.626,-2.775) (4.634,-2.825) (4.641,-2.875) (4.647,-2.925) (4.653,-2.975) (4.658,-3.025) (4.662,-3.075) (4.665,-3.125) (4.668,-3.175) (4.670,-3.225) (4.671,-3.275) (4.671,-3.325) (4.671,-3.375) (4.670,-3.425) (4.668,-3.475) (4.666,-3.525) (4.663,-3.575) (4.659,-3.625) (4.655,-3.675) (4.650,-3.725) (4.644,-3.775) (4.637,-3.825) (4.630,-3.875) (4.622,-3.925) (4.613,-3.975) (4.604,-4.025) (4.594,-4.075) (4.583,-4.125) (4.572,-4.175) (4.560,-4.225) (4.547,-4.275) (4.533,-4.325) (4.519,-4.375) (4.504,-4.425) (4.488,-4.475) (4.472,-4.525) (4.454,-4.575) (4.436,-4.625) (4.418,-4.675) (4.398,-4.725) (4.378,-4.775) (4.357,-4.825) (4.335,-4.875) (4.312,-4.925) (4.288,-4.975) (4.264,-5.025) (4.239,-5.075) (4.213,-5.125) (4.186,-5.175) (4.158,-5.225) (4.130,-5.275) (4.100,-5.325) (4.069,-5.375) (4.038,-5.425) (4.005,-5.475) (3.975,-5.520) (3.937,-5.575) (3.902,-5.625) (3.865,-5.675) (3.827,-5.725) (3.789,-5.775) (3.750,-5.823) (3.707,-5.875) (3.665,-5.925) (3.625,-5.970) (3.576,-6.025) (3.529,-6.075) (3.481,-6.125) (3.432,-6.175) (3.381,-6.225) (3.328,-6.275) (3.275,-6.324) (3.225,-6.368) (3.175,-6.412) (3.129,-6.450) (3.075,-6.495) (3.025,-6.534) (2.975,-6.573) (2.925,-6.611) (2.875,-6.647) (2.825,-6.683) (2.775,-6.717) (2.726,-6.750) (2.675,-6.784) (2.625,-6.816) (2.575,-6.846) (2.528,-6.875) (2.475,-6.906) (2.425,-6.935) (2.375,-6.962) (2.325,-6.989) (2.275,-7.015) (2.225,-7.041) (2.175,-7.066) (2.125,-7.090) (2.075,-7.113) (2.025,-7.136) (1.975,-7.158) (1.925,-7.179) (1.875,-7.200) (1.825,-7.220) (1.775,-7.239) (1.725,-7.258) (1.678,-7.275) (1.625,-7.294) (1.575,-7.311) (1.525,-7.327) (1.475,-7.343) (1.425,-7.358) (1.375,-7.372) (1.325,-7.386) (1.275,-7.400) (1.225,-7.413) (1.176,-7.425) (1.125,-7.437) (1.075,-7.448) (1.025,-7.459) (0.975,-7.469) (0.925,-7.479) (0.875,-7.488) (0.825,-7.497) (0.775,-7.505) (0.725,-7.512) (0.675,-7.519) (0.625,-7.526) (0.575,-7.532) (0.525,-7.538) (0.475,-7.543) (0.425,-7.547) (0.375,-7.551) (0.325,-7.555) (0.275,-7.558) (0.225,-7.560) (0.175,-7.562) (0.125,-7.564) (0.075,-7.565) (0.025,-7.565) (-0.050,-7.565) (-0.100,-7.564) (-0.150,-7.563) (-0.200,-7.561) (-0.250,-7.559) (-0.300,-7.556) (-0.350,-7.553) (-0.400,-7.549) (-0.450,-7.545) (-0.500,-7.540) (-0.550,-7.535) (-0.600,-7.529) (-0.650,-7.523) (-0.700,-7.516) (-0.750,-7.509) (-0.800,-7.501) (-0.850,-7.492) (-0.900,-7.484) (-0.950,-7.474) (-1.000,-7.464) (-1.050,-7.454) (-1.100,-7.443) (-1.150,-7.431) (-1.200,-7.419) (-1.250,-7.406) (-1.300,-7.393) (-1.350,-7.380) (-1.400,-7.365) (-1.450,-7.350) (-1.500,-7.335) (-1.550,-7.319) (-1.600,-7.302) (-1.650,-7.285) (-1.700,-7.267) (-1.750,-7.249) (-1.800,-7.229) (-1.850,-7.210) (-1.900,-7.189) (-1.950,-7.168) (-2.000,-7.147) (-2.049,-7.125) (-2.100,-7.101) (-2.150,-7.078) (-2.200,-7.053) (-2.250,-7.028) (-2.300,-7.002) (-2.350,-6.976) (-2.397,-6.950) (-2.442,-6.925) (-2.485,-6.900) (-2.550,-6.862) (-2.600,-6.831) (-2.649,-6.800) (-2.700,-6.767) (-2.750,-6.734) (-2.800,-6.700) (-2.850,-6.665) (-2.900,-6.629) (-2.950,-6.592) (-3.000,-6.554) (-3.050,-6.515) (-3.099,-6.475) (-3.150,-6.433) (-3.188,-6.400) (-3.246,-6.350) (-3.300,-6.301) (-3.350,-6.254) (-3.400,-6.206) (-3.450,-6.157) (-3.500,-6.106) (-3.550,-6.053) (-3.599,-6.000) (-3.643,-5.950) (-3.686,-5.900) (-3.725,-5.854) (-3.769,-5.800) (-3.800,-5.761) (-3.847,-5.700) (-3.884,-5.650) (-3.920,-5.600) (-3.950,-5.557) (-3.989,-5.500) (-4.022,-5.450) (-4.050,-5.406) (-4.085,-5.350) (-4.115,-5.300) (-4.144,-5.250) (-4.172,-5.200) (-4.200,-5.150) (-4.225,-5.102) (-4.250,-5.053) (-4.275,-5.003) (-4.300,-4.951) (-4.323,-4.900) (-4.346,-4.850) (-4.367,-4.800) (-4.388,-4.750) (-4.408,-4.700) (-4.427,-4.650) (-4.445,-4.600) (-4.463,-4.550) (-4.480,-4.500) (-4.496,-4.450) (-4.511,-4.400) (-4.526,-4.350) (-4.540,-4.300) (-4.553,-4.250) (-4.566,-4.200) (-4.578,-4.150) (-4.589,-4.100) (-4.599,-4.050) (-4.609,-4.000) (-4.618,-3.950) (-4.626,-3.900) (-4.634,-3.850) (-4.641,-3.800) (-4.647,-3.750) (-4.652,-3.700) (-4.657,-3.650) (-4.661,-3.600) (-4.665,-3.550) (-4.667,-3.500) (-4.669,-3.450) (-4.671,-3.400) (-4.671,-3.350) (-4.671,-3.300) (-4.670,-3.250) (-4.669,-3.200) (-4.667,-3.150) (-4.664,-3.100) (-4.660,-3.050) (-4.655,-3.000) (-4.650,-2.950) (-4.644,-2.900) (-4.638,-2.850) (-4.630,-2.800) (-4.622,-2.750) (-4.613,-2.700) (-4.603,-2.650) (-4.592,-2.600) (-4.581,-2.550) (-4.569,-2.500) (-4.555,-2.450) (-4.541,-2.400) (-4.526,-2.350) (-4.511,-2.300) (-4.494,-2.250) (-4.476,-2.200) (-4.458,-2.150) (-4.438,-2.100) (-4.418,-2.050) (-4.396,-2.000) (-4.375,-1.953) (-4.350,-1.900) (-4.325,-1.850) (-4.300,-1.802) (-4.275,-1.755) (-4.250,-1.711) (-4.214,-1.650) (-4.183,-1.600) (-4.151,-1.550) (-4.117,-1.500) (-4.082,-1.450) (-4.050,-1.406) (-4.007,-1.350) (-3.975,-1.310) (-3.925,-1.250) (-3.882,-1.200) (-3.836,-1.150) (-3.800,-1.112) (-3.764,-1.075) (-3.725,-1.037) (-3.686,-1.000) (-3.631,-0.950) (-3.575,-0.902) (-3.525,-0.860) (-3.481,-0.825) (-3.425,-0.782) (-3.381,-0.750) (-3.325,-0.711) (-3.275,-0.677) (-3.225,-0.645) (-3.175,-0.614) (-3.125,-0.584) (-3.075,-0.555) (-3.025,-0.527) (-2.975,-0.501) (-2.925,-0.476) (-2.875,-0.451) (-2.825,-0.428) (-2.775,-0.405) (-2.725,-0.383) (-2.675,-0.363) (-2.625,-0.343) (-2.575,-0.324) (-2.525,-0.305) (-2.475,-0.288) (-2.425,-0.271) (-2.375,-0.255) (-2.325,-0.239) (-2.276,-0.225) (-2.225,-0.211) (-2.175,-0.197) (-2.125,-0.184) (-2.075,-0.172) (-2.025,-0.161) (-1.975,-0.149) (-1.925,-0.139) (-1.875,-0.129) (-1.825,-0.120) (-1.775,-0.111) (-1.725,-0.102) (-1.675,-0.094) (-1.625,-0.087) (-1.575,-0.079) (-1.525,-0.073) (-1.475,-0.066) (-1.425,-0.061) (-1.375,-0.055) (-1.325,-0.050) (-1.275,-0.045) (-1.225,-0.041) (-1.175,-0.037) (-1.125,-0.033) (-1.075,-0.029) (-1.025,-0.026) (-0.975,-0.023) (-0.925,-0.020) (-0.875,-0.018) (-0.825,-0.016) (-0.775,-0.014) (-0.725,-0.012) (-0.675,-0.010) (-0.625,-0.009) (-0.575,-0.008) (-0.525,-0.007) (-0.475,-0.006) (-0.425,-0.005) (-0.375,-0.004) (-0.325,-0.004) (-0.275,-0.003) (-0.250,-0.003)}
\EQ{(0.250,-0.002) (0.300,-0.002) (0.350,-0.003) (0.400,-0.003) (0.450,-0.004) (0.500,-0.004) (0.550,-0.005) (0.600,-0.006) (0.650,-0.006) (0.700,-0.007) (0.750,-0.009) (0.800,-0.010) (0.850,-0.011) (0.900,-0.013) (0.950,-0.014) (1.000,-0.016) (1.050,-0.018) (1.100,-0.021) (1.150,-0.023) (1.200,-0.026) (1.250,-0.029) (1.300,-0.032) (1.350,-0.035) (1.400,-0.038) (1.450,-0.042) (1.500,-0.046) (1.550,-0.051) (1.600,-0.055) (1.650,-0.060) (1.700,-0.065) (1.750,-0.071) (1.800,-0.076) (1.850,-0.083) (1.900,-0.089) (1.950,-0.096) (2.000,-0.103) (2.050,-0.110) (2.100,-0.118) (2.150,-0.126) (2.200,-0.135) (2.250,-0.144) (2.300,-0.153) (2.350,-0.163) (2.400,-0.174) (2.450,-0.184) (2.500,-0.196) (2.550,-0.207) (2.600,-0.219) (2.650,-0.232) (2.700,-0.245) (2.750,-0.259) (2.800,-0.273) (2.850,-0.287) (2.900,-0.303) (2.950,-0.319) (3.000,-0.335) (3.050,-0.352) (3.100,-0.369) (3.150,-0.388) (3.200,-0.407) (3.247,-0.425) (3.300,-0.446) (3.350,-0.467) (3.400,-0.489) (3.450,-0.511) (3.500,-0.534) (3.550,-0.558) (3.600,-0.583) (3.650,-0.608) (3.700,-0.635) (3.750,-0.662) (3.800,-0.690) (3.850,-0.719) (3.900,-0.749) (3.950,-0.780) (4.000,-0.812) (4.050,-0.846) (4.093,-0.875) (4.150,-0.915) (4.197,-0.950) (4.250,-0.990) (4.295,-1.025) (4.350,-1.070) (4.400,-1.112) (4.444,-1.150) (4.500,-1.200) (4.550,-1.247) (4.600,-1.295) (4.650,-1.346) (4.700,-1.398) (4.748,-1.450) (4.792,-1.500) (4.825,-1.539) (4.875,-1.599) (4.916,-1.650) (4.950,-1.695) (4.991,-1.750) (5.025,-1.797) (5.061,-1.850) (5.095,-1.900) (5.125,-1.947) (5.158,-2.000) (5.188,-2.050) (5.217,-2.100) (5.245,-2.150) (5.272,-2.200) (5.298,-2.250) (5.323,-2.300) (5.347,-2.350) (5.370,-2.400) (5.393,-2.450) (5.414,-2.500) (5.435,-2.550) (5.455,-2.600) (5.474,-2.650) (5.492,-2.700) (5.510,-2.750) (5.526,-2.800) (5.542,-2.850) (5.558,-2.900) (5.572,-2.950) (5.586,-3.000) (5.599,-3.050) (5.612,-3.100) (5.623,-3.150) (5.634,-3.200) (5.645,-3.250) (5.655,-3.300) (5.664,-3.350) (5.672,-3.400) (5.680,-3.450) (5.687,-3.500) (5.694,-3.550) (5.700,-3.600) (5.705,-3.650) (5.710,-3.700) (5.714,-3.750) (5.717,-3.800) (5.720,-3.850) (5.722,-3.900) (5.724,-3.950) (5.725,-4.000) (5.726,-4.050) (5.726,-4.100) (5.725,-4.150) (5.724,-4.200) (5.722,-4.250) (5.720,-4.300) (5.717,-4.350) (5.714,-4.400) (5.710,-4.450) (5.705,-4.500) (5.700,-4.550) (5.694,-4.600) (5.688,-4.650) (5.681,-4.700) (5.673,-4.750) (5.665,-4.800) (5.657,-4.850) (5.648,-4.900) (5.638,-4.950) (5.628,-5.000) (5.617,-5.050) (5.605,-5.100) (5.593,-5.150) (5.581,-5.200) (5.567,-5.250) (5.553,-5.300) (5.539,-5.350) (5.524,-5.400) (5.509,-5.450) (5.492,-5.500) (5.476,-5.550) (5.458,-5.600) (5.440,-5.650) (5.422,-5.700) (5.402,-5.750) (5.382,-5.800) (5.362,-5.850) (5.341,-5.900) (5.319,-5.950) (5.296,-6.000) (5.275,-6.046) (5.250,-6.099) (5.225,-6.150) (5.200,-6.200) (5.175,-6.248) (5.150,-6.295) (5.125,-6.341) (5.100,-6.386) (5.075,-6.430) (5.049,-6.475) (5.019,-6.525) (4.988,-6.575) (4.957,-6.625) (4.925,-6.674) (4.892,-6.725) (4.858,-6.775) (4.825,-6.822) (4.787,-6.875) (4.751,-6.925) (4.714,-6.975) (4.675,-7.025) (4.636,-7.075) (4.600,-7.120) (4.555,-7.175) (4.513,-7.225) (4.475,-7.269) (4.425,-7.325) (4.380,-7.375) (4.333,-7.425) (4.286,-7.475) (4.250,-7.512) (4.212,-7.550) (4.175,-7.586) (4.135,-7.625) (4.082,-7.675) (4.028,-7.725) (3.975,-7.773) (3.925,-7.817) (3.886,-7.850) (3.827,-7.900) (3.775,-7.942) (3.734,-7.975) (3.675,-8.021) (3.625,-8.059) (3.575,-8.097) (3.525,-8.133) (3.475,-8.169) (3.430,-8.200) (3.375,-8.238) (3.325,-8.271) (3.281,-8.300) (3.225,-8.335) (3.175,-8.367) (3.125,-8.397) (3.078,-8.425) (3.025,-8.456) (2.975,-8.484) (2.925,-8.512) (2.875,-8.540) (2.825,-8.566) (2.775,-8.592) (2.725,-8.617) (2.675,-8.642) (2.625,-8.666) (2.575,-8.690) (2.525,-8.713) (2.475,-8.736) (2.425,-8.758) (2.375,-8.779) (2.325,-8.800) (2.275,-8.820) (2.225,-8.840) (2.175,-8.859) (2.125,-8.878) (2.075,-8.897) (2.025,-8.914) (1.975,-8.932) (1.925,-8.948) (1.875,-8.965) (1.825,-8.981) (1.775,-8.996) (1.725,-9.011) (1.677,-9.025) (1.625,-9.039) (1.575,-9.053) (1.525,-9.066) (1.475,-9.079) (1.425,-9.091) (1.375,-9.103) (1.325,-9.114) (1.275,-9.125) (1.225,-9.135) (1.175,-9.145) (1.125,-9.155) (1.075,-9.164) (1.025,-9.172) (0.975,-9.181) (0.925,-9.189) (0.875,-9.196) (0.825,-9.203) (0.775,-9.209) (0.725,-9.216) (0.675,-9.221) (0.625,-9.227) (0.575,-9.232) (0.525,-9.236) (0.475,-9.240) (0.425,-9.244) (0.375,-9.247) (0.325,-9.250) (0.275,-9.252) (0.225,-9.254) (0.175,-9.256) (0.125,-9.257) (0.075,-9.258) (0.025,-9.258) (-0.050,-9.258) (-0.100,-9.258) (-0.150,-9.257) (-0.200,-9.255) (-0.250,-9.253) (-0.300,-9.251) (-0.350,-9.249) (-0.400,-9.246) (-0.450,-9.242) (-0.500,-9.238) (-0.550,-9.234) (-0.600,-9.229) (-0.650,-9.224) (-0.700,-9.219) (-0.750,-9.213) (-0.800,-9.206) (-0.850,-9.199) (-0.900,-9.192) (-0.950,-9.185) (-1.000,-9.177) (-1.050,-9.168) (-1.100,-9.159) (-1.149,-9.150) (-1.200,-9.140) (-1.250,-9.130) (-1.300,-9.119) (-1.350,-9.108) (-1.400,-9.097) (-1.450,-9.085) (-1.500,-9.072) (-1.550,-9.060) (-1.600,-9.046) (-1.650,-9.033) (-1.700,-9.018) (-1.750,-9.004) (-1.800,-8.988) (-1.850,-8.973) (-1.900,-8.957) (-1.950,-8.940) (-2.000,-8.923) (-2.050,-8.906) (-2.100,-8.887) (-2.150,-8.869) (-2.199,-8.850) (-2.250,-8.830) (-2.300,-8.810) (-2.350,-8.790) (-2.400,-8.768) (-2.450,-8.747) (-2.499,-8.725) (-2.550,-8.702) (-2.600,-8.678) (-2.650,-8.654) (-2.700,-8.630) (-2.750,-8.605) (-2.800,-8.579) (-2.850,-8.553) (-2.900,-8.526) (-2.947,-8.500) (-2.992,-8.475) (-3.035,-8.450) (-3.100,-8.412) (-3.150,-8.382) (-3.200,-8.351) (-3.250,-8.320) (-3.300,-8.287) (-3.350,-8.254) (-3.394,-8.225) (-3.450,-8.186) (-3.500,-8.151) (-3.550,-8.115) (-3.600,-8.078) (-3.650,-8.040) (-3.700,-8.002) (-3.750,-7.962) (-3.796,-7.925) (-3.850,-7.881) (-3.900,-7.838) (-3.944,-7.800) (-4.000,-7.750) (-4.050,-7.705) (-4.100,-7.658) (-4.150,-7.611) (-4.187,-7.575) (-4.225,-7.537) (-4.261,-7.500) (-4.300,-7.460) (-4.350,-7.407) (-4.400,-7.353) (-4.448,-7.300) (-4.491,-7.250) (-4.525,-7.211) (-4.575,-7.151) (-4.616,-7.100) (-4.650,-7.058) (-4.695,-7.000) (-4.725,-6.960) (-4.769,-6.900) (-4.800,-6.857) (-4.840,-6.800) (-4.875,-6.750) (-4.908,-6.700) (-4.941,-6.650) (-4.973,-6.600) (-5.000,-6.556) (-5.034,-6.500) (-5.063,-6.450) (-5.092,-6.400) (-5.120,-6.350) (-5.147,-6.300) (-5.174,-6.250) (-5.200,-6.200) (-5.225,-6.150) (-5.250,-6.100) (-5.273,-6.050) (-5.296,-6.000) (-5.319,-5.950) (-5.341,-5.900) (-5.362,-5.850) (-5.382,-5.800) (-5.402,-5.750) (-5.422,-5.700) (-5.440,-5.650) (-5.458,-5.600) (-5.475,-5.552) (-5.492,-5.500) (-5.509,-5.450) (-5.524,-5.400) (-5.539,-5.350) (-5.553,-5.300) (-5.567,-5.250) (-5.581,-5.200) (-5.593,-5.150) (-5.605,-5.100) (-5.617,-5.050) (-5.628,-5.000) (-5.638,-4.950) (-5.648,-4.900) (-5.657,-4.850) (-5.665,-4.800) (-5.673,-4.750) (-5.681,-4.700) (-5.688,-4.650) (-5.694,-4.600) (-5.700,-4.550) (-5.705,-4.500) (-5.710,-4.450) (-5.714,-4.400) (-5.717,-4.350) (-5.720,-4.300) (-5.722,-4.250) (-5.724,-4.200) (-5.725,-4.150) (-5.726,-4.100) (-5.726,-4.050) (-5.725,-4.000) (-5.724,-3.950) (-5.722,-3.900) (-5.720,-3.850) (-5.717,-3.800) (-5.714,-3.750) (-5.710,-3.700) (-5.705,-3.650) (-5.700,-3.600) (-5.694,-3.550) (-5.687,-3.500) (-5.680,-3.450) (-5.672,-3.400) (-5.664,-3.350) (-5.655,-3.300) (-5.645,-3.250) (-5.634,-3.200) (-5.623,-3.150) (-5.612,-3.100) (-5.599,-3.050) (-5.586,-3.000) (-5.572,-2.950) (-5.558,-2.900) (-5.542,-2.850) (-5.526,-2.800) (-5.510,-2.750) (-5.492,-2.700) (-5.474,-2.650) (-5.455,-2.600) (-5.435,-2.550) (-5.414,-2.500) (-5.393,-2.450) (-5.370,-2.400) (-5.347,-2.350) (-5.325,-2.304) (-5.300,-2.254) (-5.275,-2.206) (-5.250,-2.159) (-5.225,-2.114) (-5.200,-2.071) (-5.173,-2.025) (-5.143,-1.975) (-5.111,-1.925) (-5.078,-1.875) (-5.050,-1.833) (-5.009,-1.775) (-4.975,-1.728) (-4.935,-1.675) (-4.900,-1.630) (-4.855,-1.575) (-4.814,-1.525) (-4.775,-1.481) (-4.725,-1.425) (-4.678,-1.375) (-4.630,-1.325) (-4.579,-1.275) (-4.527,-1.225) (-4.475,-1.177) (-4.425,-1.133) (-4.386,-1.100) (-4.326,-1.050) (-4.275,-1.009) (-4.231,-0.975) (-4.175,-0.933) (-4.129,-0.900) (-4.075,-0.863) (-4.025,-0.829) (-3.981,-0.800) (-3.925,-0.765) (-3.875,-0.734) (-3.825,-0.705) (-3.775,-0.676) (-3.728,-0.650) (-3.682,-0.625) (-3.634,-0.600) (-3.584,-0.575) (-3.533,-0.550) (-3.480,-0.525) (-3.425,-0.500) (-3.375,-0.478) (-3.325,-0.457) (-3.275,-0.436) (-3.225,-0.416) (-3.175,-0.397) (-3.125,-0.379) (-3.075,-0.361) (-3.025,-0.343) (-2.975,-0.327) (-2.925,-0.311) (-2.875,-0.295) (-2.825,-0.280) (-2.775,-0.266) (-2.725,-0.252) (-2.675,-0.238) (-2.625,-0.226) (-2.575,-0.213) (-2.525,-0.201) (-2.475,-0.190) (-2.425,-0.179) (-2.375,-0.168) (-2.325,-0.158) (-2.275,-0.149) (-2.225,-0.139) (-2.175,-0.131) (-2.125,-0.122) (-2.075,-0.114) (-2.025,-0.107) (-1.975,-0.099) (-1.925,-0.092) (-1.875,-0.086) (-1.825,-0.079) (-1.775,-0.074) (-1.725,-0.068) (-1.675,-0.063) (-1.625,-0.058) (-1.575,-0.053) (-1.525,-0.048) (-1.475,-0.044) (-1.425,-0.040) (-1.375,-0.037) (-1.325,-0.033) (-1.275,-0.030) (-1.225,-0.027) (-1.175,-0.024) (-1.125,-0.022) (-1.075,-0.020) (-1.025,-0.017) (-0.975,-0.015) (-0.925,-0.014) (-0.875,-0.012) (-0.825,-0.011) (-0.775,-0.009) (-0.725,-0.008) (-0.675,-0.007) (-0.625,-0.006) (-0.575,-0.005) (-0.525,-0.004) (-0.475,-0.004) (-0.425,-0.003) (-0.375,-0.003) (-0.325,-0.002) (-0.275,-0.002) (-0.250,-0.002)}
\EQ{(-6.558,-10.000) (-6.597,-9.950) (-6.636,-9.900) (-6.674,-9.850) (-6.712,-9.800) (-6.748,-9.750) (-6.785,-9.700) (-6.820,-9.650) (-6.850,-9.607) (-6.889,-9.550) (-6.923,-9.500) (-6.956,-9.450) (-6.989,-9.400) (-7.021,-9.350) (-7.050,-9.304) (-7.083,-9.250) (-7.114,-9.200) (-7.144,-9.150) (-7.173,-9.100) (-7.200,-9.053) (-7.225,-9.009) (-7.250,-8.964) (-7.275,-8.918) (-7.300,-8.872) (-7.325,-8.825) (-7.350,-8.776) (-7.375,-8.727) (-7.400,-8.677) (-7.425,-8.626) (-7.449,-8.575) (-7.473,-8.525) (-7.496,-8.475) (-7.519,-8.425) (-7.541,-8.375) (-7.562,-8.325) (-7.584,-8.275) (-7.604,-8.225) (-7.625,-8.175) (-7.645,-8.125) (-7.664,-8.075) (-7.683,-8.025) (-7.701,-7.975) (-7.719,-7.925) (-7.737,-7.875) (-7.754,-7.825) (-7.771,-7.775) (-7.787,-7.725) (-7.803,-7.675) (-7.819,-7.625) (-7.834,-7.575) (-7.849,-7.525) (-7.863,-7.475) (-7.877,-7.425) (-7.890,-7.375) (-7.903,-7.325) (-7.916,-7.275) (-7.928,-7.225) (-7.940,-7.175) (-7.951,-7.125) (-7.962,-7.075) (-7.972,-7.025) (-7.983,-6.975) (-7.992,-6.925) (-8.002,-6.875) (-8.010,-6.825) (-8.019,-6.775) (-8.027,-6.725) (-8.035,-6.675) (-8.042,-6.625) (-8.049,-6.575) (-8.055,-6.525) (-8.062,-6.475) (-8.067,-6.425) (-8.073,-6.375) (-8.077,-6.325) (-8.082,-6.275) (-8.086,-6.225) (-8.090,-6.175) (-8.093,-6.125) (-8.096,-6.075) (-8.098,-6.025) (-8.100,-5.975) (-8.102,-5.925) (-8.103,-5.875) (-8.104,-5.825) (-8.105,-5.775) (-8.105,-5.725) (-8.104,-5.675) (-8.103,-5.625) (-8.102,-5.575) (-8.100,-5.525) (-8.098,-5.475) (-8.096,-5.425) (-8.093,-5.375) (-8.090,-5.325) (-8.086,-5.275) (-8.082,-5.225) (-8.077,-5.175) (-8.072,-5.125) (-8.067,-5.075) (-8.061,-5.025) (-8.055,-4.975) (-8.048,-4.925) (-8.041,-4.875) (-8.033,-4.825) (-8.025,-4.775) (-8.016,-4.725) (-8.007,-4.675) (-7.998,-4.625) (-7.988,-4.575) (-7.977,-4.525) (-7.967,-4.475) (-7.955,-4.425) (-7.943,-4.375) (-7.931,-4.325) (-7.918,-4.275) (-7.905,-4.225) (-7.891,-4.175) (-7.877,-4.125) (-7.862,-4.075) (-7.847,-4.025) (-7.831,-3.975) (-7.815,-3.925) (-7.798,-3.875) (-7.780,-3.825) (-7.762,-3.775) (-7.744,-3.725) (-7.725,-3.676) (-7.705,-3.625) (-7.685,-3.575) (-7.664,-3.525) (-7.643,-3.475) (-7.621,-3.425) (-7.600,-3.379) (-7.575,-3.325) (-7.551,-3.275) (-7.526,-3.225) (-7.501,-3.175) (-7.475,-3.125) (-7.450,-3.077) (-7.425,-3.031) (-7.400,-2.986) (-7.375,-2.942) (-7.336,-2.875) (-7.306,-2.825) (-7.275,-2.775) (-7.243,-2.725) (-7.211,-2.675) (-7.178,-2.625) (-7.144,-2.575) (-7.109,-2.525) (-7.075,-2.478) (-7.036,-2.425) (-7.000,-2.377) (-6.960,-2.325) (-6.925,-2.281) (-6.880,-2.225) (-6.838,-2.175) (-6.800,-2.130) (-6.752,-2.075) (-6.707,-2.025) (-6.661,-1.975) (-6.625,-1.937) (-6.590,-1.900) (-6.550,-1.860) (-6.500,-1.810) (-6.464,-1.775) (-6.425,-1.738) (-6.384,-1.700) (-6.329,-1.650) (-6.275,-1.602) (-6.225,-1.559) (-6.185,-1.525) (-6.125,-1.476) (-6.075,-1.435) (-6.030,-1.400) (-5.975,-1.358) (-5.931,-1.325) (-5.875,-1.285) (-5.826,-1.250) (-5.775,-1.214) (-5.725,-1.180) (-5.679,-1.150) (-5.625,-1.115) (-5.575,-1.083) (-5.525,-1.052) (-5.480,-1.025) (-5.425,-0.993) (-5.375,-0.964) (-5.325,-0.936) (-5.275,-0.908) (-5.225,-0.881) (-5.175,-0.855) (-5.125,-0.830) (-5.075,-0.805) (-5.025,-0.780) (-4.975,-0.756) (-4.925,-0.733) (-4.875,-0.710) (-4.825,-0.688) (-4.775,-0.666) (-4.725,-0.645) (-4.678,-0.625) (-4.625,-0.604) (-4.575,-0.584) (-4.525,-0.565) (-4.475,-0.546) (-4.425,-0.527) (-4.375,-0.509) (-4.325,-0.492) (-4.276,-0.475) (-4.225,-0.458) (-4.175,-0.442) (-4.125,-0.426) (-4.075,-0.411) (-4.025,-0.396) (-3.975,-0.381) (-3.925,-0.367) (-3.875,-0.353) (-3.825,-0.340) (-3.775,-0.326) (-3.725,-0.314) (-3.675,-0.301) (-3.625,-0.289) (-3.575,-0.278) (-3.525,-0.266) (-3.475,-0.255) (-3.425,-0.244) (-3.375,-0.234) (-3.325,-0.224) (-3.275,-0.214) (-3.225,-0.204) (-3.175,-0.195) (-3.125,-0.186) (-3.075,-0.178) (-3.025,-0.169) (-2.975,-0.161) (-2.925,-0.153) (-2.875,-0.146) (-2.825,-0.139) (-2.775,-0.132) (-2.725,-0.125) (-2.675,-0.118) (-2.625,-0.112) (-2.575,-0.106) (-2.525,-0.100) (-2.475,-0.094) (-2.425,-0.089) (-2.375,-0.084) (-2.325,-0.079) (-2.275,-0.074) (-2.225,-0.069) (-2.175,-0.065) (-2.125,-0.061) (-2.075,-0.057) (-2.025,-0.053) (-1.975,-0.050) (-1.925,-0.046) (-1.875,-0.043) (-1.825,-0.040) (-1.775,-0.037) (-1.725,-0.034) (-1.675,-0.031) (-1.625,-0.029) (-1.575,-0.026) (-1.525,-0.024) (-1.475,-0.022) (-1.425,-0.020) (-1.375,-0.018) (-1.325,-0.017) (-1.275,-0.015) (-1.225,-0.014) (-1.175,-0.012) (-1.125,-0.011) (-1.075,-0.010) (-1.025,-0.009) (-0.975,-0.008) (-0.925,-0.007) (-0.875,-0.006) (-0.825,-0.005) (-0.775,-0.005) (-0.725,-0.004) (-0.675,-0.003) (-0.625,-0.003) (-0.575,-0.003) (-0.525,-0.002) (-0.475,-0.002) (-0.425,-0.002) (-0.375,-0.001) (-0.325,-0.001) (-0.275,-0.001) (-0.250,-0.001)}
\EQ{(0.250,-0.001) (0.300,-0.001) (0.350,-0.001) (0.400,-0.002) (0.450,-0.002) (0.500,-0.002) (0.550,-0.002) (0.600,-0.003) (0.650,-0.003) (0.700,-0.004) (0.750,-0.004) (0.800,-0.005) (0.850,-0.006) (0.900,-0.006) (0.950,-0.007) (1.000,-0.008) (1.050,-0.009) (1.100,-0.010) (1.150,-0.012) (1.200,-0.013) (1.250,-0.014) (1.300,-0.016) (1.350,-0.017) (1.400,-0.019) (1.450,-0.021) (1.500,-0.023) (1.550,-0.025) (1.600,-0.028) (1.650,-0.030) (1.700,-0.033) (1.750,-0.035) (1.800,-0.038) (1.850,-0.041) (1.900,-0.044) (1.950,-0.048) (2.000,-0.051) (2.050,-0.055) (2.100,-0.059) (2.150,-0.063) (2.200,-0.067) (2.250,-0.072) (2.300,-0.076) (2.350,-0.081) (2.400,-0.086) (2.450,-0.092) (2.500,-0.097) (2.550,-0.103) (2.600,-0.109) (2.650,-0.115) (2.700,-0.121) (2.750,-0.128) (2.800,-0.135) (2.850,-0.142) (2.900,-0.150) (2.950,-0.157) (3.000,-0.165) (3.050,-0.173) (3.100,-0.182) (3.150,-0.191) (3.200,-0.200) (3.250,-0.209) (3.300,-0.219) (3.350,-0.229) (3.400,-0.239) (3.450,-0.250) (3.500,-0.261) (3.550,-0.272) (3.600,-0.283) (3.650,-0.295) (3.700,-0.307) (3.750,-0.320) (3.800,-0.333) (3.850,-0.346) (3.900,-0.360) (3.950,-0.374) (4.000,-0.388) (4.050,-0.403) (4.100,-0.418) (4.150,-0.434) (4.200,-0.450) (4.250,-0.466) (4.300,-0.483) (4.348,-0.500) (4.400,-0.518) (4.450,-0.536) (4.500,-0.555) (4.550,-0.574) (4.600,-0.594) (4.650,-0.614) (4.700,-0.634) (4.750,-0.655) (4.796,-0.675) (4.850,-0.699) (4.900,-0.721) (4.950,-0.744) (5.000,-0.768) (5.050,-0.792) (5.100,-0.817) (5.150,-0.842) (5.200,-0.868) (5.250,-0.895) (5.300,-0.922) (5.350,-0.950) (5.394,-0.975) (5.450,-1.007) (5.500,-1.037) (5.550,-1.068) (5.600,-1.099) (5.650,-1.131) (5.700,-1.164) (5.750,-1.197) (5.800,-1.232) (5.850,-1.267) (5.896,-1.300) (5.950,-1.339) (5.997,-1.375) (6.050,-1.416) (6.093,-1.450) (6.150,-1.496) (6.200,-1.537) (6.244,-1.575) (6.300,-1.624) (6.350,-1.669) (6.400,-1.715) (6.438,-1.750) (6.475,-1.786) (6.515,-1.825) (6.565,-1.875) (6.600,-1.911) (6.637,-1.950) (6.675,-1.990) (6.725,-2.045) (6.774,-2.100) (6.817,-2.150) (6.850,-2.189) (6.900,-2.250) (6.940,-2.300) (6.975,-2.344) (7.018,-2.400) (7.050,-2.444) (7.091,-2.500) (7.125,-2.548) (7.161,-2.600) (7.194,-2.650) (7.225,-2.697) (7.259,-2.750) (7.291,-2.800) (7.321,-2.850) (7.350,-2.899) (7.380,-2.950) (7.408,-3.000) (7.435,-3.050) (7.462,-3.100) (7.488,-3.150) (7.514,-3.200) (7.539,-3.250) (7.563,-3.300) (7.587,-3.350) (7.609,-3.400) (7.632,-3.450) (7.653,-3.500) (7.675,-3.550) (7.695,-3.600) (7.715,-3.650) (7.734,-3.700) (7.753,-3.750) (7.771,-3.800) (7.789,-3.850) (7.806,-3.900) (7.823,-3.950) (7.839,-4.000) (7.855,-4.050) (7.870,-4.100) (7.884,-4.150) (7.898,-4.200) (7.912,-4.250) (7.925,-4.300) (7.937,-4.350) (7.949,-4.400) (7.961,-4.450) (7.972,-4.500) (7.983,-4.550) (7.993,-4.600) (8.003,-4.650) (8.012,-4.700) (8.021,-4.750) (8.029,-4.800) (8.037,-4.850) (8.044,-4.900) (8.051,-4.950) (8.058,-5.000) (8.064,-5.050) (8.070,-5.100) (8.075,-5.150) (8.080,-5.200) (8.084,-5.250) (8.088,-5.300) (8.091,-5.350) (8.095,-5.400) (8.097,-5.450) (8.099,-5.500) (8.101,-5.550) (8.103,-5.600) (8.104,-5.650) (8.104,-5.700) (8.105,-5.750) (8.104,-5.800) (8.104,-5.850) (8.103,-5.900) (8.101,-5.950) (8.099,-6.000) (8.097,-6.050) (8.094,-6.100) (8.091,-6.150) (8.088,-6.200) (8.084,-6.250) (8.080,-6.300) (8.075,-6.350) (8.070,-6.400) (8.064,-6.450) (8.059,-6.500) (8.052,-6.550) (8.046,-6.600) (8.039,-6.650) (8.031,-6.700) (8.023,-6.750) (8.015,-6.800) (8.006,-6.850) (7.997,-6.900) (7.987,-6.950) (7.978,-7.000) (7.967,-7.050) (7.956,-7.100) (7.945,-7.150) (7.934,-7.200) (7.922,-7.250) (7.909,-7.300) (7.897,-7.350) (7.883,-7.400) (7.870,-7.450) (7.856,-7.500) (7.841,-7.550) (7.826,-7.600) (7.811,-7.650) (7.795,-7.700) (7.779,-7.750) (7.763,-7.800) (7.746,-7.850) (7.728,-7.900) (7.710,-7.950) (7.692,-8.000) (7.673,-8.050) (7.654,-8.100) (7.635,-8.150) (7.615,-8.200) (7.594,-8.250) (7.573,-8.300) (7.552,-8.350) (7.530,-8.400) (7.507,-8.450) (7.485,-8.500) (7.461,-8.550) (7.438,-8.600) (7.413,-8.650) (7.389,-8.700) (7.363,-8.750) (7.338,-8.800) (7.311,-8.850) (7.285,-8.900) (7.258,-8.950) (7.230,-9.000) (7.202,-9.050) (7.175,-9.096) (7.144,-9.150) (7.114,-9.200) (7.083,-9.250) (7.052,-9.300) (7.025,-9.344) (6.989,-9.400) (6.956,-9.450) (6.925,-9.497) (6.889,-9.550) (6.855,-9.600) (6.825,-9.643) (6.785,-9.700) (6.750,-9.748) (6.712,-9.800) (6.675,-9.849) (6.636,-9.900) (6.600,-9.947) (6.558,-10.000)}
\EQ{(0.250,0.001) (0.300,0.001) (0.350,0.001) (0.400,0.002) (0.450,0.002) (0.500,0.002) (0.550,0.002) (0.600,0.003) (0.650,0.003) (0.700,0.004) (0.750,0.004) (0.800,0.005) (0.850,0.006) (0.900,0.006) (0.950,0.007) (1.000,0.008) (1.050,0.009) (1.100,0.010) (1.150,0.012) (1.200,0.013) (1.250,0.014) (1.300,0.016) (1.350,0.017) (1.400,0.019) (1.450,0.021) (1.500,0.023) (1.550,0.025) (1.600,0.028) (1.650,0.030) (1.700,0.033) (1.750,0.035) (1.800,0.038) (1.850,0.041) (1.900,0.044) (1.950,0.048) (2.000,0.051) (2.050,0.055) (2.100,0.059) (2.150,0.063) (2.200,0.067) (2.250,0.072) (2.300,0.076) (2.350,0.081) (2.400,0.086) (2.450,0.092) (2.500,0.097) (2.550,0.103) (2.600,0.109) (2.650,0.115) (2.700,0.121) (2.750,0.128) (2.800,0.135) (2.850,0.142) (2.900,0.150) (2.950,0.157) (3.000,0.165) (3.050,0.173) (3.100,0.182) (3.150,0.191) (3.200,0.200) (3.250,0.209) (3.300,0.219) (3.350,0.229) (3.400,0.239) (3.450,0.250) (3.500,0.261) (3.550,0.272) (3.600,0.283) (3.650,0.295) (3.700,0.307) (3.750,0.320) (3.800,0.333) (3.850,0.346) (3.900,0.360) (3.950,0.374) (4.000,0.388) (4.050,0.403) (4.100,0.418) (4.150,0.434) (4.200,0.450) (4.250,0.466) (4.300,0.483) (4.348,0.500) (4.400,0.518) (4.450,0.536) (4.500,0.555) (4.550,0.574) (4.600,0.594) (4.650,0.614) (4.700,0.634) (4.750,0.655) (4.796,0.675) (4.850,0.699) (4.900,0.721) (4.950,0.744) (5.000,0.768) (5.050,0.792) (5.100,0.817) (5.150,0.842) (5.200,0.868) (5.250,0.895) (5.300,0.922) (5.350,0.950) (5.394,0.975) (5.450,1.007) (5.500,1.037) (5.550,1.068) (5.600,1.099) (5.650,1.131) (5.700,1.164) (5.750,1.197) (5.800,1.232) (5.850,1.267) (5.896,1.300) (5.950,1.339) (5.997,1.375) (6.050,1.416) (6.093,1.450) (6.150,1.496) (6.200,1.537) (6.244,1.575) (6.300,1.624) (6.350,1.669) (6.400,1.715) (6.438,1.750) (6.475,1.786) (6.515,1.825) (6.565,1.875) (6.600,1.911) (6.637,1.950) (6.675,1.990) (6.725,2.045) (6.774,2.100) (6.817,2.150) (6.850,2.189) (6.900,2.250) (6.940,2.300) (6.975,2.344) (7.018,2.400) (7.050,2.444) (7.091,2.500) (7.125,2.548) (7.161,2.600) (7.194,2.650) (7.225,2.697) (7.259,2.750) (7.291,2.800) (7.321,2.850) (7.350,2.899) (7.380,2.950) (7.408,3.000) (7.435,3.050) (7.462,3.100) (7.488,3.150) (7.514,3.200) (7.539,3.250) (7.563,3.300) (7.587,3.350) (7.609,3.400) (7.632,3.450) (7.653,3.500) (7.675,3.550) (7.695,3.600) (7.715,3.650) (7.734,3.700) (7.753,3.750) (7.771,3.800) (7.789,3.850) (7.806,3.900) (7.823,3.950) (7.839,4.000) (7.855,4.050) (7.870,4.100) (7.884,4.150) (7.898,4.200) (7.912,4.250) (7.925,4.300) (7.937,4.350) (7.949,4.400) (7.961,4.450) (7.972,4.500) (7.983,4.550) (7.993,4.600) (8.003,4.650) (8.012,4.700) (8.021,4.750) (8.029,4.800) (8.037,4.850) (8.044,4.900) (8.051,4.950) (8.058,5.000) (8.064,5.050) (8.070,5.100) (8.075,5.150) (8.080,5.200) (8.084,5.250) (8.088,5.300) (8.091,5.350) (8.095,5.400) (8.097,5.450) (8.099,5.500) (8.101,5.550) (8.103,5.600) (8.104,5.650) (8.104,5.700) (8.105,5.750) (8.104,5.800) (8.104,5.850) (8.103,5.900) (8.101,5.950) (8.099,6.000) (8.097,6.050) (8.094,6.100) (8.091,6.150) (8.088,6.200) (8.084,6.250) (8.080,6.300) (8.075,6.350) (8.070,6.400) (8.064,6.450) (8.059,6.500) (8.052,6.550) (8.046,6.600) (8.039,6.650) (8.031,6.700) (8.023,6.750) (8.015,6.800) (8.006,6.850) (7.997,6.900) (7.987,6.950) (7.978,7.000) (7.967,7.050) (7.956,7.100) (7.945,7.150) (7.934,7.200) (7.922,7.250) (7.909,7.300) (7.897,7.350) (7.883,7.400) (7.870,7.450) (7.856,7.500) (7.841,7.550) (7.826,7.600) (7.811,7.650) (7.795,7.700) (7.779,7.750) (7.763,7.800) (7.746,7.850) (7.728,7.900) (7.710,7.950) (7.692,8.000) (7.673,8.050) (7.654,8.100) (7.635,8.150) (7.615,8.200) (7.594,8.250) (7.573,8.300) (7.552,8.350) (7.530,8.400) (7.507,8.450) (7.485,8.500) (7.461,8.550) (7.438,8.600) (7.413,8.650) (7.389,8.700) (7.363,8.750) (7.338,8.800) (7.311,8.850) (7.285,8.900) (7.258,8.950) (7.230,9.000) (7.202,9.050) (7.175,9.096) (7.144,9.150) (7.114,9.200) (7.083,9.250) (7.052,9.300) (7.025,9.344) (6.989,9.400) (6.956,9.450) (6.925,9.497) (6.889,9.550) (6.855,9.600) (6.825,9.643) (6.785,9.700) (6.750,9.748) (6.712,9.800) (6.675,9.849) (6.636,9.900) (6.600,9.947) (6.558,10.000)}
\EQ{(-6.558,10.000) (-6.597,9.950) (-6.636,9.900) (-6.674,9.850) (-6.712,9.800) (-6.748,9.750) (-6.785,9.700) (-6.820,9.650) (-6.850,9.607) (-6.889,9.550) (-6.923,9.500) (-6.956,9.450) (-6.989,9.400) (-7.021,9.350) (-7.050,9.304) (-7.083,9.250) (-7.114,9.200) (-7.144,9.150) (-7.173,9.100) (-7.200,9.053) (-7.225,9.009) (-7.250,8.964) (-7.275,8.918) (-7.300,8.872) (-7.325,8.825) (-7.350,8.776) (-7.375,8.727) (-7.400,8.677) (-7.425,8.626) (-7.449,8.575) (-7.473,8.525) (-7.496,8.475) (-7.519,8.425) (-7.541,8.375) (-7.562,8.325) (-7.584,8.275) (-7.604,8.225) (-7.625,8.175) (-7.645,8.125) (-7.664,8.075) (-7.683,8.025) (-7.701,7.975) (-7.719,7.925) (-7.737,7.875) (-7.754,7.825) (-7.771,7.775) (-7.787,7.725) (-7.803,7.675) (-7.819,7.625) (-7.834,7.575) (-7.849,7.525) (-7.863,7.475) (-7.877,7.425) (-7.890,7.375) (-7.903,7.325) (-7.916,7.275) (-7.928,7.225) (-7.940,7.175) (-7.951,7.125) (-7.962,7.075) (-7.972,7.025) (-7.983,6.975) (-7.992,6.925) (-8.002,6.875) (-8.010,6.825) (-8.019,6.775) (-8.027,6.725) (-8.035,6.675) (-8.042,6.625) (-8.049,6.575) (-8.055,6.525) (-8.062,6.475) (-8.067,6.425) (-8.073,6.375) (-8.077,6.325) (-8.082,6.275) (-8.086,6.225) (-8.090,6.175) (-8.093,6.125) (-8.096,6.075) (-8.098,6.025) (-8.100,5.975) (-8.102,5.925) (-8.103,5.875) (-8.104,5.825) (-8.105,5.775) (-8.105,5.725) (-8.104,5.675) (-8.103,5.625) (-8.102,5.575) (-8.100,5.525) (-8.098,5.475) (-8.096,5.425) (-8.093,5.375) (-8.090,5.325) (-8.086,5.275) (-8.082,5.225) (-8.077,5.175) (-8.072,5.125) (-8.067,5.075) (-8.061,5.025) (-8.055,4.975) (-8.048,4.925) (-8.041,4.875) (-8.033,4.825) (-8.025,4.775) (-8.016,4.725) (-8.007,4.675) (-7.998,4.625) (-7.988,4.575) (-7.977,4.525) (-7.967,4.475) (-7.955,4.425) (-7.943,4.375) (-7.931,4.325) (-7.918,4.275) (-7.905,4.225) (-7.891,4.175) (-7.877,4.125) (-7.862,4.075) (-7.847,4.025) (-7.831,3.975) (-7.815,3.925) (-7.798,3.875) (-7.780,3.825) (-7.762,3.775) (-7.744,3.725) (-7.725,3.676) (-7.705,3.625) (-7.685,3.575) (-7.664,3.525) (-7.643,3.475) (-7.621,3.425) (-7.600,3.379) (-7.575,3.325) (-7.551,3.275) (-7.526,3.225) (-7.501,3.175) (-7.475,3.125) (-7.450,3.077) (-7.425,3.031) (-7.400,2.986) (-7.375,2.942) (-7.336,2.875) (-7.306,2.825) (-7.275,2.775) (-7.243,2.725) (-7.211,2.675) (-7.178,2.625) (-7.144,2.575) (-7.109,2.525) (-7.075,2.478) (-7.036,2.425) (-7.000,2.377) (-6.960,2.325) (-6.925,2.281) (-6.880,2.225) (-6.838,2.175) (-6.800,2.130) (-6.752,2.075) (-6.707,2.025) (-6.661,1.975) (-6.625,1.937) (-6.590,1.900) (-6.550,1.860) (-6.500,1.810) (-6.464,1.775) (-6.425,1.738) (-6.384,1.700) (-6.329,1.650) (-6.275,1.602) (-6.225,1.559) (-6.185,1.525) (-6.125,1.476) (-6.075,1.435) (-6.030,1.400) (-5.975,1.358) (-5.931,1.325) (-5.875,1.285) (-5.826,1.250) (-5.775,1.214) (-5.725,1.180) (-5.679,1.150) (-5.625,1.115) (-5.575,1.083) (-5.525,1.052) (-5.480,1.025) (-5.425,0.993) (-5.375,0.964) (-5.325,0.936) (-5.275,0.908) (-5.225,0.881) (-5.175,0.855) (-5.125,0.830) (-5.075,0.805) (-5.025,0.780) (-4.975,0.756) (-4.925,0.733) (-4.875,0.710) (-4.825,0.688) (-4.775,0.666) (-4.725,0.645) (-4.678,0.625) (-4.625,0.604) (-4.575,0.584) (-4.525,0.565) (-4.475,0.546) (-4.425,0.527) (-4.375,0.509) (-4.325,0.492) (-4.276,0.475) (-4.225,0.458) (-4.175,0.442) (-4.125,0.426) (-4.075,0.411) (-4.025,0.396) (-3.975,0.381) (-3.925,0.367) (-3.875,0.353) (-3.825,0.340) (-3.775,0.326) (-3.725,0.314) (-3.675,0.301) (-3.625,0.289) (-3.575,0.278) (-3.525,0.266) (-3.475,0.255) (-3.425,0.244) (-3.375,0.234) (-3.325,0.224) (-3.275,0.214) (-3.225,0.204) (-3.175,0.195) (-3.125,0.186) (-3.075,0.178) (-3.025,0.169) (-2.975,0.161) (-2.925,0.153) (-2.875,0.146) (-2.825,0.139) (-2.775,0.132) (-2.725,0.125) (-2.675,0.118) (-2.625,0.112) (-2.575,0.106) (-2.525,0.100) (-2.475,0.094) (-2.425,0.089) (-2.375,0.084) (-2.325,0.079) (-2.275,0.074) (-2.225,0.069) (-2.175,0.065) (-2.125,0.061) (-2.075,0.057) (-2.025,0.053) (-1.975,0.050) (-1.925,0.046) (-1.875,0.043) (-1.825,0.040) (-1.775,0.037) (-1.725,0.034) (-1.675,0.031) (-1.625,0.029) (-1.575,0.026) (-1.525,0.024) (-1.475,0.022) (-1.425,0.020) (-1.375,0.018) (-1.325,0.017) (-1.275,0.015) (-1.225,0.014) (-1.175,0.012) (-1.125,0.011) (-1.075,0.010) (-1.025,0.009) (-0.975,0.008) (-0.925,0.007) (-0.875,0.006) (-0.825,0.005) (-0.775,0.005) (-0.725,0.004) (-0.675,0.003) (-0.625,0.003) (-0.575,0.003) (-0.525,0.002) (-0.475,0.002) (-0.425,0.002) (-0.375,0.001) (-0.325,0.001) (-0.275,0.001) (-0.250,0.001)}
\EQ{(0.250,0.002) (0.300,0.002) (0.350,0.003) (0.400,0.003) (0.450,0.004) (0.500,0.004) (0.550,0.005) (0.600,0.006) (0.650,0.006) (0.700,0.007) (0.750,0.009) (0.800,0.010) (0.850,0.011) (0.900,0.013) (0.950,0.014) (1.000,0.016) (1.050,0.018) (1.100,0.021) (1.150,0.023) (1.200,0.026) (1.250,0.029) (1.300,0.032) (1.350,0.035) (1.400,0.038) (1.450,0.042) (1.500,0.046) (1.550,0.051) (1.600,0.055) (1.650,0.060) (1.700,0.065) (1.750,0.071) (1.800,0.076) (1.850,0.083) (1.900,0.089) (1.950,0.096) (2.000,0.103) (2.050,0.110) (2.100,0.118) (2.150,0.126) (2.200,0.135) (2.250,0.144) (2.300,0.153) (2.350,0.163) (2.400,0.174) (2.450,0.184) (2.500,0.196) (2.550,0.207) (2.600,0.219) (2.650,0.232) (2.700,0.245) (2.750,0.259) (2.800,0.273) (2.850,0.287) (2.900,0.303) (2.950,0.319) (3.000,0.335) (3.050,0.352) (3.100,0.369) (3.150,0.388) (3.200,0.407) (3.247,0.425) (3.300,0.446) (3.350,0.467) (3.400,0.489) (3.450,0.511) (3.500,0.534) (3.550,0.558) (3.600,0.583) (3.650,0.608) (3.700,0.635) (3.750,0.662) (3.800,0.690) (3.850,0.719) (3.900,0.749) (3.950,0.780) (4.000,0.812) (4.050,0.846) (4.093,0.875) (4.150,0.915) (4.197,0.950) (4.250,0.990) (4.295,1.025) (4.350,1.070) (4.400,1.112) (4.444,1.150) (4.500,1.200) (4.550,1.247) (4.600,1.295) (4.650,1.346) (4.700,1.398) (4.748,1.450) (4.792,1.500) (4.825,1.539) (4.875,1.599) (4.916,1.650) (4.950,1.695) (4.991,1.750) (5.025,1.797) (5.061,1.850) (5.095,1.900) (5.125,1.947) (5.158,2.000) (5.188,2.050) (5.217,2.100) (5.245,2.150) (5.272,2.200) (5.298,2.250) (5.323,2.300) (5.347,2.350) (5.370,2.400) (5.393,2.450) (5.414,2.500) (5.435,2.550) (5.455,2.600) (5.474,2.650) (5.492,2.700) (5.510,2.750) (5.526,2.800) (5.542,2.850) (5.558,2.900) (5.572,2.950) (5.586,3.000) (5.599,3.050) (5.612,3.100) (5.623,3.150) (5.634,3.200) (5.645,3.250) (5.655,3.300) (5.664,3.350) (5.672,3.400) (5.680,3.450) (5.687,3.500) (5.694,3.550) (5.700,3.600) (5.705,3.650) (5.710,3.700) (5.714,3.750) (5.717,3.800) (5.720,3.850) (5.722,3.900) (5.724,3.950) (5.725,4.000) (5.726,4.050) (5.726,4.100) (5.725,4.150) (5.724,4.200) (5.722,4.250) (5.720,4.300) (5.717,4.350) (5.714,4.400) (5.710,4.450) (5.705,4.500) (5.700,4.550) (5.694,4.600) (5.688,4.650) (5.681,4.700) (5.673,4.750) (5.665,4.800) (5.657,4.850) (5.648,4.900) (5.638,4.950) (5.628,5.000) (5.617,5.050) (5.605,5.100) (5.593,5.150) (5.581,5.200) (5.567,5.250) (5.553,5.300) (5.539,5.350) (5.524,5.400) (5.509,5.450) (5.492,5.500) (5.476,5.550) (5.458,5.600) (5.440,5.650) (5.422,5.700) (5.402,5.750) (5.382,5.800) (5.362,5.850) (5.341,5.900) (5.319,5.950) (5.296,6.000) (5.275,6.046) (5.250,6.099) (5.225,6.150) (5.200,6.200) (5.175,6.248) (5.150,6.295) (5.125,6.341) (5.100,6.386) (5.075,6.430) (5.049,6.475) (5.019,6.525) (4.988,6.575) (4.957,6.625) (4.925,6.674) (4.892,6.725) (4.858,6.775) (4.825,6.822) (4.787,6.875) (4.751,6.925) (4.714,6.975) (4.675,7.025) (4.636,7.075) (4.600,7.120) (4.555,7.175) (4.513,7.225) (4.475,7.269) (4.425,7.325) (4.380,7.375) (4.333,7.425) (4.286,7.475) (4.250,7.512) (4.212,7.550) (4.175,7.586) (4.135,7.625) (4.082,7.675) (4.028,7.725) (3.975,7.773) (3.925,7.817) (3.886,7.850) (3.827,7.900) (3.775,7.942) (3.734,7.975) (3.675,8.021) (3.625,8.059) (3.575,8.097) (3.525,8.133) (3.475,8.169) (3.430,8.200) (3.375,8.238) (3.325,8.271) (3.281,8.300) (3.225,8.335) (3.175,8.367) (3.125,8.397) (3.078,8.425) (3.025,8.456) (2.975,8.484) (2.925,8.512) (2.875,8.540) (2.825,8.566) (2.775,8.592) (2.725,8.617) (2.675,8.642) (2.625,8.666) (2.575,8.690) (2.525,8.713) (2.475,8.736) (2.425,8.758) (2.375,8.779) (2.325,8.800) (2.275,8.820) (2.225,8.840) (2.175,8.859) (2.125,8.878) (2.075,8.897) (2.025,8.914) (1.975,8.932) (1.925,8.948) (1.875,8.965) (1.825,8.981) (1.775,8.996) (1.725,9.011) (1.677,9.025) (1.625,9.039) (1.575,9.053) (1.525,9.066) (1.475,9.079) (1.425,9.091) (1.375,9.103) (1.325,9.114) (1.275,9.125) (1.225,9.135) (1.175,9.145) (1.125,9.155) (1.075,9.164) (1.025,9.172) (0.975,9.181) (0.925,9.189) (0.875,9.196) (0.825,9.203) (0.775,9.209) (0.725,9.216) (0.675,9.221) (0.625,9.227) (0.575,9.232) (0.525,9.236) (0.475,9.240) (0.425,9.244) (0.375,9.247) (0.325,9.250) (0.275,9.252) (0.225,9.254) (0.175,9.256) (0.125,9.257) (0.075,9.258) (0.025,9.258) (-0.050,9.258) (-0.100,9.258) (-0.150,9.257) (-0.200,9.255) (-0.250,9.253) (-0.300,9.251) (-0.350,9.249) (-0.400,9.246) (-0.450,9.242) (-0.500,9.238) (-0.550,9.234) (-0.600,9.229) (-0.650,9.224) (-0.700,9.219) (-0.750,9.213) (-0.800,9.206) (-0.850,9.199) (-0.900,9.192) (-0.950,9.185) (-1.000,9.177) (-1.050,9.168) (-1.100,9.159) (-1.149,9.150) (-1.200,9.140) (-1.250,9.130) (-1.300,9.119) (-1.350,9.108) (-1.400,9.097) (-1.450,9.085) (-1.500,9.072) (-1.550,9.060) (-1.600,9.046) (-1.650,9.033) (-1.700,9.018) (-1.750,9.004) (-1.800,8.988) (-1.850,8.973) (-1.900,8.957) (-1.950,8.940) (-2.000,8.923) (-2.050,8.906) (-2.100,8.887) (-2.150,8.869) (-2.199,8.850) (-2.250,8.830) (-2.300,8.810) (-2.350,8.790) (-2.400,8.768) (-2.450,8.747) (-2.499,8.725) (-2.550,8.702) (-2.600,8.678) (-2.650,8.654) (-2.700,8.630) (-2.750,8.605) (-2.800,8.579) (-2.850,8.553) (-2.900,8.526) (-2.947,8.500) (-2.992,8.475) (-3.035,8.450) (-3.100,8.412) (-3.150,8.382) (-3.200,8.351) (-3.250,8.320) (-3.300,8.287) (-3.350,8.254) (-3.394,8.225) (-3.450,8.186) (-3.500,8.151) (-3.550,8.115) (-3.600,8.078) (-3.650,8.040) (-3.700,8.002) (-3.750,7.962) (-3.796,7.925) (-3.850,7.881) (-3.900,7.838) (-3.944,7.800) (-4.000,7.750) (-4.050,7.705) (-4.100,7.658) (-4.150,7.611) (-4.187,7.575) (-4.225,7.537) (-4.261,7.500) (-4.300,7.460) (-4.350,7.407) (-4.400,7.353) (-4.448,7.300) (-4.491,7.250) (-4.525,7.211) (-4.575,7.151) (-4.616,7.100) (-4.650,7.058) (-4.695,7.000) (-4.725,6.960) (-4.769,6.900) (-4.800,6.857) (-4.840,6.800) (-4.875,6.750) (-4.908,6.700) (-4.941,6.650) (-4.973,6.600) (-5.000,6.556) (-5.034,6.500) (-5.063,6.450) (-5.092,6.400) (-5.120,6.350) (-5.147,6.300) (-5.174,6.250) (-5.200,6.200) (-5.225,6.150) (-5.250,6.100) (-5.273,6.050) (-5.296,6.000) (-5.319,5.950) (-5.341,5.900) (-5.362,5.850) (-5.382,5.800) (-5.402,5.750) (-5.422,5.700) (-5.440,5.650) (-5.458,5.600) (-5.475,5.552) (-5.492,5.500) (-5.509,5.450) (-5.524,5.400) (-5.539,5.350) (-5.553,5.300) (-5.567,5.250) (-5.581,5.200) (-5.593,5.150) (-5.605,5.100) (-5.617,5.050) (-5.628,5.000) (-5.638,4.950) (-5.648,4.900) (-5.657,4.850) (-5.665,4.800) (-5.673,4.750) (-5.681,4.700) (-5.688,4.650) (-5.694,4.600) (-5.700,4.550) (-5.705,4.500) (-5.710,4.450) (-5.714,4.400) (-5.717,4.350) (-5.720,4.300) (-5.722,4.250) (-5.724,4.200) (-5.725,4.150) (-5.726,4.100) (-5.726,4.050) (-5.725,4.000) (-5.724,3.950) (-5.722,3.900) (-5.720,3.850) (-5.717,3.800) (-5.714,3.750) (-5.710,3.700) (-5.705,3.650) (-5.700,3.600) (-5.694,3.550) (-5.687,3.500) (-5.680,3.450) (-5.672,3.400) (-5.664,3.350) (-5.655,3.300) (-5.645,3.250) (-5.634,3.200) (-5.623,3.150) (-5.612,3.100) (-5.599,3.050) (-5.586,3.000) (-5.572,2.950) (-5.558,2.900) (-5.542,2.850) (-5.526,2.800) (-5.510,2.750) (-5.492,2.700) (-5.474,2.650) (-5.455,2.600) (-5.435,2.550) (-5.414,2.500) (-5.393,2.450) (-5.370,2.400) (-5.347,2.350) (-5.325,2.304) (-5.300,2.254) (-5.275,2.206) (-5.250,2.159) (-5.225,2.114) (-5.200,2.071) (-5.173,2.025) (-5.143,1.975) (-5.111,1.925) (-5.078,1.875) (-5.050,1.833) (-5.009,1.775) (-4.975,1.728) (-4.935,1.675) (-4.900,1.630) (-4.855,1.575) (-4.814,1.525) (-4.775,1.481) (-4.725,1.425) (-4.678,1.375) (-4.630,1.325) (-4.579,1.275) (-4.527,1.225) (-4.475,1.177) (-4.425,1.133) (-4.386,1.100) (-4.326,1.050) (-4.275,1.009) (-4.231,0.975) (-4.175,0.933) (-4.129,0.900) (-4.075,0.863) (-4.025,0.829) (-3.981,0.800) (-3.925,0.765) (-3.875,0.734) (-3.825,0.705) (-3.775,0.676) (-3.728,0.650) (-3.682,0.625) (-3.634,0.600) (-3.584,0.575) (-3.533,0.550) (-3.480,0.525) (-3.425,0.500) (-3.375,0.478) (-3.325,0.457) (-3.275,0.436) (-3.225,0.416) (-3.175,0.397) (-3.125,0.379) (-3.075,0.361) (-3.025,0.343) (-2.975,0.327) (-2.925,0.311) (-2.875,0.295) (-2.825,0.280) (-2.775,0.266) (-2.725,0.252) (-2.675,0.238) (-2.625,0.226) (-2.575,0.213) (-2.525,0.201) (-2.475,0.190) (-2.425,0.179) (-2.375,0.168) (-2.325,0.158) (-2.275,0.149) (-2.225,0.139) (-2.175,0.131) (-2.125,0.122) (-2.075,0.114) (-2.025,0.107) (-1.975,0.099) (-1.925,0.092) (-1.875,0.086) (-1.825,0.079) (-1.775,0.074) (-1.725,0.068) (-1.675,0.063) (-1.625,0.058) (-1.575,0.053) (-1.525,0.048) (-1.475,0.044) (-1.425,0.040) (-1.375,0.037) (-1.325,0.033) (-1.275,0.030) (-1.225,0.027) (-1.175,0.024) (-1.125,0.022) (-1.075,0.020) (-1.025,0.017) (-0.975,0.015) (-0.925,0.014) (-0.875,0.012) (-0.825,0.011) (-0.775,0.009) (-0.725,0.008) (-0.675,0.007) (-0.625,0.006) (-0.575,0.005) (-0.525,0.004) (-0.475,0.004) (-0.425,0.003) (-0.375,0.003) (-0.325,0.002) (-0.275,0.002) (-0.250,0.002)}
\EQ{(0.250,0.003) (0.300,0.003) (0.350,0.004) (0.400,0.005) (0.450,0.005) (0.500,0.006) (0.550,0.007) (0.600,0.008) (0.650,0.010) (0.700,0.011) (0.750,0.013) (0.800,0.015) (0.850,0.017) (0.900,0.019) (0.950,0.022) (1.000,0.025) (1.050,0.028) (1.100,0.031) (1.150,0.035) (1.200,0.039) (1.250,0.043) (1.300,0.047) (1.350,0.052) (1.400,0.058) (1.450,0.063) (1.500,0.070) (1.550,0.076) (1.600,0.083) (1.650,0.090) (1.700,0.098) (1.750,0.106) (1.800,0.115) (1.850,0.124) (1.900,0.134) (1.950,0.144) (2.000,0.155) (2.050,0.166) (2.100,0.178) (2.150,0.191) (2.200,0.204) (2.250,0.218) (2.300,0.232) (2.350,0.247) (2.400,0.263) (2.450,0.279) (2.500,0.296) (2.550,0.314) (2.600,0.333) (2.650,0.353) (2.700,0.373) (2.750,0.394) (2.800,0.416) (2.850,0.439) (2.900,0.463) (2.950,0.488) (3.000,0.514) (3.050,0.541) (3.100,0.569) (3.150,0.599) (3.194,0.625) (3.250,0.661) (3.300,0.694) (3.346,0.725) (3.400,0.764) (3.448,0.800) (3.500,0.840) (3.543,0.875) (3.600,0.923) (3.650,0.967) (3.700,1.013) (3.738,1.050) (3.775,1.087) (3.812,1.125) (3.850,1.165) (3.900,1.221) (3.946,1.275) (3.987,1.325) (4.025,1.373) (4.064,1.425) (4.100,1.475) (4.134,1.525) (4.167,1.575) (4.199,1.625) (4.225,1.668) (4.258,1.725) (4.286,1.775) (4.312,1.825) (4.338,1.875) (4.362,1.925) (4.385,1.975) (4.407,2.025) (4.428,2.075) (4.448,2.125) (4.467,2.175) (4.485,2.225) (4.502,2.275) (4.519,2.325) (4.534,2.375) (4.548,2.425) (4.562,2.475) (4.575,2.525) (4.587,2.575) (4.598,2.625) (4.608,2.675) (4.617,2.725) (4.626,2.775) (4.634,2.825) (4.641,2.875) (4.647,2.925) (4.653,2.975) (4.658,3.025) (4.662,3.075) (4.665,3.125) (4.668,3.175) (4.670,3.225) (4.671,3.275) (4.671,3.325) (4.671,3.375) (4.670,3.425) (4.668,3.475) (4.666,3.525) (4.663,3.575) (4.659,3.625) (4.655,3.675) (4.650,3.725) (4.644,3.775) (4.637,3.825) (4.630,3.875) (4.622,3.925) (4.613,3.975) (4.604,4.025) (4.594,4.075) (4.583,4.125) (4.572,4.175) (4.560,4.225) (4.547,4.275) (4.533,4.325) (4.519,4.375) (4.504,4.425) (4.488,4.475) (4.472,4.525) (4.454,4.575) (4.436,4.625) (4.418,4.675) (4.398,4.725) (4.378,4.775) (4.357,4.825) (4.335,4.875) (4.312,4.925) (4.288,4.975) (4.264,5.025) (4.239,5.075) (4.213,5.125) (4.186,5.175) (4.158,5.225) (4.130,5.275) (4.100,5.325) (4.069,5.375) (4.038,5.425) (4.005,5.475) (3.975,5.520) (3.937,5.575) (3.902,5.625) (3.865,5.675) (3.827,5.725) (3.789,5.775) (3.750,5.823) (3.707,5.875) (3.665,5.925) (3.625,5.970) (3.576,6.025) (3.529,6.075) (3.481,6.125) (3.432,6.175) (3.381,6.225) (3.328,6.275) (3.275,6.324) (3.225,6.368) (3.175,6.412) (3.129,6.450) (3.075,6.495) (3.025,6.534) (2.975,6.573) (2.925,6.611) (2.875,6.647) (2.825,6.683) (2.775,6.717) (2.726,6.750) (2.675,6.784) (2.625,6.816) (2.575,6.846) (2.528,6.875) (2.475,6.906) (2.425,6.935) (2.375,6.962) (2.325,6.989) (2.275,7.015) (2.225,7.041) (2.175,7.066) (2.125,7.090) (2.075,7.113) (2.025,7.136) (1.975,7.158) (1.925,7.179) (1.875,7.200) (1.825,7.220) (1.775,7.239) (1.725,7.258) (1.678,7.275) (1.625,7.294) (1.575,7.311) (1.525,7.327) (1.475,7.343) (1.425,7.358) (1.375,7.372) (1.325,7.386) (1.275,7.400) (1.225,7.413) (1.176,7.425) (1.125,7.437) (1.075,7.448) (1.025,7.459) (0.975,7.469) (0.925,7.479) (0.875,7.488) (0.825,7.497) (0.775,7.505) (0.725,7.512) (0.675,7.519) (0.625,7.526) (0.575,7.532) (0.525,7.538) (0.475,7.543) (0.425,7.547) (0.375,7.551) (0.325,7.555) (0.275,7.558) (0.225,7.560) (0.175,7.562) (0.125,7.564) (0.075,7.565) (0.025,7.565) (-0.050,7.565) (-0.100,7.564) (-0.150,7.563) (-0.200,7.561) (-0.250,7.559) (-0.300,7.556) (-0.350,7.553) (-0.400,7.549) (-0.450,7.545) (-0.500,7.540) (-0.550,7.535) (-0.600,7.529) (-0.650,7.523) (-0.700,7.516) (-0.750,7.509) (-0.800,7.501) (-0.850,7.492) (-0.900,7.484) (-0.950,7.474) (-1.000,7.464) (-1.050,7.454) (-1.100,7.443) (-1.150,7.431) (-1.200,7.419) (-1.250,7.406) (-1.300,7.393) (-1.350,7.380) (-1.400,7.365) (-1.450,7.350) (-1.500,7.335) (-1.550,7.319) (-1.600,7.302) (-1.650,7.285) (-1.700,7.267) (-1.750,7.249) (-1.800,7.229) (-1.850,7.210) (-1.900,7.189) (-1.950,7.168) (-2.000,7.147) (-2.049,7.125) (-2.100,7.101) (-2.150,7.078) (-2.200,7.053) (-2.250,7.028) (-2.300,7.002) (-2.350,6.976) (-2.397,6.950) (-2.442,6.925) (-2.485,6.900) (-2.550,6.862) (-2.600,6.831) (-2.649,6.800) (-2.700,6.767) (-2.750,6.734) (-2.800,6.700) (-2.850,6.665) (-2.900,6.629) (-2.950,6.592) (-3.000,6.554) (-3.050,6.515) (-3.099,6.475) (-3.150,6.433) (-3.188,6.400) (-3.246,6.350) (-3.300,6.301) (-3.350,6.254) (-3.400,6.206) (-3.450,6.157) (-3.500,6.106) (-3.550,6.053) (-3.599,6.000) (-3.643,5.950) (-3.686,5.900) (-3.725,5.854) (-3.769,5.800) (-3.800,5.761) (-3.847,5.700) (-3.884,5.650) (-3.920,5.600) (-3.950,5.557) (-3.989,5.500) (-4.022,5.450) (-4.050,5.406) (-4.085,5.350) (-4.115,5.300) (-4.144,5.250) (-4.172,5.200) (-4.200,5.150) (-4.225,5.102) (-4.250,5.053) (-4.275,5.003) (-4.300,4.951) (-4.323,4.900) (-4.346,4.850) (-4.367,4.800) (-4.388,4.750) (-4.408,4.700) (-4.427,4.650) (-4.445,4.600) (-4.463,4.550) (-4.480,4.500) (-4.496,4.450) (-4.511,4.400) (-4.526,4.350) (-4.540,4.300) (-4.553,4.250) (-4.566,4.200) (-4.578,4.150) (-4.589,4.100) (-4.599,4.050) (-4.609,4.000) (-4.618,3.950) (-4.626,3.900) (-4.634,3.850) (-4.641,3.800) (-4.647,3.750) (-4.652,3.700) (-4.657,3.650) (-4.661,3.600) (-4.665,3.550) (-4.667,3.500) (-4.669,3.450) (-4.671,3.400) (-4.671,3.350) (-4.671,3.300) (-4.670,3.250) (-4.669,3.200) (-4.667,3.150) (-4.664,3.100) (-4.660,3.050) (-4.655,3.000) (-4.650,2.950) (-4.644,2.900) (-4.638,2.850) (-4.630,2.800) (-4.622,2.750) (-4.613,2.700) (-4.603,2.650) (-4.592,2.600) (-4.581,2.550) (-4.569,2.500) (-4.555,2.450) (-4.541,2.400) (-4.526,2.350) (-4.511,2.300) (-4.494,2.250) (-4.476,2.200) (-4.458,2.150) (-4.438,2.100) (-4.418,2.050) (-4.396,2.000) (-4.375,1.953) (-4.350,1.900) (-4.325,1.850) (-4.300,1.802) (-4.275,1.755) (-4.250,1.711) (-4.214,1.650) (-4.183,1.600) (-4.151,1.550) (-4.117,1.500) (-4.082,1.450) (-4.050,1.406) (-4.007,1.350) (-3.975,1.310) (-3.925,1.250) (-3.882,1.200) (-3.836,1.150) (-3.800,1.112) (-3.764,1.075) (-3.725,1.037) (-3.686,1.000) (-3.631,0.950) (-3.575,0.902) (-3.525,0.860) (-3.481,0.825) (-3.425,0.782) (-3.381,0.750) (-3.325,0.711) (-3.275,0.677) (-3.225,0.645) (-3.175,0.614) (-3.125,0.584) (-3.075,0.555) (-3.025,0.527) (-2.975,0.501) (-2.925,0.476) (-2.875,0.451) (-2.825,0.428) (-2.775,0.405) (-2.725,0.383) (-2.675,0.363) (-2.625,0.343) (-2.575,0.324) (-2.525,0.305) (-2.475,0.288) (-2.425,0.271) (-2.375,0.255) (-2.325,0.239) (-2.276,0.225) (-2.225,0.211) (-2.175,0.197) (-2.125,0.184) (-2.075,0.172) (-2.025,0.161) (-1.975,0.149) (-1.925,0.139) (-1.875,0.129) (-1.825,0.120) (-1.775,0.111) (-1.725,0.102) (-1.675,0.094) (-1.625,0.087) (-1.575,0.079) (-1.525,0.073) (-1.475,0.066) (-1.425,0.061) (-1.375,0.055) (-1.325,0.050) (-1.275,0.045) (-1.225,0.041) (-1.175,0.037) (-1.125,0.033) (-1.075,0.029) (-1.025,0.026) (-0.975,0.023) (-0.925,0.020) (-0.875,0.018) (-0.825,0.016) (-0.775,0.014) (-0.725,0.012) (-0.675,0.010) (-0.625,0.009) (-0.575,0.008) (-0.525,0.007) (-0.475,0.006) (-0.425,0.005) (-0.375,0.004) (-0.325,0.004) (-0.275,0.003) (-0.250,0.003)}
\EQ{(0.250,0.004) (0.300,0.005) (0.350,0.005) (0.400,0.006) (0.450,0.007) (0.500,0.008) (0.550,0.010) (0.600,0.011) (0.650,0.013) (0.700,0.015) (0.750,0.017) (0.800,0.020) (0.850,0.022) (0.900,0.026) (0.950,0.029) (1.000,0.033) (1.050,0.037) (1.100,0.041) (1.150,0.046) (1.200,0.052) (1.250,0.057) (1.300,0.063) (1.350,0.070) (1.400,0.077) (1.450,0.085) (1.500,0.093) (1.550,0.102) (1.600,0.111) (1.650,0.121) (1.700,0.131) (1.750,0.142) (1.800,0.154) (1.850,0.166) (1.900,0.179) (1.950,0.193) (2.000,0.208) (2.050,0.223) (2.100,0.239) (2.150,0.256) (2.200,0.274) (2.250,0.293) (2.300,0.313) (2.350,0.333) (2.400,0.355) (2.450,0.378) (2.497,0.400) (2.548,0.425) (2.596,0.450) (2.642,0.475) (2.686,0.500) (2.750,0.538) (2.800,0.570) (2.847,0.600) (2.900,0.637) (2.950,0.673) (3.000,0.711) (3.050,0.750) (3.100,0.792) (3.138,0.825) (3.193,0.875) (3.245,0.925) (3.294,0.975) (3.341,1.025) (3.375,1.063) (3.425,1.122) (3.467,1.175) (3.500,1.218) (3.542,1.275) (3.575,1.323) (3.609,1.375) (3.640,1.425) (3.670,1.475) (3.698,1.525) (3.725,1.574) (3.750,1.623) (3.775,1.675) (3.798,1.725) (3.820,1.775) (3.841,1.825) (3.860,1.875) (3.878,1.925) (3.896,1.975) (3.912,2.025) (3.927,2.075) (3.941,2.125) (3.954,2.175) (3.966,2.225) (3.977,2.275) (3.988,2.325) (3.997,2.375) (4.005,2.425) (4.013,2.475) (4.020,2.525) (4.025,2.575) (4.030,2.625) (4.034,2.675) (4.037,2.725) (4.040,2.775) (4.041,2.825) (4.042,2.875) (4.042,2.925) (4.041,2.975) (4.039,3.025) (4.037,3.075) (4.033,3.125) (4.029,3.175) (4.024,3.225) (4.018,3.275) (4.012,3.325) (4.004,3.375) (3.996,3.425) (3.987,3.475) (3.977,3.525) (3.967,3.575) (3.955,3.625) (3.943,3.675) (3.930,3.725) (3.916,3.775) (3.901,3.825) (3.886,3.875) (3.869,3.925) (3.852,3.975) (3.834,4.025) (3.815,4.075) (3.795,4.125) (3.775,4.172) (3.752,4.225) (3.729,4.275) (3.706,4.325) (3.681,4.375) (3.655,4.425) (3.629,4.475) (3.601,4.525) (3.575,4.571) (3.543,4.625) (3.512,4.675) (3.480,4.725) (3.450,4.771) (3.413,4.825) (3.378,4.875) (3.341,4.925) (3.303,4.975) (3.264,5.025) (3.225,5.073) (3.181,5.125) (3.138,5.175) (3.100,5.217) (3.050,5.271) (3.000,5.322) (2.950,5.372) (2.900,5.420) (2.850,5.467) (2.800,5.512) (2.756,5.550) (2.700,5.598) (2.650,5.638) (2.603,5.675) (2.550,5.716) (2.504,5.750) (2.450,5.789) (2.400,5.823) (2.350,5.857) (2.300,5.889) (2.250,5.921) (2.202,5.950) (2.150,5.981) (2.100,6.010) (2.050,6.037) (2.000,6.064) (1.950,6.090) (1.900,6.115) (1.850,6.139) (1.800,6.163) (1.750,6.186) (1.700,6.208) (1.650,6.229) (1.600,6.249) (1.550,6.269) (1.500,6.288) (1.450,6.306) (1.400,6.323) (1.350,6.340) (1.300,6.356) (1.250,6.372) (1.200,6.387) (1.150,6.401) (1.100,6.414) (1.050,6.427) (1.000,6.439) (0.950,6.451) (0.900,6.462) (0.850,6.472) (0.800,6.482) (0.750,6.491) (0.700,6.499) (0.650,6.507) (0.600,6.515) (0.550,6.521) (0.500,6.527) (0.450,6.533) (0.400,6.538) (0.350,6.542) (0.300,6.546) (0.250,6.549) (0.200,6.552) (0.150,6.554) (0.100,6.555) (0.050,6.556) (0.000,6.556) (-0.050,6.556) (-0.100,6.555) (-0.150,6.554) (-0.200,6.552) (-0.250,6.549) (-0.300,6.546) (-0.350,6.542) (-0.400,6.538) (-0.450,6.533) (-0.500,6.527) (-0.550,6.521) (-0.600,6.515) (-0.650,6.507) (-0.700,6.499) (-0.750,6.491) (-0.800,6.482) (-0.850,6.472) (-0.900,6.462) (-0.950,6.451) (-1.000,6.439) (-1.050,6.427) (-1.100,6.414) (-1.150,6.401) (-1.200,6.387) (-1.250,6.372) (-1.300,6.356) (-1.350,6.340) (-1.400,6.323) (-1.450,6.306) (-1.500,6.288) (-1.550,6.269) (-1.598,6.250) (-1.650,6.229) (-1.700,6.208) (-1.750,6.186) (-1.800,6.163) (-1.850,6.139) (-1.900,6.115) (-1.950,6.090) (-2.000,6.064) (-2.050,6.037) (-2.100,6.010) (-2.150,5.981) (-2.200,5.951) (-2.243,5.925) (-2.300,5.889) (-2.350,5.857) (-2.397,5.825) (-2.450,5.789) (-2.500,5.753) (-2.550,5.716) (-2.600,5.678) (-2.650,5.638) (-2.697,5.600) (-2.750,5.556) (-2.800,5.512) (-2.841,5.475) (-2.895,5.425) (-2.947,5.375) (-2.997,5.325) (-3.046,5.275) (-3.092,5.225) (-3.138,5.175) (-3.175,5.132) (-3.223,5.075) (-3.264,5.025) (-3.300,4.979) (-3.341,4.925) (-3.375,4.879) (-3.413,4.825) (-3.447,4.775) (-3.475,4.733) (-3.512,4.675) (-3.543,4.625) (-3.573,4.575) (-3.600,4.527) (-3.625,4.482) (-3.650,4.435) (-3.675,4.387) (-3.700,4.337) (-3.725,4.284) (-3.750,4.229) (-3.774,4.175) (-3.795,4.125) (-3.815,4.075) (-3.834,4.025) (-3.852,3.975) (-3.869,3.925) (-3.886,3.875) (-3.901,3.825) (-3.916,3.775) (-3.930,3.725) (-3.943,3.675) (-3.955,3.625) (-3.967,3.575) (-3.977,3.525) (-3.987,3.475) (-3.996,3.425) (-4.004,3.375) (-4.012,3.325) (-4.018,3.275) (-4.024,3.225) (-4.029,3.175) (-4.033,3.125) (-4.037,3.075) (-4.039,3.025) (-4.041,2.975) (-4.042,2.925) (-4.042,2.875) (-4.041,2.825) (-4.040,2.775) (-4.037,2.725) (-4.034,2.675) (-4.030,2.625) (-4.025,2.575) (-4.020,2.525) (-4.013,2.475) (-4.005,2.425) (-3.997,2.375) (-3.988,2.325) (-3.977,2.275) (-3.966,2.225) (-3.954,2.175) (-3.941,2.125) (-3.927,2.075) (-3.912,2.025) (-3.896,1.975) (-3.878,1.925) (-3.860,1.875) (-3.841,1.825) (-3.820,1.775) (-3.798,1.725) (-3.775,1.675) (-3.751,1.625) (-3.725,1.575) (-3.700,1.528) (-3.675,1.483) (-3.640,1.425) (-3.609,1.375) (-3.576,1.325) (-3.542,1.275) (-3.505,1.225) (-3.475,1.185) (-3.427,1.125) (-3.385,1.075) (-3.350,1.035) (-3.300,0.981) (-3.250,0.930) (-3.200,0.882) (-3.150,0.836) (-3.110,0.800) (-3.050,0.750) (-3.000,0.711) (-2.953,0.675) (-2.900,0.637) (-2.850,0.602) (-2.800,0.570) (-2.750,0.538) (-2.700,0.508) (-2.650,0.480) (-2.600,0.452) (-2.550,0.426) (-2.500,0.401) (-2.450,0.378) (-2.400,0.355) (-2.350,0.333) (-2.300,0.313) (-2.250,0.293) (-2.202,0.275) (-2.150,0.256) (-2.100,0.239) (-2.050,0.223) (-2.000,0.208) (-1.950,0.193) (-1.900,0.179) (-1.850,0.166) (-1.800,0.154) (-1.750,0.142) (-1.700,0.131) (-1.650,0.121) (-1.600,0.111) (-1.550,0.102) (-1.500,0.093) (-1.450,0.085) (-1.400,0.077) (-1.350,0.070) (-1.300,0.063) (-1.250,0.057) (-1.200,0.052) (-1.150,0.046) (-1.100,0.041) (-1.050,0.037) (-1.000,0.033) (-0.950,0.029) (-0.900,0.026) (-0.850,0.022) (-0.800,0.020) (-0.750,0.017) (-0.700,0.015) (-0.650,0.013) (-0.600,0.011) (-0.550,0.010) (-0.500,0.008) (-0.450,0.007) (-0.400,0.006) (-0.350,0.005) (-0.300,0.005) (-0.250,0.004)}
\EQ{(0.250,0.005) (0.300,0.006) (0.350,0.007) (0.400,0.008) (0.450,0.009) (0.500,0.010) (0.550,0.012) (0.600,0.014) (0.650,0.016) (0.700,0.019) (0.750,0.021) (0.800,0.025) (0.850,0.028) (0.900,0.032) (0.950,0.036) (1.000,0.041) (1.050,0.046) (1.100,0.052) (1.150,0.058) (1.200,0.064) (1.250,0.072) (1.300,0.079) (1.350,0.088) (1.400,0.097) (1.450,0.106) (1.500,0.116) (1.550,0.127) (1.600,0.139) (1.650,0.151) (1.700,0.165) (1.750,0.179) (1.800,0.193) (1.850,0.209) (1.898,0.225) (1.950,0.243) (2.000,0.262) (2.050,0.282) (2.100,0.302) (2.150,0.324) (2.200,0.347) (2.250,0.371) (2.300,0.397) (2.350,0.423) (2.397,0.450) (2.450,0.481) (2.500,0.513) (2.550,0.546) (2.593,0.575) (2.650,0.617) (2.693,0.650) (2.750,0.696) (2.800,0.739) (2.840,0.775) (2.892,0.825) (2.942,0.875) (2.988,0.925) (3.025,0.967) (3.074,1.025) (3.113,1.075) (3.150,1.125) (3.186,1.175) (3.219,1.225) (3.250,1.274) (3.280,1.325) (3.308,1.375) (3.335,1.425) (3.360,1.475) (3.384,1.525) (3.406,1.575) (3.427,1.625) (3.447,1.675) (3.465,1.725) (3.482,1.775) (3.498,1.825) (3.513,1.875) (3.527,1.925) (3.540,1.975) (3.551,2.025) (3.562,2.075) (3.571,2.125) (3.580,2.175) (3.587,2.225) (3.594,2.275) (3.599,2.325) (3.604,2.375) (3.607,2.425) (3.610,2.475) (3.611,2.525) (3.612,2.575) (3.612,2.625) (3.611,2.675) (3.609,2.725) (3.606,2.775) (3.602,2.825) (3.598,2.875) (3.592,2.925) (3.586,2.975) (3.578,3.025) (3.570,3.075) (3.561,3.125) (3.551,3.175) (3.540,3.225) (3.528,3.275) (3.515,3.325) (3.501,3.375) (3.487,3.425) (3.471,3.475) (3.455,3.525) (3.437,3.575) (3.419,3.625) (3.400,3.673) (3.379,3.725) (3.357,3.775) (3.334,3.825) (3.311,3.875) (3.286,3.925) (3.260,3.975) (3.234,4.025) (3.205,4.075) (3.176,4.125) (3.150,4.168) (3.114,4.225) (3.081,4.275) (3.050,4.321) (3.011,4.375) (2.975,4.424) (2.936,4.475) (2.900,4.520) (2.855,4.575) (2.811,4.625) (2.775,4.666) (2.725,4.720) (2.675,4.771) (2.625,4.821) (2.575,4.869) (2.525,4.915) (2.485,4.950) (2.426,5.000) (2.375,5.042) (2.334,5.075) (2.275,5.120) (2.234,5.150) (2.175,5.192) (2.127,5.225) (2.075,5.260) (2.025,5.291) (1.975,5.322) (1.928,5.350) (1.885,5.375) (1.840,5.400) (1.794,5.425) (1.746,5.450) (1.700,5.473) (1.650,5.497) (1.600,5.521) (1.550,5.543) (1.500,5.565) (1.450,5.586) (1.400,5.606) (1.350,5.625) (1.300,5.643) (1.250,5.661) (1.200,5.677) (1.150,5.693) (1.100,5.709) (1.050,5.723) (1.000,5.737) (0.950,5.750) (0.900,5.762) (0.850,5.774) (0.800,5.785) (0.750,5.795) (0.700,5.805) (0.650,5.813) (0.600,5.822) (0.550,5.829) (0.500,5.836) (0.450,5.842) (0.400,5.848) (0.350,5.853) (0.300,5.857) (0.250,5.860) (0.200,5.863) (0.150,5.866) (0.100,5.867) (0.050,5.868) (0.000,5.868) (-0.050,5.868) (-0.100,5.867) (-0.150,5.866) (-0.200,5.863) (-0.250,5.860) (-0.300,5.857) (-0.350,5.853) (-0.400,5.848) (-0.450,5.842) (-0.500,5.836) (-0.550,5.829) (-0.600,5.822) (-0.650,5.813) (-0.700,5.805) (-0.750,5.795) (-0.800,5.785) (-0.850,5.774) (-0.900,5.762) (-0.950,5.750) (-1.000,5.737) (-1.050,5.723) (-1.100,5.709) (-1.150,5.693) (-1.200,5.677) (-1.250,5.661) (-1.300,5.643) (-1.349,5.625) (-1.400,5.606) (-1.450,5.586) (-1.500,5.565) (-1.550,5.543) (-1.600,5.521) (-1.650,5.497) (-1.696,5.475) (-1.746,5.450) (-1.794,5.425) (-1.840,5.400) (-1.885,5.375) (-1.928,5.350) (-1.975,5.322) (-2.025,5.291) (-2.075,5.260) (-2.125,5.226) (-2.175,5.192) (-2.225,5.157) (-2.268,5.125) (-2.325,5.082) (-2.365,5.050) (-2.425,5.001) (-2.475,4.959) (-2.513,4.925) (-2.568,4.875) (-2.621,4.825) (-2.671,4.775) (-2.720,4.725) (-2.767,4.675) (-2.811,4.625) (-2.850,4.580) (-2.896,4.525) (-2.936,4.475) (-2.974,4.425) (-3.011,4.375) (-3.047,4.325) (-3.081,4.275) (-3.114,4.225) (-3.146,4.175) (-3.175,4.127) (-3.205,4.075) (-3.234,4.025) (-3.260,3.975) (-3.286,3.925) (-3.311,3.875) (-3.334,3.825) (-3.357,3.775) (-3.379,3.725) (-3.399,3.675) (-3.419,3.625) (-3.437,3.575) (-3.455,3.525) (-3.471,3.475) (-3.487,3.425) (-3.501,3.375) (-3.515,3.325) (-3.528,3.275) (-3.540,3.225) (-3.551,3.175) (-3.561,3.125) (-3.570,3.075) (-3.578,3.025) (-3.586,2.975) (-3.592,2.925) (-3.598,2.875) (-3.602,2.825) (-3.606,2.775) (-3.609,2.725) (-3.611,2.675) (-3.612,2.625) (-3.612,2.575) (-3.611,2.525) (-3.610,2.475) (-3.607,2.425) (-3.604,2.375) (-3.599,2.325) (-3.594,2.275) (-3.587,2.225) (-3.580,2.175) (-3.571,2.125) (-3.562,2.075) (-3.551,2.025) (-3.540,1.975) (-3.527,1.925) (-3.513,1.875) (-3.498,1.825) (-3.482,1.775) (-3.465,1.725) (-3.447,1.675) (-3.427,1.625) (-3.406,1.575) (-3.384,1.525) (-3.360,1.475) (-3.335,1.425) (-3.308,1.375) (-3.280,1.325) (-3.250,1.275) (-3.219,1.225) (-3.186,1.175) (-3.150,1.125) (-3.113,1.075) (-3.075,1.026) (-3.032,0.975) (-2.988,0.925) (-2.950,0.884) (-2.900,0.833) (-2.850,0.784) (-2.813,0.750) (-2.755,0.700) (-2.700,0.655) (-2.650,0.617) (-2.600,0.580) (-2.557,0.550) (-2.500,0.513) (-2.450,0.481) (-2.400,0.452) (-2.353,0.425) (-2.307,0.400) (-2.258,0.375) (-2.207,0.350) (-2.152,0.325) (-2.100,0.302) (-2.050,0.282) (-2.000,0.262) (-1.950,0.243) (-1.900,0.226) (-1.850,0.209) (-1.800,0.193) (-1.750,0.179) (-1.700,0.165) (-1.650,0.151) (-1.600,0.139) (-1.550,0.127) (-1.500,0.116) (-1.450,0.106) (-1.400,0.097) (-1.350,0.088) (-1.300,0.079) (-1.250,0.072) (-1.200,0.064) (-1.150,0.058) (-1.100,0.052) (-1.050,0.046) (-1.000,0.041) (-0.950,0.036) (-0.900,0.032) (-0.850,0.028) (-0.800,0.025) (-0.750,0.021) (-0.700,0.019) (-0.650,0.016) (-0.600,0.014) (-0.550,0.012) (-0.500,0.010) (-0.450,0.009) (-0.400,0.008) (-0.350,0.007) (-0.300,0.006) (-0.250,0.005)}
\EQ{(0.250,0.006) (0.300,0.007) (0.350,0.008) (0.400,0.009) (0.450,0.011) (0.500,0.012) (0.550,0.014) (0.600,0.017) (0.650,0.019) (0.700,0.022) (0.750,0.026) (0.800,0.030) (0.850,0.034) (0.900,0.038) (0.950,0.044) (1.000,0.049) (1.050,0.055) (1.100,0.062) (1.150,0.069) (1.200,0.077) (1.250,0.086) (1.300,0.095) (1.350,0.105) (1.400,0.116) (1.450,0.128) (1.500,0.140) (1.550,0.153) (1.600,0.167) (1.650,0.183) (1.700,0.199) (1.750,0.216) (1.800,0.234) (1.850,0.253) (1.900,0.273) (1.950,0.295) (2.000,0.318) (2.050,0.342) (2.100,0.367) (2.150,0.394) (2.200,0.423) (2.245,0.450) (2.300,0.485) (2.350,0.519) (2.394,0.550) (2.450,0.593) (2.491,0.625) (2.549,0.675) (2.600,0.721) (2.650,0.770) (2.700,0.821) (2.749,0.875) (2.791,0.925) (2.825,0.967) (2.869,1.025) (2.900,1.069) (2.937,1.125) (2.968,1.175) (2.998,1.225) (3.025,1.274) (3.050,1.322) (3.075,1.373) (3.099,1.425) (3.120,1.475) (3.140,1.525) (3.159,1.575) (3.176,1.625) (3.192,1.675) (3.207,1.725) (3.220,1.775) (3.233,1.825) (3.244,1.875) (3.254,1.925) (3.263,1.975) (3.270,2.025) (3.277,2.075) (3.282,2.125) (3.287,2.175) (3.290,2.225) (3.293,2.275) (3.294,2.325) (3.295,2.375) (3.294,2.425) (3.293,2.475) (3.290,2.525) (3.287,2.575) (3.282,2.625) (3.276,2.675) (3.270,2.725) (3.262,2.775) (3.254,2.825) (3.245,2.875) (3.234,2.925) (3.223,2.975) (3.210,3.025) (3.197,3.075) (3.182,3.125) (3.167,3.175) (3.150,3.225) (3.133,3.275) (3.114,3.325) (3.094,3.375) (3.073,3.425) (3.051,3.475) (3.028,3.525) (3.004,3.575) (2.978,3.625) (2.951,3.675) (2.925,3.722) (2.894,3.775) (2.864,3.825) (2.832,3.875) (2.800,3.923) (2.764,3.975) (2.727,4.025) (2.689,4.075) (2.650,4.125) (2.609,4.175) (2.575,4.214) (2.525,4.270) (2.475,4.324) (2.425,4.375) (2.375,4.424) (2.325,4.471) (2.275,4.516) (2.235,4.550) (2.175,4.600) (2.125,4.640) (2.079,4.675) (2.025,4.715) (1.975,4.750) (1.925,4.784) (1.875,4.817) (1.825,4.848) (1.781,4.875) (1.725,4.907) (1.675,4.935) (1.625,4.962) (1.575,4.988) (1.525,5.013) (1.475,5.037) (1.425,5.059) (1.375,5.081) (1.325,5.102) (1.275,5.122) (1.225,5.141) (1.175,5.159) (1.125,5.177) (1.075,5.193) (1.025,5.209) (0.975,5.223) (0.925,5.238) (0.875,5.251) (0.825,5.263) (0.775,5.275) (0.725,5.286) (0.675,5.296) (0.625,5.305) (0.575,5.314) (0.525,5.322) (0.475,5.329) (0.425,5.335) (0.375,5.341) (0.325,5.346) (0.275,5.350) (0.225,5.354) (0.175,5.357) (0.125,5.359) (0.075,5.360) (0.025,5.361) (-0.050,5.361) (-0.100,5.360) (-0.150,5.358) (-0.200,5.355) (-0.250,5.352) (-0.300,5.348) (-0.350,5.344) (-0.400,5.338) (-0.450,5.332) (-0.500,5.325) (-0.550,5.318) (-0.600,5.310) (-0.650,5.301) (-0.700,5.291) (-0.750,5.280) (-0.800,5.269) (-0.850,5.257) (-0.900,5.244) (-0.950,5.231) (-1.000,5.216) (-1.050,5.201) (-1.100,5.185) (-1.150,5.168) (-1.200,5.150) (-1.250,5.132) (-1.300,5.112) (-1.350,5.092) (-1.400,5.070) (-1.450,5.048) (-1.499,5.025) (-1.550,5.001) (-1.600,4.975) (-1.648,4.950) (-1.694,4.925) (-1.738,4.900) (-1.800,4.863) (-1.850,4.833) (-1.900,4.801) (-1.950,4.767) (-2.000,4.733) (-2.045,4.700) (-2.100,4.659) (-2.144,4.625) (-2.200,4.580) (-2.250,4.537) (-2.293,4.500) (-2.347,4.450) (-2.400,4.400) (-2.450,4.350) (-2.497,4.300) (-2.543,4.250) (-2.587,4.200) (-2.625,4.155) (-2.670,4.100) (-2.708,4.050) (-2.746,4.000) (-2.775,3.959) (-2.815,3.900) (-2.848,3.850) (-2.875,3.807) (-2.909,3.750) (-2.938,3.700) (-2.965,3.650) (-2.991,3.600) (-3.016,3.550) (-3.040,3.500) (-3.062,3.450) (-3.084,3.400) (-3.104,3.350) (-3.123,3.300) (-3.141,3.250) (-3.159,3.200) (-3.175,3.150) (-3.190,3.100) (-3.204,3.050) (-3.217,3.000) (-3.228,2.950) (-3.239,2.900) (-3.249,2.850) (-3.258,2.800) (-3.266,2.750) (-3.273,2.700) (-3.279,2.650) (-3.284,2.600) (-3.288,2.550) (-3.292,2.500) (-3.294,2.450) (-3.295,2.400) (-3.295,2.350) (-3.294,2.300) (-3.292,2.250) (-3.289,2.200) (-3.285,2.150) (-3.280,2.100) (-3.274,2.050) (-3.267,2.000) (-3.258,1.950) (-3.249,1.900) (-3.238,1.850) (-3.227,1.800) (-3.214,1.750) (-3.200,1.701) (-3.184,1.650) (-3.168,1.600) (-3.150,1.551) (-3.131,1.500) (-3.110,1.450) (-3.088,1.400) (-3.064,1.350) (-3.039,1.300) (-3.012,1.250) (-2.983,1.200) (-2.953,1.150) (-2.925,1.107) (-2.887,1.050) (-2.850,1.000) (-2.811,0.950) (-2.775,0.906) (-2.727,0.850) (-2.680,0.800) (-2.630,0.750) (-2.577,0.700) (-2.525,0.654) (-2.475,0.613) (-2.427,0.575) (-2.375,0.537) (-2.325,0.502) (-2.275,0.469) (-2.225,0.438) (-2.175,0.408) (-2.125,0.381) (-2.075,0.354) (-2.025,0.329) (-1.975,0.306) (-1.925,0.284) (-1.875,0.263) (-1.825,0.243) (-1.777,0.225) (-1.725,0.207) (-1.675,0.190) (-1.625,0.175) (-1.575,0.160) (-1.525,0.147) (-1.475,0.134) (-1.425,0.122) (-1.375,0.111) (-1.325,0.100) (-1.275,0.091) (-1.225,0.082) (-1.175,0.073) (-1.125,0.066) (-1.075,0.059) (-1.025,0.052) (-0.975,0.046) (-0.925,0.041) (-0.875,0.036) (-0.825,0.032) (-0.775,0.028) (-0.725,0.024) (-0.675,0.021) (-0.625,0.018) (-0.575,0.016) (-0.525,0.013) (-0.475,0.012) (-0.425,0.010) (-0.375,0.009) (-0.325,0.007) (-0.275,0.007) (-0.250,0.006)}
\EQ{(0.250,0.007) (0.300,0.008) (0.350,0.009) (0.400,0.011) (0.450,0.012) (0.500,0.014) (0.550,0.017) (0.600,0.020) (0.650,0.023) (0.700,0.026) (0.750,0.030) (0.800,0.034) (0.850,0.039) (0.900,0.045) (0.950,0.051) (1.000,0.057) (1.050,0.065) (1.100,0.073) (1.150,0.081) (1.200,0.090) (1.250,0.101) (1.300,0.112) (1.350,0.123) (1.400,0.136) (1.450,0.149) (1.500,0.164) (1.550,0.180) (1.600,0.196) (1.650,0.214) (1.700,0.233) (1.750,0.253) (1.800,0.275) (1.850,0.298) (1.900,0.322) (1.950,0.348) (2.000,0.375) (2.050,0.404) (2.100,0.435) (2.150,0.468) (2.196,0.500) (2.250,0.540) (2.294,0.575) (2.350,0.622) (2.400,0.667) (2.450,0.715) (2.500,0.767) (2.550,0.823) (2.594,0.875) (2.625,0.915) (2.669,0.975) (2.700,1.020) (2.735,1.075) (2.765,1.125) (2.793,1.175) (2.820,1.225) (2.844,1.275) (2.867,1.325) (2.888,1.375) (2.908,1.425) (2.926,1.475) (2.943,1.525) (2.959,1.575) (2.973,1.625) (2.986,1.675) (2.997,1.725) (3.008,1.775) (3.017,1.825) (3.024,1.875) (3.031,1.925) (3.037,1.975) (3.041,2.025) (3.044,2.075) (3.047,2.125) (3.048,2.175) (3.048,2.225) (3.047,2.275) (3.044,2.325) (3.041,2.375) (3.037,2.425) (3.032,2.475) (3.025,2.525) (3.018,2.575) (3.009,2.625) (3.000,2.675) (2.989,2.725) (2.977,2.775) (2.964,2.825) (2.950,2.875) (2.935,2.925) (2.919,2.975) (2.902,3.025) (2.883,3.075) (2.864,3.125) (2.843,3.175) (2.821,3.225) (2.798,3.275) (2.775,3.321) (2.750,3.370) (2.725,3.416) (2.700,3.460) (2.662,3.525) (2.630,3.575) (2.600,3.621) (2.563,3.675) (2.527,3.725) (2.489,3.775) (2.450,3.825) (2.408,3.875) (2.375,3.914) (2.325,3.969) (2.275,4.022) (2.225,4.072) (2.175,4.120) (2.125,4.166) (2.087,4.200) (2.028,4.250) (1.975,4.292) (1.933,4.325) (1.875,4.368) (1.829,4.400) (1.775,4.437) (1.725,4.469) (1.676,4.500) (1.625,4.530) (1.575,4.559) (1.525,4.586) (1.475,4.613) (1.425,4.638) (1.375,4.662) (1.325,4.685) (1.275,4.707) (1.225,4.727) (1.175,4.747) (1.125,4.766) (1.075,4.784) (1.025,4.801) (0.975,4.818) (0.925,4.833) (0.875,4.847) (0.825,4.861) (0.775,4.873) (0.725,4.885) (0.675,4.896) (0.625,4.907) (0.575,4.916) (0.525,4.924) (0.475,4.932) (0.425,4.939) (0.375,4.945) (0.325,4.951) (0.275,4.955) (0.225,4.959) (0.175,4.962) (0.125,4.965) (0.075,4.966) (0.025,4.967) (-0.050,4.967) (-0.100,4.965) (-0.150,4.963) (-0.200,4.961) (-0.250,4.957) (-0.300,4.953) (-0.350,4.948) (-0.400,4.942) (-0.450,4.936) (-0.500,4.928) (-0.550,4.920) (-0.600,4.911) (-0.650,4.902) (-0.700,4.891) (-0.750,4.880) (-0.800,4.867) (-0.850,4.854) (-0.900,4.840) (-0.950,4.825) (-1.000,4.810) (-1.050,4.793) (-1.100,4.775) (-1.150,4.757) (-1.200,4.738) (-1.250,4.717) (-1.300,4.696) (-1.346,4.675) (-1.400,4.650) (-1.450,4.625) (-1.499,4.600) (-1.546,4.575) (-1.591,4.550) (-1.650,4.516) (-1.700,4.485) (-1.750,4.453) (-1.793,4.425) (-1.850,4.385) (-1.899,4.350) (-1.950,4.312) (-1.997,4.275) (-2.050,4.231) (-2.100,4.189) (-2.143,4.150) (-2.197,4.100) (-2.248,4.050) (-2.296,4.000) (-2.343,3.950) (-2.387,3.900) (-2.425,3.855) (-2.470,3.800) (-2.508,3.750) (-2.545,3.700) (-2.575,3.658) (-2.614,3.600) (-2.646,3.550) (-2.675,3.503) (-2.706,3.450) (-2.734,3.400) (-2.760,3.350) (-2.786,3.300) (-2.809,3.250) (-2.832,3.200) (-2.853,3.150) (-2.874,3.100) (-2.893,3.050) (-2.911,3.000) (-2.927,2.950) (-2.943,2.900) (-2.957,2.850) (-2.971,2.800) (-2.983,2.750) (-2.994,2.700) (-3.005,2.650) (-3.014,2.600) (-3.022,2.550) (-3.029,2.500) (-3.034,2.450) (-3.039,2.400) (-3.043,2.350) (-3.046,2.300) (-3.047,2.250) (-3.048,2.200) (-3.047,2.150) (-3.046,2.100) (-3.043,2.050) (-3.039,2.000) (-3.034,1.950) (-3.028,1.900) (-3.021,1.850) (-3.012,1.800) (-3.003,1.750) (-2.992,1.700) (-2.979,1.650) (-2.966,1.600) (-2.951,1.550) (-2.935,1.500) (-2.918,1.450) (-2.898,1.400) (-2.878,1.350) (-2.856,1.300) (-2.832,1.250) (-2.807,1.200) (-2.780,1.150) (-2.751,1.100) (-2.720,1.050) (-2.686,1.000) (-2.651,0.950) (-2.613,0.900) (-2.575,0.852) (-2.530,0.800) (-2.484,0.750) (-2.435,0.700) (-2.382,0.650) (-2.325,0.601) (-2.275,0.560) (-2.230,0.525) (-2.175,0.485) (-2.125,0.451) (-2.075,0.419) (-2.025,0.389) (-1.975,0.361) (-1.925,0.335) (-1.875,0.310) (-1.825,0.286) (-1.775,0.264) (-1.725,0.243) (-1.675,0.223) (-1.625,0.205) (-1.575,0.188) (-1.525,0.172) (-1.475,0.157) (-1.425,0.143) (-1.375,0.129) (-1.325,0.117) (-1.275,0.106) (-1.225,0.095) (-1.175,0.086) (-1.125,0.077) (-1.075,0.069) (-1.025,0.061) (-0.975,0.054) (-0.925,0.048) (-0.875,0.042) (-0.825,0.037) (-0.775,0.032) (-0.725,0.028) (-0.675,0.024) (-0.625,0.021) (-0.575,0.018) (-0.525,0.016) (-0.475,0.013) (-0.425,0.012) (-0.375,0.010) (-0.325,0.009) (-0.275,0.008) (-0.250,0.007)}
\EQ{(0.250,0.008) (0.300,0.009) (0.350,0.011) (0.400,0.012) (0.450,0.014) (0.500,0.017) (0.550,0.019) (0.600,0.022) (0.650,0.026) (0.700,0.030) (0.750,0.034) (0.800,0.039) (0.850,0.045) (0.900,0.051) (0.950,0.058) (1.000,0.066) (1.050,0.074) (1.100,0.083) (1.150,0.093) (1.200,0.104) (1.250,0.115) (1.300,0.128) (1.350,0.141) (1.400,0.156) (1.450,0.172) (1.500,0.188) (1.550,0.206) (1.599,0.225) (1.650,0.246) (1.700,0.268) (1.750,0.292) (1.800,0.317) (1.850,0.344) (1.900,0.372) (1.946,0.400) (2.000,0.435) (2.050,0.470) (2.091,0.500) (2.150,0.547) (2.200,0.589) (2.240,0.625) (2.292,0.675) (2.340,0.725) (2.375,0.764) (2.425,0.824) (2.465,0.875) (2.500,0.924) (2.534,0.975) (2.566,1.025) (2.595,1.075) (2.622,1.125) (2.648,1.175) (2.671,1.225) (2.693,1.275) (2.713,1.325) (2.732,1.375) (2.749,1.425) (2.765,1.475) (2.779,1.525) (2.791,1.575) (2.803,1.625) (2.813,1.675) (2.822,1.725) (2.829,1.775) (2.835,1.825) (2.840,1.875) (2.844,1.925) (2.847,1.975) (2.848,2.025) (2.849,2.075) (2.848,2.125) (2.846,2.175) (2.843,2.225) (2.838,2.275) (2.833,2.325) (2.826,2.375) (2.818,2.425) (2.810,2.475) (2.800,2.525) (2.789,2.575) (2.776,2.625) (2.763,2.675) (2.748,2.725) (2.732,2.775) (2.715,2.825) (2.697,2.875) (2.677,2.925) (2.657,2.975) (2.635,3.025) (2.611,3.075) (2.586,3.125) (2.560,3.175) (2.533,3.225) (2.504,3.275) (2.475,3.322) (2.441,3.375) (2.407,3.425) (2.375,3.470) (2.334,3.525) (2.300,3.569) (2.254,3.625) (2.211,3.675) (2.175,3.715) (2.125,3.767) (2.075,3.818) (2.025,3.865) (1.987,3.900) (1.930,3.950) (1.875,3.996) (1.825,4.035) (1.775,4.073) (1.725,4.109) (1.675,4.143) (1.627,4.175) (1.575,4.208) (1.525,4.238) (1.475,4.266) (1.425,4.294) (1.375,4.320) (1.325,4.345) (1.275,4.369) (1.225,4.392) (1.175,4.413) (1.125,4.434) (1.075,4.453) (1.025,4.472) (0.975,4.489) (0.925,4.506) (0.875,4.521) (0.825,4.536) (0.775,4.549) (0.725,4.562) (0.675,4.574) (0.625,4.585) (0.575,4.595) (0.525,4.604) (0.475,4.612) (0.425,4.620) (0.375,4.626) (0.325,4.632) (0.275,4.637) (0.225,4.641) (0.175,4.644) (0.125,4.647) (0.075,4.649) (0.025,4.649) (-0.050,4.649) (-0.100,4.648) (-0.150,4.646) (-0.200,4.643) (-0.250,4.639) (-0.300,4.635) (-0.350,4.629) (-0.400,4.623) (-0.450,4.616) (-0.500,4.608) (-0.550,4.599) (-0.600,4.590) (-0.650,4.579) (-0.700,4.568) (-0.750,4.556) (-0.800,4.543) (-0.850,4.528) (-0.900,4.513) (-0.950,4.497) (-1.000,4.480) (-1.050,4.463) (-1.100,4.444) (-1.150,4.424) (-1.200,4.402) (-1.250,4.380) (-1.300,4.357) (-1.350,4.333) (-1.400,4.307) (-1.450,4.280) (-1.500,4.252) (-1.546,4.225) (-1.600,4.192) (-1.650,4.160) (-1.700,4.126) (-1.750,4.091) (-1.800,4.054) (-1.850,4.016) (-1.900,3.975) (-1.950,3.933) (-2.000,3.888) (-2.041,3.850) (-2.093,3.800) (-2.142,3.750) (-2.188,3.700) (-2.225,3.659) (-2.275,3.600) (-2.315,3.550) (-2.350,3.504) (-2.389,3.450) (-2.424,3.400) (-2.457,3.350) (-2.488,3.300) (-2.518,3.250) (-2.547,3.200) (-2.573,3.150) (-2.599,3.100) (-2.623,3.050) (-2.646,3.000) (-2.667,2.950) (-2.687,2.900) (-2.706,2.850) (-2.724,2.800) (-2.740,2.750) (-2.756,2.700) (-2.770,2.650) (-2.783,2.600) (-2.794,2.550) (-2.805,2.500) (-2.814,2.450) (-2.822,2.400) (-2.830,2.350) (-2.836,2.300) (-2.841,2.250) (-2.844,2.200) (-2.847,2.150) (-2.848,2.100) (-2.849,2.050) (-2.848,2.000) (-2.846,1.950) (-2.842,1.900) (-2.838,1.850) (-2.832,1.800) (-2.826,1.750) (-2.817,1.700) (-2.808,1.650) (-2.797,1.600) (-2.785,1.550) (-2.772,1.500) (-2.757,1.450) (-2.741,1.400) (-2.723,1.350) (-2.703,1.300) (-2.683,1.250) (-2.660,1.200) (-2.635,1.150) (-2.609,1.100) (-2.581,1.050) (-2.550,1.000) (-2.518,0.950) (-2.483,0.900) (-2.450,0.856) (-2.405,0.800) (-2.362,0.750) (-2.325,0.709) (-2.275,0.659) (-2.225,0.611) (-2.184,0.575) (-2.125,0.526) (-2.075,0.488) (-2.025,0.452) (-1.975,0.419) (-1.925,0.387) (-1.875,0.358) (-1.825,0.330) (-1.775,0.304) (-1.725,0.280) (-1.675,0.257) (-1.625,0.236) (-1.575,0.216) (-1.525,0.197) (-1.475,0.180) (-1.425,0.164) (-1.375,0.148) (-1.325,0.134) (-1.275,0.121) (-1.225,0.109) (-1.175,0.098) (-1.125,0.088) (-1.075,0.078) (-1.025,0.070) (-0.975,0.062) (-0.925,0.055) (-0.875,0.048) (-0.825,0.042) (-0.775,0.037) (-0.725,0.032) (-0.675,0.028) (-0.625,0.024) (-0.575,0.021) (-0.525,0.018) (-0.475,0.015) (-0.425,0.013) (-0.375,0.011) (-0.325,0.010) (-0.275,0.009) (-0.250,0.008)}
\EQ{(0.250,0.009) (0.300,0.010) (0.350,0.012) (0.400,0.014) (0.450,0.016) (0.500,0.019) (0.550,0.022) (0.600,0.025) (0.650,0.029) (0.700,0.034) (0.750,0.039) (0.800,0.044) (0.850,0.051) (0.900,0.058) (0.950,0.065) (1.000,0.074) (1.050,0.083) (1.100,0.094) (1.150,0.105) (1.200,0.117) (1.250,0.130) (1.300,0.144) (1.350,0.159) (1.400,0.176) (1.450,0.194) (1.500,0.213) (1.550,0.233) (1.600,0.255) (1.650,0.279) (1.700,0.304) (1.750,0.331) (1.800,0.360) (1.850,0.391) (1.900,0.425) (1.950,0.460) (2.000,0.498) (2.050,0.540) (2.091,0.575) (2.144,0.625) (2.193,0.675) (2.238,0.725) (2.275,0.769) (2.319,0.825) (2.350,0.868) (2.389,0.925) (2.420,0.975) (2.449,1.025) (2.475,1.074) (2.500,1.124) (2.524,1.175) (2.545,1.225) (2.564,1.275) (2.582,1.325) (2.598,1.375) (2.613,1.425) (2.626,1.475) (2.638,1.525) (2.648,1.575) (2.657,1.625) (2.665,1.675) (2.671,1.725) (2.676,1.775) (2.680,1.825) (2.682,1.875) (2.683,1.925) (2.683,1.975) (2.682,2.025) (2.680,2.075) (2.676,2.125) (2.671,2.175) (2.665,2.225) (2.658,2.275) (2.649,2.325) (2.639,2.375) (2.628,2.425) (2.616,2.475) (2.603,2.525) (2.588,2.575) (2.572,2.625) (2.555,2.675) (2.537,2.725) (2.517,2.775) (2.496,2.825) (2.475,2.871) (2.450,2.923) (2.425,2.972) (2.400,3.019) (2.375,3.063) (2.338,3.125) (2.306,3.175) (2.275,3.222) (2.238,3.275) (2.200,3.325) (2.161,3.375) (2.125,3.419) (2.077,3.475) (2.031,3.525) (1.983,3.575) (1.932,3.625) (1.878,3.675) (1.825,3.722) (1.775,3.763) (1.729,3.800) (1.675,3.841) (1.628,3.875) (1.575,3.912) (1.525,3.944) (1.476,3.975) (1.425,4.005) (1.375,4.034) (1.325,4.061) (1.275,4.087) (1.225,4.111) (1.175,4.134) (1.125,4.156) (1.075,4.177) (1.025,4.197) (0.975,4.216) (0.925,4.233) (0.875,4.250) (0.825,4.266) (0.775,4.280) (0.725,4.294) (0.675,4.306) (0.625,4.318) (0.575,4.329) (0.525,4.338) (0.475,4.347) (0.425,4.355) (0.375,4.362) (0.325,4.368) (0.275,4.374) (0.225,4.378) (0.175,4.381) (0.125,4.384) (0.075,4.386) (0.025,4.387) (-0.050,4.386) (-0.100,4.385) (-0.150,4.383) (-0.200,4.380) (-0.250,4.376) (-0.300,4.371) (-0.350,4.365) (-0.400,4.359) (-0.450,4.351) (-0.500,4.343) (-0.550,4.334) (-0.600,4.323) (-0.650,4.312) (-0.700,4.300) (-0.750,4.287) (-0.800,4.273) (-0.850,4.258) (-0.900,4.242) (-0.949,4.225) (-1.000,4.207) (-1.050,4.187) (-1.100,4.167) (-1.150,4.145) (-1.200,4.123) (-1.248,4.100) (-1.298,4.075) (-1.345,4.050) (-1.391,4.025) (-1.434,4.000) (-1.500,3.960) (-1.550,3.928) (-1.592,3.900) (-1.650,3.859) (-1.696,3.825) (-1.750,3.784) (-1.791,3.750) (-1.850,3.700) (-1.900,3.655) (-1.950,3.607) (-2.000,3.557) (-2.050,3.504) (-2.099,3.450) (-2.141,3.400) (-2.175,3.358) (-2.219,3.300) (-2.250,3.258) (-2.290,3.200) (-2.322,3.150) (-2.350,3.105) (-2.383,3.050) (-2.410,3.000) (-2.436,2.950) (-2.461,2.900) (-2.484,2.850) (-2.506,2.800) (-2.527,2.750) (-2.546,2.700) (-2.564,2.650) (-2.580,2.600) (-2.596,2.550) (-2.610,2.500) (-2.622,2.450) (-2.634,2.400) (-2.644,2.350) (-2.653,2.300) (-2.661,2.250) (-2.668,2.200) (-2.674,2.150) (-2.678,2.100) (-2.681,2.050) (-2.683,2.000) (-2.683,1.950) (-2.683,1.900) (-2.681,1.850) (-2.678,1.800) (-2.674,1.750) (-2.668,1.700) (-2.661,1.650) (-2.653,1.600) (-2.643,1.550) (-2.632,1.500) (-2.620,1.450) (-2.606,1.400) (-2.590,1.350) (-2.573,1.300) (-2.555,1.250) (-2.534,1.200) (-2.512,1.150) (-2.488,1.100) (-2.463,1.050) (-2.435,1.000) (-2.405,0.950) (-2.375,0.904) (-2.337,0.850) (-2.300,0.801) (-2.259,0.750) (-2.225,0.711) (-2.175,0.657) (-2.125,0.607) (-2.075,0.561) (-2.033,0.525) (-1.975,0.479) (-1.925,0.442) (-1.875,0.408) (-1.825,0.376) (-1.775,0.346) (-1.725,0.318) (-1.675,0.292) (-1.625,0.267) (-1.575,0.244) (-1.525,0.223) (-1.475,0.203) (-1.425,0.185) (-1.375,0.168) (-1.325,0.152) (-1.275,0.137) (-1.225,0.123) (-1.175,0.111) (-1.125,0.099) (-1.075,0.088) (-1.025,0.078) (-0.975,0.070) (-0.925,0.061) (-0.875,0.054) (-0.825,0.047) (-0.775,0.041) (-0.725,0.036) (-0.675,0.031) (-0.625,0.027) (-0.575,0.023) (-0.525,0.020) (-0.475,0.017) (-0.425,0.015) (-0.375,0.013) (-0.325,0.011) (-0.275,0.010) (-0.250,0.009)}
\EQ{(0.250,0.010) (0.300,0.012) (0.350,0.013) (0.400,0.015) (0.450,0.018) (0.500,0.021) (0.550,0.024) (0.600,0.028) (0.650,0.032) (0.700,0.037) (0.750,0.043) (0.800,0.049) (0.850,0.056) (0.900,0.064) (0.950,0.073) (1.000,0.082) (1.050,0.093) (1.100,0.104) (1.150,0.116) (1.200,0.130) (1.250,0.145) (1.300,0.161) (1.350,0.178) (1.400,0.196) (1.450,0.216) (1.500,0.238) (1.550,0.261) (1.600,0.286) (1.650,0.313) (1.700,0.341) (1.750,0.372) (1.792,0.400) (1.850,0.441) (1.894,0.475) (1.950,0.521) (2.000,0.566) (2.050,0.615) (2.100,0.668) (2.150,0.725) (2.189,0.775) (2.225,0.824) (2.260,0.875) (2.292,0.925) (2.321,0.975) (2.347,1.025) (2.372,1.075) (2.395,1.125) (2.416,1.175) (2.435,1.225) (2.452,1.275) (2.468,1.325) (2.482,1.375) (2.495,1.425) (2.506,1.475) (2.515,1.525) (2.523,1.575) (2.530,1.625) (2.535,1.675) (2.539,1.725) (2.542,1.775) (2.543,1.825) (2.543,1.875) (2.542,1.925) (2.540,1.975) (2.536,2.025) (2.531,2.075) (2.525,2.125) (2.517,2.175) (2.508,2.225) (2.498,2.275) (2.487,2.325) (2.474,2.375) (2.460,2.425) (2.444,2.475) (2.428,2.525) (2.410,2.575) (2.390,2.625) (2.369,2.675) (2.347,2.725) (2.325,2.771) (2.300,2.821) (2.275,2.867) (2.250,2.912) (2.212,2.975) (2.180,3.025) (2.150,3.069) (2.110,3.125) (2.075,3.172) (2.033,3.225) (2.000,3.264) (1.950,3.321) (1.900,3.374) (1.850,3.424) (1.800,3.472) (1.750,3.517) (1.711,3.550) (1.650,3.600) (1.600,3.638) (1.550,3.675) (1.500,3.709) (1.450,3.742) (1.400,3.774) (1.356,3.800) (1.312,3.825) (1.266,3.850) (1.218,3.875) (1.168,3.900) (1.114,3.925) (1.057,3.950) (1.000,3.974) (0.950,3.993) (0.900,4.011) (0.850,4.028) (0.800,4.044) (0.750,4.059) (0.700,4.073) (0.650,4.086) (0.600,4.098) (0.550,4.108) (0.500,4.118) (0.450,4.127) (0.400,4.135) (0.350,4.142) (0.300,4.148) (0.250,4.153) (0.200,4.157) (0.150,4.161) (0.100,4.163) (0.050,4.164) (0.000,4.165) (-0.050,4.164) (-0.100,4.163) (-0.150,4.161) (-0.200,4.157) (-0.250,4.153) (-0.300,4.148) (-0.350,4.142) (-0.400,4.135) (-0.450,4.127) (-0.500,4.118) (-0.550,4.108) (-0.600,4.098) (-0.650,4.086) (-0.700,4.073) (-0.750,4.059) (-0.800,4.044) (-0.850,4.028) (-0.900,4.011) (-0.950,3.993) (-1.000,3.974) (-1.050,3.953) (-1.100,3.931) (-1.150,3.909) (-1.200,3.884) (-1.250,3.859) (-1.300,3.832) (-1.350,3.803) (-1.398,3.775) (-1.450,3.742) (-1.500,3.709) (-1.549,3.675) (-1.600,3.638) (-1.649,3.600) (-1.700,3.559) (-1.740,3.525) (-1.796,3.475) (-1.849,3.425) (-1.899,3.375) (-1.946,3.325) (-1.990,3.275) (-2.025,3.234) (-2.072,3.175) (-2.110,3.125) (-2.146,3.075) (-2.175,3.032) (-2.212,2.975) (-2.242,2.925) (-2.271,2.875) (-2.298,2.825) (-2.323,2.775) (-2.347,2.725) (-2.369,2.675) (-2.390,2.625) (-2.410,2.575) (-2.428,2.525) (-2.444,2.475) (-2.460,2.425) (-2.474,2.375) (-2.487,2.325) (-2.498,2.275) (-2.508,2.225) (-2.517,2.175) (-2.525,2.125) (-2.531,2.075) (-2.536,2.025) (-2.540,1.975) (-2.542,1.925) (-2.543,1.875) (-2.543,1.825) (-2.542,1.775) (-2.539,1.725) (-2.535,1.675) (-2.530,1.625) (-2.523,1.575) (-2.515,1.525) (-2.506,1.475) (-2.495,1.425) (-2.482,1.375) (-2.468,1.325) (-2.452,1.275) (-2.435,1.225) (-2.416,1.175) (-2.395,1.125) (-2.372,1.075) (-2.347,1.025) (-2.321,0.975) (-2.292,0.925) (-2.260,0.875) (-2.226,0.825) (-2.189,0.775) (-2.150,0.726) (-2.107,0.675) (-2.060,0.625) (-2.010,0.575) (-1.954,0.525) (-1.900,0.480) (-1.850,0.441) (-1.800,0.406) (-1.754,0.375) (-1.700,0.341) (-1.650,0.313) (-1.600,0.286) (-1.550,0.261) (-1.500,0.238) (-1.450,0.216) (-1.400,0.196) (-1.350,0.178) (-1.300,0.161) (-1.250,0.145) (-1.200,0.130) (-1.150,0.116) (-1.100,0.104) (-1.050,0.093) (-1.000,0.082) (-0.950,0.073) (-0.900,0.064) (-0.850,0.056) (-0.800,0.049) (-0.750,0.043) (-0.700,0.037) (-0.650,0.032) (-0.600,0.028) (-0.550,0.024) (-0.500,0.021) (-0.450,0.018) (-0.400,0.015) (-0.350,0.013) (-0.300,0.012) (-0.250,0.010)}

    \draw[mmaEQ] (-10,0) -- (10,0);
    % Dipolo físico - líneas de campo
% Generado por scripts/generar_datos.py (no editar a mano).
\LC{(1.086,0.495) (1.147,0.452) (1.204,0.403) (1.254,0.348) (1.298,0.287) (1.334,0.221) (1.360,0.150) (1.376,0.077) (1.382,0.003) (1.377,-0.072) (1.361,-0.146) (1.336,-0.216) (1.301,-0.282) (1.258,-0.344) (1.207,-0.399) (1.151,-0.449) (1.090,-0.492)}
\LC{(0.450,1.109) (0.501,1.164) (0.552,1.219) (0.605,1.272) (0.658,1.325) (0.713,1.376) (0.768,1.426) (0.825,1.476) (0.882,1.524) (0.941,1.571) (1.000,1.617) (1.060,1.662) (1.121,1.706) (1.183,1.748) (1.245,1.790) (1.308,1.830) (1.372,1.869) (1.437,1.907) (1.503,1.943) (1.569,1.979) (1.636,2.013) (1.703,2.046) (1.771,2.077) (1.840,2.108) (1.909,2.137) (1.978,2.165) (2.049,2.191) (2.119,2.216) (2.190,2.240) (2.262,2.262) (2.334,2.284) (2.406,2.303) (2.479,2.322) (2.552,2.339) (2.625,2.354) (2.699,2.369) (2.773,2.381) (2.847,2.393) (2.921,2.403) (2.996,2.411) (3.071,2.418) (3.145,2.424) (3.220,2.428) (3.295,2.430) (3.370,2.431) (3.445,2.430) (3.520,2.428) (3.595,2.425) (3.670,2.420) (3.745,2.413) (3.819,2.404) (3.893,2.394) (3.968,2.383) (4.041,2.370) (4.115,2.355) (4.188,2.338) (4.261,2.320) (4.333,2.301) (4.405,2.279) (4.476,2.256) (4.547,2.231) (4.617,2.205) (4.687,2.177) (4.756,2.147) (4.824,2.115) (4.891,2.082) (4.957,2.047) (5.023,2.011) (5.087,1.973) (5.151,1.933) (5.213,1.891) (5.274,1.848) (5.335,1.803) (5.393,1.756) (5.451,1.708) (5.507,1.658) (5.561,1.606) (5.614,1.553) (5.666,1.499) (5.715,1.443) (5.763,1.385) (5.810,1.326) (5.854,1.265) (5.896,1.204) (5.937,1.140) (5.975,1.076) (6.011,1.010) (6.045,0.943) (6.077,0.875) (6.107,0.807) (6.134,0.737) (6.158,0.666) (6.181,0.594) (6.200,0.522) (6.218,0.449) (6.232,0.375) (6.244,0.301) (6.254,0.227) (6.260,0.152) (6.265,0.077) (6.266,0.002) (6.265,-0.073) (6.261,-0.148) (6.254,-0.222) (6.245,-0.297) (6.233,-0.371) (6.219,-0.444) (6.202,-0.517) (6.182,-0.590) (6.160,-0.661) (6.135,-0.732) (6.108,-0.802) (6.079,-0.871) (6.047,-0.939) (6.013,-1.006) (5.977,-1.072) (5.939,-1.136) (5.899,-1.200) (5.857,-1.262) (5.812,-1.322) (5.766,-1.381) (5.718,-1.439) (5.669,-1.495) (5.617,-1.550) (5.564,-1.603) (5.510,-1.655) (5.454,-1.705) (5.397,-1.753) (5.338,-1.800) (5.278,-1.845) (5.217,-1.888) (5.155,-1.930) (5.091,-1.970) (5.027,-2.008) (4.962,-2.045) (4.895,-2.080) (4.828,-2.113) (4.760,-2.145) (4.691,-2.175) (4.622,-2.203) (4.552,-2.230) (4.481,-2.255) (4.409,-2.278) (4.338,-2.299) (4.265,-2.319) (4.193,-2.337) (4.119,-2.354) (4.046,-2.369) (3.972,-2.382) (3.898,-2.394) (3.824,-2.404) (3.749,-2.412) (3.674,-2.419) (3.600,-2.424) (3.525,-2.428) (3.450,-2.430) (3.375,-2.431) (3.300,-2.430) (3.225,-2.428) (3.150,-2.424) (3.075,-2.418) (3.000,-2.412) (2.926,-2.403) (2.852,-2.393) (2.777,-2.382) (2.704,-2.369) (2.630,-2.355) (2.556,-2.340) (2.483,-2.323) (2.411,-2.305) (2.338,-2.285) (2.266,-2.264) (2.195,-2.241) (2.124,-2.218) (2.053,-2.193) (1.983,-2.166) (1.913,-2.138) (1.844,-2.110) (1.775,-2.079) (1.707,-2.048) (1.640,-2.015) (1.573,-1.981) (1.507,-1.946) (1.441,-1.909) (1.376,-1.871) (1.312,-1.832) (1.249,-1.792) (1.186,-1.751) (1.125,-1.708) (1.064,-1.665) (1.003,-1.620) (0.944,-1.574) (0.886,-1.527) (0.828,-1.479) (0.772,-1.430) (0.716,-1.379) (0.661,-1.328) (0.608,-1.275) (0.555,-1.222) (0.504,-1.167) (0.453,-1.112) (0.437,-1.093)}
\LC{(-0.450,1.109) (-0.501,1.164) (-0.552,1.219) (-0.605,1.272) (-0.658,1.325) (-0.713,1.376) (-0.768,1.426) (-0.825,1.476) (-0.882,1.524) (-0.941,1.571) (-1.000,1.617) (-1.060,1.662) (-1.121,1.706) (-1.183,1.748) (-1.245,1.790) (-1.308,1.830) (-1.372,1.869) (-1.437,1.907) (-1.503,1.943) (-1.569,1.979) (-1.636,2.013) (-1.703,2.046) (-1.771,2.077) (-1.840,2.108) (-1.909,2.137) (-1.978,2.165) (-2.049,2.191) (-2.119,2.216) (-2.190,2.240) (-2.262,2.262) (-2.334,2.284) (-2.406,2.303) (-2.479,2.322) (-2.552,2.339) (-2.625,2.354) (-2.699,2.369) (-2.773,2.381) (-2.847,2.393) (-2.921,2.403) (-2.996,2.411) (-3.071,2.418) (-3.145,2.424) (-3.220,2.428) (-3.295,2.430) (-3.370,2.431) (-3.445,2.430) (-3.520,2.428) (-3.595,2.425) (-3.670,2.420) (-3.745,2.413) (-3.819,2.404) (-3.893,2.394) (-3.968,2.383) (-4.041,2.370) (-4.115,2.355) (-4.188,2.338) (-4.261,2.320) (-4.333,2.301) (-4.405,2.279) (-4.476,2.256) (-4.547,2.231) (-4.617,2.205) (-4.687,2.177) (-4.756,2.147) (-4.824,2.115) (-4.891,2.082) (-4.957,2.047) (-5.023,2.011) (-5.087,1.973) (-5.151,1.933) (-5.213,1.891) (-5.274,1.848) (-5.335,1.803) (-5.393,1.756) (-5.451,1.708) (-5.507,1.658) (-5.561,1.606) (-5.614,1.553) (-5.666,1.499) (-5.715,1.443) (-5.763,1.385) (-5.810,1.326) (-5.854,1.265) (-5.896,1.204) (-5.937,1.140) (-5.975,1.076) (-6.011,1.010) (-6.045,0.943) (-6.077,0.875) (-6.107,0.807) (-6.134,0.737) (-6.158,0.666) (-6.181,0.594) (-6.200,0.522) (-6.218,0.449) (-6.232,0.375) (-6.244,0.301) (-6.254,0.227) (-6.260,0.152) (-6.265,0.077) (-6.266,0.002) (-6.265,-0.073) (-6.261,-0.148) (-6.254,-0.222) (-6.245,-0.297) (-6.233,-0.371) (-6.219,-0.444) (-6.202,-0.517) (-6.182,-0.590) (-6.160,-0.661) (-6.135,-0.732) (-6.108,-0.802) (-6.079,-0.871) (-6.047,-0.939) (-6.013,-1.006) (-5.977,-1.072) (-5.939,-1.136) (-5.899,-1.200) (-5.857,-1.262) (-5.812,-1.322) (-5.766,-1.381) (-5.718,-1.439) (-5.669,-1.495) (-5.617,-1.550) (-5.564,-1.603) (-5.510,-1.655) (-5.454,-1.705) (-5.397,-1.753) (-5.338,-1.800) (-5.278,-1.845) (-5.217,-1.888) (-5.155,-1.930) (-5.091,-1.970) (-5.027,-2.008) (-4.962,-2.045) (-4.895,-2.080) (-4.828,-2.113) (-4.760,-2.145) (-4.691,-2.175) (-4.622,-2.203) (-4.552,-2.230) (-4.481,-2.255) (-4.409,-2.278) (-4.338,-2.299) (-4.265,-2.319) (-4.193,-2.337) (-4.119,-2.354) (-4.046,-2.369) (-3.972,-2.382) (-3.898,-2.394) (-3.824,-2.404) (-3.749,-2.412) (-3.674,-2.419) (-3.600,-2.424) (-3.525,-2.428) (-3.450,-2.430) (-3.375,-2.431) (-3.300,-2.430) (-3.225,-2.428) (-3.150,-2.424) (-3.075,-2.418) (-3.000,-2.412) (-2.926,-2.403) (-2.852,-2.393) (-2.777,-2.382) (-2.704,-2.369) (-2.630,-2.355) (-2.556,-2.340) (-2.483,-2.323) (-2.411,-2.305) (-2.338,-2.285) (-2.266,-2.264) (-2.195,-2.241) (-2.124,-2.218) (-2.053,-2.193) (-1.983,-2.166) (-1.913,-2.138) (-1.844,-2.110) (-1.775,-2.079) (-1.707,-2.048) (-1.640,-2.015) (-1.573,-1.981) (-1.507,-1.946) (-1.441,-1.909) (-1.376,-1.871) (-1.312,-1.832) (-1.249,-1.792) (-1.186,-1.751) (-1.125,-1.708) (-1.064,-1.665) (-1.003,-1.620) (-0.944,-1.574) (-0.886,-1.527) (-0.828,-1.479) (-0.772,-1.430) (-0.716,-1.379) (-0.661,-1.328) (-0.608,-1.275) (-0.555,-1.222) (-0.504,-1.167) (-0.453,-1.112) (-0.437,-1.093)}
\LC{(-1.086,0.495) (-1.147,0.452) (-1.204,0.403) (-1.254,0.348) (-1.298,0.287) (-1.334,0.221) (-1.360,0.150) (-1.376,0.077) (-1.382,0.003) (-1.377,-0.072) (-1.361,-0.146) (-1.336,-0.216) (-1.301,-0.282) (-1.258,-0.344) (-1.207,-0.399) (-1.151,-0.449) (-1.090,-0.492)}
\LC{(4.895,-10.000) (4.847,-9.949) (4.799,-9.892) (4.752,-9.834) (4.704,-9.776) (4.656,-9.718) (4.609,-9.660) (4.562,-9.601) (4.515,-9.543) (4.468,-9.484) (4.421,-9.426) (4.375,-9.367) (4.329,-9.308) (4.282,-9.249) (4.236,-9.190) (4.190,-9.130) (4.145,-9.071) (4.099,-9.011) (4.053,-8.952) (4.008,-8.892) (3.963,-8.832) (3.918,-8.772) (3.873,-8.712) (3.829,-8.652) (3.784,-8.591) (3.740,-8.531) (3.696,-8.470) (3.652,-8.409) (3.608,-8.349) (3.564,-8.288) (3.521,-8.226) (3.477,-8.165) (3.434,-8.104) (3.391,-8.042) (3.348,-7.981) (3.306,-7.919) (3.263,-7.857) (3.221,-7.795) (3.179,-7.733) (3.137,-7.671) (3.095,-7.609) (3.053,-7.547) (3.012,-7.484) (2.971,-7.421) (2.929,-7.359) (2.889,-7.296) (2.848,-7.233) (2.807,-7.170) (2.767,-7.107) (2.727,-7.043) (2.687,-6.980) (2.647,-6.916) (2.607,-6.852) (2.568,-6.789) (2.529,-6.725) (2.490,-6.661) (2.451,-6.597) (2.412,-6.532) (2.374,-6.468) (2.336,-6.403) (2.298,-6.339) (2.260,-6.274) (2.222,-6.209) (2.185,-6.144) (2.147,-6.079) (2.110,-6.014) (2.073,-5.948) (2.037,-5.883) (2.000,-5.817) (1.964,-5.752) (1.928,-5.686) (1.892,-5.620) (1.857,-5.554) (1.821,-5.488) (1.786,-5.422) (1.751,-5.355) (1.717,-5.289) (1.682,-5.222) (1.648,-5.156) (1.614,-5.089) (1.580,-5.022) (1.546,-4.955) (1.513,-4.888) (1.480,-4.820) (1.447,-4.753) (1.414,-4.685) (1.382,-4.618) (1.350,-4.550) (1.318,-4.482) (1.286,-4.414) (1.255,-4.346) (1.223,-4.278) (1.192,-4.210) (1.162,-4.141) (1.131,-4.073) (1.101,-4.004) (1.071,-3.935) (1.041,-3.866) (1.012,-3.797) (0.983,-3.728) (0.954,-3.659) (0.925,-3.590) (0.897,-3.520) (0.869,-3.451) (0.841,-3.381) (0.814,-3.311) (0.786,-3.241) (0.760,-3.171) (0.733,-3.101) (0.707,-3.031) (0.680,-2.961) (0.655,-2.890) (0.629,-2.820) (0.604,-2.749) (0.579,-2.678) (0.555,-2.607) (0.531,-2.536) (0.507,-2.465) (0.483,-2.394) (0.460,-2.323) (0.437,-2.251) (0.415,-2.180) (0.393,-2.108) (0.371,-2.036) (0.349,-1.964) (0.328,-1.893) (0.307,-1.820) (0.287,-1.748) (0.267,-1.676) (0.247,-1.604) (0.228,-1.531)}
\LC{(-5.083,-10.000) (-5.034,-9.952) (-4.985,-9.896) (-4.935,-9.839) (-4.886,-9.782) (-4.838,-9.725) (-4.789,-9.668) (-4.741,-9.611) (-4.692,-9.553) (-4.644,-9.496) (-4.596,-9.438) (-4.548,-9.381) (-4.500,-9.323) (-4.453,-9.265) (-4.405,-9.207) (-4.358,-9.148) (-4.311,-9.090) (-4.264,-9.032) (-4.217,-8.973) (-4.171,-8.914) (-4.124,-8.855) (-4.078,-8.796) (-4.032,-8.737) (-3.986,-8.678) (-3.940,-8.619) (-3.894,-8.559) (-3.849,-8.499) (-3.804,-8.440) (-3.758,-8.380) (-3.713,-8.320) (-3.669,-8.260) (-3.624,-8.199) (-3.580,-8.139) (-3.535,-8.078) (-3.491,-8.018) (-3.447,-7.957) (-3.403,-7.896) (-3.360,-7.835) (-3.316,-7.774) (-3.273,-7.713) (-3.230,-7.651) (-3.187,-7.590) (-3.144,-7.528) (-3.102,-7.466) (-3.060,-7.404) (-3.017,-7.342) (-2.976,-7.280) (-2.934,-7.218) (-2.892,-7.155) (-2.851,-7.093) (-2.809,-7.030) (-2.768,-6.968) (-2.728,-6.905) (-2.687,-6.842) (-2.646,-6.778) (-2.606,-6.715) (-2.566,-6.652) (-2.526,-6.588) (-2.487,-6.525) (-2.447,-6.461) (-2.408,-6.397) (-2.369,-6.333) (-2.330,-6.269) (-2.291,-6.205) (-2.253,-6.140) (-2.214,-6.076) (-2.176,-6.011) (-2.139,-5.946) (-2.101,-5.881) (-2.064,-5.816) (-2.026,-5.751) (-1.989,-5.686) (-1.953,-5.621) (-1.916,-5.555) (-1.880,-5.490) (-1.844,-5.424) (-1.808,-5.358) (-1.772,-5.292) (-1.737,-5.226) (-1.701,-5.160) (-1.666,-5.093) (-1.632,-5.027) (-1.597,-4.960) (-1.563,-4.894) (-1.529,-4.827) (-1.495,-4.760) (-1.461,-4.693) (-1.428,-4.626) (-1.395,-4.558) (-1.362,-4.491) (-1.329,-4.423) (-1.297,-4.356) (-1.265,-4.288) (-1.233,-4.220) (-1.202,-4.152) (-1.170,-4.084) (-1.139,-4.015) (-1.108,-3.947) (-1.078,-3.879) (-1.048,-3.810) (-1.018,-3.741) (-0.988,-3.672) (-0.959,-3.603) (-0.929,-3.534) (-0.900,-3.465) (-0.872,-3.396) (-0.844,-3.326) (-0.816,-3.257) (-0.788,-3.187) (-0.760,-3.117) (-0.733,-3.047) (-0.706,-2.977) (-0.680,-2.907) (-0.654,-2.837) (-0.628,-2.766) (-0.602,-2.696) (-0.577,-2.625) (-0.552,-2.555) (-0.528,-2.484) (-0.503,-2.413) (-0.479,-2.342) (-0.456,-2.270) (-0.433,-2.199) (-0.410,-2.128) (-0.387,-2.056) (-0.365,-1.984) (-0.343,-1.913) (-0.322,-1.841) (-0.301,-1.769) (-0.280,-1.697) (-0.260,-1.624) (-0.240,-1.552)}
\LC{(10.000,-6.214) (9.948,-6.223) (9.874,-6.231) (9.799,-6.239) (9.724,-6.246) (9.650,-6.253) (9.575,-6.259) (9.500,-6.264) (9.425,-6.269) (9.350,-6.273) (9.275,-6.277) (9.200,-6.280) (9.125,-6.282) (9.050,-6.284) (8.976,-6.285) (8.901,-6.286) (8.826,-6.286) (8.751,-6.285) (8.676,-6.284) (8.601,-6.282) (8.526,-6.280) (8.451,-6.277) (8.376,-6.273) (8.301,-6.269) (8.226,-6.264) (8.151,-6.259) (8.076,-6.253) (8.002,-6.246) (7.927,-6.239) (7.852,-6.232) (7.778,-6.223) (7.703,-6.215) (7.629,-6.205) (7.555,-6.195) (7.480,-6.185) (7.406,-6.174) (7.332,-6.162) (7.258,-6.150) (7.184,-6.137) (7.110,-6.124) (7.037,-6.110) (6.963,-6.096) (6.890,-6.081) (6.816,-6.065) (6.743,-6.049) (6.670,-6.033) (6.597,-6.016) (6.524,-5.998) (6.451,-5.980) (6.379,-5.961) (6.306,-5.942) (6.234,-5.922) (6.162,-5.901) (6.089,-5.880) (6.018,-5.859) (5.946,-5.837) (5.874,-5.815) (5.803,-5.791) (5.732,-5.768) (5.661,-5.744) (5.590,-5.719) (5.519,-5.694) (5.449,-5.668) (5.378,-5.642) (5.308,-5.616) (5.239,-5.588) (5.169,-5.561) (5.099,-5.532) (5.030,-5.504) (4.961,-5.475) (4.892,-5.445) (4.824,-5.415) (4.755,-5.384) (4.687,-5.352) (4.619,-5.321) (4.551,-5.288) (4.484,-5.256) (4.417,-5.222) (4.350,-5.189) (4.283,-5.154) (4.216,-5.120) (4.150,-5.084) (4.084,-5.049) (4.019,-5.013) (3.953,-4.976) (3.888,-4.939) (3.823,-4.901) (3.759,-4.863) (3.694,-4.824) (3.630,-4.785) (3.567,-4.746) (3.503,-4.706) (3.440,-4.665) (3.377,-4.624) (3.315,-4.583) (3.253,-4.541) (3.191,-4.498) (3.129,-4.455) (3.068,-4.412) (3.007,-4.368) (2.947,-4.324) (2.886,-4.279) (2.827,-4.234) (2.767,-4.188) (2.708,-4.142) (2.649,-4.096) (2.591,-4.049) (2.533,-4.001) (2.475,-3.953) (2.418,-3.905) (2.361,-3.856) (2.304,-3.807) (2.248,-3.757) (2.192,-3.707) (2.137,-3.656) (2.082,-3.605) (2.027,-3.554) (1.973,-3.502) (1.920,-3.449) (1.866,-3.396) (1.814,-3.343) (1.761,-3.290) (1.709,-3.235) (1.658,-3.181) (1.607,-3.126) (1.556,-3.070) (1.506,-3.015) (1.457,-2.958) (1.408,-2.902) (1.359,-2.845) (1.311,-2.787) (1.263,-2.729) (1.216,-2.671) (1.170,-2.612) (1.124,-2.553) (1.078,-2.493) (1.033,-2.433) (0.989,-2.372) (0.945,-2.312) (0.931,-2.291)}
\LC{(0.308,1.617) (0.332,1.688) (0.356,1.759) (0.381,1.830) (0.406,1.901) (0.432,1.971) (0.458,2.041) (0.484,2.112) (0.511,2.182) (0.539,2.251) (0.566,2.321) (0.595,2.391) (0.623,2.460) (0.652,2.529) (0.681,2.598) (0.711,2.667) (0.741,2.736) (0.772,2.804) (0.802,2.873) (0.834,2.941) (0.865,3.009) (0.897,3.077) (0.929,3.145) (0.962,3.212) (0.995,3.280) (1.028,3.347) (1.062,3.414) (1.096,3.481) (1.130,3.547) (1.164,3.614) (1.199,3.680) (1.235,3.746) (1.270,3.813) (1.306,3.878) (1.342,3.944) (1.379,4.010) (1.415,4.075) (1.453,4.140) (1.490,4.205) (1.528,4.270) (1.566,4.335) (1.604,4.399) (1.642,4.464) (1.681,4.528) (1.720,4.592) (1.760,4.656) (1.799,4.719) (1.839,4.783) (1.879,4.846) (1.920,4.909) (1.961,4.972) (2.002,5.035) (2.043,5.098) (2.084,5.160) (2.126,5.222) (2.168,5.284) (2.211,5.346) (2.253,5.408) (2.296,5.470) (2.339,5.531) (2.382,5.592) (2.426,5.653) (2.470,5.714) (2.514,5.775) (2.558,5.836) (2.603,5.896) (2.647,5.956) (2.692,6.016) (2.738,6.076) (2.783,6.136) (2.829,6.195) (2.875,6.254) (2.921,6.313) (2.967,6.372) (3.014,6.431) (3.061,6.490) (3.108,6.548) (3.155,6.606) (3.202,6.664) (3.250,6.722) (3.298,6.780) (3.346,6.838) (3.394,6.895) (3.443,6.952) (3.492,7.009) (3.541,7.066) (3.590,7.122) (3.639,7.179) (3.689,7.235) (3.739,7.291) (3.789,7.347) (3.839,7.403) (3.889,7.458) (3.940,7.514) (3.991,7.569) (4.042,7.624) (4.093,7.679) (4.144,7.733) (4.196,7.788) (4.248,7.842) (4.299,7.896) (4.352,7.950) (4.404,8.004) (4.457,8.057) (4.509,8.111) (4.562,8.164) (4.615,8.217) (4.668,8.270) (4.722,8.322) (4.776,8.375) (4.829,8.427) (4.883,8.479) (4.937,8.531) (4.992,8.583) (5.046,8.634) (5.101,8.685) (5.156,8.737) (5.211,8.788) (5.266,8.838) (5.322,8.889) (5.377,8.939) (5.433,8.989) (5.489,9.039) (5.545,9.089) (5.601,9.139) (5.657,9.188) (5.714,9.238) (5.771,9.287) (5.828,9.336) (5.885,9.384) (5.942,9.433) (5.999,9.481) (6.057,9.529) (6.114,9.577) (6.172,9.625) (6.230,9.672) (6.288,9.720) (6.347,9.767) (6.405,9.814) (6.464,9.861) (6.522,9.907) (6.581,9.954) (6.640,10.000)}
\LC{(3.378,2.882) (3.452,2.893) (3.526,2.904) (3.601,2.913) (3.675,2.921) (3.750,2.927) (3.825,2.933) (3.900,2.937) (3.975,2.940) (4.050,2.942) (4.125,2.942) (4.200,2.941) (4.275,2.939) (4.349,2.936) (4.424,2.931) (4.499,2.925) (4.574,2.918) (4.648,2.909) (4.723,2.899) (4.797,2.888) (4.871,2.876) (4.944,2.862) (5.018,2.846) (5.091,2.830) (5.164,2.812) (5.236,2.792) (5.308,2.772) (5.380,2.750) (5.451,2.726) (5.522,2.701) (5.592,2.675) (5.662,2.647) (5.731,2.618) (5.800,2.588) (5.867,2.556) (5.935,2.523) (6.001,2.488) (6.067,2.452) (6.132,2.414) (6.196,2.376) (6.259,2.335) (6.322,2.294) (6.383,2.251) (6.444,2.206) (6.503,2.160) (6.561,2.113) (6.619,2.065) (6.675,2.015) (6.729,1.964) (6.783,1.912) (6.835,1.858) (6.886,1.803) (6.936,1.747) (6.984,1.689) (7.031,1.630) (7.076,1.570) (7.120,1.509) (7.162,1.447) (7.202,1.384) (7.241,1.320) (7.277,1.254) (7.312,1.188) (7.346,1.121) (7.377,1.053) (7.407,0.984) (7.434,0.914) (7.460,0.844) (7.483,0.772) (7.505,0.701) (7.524,0.628) (7.542,0.555) (7.557,0.482) (7.570,0.408) (7.581,0.334) (7.589,0.259) (7.596,0.184) (7.600,0.110) (7.602,0.035) (7.602,-0.040) (7.600,-0.115) (7.596,-0.190) (7.589,-0.265) (7.580,-0.339) (7.569,-0.414) (7.556,-0.487) (7.540,-0.561) (7.523,-0.634) (7.503,-0.706) (7.482,-0.778) (7.458,-0.849) (7.432,-0.920) (7.404,-0.989) (7.375,-1.058) (7.343,-1.126) (7.310,-1.193) (7.275,-1.260) (7.238,-1.325) (7.199,-1.389) (7.158,-1.452) (7.116,-1.514) (7.073,-1.575) (7.027,-1.635) (6.981,-1.693) (6.932,-1.751) (6.883,-1.807) (6.831,-1.862) (6.779,-1.916) (6.725,-1.968) (6.670,-2.019) (6.614,-2.069) (6.557,-2.117) (6.498,-2.164) (6.439,-2.210) (6.378,-2.254) (6.317,-2.297) (6.254,-2.338) (6.191,-2.379) (6.127,-2.417) (6.062,-2.455) (5.996,-2.491) (5.929,-2.525) (5.862,-2.558) (5.794,-2.590) (5.726,-2.620) (5.656,-2.649) (5.587,-2.677) (5.516,-2.703) (5.446,-2.728) (5.374,-2.751) (5.303,-2.773) (5.231,-2.794) (5.158,-2.813) (5.085,-2.831) (5.012,-2.848) (4.939,-2.863) (4.865,-2.877) (4.791,-2.889) (4.717,-2.900) (4.642,-2.910) (4.568,-2.918) (4.493,-2.926) (4.419,-2.932) (4.344,-2.936) (4.269,-2.939) (4.194,-2.941) (4.119,-2.942) (4.044,-2.942) (3.969,-2.940) (3.894,-2.937) (3.819,-2.932) (3.744,-2.927) (3.670,-2.920) (3.595,-2.912) (3.521,-2.903) (3.446,-2.892) (3.397,-2.885)}
\LC{(-1.547,2.471) (-1.606,2.518) (-1.665,2.564) (-1.724,2.609) (-1.785,2.654) (-1.846,2.698) (-1.907,2.741) (-1.969,2.783) (-2.031,2.825) (-2.094,2.866) (-2.158,2.906) (-2.222,2.945) (-2.286,2.983) (-2.351,3.021) (-2.416,3.058) (-2.482,3.094) (-2.548,3.130) (-2.615,3.164) (-2.682,3.198) (-2.749,3.231) (-2.817,3.263) (-2.885,3.294) (-2.953,3.325) (-3.022,3.354) (-3.091,3.383) (-3.161,3.411) (-3.231,3.438) (-3.301,3.465) (-3.372,3.490) (-3.442,3.515) (-3.514,3.539) (-3.585,3.562) (-3.657,3.584) (-3.729,3.605) (-3.801,3.626) (-3.873,3.645) (-3.946,3.664) (-4.019,3.681) (-4.092,3.698) (-4.165,3.714) (-4.239,3.729) (-4.312,3.743) (-4.386,3.757) (-4.460,3.769) (-4.534,3.780) (-4.608,3.791) (-4.683,3.800) (-4.757,3.809) (-4.832,3.817) (-4.907,3.824) (-4.981,3.830) (-5.056,3.834) (-5.131,3.838) (-5.206,3.841) (-5.281,3.843) (-5.356,3.844) (-5.431,3.845) (-5.506,3.844) (-5.581,3.842) (-5.656,3.839) (-5.731,3.835) (-5.806,3.830) (-5.880,3.825) (-5.955,3.818) (-6.030,3.810) (-6.104,3.801) (-6.179,3.791) (-6.253,3.780) (-6.327,3.769) (-6.401,3.756) (-6.474,3.742) (-6.548,3.727) (-6.621,3.711) (-6.694,3.694) (-6.767,3.676) (-6.839,3.656) (-6.912,3.636) (-6.984,3.615) (-7.055,3.593) (-7.126,3.569) (-7.197,3.545) (-7.268,3.519) (-7.338,3.493) (-7.408,3.465) (-7.477,3.437) (-7.546,3.407) (-7.614,3.376) (-7.682,3.344) (-7.750,3.311) (-7.816,3.277) (-7.883,3.242) (-7.948,3.206) (-8.013,3.169) (-8.078,3.130) (-8.142,3.091) (-8.205,3.050) (-8.267,3.009) (-8.329,2.966) (-8.390,2.923) (-8.450,2.878) (-8.510,2.832) (-8.568,2.785) (-8.626,2.738) (-8.683,2.689) (-8.739,2.639) (-8.794,2.588) (-8.848,2.536) (-8.901,2.483) (-8.954,2.429) (-9.005,2.374) (-9.055,2.319) (-9.104,2.262) (-9.152,2.204) (-9.199,2.146) (-9.244,2.086) (-9.289,2.026) (-9.332,1.964) (-9.374,1.902) (-9.415,1.839) (-9.454,1.775) (-9.492,1.711) (-9.529,1.646) (-9.564,1.579) (-9.598,1.513) (-9.631,1.445) (-9.662,1.377) (-9.692,1.308) (-9.720,1.239) (-9.747,1.168) (-9.772,1.098) (-9.796,1.027) (-9.818,0.955) (-9.839,0.883) (-9.857,0.810) (-9.875,0.737) (-9.891,0.664) (-9.905,0.590) (-9.917,0.516) (-9.928,0.442) (-9.937,0.368) (-9.944,0.293) (-9.950,0.218) (-9.954,0.143) (-9.957,0.068) (-9.957,-0.007) (-9.956,-0.082) (-9.954,-0.156) (-9.949,-0.231) (-9.943,-0.306) (-9.935,-0.381) (-9.926,-0.455) (-9.915,-0.529) (-9.902,-0.603) (-9.888,-0.677) (-9.872,-0.750) (-9.854,-0.823) (-9.835,-0.895) (-9.814,-0.968) (-9.792,-1.039) (-9.768,-1.110) (-9.742,-1.181) (-9.715,-1.251) (-9.687,-1.320) (-9.657,-1.389) (-9.626,-1.457) (-9.593,-1.524) (-9.558,-1.591) (-9.523,-1.657) (-9.486,-1.722) (-9.447,-1.787) (-9.408,-1.850) (-9.367,-1.913) (-9.325,-1.975) (-9.281,-2.036) (-9.236,-2.097) (-9.191,-2.156) (-9.144,-2.214) (-9.095,-2.272) (-9.046,-2.328) (-8.996,-2.384) (-8.945,-2.439) (-8.892,-2.492) (-8.839,-2.545) (-8.785,-2.597) (-8.729,-2.648) (-8.673,-2.697) (-8.616,-2.746) (-8.558,-2.794) (-8.499,-2.840) (-8.440,-2.886) (-8.379,-2.930) (-8.318,-2.974) (-8.256,-3.016) (-8.194,-3.057) (-8.131,-3.098) (-8.067,-3.137) (-8.002,-3.175) (-7.937,-3.212) (-7.871,-3.248) (-7.805,-3.283) (-7.738,-3.317) (-7.670,-3.350) (-7.602,-3.381) (-7.534,-3.412) (-7.465,-3.442) (-7.396,-3.470) (-7.326,-3.498) (-7.256,-3.524) (-7.185,-3.549) (-7.114,-3.574) (-7.043,-3.597) (-6.971,-3.619) (-6.899,-3.640) (-6.827,-3.660) (-6.754,-3.679) (-6.681,-3.697) (-6.608,-3.714) (-6.535,-3.729) (-6.462,-3.744) (-6.388,-3.758) (-6.314,-3.771) (-6.240,-3.782) (-6.166,-3.793) (-6.091,-3.803) (-6.017,-3.811) (-5.942,-3.819) (-5.867,-3.826) (-5.793,-3.831) (-5.718,-3.836) (-5.643,-3.840) (-5.568,-3.842) (-5.493,-3.844) (-5.418,-3.845) (-5.343,-3.844) (-5.268,-3.843) (-5.193,-3.841) (-5.118,-3.838) (-5.043,-3.834) (-4.968,-3.829) (-4.894,-3.823) (-4.819,-3.816) (-4.744,-3.808) (-4.670,-3.799) (-4.595,-3.789) (-4.521,-3.778) (-4.447,-3.767) (-4.373,-3.754) (-4.299,-3.741) (-4.226,-3.727) (-4.152,-3.711) (-4.079,-3.695) (-4.006,-3.678) (-3.933,-3.660) (-3.861,-3.642) (-3.788,-3.622) (-3.716,-3.601) (-3.644,-3.580) (-3.573,-3.558) (-3.501,-3.535) (-3.430,-3.511) (-3.359,-3.486) (-3.289,-3.460) (-3.219,-3.434) (-3.149,-3.406) (-3.079,-3.378) (-3.010,-3.349) (-2.941,-3.319) (-2.873,-3.289) (-2.805,-3.257) (-2.737,-3.225) (-2.670,-3.192) (-2.603,-3.158) (-2.536,-3.123) (-2.470,-3.088) (-2.405,-3.052) (-2.339,-3.015) (-2.275,-2.977) (-2.210,-2.938) (-2.147,-2.899) (-2.083,-2.859) (-2.020,-2.818) (-1.958,-2.776) (-1.896,-2.733) (-1.835,-2.690) (-1.774,-2.646) (-1.714,-2.602) (-1.654,-2.556) (-1.595,-2.510) (-1.556,-2.479)}
\LC{(-0.049,-10.000) (-0.048,-9.947) (-0.048,-9.872) (-0.047,-9.797) (-0.047,-9.722) (-0.046,-9.647) (-0.045,-9.572) (-0.045,-9.497) (-0.044,-9.422) (-0.044,-9.347) (-0.043,-9.272) (-0.043,-9.197) (-0.042,-9.122) (-0.042,-9.047) (-0.041,-8.972) (-0.041,-8.897) (-0.040,-8.822) (-0.040,-8.747) (-0.039,-8.672) (-0.039,-8.597) (-0.038,-8.522) (-0.038,-8.447) (-0.037,-8.372) (-0.037,-8.297) (-0.036,-8.222) (-0.036,-8.147) (-0.035,-8.072) (-0.035,-7.997) (-0.034,-7.922) (-0.034,-7.847) (-0.033,-7.772) (-0.033,-7.697) (-0.032,-7.622) (-0.032,-7.547) (-0.031,-7.472) (-0.031,-7.397) (-0.030,-7.322) (-0.030,-7.247) (-0.029,-7.172) (-0.029,-7.097) (-0.029,-7.022) (-0.028,-6.947) (-0.028,-6.872) (-0.027,-6.797) (-0.027,-6.722) (-0.026,-6.647) (-0.026,-6.572) (-0.025,-6.497) (-0.025,-6.422) (-0.024,-6.347) (-0.024,-6.272) (-0.024,-6.197) (-0.023,-6.122) (-0.023,-6.047) (-0.022,-5.972) (-0.022,-5.897) (-0.021,-5.822) (-0.021,-5.747) (-0.021,-5.672) (-0.020,-5.597) (-0.020,-5.522) (-0.019,-5.447) (-0.019,-5.372) (-0.019,-5.297) (-0.018,-5.222) (-0.018,-5.147) (-0.017,-5.072) (-0.017,-4.997) (-0.017,-4.922) (-0.016,-4.847) (-0.016,-4.772) (-0.016,-4.697) (-0.015,-4.622) (-0.015,-4.547) (-0.014,-4.472) (-0.014,-4.397) (-0.014,-4.322) (-0.013,-4.247) (-0.013,-4.172) (-0.013,-4.097) (-0.012,-4.022) (-0.012,-3.947) (-0.012,-3.872) (-0.011,-3.797) (-0.011,-3.722) (-0.011,-3.647) (-0.010,-3.572) (-0.010,-3.497) (-0.010,-3.422) (-0.009,-3.347) (-0.009,-3.272) (-0.009,-3.197) (-0.008,-3.122) (-0.008,-3.047) (-0.008,-2.972) (-0.007,-2.897) (-0.007,-2.822) (-0.007,-2.747) (-0.006,-2.672) (-0.006,-2.597) (-0.006,-2.522) (-0.006,-2.447) (-0.005,-2.372) (-0.005,-2.347)}
\LC{(-0.371,1.674) (-0.398,1.743) (-0.426,1.813) (-0.454,1.883) (-0.483,1.952) (-0.512,2.021) (-0.541,2.090) (-0.572,2.159) (-0.602,2.227) (-0.633,2.295) (-0.665,2.363) (-0.697,2.431) (-0.729,2.499) (-0.762,2.566) (-0.795,2.634) (-0.829,2.701) (-0.863,2.768) (-0.897,2.834) (-0.932,2.901) (-0.967,2.967) (-1.003,3.033) (-1.039,3.098) (-1.075,3.164) (-1.112,3.229) (-1.149,3.295) (-1.187,3.359) (-1.225,3.424) (-1.263,3.489) (-1.302,3.553) (-1.341,3.617) (-1.380,3.681) (-1.420,3.745) (-1.460,3.808) (-1.500,3.871) (-1.541,3.934) (-1.582,3.997) (-1.623,4.059) (-1.665,4.122) (-1.707,4.184) (-1.749,4.246) (-1.792,4.307) (-1.835,4.369) (-1.878,4.430) (-1.922,4.491) (-1.966,4.552) (-2.010,4.613) (-2.054,4.673) (-2.099,4.733) (-2.144,4.793) (-2.190,4.853) (-2.235,4.912) (-2.281,4.971) (-2.328,5.030) (-2.374,5.089) (-2.421,5.148) (-2.468,5.206) (-2.516,5.264) (-2.563,5.322) (-2.611,5.380) (-2.659,5.437) (-2.708,5.495) (-2.757,5.552) (-2.806,5.608) (-2.855,5.665) (-2.904,5.721) (-2.954,5.777) (-3.004,5.833) (-3.054,5.889) (-3.105,5.944) (-3.156,6.000) (-3.207,6.055) (-3.258,6.109) (-3.310,6.164) (-3.361,6.218) (-3.413,6.272) (-3.466,6.326) (-3.518,6.380) (-3.571,6.433) (-3.624,6.486) (-3.677,6.539) (-3.730,6.592) (-3.784,6.644) (-3.838,6.696) (-3.892,6.748) (-3.946,6.800) (-4.001,6.851) (-4.055,6.903) (-4.110,6.954) (-4.166,7.004) (-4.221,7.055) (-4.277,7.105) (-4.332,7.155) (-4.388,7.205) (-4.445,7.255) (-4.501,7.304) (-4.558,7.353) (-4.615,7.402) (-4.672,7.451) (-4.729,7.499) (-4.786,7.548) (-4.844,7.595) (-4.902,7.643) (-4.960,7.691) (-5.018,7.738) (-5.077,7.785) (-5.135,7.832) (-5.194,7.878) (-5.253,7.924) (-5.313,7.970) (-5.372,8.016) (-5.432,8.062) (-5.491,8.107) (-5.551,8.152) (-5.611,8.197) (-5.672,8.241) (-5.732,8.286) (-5.793,8.330) (-5.854,8.374) (-5.915,8.417) (-5.976,8.461) (-6.038,8.504) (-6.099,8.546) (-6.161,8.589) (-6.223,8.631) (-6.285,8.673) (-6.347,8.715) (-6.409,8.757) (-6.472,8.798) (-6.535,8.839) (-6.598,8.880) (-6.661,8.921) (-6.724,8.961) (-6.787,9.001) (-6.851,9.041) (-6.914,9.081) (-6.978,9.120) (-7.042,9.159) (-7.107,9.198) (-7.171,9.237) (-7.235,9.275) (-7.300,9.313) (-7.365,9.351) (-7.430,9.388) (-7.495,9.426) (-7.560,9.463) (-7.625,9.500) (-7.691,9.536) (-7.756,9.572) (-7.822,9.609) (-7.888,9.644) (-7.954,9.680) (-8.020,9.715) (-8.087,9.750) (-8.153,9.785) (-8.220,9.819) (-8.286,9.854) (-8.353,9.888) (-8.420,9.921) (-8.487,9.955) (-8.555,9.988) (-8.600,10.000)}
\LC{(-1.946,3.725) (-1.997,3.779) (-2.049,3.834) (-2.101,3.888) (-2.154,3.941) (-2.206,3.995) (-2.260,4.047) (-2.313,4.100) (-2.367,4.152) (-2.421,4.204) (-2.476,4.255) (-2.531,4.306) (-2.586,4.357) (-2.642,4.407) (-2.698,4.457) (-2.754,4.506) (-2.811,4.556) (-2.868,4.604) (-2.925,4.653) (-2.983,4.701) (-3.041,4.748) (-3.099,4.795) (-3.158,4.842) (-3.217,4.889) (-3.276,4.935) (-3.336,4.980) (-3.395,5.025) (-3.456,5.070) (-3.516,5.115) (-3.577,5.159) (-3.638,5.202) (-3.699,5.246) (-3.760,5.289) (-3.822,5.331) (-3.884,5.373) (-3.947,5.415) (-4.009,5.456) (-4.072,5.497) (-4.135,5.537) (-4.199,5.577) (-4.262,5.617) (-4.326,5.656) (-4.391,5.695) (-4.455,5.734) (-4.520,5.772) (-4.585,5.809) (-4.650,5.846) (-4.715,5.883) (-4.781,5.920) (-4.846,5.956) (-4.912,5.991) (-4.979,6.026) (-5.045,6.061) (-5.112,6.095) (-5.179,6.129) (-5.246,6.162) (-5.313,6.196) (-5.381,6.228) (-5.449,6.260) (-5.517,6.292) (-5.585,6.323) (-5.653,6.354) (-5.722,6.385) (-5.790,6.415) (-5.859,6.444) (-5.928,6.474) (-5.998,6.502) (-6.067,6.531) (-6.137,6.558) (-6.206,6.586) (-6.276,6.613) (-6.347,6.639) (-6.417,6.666) (-6.487,6.691) (-6.558,6.717) (-6.629,6.741) (-6.700,6.766) (-6.771,6.790) (-6.842,6.813) (-6.913,6.836) (-6.985,6.859) (-7.057,6.881) (-7.128,6.902) (-7.200,6.924) (-7.272,6.944) (-7.345,6.965) (-7.417,6.984) (-7.489,7.004) (-7.562,7.023) (-7.635,7.041) (-7.707,7.059) (-7.780,7.077) (-7.853,7.094) (-7.926,7.111) (-8.000,7.127) (-8.073,7.142) (-8.147,7.158) (-8.220,7.172) (-8.294,7.187) (-8.367,7.200) (-8.441,7.214) (-8.515,7.227) (-8.589,7.239) (-8.663,7.251) (-8.737,7.262) (-8.811,7.273) (-8.886,7.284) (-8.960,7.294) (-9.034,7.303) (-9.109,7.312) (-9.183,7.321) (-9.258,7.329) (-9.333,7.336) (-9.407,7.343) (-9.482,7.350) (-9.557,7.356) (-9.631,7.362) (-9.706,7.367) (-9.781,7.371) (-9.856,7.376) (-9.931,7.379) (-10.000,7.382)}
\LC{(-10.000,-8.148) (-9.943,-8.138) (-9.869,-8.127) (-9.795,-8.117) (-9.720,-8.105) (-9.646,-8.094) (-9.572,-8.082) (-9.498,-8.069) (-9.424,-8.056) (-9.351,-8.043) (-9.277,-8.029) (-9.203,-8.015) (-9.130,-8.000) (-9.056,-7.985) (-8.983,-7.970) (-8.910,-7.954) (-8.836,-7.938) (-8.763,-7.921) (-8.690,-7.904) (-8.617,-7.886) (-8.545,-7.868) (-8.472,-7.850) (-8.399,-7.831) (-8.327,-7.812) (-8.254,-7.792) (-8.182,-7.772) (-8.110,-7.752) (-8.038,-7.731) (-7.966,-7.710) (-7.894,-7.688) (-7.822,-7.666) (-7.751,-7.644) (-7.679,-7.621) (-7.608,-7.597) (-7.537,-7.574) (-7.466,-7.550) (-7.395,-7.525) (-7.324,-7.500) (-7.254,-7.475) (-7.183,-7.449) (-7.113,-7.423) (-7.043,-7.397) (-6.973,-7.370) (-6.903,-7.343) (-6.833,-7.315) (-6.763,-7.287) (-6.694,-7.259) (-6.625,-7.230) (-6.556,-7.201) (-6.487,-7.171) (-6.418,-7.141) (-6.349,-7.111) (-6.281,-7.080) (-6.213,-7.049) (-6.145,-7.018) (-6.077,-6.986) (-6.009,-6.954) (-5.941,-6.921) (-5.874,-6.888) (-5.807,-6.855) (-5.740,-6.821) (-5.673,-6.787) (-5.607,-6.752) (-5.540,-6.718) (-5.474,-6.682) (-5.408,-6.647) (-5.342,-6.611) (-5.277,-6.574) (-5.211,-6.538) (-5.146,-6.500) (-5.081,-6.463) (-5.016,-6.425) (-4.952,-6.387) (-4.887,-6.348) (-4.823,-6.309) (-4.760,-6.270) (-4.696,-6.230) (-4.632,-6.190) (-4.569,-6.150) (-4.506,-6.109) (-4.444,-6.068) (-4.381,-6.026) (-4.319,-5.985) (-4.257,-5.942) (-4.195,-5.900) (-4.134,-5.857) (-4.072,-5.814) (-4.011,-5.770) (-3.951,-5.726) (-3.890,-5.682) (-3.830,-5.637) (-3.770,-5.592) (-3.710,-5.547) (-3.651,-5.501) (-3.592,-5.455) (-3.533,-5.408) (-3.474,-5.362) (-3.416,-5.314) (-3.358,-5.267) (-3.300,-5.219) (-3.242,-5.171) (-3.185,-5.122) (-3.128,-5.074) (-3.072,-5.024) (-3.015,-4.975) (-2.959,-4.925) (-2.904,-4.875) (-2.848,-4.824) (-2.793,-4.773) (-2.739,-4.722) (-2.684,-4.670) (-2.630,-4.618) (-2.576,-4.566) (-2.523,-4.514) (-2.470,-4.461) (-2.417,-4.407) (-2.364,-4.354) (-2.312,-4.300) (-2.260,-4.246) (-2.209,-4.191) (-2.158,-4.136) (-2.107,-4.081) (-2.057,-4.025) (-2.007,-3.969) (-1.957,-3.913) (-1.908,-3.857) (-1.859,-3.800) (-1.811,-3.742) (-1.763,-3.685) (-1.715,-3.627) (-1.668,-3.569) (-1.621,-3.510) (-1.574,-3.451) (-1.528,-3.392) (-1.482,-3.333) (-1.437,-3.273) (-1.392,-3.213) (-1.348,-3.153) (-1.304,-3.092) (-1.260,-3.031) (-1.217,-2.969) (-1.174,-2.908) (-1.146,-2.866)}
\LC{(-5.708,3.384) (-5.781,3.371) (-5.855,3.357) (-5.928,3.341) (-6.001,3.324) (-6.074,3.306) (-6.147,3.287) (-6.219,3.267) (-6.291,3.245) (-6.362,3.223) (-6.433,3.199) (-6.504,3.174) (-6.574,3.148) (-6.644,3.120) (-6.714,3.092) (-6.782,3.062) (-6.851,3.031) (-6.918,2.999) (-6.986,2.966) (-7.052,2.931) (-7.118,2.895) (-7.183,2.858) (-7.248,2.820) (-7.312,2.781) (-7.375,2.740) (-7.437,2.699) (-7.499,2.656) (-7.560,2.612) (-7.620,2.567) (-7.679,2.521) (-7.737,2.473) (-7.794,2.424) (-7.850,2.375) (-7.905,2.324) (-7.959,2.272) (-8.012,2.219) (-8.064,2.164) (-8.114,2.109) (-8.164,2.053) (-8.212,1.995) (-8.259,1.937) (-8.305,1.877) (-8.349,1.817) (-8.392,1.756) (-8.434,1.693) (-8.474,1.630) (-8.513,1.566) (-8.551,1.501) (-8.586,1.435) (-8.621,1.368) (-8.653,1.301) (-8.685,1.233) (-8.714,1.164) (-8.742,1.094) (-8.768,1.024) (-8.793,0.953) (-8.815,0.881) (-8.836,0.809) (-8.856,0.737) (-8.873,0.664) (-8.889,0.591) (-8.902,0.517) (-8.914,0.443) (-8.925,0.368) (-8.933,0.294) (-8.939,0.219) (-8.944,0.144) (-8.947,0.069) (-8.947,-0.006) (-8.946,-0.081) (-8.943,-0.156) (-8.938,-0.230) (-8.932,-0.305) (-8.923,-0.380) (-8.913,-0.454) (-8.901,-0.528) (-8.886,-0.601) (-8.870,-0.675) (-8.853,-0.748) (-8.833,-0.820) (-8.812,-0.892) (-8.789,-0.963) (-8.764,-1.034) (-8.738,-1.104) (-8.710,-1.174) (-8.680,-1.243) (-8.649,-1.311) (-8.616,-1.378) (-8.581,-1.445) (-8.545,-1.511) (-8.507,-1.575) (-8.468,-1.640) (-8.428,-1.703) (-8.386,-1.765) (-8.343,-1.826) (-8.298,-1.886) (-8.252,-1.946) (-8.205,-2.004) (-8.157,-2.061) (-8.107,-2.117) (-8.056,-2.173) (-8.004,-2.227) (-7.951,-2.280) (-7.897,-2.331) (-7.842,-2.382) (-7.785,-2.432) (-7.728,-2.480) (-7.670,-2.528) (-7.611,-2.574) (-7.551,-2.619) (-7.490,-2.662) (-7.428,-2.705) (-7.366,-2.747) (-7.302,-2.787) (-7.238,-2.826) (-7.174,-2.864) (-7.108,-2.901) (-7.042,-2.936) (-6.976,-2.971) (-6.908,-3.004) (-6.841,-3.036) (-6.772,-3.066) (-6.703,-3.096) (-6.634,-3.124) (-6.564,-3.152) (-6.494,-3.178) (-6.423,-3.202) (-6.352,-3.226) (-6.280,-3.249) (-6.208,-3.270) (-6.136,-3.290) (-6.063,-3.309) (-5.990,-3.327) (-5.917,-3.343) (-5.844,-3.359) (-5.770,-3.373) (-5.697,-3.386)}
\LC{(10.000,-7.567) (9.949,-7.562) (9.875,-7.557) (9.800,-7.551) (9.725,-7.545) (9.650,-7.539) (9.576,-7.532) (9.501,-7.524) (9.427,-7.516) (9.352,-7.507) (9.278,-7.498) (9.203,-7.489) (9.129,-7.479) (9.055,-7.468) (8.980,-7.457) (8.906,-7.446) (8.832,-7.434) (8.758,-7.422) (8.684,-7.409) (8.610,-7.396) (8.537,-7.382) (8.463,-7.368) (8.390,-7.354) (8.316,-7.339) (8.243,-7.323) (8.169,-7.307) (8.096,-7.291) (8.023,-7.274) (7.950,-7.256) (7.877,-7.239) (7.805,-7.220) (7.732,-7.202) (7.659,-7.183) (7.587,-7.163) (7.515,-7.143) (7.443,-7.123) (7.371,-7.102) (7.299,-7.080) (7.227,-7.059) (7.155,-7.036) (7.084,-7.014) (7.012,-6.991) (6.941,-6.967) (6.870,-6.943) (6.799,-6.919) (6.728,-6.894) (6.658,-6.869) (6.587,-6.843) (6.517,-6.817) (6.447,-6.790) (6.377,-6.763) (6.307,-6.736) (6.237,-6.708) (6.168,-6.680) (6.099,-6.651) (6.030,-6.622) (5.961,-6.592) (5.892,-6.562) (5.823,-6.532) (5.755,-6.501) (5.687,-6.470) (5.619,-6.438) (5.551,-6.406) (5.483,-6.374) (5.416,-6.341) (5.349,-6.308) (5.282,-6.274) (5.215,-6.240) (5.148,-6.206) (5.082,-6.171) (5.016,-6.135) (4.950,-6.100) (4.884,-6.064) (4.818,-6.027) (4.753,-5.990) (4.688,-5.953) (4.623,-5.915) (4.559,-5.877) (4.494,-5.838) (4.430,-5.799) (4.366,-5.760) (4.303,-5.720) (4.239,-5.680) (4.176,-5.640) (4.113,-5.599) (4.051,-5.558) (3.988,-5.516) (3.947,-5.488)}
\LC{(1.522,2.462) (1.580,2.510) (1.639,2.556) (1.698,2.602) (1.758,2.647) (1.818,2.692) (1.879,2.735) (1.941,2.778) (2.003,2.821) (2.065,2.862) (2.128,2.903) (2.192,2.943) (2.256,2.982) (2.320,3.020) (2.385,3.058) (2.451,3.094) (2.516,3.130) (2.583,3.166) (2.649,3.200) (2.716,3.234) (2.784,3.267) (2.851,3.299) (2.920,3.330) (2.988,3.361) (3.057,3.390) (3.126,3.419) (3.196,3.447) (3.266,3.474) (3.336,3.501) (3.406,3.526) (3.477,3.551) (3.548,3.575) (3.620,3.598) (3.691,3.620) (3.763,3.641) (3.836,3.662) (3.908,3.681) (3.981,3.700) (4.053,3.718) (4.126,3.735) (4.200,3.751) (4.273,3.766) (4.347,3.780) (4.421,3.793) (4.495,3.806) (4.569,3.818) (4.643,3.828) (4.717,3.838) (4.792,3.847) (4.866,3.855) (4.941,3.862) (5.016,3.868) (5.091,3.873) (5.165,3.877) (5.240,3.880) (5.315,3.883) (5.390,3.884) (5.465,3.884) (5.540,3.884) (5.615,3.882) (5.690,3.880) (5.765,3.876) (5.840,3.872) (5.915,3.866) (5.990,3.860) (6.064,3.852) (6.139,3.844) (6.213,3.835) (6.287,3.824) (6.362,3.813) (6.435,3.800) (6.509,3.787) (6.583,3.772) (6.656,3.757) (6.729,3.740) (6.802,3.722) (6.875,3.704) (6.947,3.684) (7.019,3.663) (7.091,3.642) (7.163,3.619) (7.234,3.595) (7.304,3.570) (7.375,3.544) (7.445,3.517) (7.514,3.489) (7.583,3.460) (7.652,3.430) (7.720,3.399) (7.788,3.366) (7.855,3.333) (7.922,3.299) (7.988,3.263) (8.053,3.226) (8.118,3.189) (8.182,3.150) (8.246,3.110) (8.309,3.070) (8.371,3.028) (8.433,2.985) (8.494,2.941) (8.554,2.896) (8.613,2.850) (8.671,2.803) (8.729,2.755) (8.786,2.706) (8.841,2.656) (8.896,2.605) (8.950,2.553) (9.003,2.500) (9.055,2.446) (9.106,2.391) (9.156,2.335) (9.205,2.278) (9.253,2.220) (9.300,2.161) (9.345,2.102) (9.389,2.041) (9.433,1.980) (9.474,1.918) (9.515,1.855) (9.554,1.791) (9.592,1.726) (9.629,1.661) (9.665,1.595) (9.699,1.528) (9.731,1.460) (9.762,1.392) (9.792,1.323) (9.820,1.254) (9.847,1.184) (9.873,1.113) (9.896,1.042) (9.919,0.970) (9.939,0.898) (9.958,0.826) (9.976,0.753) (9.992,0.679) (10.000,0.630)}
\LC{(-3.266,2.966) (-3.339,2.982) (-3.413,2.996) (-3.486,3.010) (-3.560,3.022) (-3.634,3.034) (-3.709,3.044) (-3.783,3.053) (-3.858,3.061) (-3.932,3.067) (-4.007,3.073) (-4.082,3.077) (-4.157,3.080) (-4.232,3.082) (-4.307,3.083) (-4.382,3.082) (-4.457,3.081) (-4.532,3.078) (-4.607,3.074) (-4.682,3.068) (-4.756,3.062) (-4.831,3.054) (-4.905,3.045) (-4.980,3.034) (-5.054,3.023) (-5.128,3.010) (-5.201,2.995) (-5.275,2.980) (-5.348,2.963) (-5.420,2.945) (-5.493,2.926) (-5.565,2.905) (-5.637,2.883) (-5.708,2.859) (-5.779,2.835) (-5.849,2.809) (-5.919,2.781) (-5.988,2.753) (-6.057,2.723) (-6.125,2.691) (-6.193,2.659) (-6.259,2.625) (-6.326,2.589) (-6.391,2.552) (-6.456,2.514) (-6.519,2.475) (-6.582,2.434) (-6.645,2.392) (-6.706,2.349) (-6.766,2.304) (-6.825,2.258) (-6.884,2.211) (-6.941,2.163) (-6.997,2.113) (-7.052,2.062) (-7.105,2.009) (-7.158,1.956) (-7.209,1.901) (-7.259,1.845) (-7.308,1.788) (-7.355,1.730) (-7.401,1.670) (-7.445,1.610) (-7.487,1.548) (-7.529,1.485) (-7.568,1.422) (-7.606,1.357) (-7.642,1.291) (-7.677,1.225) (-7.710,1.157) (-7.741,1.089) (-7.770,1.020) (-7.797,0.950) (-7.822,0.879) (-7.846,0.808) (-7.867,0.736) (-7.887,0.664) (-7.904,0.591) (-7.920,0.518) (-7.933,0.444) (-7.945,0.370) (-7.954,0.295) (-7.961,0.221) (-7.967,0.146) (-7.970,0.071) (-7.971,-0.004) (-7.969,-0.079) (-7.966,-0.154) (-7.961,-0.229) (-7.953,-0.304) (-7.944,-0.378) (-7.932,-0.452) (-7.918,-0.526) (-7.903,-0.599) (-7.885,-0.672) (-7.865,-0.744) (-7.843,-0.816) (-7.820,-0.887) (-7.794,-0.958) (-7.767,-1.027) (-7.737,-1.097) (-7.706,-1.165) (-7.673,-1.232) (-7.638,-1.299) (-7.602,-1.364) (-7.564,-1.429) (-7.524,-1.492) (-7.483,-1.555) (-7.440,-1.616) (-7.395,-1.677) (-7.350,-1.736) (-7.302,-1.794) (-7.254,-1.851) (-7.204,-1.907) (-7.152,-1.962) (-7.099,-2.015) (-7.046,-2.068) (-6.991,-2.118) (-6.934,-2.168) (-6.877,-2.217) (-6.819,-2.264) (-6.759,-2.309) (-6.699,-2.354) (-6.638,-2.397) (-6.575,-2.439) (-6.512,-2.480) (-6.448,-2.519) (-6.384,-2.557) (-6.318,-2.593) (-6.252,-2.628) (-6.185,-2.662) (-6.118,-2.695) (-6.049,-2.726) (-5.981,-2.756) (-5.911,-2.784) (-5.841,-2.812) (-5.771,-2.837) (-5.700,-2.862) (-5.629,-2.885) (-5.557,-2.907) (-5.485,-2.928) (-5.412,-2.947) (-5.340,-2.965) (-5.266,-2.982) (-5.193,-2.997) (-5.119,-3.011) (-5.045,-3.024) (-4.971,-3.036) (-4.897,-3.046) (-4.823,-3.055) (-4.748,-3.063) (-4.673,-3.069) (-4.598,-3.074) (-4.524,-3.078) (-4.449,-3.081) (-4.374,-3.083) (-4.299,-3.083) (-4.224,-3.082) (-4.149,-3.080) (-4.074,-3.077) (-3.999,-3.072) (-3.924,-3.067) (-3.849,-3.060) (-3.775,-3.052) (-3.700,-3.043) (-3.626,-3.032) (-3.552,-3.021) (-3.478,-3.009) (-3.404,-2.995) (-3.331,-2.980) (-3.257,-2.964)}
\LC{(2.562,5.059) (2.612,5.115) (2.662,5.171) (2.713,5.226) (2.763,5.282) (2.814,5.337) (2.865,5.391) (2.917,5.446) (2.969,5.500) (3.021,5.554) (3.073,5.608) (3.125,5.662) (3.178,5.715) (3.231,5.768) (3.285,5.820) (3.338,5.873) (3.392,5.925) (3.446,5.977) (3.501,6.029) (3.555,6.080) (3.610,6.131) (3.665,6.182) (3.721,6.233) (3.776,6.283) (3.832,6.333) (3.888,6.383) (3.944,6.432) (4.001,6.482) (4.058,6.531) (4.115,6.579) (4.172,6.628) (4.229,6.676) (4.287,6.724) (4.345,6.772) (4.403,6.819) (4.462,6.866) (4.520,6.913) (4.579,6.960) (4.638,7.006) (4.697,7.052) (4.757,7.098) (4.816,7.143) (4.876,7.188) (4.936,7.233) (4.996,7.278) (5.057,7.322) (5.118,7.366) (5.178,7.410) (5.239,7.454) (5.301,7.497) (5.362,7.540) (5.424,7.583) (5.486,7.625) (5.548,7.667) (5.610,7.709) (5.672,7.751) (5.735,7.792) (5.798,7.833) (5.861,7.874) (5.924,7.914) (5.987,7.954) (6.051,7.994) (6.114,8.034) (6.178,8.073) (6.242,8.112) (6.307,8.151) (6.371,8.190) (6.435,8.228) (6.500,8.266) (6.565,8.303) (6.630,8.340) (6.695,8.378) (6.761,8.414) (6.826,8.451) (6.892,8.487) (6.958,8.523) (7.024,8.558) (7.090,8.594) (7.156,8.629) (7.223,8.663) (7.290,8.698) (7.356,8.732) (7.423,8.766) (7.490,8.799) (7.558,8.832) (7.625,8.865) (7.693,8.898) (7.760,8.930) (7.828,8.962) (7.896,8.994) (7.964,9.026) (8.032,9.057) (8.101,9.087) (8.169,9.118) (8.238,9.148) (8.307,9.178) (8.375,9.208) (8.444,9.237) (8.514,9.266) (8.583,9.295) (8.652,9.323) (8.722,9.352) (8.791,9.379) (8.861,9.407) (8.931,9.434) (9.001,9.461) (9.071,9.488) (9.141,9.514) (9.212,9.540) (9.282,9.566) (9.353,9.591) (9.424,9.616) (9.494,9.641) (9.565,9.665) (9.636,9.689) (9.707,9.713) (9.779,9.737) (9.850,9.760) (9.921,9.783) (9.993,9.805) (10.000,9.813)}
\LC{(1.896,1.632) (1.970,1.642) (2.045,1.649) (2.120,1.654) (2.195,1.657) (2.270,1.658) (2.345,1.656) (2.420,1.653) (2.494,1.647) (2.569,1.639) (2.643,1.628) (2.717,1.615) (2.790,1.600) (2.863,1.582) (2.936,1.562) (3.007,1.539) (3.078,1.514) (3.148,1.487) (3.216,1.457) (3.284,1.425) (3.350,1.390) (3.415,1.353) (3.479,1.313) (3.541,1.270) (3.601,1.226) (3.660,1.179) (3.716,1.129) (3.770,1.077) (3.822,1.023) (3.871,0.967) (3.918,0.908) (3.962,0.848) (4.004,0.785) (4.042,0.720) (4.077,0.654) (4.109,0.586) (4.137,0.517) (4.162,0.446) (4.184,0.374) (4.201,0.301) (4.215,0.228) (4.225,0.153) (4.231,0.079) (4.233,0.004) (4.231,-0.071) (4.226,-0.146) (4.216,-0.220) (4.203,-0.294) (4.185,-0.367) (4.165,-0.439) (4.140,-0.510) (4.112,-0.580) (4.080,-0.648) (4.045,-0.714) (4.008,-0.779) (3.967,-0.842) (3.923,-0.902) (3.876,-0.961) (3.827,-1.018) (3.775,-1.072) (3.721,-1.124) (3.665,-1.174) (3.607,-1.221) (3.547,-1.266) (3.485,-1.309) (3.422,-1.349) (3.357,-1.386) (3.291,-1.421) (3.223,-1.454) (3.154,-1.484) (3.085,-1.512) (3.014,-1.537) (2.943,-1.560) (2.870,-1.580) (2.798,-1.598) (2.724,-1.614) (2.650,-1.627) (2.576,-1.638) (2.502,-1.646) (2.427,-1.652) (2.352,-1.656) (2.277,-1.658) (2.202,-1.657) (2.127,-1.654) (2.052,-1.650) (1.978,-1.642) (1.903,-1.633) (1.878,-1.630)}
\LC{(2.741,-10.000) (2.712,-9.933) (2.682,-9.864) (2.653,-9.795) (2.624,-9.726) (2.595,-9.657) (2.566,-9.588) (2.537,-9.519) (2.508,-9.449) (2.479,-9.380) (2.451,-9.311) (2.422,-9.241) (2.394,-9.172) (2.366,-9.102) (2.338,-9.033) (2.310,-8.963) (2.282,-8.894) (2.254,-8.824) (2.226,-8.754) (2.199,-8.685) (2.171,-8.615) (2.144,-8.545) (2.117,-8.475) (2.090,-8.405) (2.063,-8.335) (2.036,-8.265) (2.009,-8.195) (1.983,-8.125) (1.956,-8.055) (1.930,-7.985) (1.903,-7.914) (1.877,-7.844) (1.851,-7.774) (1.825,-7.703) (1.799,-7.633) (1.774,-7.563) (1.748,-7.492) (1.723,-7.421) (1.697,-7.351) (1.672,-7.280) (1.647,-7.210) (1.622,-7.139) (1.597,-7.068) (1.572,-6.997) (1.548,-6.926) (1.523,-6.856) (1.499,-6.785) (1.475,-6.714) (1.451,-6.643) (1.427,-6.571) (1.403,-6.500) (1.379,-6.429) (1.356,-6.358) (1.332,-6.287) (1.309,-6.215) (1.286,-6.144) (1.263,-6.073) (1.240,-6.001) (1.217,-5.930) (1.194,-5.858) (1.172,-5.787) (1.150,-5.715) (1.127,-5.644) (1.105,-5.572) (1.083,-5.500) (1.062,-5.428) (1.040,-5.357) (1.018,-5.285) (0.997,-5.213) (0.976,-5.141) (0.955,-5.069) (0.934,-4.997) (0.913,-4.925) (0.893,-4.853) (0.872,-4.781) (0.852,-4.708) (0.832,-4.636) (0.812,-4.564) (0.792,-4.492) (0.772,-4.419) (0.753,-4.347) (0.746,-4.323)}
\LC{(-3.495,5.655) (-3.552,5.704) (-3.609,5.752) (-3.667,5.801) (-3.725,5.849) (-3.783,5.896) (-3.841,5.944) (-3.899,5.991) (-3.958,6.037) (-4.017,6.084) (-4.076,6.130) (-4.135,6.175) (-4.195,6.221) (-4.255,6.266) (-4.315,6.311) (-4.376,6.355) (-4.436,6.399) (-4.497,6.443) (-4.558,6.487) (-4.619,6.530) (-4.681,6.573) (-4.743,6.615) (-4.805,6.658) (-4.867,6.699) (-4.929,6.741) (-4.992,6.782) (-5.055,6.823) (-5.118,6.864) (-5.181,6.904) (-5.245,6.944) (-5.308,6.984) (-5.372,7.023) (-5.436,7.062) (-5.500,7.101) (-5.565,7.139) (-5.630,7.177) (-5.694,7.215) (-5.759,7.252) (-5.825,7.289) (-5.890,7.326) (-5.956,7.362) (-6.022,7.398) (-6.087,7.434) (-6.154,7.469) (-6.220,7.504) (-6.286,7.539) (-6.353,7.573) (-6.420,7.607) (-6.487,7.641) (-6.554,7.674) (-6.622,7.707) (-6.689,7.740) (-6.757,7.772) (-6.825,7.804) (-6.893,7.836) (-6.961,7.867) (-7.029,7.898) (-7.098,7.928) (-7.166,7.959) (-7.235,7.988) (-7.304,8.018) (-7.373,8.047) (-7.442,8.076) (-7.512,8.105) (-7.581,8.133) (-7.651,8.160) (-7.721,8.188) (-7.791,8.215) (-7.861,8.242) (-7.931,8.268) (-8.001,8.294) (-8.072,8.320) (-8.142,8.345) (-8.213,8.370) (-8.284,8.395) (-8.355,8.419) (-8.426,8.443) (-8.497,8.466) (-8.569,8.489) (-8.640,8.512) (-8.712,8.535) (-8.783,8.557) (-8.855,8.579) (-8.927,8.600) (-8.999,8.621) (-9.071,8.642) (-9.143,8.662) (-9.216,8.682) (-9.288,8.701) (-9.360,8.720) (-9.433,8.739) (-9.506,8.758) (-9.579,8.776) (-9.651,8.793) (-9.724,8.811) (-9.797,8.828) (-9.871,8.844) (-9.944,8.860) (-10.000,8.876)}
\LC{(8.690,2.175) (8.738,2.118) (8.785,2.059) (8.831,2.000) (8.876,1.940) (8.919,1.878) (8.961,1.816) (9.002,1.753) (9.041,1.689) (9.079,1.625) (9.115,1.559) (9.150,1.493) (9.184,1.426) (9.216,1.358) (9.247,1.289) (9.276,1.220) (9.303,1.151) (9.329,1.080) (9.354,1.009) (9.376,0.938) (9.397,0.866) (9.417,0.793) (9.435,0.720) (9.451,0.647) (9.465,0.574) (9.478,0.500) (9.489,0.425) (9.498,0.351) (9.505,0.276) (9.511,0.202) (9.515,0.127) (9.517,0.052) (9.517,-0.023) (9.516,-0.098) (9.512,-0.173) (9.507,-0.248) (9.501,-0.323) (9.492,-0.397) (9.482,-0.472) (9.470,-0.546) (9.456,-0.619) (9.441,-0.693) (9.424,-0.766) (9.405,-0.838) (9.385,-0.910) (9.363,-0.982) (9.339,-1.053) (9.313,-1.124) (9.287,-1.194) (9.258,-1.263) (9.228,-1.332) (9.196,-1.400) (9.163,-1.467) (9.129,-1.534) (9.093,-1.600) (9.056,-1.665) (9.017,-1.729) (8.977,-1.792) (8.935,-1.855) (8.892,-1.916) (8.848,-1.977) (8.803,-2.037) (8.756,-2.096) (8.709,-2.154) (8.660,-2.210) (8.610,-2.266) (8.559,-2.321) (8.507,-2.375) (8.453,-2.428) (8.399,-2.480) (8.344,-2.530) (8.288,-2.580) (8.231,-2.629) (8.173,-2.676) (8.114,-2.723) (8.054,-2.768) (7.993,-2.812) (7.932,-2.856) (7.870,-2.898) (7.807,-2.939) (7.744,-2.978) (7.679,-3.017) (7.614,-3.055) (7.549,-3.091) (7.483,-3.127) (7.416,-3.161) (7.349,-3.194) (7.281,-3.226) (7.213,-3.257) (7.144,-3.286) (7.074,-3.315) (7.005,-3.343) (6.934,-3.369) (6.864,-3.394) (6.793,-3.418) (6.721,-3.441) (6.650,-3.463) (6.578,-3.484) (6.505,-3.504) (6.433,-3.523) (6.360,-3.540) (6.287,-3.557) (6.213,-3.572) (6.140,-3.586) (6.066,-3.600) (5.992,-3.612) (5.918,-3.623) (5.843,-3.633) (5.769,-3.642) (5.694,-3.650) (5.620,-3.657) (5.545,-3.662) (5.470,-3.667) (5.395,-3.671) (5.320,-3.673) (5.245,-3.675) (5.170,-3.676) (5.095,-3.675) (5.020,-3.674) (4.945,-3.672) (4.870,-3.668) (4.795,-3.664) (4.721,-3.658) (4.646,-3.652) (4.571,-3.645) (4.497,-3.636) (4.422,-3.627) (4.348,-3.617) (4.274,-3.605) (4.200,-3.593) (4.126,-3.580) (4.052,-3.566) (3.979,-3.551) (3.906,-3.535) (3.832,-3.518) (3.760,-3.500) (3.687,-3.482) (3.615,-3.462) (3.542,-3.441) (3.471,-3.420) (3.399,-3.397) (3.328,-3.374) (3.257,-3.350) (3.186,-3.325) (3.116,-3.299) (3.046,-3.272) (2.976,-3.244) (2.907,-3.216) (2.838,-3.186) (2.769,-3.156) (2.701,-3.125) (2.633,-3.093) (2.566,-3.060) (2.499,-3.026) (2.432,-2.992) (2.366,-2.957) (2.300,-2.920) (2.235,-2.883) (2.170,-2.846) (2.106,-2.807) (2.042,-2.768) (1.979,-2.727) (1.916,-2.686) (1.854,-2.644) (1.792,-2.602) (1.731,-2.558) (1.711,-2.544)}
\LC{(9.553,-10.000) (9.483,-9.975) (9.413,-9.948) (9.343,-9.920) (9.273,-9.893) (9.203,-9.865) (9.134,-9.837) (9.064,-9.808) (8.995,-9.780) (8.926,-9.751) (8.857,-9.722) (8.788,-9.692) (8.719,-9.662) (8.651,-9.632) (8.582,-9.602) (8.514,-9.571) (8.445,-9.540) (8.377,-9.509) (8.309,-9.477) (8.241,-9.445) (8.173,-9.413) (8.106,-9.381) (8.038,-9.348) (7.971,-9.315) (7.904,-9.282) (7.836,-9.248) (7.769,-9.215) (7.703,-9.181) (7.636,-9.146) (7.569,-9.112) (7.503,-9.077) (7.437,-9.042) (7.371,-9.006) (7.305,-8.970) (7.239,-8.934) (7.173,-8.898) (7.108,-8.861) (7.043,-8.825) (6.977,-8.788) (6.912,-8.750) (6.848,-8.712) (6.783,-8.675) (6.718,-8.636) (6.654,-8.598) (6.590,-8.559) (6.526,-8.520) (6.462,-8.481) (6.398,-8.441) (6.334,-8.401) (6.271,-8.361) (6.208,-8.321) (6.145,-8.280) (6.082,-8.239) (6.019,-8.198) (5.957,-8.157) (5.894,-8.115) (5.832,-8.073) (5.770,-8.031) (5.709,-7.988) (5.647,-7.946) (5.585,-7.903) (5.524,-7.859) (5.463,-7.816) (5.402,-7.772) (5.342,-7.728) (5.281,-7.684) (5.221,-7.639) (5.161,-7.594) (5.101,-7.549) (5.041,-7.504) (4.982,-7.458) (4.922,-7.412) (4.863,-7.366) (4.804,-7.319) (4.746,-7.273) (4.687,-7.226) (4.629,-7.179) (4.571,-7.131) (4.513,-7.083) (4.455,-7.035) (4.398,-6.987) (4.341,-6.939) (4.284,-6.890) (4.227,-6.841) (4.170,-6.792) (4.114,-6.742) (4.058,-6.692) (4.002,-6.642) (3.946,-6.592) (3.891,-6.542) (3.836,-6.491) (3.781,-6.440) (3.726,-6.389) (3.671,-6.337) (3.617,-6.285) (3.563,-6.233) (3.509,-6.181) (3.456,-6.129) (3.402,-6.076) (3.349,-6.023) (3.297,-5.970) (3.244,-5.916) (3.192,-5.862) (3.140,-5.808) (3.088,-5.754) (3.036,-5.699) (2.985,-5.645) (2.934,-5.590) (2.883,-5.535) (2.833,-5.479) (2.783,-5.423) (2.733,-5.367) (2.683,-5.311) (2.634,-5.255) (2.585,-5.198) (2.536,-5.141) (2.487,-5.084) (2.439,-5.026) (2.391,-4.969) (2.343,-4.911) (2.296,-4.853) (2.249,-4.794) (2.202,-4.736) (2.156,-4.677) (2.110,-4.618) (2.064,-4.558) (2.018,-4.499) (1.973,-4.439) (1.928,-4.379) (1.883,-4.319) (1.839,-4.258) (1.795,-4.197) (1.751,-4.136) (1.708,-4.075) (1.665,-4.014) (1.622,-3.952) (1.580,-3.890) (1.538,-3.828) (1.497,-3.766) (1.483,-3.745)}
\LC{(-10.000,-5.598) (-9.942,-5.614) (-9.868,-5.630) (-9.795,-5.644) (-9.721,-5.658) (-9.647,-5.672) (-9.574,-5.685) (-9.499,-5.697) (-9.425,-5.708) (-9.351,-5.719) (-9.277,-5.729) (-9.202,-5.738) (-9.128,-5.747) (-9.053,-5.755) (-8.979,-5.762) (-8.904,-5.769) (-8.829,-5.775) (-8.754,-5.780) (-8.680,-5.785) (-8.605,-5.789) (-8.530,-5.792) (-8.455,-5.795) (-8.380,-5.797) (-8.305,-5.799) (-8.230,-5.799) (-8.155,-5.799) (-8.080,-5.799) (-8.005,-5.798) (-7.930,-5.796) (-7.855,-5.794) (-7.780,-5.790) (-7.705,-5.787) (-7.630,-5.782) (-7.555,-5.777) (-7.481,-5.772) (-7.406,-5.766) (-7.331,-5.759) (-7.257,-5.751) (-7.182,-5.743) (-7.107,-5.734) (-7.033,-5.725) (-6.959,-5.715) (-6.884,-5.705) (-6.810,-5.693) (-6.736,-5.682) (-6.662,-5.669) (-6.588,-5.656) (-6.515,-5.643) (-6.441,-5.628) (-6.367,-5.614) (-6.294,-5.598) (-6.221,-5.582) (-6.148,-5.566) (-6.075,-5.548) (-6.002,-5.531) (-5.929,-5.512) (-5.857,-5.493) (-5.784,-5.474) (-5.712,-5.454) (-5.640,-5.433) (-5.568,-5.412) (-5.496,-5.390) (-5.425,-5.367) (-5.353,-5.344) (-5.282,-5.321) (-5.211,-5.297) (-5.140,-5.272) (-5.069,-5.247) (-4.999,-5.221) (-4.929,-5.195) (-4.859,-5.168) (-4.789,-5.140) (-4.720,-5.112) (-4.650,-5.084) (-4.581,-5.054) (-4.512,-5.025) (-4.444,-4.994) (-4.375,-4.964) (-4.307,-4.932) (-4.239,-4.900) (-4.172,-4.868) (-4.104,-4.835) (-4.037,-4.802) (-3.970,-4.768) (-3.904,-4.733) (-3.837,-4.698) (-3.771,-4.662) (-3.706,-4.626) (-3.640,-4.589) (-3.575,-4.552) (-3.510,-4.514) (-3.446,-4.476) (-3.382,-4.437) (-3.318,-4.398) (-3.254,-4.358) (-3.191,-4.318) (-3.128,-4.277) (-3.065,-4.236) (-3.003,-4.194) (-2.941,-4.152) (-2.880,-4.109) (-2.819,-4.065) (-2.758,-4.022) (-2.697,-3.977) (-2.637,-3.932) (-2.577,-3.887) (-2.518,-3.841) (-2.459,-3.795)}
\LC{(10.000,-4.870) (9.938,-4.895) (9.867,-4.920) (9.796,-4.943) (9.724,-4.966) (9.653,-4.988) (9.581,-5.010) (9.509,-5.030) (9.436,-5.050) (9.364,-5.069) (9.291,-5.087) (9.218,-5.105) (9.145,-5.122) (9.072,-5.138) (8.998,-5.153) (8.925,-5.168) (8.851,-5.181) (8.777,-5.194) (8.703,-5.207) (8.629,-5.218) (8.555,-5.229) (8.480,-5.239) (8.406,-5.248) (8.332,-5.257) (8.257,-5.265) (8.182,-5.272) (8.108,-5.278) (8.033,-5.284) (7.958,-5.289) (7.883,-5.293) (7.808,-5.297) (7.733,-5.300) (7.658,-5.302) (7.583,-5.303) (7.508,-5.304) (7.433,-5.304) (7.358,-5.303) (7.283,-5.302) (7.208,-5.299) (7.133,-5.297) (7.058,-5.293) (6.984,-5.289) (6.909,-5.284) (6.834,-5.278) (6.759,-5.272) (6.684,-5.265) (6.610,-5.257) (6.535,-5.249) (6.461,-5.240) (6.387,-5.230) (6.312,-5.220) (6.238,-5.209) (6.164,-5.197) (6.090,-5.185) (6.016,-5.172) (5.942,-5.158) (5.869,-5.144) (5.795,-5.129) (5.722,-5.113) (5.649,-5.097) (5.576,-5.080) (5.503,-5.062) (5.430,-5.044) (5.357,-5.025) (5.285,-5.006) (5.213,-4.986) (5.141,-4.965) (5.069,-4.944) (4.997,-4.922) (4.926,-4.899) (4.854,-4.876) (4.783,-4.852) (4.712,-4.827) (4.642,-4.802) (4.571,-4.776) (4.501,-4.750) (4.431,-4.723)}
\LC{(-0.451,2.963) (-0.469,3.036) (-0.487,3.109) (-0.505,3.181) (-0.523,3.254) (-0.541,3.327) (-0.560,3.400) (-0.579,3.472) (-0.598,3.545) (-0.617,3.617) (-0.636,3.690) (-0.656,3.762) (-0.676,3.834) (-0.696,3.907) (-0.716,3.979) (-0.736,4.051) (-0.757,4.123) (-0.777,4.195) (-0.798,4.267) (-0.819,4.339) (-0.840,4.411) (-0.862,4.483) (-0.883,4.555) (-0.905,4.627) (-0.927,4.699) (-0.949,4.770) (-0.971,4.842) (-0.993,4.914) (-1.016,4.985) (-1.039,5.057) (-1.062,5.128) (-1.085,5.199) (-1.108,5.271) (-1.131,5.342) (-1.155,5.413) (-1.178,5.484) (-1.202,5.555) (-1.226,5.627) (-1.250,5.698) (-1.274,5.769) (-1.299,5.839) (-1.323,5.910) (-1.348,5.981) (-1.373,6.052) (-1.398,6.123) (-1.423,6.193) (-1.448,6.264) (-1.474,6.334) (-1.499,6.405) (-1.525,6.475) (-1.551,6.546) (-1.577,6.616) (-1.603,6.686) (-1.629,6.757) (-1.656,6.827) (-1.682,6.897) (-1.709,6.967) (-1.736,7.037) (-1.763,7.107) (-1.790,7.177) (-1.817,7.247) (-1.845,7.317) (-1.872,7.387) (-1.900,7.456) (-1.928,7.526) (-1.955,7.596) (-1.983,7.665) (-2.012,7.735) (-2.040,7.804) (-2.068,7.873) (-2.097,7.943) (-2.125,8.012) (-2.154,8.081) (-2.183,8.151) (-2.212,8.220) (-2.241,8.289) (-2.271,8.358) (-2.300,8.427) (-2.330,8.496) (-2.359,8.565) (-2.389,8.634) (-2.419,8.702) (-2.449,8.771) (-2.479,8.840) (-2.509,8.908) (-2.540,8.977) (-2.570,9.046) (-2.601,9.114) (-2.631,9.182) (-2.662,9.251) (-2.693,9.319) (-2.724,9.387) (-2.755,9.456) (-2.787,9.524) (-2.818,9.592) (-2.850,9.660) (-2.881,9.728) (-2.913,9.796) (-2.945,9.864) (-2.977,9.932) (-3.009,9.999) (-3.020,10.000)}
\LC{(-10.000,-9.565) (-9.952,-9.544) (-9.880,-9.523) (-9.808,-9.502) (-9.736,-9.480) (-9.664,-9.458) (-9.593,-9.436) (-9.521,-9.413) (-9.450,-9.390) (-9.379,-9.367) (-9.308,-9.343) (-9.237,-9.319) (-9.166,-9.295) (-9.095,-9.270) (-9.024,-9.245) (-8.953,-9.220) (-8.883,-9.194) (-8.813,-9.168) (-8.742,-9.142) (-8.672,-9.115) (-8.602,-9.088) (-8.532,-9.061) (-8.463,-9.034) (-8.393,-9.006) (-8.323,-8.978) (-8.254,-8.949) (-8.185,-8.920) (-8.116,-8.891) (-8.047,-8.862) (-7.978,-8.832) (-7.909,-8.802) (-7.841,-8.772) (-7.772,-8.741) (-7.704,-8.710) (-7.636,-8.679) (-7.568,-8.647) (-7.500,-8.615) (-7.432,-8.583) (-7.364,-8.550) (-7.297,-8.517) (-7.230,-8.484) (-7.163,-8.451) (-7.096,-8.417) (-7.029,-8.383) (-6.962,-8.349) (-6.896,-8.314) (-6.829,-8.279) (-6.763,-8.244) (-6.697,-8.208) (-6.631,-8.172) (-6.566,-8.136) (-6.500,-8.100) (-6.435,-8.063) (-6.370,-8.026) (-6.305,-7.988) (-6.240,-7.951) (-6.175,-7.913) (-6.111,-7.874) (-6.046,-7.836) (-5.982,-7.797) (-5.918,-7.758) (-5.854,-7.718) (-5.791,-7.678) (-5.727,-7.638) (-5.664,-7.598) (-5.601,-7.557) (-5.538,-7.516) (-5.476,-7.475) (-5.413,-7.434) (-5.351,-7.392) (-5.289,-7.350) (-5.227,-7.307) (-5.165,-7.265) (-5.104,-7.222) (-5.043,-7.178) (-4.982,-7.135) (-4.921,-7.091) (-4.860,-7.047) (-4.800,-7.002) (-4.740,-6.958) (-4.680,-6.913) (-4.620,-6.867) (-4.560,-6.822) (-4.501,-6.776) (-4.442,-6.730) (-4.383,-6.683) (-4.324,-6.637) (-4.266,-6.590) (-4.207,-6.542) (-4.149,-6.495) (-4.092,-6.447) (-4.034,-6.399) (-3.977,-6.350) (-3.920,-6.302) (-3.863,-6.253) (-3.806,-6.204) (-3.750,-6.154) (-3.694,-6.104) (-3.657,-6.071)}
\LC{(-2.932,-10.000) (-2.901,-9.953) (-2.869,-9.885) (-2.838,-9.817) (-2.807,-9.749) (-2.777,-9.680) (-2.746,-9.612) (-2.715,-9.543) (-2.685,-9.475) (-2.654,-9.406) (-2.624,-9.338) (-2.594,-9.269) (-2.564,-9.200) (-2.534,-9.132) (-2.504,-9.063) (-2.474,-8.994) (-2.445,-8.925) (-2.415,-8.856) (-2.386,-8.787) (-2.357,-8.718) (-2.327,-8.649) (-2.298,-8.580) (-2.269,-8.511) (-2.241,-8.441) (-2.212,-8.372) (-2.183,-8.303) (-2.155,-8.233) (-2.127,-8.164) (-2.099,-8.094) (-2.070,-8.025) (-2.043,-7.955) (-2.015,-7.885) (-1.987,-7.816) (-1.959,-7.746) (-1.932,-7.676) (-1.905,-7.606) (-1.878,-7.536) (-1.850,-7.467) (-1.824,-7.397) (-1.797,-7.326) (-1.770,-7.256) (-1.744,-7.186) (-1.717,-7.116) (-1.691,-7.046) (-1.665,-6.976) (-1.639,-6.905) (-1.613,-6.835) (-1.587,-6.764) (-1.561,-6.694) (-1.536,-6.623) (-1.511,-6.553) (-1.485,-6.482) (-1.460,-6.411) (-1.435,-6.341) (-1.411,-6.270) (-1.386,-6.199) (-1.361,-6.128) (-1.337,-6.057) (-1.313,-5.986) (-1.289,-5.915) (-1.265,-5.844) (-1.241,-5.773) (-1.217,-5.702) (-1.194,-5.631) (-1.170,-5.559) (-1.147,-5.488) (-1.124,-5.417) (-1.101,-5.345) (-1.079,-5.274) (-1.056,-5.202) (-1.034,-5.131) (-1.011,-5.059) (-0.989,-4.987) (-0.967,-4.916) (-0.945,-4.844) (-0.924,-4.772) (-0.902,-4.700) (-0.881,-4.628) (-0.860,-4.556) (-0.839,-4.484) (-0.818,-4.412) (-0.811,-4.388)}
\LC{(-0.002,3.003) (-0.002,3.078) (-0.002,3.153) (-0.002,3.228) (-0.003,3.303) (-0.003,3.378) (-0.003,3.453) (-0.003,3.528) (-0.003,3.603) (-0.003,3.678) (-0.003,3.753) (-0.003,3.828) (-0.003,3.903) (-0.003,3.978) (-0.003,4.053) (-0.004,4.128) (-0.004,4.203) (-0.004,4.278) (-0.004,4.353) (-0.004,4.428) (-0.004,4.503) (-0.004,4.578) (-0.004,4.653) (-0.004,4.728) (-0.004,4.803) (-0.005,4.878) (-0.005,4.953) (-0.005,5.028) (-0.005,5.103) (-0.005,5.178) (-0.005,5.253) (-0.005,5.328) (-0.005,5.403) (-0.005,5.478) (-0.006,5.553) (-0.006,5.628) (-0.006,5.703) (-0.006,5.778) (-0.006,5.853) (-0.006,5.928) (-0.006,6.003) (-0.006,6.078) (-0.007,6.153) (-0.007,6.228) (-0.007,6.303) (-0.007,6.378) (-0.007,6.453) (-0.007,6.528) (-0.007,6.603) (-0.007,6.678) (-0.007,6.753) (-0.008,6.828) (-0.008,6.903) (-0.008,6.978) (-0.008,7.053) (-0.008,7.128) (-0.008,7.203) (-0.008,7.278) (-0.009,7.353) (-0.009,7.428) (-0.009,7.503) (-0.009,7.578) (-0.009,7.653) (-0.009,7.728) (-0.009,7.803) (-0.009,7.878) (-0.010,7.953) (-0.010,8.028) (-0.010,8.103) (-0.010,8.178) (-0.010,8.253) (-0.010,8.328) (-0.010,8.403) (-0.011,8.478) (-0.011,8.553) (-0.011,8.628) (-0.011,8.703) (-0.011,8.778) (-0.011,8.853) (-0.011,8.928) (-0.012,9.003) (-0.012,9.078) (-0.012,9.153) (-0.012,9.228) (-0.012,9.303) (-0.012,9.378) (-0.012,9.453) (-0.013,9.528) (-0.013,9.603) (-0.013,9.678) (-0.013,9.753) (-0.013,9.828) (-0.013,9.903) (-0.013,9.978) (-0.014,10.000)}
\LC{(4.652,4.285) (4.725,4.303) (4.798,4.319) (4.871,4.335) (4.945,4.350) (5.018,4.365) (5.092,4.378) (5.166,4.391) (5.240,4.403) (5.314,4.415) (5.388,4.425) (5.463,4.435) (5.537,4.444) (5.612,4.452) (5.686,4.459) (5.761,4.466) (5.836,4.471) (5.911,4.476) (5.986,4.480) (6.061,4.483) (6.135,4.486) (6.210,4.487) (6.285,4.488) (6.360,4.488) (6.435,4.487) (6.510,4.485) (6.585,4.483) (6.660,4.479) (6.735,4.475) (6.810,4.470) (6.885,4.464) (6.959,4.457) (7.034,4.449) (7.109,4.441) (7.183,4.431) (7.257,4.421) (7.331,4.410) (7.405,4.398) (7.479,4.385) (7.553,4.371) (7.627,4.356) (7.700,4.341) (7.773,4.324) (7.846,4.307) (7.919,4.289) (7.991,4.269) (8.064,4.249) (8.136,4.228) (8.207,4.207) (8.279,4.184) (8.350,4.160) (8.421,4.135) (8.491,4.110) (8.561,4.084) (8.631,4.056) (8.701,4.028) (8.770,3.999) (8.839,3.969) (8.907,3.938) (8.975,3.906) (9.042,3.873) (9.109,3.839) (9.176,3.804) (9.242,3.769) (9.307,3.732) (9.372,3.695) (9.437,3.656) (9.500,3.617) (9.564,3.577) (9.627,3.536) (9.689,3.494) (9.750,3.451) (9.811,3.407) (9.871,3.362) (9.931,3.317) (9.990,3.270) (10.000,3.255)}
\LC{(-3.065,3.859) (-3.131,3.895) (-3.197,3.931) (-3.263,3.966) (-3.330,4.001) (-3.396,4.034) (-3.464,4.068) (-3.531,4.100) (-3.599,4.132) (-3.667,4.163) (-3.736,4.194) (-3.805,4.224) (-3.874,4.253) (-3.943,4.282) (-4.013,4.310) (-4.082,4.337) (-4.152,4.364) (-4.223,4.390) (-4.293,4.415) (-4.364,4.440) (-4.435,4.464) (-4.507,4.487) (-4.578,4.509) (-4.650,4.531) (-4.722,4.553) (-4.794,4.573) (-4.866,4.593) (-4.939,4.612) (-5.011,4.631) (-5.084,4.648) (-5.157,4.665) (-5.231,4.682) (-5.304,4.697) (-5.377,4.712) (-5.451,4.727) (-5.525,4.740) (-5.599,4.753) (-5.673,4.765) (-5.747,4.776) (-5.821,4.787) (-5.895,4.797) (-5.970,4.806) (-6.044,4.814) (-6.119,4.822) (-6.194,4.829) (-6.268,4.835) (-6.343,4.841) (-6.418,4.845) (-6.493,4.849) (-6.568,4.852) (-6.643,4.855) (-6.718,4.856) (-6.793,4.857) (-6.868,4.857) (-6.943,4.857) (-7.018,4.855) (-7.093,4.853) (-7.168,4.850) (-7.243,4.846) (-7.318,4.842) (-7.392,4.836) (-7.467,4.830) (-7.542,4.823) (-7.616,4.816) (-7.691,4.807) (-7.765,4.798) (-7.840,4.788) (-7.914,4.777) (-7.988,4.765) (-8.062,4.752) (-8.136,4.739) (-8.209,4.725) (-8.283,4.710) (-8.356,4.694) (-8.429,4.677) (-8.502,4.660) (-8.575,4.641) (-8.647,4.622) (-8.720,4.602) (-8.792,4.582) (-8.864,4.560) (-8.935,4.537) (-9.006,4.514) (-9.077,4.490) (-9.148,4.465) (-9.218,4.439) (-9.289,4.412) (-9.358,4.385) (-9.428,4.356) (-9.497,4.327) (-9.566,4.297) (-9.634,4.266) (-9.702,4.234) (-9.769,4.202) (-9.836,4.168) (-9.903,4.134) (-9.969,4.099) (-10.000,4.075)}
\LC{(2.988,4.855) (3.044,4.904) (3.102,4.953) (3.159,5.001) (3.217,5.049) (3.275,5.096) (3.333,5.143) (3.392,5.190) (3.451,5.237) (3.510,5.283) (3.570,5.328) (3.629,5.374) (3.689,5.418) (3.750,5.463) (3.810,5.507) (3.871,5.551) (3.932,5.594) (3.994,5.638) (4.055,5.680) (4.117,5.723) (4.179,5.765) (4.242,5.806) (4.305,5.847) (4.367,5.888) (4.431,5.929) (4.494,5.969) (4.558,6.008) (4.622,6.048) (4.686,6.087) (4.750,6.125) (4.814,6.163) (4.879,6.201) (4.944,6.239) (5.010,6.276) (5.075,6.312) (5.141,6.349) (5.206,6.384) (5.273,6.420) (5.339,6.455) (5.405,6.490) (5.472,6.524) (5.539,6.558) (5.606,6.591) (5.673,6.624) (5.741,6.657) (5.808,6.690) (5.876,6.721) (5.944,6.753) (6.013,6.784) (6.081,6.815) (6.150,6.845) (6.218,6.875) (6.287,6.905) (6.356,6.934) (6.426,6.963) (6.495,6.991) (6.565,7.019) (6.635,7.046) (6.704,7.073) (6.775,7.100) (6.845,7.126) (6.915,7.152) (6.986,7.178) (7.056,7.203) (7.127,7.227) (7.198,7.252) (7.269,7.275) (7.341,7.299) (7.412,7.322) (7.484,7.344) (7.555,7.366) (7.627,7.388) (7.699,7.409) (7.771,7.430) (7.843,7.451) (7.915,7.471) (7.988,7.490) (8.060,7.509) (8.133,7.528) (8.206,7.546) (8.279,7.564) (8.352,7.582) (8.425,7.599) (8.498,7.615) (8.571,7.631) (8.644,7.647) (8.718,7.662) (8.791,7.677) (8.865,7.691) (8.939,7.705) (9.012,7.719) (9.086,7.732) (9.160,7.744) (9.234,7.757) (9.308,7.768) (9.382,7.780) (9.457,7.790) (9.531,7.801) (9.605,7.811) (9.680,7.820) (9.754,7.829) (9.829,7.838) (9.903,7.846) (9.978,7.853) (10.000,7.856)}
\LC{(-1.438,1.372) (-1.512,1.387) (-1.586,1.399) (-1.660,1.408) (-1.735,1.415) (-1.810,1.420) (-1.885,1.422) (-1.960,1.422) (-2.035,1.419) (-2.110,1.413) (-2.184,1.405) (-2.258,1.394) (-2.332,1.380) (-2.405,1.364) (-2.478,1.345) (-2.549,1.323) (-2.620,1.298) (-2.690,1.270) (-2.758,1.239) (-2.825,1.206) (-2.891,1.169) (-2.955,1.130) (-3.017,1.088) (-3.077,1.043) (-3.134,0.995) (-3.190,0.945) (-3.243,0.891) (-3.293,0.835) (-3.340,0.777) (-3.383,0.716) (-3.424,0.653) (-3.461,0.588) (-3.494,0.520) (-3.523,0.451) (-3.548,0.381) (-3.569,0.309) (-3.586,0.236) (-3.598,0.162) (-3.605,0.087) (-3.608,0.012) (-3.607,-0.063) (-3.601,-0.138) (-3.590,-0.212) (-3.575,-0.285) (-3.555,-0.358) (-3.532,-0.429) (-3.504,-0.498) (-3.472,-0.566) (-3.436,-0.632) (-3.397,-0.696) (-3.354,-0.758) (-3.308,-0.817) (-3.259,-0.874) (-3.207,-0.928) (-3.153,-0.979) (-3.095,-1.028) (-3.036,-1.074) (-2.975,-1.117) (-2.911,-1.157) (-2.846,-1.194) (-2.780,-1.229) (-2.712,-1.260) (-2.643,-1.289) (-2.572,-1.315) (-2.501,-1.338) (-2.429,-1.358) (-2.356,-1.375) (-2.282,-1.390) (-2.208,-1.402) (-2.133,-1.411) (-2.059,-1.417) (-1.984,-1.421) (-1.909,-1.422) (-1.834,-1.421) (-1.759,-1.417) (-1.684,-1.411) (-1.610,-1.402) (-1.536,-1.391) (-1.462,-1.377) (-1.437,-1.372)}
\LC{(-2.727,1.915) (-2.802,1.912) (-2.876,1.906) (-2.951,1.899) (-3.025,1.890) (-3.100,1.878) (-3.173,1.865) (-3.247,1.850) (-3.320,1.832) (-3.392,1.813) (-3.464,1.791) (-3.535,1.767) (-3.605,1.741) (-3.675,1.713) (-3.744,1.683) (-3.811,1.650) (-3.878,1.616) (-3.943,1.579) (-4.008,1.541) (-4.070,1.500) (-4.132,1.457) (-4.192,1.411) (-4.250,1.364) (-4.306,1.315) (-4.361,1.264) (-4.414,1.210) (-4.465,1.155) (-4.513,1.098) (-4.559,1.039) (-4.603,0.978) (-4.644,0.915) (-4.683,0.851) (-4.719,0.785) (-4.753,0.718) (-4.783,0.650) (-4.811,0.580) (-4.835,0.509) (-4.856,0.437) (-4.874,0.364) (-4.889,0.291) (-4.900,0.217) (-4.909,0.142) (-4.913,0.067) (-4.915,-0.008) (-4.913,-0.083) (-4.907,-0.158) (-4.898,-0.232) (-4.886,-0.306) (-4.871,-0.379) (-4.852,-0.452) (-4.830,-0.524) (-4.805,-0.594) (-4.777,-0.664) (-4.746,-0.732) (-4.712,-0.799) (-4.675,-0.865) (-4.636,-0.928) (-4.594,-0.991) (-4.550,-1.051) (-4.503,-1.110) (-4.454,-1.167) (-4.403,-1.221) (-4.350,-1.274) (-4.295,-1.325) (-4.238,-1.374) (-4.179,-1.421) (-4.119,-1.466) (-4.057,-1.508) (-3.994,-1.549) (-3.930,-1.587) (-3.864,-1.623) (-3.797,-1.657) (-3.730,-1.689) (-3.661,-1.719) (-3.591,-1.747) (-3.520,-1.772) (-3.449,-1.795) (-3.377,-1.817) (-3.305,-1.836) (-3.232,-1.853) (-3.158,-1.868) (-3.084,-1.881) (-3.010,-1.892) (-2.936,-1.901) (-2.861,-1.907) (-2.786,-1.912) (-2.736,-1.915)}
\LC{(-4.344,5.136) (-4.411,5.170) (-4.479,5.202) (-4.547,5.235) (-4.615,5.266) (-4.683,5.298) (-4.751,5.328) (-4.820,5.358) (-4.889,5.388) (-4.958,5.417) (-5.027,5.446) (-5.097,5.474) (-5.166,5.502) (-5.236,5.529) (-5.307,5.555) (-5.377,5.581) (-5.447,5.607) (-5.518,5.632) (-5.589,5.656) (-5.660,5.680) (-5.731,5.703) (-5.803,5.726) (-5.874,5.749) (-5.946,5.770) (-6.018,5.792) (-6.090,5.812) (-6.162,5.832) (-6.235,5.852) (-6.307,5.871) (-6.380,5.890) (-6.453,5.907) (-6.526,5.925) (-6.599,5.942) (-6.672,5.958) (-6.745,5.974) (-6.819,5.989) (-6.893,6.003) (-6.966,6.018) (-7.040,6.031) (-7.114,6.044) (-7.188,6.056) (-7.262,6.068) (-7.336,6.079) (-7.410,6.090) (-7.485,6.100) (-7.559,6.109) (-7.633,6.118) (-7.708,6.127) (-7.783,6.134) (-7.857,6.141) (-7.932,6.148) (-8.007,6.154) (-8.082,6.159) (-8.156,6.164) (-8.231,6.168) (-8.306,6.172) (-8.381,6.175) (-8.456,6.177) (-8.531,6.179) (-8.606,6.180) (-8.681,6.181) (-8.756,6.181) (-8.831,6.180) (-8.906,6.179) (-8.981,6.177) (-9.056,6.175) (-9.131,6.172) (-9.206,6.168) (-9.281,6.164) (-9.356,6.159) (-9.430,6.154) (-9.505,6.147) (-9.580,6.141) (-9.654,6.133) (-9.729,6.125) (-9.804,6.117) (-9.878,6.107) (-9.952,6.097) (-10.000,6.090)}
\LC{(2.341,3.550) (2.401,3.596) (2.461,3.641) (2.521,3.685) (2.582,3.729) (2.643,3.772) (2.705,3.815) (2.767,3.857) (2.829,3.899) (2.892,3.940) (2.955,3.980) (3.019,4.020) (3.082,4.060) (3.146,4.098) (3.211,4.137) (3.276,4.174) (3.341,4.212) (3.406,4.248) (3.472,4.284) (3.538,4.320) (3.605,4.354) (3.672,4.389) (3.739,4.422) (3.806,4.455) (3.873,4.488) (3.941,4.520) (4.009,4.551) (4.078,4.582) (4.147,4.612) (4.216,4.641) (4.285,4.670) (4.354,4.699) (4.424,4.726) (4.494,4.753) (4.564,4.780) (4.634,4.806) (4.705,4.831) (4.776,4.856) (4.847,4.880) (4.918,4.903) (4.990,4.926) (5.061,4.948) (5.133,4.970) (5.205,4.991) (5.277,5.011) (5.350,5.031) (5.422,5.050) (5.495,5.068) (5.568,5.086) (5.641,5.103) (5.714,5.119) (5.787,5.135) (5.861,5.150) (5.934,5.165) (6.008,5.179) (6.082,5.192) (6.156,5.205) (6.230,5.216) (6.304,5.228) (6.378,5.238) (6.453,5.248) (6.527,5.257) (6.602,5.266) (6.676,5.274) (6.751,5.281) (6.825,5.287) (6.900,5.293) (6.975,5.298) (7.050,5.303) (7.125,5.306) (7.200,5.310) (7.275,5.312) (7.350,5.314) (7.425,5.315) (7.500,5.315) (7.575,5.314) (7.650,5.313) (7.725,5.311) (7.800,5.309) (7.875,5.306) (7.949,5.302) (8.024,5.297) (8.099,5.291) (8.174,5.285) (8.249,5.278) (8.323,5.271) (8.398,5.262) (8.472,5.253) (8.546,5.243) (8.621,5.233) (8.695,5.222) (8.769,5.210) (8.843,5.197) (8.917,5.183) (8.990,5.169) (9.064,5.154) (9.137,5.138) (9.210,5.122) (9.283,5.105) (9.356,5.086) (9.429,5.068) (9.501,5.048) (9.573,5.028) (9.645,5.007) (9.717,4.985) (9.788,4.962) (9.860,4.939) (9.931,4.915) (10.000,4.890)}
\LC{(1.360,4.868) (1.391,4.936) (1.421,5.005) (1.452,5.073) (1.483,5.141) (1.515,5.209) (1.546,5.277) (1.578,5.345) (1.610,5.413) (1.642,5.481) (1.674,5.549) (1.707,5.616) (1.740,5.684) (1.772,5.751) (1.806,5.818) (1.839,5.886) (1.872,5.953) (1.906,6.020) (1.940,6.087) (1.974,6.153) (2.008,6.220) (2.043,6.287) (2.077,6.353) (2.112,6.420) (2.147,6.486) (2.182,6.552) (2.218,6.618) (2.253,6.684) (2.289,6.750) (2.325,6.816) (2.361,6.882) (2.397,6.948) (2.434,7.013) (2.470,7.079) (2.507,7.144) (2.544,7.209) (2.581,7.274) (2.619,7.340) (2.656,7.404) (2.694,7.469) (2.731,7.534) (2.769,7.599) (2.808,7.663) (2.846,7.728) (2.884,7.792) (2.923,7.857) (2.962,7.921) (3.001,7.985) (3.040,8.049) (3.079,8.113) (3.119,8.176) (3.158,8.240) (3.198,8.304) (3.238,8.367) (3.278,8.431) (3.318,8.494) (3.359,8.557) (3.399,8.620) (3.440,8.683) (3.481,8.746) (3.522,8.809) (3.563,8.872) (3.604,8.934) (3.646,8.997) (3.687,9.059) (3.729,9.121) (3.771,9.184) (3.813,9.246) (3.855,9.308) (3.897,9.370) (3.940,9.431) (3.983,9.493) (4.025,9.555) (4.068,9.616) (4.111,9.678) (4.154,9.739) (4.198,9.800) (4.241,9.861) (4.285,9.922) (4.329,9.983) (4.343,10.000)}
\LC{(-1.409,4.730) (-1.442,4.797) (-1.474,4.865) (-1.507,4.932) (-1.540,5.000) (-1.573,5.067) (-1.606,5.134) (-1.640,5.201) (-1.674,5.268) (-1.708,5.335) (-1.742,5.402) (-1.777,5.468) (-1.812,5.535) (-1.846,5.601) (-1.882,5.667) (-1.917,5.733) (-1.953,5.799) (-1.988,5.865) (-2.024,5.931) (-2.060,5.997) (-2.097,6.062) (-2.133,6.128) (-2.170,6.193) (-2.207,6.258) (-2.244,6.324) (-2.282,6.389) (-2.319,6.454) (-2.357,6.518) (-2.395,6.583) (-2.433,6.648) (-2.472,6.712) (-2.510,6.776) (-2.549,6.841) (-2.588,6.905) (-2.627,6.969) (-2.666,7.033) (-2.706,7.096) (-2.745,7.160) (-2.785,7.224) (-2.825,7.287) (-2.865,7.350) (-2.906,7.414) (-2.946,7.477) (-2.987,7.540) (-3.028,7.603) (-3.069,7.665) (-3.110,7.728) (-3.152,7.790) (-3.193,7.853) (-3.235,7.915) (-3.277,7.977) (-3.319,8.039) (-3.361,8.101) (-3.404,8.163) (-3.446,8.225) (-3.489,8.287) (-3.532,8.348) (-3.575,8.409) (-3.618,8.471) (-3.662,8.532) (-3.705,8.593) (-3.749,8.654) (-3.793,8.715) (-3.837,8.775) (-3.881,8.836) (-3.926,8.896) (-3.970,8.957) (-4.015,9.017) (-4.060,9.077) (-4.105,9.137) (-4.150,9.197) (-4.195,9.256) (-4.241,9.316) (-4.287,9.376) (-4.332,9.435) (-4.378,9.494) (-4.424,9.553) (-4.471,9.613) (-4.517,9.671) (-4.564,9.730) (-4.610,9.789) (-4.657,9.848) (-4.704,9.906) (-4.751,9.964) (-4.783,10.000)}
\LC{(-10.000,-3.295) (-9.955,-3.342) (-9.896,-3.387) (-9.835,-3.432) (-9.774,-3.475) (-9.713,-3.518) (-9.650,-3.560) (-9.588,-3.601) (-9.524,-3.641) (-9.460,-3.680) (-9.396,-3.718) (-9.330,-3.755) (-9.265,-3.791) (-9.199,-3.827) (-9.132,-3.861) (-9.065,-3.895) (-8.998,-3.928) (-8.930,-3.959) (-8.861,-3.990) (-8.792,-4.020) (-8.723,-4.049) (-8.654,-4.077) (-8.584,-4.104) (-8.514,-4.131) (-8.443,-4.156) (-8.372,-4.181) (-8.301,-4.204) (-8.229,-4.227) (-8.158,-4.248) (-8.086,-4.269) (-8.013,-4.289) (-7.941,-4.308) (-7.868,-4.326) (-7.795,-4.344) (-7.722,-4.360) (-7.648,-4.376) (-7.575,-4.390) (-7.501,-4.404) (-7.427,-4.417) (-7.353,-4.429) (-7.279,-4.440) (-7.205,-4.450) (-7.130,-4.459) (-7.056,-4.468) (-6.981,-4.476) (-6.906,-4.482) (-6.832,-4.488) (-6.757,-4.493) (-6.682,-4.498) (-6.607,-4.501) (-6.532,-4.503) (-6.457,-4.505) (-6.382,-4.506) (-6.307,-4.506) (-6.232,-4.505) (-6.157,-4.504) (-6.082,-4.501) (-6.007,-4.498) (-5.932,-4.494) (-5.858,-4.489) (-5.783,-4.483) (-5.708,-4.477) (-5.633,-4.469) (-5.559,-4.461) (-5.484,-4.452) (-5.410,-4.443) (-5.336,-4.432) (-5.262,-4.421) (-5.188,-4.409) (-5.114,-4.396) (-5.040,-4.382) (-4.966,-4.368) (-4.893,-4.353) (-4.820,-4.337) (-4.746,-4.320) (-4.674,-4.303) (-4.601,-4.285) (-4.528,-4.266) (-4.456,-4.246) (-4.384,-4.226) (-4.312,-4.204) (-4.288,-4.197)}
\LC{(0.452,3.470) (0.467,3.543) (0.482,3.617) (0.497,3.690) (0.512,3.764) (0.528,3.837) (0.544,3.910) (0.560,3.984) (0.576,4.057) (0.592,4.130) (0.608,4.203) (0.624,4.277) (0.641,4.350) (0.658,4.423) (0.674,4.496) (0.691,4.569) (0.709,4.642) (0.726,4.715) (0.743,4.788) (0.761,4.861) (0.778,4.934) (0.796,5.007) (0.814,5.080) (0.832,5.152) (0.850,5.225) (0.868,5.298) (0.887,5.371) (0.905,5.443) (0.924,5.516) (0.943,5.588) (0.962,5.661) (0.981,5.734) (1.000,5.806) (1.019,5.879) (1.038,5.951) (1.058,6.023) (1.078,6.096) (1.097,6.168) (1.117,6.241) (1.137,6.313) (1.157,6.385) (1.177,6.457) (1.198,6.530) (1.218,6.602) (1.239,6.674) (1.259,6.746) (1.280,6.818) (1.301,6.890) (1.322,6.962) (1.343,7.034) (1.364,7.106) (1.385,7.178) (1.407,7.250) (1.428,7.322) (1.450,7.393) (1.472,7.465) (1.493,7.537) (1.515,7.609) (1.537,7.680) (1.560,7.752) (1.582,7.824) (1.604,7.895) (1.627,7.967) (1.649,8.038) (1.672,8.110) (1.695,8.181) (1.717,8.253) (1.740,8.324) (1.763,8.396) (1.787,8.467) (1.810,8.538) (1.833,8.609) (1.857,8.681) (1.880,8.752) (1.904,8.823) (1.927,8.894) (1.951,8.965) (1.975,9.036) (1.999,9.108) (2.023,9.179) (2.048,9.250) (2.072,9.320) (2.096,9.391) (2.121,9.462) (2.145,9.533) (2.170,9.604) (2.195,9.675) (2.220,9.746) (2.245,9.816) (2.270,9.887) (2.295,9.958) (2.311,10.000)}
\LC{(10.000,-9.011) (9.930,-8.994) (9.857,-8.977) (9.784,-8.959) (9.711,-8.941) (9.638,-8.923) (9.566,-8.904) (9.493,-8.885) (9.421,-8.866) (9.348,-8.846) (9.276,-8.826) (9.204,-8.805) (9.132,-8.784) (9.060,-8.763) (8.988,-8.741) (8.916,-8.719) (8.845,-8.697) (8.773,-8.674) (8.702,-8.651) (8.631,-8.628) (8.560,-8.604) (8.489,-8.580) (8.418,-8.556) (8.347,-8.531) (8.276,-8.506) (8.206,-8.480) (8.135,-8.454) (8.065,-8.428) (7.995,-8.402) (7.925,-8.375) (7.855,-8.348) (7.785,-8.320) (7.716,-8.292) (7.646,-8.264) (7.577,-8.235) (7.508,-8.206) (7.439,-8.177) (7.370,-8.147) (7.301,-8.117) (7.232,-8.087) (7.164,-8.056) (7.096,-8.025) (7.027,-7.994) (6.959,-7.963) (6.892,-7.931) (6.824,-7.898) (6.756,-7.866) (6.689,-7.833) (6.622,-7.799) (6.555,-7.766) (6.488,-7.732) (6.421,-7.697) (6.355,-7.663) (6.288,-7.628) (6.222,-7.592) (6.156,-7.557) (6.090,-7.521) (6.025,-7.484) (5.959,-7.448) (5.894,-7.411) (5.829,-7.374) (5.764,-7.336) (5.699,-7.298) (5.635,-7.260) (5.571,-7.221) (5.506,-7.182) (5.442,-7.143) (5.379,-7.104) (5.315,-7.064) (5.252,-7.024) (5.189,-6.983) (5.126,-6.942) (5.063,-6.901) (5.001,-6.860) (4.938,-6.818) (4.876,-6.776) (4.814,-6.734) (4.753,-6.691) (4.712,-6.662)}
\LC{(4.295,5.298) (4.361,5.334) (4.427,5.369) (4.494,5.404) (4.561,5.438) (4.628,5.472) (4.695,5.505) (4.762,5.538) (4.830,5.570) (4.898,5.602) (4.966,5.634) (5.034,5.664) (5.103,5.695) (5.172,5.725) (5.240,5.754) (5.310,5.783) (5.379,5.812) (5.449,5.840) (5.518,5.867) (5.588,5.894) (5.658,5.921) (5.729,5.947) (5.799,5.972) (5.870,5.998) (5.941,6.022) (6.012,6.046) (6.083,6.070) (6.154,6.093) (6.226,6.115) (6.298,6.137) (6.370,6.159) (6.442,6.180) (6.514,6.201) (6.586,6.221) (6.658,6.240) (6.731,6.259) (6.804,6.278) (6.876,6.295) (6.949,6.313) (7.022,6.330) (7.096,6.346) (7.169,6.362) (7.242,6.377) (7.316,6.392) (7.389,6.407) (7.463,6.420) (7.537,6.434) (7.611,6.446) (7.685,6.458) (7.759,6.470) (7.833,6.481) (7.907,6.492) (7.982,6.502) (8.056,6.511) (8.131,6.520) (8.205,6.529) (8.280,6.536) (8.354,6.544) (8.429,6.551) (8.504,6.557) (8.579,6.562) (8.653,6.567) (8.728,6.572) (8.803,6.576) (8.878,6.579) (8.953,6.582) (9.028,6.585) (9.103,6.586) (9.178,6.588) (9.253,6.588) (9.328,6.588) (9.403,6.588) (9.478,6.587) (9.553,6.585) (9.628,6.583) (9.703,6.580) (9.778,6.576) (9.853,6.572) (9.928,6.568) (10.000,6.563)}
\LC{(0.445,6.219) (0.453,6.293) (0.461,6.368) (0.470,6.443) (0.478,6.517) (0.486,6.592) (0.495,6.666) (0.503,6.741) (0.511,6.815) (0.520,6.890) (0.528,6.964) (0.537,7.039) (0.546,7.113) (0.554,7.188) (0.563,7.262) (0.572,7.337) (0.581,7.411) (0.590,7.486) (0.598,7.560) (0.607,7.635) (0.616,7.709) (0.625,7.784) (0.634,7.858) (0.644,7.932) (0.653,8.007) (0.662,8.081) (0.671,8.156) (0.680,8.230) (0.690,8.305) (0.699,8.379) (0.709,8.453) (0.718,8.528) (0.728,8.602) (0.737,8.677) (0.747,8.751) (0.756,8.825) (0.766,8.900) (0.776,8.974) (0.785,9.048) (0.795,9.123) (0.805,9.197) (0.815,9.272) (0.825,9.346) (0.835,9.420) (0.845,9.495) (0.855,9.569) (0.865,9.643) (0.875,9.718) (0.885,9.792) (0.895,9.866) (0.905,9.940) (0.915,10.000)}
\LC{(-0.456,5.643) (-0.465,5.718) (-0.474,5.792) (-0.483,5.867) (-0.493,5.941) (-0.502,6.015) (-0.512,6.090) (-0.521,6.164) (-0.531,6.239) (-0.540,6.313) (-0.550,6.387) (-0.560,6.462) (-0.570,6.536) (-0.579,6.610) (-0.589,6.685) (-0.599,6.759) (-0.609,6.833) (-0.619,6.908) (-0.629,6.982) (-0.640,7.056) (-0.650,7.131) (-0.660,7.205) (-0.671,7.279) (-0.681,7.353) (-0.691,7.428) (-0.702,7.502) (-0.712,7.576) (-0.723,7.650) (-0.734,7.725) (-0.744,7.799) (-0.755,7.873) (-0.766,7.947) (-0.777,8.022) (-0.788,8.096) (-0.799,8.170) (-0.810,8.244) (-0.821,8.318) (-0.832,8.393) (-0.843,8.467) (-0.854,8.541) (-0.866,8.615) (-0.877,8.689) (-0.888,8.763) (-0.900,8.837) (-0.911,8.912) (-0.923,8.986) (-0.934,9.060) (-0.946,9.134) (-0.957,9.208) (-0.969,9.282) (-0.981,9.356) (-0.993,9.430) (-1.004,9.504) (-1.016,9.578) (-1.028,9.652) (-1.040,9.726) (-1.052,9.800) (-1.064,9.874) (-1.076,9.948) (-1.088,10.000)}
\LC{(1.026,1.052) (1.099,1.067) (1.173,1.078) (1.248,1.087) (1.322,1.093) (1.397,1.095) (1.472,1.094) (1.547,1.090) (1.622,1.082) (1.696,1.071) (1.770,1.057) (1.842,1.039) (1.914,1.017) (1.985,0.992) (2.054,0.963) (2.121,0.930) (2.187,0.893) (2.250,0.853) (2.311,0.809) (2.369,0.761) (2.423,0.710) (2.475,0.655) (2.522,0.597) (2.566,0.536) (2.605,0.472) (2.639,0.405) (2.668,0.336) (2.692,0.265) (2.710,0.192) (2.722,0.119) (2.729,0.044) (2.729,-0.031) (2.724,-0.106) (2.712,-0.180) (2.695,-0.253) (2.672,-0.324) (2.644,-0.394) (2.611,-0.461) (2.573,-0.526) (2.530,-0.587) (2.483,-0.646) (2.432,-0.701) (2.378,-0.753) (2.321,-0.801) (2.260,-0.846) (2.198,-0.887) (2.132,-0.924) (2.065,-0.957) (1.996,-0.987) (1.926,-1.013) (1.855,-1.035) (1.782,-1.054) (1.709,-1.069) (1.634,-1.081) (1.560,-1.089) (1.485,-1.094) (1.410,-1.095) (1.335,-1.093) (1.260,-1.088) (1.186,-1.080) (1.112,-1.069) (1.038,-1.055)}
\LC{(1.474,-10.000) (1.458,-9.936) (1.442,-9.863) (1.425,-9.789) (1.409,-9.716) (1.393,-9.643) (1.377,-9.570) (1.361,-9.496) (1.345,-9.423) (1.329,-9.350) (1.313,-9.276) (1.298,-9.203) (1.282,-9.130) (1.266,-9.056) (1.251,-8.983) (1.235,-8.910) (1.220,-8.836) (1.204,-8.763) (1.189,-8.689) (1.174,-8.616) (1.159,-8.543) (1.143,-8.469) (1.128,-8.396) (1.113,-8.322) (1.098,-8.249) (1.084,-8.175) (1.069,-8.102) (1.054,-8.028) (1.039,-7.955) (1.025,-7.881) (1.010,-7.807) (0.996,-7.734) (0.982,-7.660) (0.967,-7.587) (0.953,-7.513) (0.939,-7.439) (0.925,-7.366) (0.911,-7.292) (0.897,-7.218) (0.883,-7.145) (0.869,-7.071) (0.855,-6.997) (0.842,-6.923) (0.828,-6.850) (0.814,-6.776) (0.801,-6.702) (0.788,-6.628) (0.774,-6.554) (0.761,-6.481) (0.748,-6.407) (0.735,-6.333) (0.722,-6.259) (0.709,-6.185) (0.696,-6.111) (0.683,-6.037) (0.670,-5.963) (0.658,-5.890) (0.645,-5.816) (0.633,-5.742) (0.620,-5.668) (0.608,-5.594) (0.596,-5.520) (0.584,-5.446) (0.572,-5.372)}
\LC{(3.662,1.966) (3.733,1.944) (3.805,1.921) (3.876,1.896) (3.946,1.869) (4.015,1.840) (4.083,1.809) (4.150,1.776) (4.217,1.741) (4.282,1.705) (4.346,1.666) (4.409,1.625) (4.471,1.582) (4.531,1.537) (4.589,1.490) (4.646,1.442) (4.702,1.391) (4.755,1.339) (4.807,1.285) (4.857,1.229) (4.905,1.171) (4.951,1.111) (4.994,1.050) (5.035,0.987) (5.074,0.923) (5.110,0.858) (5.144,0.791) (5.175,0.722) (5.203,0.653) (5.229,0.582) (5.252,0.511) (5.271,0.439) (5.288,0.366) (5.302,0.292) (5.313,0.218) (5.320,0.143) (5.324,0.068) (5.326,-0.007) (5.324,-0.082) (5.319,-0.157) (5.311,-0.231) (5.300,-0.305) (5.285,-0.379) (5.268,-0.452) (5.248,-0.524) (5.225,-0.596) (5.198,-0.666) (5.170,-0.735) (5.138,-0.803) (5.104,-0.870) (5.067,-0.935) (5.028,-0.999) (4.986,-1.062) (4.942,-1.122) (4.896,-1.182) (4.848,-1.239) (4.798,-1.295) (4.746,-1.349) (4.692,-1.401) (4.636,-1.451) (4.579,-1.499) (4.520,-1.545) (4.459,-1.590) (4.398,-1.632) (4.334,-1.673) (4.270,-1.711) (4.205,-1.748) (4.138,-1.783) (4.070,-1.815) (4.002,-1.846) (3.933,-1.874) (3.863,-1.901) (3.792,-1.926) (3.720,-1.948) (3.648,-1.969)}
\LC{(-7.036,-10.000) (-6.975,-9.958) (-6.914,-9.914) (-6.854,-9.870) (-6.793,-9.826) (-6.732,-9.781) (-6.672,-9.737) (-6.612,-9.692) (-6.552,-9.647) (-6.492,-9.602) (-6.432,-9.556) (-6.373,-9.511) (-6.314,-9.465) (-6.254,-9.419) (-6.195,-9.373) (-6.136,-9.326) (-6.078,-9.280) (-6.019,-9.233) (-5.961,-9.186) (-5.902,-9.139) (-5.844,-9.091) (-5.786,-9.044) (-5.728,-8.996) (-5.671,-8.948) (-5.613,-8.899) (-5.556,-8.851) (-5.499,-8.802) (-5.442,-8.754) (-5.385,-8.705) (-5.329,-8.655) (-5.273,-8.606) (-5.216,-8.556) (-5.160,-8.506) (-5.104,-8.456) (-5.049,-8.406) (-4.993,-8.356) (-4.938,-8.305) (-4.883,-8.254) (-4.828,-8.203) (-4.773,-8.152) (-4.718,-8.100) (-4.664,-8.049) (-4.610,-7.997) (-4.556,-7.945) (-4.502,-7.893) (-4.448,-7.840) (-4.395,-7.788) (-4.342,-7.735) (-4.288,-7.682) (-4.236,-7.629) (-4.183,-7.575) (-4.130,-7.522) (-4.078,-7.468) (-4.026,-7.414) (-3.974,-7.360) (-3.923,-7.306) (-3.871,-7.251) (-3.820,-7.196) (-3.769,-7.141) (-3.718,-7.086) (-3.667,-7.031) (-3.617,-6.975) (-3.567,-6.920) (-3.517,-6.864) (-3.467,-6.808) (-3.417,-6.751) (-3.368,-6.695) (-3.319,-6.638) (-3.270,-6.581) (-3.221,-6.524) (-3.173,-6.467) (-3.125,-6.410) (-3.077,-6.352) (-3.029,-6.294) (-2.981,-6.236) (-2.934,-6.178) (-2.887,-6.120) (-2.840,-6.061) (-2.793,-6.002) (-2.747,-5.943) (-2.701,-5.884) (-2.655,-5.825) (-2.609,-5.765) (-2.564,-5.706) (-2.519,-5.646) (-2.474,-5.586) (-2.429,-5.525) (-2.385,-5.465)}
\LC{(10.000,-2.634) (9.966,-2.689) (9.914,-2.743) (9.861,-2.796) (9.807,-2.848) (9.753,-2.900) (9.697,-2.950) (9.641,-3.000) (9.584,-3.048) (9.526,-3.096) (9.468,-3.143) (9.408,-3.189) (9.348,-3.234) (9.287,-3.278) (9.226,-3.321) (9.164,-3.363) (9.101,-3.404) (9.038,-3.444) (8.974,-3.483) (8.909,-3.522) (8.844,-3.559) (8.779,-3.595) (8.713,-3.631) (8.646,-3.665) (8.579,-3.698) (8.511,-3.731) (8.443,-3.762) (8.375,-3.793) (8.306,-3.823) (8.236,-3.851) (8.167,-3.879) (8.097,-3.906) (8.026,-3.931) (7.955,-3.956) (7.884,-3.980) (7.813,-4.003) (7.741,-4.025) (7.669,-4.046) (7.597,-4.066) (7.524,-4.085) (7.451,-4.103) (7.378,-4.120) (7.305,-4.136) (7.232,-4.151) (7.158,-4.166) (7.084,-4.179) (7.011,-4.192) (6.936,-4.203) (6.862,-4.214) (6.788,-4.224) (6.713,-4.233) (6.639,-4.241) (6.564,-4.248) (6.489,-4.254) (6.415,-4.259) (6.340,-4.263) (6.265,-4.267) (6.190,-4.269) (6.115,-4.271) (6.040,-4.272) (5.965,-4.272) (5.890,-4.271) (5.815,-4.269) (5.740,-4.266) (5.665,-4.263) (5.590,-4.258) (5.515,-4.253) (5.441,-4.247) (5.366,-4.240) (5.291,-4.232) (5.217,-4.224) (5.142,-4.214) (5.068,-4.204) (4.994,-4.193) (4.920,-4.181) (4.846,-4.168) (4.772,-4.154) (4.699,-4.140) (4.625,-4.125) (4.552,-4.109) (4.479,-4.092) (4.455,-4.086)}
\LC{(-10.000,-7.077) (-9.926,-7.077) (-9.851,-7.077) (-9.776,-7.075) (-9.701,-7.074) (-9.626,-7.071) (-9.552,-7.068) (-9.477,-7.065) (-9.402,-7.061) (-9.327,-7.057) (-9.252,-7.052) (-9.177,-7.047) (-9.102,-7.041) (-9.028,-7.034) (-8.953,-7.027) (-8.878,-7.020) (-8.804,-7.012) (-8.729,-7.003) (-8.655,-6.994) (-8.580,-6.985) (-8.506,-6.975) (-8.432,-6.964) (-8.358,-6.953) (-8.284,-6.942) (-8.210,-6.930) (-8.136,-6.918) (-8.062,-6.905) (-7.988,-6.891) (-7.914,-6.877) (-7.841,-6.863) (-7.767,-6.848) (-7.694,-6.832) (-7.620,-6.816) (-7.547,-6.800) (-7.474,-6.783) (-7.401,-6.766) (-7.328,-6.748) (-7.256,-6.729) (-7.183,-6.711) (-7.111,-6.691) (-7.038,-6.671) (-6.966,-6.651) (-6.894,-6.630) (-6.822,-6.609) (-6.750,-6.588) (-6.679,-6.565) (-6.607,-6.543) (-6.536,-6.520) (-6.465,-6.496) (-6.393,-6.472) (-6.323,-6.448) (-6.252,-6.423) (-6.181,-6.397) (-6.111,-6.371) (-6.041,-6.345) (-5.971,-6.318) (-5.901,-6.291) (-5.831,-6.263) (-5.762,-6.235) (-5.692,-6.206) (-5.623,-6.177) (-5.554,-6.148) (-5.485,-6.118) (-5.417,-6.087) (-5.349,-6.057) (-5.280,-6.025) (-5.212,-5.993) (-5.145,-5.961) (-5.077,-5.929) (-5.010,-5.896)}
\LC{(-1.487,0.918) (-1.560,0.900) (-1.632,0.878) (-1.702,0.852) (-1.771,0.822) (-1.837,0.787) (-1.901,0.748) (-1.962,0.705) (-2.021,0.657) (-2.075,0.606) (-2.125,0.550) (-2.171,0.491) (-2.212,0.428) (-2.247,0.362) (-2.277,0.293) (-2.301,0.222) (-2.318,0.149) (-2.328,0.075) (-2.331,-0.000) (-2.328,-0.075) (-2.318,-0.149) (-2.300,-0.222) (-2.277,-0.294) (-2.247,-0.362) (-2.212,-0.428) (-2.171,-0.491) (-2.125,-0.551) (-2.075,-0.606) (-2.020,-0.658) (-1.962,-0.705) (-1.901,-0.748) (-1.837,-0.787) (-1.770,-0.822) (-1.702,-0.852) (-1.631,-0.879) (-1.560,-0.901) (-1.487,-0.918)}
\LC{(6.535,-10.000) (6.476,-9.962) (6.418,-9.915) (6.360,-9.867) (6.302,-9.820) (6.244,-9.772) (6.186,-9.724) (6.129,-9.676) (6.071,-9.628) (6.014,-9.579) (5.957,-9.531) (5.900,-9.482) (5.843,-9.433) (5.787,-9.384) (5.730,-9.334) (5.674,-9.285) (5.618,-9.235) (5.562,-9.185) (5.506,-9.135) (5.451,-9.084) (5.395,-9.034) (5.340,-8.983) (5.285,-8.932) (5.230,-8.881) (5.175,-8.830) (5.120,-8.779) (5.066,-8.727) (5.012,-8.675) (4.958,-8.623) (4.904,-8.571) (4.850,-8.519) (4.796,-8.467) (4.743,-8.414) (4.690,-8.361) (4.637,-8.308) (4.584,-8.255) (4.531,-8.201) (4.479,-8.148) (4.426,-8.094) (4.374,-8.040) (4.322,-7.986) (4.271,-7.932) (4.219,-7.877) (4.168,-7.823) (4.117,-7.768) (4.066,-7.713) (4.015,-7.657) (3.964,-7.602) (3.914,-7.547) (3.864,-7.491) (3.814,-7.435) (3.764,-7.379) (3.714,-7.323) (3.665,-7.266) (3.616,-7.210) (3.567,-7.153) (3.518,-7.096) (3.469,-7.039) (3.421,-6.981) (3.373,-6.924) (3.325,-6.866) (3.277,-6.808) (3.229,-6.750) (3.182,-6.692) (3.135,-6.634) (3.088,-6.575) (3.041,-6.517) (2.995,-6.458) (2.949,-6.399) (2.903,-6.339) (2.857,-6.280) (2.811,-6.220) (2.766,-6.161) (2.721,-6.101) (2.676,-6.041) (2.632,-5.980) (2.587,-5.920) (2.543,-5.859)}
\LC{(-2.806,6.393) (-2.851,6.453) (-2.896,6.513) (-2.941,6.573) (-2.986,6.633) (-3.031,6.693) (-3.077,6.752) (-3.123,6.812) (-3.169,6.871) (-3.215,6.930) (-3.261,6.989) (-3.308,7.048) (-3.355,7.106) (-3.402,7.165) (-3.449,7.223) (-3.496,7.281) (-3.544,7.339) (-3.592,7.397) (-3.640,7.454) (-3.688,7.512) (-3.736,7.569) (-3.785,7.626) (-3.834,7.683) (-3.883,7.740) (-3.932,7.797) (-3.981,7.853) (-4.031,7.909) (-4.081,7.966) (-4.131,8.021) (-4.181,8.077) (-4.231,8.133) (-4.282,8.188) (-4.332,8.244) (-4.383,8.299) (-4.434,8.354) (-4.485,8.409) (-4.537,8.463) (-4.588,8.518) (-4.640,8.572) (-4.692,8.626) (-4.744,8.680) (-4.796,8.734) (-4.849,8.788) (-4.901,8.841) (-4.954,8.894) (-5.007,8.947) (-5.060,9.000) (-5.113,9.053) (-5.167,9.106) (-5.220,9.158) (-5.274,9.211) (-5.328,9.263) (-5.382,9.315) (-5.436,9.366) (-5.491,9.418) (-5.545,9.470) (-5.600,9.521) (-5.655,9.572) (-5.710,9.623) (-5.765,9.674) (-5.821,9.724) (-5.876,9.775) (-5.932,9.825) (-5.988,9.875) (-6.044,9.925) (-6.100,9.975) (-6.137,10.000)}
\LC{(-1.501,-10.000) (-1.485,-9.935) (-1.468,-9.862) (-1.451,-9.789) (-1.435,-9.716) (-1.418,-9.643) (-1.402,-9.569) (-1.386,-9.496) (-1.370,-9.423) (-1.353,-9.350) (-1.337,-9.276) (-1.321,-9.203) (-1.305,-9.130) (-1.289,-9.057) (-1.274,-8.983) (-1.258,-8.910) (-1.242,-8.837) (-1.226,-8.763) (-1.211,-8.690) (-1.195,-8.617) (-1.180,-8.543) (-1.164,-8.470) (-1.149,-8.396) (-1.134,-8.323) (-1.119,-8.249) (-1.104,-8.176) (-1.089,-8.103) (-1.074,-8.029) (-1.059,-7.956) (-1.044,-7.882) (-1.029,-7.808) (-1.014,-7.735) (-1.000,-7.661) (-0.985,-7.588) (-0.971,-7.514) (-0.956,-7.441) (-0.942,-7.367) (-0.928,-7.293) (-0.913,-7.220) (-0.899,-7.146) (-0.885,-7.072) (-0.871,-6.999) (-0.857,-6.925) (-0.843,-6.851) (-0.830,-6.778) (-0.816,-6.704) (-0.802,-6.630) (-0.789,-6.556) (-0.775,-6.483) (-0.762,-6.409) (-0.748,-6.335) (-0.735,-6.261) (-0.722,-6.187) (-0.709,-6.113) (-0.696,-6.040) (-0.683,-5.966) (-0.670,-5.892) (-0.657,-5.818) (-0.645,-5.744) (-0.632,-5.670) (-0.620,-5.596) (-0.607,-5.522) (-0.595,-5.448) (-0.582,-5.374) (-0.570,-5.300) (-0.558,-5.226) (-0.546,-5.152) (-0.534,-5.078) (-0.522,-5.004)}
\LC{(-5.630,4.299) (-5.705,4.304) (-5.780,4.308) (-5.855,4.311) (-5.929,4.313) (-6.004,4.314) (-6.079,4.314) (-6.154,4.314) (-6.229,4.312) (-6.304,4.310) (-6.379,4.307) (-6.454,4.303) (-6.529,4.298) (-6.604,4.292) (-6.679,4.285) (-6.753,4.278) (-6.828,4.269) (-6.902,4.260) (-6.976,4.249) (-7.051,4.238) (-7.125,4.226) (-7.198,4.213) (-7.272,4.199) (-7.346,4.184) (-7.419,4.168) (-7.492,4.151) (-7.565,4.134) (-7.638,4.115) (-7.710,4.096) (-7.782,4.075) (-7.854,4.054) (-7.926,4.031) (-7.997,4.008) (-8.068,3.984) (-8.138,3.959) (-8.209,3.932) (-8.279,3.905) (-8.348,3.877) (-8.417,3.848) (-8.486,3.818) (-8.554,3.787) (-8.622,3.755) (-8.690,3.723) (-8.757,3.689) (-8.823,3.654) (-8.889,3.618) (-8.955,3.582) (-9.019,3.544) (-9.084,3.506) (-9.148,3.466) (-9.211,3.426) (-9.273,3.384) (-9.335,3.342) (-9.396,3.299) (-9.457,3.254) (-9.517,3.209) (-9.576,3.163) (-9.634,3.116) (-9.692,3.068) (-9.749,3.019) (-9.805,2.969) (-9.860,2.919) (-9.915,2.867) (-9.968,2.815) (-10.000,2.779)}

    % Dipolo físico - puntas de flecha
% Generado por scripts/generar_datos.py (no editar a mano).
\FL{1.382}{0.003}{-89.7}
\FL{1.658}{2.024}{26.2}
\FL{4.571}{2.223}{-20.4}
\FL{6.266}{-0.023}{-90.6}
\FL{4.528}{-2.238}{-160.4}
\FL{1.617}{-2.004}{153.2}
\FL{-1.658}{2.024}{153.8}
\FL{-4.571}{2.223}{-159.6}
\FL{-6.266}{-0.023}{-89.4}
\FL{-4.528}{-2.238}{-19.6}
\FL{-1.617}{-2.004}{26.8}
\FL{-1.382}{0.003}{-90.3}
\FL{3.873}{-8.712}{126.6}
\FL{2.098}{-5.992}{119.5}
\FL{0.715}{-3.054}{110.5}
\FL{-4.032}{-8.737}{52.1}
\FL{-2.189}{-6.033}{59.5}
\FL{-0.751}{-3.094}{68.8}
\FL{8.276}{-6.267}{176.4}
\FL{4.938}{-5.465}{156.7}
\FL{2.064}{-3.588}{136.7}
\FL{0.802}{2.873}{65.6}
\FL{2.098}{5.181}{56.2}
\FL{3.739}{7.291}{48.3}
\FL{5.620}{9.155}{41.3}
\FL{4.821}{2.884}{-9.5}
\FL{7.202}{1.384}{-58.2}
\FL{7.199}{-1.389}{-121.9}
\FL{4.816}{-2.885}{-170.6}
\FL{-2.682}{3.198}{153.6}
\FL{-5.281}{3.843}{178.8}
\FL{-7.883}{3.242}{-151.6}
\FL{-9.692}{1.308}{-112.7}
\FL{-9.687}{-1.320}{-67.1}
\FL{-7.871}{-3.248}{-28.2}
\FL{-5.268}{-3.843}{1.3}
\FL{-2.670}{-3.192}{26.5}
\FL{-0.035}{-8.072}{89.6}
\FL{-0.013}{-4.247}{89.7}
\FL{-1.003}{3.033}{118.6}
\FL{-2.708}{5.495}{130.4}
\FL{-4.844}{7.595}{140.4}
\FL{-7.300}{9.313}{149.6}
\FL{-3.060}{4.764}{141.0}
\FL{-5.630}{6.344}{155.7}
\FL{-8.515}{7.227}{170.3}
\FL{-8.690}{-7.904}{13.4}
\FL{-6.190}{-7.039}{24.8}
\FL{-3.910}{-5.697}{36.3}
\FL{-1.957}{-3.913}{48.8}
\FL{-7.291}{2.794}{-148.2}
\FL{-8.947}{-0.031}{-89.4}
\FL{-7.260}{-2.813}{-31.3}
\FL{8.390}{-7.354}{168.6}
\FL{5.326}{-6.297}{153.4}
\FL{2.988}{3.361}{23.6}
\FL{6.362}{3.813}{-9.2}
\FL{9.300}{2.161}{-52.1}
\FL{-4.856}{3.051}{-173.2}
\FL{-7.529}{1.485}{-122.6}
\FL{-7.511}{-1.513}{-56.8}
\FL{-4.823}{-3.055}{-6.4}
\FL{3.610}{6.131}{42.8}
\FL{5.966}{7.941}{32.3}
\FL{8.606}{9.304}{22.3}
\FL{3.458}{1.326}{-32.4}
\FL{3.422}{-1.349}{-148.8}
\FL{2.162}{-8.591}{111.4}
\FL{1.164}{-5.763}{107.4}
\FL{-4.971}{6.769}{146.7}
\FL{-8.260}{8.387}{160.9}
\FL{9.417}{0.793}{-75.7}
\FL{8.741}{-2.115}{-129.3}
\FL{6.090}{-3.595}{-170.0}
\FL{3.046}{-3.272}{158.7}
\FL{7.948}{-9.304}{153.7}
\FL{5.021}{-7.488}{142.5}
\FL{2.520}{-5.122}{130.4}
\FL{-8.680}{-5.785}{-3.3}
\FL{-6.026}{-5.537}{13.8}
\FL{-3.553}{-4.540}{30.1}
\FL{8.604}{-5.222}{-171.6}
\FL{5.771}{-5.124}{168.1}
\FL{-0.956}{4.794}{107.2}
\FL{-2.251}{8.312}{113.0}
\FL{-8.277}{-8.959}{22.4}
\FL{-5.083}{-7.207}{35.2}
\FL{-2.318}{-8.626}{67.2}
\FL{-1.257}{-5.820}{71.5}
\FL{-0.004}{4.753}{90.1}
\FL{-0.010}{8.253}{90.1}
\FL{6.061}{4.483}{2.2}
\FL{8.816}{3.979}{-23.8}
\FL{-4.746}{4.560}{164.0}
\FL{-8.332}{4.699}{-167.7}
\FL{4.558}{6.008}{31.8}
\FL{8.109}{7.522}{14.4}
\FL{-2.869}{1.182}{-150.5}
\FL{-2.846}{-1.194}{-28.6}
\FL{-4.231}{1.380}{-140.6}
\FL{-4.219}{-1.390}{-39.0}
\FL{-5.708}{5.696}{161.9}
\FL{-8.556}{6.180}{179.0}
\FL{3.472}{4.284}{28.4}
\FL{6.008}{5.179}{10.4}
\FL{8.695}{5.222}{-8.9}
\FL{2.008}{6.220}{62.7}
\FL{3.508}{8.788}{56.8}
\FL{-2.133}{6.128}{119.3}
\FL{-3.837}{8.775}{126.1}
\FL{-8.700}{-4.059}{-22.1}
\FL{-5.758}{-4.481}{4.9}
\FL{0.826}{5.128}{76.1}
\FL{1.763}{8.396}{72.1}
\FL{8.583}{-8.612}{161.5}
\FL{5.938}{-7.436}{150.5}
\FL{5.658}{5.921}{20.5}
\FL{8.529}{6.559}{4.4}
\FL{0.668}{8.131}{82.9}
\FL{-0.752}{7.848}{98.3}
\FL{2.730}{-0.006}{-90.4}
\FL{0.986}{-7.685}{101.0}
\FL{4.841}{1.247}{-48.7}
\FL{4.815}{-1.276}{-132.4}
\FL{-5.728}{-8.996}{39.7}
\FL{-3.401}{-6.732}{48.8}
\FL{8.823}{-3.571}{-150.9}
\FL{5.940}{-4.272}{179.4}
\FL{-7.426}{-6.771}{13.5}
\FL{-2.331}{-0.025}{-88.2}
\FL{5.432}{-9.068}{137.7}
\FL{3.421}{-6.981}{130.1}
\FL{-4.383}{8.299}{132.8}
\FL{-0.966}{-7.490}{78.9}
\FL{-7.997}{4.008}{-161.5}
}
  \fill[mmarojo] (0,0.5) circle (0.2) (0,-0.5) circle (0.2);
  \node[font=\scriptsize, text=mmarojo, fill=white, fill opacity=0.85, text opacity=1, inner sep=0.6pt, anchor=west] at (0.4,0.62) {$+q$};
  \node[font=\scriptsize, text=mmarojo, fill=white, fill opacity=0.85, text opacity=1, inner sep=0.6pt, anchor=west] at (0.4,-0.62) {$-q$};
  \marcoMMA{10}{5}{$x/\ell$}{$z/\ell$}
\end{tikzpicture}\\[1pt]{\small\scshape (b)\; dipolo físico}
\end{minipage}
\caption{Equipotenciales (gris) y líneas de campo (azul) de (a)~un dipolo puntual $\vp=q\ell\,\uz$ y (b)~un dipolo físico con $+q$ en
$z=\ell/2$ y $-q$ en $z=-\ell/2$, ambos con los mismos niveles de potencial. Fuera del disco central, de radio $1{,}2\,\ell$, las dos
representaciones coinciden a simple vista; las pequeñas diferencias se concentran cerca de las cargas y en los «haces» de equipotenciales que
convergen hacia el plano ecuatorial. Recreación de~\cite{FIGURAS-T2}.}
\label{fig:fisico-puntual}
\end{figure}


\begin{observacion}[Lectura de la figura]\label{obs:discrepancia}
La \cref{fig:discrepancia} representa $\varepsilon$ a lo largo de tres semirrectas que parten del origen, $\theta=0$, $\pi/4$ y $\pi/2$. Cerca del
dipolo ($r\sim\ell$) la discrepancia es de decenas por ciento, pero cae rápidamente: en el peor de los casos (sobre el eje) es del $2\%$ a
$r=5\ell$ y del $0{,}5\%$ a $r=10\ell$. En todas las direcciones decrece como $(\ell/r)^{2}$, aunque con un coeficiente que depende de $\theta$. La
razón la veremos en la \cref{sec:multipolar}: el dipolo físico es un dipolo puntual \emph{más} una serie de multipolos de orden superior
(los de orden par se anulan por simetría), el primero de los cuales, el octupolar $\ell=3$, pesa $(\ell/r)^{2}$ respecto del dipolar. Haciendo la
cuenta con el polinomio de Legendre $\Pl{3}$ (\cref{def:legendre}) se obtiene
\[
  \varepsilon(r,\theta)=\frac{\ell^{2}}{r^{2}}\,\frac{\sqrt{\Pl{3}(\cos\theta)^{2}+\frac{1}{16}\sen^{2}\theta\,\Pl{3}'(\cos\theta)^{2}}}{\sqrt{1+3\cos^{2}\theta}}
  +\ord\Bigl(\frac{\ell^{4}}{r^{4}}\Bigr),
\]
cuyo coeficiente vale $1/2$, $0{,}275$ y $3/8$ para $\theta=0$, $\pi/4$ y $\pi/2$, de acuerdo con \cref{eq:discrepancia} y con la figura.
\end{observacion}

% =========================================================================
% Figura (clon de «Tema 2. Figuras», diapositiva «Discrepancia entre dipolo
% físico y puntual»): semirrectas theta = 0, pi/4, pi/2 sobre el campo del
% dipolo puntual y discrepancia relativa |E_p - E_f|/|E_f| (en %) a lo largo
% de ellas. Datos: figuras/datos/discrepancia.dat (exactos, no ajustados).
% =========================================================================
\begin{figure}[t]
\centering
\begin{minipage}[c]{0.5\textwidth}\centering
\begin{tikzpicture}[x=0.31cm, y=0.31cm]
  \panelMMA{10}{mmafondo}{1.2}{%
    % Dipolo puntual - equipotenciales
% Generado por scripts/generar_datos.py (no editar a mano).
\EQ{(1.224,-0.476) (1.250,-0.522) (1.275,-0.574) (1.296,-0.625) (1.313,-0.675) (1.326,-0.725) (1.337,-0.775) (1.345,-0.825) (1.350,-0.875) (1.353,-0.925) (1.354,-0.975) (1.352,-1.025) (1.347,-1.075) (1.340,-1.125) (1.331,-1.175) (1.319,-1.225) (1.305,-1.275) (1.288,-1.325) (1.269,-1.375) (1.246,-1.425) (1.221,-1.475) (1.194,-1.525) (1.162,-1.575) (1.128,-1.625) (1.089,-1.675) (1.050,-1.721) (1.000,-1.774) (0.950,-1.822) (0.900,-1.865) (0.856,-1.900) (0.800,-1.940) (0.750,-1.972) (0.702,-2.000) (0.656,-2.025) (0.604,-2.050) (0.550,-2.074) (0.500,-2.093) (0.450,-2.111) (0.400,-2.126) (0.350,-2.139) (0.300,-2.151) (0.250,-2.161) (0.200,-2.168) (0.150,-2.174) (0.100,-2.179) (0.050,-2.181) (0.000,-2.182) (-0.050,-2.181) (-0.100,-2.179) (-0.150,-2.174) (-0.200,-2.168) (-0.250,-2.161) (-0.300,-2.151) (-0.350,-2.139) (-0.400,-2.126) (-0.450,-2.111) (-0.500,-2.093) (-0.547,-2.075) (-0.600,-2.052) (-0.650,-2.028) (-0.700,-2.001) (-0.745,-1.975) (-0.800,-1.940) (-0.850,-1.904) (-0.900,-1.865) (-0.947,-1.825) (-0.999,-1.775) (-1.047,-1.725) (-1.089,-1.675) (-1.125,-1.629) (-1.162,-1.575) (-1.194,-1.525) (-1.221,-1.475) (-1.246,-1.425) (-1.269,-1.375) (-1.288,-1.325) (-1.305,-1.275) (-1.319,-1.225) (-1.331,-1.175) (-1.340,-1.125) (-1.347,-1.075) (-1.352,-1.025) (-1.354,-0.975) (-1.353,-0.925) (-1.350,-0.875) (-1.345,-0.825) (-1.337,-0.775) (-1.326,-0.725) (-1.313,-0.675) (-1.296,-0.625) (-1.275,-0.575) (-1.252,-0.525) (-1.225,-0.477)}
\EQ{(1.229,-0.446) (1.263,-0.500) (1.290,-0.550) (1.313,-0.600) (1.333,-0.650) (1.349,-0.700) (1.362,-0.750) (1.373,-0.800) (1.381,-0.850) (1.386,-0.900) (1.389,-0.950) (1.389,-1.000) (1.387,-1.050) (1.383,-1.100) (1.376,-1.150) (1.367,-1.200) (1.355,-1.250) (1.341,-1.300) (1.325,-1.349) (1.306,-1.400) (1.284,-1.450) (1.260,-1.500) (1.233,-1.550) (1.203,-1.600) (1.175,-1.642) (1.132,-1.700) (1.100,-1.740) (1.050,-1.796) (1.000,-1.846) (0.950,-1.891) (0.909,-1.925) (0.850,-1.970) (0.806,-2.000) (0.750,-2.035) (0.700,-2.063) (0.650,-2.089) (0.600,-2.112) (0.550,-2.133) (0.500,-2.152) (0.450,-2.169) (0.400,-2.184) (0.350,-2.197) (0.300,-2.208) (0.250,-2.218) (0.200,-2.225) (0.150,-2.231) (0.100,-2.236) (0.050,-2.238) (0.000,-2.239) (-0.050,-2.238) (-0.100,-2.236) (-0.150,-2.231) (-0.200,-2.225) (-0.250,-2.218) (-0.300,-2.208) (-0.350,-2.197) (-0.400,-2.184) (-0.450,-2.169) (-0.500,-2.152) (-0.550,-2.133) (-0.600,-2.112) (-0.650,-2.089) (-0.700,-2.063) (-0.750,-2.035) (-0.800,-2.004) (-0.843,-1.975) (-0.900,-1.932) (-0.940,-1.900) (-0.996,-1.850) (-1.046,-1.800) (-1.091,-1.750) (-1.125,-1.709) (-1.169,-1.650) (-1.200,-1.604) (-1.233,-1.550) (-1.260,-1.500) (-1.284,-1.450) (-1.306,-1.400) (-1.325,-1.350) (-1.341,-1.300) (-1.355,-1.250) (-1.367,-1.200) (-1.376,-1.150) (-1.383,-1.100) (-1.387,-1.050) (-1.389,-1.000) (-1.389,-0.950) (-1.386,-0.900) (-1.381,-0.850) (-1.373,-0.800) (-1.362,-0.750) (-1.349,-0.700) (-1.333,-0.650) (-1.313,-0.600) (-1.290,-0.550) (-1.263,-0.500) (-1.232,-0.450) (-1.229,-0.446)}
\EQ{(1.244,-0.431) (1.274,-0.475) (1.300,-0.518) (1.325,-0.565) (1.350,-0.619) (1.372,-0.675) (1.388,-0.725) (1.400,-0.775) (1.411,-0.825) (1.419,-0.875) (1.424,-0.925) (1.427,-0.975) (1.427,-1.025) (1.425,-1.075) (1.421,-1.125) (1.414,-1.175) (1.406,-1.225) (1.395,-1.275) (1.381,-1.325) (1.365,-1.375) (1.347,-1.425) (1.326,-1.475) (1.303,-1.525) (1.277,-1.575) (1.250,-1.621) (1.215,-1.675) (1.180,-1.725) (1.141,-1.775) (1.100,-1.822) (1.050,-1.874) (1.000,-1.922) (0.950,-1.965) (0.906,-2.000) (0.850,-2.040) (0.800,-2.073) (0.755,-2.100) (0.709,-2.125) (0.660,-2.150) (0.605,-2.175) (0.550,-2.198) (0.500,-2.216) (0.450,-2.233) (0.400,-2.247) (0.350,-2.260) (0.300,-2.271) (0.250,-2.280) (0.200,-2.287) (0.150,-2.293) (0.100,-2.297) (0.050,-2.299) (0.000,-2.300) (-0.050,-2.299) (-0.100,-2.297) (-0.150,-2.293) (-0.200,-2.287) (-0.250,-2.280) (-0.300,-2.271) (-0.350,-2.260) (-0.400,-2.247) (-0.450,-2.233) (-0.500,-2.216) (-0.550,-2.198) (-0.600,-2.177) (-0.650,-2.155) (-0.700,-2.130) (-0.750,-2.103) (-0.796,-2.075) (-0.850,-2.040) (-0.900,-2.004) (-0.950,-1.965) (-0.997,-1.925) (-1.049,-1.875) (-1.097,-1.825) (-1.141,-1.775) (-1.175,-1.731) (-1.215,-1.675) (-1.248,-1.625) (-1.275,-1.578) (-1.300,-1.531) (-1.325,-1.478) (-1.347,-1.425) (-1.365,-1.375) (-1.381,-1.325) (-1.395,-1.275) (-1.406,-1.225) (-1.414,-1.175) (-1.421,-1.125) (-1.425,-1.075) (-1.427,-1.025) (-1.427,-0.975) (-1.424,-0.925) (-1.419,-0.875) (-1.411,-0.825) (-1.400,-0.775) (-1.388,-0.725) (-1.372,-0.675) (-1.352,-0.625) (-1.330,-0.575) (-1.304,-0.525) (-1.275,-0.477) (-1.244,-0.431)}
\EQ{(1.250,-0.405) (1.285,-0.450) (1.318,-0.500) (1.347,-0.550) (1.373,-0.600) (1.395,-0.650) (1.413,-0.700) (1.429,-0.750) (1.442,-0.800) (1.452,-0.850) (1.460,-0.900) (1.465,-0.950) (1.468,-1.000) (1.468,-1.050) (1.467,-1.100) (1.463,-1.150) (1.457,-1.200) (1.448,-1.250) (1.438,-1.300) (1.425,-1.350) (1.410,-1.400) (1.392,-1.450) (1.372,-1.500) (1.350,-1.550) (1.325,-1.600) (1.300,-1.645) (1.266,-1.700) (1.232,-1.750) (1.200,-1.793) (1.154,-1.850) (1.109,-1.900) (1.059,-1.950) (1.004,-2.000) (0.950,-2.044) (0.909,-2.075) (0.850,-2.116) (0.800,-2.147) (0.751,-2.175) (0.704,-2.200) (0.652,-2.225) (0.600,-2.248) (0.550,-2.268) (0.500,-2.285) (0.450,-2.301) (0.400,-2.315) (0.350,-2.328) (0.300,-2.338) (0.250,-2.347) (0.200,-2.354) (0.150,-2.360) (0.100,-2.364) (0.050,-2.366) (0.000,-2.367) (-0.050,-2.366) (-0.100,-2.364) (-0.150,-2.360) (-0.200,-2.354) (-0.250,-2.347) (-0.300,-2.338) (-0.350,-2.328) (-0.400,-2.315) (-0.450,-2.301) (-0.500,-2.285) (-0.550,-2.268) (-0.600,-2.248) (-0.650,-2.226) (-0.700,-2.202) (-0.750,-2.176) (-0.795,-2.150) (-0.850,-2.116) (-0.900,-2.081) (-0.942,-2.050) (-1.000,-2.003) (-1.050,-1.958) (-1.100,-1.909) (-1.150,-1.854) (-1.195,-1.800) (-1.225,-1.760) (-1.266,-1.700) (-1.297,-1.650) (-1.325,-1.600) (-1.350,-1.550) (-1.372,-1.500) (-1.392,-1.450) (-1.410,-1.400) (-1.425,-1.350) (-1.438,-1.300) (-1.448,-1.250) (-1.457,-1.200) (-1.463,-1.150) (-1.467,-1.100) (-1.468,-1.050) (-1.468,-1.000) (-1.465,-0.950) (-1.460,-0.900) (-1.452,-0.850) (-1.442,-0.800) (-1.429,-0.750) (-1.413,-0.700) (-1.395,-0.650) (-1.373,-0.600) (-1.350,-0.555) (-1.325,-0.512) (-1.285,-0.450) (-1.250,-0.405)}
\EQ{(1.251,-0.374) (1.295,-0.425) (1.325,-0.465) (1.365,-0.525) (1.393,-0.575) (1.418,-0.625) (1.440,-0.675) (1.458,-0.725) (1.473,-0.775) (1.486,-0.825) (1.496,-0.875) (1.504,-0.925) (1.510,-0.975) (1.513,-1.025) (1.514,-1.075) (1.512,-1.125) (1.509,-1.175) (1.503,-1.225) (1.495,-1.275) (1.486,-1.325) (1.473,-1.375) (1.459,-1.425) (1.442,-1.475) (1.423,-1.525) (1.402,-1.575) (1.378,-1.625) (1.352,-1.675) (1.325,-1.721) (1.290,-1.775) (1.255,-1.825) (1.216,-1.875) (1.175,-1.923) (1.127,-1.975) (1.075,-2.025) (1.025,-2.070) (0.975,-2.110) (0.925,-2.147) (0.875,-2.182) (0.825,-2.213) (0.775,-2.242) (0.725,-2.268) (0.675,-2.292) (0.625,-2.314) (0.575,-2.334) (0.525,-2.352) (0.475,-2.369) (0.425,-2.383) (0.375,-2.396) (0.325,-2.407) (0.275,-2.416) (0.225,-2.424) (0.175,-2.430) (0.125,-2.435) (0.075,-2.438) (0.025,-2.440) (-0.050,-2.439) (-0.100,-2.437) (-0.150,-2.433) (-0.200,-2.427) (-0.250,-2.420) (-0.300,-2.412) (-0.350,-2.402) (-0.400,-2.390) (-0.450,-2.376) (-0.500,-2.361) (-0.550,-2.344) (-0.599,-2.325) (-0.650,-2.304) (-0.700,-2.280) (-0.750,-2.255) (-0.800,-2.228) (-0.846,-2.200) (-0.900,-2.165) (-0.950,-2.129) (-1.000,-2.090) (-1.048,-2.050) (-1.100,-2.002) (-1.150,-1.951) (-1.195,-1.900) (-1.236,-1.850) (-1.273,-1.800) (-1.307,-1.750) (-1.337,-1.700) (-1.365,-1.650) (-1.390,-1.600) (-1.413,-1.550) (-1.433,-1.500) (-1.450,-1.453) (-1.466,-1.400) (-1.480,-1.350) (-1.491,-1.300) (-1.500,-1.250) (-1.506,-1.200) (-1.511,-1.150) (-1.513,-1.100) (-1.514,-1.050) (-1.512,-1.000) (-1.507,-0.950) (-1.501,-0.900) (-1.492,-0.850) (-1.480,-0.800) (-1.466,-0.750) (-1.449,-0.700) (-1.429,-0.650) (-1.406,-0.600) (-1.379,-0.550) (-1.350,-0.502) (-1.314,-0.450) (-1.275,-0.401) (-1.251,-0.374)}
\EQ{(1.266,-0.359) (1.305,-0.400) (1.346,-0.450) (1.382,-0.500) (1.414,-0.550) (1.442,-0.600) (1.467,-0.650) (1.488,-0.700) (1.506,-0.750) (1.522,-0.800) (1.534,-0.850) (1.545,-0.900) (1.553,-0.950) (1.559,-1.000) (1.562,-1.050) (1.563,-1.100) (1.563,-1.150) (1.560,-1.200) (1.555,-1.250) (1.547,-1.300) (1.538,-1.350) (1.527,-1.400) (1.513,-1.450) (1.498,-1.500) (1.480,-1.550) (1.460,-1.600) (1.437,-1.650) (1.412,-1.700) (1.385,-1.750) (1.354,-1.800) (1.325,-1.844) (1.285,-1.900) (1.250,-1.944) (1.201,-2.000) (1.154,-2.050) (1.101,-2.100) (1.050,-2.145) (1.000,-2.185) (0.950,-2.222) (0.908,-2.250) (0.850,-2.287) (0.800,-2.315) (0.750,-2.342) (0.700,-2.366) (0.650,-2.388) (0.600,-2.409) (0.550,-2.427) (0.500,-2.443) (0.450,-2.458) (0.400,-2.471) (0.350,-2.483) (0.300,-2.493) (0.250,-2.501) (0.200,-2.508) (0.150,-2.513) (0.100,-2.517) (0.050,-2.519) (0.000,-2.520) (-0.050,-2.519) (-0.100,-2.517) (-0.150,-2.513) (-0.200,-2.508) (-0.250,-2.501) (-0.300,-2.493) (-0.350,-2.483) (-0.400,-2.471) (-0.450,-2.458) (-0.500,-2.443) (-0.550,-2.427) (-0.600,-2.409) (-0.650,-2.388) (-0.700,-2.366) (-0.750,-2.342) (-0.800,-2.315) (-0.850,-2.287) (-0.900,-2.256) (-0.945,-2.225) (-1.000,-2.185) (-1.044,-2.150) (-1.100,-2.101) (-1.150,-2.054) (-1.200,-2.001) (-1.245,-1.950) (-1.285,-1.900) (-1.321,-1.850) (-1.350,-1.807) (-1.385,-1.750) (-1.412,-1.700) (-1.437,-1.650) (-1.460,-1.600) (-1.480,-1.550) (-1.498,-1.500) (-1.513,-1.450) (-1.527,-1.400) (-1.538,-1.350) (-1.547,-1.300) (-1.555,-1.250) (-1.560,-1.200) (-1.563,-1.150) (-1.563,-1.100) (-1.562,-1.050) (-1.559,-1.000) (-1.553,-0.950) (-1.545,-0.900) (-1.534,-0.850) (-1.522,-0.800) (-1.506,-0.750) (-1.488,-0.700) (-1.467,-0.650) (-1.442,-0.600) (-1.414,-0.550) (-1.382,-0.500) (-1.350,-0.456) (-1.305,-0.400) (-1.266,-0.359)}
\EQ{(1.275,-0.337) (1.314,-0.375) (1.350,-0.414) (1.400,-0.475) (1.436,-0.525) (1.467,-0.575) (1.495,-0.625) (1.519,-0.675) (1.540,-0.725) (1.559,-0.775) (1.574,-0.825) (1.587,-0.875) (1.598,-0.925) (1.606,-0.975) (1.612,-1.025) (1.616,-1.075) (1.618,-1.125) (1.618,-1.175) (1.616,-1.225) (1.611,-1.275) (1.605,-1.325) (1.597,-1.375) (1.586,-1.425) (1.574,-1.475) (1.559,-1.525) (1.543,-1.575) (1.524,-1.625) (1.503,-1.675) (1.479,-1.725) (1.454,-1.775) (1.425,-1.825) (1.394,-1.875) (1.360,-1.925) (1.325,-1.973) (1.283,-2.025) (1.250,-2.063) (1.216,-2.100) (1.175,-2.141) (1.125,-2.187) (1.081,-2.225) (1.025,-2.269) (0.983,-2.300) (0.925,-2.339) (0.875,-2.370) (0.825,-2.398) (0.775,-2.425) (0.725,-2.449) (0.675,-2.471) (0.625,-2.492) (0.575,-2.510) (0.525,-2.527) (0.475,-2.542) (0.425,-2.555) (0.375,-2.567) (0.325,-2.578) (0.275,-2.586) (0.225,-2.594) (0.175,-2.599) (0.125,-2.604) (0.075,-2.607) (0.025,-2.608) (-0.050,-2.608) (-0.100,-2.605) (-0.150,-2.602) (-0.200,-2.597) (-0.250,-2.590) (-0.300,-2.582) (-0.350,-2.573) (-0.400,-2.562) (-0.450,-2.549) (-0.500,-2.535) (-0.550,-2.519) (-0.600,-2.501) (-0.650,-2.482) (-0.700,-2.460) (-0.750,-2.437) (-0.800,-2.412) (-0.850,-2.384) (-0.900,-2.355) (-0.946,-2.325) (-1.000,-2.288) (-1.050,-2.250) (-1.100,-2.209) (-1.139,-2.175) (-1.191,-2.125) (-1.239,-2.075) (-1.275,-2.034) (-1.323,-1.975) (-1.360,-1.925) (-1.394,-1.875) (-1.425,-1.825) (-1.450,-1.781) (-1.475,-1.734) (-1.500,-1.681) (-1.524,-1.625) (-1.543,-1.575) (-1.559,-1.525) (-1.574,-1.475) (-1.586,-1.425) (-1.597,-1.375) (-1.605,-1.325) (-1.611,-1.275) (-1.616,-1.225) (-1.618,-1.175) (-1.618,-1.125) (-1.616,-1.075) (-1.612,-1.025) (-1.606,-0.975) (-1.598,-0.925) (-1.587,-0.875) (-1.574,-0.825) (-1.559,-0.775) (-1.540,-0.725) (-1.519,-0.675) (-1.495,-0.625) (-1.467,-0.575) (-1.436,-0.525) (-1.400,-0.475) (-1.360,-0.425) (-1.325,-0.387) (-1.289,-0.350) (-1.275,-0.337)}
\EQ{(1.275,-0.308) (1.323,-0.350) (1.374,-0.400) (1.419,-0.450) (1.450,-0.489) (1.493,-0.550) (1.524,-0.600) (1.550,-0.647) (1.575,-0.698) (1.597,-0.750) (1.616,-0.800) (1.632,-0.850) (1.645,-0.900) (1.656,-0.950) (1.665,-1.000) (1.672,-1.050) (1.676,-1.100) (1.679,-1.150) (1.679,-1.200) (1.678,-1.250) (1.674,-1.300) (1.669,-1.350) (1.662,-1.400) (1.652,-1.450) (1.641,-1.500) (1.628,-1.550) (1.613,-1.600) (1.595,-1.650) (1.576,-1.700) (1.554,-1.750) (1.530,-1.800) (1.504,-1.850) (1.475,-1.900) (1.444,-1.950) (1.410,-2.000) (1.375,-2.047) (1.333,-2.100) (1.289,-2.150) (1.250,-2.191) (1.200,-2.240) (1.162,-2.275) (1.102,-2.325) (1.050,-2.365) (1.001,-2.400) (0.950,-2.434) (0.900,-2.464) (0.850,-2.492) (0.800,-2.518) (0.750,-2.543) (0.700,-2.565) (0.650,-2.585) (0.600,-2.604) (0.550,-2.621) (0.500,-2.636) (0.450,-2.650) (0.400,-2.662) (0.350,-2.672) (0.300,-2.682) (0.250,-2.689) (0.200,-2.696) (0.150,-2.700) (0.100,-2.704) (0.050,-2.706) (0.000,-2.707) (-0.050,-2.706) (-0.100,-2.704) (-0.150,-2.700) (-0.200,-2.696) (-0.250,-2.689) (-0.300,-2.682) (-0.350,-2.672) (-0.400,-2.662) (-0.450,-2.650) (-0.500,-2.636) (-0.550,-2.621) (-0.600,-2.604) (-0.650,-2.585) (-0.700,-2.565) (-0.750,-2.543) (-0.800,-2.518) (-0.850,-2.492) (-0.900,-2.464) (-0.950,-2.434) (-1.000,-2.401) (-1.050,-2.365) (-1.100,-2.327) (-1.150,-2.285) (-1.189,-2.250) (-1.241,-2.200) (-1.289,-2.150) (-1.325,-2.109) (-1.373,-2.050) (-1.410,-2.000) (-1.444,-1.950) (-1.475,-1.900) (-1.504,-1.850) (-1.530,-1.800) (-1.554,-1.750) (-1.575,-1.702) (-1.595,-1.650) (-1.613,-1.600) (-1.628,-1.550) (-1.641,-1.500) (-1.652,-1.450) (-1.662,-1.400) (-1.669,-1.350) (-1.674,-1.300) (-1.678,-1.250) (-1.679,-1.200) (-1.679,-1.150) (-1.676,-1.100) (-1.672,-1.050) (-1.665,-1.000) (-1.656,-0.950) (-1.645,-0.900) (-1.632,-0.850) (-1.616,-0.800) (-1.597,-0.750) (-1.576,-0.700) (-1.552,-0.650) (-1.525,-0.601) (-1.493,-0.550) (-1.458,-0.500) (-1.425,-0.458) (-1.375,-0.401) (-1.325,-0.352) (-1.275,-0.308)}
\EQ{(1.275,-0.280) (1.325,-0.319) (1.375,-0.362) (1.414,-0.400) (1.450,-0.438) (1.500,-0.497) (1.539,-0.550) (1.572,-0.600) (1.600,-0.648) (1.625,-0.696) (1.650,-0.750) (1.670,-0.800) (1.688,-0.850) (1.703,-0.900) (1.716,-0.950) (1.726,-1.000) (1.734,-1.050) (1.741,-1.100) (1.745,-1.150) (1.747,-1.200) (1.748,-1.250) (1.746,-1.300) (1.743,-1.350) (1.738,-1.400) (1.731,-1.450) (1.722,-1.500) (1.711,-1.550) (1.698,-1.600) (1.683,-1.650) (1.667,-1.700) (1.648,-1.750) (1.627,-1.800) (1.604,-1.850) (1.579,-1.900) (1.552,-1.950) (1.525,-1.995) (1.489,-2.050) (1.454,-2.100) (1.416,-2.150) (1.375,-2.200) (1.330,-2.250) (1.282,-2.300) (1.229,-2.350) (1.175,-2.397) (1.125,-2.437) (1.075,-2.475) (1.025,-2.509) (0.975,-2.542) (0.925,-2.571) (0.875,-2.599) (0.825,-2.625) (0.775,-2.649) (0.725,-2.671) (0.675,-2.691) (0.625,-2.710) (0.575,-2.727) (0.525,-2.742) (0.475,-2.756) (0.425,-2.768) (0.375,-2.779) (0.325,-2.789) (0.275,-2.797) (0.225,-2.804) (0.175,-2.809) (0.125,-2.813) (0.075,-2.816) (0.025,-2.817) (-0.050,-2.817) (-0.100,-2.815) (-0.150,-2.811) (-0.200,-2.807) (-0.250,-2.800) (-0.300,-2.793) (-0.350,-2.784) (-0.400,-2.774) (-0.450,-2.762) (-0.500,-2.749) (-0.550,-2.735) (-0.600,-2.719) (-0.650,-2.701) (-0.700,-2.681) (-0.750,-2.660) (-0.800,-2.637) (-0.850,-2.612) (-0.900,-2.586) (-0.950,-2.557) (-1.000,-2.526) (-1.050,-2.492) (-1.100,-2.456) (-1.141,-2.425) (-1.200,-2.376) (-1.250,-2.331) (-1.300,-2.282) (-1.350,-2.228) (-1.396,-2.175) (-1.436,-2.125) (-1.472,-2.075) (-1.506,-2.025) (-1.537,-1.975) (-1.566,-1.925) (-1.592,-1.875) (-1.616,-1.825) (-1.638,-1.775) (-1.658,-1.725) (-1.675,-1.675) (-1.691,-1.625) (-1.705,-1.575) (-1.716,-1.525) (-1.726,-1.475) (-1.734,-1.425) (-1.741,-1.375) (-1.745,-1.325) (-1.747,-1.275) (-1.748,-1.225) (-1.746,-1.175) (-1.743,-1.125) (-1.738,-1.075) (-1.730,-1.025) (-1.721,-0.975) (-1.709,-0.925) (-1.695,-0.875) (-1.679,-0.825) (-1.660,-0.775) (-1.639,-0.725) (-1.614,-0.675) (-1.587,-0.625) (-1.556,-0.575) (-1.525,-0.530) (-1.482,-0.475) (-1.438,-0.425) (-1.400,-0.386) (-1.362,-0.350) (-1.301,-0.300) (-1.275,-0.280)}
\EQ{(1.288,-0.262) (1.342,-0.300) (1.400,-0.347) (1.450,-0.391) (1.500,-0.442) (1.550,-0.499) (1.590,-0.550) (1.625,-0.600) (1.656,-0.650) (1.684,-0.700) (1.708,-0.750) (1.730,-0.800) (1.750,-0.850) (1.767,-0.900) (1.781,-0.950) (1.794,-1.000) (1.804,-1.050) (1.812,-1.100) (1.818,-1.150) (1.823,-1.200) (1.825,-1.250) (1.825,-1.300) (1.824,-1.350) (1.821,-1.400) (1.816,-1.450) (1.810,-1.500) (1.801,-1.550) (1.791,-1.600) (1.779,-1.650) (1.765,-1.700) (1.750,-1.748) (1.732,-1.800) (1.712,-1.850) (1.691,-1.900) (1.667,-1.950) (1.641,-2.000) (1.613,-2.050) (1.583,-2.100) (1.550,-2.149) (1.514,-2.200) (1.476,-2.250) (1.434,-2.300) (1.400,-2.339) (1.366,-2.375) (1.325,-2.416) (1.289,-2.450) (1.250,-2.484) (1.201,-2.525) (1.150,-2.565) (1.101,-2.600) (1.050,-2.634) (1.000,-2.666) (0.950,-2.695) (0.900,-2.722) (0.850,-2.748) (0.800,-2.771) (0.750,-2.793) (0.700,-2.813) (0.650,-2.831) (0.600,-2.848) (0.550,-2.864) (0.500,-2.878) (0.450,-2.890) (0.400,-2.901) (0.350,-2.911) (0.300,-2.919) (0.250,-2.926) (0.200,-2.932) (0.150,-2.937) (0.100,-2.940) (0.050,-2.942) (0.000,-2.943) (-0.050,-2.942) (-0.100,-2.940) (-0.150,-2.937) (-0.200,-2.932) (-0.250,-2.926) (-0.300,-2.919) (-0.350,-2.911) (-0.400,-2.901) (-0.450,-2.890) (-0.500,-2.878) (-0.550,-2.864) (-0.600,-2.848) (-0.650,-2.831) (-0.700,-2.813) (-0.750,-2.793) (-0.800,-2.771) (-0.850,-2.748) (-0.900,-2.722) (-0.950,-2.695) (-1.000,-2.666) (-1.050,-2.634) (-1.100,-2.601) (-1.150,-2.565) (-1.200,-2.526) (-1.250,-2.484) (-1.289,-2.450) (-1.325,-2.416) (-1.366,-2.375) (-1.400,-2.339) (-1.434,-2.300) (-1.475,-2.251) (-1.514,-2.200) (-1.550,-2.150) (-1.583,-2.100) (-1.613,-2.050) (-1.641,-2.000) (-1.667,-1.950) (-1.691,-1.900) (-1.712,-1.850) (-1.732,-1.800) (-1.749,-1.750) (-1.765,-1.700) (-1.779,-1.650) (-1.791,-1.600) (-1.801,-1.550) (-1.810,-1.500) (-1.816,-1.450) (-1.821,-1.400) (-1.824,-1.350) (-1.825,-1.300) (-1.825,-1.250) (-1.823,-1.200) (-1.818,-1.150) (-1.812,-1.100) (-1.804,-1.050) (-1.794,-1.000) (-1.781,-0.950) (-1.767,-0.900) (-1.750,-0.851) (-1.730,-0.800) (-1.708,-0.750) (-1.684,-0.700) (-1.656,-0.650) (-1.625,-0.601) (-1.590,-0.550) (-1.551,-0.500) (-1.508,-0.450) (-1.459,-0.400) (-1.404,-0.350) (-1.350,-0.307) (-1.307,-0.275) (-1.288,-0.262)}
\EQ{(1.300,-0.243) (1.350,-0.275) (1.400,-0.310) (1.450,-0.348) (1.500,-0.391) (1.550,-0.439) (1.600,-0.493) (1.647,-0.550) (1.684,-0.600) (1.717,-0.650) (1.747,-0.700) (1.774,-0.750) (1.798,-0.800) (1.819,-0.850) (1.838,-0.900) (1.855,-0.950) (1.869,-1.000) (1.881,-1.050) (1.892,-1.100) (1.900,-1.150) (1.906,-1.200) (1.911,-1.250) (1.914,-1.300) (1.915,-1.350) (1.914,-1.400) (1.911,-1.450) (1.907,-1.500) (1.901,-1.550) (1.894,-1.600) (1.885,-1.650) (1.873,-1.700) (1.861,-1.750) (1.846,-1.800) (1.830,-1.850) (1.812,-1.900) (1.792,-1.950) (1.770,-2.000) (1.746,-2.050) (1.720,-2.100) (1.692,-2.150) (1.661,-2.200) (1.629,-2.250) (1.593,-2.300) (1.555,-2.350) (1.515,-2.400) (1.475,-2.445) (1.425,-2.498) (1.375,-2.548) (1.325,-2.593) (1.275,-2.636) (1.226,-2.675) (1.175,-2.713) (1.125,-2.747) (1.083,-2.775) (1.025,-2.810) (0.975,-2.838) (0.925,-2.865) (0.875,-2.890) (0.825,-2.913) (0.775,-2.934) (0.725,-2.954) (0.675,-2.972) (0.625,-2.989) (0.575,-3.004) (0.525,-3.018) (0.475,-3.030) (0.425,-3.042) (0.375,-3.052) (0.325,-3.060) (0.275,-3.068) (0.225,-3.074) (0.175,-3.079) (0.125,-3.082) (0.075,-3.085) (0.025,-3.086) (-0.050,-3.086) (-0.100,-3.084) (-0.150,-3.081) (-0.200,-3.076) (-0.250,-3.071) (-0.300,-3.064) (-0.350,-3.056) (-0.400,-3.047) (-0.450,-3.036) (-0.500,-3.024) (-0.550,-3.011) (-0.600,-2.996) (-0.650,-2.980) (-0.700,-2.963) (-0.750,-2.944) (-0.800,-2.923) (-0.850,-2.901) (-0.900,-2.877) (-0.950,-2.852) (-0.999,-2.825) (-1.050,-2.795) (-1.100,-2.764) (-1.150,-2.730) (-1.192,-2.700) (-1.250,-2.656) (-1.300,-2.615) (-1.345,-2.575) (-1.398,-2.525) (-1.448,-2.475) (-1.493,-2.425) (-1.535,-2.375) (-1.575,-2.325) (-1.611,-2.275) (-1.645,-2.225) (-1.675,-2.178) (-1.706,-2.125) (-1.733,-2.075) (-1.758,-2.025) (-1.781,-1.975) (-1.802,-1.925) (-1.821,-1.875) (-1.838,-1.825) (-1.854,-1.775) (-1.867,-1.725) (-1.879,-1.675) (-1.889,-1.625) (-1.898,-1.575) (-1.905,-1.525) (-1.910,-1.475) (-1.913,-1.425) (-1.915,-1.375) (-1.914,-1.325) (-1.913,-1.275) (-1.909,-1.225) (-1.903,-1.175) (-1.896,-1.125) (-1.887,-1.075) (-1.875,-1.025) (-1.862,-0.975) (-1.847,-0.925) (-1.829,-0.875) (-1.809,-0.825) (-1.786,-0.775) (-1.761,-0.725) (-1.733,-0.675) (-1.701,-0.625) (-1.666,-0.575) (-1.627,-0.525) (-1.584,-0.475) (-1.536,-0.425) (-1.482,-0.375) (-1.425,-0.328) (-1.375,-0.292) (-1.325,-0.258) (-1.300,-0.243)}
\EQ{(1.300,-0.216) (1.350,-0.244) (1.400,-0.274) (1.450,-0.308) (1.500,-0.345) (1.550,-0.385) (1.595,-0.425) (1.646,-0.475) (1.692,-0.525) (1.725,-0.565) (1.770,-0.625) (1.800,-0.669) (1.834,-0.725) (1.862,-0.775) (1.887,-0.825) (1.909,-0.875) (1.929,-0.925) (1.947,-0.975) (1.962,-1.025) (1.975,-1.075) (1.987,-1.125) (1.997,-1.175) (2.004,-1.225) (2.011,-1.275) (2.015,-1.325) (2.017,-1.375) (2.018,-1.425) (2.017,-1.475) (2.015,-1.525) (2.011,-1.575) (2.005,-1.625) (1.998,-1.675) (1.989,-1.725) (1.979,-1.775) (1.967,-1.825) (1.953,-1.875) (1.937,-1.925) (1.920,-1.975) (1.901,-2.025) (1.880,-2.075) (1.858,-2.125) (1.833,-2.175) (1.807,-2.225) (1.778,-2.275) (1.750,-2.321) (1.714,-2.375) (1.679,-2.425) (1.641,-2.475) (1.601,-2.525) (1.557,-2.575) (1.510,-2.625) (1.475,-2.661) (1.434,-2.700) (1.378,-2.750) (1.325,-2.794) (1.285,-2.825) (1.225,-2.869) (1.180,-2.900) (1.125,-2.935) (1.075,-2.965) (1.025,-2.993) (0.975,-3.020) (0.925,-3.045) (0.875,-3.068) (0.825,-3.089) (0.775,-3.109) (0.725,-3.128) (0.675,-3.145) (0.625,-3.161) (0.576,-3.175) (0.525,-3.188) (0.476,-3.200) (0.425,-3.211) (0.375,-3.220) (0.325,-3.229) (0.275,-3.236) (0.225,-3.241) (0.175,-3.246) (0.125,-3.249) (0.075,-3.252) (0.025,-3.253) (-0.050,-3.252) (-0.100,-3.251) (-0.150,-3.248) (-0.200,-3.244) (-0.250,-3.239) (-0.300,-3.232) (-0.350,-3.225) (-0.400,-3.216) (-0.450,-3.206) (-0.500,-3.195) (-0.550,-3.182) (-0.600,-3.168) (-0.650,-3.153) (-0.700,-3.137) (-0.750,-3.119) (-0.798,-3.100) (-0.850,-3.079) (-0.900,-3.056) (-0.950,-3.032) (-1.000,-3.007) (-1.050,-2.980) (-1.100,-2.950) (-1.150,-2.919) (-1.200,-2.886) (-1.250,-2.851) (-1.300,-2.814) (-1.348,-2.775) (-1.400,-2.731) (-1.450,-2.685) (-1.486,-2.650) (-1.525,-2.610) (-1.575,-2.555) (-1.621,-2.500) (-1.661,-2.450) (-1.697,-2.400) (-1.731,-2.350) (-1.763,-2.300) (-1.793,-2.250) (-1.820,-2.200) (-1.846,-2.150) (-1.869,-2.100) (-1.891,-2.050) (-1.911,-2.000) (-1.929,-1.950) (-1.945,-1.900) (-1.960,-1.850) (-1.973,-1.800) (-1.984,-1.750) (-1.994,-1.700) (-2.002,-1.650) (-2.008,-1.600) (-2.013,-1.550) (-2.017,-1.500) (-2.018,-1.450) (-2.018,-1.400) (-2.016,-1.350) (-2.013,-1.300) (-2.008,-1.250) (-2.001,-1.200) (-1.992,-1.150) (-1.982,-1.100) (-1.969,-1.050) (-1.955,-1.000) (-1.938,-0.950) (-1.919,-0.900) (-1.898,-0.850) (-1.875,-0.801) (-1.850,-0.753) (-1.825,-0.710) (-1.787,-0.650) (-1.752,-0.600) (-1.713,-0.550) (-1.675,-0.506) (-1.625,-0.454) (-1.575,-0.407) (-1.525,-0.364) (-1.475,-0.326) (-1.425,-0.291) (-1.375,-0.259) (-1.325,-0.230) (-1.300,-0.216)}
\EQ{(1.300,-0.191) (1.350,-0.215) (1.400,-0.241) (1.450,-0.270) (1.499,-0.300) (1.550,-0.335) (1.600,-0.372) (1.650,-0.414) (1.691,-0.450) (1.742,-0.500) (1.788,-0.550) (1.825,-0.594) (1.868,-0.650) (1.900,-0.696) (1.934,-0.750) (1.963,-0.800) (1.989,-0.850) (2.013,-0.900) (2.034,-0.950) (2.053,-1.000) (2.070,-1.050) (2.085,-1.100) (2.098,-1.150) (2.109,-1.200) (2.118,-1.250) (2.126,-1.300) (2.132,-1.350) (2.137,-1.400) (2.139,-1.450) (2.141,-1.500) (2.140,-1.550) (2.138,-1.600) (2.135,-1.650) (2.130,-1.700) (2.123,-1.750) (2.115,-1.800) (2.106,-1.850) (2.095,-1.900) (2.082,-1.950) (2.068,-2.000) (2.052,-2.050) (2.035,-2.100) (2.016,-2.150) (1.995,-2.200) (1.972,-2.250) (1.950,-2.296) (1.925,-2.344) (1.900,-2.389) (1.864,-2.450) (1.831,-2.500) (1.800,-2.546) (1.760,-2.600) (1.725,-2.645) (1.679,-2.700) (1.634,-2.750) (1.600,-2.786) (1.562,-2.825) (1.525,-2.860) (1.481,-2.900) (1.425,-2.947) (1.375,-2.987) (1.325,-3.025) (1.275,-3.060) (1.225,-3.093) (1.175,-3.124) (1.132,-3.150) (1.087,-3.175) (1.040,-3.200) (0.990,-3.225) (0.936,-3.250) (0.879,-3.275) (0.825,-3.297) (0.775,-3.315) (0.725,-3.333) (0.675,-3.349) (0.625,-3.364) (0.575,-3.377) (0.525,-3.390) (0.475,-3.401) (0.425,-3.411) (0.375,-3.420) (0.325,-3.427) (0.275,-3.434) (0.225,-3.439) (0.175,-3.444) (0.125,-3.447) (0.075,-3.449) (0.025,-3.450) (-0.037,-3.450) (-0.100,-3.448) (-0.150,-3.445) (-0.200,-3.442) (-0.250,-3.437) (-0.300,-3.431) (-0.350,-3.424) (-0.400,-3.415) (-0.450,-3.406) (-0.500,-3.395) (-0.550,-3.384) (-0.600,-3.371) (-0.650,-3.356) (-0.700,-3.341) (-0.748,-3.325) (-0.800,-3.306) (-0.850,-3.287) (-0.900,-3.266) (-0.950,-3.244) (-1.000,-3.220) (-1.050,-3.195) (-1.100,-3.168) (-1.150,-3.139) (-1.200,-3.109) (-1.250,-3.077) (-1.300,-3.043) (-1.350,-3.006) (-1.391,-2.975) (-1.450,-2.927) (-1.500,-2.883) (-1.550,-2.836) (-1.587,-2.800) (-1.625,-2.760) (-1.675,-2.705) (-1.721,-2.650) (-1.760,-2.600) (-1.797,-2.550) (-1.831,-2.500) (-1.864,-2.450) (-1.894,-2.400) (-1.922,-2.350) (-1.948,-2.300) (-1.972,-2.250) (-1.995,-2.200) (-2.016,-2.150) (-2.035,-2.100) (-2.052,-2.050) (-2.068,-2.000) (-2.082,-1.950) (-2.095,-1.900) (-2.106,-1.850) (-2.115,-1.800) (-2.123,-1.750) (-2.130,-1.700) (-2.135,-1.650) (-2.138,-1.600) (-2.140,-1.550) (-2.141,-1.500) (-2.139,-1.450) (-2.137,-1.400) (-2.132,-1.350) (-2.126,-1.300) (-2.118,-1.250) (-2.109,-1.200) (-2.098,-1.150) (-2.085,-1.100) (-2.070,-1.050) (-2.053,-1.000) (-2.034,-0.950) (-2.013,-0.900) (-1.989,-0.850) (-1.963,-0.800) (-1.934,-0.750) (-1.903,-0.700) (-1.868,-0.650) (-1.830,-0.600) (-1.788,-0.550) (-1.750,-0.508) (-1.700,-0.459) (-1.650,-0.414) (-1.603,-0.375) (-1.550,-0.335) (-1.500,-0.301) (-1.450,-0.270) (-1.400,-0.241) (-1.350,-0.215) (-1.300,-0.191)}
\EQ{(1.300,-0.165) (1.350,-0.186) (1.400,-0.208) (1.450,-0.233) (1.500,-0.259) (1.550,-0.288) (1.600,-0.319) (1.646,-0.350) (1.700,-0.390) (1.744,-0.425) (1.800,-0.474) (1.850,-0.522) (1.900,-0.575) (1.943,-0.625) (1.975,-0.665) (2.018,-0.725) (2.050,-0.773) (2.082,-0.825) (2.109,-0.875) (2.134,-0.925) (2.157,-0.975) (2.178,-1.025) (2.197,-1.075) (2.213,-1.125) (2.228,-1.175) (2.242,-1.225) (2.253,-1.275) (2.263,-1.325) (2.271,-1.375) (2.277,-1.425) (2.282,-1.475) (2.286,-1.525) (2.288,-1.575) (2.288,-1.625) (2.288,-1.675) (2.285,-1.725) (2.281,-1.775) (2.276,-1.825) (2.269,-1.875) (2.261,-1.925) (2.252,-1.975) (2.241,-2.025) (2.228,-2.075) (2.215,-2.125) (2.199,-2.175) (2.182,-2.225) (2.164,-2.275) (2.144,-2.325) (2.122,-2.375) (2.100,-2.423) (2.075,-2.473) (2.050,-2.520) (2.025,-2.565) (1.989,-2.625) (1.956,-2.675) (1.925,-2.720) (1.885,-2.775) (1.850,-2.820) (1.805,-2.875) (1.761,-2.925) (1.725,-2.964) (1.690,-3.000) (1.650,-3.039) (1.611,-3.075) (1.554,-3.125) (1.500,-3.170) (1.450,-3.209) (1.400,-3.245) (1.358,-3.275) (1.300,-3.313) (1.250,-3.344) (1.200,-3.373) (1.152,-3.400) (1.104,-3.425) (1.053,-3.450) (1.000,-3.475) (0.950,-3.497) (0.900,-3.517) (0.850,-3.536) (0.800,-3.554) (0.750,-3.571) (0.700,-3.587) (0.650,-3.601) (0.600,-3.614) (0.550,-3.626) (0.500,-3.637) (0.450,-3.647) (0.400,-3.656) (0.350,-3.664) (0.300,-3.670) (0.250,-3.676) (0.200,-3.680) (0.150,-3.684) (0.100,-3.687) (0.050,-3.688) (0.000,-3.689) (-0.050,-3.688) (-0.100,-3.687) (-0.150,-3.684) (-0.200,-3.680) (-0.250,-3.676) (-0.300,-3.670) (-0.350,-3.664) (-0.400,-3.656) (-0.450,-3.647) (-0.500,-3.637) (-0.550,-3.626) (-0.600,-3.614) (-0.650,-3.601) (-0.700,-3.587) (-0.750,-3.571) (-0.800,-3.554) (-0.850,-3.536) (-0.900,-3.517) (-0.950,-3.497) (-1.000,-3.475) (-1.050,-3.452) (-1.100,-3.427) (-1.150,-3.401) (-1.197,-3.375) (-1.250,-3.344) (-1.300,-3.313) (-1.350,-3.280) (-1.394,-3.250) (-1.450,-3.209) (-1.494,-3.175) (-1.550,-3.129) (-1.600,-3.085) (-1.638,-3.050) (-1.675,-3.015) (-1.714,-2.975) (-1.750,-2.937) (-1.800,-2.881) (-1.846,-2.825) (-1.885,-2.775) (-1.922,-2.725) (-1.956,-2.675) (-1.989,-2.625) (-2.019,-2.575) (-2.047,-2.525) (-2.074,-2.475) (-2.099,-2.425) (-2.122,-2.375) (-2.144,-2.325) (-2.164,-2.275) (-2.182,-2.225) (-2.199,-2.175) (-2.215,-2.125) (-2.228,-2.075) (-2.241,-2.025) (-2.252,-1.975) (-2.261,-1.925) (-2.269,-1.875) (-2.276,-1.825) (-2.281,-1.775) (-2.285,-1.725) (-2.288,-1.675) (-2.288,-1.625) (-2.288,-1.575) (-2.286,-1.525) (-2.282,-1.475) (-2.277,-1.425) (-2.271,-1.375) (-2.263,-1.325) (-2.253,-1.275) (-2.242,-1.225) (-2.228,-1.175) (-2.213,-1.125) (-2.197,-1.075) (-2.178,-1.025) (-2.157,-0.975) (-2.134,-0.925) (-2.109,-0.875) (-2.082,-0.825) (-2.051,-0.775) (-2.018,-0.725) (-1.982,-0.675) (-1.950,-0.634) (-1.900,-0.575) (-1.853,-0.525) (-1.801,-0.475) (-1.750,-0.430) (-1.700,-0.390) (-1.650,-0.353) (-1.600,-0.319) (-1.550,-0.288) (-1.500,-0.259) (-1.450,-0.233) (-1.400,-0.208) (-1.350,-0.186) (-1.300,-0.165)}
\EQ{(1.300,-0.141) (1.350,-0.158) (1.400,-0.177) (1.450,-0.197) (1.500,-0.220) (1.550,-0.243) (1.600,-0.269) (1.650,-0.297) (1.697,-0.325) (1.750,-0.359) (1.800,-0.394) (1.841,-0.425) (1.900,-0.473) (1.950,-0.517) (2.000,-0.566) (2.050,-0.619) (2.099,-0.675) (2.138,-0.725) (2.174,-0.775) (2.208,-0.825) (2.239,-0.875) (2.267,-0.925) (2.293,-0.975) (2.317,-1.025) (2.339,-1.075) (2.359,-1.125) (2.377,-1.175) (2.394,-1.225) (2.409,-1.275) (2.422,-1.325) (2.433,-1.375) (2.443,-1.425) (2.451,-1.475) (2.458,-1.525) (2.464,-1.575) (2.468,-1.625) (2.470,-1.675) (2.472,-1.725) (2.472,-1.775) (2.470,-1.825) (2.467,-1.875) (2.463,-1.925) (2.458,-1.975) (2.451,-2.025) (2.443,-2.075) (2.434,-2.125) (2.423,-2.175) (2.411,-2.225) (2.398,-2.275) (2.383,-2.325) (2.367,-2.375) (2.350,-2.423) (2.330,-2.475) (2.310,-2.525) (2.288,-2.575) (2.265,-2.625) (2.240,-2.675) (2.213,-2.725) (2.184,-2.775) (2.154,-2.825) (2.125,-2.871) (2.089,-2.925) (2.053,-2.975) (2.015,-3.025) (1.975,-3.075) (1.933,-3.125) (1.888,-3.175) (1.850,-3.215) (1.800,-3.265) (1.763,-3.300) (1.725,-3.335) (1.679,-3.375) (1.625,-3.420) (1.575,-3.459) (1.525,-3.496) (1.475,-3.531) (1.425,-3.565) (1.375,-3.597) (1.328,-3.625) (1.285,-3.650) (1.239,-3.675) (1.191,-3.700) (1.141,-3.725) (1.087,-3.750) (1.031,-3.775) (0.975,-3.798) (0.925,-3.817) (0.875,-3.835) (0.825,-3.852) (0.775,-3.868) (0.725,-3.883) (0.675,-3.897) (0.625,-3.909) (0.575,-3.921) (0.525,-3.932) (0.475,-3.941) (0.425,-3.950) (0.375,-3.958) (0.325,-3.964) (0.275,-3.970) (0.225,-3.975) (0.175,-3.978) (0.125,-3.981) (0.075,-3.983) (0.025,-3.984) (-0.050,-3.984) (-0.100,-3.982) (-0.150,-3.980) (-0.200,-3.977) (-0.250,-3.972) (-0.300,-3.967) (-0.350,-3.961) (-0.400,-3.954) (-0.450,-3.946) (-0.500,-3.937) (-0.550,-3.926) (-0.600,-3.915) (-0.650,-3.903) (-0.700,-3.890) (-0.750,-3.876) (-0.800,-3.860) (-0.850,-3.844) (-0.900,-3.826) (-0.950,-3.808) (-1.000,-3.788) (-1.050,-3.767) (-1.100,-3.744) (-1.150,-3.721) (-1.200,-3.696) (-1.250,-3.669) (-1.300,-3.641) (-1.350,-3.612) (-1.400,-3.581) (-1.448,-3.550) (-1.500,-3.514) (-1.550,-3.478) (-1.600,-3.440) (-1.650,-3.400) (-1.700,-3.357) (-1.750,-3.312) (-1.790,-3.275) (-1.840,-3.225) (-1.875,-3.189) (-1.910,-3.150) (-1.950,-3.105) (-1.995,-3.050) (-2.034,-3.000) (-2.071,-2.950) (-2.100,-2.908) (-2.139,-2.850) (-2.170,-2.800) (-2.199,-2.750) (-2.225,-2.703) (-2.250,-2.654) (-2.275,-2.603) (-2.299,-2.550) (-2.320,-2.500) (-2.340,-2.450) (-2.358,-2.400) (-2.375,-2.351) (-2.391,-2.300) (-2.405,-2.250) (-2.417,-2.200) (-2.429,-2.150) (-2.439,-2.100) (-2.447,-2.050) (-2.455,-2.000) (-2.461,-1.950) (-2.466,-1.900) (-2.469,-1.850) (-2.471,-1.800) (-2.472,-1.750) (-2.471,-1.700) (-2.469,-1.650) (-2.466,-1.600) (-2.461,-1.550) (-2.455,-1.500) (-2.447,-1.450) (-2.438,-1.400) (-2.427,-1.350) (-2.415,-1.300) (-2.401,-1.250) (-2.386,-1.200) (-2.369,-1.150) (-2.350,-1.101) (-2.329,-1.050) (-2.306,-1.000) (-2.281,-0.950) (-2.253,-0.900) (-2.225,-0.852) (-2.192,-0.800) (-2.157,-0.750) (-2.125,-0.708) (-2.078,-0.650) (-2.033,-0.600) (-1.984,-0.550) (-1.931,-0.500) (-1.875,-0.452) (-1.825,-0.413) (-1.775,-0.376) (-1.725,-0.343) (-1.675,-0.312) (-1.625,-0.283) (-1.575,-0.256) (-1.525,-0.231) (-1.475,-0.208) (-1.425,-0.187) (-1.375,-0.167) (-1.327,-0.150) (-1.300,-0.141)}
\EQ{(1.300,-0.117) (1.350,-0.131) (1.400,-0.146) (1.450,-0.163) (1.500,-0.181) (1.549,-0.200) (1.600,-0.221) (1.650,-0.244) (1.700,-0.268) (1.750,-0.293) (1.800,-0.321) (1.849,-0.350) (1.900,-0.382) (1.950,-0.416) (1.997,-0.450) (2.050,-0.492) (2.090,-0.525) (2.145,-0.575) (2.196,-0.625) (2.243,-0.675) (2.287,-0.725) (2.325,-0.772) (2.365,-0.825) (2.400,-0.875) (2.432,-0.925) (2.462,-0.975) (2.490,-1.025) (2.516,-1.075) (2.540,-1.125) (2.562,-1.175) (2.582,-1.225) (2.600,-1.275) (2.617,-1.325) (2.632,-1.375) (2.646,-1.425) (2.658,-1.475) (2.669,-1.525) (2.679,-1.575) (2.687,-1.625) (2.693,-1.675) (2.699,-1.725) (2.703,-1.775) (2.706,-1.825) (2.707,-1.875) (2.708,-1.925) (2.707,-1.975) (2.705,-2.025) (2.701,-2.075) (2.697,-2.125) (2.691,-2.175) (2.684,-2.225) (2.676,-2.275) (2.667,-2.325) (2.656,-2.375) (2.644,-2.425) (2.631,-2.475) (2.617,-2.525) (2.602,-2.575) (2.585,-2.625) (2.567,-2.675) (2.547,-2.725) (2.527,-2.775) (2.504,-2.825) (2.481,-2.875) (2.456,-2.925) (2.430,-2.975) (2.402,-3.025) (2.375,-3.070) (2.341,-3.125) (2.308,-3.175) (2.275,-3.223) (2.237,-3.275) (2.200,-3.323) (2.158,-3.375) (2.125,-3.414) (2.075,-3.470) (2.025,-3.523) (1.975,-3.574) (1.925,-3.621) (1.875,-3.667) (1.836,-3.700) (1.776,-3.750) (1.725,-3.790) (1.677,-3.825) (1.625,-3.862) (1.575,-3.896) (1.530,-3.925) (1.475,-3.959) (1.425,-3.988) (1.375,-4.016) (1.325,-4.043) (1.275,-4.068) (1.225,-4.092) (1.175,-4.115) (1.125,-4.137) (1.075,-4.158) (1.025,-4.177) (0.975,-4.196) (0.925,-4.213) (0.875,-4.229) (0.825,-4.245) (0.775,-4.259) (0.725,-4.272) (0.675,-4.285) (0.625,-4.296) (0.575,-4.307) (0.525,-4.317) (0.475,-4.325) (0.425,-4.333) (0.375,-4.340) (0.325,-4.346) (0.275,-4.351) (0.225,-4.356) (0.175,-4.359) (0.125,-4.362) (0.075,-4.363) (0.025,-4.364) (-0.050,-4.364) (-0.100,-4.363) (-0.150,-4.361) (-0.200,-4.358) (-0.250,-4.354) (-0.300,-4.349) (-0.350,-4.343) (-0.400,-4.337) (-0.450,-4.329) (-0.500,-4.321) (-0.550,-4.312) (-0.600,-4.302) (-0.650,-4.291) (-0.700,-4.279) (-0.750,-4.266) (-0.800,-4.252) (-0.850,-4.237) (-0.900,-4.221) (-0.950,-4.204) (-1.000,-4.187) (-1.050,-4.168) (-1.100,-4.147) (-1.150,-4.126) (-1.200,-4.104) (-1.250,-4.080) (-1.300,-4.056) (-1.350,-4.030) (-1.400,-4.002) (-1.448,-3.975) (-1.500,-3.944) (-1.550,-3.912) (-1.600,-3.879) (-1.642,-3.850) (-1.700,-3.808) (-1.744,-3.775) (-1.800,-3.730) (-1.850,-3.688) (-1.894,-3.650) (-1.948,-3.600) (-1.999,-3.550) (-2.047,-3.500) (-2.093,-3.450) (-2.137,-3.400) (-2.175,-3.354) (-2.218,-3.300) (-2.250,-3.257) (-2.291,-3.200) (-2.325,-3.150) (-2.357,-3.100) (-2.387,-3.050) (-2.416,-3.000) (-2.443,-2.950) (-2.469,-2.900) (-2.493,-2.850) (-2.516,-2.800) (-2.537,-2.750) (-2.557,-2.700) (-2.575,-2.653) (-2.593,-2.600) (-2.610,-2.550) (-2.624,-2.500) (-2.638,-2.450) (-2.650,-2.400) (-2.662,-2.350) (-2.672,-2.300) (-2.680,-2.250) (-2.688,-2.200) (-2.694,-2.150) (-2.699,-2.100) (-2.703,-2.050) (-2.706,-2.000) (-2.707,-1.950) (-2.708,-1.900) (-2.707,-1.850) (-2.704,-1.800) (-2.701,-1.750) (-2.696,-1.700) (-2.690,-1.650) (-2.683,-1.600) (-2.674,-1.550) (-2.664,-1.500) (-2.652,-1.450) (-2.639,-1.400) (-2.625,-1.351) (-2.609,-1.300) (-2.591,-1.250) (-2.572,-1.200) (-2.551,-1.150) (-2.528,-1.100) (-2.503,-1.050) (-2.476,-1.000) (-2.450,-0.954) (-2.416,-0.900) (-2.383,-0.850) (-2.350,-0.805) (-2.308,-0.750) (-2.275,-0.711) (-2.225,-0.655) (-2.175,-0.604) (-2.125,-0.557) (-2.075,-0.513) (-2.029,-0.475) (-1.975,-0.434) (-1.927,-0.400) (-1.875,-0.366) (-1.825,-0.335) (-1.775,-0.307) (-1.725,-0.280) (-1.675,-0.255) (-1.625,-0.232) (-1.575,-0.211) (-1.525,-0.191) (-1.475,-0.172) (-1.425,-0.155) (-1.375,-0.139) (-1.325,-0.124) (-1.300,-0.117)}
\EQ{(1.300,-0.093) (1.350,-0.104) (1.400,-0.116) (1.450,-0.130) (1.500,-0.144) (1.550,-0.159) (1.600,-0.175) (1.650,-0.193) (1.700,-0.211) (1.750,-0.231) (1.800,-0.252) (1.850,-0.275) (1.900,-0.299) (1.950,-0.324) (1.997,-0.350) (2.041,-0.375) (2.100,-0.412) (2.150,-0.444) (2.194,-0.475) (2.250,-0.517) (2.292,-0.550) (2.350,-0.599) (2.400,-0.645) (2.450,-0.693) (2.500,-0.746) (2.548,-0.800) (2.589,-0.850) (2.625,-0.896) (2.664,-0.950) (2.697,-1.000) (2.725,-1.044) (2.759,-1.100) (2.786,-1.150) (2.812,-1.200) (2.836,-1.250) (2.858,-1.300) (2.879,-1.350) (2.898,-1.400) (2.916,-1.450) (2.932,-1.500) (2.947,-1.550) (2.960,-1.600) (2.972,-1.650) (2.983,-1.700) (2.993,-1.750) (3.001,-1.800) (3.008,-1.850) (3.014,-1.900) (3.019,-1.950) (3.023,-2.000) (3.025,-2.050) (3.027,-2.100) (3.027,-2.150) (3.026,-2.200) (3.025,-2.250) (3.022,-2.300) (3.018,-2.350) (3.013,-2.400) (3.006,-2.450) (2.999,-2.500) (2.991,-2.550) (2.981,-2.600) (2.971,-2.650) (2.959,-2.700) (2.947,-2.750) (2.933,-2.800) (2.918,-2.850) (2.902,-2.900) (2.885,-2.950) (2.866,-3.000) (2.847,-3.050) (2.826,-3.100) (2.804,-3.150) (2.781,-3.200) (2.756,-3.250) (2.730,-3.300) (2.703,-3.350) (2.675,-3.399) (2.645,-3.450) (2.613,-3.500) (2.580,-3.550) (2.550,-3.594) (2.509,-3.650) (2.475,-3.696) (2.432,-3.750) (2.400,-3.789) (2.350,-3.846) (2.301,-3.900) (2.253,-3.950) (2.203,-4.000) (2.151,-4.050) (2.100,-4.096) (2.050,-4.139) (2.007,-4.175) (1.950,-4.220) (1.900,-4.258) (1.850,-4.295) (1.807,-4.325) (1.750,-4.363) (1.700,-4.395) (1.651,-4.425) (1.600,-4.455) (1.550,-4.483) (1.500,-4.510) (1.450,-4.536) (1.400,-4.561) (1.350,-4.585) (1.300,-4.607) (1.250,-4.629) (1.200,-4.649) (1.150,-4.669) (1.100,-4.687) (1.050,-4.705) (1.000,-4.722) (0.950,-4.737) (0.900,-4.752) (0.850,-4.766) (0.800,-4.779) (0.750,-4.792) (0.700,-4.803) (0.650,-4.814) (0.600,-4.824) (0.550,-4.833) (0.500,-4.841) (0.450,-4.848) (0.400,-4.855) (0.350,-4.861) (0.300,-4.866) (0.250,-4.870) (0.200,-4.873) (0.150,-4.876) (0.100,-4.878) (0.050,-4.879) (0.000,-4.880) (-0.050,-4.879) (-0.100,-4.878) (-0.150,-4.876) (-0.200,-4.873) (-0.250,-4.870) (-0.300,-4.866) (-0.350,-4.861) (-0.400,-4.855) (-0.450,-4.848) (-0.500,-4.841) (-0.550,-4.833) (-0.600,-4.824) (-0.650,-4.814) (-0.700,-4.803) (-0.750,-4.792) (-0.800,-4.779) (-0.850,-4.766) (-0.900,-4.752) (-0.950,-4.737) (-1.000,-4.722) (-1.050,-4.705) (-1.100,-4.687) (-1.150,-4.669) (-1.198,-4.650) (-1.250,-4.629) (-1.300,-4.607) (-1.350,-4.585) (-1.400,-4.561) (-1.450,-4.536) (-1.500,-4.510) (-1.550,-4.483) (-1.600,-4.455) (-1.650,-4.426) (-1.700,-4.395) (-1.750,-4.363) (-1.800,-4.330) (-1.843,-4.300) (-1.900,-4.258) (-1.944,-4.225) (-2.000,-4.181) (-2.050,-4.139) (-2.095,-4.100) (-2.150,-4.051) (-2.200,-4.003) (-2.250,-3.954) (-2.300,-3.901) (-2.347,-3.850) (-2.390,-3.800) (-2.425,-3.758) (-2.472,-3.700) (-2.509,-3.650) (-2.546,-3.600) (-2.575,-3.558) (-2.613,-3.500) (-2.645,-3.450) (-2.675,-3.400) (-2.700,-3.356) (-2.725,-3.310) (-2.750,-3.262) (-2.775,-3.212) (-2.800,-3.159) (-2.825,-3.102) (-2.847,-3.050) (-2.866,-3.000) (-2.885,-2.950) (-2.902,-2.900) (-2.918,-2.850) (-2.933,-2.800) (-2.947,-2.750) (-2.959,-2.700) (-2.971,-2.650) (-2.981,-2.600) (-2.991,-2.550) (-2.999,-2.500) (-3.006,-2.450) (-3.013,-2.400) (-3.018,-2.350) (-3.022,-2.300) (-3.025,-2.250) (-3.026,-2.200) (-3.027,-2.150) (-3.027,-2.100) (-3.025,-2.050) (-3.023,-2.000) (-3.019,-1.950) (-3.014,-1.900) (-3.008,-1.850) (-3.001,-1.800) (-2.993,-1.750) (-2.983,-1.700) (-2.972,-1.650) (-2.960,-1.600) (-2.947,-1.550) (-2.932,-1.500) (-2.916,-1.450) (-2.898,-1.400) (-2.879,-1.350) (-2.858,-1.300) (-2.836,-1.250) (-2.812,-1.200) (-2.786,-1.150) (-2.759,-1.100) (-2.729,-1.050) (-2.700,-1.004) (-2.664,-0.950) (-2.628,-0.900) (-2.589,-0.850) (-2.550,-0.802) (-2.504,-0.750) (-2.457,-0.700) (-2.406,-0.650) (-2.351,-0.600) (-2.300,-0.557) (-2.261,-0.525) (-2.200,-0.479) (-2.158,-0.450) (-2.100,-0.412) (-2.050,-0.381) (-2.000,-0.352) (-1.951,-0.325) (-1.902,-0.300) (-1.850,-0.275) (-1.800,-0.252) (-1.750,-0.231) (-1.700,-0.211) (-1.650,-0.193) (-1.600,-0.175) (-1.550,-0.159) (-1.500,-0.144) (-1.450,-0.130) (-1.400,-0.116) (-1.350,-0.104) (-1.300,-0.093)}
\EQ{(1.300,-0.070) (1.350,-0.078) (1.400,-0.087) (1.450,-0.097) (1.500,-0.107) (1.550,-0.118) (1.600,-0.130) (1.650,-0.143) (1.700,-0.157) (1.750,-0.171) (1.800,-0.187) (1.850,-0.203) (1.900,-0.220) (1.950,-0.239) (2.000,-0.258) (2.050,-0.279) (2.098,-0.300) (2.150,-0.324) (2.200,-0.348) (2.250,-0.374) (2.298,-0.400) (2.342,-0.425) (2.400,-0.460) (2.450,-0.492) (2.500,-0.525) (2.550,-0.561) (2.600,-0.598) (2.650,-0.638) (2.694,-0.675) (2.750,-0.724) (2.800,-0.772) (2.850,-0.822) (2.899,-0.875) (2.943,-0.925) (2.985,-0.975) (3.024,-1.025) (3.060,-1.075) (3.095,-1.125) (3.125,-1.171) (3.159,-1.225) (3.188,-1.275) (3.216,-1.325) (3.241,-1.375) (3.266,-1.425) (3.289,-1.475) (3.310,-1.525) (3.330,-1.575) (3.349,-1.625) (3.367,-1.675) (3.383,-1.725) (3.398,-1.775) (3.412,-1.825) (3.425,-1.875) (3.437,-1.925) (3.447,-1.975) (3.456,-2.025) (3.465,-2.075) (3.472,-2.125) (3.479,-2.175) (3.484,-2.225) (3.488,-2.275) (3.491,-2.325) (3.494,-2.375) (3.495,-2.425) (3.496,-2.475) (3.495,-2.525) (3.494,-2.575) (3.491,-2.625) (3.488,-2.675) (3.483,-2.725) (3.478,-2.775) (3.472,-2.825) (3.465,-2.875) (3.457,-2.925) (3.448,-2.975) (3.438,-3.025) (3.427,-3.075) (3.415,-3.125) (3.403,-3.175) (3.389,-3.225) (3.374,-3.275) (3.359,-3.325) (3.342,-3.375) (3.325,-3.423) (3.306,-3.475) (3.286,-3.525) (3.265,-3.575) (3.243,-3.625) (3.220,-3.675) (3.196,-3.725) (3.171,-3.775) (3.145,-3.825) (3.118,-3.875) (3.089,-3.925) (3.059,-3.975) (3.028,-4.025) (3.000,-4.068) (2.962,-4.125) (2.927,-4.175) (2.890,-4.225) (2.852,-4.275) (2.812,-4.325) (2.775,-4.370) (2.728,-4.425) (2.683,-4.475) (2.637,-4.525) (2.600,-4.563) (2.563,-4.600) (2.525,-4.637) (2.485,-4.675) (2.430,-4.725) (2.375,-4.773) (2.325,-4.815) (2.281,-4.850) (2.225,-4.894) (2.183,-4.925) (2.125,-4.967) (2.078,-5.000) (2.025,-5.035) (1.975,-5.067) (1.925,-5.098) (1.881,-5.125) (1.825,-5.157) (1.775,-5.185) (1.725,-5.212) (1.675,-5.238) (1.625,-5.262) (1.575,-5.286) (1.525,-5.309) (1.475,-5.331) (1.425,-5.352) (1.375,-5.373) (1.325,-5.392) (1.275,-5.411) (1.225,-5.428) (1.175,-5.445) (1.125,-5.462) (1.075,-5.477) (1.025,-5.492) (0.975,-5.505) (0.925,-5.519) (0.875,-5.531) (0.825,-5.543) (0.775,-5.553) (0.725,-5.564) (0.675,-5.573) (0.625,-5.582) (0.575,-5.590) (0.525,-5.597) (0.475,-5.604) (0.425,-5.610) (0.375,-5.616) (0.325,-5.620) (0.275,-5.624) (0.225,-5.628) (0.175,-5.630) (0.125,-5.632) (0.075,-5.634) (0.025,-5.634) (-0.050,-5.634) (-0.100,-5.633) (-0.150,-5.631) (-0.200,-5.629) (-0.250,-5.626) (-0.300,-5.622) (-0.350,-5.618) (-0.400,-5.613) (-0.450,-5.607) (-0.500,-5.601) (-0.550,-5.594) (-0.600,-5.586) (-0.650,-5.578) (-0.700,-5.569) (-0.750,-5.559) (-0.800,-5.548) (-0.850,-5.537) (-0.899,-5.525) (-0.950,-5.512) (-1.000,-5.499) (-1.050,-5.484) (-1.100,-5.469) (-1.150,-5.453) (-1.200,-5.437) (-1.250,-5.420) (-1.300,-5.401) (-1.350,-5.382) (-1.400,-5.363) (-1.450,-5.342) (-1.500,-5.320) (-1.550,-5.298) (-1.599,-5.275) (-1.650,-5.250) (-1.700,-5.225) (-1.747,-5.200) (-1.793,-5.175) (-1.838,-5.150) (-1.900,-5.114) (-1.950,-5.083) (-2.000,-5.051) (-2.050,-5.019) (-2.100,-4.985) (-2.149,-4.950) (-2.200,-4.913) (-2.249,-4.875) (-2.300,-4.835) (-2.343,-4.800) (-2.400,-4.751) (-2.450,-4.707) (-2.500,-4.661) (-2.538,-4.625) (-2.575,-4.588) (-2.613,-4.550) (-2.650,-4.511) (-2.700,-4.457) (-2.750,-4.400) (-2.792,-4.350) (-2.825,-4.309) (-2.871,-4.250) (-2.908,-4.200) (-2.944,-4.150) (-2.975,-4.106) (-3.012,-4.050) (-3.044,-4.000) (-3.074,-3.950) (-3.100,-3.906) (-3.125,-3.862) (-3.150,-3.816) (-3.175,-3.768) (-3.200,-3.718) (-3.225,-3.665) (-3.250,-3.610) (-3.275,-3.551) (-3.296,-3.500) (-3.315,-3.450) (-3.333,-3.400) (-3.350,-3.351) (-3.366,-3.300) (-3.382,-3.250) (-3.396,-3.200) (-3.409,-3.150) (-3.421,-3.100) (-3.433,-3.050) (-3.443,-3.000) (-3.453,-2.950) (-3.461,-2.900) (-3.469,-2.850) (-3.475,-2.800) (-3.481,-2.750) (-3.486,-2.700) (-3.490,-2.650) (-3.492,-2.600) (-3.494,-2.550) (-3.495,-2.500) (-3.496,-2.450) (-3.495,-2.400) (-3.493,-2.350) (-3.490,-2.300) (-3.486,-2.250) (-3.481,-2.200) (-3.475,-2.150) (-3.469,-2.100) (-3.461,-2.050) (-3.452,-2.000) (-3.442,-1.950) (-3.431,-1.900) (-3.419,-1.850) (-3.405,-1.800) (-3.391,-1.750) (-3.375,-1.700) (-3.358,-1.650) (-3.340,-1.600) (-3.321,-1.550) (-3.300,-1.501) (-3.277,-1.450) (-3.254,-1.400) (-3.229,-1.350) (-3.202,-1.300) (-3.175,-1.252) (-3.144,-1.200) (-3.112,-1.150) (-3.078,-1.100) (-3.042,-1.050) (-3.004,-1.000) (-2.964,-0.950) (-2.925,-0.904) (-2.876,-0.850) (-2.828,-0.800) (-2.777,-0.750) (-2.725,-0.702) (-2.675,-0.659) (-2.634,-0.625) (-2.575,-0.579) (-2.525,-0.543) (-2.475,-0.508) (-2.425,-0.475) (-2.375,-0.444) (-2.325,-0.415) (-2.275,-0.387) (-2.225,-0.361) (-2.175,-0.336) (-2.125,-0.312) (-2.075,-0.290) (-2.025,-0.269) (-1.975,-0.248) (-1.925,-0.230) (-1.875,-0.212) (-1.825,-0.195) (-1.775,-0.179) (-1.725,-0.164) (-1.676,-0.150) (-1.625,-0.137) (-1.575,-0.124) (-1.525,-0.113) (-1.475,-0.102) (-1.425,-0.092) (-1.375,-0.082) (-1.325,-0.074) (-1.300,-0.070)}
\EQ{(1.300,-0.046) (1.350,-0.052) (1.400,-0.058) (1.450,-0.064) (1.500,-0.071) (1.550,-0.079) (1.600,-0.086) (1.650,-0.095) (1.700,-0.104) (1.750,-0.113) (1.800,-0.123) (1.850,-0.134) (1.900,-0.145) (1.950,-0.157) (2.000,-0.170) (2.050,-0.183) (2.100,-0.197) (2.150,-0.212) (2.200,-0.227) (2.250,-0.243) (2.300,-0.260) (2.350,-0.278) (2.400,-0.297) (2.450,-0.317) (2.500,-0.337) (2.550,-0.359) (2.600,-0.381) (2.650,-0.405) (2.700,-0.429) (2.750,-0.455) (2.800,-0.482) (2.850,-0.510) (2.900,-0.539) (2.950,-0.570) (2.998,-0.600) (3.050,-0.635) (3.100,-0.670) (3.141,-0.700) (3.200,-0.745) (3.250,-0.785) (3.298,-0.825) (3.350,-0.871) (3.400,-0.917) (3.450,-0.966) (3.500,-1.017) (3.550,-1.071) (3.598,-1.125) (3.640,-1.175) (3.675,-1.219) (3.718,-1.275) (3.750,-1.319) (3.789,-1.375) (3.822,-1.425) (3.850,-1.470) (3.883,-1.525) (3.912,-1.575) (3.939,-1.625) (3.965,-1.675) (3.990,-1.725) (4.014,-1.775) (4.036,-1.825) (4.057,-1.875) (4.077,-1.925) (4.096,-1.975) (4.114,-2.025) (4.131,-2.075) (4.147,-2.125) (4.162,-2.175) (4.176,-2.225) (4.189,-2.275) (4.201,-2.325) (4.213,-2.375) (4.223,-2.425) (4.232,-2.475) (4.241,-2.525) (4.249,-2.575) (4.256,-2.625) (4.262,-2.675) (4.267,-2.725) (4.271,-2.775) (4.275,-2.825) (4.278,-2.875) (4.280,-2.925) (4.281,-2.975) (4.281,-3.025) (4.281,-3.075) (4.280,-3.125) (4.278,-3.175) (4.275,-3.225) (4.272,-3.275) (4.268,-3.325) (4.263,-3.375) (4.257,-3.425) (4.250,-3.475) (4.243,-3.525) (4.235,-3.575) (4.226,-3.625) (4.217,-3.675) (4.206,-3.725) (4.195,-3.775) (4.184,-3.825) (4.171,-3.875) (4.157,-3.925) (4.143,-3.975) (4.128,-4.025) (4.112,-4.075) (4.096,-4.125) (4.078,-4.175) (4.060,-4.225) (4.041,-4.275) (4.021,-4.325) (4.000,-4.375) (3.979,-4.425) (3.956,-4.475) (3.933,-4.525) (3.908,-4.575) (3.883,-4.625) (3.857,-4.675) (3.830,-4.725) (3.802,-4.775) (3.775,-4.821) (3.743,-4.875) (3.711,-4.925) (3.679,-4.975) (3.650,-5.019) (3.611,-5.075) (3.576,-5.125) (3.539,-5.175) (3.501,-5.225) (3.462,-5.275) (3.425,-5.321) (3.380,-5.375) (3.336,-5.425) (3.300,-5.466) (3.250,-5.520) (3.200,-5.573) (3.150,-5.623) (3.100,-5.672) (3.050,-5.720) (3.000,-5.766) (2.961,-5.800) (2.904,-5.850) (2.850,-5.895) (2.800,-5.935) (2.750,-5.974) (2.700,-6.013) (2.650,-6.049) (2.600,-6.085) (2.550,-6.120) (2.505,-6.150) (2.450,-6.186) (2.400,-6.218) (2.350,-6.249) (2.306,-6.275) (2.250,-6.308) (2.200,-6.336) (2.150,-6.363) (2.100,-6.390) (2.050,-6.415) (2.000,-6.440) (1.950,-6.464) (1.900,-6.488) (1.850,-6.510) (1.800,-6.532) (1.750,-6.553) (1.700,-6.574) (1.650,-6.593) (1.600,-6.612) (1.550,-6.631) (1.500,-6.648) (1.450,-6.665) (1.400,-6.682) (1.350,-6.698) (1.300,-6.713) (1.250,-6.727) (1.200,-6.741) (1.150,-6.754) (1.100,-6.767) (1.050,-6.779) (1.000,-6.791) (0.950,-6.801) (0.900,-6.812) (0.850,-6.821) (0.800,-6.831) (0.750,-6.839) (0.700,-6.847) (0.650,-6.854) (0.600,-6.861) (0.550,-6.868) (0.500,-6.873) (0.450,-6.879) (0.400,-6.883) (0.350,-6.887) (0.300,-6.891) (0.250,-6.894) (0.200,-6.896) (0.150,-6.898) (0.100,-6.900) (0.050,-6.900) (0.000,-6.901) (-0.050,-6.900) (-0.100,-6.900) (-0.150,-6.898) (-0.200,-6.896) (-0.250,-6.894) (-0.300,-6.891) (-0.350,-6.887) (-0.400,-6.883) (-0.450,-6.879) (-0.500,-6.873) (-0.550,-6.868) (-0.600,-6.861) (-0.650,-6.854) (-0.700,-6.847) (-0.750,-6.839) (-0.800,-6.831) (-0.850,-6.821) (-0.900,-6.812) (-0.950,-6.801) (-1.000,-6.791) (-1.050,-6.779) (-1.100,-6.767) (-1.150,-6.754) (-1.200,-6.741) (-1.250,-6.727) (-1.300,-6.713) (-1.350,-6.698) (-1.400,-6.682) (-1.450,-6.665) (-1.500,-6.648) (-1.550,-6.631) (-1.600,-6.612) (-1.650,-6.593) (-1.697,-6.575) (-1.750,-6.553) (-1.800,-6.532) (-1.850,-6.510) (-1.900,-6.488) (-1.950,-6.464) (-2.000,-6.440) (-2.050,-6.415) (-2.100,-6.390) (-2.150,-6.363) (-2.200,-6.336) (-2.250,-6.308) (-2.300,-6.279) (-2.348,-6.250) (-2.400,-6.218) (-2.450,-6.186) (-2.500,-6.154) (-2.543,-6.125) (-2.600,-6.085) (-2.649,-6.050) (-2.700,-6.013) (-2.749,-5.975) (-2.800,-5.935) (-2.844,-5.900) (-2.900,-5.853) (-2.950,-5.810) (-2.990,-5.775) (-3.044,-5.725) (-3.097,-5.675) (-3.148,-5.625) (-3.198,-5.575) (-3.246,-5.525) (-3.292,-5.475) (-3.336,-5.425) (-3.375,-5.380) (-3.421,-5.325) (-3.462,-5.275) (-3.500,-5.227) (-3.539,-5.175) (-3.575,-5.126) (-3.611,-5.075) (-3.646,-5.025) (-3.675,-4.981) (-3.711,-4.925) (-3.743,-4.875) (-3.773,-4.825) (-3.800,-4.778) (-3.825,-4.734) (-3.850,-4.688) (-3.875,-4.641) (-3.900,-4.592) (-3.925,-4.541) (-3.950,-4.488) (-3.975,-4.433) (-4.000,-4.376) (-4.021,-4.325) (-4.041,-4.275) (-4.060,-4.225) (-4.078,-4.175) (-4.096,-4.125) (-4.112,-4.075) (-4.128,-4.025) (-4.143,-3.975) (-4.157,-3.925) (-4.171,-3.875) (-4.184,-3.825) (-4.195,-3.775) (-4.206,-3.725) (-4.217,-3.675) (-4.226,-3.625) (-4.235,-3.575) (-4.243,-3.525) (-4.250,-3.475) (-4.257,-3.425) (-4.263,-3.375) (-4.268,-3.325) (-4.272,-3.275) (-4.275,-3.225) (-4.278,-3.175) (-4.280,-3.125) (-4.281,-3.075) (-4.281,-3.025) (-4.281,-2.975) (-4.280,-2.925) (-4.278,-2.875) (-4.275,-2.825) (-4.271,-2.775) (-4.267,-2.725) (-4.262,-2.675) (-4.256,-2.625) (-4.249,-2.575) (-4.241,-2.525) (-4.232,-2.475) (-4.223,-2.425) (-4.213,-2.375) (-4.201,-2.325) (-4.189,-2.275) (-4.176,-2.225) (-4.162,-2.175) (-4.147,-2.125) (-4.131,-2.075) (-4.114,-2.025) (-4.096,-1.975) (-4.077,-1.925) (-4.057,-1.875) (-4.036,-1.825) (-4.014,-1.775) (-3.990,-1.725) (-3.965,-1.675) (-3.939,-1.625) (-3.912,-1.575) (-3.883,-1.525) (-3.853,-1.475) (-3.825,-1.430) (-3.789,-1.375) (-3.754,-1.325) (-3.718,-1.275) (-3.680,-1.225) (-3.640,-1.175) (-3.600,-1.127) (-3.554,-1.075) (-3.508,-1.025) (-3.459,-0.975) (-3.408,-0.925) (-3.355,-0.875) (-3.300,-0.827) (-3.250,-0.785) (-3.207,-0.750) (-3.150,-0.707) (-3.107,-0.675) (-3.050,-0.635) (-3.000,-0.602) (-2.950,-0.570) (-2.900,-0.539) (-2.850,-0.510) (-2.800,-0.482) (-2.750,-0.455) (-2.700,-0.429) (-2.650,-0.405) (-2.600,-0.381) (-2.550,-0.359) (-2.500,-0.337) (-2.450,-0.317) (-2.400,-0.297) (-2.350,-0.278) (-2.300,-0.260) (-2.250,-0.243) (-2.200,-0.227) (-2.150,-0.212) (-2.100,-0.197) (-2.050,-0.183) (-2.000,-0.170) (-1.950,-0.157) (-1.900,-0.145) (-1.850,-0.134) (-1.800,-0.123) (-1.750,-0.113) (-1.700,-0.104) (-1.650,-0.095) (-1.600,-0.086) (-1.550,-0.079) (-1.500,-0.071) (-1.450,-0.064) (-1.400,-0.058) (-1.350,-0.052) (-1.300,-0.046)}
\EQ{(1.300,-0.023) (1.350,-0.026) (1.400,-0.029) (1.450,-0.032) (1.500,-0.035) (1.550,-0.039) (1.600,-0.043) (1.650,-0.047) (1.700,-0.052) (1.750,-0.056) (1.800,-0.061) (1.850,-0.067) (1.900,-0.072) (1.950,-0.078) (2.000,-0.084) (2.050,-0.091) (2.100,-0.098) (2.150,-0.105) (2.200,-0.112) (2.250,-0.120) (2.300,-0.128) (2.350,-0.137) (2.400,-0.146) (2.450,-0.155) (2.500,-0.165) (2.550,-0.175) (2.600,-0.186) (2.650,-0.197) (2.700,-0.209) (2.750,-0.220) (2.800,-0.233) (2.850,-0.246) (2.900,-0.259) (2.950,-0.273) (3.000,-0.287) (3.050,-0.302) (3.100,-0.318) (3.150,-0.334) (3.199,-0.350) (3.250,-0.367) (3.300,-0.385) (3.350,-0.403) (3.400,-0.422) (3.450,-0.442) (3.500,-0.462) (3.550,-0.483) (3.600,-0.504) (3.646,-0.525) (3.700,-0.550) (3.750,-0.573) (3.800,-0.598) (3.850,-0.623) (3.900,-0.649) (3.949,-0.675) (3.994,-0.700) (4.038,-0.725) (4.100,-0.761) (4.150,-0.792) (4.200,-0.823) (4.250,-0.856) (4.300,-0.889) (4.350,-0.923) (4.400,-0.959) (4.450,-0.996) (4.500,-1.034) (4.550,-1.073) (4.600,-1.113) (4.644,-1.150) (4.700,-1.198) (4.750,-1.243) (4.800,-1.289) (4.838,-1.325) (4.875,-1.362) (4.913,-1.400) (4.950,-1.438) (4.985,-1.475) (5.025,-1.519) (5.074,-1.575) (5.117,-1.625) (5.150,-1.665) (5.198,-1.725) (5.237,-1.775) (5.274,-1.825) (5.310,-1.875) (5.345,-1.925) (5.375,-1.970) (5.411,-2.025) (5.443,-2.075) (5.473,-2.125) (5.500,-2.171) (5.525,-2.214) (5.550,-2.259) (5.575,-2.306) (5.598,-2.350) (5.623,-2.400) (5.648,-2.450) (5.671,-2.500) (5.694,-2.550) (5.716,-2.600) (5.737,-2.650) (5.757,-2.700) (5.777,-2.750) (5.796,-2.800) (5.814,-2.850) (5.831,-2.900) (5.848,-2.950) (5.863,-3.000) (5.879,-3.050) (5.893,-3.100) (5.907,-3.150) (5.920,-3.200) (5.933,-3.250) (5.944,-3.300) (5.956,-3.350) (5.966,-3.400) (5.976,-3.450) (5.985,-3.500) (5.994,-3.550) (6.002,-3.600) (6.010,-3.650) (6.017,-3.700) (6.023,-3.750) (6.029,-3.800) (6.034,-3.850) (6.038,-3.900) (6.042,-3.950) (6.046,-4.000) (6.049,-4.050) (6.051,-4.100) (6.053,-4.150) (6.054,-4.200) (6.054,-4.250) (6.054,-4.300) (6.054,-4.350) (6.053,-4.400) (6.051,-4.450) (6.049,-4.500) (6.047,-4.550) (6.043,-4.600) (6.040,-4.650) (6.035,-4.700) (6.031,-4.750) (6.025,-4.800) (6.019,-4.850) (6.013,-4.900) (6.006,-4.950) (5.998,-5.000) (5.990,-5.050) (5.982,-5.100) (5.973,-5.150) (5.963,-5.200) (5.953,-5.250) (5.942,-5.300) (5.931,-5.350) (5.919,-5.400) (5.906,-5.450) (5.893,-5.500) (5.880,-5.550) (5.866,-5.600) (5.851,-5.650) (5.836,-5.700) (5.820,-5.750) (5.804,-5.800) (5.787,-5.850) (5.769,-5.900) (5.751,-5.950) (5.732,-6.000) (5.713,-6.050) (5.693,-6.100) (5.673,-6.150) (5.652,-6.200) (5.630,-6.250) (5.608,-6.300) (5.585,-6.350) (5.561,-6.400) (5.537,-6.450) (5.512,-6.500) (5.487,-6.550) (5.461,-6.600) (5.434,-6.650) (5.406,-6.700) (5.378,-6.750) (5.350,-6.799) (5.320,-6.850) (5.289,-6.900) (5.258,-6.950) (5.226,-7.000) (5.194,-7.050) (5.160,-7.100) (5.126,-7.150) (5.091,-7.200) (5.056,-7.250) (5.025,-7.292) (4.981,-7.350) (4.950,-7.391) (4.904,-7.450) (4.864,-7.500) (4.825,-7.547) (4.781,-7.600) (4.738,-7.650) (4.700,-7.693) (4.650,-7.749) (4.603,-7.800) (4.555,-7.850) (4.507,-7.900) (4.457,-7.950) (4.407,-8.000) (4.355,-8.050) (4.302,-8.100) (4.250,-8.147) (4.200,-8.192) (4.162,-8.225) (4.104,-8.275) (4.050,-8.320) (4.000,-8.361) (3.952,-8.400) (3.900,-8.441) (3.855,-8.475) (3.800,-8.517) (3.755,-8.550) (3.700,-8.589) (3.650,-8.625) (3.600,-8.659) (3.550,-8.693) (3.501,-8.725) (3.450,-8.758) (3.400,-8.790) (3.350,-8.821) (3.302,-8.850) (3.250,-8.881) (3.200,-8.910) (3.150,-8.939) (3.100,-8.967) (3.050,-8.994) (3.000,-9.021) (2.950,-9.047) (2.900,-9.072) (2.850,-9.097) (2.800,-9.122) (2.750,-9.146) (2.700,-9.169) (2.650,-9.192) (2.600,-9.214) (2.550,-9.236) (2.500,-9.257) (2.450,-9.278) (2.400,-9.298) (2.350,-9.318) (2.300,-9.337) (2.250,-9.356) (2.200,-9.374) (2.150,-9.392) (2.100,-9.410) (2.050,-9.427) (2.000,-9.443) (1.950,-9.459) (1.900,-9.475) (1.850,-9.490) (1.800,-9.504) (1.750,-9.519) (1.700,-9.533) (1.650,-9.546) (1.600,-9.559) (1.550,-9.571) (1.500,-9.583) (1.450,-9.595) (1.400,-9.606) (1.350,-9.617) (1.300,-9.628) (1.250,-9.638) (1.200,-9.647) (1.150,-9.656) (1.100,-9.665) (1.050,-9.674) (1.000,-9.682) (0.950,-9.689) (0.900,-9.696) (0.850,-9.703) (0.800,-9.710) (0.750,-9.716) (0.700,-9.721) (0.650,-9.726) (0.600,-9.731) (0.550,-9.736) (0.500,-9.740) (0.450,-9.743) (0.400,-9.747) (0.350,-9.750) (0.300,-9.752) (0.250,-9.754) (0.200,-9.756) (0.150,-9.757) (0.100,-9.758) (0.050,-9.759) (0.000,-9.759) (-0.050,-9.759) (-0.100,-9.758) (-0.150,-9.757) (-0.200,-9.756) (-0.250,-9.754) (-0.300,-9.752) (-0.350,-9.750) (-0.400,-9.747) (-0.450,-9.743) (-0.500,-9.740) (-0.550,-9.736) (-0.600,-9.731) (-0.650,-9.726) (-0.700,-9.721) (-0.750,-9.716) (-0.800,-9.710) (-0.850,-9.703) (-0.900,-9.696) (-0.950,-9.689) (-1.000,-9.682) (-1.050,-9.674) (-1.100,-9.665) (-1.150,-9.656) (-1.200,-9.647) (-1.250,-9.638) (-1.300,-9.628) (-1.350,-9.617) (-1.400,-9.606) (-1.450,-9.595) (-1.500,-9.583) (-1.550,-9.571) (-1.600,-9.559) (-1.650,-9.546) (-1.700,-9.533) (-1.750,-9.519) (-1.800,-9.504) (-1.850,-9.490) (-1.899,-9.475) (-1.950,-9.459) (-2.000,-9.443) (-2.050,-9.427) (-2.100,-9.410) (-2.150,-9.392) (-2.198,-9.375) (-2.250,-9.356) (-2.300,-9.337) (-2.350,-9.318) (-2.400,-9.298) (-2.450,-9.278) (-2.500,-9.257) (-2.550,-9.236) (-2.600,-9.214) (-2.650,-9.192) (-2.700,-9.169) (-2.750,-9.146) (-2.800,-9.122) (-2.850,-9.097) (-2.900,-9.072) (-2.950,-9.047) (-3.000,-9.021) (-3.050,-8.994) (-3.100,-8.967) (-3.150,-8.939) (-3.200,-8.910) (-3.250,-8.881) (-3.300,-8.851) (-3.344,-8.825) (-3.400,-8.790) (-3.450,-8.758) (-3.500,-8.726) (-3.550,-8.693) (-3.600,-8.659) (-3.650,-8.625) (-3.700,-8.589) (-3.750,-8.553) (-3.800,-8.517) (-3.850,-8.479) (-3.900,-8.441) (-3.950,-8.401) (-4.000,-8.361) (-4.045,-8.325) (-4.100,-8.279) (-4.150,-8.236) (-4.191,-8.200) (-4.247,-8.150) (-4.300,-8.101) (-4.350,-8.055) (-4.400,-8.007) (-4.450,-7.957) (-4.500,-7.907) (-4.550,-7.856) (-4.600,-7.803) (-4.649,-7.750) (-4.694,-7.700) (-4.738,-7.650) (-4.775,-7.607) (-4.823,-7.550) (-4.864,-7.500) (-4.900,-7.455) (-4.943,-7.400) (-4.975,-7.359) (-5.019,-7.300) (-5.050,-7.258) (-5.091,-7.200) (-5.125,-7.152) (-5.160,-7.100) (-5.194,-7.050) (-5.225,-7.002) (-5.258,-6.950) (-5.289,-6.900) (-5.320,-6.850) (-5.349,-6.800) (-5.375,-6.755) (-5.400,-6.711) (-5.425,-6.666) (-5.450,-6.620) (-5.474,-6.575) (-5.500,-6.525) (-5.525,-6.475) (-5.549,-6.425) (-5.573,-6.375) (-5.596,-6.325) (-5.619,-6.275) (-5.641,-6.225) (-5.662,-6.175) (-5.683,-6.125) (-5.703,-6.075) (-5.723,-6.025) (-5.742,-5.975) (-5.760,-5.925) (-5.778,-5.875) (-5.795,-5.825) (-5.812,-5.775) (-5.828,-5.725) (-5.843,-5.675) (-5.858,-5.625) (-5.873,-5.575) (-5.887,-5.525) (-5.900,-5.475) (-5.913,-5.425) (-5.925,-5.375) (-5.936,-5.325) (-5.947,-5.275) (-5.958,-5.225) (-5.968,-5.175) (-5.977,-5.125) (-5.986,-5.075) (-5.994,-5.025) (-6.002,-4.975) (-6.009,-4.925) (-6.016,-4.875) (-6.022,-4.825) (-6.028,-4.775) (-6.033,-4.725) (-6.038,-4.675) (-6.042,-4.625) (-6.045,-4.575) (-6.048,-4.525) (-6.050,-4.475) (-6.052,-4.425) (-6.054,-4.375) (-6.054,-4.325) (-6.055,-4.275) (-6.054,-4.225) (-6.053,-4.175) (-6.052,-4.125) (-6.050,-4.075) (-6.047,-4.025) (-6.044,-3.975) (-6.040,-3.925) (-6.036,-3.875) (-6.031,-3.825) (-6.026,-3.775) (-6.020,-3.725) (-6.013,-3.675) (-6.006,-3.625) (-5.998,-3.575) (-5.990,-3.525) (-5.981,-3.475) (-5.971,-3.425) (-5.961,-3.375) (-5.950,-3.325) (-5.939,-3.275) (-5.926,-3.225) (-5.914,-3.175) (-5.900,-3.125) (-5.886,-3.075) (-5.871,-3.025) (-5.856,-2.975) (-5.839,-2.925) (-5.822,-2.875) (-5.805,-2.825) (-5.786,-2.775) (-5.767,-2.725) (-5.747,-2.675) (-5.727,-2.625) (-5.705,-2.575) (-5.683,-2.525) (-5.660,-2.475) (-5.636,-2.425) (-5.611,-2.375) (-5.585,-2.325) (-5.559,-2.275) (-5.531,-2.225) (-5.503,-2.175) (-5.475,-2.128) (-5.443,-2.075) (-5.411,-2.025) (-5.378,-1.975) (-5.350,-1.933) (-5.310,-1.875) (-5.275,-1.827) (-5.237,-1.775) (-5.200,-1.727) (-5.158,-1.675) (-5.125,-1.634) (-5.075,-1.576) (-5.030,-1.525) (-4.985,-1.475) (-4.950,-1.438) (-4.913,-1.400) (-4.875,-1.362) (-4.838,-1.325) (-4.800,-1.289) (-4.758,-1.250) (-4.702,-1.200) (-4.650,-1.155) (-4.600,-1.113) (-4.553,-1.075) (-4.500,-1.034) (-4.456,-1.000) (-4.400,-0.959) (-4.352,-0.925) (-4.300,-0.889) (-4.250,-0.856) (-4.203,-0.825) (-4.150,-0.792) (-4.100,-0.761) (-4.050,-0.732) (-4.000,-0.703) (-3.950,-0.676) (-3.902,-0.650) (-3.854,-0.625) (-3.805,-0.600) (-3.754,-0.575) (-3.701,-0.550) (-3.650,-0.527) (-3.600,-0.504) (-3.550,-0.483) (-3.500,-0.462) (-3.450,-0.442) (-3.400,-0.422) (-3.350,-0.403) (-3.300,-0.385) (-3.250,-0.367) (-3.200,-0.350) (-3.150,-0.334) (-3.100,-0.318) (-3.050,-0.302) (-3.000,-0.287) (-2.950,-0.273) (-2.900,-0.259) (-2.850,-0.246) (-2.800,-0.233) (-2.750,-0.220) (-2.700,-0.209) (-2.650,-0.197) (-2.600,-0.186) (-2.550,-0.175) (-2.500,-0.165) (-2.450,-0.155) (-2.400,-0.146) (-2.350,-0.137) (-2.300,-0.128) (-2.250,-0.120) (-2.200,-0.112) (-2.150,-0.105) (-2.100,-0.098) (-2.050,-0.091) (-2.000,-0.084) (-1.950,-0.078) (-1.900,-0.072) (-1.850,-0.067) (-1.800,-0.061) (-1.750,-0.056) (-1.700,-0.052) (-1.650,-0.047) (-1.600,-0.043) (-1.550,-0.039) (-1.500,-0.035) (-1.450,-0.032) (-1.400,-0.029) (-1.350,-0.026) (-1.302,-0.023)}
\EQ{(1.300,0.023) (1.350,0.026) (1.400,0.029) (1.450,0.032) (1.500,0.035) (1.550,0.039) (1.600,0.043) (1.650,0.047) (1.700,0.052) (1.750,0.056) (1.800,0.061) (1.850,0.067) (1.900,0.072) (1.950,0.078) (2.000,0.084) (2.050,0.091) (2.100,0.098) (2.150,0.105) (2.200,0.112) (2.250,0.120) (2.300,0.128) (2.350,0.137) (2.400,0.146) (2.450,0.155) (2.500,0.165) (2.550,0.175) (2.600,0.186) (2.650,0.197) (2.700,0.209) (2.750,0.220) (2.800,0.233) (2.850,0.246) (2.900,0.259) (2.950,0.273) (3.000,0.287) (3.050,0.302) (3.100,0.318) (3.150,0.334) (3.199,0.350) (3.250,0.367) (3.300,0.385) (3.350,0.403) (3.400,0.422) (3.450,0.442) (3.500,0.462) (3.550,0.483) (3.600,0.504) (3.646,0.525) (3.700,0.550) (3.750,0.573) (3.800,0.598) (3.850,0.623) (3.900,0.649) (3.949,0.675) (3.994,0.700) (4.038,0.725) (4.100,0.761) (4.150,0.792) (4.200,0.823) (4.250,0.856) (4.300,0.889) (4.350,0.923) (4.400,0.959) (4.450,0.996) (4.500,1.034) (4.550,1.073) (4.600,1.113) (4.644,1.150) (4.700,1.198) (4.750,1.243) (4.800,1.289) (4.838,1.325) (4.875,1.362) (4.913,1.400) (4.950,1.438) (4.985,1.475) (5.025,1.519) (5.074,1.575) (5.117,1.625) (5.150,1.665) (5.198,1.725) (5.237,1.775) (5.274,1.825) (5.310,1.875) (5.345,1.925) (5.375,1.970) (5.411,2.025) (5.443,2.075) (5.473,2.125) (5.500,2.171) (5.525,2.214) (5.550,2.259) (5.575,2.306) (5.598,2.350) (5.623,2.400) (5.648,2.450) (5.671,2.500) (5.694,2.550) (5.716,2.600) (5.737,2.650) (5.757,2.700) (5.777,2.750) (5.796,2.800) (5.814,2.850) (5.831,2.900) (5.848,2.950) (5.863,3.000) (5.879,3.050) (5.893,3.100) (5.907,3.150) (5.920,3.200) (5.933,3.250) (5.944,3.300) (5.956,3.350) (5.966,3.400) (5.976,3.450) (5.985,3.500) (5.994,3.550) (6.002,3.600) (6.010,3.650) (6.017,3.700) (6.023,3.750) (6.029,3.800) (6.034,3.850) (6.038,3.900) (6.042,3.950) (6.046,4.000) (6.049,4.050) (6.051,4.100) (6.053,4.150) (6.054,4.200) (6.054,4.250) (6.054,4.300) (6.054,4.350) (6.053,4.400) (6.051,4.450) (6.049,4.500) (6.047,4.550) (6.043,4.600) (6.040,4.650) (6.035,4.700) (6.031,4.750) (6.025,4.800) (6.019,4.850) (6.013,4.900) (6.006,4.950) (5.998,5.000) (5.990,5.050) (5.982,5.100) (5.973,5.150) (5.963,5.200) (5.953,5.250) (5.942,5.300) (5.931,5.350) (5.919,5.400) (5.906,5.450) (5.893,5.500) (5.880,5.550) (5.866,5.600) (5.851,5.650) (5.836,5.700) (5.820,5.750) (5.804,5.800) (5.787,5.850) (5.769,5.900) (5.751,5.950) (5.732,6.000) (5.713,6.050) (5.693,6.100) (5.673,6.150) (5.652,6.200) (5.630,6.250) (5.608,6.300) (5.585,6.350) (5.561,6.400) (5.537,6.450) (5.512,6.500) (5.487,6.550) (5.461,6.600) (5.434,6.650) (5.406,6.700) (5.378,6.750) (5.350,6.799) (5.320,6.850) (5.289,6.900) (5.258,6.950) (5.226,7.000) (5.194,7.050) (5.160,7.100) (5.126,7.150) (5.091,7.200) (5.056,7.250) (5.025,7.292) (4.981,7.350) (4.950,7.391) (4.904,7.450) (4.864,7.500) (4.825,7.547) (4.781,7.600) (4.738,7.650) (4.700,7.693) (4.650,7.749) (4.603,7.800) (4.555,7.850) (4.507,7.900) (4.457,7.950) (4.407,8.000) (4.355,8.050) (4.302,8.100) (4.250,8.147) (4.200,8.192) (4.162,8.225) (4.104,8.275) (4.050,8.320) (4.000,8.361) (3.952,8.400) (3.900,8.441) (3.855,8.475) (3.800,8.517) (3.755,8.550) (3.700,8.589) (3.650,8.625) (3.600,8.659) (3.550,8.693) (3.501,8.725) (3.450,8.758) (3.400,8.790) (3.350,8.821) (3.302,8.850) (3.250,8.881) (3.200,8.910) (3.150,8.939) (3.100,8.967) (3.050,8.994) (3.000,9.021) (2.950,9.047) (2.900,9.072) (2.850,9.097) (2.800,9.122) (2.750,9.146) (2.700,9.169) (2.650,9.192) (2.600,9.214) (2.550,9.236) (2.500,9.257) (2.450,9.278) (2.400,9.298) (2.350,9.318) (2.300,9.337) (2.250,9.356) (2.200,9.374) (2.150,9.392) (2.100,9.410) (2.050,9.427) (2.000,9.443) (1.950,9.459) (1.900,9.475) (1.850,9.490) (1.800,9.504) (1.750,9.519) (1.700,9.533) (1.650,9.546) (1.600,9.559) (1.550,9.571) (1.500,9.583) (1.450,9.595) (1.400,9.606) (1.350,9.617) (1.300,9.628) (1.250,9.638) (1.200,9.647) (1.150,9.656) (1.100,9.665) (1.050,9.674) (1.000,9.682) (0.950,9.689) (0.900,9.696) (0.850,9.703) (0.800,9.710) (0.750,9.716) (0.700,9.721) (0.650,9.726) (0.600,9.731) (0.550,9.736) (0.500,9.740) (0.450,9.743) (0.400,9.747) (0.350,9.750) (0.300,9.752) (0.250,9.754) (0.200,9.756) (0.150,9.757) (0.100,9.758) (0.050,9.759) (0.000,9.759) (-0.050,9.759) (-0.100,9.758) (-0.150,9.757) (-0.200,9.756) (-0.250,9.754) (-0.300,9.752) (-0.350,9.750) (-0.400,9.747) (-0.450,9.743) (-0.500,9.740) (-0.550,9.736) (-0.600,9.731) (-0.650,9.726) (-0.700,9.721) (-0.750,9.716) (-0.800,9.710) (-0.850,9.703) (-0.900,9.696) (-0.950,9.689) (-1.000,9.682) (-1.050,9.674) (-1.100,9.665) (-1.150,9.656) (-1.200,9.647) (-1.250,9.638) (-1.300,9.628) (-1.350,9.617) (-1.400,9.606) (-1.450,9.595) (-1.500,9.583) (-1.550,9.571) (-1.600,9.559) (-1.650,9.546) (-1.700,9.533) (-1.750,9.519) (-1.800,9.504) (-1.850,9.490) (-1.899,9.475) (-1.950,9.459) (-2.000,9.443) (-2.050,9.427) (-2.100,9.410) (-2.150,9.392) (-2.198,9.375) (-2.250,9.356) (-2.300,9.337) (-2.350,9.318) (-2.400,9.298) (-2.450,9.278) (-2.500,9.257) (-2.550,9.236) (-2.600,9.214) (-2.650,9.192) (-2.700,9.169) (-2.750,9.146) (-2.800,9.122) (-2.850,9.097) (-2.900,9.072) (-2.950,9.047) (-3.000,9.021) (-3.050,8.994) (-3.100,8.967) (-3.150,8.939) (-3.200,8.910) (-3.250,8.881) (-3.300,8.851) (-3.344,8.825) (-3.400,8.790) (-3.450,8.758) (-3.500,8.726) (-3.550,8.693) (-3.600,8.659) (-3.650,8.625) (-3.700,8.589) (-3.750,8.553) (-3.800,8.517) (-3.850,8.479) (-3.900,8.441) (-3.950,8.401) (-4.000,8.361) (-4.045,8.325) (-4.100,8.279) (-4.150,8.236) (-4.191,8.200) (-4.247,8.150) (-4.300,8.101) (-4.350,8.055) (-4.400,8.007) (-4.450,7.957) (-4.500,7.907) (-4.550,7.856) (-4.600,7.803) (-4.649,7.750) (-4.694,7.700) (-4.738,7.650) (-4.775,7.607) (-4.823,7.550) (-4.864,7.500) (-4.900,7.455) (-4.943,7.400) (-4.975,7.359) (-5.019,7.300) (-5.050,7.258) (-5.091,7.200) (-5.125,7.152) (-5.160,7.100) (-5.194,7.050) (-5.225,7.002) (-5.258,6.950) (-5.289,6.900) (-5.320,6.850) (-5.349,6.800) (-5.375,6.755) (-5.400,6.711) (-5.425,6.666) (-5.450,6.620) (-5.474,6.575) (-5.500,6.525) (-5.525,6.475) (-5.549,6.425) (-5.573,6.375) (-5.596,6.325) (-5.619,6.275) (-5.641,6.225) (-5.662,6.175) (-5.683,6.125) (-5.703,6.075) (-5.723,6.025) (-5.742,5.975) (-5.760,5.925) (-5.778,5.875) (-5.795,5.825) (-5.812,5.775) (-5.828,5.725) (-5.843,5.675) (-5.858,5.625) (-5.873,5.575) (-5.887,5.525) (-5.900,5.475) (-5.913,5.425) (-5.925,5.375) (-5.936,5.325) (-5.947,5.275) (-5.958,5.225) (-5.968,5.175) (-5.977,5.125) (-5.986,5.075) (-5.994,5.025) (-6.002,4.975) (-6.009,4.925) (-6.016,4.875) (-6.022,4.825) (-6.028,4.775) (-6.033,4.725) (-6.038,4.675) (-6.042,4.625) (-6.045,4.575) (-6.048,4.525) (-6.050,4.475) (-6.052,4.425) (-6.054,4.375) (-6.054,4.325) (-6.055,4.275) (-6.054,4.225) (-6.053,4.175) (-6.052,4.125) (-6.050,4.075) (-6.047,4.025) (-6.044,3.975) (-6.040,3.925) (-6.036,3.875) (-6.031,3.825) (-6.026,3.775) (-6.020,3.725) (-6.013,3.675) (-6.006,3.625) (-5.998,3.575) (-5.990,3.525) (-5.981,3.475) (-5.971,3.425) (-5.961,3.375) (-5.950,3.325) (-5.939,3.275) (-5.926,3.225) (-5.914,3.175) (-5.900,3.125) (-5.886,3.075) (-5.871,3.025) (-5.856,2.975) (-5.839,2.925) (-5.822,2.875) (-5.805,2.825) (-5.786,2.775) (-5.767,2.725) (-5.747,2.675) (-5.727,2.625) (-5.705,2.575) (-5.683,2.525) (-5.660,2.475) (-5.636,2.425) (-5.611,2.375) (-5.585,2.325) (-5.559,2.275) (-5.531,2.225) (-5.503,2.175) (-5.475,2.128) (-5.443,2.075) (-5.411,2.025) (-5.378,1.975) (-5.350,1.933) (-5.310,1.875) (-5.275,1.827) (-5.237,1.775) (-5.200,1.727) (-5.158,1.675) (-5.125,1.634) (-5.075,1.576) (-5.030,1.525) (-4.985,1.475) (-4.950,1.438) (-4.913,1.400) (-4.875,1.362) (-4.838,1.325) (-4.800,1.289) (-4.758,1.250) (-4.702,1.200) (-4.650,1.155) (-4.600,1.113) (-4.553,1.075) (-4.500,1.034) (-4.456,1.000) (-4.400,0.959) (-4.352,0.925) (-4.300,0.889) (-4.250,0.856) (-4.203,0.825) (-4.150,0.792) (-4.100,0.761) (-4.050,0.732) (-4.000,0.703) (-3.950,0.676) (-3.902,0.650) (-3.854,0.625) (-3.805,0.600) (-3.754,0.575) (-3.701,0.550) (-3.650,0.527) (-3.600,0.504) (-3.550,0.483) (-3.500,0.462) (-3.450,0.442) (-3.400,0.422) (-3.350,0.403) (-3.300,0.385) (-3.250,0.367) (-3.200,0.350) (-3.150,0.334) (-3.100,0.318) (-3.050,0.302) (-3.000,0.287) (-2.950,0.273) (-2.900,0.259) (-2.850,0.246) (-2.800,0.233) (-2.750,0.220) (-2.700,0.209) (-2.650,0.197) (-2.600,0.186) (-2.550,0.175) (-2.500,0.165) (-2.450,0.155) (-2.400,0.146) (-2.350,0.137) (-2.300,0.128) (-2.250,0.120) (-2.200,0.112) (-2.150,0.105) (-2.100,0.098) (-2.050,0.091) (-2.000,0.084) (-1.950,0.078) (-1.900,0.072) (-1.850,0.067) (-1.800,0.061) (-1.750,0.056) (-1.700,0.052) (-1.650,0.047) (-1.600,0.043) (-1.550,0.039) (-1.500,0.035) (-1.450,0.032) (-1.400,0.029) (-1.350,0.026) (-1.302,0.023)}
\EQ{(1.300,0.046) (1.350,0.052) (1.400,0.058) (1.450,0.064) (1.500,0.071) (1.550,0.079) (1.600,0.086) (1.650,0.095) (1.700,0.104) (1.750,0.113) (1.800,0.123) (1.850,0.134) (1.900,0.145) (1.950,0.157) (2.000,0.170) (2.050,0.183) (2.100,0.197) (2.150,0.212) (2.200,0.227) (2.250,0.243) (2.300,0.260) (2.350,0.278) (2.400,0.297) (2.450,0.317) (2.500,0.337) (2.550,0.359) (2.600,0.381) (2.650,0.405) (2.700,0.429) (2.750,0.455) (2.800,0.482) (2.850,0.510) (2.900,0.539) (2.950,0.570) (2.998,0.600) (3.050,0.635) (3.100,0.670) (3.141,0.700) (3.200,0.745) (3.250,0.785) (3.298,0.825) (3.350,0.871) (3.400,0.917) (3.450,0.966) (3.500,1.017) (3.550,1.071) (3.598,1.125) (3.640,1.175) (3.675,1.219) (3.718,1.275) (3.750,1.319) (3.789,1.375) (3.822,1.425) (3.850,1.470) (3.883,1.525) (3.912,1.575) (3.939,1.625) (3.965,1.675) (3.990,1.725) (4.014,1.775) (4.036,1.825) (4.057,1.875) (4.077,1.925) (4.096,1.975) (4.114,2.025) (4.131,2.075) (4.147,2.125) (4.162,2.175) (4.176,2.225) (4.189,2.275) (4.201,2.325) (4.213,2.375) (4.223,2.425) (4.232,2.475) (4.241,2.525) (4.249,2.575) (4.256,2.625) (4.262,2.675) (4.267,2.725) (4.271,2.775) (4.275,2.825) (4.278,2.875) (4.280,2.925) (4.281,2.975) (4.281,3.025) (4.281,3.075) (4.280,3.125) (4.278,3.175) (4.275,3.225) (4.272,3.275) (4.268,3.325) (4.263,3.375) (4.257,3.425) (4.250,3.475) (4.243,3.525) (4.235,3.575) (4.226,3.625) (4.217,3.675) (4.206,3.725) (4.195,3.775) (4.184,3.825) (4.171,3.875) (4.157,3.925) (4.143,3.975) (4.128,4.025) (4.112,4.075) (4.096,4.125) (4.078,4.175) (4.060,4.225) (4.041,4.275) (4.021,4.325) (4.000,4.375) (3.979,4.425) (3.956,4.475) (3.933,4.525) (3.908,4.575) (3.883,4.625) (3.857,4.675) (3.830,4.725) (3.802,4.775) (3.775,4.821) (3.743,4.875) (3.711,4.925) (3.679,4.975) (3.650,5.019) (3.611,5.075) (3.576,5.125) (3.539,5.175) (3.501,5.225) (3.462,5.275) (3.425,5.321) (3.380,5.375) (3.336,5.425) (3.300,5.466) (3.250,5.520) (3.200,5.573) (3.150,5.623) (3.100,5.672) (3.050,5.720) (3.000,5.766) (2.961,5.800) (2.904,5.850) (2.850,5.895) (2.800,5.935) (2.750,5.974) (2.700,6.013) (2.650,6.049) (2.600,6.085) (2.550,6.120) (2.505,6.150) (2.450,6.186) (2.400,6.218) (2.350,6.249) (2.306,6.275) (2.250,6.308) (2.200,6.336) (2.150,6.363) (2.100,6.390) (2.050,6.415) (2.000,6.440) (1.950,6.464) (1.900,6.488) (1.850,6.510) (1.800,6.532) (1.750,6.553) (1.700,6.574) (1.650,6.593) (1.600,6.612) (1.550,6.631) (1.500,6.648) (1.450,6.665) (1.400,6.682) (1.350,6.698) (1.300,6.713) (1.250,6.727) (1.200,6.741) (1.150,6.754) (1.100,6.767) (1.050,6.779) (1.000,6.791) (0.950,6.801) (0.900,6.812) (0.850,6.821) (0.800,6.831) (0.750,6.839) (0.700,6.847) (0.650,6.854) (0.600,6.861) (0.550,6.868) (0.500,6.873) (0.450,6.879) (0.400,6.883) (0.350,6.887) (0.300,6.891) (0.250,6.894) (0.200,6.896) (0.150,6.898) (0.100,6.900) (0.050,6.900) (0.000,6.901) (-0.050,6.900) (-0.100,6.900) (-0.150,6.898) (-0.200,6.896) (-0.250,6.894) (-0.300,6.891) (-0.350,6.887) (-0.400,6.883) (-0.450,6.879) (-0.500,6.873) (-0.550,6.868) (-0.600,6.861) (-0.650,6.854) (-0.700,6.847) (-0.750,6.839) (-0.800,6.831) (-0.850,6.821) (-0.900,6.812) (-0.950,6.801) (-1.000,6.791) (-1.050,6.779) (-1.100,6.767) (-1.150,6.754) (-1.200,6.741) (-1.250,6.727) (-1.300,6.713) (-1.350,6.698) (-1.400,6.682) (-1.450,6.665) (-1.500,6.648) (-1.550,6.631) (-1.600,6.612) (-1.650,6.593) (-1.697,6.575) (-1.750,6.553) (-1.800,6.532) (-1.850,6.510) (-1.900,6.488) (-1.950,6.464) (-2.000,6.440) (-2.050,6.415) (-2.100,6.390) (-2.150,6.363) (-2.200,6.336) (-2.250,6.308) (-2.300,6.279) (-2.348,6.250) (-2.400,6.218) (-2.450,6.186) (-2.500,6.154) (-2.543,6.125) (-2.600,6.085) (-2.649,6.050) (-2.700,6.013) (-2.749,5.975) (-2.800,5.935) (-2.844,5.900) (-2.900,5.853) (-2.950,5.810) (-2.990,5.775) (-3.044,5.725) (-3.097,5.675) (-3.148,5.625) (-3.198,5.575) (-3.246,5.525) (-3.292,5.475) (-3.336,5.425) (-3.375,5.380) (-3.421,5.325) (-3.462,5.275) (-3.500,5.227) (-3.539,5.175) (-3.575,5.126) (-3.611,5.075) (-3.646,5.025) (-3.675,4.981) (-3.711,4.925) (-3.743,4.875) (-3.773,4.825) (-3.800,4.778) (-3.825,4.734) (-3.850,4.688) (-3.875,4.641) (-3.900,4.592) (-3.925,4.541) (-3.950,4.488) (-3.975,4.433) (-4.000,4.376) (-4.021,4.325) (-4.041,4.275) (-4.060,4.225) (-4.078,4.175) (-4.096,4.125) (-4.112,4.075) (-4.128,4.025) (-4.143,3.975) (-4.157,3.925) (-4.171,3.875) (-4.184,3.825) (-4.195,3.775) (-4.206,3.725) (-4.217,3.675) (-4.226,3.625) (-4.235,3.575) (-4.243,3.525) (-4.250,3.475) (-4.257,3.425) (-4.263,3.375) (-4.268,3.325) (-4.272,3.275) (-4.275,3.225) (-4.278,3.175) (-4.280,3.125) (-4.281,3.075) (-4.281,3.025) (-4.281,2.975) (-4.280,2.925) (-4.278,2.875) (-4.275,2.825) (-4.271,2.775) (-4.267,2.725) (-4.262,2.675) (-4.256,2.625) (-4.249,2.575) (-4.241,2.525) (-4.232,2.475) (-4.223,2.425) (-4.213,2.375) (-4.201,2.325) (-4.189,2.275) (-4.176,2.225) (-4.162,2.175) (-4.147,2.125) (-4.131,2.075) (-4.114,2.025) (-4.096,1.975) (-4.077,1.925) (-4.057,1.875) (-4.036,1.825) (-4.014,1.775) (-3.990,1.725) (-3.965,1.675) (-3.939,1.625) (-3.912,1.575) (-3.883,1.525) (-3.853,1.475) (-3.825,1.430) (-3.789,1.375) (-3.754,1.325) (-3.718,1.275) (-3.680,1.225) (-3.640,1.175) (-3.600,1.127) (-3.554,1.075) (-3.508,1.025) (-3.459,0.975) (-3.408,0.925) (-3.355,0.875) (-3.300,0.827) (-3.250,0.785) (-3.207,0.750) (-3.150,0.707) (-3.107,0.675) (-3.050,0.635) (-3.000,0.602) (-2.950,0.570) (-2.900,0.539) (-2.850,0.510) (-2.800,0.482) (-2.750,0.455) (-2.700,0.429) (-2.650,0.405) (-2.600,0.381) (-2.550,0.359) (-2.500,0.337) (-2.450,0.317) (-2.400,0.297) (-2.350,0.278) (-2.300,0.260) (-2.250,0.243) (-2.200,0.227) (-2.150,0.212) (-2.100,0.197) (-2.050,0.183) (-2.000,0.170) (-1.950,0.157) (-1.900,0.145) (-1.850,0.134) (-1.800,0.123) (-1.750,0.113) (-1.700,0.104) (-1.650,0.095) (-1.600,0.086) (-1.550,0.079) (-1.500,0.071) (-1.450,0.064) (-1.400,0.058) (-1.350,0.052) (-1.300,0.046)}
\EQ{(1.300,0.070) (1.350,0.078) (1.400,0.087) (1.450,0.097) (1.500,0.107) (1.550,0.118) (1.600,0.130) (1.650,0.143) (1.700,0.157) (1.750,0.171) (1.800,0.187) (1.850,0.203) (1.900,0.220) (1.950,0.239) (2.000,0.258) (2.050,0.279) (2.098,0.300) (2.150,0.324) (2.200,0.348) (2.250,0.374) (2.298,0.400) (2.342,0.425) (2.400,0.460) (2.450,0.492) (2.500,0.525) (2.550,0.561) (2.600,0.598) (2.650,0.638) (2.694,0.675) (2.750,0.724) (2.800,0.772) (2.850,0.822) (2.899,0.875) (2.943,0.925) (2.985,0.975) (3.024,1.025) (3.060,1.075) (3.095,1.125) (3.125,1.171) (3.159,1.225) (3.188,1.275) (3.216,1.325) (3.241,1.375) (3.266,1.425) (3.289,1.475) (3.310,1.525) (3.330,1.575) (3.349,1.625) (3.367,1.675) (3.383,1.725) (3.398,1.775) (3.412,1.825) (3.425,1.875) (3.437,1.925) (3.447,1.975) (3.456,2.025) (3.465,2.075) (3.472,2.125) (3.479,2.175) (3.484,2.225) (3.488,2.275) (3.491,2.325) (3.494,2.375) (3.495,2.425) (3.496,2.475) (3.495,2.525) (3.494,2.575) (3.491,2.625) (3.488,2.675) (3.483,2.725) (3.478,2.775) (3.472,2.825) (3.465,2.875) (3.457,2.925) (3.448,2.975) (3.438,3.025) (3.427,3.075) (3.415,3.125) (3.403,3.175) (3.389,3.225) (3.374,3.275) (3.359,3.325) (3.342,3.375) (3.325,3.423) (3.306,3.475) (3.286,3.525) (3.265,3.575) (3.243,3.625) (3.220,3.675) (3.196,3.725) (3.171,3.775) (3.145,3.825) (3.118,3.875) (3.089,3.925) (3.059,3.975) (3.028,4.025) (3.000,4.068) (2.962,4.125) (2.927,4.175) (2.890,4.225) (2.852,4.275) (2.812,4.325) (2.775,4.370) (2.728,4.425) (2.683,4.475) (2.637,4.525) (2.600,4.563) (2.563,4.600) (2.525,4.637) (2.485,4.675) (2.430,4.725) (2.375,4.773) (2.325,4.815) (2.281,4.850) (2.225,4.894) (2.183,4.925) (2.125,4.967) (2.078,5.000) (2.025,5.035) (1.975,5.067) (1.925,5.098) (1.881,5.125) (1.825,5.157) (1.775,5.185) (1.725,5.212) (1.675,5.238) (1.625,5.262) (1.575,5.286) (1.525,5.309) (1.475,5.331) (1.425,5.352) (1.375,5.373) (1.325,5.392) (1.275,5.411) (1.225,5.428) (1.175,5.445) (1.125,5.462) (1.075,5.477) (1.025,5.492) (0.975,5.505) (0.925,5.519) (0.875,5.531) (0.825,5.543) (0.775,5.553) (0.725,5.564) (0.675,5.573) (0.625,5.582) (0.575,5.590) (0.525,5.597) (0.475,5.604) (0.425,5.610) (0.375,5.616) (0.325,5.620) (0.275,5.624) (0.225,5.628) (0.175,5.630) (0.125,5.632) (0.075,5.634) (0.025,5.634) (-0.050,5.634) (-0.100,5.633) (-0.150,5.631) (-0.200,5.629) (-0.250,5.626) (-0.300,5.622) (-0.350,5.618) (-0.400,5.613) (-0.450,5.607) (-0.500,5.601) (-0.550,5.594) (-0.600,5.586) (-0.650,5.578) (-0.700,5.569) (-0.750,5.559) (-0.800,5.548) (-0.850,5.537) (-0.899,5.525) (-0.950,5.512) (-1.000,5.499) (-1.050,5.484) (-1.100,5.469) (-1.150,5.453) (-1.200,5.437) (-1.250,5.420) (-1.300,5.401) (-1.350,5.382) (-1.400,5.363) (-1.450,5.342) (-1.500,5.320) (-1.550,5.298) (-1.599,5.275) (-1.650,5.250) (-1.700,5.225) (-1.747,5.200) (-1.793,5.175) (-1.838,5.150) (-1.900,5.114) (-1.950,5.083) (-2.000,5.051) (-2.050,5.019) (-2.100,4.985) (-2.149,4.950) (-2.200,4.913) (-2.249,4.875) (-2.300,4.835) (-2.343,4.800) (-2.400,4.751) (-2.450,4.707) (-2.500,4.661) (-2.538,4.625) (-2.575,4.588) (-2.613,4.550) (-2.650,4.511) (-2.700,4.457) (-2.750,4.400) (-2.792,4.350) (-2.825,4.309) (-2.871,4.250) (-2.908,4.200) (-2.944,4.150) (-2.975,4.106) (-3.012,4.050) (-3.044,4.000) (-3.074,3.950) (-3.100,3.906) (-3.125,3.862) (-3.150,3.816) (-3.175,3.768) (-3.200,3.718) (-3.225,3.665) (-3.250,3.610) (-3.275,3.551) (-3.296,3.500) (-3.315,3.450) (-3.333,3.400) (-3.350,3.351) (-3.366,3.300) (-3.382,3.250) (-3.396,3.200) (-3.409,3.150) (-3.421,3.100) (-3.433,3.050) (-3.443,3.000) (-3.453,2.950) (-3.461,2.900) (-3.469,2.850) (-3.475,2.800) (-3.481,2.750) (-3.486,2.700) (-3.490,2.650) (-3.492,2.600) (-3.494,2.550) (-3.495,2.500) (-3.496,2.450) (-3.495,2.400) (-3.493,2.350) (-3.490,2.300) (-3.486,2.250) (-3.481,2.200) (-3.475,2.150) (-3.469,2.100) (-3.461,2.050) (-3.452,2.000) (-3.442,1.950) (-3.431,1.900) (-3.419,1.850) (-3.405,1.800) (-3.391,1.750) (-3.375,1.700) (-3.358,1.650) (-3.340,1.600) (-3.321,1.550) (-3.300,1.501) (-3.277,1.450) (-3.254,1.400) (-3.229,1.350) (-3.202,1.300) (-3.175,1.252) (-3.144,1.200) (-3.112,1.150) (-3.078,1.100) (-3.042,1.050) (-3.004,1.000) (-2.964,0.950) (-2.925,0.904) (-2.876,0.850) (-2.828,0.800) (-2.777,0.750) (-2.725,0.702) (-2.675,0.659) (-2.634,0.625) (-2.575,0.579) (-2.525,0.543) (-2.475,0.508) (-2.425,0.475) (-2.375,0.444) (-2.325,0.415) (-2.275,0.387) (-2.225,0.361) (-2.175,0.336) (-2.125,0.312) (-2.075,0.290) (-2.025,0.269) (-1.975,0.248) (-1.925,0.230) (-1.875,0.212) (-1.825,0.195) (-1.775,0.179) (-1.725,0.164) (-1.676,0.150) (-1.625,0.137) (-1.575,0.124) (-1.525,0.113) (-1.475,0.102) (-1.425,0.092) (-1.375,0.082) (-1.325,0.074) (-1.300,0.070)}
\EQ{(1.300,0.093) (1.350,0.104) (1.400,0.116) (1.450,0.130) (1.500,0.144) (1.550,0.159) (1.600,0.175) (1.650,0.193) (1.700,0.211) (1.750,0.231) (1.800,0.252) (1.850,0.275) (1.900,0.299) (1.950,0.324) (1.997,0.350) (2.041,0.375) (2.100,0.412) (2.150,0.444) (2.194,0.475) (2.250,0.517) (2.292,0.550) (2.350,0.599) (2.400,0.645) (2.450,0.693) (2.500,0.746) (2.548,0.800) (2.589,0.850) (2.625,0.896) (2.664,0.950) (2.697,1.000) (2.725,1.044) (2.759,1.100) (2.786,1.150) (2.812,1.200) (2.836,1.250) (2.858,1.300) (2.879,1.350) (2.898,1.400) (2.916,1.450) (2.932,1.500) (2.947,1.550) (2.960,1.600) (2.972,1.650) (2.983,1.700) (2.993,1.750) (3.001,1.800) (3.008,1.850) (3.014,1.900) (3.019,1.950) (3.023,2.000) (3.025,2.050) (3.027,2.100) (3.027,2.150) (3.026,2.200) (3.025,2.250) (3.022,2.300) (3.018,2.350) (3.013,2.400) (3.006,2.450) (2.999,2.500) (2.991,2.550) (2.981,2.600) (2.971,2.650) (2.959,2.700) (2.947,2.750) (2.933,2.800) (2.918,2.850) (2.902,2.900) (2.885,2.950) (2.866,3.000) (2.847,3.050) (2.826,3.100) (2.804,3.150) (2.781,3.200) (2.756,3.250) (2.730,3.300) (2.703,3.350) (2.675,3.399) (2.645,3.450) (2.613,3.500) (2.580,3.550) (2.550,3.594) (2.509,3.650) (2.475,3.696) (2.432,3.750) (2.400,3.789) (2.350,3.846) (2.301,3.900) (2.253,3.950) (2.203,4.000) (2.151,4.050) (2.100,4.096) (2.050,4.139) (2.007,4.175) (1.950,4.220) (1.900,4.258) (1.850,4.295) (1.807,4.325) (1.750,4.363) (1.700,4.395) (1.651,4.425) (1.600,4.455) (1.550,4.483) (1.500,4.510) (1.450,4.536) (1.400,4.561) (1.350,4.585) (1.300,4.607) (1.250,4.629) (1.200,4.649) (1.150,4.669) (1.100,4.687) (1.050,4.705) (1.000,4.722) (0.950,4.737) (0.900,4.752) (0.850,4.766) (0.800,4.779) (0.750,4.792) (0.700,4.803) (0.650,4.814) (0.600,4.824) (0.550,4.833) (0.500,4.841) (0.450,4.848) (0.400,4.855) (0.350,4.861) (0.300,4.866) (0.250,4.870) (0.200,4.873) (0.150,4.876) (0.100,4.878) (0.050,4.879) (0.000,4.880) (-0.050,4.879) (-0.100,4.878) (-0.150,4.876) (-0.200,4.873) (-0.250,4.870) (-0.300,4.866) (-0.350,4.861) (-0.400,4.855) (-0.450,4.848) (-0.500,4.841) (-0.550,4.833) (-0.600,4.824) (-0.650,4.814) (-0.700,4.803) (-0.750,4.792) (-0.800,4.779) (-0.850,4.766) (-0.900,4.752) (-0.950,4.737) (-1.000,4.722) (-1.050,4.705) (-1.100,4.687) (-1.150,4.669) (-1.198,4.650) (-1.250,4.629) (-1.300,4.607) (-1.350,4.585) (-1.400,4.561) (-1.450,4.536) (-1.500,4.510) (-1.550,4.483) (-1.600,4.455) (-1.650,4.426) (-1.700,4.395) (-1.750,4.363) (-1.800,4.330) (-1.843,4.300) (-1.900,4.258) (-1.944,4.225) (-2.000,4.181) (-2.050,4.139) (-2.095,4.100) (-2.150,4.051) (-2.200,4.003) (-2.250,3.954) (-2.300,3.901) (-2.347,3.850) (-2.390,3.800) (-2.425,3.758) (-2.472,3.700) (-2.509,3.650) (-2.546,3.600) (-2.575,3.558) (-2.613,3.500) (-2.645,3.450) (-2.675,3.400) (-2.700,3.356) (-2.725,3.310) (-2.750,3.262) (-2.775,3.212) (-2.800,3.159) (-2.825,3.102) (-2.847,3.050) (-2.866,3.000) (-2.885,2.950) (-2.902,2.900) (-2.918,2.850) (-2.933,2.800) (-2.947,2.750) (-2.959,2.700) (-2.971,2.650) (-2.981,2.600) (-2.991,2.550) (-2.999,2.500) (-3.006,2.450) (-3.013,2.400) (-3.018,2.350) (-3.022,2.300) (-3.025,2.250) (-3.026,2.200) (-3.027,2.150) (-3.027,2.100) (-3.025,2.050) (-3.023,2.000) (-3.019,1.950) (-3.014,1.900) (-3.008,1.850) (-3.001,1.800) (-2.993,1.750) (-2.983,1.700) (-2.972,1.650) (-2.960,1.600) (-2.947,1.550) (-2.932,1.500) (-2.916,1.450) (-2.898,1.400) (-2.879,1.350) (-2.858,1.300) (-2.836,1.250) (-2.812,1.200) (-2.786,1.150) (-2.759,1.100) (-2.729,1.050) (-2.700,1.004) (-2.664,0.950) (-2.628,0.900) (-2.589,0.850) (-2.550,0.802) (-2.504,0.750) (-2.457,0.700) (-2.406,0.650) (-2.351,0.600) (-2.300,0.557) (-2.261,0.525) (-2.200,0.479) (-2.158,0.450) (-2.100,0.412) (-2.050,0.381) (-2.000,0.352) (-1.951,0.325) (-1.902,0.300) (-1.850,0.275) (-1.800,0.252) (-1.750,0.231) (-1.700,0.211) (-1.650,0.193) (-1.600,0.175) (-1.550,0.159) (-1.500,0.144) (-1.450,0.130) (-1.400,0.116) (-1.350,0.104) (-1.300,0.093)}
\EQ{(1.300,0.117) (1.350,0.131) (1.400,0.146) (1.450,0.163) (1.500,0.181) (1.549,0.200) (1.600,0.221) (1.650,0.244) (1.700,0.268) (1.750,0.293) (1.800,0.321) (1.849,0.350) (1.900,0.382) (1.950,0.416) (1.997,0.450) (2.050,0.492) (2.090,0.525) (2.145,0.575) (2.196,0.625) (2.243,0.675) (2.287,0.725) (2.325,0.772) (2.365,0.825) (2.400,0.875) (2.432,0.925) (2.462,0.975) (2.490,1.025) (2.516,1.075) (2.540,1.125) (2.562,1.175) (2.582,1.225) (2.600,1.275) (2.617,1.325) (2.632,1.375) (2.646,1.425) (2.658,1.475) (2.669,1.525) (2.679,1.575) (2.687,1.625) (2.693,1.675) (2.699,1.725) (2.703,1.775) (2.706,1.825) (2.707,1.875) (2.708,1.925) (2.707,1.975) (2.705,2.025) (2.701,2.075) (2.697,2.125) (2.691,2.175) (2.684,2.225) (2.676,2.275) (2.667,2.325) (2.656,2.375) (2.644,2.425) (2.631,2.475) (2.617,2.525) (2.602,2.575) (2.585,2.625) (2.567,2.675) (2.547,2.725) (2.527,2.775) (2.504,2.825) (2.481,2.875) (2.456,2.925) (2.430,2.975) (2.402,3.025) (2.375,3.070) (2.341,3.125) (2.308,3.175) (2.275,3.223) (2.237,3.275) (2.200,3.323) (2.158,3.375) (2.125,3.414) (2.075,3.470) (2.025,3.523) (1.975,3.574) (1.925,3.621) (1.875,3.667) (1.836,3.700) (1.776,3.750) (1.725,3.790) (1.677,3.825) (1.625,3.862) (1.575,3.896) (1.530,3.925) (1.475,3.959) (1.425,3.988) (1.375,4.016) (1.325,4.043) (1.275,4.068) (1.225,4.092) (1.175,4.115) (1.125,4.137) (1.075,4.158) (1.025,4.177) (0.975,4.196) (0.925,4.213) (0.875,4.229) (0.825,4.245) (0.775,4.259) (0.725,4.272) (0.675,4.285) (0.625,4.296) (0.575,4.307) (0.525,4.317) (0.475,4.325) (0.425,4.333) (0.375,4.340) (0.325,4.346) (0.275,4.351) (0.225,4.356) (0.175,4.359) (0.125,4.362) (0.075,4.363) (0.025,4.364) (-0.050,4.364) (-0.100,4.363) (-0.150,4.361) (-0.200,4.358) (-0.250,4.354) (-0.300,4.349) (-0.350,4.343) (-0.400,4.337) (-0.450,4.329) (-0.500,4.321) (-0.550,4.312) (-0.600,4.302) (-0.650,4.291) (-0.700,4.279) (-0.750,4.266) (-0.800,4.252) (-0.850,4.237) (-0.900,4.221) (-0.950,4.204) (-1.000,4.187) (-1.050,4.168) (-1.100,4.147) (-1.150,4.126) (-1.200,4.104) (-1.250,4.080) (-1.300,4.056) (-1.350,4.030) (-1.400,4.002) (-1.448,3.975) (-1.500,3.944) (-1.550,3.912) (-1.600,3.879) (-1.642,3.850) (-1.700,3.808) (-1.744,3.775) (-1.800,3.730) (-1.850,3.688) (-1.894,3.650) (-1.948,3.600) (-1.999,3.550) (-2.047,3.500) (-2.093,3.450) (-2.137,3.400) (-2.175,3.354) (-2.218,3.300) (-2.250,3.257) (-2.291,3.200) (-2.325,3.150) (-2.357,3.100) (-2.387,3.050) (-2.416,3.000) (-2.443,2.950) (-2.469,2.900) (-2.493,2.850) (-2.516,2.800) (-2.537,2.750) (-2.557,2.700) (-2.575,2.653) (-2.593,2.600) (-2.610,2.550) (-2.624,2.500) (-2.638,2.450) (-2.650,2.400) (-2.662,2.350) (-2.672,2.300) (-2.680,2.250) (-2.688,2.200) (-2.694,2.150) (-2.699,2.100) (-2.703,2.050) (-2.706,2.000) (-2.707,1.950) (-2.708,1.900) (-2.707,1.850) (-2.704,1.800) (-2.701,1.750) (-2.696,1.700) (-2.690,1.650) (-2.683,1.600) (-2.674,1.550) (-2.664,1.500) (-2.652,1.450) (-2.639,1.400) (-2.625,1.351) (-2.609,1.300) (-2.591,1.250) (-2.572,1.200) (-2.551,1.150) (-2.528,1.100) (-2.503,1.050) (-2.476,1.000) (-2.450,0.954) (-2.416,0.900) (-2.383,0.850) (-2.350,0.805) (-2.308,0.750) (-2.275,0.711) (-2.225,0.655) (-2.175,0.604) (-2.125,0.557) (-2.075,0.513) (-2.029,0.475) (-1.975,0.434) (-1.927,0.400) (-1.875,0.366) (-1.825,0.335) (-1.775,0.307) (-1.725,0.280) (-1.675,0.255) (-1.625,0.232) (-1.575,0.211) (-1.525,0.191) (-1.475,0.172) (-1.425,0.155) (-1.375,0.139) (-1.325,0.124) (-1.300,0.117)}
\EQ{(1.300,0.141) (1.350,0.158) (1.400,0.177) (1.450,0.197) (1.500,0.220) (1.550,0.243) (1.600,0.269) (1.650,0.297) (1.697,0.325) (1.750,0.359) (1.800,0.394) (1.841,0.425) (1.900,0.473) (1.950,0.517) (2.000,0.566) (2.050,0.619) (2.099,0.675) (2.138,0.725) (2.174,0.775) (2.208,0.825) (2.239,0.875) (2.267,0.925) (2.293,0.975) (2.317,1.025) (2.339,1.075) (2.359,1.125) (2.377,1.175) (2.394,1.225) (2.409,1.275) (2.422,1.325) (2.433,1.375) (2.443,1.425) (2.451,1.475) (2.458,1.525) (2.464,1.575) (2.468,1.625) (2.470,1.675) (2.472,1.725) (2.472,1.775) (2.470,1.825) (2.467,1.875) (2.463,1.925) (2.458,1.975) (2.451,2.025) (2.443,2.075) (2.434,2.125) (2.423,2.175) (2.411,2.225) (2.398,2.275) (2.383,2.325) (2.367,2.375) (2.350,2.423) (2.330,2.475) (2.310,2.525) (2.288,2.575) (2.265,2.625) (2.240,2.675) (2.213,2.725) (2.184,2.775) (2.154,2.825) (2.125,2.871) (2.089,2.925) (2.053,2.975) (2.015,3.025) (1.975,3.075) (1.933,3.125) (1.888,3.175) (1.850,3.215) (1.800,3.265) (1.763,3.300) (1.725,3.335) (1.679,3.375) (1.625,3.420) (1.575,3.459) (1.525,3.496) (1.475,3.531) (1.425,3.565) (1.375,3.597) (1.328,3.625) (1.285,3.650) (1.239,3.675) (1.191,3.700) (1.141,3.725) (1.087,3.750) (1.031,3.775) (0.975,3.798) (0.925,3.817) (0.875,3.835) (0.825,3.852) (0.775,3.868) (0.725,3.883) (0.675,3.897) (0.625,3.909) (0.575,3.921) (0.525,3.932) (0.475,3.941) (0.425,3.950) (0.375,3.958) (0.325,3.964) (0.275,3.970) (0.225,3.975) (0.175,3.978) (0.125,3.981) (0.075,3.983) (0.025,3.984) (-0.050,3.984) (-0.100,3.982) (-0.150,3.980) (-0.200,3.977) (-0.250,3.972) (-0.300,3.967) (-0.350,3.961) (-0.400,3.954) (-0.450,3.946) (-0.500,3.937) (-0.550,3.926) (-0.600,3.915) (-0.650,3.903) (-0.700,3.890) (-0.750,3.876) (-0.800,3.860) (-0.850,3.844) (-0.900,3.826) (-0.950,3.808) (-1.000,3.788) (-1.050,3.767) (-1.100,3.744) (-1.150,3.721) (-1.200,3.696) (-1.250,3.669) (-1.300,3.641) (-1.350,3.612) (-1.400,3.581) (-1.448,3.550) (-1.500,3.514) (-1.550,3.478) (-1.600,3.440) (-1.650,3.400) (-1.700,3.357) (-1.750,3.312) (-1.790,3.275) (-1.840,3.225) (-1.875,3.189) (-1.910,3.150) (-1.950,3.105) (-1.995,3.050) (-2.034,3.000) (-2.071,2.950) (-2.100,2.908) (-2.139,2.850) (-2.170,2.800) (-2.199,2.750) (-2.225,2.703) (-2.250,2.654) (-2.275,2.603) (-2.299,2.550) (-2.320,2.500) (-2.340,2.450) (-2.358,2.400) (-2.375,2.351) (-2.391,2.300) (-2.405,2.250) (-2.417,2.200) (-2.429,2.150) (-2.439,2.100) (-2.447,2.050) (-2.455,2.000) (-2.461,1.950) (-2.466,1.900) (-2.469,1.850) (-2.471,1.800) (-2.472,1.750) (-2.471,1.700) (-2.469,1.650) (-2.466,1.600) (-2.461,1.550) (-2.455,1.500) (-2.447,1.450) (-2.438,1.400) (-2.427,1.350) (-2.415,1.300) (-2.401,1.250) (-2.386,1.200) (-2.369,1.150) (-2.350,1.101) (-2.329,1.050) (-2.306,1.000) (-2.281,0.950) (-2.253,0.900) (-2.225,0.852) (-2.192,0.800) (-2.157,0.750) (-2.125,0.708) (-2.078,0.650) (-2.033,0.600) (-1.984,0.550) (-1.931,0.500) (-1.875,0.452) (-1.825,0.413) (-1.775,0.376) (-1.725,0.343) (-1.675,0.312) (-1.625,0.283) (-1.575,0.256) (-1.525,0.231) (-1.475,0.208) (-1.425,0.187) (-1.375,0.167) (-1.327,0.150) (-1.300,0.141)}
\EQ{(1.300,0.165) (1.350,0.186) (1.400,0.208) (1.450,0.233) (1.500,0.259) (1.550,0.288) (1.600,0.319) (1.646,0.350) (1.700,0.390) (1.744,0.425) (1.800,0.474) (1.850,0.522) (1.900,0.575) (1.943,0.625) (1.975,0.665) (2.018,0.725) (2.050,0.773) (2.082,0.825) (2.109,0.875) (2.134,0.925) (2.157,0.975) (2.178,1.025) (2.197,1.075) (2.213,1.125) (2.228,1.175) (2.242,1.225) (2.253,1.275) (2.263,1.325) (2.271,1.375) (2.277,1.425) (2.282,1.475) (2.286,1.525) (2.288,1.575) (2.288,1.625) (2.288,1.675) (2.285,1.725) (2.281,1.775) (2.276,1.825) (2.269,1.875) (2.261,1.925) (2.252,1.975) (2.241,2.025) (2.228,2.075) (2.215,2.125) (2.199,2.175) (2.182,2.225) (2.164,2.275) (2.144,2.325) (2.122,2.375) (2.100,2.423) (2.075,2.473) (2.050,2.520) (2.025,2.565) (1.989,2.625) (1.956,2.675) (1.925,2.720) (1.885,2.775) (1.850,2.820) (1.805,2.875) (1.761,2.925) (1.725,2.964) (1.690,3.000) (1.650,3.039) (1.611,3.075) (1.554,3.125) (1.500,3.170) (1.450,3.209) (1.400,3.245) (1.358,3.275) (1.300,3.313) (1.250,3.344) (1.200,3.373) (1.152,3.400) (1.104,3.425) (1.053,3.450) (1.000,3.475) (0.950,3.497) (0.900,3.517) (0.850,3.536) (0.800,3.554) (0.750,3.571) (0.700,3.587) (0.650,3.601) (0.600,3.614) (0.550,3.626) (0.500,3.637) (0.450,3.647) (0.400,3.656) (0.350,3.664) (0.300,3.670) (0.250,3.676) (0.200,3.680) (0.150,3.684) (0.100,3.687) (0.050,3.688) (0.000,3.689) (-0.050,3.688) (-0.100,3.687) (-0.150,3.684) (-0.200,3.680) (-0.250,3.676) (-0.300,3.670) (-0.350,3.664) (-0.400,3.656) (-0.450,3.647) (-0.500,3.637) (-0.550,3.626) (-0.600,3.614) (-0.650,3.601) (-0.700,3.587) (-0.750,3.571) (-0.800,3.554) (-0.850,3.536) (-0.900,3.517) (-0.950,3.497) (-1.000,3.475) (-1.050,3.452) (-1.100,3.427) (-1.150,3.401) (-1.197,3.375) (-1.250,3.344) (-1.300,3.313) (-1.350,3.280) (-1.394,3.250) (-1.450,3.209) (-1.494,3.175) (-1.550,3.129) (-1.600,3.085) (-1.638,3.050) (-1.675,3.015) (-1.714,2.975) (-1.750,2.937) (-1.800,2.881) (-1.846,2.825) (-1.885,2.775) (-1.922,2.725) (-1.956,2.675) (-1.989,2.625) (-2.019,2.575) (-2.047,2.525) (-2.074,2.475) (-2.099,2.425) (-2.122,2.375) (-2.144,2.325) (-2.164,2.275) (-2.182,2.225) (-2.199,2.175) (-2.215,2.125) (-2.228,2.075) (-2.241,2.025) (-2.252,1.975) (-2.261,1.925) (-2.269,1.875) (-2.276,1.825) (-2.281,1.775) (-2.285,1.725) (-2.288,1.675) (-2.288,1.625) (-2.288,1.575) (-2.286,1.525) (-2.282,1.475) (-2.277,1.425) (-2.271,1.375) (-2.263,1.325) (-2.253,1.275) (-2.242,1.225) (-2.228,1.175) (-2.213,1.125) (-2.197,1.075) (-2.178,1.025) (-2.157,0.975) (-2.134,0.925) (-2.109,0.875) (-2.082,0.825) (-2.051,0.775) (-2.018,0.725) (-1.982,0.675) (-1.950,0.634) (-1.900,0.575) (-1.853,0.525) (-1.801,0.475) (-1.750,0.430) (-1.700,0.390) (-1.650,0.353) (-1.600,0.319) (-1.550,0.288) (-1.500,0.259) (-1.450,0.233) (-1.400,0.208) (-1.350,0.186) (-1.300,0.165)}
\EQ{(1.300,0.191) (1.350,0.215) (1.400,0.241) (1.450,0.270) (1.499,0.300) (1.550,0.335) (1.600,0.372) (1.650,0.414) (1.691,0.450) (1.742,0.500) (1.788,0.550) (1.825,0.594) (1.868,0.650) (1.900,0.696) (1.934,0.750) (1.963,0.800) (1.989,0.850) (2.013,0.900) (2.034,0.950) (2.053,1.000) (2.070,1.050) (2.085,1.100) (2.098,1.150) (2.109,1.200) (2.118,1.250) (2.126,1.300) (2.132,1.350) (2.137,1.400) (2.139,1.450) (2.141,1.500) (2.140,1.550) (2.138,1.600) (2.135,1.650) (2.130,1.700) (2.123,1.750) (2.115,1.800) (2.106,1.850) (2.095,1.900) (2.082,1.950) (2.068,2.000) (2.052,2.050) (2.035,2.100) (2.016,2.150) (1.995,2.200) (1.972,2.250) (1.950,2.296) (1.925,2.344) (1.900,2.389) (1.864,2.450) (1.831,2.500) (1.800,2.546) (1.760,2.600) (1.725,2.645) (1.679,2.700) (1.634,2.750) (1.600,2.786) (1.562,2.825) (1.525,2.860) (1.481,2.900) (1.425,2.947) (1.375,2.987) (1.325,3.025) (1.275,3.060) (1.225,3.093) (1.175,3.124) (1.132,3.150) (1.087,3.175) (1.040,3.200) (0.990,3.225) (0.936,3.250) (0.879,3.275) (0.825,3.297) (0.775,3.315) (0.725,3.333) (0.675,3.349) (0.625,3.364) (0.575,3.377) (0.525,3.390) (0.475,3.401) (0.425,3.411) (0.375,3.420) (0.325,3.427) (0.275,3.434) (0.225,3.439) (0.175,3.444) (0.125,3.447) (0.075,3.449) (0.025,3.450) (-0.037,3.450) (-0.100,3.448) (-0.150,3.445) (-0.200,3.442) (-0.250,3.437) (-0.300,3.431) (-0.350,3.424) (-0.400,3.415) (-0.450,3.406) (-0.500,3.395) (-0.550,3.384) (-0.600,3.371) (-0.650,3.356) (-0.700,3.341) (-0.748,3.325) (-0.800,3.306) (-0.850,3.287) (-0.900,3.266) (-0.950,3.244) (-1.000,3.220) (-1.050,3.195) (-1.100,3.168) (-1.150,3.139) (-1.200,3.109) (-1.250,3.077) (-1.300,3.043) (-1.350,3.006) (-1.391,2.975) (-1.450,2.927) (-1.500,2.883) (-1.550,2.836) (-1.587,2.800) (-1.625,2.760) (-1.675,2.705) (-1.721,2.650) (-1.760,2.600) (-1.797,2.550) (-1.831,2.500) (-1.864,2.450) (-1.894,2.400) (-1.922,2.350) (-1.948,2.300) (-1.972,2.250) (-1.995,2.200) (-2.016,2.150) (-2.035,2.100) (-2.052,2.050) (-2.068,2.000) (-2.082,1.950) (-2.095,1.900) (-2.106,1.850) (-2.115,1.800) (-2.123,1.750) (-2.130,1.700) (-2.135,1.650) (-2.138,1.600) (-2.140,1.550) (-2.141,1.500) (-2.139,1.450) (-2.137,1.400) (-2.132,1.350) (-2.126,1.300) (-2.118,1.250) (-2.109,1.200) (-2.098,1.150) (-2.085,1.100) (-2.070,1.050) (-2.053,1.000) (-2.034,0.950) (-2.013,0.900) (-1.989,0.850) (-1.963,0.800) (-1.934,0.750) (-1.903,0.700) (-1.868,0.650) (-1.830,0.600) (-1.788,0.550) (-1.750,0.508) (-1.700,0.459) (-1.650,0.414) (-1.603,0.375) (-1.550,0.335) (-1.500,0.301) (-1.450,0.270) (-1.400,0.241) (-1.350,0.215) (-1.300,0.191)}
\EQ{(1.300,0.216) (1.350,0.244) (1.400,0.274) (1.450,0.308) (1.500,0.345) (1.550,0.385) (1.595,0.425) (1.646,0.475) (1.692,0.525) (1.725,0.565) (1.770,0.625) (1.800,0.669) (1.834,0.725) (1.862,0.775) (1.887,0.825) (1.909,0.875) (1.929,0.925) (1.947,0.975) (1.962,1.025) (1.975,1.075) (1.987,1.125) (1.997,1.175) (2.004,1.225) (2.011,1.275) (2.015,1.325) (2.017,1.375) (2.018,1.425) (2.017,1.475) (2.015,1.525) (2.011,1.575) (2.005,1.625) (1.998,1.675) (1.989,1.725) (1.979,1.775) (1.967,1.825) (1.953,1.875) (1.937,1.925) (1.920,1.975) (1.901,2.025) (1.880,2.075) (1.858,2.125) (1.833,2.175) (1.807,2.225) (1.778,2.275) (1.750,2.321) (1.714,2.375) (1.679,2.425) (1.641,2.475) (1.601,2.525) (1.557,2.575) (1.510,2.625) (1.475,2.661) (1.434,2.700) (1.378,2.750) (1.325,2.794) (1.285,2.825) (1.225,2.869) (1.180,2.900) (1.125,2.935) (1.075,2.965) (1.025,2.993) (0.975,3.020) (0.925,3.045) (0.875,3.068) (0.825,3.089) (0.775,3.109) (0.725,3.128) (0.675,3.145) (0.625,3.161) (0.576,3.175) (0.525,3.188) (0.476,3.200) (0.425,3.211) (0.375,3.220) (0.325,3.229) (0.275,3.236) (0.225,3.241) (0.175,3.246) (0.125,3.249) (0.075,3.252) (0.025,3.253) (-0.050,3.252) (-0.100,3.251) (-0.150,3.248) (-0.200,3.244) (-0.250,3.239) (-0.300,3.232) (-0.350,3.225) (-0.400,3.216) (-0.450,3.206) (-0.500,3.195) (-0.550,3.182) (-0.600,3.168) (-0.650,3.153) (-0.700,3.137) (-0.750,3.119) (-0.798,3.100) (-0.850,3.079) (-0.900,3.056) (-0.950,3.032) (-1.000,3.007) (-1.050,2.980) (-1.100,2.950) (-1.150,2.919) (-1.200,2.886) (-1.250,2.851) (-1.300,2.814) (-1.348,2.775) (-1.400,2.731) (-1.450,2.685) (-1.486,2.650) (-1.525,2.610) (-1.575,2.555) (-1.621,2.500) (-1.661,2.450) (-1.697,2.400) (-1.731,2.350) (-1.763,2.300) (-1.793,2.250) (-1.820,2.200) (-1.846,2.150) (-1.869,2.100) (-1.891,2.050) (-1.911,2.000) (-1.929,1.950) (-1.945,1.900) (-1.960,1.850) (-1.973,1.800) (-1.984,1.750) (-1.994,1.700) (-2.002,1.650) (-2.008,1.600) (-2.013,1.550) (-2.017,1.500) (-2.018,1.450) (-2.018,1.400) (-2.016,1.350) (-2.013,1.300) (-2.008,1.250) (-2.001,1.200) (-1.992,1.150) (-1.982,1.100) (-1.969,1.050) (-1.955,1.000) (-1.938,0.950) (-1.919,0.900) (-1.898,0.850) (-1.875,0.801) (-1.850,0.753) (-1.825,0.710) (-1.787,0.650) (-1.752,0.600) (-1.713,0.550) (-1.675,0.506) (-1.625,0.454) (-1.575,0.407) (-1.525,0.364) (-1.475,0.326) (-1.425,0.291) (-1.375,0.259) (-1.325,0.230) (-1.300,0.216)}
\EQ{(1.300,0.243) (1.350,0.275) (1.400,0.310) (1.450,0.348) (1.500,0.391) (1.550,0.439) (1.600,0.493) (1.647,0.550) (1.684,0.600) (1.717,0.650) (1.747,0.700) (1.774,0.750) (1.798,0.800) (1.819,0.850) (1.838,0.900) (1.855,0.950) (1.869,1.000) (1.881,1.050) (1.892,1.100) (1.900,1.150) (1.906,1.200) (1.911,1.250) (1.914,1.300) (1.915,1.350) (1.914,1.400) (1.911,1.450) (1.907,1.500) (1.901,1.550) (1.894,1.600) (1.885,1.650) (1.873,1.700) (1.861,1.750) (1.846,1.800) (1.830,1.850) (1.812,1.900) (1.792,1.950) (1.770,2.000) (1.746,2.050) (1.720,2.100) (1.692,2.150) (1.661,2.200) (1.629,2.250) (1.593,2.300) (1.555,2.350) (1.515,2.400) (1.475,2.445) (1.425,2.498) (1.375,2.548) (1.325,2.593) (1.275,2.636) (1.226,2.675) (1.175,2.713) (1.125,2.747) (1.083,2.775) (1.025,2.810) (0.975,2.838) (0.925,2.865) (0.875,2.890) (0.825,2.913) (0.775,2.934) (0.725,2.954) (0.675,2.972) (0.625,2.989) (0.575,3.004) (0.525,3.018) (0.475,3.030) (0.425,3.042) (0.375,3.052) (0.325,3.060) (0.275,3.068) (0.225,3.074) (0.175,3.079) (0.125,3.082) (0.075,3.085) (0.025,3.086) (-0.050,3.086) (-0.100,3.084) (-0.150,3.081) (-0.200,3.076) (-0.250,3.071) (-0.300,3.064) (-0.350,3.056) (-0.400,3.047) (-0.450,3.036) (-0.500,3.024) (-0.550,3.011) (-0.600,2.996) (-0.650,2.980) (-0.700,2.963) (-0.750,2.944) (-0.800,2.923) (-0.850,2.901) (-0.900,2.877) (-0.950,2.852) (-0.999,2.825) (-1.050,2.795) (-1.100,2.764) (-1.150,2.730) (-1.192,2.700) (-1.250,2.656) (-1.300,2.615) (-1.345,2.575) (-1.398,2.525) (-1.448,2.475) (-1.493,2.425) (-1.535,2.375) (-1.575,2.325) (-1.611,2.275) (-1.645,2.225) (-1.675,2.178) (-1.706,2.125) (-1.733,2.075) (-1.758,2.025) (-1.781,1.975) (-1.802,1.925) (-1.821,1.875) (-1.838,1.825) (-1.854,1.775) (-1.867,1.725) (-1.879,1.675) (-1.889,1.625) (-1.898,1.575) (-1.905,1.525) (-1.910,1.475) (-1.913,1.425) (-1.915,1.375) (-1.914,1.325) (-1.913,1.275) (-1.909,1.225) (-1.903,1.175) (-1.896,1.125) (-1.887,1.075) (-1.875,1.025) (-1.862,0.975) (-1.847,0.925) (-1.829,0.875) (-1.809,0.825) (-1.786,0.775) (-1.761,0.725) (-1.733,0.675) (-1.701,0.625) (-1.666,0.575) (-1.627,0.525) (-1.584,0.475) (-1.536,0.425) (-1.482,0.375) (-1.425,0.328) (-1.375,0.292) (-1.325,0.258) (-1.300,0.243)}
\EQ{(1.288,0.262) (1.342,0.300) (1.400,0.347) (1.450,0.391) (1.500,0.442) (1.550,0.499) (1.590,0.550) (1.625,0.600) (1.656,0.650) (1.684,0.700) (1.708,0.750) (1.730,0.800) (1.750,0.850) (1.767,0.900) (1.781,0.950) (1.794,1.000) (1.804,1.050) (1.812,1.100) (1.818,1.150) (1.823,1.200) (1.825,1.250) (1.825,1.300) (1.824,1.350) (1.821,1.400) (1.816,1.450) (1.810,1.500) (1.801,1.550) (1.791,1.600) (1.779,1.650) (1.765,1.700) (1.750,1.748) (1.732,1.800) (1.712,1.850) (1.691,1.900) (1.667,1.950) (1.641,2.000) (1.613,2.050) (1.583,2.100) (1.550,2.149) (1.514,2.200) (1.476,2.250) (1.434,2.300) (1.400,2.339) (1.366,2.375) (1.325,2.416) (1.289,2.450) (1.250,2.484) (1.201,2.525) (1.150,2.565) (1.101,2.600) (1.050,2.634) (1.000,2.666) (0.950,2.695) (0.900,2.722) (0.850,2.748) (0.800,2.771) (0.750,2.793) (0.700,2.813) (0.650,2.831) (0.600,2.848) (0.550,2.864) (0.500,2.878) (0.450,2.890) (0.400,2.901) (0.350,2.911) (0.300,2.919) (0.250,2.926) (0.200,2.932) (0.150,2.937) (0.100,2.940) (0.050,2.942) (0.000,2.943) (-0.050,2.942) (-0.100,2.940) (-0.150,2.937) (-0.200,2.932) (-0.250,2.926) (-0.300,2.919) (-0.350,2.911) (-0.400,2.901) (-0.450,2.890) (-0.500,2.878) (-0.550,2.864) (-0.600,2.848) (-0.650,2.831) (-0.700,2.813) (-0.750,2.793) (-0.800,2.771) (-0.850,2.748) (-0.900,2.722) (-0.950,2.695) (-1.000,2.666) (-1.050,2.634) (-1.100,2.601) (-1.150,2.565) (-1.200,2.526) (-1.250,2.484) (-1.289,2.450) (-1.325,2.416) (-1.366,2.375) (-1.400,2.339) (-1.434,2.300) (-1.475,2.251) (-1.514,2.200) (-1.550,2.150) (-1.583,2.100) (-1.613,2.050) (-1.641,2.000) (-1.667,1.950) (-1.691,1.900) (-1.712,1.850) (-1.732,1.800) (-1.749,1.750) (-1.765,1.700) (-1.779,1.650) (-1.791,1.600) (-1.801,1.550) (-1.810,1.500) (-1.816,1.450) (-1.821,1.400) (-1.824,1.350) (-1.825,1.300) (-1.825,1.250) (-1.823,1.200) (-1.818,1.150) (-1.812,1.100) (-1.804,1.050) (-1.794,1.000) (-1.781,0.950) (-1.767,0.900) (-1.750,0.851) (-1.730,0.800) (-1.708,0.750) (-1.684,0.700) (-1.656,0.650) (-1.625,0.601) (-1.590,0.550) (-1.551,0.500) (-1.508,0.450) (-1.459,0.400) (-1.404,0.350) (-1.350,0.307) (-1.307,0.275) (-1.288,0.262)}
\EQ{(1.275,0.280) (1.325,0.319) (1.375,0.362) (1.414,0.400) (1.450,0.438) (1.500,0.497) (1.539,0.550) (1.572,0.600) (1.600,0.648) (1.625,0.696) (1.650,0.750) (1.670,0.800) (1.688,0.850) (1.703,0.900) (1.716,0.950) (1.726,1.000) (1.734,1.050) (1.741,1.100) (1.745,1.150) (1.747,1.200) (1.748,1.250) (1.746,1.300) (1.743,1.350) (1.738,1.400) (1.731,1.450) (1.722,1.500) (1.711,1.550) (1.698,1.600) (1.683,1.650) (1.667,1.700) (1.648,1.750) (1.627,1.800) (1.604,1.850) (1.579,1.900) (1.552,1.950) (1.525,1.995) (1.489,2.050) (1.454,2.100) (1.416,2.150) (1.375,2.200) (1.330,2.250) (1.282,2.300) (1.229,2.350) (1.175,2.397) (1.125,2.437) (1.075,2.475) (1.025,2.509) (0.975,2.542) (0.925,2.571) (0.875,2.599) (0.825,2.625) (0.775,2.649) (0.725,2.671) (0.675,2.691) (0.625,2.710) (0.575,2.727) (0.525,2.742) (0.475,2.756) (0.425,2.768) (0.375,2.779) (0.325,2.789) (0.275,2.797) (0.225,2.804) (0.175,2.809) (0.125,2.813) (0.075,2.816) (0.025,2.817) (-0.050,2.817) (-0.100,2.815) (-0.150,2.811) (-0.200,2.807) (-0.250,2.800) (-0.300,2.793) (-0.350,2.784) (-0.400,2.774) (-0.450,2.762) (-0.500,2.749) (-0.550,2.735) (-0.600,2.719) (-0.650,2.701) (-0.700,2.681) (-0.750,2.660) (-0.800,2.637) (-0.850,2.612) (-0.900,2.586) (-0.950,2.557) (-1.000,2.526) (-1.050,2.492) (-1.100,2.456) (-1.141,2.425) (-1.200,2.376) (-1.250,2.331) (-1.300,2.282) (-1.350,2.228) (-1.396,2.175) (-1.436,2.125) (-1.472,2.075) (-1.506,2.025) (-1.537,1.975) (-1.566,1.925) (-1.592,1.875) (-1.616,1.825) (-1.638,1.775) (-1.658,1.725) (-1.675,1.675) (-1.691,1.625) (-1.705,1.575) (-1.716,1.525) (-1.726,1.475) (-1.734,1.425) (-1.741,1.375) (-1.745,1.325) (-1.747,1.275) (-1.748,1.225) (-1.746,1.175) (-1.743,1.125) (-1.738,1.075) (-1.730,1.025) (-1.721,0.975) (-1.709,0.925) (-1.695,0.875) (-1.679,0.825) (-1.660,0.775) (-1.639,0.725) (-1.614,0.675) (-1.587,0.625) (-1.556,0.575) (-1.525,0.530) (-1.482,0.475) (-1.438,0.425) (-1.400,0.386) (-1.362,0.350) (-1.301,0.300) (-1.275,0.280)}
\EQ{(1.275,0.308) (1.323,0.350) (1.374,0.400) (1.419,0.450) (1.450,0.489) (1.493,0.550) (1.524,0.600) (1.550,0.647) (1.575,0.698) (1.597,0.750) (1.616,0.800) (1.632,0.850) (1.645,0.900) (1.656,0.950) (1.665,1.000) (1.672,1.050) (1.676,1.100) (1.679,1.150) (1.679,1.200) (1.678,1.250) (1.674,1.300) (1.669,1.350) (1.662,1.400) (1.652,1.450) (1.641,1.500) (1.628,1.550) (1.613,1.600) (1.595,1.650) (1.576,1.700) (1.554,1.750) (1.530,1.800) (1.504,1.850) (1.475,1.900) (1.444,1.950) (1.410,2.000) (1.375,2.047) (1.333,2.100) (1.289,2.150) (1.250,2.191) (1.200,2.240) (1.162,2.275) (1.102,2.325) (1.050,2.365) (1.001,2.400) (0.950,2.434) (0.900,2.464) (0.850,2.492) (0.800,2.518) (0.750,2.543) (0.700,2.565) (0.650,2.585) (0.600,2.604) (0.550,2.621) (0.500,2.636) (0.450,2.650) (0.400,2.662) (0.350,2.672) (0.300,2.682) (0.250,2.689) (0.200,2.696) (0.150,2.700) (0.100,2.704) (0.050,2.706) (0.000,2.707) (-0.050,2.706) (-0.100,2.704) (-0.150,2.700) (-0.200,2.696) (-0.250,2.689) (-0.300,2.682) (-0.350,2.672) (-0.400,2.662) (-0.450,2.650) (-0.500,2.636) (-0.550,2.621) (-0.600,2.604) (-0.650,2.585) (-0.700,2.565) (-0.750,2.543) (-0.800,2.518) (-0.850,2.492) (-0.900,2.464) (-0.950,2.434) (-1.000,2.401) (-1.050,2.365) (-1.100,2.327) (-1.150,2.285) (-1.189,2.250) (-1.241,2.200) (-1.289,2.150) (-1.325,2.109) (-1.373,2.050) (-1.410,2.000) (-1.444,1.950) (-1.475,1.900) (-1.504,1.850) (-1.530,1.800) (-1.554,1.750) (-1.575,1.702) (-1.595,1.650) (-1.613,1.600) (-1.628,1.550) (-1.641,1.500) (-1.652,1.450) (-1.662,1.400) (-1.669,1.350) (-1.674,1.300) (-1.678,1.250) (-1.679,1.200) (-1.679,1.150) (-1.676,1.100) (-1.672,1.050) (-1.665,1.000) (-1.656,0.950) (-1.645,0.900) (-1.632,0.850) (-1.616,0.800) (-1.597,0.750) (-1.576,0.700) (-1.552,0.650) (-1.525,0.601) (-1.493,0.550) (-1.458,0.500) (-1.425,0.458) (-1.375,0.401) (-1.325,0.352) (-1.275,0.308)}
\EQ{(1.275,0.337) (1.314,0.375) (1.350,0.414) (1.400,0.475) (1.436,0.525) (1.467,0.575) (1.495,0.625) (1.519,0.675) (1.540,0.725) (1.559,0.775) (1.574,0.825) (1.587,0.875) (1.598,0.925) (1.606,0.975) (1.612,1.025) (1.616,1.075) (1.618,1.125) (1.618,1.175) (1.616,1.225) (1.611,1.275) (1.605,1.325) (1.597,1.375) (1.586,1.425) (1.574,1.475) (1.559,1.525) (1.543,1.575) (1.524,1.625) (1.503,1.675) (1.479,1.725) (1.454,1.775) (1.425,1.825) (1.394,1.875) (1.360,1.925) (1.325,1.973) (1.283,2.025) (1.250,2.063) (1.216,2.100) (1.175,2.141) (1.125,2.187) (1.081,2.225) (1.025,2.269) (0.983,2.300) (0.925,2.339) (0.875,2.370) (0.825,2.398) (0.775,2.425) (0.725,2.449) (0.675,2.471) (0.625,2.492) (0.575,2.510) (0.525,2.527) (0.475,2.542) (0.425,2.555) (0.375,2.567) (0.325,2.578) (0.275,2.586) (0.225,2.594) (0.175,2.599) (0.125,2.604) (0.075,2.607) (0.025,2.608) (-0.050,2.608) (-0.100,2.605) (-0.150,2.602) (-0.200,2.597) (-0.250,2.590) (-0.300,2.582) (-0.350,2.573) (-0.400,2.562) (-0.450,2.549) (-0.500,2.535) (-0.550,2.519) (-0.600,2.501) (-0.650,2.482) (-0.700,2.460) (-0.750,2.437) (-0.800,2.412) (-0.850,2.384) (-0.900,2.355) (-0.946,2.325) (-1.000,2.288) (-1.050,2.250) (-1.100,2.209) (-1.139,2.175) (-1.191,2.125) (-1.239,2.075) (-1.275,2.034) (-1.323,1.975) (-1.360,1.925) (-1.394,1.875) (-1.425,1.825) (-1.450,1.781) (-1.475,1.734) (-1.500,1.681) (-1.524,1.625) (-1.543,1.575) (-1.559,1.525) (-1.574,1.475) (-1.586,1.425) (-1.597,1.375) (-1.605,1.325) (-1.611,1.275) (-1.616,1.225) (-1.618,1.175) (-1.618,1.125) (-1.616,1.075) (-1.612,1.025) (-1.606,0.975) (-1.598,0.925) (-1.587,0.875) (-1.574,0.825) (-1.559,0.775) (-1.540,0.725) (-1.519,0.675) (-1.495,0.625) (-1.467,0.575) (-1.436,0.525) (-1.400,0.475) (-1.360,0.425) (-1.325,0.387) (-1.289,0.350) (-1.275,0.337)}
\EQ{(1.266,0.359) (1.305,0.400) (1.346,0.450) (1.382,0.500) (1.414,0.550) (1.442,0.600) (1.467,0.650) (1.488,0.700) (1.506,0.750) (1.522,0.800) (1.534,0.850) (1.545,0.900) (1.553,0.950) (1.559,1.000) (1.562,1.050) (1.563,1.100) (1.563,1.150) (1.560,1.200) (1.555,1.250) (1.547,1.300) (1.538,1.350) (1.527,1.400) (1.513,1.450) (1.498,1.500) (1.480,1.550) (1.460,1.600) (1.437,1.650) (1.412,1.700) (1.385,1.750) (1.354,1.800) (1.325,1.844) (1.285,1.900) (1.250,1.944) (1.201,2.000) (1.154,2.050) (1.101,2.100) (1.050,2.145) (1.000,2.185) (0.950,2.222) (0.908,2.250) (0.850,2.287) (0.800,2.315) (0.750,2.342) (0.700,2.366) (0.650,2.388) (0.600,2.409) (0.550,2.427) (0.500,2.443) (0.450,2.458) (0.400,2.471) (0.350,2.483) (0.300,2.493) (0.250,2.501) (0.200,2.508) (0.150,2.513) (0.100,2.517) (0.050,2.519) (0.000,2.520) (-0.050,2.519) (-0.100,2.517) (-0.150,2.513) (-0.200,2.508) (-0.250,2.501) (-0.300,2.493) (-0.350,2.483) (-0.400,2.471) (-0.450,2.458) (-0.500,2.443) (-0.550,2.427) (-0.600,2.409) (-0.650,2.388) (-0.700,2.366) (-0.750,2.342) (-0.800,2.315) (-0.850,2.287) (-0.900,2.256) (-0.945,2.225) (-1.000,2.185) (-1.044,2.150) (-1.100,2.101) (-1.150,2.054) (-1.200,2.001) (-1.245,1.950) (-1.285,1.900) (-1.321,1.850) (-1.350,1.807) (-1.385,1.750) (-1.412,1.700) (-1.437,1.650) (-1.460,1.600) (-1.480,1.550) (-1.498,1.500) (-1.513,1.450) (-1.527,1.400) (-1.538,1.350) (-1.547,1.300) (-1.555,1.250) (-1.560,1.200) (-1.563,1.150) (-1.563,1.100) (-1.562,1.050) (-1.559,1.000) (-1.553,0.950) (-1.545,0.900) (-1.534,0.850) (-1.522,0.800) (-1.506,0.750) (-1.488,0.700) (-1.467,0.650) (-1.442,0.600) (-1.414,0.550) (-1.382,0.500) (-1.350,0.456) (-1.305,0.400) (-1.266,0.359)}
\EQ{(1.251,0.374) (1.295,0.425) (1.325,0.465) (1.365,0.525) (1.393,0.575) (1.418,0.625) (1.440,0.675) (1.458,0.725) (1.473,0.775) (1.486,0.825) (1.496,0.875) (1.504,0.925) (1.510,0.975) (1.513,1.025) (1.514,1.075) (1.512,1.125) (1.509,1.175) (1.503,1.225) (1.495,1.275) (1.486,1.325) (1.473,1.375) (1.459,1.425) (1.442,1.475) (1.423,1.525) (1.402,1.575) (1.378,1.625) (1.352,1.675) (1.325,1.721) (1.290,1.775) (1.255,1.825) (1.216,1.875) (1.175,1.923) (1.127,1.975) (1.075,2.025) (1.025,2.070) (0.975,2.110) (0.925,2.147) (0.875,2.182) (0.825,2.213) (0.775,2.242) (0.725,2.268) (0.675,2.292) (0.625,2.314) (0.575,2.334) (0.525,2.352) (0.475,2.369) (0.425,2.383) (0.375,2.396) (0.325,2.407) (0.275,2.416) (0.225,2.424) (0.175,2.430) (0.125,2.435) (0.075,2.438) (0.025,2.440) (-0.050,2.439) (-0.100,2.437) (-0.150,2.433) (-0.200,2.427) (-0.250,2.420) (-0.300,2.412) (-0.350,2.402) (-0.400,2.390) (-0.450,2.376) (-0.500,2.361) (-0.550,2.344) (-0.599,2.325) (-0.650,2.304) (-0.700,2.280) (-0.750,2.255) (-0.800,2.228) (-0.846,2.200) (-0.900,2.165) (-0.950,2.129) (-1.000,2.090) (-1.048,2.050) (-1.100,2.002) (-1.150,1.951) (-1.195,1.900) (-1.236,1.850) (-1.273,1.800) (-1.307,1.750) (-1.337,1.700) (-1.365,1.650) (-1.390,1.600) (-1.413,1.550) (-1.433,1.500) (-1.450,1.453) (-1.466,1.400) (-1.480,1.350) (-1.491,1.300) (-1.500,1.250) (-1.506,1.200) (-1.511,1.150) (-1.513,1.100) (-1.514,1.050) (-1.512,1.000) (-1.507,0.950) (-1.501,0.900) (-1.492,0.850) (-1.480,0.800) (-1.466,0.750) (-1.449,0.700) (-1.429,0.650) (-1.406,0.600) (-1.379,0.550) (-1.350,0.502) (-1.314,0.450) (-1.275,0.401) (-1.251,0.374)}
\EQ{(1.250,0.405) (1.285,0.450) (1.318,0.500) (1.347,0.550) (1.373,0.600) (1.395,0.650) (1.413,0.700) (1.429,0.750) (1.442,0.800) (1.452,0.850) (1.460,0.900) (1.465,0.950) (1.468,1.000) (1.468,1.050) (1.467,1.100) (1.463,1.150) (1.457,1.200) (1.448,1.250) (1.438,1.300) (1.425,1.350) (1.410,1.400) (1.392,1.450) (1.372,1.500) (1.350,1.550) (1.325,1.600) (1.300,1.645) (1.266,1.700) (1.232,1.750) (1.200,1.793) (1.154,1.850) (1.109,1.900) (1.059,1.950) (1.004,2.000) (0.950,2.044) (0.909,2.075) (0.850,2.116) (0.800,2.147) (0.751,2.175) (0.704,2.200) (0.652,2.225) (0.600,2.248) (0.550,2.268) (0.500,2.285) (0.450,2.301) (0.400,2.315) (0.350,2.328) (0.300,2.338) (0.250,2.347) (0.200,2.354) (0.150,2.360) (0.100,2.364) (0.050,2.366) (0.000,2.367) (-0.050,2.366) (-0.100,2.364) (-0.150,2.360) (-0.200,2.354) (-0.250,2.347) (-0.300,2.338) (-0.350,2.328) (-0.400,2.315) (-0.450,2.301) (-0.500,2.285) (-0.550,2.268) (-0.600,2.248) (-0.650,2.226) (-0.700,2.202) (-0.750,2.176) (-0.795,2.150) (-0.850,2.116) (-0.900,2.081) (-0.942,2.050) (-1.000,2.003) (-1.050,1.958) (-1.100,1.909) (-1.150,1.854) (-1.195,1.800) (-1.225,1.760) (-1.266,1.700) (-1.297,1.650) (-1.325,1.600) (-1.350,1.550) (-1.372,1.500) (-1.392,1.450) (-1.410,1.400) (-1.425,1.350) (-1.438,1.300) (-1.448,1.250) (-1.457,1.200) (-1.463,1.150) (-1.467,1.100) (-1.468,1.050) (-1.468,1.000) (-1.465,0.950) (-1.460,0.900) (-1.452,0.850) (-1.442,0.800) (-1.429,0.750) (-1.413,0.700) (-1.395,0.650) (-1.373,0.600) (-1.350,0.555) (-1.325,0.512) (-1.285,0.450) (-1.250,0.405)}
\EQ{(1.244,0.431) (1.274,0.475) (1.300,0.518) (1.325,0.565) (1.350,0.619) (1.372,0.675) (1.388,0.725) (1.400,0.775) (1.411,0.825) (1.419,0.875) (1.424,0.925) (1.427,0.975) (1.427,1.025) (1.425,1.075) (1.421,1.125) (1.414,1.175) (1.406,1.225) (1.395,1.275) (1.381,1.325) (1.365,1.375) (1.347,1.425) (1.326,1.475) (1.303,1.525) (1.277,1.575) (1.250,1.621) (1.215,1.675) (1.180,1.725) (1.141,1.775) (1.100,1.822) (1.050,1.874) (1.000,1.922) (0.950,1.965) (0.906,2.000) (0.850,2.040) (0.800,2.073) (0.755,2.100) (0.709,2.125) (0.660,2.150) (0.605,2.175) (0.550,2.198) (0.500,2.216) (0.450,2.233) (0.400,2.247) (0.350,2.260) (0.300,2.271) (0.250,2.280) (0.200,2.287) (0.150,2.293) (0.100,2.297) (0.050,2.299) (0.000,2.300) (-0.050,2.299) (-0.100,2.297) (-0.150,2.293) (-0.200,2.287) (-0.250,2.280) (-0.300,2.271) (-0.350,2.260) (-0.400,2.247) (-0.450,2.233) (-0.500,2.216) (-0.550,2.198) (-0.600,2.177) (-0.650,2.155) (-0.700,2.130) (-0.750,2.103) (-0.796,2.075) (-0.850,2.040) (-0.900,2.004) (-0.950,1.965) (-0.997,1.925) (-1.049,1.875) (-1.097,1.825) (-1.141,1.775) (-1.175,1.731) (-1.215,1.675) (-1.248,1.625) (-1.275,1.578) (-1.300,1.531) (-1.325,1.478) (-1.347,1.425) (-1.365,1.375) (-1.381,1.325) (-1.395,1.275) (-1.406,1.225) (-1.414,1.175) (-1.421,1.125) (-1.425,1.075) (-1.427,1.025) (-1.427,0.975) (-1.424,0.925) (-1.419,0.875) (-1.411,0.825) (-1.400,0.775) (-1.388,0.725) (-1.372,0.675) (-1.352,0.625) (-1.330,0.575) (-1.304,0.525) (-1.275,0.477) (-1.244,0.431)}
\EQ{(1.229,0.446) (1.263,0.500) (1.290,0.550) (1.313,0.600) (1.333,0.650) (1.349,0.700) (1.362,0.750) (1.373,0.800) (1.381,0.850) (1.386,0.900) (1.389,0.950) (1.389,1.000) (1.387,1.050) (1.383,1.100) (1.376,1.150) (1.367,1.200) (1.355,1.250) (1.341,1.300) (1.325,1.349) (1.306,1.400) (1.284,1.450) (1.260,1.500) (1.233,1.550) (1.203,1.600) (1.175,1.642) (1.132,1.700) (1.100,1.740) (1.050,1.796) (1.000,1.846) (0.950,1.891) (0.909,1.925) (0.850,1.970) (0.806,2.000) (0.750,2.035) (0.700,2.063) (0.650,2.089) (0.600,2.112) (0.550,2.133) (0.500,2.152) (0.450,2.169) (0.400,2.184) (0.350,2.197) (0.300,2.208) (0.250,2.218) (0.200,2.225) (0.150,2.231) (0.100,2.236) (0.050,2.238) (0.000,2.239) (-0.050,2.238) (-0.100,2.236) (-0.150,2.231) (-0.200,2.225) (-0.250,2.218) (-0.300,2.208) (-0.350,2.197) (-0.400,2.184) (-0.450,2.169) (-0.500,2.152) (-0.550,2.133) (-0.600,2.112) (-0.650,2.089) (-0.700,2.063) (-0.750,2.035) (-0.800,2.004) (-0.843,1.975) (-0.900,1.932) (-0.940,1.900) (-0.996,1.850) (-1.046,1.800) (-1.091,1.750) (-1.125,1.709) (-1.169,1.650) (-1.200,1.604) (-1.233,1.550) (-1.260,1.500) (-1.284,1.450) (-1.306,1.400) (-1.325,1.350) (-1.341,1.300) (-1.355,1.250) (-1.367,1.200) (-1.376,1.150) (-1.383,1.100) (-1.387,1.050) (-1.389,1.000) (-1.389,0.950) (-1.386,0.900) (-1.381,0.850) (-1.373,0.800) (-1.362,0.750) (-1.349,0.700) (-1.333,0.650) (-1.313,0.600) (-1.290,0.550) (-1.263,0.500) (-1.232,0.450) (-1.229,0.446)}
\EQ{(1.224,0.476) (1.250,0.522) (1.275,0.574) (1.296,0.625) (1.313,0.675) (1.326,0.725) (1.337,0.775) (1.345,0.825) (1.350,0.875) (1.353,0.925) (1.354,0.975) (1.352,1.025) (1.347,1.075) (1.340,1.125) (1.331,1.175) (1.319,1.225) (1.305,1.275) (1.288,1.325) (1.269,1.375) (1.246,1.425) (1.221,1.475) (1.194,1.525) (1.162,1.575) (1.128,1.625) (1.089,1.675) (1.050,1.721) (1.000,1.774) (0.950,1.822) (0.900,1.865) (0.856,1.900) (0.800,1.940) (0.750,1.972) (0.702,2.000) (0.656,2.025) (0.604,2.050) (0.550,2.074) (0.500,2.093) (0.450,2.111) (0.400,2.126) (0.350,2.139) (0.300,2.151) (0.250,2.161) (0.200,2.168) (0.150,2.174) (0.100,2.179) (0.050,2.181) (0.000,2.182) (-0.050,2.181) (-0.100,2.179) (-0.150,2.174) (-0.200,2.168) (-0.250,2.161) (-0.300,2.151) (-0.350,2.139) (-0.400,2.126) (-0.450,2.111) (-0.500,2.093) (-0.547,2.075) (-0.600,2.052) (-0.650,2.028) (-0.700,2.001) (-0.745,1.975) (-0.800,1.940) (-0.850,1.904) (-0.900,1.865) (-0.947,1.825) (-0.999,1.775) (-1.047,1.725) (-1.089,1.675) (-1.125,1.629) (-1.162,1.575) (-1.194,1.525) (-1.221,1.475) (-1.246,1.425) (-1.269,1.375) (-1.288,1.325) (-1.305,1.275) (-1.319,1.225) (-1.331,1.175) (-1.340,1.125) (-1.347,1.075) (-1.352,1.025) (-1.354,0.975) (-1.353,0.925) (-1.350,0.875) (-1.345,0.825) (-1.337,0.775) (-1.326,0.725) (-1.313,0.675) (-1.296,0.625) (-1.275,0.575) (-1.252,0.525) (-1.225,0.477)}

    \draw[mmaEQ] (-10,0) -- (10,0);
    % Dipolo puntual - líneas de campo
% Generado por scripts/generar_datos.py (no editar a mano).
\LC{(1.082,0.489) (1.149,0.455) (1.212,0.414) (1.269,0.365) (1.319,0.309) (1.361,0.248) (1.394,0.180) (1.417,0.109) (1.428,0.035) (1.428,-0.040) (1.415,-0.114) (1.392,-0.185) (1.358,-0.252) (1.316,-0.314) (1.265,-0.369) (1.207,-0.417) (1.144,-0.458) (1.077,-0.491)}
\LC{(0.453,1.106) (0.495,1.168) (0.539,1.229) (0.583,1.290) (0.629,1.349) (0.676,1.408) (0.724,1.465) (0.773,1.522) (0.823,1.578) (0.873,1.633) (0.925,1.688) (0.978,1.741) (1.031,1.794) (1.085,1.846) (1.141,1.896) (1.197,1.946) (1.253,1.996) (1.311,2.044) (1.369,2.091) (1.428,2.137) (1.487,2.183) (1.548,2.228) (1.609,2.271) (1.670,2.314) (1.733,2.356) (1.795,2.397) (1.859,2.437) (1.923,2.476) (1.988,2.514) (2.053,2.551) (2.118,2.587) (2.185,2.622) (2.251,2.657) (2.319,2.690) (2.386,2.722) (2.454,2.753) (2.523,2.784) (2.592,2.813) (2.661,2.841) (2.731,2.869) (2.802,2.895) (2.872,2.920) (2.943,2.945) (3.014,2.968) (3.086,2.990) (3.158,3.011) (3.230,3.031) (3.303,3.050) (3.376,3.068) (3.449,3.085) (3.522,3.101) (3.596,3.116) (3.669,3.130) (3.743,3.142) (3.817,3.154) (3.892,3.164) (3.966,3.173) (4.041,3.181) (4.115,3.188) (4.190,3.194) (4.265,3.199) (4.340,3.203) (4.415,3.205) (4.490,3.206) (4.565,3.206) (4.640,3.205) (4.715,3.203) (4.790,3.200) (4.864,3.195) (4.939,3.189) (5.014,3.182) (5.088,3.174) (5.163,3.165) (5.237,3.154) (5.311,3.142) (5.385,3.129) (5.459,3.115) (5.532,3.099) (5.605,3.082) (5.678,3.064) (5.750,3.045) (5.822,3.024) (5.894,3.002) (5.965,2.979) (6.036,2.955) (6.107,2.929) (6.177,2.902) (6.246,2.874) (6.315,2.844) (6.384,2.813) (6.451,2.781) (6.518,2.748) (6.585,2.713) (6.651,2.677) (6.716,2.640) (6.780,2.601) (6.844,2.561) (6.906,2.520) (6.968,2.478) (7.029,2.434) (7.089,2.389) (7.148,2.343) (7.207,2.296) (7.264,2.247) (7.320,2.197) (7.374,2.146) (7.428,2.094) (7.481,2.040) (7.532,1.985) (7.582,1.930) (7.631,1.873) (7.679,1.815) (7.725,1.755) (7.769,1.695) (7.813,1.634) (7.854,1.572) (7.894,1.508) (7.933,1.444) (7.970,1.379) (8.006,1.313) (8.039,1.246) (8.071,1.178) (8.102,1.109) (8.130,1.040) (8.157,0.970) (8.182,0.899) (8.205,0.828) (8.226,0.756) (8.246,0.683) (8.263,0.610) (8.278,0.537) (8.292,0.463) (8.303,0.389) (8.313,0.315) (8.320,0.240) (8.326,0.165) (8.329,0.090) (8.331,0.015) (8.330,-0.060) (8.327,-0.134) (8.323,-0.209) (8.316,-0.284) (8.307,-0.359) (8.297,-0.433) (8.284,-0.507) (8.270,-0.580) (8.253,-0.653) (8.234,-0.726) (8.214,-0.798) (8.192,-0.870) (8.168,-0.941) (8.142,-1.011) (8.114,-1.081) (8.084,-1.150) (8.053,-1.218) (8.020,-1.285) (7.985,-1.352) (7.949,-1.417) (7.911,-1.482) (7.871,-1.546) (7.830,-1.608) (7.787,-1.670) (7.743,-1.731) (7.698,-1.790) (7.651,-1.849) (7.603,-1.906) (7.553,-1.963) (7.502,-2.018) (7.450,-2.072) (7.397,-2.124) (7.342,-2.176) (7.287,-2.226) (7.230,-2.276) (7.173,-2.324) (7.114,-2.370) (7.054,-2.416) (6.994,-2.460) (6.932,-2.503) (6.870,-2.544) (6.806,-2.585) (6.742,-2.624) (6.678,-2.662) (6.612,-2.698) (6.546,-2.733) (6.479,-2.767) (6.412,-2.800) (6.343,-2.831) (6.275,-2.862) (6.206,-2.890) (6.136,-2.918) (6.066,-2.944) (5.995,-2.969) (5.924,-2.993) (5.852,-3.015) (5.780,-3.036) (5.708,-3.056) (5.635,-3.075) (5.562,-3.092) (5.489,-3.108) (5.415,-3.123) (5.342,-3.137) (5.268,-3.149) (5.194,-3.160) (5.119,-3.170) (5.045,-3.179) (4.970,-3.187) (4.895,-3.193) (4.821,-3.198) (4.746,-3.202) (4.671,-3.205) (4.596,-3.206) (4.521,-3.206) (4.446,-3.206) (4.371,-3.204) (4.296,-3.201) (4.221,-3.196) (4.146,-3.191) (4.071,-3.184) (3.997,-3.177) (3.922,-3.168) (3.848,-3.158) (3.774,-3.147) (3.700,-3.135) (3.626,-3.122) (3.552,-3.107) (3.479,-3.092) (3.406,-3.075) (3.333,-3.058) (3.260,-3.039) (3.188,-3.020) (3.116,-2.999) (3.044,-2.977) (2.973,-2.954) (2.901,-2.930) (2.831,-2.906) (2.760,-2.880) (2.690,-2.853) (2.621,-2.825) (2.551,-2.796) (2.483,-2.766) (2.414,-2.735) (2.346,-2.703) (2.279,-2.670) (2.212,-2.637) (2.146,-2.602) (2.080,-2.566) (2.014,-2.529) (1.950,-2.492) (1.885,-2.453) (1.822,-2.413) (1.758,-2.373) (1.696,-2.331) (1.634,-2.289) (1.573,-2.246) (1.512,-2.201) (1.452,-2.156) (1.393,-2.110) (1.335,-2.063) (1.277,-2.016) (1.220,-1.967) (1.164,-1.917) (1.108,-1.867) (1.053,-1.815) (1.000,-1.763) (0.947,-1.710) (0.895,-1.656) (0.843,-1.601) (0.793,-1.545) (0.744,-1.489) (0.696,-1.432) (0.648,-1.373) (0.602,-1.314) (0.557,-1.254) (0.513,-1.194) (0.470,-1.132) (0.442,-1.091)}
\LC{(-0.453,1.106) (-0.495,1.168) (-0.539,1.229) (-0.583,1.290) (-0.629,1.349) (-0.676,1.408) (-0.724,1.465) (-0.773,1.522) (-0.823,1.578) (-0.873,1.633) (-0.925,1.688) (-0.978,1.741) (-1.031,1.794) (-1.085,1.846) (-1.141,1.896) (-1.197,1.946) (-1.253,1.996) (-1.311,2.044) (-1.369,2.091) (-1.428,2.137) (-1.487,2.183) (-1.548,2.228) (-1.609,2.271) (-1.670,2.314) (-1.733,2.356) (-1.795,2.397) (-1.859,2.437) (-1.923,2.476) (-1.988,2.514) (-2.053,2.551) (-2.118,2.587) (-2.185,2.622) (-2.251,2.657) (-2.319,2.690) (-2.386,2.722) (-2.454,2.753) (-2.523,2.784) (-2.592,2.813) (-2.661,2.841) (-2.731,2.869) (-2.802,2.895) (-2.872,2.920) (-2.943,2.945) (-3.014,2.968) (-3.086,2.990) (-3.158,3.011) (-3.230,3.031) (-3.303,3.050) (-3.376,3.068) (-3.449,3.085) (-3.522,3.101) (-3.596,3.116) (-3.669,3.130) (-3.743,3.142) (-3.817,3.154) (-3.892,3.164) (-3.966,3.173) (-4.041,3.181) (-4.115,3.188) (-4.190,3.194) (-4.265,3.199) (-4.340,3.203) (-4.415,3.205) (-4.490,3.206) (-4.565,3.206) (-4.640,3.205) (-4.715,3.203) (-4.790,3.200) (-4.864,3.195) (-4.939,3.189) (-5.014,3.182) (-5.088,3.174) (-5.163,3.165) (-5.237,3.154) (-5.311,3.142) (-5.385,3.129) (-5.459,3.115) (-5.532,3.099) (-5.605,3.082) (-5.678,3.064) (-5.750,3.045) (-5.822,3.024) (-5.894,3.002) (-5.965,2.979) (-6.036,2.955) (-6.107,2.929) (-6.177,2.902) (-6.246,2.874) (-6.315,2.844) (-6.384,2.813) (-6.451,2.781) (-6.518,2.748) (-6.585,2.713) (-6.651,2.677) (-6.716,2.640) (-6.780,2.601) (-6.844,2.561) (-6.906,2.520) (-6.968,2.478) (-7.029,2.434) (-7.089,2.389) (-7.148,2.343) (-7.207,2.296) (-7.264,2.247) (-7.320,2.197) (-7.374,2.146) (-7.428,2.094) (-7.481,2.040) (-7.532,1.985) (-7.582,1.930) (-7.631,1.873) (-7.679,1.815) (-7.725,1.755) (-7.769,1.695) (-7.813,1.634) (-7.854,1.572) (-7.894,1.508) (-7.933,1.444) (-7.970,1.379) (-8.006,1.313) (-8.039,1.246) (-8.071,1.178) (-8.102,1.109) (-8.130,1.040) (-8.157,0.970) (-8.182,0.899) (-8.205,0.828) (-8.226,0.756) (-8.246,0.683) (-8.263,0.610) (-8.278,0.537) (-8.292,0.463) (-8.303,0.389) (-8.313,0.315) (-8.320,0.240) (-8.326,0.165) (-8.329,0.090) (-8.331,0.015) (-8.330,-0.060) (-8.327,-0.134) (-8.323,-0.209) (-8.316,-0.284) (-8.307,-0.359) (-8.297,-0.433) (-8.284,-0.507) (-8.270,-0.580) (-8.253,-0.653) (-8.234,-0.726) (-8.214,-0.798) (-8.192,-0.870) (-8.168,-0.941) (-8.142,-1.011) (-8.114,-1.081) (-8.084,-1.150) (-8.053,-1.218) (-8.020,-1.285) (-7.985,-1.352) (-7.949,-1.417) (-7.911,-1.482) (-7.871,-1.546) (-7.830,-1.608) (-7.787,-1.670) (-7.743,-1.731) (-7.698,-1.790) (-7.651,-1.849) (-7.603,-1.906) (-7.553,-1.963) (-7.502,-2.018) (-7.450,-2.072) (-7.397,-2.124) (-7.342,-2.176) (-7.287,-2.226) (-7.230,-2.276) (-7.173,-2.324) (-7.114,-2.370) (-7.054,-2.416) (-6.994,-2.460) (-6.932,-2.503) (-6.870,-2.544) (-6.806,-2.585) (-6.742,-2.624) (-6.678,-2.662) (-6.612,-2.698) (-6.546,-2.733) (-6.479,-2.767) (-6.412,-2.800) (-6.343,-2.831) (-6.275,-2.862) (-6.206,-2.890) (-6.136,-2.918) (-6.066,-2.944) (-5.995,-2.969) (-5.924,-2.993) (-5.852,-3.015) (-5.780,-3.036) (-5.708,-3.056) (-5.635,-3.075) (-5.562,-3.092) (-5.489,-3.108) (-5.415,-3.123) (-5.342,-3.137) (-5.268,-3.149) (-5.194,-3.160) (-5.119,-3.170) (-5.045,-3.179) (-4.970,-3.187) (-4.895,-3.193) (-4.821,-3.198) (-4.746,-3.202) (-4.671,-3.205) (-4.596,-3.206) (-4.521,-3.206) (-4.446,-3.206) (-4.371,-3.204) (-4.296,-3.201) (-4.221,-3.196) (-4.146,-3.191) (-4.071,-3.184) (-3.997,-3.177) (-3.922,-3.168) (-3.848,-3.158) (-3.774,-3.147) (-3.700,-3.135) (-3.626,-3.122) (-3.552,-3.107) (-3.479,-3.092) (-3.406,-3.075) (-3.333,-3.058) (-3.260,-3.039) (-3.188,-3.020) (-3.116,-2.999) (-3.044,-2.977) (-2.973,-2.954) (-2.901,-2.930) (-2.831,-2.906) (-2.760,-2.880) (-2.690,-2.853) (-2.621,-2.825) (-2.551,-2.796) (-2.483,-2.766) (-2.414,-2.735) (-2.346,-2.703) (-2.279,-2.670) (-2.212,-2.637) (-2.146,-2.602) (-2.080,-2.566) (-2.014,-2.529) (-1.950,-2.492) (-1.885,-2.453) (-1.822,-2.413) (-1.758,-2.373) (-1.696,-2.331) (-1.634,-2.289) (-1.573,-2.246) (-1.512,-2.201) (-1.452,-2.156) (-1.393,-2.110) (-1.335,-2.063) (-1.277,-2.016) (-1.220,-1.967) (-1.164,-1.917) (-1.108,-1.867) (-1.053,-1.815) (-1.000,-1.763) (-0.947,-1.710) (-0.895,-1.656) (-0.843,-1.601) (-0.793,-1.545) (-0.744,-1.489) (-0.696,-1.432) (-0.648,-1.373) (-0.602,-1.314) (-0.557,-1.254) (-0.513,-1.194) (-0.470,-1.132) (-0.442,-1.091)}
\LC{(-1.082,0.489) (-1.149,0.455) (-1.212,0.414) (-1.269,0.365) (-1.319,0.309) (-1.361,0.248) (-1.394,0.180) (-1.417,0.109) (-1.428,0.035) (-1.428,-0.040) (-1.415,-0.114) (-1.392,-0.185) (-1.358,-0.252) (-1.316,-0.314) (-1.265,-0.369) (-1.207,-0.417) (-1.144,-0.458) (-1.077,-0.491)}
\LC{(4.893,-10.000) (4.845,-9.951) (4.797,-9.893) (4.749,-9.835) (4.702,-9.777) (4.655,-9.719) (4.607,-9.661) (4.560,-9.602) (4.514,-9.544) (4.467,-9.485) (4.420,-9.426) (4.374,-9.368) (4.328,-9.309) (4.281,-9.249) (4.235,-9.190) (4.190,-9.131) (4.144,-9.071) (4.099,-9.012) (4.053,-8.952) (4.008,-8.892) (3.963,-8.832) (3.918,-8.772) (3.873,-8.712) (3.829,-8.651) (3.785,-8.591) (3.740,-8.530) (3.696,-8.470) (3.652,-8.409) (3.609,-8.348) (3.565,-8.287) (3.522,-8.226) (3.479,-8.164) (3.435,-8.103) (3.393,-8.041) (3.350,-7.980) (3.307,-7.918) (3.265,-7.856) (3.223,-7.794) (3.181,-7.732) (3.139,-7.670) (3.097,-7.607) (3.056,-7.545) (3.015,-7.482) (2.973,-7.419) (2.932,-7.357) (2.892,-7.294) (2.851,-7.231) (2.811,-7.167) (2.771,-7.104) (2.730,-7.041) (2.691,-6.977) (2.651,-6.913) (2.612,-6.850) (2.572,-6.786) (2.533,-6.722) (2.494,-6.658) (2.456,-6.593) (2.417,-6.529) (2.379,-6.465) (2.341,-6.400) (2.303,-6.335) (2.265,-6.270) (2.228,-6.205) (2.190,-6.140) (2.153,-6.075) (2.116,-6.010) (2.080,-5.945) (2.043,-5.879) (2.007,-5.813) (1.971,-5.748) (1.935,-5.682) (1.899,-5.616) (1.864,-5.550) (1.829,-5.483) (1.794,-5.417) (1.759,-5.351) (1.724,-5.284) (1.690,-5.217) (1.656,-5.151) (1.622,-5.084) (1.588,-5.017) (1.555,-4.950) (1.522,-4.882) (1.489,-4.815) (1.456,-4.747) (1.423,-4.680) (1.391,-4.612) (1.359,-4.544) (1.327,-4.476) (1.296,-4.408) (1.265,-4.340) (1.234,-4.272) (1.203,-4.203) (1.172,-4.135) (1.142,-4.066) (1.112,-3.998) (1.082,-3.929) (1.053,-3.860) (1.024,-3.791) (0.995,-3.721) (0.966,-3.652) (0.938,-3.583) (0.910,-3.513) (0.882,-3.444) (0.854,-3.374) (0.827,-3.304) (0.800,-3.234) (0.773,-3.164) (0.747,-3.094) (0.721,-3.023) (0.695,-2.953) (0.670,-2.882) (0.645,-2.812) (0.620,-2.741) (0.595,-2.670) (0.571,-2.599) (0.547,-2.528) (0.524,-2.457) (0.501,-2.385) (0.478,-2.314) (0.455,-2.242) (0.433,-2.171) (0.412,-2.099) (0.390,-2.027) (0.369,-1.955) (0.349,-1.883) (0.328,-1.811) (0.309,-1.738) (0.289,-1.666) (0.283,-1.642)}
\LC{(-5.042,-10.000) (-4.993,-9.951) (-4.944,-9.894) (-4.895,-9.837) (-4.847,-9.780) (-4.798,-9.722) (-4.750,-9.665) (-4.702,-9.607) (-4.654,-9.550) (-4.606,-9.492) (-4.558,-9.434) (-4.511,-9.376) (-4.463,-9.318) (-4.416,-9.259) (-4.369,-9.201) (-4.322,-9.142) (-4.276,-9.084) (-4.229,-9.025) (-4.183,-8.966) (-4.136,-8.907) (-4.090,-8.848) (-4.044,-8.789) (-3.999,-8.729) (-3.953,-8.670) (-3.908,-8.610) (-3.862,-8.550) (-3.817,-8.490) (-3.772,-8.430) (-3.728,-8.370) (-3.683,-8.310) (-3.639,-8.249) (-3.594,-8.189) (-3.550,-8.128) (-3.506,-8.067) (-3.463,-8.006) (-3.419,-7.945) (-3.376,-7.884) (-3.333,-7.823) (-3.289,-7.761) (-3.247,-7.700) (-3.204,-7.638) (-3.161,-7.576) (-3.119,-7.514) (-3.077,-7.452) (-3.035,-7.390) (-2.993,-7.328) (-2.952,-7.265) (-2.910,-7.203) (-2.869,-7.140) (-2.828,-7.077) (-2.787,-7.015) (-2.747,-6.951) (-2.706,-6.888) (-2.666,-6.825) (-2.626,-6.762) (-2.586,-6.698) (-2.546,-6.634) (-2.507,-6.571) (-2.467,-6.507) (-2.428,-6.443) (-2.389,-6.379) (-2.351,-6.314) (-2.312,-6.250) (-2.274,-6.186) (-2.236,-6.121) (-2.198,-6.056) (-2.160,-5.991) (-2.123,-5.926) (-2.086,-5.861) (-2.049,-5.796) (-2.012,-5.731) (-1.975,-5.665) (-1.939,-5.600) (-1.903,-5.534) (-1.867,-5.468) (-1.831,-5.402) (-1.796,-5.336) (-1.760,-5.270) (-1.725,-5.204) (-1.690,-5.137) (-1.656,-5.071) (-1.621,-5.004) (-1.587,-4.937) (-1.553,-4.870) (-1.520,-4.803) (-1.486,-4.736) (-1.453,-4.669) (-1.420,-4.601) (-1.388,-4.534) (-1.355,-4.466) (-1.323,-4.399) (-1.291,-4.331) (-1.260,-4.263) (-1.228,-4.195) (-1.197,-4.126) (-1.166,-4.058) (-1.135,-3.990) (-1.105,-3.921) (-1.075,-3.852) (-1.045,-3.783) (-1.016,-3.714) (-0.986,-3.645) (-0.957,-3.576) (-0.929,-3.507) (-0.900,-3.438) (-0.872,-3.368) (-0.844,-3.298) (-0.817,-3.229) (-0.790,-3.159) (-0.763,-3.089) (-0.736,-3.019) (-0.710,-2.948) (-0.684,-2.878) (-0.658,-2.808) (-0.633,-2.737) (-0.608,-2.666) (-0.583,-2.595) (-0.559,-2.524) (-0.535,-2.453) (-0.511,-2.382) (-0.488,-2.311) (-0.465,-2.240) (-0.442,-2.168) (-0.420,-2.096) (-0.398,-2.025) (-0.377,-1.953) (-0.356,-1.881) (-0.335,-1.809) (-0.315,-1.736) (-0.295,-1.664)}
\LC{(10.000,-6.289) (9.931,-6.297) (9.856,-6.305) (9.782,-6.311) (9.707,-6.317) (9.632,-6.323) (9.557,-6.328) (9.482,-6.332) (9.407,-6.336) (9.332,-6.339) (9.257,-6.342) (9.182,-6.344) (9.107,-6.345) (9.032,-6.346) (8.957,-6.346) (8.882,-6.346) (8.808,-6.344) (8.733,-6.343) (8.658,-6.341) (8.583,-6.338) (8.508,-6.335) (8.433,-6.331) (8.358,-6.326) (8.283,-6.321) (8.208,-6.315) (8.134,-6.309) (8.059,-6.302) (7.984,-6.295) (7.910,-6.287) (7.835,-6.278) (7.761,-6.269) (7.686,-6.260) (7.612,-6.249) (7.538,-6.239) (7.464,-6.227) (7.390,-6.215) (7.316,-6.203) (7.242,-6.190) (7.168,-6.176) (7.094,-6.162) (7.021,-6.148) (6.947,-6.133) (6.874,-6.117) (6.801,-6.100) (6.728,-6.084) (6.655,-6.066) (6.582,-6.048) (6.509,-6.030) (6.437,-6.011) (6.364,-5.991) (6.292,-5.971) (6.220,-5.951) (6.148,-5.930) (6.076,-5.908) (6.004,-5.886) (5.933,-5.863) (5.862,-5.840) (5.791,-5.816) (5.720,-5.792) (5.649,-5.767) (5.578,-5.742) (5.508,-5.716) (5.438,-5.690) (5.367,-5.663) (5.298,-5.636) (5.228,-5.608) (5.159,-5.579) (5.089,-5.550) (5.020,-5.521) (4.952,-5.491) (4.883,-5.461) (4.815,-5.430) (4.747,-5.399) (4.679,-5.367) (4.611,-5.334) (4.544,-5.301) (4.476,-5.268) (4.410,-5.234) (4.343,-5.200) (4.276,-5.165) (4.210,-5.130) (4.144,-5.094) (4.079,-5.058) (4.013,-5.021) (3.948,-4.984) (3.883,-4.946) (3.819,-4.908) (3.755,-4.869) (3.691,-4.830) (3.627,-4.790) (3.564,-4.750) (3.501,-4.709) (3.438,-4.668) (3.375,-4.627) (3.313,-4.585) (3.251,-4.542) (3.190,-4.499) (3.129,-4.456) (3.068,-4.412) (3.008,-4.368) (2.947,-4.323) (2.888,-4.278) (2.828,-4.232) (2.769,-4.186) (2.710,-4.139) (2.652,-4.092) (2.594,-4.045) (2.536,-3.997) (2.479,-3.948) (2.422,-3.899) (2.366,-3.850) (2.309,-3.800) (2.254,-3.750) (2.198,-3.699) (2.144,-3.648) (2.089,-3.597) (2.035,-3.545) (1.981,-3.492) (1.928,-3.439) (1.875,-3.386) (1.823,-3.332) (1.771,-3.278) (1.720,-3.223) (1.669,-3.168) (1.618,-3.113) (1.568,-3.057) (1.519,-3.001) (1.470,-2.944) (1.421,-2.887) (1.373,-2.829) (1.326,-2.771) (1.279,-2.713) (1.233,-2.654) (1.187,-2.594) (1.141,-2.535) (1.097,-2.474) (1.052,-2.414)}
\LC{(-2.969,2.475) (-3.043,2.485) (-3.118,2.495) (-3.192,2.503) (-3.267,2.509) (-3.342,2.515) (-3.417,2.518) (-3.492,2.520) (-3.567,2.521) (-3.642,2.520) (-3.717,2.518) (-3.792,2.514) (-3.866,2.509) (-3.941,2.502) (-4.016,2.494) (-4.090,2.484) (-4.164,2.472) (-4.238,2.459) (-4.312,2.445) (-4.385,2.428) (-4.458,2.411) (-4.530,2.391) (-4.602,2.370) (-4.674,2.348) (-4.745,2.323) (-4.815,2.298) (-4.885,2.270) (-4.954,2.241) (-5.022,2.210) (-5.090,2.178) (-5.157,2.144) (-5.223,2.108) (-5.288,2.071) (-5.352,2.032) (-5.415,1.991) (-5.477,1.949) (-5.538,1.905) (-5.598,1.860) (-5.656,1.813) (-5.713,1.765) (-5.769,1.714) (-5.823,1.663) (-5.876,1.610) (-5.928,1.555) (-5.977,1.499) (-6.025,1.441) (-6.072,1.382) (-6.116,1.322) (-6.159,1.260) (-6.200,1.197) (-6.239,1.133) (-6.275,1.068) (-6.310,1.001) (-6.343,0.934) (-6.373,0.865) (-6.401,0.796) (-6.427,0.725) (-6.450,0.654) (-6.471,0.582) (-6.490,0.509) (-6.506,0.436) (-6.520,0.362) (-6.531,0.288) (-6.539,0.214) (-6.545,0.139) (-6.549,0.064) (-6.550,-0.011) (-6.548,-0.086) (-6.544,-0.161) (-6.537,-0.236) (-6.528,-0.310) (-6.516,-0.384) (-6.502,-0.458) (-6.485,-0.531) (-6.465,-0.603) (-6.444,-0.675) (-6.419,-0.746) (-6.393,-0.816) (-6.364,-0.885) (-6.333,-0.954) (-6.300,-1.021) (-6.265,-1.087) (-6.227,-1.152) (-6.188,-1.216) (-6.147,-1.279) (-6.103,-1.340) (-6.058,-1.400) (-6.012,-1.458) (-5.963,-1.515) (-5.913,-1.571) (-5.861,-1.625) (-5.808,-1.678) (-5.753,-1.729) (-5.697,-1.779) (-5.639,-1.827) (-5.580,-1.874) (-5.520,-1.918) (-5.459,-1.962) (-5.397,-2.004) (-5.333,-2.044) (-5.269,-2.082) (-5.203,-2.119) (-5.137,-2.154) (-5.070,-2.188) (-5.002,-2.219) (-4.934,-2.250) (-4.864,-2.278) (-4.794,-2.305) (-4.724,-2.331) (-4.653,-2.355) (-4.581,-2.377) (-4.509,-2.397) (-4.436,-2.416) (-4.363,-2.433) (-4.290,-2.449) (-4.216,-2.463) (-4.142,-2.476) (-4.068,-2.487) (-3.994,-2.496) (-3.919,-2.504) (-3.845,-2.511) (-3.770,-2.515) (-3.695,-2.519) (-3.620,-2.521) (-3.545,-2.521) (-3.470,-2.520) (-3.395,-2.517) (-3.320,-2.513) (-3.245,-2.508) (-3.170,-2.501) (-3.096,-2.492) (-3.022,-2.482) (-2.972,-2.475)}
\LC{(0.389,1.769) (0.413,1.840) (0.437,1.911) (0.463,1.982) (0.488,2.052) (0.514,2.123) (0.541,2.193) (0.568,2.263) (0.595,2.333) (0.623,2.402) (0.651,2.472) (0.680,2.541) (0.709,2.610) (0.738,2.679) (0.768,2.748) (0.798,2.817) (0.829,2.885) (0.860,2.954) (0.891,3.022) (0.923,3.090) (0.955,3.158) (0.987,3.225) (1.020,3.293) (1.053,3.360) (1.086,3.427) (1.120,3.494) (1.154,3.561) (1.189,3.628) (1.223,3.694) (1.258,3.760) (1.294,3.827) (1.329,3.893) (1.365,3.958) (1.402,4.024) (1.438,4.089) (1.475,4.155) (1.512,4.220) (1.550,4.285) (1.588,4.350) (1.626,4.414) (1.664,4.479) (1.703,4.543) (1.742,4.607) (1.781,4.671) (1.821,4.735) (1.860,4.798) (1.900,4.862) (1.941,4.925) (1.981,4.988) (2.022,5.051) (2.063,5.113) (2.105,5.176) (2.146,5.238) (2.188,5.301) (2.230,5.363) (2.273,5.424) (2.316,5.486) (2.358,5.548) (2.402,5.609) (2.445,5.670) (2.489,5.731) (2.533,5.792) (2.577,5.853) (2.621,5.913) (2.666,5.973) (2.711,6.033) (2.756,6.093) (2.801,6.153) (2.847,6.213) (2.892,6.272) (2.938,6.331) (2.985,6.390) (3.031,6.449) (3.078,6.508) (3.125,6.566) (3.172,6.625) (3.219,6.683) (3.267,6.741) (3.315,6.799) (3.363,6.856) (3.411,6.914) (3.459,6.971) (3.508,7.028) (3.557,7.085) (3.606,7.142) (3.655,7.198) (3.705,7.255) (3.754,7.311) (3.804,7.367) (3.854,7.423) (3.905,7.478) (3.955,7.534) (4.006,7.589) (4.057,7.644) (4.108,7.699) (4.159,7.754) (4.210,7.808) (4.262,7.863) (4.314,7.917) (4.366,7.971) (4.418,8.025) (4.471,8.078) (4.523,8.132) (4.576,8.185) (4.629,8.238) (4.682,8.291) (4.735,8.344) (4.789,8.396) (4.843,8.449) (4.896,8.501) (4.950,8.553) (5.005,8.605) (5.059,8.656) (5.114,8.708) (5.168,8.759) (5.223,8.810) (5.278,8.861) (5.334,8.912) (5.389,8.962) (5.445,9.013) (5.500,9.063) (5.556,9.113) (5.613,9.163) (5.669,9.212) (5.725,9.262) (5.782,9.311) (5.839,9.360) (5.895,9.409) (5.953,9.457) (6.010,9.506) (6.067,9.554) (6.125,9.602) (6.183,9.650) (6.240,9.698) (6.298,9.745) (6.357,9.793) (6.415,9.840) (6.473,9.887) (6.532,9.933) (6.591,9.980) (6.630,10.000)}
\LC{(-3.797,3.606) (-3.869,3.626) (-3.942,3.645) (-4.014,3.663) (-4.087,3.680) (-4.161,3.696) (-4.234,3.712) (-4.307,3.726) (-4.381,3.740) (-4.455,3.753) (-4.529,3.765) (-4.603,3.775) (-4.678,3.785) (-4.752,3.794) (-4.827,3.802) (-4.901,3.810) (-4.976,3.816) (-5.051,3.821) (-5.126,3.825) (-5.201,3.829) (-5.276,3.831) (-5.351,3.832) (-5.426,3.833) (-5.501,3.832) (-5.576,3.831) (-5.651,3.828) (-5.726,3.825) (-5.800,3.820) (-5.875,3.815) (-5.950,3.808) (-6.025,3.801) (-6.099,3.792) (-6.173,3.782) (-6.248,3.772) (-6.322,3.760) (-6.396,3.748) (-6.470,3.734) (-6.543,3.719) (-6.616,3.703) (-6.689,3.687) (-6.762,3.669) (-6.835,3.650) (-6.907,3.630) (-6.979,3.609) (-7.051,3.587) (-7.122,3.564) (-7.193,3.540) (-7.264,3.514) (-7.334,3.488) (-7.404,3.461) (-7.473,3.432) (-7.542,3.402) (-7.611,3.372) (-7.679,3.340) (-7.746,3.307) (-7.813,3.273) (-7.879,3.239) (-7.945,3.202) (-8.010,3.165) (-8.075,3.127) (-8.139,3.088) (-8.202,3.048) (-8.265,3.006) (-8.326,2.964) (-8.387,2.920) (-8.448,2.876) (-8.507,2.830) (-8.566,2.784) (-8.624,2.736) (-8.681,2.687) (-8.737,2.637) (-8.792,2.587) (-8.846,2.535) (-8.900,2.482) (-8.952,2.428) (-9.003,2.373) (-9.054,2.318) (-9.103,2.261) (-9.151,2.203) (-9.198,2.145) (-9.243,2.085) (-9.288,2.025) (-9.331,1.964) (-9.373,1.902) (-9.414,1.839) (-9.454,1.775) (-9.492,1.711) (-9.529,1.645) (-9.564,1.579) (-9.598,1.513) (-9.631,1.445) (-9.662,1.377) (-9.692,1.308) (-9.720,1.239) (-9.747,1.169) (-9.772,1.098) (-9.796,1.027) (-9.818,0.955) (-9.839,0.883) (-9.858,0.810) (-9.875,0.737) (-9.891,0.664) (-9.905,0.590) (-9.918,0.517) (-9.928,0.442) (-9.938,0.368) (-9.945,0.293) (-9.951,0.218) (-9.955,0.144) (-9.957,0.069) (-9.958,-0.006) (-9.957,-0.081) (-9.954,-0.156) (-9.950,-0.231) (-9.944,-0.306) (-9.936,-0.381) (-9.927,-0.455) (-9.916,-0.529) (-9.903,-0.603) (-9.888,-0.677) (-9.872,-0.750) (-9.855,-0.823) (-9.836,-0.895) (-9.815,-0.967) (-9.792,-1.039) (-9.768,-1.110) (-9.743,-1.180) (-9.716,-1.250) (-9.687,-1.320) (-9.657,-1.388) (-9.626,-1.457) (-9.593,-1.524) (-9.558,-1.591) (-9.522,-1.656) (-9.485,-1.722) (-9.447,-1.786) (-9.407,-1.850) (-9.366,-1.912) (-9.324,-1.974) (-9.280,-2.035) (-9.236,-2.096) (-9.190,-2.155) (-9.143,-2.213) (-9.094,-2.271) (-9.045,-2.327) (-8.995,-2.383) (-8.943,-2.437) (-8.891,-2.491) (-8.837,-2.544) (-8.783,-2.595) (-8.728,-2.646) (-8.671,-2.695) (-8.614,-2.744) (-8.556,-2.792) (-8.497,-2.838) (-8.438,-2.883) (-8.377,-2.928) (-8.316,-2.971) (-8.254,-3.013) (-8.191,-3.055) (-8.128,-3.095) (-8.064,-3.134) (-7.999,-3.172) (-7.934,-3.209) (-7.868,-3.245) (-7.802,-3.279) (-7.735,-3.313) (-7.667,-3.346) (-7.599,-3.377) (-7.530,-3.408) (-7.461,-3.437) (-7.392,-3.465) (-7.322,-3.492) (-7.252,-3.519) (-7.181,-3.544) (-7.110,-3.568) (-7.039,-3.591) (-6.967,-3.613) (-6.895,-3.633) (-6.823,-3.653) (-6.750,-3.672) (-6.677,-3.690) (-6.604,-3.706) (-6.531,-3.722) (-6.457,-3.736) (-6.383,-3.750) (-6.309,-3.762) (-6.235,-3.774) (-6.161,-3.784) (-6.086,-3.794) (-6.012,-3.802) (-5.937,-3.809) (-5.863,-3.816) (-5.788,-3.821) (-5.713,-3.825) (-5.638,-3.829) (-5.563,-3.831) (-5.488,-3.832) (-5.413,-3.833) (-5.338,-3.832) (-5.263,-3.831) (-5.188,-3.828) (-5.113,-3.825) (-5.038,-3.820) (-4.964,-3.815) (-4.889,-3.808) (-4.814,-3.801) (-4.740,-3.793) (-4.665,-3.784) (-4.591,-3.774) (-4.517,-3.763) (-4.443,-3.751) (-4.369,-3.738) (-4.295,-3.724) (-4.221,-3.709) (-4.148,-3.694) (-4.075,-3.677) (-4.002,-3.660) (-3.929,-3.641) (-3.857,-3.622) (-3.809,-3.609)}
\LC{(-0.048,-10.000) (-0.048,-9.947) (-0.047,-9.872) (-0.046,-9.797) (-0.046,-9.722) (-0.045,-9.647) (-0.045,-9.572) (-0.044,-9.497) (-0.044,-9.422) (-0.043,-9.347) (-0.043,-9.272) (-0.042,-9.197) (-0.042,-9.122) (-0.041,-9.047) (-0.041,-8.972) (-0.040,-8.897) (-0.040,-8.822) (-0.039,-8.747) (-0.039,-8.672) (-0.038,-8.597) (-0.038,-8.522) (-0.037,-8.447) (-0.037,-8.372) (-0.036,-8.297) (-0.036,-8.222) (-0.035,-8.147) (-0.035,-8.072) (-0.034,-7.997) (-0.034,-7.922) (-0.033,-7.847) (-0.033,-7.772) (-0.032,-7.697) (-0.032,-7.622) (-0.031,-7.547) (-0.031,-7.472) (-0.030,-7.397) (-0.030,-7.322) (-0.030,-7.247) (-0.029,-7.172) (-0.029,-7.097) (-0.028,-7.022) (-0.028,-6.947) (-0.027,-6.872) (-0.027,-6.797) (-0.026,-6.722) (-0.026,-6.647) (-0.026,-6.572) (-0.025,-6.497) (-0.025,-6.422) (-0.024,-6.347) (-0.024,-6.272) (-0.023,-6.197) (-0.023,-6.122) (-0.023,-6.047) (-0.022,-5.972) (-0.022,-5.897) (-0.021,-5.822) (-0.021,-5.747) (-0.020,-5.672) (-0.020,-5.597) (-0.020,-5.522) (-0.019,-5.447) (-0.019,-5.372) (-0.018,-5.297) (-0.018,-5.222) (-0.018,-5.147) (-0.017,-5.072) (-0.017,-4.997) (-0.017,-4.922) (-0.016,-4.847) (-0.016,-4.772) (-0.015,-4.697) (-0.015,-4.622) (-0.015,-4.547) (-0.014,-4.472) (-0.014,-4.397) (-0.014,-4.322) (-0.013,-4.247) (-0.013,-4.172) (-0.013,-4.097) (-0.012,-4.022) (-0.012,-3.947) (-0.012,-3.872) (-0.011,-3.797) (-0.011,-3.722) (-0.011,-3.647) (-0.010,-3.572) (-0.010,-3.497) (-0.010,-3.422) (-0.009,-3.347) (-0.009,-3.272) (-0.009,-3.197) (-0.008,-3.122) (-0.008,-3.047) (-0.008,-2.972) (-0.007,-2.897) (-0.007,-2.822) (-0.007,-2.747) (-0.007,-2.672) (-0.006,-2.597) (-0.006,-2.522) (-0.006,-2.447) (-0.006,-2.372) (-0.005,-2.297)}
\LC{(-0.488,1.887) (-0.516,1.956) (-0.545,2.025) (-0.574,2.094) (-0.604,2.163) (-0.634,2.232) (-0.665,2.300) (-0.696,2.369) (-0.727,2.437) (-0.759,2.504) (-0.792,2.572) (-0.825,2.639) (-0.858,2.707) (-0.892,2.774) (-0.926,2.840) (-0.961,2.907) (-0.996,2.973) (-1.031,3.039) (-1.067,3.105) (-1.103,3.171) (-1.140,3.236) (-1.177,3.301) (-1.214,3.366) (-1.252,3.431) (-1.290,3.496) (-1.329,3.560) (-1.368,3.624) (-1.407,3.688) (-1.446,3.752) (-1.486,3.816) (-1.526,3.879) (-1.567,3.942) (-1.608,4.005) (-1.649,4.067) (-1.691,4.130) (-1.733,4.192) (-1.775,4.254) (-1.818,4.316) (-1.860,4.377) (-1.904,4.439) (-1.947,4.500) (-1.991,4.560) (-2.035,4.621) (-2.080,4.682) (-2.124,4.742) (-2.169,4.802) (-2.215,4.861) (-2.260,4.921) (-2.306,4.980) (-2.352,5.039) (-2.399,5.098) (-2.446,5.157) (-2.493,5.215) (-2.540,5.273) (-2.588,5.331) (-2.635,5.389) (-2.684,5.447) (-2.732,5.504) (-2.781,5.561) (-2.830,5.618) (-2.879,5.674) (-2.928,5.731) (-2.978,5.787) (-3.028,5.843) (-3.078,5.899) (-3.129,5.954) (-3.179,6.009) (-3.230,6.064) (-3.282,6.119) (-3.333,6.174) (-3.385,6.228) (-3.437,6.282) (-3.489,6.336) (-3.542,6.389) (-3.594,6.443) (-3.647,6.496) (-3.700,6.549) (-3.754,6.602) (-3.807,6.654) (-3.861,6.706) (-3.915,6.758) (-3.969,6.810) (-4.024,6.862) (-4.079,6.913) (-4.134,6.964) (-4.189,7.015) (-4.244,7.065) (-4.300,7.116) (-4.355,7.166) (-4.411,7.216) (-4.468,7.265) (-4.524,7.315) (-4.581,7.364) (-4.638,7.413) (-4.695,7.461) (-4.752,7.510) (-4.809,7.558) (-4.867,7.606) (-4.925,7.654) (-4.983,7.701) (-5.041,7.749) (-5.100,7.796) (-5.158,7.842) (-5.217,7.889) (-5.276,7.935) (-5.335,7.981) (-5.395,8.027) (-5.454,8.073) (-5.514,8.118) (-5.574,8.163) (-5.634,8.208) (-5.694,8.252) (-5.755,8.297) (-5.816,8.341) (-5.876,8.385) (-5.937,8.428) (-5.999,8.472) (-6.060,8.515) (-6.122,8.558) (-6.183,8.600) (-6.245,8.643) (-6.307,8.685) (-6.369,8.727) (-6.432,8.768) (-6.494,8.810) (-6.557,8.851) (-6.620,8.892) (-6.683,8.932) (-6.746,8.973) (-6.810,9.013) (-6.873,9.053) (-6.937,9.092) (-7.001,9.132) (-7.065,9.171) (-7.129,9.210) (-7.193,9.248) (-7.257,9.287) (-7.322,9.325) (-7.387,9.363) (-7.452,9.400) (-7.517,9.437) (-7.582,9.475) (-7.647,9.511) (-7.713,9.548) (-7.778,9.584) (-7.844,9.620) (-7.910,9.656) (-7.976,9.692) (-8.042,9.727) (-8.109,9.762) (-8.175,9.797) (-8.242,9.831) (-8.308,9.866) (-8.375,9.900) (-8.442,9.933) (-8.509,9.967) (-8.577,10.000) (-8.599,10.000)}
\LC{(-1.957,3.712) (-2.008,3.767) (-2.059,3.821) (-2.111,3.876) (-2.163,3.930) (-2.215,3.983) (-2.268,4.036) (-2.322,4.089) (-2.375,4.142) (-2.429,4.194) (-2.484,4.245) (-2.538,4.297) (-2.593,4.348) (-2.649,4.398) (-2.704,4.448) (-2.761,4.498) (-2.817,4.547) (-2.874,4.596) (-2.931,4.645) (-2.988,4.693) (-3.046,4.741) (-3.104,4.789) (-3.162,4.836) (-3.221,4.883) (-3.280,4.929) (-3.339,4.975) (-3.399,5.020) (-3.459,5.065) (-3.519,5.110) (-3.579,5.155) (-3.640,5.199) (-3.701,5.242) (-3.763,5.285) (-3.824,5.328) (-3.886,5.370) (-3.948,5.412) (-4.011,5.454) (-4.073,5.495) (-4.136,5.536) (-4.199,5.576) (-4.263,5.616) (-4.327,5.656) (-4.391,5.695) (-4.455,5.734) (-4.519,5.772) (-4.584,5.810) (-4.649,5.848) (-4.714,5.885) (-4.780,5.921) (-4.845,5.958) (-4.911,5.994) (-4.977,6.029) (-5.043,6.064) (-5.110,6.099) (-5.177,6.133) (-5.244,6.167) (-5.311,6.200) (-5.378,6.233) (-5.446,6.265) (-5.514,6.298) (-5.582,6.329) (-5.650,6.360) (-5.718,6.391) (-5.787,6.422) (-5.855,6.452) (-5.924,6.481) (-5.994,6.510) (-6.063,6.539) (-6.132,6.567) (-6.202,6.595) (-6.272,6.622) (-6.342,6.649) (-6.412,6.676) (-6.482,6.702) (-6.553,6.727) (-6.623,6.753) (-6.694,6.777) (-6.765,6.802) (-6.836,6.825) (-6.908,6.849) (-6.979,6.872) (-7.051,6.894) (-7.122,6.916) (-7.194,6.938) (-7.266,6.959) (-7.338,6.980) (-7.410,7.000) (-7.483,7.020) (-7.555,7.039) (-7.628,7.058) (-7.700,7.076) (-7.773,7.094) (-7.846,7.112) (-7.919,7.129) (-7.992,7.145) (-8.066,7.161) (-8.139,7.177) (-8.212,7.192) (-8.286,7.207) (-8.360,7.221) (-8.433,7.235) (-8.507,7.248) (-8.581,7.261) (-8.655,7.273) (-8.729,7.285) (-8.803,7.297) (-8.877,7.308) (-8.952,7.318) (-9.026,7.328) (-9.100,7.337) (-9.175,7.346) (-9.249,7.355) (-9.324,7.363) (-9.399,7.371) (-9.473,7.378) (-9.548,7.384) (-9.623,7.390) (-9.697,7.396) (-9.772,7.401) (-9.847,7.406) (-9.922,7.410) (-9.997,7.413) (-10.000,7.414)}
\LC{(-10.000,-8.150) (-9.943,-8.140) (-9.868,-8.130) (-9.794,-8.119) (-9.720,-8.107) (-9.646,-8.095) (-9.572,-8.083) (-9.498,-8.070) (-9.424,-8.057) (-9.350,-8.044) (-9.277,-8.030) (-9.203,-8.016) (-9.130,-8.001) (-9.056,-7.986) (-8.983,-7.970) (-8.910,-7.954) (-8.836,-7.938) (-8.763,-7.921) (-8.690,-7.903) (-8.617,-7.886) (-8.545,-7.868) (-8.472,-7.849) (-8.399,-7.830) (-8.327,-7.811) (-8.255,-7.791) (-8.182,-7.771) (-8.110,-7.750) (-8.038,-7.729) (-7.966,-7.708) (-7.895,-7.686) (-7.823,-7.664) (-7.751,-7.641) (-7.680,-7.618) (-7.609,-7.595) (-7.538,-7.571) (-7.467,-7.547) (-7.396,-7.522) (-7.325,-7.497) (-7.255,-7.472) (-7.184,-7.446) (-7.114,-7.420) (-7.044,-7.393) (-6.974,-7.366) (-6.904,-7.339) (-6.834,-7.311) (-6.765,-7.283) (-6.695,-7.255) (-6.626,-7.226) (-6.557,-7.196) (-6.488,-7.167) (-6.420,-7.137) (-6.351,-7.106) (-6.283,-7.075) (-6.214,-7.044) (-6.146,-7.012) (-6.079,-6.980) (-6.011,-6.948) (-5.944,-6.915) (-5.876,-6.882) (-5.809,-6.849) (-5.742,-6.815) (-5.676,-6.780) (-5.609,-6.746) (-5.543,-6.711) (-5.477,-6.675) (-5.411,-6.640) (-5.345,-6.603) (-5.279,-6.567) (-5.214,-6.530) (-5.149,-6.493) (-5.084,-6.455) (-5.020,-6.417) (-4.955,-6.379) (-4.891,-6.340) (-4.827,-6.301) (-4.763,-6.261) (-4.700,-6.222) (-4.636,-6.181) (-4.573,-6.141) (-4.510,-6.100) (-4.448,-6.059) (-4.385,-6.017) (-4.323,-5.975) (-4.261,-5.933) (-4.200,-5.890) (-4.138,-5.847) (-4.077,-5.803) (-4.016,-5.760) (-3.956,-5.715) (-3.895,-5.671) (-3.835,-5.626) (-3.775,-5.581) (-3.716,-5.535) (-3.656,-5.489) (-3.597,-5.443) (-3.539,-5.397) (-3.480,-5.350) (-3.422,-5.302) (-3.364,-5.255) (-3.306,-5.207) (-3.249,-5.158) (-3.192,-5.110) (-3.135,-5.061) (-3.079,-5.011) (-3.023,-4.961) (-2.967,-4.911) (-2.911,-4.861) (-2.856,-4.810) (-2.801,-4.759) (-2.747,-4.708) (-2.692,-4.656) (-2.638,-4.604) (-2.585,-4.551) (-2.531,-4.499) (-2.479,-4.445) (-2.426,-4.392) (-2.374,-4.338) (-2.322,-4.284) (-2.270,-4.230) (-2.219,-4.175) (-2.168,-4.120) (-2.118,-4.064) (-2.068,-4.008) (-2.018,-3.952) (-1.968,-3.896) (-1.919,-3.839) (-1.871,-3.782) (-1.823,-3.724) (-1.775,-3.667) (-1.727,-3.609) (-1.680,-3.550) (-1.634,-3.491) (-1.587,-3.432) (-1.542,-3.373) (-1.496,-3.313) (-1.451,-3.253) (-1.407,-3.193) (-1.363,-3.132) (-1.319,-3.071) (-1.276,-3.010) (-1.233,-2.948) (-1.191,-2.886) (-1.149,-2.824)}
\LC{(3.767,2.687) (3.842,2.687) (3.917,2.685) (3.992,2.683) (4.066,2.678) (4.141,2.673) (4.216,2.665) (4.290,2.657) (4.365,2.647) (4.439,2.635) (4.513,2.622) (4.586,2.608) (4.660,2.592) (4.732,2.574) (4.805,2.555) (4.877,2.535) (4.949,2.513) (5.020,2.489) (5.091,2.464) (5.161,2.438) (5.231,2.410) (5.299,2.380) (5.368,2.349) (5.435,2.316) (5.502,2.282) (5.568,2.246) (5.633,2.209) (5.697,2.170) (5.760,2.130) (5.823,2.088) (5.884,2.045) (5.944,2.000) (6.003,1.954) (6.061,1.906) (6.117,1.856) (6.172,1.806) (6.226,1.753) (6.279,1.700) (6.330,1.645) (6.379,1.588) (6.427,1.531) (6.473,1.472) (6.518,1.411) (6.560,1.350) (6.601,1.287) (6.641,1.223) (6.678,1.158) (6.713,1.092) (6.747,1.025) (6.778,0.956) (6.807,0.887) (6.834,0.817) (6.859,0.747) (6.882,0.675) (6.902,0.603) (6.920,0.530) (6.936,0.457) (6.950,0.383) (6.961,0.309) (6.970,0.235) (6.976,0.160) (6.980,0.085) (6.981,0.010) (6.980,-0.065) (6.977,-0.140) (6.971,-0.215) (6.963,-0.289) (6.953,-0.364) (6.940,-0.437) (6.925,-0.511) (6.907,-0.584) (6.887,-0.656) (6.865,-0.728) (6.841,-0.799) (6.814,-0.869) (6.786,-0.938) (6.755,-1.007) (6.722,-1.074) (6.687,-1.141) (6.651,-1.206) (6.612,-1.270) (6.571,-1.333) (6.529,-1.395) (6.485,-1.456) (6.439,-1.515) (6.392,-1.573) (6.343,-1.630) (6.292,-1.685) (6.240,-1.739) (6.187,-1.792) (6.132,-1.843) (6.076,-1.893) (6.018,-1.941) (5.960,-1.988) (5.900,-2.033) (5.839,-2.077) (5.777,-2.119) (5.714,-2.160) (5.650,-2.199) (5.585,-2.237) (5.519,-2.273) (5.453,-2.308) (5.386,-2.341) (5.318,-2.372) (5.249,-2.402) (5.179,-2.431) (5.109,-2.458) (5.039,-2.483) (4.968,-2.507) (4.896,-2.529) (4.824,-2.550) (4.752,-2.569) (4.679,-2.587) (4.606,-2.604) (4.532,-2.618) (4.458,-2.632) (4.384,-2.644) (4.310,-2.654) (4.236,-2.663) (4.161,-2.671) (4.086,-2.677) (4.011,-2.682) (3.936,-2.685) (3.861,-2.687) (3.786,-2.687) (3.761,-2.687)}
\LC{(10.000,-7.569) (9.949,-7.564) (9.875,-7.558) (9.800,-7.553) (9.725,-7.546) (9.650,-7.539) (9.576,-7.532) (9.501,-7.524) (9.427,-7.516) (9.352,-7.508) (9.278,-7.498) (9.203,-7.489) (9.129,-7.479) (9.055,-7.468) (8.980,-7.457) (8.906,-7.445) (8.832,-7.434) (8.758,-7.421) (8.684,-7.408) (8.611,-7.395) (8.537,-7.381) (8.463,-7.367) (8.390,-7.352) (8.316,-7.337) (8.243,-7.321) (8.170,-7.305) (8.097,-7.288) (8.024,-7.271) (7.951,-7.254) (7.878,-7.236) (7.805,-7.218) (7.733,-7.199) (7.660,-7.179) (7.588,-7.160) (7.515,-7.140) (7.443,-7.119) (7.371,-7.098) (7.299,-7.076) (7.228,-7.054) (7.156,-7.032) (7.085,-7.009) (7.013,-6.986) (6.942,-6.962) (6.871,-6.938) (6.800,-6.914) (6.730,-6.889) (6.659,-6.863) (6.589,-6.837) (6.518,-6.811) (6.448,-6.784) (6.378,-6.757) (6.309,-6.730) (6.239,-6.702) (6.170,-6.673) (6.100,-6.645) (6.031,-6.615) (5.963,-6.586) (5.894,-6.555) (5.825,-6.525) (5.757,-6.494) (5.689,-6.463) (5.621,-6.431) (5.553,-6.399) (5.486,-6.366) (5.418,-6.333) (5.351,-6.300) (5.284,-6.266) (5.218,-6.232) (5.151,-6.197) (5.085,-6.162) (5.019,-6.126) (4.953,-6.091) (4.887,-6.054) (4.822,-6.018) (4.757,-5.980) (4.692,-5.943) (4.627,-5.905) (4.562,-5.867) (4.498,-5.828) (4.434,-5.789) (4.370,-5.749) (4.307,-5.710) (4.244,-5.669) (4.181,-5.629) (4.139,-5.601)}
\LC{(3.620,3.580) (3.691,3.603) (3.763,3.624) (3.835,3.645) (3.908,3.665) (3.980,3.684) (4.053,3.703) (4.126,3.720) (4.199,3.736) (4.272,3.752) (4.346,3.767) (4.420,3.781) (4.494,3.793) (4.568,3.805) (4.642,3.816) (4.716,3.827) (4.790,3.836) (4.865,3.844) (4.940,3.851) (5.014,3.858) (5.089,3.863) (5.164,3.868) (5.239,3.871) (5.314,3.874) (5.389,3.876) (5.464,3.876) (5.539,3.876) (5.614,3.875) (5.689,3.873) (5.764,3.870) (5.839,3.865) (5.914,3.860) (5.988,3.854) (6.063,3.847) (6.137,3.839) (6.212,3.830) (6.286,3.820) (6.360,3.808) (6.434,3.796) (6.508,3.783) (6.582,3.769) (6.655,3.754) (6.728,3.737) (6.801,3.720) (6.874,3.702) (6.947,3.682) (7.019,3.662) (7.091,3.640) (7.162,3.618) (7.233,3.594) (7.304,3.570) (7.375,3.544) (7.445,3.517) (7.514,3.489) (7.584,3.460) (7.652,3.431) (7.721,3.400) (7.788,3.367) (7.856,3.334) (7.922,3.300) (7.989,3.265) (8.054,3.229) (8.119,3.191) (8.184,3.153) (8.247,3.113) (8.311,3.073) (8.373,3.031) (8.435,2.988) (8.496,2.945) (8.556,2.900) (8.615,2.854) (8.674,2.807) (8.731,2.759) (8.788,2.710) (8.844,2.660) (8.899,2.610) (8.953,2.558) (9.007,2.505) (9.059,2.451) (9.110,2.396) (9.160,2.340) (9.209,2.284) (9.257,2.226) (9.304,2.167) (9.350,2.108) (9.394,2.048) (9.438,1.986) (9.480,1.924) (9.520,1.861) (9.560,1.798) (9.598,1.733) (9.635,1.668) (9.671,1.602) (9.705,1.535) (9.738,1.468) (9.769,1.400) (9.799,1.331) (9.828,1.261) (9.855,1.191) (9.880,1.121) (9.904,1.050) (9.927,0.978) (9.947,0.906) (9.967,0.834) (9.984,0.761) (10.000,0.688)}
\LC{(2.538,5.010) (2.587,5.067) (2.637,5.123) (2.687,5.179) (2.737,5.234) (2.788,5.290) (2.839,5.345) (2.890,5.400) (2.941,5.454) (2.993,5.508) (3.045,5.563) (3.097,5.616) (3.150,5.670) (3.203,5.723) (3.256,5.776) (3.309,5.829) (3.363,5.882) (3.416,5.934) (3.470,5.986) (3.525,6.037) (3.579,6.089) (3.634,6.140) (3.689,6.191) (3.745,6.242) (3.800,6.292) (3.856,6.342) (3.912,6.392) (3.968,6.442) (4.025,6.491) (4.081,6.540) (4.138,6.589) (4.196,6.637) (4.253,6.685) (4.311,6.733) (4.369,6.781) (4.427,6.829) (4.485,6.876) (4.544,6.923) (4.602,6.969) (4.661,7.016) (4.720,7.062) (4.780,7.107) (4.840,7.153) (4.899,7.198) (4.959,7.243) (5.020,7.288) (5.080,7.332) (5.141,7.376) (5.202,7.420) (5.263,7.464) (5.324,7.507) (5.385,7.550) (5.447,7.593) (5.509,7.635) (5.571,7.677) (5.633,7.719) (5.695,7.761) (5.758,7.802) (5.821,7.843) (5.884,7.884) (5.947,7.925) (6.010,7.965) (6.074,8.005) (6.137,8.044) (6.201,8.084) (6.265,8.123) (6.329,8.162) (6.394,8.200) (6.458,8.238) (6.523,8.276) (6.588,8.314) (6.653,8.351) (6.718,8.388) (6.783,8.425) (6.849,8.462) (6.915,8.498) (6.980,8.534) (7.046,8.569) (7.113,8.605) (7.179,8.640) (7.245,8.674) (7.312,8.709) (7.379,8.743) (7.446,8.777) (7.513,8.810) (7.580,8.844) (7.647,8.877) (7.715,8.909) (7.782,8.942) (7.850,8.974) (7.918,9.005) (7.986,9.037) (8.054,9.068) (8.123,9.099) (8.191,9.130) (8.260,9.160) (8.329,9.190) (8.398,9.220) (8.467,9.249) (8.536,9.278) (8.605,9.307) (8.674,9.335) (8.744,9.363) (8.813,9.391) (8.883,9.419) (8.953,9.446) (9.023,9.473) (9.093,9.500) (9.163,9.526) (9.234,9.552) (9.304,9.578) (9.375,9.603) (9.446,9.628) (9.516,9.653) (9.587,9.677) (9.658,9.702) (9.729,9.725) (9.801,9.749) (9.872,9.772) (9.943,9.795) (10.000,9.818)}
\LC{(1.312,1.428) (1.381,1.456) (1.451,1.482) (1.523,1.506) (1.594,1.528) (1.667,1.548) (1.739,1.566) (1.813,1.582) (1.886,1.595) (1.961,1.607) (2.035,1.616) (2.110,1.623) (2.184,1.628) (2.259,1.631) (2.334,1.631) (2.409,1.629) (2.484,1.625) (2.559,1.618) (2.633,1.609) (2.708,1.597) (2.781,1.583) (2.854,1.567) (2.927,1.548) (2.999,1.527) (3.070,1.503) (3.140,1.477) (3.209,1.448) (3.278,1.417) (3.345,1.383) (3.410,1.346) (3.474,1.307) (3.537,1.266) (3.597,1.222) (3.656,1.176) (3.713,1.127) (3.768,1.075) (3.820,1.022) (3.870,0.966) (3.918,0.908) (3.963,0.848) (4.004,0.785) (4.043,0.721) (4.079,0.655) (4.111,0.588) (4.140,0.518) (4.165,0.448) (4.187,0.376) (4.205,0.303) (4.219,0.230) (4.229,0.155) (4.235,0.080) (4.238,0.006) (4.236,-0.069) (4.230,-0.144) (4.221,-0.219) (4.207,-0.292) (4.190,-0.365) (4.169,-0.437) (4.144,-0.508) (4.116,-0.577) (4.084,-0.645) (4.049,-0.712) (4.010,-0.776) (3.969,-0.839) (3.925,-0.899) (3.878,-0.958) (3.828,-1.014) (3.776,-1.068) (3.721,-1.119) (3.665,-1.169) (3.606,-1.215) (3.546,-1.260) (3.483,-1.301) (3.420,-1.341) (3.354,-1.378) (3.288,-1.412) (3.220,-1.444) (3.151,-1.473) (3.080,-1.499) (3.009,-1.524) (2.938,-1.545) (2.865,-1.564) (2.792,-1.581) (2.718,-1.595) (2.644,-1.607) (2.570,-1.617) (2.495,-1.624) (2.420,-1.628) (2.345,-1.631) (2.270,-1.631) (2.195,-1.629) (2.121,-1.624) (2.046,-1.617) (1.971,-1.608) (1.897,-1.597) (1.824,-1.584) (1.750,-1.568) (1.677,-1.551) (1.605,-1.531) (1.533,-1.510) (1.462,-1.486) (1.391,-1.460) (1.322,-1.432) (1.299,-1.423)}
\LC{(2.736,-10.000) (2.707,-9.935) (2.678,-9.866) (2.648,-9.797) (2.619,-9.728) (2.591,-9.659) (2.562,-9.589) (2.533,-9.520) (2.504,-9.451) (2.476,-9.382) (2.447,-9.312) (2.419,-9.243) (2.391,-9.173) (2.363,-9.104) (2.335,-9.034) (2.307,-8.964) (2.279,-8.895) (2.252,-8.825) (2.224,-8.755) (2.197,-8.685) (2.169,-8.616) (2.142,-8.546) (2.115,-8.476) (2.088,-8.406) (2.061,-8.336) (2.034,-8.266) (2.008,-8.196) (1.981,-8.125) (1.955,-8.055) (1.929,-7.985) (1.902,-7.915) (1.876,-7.844) (1.851,-7.774) (1.825,-7.704) (1.799,-7.633) (1.773,-7.563) (1.748,-7.492) (1.723,-7.421) (1.697,-7.351) (1.672,-7.280) (1.647,-7.209) (1.623,-7.139) (1.598,-7.068) (1.573,-6.997) (1.549,-6.926) (1.524,-6.855) (1.500,-6.784) (1.476,-6.713) (1.452,-6.642) (1.428,-6.571) (1.405,-6.500) (1.381,-6.429) (1.358,-6.357) (1.334,-6.286) (1.311,-6.215) (1.288,-6.143) (1.265,-6.072) (1.243,-6.000) (1.220,-5.929) (1.197,-5.857) (1.175,-5.786) (1.153,-5.714) (1.131,-5.642) (1.109,-5.571) (1.087,-5.499) (1.066,-5.427) (1.044,-5.355) (1.023,-5.283) (1.001,-5.211) (0.980,-5.139) (0.959,-5.067) (0.939,-4.995) (0.918,-4.923) (0.898,-4.851) (0.877,-4.779) (0.857,-4.707) (0.837,-4.634) (0.817,-4.562) (0.798,-4.490) (0.778,-4.417) (0.759,-4.345) (0.740,-4.272)}
\LC{(-3.539,5.679) (-3.596,5.728) (-3.653,5.776) (-3.711,5.824) (-3.768,5.872) (-3.826,5.920) (-3.885,5.967) (-3.943,6.014) (-4.002,6.061) (-4.061,6.107) (-4.120,6.153) (-4.179,6.199) (-4.239,6.244) (-4.299,6.289) (-4.359,6.334) (-4.420,6.378) (-4.480,6.422) (-4.541,6.466) (-4.602,6.509) (-4.664,6.553) (-4.725,6.595) (-4.787,6.638) (-4.849,6.680) (-4.911,6.722) (-4.974,6.764) (-5.037,6.805) (-5.099,6.846) (-5.162,6.886) (-5.226,6.926) (-5.289,6.966) (-5.353,7.006) (-5.417,7.045) (-5.481,7.084) (-5.545,7.123) (-5.610,7.161) (-5.674,7.199) (-5.739,7.237) (-5.804,7.274) (-5.870,7.311) (-5.935,7.347) (-6.001,7.384) (-6.067,7.420) (-6.133,7.455) (-6.199,7.490) (-6.265,7.525) (-6.332,7.560) (-6.398,7.594) (-6.465,7.628) (-6.532,7.662) (-6.600,7.695) (-6.667,7.728) (-6.735,7.761) (-6.802,7.793) (-6.870,7.825) (-6.938,7.856) (-7.006,7.887) (-7.075,7.918) (-7.143,7.949) (-7.212,7.979) (-7.281,8.009) (-7.350,8.038) (-7.419,8.067) (-7.488,8.096) (-7.558,8.125) (-7.627,8.153) (-7.697,8.180) (-7.767,8.208) (-7.837,8.235) (-7.907,8.261) (-7.977,8.288) (-8.047,8.314) (-8.118,8.339) (-8.188,8.365) (-8.259,8.389) (-8.330,8.414) (-8.401,8.438) (-8.472,8.462) (-8.543,8.485) (-8.615,8.509) (-8.686,8.531) (-8.758,8.554) (-8.829,8.576) (-8.901,8.597) (-8.973,8.619) (-9.045,8.640) (-9.117,8.660) (-9.190,8.680) (-9.262,8.700) (-9.334,8.720) (-9.407,8.739) (-9.479,8.757) (-9.552,8.776) (-9.625,8.794) (-9.698,8.811) (-9.771,8.829) (-9.844,8.845) (-9.917,8.862) (-9.990,8.878) (-10.000,8.883)}
\LC{(8.623,2.249) (8.673,2.193) (8.722,2.136) (8.769,2.078) (8.815,2.019) (8.860,1.959) (8.904,1.898) (8.947,1.836) (8.988,1.773) (9.028,1.710) (9.066,1.646) (9.103,1.580) (9.139,1.514) (9.173,1.448) (9.206,1.380) (9.237,1.312) (9.267,1.243) (9.295,1.174) (9.321,1.103) (9.346,1.033) (9.370,0.961) (9.391,0.890) (9.411,0.817) (9.430,0.745) (9.446,0.671) (9.461,0.598) (9.474,0.524) (9.486,0.450) (9.496,0.376) (9.504,0.301) (9.510,0.226) (9.514,0.151) (9.517,0.076) (9.518,0.001) (9.517,-0.073) (9.514,-0.148) (9.510,-0.223) (9.504,-0.298) (9.496,-0.373) (9.486,-0.447) (9.475,-0.521) (9.462,-0.595) (9.447,-0.668) (9.430,-0.742) (9.412,-0.814) (9.392,-0.887) (9.370,-0.958) (9.347,-1.030) (9.322,-1.101) (9.296,-1.171) (9.268,-1.240) (9.238,-1.309) (9.207,-1.377) (9.174,-1.445) (9.140,-1.512) (9.105,-1.578) (9.068,-1.643) (9.029,-1.707) (8.990,-1.771) (8.948,-1.834) (8.906,-1.896) (8.862,-1.956) (8.817,-2.016) (8.771,-2.075) (8.724,-2.134) (8.675,-2.191) (8.625,-2.247) (8.574,-2.302) (8.522,-2.356) (8.469,-2.409) (8.415,-2.461) (8.360,-2.512) (8.304,-2.562) (8.248,-2.611) (8.190,-2.659) (8.131,-2.705) (8.071,-2.751) (8.011,-2.795) (7.950,-2.839) (7.888,-2.881) (7.825,-2.922) (7.762,-2.962) (7.698,-3.001) (7.633,-3.039) (7.567,-3.076) (7.501,-3.111) (7.435,-3.146) (7.368,-3.179) (7.300,-3.211) (7.232,-3.242) (7.163,-3.272) (7.094,-3.301) (7.024,-3.328) (6.954,-3.355) (6.883,-3.380) (6.812,-3.405) (6.741,-3.428) (6.669,-3.450) (6.597,-3.471) (6.525,-3.491) (6.452,-3.509) (6.379,-3.527) (6.306,-3.544) (6.233,-3.559) (6.159,-3.574) (6.085,-3.587) (6.011,-3.599) (5.937,-3.610) (5.863,-3.620) (5.788,-3.629) (5.714,-3.637) (5.639,-3.644) (5.564,-3.650) (5.490,-3.655) (5.415,-3.658) (5.340,-3.661) (5.265,-3.663) (5.190,-3.663) (5.115,-3.663) (5.040,-3.662) (4.965,-3.659) (4.890,-3.656) (4.815,-3.651)}
\LC{(9.552,-10.000) (9.482,-9.976) (9.412,-9.949) (9.342,-9.921) (9.273,-9.894) (9.203,-9.866) (9.134,-9.837) (9.064,-9.809) (8.995,-9.780) (8.926,-9.751) (8.857,-9.722) (8.788,-9.692) (8.719,-9.662) (8.651,-9.632) (8.582,-9.601) (8.514,-9.571) (8.445,-9.539) (8.377,-9.508) (8.309,-9.476) (8.241,-9.444) (8.174,-9.412) (8.106,-9.380) (8.039,-9.347) (7.971,-9.314) (7.904,-9.281) (7.837,-9.247) (7.770,-9.213) (7.703,-9.179) (7.637,-9.144) (7.570,-9.110) (7.504,-9.075) (7.438,-9.039) (7.372,-9.004) (7.306,-8.968) (7.240,-8.932) (7.175,-8.895) (7.109,-8.859) (7.044,-8.822) (6.979,-8.785) (6.914,-8.747) (6.849,-8.709) (6.784,-8.671) (6.720,-8.633) (6.656,-8.594) (6.592,-8.556) (6.528,-8.516) (6.464,-8.477) (6.400,-8.437) (6.337,-8.398) (6.273,-8.357) (6.210,-8.317) (6.147,-8.276) (6.084,-8.235) (6.022,-8.194) (5.959,-8.152) (5.897,-8.110) (5.835,-8.068) (5.773,-8.026) (5.711,-7.983) (5.650,-7.941) (5.588,-7.897) (5.527,-7.854) (5.466,-7.810) (5.405,-7.766) (5.345,-7.722) (5.284,-7.678) (5.224,-7.633) (5.164,-7.588) (5.104,-7.543) (5.045,-7.497) (4.985,-7.452) (4.926,-7.406) (4.867,-7.359) (4.808,-7.313) (4.750,-7.266) (4.691,-7.219) (4.633,-7.172) (4.575,-7.124) (4.517,-7.076) (4.460,-7.028) (4.402,-6.980) (4.345,-6.931) (4.288,-6.882) (4.232,-6.833) (4.175,-6.784) (4.119,-6.734) (4.063,-6.684) (4.007,-6.634) (3.952,-6.584) (3.896,-6.533) (3.841,-6.482) (3.786,-6.431) (3.732,-6.380) (3.677,-6.328) (3.623,-6.276) (3.569,-6.224) (3.515,-6.172) (3.462,-6.119) (3.409,-6.066) (3.356,-6.013) (3.303,-5.960) (3.251,-5.906) (3.199,-5.852) (3.147,-5.798) (3.095,-5.744) (3.044,-5.689) (2.993,-5.634) (2.942,-5.579) (2.891,-5.524) (2.841,-5.468) (2.791,-5.412) (2.741,-5.356) (2.691,-5.300) (2.642,-5.243) (2.593,-5.187) (2.545,-5.130) (2.496,-5.072) (2.448,-5.015) (2.400,-4.957) (2.353,-4.899) (2.306,-4.841) (2.259,-4.782) (2.212,-4.723) (2.166,-4.664) (2.120,-4.605) (2.074,-4.546) (2.029,-4.486) (1.984,-4.426) (1.939,-4.366) (1.894,-4.305) (1.850,-4.245) (1.807,-4.184) (1.763,-4.123) (1.720,-4.061) (1.677,-4.000) (1.635,-3.938) (1.593,-3.876) (1.551,-3.813) (1.510,-3.751) (1.482,-3.709)}
\LC{(-10.000,-5.600) (-9.942,-5.616) (-9.869,-5.631) (-9.795,-5.646) (-9.721,-5.659) (-9.648,-5.673) (-9.574,-5.685) (-9.500,-5.697) (-9.425,-5.708) (-9.351,-5.719) (-9.277,-5.729) (-9.202,-5.738) (-9.128,-5.747) (-9.053,-5.754) (-8.979,-5.762) (-8.904,-5.768) (-8.829,-5.774) (-8.754,-5.779) (-8.679,-5.784) (-8.605,-5.788) (-8.530,-5.791) (-8.455,-5.793) (-8.380,-5.795) (-8.305,-5.796) (-8.230,-5.797) (-8.155,-5.797) (-8.080,-5.796) (-8.005,-5.795) (-7.930,-5.793) (-7.855,-5.790) (-7.780,-5.787) (-7.705,-5.783) (-7.630,-5.779) (-7.555,-5.773) (-7.480,-5.768) (-7.406,-5.761) (-7.331,-5.754) (-7.256,-5.747) (-7.182,-5.738) (-7.107,-5.729) (-7.033,-5.720) (-6.959,-5.710) (-6.885,-5.699) (-6.810,-5.688) (-6.736,-5.676) (-6.662,-5.663) (-6.589,-5.650) (-6.515,-5.636) (-6.441,-5.622) (-6.368,-5.607) (-6.294,-5.591) (-6.221,-5.575) (-6.148,-5.558) (-6.075,-5.541) (-6.002,-5.523) (-5.930,-5.504) (-5.857,-5.485) (-5.785,-5.465) (-5.713,-5.445) (-5.641,-5.424) (-5.569,-5.403) (-5.497,-5.380) (-5.426,-5.358) (-5.354,-5.335) (-5.283,-5.311) (-5.212,-5.287) (-5.141,-5.262) (-5.071,-5.236) (-5.001,-5.210) (-4.930,-5.184) (-4.860,-5.156) (-4.791,-5.129) (-4.721,-5.100) (-4.652,-5.072) (-4.583,-5.042) (-4.514,-5.012) (-4.446,-4.982) (-4.377,-4.951) (-4.309,-4.919) (-4.242,-4.887) (-4.174,-4.854) (-4.107,-4.821) (-4.040,-4.787) (-3.973,-4.753) (-3.907,-4.718) (-3.841,-4.683) (-3.775,-4.647) (-3.709,-4.611) (-3.644,-4.574) (-3.579,-4.536) (-3.514,-4.498) (-3.450,-4.460) (-3.386,-4.421) (-3.322,-4.381) (-3.259,-4.341) (-3.196,-4.301) (-3.133,-4.259) (-3.071,-4.218) (-3.009,-4.176) (-2.947,-4.133) (-2.886,-4.090) (-2.825,-4.046) (-2.764,-4.002) (-2.704,-3.957) (-2.644,-3.912) (-2.584,-3.867) (-2.525,-3.820) (-2.467,-3.774)}
\LC{(10.000,-4.876) (9.939,-4.901) (9.868,-4.925) (9.797,-4.949) (9.725,-4.972) (9.653,-4.993) (9.581,-5.015) (9.509,-5.035) (9.437,-5.055) (9.364,-5.073) (9.292,-5.092) (9.219,-5.109) (9.145,-5.125) (9.072,-5.141) (8.999,-5.156) (8.925,-5.171) (8.851,-5.184) (8.777,-5.197) (8.703,-5.209) (8.629,-5.220) (8.555,-5.231) (8.481,-5.241) (8.406,-5.250) (8.332,-5.258) (8.257,-5.266) (8.182,-5.273) (8.108,-5.279) (8.033,-5.285) (7.958,-5.289) (7.883,-5.293) (7.808,-5.297) (7.733,-5.299) (7.658,-5.301) (7.583,-5.302) (7.508,-5.303) (7.433,-5.302) (7.358,-5.301) (7.283,-5.300) (7.208,-5.297) (7.133,-5.294) (7.058,-5.290) (6.984,-5.286) (6.909,-5.281) (6.834,-5.275) (6.759,-5.268) (6.685,-5.261) (6.610,-5.253) (6.535,-5.245) (6.461,-5.236) (6.387,-5.226) (6.312,-5.215) (6.238,-5.204) (6.164,-5.192) (6.090,-5.179) (6.017,-5.166) (5.943,-5.152) (5.869,-5.137) (5.796,-5.122) (5.723,-5.106) (5.649,-5.090) (5.576,-5.073) (5.504,-5.055) (5.431,-5.036) (5.358,-5.017) (5.286,-4.998) (5.214,-4.977) (5.142,-4.956) (5.070,-4.934) (4.998,-4.912) (4.927,-4.889) (4.856,-4.866) (4.785,-4.842) (4.714,-4.817) (4.643,-4.791) (4.573,-4.765) (4.503,-4.739) (4.433,-4.711) (4.363,-4.684) (4.294,-4.655) (4.248,-4.636)}
\LC{(-0.445,2.886) (-0.462,2.959) (-0.479,3.032) (-0.497,3.105) (-0.515,3.178) (-0.533,3.251) (-0.551,3.323) (-0.570,3.396) (-0.588,3.469) (-0.607,3.541) (-0.626,3.614) (-0.645,3.686) (-0.665,3.759) (-0.684,3.831) (-0.704,3.904) (-0.724,3.976) (-0.744,4.048) (-0.765,4.120) (-0.785,4.192) (-0.806,4.265) (-0.827,4.337) (-0.848,4.409) (-0.869,4.481) (-0.891,4.552) (-0.912,4.624) (-0.934,4.696) (-0.956,4.768) (-0.978,4.839) (-1.000,4.911) (-1.023,4.983) (-1.045,5.054) (-1.068,5.126) (-1.091,5.197) (-1.114,5.268) (-1.137,5.340) (-1.160,5.411) (-1.184,5.482) (-1.208,5.553) (-1.231,5.624) (-1.255,5.695) (-1.280,5.766) (-1.304,5.837) (-1.328,5.908) (-1.353,5.979) (-1.378,6.050) (-1.403,6.121) (-1.428,6.191) (-1.453,6.262) (-1.478,6.333) (-1.504,6.403) (-1.529,6.474) (-1.555,6.544) (-1.581,6.615) (-1.607,6.685) (-1.633,6.755) (-1.659,6.825) (-1.686,6.896) (-1.712,6.966) (-1.739,7.036) (-1.766,7.106) (-1.793,7.176) (-1.820,7.246) (-1.847,7.316) (-1.875,7.385) (-1.902,7.455) (-1.930,7.525) (-1.958,7.595) (-1.986,7.664) (-2.014,7.734) (-2.042,7.803) (-2.070,7.873) (-2.099,7.942) (-2.127,8.011) (-2.156,8.081) (-2.185,8.150) (-2.214,8.219) (-2.243,8.288) (-2.272,8.357) (-2.301,8.426) (-2.331,8.495) (-2.360,8.564) (-2.390,8.633) (-2.420,8.702) (-2.450,8.771) (-2.480,8.840) (-2.510,8.908) (-2.540,8.977) (-2.570,9.045) (-2.601,9.114) (-2.632,9.182) (-2.662,9.251) (-2.693,9.319) (-2.724,9.387) (-2.755,9.456) (-2.786,9.524) (-2.818,9.592) (-2.849,9.660) (-2.881,9.728) (-2.912,9.796) (-2.944,9.864) (-2.976,9.932) (-3.008,10.000) (-3.019,10.000)}
\LC{(-10.000,-9.569) (-9.950,-9.548) (-9.878,-9.527) (-9.807,-9.505) (-9.735,-9.484) (-9.663,-9.461) (-9.592,-9.439) (-9.520,-9.416) (-9.449,-9.393) (-9.378,-9.369) (-9.307,-9.346) (-9.236,-9.321) (-9.165,-9.297) (-9.094,-9.272) (-9.023,-9.247) (-8.953,-9.222) (-8.882,-9.196) (-8.812,-9.170) (-8.742,-9.143) (-8.672,-9.117) (-8.602,-9.090) (-8.532,-9.062) (-8.462,-9.035) (-8.393,-9.007) (-8.323,-8.978) (-8.254,-8.950) (-8.185,-8.921) (-8.116,-8.892) (-8.047,-8.862) (-7.978,-8.832) (-7.909,-8.802) (-7.841,-8.771) (-7.772,-8.741) (-7.704,-8.709) (-7.636,-8.678) (-7.568,-8.646) (-7.500,-8.614) (-7.432,-8.582) (-7.365,-8.549) (-7.298,-8.516) (-7.230,-8.483) (-7.163,-8.449) (-7.096,-8.415) (-7.030,-8.381) (-6.963,-8.347) (-6.897,-8.312) (-6.830,-8.277) (-6.764,-8.241) (-6.698,-8.206) (-6.633,-8.170) (-6.567,-8.133) (-6.501,-8.097) (-6.436,-8.060) (-6.371,-8.023) (-6.306,-7.985) (-6.241,-7.947) (-6.177,-7.909) (-6.112,-7.871) (-6.048,-7.832) (-5.984,-7.793) (-5.920,-7.754) (-5.857,-7.714) (-5.793,-7.674) (-5.730,-7.634) (-5.667,-7.593) (-5.604,-7.553) (-5.541,-7.512) (-5.478,-7.470) (-5.416,-7.429) (-5.354,-7.387) (-5.292,-7.344) (-5.230,-7.302) (-5.168,-7.259) (-5.107,-7.216) (-5.046,-7.173) (-4.985,-7.129) (-4.924,-7.085) (-4.864,-7.041) (-4.803,-6.996) (-4.743,-6.951) (-4.683,-6.906) (-4.624,-6.861) (-4.564,-6.815) (-4.505,-6.769) (-4.446,-6.723) (-4.387,-6.676) (-4.328,-6.629) (-4.270,-6.582) (-4.212,-6.535) (-4.154,-6.487) (-4.096,-6.439) (-4.039,-6.391) (-3.982,-6.342) (-3.925,-6.294) (-3.868,-6.245) (-3.812,-6.195) (-3.755,-6.146) (-3.699,-6.096) (-3.644,-6.045)}
\LC{(-2.919,-10.000) (-2.888,-9.932) (-2.857,-9.864) (-2.826,-9.795) (-2.795,-9.727) (-2.764,-9.658) (-2.734,-9.590) (-2.703,-9.521) (-2.673,-9.453) (-2.642,-9.384) (-2.612,-9.316) (-2.582,-9.247) (-2.552,-9.178) (-2.523,-9.109) (-2.493,-9.040) (-2.463,-8.971) (-2.434,-8.902) (-2.405,-8.833) (-2.375,-8.764) (-2.346,-8.695) (-2.317,-8.626) (-2.288,-8.557) (-2.260,-8.488) (-2.231,-8.418) (-2.203,-8.349) (-2.174,-8.279) (-2.146,-8.210) (-2.118,-8.140) (-2.090,-8.071) (-2.062,-8.001) (-2.034,-7.932) (-2.006,-7.862) (-1.979,-7.792) (-1.951,-7.722) (-1.924,-7.652) (-1.897,-7.583) (-1.870,-7.513) (-1.843,-7.443) (-1.816,-7.373) (-1.790,-7.302) (-1.763,-7.232) (-1.737,-7.162) (-1.710,-7.092) (-1.684,-7.021) (-1.658,-6.951) (-1.632,-6.881) (-1.607,-6.810) (-1.581,-6.740) (-1.556,-6.669) (-1.530,-6.599) (-1.505,-6.528) (-1.480,-6.457) (-1.455,-6.387) (-1.431,-6.316) (-1.406,-6.245) (-1.381,-6.174) (-1.357,-6.103) (-1.333,-6.032) (-1.309,-5.961) (-1.285,-5.890) (-1.261,-5.819) (-1.238,-5.748) (-1.214,-5.676) (-1.191,-5.605) (-1.168,-5.534) (-1.145,-5.462) (-1.122,-5.391) (-1.099,-5.320) (-1.076,-5.248) (-1.054,-5.176) (-1.032,-5.105) (-1.010,-5.033) (-0.988,-4.961) (-0.966,-4.890) (-0.944,-4.818) (-0.923,-4.746) (-0.902,-4.674) (-0.880,-4.602) (-0.859,-4.530) (-0.839,-4.458) (-0.818,-4.386)}
\LC{(2.151,2.078) (2.223,2.099) (2.296,2.118) (2.368,2.135) (2.442,2.152) (2.515,2.166) (2.589,2.179) (2.663,2.191) (2.738,2.200) (2.812,2.209) (2.887,2.215) (2.962,2.220) (3.037,2.223) (3.112,2.225) (3.187,2.225) (3.262,2.223) (3.336,2.220) (3.411,2.215) (3.486,2.208) (3.560,2.199) (3.635,2.189) (3.709,2.176) (3.782,2.163) (3.856,2.147) (3.929,2.129) (4.001,2.110) (4.073,2.089) (4.144,2.066) (4.215,2.041) (4.285,2.014) (4.355,1.986) (4.423,1.955) (4.491,1.923) (4.558,1.889) (4.624,1.853) (4.688,1.815) (4.752,1.776) (4.815,1.734) (4.876,1.691) (4.936,1.646) (4.994,1.599) (5.051,1.550) (5.107,1.500) (5.160,1.447) (5.213,1.393) (5.263,1.338) (5.311,1.281) (5.358,1.222) (5.402,1.161) (5.445,1.099) (5.485,1.036) (5.523,0.971) (5.558,0.905) (5.592,0.838) (5.622,0.770) (5.651,0.700) (5.676,0.630) (5.699,0.558) (5.719,0.486) (5.736,0.413) (5.751,0.340) (5.763,0.266) (5.772,0.191) (5.778,0.116) (5.781,0.041) (5.781,-0.034) (5.778,-0.109) (5.772,-0.183) (5.764,-0.258) (5.752,-0.332) (5.738,-0.406) (5.721,-0.479) (5.701,-0.551) (5.679,-0.622) (5.653,-0.693) (5.625,-0.763) (5.595,-0.831) (5.562,-0.899) (5.527,-0.965) (5.489,-1.030) (5.449,-1.093) (5.407,-1.155) (5.363,-1.216) (5.316,-1.275) (5.268,-1.332) (5.218,-1.388) (5.166,-1.442) (5.112,-1.494) (5.057,-1.545) (5.000,-1.594) (4.942,-1.641) (4.882,-1.686) (4.821,-1.730) (4.759,-1.771) (4.695,-1.811) (4.630,-1.849) (4.565,-1.885) (4.498,-1.920) (4.430,-1.952) (4.362,-1.983) (4.293,-2.011) (4.222,-2.038) (4.152,-2.063) (4.080,-2.086) (4.009,-2.108) (3.936,-2.127) (3.863,-2.145) (3.790,-2.161) (3.716,-2.175) (3.642,-2.187) (3.568,-2.198) (3.494,-2.207) (3.419,-2.214) (3.344,-2.219) (3.269,-2.223) (3.194,-2.225) (3.119,-2.225) (3.044,-2.224) (2.969,-2.220) (2.895,-2.216) (2.820,-2.209) (2.745,-2.201) (2.671,-2.192) (2.597,-2.180) (2.523,-2.167) (2.449,-2.153) (2.376,-2.137) (2.303,-2.120) (2.230,-2.101) (2.158,-2.080)}
\LC{(-0.002,2.978) (-0.002,3.053) (-0.002,3.128) (-0.002,3.203) (-0.003,3.278) (-0.003,3.353) (-0.003,3.428) (-0.003,3.503) (-0.003,3.578) (-0.003,3.653) (-0.003,3.728) (-0.003,3.803) (-0.003,3.878) (-0.003,3.953) (-0.003,4.028) (-0.004,4.103) (-0.004,4.178) (-0.004,4.253) (-0.004,4.328) (-0.004,4.403) (-0.004,4.478) (-0.004,4.553) (-0.004,4.628) (-0.004,4.703) (-0.004,4.778) (-0.005,4.853) (-0.005,4.928) (-0.005,5.003) (-0.005,5.078) (-0.005,5.153) (-0.005,5.228) (-0.005,5.303) (-0.005,5.378) (-0.005,5.453) (-0.006,5.528) (-0.006,5.603) (-0.006,5.678) (-0.006,5.753) (-0.006,5.828) (-0.006,5.903) (-0.006,5.978) (-0.006,6.053) (-0.006,6.128) (-0.007,6.203) (-0.007,6.278) (-0.007,6.353) (-0.007,6.428) (-0.007,6.503) (-0.007,6.578) (-0.007,6.653) (-0.007,6.728) (-0.008,6.803) (-0.008,6.878) (-0.008,6.953) (-0.008,7.028) (-0.008,7.103) (-0.008,7.178) (-0.008,7.253) (-0.008,7.328) (-0.009,7.403) (-0.009,7.478) (-0.009,7.553) (-0.009,7.628) (-0.009,7.703) (-0.009,7.778) (-0.009,7.853) (-0.010,7.928) (-0.010,8.003) (-0.010,8.078) (-0.010,8.153) (-0.010,8.228) (-0.010,8.303) (-0.010,8.378) (-0.011,8.453) (-0.011,8.528) (-0.011,8.603) (-0.011,8.678) (-0.011,8.753) (-0.011,8.828) (-0.011,8.903) (-0.012,8.978) (-0.012,9.053) (-0.012,9.128) (-0.012,9.203) (-0.012,9.278) (-0.012,9.353) (-0.012,9.428) (-0.013,9.503) (-0.013,9.578) (-0.013,9.653) (-0.013,9.728) (-0.013,9.803) (-0.013,9.878) (-0.013,9.953) (-0.014,10.000)}
\LC{(4.675,4.277) (4.748,4.294) (4.821,4.311) (4.894,4.327) (4.967,4.342) (5.041,4.357) (5.115,4.370) (5.189,4.383) (5.263,4.395) (5.337,4.406) (5.411,4.417) (5.485,4.427) (5.560,4.435) (5.634,4.444) (5.709,4.451) (5.784,4.457) (5.859,4.463) (5.933,4.468) (6.008,4.472) (6.083,4.475) (6.158,4.477) (6.233,4.479) (6.308,4.480) (6.383,4.480) (6.458,4.479) (6.533,4.477) (6.608,4.474) (6.683,4.471) (6.758,4.466) (6.833,4.461) (6.907,4.455) (6.982,4.448) (7.057,4.440) (7.131,4.432) (7.206,4.422) (7.280,4.412) (7.354,4.401) (7.428,4.388) (7.502,4.375) (7.576,4.361) (7.649,4.347) (7.722,4.331) (7.796,4.314) (7.869,4.297) (7.941,4.278) (8.014,4.259) (8.086,4.239) (8.158,4.218) (8.230,4.196) (8.301,4.173) (8.372,4.149) (8.443,4.124) (8.513,4.099) (8.584,4.072) (8.653,4.045) (8.723,4.016) (8.792,3.987) (8.860,3.957) (8.929,3.926) (8.996,3.893) (9.064,3.860) (9.131,3.827) (9.197,3.792) (9.263,3.756) (9.328,3.719) (9.393,3.682) (9.458,3.643) (9.521,3.604) (9.585,3.563) (9.647,3.522) (9.709,3.480) (9.771,3.437) (9.831,3.393) (9.891,3.348) (9.951,3.302) (10.000,3.255)}
\LC{(-1.771,2.911) (-1.827,2.961) (-1.884,3.010) (-1.941,3.058) (-1.999,3.106) (-2.057,3.154) (-2.116,3.200) (-2.175,3.246) (-2.234,3.292) (-2.295,3.337) (-2.355,3.381) (-2.416,3.425) (-2.477,3.468) (-2.539,3.511) (-2.601,3.552) (-2.664,3.594) (-2.727,3.634) (-2.790,3.674) (-2.854,3.714) (-2.918,3.753) (-2.983,3.791) (-3.048,3.828) (-3.113,3.865) (-3.179,3.902) (-3.245,3.937) (-3.311,3.972) (-3.378,4.007) (-3.445,4.040) (-3.512,4.074) (-3.580,4.106) (-3.648,4.138) (-3.716,4.169) (-3.784,4.199) (-3.853,4.229) (-3.922,4.258) (-3.992,4.287) (-4.061,4.315) (-4.131,4.342) (-4.201,4.369) (-4.272,4.394) (-4.342,4.420) (-4.413,4.444) (-4.484,4.468) (-4.556,4.491) (-4.627,4.514) (-4.699,4.535) (-4.771,4.556) (-4.843,4.577) (-4.915,4.596) (-4.988,4.615) (-5.061,4.634) (-5.134,4.651) (-5.207,4.668) (-5.280,4.684) (-5.353,4.700) (-5.427,4.715) (-5.501,4.729) (-5.574,4.742) (-5.648,4.754) (-5.722,4.766) (-5.797,4.777) (-5.871,4.788) (-5.945,4.798) (-6.020,4.806) (-6.094,4.815) (-6.169,4.822) (-6.244,4.829) (-6.318,4.835) (-6.393,4.840) (-6.468,4.844) (-6.543,4.848) (-6.618,4.851) (-6.693,4.853) (-6.768,4.855) (-6.843,4.855) (-6.918,4.855) (-6.993,4.854) (-7.068,4.852) (-7.143,4.850) (-7.218,4.847) (-7.293,4.843) (-7.367,4.838) (-7.442,4.832) (-7.517,4.826) (-7.592,4.819) (-7.666,4.811) (-7.741,4.802) (-7.815,4.792) (-7.889,4.782) (-7.963,4.770) (-8.037,4.758) (-8.111,4.746) (-8.185,4.732) (-8.259,4.717) (-8.332,4.702) (-8.405,4.686) (-8.478,4.669) (-8.551,4.651) (-8.624,4.633) (-8.696,4.613) (-8.769,4.593) (-8.840,4.572) (-8.912,4.550) (-8.984,4.527) (-9.055,4.503) (-9.126,4.479) (-9.196,4.454) (-9.267,4.427) (-9.337,4.400) (-9.406,4.372) (-9.475,4.344) (-9.544,4.314) (-9.613,4.284) (-9.681,4.253) (-9.749,4.220) (-9.816,4.187) (-9.883,4.154) (-9.950,4.119) (-10.000,4.084)}
\LC{(2.995,4.844) (3.051,4.893) (3.108,4.942) (3.166,4.990) (3.223,5.038) (3.281,5.086) (3.339,5.133) (3.398,5.180) (3.456,5.227) (3.516,5.273) (3.575,5.319) (3.635,5.364) (3.694,5.410) (3.755,5.454) (3.815,5.499) (3.876,5.543) (3.937,5.586) (3.998,5.630) (4.060,5.673) (4.121,5.715) (4.183,5.757) (4.246,5.799) (4.308,5.840) (4.371,5.881) (4.434,5.922) (4.497,5.962) (4.561,6.002) (4.624,6.042) (4.688,6.081) (4.753,6.120) (4.817,6.158) (4.882,6.196) (4.947,6.234) (5.012,6.271) (5.077,6.308) (5.143,6.344) (5.208,6.380) (5.274,6.416) (5.341,6.451) (5.407,6.486) (5.474,6.521) (5.540,6.555) (5.607,6.588) (5.675,6.622) (5.742,6.655) (5.809,6.687) (5.877,6.719) (5.945,6.751) (6.013,6.782) (6.082,6.813) (6.150,6.844) (6.219,6.874) (6.288,6.904) (6.357,6.933) (6.426,6.962) (6.495,6.991) (6.565,7.019) (6.634,7.046) (6.704,7.074) (6.774,7.101) (6.845,7.127) (6.915,7.153) (6.985,7.179) (7.056,7.204) (7.127,7.229) (7.198,7.253) (7.269,7.277) (7.340,7.301) (7.411,7.324) (7.483,7.347) (7.554,7.369) (7.626,7.391) (7.698,7.413) (7.770,7.434) (7.842,7.454) (7.914,7.474) (7.987,7.494) (8.059,7.514) (8.132,7.533) (8.204,7.551) (8.277,7.569) (8.350,7.587) (8.423,7.604) (8.496,7.621) (8.569,7.637) (8.642,7.653) (8.716,7.668) (8.789,7.683) (8.863,7.698) (8.937,7.712) (9.010,7.726) (9.084,7.739) (9.158,7.752) (9.232,7.764) (9.306,7.776) (9.380,7.788) (9.454,7.799) (9.529,7.810) (9.603,7.820) (9.677,7.829) (9.752,7.839) (9.826,7.847) (9.901,7.856) (9.975,7.864) (10.000,7.866)}
\LC{(-1.118,1.217) (-1.188,1.245) (-1.258,1.270) (-1.330,1.293) (-1.402,1.314) (-1.474,1.332) (-1.548,1.348) (-1.622,1.362) (-1.696,1.372) (-1.770,1.380) (-1.845,1.386) (-1.920,1.389) (-1.995,1.389) (-2.070,1.387) (-2.145,1.382) (-2.219,1.374) (-2.294,1.363) (-2.367,1.349) (-2.440,1.333) (-2.513,1.313) (-2.585,1.291) (-2.655,1.266) (-2.725,1.238) (-2.793,1.207) (-2.860,1.172) (-2.925,1.135) (-2.988,1.095) (-3.050,1.052) (-3.109,1.006) (-3.166,0.958) (-3.220,0.906) (-3.272,0.852) (-3.321,0.795) (-3.366,0.735) (-3.409,0.673) (-3.447,0.609) (-3.482,0.543) (-3.514,0.474) (-3.541,0.405) (-3.563,0.333) (-3.582,0.260) (-3.596,0.187) (-3.605,0.112) (-3.609,0.037) (-3.609,-0.038) (-3.605,-0.112) (-3.596,-0.187) (-3.582,-0.261) (-3.563,-0.333) (-3.540,-0.405) (-3.513,-0.475) (-3.482,-0.543) (-3.447,-0.609) (-3.409,-0.673) (-3.366,-0.735) (-3.321,-0.795) (-3.272,-0.852) (-3.220,-0.906) (-3.166,-0.958) (-3.109,-1.006) (-3.049,-1.052) (-2.988,-1.095) (-2.925,-1.135) (-2.859,-1.172) (-2.793,-1.207) (-2.724,-1.238) (-2.655,-1.266) (-2.584,-1.291) (-2.513,-1.313) (-2.440,-1.333) (-2.367,-1.349) (-2.293,-1.363) (-2.219,-1.374) (-2.145,-1.382) (-2.070,-1.387) (-1.995,-1.389) (-1.920,-1.389) (-1.845,-1.386) (-1.770,-1.380) (-1.696,-1.372) (-1.621,-1.361) (-1.548,-1.348) (-1.474,-1.332) (-1.401,-1.314) (-1.329,-1.293) (-1.258,-1.270) (-1.188,-1.245) (-1.118,-1.217)}
\LC{(-10.000,-2.969) (-9.946,-3.019) (-9.890,-3.069) (-9.833,-3.118) (-9.775,-3.165) (-9.716,-3.212) (-9.657,-3.258) (-9.597,-3.303) (-9.537,-3.347) (-9.475,-3.391) (-9.413,-3.433) (-9.351,-3.474) (-9.287,-3.515) (-9.224,-3.554) (-9.159,-3.593) (-9.094,-3.630) (-9.029,-3.667) (-8.963,-3.702) (-8.897,-3.737) (-8.830,-3.771) (-8.762,-3.804) (-8.694,-3.836) (-8.626,-3.867) (-8.557,-3.897) (-8.488,-3.926) (-8.419,-3.954) (-8.349,-3.981) (-8.278,-4.007) (-8.208,-4.033) (-8.137,-4.057) (-8.066,-4.080) (-7.994,-4.103) (-7.922,-4.124) (-7.850,-4.145) (-7.778,-4.165) (-7.705,-4.184) (-7.632,-4.201) (-7.559,-4.218) (-7.486,-4.234) (-7.413,-4.249) (-7.339,-4.264) (-7.265,-4.277) (-7.191,-4.289) (-7.117,-4.301) (-7.043,-4.311) (-6.968,-4.321) (-6.894,-4.330) (-6.819,-4.338) (-6.745,-4.345) (-6.670,-4.351) (-6.595,-4.356) (-6.520,-4.361) (-6.445,-4.364) (-6.370,-4.367) (-6.295,-4.368) (-6.220,-4.369) (-6.145,-4.369) (-6.070,-4.369) (-5.995,-4.367) (-5.920,-4.364) (-5.846,-4.361) (-5.771,-4.357) (-5.696,-4.352) (-5.621,-4.346) (-5.546,-4.339) (-5.472,-4.332) (-5.397,-4.323) (-5.323,-4.314) (-5.248,-4.304) (-5.174,-4.293) (-5.100,-4.282) (-5.051,-4.273)}
\LC{(-2.119,1.838) (-2.193,1.851) (-2.267,1.863) (-2.341,1.872) (-2.416,1.880) (-2.491,1.886) (-2.566,1.890) (-2.641,1.892) (-2.716,1.892) (-2.791,1.890) (-2.866,1.885) (-2.940,1.879) (-3.015,1.871) (-3.089,1.861) (-3.163,1.849) (-3.237,1.834) (-3.310,1.818) (-3.383,1.799) (-3.455,1.778) (-3.526,1.756) (-3.597,1.730) (-3.667,1.703) (-3.736,1.674) (-3.804,1.642) (-3.871,1.608) (-3.936,1.573) (-4.001,1.534) (-4.064,1.494) (-4.126,1.452) (-4.186,1.407) (-4.245,1.360) (-4.302,1.311) (-4.357,1.261) (-4.410,1.208) (-4.461,1.153) (-4.510,1.096) (-4.557,1.037) (-4.601,0.977) (-4.643,0.914) (-4.682,0.850) (-4.718,0.785) (-4.752,0.718) (-4.782,0.649) (-4.810,0.580) (-4.835,0.509) (-4.856,0.437) (-4.874,0.364) (-4.889,0.291) (-4.901,0.216) (-4.909,0.142) (-4.914,0.067) (-4.915,-0.008) (-4.913,-0.083) (-4.908,-0.158) (-4.899,-0.232) (-4.887,-0.306) (-4.871,-0.379) (-4.852,-0.452) (-4.830,-0.524) (-4.805,-0.594) (-4.776,-0.664) (-4.745,-0.732) (-4.711,-0.799) (-4.674,-0.864) (-4.634,-0.927) (-4.592,-0.989) (-4.547,-1.050) (-4.500,-1.108) (-4.451,-1.164) (-4.399,-1.219) (-4.346,-1.271) (-4.290,-1.322) (-4.233,-1.370) (-4.174,-1.417) (-4.113,-1.461) (-4.051,-1.503) (-3.987,-1.543) (-3.923,-1.580) (-3.857,-1.616) (-3.789,-1.649) (-3.721,-1.680) (-3.652,-1.709) (-3.582,-1.736) (-3.511,-1.761) (-3.440,-1.783) (-3.367,-1.803) (-3.295,-1.821) (-3.221,-1.838) (-3.148,-1.851) (-3.074,-1.863) (-2.999,-1.873) (-2.925,-1.881) (-2.850,-1.887) (-2.775,-1.890) (-2.700,-1.892) (-2.625,-1.891) (-2.550,-1.889) (-2.475,-1.885) (-2.400,-1.879) (-2.326,-1.871) (-2.251,-1.861) (-2.177,-1.849) (-2.128,-1.840)}
\LC{(-4.280,5.092) (-4.347,5.126) (-4.414,5.159) (-4.482,5.192) (-4.549,5.225) (-4.617,5.257) (-4.685,5.288) (-4.754,5.319) (-4.822,5.349) (-4.891,5.379) (-4.960,5.409) (-5.029,5.437) (-5.099,5.466) (-5.168,5.494) (-5.238,5.521) (-5.308,5.548) (-5.379,5.574) (-5.449,5.600) (-5.520,5.625) (-5.590,5.650) (-5.661,5.674) (-5.733,5.697) (-5.804,5.720) (-5.876,5.743) (-5.947,5.765) (-6.019,5.786) (-6.091,5.807) (-6.163,5.827) (-6.236,5.847) (-6.308,5.866) (-6.381,5.885) (-6.454,5.903) (-6.527,5.921) (-6.600,5.938) (-6.673,5.954) (-6.746,5.970) (-6.819,5.986) (-6.893,6.001) (-6.967,6.015) (-7.040,6.028) (-7.114,6.042) (-7.188,6.054) (-7.262,6.066) (-7.336,6.078) (-7.410,6.088) (-7.485,6.099) (-7.559,6.108) (-7.634,6.117) (-7.708,6.126) (-7.783,6.134) (-7.857,6.141) (-7.932,6.148) (-8.007,6.154) (-8.082,6.160) (-8.156,6.165) (-8.231,6.169) (-8.306,6.173) (-8.381,6.176) (-8.456,6.179) (-8.531,6.181) (-8.606,6.182) (-8.681,6.183) (-8.756,6.183) (-8.831,6.183) (-8.906,6.182) (-8.981,6.180) (-9.056,6.178) (-9.131,6.175) (-9.206,6.172) (-9.281,6.168) (-9.356,6.163) (-9.430,6.158) (-9.505,6.152) (-9.580,6.145) (-9.654,6.138) (-9.729,6.130) (-9.804,6.122) (-9.878,6.113) (-9.952,6.103) (-10.000,6.096)}
\LC{(1.461,2.722) (1.513,2.776) (1.565,2.830) (1.618,2.883) (1.671,2.936) (1.725,2.988) (1.779,3.040) (1.834,3.091) (1.890,3.142) (1.945,3.192) (2.002,3.242) (2.058,3.291) (2.115,3.339) (2.173,3.387) (2.231,3.435) (2.290,3.482) (2.349,3.528) (2.408,3.574) (2.468,3.619) (2.528,3.664) (2.588,3.708) (2.649,3.752) (2.711,3.795) (2.773,3.837) (2.835,3.879) (2.897,3.921) (2.960,3.962) (3.024,4.002) (3.087,4.042) (3.151,4.081) (3.215,4.119) (3.280,4.157) (3.345,4.195) (3.410,4.232) (3.476,4.268) (3.542,4.304) (3.608,4.339) (3.675,4.373) (3.742,4.407) (3.809,4.441) (3.876,4.473) (3.944,4.506) (4.012,4.537) (4.080,4.568) (4.149,4.599) (4.218,4.628) (4.287,4.658) (4.356,4.686) (4.426,4.714) (4.496,4.741) (4.566,4.768) (4.636,4.794) (4.706,4.820) (4.777,4.845) (4.848,4.869) (4.919,4.893) (4.991,4.916) (5.062,4.938) (5.134,4.960) (5.206,4.981) (5.278,5.002) (5.350,5.022) (5.423,5.041) (5.496,5.060) (5.568,5.078) (5.641,5.095) (5.714,5.112) (5.788,5.128) (5.861,5.143) (5.935,5.158) (6.008,5.172) (6.082,5.185) (6.156,5.198) (6.230,5.210) (6.304,5.222) (6.378,5.233) (6.453,5.243) (6.527,5.252) (6.602,5.261) (6.676,5.269) (6.751,5.276) (6.825,5.283) (6.900,5.289) (6.975,5.294) (7.050,5.299) (7.125,5.303) (7.200,5.306) (7.275,5.309) (7.350,5.311) (7.425,5.312) (7.500,5.312) (7.575,5.312) (7.650,5.311) (7.725,5.310) (7.800,5.307) (7.874,5.304) (7.949,5.300) (8.024,5.296) (8.099,5.291) (8.174,5.285) (8.249,5.278) (8.323,5.271) (8.398,5.262) (8.472,5.254) (8.547,5.244) (8.621,5.234) (8.695,5.223) (8.769,5.211) (8.843,5.198) (8.917,5.185) (8.990,5.171) (9.064,5.156) (9.137,5.140) (9.211,5.124) (9.284,5.107) (9.356,5.089) (9.429,5.071) (9.502,5.051) (9.574,5.031) (9.646,5.010) (9.718,4.989) (9.789,4.966) (9.861,4.943) (9.932,4.919) (10.000,4.894)}
\LC{(1.369,4.863) (1.399,4.932) (1.429,5.000) (1.460,5.069) (1.491,5.137) (1.522,5.205) (1.553,5.273) (1.585,5.341) (1.617,5.409) (1.649,5.477) (1.681,5.545) (1.713,5.613) (1.746,5.680) (1.779,5.748) (1.812,5.815) (1.845,5.882) (1.878,5.949) (1.912,6.017) (1.945,6.084) (1.979,6.150) (2.013,6.217) (2.048,6.284) (2.082,6.350) (2.117,6.417) (2.152,6.483) (2.187,6.550) (2.222,6.616) (2.258,6.682) (2.293,6.748) (2.329,6.814) (2.365,6.880) (2.401,6.945) (2.437,7.011) (2.474,7.077) (2.510,7.142) (2.547,7.207) (2.584,7.273) (2.621,7.338) (2.659,7.403) (2.696,7.468) (2.734,7.533) (2.772,7.597) (2.810,7.662) (2.848,7.727) (2.886,7.791) (2.925,7.855) (2.964,7.920) (3.002,7.984) (3.041,8.048) (3.081,8.112) (3.120,8.176) (3.159,8.239) (3.199,8.303) (3.239,8.367) (3.279,8.430) (3.319,8.494) (3.359,8.557) (3.399,8.620) (3.440,8.683) (3.481,8.746) (3.522,8.809) (3.563,8.872) (3.604,8.934) (3.645,8.997) (3.687,9.059) (3.728,9.122) (3.770,9.184) (3.812,9.246) (3.854,9.308) (3.896,9.370) (3.939,9.432) (3.981,9.494) (4.024,9.556) (4.067,9.617) (4.110,9.679) (4.153,9.740) (4.196,9.802) (4.239,9.863) (4.283,9.924) (4.326,9.985) (4.341,10.000)}
\LC{(-1.407,4.702) (-1.439,4.770) (-1.472,4.837) (-1.504,4.905) (-1.537,4.973) (-1.570,5.040) (-1.603,5.107) (-1.636,5.174) (-1.670,5.241) (-1.704,5.308) (-1.738,5.375) (-1.772,5.442) (-1.806,5.508) (-1.841,5.575) (-1.876,5.641) (-1.911,5.708) (-1.947,5.774) (-1.982,5.840) (-2.018,5.906) (-2.054,5.972) (-2.090,6.037) (-2.126,6.103) (-2.163,6.168) (-2.200,6.234) (-2.237,6.299) (-2.274,6.364) (-2.311,6.429) (-2.349,6.494) (-2.386,6.559) (-2.424,6.624) (-2.463,6.688) (-2.501,6.753) (-2.539,6.817) (-2.578,6.881) (-2.617,6.945) (-2.656,7.009) (-2.695,7.073) (-2.735,7.137) (-2.774,7.201) (-2.814,7.264) (-2.854,7.328) (-2.894,7.391) (-2.935,7.454) (-2.975,7.517) (-3.016,7.580) (-3.057,7.643) (-3.098,7.706) (-3.139,7.769) (-3.180,7.831) (-3.222,7.894) (-3.264,7.956) (-3.306,8.018) (-3.348,8.080) (-3.390,8.142) (-3.432,8.204) (-3.475,8.266) (-3.518,8.327) (-3.561,8.389) (-3.604,8.450) (-3.647,8.512) (-3.690,8.573) (-3.734,8.634) (-3.778,8.695) (-3.822,8.756) (-3.866,8.816) (-3.910,8.877) (-3.954,8.937) (-3.999,8.998) (-4.043,9.058) (-4.088,9.118) (-4.133,9.178) (-4.178,9.238) (-4.224,9.298) (-4.269,9.357) (-4.315,9.417) (-4.360,9.476) (-4.406,9.536) (-4.452,9.595) (-4.499,9.654) (-4.545,9.713) (-4.591,9.772) (-4.638,9.831) (-4.685,9.889) (-4.732,9.948) (-4.779,10.000)}
\LC{(0.449,3.419) (0.464,3.492) (0.479,3.566) (0.494,3.639) (0.509,3.713) (0.524,3.786) (0.540,3.860) (0.556,3.933) (0.571,4.006) (0.587,4.080) (0.603,4.153) (0.619,4.226) (0.636,4.299) (0.652,4.373) (0.669,4.446) (0.686,4.519) (0.703,4.592) (0.720,4.665) (0.737,4.738) (0.754,4.811) (0.772,4.884) (0.789,4.957) (0.807,5.030) (0.825,5.102) (0.843,5.175) (0.861,5.248) (0.879,5.321) (0.897,5.394) (0.916,5.466) (0.934,5.539) (0.953,5.612) (0.972,5.684) (0.991,5.757) (1.010,5.829) (1.029,5.902) (1.048,5.974) (1.068,6.047) (1.087,6.119) (1.107,6.191) (1.127,6.264) (1.147,6.336) (1.167,6.408) (1.187,6.481) (1.207,6.553) (1.227,6.625) (1.248,6.697) (1.268,6.769) (1.289,6.841) (1.310,6.913) (1.331,6.986) (1.352,7.057) (1.373,7.129) (1.394,7.201) (1.415,7.273) (1.437,7.345) (1.458,7.417) (1.480,7.489) (1.502,7.561) (1.524,7.632) (1.546,7.704) (1.568,7.776) (1.590,7.847) (1.612,7.919) (1.635,7.991) (1.657,8.062) (1.680,8.134) (1.702,8.205) (1.725,8.277) (1.748,8.348) (1.771,8.419) (1.794,8.491) (1.817,8.562) (1.841,8.633) (1.864,8.705) (1.887,8.776) (1.911,8.847) (1.935,8.918) (1.958,8.989) (1.982,9.060) (2.006,9.132) (2.030,9.203) (2.054,9.274) (2.079,9.345) (2.103,9.416) (2.127,9.486) (2.152,9.557) (2.176,9.628) (2.201,9.699) (2.226,9.770) (2.251,9.841) (2.276,9.911) (2.301,9.982) (2.309,10.000)}
\LC{(10.000,-9.013) (9.929,-8.996) (9.856,-8.978) (9.784,-8.960) (9.711,-8.942) (9.638,-8.924) (9.565,-8.905) (9.493,-8.886) (9.421,-8.866) (9.348,-8.846) (9.276,-8.826) (9.204,-8.805) (9.132,-8.784) (9.060,-8.763) (8.988,-8.741) (8.917,-8.719) (8.845,-8.696) (8.774,-8.674) (8.702,-8.650) (8.631,-8.627) (8.560,-8.603) (8.489,-8.579) (8.418,-8.554) (8.347,-8.529) (8.277,-8.504) (8.206,-8.478) (8.136,-8.453) (8.066,-8.426) (7.996,-8.400) (7.926,-8.373) (7.856,-8.345) (7.786,-8.318) (7.717,-8.289) (7.647,-8.261) (7.578,-8.232) (7.509,-8.203) (7.440,-8.174) (7.371,-8.144) (7.302,-8.114) (7.234,-8.084) (7.165,-8.053) (7.097,-8.022) (7.029,-7.990) (6.961,-7.959) (6.893,-7.927) (6.825,-7.894) (6.758,-7.861) (6.691,-7.828) (6.624,-7.795) (6.557,-7.761) (6.490,-7.727) (6.423,-7.692) (6.357,-7.658) (6.290,-7.623) (6.224,-7.587) (6.158,-7.551) (6.093,-7.515) (6.027,-7.479) (5.962,-7.442) (5.897,-7.405) (5.832,-7.368) (5.767,-7.330) (5.702,-7.292) (5.638,-7.254) (5.573,-7.215) (5.509,-7.176) (5.445,-7.137) (5.382,-7.097) (5.318,-7.057) (5.255,-7.017) (5.192,-6.976) (5.129,-6.935) (5.067,-6.894) (5.004,-6.852) (4.942,-6.810) (4.880,-6.768) (4.818,-6.726) (4.757,-6.683) (4.716,-6.654)}
\LC{(4.277,5.274) (4.342,5.310) (4.408,5.346) (4.475,5.381) (4.541,5.415) (4.608,5.449) (4.675,5.483) (4.742,5.516) (4.810,5.549) (4.878,5.581) (4.946,5.613) (5.014,5.644) (5.082,5.675) (5.151,5.705) (5.220,5.735) (5.289,5.764) (5.358,5.793) (5.427,5.822) (5.497,5.850) (5.567,5.877) (5.637,5.904) (5.707,5.930) (5.777,5.956) (5.848,5.982) (5.919,6.006) (5.989,6.031) (6.061,6.055) (6.132,6.078) (6.203,6.101) (6.275,6.123) (6.347,6.145) (6.418,6.167) (6.490,6.188) (6.563,6.208) (6.635,6.228) (6.707,6.247) (6.780,6.266) (6.853,6.284) (6.926,6.302) (6.999,6.319) (7.072,6.336) (7.145,6.352) (7.218,6.368) (7.292,6.383) (7.365,6.398) (7.439,6.412) (7.513,6.425) (7.587,6.439) (7.661,6.451) (7.735,6.463) (7.809,6.474) (7.883,6.485) (7.957,6.496) (8.032,6.506) (8.106,6.515) (8.180,6.524) (8.255,6.532) (8.330,6.539) (8.404,6.547) (8.479,6.553) (8.554,6.559) (8.629,6.565) (8.703,6.569) (8.778,6.574) (8.853,6.578) (8.928,6.581) (9.003,6.583) (9.078,6.586) (9.153,6.587) (9.228,6.588) (9.303,6.588) (9.378,6.588) (9.453,6.587) (9.528,6.586) (9.603,6.584) (9.678,6.582) (9.753,6.579) (9.828,6.575) (9.903,6.571) (9.977,6.566) (10.000,6.564)}
\LC{(0.444,6.194) (0.452,6.268) (0.460,6.343) (0.468,6.418) (0.477,6.492) (0.485,6.567) (0.493,6.641) (0.502,6.716) (0.510,6.790) (0.518,6.865) (0.527,6.939) (0.535,7.014) (0.544,7.088) (0.553,7.163) (0.561,7.237) (0.570,7.312) (0.579,7.386) (0.588,7.461) (0.596,7.535) (0.605,7.610) (0.614,7.684) (0.623,7.759) (0.632,7.833) (0.641,7.908) (0.650,7.982) (0.660,8.056) (0.669,8.131) (0.678,8.205) (0.687,8.280) (0.697,8.354) (0.706,8.429) (0.715,8.503) (0.725,8.577) (0.734,8.652) (0.744,8.726) (0.753,8.801) (0.763,8.875) (0.773,8.949) (0.782,9.024) (0.792,9.098) (0.802,9.172) (0.812,9.247) (0.822,9.321) (0.831,9.395) (0.841,9.470) (0.851,9.544) (0.861,9.618) (0.871,9.693) (0.881,9.767) (0.892,9.841) (0.902,9.916) (0.912,9.990) (0.915,10.000)}
\LC{(-0.458,5.643) (-0.467,5.717) (-0.476,5.792) (-0.485,5.866) (-0.495,5.941) (-0.504,6.015) (-0.513,6.090) (-0.523,6.164) (-0.532,6.238) (-0.542,6.313) (-0.552,6.387) (-0.561,6.461) (-0.571,6.536) (-0.581,6.610) (-0.591,6.685) (-0.601,6.759) (-0.611,6.833) (-0.621,6.908) (-0.631,6.982) (-0.641,7.056) (-0.651,7.130) (-0.661,7.205) (-0.672,7.279) (-0.682,7.353) (-0.692,7.428) (-0.703,7.502) (-0.713,7.576) (-0.724,7.650) (-0.734,7.725) (-0.745,7.799) (-0.756,7.873) (-0.767,7.947) (-0.777,8.022) (-0.788,8.096) (-0.799,8.170) (-0.810,8.244) (-0.821,8.318) (-0.832,8.392) (-0.843,8.467) (-0.854,8.541) (-0.866,8.615) (-0.877,8.689) (-0.888,8.763) (-0.900,8.837) (-0.911,8.912) (-0.922,8.986) (-0.934,9.060) (-0.945,9.134) (-0.957,9.208) (-0.969,9.282) (-0.980,9.356) (-0.992,9.430) (-1.004,9.504) (-1.016,9.578) (-1.028,9.652) (-1.039,9.726) (-1.051,9.801) (-1.063,9.875) (-1.076,9.949) (-1.088,10.000)}
\LC{(1.579,1.049) (1.654,1.043) (1.728,1.033) (1.802,1.020) (1.875,1.002) (1.947,0.981) (2.017,0.956) (2.086,0.927) (2.154,0.894) (2.219,0.857) (2.282,0.816) (2.342,0.771) (2.399,0.723) (2.453,0.671) (2.503,0.615) (2.549,0.556) (2.591,0.493) (2.628,0.428) (2.660,0.360) (2.686,0.290) (2.707,0.218) (2.722,0.144) (2.731,0.070) (2.733,-0.005) (2.730,-0.080) (2.720,-0.154) (2.704,-0.228) (2.683,-0.299) (2.656,-0.369) (2.623,-0.437) (2.586,-0.502) (2.544,-0.564) (2.497,-0.622) (2.446,-0.678) (2.392,-0.729) (2.334,-0.777) (2.274,-0.822) (2.210,-0.862) (2.145,-0.898) (2.077,-0.931) (2.008,-0.959) (1.937,-0.984) (1.865,-1.005) (1.792,-1.022) (1.718,-1.035) (1.644,-1.044) (1.594,-1.048)}
\LC{(1.474,-10.000) (1.457,-9.936) (1.441,-9.863) (1.425,-9.789) (1.409,-9.716) (1.393,-9.643) (1.377,-9.570) (1.361,-9.496) (1.345,-9.423) (1.329,-9.350) (1.313,-9.276) (1.298,-9.203) (1.282,-9.130) (1.266,-9.056) (1.251,-8.983) (1.235,-8.910) (1.220,-8.836) (1.205,-8.763) (1.189,-8.689) (1.174,-8.616) (1.159,-8.542) (1.144,-8.469) (1.129,-8.395) (1.114,-8.322) (1.099,-8.248) (1.084,-8.175) (1.070,-8.101) (1.055,-8.028) (1.040,-7.954) (1.026,-7.881) (1.012,-7.807) (0.997,-7.734) (0.983,-7.660) (0.969,-7.586) (0.954,-7.513) (0.940,-7.439) (0.926,-7.365) (0.912,-7.292) (0.898,-7.218) (0.884,-7.144) (0.871,-7.070) (0.857,-6.997) (0.843,-6.923) (0.830,-6.849) (0.816,-6.775) (0.803,-6.702) (0.790,-6.628) (0.776,-6.554) (0.763,-6.480) (0.750,-6.406) (0.737,-6.332) (0.724,-6.259) (0.711,-6.185) (0.699,-6.111) (0.686,-6.037) (0.673,-5.963) (0.661,-5.889) (0.648,-5.815) (0.636,-5.741) (0.623,-5.667) (0.611,-5.593) (0.599,-5.519) (0.587,-5.445) (0.575,-5.371) (0.571,-5.346)}
\LC{(-7.036,-10.000) (-6.975,-9.958) (-6.914,-9.914) (-6.853,-9.870) (-6.793,-9.826) (-6.732,-9.782) (-6.672,-9.737) (-6.612,-9.692) (-6.552,-9.647) (-6.492,-9.602) (-6.433,-9.556) (-6.373,-9.510) (-6.314,-9.464) (-6.255,-9.418) (-6.196,-9.372) (-6.137,-9.326) (-6.078,-9.279) (-6.020,-9.232) (-5.961,-9.185) (-5.903,-9.137) (-5.845,-9.090) (-5.787,-9.042) (-5.730,-8.994) (-5.672,-8.946) (-5.615,-8.898) (-5.558,-8.849) (-5.501,-8.800) (-5.444,-8.751) (-5.387,-8.702) (-5.331,-8.653) (-5.274,-8.603) (-5.218,-8.554) (-5.162,-8.504) (-5.107,-8.454) (-5.051,-8.403) (-4.996,-8.353) (-4.940,-8.302) (-4.885,-8.251) (-4.830,-8.200) (-4.776,-8.149) (-4.721,-8.097) (-4.667,-8.045) (-4.613,-7.993) (-4.559,-7.941) (-4.505,-7.889) (-4.452,-7.836) (-4.398,-7.784) (-4.345,-7.731) (-4.292,-7.678) (-4.239,-7.624) (-4.187,-7.571) (-4.134,-7.517) (-4.082,-7.463) (-4.030,-7.409) (-3.978,-7.355) (-3.927,-7.301) (-3.876,-7.246) (-3.824,-7.191) (-3.773,-7.136) (-3.723,-7.081) (-3.672,-7.025) (-3.622,-6.970) (-3.572,-6.914) (-3.522,-6.858) (-3.472,-6.802) (-3.423,-6.745) (-3.374,-6.689) (-3.325,-6.632) (-3.276,-6.575) (-3.227,-6.518) (-3.179,-6.461) (-3.131,-6.403) (-3.083,-6.345) (-3.035,-6.287) (-2.988,-6.229) (-2.941,-6.171) (-2.894,-6.113) (-2.847,-6.054) (-2.800,-5.995) (-2.754,-5.936) (-2.708,-5.877) (-2.662,-5.817) (-2.617,-5.758) (-2.572,-5.698) (-2.527,-5.638) (-2.482,-5.578) (-2.437,-5.517) (-2.393,-5.457) (-2.349,-5.396) (-2.320,-5.356)}
\LC{(10.000,-2.658) (9.952,-2.713) (9.900,-2.767) (9.847,-2.819) (9.793,-2.871) (9.738,-2.922) (9.682,-2.972) (9.625,-3.021) (9.568,-3.069) (9.509,-3.117) (9.450,-3.163) (9.391,-3.208) (9.330,-3.253) (9.269,-3.296) (9.207,-3.339) (9.145,-3.381) (9.082,-3.421) (9.018,-3.461) (8.954,-3.500) (8.889,-3.537) (8.824,-3.574) (8.758,-3.610) (8.691,-3.645) (8.625,-3.679) (8.557,-3.711) (8.489,-3.743) (8.421,-3.774) (8.352,-3.804) (8.283,-3.833) (8.214,-3.861) (8.144,-3.889) (8.073,-3.915) (8.003,-3.940) (7.932,-3.964) (7.860,-3.987) (7.789,-4.010) (7.717,-4.031) (7.645,-4.051) (7.572,-4.071) (7.500,-4.089) (7.427,-4.107) (7.354,-4.124) (7.280,-4.139) (7.207,-4.154) (7.133,-4.168) (7.059,-4.181) (6.985,-4.193) (6.911,-4.204) (6.837,-4.214) (6.762,-4.223) (6.688,-4.231) (6.613,-4.239) (6.538,-4.245) (6.463,-4.251) (6.389,-4.255) (6.314,-4.259) (6.239,-4.262) (6.164,-4.264) (6.089,-4.265) (6.014,-4.265) (5.939,-4.265) (5.864,-4.263) (5.789,-4.261) (5.714,-4.258) (5.639,-4.253) (5.564,-4.248) (5.489,-4.243) (5.415,-4.236) (5.340,-4.228) (5.266,-4.220) (5.191,-4.211) (5.117,-4.201) (5.043,-4.190) (4.969,-4.178) (4.895,-4.165) (4.821,-4.152) (4.747,-4.138) (4.674,-4.123) (4.600,-4.107) (4.527,-4.090) (4.454,-4.073) (4.382,-4.055) (4.309,-4.036) (4.237,-4.016) (4.165,-3.995) (4.093,-3.974) (4.021,-3.952) (3.950,-3.929) (3.879,-3.905) (3.808,-3.880) (3.737,-3.855) (3.667,-3.829) (3.597,-3.802) (3.527,-3.774) (3.458,-3.746) (3.388,-3.717) (3.320,-3.687) (3.251,-3.657) (3.183,-3.625) (3.115,-3.593) (3.048,-3.560) (2.981,-3.527) (2.914,-3.492) (2.848,-3.457) (2.782,-3.422) (2.716,-3.385) (2.651,-3.348) (2.586,-3.310) (2.522,-3.272)}
\LC{(-10.000,-4.265) (-9.942,-4.298) (-9.875,-4.330) (-9.806,-4.362) (-9.738,-4.392) (-9.669,-4.422) (-9.600,-4.451) (-9.530,-4.479) (-9.460,-4.506) (-9.390,-4.532) (-9.320,-4.558) (-9.249,-4.583) (-9.178,-4.606) (-9.106,-4.629) (-9.035,-4.652) (-8.963,-4.673) (-8.891,-4.694) (-8.818,-4.713) (-8.746,-4.732) (-8.673,-4.751) (-8.600,-4.768) (-8.527,-4.784) (-8.454,-4.800) (-8.380,-4.815) (-8.306,-4.829) (-8.233,-4.842) (-8.159,-4.855) (-8.085,-4.867) (-8.010,-4.878) (-7.936,-4.888) (-7.862,-4.897) (-7.787,-4.906) (-7.712,-4.913) (-7.638,-4.920) (-7.563,-4.926) (-7.488,-4.932) (-7.413,-4.936) (-7.339,-4.940) (-7.264,-4.943) (-7.189,-4.946) (-7.114,-4.947) (-7.039,-4.948) (-6.964,-4.948) (-6.889,-4.947) (-6.814,-4.946) (-6.739,-4.944) (-6.664,-4.941) (-6.589,-4.937) (-6.514,-4.933) (-6.439,-4.927) (-6.364,-4.921) (-6.290,-4.915) (-6.215,-4.907) (-6.140,-4.899) (-6.066,-4.890) (-5.992,-4.881) (-5.917,-4.870) (-5.843,-4.859) (-5.769,-4.848) (-5.695,-4.835) (-5.621,-4.822) (-5.548,-4.808) (-5.474,-4.794) (-5.401,-4.778)}
\LC{(-10.000,-7.079) (-9.926,-7.079) (-9.851,-7.078) (-9.776,-7.076) (-9.701,-7.074) (-9.626,-7.072) (-9.552,-7.069) (-9.477,-7.066) (-9.402,-7.062) (-9.327,-7.057) (-9.252,-7.052) (-9.177,-7.046) (-9.102,-7.040) (-9.028,-7.034) (-8.953,-7.027) (-8.878,-7.019) (-8.804,-7.011) (-8.729,-7.002) (-8.655,-6.993) (-8.581,-6.984) (-8.506,-6.973) (-8.432,-6.963) (-8.358,-6.952) (-8.284,-6.940) (-8.210,-6.928) (-8.136,-6.915) (-8.062,-6.902) (-7.988,-6.888) (-7.915,-6.874) (-7.841,-6.860) (-7.768,-6.845) (-7.694,-6.829) (-7.621,-6.813) (-7.548,-6.796) (-7.475,-6.779) (-7.402,-6.762) (-7.329,-6.744) (-7.256,-6.725) (-7.184,-6.706) (-7.111,-6.687) (-7.039,-6.667) (-6.967,-6.646) (-6.895,-6.625) (-6.823,-6.604) (-6.751,-6.582) (-6.680,-6.560) (-6.608,-6.537) (-6.537,-6.514) (-6.466,-6.490) (-6.395,-6.466) (-6.324,-6.441) (-6.253,-6.416) (-6.183,-6.391) (-6.112,-6.365) (-6.042,-6.338) (-5.972,-6.311) (-5.902,-6.284) (-5.833,-6.256) (-5.763,-6.227) (-5.694,-6.199) (-5.625,-6.169) (-5.556,-6.140) (-5.488,-6.110) (-5.419,-6.079) (-5.351,-6.048) (-5.283,-6.016) (-5.215,-5.985) (-5.147,-5.952) (-5.080,-5.919) (-5.013,-5.886)}
\LC{(6.534,-10.000) (6.476,-9.962) (6.418,-9.915) (6.360,-9.868) (6.302,-9.820) (6.244,-9.772) (6.186,-9.724) (6.129,-9.676) (6.072,-9.628) (6.014,-9.579) (5.957,-9.530) (5.901,-9.481) (5.844,-9.432) (5.787,-9.383) (5.731,-9.333) (5.675,-9.284) (5.619,-9.234) (5.563,-9.184) (5.507,-9.134) (5.452,-9.083) (5.396,-9.033) (5.341,-8.982) (5.286,-8.931) (5.231,-8.880) (5.177,-8.828) (5.122,-8.777) (5.068,-8.725) (5.013,-8.673) (4.959,-8.621) (4.906,-8.569) (4.852,-8.517) (4.799,-8.464) (4.745,-8.411) (4.692,-8.358) (4.639,-8.305) (4.586,-8.252) (4.534,-8.198) (4.482,-8.145) (4.429,-8.091) (4.377,-8.037) (4.325,-7.983) (4.274,-7.928) (4.222,-7.874) (4.171,-7.819) (4.120,-7.764) (4.069,-7.709) (4.019,-7.654) (3.968,-7.598) (3.918,-7.542) (3.868,-7.487) (3.818,-7.431) (3.768,-7.374) (3.719,-7.318) (3.669,-7.262) (3.620,-7.205) (3.571,-7.148) (3.523,-7.091) (3.474,-7.034) (3.426,-6.976) (3.378,-6.919) (3.330,-6.861) (3.282,-6.803) (3.235,-6.745) (3.188,-6.687) (3.141,-6.628) (3.094,-6.570) (3.047,-6.511) (3.001,-6.452) (2.955,-6.393) (2.909,-6.333) (2.863,-6.274) (2.818,-6.214) (2.773,-6.154) (2.728,-6.094) (2.683,-6.034) (2.639,-5.974) (2.594,-5.913) (2.550,-5.852) (2.521,-5.812)}
\LC{(-2.798,6.367) (-2.842,6.428) (-2.886,6.488) (-2.931,6.548) (-2.976,6.608) (-3.021,6.668) (-3.067,6.728) (-3.112,6.788) (-3.158,6.847) (-3.204,6.906) (-3.250,6.965) (-3.297,7.024) (-3.343,7.083) (-3.390,7.141) (-3.437,7.200) (-3.484,7.258) (-3.532,7.316) (-3.580,7.374) (-3.627,7.432) (-3.675,7.489) (-3.724,7.547) (-3.772,7.604) (-3.821,7.661) (-3.870,7.718) (-3.919,7.775) (-3.968,7.832) (-4.017,7.888) (-4.067,7.944) (-4.116,8.000) (-4.166,8.056) (-4.216,8.112) (-4.267,8.168) (-4.317,8.223) (-4.368,8.279) (-4.419,8.334) (-4.470,8.389) (-4.521,8.444) (-4.572,8.498) (-4.624,8.553) (-4.676,8.607) (-4.728,8.661) (-4.780,8.715) (-4.832,8.769) (-4.884,8.822) (-4.937,8.876) (-4.990,8.929) (-5.043,8.982) (-5.096,9.035) (-5.149,9.088) (-5.203,9.141) (-5.256,9.193) (-5.310,9.245) (-5.364,9.298) (-5.418,9.349) (-5.472,9.401) (-5.527,9.453) (-5.581,9.504) (-5.636,9.556) (-5.691,9.607) (-5.746,9.658) (-5.801,9.708) (-5.857,9.759) (-5.912,9.809) (-5.968,9.860) (-6.024,9.910) (-6.080,9.960) (-6.136,10.000)}
\LC{(-1.929,0.938) (-1.998,0.911) (-2.066,0.879) (-2.132,0.844) (-2.196,0.804) (-2.257,0.760) (-2.315,0.713) (-2.370,0.661) (-2.421,0.606) (-2.467,0.548) (-2.510,0.486) (-2.547,0.421) (-2.579,0.353) (-2.606,0.283) (-2.627,0.211) (-2.641,0.137) (-2.650,0.063) (-2.652,-0.012) (-2.648,-0.087) (-2.637,-0.161) (-2.621,-0.234) (-2.598,-0.306) (-2.570,-0.375) (-2.536,-0.442) (-2.497,-0.506) (-2.453,-0.567) (-2.405,-0.624) (-2.353,-0.678) (-2.297,-0.728) (-2.238,-0.775) (-2.176,-0.817) (-2.112,-0.855) (-2.045,-0.890) (-1.976,-0.920) (-1.930,-0.938)}
\LC{(-1.501,-10.000) (-1.484,-9.935) (-1.467,-9.862) (-1.451,-9.789) (-1.435,-9.716) (-1.418,-9.643) (-1.402,-9.569) (-1.386,-9.496) (-1.369,-9.423) (-1.353,-9.350) (-1.337,-9.276) (-1.321,-9.203) (-1.305,-9.130) (-1.290,-9.057) (-1.274,-8.983) (-1.258,-8.910) (-1.242,-8.837) (-1.227,-8.763) (-1.211,-8.690) (-1.196,-8.616) (-1.180,-8.543) (-1.165,-8.470) (-1.150,-8.396) (-1.135,-8.323) (-1.120,-8.249) (-1.104,-8.176) (-1.089,-8.102) (-1.075,-8.029) (-1.060,-7.955) (-1.045,-7.882) (-1.030,-7.808) (-1.016,-7.735) (-1.001,-7.661) (-0.986,-7.588) (-0.972,-7.514) (-0.958,-7.440) (-0.943,-7.367) (-0.929,-7.293) (-0.915,-7.219) (-0.901,-7.146) (-0.887,-7.072) (-0.873,-6.998) (-0.859,-6.925) (-0.845,-6.851) (-0.832,-6.777) (-0.818,-6.703) (-0.804,-6.630) (-0.791,-6.556) (-0.777,-6.482) (-0.764,-6.408) (-0.751,-6.334) (-0.738,-6.261) (-0.725,-6.187) (-0.712,-6.113) (-0.699,-6.039) (-0.686,-5.965) (-0.673,-5.891) (-0.660,-5.817) (-0.648,-5.743) (-0.635,-5.669) (-0.623,-5.596) (-0.610,-5.522) (-0.598,-5.448) (-0.586,-5.374) (-0.574,-5.300) (-0.562,-5.226) (-0.550,-5.152) (-0.538,-5.077) (-0.526,-5.003) (-0.522,-4.979)}
\LC{(-5.577,4.285) (-5.652,4.290) (-5.727,4.295) (-5.802,4.299) (-5.877,4.302) (-5.952,4.304) (-6.027,4.305) (-6.102,4.305) (-6.177,4.304) (-6.252,4.303) (-6.327,4.301) (-6.401,4.297) (-6.476,4.293) (-6.551,4.288) (-6.626,4.282) (-6.701,4.276) (-6.775,4.268) (-6.850,4.259) (-6.924,4.250) (-6.998,4.239) (-7.073,4.228) (-7.147,4.216) (-7.220,4.203) (-7.294,4.189) (-7.368,4.174) (-7.441,4.158) (-7.514,4.141) (-7.587,4.123) (-7.659,4.104) (-7.732,4.085) (-7.804,4.064) (-7.876,4.042) (-7.947,4.020) (-8.019,3.996) (-8.089,3.972) (-8.160,3.947) (-8.230,3.920) (-8.300,3.893) (-8.370,3.865) (-8.439,3.836) (-8.507,3.806) (-8.576,3.774) (-8.643,3.742) (-8.711,3.709) (-8.778,3.675) (-8.844,3.641) (-8.910,3.605) (-8.975,3.568) (-9.040,3.530) (-9.104,3.491) (-9.168,3.452) (-9.231,3.411) (-9.293,3.369) (-9.355,3.327) (-9.416,3.283) (-9.477,3.239) (-9.536,3.194) (-9.595,3.147) (-9.654,3.100) (-9.711,3.052) (-9.768,3.003) (-9.824,2.953) (-9.879,2.902) (-9.933,2.850) (-9.986,2.797) (-10.000,2.780)}

    % Dipolo puntual - puntas de flecha
% Generado por scripts/generar_datos.py (no editar a mano).
\FL{1.429}{-0.015}{-91.8}
\FL{1.467}{2.168}{37.2}
\FL{4.190}{3.194}{4.1}
\FL{6.989}{2.463}{-35.4}
\FL{8.330}{-0.035}{-90.7}
\FL{6.953}{-2.489}{-145.3}
\FL{4.146}{-3.191}{175.4}
\FL{1.433}{-2.141}{142.3}
\FL{-1.467}{2.168}{142.8}
\FL{-4.190}{3.194}{175.9}
\FL{-6.989}{2.463}{-144.6}
\FL{-8.330}{-0.035}{-89.3}
\FL{-6.953}{-2.489}{-34.7}
\FL{-4.146}{-3.191}{4.6}
\FL{-1.433}{-2.141}{37.7}
\FL{-1.429}{-0.015}{-88.2}
\FL{3.888}{-8.732}{126.6}
\FL{2.141}{-6.053}{119.5}
\FL{0.773}{-3.164}{110.7}
\FL{-4.014}{-8.749}{52.4}
\FL{-2.211}{-6.078}{59.7}
\FL{-0.799}{-3.182}{68.8}
\FL{8.283}{-6.321}{175.9}
\FL{4.997}{-5.511}{156.6}
\FL{2.162}{-3.665}{137.0}
\FL{-4.602}{2.370}{-163.1}
\FL{-6.550}{-0.011}{-89.7}
\FL{-4.581}{-2.377}{-16.5}
\FL{0.881}{2.999}{65.3}
\FL{2.174}{5.280}{56.0}
\FL{3.771}{7.329}{48.3}
\FL{5.631}{9.179}{41.4}
\FL{-5.151}{3.827}{177.3}
\FL{-7.813}{3.273}{-152.7}
\FL{-9.682}{1.331}{-113.2}
\FL{-9.677}{-1.343}{-66.6}
\FL{-7.802}{-3.279}{-27.1}
\FL{-5.138}{-3.826}{2.8}
\FL{-0.035}{-8.072}{89.6}
\FL{-0.013}{-4.222}{89.7}
\FL{-1.128}{3.214}{119.3}
\FL{-2.813}{5.599}{130.8}
\FL{-4.925}{7.654}{140.6}
\FL{-7.322}{9.325}{149.6}
\FL{-3.065}{4.757}{140.7}
\FL{-5.627}{6.350}{155.4}
\FL{-8.507}{7.248}{170.0}
\FL{-8.690}{-7.903}{13.5}
\FL{-6.169}{-7.023}{25.0}
\FL{-3.895}{-5.671}{36.6}
\FL{-1.952}{-3.877}{49.1}
\FL{5.299}{2.380}{-23.9}
\FL{6.981}{-0.015}{-90.4}
\FL{5.272}{-2.392}{-156.6}
\FL{8.439}{-7.362}{168.7}
\FL{5.486}{-6.366}{154.1}
\FL{4.915}{3.849}{5.5}
\FL{7.491}{3.499}{-21.9}
\FL{9.507}{1.882}{-57.2}
\FL{3.579}{6.089}{43.2}
\FL{5.947}{7.925}{32.5}
\FL{8.605}{9.307}{22.4}
\FL{3.187}{1.458}{-23.0}
\FL{3.151}{-1.473}{-158.1}
\FL{2.160}{-8.592}{111.3}
\FL{1.153}{-5.714}{107.2}
\FL{-4.995}{6.777}{146.6}
\FL{-8.259}{8.389}{160.8}
\FL{9.384}{0.914}{-73.4}
\FL{8.847}{-1.977}{-126.7}
\FL{6.306}{-3.544}{-167.7}
\FL{7.949}{-9.303}{153.6}
\FL{5.005}{-7.467}{142.4}
\FL{2.512}{-5.091}{130.1}
\FL{-8.679}{-5.784}{-3.2}
\FL{-6.027}{-5.529}{13.9}
\FL{-3.557}{-4.524}{30.4}
\FL{8.555}{-5.231}{-172.2}
\FL{5.625}{-5.084}{166.8}
\FL{-0.941}{4.720}{107.0}
\FL{-2.243}{8.288}{112.8}
\FL{-8.254}{-8.950}{22.6}
\FL{-5.066}{-7.187}{35.4}
\FL{-2.308}{-8.603}{67.3}
\FL{-1.253}{-5.795}{71.6}
\FL{3.758}{2.167}{-10.9}
\FL{5.781}{-0.009}{-90.3}
\FL{3.741}{-2.171}{-169.4}
\FL{-0.004}{4.753}{90.1}
\FL{-0.010}{8.253}{90.1}
\FL{6.083}{4.475}{2.1}
\FL{8.815}{3.977}{-23.7}
\FL{-2.961}{3.778}{149.5}
\FL{-5.722}{4.766}{171.2}
\FL{-8.648}{4.626}{-165.1}
\FL{4.561}{6.002}{32.0}
\FL{8.107}{7.526}{14.6}
\FL{-2.725}{1.238}{-156.7}
\FL{-2.701}{-1.247}{-22.5}
\FL{-8.897}{-3.737}{-27.2}
\FL{-6.345}{-4.367}{-1.5}
\FL{-3.980}{1.547}{-149.1}
\FL{-3.966}{-1.555}{-30.5}
\FL{-5.638}{5.666}{161.3}
\FL{-8.531}{6.181}{178.7}
\FL{2.629}{3.737}{35.5}
\FL{5.423}{5.041}{14.6}
\FL{8.497}{5.250}{-7.3}
\FL{2.013}{6.217}{62.9}
\FL{3.508}{8.788}{57.0}
\FL{-2.126}{6.103}{119.1}
\FL{-3.822}{8.756}{125.9}
\FL{0.825}{5.102}{76.2}
\FL{1.763}{8.396}{72.1}
\FL{8.584}{-8.611}{161.4}
\FL{5.940}{-7.430}{150.4}
\FL{5.637}{5.904}{20.8}
\FL{8.529}{6.557}{4.5}
\FL{0.666}{8.106}{83.0}
\FL{-0.752}{7.848}{98.2}
\FL{2.733}{-0.005}{-90.3}
\FL{0.983}{-7.660}{101.0}
\FL{-5.711}{-8.978}{39.9}
\FL{-3.341}{-6.651}{49.2}
\FL{8.932}{-3.512}{-149.6}
\FL{6.364}{-4.257}{-177.0}
\FL{3.714}{-3.846}{159.8}
\FL{-7.737}{-4.911}{-5.8}
\FL{-7.426}{-6.768}{13.6}
\FL{5.415}{-9.049}{137.6}
\FL{3.410}{-6.957}{129.9}
\FL{-4.368}{8.279}{132.6}
\FL{-2.652}{-0.012}{-89.2}
\FL{-0.967}{-7.489}{78.9}
\FL{-7.971}{4.012}{-161.9}

    \draw[mmaverde!85!black, line width=2.1pt, dash pattern=on 4pt off 2.5pt] (0,1.2) -- (0,9.4);
    \draw[mmanaranja, line width=2.1pt, dash pattern=on 4pt off 2.5pt] (0.85,0.85) -- (6.6,6.6);
    \draw[mmaazul, line width=2.1pt, dash pattern=on 4pt off 2.5pt] (1.2,0) -- (9.4,0);
    \node[fill=mmafondo, inner sep=1pt, font=\scriptsize, text=mmaverde!70!black] at (1.75,8.6) {$\theta=0$};
    \node[fill=mmafondo, inner sep=1pt, font=\scriptsize, text=mmanaranja!80!black] at (7.4,5.6) {$\theta=\pi/4$};
    \node[fill=mmafondo, inner sep=1pt, font=\scriptsize, text=mmaazul!80!black] at (7.6,1.15) {$\theta=\pi/2$};}
  \dipoloRojo{0.85}
  \marcoMMA{10}{5}{$x/\ell$}{$z/\ell$}
\end{tikzpicture}
\end{minipage}\hfill
\begin{minipage}[c]{0.48\textwidth}\centering
\begin{tikzpicture}
\begin{axis}[marcoMMA, width=6.6cm, height=4.3cm, xmin=1, xmax=10.5, ymin=0, ymax=34,
             ylabel={$\varepsilon$ (\%)}, xlabel={$r/\ell$}, ylabel style={yshift=-4pt},
             xtick={2,4,6,8,10}, ytick={0,10,20,30}, every axis plot/.style={line width=1pt}]
  \addplot[mmaazul]              table[x=r, y=e90] {figuras/datos/discrepancia.dat};
  \addplot[mmanaranja]           table[x=r, y=e45] {figuras/datos/discrepancia.dat};
  \addplot[mmaverde!85!black]    table[x=r, y=e0]  {figuras/datos/discrepancia.dat};
  \draw[pizarra, densely dashed, line width=0.6pt] (axis cs:5,-0.6) rectangle (axis cs:10.3,3.0);
  \node[font=\scriptsize, text=mmaverde!70!black, anchor=west] at (axis cs:1.75,27) {$\theta=0$};
  \node[font=\scriptsize, text=mmaazul!80!black, anchor=west] at (axis cs:2.4,15) {$\theta=\pi/2$};
  \node[font=\scriptsize, text=mmanaranja!80!black, anchor=west] at (axis cs:1.3,3.2) {$\theta=\pi/4$};
\end{axis}
\end{tikzpicture}\\[4pt]
\begin{tikzpicture}
\begin{axis}[marcoMMA, width=6.6cm, height=4.3cm, xmin=5, xmax=15, ymin=0, ymax=2.1,
             ylabel={$\varepsilon$ (\%)}, xlabel={$r/\ell$}, ylabel style={yshift=-4pt},
             xtick={6,8,...,14}, ytick={0,0.5,1,1.5,2}, every axis plot/.style={line width=1pt},
             y tick label style={/pgf/number format/.cd, fixed, precision=1, use comma}]
  \addplot[mmaazul]              table[x=r, y=e90] {figuras/datos/discrepancia.dat};
  \addplot[mmanaranja]           table[x=r, y=e45] {figuras/datos/discrepancia.dat};
  \addplot[mmaverde!85!black]    table[x=r, y=e0]  {figuras/datos/discrepancia.dat};
  \node[font=\scriptsize, text=mmaverde!70!black, anchor=west] at (axis cs:6.1,1.55) {$\theta=0$};
  \node[font=\scriptsize, text=mmaazul!80!black, anchor=south west] at (axis cs:7.4,0.78) {$\theta=\pi/2$};
  \node[font=\scriptsize, text=mmanaranja!80!black, anchor=north west] at (axis cs:8.2,0.27) {$\theta=\pi/4$};
\end{axis}
\end{tikzpicture}
\end{minipage}
\caption{Discrepancia entre el dipolo físico y el puntual. Izquierda: campo del dipolo puntual con las tres semirrectas sobre las que se
evalúa. Derecha: discrepancia relativa $\varepsilon=\norm{\vE_{p}-\vE_{f}}/\norm{\vE_{f}}$ (\cref{prop:discrepancia}) en función de la distancia; el
panel inferior amplía la región recuadrada y la prolonga hasta $r=15\ell$. La discrepancia cae como $(\ell/r)^{2}$ y es inferior al $0{,}5\%$
para $r>10\ell$ en cualquier dirección. Recreación de~\cite{FIGURAS-T2} con los valores exactos.}
\label{fig:discrepancia}
\end{figure}


% =========================================================================
\section{Distribuciones de dipolos: la polarización}\label{sec:polarizacion}
% =========================================================================
Del mismo modo que una distribución continua de carga es el paso al continuo de un sistema de cargas puntuales, $\rho=\dif q/\dVp$,
podemos considerar una distribución continua de \emph{dipolos}. Es la situación de un material dieléctrico, en el que un número enorme de
dipolos moleculares, muy próximos entre sí, ocupa una región $V'$ del espacio.

% -------------------------------------------------------------------------
\subsection{Densidad de momento dipolar}\label{sec:pol-def}
% -------------------------------------------------------------------------
\begin{definicion}[Polarización]\label{def:polarizacion}
Sea una región $V'$ con una distribución continua de dipolos. Se llama \emph{polarización} al campo vectorial
\begin{equation}\label{eq:polarizacion}
  \vP(\rp)=\frac{\dif\vp}{\dVp},\qquad\text{es decir,}\qquad \dif\vp=\vP(\rp)\,\dVp ,
\end{equation}
momento dipolar por unidad de volumen en el punto fuente $\rp$. Su unidad es $[\vP]=\mathrm{C\,m}/\mathrm{m^{3}}=\mathrm{C/m^{2}}$, la misma que la de una
densidad superficial de carga.
\end{definicion}

\begin{observacion}[Analogía con la densidad de carga]\label{obs:analogia}
\begin{center}\vspace{2pt}
\renewcommand{\arraystretch}{1.25}
\begin{tabular}{@{}lcc@{}}
  \toprule
  & \textsc{Carga puntual} & \textsc{Dipolo puntual}\\
  \midrule
  Magnitud discreta & $q$ & $\vp$\\
  Densidad continua & $\rho(\rp)=\dif q/\dVp$ & $\vP(\rp)=\dif\vp/\dVp$\\
  Unidad de la densidad & $\mathrm{C/m^{3}}$ & $\mathrm{C/m^{2}}$\\
  Potencial elemental & $\dfrac{\dif q}{4\pi\eo R}$ & $\dfrac{\dif\vp\cdot\uR}{4\pi\eo R^{2}}$\\
  \bottomrule
\end{tabular}
\end{center}
\vspace{2pt}
La polarización es una magnitud fundamental en el estudio de los materiales dieléctricos, que abordaremos más adelante en el curso.
\end{observacion}

% =========================================================================
% Figura (apuntes manuscritos, pág. 5, y de teoría, pág. II-4): región V'
% polarizada; el elemento dV' en r' contiene el dipolo dp = P dV', que crea
% en P (posición r) el potencial dphi = dp.u_R/(4 pi eps0 R^2).
% =========================================================================
\begin{figure}[ht]
\centering
\begin{tikzpicture}[x=1cm, y=1cm, >=Stealth]
  \ejesO{4.4}{6.6}{1.6}
  \regionV
  % pequeños dipolos (polarización) en el interior
  \foreach \px/\py in {1.45/2.5, 1.75/3.35, 2.55/3.55, 3.15/3.1, 3.2/2.25, 2.45/1.95, 1.75/1.95}{%
    \draw[-{Stealth[length=2.6pt,width=2.2pt]}, magacero!75, line width=0.55pt] (\px-0.12,\py-0.08) -- (\px+0.14,\py+0.1);}
  \node[etiq, text=magprof] at (0.8,3.65) {$V'$};
  % elemento de volumen
  \coordinate (Q) at (2.35,2.75);
  \coordinate (P) at (6.2,3.25);
  \draw[vect] (0,0) -- ($(Q)+(-0.1,-0.12)$) node[pos=0.5, above left=-1pt, etiq] {$\rp$};
  \cubito{Q}
  \node[etiq, font=\footnotesize, anchor=south east] at ($(Q)+(-0.1,0.12)$) {$\dVp$};
  \draw[vect] (0,0) -- (P) node[pos=0.72, below right=-2pt, etiq] {$\vc{r}$};
  \draw[vectrojo] ($(Q)+(0.22,0.06)$) -- (P) node[pos=0.6, above=2pt, etiq, text=rojoteo] {$\vc{R}=\vc{r}-\rp$};
  \fill[magprof] (P) circle (1.5pt) node[right=2pt, etiq] {$P$};
  % normal exterior
  \coordinate (B) at (3.3,1.8);
  \draw[vectverde] (B) -- ++(0.42,-0.44) node[below right=-2pt, etiq, text=verdeteo] {$\un$};
  \node[etiq, font=\footnotesize, text=magacero, anchor=south] at (2.45,3.98) {$S'=\partial V'$};
  \fill[black] (0,0) circle (1.2pt) node[below right=1pt, etiq, font=\footnotesize] {$O$};
  % recuadro con las relaciones
  \node[cuadro, anchor=north west, align=left, font=\footnotesize, inner sep=4pt] at (7.0,2.6)
    {$\dif\vp=\vP(\rp)\,\dVp$\\[4pt]
     $\dif\phi(\vc{r})=\dfrac{\dif\vp\cdot\uR}{4\pi\eo R^{2}}$\\[7pt]
     $\rhoP=-\nabp\cdot\vP$,\quad $\sigP=\vP\cdot\un$};
\end{tikzpicture}
\caption{Distribución continua de dipolos en una región $V'$ limitada por la superficie $S'$, de normal exterior $\un$. El elemento $\dVp$ situado en $\rp$
contiene el momento dipolar $\dif\vp=\vP\,\dVp$, que contribuye al potencial en el punto campo $P$ como un dipolo puntual a la distancia $R$. El potencial
total equivale al de las cargas de polarización $\rhoP$ y $\sigP$ (\cref{teo:cargas-pol}).}
\label{fig:polarizacion}
\end{figure}


% -------------------------------------------------------------------------
\subsection{Potencial de una distribución de dipolos}\label{sec:pol-potencial}
% -------------------------------------------------------------------------
Cada elemento de volumen $\dVp$ situado en $\rp$ contiene un dipolo $\dif\vp=\vP\,\dVp$ que, por la \cref{prop:dipolo-Rp}, crea en $\vc{r}$ el potencial
$\dif\phi=\dif\vp\cdot\uR/(4\pi\eo R^{2})$ (\cref{fig:polarizacion}). Por superposición,
\begin{equation}\label{eq:pot-polarizacion}
  \phi(\vc{r})=\frac{1}{4\pi\eo}\int_{V'}\frac{\vP(\rp)\cdot\uR}{R^{2}}\,\dVp=\frac{1}{4\pi\eo}\int_{V'}\frac{\vP(\rp)\cdot\vc{R}}{R^{3}}\,\dVp .
\end{equation}
Esta integral tiene una interpretación sorprendente: el potencial es el mismo que crearía una cierta distribución de \emph{carga}, en volumen y en la
superficie que limita $V'$.

\begin{lema}[Gradiente de $1/R$ respecto de las coordenadas fuente]\label{lem:grad-1R}
Para $\vc{r}\neq\rp$,
\begin{equation}\label{eq:grad-1R}
  \nabp\Bigl(\frac{1}{R}\Bigr)=\frac{\vc{R}}{R^{3}}=\frac{\uR}{R^{2}}=-\nab\Bigl(\frac{1}{R}\Bigr).
\end{equation}
\end{lema}
\begin{proof}
Es el gradiente de $R^{n}$ con $n=-1$ estudiado en el Tema~0~\cite{TEMA0}: $\nab R^{n}=nR^{n-1}\uR$, de donde $\nab(1/R)=-\uR/R^{2}$. Como
$\vc{R}=\vc{r}-\rp$, derivar respecto de $x'_j$ equivale a derivar respecto de $x_j$ cambiando el signo, luego $\nabp(1/R)=+\uR/R^{2}$.
\end{proof}

\begin{teorema}[Cargas de polarización]\label{teo:cargas-pol}
Sea $\vP$ una polarización de clase $C^{1}$ en una región $V'$ limitada por una superficie cerrada $S'=\partial V'$ con normal unitaria exterior
$\un$. El potencial \cref{eq:pot-polarizacion} es
\begin{equation}\label{eq:pot-cargas-pol}
  \phi(\vc{r})=\frac{1}{4\pi\eo}\int_{V'}\frac{\rhoP(\rp)}{R}\,\dVp+\frac{1}{4\pi\eo}\oiint_{S'}\frac{\sigP(\rp)}{R}\,\dSp,
\end{equation}
donde
\begin{equation}\label{eq:densidades-pol}
  \boxed{\;\rhoP=-\nabp\cdot\vP,\qquad \sigP=\vP\cdot\un\;}
\end{equation}
son las \emph{densidades volúmica y superficial de carga de polarización}.
\end{teorema}
\begin{proof}
Por el \cref{lem:grad-1R}, el integrando de \cref{eq:pot-polarizacion} es $\vP\cdot\nabp(1/R)$. La regla del producto para la divergencia,
$\nabp\cdot(f\vc{A})=f\,\nabp\cdot\vc{A}+\vc{A}\cdot\nabp f$, con $f=1/R$ y $\vc{A}=\vP$, da
\[
  \nabp\cdot\Bigl(\frac{\vP}{R}\Bigr)=\frac{\nabp\cdot\vP}{R}+\vP\cdot\nabp\Bigl(\frac{1}{R}\Bigr)
  \quad\Longrightarrow\quad
  \frac{\vP\cdot\uR}{R^{2}}=-\frac{\nabp\cdot\vP}{R}+\nabp\cdot\Bigl(\frac{\vP}{R}\Bigr).
\]
Integrando en $V'$,
\[
  \phi(\vc{r})=\frac{1}{4\pi\eo}\Biggl[\int_{V'}\frac{-\nabp\cdot\vP}{R}\,\dVp+\int_{V'}\nabp\cdot\Bigl(\frac{\vP}{R}\Bigr)\dVp\Biggr],
\]
y por el teorema de la divergencia (de Gauss)~\cite{TEMA0}, con $\dif\vc{S}\,'=\dSp\,\un$, la segunda integral es
$\oiint_{S'}\frac{\vP\cdot\un}{R}\,\dSp$. Basta llamar $\rhoP$ y $\sigP$ a los numeradores. (Si $\vc{r}\in V'$, la singularidad en $R=0$ es integrable y
el resultado se mantiene excluyendo una pequeña esfera alrededor de $\vc{r}$ y haciendo tender su radio a cero.)
\end{proof}

\begin{observacion}[Cargas ligadas]\label{obs:ligadas}
Comparando \cref{eq:pot-cargas-pol} con el potencial de una distribución de carga (\cref{eq:integrales}), una distribución de dipolos equivale, a efectos de
potencial y de campo, a una carga $\rhoP$ en volumen más una carga $\sigP$ en superficie. Estas cargas no son libres: proceden del desplazamiento
microscópico de las cargas ligadas de cada molécula, y por eso la terminología habitual de los dieléctricos las llama \emph{cargas ligadas}. La normal
$\un$ apunta, por convenio, hacia fuera del volumen. Físicamente, $\sigP$ es la carga que «asoma» en la superficie al no estar compensada por la
molécula vecina, y $\rhoP$ aparece donde la polarización no es uniforme (si $\vP$ es constante, $\rhoP=0$ y solo hay carga superficial).
\end{observacion}

\begin{proposicion}[La carga total de polarización es nula]\label{prop:carga-pol-nula}
\begin{equation}\label{eq:carga-pol-nula}
  Q_{P}=\int_{V'}\rhoP\,\dVp+\oiint_{S'}\sigP\,\dSp=0 .
\end{equation}
\end{proposicion}
\begin{proof}
Por el teorema de la divergencia,
\[
  Q_{P}=-\int_{V'}\nabp\cdot\vP\,\dVp+\oiint_{S'}\vP\cdot\un\,\dSp=-\oiint_{S'}\vP\cdot\dif\vc{S}\,'+\oiint_{S'}\vP\cdot\dif\vc{S}\,'=0 .
\]
Es lo esperable: cada dipolo es neutro, y por tanto también lo es cualquier conjunto de ellos.
\end{proof}

% =========================================================================
\section{El desarrollo multipolar del potencial}\label{sec:multipolar}
% =========================================================================
Volvemos al problema general: una distribución $\rho(\rp)$ localizada en un volumen $V'$, cuyo potencial (\cref{eq:pot-integral}) queremos aproximar en
puntos alejados, $r\gg r'$ (\cref{fig:multipolar-geom}). El único ingrediente del integrando que depende a la vez del punto campo y del punto
fuente es $1/R$; desarrollándolo en potencias de $r'/r$ separaremos ambas dependencias.

% =========================================================================
% Figura (apuntes manuscritos, pág. 6; teoría, pág. II-5): geometría del
% desarrollo multipolar. theta' es el ángulo entre r (punto campo) y r'
% (punto fuente); R = |r - r'|.
% =========================================================================
\begin{figure}[ht]
\centering
\begin{tikzpicture}[x=1cm, y=1cm, >=Stealth]
  \draw[eje] (0,0) -- (0,3.3) node[above, etiq] {$z$};
  \draw[eje] (0,0) -- (6.4,0) node[right, etiq] {$y$};
  \draw[eje] (0,0) -- (-1.4,-1.2) node[below left=-2pt, etiq] {$x$};
  \draw[magacero, line width=0.7pt, fill=rosaclaro!55]
    plot[smooth cycle, tension=0.8] coordinates {(-1.2,0.6) (-0.5,2.2) (0.9,2.6) (1.9,1.9) (1.6,0.5) (0.5,-0.3) (-0.8,-0.4)};
  \node[etiq, text=magprof] at (-0.95,1.9) {$V'$};
  \coordinate (O) at (0,0);
  \coordinate (Q) at (0.75,1.6);
  \coordinate (P) at (6.0,2.2);
  \draw[magprof, line width=0.5pt, fill=white] ($(Q)+(-0.12,-0.12)$) rectangle ($(Q)+(0.12,0.12)$);
  \node[etiq, font=\footnotesize, anchor=south] at ($(Q)+(0,0.16)$) {$\dif q=\rho\,\dVp$};
  \draw[vect] (0,0) -- (Q) node[pos=0.5, left=1pt, etiq] {$\rp$};
  \draw[vect] (0,0) -- (P) node[pos=0.55, below right=-2pt, etiq] {$\vc{r}$};
  \draw[vectrojo] ($(Q)+(0.13,0)$) -- (P) node[pos=0.55, above=2pt, etiq, text=rojoteo] {$\vc{R}=\vc{r}-\rp$};
  \pic[draw=magacero, line width=0.5pt, angle radius=0.85cm, "$\theta'$", angle eccentricity=1.38, font=\small] {angle=P--O--Q};
  \fill[magprof] (P) circle (1.5pt) node[right=2pt, etiq] {$P$};
  \node[cuadro, anchor=west, font=\footnotesize, align=left] at (7.0,1.1)
     {$R^{2}=r^{2}+r'^{2}-2rr'\cos\theta'$\\[3pt] $\vc{r}\cdot\rp=rr'\cos\theta'$};
\end{tikzpicture}
\caption{Geometría del desarrollo multipolar: el elemento de carga $\dif q=\rho\,\dVp$, en $\rp$, y el punto campo $P$, en $\vc{r}$, forman con el origen el ángulo
$\theta'$. Cuando $r'\ll r$, el vector $\vc{R}$ es casi paralelo a $\vc{r}$ y $1/R$ admite un desarrollo en potencias de $r'/r$.}
\label{fig:multipolar-geom}
\end{figure}


% -------------------------------------------------------------------------
\subsection{Desarrollo de \texorpdfstring{$1/R$}{1/R} y polinomios de Legendre}\label{sec:desarrollo-1R}
% -------------------------------------------------------------------------
Por el teorema del coseno, $R^{2}=\norm{\vc{r}-\rp}^{2}=r^{2}+r'^{2}-2rr'\cosp$, de modo que
\begin{equation}\label{eq:1R-coseno}
  \frac{1}{R}=\bigl(r^{2}+r'^{2}-2rr'\cosp\bigr)^{-1/2}
  =\frac{1}{r}\Bigl[1+\underbrace{\Bigl(\frac{r'}{r}\Bigr)^{2}-2\,\frac{r'}{r}\cosp}_{x}\Bigr]^{-1/2}.
\end{equation}
Si $r'/r\ll1$, también $\abs{x}\ll1$, y podemos usar el binomio de Newton.

\begin{lema}[Serie binomial de Newton]\label{lem:binomio}
Para todo $n\in\R$ y $\abs{x}<1$, $(1+x)^{n}=\sum_{k=0}^{\infty}\binom{n}{k}x^{k}$, con
$\binom{n}{k}=\frac{n(n-1)\cdots(n-k+1)}{k!}$. En particular, para $n=-\frac12$,
\begin{equation}\label{eq:binomio-12}
  (1+x)^{-1/2}=1-\frac{1}{2}\,x+\frac{3}{8}\,x^{2}-\frac{5}{16}\,x^{3}+\cdots
\end{equation}
\end{lema}
\begin{proof}
Es la serie de Taylor de $(1+x)^{n}$ en $x=0$, cuyo radio de convergencia es $1$ cuando $n\notin\N$. Los coeficientes de \cref{eq:binomio-12} son
$\binom{-1/2}{1}=-\frac12$, $\binom{-1/2}{2}=\frac{(-\frac12)(-\frac32)}{2}=\frac38$ y $\binom{-1/2}{3}=\frac{(-\frac12)(-\frac32)(-\frac52)}{6}=-\frac{5}{16}$.
\end{proof}

Sustituyendo $x$ en \cref{eq:binomio-12} y \emph{agrupando por potencias} de $r'/r$ (y no de $x$, que mezcla órdenes distintos):
\[
  \frac{1}{R}=\frac{1}{r}\Bigl\{1-\frac{1}{2}\Bigl[\Bigl(\frac{r'}{r}\Bigr)^{2}-2\,\frac{r'}{r}\cosp\Bigr]
  +\frac{3}{8}\Bigl[\Bigl(\frac{r'}{r}\Bigr)^{2}-2\,\frac{r'}{r}\cosp\Bigr]^{2}+\cdots\Bigr\}.
\]
El término de orden $r'/r$ procede solo del primer corchete: $+\frac{r'}{r}\cosp$. El de orden $(r'/r)^{2}$ recibe dos contribuciones, $-\frac12(r'/r)^{2}$
del primer corchete y $\frac38\bigl(-2\frac{r'}{r}\cosp\bigr)^{2}=\frac32(r'/r)^{2}\cos^{2}\theta'$ del segundo, que suman
$(r'/r)^{2}\,\frac12(3\cos^{2}\theta'-1)$. Por tanto,
\begin{equation}\label{eq:1R-orden2}
  \frac{1}{R}=\frac{1}{r}\Bigl[1+\frac{r'}{r}\cosp+\Bigl(\frac{r'}{r}\Bigr)^{2}\frac{1}{2}\bigl(3\cos^{2}\theta'-1\bigr)+\ord\Bigl(\frac{r'^{3}}{r^{3}}\Bigr)\Bigr].
\end{equation}
Los coeficientes de cada potencia son polinomios en $\cosp$ que reconocemos enseguida.

\begin{definicion}[Polinomios de Legendre]\label{def:legendre}
Los \emph{polinomios de Legendre} $\Pl{\ell}(x)$, $\ell\in\N$, $x\in[-1,1]$, son los coeficientes del desarrollo de la \emph{función generatriz}
\begin{equation}\label{eq:generatriz}
  g(x,t)=\frac{1}{\sqrt{1-2xt+t^{2}}}=\sum_{\ell=0}^{\infty}\Pl{\ell}(x)\,t^{\ell},\qquad \abs{t}<1 .
\end{equation}
Fueron introducidos por Legendre precisamente al estudiar la atracción gravitatoria de cuerpos extensos~\cite{LEGENDRE}, un problema matemáticamente idéntico al nuestro.
Equivalentemente, vienen dados por la fórmula de Rodrigues $\Pl{\ell}(x)=\frac{1}{2^{\ell}\ell!}\frac{\dif^{\ell}}{\dif x^{\ell}}(x^{2}-1)^{\ell}$. Los primeros son
\begin{equation}\label{eq:legendre}
  \begin{gathered}
  \Pl{0}(x)=1,\qquad \Pl{1}(x)=x,\qquad \Pl{2}(x)=\tfrac12(3x^{2}-1),\\
  \Pl{3}(x)=\tfrac12(5x^{3}-3x),\qquad \Pl{4}(x)=\tfrac18(35x^{4}-30x^{2}+3).
  \end{gathered}
\end{equation}
\end{definicion}

\begin{teorema}[Desarrollo de $1/R$ en polinomios de Legendre]\label{teo:1R-legendre}
Si $r'<r$,
\begin{equation}\label{eq:1R-legendre}
  \boxed{\;\frac{1}{R}=\frac{1}{\norm{\vc{r}-\rp}}=\sum_{\ell=0}^{\infty}\frac{r'^{\,\ell}}{r^{\ell+1}}\,\Pl{\ell}(\cosp)\;}
\end{equation}
y la serie converge absolutamente.
\end{teorema}
\begin{proof}
Por \cref{eq:1R-coseno}, $1/R=\frac1r\bigl(1-2\cosp\,t+t^{2}\bigr)^{-1/2}=\frac1r\,g(\cosp,t)$ con $t=r'/r\in[0,1)$. Basta aplicar \cref{eq:generatriz}. La convergencia
para $\abs{t}<1$ se debe a que, con $\abs{x}\le1$, las raíces de $1-2xt+t^{2}$ son $t=x\pm\mathrm{i}\sqrt{1-x^{2}}$, de módulo $1$: la función $g(x,\cdot)$ es analítica en el disco
$\abs{t}<1$ y su serie de potencias converge en él. El cálculo hasta orden dos, \cref{eq:1R-orden2}, confirma los tres primeros coeficientes.
\end{proof}

\begin{lema}[Acotación de los polinomios de Legendre]\label{lem:legendre-cota}
Para todo $\ell\in\N$ y $x\in[-1,1]$, $\abs{\Pl{\ell}(x)}\le1$, con igualdad en $x=\pm1$: $\Pl{\ell}(1)=1$ y $\Pl{\ell}(-1)=(-1)^{\ell}$.
\end{lema}
\begin{proof}
Se cumple la primera integral de Laplace~\cite{ARFKEN}, $\Pl{\ell}(\cos\theta)=\frac{1}{\pi}\int_{0}^{\pi}\bigl(\cos\theta+\mathrm{i}\sen\theta\cos u\bigr)^{\ell}\dif u$. Como
$\abs{\cos\theta+\mathrm{i}\sen\theta\cos u}=\bigl(\cos^{2}\theta+\sen^{2}\theta\cos^{2}u\bigr)^{1/2}\le1$, el integrando tiene módulo $\le1$ y
$\abs{\Pl{\ell}}\le\frac1\pi\int_0^\pi\dif u=1$. Para los valores en $\pm1$: $g(1,t)=(1-t)^{-1}=\sum t^{\ell}$ y $g(-1,t)=(1+t)^{-1}=\sum(-t)^{\ell}$.
\end{proof}

% =========================================================================
% Figura (clon de «Tema 2. Figuras», diapositiva «Polinomios de Legendre»):
% desarrollo de 1/R para r'/r = 1/3 (r = 3, r' = 1) en función de theta'.
% Arriba, cada término (r'^l / r^{l+1}) P_l(cos theta'); abajo, las sumas
% parciales frente a la función exacta 1/sqrt(10 - 6 cos theta').
% Todo se evalúa analíticamente en pgfplots.
% =========================================================================
\begin{figure}[t]
\centering
\pgfplotsset{legendreejes/.style={marcoMMA, width=8.4cm, height=4.6cm, xmin=0, xmax=6.2832,
    trig format plots=rad, domain=0:6.2832, samples=181, every axis plot/.style={line width=1.1pt},
    xtick={0,1.5708,3.1416,4.7124}, xticklabels={$0$,$\frac{\pi}{2}$,$\pi$,$\frac{3\pi}{2}$},
    xlabel={$\theta'$}, ylabel style={yshift=-3pt}}}
\begin{minipage}[c]{0.33\textwidth}\centering
{\small $\displaystyle\frac{1}{R}=\sum_{\ell=0}^{\infty}\frac{r'^{\,\ell}}{r^{\ell+1}}\,\Pl{\ell}(\cos\theta')$}\\[10pt]
\begin{tikzpicture}[x=1cm, y=1cm, >=Stealth]
  \draw[pizarra, densely dotted, line width=0.5pt] (0,0) circle (2);
  \coordinate (O) at (0,0);
  \coordinate (A) at (40:2);
  \coordinate (B) at (0.67,0);
  \draw[->, line width=1pt, black] (O) -- (A) node[pos=0.55, above left=-1pt, etiq] {$\vc{r}$};
  \draw[->, line width=1pt, mmaazul] (O) -- (B) node[below=2pt, etiq, text=mmaazul] {$\rp$};
  \draw[->, line width=1pt, mmarojo] (B) -- (A) node[pos=0.5, right=1pt, etiq, text=mmarojo] {$\vc{R}$};
  \pic[draw=mmaazul, line width=0.5pt, angle radius=0.42cm, "$\theta'$", angle eccentricity=1.7, font=\scriptsize] {angle=B--O--A};
\end{tikzpicture}\\[4pt]
{\footnotesize $r'/r=1/3$}
\end{minipage}\hfill
\begin{minipage}[c]{0.64\textwidth}\centering
\begin{tikzpicture}
\begin{axis}[legendreejes, ymin=-0.14, ymax=0.37, ylabel={$(1/R)_{\ell}$},
             ytick={-0.1,0,0.1,0.2,0.3},
             y tick label style={/pgf/number format/.cd, fixed, precision=1, use comma}]
  \addplot[grispizarra, line width=0.4pt] {0};
  \addplot[mmaazul]     {1/3};
  \addplot[mmanaranja]  {(1/9)*cos(x)};
  \addplot[mmaverde]    {(1/27)*0.5*(3*cos(x)^2-1)};
  \addplot[mmarojizo]   {(1/81)*0.5*(5*cos(x)^3-3*cos(x))};
  \node[font=\scriptsize, text=mmaazul!80!black, fill=white, inner sep=1pt] at (axis cs:3.1416,0.30) {$\ell=0$};
  \node[font=\scriptsize, text=mmanaranja!80!black, anchor=south] at (axis cs:0.55,0.12) {$\ell=1$};
  \node[font=\scriptsize, text=mmaverde!70!black, anchor=south] at (axis cs:3.1416,0.045) {$\ell=2$};
  \node[font=\scriptsize, text=mmarojizo!85!black, anchor=north] at (axis cs:4.95,-0.025) {$\ell=3$};
\end{axis}
\end{tikzpicture}\\[2pt]
\begin{tikzpicture}
\begin{axis}[legendreejes, ymin=0.22, ymax=0.52, ylabel={$\sum_{\ell\le n}(1/R)_{\ell}$},
             ytick={0.25,0.30,...,0.50},
             y tick label style={/pgf/number format/.cd, fixed, fixed zerofill, precision=2, use comma},
             legend columns=5, legend cell align=left,
             legend style={font=\scriptsize, draw=none, at={(0.5,-0.42)}, anchor=north,
                           /tikz/every even column/.append style={column sep=6pt}}]
  \addplot[mmavioleta, line width=2.6pt, opacity=0.55] {1/sqrt(10-6*cos(x))};
  \addplot[mmanaranja, line width=1.3pt] {1/3};
  \addplot[mmaverde]   {1/3+(1/9)*cos(x)};
  \addplot[mmarojizo]  {1/3+(1/9)*cos(x)+(1/27)*0.5*(3*cos(x)^2-1)};
  \addplot[mmaazul, line width=0.8pt]
     {1/3+(1/9)*cos(x)+(1/27)*0.5*(3*cos(x)^2-1)+(1/81)*0.5*(5*cos(x)^3-3*cos(x))};
  \legend{exacta, $n=0$, $n=1$, $n=2$, $n=3$}
\end{axis}
\end{tikzpicture}
\end{minipage}
\caption{Desarrollo de $1/R$ en polinomios de Legendre (\cref{teo:1R-legendre}) para $r=3$ y $r'=1$ ($r'/r=1/3$), en función del ángulo $\theta'$ entre $\vc{r}$ y
$\rp$. Arriba: los cuatro primeros términos $(1/R)_{\ell}=r'^{\,\ell}\Pl{\ell}(\cos\theta')/r^{\ell+1}$, cuya amplitud decrece en progresión geométrica de razón $1/3$.
Abajo: sumas parciales hasta orden $n$ frente a la función exacta $1/R=(10-6\cos\theta')^{-1/2}$; con $n=3$ el error máximo es ya de solo un $1{,}2\%$, y cada orden adicional lo divide aproximadamente por tres. Recreación
de~\cite{FIGURAS-T2}.}
\label{fig:legendre}
\end{figure}


\begin{observacion}[Convergencia del desarrollo]\label{obs:convergencia}
La \cref{fig:legendre} ilustra el teorema para $r'/r=1/3$. Cada término está acotado por $(r'/r)^{\ell}/r$ (\cref{lem:legendre-cota}), de modo que la serie
converge como una geométrica de razón $r'/r$: con solo cuatro términos la suma parcial es prácticamente indistinguible de la función exacta (el error es del orden del $1\%$). Cuanto más lejos esté el
punto campo, más rápida es la convergencia.
\end{observacion}

% -------------------------------------------------------------------------
\subsection{El desarrollo multipolar}\label{sec:desarrollo}
% -------------------------------------------------------------------------
\begin{teorema}[Desarrollo multipolar del potencial]\label{teo:multipolar}
Sea $\rho$ una distribución acotada contenida en la bola de radio $a$ centrada en el origen, $\rho(\rp)=0$ si $r'>a$. Para todo punto campo con $r>a$,
\begin{equation}\label{eq:desarrollo-multipolar}
  \boxed{\;\phi(\vc{r})=\sum_{\ell=0}^{\infty}\phimul{\ell}(\vc{r}),\qquad
  \phimul{\ell}(\vc{r})=\frac{1}{4\pi\eo}\,\frac{1}{r^{\ell+1}}\int_{V'}r'^{\,\ell}\,\rho(\rp)\,\Pl{\ell}(\cosp)\,\dVp\;}
\end{equation}
y la serie converge.
\end{teorema}
\begin{proof}
Sustituimos \cref{eq:1R-legendre}, válido porque $r'\le a<r$ en todo el soporte de $\rho$, en \cref{eq:pot-integral}:
\[
  \phi(\vc{r})=\frac{1}{4\pi\eo}\int_{V'}\rho(\rp)\sum_{\ell=0}^{\infty}\frac{r'^{\,\ell}}{r^{\ell+1}}\Pl{\ell}(\cosp)\,\dVp .
\]
Por el \cref{lem:legendre-cota}, el término $\ell$-ésimo de la serie está acotado por $\abs{\rho}\,a^{\ell}/r^{\ell+1}$, y $\sum_\ell a^{\ell}/r^{\ell+1}$ converge
porque $a/r<1$. Por el criterio $M$ de Weierstrass la serie converge uniformemente en $V'$, lo que permite intercambiar suma e integral; finalmente sacamos de la integral
el factor $1/r^{\ell+1}$, que no depende de las coordenadas fuente.
\end{proof}

\begin{observacion}[Lectura del desarrollo]\label{obs:lectura}
\begin{enumerate}[label=(\alph*)]
  \item \emph{Es exacto}: \cref{eq:desarrollo-multipolar} no es una aproximación, sino una reescritura del potencial válida fuera de la esfera que contiene la
        distribución. La aproximación consiste en \emph{truncarlo}.
  \item \emph{Jerarquía}: el término $\ell$-ésimo decae como $1/r^{\ell+1}$, de modo que a grandes distancias los primeros términos no nulos son los dominantes.
        Reciben los nombres de término \emph{monopolar} ($\ell=0$), \emph{dipolar} ($\ell=1$), \emph{cuadrupolar} ($\ell=2$), \emph{octupolar} ($\ell=3$), \dots: el
        término $\ell$-ésimo corresponde a un $2^{\ell}$-polo.
  \item \emph{Lo que falta}: el ángulo $\theta'$ depende tanto de $\rp$ como de $\vc{r}$, así que en cada integral siguen mezcladas las coordenadas del punto campo
        y las del punto fuente. Nuestro objetivo en lo que sigue es \emph{sacar de la integral la dependencia angular} del punto campo, para que cada término sea un
        momento de la distribución (que se calcula una sola vez) multiplicado por una función explícita de $(r,\theta,\varphi)$.
\end{enumerate}
\end{observacion}

% -------------------------------------------------------------------------
\subsection{Término monopolar}\label{sec:monopolar}
% -------------------------------------------------------------------------
Con $\ell=0$, $\Pl{0}=1$ y $r'^{\,0}=1$, la integral es la carga total:

\begin{proposicion}[Término monopolar]\label{prop:monopolar}
\begin{equation}\label{eq:monopolar}
  \boxed{\;\phimul{0}(\vc{r})=\frac{Q}{4\pi\eo\,r},\qquad Q=\int_{V'}\rho(\rp)\,\dVp\;}
\end{equation}
La carga neta $Q$ se denomina \emph{momento monopolar} (momento multipolar de orden cero) de la distribución.
\end{proposicion}

\begin{observacion}\label{obs:monopolar}
El término monopolar coincide con el potencial de una carga puntual $Q$ situada en el origen. Decae como $1/r$ y es el dominante siempre que $Q\neq0$: es la
afirmación precisa de que, \emph{suficientemente} lejos, toda distribución con carga neta parece una carga puntual. Obviamente, $Q$ no depende del origen de
coordenadas escogido.
\end{observacion}

% -------------------------------------------------------------------------
\subsection{Término dipolar}\label{sec:dipolar}
% -------------------------------------------------------------------------
Con $\ell=1$, $\Pl{1}(\cosp)=\cosp$ y
\[
  \phimul{1}(\vc{r})=\frac{1}{4\pi\eo\,r^{2}}\int_{V'}\rho(\rp)\,r'\cosp\,\dVp .
\]
La clave para separar variables es escribir $r'\cosp$ como un producto escalar: $r'\cosp=\rp\cdot\ur$, donde $\ur=\vc{r}/r$ \emph{no} depende de las coordenadas fuente
y puede salir de la integral.

\begin{proposicion}[Término dipolar]\label{prop:dipolar}
\begin{equation}\label{eq:dipolar}
  \boxed{\;\phimul{1}(\vc{r})=\frac{\vp\cdot\ur}{4\pi\eo\,r^{2}},\qquad \vp=\int_{V'}\rp\,\rho(\rp)\,\dVp\;}
\end{equation}
El vector $\vp$ es el \emph{momento dipolar} (momento multipolar de orden uno) de la distribución.
\end{proposicion}
\begin{proof}
$\int_{V'}\rho\,r'\cosp\,\dVp=\int_{V'}\rho\,(\rp\cdot\ur)\,\dVp=\bigl(\int_{V'}\rho\,\rp\,\dVp\bigr)\cdot\ur=\vp\cdot\ur$, por la linealidad de la integral.
\end{proof}

\begin{observacion}[Generalización del momento dipolar]\label{obs:p-general}
La \cref{eq:dipolar} coincide con el potencial de un dipolo puntual situado en el origen (\cref{prop:pot-dipolo}), y generaliza la definición $\vp=q\vl$ a
distribuciones continuas. Para un sistema de cargas puntuales, $\rho(\rp)=\sum_iq_i\,\delta^{3}(\rp-\rp_i)$, la propiedad de filtrado de la delta de
Dirac~\cite{TEMA1} da
\begin{equation}\label{eq:p-discreto}
  \vp=\int_{V'}\rp\sum_iq_i\,\delta^{3}(\rp-\rp_i)\,\dVp=\sum_{i}q_i\,\rp_i ,
\end{equation}
y para el dipolo físico, $\vp=q\,\rp_{+}-q\,\rp_{-}=q\vl$, como debía ser. Las distribuciones superficiales y lineales se tratan igual, cambiando
$\rho\,\dVp$ por $\sigma\,\dSp$ o $\lambda\,\dlp$.
\end{observacion}

A diferencia de la carga neta, el momento dipolar \emph{depende del origen} de coordenadas.

% =========================================================================
% Figura (apuntes manuscritos, pág. 7; teoría, pág. II-7): cambio de origen
% del momento dipolar. r'' = r' - OO'.
% =========================================================================
\begin{figure}[ht]
\centering
\begin{tikzpicture}[x=1cm, y=1cm, >=Stealth]
  \draw[magacero, line width=0.7pt, fill=rosaclaro!55]
    plot[smooth cycle, tension=0.8] coordinates {(-1.3,1.6) (0.2,3.1) (2.2,3.4) (3.4,2.4) (3.0,0.9) (1.4,0.3) (-0.3,0.4)};
  \node[etiq, text=magprof] at (3.55,3.15) {$V'$};
  \coordinate (O) at (0,0);
  \coordinate (Op) at (-0.75,1.45);
  \coordinate (Q) at (1.9,2.5);
  % ejes del origen O
  \draw[eje] (O) -- (0,1.0) node[above, etiq, font=\footnotesize] {$z$};
  \draw[eje] (O) -- (4.3,0) node[right, etiq] {$y$};
  \draw[eje] (O) -- (-0.9,-0.75);
  \fill[black] (O) circle (1.4pt) node[below right=1pt, etiq] {$O$};
  \fill[black] (Op) circle (1.4pt) node[left=2pt, etiq] {$O'$};
  \draw[magprof, line width=0.5pt, fill=white] ($(Q)+(-0.12,-0.12)$) rectangle ($(Q)+(0.12,0.12)$);
  \node[etiq, font=\footnotesize, anchor=south west] at ($(Q)+(0.12,0.08)$) {$\dVp$};
  \draw[vect] (O) -- ($(Q)+(-0.08,-0.12)$) node[pos=0.55, right=1pt, etiq] {$\rp$};
  \draw[vectrojo] (Op) -- ($(Q)+(-0.13,-0.02)$) node[pos=0.5, above=2pt, etiq, text=rojoteo] {$\vc{r}\,''$};
  \draw[vectverde] (O) -- (Op) node[pos=0.3, left=6pt, etiq, text=verdeteo] {$\OOp$};
  \node[cuadro, anchor=west, font=\footnotesize, align=left] at (4.9,1.8)
     {$\vc{r}\,''=\rp-\OOp$\\[3pt] $\vp_{O'}=\vp_{O}-Q\,\OOp$};
\end{tikzpicture}
\caption{Cambio de origen. Un mismo elemento de carga tiene posición $\rp$ respecto de $O$ y $\vc{r}\,''$ respecto de $O'$. El momento dipolar cambia en $-Q\,\OOp$, de modo que
no depende del origen si la distribución es neutra (\cref{prop:p-origen}).}
\label{fig:origen}
\end{figure}


\begin{proposicion}[Cambio de origen del momento dipolar]\label{prop:p-origen}
Sean $\vp_{O}$ y $\vp_{O'}$ los momentos dipolares de una misma distribución respecto de dos orígenes $O$ y $O'$ (\cref{fig:origen}). Entonces
\begin{equation}\label{eq:p-origen}
  \boxed{\;\vp_{O'}=\vp_{O}-Q\,\OOp\;}
\end{equation}
\end{proposicion}
\begin{proof}
Si $\rp$ y $\vc{r}\,''$ son las posiciones de un mismo punto fuente respecto de $O$ y de $O'$, entonces $\vc{r}\,''=\rp-\OOp$, y
\[
  \vp_{O'}=\int_{V'}\vc{r}\,''\rho\,\dVp=\int_{V'}\bigl(\rp-\OOp\bigr)\rho\,\dVp=\int_{V'}\rp\rho\,\dVp-\OOp\int_{V'}\rho\,\dVp=\vp_{O}-Q\,\OOp .\qedhere
\]
\end{proof}

\begin{corolario}[Distribuciones neutras y centro de carga]\label{cor:centro-carga}
\begin{enumerate}[label=(\roman*)]
  \item Si la carga neta es nula, $Q=0$, el momento dipolar no depende del origen. Además, en ese caso el término dipolar es la contribución dominante al potencial.
  \item Si $Q\neq0$, siempre puede escogerse un origen respecto del cual la distribución no tenga momento dipolar: basta tomar
        \begin{equation}\label{eq:centro-carga}
          \OOp=\vc{r}_{c}=\frac{\vp_{O}}{Q},
        \end{equation}
        el \emph{centro de carga} de la distribución, análogo eléctrico del centro de masas.
\end{enumerate}
\end{corolario}
\begin{proof}
(i) Es \cref{eq:p-origen} con $Q=0$; como $\phimul{0}=0$, el primer término no nulo del desarrollo es el dipolar. (ii) Con $\OOp=\vp_{O}/Q$,
\cref{eq:p-origen} da $\vp_{O'}=\vp_{O}-Q\,\vp_{O}/Q=\vc{0}$.
\end{proof}

% -------------------------------------------------------------------------
\subsection{Término cuadrupolar}\label{sec:cuadrupolar}
% -------------------------------------------------------------------------
Con $\ell=2$ y $\Pl{2}(\cosp)=\frac12(3\cos^{2}\theta'-1)$,
\begin{equation}\label{eq:cuadrupolar-integral}
  \phimul{2}(\vc{r})=\frac{1}{4\pi\eo\,r^{3}}\int_{V'}\rho(\rp)\,r'^{2}\,\Pl{2}(\cosp)\,\dVp=\frac{1}{8\pi\eo\,r^{3}}\int_{V'}\rho(\rp)\,r'^{2}\bigl(3\cos^{2}\theta'-1\bigr)\dVp .
\end{equation}
Como $\cosp$ depende de $\vc{r}$ y de $\rp$, tratamos de nuevo de separar lo que depende del punto campo de lo que depende del punto fuente.

\begin{lema}[Separación de variables en el término cuadrupolar]\label{lem:separacion}
Con los cosenos directores $\cd{j}$ del punto campo (\cref{not:puntos}),
\begin{equation}\label{eq:separacion}
  r'^{2}\bigl(3\cos^{2}\theta'-1\bigr)=\sum_{j,k}T_{jk}\,\cd{j}\cd{k},\qquad T_{jk}=3x'_jx'_k-\delt{j}{k}\,r'^{2},\qquad j,k\in\{x,y,z\}.
\end{equation}
\end{lema}
\begin{proof}
Como $r'\cosp=\ur\cdot\rp=\cd{x}x'+\cd{y}y'+\cd{z}z'$ y $\cd{x}^{2}+\cd{y}^{2}+\cd{z}^{2}=1$,
\[
  r'^{2}\bigl(3\cos^{2}\theta'-1\bigr)=3\bigl(\ur\cdot\rp\bigr)^{2}-r'^{2}=3\bigl(\cd{x}x'+\cd{y}y'+\cd{z}z'\bigr)^{2}-r'^{2}\bigl(\cd{x}^{2}+\cd{y}^{2}+\cd{z}^{2}\bigr).
\]
Desarrollando el cuadrado con $(a+b+c)^{2}=a^{2}+b^{2}+c^{2}+2ab+2ac+2bc$,
\[
  =3\bigl(\cd{x}^{2}x'^{2}+\cd{y}^{2}y'^{2}+\cd{z}^{2}z'^{2}\bigr)+6\bigl(\cd{x}\cd{y}\,x'y'+\cd{x}\cd{z}\,x'z'+\cd{y}\cd{z}\,y'z'\bigr)-r'^{2}\bigl(\cd{x}^{2}+\cd{y}^{2}+\cd{z}^{2}\bigr),
\]
y repartiendo cada término cruzado $6\,\cd{j}\cd{k}\,x'_jx'_k$ ($j\neq k$) a partes iguales entre las posiciones $(j,k)$ y $(k,j)$, se obtiene la forma cuadrática
\[
  \begin{pmatrix}\cd{x}&\cd{y}&\cd{z}\end{pmatrix}
  \begin{pmatrix}3x'^{2}-r'^{2}&3x'y'&3x'z'\\ 3x'y'&3y'^{2}-r'^{2}&3y'z'\\ 3x'z'&3y'z'&3z'^{2}-r'^{2}\end{pmatrix}
  \begin{pmatrix}\cd{x}\\ \cd{y}\\ \cd{z}\end{pmatrix}=\sum_{j,k}T_{jk}\,\cd{j}\cd{k}.\qedhere
\]
\end{proof}

La matriz $T_{jk}$ solo depende del punto fuente y los cosenos directores solo del punto campo: integrando en $V'$ conseguimos la separación buscada.

\begin{definicion}[Tensor momento cuadrupolar]\label{def:cuadrupolo}
Se llama \emph{tensor momento cuadrupolar} (o, simplemente, momento cuadrupolar) de una distribución de carga, respecto del origen $O$, a la matriz
$\Qt=(\Qjk)$ de componentes
\begin{equation}\label{eq:Qjk}
  \boxed{\;\Qjk=\int_{V'}\rho(\rp)\bigl(3x'_jx'_k-\delt{j}{k}\,r'^{2}\bigr)\dVp\;},\qquad j,k\in\{x,y,z\}.
\end{equation}
Para un sistema discreto de cargas puntuales $q_i$ en $\rp_i=(x'_i,y'_i,z'_i)$,
\begin{equation}\label{eq:Qjk-discreto}
  \Qjk=\sum_{i}q_i\bigl(3x'_{i,j}x'_{i,k}-\delt{j}{k}\,r_i'^{\,2}\bigr).
\end{equation}
\end{definicion}

\begin{proposicion}[Término cuadrupolar]\label{prop:cuadrupolar}
\begin{equation}\label{eq:cuadrupolar}
  \boxed{\;\phimul{2}(\vc{r})=\frac{1}{8\pi\eo\,r^{3}}\sum_{j,k}\Qjk\,\cd{j}\cd{k}=\frac{1}{8\pi\eo\,r^{5}}\sum_{j,k}\Qjk\,x_jx_k\;}
\end{equation}
\end{proposicion}
\begin{proof}
Se integra \cref{eq:separacion} multiplicada por $\rho$ y se sustituye en \cref{eq:cuadrupolar-integral}; los $\cd{j}$ salen de la integral porque no dependen de
$\rp$. La segunda forma se sigue de $\cd{j}=x_j/r$.
\end{proof}

\begin{observacion}[Analogía con el término dipolar]\label{obs:analogia-p-Q}
El término dipolar puede escribirse igual, $\vp\cdot\ur=\sum_jp_j\,\cd{j}$: es una forma \emph{lineal} en los cosenos directores, mientras que el cuadrupolar es una
forma \emph{cuadrática}. En general, el término $\ell$-ésimo es una forma de grado $\ell$ en los $\cd{j}$, cuyos coeficientes son los momentos de orden $\ell$ de la
distribución. El término cuadrupolar decae como $1/r^{3}$ y es el dominante cuando $Q=0$ y $\vp=\vc{0}$; su valor coincide con el potencial de un
\emph{cuadrupolo} (por ejemplo, dos dipolos de sentidos contrarios, como en la \cref{fig:cuadrupolo-plano}) situado en el origen.
\end{observacion}

\subsubsection*{Propiedades del tensor cuadrupolar}
\begin{proposicion}[Componentes independientes]\label{prop:Q-propiedades}
\begin{enumerate}[label=(\roman*)]
  \item $\Qt$ es \emph{simétrico}, $\Qjk=Q_{kj}$: de sus nueve componentes, solo seis son independientes.
  \item $\Qt$ tiene \emph{traza nula}, $\Tr\Qt=Q_{xx}+Q_{yy}+Q_{zz}=0$: esta ligadura elimina un grado de libertad más y deja cinco componentes independientes.
  \item Existe un sistema de ejes cartesianos, los \emph{ejes principales}, en el que $\Qt$ es diagonal, $\Qt=\operatorname{diag}(Q_1,Q_2,Q_3)$ con
        $Q_1+Q_2+Q_3=0$, y en él
        \begin{equation}\label{eq:Q-principal}
          \phimul{2}(\vc{r})=\frac{1}{8\pi\eo\,r^{3}}\bigl(Q_1\cd{1}^{2}+Q_2\cd{2}^{2}+Q_3\cd{3}^{2}\bigr).
        \end{equation}
\end{enumerate}
\end{proposicion}
\begin{proof}
(i) El integrando $3x'_jx'_k-\delt{j}{k}r'^{2}$ es simétrico en $j,k$. (ii) Sumando las componentes diagonales,
\[
  \Tr\Qt=\sum_{j}Q_{jj}=\int_{V'}\rho\bigl(3x'^{2}-r'^{2}+3y'^{2}-r'^{2}+3z'^{2}-r'^{2}\bigr)\dVp=\int_{V'}\rho\bigl(3r'^{2}-3r'^{2}\bigr)\dVp=0 .
\]
(iii) Toda matriz real simétrica es diagonalizable mediante una matriz ortogonal (teorema espectral): sus autovectores, ortonormales, definen los ejes
principales, y una rotación conserva la traza.
\end{proof}

\begin{observacion}[Cinco números: dos autovalores y tres ángulos]\label{obs:cinco}
El recuento de la \cref{prop:Q-propiedades} tiene una lectura geométrica: dos de los autovalores (el tercero lo fija la traza nula) y los tres ángulos que orientan los
ejes principales. En general, el momento de orden $\ell$ tiene $2\ell+1$ componentes independientes: $1$ para la carga, $3$ para el dipolo y $5$ para el
cuadrupolo.
\end{observacion}

\begin{observacion}[Relación con el tensor de inercia]\label{obs:inercia}
La estructura de $\Qt$ recuerda a la del tensor de inercia de la mecánica, $I_{jk}=\int\rho_m\bigl(r'^{2}\delt{j}{k}-x'_jx'_k\bigr)\dVp$: ambos son formas cuadráticas
construidas con los segundos momentos de la distribución. De hecho, si se sustituye la densidad de masa $\rho_m$ por la de carga $\rho$, se tiene
$\Qjk=\delt{j}{k}\Tr I-3I_{jk}$, de modo que los ejes principales de inercia y los del cuadrupolo coinciden.
\end{observacion}

\begin{proposicion}[Cambio de origen del tensor cuadrupolar]\label{prop:Q-origen}
Si $\vc{d}=\OOp$ y $Q$, $\vp$ y $\Qjk$ se refieren al origen $O$, el tensor cuadrupolar respecto de $O'$ es
\begin{equation}\label{eq:Q-origen}
  Q'_{jk}=\Qjk-3\bigl(p_jd_k+d_jp_k\bigr)+2\,\delt{j}{k}\,(\vp\cdot\vc{d})+Q\bigl(3d_jd_k-\delt{j}{k}d^{2}\bigr).
\end{equation}
En particular, si $Q=0$ y $\vp=\vc{0}$ el tensor cuadrupolar no depende del origen.
\end{proposicion}
\begin{proof}
Con $\vc{r}\,''=\rp-\vc{d}$, el integrando de $Q'_{jk}$ es
$3(x'_j-d_j)(x'_k-d_k)-\delt{j}{k}\norm{\rp-\vc{d}}^{2}=\bigl(3x'_jx'_k-\delt{j}{k}r'^{2}\bigr)-3\bigl(x'_jd_k+d_jx'_k\bigr)+2\,\delt{j}{k}\,\rp\cdot\vc{d}+3d_jd_k-\delt{j}{k}d^{2}$.
Al multiplicar por $\rho$ e integrar, $\int\rho\,x'_j=p_j$ y $\int\rho=Q$.
\end{proof}

Es un ejemplo de una regla general: \emph{el primer momento multipolar no nulo de una distribución no depende del origen} (la carga si $Q\neq0$; el dipolo si $Q=0$;
el cuadrupolo si además $\vp=\vc{0}$\dots).

\begin{observacion}[Interpretación geométrica: forma de la distribución]\label{obs:forma}
Para una distribución positiva con simetría de revolución alrededor de $z$ (\cref{sec:revolucion}), $Q_{zz}=\int\rho\,(2z'^{2}-x'^{2}-y'^{2})\dVp$ compara la extensión
de la carga a lo largo del eje con la transversal (\cref{fig:formas}):
\begin{itemize}
  \item $Q_{zz}>0$: la carga está \emph{alargada} según $z$ (forma de «cigarro», elipsoide prolato);
  \item $Q_{zz}<0$: la carga está \emph{achatada} (forma de «disco», elipsoide oblato);
  \item $Q_{zz}=0$: no hay deformación cuadrupolar (por ejemplo, simetría esférica).
\end{itemize}
Las componentes fuera de la diagonal indican que la distribución está \emph{girada} respecto de los ejes: si $Q_{xy}\neq0$, por ejemplo, la dirección de alargamiento está
orientada oblicuamente en el plano $xy$. Así se describen, por ejemplo, las deformaciones de los núcleos atómicos.
\end{observacion}

% =========================================================================
% Figura (apuntes manuscritos, pág. 12): interpretación del signo de Q_zz y
% de las componentes fuera de la diagonal. Elipses sombreadas en la paleta
% del documento (sección meridiana de la distribución de carga positiva).
% =========================================================================
\begin{figure}[ht]
\centering
\newcommand{\forma}[4]{% semieje x, semieje z, giro, rótulo
  \begin{tikzpicture}[x=1cm, y=1cm, >=Stealth]
    \draw[eje] (0,-1.55) -- (0,1.65) node[above, etiq, font=\footnotesize] {$z$};
    \draw[eje] (-1.55,0) -- (1.65,0) node[right, etiq, font=\footnotesize] {$x$};
    \shade[rotate=#3, inner color=magacero!55, outer color=rosaclaro!60, opacity=0.9] (0,0) ellipse ({#1} and {#2});
    \draw[rotate=#3, magacero, line width=0.7pt] (0,0) ellipse ({#1} and {#2});
    \draw[rotate=#3, densely dashed, magprof, line width=0.4pt] (0,{-#2-0.15}) -- (0,{#2+0.15});
    \node[font=\footnotesize, align=center] at (0,-2.05) {#4};
  \end{tikzpicture}}
\forma{0.55}{1.3}{0}{$Q_{zz}>0$: alargada\\(«cigarro»)}\hfill
\forma{1.3}{0.55}{0}{$Q_{zz}<0$: achatada\\(«disco»)}\hfill
\forma{0.95}{0.95}{0}{$Q_{zz}=0$: sin deformación\\(simetría esférica)}\hfill
\forma{0.55}{1.3}{-35}{$Q_{xz}\neq0$: girada\\(ejes principales oblicuos)}
\caption{Lectura geométrica del tensor cuadrupolar de una distribución de carga positiva. Con simetría de revolución alrededor de $z$, el signo de $Q_{zz}$ distingue una
carga alargada a lo largo del eje de una achatada; las componentes fuera de la diagonal no nulas indican que los ejes principales (a trazos) están girados respecto de los
ejes coordenados.}
\label{fig:formas}
\end{figure}


% -------------------------------------------------------------------------
\subsection{Distribuciones con simetría de revolución}\label{sec:revolucion}
% -------------------------------------------------------------------------
Cuando la distribución es invariante frente a rotaciones alrededor de un eje, el número de componentes independientes se reduce drásticamente.

\begin{teorema}[Momentos de una distribución con simetría de revolución]\label{teo:revolucion}
Sea $\rho(\rp)=\rho(r',\theta')$ independiente del ángulo azimutal $\varphi'$ (simetría de revolución alrededor del eje $z$). Entonces:
\begin{enumerate}[label=(\roman*)]
  \item el momento dipolar es paralelo al eje: $p_x=p_y=0$ y $\vp=p_z\,\uz$;
  \item el tensor cuadrupolar es diagonal y queda determinado por una sola componente:
        \begin{equation}\label{eq:Q-revolucion}
          Q_{jk}=0\ (j\neq k),\qquad Q_{xx}=Q_{yy}=-\tfrac12\,Q_{zz},\qquad Q_{zz}=\int_{V'}\rho\,\bigl(3z'^{2}-r'^{2}\bigr)\dVp ;
        \end{equation}
  \item el término cuadrupolar es
        \begin{equation}\label{eq:cuadrupolar-revolucion}
          \boxed{\;\phimul{2}(\vc{r})=\frac{Q_{zz}}{16\pi\eo\,r^{3}}\bigl(3\cos^{2}\theta-1\bigr)=\frac{Q_{zz}}{8\pi\eo\,r^{3}}\,\Pl{2}(\cos\theta)\;}
        \end{equation}
\end{enumerate}
\end{teorema}
\begin{proof}
Usamos coordenadas esféricas para el punto fuente, $x'=r'\sen\theta'\cos\varphi'$, $y'=r'\sen\theta'\sen\varphi'$, $z'=r'\cos\theta'$ y
$\dVp=r'^{2}\sen\theta'\,\dif r'\,\dif\theta'\,\dif\varphi'$. Como $\rho$ no depende de $\varphi'$, en cada integral se puede separar la integral en $\varphi'$ (teorema de Fubini),
y las únicas que aparecen son
\[
  \int_{0}^{2\pi}\cos\varphi'\,\dif\varphi'=\bigl[\sen\varphi'\bigr]_{0}^{2\pi}=0,\qquad
  \int_{0}^{2\pi}\sen\varphi'\,\dif\varphi'=\bigl[-\cos\varphi'\bigr]_{0}^{2\pi}=-1-(-1)=0,
\]
\[
  \int_{0}^{2\pi}\sen\varphi'\cos\varphi'\,\dif\varphi'=\frac12\int_{0}^{2\pi}\frac{\dif}{\dif\varphi'}\bigl(\sen^{2}\varphi'\bigr)\dif\varphi'=\frac12\bigl[\sen^{2}\varphi'\bigr]_{0}^{2\pi}=0,
\]
\[
  \int_{0}^{2\pi}\cos^{2}\varphi'\,\dif\varphi'=\int_{0}^{2\pi}\sen^{2}\varphi'\,\dif\varphi'=\pi .
\]
(i) $p_x=\int\rho\,x'\,\dVp\propto\int_0^{2\pi}\cos\varphi'\dif\varphi'=0$ y, análogamente, $p_y\propto\int_0^{2\pi}\sen\varphi'\dif\varphi'=0$.

(ii) Fuera de la diagonal, $\delt{j}{k}=0$ y $Q_{jk}=3\int\rho\,x'_jx'_k\,\dVp$. Como $x'z'=r'^{2}\sen\theta'\cos\theta'\cos\varphi'$, $y'z'=r'^{2}\sen\theta'\cos\theta'\sen\varphi'$ y
$x'y'=r'^{2}\sen^{2}\theta'\cos\varphi'\sen\varphi'$, se tiene $Q_{xz}\propto\int\cos\varphi'=0$, $Q_{yz}\propto\int\sen\varphi'=0$ y $Q_{xy}\propto\int\sen\varphi'\cos\varphi'=0$. En la
diagonal, $Q_{xx}$ y $Q_{yy}$ difieren solo en que llevan $\cos^{2}\varphi'$ o $\sen^{2}\varphi'$, cuyas integrales coinciden: $Q_{xx}=Q_{yy}$. Con la traza nula,
$2Q_{xx}+Q_{zz}=0$, de donde $Q_{xx}=Q_{yy}=-\frac12Q_{zz}$.

(iii) Sustituyendo en \cref{eq:cuadrupolar} y usando $\cd{x}^{2}+\cd{y}^{2}=1-\cd{z}^{2}$ y $\cd{z}=\cos\theta$,
\begin{align*}
  \phimul{2}&=\frac{1}{8\pi\eo r^{3}}\Bigl[Q_{xx}\bigl(\cd{x}^{2}+\cd{y}^{2}\bigr)+Q_{zz}\cd{z}^{2}\Bigr]
  =\frac{1}{8\pi\eo r^{3}}\Bigl[-\frac12Q_{zz}\bigl(1-\cd{z}^{2}\bigr)+Q_{zz}\cd{z}^{2}\Bigr]\\
  &=\frac{Q_{zz}}{8\pi\eo r^{3}}\Bigl(\frac32\cos^{2}\theta-\frac12\Bigr).\qedhere
\end{align*}
\end{proof}

\begin{observacion}[Por qué se anulan: argumento geométrico]\label{obs:anulacion}
Las integrales en $\varphi'$ se anulan por una razón de signos: $\sen\varphi'\cos\varphi'>0$ en el primer y el tercer cuadrante y $<0$ en el segundo y el cuarto, con
contribuciones que se compensan exactamente al dar la vuelta completa (también se ve escribiendo $\sen\varphi'\cos\varphi'=\frac12\sen2\varphi'$, una sinusoide completa
de integral nula en un periodo). Físicamente: para cada punto $(x',y')$ hay otro, girado, con la misma densidad y el signo opuesto de $x'$ o de $y'$; la integración en
$\varphi'$ «mata» todos los términos que contienen $\sen\varphi'$, $\cos\varphi'$ o su producto. Y los ejes $x$ e $y$ son físicamente equivalentes, de ahí $Q_{xx}=Q_{yy}$. Se
pasa así de $9$ componentes a $6$ (simetría), a $5$ (traza nula) y, con simetría de revolución, a \emph{una sola} (dos ligaduras más).
\end{observacion}

Con simetría de revolución, por tanto, los dos primeros términos no monopolares tienen la misma estructura,
\begin{equation}\label{eq:dip-cuad-legendre}
  \phimul{1}=\frac{p\cos\theta}{4\pi\eo r^{2}}=\frac{p}{4\pi\eo r^{2}}\,\Pl{1}(\cos\theta),\qquad
  \phimul{2}=\frac{Q_{zz}}{8\pi\eo r^{3}}\,\Pl{2}(\cos\theta),
\end{equation}
donde ahora $\theta$ es la colatitud del punto campo respecto del eje de simetría.

\begin{corolario}[Campo, equipotenciales y líneas del término cuadrupolar]\label{cor:campo-cuadrupolo}
Para el término \cref{eq:cuadrupolar-revolucion}:
\begin{enumerate}[label=(\roman*)]
  \item el campo es $\vE_{2}=\dfrac{3Q_{zz}}{16\pi\eo r^{4}}\Bigl[\bigl(3\cos^{2}\theta-1\bigr)\ur+2\sen\theta\cos\theta\,\uth\Bigr]$, que decae como $1/r^{4}$;
  \item las equipotenciales $\phimul{2}=\phi_0$ son las curvas $r^{3}=K\,(3\cos^{2}\theta-1)$, y la equipotencial nula es el \emph{cono nodal} $\cos^{2}\theta=1/3$,
        es decir, $\theta\simeq54{,}7^{\circ}$ y $125{,}3^{\circ}$;
  \item las líneas de campo son las curvas $r^{2}=C\sen^{2}\theta\,\abs{\cos\theta}$, más el eje $z$ y el plano ecuatorial.
\end{enumerate}
\end{corolario}
\begin{proof}
(i) $E_r=-\partial_r\phimul{2}=3\phimul{2}/r$ y $E_\theta=-\frac1r\partial_\theta\phimul{2}=-\frac{Q_{zz}}{16\pi\eo r^{4}}(-6\cos\theta\sen\theta)$.
(ii) Es inmediato. (iii) Como en la \cref{prop:lineas-dipolo},
$\frac{\dif r}{r\,\dif\theta}=\frac{E_r}{E_\theta}=\frac{3\cos^{2}\theta-1}{2\sen\theta\cos\theta}$, y derivando $\frac12\ln\bigl(\sen^{2}\theta\cos\theta\bigr)$ se comprueba que
$\frac{\dif}{\dif\theta}\,\frac12\ln\bigl(\sen^{2}\theta\cos\theta\bigr)=\frac{2\cos^{2}\theta-\sen^{2}\theta}{2\sen\theta\cos\theta}=\frac{3\cos^{2}\theta-1}{2\sen\theta\cos\theta}$.
\end{proof}

% =========================================================================
% Figura (clon de «Tema 2. Figuras», diapositiva «Equipotenciales»):
% equipotenciales de los términos monopolar (phi ~ 1/r), dipolar
% (phi ~ cos t / r^2) y cuadrupolar (phi ~ P2(cos t)/r^3), a intervalos
% iguales de potencial. Monopolo: circunferencias analíticas; dipolo y
% cuadrupolo: datos de figuras/datos/{dip,cua}_puntual_eq.tex.
% =========================================================================
\begin{figure}[ht]
\centering
\newcommand{\rotulo}[1]{{\footnotesize #1}\\[3pt]}
\renewcommand{\EQ}[1]{\draw[mmaEQsolo] plot coordinates {#1};}
\colorlet{marcoc}{grispizarra!60}
\begin{minipage}[t]{0.325\textwidth}\centering
\rotulo{Término monopolar ($\ell=0$)}
\begin{tikzpicture}[x=0.198cm, y=0.198cm]
  \panelMMA{10}{white}{0}{%
    \foreach \k in {0,...,12}{\pgfmathsetmacro{\rr}{1/(0.1+0.0194*\k)}\draw[mmaEQsolo] (0,0) circle (\rr);}}
  \marcoMMA{10}{5}{$x$ (u.\,arb.)}{$z$ (u.\,arb.)}
\end{tikzpicture}
\end{minipage}\hfill
\begin{minipage}[t]{0.325\textwidth}\centering
\rotulo{Término dipolar ($\ell=1$)}
\begin{tikzpicture}[x=0.198cm, y=0.198cm]
  \panelMMA{10}{white}{0}{%
    % Dipolo puntual - equipotenciales
% Generado por scripts/generar_datos.py (no editar a mano).
\EQ{(1.224,-0.476) (1.250,-0.522) (1.275,-0.574) (1.296,-0.625) (1.313,-0.675) (1.326,-0.725) (1.337,-0.775) (1.345,-0.825) (1.350,-0.875) (1.353,-0.925) (1.354,-0.975) (1.352,-1.025) (1.347,-1.075) (1.340,-1.125) (1.331,-1.175) (1.319,-1.225) (1.305,-1.275) (1.288,-1.325) (1.269,-1.375) (1.246,-1.425) (1.221,-1.475) (1.194,-1.525) (1.162,-1.575) (1.128,-1.625) (1.089,-1.675) (1.050,-1.721) (1.000,-1.774) (0.950,-1.822) (0.900,-1.865) (0.856,-1.900) (0.800,-1.940) (0.750,-1.972) (0.702,-2.000) (0.656,-2.025) (0.604,-2.050) (0.550,-2.074) (0.500,-2.093) (0.450,-2.111) (0.400,-2.126) (0.350,-2.139) (0.300,-2.151) (0.250,-2.161) (0.200,-2.168) (0.150,-2.174) (0.100,-2.179) (0.050,-2.181) (0.000,-2.182) (-0.050,-2.181) (-0.100,-2.179) (-0.150,-2.174) (-0.200,-2.168) (-0.250,-2.161) (-0.300,-2.151) (-0.350,-2.139) (-0.400,-2.126) (-0.450,-2.111) (-0.500,-2.093) (-0.547,-2.075) (-0.600,-2.052) (-0.650,-2.028) (-0.700,-2.001) (-0.745,-1.975) (-0.800,-1.940) (-0.850,-1.904) (-0.900,-1.865) (-0.947,-1.825) (-0.999,-1.775) (-1.047,-1.725) (-1.089,-1.675) (-1.125,-1.629) (-1.162,-1.575) (-1.194,-1.525) (-1.221,-1.475) (-1.246,-1.425) (-1.269,-1.375) (-1.288,-1.325) (-1.305,-1.275) (-1.319,-1.225) (-1.331,-1.175) (-1.340,-1.125) (-1.347,-1.075) (-1.352,-1.025) (-1.354,-0.975) (-1.353,-0.925) (-1.350,-0.875) (-1.345,-0.825) (-1.337,-0.775) (-1.326,-0.725) (-1.313,-0.675) (-1.296,-0.625) (-1.275,-0.575) (-1.252,-0.525) (-1.225,-0.477)}
\EQ{(1.229,-0.446) (1.263,-0.500) (1.290,-0.550) (1.313,-0.600) (1.333,-0.650) (1.349,-0.700) (1.362,-0.750) (1.373,-0.800) (1.381,-0.850) (1.386,-0.900) (1.389,-0.950) (1.389,-1.000) (1.387,-1.050) (1.383,-1.100) (1.376,-1.150) (1.367,-1.200) (1.355,-1.250) (1.341,-1.300) (1.325,-1.349) (1.306,-1.400) (1.284,-1.450) (1.260,-1.500) (1.233,-1.550) (1.203,-1.600) (1.175,-1.642) (1.132,-1.700) (1.100,-1.740) (1.050,-1.796) (1.000,-1.846) (0.950,-1.891) (0.909,-1.925) (0.850,-1.970) (0.806,-2.000) (0.750,-2.035) (0.700,-2.063) (0.650,-2.089) (0.600,-2.112) (0.550,-2.133) (0.500,-2.152) (0.450,-2.169) (0.400,-2.184) (0.350,-2.197) (0.300,-2.208) (0.250,-2.218) (0.200,-2.225) (0.150,-2.231) (0.100,-2.236) (0.050,-2.238) (0.000,-2.239) (-0.050,-2.238) (-0.100,-2.236) (-0.150,-2.231) (-0.200,-2.225) (-0.250,-2.218) (-0.300,-2.208) (-0.350,-2.197) (-0.400,-2.184) (-0.450,-2.169) (-0.500,-2.152) (-0.550,-2.133) (-0.600,-2.112) (-0.650,-2.089) (-0.700,-2.063) (-0.750,-2.035) (-0.800,-2.004) (-0.843,-1.975) (-0.900,-1.932) (-0.940,-1.900) (-0.996,-1.850) (-1.046,-1.800) (-1.091,-1.750) (-1.125,-1.709) (-1.169,-1.650) (-1.200,-1.604) (-1.233,-1.550) (-1.260,-1.500) (-1.284,-1.450) (-1.306,-1.400) (-1.325,-1.350) (-1.341,-1.300) (-1.355,-1.250) (-1.367,-1.200) (-1.376,-1.150) (-1.383,-1.100) (-1.387,-1.050) (-1.389,-1.000) (-1.389,-0.950) (-1.386,-0.900) (-1.381,-0.850) (-1.373,-0.800) (-1.362,-0.750) (-1.349,-0.700) (-1.333,-0.650) (-1.313,-0.600) (-1.290,-0.550) (-1.263,-0.500) (-1.232,-0.450) (-1.229,-0.446)}
\EQ{(1.244,-0.431) (1.274,-0.475) (1.300,-0.518) (1.325,-0.565) (1.350,-0.619) (1.372,-0.675) (1.388,-0.725) (1.400,-0.775) (1.411,-0.825) (1.419,-0.875) (1.424,-0.925) (1.427,-0.975) (1.427,-1.025) (1.425,-1.075) (1.421,-1.125) (1.414,-1.175) (1.406,-1.225) (1.395,-1.275) (1.381,-1.325) (1.365,-1.375) (1.347,-1.425) (1.326,-1.475) (1.303,-1.525) (1.277,-1.575) (1.250,-1.621) (1.215,-1.675) (1.180,-1.725) (1.141,-1.775) (1.100,-1.822) (1.050,-1.874) (1.000,-1.922) (0.950,-1.965) (0.906,-2.000) (0.850,-2.040) (0.800,-2.073) (0.755,-2.100) (0.709,-2.125) (0.660,-2.150) (0.605,-2.175) (0.550,-2.198) (0.500,-2.216) (0.450,-2.233) (0.400,-2.247) (0.350,-2.260) (0.300,-2.271) (0.250,-2.280) (0.200,-2.287) (0.150,-2.293) (0.100,-2.297) (0.050,-2.299) (0.000,-2.300) (-0.050,-2.299) (-0.100,-2.297) (-0.150,-2.293) (-0.200,-2.287) (-0.250,-2.280) (-0.300,-2.271) (-0.350,-2.260) (-0.400,-2.247) (-0.450,-2.233) (-0.500,-2.216) (-0.550,-2.198) (-0.600,-2.177) (-0.650,-2.155) (-0.700,-2.130) (-0.750,-2.103) (-0.796,-2.075) (-0.850,-2.040) (-0.900,-2.004) (-0.950,-1.965) (-0.997,-1.925) (-1.049,-1.875) (-1.097,-1.825) (-1.141,-1.775) (-1.175,-1.731) (-1.215,-1.675) (-1.248,-1.625) (-1.275,-1.578) (-1.300,-1.531) (-1.325,-1.478) (-1.347,-1.425) (-1.365,-1.375) (-1.381,-1.325) (-1.395,-1.275) (-1.406,-1.225) (-1.414,-1.175) (-1.421,-1.125) (-1.425,-1.075) (-1.427,-1.025) (-1.427,-0.975) (-1.424,-0.925) (-1.419,-0.875) (-1.411,-0.825) (-1.400,-0.775) (-1.388,-0.725) (-1.372,-0.675) (-1.352,-0.625) (-1.330,-0.575) (-1.304,-0.525) (-1.275,-0.477) (-1.244,-0.431)}
\EQ{(1.250,-0.405) (1.285,-0.450) (1.318,-0.500) (1.347,-0.550) (1.373,-0.600) (1.395,-0.650) (1.413,-0.700) (1.429,-0.750) (1.442,-0.800) (1.452,-0.850) (1.460,-0.900) (1.465,-0.950) (1.468,-1.000) (1.468,-1.050) (1.467,-1.100) (1.463,-1.150) (1.457,-1.200) (1.448,-1.250) (1.438,-1.300) (1.425,-1.350) (1.410,-1.400) (1.392,-1.450) (1.372,-1.500) (1.350,-1.550) (1.325,-1.600) (1.300,-1.645) (1.266,-1.700) (1.232,-1.750) (1.200,-1.793) (1.154,-1.850) (1.109,-1.900) (1.059,-1.950) (1.004,-2.000) (0.950,-2.044) (0.909,-2.075) (0.850,-2.116) (0.800,-2.147) (0.751,-2.175) (0.704,-2.200) (0.652,-2.225) (0.600,-2.248) (0.550,-2.268) (0.500,-2.285) (0.450,-2.301) (0.400,-2.315) (0.350,-2.328) (0.300,-2.338) (0.250,-2.347) (0.200,-2.354) (0.150,-2.360) (0.100,-2.364) (0.050,-2.366) (0.000,-2.367) (-0.050,-2.366) (-0.100,-2.364) (-0.150,-2.360) (-0.200,-2.354) (-0.250,-2.347) (-0.300,-2.338) (-0.350,-2.328) (-0.400,-2.315) (-0.450,-2.301) (-0.500,-2.285) (-0.550,-2.268) (-0.600,-2.248) (-0.650,-2.226) (-0.700,-2.202) (-0.750,-2.176) (-0.795,-2.150) (-0.850,-2.116) (-0.900,-2.081) (-0.942,-2.050) (-1.000,-2.003) (-1.050,-1.958) (-1.100,-1.909) (-1.150,-1.854) (-1.195,-1.800) (-1.225,-1.760) (-1.266,-1.700) (-1.297,-1.650) (-1.325,-1.600) (-1.350,-1.550) (-1.372,-1.500) (-1.392,-1.450) (-1.410,-1.400) (-1.425,-1.350) (-1.438,-1.300) (-1.448,-1.250) (-1.457,-1.200) (-1.463,-1.150) (-1.467,-1.100) (-1.468,-1.050) (-1.468,-1.000) (-1.465,-0.950) (-1.460,-0.900) (-1.452,-0.850) (-1.442,-0.800) (-1.429,-0.750) (-1.413,-0.700) (-1.395,-0.650) (-1.373,-0.600) (-1.350,-0.555) (-1.325,-0.512) (-1.285,-0.450) (-1.250,-0.405)}
\EQ{(1.251,-0.374) (1.295,-0.425) (1.325,-0.465) (1.365,-0.525) (1.393,-0.575) (1.418,-0.625) (1.440,-0.675) (1.458,-0.725) (1.473,-0.775) (1.486,-0.825) (1.496,-0.875) (1.504,-0.925) (1.510,-0.975) (1.513,-1.025) (1.514,-1.075) (1.512,-1.125) (1.509,-1.175) (1.503,-1.225) (1.495,-1.275) (1.486,-1.325) (1.473,-1.375) (1.459,-1.425) (1.442,-1.475) (1.423,-1.525) (1.402,-1.575) (1.378,-1.625) (1.352,-1.675) (1.325,-1.721) (1.290,-1.775) (1.255,-1.825) (1.216,-1.875) (1.175,-1.923) (1.127,-1.975) (1.075,-2.025) (1.025,-2.070) (0.975,-2.110) (0.925,-2.147) (0.875,-2.182) (0.825,-2.213) (0.775,-2.242) (0.725,-2.268) (0.675,-2.292) (0.625,-2.314) (0.575,-2.334) (0.525,-2.352) (0.475,-2.369) (0.425,-2.383) (0.375,-2.396) (0.325,-2.407) (0.275,-2.416) (0.225,-2.424) (0.175,-2.430) (0.125,-2.435) (0.075,-2.438) (0.025,-2.440) (-0.050,-2.439) (-0.100,-2.437) (-0.150,-2.433) (-0.200,-2.427) (-0.250,-2.420) (-0.300,-2.412) (-0.350,-2.402) (-0.400,-2.390) (-0.450,-2.376) (-0.500,-2.361) (-0.550,-2.344) (-0.599,-2.325) (-0.650,-2.304) (-0.700,-2.280) (-0.750,-2.255) (-0.800,-2.228) (-0.846,-2.200) (-0.900,-2.165) (-0.950,-2.129) (-1.000,-2.090) (-1.048,-2.050) (-1.100,-2.002) (-1.150,-1.951) (-1.195,-1.900) (-1.236,-1.850) (-1.273,-1.800) (-1.307,-1.750) (-1.337,-1.700) (-1.365,-1.650) (-1.390,-1.600) (-1.413,-1.550) (-1.433,-1.500) (-1.450,-1.453) (-1.466,-1.400) (-1.480,-1.350) (-1.491,-1.300) (-1.500,-1.250) (-1.506,-1.200) (-1.511,-1.150) (-1.513,-1.100) (-1.514,-1.050) (-1.512,-1.000) (-1.507,-0.950) (-1.501,-0.900) (-1.492,-0.850) (-1.480,-0.800) (-1.466,-0.750) (-1.449,-0.700) (-1.429,-0.650) (-1.406,-0.600) (-1.379,-0.550) (-1.350,-0.502) (-1.314,-0.450) (-1.275,-0.401) (-1.251,-0.374)}
\EQ{(1.266,-0.359) (1.305,-0.400) (1.346,-0.450) (1.382,-0.500) (1.414,-0.550) (1.442,-0.600) (1.467,-0.650) (1.488,-0.700) (1.506,-0.750) (1.522,-0.800) (1.534,-0.850) (1.545,-0.900) (1.553,-0.950) (1.559,-1.000) (1.562,-1.050) (1.563,-1.100) (1.563,-1.150) (1.560,-1.200) (1.555,-1.250) (1.547,-1.300) (1.538,-1.350) (1.527,-1.400) (1.513,-1.450) (1.498,-1.500) (1.480,-1.550) (1.460,-1.600) (1.437,-1.650) (1.412,-1.700) (1.385,-1.750) (1.354,-1.800) (1.325,-1.844) (1.285,-1.900) (1.250,-1.944) (1.201,-2.000) (1.154,-2.050) (1.101,-2.100) (1.050,-2.145) (1.000,-2.185) (0.950,-2.222) (0.908,-2.250) (0.850,-2.287) (0.800,-2.315) (0.750,-2.342) (0.700,-2.366) (0.650,-2.388) (0.600,-2.409) (0.550,-2.427) (0.500,-2.443) (0.450,-2.458) (0.400,-2.471) (0.350,-2.483) (0.300,-2.493) (0.250,-2.501) (0.200,-2.508) (0.150,-2.513) (0.100,-2.517) (0.050,-2.519) (0.000,-2.520) (-0.050,-2.519) (-0.100,-2.517) (-0.150,-2.513) (-0.200,-2.508) (-0.250,-2.501) (-0.300,-2.493) (-0.350,-2.483) (-0.400,-2.471) (-0.450,-2.458) (-0.500,-2.443) (-0.550,-2.427) (-0.600,-2.409) (-0.650,-2.388) (-0.700,-2.366) (-0.750,-2.342) (-0.800,-2.315) (-0.850,-2.287) (-0.900,-2.256) (-0.945,-2.225) (-1.000,-2.185) (-1.044,-2.150) (-1.100,-2.101) (-1.150,-2.054) (-1.200,-2.001) (-1.245,-1.950) (-1.285,-1.900) (-1.321,-1.850) (-1.350,-1.807) (-1.385,-1.750) (-1.412,-1.700) (-1.437,-1.650) (-1.460,-1.600) (-1.480,-1.550) (-1.498,-1.500) (-1.513,-1.450) (-1.527,-1.400) (-1.538,-1.350) (-1.547,-1.300) (-1.555,-1.250) (-1.560,-1.200) (-1.563,-1.150) (-1.563,-1.100) (-1.562,-1.050) (-1.559,-1.000) (-1.553,-0.950) (-1.545,-0.900) (-1.534,-0.850) (-1.522,-0.800) (-1.506,-0.750) (-1.488,-0.700) (-1.467,-0.650) (-1.442,-0.600) (-1.414,-0.550) (-1.382,-0.500) (-1.350,-0.456) (-1.305,-0.400) (-1.266,-0.359)}
\EQ{(1.275,-0.337) (1.314,-0.375) (1.350,-0.414) (1.400,-0.475) (1.436,-0.525) (1.467,-0.575) (1.495,-0.625) (1.519,-0.675) (1.540,-0.725) (1.559,-0.775) (1.574,-0.825) (1.587,-0.875) (1.598,-0.925) (1.606,-0.975) (1.612,-1.025) (1.616,-1.075) (1.618,-1.125) (1.618,-1.175) (1.616,-1.225) (1.611,-1.275) (1.605,-1.325) (1.597,-1.375) (1.586,-1.425) (1.574,-1.475) (1.559,-1.525) (1.543,-1.575) (1.524,-1.625) (1.503,-1.675) (1.479,-1.725) (1.454,-1.775) (1.425,-1.825) (1.394,-1.875) (1.360,-1.925) (1.325,-1.973) (1.283,-2.025) (1.250,-2.063) (1.216,-2.100) (1.175,-2.141) (1.125,-2.187) (1.081,-2.225) (1.025,-2.269) (0.983,-2.300) (0.925,-2.339) (0.875,-2.370) (0.825,-2.398) (0.775,-2.425) (0.725,-2.449) (0.675,-2.471) (0.625,-2.492) (0.575,-2.510) (0.525,-2.527) (0.475,-2.542) (0.425,-2.555) (0.375,-2.567) (0.325,-2.578) (0.275,-2.586) (0.225,-2.594) (0.175,-2.599) (0.125,-2.604) (0.075,-2.607) (0.025,-2.608) (-0.050,-2.608) (-0.100,-2.605) (-0.150,-2.602) (-0.200,-2.597) (-0.250,-2.590) (-0.300,-2.582) (-0.350,-2.573) (-0.400,-2.562) (-0.450,-2.549) (-0.500,-2.535) (-0.550,-2.519) (-0.600,-2.501) (-0.650,-2.482) (-0.700,-2.460) (-0.750,-2.437) (-0.800,-2.412) (-0.850,-2.384) (-0.900,-2.355) (-0.946,-2.325) (-1.000,-2.288) (-1.050,-2.250) (-1.100,-2.209) (-1.139,-2.175) (-1.191,-2.125) (-1.239,-2.075) (-1.275,-2.034) (-1.323,-1.975) (-1.360,-1.925) (-1.394,-1.875) (-1.425,-1.825) (-1.450,-1.781) (-1.475,-1.734) (-1.500,-1.681) (-1.524,-1.625) (-1.543,-1.575) (-1.559,-1.525) (-1.574,-1.475) (-1.586,-1.425) (-1.597,-1.375) (-1.605,-1.325) (-1.611,-1.275) (-1.616,-1.225) (-1.618,-1.175) (-1.618,-1.125) (-1.616,-1.075) (-1.612,-1.025) (-1.606,-0.975) (-1.598,-0.925) (-1.587,-0.875) (-1.574,-0.825) (-1.559,-0.775) (-1.540,-0.725) (-1.519,-0.675) (-1.495,-0.625) (-1.467,-0.575) (-1.436,-0.525) (-1.400,-0.475) (-1.360,-0.425) (-1.325,-0.387) (-1.289,-0.350) (-1.275,-0.337)}
\EQ{(1.275,-0.308) (1.323,-0.350) (1.374,-0.400) (1.419,-0.450) (1.450,-0.489) (1.493,-0.550) (1.524,-0.600) (1.550,-0.647) (1.575,-0.698) (1.597,-0.750) (1.616,-0.800) (1.632,-0.850) (1.645,-0.900) (1.656,-0.950) (1.665,-1.000) (1.672,-1.050) (1.676,-1.100) (1.679,-1.150) (1.679,-1.200) (1.678,-1.250) (1.674,-1.300) (1.669,-1.350) (1.662,-1.400) (1.652,-1.450) (1.641,-1.500) (1.628,-1.550) (1.613,-1.600) (1.595,-1.650) (1.576,-1.700) (1.554,-1.750) (1.530,-1.800) (1.504,-1.850) (1.475,-1.900) (1.444,-1.950) (1.410,-2.000) (1.375,-2.047) (1.333,-2.100) (1.289,-2.150) (1.250,-2.191) (1.200,-2.240) (1.162,-2.275) (1.102,-2.325) (1.050,-2.365) (1.001,-2.400) (0.950,-2.434) (0.900,-2.464) (0.850,-2.492) (0.800,-2.518) (0.750,-2.543) (0.700,-2.565) (0.650,-2.585) (0.600,-2.604) (0.550,-2.621) (0.500,-2.636) (0.450,-2.650) (0.400,-2.662) (0.350,-2.672) (0.300,-2.682) (0.250,-2.689) (0.200,-2.696) (0.150,-2.700) (0.100,-2.704) (0.050,-2.706) (0.000,-2.707) (-0.050,-2.706) (-0.100,-2.704) (-0.150,-2.700) (-0.200,-2.696) (-0.250,-2.689) (-0.300,-2.682) (-0.350,-2.672) (-0.400,-2.662) (-0.450,-2.650) (-0.500,-2.636) (-0.550,-2.621) (-0.600,-2.604) (-0.650,-2.585) (-0.700,-2.565) (-0.750,-2.543) (-0.800,-2.518) (-0.850,-2.492) (-0.900,-2.464) (-0.950,-2.434) (-1.000,-2.401) (-1.050,-2.365) (-1.100,-2.327) (-1.150,-2.285) (-1.189,-2.250) (-1.241,-2.200) (-1.289,-2.150) (-1.325,-2.109) (-1.373,-2.050) (-1.410,-2.000) (-1.444,-1.950) (-1.475,-1.900) (-1.504,-1.850) (-1.530,-1.800) (-1.554,-1.750) (-1.575,-1.702) (-1.595,-1.650) (-1.613,-1.600) (-1.628,-1.550) (-1.641,-1.500) (-1.652,-1.450) (-1.662,-1.400) (-1.669,-1.350) (-1.674,-1.300) (-1.678,-1.250) (-1.679,-1.200) (-1.679,-1.150) (-1.676,-1.100) (-1.672,-1.050) (-1.665,-1.000) (-1.656,-0.950) (-1.645,-0.900) (-1.632,-0.850) (-1.616,-0.800) (-1.597,-0.750) (-1.576,-0.700) (-1.552,-0.650) (-1.525,-0.601) (-1.493,-0.550) (-1.458,-0.500) (-1.425,-0.458) (-1.375,-0.401) (-1.325,-0.352) (-1.275,-0.308)}
\EQ{(1.275,-0.280) (1.325,-0.319) (1.375,-0.362) (1.414,-0.400) (1.450,-0.438) (1.500,-0.497) (1.539,-0.550) (1.572,-0.600) (1.600,-0.648) (1.625,-0.696) (1.650,-0.750) (1.670,-0.800) (1.688,-0.850) (1.703,-0.900) (1.716,-0.950) (1.726,-1.000) (1.734,-1.050) (1.741,-1.100) (1.745,-1.150) (1.747,-1.200) (1.748,-1.250) (1.746,-1.300) (1.743,-1.350) (1.738,-1.400) (1.731,-1.450) (1.722,-1.500) (1.711,-1.550) (1.698,-1.600) (1.683,-1.650) (1.667,-1.700) (1.648,-1.750) (1.627,-1.800) (1.604,-1.850) (1.579,-1.900) (1.552,-1.950) (1.525,-1.995) (1.489,-2.050) (1.454,-2.100) (1.416,-2.150) (1.375,-2.200) (1.330,-2.250) (1.282,-2.300) (1.229,-2.350) (1.175,-2.397) (1.125,-2.437) (1.075,-2.475) (1.025,-2.509) (0.975,-2.542) (0.925,-2.571) (0.875,-2.599) (0.825,-2.625) (0.775,-2.649) (0.725,-2.671) (0.675,-2.691) (0.625,-2.710) (0.575,-2.727) (0.525,-2.742) (0.475,-2.756) (0.425,-2.768) (0.375,-2.779) (0.325,-2.789) (0.275,-2.797) (0.225,-2.804) (0.175,-2.809) (0.125,-2.813) (0.075,-2.816) (0.025,-2.817) (-0.050,-2.817) (-0.100,-2.815) (-0.150,-2.811) (-0.200,-2.807) (-0.250,-2.800) (-0.300,-2.793) (-0.350,-2.784) (-0.400,-2.774) (-0.450,-2.762) (-0.500,-2.749) (-0.550,-2.735) (-0.600,-2.719) (-0.650,-2.701) (-0.700,-2.681) (-0.750,-2.660) (-0.800,-2.637) (-0.850,-2.612) (-0.900,-2.586) (-0.950,-2.557) (-1.000,-2.526) (-1.050,-2.492) (-1.100,-2.456) (-1.141,-2.425) (-1.200,-2.376) (-1.250,-2.331) (-1.300,-2.282) (-1.350,-2.228) (-1.396,-2.175) (-1.436,-2.125) (-1.472,-2.075) (-1.506,-2.025) (-1.537,-1.975) (-1.566,-1.925) (-1.592,-1.875) (-1.616,-1.825) (-1.638,-1.775) (-1.658,-1.725) (-1.675,-1.675) (-1.691,-1.625) (-1.705,-1.575) (-1.716,-1.525) (-1.726,-1.475) (-1.734,-1.425) (-1.741,-1.375) (-1.745,-1.325) (-1.747,-1.275) (-1.748,-1.225) (-1.746,-1.175) (-1.743,-1.125) (-1.738,-1.075) (-1.730,-1.025) (-1.721,-0.975) (-1.709,-0.925) (-1.695,-0.875) (-1.679,-0.825) (-1.660,-0.775) (-1.639,-0.725) (-1.614,-0.675) (-1.587,-0.625) (-1.556,-0.575) (-1.525,-0.530) (-1.482,-0.475) (-1.438,-0.425) (-1.400,-0.386) (-1.362,-0.350) (-1.301,-0.300) (-1.275,-0.280)}
\EQ{(1.288,-0.262) (1.342,-0.300) (1.400,-0.347) (1.450,-0.391) (1.500,-0.442) (1.550,-0.499) (1.590,-0.550) (1.625,-0.600) (1.656,-0.650) (1.684,-0.700) (1.708,-0.750) (1.730,-0.800) (1.750,-0.850) (1.767,-0.900) (1.781,-0.950) (1.794,-1.000) (1.804,-1.050) (1.812,-1.100) (1.818,-1.150) (1.823,-1.200) (1.825,-1.250) (1.825,-1.300) (1.824,-1.350) (1.821,-1.400) (1.816,-1.450) (1.810,-1.500) (1.801,-1.550) (1.791,-1.600) (1.779,-1.650) (1.765,-1.700) (1.750,-1.748) (1.732,-1.800) (1.712,-1.850) (1.691,-1.900) (1.667,-1.950) (1.641,-2.000) (1.613,-2.050) (1.583,-2.100) (1.550,-2.149) (1.514,-2.200) (1.476,-2.250) (1.434,-2.300) (1.400,-2.339) (1.366,-2.375) (1.325,-2.416) (1.289,-2.450) (1.250,-2.484) (1.201,-2.525) (1.150,-2.565) (1.101,-2.600) (1.050,-2.634) (1.000,-2.666) (0.950,-2.695) (0.900,-2.722) (0.850,-2.748) (0.800,-2.771) (0.750,-2.793) (0.700,-2.813) (0.650,-2.831) (0.600,-2.848) (0.550,-2.864) (0.500,-2.878) (0.450,-2.890) (0.400,-2.901) (0.350,-2.911) (0.300,-2.919) (0.250,-2.926) (0.200,-2.932) (0.150,-2.937) (0.100,-2.940) (0.050,-2.942) (0.000,-2.943) (-0.050,-2.942) (-0.100,-2.940) (-0.150,-2.937) (-0.200,-2.932) (-0.250,-2.926) (-0.300,-2.919) (-0.350,-2.911) (-0.400,-2.901) (-0.450,-2.890) (-0.500,-2.878) (-0.550,-2.864) (-0.600,-2.848) (-0.650,-2.831) (-0.700,-2.813) (-0.750,-2.793) (-0.800,-2.771) (-0.850,-2.748) (-0.900,-2.722) (-0.950,-2.695) (-1.000,-2.666) (-1.050,-2.634) (-1.100,-2.601) (-1.150,-2.565) (-1.200,-2.526) (-1.250,-2.484) (-1.289,-2.450) (-1.325,-2.416) (-1.366,-2.375) (-1.400,-2.339) (-1.434,-2.300) (-1.475,-2.251) (-1.514,-2.200) (-1.550,-2.150) (-1.583,-2.100) (-1.613,-2.050) (-1.641,-2.000) (-1.667,-1.950) (-1.691,-1.900) (-1.712,-1.850) (-1.732,-1.800) (-1.749,-1.750) (-1.765,-1.700) (-1.779,-1.650) (-1.791,-1.600) (-1.801,-1.550) (-1.810,-1.500) (-1.816,-1.450) (-1.821,-1.400) (-1.824,-1.350) (-1.825,-1.300) (-1.825,-1.250) (-1.823,-1.200) (-1.818,-1.150) (-1.812,-1.100) (-1.804,-1.050) (-1.794,-1.000) (-1.781,-0.950) (-1.767,-0.900) (-1.750,-0.851) (-1.730,-0.800) (-1.708,-0.750) (-1.684,-0.700) (-1.656,-0.650) (-1.625,-0.601) (-1.590,-0.550) (-1.551,-0.500) (-1.508,-0.450) (-1.459,-0.400) (-1.404,-0.350) (-1.350,-0.307) (-1.307,-0.275) (-1.288,-0.262)}
\EQ{(1.300,-0.243) (1.350,-0.275) (1.400,-0.310) (1.450,-0.348) (1.500,-0.391) (1.550,-0.439) (1.600,-0.493) (1.647,-0.550) (1.684,-0.600) (1.717,-0.650) (1.747,-0.700) (1.774,-0.750) (1.798,-0.800) (1.819,-0.850) (1.838,-0.900) (1.855,-0.950) (1.869,-1.000) (1.881,-1.050) (1.892,-1.100) (1.900,-1.150) (1.906,-1.200) (1.911,-1.250) (1.914,-1.300) (1.915,-1.350) (1.914,-1.400) (1.911,-1.450) (1.907,-1.500) (1.901,-1.550) (1.894,-1.600) (1.885,-1.650) (1.873,-1.700) (1.861,-1.750) (1.846,-1.800) (1.830,-1.850) (1.812,-1.900) (1.792,-1.950) (1.770,-2.000) (1.746,-2.050) (1.720,-2.100) (1.692,-2.150) (1.661,-2.200) (1.629,-2.250) (1.593,-2.300) (1.555,-2.350) (1.515,-2.400) (1.475,-2.445) (1.425,-2.498) (1.375,-2.548) (1.325,-2.593) (1.275,-2.636) (1.226,-2.675) (1.175,-2.713) (1.125,-2.747) (1.083,-2.775) (1.025,-2.810) (0.975,-2.838) (0.925,-2.865) (0.875,-2.890) (0.825,-2.913) (0.775,-2.934) (0.725,-2.954) (0.675,-2.972) (0.625,-2.989) (0.575,-3.004) (0.525,-3.018) (0.475,-3.030) (0.425,-3.042) (0.375,-3.052) (0.325,-3.060) (0.275,-3.068) (0.225,-3.074) (0.175,-3.079) (0.125,-3.082) (0.075,-3.085) (0.025,-3.086) (-0.050,-3.086) (-0.100,-3.084) (-0.150,-3.081) (-0.200,-3.076) (-0.250,-3.071) (-0.300,-3.064) (-0.350,-3.056) (-0.400,-3.047) (-0.450,-3.036) (-0.500,-3.024) (-0.550,-3.011) (-0.600,-2.996) (-0.650,-2.980) (-0.700,-2.963) (-0.750,-2.944) (-0.800,-2.923) (-0.850,-2.901) (-0.900,-2.877) (-0.950,-2.852) (-0.999,-2.825) (-1.050,-2.795) (-1.100,-2.764) (-1.150,-2.730) (-1.192,-2.700) (-1.250,-2.656) (-1.300,-2.615) (-1.345,-2.575) (-1.398,-2.525) (-1.448,-2.475) (-1.493,-2.425) (-1.535,-2.375) (-1.575,-2.325) (-1.611,-2.275) (-1.645,-2.225) (-1.675,-2.178) (-1.706,-2.125) (-1.733,-2.075) (-1.758,-2.025) (-1.781,-1.975) (-1.802,-1.925) (-1.821,-1.875) (-1.838,-1.825) (-1.854,-1.775) (-1.867,-1.725) (-1.879,-1.675) (-1.889,-1.625) (-1.898,-1.575) (-1.905,-1.525) (-1.910,-1.475) (-1.913,-1.425) (-1.915,-1.375) (-1.914,-1.325) (-1.913,-1.275) (-1.909,-1.225) (-1.903,-1.175) (-1.896,-1.125) (-1.887,-1.075) (-1.875,-1.025) (-1.862,-0.975) (-1.847,-0.925) (-1.829,-0.875) (-1.809,-0.825) (-1.786,-0.775) (-1.761,-0.725) (-1.733,-0.675) (-1.701,-0.625) (-1.666,-0.575) (-1.627,-0.525) (-1.584,-0.475) (-1.536,-0.425) (-1.482,-0.375) (-1.425,-0.328) (-1.375,-0.292) (-1.325,-0.258) (-1.300,-0.243)}
\EQ{(1.300,-0.216) (1.350,-0.244) (1.400,-0.274) (1.450,-0.308) (1.500,-0.345) (1.550,-0.385) (1.595,-0.425) (1.646,-0.475) (1.692,-0.525) (1.725,-0.565) (1.770,-0.625) (1.800,-0.669) (1.834,-0.725) (1.862,-0.775) (1.887,-0.825) (1.909,-0.875) (1.929,-0.925) (1.947,-0.975) (1.962,-1.025) (1.975,-1.075) (1.987,-1.125) (1.997,-1.175) (2.004,-1.225) (2.011,-1.275) (2.015,-1.325) (2.017,-1.375) (2.018,-1.425) (2.017,-1.475) (2.015,-1.525) (2.011,-1.575) (2.005,-1.625) (1.998,-1.675) (1.989,-1.725) (1.979,-1.775) (1.967,-1.825) (1.953,-1.875) (1.937,-1.925) (1.920,-1.975) (1.901,-2.025) (1.880,-2.075) (1.858,-2.125) (1.833,-2.175) (1.807,-2.225) (1.778,-2.275) (1.750,-2.321) (1.714,-2.375) (1.679,-2.425) (1.641,-2.475) (1.601,-2.525) (1.557,-2.575) (1.510,-2.625) (1.475,-2.661) (1.434,-2.700) (1.378,-2.750) (1.325,-2.794) (1.285,-2.825) (1.225,-2.869) (1.180,-2.900) (1.125,-2.935) (1.075,-2.965) (1.025,-2.993) (0.975,-3.020) (0.925,-3.045) (0.875,-3.068) (0.825,-3.089) (0.775,-3.109) (0.725,-3.128) (0.675,-3.145) (0.625,-3.161) (0.576,-3.175) (0.525,-3.188) (0.476,-3.200) (0.425,-3.211) (0.375,-3.220) (0.325,-3.229) (0.275,-3.236) (0.225,-3.241) (0.175,-3.246) (0.125,-3.249) (0.075,-3.252) (0.025,-3.253) (-0.050,-3.252) (-0.100,-3.251) (-0.150,-3.248) (-0.200,-3.244) (-0.250,-3.239) (-0.300,-3.232) (-0.350,-3.225) (-0.400,-3.216) (-0.450,-3.206) (-0.500,-3.195) (-0.550,-3.182) (-0.600,-3.168) (-0.650,-3.153) (-0.700,-3.137) (-0.750,-3.119) (-0.798,-3.100) (-0.850,-3.079) (-0.900,-3.056) (-0.950,-3.032) (-1.000,-3.007) (-1.050,-2.980) (-1.100,-2.950) (-1.150,-2.919) (-1.200,-2.886) (-1.250,-2.851) (-1.300,-2.814) (-1.348,-2.775) (-1.400,-2.731) (-1.450,-2.685) (-1.486,-2.650) (-1.525,-2.610) (-1.575,-2.555) (-1.621,-2.500) (-1.661,-2.450) (-1.697,-2.400) (-1.731,-2.350) (-1.763,-2.300) (-1.793,-2.250) (-1.820,-2.200) (-1.846,-2.150) (-1.869,-2.100) (-1.891,-2.050) (-1.911,-2.000) (-1.929,-1.950) (-1.945,-1.900) (-1.960,-1.850) (-1.973,-1.800) (-1.984,-1.750) (-1.994,-1.700) (-2.002,-1.650) (-2.008,-1.600) (-2.013,-1.550) (-2.017,-1.500) (-2.018,-1.450) (-2.018,-1.400) (-2.016,-1.350) (-2.013,-1.300) (-2.008,-1.250) (-2.001,-1.200) (-1.992,-1.150) (-1.982,-1.100) (-1.969,-1.050) (-1.955,-1.000) (-1.938,-0.950) (-1.919,-0.900) (-1.898,-0.850) (-1.875,-0.801) (-1.850,-0.753) (-1.825,-0.710) (-1.787,-0.650) (-1.752,-0.600) (-1.713,-0.550) (-1.675,-0.506) (-1.625,-0.454) (-1.575,-0.407) (-1.525,-0.364) (-1.475,-0.326) (-1.425,-0.291) (-1.375,-0.259) (-1.325,-0.230) (-1.300,-0.216)}
\EQ{(1.300,-0.191) (1.350,-0.215) (1.400,-0.241) (1.450,-0.270) (1.499,-0.300) (1.550,-0.335) (1.600,-0.372) (1.650,-0.414) (1.691,-0.450) (1.742,-0.500) (1.788,-0.550) (1.825,-0.594) (1.868,-0.650) (1.900,-0.696) (1.934,-0.750) (1.963,-0.800) (1.989,-0.850) (2.013,-0.900) (2.034,-0.950) (2.053,-1.000) (2.070,-1.050) (2.085,-1.100) (2.098,-1.150) (2.109,-1.200) (2.118,-1.250) (2.126,-1.300) (2.132,-1.350) (2.137,-1.400) (2.139,-1.450) (2.141,-1.500) (2.140,-1.550) (2.138,-1.600) (2.135,-1.650) (2.130,-1.700) (2.123,-1.750) (2.115,-1.800) (2.106,-1.850) (2.095,-1.900) (2.082,-1.950) (2.068,-2.000) (2.052,-2.050) (2.035,-2.100) (2.016,-2.150) (1.995,-2.200) (1.972,-2.250) (1.950,-2.296) (1.925,-2.344) (1.900,-2.389) (1.864,-2.450) (1.831,-2.500) (1.800,-2.546) (1.760,-2.600) (1.725,-2.645) (1.679,-2.700) (1.634,-2.750) (1.600,-2.786) (1.562,-2.825) (1.525,-2.860) (1.481,-2.900) (1.425,-2.947) (1.375,-2.987) (1.325,-3.025) (1.275,-3.060) (1.225,-3.093) (1.175,-3.124) (1.132,-3.150) (1.087,-3.175) (1.040,-3.200) (0.990,-3.225) (0.936,-3.250) (0.879,-3.275) (0.825,-3.297) (0.775,-3.315) (0.725,-3.333) (0.675,-3.349) (0.625,-3.364) (0.575,-3.377) (0.525,-3.390) (0.475,-3.401) (0.425,-3.411) (0.375,-3.420) (0.325,-3.427) (0.275,-3.434) (0.225,-3.439) (0.175,-3.444) (0.125,-3.447) (0.075,-3.449) (0.025,-3.450) (-0.037,-3.450) (-0.100,-3.448) (-0.150,-3.445) (-0.200,-3.442) (-0.250,-3.437) (-0.300,-3.431) (-0.350,-3.424) (-0.400,-3.415) (-0.450,-3.406) (-0.500,-3.395) (-0.550,-3.384) (-0.600,-3.371) (-0.650,-3.356) (-0.700,-3.341) (-0.748,-3.325) (-0.800,-3.306) (-0.850,-3.287) (-0.900,-3.266) (-0.950,-3.244) (-1.000,-3.220) (-1.050,-3.195) (-1.100,-3.168) (-1.150,-3.139) (-1.200,-3.109) (-1.250,-3.077) (-1.300,-3.043) (-1.350,-3.006) (-1.391,-2.975) (-1.450,-2.927) (-1.500,-2.883) (-1.550,-2.836) (-1.587,-2.800) (-1.625,-2.760) (-1.675,-2.705) (-1.721,-2.650) (-1.760,-2.600) (-1.797,-2.550) (-1.831,-2.500) (-1.864,-2.450) (-1.894,-2.400) (-1.922,-2.350) (-1.948,-2.300) (-1.972,-2.250) (-1.995,-2.200) (-2.016,-2.150) (-2.035,-2.100) (-2.052,-2.050) (-2.068,-2.000) (-2.082,-1.950) (-2.095,-1.900) (-2.106,-1.850) (-2.115,-1.800) (-2.123,-1.750) (-2.130,-1.700) (-2.135,-1.650) (-2.138,-1.600) (-2.140,-1.550) (-2.141,-1.500) (-2.139,-1.450) (-2.137,-1.400) (-2.132,-1.350) (-2.126,-1.300) (-2.118,-1.250) (-2.109,-1.200) (-2.098,-1.150) (-2.085,-1.100) (-2.070,-1.050) (-2.053,-1.000) (-2.034,-0.950) (-2.013,-0.900) (-1.989,-0.850) (-1.963,-0.800) (-1.934,-0.750) (-1.903,-0.700) (-1.868,-0.650) (-1.830,-0.600) (-1.788,-0.550) (-1.750,-0.508) (-1.700,-0.459) (-1.650,-0.414) (-1.603,-0.375) (-1.550,-0.335) (-1.500,-0.301) (-1.450,-0.270) (-1.400,-0.241) (-1.350,-0.215) (-1.300,-0.191)}
\EQ{(1.300,-0.165) (1.350,-0.186) (1.400,-0.208) (1.450,-0.233) (1.500,-0.259) (1.550,-0.288) (1.600,-0.319) (1.646,-0.350) (1.700,-0.390) (1.744,-0.425) (1.800,-0.474) (1.850,-0.522) (1.900,-0.575) (1.943,-0.625) (1.975,-0.665) (2.018,-0.725) (2.050,-0.773) (2.082,-0.825) (2.109,-0.875) (2.134,-0.925) (2.157,-0.975) (2.178,-1.025) (2.197,-1.075) (2.213,-1.125) (2.228,-1.175) (2.242,-1.225) (2.253,-1.275) (2.263,-1.325) (2.271,-1.375) (2.277,-1.425) (2.282,-1.475) (2.286,-1.525) (2.288,-1.575) (2.288,-1.625) (2.288,-1.675) (2.285,-1.725) (2.281,-1.775) (2.276,-1.825) (2.269,-1.875) (2.261,-1.925) (2.252,-1.975) (2.241,-2.025) (2.228,-2.075) (2.215,-2.125) (2.199,-2.175) (2.182,-2.225) (2.164,-2.275) (2.144,-2.325) (2.122,-2.375) (2.100,-2.423) (2.075,-2.473) (2.050,-2.520) (2.025,-2.565) (1.989,-2.625) (1.956,-2.675) (1.925,-2.720) (1.885,-2.775) (1.850,-2.820) (1.805,-2.875) (1.761,-2.925) (1.725,-2.964) (1.690,-3.000) (1.650,-3.039) (1.611,-3.075) (1.554,-3.125) (1.500,-3.170) (1.450,-3.209) (1.400,-3.245) (1.358,-3.275) (1.300,-3.313) (1.250,-3.344) (1.200,-3.373) (1.152,-3.400) (1.104,-3.425) (1.053,-3.450) (1.000,-3.475) (0.950,-3.497) (0.900,-3.517) (0.850,-3.536) (0.800,-3.554) (0.750,-3.571) (0.700,-3.587) (0.650,-3.601) (0.600,-3.614) (0.550,-3.626) (0.500,-3.637) (0.450,-3.647) (0.400,-3.656) (0.350,-3.664) (0.300,-3.670) (0.250,-3.676) (0.200,-3.680) (0.150,-3.684) (0.100,-3.687) (0.050,-3.688) (0.000,-3.689) (-0.050,-3.688) (-0.100,-3.687) (-0.150,-3.684) (-0.200,-3.680) (-0.250,-3.676) (-0.300,-3.670) (-0.350,-3.664) (-0.400,-3.656) (-0.450,-3.647) (-0.500,-3.637) (-0.550,-3.626) (-0.600,-3.614) (-0.650,-3.601) (-0.700,-3.587) (-0.750,-3.571) (-0.800,-3.554) (-0.850,-3.536) (-0.900,-3.517) (-0.950,-3.497) (-1.000,-3.475) (-1.050,-3.452) (-1.100,-3.427) (-1.150,-3.401) (-1.197,-3.375) (-1.250,-3.344) (-1.300,-3.313) (-1.350,-3.280) (-1.394,-3.250) (-1.450,-3.209) (-1.494,-3.175) (-1.550,-3.129) (-1.600,-3.085) (-1.638,-3.050) (-1.675,-3.015) (-1.714,-2.975) (-1.750,-2.937) (-1.800,-2.881) (-1.846,-2.825) (-1.885,-2.775) (-1.922,-2.725) (-1.956,-2.675) (-1.989,-2.625) (-2.019,-2.575) (-2.047,-2.525) (-2.074,-2.475) (-2.099,-2.425) (-2.122,-2.375) (-2.144,-2.325) (-2.164,-2.275) (-2.182,-2.225) (-2.199,-2.175) (-2.215,-2.125) (-2.228,-2.075) (-2.241,-2.025) (-2.252,-1.975) (-2.261,-1.925) (-2.269,-1.875) (-2.276,-1.825) (-2.281,-1.775) (-2.285,-1.725) (-2.288,-1.675) (-2.288,-1.625) (-2.288,-1.575) (-2.286,-1.525) (-2.282,-1.475) (-2.277,-1.425) (-2.271,-1.375) (-2.263,-1.325) (-2.253,-1.275) (-2.242,-1.225) (-2.228,-1.175) (-2.213,-1.125) (-2.197,-1.075) (-2.178,-1.025) (-2.157,-0.975) (-2.134,-0.925) (-2.109,-0.875) (-2.082,-0.825) (-2.051,-0.775) (-2.018,-0.725) (-1.982,-0.675) (-1.950,-0.634) (-1.900,-0.575) (-1.853,-0.525) (-1.801,-0.475) (-1.750,-0.430) (-1.700,-0.390) (-1.650,-0.353) (-1.600,-0.319) (-1.550,-0.288) (-1.500,-0.259) (-1.450,-0.233) (-1.400,-0.208) (-1.350,-0.186) (-1.300,-0.165)}
\EQ{(1.300,-0.141) (1.350,-0.158) (1.400,-0.177) (1.450,-0.197) (1.500,-0.220) (1.550,-0.243) (1.600,-0.269) (1.650,-0.297) (1.697,-0.325) (1.750,-0.359) (1.800,-0.394) (1.841,-0.425) (1.900,-0.473) (1.950,-0.517) (2.000,-0.566) (2.050,-0.619) (2.099,-0.675) (2.138,-0.725) (2.174,-0.775) (2.208,-0.825) (2.239,-0.875) (2.267,-0.925) (2.293,-0.975) (2.317,-1.025) (2.339,-1.075) (2.359,-1.125) (2.377,-1.175) (2.394,-1.225) (2.409,-1.275) (2.422,-1.325) (2.433,-1.375) (2.443,-1.425) (2.451,-1.475) (2.458,-1.525) (2.464,-1.575) (2.468,-1.625) (2.470,-1.675) (2.472,-1.725) (2.472,-1.775) (2.470,-1.825) (2.467,-1.875) (2.463,-1.925) (2.458,-1.975) (2.451,-2.025) (2.443,-2.075) (2.434,-2.125) (2.423,-2.175) (2.411,-2.225) (2.398,-2.275) (2.383,-2.325) (2.367,-2.375) (2.350,-2.423) (2.330,-2.475) (2.310,-2.525) (2.288,-2.575) (2.265,-2.625) (2.240,-2.675) (2.213,-2.725) (2.184,-2.775) (2.154,-2.825) (2.125,-2.871) (2.089,-2.925) (2.053,-2.975) (2.015,-3.025) (1.975,-3.075) (1.933,-3.125) (1.888,-3.175) (1.850,-3.215) (1.800,-3.265) (1.763,-3.300) (1.725,-3.335) (1.679,-3.375) (1.625,-3.420) (1.575,-3.459) (1.525,-3.496) (1.475,-3.531) (1.425,-3.565) (1.375,-3.597) (1.328,-3.625) (1.285,-3.650) (1.239,-3.675) (1.191,-3.700) (1.141,-3.725) (1.087,-3.750) (1.031,-3.775) (0.975,-3.798) (0.925,-3.817) (0.875,-3.835) (0.825,-3.852) (0.775,-3.868) (0.725,-3.883) (0.675,-3.897) (0.625,-3.909) (0.575,-3.921) (0.525,-3.932) (0.475,-3.941) (0.425,-3.950) (0.375,-3.958) (0.325,-3.964) (0.275,-3.970) (0.225,-3.975) (0.175,-3.978) (0.125,-3.981) (0.075,-3.983) (0.025,-3.984) (-0.050,-3.984) (-0.100,-3.982) (-0.150,-3.980) (-0.200,-3.977) (-0.250,-3.972) (-0.300,-3.967) (-0.350,-3.961) (-0.400,-3.954) (-0.450,-3.946) (-0.500,-3.937) (-0.550,-3.926) (-0.600,-3.915) (-0.650,-3.903) (-0.700,-3.890) (-0.750,-3.876) (-0.800,-3.860) (-0.850,-3.844) (-0.900,-3.826) (-0.950,-3.808) (-1.000,-3.788) (-1.050,-3.767) (-1.100,-3.744) (-1.150,-3.721) (-1.200,-3.696) (-1.250,-3.669) (-1.300,-3.641) (-1.350,-3.612) (-1.400,-3.581) (-1.448,-3.550) (-1.500,-3.514) (-1.550,-3.478) (-1.600,-3.440) (-1.650,-3.400) (-1.700,-3.357) (-1.750,-3.312) (-1.790,-3.275) (-1.840,-3.225) (-1.875,-3.189) (-1.910,-3.150) (-1.950,-3.105) (-1.995,-3.050) (-2.034,-3.000) (-2.071,-2.950) (-2.100,-2.908) (-2.139,-2.850) (-2.170,-2.800) (-2.199,-2.750) (-2.225,-2.703) (-2.250,-2.654) (-2.275,-2.603) (-2.299,-2.550) (-2.320,-2.500) (-2.340,-2.450) (-2.358,-2.400) (-2.375,-2.351) (-2.391,-2.300) (-2.405,-2.250) (-2.417,-2.200) (-2.429,-2.150) (-2.439,-2.100) (-2.447,-2.050) (-2.455,-2.000) (-2.461,-1.950) (-2.466,-1.900) (-2.469,-1.850) (-2.471,-1.800) (-2.472,-1.750) (-2.471,-1.700) (-2.469,-1.650) (-2.466,-1.600) (-2.461,-1.550) (-2.455,-1.500) (-2.447,-1.450) (-2.438,-1.400) (-2.427,-1.350) (-2.415,-1.300) (-2.401,-1.250) (-2.386,-1.200) (-2.369,-1.150) (-2.350,-1.101) (-2.329,-1.050) (-2.306,-1.000) (-2.281,-0.950) (-2.253,-0.900) (-2.225,-0.852) (-2.192,-0.800) (-2.157,-0.750) (-2.125,-0.708) (-2.078,-0.650) (-2.033,-0.600) (-1.984,-0.550) (-1.931,-0.500) (-1.875,-0.452) (-1.825,-0.413) (-1.775,-0.376) (-1.725,-0.343) (-1.675,-0.312) (-1.625,-0.283) (-1.575,-0.256) (-1.525,-0.231) (-1.475,-0.208) (-1.425,-0.187) (-1.375,-0.167) (-1.327,-0.150) (-1.300,-0.141)}
\EQ{(1.300,-0.117) (1.350,-0.131) (1.400,-0.146) (1.450,-0.163) (1.500,-0.181) (1.549,-0.200) (1.600,-0.221) (1.650,-0.244) (1.700,-0.268) (1.750,-0.293) (1.800,-0.321) (1.849,-0.350) (1.900,-0.382) (1.950,-0.416) (1.997,-0.450) (2.050,-0.492) (2.090,-0.525) (2.145,-0.575) (2.196,-0.625) (2.243,-0.675) (2.287,-0.725) (2.325,-0.772) (2.365,-0.825) (2.400,-0.875) (2.432,-0.925) (2.462,-0.975) (2.490,-1.025) (2.516,-1.075) (2.540,-1.125) (2.562,-1.175) (2.582,-1.225) (2.600,-1.275) (2.617,-1.325) (2.632,-1.375) (2.646,-1.425) (2.658,-1.475) (2.669,-1.525) (2.679,-1.575) (2.687,-1.625) (2.693,-1.675) (2.699,-1.725) (2.703,-1.775) (2.706,-1.825) (2.707,-1.875) (2.708,-1.925) (2.707,-1.975) (2.705,-2.025) (2.701,-2.075) (2.697,-2.125) (2.691,-2.175) (2.684,-2.225) (2.676,-2.275) (2.667,-2.325) (2.656,-2.375) (2.644,-2.425) (2.631,-2.475) (2.617,-2.525) (2.602,-2.575) (2.585,-2.625) (2.567,-2.675) (2.547,-2.725) (2.527,-2.775) (2.504,-2.825) (2.481,-2.875) (2.456,-2.925) (2.430,-2.975) (2.402,-3.025) (2.375,-3.070) (2.341,-3.125) (2.308,-3.175) (2.275,-3.223) (2.237,-3.275) (2.200,-3.323) (2.158,-3.375) (2.125,-3.414) (2.075,-3.470) (2.025,-3.523) (1.975,-3.574) (1.925,-3.621) (1.875,-3.667) (1.836,-3.700) (1.776,-3.750) (1.725,-3.790) (1.677,-3.825) (1.625,-3.862) (1.575,-3.896) (1.530,-3.925) (1.475,-3.959) (1.425,-3.988) (1.375,-4.016) (1.325,-4.043) (1.275,-4.068) (1.225,-4.092) (1.175,-4.115) (1.125,-4.137) (1.075,-4.158) (1.025,-4.177) (0.975,-4.196) (0.925,-4.213) (0.875,-4.229) (0.825,-4.245) (0.775,-4.259) (0.725,-4.272) (0.675,-4.285) (0.625,-4.296) (0.575,-4.307) (0.525,-4.317) (0.475,-4.325) (0.425,-4.333) (0.375,-4.340) (0.325,-4.346) (0.275,-4.351) (0.225,-4.356) (0.175,-4.359) (0.125,-4.362) (0.075,-4.363) (0.025,-4.364) (-0.050,-4.364) (-0.100,-4.363) (-0.150,-4.361) (-0.200,-4.358) (-0.250,-4.354) (-0.300,-4.349) (-0.350,-4.343) (-0.400,-4.337) (-0.450,-4.329) (-0.500,-4.321) (-0.550,-4.312) (-0.600,-4.302) (-0.650,-4.291) (-0.700,-4.279) (-0.750,-4.266) (-0.800,-4.252) (-0.850,-4.237) (-0.900,-4.221) (-0.950,-4.204) (-1.000,-4.187) (-1.050,-4.168) (-1.100,-4.147) (-1.150,-4.126) (-1.200,-4.104) (-1.250,-4.080) (-1.300,-4.056) (-1.350,-4.030) (-1.400,-4.002) (-1.448,-3.975) (-1.500,-3.944) (-1.550,-3.912) (-1.600,-3.879) (-1.642,-3.850) (-1.700,-3.808) (-1.744,-3.775) (-1.800,-3.730) (-1.850,-3.688) (-1.894,-3.650) (-1.948,-3.600) (-1.999,-3.550) (-2.047,-3.500) (-2.093,-3.450) (-2.137,-3.400) (-2.175,-3.354) (-2.218,-3.300) (-2.250,-3.257) (-2.291,-3.200) (-2.325,-3.150) (-2.357,-3.100) (-2.387,-3.050) (-2.416,-3.000) (-2.443,-2.950) (-2.469,-2.900) (-2.493,-2.850) (-2.516,-2.800) (-2.537,-2.750) (-2.557,-2.700) (-2.575,-2.653) (-2.593,-2.600) (-2.610,-2.550) (-2.624,-2.500) (-2.638,-2.450) (-2.650,-2.400) (-2.662,-2.350) (-2.672,-2.300) (-2.680,-2.250) (-2.688,-2.200) (-2.694,-2.150) (-2.699,-2.100) (-2.703,-2.050) (-2.706,-2.000) (-2.707,-1.950) (-2.708,-1.900) (-2.707,-1.850) (-2.704,-1.800) (-2.701,-1.750) (-2.696,-1.700) (-2.690,-1.650) (-2.683,-1.600) (-2.674,-1.550) (-2.664,-1.500) (-2.652,-1.450) (-2.639,-1.400) (-2.625,-1.351) (-2.609,-1.300) (-2.591,-1.250) (-2.572,-1.200) (-2.551,-1.150) (-2.528,-1.100) (-2.503,-1.050) (-2.476,-1.000) (-2.450,-0.954) (-2.416,-0.900) (-2.383,-0.850) (-2.350,-0.805) (-2.308,-0.750) (-2.275,-0.711) (-2.225,-0.655) (-2.175,-0.604) (-2.125,-0.557) (-2.075,-0.513) (-2.029,-0.475) (-1.975,-0.434) (-1.927,-0.400) (-1.875,-0.366) (-1.825,-0.335) (-1.775,-0.307) (-1.725,-0.280) (-1.675,-0.255) (-1.625,-0.232) (-1.575,-0.211) (-1.525,-0.191) (-1.475,-0.172) (-1.425,-0.155) (-1.375,-0.139) (-1.325,-0.124) (-1.300,-0.117)}
\EQ{(1.300,-0.093) (1.350,-0.104) (1.400,-0.116) (1.450,-0.130) (1.500,-0.144) (1.550,-0.159) (1.600,-0.175) (1.650,-0.193) (1.700,-0.211) (1.750,-0.231) (1.800,-0.252) (1.850,-0.275) (1.900,-0.299) (1.950,-0.324) (1.997,-0.350) (2.041,-0.375) (2.100,-0.412) (2.150,-0.444) (2.194,-0.475) (2.250,-0.517) (2.292,-0.550) (2.350,-0.599) (2.400,-0.645) (2.450,-0.693) (2.500,-0.746) (2.548,-0.800) (2.589,-0.850) (2.625,-0.896) (2.664,-0.950) (2.697,-1.000) (2.725,-1.044) (2.759,-1.100) (2.786,-1.150) (2.812,-1.200) (2.836,-1.250) (2.858,-1.300) (2.879,-1.350) (2.898,-1.400) (2.916,-1.450) (2.932,-1.500) (2.947,-1.550) (2.960,-1.600) (2.972,-1.650) (2.983,-1.700) (2.993,-1.750) (3.001,-1.800) (3.008,-1.850) (3.014,-1.900) (3.019,-1.950) (3.023,-2.000) (3.025,-2.050) (3.027,-2.100) (3.027,-2.150) (3.026,-2.200) (3.025,-2.250) (3.022,-2.300) (3.018,-2.350) (3.013,-2.400) (3.006,-2.450) (2.999,-2.500) (2.991,-2.550) (2.981,-2.600) (2.971,-2.650) (2.959,-2.700) (2.947,-2.750) (2.933,-2.800) (2.918,-2.850) (2.902,-2.900) (2.885,-2.950) (2.866,-3.000) (2.847,-3.050) (2.826,-3.100) (2.804,-3.150) (2.781,-3.200) (2.756,-3.250) (2.730,-3.300) (2.703,-3.350) (2.675,-3.399) (2.645,-3.450) (2.613,-3.500) (2.580,-3.550) (2.550,-3.594) (2.509,-3.650) (2.475,-3.696) (2.432,-3.750) (2.400,-3.789) (2.350,-3.846) (2.301,-3.900) (2.253,-3.950) (2.203,-4.000) (2.151,-4.050) (2.100,-4.096) (2.050,-4.139) (2.007,-4.175) (1.950,-4.220) (1.900,-4.258) (1.850,-4.295) (1.807,-4.325) (1.750,-4.363) (1.700,-4.395) (1.651,-4.425) (1.600,-4.455) (1.550,-4.483) (1.500,-4.510) (1.450,-4.536) (1.400,-4.561) (1.350,-4.585) (1.300,-4.607) (1.250,-4.629) (1.200,-4.649) (1.150,-4.669) (1.100,-4.687) (1.050,-4.705) (1.000,-4.722) (0.950,-4.737) (0.900,-4.752) (0.850,-4.766) (0.800,-4.779) (0.750,-4.792) (0.700,-4.803) (0.650,-4.814) (0.600,-4.824) (0.550,-4.833) (0.500,-4.841) (0.450,-4.848) (0.400,-4.855) (0.350,-4.861) (0.300,-4.866) (0.250,-4.870) (0.200,-4.873) (0.150,-4.876) (0.100,-4.878) (0.050,-4.879) (0.000,-4.880) (-0.050,-4.879) (-0.100,-4.878) (-0.150,-4.876) (-0.200,-4.873) (-0.250,-4.870) (-0.300,-4.866) (-0.350,-4.861) (-0.400,-4.855) (-0.450,-4.848) (-0.500,-4.841) (-0.550,-4.833) (-0.600,-4.824) (-0.650,-4.814) (-0.700,-4.803) (-0.750,-4.792) (-0.800,-4.779) (-0.850,-4.766) (-0.900,-4.752) (-0.950,-4.737) (-1.000,-4.722) (-1.050,-4.705) (-1.100,-4.687) (-1.150,-4.669) (-1.198,-4.650) (-1.250,-4.629) (-1.300,-4.607) (-1.350,-4.585) (-1.400,-4.561) (-1.450,-4.536) (-1.500,-4.510) (-1.550,-4.483) (-1.600,-4.455) (-1.650,-4.426) (-1.700,-4.395) (-1.750,-4.363) (-1.800,-4.330) (-1.843,-4.300) (-1.900,-4.258) (-1.944,-4.225) (-2.000,-4.181) (-2.050,-4.139) (-2.095,-4.100) (-2.150,-4.051) (-2.200,-4.003) (-2.250,-3.954) (-2.300,-3.901) (-2.347,-3.850) (-2.390,-3.800) (-2.425,-3.758) (-2.472,-3.700) (-2.509,-3.650) (-2.546,-3.600) (-2.575,-3.558) (-2.613,-3.500) (-2.645,-3.450) (-2.675,-3.400) (-2.700,-3.356) (-2.725,-3.310) (-2.750,-3.262) (-2.775,-3.212) (-2.800,-3.159) (-2.825,-3.102) (-2.847,-3.050) (-2.866,-3.000) (-2.885,-2.950) (-2.902,-2.900) (-2.918,-2.850) (-2.933,-2.800) (-2.947,-2.750) (-2.959,-2.700) (-2.971,-2.650) (-2.981,-2.600) (-2.991,-2.550) (-2.999,-2.500) (-3.006,-2.450) (-3.013,-2.400) (-3.018,-2.350) (-3.022,-2.300) (-3.025,-2.250) (-3.026,-2.200) (-3.027,-2.150) (-3.027,-2.100) (-3.025,-2.050) (-3.023,-2.000) (-3.019,-1.950) (-3.014,-1.900) (-3.008,-1.850) (-3.001,-1.800) (-2.993,-1.750) (-2.983,-1.700) (-2.972,-1.650) (-2.960,-1.600) (-2.947,-1.550) (-2.932,-1.500) (-2.916,-1.450) (-2.898,-1.400) (-2.879,-1.350) (-2.858,-1.300) (-2.836,-1.250) (-2.812,-1.200) (-2.786,-1.150) (-2.759,-1.100) (-2.729,-1.050) (-2.700,-1.004) (-2.664,-0.950) (-2.628,-0.900) (-2.589,-0.850) (-2.550,-0.802) (-2.504,-0.750) (-2.457,-0.700) (-2.406,-0.650) (-2.351,-0.600) (-2.300,-0.557) (-2.261,-0.525) (-2.200,-0.479) (-2.158,-0.450) (-2.100,-0.412) (-2.050,-0.381) (-2.000,-0.352) (-1.951,-0.325) (-1.902,-0.300) (-1.850,-0.275) (-1.800,-0.252) (-1.750,-0.231) (-1.700,-0.211) (-1.650,-0.193) (-1.600,-0.175) (-1.550,-0.159) (-1.500,-0.144) (-1.450,-0.130) (-1.400,-0.116) (-1.350,-0.104) (-1.300,-0.093)}
\EQ{(1.300,-0.070) (1.350,-0.078) (1.400,-0.087) (1.450,-0.097) (1.500,-0.107) (1.550,-0.118) (1.600,-0.130) (1.650,-0.143) (1.700,-0.157) (1.750,-0.171) (1.800,-0.187) (1.850,-0.203) (1.900,-0.220) (1.950,-0.239) (2.000,-0.258) (2.050,-0.279) (2.098,-0.300) (2.150,-0.324) (2.200,-0.348) (2.250,-0.374) (2.298,-0.400) (2.342,-0.425) (2.400,-0.460) (2.450,-0.492) (2.500,-0.525) (2.550,-0.561) (2.600,-0.598) (2.650,-0.638) (2.694,-0.675) (2.750,-0.724) (2.800,-0.772) (2.850,-0.822) (2.899,-0.875) (2.943,-0.925) (2.985,-0.975) (3.024,-1.025) (3.060,-1.075) (3.095,-1.125) (3.125,-1.171) (3.159,-1.225) (3.188,-1.275) (3.216,-1.325) (3.241,-1.375) (3.266,-1.425) (3.289,-1.475) (3.310,-1.525) (3.330,-1.575) (3.349,-1.625) (3.367,-1.675) (3.383,-1.725) (3.398,-1.775) (3.412,-1.825) (3.425,-1.875) (3.437,-1.925) (3.447,-1.975) (3.456,-2.025) (3.465,-2.075) (3.472,-2.125) (3.479,-2.175) (3.484,-2.225) (3.488,-2.275) (3.491,-2.325) (3.494,-2.375) (3.495,-2.425) (3.496,-2.475) (3.495,-2.525) (3.494,-2.575) (3.491,-2.625) (3.488,-2.675) (3.483,-2.725) (3.478,-2.775) (3.472,-2.825) (3.465,-2.875) (3.457,-2.925) (3.448,-2.975) (3.438,-3.025) (3.427,-3.075) (3.415,-3.125) (3.403,-3.175) (3.389,-3.225) (3.374,-3.275) (3.359,-3.325) (3.342,-3.375) (3.325,-3.423) (3.306,-3.475) (3.286,-3.525) (3.265,-3.575) (3.243,-3.625) (3.220,-3.675) (3.196,-3.725) (3.171,-3.775) (3.145,-3.825) (3.118,-3.875) (3.089,-3.925) (3.059,-3.975) (3.028,-4.025) (3.000,-4.068) (2.962,-4.125) (2.927,-4.175) (2.890,-4.225) (2.852,-4.275) (2.812,-4.325) (2.775,-4.370) (2.728,-4.425) (2.683,-4.475) (2.637,-4.525) (2.600,-4.563) (2.563,-4.600) (2.525,-4.637) (2.485,-4.675) (2.430,-4.725) (2.375,-4.773) (2.325,-4.815) (2.281,-4.850) (2.225,-4.894) (2.183,-4.925) (2.125,-4.967) (2.078,-5.000) (2.025,-5.035) (1.975,-5.067) (1.925,-5.098) (1.881,-5.125) (1.825,-5.157) (1.775,-5.185) (1.725,-5.212) (1.675,-5.238) (1.625,-5.262) (1.575,-5.286) (1.525,-5.309) (1.475,-5.331) (1.425,-5.352) (1.375,-5.373) (1.325,-5.392) (1.275,-5.411) (1.225,-5.428) (1.175,-5.445) (1.125,-5.462) (1.075,-5.477) (1.025,-5.492) (0.975,-5.505) (0.925,-5.519) (0.875,-5.531) (0.825,-5.543) (0.775,-5.553) (0.725,-5.564) (0.675,-5.573) (0.625,-5.582) (0.575,-5.590) (0.525,-5.597) (0.475,-5.604) (0.425,-5.610) (0.375,-5.616) (0.325,-5.620) (0.275,-5.624) (0.225,-5.628) (0.175,-5.630) (0.125,-5.632) (0.075,-5.634) (0.025,-5.634) (-0.050,-5.634) (-0.100,-5.633) (-0.150,-5.631) (-0.200,-5.629) (-0.250,-5.626) (-0.300,-5.622) (-0.350,-5.618) (-0.400,-5.613) (-0.450,-5.607) (-0.500,-5.601) (-0.550,-5.594) (-0.600,-5.586) (-0.650,-5.578) (-0.700,-5.569) (-0.750,-5.559) (-0.800,-5.548) (-0.850,-5.537) (-0.899,-5.525) (-0.950,-5.512) (-1.000,-5.499) (-1.050,-5.484) (-1.100,-5.469) (-1.150,-5.453) (-1.200,-5.437) (-1.250,-5.420) (-1.300,-5.401) (-1.350,-5.382) (-1.400,-5.363) (-1.450,-5.342) (-1.500,-5.320) (-1.550,-5.298) (-1.599,-5.275) (-1.650,-5.250) (-1.700,-5.225) (-1.747,-5.200) (-1.793,-5.175) (-1.838,-5.150) (-1.900,-5.114) (-1.950,-5.083) (-2.000,-5.051) (-2.050,-5.019) (-2.100,-4.985) (-2.149,-4.950) (-2.200,-4.913) (-2.249,-4.875) (-2.300,-4.835) (-2.343,-4.800) (-2.400,-4.751) (-2.450,-4.707) (-2.500,-4.661) (-2.538,-4.625) (-2.575,-4.588) (-2.613,-4.550) (-2.650,-4.511) (-2.700,-4.457) (-2.750,-4.400) (-2.792,-4.350) (-2.825,-4.309) (-2.871,-4.250) (-2.908,-4.200) (-2.944,-4.150) (-2.975,-4.106) (-3.012,-4.050) (-3.044,-4.000) (-3.074,-3.950) (-3.100,-3.906) (-3.125,-3.862) (-3.150,-3.816) (-3.175,-3.768) (-3.200,-3.718) (-3.225,-3.665) (-3.250,-3.610) (-3.275,-3.551) (-3.296,-3.500) (-3.315,-3.450) (-3.333,-3.400) (-3.350,-3.351) (-3.366,-3.300) (-3.382,-3.250) (-3.396,-3.200) (-3.409,-3.150) (-3.421,-3.100) (-3.433,-3.050) (-3.443,-3.000) (-3.453,-2.950) (-3.461,-2.900) (-3.469,-2.850) (-3.475,-2.800) (-3.481,-2.750) (-3.486,-2.700) (-3.490,-2.650) (-3.492,-2.600) (-3.494,-2.550) (-3.495,-2.500) (-3.496,-2.450) (-3.495,-2.400) (-3.493,-2.350) (-3.490,-2.300) (-3.486,-2.250) (-3.481,-2.200) (-3.475,-2.150) (-3.469,-2.100) (-3.461,-2.050) (-3.452,-2.000) (-3.442,-1.950) (-3.431,-1.900) (-3.419,-1.850) (-3.405,-1.800) (-3.391,-1.750) (-3.375,-1.700) (-3.358,-1.650) (-3.340,-1.600) (-3.321,-1.550) (-3.300,-1.501) (-3.277,-1.450) (-3.254,-1.400) (-3.229,-1.350) (-3.202,-1.300) (-3.175,-1.252) (-3.144,-1.200) (-3.112,-1.150) (-3.078,-1.100) (-3.042,-1.050) (-3.004,-1.000) (-2.964,-0.950) (-2.925,-0.904) (-2.876,-0.850) (-2.828,-0.800) (-2.777,-0.750) (-2.725,-0.702) (-2.675,-0.659) (-2.634,-0.625) (-2.575,-0.579) (-2.525,-0.543) (-2.475,-0.508) (-2.425,-0.475) (-2.375,-0.444) (-2.325,-0.415) (-2.275,-0.387) (-2.225,-0.361) (-2.175,-0.336) (-2.125,-0.312) (-2.075,-0.290) (-2.025,-0.269) (-1.975,-0.248) (-1.925,-0.230) (-1.875,-0.212) (-1.825,-0.195) (-1.775,-0.179) (-1.725,-0.164) (-1.676,-0.150) (-1.625,-0.137) (-1.575,-0.124) (-1.525,-0.113) (-1.475,-0.102) (-1.425,-0.092) (-1.375,-0.082) (-1.325,-0.074) (-1.300,-0.070)}
\EQ{(1.300,-0.046) (1.350,-0.052) (1.400,-0.058) (1.450,-0.064) (1.500,-0.071) (1.550,-0.079) (1.600,-0.086) (1.650,-0.095) (1.700,-0.104) (1.750,-0.113) (1.800,-0.123) (1.850,-0.134) (1.900,-0.145) (1.950,-0.157) (2.000,-0.170) (2.050,-0.183) (2.100,-0.197) (2.150,-0.212) (2.200,-0.227) (2.250,-0.243) (2.300,-0.260) (2.350,-0.278) (2.400,-0.297) (2.450,-0.317) (2.500,-0.337) (2.550,-0.359) (2.600,-0.381) (2.650,-0.405) (2.700,-0.429) (2.750,-0.455) (2.800,-0.482) (2.850,-0.510) (2.900,-0.539) (2.950,-0.570) (2.998,-0.600) (3.050,-0.635) (3.100,-0.670) (3.141,-0.700) (3.200,-0.745) (3.250,-0.785) (3.298,-0.825) (3.350,-0.871) (3.400,-0.917) (3.450,-0.966) (3.500,-1.017) (3.550,-1.071) (3.598,-1.125) (3.640,-1.175) (3.675,-1.219) (3.718,-1.275) (3.750,-1.319) (3.789,-1.375) (3.822,-1.425) (3.850,-1.470) (3.883,-1.525) (3.912,-1.575) (3.939,-1.625) (3.965,-1.675) (3.990,-1.725) (4.014,-1.775) (4.036,-1.825) (4.057,-1.875) (4.077,-1.925) (4.096,-1.975) (4.114,-2.025) (4.131,-2.075) (4.147,-2.125) (4.162,-2.175) (4.176,-2.225) (4.189,-2.275) (4.201,-2.325) (4.213,-2.375) (4.223,-2.425) (4.232,-2.475) (4.241,-2.525) (4.249,-2.575) (4.256,-2.625) (4.262,-2.675) (4.267,-2.725) (4.271,-2.775) (4.275,-2.825) (4.278,-2.875) (4.280,-2.925) (4.281,-2.975) (4.281,-3.025) (4.281,-3.075) (4.280,-3.125) (4.278,-3.175) (4.275,-3.225) (4.272,-3.275) (4.268,-3.325) (4.263,-3.375) (4.257,-3.425) (4.250,-3.475) (4.243,-3.525) (4.235,-3.575) (4.226,-3.625) (4.217,-3.675) (4.206,-3.725) (4.195,-3.775) (4.184,-3.825) (4.171,-3.875) (4.157,-3.925) (4.143,-3.975) (4.128,-4.025) (4.112,-4.075) (4.096,-4.125) (4.078,-4.175) (4.060,-4.225) (4.041,-4.275) (4.021,-4.325) (4.000,-4.375) (3.979,-4.425) (3.956,-4.475) (3.933,-4.525) (3.908,-4.575) (3.883,-4.625) (3.857,-4.675) (3.830,-4.725) (3.802,-4.775) (3.775,-4.821) (3.743,-4.875) (3.711,-4.925) (3.679,-4.975) (3.650,-5.019) (3.611,-5.075) (3.576,-5.125) (3.539,-5.175) (3.501,-5.225) (3.462,-5.275) (3.425,-5.321) (3.380,-5.375) (3.336,-5.425) (3.300,-5.466) (3.250,-5.520) (3.200,-5.573) (3.150,-5.623) (3.100,-5.672) (3.050,-5.720) (3.000,-5.766) (2.961,-5.800) (2.904,-5.850) (2.850,-5.895) (2.800,-5.935) (2.750,-5.974) (2.700,-6.013) (2.650,-6.049) (2.600,-6.085) (2.550,-6.120) (2.505,-6.150) (2.450,-6.186) (2.400,-6.218) (2.350,-6.249) (2.306,-6.275) (2.250,-6.308) (2.200,-6.336) (2.150,-6.363) (2.100,-6.390) (2.050,-6.415) (2.000,-6.440) (1.950,-6.464) (1.900,-6.488) (1.850,-6.510) (1.800,-6.532) (1.750,-6.553) (1.700,-6.574) (1.650,-6.593) (1.600,-6.612) (1.550,-6.631) (1.500,-6.648) (1.450,-6.665) (1.400,-6.682) (1.350,-6.698) (1.300,-6.713) (1.250,-6.727) (1.200,-6.741) (1.150,-6.754) (1.100,-6.767) (1.050,-6.779) (1.000,-6.791) (0.950,-6.801) (0.900,-6.812) (0.850,-6.821) (0.800,-6.831) (0.750,-6.839) (0.700,-6.847) (0.650,-6.854) (0.600,-6.861) (0.550,-6.868) (0.500,-6.873) (0.450,-6.879) (0.400,-6.883) (0.350,-6.887) (0.300,-6.891) (0.250,-6.894) (0.200,-6.896) (0.150,-6.898) (0.100,-6.900) (0.050,-6.900) (0.000,-6.901) (-0.050,-6.900) (-0.100,-6.900) (-0.150,-6.898) (-0.200,-6.896) (-0.250,-6.894) (-0.300,-6.891) (-0.350,-6.887) (-0.400,-6.883) (-0.450,-6.879) (-0.500,-6.873) (-0.550,-6.868) (-0.600,-6.861) (-0.650,-6.854) (-0.700,-6.847) (-0.750,-6.839) (-0.800,-6.831) (-0.850,-6.821) (-0.900,-6.812) (-0.950,-6.801) (-1.000,-6.791) (-1.050,-6.779) (-1.100,-6.767) (-1.150,-6.754) (-1.200,-6.741) (-1.250,-6.727) (-1.300,-6.713) (-1.350,-6.698) (-1.400,-6.682) (-1.450,-6.665) (-1.500,-6.648) (-1.550,-6.631) (-1.600,-6.612) (-1.650,-6.593) (-1.697,-6.575) (-1.750,-6.553) (-1.800,-6.532) (-1.850,-6.510) (-1.900,-6.488) (-1.950,-6.464) (-2.000,-6.440) (-2.050,-6.415) (-2.100,-6.390) (-2.150,-6.363) (-2.200,-6.336) (-2.250,-6.308) (-2.300,-6.279) (-2.348,-6.250) (-2.400,-6.218) (-2.450,-6.186) (-2.500,-6.154) (-2.543,-6.125) (-2.600,-6.085) (-2.649,-6.050) (-2.700,-6.013) (-2.749,-5.975) (-2.800,-5.935) (-2.844,-5.900) (-2.900,-5.853) (-2.950,-5.810) (-2.990,-5.775) (-3.044,-5.725) (-3.097,-5.675) (-3.148,-5.625) (-3.198,-5.575) (-3.246,-5.525) (-3.292,-5.475) (-3.336,-5.425) (-3.375,-5.380) (-3.421,-5.325) (-3.462,-5.275) (-3.500,-5.227) (-3.539,-5.175) (-3.575,-5.126) (-3.611,-5.075) (-3.646,-5.025) (-3.675,-4.981) (-3.711,-4.925) (-3.743,-4.875) (-3.773,-4.825) (-3.800,-4.778) (-3.825,-4.734) (-3.850,-4.688) (-3.875,-4.641) (-3.900,-4.592) (-3.925,-4.541) (-3.950,-4.488) (-3.975,-4.433) (-4.000,-4.376) (-4.021,-4.325) (-4.041,-4.275) (-4.060,-4.225) (-4.078,-4.175) (-4.096,-4.125) (-4.112,-4.075) (-4.128,-4.025) (-4.143,-3.975) (-4.157,-3.925) (-4.171,-3.875) (-4.184,-3.825) (-4.195,-3.775) (-4.206,-3.725) (-4.217,-3.675) (-4.226,-3.625) (-4.235,-3.575) (-4.243,-3.525) (-4.250,-3.475) (-4.257,-3.425) (-4.263,-3.375) (-4.268,-3.325) (-4.272,-3.275) (-4.275,-3.225) (-4.278,-3.175) (-4.280,-3.125) (-4.281,-3.075) (-4.281,-3.025) (-4.281,-2.975) (-4.280,-2.925) (-4.278,-2.875) (-4.275,-2.825) (-4.271,-2.775) (-4.267,-2.725) (-4.262,-2.675) (-4.256,-2.625) (-4.249,-2.575) (-4.241,-2.525) (-4.232,-2.475) (-4.223,-2.425) (-4.213,-2.375) (-4.201,-2.325) (-4.189,-2.275) (-4.176,-2.225) (-4.162,-2.175) (-4.147,-2.125) (-4.131,-2.075) (-4.114,-2.025) (-4.096,-1.975) (-4.077,-1.925) (-4.057,-1.875) (-4.036,-1.825) (-4.014,-1.775) (-3.990,-1.725) (-3.965,-1.675) (-3.939,-1.625) (-3.912,-1.575) (-3.883,-1.525) (-3.853,-1.475) (-3.825,-1.430) (-3.789,-1.375) (-3.754,-1.325) (-3.718,-1.275) (-3.680,-1.225) (-3.640,-1.175) (-3.600,-1.127) (-3.554,-1.075) (-3.508,-1.025) (-3.459,-0.975) (-3.408,-0.925) (-3.355,-0.875) (-3.300,-0.827) (-3.250,-0.785) (-3.207,-0.750) (-3.150,-0.707) (-3.107,-0.675) (-3.050,-0.635) (-3.000,-0.602) (-2.950,-0.570) (-2.900,-0.539) (-2.850,-0.510) (-2.800,-0.482) (-2.750,-0.455) (-2.700,-0.429) (-2.650,-0.405) (-2.600,-0.381) (-2.550,-0.359) (-2.500,-0.337) (-2.450,-0.317) (-2.400,-0.297) (-2.350,-0.278) (-2.300,-0.260) (-2.250,-0.243) (-2.200,-0.227) (-2.150,-0.212) (-2.100,-0.197) (-2.050,-0.183) (-2.000,-0.170) (-1.950,-0.157) (-1.900,-0.145) (-1.850,-0.134) (-1.800,-0.123) (-1.750,-0.113) (-1.700,-0.104) (-1.650,-0.095) (-1.600,-0.086) (-1.550,-0.079) (-1.500,-0.071) (-1.450,-0.064) (-1.400,-0.058) (-1.350,-0.052) (-1.300,-0.046)}
\EQ{(1.300,-0.023) (1.350,-0.026) (1.400,-0.029) (1.450,-0.032) (1.500,-0.035) (1.550,-0.039) (1.600,-0.043) (1.650,-0.047) (1.700,-0.052) (1.750,-0.056) (1.800,-0.061) (1.850,-0.067) (1.900,-0.072) (1.950,-0.078) (2.000,-0.084) (2.050,-0.091) (2.100,-0.098) (2.150,-0.105) (2.200,-0.112) (2.250,-0.120) (2.300,-0.128) (2.350,-0.137) (2.400,-0.146) (2.450,-0.155) (2.500,-0.165) (2.550,-0.175) (2.600,-0.186) (2.650,-0.197) (2.700,-0.209) (2.750,-0.220) (2.800,-0.233) (2.850,-0.246) (2.900,-0.259) (2.950,-0.273) (3.000,-0.287) (3.050,-0.302) (3.100,-0.318) (3.150,-0.334) (3.199,-0.350) (3.250,-0.367) (3.300,-0.385) (3.350,-0.403) (3.400,-0.422) (3.450,-0.442) (3.500,-0.462) (3.550,-0.483) (3.600,-0.504) (3.646,-0.525) (3.700,-0.550) (3.750,-0.573) (3.800,-0.598) (3.850,-0.623) (3.900,-0.649) (3.949,-0.675) (3.994,-0.700) (4.038,-0.725) (4.100,-0.761) (4.150,-0.792) (4.200,-0.823) (4.250,-0.856) (4.300,-0.889) (4.350,-0.923) (4.400,-0.959) (4.450,-0.996) (4.500,-1.034) (4.550,-1.073) (4.600,-1.113) (4.644,-1.150) (4.700,-1.198) (4.750,-1.243) (4.800,-1.289) (4.838,-1.325) (4.875,-1.362) (4.913,-1.400) (4.950,-1.438) (4.985,-1.475) (5.025,-1.519) (5.074,-1.575) (5.117,-1.625) (5.150,-1.665) (5.198,-1.725) (5.237,-1.775) (5.274,-1.825) (5.310,-1.875) (5.345,-1.925) (5.375,-1.970) (5.411,-2.025) (5.443,-2.075) (5.473,-2.125) (5.500,-2.171) (5.525,-2.214) (5.550,-2.259) (5.575,-2.306) (5.598,-2.350) (5.623,-2.400) (5.648,-2.450) (5.671,-2.500) (5.694,-2.550) (5.716,-2.600) (5.737,-2.650) (5.757,-2.700) (5.777,-2.750) (5.796,-2.800) (5.814,-2.850) (5.831,-2.900) (5.848,-2.950) (5.863,-3.000) (5.879,-3.050) (5.893,-3.100) (5.907,-3.150) (5.920,-3.200) (5.933,-3.250) (5.944,-3.300) (5.956,-3.350) (5.966,-3.400) (5.976,-3.450) (5.985,-3.500) (5.994,-3.550) (6.002,-3.600) (6.010,-3.650) (6.017,-3.700) (6.023,-3.750) (6.029,-3.800) (6.034,-3.850) (6.038,-3.900) (6.042,-3.950) (6.046,-4.000) (6.049,-4.050) (6.051,-4.100) (6.053,-4.150) (6.054,-4.200) (6.054,-4.250) (6.054,-4.300) (6.054,-4.350) (6.053,-4.400) (6.051,-4.450) (6.049,-4.500) (6.047,-4.550) (6.043,-4.600) (6.040,-4.650) (6.035,-4.700) (6.031,-4.750) (6.025,-4.800) (6.019,-4.850) (6.013,-4.900) (6.006,-4.950) (5.998,-5.000) (5.990,-5.050) (5.982,-5.100) (5.973,-5.150) (5.963,-5.200) (5.953,-5.250) (5.942,-5.300) (5.931,-5.350) (5.919,-5.400) (5.906,-5.450) (5.893,-5.500) (5.880,-5.550) (5.866,-5.600) (5.851,-5.650) (5.836,-5.700) (5.820,-5.750) (5.804,-5.800) (5.787,-5.850) (5.769,-5.900) (5.751,-5.950) (5.732,-6.000) (5.713,-6.050) (5.693,-6.100) (5.673,-6.150) (5.652,-6.200) (5.630,-6.250) (5.608,-6.300) (5.585,-6.350) (5.561,-6.400) (5.537,-6.450) (5.512,-6.500) (5.487,-6.550) (5.461,-6.600) (5.434,-6.650) (5.406,-6.700) (5.378,-6.750) (5.350,-6.799) (5.320,-6.850) (5.289,-6.900) (5.258,-6.950) (5.226,-7.000) (5.194,-7.050) (5.160,-7.100) (5.126,-7.150) (5.091,-7.200) (5.056,-7.250) (5.025,-7.292) (4.981,-7.350) (4.950,-7.391) (4.904,-7.450) (4.864,-7.500) (4.825,-7.547) (4.781,-7.600) (4.738,-7.650) (4.700,-7.693) (4.650,-7.749) (4.603,-7.800) (4.555,-7.850) (4.507,-7.900) (4.457,-7.950) (4.407,-8.000) (4.355,-8.050) (4.302,-8.100) (4.250,-8.147) (4.200,-8.192) (4.162,-8.225) (4.104,-8.275) (4.050,-8.320) (4.000,-8.361) (3.952,-8.400) (3.900,-8.441) (3.855,-8.475) (3.800,-8.517) (3.755,-8.550) (3.700,-8.589) (3.650,-8.625) (3.600,-8.659) (3.550,-8.693) (3.501,-8.725) (3.450,-8.758) (3.400,-8.790) (3.350,-8.821) (3.302,-8.850) (3.250,-8.881) (3.200,-8.910) (3.150,-8.939) (3.100,-8.967) (3.050,-8.994) (3.000,-9.021) (2.950,-9.047) (2.900,-9.072) (2.850,-9.097) (2.800,-9.122) (2.750,-9.146) (2.700,-9.169) (2.650,-9.192) (2.600,-9.214) (2.550,-9.236) (2.500,-9.257) (2.450,-9.278) (2.400,-9.298) (2.350,-9.318) (2.300,-9.337) (2.250,-9.356) (2.200,-9.374) (2.150,-9.392) (2.100,-9.410) (2.050,-9.427) (2.000,-9.443) (1.950,-9.459) (1.900,-9.475) (1.850,-9.490) (1.800,-9.504) (1.750,-9.519) (1.700,-9.533) (1.650,-9.546) (1.600,-9.559) (1.550,-9.571) (1.500,-9.583) (1.450,-9.595) (1.400,-9.606) (1.350,-9.617) (1.300,-9.628) (1.250,-9.638) (1.200,-9.647) (1.150,-9.656) (1.100,-9.665) (1.050,-9.674) (1.000,-9.682) (0.950,-9.689) (0.900,-9.696) (0.850,-9.703) (0.800,-9.710) (0.750,-9.716) (0.700,-9.721) (0.650,-9.726) (0.600,-9.731) (0.550,-9.736) (0.500,-9.740) (0.450,-9.743) (0.400,-9.747) (0.350,-9.750) (0.300,-9.752) (0.250,-9.754) (0.200,-9.756) (0.150,-9.757) (0.100,-9.758) (0.050,-9.759) (0.000,-9.759) (-0.050,-9.759) (-0.100,-9.758) (-0.150,-9.757) (-0.200,-9.756) (-0.250,-9.754) (-0.300,-9.752) (-0.350,-9.750) (-0.400,-9.747) (-0.450,-9.743) (-0.500,-9.740) (-0.550,-9.736) (-0.600,-9.731) (-0.650,-9.726) (-0.700,-9.721) (-0.750,-9.716) (-0.800,-9.710) (-0.850,-9.703) (-0.900,-9.696) (-0.950,-9.689) (-1.000,-9.682) (-1.050,-9.674) (-1.100,-9.665) (-1.150,-9.656) (-1.200,-9.647) (-1.250,-9.638) (-1.300,-9.628) (-1.350,-9.617) (-1.400,-9.606) (-1.450,-9.595) (-1.500,-9.583) (-1.550,-9.571) (-1.600,-9.559) (-1.650,-9.546) (-1.700,-9.533) (-1.750,-9.519) (-1.800,-9.504) (-1.850,-9.490) (-1.899,-9.475) (-1.950,-9.459) (-2.000,-9.443) (-2.050,-9.427) (-2.100,-9.410) (-2.150,-9.392) (-2.198,-9.375) (-2.250,-9.356) (-2.300,-9.337) (-2.350,-9.318) (-2.400,-9.298) (-2.450,-9.278) (-2.500,-9.257) (-2.550,-9.236) (-2.600,-9.214) (-2.650,-9.192) (-2.700,-9.169) (-2.750,-9.146) (-2.800,-9.122) (-2.850,-9.097) (-2.900,-9.072) (-2.950,-9.047) (-3.000,-9.021) (-3.050,-8.994) (-3.100,-8.967) (-3.150,-8.939) (-3.200,-8.910) (-3.250,-8.881) (-3.300,-8.851) (-3.344,-8.825) (-3.400,-8.790) (-3.450,-8.758) (-3.500,-8.726) (-3.550,-8.693) (-3.600,-8.659) (-3.650,-8.625) (-3.700,-8.589) (-3.750,-8.553) (-3.800,-8.517) (-3.850,-8.479) (-3.900,-8.441) (-3.950,-8.401) (-4.000,-8.361) (-4.045,-8.325) (-4.100,-8.279) (-4.150,-8.236) (-4.191,-8.200) (-4.247,-8.150) (-4.300,-8.101) (-4.350,-8.055) (-4.400,-8.007) (-4.450,-7.957) (-4.500,-7.907) (-4.550,-7.856) (-4.600,-7.803) (-4.649,-7.750) (-4.694,-7.700) (-4.738,-7.650) (-4.775,-7.607) (-4.823,-7.550) (-4.864,-7.500) (-4.900,-7.455) (-4.943,-7.400) (-4.975,-7.359) (-5.019,-7.300) (-5.050,-7.258) (-5.091,-7.200) (-5.125,-7.152) (-5.160,-7.100) (-5.194,-7.050) (-5.225,-7.002) (-5.258,-6.950) (-5.289,-6.900) (-5.320,-6.850) (-5.349,-6.800) (-5.375,-6.755) (-5.400,-6.711) (-5.425,-6.666) (-5.450,-6.620) (-5.474,-6.575) (-5.500,-6.525) (-5.525,-6.475) (-5.549,-6.425) (-5.573,-6.375) (-5.596,-6.325) (-5.619,-6.275) (-5.641,-6.225) (-5.662,-6.175) (-5.683,-6.125) (-5.703,-6.075) (-5.723,-6.025) (-5.742,-5.975) (-5.760,-5.925) (-5.778,-5.875) (-5.795,-5.825) (-5.812,-5.775) (-5.828,-5.725) (-5.843,-5.675) (-5.858,-5.625) (-5.873,-5.575) (-5.887,-5.525) (-5.900,-5.475) (-5.913,-5.425) (-5.925,-5.375) (-5.936,-5.325) (-5.947,-5.275) (-5.958,-5.225) (-5.968,-5.175) (-5.977,-5.125) (-5.986,-5.075) (-5.994,-5.025) (-6.002,-4.975) (-6.009,-4.925) (-6.016,-4.875) (-6.022,-4.825) (-6.028,-4.775) (-6.033,-4.725) (-6.038,-4.675) (-6.042,-4.625) (-6.045,-4.575) (-6.048,-4.525) (-6.050,-4.475) (-6.052,-4.425) (-6.054,-4.375) (-6.054,-4.325) (-6.055,-4.275) (-6.054,-4.225) (-6.053,-4.175) (-6.052,-4.125) (-6.050,-4.075) (-6.047,-4.025) (-6.044,-3.975) (-6.040,-3.925) (-6.036,-3.875) (-6.031,-3.825) (-6.026,-3.775) (-6.020,-3.725) (-6.013,-3.675) (-6.006,-3.625) (-5.998,-3.575) (-5.990,-3.525) (-5.981,-3.475) (-5.971,-3.425) (-5.961,-3.375) (-5.950,-3.325) (-5.939,-3.275) (-5.926,-3.225) (-5.914,-3.175) (-5.900,-3.125) (-5.886,-3.075) (-5.871,-3.025) (-5.856,-2.975) (-5.839,-2.925) (-5.822,-2.875) (-5.805,-2.825) (-5.786,-2.775) (-5.767,-2.725) (-5.747,-2.675) (-5.727,-2.625) (-5.705,-2.575) (-5.683,-2.525) (-5.660,-2.475) (-5.636,-2.425) (-5.611,-2.375) (-5.585,-2.325) (-5.559,-2.275) (-5.531,-2.225) (-5.503,-2.175) (-5.475,-2.128) (-5.443,-2.075) (-5.411,-2.025) (-5.378,-1.975) (-5.350,-1.933) (-5.310,-1.875) (-5.275,-1.827) (-5.237,-1.775) (-5.200,-1.727) (-5.158,-1.675) (-5.125,-1.634) (-5.075,-1.576) (-5.030,-1.525) (-4.985,-1.475) (-4.950,-1.438) (-4.913,-1.400) (-4.875,-1.362) (-4.838,-1.325) (-4.800,-1.289) (-4.758,-1.250) (-4.702,-1.200) (-4.650,-1.155) (-4.600,-1.113) (-4.553,-1.075) (-4.500,-1.034) (-4.456,-1.000) (-4.400,-0.959) (-4.352,-0.925) (-4.300,-0.889) (-4.250,-0.856) (-4.203,-0.825) (-4.150,-0.792) (-4.100,-0.761) (-4.050,-0.732) (-4.000,-0.703) (-3.950,-0.676) (-3.902,-0.650) (-3.854,-0.625) (-3.805,-0.600) (-3.754,-0.575) (-3.701,-0.550) (-3.650,-0.527) (-3.600,-0.504) (-3.550,-0.483) (-3.500,-0.462) (-3.450,-0.442) (-3.400,-0.422) (-3.350,-0.403) (-3.300,-0.385) (-3.250,-0.367) (-3.200,-0.350) (-3.150,-0.334) (-3.100,-0.318) (-3.050,-0.302) (-3.000,-0.287) (-2.950,-0.273) (-2.900,-0.259) (-2.850,-0.246) (-2.800,-0.233) (-2.750,-0.220) (-2.700,-0.209) (-2.650,-0.197) (-2.600,-0.186) (-2.550,-0.175) (-2.500,-0.165) (-2.450,-0.155) (-2.400,-0.146) (-2.350,-0.137) (-2.300,-0.128) (-2.250,-0.120) (-2.200,-0.112) (-2.150,-0.105) (-2.100,-0.098) (-2.050,-0.091) (-2.000,-0.084) (-1.950,-0.078) (-1.900,-0.072) (-1.850,-0.067) (-1.800,-0.061) (-1.750,-0.056) (-1.700,-0.052) (-1.650,-0.047) (-1.600,-0.043) (-1.550,-0.039) (-1.500,-0.035) (-1.450,-0.032) (-1.400,-0.029) (-1.350,-0.026) (-1.302,-0.023)}
\EQ{(1.300,0.023) (1.350,0.026) (1.400,0.029) (1.450,0.032) (1.500,0.035) (1.550,0.039) (1.600,0.043) (1.650,0.047) (1.700,0.052) (1.750,0.056) (1.800,0.061) (1.850,0.067) (1.900,0.072) (1.950,0.078) (2.000,0.084) (2.050,0.091) (2.100,0.098) (2.150,0.105) (2.200,0.112) (2.250,0.120) (2.300,0.128) (2.350,0.137) (2.400,0.146) (2.450,0.155) (2.500,0.165) (2.550,0.175) (2.600,0.186) (2.650,0.197) (2.700,0.209) (2.750,0.220) (2.800,0.233) (2.850,0.246) (2.900,0.259) (2.950,0.273) (3.000,0.287) (3.050,0.302) (3.100,0.318) (3.150,0.334) (3.199,0.350) (3.250,0.367) (3.300,0.385) (3.350,0.403) (3.400,0.422) (3.450,0.442) (3.500,0.462) (3.550,0.483) (3.600,0.504) (3.646,0.525) (3.700,0.550) (3.750,0.573) (3.800,0.598) (3.850,0.623) (3.900,0.649) (3.949,0.675) (3.994,0.700) (4.038,0.725) (4.100,0.761) (4.150,0.792) (4.200,0.823) (4.250,0.856) (4.300,0.889) (4.350,0.923) (4.400,0.959) (4.450,0.996) (4.500,1.034) (4.550,1.073) (4.600,1.113) (4.644,1.150) (4.700,1.198) (4.750,1.243) (4.800,1.289) (4.838,1.325) (4.875,1.362) (4.913,1.400) (4.950,1.438) (4.985,1.475) (5.025,1.519) (5.074,1.575) (5.117,1.625) (5.150,1.665) (5.198,1.725) (5.237,1.775) (5.274,1.825) (5.310,1.875) (5.345,1.925) (5.375,1.970) (5.411,2.025) (5.443,2.075) (5.473,2.125) (5.500,2.171) (5.525,2.214) (5.550,2.259) (5.575,2.306) (5.598,2.350) (5.623,2.400) (5.648,2.450) (5.671,2.500) (5.694,2.550) (5.716,2.600) (5.737,2.650) (5.757,2.700) (5.777,2.750) (5.796,2.800) (5.814,2.850) (5.831,2.900) (5.848,2.950) (5.863,3.000) (5.879,3.050) (5.893,3.100) (5.907,3.150) (5.920,3.200) (5.933,3.250) (5.944,3.300) (5.956,3.350) (5.966,3.400) (5.976,3.450) (5.985,3.500) (5.994,3.550) (6.002,3.600) (6.010,3.650) (6.017,3.700) (6.023,3.750) (6.029,3.800) (6.034,3.850) (6.038,3.900) (6.042,3.950) (6.046,4.000) (6.049,4.050) (6.051,4.100) (6.053,4.150) (6.054,4.200) (6.054,4.250) (6.054,4.300) (6.054,4.350) (6.053,4.400) (6.051,4.450) (6.049,4.500) (6.047,4.550) (6.043,4.600) (6.040,4.650) (6.035,4.700) (6.031,4.750) (6.025,4.800) (6.019,4.850) (6.013,4.900) (6.006,4.950) (5.998,5.000) (5.990,5.050) (5.982,5.100) (5.973,5.150) (5.963,5.200) (5.953,5.250) (5.942,5.300) (5.931,5.350) (5.919,5.400) (5.906,5.450) (5.893,5.500) (5.880,5.550) (5.866,5.600) (5.851,5.650) (5.836,5.700) (5.820,5.750) (5.804,5.800) (5.787,5.850) (5.769,5.900) (5.751,5.950) (5.732,6.000) (5.713,6.050) (5.693,6.100) (5.673,6.150) (5.652,6.200) (5.630,6.250) (5.608,6.300) (5.585,6.350) (5.561,6.400) (5.537,6.450) (5.512,6.500) (5.487,6.550) (5.461,6.600) (5.434,6.650) (5.406,6.700) (5.378,6.750) (5.350,6.799) (5.320,6.850) (5.289,6.900) (5.258,6.950) (5.226,7.000) (5.194,7.050) (5.160,7.100) (5.126,7.150) (5.091,7.200) (5.056,7.250) (5.025,7.292) (4.981,7.350) (4.950,7.391) (4.904,7.450) (4.864,7.500) (4.825,7.547) (4.781,7.600) (4.738,7.650) (4.700,7.693) (4.650,7.749) (4.603,7.800) (4.555,7.850) (4.507,7.900) (4.457,7.950) (4.407,8.000) (4.355,8.050) (4.302,8.100) (4.250,8.147) (4.200,8.192) (4.162,8.225) (4.104,8.275) (4.050,8.320) (4.000,8.361) (3.952,8.400) (3.900,8.441) (3.855,8.475) (3.800,8.517) (3.755,8.550) (3.700,8.589) (3.650,8.625) (3.600,8.659) (3.550,8.693) (3.501,8.725) (3.450,8.758) (3.400,8.790) (3.350,8.821) (3.302,8.850) (3.250,8.881) (3.200,8.910) (3.150,8.939) (3.100,8.967) (3.050,8.994) (3.000,9.021) (2.950,9.047) (2.900,9.072) (2.850,9.097) (2.800,9.122) (2.750,9.146) (2.700,9.169) (2.650,9.192) (2.600,9.214) (2.550,9.236) (2.500,9.257) (2.450,9.278) (2.400,9.298) (2.350,9.318) (2.300,9.337) (2.250,9.356) (2.200,9.374) (2.150,9.392) (2.100,9.410) (2.050,9.427) (2.000,9.443) (1.950,9.459) (1.900,9.475) (1.850,9.490) (1.800,9.504) (1.750,9.519) (1.700,9.533) (1.650,9.546) (1.600,9.559) (1.550,9.571) (1.500,9.583) (1.450,9.595) (1.400,9.606) (1.350,9.617) (1.300,9.628) (1.250,9.638) (1.200,9.647) (1.150,9.656) (1.100,9.665) (1.050,9.674) (1.000,9.682) (0.950,9.689) (0.900,9.696) (0.850,9.703) (0.800,9.710) (0.750,9.716) (0.700,9.721) (0.650,9.726) (0.600,9.731) (0.550,9.736) (0.500,9.740) (0.450,9.743) (0.400,9.747) (0.350,9.750) (0.300,9.752) (0.250,9.754) (0.200,9.756) (0.150,9.757) (0.100,9.758) (0.050,9.759) (0.000,9.759) (-0.050,9.759) (-0.100,9.758) (-0.150,9.757) (-0.200,9.756) (-0.250,9.754) (-0.300,9.752) (-0.350,9.750) (-0.400,9.747) (-0.450,9.743) (-0.500,9.740) (-0.550,9.736) (-0.600,9.731) (-0.650,9.726) (-0.700,9.721) (-0.750,9.716) (-0.800,9.710) (-0.850,9.703) (-0.900,9.696) (-0.950,9.689) (-1.000,9.682) (-1.050,9.674) (-1.100,9.665) (-1.150,9.656) (-1.200,9.647) (-1.250,9.638) (-1.300,9.628) (-1.350,9.617) (-1.400,9.606) (-1.450,9.595) (-1.500,9.583) (-1.550,9.571) (-1.600,9.559) (-1.650,9.546) (-1.700,9.533) (-1.750,9.519) (-1.800,9.504) (-1.850,9.490) (-1.899,9.475) (-1.950,9.459) (-2.000,9.443) (-2.050,9.427) (-2.100,9.410) (-2.150,9.392) (-2.198,9.375) (-2.250,9.356) (-2.300,9.337) (-2.350,9.318) (-2.400,9.298) (-2.450,9.278) (-2.500,9.257) (-2.550,9.236) (-2.600,9.214) (-2.650,9.192) (-2.700,9.169) (-2.750,9.146) (-2.800,9.122) (-2.850,9.097) (-2.900,9.072) (-2.950,9.047) (-3.000,9.021) (-3.050,8.994) (-3.100,8.967) (-3.150,8.939) (-3.200,8.910) (-3.250,8.881) (-3.300,8.851) (-3.344,8.825) (-3.400,8.790) (-3.450,8.758) (-3.500,8.726) (-3.550,8.693) (-3.600,8.659) (-3.650,8.625) (-3.700,8.589) (-3.750,8.553) (-3.800,8.517) (-3.850,8.479) (-3.900,8.441) (-3.950,8.401) (-4.000,8.361) (-4.045,8.325) (-4.100,8.279) (-4.150,8.236) (-4.191,8.200) (-4.247,8.150) (-4.300,8.101) (-4.350,8.055) (-4.400,8.007) (-4.450,7.957) (-4.500,7.907) (-4.550,7.856) (-4.600,7.803) (-4.649,7.750) (-4.694,7.700) (-4.738,7.650) (-4.775,7.607) (-4.823,7.550) (-4.864,7.500) (-4.900,7.455) (-4.943,7.400) (-4.975,7.359) (-5.019,7.300) (-5.050,7.258) (-5.091,7.200) (-5.125,7.152) (-5.160,7.100) (-5.194,7.050) (-5.225,7.002) (-5.258,6.950) (-5.289,6.900) (-5.320,6.850) (-5.349,6.800) (-5.375,6.755) (-5.400,6.711) (-5.425,6.666) (-5.450,6.620) (-5.474,6.575) (-5.500,6.525) (-5.525,6.475) (-5.549,6.425) (-5.573,6.375) (-5.596,6.325) (-5.619,6.275) (-5.641,6.225) (-5.662,6.175) (-5.683,6.125) (-5.703,6.075) (-5.723,6.025) (-5.742,5.975) (-5.760,5.925) (-5.778,5.875) (-5.795,5.825) (-5.812,5.775) (-5.828,5.725) (-5.843,5.675) (-5.858,5.625) (-5.873,5.575) (-5.887,5.525) (-5.900,5.475) (-5.913,5.425) (-5.925,5.375) (-5.936,5.325) (-5.947,5.275) (-5.958,5.225) (-5.968,5.175) (-5.977,5.125) (-5.986,5.075) (-5.994,5.025) (-6.002,4.975) (-6.009,4.925) (-6.016,4.875) (-6.022,4.825) (-6.028,4.775) (-6.033,4.725) (-6.038,4.675) (-6.042,4.625) (-6.045,4.575) (-6.048,4.525) (-6.050,4.475) (-6.052,4.425) (-6.054,4.375) (-6.054,4.325) (-6.055,4.275) (-6.054,4.225) (-6.053,4.175) (-6.052,4.125) (-6.050,4.075) (-6.047,4.025) (-6.044,3.975) (-6.040,3.925) (-6.036,3.875) (-6.031,3.825) (-6.026,3.775) (-6.020,3.725) (-6.013,3.675) (-6.006,3.625) (-5.998,3.575) (-5.990,3.525) (-5.981,3.475) (-5.971,3.425) (-5.961,3.375) (-5.950,3.325) (-5.939,3.275) (-5.926,3.225) (-5.914,3.175) (-5.900,3.125) (-5.886,3.075) (-5.871,3.025) (-5.856,2.975) (-5.839,2.925) (-5.822,2.875) (-5.805,2.825) (-5.786,2.775) (-5.767,2.725) (-5.747,2.675) (-5.727,2.625) (-5.705,2.575) (-5.683,2.525) (-5.660,2.475) (-5.636,2.425) (-5.611,2.375) (-5.585,2.325) (-5.559,2.275) (-5.531,2.225) (-5.503,2.175) (-5.475,2.128) (-5.443,2.075) (-5.411,2.025) (-5.378,1.975) (-5.350,1.933) (-5.310,1.875) (-5.275,1.827) (-5.237,1.775) (-5.200,1.727) (-5.158,1.675) (-5.125,1.634) (-5.075,1.576) (-5.030,1.525) (-4.985,1.475) (-4.950,1.438) (-4.913,1.400) (-4.875,1.362) (-4.838,1.325) (-4.800,1.289) (-4.758,1.250) (-4.702,1.200) (-4.650,1.155) (-4.600,1.113) (-4.553,1.075) (-4.500,1.034) (-4.456,1.000) (-4.400,0.959) (-4.352,0.925) (-4.300,0.889) (-4.250,0.856) (-4.203,0.825) (-4.150,0.792) (-4.100,0.761) (-4.050,0.732) (-4.000,0.703) (-3.950,0.676) (-3.902,0.650) (-3.854,0.625) (-3.805,0.600) (-3.754,0.575) (-3.701,0.550) (-3.650,0.527) (-3.600,0.504) (-3.550,0.483) (-3.500,0.462) (-3.450,0.442) (-3.400,0.422) (-3.350,0.403) (-3.300,0.385) (-3.250,0.367) (-3.200,0.350) (-3.150,0.334) (-3.100,0.318) (-3.050,0.302) (-3.000,0.287) (-2.950,0.273) (-2.900,0.259) (-2.850,0.246) (-2.800,0.233) (-2.750,0.220) (-2.700,0.209) (-2.650,0.197) (-2.600,0.186) (-2.550,0.175) (-2.500,0.165) (-2.450,0.155) (-2.400,0.146) (-2.350,0.137) (-2.300,0.128) (-2.250,0.120) (-2.200,0.112) (-2.150,0.105) (-2.100,0.098) (-2.050,0.091) (-2.000,0.084) (-1.950,0.078) (-1.900,0.072) (-1.850,0.067) (-1.800,0.061) (-1.750,0.056) (-1.700,0.052) (-1.650,0.047) (-1.600,0.043) (-1.550,0.039) (-1.500,0.035) (-1.450,0.032) (-1.400,0.029) (-1.350,0.026) (-1.302,0.023)}
\EQ{(1.300,0.046) (1.350,0.052) (1.400,0.058) (1.450,0.064) (1.500,0.071) (1.550,0.079) (1.600,0.086) (1.650,0.095) (1.700,0.104) (1.750,0.113) (1.800,0.123) (1.850,0.134) (1.900,0.145) (1.950,0.157) (2.000,0.170) (2.050,0.183) (2.100,0.197) (2.150,0.212) (2.200,0.227) (2.250,0.243) (2.300,0.260) (2.350,0.278) (2.400,0.297) (2.450,0.317) (2.500,0.337) (2.550,0.359) (2.600,0.381) (2.650,0.405) (2.700,0.429) (2.750,0.455) (2.800,0.482) (2.850,0.510) (2.900,0.539) (2.950,0.570) (2.998,0.600) (3.050,0.635) (3.100,0.670) (3.141,0.700) (3.200,0.745) (3.250,0.785) (3.298,0.825) (3.350,0.871) (3.400,0.917) (3.450,0.966) (3.500,1.017) (3.550,1.071) (3.598,1.125) (3.640,1.175) (3.675,1.219) (3.718,1.275) (3.750,1.319) (3.789,1.375) (3.822,1.425) (3.850,1.470) (3.883,1.525) (3.912,1.575) (3.939,1.625) (3.965,1.675) (3.990,1.725) (4.014,1.775) (4.036,1.825) (4.057,1.875) (4.077,1.925) (4.096,1.975) (4.114,2.025) (4.131,2.075) (4.147,2.125) (4.162,2.175) (4.176,2.225) (4.189,2.275) (4.201,2.325) (4.213,2.375) (4.223,2.425) (4.232,2.475) (4.241,2.525) (4.249,2.575) (4.256,2.625) (4.262,2.675) (4.267,2.725) (4.271,2.775) (4.275,2.825) (4.278,2.875) (4.280,2.925) (4.281,2.975) (4.281,3.025) (4.281,3.075) (4.280,3.125) (4.278,3.175) (4.275,3.225) (4.272,3.275) (4.268,3.325) (4.263,3.375) (4.257,3.425) (4.250,3.475) (4.243,3.525) (4.235,3.575) (4.226,3.625) (4.217,3.675) (4.206,3.725) (4.195,3.775) (4.184,3.825) (4.171,3.875) (4.157,3.925) (4.143,3.975) (4.128,4.025) (4.112,4.075) (4.096,4.125) (4.078,4.175) (4.060,4.225) (4.041,4.275) (4.021,4.325) (4.000,4.375) (3.979,4.425) (3.956,4.475) (3.933,4.525) (3.908,4.575) (3.883,4.625) (3.857,4.675) (3.830,4.725) (3.802,4.775) (3.775,4.821) (3.743,4.875) (3.711,4.925) (3.679,4.975) (3.650,5.019) (3.611,5.075) (3.576,5.125) (3.539,5.175) (3.501,5.225) (3.462,5.275) (3.425,5.321) (3.380,5.375) (3.336,5.425) (3.300,5.466) (3.250,5.520) (3.200,5.573) (3.150,5.623) (3.100,5.672) (3.050,5.720) (3.000,5.766) (2.961,5.800) (2.904,5.850) (2.850,5.895) (2.800,5.935) (2.750,5.974) (2.700,6.013) (2.650,6.049) (2.600,6.085) (2.550,6.120) (2.505,6.150) (2.450,6.186) (2.400,6.218) (2.350,6.249) (2.306,6.275) (2.250,6.308) (2.200,6.336) (2.150,6.363) (2.100,6.390) (2.050,6.415) (2.000,6.440) (1.950,6.464) (1.900,6.488) (1.850,6.510) (1.800,6.532) (1.750,6.553) (1.700,6.574) (1.650,6.593) (1.600,6.612) (1.550,6.631) (1.500,6.648) (1.450,6.665) (1.400,6.682) (1.350,6.698) (1.300,6.713) (1.250,6.727) (1.200,6.741) (1.150,6.754) (1.100,6.767) (1.050,6.779) (1.000,6.791) (0.950,6.801) (0.900,6.812) (0.850,6.821) (0.800,6.831) (0.750,6.839) (0.700,6.847) (0.650,6.854) (0.600,6.861) (0.550,6.868) (0.500,6.873) (0.450,6.879) (0.400,6.883) (0.350,6.887) (0.300,6.891) (0.250,6.894) (0.200,6.896) (0.150,6.898) (0.100,6.900) (0.050,6.900) (0.000,6.901) (-0.050,6.900) (-0.100,6.900) (-0.150,6.898) (-0.200,6.896) (-0.250,6.894) (-0.300,6.891) (-0.350,6.887) (-0.400,6.883) (-0.450,6.879) (-0.500,6.873) (-0.550,6.868) (-0.600,6.861) (-0.650,6.854) (-0.700,6.847) (-0.750,6.839) (-0.800,6.831) (-0.850,6.821) (-0.900,6.812) (-0.950,6.801) (-1.000,6.791) (-1.050,6.779) (-1.100,6.767) (-1.150,6.754) (-1.200,6.741) (-1.250,6.727) (-1.300,6.713) (-1.350,6.698) (-1.400,6.682) (-1.450,6.665) (-1.500,6.648) (-1.550,6.631) (-1.600,6.612) (-1.650,6.593) (-1.697,6.575) (-1.750,6.553) (-1.800,6.532) (-1.850,6.510) (-1.900,6.488) (-1.950,6.464) (-2.000,6.440) (-2.050,6.415) (-2.100,6.390) (-2.150,6.363) (-2.200,6.336) (-2.250,6.308) (-2.300,6.279) (-2.348,6.250) (-2.400,6.218) (-2.450,6.186) (-2.500,6.154) (-2.543,6.125) (-2.600,6.085) (-2.649,6.050) (-2.700,6.013) (-2.749,5.975) (-2.800,5.935) (-2.844,5.900) (-2.900,5.853) (-2.950,5.810) (-2.990,5.775) (-3.044,5.725) (-3.097,5.675) (-3.148,5.625) (-3.198,5.575) (-3.246,5.525) (-3.292,5.475) (-3.336,5.425) (-3.375,5.380) (-3.421,5.325) (-3.462,5.275) (-3.500,5.227) (-3.539,5.175) (-3.575,5.126) (-3.611,5.075) (-3.646,5.025) (-3.675,4.981) (-3.711,4.925) (-3.743,4.875) (-3.773,4.825) (-3.800,4.778) (-3.825,4.734) (-3.850,4.688) (-3.875,4.641) (-3.900,4.592) (-3.925,4.541) (-3.950,4.488) (-3.975,4.433) (-4.000,4.376) (-4.021,4.325) (-4.041,4.275) (-4.060,4.225) (-4.078,4.175) (-4.096,4.125) (-4.112,4.075) (-4.128,4.025) (-4.143,3.975) (-4.157,3.925) (-4.171,3.875) (-4.184,3.825) (-4.195,3.775) (-4.206,3.725) (-4.217,3.675) (-4.226,3.625) (-4.235,3.575) (-4.243,3.525) (-4.250,3.475) (-4.257,3.425) (-4.263,3.375) (-4.268,3.325) (-4.272,3.275) (-4.275,3.225) (-4.278,3.175) (-4.280,3.125) (-4.281,3.075) (-4.281,3.025) (-4.281,2.975) (-4.280,2.925) (-4.278,2.875) (-4.275,2.825) (-4.271,2.775) (-4.267,2.725) (-4.262,2.675) (-4.256,2.625) (-4.249,2.575) (-4.241,2.525) (-4.232,2.475) (-4.223,2.425) (-4.213,2.375) (-4.201,2.325) (-4.189,2.275) (-4.176,2.225) (-4.162,2.175) (-4.147,2.125) (-4.131,2.075) (-4.114,2.025) (-4.096,1.975) (-4.077,1.925) (-4.057,1.875) (-4.036,1.825) (-4.014,1.775) (-3.990,1.725) (-3.965,1.675) (-3.939,1.625) (-3.912,1.575) (-3.883,1.525) (-3.853,1.475) (-3.825,1.430) (-3.789,1.375) (-3.754,1.325) (-3.718,1.275) (-3.680,1.225) (-3.640,1.175) (-3.600,1.127) (-3.554,1.075) (-3.508,1.025) (-3.459,0.975) (-3.408,0.925) (-3.355,0.875) (-3.300,0.827) (-3.250,0.785) (-3.207,0.750) (-3.150,0.707) (-3.107,0.675) (-3.050,0.635) (-3.000,0.602) (-2.950,0.570) (-2.900,0.539) (-2.850,0.510) (-2.800,0.482) (-2.750,0.455) (-2.700,0.429) (-2.650,0.405) (-2.600,0.381) (-2.550,0.359) (-2.500,0.337) (-2.450,0.317) (-2.400,0.297) (-2.350,0.278) (-2.300,0.260) (-2.250,0.243) (-2.200,0.227) (-2.150,0.212) (-2.100,0.197) (-2.050,0.183) (-2.000,0.170) (-1.950,0.157) (-1.900,0.145) (-1.850,0.134) (-1.800,0.123) (-1.750,0.113) (-1.700,0.104) (-1.650,0.095) (-1.600,0.086) (-1.550,0.079) (-1.500,0.071) (-1.450,0.064) (-1.400,0.058) (-1.350,0.052) (-1.300,0.046)}
\EQ{(1.300,0.070) (1.350,0.078) (1.400,0.087) (1.450,0.097) (1.500,0.107) (1.550,0.118) (1.600,0.130) (1.650,0.143) (1.700,0.157) (1.750,0.171) (1.800,0.187) (1.850,0.203) (1.900,0.220) (1.950,0.239) (2.000,0.258) (2.050,0.279) (2.098,0.300) (2.150,0.324) (2.200,0.348) (2.250,0.374) (2.298,0.400) (2.342,0.425) (2.400,0.460) (2.450,0.492) (2.500,0.525) (2.550,0.561) (2.600,0.598) (2.650,0.638) (2.694,0.675) (2.750,0.724) (2.800,0.772) (2.850,0.822) (2.899,0.875) (2.943,0.925) (2.985,0.975) (3.024,1.025) (3.060,1.075) (3.095,1.125) (3.125,1.171) (3.159,1.225) (3.188,1.275) (3.216,1.325) (3.241,1.375) (3.266,1.425) (3.289,1.475) (3.310,1.525) (3.330,1.575) (3.349,1.625) (3.367,1.675) (3.383,1.725) (3.398,1.775) (3.412,1.825) (3.425,1.875) (3.437,1.925) (3.447,1.975) (3.456,2.025) (3.465,2.075) (3.472,2.125) (3.479,2.175) (3.484,2.225) (3.488,2.275) (3.491,2.325) (3.494,2.375) (3.495,2.425) (3.496,2.475) (3.495,2.525) (3.494,2.575) (3.491,2.625) (3.488,2.675) (3.483,2.725) (3.478,2.775) (3.472,2.825) (3.465,2.875) (3.457,2.925) (3.448,2.975) (3.438,3.025) (3.427,3.075) (3.415,3.125) (3.403,3.175) (3.389,3.225) (3.374,3.275) (3.359,3.325) (3.342,3.375) (3.325,3.423) (3.306,3.475) (3.286,3.525) (3.265,3.575) (3.243,3.625) (3.220,3.675) (3.196,3.725) (3.171,3.775) (3.145,3.825) (3.118,3.875) (3.089,3.925) (3.059,3.975) (3.028,4.025) (3.000,4.068) (2.962,4.125) (2.927,4.175) (2.890,4.225) (2.852,4.275) (2.812,4.325) (2.775,4.370) (2.728,4.425) (2.683,4.475) (2.637,4.525) (2.600,4.563) (2.563,4.600) (2.525,4.637) (2.485,4.675) (2.430,4.725) (2.375,4.773) (2.325,4.815) (2.281,4.850) (2.225,4.894) (2.183,4.925) (2.125,4.967) (2.078,5.000) (2.025,5.035) (1.975,5.067) (1.925,5.098) (1.881,5.125) (1.825,5.157) (1.775,5.185) (1.725,5.212) (1.675,5.238) (1.625,5.262) (1.575,5.286) (1.525,5.309) (1.475,5.331) (1.425,5.352) (1.375,5.373) (1.325,5.392) (1.275,5.411) (1.225,5.428) (1.175,5.445) (1.125,5.462) (1.075,5.477) (1.025,5.492) (0.975,5.505) (0.925,5.519) (0.875,5.531) (0.825,5.543) (0.775,5.553) (0.725,5.564) (0.675,5.573) (0.625,5.582) (0.575,5.590) (0.525,5.597) (0.475,5.604) (0.425,5.610) (0.375,5.616) (0.325,5.620) (0.275,5.624) (0.225,5.628) (0.175,5.630) (0.125,5.632) (0.075,5.634) (0.025,5.634) (-0.050,5.634) (-0.100,5.633) (-0.150,5.631) (-0.200,5.629) (-0.250,5.626) (-0.300,5.622) (-0.350,5.618) (-0.400,5.613) (-0.450,5.607) (-0.500,5.601) (-0.550,5.594) (-0.600,5.586) (-0.650,5.578) (-0.700,5.569) (-0.750,5.559) (-0.800,5.548) (-0.850,5.537) (-0.899,5.525) (-0.950,5.512) (-1.000,5.499) (-1.050,5.484) (-1.100,5.469) (-1.150,5.453) (-1.200,5.437) (-1.250,5.420) (-1.300,5.401) (-1.350,5.382) (-1.400,5.363) (-1.450,5.342) (-1.500,5.320) (-1.550,5.298) (-1.599,5.275) (-1.650,5.250) (-1.700,5.225) (-1.747,5.200) (-1.793,5.175) (-1.838,5.150) (-1.900,5.114) (-1.950,5.083) (-2.000,5.051) (-2.050,5.019) (-2.100,4.985) (-2.149,4.950) (-2.200,4.913) (-2.249,4.875) (-2.300,4.835) (-2.343,4.800) (-2.400,4.751) (-2.450,4.707) (-2.500,4.661) (-2.538,4.625) (-2.575,4.588) (-2.613,4.550) (-2.650,4.511) (-2.700,4.457) (-2.750,4.400) (-2.792,4.350) (-2.825,4.309) (-2.871,4.250) (-2.908,4.200) (-2.944,4.150) (-2.975,4.106) (-3.012,4.050) (-3.044,4.000) (-3.074,3.950) (-3.100,3.906) (-3.125,3.862) (-3.150,3.816) (-3.175,3.768) (-3.200,3.718) (-3.225,3.665) (-3.250,3.610) (-3.275,3.551) (-3.296,3.500) (-3.315,3.450) (-3.333,3.400) (-3.350,3.351) (-3.366,3.300) (-3.382,3.250) (-3.396,3.200) (-3.409,3.150) (-3.421,3.100) (-3.433,3.050) (-3.443,3.000) (-3.453,2.950) (-3.461,2.900) (-3.469,2.850) (-3.475,2.800) (-3.481,2.750) (-3.486,2.700) (-3.490,2.650) (-3.492,2.600) (-3.494,2.550) (-3.495,2.500) (-3.496,2.450) (-3.495,2.400) (-3.493,2.350) (-3.490,2.300) (-3.486,2.250) (-3.481,2.200) (-3.475,2.150) (-3.469,2.100) (-3.461,2.050) (-3.452,2.000) (-3.442,1.950) (-3.431,1.900) (-3.419,1.850) (-3.405,1.800) (-3.391,1.750) (-3.375,1.700) (-3.358,1.650) (-3.340,1.600) (-3.321,1.550) (-3.300,1.501) (-3.277,1.450) (-3.254,1.400) (-3.229,1.350) (-3.202,1.300) (-3.175,1.252) (-3.144,1.200) (-3.112,1.150) (-3.078,1.100) (-3.042,1.050) (-3.004,1.000) (-2.964,0.950) (-2.925,0.904) (-2.876,0.850) (-2.828,0.800) (-2.777,0.750) (-2.725,0.702) (-2.675,0.659) (-2.634,0.625) (-2.575,0.579) (-2.525,0.543) (-2.475,0.508) (-2.425,0.475) (-2.375,0.444) (-2.325,0.415) (-2.275,0.387) (-2.225,0.361) (-2.175,0.336) (-2.125,0.312) (-2.075,0.290) (-2.025,0.269) (-1.975,0.248) (-1.925,0.230) (-1.875,0.212) (-1.825,0.195) (-1.775,0.179) (-1.725,0.164) (-1.676,0.150) (-1.625,0.137) (-1.575,0.124) (-1.525,0.113) (-1.475,0.102) (-1.425,0.092) (-1.375,0.082) (-1.325,0.074) (-1.300,0.070)}
\EQ{(1.300,0.093) (1.350,0.104) (1.400,0.116) (1.450,0.130) (1.500,0.144) (1.550,0.159) (1.600,0.175) (1.650,0.193) (1.700,0.211) (1.750,0.231) (1.800,0.252) (1.850,0.275) (1.900,0.299) (1.950,0.324) (1.997,0.350) (2.041,0.375) (2.100,0.412) (2.150,0.444) (2.194,0.475) (2.250,0.517) (2.292,0.550) (2.350,0.599) (2.400,0.645) (2.450,0.693) (2.500,0.746) (2.548,0.800) (2.589,0.850) (2.625,0.896) (2.664,0.950) (2.697,1.000) (2.725,1.044) (2.759,1.100) (2.786,1.150) (2.812,1.200) (2.836,1.250) (2.858,1.300) (2.879,1.350) (2.898,1.400) (2.916,1.450) (2.932,1.500) (2.947,1.550) (2.960,1.600) (2.972,1.650) (2.983,1.700) (2.993,1.750) (3.001,1.800) (3.008,1.850) (3.014,1.900) (3.019,1.950) (3.023,2.000) (3.025,2.050) (3.027,2.100) (3.027,2.150) (3.026,2.200) (3.025,2.250) (3.022,2.300) (3.018,2.350) (3.013,2.400) (3.006,2.450) (2.999,2.500) (2.991,2.550) (2.981,2.600) (2.971,2.650) (2.959,2.700) (2.947,2.750) (2.933,2.800) (2.918,2.850) (2.902,2.900) (2.885,2.950) (2.866,3.000) (2.847,3.050) (2.826,3.100) (2.804,3.150) (2.781,3.200) (2.756,3.250) (2.730,3.300) (2.703,3.350) (2.675,3.399) (2.645,3.450) (2.613,3.500) (2.580,3.550) (2.550,3.594) (2.509,3.650) (2.475,3.696) (2.432,3.750) (2.400,3.789) (2.350,3.846) (2.301,3.900) (2.253,3.950) (2.203,4.000) (2.151,4.050) (2.100,4.096) (2.050,4.139) (2.007,4.175) (1.950,4.220) (1.900,4.258) (1.850,4.295) (1.807,4.325) (1.750,4.363) (1.700,4.395) (1.651,4.425) (1.600,4.455) (1.550,4.483) (1.500,4.510) (1.450,4.536) (1.400,4.561) (1.350,4.585) (1.300,4.607) (1.250,4.629) (1.200,4.649) (1.150,4.669) (1.100,4.687) (1.050,4.705) (1.000,4.722) (0.950,4.737) (0.900,4.752) (0.850,4.766) (0.800,4.779) (0.750,4.792) (0.700,4.803) (0.650,4.814) (0.600,4.824) (0.550,4.833) (0.500,4.841) (0.450,4.848) (0.400,4.855) (0.350,4.861) (0.300,4.866) (0.250,4.870) (0.200,4.873) (0.150,4.876) (0.100,4.878) (0.050,4.879) (0.000,4.880) (-0.050,4.879) (-0.100,4.878) (-0.150,4.876) (-0.200,4.873) (-0.250,4.870) (-0.300,4.866) (-0.350,4.861) (-0.400,4.855) (-0.450,4.848) (-0.500,4.841) (-0.550,4.833) (-0.600,4.824) (-0.650,4.814) (-0.700,4.803) (-0.750,4.792) (-0.800,4.779) (-0.850,4.766) (-0.900,4.752) (-0.950,4.737) (-1.000,4.722) (-1.050,4.705) (-1.100,4.687) (-1.150,4.669) (-1.198,4.650) (-1.250,4.629) (-1.300,4.607) (-1.350,4.585) (-1.400,4.561) (-1.450,4.536) (-1.500,4.510) (-1.550,4.483) (-1.600,4.455) (-1.650,4.426) (-1.700,4.395) (-1.750,4.363) (-1.800,4.330) (-1.843,4.300) (-1.900,4.258) (-1.944,4.225) (-2.000,4.181) (-2.050,4.139) (-2.095,4.100) (-2.150,4.051) (-2.200,4.003) (-2.250,3.954) (-2.300,3.901) (-2.347,3.850) (-2.390,3.800) (-2.425,3.758) (-2.472,3.700) (-2.509,3.650) (-2.546,3.600) (-2.575,3.558) (-2.613,3.500) (-2.645,3.450) (-2.675,3.400) (-2.700,3.356) (-2.725,3.310) (-2.750,3.262) (-2.775,3.212) (-2.800,3.159) (-2.825,3.102) (-2.847,3.050) (-2.866,3.000) (-2.885,2.950) (-2.902,2.900) (-2.918,2.850) (-2.933,2.800) (-2.947,2.750) (-2.959,2.700) (-2.971,2.650) (-2.981,2.600) (-2.991,2.550) (-2.999,2.500) (-3.006,2.450) (-3.013,2.400) (-3.018,2.350) (-3.022,2.300) (-3.025,2.250) (-3.026,2.200) (-3.027,2.150) (-3.027,2.100) (-3.025,2.050) (-3.023,2.000) (-3.019,1.950) (-3.014,1.900) (-3.008,1.850) (-3.001,1.800) (-2.993,1.750) (-2.983,1.700) (-2.972,1.650) (-2.960,1.600) (-2.947,1.550) (-2.932,1.500) (-2.916,1.450) (-2.898,1.400) (-2.879,1.350) (-2.858,1.300) (-2.836,1.250) (-2.812,1.200) (-2.786,1.150) (-2.759,1.100) (-2.729,1.050) (-2.700,1.004) (-2.664,0.950) (-2.628,0.900) (-2.589,0.850) (-2.550,0.802) (-2.504,0.750) (-2.457,0.700) (-2.406,0.650) (-2.351,0.600) (-2.300,0.557) (-2.261,0.525) (-2.200,0.479) (-2.158,0.450) (-2.100,0.412) (-2.050,0.381) (-2.000,0.352) (-1.951,0.325) (-1.902,0.300) (-1.850,0.275) (-1.800,0.252) (-1.750,0.231) (-1.700,0.211) (-1.650,0.193) (-1.600,0.175) (-1.550,0.159) (-1.500,0.144) (-1.450,0.130) (-1.400,0.116) (-1.350,0.104) (-1.300,0.093)}
\EQ{(1.300,0.117) (1.350,0.131) (1.400,0.146) (1.450,0.163) (1.500,0.181) (1.549,0.200) (1.600,0.221) (1.650,0.244) (1.700,0.268) (1.750,0.293) (1.800,0.321) (1.849,0.350) (1.900,0.382) (1.950,0.416) (1.997,0.450) (2.050,0.492) (2.090,0.525) (2.145,0.575) (2.196,0.625) (2.243,0.675) (2.287,0.725) (2.325,0.772) (2.365,0.825) (2.400,0.875) (2.432,0.925) (2.462,0.975) (2.490,1.025) (2.516,1.075) (2.540,1.125) (2.562,1.175) (2.582,1.225) (2.600,1.275) (2.617,1.325) (2.632,1.375) (2.646,1.425) (2.658,1.475) (2.669,1.525) (2.679,1.575) (2.687,1.625) (2.693,1.675) (2.699,1.725) (2.703,1.775) (2.706,1.825) (2.707,1.875) (2.708,1.925) (2.707,1.975) (2.705,2.025) (2.701,2.075) (2.697,2.125) (2.691,2.175) (2.684,2.225) (2.676,2.275) (2.667,2.325) (2.656,2.375) (2.644,2.425) (2.631,2.475) (2.617,2.525) (2.602,2.575) (2.585,2.625) (2.567,2.675) (2.547,2.725) (2.527,2.775) (2.504,2.825) (2.481,2.875) (2.456,2.925) (2.430,2.975) (2.402,3.025) (2.375,3.070) (2.341,3.125) (2.308,3.175) (2.275,3.223) (2.237,3.275) (2.200,3.323) (2.158,3.375) (2.125,3.414) (2.075,3.470) (2.025,3.523) (1.975,3.574) (1.925,3.621) (1.875,3.667) (1.836,3.700) (1.776,3.750) (1.725,3.790) (1.677,3.825) (1.625,3.862) (1.575,3.896) (1.530,3.925) (1.475,3.959) (1.425,3.988) (1.375,4.016) (1.325,4.043) (1.275,4.068) (1.225,4.092) (1.175,4.115) (1.125,4.137) (1.075,4.158) (1.025,4.177) (0.975,4.196) (0.925,4.213) (0.875,4.229) (0.825,4.245) (0.775,4.259) (0.725,4.272) (0.675,4.285) (0.625,4.296) (0.575,4.307) (0.525,4.317) (0.475,4.325) (0.425,4.333) (0.375,4.340) (0.325,4.346) (0.275,4.351) (0.225,4.356) (0.175,4.359) (0.125,4.362) (0.075,4.363) (0.025,4.364) (-0.050,4.364) (-0.100,4.363) (-0.150,4.361) (-0.200,4.358) (-0.250,4.354) (-0.300,4.349) (-0.350,4.343) (-0.400,4.337) (-0.450,4.329) (-0.500,4.321) (-0.550,4.312) (-0.600,4.302) (-0.650,4.291) (-0.700,4.279) (-0.750,4.266) (-0.800,4.252) (-0.850,4.237) (-0.900,4.221) (-0.950,4.204) (-1.000,4.187) (-1.050,4.168) (-1.100,4.147) (-1.150,4.126) (-1.200,4.104) (-1.250,4.080) (-1.300,4.056) (-1.350,4.030) (-1.400,4.002) (-1.448,3.975) (-1.500,3.944) (-1.550,3.912) (-1.600,3.879) (-1.642,3.850) (-1.700,3.808) (-1.744,3.775) (-1.800,3.730) (-1.850,3.688) (-1.894,3.650) (-1.948,3.600) (-1.999,3.550) (-2.047,3.500) (-2.093,3.450) (-2.137,3.400) (-2.175,3.354) (-2.218,3.300) (-2.250,3.257) (-2.291,3.200) (-2.325,3.150) (-2.357,3.100) (-2.387,3.050) (-2.416,3.000) (-2.443,2.950) (-2.469,2.900) (-2.493,2.850) (-2.516,2.800) (-2.537,2.750) (-2.557,2.700) (-2.575,2.653) (-2.593,2.600) (-2.610,2.550) (-2.624,2.500) (-2.638,2.450) (-2.650,2.400) (-2.662,2.350) (-2.672,2.300) (-2.680,2.250) (-2.688,2.200) (-2.694,2.150) (-2.699,2.100) (-2.703,2.050) (-2.706,2.000) (-2.707,1.950) (-2.708,1.900) (-2.707,1.850) (-2.704,1.800) (-2.701,1.750) (-2.696,1.700) (-2.690,1.650) (-2.683,1.600) (-2.674,1.550) (-2.664,1.500) (-2.652,1.450) (-2.639,1.400) (-2.625,1.351) (-2.609,1.300) (-2.591,1.250) (-2.572,1.200) (-2.551,1.150) (-2.528,1.100) (-2.503,1.050) (-2.476,1.000) (-2.450,0.954) (-2.416,0.900) (-2.383,0.850) (-2.350,0.805) (-2.308,0.750) (-2.275,0.711) (-2.225,0.655) (-2.175,0.604) (-2.125,0.557) (-2.075,0.513) (-2.029,0.475) (-1.975,0.434) (-1.927,0.400) (-1.875,0.366) (-1.825,0.335) (-1.775,0.307) (-1.725,0.280) (-1.675,0.255) (-1.625,0.232) (-1.575,0.211) (-1.525,0.191) (-1.475,0.172) (-1.425,0.155) (-1.375,0.139) (-1.325,0.124) (-1.300,0.117)}
\EQ{(1.300,0.141) (1.350,0.158) (1.400,0.177) (1.450,0.197) (1.500,0.220) (1.550,0.243) (1.600,0.269) (1.650,0.297) (1.697,0.325) (1.750,0.359) (1.800,0.394) (1.841,0.425) (1.900,0.473) (1.950,0.517) (2.000,0.566) (2.050,0.619) (2.099,0.675) (2.138,0.725) (2.174,0.775) (2.208,0.825) (2.239,0.875) (2.267,0.925) (2.293,0.975) (2.317,1.025) (2.339,1.075) (2.359,1.125) (2.377,1.175) (2.394,1.225) (2.409,1.275) (2.422,1.325) (2.433,1.375) (2.443,1.425) (2.451,1.475) (2.458,1.525) (2.464,1.575) (2.468,1.625) (2.470,1.675) (2.472,1.725) (2.472,1.775) (2.470,1.825) (2.467,1.875) (2.463,1.925) (2.458,1.975) (2.451,2.025) (2.443,2.075) (2.434,2.125) (2.423,2.175) (2.411,2.225) (2.398,2.275) (2.383,2.325) (2.367,2.375) (2.350,2.423) (2.330,2.475) (2.310,2.525) (2.288,2.575) (2.265,2.625) (2.240,2.675) (2.213,2.725) (2.184,2.775) (2.154,2.825) (2.125,2.871) (2.089,2.925) (2.053,2.975) (2.015,3.025) (1.975,3.075) (1.933,3.125) (1.888,3.175) (1.850,3.215) (1.800,3.265) (1.763,3.300) (1.725,3.335) (1.679,3.375) (1.625,3.420) (1.575,3.459) (1.525,3.496) (1.475,3.531) (1.425,3.565) (1.375,3.597) (1.328,3.625) (1.285,3.650) (1.239,3.675) (1.191,3.700) (1.141,3.725) (1.087,3.750) (1.031,3.775) (0.975,3.798) (0.925,3.817) (0.875,3.835) (0.825,3.852) (0.775,3.868) (0.725,3.883) (0.675,3.897) (0.625,3.909) (0.575,3.921) (0.525,3.932) (0.475,3.941) (0.425,3.950) (0.375,3.958) (0.325,3.964) (0.275,3.970) (0.225,3.975) (0.175,3.978) (0.125,3.981) (0.075,3.983) (0.025,3.984) (-0.050,3.984) (-0.100,3.982) (-0.150,3.980) (-0.200,3.977) (-0.250,3.972) (-0.300,3.967) (-0.350,3.961) (-0.400,3.954) (-0.450,3.946) (-0.500,3.937) (-0.550,3.926) (-0.600,3.915) (-0.650,3.903) (-0.700,3.890) (-0.750,3.876) (-0.800,3.860) (-0.850,3.844) (-0.900,3.826) (-0.950,3.808) (-1.000,3.788) (-1.050,3.767) (-1.100,3.744) (-1.150,3.721) (-1.200,3.696) (-1.250,3.669) (-1.300,3.641) (-1.350,3.612) (-1.400,3.581) (-1.448,3.550) (-1.500,3.514) (-1.550,3.478) (-1.600,3.440) (-1.650,3.400) (-1.700,3.357) (-1.750,3.312) (-1.790,3.275) (-1.840,3.225) (-1.875,3.189) (-1.910,3.150) (-1.950,3.105) (-1.995,3.050) (-2.034,3.000) (-2.071,2.950) (-2.100,2.908) (-2.139,2.850) (-2.170,2.800) (-2.199,2.750) (-2.225,2.703) (-2.250,2.654) (-2.275,2.603) (-2.299,2.550) (-2.320,2.500) (-2.340,2.450) (-2.358,2.400) (-2.375,2.351) (-2.391,2.300) (-2.405,2.250) (-2.417,2.200) (-2.429,2.150) (-2.439,2.100) (-2.447,2.050) (-2.455,2.000) (-2.461,1.950) (-2.466,1.900) (-2.469,1.850) (-2.471,1.800) (-2.472,1.750) (-2.471,1.700) (-2.469,1.650) (-2.466,1.600) (-2.461,1.550) (-2.455,1.500) (-2.447,1.450) (-2.438,1.400) (-2.427,1.350) (-2.415,1.300) (-2.401,1.250) (-2.386,1.200) (-2.369,1.150) (-2.350,1.101) (-2.329,1.050) (-2.306,1.000) (-2.281,0.950) (-2.253,0.900) (-2.225,0.852) (-2.192,0.800) (-2.157,0.750) (-2.125,0.708) (-2.078,0.650) (-2.033,0.600) (-1.984,0.550) (-1.931,0.500) (-1.875,0.452) (-1.825,0.413) (-1.775,0.376) (-1.725,0.343) (-1.675,0.312) (-1.625,0.283) (-1.575,0.256) (-1.525,0.231) (-1.475,0.208) (-1.425,0.187) (-1.375,0.167) (-1.327,0.150) (-1.300,0.141)}
\EQ{(1.300,0.165) (1.350,0.186) (1.400,0.208) (1.450,0.233) (1.500,0.259) (1.550,0.288) (1.600,0.319) (1.646,0.350) (1.700,0.390) (1.744,0.425) (1.800,0.474) (1.850,0.522) (1.900,0.575) (1.943,0.625) (1.975,0.665) (2.018,0.725) (2.050,0.773) (2.082,0.825) (2.109,0.875) (2.134,0.925) (2.157,0.975) (2.178,1.025) (2.197,1.075) (2.213,1.125) (2.228,1.175) (2.242,1.225) (2.253,1.275) (2.263,1.325) (2.271,1.375) (2.277,1.425) (2.282,1.475) (2.286,1.525) (2.288,1.575) (2.288,1.625) (2.288,1.675) (2.285,1.725) (2.281,1.775) (2.276,1.825) (2.269,1.875) (2.261,1.925) (2.252,1.975) (2.241,2.025) (2.228,2.075) (2.215,2.125) (2.199,2.175) (2.182,2.225) (2.164,2.275) (2.144,2.325) (2.122,2.375) (2.100,2.423) (2.075,2.473) (2.050,2.520) (2.025,2.565) (1.989,2.625) (1.956,2.675) (1.925,2.720) (1.885,2.775) (1.850,2.820) (1.805,2.875) (1.761,2.925) (1.725,2.964) (1.690,3.000) (1.650,3.039) (1.611,3.075) (1.554,3.125) (1.500,3.170) (1.450,3.209) (1.400,3.245) (1.358,3.275) (1.300,3.313) (1.250,3.344) (1.200,3.373) (1.152,3.400) (1.104,3.425) (1.053,3.450) (1.000,3.475) (0.950,3.497) (0.900,3.517) (0.850,3.536) (0.800,3.554) (0.750,3.571) (0.700,3.587) (0.650,3.601) (0.600,3.614) (0.550,3.626) (0.500,3.637) (0.450,3.647) (0.400,3.656) (0.350,3.664) (0.300,3.670) (0.250,3.676) (0.200,3.680) (0.150,3.684) (0.100,3.687) (0.050,3.688) (0.000,3.689) (-0.050,3.688) (-0.100,3.687) (-0.150,3.684) (-0.200,3.680) (-0.250,3.676) (-0.300,3.670) (-0.350,3.664) (-0.400,3.656) (-0.450,3.647) (-0.500,3.637) (-0.550,3.626) (-0.600,3.614) (-0.650,3.601) (-0.700,3.587) (-0.750,3.571) (-0.800,3.554) (-0.850,3.536) (-0.900,3.517) (-0.950,3.497) (-1.000,3.475) (-1.050,3.452) (-1.100,3.427) (-1.150,3.401) (-1.197,3.375) (-1.250,3.344) (-1.300,3.313) (-1.350,3.280) (-1.394,3.250) (-1.450,3.209) (-1.494,3.175) (-1.550,3.129) (-1.600,3.085) (-1.638,3.050) (-1.675,3.015) (-1.714,2.975) (-1.750,2.937) (-1.800,2.881) (-1.846,2.825) (-1.885,2.775) (-1.922,2.725) (-1.956,2.675) (-1.989,2.625) (-2.019,2.575) (-2.047,2.525) (-2.074,2.475) (-2.099,2.425) (-2.122,2.375) (-2.144,2.325) (-2.164,2.275) (-2.182,2.225) (-2.199,2.175) (-2.215,2.125) (-2.228,2.075) (-2.241,2.025) (-2.252,1.975) (-2.261,1.925) (-2.269,1.875) (-2.276,1.825) (-2.281,1.775) (-2.285,1.725) (-2.288,1.675) (-2.288,1.625) (-2.288,1.575) (-2.286,1.525) (-2.282,1.475) (-2.277,1.425) (-2.271,1.375) (-2.263,1.325) (-2.253,1.275) (-2.242,1.225) (-2.228,1.175) (-2.213,1.125) (-2.197,1.075) (-2.178,1.025) (-2.157,0.975) (-2.134,0.925) (-2.109,0.875) (-2.082,0.825) (-2.051,0.775) (-2.018,0.725) (-1.982,0.675) (-1.950,0.634) (-1.900,0.575) (-1.853,0.525) (-1.801,0.475) (-1.750,0.430) (-1.700,0.390) (-1.650,0.353) (-1.600,0.319) (-1.550,0.288) (-1.500,0.259) (-1.450,0.233) (-1.400,0.208) (-1.350,0.186) (-1.300,0.165)}
\EQ{(1.300,0.191) (1.350,0.215) (1.400,0.241) (1.450,0.270) (1.499,0.300) (1.550,0.335) (1.600,0.372) (1.650,0.414) (1.691,0.450) (1.742,0.500) (1.788,0.550) (1.825,0.594) (1.868,0.650) (1.900,0.696) (1.934,0.750) (1.963,0.800) (1.989,0.850) (2.013,0.900) (2.034,0.950) (2.053,1.000) (2.070,1.050) (2.085,1.100) (2.098,1.150) (2.109,1.200) (2.118,1.250) (2.126,1.300) (2.132,1.350) (2.137,1.400) (2.139,1.450) (2.141,1.500) (2.140,1.550) (2.138,1.600) (2.135,1.650) (2.130,1.700) (2.123,1.750) (2.115,1.800) (2.106,1.850) (2.095,1.900) (2.082,1.950) (2.068,2.000) (2.052,2.050) (2.035,2.100) (2.016,2.150) (1.995,2.200) (1.972,2.250) (1.950,2.296) (1.925,2.344) (1.900,2.389) (1.864,2.450) (1.831,2.500) (1.800,2.546) (1.760,2.600) (1.725,2.645) (1.679,2.700) (1.634,2.750) (1.600,2.786) (1.562,2.825) (1.525,2.860) (1.481,2.900) (1.425,2.947) (1.375,2.987) (1.325,3.025) (1.275,3.060) (1.225,3.093) (1.175,3.124) (1.132,3.150) (1.087,3.175) (1.040,3.200) (0.990,3.225) (0.936,3.250) (0.879,3.275) (0.825,3.297) (0.775,3.315) (0.725,3.333) (0.675,3.349) (0.625,3.364) (0.575,3.377) (0.525,3.390) (0.475,3.401) (0.425,3.411) (0.375,3.420) (0.325,3.427) (0.275,3.434) (0.225,3.439) (0.175,3.444) (0.125,3.447) (0.075,3.449) (0.025,3.450) (-0.037,3.450) (-0.100,3.448) (-0.150,3.445) (-0.200,3.442) (-0.250,3.437) (-0.300,3.431) (-0.350,3.424) (-0.400,3.415) (-0.450,3.406) (-0.500,3.395) (-0.550,3.384) (-0.600,3.371) (-0.650,3.356) (-0.700,3.341) (-0.748,3.325) (-0.800,3.306) (-0.850,3.287) (-0.900,3.266) (-0.950,3.244) (-1.000,3.220) (-1.050,3.195) (-1.100,3.168) (-1.150,3.139) (-1.200,3.109) (-1.250,3.077) (-1.300,3.043) (-1.350,3.006) (-1.391,2.975) (-1.450,2.927) (-1.500,2.883) (-1.550,2.836) (-1.587,2.800) (-1.625,2.760) (-1.675,2.705) (-1.721,2.650) (-1.760,2.600) (-1.797,2.550) (-1.831,2.500) (-1.864,2.450) (-1.894,2.400) (-1.922,2.350) (-1.948,2.300) (-1.972,2.250) (-1.995,2.200) (-2.016,2.150) (-2.035,2.100) (-2.052,2.050) (-2.068,2.000) (-2.082,1.950) (-2.095,1.900) (-2.106,1.850) (-2.115,1.800) (-2.123,1.750) (-2.130,1.700) (-2.135,1.650) (-2.138,1.600) (-2.140,1.550) (-2.141,1.500) (-2.139,1.450) (-2.137,1.400) (-2.132,1.350) (-2.126,1.300) (-2.118,1.250) (-2.109,1.200) (-2.098,1.150) (-2.085,1.100) (-2.070,1.050) (-2.053,1.000) (-2.034,0.950) (-2.013,0.900) (-1.989,0.850) (-1.963,0.800) (-1.934,0.750) (-1.903,0.700) (-1.868,0.650) (-1.830,0.600) (-1.788,0.550) (-1.750,0.508) (-1.700,0.459) (-1.650,0.414) (-1.603,0.375) (-1.550,0.335) (-1.500,0.301) (-1.450,0.270) (-1.400,0.241) (-1.350,0.215) (-1.300,0.191)}
\EQ{(1.300,0.216) (1.350,0.244) (1.400,0.274) (1.450,0.308) (1.500,0.345) (1.550,0.385) (1.595,0.425) (1.646,0.475) (1.692,0.525) (1.725,0.565) (1.770,0.625) (1.800,0.669) (1.834,0.725) (1.862,0.775) (1.887,0.825) (1.909,0.875) (1.929,0.925) (1.947,0.975) (1.962,1.025) (1.975,1.075) (1.987,1.125) (1.997,1.175) (2.004,1.225) (2.011,1.275) (2.015,1.325) (2.017,1.375) (2.018,1.425) (2.017,1.475) (2.015,1.525) (2.011,1.575) (2.005,1.625) (1.998,1.675) (1.989,1.725) (1.979,1.775) (1.967,1.825) (1.953,1.875) (1.937,1.925) (1.920,1.975) (1.901,2.025) (1.880,2.075) (1.858,2.125) (1.833,2.175) (1.807,2.225) (1.778,2.275) (1.750,2.321) (1.714,2.375) (1.679,2.425) (1.641,2.475) (1.601,2.525) (1.557,2.575) (1.510,2.625) (1.475,2.661) (1.434,2.700) (1.378,2.750) (1.325,2.794) (1.285,2.825) (1.225,2.869) (1.180,2.900) (1.125,2.935) (1.075,2.965) (1.025,2.993) (0.975,3.020) (0.925,3.045) (0.875,3.068) (0.825,3.089) (0.775,3.109) (0.725,3.128) (0.675,3.145) (0.625,3.161) (0.576,3.175) (0.525,3.188) (0.476,3.200) (0.425,3.211) (0.375,3.220) (0.325,3.229) (0.275,3.236) (0.225,3.241) (0.175,3.246) (0.125,3.249) (0.075,3.252) (0.025,3.253) (-0.050,3.252) (-0.100,3.251) (-0.150,3.248) (-0.200,3.244) (-0.250,3.239) (-0.300,3.232) (-0.350,3.225) (-0.400,3.216) (-0.450,3.206) (-0.500,3.195) (-0.550,3.182) (-0.600,3.168) (-0.650,3.153) (-0.700,3.137) (-0.750,3.119) (-0.798,3.100) (-0.850,3.079) (-0.900,3.056) (-0.950,3.032) (-1.000,3.007) (-1.050,2.980) (-1.100,2.950) (-1.150,2.919) (-1.200,2.886) (-1.250,2.851) (-1.300,2.814) (-1.348,2.775) (-1.400,2.731) (-1.450,2.685) (-1.486,2.650) (-1.525,2.610) (-1.575,2.555) (-1.621,2.500) (-1.661,2.450) (-1.697,2.400) (-1.731,2.350) (-1.763,2.300) (-1.793,2.250) (-1.820,2.200) (-1.846,2.150) (-1.869,2.100) (-1.891,2.050) (-1.911,2.000) (-1.929,1.950) (-1.945,1.900) (-1.960,1.850) (-1.973,1.800) (-1.984,1.750) (-1.994,1.700) (-2.002,1.650) (-2.008,1.600) (-2.013,1.550) (-2.017,1.500) (-2.018,1.450) (-2.018,1.400) (-2.016,1.350) (-2.013,1.300) (-2.008,1.250) (-2.001,1.200) (-1.992,1.150) (-1.982,1.100) (-1.969,1.050) (-1.955,1.000) (-1.938,0.950) (-1.919,0.900) (-1.898,0.850) (-1.875,0.801) (-1.850,0.753) (-1.825,0.710) (-1.787,0.650) (-1.752,0.600) (-1.713,0.550) (-1.675,0.506) (-1.625,0.454) (-1.575,0.407) (-1.525,0.364) (-1.475,0.326) (-1.425,0.291) (-1.375,0.259) (-1.325,0.230) (-1.300,0.216)}
\EQ{(1.300,0.243) (1.350,0.275) (1.400,0.310) (1.450,0.348) (1.500,0.391) (1.550,0.439) (1.600,0.493) (1.647,0.550) (1.684,0.600) (1.717,0.650) (1.747,0.700) (1.774,0.750) (1.798,0.800) (1.819,0.850) (1.838,0.900) (1.855,0.950) (1.869,1.000) (1.881,1.050) (1.892,1.100) (1.900,1.150) (1.906,1.200) (1.911,1.250) (1.914,1.300) (1.915,1.350) (1.914,1.400) (1.911,1.450) (1.907,1.500) (1.901,1.550) (1.894,1.600) (1.885,1.650) (1.873,1.700) (1.861,1.750) (1.846,1.800) (1.830,1.850) (1.812,1.900) (1.792,1.950) (1.770,2.000) (1.746,2.050) (1.720,2.100) (1.692,2.150) (1.661,2.200) (1.629,2.250) (1.593,2.300) (1.555,2.350) (1.515,2.400) (1.475,2.445) (1.425,2.498) (1.375,2.548) (1.325,2.593) (1.275,2.636) (1.226,2.675) (1.175,2.713) (1.125,2.747) (1.083,2.775) (1.025,2.810) (0.975,2.838) (0.925,2.865) (0.875,2.890) (0.825,2.913) (0.775,2.934) (0.725,2.954) (0.675,2.972) (0.625,2.989) (0.575,3.004) (0.525,3.018) (0.475,3.030) (0.425,3.042) (0.375,3.052) (0.325,3.060) (0.275,3.068) (0.225,3.074) (0.175,3.079) (0.125,3.082) (0.075,3.085) (0.025,3.086) (-0.050,3.086) (-0.100,3.084) (-0.150,3.081) (-0.200,3.076) (-0.250,3.071) (-0.300,3.064) (-0.350,3.056) (-0.400,3.047) (-0.450,3.036) (-0.500,3.024) (-0.550,3.011) (-0.600,2.996) (-0.650,2.980) (-0.700,2.963) (-0.750,2.944) (-0.800,2.923) (-0.850,2.901) (-0.900,2.877) (-0.950,2.852) (-0.999,2.825) (-1.050,2.795) (-1.100,2.764) (-1.150,2.730) (-1.192,2.700) (-1.250,2.656) (-1.300,2.615) (-1.345,2.575) (-1.398,2.525) (-1.448,2.475) (-1.493,2.425) (-1.535,2.375) (-1.575,2.325) (-1.611,2.275) (-1.645,2.225) (-1.675,2.178) (-1.706,2.125) (-1.733,2.075) (-1.758,2.025) (-1.781,1.975) (-1.802,1.925) (-1.821,1.875) (-1.838,1.825) (-1.854,1.775) (-1.867,1.725) (-1.879,1.675) (-1.889,1.625) (-1.898,1.575) (-1.905,1.525) (-1.910,1.475) (-1.913,1.425) (-1.915,1.375) (-1.914,1.325) (-1.913,1.275) (-1.909,1.225) (-1.903,1.175) (-1.896,1.125) (-1.887,1.075) (-1.875,1.025) (-1.862,0.975) (-1.847,0.925) (-1.829,0.875) (-1.809,0.825) (-1.786,0.775) (-1.761,0.725) (-1.733,0.675) (-1.701,0.625) (-1.666,0.575) (-1.627,0.525) (-1.584,0.475) (-1.536,0.425) (-1.482,0.375) (-1.425,0.328) (-1.375,0.292) (-1.325,0.258) (-1.300,0.243)}
\EQ{(1.288,0.262) (1.342,0.300) (1.400,0.347) (1.450,0.391) (1.500,0.442) (1.550,0.499) (1.590,0.550) (1.625,0.600) (1.656,0.650) (1.684,0.700) (1.708,0.750) (1.730,0.800) (1.750,0.850) (1.767,0.900) (1.781,0.950) (1.794,1.000) (1.804,1.050) (1.812,1.100) (1.818,1.150) (1.823,1.200) (1.825,1.250) (1.825,1.300) (1.824,1.350) (1.821,1.400) (1.816,1.450) (1.810,1.500) (1.801,1.550) (1.791,1.600) (1.779,1.650) (1.765,1.700) (1.750,1.748) (1.732,1.800) (1.712,1.850) (1.691,1.900) (1.667,1.950) (1.641,2.000) (1.613,2.050) (1.583,2.100) (1.550,2.149) (1.514,2.200) (1.476,2.250) (1.434,2.300) (1.400,2.339) (1.366,2.375) (1.325,2.416) (1.289,2.450) (1.250,2.484) (1.201,2.525) (1.150,2.565) (1.101,2.600) (1.050,2.634) (1.000,2.666) (0.950,2.695) (0.900,2.722) (0.850,2.748) (0.800,2.771) (0.750,2.793) (0.700,2.813) (0.650,2.831) (0.600,2.848) (0.550,2.864) (0.500,2.878) (0.450,2.890) (0.400,2.901) (0.350,2.911) (0.300,2.919) (0.250,2.926) (0.200,2.932) (0.150,2.937) (0.100,2.940) (0.050,2.942) (0.000,2.943) (-0.050,2.942) (-0.100,2.940) (-0.150,2.937) (-0.200,2.932) (-0.250,2.926) (-0.300,2.919) (-0.350,2.911) (-0.400,2.901) (-0.450,2.890) (-0.500,2.878) (-0.550,2.864) (-0.600,2.848) (-0.650,2.831) (-0.700,2.813) (-0.750,2.793) (-0.800,2.771) (-0.850,2.748) (-0.900,2.722) (-0.950,2.695) (-1.000,2.666) (-1.050,2.634) (-1.100,2.601) (-1.150,2.565) (-1.200,2.526) (-1.250,2.484) (-1.289,2.450) (-1.325,2.416) (-1.366,2.375) (-1.400,2.339) (-1.434,2.300) (-1.475,2.251) (-1.514,2.200) (-1.550,2.150) (-1.583,2.100) (-1.613,2.050) (-1.641,2.000) (-1.667,1.950) (-1.691,1.900) (-1.712,1.850) (-1.732,1.800) (-1.749,1.750) (-1.765,1.700) (-1.779,1.650) (-1.791,1.600) (-1.801,1.550) (-1.810,1.500) (-1.816,1.450) (-1.821,1.400) (-1.824,1.350) (-1.825,1.300) (-1.825,1.250) (-1.823,1.200) (-1.818,1.150) (-1.812,1.100) (-1.804,1.050) (-1.794,1.000) (-1.781,0.950) (-1.767,0.900) (-1.750,0.851) (-1.730,0.800) (-1.708,0.750) (-1.684,0.700) (-1.656,0.650) (-1.625,0.601) (-1.590,0.550) (-1.551,0.500) (-1.508,0.450) (-1.459,0.400) (-1.404,0.350) (-1.350,0.307) (-1.307,0.275) (-1.288,0.262)}
\EQ{(1.275,0.280) (1.325,0.319) (1.375,0.362) (1.414,0.400) (1.450,0.438) (1.500,0.497) (1.539,0.550) (1.572,0.600) (1.600,0.648) (1.625,0.696) (1.650,0.750) (1.670,0.800) (1.688,0.850) (1.703,0.900) (1.716,0.950) (1.726,1.000) (1.734,1.050) (1.741,1.100) (1.745,1.150) (1.747,1.200) (1.748,1.250) (1.746,1.300) (1.743,1.350) (1.738,1.400) (1.731,1.450) (1.722,1.500) (1.711,1.550) (1.698,1.600) (1.683,1.650) (1.667,1.700) (1.648,1.750) (1.627,1.800) (1.604,1.850) (1.579,1.900) (1.552,1.950) (1.525,1.995) (1.489,2.050) (1.454,2.100) (1.416,2.150) (1.375,2.200) (1.330,2.250) (1.282,2.300) (1.229,2.350) (1.175,2.397) (1.125,2.437) (1.075,2.475) (1.025,2.509) (0.975,2.542) (0.925,2.571) (0.875,2.599) (0.825,2.625) (0.775,2.649) (0.725,2.671) (0.675,2.691) (0.625,2.710) (0.575,2.727) (0.525,2.742) (0.475,2.756) (0.425,2.768) (0.375,2.779) (0.325,2.789) (0.275,2.797) (0.225,2.804) (0.175,2.809) (0.125,2.813) (0.075,2.816) (0.025,2.817) (-0.050,2.817) (-0.100,2.815) (-0.150,2.811) (-0.200,2.807) (-0.250,2.800) (-0.300,2.793) (-0.350,2.784) (-0.400,2.774) (-0.450,2.762) (-0.500,2.749) (-0.550,2.735) (-0.600,2.719) (-0.650,2.701) (-0.700,2.681) (-0.750,2.660) (-0.800,2.637) (-0.850,2.612) (-0.900,2.586) (-0.950,2.557) (-1.000,2.526) (-1.050,2.492) (-1.100,2.456) (-1.141,2.425) (-1.200,2.376) (-1.250,2.331) (-1.300,2.282) (-1.350,2.228) (-1.396,2.175) (-1.436,2.125) (-1.472,2.075) (-1.506,2.025) (-1.537,1.975) (-1.566,1.925) (-1.592,1.875) (-1.616,1.825) (-1.638,1.775) (-1.658,1.725) (-1.675,1.675) (-1.691,1.625) (-1.705,1.575) (-1.716,1.525) (-1.726,1.475) (-1.734,1.425) (-1.741,1.375) (-1.745,1.325) (-1.747,1.275) (-1.748,1.225) (-1.746,1.175) (-1.743,1.125) (-1.738,1.075) (-1.730,1.025) (-1.721,0.975) (-1.709,0.925) (-1.695,0.875) (-1.679,0.825) (-1.660,0.775) (-1.639,0.725) (-1.614,0.675) (-1.587,0.625) (-1.556,0.575) (-1.525,0.530) (-1.482,0.475) (-1.438,0.425) (-1.400,0.386) (-1.362,0.350) (-1.301,0.300) (-1.275,0.280)}
\EQ{(1.275,0.308) (1.323,0.350) (1.374,0.400) (1.419,0.450) (1.450,0.489) (1.493,0.550) (1.524,0.600) (1.550,0.647) (1.575,0.698) (1.597,0.750) (1.616,0.800) (1.632,0.850) (1.645,0.900) (1.656,0.950) (1.665,1.000) (1.672,1.050) (1.676,1.100) (1.679,1.150) (1.679,1.200) (1.678,1.250) (1.674,1.300) (1.669,1.350) (1.662,1.400) (1.652,1.450) (1.641,1.500) (1.628,1.550) (1.613,1.600) (1.595,1.650) (1.576,1.700) (1.554,1.750) (1.530,1.800) (1.504,1.850) (1.475,1.900) (1.444,1.950) (1.410,2.000) (1.375,2.047) (1.333,2.100) (1.289,2.150) (1.250,2.191) (1.200,2.240) (1.162,2.275) (1.102,2.325) (1.050,2.365) (1.001,2.400) (0.950,2.434) (0.900,2.464) (0.850,2.492) (0.800,2.518) (0.750,2.543) (0.700,2.565) (0.650,2.585) (0.600,2.604) (0.550,2.621) (0.500,2.636) (0.450,2.650) (0.400,2.662) (0.350,2.672) (0.300,2.682) (0.250,2.689) (0.200,2.696) (0.150,2.700) (0.100,2.704) (0.050,2.706) (0.000,2.707) (-0.050,2.706) (-0.100,2.704) (-0.150,2.700) (-0.200,2.696) (-0.250,2.689) (-0.300,2.682) (-0.350,2.672) (-0.400,2.662) (-0.450,2.650) (-0.500,2.636) (-0.550,2.621) (-0.600,2.604) (-0.650,2.585) (-0.700,2.565) (-0.750,2.543) (-0.800,2.518) (-0.850,2.492) (-0.900,2.464) (-0.950,2.434) (-1.000,2.401) (-1.050,2.365) (-1.100,2.327) (-1.150,2.285) (-1.189,2.250) (-1.241,2.200) (-1.289,2.150) (-1.325,2.109) (-1.373,2.050) (-1.410,2.000) (-1.444,1.950) (-1.475,1.900) (-1.504,1.850) (-1.530,1.800) (-1.554,1.750) (-1.575,1.702) (-1.595,1.650) (-1.613,1.600) (-1.628,1.550) (-1.641,1.500) (-1.652,1.450) (-1.662,1.400) (-1.669,1.350) (-1.674,1.300) (-1.678,1.250) (-1.679,1.200) (-1.679,1.150) (-1.676,1.100) (-1.672,1.050) (-1.665,1.000) (-1.656,0.950) (-1.645,0.900) (-1.632,0.850) (-1.616,0.800) (-1.597,0.750) (-1.576,0.700) (-1.552,0.650) (-1.525,0.601) (-1.493,0.550) (-1.458,0.500) (-1.425,0.458) (-1.375,0.401) (-1.325,0.352) (-1.275,0.308)}
\EQ{(1.275,0.337) (1.314,0.375) (1.350,0.414) (1.400,0.475) (1.436,0.525) (1.467,0.575) (1.495,0.625) (1.519,0.675) (1.540,0.725) (1.559,0.775) (1.574,0.825) (1.587,0.875) (1.598,0.925) (1.606,0.975) (1.612,1.025) (1.616,1.075) (1.618,1.125) (1.618,1.175) (1.616,1.225) (1.611,1.275) (1.605,1.325) (1.597,1.375) (1.586,1.425) (1.574,1.475) (1.559,1.525) (1.543,1.575) (1.524,1.625) (1.503,1.675) (1.479,1.725) (1.454,1.775) (1.425,1.825) (1.394,1.875) (1.360,1.925) (1.325,1.973) (1.283,2.025) (1.250,2.063) (1.216,2.100) (1.175,2.141) (1.125,2.187) (1.081,2.225) (1.025,2.269) (0.983,2.300) (0.925,2.339) (0.875,2.370) (0.825,2.398) (0.775,2.425) (0.725,2.449) (0.675,2.471) (0.625,2.492) (0.575,2.510) (0.525,2.527) (0.475,2.542) (0.425,2.555) (0.375,2.567) (0.325,2.578) (0.275,2.586) (0.225,2.594) (0.175,2.599) (0.125,2.604) (0.075,2.607) (0.025,2.608) (-0.050,2.608) (-0.100,2.605) (-0.150,2.602) (-0.200,2.597) (-0.250,2.590) (-0.300,2.582) (-0.350,2.573) (-0.400,2.562) (-0.450,2.549) (-0.500,2.535) (-0.550,2.519) (-0.600,2.501) (-0.650,2.482) (-0.700,2.460) (-0.750,2.437) (-0.800,2.412) (-0.850,2.384) (-0.900,2.355) (-0.946,2.325) (-1.000,2.288) (-1.050,2.250) (-1.100,2.209) (-1.139,2.175) (-1.191,2.125) (-1.239,2.075) (-1.275,2.034) (-1.323,1.975) (-1.360,1.925) (-1.394,1.875) (-1.425,1.825) (-1.450,1.781) (-1.475,1.734) (-1.500,1.681) (-1.524,1.625) (-1.543,1.575) (-1.559,1.525) (-1.574,1.475) (-1.586,1.425) (-1.597,1.375) (-1.605,1.325) (-1.611,1.275) (-1.616,1.225) (-1.618,1.175) (-1.618,1.125) (-1.616,1.075) (-1.612,1.025) (-1.606,0.975) (-1.598,0.925) (-1.587,0.875) (-1.574,0.825) (-1.559,0.775) (-1.540,0.725) (-1.519,0.675) (-1.495,0.625) (-1.467,0.575) (-1.436,0.525) (-1.400,0.475) (-1.360,0.425) (-1.325,0.387) (-1.289,0.350) (-1.275,0.337)}
\EQ{(1.266,0.359) (1.305,0.400) (1.346,0.450) (1.382,0.500) (1.414,0.550) (1.442,0.600) (1.467,0.650) (1.488,0.700) (1.506,0.750) (1.522,0.800) (1.534,0.850) (1.545,0.900) (1.553,0.950) (1.559,1.000) (1.562,1.050) (1.563,1.100) (1.563,1.150) (1.560,1.200) (1.555,1.250) (1.547,1.300) (1.538,1.350) (1.527,1.400) (1.513,1.450) (1.498,1.500) (1.480,1.550) (1.460,1.600) (1.437,1.650) (1.412,1.700) (1.385,1.750) (1.354,1.800) (1.325,1.844) (1.285,1.900) (1.250,1.944) (1.201,2.000) (1.154,2.050) (1.101,2.100) (1.050,2.145) (1.000,2.185) (0.950,2.222) (0.908,2.250) (0.850,2.287) (0.800,2.315) (0.750,2.342) (0.700,2.366) (0.650,2.388) (0.600,2.409) (0.550,2.427) (0.500,2.443) (0.450,2.458) (0.400,2.471) (0.350,2.483) (0.300,2.493) (0.250,2.501) (0.200,2.508) (0.150,2.513) (0.100,2.517) (0.050,2.519) (0.000,2.520) (-0.050,2.519) (-0.100,2.517) (-0.150,2.513) (-0.200,2.508) (-0.250,2.501) (-0.300,2.493) (-0.350,2.483) (-0.400,2.471) (-0.450,2.458) (-0.500,2.443) (-0.550,2.427) (-0.600,2.409) (-0.650,2.388) (-0.700,2.366) (-0.750,2.342) (-0.800,2.315) (-0.850,2.287) (-0.900,2.256) (-0.945,2.225) (-1.000,2.185) (-1.044,2.150) (-1.100,2.101) (-1.150,2.054) (-1.200,2.001) (-1.245,1.950) (-1.285,1.900) (-1.321,1.850) (-1.350,1.807) (-1.385,1.750) (-1.412,1.700) (-1.437,1.650) (-1.460,1.600) (-1.480,1.550) (-1.498,1.500) (-1.513,1.450) (-1.527,1.400) (-1.538,1.350) (-1.547,1.300) (-1.555,1.250) (-1.560,1.200) (-1.563,1.150) (-1.563,1.100) (-1.562,1.050) (-1.559,1.000) (-1.553,0.950) (-1.545,0.900) (-1.534,0.850) (-1.522,0.800) (-1.506,0.750) (-1.488,0.700) (-1.467,0.650) (-1.442,0.600) (-1.414,0.550) (-1.382,0.500) (-1.350,0.456) (-1.305,0.400) (-1.266,0.359)}
\EQ{(1.251,0.374) (1.295,0.425) (1.325,0.465) (1.365,0.525) (1.393,0.575) (1.418,0.625) (1.440,0.675) (1.458,0.725) (1.473,0.775) (1.486,0.825) (1.496,0.875) (1.504,0.925) (1.510,0.975) (1.513,1.025) (1.514,1.075) (1.512,1.125) (1.509,1.175) (1.503,1.225) (1.495,1.275) (1.486,1.325) (1.473,1.375) (1.459,1.425) (1.442,1.475) (1.423,1.525) (1.402,1.575) (1.378,1.625) (1.352,1.675) (1.325,1.721) (1.290,1.775) (1.255,1.825) (1.216,1.875) (1.175,1.923) (1.127,1.975) (1.075,2.025) (1.025,2.070) (0.975,2.110) (0.925,2.147) (0.875,2.182) (0.825,2.213) (0.775,2.242) (0.725,2.268) (0.675,2.292) (0.625,2.314) (0.575,2.334) (0.525,2.352) (0.475,2.369) (0.425,2.383) (0.375,2.396) (0.325,2.407) (0.275,2.416) (0.225,2.424) (0.175,2.430) (0.125,2.435) (0.075,2.438) (0.025,2.440) (-0.050,2.439) (-0.100,2.437) (-0.150,2.433) (-0.200,2.427) (-0.250,2.420) (-0.300,2.412) (-0.350,2.402) (-0.400,2.390) (-0.450,2.376) (-0.500,2.361) (-0.550,2.344) (-0.599,2.325) (-0.650,2.304) (-0.700,2.280) (-0.750,2.255) (-0.800,2.228) (-0.846,2.200) (-0.900,2.165) (-0.950,2.129) (-1.000,2.090) (-1.048,2.050) (-1.100,2.002) (-1.150,1.951) (-1.195,1.900) (-1.236,1.850) (-1.273,1.800) (-1.307,1.750) (-1.337,1.700) (-1.365,1.650) (-1.390,1.600) (-1.413,1.550) (-1.433,1.500) (-1.450,1.453) (-1.466,1.400) (-1.480,1.350) (-1.491,1.300) (-1.500,1.250) (-1.506,1.200) (-1.511,1.150) (-1.513,1.100) (-1.514,1.050) (-1.512,1.000) (-1.507,0.950) (-1.501,0.900) (-1.492,0.850) (-1.480,0.800) (-1.466,0.750) (-1.449,0.700) (-1.429,0.650) (-1.406,0.600) (-1.379,0.550) (-1.350,0.502) (-1.314,0.450) (-1.275,0.401) (-1.251,0.374)}
\EQ{(1.250,0.405) (1.285,0.450) (1.318,0.500) (1.347,0.550) (1.373,0.600) (1.395,0.650) (1.413,0.700) (1.429,0.750) (1.442,0.800) (1.452,0.850) (1.460,0.900) (1.465,0.950) (1.468,1.000) (1.468,1.050) (1.467,1.100) (1.463,1.150) (1.457,1.200) (1.448,1.250) (1.438,1.300) (1.425,1.350) (1.410,1.400) (1.392,1.450) (1.372,1.500) (1.350,1.550) (1.325,1.600) (1.300,1.645) (1.266,1.700) (1.232,1.750) (1.200,1.793) (1.154,1.850) (1.109,1.900) (1.059,1.950) (1.004,2.000) (0.950,2.044) (0.909,2.075) (0.850,2.116) (0.800,2.147) (0.751,2.175) (0.704,2.200) (0.652,2.225) (0.600,2.248) (0.550,2.268) (0.500,2.285) (0.450,2.301) (0.400,2.315) (0.350,2.328) (0.300,2.338) (0.250,2.347) (0.200,2.354) (0.150,2.360) (0.100,2.364) (0.050,2.366) (0.000,2.367) (-0.050,2.366) (-0.100,2.364) (-0.150,2.360) (-0.200,2.354) (-0.250,2.347) (-0.300,2.338) (-0.350,2.328) (-0.400,2.315) (-0.450,2.301) (-0.500,2.285) (-0.550,2.268) (-0.600,2.248) (-0.650,2.226) (-0.700,2.202) (-0.750,2.176) (-0.795,2.150) (-0.850,2.116) (-0.900,2.081) (-0.942,2.050) (-1.000,2.003) (-1.050,1.958) (-1.100,1.909) (-1.150,1.854) (-1.195,1.800) (-1.225,1.760) (-1.266,1.700) (-1.297,1.650) (-1.325,1.600) (-1.350,1.550) (-1.372,1.500) (-1.392,1.450) (-1.410,1.400) (-1.425,1.350) (-1.438,1.300) (-1.448,1.250) (-1.457,1.200) (-1.463,1.150) (-1.467,1.100) (-1.468,1.050) (-1.468,1.000) (-1.465,0.950) (-1.460,0.900) (-1.452,0.850) (-1.442,0.800) (-1.429,0.750) (-1.413,0.700) (-1.395,0.650) (-1.373,0.600) (-1.350,0.555) (-1.325,0.512) (-1.285,0.450) (-1.250,0.405)}
\EQ{(1.244,0.431) (1.274,0.475) (1.300,0.518) (1.325,0.565) (1.350,0.619) (1.372,0.675) (1.388,0.725) (1.400,0.775) (1.411,0.825) (1.419,0.875) (1.424,0.925) (1.427,0.975) (1.427,1.025) (1.425,1.075) (1.421,1.125) (1.414,1.175) (1.406,1.225) (1.395,1.275) (1.381,1.325) (1.365,1.375) (1.347,1.425) (1.326,1.475) (1.303,1.525) (1.277,1.575) (1.250,1.621) (1.215,1.675) (1.180,1.725) (1.141,1.775) (1.100,1.822) (1.050,1.874) (1.000,1.922) (0.950,1.965) (0.906,2.000) (0.850,2.040) (0.800,2.073) (0.755,2.100) (0.709,2.125) (0.660,2.150) (0.605,2.175) (0.550,2.198) (0.500,2.216) (0.450,2.233) (0.400,2.247) (0.350,2.260) (0.300,2.271) (0.250,2.280) (0.200,2.287) (0.150,2.293) (0.100,2.297) (0.050,2.299) (0.000,2.300) (-0.050,2.299) (-0.100,2.297) (-0.150,2.293) (-0.200,2.287) (-0.250,2.280) (-0.300,2.271) (-0.350,2.260) (-0.400,2.247) (-0.450,2.233) (-0.500,2.216) (-0.550,2.198) (-0.600,2.177) (-0.650,2.155) (-0.700,2.130) (-0.750,2.103) (-0.796,2.075) (-0.850,2.040) (-0.900,2.004) (-0.950,1.965) (-0.997,1.925) (-1.049,1.875) (-1.097,1.825) (-1.141,1.775) (-1.175,1.731) (-1.215,1.675) (-1.248,1.625) (-1.275,1.578) (-1.300,1.531) (-1.325,1.478) (-1.347,1.425) (-1.365,1.375) (-1.381,1.325) (-1.395,1.275) (-1.406,1.225) (-1.414,1.175) (-1.421,1.125) (-1.425,1.075) (-1.427,1.025) (-1.427,0.975) (-1.424,0.925) (-1.419,0.875) (-1.411,0.825) (-1.400,0.775) (-1.388,0.725) (-1.372,0.675) (-1.352,0.625) (-1.330,0.575) (-1.304,0.525) (-1.275,0.477) (-1.244,0.431)}
\EQ{(1.229,0.446) (1.263,0.500) (1.290,0.550) (1.313,0.600) (1.333,0.650) (1.349,0.700) (1.362,0.750) (1.373,0.800) (1.381,0.850) (1.386,0.900) (1.389,0.950) (1.389,1.000) (1.387,1.050) (1.383,1.100) (1.376,1.150) (1.367,1.200) (1.355,1.250) (1.341,1.300) (1.325,1.349) (1.306,1.400) (1.284,1.450) (1.260,1.500) (1.233,1.550) (1.203,1.600) (1.175,1.642) (1.132,1.700) (1.100,1.740) (1.050,1.796) (1.000,1.846) (0.950,1.891) (0.909,1.925) (0.850,1.970) (0.806,2.000) (0.750,2.035) (0.700,2.063) (0.650,2.089) (0.600,2.112) (0.550,2.133) (0.500,2.152) (0.450,2.169) (0.400,2.184) (0.350,2.197) (0.300,2.208) (0.250,2.218) (0.200,2.225) (0.150,2.231) (0.100,2.236) (0.050,2.238) (0.000,2.239) (-0.050,2.238) (-0.100,2.236) (-0.150,2.231) (-0.200,2.225) (-0.250,2.218) (-0.300,2.208) (-0.350,2.197) (-0.400,2.184) (-0.450,2.169) (-0.500,2.152) (-0.550,2.133) (-0.600,2.112) (-0.650,2.089) (-0.700,2.063) (-0.750,2.035) (-0.800,2.004) (-0.843,1.975) (-0.900,1.932) (-0.940,1.900) (-0.996,1.850) (-1.046,1.800) (-1.091,1.750) (-1.125,1.709) (-1.169,1.650) (-1.200,1.604) (-1.233,1.550) (-1.260,1.500) (-1.284,1.450) (-1.306,1.400) (-1.325,1.350) (-1.341,1.300) (-1.355,1.250) (-1.367,1.200) (-1.376,1.150) (-1.383,1.100) (-1.387,1.050) (-1.389,1.000) (-1.389,0.950) (-1.386,0.900) (-1.381,0.850) (-1.373,0.800) (-1.362,0.750) (-1.349,0.700) (-1.333,0.650) (-1.313,0.600) (-1.290,0.550) (-1.263,0.500) (-1.232,0.450) (-1.229,0.446)}
\EQ{(1.224,0.476) (1.250,0.522) (1.275,0.574) (1.296,0.625) (1.313,0.675) (1.326,0.725) (1.337,0.775) (1.345,0.825) (1.350,0.875) (1.353,0.925) (1.354,0.975) (1.352,1.025) (1.347,1.075) (1.340,1.125) (1.331,1.175) (1.319,1.225) (1.305,1.275) (1.288,1.325) (1.269,1.375) (1.246,1.425) (1.221,1.475) (1.194,1.525) (1.162,1.575) (1.128,1.625) (1.089,1.675) (1.050,1.721) (1.000,1.774) (0.950,1.822) (0.900,1.865) (0.856,1.900) (0.800,1.940) (0.750,1.972) (0.702,2.000) (0.656,2.025) (0.604,2.050) (0.550,2.074) (0.500,2.093) (0.450,2.111) (0.400,2.126) (0.350,2.139) (0.300,2.151) (0.250,2.161) (0.200,2.168) (0.150,2.174) (0.100,2.179) (0.050,2.181) (0.000,2.182) (-0.050,2.181) (-0.100,2.179) (-0.150,2.174) (-0.200,2.168) (-0.250,2.161) (-0.300,2.151) (-0.350,2.139) (-0.400,2.126) (-0.450,2.111) (-0.500,2.093) (-0.547,2.075) (-0.600,2.052) (-0.650,2.028) (-0.700,2.001) (-0.745,1.975) (-0.800,1.940) (-0.850,1.904) (-0.900,1.865) (-0.947,1.825) (-0.999,1.775) (-1.047,1.725) (-1.089,1.675) (-1.125,1.629) (-1.162,1.575) (-1.194,1.525) (-1.221,1.475) (-1.246,1.425) (-1.269,1.375) (-1.288,1.325) (-1.305,1.275) (-1.319,1.225) (-1.331,1.175) (-1.340,1.125) (-1.347,1.075) (-1.352,1.025) (-1.354,0.975) (-1.353,0.925) (-1.350,0.875) (-1.345,0.825) (-1.337,0.775) (-1.326,0.725) (-1.313,0.675) (-1.296,0.625) (-1.275,0.575) (-1.252,0.525) (-1.225,0.477)}

    \draw[mmaEQsolo] (-10,0) -- (-1.2,0) (1.2,0) -- (10,0);}
  \marcoMMA{10}{5}{$x$ (u.\,arb.)}{$z$ (u.\,arb.)}
\end{tikzpicture}
\end{minipage}\hfill
\begin{minipage}[t]{0.325\textwidth}\centering
\rotulo{Término cuadrupolar ($\ell=2$)}
\begin{tikzpicture}[x=0.198cm, y=0.198cm]
  \panelMMA{10}{white}{0}{% Cuadrupolo puntual - equipotenciales
% Generado por scripts/generar_datos.py (no editar a mano).
\EQ{(-0.220,-0.155) (-0.275,-0.195) (-0.319,-0.225) (-0.375,-0.266) (-0.425,-0.300) (-0.475,-0.336) (-0.525,-0.371) (-0.568,-0.400) (-0.625,-0.441) (-0.675,-0.475) (-0.725,-0.510) (-0.775,-0.544) (-0.821,-0.575) (-0.875,-0.612) (-0.925,-0.646) (-0.969,-0.675) (-1.025,-0.713) (-1.075,-0.745) (-1.121,-0.775) (-1.175,-0.810) (-1.225,-0.841) (-1.275,-0.872) (-1.321,-0.900) (-1.375,-0.933) (-1.425,-0.963) (-1.475,-0.991) (-1.525,-1.020) (-1.575,-1.047) (-1.625,-1.074) (-1.674,-1.100) (-1.723,-1.125) (-1.774,-1.150) (-1.825,-1.175) (-1.875,-1.198) (-1.925,-1.220) (-1.975,-1.241) (-2.025,-1.262) (-2.075,-1.281) (-2.125,-1.300) (-2.175,-1.317) (-2.225,-1.334) (-2.275,-1.349) (-2.325,-1.364) (-2.375,-1.377) (-2.425,-1.389) (-2.475,-1.400) (-2.525,-1.410) (-2.575,-1.419) (-2.625,-1.426) (-2.675,-1.433) (-2.725,-1.438) (-2.775,-1.441) (-2.825,-1.444) (-2.875,-1.445) (-2.925,-1.445) (-2.975,-1.443) (-3.025,-1.440) (-3.075,-1.435) (-3.125,-1.429) (-3.175,-1.421) (-3.225,-1.412) (-3.275,-1.401) (-3.325,-1.388) (-3.375,-1.373) (-3.425,-1.357) (-3.475,-1.338) (-3.525,-1.318) (-3.575,-1.295) (-3.625,-1.270) (-3.675,-1.242) (-3.725,-1.212) (-3.775,-1.179) (-3.825,-1.142) (-3.875,-1.102) (-3.925,-1.058) (-3.975,-1.010) (-4.025,-0.956) (-4.072,-0.900) (-4.110,-0.850) (-4.145,-0.800) (-4.175,-0.753) (-4.206,-0.700) (-4.232,-0.650) (-4.256,-0.600) (-4.277,-0.550) (-4.297,-0.500) (-4.314,-0.450) (-4.329,-0.400) (-4.343,-0.350) (-4.354,-0.300) (-4.364,-0.250) (-4.371,-0.200) (-4.377,-0.150) (-4.382,-0.100) (-4.384,-0.050) (-4.385,0.000) (-4.384,0.050) (-4.382,0.100) (-4.377,0.150) (-4.371,0.200) (-4.364,0.250) (-4.354,0.300) (-4.343,0.350) (-4.329,0.400) (-4.314,0.450) (-4.297,0.500) (-4.277,0.550) (-4.256,0.600) (-4.232,0.650) (-4.206,0.700) (-4.177,0.750) (-4.150,0.792) (-4.110,0.850) (-4.075,0.896) (-4.030,0.950) (-3.985,1.000) (-3.934,1.050) (-3.878,1.100) (-3.825,1.142) (-3.780,1.175) (-3.725,1.212) (-3.675,1.242) (-3.625,1.270) (-3.575,1.295) (-3.525,1.318) (-3.475,1.338) (-3.425,1.357) (-3.375,1.373) (-3.325,1.388) (-3.275,1.401) (-3.225,1.412) (-3.175,1.421) (-3.125,1.429) (-3.075,1.435) (-3.025,1.440) (-2.975,1.443) (-2.925,1.445) (-2.875,1.445) (-2.825,1.444) (-2.775,1.441) (-2.725,1.438) (-2.675,1.433) (-2.625,1.426) (-2.575,1.419) (-2.525,1.410) (-2.475,1.400) (-2.425,1.389) (-2.375,1.377) (-2.325,1.364) (-2.275,1.349) (-2.225,1.334) (-2.175,1.317) (-2.126,1.300) (-2.075,1.281) (-2.025,1.262) (-1.975,1.241) (-1.925,1.220) (-1.875,1.198) (-1.826,1.175) (-1.775,1.151) (-1.725,1.126) (-1.675,1.100) (-1.627,1.075) (-1.580,1.050) (-1.535,1.025) (-1.491,1.000) (-1.447,0.975) (-1.400,0.948) (-1.350,0.918) (-1.300,0.888) (-1.250,0.857) (-1.200,0.825) (-1.150,0.794) (-1.100,0.762) (-1.050,0.729) (-1.007,0.700) (-0.950,0.663) (-0.900,0.629) (-0.858,0.600) (-0.800,0.561) (-0.750,0.527) (-0.700,0.493) (-0.650,0.458) (-0.604,0.425) (-0.550,0.388) (-0.500,0.353) (-0.450,0.318) (-0.400,0.283) (-0.354,0.250) (-0.300,0.213) (-0.250,0.177) (-0.220,0.155)}
\EQ{(0.206,0.144) (0.248,0.175) (0.300,0.213) (0.350,0.248) (0.400,0.283) (0.450,0.318) (0.496,0.350) (0.550,0.388) (0.600,0.423) (0.650,0.458) (0.700,0.493) (0.748,0.525) (0.800,0.561) (0.850,0.595) (0.895,0.625) (0.950,0.663) (1.000,0.696) (1.045,0.725) (1.100,0.762) (1.150,0.794) (1.200,0.825) (1.250,0.857) (1.300,0.888) (1.350,0.918) (1.400,0.948) (1.447,0.975) (1.491,1.000) (1.535,1.025) (1.580,1.050) (1.625,1.074) (1.674,1.100) (1.723,1.125) (1.774,1.150) (1.825,1.175) (1.875,1.198) (1.925,1.220) (1.975,1.241) (2.025,1.262) (2.075,1.281) (2.125,1.300) (2.175,1.317) (2.225,1.334) (2.275,1.349) (2.325,1.364) (2.375,1.377) (2.425,1.389) (2.475,1.400) (2.525,1.410) (2.575,1.419) (2.625,1.426) (2.675,1.433) (2.725,1.438) (2.775,1.441) (2.825,1.444) (2.875,1.445) (2.925,1.445) (2.975,1.443) (3.025,1.440) (3.075,1.435) (3.125,1.429) (3.175,1.421) (3.225,1.412) (3.275,1.401) (3.325,1.388) (3.375,1.373) (3.425,1.357) (3.475,1.338) (3.525,1.318) (3.575,1.295) (3.625,1.270) (3.675,1.242) (3.725,1.212) (3.775,1.179) (3.825,1.142) (3.875,1.102) (3.925,1.058) (3.975,1.010) (4.025,0.956) (4.072,0.900) (4.110,0.850) (4.145,0.800) (4.175,0.753) (4.206,0.700) (4.232,0.650) (4.256,0.600) (4.277,0.550) (4.297,0.500) (4.314,0.450) (4.329,0.400) (4.343,0.350) (4.354,0.300) (4.364,0.250) (4.371,0.200) (4.377,0.150) (4.382,0.100) (4.384,0.050) (4.385,0.000) (4.384,-0.050) (4.382,-0.100) (4.377,-0.150) (4.371,-0.200) (4.364,-0.250) (4.354,-0.300) (4.343,-0.350) (4.329,-0.400) (4.314,-0.450) (4.297,-0.500) (4.277,-0.550) (4.256,-0.600) (4.232,-0.650) (4.206,-0.700) (4.177,-0.750) (4.150,-0.792) (4.110,-0.850) (4.075,-0.896) (4.030,-0.950) (3.985,-1.000) (3.934,-1.050) (3.878,-1.100) (3.825,-1.142) (3.780,-1.175) (3.725,-1.212) (3.675,-1.242) (3.625,-1.270) (3.575,-1.295) (3.525,-1.318) (3.475,-1.338) (3.425,-1.357) (3.375,-1.373) (3.325,-1.388) (3.275,-1.401) (3.225,-1.412) (3.175,-1.421) (3.125,-1.429) (3.075,-1.435) (3.025,-1.440) (2.975,-1.443) (2.925,-1.445) (2.875,-1.445) (2.825,-1.444) (2.775,-1.441) (2.725,-1.438) (2.675,-1.433) (2.625,-1.426) (2.575,-1.419) (2.525,-1.410) (2.475,-1.400) (2.425,-1.389) (2.375,-1.377) (2.325,-1.364) (2.275,-1.349) (2.225,-1.334) (2.175,-1.317) (2.126,-1.300) (2.075,-1.281) (2.025,-1.262) (1.975,-1.241) (1.925,-1.220) (1.875,-1.198) (1.826,-1.175) (1.775,-1.151) (1.725,-1.126) (1.675,-1.100) (1.627,-1.075) (1.580,-1.050) (1.535,-1.025) (1.491,-1.000) (1.447,-0.975) (1.400,-0.948) (1.350,-0.918) (1.300,-0.888) (1.250,-0.857) (1.200,-0.825) (1.150,-0.794) (1.100,-0.762) (1.050,-0.729) (1.007,-0.700) (0.950,-0.663) (0.900,-0.629) (0.858,-0.600) (0.800,-0.561) (0.750,-0.527) (0.700,-0.493) (0.650,-0.458) (0.604,-0.425) (0.550,-0.388) (0.500,-0.353) (0.450,-0.318) (0.400,-0.283) (0.354,-0.250) (0.300,-0.213) (0.250,-0.177) (0.206,-0.144)}
\EQ{(-0.220,-0.155) (-0.275,-0.195) (-0.319,-0.225) (-0.375,-0.266) (-0.425,-0.300) (-0.475,-0.336) (-0.525,-0.371) (-0.568,-0.400) (-0.625,-0.441) (-0.675,-0.475) (-0.725,-0.510) (-0.775,-0.545) (-0.820,-0.575) (-0.875,-0.613) (-0.925,-0.647) (-0.968,-0.675) (-1.025,-0.714) (-1.075,-0.747) (-1.119,-0.775) (-1.175,-0.812) (-1.225,-0.844) (-1.275,-0.875) (-1.325,-0.906) (-1.375,-0.937) (-1.425,-0.967) (-1.475,-0.996) (-1.525,-1.025) (-1.569,-1.050) (-1.614,-1.075) (-1.660,-1.100) (-1.707,-1.125) (-1.755,-1.150) (-1.800,-1.173) (-1.850,-1.197) (-1.900,-1.221) (-1.950,-1.244) (-2.000,-1.266) (-2.050,-1.288) (-2.100,-1.308) (-2.150,-1.327) (-2.200,-1.346) (-2.250,-1.363) (-2.300,-1.380) (-2.350,-1.396) (-2.400,-1.410) (-2.450,-1.424) (-2.500,-1.436) (-2.550,-1.447) (-2.600,-1.458) (-2.650,-1.467) (-2.700,-1.475) (-2.750,-1.481) (-2.800,-1.487) (-2.850,-1.491) (-2.900,-1.494) (-2.950,-1.496) (-3.000,-1.497) (-3.050,-1.496) (-3.100,-1.493) (-3.150,-1.490) (-3.200,-1.485) (-3.250,-1.478) (-3.300,-1.470) (-3.350,-1.460) (-3.400,-1.449) (-3.450,-1.436) (-3.500,-1.421) (-3.550,-1.405) (-3.600,-1.386) (-3.650,-1.365) (-3.700,-1.343) (-3.750,-1.318) (-3.800,-1.290) (-3.850,-1.260) (-3.900,-1.228) (-3.950,-1.192) (-4.000,-1.153) (-4.050,-1.110) (-4.088,-1.075) (-4.125,-1.038) (-4.160,-1.000) (-4.200,-0.954) (-4.243,-0.900) (-4.275,-0.855) (-4.312,-0.800) (-4.342,-0.750) (-4.369,-0.700) (-4.395,-0.650) (-4.417,-0.600) (-4.438,-0.550) (-4.457,-0.500) (-4.473,-0.450) (-4.488,-0.400) (-4.501,-0.350) (-4.512,-0.300) (-4.521,-0.250) (-4.529,-0.200) (-4.534,-0.150) (-4.539,-0.100) (-4.541,-0.050) (-4.542,0.000) (-4.541,0.050) (-4.539,0.100) (-4.534,0.150) (-4.529,0.200) (-4.521,0.250) (-4.512,0.300) (-4.501,0.350) (-4.488,0.400) (-4.473,0.450) (-4.457,0.500) (-4.438,0.550) (-4.417,0.600) (-4.395,0.650) (-4.369,0.700) (-4.342,0.750) (-4.312,0.800) (-4.279,0.850) (-4.243,0.900) (-4.203,0.950) (-4.160,1.000) (-4.125,1.038) (-4.088,1.075) (-4.050,1.110) (-4.003,1.150) (-3.950,1.192) (-3.904,1.225) (-3.850,1.260) (-3.800,1.290) (-3.750,1.318) (-3.700,1.343) (-3.650,1.365) (-3.600,1.386) (-3.550,1.405) (-3.500,1.421) (-3.450,1.436) (-3.400,1.449) (-3.350,1.460) (-3.300,1.470) (-3.250,1.478) (-3.200,1.485) (-3.150,1.490) (-3.100,1.493) (-3.050,1.496) (-3.000,1.497) (-2.950,1.496) (-2.900,1.494) (-2.850,1.491) (-2.800,1.487) (-2.750,1.481) (-2.700,1.475) (-2.650,1.467) (-2.600,1.458) (-2.550,1.447) (-2.500,1.436) (-2.450,1.424) (-2.400,1.410) (-2.350,1.396) (-2.300,1.380) (-2.250,1.363) (-2.200,1.346) (-2.150,1.327) (-2.100,1.308) (-2.050,1.288) (-2.000,1.266) (-1.950,1.244) (-1.900,1.221) (-1.850,1.197) (-1.804,1.175) (-1.755,1.150) (-1.707,1.125) (-1.660,1.100) (-1.614,1.075) (-1.569,1.050) (-1.525,1.025) (-1.482,1.000) (-1.425,0.967) (-1.375,0.937) (-1.325,0.906) (-1.275,0.875) (-1.225,0.844) (-1.175,0.812) (-1.125,0.779) (-1.081,0.750) (-1.025,0.714) (-0.975,0.680) (-0.930,0.650) (-0.875,0.613) (-0.825,0.579) (-0.784,0.550) (-0.725,0.510) (-0.675,0.475) (-0.625,0.441) (-0.575,0.406) (-0.532,0.375) (-0.475,0.336) (-0.425,0.300) (-0.375,0.266) (-0.325,0.230) (-0.284,0.200) (-0.225,0.160) (-0.220,0.155)}
\EQ{(0.206,0.144) (0.248,0.175) (0.300,0.213) (0.350,0.248) (0.400,0.283) (0.450,0.318) (0.496,0.350) (0.550,0.388) (0.600,0.423) (0.650,0.458) (0.700,0.493) (0.747,0.525) (0.800,0.562) (0.850,0.596) (0.894,0.625) (0.950,0.664) (1.000,0.697) (1.043,0.725) (1.100,0.763) (1.150,0.796) (1.196,0.825) (1.250,0.859) (1.300,0.891) (1.350,0.921) (1.397,0.950) (1.450,0.982) (1.500,1.011) (1.550,1.040) (1.600,1.068) (1.650,1.095) (1.700,1.122) (1.750,1.148) (1.800,1.173) (1.850,1.197) (1.900,1.221) (1.950,1.244) (2.000,1.266) (2.050,1.288) (2.100,1.308) (2.150,1.327) (2.200,1.346) (2.250,1.363) (2.300,1.380) (2.350,1.396) (2.400,1.410) (2.450,1.424) (2.500,1.436) (2.550,1.447) (2.600,1.458) (2.650,1.467) (2.700,1.475) (2.750,1.481) (2.800,1.487) (2.850,1.491) (2.900,1.494) (2.950,1.496) (3.000,1.497) (3.050,1.496) (3.100,1.493) (3.150,1.490) (3.200,1.485) (3.250,1.478) (3.300,1.470) (3.350,1.460) (3.400,1.449) (3.450,1.436) (3.500,1.421) (3.550,1.405) (3.600,1.386) (3.650,1.365) (3.700,1.343) (3.750,1.318) (3.800,1.290) (3.850,1.260) (3.900,1.228) (3.950,1.192) (4.000,1.153) (4.050,1.110) (4.088,1.075) (4.125,1.038) (4.160,1.000) (4.200,0.954) (4.243,0.900) (4.275,0.855) (4.312,0.800) (4.342,0.750) (4.369,0.700) (4.395,0.650) (4.417,0.600) (4.438,0.550) (4.457,0.500) (4.473,0.450) (4.488,0.400) (4.501,0.350) (4.512,0.300) (4.521,0.250) (4.529,0.200) (4.534,0.150) (4.539,0.100) (4.541,0.050) (4.542,0.000) (4.541,-0.050) (4.539,-0.100) (4.534,-0.150) (4.529,-0.200) (4.521,-0.250) (4.512,-0.300) (4.501,-0.350) (4.488,-0.400) (4.473,-0.450) (4.457,-0.500) (4.438,-0.550) (4.417,-0.600) (4.395,-0.650) (4.369,-0.700) (4.342,-0.750) (4.312,-0.800) (4.279,-0.850) (4.243,-0.900) (4.203,-0.950) (4.160,-1.000) (4.125,-1.038) (4.088,-1.075) (4.050,-1.110) (4.003,-1.150) (3.950,-1.192) (3.904,-1.225) (3.850,-1.260) (3.800,-1.290) (3.750,-1.318) (3.700,-1.343) (3.650,-1.365) (3.600,-1.386) (3.550,-1.405) (3.500,-1.421) (3.450,-1.436) (3.400,-1.449) (3.350,-1.460) (3.300,-1.470) (3.250,-1.478) (3.200,-1.485) (3.150,-1.490) (3.100,-1.493) (3.050,-1.496) (3.000,-1.497) (2.950,-1.496) (2.900,-1.494) (2.850,-1.491) (2.800,-1.487) (2.750,-1.481) (2.700,-1.475) (2.650,-1.467) (2.600,-1.458) (2.550,-1.447) (2.500,-1.436) (2.450,-1.424) (2.400,-1.410) (2.350,-1.396) (2.300,-1.380) (2.250,-1.363) (2.200,-1.346) (2.150,-1.327) (2.100,-1.308) (2.050,-1.288) (2.000,-1.266) (1.950,-1.244) (1.900,-1.221) (1.850,-1.197) (1.804,-1.175) (1.755,-1.150) (1.707,-1.125) (1.660,-1.100) (1.614,-1.075) (1.569,-1.050) (1.525,-1.025) (1.482,-1.000) (1.425,-0.967) (1.375,-0.937) (1.325,-0.906) (1.275,-0.875) (1.225,-0.844) (1.175,-0.812) (1.125,-0.779) (1.081,-0.750) (1.025,-0.714) (0.975,-0.680) (0.930,-0.650) (0.875,-0.613) (0.825,-0.579) (0.784,-0.550) (0.725,-0.510) (0.675,-0.475) (0.625,-0.441) (0.575,-0.406) (0.532,-0.375) (0.475,-0.336) (0.425,-0.300) (0.375,-0.266) (0.325,-0.230) (0.284,-0.200) (0.225,-0.160) (0.206,-0.144)}
\EQ{(-0.220,-0.155) (-0.275,-0.195) (-0.319,-0.225) (-0.375,-0.266) (-0.425,-0.300) (-0.475,-0.336) (-0.525,-0.371) (-0.568,-0.400) (-0.625,-0.441) (-0.674,-0.475) (-0.725,-0.510) (-0.775,-0.545) (-0.819,-0.575) (-0.875,-0.614) (-0.925,-0.647) (-0.975,-0.681) (-1.025,-0.715) (-1.075,-0.748) (-1.125,-0.781) (-1.175,-0.814) (-1.225,-0.846) (-1.271,-0.875) (-1.325,-0.910) (-1.375,-0.941) (-1.425,-0.971) (-1.473,-1.000) (-1.525,-1.031) (-1.575,-1.060) (-1.625,-1.089) (-1.675,-1.116) (-1.725,-1.144) (-1.775,-1.170) (-1.825,-1.196) (-1.875,-1.221) (-1.925,-1.246) (-1.975,-1.270) (-2.025,-1.293) (-2.075,-1.315) (-2.125,-1.336) (-2.175,-1.357) (-2.222,-1.375) (-2.275,-1.395) (-2.325,-1.413) (-2.375,-1.430) (-2.425,-1.446) (-2.475,-1.461) (-2.525,-1.474) (-2.575,-1.487) (-2.625,-1.499) (-2.675,-1.510) (-2.725,-1.520) (-2.775,-1.529) (-2.825,-1.536) (-2.875,-1.542) (-2.925,-1.548) (-2.975,-1.552) (-3.025,-1.555) (-3.075,-1.556) (-3.125,-1.556) (-3.175,-1.556) (-3.225,-1.553) (-3.275,-1.550) (-3.325,-1.545) (-3.375,-1.538) (-3.425,-1.530) (-3.475,-1.521) (-3.525,-1.510) (-3.575,-1.497) (-3.625,-1.483) (-3.675,-1.467) (-3.725,-1.449) (-3.775,-1.429) (-3.825,-1.407) (-3.875,-1.383) (-3.925,-1.357) (-3.975,-1.329) (-4.021,-1.300) (-4.075,-1.264) (-4.125,-1.227) (-4.175,-1.187) (-4.217,-1.150) (-4.270,-1.100) (-4.318,-1.050) (-4.362,-1.000) (-4.400,-0.952) (-4.439,-0.900) (-4.473,-0.850) (-4.500,-0.807) (-4.533,-0.750) (-4.559,-0.700) (-4.583,-0.650) (-4.605,-0.600) (-4.624,-0.550) (-4.642,-0.500) (-4.658,-0.450) (-4.672,-0.400) (-4.684,-0.350) (-4.695,-0.300) (-4.704,-0.250) (-4.711,-0.200) (-4.717,-0.150) (-4.720,-0.100) (-4.723,-0.050) (-4.724,0.000) (-4.723,0.050) (-4.720,0.100) (-4.717,0.150) (-4.711,0.200) (-4.704,0.250) (-4.695,0.300) (-4.684,0.350) (-4.672,0.400) (-4.658,0.450) (-4.642,0.500) (-4.625,0.548) (-4.605,0.600) (-4.583,0.650) (-4.559,0.700) (-4.533,0.750) (-4.504,0.800) (-4.475,0.847) (-4.439,0.900) (-4.402,0.950) (-4.362,1.000) (-4.325,1.042) (-4.275,1.095) (-4.225,1.143) (-4.175,1.187) (-4.127,1.225) (-4.075,1.264) (-4.025,1.298) (-3.981,1.325) (-3.925,1.357) (-3.875,1.383) (-3.825,1.407) (-3.775,1.429) (-3.725,1.449) (-3.675,1.467) (-3.625,1.483) (-3.575,1.497) (-3.525,1.510) (-3.475,1.521) (-3.425,1.530) (-3.375,1.538) (-3.325,1.545) (-3.275,1.550) (-3.225,1.553) (-3.175,1.556) (-3.125,1.556) (-3.075,1.556) (-3.025,1.555) (-2.975,1.552) (-2.925,1.548) (-2.875,1.542) (-2.825,1.536) (-2.775,1.529) (-2.725,1.520) (-2.675,1.510) (-2.625,1.499) (-2.575,1.487) (-2.525,1.474) (-2.475,1.461) (-2.425,1.446) (-2.375,1.430) (-2.325,1.413) (-2.275,1.395) (-2.225,1.376) (-2.175,1.357) (-2.125,1.336) (-2.075,1.315) (-2.025,1.293) (-1.975,1.270) (-1.925,1.246) (-1.875,1.221) (-1.825,1.196) (-1.775,1.170) (-1.725,1.144) (-1.675,1.116) (-1.625,1.089) (-1.575,1.060) (-1.525,1.031) (-1.475,1.001) (-1.432,0.975) (-1.375,0.941) (-1.325,0.910) (-1.275,0.878) (-1.232,0.850) (-1.175,0.814) (-1.125,0.781) (-1.078,0.750) (-1.025,0.715) (-0.975,0.681) (-0.929,0.650) (-0.875,0.614) (-0.825,0.579) (-0.783,0.550) (-0.725,0.510) (-0.675,0.475) (-0.625,0.441) (-0.575,0.406) (-0.532,0.375) (-0.475,0.336) (-0.425,0.300) (-0.375,0.266) (-0.325,0.230) (-0.284,0.200) (-0.225,0.160) (-0.220,0.155)}
\EQ{(0.206,0.144) (0.248,0.175) (0.300,0.213) (0.350,0.248) (0.400,0.283) (0.450,0.318) (0.496,0.350) (0.550,0.389) (0.600,0.423) (0.650,0.458) (0.700,0.493) (0.747,0.525) (0.800,0.562) (0.850,0.596) (0.893,0.625) (0.950,0.665) (1.000,0.698) (1.050,0.732) (1.100,0.765) (1.150,0.797) (1.193,0.825) (1.250,0.862) (1.300,0.894) (1.350,0.925) (1.400,0.956) (1.450,0.986) (1.500,1.016) (1.550,1.046) (1.600,1.074) (1.646,1.100) (1.691,1.125) (1.737,1.150) (1.785,1.175) (1.833,1.200) (1.883,1.225) (1.934,1.250) (1.987,1.275) (2.042,1.300) (2.099,1.325) (2.150,1.346) (2.200,1.367) (2.250,1.386) (2.300,1.404) (2.350,1.421) (2.400,1.438) (2.450,1.453) (2.500,1.468) (2.550,1.481) (2.600,1.493) (2.650,1.505) (2.700,1.515) (2.750,1.524) (2.800,1.532) (2.850,1.539) (2.900,1.545) (2.950,1.550) (3.000,1.553) (3.050,1.556) (3.100,1.556) (3.150,1.556) (3.200,1.555) (3.250,1.552) (3.300,1.547) (3.350,1.542) (3.400,1.534) (3.450,1.526) (3.500,1.515) (3.550,1.504) (3.600,1.490) (3.650,1.475) (3.700,1.458) (3.750,1.439) (3.800,1.418) (3.850,1.396) (3.900,1.371) (3.950,1.343) (4.000,1.313) (4.050,1.281) (4.094,1.250) (4.150,1.207) (4.189,1.175) (4.244,1.125) (4.294,1.075) (4.340,1.025) (4.375,0.984) (4.421,0.925) (4.450,0.884) (4.489,0.825) (4.519,0.775) (4.546,0.725) (4.571,0.675) (4.594,0.625) (4.615,0.575) (4.633,0.525) (4.650,0.475) (4.665,0.425) (4.678,0.375) (4.690,0.325) (4.699,0.275) (4.708,0.225) (4.714,0.175) (4.719,0.125) (4.722,0.075) (4.723,0.025) (4.723,-0.050) (4.720,-0.100) (4.717,-0.150) (4.711,-0.200) (4.704,-0.250) (4.695,-0.300) (4.684,-0.350) (4.672,-0.400) (4.658,-0.450) (4.642,-0.500) (4.625,-0.548) (4.605,-0.600) (4.583,-0.650) (4.559,-0.700) (4.533,-0.750) (4.504,-0.800) (4.475,-0.847) (4.439,-0.900) (4.402,-0.950) (4.362,-1.000) (4.325,-1.042) (4.275,-1.095) (4.225,-1.143) (4.175,-1.187) (4.127,-1.225) (4.075,-1.264) (4.025,-1.298) (3.981,-1.325) (3.925,-1.357) (3.875,-1.383) (3.825,-1.407) (3.775,-1.429) (3.725,-1.449) (3.675,-1.467) (3.625,-1.483) (3.575,-1.497) (3.525,-1.510) (3.475,-1.521) (3.425,-1.530) (3.375,-1.538) (3.325,-1.545) (3.275,-1.550) (3.225,-1.553) (3.175,-1.556) (3.125,-1.556) (3.075,-1.556) (3.025,-1.555) (2.975,-1.552) (2.925,-1.548) (2.875,-1.542) (2.825,-1.536) (2.775,-1.529) (2.725,-1.520) (2.675,-1.510) (2.625,-1.499) (2.575,-1.487) (2.525,-1.474) (2.475,-1.461) (2.425,-1.446) (2.375,-1.430) (2.325,-1.413) (2.275,-1.395) (2.225,-1.376) (2.175,-1.357) (2.125,-1.336) (2.075,-1.315) (2.025,-1.293) (1.975,-1.270) (1.925,-1.246) (1.875,-1.221) (1.825,-1.196) (1.775,-1.170) (1.725,-1.144) (1.675,-1.116) (1.625,-1.089) (1.575,-1.060) (1.525,-1.031) (1.475,-1.001) (1.432,-0.975) (1.375,-0.941) (1.325,-0.910) (1.275,-0.878) (1.232,-0.850) (1.175,-0.814) (1.125,-0.781) (1.078,-0.750) (1.025,-0.715) (0.975,-0.681) (0.929,-0.650) (0.875,-0.614) (0.825,-0.579) (0.783,-0.550) (0.725,-0.510) (0.675,-0.475) (0.625,-0.441) (0.575,-0.406) (0.532,-0.375) (0.475,-0.336) (0.425,-0.300) (0.375,-0.266) (0.325,-0.230) (0.284,-0.200) (0.225,-0.160) (0.206,-0.144)}
\EQ{(-0.220,-0.155) (-0.275,-0.195) (-0.319,-0.225) (-0.375,-0.266) (-0.425,-0.300) (-0.475,-0.336) (-0.525,-0.371) (-0.568,-0.400) (-0.625,-0.441) (-0.674,-0.475) (-0.725,-0.511) (-0.775,-0.545) (-0.819,-0.575) (-0.875,-0.614) (-0.925,-0.648) (-0.975,-0.682) (-1.025,-0.716) (-1.075,-0.750) (-1.125,-0.783) (-1.175,-0.816) (-1.225,-0.848) (-1.275,-0.881) (-1.325,-0.913) (-1.375,-0.944) (-1.424,-0.975) (-1.475,-1.006) (-1.525,-1.037) (-1.575,-1.066) (-1.625,-1.096) (-1.675,-1.124) (-1.721,-1.150) (-1.766,-1.175) (-1.812,-1.200) (-1.859,-1.225) (-1.907,-1.250) (-1.957,-1.275) (-2.008,-1.300) (-2.060,-1.325) (-2.114,-1.350) (-2.171,-1.375) (-2.225,-1.398) (-2.275,-1.419) (-2.325,-1.438) (-2.375,-1.457) (-2.425,-1.475) (-2.475,-1.492) (-2.525,-1.508) (-2.575,-1.523) (-2.625,-1.538) (-2.675,-1.551) (-2.725,-1.563) (-2.775,-1.574) (-2.825,-1.585) (-2.875,-1.594) (-2.925,-1.602) (-2.975,-1.609) (-3.025,-1.615) (-3.075,-1.620) (-3.125,-1.623) (-3.175,-1.626) (-3.225,-1.627) (-3.275,-1.627) (-3.325,-1.626) (-3.375,-1.624) (-3.425,-1.620) (-3.475,-1.615) (-3.525,-1.609) (-3.575,-1.601) (-3.625,-1.592) (-3.675,-1.581) (-3.725,-1.569) (-3.775,-1.555) (-3.825,-1.539) (-3.875,-1.522) (-3.925,-1.503) (-3.975,-1.482) (-4.025,-1.459) (-4.075,-1.434) (-4.125,-1.407) (-4.175,-1.378) (-4.218,-1.350) (-4.275,-1.311) (-4.323,-1.275) (-4.375,-1.232) (-4.425,-1.188) (-4.464,-1.150) (-4.500,-1.113) (-4.535,-1.075) (-4.575,-1.028) (-4.616,-0.975) (-4.650,-0.928) (-4.685,-0.875) (-4.716,-0.825) (-4.744,-0.775) (-4.770,-0.725) (-4.793,-0.675) (-4.815,-0.625) (-4.835,-0.575) (-4.853,-0.525) (-4.869,-0.475) (-4.883,-0.425) (-4.895,-0.375) (-4.906,-0.325) (-4.916,-0.275) (-4.923,-0.225) (-4.929,-0.175) (-4.934,-0.125) (-4.937,-0.075) (-4.938,-0.025) (-4.938,0.050) (-4.936,0.100) (-4.932,0.150) (-4.926,0.200) (-4.920,0.250) (-4.911,0.300) (-4.901,0.350) (-4.889,0.400) (-4.876,0.450) (-4.861,0.500) (-4.844,0.550) (-4.825,0.600) (-4.804,0.650) (-4.782,0.700) (-4.757,0.750) (-4.730,0.800) (-4.701,0.850) (-4.669,0.900) (-4.635,0.950) (-4.600,0.996) (-4.556,1.050) (-4.512,1.100) (-4.475,1.139) (-4.438,1.175) (-4.400,1.210) (-4.354,1.250) (-4.300,1.292) (-4.255,1.325) (-4.200,1.362) (-4.150,1.393) (-4.100,1.421) (-4.050,1.447) (-4.000,1.471) (-3.950,1.493) (-3.900,1.513) (-3.850,1.531) (-3.800,1.547) (-3.750,1.562) (-3.700,1.575) (-3.650,1.586) (-3.600,1.596) (-3.550,1.605) (-3.500,1.612) (-3.450,1.618) (-3.400,1.622) (-3.350,1.625) (-3.300,1.627) (-3.250,1.627) (-3.200,1.627) (-3.150,1.625) (-3.100,1.622) (-3.050,1.617) (-3.000,1.612) (-2.950,1.605) (-2.900,1.598) (-2.850,1.589) (-2.800,1.580) (-2.750,1.569) (-2.700,1.557) (-2.650,1.544) (-2.600,1.531) (-2.550,1.516) (-2.500,1.500) (-2.450,1.484) (-2.400,1.466) (-2.350,1.448) (-2.300,1.429) (-2.250,1.409) (-2.200,1.388) (-2.150,1.366) (-2.100,1.344) (-2.050,1.320) (-2.000,1.296) (-1.950,1.272) (-1.900,1.246) (-1.850,1.220) (-1.800,1.194) (-1.750,1.167) (-1.700,1.139) (-1.650,1.110) (-1.600,1.081) (-1.550,1.051) (-1.506,1.025) (-1.450,0.991) (-1.400,0.960) (-1.350,0.929) (-1.305,0.900) (-1.250,0.865) (-1.200,0.832) (-1.151,0.800) (-1.100,0.766) (-1.050,0.733) (-1.001,0.700) (-0.950,0.665) (-0.900,0.631) (-0.855,0.600) (-0.800,0.563) (-0.750,0.528) (-0.700,0.493) (-0.650,0.459) (-0.603,0.425) (-0.550,0.389) (-0.500,0.353) (-0.450,0.318) (-0.400,0.283) (-0.354,0.250) (-0.300,0.213) (-0.250,0.177) (-0.220,0.155)}
\EQ{(0.206,0.144) (0.248,0.175) (0.300,0.213) (0.350,0.248) (0.400,0.283) (0.450,0.318) (0.496,0.350) (0.550,0.389) (0.600,0.423) (0.650,0.459) (0.700,0.493) (0.746,0.525) (0.800,0.563) (0.850,0.597) (0.891,0.625) (0.950,0.665) (1.000,0.699) (1.050,0.733) (1.100,0.766) (1.150,0.799) (1.200,0.832) (1.250,0.865) (1.300,0.897) (1.345,0.925) (1.400,0.960) (1.450,0.991) (1.500,1.021) (1.548,1.050) (1.600,1.081) (1.650,1.110) (1.700,1.139) (1.750,1.167) (1.800,1.194) (1.850,1.220) (1.900,1.246) (1.950,1.272) (2.000,1.296) (2.050,1.320) (2.100,1.344) (2.150,1.366) (2.200,1.388) (2.250,1.409) (2.300,1.429) (2.350,1.448) (2.400,1.466) (2.450,1.484) (2.500,1.500) (2.550,1.516) (2.600,1.531) (2.650,1.544) (2.700,1.557) (2.750,1.569) (2.800,1.580) (2.850,1.589) (2.900,1.598) (2.950,1.605) (3.000,1.612) (3.050,1.617) (3.100,1.622) (3.150,1.625) (3.200,1.627) (3.250,1.627) (3.300,1.627) (3.350,1.625) (3.400,1.622) (3.450,1.618) (3.500,1.612) (3.550,1.605) (3.600,1.596) (3.650,1.586) (3.700,1.575) (3.750,1.562) (3.800,1.547) (3.850,1.531) (3.900,1.513) (3.950,1.493) (4.000,1.471) (4.050,1.447) (4.100,1.421) (4.150,1.393) (4.200,1.362) (4.250,1.328) (4.300,1.292) (4.350,1.253) (4.400,1.210) (4.438,1.175) (4.475,1.139) (4.512,1.100) (4.550,1.058) (4.597,1.000) (4.635,0.950) (4.669,0.900) (4.700,0.851) (4.730,0.800) (4.757,0.750) (4.782,0.700) (4.804,0.650) (4.825,0.600) (4.844,0.550) (4.861,0.500) (4.876,0.450) (4.889,0.400) (4.901,0.350) (4.911,0.300) (4.920,0.250) (4.926,0.200) (4.932,0.150) (4.936,0.100) (4.938,0.050) (4.939,0.000) (4.938,-0.050) (4.936,-0.100) (4.932,-0.150) (4.926,-0.200) (4.920,-0.250) (4.911,-0.300) (4.901,-0.350) (4.889,-0.400) (4.876,-0.450) (4.861,-0.500) (4.844,-0.550) (4.825,-0.600) (4.804,-0.650) (4.782,-0.700) (4.757,-0.750) (4.730,-0.800) (4.701,-0.850) (4.669,-0.900) (4.635,-0.950) (4.600,-0.996) (4.556,-1.050) (4.512,-1.100) (4.475,-1.139) (4.438,-1.175) (4.400,-1.210) (4.354,-1.250) (4.300,-1.292) (4.255,-1.325) (4.200,-1.362) (4.150,-1.393) (4.100,-1.421) (4.050,-1.447) (4.000,-1.471) (3.950,-1.493) (3.900,-1.513) (3.850,-1.531) (3.800,-1.547) (3.750,-1.562) (3.700,-1.575) (3.650,-1.586) (3.600,-1.596) (3.550,-1.605) (3.500,-1.612) (3.450,-1.618) (3.400,-1.622) (3.350,-1.625) (3.300,-1.627) (3.250,-1.627) (3.200,-1.627) (3.150,-1.625) (3.100,-1.622) (3.050,-1.617) (3.000,-1.612) (2.950,-1.605) (2.900,-1.598) (2.850,-1.589) (2.800,-1.580) (2.750,-1.569) (2.700,-1.557) (2.650,-1.544) (2.600,-1.531) (2.550,-1.516) (2.500,-1.500) (2.450,-1.484) (2.400,-1.466) (2.350,-1.448) (2.300,-1.429) (2.250,-1.409) (2.200,-1.388) (2.150,-1.366) (2.100,-1.344) (2.050,-1.320) (2.000,-1.296) (1.950,-1.272) (1.900,-1.246) (1.850,-1.220) (1.800,-1.194) (1.750,-1.167) (1.700,-1.139) (1.650,-1.110) (1.600,-1.081) (1.550,-1.051) (1.506,-1.025) (1.450,-0.991) (1.400,-0.960) (1.350,-0.929) (1.305,-0.900) (1.250,-0.865) (1.200,-0.832) (1.151,-0.800) (1.100,-0.766) (1.050,-0.733) (1.001,-0.700) (0.950,-0.665) (0.900,-0.631) (0.855,-0.600) (0.800,-0.563) (0.750,-0.528) (0.700,-0.493) (0.650,-0.459) (0.603,-0.425) (0.550,-0.389) (0.500,-0.353) (0.450,-0.318) (0.400,-0.283) (0.354,-0.250) (0.300,-0.213) (0.250,-0.177) (0.206,-0.144)}
\EQ{(-0.220,-0.155) (-0.275,-0.195) (-0.319,-0.225) (-0.375,-0.266) (-0.425,-0.300) (-0.475,-0.336) (-0.525,-0.371) (-0.567,-0.400) (-0.625,-0.441) (-0.674,-0.475) (-0.725,-0.511) (-0.775,-0.546) (-0.818,-0.575) (-0.875,-0.615) (-0.925,-0.649) (-0.975,-0.683) (-1.025,-0.717) (-1.074,-0.750) (-1.125,-0.785) (-1.175,-0.818) (-1.224,-0.850) (-1.275,-0.884) (-1.325,-0.916) (-1.375,-0.948) (-1.417,-0.975) (-1.475,-1.011) (-1.525,-1.042) (-1.575,-1.073) (-1.620,-1.100) (-1.675,-1.133) (-1.725,-1.162) (-1.775,-1.190) (-1.825,-1.218) (-1.875,-1.246) (-1.925,-1.273) (-1.975,-1.299) (-2.025,-1.325) (-2.075,-1.350) (-2.125,-1.374) (-2.175,-1.398) (-2.225,-1.421) (-2.275,-1.443) (-2.325,-1.465) (-2.375,-1.486) (-2.425,-1.505) (-2.475,-1.525) (-2.525,-1.543) (-2.575,-1.560) (-2.625,-1.577) (-2.675,-1.593) (-2.725,-1.608) (-2.775,-1.622) (-2.825,-1.635) (-2.875,-1.647) (-2.925,-1.658) (-2.975,-1.668) (-3.025,-1.677) (-3.075,-1.686) (-3.125,-1.693) (-3.175,-1.699) (-3.225,-1.704) (-3.275,-1.708) (-3.325,-1.711) (-3.375,-1.713) (-3.425,-1.713) (-3.475,-1.713) (-3.525,-1.711) (-3.575,-1.708) (-3.625,-1.704) (-3.675,-1.698) (-3.725,-1.691) (-3.775,-1.683) (-3.825,-1.674) (-3.875,-1.663) (-3.925,-1.650) (-3.975,-1.636) (-4.025,-1.621) (-4.075,-1.604) (-4.125,-1.585) (-4.175,-1.564) (-4.225,-1.542) (-4.275,-1.518) (-4.325,-1.491) (-4.375,-1.462) (-4.425,-1.431) (-4.472,-1.400) (-4.525,-1.362) (-4.572,-1.325) (-4.625,-1.280) (-4.675,-1.234) (-4.725,-1.185) (-4.775,-1.130) (-4.822,-1.075) (-4.861,-1.025) (-4.897,-0.975) (-4.925,-0.932) (-4.961,-0.875) (-4.989,-0.825) (-5.015,-0.775) (-5.040,-0.725) (-5.062,-0.675) (-5.082,-0.625) (-5.100,-0.577) (-5.117,-0.525) (-5.133,-0.475) (-5.146,-0.425) (-5.158,-0.375) (-5.168,-0.325) (-5.177,-0.275) (-5.184,-0.225) (-5.190,-0.175) (-5.195,-0.125) (-5.197,-0.075) (-5.199,-0.025) (-5.198,0.050) (-5.196,0.100) (-5.193,0.150) (-5.187,0.200) (-5.181,0.250) (-5.173,0.300) (-5.163,0.350) (-5.152,0.400) (-5.140,0.450) (-5.125,0.500) (-5.109,0.550) (-5.092,0.600) (-5.072,0.650) (-5.051,0.700) (-5.028,0.750) (-5.002,0.800) (-4.975,0.850) (-4.946,0.900) (-4.914,0.950) (-4.879,1.000) (-4.842,1.050) (-4.801,1.100) (-4.758,1.150) (-4.710,1.200) (-4.659,1.250) (-4.602,1.300) (-4.550,1.343) (-4.507,1.375) (-4.450,1.415) (-4.400,1.447) (-4.353,1.475) (-4.309,1.500) (-4.260,1.525) (-4.208,1.550) (-4.150,1.575) (-4.100,1.595) (-4.050,1.613) (-4.000,1.629) (-3.950,1.644) (-3.900,1.657) (-3.850,1.668) (-3.800,1.679) (-3.750,1.687) (-3.700,1.695) (-3.650,1.701) (-3.600,1.706) (-3.550,1.710) (-3.500,1.712) (-3.450,1.713) (-3.400,1.713) (-3.350,1.712) (-3.300,1.710) (-3.250,1.706) (-3.200,1.702) (-3.150,1.696) (-3.100,1.689) (-3.050,1.682) (-3.000,1.673) (-2.950,1.663) (-2.900,1.653) (-2.850,1.641) (-2.800,1.628) (-2.750,1.615) (-2.700,1.600) (-2.650,1.585) (-2.600,1.569) (-2.550,1.552) (-2.500,1.534) (-2.450,1.515) (-2.400,1.496) (-2.350,1.475) (-2.300,1.454) (-2.250,1.432) (-2.200,1.410) (-2.150,1.386) (-2.100,1.362) (-2.050,1.338) (-2.000,1.312) (-1.950,1.286) (-1.900,1.259) (-1.850,1.232) (-1.800,1.204) (-1.750,1.176) (-1.705,1.150) (-1.650,1.118) (-1.600,1.088) (-1.550,1.058) (-1.500,1.027) (-1.457,1.000) (-1.400,0.964) (-1.350,0.932) (-1.300,0.900) (-1.250,0.867) (-1.200,0.835) (-1.150,0.801) (-1.100,0.768) (-1.050,0.734) (-1.000,0.700) (-0.950,0.666) (-0.900,0.632) (-0.854,0.600) (-0.800,0.563) (-0.750,0.528) (-0.700,0.494) (-0.650,0.459) (-0.603,0.425) (-0.550,0.389) (-0.500,0.353) (-0.450,0.318) (-0.400,0.283) (-0.354,0.250) (-0.300,0.213) (-0.250,0.177) (-0.220,0.155)}
\EQ{(0.206,0.144) (0.248,0.175) (0.300,0.213) (0.350,0.248) (0.400,0.283) (0.450,0.318) (0.496,0.350) (0.550,0.389) (0.600,0.423) (0.650,0.459) (0.700,0.494) (0.746,0.525) (0.800,0.563) (0.850,0.598) (0.900,0.632) (0.950,0.666) (1.000,0.700) (1.050,0.734) (1.100,0.768) (1.148,0.800) (1.200,0.835) (1.250,0.867) (1.300,0.900) (1.350,0.932) (1.400,0.964) (1.450,0.996) (1.497,1.025) (1.550,1.058) (1.600,1.088) (1.650,1.118) (1.700,1.147) (1.748,1.175) (1.792,1.200) (1.837,1.225) (1.883,1.250) (1.929,1.275) (1.975,1.299) (2.025,1.325) (2.075,1.350) (2.125,1.374) (2.175,1.398) (2.225,1.421) (2.275,1.443) (2.325,1.465) (2.375,1.486) (2.425,1.505) (2.475,1.525) (2.525,1.543) (2.575,1.560) (2.625,1.577) (2.675,1.593) (2.725,1.608) (2.775,1.622) (2.825,1.635) (2.875,1.647) (2.925,1.658) (2.975,1.668) (3.025,1.677) (3.075,1.686) (3.125,1.693) (3.175,1.699) (3.225,1.704) (3.275,1.708) (3.325,1.711) (3.375,1.713) (3.425,1.713) (3.475,1.713) (3.525,1.711) (3.575,1.708) (3.625,1.704) (3.675,1.698) (3.725,1.691) (3.775,1.683) (3.825,1.674) (3.875,1.663) (3.925,1.650) (3.975,1.636) (4.025,1.621) (4.075,1.604) (4.125,1.585) (4.175,1.564) (4.225,1.542) (4.275,1.518) (4.325,1.491) (4.375,1.462) (4.425,1.431) (4.472,1.400) (4.525,1.362) (4.572,1.325) (4.625,1.280) (4.675,1.234) (4.725,1.185) (4.775,1.130) (4.822,1.075) (4.861,1.025) (4.897,0.975) (4.925,0.932) (4.961,0.875) (4.989,0.825) (5.015,0.775) (5.040,0.725) (5.062,0.675) (5.082,0.625) (5.100,0.577) (5.117,0.525) (5.133,0.475) (5.146,0.425) (5.158,0.375) (5.168,0.325) (5.177,0.275) (5.184,0.225) (5.190,0.175) (5.195,0.125) (5.197,0.075) (5.199,0.025) (5.198,-0.050) (5.196,-0.100) (5.193,-0.150) (5.187,-0.200) (5.181,-0.250) (5.173,-0.300) (5.163,-0.350) (5.152,-0.400) (5.140,-0.450) (5.125,-0.500) (5.109,-0.550) (5.092,-0.600) (5.072,-0.650) (5.051,-0.700) (5.028,-0.750) (5.002,-0.800) (4.975,-0.850) (4.946,-0.900) (4.914,-0.950) (4.879,-1.000) (4.842,-1.050) (4.801,-1.100) (4.758,-1.150) (4.710,-1.200) (4.659,-1.250) (4.602,-1.300) (4.550,-1.343) (4.507,-1.375) (4.450,-1.415) (4.400,-1.447) (4.353,-1.475) (4.309,-1.500) (4.260,-1.525) (4.208,-1.550) (4.150,-1.575) (4.100,-1.595) (4.050,-1.613) (4.000,-1.629) (3.950,-1.644) (3.900,-1.657) (3.850,-1.668) (3.800,-1.679) (3.750,-1.687) (3.700,-1.695) (3.650,-1.701) (3.600,-1.706) (3.550,-1.710) (3.500,-1.712) (3.450,-1.713) (3.400,-1.713) (3.350,-1.712) (3.300,-1.710) (3.250,-1.706) (3.200,-1.702) (3.150,-1.696) (3.100,-1.689) (3.050,-1.682) (3.000,-1.673) (2.950,-1.663) (2.900,-1.653) (2.850,-1.641) (2.800,-1.628) (2.750,-1.615) (2.700,-1.600) (2.650,-1.585) (2.600,-1.569) (2.550,-1.552) (2.500,-1.534) (2.450,-1.515) (2.400,-1.496) (2.350,-1.475) (2.300,-1.454) (2.250,-1.432) (2.200,-1.410) (2.150,-1.386) (2.100,-1.362) (2.050,-1.338) (2.000,-1.312) (1.950,-1.286) (1.900,-1.259) (1.850,-1.232) (1.800,-1.204) (1.750,-1.176) (1.705,-1.150) (1.650,-1.118) (1.600,-1.088) (1.550,-1.058) (1.500,-1.027) (1.457,-1.000) (1.400,-0.964) (1.350,-0.932) (1.300,-0.900) (1.250,-0.867) (1.200,-0.835) (1.150,-0.801) (1.100,-0.768) (1.050,-0.734) (1.000,-0.700) (0.950,-0.666) (0.900,-0.632) (0.854,-0.600) (0.800,-0.563) (0.750,-0.528) (0.700,-0.494) (0.650,-0.459) (0.603,-0.425) (0.550,-0.389) (0.500,-0.353) (0.450,-0.318) (0.400,-0.283) (0.354,-0.250) (0.300,-0.213) (0.250,-0.177) (0.206,-0.144)}
\EQ{(-0.220,-0.155) (-0.275,-0.195) (-0.319,-0.225) (-0.375,-0.266) (-0.425,-0.300) (-0.475,-0.336) (-0.525,-0.371) (-0.567,-0.400) (-0.625,-0.441) (-0.674,-0.475) (-0.725,-0.511) (-0.775,-0.546) (-0.817,-0.575) (-0.875,-0.616) (-0.925,-0.650) (-0.975,-0.684) (-1.025,-0.719) (-1.072,-0.750) (-1.125,-0.787) (-1.175,-0.820) (-1.220,-0.850) (-1.275,-0.887) (-1.325,-0.920) (-1.372,-0.950) (-1.425,-0.985) (-1.475,-1.017) (-1.525,-1.048) (-1.568,-1.075) (-1.625,-1.110) (-1.675,-1.141) (-1.725,-1.171) (-1.774,-1.200) (-1.825,-1.230) (-1.875,-1.259) (-1.925,-1.287) (-1.975,-1.314) (-2.025,-1.342) (-2.075,-1.368) (-2.125,-1.394) (-2.175,-1.420) (-2.225,-1.444) (-2.275,-1.468) (-2.325,-1.492) (-2.375,-1.515) (-2.425,-1.537) (-2.475,-1.558) (-2.525,-1.579) (-2.575,-1.599) (-2.625,-1.618) (-2.675,-1.637) (-2.725,-1.654) (-2.775,-1.671) (-2.825,-1.687) (-2.875,-1.702) (-2.925,-1.717) (-2.975,-1.730) (-3.025,-1.743) (-3.075,-1.755) (-3.125,-1.765) (-3.175,-1.775) (-3.225,-1.784) (-3.275,-1.792) (-3.325,-1.799) (-3.375,-1.805) (-3.425,-1.811) (-3.475,-1.815) (-3.525,-1.818) (-3.575,-1.820) (-3.625,-1.820) (-3.675,-1.820) (-3.725,-1.819) (-3.775,-1.816) (-3.825,-1.813) (-3.875,-1.808) (-3.925,-1.802) (-3.975,-1.795) (-4.025,-1.786) (-4.075,-1.776) (-4.125,-1.765) (-4.175,-1.753) (-4.225,-1.739) (-4.275,-1.723) (-4.325,-1.706) (-4.375,-1.688) (-4.425,-1.667) (-4.475,-1.645) (-4.525,-1.622) (-4.575,-1.596) (-4.625,-1.568) (-4.675,-1.538) (-4.725,-1.506) (-4.770,-1.475) (-4.825,-1.435) (-4.868,-1.400) (-4.925,-1.352) (-4.975,-1.305) (-5.025,-1.255) (-5.075,-1.200) (-5.118,-1.150) (-5.150,-1.109) (-5.194,-1.050) (-5.225,-1.004) (-5.259,-0.950) (-5.289,-0.900) (-5.316,-0.850) (-5.341,-0.800) (-5.365,-0.750) (-5.386,-0.700) (-5.406,-0.650) (-5.424,-0.600) (-5.441,-0.550) (-5.456,-0.500) (-5.469,-0.450) (-5.481,-0.400) (-5.491,-0.350) (-5.500,-0.300) (-5.508,-0.250) (-5.514,-0.200) (-5.519,-0.150) (-5.522,-0.100) (-5.524,-0.050) (-5.525,0.000) (-5.524,0.050) (-5.522,0.100) (-5.519,0.150) (-5.514,0.200) (-5.508,0.250) (-5.500,0.300) (-5.491,0.350) (-5.481,0.400) (-5.469,0.450) (-5.456,0.500) (-5.441,0.550) (-5.424,0.600) (-5.406,0.650) (-5.386,0.700) (-5.365,0.750) (-5.341,0.800) (-5.316,0.850) (-5.289,0.900) (-5.259,0.950) (-5.228,1.000) (-5.194,1.050) (-5.157,1.100) (-5.125,1.141) (-5.075,1.200) (-5.030,1.250) (-4.980,1.300) (-4.927,1.350) (-4.875,1.395) (-4.825,1.435) (-4.775,1.472) (-4.725,1.506) (-4.675,1.538) (-4.625,1.568) (-4.575,1.596) (-4.525,1.622) (-4.475,1.645) (-4.425,1.667) (-4.375,1.688) (-4.325,1.706) (-4.275,1.723) (-4.225,1.739) (-4.175,1.753) (-4.125,1.765) (-4.075,1.776) (-4.025,1.786) (-3.975,1.795) (-3.925,1.802) (-3.875,1.808) (-3.825,1.813) (-3.775,1.816) (-3.725,1.819) (-3.675,1.820) (-3.625,1.820) (-3.575,1.820) (-3.525,1.818) (-3.475,1.815) (-3.425,1.811) (-3.375,1.805) (-3.325,1.799) (-3.275,1.792) (-3.225,1.784) (-3.175,1.775) (-3.125,1.765) (-3.075,1.755) (-3.025,1.743) (-2.975,1.730) (-2.925,1.717) (-2.875,1.702) (-2.825,1.687) (-2.775,1.671) (-2.725,1.654) (-2.675,1.637) (-2.625,1.618) (-2.578,1.600) (-2.525,1.579) (-2.475,1.558) (-2.425,1.537) (-2.375,1.515) (-2.325,1.492) (-2.275,1.468) (-2.225,1.444) (-2.175,1.420) (-2.125,1.394) (-2.075,1.368) (-2.025,1.342) (-1.975,1.314) (-1.925,1.287) (-1.875,1.259) (-1.825,1.230) (-1.775,1.201) (-1.725,1.171) (-1.675,1.141) (-1.625,1.110) (-1.575,1.079) (-1.528,1.050) (-1.475,1.017) (-1.425,0.985) (-1.375,0.952) (-1.325,0.920) (-1.275,0.887) (-1.225,0.854) (-1.183,0.825) (-1.125,0.787) (-1.075,0.753) (-1.025,0.719) (-0.975,0.684) (-0.925,0.650) (-0.875,0.616) (-0.825,0.581) (-0.781,0.550) (-0.725,0.511) (-0.675,0.476) (-0.625,0.441) (-0.575,0.406) (-0.532,0.375) (-0.475,0.336) (-0.425,0.300) (-0.375,0.266) (-0.325,0.230) (-0.284,0.200) (-0.225,0.160) (-0.220,0.155)}
\EQ{(0.206,0.144) (0.248,0.175) (0.300,0.213) (0.350,0.248) (0.400,0.283) (0.450,0.318) (0.496,0.350) (0.550,0.389) (0.600,0.424) (0.650,0.459) (0.700,0.494) (0.745,0.525) (0.800,0.564) (0.850,0.598) (0.900,0.633) (0.950,0.667) (0.998,0.700) (1.050,0.736) (1.100,0.770) (1.145,0.800) (1.200,0.837) (1.250,0.870) (1.296,0.900) (1.350,0.936) (1.400,0.968) (1.449,1.000) (1.500,1.032) (1.550,1.064) (1.600,1.095) (1.649,1.125) (1.700,1.156) (1.750,1.186) (1.800,1.215) (1.850,1.244) (1.900,1.273) (1.949,1.300) (1.995,1.325) (2.041,1.350) (2.088,1.375) (2.137,1.400) (2.186,1.425) (2.237,1.450) (2.289,1.475) (2.343,1.500) (2.398,1.525) (2.450,1.548) (2.500,1.569) (2.550,1.589) (2.600,1.609) (2.650,1.628) (2.700,1.646) (2.750,1.663) (2.800,1.679) (2.850,1.695) (2.900,1.710) (2.950,1.724) (3.000,1.737) (3.050,1.749) (3.100,1.760) (3.150,1.771) (3.200,1.780) (3.250,1.788) (3.300,1.796) (3.350,1.803) (3.400,1.808) (3.450,1.813) (3.500,1.816) (3.550,1.819) (3.600,1.820) (3.650,1.820) (3.700,1.820) (3.750,1.818) (3.800,1.815) (3.850,1.811) (3.900,1.805) (3.950,1.799) (4.000,1.791) (4.050,1.782) (4.100,1.771) (4.150,1.759) (4.200,1.746) (4.250,1.731) (4.300,1.715) (4.350,1.697) (4.400,1.678) (4.450,1.657) (4.500,1.634) (4.550,1.609) (4.600,1.582) (4.650,1.553) (4.696,1.525) (4.750,1.489) (4.800,1.453) (4.850,1.415) (4.898,1.375) (4.950,1.329) (5.000,1.280) (5.050,1.228) (5.097,1.175) (5.138,1.125) (5.175,1.076) (5.211,1.025) (5.244,0.975) (5.274,0.925) (5.300,0.880) (5.325,0.832) (5.350,0.782) (5.375,0.726) (5.396,0.675) (5.415,0.625) (5.433,0.575) (5.448,0.525) (5.462,0.475) (5.475,0.425) (5.486,0.375) (5.496,0.325) (5.504,0.275) (5.511,0.225) (5.516,0.175) (5.521,0.125) (5.523,0.075) (5.525,0.025) (5.524,-0.050) (5.522,-0.100) (5.519,-0.150) (5.514,-0.200) (5.508,-0.250) (5.500,-0.300) (5.491,-0.350) (5.481,-0.400) (5.469,-0.450) (5.456,-0.500) (5.441,-0.550) (5.424,-0.600) (5.406,-0.650) (5.386,-0.700) (5.365,-0.750) (5.341,-0.800) (5.316,-0.850) (5.289,-0.900) (5.259,-0.950) (5.228,-1.000) (5.194,-1.050) (5.157,-1.100) (5.125,-1.141) (5.075,-1.200) (5.030,-1.250) (4.980,-1.300) (4.927,-1.350) (4.875,-1.395) (4.825,-1.435) (4.775,-1.472) (4.725,-1.506) (4.675,-1.538) (4.625,-1.568) (4.575,-1.596) (4.525,-1.622) (4.475,-1.645) (4.425,-1.667) (4.375,-1.688) (4.325,-1.706) (4.275,-1.723) (4.225,-1.739) (4.175,-1.753) (4.125,-1.765) (4.075,-1.776) (4.025,-1.786) (3.975,-1.795) (3.925,-1.802) (3.875,-1.808) (3.825,-1.813) (3.775,-1.816) (3.725,-1.819) (3.675,-1.820) (3.625,-1.820) (3.575,-1.820) (3.525,-1.818) (3.475,-1.815) (3.425,-1.811) (3.375,-1.805) (3.325,-1.799) (3.275,-1.792) (3.225,-1.784) (3.175,-1.775) (3.125,-1.765) (3.075,-1.755) (3.025,-1.743) (2.975,-1.730) (2.925,-1.717) (2.875,-1.702) (2.825,-1.687) (2.775,-1.671) (2.725,-1.654) (2.675,-1.637) (2.625,-1.618) (2.578,-1.600) (2.525,-1.579) (2.475,-1.558) (2.425,-1.537) (2.375,-1.515) (2.325,-1.492) (2.275,-1.468) (2.225,-1.444) (2.175,-1.420) (2.125,-1.394) (2.075,-1.368) (2.025,-1.342) (1.975,-1.314) (1.925,-1.287) (1.875,-1.259) (1.825,-1.230) (1.775,-1.201) (1.725,-1.171) (1.675,-1.141) (1.625,-1.110) (1.575,-1.079) (1.528,-1.050) (1.475,-1.017) (1.425,-0.985) (1.375,-0.952) (1.325,-0.920) (1.275,-0.887) (1.225,-0.854) (1.183,-0.825) (1.125,-0.787) (1.075,-0.753) (1.025,-0.719) (0.975,-0.684) (0.925,-0.650) (0.875,-0.616) (0.825,-0.581) (0.781,-0.550) (0.725,-0.511) (0.675,-0.476) (0.625,-0.441) (0.575,-0.406) (0.532,-0.375) (0.475,-0.336) (0.425,-0.300) (0.375,-0.266) (0.325,-0.230) (0.284,-0.200) (0.225,-0.160) (0.206,-0.144)}
\EQ{(-0.220,-0.155) (-0.275,-0.195) (-0.319,-0.225) (-0.375,-0.266) (-0.425,-0.300) (-0.475,-0.336) (-0.525,-0.371) (-0.567,-0.400) (-0.625,-0.442) (-0.673,-0.475) (-0.725,-0.512) (-0.775,-0.547) (-0.817,-0.575) (-0.875,-0.616) (-0.924,-0.650) (-0.975,-0.686) (-1.025,-0.720) (-1.069,-0.750) (-1.125,-0.788) (-1.175,-0.822) (-1.225,-0.856) (-1.275,-0.890) (-1.325,-0.923) (-1.375,-0.956) (-1.425,-0.989) (-1.475,-1.022) (-1.519,-1.050) (-1.575,-1.086) (-1.625,-1.118) (-1.675,-1.149) (-1.725,-1.180) (-1.775,-1.211) (-1.825,-1.242) (-1.875,-1.271) (-1.924,-1.300) (-1.975,-1.330) (-2.025,-1.359) (-2.075,-1.387) (-2.125,-1.415) (-2.175,-1.442) (-2.225,-1.468) (-2.275,-1.495) (-2.325,-1.520) (-2.375,-1.545) (-2.425,-1.570) (-2.475,-1.593) (-2.525,-1.617) (-2.575,-1.639) (-2.625,-1.661) (-2.675,-1.682) (-2.725,-1.703) (-2.775,-1.723) (-2.825,-1.742) (-2.875,-1.761) (-2.925,-1.778) (-2.975,-1.796) (-3.025,-1.812) (-3.075,-1.827) (-3.125,-1.842) (-3.175,-1.856) (-3.225,-1.869) (-3.275,-1.881) (-3.325,-1.893) (-3.375,-1.904) (-3.425,-1.913) (-3.475,-1.922) (-3.525,-1.930) (-3.575,-1.937) (-3.625,-1.944) (-3.675,-1.949) (-3.725,-1.953) (-3.775,-1.957) (-3.825,-1.959) (-3.875,-1.961) (-3.925,-1.961) (-3.975,-1.960) (-4.025,-1.959) (-4.075,-1.956) (-4.125,-1.952) (-4.175,-1.948) (-4.225,-1.942) (-4.275,-1.934) (-4.325,-1.926) (-4.375,-1.917) (-4.425,-1.906) (-4.475,-1.894) (-4.525,-1.880) (-4.575,-1.866) (-4.624,-1.850) (-4.675,-1.832) (-4.725,-1.813) (-4.775,-1.793) (-4.825,-1.770) (-4.875,-1.746) (-4.925,-1.721) (-4.975,-1.693) (-5.025,-1.664) (-5.075,-1.632) (-5.123,-1.600) (-5.175,-1.562) (-5.223,-1.525) (-5.275,-1.482) (-5.325,-1.438) (-5.365,-1.400) (-5.415,-1.350) (-5.450,-1.312) (-5.500,-1.254) (-5.544,-1.200) (-5.575,-1.158) (-5.616,-1.100) (-5.649,-1.050) (-5.675,-1.007) (-5.708,-0.950) (-5.735,-0.900) (-5.759,-0.850) (-5.782,-0.800) (-5.804,-0.750) (-5.823,-0.700) (-5.842,-0.650) (-5.858,-0.600) (-5.874,-0.550) (-5.887,-0.500) (-5.900,-0.450) (-5.911,-0.400) (-5.920,-0.350) (-5.929,-0.300) (-5.936,-0.250) (-5.941,-0.200) (-5.946,-0.150) (-5.949,-0.100) (-5.951,-0.050) (-5.951,0.000) (-5.951,0.050) (-5.949,0.100) (-5.946,0.150) (-5.941,0.200) (-5.936,0.250) (-5.929,0.300) (-5.920,0.350) (-5.911,0.400) (-5.900,0.450) (-5.887,0.500) (-5.874,0.550) (-5.858,0.600) (-5.842,0.650) (-5.823,0.700) (-5.804,0.750) (-5.782,0.800) (-5.759,0.850) (-5.735,0.900) (-5.708,0.950) (-5.679,1.000) (-5.650,1.048) (-5.616,1.100) (-5.581,1.150) (-5.550,1.192) (-5.504,1.250) (-5.461,1.300) (-5.425,1.339) (-5.375,1.390) (-5.325,1.438) (-5.284,1.475) (-5.225,1.524) (-5.175,1.562) (-5.125,1.599) (-5.075,1.632) (-5.025,1.664) (-4.975,1.693) (-4.925,1.721) (-4.875,1.746) (-4.825,1.770) (-4.775,1.793) (-4.725,1.813) (-4.675,1.832) (-4.625,1.850) (-4.575,1.866) (-4.525,1.880) (-4.475,1.894) (-4.425,1.906) (-4.375,1.917) (-4.325,1.926) (-4.275,1.934) (-4.225,1.942) (-4.175,1.948) (-4.125,1.952) (-4.075,1.956) (-4.025,1.959) (-3.975,1.960) (-3.925,1.961) (-3.875,1.961) (-3.825,1.959) (-3.775,1.957) (-3.725,1.953) (-3.675,1.949) (-3.625,1.944) (-3.575,1.937) (-3.525,1.930) (-3.475,1.922) (-3.425,1.913) (-3.375,1.904) (-3.325,1.893) (-3.275,1.881) (-3.225,1.869) (-3.175,1.856) (-3.125,1.842) (-3.075,1.827) (-3.025,1.812) (-2.975,1.796) (-2.925,1.778) (-2.875,1.761) (-2.825,1.742) (-2.775,1.723) (-2.725,1.703) (-2.675,1.682) (-2.625,1.661) (-2.575,1.639) (-2.525,1.617) (-2.475,1.593) (-2.425,1.570) (-2.375,1.545) (-2.325,1.520) (-2.275,1.495) (-2.225,1.468) (-2.175,1.442) (-2.125,1.415) (-2.075,1.387) (-2.025,1.359) (-1.975,1.330) (-1.925,1.301) (-1.881,1.275) (-1.825,1.242) (-1.775,1.211) (-1.725,1.180) (-1.676,1.150) (-1.625,1.118) (-1.575,1.086) (-1.525,1.054) (-1.480,1.025) (-1.425,0.989) (-1.375,0.956) (-1.328,0.925) (-1.275,0.890) (-1.225,0.856) (-1.179,0.825) (-1.125,0.788) (-1.075,0.754) (-1.033,0.725) (-0.975,0.686) (-0.925,0.651) (-0.875,0.616) (-0.825,0.581) (-0.781,0.550) (-0.725,0.512) (-0.675,0.476) (-0.625,0.442) (-0.575,0.406) (-0.531,0.375) (-0.475,0.336) (-0.425,0.300) (-0.375,0.266) (-0.325,0.230) (-0.284,0.200) (-0.225,0.160) (-0.220,0.155)}
\EQ{(0.206,0.144) (0.248,0.175) (0.300,0.213) (0.350,0.248) (0.400,0.283) (0.450,0.318) (0.496,0.350) (0.550,0.389) (0.600,0.424) (0.650,0.459) (0.700,0.494) (0.745,0.525) (0.800,0.564) (0.850,0.599) (0.900,0.634) (0.950,0.668) (0.997,0.700) (1.050,0.737) (1.100,0.771) (1.143,0.800) (1.200,0.839) (1.250,0.873) (1.300,0.906) (1.350,0.940) (1.400,0.973) (1.442,1.000) (1.500,1.038) (1.550,1.070) (1.597,1.100) (1.650,1.134) (1.700,1.165) (1.750,1.196) (1.798,1.225) (1.850,1.257) (1.900,1.286) (1.950,1.316) (2.000,1.344) (2.050,1.373) (2.099,1.400) (2.144,1.425) (2.191,1.450) (2.238,1.475) (2.286,1.500) (2.335,1.525) (2.385,1.550) (2.437,1.575) (2.489,1.600) (2.544,1.625) (2.600,1.650) (2.650,1.672) (2.700,1.693) (2.750,1.713) (2.800,1.733) (2.850,1.751) (2.900,1.770) (2.950,1.787) (3.000,1.804) (3.050,1.820) (3.100,1.835) (3.150,1.849) (3.200,1.863) (3.249,1.875) (3.300,1.887) (3.350,1.898) (3.400,1.909) (3.450,1.918) (3.500,1.926) (3.550,1.934) (3.600,1.941) (3.650,1.946) (3.700,1.951) (3.750,1.955) (3.800,1.958) (3.850,1.960) (3.900,1.961) (3.950,1.961) (4.000,1.960) (4.050,1.958) (4.100,1.954) (4.150,1.950) (4.200,1.945) (4.250,1.938) (4.300,1.930) (4.350,1.922) (4.400,1.911) (4.450,1.900) (4.500,1.887) (4.550,1.873) (4.600,1.858) (4.650,1.841) (4.700,1.823) (4.750,1.803) (4.800,1.782) (4.850,1.759) (4.900,1.734) (4.950,1.707) (5.000,1.679) (5.047,1.650) (5.100,1.616) (5.150,1.581) (5.191,1.550) (5.250,1.503) (5.300,1.461) (5.339,1.425) (5.390,1.375) (5.438,1.325) (5.475,1.284) (5.524,1.225) (5.563,1.175) (5.599,1.125) (5.633,1.075) (5.664,1.025) (5.694,0.975) (5.721,0.925) (5.747,0.875) (5.771,0.825) (5.793,0.775) (5.814,0.725) (5.833,0.675) (5.850,0.625) (5.866,0.575) (5.881,0.525) (5.894,0.475) (5.905,0.425) (5.916,0.375) (5.925,0.325) (5.932,0.275) (5.939,0.225) (5.944,0.175) (5.947,0.125) (5.950,0.075) (5.951,0.025) (5.951,-0.050) (5.949,-0.100) (5.946,-0.150) (5.941,-0.200) (5.936,-0.250) (5.929,-0.300) (5.920,-0.350) (5.911,-0.400) (5.900,-0.450) (5.887,-0.500) (5.874,-0.550) (5.858,-0.600) (5.842,-0.650) (5.823,-0.700) (5.804,-0.750) (5.782,-0.800) (5.759,-0.850) (5.735,-0.900) (5.708,-0.950) (5.679,-1.000) (5.650,-1.048) (5.616,-1.100) (5.581,-1.150) (5.550,-1.192) (5.504,-1.250) (5.461,-1.300) (5.425,-1.339) (5.375,-1.390) (5.325,-1.438) (5.284,-1.475) (5.225,-1.524) (5.175,-1.562) (5.125,-1.599) (5.075,-1.632) (5.025,-1.664) (4.975,-1.693) (4.925,-1.721) (4.875,-1.746) (4.825,-1.770) (4.775,-1.793) (4.725,-1.813) (4.675,-1.832) (4.625,-1.850) (4.575,-1.866) (4.525,-1.880) (4.475,-1.894) (4.425,-1.906) (4.375,-1.917) (4.325,-1.926) (4.275,-1.934) (4.225,-1.942) (4.175,-1.948) (4.125,-1.952) (4.075,-1.956) (4.025,-1.959) (3.975,-1.960) (3.925,-1.961) (3.875,-1.961) (3.825,-1.959) (3.775,-1.957) (3.725,-1.953) (3.675,-1.949) (3.625,-1.944) (3.575,-1.937) (3.525,-1.930) (3.475,-1.922) (3.425,-1.913) (3.375,-1.904) (3.325,-1.893) (3.275,-1.881) (3.225,-1.869) (3.175,-1.856) (3.125,-1.842) (3.075,-1.827) (3.025,-1.812) (2.975,-1.796) (2.925,-1.778) (2.875,-1.761) (2.825,-1.742) (2.775,-1.723) (2.725,-1.703) (2.675,-1.682) (2.625,-1.661) (2.575,-1.639) (2.525,-1.617) (2.475,-1.593) (2.425,-1.570) (2.375,-1.545) (2.325,-1.520) (2.275,-1.495) (2.225,-1.468) (2.175,-1.442) (2.125,-1.415) (2.075,-1.387) (2.025,-1.359) (1.975,-1.330) (1.925,-1.301) (1.881,-1.275) (1.825,-1.242) (1.775,-1.211) (1.725,-1.180) (1.676,-1.150) (1.625,-1.118) (1.575,-1.086) (1.525,-1.054) (1.480,-1.025) (1.425,-0.989) (1.375,-0.956) (1.328,-0.925) (1.275,-0.890) (1.225,-0.856) (1.179,-0.825) (1.125,-0.788) (1.075,-0.754) (1.033,-0.725) (0.975,-0.686) (0.925,-0.651) (0.875,-0.616) (0.825,-0.581) (0.781,-0.550) (0.725,-0.512) (0.675,-0.476) (0.625,-0.442) (0.575,-0.406) (0.531,-0.375) (0.475,-0.336) (0.425,-0.300) (0.375,-0.266) (0.325,-0.230) (0.284,-0.200) (0.225,-0.160) (0.206,-0.144)}
\EQ{(-0.220,-0.155) (-0.275,-0.195) (-0.319,-0.225) (-0.375,-0.266) (-0.425,-0.300) (-0.475,-0.336) (-0.525,-0.371) (-0.567,-0.400) (-0.625,-0.442) (-0.673,-0.475) (-0.725,-0.512) (-0.775,-0.547) (-0.825,-0.582) (-0.875,-0.617) (-0.923,-0.650) (-0.975,-0.687) (-1.025,-0.721) (-1.067,-0.750) (-1.125,-0.790) (-1.175,-0.824) (-1.225,-0.859) (-1.275,-0.893) (-1.323,-0.925) (-1.375,-0.960) (-1.425,-0.994) (-1.472,-1.025) (-1.525,-1.060) (-1.575,-1.093) (-1.624,-1.125) (-1.675,-1.158) (-1.725,-1.190) (-1.775,-1.222) (-1.820,-1.250) (-1.875,-1.285) (-1.925,-1.316) (-1.975,-1.346) (-2.023,-1.375) (-2.075,-1.406) (-2.125,-1.436) (-2.175,-1.465) (-2.225,-1.493) (-2.275,-1.521) (-2.325,-1.549) (-2.372,-1.575) (-2.419,-1.600) (-2.466,-1.625) (-2.514,-1.650) (-2.563,-1.675) (-2.613,-1.700) (-2.664,-1.725) (-2.716,-1.750) (-2.770,-1.775) (-2.825,-1.800) (-2.875,-1.822) (-2.925,-1.844) (-2.975,-1.865) (-3.025,-1.885) (-3.075,-1.904) (-3.125,-1.923) (-3.175,-1.942) (-3.225,-1.959) (-3.275,-1.976) (-3.325,-1.993) (-3.375,-2.008) (-3.425,-2.023) (-3.475,-2.037) (-3.525,-2.051) (-3.575,-2.064) (-3.625,-2.076) (-3.675,-2.087) (-3.725,-2.097) (-3.775,-2.107) (-3.825,-2.116) (-3.875,-2.124) (-3.925,-2.131) (-3.975,-2.138) (-4.025,-2.143) (-4.075,-2.148) (-4.125,-2.152) (-4.175,-2.155) (-4.225,-2.157) (-4.275,-2.158) (-4.325,-2.158) (-4.375,-2.158) (-4.425,-2.156) (-4.475,-2.154) (-4.525,-2.150) (-4.575,-2.146) (-4.625,-2.140) (-4.675,-2.134) (-4.725,-2.126) (-4.775,-2.117) (-4.825,-2.108) (-4.875,-2.097) (-4.925,-2.085) (-4.975,-2.071) (-5.025,-2.057) (-5.075,-2.041) (-5.125,-2.024) (-5.175,-2.006) (-5.225,-1.986) (-5.275,-1.965) (-5.325,-1.942) (-5.375,-1.918) (-5.425,-1.892) (-5.475,-1.864) (-5.525,-1.835) (-5.575,-1.804) (-5.618,-1.775) (-5.675,-1.735) (-5.722,-1.700) (-5.775,-1.658) (-5.814,-1.625) (-5.869,-1.575) (-5.921,-1.525) (-5.970,-1.475) (-6.016,-1.425) (-6.050,-1.385) (-6.098,-1.325) (-6.136,-1.275) (-6.172,-1.225) (-6.200,-1.183) (-6.236,-1.125) (-6.266,-1.075) (-6.294,-1.025) (-6.320,-0.975) (-6.344,-0.925) (-6.367,-0.875) (-6.388,-0.825) (-6.408,-0.775) (-6.426,-0.725) (-6.443,-0.675) (-6.459,-0.625) (-6.473,-0.575) (-6.486,-0.525) (-6.498,-0.475) (-6.509,-0.425) (-6.518,-0.375) (-6.526,-0.325) (-6.533,-0.275) (-6.539,-0.225) (-6.543,-0.175) (-6.547,-0.125) (-6.549,-0.075) (-6.550,-0.025) (-6.550,0.037) (-6.548,0.100) (-6.545,0.150) (-6.541,0.200) (-6.536,0.250) (-6.530,0.300) (-6.522,0.350) (-6.513,0.400) (-6.503,0.450) (-6.492,0.500) (-6.480,0.550) (-6.466,0.600) (-6.451,0.650) (-6.435,0.700) (-6.417,0.750) (-6.398,0.800) (-6.378,0.850) (-6.356,0.900) (-6.332,0.950) (-6.307,1.000) (-6.280,1.050) (-6.251,1.100) (-6.225,1.143) (-6.189,1.200) (-6.154,1.250) (-6.118,1.300) (-6.079,1.350) (-6.037,1.400) (-6.000,1.442) (-5.950,1.496) (-5.900,1.546) (-5.850,1.593) (-5.800,1.637) (-5.754,1.675) (-5.700,1.717) (-5.654,1.750) (-5.600,1.787) (-5.550,1.819) (-5.500,1.850) (-5.456,1.875) (-5.409,1.900) (-5.360,1.925) (-5.307,1.950) (-5.251,1.975) (-5.200,1.996) (-5.150,2.015) (-5.100,2.033) (-5.050,2.049) (-5.000,2.064) (-4.950,2.078) (-4.900,2.091) (-4.850,2.102) (-4.800,2.113) (-4.750,2.122) (-4.700,2.130) (-4.650,2.137) (-4.600,2.143) (-4.550,2.148) (-4.500,2.152) (-4.450,2.155) (-4.400,2.157) (-4.350,2.158) (-4.300,2.158) (-4.250,2.158) (-4.200,2.156) (-4.150,2.153) (-4.100,2.150) (-4.050,2.146) (-4.000,2.140) (-3.950,2.134) (-3.900,2.128) (-3.850,2.120) (-3.800,2.111) (-3.750,2.102) (-3.700,2.092) (-3.650,2.081) (-3.600,2.070) (-3.550,2.057) (-3.500,2.044) (-3.450,2.030) (-3.400,2.016) (-3.350,2.001) (-3.300,1.985) (-3.250,1.968) (-3.200,1.951) (-3.150,1.933) (-3.100,1.914) (-3.050,1.895) (-3.001,1.875) (-2.950,1.854) (-2.900,1.833) (-2.850,1.811) (-2.800,1.789) (-2.750,1.766) (-2.700,1.742) (-2.650,1.718) (-2.600,1.694) (-2.550,1.669) (-2.500,1.643) (-2.450,1.617) (-2.400,1.590) (-2.350,1.563) (-2.300,1.535) (-2.250,1.507) (-2.200,1.479) (-2.150,1.450) (-2.100,1.421) (-2.050,1.391) (-2.000,1.361) (-1.950,1.331) (-1.900,1.300) (-1.850,1.269) (-1.800,1.238) (-1.750,1.206) (-1.702,1.175) (-1.650,1.142) (-1.600,1.109) (-1.550,1.076) (-1.500,1.044) (-1.450,1.010) (-1.400,0.977) (-1.350,0.943) (-1.300,0.910) (-1.250,0.876) (-1.200,0.842) (-1.150,0.807) (-1.103,0.775) (-1.050,0.739) (-1.000,0.704) (-0.950,0.669) (-0.900,0.634) (-0.851,0.600) (-0.800,0.565) (-0.750,0.529) (-0.709,0.500) (-0.650,0.459) (-0.602,0.425) (-0.550,0.389) (-0.500,0.354) (-0.450,0.319) (-0.400,0.283) (-0.354,0.250) (-0.300,0.213) (-0.250,0.177) (-0.220,0.155)}
\EQ{(0.206,0.144) (0.248,0.175) (0.300,0.213) (0.350,0.248) (0.400,0.283) (0.450,0.319) (0.496,0.350) (0.550,0.389) (0.600,0.424) (0.650,0.459) (0.700,0.494) (0.744,0.525) (0.800,0.565) (0.850,0.599) (0.900,0.634) (0.950,0.669) (0.995,0.700) (1.050,0.739) (1.100,0.773) (1.150,0.807) (1.200,0.842) (1.249,0.875) (1.300,0.910) (1.350,0.943) (1.397,0.975) (1.450,1.010) (1.500,1.044) (1.548,1.075) (1.600,1.109) (1.650,1.142) (1.700,1.174) (1.750,1.206) (1.800,1.238) (1.850,1.269) (1.900,1.300) (1.950,1.331) (2.000,1.361) (2.050,1.391) (2.100,1.421) (2.150,1.450) (2.193,1.475) (2.237,1.500) (2.282,1.525) (2.327,1.550) (2.372,1.575) (2.419,1.600) (2.466,1.625) (2.514,1.650) (2.563,1.675) (2.613,1.700) (2.664,1.725) (2.716,1.750) (2.770,1.775) (2.825,1.800) (2.875,1.822) (2.925,1.844) (2.975,1.865) (3.025,1.885) (3.075,1.904) (3.125,1.923) (3.175,1.942) (3.225,1.959) (3.275,1.976) (3.325,1.993) (3.375,2.008) (3.425,2.023) (3.475,2.037) (3.525,2.051) (3.575,2.064) (3.625,2.076) (3.675,2.087) (3.725,2.097) (3.775,2.107) (3.825,2.116) (3.875,2.124) (3.925,2.131) (3.975,2.138) (4.025,2.143) (4.075,2.148) (4.125,2.152) (4.175,2.155) (4.225,2.157) (4.275,2.158) (4.325,2.158) (4.375,2.158) (4.425,2.156) (4.475,2.154) (4.525,2.150) (4.575,2.146) (4.625,2.140) (4.675,2.134) (4.725,2.126) (4.775,2.117) (4.825,2.108) (4.875,2.097) (4.925,2.085) (4.975,2.071) (5.025,2.057) (5.075,2.041) (5.125,2.024) (5.175,2.006) (5.225,1.986) (5.275,1.965) (5.325,1.942) (5.375,1.918) (5.425,1.892) (5.475,1.864) (5.525,1.835) (5.575,1.804) (5.618,1.775) (5.675,1.735) (5.722,1.700) (5.775,1.658) (5.814,1.625) (5.869,1.575) (5.921,1.525) (5.970,1.475) (6.016,1.425) (6.050,1.385) (6.098,1.325) (6.136,1.275) (6.172,1.225) (6.200,1.183) (6.236,1.125) (6.266,1.075) (6.294,1.025) (6.320,0.975) (6.344,0.925) (6.367,0.875) (6.388,0.825) (6.408,0.775) (6.426,0.725) (6.443,0.675) (6.459,0.625) (6.473,0.575) (6.486,0.525) (6.498,0.475) (6.509,0.425) (6.518,0.375) (6.526,0.325) (6.533,0.275) (6.539,0.225) (6.543,0.175) (6.547,0.125) (6.549,0.075) (6.550,0.025) (6.550,-0.037) (6.548,-0.100) (6.545,-0.150) (6.541,-0.200) (6.536,-0.250) (6.530,-0.300) (6.522,-0.350) (6.513,-0.400) (6.503,-0.450) (6.492,-0.500) (6.480,-0.550) (6.466,-0.600) (6.451,-0.650) (6.435,-0.700) (6.417,-0.750) (6.398,-0.800) (6.378,-0.850) (6.356,-0.900) (6.332,-0.950) (6.307,-1.000) (6.280,-1.050) (6.251,-1.100) (6.225,-1.143) (6.189,-1.200) (6.154,-1.250) (6.118,-1.300) (6.079,-1.350) (6.037,-1.400) (6.000,-1.442) (5.950,-1.496) (5.900,-1.546) (5.850,-1.593) (5.800,-1.637) (5.754,-1.675) (5.700,-1.717) (5.654,-1.750) (5.600,-1.787) (5.550,-1.819) (5.500,-1.850) (5.456,-1.875) (5.409,-1.900) (5.360,-1.925) (5.307,-1.950) (5.251,-1.975) (5.200,-1.996) (5.150,-2.015) (5.100,-2.033) (5.050,-2.049) (5.000,-2.064) (4.950,-2.078) (4.900,-2.091) (4.850,-2.102) (4.800,-2.113) (4.750,-2.122) (4.700,-2.130) (4.650,-2.137) (4.600,-2.143) (4.550,-2.148) (4.500,-2.152) (4.450,-2.155) (4.400,-2.157) (4.350,-2.158) (4.300,-2.158) (4.250,-2.158) (4.200,-2.156) (4.150,-2.153) (4.100,-2.150) (4.050,-2.146) (4.000,-2.140) (3.950,-2.134) (3.900,-2.128) (3.850,-2.120) (3.800,-2.111) (3.750,-2.102) (3.700,-2.092) (3.650,-2.081) (3.600,-2.070) (3.550,-2.057) (3.500,-2.044) (3.450,-2.030) (3.400,-2.016) (3.350,-2.001) (3.300,-1.985) (3.250,-1.968) (3.200,-1.951) (3.150,-1.933) (3.100,-1.914) (3.050,-1.895) (3.001,-1.875) (2.950,-1.854) (2.900,-1.833) (2.850,-1.811) (2.800,-1.789) (2.750,-1.766) (2.700,-1.742) (2.650,-1.718) (2.600,-1.694) (2.550,-1.669) (2.500,-1.643) (2.450,-1.617) (2.400,-1.590) (2.350,-1.563) (2.300,-1.535) (2.250,-1.507) (2.200,-1.479) (2.150,-1.450) (2.100,-1.421) (2.050,-1.391) (2.000,-1.361) (1.950,-1.331) (1.900,-1.300) (1.850,-1.269) (1.800,-1.238) (1.750,-1.206) (1.702,-1.175) (1.650,-1.142) (1.600,-1.109) (1.550,-1.076) (1.500,-1.044) (1.450,-1.010) (1.400,-0.977) (1.350,-0.943) (1.300,-0.910) (1.250,-0.876) (1.200,-0.842) (1.150,-0.807) (1.103,-0.775) (1.050,-0.739) (1.000,-0.704) (0.950,-0.669) (0.900,-0.634) (0.851,-0.600) (0.800,-0.565) (0.750,-0.529) (0.709,-0.500) (0.650,-0.459) (0.602,-0.425) (0.550,-0.389) (0.500,-0.354) (0.450,-0.319) (0.400,-0.283) (0.354,-0.250) (0.300,-0.213) (0.250,-0.177) (0.206,-0.144)}
\EQ{(-0.220,-0.155) (-0.275,-0.195) (-0.319,-0.225) (-0.375,-0.266) (-0.424,-0.300) (-0.475,-0.336) (-0.525,-0.371) (-0.567,-0.400) (-0.625,-0.442) (-0.673,-0.475) (-0.725,-0.512) (-0.775,-0.547) (-0.825,-0.583) (-0.875,-0.618) (-0.922,-0.650) (-0.975,-0.688) (-1.025,-0.722) (-1.075,-0.757) (-1.125,-0.792) (-1.173,-0.825) (-1.225,-0.861) (-1.275,-0.896) (-1.318,-0.925) (-1.375,-0.964) (-1.425,-0.998) (-1.475,-1.032) (-1.525,-1.066) (-1.575,-1.100) (-1.625,-1.133) (-1.675,-1.167) (-1.725,-1.200) (-1.775,-1.233) (-1.825,-1.266) (-1.875,-1.298) (-1.925,-1.330) (-1.975,-1.362) (-2.025,-1.394) (-2.074,-1.425) (-2.125,-1.457) (-2.175,-1.488) (-2.225,-1.519) (-2.275,-1.549) (-2.325,-1.580) (-2.375,-1.609) (-2.425,-1.639) (-2.475,-1.668) (-2.525,-1.697) (-2.575,-1.725) (-2.619,-1.750) (-2.665,-1.775) (-2.710,-1.800) (-2.757,-1.825) (-2.803,-1.850) (-2.850,-1.874) (-2.900,-1.900) (-2.949,-1.925) (-2.999,-1.950) (-3.050,-1.975) (-3.100,-1.999) (-3.150,-2.022) (-3.200,-2.045) (-3.250,-2.068) (-3.300,-2.090) (-3.350,-2.111) (-3.400,-2.132) (-3.450,-2.152) (-3.500,-2.172) (-3.550,-2.191) (-3.600,-2.210) (-3.650,-2.228) (-3.700,-2.246) (-3.750,-2.263) (-3.800,-2.279) (-3.850,-2.295) (-3.900,-2.310) (-3.950,-2.324) (-4.000,-2.338) (-4.050,-2.352) (-4.100,-2.364) (-4.150,-2.376) (-4.200,-2.387) (-4.250,-2.398) (-4.300,-2.408) (-4.350,-2.417) (-4.400,-2.425) (-4.450,-2.433) (-4.500,-2.440) (-4.550,-2.447) (-4.600,-2.452) (-4.650,-2.457) (-4.700,-2.462) (-4.750,-2.465) (-4.800,-2.468) (-4.850,-2.469) (-4.900,-2.470) (-4.950,-2.471) (-5.000,-2.470) (-5.050,-2.469) (-5.100,-2.467) (-5.150,-2.464) (-5.200,-2.460) (-5.250,-2.455) (-5.300,-2.449) (-5.350,-2.443) (-5.400,-2.435) (-5.450,-2.427) (-5.500,-2.417) (-5.550,-2.407) (-5.600,-2.395) (-5.650,-2.383) (-5.700,-2.370) (-5.750,-2.355) (-5.800,-2.339) (-5.850,-2.323) (-5.900,-2.305) (-5.950,-2.286) (-6.000,-2.265) (-6.050,-2.244) (-6.100,-2.221) (-6.150,-2.196) (-6.200,-2.171) (-6.250,-2.144) (-6.300,-2.115) (-6.350,-2.085) (-6.400,-2.053) (-6.450,-2.019) (-6.500,-1.983) (-6.545,-1.950) (-6.600,-1.906) (-6.650,-1.864) (-6.695,-1.825) (-6.748,-1.775) (-6.799,-1.725) (-6.846,-1.675) (-6.892,-1.625) (-6.925,-1.586) (-6.974,-1.525) (-7.013,-1.475) (-7.049,-1.425) (-7.083,-1.375) (-7.116,-1.325) (-7.147,-1.275) (-7.175,-1.226) (-7.200,-1.181) (-7.225,-1.134) (-7.250,-1.084) (-7.275,-1.031) (-7.300,-0.975) (-7.320,-0.925) (-7.340,-0.875) (-7.358,-0.825) (-7.375,-0.775) (-7.391,-0.725) (-7.405,-0.675) (-7.419,-0.625) (-7.431,-0.575) (-7.443,-0.525) (-7.453,-0.475) (-7.462,-0.425) (-7.470,-0.375) (-7.477,-0.325) (-7.483,-0.275) (-7.488,-0.225) (-7.492,-0.175) (-7.495,-0.125) (-7.497,-0.075) (-7.498,-0.025) (-7.498,0.050) (-7.496,0.100) (-7.494,0.150) (-7.490,0.200) (-7.486,0.250) (-7.480,0.300) (-7.474,0.350) (-7.466,0.400) (-7.457,0.450) (-7.448,0.500) (-7.437,0.550) (-7.425,0.600) (-7.412,0.650) (-7.398,0.700) (-7.383,0.750) (-7.367,0.800) (-7.349,0.850) (-7.330,0.900) (-7.310,0.950) (-7.289,1.000) (-7.266,1.050) (-7.242,1.100) (-7.217,1.150) (-7.190,1.200) (-7.161,1.250) (-7.131,1.300) (-7.100,1.350) (-7.066,1.400) (-7.031,1.450) (-7.000,1.492) (-6.955,1.550) (-6.913,1.600) (-6.875,1.644) (-6.825,1.698) (-6.775,1.749) (-6.725,1.797) (-6.675,1.843) (-6.625,1.886) (-6.577,1.925) (-6.525,1.965) (-6.477,2.000) (-6.425,2.036) (-6.375,2.069) (-6.325,2.100) (-6.275,2.129) (-6.225,2.157) (-6.175,2.184) (-6.125,2.209) (-6.075,2.232) (-6.025,2.255) (-5.977,2.275) (-5.925,2.295) (-5.875,2.314) (-5.825,2.331) (-5.775,2.347) (-5.725,2.362) (-5.675,2.376) (-5.625,2.389) (-5.575,2.401) (-5.525,2.412) (-5.475,2.422) (-5.425,2.431) (-5.375,2.439) (-5.325,2.446) (-5.275,2.452) (-5.225,2.457) (-5.175,2.462) (-5.125,2.465) (-5.075,2.468) (-5.025,2.470) (-4.975,2.471) (-4.925,2.471) (-4.875,2.470) (-4.825,2.469) (-4.775,2.466) (-4.725,2.463) (-4.675,2.460) (-4.625,2.455) (-4.575,2.450) (-4.525,2.444) (-4.475,2.437) (-4.425,2.430) (-4.375,2.421) (-4.325,2.412) (-4.275,2.403) (-4.225,2.393) (-4.175,2.382) (-4.125,2.370) (-4.075,2.358) (-4.025,2.345) (-3.975,2.331) (-3.925,2.317) (-3.875,2.302) (-3.825,2.287) (-3.775,2.271) (-3.725,2.254) (-3.675,2.237) (-3.625,2.219) (-3.575,2.201) (-3.525,2.182) (-3.475,2.162) (-3.425,2.142) (-3.375,2.122) (-3.325,2.100) (-3.275,2.079) (-3.225,2.057) (-3.175,2.034) (-3.125,2.011) (-3.075,1.987) (-3.025,1.963) (-2.975,1.938) (-2.925,1.913) (-2.875,1.887) (-2.825,1.861) (-2.775,1.835) (-2.725,1.808) (-2.675,1.781) (-2.625,1.753) (-2.575,1.725) (-2.531,1.700) (-2.487,1.675) (-2.425,1.639) (-2.375,1.609) (-2.325,1.580) (-2.276,1.550) (-2.225,1.519) (-2.175,1.488) (-2.125,1.457) (-2.075,1.426) (-2.025,1.394) (-1.975,1.362) (-1.925,1.330) (-1.878,1.300) (-1.825,1.266) (-1.775,1.233) (-1.725,1.200) (-1.675,1.167) (-1.625,1.133) (-1.576,1.100) (-1.525,1.066) (-1.475,1.032) (-1.428,1.000) (-1.375,0.964) (-1.325,0.930) (-1.282,0.900) (-1.225,0.861) (-1.175,0.827) (-1.125,0.792) (-1.075,0.757) (-1.029,0.725) (-0.975,0.688) (-0.925,0.652) (-0.875,0.618) (-0.825,0.583) (-0.779,0.550) (-0.725,0.512) (-0.675,0.477) (-0.625,0.442) (-0.575,0.407) (-0.531,0.375) (-0.475,0.336) (-0.425,0.300) (-0.375,0.266) (-0.325,0.230) (-0.284,0.200) (-0.225,0.160) (-0.220,0.155)}
\EQ{(0.206,0.144) (0.248,0.175) (0.300,0.213) (0.350,0.248) (0.400,0.283) (0.450,0.319) (0.496,0.350) (0.550,0.389) (0.600,0.424) (0.650,0.460) (0.700,0.495) (0.744,0.525) (0.800,0.565) (0.850,0.600) (0.900,0.635) (0.950,0.670) (0.993,0.700) (1.050,0.740) (1.100,0.774) (1.150,0.809) (1.200,0.844) (1.245,0.875) (1.300,0.913) (1.350,0.947) (1.400,0.981) (1.450,1.015) (1.500,1.049) (1.550,1.083) (1.600,1.117) (1.650,1.150) (1.700,1.183) (1.750,1.216) (1.800,1.249) (1.850,1.282) (1.900,1.314) (1.950,1.346) (1.995,1.375) (2.050,1.410) (2.100,1.442) (2.150,1.473) (2.194,1.500) (2.250,1.534) (2.300,1.565) (2.350,1.594) (2.400,1.624) (2.444,1.650) (2.500,1.682) (2.550,1.711) (2.600,1.739) (2.650,1.767) (2.700,1.795) (2.750,1.822) (2.800,1.848) (2.850,1.874) (2.900,1.900) (2.949,1.925) (2.999,1.950) (3.050,1.975) (3.100,1.999) (3.150,2.022) (3.200,2.045) (3.250,2.068) (3.300,2.090) (3.350,2.111) (3.400,2.132) (3.450,2.152) (3.500,2.172) (3.550,2.191) (3.600,2.210) (3.650,2.228) (3.700,2.246) (3.750,2.263) (3.800,2.279) (3.850,2.295) (3.900,2.310) (3.950,2.324) (4.000,2.338) (4.050,2.352) (4.100,2.364) (4.150,2.376) (4.200,2.387) (4.250,2.398) (4.300,2.408) (4.350,2.417) (4.400,2.425) (4.450,2.433) (4.500,2.440) (4.550,2.447) (4.600,2.452) (4.650,2.457) (4.700,2.462) (4.750,2.465) (4.800,2.468) (4.850,2.469) (4.900,2.470) (4.950,2.471) (5.000,2.470) (5.050,2.469) (5.100,2.467) (5.150,2.464) (5.200,2.460) (5.250,2.455) (5.300,2.449) (5.350,2.443) (5.400,2.435) (5.450,2.427) (5.500,2.417) (5.550,2.407) (5.600,2.395) (5.650,2.383) (5.700,2.370) (5.750,2.355) (5.800,2.339) (5.850,2.323) (5.900,2.305) (5.950,2.286) (6.000,2.265) (6.050,2.244) (6.100,2.221) (6.150,2.196) (6.200,2.171) (6.250,2.144) (6.300,2.115) (6.350,2.085) (6.400,2.053) (6.450,2.019) (6.500,1.983) (6.545,1.950) (6.600,1.906) (6.650,1.864) (6.695,1.825) (6.748,1.775) (6.799,1.725) (6.846,1.675) (6.892,1.625) (6.925,1.586) (6.974,1.525) (7.013,1.475) (7.049,1.425) (7.083,1.375) (7.116,1.325) (7.147,1.275) (7.175,1.226) (7.200,1.181) (7.225,1.134) (7.250,1.084) (7.275,1.031) (7.300,0.975) (7.320,0.925) (7.340,0.875) (7.358,0.825) (7.375,0.775) (7.391,0.725) (7.405,0.675) (7.419,0.625) (7.431,0.575) (7.443,0.525) (7.453,0.475) (7.462,0.425) (7.470,0.375) (7.477,0.325) (7.483,0.275) (7.488,0.225) (7.492,0.175) (7.495,0.125) (7.497,0.075) (7.498,0.025) (7.498,-0.050) (7.496,-0.100) (7.494,-0.150) (7.490,-0.200) (7.486,-0.250) (7.480,-0.300) (7.474,-0.350) (7.466,-0.400) (7.457,-0.450) (7.448,-0.500) (7.437,-0.550) (7.425,-0.600) (7.412,-0.650) (7.398,-0.700) (7.383,-0.750) (7.367,-0.800) (7.349,-0.850) (7.330,-0.900) (7.310,-0.950) (7.289,-1.000) (7.266,-1.050) (7.242,-1.100) (7.217,-1.150) (7.190,-1.200) (7.161,-1.250) (7.131,-1.300) (7.100,-1.350) (7.066,-1.400) (7.031,-1.450) (7.000,-1.492) (6.955,-1.550) (6.913,-1.600) (6.875,-1.644) (6.825,-1.698) (6.775,-1.749) (6.725,-1.797) (6.675,-1.843) (6.625,-1.886) (6.577,-1.925) (6.525,-1.965) (6.477,-2.000) (6.425,-2.036) (6.375,-2.069) (6.325,-2.100) (6.275,-2.129) (6.225,-2.157) (6.175,-2.184) (6.125,-2.209) (6.075,-2.232) (6.025,-2.255) (5.977,-2.275) (5.925,-2.295) (5.875,-2.314) (5.825,-2.331) (5.775,-2.347) (5.725,-2.362) (5.675,-2.376) (5.625,-2.389) (5.575,-2.401) (5.525,-2.412) (5.475,-2.422) (5.425,-2.431) (5.375,-2.439) (5.325,-2.446) (5.275,-2.452) (5.225,-2.457) (5.175,-2.462) (5.125,-2.465) (5.075,-2.468) (5.025,-2.470) (4.975,-2.471) (4.925,-2.471) (4.875,-2.470) (4.825,-2.469) (4.775,-2.466) (4.725,-2.463) (4.675,-2.460) (4.625,-2.455) (4.575,-2.450) (4.525,-2.444) (4.475,-2.437) (4.425,-2.430) (4.375,-2.421) (4.325,-2.412) (4.275,-2.403) (4.225,-2.393) (4.175,-2.382) (4.125,-2.370) (4.075,-2.358) (4.025,-2.345) (3.975,-2.331) (3.925,-2.317) (3.875,-2.302) (3.825,-2.287) (3.775,-2.271) (3.725,-2.254) (3.675,-2.237) (3.625,-2.219) (3.575,-2.201) (3.525,-2.182) (3.475,-2.162) (3.425,-2.142) (3.375,-2.122) (3.325,-2.100) (3.275,-2.079) (3.225,-2.057) (3.175,-2.034) (3.125,-2.011) (3.075,-1.987) (3.025,-1.963) (2.975,-1.938) (2.925,-1.913) (2.875,-1.887) (2.825,-1.861) (2.775,-1.835) (2.725,-1.808) (2.675,-1.781) (2.625,-1.753) (2.575,-1.725) (2.531,-1.700) (2.487,-1.675) (2.425,-1.639) (2.375,-1.609) (2.325,-1.580) (2.276,-1.550) (2.225,-1.519) (2.175,-1.488) (2.125,-1.457) (2.075,-1.426) (2.025,-1.394) (1.975,-1.362) (1.925,-1.330) (1.878,-1.300) (1.825,-1.266) (1.775,-1.233) (1.725,-1.200) (1.675,-1.167) (1.625,-1.133) (1.576,-1.100) (1.525,-1.066) (1.475,-1.032) (1.428,-1.000) (1.375,-0.964) (1.325,-0.930) (1.282,-0.900) (1.225,-0.861) (1.175,-0.827) (1.125,-0.792) (1.075,-0.757) (1.029,-0.725) (0.975,-0.688) (0.925,-0.652) (0.875,-0.618) (0.825,-0.583) (0.779,-0.550) (0.725,-0.512) (0.675,-0.477) (0.625,-0.442) (0.575,-0.407) (0.531,-0.375) (0.475,-0.336) (0.425,-0.300) (0.375,-0.266) (0.325,-0.230) (0.284,-0.200) (0.225,-0.160) (0.206,-0.144)}
\EQ{(-0.220,-0.155) (-0.275,-0.195) (-0.319,-0.225) (-0.375,-0.266) (-0.424,-0.300) (-0.475,-0.336) (-0.525,-0.371) (-0.567,-0.400) (-0.625,-0.442) (-0.672,-0.475) (-0.725,-0.513) (-0.775,-0.548) (-0.825,-0.583) (-0.875,-0.618) (-0.921,-0.650) (-0.975,-0.689) (-1.025,-0.724) (-1.075,-0.759) (-1.125,-0.794) (-1.170,-0.825) (-1.225,-0.864) (-1.275,-0.899) (-1.325,-0.934) (-1.375,-0.968) (-1.421,-1.000) (-1.475,-1.038) (-1.525,-1.072) (-1.575,-1.107) (-1.625,-1.141) (-1.674,-1.175) (-1.725,-1.210) (-1.775,-1.244) (-1.821,-1.275) (-1.875,-1.312) (-1.925,-1.346) (-1.969,-1.375) (-2.025,-1.413) (-2.075,-1.446) (-2.118,-1.475) (-2.175,-1.513) (-2.225,-1.546) (-2.270,-1.575) (-2.325,-1.611) (-2.375,-1.644) (-2.424,-1.675) (-2.475,-1.708) (-2.525,-1.740) (-2.575,-1.772) (-2.620,-1.800) (-2.675,-1.834) (-2.725,-1.866) (-2.775,-1.896) (-2.822,-1.925) (-2.875,-1.957) (-2.925,-1.988) (-2.975,-2.017) (-3.025,-2.047) (-3.073,-2.075) (-3.125,-2.105) (-3.175,-2.134) (-3.225,-2.163) (-3.275,-2.191) (-3.325,-2.219) (-3.375,-2.246) (-3.425,-2.274) (-3.474,-2.300) (-3.521,-2.325) (-3.569,-2.350) (-3.617,-2.375) (-3.665,-2.400) (-3.715,-2.425) (-3.765,-2.450) (-3.816,-2.475) (-3.867,-2.500) (-3.920,-2.525) (-3.973,-2.550) (-4.025,-2.574) (-4.075,-2.596) (-4.125,-2.619) (-4.175,-2.640) (-4.225,-2.662) (-4.275,-2.683) (-4.325,-2.703) (-4.375,-2.723) (-4.425,-2.743) (-4.475,-2.762) (-4.525,-2.781) (-4.575,-2.799) (-4.625,-2.817) (-4.675,-2.834) (-4.725,-2.851) (-4.775,-2.867) (-4.825,-2.883) (-4.875,-2.899) (-4.925,-2.914) (-4.975,-2.928) (-5.025,-2.942) (-5.075,-2.956) (-5.125,-2.968) (-5.175,-2.981) (-5.225,-2.993) (-5.275,-3.004) (-5.325,-3.015) (-5.375,-3.025) (-5.425,-3.035) (-5.475,-3.044) (-5.525,-3.053) (-5.575,-3.061) (-5.625,-3.069) (-5.675,-3.075) (-5.725,-3.082) (-5.775,-3.088) (-5.825,-3.093) (-5.875,-3.098) (-5.925,-3.102) (-5.975,-3.105) (-6.025,-3.108) (-6.075,-3.110) (-6.125,-3.112) (-6.175,-3.113) (-6.225,-3.113) (-6.275,-3.113) (-6.325,-3.112) (-6.375,-3.110) (-6.425,-3.108) (-6.475,-3.105) (-6.525,-3.101) (-6.575,-3.097) (-6.625,-3.092) (-6.675,-3.086) (-6.725,-3.080) (-6.775,-3.072) (-6.825,-3.065) (-6.875,-3.056) (-6.925,-3.046) (-6.975,-3.036) (-7.025,-3.025) (-7.075,-3.013) (-7.125,-3.001) (-7.175,-2.987) (-7.225,-2.973) (-7.275,-2.958) (-7.325,-2.942) (-7.375,-2.925) (-7.425,-2.907) (-7.475,-2.888) (-7.525,-2.868) (-7.575,-2.848) (-7.625,-2.826) (-7.675,-2.803) (-7.725,-2.779) (-7.775,-2.754) (-7.825,-2.728) (-7.875,-2.701) (-7.920,-2.675) (-7.975,-2.642) (-8.025,-2.611) (-8.075,-2.579) (-8.117,-2.550) (-8.175,-2.509) (-8.222,-2.475) (-8.275,-2.434) (-8.318,-2.400) (-8.375,-2.352) (-8.425,-2.308) (-8.475,-2.262) (-8.514,-2.225) (-8.550,-2.189) (-8.589,-2.150) (-8.625,-2.111) (-8.675,-2.056) (-8.723,-2.000) (-8.765,-1.950) (-8.800,-1.905) (-8.842,-1.850) (-8.875,-1.804) (-8.912,-1.750) (-8.946,-1.700) (-8.975,-1.654) (-9.008,-1.600) (-9.037,-1.550) (-9.065,-1.500) (-9.092,-1.450) (-9.118,-1.400) (-9.142,-1.350) (-9.165,-1.300) (-9.188,-1.250) (-9.209,-1.200) (-9.229,-1.150) (-9.248,-1.100) (-9.267,-1.050) (-9.284,-1.000) (-9.300,-0.951) (-9.316,-0.900) (-9.330,-0.850) (-9.344,-0.800) (-9.357,-0.750) (-9.368,-0.700) (-9.379,-0.650) (-9.390,-0.600) (-9.399,-0.550) (-9.407,-0.500) (-9.415,-0.450) (-9.422,-0.400) (-9.428,-0.350) (-9.433,-0.300) (-9.437,-0.250) (-9.441,-0.200) (-9.444,-0.150) (-9.446,-0.100) (-9.447,-0.050) (-9.447,0.000) (-9.447,0.050) (-9.446,0.100) (-9.444,0.150) (-9.441,0.200) (-9.437,0.250) (-9.433,0.300) (-9.428,0.350) (-9.422,0.400) (-9.415,0.450) (-9.407,0.500) (-9.399,0.550) (-9.390,0.600) (-9.379,0.650) (-9.368,0.700) (-9.357,0.750) (-9.344,0.800) (-9.330,0.850) (-9.316,0.900) (-9.300,0.950) (-9.284,1.000) (-9.267,1.050) (-9.248,1.100) (-9.229,1.150) (-9.209,1.200) (-9.188,1.250) (-9.165,1.300) (-9.142,1.350) (-9.118,1.400) (-9.092,1.450) (-9.065,1.500) (-9.037,1.550) (-9.008,1.600) (-8.978,1.650) (-8.950,1.693) (-8.912,1.750) (-8.878,1.800) (-8.842,1.850) (-8.804,1.900) (-8.765,1.950) (-8.725,1.998) (-8.680,2.050) (-8.636,2.100) (-8.600,2.138) (-8.564,2.175) (-8.525,2.214) (-8.488,2.250) (-8.450,2.285) (-8.406,2.325) (-8.350,2.373) (-8.300,2.414) (-8.255,2.450) (-8.200,2.491) (-8.153,2.525) (-8.100,2.562) (-8.050,2.595) (-8.003,2.625) (-7.950,2.657) (-7.900,2.686) (-7.850,2.714) (-7.800,2.741) (-7.750,2.767) (-7.700,2.791) (-7.650,2.814) (-7.600,2.837) (-7.550,2.858) (-7.500,2.878) (-7.450,2.898) (-7.400,2.916) (-7.350,2.933) (-7.300,2.950) (-7.250,2.966) (-7.200,2.980) (-7.150,2.994) (-7.100,3.007) (-7.050,3.019) (-7.000,3.031) (-6.950,3.041) (-6.900,3.051) (-6.850,3.060) (-6.800,3.069) (-6.750,3.076) (-6.700,3.083) (-6.650,3.089) (-6.600,3.094) (-6.550,3.099) (-6.500,3.103) (-6.450,3.106) (-6.400,3.109) (-6.350,3.111) (-6.300,3.112) (-6.250,3.113) (-6.200,3.113) (-6.150,3.112) (-6.100,3.111) (-6.050,3.109) (-6.000,3.107) (-5.950,3.103) (-5.900,3.100) (-5.850,3.095) (-5.800,3.090) (-5.750,3.085) (-5.700,3.079) (-5.650,3.072) (-5.600,3.065) (-5.550,3.057) (-5.500,3.049) (-5.450,3.040) (-5.400,3.030) (-5.350,3.020) (-5.300,3.010) (-5.250,2.998) (-5.200,2.987) (-5.151,2.975) (-5.100,2.962) (-5.050,2.949) (-5.000,2.935) (-4.950,2.921) (-4.900,2.906) (-4.850,2.891) (-4.800,2.875) (-4.750,2.859) (-4.700,2.843) (-4.650,2.825) (-4.600,2.808) (-4.550,2.790) (-4.500,2.771) (-4.450,2.752) (-4.400,2.733) (-4.350,2.713) (-4.300,2.693) (-4.250,2.672) (-4.200,2.651) (-4.150,2.629) (-4.100,2.608) (-4.050,2.585) (-4.000,2.562) (-3.950,2.539) (-3.900,2.516) (-3.850,2.492) (-3.800,2.467) (-3.750,2.443) (-3.700,2.418) (-3.650,2.392) (-3.600,2.366) (-3.550,2.340) (-3.500,2.314) (-3.450,2.287) (-3.400,2.260) (-3.350,2.233) (-3.300,2.205) (-3.250,2.177) (-3.203,2.150) (-3.159,2.125) (-3.116,2.100) (-3.050,2.062) (-3.000,2.032) (-2.950,2.002) (-2.904,1.975) (-2.850,1.942) (-2.800,1.912) (-2.750,1.881) (-2.700,1.850) (-2.650,1.819) (-2.600,1.787) (-2.550,1.756) (-2.502,1.725) (-2.450,1.692) (-2.400,1.660) (-2.350,1.627) (-2.300,1.595) (-2.250,1.562) (-2.200,1.529) (-2.156,1.500) (-2.100,1.463) (-2.050,1.430) (-2.006,1.400) (-1.950,1.363) (-1.900,1.329) (-1.858,1.300) (-1.800,1.261) (-1.750,1.227) (-1.700,1.193) (-1.650,1.158) (-1.602,1.125) (-1.550,1.090) (-1.500,1.055) (-1.457,1.025) (-1.400,0.986) (-1.350,0.951) (-1.300,0.916) (-1.250,0.881) (-1.206,0.850) (-1.150,0.811) (-1.100,0.776) (-1.050,0.741) (-1.000,0.706) (-0.956,0.675) (-0.900,0.636) (-0.850,0.600) (-0.800,0.566) (-0.750,0.530) (-0.708,0.500) (-0.650,0.460) (-0.601,0.425) (-0.550,0.389) (-0.500,0.354) (-0.450,0.319) (-0.400,0.283) (-0.354,0.250) (-0.300,0.213) (-0.250,0.177) (-0.220,0.155)}
\EQ{(0.206,0.144) (0.248,0.175) (0.300,0.213) (0.350,0.248) (0.400,0.283) (0.450,0.319) (0.496,0.350) (0.550,0.389) (0.600,0.424) (0.650,0.460) (0.700,0.495) (0.743,0.525) (0.800,0.566) (0.849,0.600) (0.900,0.636) (0.950,0.671) (0.992,0.700) (1.050,0.741) (1.098,0.775) (1.150,0.811) (1.200,0.846) (1.242,0.875) (1.300,0.916) (1.349,0.950) (1.400,0.986) (1.450,1.020) (1.493,1.050) (1.550,1.090) (1.600,1.124) (1.650,1.158) (1.700,1.193) (1.748,1.225) (1.800,1.261) (1.850,1.295) (1.895,1.325) (1.950,1.363) (2.000,1.396) (2.043,1.425) (2.100,1.463) (2.150,1.496) (2.194,1.525) (2.250,1.562) (2.300,1.595) (2.347,1.625) (2.400,1.660) (2.450,1.692) (2.500,1.724) (2.550,1.756) (2.600,1.787) (2.650,1.819) (2.700,1.850) (2.750,1.881) (2.800,1.912) (2.850,1.942) (2.900,1.972) (2.946,2.000) (3.000,2.032) (3.050,2.062) (3.100,2.091) (3.150,2.120) (3.200,2.148) (3.247,2.175) (3.292,2.200) (3.337,2.225) (3.382,2.250) (3.428,2.275) (3.474,2.300) (3.521,2.325) (3.569,2.350) (3.617,2.375) (3.665,2.400) (3.715,2.425) (3.765,2.450) (3.816,2.475) (3.867,2.500) (3.920,2.525) (3.973,2.550) (4.025,2.574) (4.075,2.596) (4.125,2.619) (4.175,2.640) (4.225,2.662) (4.275,2.683) (4.325,2.703) (4.375,2.723) (4.425,2.743) (4.475,2.762) (4.525,2.781) (4.575,2.799) (4.625,2.817) (4.675,2.834) (4.725,2.851) (4.775,2.867) (4.825,2.883) (4.875,2.899) (4.925,2.914) (4.975,2.928) (5.025,2.942) (5.075,2.956) (5.125,2.968) (5.175,2.981) (5.225,2.993) (5.275,3.004) (5.325,3.015) (5.375,3.025) (5.425,3.035) (5.475,3.044) (5.525,3.053) (5.575,3.061) (5.625,3.069) (5.675,3.075) (5.725,3.082) (5.775,3.088) (5.825,3.093) (5.875,3.098) (5.925,3.102) (5.975,3.105) (6.025,3.108) (6.075,3.110) (6.125,3.112) (6.175,3.113) (6.225,3.113) (6.275,3.113) (6.325,3.112) (6.375,3.110) (6.425,3.108) (6.475,3.105) (6.525,3.101) (6.575,3.097) (6.625,3.092) (6.675,3.086) (6.725,3.080) (6.775,3.072) (6.825,3.065) (6.875,3.056) (6.925,3.046) (6.975,3.036) (7.025,3.025) (7.075,3.013) (7.125,3.001) (7.175,2.987) (7.225,2.973) (7.275,2.958) (7.325,2.942) (7.375,2.925) (7.425,2.907) (7.475,2.888) (7.525,2.868) (7.575,2.848) (7.625,2.826) (7.675,2.803) (7.725,2.779) (7.775,2.754) (7.825,2.728) (7.875,2.701) (7.920,2.675) (7.975,2.642) (8.025,2.611) (8.075,2.579) (8.117,2.550) (8.175,2.509) (8.222,2.475) (8.275,2.434) (8.318,2.400) (8.375,2.352) (8.425,2.308) (8.475,2.262) (8.514,2.225) (8.550,2.189) (8.589,2.150) (8.625,2.111) (8.675,2.056) (8.723,2.000) (8.765,1.950) (8.800,1.905) (8.842,1.850) (8.875,1.804) (8.912,1.750) (8.946,1.700) (8.975,1.654) (9.008,1.600) (9.037,1.550) (9.065,1.500) (9.092,1.450) (9.118,1.400) (9.142,1.350) (9.165,1.300) (9.188,1.250) (9.209,1.200) (9.229,1.150) (9.248,1.100) (9.267,1.050) (9.284,1.000) (9.300,0.951) (9.316,0.900) (9.330,0.850) (9.344,0.800) (9.357,0.750) (9.368,0.700) (9.379,0.650) (9.390,0.600) (9.399,0.550) (9.407,0.500) (9.415,0.450) (9.422,0.400) (9.428,0.350) (9.433,0.300) (9.437,0.250) (9.441,0.200) (9.444,0.150) (9.446,0.100) (9.447,0.050) (9.447,0.000) (9.447,-0.050) (9.446,-0.100) (9.444,-0.150) (9.441,-0.200) (9.437,-0.250) (9.433,-0.300) (9.428,-0.350) (9.422,-0.400) (9.415,-0.450) (9.407,-0.500) (9.399,-0.550) (9.390,-0.600) (9.379,-0.650) (9.368,-0.700) (9.357,-0.750) (9.344,-0.800) (9.330,-0.850) (9.316,-0.900) (9.300,-0.950) (9.284,-1.000) (9.267,-1.050) (9.248,-1.100) (9.229,-1.150) (9.209,-1.200) (9.188,-1.250) (9.165,-1.300) (9.142,-1.350) (9.118,-1.400) (9.092,-1.450) (9.065,-1.500) (9.037,-1.550) (9.008,-1.600) (8.978,-1.650) (8.950,-1.693) (8.912,-1.750) (8.878,-1.800) (8.842,-1.850) (8.804,-1.900) (8.765,-1.950) (8.725,-1.998) (8.680,-2.050) (8.636,-2.100) (8.600,-2.138) (8.564,-2.175) (8.525,-2.214) (8.488,-2.250) (8.450,-2.285) (8.406,-2.325) (8.350,-2.373) (8.300,-2.414) (8.255,-2.450) (8.200,-2.491) (8.153,-2.525) (8.100,-2.562) (8.050,-2.595) (8.003,-2.625) (7.950,-2.657) (7.900,-2.686) (7.850,-2.714) (7.800,-2.741) (7.750,-2.767) (7.700,-2.791) (7.650,-2.814) (7.600,-2.837) (7.550,-2.858) (7.500,-2.878) (7.450,-2.898) (7.400,-2.916) (7.350,-2.933) (7.300,-2.950) (7.250,-2.966) (7.200,-2.980) (7.150,-2.994) (7.100,-3.007) (7.050,-3.019) (7.000,-3.031) (6.950,-3.041) (6.900,-3.051) (6.850,-3.060) (6.800,-3.069) (6.750,-3.076) (6.700,-3.083) (6.650,-3.089) (6.600,-3.094) (6.550,-3.099) (6.500,-3.103) (6.450,-3.106) (6.400,-3.109) (6.350,-3.111) (6.300,-3.112) (6.250,-3.113) (6.200,-3.113) (6.150,-3.112) (6.100,-3.111) (6.050,-3.109) (6.000,-3.107) (5.950,-3.103) (5.900,-3.100) (5.850,-3.095) (5.800,-3.090) (5.750,-3.085) (5.700,-3.079) (5.650,-3.072) (5.600,-3.065) (5.550,-3.057) (5.500,-3.049) (5.450,-3.040) (5.400,-3.030) (5.350,-3.020) (5.300,-3.010) (5.250,-2.998) (5.200,-2.987) (5.151,-2.975) (5.100,-2.962) (5.050,-2.949) (5.000,-2.935) (4.950,-2.921) (4.900,-2.906) (4.850,-2.891) (4.800,-2.875) (4.750,-2.859) (4.700,-2.843) (4.650,-2.825) (4.600,-2.808) (4.550,-2.790) (4.500,-2.771) (4.450,-2.752) (4.400,-2.733) (4.350,-2.713) (4.300,-2.693) (4.250,-2.672) (4.200,-2.651) (4.150,-2.629) (4.100,-2.608) (4.050,-2.585) (4.000,-2.562) (3.950,-2.539) (3.900,-2.516) (3.850,-2.492) (3.800,-2.467) (3.750,-2.443) (3.700,-2.418) (3.650,-2.392) (3.600,-2.366) (3.550,-2.340) (3.500,-2.314) (3.450,-2.287) (3.400,-2.260) (3.350,-2.233) (3.300,-2.205) (3.250,-2.177) (3.203,-2.150) (3.159,-2.125) (3.116,-2.100) (3.050,-2.062) (3.000,-2.032) (2.950,-2.002) (2.904,-1.975) (2.850,-1.942) (2.800,-1.912) (2.750,-1.881) (2.700,-1.850) (2.650,-1.819) (2.600,-1.787) (2.550,-1.756) (2.502,-1.725) (2.450,-1.692) (2.400,-1.660) (2.350,-1.627) (2.300,-1.595) (2.250,-1.562) (2.200,-1.529) (2.156,-1.500) (2.100,-1.463) (2.050,-1.430) (2.006,-1.400) (1.950,-1.363) (1.900,-1.329) (1.858,-1.300) (1.800,-1.261) (1.750,-1.227) (1.700,-1.193) (1.650,-1.158) (1.602,-1.125) (1.550,-1.090) (1.500,-1.055) (1.457,-1.025) (1.400,-0.986) (1.350,-0.951) (1.300,-0.916) (1.250,-0.881) (1.206,-0.850) (1.150,-0.811) (1.100,-0.776) (1.050,-0.741) (1.000,-0.706) (0.956,-0.675) (0.900,-0.636) (0.850,-0.600) (0.800,-0.566) (0.750,-0.530) (0.708,-0.500) (0.650,-0.460) (0.601,-0.425) (0.550,-0.389) (0.500,-0.354) (0.450,-0.319) (0.400,-0.283) (0.354,-0.250) (0.300,-0.213) (0.250,-0.177) (0.206,-0.144)}
\EQ{(10.000,-7.071) (9.950,-7.036) (9.900,-7.000) (9.850,-6.965) (9.800,-6.930) (9.758,-6.900) (9.700,-6.859) (9.652,-6.825) (9.600,-6.788) (9.550,-6.753) (9.500,-6.718) (9.450,-6.682) (9.405,-6.650) (9.350,-6.611) (9.300,-6.576) (9.250,-6.541) (9.200,-6.505) (9.157,-6.475) (9.100,-6.435) (9.051,-6.400) (9.000,-6.364) (8.950,-6.329) (8.900,-6.293) (8.850,-6.258) (8.804,-6.225) (8.750,-6.187) (8.700,-6.152) (8.650,-6.117) (8.600,-6.081) (8.556,-6.050) (8.500,-6.010) (8.450,-5.975) (8.400,-5.940) (8.350,-5.904) (8.309,-5.875) (8.250,-5.834) (8.202,-5.800) (8.150,-5.763) (8.100,-5.728) (8.050,-5.692) (8.000,-5.657) (7.955,-5.625) (7.900,-5.586) (7.850,-5.551) (7.800,-5.515) (7.750,-5.480) (7.708,-5.450) (7.650,-5.409) (7.601,-5.375) (7.550,-5.339) (7.500,-5.303) (7.450,-5.268) (7.400,-5.233) (7.354,-5.200) (7.300,-5.162) (7.250,-5.127) (7.200,-5.091) (7.150,-5.056) (7.106,-5.025) (7.050,-4.985) (7.000,-4.950) (6.950,-4.914) (6.900,-4.879) (6.859,-4.850) (6.800,-4.808) (6.753,-4.775) (6.700,-4.738) (6.650,-4.702) (6.600,-4.667) (6.550,-4.632) (6.505,-4.600) (6.450,-4.561) (6.400,-4.525) (6.350,-4.490) (6.300,-4.455) (6.258,-4.425) (6.200,-4.384) (6.152,-4.350) (6.100,-4.313) (6.050,-4.278) (6.000,-4.243) (5.950,-4.207) (5.904,-4.175) (5.850,-4.137) (5.800,-4.101) (5.750,-4.066) (5.700,-4.031) (5.657,-4.000) (5.600,-3.960) (5.551,-3.925) (5.500,-3.889) (5.450,-3.854) (5.400,-3.818) (5.350,-3.783) (5.303,-3.750) (5.250,-3.712) (5.200,-3.677) (5.150,-3.642) (5.100,-3.606) (5.056,-3.575) (5.000,-3.536) (4.950,-3.500) (4.900,-3.465) (4.850,-3.429) (4.808,-3.400) (4.750,-3.359) (4.702,-3.325) (4.650,-3.288) (4.600,-3.253) (4.550,-3.217) (4.500,-3.182) (4.455,-3.150) (4.400,-3.111) (4.350,-3.076) (4.300,-3.041) (4.250,-3.005) (4.207,-2.975) (4.150,-2.935) (4.101,-2.900) (4.050,-2.864) (4.000,-2.828) (3.950,-2.793) (3.900,-2.758) (3.854,-2.725) (3.800,-2.687) (3.750,-2.652) (3.700,-2.616) (3.650,-2.581) (3.606,-2.550) (3.550,-2.510) (3.500,-2.475) (3.450,-2.440) (3.400,-2.404) (3.359,-2.375) (3.300,-2.334) (3.253,-2.300) (3.200,-2.263) (3.150,-2.227) (3.100,-2.192) (3.050,-2.157) (3.005,-2.125) (2.950,-2.086) (2.900,-2.051) (2.850,-2.015) (2.800,-1.980) (2.758,-1.950) (2.700,-1.909) (2.652,-1.875) (2.600,-1.839) (2.550,-1.803) (2.500,-1.768) (2.450,-1.732) (2.404,-1.700) (2.350,-1.662) (2.300,-1.626) (2.250,-1.591) (2.200,-1.556) (2.157,-1.525) (2.100,-1.485) (2.051,-1.450) (2.000,-1.414) (1.950,-1.379) (1.900,-1.344) (1.850,-1.308) (1.803,-1.275) (1.750,-1.238) (1.700,-1.202) (1.650,-1.167) (1.600,-1.131) (1.556,-1.100) (1.500,-1.061) (1.450,-1.025) (1.400,-0.990) (1.350,-0.955) (1.308,-0.925) (1.250,-0.884) (1.202,-0.850) (1.150,-0.813) (1.100,-0.778) (1.050,-0.743) (1.000,-0.707) (0.955,-0.675) (0.900,-0.637) (0.850,-0.601) (0.800,-0.566) (0.750,-0.531) (0.708,-0.500) (0.650,-0.460) (0.601,-0.425) (0.550,-0.389) (0.500,-0.354) (0.450,-0.319) (0.400,-0.283) (0.354,-0.250) (0.300,-0.213) (0.250,-0.177) (0.206,-0.144)}
\EQ{(-0.220,-0.155) (-0.275,-0.195) (-0.319,-0.225) (-0.375,-0.266) (-0.424,-0.300) (-0.475,-0.336) (-0.525,-0.371) (-0.566,-0.400) (-0.625,-0.442) (-0.672,-0.475) (-0.725,-0.513) (-0.775,-0.548) (-0.825,-0.584) (-0.875,-0.619) (-0.920,-0.650) (-0.975,-0.690) (-1.025,-0.725) (-1.075,-0.760) (-1.125,-0.796) (-1.167,-0.825) (-1.225,-0.866) (-1.273,-0.900) (-1.325,-0.937) (-1.375,-0.972) (-1.425,-1.008) (-1.475,-1.043) (-1.520,-1.075) (-1.575,-1.114) (-1.625,-1.149) (-1.675,-1.185) (-1.725,-1.220) (-1.768,-1.250) (-1.825,-1.291) (-1.874,-1.325) (-1.925,-1.361) (-1.975,-1.397) (-2.025,-1.432) (-2.075,-1.467) (-2.121,-1.500) (-2.175,-1.538) (-2.225,-1.573) (-2.275,-1.609) (-2.325,-1.644) (-2.369,-1.675) (-2.425,-1.715) (-2.475,-1.750) (-2.525,-1.786) (-2.575,-1.821) (-2.616,-1.850) (-2.675,-1.892) (-2.722,-1.925) (-2.775,-1.962) (-2.825,-1.998) (-2.875,-2.033) (-2.925,-2.068) (-2.970,-2.100) (-3.025,-2.139) (-3.075,-2.174) (-3.125,-2.210) (-3.175,-2.245) (-3.217,-2.275) (-3.275,-2.316) (-3.323,-2.350) (-3.375,-2.387) (-3.425,-2.422) (-3.475,-2.457) (-3.525,-2.493) (-3.571,-2.525) (-3.625,-2.563) (-3.675,-2.599) (-3.725,-2.634) (-3.775,-2.669) (-3.818,-2.700) (-3.875,-2.740) (-3.924,-2.775) (-3.975,-2.811) (-4.025,-2.846) (-4.066,-2.875) (-4.125,-2.917) (-4.172,-2.950) (-4.225,-2.988) (-4.275,-3.023) (-4.325,-3.058) (-4.375,-3.094) (-4.419,-3.125) (-4.475,-3.164) (-4.525,-3.200) (-4.575,-3.235) (-4.625,-3.270) (-4.667,-3.300) (-4.725,-3.341) (-4.773,-3.375) (-4.825,-3.412) (-4.875,-3.447) (-4.925,-3.483) (-4.975,-3.518) (-5.021,-3.550) (-5.075,-3.589) (-5.125,-3.624) (-5.175,-3.659) (-5.225,-3.695) (-5.268,-3.725) (-5.325,-3.765) (-5.374,-3.800) (-5.425,-3.836) (-5.475,-3.871) (-5.525,-3.907) (-5.575,-3.942) (-5.622,-3.975) (-5.675,-4.013) (-5.725,-4.048) (-5.775,-4.084) (-5.825,-4.119) (-5.869,-4.150) (-5.925,-4.190) (-5.975,-4.225) (-6.025,-4.260) (-6.075,-4.296) (-6.117,-4.325) (-6.175,-4.366) (-6.223,-4.400) (-6.275,-4.437) (-6.325,-4.472) (-6.375,-4.508) (-6.425,-4.543) (-6.470,-4.575) (-6.525,-4.614) (-6.575,-4.649) (-6.625,-4.685) (-6.675,-4.720) (-6.718,-4.750) (-6.775,-4.791) (-6.824,-4.825) (-6.875,-4.861) (-6.925,-4.897) (-6.975,-4.932) (-7.025,-4.967) (-7.071,-5.000) (-7.125,-5.038) (-7.175,-5.073) (-7.225,-5.109) (-7.275,-5.144) (-7.319,-5.175) (-7.375,-5.215) (-7.425,-5.250) (-7.475,-5.286) (-7.525,-5.321) (-7.566,-5.350) (-7.625,-5.392) (-7.672,-5.425) (-7.725,-5.462) (-7.775,-5.498) (-7.825,-5.533) (-7.875,-5.568) (-7.920,-5.600) (-7.975,-5.639) (-8.025,-5.675) (-8.075,-5.710) (-8.125,-5.745) (-8.167,-5.775) (-8.225,-5.816) (-8.273,-5.850) (-8.325,-5.887) (-8.375,-5.922) (-8.425,-5.957) (-8.475,-5.993) (-8.521,-6.025) (-8.575,-6.063) (-8.625,-6.099) (-8.675,-6.134) (-8.725,-6.170) (-8.768,-6.200) (-8.825,-6.240) (-8.874,-6.275) (-8.925,-6.311) (-8.975,-6.346) (-9.025,-6.382) (-9.075,-6.417) (-9.122,-6.450) (-9.175,-6.488) (-9.225,-6.523) (-9.275,-6.558) (-9.325,-6.594) (-9.369,-6.625) (-9.425,-6.665) (-9.475,-6.700) (-9.525,-6.735) (-9.575,-6.771) (-9.617,-6.800) (-9.675,-6.841) (-9.723,-6.875) (-9.775,-6.912) (-9.825,-6.947) (-9.875,-6.983) (-9.925,-7.018) (-9.970,-7.050) (-10.000,-7.071)}
\EQ{(0.206,0.144) (0.248,0.175) (0.300,0.213) (0.350,0.248) (0.400,0.283) (0.450,0.319) (0.496,0.350) (0.550,0.389) (0.600,0.424) (0.650,0.460) (0.700,0.495) (0.743,0.525) (0.800,0.566) (0.849,0.600) (0.900,0.637) (0.950,0.672) (1.000,0.707) (1.050,0.743) (1.096,0.775) (1.150,0.813) (1.200,0.849) (1.250,0.884) (1.300,0.919) (1.344,0.950) (1.400,0.990) (1.450,1.025) (1.500,1.061) (1.550,1.096) (1.591,1.125) (1.650,1.167) (1.697,1.200) (1.750,1.238) (1.800,1.273) (1.850,1.308) (1.900,1.344) (1.945,1.375) (2.000,1.414) (2.050,1.450) (2.100,1.485) (2.150,1.520) (2.192,1.550) (2.250,1.591) (2.298,1.625) (2.350,1.662) (2.400,1.697) (2.450,1.732) (2.500,1.768) (2.546,1.800) (2.600,1.839) (2.650,1.874) (2.700,1.909) (2.750,1.945) (2.793,1.975) (2.850,2.015) (2.899,2.050) (2.950,2.086) (3.000,2.121) (3.050,2.157) (3.100,2.192) (3.147,2.225) (3.200,2.263) (3.250,2.298) (3.300,2.334) (3.350,2.369) (3.394,2.400) (3.450,2.440) (3.500,2.475) (3.550,2.510) (3.600,2.546) (3.642,2.575) (3.700,2.616) (3.748,2.650) (3.800,2.687) (3.850,2.722) (3.900,2.758) (3.950,2.793) (3.995,2.825) (4.050,2.864) (4.100,2.899) (4.150,2.935) (4.200,2.970) (4.243,3.000) (4.300,3.041) (4.349,3.075) (4.400,3.111) (4.450,3.147) (4.500,3.182) (4.550,3.217) (4.596,3.250) (4.650,3.288) (4.700,3.323) (4.750,3.359) (4.800,3.394) (4.844,3.425) (4.900,3.465) (4.950,3.500) (5.000,3.536) (5.050,3.571) (5.091,3.600) (5.150,3.642) (5.197,3.675) (5.250,3.712) (5.300,3.748) (5.350,3.783) (5.400,3.818) (5.445,3.850) (5.500,3.889) (5.550,3.924) (5.600,3.960) (5.650,3.995) (5.692,4.025) (5.750,4.066) (5.798,4.100) (5.850,4.137) (5.900,4.172) (5.950,4.207) (6.000,4.243) (6.046,4.275) (6.100,4.313) (6.150,4.349) (6.200,4.384) (6.250,4.419) (6.293,4.450) (6.350,4.490) (6.399,4.525) (6.450,4.561) (6.500,4.596) (6.550,4.632) (6.600,4.667) (6.647,4.700) (6.700,4.738) (6.750,4.773) (6.800,4.808) (6.850,4.844) (6.894,4.875) (6.950,4.914) (7.000,4.950) (7.050,4.985) (7.100,5.020) (7.142,5.050) (7.200,5.091) (7.248,5.125) (7.300,5.162) (7.350,5.197) (7.400,5.233) (7.450,5.268) (7.495,5.300) (7.550,5.339) (7.600,5.374) (7.650,5.409) (7.700,5.445) (7.743,5.475) (7.800,5.515) (7.849,5.550) (7.900,5.586) (7.950,5.622) (8.000,5.657) (8.050,5.692) (8.096,5.725) (8.150,5.763) (8.200,5.798) (8.250,5.834) (8.300,5.869) (8.344,5.900) (8.400,5.940) (8.450,5.975) (8.500,6.010) (8.550,6.046) (8.591,6.075) (8.650,6.117) (8.697,6.150) (8.750,6.187) (8.800,6.223) (8.850,6.258) (8.900,6.293) (8.945,6.325) (9.000,6.364) (9.050,6.399) (9.100,6.435) (9.150,6.470) (9.192,6.500) (9.250,6.541) (9.298,6.575) (9.350,6.611) (9.400,6.647) (9.450,6.682) (9.500,6.718) (9.546,6.750) (9.600,6.788) (9.650,6.824) (9.700,6.859) (9.750,6.894) (9.793,6.925) (9.850,6.965) (9.899,7.000) (9.950,7.036) (10.000,7.071)}
\EQ{(-10.000,7.071) (-9.950,7.036) (-9.900,7.000) (-9.850,6.965) (-9.800,6.930) (-9.758,6.900) (-9.700,6.859) (-9.652,6.825) (-9.600,6.788) (-9.550,6.753) (-9.500,6.718) (-9.450,6.682) (-9.405,6.650) (-9.350,6.611) (-9.300,6.576) (-9.250,6.541) (-9.200,6.505) (-9.157,6.475) (-9.100,6.435) (-9.051,6.400) (-9.000,6.364) (-8.950,6.329) (-8.900,6.293) (-8.850,6.258) (-8.804,6.225) (-8.750,6.187) (-8.700,6.152) (-8.650,6.117) (-8.600,6.081) (-8.556,6.050) (-8.500,6.010) (-8.450,5.975) (-8.400,5.940) (-8.350,5.904) (-8.309,5.875) (-8.250,5.834) (-8.202,5.800) (-8.150,5.763) (-8.100,5.728) (-8.050,5.692) (-8.000,5.657) (-7.955,5.625) (-7.900,5.586) (-7.850,5.551) (-7.800,5.515) (-7.750,5.480) (-7.708,5.450) (-7.650,5.409) (-7.601,5.375) (-7.550,5.339) (-7.500,5.303) (-7.450,5.268) (-7.400,5.233) (-7.354,5.200) (-7.300,5.162) (-7.250,5.127) (-7.200,5.091) (-7.150,5.056) (-7.106,5.025) (-7.050,4.985) (-7.000,4.950) (-6.950,4.914) (-6.900,4.879) (-6.859,4.850) (-6.800,4.808) (-6.753,4.775) (-6.700,4.738) (-6.650,4.702) (-6.600,4.667) (-6.550,4.632) (-6.505,4.600) (-6.450,4.561) (-6.400,4.525) (-6.350,4.490) (-6.300,4.455) (-6.258,4.425) (-6.200,4.384) (-6.152,4.350) (-6.100,4.313) (-6.050,4.278) (-6.000,4.243) (-5.950,4.207) (-5.904,4.175) (-5.850,4.137) (-5.800,4.101) (-5.750,4.066) (-5.700,4.031) (-5.657,4.000) (-5.600,3.960) (-5.551,3.925) (-5.500,3.889) (-5.450,3.854) (-5.400,3.818) (-5.350,3.783) (-5.303,3.750) (-5.250,3.712) (-5.200,3.677) (-5.150,3.642) (-5.100,3.606) (-5.056,3.575) (-5.000,3.536) (-4.950,3.500) (-4.900,3.465) (-4.850,3.429) (-4.808,3.400) (-4.750,3.359) (-4.702,3.325) (-4.650,3.288) (-4.600,3.253) (-4.550,3.217) (-4.500,3.182) (-4.455,3.150) (-4.400,3.111) (-4.350,3.076) (-4.300,3.041) (-4.250,3.005) (-4.207,2.975) (-4.150,2.935) (-4.101,2.900) (-4.050,2.864) (-4.000,2.828) (-3.950,2.793) (-3.900,2.758) (-3.854,2.725) (-3.800,2.687) (-3.750,2.652) (-3.700,2.616) (-3.650,2.581) (-3.606,2.550) (-3.550,2.510) (-3.500,2.475) (-3.450,2.440) (-3.400,2.404) (-3.359,2.375) (-3.300,2.334) (-3.253,2.300) (-3.200,2.263) (-3.150,2.227) (-3.100,2.192) (-3.050,2.157) (-3.005,2.125) (-2.950,2.086) (-2.900,2.051) (-2.850,2.015) (-2.800,1.980) (-2.758,1.950) (-2.700,1.909) (-2.652,1.875) (-2.600,1.839) (-2.550,1.803) (-2.500,1.768) (-2.450,1.732) (-2.404,1.700) (-2.350,1.662) (-2.300,1.626) (-2.250,1.591) (-2.200,1.556) (-2.157,1.525) (-2.100,1.485) (-2.051,1.450) (-2.000,1.414) (-1.950,1.379) (-1.900,1.344) (-1.850,1.308) (-1.803,1.275) (-1.750,1.238) (-1.700,1.202) (-1.650,1.167) (-1.600,1.131) (-1.556,1.100) (-1.500,1.061) (-1.450,1.025) (-1.400,0.990) (-1.350,0.955) (-1.308,0.925) (-1.250,0.884) (-1.202,0.850) (-1.150,0.813) (-1.100,0.778) (-1.050,0.743) (-1.000,0.707) (-0.955,0.675) (-0.900,0.637) (-0.850,0.601) (-0.800,0.566) (-0.750,0.531) (-0.708,0.500) (-0.650,0.460) (-0.601,0.425) (-0.550,0.389) (-0.500,0.354) (-0.450,0.319) (-0.400,0.283) (-0.354,0.250) (-0.300,0.213) (-0.250,0.177) (-0.220,0.155)}
\EQ{(4.306,-10.000) (4.350,-9.950) (4.392,-9.900) (4.425,-9.860) (4.473,-9.800) (4.513,-9.750) (4.550,-9.701) (4.588,-9.650) (4.624,-9.600) (4.660,-9.550) (4.695,-9.500) (4.725,-9.455) (4.761,-9.400) (4.793,-9.350) (4.824,-9.300) (4.855,-9.250) (4.884,-9.200) (4.913,-9.150) (4.941,-9.100) (4.969,-9.050) (4.995,-9.000) (5.021,-8.950) (5.046,-8.900) (5.071,-8.850) (5.095,-8.800) (5.118,-8.750) (5.140,-8.700) (5.162,-8.650) (5.183,-8.600) (5.203,-8.550) (5.223,-8.500) (5.242,-8.450) (5.261,-8.400) (5.279,-8.350) (5.296,-8.300) (5.313,-8.250) (5.329,-8.200) (5.344,-8.150) (5.359,-8.100) (5.373,-8.050) (5.387,-8.000) (5.400,-7.951) (5.413,-7.900) (5.425,-7.850) (5.436,-7.800) (5.447,-7.750) (5.457,-7.700) (5.467,-7.650) (5.476,-7.600) (5.484,-7.550) (5.492,-7.500) (5.500,-7.450) (5.507,-7.400) (5.513,-7.350) (5.519,-7.300) (5.524,-7.250) (5.529,-7.200) (5.533,-7.150) (5.537,-7.100) (5.540,-7.050) (5.542,-7.000) (5.544,-6.950) (5.545,-6.900) (5.546,-6.850) (5.547,-6.800) (5.546,-6.750) (5.546,-6.700) (5.544,-6.650) (5.543,-6.600) (5.540,-6.550) (5.537,-6.500) (5.534,-6.450) (5.530,-6.400) (5.525,-6.350) (5.520,-6.300) (5.515,-6.250) (5.509,-6.200) (5.502,-6.150) (5.495,-6.100) (5.487,-6.050) (5.478,-6.000) (5.469,-5.950) (5.460,-5.900) (5.450,-5.850) (5.439,-5.800) (5.428,-5.750) (5.416,-5.700) (5.404,-5.650) (5.391,-5.600) (5.378,-5.550) (5.363,-5.500) (5.349,-5.450) (5.334,-5.400) (5.318,-5.350) (5.301,-5.300) (5.284,-5.250) (5.267,-5.200) (5.249,-5.150) (5.230,-5.100) (5.210,-5.050) (5.190,-5.000) (5.170,-4.950) (5.150,-4.903) (5.127,-4.850) (5.104,-4.800) (5.081,-4.750) (5.057,-4.700) (5.033,-4.650) (5.008,-4.600) (4.982,-4.550) (4.956,-4.500) (4.929,-4.450) (4.901,-4.400) (4.875,-4.354) (4.844,-4.300) (4.814,-4.250) (4.784,-4.200) (4.753,-4.150) (4.725,-4.107) (4.688,-4.050) (4.655,-4.000) (4.625,-3.955) (4.587,-3.900) (4.551,-3.850) (4.516,-3.800) (4.479,-3.750) (4.441,-3.700) (4.403,-3.650) (4.364,-3.600) (4.325,-3.551) (4.284,-3.500) (4.250,-3.458) (4.201,-3.400) (4.159,-3.350) (4.125,-3.311) (4.075,-3.254) (4.026,-3.200) (3.981,-3.150) (3.935,-3.100) (3.900,-3.063) (3.864,-3.025) (3.825,-2.985) (3.775,-2.934) (3.725,-2.883) (3.675,-2.834) (3.625,-2.785) (3.575,-2.736) (3.537,-2.700) (3.500,-2.665) (3.457,-2.625) (3.403,-2.575) (3.350,-2.526) (3.300,-2.481) (3.250,-2.437) (3.209,-2.400) (3.151,-2.350) (3.100,-2.306) (3.050,-2.263) (3.006,-2.225) (2.950,-2.178) (2.900,-2.137) (2.856,-2.100) (2.800,-2.054) (2.750,-2.014) (2.702,-1.975) (2.650,-1.933) (2.609,-1.900) (2.550,-1.854) (2.500,-1.814) (2.450,-1.775) (2.400,-1.737) (2.353,-1.700) (2.300,-1.660) (2.255,-1.625) (2.200,-1.583) (2.156,-1.550) (2.100,-1.508) (2.056,-1.475) (2.000,-1.433) (1.956,-1.400) (1.900,-1.359) (1.854,-1.325) (1.800,-1.285) (1.752,-1.250) (1.700,-1.212) (1.650,-1.175) (1.600,-1.139) (1.550,-1.103) (1.500,-1.067) (1.450,-1.031) (1.408,-1.000) (1.350,-0.959) (1.304,-0.925) (1.250,-0.887) (1.200,-0.851) (1.150,-0.815) (1.100,-0.780) (1.059,-0.750) (1.000,-0.708) (0.953,-0.675) (0.900,-0.637) (0.850,-0.602) (0.800,-0.566) (0.750,-0.531) (0.707,-0.500) (0.650,-0.460) (0.601,-0.425) (0.550,-0.389) (0.500,-0.354) (0.450,-0.319) (0.400,-0.283) (0.354,-0.250) (0.300,-0.213) (0.250,-0.177) (0.206,-0.144)}
\EQ{(-0.220,-0.155) (-0.275,-0.195) (-0.319,-0.225) (-0.375,-0.266) (-0.424,-0.300) (-0.475,-0.336) (-0.525,-0.372) (-0.566,-0.400) (-0.625,-0.442) (-0.672,-0.475) (-0.725,-0.513) (-0.775,-0.549) (-0.825,-0.584) (-0.875,-0.620) (-0.918,-0.650) (-0.975,-0.691) (-1.024,-0.725) (-1.075,-0.762) (-1.125,-0.797) (-1.175,-0.833) (-1.225,-0.869) (-1.269,-0.900) (-1.325,-0.941) (-1.373,-0.975) (-1.425,-1.013) (-1.475,-1.049) (-1.525,-1.085) (-1.575,-1.121) (-1.625,-1.157) (-1.675,-1.194) (-1.718,-1.225) (-1.775,-1.267) (-1.820,-1.300) (-1.875,-1.340) (-1.922,-1.375) (-1.975,-1.415) (-2.023,-1.450) (-2.075,-1.489) (-2.123,-1.525) (-2.175,-1.564) (-2.222,-1.600) (-2.275,-1.640) (-2.320,-1.675) (-2.375,-1.717) (-2.418,-1.750) (-2.475,-1.795) (-2.525,-1.834) (-2.575,-1.873) (-2.625,-1.913) (-2.671,-1.950) (-2.725,-1.993) (-2.764,-2.025) (-2.825,-2.075) (-2.875,-2.116) (-2.916,-2.150) (-2.975,-2.199) (-3.025,-2.241) (-3.065,-2.275) (-3.123,-2.325) (-3.175,-2.371) (-3.225,-2.415) (-3.265,-2.450) (-3.321,-2.500) (-3.375,-2.549) (-3.425,-2.595) (-3.475,-2.642) (-3.525,-2.689) (-3.563,-2.725) (-3.600,-2.760) (-3.641,-2.800) (-3.692,-2.850) (-3.742,-2.900) (-3.791,-2.950) (-3.840,-3.000) (-3.875,-3.037) (-3.911,-3.075) (-3.950,-3.117) (-4.000,-3.171) (-4.049,-3.225) (-4.093,-3.275) (-4.137,-3.325) (-4.175,-3.369) (-4.222,-3.425) (-4.264,-3.475) (-4.300,-3.519) (-4.345,-3.575) (-4.384,-3.625) (-4.422,-3.675) (-4.460,-3.725) (-4.497,-3.775) (-4.534,-3.825) (-4.569,-3.875) (-4.600,-3.919) (-4.638,-3.975) (-4.672,-4.025) (-4.700,-4.068) (-4.737,-4.125) (-4.768,-4.175) (-4.799,-4.225) (-4.825,-4.269) (-4.858,-4.325) (-4.887,-4.375) (-4.915,-4.425) (-4.942,-4.475) (-4.969,-4.525) (-4.995,-4.575) (-5.020,-4.625) (-5.045,-4.675) (-5.069,-4.725) (-5.093,-4.775) (-5.116,-4.825) (-5.138,-4.875) (-5.159,-4.925) (-5.180,-4.975) (-5.200,-5.024) (-5.220,-5.075) (-5.239,-5.125) (-5.258,-5.175) (-5.275,-5.223) (-5.293,-5.275) (-5.310,-5.325) (-5.325,-5.373) (-5.341,-5.425) (-5.356,-5.475) (-5.371,-5.525) (-5.384,-5.575) (-5.398,-5.625) (-5.410,-5.675) (-5.422,-5.725) (-5.434,-5.775) (-5.445,-5.825) (-5.455,-5.875) (-5.465,-5.925) (-5.474,-5.975) (-5.483,-6.025) (-5.491,-6.075) (-5.498,-6.125) (-5.505,-6.175) (-5.512,-6.225) (-5.518,-6.275) (-5.523,-6.325) (-5.528,-6.375) (-5.532,-6.425) (-5.536,-6.475) (-5.539,-6.525) (-5.542,-6.575) (-5.544,-6.625) (-5.545,-6.675) (-5.546,-6.725) (-5.547,-6.775) (-5.547,-6.825) (-5.546,-6.875) (-5.545,-6.925) (-5.543,-6.975) (-5.541,-7.025) (-5.538,-7.075) (-5.535,-7.125) (-5.531,-7.175) (-5.526,-7.225) (-5.522,-7.275) (-5.516,-7.325) (-5.510,-7.375) (-5.503,-7.425) (-5.496,-7.475) (-5.489,-7.525) (-5.480,-7.575) (-5.471,-7.625) (-5.462,-7.675) (-5.452,-7.725) (-5.442,-7.775) (-5.430,-7.825) (-5.419,-7.875) (-5.407,-7.925) (-5.394,-7.975) (-5.380,-8.025) (-5.366,-8.075) (-5.352,-8.125) (-5.337,-8.175) (-5.321,-8.225) (-5.305,-8.275) (-5.288,-8.325) (-5.270,-8.375) (-5.252,-8.425) (-5.233,-8.475) (-5.213,-8.525) (-5.193,-8.575) (-5.173,-8.625) (-5.151,-8.675) (-5.129,-8.725) (-5.106,-8.775) (-5.083,-8.825) (-5.059,-8.875) (-5.034,-8.925) (-5.008,-8.975) (-4.982,-9.025) (-4.955,-9.075) (-4.927,-9.125) (-4.900,-9.173) (-4.870,-9.225) (-4.840,-9.275) (-4.809,-9.325) (-4.777,-9.375) (-4.750,-9.417) (-4.711,-9.475) (-4.677,-9.525) (-4.642,-9.575) (-4.606,-9.625) (-4.575,-9.668) (-4.532,-9.725) (-4.500,-9.766) (-4.453,-9.825) (-4.413,-9.875) (-4.375,-9.920) (-4.328,-9.975) (-4.306,-10.000)}
\EQ{(0.206,0.144) (0.248,0.175) (0.300,0.213) (0.350,0.248) (0.400,0.283) (0.450,0.319) (0.495,0.350) (0.550,0.389) (0.600,0.424) (0.650,0.460) (0.700,0.495) (0.742,0.525) (0.800,0.566) (0.848,0.600) (0.900,0.637) (0.950,0.673) (1.000,0.708) (1.050,0.744) (1.094,0.775) (1.150,0.815) (1.199,0.850) (1.250,0.887) (1.300,0.923) (1.350,0.959) (1.400,0.995) (1.443,1.025) (1.500,1.067) (1.546,1.100) (1.600,1.139) (1.650,1.175) (1.700,1.212) (1.750,1.248) (1.800,1.285) (1.850,1.322) (1.900,1.359) (1.950,1.396) (2.000,1.433) (2.050,1.470) (2.100,1.508) (2.150,1.546) (2.200,1.583) (2.250,1.621) (2.300,1.660) (2.350,1.698) (2.400,1.737) (2.450,1.775) (2.500,1.814) (2.546,1.850) (2.600,1.893) (2.640,1.925) (2.700,1.973) (2.750,2.014) (2.795,2.050) (2.850,2.095) (2.900,2.137) (2.946,2.175) (3.000,2.220) (3.050,2.263) (3.094,2.300) (3.150,2.349) (3.200,2.393) (3.250,2.437) (3.293,2.475) (3.349,2.525) (3.400,2.572) (3.450,2.618) (3.500,2.665) (3.537,2.700) (3.575,2.736) (3.615,2.775) (3.666,2.825) (3.717,2.875) (3.767,2.925) (3.815,2.975) (3.864,3.025) (3.900,3.063) (3.935,3.100) (3.975,3.144) (4.025,3.198) (4.071,3.250) (4.115,3.300) (4.150,3.340) (4.200,3.399) (4.243,3.450) (4.275,3.489) (4.325,3.550) (4.364,3.600) (4.400,3.646) (4.441,3.700) (4.475,3.745) (4.516,3.800) (4.550,3.848) (4.587,3.900) (4.621,3.950) (4.650,3.992) (4.688,4.050) (4.721,4.100) (4.750,4.146) (4.784,4.200) (4.814,4.250) (4.844,4.300) (4.873,4.350) (4.900,4.398) (4.925,4.443) (4.950,4.489) (4.975,4.536) (5.000,4.585) (5.025,4.634) (5.050,4.685) (5.075,4.737) (5.100,4.791) (5.125,4.846) (5.149,4.900) (5.170,4.950) (5.190,5.000) (5.210,5.050) (5.230,5.100) (5.249,5.150) (5.267,5.200) (5.284,5.250) (5.301,5.300) (5.318,5.350) (5.334,5.400) (5.349,5.450) (5.363,5.500) (5.378,5.550) (5.391,5.600) (5.404,5.650) (5.416,5.700) (5.428,5.750) (5.439,5.800) (5.450,5.850) (5.460,5.900) (5.469,5.950) (5.478,6.000) (5.487,6.050) (5.495,6.100) (5.502,6.150) (5.509,6.200) (5.515,6.250) (5.520,6.300) (5.525,6.350) (5.530,6.400) (5.534,6.450) (5.537,6.500) (5.540,6.550) (5.543,6.600) (5.544,6.650) (5.546,6.700) (5.546,6.750) (5.547,6.800) (5.546,6.850) (5.545,6.900) (5.544,6.950) (5.542,7.000) (5.540,7.050) (5.537,7.100) (5.533,7.150) (5.529,7.200) (5.524,7.250) (5.519,7.300) (5.513,7.350) (5.507,7.400) (5.500,7.450) (5.492,7.500) (5.484,7.550) (5.476,7.600) (5.467,7.650) (5.457,7.700) (5.447,7.750) (5.436,7.800) (5.425,7.850) (5.413,7.900) (5.400,7.950) (5.387,8.000) (5.373,8.050) (5.359,8.100) (5.344,8.150) (5.329,8.200) (5.313,8.250) (5.296,8.300) (5.279,8.350) (5.261,8.400) (5.242,8.450) (5.223,8.500) (5.203,8.550) (5.183,8.600) (5.162,8.650) (5.140,8.700) (5.118,8.750) (5.095,8.800) (5.071,8.850) (5.046,8.900) (5.021,8.950) (4.995,9.000) (4.969,9.050) (4.941,9.100) (4.913,9.150) (4.884,9.200) (4.855,9.250) (4.825,9.299) (4.793,9.350) (4.761,9.400) (4.728,9.450) (4.700,9.492) (4.660,9.550) (4.625,9.599) (4.588,9.650) (4.551,9.700) (4.513,9.750) (4.475,9.798) (4.433,9.850) (4.400,9.890) (4.350,9.949) (4.306,10.000)}
\EQ{(-4.306,10.000) (-4.350,9.950) (-4.392,9.900) (-4.425,9.860) (-4.473,9.800) (-4.513,9.750) (-4.550,9.701) (-4.588,9.650) (-4.624,9.600) (-4.660,9.550) (-4.695,9.500) (-4.725,9.455) (-4.761,9.400) (-4.793,9.350) (-4.824,9.300) (-4.855,9.250) (-4.884,9.200) (-4.913,9.150) (-4.941,9.100) (-4.969,9.050) (-4.995,9.000) (-5.021,8.950) (-5.046,8.900) (-5.071,8.850) (-5.095,8.800) (-5.118,8.750) (-5.140,8.700) (-5.162,8.650) (-5.183,8.600) (-5.203,8.550) (-5.223,8.500) (-5.242,8.450) (-5.261,8.400) (-5.279,8.350) (-5.296,8.300) (-5.313,8.250) (-5.329,8.200) (-5.344,8.150) (-5.359,8.100) (-5.373,8.050) (-5.387,8.000) (-5.400,7.951) (-5.413,7.900) (-5.425,7.850) (-5.436,7.800) (-5.447,7.750) (-5.457,7.700) (-5.467,7.650) (-5.476,7.600) (-5.484,7.550) (-5.492,7.500) (-5.500,7.450) (-5.507,7.400) (-5.513,7.350) (-5.519,7.300) (-5.524,7.250) (-5.529,7.200) (-5.533,7.150) (-5.537,7.100) (-5.540,7.050) (-5.542,7.000) (-5.544,6.950) (-5.545,6.900) (-5.546,6.850) (-5.547,6.800) (-5.546,6.750) (-5.546,6.700) (-5.544,6.650) (-5.543,6.600) (-5.540,6.550) (-5.537,6.500) (-5.534,6.450) (-5.530,6.400) (-5.525,6.350) (-5.520,6.300) (-5.515,6.250) (-5.509,6.200) (-5.502,6.150) (-5.495,6.100) (-5.487,6.050) (-5.478,6.000) (-5.469,5.950) (-5.460,5.900) (-5.450,5.850) (-5.439,5.800) (-5.428,5.750) (-5.416,5.700) (-5.404,5.650) (-5.391,5.600) (-5.378,5.550) (-5.363,5.500) (-5.349,5.450) (-5.334,5.400) (-5.318,5.350) (-5.301,5.300) (-5.284,5.250) (-5.267,5.200) (-5.249,5.150) (-5.230,5.100) (-5.210,5.050) (-5.190,5.000) (-5.170,4.950) (-5.150,4.903) (-5.127,4.850) (-5.104,4.800) (-5.081,4.750) (-5.057,4.700) (-5.033,4.650) (-5.008,4.600) (-4.982,4.550) (-4.956,4.500) (-4.929,4.450) (-4.901,4.400) (-4.875,4.354) (-4.844,4.300) (-4.814,4.250) (-4.784,4.200) (-4.753,4.150) (-4.725,4.107) (-4.688,4.050) (-4.655,4.000) (-4.625,3.955) (-4.587,3.900) (-4.551,3.850) (-4.516,3.800) (-4.479,3.750) (-4.441,3.700) (-4.403,3.650) (-4.364,3.600) (-4.325,3.551) (-4.284,3.500) (-4.250,3.458) (-4.201,3.400) (-4.159,3.350) (-4.125,3.311) (-4.075,3.254) (-4.026,3.200) (-3.981,3.150) (-3.935,3.100) (-3.900,3.063) (-3.864,3.025) (-3.825,2.985) (-3.775,2.934) (-3.725,2.883) (-3.675,2.834) (-3.625,2.785) (-3.575,2.736) (-3.537,2.700) (-3.500,2.665) (-3.457,2.625) (-3.403,2.575) (-3.350,2.526) (-3.300,2.481) (-3.250,2.437) (-3.209,2.400) (-3.151,2.350) (-3.100,2.306) (-3.050,2.263) (-3.006,2.225) (-2.950,2.178) (-2.900,2.137) (-2.856,2.100) (-2.800,2.054) (-2.750,2.014) (-2.702,1.975) (-2.650,1.933) (-2.609,1.900) (-2.550,1.854) (-2.500,1.814) (-2.450,1.775) (-2.400,1.737) (-2.353,1.700) (-2.300,1.660) (-2.255,1.625) (-2.200,1.583) (-2.156,1.550) (-2.100,1.508) (-2.056,1.475) (-2.000,1.433) (-1.956,1.400) (-1.900,1.359) (-1.854,1.325) (-1.800,1.285) (-1.752,1.250) (-1.700,1.212) (-1.650,1.175) (-1.600,1.139) (-1.550,1.103) (-1.500,1.067) (-1.450,1.031) (-1.408,1.000) (-1.350,0.959) (-1.304,0.925) (-1.250,0.887) (-1.200,0.851) (-1.150,0.815) (-1.100,0.780) (-1.059,0.750) (-1.000,0.708) (-0.953,0.675) (-0.900,0.637) (-0.850,0.602) (-0.800,0.566) (-0.750,0.531) (-0.707,0.500) (-0.650,0.460) (-0.601,0.425) (-0.550,0.389) (-0.500,0.354) (-0.450,0.319) (-0.400,0.283) (-0.354,0.250) (-0.300,0.213) (-0.250,0.177) (-0.220,0.155)}
\EQ{(-0.220,-0.155) (-0.275,-0.195) (-0.319,-0.225) (-0.375,-0.266) (-0.424,-0.300) (-0.475,-0.337) (-0.525,-0.372) (-0.575,-0.407) (-0.625,-0.443) (-0.671,-0.475) (-0.725,-0.514) (-0.775,-0.549) (-0.825,-0.585) (-0.875,-0.620) (-0.917,-0.650) (-0.975,-0.692) (-1.022,-0.725) (-1.075,-0.763) (-1.125,-0.799) (-1.175,-0.835) (-1.225,-0.872) (-1.275,-0.908) (-1.325,-0.944) (-1.368,-0.975) (-1.425,-1.017) (-1.470,-1.050) (-1.525,-1.091) (-1.571,-1.125) (-1.625,-1.166) (-1.671,-1.200) (-1.725,-1.241) (-1.770,-1.275) (-1.825,-1.317) (-1.869,-1.350) (-1.925,-1.394) (-1.965,-1.425) (-2.025,-1.472) (-2.075,-1.512) (-2.123,-1.550) (-2.175,-1.592) (-2.216,-1.625) (-2.275,-1.674) (-2.325,-1.715) (-2.367,-1.750) (-2.425,-1.799) (-2.475,-1.843) (-2.525,-1.886) (-2.570,-1.925) (-2.625,-1.975) (-2.675,-2.020) (-2.725,-2.066) (-2.762,-2.100) (-2.800,-2.136) (-2.841,-2.175) (-2.893,-2.225) (-2.943,-2.275) (-2.993,-2.325) (-3.042,-2.375) (-3.090,-2.425) (-3.125,-2.463) (-3.160,-2.500) (-3.200,-2.545) (-3.249,-2.600) (-3.293,-2.650) (-3.336,-2.700) (-3.375,-2.747) (-3.418,-2.800) (-3.450,-2.840) (-3.497,-2.900) (-3.535,-2.950) (-3.572,-3.000) (-3.608,-3.050) (-3.644,-3.100) (-3.675,-3.146) (-3.712,-3.200) (-3.744,-3.250) (-3.775,-3.299) (-3.807,-3.350) (-3.837,-3.400) (-3.866,-3.450) (-3.894,-3.500) (-3.922,-3.550) (-3.948,-3.600) (-3.974,-3.650) (-3.999,-3.700) (-4.023,-3.750) (-4.047,-3.800) (-4.069,-3.850) (-4.091,-3.900) (-4.112,-3.950) (-4.132,-4.000) (-4.152,-4.050) (-4.170,-4.100) (-4.188,-4.150) (-4.206,-4.200) (-4.222,-4.250) (-4.238,-4.300) (-4.253,-4.350) (-4.267,-4.400) (-4.280,-4.450) (-4.293,-4.500) (-4.305,-4.550) (-4.316,-4.600) (-4.327,-4.650) (-4.337,-4.700) (-4.346,-4.750) (-4.354,-4.800) (-4.362,-4.850) (-4.369,-4.900) (-4.376,-4.950) (-4.381,-5.000) (-4.386,-5.050) (-4.391,-5.100) (-4.394,-5.150) (-4.397,-5.200) (-4.400,-5.250) (-4.401,-5.300) (-4.402,-5.350) (-4.402,-5.400) (-4.402,-5.450) (-4.401,-5.500) (-4.399,-5.550) (-4.396,-5.600) (-4.393,-5.650) (-4.389,-5.700) (-4.385,-5.750) (-4.380,-5.800) (-4.374,-5.850) (-4.367,-5.900) (-4.360,-5.950) (-4.352,-6.000) (-4.343,-6.050) (-4.334,-6.100) (-4.323,-6.150) (-4.313,-6.200) (-4.301,-6.250) (-4.289,-6.300) (-4.276,-6.350) (-4.262,-6.400) (-4.247,-6.450) (-4.232,-6.500) (-4.216,-6.550) (-4.200,-6.598) (-4.182,-6.650) (-4.163,-6.700) (-4.144,-6.750) (-4.125,-6.798) (-4.104,-6.850) (-4.082,-6.900) (-4.060,-6.950) (-4.036,-7.000) (-4.012,-7.050) (-3.987,-7.100) (-3.961,-7.150) (-3.934,-7.200) (-3.907,-7.250) (-3.878,-7.300) (-3.850,-7.347) (-3.818,-7.400) (-3.786,-7.450) (-3.753,-7.500) (-3.725,-7.542) (-3.684,-7.600) (-3.650,-7.648) (-3.611,-7.700) (-3.575,-7.747) (-3.533,-7.800) (-3.500,-7.841) (-3.450,-7.900) (-3.406,-7.950) (-3.361,-8.000) (-3.325,-8.039) (-3.275,-8.092) (-3.225,-8.142) (-3.175,-8.191) (-3.125,-8.238) (-3.085,-8.275) (-3.029,-8.325) (-2.975,-8.372) (-2.925,-8.414) (-2.880,-8.450) (-2.825,-8.494) (-2.784,-8.525) (-2.725,-8.569) (-2.682,-8.600) (-2.625,-8.640) (-2.575,-8.674) (-2.525,-8.707) (-2.475,-8.739) (-2.425,-8.770) (-2.375,-8.800) (-2.325,-8.829) (-2.275,-8.858) (-2.225,-8.886) (-2.175,-8.912) (-2.125,-8.938) (-2.075,-8.964) (-2.025,-8.988) (-1.975,-9.012) (-1.925,-9.035) (-1.875,-9.057) (-1.825,-9.078) (-1.775,-9.099) (-1.725,-9.119) (-1.675,-9.139) (-1.625,-9.158) (-1.577,-9.175) (-1.525,-9.193) (-1.475,-9.210) (-1.425,-9.226) (-1.375,-9.242) (-1.325,-9.257) (-1.275,-9.272) (-1.225,-9.285) (-1.175,-9.298) (-1.125,-9.311) (-1.075,-9.323) (-1.025,-9.335) (-0.975,-9.345) (-0.925,-9.356) (-0.875,-9.365) (-0.825,-9.375) (-0.775,-9.383) (-0.725,-9.391) (-0.675,-9.399) (-0.625,-9.406) (-0.575,-9.412) (-0.525,-9.418) (-0.475,-9.423) (-0.425,-9.428) (-0.375,-9.432) (-0.325,-9.436) (-0.275,-9.439) (-0.225,-9.442) (-0.175,-9.444) (-0.125,-9.446) (-0.075,-9.447) (-0.025,-9.447) (0.050,-9.447) (0.100,-9.446) (0.150,-9.445) (0.200,-9.443) (0.250,-9.441) (0.300,-9.438) (0.350,-9.434) (0.400,-9.430) (0.450,-9.426) (0.500,-9.421) (0.550,-9.415) (0.600,-9.409) (0.650,-9.402) (0.700,-9.395) (0.750,-9.387) (0.800,-9.379) (0.850,-9.370) (0.900,-9.361) (0.950,-9.351) (1.000,-9.340) (1.050,-9.329) (1.100,-9.317) (1.150,-9.305) (1.200,-9.292) (1.250,-9.278) (1.300,-9.264) (1.349,-9.250) (1.400,-9.234) (1.450,-9.218) (1.500,-9.202) (1.550,-9.185) (1.600,-9.167) (1.650,-9.148) (1.700,-9.129) (1.750,-9.109) (1.800,-9.089) (1.850,-9.068) (1.900,-9.046) (1.946,-9.025) (2.000,-9.000) (2.050,-8.976) (2.100,-8.951) (2.150,-8.925) (2.198,-8.900) (2.244,-8.875) (2.289,-8.850) (2.333,-8.825) (2.400,-8.785) (2.450,-8.755) (2.497,-8.725) (2.550,-8.691) (2.600,-8.657) (2.647,-8.625) (2.700,-8.587) (2.750,-8.551) (2.800,-8.513) (2.849,-8.475) (2.900,-8.434) (2.942,-8.400) (3.000,-8.351) (3.050,-8.307) (3.100,-8.262) (3.139,-8.225) (3.192,-8.175) (3.242,-8.125) (3.291,-8.075) (3.338,-8.025) (3.375,-7.985) (3.425,-7.929) (3.471,-7.875) (3.513,-7.825) (3.550,-7.779) (3.592,-7.725) (3.625,-7.682) (3.667,-7.625) (3.700,-7.578) (3.736,-7.525) (3.770,-7.475) (3.800,-7.428) (3.833,-7.375) (3.863,-7.325) (3.892,-7.275) (3.921,-7.225) (3.948,-7.175) (3.974,-7.125) (4.000,-7.075) (4.024,-7.025) (4.048,-6.975) (4.071,-6.925) (4.093,-6.875) (4.114,-6.825) (4.134,-6.775) (4.154,-6.725) (4.173,-6.675) (4.191,-6.625) (4.208,-6.575) (4.224,-6.525) (4.240,-6.475) (4.255,-6.425) (4.269,-6.375) (4.282,-6.325) (4.295,-6.275) (4.307,-6.225) (4.318,-6.175) (4.329,-6.125) (4.338,-6.075) (4.347,-6.025) (4.356,-5.975) (4.364,-5.925) (4.370,-5.875) (4.377,-5.825) (4.382,-5.775) (4.387,-5.725) (4.391,-5.675) (4.395,-5.625) (4.398,-5.575) (4.400,-5.525) (4.401,-5.475) (4.402,-5.425) (4.402,-5.375) (4.402,-5.325) (4.401,-5.275) (4.399,-5.225) (4.396,-5.175) (4.393,-5.125) (4.389,-5.075) (4.384,-5.025) (4.379,-4.975) (4.373,-4.925) (4.366,-4.875) (4.358,-4.825) (4.350,-4.775) (4.342,-4.725) (4.332,-4.675) (4.322,-4.625) (4.311,-4.575) (4.299,-4.525) (4.287,-4.475) (4.274,-4.425) (4.260,-4.375) (4.245,-4.325) (4.230,-4.275) (4.214,-4.225) (4.197,-4.175) (4.180,-4.125) (4.161,-4.075) (4.142,-4.025) (4.122,-3.975) (4.102,-3.925) (4.080,-3.875) (4.058,-3.825) (4.035,-3.775) (4.011,-3.725) (3.987,-3.675) (3.961,-3.625) (3.935,-3.575) (3.908,-3.525) (3.880,-3.475) (3.852,-3.425) (3.825,-3.380) (3.792,-3.325) (3.760,-3.275) (3.728,-3.225) (3.700,-3.183) (3.661,-3.125) (3.626,-3.075) (3.590,-3.025) (3.554,-2.975) (3.516,-2.925) (3.478,-2.875) (3.438,-2.825) (3.400,-2.778) (3.357,-2.725) (3.315,-2.675) (3.275,-2.629) (3.227,-2.575) (3.182,-2.525) (3.136,-2.475) (3.100,-2.436) (3.050,-2.384) (3.000,-2.332) (2.950,-2.282) (2.900,-2.232) (2.850,-2.184) (2.800,-2.136) (2.762,-2.100) (2.725,-2.066) (2.681,-2.025) (2.626,-1.975) (2.575,-1.930) (2.525,-1.886) (2.484,-1.850) (2.426,-1.800) (2.375,-1.757) (2.325,-1.715) (2.277,-1.675) (2.225,-1.633) (2.185,-1.600) (2.125,-1.552) (2.075,-1.512) (2.029,-1.475) (1.975,-1.433) (1.933,-1.400) (1.875,-1.355) (1.825,-1.317) (1.775,-1.279) (1.725,-1.241) (1.675,-1.203) (1.625,-1.166) (1.575,-1.128) (1.525,-1.091) (1.475,-1.054) (1.425,-1.017) (1.375,-0.981) (1.333,-0.950) (1.275,-0.908) (1.230,-0.875) (1.175,-0.835) (1.126,-0.800) (1.075,-0.763) (1.025,-0.727) (0.975,-0.692) (0.925,-0.656) (0.882,-0.625) (0.825,-0.585) (0.777,-0.550) (0.725,-0.514) (0.675,-0.478) (0.625,-0.443) (0.575,-0.407) (0.531,-0.375) (0.475,-0.337) (0.425,-0.301) (0.375,-0.266) (0.325,-0.230) (0.284,-0.200) (0.225,-0.160) (0.206,-0.144)}
\EQ{(0.206,0.144) (0.248,0.175) (0.300,0.213) (0.350,0.248) (0.400,0.283) (0.450,0.319) (0.495,0.350) (0.550,0.390) (0.600,0.425) (0.650,0.460) (0.700,0.496) (0.742,0.525) (0.800,0.567) (0.847,0.600) (0.900,0.638) (0.950,0.674) (1.000,0.710) (1.050,0.745) (1.092,0.775) (1.150,0.817) (1.196,0.850) (1.250,0.890) (1.299,0.925) (1.350,0.963) (1.400,0.999) (1.450,1.036) (1.500,1.073) (1.550,1.110) (1.600,1.147) (1.650,1.184) (1.700,1.222) (1.750,1.260) (1.800,1.298) (1.850,1.336) (1.900,1.374) (1.950,1.413) (1.997,1.450) (2.050,1.492) (2.092,1.525) (2.150,1.572) (2.200,1.612) (2.246,1.650) (2.300,1.694) (2.350,1.736) (2.396,1.775) (2.450,1.821) (2.500,1.864) (2.541,1.900) (2.598,1.950) (2.650,1.997) (2.700,2.043) (2.750,2.089) (2.789,2.125) (2.825,2.160) (2.867,2.200) (2.918,2.250) (2.968,2.300) (3.017,2.350) (3.066,2.400) (3.113,2.450) (3.150,2.490) (3.200,2.545) (3.249,2.600) (3.293,2.650) (3.336,2.700) (3.375,2.747) (3.418,2.800) (3.450,2.840) (3.497,2.900) (3.535,2.950) (3.572,3.000) (3.608,3.050) (3.644,3.100) (3.675,3.146) (3.712,3.200) (3.744,3.250) (3.775,3.299) (3.807,3.350) (3.837,3.400) (3.866,3.450) (3.894,3.500) (3.922,3.550) (3.948,3.600) (3.974,3.650) (3.999,3.700) (4.023,3.750) (4.047,3.800) (4.069,3.850) (4.091,3.900) (4.112,3.950) (4.132,4.000) (4.152,4.050) (4.170,4.100) (4.188,4.150) (4.206,4.200) (4.222,4.250) (4.238,4.300) (4.253,4.350) (4.267,4.400) (4.280,4.450) (4.293,4.500) (4.305,4.550) (4.316,4.600) (4.327,4.650) (4.337,4.700) (4.346,4.750) (4.354,4.800) (4.362,4.850) (4.369,4.900) (4.376,4.950) (4.381,5.000) (4.386,5.050) (4.391,5.100) (4.394,5.150) (4.397,5.200) (4.400,5.250) (4.401,5.300) (4.402,5.350) (4.402,5.400) (4.402,5.450) (4.401,5.500) (4.399,5.550) (4.396,5.600) (4.393,5.650) (4.389,5.700) (4.385,5.750) (4.380,5.800) (4.374,5.850) (4.367,5.900) (4.360,5.950) (4.352,6.000) (4.343,6.050) (4.334,6.100) (4.323,6.150) (4.313,6.200) (4.301,6.250) (4.289,6.300) (4.276,6.350) (4.262,6.400) (4.247,6.450) (4.232,6.500) (4.216,6.550) (4.200,6.598) (4.182,6.650) (4.163,6.700) (4.144,6.750) (4.125,6.798) (4.104,6.850) (4.082,6.900) (4.060,6.950) (4.036,7.000) (4.012,7.050) (3.987,7.100) (3.961,7.150) (3.934,7.200) (3.907,7.250) (3.878,7.300) (3.850,7.347) (3.818,7.400) (3.786,7.450) (3.753,7.500) (3.725,7.542) (3.684,7.600) (3.650,7.648) (3.611,7.700) (3.575,7.747) (3.533,7.800) (3.500,7.841) (3.450,7.900) (3.406,7.950) (3.361,8.000) (3.325,8.039) (3.275,8.092) (3.225,8.142) (3.175,8.191) (3.125,8.238) (3.085,8.275) (3.029,8.325) (2.975,8.372) (2.925,8.414) (2.880,8.450) (2.825,8.494) (2.784,8.525) (2.725,8.569) (2.682,8.600) (2.625,8.640) (2.575,8.674) (2.525,8.707) (2.475,8.739) (2.425,8.770) (2.375,8.800) (2.325,8.829) (2.275,8.858) (2.225,8.886) (2.175,8.912) (2.125,8.938) (2.075,8.964) (2.025,8.988) (1.975,9.012) (1.925,9.035) (1.875,9.057) (1.825,9.078) (1.775,9.099) (1.725,9.119) (1.675,9.139) (1.625,9.158) (1.577,9.175) (1.525,9.193) (1.475,9.210) (1.425,9.226) (1.375,9.242) (1.325,9.257) (1.275,9.272) (1.225,9.285) (1.175,9.298) (1.125,9.311) (1.075,9.323) (1.025,9.335) (0.975,9.345) (0.925,9.356) (0.875,9.365) (0.825,9.375) (0.775,9.383) (0.725,9.391) (0.675,9.399) (0.625,9.406) (0.575,9.412) (0.525,9.418) (0.475,9.423) (0.425,9.428) (0.375,9.432) (0.325,9.436) (0.275,9.439) (0.225,9.442) (0.175,9.444) (0.125,9.446) (0.075,9.447) (0.025,9.447) (-0.050,9.447) (-0.100,9.446) (-0.150,9.445) (-0.200,9.443) (-0.250,9.441) (-0.300,9.438) (-0.350,9.434) (-0.400,9.430) (-0.450,9.426) (-0.500,9.421) (-0.550,9.415) (-0.600,9.409) (-0.650,9.402) (-0.700,9.395) (-0.750,9.387) (-0.800,9.379) (-0.850,9.370) (-0.900,9.361) (-0.950,9.351) (-1.000,9.340) (-1.050,9.329) (-1.100,9.317) (-1.150,9.305) (-1.200,9.292) (-1.250,9.278) (-1.300,9.264) (-1.349,9.250) (-1.400,9.234) (-1.450,9.218) (-1.500,9.202) (-1.550,9.185) (-1.600,9.167) (-1.650,9.148) (-1.700,9.129) (-1.750,9.109) (-1.800,9.089) (-1.850,9.068) (-1.900,9.046) (-1.946,9.025) (-2.000,9.000) (-2.050,8.976) (-2.100,8.951) (-2.150,8.925) (-2.198,8.900) (-2.244,8.875) (-2.289,8.850) (-2.333,8.825) (-2.400,8.785) (-2.450,8.755) (-2.497,8.725) (-2.550,8.691) (-2.600,8.657) (-2.647,8.625) (-2.700,8.587) (-2.750,8.551) (-2.800,8.513) (-2.849,8.475) (-2.900,8.434) (-2.942,8.400) (-3.000,8.351) (-3.050,8.307) (-3.100,8.262) (-3.139,8.225) (-3.192,8.175) (-3.242,8.125) (-3.291,8.075) (-3.338,8.025) (-3.375,7.985) (-3.425,7.929) (-3.471,7.875) (-3.513,7.825) (-3.550,7.779) (-3.592,7.725) (-3.625,7.682) (-3.667,7.625) (-3.700,7.578) (-3.736,7.525) (-3.770,7.475) (-3.800,7.428) (-3.833,7.375) (-3.863,7.325) (-3.892,7.275) (-3.921,7.225) (-3.948,7.175) (-3.974,7.125) (-4.000,7.075) (-4.024,7.025) (-4.048,6.975) (-4.071,6.925) (-4.093,6.875) (-4.114,6.825) (-4.134,6.775) (-4.154,6.725) (-4.173,6.675) (-4.191,6.625) (-4.208,6.575) (-4.224,6.525) (-4.240,6.475) (-4.255,6.425) (-4.269,6.375) (-4.282,6.325) (-4.295,6.275) (-4.307,6.225) (-4.318,6.175) (-4.329,6.125) (-4.338,6.075) (-4.347,6.025) (-4.356,5.975) (-4.364,5.925) (-4.370,5.875) (-4.377,5.825) (-4.382,5.775) (-4.387,5.725) (-4.391,5.675) (-4.395,5.625) (-4.398,5.575) (-4.400,5.525) (-4.401,5.475) (-4.402,5.425) (-4.402,5.375) (-4.402,5.325) (-4.401,5.275) (-4.399,5.225) (-4.396,5.175) (-4.393,5.125) (-4.389,5.075) (-4.384,5.025) (-4.379,4.975) (-4.373,4.925) (-4.366,4.875) (-4.358,4.825) (-4.350,4.775) (-4.342,4.725) (-4.332,4.675) (-4.322,4.625) (-4.311,4.575) (-4.299,4.525) (-4.287,4.475) (-4.274,4.425) (-4.260,4.375) (-4.245,4.325) (-4.230,4.275) (-4.214,4.225) (-4.197,4.175) (-4.180,4.125) (-4.161,4.075) (-4.142,4.025) (-4.122,3.975) (-4.102,3.925) (-4.080,3.875) (-4.058,3.825) (-4.035,3.775) (-4.011,3.725) (-3.987,3.675) (-3.961,3.625) (-3.935,3.575) (-3.908,3.525) (-3.880,3.475) (-3.852,3.425) (-3.825,3.380) (-3.792,3.325) (-3.760,3.275) (-3.728,3.225) (-3.700,3.183) (-3.661,3.125) (-3.626,3.075) (-3.590,3.025) (-3.554,2.975) (-3.516,2.925) (-3.478,2.875) (-3.438,2.825) (-3.400,2.778) (-3.357,2.725) (-3.315,2.675) (-3.275,2.629) (-3.227,2.575) (-3.182,2.525) (-3.136,2.475) (-3.100,2.436) (-3.050,2.384) (-3.000,2.332) (-2.950,2.282) (-2.900,2.232) (-2.850,2.184) (-2.800,2.136) (-2.762,2.100) (-2.725,2.066) (-2.681,2.025) (-2.626,1.975) (-2.575,1.930) (-2.525,1.886) (-2.484,1.850) (-2.426,1.800) (-2.375,1.757) (-2.325,1.715) (-2.277,1.675) (-2.225,1.633) (-2.185,1.600) (-2.125,1.552) (-2.075,1.512) (-2.029,1.475) (-1.975,1.433) (-1.933,1.400) (-1.875,1.355) (-1.825,1.317) (-1.775,1.279) (-1.725,1.241) (-1.675,1.203) (-1.625,1.166) (-1.575,1.128) (-1.525,1.091) (-1.475,1.054) (-1.425,1.017) (-1.375,0.981) (-1.333,0.950) (-1.275,0.908) (-1.230,0.875) (-1.175,0.835) (-1.126,0.800) (-1.075,0.763) (-1.025,0.727) (-0.975,0.692) (-0.925,0.656) (-0.882,0.625) (-0.825,0.585) (-0.777,0.550) (-0.725,0.514) (-0.675,0.478) (-0.625,0.443) (-0.575,0.407) (-0.531,0.375) (-0.475,0.337) (-0.425,0.301) (-0.375,0.266) (-0.325,0.230) (-0.284,0.200) (-0.225,0.160) (-0.220,0.155)}
\EQ{(-0.220,-0.155) (-0.275,-0.195) (-0.319,-0.225) (-0.375,-0.266) (-0.424,-0.300) (-0.475,-0.337) (-0.525,-0.372) (-0.575,-0.407) (-0.625,-0.443) (-0.671,-0.475) (-0.725,-0.514) (-0.775,-0.549) (-0.825,-0.585) (-0.875,-0.621) (-0.916,-0.650) (-0.975,-0.693) (-1.020,-0.725) (-1.075,-0.765) (-1.124,-0.800) (-1.175,-0.838) (-1.225,-0.874) (-1.275,-0.911) (-1.325,-0.948) (-1.375,-0.985) (-1.425,-1.022) (-1.475,-1.060) (-1.525,-1.098) (-1.575,-1.136) (-1.625,-1.174) (-1.675,-1.213) (-1.723,-1.250) (-1.775,-1.291) (-1.818,-1.325) (-1.875,-1.370) (-1.925,-1.411) (-1.973,-1.450) (-2.025,-1.493) (-2.075,-1.535) (-2.122,-1.575) (-2.175,-1.621) (-2.225,-1.664) (-2.265,-1.700) (-2.321,-1.750) (-2.375,-1.800) (-2.425,-1.846) (-2.475,-1.894) (-2.525,-1.942) (-2.575,-1.991) (-2.625,-2.042) (-2.675,-2.093) (-2.725,-2.146) (-2.775,-2.200) (-2.820,-2.250) (-2.864,-2.300) (-2.900,-2.342) (-2.949,-2.400) (-2.989,-2.450) (-3.025,-2.495) (-3.068,-2.550) (-3.100,-2.593) (-3.142,-2.650) (-3.175,-2.697) (-3.212,-2.750) (-3.245,-2.800) (-3.275,-2.846) (-3.309,-2.900) (-3.340,-2.950) (-3.369,-3.000) (-3.398,-3.050) (-3.425,-3.099) (-3.450,-3.146) (-3.475,-3.195) (-3.500,-3.245) (-3.525,-3.298) (-3.549,-3.350) (-3.571,-3.400) (-3.592,-3.450) (-3.612,-3.500) (-3.631,-3.550) (-3.650,-3.600) (-3.667,-3.650) (-3.684,-3.700) (-3.700,-3.750) (-3.715,-3.800) (-3.729,-3.850) (-3.742,-3.900) (-3.755,-3.950) (-3.767,-4.000) (-3.777,-4.050) (-3.787,-4.100) (-3.797,-4.150) (-3.805,-4.200) (-3.813,-4.250) (-3.820,-4.300) (-3.826,-4.350) (-3.831,-4.400) (-3.835,-4.450) (-3.839,-4.500) (-3.842,-4.550) (-3.844,-4.600) (-3.845,-4.650) (-3.846,-4.700) (-3.846,-4.750) (-3.845,-4.800) (-3.843,-4.850) (-3.840,-4.900) (-3.837,-4.950) (-3.833,-5.000) (-3.828,-5.050) (-3.822,-5.100) (-3.816,-5.150) (-3.808,-5.200) (-3.800,-5.250) (-3.791,-5.300) (-3.782,-5.350) (-3.771,-5.400) (-3.760,-5.450) (-3.747,-5.500) (-3.734,-5.550) (-3.721,-5.600) (-3.706,-5.650) (-3.690,-5.700) (-3.674,-5.750) (-3.656,-5.800) (-3.638,-5.850) (-3.619,-5.900) (-3.600,-5.948) (-3.578,-6.000) (-3.556,-6.050) (-3.533,-6.100) (-3.509,-6.150) (-3.484,-6.200) (-3.458,-6.250) (-3.431,-6.300) (-3.403,-6.350) (-3.375,-6.399) (-3.344,-6.450) (-3.313,-6.500) (-3.280,-6.550) (-3.250,-6.594) (-3.211,-6.650) (-3.175,-6.699) (-3.137,-6.750) (-3.100,-6.797) (-3.057,-6.850) (-3.015,-6.900) (-2.975,-6.946) (-2.926,-7.000) (-2.879,-7.050) (-2.830,-7.100) (-2.779,-7.150) (-2.727,-7.200) (-2.675,-7.247) (-2.625,-7.291) (-2.585,-7.325) (-2.525,-7.375) (-2.475,-7.414) (-2.429,-7.450) (-2.375,-7.490) (-2.326,-7.525) (-2.275,-7.560) (-2.225,-7.594) (-2.176,-7.625) (-2.125,-7.657) (-2.075,-7.688) (-2.025,-7.717) (-1.975,-7.745) (-1.925,-7.772) (-1.875,-7.799) (-1.825,-7.824) (-1.775,-7.849) (-1.725,-7.873) (-1.675,-7.896) (-1.625,-7.918) (-1.575,-7.939) (-1.525,-7.959) (-1.475,-7.979) (-1.425,-7.998) (-1.375,-8.016) (-1.325,-8.034) (-1.276,-8.050) (-1.225,-8.066) (-1.175,-8.082) (-1.125,-8.096) (-1.075,-8.110) (-1.025,-8.123) (-0.975,-8.136) (-0.925,-8.148) (-0.875,-8.159) (-0.825,-8.170) (-0.775,-8.179) (-0.725,-8.189) (-0.675,-8.197) (-0.625,-8.205) (-0.575,-8.213) (-0.525,-8.219) (-0.475,-8.225) (-0.425,-8.231) (-0.375,-8.236) (-0.325,-8.240) (-0.275,-8.244) (-0.225,-8.247) (-0.175,-8.249) (-0.125,-8.251) (-0.075,-8.252) (-0.025,-8.253) (0.050,-8.253) (0.100,-8.252) (0.150,-8.250) (0.200,-8.248) (0.250,-8.245) (0.300,-8.242) (0.350,-8.238) (0.400,-8.234) (0.450,-8.228) (0.500,-8.223) (0.550,-8.216) (0.600,-8.209) (0.650,-8.201) (0.700,-8.193) (0.750,-8.184) (0.800,-8.175) (0.850,-8.164) (0.900,-8.153) (0.950,-8.142) (1.000,-8.130) (1.050,-8.117) (1.100,-8.103) (1.150,-8.089) (1.200,-8.074) (1.250,-8.058) (1.300,-8.042) (1.350,-8.025) (1.400,-8.007) (1.450,-7.989) (1.500,-7.969) (1.548,-7.950) (1.600,-7.928) (1.650,-7.907) (1.700,-7.884) (1.750,-7.861) (1.800,-7.837) (1.850,-7.812) (1.900,-7.786) (1.950,-7.759) (2.000,-7.731) (2.050,-7.702) (2.096,-7.675) (2.150,-7.642) (2.200,-7.610) (2.250,-7.577) (2.300,-7.543) (2.350,-7.508) (2.395,-7.475) (2.450,-7.434) (2.493,-7.400) (2.550,-7.354) (2.600,-7.313) (2.644,-7.275) (2.700,-7.225) (2.750,-7.178) (2.800,-7.130) (2.850,-7.080) (2.900,-7.028) (2.949,-6.975) (2.993,-6.925) (3.036,-6.875) (3.075,-6.828) (3.117,-6.775) (3.150,-6.733) (3.193,-6.675) (3.225,-6.630) (3.263,-6.575) (3.296,-6.525) (3.325,-6.480) (3.359,-6.425) (3.389,-6.375) (3.418,-6.325) (3.445,-6.275) (3.472,-6.225) (3.497,-6.175) (3.521,-6.125) (3.545,-6.075) (3.567,-6.025) (3.589,-5.975) (3.609,-5.925) (3.629,-5.875) (3.647,-5.825) (3.665,-5.775) (3.682,-5.725) (3.698,-5.675) (3.713,-5.625) (3.728,-5.575) (3.741,-5.525) (3.754,-5.475) (3.765,-5.425) (3.776,-5.375) (3.787,-5.325) (3.796,-5.275) (3.804,-5.225) (3.812,-5.175) (3.819,-5.125) (3.825,-5.075) (3.830,-5.025) (3.835,-4.975) (3.839,-4.925) (3.842,-4.875) (3.844,-4.825) (3.845,-4.775) (3.846,-4.725) (3.846,-4.675) (3.845,-4.625) (3.843,-4.575) (3.841,-4.525) (3.837,-4.475) (3.833,-4.425) (3.828,-4.375) (3.823,-4.325) (3.816,-4.275) (3.809,-4.225) (3.801,-4.175) (3.792,-4.125) (3.783,-4.075) (3.772,-4.025) (3.761,-3.975) (3.749,-3.925) (3.736,-3.875) (3.722,-3.825) (3.708,-3.775) (3.692,-3.725) (3.676,-3.675) (3.659,-3.625) (3.641,-3.575) (3.622,-3.525) (3.602,-3.475) (3.582,-3.425) (3.560,-3.375) (3.538,-3.325) (3.514,-3.275) (3.490,-3.225) (3.465,-3.175) (3.439,-3.125) (3.412,-3.075) (3.384,-3.025) (3.355,-2.975) (3.325,-2.926) (3.294,-2.875) (3.262,-2.825) (3.229,-2.775) (3.200,-2.733) (3.160,-2.675) (3.125,-2.627) (3.087,-2.575) (3.050,-2.527) (3.009,-2.475) (2.975,-2.432) (2.928,-2.375) (2.886,-2.325) (2.850,-2.284) (2.800,-2.228) (2.752,-2.175) (2.705,-2.125) (2.657,-2.075) (2.609,-2.025) (2.559,-1.975) (2.508,-1.925) (2.456,-1.875) (2.403,-1.825) (2.350,-1.777) (2.300,-1.731) (2.250,-1.687) (2.209,-1.650) (2.151,-1.600) (2.100,-1.556) (2.050,-1.514) (2.003,-1.475) (1.950,-1.431) (1.900,-1.391) (1.850,-1.350) (1.800,-1.311) (1.755,-1.275) (1.700,-1.232) (1.659,-1.200) (1.600,-1.155) (1.550,-1.117) (1.500,-1.079) (1.450,-1.041) (1.400,-1.004) (1.350,-0.966) (1.300,-0.929) (1.250,-0.893) (1.200,-0.856) (1.158,-0.825) (1.100,-0.783) (1.055,-0.750) (1.000,-0.711) (0.951,-0.675) (0.900,-0.639) (0.850,-0.603) (0.800,-0.567) (0.750,-0.532) (0.706,-0.500) (0.650,-0.461) (0.600,-0.425) (0.550,-0.390) (0.500,-0.354) (0.450,-0.319) (0.400,-0.284) (0.354,-0.250) (0.300,-0.213) (0.250,-0.177) (0.206,-0.144)}
\EQ{(0.206,0.144) (0.248,0.175) (0.300,0.213) (0.350,0.248) (0.400,0.284) (0.450,0.319) (0.495,0.350) (0.550,0.390) (0.600,0.425) (0.650,0.461) (0.700,0.496) (0.742,0.525) (0.800,0.567) (0.846,0.600) (0.900,0.639) (0.950,0.675) (1.000,0.711) (1.050,0.747) (1.100,0.783) (1.150,0.819) (1.192,0.850) (1.250,0.893) (1.294,0.925) (1.350,0.966) (1.396,1.000) (1.450,1.041) (1.495,1.075) (1.550,1.117) (1.594,1.150) (1.650,1.193) (1.691,1.225) (1.750,1.271) (1.800,1.311) (1.850,1.350) (1.900,1.391) (1.942,1.425) (2.000,1.472) (2.050,1.514) (2.093,1.550) (2.150,1.599) (2.200,1.643) (2.250,1.687) (2.293,1.725) (2.348,1.775) (2.400,1.823) (2.450,1.870) (2.500,1.918) (2.550,1.967) (2.600,2.017) (2.650,2.068) (2.700,2.120) (2.750,2.173) (2.797,2.225) (2.842,2.275) (2.875,2.313) (2.925,2.372) (2.969,2.425) (3.009,2.475) (3.049,2.525) (3.087,2.575) (3.124,2.625) (3.160,2.675) (3.195,2.725) (3.225,2.770) (3.262,2.825) (3.294,2.875) (3.325,2.925) (3.355,2.975) (3.384,3.025) (3.412,3.075) (3.439,3.125) (3.465,3.175) (3.490,3.225) (3.514,3.275) (3.538,3.325) (3.560,3.375) (3.582,3.425) (3.602,3.475) (3.622,3.525) (3.641,3.575) (3.659,3.625) (3.675,3.673) (3.692,3.725) (3.708,3.775) (3.722,3.825) (3.736,3.875) (3.749,3.925) (3.761,3.975) (3.772,4.025) (3.783,4.075) (3.792,4.125) (3.801,4.175) (3.809,4.225) (3.816,4.275) (3.823,4.325) (3.828,4.375) (3.833,4.425) (3.837,4.475) (3.841,4.525) (3.843,4.575) (3.845,4.625) (3.846,4.675) (3.846,4.725) (3.845,4.775) (3.844,4.825) (3.842,4.875) (3.839,4.925) (3.835,4.975) (3.830,5.025) (3.825,5.075) (3.819,5.125) (3.812,5.175) (3.804,5.225) (3.796,5.275) (3.787,5.325) (3.776,5.375) (3.765,5.425) (3.754,5.475) (3.741,5.525) (3.728,5.575) (3.713,5.625) (3.698,5.675) (3.682,5.725) (3.665,5.775) (3.647,5.825) (3.629,5.875) (3.609,5.925) (3.589,5.975) (3.567,6.025) (3.545,6.075) (3.521,6.125) (3.497,6.175) (3.472,6.225) (3.445,6.275) (3.418,6.325) (3.389,6.375) (3.359,6.425) (3.328,6.475) (3.300,6.519) (3.263,6.575) (3.229,6.625) (3.193,6.675) (3.156,6.725) (3.125,6.765) (3.077,6.825) (3.036,6.875) (3.000,6.917) (2.950,6.974) (2.903,7.025) (2.855,7.075) (2.805,7.125) (2.753,7.175) (2.700,7.225) (2.650,7.269) (2.600,7.313) (2.555,7.350) (2.500,7.395) (2.450,7.434) (2.400,7.471) (2.350,7.508) (2.300,7.543) (2.253,7.575) (2.200,7.610) (2.150,7.642) (2.100,7.673) (2.054,7.700) (2.000,7.731) (1.950,7.759) (1.900,7.786) (1.850,7.812) (1.800,7.837) (1.750,7.861) (1.700,7.884) (1.650,7.907) (1.600,7.928) (1.550,7.949) (1.500,7.969) (1.450,7.989) (1.400,8.007) (1.350,8.025) (1.300,8.042) (1.250,8.058) (1.200,8.074) (1.150,8.089) (1.100,8.103) (1.050,8.117) (1.000,8.130) (0.950,8.142) (0.900,8.153) (0.850,8.164) (0.800,8.175) (0.750,8.184) (0.700,8.193) (0.650,8.201) (0.600,8.209) (0.550,8.216) (0.500,8.223) (0.450,8.228) (0.400,8.234) (0.350,8.238) (0.300,8.242) (0.250,8.245) (0.200,8.248) (0.150,8.250) (0.100,8.252) (0.050,8.253) (0.000,8.253) (-0.050,8.253) (-0.100,8.252) (-0.150,8.250) (-0.200,8.248) (-0.250,8.245) (-0.300,8.242) (-0.350,8.238) (-0.400,8.234) (-0.450,8.228) (-0.500,8.223) (-0.550,8.216) (-0.600,8.209) (-0.650,8.201) (-0.700,8.193) (-0.750,8.184) (-0.800,8.175) (-0.850,8.164) (-0.900,8.153) (-0.950,8.142) (-1.000,8.130) (-1.050,8.117) (-1.100,8.103) (-1.150,8.089) (-1.200,8.074) (-1.250,8.058) (-1.300,8.042) (-1.350,8.025) (-1.400,8.007) (-1.450,7.989) (-1.500,7.969) (-1.548,7.950) (-1.600,7.928) (-1.650,7.907) (-1.700,7.884) (-1.750,7.861) (-1.800,7.837) (-1.850,7.812) (-1.900,7.786) (-1.950,7.759) (-2.000,7.731) (-2.050,7.702) (-2.096,7.675) (-2.150,7.642) (-2.200,7.610) (-2.250,7.577) (-2.300,7.543) (-2.350,7.508) (-2.395,7.475) (-2.450,7.434) (-2.493,7.400) (-2.550,7.354) (-2.600,7.313) (-2.644,7.275) (-2.700,7.225) (-2.750,7.178) (-2.800,7.130) (-2.850,7.080) (-2.900,7.028) (-2.949,6.975) (-2.993,6.925) (-3.036,6.875) (-3.075,6.828) (-3.117,6.775) (-3.150,6.733) (-3.193,6.675) (-3.225,6.630) (-3.263,6.575) (-3.296,6.525) (-3.325,6.480) (-3.359,6.425) (-3.389,6.375) (-3.418,6.325) (-3.445,6.275) (-3.472,6.225) (-3.497,6.175) (-3.521,6.125) (-3.545,6.075) (-3.567,6.025) (-3.589,5.975) (-3.609,5.925) (-3.629,5.875) (-3.647,5.825) (-3.665,5.775) (-3.682,5.725) (-3.698,5.675) (-3.713,5.625) (-3.728,5.575) (-3.741,5.525) (-3.754,5.475) (-3.765,5.425) (-3.776,5.375) (-3.787,5.325) (-3.796,5.275) (-3.804,5.225) (-3.812,5.175) (-3.819,5.125) (-3.825,5.075) (-3.830,5.025) (-3.835,4.975) (-3.839,4.925) (-3.842,4.875) (-3.844,4.825) (-3.845,4.775) (-3.846,4.725) (-3.846,4.675) (-3.845,4.625) (-3.843,4.575) (-3.841,4.525) (-3.837,4.475) (-3.833,4.425) (-3.828,4.375) (-3.823,4.325) (-3.816,4.275) (-3.809,4.225) (-3.801,4.175) (-3.792,4.125) (-3.783,4.075) (-3.772,4.025) (-3.761,3.975) (-3.749,3.925) (-3.736,3.875) (-3.722,3.825) (-3.708,3.775) (-3.692,3.725) (-3.676,3.675) (-3.659,3.625) (-3.641,3.575) (-3.622,3.525) (-3.602,3.475) (-3.582,3.425) (-3.560,3.375) (-3.538,3.325) (-3.514,3.275) (-3.490,3.225) (-3.465,3.175) (-3.439,3.125) (-3.412,3.075) (-3.384,3.025) (-3.355,2.975) (-3.325,2.926) (-3.294,2.875) (-3.262,2.825) (-3.229,2.775) (-3.200,2.733) (-3.160,2.675) (-3.125,2.627) (-3.087,2.575) (-3.050,2.527) (-3.009,2.475) (-2.975,2.432) (-2.928,2.375) (-2.886,2.325) (-2.850,2.284) (-2.800,2.228) (-2.752,2.175) (-2.705,2.125) (-2.657,2.075) (-2.609,2.025) (-2.559,1.975) (-2.508,1.925) (-2.456,1.875) (-2.403,1.825) (-2.350,1.777) (-2.300,1.731) (-2.250,1.687) (-2.209,1.650) (-2.151,1.600) (-2.100,1.556) (-2.050,1.514) (-2.003,1.475) (-1.950,1.431) (-1.900,1.391) (-1.850,1.350) (-1.800,1.311) (-1.755,1.275) (-1.700,1.232) (-1.659,1.200) (-1.600,1.155) (-1.550,1.117) (-1.500,1.079) (-1.450,1.041) (-1.400,1.004) (-1.350,0.966) (-1.300,0.929) (-1.250,0.893) (-1.200,0.856) (-1.158,0.825) (-1.100,0.783) (-1.055,0.750) (-1.000,0.711) (-0.951,0.675) (-0.900,0.639) (-0.850,0.603) (-0.800,0.567) (-0.750,0.532) (-0.706,0.500) (-0.650,0.461) (-0.600,0.425) (-0.550,0.390) (-0.500,0.354) (-0.450,0.319) (-0.400,0.284) (-0.354,0.250) (-0.300,0.213) (-0.250,0.177) (-0.220,0.155)}
\EQ{(-0.220,-0.155) (-0.275,-0.195) (-0.319,-0.225) (-0.375,-0.266) (-0.424,-0.300) (-0.475,-0.337) (-0.525,-0.372) (-0.575,-0.407) (-0.625,-0.443) (-0.671,-0.475) (-0.725,-0.514) (-0.775,-0.550) (-0.825,-0.586) (-0.875,-0.622) (-0.925,-0.658) (-0.975,-0.694) (-1.019,-0.725) (-1.075,-0.767) (-1.121,-0.800) (-1.175,-0.840) (-1.223,-0.875) (-1.275,-0.914) (-1.323,-0.950) (-1.375,-0.989) (-1.422,-1.025) (-1.475,-1.066) (-1.520,-1.100) (-1.575,-1.143) (-1.616,-1.175) (-1.675,-1.222) (-1.725,-1.263) (-1.771,-1.300) (-1.825,-1.344) (-1.875,-1.386) (-1.921,-1.425) (-1.975,-1.472) (-2.025,-1.515) (-2.064,-1.550) (-2.120,-1.600) (-2.174,-1.650) (-2.225,-1.698) (-2.275,-1.746) (-2.325,-1.795) (-2.375,-1.845) (-2.425,-1.896) (-2.475,-1.949) (-2.522,-2.000) (-2.567,-2.050) (-2.611,-2.100) (-2.650,-2.146) (-2.695,-2.200) (-2.735,-2.250) (-2.774,-2.300) (-2.812,-2.350) (-2.849,-2.400) (-2.885,-2.450) (-2.919,-2.500) (-2.950,-2.546) (-2.985,-2.600) (-3.016,-2.650) (-3.046,-2.700) (-3.075,-2.750) (-3.100,-2.794) (-3.125,-2.841) (-3.150,-2.889) (-3.175,-2.939) (-3.200,-2.991) (-3.225,-3.045) (-3.249,-3.100) (-3.270,-3.150) (-3.290,-3.200) (-3.309,-3.250) (-3.326,-3.300) (-3.343,-3.350) (-3.359,-3.400) (-3.374,-3.450) (-3.389,-3.500) (-3.402,-3.550) (-3.414,-3.600) (-3.426,-3.650) (-3.436,-3.700) (-3.446,-3.750) (-3.455,-3.800) (-3.462,-3.850) (-3.469,-3.900) (-3.476,-3.950) (-3.481,-4.000) (-3.485,-4.050) (-3.489,-4.100) (-3.491,-4.150) (-3.493,-4.200) (-3.494,-4.250) (-3.494,-4.300) (-3.493,-4.350) (-3.492,-4.400) (-3.489,-4.450) (-3.486,-4.500) (-3.482,-4.550) (-3.477,-4.600) (-3.471,-4.650) (-3.464,-4.700) (-3.456,-4.750) (-3.447,-4.800) (-3.438,-4.850) (-3.428,-4.900) (-3.416,-4.950) (-3.404,-5.000) (-3.391,-5.050) (-3.377,-5.100) (-3.362,-5.150) (-3.346,-5.200) (-3.329,-5.250) (-3.311,-5.300) (-3.292,-5.350) (-3.272,-5.400) (-3.251,-5.450) (-3.229,-5.500) (-3.206,-5.550) (-3.182,-5.600) (-3.157,-5.650) (-3.131,-5.700) (-3.103,-5.750) (-3.075,-5.799) (-3.044,-5.850) (-3.013,-5.900) (-2.981,-5.950) (-2.950,-5.995) (-2.912,-6.050) (-2.875,-6.100) (-2.837,-6.150) (-2.800,-6.196) (-2.755,-6.250) (-2.712,-6.300) (-2.675,-6.342) (-2.625,-6.396) (-2.575,-6.447) (-2.525,-6.496) (-2.475,-6.544) (-2.425,-6.589) (-2.384,-6.625) (-2.325,-6.675) (-2.275,-6.716) (-2.231,-6.750) (-2.175,-6.792) (-2.130,-6.825) (-2.075,-6.863) (-2.025,-6.897) (-1.982,-6.925) (-1.925,-6.961) (-1.875,-6.991) (-1.825,-7.020) (-1.775,-7.048) (-1.725,-7.074) (-1.676,-7.100) (-1.625,-7.125) (-1.575,-7.149) (-1.525,-7.172) (-1.475,-7.194) (-1.425,-7.215) (-1.375,-7.236) (-1.325,-7.255) (-1.275,-7.274) (-1.225,-7.292) (-1.175,-7.309) (-1.125,-7.325) (-1.075,-7.340) (-1.025,-7.355) (-0.975,-7.369) (-0.925,-7.382) (-0.875,-7.395) (-0.825,-7.406) (-0.775,-7.417) (-0.725,-7.427) (-0.675,-7.437) (-0.625,-7.446) (-0.575,-7.454) (-0.525,-7.461) (-0.475,-7.468) (-0.425,-7.474) (-0.375,-7.479) (-0.325,-7.484) (-0.275,-7.488) (-0.225,-7.492) (-0.175,-7.494) (-0.125,-7.496) (-0.075,-7.498) (-0.025,-7.498) (0.050,-7.498) (0.100,-7.497) (0.150,-7.495) (0.200,-7.493) (0.250,-7.490) (0.300,-7.486) (0.350,-7.482) (0.400,-7.477) (0.450,-7.471) (0.500,-7.465) (0.550,-7.458) (0.599,-7.450) (0.650,-7.441) (0.700,-7.432) (0.750,-7.422) (0.800,-7.412) (0.850,-7.400) (0.900,-7.388) (0.950,-7.376) (1.000,-7.362) (1.050,-7.348) (1.100,-7.333) (1.150,-7.317) (1.200,-7.300) (1.250,-7.283) (1.300,-7.265) (1.350,-7.246) (1.400,-7.226) (1.450,-7.205) (1.500,-7.183) (1.550,-7.161) (1.600,-7.137) (1.650,-7.113) (1.700,-7.088) (1.750,-7.061) (1.800,-7.034) (1.850,-7.005) (1.900,-6.976) (1.950,-6.945) (2.000,-6.913) (2.050,-6.880) (2.094,-6.850) (2.150,-6.811) (2.198,-6.775) (2.250,-6.735) (2.295,-6.700) (2.350,-6.654) (2.400,-6.611) (2.441,-6.575) (2.495,-6.525) (2.547,-6.475) (2.597,-6.425) (2.644,-6.375) (2.690,-6.325) (2.725,-6.286) (2.775,-6.227) (2.817,-6.175) (2.850,-6.133) (2.893,-6.075) (2.925,-6.031) (2.964,-5.975) (2.997,-5.925) (3.025,-5.881) (3.060,-5.825) (3.089,-5.775) (3.117,-5.725) (3.144,-5.675) (3.170,-5.625) (3.195,-5.575) (3.218,-5.525) (3.241,-5.475) (3.262,-5.425) (3.282,-5.375) (3.302,-5.325) (3.320,-5.275) (3.338,-5.225) (3.354,-5.175) (3.370,-5.125) (3.384,-5.075) (3.398,-5.025) (3.410,-4.975) (3.422,-4.925) (3.433,-4.875) (3.443,-4.825) (3.452,-4.775) (3.460,-4.725) (3.467,-4.675) (3.474,-4.625) (3.479,-4.575) (3.484,-4.525) (3.488,-4.475) (3.491,-4.425) (3.493,-4.375) (3.494,-4.325) (3.494,-4.275) (3.494,-4.225) (3.492,-4.175) (3.490,-4.125) (3.487,-4.075) (3.483,-4.025) (3.478,-3.975) (3.473,-3.925) (3.466,-3.875) (3.459,-3.825) (3.450,-3.775) (3.441,-3.725) (3.431,-3.675) (3.420,-3.625) (3.408,-3.575) (3.395,-3.525) (3.382,-3.475) (3.367,-3.425) (3.352,-3.375) (3.335,-3.325) (3.318,-3.275) (3.300,-3.227) (3.280,-3.175) (3.260,-3.125) (3.238,-3.075) (3.216,-3.025) (3.193,-2.975) (3.168,-2.925) (3.143,-2.875) (3.117,-2.825) (3.089,-2.775) (3.061,-2.725) (3.031,-2.675) (3.001,-2.625) (2.969,-2.575) (2.936,-2.525) (2.902,-2.475) (2.867,-2.425) (2.831,-2.375) (2.800,-2.334) (2.755,-2.275) (2.715,-2.225) (2.675,-2.176) (2.632,-2.125) (2.589,-2.075) (2.550,-2.031) (2.500,-1.976) (2.452,-1.925) (2.404,-1.875) (2.355,-1.825) (2.305,-1.775) (2.254,-1.725) (2.201,-1.675) (2.150,-1.628) (2.100,-1.582) (2.050,-1.537) (2.008,-1.500) (1.950,-1.450) (1.900,-1.407) (1.850,-1.365) (1.802,-1.325) (1.750,-1.283) (1.710,-1.250) (1.650,-1.202) (1.600,-1.163) (1.552,-1.125) (1.500,-1.085) (1.455,-1.050) (1.400,-1.008) (1.357,-0.975) (1.300,-0.933) (1.257,-0.900) (1.200,-0.858) (1.155,-0.825) (1.100,-0.785) (1.053,-0.750) (1.000,-0.712) (0.950,-0.676) (0.900,-0.640) (0.850,-0.604) (0.800,-0.568) (0.750,-0.532) (0.706,-0.500) (0.650,-0.461) (0.600,-0.425) (0.550,-0.390) (0.500,-0.354) (0.450,-0.319) (0.400,-0.284) (0.354,-0.250) (0.300,-0.213) (0.250,-0.177) (0.206,-0.144)}
\EQ{(0.206,0.144) (0.248,0.175) (0.300,0.213) (0.350,0.248) (0.400,0.284) (0.450,0.319) (0.495,0.350) (0.550,0.390) (0.600,0.425) (0.650,0.461) (0.700,0.496) (0.741,0.525) (0.800,0.568) (0.845,0.600) (0.900,0.640) (0.949,0.675) (1.000,0.712) (1.050,0.748) (1.100,0.785) (1.150,0.821) (1.200,0.858) (1.250,0.895) (1.300,0.933) (1.350,0.970) (1.400,1.008) (1.450,1.046) (1.500,1.085) (1.550,1.124) (1.600,1.163) (1.647,1.200) (1.700,1.242) (1.741,1.275) (1.800,1.324) (1.850,1.365) (1.892,1.400) (1.950,1.450) (2.000,1.493) (2.050,1.537) (2.092,1.575) (2.147,1.625) (2.200,1.674) (2.250,1.722) (2.300,1.770) (2.350,1.820) (2.400,1.871) (2.450,1.923) (2.499,1.975) (2.545,2.025) (2.589,2.075) (2.625,2.117) (2.674,2.175) (2.715,2.225) (2.750,2.269) (2.793,2.325) (2.825,2.367) (2.867,2.425) (2.900,2.472) (2.936,2.525) (2.969,2.575) (3.000,2.624) (3.031,2.675) (3.061,2.725) (3.089,2.775) (3.117,2.825) (3.143,2.875) (3.168,2.925) (3.193,2.975) (3.216,3.025) (3.238,3.075) (3.260,3.125) (3.280,3.175) (3.299,3.225) (3.318,3.275) (3.335,3.325) (3.352,3.375) (3.367,3.425) (3.382,3.475) (3.395,3.525) (3.408,3.575) (3.420,3.625) (3.431,3.675) (3.441,3.725) (3.450,3.775) (3.459,3.825) (3.466,3.875) (3.473,3.925) (3.478,3.975) (3.483,4.025) (3.487,4.075) (3.490,4.125) (3.492,4.175) (3.494,4.225) (3.494,4.275) (3.494,4.325) (3.493,4.375) (3.491,4.425) (3.488,4.475) (3.484,4.525) (3.479,4.575) (3.474,4.625) (3.467,4.675) (3.460,4.725) (3.452,4.775) (3.443,4.825) (3.433,4.875) (3.422,4.925) (3.410,4.975) (3.398,5.025) (3.384,5.075) (3.370,5.125) (3.354,5.175) (3.338,5.225) (3.320,5.275) (3.302,5.325) (3.282,5.375) (3.262,5.425) (3.241,5.475) (3.218,5.525) (3.195,5.575) (3.170,5.625) (3.144,5.675) (3.117,5.725) (3.089,5.775) (3.060,5.825) (3.029,5.875) (3.000,5.921) (2.964,5.975) (2.929,6.025) (2.900,6.066) (2.856,6.125) (2.825,6.165) (2.776,6.225) (2.734,6.275) (2.700,6.314) (2.650,6.369) (2.600,6.422) (2.550,6.472) (2.500,6.520) (2.450,6.567) (2.400,6.611) (2.355,6.650) (2.300,6.696) (2.250,6.735) (2.200,6.774) (2.150,6.811) (2.100,6.846) (2.058,6.875) (2.000,6.913) (1.950,6.945) (1.901,6.975) (1.850,7.005) (1.800,7.034) (1.750,7.061) (1.700,7.088) (1.650,7.113) (1.600,7.137) (1.550,7.161) (1.500,7.183) (1.450,7.205) (1.402,7.225) (1.350,7.246) (1.300,7.265) (1.250,7.283) (1.201,7.300) (1.150,7.317) (1.100,7.333) (1.050,7.348) (1.000,7.362) (0.950,7.376) (0.900,7.388) (0.850,7.400) (0.800,7.412) (0.750,7.422) (0.700,7.432) (0.650,7.441) (0.600,7.450) (0.550,7.458) (0.500,7.465) (0.450,7.471) (0.400,7.477) (0.350,7.482) (0.300,7.486) (0.250,7.490) (0.200,7.493) (0.150,7.495) (0.100,7.497) (0.050,7.498) (0.000,7.498) (-0.050,7.498) (-0.100,7.497) (-0.150,7.495) (-0.200,7.493) (-0.250,7.490) (-0.300,7.486) (-0.350,7.482) (-0.400,7.477) (-0.450,7.471) (-0.500,7.465) (-0.550,7.458) (-0.599,7.450) (-0.650,7.441) (-0.700,7.432) (-0.750,7.422) (-0.800,7.412) (-0.850,7.400) (-0.900,7.388) (-0.950,7.376) (-1.000,7.362) (-1.050,7.348) (-1.100,7.333) (-1.150,7.317) (-1.200,7.300) (-1.250,7.283) (-1.300,7.265) (-1.350,7.246) (-1.400,7.226) (-1.450,7.205) (-1.500,7.183) (-1.550,7.161) (-1.600,7.137) (-1.650,7.113) (-1.700,7.088) (-1.750,7.061) (-1.800,7.034) (-1.850,7.005) (-1.900,6.976) (-1.950,6.945) (-2.000,6.913) (-2.050,6.880) (-2.094,6.850) (-2.150,6.811) (-2.198,6.775) (-2.250,6.735) (-2.295,6.700) (-2.350,6.654) (-2.400,6.611) (-2.441,6.575) (-2.495,6.525) (-2.547,6.475) (-2.597,6.425) (-2.644,6.375) (-2.690,6.325) (-2.725,6.286) (-2.775,6.227) (-2.817,6.175) (-2.850,6.133) (-2.893,6.075) (-2.925,6.031) (-2.964,5.975) (-2.997,5.925) (-3.025,5.881) (-3.060,5.825) (-3.089,5.775) (-3.117,5.725) (-3.144,5.675) (-3.170,5.625) (-3.195,5.575) (-3.218,5.525) (-3.241,5.475) (-3.262,5.425) (-3.282,5.375) (-3.302,5.325) (-3.320,5.275) (-3.338,5.225) (-3.354,5.175) (-3.370,5.125) (-3.384,5.075) (-3.398,5.025) (-3.410,4.975) (-3.422,4.925) (-3.433,4.875) (-3.443,4.825) (-3.452,4.775) (-3.460,4.725) (-3.467,4.675) (-3.474,4.625) (-3.479,4.575) (-3.484,4.525) (-3.488,4.475) (-3.491,4.425) (-3.493,4.375) (-3.494,4.325) (-3.494,4.275) (-3.494,4.225) (-3.492,4.175) (-3.490,4.125) (-3.487,4.075) (-3.483,4.025) (-3.478,3.975) (-3.473,3.925) (-3.466,3.875) (-3.459,3.825) (-3.450,3.775) (-3.441,3.725) (-3.431,3.675) (-3.420,3.625) (-3.408,3.575) (-3.395,3.525) (-3.382,3.475) (-3.367,3.425) (-3.352,3.375) (-3.335,3.325) (-3.318,3.275) (-3.300,3.227) (-3.280,3.175) (-3.260,3.125) (-3.238,3.075) (-3.216,3.025) (-3.193,2.975) (-3.168,2.925) (-3.143,2.875) (-3.117,2.825) (-3.089,2.775) (-3.061,2.725) (-3.031,2.675) (-3.001,2.625) (-2.969,2.575) (-2.936,2.525) (-2.902,2.475) (-2.867,2.425) (-2.831,2.375) (-2.800,2.334) (-2.755,2.275) (-2.715,2.225) (-2.675,2.176) (-2.632,2.125) (-2.589,2.075) (-2.550,2.031) (-2.500,1.976) (-2.452,1.925) (-2.404,1.875) (-2.355,1.825) (-2.305,1.775) (-2.254,1.725) (-2.201,1.675) (-2.150,1.628) (-2.100,1.582) (-2.050,1.537) (-2.008,1.500) (-1.950,1.450) (-1.900,1.407) (-1.850,1.365) (-1.802,1.325) (-1.750,1.283) (-1.710,1.250) (-1.650,1.202) (-1.600,1.163) (-1.552,1.125) (-1.500,1.085) (-1.455,1.050) (-1.400,1.008) (-1.357,0.975) (-1.300,0.933) (-1.257,0.900) (-1.200,0.858) (-1.155,0.825) (-1.100,0.785) (-1.053,0.750) (-1.000,0.712) (-0.950,0.676) (-0.900,0.640) (-0.850,0.604) (-0.800,0.568) (-0.750,0.532) (-0.706,0.500) (-0.650,0.461) (-0.600,0.425) (-0.550,0.390) (-0.500,0.354) (-0.450,0.319) (-0.400,0.284) (-0.354,0.250) (-0.300,0.213) (-0.250,0.177) (-0.220,0.155)}
\EQ{(-0.220,-0.155) (-0.275,-0.195) (-0.319,-0.225) (-0.375,-0.266) (-0.424,-0.300) (-0.475,-0.337) (-0.525,-0.372) (-0.575,-0.408) (-0.625,-0.443) (-0.671,-0.475) (-0.725,-0.515) (-0.775,-0.550) (-0.825,-0.586) (-0.875,-0.622) (-0.925,-0.659) (-0.975,-0.695) (-1.017,-0.725) (-1.075,-0.768) (-1.119,-0.800) (-1.175,-0.842) (-1.219,-0.875) (-1.275,-0.917) (-1.319,-0.950) (-1.375,-0.994) (-1.416,-1.025) (-1.475,-1.071) (-1.525,-1.111) (-1.574,-1.150) (-1.625,-1.191) (-1.667,-1.225) (-1.725,-1.274) (-1.775,-1.316) (-1.815,-1.350) (-1.872,-1.400) (-1.925,-1.447) (-1.975,-1.492) (-2.025,-1.538) (-2.064,-1.575) (-2.117,-1.625) (-2.167,-1.675) (-2.217,-1.725) (-2.265,-1.775) (-2.300,-1.812) (-2.335,-1.850) (-2.375,-1.894) (-2.424,-1.950) (-2.467,-2.000) (-2.500,-2.040) (-2.548,-2.100) (-2.587,-2.150) (-2.624,-2.200) (-2.661,-2.250) (-2.696,-2.300) (-2.725,-2.343) (-2.762,-2.400) (-2.794,-2.450) (-2.824,-2.500) (-2.850,-2.545) (-2.875,-2.589) (-2.900,-2.635) (-2.925,-2.683) (-2.950,-2.733) (-2.975,-2.786) (-3.000,-2.841) (-3.025,-2.899) (-3.046,-2.950) (-3.065,-3.000) (-3.083,-3.050) (-3.100,-3.099) (-3.117,-3.150) (-3.132,-3.200) (-3.146,-3.250) (-3.159,-3.300) (-3.171,-3.350) (-3.183,-3.400) (-3.193,-3.450) (-3.202,-3.500) (-3.211,-3.550) (-3.218,-3.600) (-3.224,-3.650) (-3.230,-3.700) (-3.235,-3.750) (-3.238,-3.800) (-3.241,-3.850) (-3.243,-3.900) (-3.244,-3.950) (-3.244,-4.000) (-3.243,-4.050) (-3.241,-4.100) (-3.238,-4.150) (-3.234,-4.200) (-3.229,-4.250) (-3.224,-4.300) (-3.217,-4.350) (-3.210,-4.400) (-3.201,-4.450) (-3.192,-4.500) (-3.182,-4.550) (-3.170,-4.600) (-3.158,-4.650) (-3.145,-4.700) (-3.130,-4.750) (-3.115,-4.800) (-3.099,-4.850) (-3.081,-4.900) (-3.063,-4.950) (-3.043,-5.000) (-3.022,-5.050) (-3.001,-5.100) (-2.978,-5.150) (-2.954,-5.200) (-2.928,-5.250) (-2.902,-5.300) (-2.875,-5.348) (-2.845,-5.400) (-2.814,-5.450) (-2.782,-5.500) (-2.750,-5.549) (-2.715,-5.600) (-2.678,-5.650) (-2.640,-5.700) (-2.601,-5.750) (-2.560,-5.800) (-2.525,-5.840) (-2.475,-5.896) (-2.425,-5.949) (-2.375,-6.000) (-2.325,-6.049) (-2.275,-6.096) (-2.225,-6.140) (-2.184,-6.175) (-2.125,-6.224) (-2.075,-6.263) (-2.026,-6.300) (-1.975,-6.337) (-1.925,-6.372) (-1.875,-6.406) (-1.825,-6.438) (-1.775,-6.469) (-1.725,-6.499) (-1.679,-6.525) (-1.634,-6.550) (-1.587,-6.575) (-1.538,-6.600) (-1.487,-6.625) (-1.433,-6.650) (-1.377,-6.675) (-1.325,-6.697) (-1.275,-6.718) (-1.225,-6.737) (-1.175,-6.756) (-1.125,-6.773) (-1.075,-6.790) (-1.025,-6.806) (-0.975,-6.821) (-0.925,-6.835) (-0.875,-6.849) (-0.825,-6.861) (-0.775,-6.873) (-0.725,-6.884) (-0.675,-6.895) (-0.625,-6.904) (-0.575,-6.913) (-0.525,-6.921) (-0.475,-6.928) (-0.425,-6.935) (-0.375,-6.941) (-0.325,-6.946) (-0.275,-6.950) (-0.225,-6.954) (-0.175,-6.956) (-0.125,-6.959) (-0.075,-6.960) (-0.025,-6.961) (0.050,-6.960) (0.100,-6.959) (0.150,-6.958) (0.200,-6.955) (0.250,-6.952) (0.300,-6.948) (0.350,-6.943) (0.400,-6.938) (0.450,-6.932) (0.500,-6.925) (0.550,-6.917) (0.600,-6.909) (0.650,-6.899) (0.700,-6.890) (0.750,-6.879) (0.800,-6.867) (0.850,-6.855) (0.900,-6.842) (0.950,-6.828) (1.000,-6.814) (1.050,-6.798) (1.100,-6.782) (1.150,-6.765) (1.200,-6.747) (1.250,-6.728) (1.300,-6.708) (1.350,-6.687) (1.400,-6.665) (1.450,-6.642) (1.500,-6.619) (1.550,-6.594) (1.600,-6.568) (1.650,-6.541) (1.700,-6.513) (1.750,-6.484) (1.800,-6.454) (1.845,-6.425) (1.900,-6.389) (1.950,-6.355) (1.992,-6.325) (2.050,-6.282) (2.092,-6.250) (2.150,-6.204) (2.200,-6.162) (2.242,-6.125) (2.297,-6.075) (2.350,-6.025) (2.400,-5.975) (2.448,-5.925) (2.494,-5.875) (2.538,-5.825) (2.575,-5.782) (2.621,-5.725) (2.660,-5.675) (2.697,-5.625) (2.732,-5.575) (2.766,-5.525) (2.799,-5.475) (2.830,-5.425) (2.859,-5.375) (2.888,-5.325) (2.915,-5.275) (2.941,-5.225) (2.966,-5.175) (2.989,-5.125) (3.012,-5.075) (3.033,-5.025) (3.053,-4.975) (3.072,-4.925) (3.090,-4.875) (3.107,-4.825) (3.123,-4.775) (3.138,-4.725) (3.151,-4.675) (3.164,-4.625) (3.176,-4.575) (3.187,-4.525) (3.197,-4.475) (3.206,-4.425) (3.214,-4.375) (3.221,-4.325) (3.227,-4.275) (3.232,-4.225) (3.236,-4.175) (3.239,-4.125) (3.242,-4.075) (3.243,-4.025) (3.244,-3.975) (3.243,-3.925) (3.242,-3.875) (3.240,-3.825) (3.237,-3.775) (3.232,-3.725) (3.227,-3.675) (3.221,-3.625) (3.214,-3.575) (3.207,-3.525) (3.198,-3.475) (3.188,-3.425) (3.177,-3.375) (3.165,-3.325) (3.153,-3.275) (3.139,-3.225) (3.124,-3.175) (3.109,-3.125) (3.092,-3.075) (3.075,-3.027) (3.056,-2.975) (3.036,-2.925) (3.015,-2.875) (2.993,-2.825) (2.970,-2.775) (2.946,-2.725) (2.921,-2.675) (2.895,-2.625) (2.867,-2.575) (2.839,-2.525) (2.809,-2.475) (2.778,-2.425) (2.750,-2.381) (2.713,-2.325) (2.678,-2.275) (2.643,-2.225) (2.606,-2.175) (2.575,-2.135) (2.528,-2.075) (2.488,-2.025) (2.450,-1.980) (2.402,-1.925) (2.358,-1.875) (2.312,-1.825) (2.275,-1.785) (2.225,-1.733) (2.175,-1.683) (2.125,-1.633) (2.075,-1.585) (2.038,-1.550) (2.000,-1.515) (1.957,-1.475) (1.901,-1.425) (1.850,-1.381) (1.800,-1.337) (1.756,-1.300) (1.700,-1.253) (1.650,-1.212) (1.605,-1.175) (1.550,-1.131) (1.500,-1.091) (1.450,-1.052) (1.400,-1.013) (1.351,-0.975) (1.300,-0.936) (1.252,-0.900) (1.200,-0.861) (1.152,-0.825) (1.100,-0.787) (1.051,-0.750) (1.000,-0.713) (0.950,-0.677) (0.900,-0.640) (0.850,-0.604) (0.800,-0.568) (0.750,-0.532) (0.705,-0.500) (0.650,-0.461) (0.600,-0.425) (0.550,-0.390) (0.500,-0.354) (0.450,-0.319) (0.400,-0.284) (0.354,-0.250) (0.300,-0.213) (0.250,-0.177) (0.206,-0.144)}
\EQ{(0.206,0.144) (0.248,0.175) (0.300,0.213) (0.350,0.248) (0.400,0.284) (0.450,0.319) (0.495,0.350) (0.550,0.390) (0.600,0.425) (0.650,0.461) (0.700,0.497) (0.750,0.532) (0.800,0.568) (0.845,0.600) (0.900,0.640) (0.948,0.675) (1.000,0.713) (1.050,0.750) (1.100,0.787) (1.150,0.823) (1.200,0.861) (1.250,0.898) (1.300,0.936) (1.350,0.974) (1.400,1.013) (1.448,1.050) (1.500,1.091) (1.543,1.125) (1.600,1.171) (1.650,1.212) (1.697,1.250) (1.750,1.295) (1.800,1.337) (1.844,1.375) (1.900,1.424) (1.950,1.469) (2.000,1.515) (2.038,1.550) (2.075,1.585) (2.117,1.625) (2.167,1.675) (2.217,1.725) (2.265,1.775) (2.300,1.812) (2.335,1.850) (2.375,1.894) (2.424,1.950) (2.467,2.000) (2.500,2.040) (2.548,2.100) (2.587,2.150) (2.624,2.200) (2.661,2.250) (2.696,2.300) (2.725,2.343) (2.762,2.400) (2.794,2.450) (2.824,2.500) (2.850,2.545) (2.875,2.589) (2.900,2.635) (2.925,2.683) (2.950,2.733) (2.975,2.786) (3.000,2.841) (3.025,2.899) (3.046,2.950) (3.065,3.000) (3.083,3.050) (3.100,3.099) (3.117,3.150) (3.132,3.200) (3.146,3.250) (3.159,3.300) (3.171,3.350) (3.183,3.400) (3.193,3.450) (3.202,3.500) (3.211,3.550) (3.218,3.600) (3.224,3.650) (3.230,3.700) (3.235,3.750) (3.238,3.800) (3.241,3.850) (3.243,3.900) (3.244,3.950) (3.244,4.000) (3.243,4.050) (3.241,4.100) (3.238,4.150) (3.234,4.200) (3.229,4.250) (3.224,4.300) (3.217,4.350) (3.210,4.400) (3.201,4.450) (3.192,4.500) (3.182,4.550) (3.170,4.600) (3.158,4.650) (3.145,4.700) (3.130,4.750) (3.115,4.800) (3.099,4.850) (3.081,4.900) (3.063,4.950) (3.043,5.000) (3.022,5.050) (3.001,5.100) (2.978,5.150) (2.954,5.200) (2.928,5.250) (2.902,5.300) (2.875,5.348) (2.845,5.400) (2.814,5.450) (2.782,5.500) (2.750,5.549) (2.715,5.600) (2.678,5.650) (2.640,5.700) (2.601,5.750) (2.560,5.800) (2.525,5.840) (2.475,5.896) (2.425,5.949) (2.375,6.000) (2.325,6.049) (2.275,6.096) (2.225,6.140) (2.184,6.175) (2.125,6.224) (2.075,6.263) (2.026,6.300) (1.975,6.337) (1.925,6.372) (1.875,6.406) (1.825,6.438) (1.775,6.469) (1.725,6.499) (1.679,6.525) (1.634,6.550) (1.587,6.575) (1.538,6.600) (1.487,6.625) (1.433,6.650) (1.377,6.675) (1.325,6.697) (1.275,6.718) (1.225,6.737) (1.175,6.756) (1.125,6.773) (1.075,6.790) (1.025,6.806) (0.975,6.821) (0.925,6.835) (0.875,6.849) (0.825,6.861) (0.775,6.873) (0.725,6.884) (0.675,6.895) (0.625,6.904) (0.575,6.913) (0.525,6.921) (0.475,6.928) (0.425,6.935) (0.375,6.941) (0.325,6.946) (0.275,6.950) (0.225,6.954) (0.175,6.956) (0.125,6.959) (0.075,6.960) (0.025,6.961) (-0.050,6.960) (-0.100,6.959) (-0.150,6.958) (-0.200,6.955) (-0.250,6.952) (-0.300,6.948) (-0.350,6.943) (-0.400,6.938) (-0.450,6.932) (-0.500,6.925) (-0.550,6.917) (-0.600,6.909) (-0.650,6.899) (-0.700,6.890) (-0.750,6.879) (-0.800,6.867) (-0.850,6.855) (-0.900,6.842) (-0.950,6.828) (-1.000,6.814) (-1.050,6.798) (-1.100,6.782) (-1.150,6.765) (-1.200,6.747) (-1.250,6.728) (-1.300,6.708) (-1.350,6.687) (-1.400,6.665) (-1.450,6.642) (-1.500,6.619) (-1.550,6.594) (-1.600,6.568) (-1.650,6.541) (-1.700,6.513) (-1.750,6.484) (-1.800,6.454) (-1.845,6.425) (-1.900,6.389) (-1.950,6.355) (-1.992,6.325) (-2.050,6.282) (-2.092,6.250) (-2.150,6.204) (-2.200,6.162) (-2.242,6.125) (-2.297,6.075) (-2.350,6.025) (-2.400,5.975) (-2.448,5.925) (-2.494,5.875) (-2.538,5.825) (-2.575,5.782) (-2.621,5.725) (-2.660,5.675) (-2.697,5.625) (-2.732,5.575) (-2.766,5.525) (-2.799,5.475) (-2.830,5.425) (-2.859,5.375) (-2.888,5.325) (-2.915,5.275) (-2.941,5.225) (-2.966,5.175) (-2.989,5.125) (-3.012,5.075) (-3.033,5.025) (-3.053,4.975) (-3.072,4.925) (-3.090,4.875) (-3.107,4.825) (-3.123,4.775) (-3.138,4.725) (-3.151,4.675) (-3.164,4.625) (-3.176,4.575) (-3.187,4.525) (-3.197,4.475) (-3.206,4.425) (-3.214,4.375) (-3.221,4.325) (-3.227,4.275) (-3.232,4.225) (-3.236,4.175) (-3.239,4.125) (-3.242,4.075) (-3.243,4.025) (-3.244,3.975) (-3.243,3.925) (-3.242,3.875) (-3.240,3.825) (-3.237,3.775) (-3.232,3.725) (-3.227,3.675) (-3.221,3.625) (-3.214,3.575) (-3.207,3.525) (-3.198,3.475) (-3.188,3.425) (-3.177,3.375) (-3.165,3.325) (-3.153,3.275) (-3.139,3.225) (-3.124,3.175) (-3.109,3.125) (-3.092,3.075) (-3.075,3.027) (-3.056,2.975) (-3.036,2.925) (-3.015,2.875) (-2.993,2.825) (-2.970,2.775) (-2.946,2.725) (-2.921,2.675) (-2.895,2.625) (-2.867,2.575) (-2.839,2.525) (-2.809,2.475) (-2.778,2.425) (-2.750,2.381) (-2.713,2.325) (-2.678,2.275) (-2.643,2.225) (-2.606,2.175) (-2.575,2.135) (-2.528,2.075) (-2.488,2.025) (-2.450,1.980) (-2.402,1.925) (-2.358,1.875) (-2.312,1.825) (-2.275,1.785) (-2.225,1.733) (-2.175,1.683) (-2.125,1.633) (-2.075,1.585) (-2.038,1.550) (-2.000,1.515) (-1.957,1.475) (-1.901,1.425) (-1.850,1.381) (-1.800,1.337) (-1.756,1.300) (-1.700,1.253) (-1.650,1.212) (-1.605,1.175) (-1.550,1.131) (-1.500,1.091) (-1.450,1.052) (-1.400,1.013) (-1.351,0.975) (-1.300,0.936) (-1.252,0.900) (-1.200,0.861) (-1.152,0.825) (-1.100,0.787) (-1.051,0.750) (-1.000,0.713) (-0.950,0.677) (-0.900,0.640) (-0.850,0.604) (-0.800,0.568) (-0.750,0.532) (-0.705,0.500) (-0.650,0.461) (-0.600,0.425) (-0.550,0.390) (-0.500,0.354) (-0.450,0.319) (-0.400,0.284) (-0.354,0.250) (-0.300,0.213) (-0.250,0.177) (-0.220,0.155)}
\EQ{(-0.220,-0.155) (-0.275,-0.195) (-0.319,-0.225) (-0.375,-0.266) (-0.424,-0.300) (-0.475,-0.337) (-0.525,-0.372) (-0.575,-0.408) (-0.625,-0.443) (-0.670,-0.475) (-0.725,-0.515) (-0.774,-0.550) (-0.825,-0.587) (-0.875,-0.623) (-0.925,-0.659) (-0.975,-0.696) (-1.025,-0.733) (-1.075,-0.770) (-1.116,-0.800) (-1.175,-0.844) (-1.216,-0.875) (-1.275,-0.920) (-1.325,-0.959) (-1.375,-0.998) (-1.425,-1.037) (-1.472,-1.075) (-1.525,-1.118) (-1.565,-1.150) (-1.625,-1.200) (-1.675,-1.242) (-1.713,-1.275) (-1.771,-1.325) (-1.825,-1.374) (-1.875,-1.419) (-1.925,-1.466) (-1.962,-1.500) (-2.000,-1.537) (-2.038,-1.575) (-2.075,-1.612) (-2.112,-1.650) (-2.150,-1.690) (-2.200,-1.744) (-2.250,-1.799) (-2.294,-1.850) (-2.336,-1.900) (-2.375,-1.948) (-2.417,-2.000) (-2.450,-2.044) (-2.491,-2.100) (-2.525,-2.147) (-2.561,-2.200) (-2.594,-2.250) (-2.625,-2.299) (-2.656,-2.350) (-2.685,-2.400) (-2.713,-2.450) (-2.740,-2.500) (-2.765,-2.550) (-2.790,-2.600) (-2.813,-2.650) (-2.835,-2.700) (-2.856,-2.750) (-2.875,-2.799) (-2.894,-2.850) (-2.912,-2.900) (-2.928,-2.950) (-2.944,-3.000) (-2.958,-3.050) (-2.972,-3.100) (-2.984,-3.150) (-2.995,-3.200) (-3.005,-3.250) (-3.015,-3.300) (-3.023,-3.350) (-3.030,-3.400) (-3.036,-3.450) (-3.041,-3.500) (-3.045,-3.550) (-3.049,-3.600) (-3.051,-3.650) (-3.052,-3.700) (-3.052,-3.750) (-3.052,-3.800) (-3.050,-3.850) (-3.047,-3.900) (-3.044,-3.950) (-3.039,-4.000) (-3.033,-4.050) (-3.027,-4.100) (-3.019,-4.150) (-3.010,-4.200) (-3.001,-4.250) (-2.990,-4.300) (-2.978,-4.350) (-2.965,-4.400) (-2.951,-4.450) (-2.937,-4.500) (-2.921,-4.550) (-2.903,-4.600) (-2.885,-4.650) (-2.866,-4.700) (-2.845,-4.750) (-2.825,-4.796) (-2.800,-4.850) (-2.776,-4.900) (-2.751,-4.950) (-2.725,-4.998) (-2.700,-5.042) (-2.666,-5.100) (-2.635,-5.150) (-2.602,-5.200) (-2.568,-5.250) (-2.533,-5.300) (-2.500,-5.344) (-2.456,-5.400) (-2.415,-5.450) (-2.375,-5.497) (-2.328,-5.550) (-2.281,-5.600) (-2.231,-5.650) (-2.180,-5.700) (-2.125,-5.750) (-2.075,-5.794) (-2.025,-5.836) (-1.977,-5.875) (-1.925,-5.915) (-1.878,-5.950) (-1.825,-5.987) (-1.775,-6.021) (-1.731,-6.050) (-1.675,-6.085) (-1.625,-6.115) (-1.575,-6.143) (-1.525,-6.170) (-1.475,-6.197) (-1.425,-6.222) (-1.375,-6.246) (-1.325,-6.269) (-1.275,-6.291) (-1.225,-6.311) (-1.175,-6.331) (-1.126,-6.350) (-1.075,-6.368) (-1.025,-6.385) (-0.975,-6.401) (-0.925,-6.417) (-0.875,-6.431) (-0.825,-6.445) (-0.775,-6.457) (-0.725,-6.469) (-0.675,-6.480) (-0.625,-6.490) (-0.575,-6.499) (-0.525,-6.508) (-0.475,-6.516) (-0.425,-6.523) (-0.375,-6.529) (-0.325,-6.534) (-0.275,-6.539) (-0.225,-6.543) (-0.175,-6.546) (-0.125,-6.548) (-0.075,-6.549) (-0.025,-6.550) (0.047,-6.550) (0.100,-6.549) (0.150,-6.547) (0.200,-6.544) (0.250,-6.541) (0.300,-6.537) (0.350,-6.532) (0.400,-6.526) (0.450,-6.519) (0.500,-6.512) (0.550,-6.504) (0.600,-6.495) (0.650,-6.485) (0.700,-6.475) (0.750,-6.463) (0.800,-6.451) (0.850,-6.438) (0.900,-6.424) (0.950,-6.409) (1.000,-6.393) (1.050,-6.377) (1.100,-6.359) (1.150,-6.341) (1.200,-6.322) (1.250,-6.301) (1.300,-6.280) (1.350,-6.257) (1.400,-6.234) (1.450,-6.209) (1.500,-6.184) (1.550,-6.157) (1.600,-6.129) (1.650,-6.100) (1.700,-6.069) (1.750,-6.038) (1.800,-6.004) (1.842,-5.975) (1.900,-5.934) (1.945,-5.900) (2.000,-5.857) (2.050,-5.816) (2.097,-5.775) (2.150,-5.728) (2.200,-5.681) (2.250,-5.632) (2.300,-5.580) (2.350,-5.525) (2.394,-5.475) (2.436,-5.425) (2.475,-5.376) (2.514,-5.325) (2.550,-5.276) (2.586,-5.225) (2.619,-5.175) (2.650,-5.126) (2.681,-5.075) (2.710,-5.025) (2.737,-4.975) (2.763,-4.925) (2.788,-4.875) (2.812,-4.825) (2.834,-4.775) (2.856,-4.725) (2.875,-4.676) (2.894,-4.625) (2.912,-4.575) (2.929,-4.525) (2.944,-4.475) (2.959,-4.425) (2.972,-4.375) (2.984,-4.325) (2.995,-4.275) (3.006,-4.225) (3.015,-4.175) (3.023,-4.125) (3.030,-4.075) (3.036,-4.025) (3.041,-3.975) (3.046,-3.925) (3.049,-3.875) (3.051,-3.825) (3.052,-3.775) (3.052,-3.725) (3.052,-3.675) (3.050,-3.625) (3.047,-3.575) (3.044,-3.525) (3.039,-3.475) (3.033,-3.425) (3.026,-3.375) (3.019,-3.325) (3.010,-3.275) (3.000,-3.225) (2.990,-3.175) (2.978,-3.125) (2.965,-3.075) (2.951,-3.025) (2.936,-2.975) (2.920,-2.925) (2.903,-2.875) (2.885,-2.825) (2.866,-2.775) (2.845,-2.725) (2.825,-2.677) (2.801,-2.625) (2.778,-2.575) (2.753,-2.525) (2.726,-2.475) (2.700,-2.427) (2.675,-2.383) (2.641,-2.325) (2.610,-2.275) (2.578,-2.225) (2.544,-2.175) (2.509,-2.125) (2.475,-2.078) (2.436,-2.025) (2.400,-1.979) (2.357,-1.925) (2.315,-1.875) (2.275,-1.828) (2.228,-1.775) (2.183,-1.725) (2.136,-1.675) (2.100,-1.638) (2.063,-1.600) (2.025,-1.562) (1.987,-1.525) (1.950,-1.489) (1.909,-1.450) (1.854,-1.400) (1.800,-1.351) (1.750,-1.307) (1.700,-1.264) (1.655,-1.225) (1.600,-1.179) (1.550,-1.138) (1.504,-1.100) (1.450,-1.057) (1.410,-1.025) (1.350,-0.978) (1.300,-0.940) (1.250,-0.901) (1.200,-0.863) (1.150,-0.826) (1.100,-0.788) (1.050,-0.751) (1.000,-0.714) (0.950,-0.678) (0.900,-0.641) (0.850,-0.605) (0.809,-0.575) (0.750,-0.533) (0.705,-0.500) (0.650,-0.461) (0.600,-0.425) (0.550,-0.390) (0.500,-0.354) (0.450,-0.319) (0.400,-0.284) (0.354,-0.250) (0.300,-0.213) (0.250,-0.177) (0.206,-0.144)}
\EQ{(0.206,0.144) (0.248,0.175) (0.300,0.213) (0.350,0.248) (0.400,0.284) (0.450,0.319) (0.495,0.350) (0.550,0.390) (0.600,0.425) (0.650,0.461) (0.700,0.497) (0.750,0.533) (0.800,0.569) (0.844,0.600) (0.900,0.641) (0.947,0.675) (1.000,0.714) (1.049,0.750) (1.100,0.788) (1.149,0.825) (1.200,0.863) (1.249,0.900) (1.300,0.940) (1.346,0.975) (1.400,1.018) (1.441,1.050) (1.500,1.097) (1.550,1.138) (1.595,1.175) (1.650,1.221) (1.700,1.264) (1.742,1.300) (1.799,1.350) (1.850,1.396) (1.900,1.442) (1.950,1.489) (1.987,1.525) (2.025,1.562) (2.063,1.600) (2.100,1.638) (2.136,1.675) (2.175,1.717) (2.225,1.771) (2.273,1.825) (2.315,1.875) (2.350,1.917) (2.397,1.975) (2.436,2.025) (2.473,2.075) (2.509,2.125) (2.544,2.175) (2.575,2.221) (2.610,2.275) (2.641,2.325) (2.671,2.375) (2.699,2.425) (2.725,2.472) (2.750,2.520) (2.775,2.570) (2.800,2.622) (2.824,2.675) (2.845,2.725) (2.866,2.775) (2.885,2.825) (2.903,2.875) (2.920,2.925) (2.936,2.975) (2.951,3.025) (2.965,3.075) (2.978,3.125) (2.990,3.175) (3.000,3.225) (3.010,3.275) (3.019,3.325) (3.026,3.375) (3.033,3.425) (3.039,3.475) (3.044,3.525) (3.047,3.575) (3.050,3.625) (3.052,3.675) (3.052,3.725) (3.052,3.775) (3.051,3.825) (3.049,3.875) (3.046,3.925) (3.041,3.975) (3.036,4.025) (3.030,4.075) (3.023,4.125) (3.015,4.175) (3.006,4.225) (2.995,4.275) (2.984,4.325) (2.972,4.375) (2.959,4.425) (2.944,4.475) (2.929,4.525) (2.912,4.575) (2.894,4.625) (2.876,4.675) (2.856,4.725) (2.834,4.775) (2.812,4.825) (2.788,4.875) (2.763,4.925) (2.737,4.975) (2.710,5.025) (2.681,5.075) (2.651,5.125) (2.619,5.175) (2.586,5.225) (2.551,5.275) (2.514,5.325) (2.476,5.375) (2.436,5.425) (2.400,5.468) (2.350,5.525) (2.305,5.575) (2.256,5.625) (2.206,5.675) (2.153,5.725) (2.100,5.773) (2.050,5.816) (2.008,5.850) (1.950,5.896) (1.900,5.934) (1.850,5.970) (1.806,6.000) (1.750,6.038) (1.700,6.069) (1.650,6.100) (1.600,6.129) (1.550,6.157) (1.500,6.184) (1.450,6.209) (1.400,6.234) (1.350,6.257) (1.300,6.280) (1.253,6.300) (1.200,6.322) (1.150,6.341) (1.100,6.359) (1.050,6.377) (1.000,6.393) (0.950,6.409) (0.900,6.424) (0.850,6.438) (0.800,6.451) (0.750,6.463) (0.700,6.475) (0.650,6.485) (0.600,6.495) (0.550,6.504) (0.500,6.512) (0.450,6.519) (0.400,6.526) (0.350,6.532) (0.300,6.537) (0.250,6.541) (0.200,6.544) (0.150,6.547) (0.100,6.549) (0.050,6.550) (0.000,6.550) (-0.050,6.550) (-0.100,6.549) (-0.150,6.547) (-0.200,6.544) (-0.250,6.541) (-0.300,6.537) (-0.350,6.532) (-0.400,6.526) (-0.450,6.519) (-0.500,6.512) (-0.550,6.504) (-0.600,6.495) (-0.650,6.485) (-0.700,6.475) (-0.750,6.463) (-0.800,6.451) (-0.850,6.438) (-0.900,6.424) (-0.950,6.409) (-1.000,6.393) (-1.050,6.377) (-1.100,6.359) (-1.150,6.341) (-1.200,6.322) (-1.250,6.301) (-1.300,6.280) (-1.350,6.257) (-1.400,6.234) (-1.450,6.209) (-1.500,6.184) (-1.550,6.157) (-1.600,6.129) (-1.650,6.100) (-1.700,6.069) (-1.750,6.038) (-1.800,6.004) (-1.842,5.975) (-1.900,5.934) (-1.945,5.900) (-2.000,5.857) (-2.050,5.816) (-2.097,5.775) (-2.150,5.728) (-2.200,5.681) (-2.250,5.632) (-2.300,5.580) (-2.350,5.525) (-2.394,5.475) (-2.436,5.425) (-2.475,5.376) (-2.514,5.325) (-2.550,5.276) (-2.586,5.225) (-2.619,5.175) (-2.650,5.126) (-2.681,5.075) (-2.710,5.025) (-2.737,4.975) (-2.763,4.925) (-2.788,4.875) (-2.812,4.825) (-2.834,4.775) (-2.856,4.725) (-2.875,4.676) (-2.894,4.625) (-2.912,4.575) (-2.929,4.525) (-2.944,4.475) (-2.959,4.425) (-2.972,4.375) (-2.984,4.325) (-2.995,4.275) (-3.006,4.225) (-3.015,4.175) (-3.023,4.125) (-3.030,4.075) (-3.036,4.025) (-3.041,3.975) (-3.046,3.925) (-3.049,3.875) (-3.051,3.825) (-3.052,3.775) (-3.052,3.725) (-3.052,3.675) (-3.050,3.625) (-3.047,3.575) (-3.044,3.525) (-3.039,3.475) (-3.033,3.425) (-3.026,3.375) (-3.019,3.325) (-3.010,3.275) (-3.000,3.225) (-2.990,3.175) (-2.978,3.125) (-2.965,3.075) (-2.951,3.025) (-2.936,2.975) (-2.920,2.925) (-2.903,2.875) (-2.885,2.825) (-2.866,2.775) (-2.845,2.725) (-2.825,2.677) (-2.801,2.625) (-2.778,2.575) (-2.753,2.525) (-2.726,2.475) (-2.700,2.427) (-2.675,2.383) (-2.641,2.325) (-2.610,2.275) (-2.578,2.225) (-2.544,2.175) (-2.509,2.125) (-2.475,2.078) (-2.436,2.025) (-2.400,1.979) (-2.357,1.925) (-2.315,1.875) (-2.275,1.828) (-2.228,1.775) (-2.183,1.725) (-2.136,1.675) (-2.100,1.638) (-2.063,1.600) (-2.025,1.562) (-1.987,1.525) (-1.950,1.489) (-1.909,1.450) (-1.854,1.400) (-1.800,1.351) (-1.750,1.307) (-1.700,1.264) (-1.655,1.225) (-1.600,1.179) (-1.550,1.138) (-1.504,1.100) (-1.450,1.057) (-1.410,1.025) (-1.350,0.978) (-1.300,0.940) (-1.250,0.901) (-1.200,0.863) (-1.150,0.826) (-1.100,0.788) (-1.050,0.751) (-1.000,0.714) (-0.950,0.678) (-0.900,0.641) (-0.850,0.605) (-0.809,0.575) (-0.750,0.533) (-0.705,0.500) (-0.650,0.461) (-0.600,0.425) (-0.550,0.390) (-0.500,0.354) (-0.450,0.319) (-0.400,0.284) (-0.354,0.250) (-0.300,0.213) (-0.250,0.177) (-0.220,0.155)}
\EQ{(-0.220,-0.155) (-0.275,-0.195) (-0.319,-0.225) (-0.375,-0.266) (-0.424,-0.300) (-0.475,-0.337) (-0.525,-0.372) (-0.575,-0.408) (-0.625,-0.444) (-0.670,-0.475) (-0.725,-0.515) (-0.774,-0.550) (-0.825,-0.587) (-0.875,-0.624) (-0.925,-0.660) (-0.975,-0.697) (-1.025,-0.734) (-1.075,-0.771) (-1.125,-0.809) (-1.175,-0.847) (-1.225,-0.885) (-1.275,-0.924) (-1.325,-0.963) (-1.372,-1.000) (-1.425,-1.043) (-1.465,-1.075) (-1.525,-1.124) (-1.575,-1.166) (-1.615,-1.200) (-1.672,-1.250) (-1.725,-1.297) (-1.775,-1.342) (-1.825,-1.389) (-1.863,-1.425) (-1.900,-1.461) (-1.940,-1.500) (-1.990,-1.550) (-2.025,-1.587) (-2.061,-1.625) (-2.100,-1.668) (-2.150,-1.724) (-2.194,-1.775) (-2.236,-1.825) (-2.275,-1.874) (-2.315,-1.925) (-2.350,-1.972) (-2.388,-2.025) (-2.423,-2.075) (-2.456,-2.125) (-2.488,-2.175) (-2.519,-2.225) (-2.548,-2.275) (-2.575,-2.324) (-2.600,-2.370) (-2.625,-2.419) (-2.650,-2.470) (-2.675,-2.524) (-2.697,-2.575) (-2.718,-2.625) (-2.737,-2.675) (-2.756,-2.725) (-2.773,-2.775) (-2.789,-2.825) (-2.804,-2.875) (-2.818,-2.925) (-2.830,-2.975) (-2.842,-3.025) (-2.852,-3.075) (-2.862,-3.125) (-2.870,-3.175) (-2.877,-3.225) (-2.884,-3.275) (-2.889,-3.325) (-2.893,-3.375) (-2.896,-3.425) (-2.898,-3.475) (-2.899,-3.525) (-2.899,-3.575) (-2.898,-3.625) (-2.896,-3.675) (-2.893,-3.725) (-2.889,-3.775) (-2.884,-3.825) (-2.878,-3.875) (-2.871,-3.925) (-2.862,-3.975) (-2.853,-4.025) (-2.842,-4.075) (-2.831,-4.125) (-2.818,-4.175) (-2.804,-4.225) (-2.789,-4.275) (-2.773,-4.325) (-2.756,-4.375) (-2.738,-4.425) (-2.718,-4.475) (-2.697,-4.525) (-2.675,-4.575) (-2.652,-4.625) (-2.627,-4.675) (-2.601,-4.725) (-2.575,-4.772) (-2.550,-4.815) (-2.514,-4.875) (-2.482,-4.925) (-2.450,-4.972) (-2.413,-5.025) (-2.376,-5.075) (-2.337,-5.125) (-2.300,-5.170) (-2.253,-5.225) (-2.208,-5.275) (-2.161,-5.325) (-2.125,-5.362) (-2.086,-5.400) (-2.032,-5.450) (-1.976,-5.500) (-1.925,-5.543) (-1.885,-5.575) (-1.825,-5.621) (-1.775,-5.658) (-1.725,-5.693) (-1.678,-5.725) (-1.625,-5.759) (-1.575,-5.790) (-1.525,-5.819) (-1.475,-5.847) (-1.425,-5.874) (-1.375,-5.900) (-1.325,-5.924) (-1.275,-5.947) (-1.225,-5.970) (-1.175,-5.991) (-1.125,-6.011) (-1.075,-6.030) (-1.025,-6.048) (-0.975,-6.065) (-0.925,-6.081) (-0.875,-6.096) (-0.825,-6.111) (-0.775,-6.124) (-0.725,-6.137) (-0.675,-6.148) (-0.625,-6.159) (-0.575,-6.169) (-0.525,-6.178) (-0.475,-6.186) (-0.425,-6.193) (-0.375,-6.200) (-0.325,-6.205) (-0.275,-6.210) (-0.225,-6.214) (-0.175,-6.217) (-0.125,-6.220) (-0.075,-6.221) (-0.025,-6.222) (0.050,-6.222) (0.100,-6.221) (0.150,-6.219) (0.200,-6.216) (0.250,-6.212) (0.300,-6.208) (0.350,-6.203) (0.400,-6.196) (0.450,-6.190) (0.500,-6.182) (0.550,-6.173) (0.600,-6.164) (0.650,-6.154) (0.700,-6.142) (0.750,-6.130) (0.800,-6.117) (0.850,-6.104) (0.900,-6.089) (0.950,-6.073) (1.000,-6.057) (1.050,-6.039) (1.100,-6.020) (1.150,-6.001) (1.200,-5.980) (1.250,-5.959) (1.300,-5.936) (1.350,-5.912) (1.400,-5.887) (1.450,-5.861) (1.500,-5.833) (1.550,-5.805) (1.599,-5.775) (1.650,-5.743) (1.700,-5.710) (1.750,-5.676) (1.800,-5.640) (1.850,-5.602) (1.900,-5.563) (1.946,-5.525) (2.000,-5.479) (2.050,-5.434) (2.100,-5.386) (2.137,-5.350) (2.175,-5.311) (2.225,-5.257) (2.275,-5.200) (2.317,-5.150) (2.350,-5.108) (2.394,-5.050) (2.425,-5.008) (2.465,-4.950) (2.498,-4.900) (2.525,-4.857) (2.559,-4.800) (2.587,-4.750) (2.614,-4.700) (2.639,-4.650) (2.663,-4.600) (2.686,-4.550) (2.708,-4.500) (2.728,-4.450) (2.747,-4.400) (2.765,-4.350) (2.781,-4.300) (2.797,-4.250) (2.811,-4.200) (2.825,-4.150) (2.837,-4.100) (2.848,-4.050) (2.858,-4.000) (2.867,-3.950) (2.874,-3.900) (2.881,-3.850) (2.887,-3.800) (2.891,-3.750) (2.895,-3.700) (2.898,-3.650) (2.899,-3.600) (2.900,-3.550) (2.899,-3.500) (2.897,-3.450) (2.895,-3.400) (2.891,-3.350) (2.887,-3.300) (2.881,-3.250) (2.874,-3.200) (2.866,-3.150) (2.857,-3.100) (2.847,-3.050) (2.836,-3.000) (2.824,-2.950) (2.811,-2.900) (2.796,-2.850) (2.781,-2.800) (2.764,-2.750) (2.747,-2.700) (2.728,-2.650) (2.708,-2.600) (2.687,-2.550) (2.664,-2.500) (2.640,-2.450) (2.616,-2.400) (2.589,-2.350) (2.562,-2.300) (2.533,-2.250) (2.503,-2.200) (2.475,-2.155) (2.440,-2.100) (2.406,-2.050) (2.375,-2.007) (2.334,-1.950) (2.300,-1.906) (2.256,-1.850) (2.215,-1.800) (2.175,-1.752) (2.129,-1.700) (2.084,-1.650) (2.050,-1.613) (2.014,-1.575) (1.975,-1.535) (1.925,-1.485) (1.889,-1.450) (1.850,-1.413) (1.811,-1.375) (1.756,-1.325) (1.700,-1.275) (1.650,-1.231) (1.600,-1.188) (1.556,-1.150) (1.500,-1.104) (1.450,-1.063) (1.404,-1.025) (1.350,-0.983) (1.309,-0.950) (1.250,-0.904) (1.200,-0.866) (1.150,-0.828) (1.100,-0.790) (1.050,-0.753) (1.000,-0.716) (0.950,-0.679) (0.900,-0.642) (0.850,-0.606) (0.809,-0.575) (0.750,-0.533) (0.705,-0.500) (0.650,-0.461) (0.600,-0.425) (0.550,-0.390) (0.500,-0.354) (0.450,-0.319) (0.400,-0.284) (0.354,-0.250) (0.300,-0.213) (0.250,-0.177) (0.206,-0.144)}
\EQ{(0.206,0.144) (0.248,0.175) (0.300,0.213) (0.350,0.248) (0.400,0.284) (0.450,0.319) (0.495,0.350) (0.550,0.390) (0.600,0.425) (0.650,0.461) (0.700,0.497) (0.750,0.533) (0.800,0.569) (0.843,0.600) (0.900,0.642) (0.946,0.675) (1.000,0.716) (1.047,0.750) (1.100,0.790) (1.147,0.825) (1.200,0.866) (1.245,0.900) (1.300,0.943) (1.341,0.975) (1.400,1.022) (1.450,1.063) (1.496,1.100) (1.550,1.145) (1.600,1.188) (1.644,1.225) (1.700,1.275) (1.750,1.320) (1.800,1.366) (1.837,1.400) (1.875,1.436) (1.915,1.475) (1.965,1.525) (2.000,1.561) (2.038,1.600) (2.075,1.640) (2.125,1.695) (2.173,1.750) (2.215,1.800) (2.250,1.843) (2.296,1.900) (2.334,1.950) (2.370,2.000) (2.400,2.042) (2.440,2.100) (2.472,2.150) (2.500,2.195) (2.533,2.250) (2.562,2.300) (2.589,2.350) (2.616,2.400) (2.640,2.450) (2.664,2.500) (2.687,2.550) (2.708,2.600) (2.728,2.650) (2.747,2.700) (2.764,2.750) (2.781,2.800) (2.796,2.850) (2.811,2.900) (2.824,2.950) (2.836,3.000) (2.847,3.050) (2.857,3.100) (2.866,3.150) (2.874,3.200) (2.881,3.250) (2.887,3.300) (2.891,3.350) (2.895,3.400) (2.897,3.450) (2.899,3.500) (2.900,3.550) (2.899,3.600) (2.898,3.650) (2.895,3.700) (2.891,3.750) (2.887,3.800) (2.881,3.850) (2.874,3.900) (2.867,3.950) (2.858,4.000) (2.848,4.050) (2.837,4.100) (2.825,4.150) (2.811,4.200) (2.797,4.250) (2.781,4.300) (2.765,4.350) (2.747,4.400) (2.728,4.450) (2.708,4.500) (2.686,4.550) (2.663,4.600) (2.639,4.650) (2.614,4.700) (2.587,4.750) (2.559,4.800) (2.529,4.850) (2.500,4.897) (2.465,4.950) (2.431,5.000) (2.400,5.043) (2.357,5.100) (2.325,5.140) (2.275,5.200) (2.231,5.250) (2.185,5.300) (2.150,5.337) (2.112,5.375) (2.075,5.410) (2.032,5.450) (1.976,5.500) (1.925,5.543) (1.885,5.575) (1.825,5.621) (1.775,5.658) (1.725,5.693) (1.678,5.725) (1.625,5.759) (1.575,5.790) (1.525,5.819) (1.475,5.847) (1.425,5.874) (1.375,5.900) (1.325,5.924) (1.275,5.947) (1.225,5.970) (1.175,5.991) (1.125,6.011) (1.075,6.030) (1.025,6.048) (0.975,6.065) (0.925,6.081) (0.875,6.096) (0.825,6.111) (0.775,6.124) (0.725,6.137) (0.675,6.148) (0.625,6.159) (0.575,6.169) (0.525,6.178) (0.475,6.186) (0.425,6.193) (0.375,6.200) (0.325,6.205) (0.275,6.210) (0.225,6.214) (0.175,6.217) (0.125,6.220) (0.075,6.221) (0.025,6.222) (-0.050,6.222) (-0.100,6.221) (-0.150,6.219) (-0.200,6.216) (-0.250,6.212) (-0.300,6.208) (-0.350,6.203) (-0.400,6.196) (-0.450,6.190) (-0.500,6.182) (-0.550,6.173) (-0.600,6.164) (-0.650,6.154) (-0.700,6.142) (-0.750,6.130) (-0.800,6.117) (-0.850,6.104) (-0.900,6.089) (-0.950,6.073) (-1.000,6.057) (-1.050,6.039) (-1.100,6.020) (-1.150,6.001) (-1.200,5.980) (-1.250,5.959) (-1.300,5.936) (-1.350,5.912) (-1.400,5.887) (-1.450,5.861) (-1.500,5.833) (-1.550,5.805) (-1.599,5.775) (-1.650,5.743) (-1.700,5.710) (-1.750,5.676) (-1.800,5.640) (-1.850,5.602) (-1.900,5.563) (-1.946,5.525) (-2.000,5.479) (-2.050,5.434) (-2.100,5.386) (-2.137,5.350) (-2.175,5.311) (-2.225,5.257) (-2.275,5.200) (-2.317,5.150) (-2.350,5.108) (-2.394,5.050) (-2.425,5.008) (-2.465,4.950) (-2.498,4.900) (-2.525,4.857) (-2.559,4.800) (-2.587,4.750) (-2.614,4.700) (-2.639,4.650) (-2.663,4.600) (-2.686,4.550) (-2.708,4.500) (-2.728,4.450) (-2.747,4.400) (-2.765,4.350) (-2.781,4.300) (-2.797,4.250) (-2.811,4.200) (-2.825,4.150) (-2.837,4.100) (-2.848,4.050) (-2.858,4.000) (-2.867,3.950) (-2.874,3.900) (-2.881,3.850) (-2.887,3.800) (-2.891,3.750) (-2.895,3.700) (-2.898,3.650) (-2.899,3.600) (-2.900,3.550) (-2.899,3.500) (-2.897,3.450) (-2.895,3.400) (-2.891,3.350) (-2.887,3.300) (-2.881,3.250) (-2.874,3.200) (-2.866,3.150) (-2.857,3.100) (-2.847,3.050) (-2.836,3.000) (-2.824,2.950) (-2.811,2.900) (-2.796,2.850) (-2.781,2.800) (-2.764,2.750) (-2.747,2.700) (-2.728,2.650) (-2.708,2.600) (-2.687,2.550) (-2.664,2.500) (-2.640,2.450) (-2.616,2.400) (-2.589,2.350) (-2.562,2.300) (-2.533,2.250) (-2.503,2.200) (-2.475,2.155) (-2.440,2.100) (-2.406,2.050) (-2.375,2.007) (-2.334,1.950) (-2.300,1.906) (-2.256,1.850) (-2.215,1.800) (-2.175,1.752) (-2.129,1.700) (-2.084,1.650) (-2.050,1.613) (-2.014,1.575) (-1.975,1.535) (-1.925,1.485) (-1.889,1.450) (-1.850,1.413) (-1.811,1.375) (-1.756,1.325) (-1.700,1.275) (-1.650,1.231) (-1.600,1.188) (-1.556,1.150) (-1.500,1.104) (-1.450,1.063) (-1.404,1.025) (-1.350,0.983) (-1.309,0.950) (-1.250,0.904) (-1.200,0.866) (-1.150,0.828) (-1.100,0.790) (-1.050,0.753) (-1.000,0.716) (-0.950,0.679) (-0.900,0.642) (-0.850,0.606) (-0.809,0.575) (-0.750,0.533) (-0.705,0.500) (-0.650,0.461) (-0.600,0.425) (-0.550,0.390) (-0.500,0.354) (-0.450,0.319) (-0.400,0.284) (-0.354,0.250) (-0.300,0.213) (-0.250,0.177) (-0.220,0.155)}
\EQ{(-0.220,-0.155) (-0.275,-0.195) (-0.319,-0.225) (-0.375,-0.266) (-0.424,-0.300) (-0.475,-0.337) (-0.525,-0.372) (-0.575,-0.408) (-0.625,-0.444) (-0.670,-0.475) (-0.725,-0.516) (-0.773,-0.550) (-0.825,-0.588) (-0.875,-0.624) (-0.925,-0.661) (-0.975,-0.698) (-1.025,-0.735) (-1.075,-0.773) (-1.125,-0.811) (-1.175,-0.849) (-1.225,-0.888) (-1.273,-0.925) (-1.325,-0.967) (-1.367,-1.000) (-1.425,-1.048) (-1.475,-1.089) (-1.518,-1.125) (-1.575,-1.174) (-1.625,-1.218) (-1.675,-1.263) (-1.715,-1.300) (-1.769,-1.350) (-1.821,-1.400) (-1.871,-1.450) (-1.920,-1.500) (-1.967,-1.550) (-2.013,-1.600) (-2.050,-1.641) (-2.100,-1.699) (-2.142,-1.750) (-2.175,-1.791) (-2.221,-1.850) (-2.258,-1.900) (-2.293,-1.950) (-2.325,-1.996) (-2.361,-2.050) (-2.392,-2.100) (-2.422,-2.150) (-2.450,-2.199) (-2.475,-2.245) (-2.500,-2.292) (-2.525,-2.343) (-2.550,-2.396) (-2.574,-2.450) (-2.595,-2.500) (-2.615,-2.550) (-2.633,-2.600) (-2.650,-2.648) (-2.667,-2.700) (-2.682,-2.750) (-2.696,-2.800) (-2.708,-2.850) (-2.720,-2.900) (-2.730,-2.950) (-2.739,-3.000) (-2.747,-3.050) (-2.754,-3.100) (-2.760,-3.150) (-2.765,-3.200) (-2.769,-3.250) (-2.771,-3.300) (-2.773,-3.350) (-2.773,-3.400) (-2.773,-3.450) (-2.771,-3.500) (-2.768,-3.550) (-2.764,-3.600) (-2.759,-3.650) (-2.753,-3.700) (-2.746,-3.750) (-2.738,-3.800) (-2.729,-3.850) (-2.718,-3.900) (-2.706,-3.950) (-2.694,-4.000) (-2.680,-4.050) (-2.664,-4.100) (-2.648,-4.150) (-2.630,-4.200) (-2.612,-4.250) (-2.592,-4.300) (-2.570,-4.350) (-2.547,-4.400) (-2.525,-4.446) (-2.500,-4.496) (-2.475,-4.542) (-2.450,-4.587) (-2.412,-4.650) (-2.381,-4.700) (-2.350,-4.746) (-2.312,-4.800) (-2.275,-4.850) (-2.237,-4.900) (-2.200,-4.945) (-2.153,-5.000) (-2.108,-5.050) (-2.061,-5.100) (-2.025,-5.136) (-1.984,-5.175) (-1.930,-5.225) (-1.875,-5.273) (-1.825,-5.315) (-1.781,-5.350) (-1.725,-5.392) (-1.679,-5.425) (-1.625,-5.462) (-1.575,-5.495) (-1.527,-5.525) (-1.475,-5.556) (-1.425,-5.585) (-1.375,-5.612) (-1.325,-5.638) (-1.275,-5.662) (-1.225,-5.686) (-1.175,-5.708) (-1.125,-5.729) (-1.075,-5.749) (-1.025,-5.769) (-0.975,-5.787) (-0.925,-5.804) (-0.875,-5.819) (-0.825,-5.834) (-0.775,-5.848) (-0.725,-5.862) (-0.675,-5.874) (-0.625,-5.885) (-0.575,-5.895) (-0.525,-5.905) (-0.475,-5.913) (-0.425,-5.921) (-0.375,-5.928) (-0.325,-5.934) (-0.275,-5.939) (-0.225,-5.943) (-0.175,-5.946) (-0.125,-5.949) (-0.075,-5.950) (-0.025,-5.951) (0.050,-5.951) (0.100,-5.950) (0.150,-5.948) (0.200,-5.945) (0.250,-5.941) (0.300,-5.936) (0.350,-5.931) (0.400,-5.924) (0.450,-5.917) (0.500,-5.909) (0.550,-5.900) (0.600,-5.890) (0.650,-5.879) (0.700,-5.868) (0.750,-5.855) (0.800,-5.842) (0.850,-5.827) (0.900,-5.812) (0.950,-5.795) (1.000,-5.778) (1.050,-5.759) (1.100,-5.740) (1.150,-5.719) (1.200,-5.697) (1.249,-5.675) (1.300,-5.650) (1.350,-5.625) (1.397,-5.600) (1.442,-5.575) (1.486,-5.550) (1.550,-5.511) (1.600,-5.479) (1.643,-5.450) (1.700,-5.410) (1.748,-5.375) (1.800,-5.335) (1.843,-5.300) (1.900,-5.252) (1.950,-5.207) (2.000,-5.160) (2.036,-5.125) (2.075,-5.085) (2.125,-5.031) (2.175,-4.975) (2.217,-4.925) (2.250,-4.883) (2.294,-4.825) (2.325,-4.782) (2.364,-4.725) (2.397,-4.675) (2.425,-4.629) (2.457,-4.575) (2.484,-4.525) (2.511,-4.475) (2.535,-4.425) (2.559,-4.375) (2.581,-4.325) (2.602,-4.275) (2.621,-4.225) (2.639,-4.175) (2.656,-4.125) (2.672,-4.075) (2.687,-4.025) (2.700,-3.975) (2.712,-3.925) (2.723,-3.875) (2.733,-3.825) (2.742,-3.775) (2.750,-3.725) (2.757,-3.675) (2.762,-3.625) (2.766,-3.575) (2.770,-3.525) (2.772,-3.475) (2.773,-3.425) (2.773,-3.375) (2.772,-3.325) (2.770,-3.275) (2.767,-3.225) (2.763,-3.175) (2.757,-3.125) (2.751,-3.075) (2.743,-3.025) (2.735,-2.975) (2.725,-2.926) (2.714,-2.875) (2.702,-2.825) (2.689,-2.775) (2.674,-2.725) (2.659,-2.675) (2.642,-2.625) (2.625,-2.577) (2.605,-2.525) (2.585,-2.475) (2.563,-2.425) (2.541,-2.375) (2.517,-2.325) (2.491,-2.275) (2.464,-2.225) (2.436,-2.175) (2.407,-2.125) (2.376,-2.075) (2.344,-2.025) (2.311,-1.975) (2.276,-1.925) (2.240,-1.875) (2.202,-1.825) (2.162,-1.775) (2.125,-1.729) (2.079,-1.675) (2.036,-1.625) (2.000,-1.586) (1.950,-1.532) (1.900,-1.480) (1.850,-1.429) (1.800,-1.380) (1.750,-1.333) (1.700,-1.286) (1.661,-1.250) (1.604,-1.200) (1.550,-1.153) (1.500,-1.110) (1.458,-1.075) (1.400,-1.027) (1.350,-0.987) (1.304,-0.950) (1.250,-0.907) (1.209,-0.875) (1.150,-0.830) (1.100,-0.792) (1.050,-0.754) (1.000,-0.717) (0.950,-0.680) (0.900,-0.643) (0.850,-0.606) (0.808,-0.575) (0.750,-0.534) (0.704,-0.500) (0.650,-0.462) (0.600,-0.425) (0.550,-0.390) (0.500,-0.354) (0.450,-0.319) (0.400,-0.284) (0.354,-0.250) (0.300,-0.213) (0.250,-0.177) (0.206,-0.144)}
\EQ{(0.206,0.144) (0.248,0.175) (0.300,0.213) (0.350,0.248) (0.400,0.284) (0.450,0.319) (0.495,0.350) (0.550,0.390) (0.599,0.425) (0.650,0.462) (0.700,0.497) (0.750,0.534) (0.800,0.570) (0.842,0.600) (0.900,0.643) (0.944,0.675) (1.000,0.717) (1.045,0.750) (1.100,0.792) (1.144,0.825) (1.200,0.868) (1.241,0.900) (1.300,0.947) (1.350,0.987) (1.398,1.025) (1.450,1.068) (1.500,1.110) (1.547,1.150) (1.600,1.196) (1.650,1.241) (1.688,1.275) (1.725,1.309) (1.769,1.350) (1.821,1.400) (1.871,1.450) (1.920,1.500) (1.967,1.550) (2.013,1.600) (2.050,1.641) (2.100,1.699) (2.142,1.750) (2.175,1.791) (2.221,1.850) (2.258,1.900) (2.293,1.950) (2.325,1.996) (2.361,2.050) (2.392,2.100) (2.422,2.150) (2.450,2.199) (2.475,2.245) (2.500,2.292) (2.525,2.343) (2.550,2.396) (2.574,2.450) (2.595,2.500) (2.615,2.550) (2.633,2.600) (2.650,2.648) (2.667,2.700) (2.682,2.750) (2.696,2.800) (2.708,2.850) (2.720,2.900) (2.730,2.950) (2.739,3.000) (2.747,3.050) (2.754,3.100) (2.760,3.150) (2.765,3.200) (2.769,3.250) (2.771,3.300) (2.773,3.350) (2.773,3.400) (2.773,3.450) (2.771,3.500) (2.768,3.550) (2.764,3.600) (2.759,3.650) (2.753,3.700) (2.746,3.750) (2.738,3.800) (2.729,3.850) (2.718,3.900) (2.706,3.950) (2.694,4.000) (2.680,4.050) (2.664,4.100) (2.648,4.150) (2.630,4.200) (2.612,4.250) (2.592,4.300) (2.570,4.350) (2.547,4.400) (2.525,4.446) (2.500,4.496) (2.475,4.542) (2.450,4.587) (2.412,4.650) (2.381,4.700) (2.350,4.746) (2.312,4.800) (2.275,4.850) (2.237,4.900) (2.200,4.945) (2.153,5.000) (2.108,5.050) (2.061,5.100) (2.025,5.136) (1.984,5.175) (1.930,5.225) (1.875,5.273) (1.825,5.315) (1.781,5.350) (1.725,5.392) (1.679,5.425) (1.625,5.462) (1.575,5.495) (1.527,5.525) (1.475,5.556) (1.425,5.585) (1.375,5.612) (1.325,5.638) (1.275,5.662) (1.225,5.686) (1.175,5.708) (1.125,5.729) (1.075,5.749) (1.025,5.769) (0.975,5.787) (0.925,5.804) (0.875,5.819) (0.825,5.834) (0.775,5.848) (0.725,5.862) (0.675,5.874) (0.625,5.885) (0.575,5.895) (0.525,5.905) (0.475,5.913) (0.425,5.921) (0.375,5.928) (0.325,5.934) (0.275,5.939) (0.225,5.943) (0.175,5.946) (0.125,5.949) (0.075,5.950) (0.025,5.951) (-0.050,5.951) (-0.100,5.950) (-0.150,5.948) (-0.200,5.945) (-0.250,5.941) (-0.300,5.936) (-0.350,5.931) (-0.400,5.924) (-0.450,5.917) (-0.500,5.909) (-0.550,5.900) (-0.600,5.890) (-0.650,5.879) (-0.700,5.868) (-0.750,5.855) (-0.800,5.842) (-0.850,5.827) (-0.900,5.812) (-0.950,5.795) (-1.000,5.778) (-1.050,5.759) (-1.100,5.740) (-1.150,5.719) (-1.200,5.697) (-1.249,5.675) (-1.300,5.650) (-1.350,5.625) (-1.397,5.600) (-1.442,5.575) (-1.486,5.550) (-1.550,5.511) (-1.600,5.479) (-1.643,5.450) (-1.700,5.410) (-1.748,5.375) (-1.800,5.335) (-1.843,5.300) (-1.900,5.252) (-1.950,5.207) (-2.000,5.160) (-2.036,5.125) (-2.075,5.085) (-2.125,5.031) (-2.175,4.975) (-2.217,4.925) (-2.250,4.883) (-2.294,4.825) (-2.325,4.782) (-2.364,4.725) (-2.397,4.675) (-2.425,4.629) (-2.457,4.575) (-2.484,4.525) (-2.511,4.475) (-2.535,4.425) (-2.559,4.375) (-2.581,4.325) (-2.602,4.275) (-2.621,4.225) (-2.639,4.175) (-2.656,4.125) (-2.672,4.075) (-2.687,4.025) (-2.700,3.975) (-2.712,3.925) (-2.723,3.875) (-2.733,3.825) (-2.742,3.775) (-2.750,3.725) (-2.757,3.675) (-2.762,3.625) (-2.766,3.575) (-2.770,3.525) (-2.772,3.475) (-2.773,3.425) (-2.773,3.375) (-2.772,3.325) (-2.770,3.275) (-2.767,3.225) (-2.763,3.175) (-2.757,3.125) (-2.751,3.075) (-2.743,3.025) (-2.735,2.975) (-2.725,2.926) (-2.714,2.875) (-2.702,2.825) (-2.689,2.775) (-2.674,2.725) (-2.659,2.675) (-2.642,2.625) (-2.625,2.577) (-2.605,2.525) (-2.585,2.475) (-2.563,2.425) (-2.541,2.375) (-2.517,2.325) (-2.491,2.275) (-2.464,2.225) (-2.436,2.175) (-2.407,2.125) (-2.376,2.075) (-2.344,2.025) (-2.311,1.975) (-2.276,1.925) (-2.240,1.875) (-2.202,1.825) (-2.162,1.775) (-2.125,1.729) (-2.079,1.675) (-2.036,1.625) (-2.000,1.586) (-1.950,1.532) (-1.900,1.480) (-1.850,1.429) (-1.800,1.380) (-1.750,1.333) (-1.700,1.286) (-1.661,1.250) (-1.604,1.200) (-1.550,1.153) (-1.500,1.110) (-1.458,1.075) (-1.400,1.027) (-1.350,0.987) (-1.304,0.950) (-1.250,0.907) (-1.209,0.875) (-1.150,0.830) (-1.100,0.792) (-1.050,0.754) (-1.000,0.717) (-0.950,0.680) (-0.900,0.643) (-0.850,0.606) (-0.808,0.575) (-0.750,0.534) (-0.704,0.500) (-0.650,0.462) (-0.600,0.425) (-0.550,0.390) (-0.500,0.354) (-0.450,0.319) (-0.400,0.284) (-0.354,0.250) (-0.300,0.213) (-0.250,0.177) (-0.220,0.155)}
\EQ{(-0.220,-0.155) (-0.275,-0.195) (-0.319,-0.225) (-0.375,-0.266) (-0.424,-0.300) (-0.475,-0.337) (-0.525,-0.372) (-0.575,-0.408) (-0.625,-0.444) (-0.669,-0.475) (-0.725,-0.516) (-0.773,-0.550) (-0.825,-0.589) (-0.875,-0.625) (-0.925,-0.662) (-0.975,-0.699) (-1.025,-0.737) (-1.075,-0.774) (-1.125,-0.813) (-1.173,-0.850) (-1.225,-0.891) (-1.269,-0.925) (-1.325,-0.970) (-1.375,-1.011) (-1.422,-1.050) (-1.475,-1.095) (-1.525,-1.139) (-1.567,-1.175) (-1.622,-1.225) (-1.675,-1.274) (-1.725,-1.322) (-1.775,-1.370) (-1.825,-1.421) (-1.875,-1.473) (-1.924,-1.525) (-1.969,-1.575) (-2.012,-1.625) (-2.050,-1.671) (-2.094,-1.725) (-2.125,-1.766) (-2.170,-1.825) (-2.200,-1.868) (-2.239,-1.925) (-2.272,-1.975) (-2.300,-2.020) (-2.333,-2.075) (-2.362,-2.125) (-2.389,-2.175) (-2.414,-2.225) (-2.439,-2.275) (-2.462,-2.325) (-2.483,-2.375) (-2.504,-2.425) (-2.523,-2.475) (-2.541,-2.525) (-2.558,-2.575) (-2.573,-2.625) (-2.587,-2.675) (-2.600,-2.725) (-2.612,-2.775) (-2.623,-2.825) (-2.632,-2.875) (-2.640,-2.925) (-2.647,-2.975) (-2.654,-3.025) (-2.658,-3.075) (-2.662,-3.125) (-2.665,-3.175) (-2.666,-3.225) (-2.667,-3.275) (-2.666,-3.325) (-2.664,-3.375) (-2.661,-3.425) (-2.657,-3.475) (-2.651,-3.525) (-2.645,-3.575) (-2.637,-3.625) (-2.629,-3.675) (-2.619,-3.725) (-2.608,-3.775) (-2.595,-3.825) (-2.582,-3.875) (-2.567,-3.925) (-2.551,-3.975) (-2.534,-4.025) (-2.515,-4.075) (-2.496,-4.125) (-2.475,-4.174) (-2.452,-4.225) (-2.428,-4.275) (-2.402,-4.325) (-2.375,-4.375) (-2.350,-4.420) (-2.317,-4.475) (-2.285,-4.525) (-2.251,-4.575) (-2.216,-4.625) (-2.179,-4.675) (-2.140,-4.725) (-2.100,-4.773) (-2.055,-4.825) (-2.009,-4.875) (-1.960,-4.925) (-1.909,-4.975) (-1.855,-5.025) (-1.800,-5.073) (-1.750,-5.114) (-1.704,-5.150) (-1.650,-5.191) (-1.602,-5.225) (-1.550,-5.260) (-1.500,-5.293) (-1.450,-5.323) (-1.405,-5.350) (-1.361,-5.375) (-1.314,-5.400) (-1.265,-5.425) (-1.214,-5.450) (-1.160,-5.475) (-1.103,-5.500) (-1.050,-5.522) (-1.000,-5.541) (-0.950,-5.559) (-0.900,-5.577) (-0.850,-5.593) (-0.800,-5.608) (-0.750,-5.622) (-0.700,-5.635) (-0.650,-5.647) (-0.600,-5.659) (-0.550,-5.669) (-0.500,-5.678) (-0.450,-5.687) (-0.400,-5.694) (-0.350,-5.701) (-0.300,-5.707) (-0.250,-5.711) (-0.200,-5.715) (-0.150,-5.718) (-0.100,-5.721) (-0.050,-5.722) (0.000,-5.722) (0.050,-5.722) (0.100,-5.721) (0.150,-5.718) (0.200,-5.715) (0.250,-5.711) (0.300,-5.707) (0.350,-5.701) (0.400,-5.694) (0.450,-5.687) (0.500,-5.678) (0.550,-5.669) (0.600,-5.659) (0.650,-5.647) (0.700,-5.635) (0.750,-5.622) (0.800,-5.608) (0.850,-5.593) (0.900,-5.577) (0.950,-5.559) (1.000,-5.541) (1.050,-5.522) (1.100,-5.501) (1.150,-5.480) (1.200,-5.457) (1.250,-5.433) (1.300,-5.407) (1.350,-5.381) (1.400,-5.353) (1.447,-5.325) (1.500,-5.293) (1.550,-5.260) (1.600,-5.226) (1.650,-5.191) (1.700,-5.153) (1.750,-5.114) (1.797,-5.075) (1.850,-5.029) (1.900,-4.983) (1.950,-4.935) (2.000,-4.884) (2.050,-4.830) (2.098,-4.775) (2.140,-4.725) (2.175,-4.680) (2.216,-4.625) (2.250,-4.577) (2.285,-4.525) (2.317,-4.475) (2.347,-4.425) (2.375,-4.376) (2.400,-4.330) (2.425,-4.281) (2.450,-4.229) (2.474,-4.175) (2.496,-4.125) (2.515,-4.075) (2.534,-4.025) (2.551,-3.975) (2.567,-3.925) (2.582,-3.875) (2.595,-3.825) (2.608,-3.775) (2.619,-3.725) (2.629,-3.675) (2.637,-3.625) (2.645,-3.575) (2.651,-3.525) (2.657,-3.475) (2.661,-3.425) (2.664,-3.375) (2.666,-3.325) (2.667,-3.275) (2.666,-3.225) (2.665,-3.175) (2.662,-3.125) (2.658,-3.075) (2.654,-3.025) (2.647,-2.975) (2.640,-2.925) (2.632,-2.875) (2.623,-2.825) (2.612,-2.775) (2.600,-2.725) (2.587,-2.675) (2.573,-2.625) (2.558,-2.575) (2.541,-2.525) (2.523,-2.475) (2.504,-2.425) (2.483,-2.375) (2.462,-2.325) (2.439,-2.275) (2.414,-2.225) (2.389,-2.175) (2.362,-2.125) (2.333,-2.075) (2.303,-2.025) (2.275,-1.980) (2.239,-1.925) (2.205,-1.875) (2.175,-1.833) (2.132,-1.775) (2.100,-1.733) (2.053,-1.675) (2.012,-1.625) (1.975,-1.583) (1.925,-1.527) (1.877,-1.475) (1.829,-1.425) (1.780,-1.375) (1.729,-1.325) (1.676,-1.275) (1.625,-1.228) (1.575,-1.183) (1.525,-1.139) (1.481,-1.100) (1.425,-1.053) (1.375,-1.011) (1.331,-0.975) (1.275,-0.930) (1.225,-0.891) (1.175,-0.851) (1.125,-0.813) (1.076,-0.775) (1.025,-0.737) (0.976,-0.700) (0.925,-0.662) (0.875,-0.625) (0.825,-0.589) (0.775,-0.552) (0.725,-0.516) (0.675,-0.480) (0.635,-0.450) (0.575,-0.408) (0.530,-0.375) (0.475,-0.337) (0.425,-0.301) (0.375,-0.266) (0.325,-0.230) (0.284,-0.200) (0.225,-0.160) (0.206,-0.144)}
\EQ{(0.206,0.144) (0.248,0.175) (0.300,0.213) (0.350,0.248) (0.400,0.284) (0.450,0.319) (0.495,0.350) (0.550,0.390) (0.599,0.425) (0.650,0.462) (0.700,0.498) (0.750,0.534) (0.800,0.570) (0.842,0.600) (0.900,0.644) (0.943,0.675) (1.000,0.718) (1.043,0.750) (1.100,0.794) (1.141,0.825) (1.200,0.871) (1.250,0.910) (1.300,0.950) (1.350,0.991) (1.392,1.025) (1.450,1.074) (1.500,1.117) (1.538,1.150) (1.595,1.200) (1.649,1.250) (1.700,1.298) (1.750,1.346) (1.800,1.395) (1.850,1.447) (1.900,1.499) (1.946,1.550) (1.990,1.600) (2.025,1.641) (2.074,1.700) (2.113,1.750) (2.150,1.799) (2.188,1.850) (2.222,1.900) (2.256,1.950) (2.288,2.000) (2.318,2.050) (2.348,2.100) (2.375,2.149) (2.400,2.197) (2.425,2.247) (2.450,2.299) (2.473,2.350) (2.494,2.400) (2.514,2.450) (2.532,2.500) (2.549,2.550) (2.565,2.600) (2.580,2.650) (2.594,2.700) (2.606,2.750) (2.617,2.800) (2.627,2.850) (2.636,2.900) (2.644,2.950) (2.651,3.000) (2.656,3.050) (2.660,3.100) (2.664,3.150) (2.666,3.200) (2.667,3.250) (2.666,3.300) (2.665,3.350) (2.663,3.400) (2.659,3.450) (2.654,3.500) (2.648,3.550) (2.641,3.600) (2.633,3.650) (2.624,3.700) (2.613,3.750) (2.602,3.800) (2.589,3.850) (2.575,3.899) (2.559,3.950) (2.543,4.000) (2.525,4.050) (2.506,4.100) (2.485,4.150) (2.463,4.200) (2.440,4.250) (2.415,4.300) (2.389,4.350) (2.361,4.400) (2.332,4.450) (2.301,4.500) (2.268,4.550) (2.234,4.600) (2.200,4.647) (2.160,4.700) (2.125,4.743) (2.077,4.800) (2.032,4.850) (1.985,4.900) (1.935,4.950) (1.882,5.000) (1.826,5.050) (1.775,5.094) (1.725,5.134) (1.675,5.172) (1.625,5.209) (1.575,5.244) (1.528,5.275) (1.475,5.308) (1.425,5.338) (1.375,5.367) (1.325,5.394) (1.275,5.420) (1.225,5.445) (1.175,5.468) (1.125,5.491) (1.075,5.512) (1.025,5.532) (0.976,5.550) (0.925,5.568) (0.875,5.585) (0.826,5.600) (0.775,5.615) (0.725,5.629) (0.675,5.641) (0.625,5.653) (0.575,5.664) (0.525,5.674) (0.475,5.683) (0.425,5.691) (0.375,5.698) (0.325,5.704) (0.275,5.709) (0.225,5.713) (0.175,5.717) (0.125,5.720) (0.075,5.721) (0.025,5.722) (-0.050,5.722) (-0.100,5.721) (-0.150,5.718) (-0.200,5.715) (-0.250,5.711) (-0.300,5.707) (-0.350,5.701) (-0.400,5.694) (-0.450,5.687) (-0.500,5.678) (-0.550,5.669) (-0.600,5.659) (-0.650,5.647) (-0.700,5.635) (-0.750,5.622) (-0.800,5.608) (-0.850,5.593) (-0.900,5.577) (-0.950,5.559) (-1.000,5.541) (-1.050,5.522) (-1.100,5.501) (-1.150,5.480) (-1.200,5.457) (-1.250,5.433) (-1.300,5.407) (-1.350,5.381) (-1.400,5.353) (-1.447,5.325) (-1.500,5.293) (-1.550,5.260) (-1.600,5.226) (-1.650,5.191) (-1.700,5.153) (-1.750,5.114) (-1.797,5.075) (-1.850,5.029) (-1.900,4.983) (-1.950,4.935) (-2.000,4.884) (-2.050,4.830) (-2.098,4.775) (-2.140,4.725) (-2.175,4.680) (-2.216,4.625) (-2.250,4.577) (-2.285,4.525) (-2.317,4.475) (-2.347,4.425) (-2.375,4.376) (-2.400,4.330) (-2.425,4.281) (-2.450,4.229) (-2.474,4.175) (-2.496,4.125) (-2.515,4.075) (-2.534,4.025) (-2.551,3.975) (-2.567,3.925) (-2.582,3.875) (-2.595,3.825) (-2.608,3.775) (-2.619,3.725) (-2.629,3.675) (-2.637,3.625) (-2.645,3.575) (-2.651,3.525) (-2.657,3.475) (-2.661,3.425) (-2.664,3.375) (-2.666,3.325) (-2.667,3.275) (-2.666,3.225) (-2.665,3.175) (-2.662,3.125) (-2.658,3.075) (-2.654,3.025) (-2.647,2.975) (-2.640,2.925) (-2.632,2.875) (-2.623,2.825) (-2.612,2.775) (-2.600,2.725) (-2.587,2.675) (-2.573,2.625) (-2.558,2.575) (-2.541,2.525) (-2.523,2.475) (-2.504,2.425) (-2.483,2.375) (-2.462,2.325) (-2.439,2.275) (-2.414,2.225) (-2.389,2.175) (-2.362,2.125) (-2.333,2.075) (-2.303,2.025) (-2.275,1.980) (-2.239,1.925) (-2.205,1.875) (-2.175,1.833) (-2.132,1.775) (-2.100,1.733) (-2.053,1.675) (-2.012,1.625) (-1.975,1.583) (-1.925,1.527) (-1.877,1.475) (-1.829,1.425) (-1.780,1.375) (-1.729,1.325) (-1.676,1.275) (-1.625,1.228) (-1.575,1.183) (-1.525,1.139) (-1.481,1.100) (-1.425,1.053) (-1.375,1.011) (-1.331,0.975) (-1.275,0.930) (-1.225,0.891) (-1.175,0.851) (-1.125,0.813) (-1.076,0.775) (-1.025,0.737) (-0.976,0.700) (-0.925,0.662) (-0.875,0.625) (-0.825,0.589) (-0.775,0.552) (-0.725,0.516) (-0.675,0.480) (-0.635,0.450) (-0.575,0.408) (-0.530,0.375) (-0.475,0.337) (-0.425,0.301) (-0.375,0.266) (-0.325,0.230) (-0.284,0.200) (-0.225,0.160) (-0.220,0.155)}
\EQ{(-0.220,-0.155) (-0.275,-0.195) (-0.319,-0.225) (-0.375,-0.266) (-0.424,-0.300) (-0.475,-0.337) (-0.525,-0.372) (-0.575,-0.408) (-0.625,-0.444) (-0.669,-0.475) (-0.725,-0.516) (-0.772,-0.550) (-0.825,-0.589) (-0.874,-0.625) (-0.925,-0.663) (-0.975,-0.700) (-1.025,-0.738) (-1.074,-0.775) (-1.125,-0.815) (-1.171,-0.850) (-1.225,-0.893) (-1.265,-0.925) (-1.325,-0.974) (-1.375,-1.016) (-1.416,-1.050) (-1.473,-1.100) (-1.525,-1.146) (-1.575,-1.191) (-1.612,-1.225) (-1.650,-1.261) (-1.691,-1.300) (-1.741,-1.350) (-1.790,-1.400) (-1.825,-1.438) (-1.860,-1.475) (-1.900,-1.520) (-1.947,-1.575) (-1.989,-1.625) (-2.025,-1.670) (-2.067,-1.725) (-2.100,-1.770) (-2.139,-1.825) (-2.173,-1.875) (-2.205,-1.925) (-2.236,-1.975) (-2.265,-2.025) (-2.293,-2.075) (-2.319,-2.125) (-2.344,-2.175) (-2.368,-2.225) (-2.390,-2.275) (-2.411,-2.325) (-2.431,-2.375) (-2.449,-2.425) (-2.466,-2.475) (-2.481,-2.525) (-2.496,-2.575) (-2.509,-2.625) (-2.521,-2.675) (-2.532,-2.725) (-2.541,-2.775) (-2.549,-2.825) (-2.557,-2.875) (-2.562,-2.925) (-2.567,-2.975) (-2.571,-3.025) (-2.573,-3.075) (-2.574,-3.125) (-2.574,-3.175) (-2.573,-3.225) (-2.571,-3.275) (-2.568,-3.325) (-2.563,-3.375) (-2.557,-3.425) (-2.550,-3.475) (-2.542,-3.525) (-2.533,-3.575) (-2.522,-3.625) (-2.510,-3.675) (-2.497,-3.725) (-2.483,-3.775) (-2.467,-3.825) (-2.450,-3.875) (-2.432,-3.925) (-2.413,-3.975) (-2.392,-4.025) (-2.369,-4.075) (-2.345,-4.125) (-2.320,-4.175) (-2.293,-4.225) (-2.264,-4.275) (-2.234,-4.325) (-2.202,-4.375) (-2.168,-4.425) (-2.133,-4.475) (-2.100,-4.518) (-2.055,-4.575) (-2.013,-4.625) (-1.975,-4.668) (-1.925,-4.722) (-1.875,-4.772) (-1.825,-4.820) (-1.775,-4.866) (-1.735,-4.900) (-1.675,-4.949) (-1.625,-4.988) (-1.575,-5.025) (-1.525,-5.060) (-1.475,-5.093) (-1.425,-5.124) (-1.375,-5.155) (-1.325,-5.183) (-1.275,-5.210) (-1.225,-5.236) (-1.175,-5.261) (-1.125,-5.284) (-1.075,-5.306) (-1.025,-5.327) (-0.975,-5.346) (-0.925,-5.365) (-0.875,-5.382) (-0.825,-5.398) (-0.775,-5.414) (-0.725,-5.428) (-0.675,-5.441) (-0.625,-5.453) (-0.575,-5.464) (-0.525,-5.474) (-0.475,-5.484) (-0.425,-5.492) (-0.375,-5.499) (-0.325,-5.506) (-0.275,-5.511) (-0.225,-5.516) (-0.175,-5.519) (-0.125,-5.522) (-0.075,-5.524) (-0.025,-5.525) (0.050,-5.524) (0.100,-5.523) (0.150,-5.521) (0.200,-5.518) (0.250,-5.513) (0.300,-5.508) (0.350,-5.503) (0.400,-5.496) (0.450,-5.488) (0.500,-5.479) (0.550,-5.469) (0.600,-5.459) (0.650,-5.447) (0.700,-5.434) (0.750,-5.421) (0.800,-5.406) (0.850,-5.390) (0.900,-5.374) (0.950,-5.356) (1.000,-5.337) (1.050,-5.316) (1.100,-5.295) (1.150,-5.272) (1.197,-5.250) (1.247,-5.225) (1.294,-5.200) (1.340,-5.175) (1.400,-5.140) (1.450,-5.109) (1.500,-5.076) (1.550,-5.042) (1.600,-5.006) (1.642,-4.975) (1.700,-4.929) (1.750,-4.887) (1.792,-4.850) (1.846,-4.800) (1.898,-4.750) (1.946,-4.700) (1.991,-4.650) (2.025,-4.611) (2.075,-4.550) (2.114,-4.500) (2.150,-4.451) (2.185,-4.400) (2.218,-4.350) (2.249,-4.300) (2.275,-4.257) (2.300,-4.212) (2.325,-4.165) (2.350,-4.115) (2.375,-4.062) (2.400,-4.005) (2.423,-3.950) (2.442,-3.900) (2.459,-3.850) (2.475,-3.801) (2.490,-3.750) (2.504,-3.700) (2.516,-3.650) (2.528,-3.600) (2.538,-3.550) (2.546,-3.500) (2.554,-3.450) (2.560,-3.400) (2.566,-3.350) (2.570,-3.300) (2.572,-3.250) (2.574,-3.200) (2.575,-3.150) (2.574,-3.100) (2.572,-3.050) (2.569,-3.000) (2.565,-2.950) (2.560,-2.900) (2.553,-2.850) (2.545,-2.800) (2.537,-2.750) (2.526,-2.700) (2.515,-2.650) (2.502,-2.600) (2.489,-2.550) (2.474,-2.500) (2.457,-2.450) (2.440,-2.400) (2.421,-2.350) (2.401,-2.300) (2.379,-2.250) (2.356,-2.200) (2.332,-2.150) (2.306,-2.100) (2.279,-2.050) (2.251,-2.000) (2.221,-1.950) (2.189,-1.900) (2.156,-1.850) (2.125,-1.805) (2.086,-1.750) (2.050,-1.702) (2.009,-1.650) (1.975,-1.608) (1.926,-1.550) (1.882,-1.500) (1.837,-1.450) (1.800,-1.411) (1.765,-1.375) (1.725,-1.334) (1.675,-1.285) (1.639,-1.250) (1.600,-1.214) (1.558,-1.175) (1.502,-1.125) (1.450,-1.080) (1.400,-1.037) (1.356,-1.000) (1.300,-0.954) (1.250,-0.913) (1.202,-0.875) (1.150,-0.834) (1.107,-0.800) (1.050,-0.757) (1.008,-0.725) (0.950,-0.682) (0.908,-0.650) (0.850,-0.607) (0.807,-0.575) (0.750,-0.534) (0.703,-0.500) (0.650,-0.462) (0.600,-0.426) (0.550,-0.390) (0.500,-0.355) (0.450,-0.319) (0.400,-0.284) (0.354,-0.250) (0.300,-0.213) (0.250,-0.177) (0.206,-0.144)}
\EQ{(0.206,0.144) (0.248,0.175) (0.300,0.213) (0.350,0.248) (0.400,0.284) (0.450,0.319) (0.495,0.350) (0.550,0.390) (0.599,0.425) (0.650,0.462) (0.700,0.498) (0.750,0.534) (0.800,0.571) (0.841,0.600) (0.900,0.644) (0.942,0.675) (1.000,0.719) (1.041,0.750) (1.100,0.795) (1.150,0.834) (1.200,0.873) (1.250,0.913) (1.296,0.950) (1.350,0.995) (1.400,1.037) (1.445,1.075) (1.500,1.123) (1.550,1.168) (1.600,1.214) (1.639,1.250) (1.675,1.285) (1.716,1.325) (1.765,1.375) (1.800,1.411) (1.837,1.450) (1.875,1.492) (1.925,1.549) (1.968,1.600) (2.000,1.639) (2.048,1.700) (2.086,1.750) (2.122,1.800) (2.156,1.850) (2.189,1.900) (2.221,1.950) (2.250,1.999) (2.275,2.043) (2.300,2.088) (2.325,2.136) (2.350,2.187) (2.375,2.241) (2.400,2.299) (2.421,2.350) (2.440,2.400) (2.457,2.450) (2.474,2.500) (2.489,2.550) (2.502,2.600) (2.515,2.650) (2.526,2.700) (2.537,2.750) (2.545,2.800) (2.553,2.850) (2.560,2.900) (2.565,2.950) (2.569,3.000) (2.572,3.050) (2.574,3.100) (2.575,3.150) (2.574,3.200) (2.572,3.250) (2.570,3.300) (2.566,3.350) (2.560,3.400) (2.554,3.450) (2.546,3.500) (2.538,3.550) (2.528,3.600) (2.516,3.650) (2.504,3.700) (2.490,3.750) (2.475,3.800) (2.459,3.850) (2.442,3.900) (2.423,3.950) (2.402,4.000) (2.381,4.050) (2.357,4.100) (2.333,4.150) (2.307,4.200) (2.279,4.250) (2.250,4.299) (2.218,4.350) (2.185,4.400) (2.151,4.450) (2.114,4.500) (2.075,4.550) (2.034,4.600) (2.000,4.640) (1.950,4.695) (1.900,4.747) (1.850,4.797) (1.800,4.843) (1.750,4.887) (1.705,4.925) (1.650,4.969) (1.609,5.000) (1.550,5.042) (1.502,5.075) (1.450,5.109) (1.400,5.140) (1.350,5.169) (1.300,5.197) (1.250,5.223) (1.200,5.249) (1.150,5.272) (1.100,5.295) (1.050,5.316) (1.000,5.337) (0.950,5.356) (0.900,5.374) (0.850,5.390) (0.800,5.406) (0.750,5.421) (0.700,5.434) (0.650,5.447) (0.600,5.459) (0.550,5.469) (0.500,5.479) (0.450,5.488) (0.400,5.496) (0.350,5.503) (0.300,5.508) (0.250,5.513) (0.200,5.518) (0.150,5.521) (0.100,5.523) (0.050,5.524) (0.000,5.525) (-0.050,5.524) (-0.100,5.523) (-0.150,5.521) (-0.200,5.518) (-0.250,5.513) (-0.300,5.508) (-0.350,5.503) (-0.400,5.496) (-0.450,5.488) (-0.500,5.479) (-0.550,5.469) (-0.600,5.459) (-0.650,5.447) (-0.700,5.434) (-0.750,5.421) (-0.800,5.406) (-0.850,5.390) (-0.900,5.374) (-0.950,5.356) (-1.000,5.337) (-1.050,5.316) (-1.100,5.295) (-1.150,5.272) (-1.197,5.250) (-1.247,5.225) (-1.294,5.200) (-1.340,5.175) (-1.400,5.140) (-1.450,5.109) (-1.500,5.076) (-1.550,5.042) (-1.600,5.006) (-1.642,4.975) (-1.700,4.929) (-1.750,4.887) (-1.792,4.850) (-1.846,4.800) (-1.898,4.750) (-1.946,4.700) (-1.991,4.650) (-2.025,4.611) (-2.075,4.550) (-2.114,4.500) (-2.150,4.451) (-2.185,4.400) (-2.218,4.350) (-2.249,4.300) (-2.275,4.257) (-2.300,4.212) (-2.325,4.165) (-2.350,4.115) (-2.375,4.062) (-2.400,4.005) (-2.423,3.950) (-2.442,3.900) (-2.459,3.850) (-2.475,3.801) (-2.490,3.750) (-2.504,3.700) (-2.516,3.650) (-2.528,3.600) (-2.538,3.550) (-2.546,3.500) (-2.554,3.450) (-2.560,3.400) (-2.566,3.350) (-2.570,3.300) (-2.572,3.250) (-2.574,3.200) (-2.575,3.150) (-2.574,3.100) (-2.572,3.050) (-2.569,3.000) (-2.565,2.950) (-2.560,2.900) (-2.553,2.850) (-2.545,2.800) (-2.537,2.750) (-2.526,2.700) (-2.515,2.650) (-2.502,2.600) (-2.489,2.550) (-2.474,2.500) (-2.457,2.450) (-2.440,2.400) (-2.421,2.350) (-2.401,2.300) (-2.379,2.250) (-2.356,2.200) (-2.332,2.150) (-2.306,2.100) (-2.279,2.050) (-2.251,2.000) (-2.221,1.950) (-2.189,1.900) (-2.156,1.850) (-2.125,1.805) (-2.086,1.750) (-2.050,1.702) (-2.009,1.650) (-1.975,1.608) (-1.926,1.550) (-1.882,1.500) (-1.837,1.450) (-1.800,1.411) (-1.765,1.375) (-1.725,1.334) (-1.675,1.285) (-1.639,1.250) (-1.600,1.214) (-1.558,1.175) (-1.502,1.125) (-1.450,1.080) (-1.400,1.037) (-1.356,1.000) (-1.300,0.954) (-1.250,0.913) (-1.202,0.875) (-1.150,0.834) (-1.107,0.800) (-1.050,0.757) (-1.008,0.725) (-0.950,0.682) (-0.908,0.650) (-0.850,0.607) (-0.807,0.575) (-0.750,0.534) (-0.703,0.500) (-0.650,0.462) (-0.600,0.426) (-0.550,0.390) (-0.500,0.355) (-0.450,0.319) (-0.400,0.284) (-0.354,0.250) (-0.300,0.213) (-0.250,0.177) (-0.220,0.155)}
}
  \marcoMMA{10}{5}{$x$ (u.\,arb.)}{$z$ (u.\,arb.)}
\end{tikzpicture}
\end{minipage}
\caption{Equipotenciales, a intervalos iguales de potencial, de los tres primeros términos del desarrollo multipolar con simetría de revolución alrededor de $z$:
monopolar (esferas $r=\text{cte}$, que se apiñan hacia el centro porque $\phi\propto1/r$), dipolar ($r^{2}\propto\cos\theta$, dos lóbulos de signo opuesto separados por
el plano $z=0$) y cuadrupolar ($r^{3}\propto3\cos^{2}\theta-1$, cuatro lóbulos separados por el cono nodal $\cos^{2}\theta=1/3$, las rectas oblicuas, hacia el que convergen las curvas; la región en blanco, con forma de trébol, es la de potencial más intenso que el último nivel trazado). Recreación
de~\cite{FIGURAS-T2}.}
\label{fig:equipotenciales}
\end{figure}


La \cref{fig:equipotenciales} compara las equipotenciales de los tres primeros términos: esferas para el monopolo, los dos lóbulos del dipolo y los cuatro
del cuadrupolo, separados por el cono nodal. La \cref{fig:cuadrupolo-mma} completa la descripción del término cuadrupolar con sus líneas de campo y su campo
vectorial, y la \cref{fig:cuadrupolo-fp} lo compara con un cuadrupolo físico de tres cargas, que estudiamos a continuación.

% =========================================================================
% Figura (clon de «Tema 2. Figuras», diapositivas «Término cuadrupolar»):
% (a) equipotenciales, (b) equipotenciales y líneas de campo y (c) campo
% eléctrico del término cuadrupolar con simetría de revolución,
% phi_2 ~ Q_zz P2(cos t)/r^3, Q_zz > 0. Datos: figuras/datos/cua_puntual_*.
% =========================================================================
\begin{figure}[t]
\centering
\newcommand{\icono}{% dos dipolos opuestos unidos por su carga negativa
  \draw[-{Stealth[length=2.8pt,width=2.6pt]}, mmarojo, line width=0.9pt] (0,0) -- (0,0.95);
  \draw[-{Stealth[length=2.8pt,width=2.6pt]}, mmarojo, line width=0.9pt] (0,0) -- (0,-0.95);
  \fill[mmarojo] (0,0) circle (0.2);}
\begin{minipage}[t]{0.325\textwidth}\centering
{\footnotesize Equipotenciales}\\[3pt]
\begin{tikzpicture}[x=0.198cm, y=0.198cm]
  \renewcommand{\EQ}[1]{\draw[mmaEQsolo] plot coordinates {#1};}%
  \colorlet{marcoc}{grispizarra!60}
  \panelMMA{10}{white}{0}{% Cuadrupolo puntual - equipotenciales
% Generado por scripts/generar_datos.py (no editar a mano).
\EQ{(-0.220,-0.155) (-0.275,-0.195) (-0.319,-0.225) (-0.375,-0.266) (-0.425,-0.300) (-0.475,-0.336) (-0.525,-0.371) (-0.568,-0.400) (-0.625,-0.441) (-0.675,-0.475) (-0.725,-0.510) (-0.775,-0.544) (-0.821,-0.575) (-0.875,-0.612) (-0.925,-0.646) (-0.969,-0.675) (-1.025,-0.713) (-1.075,-0.745) (-1.121,-0.775) (-1.175,-0.810) (-1.225,-0.841) (-1.275,-0.872) (-1.321,-0.900) (-1.375,-0.933) (-1.425,-0.963) (-1.475,-0.991) (-1.525,-1.020) (-1.575,-1.047) (-1.625,-1.074) (-1.674,-1.100) (-1.723,-1.125) (-1.774,-1.150) (-1.825,-1.175) (-1.875,-1.198) (-1.925,-1.220) (-1.975,-1.241) (-2.025,-1.262) (-2.075,-1.281) (-2.125,-1.300) (-2.175,-1.317) (-2.225,-1.334) (-2.275,-1.349) (-2.325,-1.364) (-2.375,-1.377) (-2.425,-1.389) (-2.475,-1.400) (-2.525,-1.410) (-2.575,-1.419) (-2.625,-1.426) (-2.675,-1.433) (-2.725,-1.438) (-2.775,-1.441) (-2.825,-1.444) (-2.875,-1.445) (-2.925,-1.445) (-2.975,-1.443) (-3.025,-1.440) (-3.075,-1.435) (-3.125,-1.429) (-3.175,-1.421) (-3.225,-1.412) (-3.275,-1.401) (-3.325,-1.388) (-3.375,-1.373) (-3.425,-1.357) (-3.475,-1.338) (-3.525,-1.318) (-3.575,-1.295) (-3.625,-1.270) (-3.675,-1.242) (-3.725,-1.212) (-3.775,-1.179) (-3.825,-1.142) (-3.875,-1.102) (-3.925,-1.058) (-3.975,-1.010) (-4.025,-0.956) (-4.072,-0.900) (-4.110,-0.850) (-4.145,-0.800) (-4.175,-0.753) (-4.206,-0.700) (-4.232,-0.650) (-4.256,-0.600) (-4.277,-0.550) (-4.297,-0.500) (-4.314,-0.450) (-4.329,-0.400) (-4.343,-0.350) (-4.354,-0.300) (-4.364,-0.250) (-4.371,-0.200) (-4.377,-0.150) (-4.382,-0.100) (-4.384,-0.050) (-4.385,0.000) (-4.384,0.050) (-4.382,0.100) (-4.377,0.150) (-4.371,0.200) (-4.364,0.250) (-4.354,0.300) (-4.343,0.350) (-4.329,0.400) (-4.314,0.450) (-4.297,0.500) (-4.277,0.550) (-4.256,0.600) (-4.232,0.650) (-4.206,0.700) (-4.177,0.750) (-4.150,0.792) (-4.110,0.850) (-4.075,0.896) (-4.030,0.950) (-3.985,1.000) (-3.934,1.050) (-3.878,1.100) (-3.825,1.142) (-3.780,1.175) (-3.725,1.212) (-3.675,1.242) (-3.625,1.270) (-3.575,1.295) (-3.525,1.318) (-3.475,1.338) (-3.425,1.357) (-3.375,1.373) (-3.325,1.388) (-3.275,1.401) (-3.225,1.412) (-3.175,1.421) (-3.125,1.429) (-3.075,1.435) (-3.025,1.440) (-2.975,1.443) (-2.925,1.445) (-2.875,1.445) (-2.825,1.444) (-2.775,1.441) (-2.725,1.438) (-2.675,1.433) (-2.625,1.426) (-2.575,1.419) (-2.525,1.410) (-2.475,1.400) (-2.425,1.389) (-2.375,1.377) (-2.325,1.364) (-2.275,1.349) (-2.225,1.334) (-2.175,1.317) (-2.126,1.300) (-2.075,1.281) (-2.025,1.262) (-1.975,1.241) (-1.925,1.220) (-1.875,1.198) (-1.826,1.175) (-1.775,1.151) (-1.725,1.126) (-1.675,1.100) (-1.627,1.075) (-1.580,1.050) (-1.535,1.025) (-1.491,1.000) (-1.447,0.975) (-1.400,0.948) (-1.350,0.918) (-1.300,0.888) (-1.250,0.857) (-1.200,0.825) (-1.150,0.794) (-1.100,0.762) (-1.050,0.729) (-1.007,0.700) (-0.950,0.663) (-0.900,0.629) (-0.858,0.600) (-0.800,0.561) (-0.750,0.527) (-0.700,0.493) (-0.650,0.458) (-0.604,0.425) (-0.550,0.388) (-0.500,0.353) (-0.450,0.318) (-0.400,0.283) (-0.354,0.250) (-0.300,0.213) (-0.250,0.177) (-0.220,0.155)}
\EQ{(0.206,0.144) (0.248,0.175) (0.300,0.213) (0.350,0.248) (0.400,0.283) (0.450,0.318) (0.496,0.350) (0.550,0.388) (0.600,0.423) (0.650,0.458) (0.700,0.493) (0.748,0.525) (0.800,0.561) (0.850,0.595) (0.895,0.625) (0.950,0.663) (1.000,0.696) (1.045,0.725) (1.100,0.762) (1.150,0.794) (1.200,0.825) (1.250,0.857) (1.300,0.888) (1.350,0.918) (1.400,0.948) (1.447,0.975) (1.491,1.000) (1.535,1.025) (1.580,1.050) (1.625,1.074) (1.674,1.100) (1.723,1.125) (1.774,1.150) (1.825,1.175) (1.875,1.198) (1.925,1.220) (1.975,1.241) (2.025,1.262) (2.075,1.281) (2.125,1.300) (2.175,1.317) (2.225,1.334) (2.275,1.349) (2.325,1.364) (2.375,1.377) (2.425,1.389) (2.475,1.400) (2.525,1.410) (2.575,1.419) (2.625,1.426) (2.675,1.433) (2.725,1.438) (2.775,1.441) (2.825,1.444) (2.875,1.445) (2.925,1.445) (2.975,1.443) (3.025,1.440) (3.075,1.435) (3.125,1.429) (3.175,1.421) (3.225,1.412) (3.275,1.401) (3.325,1.388) (3.375,1.373) (3.425,1.357) (3.475,1.338) (3.525,1.318) (3.575,1.295) (3.625,1.270) (3.675,1.242) (3.725,1.212) (3.775,1.179) (3.825,1.142) (3.875,1.102) (3.925,1.058) (3.975,1.010) (4.025,0.956) (4.072,0.900) (4.110,0.850) (4.145,0.800) (4.175,0.753) (4.206,0.700) (4.232,0.650) (4.256,0.600) (4.277,0.550) (4.297,0.500) (4.314,0.450) (4.329,0.400) (4.343,0.350) (4.354,0.300) (4.364,0.250) (4.371,0.200) (4.377,0.150) (4.382,0.100) (4.384,0.050) (4.385,0.000) (4.384,-0.050) (4.382,-0.100) (4.377,-0.150) (4.371,-0.200) (4.364,-0.250) (4.354,-0.300) (4.343,-0.350) (4.329,-0.400) (4.314,-0.450) (4.297,-0.500) (4.277,-0.550) (4.256,-0.600) (4.232,-0.650) (4.206,-0.700) (4.177,-0.750) (4.150,-0.792) (4.110,-0.850) (4.075,-0.896) (4.030,-0.950) (3.985,-1.000) (3.934,-1.050) (3.878,-1.100) (3.825,-1.142) (3.780,-1.175) (3.725,-1.212) (3.675,-1.242) (3.625,-1.270) (3.575,-1.295) (3.525,-1.318) (3.475,-1.338) (3.425,-1.357) (3.375,-1.373) (3.325,-1.388) (3.275,-1.401) (3.225,-1.412) (3.175,-1.421) (3.125,-1.429) (3.075,-1.435) (3.025,-1.440) (2.975,-1.443) (2.925,-1.445) (2.875,-1.445) (2.825,-1.444) (2.775,-1.441) (2.725,-1.438) (2.675,-1.433) (2.625,-1.426) (2.575,-1.419) (2.525,-1.410) (2.475,-1.400) (2.425,-1.389) (2.375,-1.377) (2.325,-1.364) (2.275,-1.349) (2.225,-1.334) (2.175,-1.317) (2.126,-1.300) (2.075,-1.281) (2.025,-1.262) (1.975,-1.241) (1.925,-1.220) (1.875,-1.198) (1.826,-1.175) (1.775,-1.151) (1.725,-1.126) (1.675,-1.100) (1.627,-1.075) (1.580,-1.050) (1.535,-1.025) (1.491,-1.000) (1.447,-0.975) (1.400,-0.948) (1.350,-0.918) (1.300,-0.888) (1.250,-0.857) (1.200,-0.825) (1.150,-0.794) (1.100,-0.762) (1.050,-0.729) (1.007,-0.700) (0.950,-0.663) (0.900,-0.629) (0.858,-0.600) (0.800,-0.561) (0.750,-0.527) (0.700,-0.493) (0.650,-0.458) (0.604,-0.425) (0.550,-0.388) (0.500,-0.353) (0.450,-0.318) (0.400,-0.283) (0.354,-0.250) (0.300,-0.213) (0.250,-0.177) (0.206,-0.144)}
\EQ{(-0.220,-0.155) (-0.275,-0.195) (-0.319,-0.225) (-0.375,-0.266) (-0.425,-0.300) (-0.475,-0.336) (-0.525,-0.371) (-0.568,-0.400) (-0.625,-0.441) (-0.675,-0.475) (-0.725,-0.510) (-0.775,-0.545) (-0.820,-0.575) (-0.875,-0.613) (-0.925,-0.647) (-0.968,-0.675) (-1.025,-0.714) (-1.075,-0.747) (-1.119,-0.775) (-1.175,-0.812) (-1.225,-0.844) (-1.275,-0.875) (-1.325,-0.906) (-1.375,-0.937) (-1.425,-0.967) (-1.475,-0.996) (-1.525,-1.025) (-1.569,-1.050) (-1.614,-1.075) (-1.660,-1.100) (-1.707,-1.125) (-1.755,-1.150) (-1.800,-1.173) (-1.850,-1.197) (-1.900,-1.221) (-1.950,-1.244) (-2.000,-1.266) (-2.050,-1.288) (-2.100,-1.308) (-2.150,-1.327) (-2.200,-1.346) (-2.250,-1.363) (-2.300,-1.380) (-2.350,-1.396) (-2.400,-1.410) (-2.450,-1.424) (-2.500,-1.436) (-2.550,-1.447) (-2.600,-1.458) (-2.650,-1.467) (-2.700,-1.475) (-2.750,-1.481) (-2.800,-1.487) (-2.850,-1.491) (-2.900,-1.494) (-2.950,-1.496) (-3.000,-1.497) (-3.050,-1.496) (-3.100,-1.493) (-3.150,-1.490) (-3.200,-1.485) (-3.250,-1.478) (-3.300,-1.470) (-3.350,-1.460) (-3.400,-1.449) (-3.450,-1.436) (-3.500,-1.421) (-3.550,-1.405) (-3.600,-1.386) (-3.650,-1.365) (-3.700,-1.343) (-3.750,-1.318) (-3.800,-1.290) (-3.850,-1.260) (-3.900,-1.228) (-3.950,-1.192) (-4.000,-1.153) (-4.050,-1.110) (-4.088,-1.075) (-4.125,-1.038) (-4.160,-1.000) (-4.200,-0.954) (-4.243,-0.900) (-4.275,-0.855) (-4.312,-0.800) (-4.342,-0.750) (-4.369,-0.700) (-4.395,-0.650) (-4.417,-0.600) (-4.438,-0.550) (-4.457,-0.500) (-4.473,-0.450) (-4.488,-0.400) (-4.501,-0.350) (-4.512,-0.300) (-4.521,-0.250) (-4.529,-0.200) (-4.534,-0.150) (-4.539,-0.100) (-4.541,-0.050) (-4.542,0.000) (-4.541,0.050) (-4.539,0.100) (-4.534,0.150) (-4.529,0.200) (-4.521,0.250) (-4.512,0.300) (-4.501,0.350) (-4.488,0.400) (-4.473,0.450) (-4.457,0.500) (-4.438,0.550) (-4.417,0.600) (-4.395,0.650) (-4.369,0.700) (-4.342,0.750) (-4.312,0.800) (-4.279,0.850) (-4.243,0.900) (-4.203,0.950) (-4.160,1.000) (-4.125,1.038) (-4.088,1.075) (-4.050,1.110) (-4.003,1.150) (-3.950,1.192) (-3.904,1.225) (-3.850,1.260) (-3.800,1.290) (-3.750,1.318) (-3.700,1.343) (-3.650,1.365) (-3.600,1.386) (-3.550,1.405) (-3.500,1.421) (-3.450,1.436) (-3.400,1.449) (-3.350,1.460) (-3.300,1.470) (-3.250,1.478) (-3.200,1.485) (-3.150,1.490) (-3.100,1.493) (-3.050,1.496) (-3.000,1.497) (-2.950,1.496) (-2.900,1.494) (-2.850,1.491) (-2.800,1.487) (-2.750,1.481) (-2.700,1.475) (-2.650,1.467) (-2.600,1.458) (-2.550,1.447) (-2.500,1.436) (-2.450,1.424) (-2.400,1.410) (-2.350,1.396) (-2.300,1.380) (-2.250,1.363) (-2.200,1.346) (-2.150,1.327) (-2.100,1.308) (-2.050,1.288) (-2.000,1.266) (-1.950,1.244) (-1.900,1.221) (-1.850,1.197) (-1.804,1.175) (-1.755,1.150) (-1.707,1.125) (-1.660,1.100) (-1.614,1.075) (-1.569,1.050) (-1.525,1.025) (-1.482,1.000) (-1.425,0.967) (-1.375,0.937) (-1.325,0.906) (-1.275,0.875) (-1.225,0.844) (-1.175,0.812) (-1.125,0.779) (-1.081,0.750) (-1.025,0.714) (-0.975,0.680) (-0.930,0.650) (-0.875,0.613) (-0.825,0.579) (-0.784,0.550) (-0.725,0.510) (-0.675,0.475) (-0.625,0.441) (-0.575,0.406) (-0.532,0.375) (-0.475,0.336) (-0.425,0.300) (-0.375,0.266) (-0.325,0.230) (-0.284,0.200) (-0.225,0.160) (-0.220,0.155)}
\EQ{(0.206,0.144) (0.248,0.175) (0.300,0.213) (0.350,0.248) (0.400,0.283) (0.450,0.318) (0.496,0.350) (0.550,0.388) (0.600,0.423) (0.650,0.458) (0.700,0.493) (0.747,0.525) (0.800,0.562) (0.850,0.596) (0.894,0.625) (0.950,0.664) (1.000,0.697) (1.043,0.725) (1.100,0.763) (1.150,0.796) (1.196,0.825) (1.250,0.859) (1.300,0.891) (1.350,0.921) (1.397,0.950) (1.450,0.982) (1.500,1.011) (1.550,1.040) (1.600,1.068) (1.650,1.095) (1.700,1.122) (1.750,1.148) (1.800,1.173) (1.850,1.197) (1.900,1.221) (1.950,1.244) (2.000,1.266) (2.050,1.288) (2.100,1.308) (2.150,1.327) (2.200,1.346) (2.250,1.363) (2.300,1.380) (2.350,1.396) (2.400,1.410) (2.450,1.424) (2.500,1.436) (2.550,1.447) (2.600,1.458) (2.650,1.467) (2.700,1.475) (2.750,1.481) (2.800,1.487) (2.850,1.491) (2.900,1.494) (2.950,1.496) (3.000,1.497) (3.050,1.496) (3.100,1.493) (3.150,1.490) (3.200,1.485) (3.250,1.478) (3.300,1.470) (3.350,1.460) (3.400,1.449) (3.450,1.436) (3.500,1.421) (3.550,1.405) (3.600,1.386) (3.650,1.365) (3.700,1.343) (3.750,1.318) (3.800,1.290) (3.850,1.260) (3.900,1.228) (3.950,1.192) (4.000,1.153) (4.050,1.110) (4.088,1.075) (4.125,1.038) (4.160,1.000) (4.200,0.954) (4.243,0.900) (4.275,0.855) (4.312,0.800) (4.342,0.750) (4.369,0.700) (4.395,0.650) (4.417,0.600) (4.438,0.550) (4.457,0.500) (4.473,0.450) (4.488,0.400) (4.501,0.350) (4.512,0.300) (4.521,0.250) (4.529,0.200) (4.534,0.150) (4.539,0.100) (4.541,0.050) (4.542,0.000) (4.541,-0.050) (4.539,-0.100) (4.534,-0.150) (4.529,-0.200) (4.521,-0.250) (4.512,-0.300) (4.501,-0.350) (4.488,-0.400) (4.473,-0.450) (4.457,-0.500) (4.438,-0.550) (4.417,-0.600) (4.395,-0.650) (4.369,-0.700) (4.342,-0.750) (4.312,-0.800) (4.279,-0.850) (4.243,-0.900) (4.203,-0.950) (4.160,-1.000) (4.125,-1.038) (4.088,-1.075) (4.050,-1.110) (4.003,-1.150) (3.950,-1.192) (3.904,-1.225) (3.850,-1.260) (3.800,-1.290) (3.750,-1.318) (3.700,-1.343) (3.650,-1.365) (3.600,-1.386) (3.550,-1.405) (3.500,-1.421) (3.450,-1.436) (3.400,-1.449) (3.350,-1.460) (3.300,-1.470) (3.250,-1.478) (3.200,-1.485) (3.150,-1.490) (3.100,-1.493) (3.050,-1.496) (3.000,-1.497) (2.950,-1.496) (2.900,-1.494) (2.850,-1.491) (2.800,-1.487) (2.750,-1.481) (2.700,-1.475) (2.650,-1.467) (2.600,-1.458) (2.550,-1.447) (2.500,-1.436) (2.450,-1.424) (2.400,-1.410) (2.350,-1.396) (2.300,-1.380) (2.250,-1.363) (2.200,-1.346) (2.150,-1.327) (2.100,-1.308) (2.050,-1.288) (2.000,-1.266) (1.950,-1.244) (1.900,-1.221) (1.850,-1.197) (1.804,-1.175) (1.755,-1.150) (1.707,-1.125) (1.660,-1.100) (1.614,-1.075) (1.569,-1.050) (1.525,-1.025) (1.482,-1.000) (1.425,-0.967) (1.375,-0.937) (1.325,-0.906) (1.275,-0.875) (1.225,-0.844) (1.175,-0.812) (1.125,-0.779) (1.081,-0.750) (1.025,-0.714) (0.975,-0.680) (0.930,-0.650) (0.875,-0.613) (0.825,-0.579) (0.784,-0.550) (0.725,-0.510) (0.675,-0.475) (0.625,-0.441) (0.575,-0.406) (0.532,-0.375) (0.475,-0.336) (0.425,-0.300) (0.375,-0.266) (0.325,-0.230) (0.284,-0.200) (0.225,-0.160) (0.206,-0.144)}
\EQ{(-0.220,-0.155) (-0.275,-0.195) (-0.319,-0.225) (-0.375,-0.266) (-0.425,-0.300) (-0.475,-0.336) (-0.525,-0.371) (-0.568,-0.400) (-0.625,-0.441) (-0.674,-0.475) (-0.725,-0.510) (-0.775,-0.545) (-0.819,-0.575) (-0.875,-0.614) (-0.925,-0.647) (-0.975,-0.681) (-1.025,-0.715) (-1.075,-0.748) (-1.125,-0.781) (-1.175,-0.814) (-1.225,-0.846) (-1.271,-0.875) (-1.325,-0.910) (-1.375,-0.941) (-1.425,-0.971) (-1.473,-1.000) (-1.525,-1.031) (-1.575,-1.060) (-1.625,-1.089) (-1.675,-1.116) (-1.725,-1.144) (-1.775,-1.170) (-1.825,-1.196) (-1.875,-1.221) (-1.925,-1.246) (-1.975,-1.270) (-2.025,-1.293) (-2.075,-1.315) (-2.125,-1.336) (-2.175,-1.357) (-2.222,-1.375) (-2.275,-1.395) (-2.325,-1.413) (-2.375,-1.430) (-2.425,-1.446) (-2.475,-1.461) (-2.525,-1.474) (-2.575,-1.487) (-2.625,-1.499) (-2.675,-1.510) (-2.725,-1.520) (-2.775,-1.529) (-2.825,-1.536) (-2.875,-1.542) (-2.925,-1.548) (-2.975,-1.552) (-3.025,-1.555) (-3.075,-1.556) (-3.125,-1.556) (-3.175,-1.556) (-3.225,-1.553) (-3.275,-1.550) (-3.325,-1.545) (-3.375,-1.538) (-3.425,-1.530) (-3.475,-1.521) (-3.525,-1.510) (-3.575,-1.497) (-3.625,-1.483) (-3.675,-1.467) (-3.725,-1.449) (-3.775,-1.429) (-3.825,-1.407) (-3.875,-1.383) (-3.925,-1.357) (-3.975,-1.329) (-4.021,-1.300) (-4.075,-1.264) (-4.125,-1.227) (-4.175,-1.187) (-4.217,-1.150) (-4.270,-1.100) (-4.318,-1.050) (-4.362,-1.000) (-4.400,-0.952) (-4.439,-0.900) (-4.473,-0.850) (-4.500,-0.807) (-4.533,-0.750) (-4.559,-0.700) (-4.583,-0.650) (-4.605,-0.600) (-4.624,-0.550) (-4.642,-0.500) (-4.658,-0.450) (-4.672,-0.400) (-4.684,-0.350) (-4.695,-0.300) (-4.704,-0.250) (-4.711,-0.200) (-4.717,-0.150) (-4.720,-0.100) (-4.723,-0.050) (-4.724,0.000) (-4.723,0.050) (-4.720,0.100) (-4.717,0.150) (-4.711,0.200) (-4.704,0.250) (-4.695,0.300) (-4.684,0.350) (-4.672,0.400) (-4.658,0.450) (-4.642,0.500) (-4.625,0.548) (-4.605,0.600) (-4.583,0.650) (-4.559,0.700) (-4.533,0.750) (-4.504,0.800) (-4.475,0.847) (-4.439,0.900) (-4.402,0.950) (-4.362,1.000) (-4.325,1.042) (-4.275,1.095) (-4.225,1.143) (-4.175,1.187) (-4.127,1.225) (-4.075,1.264) (-4.025,1.298) (-3.981,1.325) (-3.925,1.357) (-3.875,1.383) (-3.825,1.407) (-3.775,1.429) (-3.725,1.449) (-3.675,1.467) (-3.625,1.483) (-3.575,1.497) (-3.525,1.510) (-3.475,1.521) (-3.425,1.530) (-3.375,1.538) (-3.325,1.545) (-3.275,1.550) (-3.225,1.553) (-3.175,1.556) (-3.125,1.556) (-3.075,1.556) (-3.025,1.555) (-2.975,1.552) (-2.925,1.548) (-2.875,1.542) (-2.825,1.536) (-2.775,1.529) (-2.725,1.520) (-2.675,1.510) (-2.625,1.499) (-2.575,1.487) (-2.525,1.474) (-2.475,1.461) (-2.425,1.446) (-2.375,1.430) (-2.325,1.413) (-2.275,1.395) (-2.225,1.376) (-2.175,1.357) (-2.125,1.336) (-2.075,1.315) (-2.025,1.293) (-1.975,1.270) (-1.925,1.246) (-1.875,1.221) (-1.825,1.196) (-1.775,1.170) (-1.725,1.144) (-1.675,1.116) (-1.625,1.089) (-1.575,1.060) (-1.525,1.031) (-1.475,1.001) (-1.432,0.975) (-1.375,0.941) (-1.325,0.910) (-1.275,0.878) (-1.232,0.850) (-1.175,0.814) (-1.125,0.781) (-1.078,0.750) (-1.025,0.715) (-0.975,0.681) (-0.929,0.650) (-0.875,0.614) (-0.825,0.579) (-0.783,0.550) (-0.725,0.510) (-0.675,0.475) (-0.625,0.441) (-0.575,0.406) (-0.532,0.375) (-0.475,0.336) (-0.425,0.300) (-0.375,0.266) (-0.325,0.230) (-0.284,0.200) (-0.225,0.160) (-0.220,0.155)}
\EQ{(0.206,0.144) (0.248,0.175) (0.300,0.213) (0.350,0.248) (0.400,0.283) (0.450,0.318) (0.496,0.350) (0.550,0.389) (0.600,0.423) (0.650,0.458) (0.700,0.493) (0.747,0.525) (0.800,0.562) (0.850,0.596) (0.893,0.625) (0.950,0.665) (1.000,0.698) (1.050,0.732) (1.100,0.765) (1.150,0.797) (1.193,0.825) (1.250,0.862) (1.300,0.894) (1.350,0.925) (1.400,0.956) (1.450,0.986) (1.500,1.016) (1.550,1.046) (1.600,1.074) (1.646,1.100) (1.691,1.125) (1.737,1.150) (1.785,1.175) (1.833,1.200) (1.883,1.225) (1.934,1.250) (1.987,1.275) (2.042,1.300) (2.099,1.325) (2.150,1.346) (2.200,1.367) (2.250,1.386) (2.300,1.404) (2.350,1.421) (2.400,1.438) (2.450,1.453) (2.500,1.468) (2.550,1.481) (2.600,1.493) (2.650,1.505) (2.700,1.515) (2.750,1.524) (2.800,1.532) (2.850,1.539) (2.900,1.545) (2.950,1.550) (3.000,1.553) (3.050,1.556) (3.100,1.556) (3.150,1.556) (3.200,1.555) (3.250,1.552) (3.300,1.547) (3.350,1.542) (3.400,1.534) (3.450,1.526) (3.500,1.515) (3.550,1.504) (3.600,1.490) (3.650,1.475) (3.700,1.458) (3.750,1.439) (3.800,1.418) (3.850,1.396) (3.900,1.371) (3.950,1.343) (4.000,1.313) (4.050,1.281) (4.094,1.250) (4.150,1.207) (4.189,1.175) (4.244,1.125) (4.294,1.075) (4.340,1.025) (4.375,0.984) (4.421,0.925) (4.450,0.884) (4.489,0.825) (4.519,0.775) (4.546,0.725) (4.571,0.675) (4.594,0.625) (4.615,0.575) (4.633,0.525) (4.650,0.475) (4.665,0.425) (4.678,0.375) (4.690,0.325) (4.699,0.275) (4.708,0.225) (4.714,0.175) (4.719,0.125) (4.722,0.075) (4.723,0.025) (4.723,-0.050) (4.720,-0.100) (4.717,-0.150) (4.711,-0.200) (4.704,-0.250) (4.695,-0.300) (4.684,-0.350) (4.672,-0.400) (4.658,-0.450) (4.642,-0.500) (4.625,-0.548) (4.605,-0.600) (4.583,-0.650) (4.559,-0.700) (4.533,-0.750) (4.504,-0.800) (4.475,-0.847) (4.439,-0.900) (4.402,-0.950) (4.362,-1.000) (4.325,-1.042) (4.275,-1.095) (4.225,-1.143) (4.175,-1.187) (4.127,-1.225) (4.075,-1.264) (4.025,-1.298) (3.981,-1.325) (3.925,-1.357) (3.875,-1.383) (3.825,-1.407) (3.775,-1.429) (3.725,-1.449) (3.675,-1.467) (3.625,-1.483) (3.575,-1.497) (3.525,-1.510) (3.475,-1.521) (3.425,-1.530) (3.375,-1.538) (3.325,-1.545) (3.275,-1.550) (3.225,-1.553) (3.175,-1.556) (3.125,-1.556) (3.075,-1.556) (3.025,-1.555) (2.975,-1.552) (2.925,-1.548) (2.875,-1.542) (2.825,-1.536) (2.775,-1.529) (2.725,-1.520) (2.675,-1.510) (2.625,-1.499) (2.575,-1.487) (2.525,-1.474) (2.475,-1.461) (2.425,-1.446) (2.375,-1.430) (2.325,-1.413) (2.275,-1.395) (2.225,-1.376) (2.175,-1.357) (2.125,-1.336) (2.075,-1.315) (2.025,-1.293) (1.975,-1.270) (1.925,-1.246) (1.875,-1.221) (1.825,-1.196) (1.775,-1.170) (1.725,-1.144) (1.675,-1.116) (1.625,-1.089) (1.575,-1.060) (1.525,-1.031) (1.475,-1.001) (1.432,-0.975) (1.375,-0.941) (1.325,-0.910) (1.275,-0.878) (1.232,-0.850) (1.175,-0.814) (1.125,-0.781) (1.078,-0.750) (1.025,-0.715) (0.975,-0.681) (0.929,-0.650) (0.875,-0.614) (0.825,-0.579) (0.783,-0.550) (0.725,-0.510) (0.675,-0.475) (0.625,-0.441) (0.575,-0.406) (0.532,-0.375) (0.475,-0.336) (0.425,-0.300) (0.375,-0.266) (0.325,-0.230) (0.284,-0.200) (0.225,-0.160) (0.206,-0.144)}
\EQ{(-0.220,-0.155) (-0.275,-0.195) (-0.319,-0.225) (-0.375,-0.266) (-0.425,-0.300) (-0.475,-0.336) (-0.525,-0.371) (-0.568,-0.400) (-0.625,-0.441) (-0.674,-0.475) (-0.725,-0.511) (-0.775,-0.545) (-0.819,-0.575) (-0.875,-0.614) (-0.925,-0.648) (-0.975,-0.682) (-1.025,-0.716) (-1.075,-0.750) (-1.125,-0.783) (-1.175,-0.816) (-1.225,-0.848) (-1.275,-0.881) (-1.325,-0.913) (-1.375,-0.944) (-1.424,-0.975) (-1.475,-1.006) (-1.525,-1.037) (-1.575,-1.066) (-1.625,-1.096) (-1.675,-1.124) (-1.721,-1.150) (-1.766,-1.175) (-1.812,-1.200) (-1.859,-1.225) (-1.907,-1.250) (-1.957,-1.275) (-2.008,-1.300) (-2.060,-1.325) (-2.114,-1.350) (-2.171,-1.375) (-2.225,-1.398) (-2.275,-1.419) (-2.325,-1.438) (-2.375,-1.457) (-2.425,-1.475) (-2.475,-1.492) (-2.525,-1.508) (-2.575,-1.523) (-2.625,-1.538) (-2.675,-1.551) (-2.725,-1.563) (-2.775,-1.574) (-2.825,-1.585) (-2.875,-1.594) (-2.925,-1.602) (-2.975,-1.609) (-3.025,-1.615) (-3.075,-1.620) (-3.125,-1.623) (-3.175,-1.626) (-3.225,-1.627) (-3.275,-1.627) (-3.325,-1.626) (-3.375,-1.624) (-3.425,-1.620) (-3.475,-1.615) (-3.525,-1.609) (-3.575,-1.601) (-3.625,-1.592) (-3.675,-1.581) (-3.725,-1.569) (-3.775,-1.555) (-3.825,-1.539) (-3.875,-1.522) (-3.925,-1.503) (-3.975,-1.482) (-4.025,-1.459) (-4.075,-1.434) (-4.125,-1.407) (-4.175,-1.378) (-4.218,-1.350) (-4.275,-1.311) (-4.323,-1.275) (-4.375,-1.232) (-4.425,-1.188) (-4.464,-1.150) (-4.500,-1.113) (-4.535,-1.075) (-4.575,-1.028) (-4.616,-0.975) (-4.650,-0.928) (-4.685,-0.875) (-4.716,-0.825) (-4.744,-0.775) (-4.770,-0.725) (-4.793,-0.675) (-4.815,-0.625) (-4.835,-0.575) (-4.853,-0.525) (-4.869,-0.475) (-4.883,-0.425) (-4.895,-0.375) (-4.906,-0.325) (-4.916,-0.275) (-4.923,-0.225) (-4.929,-0.175) (-4.934,-0.125) (-4.937,-0.075) (-4.938,-0.025) (-4.938,0.050) (-4.936,0.100) (-4.932,0.150) (-4.926,0.200) (-4.920,0.250) (-4.911,0.300) (-4.901,0.350) (-4.889,0.400) (-4.876,0.450) (-4.861,0.500) (-4.844,0.550) (-4.825,0.600) (-4.804,0.650) (-4.782,0.700) (-4.757,0.750) (-4.730,0.800) (-4.701,0.850) (-4.669,0.900) (-4.635,0.950) (-4.600,0.996) (-4.556,1.050) (-4.512,1.100) (-4.475,1.139) (-4.438,1.175) (-4.400,1.210) (-4.354,1.250) (-4.300,1.292) (-4.255,1.325) (-4.200,1.362) (-4.150,1.393) (-4.100,1.421) (-4.050,1.447) (-4.000,1.471) (-3.950,1.493) (-3.900,1.513) (-3.850,1.531) (-3.800,1.547) (-3.750,1.562) (-3.700,1.575) (-3.650,1.586) (-3.600,1.596) (-3.550,1.605) (-3.500,1.612) (-3.450,1.618) (-3.400,1.622) (-3.350,1.625) (-3.300,1.627) (-3.250,1.627) (-3.200,1.627) (-3.150,1.625) (-3.100,1.622) (-3.050,1.617) (-3.000,1.612) (-2.950,1.605) (-2.900,1.598) (-2.850,1.589) (-2.800,1.580) (-2.750,1.569) (-2.700,1.557) (-2.650,1.544) (-2.600,1.531) (-2.550,1.516) (-2.500,1.500) (-2.450,1.484) (-2.400,1.466) (-2.350,1.448) (-2.300,1.429) (-2.250,1.409) (-2.200,1.388) (-2.150,1.366) (-2.100,1.344) (-2.050,1.320) (-2.000,1.296) (-1.950,1.272) (-1.900,1.246) (-1.850,1.220) (-1.800,1.194) (-1.750,1.167) (-1.700,1.139) (-1.650,1.110) (-1.600,1.081) (-1.550,1.051) (-1.506,1.025) (-1.450,0.991) (-1.400,0.960) (-1.350,0.929) (-1.305,0.900) (-1.250,0.865) (-1.200,0.832) (-1.151,0.800) (-1.100,0.766) (-1.050,0.733) (-1.001,0.700) (-0.950,0.665) (-0.900,0.631) (-0.855,0.600) (-0.800,0.563) (-0.750,0.528) (-0.700,0.493) (-0.650,0.459) (-0.603,0.425) (-0.550,0.389) (-0.500,0.353) (-0.450,0.318) (-0.400,0.283) (-0.354,0.250) (-0.300,0.213) (-0.250,0.177) (-0.220,0.155)}
\EQ{(0.206,0.144) (0.248,0.175) (0.300,0.213) (0.350,0.248) (0.400,0.283) (0.450,0.318) (0.496,0.350) (0.550,0.389) (0.600,0.423) (0.650,0.459) (0.700,0.493) (0.746,0.525) (0.800,0.563) (0.850,0.597) (0.891,0.625) (0.950,0.665) (1.000,0.699) (1.050,0.733) (1.100,0.766) (1.150,0.799) (1.200,0.832) (1.250,0.865) (1.300,0.897) (1.345,0.925) (1.400,0.960) (1.450,0.991) (1.500,1.021) (1.548,1.050) (1.600,1.081) (1.650,1.110) (1.700,1.139) (1.750,1.167) (1.800,1.194) (1.850,1.220) (1.900,1.246) (1.950,1.272) (2.000,1.296) (2.050,1.320) (2.100,1.344) (2.150,1.366) (2.200,1.388) (2.250,1.409) (2.300,1.429) (2.350,1.448) (2.400,1.466) (2.450,1.484) (2.500,1.500) (2.550,1.516) (2.600,1.531) (2.650,1.544) (2.700,1.557) (2.750,1.569) (2.800,1.580) (2.850,1.589) (2.900,1.598) (2.950,1.605) (3.000,1.612) (3.050,1.617) (3.100,1.622) (3.150,1.625) (3.200,1.627) (3.250,1.627) (3.300,1.627) (3.350,1.625) (3.400,1.622) (3.450,1.618) (3.500,1.612) (3.550,1.605) (3.600,1.596) (3.650,1.586) (3.700,1.575) (3.750,1.562) (3.800,1.547) (3.850,1.531) (3.900,1.513) (3.950,1.493) (4.000,1.471) (4.050,1.447) (4.100,1.421) (4.150,1.393) (4.200,1.362) (4.250,1.328) (4.300,1.292) (4.350,1.253) (4.400,1.210) (4.438,1.175) (4.475,1.139) (4.512,1.100) (4.550,1.058) (4.597,1.000) (4.635,0.950) (4.669,0.900) (4.700,0.851) (4.730,0.800) (4.757,0.750) (4.782,0.700) (4.804,0.650) (4.825,0.600) (4.844,0.550) (4.861,0.500) (4.876,0.450) (4.889,0.400) (4.901,0.350) (4.911,0.300) (4.920,0.250) (4.926,0.200) (4.932,0.150) (4.936,0.100) (4.938,0.050) (4.939,0.000) (4.938,-0.050) (4.936,-0.100) (4.932,-0.150) (4.926,-0.200) (4.920,-0.250) (4.911,-0.300) (4.901,-0.350) (4.889,-0.400) (4.876,-0.450) (4.861,-0.500) (4.844,-0.550) (4.825,-0.600) (4.804,-0.650) (4.782,-0.700) (4.757,-0.750) (4.730,-0.800) (4.701,-0.850) (4.669,-0.900) (4.635,-0.950) (4.600,-0.996) (4.556,-1.050) (4.512,-1.100) (4.475,-1.139) (4.438,-1.175) (4.400,-1.210) (4.354,-1.250) (4.300,-1.292) (4.255,-1.325) (4.200,-1.362) (4.150,-1.393) (4.100,-1.421) (4.050,-1.447) (4.000,-1.471) (3.950,-1.493) (3.900,-1.513) (3.850,-1.531) (3.800,-1.547) (3.750,-1.562) (3.700,-1.575) (3.650,-1.586) (3.600,-1.596) (3.550,-1.605) (3.500,-1.612) (3.450,-1.618) (3.400,-1.622) (3.350,-1.625) (3.300,-1.627) (3.250,-1.627) (3.200,-1.627) (3.150,-1.625) (3.100,-1.622) (3.050,-1.617) (3.000,-1.612) (2.950,-1.605) (2.900,-1.598) (2.850,-1.589) (2.800,-1.580) (2.750,-1.569) (2.700,-1.557) (2.650,-1.544) (2.600,-1.531) (2.550,-1.516) (2.500,-1.500) (2.450,-1.484) (2.400,-1.466) (2.350,-1.448) (2.300,-1.429) (2.250,-1.409) (2.200,-1.388) (2.150,-1.366) (2.100,-1.344) (2.050,-1.320) (2.000,-1.296) (1.950,-1.272) (1.900,-1.246) (1.850,-1.220) (1.800,-1.194) (1.750,-1.167) (1.700,-1.139) (1.650,-1.110) (1.600,-1.081) (1.550,-1.051) (1.506,-1.025) (1.450,-0.991) (1.400,-0.960) (1.350,-0.929) (1.305,-0.900) (1.250,-0.865) (1.200,-0.832) (1.151,-0.800) (1.100,-0.766) (1.050,-0.733) (1.001,-0.700) (0.950,-0.665) (0.900,-0.631) (0.855,-0.600) (0.800,-0.563) (0.750,-0.528) (0.700,-0.493) (0.650,-0.459) (0.603,-0.425) (0.550,-0.389) (0.500,-0.353) (0.450,-0.318) (0.400,-0.283) (0.354,-0.250) (0.300,-0.213) (0.250,-0.177) (0.206,-0.144)}
\EQ{(-0.220,-0.155) (-0.275,-0.195) (-0.319,-0.225) (-0.375,-0.266) (-0.425,-0.300) (-0.475,-0.336) (-0.525,-0.371) (-0.567,-0.400) (-0.625,-0.441) (-0.674,-0.475) (-0.725,-0.511) (-0.775,-0.546) (-0.818,-0.575) (-0.875,-0.615) (-0.925,-0.649) (-0.975,-0.683) (-1.025,-0.717) (-1.074,-0.750) (-1.125,-0.785) (-1.175,-0.818) (-1.224,-0.850) (-1.275,-0.884) (-1.325,-0.916) (-1.375,-0.948) (-1.417,-0.975) (-1.475,-1.011) (-1.525,-1.042) (-1.575,-1.073) (-1.620,-1.100) (-1.675,-1.133) (-1.725,-1.162) (-1.775,-1.190) (-1.825,-1.218) (-1.875,-1.246) (-1.925,-1.273) (-1.975,-1.299) (-2.025,-1.325) (-2.075,-1.350) (-2.125,-1.374) (-2.175,-1.398) (-2.225,-1.421) (-2.275,-1.443) (-2.325,-1.465) (-2.375,-1.486) (-2.425,-1.505) (-2.475,-1.525) (-2.525,-1.543) (-2.575,-1.560) (-2.625,-1.577) (-2.675,-1.593) (-2.725,-1.608) (-2.775,-1.622) (-2.825,-1.635) (-2.875,-1.647) (-2.925,-1.658) (-2.975,-1.668) (-3.025,-1.677) (-3.075,-1.686) (-3.125,-1.693) (-3.175,-1.699) (-3.225,-1.704) (-3.275,-1.708) (-3.325,-1.711) (-3.375,-1.713) (-3.425,-1.713) (-3.475,-1.713) (-3.525,-1.711) (-3.575,-1.708) (-3.625,-1.704) (-3.675,-1.698) (-3.725,-1.691) (-3.775,-1.683) (-3.825,-1.674) (-3.875,-1.663) (-3.925,-1.650) (-3.975,-1.636) (-4.025,-1.621) (-4.075,-1.604) (-4.125,-1.585) (-4.175,-1.564) (-4.225,-1.542) (-4.275,-1.518) (-4.325,-1.491) (-4.375,-1.462) (-4.425,-1.431) (-4.472,-1.400) (-4.525,-1.362) (-4.572,-1.325) (-4.625,-1.280) (-4.675,-1.234) (-4.725,-1.185) (-4.775,-1.130) (-4.822,-1.075) (-4.861,-1.025) (-4.897,-0.975) (-4.925,-0.932) (-4.961,-0.875) (-4.989,-0.825) (-5.015,-0.775) (-5.040,-0.725) (-5.062,-0.675) (-5.082,-0.625) (-5.100,-0.577) (-5.117,-0.525) (-5.133,-0.475) (-5.146,-0.425) (-5.158,-0.375) (-5.168,-0.325) (-5.177,-0.275) (-5.184,-0.225) (-5.190,-0.175) (-5.195,-0.125) (-5.197,-0.075) (-5.199,-0.025) (-5.198,0.050) (-5.196,0.100) (-5.193,0.150) (-5.187,0.200) (-5.181,0.250) (-5.173,0.300) (-5.163,0.350) (-5.152,0.400) (-5.140,0.450) (-5.125,0.500) (-5.109,0.550) (-5.092,0.600) (-5.072,0.650) (-5.051,0.700) (-5.028,0.750) (-5.002,0.800) (-4.975,0.850) (-4.946,0.900) (-4.914,0.950) (-4.879,1.000) (-4.842,1.050) (-4.801,1.100) (-4.758,1.150) (-4.710,1.200) (-4.659,1.250) (-4.602,1.300) (-4.550,1.343) (-4.507,1.375) (-4.450,1.415) (-4.400,1.447) (-4.353,1.475) (-4.309,1.500) (-4.260,1.525) (-4.208,1.550) (-4.150,1.575) (-4.100,1.595) (-4.050,1.613) (-4.000,1.629) (-3.950,1.644) (-3.900,1.657) (-3.850,1.668) (-3.800,1.679) (-3.750,1.687) (-3.700,1.695) (-3.650,1.701) (-3.600,1.706) (-3.550,1.710) (-3.500,1.712) (-3.450,1.713) (-3.400,1.713) (-3.350,1.712) (-3.300,1.710) (-3.250,1.706) (-3.200,1.702) (-3.150,1.696) (-3.100,1.689) (-3.050,1.682) (-3.000,1.673) (-2.950,1.663) (-2.900,1.653) (-2.850,1.641) (-2.800,1.628) (-2.750,1.615) (-2.700,1.600) (-2.650,1.585) (-2.600,1.569) (-2.550,1.552) (-2.500,1.534) (-2.450,1.515) (-2.400,1.496) (-2.350,1.475) (-2.300,1.454) (-2.250,1.432) (-2.200,1.410) (-2.150,1.386) (-2.100,1.362) (-2.050,1.338) (-2.000,1.312) (-1.950,1.286) (-1.900,1.259) (-1.850,1.232) (-1.800,1.204) (-1.750,1.176) (-1.705,1.150) (-1.650,1.118) (-1.600,1.088) (-1.550,1.058) (-1.500,1.027) (-1.457,1.000) (-1.400,0.964) (-1.350,0.932) (-1.300,0.900) (-1.250,0.867) (-1.200,0.835) (-1.150,0.801) (-1.100,0.768) (-1.050,0.734) (-1.000,0.700) (-0.950,0.666) (-0.900,0.632) (-0.854,0.600) (-0.800,0.563) (-0.750,0.528) (-0.700,0.494) (-0.650,0.459) (-0.603,0.425) (-0.550,0.389) (-0.500,0.353) (-0.450,0.318) (-0.400,0.283) (-0.354,0.250) (-0.300,0.213) (-0.250,0.177) (-0.220,0.155)}
\EQ{(0.206,0.144) (0.248,0.175) (0.300,0.213) (0.350,0.248) (0.400,0.283) (0.450,0.318) (0.496,0.350) (0.550,0.389) (0.600,0.423) (0.650,0.459) (0.700,0.494) (0.746,0.525) (0.800,0.563) (0.850,0.598) (0.900,0.632) (0.950,0.666) (1.000,0.700) (1.050,0.734) (1.100,0.768) (1.148,0.800) (1.200,0.835) (1.250,0.867) (1.300,0.900) (1.350,0.932) (1.400,0.964) (1.450,0.996) (1.497,1.025) (1.550,1.058) (1.600,1.088) (1.650,1.118) (1.700,1.147) (1.748,1.175) (1.792,1.200) (1.837,1.225) (1.883,1.250) (1.929,1.275) (1.975,1.299) (2.025,1.325) (2.075,1.350) (2.125,1.374) (2.175,1.398) (2.225,1.421) (2.275,1.443) (2.325,1.465) (2.375,1.486) (2.425,1.505) (2.475,1.525) (2.525,1.543) (2.575,1.560) (2.625,1.577) (2.675,1.593) (2.725,1.608) (2.775,1.622) (2.825,1.635) (2.875,1.647) (2.925,1.658) (2.975,1.668) (3.025,1.677) (3.075,1.686) (3.125,1.693) (3.175,1.699) (3.225,1.704) (3.275,1.708) (3.325,1.711) (3.375,1.713) (3.425,1.713) (3.475,1.713) (3.525,1.711) (3.575,1.708) (3.625,1.704) (3.675,1.698) (3.725,1.691) (3.775,1.683) (3.825,1.674) (3.875,1.663) (3.925,1.650) (3.975,1.636) (4.025,1.621) (4.075,1.604) (4.125,1.585) (4.175,1.564) (4.225,1.542) (4.275,1.518) (4.325,1.491) (4.375,1.462) (4.425,1.431) (4.472,1.400) (4.525,1.362) (4.572,1.325) (4.625,1.280) (4.675,1.234) (4.725,1.185) (4.775,1.130) (4.822,1.075) (4.861,1.025) (4.897,0.975) (4.925,0.932) (4.961,0.875) (4.989,0.825) (5.015,0.775) (5.040,0.725) (5.062,0.675) (5.082,0.625) (5.100,0.577) (5.117,0.525) (5.133,0.475) (5.146,0.425) (5.158,0.375) (5.168,0.325) (5.177,0.275) (5.184,0.225) (5.190,0.175) (5.195,0.125) (5.197,0.075) (5.199,0.025) (5.198,-0.050) (5.196,-0.100) (5.193,-0.150) (5.187,-0.200) (5.181,-0.250) (5.173,-0.300) (5.163,-0.350) (5.152,-0.400) (5.140,-0.450) (5.125,-0.500) (5.109,-0.550) (5.092,-0.600) (5.072,-0.650) (5.051,-0.700) (5.028,-0.750) (5.002,-0.800) (4.975,-0.850) (4.946,-0.900) (4.914,-0.950) (4.879,-1.000) (4.842,-1.050) (4.801,-1.100) (4.758,-1.150) (4.710,-1.200) (4.659,-1.250) (4.602,-1.300) (4.550,-1.343) (4.507,-1.375) (4.450,-1.415) (4.400,-1.447) (4.353,-1.475) (4.309,-1.500) (4.260,-1.525) (4.208,-1.550) (4.150,-1.575) (4.100,-1.595) (4.050,-1.613) (4.000,-1.629) (3.950,-1.644) (3.900,-1.657) (3.850,-1.668) (3.800,-1.679) (3.750,-1.687) (3.700,-1.695) (3.650,-1.701) (3.600,-1.706) (3.550,-1.710) (3.500,-1.712) (3.450,-1.713) (3.400,-1.713) (3.350,-1.712) (3.300,-1.710) (3.250,-1.706) (3.200,-1.702) (3.150,-1.696) (3.100,-1.689) (3.050,-1.682) (3.000,-1.673) (2.950,-1.663) (2.900,-1.653) (2.850,-1.641) (2.800,-1.628) (2.750,-1.615) (2.700,-1.600) (2.650,-1.585) (2.600,-1.569) (2.550,-1.552) (2.500,-1.534) (2.450,-1.515) (2.400,-1.496) (2.350,-1.475) (2.300,-1.454) (2.250,-1.432) (2.200,-1.410) (2.150,-1.386) (2.100,-1.362) (2.050,-1.338) (2.000,-1.312) (1.950,-1.286) (1.900,-1.259) (1.850,-1.232) (1.800,-1.204) (1.750,-1.176) (1.705,-1.150) (1.650,-1.118) (1.600,-1.088) (1.550,-1.058) (1.500,-1.027) (1.457,-1.000) (1.400,-0.964) (1.350,-0.932) (1.300,-0.900) (1.250,-0.867) (1.200,-0.835) (1.150,-0.801) (1.100,-0.768) (1.050,-0.734) (1.000,-0.700) (0.950,-0.666) (0.900,-0.632) (0.854,-0.600) (0.800,-0.563) (0.750,-0.528) (0.700,-0.494) (0.650,-0.459) (0.603,-0.425) (0.550,-0.389) (0.500,-0.353) (0.450,-0.318) (0.400,-0.283) (0.354,-0.250) (0.300,-0.213) (0.250,-0.177) (0.206,-0.144)}
\EQ{(-0.220,-0.155) (-0.275,-0.195) (-0.319,-0.225) (-0.375,-0.266) (-0.425,-0.300) (-0.475,-0.336) (-0.525,-0.371) (-0.567,-0.400) (-0.625,-0.441) (-0.674,-0.475) (-0.725,-0.511) (-0.775,-0.546) (-0.817,-0.575) (-0.875,-0.616) (-0.925,-0.650) (-0.975,-0.684) (-1.025,-0.719) (-1.072,-0.750) (-1.125,-0.787) (-1.175,-0.820) (-1.220,-0.850) (-1.275,-0.887) (-1.325,-0.920) (-1.372,-0.950) (-1.425,-0.985) (-1.475,-1.017) (-1.525,-1.048) (-1.568,-1.075) (-1.625,-1.110) (-1.675,-1.141) (-1.725,-1.171) (-1.774,-1.200) (-1.825,-1.230) (-1.875,-1.259) (-1.925,-1.287) (-1.975,-1.314) (-2.025,-1.342) (-2.075,-1.368) (-2.125,-1.394) (-2.175,-1.420) (-2.225,-1.444) (-2.275,-1.468) (-2.325,-1.492) (-2.375,-1.515) (-2.425,-1.537) (-2.475,-1.558) (-2.525,-1.579) (-2.575,-1.599) (-2.625,-1.618) (-2.675,-1.637) (-2.725,-1.654) (-2.775,-1.671) (-2.825,-1.687) (-2.875,-1.702) (-2.925,-1.717) (-2.975,-1.730) (-3.025,-1.743) (-3.075,-1.755) (-3.125,-1.765) (-3.175,-1.775) (-3.225,-1.784) (-3.275,-1.792) (-3.325,-1.799) (-3.375,-1.805) (-3.425,-1.811) (-3.475,-1.815) (-3.525,-1.818) (-3.575,-1.820) (-3.625,-1.820) (-3.675,-1.820) (-3.725,-1.819) (-3.775,-1.816) (-3.825,-1.813) (-3.875,-1.808) (-3.925,-1.802) (-3.975,-1.795) (-4.025,-1.786) (-4.075,-1.776) (-4.125,-1.765) (-4.175,-1.753) (-4.225,-1.739) (-4.275,-1.723) (-4.325,-1.706) (-4.375,-1.688) (-4.425,-1.667) (-4.475,-1.645) (-4.525,-1.622) (-4.575,-1.596) (-4.625,-1.568) (-4.675,-1.538) (-4.725,-1.506) (-4.770,-1.475) (-4.825,-1.435) (-4.868,-1.400) (-4.925,-1.352) (-4.975,-1.305) (-5.025,-1.255) (-5.075,-1.200) (-5.118,-1.150) (-5.150,-1.109) (-5.194,-1.050) (-5.225,-1.004) (-5.259,-0.950) (-5.289,-0.900) (-5.316,-0.850) (-5.341,-0.800) (-5.365,-0.750) (-5.386,-0.700) (-5.406,-0.650) (-5.424,-0.600) (-5.441,-0.550) (-5.456,-0.500) (-5.469,-0.450) (-5.481,-0.400) (-5.491,-0.350) (-5.500,-0.300) (-5.508,-0.250) (-5.514,-0.200) (-5.519,-0.150) (-5.522,-0.100) (-5.524,-0.050) (-5.525,0.000) (-5.524,0.050) (-5.522,0.100) (-5.519,0.150) (-5.514,0.200) (-5.508,0.250) (-5.500,0.300) (-5.491,0.350) (-5.481,0.400) (-5.469,0.450) (-5.456,0.500) (-5.441,0.550) (-5.424,0.600) (-5.406,0.650) (-5.386,0.700) (-5.365,0.750) (-5.341,0.800) (-5.316,0.850) (-5.289,0.900) (-5.259,0.950) (-5.228,1.000) (-5.194,1.050) (-5.157,1.100) (-5.125,1.141) (-5.075,1.200) (-5.030,1.250) (-4.980,1.300) (-4.927,1.350) (-4.875,1.395) (-4.825,1.435) (-4.775,1.472) (-4.725,1.506) (-4.675,1.538) (-4.625,1.568) (-4.575,1.596) (-4.525,1.622) (-4.475,1.645) (-4.425,1.667) (-4.375,1.688) (-4.325,1.706) (-4.275,1.723) (-4.225,1.739) (-4.175,1.753) (-4.125,1.765) (-4.075,1.776) (-4.025,1.786) (-3.975,1.795) (-3.925,1.802) (-3.875,1.808) (-3.825,1.813) (-3.775,1.816) (-3.725,1.819) (-3.675,1.820) (-3.625,1.820) (-3.575,1.820) (-3.525,1.818) (-3.475,1.815) (-3.425,1.811) (-3.375,1.805) (-3.325,1.799) (-3.275,1.792) (-3.225,1.784) (-3.175,1.775) (-3.125,1.765) (-3.075,1.755) (-3.025,1.743) (-2.975,1.730) (-2.925,1.717) (-2.875,1.702) (-2.825,1.687) (-2.775,1.671) (-2.725,1.654) (-2.675,1.637) (-2.625,1.618) (-2.578,1.600) (-2.525,1.579) (-2.475,1.558) (-2.425,1.537) (-2.375,1.515) (-2.325,1.492) (-2.275,1.468) (-2.225,1.444) (-2.175,1.420) (-2.125,1.394) (-2.075,1.368) (-2.025,1.342) (-1.975,1.314) (-1.925,1.287) (-1.875,1.259) (-1.825,1.230) (-1.775,1.201) (-1.725,1.171) (-1.675,1.141) (-1.625,1.110) (-1.575,1.079) (-1.528,1.050) (-1.475,1.017) (-1.425,0.985) (-1.375,0.952) (-1.325,0.920) (-1.275,0.887) (-1.225,0.854) (-1.183,0.825) (-1.125,0.787) (-1.075,0.753) (-1.025,0.719) (-0.975,0.684) (-0.925,0.650) (-0.875,0.616) (-0.825,0.581) (-0.781,0.550) (-0.725,0.511) (-0.675,0.476) (-0.625,0.441) (-0.575,0.406) (-0.532,0.375) (-0.475,0.336) (-0.425,0.300) (-0.375,0.266) (-0.325,0.230) (-0.284,0.200) (-0.225,0.160) (-0.220,0.155)}
\EQ{(0.206,0.144) (0.248,0.175) (0.300,0.213) (0.350,0.248) (0.400,0.283) (0.450,0.318) (0.496,0.350) (0.550,0.389) (0.600,0.424) (0.650,0.459) (0.700,0.494) (0.745,0.525) (0.800,0.564) (0.850,0.598) (0.900,0.633) (0.950,0.667) (0.998,0.700) (1.050,0.736) (1.100,0.770) (1.145,0.800) (1.200,0.837) (1.250,0.870) (1.296,0.900) (1.350,0.936) (1.400,0.968) (1.449,1.000) (1.500,1.032) (1.550,1.064) (1.600,1.095) (1.649,1.125) (1.700,1.156) (1.750,1.186) (1.800,1.215) (1.850,1.244) (1.900,1.273) (1.949,1.300) (1.995,1.325) (2.041,1.350) (2.088,1.375) (2.137,1.400) (2.186,1.425) (2.237,1.450) (2.289,1.475) (2.343,1.500) (2.398,1.525) (2.450,1.548) (2.500,1.569) (2.550,1.589) (2.600,1.609) (2.650,1.628) (2.700,1.646) (2.750,1.663) (2.800,1.679) (2.850,1.695) (2.900,1.710) (2.950,1.724) (3.000,1.737) (3.050,1.749) (3.100,1.760) (3.150,1.771) (3.200,1.780) (3.250,1.788) (3.300,1.796) (3.350,1.803) (3.400,1.808) (3.450,1.813) (3.500,1.816) (3.550,1.819) (3.600,1.820) (3.650,1.820) (3.700,1.820) (3.750,1.818) (3.800,1.815) (3.850,1.811) (3.900,1.805) (3.950,1.799) (4.000,1.791) (4.050,1.782) (4.100,1.771) (4.150,1.759) (4.200,1.746) (4.250,1.731) (4.300,1.715) (4.350,1.697) (4.400,1.678) (4.450,1.657) (4.500,1.634) (4.550,1.609) (4.600,1.582) (4.650,1.553) (4.696,1.525) (4.750,1.489) (4.800,1.453) (4.850,1.415) (4.898,1.375) (4.950,1.329) (5.000,1.280) (5.050,1.228) (5.097,1.175) (5.138,1.125) (5.175,1.076) (5.211,1.025) (5.244,0.975) (5.274,0.925) (5.300,0.880) (5.325,0.832) (5.350,0.782) (5.375,0.726) (5.396,0.675) (5.415,0.625) (5.433,0.575) (5.448,0.525) (5.462,0.475) (5.475,0.425) (5.486,0.375) (5.496,0.325) (5.504,0.275) (5.511,0.225) (5.516,0.175) (5.521,0.125) (5.523,0.075) (5.525,0.025) (5.524,-0.050) (5.522,-0.100) (5.519,-0.150) (5.514,-0.200) (5.508,-0.250) (5.500,-0.300) (5.491,-0.350) (5.481,-0.400) (5.469,-0.450) (5.456,-0.500) (5.441,-0.550) (5.424,-0.600) (5.406,-0.650) (5.386,-0.700) (5.365,-0.750) (5.341,-0.800) (5.316,-0.850) (5.289,-0.900) (5.259,-0.950) (5.228,-1.000) (5.194,-1.050) (5.157,-1.100) (5.125,-1.141) (5.075,-1.200) (5.030,-1.250) (4.980,-1.300) (4.927,-1.350) (4.875,-1.395) (4.825,-1.435) (4.775,-1.472) (4.725,-1.506) (4.675,-1.538) (4.625,-1.568) (4.575,-1.596) (4.525,-1.622) (4.475,-1.645) (4.425,-1.667) (4.375,-1.688) (4.325,-1.706) (4.275,-1.723) (4.225,-1.739) (4.175,-1.753) (4.125,-1.765) (4.075,-1.776) (4.025,-1.786) (3.975,-1.795) (3.925,-1.802) (3.875,-1.808) (3.825,-1.813) (3.775,-1.816) (3.725,-1.819) (3.675,-1.820) (3.625,-1.820) (3.575,-1.820) (3.525,-1.818) (3.475,-1.815) (3.425,-1.811) (3.375,-1.805) (3.325,-1.799) (3.275,-1.792) (3.225,-1.784) (3.175,-1.775) (3.125,-1.765) (3.075,-1.755) (3.025,-1.743) (2.975,-1.730) (2.925,-1.717) (2.875,-1.702) (2.825,-1.687) (2.775,-1.671) (2.725,-1.654) (2.675,-1.637) (2.625,-1.618) (2.578,-1.600) (2.525,-1.579) (2.475,-1.558) (2.425,-1.537) (2.375,-1.515) (2.325,-1.492) (2.275,-1.468) (2.225,-1.444) (2.175,-1.420) (2.125,-1.394) (2.075,-1.368) (2.025,-1.342) (1.975,-1.314) (1.925,-1.287) (1.875,-1.259) (1.825,-1.230) (1.775,-1.201) (1.725,-1.171) (1.675,-1.141) (1.625,-1.110) (1.575,-1.079) (1.528,-1.050) (1.475,-1.017) (1.425,-0.985) (1.375,-0.952) (1.325,-0.920) (1.275,-0.887) (1.225,-0.854) (1.183,-0.825) (1.125,-0.787) (1.075,-0.753) (1.025,-0.719) (0.975,-0.684) (0.925,-0.650) (0.875,-0.616) (0.825,-0.581) (0.781,-0.550) (0.725,-0.511) (0.675,-0.476) (0.625,-0.441) (0.575,-0.406) (0.532,-0.375) (0.475,-0.336) (0.425,-0.300) (0.375,-0.266) (0.325,-0.230) (0.284,-0.200) (0.225,-0.160) (0.206,-0.144)}
\EQ{(-0.220,-0.155) (-0.275,-0.195) (-0.319,-0.225) (-0.375,-0.266) (-0.425,-0.300) (-0.475,-0.336) (-0.525,-0.371) (-0.567,-0.400) (-0.625,-0.442) (-0.673,-0.475) (-0.725,-0.512) (-0.775,-0.547) (-0.817,-0.575) (-0.875,-0.616) (-0.924,-0.650) (-0.975,-0.686) (-1.025,-0.720) (-1.069,-0.750) (-1.125,-0.788) (-1.175,-0.822) (-1.225,-0.856) (-1.275,-0.890) (-1.325,-0.923) (-1.375,-0.956) (-1.425,-0.989) (-1.475,-1.022) (-1.519,-1.050) (-1.575,-1.086) (-1.625,-1.118) (-1.675,-1.149) (-1.725,-1.180) (-1.775,-1.211) (-1.825,-1.242) (-1.875,-1.271) (-1.924,-1.300) (-1.975,-1.330) (-2.025,-1.359) (-2.075,-1.387) (-2.125,-1.415) (-2.175,-1.442) (-2.225,-1.468) (-2.275,-1.495) (-2.325,-1.520) (-2.375,-1.545) (-2.425,-1.570) (-2.475,-1.593) (-2.525,-1.617) (-2.575,-1.639) (-2.625,-1.661) (-2.675,-1.682) (-2.725,-1.703) (-2.775,-1.723) (-2.825,-1.742) (-2.875,-1.761) (-2.925,-1.778) (-2.975,-1.796) (-3.025,-1.812) (-3.075,-1.827) (-3.125,-1.842) (-3.175,-1.856) (-3.225,-1.869) (-3.275,-1.881) (-3.325,-1.893) (-3.375,-1.904) (-3.425,-1.913) (-3.475,-1.922) (-3.525,-1.930) (-3.575,-1.937) (-3.625,-1.944) (-3.675,-1.949) (-3.725,-1.953) (-3.775,-1.957) (-3.825,-1.959) (-3.875,-1.961) (-3.925,-1.961) (-3.975,-1.960) (-4.025,-1.959) (-4.075,-1.956) (-4.125,-1.952) (-4.175,-1.948) (-4.225,-1.942) (-4.275,-1.934) (-4.325,-1.926) (-4.375,-1.917) (-4.425,-1.906) (-4.475,-1.894) (-4.525,-1.880) (-4.575,-1.866) (-4.624,-1.850) (-4.675,-1.832) (-4.725,-1.813) (-4.775,-1.793) (-4.825,-1.770) (-4.875,-1.746) (-4.925,-1.721) (-4.975,-1.693) (-5.025,-1.664) (-5.075,-1.632) (-5.123,-1.600) (-5.175,-1.562) (-5.223,-1.525) (-5.275,-1.482) (-5.325,-1.438) (-5.365,-1.400) (-5.415,-1.350) (-5.450,-1.312) (-5.500,-1.254) (-5.544,-1.200) (-5.575,-1.158) (-5.616,-1.100) (-5.649,-1.050) (-5.675,-1.007) (-5.708,-0.950) (-5.735,-0.900) (-5.759,-0.850) (-5.782,-0.800) (-5.804,-0.750) (-5.823,-0.700) (-5.842,-0.650) (-5.858,-0.600) (-5.874,-0.550) (-5.887,-0.500) (-5.900,-0.450) (-5.911,-0.400) (-5.920,-0.350) (-5.929,-0.300) (-5.936,-0.250) (-5.941,-0.200) (-5.946,-0.150) (-5.949,-0.100) (-5.951,-0.050) (-5.951,0.000) (-5.951,0.050) (-5.949,0.100) (-5.946,0.150) (-5.941,0.200) (-5.936,0.250) (-5.929,0.300) (-5.920,0.350) (-5.911,0.400) (-5.900,0.450) (-5.887,0.500) (-5.874,0.550) (-5.858,0.600) (-5.842,0.650) (-5.823,0.700) (-5.804,0.750) (-5.782,0.800) (-5.759,0.850) (-5.735,0.900) (-5.708,0.950) (-5.679,1.000) (-5.650,1.048) (-5.616,1.100) (-5.581,1.150) (-5.550,1.192) (-5.504,1.250) (-5.461,1.300) (-5.425,1.339) (-5.375,1.390) (-5.325,1.438) (-5.284,1.475) (-5.225,1.524) (-5.175,1.562) (-5.125,1.599) (-5.075,1.632) (-5.025,1.664) (-4.975,1.693) (-4.925,1.721) (-4.875,1.746) (-4.825,1.770) (-4.775,1.793) (-4.725,1.813) (-4.675,1.832) (-4.625,1.850) (-4.575,1.866) (-4.525,1.880) (-4.475,1.894) (-4.425,1.906) (-4.375,1.917) (-4.325,1.926) (-4.275,1.934) (-4.225,1.942) (-4.175,1.948) (-4.125,1.952) (-4.075,1.956) (-4.025,1.959) (-3.975,1.960) (-3.925,1.961) (-3.875,1.961) (-3.825,1.959) (-3.775,1.957) (-3.725,1.953) (-3.675,1.949) (-3.625,1.944) (-3.575,1.937) (-3.525,1.930) (-3.475,1.922) (-3.425,1.913) (-3.375,1.904) (-3.325,1.893) (-3.275,1.881) (-3.225,1.869) (-3.175,1.856) (-3.125,1.842) (-3.075,1.827) (-3.025,1.812) (-2.975,1.796) (-2.925,1.778) (-2.875,1.761) (-2.825,1.742) (-2.775,1.723) (-2.725,1.703) (-2.675,1.682) (-2.625,1.661) (-2.575,1.639) (-2.525,1.617) (-2.475,1.593) (-2.425,1.570) (-2.375,1.545) (-2.325,1.520) (-2.275,1.495) (-2.225,1.468) (-2.175,1.442) (-2.125,1.415) (-2.075,1.387) (-2.025,1.359) (-1.975,1.330) (-1.925,1.301) (-1.881,1.275) (-1.825,1.242) (-1.775,1.211) (-1.725,1.180) (-1.676,1.150) (-1.625,1.118) (-1.575,1.086) (-1.525,1.054) (-1.480,1.025) (-1.425,0.989) (-1.375,0.956) (-1.328,0.925) (-1.275,0.890) (-1.225,0.856) (-1.179,0.825) (-1.125,0.788) (-1.075,0.754) (-1.033,0.725) (-0.975,0.686) (-0.925,0.651) (-0.875,0.616) (-0.825,0.581) (-0.781,0.550) (-0.725,0.512) (-0.675,0.476) (-0.625,0.442) (-0.575,0.406) (-0.531,0.375) (-0.475,0.336) (-0.425,0.300) (-0.375,0.266) (-0.325,0.230) (-0.284,0.200) (-0.225,0.160) (-0.220,0.155)}
\EQ{(0.206,0.144) (0.248,0.175) (0.300,0.213) (0.350,0.248) (0.400,0.283) (0.450,0.318) (0.496,0.350) (0.550,0.389) (0.600,0.424) (0.650,0.459) (0.700,0.494) (0.745,0.525) (0.800,0.564) (0.850,0.599) (0.900,0.634) (0.950,0.668) (0.997,0.700) (1.050,0.737) (1.100,0.771) (1.143,0.800) (1.200,0.839) (1.250,0.873) (1.300,0.906) (1.350,0.940) (1.400,0.973) (1.442,1.000) (1.500,1.038) (1.550,1.070) (1.597,1.100) (1.650,1.134) (1.700,1.165) (1.750,1.196) (1.798,1.225) (1.850,1.257) (1.900,1.286) (1.950,1.316) (2.000,1.344) (2.050,1.373) (2.099,1.400) (2.144,1.425) (2.191,1.450) (2.238,1.475) (2.286,1.500) (2.335,1.525) (2.385,1.550) (2.437,1.575) (2.489,1.600) (2.544,1.625) (2.600,1.650) (2.650,1.672) (2.700,1.693) (2.750,1.713) (2.800,1.733) (2.850,1.751) (2.900,1.770) (2.950,1.787) (3.000,1.804) (3.050,1.820) (3.100,1.835) (3.150,1.849) (3.200,1.863) (3.249,1.875) (3.300,1.887) (3.350,1.898) (3.400,1.909) (3.450,1.918) (3.500,1.926) (3.550,1.934) (3.600,1.941) (3.650,1.946) (3.700,1.951) (3.750,1.955) (3.800,1.958) (3.850,1.960) (3.900,1.961) (3.950,1.961) (4.000,1.960) (4.050,1.958) (4.100,1.954) (4.150,1.950) (4.200,1.945) (4.250,1.938) (4.300,1.930) (4.350,1.922) (4.400,1.911) (4.450,1.900) (4.500,1.887) (4.550,1.873) (4.600,1.858) (4.650,1.841) (4.700,1.823) (4.750,1.803) (4.800,1.782) (4.850,1.759) (4.900,1.734) (4.950,1.707) (5.000,1.679) (5.047,1.650) (5.100,1.616) (5.150,1.581) (5.191,1.550) (5.250,1.503) (5.300,1.461) (5.339,1.425) (5.390,1.375) (5.438,1.325) (5.475,1.284) (5.524,1.225) (5.563,1.175) (5.599,1.125) (5.633,1.075) (5.664,1.025) (5.694,0.975) (5.721,0.925) (5.747,0.875) (5.771,0.825) (5.793,0.775) (5.814,0.725) (5.833,0.675) (5.850,0.625) (5.866,0.575) (5.881,0.525) (5.894,0.475) (5.905,0.425) (5.916,0.375) (5.925,0.325) (5.932,0.275) (5.939,0.225) (5.944,0.175) (5.947,0.125) (5.950,0.075) (5.951,0.025) (5.951,-0.050) (5.949,-0.100) (5.946,-0.150) (5.941,-0.200) (5.936,-0.250) (5.929,-0.300) (5.920,-0.350) (5.911,-0.400) (5.900,-0.450) (5.887,-0.500) (5.874,-0.550) (5.858,-0.600) (5.842,-0.650) (5.823,-0.700) (5.804,-0.750) (5.782,-0.800) (5.759,-0.850) (5.735,-0.900) (5.708,-0.950) (5.679,-1.000) (5.650,-1.048) (5.616,-1.100) (5.581,-1.150) (5.550,-1.192) (5.504,-1.250) (5.461,-1.300) (5.425,-1.339) (5.375,-1.390) (5.325,-1.438) (5.284,-1.475) (5.225,-1.524) (5.175,-1.562) (5.125,-1.599) (5.075,-1.632) (5.025,-1.664) (4.975,-1.693) (4.925,-1.721) (4.875,-1.746) (4.825,-1.770) (4.775,-1.793) (4.725,-1.813) (4.675,-1.832) (4.625,-1.850) (4.575,-1.866) (4.525,-1.880) (4.475,-1.894) (4.425,-1.906) (4.375,-1.917) (4.325,-1.926) (4.275,-1.934) (4.225,-1.942) (4.175,-1.948) (4.125,-1.952) (4.075,-1.956) (4.025,-1.959) (3.975,-1.960) (3.925,-1.961) (3.875,-1.961) (3.825,-1.959) (3.775,-1.957) (3.725,-1.953) (3.675,-1.949) (3.625,-1.944) (3.575,-1.937) (3.525,-1.930) (3.475,-1.922) (3.425,-1.913) (3.375,-1.904) (3.325,-1.893) (3.275,-1.881) (3.225,-1.869) (3.175,-1.856) (3.125,-1.842) (3.075,-1.827) (3.025,-1.812) (2.975,-1.796) (2.925,-1.778) (2.875,-1.761) (2.825,-1.742) (2.775,-1.723) (2.725,-1.703) (2.675,-1.682) (2.625,-1.661) (2.575,-1.639) (2.525,-1.617) (2.475,-1.593) (2.425,-1.570) (2.375,-1.545) (2.325,-1.520) (2.275,-1.495) (2.225,-1.468) (2.175,-1.442) (2.125,-1.415) (2.075,-1.387) (2.025,-1.359) (1.975,-1.330) (1.925,-1.301) (1.881,-1.275) (1.825,-1.242) (1.775,-1.211) (1.725,-1.180) (1.676,-1.150) (1.625,-1.118) (1.575,-1.086) (1.525,-1.054) (1.480,-1.025) (1.425,-0.989) (1.375,-0.956) (1.328,-0.925) (1.275,-0.890) (1.225,-0.856) (1.179,-0.825) (1.125,-0.788) (1.075,-0.754) (1.033,-0.725) (0.975,-0.686) (0.925,-0.651) (0.875,-0.616) (0.825,-0.581) (0.781,-0.550) (0.725,-0.512) (0.675,-0.476) (0.625,-0.442) (0.575,-0.406) (0.531,-0.375) (0.475,-0.336) (0.425,-0.300) (0.375,-0.266) (0.325,-0.230) (0.284,-0.200) (0.225,-0.160) (0.206,-0.144)}
\EQ{(-0.220,-0.155) (-0.275,-0.195) (-0.319,-0.225) (-0.375,-0.266) (-0.425,-0.300) (-0.475,-0.336) (-0.525,-0.371) (-0.567,-0.400) (-0.625,-0.442) (-0.673,-0.475) (-0.725,-0.512) (-0.775,-0.547) (-0.825,-0.582) (-0.875,-0.617) (-0.923,-0.650) (-0.975,-0.687) (-1.025,-0.721) (-1.067,-0.750) (-1.125,-0.790) (-1.175,-0.824) (-1.225,-0.859) (-1.275,-0.893) (-1.323,-0.925) (-1.375,-0.960) (-1.425,-0.994) (-1.472,-1.025) (-1.525,-1.060) (-1.575,-1.093) (-1.624,-1.125) (-1.675,-1.158) (-1.725,-1.190) (-1.775,-1.222) (-1.820,-1.250) (-1.875,-1.285) (-1.925,-1.316) (-1.975,-1.346) (-2.023,-1.375) (-2.075,-1.406) (-2.125,-1.436) (-2.175,-1.465) (-2.225,-1.493) (-2.275,-1.521) (-2.325,-1.549) (-2.372,-1.575) (-2.419,-1.600) (-2.466,-1.625) (-2.514,-1.650) (-2.563,-1.675) (-2.613,-1.700) (-2.664,-1.725) (-2.716,-1.750) (-2.770,-1.775) (-2.825,-1.800) (-2.875,-1.822) (-2.925,-1.844) (-2.975,-1.865) (-3.025,-1.885) (-3.075,-1.904) (-3.125,-1.923) (-3.175,-1.942) (-3.225,-1.959) (-3.275,-1.976) (-3.325,-1.993) (-3.375,-2.008) (-3.425,-2.023) (-3.475,-2.037) (-3.525,-2.051) (-3.575,-2.064) (-3.625,-2.076) (-3.675,-2.087) (-3.725,-2.097) (-3.775,-2.107) (-3.825,-2.116) (-3.875,-2.124) (-3.925,-2.131) (-3.975,-2.138) (-4.025,-2.143) (-4.075,-2.148) (-4.125,-2.152) (-4.175,-2.155) (-4.225,-2.157) (-4.275,-2.158) (-4.325,-2.158) (-4.375,-2.158) (-4.425,-2.156) (-4.475,-2.154) (-4.525,-2.150) (-4.575,-2.146) (-4.625,-2.140) (-4.675,-2.134) (-4.725,-2.126) (-4.775,-2.117) (-4.825,-2.108) (-4.875,-2.097) (-4.925,-2.085) (-4.975,-2.071) (-5.025,-2.057) (-5.075,-2.041) (-5.125,-2.024) (-5.175,-2.006) (-5.225,-1.986) (-5.275,-1.965) (-5.325,-1.942) (-5.375,-1.918) (-5.425,-1.892) (-5.475,-1.864) (-5.525,-1.835) (-5.575,-1.804) (-5.618,-1.775) (-5.675,-1.735) (-5.722,-1.700) (-5.775,-1.658) (-5.814,-1.625) (-5.869,-1.575) (-5.921,-1.525) (-5.970,-1.475) (-6.016,-1.425) (-6.050,-1.385) (-6.098,-1.325) (-6.136,-1.275) (-6.172,-1.225) (-6.200,-1.183) (-6.236,-1.125) (-6.266,-1.075) (-6.294,-1.025) (-6.320,-0.975) (-6.344,-0.925) (-6.367,-0.875) (-6.388,-0.825) (-6.408,-0.775) (-6.426,-0.725) (-6.443,-0.675) (-6.459,-0.625) (-6.473,-0.575) (-6.486,-0.525) (-6.498,-0.475) (-6.509,-0.425) (-6.518,-0.375) (-6.526,-0.325) (-6.533,-0.275) (-6.539,-0.225) (-6.543,-0.175) (-6.547,-0.125) (-6.549,-0.075) (-6.550,-0.025) (-6.550,0.037) (-6.548,0.100) (-6.545,0.150) (-6.541,0.200) (-6.536,0.250) (-6.530,0.300) (-6.522,0.350) (-6.513,0.400) (-6.503,0.450) (-6.492,0.500) (-6.480,0.550) (-6.466,0.600) (-6.451,0.650) (-6.435,0.700) (-6.417,0.750) (-6.398,0.800) (-6.378,0.850) (-6.356,0.900) (-6.332,0.950) (-6.307,1.000) (-6.280,1.050) (-6.251,1.100) (-6.225,1.143) (-6.189,1.200) (-6.154,1.250) (-6.118,1.300) (-6.079,1.350) (-6.037,1.400) (-6.000,1.442) (-5.950,1.496) (-5.900,1.546) (-5.850,1.593) (-5.800,1.637) (-5.754,1.675) (-5.700,1.717) (-5.654,1.750) (-5.600,1.787) (-5.550,1.819) (-5.500,1.850) (-5.456,1.875) (-5.409,1.900) (-5.360,1.925) (-5.307,1.950) (-5.251,1.975) (-5.200,1.996) (-5.150,2.015) (-5.100,2.033) (-5.050,2.049) (-5.000,2.064) (-4.950,2.078) (-4.900,2.091) (-4.850,2.102) (-4.800,2.113) (-4.750,2.122) (-4.700,2.130) (-4.650,2.137) (-4.600,2.143) (-4.550,2.148) (-4.500,2.152) (-4.450,2.155) (-4.400,2.157) (-4.350,2.158) (-4.300,2.158) (-4.250,2.158) (-4.200,2.156) (-4.150,2.153) (-4.100,2.150) (-4.050,2.146) (-4.000,2.140) (-3.950,2.134) (-3.900,2.128) (-3.850,2.120) (-3.800,2.111) (-3.750,2.102) (-3.700,2.092) (-3.650,2.081) (-3.600,2.070) (-3.550,2.057) (-3.500,2.044) (-3.450,2.030) (-3.400,2.016) (-3.350,2.001) (-3.300,1.985) (-3.250,1.968) (-3.200,1.951) (-3.150,1.933) (-3.100,1.914) (-3.050,1.895) (-3.001,1.875) (-2.950,1.854) (-2.900,1.833) (-2.850,1.811) (-2.800,1.789) (-2.750,1.766) (-2.700,1.742) (-2.650,1.718) (-2.600,1.694) (-2.550,1.669) (-2.500,1.643) (-2.450,1.617) (-2.400,1.590) (-2.350,1.563) (-2.300,1.535) (-2.250,1.507) (-2.200,1.479) (-2.150,1.450) (-2.100,1.421) (-2.050,1.391) (-2.000,1.361) (-1.950,1.331) (-1.900,1.300) (-1.850,1.269) (-1.800,1.238) (-1.750,1.206) (-1.702,1.175) (-1.650,1.142) (-1.600,1.109) (-1.550,1.076) (-1.500,1.044) (-1.450,1.010) (-1.400,0.977) (-1.350,0.943) (-1.300,0.910) (-1.250,0.876) (-1.200,0.842) (-1.150,0.807) (-1.103,0.775) (-1.050,0.739) (-1.000,0.704) (-0.950,0.669) (-0.900,0.634) (-0.851,0.600) (-0.800,0.565) (-0.750,0.529) (-0.709,0.500) (-0.650,0.459) (-0.602,0.425) (-0.550,0.389) (-0.500,0.354) (-0.450,0.319) (-0.400,0.283) (-0.354,0.250) (-0.300,0.213) (-0.250,0.177) (-0.220,0.155)}
\EQ{(0.206,0.144) (0.248,0.175) (0.300,0.213) (0.350,0.248) (0.400,0.283) (0.450,0.319) (0.496,0.350) (0.550,0.389) (0.600,0.424) (0.650,0.459) (0.700,0.494) (0.744,0.525) (0.800,0.565) (0.850,0.599) (0.900,0.634) (0.950,0.669) (0.995,0.700) (1.050,0.739) (1.100,0.773) (1.150,0.807) (1.200,0.842) (1.249,0.875) (1.300,0.910) (1.350,0.943) (1.397,0.975) (1.450,1.010) (1.500,1.044) (1.548,1.075) (1.600,1.109) (1.650,1.142) (1.700,1.174) (1.750,1.206) (1.800,1.238) (1.850,1.269) (1.900,1.300) (1.950,1.331) (2.000,1.361) (2.050,1.391) (2.100,1.421) (2.150,1.450) (2.193,1.475) (2.237,1.500) (2.282,1.525) (2.327,1.550) (2.372,1.575) (2.419,1.600) (2.466,1.625) (2.514,1.650) (2.563,1.675) (2.613,1.700) (2.664,1.725) (2.716,1.750) (2.770,1.775) (2.825,1.800) (2.875,1.822) (2.925,1.844) (2.975,1.865) (3.025,1.885) (3.075,1.904) (3.125,1.923) (3.175,1.942) (3.225,1.959) (3.275,1.976) (3.325,1.993) (3.375,2.008) (3.425,2.023) (3.475,2.037) (3.525,2.051) (3.575,2.064) (3.625,2.076) (3.675,2.087) (3.725,2.097) (3.775,2.107) (3.825,2.116) (3.875,2.124) (3.925,2.131) (3.975,2.138) (4.025,2.143) (4.075,2.148) (4.125,2.152) (4.175,2.155) (4.225,2.157) (4.275,2.158) (4.325,2.158) (4.375,2.158) (4.425,2.156) (4.475,2.154) (4.525,2.150) (4.575,2.146) (4.625,2.140) (4.675,2.134) (4.725,2.126) (4.775,2.117) (4.825,2.108) (4.875,2.097) (4.925,2.085) (4.975,2.071) (5.025,2.057) (5.075,2.041) (5.125,2.024) (5.175,2.006) (5.225,1.986) (5.275,1.965) (5.325,1.942) (5.375,1.918) (5.425,1.892) (5.475,1.864) (5.525,1.835) (5.575,1.804) (5.618,1.775) (5.675,1.735) (5.722,1.700) (5.775,1.658) (5.814,1.625) (5.869,1.575) (5.921,1.525) (5.970,1.475) (6.016,1.425) (6.050,1.385) (6.098,1.325) (6.136,1.275) (6.172,1.225) (6.200,1.183) (6.236,1.125) (6.266,1.075) (6.294,1.025) (6.320,0.975) (6.344,0.925) (6.367,0.875) (6.388,0.825) (6.408,0.775) (6.426,0.725) (6.443,0.675) (6.459,0.625) (6.473,0.575) (6.486,0.525) (6.498,0.475) (6.509,0.425) (6.518,0.375) (6.526,0.325) (6.533,0.275) (6.539,0.225) (6.543,0.175) (6.547,0.125) (6.549,0.075) (6.550,0.025) (6.550,-0.037) (6.548,-0.100) (6.545,-0.150) (6.541,-0.200) (6.536,-0.250) (6.530,-0.300) (6.522,-0.350) (6.513,-0.400) (6.503,-0.450) (6.492,-0.500) (6.480,-0.550) (6.466,-0.600) (6.451,-0.650) (6.435,-0.700) (6.417,-0.750) (6.398,-0.800) (6.378,-0.850) (6.356,-0.900) (6.332,-0.950) (6.307,-1.000) (6.280,-1.050) (6.251,-1.100) (6.225,-1.143) (6.189,-1.200) (6.154,-1.250) (6.118,-1.300) (6.079,-1.350) (6.037,-1.400) (6.000,-1.442) (5.950,-1.496) (5.900,-1.546) (5.850,-1.593) (5.800,-1.637) (5.754,-1.675) (5.700,-1.717) (5.654,-1.750) (5.600,-1.787) (5.550,-1.819) (5.500,-1.850) (5.456,-1.875) (5.409,-1.900) (5.360,-1.925) (5.307,-1.950) (5.251,-1.975) (5.200,-1.996) (5.150,-2.015) (5.100,-2.033) (5.050,-2.049) (5.000,-2.064) (4.950,-2.078) (4.900,-2.091) (4.850,-2.102) (4.800,-2.113) (4.750,-2.122) (4.700,-2.130) (4.650,-2.137) (4.600,-2.143) (4.550,-2.148) (4.500,-2.152) (4.450,-2.155) (4.400,-2.157) (4.350,-2.158) (4.300,-2.158) (4.250,-2.158) (4.200,-2.156) (4.150,-2.153) (4.100,-2.150) (4.050,-2.146) (4.000,-2.140) (3.950,-2.134) (3.900,-2.128) (3.850,-2.120) (3.800,-2.111) (3.750,-2.102) (3.700,-2.092) (3.650,-2.081) (3.600,-2.070) (3.550,-2.057) (3.500,-2.044) (3.450,-2.030) (3.400,-2.016) (3.350,-2.001) (3.300,-1.985) (3.250,-1.968) (3.200,-1.951) (3.150,-1.933) (3.100,-1.914) (3.050,-1.895) (3.001,-1.875) (2.950,-1.854) (2.900,-1.833) (2.850,-1.811) (2.800,-1.789) (2.750,-1.766) (2.700,-1.742) (2.650,-1.718) (2.600,-1.694) (2.550,-1.669) (2.500,-1.643) (2.450,-1.617) (2.400,-1.590) (2.350,-1.563) (2.300,-1.535) (2.250,-1.507) (2.200,-1.479) (2.150,-1.450) (2.100,-1.421) (2.050,-1.391) (2.000,-1.361) (1.950,-1.331) (1.900,-1.300) (1.850,-1.269) (1.800,-1.238) (1.750,-1.206) (1.702,-1.175) (1.650,-1.142) (1.600,-1.109) (1.550,-1.076) (1.500,-1.044) (1.450,-1.010) (1.400,-0.977) (1.350,-0.943) (1.300,-0.910) (1.250,-0.876) (1.200,-0.842) (1.150,-0.807) (1.103,-0.775) (1.050,-0.739) (1.000,-0.704) (0.950,-0.669) (0.900,-0.634) (0.851,-0.600) (0.800,-0.565) (0.750,-0.529) (0.709,-0.500) (0.650,-0.459) (0.602,-0.425) (0.550,-0.389) (0.500,-0.354) (0.450,-0.319) (0.400,-0.283) (0.354,-0.250) (0.300,-0.213) (0.250,-0.177) (0.206,-0.144)}
\EQ{(-0.220,-0.155) (-0.275,-0.195) (-0.319,-0.225) (-0.375,-0.266) (-0.424,-0.300) (-0.475,-0.336) (-0.525,-0.371) (-0.567,-0.400) (-0.625,-0.442) (-0.673,-0.475) (-0.725,-0.512) (-0.775,-0.547) (-0.825,-0.583) (-0.875,-0.618) (-0.922,-0.650) (-0.975,-0.688) (-1.025,-0.722) (-1.075,-0.757) (-1.125,-0.792) (-1.173,-0.825) (-1.225,-0.861) (-1.275,-0.896) (-1.318,-0.925) (-1.375,-0.964) (-1.425,-0.998) (-1.475,-1.032) (-1.525,-1.066) (-1.575,-1.100) (-1.625,-1.133) (-1.675,-1.167) (-1.725,-1.200) (-1.775,-1.233) (-1.825,-1.266) (-1.875,-1.298) (-1.925,-1.330) (-1.975,-1.362) (-2.025,-1.394) (-2.074,-1.425) (-2.125,-1.457) (-2.175,-1.488) (-2.225,-1.519) (-2.275,-1.549) (-2.325,-1.580) (-2.375,-1.609) (-2.425,-1.639) (-2.475,-1.668) (-2.525,-1.697) (-2.575,-1.725) (-2.619,-1.750) (-2.665,-1.775) (-2.710,-1.800) (-2.757,-1.825) (-2.803,-1.850) (-2.850,-1.874) (-2.900,-1.900) (-2.949,-1.925) (-2.999,-1.950) (-3.050,-1.975) (-3.100,-1.999) (-3.150,-2.022) (-3.200,-2.045) (-3.250,-2.068) (-3.300,-2.090) (-3.350,-2.111) (-3.400,-2.132) (-3.450,-2.152) (-3.500,-2.172) (-3.550,-2.191) (-3.600,-2.210) (-3.650,-2.228) (-3.700,-2.246) (-3.750,-2.263) (-3.800,-2.279) (-3.850,-2.295) (-3.900,-2.310) (-3.950,-2.324) (-4.000,-2.338) (-4.050,-2.352) (-4.100,-2.364) (-4.150,-2.376) (-4.200,-2.387) (-4.250,-2.398) (-4.300,-2.408) (-4.350,-2.417) (-4.400,-2.425) (-4.450,-2.433) (-4.500,-2.440) (-4.550,-2.447) (-4.600,-2.452) (-4.650,-2.457) (-4.700,-2.462) (-4.750,-2.465) (-4.800,-2.468) (-4.850,-2.469) (-4.900,-2.470) (-4.950,-2.471) (-5.000,-2.470) (-5.050,-2.469) (-5.100,-2.467) (-5.150,-2.464) (-5.200,-2.460) (-5.250,-2.455) (-5.300,-2.449) (-5.350,-2.443) (-5.400,-2.435) (-5.450,-2.427) (-5.500,-2.417) (-5.550,-2.407) (-5.600,-2.395) (-5.650,-2.383) (-5.700,-2.370) (-5.750,-2.355) (-5.800,-2.339) (-5.850,-2.323) (-5.900,-2.305) (-5.950,-2.286) (-6.000,-2.265) (-6.050,-2.244) (-6.100,-2.221) (-6.150,-2.196) (-6.200,-2.171) (-6.250,-2.144) (-6.300,-2.115) (-6.350,-2.085) (-6.400,-2.053) (-6.450,-2.019) (-6.500,-1.983) (-6.545,-1.950) (-6.600,-1.906) (-6.650,-1.864) (-6.695,-1.825) (-6.748,-1.775) (-6.799,-1.725) (-6.846,-1.675) (-6.892,-1.625) (-6.925,-1.586) (-6.974,-1.525) (-7.013,-1.475) (-7.049,-1.425) (-7.083,-1.375) (-7.116,-1.325) (-7.147,-1.275) (-7.175,-1.226) (-7.200,-1.181) (-7.225,-1.134) (-7.250,-1.084) (-7.275,-1.031) (-7.300,-0.975) (-7.320,-0.925) (-7.340,-0.875) (-7.358,-0.825) (-7.375,-0.775) (-7.391,-0.725) (-7.405,-0.675) (-7.419,-0.625) (-7.431,-0.575) (-7.443,-0.525) (-7.453,-0.475) (-7.462,-0.425) (-7.470,-0.375) (-7.477,-0.325) (-7.483,-0.275) (-7.488,-0.225) (-7.492,-0.175) (-7.495,-0.125) (-7.497,-0.075) (-7.498,-0.025) (-7.498,0.050) (-7.496,0.100) (-7.494,0.150) (-7.490,0.200) (-7.486,0.250) (-7.480,0.300) (-7.474,0.350) (-7.466,0.400) (-7.457,0.450) (-7.448,0.500) (-7.437,0.550) (-7.425,0.600) (-7.412,0.650) (-7.398,0.700) (-7.383,0.750) (-7.367,0.800) (-7.349,0.850) (-7.330,0.900) (-7.310,0.950) (-7.289,1.000) (-7.266,1.050) (-7.242,1.100) (-7.217,1.150) (-7.190,1.200) (-7.161,1.250) (-7.131,1.300) (-7.100,1.350) (-7.066,1.400) (-7.031,1.450) (-7.000,1.492) (-6.955,1.550) (-6.913,1.600) (-6.875,1.644) (-6.825,1.698) (-6.775,1.749) (-6.725,1.797) (-6.675,1.843) (-6.625,1.886) (-6.577,1.925) (-6.525,1.965) (-6.477,2.000) (-6.425,2.036) (-6.375,2.069) (-6.325,2.100) (-6.275,2.129) (-6.225,2.157) (-6.175,2.184) (-6.125,2.209) (-6.075,2.232) (-6.025,2.255) (-5.977,2.275) (-5.925,2.295) (-5.875,2.314) (-5.825,2.331) (-5.775,2.347) (-5.725,2.362) (-5.675,2.376) (-5.625,2.389) (-5.575,2.401) (-5.525,2.412) (-5.475,2.422) (-5.425,2.431) (-5.375,2.439) (-5.325,2.446) (-5.275,2.452) (-5.225,2.457) (-5.175,2.462) (-5.125,2.465) (-5.075,2.468) (-5.025,2.470) (-4.975,2.471) (-4.925,2.471) (-4.875,2.470) (-4.825,2.469) (-4.775,2.466) (-4.725,2.463) (-4.675,2.460) (-4.625,2.455) (-4.575,2.450) (-4.525,2.444) (-4.475,2.437) (-4.425,2.430) (-4.375,2.421) (-4.325,2.412) (-4.275,2.403) (-4.225,2.393) (-4.175,2.382) (-4.125,2.370) (-4.075,2.358) (-4.025,2.345) (-3.975,2.331) (-3.925,2.317) (-3.875,2.302) (-3.825,2.287) (-3.775,2.271) (-3.725,2.254) (-3.675,2.237) (-3.625,2.219) (-3.575,2.201) (-3.525,2.182) (-3.475,2.162) (-3.425,2.142) (-3.375,2.122) (-3.325,2.100) (-3.275,2.079) (-3.225,2.057) (-3.175,2.034) (-3.125,2.011) (-3.075,1.987) (-3.025,1.963) (-2.975,1.938) (-2.925,1.913) (-2.875,1.887) (-2.825,1.861) (-2.775,1.835) (-2.725,1.808) (-2.675,1.781) (-2.625,1.753) (-2.575,1.725) (-2.531,1.700) (-2.487,1.675) (-2.425,1.639) (-2.375,1.609) (-2.325,1.580) (-2.276,1.550) (-2.225,1.519) (-2.175,1.488) (-2.125,1.457) (-2.075,1.426) (-2.025,1.394) (-1.975,1.362) (-1.925,1.330) (-1.878,1.300) (-1.825,1.266) (-1.775,1.233) (-1.725,1.200) (-1.675,1.167) (-1.625,1.133) (-1.576,1.100) (-1.525,1.066) (-1.475,1.032) (-1.428,1.000) (-1.375,0.964) (-1.325,0.930) (-1.282,0.900) (-1.225,0.861) (-1.175,0.827) (-1.125,0.792) (-1.075,0.757) (-1.029,0.725) (-0.975,0.688) (-0.925,0.652) (-0.875,0.618) (-0.825,0.583) (-0.779,0.550) (-0.725,0.512) (-0.675,0.477) (-0.625,0.442) (-0.575,0.407) (-0.531,0.375) (-0.475,0.336) (-0.425,0.300) (-0.375,0.266) (-0.325,0.230) (-0.284,0.200) (-0.225,0.160) (-0.220,0.155)}
\EQ{(0.206,0.144) (0.248,0.175) (0.300,0.213) (0.350,0.248) (0.400,0.283) (0.450,0.319) (0.496,0.350) (0.550,0.389) (0.600,0.424) (0.650,0.460) (0.700,0.495) (0.744,0.525) (0.800,0.565) (0.850,0.600) (0.900,0.635) (0.950,0.670) (0.993,0.700) (1.050,0.740) (1.100,0.774) (1.150,0.809) (1.200,0.844) (1.245,0.875) (1.300,0.913) (1.350,0.947) (1.400,0.981) (1.450,1.015) (1.500,1.049) (1.550,1.083) (1.600,1.117) (1.650,1.150) (1.700,1.183) (1.750,1.216) (1.800,1.249) (1.850,1.282) (1.900,1.314) (1.950,1.346) (1.995,1.375) (2.050,1.410) (2.100,1.442) (2.150,1.473) (2.194,1.500) (2.250,1.534) (2.300,1.565) (2.350,1.594) (2.400,1.624) (2.444,1.650) (2.500,1.682) (2.550,1.711) (2.600,1.739) (2.650,1.767) (2.700,1.795) (2.750,1.822) (2.800,1.848) (2.850,1.874) (2.900,1.900) (2.949,1.925) (2.999,1.950) (3.050,1.975) (3.100,1.999) (3.150,2.022) (3.200,2.045) (3.250,2.068) (3.300,2.090) (3.350,2.111) (3.400,2.132) (3.450,2.152) (3.500,2.172) (3.550,2.191) (3.600,2.210) (3.650,2.228) (3.700,2.246) (3.750,2.263) (3.800,2.279) (3.850,2.295) (3.900,2.310) (3.950,2.324) (4.000,2.338) (4.050,2.352) (4.100,2.364) (4.150,2.376) (4.200,2.387) (4.250,2.398) (4.300,2.408) (4.350,2.417) (4.400,2.425) (4.450,2.433) (4.500,2.440) (4.550,2.447) (4.600,2.452) (4.650,2.457) (4.700,2.462) (4.750,2.465) (4.800,2.468) (4.850,2.469) (4.900,2.470) (4.950,2.471) (5.000,2.470) (5.050,2.469) (5.100,2.467) (5.150,2.464) (5.200,2.460) (5.250,2.455) (5.300,2.449) (5.350,2.443) (5.400,2.435) (5.450,2.427) (5.500,2.417) (5.550,2.407) (5.600,2.395) (5.650,2.383) (5.700,2.370) (5.750,2.355) (5.800,2.339) (5.850,2.323) (5.900,2.305) (5.950,2.286) (6.000,2.265) (6.050,2.244) (6.100,2.221) (6.150,2.196) (6.200,2.171) (6.250,2.144) (6.300,2.115) (6.350,2.085) (6.400,2.053) (6.450,2.019) (6.500,1.983) (6.545,1.950) (6.600,1.906) (6.650,1.864) (6.695,1.825) (6.748,1.775) (6.799,1.725) (6.846,1.675) (6.892,1.625) (6.925,1.586) (6.974,1.525) (7.013,1.475) (7.049,1.425) (7.083,1.375) (7.116,1.325) (7.147,1.275) (7.175,1.226) (7.200,1.181) (7.225,1.134) (7.250,1.084) (7.275,1.031) (7.300,0.975) (7.320,0.925) (7.340,0.875) (7.358,0.825) (7.375,0.775) (7.391,0.725) (7.405,0.675) (7.419,0.625) (7.431,0.575) (7.443,0.525) (7.453,0.475) (7.462,0.425) (7.470,0.375) (7.477,0.325) (7.483,0.275) (7.488,0.225) (7.492,0.175) (7.495,0.125) (7.497,0.075) (7.498,0.025) (7.498,-0.050) (7.496,-0.100) (7.494,-0.150) (7.490,-0.200) (7.486,-0.250) (7.480,-0.300) (7.474,-0.350) (7.466,-0.400) (7.457,-0.450) (7.448,-0.500) (7.437,-0.550) (7.425,-0.600) (7.412,-0.650) (7.398,-0.700) (7.383,-0.750) (7.367,-0.800) (7.349,-0.850) (7.330,-0.900) (7.310,-0.950) (7.289,-1.000) (7.266,-1.050) (7.242,-1.100) (7.217,-1.150) (7.190,-1.200) (7.161,-1.250) (7.131,-1.300) (7.100,-1.350) (7.066,-1.400) (7.031,-1.450) (7.000,-1.492) (6.955,-1.550) (6.913,-1.600) (6.875,-1.644) (6.825,-1.698) (6.775,-1.749) (6.725,-1.797) (6.675,-1.843) (6.625,-1.886) (6.577,-1.925) (6.525,-1.965) (6.477,-2.000) (6.425,-2.036) (6.375,-2.069) (6.325,-2.100) (6.275,-2.129) (6.225,-2.157) (6.175,-2.184) (6.125,-2.209) (6.075,-2.232) (6.025,-2.255) (5.977,-2.275) (5.925,-2.295) (5.875,-2.314) (5.825,-2.331) (5.775,-2.347) (5.725,-2.362) (5.675,-2.376) (5.625,-2.389) (5.575,-2.401) (5.525,-2.412) (5.475,-2.422) (5.425,-2.431) (5.375,-2.439) (5.325,-2.446) (5.275,-2.452) (5.225,-2.457) (5.175,-2.462) (5.125,-2.465) (5.075,-2.468) (5.025,-2.470) (4.975,-2.471) (4.925,-2.471) (4.875,-2.470) (4.825,-2.469) (4.775,-2.466) (4.725,-2.463) (4.675,-2.460) (4.625,-2.455) (4.575,-2.450) (4.525,-2.444) (4.475,-2.437) (4.425,-2.430) (4.375,-2.421) (4.325,-2.412) (4.275,-2.403) (4.225,-2.393) (4.175,-2.382) (4.125,-2.370) (4.075,-2.358) (4.025,-2.345) (3.975,-2.331) (3.925,-2.317) (3.875,-2.302) (3.825,-2.287) (3.775,-2.271) (3.725,-2.254) (3.675,-2.237) (3.625,-2.219) (3.575,-2.201) (3.525,-2.182) (3.475,-2.162) (3.425,-2.142) (3.375,-2.122) (3.325,-2.100) (3.275,-2.079) (3.225,-2.057) (3.175,-2.034) (3.125,-2.011) (3.075,-1.987) (3.025,-1.963) (2.975,-1.938) (2.925,-1.913) (2.875,-1.887) (2.825,-1.861) (2.775,-1.835) (2.725,-1.808) (2.675,-1.781) (2.625,-1.753) (2.575,-1.725) (2.531,-1.700) (2.487,-1.675) (2.425,-1.639) (2.375,-1.609) (2.325,-1.580) (2.276,-1.550) (2.225,-1.519) (2.175,-1.488) (2.125,-1.457) (2.075,-1.426) (2.025,-1.394) (1.975,-1.362) (1.925,-1.330) (1.878,-1.300) (1.825,-1.266) (1.775,-1.233) (1.725,-1.200) (1.675,-1.167) (1.625,-1.133) (1.576,-1.100) (1.525,-1.066) (1.475,-1.032) (1.428,-1.000) (1.375,-0.964) (1.325,-0.930) (1.282,-0.900) (1.225,-0.861) (1.175,-0.827) (1.125,-0.792) (1.075,-0.757) (1.029,-0.725) (0.975,-0.688) (0.925,-0.652) (0.875,-0.618) (0.825,-0.583) (0.779,-0.550) (0.725,-0.512) (0.675,-0.477) (0.625,-0.442) (0.575,-0.407) (0.531,-0.375) (0.475,-0.336) (0.425,-0.300) (0.375,-0.266) (0.325,-0.230) (0.284,-0.200) (0.225,-0.160) (0.206,-0.144)}
\EQ{(-0.220,-0.155) (-0.275,-0.195) (-0.319,-0.225) (-0.375,-0.266) (-0.424,-0.300) (-0.475,-0.336) (-0.525,-0.371) (-0.567,-0.400) (-0.625,-0.442) (-0.672,-0.475) (-0.725,-0.513) (-0.775,-0.548) (-0.825,-0.583) (-0.875,-0.618) (-0.921,-0.650) (-0.975,-0.689) (-1.025,-0.724) (-1.075,-0.759) (-1.125,-0.794) (-1.170,-0.825) (-1.225,-0.864) (-1.275,-0.899) (-1.325,-0.934) (-1.375,-0.968) (-1.421,-1.000) (-1.475,-1.038) (-1.525,-1.072) (-1.575,-1.107) (-1.625,-1.141) (-1.674,-1.175) (-1.725,-1.210) (-1.775,-1.244) (-1.821,-1.275) (-1.875,-1.312) (-1.925,-1.346) (-1.969,-1.375) (-2.025,-1.413) (-2.075,-1.446) (-2.118,-1.475) (-2.175,-1.513) (-2.225,-1.546) (-2.270,-1.575) (-2.325,-1.611) (-2.375,-1.644) (-2.424,-1.675) (-2.475,-1.708) (-2.525,-1.740) (-2.575,-1.772) (-2.620,-1.800) (-2.675,-1.834) (-2.725,-1.866) (-2.775,-1.896) (-2.822,-1.925) (-2.875,-1.957) (-2.925,-1.988) (-2.975,-2.017) (-3.025,-2.047) (-3.073,-2.075) (-3.125,-2.105) (-3.175,-2.134) (-3.225,-2.163) (-3.275,-2.191) (-3.325,-2.219) (-3.375,-2.246) (-3.425,-2.274) (-3.474,-2.300) (-3.521,-2.325) (-3.569,-2.350) (-3.617,-2.375) (-3.665,-2.400) (-3.715,-2.425) (-3.765,-2.450) (-3.816,-2.475) (-3.867,-2.500) (-3.920,-2.525) (-3.973,-2.550) (-4.025,-2.574) (-4.075,-2.596) (-4.125,-2.619) (-4.175,-2.640) (-4.225,-2.662) (-4.275,-2.683) (-4.325,-2.703) (-4.375,-2.723) (-4.425,-2.743) (-4.475,-2.762) (-4.525,-2.781) (-4.575,-2.799) (-4.625,-2.817) (-4.675,-2.834) (-4.725,-2.851) (-4.775,-2.867) (-4.825,-2.883) (-4.875,-2.899) (-4.925,-2.914) (-4.975,-2.928) (-5.025,-2.942) (-5.075,-2.956) (-5.125,-2.968) (-5.175,-2.981) (-5.225,-2.993) (-5.275,-3.004) (-5.325,-3.015) (-5.375,-3.025) (-5.425,-3.035) (-5.475,-3.044) (-5.525,-3.053) (-5.575,-3.061) (-5.625,-3.069) (-5.675,-3.075) (-5.725,-3.082) (-5.775,-3.088) (-5.825,-3.093) (-5.875,-3.098) (-5.925,-3.102) (-5.975,-3.105) (-6.025,-3.108) (-6.075,-3.110) (-6.125,-3.112) (-6.175,-3.113) (-6.225,-3.113) (-6.275,-3.113) (-6.325,-3.112) (-6.375,-3.110) (-6.425,-3.108) (-6.475,-3.105) (-6.525,-3.101) (-6.575,-3.097) (-6.625,-3.092) (-6.675,-3.086) (-6.725,-3.080) (-6.775,-3.072) (-6.825,-3.065) (-6.875,-3.056) (-6.925,-3.046) (-6.975,-3.036) (-7.025,-3.025) (-7.075,-3.013) (-7.125,-3.001) (-7.175,-2.987) (-7.225,-2.973) (-7.275,-2.958) (-7.325,-2.942) (-7.375,-2.925) (-7.425,-2.907) (-7.475,-2.888) (-7.525,-2.868) (-7.575,-2.848) (-7.625,-2.826) (-7.675,-2.803) (-7.725,-2.779) (-7.775,-2.754) (-7.825,-2.728) (-7.875,-2.701) (-7.920,-2.675) (-7.975,-2.642) (-8.025,-2.611) (-8.075,-2.579) (-8.117,-2.550) (-8.175,-2.509) (-8.222,-2.475) (-8.275,-2.434) (-8.318,-2.400) (-8.375,-2.352) (-8.425,-2.308) (-8.475,-2.262) (-8.514,-2.225) (-8.550,-2.189) (-8.589,-2.150) (-8.625,-2.111) (-8.675,-2.056) (-8.723,-2.000) (-8.765,-1.950) (-8.800,-1.905) (-8.842,-1.850) (-8.875,-1.804) (-8.912,-1.750) (-8.946,-1.700) (-8.975,-1.654) (-9.008,-1.600) (-9.037,-1.550) (-9.065,-1.500) (-9.092,-1.450) (-9.118,-1.400) (-9.142,-1.350) (-9.165,-1.300) (-9.188,-1.250) (-9.209,-1.200) (-9.229,-1.150) (-9.248,-1.100) (-9.267,-1.050) (-9.284,-1.000) (-9.300,-0.951) (-9.316,-0.900) (-9.330,-0.850) (-9.344,-0.800) (-9.357,-0.750) (-9.368,-0.700) (-9.379,-0.650) (-9.390,-0.600) (-9.399,-0.550) (-9.407,-0.500) (-9.415,-0.450) (-9.422,-0.400) (-9.428,-0.350) (-9.433,-0.300) (-9.437,-0.250) (-9.441,-0.200) (-9.444,-0.150) (-9.446,-0.100) (-9.447,-0.050) (-9.447,0.000) (-9.447,0.050) (-9.446,0.100) (-9.444,0.150) (-9.441,0.200) (-9.437,0.250) (-9.433,0.300) (-9.428,0.350) (-9.422,0.400) (-9.415,0.450) (-9.407,0.500) (-9.399,0.550) (-9.390,0.600) (-9.379,0.650) (-9.368,0.700) (-9.357,0.750) (-9.344,0.800) (-9.330,0.850) (-9.316,0.900) (-9.300,0.950) (-9.284,1.000) (-9.267,1.050) (-9.248,1.100) (-9.229,1.150) (-9.209,1.200) (-9.188,1.250) (-9.165,1.300) (-9.142,1.350) (-9.118,1.400) (-9.092,1.450) (-9.065,1.500) (-9.037,1.550) (-9.008,1.600) (-8.978,1.650) (-8.950,1.693) (-8.912,1.750) (-8.878,1.800) (-8.842,1.850) (-8.804,1.900) (-8.765,1.950) (-8.725,1.998) (-8.680,2.050) (-8.636,2.100) (-8.600,2.138) (-8.564,2.175) (-8.525,2.214) (-8.488,2.250) (-8.450,2.285) (-8.406,2.325) (-8.350,2.373) (-8.300,2.414) (-8.255,2.450) (-8.200,2.491) (-8.153,2.525) (-8.100,2.562) (-8.050,2.595) (-8.003,2.625) (-7.950,2.657) (-7.900,2.686) (-7.850,2.714) (-7.800,2.741) (-7.750,2.767) (-7.700,2.791) (-7.650,2.814) (-7.600,2.837) (-7.550,2.858) (-7.500,2.878) (-7.450,2.898) (-7.400,2.916) (-7.350,2.933) (-7.300,2.950) (-7.250,2.966) (-7.200,2.980) (-7.150,2.994) (-7.100,3.007) (-7.050,3.019) (-7.000,3.031) (-6.950,3.041) (-6.900,3.051) (-6.850,3.060) (-6.800,3.069) (-6.750,3.076) (-6.700,3.083) (-6.650,3.089) (-6.600,3.094) (-6.550,3.099) (-6.500,3.103) (-6.450,3.106) (-6.400,3.109) (-6.350,3.111) (-6.300,3.112) (-6.250,3.113) (-6.200,3.113) (-6.150,3.112) (-6.100,3.111) (-6.050,3.109) (-6.000,3.107) (-5.950,3.103) (-5.900,3.100) (-5.850,3.095) (-5.800,3.090) (-5.750,3.085) (-5.700,3.079) (-5.650,3.072) (-5.600,3.065) (-5.550,3.057) (-5.500,3.049) (-5.450,3.040) (-5.400,3.030) (-5.350,3.020) (-5.300,3.010) (-5.250,2.998) (-5.200,2.987) (-5.151,2.975) (-5.100,2.962) (-5.050,2.949) (-5.000,2.935) (-4.950,2.921) (-4.900,2.906) (-4.850,2.891) (-4.800,2.875) (-4.750,2.859) (-4.700,2.843) (-4.650,2.825) (-4.600,2.808) (-4.550,2.790) (-4.500,2.771) (-4.450,2.752) (-4.400,2.733) (-4.350,2.713) (-4.300,2.693) (-4.250,2.672) (-4.200,2.651) (-4.150,2.629) (-4.100,2.608) (-4.050,2.585) (-4.000,2.562) (-3.950,2.539) (-3.900,2.516) (-3.850,2.492) (-3.800,2.467) (-3.750,2.443) (-3.700,2.418) (-3.650,2.392) (-3.600,2.366) (-3.550,2.340) (-3.500,2.314) (-3.450,2.287) (-3.400,2.260) (-3.350,2.233) (-3.300,2.205) (-3.250,2.177) (-3.203,2.150) (-3.159,2.125) (-3.116,2.100) (-3.050,2.062) (-3.000,2.032) (-2.950,2.002) (-2.904,1.975) (-2.850,1.942) (-2.800,1.912) (-2.750,1.881) (-2.700,1.850) (-2.650,1.819) (-2.600,1.787) (-2.550,1.756) (-2.502,1.725) (-2.450,1.692) (-2.400,1.660) (-2.350,1.627) (-2.300,1.595) (-2.250,1.562) (-2.200,1.529) (-2.156,1.500) (-2.100,1.463) (-2.050,1.430) (-2.006,1.400) (-1.950,1.363) (-1.900,1.329) (-1.858,1.300) (-1.800,1.261) (-1.750,1.227) (-1.700,1.193) (-1.650,1.158) (-1.602,1.125) (-1.550,1.090) (-1.500,1.055) (-1.457,1.025) (-1.400,0.986) (-1.350,0.951) (-1.300,0.916) (-1.250,0.881) (-1.206,0.850) (-1.150,0.811) (-1.100,0.776) (-1.050,0.741) (-1.000,0.706) (-0.956,0.675) (-0.900,0.636) (-0.850,0.600) (-0.800,0.566) (-0.750,0.530) (-0.708,0.500) (-0.650,0.460) (-0.601,0.425) (-0.550,0.389) (-0.500,0.354) (-0.450,0.319) (-0.400,0.283) (-0.354,0.250) (-0.300,0.213) (-0.250,0.177) (-0.220,0.155)}
\EQ{(0.206,0.144) (0.248,0.175) (0.300,0.213) (0.350,0.248) (0.400,0.283) (0.450,0.319) (0.496,0.350) (0.550,0.389) (0.600,0.424) (0.650,0.460) (0.700,0.495) (0.743,0.525) (0.800,0.566) (0.849,0.600) (0.900,0.636) (0.950,0.671) (0.992,0.700) (1.050,0.741) (1.098,0.775) (1.150,0.811) (1.200,0.846) (1.242,0.875) (1.300,0.916) (1.349,0.950) (1.400,0.986) (1.450,1.020) (1.493,1.050) (1.550,1.090) (1.600,1.124) (1.650,1.158) (1.700,1.193) (1.748,1.225) (1.800,1.261) (1.850,1.295) (1.895,1.325) (1.950,1.363) (2.000,1.396) (2.043,1.425) (2.100,1.463) (2.150,1.496) (2.194,1.525) (2.250,1.562) (2.300,1.595) (2.347,1.625) (2.400,1.660) (2.450,1.692) (2.500,1.724) (2.550,1.756) (2.600,1.787) (2.650,1.819) (2.700,1.850) (2.750,1.881) (2.800,1.912) (2.850,1.942) (2.900,1.972) (2.946,2.000) (3.000,2.032) (3.050,2.062) (3.100,2.091) (3.150,2.120) (3.200,2.148) (3.247,2.175) (3.292,2.200) (3.337,2.225) (3.382,2.250) (3.428,2.275) (3.474,2.300) (3.521,2.325) (3.569,2.350) (3.617,2.375) (3.665,2.400) (3.715,2.425) (3.765,2.450) (3.816,2.475) (3.867,2.500) (3.920,2.525) (3.973,2.550) (4.025,2.574) (4.075,2.596) (4.125,2.619) (4.175,2.640) (4.225,2.662) (4.275,2.683) (4.325,2.703) (4.375,2.723) (4.425,2.743) (4.475,2.762) (4.525,2.781) (4.575,2.799) (4.625,2.817) (4.675,2.834) (4.725,2.851) (4.775,2.867) (4.825,2.883) (4.875,2.899) (4.925,2.914) (4.975,2.928) (5.025,2.942) (5.075,2.956) (5.125,2.968) (5.175,2.981) (5.225,2.993) (5.275,3.004) (5.325,3.015) (5.375,3.025) (5.425,3.035) (5.475,3.044) (5.525,3.053) (5.575,3.061) (5.625,3.069) (5.675,3.075) (5.725,3.082) (5.775,3.088) (5.825,3.093) (5.875,3.098) (5.925,3.102) (5.975,3.105) (6.025,3.108) (6.075,3.110) (6.125,3.112) (6.175,3.113) (6.225,3.113) (6.275,3.113) (6.325,3.112) (6.375,3.110) (6.425,3.108) (6.475,3.105) (6.525,3.101) (6.575,3.097) (6.625,3.092) (6.675,3.086) (6.725,3.080) (6.775,3.072) (6.825,3.065) (6.875,3.056) (6.925,3.046) (6.975,3.036) (7.025,3.025) (7.075,3.013) (7.125,3.001) (7.175,2.987) (7.225,2.973) (7.275,2.958) (7.325,2.942) (7.375,2.925) (7.425,2.907) (7.475,2.888) (7.525,2.868) (7.575,2.848) (7.625,2.826) (7.675,2.803) (7.725,2.779) (7.775,2.754) (7.825,2.728) (7.875,2.701) (7.920,2.675) (7.975,2.642) (8.025,2.611) (8.075,2.579) (8.117,2.550) (8.175,2.509) (8.222,2.475) (8.275,2.434) (8.318,2.400) (8.375,2.352) (8.425,2.308) (8.475,2.262) (8.514,2.225) (8.550,2.189) (8.589,2.150) (8.625,2.111) (8.675,2.056) (8.723,2.000) (8.765,1.950) (8.800,1.905) (8.842,1.850) (8.875,1.804) (8.912,1.750) (8.946,1.700) (8.975,1.654) (9.008,1.600) (9.037,1.550) (9.065,1.500) (9.092,1.450) (9.118,1.400) (9.142,1.350) (9.165,1.300) (9.188,1.250) (9.209,1.200) (9.229,1.150) (9.248,1.100) (9.267,1.050) (9.284,1.000) (9.300,0.951) (9.316,0.900) (9.330,0.850) (9.344,0.800) (9.357,0.750) (9.368,0.700) (9.379,0.650) (9.390,0.600) (9.399,0.550) (9.407,0.500) (9.415,0.450) (9.422,0.400) (9.428,0.350) (9.433,0.300) (9.437,0.250) (9.441,0.200) (9.444,0.150) (9.446,0.100) (9.447,0.050) (9.447,0.000) (9.447,-0.050) (9.446,-0.100) (9.444,-0.150) (9.441,-0.200) (9.437,-0.250) (9.433,-0.300) (9.428,-0.350) (9.422,-0.400) (9.415,-0.450) (9.407,-0.500) (9.399,-0.550) (9.390,-0.600) (9.379,-0.650) (9.368,-0.700) (9.357,-0.750) (9.344,-0.800) (9.330,-0.850) (9.316,-0.900) (9.300,-0.950) (9.284,-1.000) (9.267,-1.050) (9.248,-1.100) (9.229,-1.150) (9.209,-1.200) (9.188,-1.250) (9.165,-1.300) (9.142,-1.350) (9.118,-1.400) (9.092,-1.450) (9.065,-1.500) (9.037,-1.550) (9.008,-1.600) (8.978,-1.650) (8.950,-1.693) (8.912,-1.750) (8.878,-1.800) (8.842,-1.850) (8.804,-1.900) (8.765,-1.950) (8.725,-1.998) (8.680,-2.050) (8.636,-2.100) (8.600,-2.138) (8.564,-2.175) (8.525,-2.214) (8.488,-2.250) (8.450,-2.285) (8.406,-2.325) (8.350,-2.373) (8.300,-2.414) (8.255,-2.450) (8.200,-2.491) (8.153,-2.525) (8.100,-2.562) (8.050,-2.595) (8.003,-2.625) (7.950,-2.657) (7.900,-2.686) (7.850,-2.714) (7.800,-2.741) (7.750,-2.767) (7.700,-2.791) (7.650,-2.814) (7.600,-2.837) (7.550,-2.858) (7.500,-2.878) (7.450,-2.898) (7.400,-2.916) (7.350,-2.933) (7.300,-2.950) (7.250,-2.966) (7.200,-2.980) (7.150,-2.994) (7.100,-3.007) (7.050,-3.019) (7.000,-3.031) (6.950,-3.041) (6.900,-3.051) (6.850,-3.060) (6.800,-3.069) (6.750,-3.076) (6.700,-3.083) (6.650,-3.089) (6.600,-3.094) (6.550,-3.099) (6.500,-3.103) (6.450,-3.106) (6.400,-3.109) (6.350,-3.111) (6.300,-3.112) (6.250,-3.113) (6.200,-3.113) (6.150,-3.112) (6.100,-3.111) (6.050,-3.109) (6.000,-3.107) (5.950,-3.103) (5.900,-3.100) (5.850,-3.095) (5.800,-3.090) (5.750,-3.085) (5.700,-3.079) (5.650,-3.072) (5.600,-3.065) (5.550,-3.057) (5.500,-3.049) (5.450,-3.040) (5.400,-3.030) (5.350,-3.020) (5.300,-3.010) (5.250,-2.998) (5.200,-2.987) (5.151,-2.975) (5.100,-2.962) (5.050,-2.949) (5.000,-2.935) (4.950,-2.921) (4.900,-2.906) (4.850,-2.891) (4.800,-2.875) (4.750,-2.859) (4.700,-2.843) (4.650,-2.825) (4.600,-2.808) (4.550,-2.790) (4.500,-2.771) (4.450,-2.752) (4.400,-2.733) (4.350,-2.713) (4.300,-2.693) (4.250,-2.672) (4.200,-2.651) (4.150,-2.629) (4.100,-2.608) (4.050,-2.585) (4.000,-2.562) (3.950,-2.539) (3.900,-2.516) (3.850,-2.492) (3.800,-2.467) (3.750,-2.443) (3.700,-2.418) (3.650,-2.392) (3.600,-2.366) (3.550,-2.340) (3.500,-2.314) (3.450,-2.287) (3.400,-2.260) (3.350,-2.233) (3.300,-2.205) (3.250,-2.177) (3.203,-2.150) (3.159,-2.125) (3.116,-2.100) (3.050,-2.062) (3.000,-2.032) (2.950,-2.002) (2.904,-1.975) (2.850,-1.942) (2.800,-1.912) (2.750,-1.881) (2.700,-1.850) (2.650,-1.819) (2.600,-1.787) (2.550,-1.756) (2.502,-1.725) (2.450,-1.692) (2.400,-1.660) (2.350,-1.627) (2.300,-1.595) (2.250,-1.562) (2.200,-1.529) (2.156,-1.500) (2.100,-1.463) (2.050,-1.430) (2.006,-1.400) (1.950,-1.363) (1.900,-1.329) (1.858,-1.300) (1.800,-1.261) (1.750,-1.227) (1.700,-1.193) (1.650,-1.158) (1.602,-1.125) (1.550,-1.090) (1.500,-1.055) (1.457,-1.025) (1.400,-0.986) (1.350,-0.951) (1.300,-0.916) (1.250,-0.881) (1.206,-0.850) (1.150,-0.811) (1.100,-0.776) (1.050,-0.741) (1.000,-0.706) (0.956,-0.675) (0.900,-0.636) (0.850,-0.600) (0.800,-0.566) (0.750,-0.530) (0.708,-0.500) (0.650,-0.460) (0.601,-0.425) (0.550,-0.389) (0.500,-0.354) (0.450,-0.319) (0.400,-0.283) (0.354,-0.250) (0.300,-0.213) (0.250,-0.177) (0.206,-0.144)}
\EQ{(10.000,-7.071) (9.950,-7.036) (9.900,-7.000) (9.850,-6.965) (9.800,-6.930) (9.758,-6.900) (9.700,-6.859) (9.652,-6.825) (9.600,-6.788) (9.550,-6.753) (9.500,-6.718) (9.450,-6.682) (9.405,-6.650) (9.350,-6.611) (9.300,-6.576) (9.250,-6.541) (9.200,-6.505) (9.157,-6.475) (9.100,-6.435) (9.051,-6.400) (9.000,-6.364) (8.950,-6.329) (8.900,-6.293) (8.850,-6.258) (8.804,-6.225) (8.750,-6.187) (8.700,-6.152) (8.650,-6.117) (8.600,-6.081) (8.556,-6.050) (8.500,-6.010) (8.450,-5.975) (8.400,-5.940) (8.350,-5.904) (8.309,-5.875) (8.250,-5.834) (8.202,-5.800) (8.150,-5.763) (8.100,-5.728) (8.050,-5.692) (8.000,-5.657) (7.955,-5.625) (7.900,-5.586) (7.850,-5.551) (7.800,-5.515) (7.750,-5.480) (7.708,-5.450) (7.650,-5.409) (7.601,-5.375) (7.550,-5.339) (7.500,-5.303) (7.450,-5.268) (7.400,-5.233) (7.354,-5.200) (7.300,-5.162) (7.250,-5.127) (7.200,-5.091) (7.150,-5.056) (7.106,-5.025) (7.050,-4.985) (7.000,-4.950) (6.950,-4.914) (6.900,-4.879) (6.859,-4.850) (6.800,-4.808) (6.753,-4.775) (6.700,-4.738) (6.650,-4.702) (6.600,-4.667) (6.550,-4.632) (6.505,-4.600) (6.450,-4.561) (6.400,-4.525) (6.350,-4.490) (6.300,-4.455) (6.258,-4.425) (6.200,-4.384) (6.152,-4.350) (6.100,-4.313) (6.050,-4.278) (6.000,-4.243) (5.950,-4.207) (5.904,-4.175) (5.850,-4.137) (5.800,-4.101) (5.750,-4.066) (5.700,-4.031) (5.657,-4.000) (5.600,-3.960) (5.551,-3.925) (5.500,-3.889) (5.450,-3.854) (5.400,-3.818) (5.350,-3.783) (5.303,-3.750) (5.250,-3.712) (5.200,-3.677) (5.150,-3.642) (5.100,-3.606) (5.056,-3.575) (5.000,-3.536) (4.950,-3.500) (4.900,-3.465) (4.850,-3.429) (4.808,-3.400) (4.750,-3.359) (4.702,-3.325) (4.650,-3.288) (4.600,-3.253) (4.550,-3.217) (4.500,-3.182) (4.455,-3.150) (4.400,-3.111) (4.350,-3.076) (4.300,-3.041) (4.250,-3.005) (4.207,-2.975) (4.150,-2.935) (4.101,-2.900) (4.050,-2.864) (4.000,-2.828) (3.950,-2.793) (3.900,-2.758) (3.854,-2.725) (3.800,-2.687) (3.750,-2.652) (3.700,-2.616) (3.650,-2.581) (3.606,-2.550) (3.550,-2.510) (3.500,-2.475) (3.450,-2.440) (3.400,-2.404) (3.359,-2.375) (3.300,-2.334) (3.253,-2.300) (3.200,-2.263) (3.150,-2.227) (3.100,-2.192) (3.050,-2.157) (3.005,-2.125) (2.950,-2.086) (2.900,-2.051) (2.850,-2.015) (2.800,-1.980) (2.758,-1.950) (2.700,-1.909) (2.652,-1.875) (2.600,-1.839) (2.550,-1.803) (2.500,-1.768) (2.450,-1.732) (2.404,-1.700) (2.350,-1.662) (2.300,-1.626) (2.250,-1.591) (2.200,-1.556) (2.157,-1.525) (2.100,-1.485) (2.051,-1.450) (2.000,-1.414) (1.950,-1.379) (1.900,-1.344) (1.850,-1.308) (1.803,-1.275) (1.750,-1.238) (1.700,-1.202) (1.650,-1.167) (1.600,-1.131) (1.556,-1.100) (1.500,-1.061) (1.450,-1.025) (1.400,-0.990) (1.350,-0.955) (1.308,-0.925) (1.250,-0.884) (1.202,-0.850) (1.150,-0.813) (1.100,-0.778) (1.050,-0.743) (1.000,-0.707) (0.955,-0.675) (0.900,-0.637) (0.850,-0.601) (0.800,-0.566) (0.750,-0.531) (0.708,-0.500) (0.650,-0.460) (0.601,-0.425) (0.550,-0.389) (0.500,-0.354) (0.450,-0.319) (0.400,-0.283) (0.354,-0.250) (0.300,-0.213) (0.250,-0.177) (0.206,-0.144)}
\EQ{(-0.220,-0.155) (-0.275,-0.195) (-0.319,-0.225) (-0.375,-0.266) (-0.424,-0.300) (-0.475,-0.336) (-0.525,-0.371) (-0.566,-0.400) (-0.625,-0.442) (-0.672,-0.475) (-0.725,-0.513) (-0.775,-0.548) (-0.825,-0.584) (-0.875,-0.619) (-0.920,-0.650) (-0.975,-0.690) (-1.025,-0.725) (-1.075,-0.760) (-1.125,-0.796) (-1.167,-0.825) (-1.225,-0.866) (-1.273,-0.900) (-1.325,-0.937) (-1.375,-0.972) (-1.425,-1.008) (-1.475,-1.043) (-1.520,-1.075) (-1.575,-1.114) (-1.625,-1.149) (-1.675,-1.185) (-1.725,-1.220) (-1.768,-1.250) (-1.825,-1.291) (-1.874,-1.325) (-1.925,-1.361) (-1.975,-1.397) (-2.025,-1.432) (-2.075,-1.467) (-2.121,-1.500) (-2.175,-1.538) (-2.225,-1.573) (-2.275,-1.609) (-2.325,-1.644) (-2.369,-1.675) (-2.425,-1.715) (-2.475,-1.750) (-2.525,-1.786) (-2.575,-1.821) (-2.616,-1.850) (-2.675,-1.892) (-2.722,-1.925) (-2.775,-1.962) (-2.825,-1.998) (-2.875,-2.033) (-2.925,-2.068) (-2.970,-2.100) (-3.025,-2.139) (-3.075,-2.174) (-3.125,-2.210) (-3.175,-2.245) (-3.217,-2.275) (-3.275,-2.316) (-3.323,-2.350) (-3.375,-2.387) (-3.425,-2.422) (-3.475,-2.457) (-3.525,-2.493) (-3.571,-2.525) (-3.625,-2.563) (-3.675,-2.599) (-3.725,-2.634) (-3.775,-2.669) (-3.818,-2.700) (-3.875,-2.740) (-3.924,-2.775) (-3.975,-2.811) (-4.025,-2.846) (-4.066,-2.875) (-4.125,-2.917) (-4.172,-2.950) (-4.225,-2.988) (-4.275,-3.023) (-4.325,-3.058) (-4.375,-3.094) (-4.419,-3.125) (-4.475,-3.164) (-4.525,-3.200) (-4.575,-3.235) (-4.625,-3.270) (-4.667,-3.300) (-4.725,-3.341) (-4.773,-3.375) (-4.825,-3.412) (-4.875,-3.447) (-4.925,-3.483) (-4.975,-3.518) (-5.021,-3.550) (-5.075,-3.589) (-5.125,-3.624) (-5.175,-3.659) (-5.225,-3.695) (-5.268,-3.725) (-5.325,-3.765) (-5.374,-3.800) (-5.425,-3.836) (-5.475,-3.871) (-5.525,-3.907) (-5.575,-3.942) (-5.622,-3.975) (-5.675,-4.013) (-5.725,-4.048) (-5.775,-4.084) (-5.825,-4.119) (-5.869,-4.150) (-5.925,-4.190) (-5.975,-4.225) (-6.025,-4.260) (-6.075,-4.296) (-6.117,-4.325) (-6.175,-4.366) (-6.223,-4.400) (-6.275,-4.437) (-6.325,-4.472) (-6.375,-4.508) (-6.425,-4.543) (-6.470,-4.575) (-6.525,-4.614) (-6.575,-4.649) (-6.625,-4.685) (-6.675,-4.720) (-6.718,-4.750) (-6.775,-4.791) (-6.824,-4.825) (-6.875,-4.861) (-6.925,-4.897) (-6.975,-4.932) (-7.025,-4.967) (-7.071,-5.000) (-7.125,-5.038) (-7.175,-5.073) (-7.225,-5.109) (-7.275,-5.144) (-7.319,-5.175) (-7.375,-5.215) (-7.425,-5.250) (-7.475,-5.286) (-7.525,-5.321) (-7.566,-5.350) (-7.625,-5.392) (-7.672,-5.425) (-7.725,-5.462) (-7.775,-5.498) (-7.825,-5.533) (-7.875,-5.568) (-7.920,-5.600) (-7.975,-5.639) (-8.025,-5.675) (-8.075,-5.710) (-8.125,-5.745) (-8.167,-5.775) (-8.225,-5.816) (-8.273,-5.850) (-8.325,-5.887) (-8.375,-5.922) (-8.425,-5.957) (-8.475,-5.993) (-8.521,-6.025) (-8.575,-6.063) (-8.625,-6.099) (-8.675,-6.134) (-8.725,-6.170) (-8.768,-6.200) (-8.825,-6.240) (-8.874,-6.275) (-8.925,-6.311) (-8.975,-6.346) (-9.025,-6.382) (-9.075,-6.417) (-9.122,-6.450) (-9.175,-6.488) (-9.225,-6.523) (-9.275,-6.558) (-9.325,-6.594) (-9.369,-6.625) (-9.425,-6.665) (-9.475,-6.700) (-9.525,-6.735) (-9.575,-6.771) (-9.617,-6.800) (-9.675,-6.841) (-9.723,-6.875) (-9.775,-6.912) (-9.825,-6.947) (-9.875,-6.983) (-9.925,-7.018) (-9.970,-7.050) (-10.000,-7.071)}
\EQ{(0.206,0.144) (0.248,0.175) (0.300,0.213) (0.350,0.248) (0.400,0.283) (0.450,0.319) (0.496,0.350) (0.550,0.389) (0.600,0.424) (0.650,0.460) (0.700,0.495) (0.743,0.525) (0.800,0.566) (0.849,0.600) (0.900,0.637) (0.950,0.672) (1.000,0.707) (1.050,0.743) (1.096,0.775) (1.150,0.813) (1.200,0.849) (1.250,0.884) (1.300,0.919) (1.344,0.950) (1.400,0.990) (1.450,1.025) (1.500,1.061) (1.550,1.096) (1.591,1.125) (1.650,1.167) (1.697,1.200) (1.750,1.238) (1.800,1.273) (1.850,1.308) (1.900,1.344) (1.945,1.375) (2.000,1.414) (2.050,1.450) (2.100,1.485) (2.150,1.520) (2.192,1.550) (2.250,1.591) (2.298,1.625) (2.350,1.662) (2.400,1.697) (2.450,1.732) (2.500,1.768) (2.546,1.800) (2.600,1.839) (2.650,1.874) (2.700,1.909) (2.750,1.945) (2.793,1.975) (2.850,2.015) (2.899,2.050) (2.950,2.086) (3.000,2.121) (3.050,2.157) (3.100,2.192) (3.147,2.225) (3.200,2.263) (3.250,2.298) (3.300,2.334) (3.350,2.369) (3.394,2.400) (3.450,2.440) (3.500,2.475) (3.550,2.510) (3.600,2.546) (3.642,2.575) (3.700,2.616) (3.748,2.650) (3.800,2.687) (3.850,2.722) (3.900,2.758) (3.950,2.793) (3.995,2.825) (4.050,2.864) (4.100,2.899) (4.150,2.935) (4.200,2.970) (4.243,3.000) (4.300,3.041) (4.349,3.075) (4.400,3.111) (4.450,3.147) (4.500,3.182) (4.550,3.217) (4.596,3.250) (4.650,3.288) (4.700,3.323) (4.750,3.359) (4.800,3.394) (4.844,3.425) (4.900,3.465) (4.950,3.500) (5.000,3.536) (5.050,3.571) (5.091,3.600) (5.150,3.642) (5.197,3.675) (5.250,3.712) (5.300,3.748) (5.350,3.783) (5.400,3.818) (5.445,3.850) (5.500,3.889) (5.550,3.924) (5.600,3.960) (5.650,3.995) (5.692,4.025) (5.750,4.066) (5.798,4.100) (5.850,4.137) (5.900,4.172) (5.950,4.207) (6.000,4.243) (6.046,4.275) (6.100,4.313) (6.150,4.349) (6.200,4.384) (6.250,4.419) (6.293,4.450) (6.350,4.490) (6.399,4.525) (6.450,4.561) (6.500,4.596) (6.550,4.632) (6.600,4.667) (6.647,4.700) (6.700,4.738) (6.750,4.773) (6.800,4.808) (6.850,4.844) (6.894,4.875) (6.950,4.914) (7.000,4.950) (7.050,4.985) (7.100,5.020) (7.142,5.050) (7.200,5.091) (7.248,5.125) (7.300,5.162) (7.350,5.197) (7.400,5.233) (7.450,5.268) (7.495,5.300) (7.550,5.339) (7.600,5.374) (7.650,5.409) (7.700,5.445) (7.743,5.475) (7.800,5.515) (7.849,5.550) (7.900,5.586) (7.950,5.622) (8.000,5.657) (8.050,5.692) (8.096,5.725) (8.150,5.763) (8.200,5.798) (8.250,5.834) (8.300,5.869) (8.344,5.900) (8.400,5.940) (8.450,5.975) (8.500,6.010) (8.550,6.046) (8.591,6.075) (8.650,6.117) (8.697,6.150) (8.750,6.187) (8.800,6.223) (8.850,6.258) (8.900,6.293) (8.945,6.325) (9.000,6.364) (9.050,6.399) (9.100,6.435) (9.150,6.470) (9.192,6.500) (9.250,6.541) (9.298,6.575) (9.350,6.611) (9.400,6.647) (9.450,6.682) (9.500,6.718) (9.546,6.750) (9.600,6.788) (9.650,6.824) (9.700,6.859) (9.750,6.894) (9.793,6.925) (9.850,6.965) (9.899,7.000) (9.950,7.036) (10.000,7.071)}
\EQ{(-10.000,7.071) (-9.950,7.036) (-9.900,7.000) (-9.850,6.965) (-9.800,6.930) (-9.758,6.900) (-9.700,6.859) (-9.652,6.825) (-9.600,6.788) (-9.550,6.753) (-9.500,6.718) (-9.450,6.682) (-9.405,6.650) (-9.350,6.611) (-9.300,6.576) (-9.250,6.541) (-9.200,6.505) (-9.157,6.475) (-9.100,6.435) (-9.051,6.400) (-9.000,6.364) (-8.950,6.329) (-8.900,6.293) (-8.850,6.258) (-8.804,6.225) (-8.750,6.187) (-8.700,6.152) (-8.650,6.117) (-8.600,6.081) (-8.556,6.050) (-8.500,6.010) (-8.450,5.975) (-8.400,5.940) (-8.350,5.904) (-8.309,5.875) (-8.250,5.834) (-8.202,5.800) (-8.150,5.763) (-8.100,5.728) (-8.050,5.692) (-8.000,5.657) (-7.955,5.625) (-7.900,5.586) (-7.850,5.551) (-7.800,5.515) (-7.750,5.480) (-7.708,5.450) (-7.650,5.409) (-7.601,5.375) (-7.550,5.339) (-7.500,5.303) (-7.450,5.268) (-7.400,5.233) (-7.354,5.200) (-7.300,5.162) (-7.250,5.127) (-7.200,5.091) (-7.150,5.056) (-7.106,5.025) (-7.050,4.985) (-7.000,4.950) (-6.950,4.914) (-6.900,4.879) (-6.859,4.850) (-6.800,4.808) (-6.753,4.775) (-6.700,4.738) (-6.650,4.702) (-6.600,4.667) (-6.550,4.632) (-6.505,4.600) (-6.450,4.561) (-6.400,4.525) (-6.350,4.490) (-6.300,4.455) (-6.258,4.425) (-6.200,4.384) (-6.152,4.350) (-6.100,4.313) (-6.050,4.278) (-6.000,4.243) (-5.950,4.207) (-5.904,4.175) (-5.850,4.137) (-5.800,4.101) (-5.750,4.066) (-5.700,4.031) (-5.657,4.000) (-5.600,3.960) (-5.551,3.925) (-5.500,3.889) (-5.450,3.854) (-5.400,3.818) (-5.350,3.783) (-5.303,3.750) (-5.250,3.712) (-5.200,3.677) (-5.150,3.642) (-5.100,3.606) (-5.056,3.575) (-5.000,3.536) (-4.950,3.500) (-4.900,3.465) (-4.850,3.429) (-4.808,3.400) (-4.750,3.359) (-4.702,3.325) (-4.650,3.288) (-4.600,3.253) (-4.550,3.217) (-4.500,3.182) (-4.455,3.150) (-4.400,3.111) (-4.350,3.076) (-4.300,3.041) (-4.250,3.005) (-4.207,2.975) (-4.150,2.935) (-4.101,2.900) (-4.050,2.864) (-4.000,2.828) (-3.950,2.793) (-3.900,2.758) (-3.854,2.725) (-3.800,2.687) (-3.750,2.652) (-3.700,2.616) (-3.650,2.581) (-3.606,2.550) (-3.550,2.510) (-3.500,2.475) (-3.450,2.440) (-3.400,2.404) (-3.359,2.375) (-3.300,2.334) (-3.253,2.300) (-3.200,2.263) (-3.150,2.227) (-3.100,2.192) (-3.050,2.157) (-3.005,2.125) (-2.950,2.086) (-2.900,2.051) (-2.850,2.015) (-2.800,1.980) (-2.758,1.950) (-2.700,1.909) (-2.652,1.875) (-2.600,1.839) (-2.550,1.803) (-2.500,1.768) (-2.450,1.732) (-2.404,1.700) (-2.350,1.662) (-2.300,1.626) (-2.250,1.591) (-2.200,1.556) (-2.157,1.525) (-2.100,1.485) (-2.051,1.450) (-2.000,1.414) (-1.950,1.379) (-1.900,1.344) (-1.850,1.308) (-1.803,1.275) (-1.750,1.238) (-1.700,1.202) (-1.650,1.167) (-1.600,1.131) (-1.556,1.100) (-1.500,1.061) (-1.450,1.025) (-1.400,0.990) (-1.350,0.955) (-1.308,0.925) (-1.250,0.884) (-1.202,0.850) (-1.150,0.813) (-1.100,0.778) (-1.050,0.743) (-1.000,0.707) (-0.955,0.675) (-0.900,0.637) (-0.850,0.601) (-0.800,0.566) (-0.750,0.531) (-0.708,0.500) (-0.650,0.460) (-0.601,0.425) (-0.550,0.389) (-0.500,0.354) (-0.450,0.319) (-0.400,0.283) (-0.354,0.250) (-0.300,0.213) (-0.250,0.177) (-0.220,0.155)}
\EQ{(4.306,-10.000) (4.350,-9.950) (4.392,-9.900) (4.425,-9.860) (4.473,-9.800) (4.513,-9.750) (4.550,-9.701) (4.588,-9.650) (4.624,-9.600) (4.660,-9.550) (4.695,-9.500) (4.725,-9.455) (4.761,-9.400) (4.793,-9.350) (4.824,-9.300) (4.855,-9.250) (4.884,-9.200) (4.913,-9.150) (4.941,-9.100) (4.969,-9.050) (4.995,-9.000) (5.021,-8.950) (5.046,-8.900) (5.071,-8.850) (5.095,-8.800) (5.118,-8.750) (5.140,-8.700) (5.162,-8.650) (5.183,-8.600) (5.203,-8.550) (5.223,-8.500) (5.242,-8.450) (5.261,-8.400) (5.279,-8.350) (5.296,-8.300) (5.313,-8.250) (5.329,-8.200) (5.344,-8.150) (5.359,-8.100) (5.373,-8.050) (5.387,-8.000) (5.400,-7.951) (5.413,-7.900) (5.425,-7.850) (5.436,-7.800) (5.447,-7.750) (5.457,-7.700) (5.467,-7.650) (5.476,-7.600) (5.484,-7.550) (5.492,-7.500) (5.500,-7.450) (5.507,-7.400) (5.513,-7.350) (5.519,-7.300) (5.524,-7.250) (5.529,-7.200) (5.533,-7.150) (5.537,-7.100) (5.540,-7.050) (5.542,-7.000) (5.544,-6.950) (5.545,-6.900) (5.546,-6.850) (5.547,-6.800) (5.546,-6.750) (5.546,-6.700) (5.544,-6.650) (5.543,-6.600) (5.540,-6.550) (5.537,-6.500) (5.534,-6.450) (5.530,-6.400) (5.525,-6.350) (5.520,-6.300) (5.515,-6.250) (5.509,-6.200) (5.502,-6.150) (5.495,-6.100) (5.487,-6.050) (5.478,-6.000) (5.469,-5.950) (5.460,-5.900) (5.450,-5.850) (5.439,-5.800) (5.428,-5.750) (5.416,-5.700) (5.404,-5.650) (5.391,-5.600) (5.378,-5.550) (5.363,-5.500) (5.349,-5.450) (5.334,-5.400) (5.318,-5.350) (5.301,-5.300) (5.284,-5.250) (5.267,-5.200) (5.249,-5.150) (5.230,-5.100) (5.210,-5.050) (5.190,-5.000) (5.170,-4.950) (5.150,-4.903) (5.127,-4.850) (5.104,-4.800) (5.081,-4.750) (5.057,-4.700) (5.033,-4.650) (5.008,-4.600) (4.982,-4.550) (4.956,-4.500) (4.929,-4.450) (4.901,-4.400) (4.875,-4.354) (4.844,-4.300) (4.814,-4.250) (4.784,-4.200) (4.753,-4.150) (4.725,-4.107) (4.688,-4.050) (4.655,-4.000) (4.625,-3.955) (4.587,-3.900) (4.551,-3.850) (4.516,-3.800) (4.479,-3.750) (4.441,-3.700) (4.403,-3.650) (4.364,-3.600) (4.325,-3.551) (4.284,-3.500) (4.250,-3.458) (4.201,-3.400) (4.159,-3.350) (4.125,-3.311) (4.075,-3.254) (4.026,-3.200) (3.981,-3.150) (3.935,-3.100) (3.900,-3.063) (3.864,-3.025) (3.825,-2.985) (3.775,-2.934) (3.725,-2.883) (3.675,-2.834) (3.625,-2.785) (3.575,-2.736) (3.537,-2.700) (3.500,-2.665) (3.457,-2.625) (3.403,-2.575) (3.350,-2.526) (3.300,-2.481) (3.250,-2.437) (3.209,-2.400) (3.151,-2.350) (3.100,-2.306) (3.050,-2.263) (3.006,-2.225) (2.950,-2.178) (2.900,-2.137) (2.856,-2.100) (2.800,-2.054) (2.750,-2.014) (2.702,-1.975) (2.650,-1.933) (2.609,-1.900) (2.550,-1.854) (2.500,-1.814) (2.450,-1.775) (2.400,-1.737) (2.353,-1.700) (2.300,-1.660) (2.255,-1.625) (2.200,-1.583) (2.156,-1.550) (2.100,-1.508) (2.056,-1.475) (2.000,-1.433) (1.956,-1.400) (1.900,-1.359) (1.854,-1.325) (1.800,-1.285) (1.752,-1.250) (1.700,-1.212) (1.650,-1.175) (1.600,-1.139) (1.550,-1.103) (1.500,-1.067) (1.450,-1.031) (1.408,-1.000) (1.350,-0.959) (1.304,-0.925) (1.250,-0.887) (1.200,-0.851) (1.150,-0.815) (1.100,-0.780) (1.059,-0.750) (1.000,-0.708) (0.953,-0.675) (0.900,-0.637) (0.850,-0.602) (0.800,-0.566) (0.750,-0.531) (0.707,-0.500) (0.650,-0.460) (0.601,-0.425) (0.550,-0.389) (0.500,-0.354) (0.450,-0.319) (0.400,-0.283) (0.354,-0.250) (0.300,-0.213) (0.250,-0.177) (0.206,-0.144)}
\EQ{(-0.220,-0.155) (-0.275,-0.195) (-0.319,-0.225) (-0.375,-0.266) (-0.424,-0.300) (-0.475,-0.336) (-0.525,-0.372) (-0.566,-0.400) (-0.625,-0.442) (-0.672,-0.475) (-0.725,-0.513) (-0.775,-0.549) (-0.825,-0.584) (-0.875,-0.620) (-0.918,-0.650) (-0.975,-0.691) (-1.024,-0.725) (-1.075,-0.762) (-1.125,-0.797) (-1.175,-0.833) (-1.225,-0.869) (-1.269,-0.900) (-1.325,-0.941) (-1.373,-0.975) (-1.425,-1.013) (-1.475,-1.049) (-1.525,-1.085) (-1.575,-1.121) (-1.625,-1.157) (-1.675,-1.194) (-1.718,-1.225) (-1.775,-1.267) (-1.820,-1.300) (-1.875,-1.340) (-1.922,-1.375) (-1.975,-1.415) (-2.023,-1.450) (-2.075,-1.489) (-2.123,-1.525) (-2.175,-1.564) (-2.222,-1.600) (-2.275,-1.640) (-2.320,-1.675) (-2.375,-1.717) (-2.418,-1.750) (-2.475,-1.795) (-2.525,-1.834) (-2.575,-1.873) (-2.625,-1.913) (-2.671,-1.950) (-2.725,-1.993) (-2.764,-2.025) (-2.825,-2.075) (-2.875,-2.116) (-2.916,-2.150) (-2.975,-2.199) (-3.025,-2.241) (-3.065,-2.275) (-3.123,-2.325) (-3.175,-2.371) (-3.225,-2.415) (-3.265,-2.450) (-3.321,-2.500) (-3.375,-2.549) (-3.425,-2.595) (-3.475,-2.642) (-3.525,-2.689) (-3.563,-2.725) (-3.600,-2.760) (-3.641,-2.800) (-3.692,-2.850) (-3.742,-2.900) (-3.791,-2.950) (-3.840,-3.000) (-3.875,-3.037) (-3.911,-3.075) (-3.950,-3.117) (-4.000,-3.171) (-4.049,-3.225) (-4.093,-3.275) (-4.137,-3.325) (-4.175,-3.369) (-4.222,-3.425) (-4.264,-3.475) (-4.300,-3.519) (-4.345,-3.575) (-4.384,-3.625) (-4.422,-3.675) (-4.460,-3.725) (-4.497,-3.775) (-4.534,-3.825) (-4.569,-3.875) (-4.600,-3.919) (-4.638,-3.975) (-4.672,-4.025) (-4.700,-4.068) (-4.737,-4.125) (-4.768,-4.175) (-4.799,-4.225) (-4.825,-4.269) (-4.858,-4.325) (-4.887,-4.375) (-4.915,-4.425) (-4.942,-4.475) (-4.969,-4.525) (-4.995,-4.575) (-5.020,-4.625) (-5.045,-4.675) (-5.069,-4.725) (-5.093,-4.775) (-5.116,-4.825) (-5.138,-4.875) (-5.159,-4.925) (-5.180,-4.975) (-5.200,-5.024) (-5.220,-5.075) (-5.239,-5.125) (-5.258,-5.175) (-5.275,-5.223) (-5.293,-5.275) (-5.310,-5.325) (-5.325,-5.373) (-5.341,-5.425) (-5.356,-5.475) (-5.371,-5.525) (-5.384,-5.575) (-5.398,-5.625) (-5.410,-5.675) (-5.422,-5.725) (-5.434,-5.775) (-5.445,-5.825) (-5.455,-5.875) (-5.465,-5.925) (-5.474,-5.975) (-5.483,-6.025) (-5.491,-6.075) (-5.498,-6.125) (-5.505,-6.175) (-5.512,-6.225) (-5.518,-6.275) (-5.523,-6.325) (-5.528,-6.375) (-5.532,-6.425) (-5.536,-6.475) (-5.539,-6.525) (-5.542,-6.575) (-5.544,-6.625) (-5.545,-6.675) (-5.546,-6.725) (-5.547,-6.775) (-5.547,-6.825) (-5.546,-6.875) (-5.545,-6.925) (-5.543,-6.975) (-5.541,-7.025) (-5.538,-7.075) (-5.535,-7.125) (-5.531,-7.175) (-5.526,-7.225) (-5.522,-7.275) (-5.516,-7.325) (-5.510,-7.375) (-5.503,-7.425) (-5.496,-7.475) (-5.489,-7.525) (-5.480,-7.575) (-5.471,-7.625) (-5.462,-7.675) (-5.452,-7.725) (-5.442,-7.775) (-5.430,-7.825) (-5.419,-7.875) (-5.407,-7.925) (-5.394,-7.975) (-5.380,-8.025) (-5.366,-8.075) (-5.352,-8.125) (-5.337,-8.175) (-5.321,-8.225) (-5.305,-8.275) (-5.288,-8.325) (-5.270,-8.375) (-5.252,-8.425) (-5.233,-8.475) (-5.213,-8.525) (-5.193,-8.575) (-5.173,-8.625) (-5.151,-8.675) (-5.129,-8.725) (-5.106,-8.775) (-5.083,-8.825) (-5.059,-8.875) (-5.034,-8.925) (-5.008,-8.975) (-4.982,-9.025) (-4.955,-9.075) (-4.927,-9.125) (-4.900,-9.173) (-4.870,-9.225) (-4.840,-9.275) (-4.809,-9.325) (-4.777,-9.375) (-4.750,-9.417) (-4.711,-9.475) (-4.677,-9.525) (-4.642,-9.575) (-4.606,-9.625) (-4.575,-9.668) (-4.532,-9.725) (-4.500,-9.766) (-4.453,-9.825) (-4.413,-9.875) (-4.375,-9.920) (-4.328,-9.975) (-4.306,-10.000)}
\EQ{(0.206,0.144) (0.248,0.175) (0.300,0.213) (0.350,0.248) (0.400,0.283) (0.450,0.319) (0.495,0.350) (0.550,0.389) (0.600,0.424) (0.650,0.460) (0.700,0.495) (0.742,0.525) (0.800,0.566) (0.848,0.600) (0.900,0.637) (0.950,0.673) (1.000,0.708) (1.050,0.744) (1.094,0.775) (1.150,0.815) (1.199,0.850) (1.250,0.887) (1.300,0.923) (1.350,0.959) (1.400,0.995) (1.443,1.025) (1.500,1.067) (1.546,1.100) (1.600,1.139) (1.650,1.175) (1.700,1.212) (1.750,1.248) (1.800,1.285) (1.850,1.322) (1.900,1.359) (1.950,1.396) (2.000,1.433) (2.050,1.470) (2.100,1.508) (2.150,1.546) (2.200,1.583) (2.250,1.621) (2.300,1.660) (2.350,1.698) (2.400,1.737) (2.450,1.775) (2.500,1.814) (2.546,1.850) (2.600,1.893) (2.640,1.925) (2.700,1.973) (2.750,2.014) (2.795,2.050) (2.850,2.095) (2.900,2.137) (2.946,2.175) (3.000,2.220) (3.050,2.263) (3.094,2.300) (3.150,2.349) (3.200,2.393) (3.250,2.437) (3.293,2.475) (3.349,2.525) (3.400,2.572) (3.450,2.618) (3.500,2.665) (3.537,2.700) (3.575,2.736) (3.615,2.775) (3.666,2.825) (3.717,2.875) (3.767,2.925) (3.815,2.975) (3.864,3.025) (3.900,3.063) (3.935,3.100) (3.975,3.144) (4.025,3.198) (4.071,3.250) (4.115,3.300) (4.150,3.340) (4.200,3.399) (4.243,3.450) (4.275,3.489) (4.325,3.550) (4.364,3.600) (4.400,3.646) (4.441,3.700) (4.475,3.745) (4.516,3.800) (4.550,3.848) (4.587,3.900) (4.621,3.950) (4.650,3.992) (4.688,4.050) (4.721,4.100) (4.750,4.146) (4.784,4.200) (4.814,4.250) (4.844,4.300) (4.873,4.350) (4.900,4.398) (4.925,4.443) (4.950,4.489) (4.975,4.536) (5.000,4.585) (5.025,4.634) (5.050,4.685) (5.075,4.737) (5.100,4.791) (5.125,4.846) (5.149,4.900) (5.170,4.950) (5.190,5.000) (5.210,5.050) (5.230,5.100) (5.249,5.150) (5.267,5.200) (5.284,5.250) (5.301,5.300) (5.318,5.350) (5.334,5.400) (5.349,5.450) (5.363,5.500) (5.378,5.550) (5.391,5.600) (5.404,5.650) (5.416,5.700) (5.428,5.750) (5.439,5.800) (5.450,5.850) (5.460,5.900) (5.469,5.950) (5.478,6.000) (5.487,6.050) (5.495,6.100) (5.502,6.150) (5.509,6.200) (5.515,6.250) (5.520,6.300) (5.525,6.350) (5.530,6.400) (5.534,6.450) (5.537,6.500) (5.540,6.550) (5.543,6.600) (5.544,6.650) (5.546,6.700) (5.546,6.750) (5.547,6.800) (5.546,6.850) (5.545,6.900) (5.544,6.950) (5.542,7.000) (5.540,7.050) (5.537,7.100) (5.533,7.150) (5.529,7.200) (5.524,7.250) (5.519,7.300) (5.513,7.350) (5.507,7.400) (5.500,7.450) (5.492,7.500) (5.484,7.550) (5.476,7.600) (5.467,7.650) (5.457,7.700) (5.447,7.750) (5.436,7.800) (5.425,7.850) (5.413,7.900) (5.400,7.950) (5.387,8.000) (5.373,8.050) (5.359,8.100) (5.344,8.150) (5.329,8.200) (5.313,8.250) (5.296,8.300) (5.279,8.350) (5.261,8.400) (5.242,8.450) (5.223,8.500) (5.203,8.550) (5.183,8.600) (5.162,8.650) (5.140,8.700) (5.118,8.750) (5.095,8.800) (5.071,8.850) (5.046,8.900) (5.021,8.950) (4.995,9.000) (4.969,9.050) (4.941,9.100) (4.913,9.150) (4.884,9.200) (4.855,9.250) (4.825,9.299) (4.793,9.350) (4.761,9.400) (4.728,9.450) (4.700,9.492) (4.660,9.550) (4.625,9.599) (4.588,9.650) (4.551,9.700) (4.513,9.750) (4.475,9.798) (4.433,9.850) (4.400,9.890) (4.350,9.949) (4.306,10.000)}
\EQ{(-4.306,10.000) (-4.350,9.950) (-4.392,9.900) (-4.425,9.860) (-4.473,9.800) (-4.513,9.750) (-4.550,9.701) (-4.588,9.650) (-4.624,9.600) (-4.660,9.550) (-4.695,9.500) (-4.725,9.455) (-4.761,9.400) (-4.793,9.350) (-4.824,9.300) (-4.855,9.250) (-4.884,9.200) (-4.913,9.150) (-4.941,9.100) (-4.969,9.050) (-4.995,9.000) (-5.021,8.950) (-5.046,8.900) (-5.071,8.850) (-5.095,8.800) (-5.118,8.750) (-5.140,8.700) (-5.162,8.650) (-5.183,8.600) (-5.203,8.550) (-5.223,8.500) (-5.242,8.450) (-5.261,8.400) (-5.279,8.350) (-5.296,8.300) (-5.313,8.250) (-5.329,8.200) (-5.344,8.150) (-5.359,8.100) (-5.373,8.050) (-5.387,8.000) (-5.400,7.951) (-5.413,7.900) (-5.425,7.850) (-5.436,7.800) (-5.447,7.750) (-5.457,7.700) (-5.467,7.650) (-5.476,7.600) (-5.484,7.550) (-5.492,7.500) (-5.500,7.450) (-5.507,7.400) (-5.513,7.350) (-5.519,7.300) (-5.524,7.250) (-5.529,7.200) (-5.533,7.150) (-5.537,7.100) (-5.540,7.050) (-5.542,7.000) (-5.544,6.950) (-5.545,6.900) (-5.546,6.850) (-5.547,6.800) (-5.546,6.750) (-5.546,6.700) (-5.544,6.650) (-5.543,6.600) (-5.540,6.550) (-5.537,6.500) (-5.534,6.450) (-5.530,6.400) (-5.525,6.350) (-5.520,6.300) (-5.515,6.250) (-5.509,6.200) (-5.502,6.150) (-5.495,6.100) (-5.487,6.050) (-5.478,6.000) (-5.469,5.950) (-5.460,5.900) (-5.450,5.850) (-5.439,5.800) (-5.428,5.750) (-5.416,5.700) (-5.404,5.650) (-5.391,5.600) (-5.378,5.550) (-5.363,5.500) (-5.349,5.450) (-5.334,5.400) (-5.318,5.350) (-5.301,5.300) (-5.284,5.250) (-5.267,5.200) (-5.249,5.150) (-5.230,5.100) (-5.210,5.050) (-5.190,5.000) (-5.170,4.950) (-5.150,4.903) (-5.127,4.850) (-5.104,4.800) (-5.081,4.750) (-5.057,4.700) (-5.033,4.650) (-5.008,4.600) (-4.982,4.550) (-4.956,4.500) (-4.929,4.450) (-4.901,4.400) (-4.875,4.354) (-4.844,4.300) (-4.814,4.250) (-4.784,4.200) (-4.753,4.150) (-4.725,4.107) (-4.688,4.050) (-4.655,4.000) (-4.625,3.955) (-4.587,3.900) (-4.551,3.850) (-4.516,3.800) (-4.479,3.750) (-4.441,3.700) (-4.403,3.650) (-4.364,3.600) (-4.325,3.551) (-4.284,3.500) (-4.250,3.458) (-4.201,3.400) (-4.159,3.350) (-4.125,3.311) (-4.075,3.254) (-4.026,3.200) (-3.981,3.150) (-3.935,3.100) (-3.900,3.063) (-3.864,3.025) (-3.825,2.985) (-3.775,2.934) (-3.725,2.883) (-3.675,2.834) (-3.625,2.785) (-3.575,2.736) (-3.537,2.700) (-3.500,2.665) (-3.457,2.625) (-3.403,2.575) (-3.350,2.526) (-3.300,2.481) (-3.250,2.437) (-3.209,2.400) (-3.151,2.350) (-3.100,2.306) (-3.050,2.263) (-3.006,2.225) (-2.950,2.178) (-2.900,2.137) (-2.856,2.100) (-2.800,2.054) (-2.750,2.014) (-2.702,1.975) (-2.650,1.933) (-2.609,1.900) (-2.550,1.854) (-2.500,1.814) (-2.450,1.775) (-2.400,1.737) (-2.353,1.700) (-2.300,1.660) (-2.255,1.625) (-2.200,1.583) (-2.156,1.550) (-2.100,1.508) (-2.056,1.475) (-2.000,1.433) (-1.956,1.400) (-1.900,1.359) (-1.854,1.325) (-1.800,1.285) (-1.752,1.250) (-1.700,1.212) (-1.650,1.175) (-1.600,1.139) (-1.550,1.103) (-1.500,1.067) (-1.450,1.031) (-1.408,1.000) (-1.350,0.959) (-1.304,0.925) (-1.250,0.887) (-1.200,0.851) (-1.150,0.815) (-1.100,0.780) (-1.059,0.750) (-1.000,0.708) (-0.953,0.675) (-0.900,0.637) (-0.850,0.602) (-0.800,0.566) (-0.750,0.531) (-0.707,0.500) (-0.650,0.460) (-0.601,0.425) (-0.550,0.389) (-0.500,0.354) (-0.450,0.319) (-0.400,0.283) (-0.354,0.250) (-0.300,0.213) (-0.250,0.177) (-0.220,0.155)}
\EQ{(-0.220,-0.155) (-0.275,-0.195) (-0.319,-0.225) (-0.375,-0.266) (-0.424,-0.300) (-0.475,-0.337) (-0.525,-0.372) (-0.575,-0.407) (-0.625,-0.443) (-0.671,-0.475) (-0.725,-0.514) (-0.775,-0.549) (-0.825,-0.585) (-0.875,-0.620) (-0.917,-0.650) (-0.975,-0.692) (-1.022,-0.725) (-1.075,-0.763) (-1.125,-0.799) (-1.175,-0.835) (-1.225,-0.872) (-1.275,-0.908) (-1.325,-0.944) (-1.368,-0.975) (-1.425,-1.017) (-1.470,-1.050) (-1.525,-1.091) (-1.571,-1.125) (-1.625,-1.166) (-1.671,-1.200) (-1.725,-1.241) (-1.770,-1.275) (-1.825,-1.317) (-1.869,-1.350) (-1.925,-1.394) (-1.965,-1.425) (-2.025,-1.472) (-2.075,-1.512) (-2.123,-1.550) (-2.175,-1.592) (-2.216,-1.625) (-2.275,-1.674) (-2.325,-1.715) (-2.367,-1.750) (-2.425,-1.799) (-2.475,-1.843) (-2.525,-1.886) (-2.570,-1.925) (-2.625,-1.975) (-2.675,-2.020) (-2.725,-2.066) (-2.762,-2.100) (-2.800,-2.136) (-2.841,-2.175) (-2.893,-2.225) (-2.943,-2.275) (-2.993,-2.325) (-3.042,-2.375) (-3.090,-2.425) (-3.125,-2.463) (-3.160,-2.500) (-3.200,-2.545) (-3.249,-2.600) (-3.293,-2.650) (-3.336,-2.700) (-3.375,-2.747) (-3.418,-2.800) (-3.450,-2.840) (-3.497,-2.900) (-3.535,-2.950) (-3.572,-3.000) (-3.608,-3.050) (-3.644,-3.100) (-3.675,-3.146) (-3.712,-3.200) (-3.744,-3.250) (-3.775,-3.299) (-3.807,-3.350) (-3.837,-3.400) (-3.866,-3.450) (-3.894,-3.500) (-3.922,-3.550) (-3.948,-3.600) (-3.974,-3.650) (-3.999,-3.700) (-4.023,-3.750) (-4.047,-3.800) (-4.069,-3.850) (-4.091,-3.900) (-4.112,-3.950) (-4.132,-4.000) (-4.152,-4.050) (-4.170,-4.100) (-4.188,-4.150) (-4.206,-4.200) (-4.222,-4.250) (-4.238,-4.300) (-4.253,-4.350) (-4.267,-4.400) (-4.280,-4.450) (-4.293,-4.500) (-4.305,-4.550) (-4.316,-4.600) (-4.327,-4.650) (-4.337,-4.700) (-4.346,-4.750) (-4.354,-4.800) (-4.362,-4.850) (-4.369,-4.900) (-4.376,-4.950) (-4.381,-5.000) (-4.386,-5.050) (-4.391,-5.100) (-4.394,-5.150) (-4.397,-5.200) (-4.400,-5.250) (-4.401,-5.300) (-4.402,-5.350) (-4.402,-5.400) (-4.402,-5.450) (-4.401,-5.500) (-4.399,-5.550) (-4.396,-5.600) (-4.393,-5.650) (-4.389,-5.700) (-4.385,-5.750) (-4.380,-5.800) (-4.374,-5.850) (-4.367,-5.900) (-4.360,-5.950) (-4.352,-6.000) (-4.343,-6.050) (-4.334,-6.100) (-4.323,-6.150) (-4.313,-6.200) (-4.301,-6.250) (-4.289,-6.300) (-4.276,-6.350) (-4.262,-6.400) (-4.247,-6.450) (-4.232,-6.500) (-4.216,-6.550) (-4.200,-6.598) (-4.182,-6.650) (-4.163,-6.700) (-4.144,-6.750) (-4.125,-6.798) (-4.104,-6.850) (-4.082,-6.900) (-4.060,-6.950) (-4.036,-7.000) (-4.012,-7.050) (-3.987,-7.100) (-3.961,-7.150) (-3.934,-7.200) (-3.907,-7.250) (-3.878,-7.300) (-3.850,-7.347) (-3.818,-7.400) (-3.786,-7.450) (-3.753,-7.500) (-3.725,-7.542) (-3.684,-7.600) (-3.650,-7.648) (-3.611,-7.700) (-3.575,-7.747) (-3.533,-7.800) (-3.500,-7.841) (-3.450,-7.900) (-3.406,-7.950) (-3.361,-8.000) (-3.325,-8.039) (-3.275,-8.092) (-3.225,-8.142) (-3.175,-8.191) (-3.125,-8.238) (-3.085,-8.275) (-3.029,-8.325) (-2.975,-8.372) (-2.925,-8.414) (-2.880,-8.450) (-2.825,-8.494) (-2.784,-8.525) (-2.725,-8.569) (-2.682,-8.600) (-2.625,-8.640) (-2.575,-8.674) (-2.525,-8.707) (-2.475,-8.739) (-2.425,-8.770) (-2.375,-8.800) (-2.325,-8.829) (-2.275,-8.858) (-2.225,-8.886) (-2.175,-8.912) (-2.125,-8.938) (-2.075,-8.964) (-2.025,-8.988) (-1.975,-9.012) (-1.925,-9.035) (-1.875,-9.057) (-1.825,-9.078) (-1.775,-9.099) (-1.725,-9.119) (-1.675,-9.139) (-1.625,-9.158) (-1.577,-9.175) (-1.525,-9.193) (-1.475,-9.210) (-1.425,-9.226) (-1.375,-9.242) (-1.325,-9.257) (-1.275,-9.272) (-1.225,-9.285) (-1.175,-9.298) (-1.125,-9.311) (-1.075,-9.323) (-1.025,-9.335) (-0.975,-9.345) (-0.925,-9.356) (-0.875,-9.365) (-0.825,-9.375) (-0.775,-9.383) (-0.725,-9.391) (-0.675,-9.399) (-0.625,-9.406) (-0.575,-9.412) (-0.525,-9.418) (-0.475,-9.423) (-0.425,-9.428) (-0.375,-9.432) (-0.325,-9.436) (-0.275,-9.439) (-0.225,-9.442) (-0.175,-9.444) (-0.125,-9.446) (-0.075,-9.447) (-0.025,-9.447) (0.050,-9.447) (0.100,-9.446) (0.150,-9.445) (0.200,-9.443) (0.250,-9.441) (0.300,-9.438) (0.350,-9.434) (0.400,-9.430) (0.450,-9.426) (0.500,-9.421) (0.550,-9.415) (0.600,-9.409) (0.650,-9.402) (0.700,-9.395) (0.750,-9.387) (0.800,-9.379) (0.850,-9.370) (0.900,-9.361) (0.950,-9.351) (1.000,-9.340) (1.050,-9.329) (1.100,-9.317) (1.150,-9.305) (1.200,-9.292) (1.250,-9.278) (1.300,-9.264) (1.349,-9.250) (1.400,-9.234) (1.450,-9.218) (1.500,-9.202) (1.550,-9.185) (1.600,-9.167) (1.650,-9.148) (1.700,-9.129) (1.750,-9.109) (1.800,-9.089) (1.850,-9.068) (1.900,-9.046) (1.946,-9.025) (2.000,-9.000) (2.050,-8.976) (2.100,-8.951) (2.150,-8.925) (2.198,-8.900) (2.244,-8.875) (2.289,-8.850) (2.333,-8.825) (2.400,-8.785) (2.450,-8.755) (2.497,-8.725) (2.550,-8.691) (2.600,-8.657) (2.647,-8.625) (2.700,-8.587) (2.750,-8.551) (2.800,-8.513) (2.849,-8.475) (2.900,-8.434) (2.942,-8.400) (3.000,-8.351) (3.050,-8.307) (3.100,-8.262) (3.139,-8.225) (3.192,-8.175) (3.242,-8.125) (3.291,-8.075) (3.338,-8.025) (3.375,-7.985) (3.425,-7.929) (3.471,-7.875) (3.513,-7.825) (3.550,-7.779) (3.592,-7.725) (3.625,-7.682) (3.667,-7.625) (3.700,-7.578) (3.736,-7.525) (3.770,-7.475) (3.800,-7.428) (3.833,-7.375) (3.863,-7.325) (3.892,-7.275) (3.921,-7.225) (3.948,-7.175) (3.974,-7.125) (4.000,-7.075) (4.024,-7.025) (4.048,-6.975) (4.071,-6.925) (4.093,-6.875) (4.114,-6.825) (4.134,-6.775) (4.154,-6.725) (4.173,-6.675) (4.191,-6.625) (4.208,-6.575) (4.224,-6.525) (4.240,-6.475) (4.255,-6.425) (4.269,-6.375) (4.282,-6.325) (4.295,-6.275) (4.307,-6.225) (4.318,-6.175) (4.329,-6.125) (4.338,-6.075) (4.347,-6.025) (4.356,-5.975) (4.364,-5.925) (4.370,-5.875) (4.377,-5.825) (4.382,-5.775) (4.387,-5.725) (4.391,-5.675) (4.395,-5.625) (4.398,-5.575) (4.400,-5.525) (4.401,-5.475) (4.402,-5.425) (4.402,-5.375) (4.402,-5.325) (4.401,-5.275) (4.399,-5.225) (4.396,-5.175) (4.393,-5.125) (4.389,-5.075) (4.384,-5.025) (4.379,-4.975) (4.373,-4.925) (4.366,-4.875) (4.358,-4.825) (4.350,-4.775) (4.342,-4.725) (4.332,-4.675) (4.322,-4.625) (4.311,-4.575) (4.299,-4.525) (4.287,-4.475) (4.274,-4.425) (4.260,-4.375) (4.245,-4.325) (4.230,-4.275) (4.214,-4.225) (4.197,-4.175) (4.180,-4.125) (4.161,-4.075) (4.142,-4.025) (4.122,-3.975) (4.102,-3.925) (4.080,-3.875) (4.058,-3.825) (4.035,-3.775) (4.011,-3.725) (3.987,-3.675) (3.961,-3.625) (3.935,-3.575) (3.908,-3.525) (3.880,-3.475) (3.852,-3.425) (3.825,-3.380) (3.792,-3.325) (3.760,-3.275) (3.728,-3.225) (3.700,-3.183) (3.661,-3.125) (3.626,-3.075) (3.590,-3.025) (3.554,-2.975) (3.516,-2.925) (3.478,-2.875) (3.438,-2.825) (3.400,-2.778) (3.357,-2.725) (3.315,-2.675) (3.275,-2.629) (3.227,-2.575) (3.182,-2.525) (3.136,-2.475) (3.100,-2.436) (3.050,-2.384) (3.000,-2.332) (2.950,-2.282) (2.900,-2.232) (2.850,-2.184) (2.800,-2.136) (2.762,-2.100) (2.725,-2.066) (2.681,-2.025) (2.626,-1.975) (2.575,-1.930) (2.525,-1.886) (2.484,-1.850) (2.426,-1.800) (2.375,-1.757) (2.325,-1.715) (2.277,-1.675) (2.225,-1.633) (2.185,-1.600) (2.125,-1.552) (2.075,-1.512) (2.029,-1.475) (1.975,-1.433) (1.933,-1.400) (1.875,-1.355) (1.825,-1.317) (1.775,-1.279) (1.725,-1.241) (1.675,-1.203) (1.625,-1.166) (1.575,-1.128) (1.525,-1.091) (1.475,-1.054) (1.425,-1.017) (1.375,-0.981) (1.333,-0.950) (1.275,-0.908) (1.230,-0.875) (1.175,-0.835) (1.126,-0.800) (1.075,-0.763) (1.025,-0.727) (0.975,-0.692) (0.925,-0.656) (0.882,-0.625) (0.825,-0.585) (0.777,-0.550) (0.725,-0.514) (0.675,-0.478) (0.625,-0.443) (0.575,-0.407) (0.531,-0.375) (0.475,-0.337) (0.425,-0.301) (0.375,-0.266) (0.325,-0.230) (0.284,-0.200) (0.225,-0.160) (0.206,-0.144)}
\EQ{(0.206,0.144) (0.248,0.175) (0.300,0.213) (0.350,0.248) (0.400,0.283) (0.450,0.319) (0.495,0.350) (0.550,0.390) (0.600,0.425) (0.650,0.460) (0.700,0.496) (0.742,0.525) (0.800,0.567) (0.847,0.600) (0.900,0.638) (0.950,0.674) (1.000,0.710) (1.050,0.745) (1.092,0.775) (1.150,0.817) (1.196,0.850) (1.250,0.890) (1.299,0.925) (1.350,0.963) (1.400,0.999) (1.450,1.036) (1.500,1.073) (1.550,1.110) (1.600,1.147) (1.650,1.184) (1.700,1.222) (1.750,1.260) (1.800,1.298) (1.850,1.336) (1.900,1.374) (1.950,1.413) (1.997,1.450) (2.050,1.492) (2.092,1.525) (2.150,1.572) (2.200,1.612) (2.246,1.650) (2.300,1.694) (2.350,1.736) (2.396,1.775) (2.450,1.821) (2.500,1.864) (2.541,1.900) (2.598,1.950) (2.650,1.997) (2.700,2.043) (2.750,2.089) (2.789,2.125) (2.825,2.160) (2.867,2.200) (2.918,2.250) (2.968,2.300) (3.017,2.350) (3.066,2.400) (3.113,2.450) (3.150,2.490) (3.200,2.545) (3.249,2.600) (3.293,2.650) (3.336,2.700) (3.375,2.747) (3.418,2.800) (3.450,2.840) (3.497,2.900) (3.535,2.950) (3.572,3.000) (3.608,3.050) (3.644,3.100) (3.675,3.146) (3.712,3.200) (3.744,3.250) (3.775,3.299) (3.807,3.350) (3.837,3.400) (3.866,3.450) (3.894,3.500) (3.922,3.550) (3.948,3.600) (3.974,3.650) (3.999,3.700) (4.023,3.750) (4.047,3.800) (4.069,3.850) (4.091,3.900) (4.112,3.950) (4.132,4.000) (4.152,4.050) (4.170,4.100) (4.188,4.150) (4.206,4.200) (4.222,4.250) (4.238,4.300) (4.253,4.350) (4.267,4.400) (4.280,4.450) (4.293,4.500) (4.305,4.550) (4.316,4.600) (4.327,4.650) (4.337,4.700) (4.346,4.750) (4.354,4.800) (4.362,4.850) (4.369,4.900) (4.376,4.950) (4.381,5.000) (4.386,5.050) (4.391,5.100) (4.394,5.150) (4.397,5.200) (4.400,5.250) (4.401,5.300) (4.402,5.350) (4.402,5.400) (4.402,5.450) (4.401,5.500) (4.399,5.550) (4.396,5.600) (4.393,5.650) (4.389,5.700) (4.385,5.750) (4.380,5.800) (4.374,5.850) (4.367,5.900) (4.360,5.950) (4.352,6.000) (4.343,6.050) (4.334,6.100) (4.323,6.150) (4.313,6.200) (4.301,6.250) (4.289,6.300) (4.276,6.350) (4.262,6.400) (4.247,6.450) (4.232,6.500) (4.216,6.550) (4.200,6.598) (4.182,6.650) (4.163,6.700) (4.144,6.750) (4.125,6.798) (4.104,6.850) (4.082,6.900) (4.060,6.950) (4.036,7.000) (4.012,7.050) (3.987,7.100) (3.961,7.150) (3.934,7.200) (3.907,7.250) (3.878,7.300) (3.850,7.347) (3.818,7.400) (3.786,7.450) (3.753,7.500) (3.725,7.542) (3.684,7.600) (3.650,7.648) (3.611,7.700) (3.575,7.747) (3.533,7.800) (3.500,7.841) (3.450,7.900) (3.406,7.950) (3.361,8.000) (3.325,8.039) (3.275,8.092) (3.225,8.142) (3.175,8.191) (3.125,8.238) (3.085,8.275) (3.029,8.325) (2.975,8.372) (2.925,8.414) (2.880,8.450) (2.825,8.494) (2.784,8.525) (2.725,8.569) (2.682,8.600) (2.625,8.640) (2.575,8.674) (2.525,8.707) (2.475,8.739) (2.425,8.770) (2.375,8.800) (2.325,8.829) (2.275,8.858) (2.225,8.886) (2.175,8.912) (2.125,8.938) (2.075,8.964) (2.025,8.988) (1.975,9.012) (1.925,9.035) (1.875,9.057) (1.825,9.078) (1.775,9.099) (1.725,9.119) (1.675,9.139) (1.625,9.158) (1.577,9.175) (1.525,9.193) (1.475,9.210) (1.425,9.226) (1.375,9.242) (1.325,9.257) (1.275,9.272) (1.225,9.285) (1.175,9.298) (1.125,9.311) (1.075,9.323) (1.025,9.335) (0.975,9.345) (0.925,9.356) (0.875,9.365) (0.825,9.375) (0.775,9.383) (0.725,9.391) (0.675,9.399) (0.625,9.406) (0.575,9.412) (0.525,9.418) (0.475,9.423) (0.425,9.428) (0.375,9.432) (0.325,9.436) (0.275,9.439) (0.225,9.442) (0.175,9.444) (0.125,9.446) (0.075,9.447) (0.025,9.447) (-0.050,9.447) (-0.100,9.446) (-0.150,9.445) (-0.200,9.443) (-0.250,9.441) (-0.300,9.438) (-0.350,9.434) (-0.400,9.430) (-0.450,9.426) (-0.500,9.421) (-0.550,9.415) (-0.600,9.409) (-0.650,9.402) (-0.700,9.395) (-0.750,9.387) (-0.800,9.379) (-0.850,9.370) (-0.900,9.361) (-0.950,9.351) (-1.000,9.340) (-1.050,9.329) (-1.100,9.317) (-1.150,9.305) (-1.200,9.292) (-1.250,9.278) (-1.300,9.264) (-1.349,9.250) (-1.400,9.234) (-1.450,9.218) (-1.500,9.202) (-1.550,9.185) (-1.600,9.167) (-1.650,9.148) (-1.700,9.129) (-1.750,9.109) (-1.800,9.089) (-1.850,9.068) (-1.900,9.046) (-1.946,9.025) (-2.000,9.000) (-2.050,8.976) (-2.100,8.951) (-2.150,8.925) (-2.198,8.900) (-2.244,8.875) (-2.289,8.850) (-2.333,8.825) (-2.400,8.785) (-2.450,8.755) (-2.497,8.725) (-2.550,8.691) (-2.600,8.657) (-2.647,8.625) (-2.700,8.587) (-2.750,8.551) (-2.800,8.513) (-2.849,8.475) (-2.900,8.434) (-2.942,8.400) (-3.000,8.351) (-3.050,8.307) (-3.100,8.262) (-3.139,8.225) (-3.192,8.175) (-3.242,8.125) (-3.291,8.075) (-3.338,8.025) (-3.375,7.985) (-3.425,7.929) (-3.471,7.875) (-3.513,7.825) (-3.550,7.779) (-3.592,7.725) (-3.625,7.682) (-3.667,7.625) (-3.700,7.578) (-3.736,7.525) (-3.770,7.475) (-3.800,7.428) (-3.833,7.375) (-3.863,7.325) (-3.892,7.275) (-3.921,7.225) (-3.948,7.175) (-3.974,7.125) (-4.000,7.075) (-4.024,7.025) (-4.048,6.975) (-4.071,6.925) (-4.093,6.875) (-4.114,6.825) (-4.134,6.775) (-4.154,6.725) (-4.173,6.675) (-4.191,6.625) (-4.208,6.575) (-4.224,6.525) (-4.240,6.475) (-4.255,6.425) (-4.269,6.375) (-4.282,6.325) (-4.295,6.275) (-4.307,6.225) (-4.318,6.175) (-4.329,6.125) (-4.338,6.075) (-4.347,6.025) (-4.356,5.975) (-4.364,5.925) (-4.370,5.875) (-4.377,5.825) (-4.382,5.775) (-4.387,5.725) (-4.391,5.675) (-4.395,5.625) (-4.398,5.575) (-4.400,5.525) (-4.401,5.475) (-4.402,5.425) (-4.402,5.375) (-4.402,5.325) (-4.401,5.275) (-4.399,5.225) (-4.396,5.175) (-4.393,5.125) (-4.389,5.075) (-4.384,5.025) (-4.379,4.975) (-4.373,4.925) (-4.366,4.875) (-4.358,4.825) (-4.350,4.775) (-4.342,4.725) (-4.332,4.675) (-4.322,4.625) (-4.311,4.575) (-4.299,4.525) (-4.287,4.475) (-4.274,4.425) (-4.260,4.375) (-4.245,4.325) (-4.230,4.275) (-4.214,4.225) (-4.197,4.175) (-4.180,4.125) (-4.161,4.075) (-4.142,4.025) (-4.122,3.975) (-4.102,3.925) (-4.080,3.875) (-4.058,3.825) (-4.035,3.775) (-4.011,3.725) (-3.987,3.675) (-3.961,3.625) (-3.935,3.575) (-3.908,3.525) (-3.880,3.475) (-3.852,3.425) (-3.825,3.380) (-3.792,3.325) (-3.760,3.275) (-3.728,3.225) (-3.700,3.183) (-3.661,3.125) (-3.626,3.075) (-3.590,3.025) (-3.554,2.975) (-3.516,2.925) (-3.478,2.875) (-3.438,2.825) (-3.400,2.778) (-3.357,2.725) (-3.315,2.675) (-3.275,2.629) (-3.227,2.575) (-3.182,2.525) (-3.136,2.475) (-3.100,2.436) (-3.050,2.384) (-3.000,2.332) (-2.950,2.282) (-2.900,2.232) (-2.850,2.184) (-2.800,2.136) (-2.762,2.100) (-2.725,2.066) (-2.681,2.025) (-2.626,1.975) (-2.575,1.930) (-2.525,1.886) (-2.484,1.850) (-2.426,1.800) (-2.375,1.757) (-2.325,1.715) (-2.277,1.675) (-2.225,1.633) (-2.185,1.600) (-2.125,1.552) (-2.075,1.512) (-2.029,1.475) (-1.975,1.433) (-1.933,1.400) (-1.875,1.355) (-1.825,1.317) (-1.775,1.279) (-1.725,1.241) (-1.675,1.203) (-1.625,1.166) (-1.575,1.128) (-1.525,1.091) (-1.475,1.054) (-1.425,1.017) (-1.375,0.981) (-1.333,0.950) (-1.275,0.908) (-1.230,0.875) (-1.175,0.835) (-1.126,0.800) (-1.075,0.763) (-1.025,0.727) (-0.975,0.692) (-0.925,0.656) (-0.882,0.625) (-0.825,0.585) (-0.777,0.550) (-0.725,0.514) (-0.675,0.478) (-0.625,0.443) (-0.575,0.407) (-0.531,0.375) (-0.475,0.337) (-0.425,0.301) (-0.375,0.266) (-0.325,0.230) (-0.284,0.200) (-0.225,0.160) (-0.220,0.155)}
\EQ{(-0.220,-0.155) (-0.275,-0.195) (-0.319,-0.225) (-0.375,-0.266) (-0.424,-0.300) (-0.475,-0.337) (-0.525,-0.372) (-0.575,-0.407) (-0.625,-0.443) (-0.671,-0.475) (-0.725,-0.514) (-0.775,-0.549) (-0.825,-0.585) (-0.875,-0.621) (-0.916,-0.650) (-0.975,-0.693) (-1.020,-0.725) (-1.075,-0.765) (-1.124,-0.800) (-1.175,-0.838) (-1.225,-0.874) (-1.275,-0.911) (-1.325,-0.948) (-1.375,-0.985) (-1.425,-1.022) (-1.475,-1.060) (-1.525,-1.098) (-1.575,-1.136) (-1.625,-1.174) (-1.675,-1.213) (-1.723,-1.250) (-1.775,-1.291) (-1.818,-1.325) (-1.875,-1.370) (-1.925,-1.411) (-1.973,-1.450) (-2.025,-1.493) (-2.075,-1.535) (-2.122,-1.575) (-2.175,-1.621) (-2.225,-1.664) (-2.265,-1.700) (-2.321,-1.750) (-2.375,-1.800) (-2.425,-1.846) (-2.475,-1.894) (-2.525,-1.942) (-2.575,-1.991) (-2.625,-2.042) (-2.675,-2.093) (-2.725,-2.146) (-2.775,-2.200) (-2.820,-2.250) (-2.864,-2.300) (-2.900,-2.342) (-2.949,-2.400) (-2.989,-2.450) (-3.025,-2.495) (-3.068,-2.550) (-3.100,-2.593) (-3.142,-2.650) (-3.175,-2.697) (-3.212,-2.750) (-3.245,-2.800) (-3.275,-2.846) (-3.309,-2.900) (-3.340,-2.950) (-3.369,-3.000) (-3.398,-3.050) (-3.425,-3.099) (-3.450,-3.146) (-3.475,-3.195) (-3.500,-3.245) (-3.525,-3.298) (-3.549,-3.350) (-3.571,-3.400) (-3.592,-3.450) (-3.612,-3.500) (-3.631,-3.550) (-3.650,-3.600) (-3.667,-3.650) (-3.684,-3.700) (-3.700,-3.750) (-3.715,-3.800) (-3.729,-3.850) (-3.742,-3.900) (-3.755,-3.950) (-3.767,-4.000) (-3.777,-4.050) (-3.787,-4.100) (-3.797,-4.150) (-3.805,-4.200) (-3.813,-4.250) (-3.820,-4.300) (-3.826,-4.350) (-3.831,-4.400) (-3.835,-4.450) (-3.839,-4.500) (-3.842,-4.550) (-3.844,-4.600) (-3.845,-4.650) (-3.846,-4.700) (-3.846,-4.750) (-3.845,-4.800) (-3.843,-4.850) (-3.840,-4.900) (-3.837,-4.950) (-3.833,-5.000) (-3.828,-5.050) (-3.822,-5.100) (-3.816,-5.150) (-3.808,-5.200) (-3.800,-5.250) (-3.791,-5.300) (-3.782,-5.350) (-3.771,-5.400) (-3.760,-5.450) (-3.747,-5.500) (-3.734,-5.550) (-3.721,-5.600) (-3.706,-5.650) (-3.690,-5.700) (-3.674,-5.750) (-3.656,-5.800) (-3.638,-5.850) (-3.619,-5.900) (-3.600,-5.948) (-3.578,-6.000) (-3.556,-6.050) (-3.533,-6.100) (-3.509,-6.150) (-3.484,-6.200) (-3.458,-6.250) (-3.431,-6.300) (-3.403,-6.350) (-3.375,-6.399) (-3.344,-6.450) (-3.313,-6.500) (-3.280,-6.550) (-3.250,-6.594) (-3.211,-6.650) (-3.175,-6.699) (-3.137,-6.750) (-3.100,-6.797) (-3.057,-6.850) (-3.015,-6.900) (-2.975,-6.946) (-2.926,-7.000) (-2.879,-7.050) (-2.830,-7.100) (-2.779,-7.150) (-2.727,-7.200) (-2.675,-7.247) (-2.625,-7.291) (-2.585,-7.325) (-2.525,-7.375) (-2.475,-7.414) (-2.429,-7.450) (-2.375,-7.490) (-2.326,-7.525) (-2.275,-7.560) (-2.225,-7.594) (-2.176,-7.625) (-2.125,-7.657) (-2.075,-7.688) (-2.025,-7.717) (-1.975,-7.745) (-1.925,-7.772) (-1.875,-7.799) (-1.825,-7.824) (-1.775,-7.849) (-1.725,-7.873) (-1.675,-7.896) (-1.625,-7.918) (-1.575,-7.939) (-1.525,-7.959) (-1.475,-7.979) (-1.425,-7.998) (-1.375,-8.016) (-1.325,-8.034) (-1.276,-8.050) (-1.225,-8.066) (-1.175,-8.082) (-1.125,-8.096) (-1.075,-8.110) (-1.025,-8.123) (-0.975,-8.136) (-0.925,-8.148) (-0.875,-8.159) (-0.825,-8.170) (-0.775,-8.179) (-0.725,-8.189) (-0.675,-8.197) (-0.625,-8.205) (-0.575,-8.213) (-0.525,-8.219) (-0.475,-8.225) (-0.425,-8.231) (-0.375,-8.236) (-0.325,-8.240) (-0.275,-8.244) (-0.225,-8.247) (-0.175,-8.249) (-0.125,-8.251) (-0.075,-8.252) (-0.025,-8.253) (0.050,-8.253) (0.100,-8.252) (0.150,-8.250) (0.200,-8.248) (0.250,-8.245) (0.300,-8.242) (0.350,-8.238) (0.400,-8.234) (0.450,-8.228) (0.500,-8.223) (0.550,-8.216) (0.600,-8.209) (0.650,-8.201) (0.700,-8.193) (0.750,-8.184) (0.800,-8.175) (0.850,-8.164) (0.900,-8.153) (0.950,-8.142) (1.000,-8.130) (1.050,-8.117) (1.100,-8.103) (1.150,-8.089) (1.200,-8.074) (1.250,-8.058) (1.300,-8.042) (1.350,-8.025) (1.400,-8.007) (1.450,-7.989) (1.500,-7.969) (1.548,-7.950) (1.600,-7.928) (1.650,-7.907) (1.700,-7.884) (1.750,-7.861) (1.800,-7.837) (1.850,-7.812) (1.900,-7.786) (1.950,-7.759) (2.000,-7.731) (2.050,-7.702) (2.096,-7.675) (2.150,-7.642) (2.200,-7.610) (2.250,-7.577) (2.300,-7.543) (2.350,-7.508) (2.395,-7.475) (2.450,-7.434) (2.493,-7.400) (2.550,-7.354) (2.600,-7.313) (2.644,-7.275) (2.700,-7.225) (2.750,-7.178) (2.800,-7.130) (2.850,-7.080) (2.900,-7.028) (2.949,-6.975) (2.993,-6.925) (3.036,-6.875) (3.075,-6.828) (3.117,-6.775) (3.150,-6.733) (3.193,-6.675) (3.225,-6.630) (3.263,-6.575) (3.296,-6.525) (3.325,-6.480) (3.359,-6.425) (3.389,-6.375) (3.418,-6.325) (3.445,-6.275) (3.472,-6.225) (3.497,-6.175) (3.521,-6.125) (3.545,-6.075) (3.567,-6.025) (3.589,-5.975) (3.609,-5.925) (3.629,-5.875) (3.647,-5.825) (3.665,-5.775) (3.682,-5.725) (3.698,-5.675) (3.713,-5.625) (3.728,-5.575) (3.741,-5.525) (3.754,-5.475) (3.765,-5.425) (3.776,-5.375) (3.787,-5.325) (3.796,-5.275) (3.804,-5.225) (3.812,-5.175) (3.819,-5.125) (3.825,-5.075) (3.830,-5.025) (3.835,-4.975) (3.839,-4.925) (3.842,-4.875) (3.844,-4.825) (3.845,-4.775) (3.846,-4.725) (3.846,-4.675) (3.845,-4.625) (3.843,-4.575) (3.841,-4.525) (3.837,-4.475) (3.833,-4.425) (3.828,-4.375) (3.823,-4.325) (3.816,-4.275) (3.809,-4.225) (3.801,-4.175) (3.792,-4.125) (3.783,-4.075) (3.772,-4.025) (3.761,-3.975) (3.749,-3.925) (3.736,-3.875) (3.722,-3.825) (3.708,-3.775) (3.692,-3.725) (3.676,-3.675) (3.659,-3.625) (3.641,-3.575) (3.622,-3.525) (3.602,-3.475) (3.582,-3.425) (3.560,-3.375) (3.538,-3.325) (3.514,-3.275) (3.490,-3.225) (3.465,-3.175) (3.439,-3.125) (3.412,-3.075) (3.384,-3.025) (3.355,-2.975) (3.325,-2.926) (3.294,-2.875) (3.262,-2.825) (3.229,-2.775) (3.200,-2.733) (3.160,-2.675) (3.125,-2.627) (3.087,-2.575) (3.050,-2.527) (3.009,-2.475) (2.975,-2.432) (2.928,-2.375) (2.886,-2.325) (2.850,-2.284) (2.800,-2.228) (2.752,-2.175) (2.705,-2.125) (2.657,-2.075) (2.609,-2.025) (2.559,-1.975) (2.508,-1.925) (2.456,-1.875) (2.403,-1.825) (2.350,-1.777) (2.300,-1.731) (2.250,-1.687) (2.209,-1.650) (2.151,-1.600) (2.100,-1.556) (2.050,-1.514) (2.003,-1.475) (1.950,-1.431) (1.900,-1.391) (1.850,-1.350) (1.800,-1.311) (1.755,-1.275) (1.700,-1.232) (1.659,-1.200) (1.600,-1.155) (1.550,-1.117) (1.500,-1.079) (1.450,-1.041) (1.400,-1.004) (1.350,-0.966) (1.300,-0.929) (1.250,-0.893) (1.200,-0.856) (1.158,-0.825) (1.100,-0.783) (1.055,-0.750) (1.000,-0.711) (0.951,-0.675) (0.900,-0.639) (0.850,-0.603) (0.800,-0.567) (0.750,-0.532) (0.706,-0.500) (0.650,-0.461) (0.600,-0.425) (0.550,-0.390) (0.500,-0.354) (0.450,-0.319) (0.400,-0.284) (0.354,-0.250) (0.300,-0.213) (0.250,-0.177) (0.206,-0.144)}
\EQ{(0.206,0.144) (0.248,0.175) (0.300,0.213) (0.350,0.248) (0.400,0.284) (0.450,0.319) (0.495,0.350) (0.550,0.390) (0.600,0.425) (0.650,0.461) (0.700,0.496) (0.742,0.525) (0.800,0.567) (0.846,0.600) (0.900,0.639) (0.950,0.675) (1.000,0.711) (1.050,0.747) (1.100,0.783) (1.150,0.819) (1.192,0.850) (1.250,0.893) (1.294,0.925) (1.350,0.966) (1.396,1.000) (1.450,1.041) (1.495,1.075) (1.550,1.117) (1.594,1.150) (1.650,1.193) (1.691,1.225) (1.750,1.271) (1.800,1.311) (1.850,1.350) (1.900,1.391) (1.942,1.425) (2.000,1.472) (2.050,1.514) (2.093,1.550) (2.150,1.599) (2.200,1.643) (2.250,1.687) (2.293,1.725) (2.348,1.775) (2.400,1.823) (2.450,1.870) (2.500,1.918) (2.550,1.967) (2.600,2.017) (2.650,2.068) (2.700,2.120) (2.750,2.173) (2.797,2.225) (2.842,2.275) (2.875,2.313) (2.925,2.372) (2.969,2.425) (3.009,2.475) (3.049,2.525) (3.087,2.575) (3.124,2.625) (3.160,2.675) (3.195,2.725) (3.225,2.770) (3.262,2.825) (3.294,2.875) (3.325,2.925) (3.355,2.975) (3.384,3.025) (3.412,3.075) (3.439,3.125) (3.465,3.175) (3.490,3.225) (3.514,3.275) (3.538,3.325) (3.560,3.375) (3.582,3.425) (3.602,3.475) (3.622,3.525) (3.641,3.575) (3.659,3.625) (3.675,3.673) (3.692,3.725) (3.708,3.775) (3.722,3.825) (3.736,3.875) (3.749,3.925) (3.761,3.975) (3.772,4.025) (3.783,4.075) (3.792,4.125) (3.801,4.175) (3.809,4.225) (3.816,4.275) (3.823,4.325) (3.828,4.375) (3.833,4.425) (3.837,4.475) (3.841,4.525) (3.843,4.575) (3.845,4.625) (3.846,4.675) (3.846,4.725) (3.845,4.775) (3.844,4.825) (3.842,4.875) (3.839,4.925) (3.835,4.975) (3.830,5.025) (3.825,5.075) (3.819,5.125) (3.812,5.175) (3.804,5.225) (3.796,5.275) (3.787,5.325) (3.776,5.375) (3.765,5.425) (3.754,5.475) (3.741,5.525) (3.728,5.575) (3.713,5.625) (3.698,5.675) (3.682,5.725) (3.665,5.775) (3.647,5.825) (3.629,5.875) (3.609,5.925) (3.589,5.975) (3.567,6.025) (3.545,6.075) (3.521,6.125) (3.497,6.175) (3.472,6.225) (3.445,6.275) (3.418,6.325) (3.389,6.375) (3.359,6.425) (3.328,6.475) (3.300,6.519) (3.263,6.575) (3.229,6.625) (3.193,6.675) (3.156,6.725) (3.125,6.765) (3.077,6.825) (3.036,6.875) (3.000,6.917) (2.950,6.974) (2.903,7.025) (2.855,7.075) (2.805,7.125) (2.753,7.175) (2.700,7.225) (2.650,7.269) (2.600,7.313) (2.555,7.350) (2.500,7.395) (2.450,7.434) (2.400,7.471) (2.350,7.508) (2.300,7.543) (2.253,7.575) (2.200,7.610) (2.150,7.642) (2.100,7.673) (2.054,7.700) (2.000,7.731) (1.950,7.759) (1.900,7.786) (1.850,7.812) (1.800,7.837) (1.750,7.861) (1.700,7.884) (1.650,7.907) (1.600,7.928) (1.550,7.949) (1.500,7.969) (1.450,7.989) (1.400,8.007) (1.350,8.025) (1.300,8.042) (1.250,8.058) (1.200,8.074) (1.150,8.089) (1.100,8.103) (1.050,8.117) (1.000,8.130) (0.950,8.142) (0.900,8.153) (0.850,8.164) (0.800,8.175) (0.750,8.184) (0.700,8.193) (0.650,8.201) (0.600,8.209) (0.550,8.216) (0.500,8.223) (0.450,8.228) (0.400,8.234) (0.350,8.238) (0.300,8.242) (0.250,8.245) (0.200,8.248) (0.150,8.250) (0.100,8.252) (0.050,8.253) (0.000,8.253) (-0.050,8.253) (-0.100,8.252) (-0.150,8.250) (-0.200,8.248) (-0.250,8.245) (-0.300,8.242) (-0.350,8.238) (-0.400,8.234) (-0.450,8.228) (-0.500,8.223) (-0.550,8.216) (-0.600,8.209) (-0.650,8.201) (-0.700,8.193) (-0.750,8.184) (-0.800,8.175) (-0.850,8.164) (-0.900,8.153) (-0.950,8.142) (-1.000,8.130) (-1.050,8.117) (-1.100,8.103) (-1.150,8.089) (-1.200,8.074) (-1.250,8.058) (-1.300,8.042) (-1.350,8.025) (-1.400,8.007) (-1.450,7.989) (-1.500,7.969) (-1.548,7.950) (-1.600,7.928) (-1.650,7.907) (-1.700,7.884) (-1.750,7.861) (-1.800,7.837) (-1.850,7.812) (-1.900,7.786) (-1.950,7.759) (-2.000,7.731) (-2.050,7.702) (-2.096,7.675) (-2.150,7.642) (-2.200,7.610) (-2.250,7.577) (-2.300,7.543) (-2.350,7.508) (-2.395,7.475) (-2.450,7.434) (-2.493,7.400) (-2.550,7.354) (-2.600,7.313) (-2.644,7.275) (-2.700,7.225) (-2.750,7.178) (-2.800,7.130) (-2.850,7.080) (-2.900,7.028) (-2.949,6.975) (-2.993,6.925) (-3.036,6.875) (-3.075,6.828) (-3.117,6.775) (-3.150,6.733) (-3.193,6.675) (-3.225,6.630) (-3.263,6.575) (-3.296,6.525) (-3.325,6.480) (-3.359,6.425) (-3.389,6.375) (-3.418,6.325) (-3.445,6.275) (-3.472,6.225) (-3.497,6.175) (-3.521,6.125) (-3.545,6.075) (-3.567,6.025) (-3.589,5.975) (-3.609,5.925) (-3.629,5.875) (-3.647,5.825) (-3.665,5.775) (-3.682,5.725) (-3.698,5.675) (-3.713,5.625) (-3.728,5.575) (-3.741,5.525) (-3.754,5.475) (-3.765,5.425) (-3.776,5.375) (-3.787,5.325) (-3.796,5.275) (-3.804,5.225) (-3.812,5.175) (-3.819,5.125) (-3.825,5.075) (-3.830,5.025) (-3.835,4.975) (-3.839,4.925) (-3.842,4.875) (-3.844,4.825) (-3.845,4.775) (-3.846,4.725) (-3.846,4.675) (-3.845,4.625) (-3.843,4.575) (-3.841,4.525) (-3.837,4.475) (-3.833,4.425) (-3.828,4.375) (-3.823,4.325) (-3.816,4.275) (-3.809,4.225) (-3.801,4.175) (-3.792,4.125) (-3.783,4.075) (-3.772,4.025) (-3.761,3.975) (-3.749,3.925) (-3.736,3.875) (-3.722,3.825) (-3.708,3.775) (-3.692,3.725) (-3.676,3.675) (-3.659,3.625) (-3.641,3.575) (-3.622,3.525) (-3.602,3.475) (-3.582,3.425) (-3.560,3.375) (-3.538,3.325) (-3.514,3.275) (-3.490,3.225) (-3.465,3.175) (-3.439,3.125) (-3.412,3.075) (-3.384,3.025) (-3.355,2.975) (-3.325,2.926) (-3.294,2.875) (-3.262,2.825) (-3.229,2.775) (-3.200,2.733) (-3.160,2.675) (-3.125,2.627) (-3.087,2.575) (-3.050,2.527) (-3.009,2.475) (-2.975,2.432) (-2.928,2.375) (-2.886,2.325) (-2.850,2.284) (-2.800,2.228) (-2.752,2.175) (-2.705,2.125) (-2.657,2.075) (-2.609,2.025) (-2.559,1.975) (-2.508,1.925) (-2.456,1.875) (-2.403,1.825) (-2.350,1.777) (-2.300,1.731) (-2.250,1.687) (-2.209,1.650) (-2.151,1.600) (-2.100,1.556) (-2.050,1.514) (-2.003,1.475) (-1.950,1.431) (-1.900,1.391) (-1.850,1.350) (-1.800,1.311) (-1.755,1.275) (-1.700,1.232) (-1.659,1.200) (-1.600,1.155) (-1.550,1.117) (-1.500,1.079) (-1.450,1.041) (-1.400,1.004) (-1.350,0.966) (-1.300,0.929) (-1.250,0.893) (-1.200,0.856) (-1.158,0.825) (-1.100,0.783) (-1.055,0.750) (-1.000,0.711) (-0.951,0.675) (-0.900,0.639) (-0.850,0.603) (-0.800,0.567) (-0.750,0.532) (-0.706,0.500) (-0.650,0.461) (-0.600,0.425) (-0.550,0.390) (-0.500,0.354) (-0.450,0.319) (-0.400,0.284) (-0.354,0.250) (-0.300,0.213) (-0.250,0.177) (-0.220,0.155)}
\EQ{(-0.220,-0.155) (-0.275,-0.195) (-0.319,-0.225) (-0.375,-0.266) (-0.424,-0.300) (-0.475,-0.337) (-0.525,-0.372) (-0.575,-0.407) (-0.625,-0.443) (-0.671,-0.475) (-0.725,-0.514) (-0.775,-0.550) (-0.825,-0.586) (-0.875,-0.622) (-0.925,-0.658) (-0.975,-0.694) (-1.019,-0.725) (-1.075,-0.767) (-1.121,-0.800) (-1.175,-0.840) (-1.223,-0.875) (-1.275,-0.914) (-1.323,-0.950) (-1.375,-0.989) (-1.422,-1.025) (-1.475,-1.066) (-1.520,-1.100) (-1.575,-1.143) (-1.616,-1.175) (-1.675,-1.222) (-1.725,-1.263) (-1.771,-1.300) (-1.825,-1.344) (-1.875,-1.386) (-1.921,-1.425) (-1.975,-1.472) (-2.025,-1.515) (-2.064,-1.550) (-2.120,-1.600) (-2.174,-1.650) (-2.225,-1.698) (-2.275,-1.746) (-2.325,-1.795) (-2.375,-1.845) (-2.425,-1.896) (-2.475,-1.949) (-2.522,-2.000) (-2.567,-2.050) (-2.611,-2.100) (-2.650,-2.146) (-2.695,-2.200) (-2.735,-2.250) (-2.774,-2.300) (-2.812,-2.350) (-2.849,-2.400) (-2.885,-2.450) (-2.919,-2.500) (-2.950,-2.546) (-2.985,-2.600) (-3.016,-2.650) (-3.046,-2.700) (-3.075,-2.750) (-3.100,-2.794) (-3.125,-2.841) (-3.150,-2.889) (-3.175,-2.939) (-3.200,-2.991) (-3.225,-3.045) (-3.249,-3.100) (-3.270,-3.150) (-3.290,-3.200) (-3.309,-3.250) (-3.326,-3.300) (-3.343,-3.350) (-3.359,-3.400) (-3.374,-3.450) (-3.389,-3.500) (-3.402,-3.550) (-3.414,-3.600) (-3.426,-3.650) (-3.436,-3.700) (-3.446,-3.750) (-3.455,-3.800) (-3.462,-3.850) (-3.469,-3.900) (-3.476,-3.950) (-3.481,-4.000) (-3.485,-4.050) (-3.489,-4.100) (-3.491,-4.150) (-3.493,-4.200) (-3.494,-4.250) (-3.494,-4.300) (-3.493,-4.350) (-3.492,-4.400) (-3.489,-4.450) (-3.486,-4.500) (-3.482,-4.550) (-3.477,-4.600) (-3.471,-4.650) (-3.464,-4.700) (-3.456,-4.750) (-3.447,-4.800) (-3.438,-4.850) (-3.428,-4.900) (-3.416,-4.950) (-3.404,-5.000) (-3.391,-5.050) (-3.377,-5.100) (-3.362,-5.150) (-3.346,-5.200) (-3.329,-5.250) (-3.311,-5.300) (-3.292,-5.350) (-3.272,-5.400) (-3.251,-5.450) (-3.229,-5.500) (-3.206,-5.550) (-3.182,-5.600) (-3.157,-5.650) (-3.131,-5.700) (-3.103,-5.750) (-3.075,-5.799) (-3.044,-5.850) (-3.013,-5.900) (-2.981,-5.950) (-2.950,-5.995) (-2.912,-6.050) (-2.875,-6.100) (-2.837,-6.150) (-2.800,-6.196) (-2.755,-6.250) (-2.712,-6.300) (-2.675,-6.342) (-2.625,-6.396) (-2.575,-6.447) (-2.525,-6.496) (-2.475,-6.544) (-2.425,-6.589) (-2.384,-6.625) (-2.325,-6.675) (-2.275,-6.716) (-2.231,-6.750) (-2.175,-6.792) (-2.130,-6.825) (-2.075,-6.863) (-2.025,-6.897) (-1.982,-6.925) (-1.925,-6.961) (-1.875,-6.991) (-1.825,-7.020) (-1.775,-7.048) (-1.725,-7.074) (-1.676,-7.100) (-1.625,-7.125) (-1.575,-7.149) (-1.525,-7.172) (-1.475,-7.194) (-1.425,-7.215) (-1.375,-7.236) (-1.325,-7.255) (-1.275,-7.274) (-1.225,-7.292) (-1.175,-7.309) (-1.125,-7.325) (-1.075,-7.340) (-1.025,-7.355) (-0.975,-7.369) (-0.925,-7.382) (-0.875,-7.395) (-0.825,-7.406) (-0.775,-7.417) (-0.725,-7.427) (-0.675,-7.437) (-0.625,-7.446) (-0.575,-7.454) (-0.525,-7.461) (-0.475,-7.468) (-0.425,-7.474) (-0.375,-7.479) (-0.325,-7.484) (-0.275,-7.488) (-0.225,-7.492) (-0.175,-7.494) (-0.125,-7.496) (-0.075,-7.498) (-0.025,-7.498) (0.050,-7.498) (0.100,-7.497) (0.150,-7.495) (0.200,-7.493) (0.250,-7.490) (0.300,-7.486) (0.350,-7.482) (0.400,-7.477) (0.450,-7.471) (0.500,-7.465) (0.550,-7.458) (0.599,-7.450) (0.650,-7.441) (0.700,-7.432) (0.750,-7.422) (0.800,-7.412) (0.850,-7.400) (0.900,-7.388) (0.950,-7.376) (1.000,-7.362) (1.050,-7.348) (1.100,-7.333) (1.150,-7.317) (1.200,-7.300) (1.250,-7.283) (1.300,-7.265) (1.350,-7.246) (1.400,-7.226) (1.450,-7.205) (1.500,-7.183) (1.550,-7.161) (1.600,-7.137) (1.650,-7.113) (1.700,-7.088) (1.750,-7.061) (1.800,-7.034) (1.850,-7.005) (1.900,-6.976) (1.950,-6.945) (2.000,-6.913) (2.050,-6.880) (2.094,-6.850) (2.150,-6.811) (2.198,-6.775) (2.250,-6.735) (2.295,-6.700) (2.350,-6.654) (2.400,-6.611) (2.441,-6.575) (2.495,-6.525) (2.547,-6.475) (2.597,-6.425) (2.644,-6.375) (2.690,-6.325) (2.725,-6.286) (2.775,-6.227) (2.817,-6.175) (2.850,-6.133) (2.893,-6.075) (2.925,-6.031) (2.964,-5.975) (2.997,-5.925) (3.025,-5.881) (3.060,-5.825) (3.089,-5.775) (3.117,-5.725) (3.144,-5.675) (3.170,-5.625) (3.195,-5.575) (3.218,-5.525) (3.241,-5.475) (3.262,-5.425) (3.282,-5.375) (3.302,-5.325) (3.320,-5.275) (3.338,-5.225) (3.354,-5.175) (3.370,-5.125) (3.384,-5.075) (3.398,-5.025) (3.410,-4.975) (3.422,-4.925) (3.433,-4.875) (3.443,-4.825) (3.452,-4.775) (3.460,-4.725) (3.467,-4.675) (3.474,-4.625) (3.479,-4.575) (3.484,-4.525) (3.488,-4.475) (3.491,-4.425) (3.493,-4.375) (3.494,-4.325) (3.494,-4.275) (3.494,-4.225) (3.492,-4.175) (3.490,-4.125) (3.487,-4.075) (3.483,-4.025) (3.478,-3.975) (3.473,-3.925) (3.466,-3.875) (3.459,-3.825) (3.450,-3.775) (3.441,-3.725) (3.431,-3.675) (3.420,-3.625) (3.408,-3.575) (3.395,-3.525) (3.382,-3.475) (3.367,-3.425) (3.352,-3.375) (3.335,-3.325) (3.318,-3.275) (3.300,-3.227) (3.280,-3.175) (3.260,-3.125) (3.238,-3.075) (3.216,-3.025) (3.193,-2.975) (3.168,-2.925) (3.143,-2.875) (3.117,-2.825) (3.089,-2.775) (3.061,-2.725) (3.031,-2.675) (3.001,-2.625) (2.969,-2.575) (2.936,-2.525) (2.902,-2.475) (2.867,-2.425) (2.831,-2.375) (2.800,-2.334) (2.755,-2.275) (2.715,-2.225) (2.675,-2.176) (2.632,-2.125) (2.589,-2.075) (2.550,-2.031) (2.500,-1.976) (2.452,-1.925) (2.404,-1.875) (2.355,-1.825) (2.305,-1.775) (2.254,-1.725) (2.201,-1.675) (2.150,-1.628) (2.100,-1.582) (2.050,-1.537) (2.008,-1.500) (1.950,-1.450) (1.900,-1.407) (1.850,-1.365) (1.802,-1.325) (1.750,-1.283) (1.710,-1.250) (1.650,-1.202) (1.600,-1.163) (1.552,-1.125) (1.500,-1.085) (1.455,-1.050) (1.400,-1.008) (1.357,-0.975) (1.300,-0.933) (1.257,-0.900) (1.200,-0.858) (1.155,-0.825) (1.100,-0.785) (1.053,-0.750) (1.000,-0.712) (0.950,-0.676) (0.900,-0.640) (0.850,-0.604) (0.800,-0.568) (0.750,-0.532) (0.706,-0.500) (0.650,-0.461) (0.600,-0.425) (0.550,-0.390) (0.500,-0.354) (0.450,-0.319) (0.400,-0.284) (0.354,-0.250) (0.300,-0.213) (0.250,-0.177) (0.206,-0.144)}
\EQ{(0.206,0.144) (0.248,0.175) (0.300,0.213) (0.350,0.248) (0.400,0.284) (0.450,0.319) (0.495,0.350) (0.550,0.390) (0.600,0.425) (0.650,0.461) (0.700,0.496) (0.741,0.525) (0.800,0.568) (0.845,0.600) (0.900,0.640) (0.949,0.675) (1.000,0.712) (1.050,0.748) (1.100,0.785) (1.150,0.821) (1.200,0.858) (1.250,0.895) (1.300,0.933) (1.350,0.970) (1.400,1.008) (1.450,1.046) (1.500,1.085) (1.550,1.124) (1.600,1.163) (1.647,1.200) (1.700,1.242) (1.741,1.275) (1.800,1.324) (1.850,1.365) (1.892,1.400) (1.950,1.450) (2.000,1.493) (2.050,1.537) (2.092,1.575) (2.147,1.625) (2.200,1.674) (2.250,1.722) (2.300,1.770) (2.350,1.820) (2.400,1.871) (2.450,1.923) (2.499,1.975) (2.545,2.025) (2.589,2.075) (2.625,2.117) (2.674,2.175) (2.715,2.225) (2.750,2.269) (2.793,2.325) (2.825,2.367) (2.867,2.425) (2.900,2.472) (2.936,2.525) (2.969,2.575) (3.000,2.624) (3.031,2.675) (3.061,2.725) (3.089,2.775) (3.117,2.825) (3.143,2.875) (3.168,2.925) (3.193,2.975) (3.216,3.025) (3.238,3.075) (3.260,3.125) (3.280,3.175) (3.299,3.225) (3.318,3.275) (3.335,3.325) (3.352,3.375) (3.367,3.425) (3.382,3.475) (3.395,3.525) (3.408,3.575) (3.420,3.625) (3.431,3.675) (3.441,3.725) (3.450,3.775) (3.459,3.825) (3.466,3.875) (3.473,3.925) (3.478,3.975) (3.483,4.025) (3.487,4.075) (3.490,4.125) (3.492,4.175) (3.494,4.225) (3.494,4.275) (3.494,4.325) (3.493,4.375) (3.491,4.425) (3.488,4.475) (3.484,4.525) (3.479,4.575) (3.474,4.625) (3.467,4.675) (3.460,4.725) (3.452,4.775) (3.443,4.825) (3.433,4.875) (3.422,4.925) (3.410,4.975) (3.398,5.025) (3.384,5.075) (3.370,5.125) (3.354,5.175) (3.338,5.225) (3.320,5.275) (3.302,5.325) (3.282,5.375) (3.262,5.425) (3.241,5.475) (3.218,5.525) (3.195,5.575) (3.170,5.625) (3.144,5.675) (3.117,5.725) (3.089,5.775) (3.060,5.825) (3.029,5.875) (3.000,5.921) (2.964,5.975) (2.929,6.025) (2.900,6.066) (2.856,6.125) (2.825,6.165) (2.776,6.225) (2.734,6.275) (2.700,6.314) (2.650,6.369) (2.600,6.422) (2.550,6.472) (2.500,6.520) (2.450,6.567) (2.400,6.611) (2.355,6.650) (2.300,6.696) (2.250,6.735) (2.200,6.774) (2.150,6.811) (2.100,6.846) (2.058,6.875) (2.000,6.913) (1.950,6.945) (1.901,6.975) (1.850,7.005) (1.800,7.034) (1.750,7.061) (1.700,7.088) (1.650,7.113) (1.600,7.137) (1.550,7.161) (1.500,7.183) (1.450,7.205) (1.402,7.225) (1.350,7.246) (1.300,7.265) (1.250,7.283) (1.201,7.300) (1.150,7.317) (1.100,7.333) (1.050,7.348) (1.000,7.362) (0.950,7.376) (0.900,7.388) (0.850,7.400) (0.800,7.412) (0.750,7.422) (0.700,7.432) (0.650,7.441) (0.600,7.450) (0.550,7.458) (0.500,7.465) (0.450,7.471) (0.400,7.477) (0.350,7.482) (0.300,7.486) (0.250,7.490) (0.200,7.493) (0.150,7.495) (0.100,7.497) (0.050,7.498) (0.000,7.498) (-0.050,7.498) (-0.100,7.497) (-0.150,7.495) (-0.200,7.493) (-0.250,7.490) (-0.300,7.486) (-0.350,7.482) (-0.400,7.477) (-0.450,7.471) (-0.500,7.465) (-0.550,7.458) (-0.599,7.450) (-0.650,7.441) (-0.700,7.432) (-0.750,7.422) (-0.800,7.412) (-0.850,7.400) (-0.900,7.388) (-0.950,7.376) (-1.000,7.362) (-1.050,7.348) (-1.100,7.333) (-1.150,7.317) (-1.200,7.300) (-1.250,7.283) (-1.300,7.265) (-1.350,7.246) (-1.400,7.226) (-1.450,7.205) (-1.500,7.183) (-1.550,7.161) (-1.600,7.137) (-1.650,7.113) (-1.700,7.088) (-1.750,7.061) (-1.800,7.034) (-1.850,7.005) (-1.900,6.976) (-1.950,6.945) (-2.000,6.913) (-2.050,6.880) (-2.094,6.850) (-2.150,6.811) (-2.198,6.775) (-2.250,6.735) (-2.295,6.700) (-2.350,6.654) (-2.400,6.611) (-2.441,6.575) (-2.495,6.525) (-2.547,6.475) (-2.597,6.425) (-2.644,6.375) (-2.690,6.325) (-2.725,6.286) (-2.775,6.227) (-2.817,6.175) (-2.850,6.133) (-2.893,6.075) (-2.925,6.031) (-2.964,5.975) (-2.997,5.925) (-3.025,5.881) (-3.060,5.825) (-3.089,5.775) (-3.117,5.725) (-3.144,5.675) (-3.170,5.625) (-3.195,5.575) (-3.218,5.525) (-3.241,5.475) (-3.262,5.425) (-3.282,5.375) (-3.302,5.325) (-3.320,5.275) (-3.338,5.225) (-3.354,5.175) (-3.370,5.125) (-3.384,5.075) (-3.398,5.025) (-3.410,4.975) (-3.422,4.925) (-3.433,4.875) (-3.443,4.825) (-3.452,4.775) (-3.460,4.725) (-3.467,4.675) (-3.474,4.625) (-3.479,4.575) (-3.484,4.525) (-3.488,4.475) (-3.491,4.425) (-3.493,4.375) (-3.494,4.325) (-3.494,4.275) (-3.494,4.225) (-3.492,4.175) (-3.490,4.125) (-3.487,4.075) (-3.483,4.025) (-3.478,3.975) (-3.473,3.925) (-3.466,3.875) (-3.459,3.825) (-3.450,3.775) (-3.441,3.725) (-3.431,3.675) (-3.420,3.625) (-3.408,3.575) (-3.395,3.525) (-3.382,3.475) (-3.367,3.425) (-3.352,3.375) (-3.335,3.325) (-3.318,3.275) (-3.300,3.227) (-3.280,3.175) (-3.260,3.125) (-3.238,3.075) (-3.216,3.025) (-3.193,2.975) (-3.168,2.925) (-3.143,2.875) (-3.117,2.825) (-3.089,2.775) (-3.061,2.725) (-3.031,2.675) (-3.001,2.625) (-2.969,2.575) (-2.936,2.525) (-2.902,2.475) (-2.867,2.425) (-2.831,2.375) (-2.800,2.334) (-2.755,2.275) (-2.715,2.225) (-2.675,2.176) (-2.632,2.125) (-2.589,2.075) (-2.550,2.031) (-2.500,1.976) (-2.452,1.925) (-2.404,1.875) (-2.355,1.825) (-2.305,1.775) (-2.254,1.725) (-2.201,1.675) (-2.150,1.628) (-2.100,1.582) (-2.050,1.537) (-2.008,1.500) (-1.950,1.450) (-1.900,1.407) (-1.850,1.365) (-1.802,1.325) (-1.750,1.283) (-1.710,1.250) (-1.650,1.202) (-1.600,1.163) (-1.552,1.125) (-1.500,1.085) (-1.455,1.050) (-1.400,1.008) (-1.357,0.975) (-1.300,0.933) (-1.257,0.900) (-1.200,0.858) (-1.155,0.825) (-1.100,0.785) (-1.053,0.750) (-1.000,0.712) (-0.950,0.676) (-0.900,0.640) (-0.850,0.604) (-0.800,0.568) (-0.750,0.532) (-0.706,0.500) (-0.650,0.461) (-0.600,0.425) (-0.550,0.390) (-0.500,0.354) (-0.450,0.319) (-0.400,0.284) (-0.354,0.250) (-0.300,0.213) (-0.250,0.177) (-0.220,0.155)}
\EQ{(-0.220,-0.155) (-0.275,-0.195) (-0.319,-0.225) (-0.375,-0.266) (-0.424,-0.300) (-0.475,-0.337) (-0.525,-0.372) (-0.575,-0.408) (-0.625,-0.443) (-0.671,-0.475) (-0.725,-0.515) (-0.775,-0.550) (-0.825,-0.586) (-0.875,-0.622) (-0.925,-0.659) (-0.975,-0.695) (-1.017,-0.725) (-1.075,-0.768) (-1.119,-0.800) (-1.175,-0.842) (-1.219,-0.875) (-1.275,-0.917) (-1.319,-0.950) (-1.375,-0.994) (-1.416,-1.025) (-1.475,-1.071) (-1.525,-1.111) (-1.574,-1.150) (-1.625,-1.191) (-1.667,-1.225) (-1.725,-1.274) (-1.775,-1.316) (-1.815,-1.350) (-1.872,-1.400) (-1.925,-1.447) (-1.975,-1.492) (-2.025,-1.538) (-2.064,-1.575) (-2.117,-1.625) (-2.167,-1.675) (-2.217,-1.725) (-2.265,-1.775) (-2.300,-1.812) (-2.335,-1.850) (-2.375,-1.894) (-2.424,-1.950) (-2.467,-2.000) (-2.500,-2.040) (-2.548,-2.100) (-2.587,-2.150) (-2.624,-2.200) (-2.661,-2.250) (-2.696,-2.300) (-2.725,-2.343) (-2.762,-2.400) (-2.794,-2.450) (-2.824,-2.500) (-2.850,-2.545) (-2.875,-2.589) (-2.900,-2.635) (-2.925,-2.683) (-2.950,-2.733) (-2.975,-2.786) (-3.000,-2.841) (-3.025,-2.899) (-3.046,-2.950) (-3.065,-3.000) (-3.083,-3.050) (-3.100,-3.099) (-3.117,-3.150) (-3.132,-3.200) (-3.146,-3.250) (-3.159,-3.300) (-3.171,-3.350) (-3.183,-3.400) (-3.193,-3.450) (-3.202,-3.500) (-3.211,-3.550) (-3.218,-3.600) (-3.224,-3.650) (-3.230,-3.700) (-3.235,-3.750) (-3.238,-3.800) (-3.241,-3.850) (-3.243,-3.900) (-3.244,-3.950) (-3.244,-4.000) (-3.243,-4.050) (-3.241,-4.100) (-3.238,-4.150) (-3.234,-4.200) (-3.229,-4.250) (-3.224,-4.300) (-3.217,-4.350) (-3.210,-4.400) (-3.201,-4.450) (-3.192,-4.500) (-3.182,-4.550) (-3.170,-4.600) (-3.158,-4.650) (-3.145,-4.700) (-3.130,-4.750) (-3.115,-4.800) (-3.099,-4.850) (-3.081,-4.900) (-3.063,-4.950) (-3.043,-5.000) (-3.022,-5.050) (-3.001,-5.100) (-2.978,-5.150) (-2.954,-5.200) (-2.928,-5.250) (-2.902,-5.300) (-2.875,-5.348) (-2.845,-5.400) (-2.814,-5.450) (-2.782,-5.500) (-2.750,-5.549) (-2.715,-5.600) (-2.678,-5.650) (-2.640,-5.700) (-2.601,-5.750) (-2.560,-5.800) (-2.525,-5.840) (-2.475,-5.896) (-2.425,-5.949) (-2.375,-6.000) (-2.325,-6.049) (-2.275,-6.096) (-2.225,-6.140) (-2.184,-6.175) (-2.125,-6.224) (-2.075,-6.263) (-2.026,-6.300) (-1.975,-6.337) (-1.925,-6.372) (-1.875,-6.406) (-1.825,-6.438) (-1.775,-6.469) (-1.725,-6.499) (-1.679,-6.525) (-1.634,-6.550) (-1.587,-6.575) (-1.538,-6.600) (-1.487,-6.625) (-1.433,-6.650) (-1.377,-6.675) (-1.325,-6.697) (-1.275,-6.718) (-1.225,-6.737) (-1.175,-6.756) (-1.125,-6.773) (-1.075,-6.790) (-1.025,-6.806) (-0.975,-6.821) (-0.925,-6.835) (-0.875,-6.849) (-0.825,-6.861) (-0.775,-6.873) (-0.725,-6.884) (-0.675,-6.895) (-0.625,-6.904) (-0.575,-6.913) (-0.525,-6.921) (-0.475,-6.928) (-0.425,-6.935) (-0.375,-6.941) (-0.325,-6.946) (-0.275,-6.950) (-0.225,-6.954) (-0.175,-6.956) (-0.125,-6.959) (-0.075,-6.960) (-0.025,-6.961) (0.050,-6.960) (0.100,-6.959) (0.150,-6.958) (0.200,-6.955) (0.250,-6.952) (0.300,-6.948) (0.350,-6.943) (0.400,-6.938) (0.450,-6.932) (0.500,-6.925) (0.550,-6.917) (0.600,-6.909) (0.650,-6.899) (0.700,-6.890) (0.750,-6.879) (0.800,-6.867) (0.850,-6.855) (0.900,-6.842) (0.950,-6.828) (1.000,-6.814) (1.050,-6.798) (1.100,-6.782) (1.150,-6.765) (1.200,-6.747) (1.250,-6.728) (1.300,-6.708) (1.350,-6.687) (1.400,-6.665) (1.450,-6.642) (1.500,-6.619) (1.550,-6.594) (1.600,-6.568) (1.650,-6.541) (1.700,-6.513) (1.750,-6.484) (1.800,-6.454) (1.845,-6.425) (1.900,-6.389) (1.950,-6.355) (1.992,-6.325) (2.050,-6.282) (2.092,-6.250) (2.150,-6.204) (2.200,-6.162) (2.242,-6.125) (2.297,-6.075) (2.350,-6.025) (2.400,-5.975) (2.448,-5.925) (2.494,-5.875) (2.538,-5.825) (2.575,-5.782) (2.621,-5.725) (2.660,-5.675) (2.697,-5.625) (2.732,-5.575) (2.766,-5.525) (2.799,-5.475) (2.830,-5.425) (2.859,-5.375) (2.888,-5.325) (2.915,-5.275) (2.941,-5.225) (2.966,-5.175) (2.989,-5.125) (3.012,-5.075) (3.033,-5.025) (3.053,-4.975) (3.072,-4.925) (3.090,-4.875) (3.107,-4.825) (3.123,-4.775) (3.138,-4.725) (3.151,-4.675) (3.164,-4.625) (3.176,-4.575) (3.187,-4.525) (3.197,-4.475) (3.206,-4.425) (3.214,-4.375) (3.221,-4.325) (3.227,-4.275) (3.232,-4.225) (3.236,-4.175) (3.239,-4.125) (3.242,-4.075) (3.243,-4.025) (3.244,-3.975) (3.243,-3.925) (3.242,-3.875) (3.240,-3.825) (3.237,-3.775) (3.232,-3.725) (3.227,-3.675) (3.221,-3.625) (3.214,-3.575) (3.207,-3.525) (3.198,-3.475) (3.188,-3.425) (3.177,-3.375) (3.165,-3.325) (3.153,-3.275) (3.139,-3.225) (3.124,-3.175) (3.109,-3.125) (3.092,-3.075) (3.075,-3.027) (3.056,-2.975) (3.036,-2.925) (3.015,-2.875) (2.993,-2.825) (2.970,-2.775) (2.946,-2.725) (2.921,-2.675) (2.895,-2.625) (2.867,-2.575) (2.839,-2.525) (2.809,-2.475) (2.778,-2.425) (2.750,-2.381) (2.713,-2.325) (2.678,-2.275) (2.643,-2.225) (2.606,-2.175) (2.575,-2.135) (2.528,-2.075) (2.488,-2.025) (2.450,-1.980) (2.402,-1.925) (2.358,-1.875) (2.312,-1.825) (2.275,-1.785) (2.225,-1.733) (2.175,-1.683) (2.125,-1.633) (2.075,-1.585) (2.038,-1.550) (2.000,-1.515) (1.957,-1.475) (1.901,-1.425) (1.850,-1.381) (1.800,-1.337) (1.756,-1.300) (1.700,-1.253) (1.650,-1.212) (1.605,-1.175) (1.550,-1.131) (1.500,-1.091) (1.450,-1.052) (1.400,-1.013) (1.351,-0.975) (1.300,-0.936) (1.252,-0.900) (1.200,-0.861) (1.152,-0.825) (1.100,-0.787) (1.051,-0.750) (1.000,-0.713) (0.950,-0.677) (0.900,-0.640) (0.850,-0.604) (0.800,-0.568) (0.750,-0.532) (0.705,-0.500) (0.650,-0.461) (0.600,-0.425) (0.550,-0.390) (0.500,-0.354) (0.450,-0.319) (0.400,-0.284) (0.354,-0.250) (0.300,-0.213) (0.250,-0.177) (0.206,-0.144)}
\EQ{(0.206,0.144) (0.248,0.175) (0.300,0.213) (0.350,0.248) (0.400,0.284) (0.450,0.319) (0.495,0.350) (0.550,0.390) (0.600,0.425) (0.650,0.461) (0.700,0.497) (0.750,0.532) (0.800,0.568) (0.845,0.600) (0.900,0.640) (0.948,0.675) (1.000,0.713) (1.050,0.750) (1.100,0.787) (1.150,0.823) (1.200,0.861) (1.250,0.898) (1.300,0.936) (1.350,0.974) (1.400,1.013) (1.448,1.050) (1.500,1.091) (1.543,1.125) (1.600,1.171) (1.650,1.212) (1.697,1.250) (1.750,1.295) (1.800,1.337) (1.844,1.375) (1.900,1.424) (1.950,1.469) (2.000,1.515) (2.038,1.550) (2.075,1.585) (2.117,1.625) (2.167,1.675) (2.217,1.725) (2.265,1.775) (2.300,1.812) (2.335,1.850) (2.375,1.894) (2.424,1.950) (2.467,2.000) (2.500,2.040) (2.548,2.100) (2.587,2.150) (2.624,2.200) (2.661,2.250) (2.696,2.300) (2.725,2.343) (2.762,2.400) (2.794,2.450) (2.824,2.500) (2.850,2.545) (2.875,2.589) (2.900,2.635) (2.925,2.683) (2.950,2.733) (2.975,2.786) (3.000,2.841) (3.025,2.899) (3.046,2.950) (3.065,3.000) (3.083,3.050) (3.100,3.099) (3.117,3.150) (3.132,3.200) (3.146,3.250) (3.159,3.300) (3.171,3.350) (3.183,3.400) (3.193,3.450) (3.202,3.500) (3.211,3.550) (3.218,3.600) (3.224,3.650) (3.230,3.700) (3.235,3.750) (3.238,3.800) (3.241,3.850) (3.243,3.900) (3.244,3.950) (3.244,4.000) (3.243,4.050) (3.241,4.100) (3.238,4.150) (3.234,4.200) (3.229,4.250) (3.224,4.300) (3.217,4.350) (3.210,4.400) (3.201,4.450) (3.192,4.500) (3.182,4.550) (3.170,4.600) (3.158,4.650) (3.145,4.700) (3.130,4.750) (3.115,4.800) (3.099,4.850) (3.081,4.900) (3.063,4.950) (3.043,5.000) (3.022,5.050) (3.001,5.100) (2.978,5.150) (2.954,5.200) (2.928,5.250) (2.902,5.300) (2.875,5.348) (2.845,5.400) (2.814,5.450) (2.782,5.500) (2.750,5.549) (2.715,5.600) (2.678,5.650) (2.640,5.700) (2.601,5.750) (2.560,5.800) (2.525,5.840) (2.475,5.896) (2.425,5.949) (2.375,6.000) (2.325,6.049) (2.275,6.096) (2.225,6.140) (2.184,6.175) (2.125,6.224) (2.075,6.263) (2.026,6.300) (1.975,6.337) (1.925,6.372) (1.875,6.406) (1.825,6.438) (1.775,6.469) (1.725,6.499) (1.679,6.525) (1.634,6.550) (1.587,6.575) (1.538,6.600) (1.487,6.625) (1.433,6.650) (1.377,6.675) (1.325,6.697) (1.275,6.718) (1.225,6.737) (1.175,6.756) (1.125,6.773) (1.075,6.790) (1.025,6.806) (0.975,6.821) (0.925,6.835) (0.875,6.849) (0.825,6.861) (0.775,6.873) (0.725,6.884) (0.675,6.895) (0.625,6.904) (0.575,6.913) (0.525,6.921) (0.475,6.928) (0.425,6.935) (0.375,6.941) (0.325,6.946) (0.275,6.950) (0.225,6.954) (0.175,6.956) (0.125,6.959) (0.075,6.960) (0.025,6.961) (-0.050,6.960) (-0.100,6.959) (-0.150,6.958) (-0.200,6.955) (-0.250,6.952) (-0.300,6.948) (-0.350,6.943) (-0.400,6.938) (-0.450,6.932) (-0.500,6.925) (-0.550,6.917) (-0.600,6.909) (-0.650,6.899) (-0.700,6.890) (-0.750,6.879) (-0.800,6.867) (-0.850,6.855) (-0.900,6.842) (-0.950,6.828) (-1.000,6.814) (-1.050,6.798) (-1.100,6.782) (-1.150,6.765) (-1.200,6.747) (-1.250,6.728) (-1.300,6.708) (-1.350,6.687) (-1.400,6.665) (-1.450,6.642) (-1.500,6.619) (-1.550,6.594) (-1.600,6.568) (-1.650,6.541) (-1.700,6.513) (-1.750,6.484) (-1.800,6.454) (-1.845,6.425) (-1.900,6.389) (-1.950,6.355) (-1.992,6.325) (-2.050,6.282) (-2.092,6.250) (-2.150,6.204) (-2.200,6.162) (-2.242,6.125) (-2.297,6.075) (-2.350,6.025) (-2.400,5.975) (-2.448,5.925) (-2.494,5.875) (-2.538,5.825) (-2.575,5.782) (-2.621,5.725) (-2.660,5.675) (-2.697,5.625) (-2.732,5.575) (-2.766,5.525) (-2.799,5.475) (-2.830,5.425) (-2.859,5.375) (-2.888,5.325) (-2.915,5.275) (-2.941,5.225) (-2.966,5.175) (-2.989,5.125) (-3.012,5.075) (-3.033,5.025) (-3.053,4.975) (-3.072,4.925) (-3.090,4.875) (-3.107,4.825) (-3.123,4.775) (-3.138,4.725) (-3.151,4.675) (-3.164,4.625) (-3.176,4.575) (-3.187,4.525) (-3.197,4.475) (-3.206,4.425) (-3.214,4.375) (-3.221,4.325) (-3.227,4.275) (-3.232,4.225) (-3.236,4.175) (-3.239,4.125) (-3.242,4.075) (-3.243,4.025) (-3.244,3.975) (-3.243,3.925) (-3.242,3.875) (-3.240,3.825) (-3.237,3.775) (-3.232,3.725) (-3.227,3.675) (-3.221,3.625) (-3.214,3.575) (-3.207,3.525) (-3.198,3.475) (-3.188,3.425) (-3.177,3.375) (-3.165,3.325) (-3.153,3.275) (-3.139,3.225) (-3.124,3.175) (-3.109,3.125) (-3.092,3.075) (-3.075,3.027) (-3.056,2.975) (-3.036,2.925) (-3.015,2.875) (-2.993,2.825) (-2.970,2.775) (-2.946,2.725) (-2.921,2.675) (-2.895,2.625) (-2.867,2.575) (-2.839,2.525) (-2.809,2.475) (-2.778,2.425) (-2.750,2.381) (-2.713,2.325) (-2.678,2.275) (-2.643,2.225) (-2.606,2.175) (-2.575,2.135) (-2.528,2.075) (-2.488,2.025) (-2.450,1.980) (-2.402,1.925) (-2.358,1.875) (-2.312,1.825) (-2.275,1.785) (-2.225,1.733) (-2.175,1.683) (-2.125,1.633) (-2.075,1.585) (-2.038,1.550) (-2.000,1.515) (-1.957,1.475) (-1.901,1.425) (-1.850,1.381) (-1.800,1.337) (-1.756,1.300) (-1.700,1.253) (-1.650,1.212) (-1.605,1.175) (-1.550,1.131) (-1.500,1.091) (-1.450,1.052) (-1.400,1.013) (-1.351,0.975) (-1.300,0.936) (-1.252,0.900) (-1.200,0.861) (-1.152,0.825) (-1.100,0.787) (-1.051,0.750) (-1.000,0.713) (-0.950,0.677) (-0.900,0.640) (-0.850,0.604) (-0.800,0.568) (-0.750,0.532) (-0.705,0.500) (-0.650,0.461) (-0.600,0.425) (-0.550,0.390) (-0.500,0.354) (-0.450,0.319) (-0.400,0.284) (-0.354,0.250) (-0.300,0.213) (-0.250,0.177) (-0.220,0.155)}
\EQ{(-0.220,-0.155) (-0.275,-0.195) (-0.319,-0.225) (-0.375,-0.266) (-0.424,-0.300) (-0.475,-0.337) (-0.525,-0.372) (-0.575,-0.408) (-0.625,-0.443) (-0.670,-0.475) (-0.725,-0.515) (-0.774,-0.550) (-0.825,-0.587) (-0.875,-0.623) (-0.925,-0.659) (-0.975,-0.696) (-1.025,-0.733) (-1.075,-0.770) (-1.116,-0.800) (-1.175,-0.844) (-1.216,-0.875) (-1.275,-0.920) (-1.325,-0.959) (-1.375,-0.998) (-1.425,-1.037) (-1.472,-1.075) (-1.525,-1.118) (-1.565,-1.150) (-1.625,-1.200) (-1.675,-1.242) (-1.713,-1.275) (-1.771,-1.325) (-1.825,-1.374) (-1.875,-1.419) (-1.925,-1.466) (-1.962,-1.500) (-2.000,-1.537) (-2.038,-1.575) (-2.075,-1.612) (-2.112,-1.650) (-2.150,-1.690) (-2.200,-1.744) (-2.250,-1.799) (-2.294,-1.850) (-2.336,-1.900) (-2.375,-1.948) (-2.417,-2.000) (-2.450,-2.044) (-2.491,-2.100) (-2.525,-2.147) (-2.561,-2.200) (-2.594,-2.250) (-2.625,-2.299) (-2.656,-2.350) (-2.685,-2.400) (-2.713,-2.450) (-2.740,-2.500) (-2.765,-2.550) (-2.790,-2.600) (-2.813,-2.650) (-2.835,-2.700) (-2.856,-2.750) (-2.875,-2.799) (-2.894,-2.850) (-2.912,-2.900) (-2.928,-2.950) (-2.944,-3.000) (-2.958,-3.050) (-2.972,-3.100) (-2.984,-3.150) (-2.995,-3.200) (-3.005,-3.250) (-3.015,-3.300) (-3.023,-3.350) (-3.030,-3.400) (-3.036,-3.450) (-3.041,-3.500) (-3.045,-3.550) (-3.049,-3.600) (-3.051,-3.650) (-3.052,-3.700) (-3.052,-3.750) (-3.052,-3.800) (-3.050,-3.850) (-3.047,-3.900) (-3.044,-3.950) (-3.039,-4.000) (-3.033,-4.050) (-3.027,-4.100) (-3.019,-4.150) (-3.010,-4.200) (-3.001,-4.250) (-2.990,-4.300) (-2.978,-4.350) (-2.965,-4.400) (-2.951,-4.450) (-2.937,-4.500) (-2.921,-4.550) (-2.903,-4.600) (-2.885,-4.650) (-2.866,-4.700) (-2.845,-4.750) (-2.825,-4.796) (-2.800,-4.850) (-2.776,-4.900) (-2.751,-4.950) (-2.725,-4.998) (-2.700,-5.042) (-2.666,-5.100) (-2.635,-5.150) (-2.602,-5.200) (-2.568,-5.250) (-2.533,-5.300) (-2.500,-5.344) (-2.456,-5.400) (-2.415,-5.450) (-2.375,-5.497) (-2.328,-5.550) (-2.281,-5.600) (-2.231,-5.650) (-2.180,-5.700) (-2.125,-5.750) (-2.075,-5.794) (-2.025,-5.836) (-1.977,-5.875) (-1.925,-5.915) (-1.878,-5.950) (-1.825,-5.987) (-1.775,-6.021) (-1.731,-6.050) (-1.675,-6.085) (-1.625,-6.115) (-1.575,-6.143) (-1.525,-6.170) (-1.475,-6.197) (-1.425,-6.222) (-1.375,-6.246) (-1.325,-6.269) (-1.275,-6.291) (-1.225,-6.311) (-1.175,-6.331) (-1.126,-6.350) (-1.075,-6.368) (-1.025,-6.385) (-0.975,-6.401) (-0.925,-6.417) (-0.875,-6.431) (-0.825,-6.445) (-0.775,-6.457) (-0.725,-6.469) (-0.675,-6.480) (-0.625,-6.490) (-0.575,-6.499) (-0.525,-6.508) (-0.475,-6.516) (-0.425,-6.523) (-0.375,-6.529) (-0.325,-6.534) (-0.275,-6.539) (-0.225,-6.543) (-0.175,-6.546) (-0.125,-6.548) (-0.075,-6.549) (-0.025,-6.550) (0.047,-6.550) (0.100,-6.549) (0.150,-6.547) (0.200,-6.544) (0.250,-6.541) (0.300,-6.537) (0.350,-6.532) (0.400,-6.526) (0.450,-6.519) (0.500,-6.512) (0.550,-6.504) (0.600,-6.495) (0.650,-6.485) (0.700,-6.475) (0.750,-6.463) (0.800,-6.451) (0.850,-6.438) (0.900,-6.424) (0.950,-6.409) (1.000,-6.393) (1.050,-6.377) (1.100,-6.359) (1.150,-6.341) (1.200,-6.322) (1.250,-6.301) (1.300,-6.280) (1.350,-6.257) (1.400,-6.234) (1.450,-6.209) (1.500,-6.184) (1.550,-6.157) (1.600,-6.129) (1.650,-6.100) (1.700,-6.069) (1.750,-6.038) (1.800,-6.004) (1.842,-5.975) (1.900,-5.934) (1.945,-5.900) (2.000,-5.857) (2.050,-5.816) (2.097,-5.775) (2.150,-5.728) (2.200,-5.681) (2.250,-5.632) (2.300,-5.580) (2.350,-5.525) (2.394,-5.475) (2.436,-5.425) (2.475,-5.376) (2.514,-5.325) (2.550,-5.276) (2.586,-5.225) (2.619,-5.175) (2.650,-5.126) (2.681,-5.075) (2.710,-5.025) (2.737,-4.975) (2.763,-4.925) (2.788,-4.875) (2.812,-4.825) (2.834,-4.775) (2.856,-4.725) (2.875,-4.676) (2.894,-4.625) (2.912,-4.575) (2.929,-4.525) (2.944,-4.475) (2.959,-4.425) (2.972,-4.375) (2.984,-4.325) (2.995,-4.275) (3.006,-4.225) (3.015,-4.175) (3.023,-4.125) (3.030,-4.075) (3.036,-4.025) (3.041,-3.975) (3.046,-3.925) (3.049,-3.875) (3.051,-3.825) (3.052,-3.775) (3.052,-3.725) (3.052,-3.675) (3.050,-3.625) (3.047,-3.575) (3.044,-3.525) (3.039,-3.475) (3.033,-3.425) (3.026,-3.375) (3.019,-3.325) (3.010,-3.275) (3.000,-3.225) (2.990,-3.175) (2.978,-3.125) (2.965,-3.075) (2.951,-3.025) (2.936,-2.975) (2.920,-2.925) (2.903,-2.875) (2.885,-2.825) (2.866,-2.775) (2.845,-2.725) (2.825,-2.677) (2.801,-2.625) (2.778,-2.575) (2.753,-2.525) (2.726,-2.475) (2.700,-2.427) (2.675,-2.383) (2.641,-2.325) (2.610,-2.275) (2.578,-2.225) (2.544,-2.175) (2.509,-2.125) (2.475,-2.078) (2.436,-2.025) (2.400,-1.979) (2.357,-1.925) (2.315,-1.875) (2.275,-1.828) (2.228,-1.775) (2.183,-1.725) (2.136,-1.675) (2.100,-1.638) (2.063,-1.600) (2.025,-1.562) (1.987,-1.525) (1.950,-1.489) (1.909,-1.450) (1.854,-1.400) (1.800,-1.351) (1.750,-1.307) (1.700,-1.264) (1.655,-1.225) (1.600,-1.179) (1.550,-1.138) (1.504,-1.100) (1.450,-1.057) (1.410,-1.025) (1.350,-0.978) (1.300,-0.940) (1.250,-0.901) (1.200,-0.863) (1.150,-0.826) (1.100,-0.788) (1.050,-0.751) (1.000,-0.714) (0.950,-0.678) (0.900,-0.641) (0.850,-0.605) (0.809,-0.575) (0.750,-0.533) (0.705,-0.500) (0.650,-0.461) (0.600,-0.425) (0.550,-0.390) (0.500,-0.354) (0.450,-0.319) (0.400,-0.284) (0.354,-0.250) (0.300,-0.213) (0.250,-0.177) (0.206,-0.144)}
\EQ{(0.206,0.144) (0.248,0.175) (0.300,0.213) (0.350,0.248) (0.400,0.284) (0.450,0.319) (0.495,0.350) (0.550,0.390) (0.600,0.425) (0.650,0.461) (0.700,0.497) (0.750,0.533) (0.800,0.569) (0.844,0.600) (0.900,0.641) (0.947,0.675) (1.000,0.714) (1.049,0.750) (1.100,0.788) (1.149,0.825) (1.200,0.863) (1.249,0.900) (1.300,0.940) (1.346,0.975) (1.400,1.018) (1.441,1.050) (1.500,1.097) (1.550,1.138) (1.595,1.175) (1.650,1.221) (1.700,1.264) (1.742,1.300) (1.799,1.350) (1.850,1.396) (1.900,1.442) (1.950,1.489) (1.987,1.525) (2.025,1.562) (2.063,1.600) (2.100,1.638) (2.136,1.675) (2.175,1.717) (2.225,1.771) (2.273,1.825) (2.315,1.875) (2.350,1.917) (2.397,1.975) (2.436,2.025) (2.473,2.075) (2.509,2.125) (2.544,2.175) (2.575,2.221) (2.610,2.275) (2.641,2.325) (2.671,2.375) (2.699,2.425) (2.725,2.472) (2.750,2.520) (2.775,2.570) (2.800,2.622) (2.824,2.675) (2.845,2.725) (2.866,2.775) (2.885,2.825) (2.903,2.875) (2.920,2.925) (2.936,2.975) (2.951,3.025) (2.965,3.075) (2.978,3.125) (2.990,3.175) (3.000,3.225) (3.010,3.275) (3.019,3.325) (3.026,3.375) (3.033,3.425) (3.039,3.475) (3.044,3.525) (3.047,3.575) (3.050,3.625) (3.052,3.675) (3.052,3.725) (3.052,3.775) (3.051,3.825) (3.049,3.875) (3.046,3.925) (3.041,3.975) (3.036,4.025) (3.030,4.075) (3.023,4.125) (3.015,4.175) (3.006,4.225) (2.995,4.275) (2.984,4.325) (2.972,4.375) (2.959,4.425) (2.944,4.475) (2.929,4.525) (2.912,4.575) (2.894,4.625) (2.876,4.675) (2.856,4.725) (2.834,4.775) (2.812,4.825) (2.788,4.875) (2.763,4.925) (2.737,4.975) (2.710,5.025) (2.681,5.075) (2.651,5.125) (2.619,5.175) (2.586,5.225) (2.551,5.275) (2.514,5.325) (2.476,5.375) (2.436,5.425) (2.400,5.468) (2.350,5.525) (2.305,5.575) (2.256,5.625) (2.206,5.675) (2.153,5.725) (2.100,5.773) (2.050,5.816) (2.008,5.850) (1.950,5.896) (1.900,5.934) (1.850,5.970) (1.806,6.000) (1.750,6.038) (1.700,6.069) (1.650,6.100) (1.600,6.129) (1.550,6.157) (1.500,6.184) (1.450,6.209) (1.400,6.234) (1.350,6.257) (1.300,6.280) (1.253,6.300) (1.200,6.322) (1.150,6.341) (1.100,6.359) (1.050,6.377) (1.000,6.393) (0.950,6.409) (0.900,6.424) (0.850,6.438) (0.800,6.451) (0.750,6.463) (0.700,6.475) (0.650,6.485) (0.600,6.495) (0.550,6.504) (0.500,6.512) (0.450,6.519) (0.400,6.526) (0.350,6.532) (0.300,6.537) (0.250,6.541) (0.200,6.544) (0.150,6.547) (0.100,6.549) (0.050,6.550) (0.000,6.550) (-0.050,6.550) (-0.100,6.549) (-0.150,6.547) (-0.200,6.544) (-0.250,6.541) (-0.300,6.537) (-0.350,6.532) (-0.400,6.526) (-0.450,6.519) (-0.500,6.512) (-0.550,6.504) (-0.600,6.495) (-0.650,6.485) (-0.700,6.475) (-0.750,6.463) (-0.800,6.451) (-0.850,6.438) (-0.900,6.424) (-0.950,6.409) (-1.000,6.393) (-1.050,6.377) (-1.100,6.359) (-1.150,6.341) (-1.200,6.322) (-1.250,6.301) (-1.300,6.280) (-1.350,6.257) (-1.400,6.234) (-1.450,6.209) (-1.500,6.184) (-1.550,6.157) (-1.600,6.129) (-1.650,6.100) (-1.700,6.069) (-1.750,6.038) (-1.800,6.004) (-1.842,5.975) (-1.900,5.934) (-1.945,5.900) (-2.000,5.857) (-2.050,5.816) (-2.097,5.775) (-2.150,5.728) (-2.200,5.681) (-2.250,5.632) (-2.300,5.580) (-2.350,5.525) (-2.394,5.475) (-2.436,5.425) (-2.475,5.376) (-2.514,5.325) (-2.550,5.276) (-2.586,5.225) (-2.619,5.175) (-2.650,5.126) (-2.681,5.075) (-2.710,5.025) (-2.737,4.975) (-2.763,4.925) (-2.788,4.875) (-2.812,4.825) (-2.834,4.775) (-2.856,4.725) (-2.875,4.676) (-2.894,4.625) (-2.912,4.575) (-2.929,4.525) (-2.944,4.475) (-2.959,4.425) (-2.972,4.375) (-2.984,4.325) (-2.995,4.275) (-3.006,4.225) (-3.015,4.175) (-3.023,4.125) (-3.030,4.075) (-3.036,4.025) (-3.041,3.975) (-3.046,3.925) (-3.049,3.875) (-3.051,3.825) (-3.052,3.775) (-3.052,3.725) (-3.052,3.675) (-3.050,3.625) (-3.047,3.575) (-3.044,3.525) (-3.039,3.475) (-3.033,3.425) (-3.026,3.375) (-3.019,3.325) (-3.010,3.275) (-3.000,3.225) (-2.990,3.175) (-2.978,3.125) (-2.965,3.075) (-2.951,3.025) (-2.936,2.975) (-2.920,2.925) (-2.903,2.875) (-2.885,2.825) (-2.866,2.775) (-2.845,2.725) (-2.825,2.677) (-2.801,2.625) (-2.778,2.575) (-2.753,2.525) (-2.726,2.475) (-2.700,2.427) (-2.675,2.383) (-2.641,2.325) (-2.610,2.275) (-2.578,2.225) (-2.544,2.175) (-2.509,2.125) (-2.475,2.078) (-2.436,2.025) (-2.400,1.979) (-2.357,1.925) (-2.315,1.875) (-2.275,1.828) (-2.228,1.775) (-2.183,1.725) (-2.136,1.675) (-2.100,1.638) (-2.063,1.600) (-2.025,1.562) (-1.987,1.525) (-1.950,1.489) (-1.909,1.450) (-1.854,1.400) (-1.800,1.351) (-1.750,1.307) (-1.700,1.264) (-1.655,1.225) (-1.600,1.179) (-1.550,1.138) (-1.504,1.100) (-1.450,1.057) (-1.410,1.025) (-1.350,0.978) (-1.300,0.940) (-1.250,0.901) (-1.200,0.863) (-1.150,0.826) (-1.100,0.788) (-1.050,0.751) (-1.000,0.714) (-0.950,0.678) (-0.900,0.641) (-0.850,0.605) (-0.809,0.575) (-0.750,0.533) (-0.705,0.500) (-0.650,0.461) (-0.600,0.425) (-0.550,0.390) (-0.500,0.354) (-0.450,0.319) (-0.400,0.284) (-0.354,0.250) (-0.300,0.213) (-0.250,0.177) (-0.220,0.155)}
\EQ{(-0.220,-0.155) (-0.275,-0.195) (-0.319,-0.225) (-0.375,-0.266) (-0.424,-0.300) (-0.475,-0.337) (-0.525,-0.372) (-0.575,-0.408) (-0.625,-0.444) (-0.670,-0.475) (-0.725,-0.515) (-0.774,-0.550) (-0.825,-0.587) (-0.875,-0.624) (-0.925,-0.660) (-0.975,-0.697) (-1.025,-0.734) (-1.075,-0.771) (-1.125,-0.809) (-1.175,-0.847) (-1.225,-0.885) (-1.275,-0.924) (-1.325,-0.963) (-1.372,-1.000) (-1.425,-1.043) (-1.465,-1.075) (-1.525,-1.124) (-1.575,-1.166) (-1.615,-1.200) (-1.672,-1.250) (-1.725,-1.297) (-1.775,-1.342) (-1.825,-1.389) (-1.863,-1.425) (-1.900,-1.461) (-1.940,-1.500) (-1.990,-1.550) (-2.025,-1.587) (-2.061,-1.625) (-2.100,-1.668) (-2.150,-1.724) (-2.194,-1.775) (-2.236,-1.825) (-2.275,-1.874) (-2.315,-1.925) (-2.350,-1.972) (-2.388,-2.025) (-2.423,-2.075) (-2.456,-2.125) (-2.488,-2.175) (-2.519,-2.225) (-2.548,-2.275) (-2.575,-2.324) (-2.600,-2.370) (-2.625,-2.419) (-2.650,-2.470) (-2.675,-2.524) (-2.697,-2.575) (-2.718,-2.625) (-2.737,-2.675) (-2.756,-2.725) (-2.773,-2.775) (-2.789,-2.825) (-2.804,-2.875) (-2.818,-2.925) (-2.830,-2.975) (-2.842,-3.025) (-2.852,-3.075) (-2.862,-3.125) (-2.870,-3.175) (-2.877,-3.225) (-2.884,-3.275) (-2.889,-3.325) (-2.893,-3.375) (-2.896,-3.425) (-2.898,-3.475) (-2.899,-3.525) (-2.899,-3.575) (-2.898,-3.625) (-2.896,-3.675) (-2.893,-3.725) (-2.889,-3.775) (-2.884,-3.825) (-2.878,-3.875) (-2.871,-3.925) (-2.862,-3.975) (-2.853,-4.025) (-2.842,-4.075) (-2.831,-4.125) (-2.818,-4.175) (-2.804,-4.225) (-2.789,-4.275) (-2.773,-4.325) (-2.756,-4.375) (-2.738,-4.425) (-2.718,-4.475) (-2.697,-4.525) (-2.675,-4.575) (-2.652,-4.625) (-2.627,-4.675) (-2.601,-4.725) (-2.575,-4.772) (-2.550,-4.815) (-2.514,-4.875) (-2.482,-4.925) (-2.450,-4.972) (-2.413,-5.025) (-2.376,-5.075) (-2.337,-5.125) (-2.300,-5.170) (-2.253,-5.225) (-2.208,-5.275) (-2.161,-5.325) (-2.125,-5.362) (-2.086,-5.400) (-2.032,-5.450) (-1.976,-5.500) (-1.925,-5.543) (-1.885,-5.575) (-1.825,-5.621) (-1.775,-5.658) (-1.725,-5.693) (-1.678,-5.725) (-1.625,-5.759) (-1.575,-5.790) (-1.525,-5.819) (-1.475,-5.847) (-1.425,-5.874) (-1.375,-5.900) (-1.325,-5.924) (-1.275,-5.947) (-1.225,-5.970) (-1.175,-5.991) (-1.125,-6.011) (-1.075,-6.030) (-1.025,-6.048) (-0.975,-6.065) (-0.925,-6.081) (-0.875,-6.096) (-0.825,-6.111) (-0.775,-6.124) (-0.725,-6.137) (-0.675,-6.148) (-0.625,-6.159) (-0.575,-6.169) (-0.525,-6.178) (-0.475,-6.186) (-0.425,-6.193) (-0.375,-6.200) (-0.325,-6.205) (-0.275,-6.210) (-0.225,-6.214) (-0.175,-6.217) (-0.125,-6.220) (-0.075,-6.221) (-0.025,-6.222) (0.050,-6.222) (0.100,-6.221) (0.150,-6.219) (0.200,-6.216) (0.250,-6.212) (0.300,-6.208) (0.350,-6.203) (0.400,-6.196) (0.450,-6.190) (0.500,-6.182) (0.550,-6.173) (0.600,-6.164) (0.650,-6.154) (0.700,-6.142) (0.750,-6.130) (0.800,-6.117) (0.850,-6.104) (0.900,-6.089) (0.950,-6.073) (1.000,-6.057) (1.050,-6.039) (1.100,-6.020) (1.150,-6.001) (1.200,-5.980) (1.250,-5.959) (1.300,-5.936) (1.350,-5.912) (1.400,-5.887) (1.450,-5.861) (1.500,-5.833) (1.550,-5.805) (1.599,-5.775) (1.650,-5.743) (1.700,-5.710) (1.750,-5.676) (1.800,-5.640) (1.850,-5.602) (1.900,-5.563) (1.946,-5.525) (2.000,-5.479) (2.050,-5.434) (2.100,-5.386) (2.137,-5.350) (2.175,-5.311) (2.225,-5.257) (2.275,-5.200) (2.317,-5.150) (2.350,-5.108) (2.394,-5.050) (2.425,-5.008) (2.465,-4.950) (2.498,-4.900) (2.525,-4.857) (2.559,-4.800) (2.587,-4.750) (2.614,-4.700) (2.639,-4.650) (2.663,-4.600) (2.686,-4.550) (2.708,-4.500) (2.728,-4.450) (2.747,-4.400) (2.765,-4.350) (2.781,-4.300) (2.797,-4.250) (2.811,-4.200) (2.825,-4.150) (2.837,-4.100) (2.848,-4.050) (2.858,-4.000) (2.867,-3.950) (2.874,-3.900) (2.881,-3.850) (2.887,-3.800) (2.891,-3.750) (2.895,-3.700) (2.898,-3.650) (2.899,-3.600) (2.900,-3.550) (2.899,-3.500) (2.897,-3.450) (2.895,-3.400) (2.891,-3.350) (2.887,-3.300) (2.881,-3.250) (2.874,-3.200) (2.866,-3.150) (2.857,-3.100) (2.847,-3.050) (2.836,-3.000) (2.824,-2.950) (2.811,-2.900) (2.796,-2.850) (2.781,-2.800) (2.764,-2.750) (2.747,-2.700) (2.728,-2.650) (2.708,-2.600) (2.687,-2.550) (2.664,-2.500) (2.640,-2.450) (2.616,-2.400) (2.589,-2.350) (2.562,-2.300) (2.533,-2.250) (2.503,-2.200) (2.475,-2.155) (2.440,-2.100) (2.406,-2.050) (2.375,-2.007) (2.334,-1.950) (2.300,-1.906) (2.256,-1.850) (2.215,-1.800) (2.175,-1.752) (2.129,-1.700) (2.084,-1.650) (2.050,-1.613) (2.014,-1.575) (1.975,-1.535) (1.925,-1.485) (1.889,-1.450) (1.850,-1.413) (1.811,-1.375) (1.756,-1.325) (1.700,-1.275) (1.650,-1.231) (1.600,-1.188) (1.556,-1.150) (1.500,-1.104) (1.450,-1.063) (1.404,-1.025) (1.350,-0.983) (1.309,-0.950) (1.250,-0.904) (1.200,-0.866) (1.150,-0.828) (1.100,-0.790) (1.050,-0.753) (1.000,-0.716) (0.950,-0.679) (0.900,-0.642) (0.850,-0.606) (0.809,-0.575) (0.750,-0.533) (0.705,-0.500) (0.650,-0.461) (0.600,-0.425) (0.550,-0.390) (0.500,-0.354) (0.450,-0.319) (0.400,-0.284) (0.354,-0.250) (0.300,-0.213) (0.250,-0.177) (0.206,-0.144)}
\EQ{(0.206,0.144) (0.248,0.175) (0.300,0.213) (0.350,0.248) (0.400,0.284) (0.450,0.319) (0.495,0.350) (0.550,0.390) (0.600,0.425) (0.650,0.461) (0.700,0.497) (0.750,0.533) (0.800,0.569) (0.843,0.600) (0.900,0.642) (0.946,0.675) (1.000,0.716) (1.047,0.750) (1.100,0.790) (1.147,0.825) (1.200,0.866) (1.245,0.900) (1.300,0.943) (1.341,0.975) (1.400,1.022) (1.450,1.063) (1.496,1.100) (1.550,1.145) (1.600,1.188) (1.644,1.225) (1.700,1.275) (1.750,1.320) (1.800,1.366) (1.837,1.400) (1.875,1.436) (1.915,1.475) (1.965,1.525) (2.000,1.561) (2.038,1.600) (2.075,1.640) (2.125,1.695) (2.173,1.750) (2.215,1.800) (2.250,1.843) (2.296,1.900) (2.334,1.950) (2.370,2.000) (2.400,2.042) (2.440,2.100) (2.472,2.150) (2.500,2.195) (2.533,2.250) (2.562,2.300) (2.589,2.350) (2.616,2.400) (2.640,2.450) (2.664,2.500) (2.687,2.550) (2.708,2.600) (2.728,2.650) (2.747,2.700) (2.764,2.750) (2.781,2.800) (2.796,2.850) (2.811,2.900) (2.824,2.950) (2.836,3.000) (2.847,3.050) (2.857,3.100) (2.866,3.150) (2.874,3.200) (2.881,3.250) (2.887,3.300) (2.891,3.350) (2.895,3.400) (2.897,3.450) (2.899,3.500) (2.900,3.550) (2.899,3.600) (2.898,3.650) (2.895,3.700) (2.891,3.750) (2.887,3.800) (2.881,3.850) (2.874,3.900) (2.867,3.950) (2.858,4.000) (2.848,4.050) (2.837,4.100) (2.825,4.150) (2.811,4.200) (2.797,4.250) (2.781,4.300) (2.765,4.350) (2.747,4.400) (2.728,4.450) (2.708,4.500) (2.686,4.550) (2.663,4.600) (2.639,4.650) (2.614,4.700) (2.587,4.750) (2.559,4.800) (2.529,4.850) (2.500,4.897) (2.465,4.950) (2.431,5.000) (2.400,5.043) (2.357,5.100) (2.325,5.140) (2.275,5.200) (2.231,5.250) (2.185,5.300) (2.150,5.337) (2.112,5.375) (2.075,5.410) (2.032,5.450) (1.976,5.500) (1.925,5.543) (1.885,5.575) (1.825,5.621) (1.775,5.658) (1.725,5.693) (1.678,5.725) (1.625,5.759) (1.575,5.790) (1.525,5.819) (1.475,5.847) (1.425,5.874) (1.375,5.900) (1.325,5.924) (1.275,5.947) (1.225,5.970) (1.175,5.991) (1.125,6.011) (1.075,6.030) (1.025,6.048) (0.975,6.065) (0.925,6.081) (0.875,6.096) (0.825,6.111) (0.775,6.124) (0.725,6.137) (0.675,6.148) (0.625,6.159) (0.575,6.169) (0.525,6.178) (0.475,6.186) (0.425,6.193) (0.375,6.200) (0.325,6.205) (0.275,6.210) (0.225,6.214) (0.175,6.217) (0.125,6.220) (0.075,6.221) (0.025,6.222) (-0.050,6.222) (-0.100,6.221) (-0.150,6.219) (-0.200,6.216) (-0.250,6.212) (-0.300,6.208) (-0.350,6.203) (-0.400,6.196) (-0.450,6.190) (-0.500,6.182) (-0.550,6.173) (-0.600,6.164) (-0.650,6.154) (-0.700,6.142) (-0.750,6.130) (-0.800,6.117) (-0.850,6.104) (-0.900,6.089) (-0.950,6.073) (-1.000,6.057) (-1.050,6.039) (-1.100,6.020) (-1.150,6.001) (-1.200,5.980) (-1.250,5.959) (-1.300,5.936) (-1.350,5.912) (-1.400,5.887) (-1.450,5.861) (-1.500,5.833) (-1.550,5.805) (-1.599,5.775) (-1.650,5.743) (-1.700,5.710) (-1.750,5.676) (-1.800,5.640) (-1.850,5.602) (-1.900,5.563) (-1.946,5.525) (-2.000,5.479) (-2.050,5.434) (-2.100,5.386) (-2.137,5.350) (-2.175,5.311) (-2.225,5.257) (-2.275,5.200) (-2.317,5.150) (-2.350,5.108) (-2.394,5.050) (-2.425,5.008) (-2.465,4.950) (-2.498,4.900) (-2.525,4.857) (-2.559,4.800) (-2.587,4.750) (-2.614,4.700) (-2.639,4.650) (-2.663,4.600) (-2.686,4.550) (-2.708,4.500) (-2.728,4.450) (-2.747,4.400) (-2.765,4.350) (-2.781,4.300) (-2.797,4.250) (-2.811,4.200) (-2.825,4.150) (-2.837,4.100) (-2.848,4.050) (-2.858,4.000) (-2.867,3.950) (-2.874,3.900) (-2.881,3.850) (-2.887,3.800) (-2.891,3.750) (-2.895,3.700) (-2.898,3.650) (-2.899,3.600) (-2.900,3.550) (-2.899,3.500) (-2.897,3.450) (-2.895,3.400) (-2.891,3.350) (-2.887,3.300) (-2.881,3.250) (-2.874,3.200) (-2.866,3.150) (-2.857,3.100) (-2.847,3.050) (-2.836,3.000) (-2.824,2.950) (-2.811,2.900) (-2.796,2.850) (-2.781,2.800) (-2.764,2.750) (-2.747,2.700) (-2.728,2.650) (-2.708,2.600) (-2.687,2.550) (-2.664,2.500) (-2.640,2.450) (-2.616,2.400) (-2.589,2.350) (-2.562,2.300) (-2.533,2.250) (-2.503,2.200) (-2.475,2.155) (-2.440,2.100) (-2.406,2.050) (-2.375,2.007) (-2.334,1.950) (-2.300,1.906) (-2.256,1.850) (-2.215,1.800) (-2.175,1.752) (-2.129,1.700) (-2.084,1.650) (-2.050,1.613) (-2.014,1.575) (-1.975,1.535) (-1.925,1.485) (-1.889,1.450) (-1.850,1.413) (-1.811,1.375) (-1.756,1.325) (-1.700,1.275) (-1.650,1.231) (-1.600,1.188) (-1.556,1.150) (-1.500,1.104) (-1.450,1.063) (-1.404,1.025) (-1.350,0.983) (-1.309,0.950) (-1.250,0.904) (-1.200,0.866) (-1.150,0.828) (-1.100,0.790) (-1.050,0.753) (-1.000,0.716) (-0.950,0.679) (-0.900,0.642) (-0.850,0.606) (-0.809,0.575) (-0.750,0.533) (-0.705,0.500) (-0.650,0.461) (-0.600,0.425) (-0.550,0.390) (-0.500,0.354) (-0.450,0.319) (-0.400,0.284) (-0.354,0.250) (-0.300,0.213) (-0.250,0.177) (-0.220,0.155)}
\EQ{(-0.220,-0.155) (-0.275,-0.195) (-0.319,-0.225) (-0.375,-0.266) (-0.424,-0.300) (-0.475,-0.337) (-0.525,-0.372) (-0.575,-0.408) (-0.625,-0.444) (-0.670,-0.475) (-0.725,-0.516) (-0.773,-0.550) (-0.825,-0.588) (-0.875,-0.624) (-0.925,-0.661) (-0.975,-0.698) (-1.025,-0.735) (-1.075,-0.773) (-1.125,-0.811) (-1.175,-0.849) (-1.225,-0.888) (-1.273,-0.925) (-1.325,-0.967) (-1.367,-1.000) (-1.425,-1.048) (-1.475,-1.089) (-1.518,-1.125) (-1.575,-1.174) (-1.625,-1.218) (-1.675,-1.263) (-1.715,-1.300) (-1.769,-1.350) (-1.821,-1.400) (-1.871,-1.450) (-1.920,-1.500) (-1.967,-1.550) (-2.013,-1.600) (-2.050,-1.641) (-2.100,-1.699) (-2.142,-1.750) (-2.175,-1.791) (-2.221,-1.850) (-2.258,-1.900) (-2.293,-1.950) (-2.325,-1.996) (-2.361,-2.050) (-2.392,-2.100) (-2.422,-2.150) (-2.450,-2.199) (-2.475,-2.245) (-2.500,-2.292) (-2.525,-2.343) (-2.550,-2.396) (-2.574,-2.450) (-2.595,-2.500) (-2.615,-2.550) (-2.633,-2.600) (-2.650,-2.648) (-2.667,-2.700) (-2.682,-2.750) (-2.696,-2.800) (-2.708,-2.850) (-2.720,-2.900) (-2.730,-2.950) (-2.739,-3.000) (-2.747,-3.050) (-2.754,-3.100) (-2.760,-3.150) (-2.765,-3.200) (-2.769,-3.250) (-2.771,-3.300) (-2.773,-3.350) (-2.773,-3.400) (-2.773,-3.450) (-2.771,-3.500) (-2.768,-3.550) (-2.764,-3.600) (-2.759,-3.650) (-2.753,-3.700) (-2.746,-3.750) (-2.738,-3.800) (-2.729,-3.850) (-2.718,-3.900) (-2.706,-3.950) (-2.694,-4.000) (-2.680,-4.050) (-2.664,-4.100) (-2.648,-4.150) (-2.630,-4.200) (-2.612,-4.250) (-2.592,-4.300) (-2.570,-4.350) (-2.547,-4.400) (-2.525,-4.446) (-2.500,-4.496) (-2.475,-4.542) (-2.450,-4.587) (-2.412,-4.650) (-2.381,-4.700) (-2.350,-4.746) (-2.312,-4.800) (-2.275,-4.850) (-2.237,-4.900) (-2.200,-4.945) (-2.153,-5.000) (-2.108,-5.050) (-2.061,-5.100) (-2.025,-5.136) (-1.984,-5.175) (-1.930,-5.225) (-1.875,-5.273) (-1.825,-5.315) (-1.781,-5.350) (-1.725,-5.392) (-1.679,-5.425) (-1.625,-5.462) (-1.575,-5.495) (-1.527,-5.525) (-1.475,-5.556) (-1.425,-5.585) (-1.375,-5.612) (-1.325,-5.638) (-1.275,-5.662) (-1.225,-5.686) (-1.175,-5.708) (-1.125,-5.729) (-1.075,-5.749) (-1.025,-5.769) (-0.975,-5.787) (-0.925,-5.804) (-0.875,-5.819) (-0.825,-5.834) (-0.775,-5.848) (-0.725,-5.862) (-0.675,-5.874) (-0.625,-5.885) (-0.575,-5.895) (-0.525,-5.905) (-0.475,-5.913) (-0.425,-5.921) (-0.375,-5.928) (-0.325,-5.934) (-0.275,-5.939) (-0.225,-5.943) (-0.175,-5.946) (-0.125,-5.949) (-0.075,-5.950) (-0.025,-5.951) (0.050,-5.951) (0.100,-5.950) (0.150,-5.948) (0.200,-5.945) (0.250,-5.941) (0.300,-5.936) (0.350,-5.931) (0.400,-5.924) (0.450,-5.917) (0.500,-5.909) (0.550,-5.900) (0.600,-5.890) (0.650,-5.879) (0.700,-5.868) (0.750,-5.855) (0.800,-5.842) (0.850,-5.827) (0.900,-5.812) (0.950,-5.795) (1.000,-5.778) (1.050,-5.759) (1.100,-5.740) (1.150,-5.719) (1.200,-5.697) (1.249,-5.675) (1.300,-5.650) (1.350,-5.625) (1.397,-5.600) (1.442,-5.575) (1.486,-5.550) (1.550,-5.511) (1.600,-5.479) (1.643,-5.450) (1.700,-5.410) (1.748,-5.375) (1.800,-5.335) (1.843,-5.300) (1.900,-5.252) (1.950,-5.207) (2.000,-5.160) (2.036,-5.125) (2.075,-5.085) (2.125,-5.031) (2.175,-4.975) (2.217,-4.925) (2.250,-4.883) (2.294,-4.825) (2.325,-4.782) (2.364,-4.725) (2.397,-4.675) (2.425,-4.629) (2.457,-4.575) (2.484,-4.525) (2.511,-4.475) (2.535,-4.425) (2.559,-4.375) (2.581,-4.325) (2.602,-4.275) (2.621,-4.225) (2.639,-4.175) (2.656,-4.125) (2.672,-4.075) (2.687,-4.025) (2.700,-3.975) (2.712,-3.925) (2.723,-3.875) (2.733,-3.825) (2.742,-3.775) (2.750,-3.725) (2.757,-3.675) (2.762,-3.625) (2.766,-3.575) (2.770,-3.525) (2.772,-3.475) (2.773,-3.425) (2.773,-3.375) (2.772,-3.325) (2.770,-3.275) (2.767,-3.225) (2.763,-3.175) (2.757,-3.125) (2.751,-3.075) (2.743,-3.025) (2.735,-2.975) (2.725,-2.926) (2.714,-2.875) (2.702,-2.825) (2.689,-2.775) (2.674,-2.725) (2.659,-2.675) (2.642,-2.625) (2.625,-2.577) (2.605,-2.525) (2.585,-2.475) (2.563,-2.425) (2.541,-2.375) (2.517,-2.325) (2.491,-2.275) (2.464,-2.225) (2.436,-2.175) (2.407,-2.125) (2.376,-2.075) (2.344,-2.025) (2.311,-1.975) (2.276,-1.925) (2.240,-1.875) (2.202,-1.825) (2.162,-1.775) (2.125,-1.729) (2.079,-1.675) (2.036,-1.625) (2.000,-1.586) (1.950,-1.532) (1.900,-1.480) (1.850,-1.429) (1.800,-1.380) (1.750,-1.333) (1.700,-1.286) (1.661,-1.250) (1.604,-1.200) (1.550,-1.153) (1.500,-1.110) (1.458,-1.075) (1.400,-1.027) (1.350,-0.987) (1.304,-0.950) (1.250,-0.907) (1.209,-0.875) (1.150,-0.830) (1.100,-0.792) (1.050,-0.754) (1.000,-0.717) (0.950,-0.680) (0.900,-0.643) (0.850,-0.606) (0.808,-0.575) (0.750,-0.534) (0.704,-0.500) (0.650,-0.462) (0.600,-0.425) (0.550,-0.390) (0.500,-0.354) (0.450,-0.319) (0.400,-0.284) (0.354,-0.250) (0.300,-0.213) (0.250,-0.177) (0.206,-0.144)}
\EQ{(0.206,0.144) (0.248,0.175) (0.300,0.213) (0.350,0.248) (0.400,0.284) (0.450,0.319) (0.495,0.350) (0.550,0.390) (0.599,0.425) (0.650,0.462) (0.700,0.497) (0.750,0.534) (0.800,0.570) (0.842,0.600) (0.900,0.643) (0.944,0.675) (1.000,0.717) (1.045,0.750) (1.100,0.792) (1.144,0.825) (1.200,0.868) (1.241,0.900) (1.300,0.947) (1.350,0.987) (1.398,1.025) (1.450,1.068) (1.500,1.110) (1.547,1.150) (1.600,1.196) (1.650,1.241) (1.688,1.275) (1.725,1.309) (1.769,1.350) (1.821,1.400) (1.871,1.450) (1.920,1.500) (1.967,1.550) (2.013,1.600) (2.050,1.641) (2.100,1.699) (2.142,1.750) (2.175,1.791) (2.221,1.850) (2.258,1.900) (2.293,1.950) (2.325,1.996) (2.361,2.050) (2.392,2.100) (2.422,2.150) (2.450,2.199) (2.475,2.245) (2.500,2.292) (2.525,2.343) (2.550,2.396) (2.574,2.450) (2.595,2.500) (2.615,2.550) (2.633,2.600) (2.650,2.648) (2.667,2.700) (2.682,2.750) (2.696,2.800) (2.708,2.850) (2.720,2.900) (2.730,2.950) (2.739,3.000) (2.747,3.050) (2.754,3.100) (2.760,3.150) (2.765,3.200) (2.769,3.250) (2.771,3.300) (2.773,3.350) (2.773,3.400) (2.773,3.450) (2.771,3.500) (2.768,3.550) (2.764,3.600) (2.759,3.650) (2.753,3.700) (2.746,3.750) (2.738,3.800) (2.729,3.850) (2.718,3.900) (2.706,3.950) (2.694,4.000) (2.680,4.050) (2.664,4.100) (2.648,4.150) (2.630,4.200) (2.612,4.250) (2.592,4.300) (2.570,4.350) (2.547,4.400) (2.525,4.446) (2.500,4.496) (2.475,4.542) (2.450,4.587) (2.412,4.650) (2.381,4.700) (2.350,4.746) (2.312,4.800) (2.275,4.850) (2.237,4.900) (2.200,4.945) (2.153,5.000) (2.108,5.050) (2.061,5.100) (2.025,5.136) (1.984,5.175) (1.930,5.225) (1.875,5.273) (1.825,5.315) (1.781,5.350) (1.725,5.392) (1.679,5.425) (1.625,5.462) (1.575,5.495) (1.527,5.525) (1.475,5.556) (1.425,5.585) (1.375,5.612) (1.325,5.638) (1.275,5.662) (1.225,5.686) (1.175,5.708) (1.125,5.729) (1.075,5.749) (1.025,5.769) (0.975,5.787) (0.925,5.804) (0.875,5.819) (0.825,5.834) (0.775,5.848) (0.725,5.862) (0.675,5.874) (0.625,5.885) (0.575,5.895) (0.525,5.905) (0.475,5.913) (0.425,5.921) (0.375,5.928) (0.325,5.934) (0.275,5.939) (0.225,5.943) (0.175,5.946) (0.125,5.949) (0.075,5.950) (0.025,5.951) (-0.050,5.951) (-0.100,5.950) (-0.150,5.948) (-0.200,5.945) (-0.250,5.941) (-0.300,5.936) (-0.350,5.931) (-0.400,5.924) (-0.450,5.917) (-0.500,5.909) (-0.550,5.900) (-0.600,5.890) (-0.650,5.879) (-0.700,5.868) (-0.750,5.855) (-0.800,5.842) (-0.850,5.827) (-0.900,5.812) (-0.950,5.795) (-1.000,5.778) (-1.050,5.759) (-1.100,5.740) (-1.150,5.719) (-1.200,5.697) (-1.249,5.675) (-1.300,5.650) (-1.350,5.625) (-1.397,5.600) (-1.442,5.575) (-1.486,5.550) (-1.550,5.511) (-1.600,5.479) (-1.643,5.450) (-1.700,5.410) (-1.748,5.375) (-1.800,5.335) (-1.843,5.300) (-1.900,5.252) (-1.950,5.207) (-2.000,5.160) (-2.036,5.125) (-2.075,5.085) (-2.125,5.031) (-2.175,4.975) (-2.217,4.925) (-2.250,4.883) (-2.294,4.825) (-2.325,4.782) (-2.364,4.725) (-2.397,4.675) (-2.425,4.629) (-2.457,4.575) (-2.484,4.525) (-2.511,4.475) (-2.535,4.425) (-2.559,4.375) (-2.581,4.325) (-2.602,4.275) (-2.621,4.225) (-2.639,4.175) (-2.656,4.125) (-2.672,4.075) (-2.687,4.025) (-2.700,3.975) (-2.712,3.925) (-2.723,3.875) (-2.733,3.825) (-2.742,3.775) (-2.750,3.725) (-2.757,3.675) (-2.762,3.625) (-2.766,3.575) (-2.770,3.525) (-2.772,3.475) (-2.773,3.425) (-2.773,3.375) (-2.772,3.325) (-2.770,3.275) (-2.767,3.225) (-2.763,3.175) (-2.757,3.125) (-2.751,3.075) (-2.743,3.025) (-2.735,2.975) (-2.725,2.926) (-2.714,2.875) (-2.702,2.825) (-2.689,2.775) (-2.674,2.725) (-2.659,2.675) (-2.642,2.625) (-2.625,2.577) (-2.605,2.525) (-2.585,2.475) (-2.563,2.425) (-2.541,2.375) (-2.517,2.325) (-2.491,2.275) (-2.464,2.225) (-2.436,2.175) (-2.407,2.125) (-2.376,2.075) (-2.344,2.025) (-2.311,1.975) (-2.276,1.925) (-2.240,1.875) (-2.202,1.825) (-2.162,1.775) (-2.125,1.729) (-2.079,1.675) (-2.036,1.625) (-2.000,1.586) (-1.950,1.532) (-1.900,1.480) (-1.850,1.429) (-1.800,1.380) (-1.750,1.333) (-1.700,1.286) (-1.661,1.250) (-1.604,1.200) (-1.550,1.153) (-1.500,1.110) (-1.458,1.075) (-1.400,1.027) (-1.350,0.987) (-1.304,0.950) (-1.250,0.907) (-1.209,0.875) (-1.150,0.830) (-1.100,0.792) (-1.050,0.754) (-1.000,0.717) (-0.950,0.680) (-0.900,0.643) (-0.850,0.606) (-0.808,0.575) (-0.750,0.534) (-0.704,0.500) (-0.650,0.462) (-0.600,0.425) (-0.550,0.390) (-0.500,0.354) (-0.450,0.319) (-0.400,0.284) (-0.354,0.250) (-0.300,0.213) (-0.250,0.177) (-0.220,0.155)}
\EQ{(-0.220,-0.155) (-0.275,-0.195) (-0.319,-0.225) (-0.375,-0.266) (-0.424,-0.300) (-0.475,-0.337) (-0.525,-0.372) (-0.575,-0.408) (-0.625,-0.444) (-0.669,-0.475) (-0.725,-0.516) (-0.773,-0.550) (-0.825,-0.589) (-0.875,-0.625) (-0.925,-0.662) (-0.975,-0.699) (-1.025,-0.737) (-1.075,-0.774) (-1.125,-0.813) (-1.173,-0.850) (-1.225,-0.891) (-1.269,-0.925) (-1.325,-0.970) (-1.375,-1.011) (-1.422,-1.050) (-1.475,-1.095) (-1.525,-1.139) (-1.567,-1.175) (-1.622,-1.225) (-1.675,-1.274) (-1.725,-1.322) (-1.775,-1.370) (-1.825,-1.421) (-1.875,-1.473) (-1.924,-1.525) (-1.969,-1.575) (-2.012,-1.625) (-2.050,-1.671) (-2.094,-1.725) (-2.125,-1.766) (-2.170,-1.825) (-2.200,-1.868) (-2.239,-1.925) (-2.272,-1.975) (-2.300,-2.020) (-2.333,-2.075) (-2.362,-2.125) (-2.389,-2.175) (-2.414,-2.225) (-2.439,-2.275) (-2.462,-2.325) (-2.483,-2.375) (-2.504,-2.425) (-2.523,-2.475) (-2.541,-2.525) (-2.558,-2.575) (-2.573,-2.625) (-2.587,-2.675) (-2.600,-2.725) (-2.612,-2.775) (-2.623,-2.825) (-2.632,-2.875) (-2.640,-2.925) (-2.647,-2.975) (-2.654,-3.025) (-2.658,-3.075) (-2.662,-3.125) (-2.665,-3.175) (-2.666,-3.225) (-2.667,-3.275) (-2.666,-3.325) (-2.664,-3.375) (-2.661,-3.425) (-2.657,-3.475) (-2.651,-3.525) (-2.645,-3.575) (-2.637,-3.625) (-2.629,-3.675) (-2.619,-3.725) (-2.608,-3.775) (-2.595,-3.825) (-2.582,-3.875) (-2.567,-3.925) (-2.551,-3.975) (-2.534,-4.025) (-2.515,-4.075) (-2.496,-4.125) (-2.475,-4.174) (-2.452,-4.225) (-2.428,-4.275) (-2.402,-4.325) (-2.375,-4.375) (-2.350,-4.420) (-2.317,-4.475) (-2.285,-4.525) (-2.251,-4.575) (-2.216,-4.625) (-2.179,-4.675) (-2.140,-4.725) (-2.100,-4.773) (-2.055,-4.825) (-2.009,-4.875) (-1.960,-4.925) (-1.909,-4.975) (-1.855,-5.025) (-1.800,-5.073) (-1.750,-5.114) (-1.704,-5.150) (-1.650,-5.191) (-1.602,-5.225) (-1.550,-5.260) (-1.500,-5.293) (-1.450,-5.323) (-1.405,-5.350) (-1.361,-5.375) (-1.314,-5.400) (-1.265,-5.425) (-1.214,-5.450) (-1.160,-5.475) (-1.103,-5.500) (-1.050,-5.522) (-1.000,-5.541) (-0.950,-5.559) (-0.900,-5.577) (-0.850,-5.593) (-0.800,-5.608) (-0.750,-5.622) (-0.700,-5.635) (-0.650,-5.647) (-0.600,-5.659) (-0.550,-5.669) (-0.500,-5.678) (-0.450,-5.687) (-0.400,-5.694) (-0.350,-5.701) (-0.300,-5.707) (-0.250,-5.711) (-0.200,-5.715) (-0.150,-5.718) (-0.100,-5.721) (-0.050,-5.722) (0.000,-5.722) (0.050,-5.722) (0.100,-5.721) (0.150,-5.718) (0.200,-5.715) (0.250,-5.711) (0.300,-5.707) (0.350,-5.701) (0.400,-5.694) (0.450,-5.687) (0.500,-5.678) (0.550,-5.669) (0.600,-5.659) (0.650,-5.647) (0.700,-5.635) (0.750,-5.622) (0.800,-5.608) (0.850,-5.593) (0.900,-5.577) (0.950,-5.559) (1.000,-5.541) (1.050,-5.522) (1.100,-5.501) (1.150,-5.480) (1.200,-5.457) (1.250,-5.433) (1.300,-5.407) (1.350,-5.381) (1.400,-5.353) (1.447,-5.325) (1.500,-5.293) (1.550,-5.260) (1.600,-5.226) (1.650,-5.191) (1.700,-5.153) (1.750,-5.114) (1.797,-5.075) (1.850,-5.029) (1.900,-4.983) (1.950,-4.935) (2.000,-4.884) (2.050,-4.830) (2.098,-4.775) (2.140,-4.725) (2.175,-4.680) (2.216,-4.625) (2.250,-4.577) (2.285,-4.525) (2.317,-4.475) (2.347,-4.425) (2.375,-4.376) (2.400,-4.330) (2.425,-4.281) (2.450,-4.229) (2.474,-4.175) (2.496,-4.125) (2.515,-4.075) (2.534,-4.025) (2.551,-3.975) (2.567,-3.925) (2.582,-3.875) (2.595,-3.825) (2.608,-3.775) (2.619,-3.725) (2.629,-3.675) (2.637,-3.625) (2.645,-3.575) (2.651,-3.525) (2.657,-3.475) (2.661,-3.425) (2.664,-3.375) (2.666,-3.325) (2.667,-3.275) (2.666,-3.225) (2.665,-3.175) (2.662,-3.125) (2.658,-3.075) (2.654,-3.025) (2.647,-2.975) (2.640,-2.925) (2.632,-2.875) (2.623,-2.825) (2.612,-2.775) (2.600,-2.725) (2.587,-2.675) (2.573,-2.625) (2.558,-2.575) (2.541,-2.525) (2.523,-2.475) (2.504,-2.425) (2.483,-2.375) (2.462,-2.325) (2.439,-2.275) (2.414,-2.225) (2.389,-2.175) (2.362,-2.125) (2.333,-2.075) (2.303,-2.025) (2.275,-1.980) (2.239,-1.925) (2.205,-1.875) (2.175,-1.833) (2.132,-1.775) (2.100,-1.733) (2.053,-1.675) (2.012,-1.625) (1.975,-1.583) (1.925,-1.527) (1.877,-1.475) (1.829,-1.425) (1.780,-1.375) (1.729,-1.325) (1.676,-1.275) (1.625,-1.228) (1.575,-1.183) (1.525,-1.139) (1.481,-1.100) (1.425,-1.053) (1.375,-1.011) (1.331,-0.975) (1.275,-0.930) (1.225,-0.891) (1.175,-0.851) (1.125,-0.813) (1.076,-0.775) (1.025,-0.737) (0.976,-0.700) (0.925,-0.662) (0.875,-0.625) (0.825,-0.589) (0.775,-0.552) (0.725,-0.516) (0.675,-0.480) (0.635,-0.450) (0.575,-0.408) (0.530,-0.375) (0.475,-0.337) (0.425,-0.301) (0.375,-0.266) (0.325,-0.230) (0.284,-0.200) (0.225,-0.160) (0.206,-0.144)}
\EQ{(0.206,0.144) (0.248,0.175) (0.300,0.213) (0.350,0.248) (0.400,0.284) (0.450,0.319) (0.495,0.350) (0.550,0.390) (0.599,0.425) (0.650,0.462) (0.700,0.498) (0.750,0.534) (0.800,0.570) (0.842,0.600) (0.900,0.644) (0.943,0.675) (1.000,0.718) (1.043,0.750) (1.100,0.794) (1.141,0.825) (1.200,0.871) (1.250,0.910) (1.300,0.950) (1.350,0.991) (1.392,1.025) (1.450,1.074) (1.500,1.117) (1.538,1.150) (1.595,1.200) (1.649,1.250) (1.700,1.298) (1.750,1.346) (1.800,1.395) (1.850,1.447) (1.900,1.499) (1.946,1.550) (1.990,1.600) (2.025,1.641) (2.074,1.700) (2.113,1.750) (2.150,1.799) (2.188,1.850) (2.222,1.900) (2.256,1.950) (2.288,2.000) (2.318,2.050) (2.348,2.100) (2.375,2.149) (2.400,2.197) (2.425,2.247) (2.450,2.299) (2.473,2.350) (2.494,2.400) (2.514,2.450) (2.532,2.500) (2.549,2.550) (2.565,2.600) (2.580,2.650) (2.594,2.700) (2.606,2.750) (2.617,2.800) (2.627,2.850) (2.636,2.900) (2.644,2.950) (2.651,3.000) (2.656,3.050) (2.660,3.100) (2.664,3.150) (2.666,3.200) (2.667,3.250) (2.666,3.300) (2.665,3.350) (2.663,3.400) (2.659,3.450) (2.654,3.500) (2.648,3.550) (2.641,3.600) (2.633,3.650) (2.624,3.700) (2.613,3.750) (2.602,3.800) (2.589,3.850) (2.575,3.899) (2.559,3.950) (2.543,4.000) (2.525,4.050) (2.506,4.100) (2.485,4.150) (2.463,4.200) (2.440,4.250) (2.415,4.300) (2.389,4.350) (2.361,4.400) (2.332,4.450) (2.301,4.500) (2.268,4.550) (2.234,4.600) (2.200,4.647) (2.160,4.700) (2.125,4.743) (2.077,4.800) (2.032,4.850) (1.985,4.900) (1.935,4.950) (1.882,5.000) (1.826,5.050) (1.775,5.094) (1.725,5.134) (1.675,5.172) (1.625,5.209) (1.575,5.244) (1.528,5.275) (1.475,5.308) (1.425,5.338) (1.375,5.367) (1.325,5.394) (1.275,5.420) (1.225,5.445) (1.175,5.468) (1.125,5.491) (1.075,5.512) (1.025,5.532) (0.976,5.550) (0.925,5.568) (0.875,5.585) (0.826,5.600) (0.775,5.615) (0.725,5.629) (0.675,5.641) (0.625,5.653) (0.575,5.664) (0.525,5.674) (0.475,5.683) (0.425,5.691) (0.375,5.698) (0.325,5.704) (0.275,5.709) (0.225,5.713) (0.175,5.717) (0.125,5.720) (0.075,5.721) (0.025,5.722) (-0.050,5.722) (-0.100,5.721) (-0.150,5.718) (-0.200,5.715) (-0.250,5.711) (-0.300,5.707) (-0.350,5.701) (-0.400,5.694) (-0.450,5.687) (-0.500,5.678) (-0.550,5.669) (-0.600,5.659) (-0.650,5.647) (-0.700,5.635) (-0.750,5.622) (-0.800,5.608) (-0.850,5.593) (-0.900,5.577) (-0.950,5.559) (-1.000,5.541) (-1.050,5.522) (-1.100,5.501) (-1.150,5.480) (-1.200,5.457) (-1.250,5.433) (-1.300,5.407) (-1.350,5.381) (-1.400,5.353) (-1.447,5.325) (-1.500,5.293) (-1.550,5.260) (-1.600,5.226) (-1.650,5.191) (-1.700,5.153) (-1.750,5.114) (-1.797,5.075) (-1.850,5.029) (-1.900,4.983) (-1.950,4.935) (-2.000,4.884) (-2.050,4.830) (-2.098,4.775) (-2.140,4.725) (-2.175,4.680) (-2.216,4.625) (-2.250,4.577) (-2.285,4.525) (-2.317,4.475) (-2.347,4.425) (-2.375,4.376) (-2.400,4.330) (-2.425,4.281) (-2.450,4.229) (-2.474,4.175) (-2.496,4.125) (-2.515,4.075) (-2.534,4.025) (-2.551,3.975) (-2.567,3.925) (-2.582,3.875) (-2.595,3.825) (-2.608,3.775) (-2.619,3.725) (-2.629,3.675) (-2.637,3.625) (-2.645,3.575) (-2.651,3.525) (-2.657,3.475) (-2.661,3.425) (-2.664,3.375) (-2.666,3.325) (-2.667,3.275) (-2.666,3.225) (-2.665,3.175) (-2.662,3.125) (-2.658,3.075) (-2.654,3.025) (-2.647,2.975) (-2.640,2.925) (-2.632,2.875) (-2.623,2.825) (-2.612,2.775) (-2.600,2.725) (-2.587,2.675) (-2.573,2.625) (-2.558,2.575) (-2.541,2.525) (-2.523,2.475) (-2.504,2.425) (-2.483,2.375) (-2.462,2.325) (-2.439,2.275) (-2.414,2.225) (-2.389,2.175) (-2.362,2.125) (-2.333,2.075) (-2.303,2.025) (-2.275,1.980) (-2.239,1.925) (-2.205,1.875) (-2.175,1.833) (-2.132,1.775) (-2.100,1.733) (-2.053,1.675) (-2.012,1.625) (-1.975,1.583) (-1.925,1.527) (-1.877,1.475) (-1.829,1.425) (-1.780,1.375) (-1.729,1.325) (-1.676,1.275) (-1.625,1.228) (-1.575,1.183) (-1.525,1.139) (-1.481,1.100) (-1.425,1.053) (-1.375,1.011) (-1.331,0.975) (-1.275,0.930) (-1.225,0.891) (-1.175,0.851) (-1.125,0.813) (-1.076,0.775) (-1.025,0.737) (-0.976,0.700) (-0.925,0.662) (-0.875,0.625) (-0.825,0.589) (-0.775,0.552) (-0.725,0.516) (-0.675,0.480) (-0.635,0.450) (-0.575,0.408) (-0.530,0.375) (-0.475,0.337) (-0.425,0.301) (-0.375,0.266) (-0.325,0.230) (-0.284,0.200) (-0.225,0.160) (-0.220,0.155)}
\EQ{(-0.220,-0.155) (-0.275,-0.195) (-0.319,-0.225) (-0.375,-0.266) (-0.424,-0.300) (-0.475,-0.337) (-0.525,-0.372) (-0.575,-0.408) (-0.625,-0.444) (-0.669,-0.475) (-0.725,-0.516) (-0.772,-0.550) (-0.825,-0.589) (-0.874,-0.625) (-0.925,-0.663) (-0.975,-0.700) (-1.025,-0.738) (-1.074,-0.775) (-1.125,-0.815) (-1.171,-0.850) (-1.225,-0.893) (-1.265,-0.925) (-1.325,-0.974) (-1.375,-1.016) (-1.416,-1.050) (-1.473,-1.100) (-1.525,-1.146) (-1.575,-1.191) (-1.612,-1.225) (-1.650,-1.261) (-1.691,-1.300) (-1.741,-1.350) (-1.790,-1.400) (-1.825,-1.438) (-1.860,-1.475) (-1.900,-1.520) (-1.947,-1.575) (-1.989,-1.625) (-2.025,-1.670) (-2.067,-1.725) (-2.100,-1.770) (-2.139,-1.825) (-2.173,-1.875) (-2.205,-1.925) (-2.236,-1.975) (-2.265,-2.025) (-2.293,-2.075) (-2.319,-2.125) (-2.344,-2.175) (-2.368,-2.225) (-2.390,-2.275) (-2.411,-2.325) (-2.431,-2.375) (-2.449,-2.425) (-2.466,-2.475) (-2.481,-2.525) (-2.496,-2.575) (-2.509,-2.625) (-2.521,-2.675) (-2.532,-2.725) (-2.541,-2.775) (-2.549,-2.825) (-2.557,-2.875) (-2.562,-2.925) (-2.567,-2.975) (-2.571,-3.025) (-2.573,-3.075) (-2.574,-3.125) (-2.574,-3.175) (-2.573,-3.225) (-2.571,-3.275) (-2.568,-3.325) (-2.563,-3.375) (-2.557,-3.425) (-2.550,-3.475) (-2.542,-3.525) (-2.533,-3.575) (-2.522,-3.625) (-2.510,-3.675) (-2.497,-3.725) (-2.483,-3.775) (-2.467,-3.825) (-2.450,-3.875) (-2.432,-3.925) (-2.413,-3.975) (-2.392,-4.025) (-2.369,-4.075) (-2.345,-4.125) (-2.320,-4.175) (-2.293,-4.225) (-2.264,-4.275) (-2.234,-4.325) (-2.202,-4.375) (-2.168,-4.425) (-2.133,-4.475) (-2.100,-4.518) (-2.055,-4.575) (-2.013,-4.625) (-1.975,-4.668) (-1.925,-4.722) (-1.875,-4.772) (-1.825,-4.820) (-1.775,-4.866) (-1.735,-4.900) (-1.675,-4.949) (-1.625,-4.988) (-1.575,-5.025) (-1.525,-5.060) (-1.475,-5.093) (-1.425,-5.124) (-1.375,-5.155) (-1.325,-5.183) (-1.275,-5.210) (-1.225,-5.236) (-1.175,-5.261) (-1.125,-5.284) (-1.075,-5.306) (-1.025,-5.327) (-0.975,-5.346) (-0.925,-5.365) (-0.875,-5.382) (-0.825,-5.398) (-0.775,-5.414) (-0.725,-5.428) (-0.675,-5.441) (-0.625,-5.453) (-0.575,-5.464) (-0.525,-5.474) (-0.475,-5.484) (-0.425,-5.492) (-0.375,-5.499) (-0.325,-5.506) (-0.275,-5.511) (-0.225,-5.516) (-0.175,-5.519) (-0.125,-5.522) (-0.075,-5.524) (-0.025,-5.525) (0.050,-5.524) (0.100,-5.523) (0.150,-5.521) (0.200,-5.518) (0.250,-5.513) (0.300,-5.508) (0.350,-5.503) (0.400,-5.496) (0.450,-5.488) (0.500,-5.479) (0.550,-5.469) (0.600,-5.459) (0.650,-5.447) (0.700,-5.434) (0.750,-5.421) (0.800,-5.406) (0.850,-5.390) (0.900,-5.374) (0.950,-5.356) (1.000,-5.337) (1.050,-5.316) (1.100,-5.295) (1.150,-5.272) (1.197,-5.250) (1.247,-5.225) (1.294,-5.200) (1.340,-5.175) (1.400,-5.140) (1.450,-5.109) (1.500,-5.076) (1.550,-5.042) (1.600,-5.006) (1.642,-4.975) (1.700,-4.929) (1.750,-4.887) (1.792,-4.850) (1.846,-4.800) (1.898,-4.750) (1.946,-4.700) (1.991,-4.650) (2.025,-4.611) (2.075,-4.550) (2.114,-4.500) (2.150,-4.451) (2.185,-4.400) (2.218,-4.350) (2.249,-4.300) (2.275,-4.257) (2.300,-4.212) (2.325,-4.165) (2.350,-4.115) (2.375,-4.062) (2.400,-4.005) (2.423,-3.950) (2.442,-3.900) (2.459,-3.850) (2.475,-3.801) (2.490,-3.750) (2.504,-3.700) (2.516,-3.650) (2.528,-3.600) (2.538,-3.550) (2.546,-3.500) (2.554,-3.450) (2.560,-3.400) (2.566,-3.350) (2.570,-3.300) (2.572,-3.250) (2.574,-3.200) (2.575,-3.150) (2.574,-3.100) (2.572,-3.050) (2.569,-3.000) (2.565,-2.950) (2.560,-2.900) (2.553,-2.850) (2.545,-2.800) (2.537,-2.750) (2.526,-2.700) (2.515,-2.650) (2.502,-2.600) (2.489,-2.550) (2.474,-2.500) (2.457,-2.450) (2.440,-2.400) (2.421,-2.350) (2.401,-2.300) (2.379,-2.250) (2.356,-2.200) (2.332,-2.150) (2.306,-2.100) (2.279,-2.050) (2.251,-2.000) (2.221,-1.950) (2.189,-1.900) (2.156,-1.850) (2.125,-1.805) (2.086,-1.750) (2.050,-1.702) (2.009,-1.650) (1.975,-1.608) (1.926,-1.550) (1.882,-1.500) (1.837,-1.450) (1.800,-1.411) (1.765,-1.375) (1.725,-1.334) (1.675,-1.285) (1.639,-1.250) (1.600,-1.214) (1.558,-1.175) (1.502,-1.125) (1.450,-1.080) (1.400,-1.037) (1.356,-1.000) (1.300,-0.954) (1.250,-0.913) (1.202,-0.875) (1.150,-0.834) (1.107,-0.800) (1.050,-0.757) (1.008,-0.725) (0.950,-0.682) (0.908,-0.650) (0.850,-0.607) (0.807,-0.575) (0.750,-0.534) (0.703,-0.500) (0.650,-0.462) (0.600,-0.426) (0.550,-0.390) (0.500,-0.355) (0.450,-0.319) (0.400,-0.284) (0.354,-0.250) (0.300,-0.213) (0.250,-0.177) (0.206,-0.144)}
\EQ{(0.206,0.144) (0.248,0.175) (0.300,0.213) (0.350,0.248) (0.400,0.284) (0.450,0.319) (0.495,0.350) (0.550,0.390) (0.599,0.425) (0.650,0.462) (0.700,0.498) (0.750,0.534) (0.800,0.571) (0.841,0.600) (0.900,0.644) (0.942,0.675) (1.000,0.719) (1.041,0.750) (1.100,0.795) (1.150,0.834) (1.200,0.873) (1.250,0.913) (1.296,0.950) (1.350,0.995) (1.400,1.037) (1.445,1.075) (1.500,1.123) (1.550,1.168) (1.600,1.214) (1.639,1.250) (1.675,1.285) (1.716,1.325) (1.765,1.375) (1.800,1.411) (1.837,1.450) (1.875,1.492) (1.925,1.549) (1.968,1.600) (2.000,1.639) (2.048,1.700) (2.086,1.750) (2.122,1.800) (2.156,1.850) (2.189,1.900) (2.221,1.950) (2.250,1.999) (2.275,2.043) (2.300,2.088) (2.325,2.136) (2.350,2.187) (2.375,2.241) (2.400,2.299) (2.421,2.350) (2.440,2.400) (2.457,2.450) (2.474,2.500) (2.489,2.550) (2.502,2.600) (2.515,2.650) (2.526,2.700) (2.537,2.750) (2.545,2.800) (2.553,2.850) (2.560,2.900) (2.565,2.950) (2.569,3.000) (2.572,3.050) (2.574,3.100) (2.575,3.150) (2.574,3.200) (2.572,3.250) (2.570,3.300) (2.566,3.350) (2.560,3.400) (2.554,3.450) (2.546,3.500) (2.538,3.550) (2.528,3.600) (2.516,3.650) (2.504,3.700) (2.490,3.750) (2.475,3.800) (2.459,3.850) (2.442,3.900) (2.423,3.950) (2.402,4.000) (2.381,4.050) (2.357,4.100) (2.333,4.150) (2.307,4.200) (2.279,4.250) (2.250,4.299) (2.218,4.350) (2.185,4.400) (2.151,4.450) (2.114,4.500) (2.075,4.550) (2.034,4.600) (2.000,4.640) (1.950,4.695) (1.900,4.747) (1.850,4.797) (1.800,4.843) (1.750,4.887) (1.705,4.925) (1.650,4.969) (1.609,5.000) (1.550,5.042) (1.502,5.075) (1.450,5.109) (1.400,5.140) (1.350,5.169) (1.300,5.197) (1.250,5.223) (1.200,5.249) (1.150,5.272) (1.100,5.295) (1.050,5.316) (1.000,5.337) (0.950,5.356) (0.900,5.374) (0.850,5.390) (0.800,5.406) (0.750,5.421) (0.700,5.434) (0.650,5.447) (0.600,5.459) (0.550,5.469) (0.500,5.479) (0.450,5.488) (0.400,5.496) (0.350,5.503) (0.300,5.508) (0.250,5.513) (0.200,5.518) (0.150,5.521) (0.100,5.523) (0.050,5.524) (0.000,5.525) (-0.050,5.524) (-0.100,5.523) (-0.150,5.521) (-0.200,5.518) (-0.250,5.513) (-0.300,5.508) (-0.350,5.503) (-0.400,5.496) (-0.450,5.488) (-0.500,5.479) (-0.550,5.469) (-0.600,5.459) (-0.650,5.447) (-0.700,5.434) (-0.750,5.421) (-0.800,5.406) (-0.850,5.390) (-0.900,5.374) (-0.950,5.356) (-1.000,5.337) (-1.050,5.316) (-1.100,5.295) (-1.150,5.272) (-1.197,5.250) (-1.247,5.225) (-1.294,5.200) (-1.340,5.175) (-1.400,5.140) (-1.450,5.109) (-1.500,5.076) (-1.550,5.042) (-1.600,5.006) (-1.642,4.975) (-1.700,4.929) (-1.750,4.887) (-1.792,4.850) (-1.846,4.800) (-1.898,4.750) (-1.946,4.700) (-1.991,4.650) (-2.025,4.611) (-2.075,4.550) (-2.114,4.500) (-2.150,4.451) (-2.185,4.400) (-2.218,4.350) (-2.249,4.300) (-2.275,4.257) (-2.300,4.212) (-2.325,4.165) (-2.350,4.115) (-2.375,4.062) (-2.400,4.005) (-2.423,3.950) (-2.442,3.900) (-2.459,3.850) (-2.475,3.801) (-2.490,3.750) (-2.504,3.700) (-2.516,3.650) (-2.528,3.600) (-2.538,3.550) (-2.546,3.500) (-2.554,3.450) (-2.560,3.400) (-2.566,3.350) (-2.570,3.300) (-2.572,3.250) (-2.574,3.200) (-2.575,3.150) (-2.574,3.100) (-2.572,3.050) (-2.569,3.000) (-2.565,2.950) (-2.560,2.900) (-2.553,2.850) (-2.545,2.800) (-2.537,2.750) (-2.526,2.700) (-2.515,2.650) (-2.502,2.600) (-2.489,2.550) (-2.474,2.500) (-2.457,2.450) (-2.440,2.400) (-2.421,2.350) (-2.401,2.300) (-2.379,2.250) (-2.356,2.200) (-2.332,2.150) (-2.306,2.100) (-2.279,2.050) (-2.251,2.000) (-2.221,1.950) (-2.189,1.900) (-2.156,1.850) (-2.125,1.805) (-2.086,1.750) (-2.050,1.702) (-2.009,1.650) (-1.975,1.608) (-1.926,1.550) (-1.882,1.500) (-1.837,1.450) (-1.800,1.411) (-1.765,1.375) (-1.725,1.334) (-1.675,1.285) (-1.639,1.250) (-1.600,1.214) (-1.558,1.175) (-1.502,1.125) (-1.450,1.080) (-1.400,1.037) (-1.356,1.000) (-1.300,0.954) (-1.250,0.913) (-1.202,0.875) (-1.150,0.834) (-1.107,0.800) (-1.050,0.757) (-1.008,0.725) (-0.950,0.682) (-0.908,0.650) (-0.850,0.607) (-0.807,0.575) (-0.750,0.534) (-0.703,0.500) (-0.650,0.462) (-0.600,0.426) (-0.550,0.390) (-0.500,0.355) (-0.450,0.319) (-0.400,0.284) (-0.354,0.250) (-0.300,0.213) (-0.250,0.177) (-0.220,0.155)}
}
  \icono
  \marcoMMA{10}{5}{$x$ (u.\,arb.)}{$z$ (u.\,arb.)}
\end{tikzpicture}\\[1pt]{\small\scshape (a)}
\end{minipage}\hfill
\begin{minipage}[t]{0.325\textwidth}\centering
{\footnotesize Equipotenciales y líneas de campo}\\[3pt]
\begin{tikzpicture}[x=0.198cm, y=0.198cm]
  \panelMMA{10}{mmafondo}{1.2}{%
    % Cuadrupolo puntual - equipotenciales
% Generado por scripts/generar_datos.py (no editar a mano).
\EQ{(-0.220,-0.155) (-0.275,-0.195) (-0.319,-0.225) (-0.375,-0.266) (-0.425,-0.300) (-0.475,-0.336) (-0.525,-0.371) (-0.568,-0.400) (-0.625,-0.441) (-0.675,-0.475) (-0.725,-0.510) (-0.775,-0.544) (-0.821,-0.575) (-0.875,-0.612) (-0.925,-0.646) (-0.969,-0.675) (-1.025,-0.713) (-1.075,-0.745) (-1.121,-0.775) (-1.175,-0.810) (-1.225,-0.841) (-1.275,-0.872) (-1.321,-0.900) (-1.375,-0.933) (-1.425,-0.963) (-1.475,-0.991) (-1.525,-1.020) (-1.575,-1.047) (-1.625,-1.074) (-1.674,-1.100) (-1.723,-1.125) (-1.774,-1.150) (-1.825,-1.175) (-1.875,-1.198) (-1.925,-1.220) (-1.975,-1.241) (-2.025,-1.262) (-2.075,-1.281) (-2.125,-1.300) (-2.175,-1.317) (-2.225,-1.334) (-2.275,-1.349) (-2.325,-1.364) (-2.375,-1.377) (-2.425,-1.389) (-2.475,-1.400) (-2.525,-1.410) (-2.575,-1.419) (-2.625,-1.426) (-2.675,-1.433) (-2.725,-1.438) (-2.775,-1.441) (-2.825,-1.444) (-2.875,-1.445) (-2.925,-1.445) (-2.975,-1.443) (-3.025,-1.440) (-3.075,-1.435) (-3.125,-1.429) (-3.175,-1.421) (-3.225,-1.412) (-3.275,-1.401) (-3.325,-1.388) (-3.375,-1.373) (-3.425,-1.357) (-3.475,-1.338) (-3.525,-1.318) (-3.575,-1.295) (-3.625,-1.270) (-3.675,-1.242) (-3.725,-1.212) (-3.775,-1.179) (-3.825,-1.142) (-3.875,-1.102) (-3.925,-1.058) (-3.975,-1.010) (-4.025,-0.956) (-4.072,-0.900) (-4.110,-0.850) (-4.145,-0.800) (-4.175,-0.753) (-4.206,-0.700) (-4.232,-0.650) (-4.256,-0.600) (-4.277,-0.550) (-4.297,-0.500) (-4.314,-0.450) (-4.329,-0.400) (-4.343,-0.350) (-4.354,-0.300) (-4.364,-0.250) (-4.371,-0.200) (-4.377,-0.150) (-4.382,-0.100) (-4.384,-0.050) (-4.385,0.000) (-4.384,0.050) (-4.382,0.100) (-4.377,0.150) (-4.371,0.200) (-4.364,0.250) (-4.354,0.300) (-4.343,0.350) (-4.329,0.400) (-4.314,0.450) (-4.297,0.500) (-4.277,0.550) (-4.256,0.600) (-4.232,0.650) (-4.206,0.700) (-4.177,0.750) (-4.150,0.792) (-4.110,0.850) (-4.075,0.896) (-4.030,0.950) (-3.985,1.000) (-3.934,1.050) (-3.878,1.100) (-3.825,1.142) (-3.780,1.175) (-3.725,1.212) (-3.675,1.242) (-3.625,1.270) (-3.575,1.295) (-3.525,1.318) (-3.475,1.338) (-3.425,1.357) (-3.375,1.373) (-3.325,1.388) (-3.275,1.401) (-3.225,1.412) (-3.175,1.421) (-3.125,1.429) (-3.075,1.435) (-3.025,1.440) (-2.975,1.443) (-2.925,1.445) (-2.875,1.445) (-2.825,1.444) (-2.775,1.441) (-2.725,1.438) (-2.675,1.433) (-2.625,1.426) (-2.575,1.419) (-2.525,1.410) (-2.475,1.400) (-2.425,1.389) (-2.375,1.377) (-2.325,1.364) (-2.275,1.349) (-2.225,1.334) (-2.175,1.317) (-2.126,1.300) (-2.075,1.281) (-2.025,1.262) (-1.975,1.241) (-1.925,1.220) (-1.875,1.198) (-1.826,1.175) (-1.775,1.151) (-1.725,1.126) (-1.675,1.100) (-1.627,1.075) (-1.580,1.050) (-1.535,1.025) (-1.491,1.000) (-1.447,0.975) (-1.400,0.948) (-1.350,0.918) (-1.300,0.888) (-1.250,0.857) (-1.200,0.825) (-1.150,0.794) (-1.100,0.762) (-1.050,0.729) (-1.007,0.700) (-0.950,0.663) (-0.900,0.629) (-0.858,0.600) (-0.800,0.561) (-0.750,0.527) (-0.700,0.493) (-0.650,0.458) (-0.604,0.425) (-0.550,0.388) (-0.500,0.353) (-0.450,0.318) (-0.400,0.283) (-0.354,0.250) (-0.300,0.213) (-0.250,0.177) (-0.220,0.155)}
\EQ{(0.206,0.144) (0.248,0.175) (0.300,0.213) (0.350,0.248) (0.400,0.283) (0.450,0.318) (0.496,0.350) (0.550,0.388) (0.600,0.423) (0.650,0.458) (0.700,0.493) (0.748,0.525) (0.800,0.561) (0.850,0.595) (0.895,0.625) (0.950,0.663) (1.000,0.696) (1.045,0.725) (1.100,0.762) (1.150,0.794) (1.200,0.825) (1.250,0.857) (1.300,0.888) (1.350,0.918) (1.400,0.948) (1.447,0.975) (1.491,1.000) (1.535,1.025) (1.580,1.050) (1.625,1.074) (1.674,1.100) (1.723,1.125) (1.774,1.150) (1.825,1.175) (1.875,1.198) (1.925,1.220) (1.975,1.241) (2.025,1.262) (2.075,1.281) (2.125,1.300) (2.175,1.317) (2.225,1.334) (2.275,1.349) (2.325,1.364) (2.375,1.377) (2.425,1.389) (2.475,1.400) (2.525,1.410) (2.575,1.419) (2.625,1.426) (2.675,1.433) (2.725,1.438) (2.775,1.441) (2.825,1.444) (2.875,1.445) (2.925,1.445) (2.975,1.443) (3.025,1.440) (3.075,1.435) (3.125,1.429) (3.175,1.421) (3.225,1.412) (3.275,1.401) (3.325,1.388) (3.375,1.373) (3.425,1.357) (3.475,1.338) (3.525,1.318) (3.575,1.295) (3.625,1.270) (3.675,1.242) (3.725,1.212) (3.775,1.179) (3.825,1.142) (3.875,1.102) (3.925,1.058) (3.975,1.010) (4.025,0.956) (4.072,0.900) (4.110,0.850) (4.145,0.800) (4.175,0.753) (4.206,0.700) (4.232,0.650) (4.256,0.600) (4.277,0.550) (4.297,0.500) (4.314,0.450) (4.329,0.400) (4.343,0.350) (4.354,0.300) (4.364,0.250) (4.371,0.200) (4.377,0.150) (4.382,0.100) (4.384,0.050) (4.385,0.000) (4.384,-0.050) (4.382,-0.100) (4.377,-0.150) (4.371,-0.200) (4.364,-0.250) (4.354,-0.300) (4.343,-0.350) (4.329,-0.400) (4.314,-0.450) (4.297,-0.500) (4.277,-0.550) (4.256,-0.600) (4.232,-0.650) (4.206,-0.700) (4.177,-0.750) (4.150,-0.792) (4.110,-0.850) (4.075,-0.896) (4.030,-0.950) (3.985,-1.000) (3.934,-1.050) (3.878,-1.100) (3.825,-1.142) (3.780,-1.175) (3.725,-1.212) (3.675,-1.242) (3.625,-1.270) (3.575,-1.295) (3.525,-1.318) (3.475,-1.338) (3.425,-1.357) (3.375,-1.373) (3.325,-1.388) (3.275,-1.401) (3.225,-1.412) (3.175,-1.421) (3.125,-1.429) (3.075,-1.435) (3.025,-1.440) (2.975,-1.443) (2.925,-1.445) (2.875,-1.445) (2.825,-1.444) (2.775,-1.441) (2.725,-1.438) (2.675,-1.433) (2.625,-1.426) (2.575,-1.419) (2.525,-1.410) (2.475,-1.400) (2.425,-1.389) (2.375,-1.377) (2.325,-1.364) (2.275,-1.349) (2.225,-1.334) (2.175,-1.317) (2.126,-1.300) (2.075,-1.281) (2.025,-1.262) (1.975,-1.241) (1.925,-1.220) (1.875,-1.198) (1.826,-1.175) (1.775,-1.151) (1.725,-1.126) (1.675,-1.100) (1.627,-1.075) (1.580,-1.050) (1.535,-1.025) (1.491,-1.000) (1.447,-0.975) (1.400,-0.948) (1.350,-0.918) (1.300,-0.888) (1.250,-0.857) (1.200,-0.825) (1.150,-0.794) (1.100,-0.762) (1.050,-0.729) (1.007,-0.700) (0.950,-0.663) (0.900,-0.629) (0.858,-0.600) (0.800,-0.561) (0.750,-0.527) (0.700,-0.493) (0.650,-0.458) (0.604,-0.425) (0.550,-0.388) (0.500,-0.353) (0.450,-0.318) (0.400,-0.283) (0.354,-0.250) (0.300,-0.213) (0.250,-0.177) (0.206,-0.144)}
\EQ{(-0.220,-0.155) (-0.275,-0.195) (-0.319,-0.225) (-0.375,-0.266) (-0.425,-0.300) (-0.475,-0.336) (-0.525,-0.371) (-0.568,-0.400) (-0.625,-0.441) (-0.675,-0.475) (-0.725,-0.510) (-0.775,-0.545) (-0.820,-0.575) (-0.875,-0.613) (-0.925,-0.647) (-0.968,-0.675) (-1.025,-0.714) (-1.075,-0.747) (-1.119,-0.775) (-1.175,-0.812) (-1.225,-0.844) (-1.275,-0.875) (-1.325,-0.906) (-1.375,-0.937) (-1.425,-0.967) (-1.475,-0.996) (-1.525,-1.025) (-1.569,-1.050) (-1.614,-1.075) (-1.660,-1.100) (-1.707,-1.125) (-1.755,-1.150) (-1.800,-1.173) (-1.850,-1.197) (-1.900,-1.221) (-1.950,-1.244) (-2.000,-1.266) (-2.050,-1.288) (-2.100,-1.308) (-2.150,-1.327) (-2.200,-1.346) (-2.250,-1.363) (-2.300,-1.380) (-2.350,-1.396) (-2.400,-1.410) (-2.450,-1.424) (-2.500,-1.436) (-2.550,-1.447) (-2.600,-1.458) (-2.650,-1.467) (-2.700,-1.475) (-2.750,-1.481) (-2.800,-1.487) (-2.850,-1.491) (-2.900,-1.494) (-2.950,-1.496) (-3.000,-1.497) (-3.050,-1.496) (-3.100,-1.493) (-3.150,-1.490) (-3.200,-1.485) (-3.250,-1.478) (-3.300,-1.470) (-3.350,-1.460) (-3.400,-1.449) (-3.450,-1.436) (-3.500,-1.421) (-3.550,-1.405) (-3.600,-1.386) (-3.650,-1.365) (-3.700,-1.343) (-3.750,-1.318) (-3.800,-1.290) (-3.850,-1.260) (-3.900,-1.228) (-3.950,-1.192) (-4.000,-1.153) (-4.050,-1.110) (-4.088,-1.075) (-4.125,-1.038) (-4.160,-1.000) (-4.200,-0.954) (-4.243,-0.900) (-4.275,-0.855) (-4.312,-0.800) (-4.342,-0.750) (-4.369,-0.700) (-4.395,-0.650) (-4.417,-0.600) (-4.438,-0.550) (-4.457,-0.500) (-4.473,-0.450) (-4.488,-0.400) (-4.501,-0.350) (-4.512,-0.300) (-4.521,-0.250) (-4.529,-0.200) (-4.534,-0.150) (-4.539,-0.100) (-4.541,-0.050) (-4.542,0.000) (-4.541,0.050) (-4.539,0.100) (-4.534,0.150) (-4.529,0.200) (-4.521,0.250) (-4.512,0.300) (-4.501,0.350) (-4.488,0.400) (-4.473,0.450) (-4.457,0.500) (-4.438,0.550) (-4.417,0.600) (-4.395,0.650) (-4.369,0.700) (-4.342,0.750) (-4.312,0.800) (-4.279,0.850) (-4.243,0.900) (-4.203,0.950) (-4.160,1.000) (-4.125,1.038) (-4.088,1.075) (-4.050,1.110) (-4.003,1.150) (-3.950,1.192) (-3.904,1.225) (-3.850,1.260) (-3.800,1.290) (-3.750,1.318) (-3.700,1.343) (-3.650,1.365) (-3.600,1.386) (-3.550,1.405) (-3.500,1.421) (-3.450,1.436) (-3.400,1.449) (-3.350,1.460) (-3.300,1.470) (-3.250,1.478) (-3.200,1.485) (-3.150,1.490) (-3.100,1.493) (-3.050,1.496) (-3.000,1.497) (-2.950,1.496) (-2.900,1.494) (-2.850,1.491) (-2.800,1.487) (-2.750,1.481) (-2.700,1.475) (-2.650,1.467) (-2.600,1.458) (-2.550,1.447) (-2.500,1.436) (-2.450,1.424) (-2.400,1.410) (-2.350,1.396) (-2.300,1.380) (-2.250,1.363) (-2.200,1.346) (-2.150,1.327) (-2.100,1.308) (-2.050,1.288) (-2.000,1.266) (-1.950,1.244) (-1.900,1.221) (-1.850,1.197) (-1.804,1.175) (-1.755,1.150) (-1.707,1.125) (-1.660,1.100) (-1.614,1.075) (-1.569,1.050) (-1.525,1.025) (-1.482,1.000) (-1.425,0.967) (-1.375,0.937) (-1.325,0.906) (-1.275,0.875) (-1.225,0.844) (-1.175,0.812) (-1.125,0.779) (-1.081,0.750) (-1.025,0.714) (-0.975,0.680) (-0.930,0.650) (-0.875,0.613) (-0.825,0.579) (-0.784,0.550) (-0.725,0.510) (-0.675,0.475) (-0.625,0.441) (-0.575,0.406) (-0.532,0.375) (-0.475,0.336) (-0.425,0.300) (-0.375,0.266) (-0.325,0.230) (-0.284,0.200) (-0.225,0.160) (-0.220,0.155)}
\EQ{(0.206,0.144) (0.248,0.175) (0.300,0.213) (0.350,0.248) (0.400,0.283) (0.450,0.318) (0.496,0.350) (0.550,0.388) (0.600,0.423) (0.650,0.458) (0.700,0.493) (0.747,0.525) (0.800,0.562) (0.850,0.596) (0.894,0.625) (0.950,0.664) (1.000,0.697) (1.043,0.725) (1.100,0.763) (1.150,0.796) (1.196,0.825) (1.250,0.859) (1.300,0.891) (1.350,0.921) (1.397,0.950) (1.450,0.982) (1.500,1.011) (1.550,1.040) (1.600,1.068) (1.650,1.095) (1.700,1.122) (1.750,1.148) (1.800,1.173) (1.850,1.197) (1.900,1.221) (1.950,1.244) (2.000,1.266) (2.050,1.288) (2.100,1.308) (2.150,1.327) (2.200,1.346) (2.250,1.363) (2.300,1.380) (2.350,1.396) (2.400,1.410) (2.450,1.424) (2.500,1.436) (2.550,1.447) (2.600,1.458) (2.650,1.467) (2.700,1.475) (2.750,1.481) (2.800,1.487) (2.850,1.491) (2.900,1.494) (2.950,1.496) (3.000,1.497) (3.050,1.496) (3.100,1.493) (3.150,1.490) (3.200,1.485) (3.250,1.478) (3.300,1.470) (3.350,1.460) (3.400,1.449) (3.450,1.436) (3.500,1.421) (3.550,1.405) (3.600,1.386) (3.650,1.365) (3.700,1.343) (3.750,1.318) (3.800,1.290) (3.850,1.260) (3.900,1.228) (3.950,1.192) (4.000,1.153) (4.050,1.110) (4.088,1.075) (4.125,1.038) (4.160,1.000) (4.200,0.954) (4.243,0.900) (4.275,0.855) (4.312,0.800) (4.342,0.750) (4.369,0.700) (4.395,0.650) (4.417,0.600) (4.438,0.550) (4.457,0.500) (4.473,0.450) (4.488,0.400) (4.501,0.350) (4.512,0.300) (4.521,0.250) (4.529,0.200) (4.534,0.150) (4.539,0.100) (4.541,0.050) (4.542,0.000) (4.541,-0.050) (4.539,-0.100) (4.534,-0.150) (4.529,-0.200) (4.521,-0.250) (4.512,-0.300) (4.501,-0.350) (4.488,-0.400) (4.473,-0.450) (4.457,-0.500) (4.438,-0.550) (4.417,-0.600) (4.395,-0.650) (4.369,-0.700) (4.342,-0.750) (4.312,-0.800) (4.279,-0.850) (4.243,-0.900) (4.203,-0.950) (4.160,-1.000) (4.125,-1.038) (4.088,-1.075) (4.050,-1.110) (4.003,-1.150) (3.950,-1.192) (3.904,-1.225) (3.850,-1.260) (3.800,-1.290) (3.750,-1.318) (3.700,-1.343) (3.650,-1.365) (3.600,-1.386) (3.550,-1.405) (3.500,-1.421) (3.450,-1.436) (3.400,-1.449) (3.350,-1.460) (3.300,-1.470) (3.250,-1.478) (3.200,-1.485) (3.150,-1.490) (3.100,-1.493) (3.050,-1.496) (3.000,-1.497) (2.950,-1.496) (2.900,-1.494) (2.850,-1.491) (2.800,-1.487) (2.750,-1.481) (2.700,-1.475) (2.650,-1.467) (2.600,-1.458) (2.550,-1.447) (2.500,-1.436) (2.450,-1.424) (2.400,-1.410) (2.350,-1.396) (2.300,-1.380) (2.250,-1.363) (2.200,-1.346) (2.150,-1.327) (2.100,-1.308) (2.050,-1.288) (2.000,-1.266) (1.950,-1.244) (1.900,-1.221) (1.850,-1.197) (1.804,-1.175) (1.755,-1.150) (1.707,-1.125) (1.660,-1.100) (1.614,-1.075) (1.569,-1.050) (1.525,-1.025) (1.482,-1.000) (1.425,-0.967) (1.375,-0.937) (1.325,-0.906) (1.275,-0.875) (1.225,-0.844) (1.175,-0.812) (1.125,-0.779) (1.081,-0.750) (1.025,-0.714) (0.975,-0.680) (0.930,-0.650) (0.875,-0.613) (0.825,-0.579) (0.784,-0.550) (0.725,-0.510) (0.675,-0.475) (0.625,-0.441) (0.575,-0.406) (0.532,-0.375) (0.475,-0.336) (0.425,-0.300) (0.375,-0.266) (0.325,-0.230) (0.284,-0.200) (0.225,-0.160) (0.206,-0.144)}
\EQ{(-0.220,-0.155) (-0.275,-0.195) (-0.319,-0.225) (-0.375,-0.266) (-0.425,-0.300) (-0.475,-0.336) (-0.525,-0.371) (-0.568,-0.400) (-0.625,-0.441) (-0.674,-0.475) (-0.725,-0.510) (-0.775,-0.545) (-0.819,-0.575) (-0.875,-0.614) (-0.925,-0.647) (-0.975,-0.681) (-1.025,-0.715) (-1.075,-0.748) (-1.125,-0.781) (-1.175,-0.814) (-1.225,-0.846) (-1.271,-0.875) (-1.325,-0.910) (-1.375,-0.941) (-1.425,-0.971) (-1.473,-1.000) (-1.525,-1.031) (-1.575,-1.060) (-1.625,-1.089) (-1.675,-1.116) (-1.725,-1.144) (-1.775,-1.170) (-1.825,-1.196) (-1.875,-1.221) (-1.925,-1.246) (-1.975,-1.270) (-2.025,-1.293) (-2.075,-1.315) (-2.125,-1.336) (-2.175,-1.357) (-2.222,-1.375) (-2.275,-1.395) (-2.325,-1.413) (-2.375,-1.430) (-2.425,-1.446) (-2.475,-1.461) (-2.525,-1.474) (-2.575,-1.487) (-2.625,-1.499) (-2.675,-1.510) (-2.725,-1.520) (-2.775,-1.529) (-2.825,-1.536) (-2.875,-1.542) (-2.925,-1.548) (-2.975,-1.552) (-3.025,-1.555) (-3.075,-1.556) (-3.125,-1.556) (-3.175,-1.556) (-3.225,-1.553) (-3.275,-1.550) (-3.325,-1.545) (-3.375,-1.538) (-3.425,-1.530) (-3.475,-1.521) (-3.525,-1.510) (-3.575,-1.497) (-3.625,-1.483) (-3.675,-1.467) (-3.725,-1.449) (-3.775,-1.429) (-3.825,-1.407) (-3.875,-1.383) (-3.925,-1.357) (-3.975,-1.329) (-4.021,-1.300) (-4.075,-1.264) (-4.125,-1.227) (-4.175,-1.187) (-4.217,-1.150) (-4.270,-1.100) (-4.318,-1.050) (-4.362,-1.000) (-4.400,-0.952) (-4.439,-0.900) (-4.473,-0.850) (-4.500,-0.807) (-4.533,-0.750) (-4.559,-0.700) (-4.583,-0.650) (-4.605,-0.600) (-4.624,-0.550) (-4.642,-0.500) (-4.658,-0.450) (-4.672,-0.400) (-4.684,-0.350) (-4.695,-0.300) (-4.704,-0.250) (-4.711,-0.200) (-4.717,-0.150) (-4.720,-0.100) (-4.723,-0.050) (-4.724,0.000) (-4.723,0.050) (-4.720,0.100) (-4.717,0.150) (-4.711,0.200) (-4.704,0.250) (-4.695,0.300) (-4.684,0.350) (-4.672,0.400) (-4.658,0.450) (-4.642,0.500) (-4.625,0.548) (-4.605,0.600) (-4.583,0.650) (-4.559,0.700) (-4.533,0.750) (-4.504,0.800) (-4.475,0.847) (-4.439,0.900) (-4.402,0.950) (-4.362,1.000) (-4.325,1.042) (-4.275,1.095) (-4.225,1.143) (-4.175,1.187) (-4.127,1.225) (-4.075,1.264) (-4.025,1.298) (-3.981,1.325) (-3.925,1.357) (-3.875,1.383) (-3.825,1.407) (-3.775,1.429) (-3.725,1.449) (-3.675,1.467) (-3.625,1.483) (-3.575,1.497) (-3.525,1.510) (-3.475,1.521) (-3.425,1.530) (-3.375,1.538) (-3.325,1.545) (-3.275,1.550) (-3.225,1.553) (-3.175,1.556) (-3.125,1.556) (-3.075,1.556) (-3.025,1.555) (-2.975,1.552) (-2.925,1.548) (-2.875,1.542) (-2.825,1.536) (-2.775,1.529) (-2.725,1.520) (-2.675,1.510) (-2.625,1.499) (-2.575,1.487) (-2.525,1.474) (-2.475,1.461) (-2.425,1.446) (-2.375,1.430) (-2.325,1.413) (-2.275,1.395) (-2.225,1.376) (-2.175,1.357) (-2.125,1.336) (-2.075,1.315) (-2.025,1.293) (-1.975,1.270) (-1.925,1.246) (-1.875,1.221) (-1.825,1.196) (-1.775,1.170) (-1.725,1.144) (-1.675,1.116) (-1.625,1.089) (-1.575,1.060) (-1.525,1.031) (-1.475,1.001) (-1.432,0.975) (-1.375,0.941) (-1.325,0.910) (-1.275,0.878) (-1.232,0.850) (-1.175,0.814) (-1.125,0.781) (-1.078,0.750) (-1.025,0.715) (-0.975,0.681) (-0.929,0.650) (-0.875,0.614) (-0.825,0.579) (-0.783,0.550) (-0.725,0.510) (-0.675,0.475) (-0.625,0.441) (-0.575,0.406) (-0.532,0.375) (-0.475,0.336) (-0.425,0.300) (-0.375,0.266) (-0.325,0.230) (-0.284,0.200) (-0.225,0.160) (-0.220,0.155)}
\EQ{(0.206,0.144) (0.248,0.175) (0.300,0.213) (0.350,0.248) (0.400,0.283) (0.450,0.318) (0.496,0.350) (0.550,0.389) (0.600,0.423) (0.650,0.458) (0.700,0.493) (0.747,0.525) (0.800,0.562) (0.850,0.596) (0.893,0.625) (0.950,0.665) (1.000,0.698) (1.050,0.732) (1.100,0.765) (1.150,0.797) (1.193,0.825) (1.250,0.862) (1.300,0.894) (1.350,0.925) (1.400,0.956) (1.450,0.986) (1.500,1.016) (1.550,1.046) (1.600,1.074) (1.646,1.100) (1.691,1.125) (1.737,1.150) (1.785,1.175) (1.833,1.200) (1.883,1.225) (1.934,1.250) (1.987,1.275) (2.042,1.300) (2.099,1.325) (2.150,1.346) (2.200,1.367) (2.250,1.386) (2.300,1.404) (2.350,1.421) (2.400,1.438) (2.450,1.453) (2.500,1.468) (2.550,1.481) (2.600,1.493) (2.650,1.505) (2.700,1.515) (2.750,1.524) (2.800,1.532) (2.850,1.539) (2.900,1.545) (2.950,1.550) (3.000,1.553) (3.050,1.556) (3.100,1.556) (3.150,1.556) (3.200,1.555) (3.250,1.552) (3.300,1.547) (3.350,1.542) (3.400,1.534) (3.450,1.526) (3.500,1.515) (3.550,1.504) (3.600,1.490) (3.650,1.475) (3.700,1.458) (3.750,1.439) (3.800,1.418) (3.850,1.396) (3.900,1.371) (3.950,1.343) (4.000,1.313) (4.050,1.281) (4.094,1.250) (4.150,1.207) (4.189,1.175) (4.244,1.125) (4.294,1.075) (4.340,1.025) (4.375,0.984) (4.421,0.925) (4.450,0.884) (4.489,0.825) (4.519,0.775) (4.546,0.725) (4.571,0.675) (4.594,0.625) (4.615,0.575) (4.633,0.525) (4.650,0.475) (4.665,0.425) (4.678,0.375) (4.690,0.325) (4.699,0.275) (4.708,0.225) (4.714,0.175) (4.719,0.125) (4.722,0.075) (4.723,0.025) (4.723,-0.050) (4.720,-0.100) (4.717,-0.150) (4.711,-0.200) (4.704,-0.250) (4.695,-0.300) (4.684,-0.350) (4.672,-0.400) (4.658,-0.450) (4.642,-0.500) (4.625,-0.548) (4.605,-0.600) (4.583,-0.650) (4.559,-0.700) (4.533,-0.750) (4.504,-0.800) (4.475,-0.847) (4.439,-0.900) (4.402,-0.950) (4.362,-1.000) (4.325,-1.042) (4.275,-1.095) (4.225,-1.143) (4.175,-1.187) (4.127,-1.225) (4.075,-1.264) (4.025,-1.298) (3.981,-1.325) (3.925,-1.357) (3.875,-1.383) (3.825,-1.407) (3.775,-1.429) (3.725,-1.449) (3.675,-1.467) (3.625,-1.483) (3.575,-1.497) (3.525,-1.510) (3.475,-1.521) (3.425,-1.530) (3.375,-1.538) (3.325,-1.545) (3.275,-1.550) (3.225,-1.553) (3.175,-1.556) (3.125,-1.556) (3.075,-1.556) (3.025,-1.555) (2.975,-1.552) (2.925,-1.548) (2.875,-1.542) (2.825,-1.536) (2.775,-1.529) (2.725,-1.520) (2.675,-1.510) (2.625,-1.499) (2.575,-1.487) (2.525,-1.474) (2.475,-1.461) (2.425,-1.446) (2.375,-1.430) (2.325,-1.413) (2.275,-1.395) (2.225,-1.376) (2.175,-1.357) (2.125,-1.336) (2.075,-1.315) (2.025,-1.293) (1.975,-1.270) (1.925,-1.246) (1.875,-1.221) (1.825,-1.196) (1.775,-1.170) (1.725,-1.144) (1.675,-1.116) (1.625,-1.089) (1.575,-1.060) (1.525,-1.031) (1.475,-1.001) (1.432,-0.975) (1.375,-0.941) (1.325,-0.910) (1.275,-0.878) (1.232,-0.850) (1.175,-0.814) (1.125,-0.781) (1.078,-0.750) (1.025,-0.715) (0.975,-0.681) (0.929,-0.650) (0.875,-0.614) (0.825,-0.579) (0.783,-0.550) (0.725,-0.510) (0.675,-0.475) (0.625,-0.441) (0.575,-0.406) (0.532,-0.375) (0.475,-0.336) (0.425,-0.300) (0.375,-0.266) (0.325,-0.230) (0.284,-0.200) (0.225,-0.160) (0.206,-0.144)}
\EQ{(-0.220,-0.155) (-0.275,-0.195) (-0.319,-0.225) (-0.375,-0.266) (-0.425,-0.300) (-0.475,-0.336) (-0.525,-0.371) (-0.568,-0.400) (-0.625,-0.441) (-0.674,-0.475) (-0.725,-0.511) (-0.775,-0.545) (-0.819,-0.575) (-0.875,-0.614) (-0.925,-0.648) (-0.975,-0.682) (-1.025,-0.716) (-1.075,-0.750) (-1.125,-0.783) (-1.175,-0.816) (-1.225,-0.848) (-1.275,-0.881) (-1.325,-0.913) (-1.375,-0.944) (-1.424,-0.975) (-1.475,-1.006) (-1.525,-1.037) (-1.575,-1.066) (-1.625,-1.096) (-1.675,-1.124) (-1.721,-1.150) (-1.766,-1.175) (-1.812,-1.200) (-1.859,-1.225) (-1.907,-1.250) (-1.957,-1.275) (-2.008,-1.300) (-2.060,-1.325) (-2.114,-1.350) (-2.171,-1.375) (-2.225,-1.398) (-2.275,-1.419) (-2.325,-1.438) (-2.375,-1.457) (-2.425,-1.475) (-2.475,-1.492) (-2.525,-1.508) (-2.575,-1.523) (-2.625,-1.538) (-2.675,-1.551) (-2.725,-1.563) (-2.775,-1.574) (-2.825,-1.585) (-2.875,-1.594) (-2.925,-1.602) (-2.975,-1.609) (-3.025,-1.615) (-3.075,-1.620) (-3.125,-1.623) (-3.175,-1.626) (-3.225,-1.627) (-3.275,-1.627) (-3.325,-1.626) (-3.375,-1.624) (-3.425,-1.620) (-3.475,-1.615) (-3.525,-1.609) (-3.575,-1.601) (-3.625,-1.592) (-3.675,-1.581) (-3.725,-1.569) (-3.775,-1.555) (-3.825,-1.539) (-3.875,-1.522) (-3.925,-1.503) (-3.975,-1.482) (-4.025,-1.459) (-4.075,-1.434) (-4.125,-1.407) (-4.175,-1.378) (-4.218,-1.350) (-4.275,-1.311) (-4.323,-1.275) (-4.375,-1.232) (-4.425,-1.188) (-4.464,-1.150) (-4.500,-1.113) (-4.535,-1.075) (-4.575,-1.028) (-4.616,-0.975) (-4.650,-0.928) (-4.685,-0.875) (-4.716,-0.825) (-4.744,-0.775) (-4.770,-0.725) (-4.793,-0.675) (-4.815,-0.625) (-4.835,-0.575) (-4.853,-0.525) (-4.869,-0.475) (-4.883,-0.425) (-4.895,-0.375) (-4.906,-0.325) (-4.916,-0.275) (-4.923,-0.225) (-4.929,-0.175) (-4.934,-0.125) (-4.937,-0.075) (-4.938,-0.025) (-4.938,0.050) (-4.936,0.100) (-4.932,0.150) (-4.926,0.200) (-4.920,0.250) (-4.911,0.300) (-4.901,0.350) (-4.889,0.400) (-4.876,0.450) (-4.861,0.500) (-4.844,0.550) (-4.825,0.600) (-4.804,0.650) (-4.782,0.700) (-4.757,0.750) (-4.730,0.800) (-4.701,0.850) (-4.669,0.900) (-4.635,0.950) (-4.600,0.996) (-4.556,1.050) (-4.512,1.100) (-4.475,1.139) (-4.438,1.175) (-4.400,1.210) (-4.354,1.250) (-4.300,1.292) (-4.255,1.325) (-4.200,1.362) (-4.150,1.393) (-4.100,1.421) (-4.050,1.447) (-4.000,1.471) (-3.950,1.493) (-3.900,1.513) (-3.850,1.531) (-3.800,1.547) (-3.750,1.562) (-3.700,1.575) (-3.650,1.586) (-3.600,1.596) (-3.550,1.605) (-3.500,1.612) (-3.450,1.618) (-3.400,1.622) (-3.350,1.625) (-3.300,1.627) (-3.250,1.627) (-3.200,1.627) (-3.150,1.625) (-3.100,1.622) (-3.050,1.617) (-3.000,1.612) (-2.950,1.605) (-2.900,1.598) (-2.850,1.589) (-2.800,1.580) (-2.750,1.569) (-2.700,1.557) (-2.650,1.544) (-2.600,1.531) (-2.550,1.516) (-2.500,1.500) (-2.450,1.484) (-2.400,1.466) (-2.350,1.448) (-2.300,1.429) (-2.250,1.409) (-2.200,1.388) (-2.150,1.366) (-2.100,1.344) (-2.050,1.320) (-2.000,1.296) (-1.950,1.272) (-1.900,1.246) (-1.850,1.220) (-1.800,1.194) (-1.750,1.167) (-1.700,1.139) (-1.650,1.110) (-1.600,1.081) (-1.550,1.051) (-1.506,1.025) (-1.450,0.991) (-1.400,0.960) (-1.350,0.929) (-1.305,0.900) (-1.250,0.865) (-1.200,0.832) (-1.151,0.800) (-1.100,0.766) (-1.050,0.733) (-1.001,0.700) (-0.950,0.665) (-0.900,0.631) (-0.855,0.600) (-0.800,0.563) (-0.750,0.528) (-0.700,0.493) (-0.650,0.459) (-0.603,0.425) (-0.550,0.389) (-0.500,0.353) (-0.450,0.318) (-0.400,0.283) (-0.354,0.250) (-0.300,0.213) (-0.250,0.177) (-0.220,0.155)}
\EQ{(0.206,0.144) (0.248,0.175) (0.300,0.213) (0.350,0.248) (0.400,0.283) (0.450,0.318) (0.496,0.350) (0.550,0.389) (0.600,0.423) (0.650,0.459) (0.700,0.493) (0.746,0.525) (0.800,0.563) (0.850,0.597) (0.891,0.625) (0.950,0.665) (1.000,0.699) (1.050,0.733) (1.100,0.766) (1.150,0.799) (1.200,0.832) (1.250,0.865) (1.300,0.897) (1.345,0.925) (1.400,0.960) (1.450,0.991) (1.500,1.021) (1.548,1.050) (1.600,1.081) (1.650,1.110) (1.700,1.139) (1.750,1.167) (1.800,1.194) (1.850,1.220) (1.900,1.246) (1.950,1.272) (2.000,1.296) (2.050,1.320) (2.100,1.344) (2.150,1.366) (2.200,1.388) (2.250,1.409) (2.300,1.429) (2.350,1.448) (2.400,1.466) (2.450,1.484) (2.500,1.500) (2.550,1.516) (2.600,1.531) (2.650,1.544) (2.700,1.557) (2.750,1.569) (2.800,1.580) (2.850,1.589) (2.900,1.598) (2.950,1.605) (3.000,1.612) (3.050,1.617) (3.100,1.622) (3.150,1.625) (3.200,1.627) (3.250,1.627) (3.300,1.627) (3.350,1.625) (3.400,1.622) (3.450,1.618) (3.500,1.612) (3.550,1.605) (3.600,1.596) (3.650,1.586) (3.700,1.575) (3.750,1.562) (3.800,1.547) (3.850,1.531) (3.900,1.513) (3.950,1.493) (4.000,1.471) (4.050,1.447) (4.100,1.421) (4.150,1.393) (4.200,1.362) (4.250,1.328) (4.300,1.292) (4.350,1.253) (4.400,1.210) (4.438,1.175) (4.475,1.139) (4.512,1.100) (4.550,1.058) (4.597,1.000) (4.635,0.950) (4.669,0.900) (4.700,0.851) (4.730,0.800) (4.757,0.750) (4.782,0.700) (4.804,0.650) (4.825,0.600) (4.844,0.550) (4.861,0.500) (4.876,0.450) (4.889,0.400) (4.901,0.350) (4.911,0.300) (4.920,0.250) (4.926,0.200) (4.932,0.150) (4.936,0.100) (4.938,0.050) (4.939,0.000) (4.938,-0.050) (4.936,-0.100) (4.932,-0.150) (4.926,-0.200) (4.920,-0.250) (4.911,-0.300) (4.901,-0.350) (4.889,-0.400) (4.876,-0.450) (4.861,-0.500) (4.844,-0.550) (4.825,-0.600) (4.804,-0.650) (4.782,-0.700) (4.757,-0.750) (4.730,-0.800) (4.701,-0.850) (4.669,-0.900) (4.635,-0.950) (4.600,-0.996) (4.556,-1.050) (4.512,-1.100) (4.475,-1.139) (4.438,-1.175) (4.400,-1.210) (4.354,-1.250) (4.300,-1.292) (4.255,-1.325) (4.200,-1.362) (4.150,-1.393) (4.100,-1.421) (4.050,-1.447) (4.000,-1.471) (3.950,-1.493) (3.900,-1.513) (3.850,-1.531) (3.800,-1.547) (3.750,-1.562) (3.700,-1.575) (3.650,-1.586) (3.600,-1.596) (3.550,-1.605) (3.500,-1.612) (3.450,-1.618) (3.400,-1.622) (3.350,-1.625) (3.300,-1.627) (3.250,-1.627) (3.200,-1.627) (3.150,-1.625) (3.100,-1.622) (3.050,-1.617) (3.000,-1.612) (2.950,-1.605) (2.900,-1.598) (2.850,-1.589) (2.800,-1.580) (2.750,-1.569) (2.700,-1.557) (2.650,-1.544) (2.600,-1.531) (2.550,-1.516) (2.500,-1.500) (2.450,-1.484) (2.400,-1.466) (2.350,-1.448) (2.300,-1.429) (2.250,-1.409) (2.200,-1.388) (2.150,-1.366) (2.100,-1.344) (2.050,-1.320) (2.000,-1.296) (1.950,-1.272) (1.900,-1.246) (1.850,-1.220) (1.800,-1.194) (1.750,-1.167) (1.700,-1.139) (1.650,-1.110) (1.600,-1.081) (1.550,-1.051) (1.506,-1.025) (1.450,-0.991) (1.400,-0.960) (1.350,-0.929) (1.305,-0.900) (1.250,-0.865) (1.200,-0.832) (1.151,-0.800) (1.100,-0.766) (1.050,-0.733) (1.001,-0.700) (0.950,-0.665) (0.900,-0.631) (0.855,-0.600) (0.800,-0.563) (0.750,-0.528) (0.700,-0.493) (0.650,-0.459) (0.603,-0.425) (0.550,-0.389) (0.500,-0.353) (0.450,-0.318) (0.400,-0.283) (0.354,-0.250) (0.300,-0.213) (0.250,-0.177) (0.206,-0.144)}
\EQ{(-0.220,-0.155) (-0.275,-0.195) (-0.319,-0.225) (-0.375,-0.266) (-0.425,-0.300) (-0.475,-0.336) (-0.525,-0.371) (-0.567,-0.400) (-0.625,-0.441) (-0.674,-0.475) (-0.725,-0.511) (-0.775,-0.546) (-0.818,-0.575) (-0.875,-0.615) (-0.925,-0.649) (-0.975,-0.683) (-1.025,-0.717) (-1.074,-0.750) (-1.125,-0.785) (-1.175,-0.818) (-1.224,-0.850) (-1.275,-0.884) (-1.325,-0.916) (-1.375,-0.948) (-1.417,-0.975) (-1.475,-1.011) (-1.525,-1.042) (-1.575,-1.073) (-1.620,-1.100) (-1.675,-1.133) (-1.725,-1.162) (-1.775,-1.190) (-1.825,-1.218) (-1.875,-1.246) (-1.925,-1.273) (-1.975,-1.299) (-2.025,-1.325) (-2.075,-1.350) (-2.125,-1.374) (-2.175,-1.398) (-2.225,-1.421) (-2.275,-1.443) (-2.325,-1.465) (-2.375,-1.486) (-2.425,-1.505) (-2.475,-1.525) (-2.525,-1.543) (-2.575,-1.560) (-2.625,-1.577) (-2.675,-1.593) (-2.725,-1.608) (-2.775,-1.622) (-2.825,-1.635) (-2.875,-1.647) (-2.925,-1.658) (-2.975,-1.668) (-3.025,-1.677) (-3.075,-1.686) (-3.125,-1.693) (-3.175,-1.699) (-3.225,-1.704) (-3.275,-1.708) (-3.325,-1.711) (-3.375,-1.713) (-3.425,-1.713) (-3.475,-1.713) (-3.525,-1.711) (-3.575,-1.708) (-3.625,-1.704) (-3.675,-1.698) (-3.725,-1.691) (-3.775,-1.683) (-3.825,-1.674) (-3.875,-1.663) (-3.925,-1.650) (-3.975,-1.636) (-4.025,-1.621) (-4.075,-1.604) (-4.125,-1.585) (-4.175,-1.564) (-4.225,-1.542) (-4.275,-1.518) (-4.325,-1.491) (-4.375,-1.462) (-4.425,-1.431) (-4.472,-1.400) (-4.525,-1.362) (-4.572,-1.325) (-4.625,-1.280) (-4.675,-1.234) (-4.725,-1.185) (-4.775,-1.130) (-4.822,-1.075) (-4.861,-1.025) (-4.897,-0.975) (-4.925,-0.932) (-4.961,-0.875) (-4.989,-0.825) (-5.015,-0.775) (-5.040,-0.725) (-5.062,-0.675) (-5.082,-0.625) (-5.100,-0.577) (-5.117,-0.525) (-5.133,-0.475) (-5.146,-0.425) (-5.158,-0.375) (-5.168,-0.325) (-5.177,-0.275) (-5.184,-0.225) (-5.190,-0.175) (-5.195,-0.125) (-5.197,-0.075) (-5.199,-0.025) (-5.198,0.050) (-5.196,0.100) (-5.193,0.150) (-5.187,0.200) (-5.181,0.250) (-5.173,0.300) (-5.163,0.350) (-5.152,0.400) (-5.140,0.450) (-5.125,0.500) (-5.109,0.550) (-5.092,0.600) (-5.072,0.650) (-5.051,0.700) (-5.028,0.750) (-5.002,0.800) (-4.975,0.850) (-4.946,0.900) (-4.914,0.950) (-4.879,1.000) (-4.842,1.050) (-4.801,1.100) (-4.758,1.150) (-4.710,1.200) (-4.659,1.250) (-4.602,1.300) (-4.550,1.343) (-4.507,1.375) (-4.450,1.415) (-4.400,1.447) (-4.353,1.475) (-4.309,1.500) (-4.260,1.525) (-4.208,1.550) (-4.150,1.575) (-4.100,1.595) (-4.050,1.613) (-4.000,1.629) (-3.950,1.644) (-3.900,1.657) (-3.850,1.668) (-3.800,1.679) (-3.750,1.687) (-3.700,1.695) (-3.650,1.701) (-3.600,1.706) (-3.550,1.710) (-3.500,1.712) (-3.450,1.713) (-3.400,1.713) (-3.350,1.712) (-3.300,1.710) (-3.250,1.706) (-3.200,1.702) (-3.150,1.696) (-3.100,1.689) (-3.050,1.682) (-3.000,1.673) (-2.950,1.663) (-2.900,1.653) (-2.850,1.641) (-2.800,1.628) (-2.750,1.615) (-2.700,1.600) (-2.650,1.585) (-2.600,1.569) (-2.550,1.552) (-2.500,1.534) (-2.450,1.515) (-2.400,1.496) (-2.350,1.475) (-2.300,1.454) (-2.250,1.432) (-2.200,1.410) (-2.150,1.386) (-2.100,1.362) (-2.050,1.338) (-2.000,1.312) (-1.950,1.286) (-1.900,1.259) (-1.850,1.232) (-1.800,1.204) (-1.750,1.176) (-1.705,1.150) (-1.650,1.118) (-1.600,1.088) (-1.550,1.058) (-1.500,1.027) (-1.457,1.000) (-1.400,0.964) (-1.350,0.932) (-1.300,0.900) (-1.250,0.867) (-1.200,0.835) (-1.150,0.801) (-1.100,0.768) (-1.050,0.734) (-1.000,0.700) (-0.950,0.666) (-0.900,0.632) (-0.854,0.600) (-0.800,0.563) (-0.750,0.528) (-0.700,0.494) (-0.650,0.459) (-0.603,0.425) (-0.550,0.389) (-0.500,0.353) (-0.450,0.318) (-0.400,0.283) (-0.354,0.250) (-0.300,0.213) (-0.250,0.177) (-0.220,0.155)}
\EQ{(0.206,0.144) (0.248,0.175) (0.300,0.213) (0.350,0.248) (0.400,0.283) (0.450,0.318) (0.496,0.350) (0.550,0.389) (0.600,0.423) (0.650,0.459) (0.700,0.494) (0.746,0.525) (0.800,0.563) (0.850,0.598) (0.900,0.632) (0.950,0.666) (1.000,0.700) (1.050,0.734) (1.100,0.768) (1.148,0.800) (1.200,0.835) (1.250,0.867) (1.300,0.900) (1.350,0.932) (1.400,0.964) (1.450,0.996) (1.497,1.025) (1.550,1.058) (1.600,1.088) (1.650,1.118) (1.700,1.147) (1.748,1.175) (1.792,1.200) (1.837,1.225) (1.883,1.250) (1.929,1.275) (1.975,1.299) (2.025,1.325) (2.075,1.350) (2.125,1.374) (2.175,1.398) (2.225,1.421) (2.275,1.443) (2.325,1.465) (2.375,1.486) (2.425,1.505) (2.475,1.525) (2.525,1.543) (2.575,1.560) (2.625,1.577) (2.675,1.593) (2.725,1.608) (2.775,1.622) (2.825,1.635) (2.875,1.647) (2.925,1.658) (2.975,1.668) (3.025,1.677) (3.075,1.686) (3.125,1.693) (3.175,1.699) (3.225,1.704) (3.275,1.708) (3.325,1.711) (3.375,1.713) (3.425,1.713) (3.475,1.713) (3.525,1.711) (3.575,1.708) (3.625,1.704) (3.675,1.698) (3.725,1.691) (3.775,1.683) (3.825,1.674) (3.875,1.663) (3.925,1.650) (3.975,1.636) (4.025,1.621) (4.075,1.604) (4.125,1.585) (4.175,1.564) (4.225,1.542) (4.275,1.518) (4.325,1.491) (4.375,1.462) (4.425,1.431) (4.472,1.400) (4.525,1.362) (4.572,1.325) (4.625,1.280) (4.675,1.234) (4.725,1.185) (4.775,1.130) (4.822,1.075) (4.861,1.025) (4.897,0.975) (4.925,0.932) (4.961,0.875) (4.989,0.825) (5.015,0.775) (5.040,0.725) (5.062,0.675) (5.082,0.625) (5.100,0.577) (5.117,0.525) (5.133,0.475) (5.146,0.425) (5.158,0.375) (5.168,0.325) (5.177,0.275) (5.184,0.225) (5.190,0.175) (5.195,0.125) (5.197,0.075) (5.199,0.025) (5.198,-0.050) (5.196,-0.100) (5.193,-0.150) (5.187,-0.200) (5.181,-0.250) (5.173,-0.300) (5.163,-0.350) (5.152,-0.400) (5.140,-0.450) (5.125,-0.500) (5.109,-0.550) (5.092,-0.600) (5.072,-0.650) (5.051,-0.700) (5.028,-0.750) (5.002,-0.800) (4.975,-0.850) (4.946,-0.900) (4.914,-0.950) (4.879,-1.000) (4.842,-1.050) (4.801,-1.100) (4.758,-1.150) (4.710,-1.200) (4.659,-1.250) (4.602,-1.300) (4.550,-1.343) (4.507,-1.375) (4.450,-1.415) (4.400,-1.447) (4.353,-1.475) (4.309,-1.500) (4.260,-1.525) (4.208,-1.550) (4.150,-1.575) (4.100,-1.595) (4.050,-1.613) (4.000,-1.629) (3.950,-1.644) (3.900,-1.657) (3.850,-1.668) (3.800,-1.679) (3.750,-1.687) (3.700,-1.695) (3.650,-1.701) (3.600,-1.706) (3.550,-1.710) (3.500,-1.712) (3.450,-1.713) (3.400,-1.713) (3.350,-1.712) (3.300,-1.710) (3.250,-1.706) (3.200,-1.702) (3.150,-1.696) (3.100,-1.689) (3.050,-1.682) (3.000,-1.673) (2.950,-1.663) (2.900,-1.653) (2.850,-1.641) (2.800,-1.628) (2.750,-1.615) (2.700,-1.600) (2.650,-1.585) (2.600,-1.569) (2.550,-1.552) (2.500,-1.534) (2.450,-1.515) (2.400,-1.496) (2.350,-1.475) (2.300,-1.454) (2.250,-1.432) (2.200,-1.410) (2.150,-1.386) (2.100,-1.362) (2.050,-1.338) (2.000,-1.312) (1.950,-1.286) (1.900,-1.259) (1.850,-1.232) (1.800,-1.204) (1.750,-1.176) (1.705,-1.150) (1.650,-1.118) (1.600,-1.088) (1.550,-1.058) (1.500,-1.027) (1.457,-1.000) (1.400,-0.964) (1.350,-0.932) (1.300,-0.900) (1.250,-0.867) (1.200,-0.835) (1.150,-0.801) (1.100,-0.768) (1.050,-0.734) (1.000,-0.700) (0.950,-0.666) (0.900,-0.632) (0.854,-0.600) (0.800,-0.563) (0.750,-0.528) (0.700,-0.494) (0.650,-0.459) (0.603,-0.425) (0.550,-0.389) (0.500,-0.353) (0.450,-0.318) (0.400,-0.283) (0.354,-0.250) (0.300,-0.213) (0.250,-0.177) (0.206,-0.144)}
\EQ{(-0.220,-0.155) (-0.275,-0.195) (-0.319,-0.225) (-0.375,-0.266) (-0.425,-0.300) (-0.475,-0.336) (-0.525,-0.371) (-0.567,-0.400) (-0.625,-0.441) (-0.674,-0.475) (-0.725,-0.511) (-0.775,-0.546) (-0.817,-0.575) (-0.875,-0.616) (-0.925,-0.650) (-0.975,-0.684) (-1.025,-0.719) (-1.072,-0.750) (-1.125,-0.787) (-1.175,-0.820) (-1.220,-0.850) (-1.275,-0.887) (-1.325,-0.920) (-1.372,-0.950) (-1.425,-0.985) (-1.475,-1.017) (-1.525,-1.048) (-1.568,-1.075) (-1.625,-1.110) (-1.675,-1.141) (-1.725,-1.171) (-1.774,-1.200) (-1.825,-1.230) (-1.875,-1.259) (-1.925,-1.287) (-1.975,-1.314) (-2.025,-1.342) (-2.075,-1.368) (-2.125,-1.394) (-2.175,-1.420) (-2.225,-1.444) (-2.275,-1.468) (-2.325,-1.492) (-2.375,-1.515) (-2.425,-1.537) (-2.475,-1.558) (-2.525,-1.579) (-2.575,-1.599) (-2.625,-1.618) (-2.675,-1.637) (-2.725,-1.654) (-2.775,-1.671) (-2.825,-1.687) (-2.875,-1.702) (-2.925,-1.717) (-2.975,-1.730) (-3.025,-1.743) (-3.075,-1.755) (-3.125,-1.765) (-3.175,-1.775) (-3.225,-1.784) (-3.275,-1.792) (-3.325,-1.799) (-3.375,-1.805) (-3.425,-1.811) (-3.475,-1.815) (-3.525,-1.818) (-3.575,-1.820) (-3.625,-1.820) (-3.675,-1.820) (-3.725,-1.819) (-3.775,-1.816) (-3.825,-1.813) (-3.875,-1.808) (-3.925,-1.802) (-3.975,-1.795) (-4.025,-1.786) (-4.075,-1.776) (-4.125,-1.765) (-4.175,-1.753) (-4.225,-1.739) (-4.275,-1.723) (-4.325,-1.706) (-4.375,-1.688) (-4.425,-1.667) (-4.475,-1.645) (-4.525,-1.622) (-4.575,-1.596) (-4.625,-1.568) (-4.675,-1.538) (-4.725,-1.506) (-4.770,-1.475) (-4.825,-1.435) (-4.868,-1.400) (-4.925,-1.352) (-4.975,-1.305) (-5.025,-1.255) (-5.075,-1.200) (-5.118,-1.150) (-5.150,-1.109) (-5.194,-1.050) (-5.225,-1.004) (-5.259,-0.950) (-5.289,-0.900) (-5.316,-0.850) (-5.341,-0.800) (-5.365,-0.750) (-5.386,-0.700) (-5.406,-0.650) (-5.424,-0.600) (-5.441,-0.550) (-5.456,-0.500) (-5.469,-0.450) (-5.481,-0.400) (-5.491,-0.350) (-5.500,-0.300) (-5.508,-0.250) (-5.514,-0.200) (-5.519,-0.150) (-5.522,-0.100) (-5.524,-0.050) (-5.525,0.000) (-5.524,0.050) (-5.522,0.100) (-5.519,0.150) (-5.514,0.200) (-5.508,0.250) (-5.500,0.300) (-5.491,0.350) (-5.481,0.400) (-5.469,0.450) (-5.456,0.500) (-5.441,0.550) (-5.424,0.600) (-5.406,0.650) (-5.386,0.700) (-5.365,0.750) (-5.341,0.800) (-5.316,0.850) (-5.289,0.900) (-5.259,0.950) (-5.228,1.000) (-5.194,1.050) (-5.157,1.100) (-5.125,1.141) (-5.075,1.200) (-5.030,1.250) (-4.980,1.300) (-4.927,1.350) (-4.875,1.395) (-4.825,1.435) (-4.775,1.472) (-4.725,1.506) (-4.675,1.538) (-4.625,1.568) (-4.575,1.596) (-4.525,1.622) (-4.475,1.645) (-4.425,1.667) (-4.375,1.688) (-4.325,1.706) (-4.275,1.723) (-4.225,1.739) (-4.175,1.753) (-4.125,1.765) (-4.075,1.776) (-4.025,1.786) (-3.975,1.795) (-3.925,1.802) (-3.875,1.808) (-3.825,1.813) (-3.775,1.816) (-3.725,1.819) (-3.675,1.820) (-3.625,1.820) (-3.575,1.820) (-3.525,1.818) (-3.475,1.815) (-3.425,1.811) (-3.375,1.805) (-3.325,1.799) (-3.275,1.792) (-3.225,1.784) (-3.175,1.775) (-3.125,1.765) (-3.075,1.755) (-3.025,1.743) (-2.975,1.730) (-2.925,1.717) (-2.875,1.702) (-2.825,1.687) (-2.775,1.671) (-2.725,1.654) (-2.675,1.637) (-2.625,1.618) (-2.578,1.600) (-2.525,1.579) (-2.475,1.558) (-2.425,1.537) (-2.375,1.515) (-2.325,1.492) (-2.275,1.468) (-2.225,1.444) (-2.175,1.420) (-2.125,1.394) (-2.075,1.368) (-2.025,1.342) (-1.975,1.314) (-1.925,1.287) (-1.875,1.259) (-1.825,1.230) (-1.775,1.201) (-1.725,1.171) (-1.675,1.141) (-1.625,1.110) (-1.575,1.079) (-1.528,1.050) (-1.475,1.017) (-1.425,0.985) (-1.375,0.952) (-1.325,0.920) (-1.275,0.887) (-1.225,0.854) (-1.183,0.825) (-1.125,0.787) (-1.075,0.753) (-1.025,0.719) (-0.975,0.684) (-0.925,0.650) (-0.875,0.616) (-0.825,0.581) (-0.781,0.550) (-0.725,0.511) (-0.675,0.476) (-0.625,0.441) (-0.575,0.406) (-0.532,0.375) (-0.475,0.336) (-0.425,0.300) (-0.375,0.266) (-0.325,0.230) (-0.284,0.200) (-0.225,0.160) (-0.220,0.155)}
\EQ{(0.206,0.144) (0.248,0.175) (0.300,0.213) (0.350,0.248) (0.400,0.283) (0.450,0.318) (0.496,0.350) (0.550,0.389) (0.600,0.424) (0.650,0.459) (0.700,0.494) (0.745,0.525) (0.800,0.564) (0.850,0.598) (0.900,0.633) (0.950,0.667) (0.998,0.700) (1.050,0.736) (1.100,0.770) (1.145,0.800) (1.200,0.837) (1.250,0.870) (1.296,0.900) (1.350,0.936) (1.400,0.968) (1.449,1.000) (1.500,1.032) (1.550,1.064) (1.600,1.095) (1.649,1.125) (1.700,1.156) (1.750,1.186) (1.800,1.215) (1.850,1.244) (1.900,1.273) (1.949,1.300) (1.995,1.325) (2.041,1.350) (2.088,1.375) (2.137,1.400) (2.186,1.425) (2.237,1.450) (2.289,1.475) (2.343,1.500) (2.398,1.525) (2.450,1.548) (2.500,1.569) (2.550,1.589) (2.600,1.609) (2.650,1.628) (2.700,1.646) (2.750,1.663) (2.800,1.679) (2.850,1.695) (2.900,1.710) (2.950,1.724) (3.000,1.737) (3.050,1.749) (3.100,1.760) (3.150,1.771) (3.200,1.780) (3.250,1.788) (3.300,1.796) (3.350,1.803) (3.400,1.808) (3.450,1.813) (3.500,1.816) (3.550,1.819) (3.600,1.820) (3.650,1.820) (3.700,1.820) (3.750,1.818) (3.800,1.815) (3.850,1.811) (3.900,1.805) (3.950,1.799) (4.000,1.791) (4.050,1.782) (4.100,1.771) (4.150,1.759) (4.200,1.746) (4.250,1.731) (4.300,1.715) (4.350,1.697) (4.400,1.678) (4.450,1.657) (4.500,1.634) (4.550,1.609) (4.600,1.582) (4.650,1.553) (4.696,1.525) (4.750,1.489) (4.800,1.453) (4.850,1.415) (4.898,1.375) (4.950,1.329) (5.000,1.280) (5.050,1.228) (5.097,1.175) (5.138,1.125) (5.175,1.076) (5.211,1.025) (5.244,0.975) (5.274,0.925) (5.300,0.880) (5.325,0.832) (5.350,0.782) (5.375,0.726) (5.396,0.675) (5.415,0.625) (5.433,0.575) (5.448,0.525) (5.462,0.475) (5.475,0.425) (5.486,0.375) (5.496,0.325) (5.504,0.275) (5.511,0.225) (5.516,0.175) (5.521,0.125) (5.523,0.075) (5.525,0.025) (5.524,-0.050) (5.522,-0.100) (5.519,-0.150) (5.514,-0.200) (5.508,-0.250) (5.500,-0.300) (5.491,-0.350) (5.481,-0.400) (5.469,-0.450) (5.456,-0.500) (5.441,-0.550) (5.424,-0.600) (5.406,-0.650) (5.386,-0.700) (5.365,-0.750) (5.341,-0.800) (5.316,-0.850) (5.289,-0.900) (5.259,-0.950) (5.228,-1.000) (5.194,-1.050) (5.157,-1.100) (5.125,-1.141) (5.075,-1.200) (5.030,-1.250) (4.980,-1.300) (4.927,-1.350) (4.875,-1.395) (4.825,-1.435) (4.775,-1.472) (4.725,-1.506) (4.675,-1.538) (4.625,-1.568) (4.575,-1.596) (4.525,-1.622) (4.475,-1.645) (4.425,-1.667) (4.375,-1.688) (4.325,-1.706) (4.275,-1.723) (4.225,-1.739) (4.175,-1.753) (4.125,-1.765) (4.075,-1.776) (4.025,-1.786) (3.975,-1.795) (3.925,-1.802) (3.875,-1.808) (3.825,-1.813) (3.775,-1.816) (3.725,-1.819) (3.675,-1.820) (3.625,-1.820) (3.575,-1.820) (3.525,-1.818) (3.475,-1.815) (3.425,-1.811) (3.375,-1.805) (3.325,-1.799) (3.275,-1.792) (3.225,-1.784) (3.175,-1.775) (3.125,-1.765) (3.075,-1.755) (3.025,-1.743) (2.975,-1.730) (2.925,-1.717) (2.875,-1.702) (2.825,-1.687) (2.775,-1.671) (2.725,-1.654) (2.675,-1.637) (2.625,-1.618) (2.578,-1.600) (2.525,-1.579) (2.475,-1.558) (2.425,-1.537) (2.375,-1.515) (2.325,-1.492) (2.275,-1.468) (2.225,-1.444) (2.175,-1.420) (2.125,-1.394) (2.075,-1.368) (2.025,-1.342) (1.975,-1.314) (1.925,-1.287) (1.875,-1.259) (1.825,-1.230) (1.775,-1.201) (1.725,-1.171) (1.675,-1.141) (1.625,-1.110) (1.575,-1.079) (1.528,-1.050) (1.475,-1.017) (1.425,-0.985) (1.375,-0.952) (1.325,-0.920) (1.275,-0.887) (1.225,-0.854) (1.183,-0.825) (1.125,-0.787) (1.075,-0.753) (1.025,-0.719) (0.975,-0.684) (0.925,-0.650) (0.875,-0.616) (0.825,-0.581) (0.781,-0.550) (0.725,-0.511) (0.675,-0.476) (0.625,-0.441) (0.575,-0.406) (0.532,-0.375) (0.475,-0.336) (0.425,-0.300) (0.375,-0.266) (0.325,-0.230) (0.284,-0.200) (0.225,-0.160) (0.206,-0.144)}
\EQ{(-0.220,-0.155) (-0.275,-0.195) (-0.319,-0.225) (-0.375,-0.266) (-0.425,-0.300) (-0.475,-0.336) (-0.525,-0.371) (-0.567,-0.400) (-0.625,-0.442) (-0.673,-0.475) (-0.725,-0.512) (-0.775,-0.547) (-0.817,-0.575) (-0.875,-0.616) (-0.924,-0.650) (-0.975,-0.686) (-1.025,-0.720) (-1.069,-0.750) (-1.125,-0.788) (-1.175,-0.822) (-1.225,-0.856) (-1.275,-0.890) (-1.325,-0.923) (-1.375,-0.956) (-1.425,-0.989) (-1.475,-1.022) (-1.519,-1.050) (-1.575,-1.086) (-1.625,-1.118) (-1.675,-1.149) (-1.725,-1.180) (-1.775,-1.211) (-1.825,-1.242) (-1.875,-1.271) (-1.924,-1.300) (-1.975,-1.330) (-2.025,-1.359) (-2.075,-1.387) (-2.125,-1.415) (-2.175,-1.442) (-2.225,-1.468) (-2.275,-1.495) (-2.325,-1.520) (-2.375,-1.545) (-2.425,-1.570) (-2.475,-1.593) (-2.525,-1.617) (-2.575,-1.639) (-2.625,-1.661) (-2.675,-1.682) (-2.725,-1.703) (-2.775,-1.723) (-2.825,-1.742) (-2.875,-1.761) (-2.925,-1.778) (-2.975,-1.796) (-3.025,-1.812) (-3.075,-1.827) (-3.125,-1.842) (-3.175,-1.856) (-3.225,-1.869) (-3.275,-1.881) (-3.325,-1.893) (-3.375,-1.904) (-3.425,-1.913) (-3.475,-1.922) (-3.525,-1.930) (-3.575,-1.937) (-3.625,-1.944) (-3.675,-1.949) (-3.725,-1.953) (-3.775,-1.957) (-3.825,-1.959) (-3.875,-1.961) (-3.925,-1.961) (-3.975,-1.960) (-4.025,-1.959) (-4.075,-1.956) (-4.125,-1.952) (-4.175,-1.948) (-4.225,-1.942) (-4.275,-1.934) (-4.325,-1.926) (-4.375,-1.917) (-4.425,-1.906) (-4.475,-1.894) (-4.525,-1.880) (-4.575,-1.866) (-4.624,-1.850) (-4.675,-1.832) (-4.725,-1.813) (-4.775,-1.793) (-4.825,-1.770) (-4.875,-1.746) (-4.925,-1.721) (-4.975,-1.693) (-5.025,-1.664) (-5.075,-1.632) (-5.123,-1.600) (-5.175,-1.562) (-5.223,-1.525) (-5.275,-1.482) (-5.325,-1.438) (-5.365,-1.400) (-5.415,-1.350) (-5.450,-1.312) (-5.500,-1.254) (-5.544,-1.200) (-5.575,-1.158) (-5.616,-1.100) (-5.649,-1.050) (-5.675,-1.007) (-5.708,-0.950) (-5.735,-0.900) (-5.759,-0.850) (-5.782,-0.800) (-5.804,-0.750) (-5.823,-0.700) (-5.842,-0.650) (-5.858,-0.600) (-5.874,-0.550) (-5.887,-0.500) (-5.900,-0.450) (-5.911,-0.400) (-5.920,-0.350) (-5.929,-0.300) (-5.936,-0.250) (-5.941,-0.200) (-5.946,-0.150) (-5.949,-0.100) (-5.951,-0.050) (-5.951,0.000) (-5.951,0.050) (-5.949,0.100) (-5.946,0.150) (-5.941,0.200) (-5.936,0.250) (-5.929,0.300) (-5.920,0.350) (-5.911,0.400) (-5.900,0.450) (-5.887,0.500) (-5.874,0.550) (-5.858,0.600) (-5.842,0.650) (-5.823,0.700) (-5.804,0.750) (-5.782,0.800) (-5.759,0.850) (-5.735,0.900) (-5.708,0.950) (-5.679,1.000) (-5.650,1.048) (-5.616,1.100) (-5.581,1.150) (-5.550,1.192) (-5.504,1.250) (-5.461,1.300) (-5.425,1.339) (-5.375,1.390) (-5.325,1.438) (-5.284,1.475) (-5.225,1.524) (-5.175,1.562) (-5.125,1.599) (-5.075,1.632) (-5.025,1.664) (-4.975,1.693) (-4.925,1.721) (-4.875,1.746) (-4.825,1.770) (-4.775,1.793) (-4.725,1.813) (-4.675,1.832) (-4.625,1.850) (-4.575,1.866) (-4.525,1.880) (-4.475,1.894) (-4.425,1.906) (-4.375,1.917) (-4.325,1.926) (-4.275,1.934) (-4.225,1.942) (-4.175,1.948) (-4.125,1.952) (-4.075,1.956) (-4.025,1.959) (-3.975,1.960) (-3.925,1.961) (-3.875,1.961) (-3.825,1.959) (-3.775,1.957) (-3.725,1.953) (-3.675,1.949) (-3.625,1.944) (-3.575,1.937) (-3.525,1.930) (-3.475,1.922) (-3.425,1.913) (-3.375,1.904) (-3.325,1.893) (-3.275,1.881) (-3.225,1.869) (-3.175,1.856) (-3.125,1.842) (-3.075,1.827) (-3.025,1.812) (-2.975,1.796) (-2.925,1.778) (-2.875,1.761) (-2.825,1.742) (-2.775,1.723) (-2.725,1.703) (-2.675,1.682) (-2.625,1.661) (-2.575,1.639) (-2.525,1.617) (-2.475,1.593) (-2.425,1.570) (-2.375,1.545) (-2.325,1.520) (-2.275,1.495) (-2.225,1.468) (-2.175,1.442) (-2.125,1.415) (-2.075,1.387) (-2.025,1.359) (-1.975,1.330) (-1.925,1.301) (-1.881,1.275) (-1.825,1.242) (-1.775,1.211) (-1.725,1.180) (-1.676,1.150) (-1.625,1.118) (-1.575,1.086) (-1.525,1.054) (-1.480,1.025) (-1.425,0.989) (-1.375,0.956) (-1.328,0.925) (-1.275,0.890) (-1.225,0.856) (-1.179,0.825) (-1.125,0.788) (-1.075,0.754) (-1.033,0.725) (-0.975,0.686) (-0.925,0.651) (-0.875,0.616) (-0.825,0.581) (-0.781,0.550) (-0.725,0.512) (-0.675,0.476) (-0.625,0.442) (-0.575,0.406) (-0.531,0.375) (-0.475,0.336) (-0.425,0.300) (-0.375,0.266) (-0.325,0.230) (-0.284,0.200) (-0.225,0.160) (-0.220,0.155)}
\EQ{(0.206,0.144) (0.248,0.175) (0.300,0.213) (0.350,0.248) (0.400,0.283) (0.450,0.318) (0.496,0.350) (0.550,0.389) (0.600,0.424) (0.650,0.459) (0.700,0.494) (0.745,0.525) (0.800,0.564) (0.850,0.599) (0.900,0.634) (0.950,0.668) (0.997,0.700) (1.050,0.737) (1.100,0.771) (1.143,0.800) (1.200,0.839) (1.250,0.873) (1.300,0.906) (1.350,0.940) (1.400,0.973) (1.442,1.000) (1.500,1.038) (1.550,1.070) (1.597,1.100) (1.650,1.134) (1.700,1.165) (1.750,1.196) (1.798,1.225) (1.850,1.257) (1.900,1.286) (1.950,1.316) (2.000,1.344) (2.050,1.373) (2.099,1.400) (2.144,1.425) (2.191,1.450) (2.238,1.475) (2.286,1.500) (2.335,1.525) (2.385,1.550) (2.437,1.575) (2.489,1.600) (2.544,1.625) (2.600,1.650) (2.650,1.672) (2.700,1.693) (2.750,1.713) (2.800,1.733) (2.850,1.751) (2.900,1.770) (2.950,1.787) (3.000,1.804) (3.050,1.820) (3.100,1.835) (3.150,1.849) (3.200,1.863) (3.249,1.875) (3.300,1.887) (3.350,1.898) (3.400,1.909) (3.450,1.918) (3.500,1.926) (3.550,1.934) (3.600,1.941) (3.650,1.946) (3.700,1.951) (3.750,1.955) (3.800,1.958) (3.850,1.960) (3.900,1.961) (3.950,1.961) (4.000,1.960) (4.050,1.958) (4.100,1.954) (4.150,1.950) (4.200,1.945) (4.250,1.938) (4.300,1.930) (4.350,1.922) (4.400,1.911) (4.450,1.900) (4.500,1.887) (4.550,1.873) (4.600,1.858) (4.650,1.841) (4.700,1.823) (4.750,1.803) (4.800,1.782) (4.850,1.759) (4.900,1.734) (4.950,1.707) (5.000,1.679) (5.047,1.650) (5.100,1.616) (5.150,1.581) (5.191,1.550) (5.250,1.503) (5.300,1.461) (5.339,1.425) (5.390,1.375) (5.438,1.325) (5.475,1.284) (5.524,1.225) (5.563,1.175) (5.599,1.125) (5.633,1.075) (5.664,1.025) (5.694,0.975) (5.721,0.925) (5.747,0.875) (5.771,0.825) (5.793,0.775) (5.814,0.725) (5.833,0.675) (5.850,0.625) (5.866,0.575) (5.881,0.525) (5.894,0.475) (5.905,0.425) (5.916,0.375) (5.925,0.325) (5.932,0.275) (5.939,0.225) (5.944,0.175) (5.947,0.125) (5.950,0.075) (5.951,0.025) (5.951,-0.050) (5.949,-0.100) (5.946,-0.150) (5.941,-0.200) (5.936,-0.250) (5.929,-0.300) (5.920,-0.350) (5.911,-0.400) (5.900,-0.450) (5.887,-0.500) (5.874,-0.550) (5.858,-0.600) (5.842,-0.650) (5.823,-0.700) (5.804,-0.750) (5.782,-0.800) (5.759,-0.850) (5.735,-0.900) (5.708,-0.950) (5.679,-1.000) (5.650,-1.048) (5.616,-1.100) (5.581,-1.150) (5.550,-1.192) (5.504,-1.250) (5.461,-1.300) (5.425,-1.339) (5.375,-1.390) (5.325,-1.438) (5.284,-1.475) (5.225,-1.524) (5.175,-1.562) (5.125,-1.599) (5.075,-1.632) (5.025,-1.664) (4.975,-1.693) (4.925,-1.721) (4.875,-1.746) (4.825,-1.770) (4.775,-1.793) (4.725,-1.813) (4.675,-1.832) (4.625,-1.850) (4.575,-1.866) (4.525,-1.880) (4.475,-1.894) (4.425,-1.906) (4.375,-1.917) (4.325,-1.926) (4.275,-1.934) (4.225,-1.942) (4.175,-1.948) (4.125,-1.952) (4.075,-1.956) (4.025,-1.959) (3.975,-1.960) (3.925,-1.961) (3.875,-1.961) (3.825,-1.959) (3.775,-1.957) (3.725,-1.953) (3.675,-1.949) (3.625,-1.944) (3.575,-1.937) (3.525,-1.930) (3.475,-1.922) (3.425,-1.913) (3.375,-1.904) (3.325,-1.893) (3.275,-1.881) (3.225,-1.869) (3.175,-1.856) (3.125,-1.842) (3.075,-1.827) (3.025,-1.812) (2.975,-1.796) (2.925,-1.778) (2.875,-1.761) (2.825,-1.742) (2.775,-1.723) (2.725,-1.703) (2.675,-1.682) (2.625,-1.661) (2.575,-1.639) (2.525,-1.617) (2.475,-1.593) (2.425,-1.570) (2.375,-1.545) (2.325,-1.520) (2.275,-1.495) (2.225,-1.468) (2.175,-1.442) (2.125,-1.415) (2.075,-1.387) (2.025,-1.359) (1.975,-1.330) (1.925,-1.301) (1.881,-1.275) (1.825,-1.242) (1.775,-1.211) (1.725,-1.180) (1.676,-1.150) (1.625,-1.118) (1.575,-1.086) (1.525,-1.054) (1.480,-1.025) (1.425,-0.989) (1.375,-0.956) (1.328,-0.925) (1.275,-0.890) (1.225,-0.856) (1.179,-0.825) (1.125,-0.788) (1.075,-0.754) (1.033,-0.725) (0.975,-0.686) (0.925,-0.651) (0.875,-0.616) (0.825,-0.581) (0.781,-0.550) (0.725,-0.512) (0.675,-0.476) (0.625,-0.442) (0.575,-0.406) (0.531,-0.375) (0.475,-0.336) (0.425,-0.300) (0.375,-0.266) (0.325,-0.230) (0.284,-0.200) (0.225,-0.160) (0.206,-0.144)}
\EQ{(-0.220,-0.155) (-0.275,-0.195) (-0.319,-0.225) (-0.375,-0.266) (-0.425,-0.300) (-0.475,-0.336) (-0.525,-0.371) (-0.567,-0.400) (-0.625,-0.442) (-0.673,-0.475) (-0.725,-0.512) (-0.775,-0.547) (-0.825,-0.582) (-0.875,-0.617) (-0.923,-0.650) (-0.975,-0.687) (-1.025,-0.721) (-1.067,-0.750) (-1.125,-0.790) (-1.175,-0.824) (-1.225,-0.859) (-1.275,-0.893) (-1.323,-0.925) (-1.375,-0.960) (-1.425,-0.994) (-1.472,-1.025) (-1.525,-1.060) (-1.575,-1.093) (-1.624,-1.125) (-1.675,-1.158) (-1.725,-1.190) (-1.775,-1.222) (-1.820,-1.250) (-1.875,-1.285) (-1.925,-1.316) (-1.975,-1.346) (-2.023,-1.375) (-2.075,-1.406) (-2.125,-1.436) (-2.175,-1.465) (-2.225,-1.493) (-2.275,-1.521) (-2.325,-1.549) (-2.372,-1.575) (-2.419,-1.600) (-2.466,-1.625) (-2.514,-1.650) (-2.563,-1.675) (-2.613,-1.700) (-2.664,-1.725) (-2.716,-1.750) (-2.770,-1.775) (-2.825,-1.800) (-2.875,-1.822) (-2.925,-1.844) (-2.975,-1.865) (-3.025,-1.885) (-3.075,-1.904) (-3.125,-1.923) (-3.175,-1.942) (-3.225,-1.959) (-3.275,-1.976) (-3.325,-1.993) (-3.375,-2.008) (-3.425,-2.023) (-3.475,-2.037) (-3.525,-2.051) (-3.575,-2.064) (-3.625,-2.076) (-3.675,-2.087) (-3.725,-2.097) (-3.775,-2.107) (-3.825,-2.116) (-3.875,-2.124) (-3.925,-2.131) (-3.975,-2.138) (-4.025,-2.143) (-4.075,-2.148) (-4.125,-2.152) (-4.175,-2.155) (-4.225,-2.157) (-4.275,-2.158) (-4.325,-2.158) (-4.375,-2.158) (-4.425,-2.156) (-4.475,-2.154) (-4.525,-2.150) (-4.575,-2.146) (-4.625,-2.140) (-4.675,-2.134) (-4.725,-2.126) (-4.775,-2.117) (-4.825,-2.108) (-4.875,-2.097) (-4.925,-2.085) (-4.975,-2.071) (-5.025,-2.057) (-5.075,-2.041) (-5.125,-2.024) (-5.175,-2.006) (-5.225,-1.986) (-5.275,-1.965) (-5.325,-1.942) (-5.375,-1.918) (-5.425,-1.892) (-5.475,-1.864) (-5.525,-1.835) (-5.575,-1.804) (-5.618,-1.775) (-5.675,-1.735) (-5.722,-1.700) (-5.775,-1.658) (-5.814,-1.625) (-5.869,-1.575) (-5.921,-1.525) (-5.970,-1.475) (-6.016,-1.425) (-6.050,-1.385) (-6.098,-1.325) (-6.136,-1.275) (-6.172,-1.225) (-6.200,-1.183) (-6.236,-1.125) (-6.266,-1.075) (-6.294,-1.025) (-6.320,-0.975) (-6.344,-0.925) (-6.367,-0.875) (-6.388,-0.825) (-6.408,-0.775) (-6.426,-0.725) (-6.443,-0.675) (-6.459,-0.625) (-6.473,-0.575) (-6.486,-0.525) (-6.498,-0.475) (-6.509,-0.425) (-6.518,-0.375) (-6.526,-0.325) (-6.533,-0.275) (-6.539,-0.225) (-6.543,-0.175) (-6.547,-0.125) (-6.549,-0.075) (-6.550,-0.025) (-6.550,0.037) (-6.548,0.100) (-6.545,0.150) (-6.541,0.200) (-6.536,0.250) (-6.530,0.300) (-6.522,0.350) (-6.513,0.400) (-6.503,0.450) (-6.492,0.500) (-6.480,0.550) (-6.466,0.600) (-6.451,0.650) (-6.435,0.700) (-6.417,0.750) (-6.398,0.800) (-6.378,0.850) (-6.356,0.900) (-6.332,0.950) (-6.307,1.000) (-6.280,1.050) (-6.251,1.100) (-6.225,1.143) (-6.189,1.200) (-6.154,1.250) (-6.118,1.300) (-6.079,1.350) (-6.037,1.400) (-6.000,1.442) (-5.950,1.496) (-5.900,1.546) (-5.850,1.593) (-5.800,1.637) (-5.754,1.675) (-5.700,1.717) (-5.654,1.750) (-5.600,1.787) (-5.550,1.819) (-5.500,1.850) (-5.456,1.875) (-5.409,1.900) (-5.360,1.925) (-5.307,1.950) (-5.251,1.975) (-5.200,1.996) (-5.150,2.015) (-5.100,2.033) (-5.050,2.049) (-5.000,2.064) (-4.950,2.078) (-4.900,2.091) (-4.850,2.102) (-4.800,2.113) (-4.750,2.122) (-4.700,2.130) (-4.650,2.137) (-4.600,2.143) (-4.550,2.148) (-4.500,2.152) (-4.450,2.155) (-4.400,2.157) (-4.350,2.158) (-4.300,2.158) (-4.250,2.158) (-4.200,2.156) (-4.150,2.153) (-4.100,2.150) (-4.050,2.146) (-4.000,2.140) (-3.950,2.134) (-3.900,2.128) (-3.850,2.120) (-3.800,2.111) (-3.750,2.102) (-3.700,2.092) (-3.650,2.081) (-3.600,2.070) (-3.550,2.057) (-3.500,2.044) (-3.450,2.030) (-3.400,2.016) (-3.350,2.001) (-3.300,1.985) (-3.250,1.968) (-3.200,1.951) (-3.150,1.933) (-3.100,1.914) (-3.050,1.895) (-3.001,1.875) (-2.950,1.854) (-2.900,1.833) (-2.850,1.811) (-2.800,1.789) (-2.750,1.766) (-2.700,1.742) (-2.650,1.718) (-2.600,1.694) (-2.550,1.669) (-2.500,1.643) (-2.450,1.617) (-2.400,1.590) (-2.350,1.563) (-2.300,1.535) (-2.250,1.507) (-2.200,1.479) (-2.150,1.450) (-2.100,1.421) (-2.050,1.391) (-2.000,1.361) (-1.950,1.331) (-1.900,1.300) (-1.850,1.269) (-1.800,1.238) (-1.750,1.206) (-1.702,1.175) (-1.650,1.142) (-1.600,1.109) (-1.550,1.076) (-1.500,1.044) (-1.450,1.010) (-1.400,0.977) (-1.350,0.943) (-1.300,0.910) (-1.250,0.876) (-1.200,0.842) (-1.150,0.807) (-1.103,0.775) (-1.050,0.739) (-1.000,0.704) (-0.950,0.669) (-0.900,0.634) (-0.851,0.600) (-0.800,0.565) (-0.750,0.529) (-0.709,0.500) (-0.650,0.459) (-0.602,0.425) (-0.550,0.389) (-0.500,0.354) (-0.450,0.319) (-0.400,0.283) (-0.354,0.250) (-0.300,0.213) (-0.250,0.177) (-0.220,0.155)}
\EQ{(0.206,0.144) (0.248,0.175) (0.300,0.213) (0.350,0.248) (0.400,0.283) (0.450,0.319) (0.496,0.350) (0.550,0.389) (0.600,0.424) (0.650,0.459) (0.700,0.494) (0.744,0.525) (0.800,0.565) (0.850,0.599) (0.900,0.634) (0.950,0.669) (0.995,0.700) (1.050,0.739) (1.100,0.773) (1.150,0.807) (1.200,0.842) (1.249,0.875) (1.300,0.910) (1.350,0.943) (1.397,0.975) (1.450,1.010) (1.500,1.044) (1.548,1.075) (1.600,1.109) (1.650,1.142) (1.700,1.174) (1.750,1.206) (1.800,1.238) (1.850,1.269) (1.900,1.300) (1.950,1.331) (2.000,1.361) (2.050,1.391) (2.100,1.421) (2.150,1.450) (2.193,1.475) (2.237,1.500) (2.282,1.525) (2.327,1.550) (2.372,1.575) (2.419,1.600) (2.466,1.625) (2.514,1.650) (2.563,1.675) (2.613,1.700) (2.664,1.725) (2.716,1.750) (2.770,1.775) (2.825,1.800) (2.875,1.822) (2.925,1.844) (2.975,1.865) (3.025,1.885) (3.075,1.904) (3.125,1.923) (3.175,1.942) (3.225,1.959) (3.275,1.976) (3.325,1.993) (3.375,2.008) (3.425,2.023) (3.475,2.037) (3.525,2.051) (3.575,2.064) (3.625,2.076) (3.675,2.087) (3.725,2.097) (3.775,2.107) (3.825,2.116) (3.875,2.124) (3.925,2.131) (3.975,2.138) (4.025,2.143) (4.075,2.148) (4.125,2.152) (4.175,2.155) (4.225,2.157) (4.275,2.158) (4.325,2.158) (4.375,2.158) (4.425,2.156) (4.475,2.154) (4.525,2.150) (4.575,2.146) (4.625,2.140) (4.675,2.134) (4.725,2.126) (4.775,2.117) (4.825,2.108) (4.875,2.097) (4.925,2.085) (4.975,2.071) (5.025,2.057) (5.075,2.041) (5.125,2.024) (5.175,2.006) (5.225,1.986) (5.275,1.965) (5.325,1.942) (5.375,1.918) (5.425,1.892) (5.475,1.864) (5.525,1.835) (5.575,1.804) (5.618,1.775) (5.675,1.735) (5.722,1.700) (5.775,1.658) (5.814,1.625) (5.869,1.575) (5.921,1.525) (5.970,1.475) (6.016,1.425) (6.050,1.385) (6.098,1.325) (6.136,1.275) (6.172,1.225) (6.200,1.183) (6.236,1.125) (6.266,1.075) (6.294,1.025) (6.320,0.975) (6.344,0.925) (6.367,0.875) (6.388,0.825) (6.408,0.775) (6.426,0.725) (6.443,0.675) (6.459,0.625) (6.473,0.575) (6.486,0.525) (6.498,0.475) (6.509,0.425) (6.518,0.375) (6.526,0.325) (6.533,0.275) (6.539,0.225) (6.543,0.175) (6.547,0.125) (6.549,0.075) (6.550,0.025) (6.550,-0.037) (6.548,-0.100) (6.545,-0.150) (6.541,-0.200) (6.536,-0.250) (6.530,-0.300) (6.522,-0.350) (6.513,-0.400) (6.503,-0.450) (6.492,-0.500) (6.480,-0.550) (6.466,-0.600) (6.451,-0.650) (6.435,-0.700) (6.417,-0.750) (6.398,-0.800) (6.378,-0.850) (6.356,-0.900) (6.332,-0.950) (6.307,-1.000) (6.280,-1.050) (6.251,-1.100) (6.225,-1.143) (6.189,-1.200) (6.154,-1.250) (6.118,-1.300) (6.079,-1.350) (6.037,-1.400) (6.000,-1.442) (5.950,-1.496) (5.900,-1.546) (5.850,-1.593) (5.800,-1.637) (5.754,-1.675) (5.700,-1.717) (5.654,-1.750) (5.600,-1.787) (5.550,-1.819) (5.500,-1.850) (5.456,-1.875) (5.409,-1.900) (5.360,-1.925) (5.307,-1.950) (5.251,-1.975) (5.200,-1.996) (5.150,-2.015) (5.100,-2.033) (5.050,-2.049) (5.000,-2.064) (4.950,-2.078) (4.900,-2.091) (4.850,-2.102) (4.800,-2.113) (4.750,-2.122) (4.700,-2.130) (4.650,-2.137) (4.600,-2.143) (4.550,-2.148) (4.500,-2.152) (4.450,-2.155) (4.400,-2.157) (4.350,-2.158) (4.300,-2.158) (4.250,-2.158) (4.200,-2.156) (4.150,-2.153) (4.100,-2.150) (4.050,-2.146) (4.000,-2.140) (3.950,-2.134) (3.900,-2.128) (3.850,-2.120) (3.800,-2.111) (3.750,-2.102) (3.700,-2.092) (3.650,-2.081) (3.600,-2.070) (3.550,-2.057) (3.500,-2.044) (3.450,-2.030) (3.400,-2.016) (3.350,-2.001) (3.300,-1.985) (3.250,-1.968) (3.200,-1.951) (3.150,-1.933) (3.100,-1.914) (3.050,-1.895) (3.001,-1.875) (2.950,-1.854) (2.900,-1.833) (2.850,-1.811) (2.800,-1.789) (2.750,-1.766) (2.700,-1.742) (2.650,-1.718) (2.600,-1.694) (2.550,-1.669) (2.500,-1.643) (2.450,-1.617) (2.400,-1.590) (2.350,-1.563) (2.300,-1.535) (2.250,-1.507) (2.200,-1.479) (2.150,-1.450) (2.100,-1.421) (2.050,-1.391) (2.000,-1.361) (1.950,-1.331) (1.900,-1.300) (1.850,-1.269) (1.800,-1.238) (1.750,-1.206) (1.702,-1.175) (1.650,-1.142) (1.600,-1.109) (1.550,-1.076) (1.500,-1.044) (1.450,-1.010) (1.400,-0.977) (1.350,-0.943) (1.300,-0.910) (1.250,-0.876) (1.200,-0.842) (1.150,-0.807) (1.103,-0.775) (1.050,-0.739) (1.000,-0.704) (0.950,-0.669) (0.900,-0.634) (0.851,-0.600) (0.800,-0.565) (0.750,-0.529) (0.709,-0.500) (0.650,-0.459) (0.602,-0.425) (0.550,-0.389) (0.500,-0.354) (0.450,-0.319) (0.400,-0.283) (0.354,-0.250) (0.300,-0.213) (0.250,-0.177) (0.206,-0.144)}
\EQ{(-0.220,-0.155) (-0.275,-0.195) (-0.319,-0.225) (-0.375,-0.266) (-0.424,-0.300) (-0.475,-0.336) (-0.525,-0.371) (-0.567,-0.400) (-0.625,-0.442) (-0.673,-0.475) (-0.725,-0.512) (-0.775,-0.547) (-0.825,-0.583) (-0.875,-0.618) (-0.922,-0.650) (-0.975,-0.688) (-1.025,-0.722) (-1.075,-0.757) (-1.125,-0.792) (-1.173,-0.825) (-1.225,-0.861) (-1.275,-0.896) (-1.318,-0.925) (-1.375,-0.964) (-1.425,-0.998) (-1.475,-1.032) (-1.525,-1.066) (-1.575,-1.100) (-1.625,-1.133) (-1.675,-1.167) (-1.725,-1.200) (-1.775,-1.233) (-1.825,-1.266) (-1.875,-1.298) (-1.925,-1.330) (-1.975,-1.362) (-2.025,-1.394) (-2.074,-1.425) (-2.125,-1.457) (-2.175,-1.488) (-2.225,-1.519) (-2.275,-1.549) (-2.325,-1.580) (-2.375,-1.609) (-2.425,-1.639) (-2.475,-1.668) (-2.525,-1.697) (-2.575,-1.725) (-2.619,-1.750) (-2.665,-1.775) (-2.710,-1.800) (-2.757,-1.825) (-2.803,-1.850) (-2.850,-1.874) (-2.900,-1.900) (-2.949,-1.925) (-2.999,-1.950) (-3.050,-1.975) (-3.100,-1.999) (-3.150,-2.022) (-3.200,-2.045) (-3.250,-2.068) (-3.300,-2.090) (-3.350,-2.111) (-3.400,-2.132) (-3.450,-2.152) (-3.500,-2.172) (-3.550,-2.191) (-3.600,-2.210) (-3.650,-2.228) (-3.700,-2.246) (-3.750,-2.263) (-3.800,-2.279) (-3.850,-2.295) (-3.900,-2.310) (-3.950,-2.324) (-4.000,-2.338) (-4.050,-2.352) (-4.100,-2.364) (-4.150,-2.376) (-4.200,-2.387) (-4.250,-2.398) (-4.300,-2.408) (-4.350,-2.417) (-4.400,-2.425) (-4.450,-2.433) (-4.500,-2.440) (-4.550,-2.447) (-4.600,-2.452) (-4.650,-2.457) (-4.700,-2.462) (-4.750,-2.465) (-4.800,-2.468) (-4.850,-2.469) (-4.900,-2.470) (-4.950,-2.471) (-5.000,-2.470) (-5.050,-2.469) (-5.100,-2.467) (-5.150,-2.464) (-5.200,-2.460) (-5.250,-2.455) (-5.300,-2.449) (-5.350,-2.443) (-5.400,-2.435) (-5.450,-2.427) (-5.500,-2.417) (-5.550,-2.407) (-5.600,-2.395) (-5.650,-2.383) (-5.700,-2.370) (-5.750,-2.355) (-5.800,-2.339) (-5.850,-2.323) (-5.900,-2.305) (-5.950,-2.286) (-6.000,-2.265) (-6.050,-2.244) (-6.100,-2.221) (-6.150,-2.196) (-6.200,-2.171) (-6.250,-2.144) (-6.300,-2.115) (-6.350,-2.085) (-6.400,-2.053) (-6.450,-2.019) (-6.500,-1.983) (-6.545,-1.950) (-6.600,-1.906) (-6.650,-1.864) (-6.695,-1.825) (-6.748,-1.775) (-6.799,-1.725) (-6.846,-1.675) (-6.892,-1.625) (-6.925,-1.586) (-6.974,-1.525) (-7.013,-1.475) (-7.049,-1.425) (-7.083,-1.375) (-7.116,-1.325) (-7.147,-1.275) (-7.175,-1.226) (-7.200,-1.181) (-7.225,-1.134) (-7.250,-1.084) (-7.275,-1.031) (-7.300,-0.975) (-7.320,-0.925) (-7.340,-0.875) (-7.358,-0.825) (-7.375,-0.775) (-7.391,-0.725) (-7.405,-0.675) (-7.419,-0.625) (-7.431,-0.575) (-7.443,-0.525) (-7.453,-0.475) (-7.462,-0.425) (-7.470,-0.375) (-7.477,-0.325) (-7.483,-0.275) (-7.488,-0.225) (-7.492,-0.175) (-7.495,-0.125) (-7.497,-0.075) (-7.498,-0.025) (-7.498,0.050) (-7.496,0.100) (-7.494,0.150) (-7.490,0.200) (-7.486,0.250) (-7.480,0.300) (-7.474,0.350) (-7.466,0.400) (-7.457,0.450) (-7.448,0.500) (-7.437,0.550) (-7.425,0.600) (-7.412,0.650) (-7.398,0.700) (-7.383,0.750) (-7.367,0.800) (-7.349,0.850) (-7.330,0.900) (-7.310,0.950) (-7.289,1.000) (-7.266,1.050) (-7.242,1.100) (-7.217,1.150) (-7.190,1.200) (-7.161,1.250) (-7.131,1.300) (-7.100,1.350) (-7.066,1.400) (-7.031,1.450) (-7.000,1.492) (-6.955,1.550) (-6.913,1.600) (-6.875,1.644) (-6.825,1.698) (-6.775,1.749) (-6.725,1.797) (-6.675,1.843) (-6.625,1.886) (-6.577,1.925) (-6.525,1.965) (-6.477,2.000) (-6.425,2.036) (-6.375,2.069) (-6.325,2.100) (-6.275,2.129) (-6.225,2.157) (-6.175,2.184) (-6.125,2.209) (-6.075,2.232) (-6.025,2.255) (-5.977,2.275) (-5.925,2.295) (-5.875,2.314) (-5.825,2.331) (-5.775,2.347) (-5.725,2.362) (-5.675,2.376) (-5.625,2.389) (-5.575,2.401) (-5.525,2.412) (-5.475,2.422) (-5.425,2.431) (-5.375,2.439) (-5.325,2.446) (-5.275,2.452) (-5.225,2.457) (-5.175,2.462) (-5.125,2.465) (-5.075,2.468) (-5.025,2.470) (-4.975,2.471) (-4.925,2.471) (-4.875,2.470) (-4.825,2.469) (-4.775,2.466) (-4.725,2.463) (-4.675,2.460) (-4.625,2.455) (-4.575,2.450) (-4.525,2.444) (-4.475,2.437) (-4.425,2.430) (-4.375,2.421) (-4.325,2.412) (-4.275,2.403) (-4.225,2.393) (-4.175,2.382) (-4.125,2.370) (-4.075,2.358) (-4.025,2.345) (-3.975,2.331) (-3.925,2.317) (-3.875,2.302) (-3.825,2.287) (-3.775,2.271) (-3.725,2.254) (-3.675,2.237) (-3.625,2.219) (-3.575,2.201) (-3.525,2.182) (-3.475,2.162) (-3.425,2.142) (-3.375,2.122) (-3.325,2.100) (-3.275,2.079) (-3.225,2.057) (-3.175,2.034) (-3.125,2.011) (-3.075,1.987) (-3.025,1.963) (-2.975,1.938) (-2.925,1.913) (-2.875,1.887) (-2.825,1.861) (-2.775,1.835) (-2.725,1.808) (-2.675,1.781) (-2.625,1.753) (-2.575,1.725) (-2.531,1.700) (-2.487,1.675) (-2.425,1.639) (-2.375,1.609) (-2.325,1.580) (-2.276,1.550) (-2.225,1.519) (-2.175,1.488) (-2.125,1.457) (-2.075,1.426) (-2.025,1.394) (-1.975,1.362) (-1.925,1.330) (-1.878,1.300) (-1.825,1.266) (-1.775,1.233) (-1.725,1.200) (-1.675,1.167) (-1.625,1.133) (-1.576,1.100) (-1.525,1.066) (-1.475,1.032) (-1.428,1.000) (-1.375,0.964) (-1.325,0.930) (-1.282,0.900) (-1.225,0.861) (-1.175,0.827) (-1.125,0.792) (-1.075,0.757) (-1.029,0.725) (-0.975,0.688) (-0.925,0.652) (-0.875,0.618) (-0.825,0.583) (-0.779,0.550) (-0.725,0.512) (-0.675,0.477) (-0.625,0.442) (-0.575,0.407) (-0.531,0.375) (-0.475,0.336) (-0.425,0.300) (-0.375,0.266) (-0.325,0.230) (-0.284,0.200) (-0.225,0.160) (-0.220,0.155)}
\EQ{(0.206,0.144) (0.248,0.175) (0.300,0.213) (0.350,0.248) (0.400,0.283) (0.450,0.319) (0.496,0.350) (0.550,0.389) (0.600,0.424) (0.650,0.460) (0.700,0.495) (0.744,0.525) (0.800,0.565) (0.850,0.600) (0.900,0.635) (0.950,0.670) (0.993,0.700) (1.050,0.740) (1.100,0.774) (1.150,0.809) (1.200,0.844) (1.245,0.875) (1.300,0.913) (1.350,0.947) (1.400,0.981) (1.450,1.015) (1.500,1.049) (1.550,1.083) (1.600,1.117) (1.650,1.150) (1.700,1.183) (1.750,1.216) (1.800,1.249) (1.850,1.282) (1.900,1.314) (1.950,1.346) (1.995,1.375) (2.050,1.410) (2.100,1.442) (2.150,1.473) (2.194,1.500) (2.250,1.534) (2.300,1.565) (2.350,1.594) (2.400,1.624) (2.444,1.650) (2.500,1.682) (2.550,1.711) (2.600,1.739) (2.650,1.767) (2.700,1.795) (2.750,1.822) (2.800,1.848) (2.850,1.874) (2.900,1.900) (2.949,1.925) (2.999,1.950) (3.050,1.975) (3.100,1.999) (3.150,2.022) (3.200,2.045) (3.250,2.068) (3.300,2.090) (3.350,2.111) (3.400,2.132) (3.450,2.152) (3.500,2.172) (3.550,2.191) (3.600,2.210) (3.650,2.228) (3.700,2.246) (3.750,2.263) (3.800,2.279) (3.850,2.295) (3.900,2.310) (3.950,2.324) (4.000,2.338) (4.050,2.352) (4.100,2.364) (4.150,2.376) (4.200,2.387) (4.250,2.398) (4.300,2.408) (4.350,2.417) (4.400,2.425) (4.450,2.433) (4.500,2.440) (4.550,2.447) (4.600,2.452) (4.650,2.457) (4.700,2.462) (4.750,2.465) (4.800,2.468) (4.850,2.469) (4.900,2.470) (4.950,2.471) (5.000,2.470) (5.050,2.469) (5.100,2.467) (5.150,2.464) (5.200,2.460) (5.250,2.455) (5.300,2.449) (5.350,2.443) (5.400,2.435) (5.450,2.427) (5.500,2.417) (5.550,2.407) (5.600,2.395) (5.650,2.383) (5.700,2.370) (5.750,2.355) (5.800,2.339) (5.850,2.323) (5.900,2.305) (5.950,2.286) (6.000,2.265) (6.050,2.244) (6.100,2.221) (6.150,2.196) (6.200,2.171) (6.250,2.144) (6.300,2.115) (6.350,2.085) (6.400,2.053) (6.450,2.019) (6.500,1.983) (6.545,1.950) (6.600,1.906) (6.650,1.864) (6.695,1.825) (6.748,1.775) (6.799,1.725) (6.846,1.675) (6.892,1.625) (6.925,1.586) (6.974,1.525) (7.013,1.475) (7.049,1.425) (7.083,1.375) (7.116,1.325) (7.147,1.275) (7.175,1.226) (7.200,1.181) (7.225,1.134) (7.250,1.084) (7.275,1.031) (7.300,0.975) (7.320,0.925) (7.340,0.875) (7.358,0.825) (7.375,0.775) (7.391,0.725) (7.405,0.675) (7.419,0.625) (7.431,0.575) (7.443,0.525) (7.453,0.475) (7.462,0.425) (7.470,0.375) (7.477,0.325) (7.483,0.275) (7.488,0.225) (7.492,0.175) (7.495,0.125) (7.497,0.075) (7.498,0.025) (7.498,-0.050) (7.496,-0.100) (7.494,-0.150) (7.490,-0.200) (7.486,-0.250) (7.480,-0.300) (7.474,-0.350) (7.466,-0.400) (7.457,-0.450) (7.448,-0.500) (7.437,-0.550) (7.425,-0.600) (7.412,-0.650) (7.398,-0.700) (7.383,-0.750) (7.367,-0.800) (7.349,-0.850) (7.330,-0.900) (7.310,-0.950) (7.289,-1.000) (7.266,-1.050) (7.242,-1.100) (7.217,-1.150) (7.190,-1.200) (7.161,-1.250) (7.131,-1.300) (7.100,-1.350) (7.066,-1.400) (7.031,-1.450) (7.000,-1.492) (6.955,-1.550) (6.913,-1.600) (6.875,-1.644) (6.825,-1.698) (6.775,-1.749) (6.725,-1.797) (6.675,-1.843) (6.625,-1.886) (6.577,-1.925) (6.525,-1.965) (6.477,-2.000) (6.425,-2.036) (6.375,-2.069) (6.325,-2.100) (6.275,-2.129) (6.225,-2.157) (6.175,-2.184) (6.125,-2.209) (6.075,-2.232) (6.025,-2.255) (5.977,-2.275) (5.925,-2.295) (5.875,-2.314) (5.825,-2.331) (5.775,-2.347) (5.725,-2.362) (5.675,-2.376) (5.625,-2.389) (5.575,-2.401) (5.525,-2.412) (5.475,-2.422) (5.425,-2.431) (5.375,-2.439) (5.325,-2.446) (5.275,-2.452) (5.225,-2.457) (5.175,-2.462) (5.125,-2.465) (5.075,-2.468) (5.025,-2.470) (4.975,-2.471) (4.925,-2.471) (4.875,-2.470) (4.825,-2.469) (4.775,-2.466) (4.725,-2.463) (4.675,-2.460) (4.625,-2.455) (4.575,-2.450) (4.525,-2.444) (4.475,-2.437) (4.425,-2.430) (4.375,-2.421) (4.325,-2.412) (4.275,-2.403) (4.225,-2.393) (4.175,-2.382) (4.125,-2.370) (4.075,-2.358) (4.025,-2.345) (3.975,-2.331) (3.925,-2.317) (3.875,-2.302) (3.825,-2.287) (3.775,-2.271) (3.725,-2.254) (3.675,-2.237) (3.625,-2.219) (3.575,-2.201) (3.525,-2.182) (3.475,-2.162) (3.425,-2.142) (3.375,-2.122) (3.325,-2.100) (3.275,-2.079) (3.225,-2.057) (3.175,-2.034) (3.125,-2.011) (3.075,-1.987) (3.025,-1.963) (2.975,-1.938) (2.925,-1.913) (2.875,-1.887) (2.825,-1.861) (2.775,-1.835) (2.725,-1.808) (2.675,-1.781) (2.625,-1.753) (2.575,-1.725) (2.531,-1.700) (2.487,-1.675) (2.425,-1.639) (2.375,-1.609) (2.325,-1.580) (2.276,-1.550) (2.225,-1.519) (2.175,-1.488) (2.125,-1.457) (2.075,-1.426) (2.025,-1.394) (1.975,-1.362) (1.925,-1.330) (1.878,-1.300) (1.825,-1.266) (1.775,-1.233) (1.725,-1.200) (1.675,-1.167) (1.625,-1.133) (1.576,-1.100) (1.525,-1.066) (1.475,-1.032) (1.428,-1.000) (1.375,-0.964) (1.325,-0.930) (1.282,-0.900) (1.225,-0.861) (1.175,-0.827) (1.125,-0.792) (1.075,-0.757) (1.029,-0.725) (0.975,-0.688) (0.925,-0.652) (0.875,-0.618) (0.825,-0.583) (0.779,-0.550) (0.725,-0.512) (0.675,-0.477) (0.625,-0.442) (0.575,-0.407) (0.531,-0.375) (0.475,-0.336) (0.425,-0.300) (0.375,-0.266) (0.325,-0.230) (0.284,-0.200) (0.225,-0.160) (0.206,-0.144)}
\EQ{(-0.220,-0.155) (-0.275,-0.195) (-0.319,-0.225) (-0.375,-0.266) (-0.424,-0.300) (-0.475,-0.336) (-0.525,-0.371) (-0.567,-0.400) (-0.625,-0.442) (-0.672,-0.475) (-0.725,-0.513) (-0.775,-0.548) (-0.825,-0.583) (-0.875,-0.618) (-0.921,-0.650) (-0.975,-0.689) (-1.025,-0.724) (-1.075,-0.759) (-1.125,-0.794) (-1.170,-0.825) (-1.225,-0.864) (-1.275,-0.899) (-1.325,-0.934) (-1.375,-0.968) (-1.421,-1.000) (-1.475,-1.038) (-1.525,-1.072) (-1.575,-1.107) (-1.625,-1.141) (-1.674,-1.175) (-1.725,-1.210) (-1.775,-1.244) (-1.821,-1.275) (-1.875,-1.312) (-1.925,-1.346) (-1.969,-1.375) (-2.025,-1.413) (-2.075,-1.446) (-2.118,-1.475) (-2.175,-1.513) (-2.225,-1.546) (-2.270,-1.575) (-2.325,-1.611) (-2.375,-1.644) (-2.424,-1.675) (-2.475,-1.708) (-2.525,-1.740) (-2.575,-1.772) (-2.620,-1.800) (-2.675,-1.834) (-2.725,-1.866) (-2.775,-1.896) (-2.822,-1.925) (-2.875,-1.957) (-2.925,-1.988) (-2.975,-2.017) (-3.025,-2.047) (-3.073,-2.075) (-3.125,-2.105) (-3.175,-2.134) (-3.225,-2.163) (-3.275,-2.191) (-3.325,-2.219) (-3.375,-2.246) (-3.425,-2.274) (-3.474,-2.300) (-3.521,-2.325) (-3.569,-2.350) (-3.617,-2.375) (-3.665,-2.400) (-3.715,-2.425) (-3.765,-2.450) (-3.816,-2.475) (-3.867,-2.500) (-3.920,-2.525) (-3.973,-2.550) (-4.025,-2.574) (-4.075,-2.596) (-4.125,-2.619) (-4.175,-2.640) (-4.225,-2.662) (-4.275,-2.683) (-4.325,-2.703) (-4.375,-2.723) (-4.425,-2.743) (-4.475,-2.762) (-4.525,-2.781) (-4.575,-2.799) (-4.625,-2.817) (-4.675,-2.834) (-4.725,-2.851) (-4.775,-2.867) (-4.825,-2.883) (-4.875,-2.899) (-4.925,-2.914) (-4.975,-2.928) (-5.025,-2.942) (-5.075,-2.956) (-5.125,-2.968) (-5.175,-2.981) (-5.225,-2.993) (-5.275,-3.004) (-5.325,-3.015) (-5.375,-3.025) (-5.425,-3.035) (-5.475,-3.044) (-5.525,-3.053) (-5.575,-3.061) (-5.625,-3.069) (-5.675,-3.075) (-5.725,-3.082) (-5.775,-3.088) (-5.825,-3.093) (-5.875,-3.098) (-5.925,-3.102) (-5.975,-3.105) (-6.025,-3.108) (-6.075,-3.110) (-6.125,-3.112) (-6.175,-3.113) (-6.225,-3.113) (-6.275,-3.113) (-6.325,-3.112) (-6.375,-3.110) (-6.425,-3.108) (-6.475,-3.105) (-6.525,-3.101) (-6.575,-3.097) (-6.625,-3.092) (-6.675,-3.086) (-6.725,-3.080) (-6.775,-3.072) (-6.825,-3.065) (-6.875,-3.056) (-6.925,-3.046) (-6.975,-3.036) (-7.025,-3.025) (-7.075,-3.013) (-7.125,-3.001) (-7.175,-2.987) (-7.225,-2.973) (-7.275,-2.958) (-7.325,-2.942) (-7.375,-2.925) (-7.425,-2.907) (-7.475,-2.888) (-7.525,-2.868) (-7.575,-2.848) (-7.625,-2.826) (-7.675,-2.803) (-7.725,-2.779) (-7.775,-2.754) (-7.825,-2.728) (-7.875,-2.701) (-7.920,-2.675) (-7.975,-2.642) (-8.025,-2.611) (-8.075,-2.579) (-8.117,-2.550) (-8.175,-2.509) (-8.222,-2.475) (-8.275,-2.434) (-8.318,-2.400) (-8.375,-2.352) (-8.425,-2.308) (-8.475,-2.262) (-8.514,-2.225) (-8.550,-2.189) (-8.589,-2.150) (-8.625,-2.111) (-8.675,-2.056) (-8.723,-2.000) (-8.765,-1.950) (-8.800,-1.905) (-8.842,-1.850) (-8.875,-1.804) (-8.912,-1.750) (-8.946,-1.700) (-8.975,-1.654) (-9.008,-1.600) (-9.037,-1.550) (-9.065,-1.500) (-9.092,-1.450) (-9.118,-1.400) (-9.142,-1.350) (-9.165,-1.300) (-9.188,-1.250) (-9.209,-1.200) (-9.229,-1.150) (-9.248,-1.100) (-9.267,-1.050) (-9.284,-1.000) (-9.300,-0.951) (-9.316,-0.900) (-9.330,-0.850) (-9.344,-0.800) (-9.357,-0.750) (-9.368,-0.700) (-9.379,-0.650) (-9.390,-0.600) (-9.399,-0.550) (-9.407,-0.500) (-9.415,-0.450) (-9.422,-0.400) (-9.428,-0.350) (-9.433,-0.300) (-9.437,-0.250) (-9.441,-0.200) (-9.444,-0.150) (-9.446,-0.100) (-9.447,-0.050) (-9.447,0.000) (-9.447,0.050) (-9.446,0.100) (-9.444,0.150) (-9.441,0.200) (-9.437,0.250) (-9.433,0.300) (-9.428,0.350) (-9.422,0.400) (-9.415,0.450) (-9.407,0.500) (-9.399,0.550) (-9.390,0.600) (-9.379,0.650) (-9.368,0.700) (-9.357,0.750) (-9.344,0.800) (-9.330,0.850) (-9.316,0.900) (-9.300,0.950) (-9.284,1.000) (-9.267,1.050) (-9.248,1.100) (-9.229,1.150) (-9.209,1.200) (-9.188,1.250) (-9.165,1.300) (-9.142,1.350) (-9.118,1.400) (-9.092,1.450) (-9.065,1.500) (-9.037,1.550) (-9.008,1.600) (-8.978,1.650) (-8.950,1.693) (-8.912,1.750) (-8.878,1.800) (-8.842,1.850) (-8.804,1.900) (-8.765,1.950) (-8.725,1.998) (-8.680,2.050) (-8.636,2.100) (-8.600,2.138) (-8.564,2.175) (-8.525,2.214) (-8.488,2.250) (-8.450,2.285) (-8.406,2.325) (-8.350,2.373) (-8.300,2.414) (-8.255,2.450) (-8.200,2.491) (-8.153,2.525) (-8.100,2.562) (-8.050,2.595) (-8.003,2.625) (-7.950,2.657) (-7.900,2.686) (-7.850,2.714) (-7.800,2.741) (-7.750,2.767) (-7.700,2.791) (-7.650,2.814) (-7.600,2.837) (-7.550,2.858) (-7.500,2.878) (-7.450,2.898) (-7.400,2.916) (-7.350,2.933) (-7.300,2.950) (-7.250,2.966) (-7.200,2.980) (-7.150,2.994) (-7.100,3.007) (-7.050,3.019) (-7.000,3.031) (-6.950,3.041) (-6.900,3.051) (-6.850,3.060) (-6.800,3.069) (-6.750,3.076) (-6.700,3.083) (-6.650,3.089) (-6.600,3.094) (-6.550,3.099) (-6.500,3.103) (-6.450,3.106) (-6.400,3.109) (-6.350,3.111) (-6.300,3.112) (-6.250,3.113) (-6.200,3.113) (-6.150,3.112) (-6.100,3.111) (-6.050,3.109) (-6.000,3.107) (-5.950,3.103) (-5.900,3.100) (-5.850,3.095) (-5.800,3.090) (-5.750,3.085) (-5.700,3.079) (-5.650,3.072) (-5.600,3.065) (-5.550,3.057) (-5.500,3.049) (-5.450,3.040) (-5.400,3.030) (-5.350,3.020) (-5.300,3.010) (-5.250,2.998) (-5.200,2.987) (-5.151,2.975) (-5.100,2.962) (-5.050,2.949) (-5.000,2.935) (-4.950,2.921) (-4.900,2.906) (-4.850,2.891) (-4.800,2.875) (-4.750,2.859) (-4.700,2.843) (-4.650,2.825) (-4.600,2.808) (-4.550,2.790) (-4.500,2.771) (-4.450,2.752) (-4.400,2.733) (-4.350,2.713) (-4.300,2.693) (-4.250,2.672) (-4.200,2.651) (-4.150,2.629) (-4.100,2.608) (-4.050,2.585) (-4.000,2.562) (-3.950,2.539) (-3.900,2.516) (-3.850,2.492) (-3.800,2.467) (-3.750,2.443) (-3.700,2.418) (-3.650,2.392) (-3.600,2.366) (-3.550,2.340) (-3.500,2.314) (-3.450,2.287) (-3.400,2.260) (-3.350,2.233) (-3.300,2.205) (-3.250,2.177) (-3.203,2.150) (-3.159,2.125) (-3.116,2.100) (-3.050,2.062) (-3.000,2.032) (-2.950,2.002) (-2.904,1.975) (-2.850,1.942) (-2.800,1.912) (-2.750,1.881) (-2.700,1.850) (-2.650,1.819) (-2.600,1.787) (-2.550,1.756) (-2.502,1.725) (-2.450,1.692) (-2.400,1.660) (-2.350,1.627) (-2.300,1.595) (-2.250,1.562) (-2.200,1.529) (-2.156,1.500) (-2.100,1.463) (-2.050,1.430) (-2.006,1.400) (-1.950,1.363) (-1.900,1.329) (-1.858,1.300) (-1.800,1.261) (-1.750,1.227) (-1.700,1.193) (-1.650,1.158) (-1.602,1.125) (-1.550,1.090) (-1.500,1.055) (-1.457,1.025) (-1.400,0.986) (-1.350,0.951) (-1.300,0.916) (-1.250,0.881) (-1.206,0.850) (-1.150,0.811) (-1.100,0.776) (-1.050,0.741) (-1.000,0.706) (-0.956,0.675) (-0.900,0.636) (-0.850,0.600) (-0.800,0.566) (-0.750,0.530) (-0.708,0.500) (-0.650,0.460) (-0.601,0.425) (-0.550,0.389) (-0.500,0.354) (-0.450,0.319) (-0.400,0.283) (-0.354,0.250) (-0.300,0.213) (-0.250,0.177) (-0.220,0.155)}
\EQ{(0.206,0.144) (0.248,0.175) (0.300,0.213) (0.350,0.248) (0.400,0.283) (0.450,0.319) (0.496,0.350) (0.550,0.389) (0.600,0.424) (0.650,0.460) (0.700,0.495) (0.743,0.525) (0.800,0.566) (0.849,0.600) (0.900,0.636) (0.950,0.671) (0.992,0.700) (1.050,0.741) (1.098,0.775) (1.150,0.811) (1.200,0.846) (1.242,0.875) (1.300,0.916) (1.349,0.950) (1.400,0.986) (1.450,1.020) (1.493,1.050) (1.550,1.090) (1.600,1.124) (1.650,1.158) (1.700,1.193) (1.748,1.225) (1.800,1.261) (1.850,1.295) (1.895,1.325) (1.950,1.363) (2.000,1.396) (2.043,1.425) (2.100,1.463) (2.150,1.496) (2.194,1.525) (2.250,1.562) (2.300,1.595) (2.347,1.625) (2.400,1.660) (2.450,1.692) (2.500,1.724) (2.550,1.756) (2.600,1.787) (2.650,1.819) (2.700,1.850) (2.750,1.881) (2.800,1.912) (2.850,1.942) (2.900,1.972) (2.946,2.000) (3.000,2.032) (3.050,2.062) (3.100,2.091) (3.150,2.120) (3.200,2.148) (3.247,2.175) (3.292,2.200) (3.337,2.225) (3.382,2.250) (3.428,2.275) (3.474,2.300) (3.521,2.325) (3.569,2.350) (3.617,2.375) (3.665,2.400) (3.715,2.425) (3.765,2.450) (3.816,2.475) (3.867,2.500) (3.920,2.525) (3.973,2.550) (4.025,2.574) (4.075,2.596) (4.125,2.619) (4.175,2.640) (4.225,2.662) (4.275,2.683) (4.325,2.703) (4.375,2.723) (4.425,2.743) (4.475,2.762) (4.525,2.781) (4.575,2.799) (4.625,2.817) (4.675,2.834) (4.725,2.851) (4.775,2.867) (4.825,2.883) (4.875,2.899) (4.925,2.914) (4.975,2.928) (5.025,2.942) (5.075,2.956) (5.125,2.968) (5.175,2.981) (5.225,2.993) (5.275,3.004) (5.325,3.015) (5.375,3.025) (5.425,3.035) (5.475,3.044) (5.525,3.053) (5.575,3.061) (5.625,3.069) (5.675,3.075) (5.725,3.082) (5.775,3.088) (5.825,3.093) (5.875,3.098) (5.925,3.102) (5.975,3.105) (6.025,3.108) (6.075,3.110) (6.125,3.112) (6.175,3.113) (6.225,3.113) (6.275,3.113) (6.325,3.112) (6.375,3.110) (6.425,3.108) (6.475,3.105) (6.525,3.101) (6.575,3.097) (6.625,3.092) (6.675,3.086) (6.725,3.080) (6.775,3.072) (6.825,3.065) (6.875,3.056) (6.925,3.046) (6.975,3.036) (7.025,3.025) (7.075,3.013) (7.125,3.001) (7.175,2.987) (7.225,2.973) (7.275,2.958) (7.325,2.942) (7.375,2.925) (7.425,2.907) (7.475,2.888) (7.525,2.868) (7.575,2.848) (7.625,2.826) (7.675,2.803) (7.725,2.779) (7.775,2.754) (7.825,2.728) (7.875,2.701) (7.920,2.675) (7.975,2.642) (8.025,2.611) (8.075,2.579) (8.117,2.550) (8.175,2.509) (8.222,2.475) (8.275,2.434) (8.318,2.400) (8.375,2.352) (8.425,2.308) (8.475,2.262) (8.514,2.225) (8.550,2.189) (8.589,2.150) (8.625,2.111) (8.675,2.056) (8.723,2.000) (8.765,1.950) (8.800,1.905) (8.842,1.850) (8.875,1.804) (8.912,1.750) (8.946,1.700) (8.975,1.654) (9.008,1.600) (9.037,1.550) (9.065,1.500) (9.092,1.450) (9.118,1.400) (9.142,1.350) (9.165,1.300) (9.188,1.250) (9.209,1.200) (9.229,1.150) (9.248,1.100) (9.267,1.050) (9.284,1.000) (9.300,0.951) (9.316,0.900) (9.330,0.850) (9.344,0.800) (9.357,0.750) (9.368,0.700) (9.379,0.650) (9.390,0.600) (9.399,0.550) (9.407,0.500) (9.415,0.450) (9.422,0.400) (9.428,0.350) (9.433,0.300) (9.437,0.250) (9.441,0.200) (9.444,0.150) (9.446,0.100) (9.447,0.050) (9.447,0.000) (9.447,-0.050) (9.446,-0.100) (9.444,-0.150) (9.441,-0.200) (9.437,-0.250) (9.433,-0.300) (9.428,-0.350) (9.422,-0.400) (9.415,-0.450) (9.407,-0.500) (9.399,-0.550) (9.390,-0.600) (9.379,-0.650) (9.368,-0.700) (9.357,-0.750) (9.344,-0.800) (9.330,-0.850) (9.316,-0.900) (9.300,-0.950) (9.284,-1.000) (9.267,-1.050) (9.248,-1.100) (9.229,-1.150) (9.209,-1.200) (9.188,-1.250) (9.165,-1.300) (9.142,-1.350) (9.118,-1.400) (9.092,-1.450) (9.065,-1.500) (9.037,-1.550) (9.008,-1.600) (8.978,-1.650) (8.950,-1.693) (8.912,-1.750) (8.878,-1.800) (8.842,-1.850) (8.804,-1.900) (8.765,-1.950) (8.725,-1.998) (8.680,-2.050) (8.636,-2.100) (8.600,-2.138) (8.564,-2.175) (8.525,-2.214) (8.488,-2.250) (8.450,-2.285) (8.406,-2.325) (8.350,-2.373) (8.300,-2.414) (8.255,-2.450) (8.200,-2.491) (8.153,-2.525) (8.100,-2.562) (8.050,-2.595) (8.003,-2.625) (7.950,-2.657) (7.900,-2.686) (7.850,-2.714) (7.800,-2.741) (7.750,-2.767) (7.700,-2.791) (7.650,-2.814) (7.600,-2.837) (7.550,-2.858) (7.500,-2.878) (7.450,-2.898) (7.400,-2.916) (7.350,-2.933) (7.300,-2.950) (7.250,-2.966) (7.200,-2.980) (7.150,-2.994) (7.100,-3.007) (7.050,-3.019) (7.000,-3.031) (6.950,-3.041) (6.900,-3.051) (6.850,-3.060) (6.800,-3.069) (6.750,-3.076) (6.700,-3.083) (6.650,-3.089) (6.600,-3.094) (6.550,-3.099) (6.500,-3.103) (6.450,-3.106) (6.400,-3.109) (6.350,-3.111) (6.300,-3.112) (6.250,-3.113) (6.200,-3.113) (6.150,-3.112) (6.100,-3.111) (6.050,-3.109) (6.000,-3.107) (5.950,-3.103) (5.900,-3.100) (5.850,-3.095) (5.800,-3.090) (5.750,-3.085) (5.700,-3.079) (5.650,-3.072) (5.600,-3.065) (5.550,-3.057) (5.500,-3.049) (5.450,-3.040) (5.400,-3.030) (5.350,-3.020) (5.300,-3.010) (5.250,-2.998) (5.200,-2.987) (5.151,-2.975) (5.100,-2.962) (5.050,-2.949) (5.000,-2.935) (4.950,-2.921) (4.900,-2.906) (4.850,-2.891) (4.800,-2.875) (4.750,-2.859) (4.700,-2.843) (4.650,-2.825) (4.600,-2.808) (4.550,-2.790) (4.500,-2.771) (4.450,-2.752) (4.400,-2.733) (4.350,-2.713) (4.300,-2.693) (4.250,-2.672) (4.200,-2.651) (4.150,-2.629) (4.100,-2.608) (4.050,-2.585) (4.000,-2.562) (3.950,-2.539) (3.900,-2.516) (3.850,-2.492) (3.800,-2.467) (3.750,-2.443) (3.700,-2.418) (3.650,-2.392) (3.600,-2.366) (3.550,-2.340) (3.500,-2.314) (3.450,-2.287) (3.400,-2.260) (3.350,-2.233) (3.300,-2.205) (3.250,-2.177) (3.203,-2.150) (3.159,-2.125) (3.116,-2.100) (3.050,-2.062) (3.000,-2.032) (2.950,-2.002) (2.904,-1.975) (2.850,-1.942) (2.800,-1.912) (2.750,-1.881) (2.700,-1.850) (2.650,-1.819) (2.600,-1.787) (2.550,-1.756) (2.502,-1.725) (2.450,-1.692) (2.400,-1.660) (2.350,-1.627) (2.300,-1.595) (2.250,-1.562) (2.200,-1.529) (2.156,-1.500) (2.100,-1.463) (2.050,-1.430) (2.006,-1.400) (1.950,-1.363) (1.900,-1.329) (1.858,-1.300) (1.800,-1.261) (1.750,-1.227) (1.700,-1.193) (1.650,-1.158) (1.602,-1.125) (1.550,-1.090) (1.500,-1.055) (1.457,-1.025) (1.400,-0.986) (1.350,-0.951) (1.300,-0.916) (1.250,-0.881) (1.206,-0.850) (1.150,-0.811) (1.100,-0.776) (1.050,-0.741) (1.000,-0.706) (0.956,-0.675) (0.900,-0.636) (0.850,-0.600) (0.800,-0.566) (0.750,-0.530) (0.708,-0.500) (0.650,-0.460) (0.601,-0.425) (0.550,-0.389) (0.500,-0.354) (0.450,-0.319) (0.400,-0.283) (0.354,-0.250) (0.300,-0.213) (0.250,-0.177) (0.206,-0.144)}
\EQ{(10.000,-7.071) (9.950,-7.036) (9.900,-7.000) (9.850,-6.965) (9.800,-6.930) (9.758,-6.900) (9.700,-6.859) (9.652,-6.825) (9.600,-6.788) (9.550,-6.753) (9.500,-6.718) (9.450,-6.682) (9.405,-6.650) (9.350,-6.611) (9.300,-6.576) (9.250,-6.541) (9.200,-6.505) (9.157,-6.475) (9.100,-6.435) (9.051,-6.400) (9.000,-6.364) (8.950,-6.329) (8.900,-6.293) (8.850,-6.258) (8.804,-6.225) (8.750,-6.187) (8.700,-6.152) (8.650,-6.117) (8.600,-6.081) (8.556,-6.050) (8.500,-6.010) (8.450,-5.975) (8.400,-5.940) (8.350,-5.904) (8.309,-5.875) (8.250,-5.834) (8.202,-5.800) (8.150,-5.763) (8.100,-5.728) (8.050,-5.692) (8.000,-5.657) (7.955,-5.625) (7.900,-5.586) (7.850,-5.551) (7.800,-5.515) (7.750,-5.480) (7.708,-5.450) (7.650,-5.409) (7.601,-5.375) (7.550,-5.339) (7.500,-5.303) (7.450,-5.268) (7.400,-5.233) (7.354,-5.200) (7.300,-5.162) (7.250,-5.127) (7.200,-5.091) (7.150,-5.056) (7.106,-5.025) (7.050,-4.985) (7.000,-4.950) (6.950,-4.914) (6.900,-4.879) (6.859,-4.850) (6.800,-4.808) (6.753,-4.775) (6.700,-4.738) (6.650,-4.702) (6.600,-4.667) (6.550,-4.632) (6.505,-4.600) (6.450,-4.561) (6.400,-4.525) (6.350,-4.490) (6.300,-4.455) (6.258,-4.425) (6.200,-4.384) (6.152,-4.350) (6.100,-4.313) (6.050,-4.278) (6.000,-4.243) (5.950,-4.207) (5.904,-4.175) (5.850,-4.137) (5.800,-4.101) (5.750,-4.066) (5.700,-4.031) (5.657,-4.000) (5.600,-3.960) (5.551,-3.925) (5.500,-3.889) (5.450,-3.854) (5.400,-3.818) (5.350,-3.783) (5.303,-3.750) (5.250,-3.712) (5.200,-3.677) (5.150,-3.642) (5.100,-3.606) (5.056,-3.575) (5.000,-3.536) (4.950,-3.500) (4.900,-3.465) (4.850,-3.429) (4.808,-3.400) (4.750,-3.359) (4.702,-3.325) (4.650,-3.288) (4.600,-3.253) (4.550,-3.217) (4.500,-3.182) (4.455,-3.150) (4.400,-3.111) (4.350,-3.076) (4.300,-3.041) (4.250,-3.005) (4.207,-2.975) (4.150,-2.935) (4.101,-2.900) (4.050,-2.864) (4.000,-2.828) (3.950,-2.793) (3.900,-2.758) (3.854,-2.725) (3.800,-2.687) (3.750,-2.652) (3.700,-2.616) (3.650,-2.581) (3.606,-2.550) (3.550,-2.510) (3.500,-2.475) (3.450,-2.440) (3.400,-2.404) (3.359,-2.375) (3.300,-2.334) (3.253,-2.300) (3.200,-2.263) (3.150,-2.227) (3.100,-2.192) (3.050,-2.157) (3.005,-2.125) (2.950,-2.086) (2.900,-2.051) (2.850,-2.015) (2.800,-1.980) (2.758,-1.950) (2.700,-1.909) (2.652,-1.875) (2.600,-1.839) (2.550,-1.803) (2.500,-1.768) (2.450,-1.732) (2.404,-1.700) (2.350,-1.662) (2.300,-1.626) (2.250,-1.591) (2.200,-1.556) (2.157,-1.525) (2.100,-1.485) (2.051,-1.450) (2.000,-1.414) (1.950,-1.379) (1.900,-1.344) (1.850,-1.308) (1.803,-1.275) (1.750,-1.238) (1.700,-1.202) (1.650,-1.167) (1.600,-1.131) (1.556,-1.100) (1.500,-1.061) (1.450,-1.025) (1.400,-0.990) (1.350,-0.955) (1.308,-0.925) (1.250,-0.884) (1.202,-0.850) (1.150,-0.813) (1.100,-0.778) (1.050,-0.743) (1.000,-0.707) (0.955,-0.675) (0.900,-0.637) (0.850,-0.601) (0.800,-0.566) (0.750,-0.531) (0.708,-0.500) (0.650,-0.460) (0.601,-0.425) (0.550,-0.389) (0.500,-0.354) (0.450,-0.319) (0.400,-0.283) (0.354,-0.250) (0.300,-0.213) (0.250,-0.177) (0.206,-0.144)}
\EQ{(-0.220,-0.155) (-0.275,-0.195) (-0.319,-0.225) (-0.375,-0.266) (-0.424,-0.300) (-0.475,-0.336) (-0.525,-0.371) (-0.566,-0.400) (-0.625,-0.442) (-0.672,-0.475) (-0.725,-0.513) (-0.775,-0.548) (-0.825,-0.584) (-0.875,-0.619) (-0.920,-0.650) (-0.975,-0.690) (-1.025,-0.725) (-1.075,-0.760) (-1.125,-0.796) (-1.167,-0.825) (-1.225,-0.866) (-1.273,-0.900) (-1.325,-0.937) (-1.375,-0.972) (-1.425,-1.008) (-1.475,-1.043) (-1.520,-1.075) (-1.575,-1.114) (-1.625,-1.149) (-1.675,-1.185) (-1.725,-1.220) (-1.768,-1.250) (-1.825,-1.291) (-1.874,-1.325) (-1.925,-1.361) (-1.975,-1.397) (-2.025,-1.432) (-2.075,-1.467) (-2.121,-1.500) (-2.175,-1.538) (-2.225,-1.573) (-2.275,-1.609) (-2.325,-1.644) (-2.369,-1.675) (-2.425,-1.715) (-2.475,-1.750) (-2.525,-1.786) (-2.575,-1.821) (-2.616,-1.850) (-2.675,-1.892) (-2.722,-1.925) (-2.775,-1.962) (-2.825,-1.998) (-2.875,-2.033) (-2.925,-2.068) (-2.970,-2.100) (-3.025,-2.139) (-3.075,-2.174) (-3.125,-2.210) (-3.175,-2.245) (-3.217,-2.275) (-3.275,-2.316) (-3.323,-2.350) (-3.375,-2.387) (-3.425,-2.422) (-3.475,-2.457) (-3.525,-2.493) (-3.571,-2.525) (-3.625,-2.563) (-3.675,-2.599) (-3.725,-2.634) (-3.775,-2.669) (-3.818,-2.700) (-3.875,-2.740) (-3.924,-2.775) (-3.975,-2.811) (-4.025,-2.846) (-4.066,-2.875) (-4.125,-2.917) (-4.172,-2.950) (-4.225,-2.988) (-4.275,-3.023) (-4.325,-3.058) (-4.375,-3.094) (-4.419,-3.125) (-4.475,-3.164) (-4.525,-3.200) (-4.575,-3.235) (-4.625,-3.270) (-4.667,-3.300) (-4.725,-3.341) (-4.773,-3.375) (-4.825,-3.412) (-4.875,-3.447) (-4.925,-3.483) (-4.975,-3.518) (-5.021,-3.550) (-5.075,-3.589) (-5.125,-3.624) (-5.175,-3.659) (-5.225,-3.695) (-5.268,-3.725) (-5.325,-3.765) (-5.374,-3.800) (-5.425,-3.836) (-5.475,-3.871) (-5.525,-3.907) (-5.575,-3.942) (-5.622,-3.975) (-5.675,-4.013) (-5.725,-4.048) (-5.775,-4.084) (-5.825,-4.119) (-5.869,-4.150) (-5.925,-4.190) (-5.975,-4.225) (-6.025,-4.260) (-6.075,-4.296) (-6.117,-4.325) (-6.175,-4.366) (-6.223,-4.400) (-6.275,-4.437) (-6.325,-4.472) (-6.375,-4.508) (-6.425,-4.543) (-6.470,-4.575) (-6.525,-4.614) (-6.575,-4.649) (-6.625,-4.685) (-6.675,-4.720) (-6.718,-4.750) (-6.775,-4.791) (-6.824,-4.825) (-6.875,-4.861) (-6.925,-4.897) (-6.975,-4.932) (-7.025,-4.967) (-7.071,-5.000) (-7.125,-5.038) (-7.175,-5.073) (-7.225,-5.109) (-7.275,-5.144) (-7.319,-5.175) (-7.375,-5.215) (-7.425,-5.250) (-7.475,-5.286) (-7.525,-5.321) (-7.566,-5.350) (-7.625,-5.392) (-7.672,-5.425) (-7.725,-5.462) (-7.775,-5.498) (-7.825,-5.533) (-7.875,-5.568) (-7.920,-5.600) (-7.975,-5.639) (-8.025,-5.675) (-8.075,-5.710) (-8.125,-5.745) (-8.167,-5.775) (-8.225,-5.816) (-8.273,-5.850) (-8.325,-5.887) (-8.375,-5.922) (-8.425,-5.957) (-8.475,-5.993) (-8.521,-6.025) (-8.575,-6.063) (-8.625,-6.099) (-8.675,-6.134) (-8.725,-6.170) (-8.768,-6.200) (-8.825,-6.240) (-8.874,-6.275) (-8.925,-6.311) (-8.975,-6.346) (-9.025,-6.382) (-9.075,-6.417) (-9.122,-6.450) (-9.175,-6.488) (-9.225,-6.523) (-9.275,-6.558) (-9.325,-6.594) (-9.369,-6.625) (-9.425,-6.665) (-9.475,-6.700) (-9.525,-6.735) (-9.575,-6.771) (-9.617,-6.800) (-9.675,-6.841) (-9.723,-6.875) (-9.775,-6.912) (-9.825,-6.947) (-9.875,-6.983) (-9.925,-7.018) (-9.970,-7.050) (-10.000,-7.071)}
\EQ{(0.206,0.144) (0.248,0.175) (0.300,0.213) (0.350,0.248) (0.400,0.283) (0.450,0.319) (0.496,0.350) (0.550,0.389) (0.600,0.424) (0.650,0.460) (0.700,0.495) (0.743,0.525) (0.800,0.566) (0.849,0.600) (0.900,0.637) (0.950,0.672) (1.000,0.707) (1.050,0.743) (1.096,0.775) (1.150,0.813) (1.200,0.849) (1.250,0.884) (1.300,0.919) (1.344,0.950) (1.400,0.990) (1.450,1.025) (1.500,1.061) (1.550,1.096) (1.591,1.125) (1.650,1.167) (1.697,1.200) (1.750,1.238) (1.800,1.273) (1.850,1.308) (1.900,1.344) (1.945,1.375) (2.000,1.414) (2.050,1.450) (2.100,1.485) (2.150,1.520) (2.192,1.550) (2.250,1.591) (2.298,1.625) (2.350,1.662) (2.400,1.697) (2.450,1.732) (2.500,1.768) (2.546,1.800) (2.600,1.839) (2.650,1.874) (2.700,1.909) (2.750,1.945) (2.793,1.975) (2.850,2.015) (2.899,2.050) (2.950,2.086) (3.000,2.121) (3.050,2.157) (3.100,2.192) (3.147,2.225) (3.200,2.263) (3.250,2.298) (3.300,2.334) (3.350,2.369) (3.394,2.400) (3.450,2.440) (3.500,2.475) (3.550,2.510) (3.600,2.546) (3.642,2.575) (3.700,2.616) (3.748,2.650) (3.800,2.687) (3.850,2.722) (3.900,2.758) (3.950,2.793) (3.995,2.825) (4.050,2.864) (4.100,2.899) (4.150,2.935) (4.200,2.970) (4.243,3.000) (4.300,3.041) (4.349,3.075) (4.400,3.111) (4.450,3.147) (4.500,3.182) (4.550,3.217) (4.596,3.250) (4.650,3.288) (4.700,3.323) (4.750,3.359) (4.800,3.394) (4.844,3.425) (4.900,3.465) (4.950,3.500) (5.000,3.536) (5.050,3.571) (5.091,3.600) (5.150,3.642) (5.197,3.675) (5.250,3.712) (5.300,3.748) (5.350,3.783) (5.400,3.818) (5.445,3.850) (5.500,3.889) (5.550,3.924) (5.600,3.960) (5.650,3.995) (5.692,4.025) (5.750,4.066) (5.798,4.100) (5.850,4.137) (5.900,4.172) (5.950,4.207) (6.000,4.243) (6.046,4.275) (6.100,4.313) (6.150,4.349) (6.200,4.384) (6.250,4.419) (6.293,4.450) (6.350,4.490) (6.399,4.525) (6.450,4.561) (6.500,4.596) (6.550,4.632) (6.600,4.667) (6.647,4.700) (6.700,4.738) (6.750,4.773) (6.800,4.808) (6.850,4.844) (6.894,4.875) (6.950,4.914) (7.000,4.950) (7.050,4.985) (7.100,5.020) (7.142,5.050) (7.200,5.091) (7.248,5.125) (7.300,5.162) (7.350,5.197) (7.400,5.233) (7.450,5.268) (7.495,5.300) (7.550,5.339) (7.600,5.374) (7.650,5.409) (7.700,5.445) (7.743,5.475) (7.800,5.515) (7.849,5.550) (7.900,5.586) (7.950,5.622) (8.000,5.657) (8.050,5.692) (8.096,5.725) (8.150,5.763) (8.200,5.798) (8.250,5.834) (8.300,5.869) (8.344,5.900) (8.400,5.940) (8.450,5.975) (8.500,6.010) (8.550,6.046) (8.591,6.075) (8.650,6.117) (8.697,6.150) (8.750,6.187) (8.800,6.223) (8.850,6.258) (8.900,6.293) (8.945,6.325) (9.000,6.364) (9.050,6.399) (9.100,6.435) (9.150,6.470) (9.192,6.500) (9.250,6.541) (9.298,6.575) (9.350,6.611) (9.400,6.647) (9.450,6.682) (9.500,6.718) (9.546,6.750) (9.600,6.788) (9.650,6.824) (9.700,6.859) (9.750,6.894) (9.793,6.925) (9.850,6.965) (9.899,7.000) (9.950,7.036) (10.000,7.071)}
\EQ{(-10.000,7.071) (-9.950,7.036) (-9.900,7.000) (-9.850,6.965) (-9.800,6.930) (-9.758,6.900) (-9.700,6.859) (-9.652,6.825) (-9.600,6.788) (-9.550,6.753) (-9.500,6.718) (-9.450,6.682) (-9.405,6.650) (-9.350,6.611) (-9.300,6.576) (-9.250,6.541) (-9.200,6.505) (-9.157,6.475) (-9.100,6.435) (-9.051,6.400) (-9.000,6.364) (-8.950,6.329) (-8.900,6.293) (-8.850,6.258) (-8.804,6.225) (-8.750,6.187) (-8.700,6.152) (-8.650,6.117) (-8.600,6.081) (-8.556,6.050) (-8.500,6.010) (-8.450,5.975) (-8.400,5.940) (-8.350,5.904) (-8.309,5.875) (-8.250,5.834) (-8.202,5.800) (-8.150,5.763) (-8.100,5.728) (-8.050,5.692) (-8.000,5.657) (-7.955,5.625) (-7.900,5.586) (-7.850,5.551) (-7.800,5.515) (-7.750,5.480) (-7.708,5.450) (-7.650,5.409) (-7.601,5.375) (-7.550,5.339) (-7.500,5.303) (-7.450,5.268) (-7.400,5.233) (-7.354,5.200) (-7.300,5.162) (-7.250,5.127) (-7.200,5.091) (-7.150,5.056) (-7.106,5.025) (-7.050,4.985) (-7.000,4.950) (-6.950,4.914) (-6.900,4.879) (-6.859,4.850) (-6.800,4.808) (-6.753,4.775) (-6.700,4.738) (-6.650,4.702) (-6.600,4.667) (-6.550,4.632) (-6.505,4.600) (-6.450,4.561) (-6.400,4.525) (-6.350,4.490) (-6.300,4.455) (-6.258,4.425) (-6.200,4.384) (-6.152,4.350) (-6.100,4.313) (-6.050,4.278) (-6.000,4.243) (-5.950,4.207) (-5.904,4.175) (-5.850,4.137) (-5.800,4.101) (-5.750,4.066) (-5.700,4.031) (-5.657,4.000) (-5.600,3.960) (-5.551,3.925) (-5.500,3.889) (-5.450,3.854) (-5.400,3.818) (-5.350,3.783) (-5.303,3.750) (-5.250,3.712) (-5.200,3.677) (-5.150,3.642) (-5.100,3.606) (-5.056,3.575) (-5.000,3.536) (-4.950,3.500) (-4.900,3.465) (-4.850,3.429) (-4.808,3.400) (-4.750,3.359) (-4.702,3.325) (-4.650,3.288) (-4.600,3.253) (-4.550,3.217) (-4.500,3.182) (-4.455,3.150) (-4.400,3.111) (-4.350,3.076) (-4.300,3.041) (-4.250,3.005) (-4.207,2.975) (-4.150,2.935) (-4.101,2.900) (-4.050,2.864) (-4.000,2.828) (-3.950,2.793) (-3.900,2.758) (-3.854,2.725) (-3.800,2.687) (-3.750,2.652) (-3.700,2.616) (-3.650,2.581) (-3.606,2.550) (-3.550,2.510) (-3.500,2.475) (-3.450,2.440) (-3.400,2.404) (-3.359,2.375) (-3.300,2.334) (-3.253,2.300) (-3.200,2.263) (-3.150,2.227) (-3.100,2.192) (-3.050,2.157) (-3.005,2.125) (-2.950,2.086) (-2.900,2.051) (-2.850,2.015) (-2.800,1.980) (-2.758,1.950) (-2.700,1.909) (-2.652,1.875) (-2.600,1.839) (-2.550,1.803) (-2.500,1.768) (-2.450,1.732) (-2.404,1.700) (-2.350,1.662) (-2.300,1.626) (-2.250,1.591) (-2.200,1.556) (-2.157,1.525) (-2.100,1.485) (-2.051,1.450) (-2.000,1.414) (-1.950,1.379) (-1.900,1.344) (-1.850,1.308) (-1.803,1.275) (-1.750,1.238) (-1.700,1.202) (-1.650,1.167) (-1.600,1.131) (-1.556,1.100) (-1.500,1.061) (-1.450,1.025) (-1.400,0.990) (-1.350,0.955) (-1.308,0.925) (-1.250,0.884) (-1.202,0.850) (-1.150,0.813) (-1.100,0.778) (-1.050,0.743) (-1.000,0.707) (-0.955,0.675) (-0.900,0.637) (-0.850,0.601) (-0.800,0.566) (-0.750,0.531) (-0.708,0.500) (-0.650,0.460) (-0.601,0.425) (-0.550,0.389) (-0.500,0.354) (-0.450,0.319) (-0.400,0.283) (-0.354,0.250) (-0.300,0.213) (-0.250,0.177) (-0.220,0.155)}
\EQ{(4.306,-10.000) (4.350,-9.950) (4.392,-9.900) (4.425,-9.860) (4.473,-9.800) (4.513,-9.750) (4.550,-9.701) (4.588,-9.650) (4.624,-9.600) (4.660,-9.550) (4.695,-9.500) (4.725,-9.455) (4.761,-9.400) (4.793,-9.350) (4.824,-9.300) (4.855,-9.250) (4.884,-9.200) (4.913,-9.150) (4.941,-9.100) (4.969,-9.050) (4.995,-9.000) (5.021,-8.950) (5.046,-8.900) (5.071,-8.850) (5.095,-8.800) (5.118,-8.750) (5.140,-8.700) (5.162,-8.650) (5.183,-8.600) (5.203,-8.550) (5.223,-8.500) (5.242,-8.450) (5.261,-8.400) (5.279,-8.350) (5.296,-8.300) (5.313,-8.250) (5.329,-8.200) (5.344,-8.150) (5.359,-8.100) (5.373,-8.050) (5.387,-8.000) (5.400,-7.951) (5.413,-7.900) (5.425,-7.850) (5.436,-7.800) (5.447,-7.750) (5.457,-7.700) (5.467,-7.650) (5.476,-7.600) (5.484,-7.550) (5.492,-7.500) (5.500,-7.450) (5.507,-7.400) (5.513,-7.350) (5.519,-7.300) (5.524,-7.250) (5.529,-7.200) (5.533,-7.150) (5.537,-7.100) (5.540,-7.050) (5.542,-7.000) (5.544,-6.950) (5.545,-6.900) (5.546,-6.850) (5.547,-6.800) (5.546,-6.750) (5.546,-6.700) (5.544,-6.650) (5.543,-6.600) (5.540,-6.550) (5.537,-6.500) (5.534,-6.450) (5.530,-6.400) (5.525,-6.350) (5.520,-6.300) (5.515,-6.250) (5.509,-6.200) (5.502,-6.150) (5.495,-6.100) (5.487,-6.050) (5.478,-6.000) (5.469,-5.950) (5.460,-5.900) (5.450,-5.850) (5.439,-5.800) (5.428,-5.750) (5.416,-5.700) (5.404,-5.650) (5.391,-5.600) (5.378,-5.550) (5.363,-5.500) (5.349,-5.450) (5.334,-5.400) (5.318,-5.350) (5.301,-5.300) (5.284,-5.250) (5.267,-5.200) (5.249,-5.150) (5.230,-5.100) (5.210,-5.050) (5.190,-5.000) (5.170,-4.950) (5.150,-4.903) (5.127,-4.850) (5.104,-4.800) (5.081,-4.750) (5.057,-4.700) (5.033,-4.650) (5.008,-4.600) (4.982,-4.550) (4.956,-4.500) (4.929,-4.450) (4.901,-4.400) (4.875,-4.354) (4.844,-4.300) (4.814,-4.250) (4.784,-4.200) (4.753,-4.150) (4.725,-4.107) (4.688,-4.050) (4.655,-4.000) (4.625,-3.955) (4.587,-3.900) (4.551,-3.850) (4.516,-3.800) (4.479,-3.750) (4.441,-3.700) (4.403,-3.650) (4.364,-3.600) (4.325,-3.551) (4.284,-3.500) (4.250,-3.458) (4.201,-3.400) (4.159,-3.350) (4.125,-3.311) (4.075,-3.254) (4.026,-3.200) (3.981,-3.150) (3.935,-3.100) (3.900,-3.063) (3.864,-3.025) (3.825,-2.985) (3.775,-2.934) (3.725,-2.883) (3.675,-2.834) (3.625,-2.785) (3.575,-2.736) (3.537,-2.700) (3.500,-2.665) (3.457,-2.625) (3.403,-2.575) (3.350,-2.526) (3.300,-2.481) (3.250,-2.437) (3.209,-2.400) (3.151,-2.350) (3.100,-2.306) (3.050,-2.263) (3.006,-2.225) (2.950,-2.178) (2.900,-2.137) (2.856,-2.100) (2.800,-2.054) (2.750,-2.014) (2.702,-1.975) (2.650,-1.933) (2.609,-1.900) (2.550,-1.854) (2.500,-1.814) (2.450,-1.775) (2.400,-1.737) (2.353,-1.700) (2.300,-1.660) (2.255,-1.625) (2.200,-1.583) (2.156,-1.550) (2.100,-1.508) (2.056,-1.475) (2.000,-1.433) (1.956,-1.400) (1.900,-1.359) (1.854,-1.325) (1.800,-1.285) (1.752,-1.250) (1.700,-1.212) (1.650,-1.175) (1.600,-1.139) (1.550,-1.103) (1.500,-1.067) (1.450,-1.031) (1.408,-1.000) (1.350,-0.959) (1.304,-0.925) (1.250,-0.887) (1.200,-0.851) (1.150,-0.815) (1.100,-0.780) (1.059,-0.750) (1.000,-0.708) (0.953,-0.675) (0.900,-0.637) (0.850,-0.602) (0.800,-0.566) (0.750,-0.531) (0.707,-0.500) (0.650,-0.460) (0.601,-0.425) (0.550,-0.389) (0.500,-0.354) (0.450,-0.319) (0.400,-0.283) (0.354,-0.250) (0.300,-0.213) (0.250,-0.177) (0.206,-0.144)}
\EQ{(-0.220,-0.155) (-0.275,-0.195) (-0.319,-0.225) (-0.375,-0.266) (-0.424,-0.300) (-0.475,-0.336) (-0.525,-0.372) (-0.566,-0.400) (-0.625,-0.442) (-0.672,-0.475) (-0.725,-0.513) (-0.775,-0.549) (-0.825,-0.584) (-0.875,-0.620) (-0.918,-0.650) (-0.975,-0.691) (-1.024,-0.725) (-1.075,-0.762) (-1.125,-0.797) (-1.175,-0.833) (-1.225,-0.869) (-1.269,-0.900) (-1.325,-0.941) (-1.373,-0.975) (-1.425,-1.013) (-1.475,-1.049) (-1.525,-1.085) (-1.575,-1.121) (-1.625,-1.157) (-1.675,-1.194) (-1.718,-1.225) (-1.775,-1.267) (-1.820,-1.300) (-1.875,-1.340) (-1.922,-1.375) (-1.975,-1.415) (-2.023,-1.450) (-2.075,-1.489) (-2.123,-1.525) (-2.175,-1.564) (-2.222,-1.600) (-2.275,-1.640) (-2.320,-1.675) (-2.375,-1.717) (-2.418,-1.750) (-2.475,-1.795) (-2.525,-1.834) (-2.575,-1.873) (-2.625,-1.913) (-2.671,-1.950) (-2.725,-1.993) (-2.764,-2.025) (-2.825,-2.075) (-2.875,-2.116) (-2.916,-2.150) (-2.975,-2.199) (-3.025,-2.241) (-3.065,-2.275) (-3.123,-2.325) (-3.175,-2.371) (-3.225,-2.415) (-3.265,-2.450) (-3.321,-2.500) (-3.375,-2.549) (-3.425,-2.595) (-3.475,-2.642) (-3.525,-2.689) (-3.563,-2.725) (-3.600,-2.760) (-3.641,-2.800) (-3.692,-2.850) (-3.742,-2.900) (-3.791,-2.950) (-3.840,-3.000) (-3.875,-3.037) (-3.911,-3.075) (-3.950,-3.117) (-4.000,-3.171) (-4.049,-3.225) (-4.093,-3.275) (-4.137,-3.325) (-4.175,-3.369) (-4.222,-3.425) (-4.264,-3.475) (-4.300,-3.519) (-4.345,-3.575) (-4.384,-3.625) (-4.422,-3.675) (-4.460,-3.725) (-4.497,-3.775) (-4.534,-3.825) (-4.569,-3.875) (-4.600,-3.919) (-4.638,-3.975) (-4.672,-4.025) (-4.700,-4.068) (-4.737,-4.125) (-4.768,-4.175) (-4.799,-4.225) (-4.825,-4.269) (-4.858,-4.325) (-4.887,-4.375) (-4.915,-4.425) (-4.942,-4.475) (-4.969,-4.525) (-4.995,-4.575) (-5.020,-4.625) (-5.045,-4.675) (-5.069,-4.725) (-5.093,-4.775) (-5.116,-4.825) (-5.138,-4.875) (-5.159,-4.925) (-5.180,-4.975) (-5.200,-5.024) (-5.220,-5.075) (-5.239,-5.125) (-5.258,-5.175) (-5.275,-5.223) (-5.293,-5.275) (-5.310,-5.325) (-5.325,-5.373) (-5.341,-5.425) (-5.356,-5.475) (-5.371,-5.525) (-5.384,-5.575) (-5.398,-5.625) (-5.410,-5.675) (-5.422,-5.725) (-5.434,-5.775) (-5.445,-5.825) (-5.455,-5.875) (-5.465,-5.925) (-5.474,-5.975) (-5.483,-6.025) (-5.491,-6.075) (-5.498,-6.125) (-5.505,-6.175) (-5.512,-6.225) (-5.518,-6.275) (-5.523,-6.325) (-5.528,-6.375) (-5.532,-6.425) (-5.536,-6.475) (-5.539,-6.525) (-5.542,-6.575) (-5.544,-6.625) (-5.545,-6.675) (-5.546,-6.725) (-5.547,-6.775) (-5.547,-6.825) (-5.546,-6.875) (-5.545,-6.925) (-5.543,-6.975) (-5.541,-7.025) (-5.538,-7.075) (-5.535,-7.125) (-5.531,-7.175) (-5.526,-7.225) (-5.522,-7.275) (-5.516,-7.325) (-5.510,-7.375) (-5.503,-7.425) (-5.496,-7.475) (-5.489,-7.525) (-5.480,-7.575) (-5.471,-7.625) (-5.462,-7.675) (-5.452,-7.725) (-5.442,-7.775) (-5.430,-7.825) (-5.419,-7.875) (-5.407,-7.925) (-5.394,-7.975) (-5.380,-8.025) (-5.366,-8.075) (-5.352,-8.125) (-5.337,-8.175) (-5.321,-8.225) (-5.305,-8.275) (-5.288,-8.325) (-5.270,-8.375) (-5.252,-8.425) (-5.233,-8.475) (-5.213,-8.525) (-5.193,-8.575) (-5.173,-8.625) (-5.151,-8.675) (-5.129,-8.725) (-5.106,-8.775) (-5.083,-8.825) (-5.059,-8.875) (-5.034,-8.925) (-5.008,-8.975) (-4.982,-9.025) (-4.955,-9.075) (-4.927,-9.125) (-4.900,-9.173) (-4.870,-9.225) (-4.840,-9.275) (-4.809,-9.325) (-4.777,-9.375) (-4.750,-9.417) (-4.711,-9.475) (-4.677,-9.525) (-4.642,-9.575) (-4.606,-9.625) (-4.575,-9.668) (-4.532,-9.725) (-4.500,-9.766) (-4.453,-9.825) (-4.413,-9.875) (-4.375,-9.920) (-4.328,-9.975) (-4.306,-10.000)}
\EQ{(0.206,0.144) (0.248,0.175) (0.300,0.213) (0.350,0.248) (0.400,0.283) (0.450,0.319) (0.495,0.350) (0.550,0.389) (0.600,0.424) (0.650,0.460) (0.700,0.495) (0.742,0.525) (0.800,0.566) (0.848,0.600) (0.900,0.637) (0.950,0.673) (1.000,0.708) (1.050,0.744) (1.094,0.775) (1.150,0.815) (1.199,0.850) (1.250,0.887) (1.300,0.923) (1.350,0.959) (1.400,0.995) (1.443,1.025) (1.500,1.067) (1.546,1.100) (1.600,1.139) (1.650,1.175) (1.700,1.212) (1.750,1.248) (1.800,1.285) (1.850,1.322) (1.900,1.359) (1.950,1.396) (2.000,1.433) (2.050,1.470) (2.100,1.508) (2.150,1.546) (2.200,1.583) (2.250,1.621) (2.300,1.660) (2.350,1.698) (2.400,1.737) (2.450,1.775) (2.500,1.814) (2.546,1.850) (2.600,1.893) (2.640,1.925) (2.700,1.973) (2.750,2.014) (2.795,2.050) (2.850,2.095) (2.900,2.137) (2.946,2.175) (3.000,2.220) (3.050,2.263) (3.094,2.300) (3.150,2.349) (3.200,2.393) (3.250,2.437) (3.293,2.475) (3.349,2.525) (3.400,2.572) (3.450,2.618) (3.500,2.665) (3.537,2.700) (3.575,2.736) (3.615,2.775) (3.666,2.825) (3.717,2.875) (3.767,2.925) (3.815,2.975) (3.864,3.025) (3.900,3.063) (3.935,3.100) (3.975,3.144) (4.025,3.198) (4.071,3.250) (4.115,3.300) (4.150,3.340) (4.200,3.399) (4.243,3.450) (4.275,3.489) (4.325,3.550) (4.364,3.600) (4.400,3.646) (4.441,3.700) (4.475,3.745) (4.516,3.800) (4.550,3.848) (4.587,3.900) (4.621,3.950) (4.650,3.992) (4.688,4.050) (4.721,4.100) (4.750,4.146) (4.784,4.200) (4.814,4.250) (4.844,4.300) (4.873,4.350) (4.900,4.398) (4.925,4.443) (4.950,4.489) (4.975,4.536) (5.000,4.585) (5.025,4.634) (5.050,4.685) (5.075,4.737) (5.100,4.791) (5.125,4.846) (5.149,4.900) (5.170,4.950) (5.190,5.000) (5.210,5.050) (5.230,5.100) (5.249,5.150) (5.267,5.200) (5.284,5.250) (5.301,5.300) (5.318,5.350) (5.334,5.400) (5.349,5.450) (5.363,5.500) (5.378,5.550) (5.391,5.600) (5.404,5.650) (5.416,5.700) (5.428,5.750) (5.439,5.800) (5.450,5.850) (5.460,5.900) (5.469,5.950) (5.478,6.000) (5.487,6.050) (5.495,6.100) (5.502,6.150) (5.509,6.200) (5.515,6.250) (5.520,6.300) (5.525,6.350) (5.530,6.400) (5.534,6.450) (5.537,6.500) (5.540,6.550) (5.543,6.600) (5.544,6.650) (5.546,6.700) (5.546,6.750) (5.547,6.800) (5.546,6.850) (5.545,6.900) (5.544,6.950) (5.542,7.000) (5.540,7.050) (5.537,7.100) (5.533,7.150) (5.529,7.200) (5.524,7.250) (5.519,7.300) (5.513,7.350) (5.507,7.400) (5.500,7.450) (5.492,7.500) (5.484,7.550) (5.476,7.600) (5.467,7.650) (5.457,7.700) (5.447,7.750) (5.436,7.800) (5.425,7.850) (5.413,7.900) (5.400,7.950) (5.387,8.000) (5.373,8.050) (5.359,8.100) (5.344,8.150) (5.329,8.200) (5.313,8.250) (5.296,8.300) (5.279,8.350) (5.261,8.400) (5.242,8.450) (5.223,8.500) (5.203,8.550) (5.183,8.600) (5.162,8.650) (5.140,8.700) (5.118,8.750) (5.095,8.800) (5.071,8.850) (5.046,8.900) (5.021,8.950) (4.995,9.000) (4.969,9.050) (4.941,9.100) (4.913,9.150) (4.884,9.200) (4.855,9.250) (4.825,9.299) (4.793,9.350) (4.761,9.400) (4.728,9.450) (4.700,9.492) (4.660,9.550) (4.625,9.599) (4.588,9.650) (4.551,9.700) (4.513,9.750) (4.475,9.798) (4.433,9.850) (4.400,9.890) (4.350,9.949) (4.306,10.000)}
\EQ{(-4.306,10.000) (-4.350,9.950) (-4.392,9.900) (-4.425,9.860) (-4.473,9.800) (-4.513,9.750) (-4.550,9.701) (-4.588,9.650) (-4.624,9.600) (-4.660,9.550) (-4.695,9.500) (-4.725,9.455) (-4.761,9.400) (-4.793,9.350) (-4.824,9.300) (-4.855,9.250) (-4.884,9.200) (-4.913,9.150) (-4.941,9.100) (-4.969,9.050) (-4.995,9.000) (-5.021,8.950) (-5.046,8.900) (-5.071,8.850) (-5.095,8.800) (-5.118,8.750) (-5.140,8.700) (-5.162,8.650) (-5.183,8.600) (-5.203,8.550) (-5.223,8.500) (-5.242,8.450) (-5.261,8.400) (-5.279,8.350) (-5.296,8.300) (-5.313,8.250) (-5.329,8.200) (-5.344,8.150) (-5.359,8.100) (-5.373,8.050) (-5.387,8.000) (-5.400,7.951) (-5.413,7.900) (-5.425,7.850) (-5.436,7.800) (-5.447,7.750) (-5.457,7.700) (-5.467,7.650) (-5.476,7.600) (-5.484,7.550) (-5.492,7.500) (-5.500,7.450) (-5.507,7.400) (-5.513,7.350) (-5.519,7.300) (-5.524,7.250) (-5.529,7.200) (-5.533,7.150) (-5.537,7.100) (-5.540,7.050) (-5.542,7.000) (-5.544,6.950) (-5.545,6.900) (-5.546,6.850) (-5.547,6.800) (-5.546,6.750) (-5.546,6.700) (-5.544,6.650) (-5.543,6.600) (-5.540,6.550) (-5.537,6.500) (-5.534,6.450) (-5.530,6.400) (-5.525,6.350) (-5.520,6.300) (-5.515,6.250) (-5.509,6.200) (-5.502,6.150) (-5.495,6.100) (-5.487,6.050) (-5.478,6.000) (-5.469,5.950) (-5.460,5.900) (-5.450,5.850) (-5.439,5.800) (-5.428,5.750) (-5.416,5.700) (-5.404,5.650) (-5.391,5.600) (-5.378,5.550) (-5.363,5.500) (-5.349,5.450) (-5.334,5.400) (-5.318,5.350) (-5.301,5.300) (-5.284,5.250) (-5.267,5.200) (-5.249,5.150) (-5.230,5.100) (-5.210,5.050) (-5.190,5.000) (-5.170,4.950) (-5.150,4.903) (-5.127,4.850) (-5.104,4.800) (-5.081,4.750) (-5.057,4.700) (-5.033,4.650) (-5.008,4.600) (-4.982,4.550) (-4.956,4.500) (-4.929,4.450) (-4.901,4.400) (-4.875,4.354) (-4.844,4.300) (-4.814,4.250) (-4.784,4.200) (-4.753,4.150) (-4.725,4.107) (-4.688,4.050) (-4.655,4.000) (-4.625,3.955) (-4.587,3.900) (-4.551,3.850) (-4.516,3.800) (-4.479,3.750) (-4.441,3.700) (-4.403,3.650) (-4.364,3.600) (-4.325,3.551) (-4.284,3.500) (-4.250,3.458) (-4.201,3.400) (-4.159,3.350) (-4.125,3.311) (-4.075,3.254) (-4.026,3.200) (-3.981,3.150) (-3.935,3.100) (-3.900,3.063) (-3.864,3.025) (-3.825,2.985) (-3.775,2.934) (-3.725,2.883) (-3.675,2.834) (-3.625,2.785) (-3.575,2.736) (-3.537,2.700) (-3.500,2.665) (-3.457,2.625) (-3.403,2.575) (-3.350,2.526) (-3.300,2.481) (-3.250,2.437) (-3.209,2.400) (-3.151,2.350) (-3.100,2.306) (-3.050,2.263) (-3.006,2.225) (-2.950,2.178) (-2.900,2.137) (-2.856,2.100) (-2.800,2.054) (-2.750,2.014) (-2.702,1.975) (-2.650,1.933) (-2.609,1.900) (-2.550,1.854) (-2.500,1.814) (-2.450,1.775) (-2.400,1.737) (-2.353,1.700) (-2.300,1.660) (-2.255,1.625) (-2.200,1.583) (-2.156,1.550) (-2.100,1.508) (-2.056,1.475) (-2.000,1.433) (-1.956,1.400) (-1.900,1.359) (-1.854,1.325) (-1.800,1.285) (-1.752,1.250) (-1.700,1.212) (-1.650,1.175) (-1.600,1.139) (-1.550,1.103) (-1.500,1.067) (-1.450,1.031) (-1.408,1.000) (-1.350,0.959) (-1.304,0.925) (-1.250,0.887) (-1.200,0.851) (-1.150,0.815) (-1.100,0.780) (-1.059,0.750) (-1.000,0.708) (-0.953,0.675) (-0.900,0.637) (-0.850,0.602) (-0.800,0.566) (-0.750,0.531) (-0.707,0.500) (-0.650,0.460) (-0.601,0.425) (-0.550,0.389) (-0.500,0.354) (-0.450,0.319) (-0.400,0.283) (-0.354,0.250) (-0.300,0.213) (-0.250,0.177) (-0.220,0.155)}
\EQ{(-0.220,-0.155) (-0.275,-0.195) (-0.319,-0.225) (-0.375,-0.266) (-0.424,-0.300) (-0.475,-0.337) (-0.525,-0.372) (-0.575,-0.407) (-0.625,-0.443) (-0.671,-0.475) (-0.725,-0.514) (-0.775,-0.549) (-0.825,-0.585) (-0.875,-0.620) (-0.917,-0.650) (-0.975,-0.692) (-1.022,-0.725) (-1.075,-0.763) (-1.125,-0.799) (-1.175,-0.835) (-1.225,-0.872) (-1.275,-0.908) (-1.325,-0.944) (-1.368,-0.975) (-1.425,-1.017) (-1.470,-1.050) (-1.525,-1.091) (-1.571,-1.125) (-1.625,-1.166) (-1.671,-1.200) (-1.725,-1.241) (-1.770,-1.275) (-1.825,-1.317) (-1.869,-1.350) (-1.925,-1.394) (-1.965,-1.425) (-2.025,-1.472) (-2.075,-1.512) (-2.123,-1.550) (-2.175,-1.592) (-2.216,-1.625) (-2.275,-1.674) (-2.325,-1.715) (-2.367,-1.750) (-2.425,-1.799) (-2.475,-1.843) (-2.525,-1.886) (-2.570,-1.925) (-2.625,-1.975) (-2.675,-2.020) (-2.725,-2.066) (-2.762,-2.100) (-2.800,-2.136) (-2.841,-2.175) (-2.893,-2.225) (-2.943,-2.275) (-2.993,-2.325) (-3.042,-2.375) (-3.090,-2.425) (-3.125,-2.463) (-3.160,-2.500) (-3.200,-2.545) (-3.249,-2.600) (-3.293,-2.650) (-3.336,-2.700) (-3.375,-2.747) (-3.418,-2.800) (-3.450,-2.840) (-3.497,-2.900) (-3.535,-2.950) (-3.572,-3.000) (-3.608,-3.050) (-3.644,-3.100) (-3.675,-3.146) (-3.712,-3.200) (-3.744,-3.250) (-3.775,-3.299) (-3.807,-3.350) (-3.837,-3.400) (-3.866,-3.450) (-3.894,-3.500) (-3.922,-3.550) (-3.948,-3.600) (-3.974,-3.650) (-3.999,-3.700) (-4.023,-3.750) (-4.047,-3.800) (-4.069,-3.850) (-4.091,-3.900) (-4.112,-3.950) (-4.132,-4.000) (-4.152,-4.050) (-4.170,-4.100) (-4.188,-4.150) (-4.206,-4.200) (-4.222,-4.250) (-4.238,-4.300) (-4.253,-4.350) (-4.267,-4.400) (-4.280,-4.450) (-4.293,-4.500) (-4.305,-4.550) (-4.316,-4.600) (-4.327,-4.650) (-4.337,-4.700) (-4.346,-4.750) (-4.354,-4.800) (-4.362,-4.850) (-4.369,-4.900) (-4.376,-4.950) (-4.381,-5.000) (-4.386,-5.050) (-4.391,-5.100) (-4.394,-5.150) (-4.397,-5.200) (-4.400,-5.250) (-4.401,-5.300) (-4.402,-5.350) (-4.402,-5.400) (-4.402,-5.450) (-4.401,-5.500) (-4.399,-5.550) (-4.396,-5.600) (-4.393,-5.650) (-4.389,-5.700) (-4.385,-5.750) (-4.380,-5.800) (-4.374,-5.850) (-4.367,-5.900) (-4.360,-5.950) (-4.352,-6.000) (-4.343,-6.050) (-4.334,-6.100) (-4.323,-6.150) (-4.313,-6.200) (-4.301,-6.250) (-4.289,-6.300) (-4.276,-6.350) (-4.262,-6.400) (-4.247,-6.450) (-4.232,-6.500) (-4.216,-6.550) (-4.200,-6.598) (-4.182,-6.650) (-4.163,-6.700) (-4.144,-6.750) (-4.125,-6.798) (-4.104,-6.850) (-4.082,-6.900) (-4.060,-6.950) (-4.036,-7.000) (-4.012,-7.050) (-3.987,-7.100) (-3.961,-7.150) (-3.934,-7.200) (-3.907,-7.250) (-3.878,-7.300) (-3.850,-7.347) (-3.818,-7.400) (-3.786,-7.450) (-3.753,-7.500) (-3.725,-7.542) (-3.684,-7.600) (-3.650,-7.648) (-3.611,-7.700) (-3.575,-7.747) (-3.533,-7.800) (-3.500,-7.841) (-3.450,-7.900) (-3.406,-7.950) (-3.361,-8.000) (-3.325,-8.039) (-3.275,-8.092) (-3.225,-8.142) (-3.175,-8.191) (-3.125,-8.238) (-3.085,-8.275) (-3.029,-8.325) (-2.975,-8.372) (-2.925,-8.414) (-2.880,-8.450) (-2.825,-8.494) (-2.784,-8.525) (-2.725,-8.569) (-2.682,-8.600) (-2.625,-8.640) (-2.575,-8.674) (-2.525,-8.707) (-2.475,-8.739) (-2.425,-8.770) (-2.375,-8.800) (-2.325,-8.829) (-2.275,-8.858) (-2.225,-8.886) (-2.175,-8.912) (-2.125,-8.938) (-2.075,-8.964) (-2.025,-8.988) (-1.975,-9.012) (-1.925,-9.035) (-1.875,-9.057) (-1.825,-9.078) (-1.775,-9.099) (-1.725,-9.119) (-1.675,-9.139) (-1.625,-9.158) (-1.577,-9.175) (-1.525,-9.193) (-1.475,-9.210) (-1.425,-9.226) (-1.375,-9.242) (-1.325,-9.257) (-1.275,-9.272) (-1.225,-9.285) (-1.175,-9.298) (-1.125,-9.311) (-1.075,-9.323) (-1.025,-9.335) (-0.975,-9.345) (-0.925,-9.356) (-0.875,-9.365) (-0.825,-9.375) (-0.775,-9.383) (-0.725,-9.391) (-0.675,-9.399) (-0.625,-9.406) (-0.575,-9.412) (-0.525,-9.418) (-0.475,-9.423) (-0.425,-9.428) (-0.375,-9.432) (-0.325,-9.436) (-0.275,-9.439) (-0.225,-9.442) (-0.175,-9.444) (-0.125,-9.446) (-0.075,-9.447) (-0.025,-9.447) (0.050,-9.447) (0.100,-9.446) (0.150,-9.445) (0.200,-9.443) (0.250,-9.441) (0.300,-9.438) (0.350,-9.434) (0.400,-9.430) (0.450,-9.426) (0.500,-9.421) (0.550,-9.415) (0.600,-9.409) (0.650,-9.402) (0.700,-9.395) (0.750,-9.387) (0.800,-9.379) (0.850,-9.370) (0.900,-9.361) (0.950,-9.351) (1.000,-9.340) (1.050,-9.329) (1.100,-9.317) (1.150,-9.305) (1.200,-9.292) (1.250,-9.278) (1.300,-9.264) (1.349,-9.250) (1.400,-9.234) (1.450,-9.218) (1.500,-9.202) (1.550,-9.185) (1.600,-9.167) (1.650,-9.148) (1.700,-9.129) (1.750,-9.109) (1.800,-9.089) (1.850,-9.068) (1.900,-9.046) (1.946,-9.025) (2.000,-9.000) (2.050,-8.976) (2.100,-8.951) (2.150,-8.925) (2.198,-8.900) (2.244,-8.875) (2.289,-8.850) (2.333,-8.825) (2.400,-8.785) (2.450,-8.755) (2.497,-8.725) (2.550,-8.691) (2.600,-8.657) (2.647,-8.625) (2.700,-8.587) (2.750,-8.551) (2.800,-8.513) (2.849,-8.475) (2.900,-8.434) (2.942,-8.400) (3.000,-8.351) (3.050,-8.307) (3.100,-8.262) (3.139,-8.225) (3.192,-8.175) (3.242,-8.125) (3.291,-8.075) (3.338,-8.025) (3.375,-7.985) (3.425,-7.929) (3.471,-7.875) (3.513,-7.825) (3.550,-7.779) (3.592,-7.725) (3.625,-7.682) (3.667,-7.625) (3.700,-7.578) (3.736,-7.525) (3.770,-7.475) (3.800,-7.428) (3.833,-7.375) (3.863,-7.325) (3.892,-7.275) (3.921,-7.225) (3.948,-7.175) (3.974,-7.125) (4.000,-7.075) (4.024,-7.025) (4.048,-6.975) (4.071,-6.925) (4.093,-6.875) (4.114,-6.825) (4.134,-6.775) (4.154,-6.725) (4.173,-6.675) (4.191,-6.625) (4.208,-6.575) (4.224,-6.525) (4.240,-6.475) (4.255,-6.425) (4.269,-6.375) (4.282,-6.325) (4.295,-6.275) (4.307,-6.225) (4.318,-6.175) (4.329,-6.125) (4.338,-6.075) (4.347,-6.025) (4.356,-5.975) (4.364,-5.925) (4.370,-5.875) (4.377,-5.825) (4.382,-5.775) (4.387,-5.725) (4.391,-5.675) (4.395,-5.625) (4.398,-5.575) (4.400,-5.525) (4.401,-5.475) (4.402,-5.425) (4.402,-5.375) (4.402,-5.325) (4.401,-5.275) (4.399,-5.225) (4.396,-5.175) (4.393,-5.125) (4.389,-5.075) (4.384,-5.025) (4.379,-4.975) (4.373,-4.925) (4.366,-4.875) (4.358,-4.825) (4.350,-4.775) (4.342,-4.725) (4.332,-4.675) (4.322,-4.625) (4.311,-4.575) (4.299,-4.525) (4.287,-4.475) (4.274,-4.425) (4.260,-4.375) (4.245,-4.325) (4.230,-4.275) (4.214,-4.225) (4.197,-4.175) (4.180,-4.125) (4.161,-4.075) (4.142,-4.025) (4.122,-3.975) (4.102,-3.925) (4.080,-3.875) (4.058,-3.825) (4.035,-3.775) (4.011,-3.725) (3.987,-3.675) (3.961,-3.625) (3.935,-3.575) (3.908,-3.525) (3.880,-3.475) (3.852,-3.425) (3.825,-3.380) (3.792,-3.325) (3.760,-3.275) (3.728,-3.225) (3.700,-3.183) (3.661,-3.125) (3.626,-3.075) (3.590,-3.025) (3.554,-2.975) (3.516,-2.925) (3.478,-2.875) (3.438,-2.825) (3.400,-2.778) (3.357,-2.725) (3.315,-2.675) (3.275,-2.629) (3.227,-2.575) (3.182,-2.525) (3.136,-2.475) (3.100,-2.436) (3.050,-2.384) (3.000,-2.332) (2.950,-2.282) (2.900,-2.232) (2.850,-2.184) (2.800,-2.136) (2.762,-2.100) (2.725,-2.066) (2.681,-2.025) (2.626,-1.975) (2.575,-1.930) (2.525,-1.886) (2.484,-1.850) (2.426,-1.800) (2.375,-1.757) (2.325,-1.715) (2.277,-1.675) (2.225,-1.633) (2.185,-1.600) (2.125,-1.552) (2.075,-1.512) (2.029,-1.475) (1.975,-1.433) (1.933,-1.400) (1.875,-1.355) (1.825,-1.317) (1.775,-1.279) (1.725,-1.241) (1.675,-1.203) (1.625,-1.166) (1.575,-1.128) (1.525,-1.091) (1.475,-1.054) (1.425,-1.017) (1.375,-0.981) (1.333,-0.950) (1.275,-0.908) (1.230,-0.875) (1.175,-0.835) (1.126,-0.800) (1.075,-0.763) (1.025,-0.727) (0.975,-0.692) (0.925,-0.656) (0.882,-0.625) (0.825,-0.585) (0.777,-0.550) (0.725,-0.514) (0.675,-0.478) (0.625,-0.443) (0.575,-0.407) (0.531,-0.375) (0.475,-0.337) (0.425,-0.301) (0.375,-0.266) (0.325,-0.230) (0.284,-0.200) (0.225,-0.160) (0.206,-0.144)}
\EQ{(0.206,0.144) (0.248,0.175) (0.300,0.213) (0.350,0.248) (0.400,0.283) (0.450,0.319) (0.495,0.350) (0.550,0.390) (0.600,0.425) (0.650,0.460) (0.700,0.496) (0.742,0.525) (0.800,0.567) (0.847,0.600) (0.900,0.638) (0.950,0.674) (1.000,0.710) (1.050,0.745) (1.092,0.775) (1.150,0.817) (1.196,0.850) (1.250,0.890) (1.299,0.925) (1.350,0.963) (1.400,0.999) (1.450,1.036) (1.500,1.073) (1.550,1.110) (1.600,1.147) (1.650,1.184) (1.700,1.222) (1.750,1.260) (1.800,1.298) (1.850,1.336) (1.900,1.374) (1.950,1.413) (1.997,1.450) (2.050,1.492) (2.092,1.525) (2.150,1.572) (2.200,1.612) (2.246,1.650) (2.300,1.694) (2.350,1.736) (2.396,1.775) (2.450,1.821) (2.500,1.864) (2.541,1.900) (2.598,1.950) (2.650,1.997) (2.700,2.043) (2.750,2.089) (2.789,2.125) (2.825,2.160) (2.867,2.200) (2.918,2.250) (2.968,2.300) (3.017,2.350) (3.066,2.400) (3.113,2.450) (3.150,2.490) (3.200,2.545) (3.249,2.600) (3.293,2.650) (3.336,2.700) (3.375,2.747) (3.418,2.800) (3.450,2.840) (3.497,2.900) (3.535,2.950) (3.572,3.000) (3.608,3.050) (3.644,3.100) (3.675,3.146) (3.712,3.200) (3.744,3.250) (3.775,3.299) (3.807,3.350) (3.837,3.400) (3.866,3.450) (3.894,3.500) (3.922,3.550) (3.948,3.600) (3.974,3.650) (3.999,3.700) (4.023,3.750) (4.047,3.800) (4.069,3.850) (4.091,3.900) (4.112,3.950) (4.132,4.000) (4.152,4.050) (4.170,4.100) (4.188,4.150) (4.206,4.200) (4.222,4.250) (4.238,4.300) (4.253,4.350) (4.267,4.400) (4.280,4.450) (4.293,4.500) (4.305,4.550) (4.316,4.600) (4.327,4.650) (4.337,4.700) (4.346,4.750) (4.354,4.800) (4.362,4.850) (4.369,4.900) (4.376,4.950) (4.381,5.000) (4.386,5.050) (4.391,5.100) (4.394,5.150) (4.397,5.200) (4.400,5.250) (4.401,5.300) (4.402,5.350) (4.402,5.400) (4.402,5.450) (4.401,5.500) (4.399,5.550) (4.396,5.600) (4.393,5.650) (4.389,5.700) (4.385,5.750) (4.380,5.800) (4.374,5.850) (4.367,5.900) (4.360,5.950) (4.352,6.000) (4.343,6.050) (4.334,6.100) (4.323,6.150) (4.313,6.200) (4.301,6.250) (4.289,6.300) (4.276,6.350) (4.262,6.400) (4.247,6.450) (4.232,6.500) (4.216,6.550) (4.200,6.598) (4.182,6.650) (4.163,6.700) (4.144,6.750) (4.125,6.798) (4.104,6.850) (4.082,6.900) (4.060,6.950) (4.036,7.000) (4.012,7.050) (3.987,7.100) (3.961,7.150) (3.934,7.200) (3.907,7.250) (3.878,7.300) (3.850,7.347) (3.818,7.400) (3.786,7.450) (3.753,7.500) (3.725,7.542) (3.684,7.600) (3.650,7.648) (3.611,7.700) (3.575,7.747) (3.533,7.800) (3.500,7.841) (3.450,7.900) (3.406,7.950) (3.361,8.000) (3.325,8.039) (3.275,8.092) (3.225,8.142) (3.175,8.191) (3.125,8.238) (3.085,8.275) (3.029,8.325) (2.975,8.372) (2.925,8.414) (2.880,8.450) (2.825,8.494) (2.784,8.525) (2.725,8.569) (2.682,8.600) (2.625,8.640) (2.575,8.674) (2.525,8.707) (2.475,8.739) (2.425,8.770) (2.375,8.800) (2.325,8.829) (2.275,8.858) (2.225,8.886) (2.175,8.912) (2.125,8.938) (2.075,8.964) (2.025,8.988) (1.975,9.012) (1.925,9.035) (1.875,9.057) (1.825,9.078) (1.775,9.099) (1.725,9.119) (1.675,9.139) (1.625,9.158) (1.577,9.175) (1.525,9.193) (1.475,9.210) (1.425,9.226) (1.375,9.242) (1.325,9.257) (1.275,9.272) (1.225,9.285) (1.175,9.298) (1.125,9.311) (1.075,9.323) (1.025,9.335) (0.975,9.345) (0.925,9.356) (0.875,9.365) (0.825,9.375) (0.775,9.383) (0.725,9.391) (0.675,9.399) (0.625,9.406) (0.575,9.412) (0.525,9.418) (0.475,9.423) (0.425,9.428) (0.375,9.432) (0.325,9.436) (0.275,9.439) (0.225,9.442) (0.175,9.444) (0.125,9.446) (0.075,9.447) (0.025,9.447) (-0.050,9.447) (-0.100,9.446) (-0.150,9.445) (-0.200,9.443) (-0.250,9.441) (-0.300,9.438) (-0.350,9.434) (-0.400,9.430) (-0.450,9.426) (-0.500,9.421) (-0.550,9.415) (-0.600,9.409) (-0.650,9.402) (-0.700,9.395) (-0.750,9.387) (-0.800,9.379) (-0.850,9.370) (-0.900,9.361) (-0.950,9.351) (-1.000,9.340) (-1.050,9.329) (-1.100,9.317) (-1.150,9.305) (-1.200,9.292) (-1.250,9.278) (-1.300,9.264) (-1.349,9.250) (-1.400,9.234) (-1.450,9.218) (-1.500,9.202) (-1.550,9.185) (-1.600,9.167) (-1.650,9.148) (-1.700,9.129) (-1.750,9.109) (-1.800,9.089) (-1.850,9.068) (-1.900,9.046) (-1.946,9.025) (-2.000,9.000) (-2.050,8.976) (-2.100,8.951) (-2.150,8.925) (-2.198,8.900) (-2.244,8.875) (-2.289,8.850) (-2.333,8.825) (-2.400,8.785) (-2.450,8.755) (-2.497,8.725) (-2.550,8.691) (-2.600,8.657) (-2.647,8.625) (-2.700,8.587) (-2.750,8.551) (-2.800,8.513) (-2.849,8.475) (-2.900,8.434) (-2.942,8.400) (-3.000,8.351) (-3.050,8.307) (-3.100,8.262) (-3.139,8.225) (-3.192,8.175) (-3.242,8.125) (-3.291,8.075) (-3.338,8.025) (-3.375,7.985) (-3.425,7.929) (-3.471,7.875) (-3.513,7.825) (-3.550,7.779) (-3.592,7.725) (-3.625,7.682) (-3.667,7.625) (-3.700,7.578) (-3.736,7.525) (-3.770,7.475) (-3.800,7.428) (-3.833,7.375) (-3.863,7.325) (-3.892,7.275) (-3.921,7.225) (-3.948,7.175) (-3.974,7.125) (-4.000,7.075) (-4.024,7.025) (-4.048,6.975) (-4.071,6.925) (-4.093,6.875) (-4.114,6.825) (-4.134,6.775) (-4.154,6.725) (-4.173,6.675) (-4.191,6.625) (-4.208,6.575) (-4.224,6.525) (-4.240,6.475) (-4.255,6.425) (-4.269,6.375) (-4.282,6.325) (-4.295,6.275) (-4.307,6.225) (-4.318,6.175) (-4.329,6.125) (-4.338,6.075) (-4.347,6.025) (-4.356,5.975) (-4.364,5.925) (-4.370,5.875) (-4.377,5.825) (-4.382,5.775) (-4.387,5.725) (-4.391,5.675) (-4.395,5.625) (-4.398,5.575) (-4.400,5.525) (-4.401,5.475) (-4.402,5.425) (-4.402,5.375) (-4.402,5.325) (-4.401,5.275) (-4.399,5.225) (-4.396,5.175) (-4.393,5.125) (-4.389,5.075) (-4.384,5.025) (-4.379,4.975) (-4.373,4.925) (-4.366,4.875) (-4.358,4.825) (-4.350,4.775) (-4.342,4.725) (-4.332,4.675) (-4.322,4.625) (-4.311,4.575) (-4.299,4.525) (-4.287,4.475) (-4.274,4.425) (-4.260,4.375) (-4.245,4.325) (-4.230,4.275) (-4.214,4.225) (-4.197,4.175) (-4.180,4.125) (-4.161,4.075) (-4.142,4.025) (-4.122,3.975) (-4.102,3.925) (-4.080,3.875) (-4.058,3.825) (-4.035,3.775) (-4.011,3.725) (-3.987,3.675) (-3.961,3.625) (-3.935,3.575) (-3.908,3.525) (-3.880,3.475) (-3.852,3.425) (-3.825,3.380) (-3.792,3.325) (-3.760,3.275) (-3.728,3.225) (-3.700,3.183) (-3.661,3.125) (-3.626,3.075) (-3.590,3.025) (-3.554,2.975) (-3.516,2.925) (-3.478,2.875) (-3.438,2.825) (-3.400,2.778) (-3.357,2.725) (-3.315,2.675) (-3.275,2.629) (-3.227,2.575) (-3.182,2.525) (-3.136,2.475) (-3.100,2.436) (-3.050,2.384) (-3.000,2.332) (-2.950,2.282) (-2.900,2.232) (-2.850,2.184) (-2.800,2.136) (-2.762,2.100) (-2.725,2.066) (-2.681,2.025) (-2.626,1.975) (-2.575,1.930) (-2.525,1.886) (-2.484,1.850) (-2.426,1.800) (-2.375,1.757) (-2.325,1.715) (-2.277,1.675) (-2.225,1.633) (-2.185,1.600) (-2.125,1.552) (-2.075,1.512) (-2.029,1.475) (-1.975,1.433) (-1.933,1.400) (-1.875,1.355) (-1.825,1.317) (-1.775,1.279) (-1.725,1.241) (-1.675,1.203) (-1.625,1.166) (-1.575,1.128) (-1.525,1.091) (-1.475,1.054) (-1.425,1.017) (-1.375,0.981) (-1.333,0.950) (-1.275,0.908) (-1.230,0.875) (-1.175,0.835) (-1.126,0.800) (-1.075,0.763) (-1.025,0.727) (-0.975,0.692) (-0.925,0.656) (-0.882,0.625) (-0.825,0.585) (-0.777,0.550) (-0.725,0.514) (-0.675,0.478) (-0.625,0.443) (-0.575,0.407) (-0.531,0.375) (-0.475,0.337) (-0.425,0.301) (-0.375,0.266) (-0.325,0.230) (-0.284,0.200) (-0.225,0.160) (-0.220,0.155)}
\EQ{(-0.220,-0.155) (-0.275,-0.195) (-0.319,-0.225) (-0.375,-0.266) (-0.424,-0.300) (-0.475,-0.337) (-0.525,-0.372) (-0.575,-0.407) (-0.625,-0.443) (-0.671,-0.475) (-0.725,-0.514) (-0.775,-0.549) (-0.825,-0.585) (-0.875,-0.621) (-0.916,-0.650) (-0.975,-0.693) (-1.020,-0.725) (-1.075,-0.765) (-1.124,-0.800) (-1.175,-0.838) (-1.225,-0.874) (-1.275,-0.911) (-1.325,-0.948) (-1.375,-0.985) (-1.425,-1.022) (-1.475,-1.060) (-1.525,-1.098) (-1.575,-1.136) (-1.625,-1.174) (-1.675,-1.213) (-1.723,-1.250) (-1.775,-1.291) (-1.818,-1.325) (-1.875,-1.370) (-1.925,-1.411) (-1.973,-1.450) (-2.025,-1.493) (-2.075,-1.535) (-2.122,-1.575) (-2.175,-1.621) (-2.225,-1.664) (-2.265,-1.700) (-2.321,-1.750) (-2.375,-1.800) (-2.425,-1.846) (-2.475,-1.894) (-2.525,-1.942) (-2.575,-1.991) (-2.625,-2.042) (-2.675,-2.093) (-2.725,-2.146) (-2.775,-2.200) (-2.820,-2.250) (-2.864,-2.300) (-2.900,-2.342) (-2.949,-2.400) (-2.989,-2.450) (-3.025,-2.495) (-3.068,-2.550) (-3.100,-2.593) (-3.142,-2.650) (-3.175,-2.697) (-3.212,-2.750) (-3.245,-2.800) (-3.275,-2.846) (-3.309,-2.900) (-3.340,-2.950) (-3.369,-3.000) (-3.398,-3.050) (-3.425,-3.099) (-3.450,-3.146) (-3.475,-3.195) (-3.500,-3.245) (-3.525,-3.298) (-3.549,-3.350) (-3.571,-3.400) (-3.592,-3.450) (-3.612,-3.500) (-3.631,-3.550) (-3.650,-3.600) (-3.667,-3.650) (-3.684,-3.700) (-3.700,-3.750) (-3.715,-3.800) (-3.729,-3.850) (-3.742,-3.900) (-3.755,-3.950) (-3.767,-4.000) (-3.777,-4.050) (-3.787,-4.100) (-3.797,-4.150) (-3.805,-4.200) (-3.813,-4.250) (-3.820,-4.300) (-3.826,-4.350) (-3.831,-4.400) (-3.835,-4.450) (-3.839,-4.500) (-3.842,-4.550) (-3.844,-4.600) (-3.845,-4.650) (-3.846,-4.700) (-3.846,-4.750) (-3.845,-4.800) (-3.843,-4.850) (-3.840,-4.900) (-3.837,-4.950) (-3.833,-5.000) (-3.828,-5.050) (-3.822,-5.100) (-3.816,-5.150) (-3.808,-5.200) (-3.800,-5.250) (-3.791,-5.300) (-3.782,-5.350) (-3.771,-5.400) (-3.760,-5.450) (-3.747,-5.500) (-3.734,-5.550) (-3.721,-5.600) (-3.706,-5.650) (-3.690,-5.700) (-3.674,-5.750) (-3.656,-5.800) (-3.638,-5.850) (-3.619,-5.900) (-3.600,-5.948) (-3.578,-6.000) (-3.556,-6.050) (-3.533,-6.100) (-3.509,-6.150) (-3.484,-6.200) (-3.458,-6.250) (-3.431,-6.300) (-3.403,-6.350) (-3.375,-6.399) (-3.344,-6.450) (-3.313,-6.500) (-3.280,-6.550) (-3.250,-6.594) (-3.211,-6.650) (-3.175,-6.699) (-3.137,-6.750) (-3.100,-6.797) (-3.057,-6.850) (-3.015,-6.900) (-2.975,-6.946) (-2.926,-7.000) (-2.879,-7.050) (-2.830,-7.100) (-2.779,-7.150) (-2.727,-7.200) (-2.675,-7.247) (-2.625,-7.291) (-2.585,-7.325) (-2.525,-7.375) (-2.475,-7.414) (-2.429,-7.450) (-2.375,-7.490) (-2.326,-7.525) (-2.275,-7.560) (-2.225,-7.594) (-2.176,-7.625) (-2.125,-7.657) (-2.075,-7.688) (-2.025,-7.717) (-1.975,-7.745) (-1.925,-7.772) (-1.875,-7.799) (-1.825,-7.824) (-1.775,-7.849) (-1.725,-7.873) (-1.675,-7.896) (-1.625,-7.918) (-1.575,-7.939) (-1.525,-7.959) (-1.475,-7.979) (-1.425,-7.998) (-1.375,-8.016) (-1.325,-8.034) (-1.276,-8.050) (-1.225,-8.066) (-1.175,-8.082) (-1.125,-8.096) (-1.075,-8.110) (-1.025,-8.123) (-0.975,-8.136) (-0.925,-8.148) (-0.875,-8.159) (-0.825,-8.170) (-0.775,-8.179) (-0.725,-8.189) (-0.675,-8.197) (-0.625,-8.205) (-0.575,-8.213) (-0.525,-8.219) (-0.475,-8.225) (-0.425,-8.231) (-0.375,-8.236) (-0.325,-8.240) (-0.275,-8.244) (-0.225,-8.247) (-0.175,-8.249) (-0.125,-8.251) (-0.075,-8.252) (-0.025,-8.253) (0.050,-8.253) (0.100,-8.252) (0.150,-8.250) (0.200,-8.248) (0.250,-8.245) (0.300,-8.242) (0.350,-8.238) (0.400,-8.234) (0.450,-8.228) (0.500,-8.223) (0.550,-8.216) (0.600,-8.209) (0.650,-8.201) (0.700,-8.193) (0.750,-8.184) (0.800,-8.175) (0.850,-8.164) (0.900,-8.153) (0.950,-8.142) (1.000,-8.130) (1.050,-8.117) (1.100,-8.103) (1.150,-8.089) (1.200,-8.074) (1.250,-8.058) (1.300,-8.042) (1.350,-8.025) (1.400,-8.007) (1.450,-7.989) (1.500,-7.969) (1.548,-7.950) (1.600,-7.928) (1.650,-7.907) (1.700,-7.884) (1.750,-7.861) (1.800,-7.837) (1.850,-7.812) (1.900,-7.786) (1.950,-7.759) (2.000,-7.731) (2.050,-7.702) (2.096,-7.675) (2.150,-7.642) (2.200,-7.610) (2.250,-7.577) (2.300,-7.543) (2.350,-7.508) (2.395,-7.475) (2.450,-7.434) (2.493,-7.400) (2.550,-7.354) (2.600,-7.313) (2.644,-7.275) (2.700,-7.225) (2.750,-7.178) (2.800,-7.130) (2.850,-7.080) (2.900,-7.028) (2.949,-6.975) (2.993,-6.925) (3.036,-6.875) (3.075,-6.828) (3.117,-6.775) (3.150,-6.733) (3.193,-6.675) (3.225,-6.630) (3.263,-6.575) (3.296,-6.525) (3.325,-6.480) (3.359,-6.425) (3.389,-6.375) (3.418,-6.325) (3.445,-6.275) (3.472,-6.225) (3.497,-6.175) (3.521,-6.125) (3.545,-6.075) (3.567,-6.025) (3.589,-5.975) (3.609,-5.925) (3.629,-5.875) (3.647,-5.825) (3.665,-5.775) (3.682,-5.725) (3.698,-5.675) (3.713,-5.625) (3.728,-5.575) (3.741,-5.525) (3.754,-5.475) (3.765,-5.425) (3.776,-5.375) (3.787,-5.325) (3.796,-5.275) (3.804,-5.225) (3.812,-5.175) (3.819,-5.125) (3.825,-5.075) (3.830,-5.025) (3.835,-4.975) (3.839,-4.925) (3.842,-4.875) (3.844,-4.825) (3.845,-4.775) (3.846,-4.725) (3.846,-4.675) (3.845,-4.625) (3.843,-4.575) (3.841,-4.525) (3.837,-4.475) (3.833,-4.425) (3.828,-4.375) (3.823,-4.325) (3.816,-4.275) (3.809,-4.225) (3.801,-4.175) (3.792,-4.125) (3.783,-4.075) (3.772,-4.025) (3.761,-3.975) (3.749,-3.925) (3.736,-3.875) (3.722,-3.825) (3.708,-3.775) (3.692,-3.725) (3.676,-3.675) (3.659,-3.625) (3.641,-3.575) (3.622,-3.525) (3.602,-3.475) (3.582,-3.425) (3.560,-3.375) (3.538,-3.325) (3.514,-3.275) (3.490,-3.225) (3.465,-3.175) (3.439,-3.125) (3.412,-3.075) (3.384,-3.025) (3.355,-2.975) (3.325,-2.926) (3.294,-2.875) (3.262,-2.825) (3.229,-2.775) (3.200,-2.733) (3.160,-2.675) (3.125,-2.627) (3.087,-2.575) (3.050,-2.527) (3.009,-2.475) (2.975,-2.432) (2.928,-2.375) (2.886,-2.325) (2.850,-2.284) (2.800,-2.228) (2.752,-2.175) (2.705,-2.125) (2.657,-2.075) (2.609,-2.025) (2.559,-1.975) (2.508,-1.925) (2.456,-1.875) (2.403,-1.825) (2.350,-1.777) (2.300,-1.731) (2.250,-1.687) (2.209,-1.650) (2.151,-1.600) (2.100,-1.556) (2.050,-1.514) (2.003,-1.475) (1.950,-1.431) (1.900,-1.391) (1.850,-1.350) (1.800,-1.311) (1.755,-1.275) (1.700,-1.232) (1.659,-1.200) (1.600,-1.155) (1.550,-1.117) (1.500,-1.079) (1.450,-1.041) (1.400,-1.004) (1.350,-0.966) (1.300,-0.929) (1.250,-0.893) (1.200,-0.856) (1.158,-0.825) (1.100,-0.783) (1.055,-0.750) (1.000,-0.711) (0.951,-0.675) (0.900,-0.639) (0.850,-0.603) (0.800,-0.567) (0.750,-0.532) (0.706,-0.500) (0.650,-0.461) (0.600,-0.425) (0.550,-0.390) (0.500,-0.354) (0.450,-0.319) (0.400,-0.284) (0.354,-0.250) (0.300,-0.213) (0.250,-0.177) (0.206,-0.144)}
\EQ{(0.206,0.144) (0.248,0.175) (0.300,0.213) (0.350,0.248) (0.400,0.284) (0.450,0.319) (0.495,0.350) (0.550,0.390) (0.600,0.425) (0.650,0.461) (0.700,0.496) (0.742,0.525) (0.800,0.567) (0.846,0.600) (0.900,0.639) (0.950,0.675) (1.000,0.711) (1.050,0.747) (1.100,0.783) (1.150,0.819) (1.192,0.850) (1.250,0.893) (1.294,0.925) (1.350,0.966) (1.396,1.000) (1.450,1.041) (1.495,1.075) (1.550,1.117) (1.594,1.150) (1.650,1.193) (1.691,1.225) (1.750,1.271) (1.800,1.311) (1.850,1.350) (1.900,1.391) (1.942,1.425) (2.000,1.472) (2.050,1.514) (2.093,1.550) (2.150,1.599) (2.200,1.643) (2.250,1.687) (2.293,1.725) (2.348,1.775) (2.400,1.823) (2.450,1.870) (2.500,1.918) (2.550,1.967) (2.600,2.017) (2.650,2.068) (2.700,2.120) (2.750,2.173) (2.797,2.225) (2.842,2.275) (2.875,2.313) (2.925,2.372) (2.969,2.425) (3.009,2.475) (3.049,2.525) (3.087,2.575) (3.124,2.625) (3.160,2.675) (3.195,2.725) (3.225,2.770) (3.262,2.825) (3.294,2.875) (3.325,2.925) (3.355,2.975) (3.384,3.025) (3.412,3.075) (3.439,3.125) (3.465,3.175) (3.490,3.225) (3.514,3.275) (3.538,3.325) (3.560,3.375) (3.582,3.425) (3.602,3.475) (3.622,3.525) (3.641,3.575) (3.659,3.625) (3.675,3.673) (3.692,3.725) (3.708,3.775) (3.722,3.825) (3.736,3.875) (3.749,3.925) (3.761,3.975) (3.772,4.025) (3.783,4.075) (3.792,4.125) (3.801,4.175) (3.809,4.225) (3.816,4.275) (3.823,4.325) (3.828,4.375) (3.833,4.425) (3.837,4.475) (3.841,4.525) (3.843,4.575) (3.845,4.625) (3.846,4.675) (3.846,4.725) (3.845,4.775) (3.844,4.825) (3.842,4.875) (3.839,4.925) (3.835,4.975) (3.830,5.025) (3.825,5.075) (3.819,5.125) (3.812,5.175) (3.804,5.225) (3.796,5.275) (3.787,5.325) (3.776,5.375) (3.765,5.425) (3.754,5.475) (3.741,5.525) (3.728,5.575) (3.713,5.625) (3.698,5.675) (3.682,5.725) (3.665,5.775) (3.647,5.825) (3.629,5.875) (3.609,5.925) (3.589,5.975) (3.567,6.025) (3.545,6.075) (3.521,6.125) (3.497,6.175) (3.472,6.225) (3.445,6.275) (3.418,6.325) (3.389,6.375) (3.359,6.425) (3.328,6.475) (3.300,6.519) (3.263,6.575) (3.229,6.625) (3.193,6.675) (3.156,6.725) (3.125,6.765) (3.077,6.825) (3.036,6.875) (3.000,6.917) (2.950,6.974) (2.903,7.025) (2.855,7.075) (2.805,7.125) (2.753,7.175) (2.700,7.225) (2.650,7.269) (2.600,7.313) (2.555,7.350) (2.500,7.395) (2.450,7.434) (2.400,7.471) (2.350,7.508) (2.300,7.543) (2.253,7.575) (2.200,7.610) (2.150,7.642) (2.100,7.673) (2.054,7.700) (2.000,7.731) (1.950,7.759) (1.900,7.786) (1.850,7.812) (1.800,7.837) (1.750,7.861) (1.700,7.884) (1.650,7.907) (1.600,7.928) (1.550,7.949) (1.500,7.969) (1.450,7.989) (1.400,8.007) (1.350,8.025) (1.300,8.042) (1.250,8.058) (1.200,8.074) (1.150,8.089) (1.100,8.103) (1.050,8.117) (1.000,8.130) (0.950,8.142) (0.900,8.153) (0.850,8.164) (0.800,8.175) (0.750,8.184) (0.700,8.193) (0.650,8.201) (0.600,8.209) (0.550,8.216) (0.500,8.223) (0.450,8.228) (0.400,8.234) (0.350,8.238) (0.300,8.242) (0.250,8.245) (0.200,8.248) (0.150,8.250) (0.100,8.252) (0.050,8.253) (0.000,8.253) (-0.050,8.253) (-0.100,8.252) (-0.150,8.250) (-0.200,8.248) (-0.250,8.245) (-0.300,8.242) (-0.350,8.238) (-0.400,8.234) (-0.450,8.228) (-0.500,8.223) (-0.550,8.216) (-0.600,8.209) (-0.650,8.201) (-0.700,8.193) (-0.750,8.184) (-0.800,8.175) (-0.850,8.164) (-0.900,8.153) (-0.950,8.142) (-1.000,8.130) (-1.050,8.117) (-1.100,8.103) (-1.150,8.089) (-1.200,8.074) (-1.250,8.058) (-1.300,8.042) (-1.350,8.025) (-1.400,8.007) (-1.450,7.989) (-1.500,7.969) (-1.548,7.950) (-1.600,7.928) (-1.650,7.907) (-1.700,7.884) (-1.750,7.861) (-1.800,7.837) (-1.850,7.812) (-1.900,7.786) (-1.950,7.759) (-2.000,7.731) (-2.050,7.702) (-2.096,7.675) (-2.150,7.642) (-2.200,7.610) (-2.250,7.577) (-2.300,7.543) (-2.350,7.508) (-2.395,7.475) (-2.450,7.434) (-2.493,7.400) (-2.550,7.354) (-2.600,7.313) (-2.644,7.275) (-2.700,7.225) (-2.750,7.178) (-2.800,7.130) (-2.850,7.080) (-2.900,7.028) (-2.949,6.975) (-2.993,6.925) (-3.036,6.875) (-3.075,6.828) (-3.117,6.775) (-3.150,6.733) (-3.193,6.675) (-3.225,6.630) (-3.263,6.575) (-3.296,6.525) (-3.325,6.480) (-3.359,6.425) (-3.389,6.375) (-3.418,6.325) (-3.445,6.275) (-3.472,6.225) (-3.497,6.175) (-3.521,6.125) (-3.545,6.075) (-3.567,6.025) (-3.589,5.975) (-3.609,5.925) (-3.629,5.875) (-3.647,5.825) (-3.665,5.775) (-3.682,5.725) (-3.698,5.675) (-3.713,5.625) (-3.728,5.575) (-3.741,5.525) (-3.754,5.475) (-3.765,5.425) (-3.776,5.375) (-3.787,5.325) (-3.796,5.275) (-3.804,5.225) (-3.812,5.175) (-3.819,5.125) (-3.825,5.075) (-3.830,5.025) (-3.835,4.975) (-3.839,4.925) (-3.842,4.875) (-3.844,4.825) (-3.845,4.775) (-3.846,4.725) (-3.846,4.675) (-3.845,4.625) (-3.843,4.575) (-3.841,4.525) (-3.837,4.475) (-3.833,4.425) (-3.828,4.375) (-3.823,4.325) (-3.816,4.275) (-3.809,4.225) (-3.801,4.175) (-3.792,4.125) (-3.783,4.075) (-3.772,4.025) (-3.761,3.975) (-3.749,3.925) (-3.736,3.875) (-3.722,3.825) (-3.708,3.775) (-3.692,3.725) (-3.676,3.675) (-3.659,3.625) (-3.641,3.575) (-3.622,3.525) (-3.602,3.475) (-3.582,3.425) (-3.560,3.375) (-3.538,3.325) (-3.514,3.275) (-3.490,3.225) (-3.465,3.175) (-3.439,3.125) (-3.412,3.075) (-3.384,3.025) (-3.355,2.975) (-3.325,2.926) (-3.294,2.875) (-3.262,2.825) (-3.229,2.775) (-3.200,2.733) (-3.160,2.675) (-3.125,2.627) (-3.087,2.575) (-3.050,2.527) (-3.009,2.475) (-2.975,2.432) (-2.928,2.375) (-2.886,2.325) (-2.850,2.284) (-2.800,2.228) (-2.752,2.175) (-2.705,2.125) (-2.657,2.075) (-2.609,2.025) (-2.559,1.975) (-2.508,1.925) (-2.456,1.875) (-2.403,1.825) (-2.350,1.777) (-2.300,1.731) (-2.250,1.687) (-2.209,1.650) (-2.151,1.600) (-2.100,1.556) (-2.050,1.514) (-2.003,1.475) (-1.950,1.431) (-1.900,1.391) (-1.850,1.350) (-1.800,1.311) (-1.755,1.275) (-1.700,1.232) (-1.659,1.200) (-1.600,1.155) (-1.550,1.117) (-1.500,1.079) (-1.450,1.041) (-1.400,1.004) (-1.350,0.966) (-1.300,0.929) (-1.250,0.893) (-1.200,0.856) (-1.158,0.825) (-1.100,0.783) (-1.055,0.750) (-1.000,0.711) (-0.951,0.675) (-0.900,0.639) (-0.850,0.603) (-0.800,0.567) (-0.750,0.532) (-0.706,0.500) (-0.650,0.461) (-0.600,0.425) (-0.550,0.390) (-0.500,0.354) (-0.450,0.319) (-0.400,0.284) (-0.354,0.250) (-0.300,0.213) (-0.250,0.177) (-0.220,0.155)}
\EQ{(-0.220,-0.155) (-0.275,-0.195) (-0.319,-0.225) (-0.375,-0.266) (-0.424,-0.300) (-0.475,-0.337) (-0.525,-0.372) (-0.575,-0.407) (-0.625,-0.443) (-0.671,-0.475) (-0.725,-0.514) (-0.775,-0.550) (-0.825,-0.586) (-0.875,-0.622) (-0.925,-0.658) (-0.975,-0.694) (-1.019,-0.725) (-1.075,-0.767) (-1.121,-0.800) (-1.175,-0.840) (-1.223,-0.875) (-1.275,-0.914) (-1.323,-0.950) (-1.375,-0.989) (-1.422,-1.025) (-1.475,-1.066) (-1.520,-1.100) (-1.575,-1.143) (-1.616,-1.175) (-1.675,-1.222) (-1.725,-1.263) (-1.771,-1.300) (-1.825,-1.344) (-1.875,-1.386) (-1.921,-1.425) (-1.975,-1.472) (-2.025,-1.515) (-2.064,-1.550) (-2.120,-1.600) (-2.174,-1.650) (-2.225,-1.698) (-2.275,-1.746) (-2.325,-1.795) (-2.375,-1.845) (-2.425,-1.896) (-2.475,-1.949) (-2.522,-2.000) (-2.567,-2.050) (-2.611,-2.100) (-2.650,-2.146) (-2.695,-2.200) (-2.735,-2.250) (-2.774,-2.300) (-2.812,-2.350) (-2.849,-2.400) (-2.885,-2.450) (-2.919,-2.500) (-2.950,-2.546) (-2.985,-2.600) (-3.016,-2.650) (-3.046,-2.700) (-3.075,-2.750) (-3.100,-2.794) (-3.125,-2.841) (-3.150,-2.889) (-3.175,-2.939) (-3.200,-2.991) (-3.225,-3.045) (-3.249,-3.100) (-3.270,-3.150) (-3.290,-3.200) (-3.309,-3.250) (-3.326,-3.300) (-3.343,-3.350) (-3.359,-3.400) (-3.374,-3.450) (-3.389,-3.500) (-3.402,-3.550) (-3.414,-3.600) (-3.426,-3.650) (-3.436,-3.700) (-3.446,-3.750) (-3.455,-3.800) (-3.462,-3.850) (-3.469,-3.900) (-3.476,-3.950) (-3.481,-4.000) (-3.485,-4.050) (-3.489,-4.100) (-3.491,-4.150) (-3.493,-4.200) (-3.494,-4.250) (-3.494,-4.300) (-3.493,-4.350) (-3.492,-4.400) (-3.489,-4.450) (-3.486,-4.500) (-3.482,-4.550) (-3.477,-4.600) (-3.471,-4.650) (-3.464,-4.700) (-3.456,-4.750) (-3.447,-4.800) (-3.438,-4.850) (-3.428,-4.900) (-3.416,-4.950) (-3.404,-5.000) (-3.391,-5.050) (-3.377,-5.100) (-3.362,-5.150) (-3.346,-5.200) (-3.329,-5.250) (-3.311,-5.300) (-3.292,-5.350) (-3.272,-5.400) (-3.251,-5.450) (-3.229,-5.500) (-3.206,-5.550) (-3.182,-5.600) (-3.157,-5.650) (-3.131,-5.700) (-3.103,-5.750) (-3.075,-5.799) (-3.044,-5.850) (-3.013,-5.900) (-2.981,-5.950) (-2.950,-5.995) (-2.912,-6.050) (-2.875,-6.100) (-2.837,-6.150) (-2.800,-6.196) (-2.755,-6.250) (-2.712,-6.300) (-2.675,-6.342) (-2.625,-6.396) (-2.575,-6.447) (-2.525,-6.496) (-2.475,-6.544) (-2.425,-6.589) (-2.384,-6.625) (-2.325,-6.675) (-2.275,-6.716) (-2.231,-6.750) (-2.175,-6.792) (-2.130,-6.825) (-2.075,-6.863) (-2.025,-6.897) (-1.982,-6.925) (-1.925,-6.961) (-1.875,-6.991) (-1.825,-7.020) (-1.775,-7.048) (-1.725,-7.074) (-1.676,-7.100) (-1.625,-7.125) (-1.575,-7.149) (-1.525,-7.172) (-1.475,-7.194) (-1.425,-7.215) (-1.375,-7.236) (-1.325,-7.255) (-1.275,-7.274) (-1.225,-7.292) (-1.175,-7.309) (-1.125,-7.325) (-1.075,-7.340) (-1.025,-7.355) (-0.975,-7.369) (-0.925,-7.382) (-0.875,-7.395) (-0.825,-7.406) (-0.775,-7.417) (-0.725,-7.427) (-0.675,-7.437) (-0.625,-7.446) (-0.575,-7.454) (-0.525,-7.461) (-0.475,-7.468) (-0.425,-7.474) (-0.375,-7.479) (-0.325,-7.484) (-0.275,-7.488) (-0.225,-7.492) (-0.175,-7.494) (-0.125,-7.496) (-0.075,-7.498) (-0.025,-7.498) (0.050,-7.498) (0.100,-7.497) (0.150,-7.495) (0.200,-7.493) (0.250,-7.490) (0.300,-7.486) (0.350,-7.482) (0.400,-7.477) (0.450,-7.471) (0.500,-7.465) (0.550,-7.458) (0.599,-7.450) (0.650,-7.441) (0.700,-7.432) (0.750,-7.422) (0.800,-7.412) (0.850,-7.400) (0.900,-7.388) (0.950,-7.376) (1.000,-7.362) (1.050,-7.348) (1.100,-7.333) (1.150,-7.317) (1.200,-7.300) (1.250,-7.283) (1.300,-7.265) (1.350,-7.246) (1.400,-7.226) (1.450,-7.205) (1.500,-7.183) (1.550,-7.161) (1.600,-7.137) (1.650,-7.113) (1.700,-7.088) (1.750,-7.061) (1.800,-7.034) (1.850,-7.005) (1.900,-6.976) (1.950,-6.945) (2.000,-6.913) (2.050,-6.880) (2.094,-6.850) (2.150,-6.811) (2.198,-6.775) (2.250,-6.735) (2.295,-6.700) (2.350,-6.654) (2.400,-6.611) (2.441,-6.575) (2.495,-6.525) (2.547,-6.475) (2.597,-6.425) (2.644,-6.375) (2.690,-6.325) (2.725,-6.286) (2.775,-6.227) (2.817,-6.175) (2.850,-6.133) (2.893,-6.075) (2.925,-6.031) (2.964,-5.975) (2.997,-5.925) (3.025,-5.881) (3.060,-5.825) (3.089,-5.775) (3.117,-5.725) (3.144,-5.675) (3.170,-5.625) (3.195,-5.575) (3.218,-5.525) (3.241,-5.475) (3.262,-5.425) (3.282,-5.375) (3.302,-5.325) (3.320,-5.275) (3.338,-5.225) (3.354,-5.175) (3.370,-5.125) (3.384,-5.075) (3.398,-5.025) (3.410,-4.975) (3.422,-4.925) (3.433,-4.875) (3.443,-4.825) (3.452,-4.775) (3.460,-4.725) (3.467,-4.675) (3.474,-4.625) (3.479,-4.575) (3.484,-4.525) (3.488,-4.475) (3.491,-4.425) (3.493,-4.375) (3.494,-4.325) (3.494,-4.275) (3.494,-4.225) (3.492,-4.175) (3.490,-4.125) (3.487,-4.075) (3.483,-4.025) (3.478,-3.975) (3.473,-3.925) (3.466,-3.875) (3.459,-3.825) (3.450,-3.775) (3.441,-3.725) (3.431,-3.675) (3.420,-3.625) (3.408,-3.575) (3.395,-3.525) (3.382,-3.475) (3.367,-3.425) (3.352,-3.375) (3.335,-3.325) (3.318,-3.275) (3.300,-3.227) (3.280,-3.175) (3.260,-3.125) (3.238,-3.075) (3.216,-3.025) (3.193,-2.975) (3.168,-2.925) (3.143,-2.875) (3.117,-2.825) (3.089,-2.775) (3.061,-2.725) (3.031,-2.675) (3.001,-2.625) (2.969,-2.575) (2.936,-2.525) (2.902,-2.475) (2.867,-2.425) (2.831,-2.375) (2.800,-2.334) (2.755,-2.275) (2.715,-2.225) (2.675,-2.176) (2.632,-2.125) (2.589,-2.075) (2.550,-2.031) (2.500,-1.976) (2.452,-1.925) (2.404,-1.875) (2.355,-1.825) (2.305,-1.775) (2.254,-1.725) (2.201,-1.675) (2.150,-1.628) (2.100,-1.582) (2.050,-1.537) (2.008,-1.500) (1.950,-1.450) (1.900,-1.407) (1.850,-1.365) (1.802,-1.325) (1.750,-1.283) (1.710,-1.250) (1.650,-1.202) (1.600,-1.163) (1.552,-1.125) (1.500,-1.085) (1.455,-1.050) (1.400,-1.008) (1.357,-0.975) (1.300,-0.933) (1.257,-0.900) (1.200,-0.858) (1.155,-0.825) (1.100,-0.785) (1.053,-0.750) (1.000,-0.712) (0.950,-0.676) (0.900,-0.640) (0.850,-0.604) (0.800,-0.568) (0.750,-0.532) (0.706,-0.500) (0.650,-0.461) (0.600,-0.425) (0.550,-0.390) (0.500,-0.354) (0.450,-0.319) (0.400,-0.284) (0.354,-0.250) (0.300,-0.213) (0.250,-0.177) (0.206,-0.144)}
\EQ{(0.206,0.144) (0.248,0.175) (0.300,0.213) (0.350,0.248) (0.400,0.284) (0.450,0.319) (0.495,0.350) (0.550,0.390) (0.600,0.425) (0.650,0.461) (0.700,0.496) (0.741,0.525) (0.800,0.568) (0.845,0.600) (0.900,0.640) (0.949,0.675) (1.000,0.712) (1.050,0.748) (1.100,0.785) (1.150,0.821) (1.200,0.858) (1.250,0.895) (1.300,0.933) (1.350,0.970) (1.400,1.008) (1.450,1.046) (1.500,1.085) (1.550,1.124) (1.600,1.163) (1.647,1.200) (1.700,1.242) (1.741,1.275) (1.800,1.324) (1.850,1.365) (1.892,1.400) (1.950,1.450) (2.000,1.493) (2.050,1.537) (2.092,1.575) (2.147,1.625) (2.200,1.674) (2.250,1.722) (2.300,1.770) (2.350,1.820) (2.400,1.871) (2.450,1.923) (2.499,1.975) (2.545,2.025) (2.589,2.075) (2.625,2.117) (2.674,2.175) (2.715,2.225) (2.750,2.269) (2.793,2.325) (2.825,2.367) (2.867,2.425) (2.900,2.472) (2.936,2.525) (2.969,2.575) (3.000,2.624) (3.031,2.675) (3.061,2.725) (3.089,2.775) (3.117,2.825) (3.143,2.875) (3.168,2.925) (3.193,2.975) (3.216,3.025) (3.238,3.075) (3.260,3.125) (3.280,3.175) (3.299,3.225) (3.318,3.275) (3.335,3.325) (3.352,3.375) (3.367,3.425) (3.382,3.475) (3.395,3.525) (3.408,3.575) (3.420,3.625) (3.431,3.675) (3.441,3.725) (3.450,3.775) (3.459,3.825) (3.466,3.875) (3.473,3.925) (3.478,3.975) (3.483,4.025) (3.487,4.075) (3.490,4.125) (3.492,4.175) (3.494,4.225) (3.494,4.275) (3.494,4.325) (3.493,4.375) (3.491,4.425) (3.488,4.475) (3.484,4.525) (3.479,4.575) (3.474,4.625) (3.467,4.675) (3.460,4.725) (3.452,4.775) (3.443,4.825) (3.433,4.875) (3.422,4.925) (3.410,4.975) (3.398,5.025) (3.384,5.075) (3.370,5.125) (3.354,5.175) (3.338,5.225) (3.320,5.275) (3.302,5.325) (3.282,5.375) (3.262,5.425) (3.241,5.475) (3.218,5.525) (3.195,5.575) (3.170,5.625) (3.144,5.675) (3.117,5.725) (3.089,5.775) (3.060,5.825) (3.029,5.875) (3.000,5.921) (2.964,5.975) (2.929,6.025) (2.900,6.066) (2.856,6.125) (2.825,6.165) (2.776,6.225) (2.734,6.275) (2.700,6.314) (2.650,6.369) (2.600,6.422) (2.550,6.472) (2.500,6.520) (2.450,6.567) (2.400,6.611) (2.355,6.650) (2.300,6.696) (2.250,6.735) (2.200,6.774) (2.150,6.811) (2.100,6.846) (2.058,6.875) (2.000,6.913) (1.950,6.945) (1.901,6.975) (1.850,7.005) (1.800,7.034) (1.750,7.061) (1.700,7.088) (1.650,7.113) (1.600,7.137) (1.550,7.161) (1.500,7.183) (1.450,7.205) (1.402,7.225) (1.350,7.246) (1.300,7.265) (1.250,7.283) (1.201,7.300) (1.150,7.317) (1.100,7.333) (1.050,7.348) (1.000,7.362) (0.950,7.376) (0.900,7.388) (0.850,7.400) (0.800,7.412) (0.750,7.422) (0.700,7.432) (0.650,7.441) (0.600,7.450) (0.550,7.458) (0.500,7.465) (0.450,7.471) (0.400,7.477) (0.350,7.482) (0.300,7.486) (0.250,7.490) (0.200,7.493) (0.150,7.495) (0.100,7.497) (0.050,7.498) (0.000,7.498) (-0.050,7.498) (-0.100,7.497) (-0.150,7.495) (-0.200,7.493) (-0.250,7.490) (-0.300,7.486) (-0.350,7.482) (-0.400,7.477) (-0.450,7.471) (-0.500,7.465) (-0.550,7.458) (-0.599,7.450) (-0.650,7.441) (-0.700,7.432) (-0.750,7.422) (-0.800,7.412) (-0.850,7.400) (-0.900,7.388) (-0.950,7.376) (-1.000,7.362) (-1.050,7.348) (-1.100,7.333) (-1.150,7.317) (-1.200,7.300) (-1.250,7.283) (-1.300,7.265) (-1.350,7.246) (-1.400,7.226) (-1.450,7.205) (-1.500,7.183) (-1.550,7.161) (-1.600,7.137) (-1.650,7.113) (-1.700,7.088) (-1.750,7.061) (-1.800,7.034) (-1.850,7.005) (-1.900,6.976) (-1.950,6.945) (-2.000,6.913) (-2.050,6.880) (-2.094,6.850) (-2.150,6.811) (-2.198,6.775) (-2.250,6.735) (-2.295,6.700) (-2.350,6.654) (-2.400,6.611) (-2.441,6.575) (-2.495,6.525) (-2.547,6.475) (-2.597,6.425) (-2.644,6.375) (-2.690,6.325) (-2.725,6.286) (-2.775,6.227) (-2.817,6.175) (-2.850,6.133) (-2.893,6.075) (-2.925,6.031) (-2.964,5.975) (-2.997,5.925) (-3.025,5.881) (-3.060,5.825) (-3.089,5.775) (-3.117,5.725) (-3.144,5.675) (-3.170,5.625) (-3.195,5.575) (-3.218,5.525) (-3.241,5.475) (-3.262,5.425) (-3.282,5.375) (-3.302,5.325) (-3.320,5.275) (-3.338,5.225) (-3.354,5.175) (-3.370,5.125) (-3.384,5.075) (-3.398,5.025) (-3.410,4.975) (-3.422,4.925) (-3.433,4.875) (-3.443,4.825) (-3.452,4.775) (-3.460,4.725) (-3.467,4.675) (-3.474,4.625) (-3.479,4.575) (-3.484,4.525) (-3.488,4.475) (-3.491,4.425) (-3.493,4.375) (-3.494,4.325) (-3.494,4.275) (-3.494,4.225) (-3.492,4.175) (-3.490,4.125) (-3.487,4.075) (-3.483,4.025) (-3.478,3.975) (-3.473,3.925) (-3.466,3.875) (-3.459,3.825) (-3.450,3.775) (-3.441,3.725) (-3.431,3.675) (-3.420,3.625) (-3.408,3.575) (-3.395,3.525) (-3.382,3.475) (-3.367,3.425) (-3.352,3.375) (-3.335,3.325) (-3.318,3.275) (-3.300,3.227) (-3.280,3.175) (-3.260,3.125) (-3.238,3.075) (-3.216,3.025) (-3.193,2.975) (-3.168,2.925) (-3.143,2.875) (-3.117,2.825) (-3.089,2.775) (-3.061,2.725) (-3.031,2.675) (-3.001,2.625) (-2.969,2.575) (-2.936,2.525) (-2.902,2.475) (-2.867,2.425) (-2.831,2.375) (-2.800,2.334) (-2.755,2.275) (-2.715,2.225) (-2.675,2.176) (-2.632,2.125) (-2.589,2.075) (-2.550,2.031) (-2.500,1.976) (-2.452,1.925) (-2.404,1.875) (-2.355,1.825) (-2.305,1.775) (-2.254,1.725) (-2.201,1.675) (-2.150,1.628) (-2.100,1.582) (-2.050,1.537) (-2.008,1.500) (-1.950,1.450) (-1.900,1.407) (-1.850,1.365) (-1.802,1.325) (-1.750,1.283) (-1.710,1.250) (-1.650,1.202) (-1.600,1.163) (-1.552,1.125) (-1.500,1.085) (-1.455,1.050) (-1.400,1.008) (-1.357,0.975) (-1.300,0.933) (-1.257,0.900) (-1.200,0.858) (-1.155,0.825) (-1.100,0.785) (-1.053,0.750) (-1.000,0.712) (-0.950,0.676) (-0.900,0.640) (-0.850,0.604) (-0.800,0.568) (-0.750,0.532) (-0.706,0.500) (-0.650,0.461) (-0.600,0.425) (-0.550,0.390) (-0.500,0.354) (-0.450,0.319) (-0.400,0.284) (-0.354,0.250) (-0.300,0.213) (-0.250,0.177) (-0.220,0.155)}
\EQ{(-0.220,-0.155) (-0.275,-0.195) (-0.319,-0.225) (-0.375,-0.266) (-0.424,-0.300) (-0.475,-0.337) (-0.525,-0.372) (-0.575,-0.408) (-0.625,-0.443) (-0.671,-0.475) (-0.725,-0.515) (-0.775,-0.550) (-0.825,-0.586) (-0.875,-0.622) (-0.925,-0.659) (-0.975,-0.695) (-1.017,-0.725) (-1.075,-0.768) (-1.119,-0.800) (-1.175,-0.842) (-1.219,-0.875) (-1.275,-0.917) (-1.319,-0.950) (-1.375,-0.994) (-1.416,-1.025) (-1.475,-1.071) (-1.525,-1.111) (-1.574,-1.150) (-1.625,-1.191) (-1.667,-1.225) (-1.725,-1.274) (-1.775,-1.316) (-1.815,-1.350) (-1.872,-1.400) (-1.925,-1.447) (-1.975,-1.492) (-2.025,-1.538) (-2.064,-1.575) (-2.117,-1.625) (-2.167,-1.675) (-2.217,-1.725) (-2.265,-1.775) (-2.300,-1.812) (-2.335,-1.850) (-2.375,-1.894) (-2.424,-1.950) (-2.467,-2.000) (-2.500,-2.040) (-2.548,-2.100) (-2.587,-2.150) (-2.624,-2.200) (-2.661,-2.250) (-2.696,-2.300) (-2.725,-2.343) (-2.762,-2.400) (-2.794,-2.450) (-2.824,-2.500) (-2.850,-2.545) (-2.875,-2.589) (-2.900,-2.635) (-2.925,-2.683) (-2.950,-2.733) (-2.975,-2.786) (-3.000,-2.841) (-3.025,-2.899) (-3.046,-2.950) (-3.065,-3.000) (-3.083,-3.050) (-3.100,-3.099) (-3.117,-3.150) (-3.132,-3.200) (-3.146,-3.250) (-3.159,-3.300) (-3.171,-3.350) (-3.183,-3.400) (-3.193,-3.450) (-3.202,-3.500) (-3.211,-3.550) (-3.218,-3.600) (-3.224,-3.650) (-3.230,-3.700) (-3.235,-3.750) (-3.238,-3.800) (-3.241,-3.850) (-3.243,-3.900) (-3.244,-3.950) (-3.244,-4.000) (-3.243,-4.050) (-3.241,-4.100) (-3.238,-4.150) (-3.234,-4.200) (-3.229,-4.250) (-3.224,-4.300) (-3.217,-4.350) (-3.210,-4.400) (-3.201,-4.450) (-3.192,-4.500) (-3.182,-4.550) (-3.170,-4.600) (-3.158,-4.650) (-3.145,-4.700) (-3.130,-4.750) (-3.115,-4.800) (-3.099,-4.850) (-3.081,-4.900) (-3.063,-4.950) (-3.043,-5.000) (-3.022,-5.050) (-3.001,-5.100) (-2.978,-5.150) (-2.954,-5.200) (-2.928,-5.250) (-2.902,-5.300) (-2.875,-5.348) (-2.845,-5.400) (-2.814,-5.450) (-2.782,-5.500) (-2.750,-5.549) (-2.715,-5.600) (-2.678,-5.650) (-2.640,-5.700) (-2.601,-5.750) (-2.560,-5.800) (-2.525,-5.840) (-2.475,-5.896) (-2.425,-5.949) (-2.375,-6.000) (-2.325,-6.049) (-2.275,-6.096) (-2.225,-6.140) (-2.184,-6.175) (-2.125,-6.224) (-2.075,-6.263) (-2.026,-6.300) (-1.975,-6.337) (-1.925,-6.372) (-1.875,-6.406) (-1.825,-6.438) (-1.775,-6.469) (-1.725,-6.499) (-1.679,-6.525) (-1.634,-6.550) (-1.587,-6.575) (-1.538,-6.600) (-1.487,-6.625) (-1.433,-6.650) (-1.377,-6.675) (-1.325,-6.697) (-1.275,-6.718) (-1.225,-6.737) (-1.175,-6.756) (-1.125,-6.773) (-1.075,-6.790) (-1.025,-6.806) (-0.975,-6.821) (-0.925,-6.835) (-0.875,-6.849) (-0.825,-6.861) (-0.775,-6.873) (-0.725,-6.884) (-0.675,-6.895) (-0.625,-6.904) (-0.575,-6.913) (-0.525,-6.921) (-0.475,-6.928) (-0.425,-6.935) (-0.375,-6.941) (-0.325,-6.946) (-0.275,-6.950) (-0.225,-6.954) (-0.175,-6.956) (-0.125,-6.959) (-0.075,-6.960) (-0.025,-6.961) (0.050,-6.960) (0.100,-6.959) (0.150,-6.958) (0.200,-6.955) (0.250,-6.952) (0.300,-6.948) (0.350,-6.943) (0.400,-6.938) (0.450,-6.932) (0.500,-6.925) (0.550,-6.917) (0.600,-6.909) (0.650,-6.899) (0.700,-6.890) (0.750,-6.879) (0.800,-6.867) (0.850,-6.855) (0.900,-6.842) (0.950,-6.828) (1.000,-6.814) (1.050,-6.798) (1.100,-6.782) (1.150,-6.765) (1.200,-6.747) (1.250,-6.728) (1.300,-6.708) (1.350,-6.687) (1.400,-6.665) (1.450,-6.642) (1.500,-6.619) (1.550,-6.594) (1.600,-6.568) (1.650,-6.541) (1.700,-6.513) (1.750,-6.484) (1.800,-6.454) (1.845,-6.425) (1.900,-6.389) (1.950,-6.355) (1.992,-6.325) (2.050,-6.282) (2.092,-6.250) (2.150,-6.204) (2.200,-6.162) (2.242,-6.125) (2.297,-6.075) (2.350,-6.025) (2.400,-5.975) (2.448,-5.925) (2.494,-5.875) (2.538,-5.825) (2.575,-5.782) (2.621,-5.725) (2.660,-5.675) (2.697,-5.625) (2.732,-5.575) (2.766,-5.525) (2.799,-5.475) (2.830,-5.425) (2.859,-5.375) (2.888,-5.325) (2.915,-5.275) (2.941,-5.225) (2.966,-5.175) (2.989,-5.125) (3.012,-5.075) (3.033,-5.025) (3.053,-4.975) (3.072,-4.925) (3.090,-4.875) (3.107,-4.825) (3.123,-4.775) (3.138,-4.725) (3.151,-4.675) (3.164,-4.625) (3.176,-4.575) (3.187,-4.525) (3.197,-4.475) (3.206,-4.425) (3.214,-4.375) (3.221,-4.325) (3.227,-4.275) (3.232,-4.225) (3.236,-4.175) (3.239,-4.125) (3.242,-4.075) (3.243,-4.025) (3.244,-3.975) (3.243,-3.925) (3.242,-3.875) (3.240,-3.825) (3.237,-3.775) (3.232,-3.725) (3.227,-3.675) (3.221,-3.625) (3.214,-3.575) (3.207,-3.525) (3.198,-3.475) (3.188,-3.425) (3.177,-3.375) (3.165,-3.325) (3.153,-3.275) (3.139,-3.225) (3.124,-3.175) (3.109,-3.125) (3.092,-3.075) (3.075,-3.027) (3.056,-2.975) (3.036,-2.925) (3.015,-2.875) (2.993,-2.825) (2.970,-2.775) (2.946,-2.725) (2.921,-2.675) (2.895,-2.625) (2.867,-2.575) (2.839,-2.525) (2.809,-2.475) (2.778,-2.425) (2.750,-2.381) (2.713,-2.325) (2.678,-2.275) (2.643,-2.225) (2.606,-2.175) (2.575,-2.135) (2.528,-2.075) (2.488,-2.025) (2.450,-1.980) (2.402,-1.925) (2.358,-1.875) (2.312,-1.825) (2.275,-1.785) (2.225,-1.733) (2.175,-1.683) (2.125,-1.633) (2.075,-1.585) (2.038,-1.550) (2.000,-1.515) (1.957,-1.475) (1.901,-1.425) (1.850,-1.381) (1.800,-1.337) (1.756,-1.300) (1.700,-1.253) (1.650,-1.212) (1.605,-1.175) (1.550,-1.131) (1.500,-1.091) (1.450,-1.052) (1.400,-1.013) (1.351,-0.975) (1.300,-0.936) (1.252,-0.900) (1.200,-0.861) (1.152,-0.825) (1.100,-0.787) (1.051,-0.750) (1.000,-0.713) (0.950,-0.677) (0.900,-0.640) (0.850,-0.604) (0.800,-0.568) (0.750,-0.532) (0.705,-0.500) (0.650,-0.461) (0.600,-0.425) (0.550,-0.390) (0.500,-0.354) (0.450,-0.319) (0.400,-0.284) (0.354,-0.250) (0.300,-0.213) (0.250,-0.177) (0.206,-0.144)}
\EQ{(0.206,0.144) (0.248,0.175) (0.300,0.213) (0.350,0.248) (0.400,0.284) (0.450,0.319) (0.495,0.350) (0.550,0.390) (0.600,0.425) (0.650,0.461) (0.700,0.497) (0.750,0.532) (0.800,0.568) (0.845,0.600) (0.900,0.640) (0.948,0.675) (1.000,0.713) (1.050,0.750) (1.100,0.787) (1.150,0.823) (1.200,0.861) (1.250,0.898) (1.300,0.936) (1.350,0.974) (1.400,1.013) (1.448,1.050) (1.500,1.091) (1.543,1.125) (1.600,1.171) (1.650,1.212) (1.697,1.250) (1.750,1.295) (1.800,1.337) (1.844,1.375) (1.900,1.424) (1.950,1.469) (2.000,1.515) (2.038,1.550) (2.075,1.585) (2.117,1.625) (2.167,1.675) (2.217,1.725) (2.265,1.775) (2.300,1.812) (2.335,1.850) (2.375,1.894) (2.424,1.950) (2.467,2.000) (2.500,2.040) (2.548,2.100) (2.587,2.150) (2.624,2.200) (2.661,2.250) (2.696,2.300) (2.725,2.343) (2.762,2.400) (2.794,2.450) (2.824,2.500) (2.850,2.545) (2.875,2.589) (2.900,2.635) (2.925,2.683) (2.950,2.733) (2.975,2.786) (3.000,2.841) (3.025,2.899) (3.046,2.950) (3.065,3.000) (3.083,3.050) (3.100,3.099) (3.117,3.150) (3.132,3.200) (3.146,3.250) (3.159,3.300) (3.171,3.350) (3.183,3.400) (3.193,3.450) (3.202,3.500) (3.211,3.550) (3.218,3.600) (3.224,3.650) (3.230,3.700) (3.235,3.750) (3.238,3.800) (3.241,3.850) (3.243,3.900) (3.244,3.950) (3.244,4.000) (3.243,4.050) (3.241,4.100) (3.238,4.150) (3.234,4.200) (3.229,4.250) (3.224,4.300) (3.217,4.350) (3.210,4.400) (3.201,4.450) (3.192,4.500) (3.182,4.550) (3.170,4.600) (3.158,4.650) (3.145,4.700) (3.130,4.750) (3.115,4.800) (3.099,4.850) (3.081,4.900) (3.063,4.950) (3.043,5.000) (3.022,5.050) (3.001,5.100) (2.978,5.150) (2.954,5.200) (2.928,5.250) (2.902,5.300) (2.875,5.348) (2.845,5.400) (2.814,5.450) (2.782,5.500) (2.750,5.549) (2.715,5.600) (2.678,5.650) (2.640,5.700) (2.601,5.750) (2.560,5.800) (2.525,5.840) (2.475,5.896) (2.425,5.949) (2.375,6.000) (2.325,6.049) (2.275,6.096) (2.225,6.140) (2.184,6.175) (2.125,6.224) (2.075,6.263) (2.026,6.300) (1.975,6.337) (1.925,6.372) (1.875,6.406) (1.825,6.438) (1.775,6.469) (1.725,6.499) (1.679,6.525) (1.634,6.550) (1.587,6.575) (1.538,6.600) (1.487,6.625) (1.433,6.650) (1.377,6.675) (1.325,6.697) (1.275,6.718) (1.225,6.737) (1.175,6.756) (1.125,6.773) (1.075,6.790) (1.025,6.806) (0.975,6.821) (0.925,6.835) (0.875,6.849) (0.825,6.861) (0.775,6.873) (0.725,6.884) (0.675,6.895) (0.625,6.904) (0.575,6.913) (0.525,6.921) (0.475,6.928) (0.425,6.935) (0.375,6.941) (0.325,6.946) (0.275,6.950) (0.225,6.954) (0.175,6.956) (0.125,6.959) (0.075,6.960) (0.025,6.961) (-0.050,6.960) (-0.100,6.959) (-0.150,6.958) (-0.200,6.955) (-0.250,6.952) (-0.300,6.948) (-0.350,6.943) (-0.400,6.938) (-0.450,6.932) (-0.500,6.925) (-0.550,6.917) (-0.600,6.909) (-0.650,6.899) (-0.700,6.890) (-0.750,6.879) (-0.800,6.867) (-0.850,6.855) (-0.900,6.842) (-0.950,6.828) (-1.000,6.814) (-1.050,6.798) (-1.100,6.782) (-1.150,6.765) (-1.200,6.747) (-1.250,6.728) (-1.300,6.708) (-1.350,6.687) (-1.400,6.665) (-1.450,6.642) (-1.500,6.619) (-1.550,6.594) (-1.600,6.568) (-1.650,6.541) (-1.700,6.513) (-1.750,6.484) (-1.800,6.454) (-1.845,6.425) (-1.900,6.389) (-1.950,6.355) (-1.992,6.325) (-2.050,6.282) (-2.092,6.250) (-2.150,6.204) (-2.200,6.162) (-2.242,6.125) (-2.297,6.075) (-2.350,6.025) (-2.400,5.975) (-2.448,5.925) (-2.494,5.875) (-2.538,5.825) (-2.575,5.782) (-2.621,5.725) (-2.660,5.675) (-2.697,5.625) (-2.732,5.575) (-2.766,5.525) (-2.799,5.475) (-2.830,5.425) (-2.859,5.375) (-2.888,5.325) (-2.915,5.275) (-2.941,5.225) (-2.966,5.175) (-2.989,5.125) (-3.012,5.075) (-3.033,5.025) (-3.053,4.975) (-3.072,4.925) (-3.090,4.875) (-3.107,4.825) (-3.123,4.775) (-3.138,4.725) (-3.151,4.675) (-3.164,4.625) (-3.176,4.575) (-3.187,4.525) (-3.197,4.475) (-3.206,4.425) (-3.214,4.375) (-3.221,4.325) (-3.227,4.275) (-3.232,4.225) (-3.236,4.175) (-3.239,4.125) (-3.242,4.075) (-3.243,4.025) (-3.244,3.975) (-3.243,3.925) (-3.242,3.875) (-3.240,3.825) (-3.237,3.775) (-3.232,3.725) (-3.227,3.675) (-3.221,3.625) (-3.214,3.575) (-3.207,3.525) (-3.198,3.475) (-3.188,3.425) (-3.177,3.375) (-3.165,3.325) (-3.153,3.275) (-3.139,3.225) (-3.124,3.175) (-3.109,3.125) (-3.092,3.075) (-3.075,3.027) (-3.056,2.975) (-3.036,2.925) (-3.015,2.875) (-2.993,2.825) (-2.970,2.775) (-2.946,2.725) (-2.921,2.675) (-2.895,2.625) (-2.867,2.575) (-2.839,2.525) (-2.809,2.475) (-2.778,2.425) (-2.750,2.381) (-2.713,2.325) (-2.678,2.275) (-2.643,2.225) (-2.606,2.175) (-2.575,2.135) (-2.528,2.075) (-2.488,2.025) (-2.450,1.980) (-2.402,1.925) (-2.358,1.875) (-2.312,1.825) (-2.275,1.785) (-2.225,1.733) (-2.175,1.683) (-2.125,1.633) (-2.075,1.585) (-2.038,1.550) (-2.000,1.515) (-1.957,1.475) (-1.901,1.425) (-1.850,1.381) (-1.800,1.337) (-1.756,1.300) (-1.700,1.253) (-1.650,1.212) (-1.605,1.175) (-1.550,1.131) (-1.500,1.091) (-1.450,1.052) (-1.400,1.013) (-1.351,0.975) (-1.300,0.936) (-1.252,0.900) (-1.200,0.861) (-1.152,0.825) (-1.100,0.787) (-1.051,0.750) (-1.000,0.713) (-0.950,0.677) (-0.900,0.640) (-0.850,0.604) (-0.800,0.568) (-0.750,0.532) (-0.705,0.500) (-0.650,0.461) (-0.600,0.425) (-0.550,0.390) (-0.500,0.354) (-0.450,0.319) (-0.400,0.284) (-0.354,0.250) (-0.300,0.213) (-0.250,0.177) (-0.220,0.155)}
\EQ{(-0.220,-0.155) (-0.275,-0.195) (-0.319,-0.225) (-0.375,-0.266) (-0.424,-0.300) (-0.475,-0.337) (-0.525,-0.372) (-0.575,-0.408) (-0.625,-0.443) (-0.670,-0.475) (-0.725,-0.515) (-0.774,-0.550) (-0.825,-0.587) (-0.875,-0.623) (-0.925,-0.659) (-0.975,-0.696) (-1.025,-0.733) (-1.075,-0.770) (-1.116,-0.800) (-1.175,-0.844) (-1.216,-0.875) (-1.275,-0.920) (-1.325,-0.959) (-1.375,-0.998) (-1.425,-1.037) (-1.472,-1.075) (-1.525,-1.118) (-1.565,-1.150) (-1.625,-1.200) (-1.675,-1.242) (-1.713,-1.275) (-1.771,-1.325) (-1.825,-1.374) (-1.875,-1.419) (-1.925,-1.466) (-1.962,-1.500) (-2.000,-1.537) (-2.038,-1.575) (-2.075,-1.612) (-2.112,-1.650) (-2.150,-1.690) (-2.200,-1.744) (-2.250,-1.799) (-2.294,-1.850) (-2.336,-1.900) (-2.375,-1.948) (-2.417,-2.000) (-2.450,-2.044) (-2.491,-2.100) (-2.525,-2.147) (-2.561,-2.200) (-2.594,-2.250) (-2.625,-2.299) (-2.656,-2.350) (-2.685,-2.400) (-2.713,-2.450) (-2.740,-2.500) (-2.765,-2.550) (-2.790,-2.600) (-2.813,-2.650) (-2.835,-2.700) (-2.856,-2.750) (-2.875,-2.799) (-2.894,-2.850) (-2.912,-2.900) (-2.928,-2.950) (-2.944,-3.000) (-2.958,-3.050) (-2.972,-3.100) (-2.984,-3.150) (-2.995,-3.200) (-3.005,-3.250) (-3.015,-3.300) (-3.023,-3.350) (-3.030,-3.400) (-3.036,-3.450) (-3.041,-3.500) (-3.045,-3.550) (-3.049,-3.600) (-3.051,-3.650) (-3.052,-3.700) (-3.052,-3.750) (-3.052,-3.800) (-3.050,-3.850) (-3.047,-3.900) (-3.044,-3.950) (-3.039,-4.000) (-3.033,-4.050) (-3.027,-4.100) (-3.019,-4.150) (-3.010,-4.200) (-3.001,-4.250) (-2.990,-4.300) (-2.978,-4.350) (-2.965,-4.400) (-2.951,-4.450) (-2.937,-4.500) (-2.921,-4.550) (-2.903,-4.600) (-2.885,-4.650) (-2.866,-4.700) (-2.845,-4.750) (-2.825,-4.796) (-2.800,-4.850) (-2.776,-4.900) (-2.751,-4.950) (-2.725,-4.998) (-2.700,-5.042) (-2.666,-5.100) (-2.635,-5.150) (-2.602,-5.200) (-2.568,-5.250) (-2.533,-5.300) (-2.500,-5.344) (-2.456,-5.400) (-2.415,-5.450) (-2.375,-5.497) (-2.328,-5.550) (-2.281,-5.600) (-2.231,-5.650) (-2.180,-5.700) (-2.125,-5.750) (-2.075,-5.794) (-2.025,-5.836) (-1.977,-5.875) (-1.925,-5.915) (-1.878,-5.950) (-1.825,-5.987) (-1.775,-6.021) (-1.731,-6.050) (-1.675,-6.085) (-1.625,-6.115) (-1.575,-6.143) (-1.525,-6.170) (-1.475,-6.197) (-1.425,-6.222) (-1.375,-6.246) (-1.325,-6.269) (-1.275,-6.291) (-1.225,-6.311) (-1.175,-6.331) (-1.126,-6.350) (-1.075,-6.368) (-1.025,-6.385) (-0.975,-6.401) (-0.925,-6.417) (-0.875,-6.431) (-0.825,-6.445) (-0.775,-6.457) (-0.725,-6.469) (-0.675,-6.480) (-0.625,-6.490) (-0.575,-6.499) (-0.525,-6.508) (-0.475,-6.516) (-0.425,-6.523) (-0.375,-6.529) (-0.325,-6.534) (-0.275,-6.539) (-0.225,-6.543) (-0.175,-6.546) (-0.125,-6.548) (-0.075,-6.549) (-0.025,-6.550) (0.047,-6.550) (0.100,-6.549) (0.150,-6.547) (0.200,-6.544) (0.250,-6.541) (0.300,-6.537) (0.350,-6.532) (0.400,-6.526) (0.450,-6.519) (0.500,-6.512) (0.550,-6.504) (0.600,-6.495) (0.650,-6.485) (0.700,-6.475) (0.750,-6.463) (0.800,-6.451) (0.850,-6.438) (0.900,-6.424) (0.950,-6.409) (1.000,-6.393) (1.050,-6.377) (1.100,-6.359) (1.150,-6.341) (1.200,-6.322) (1.250,-6.301) (1.300,-6.280) (1.350,-6.257) (1.400,-6.234) (1.450,-6.209) (1.500,-6.184) (1.550,-6.157) (1.600,-6.129) (1.650,-6.100) (1.700,-6.069) (1.750,-6.038) (1.800,-6.004) (1.842,-5.975) (1.900,-5.934) (1.945,-5.900) (2.000,-5.857) (2.050,-5.816) (2.097,-5.775) (2.150,-5.728) (2.200,-5.681) (2.250,-5.632) (2.300,-5.580) (2.350,-5.525) (2.394,-5.475) (2.436,-5.425) (2.475,-5.376) (2.514,-5.325) (2.550,-5.276) (2.586,-5.225) (2.619,-5.175) (2.650,-5.126) (2.681,-5.075) (2.710,-5.025) (2.737,-4.975) (2.763,-4.925) (2.788,-4.875) (2.812,-4.825) (2.834,-4.775) (2.856,-4.725) (2.875,-4.676) (2.894,-4.625) (2.912,-4.575) (2.929,-4.525) (2.944,-4.475) (2.959,-4.425) (2.972,-4.375) (2.984,-4.325) (2.995,-4.275) (3.006,-4.225) (3.015,-4.175) (3.023,-4.125) (3.030,-4.075) (3.036,-4.025) (3.041,-3.975) (3.046,-3.925) (3.049,-3.875) (3.051,-3.825) (3.052,-3.775) (3.052,-3.725) (3.052,-3.675) (3.050,-3.625) (3.047,-3.575) (3.044,-3.525) (3.039,-3.475) (3.033,-3.425) (3.026,-3.375) (3.019,-3.325) (3.010,-3.275) (3.000,-3.225) (2.990,-3.175) (2.978,-3.125) (2.965,-3.075) (2.951,-3.025) (2.936,-2.975) (2.920,-2.925) (2.903,-2.875) (2.885,-2.825) (2.866,-2.775) (2.845,-2.725) (2.825,-2.677) (2.801,-2.625) (2.778,-2.575) (2.753,-2.525) (2.726,-2.475) (2.700,-2.427) (2.675,-2.383) (2.641,-2.325) (2.610,-2.275) (2.578,-2.225) (2.544,-2.175) (2.509,-2.125) (2.475,-2.078) (2.436,-2.025) (2.400,-1.979) (2.357,-1.925) (2.315,-1.875) (2.275,-1.828) (2.228,-1.775) (2.183,-1.725) (2.136,-1.675) (2.100,-1.638) (2.063,-1.600) (2.025,-1.562) (1.987,-1.525) (1.950,-1.489) (1.909,-1.450) (1.854,-1.400) (1.800,-1.351) (1.750,-1.307) (1.700,-1.264) (1.655,-1.225) (1.600,-1.179) (1.550,-1.138) (1.504,-1.100) (1.450,-1.057) (1.410,-1.025) (1.350,-0.978) (1.300,-0.940) (1.250,-0.901) (1.200,-0.863) (1.150,-0.826) (1.100,-0.788) (1.050,-0.751) (1.000,-0.714) (0.950,-0.678) (0.900,-0.641) (0.850,-0.605) (0.809,-0.575) (0.750,-0.533) (0.705,-0.500) (0.650,-0.461) (0.600,-0.425) (0.550,-0.390) (0.500,-0.354) (0.450,-0.319) (0.400,-0.284) (0.354,-0.250) (0.300,-0.213) (0.250,-0.177) (0.206,-0.144)}
\EQ{(0.206,0.144) (0.248,0.175) (0.300,0.213) (0.350,0.248) (0.400,0.284) (0.450,0.319) (0.495,0.350) (0.550,0.390) (0.600,0.425) (0.650,0.461) (0.700,0.497) (0.750,0.533) (0.800,0.569) (0.844,0.600) (0.900,0.641) (0.947,0.675) (1.000,0.714) (1.049,0.750) (1.100,0.788) (1.149,0.825) (1.200,0.863) (1.249,0.900) (1.300,0.940) (1.346,0.975) (1.400,1.018) (1.441,1.050) (1.500,1.097) (1.550,1.138) (1.595,1.175) (1.650,1.221) (1.700,1.264) (1.742,1.300) (1.799,1.350) (1.850,1.396) (1.900,1.442) (1.950,1.489) (1.987,1.525) (2.025,1.562) (2.063,1.600) (2.100,1.638) (2.136,1.675) (2.175,1.717) (2.225,1.771) (2.273,1.825) (2.315,1.875) (2.350,1.917) (2.397,1.975) (2.436,2.025) (2.473,2.075) (2.509,2.125) (2.544,2.175) (2.575,2.221) (2.610,2.275) (2.641,2.325) (2.671,2.375) (2.699,2.425) (2.725,2.472) (2.750,2.520) (2.775,2.570) (2.800,2.622) (2.824,2.675) (2.845,2.725) (2.866,2.775) (2.885,2.825) (2.903,2.875) (2.920,2.925) (2.936,2.975) (2.951,3.025) (2.965,3.075) (2.978,3.125) (2.990,3.175) (3.000,3.225) (3.010,3.275) (3.019,3.325) (3.026,3.375) (3.033,3.425) (3.039,3.475) (3.044,3.525) (3.047,3.575) (3.050,3.625) (3.052,3.675) (3.052,3.725) (3.052,3.775) (3.051,3.825) (3.049,3.875) (3.046,3.925) (3.041,3.975) (3.036,4.025) (3.030,4.075) (3.023,4.125) (3.015,4.175) (3.006,4.225) (2.995,4.275) (2.984,4.325) (2.972,4.375) (2.959,4.425) (2.944,4.475) (2.929,4.525) (2.912,4.575) (2.894,4.625) (2.876,4.675) (2.856,4.725) (2.834,4.775) (2.812,4.825) (2.788,4.875) (2.763,4.925) (2.737,4.975) (2.710,5.025) (2.681,5.075) (2.651,5.125) (2.619,5.175) (2.586,5.225) (2.551,5.275) (2.514,5.325) (2.476,5.375) (2.436,5.425) (2.400,5.468) (2.350,5.525) (2.305,5.575) (2.256,5.625) (2.206,5.675) (2.153,5.725) (2.100,5.773) (2.050,5.816) (2.008,5.850) (1.950,5.896) (1.900,5.934) (1.850,5.970) (1.806,6.000) (1.750,6.038) (1.700,6.069) (1.650,6.100) (1.600,6.129) (1.550,6.157) (1.500,6.184) (1.450,6.209) (1.400,6.234) (1.350,6.257) (1.300,6.280) (1.253,6.300) (1.200,6.322) (1.150,6.341) (1.100,6.359) (1.050,6.377) (1.000,6.393) (0.950,6.409) (0.900,6.424) (0.850,6.438) (0.800,6.451) (0.750,6.463) (0.700,6.475) (0.650,6.485) (0.600,6.495) (0.550,6.504) (0.500,6.512) (0.450,6.519) (0.400,6.526) (0.350,6.532) (0.300,6.537) (0.250,6.541) (0.200,6.544) (0.150,6.547) (0.100,6.549) (0.050,6.550) (0.000,6.550) (-0.050,6.550) (-0.100,6.549) (-0.150,6.547) (-0.200,6.544) (-0.250,6.541) (-0.300,6.537) (-0.350,6.532) (-0.400,6.526) (-0.450,6.519) (-0.500,6.512) (-0.550,6.504) (-0.600,6.495) (-0.650,6.485) (-0.700,6.475) (-0.750,6.463) (-0.800,6.451) (-0.850,6.438) (-0.900,6.424) (-0.950,6.409) (-1.000,6.393) (-1.050,6.377) (-1.100,6.359) (-1.150,6.341) (-1.200,6.322) (-1.250,6.301) (-1.300,6.280) (-1.350,6.257) (-1.400,6.234) (-1.450,6.209) (-1.500,6.184) (-1.550,6.157) (-1.600,6.129) (-1.650,6.100) (-1.700,6.069) (-1.750,6.038) (-1.800,6.004) (-1.842,5.975) (-1.900,5.934) (-1.945,5.900) (-2.000,5.857) (-2.050,5.816) (-2.097,5.775) (-2.150,5.728) (-2.200,5.681) (-2.250,5.632) (-2.300,5.580) (-2.350,5.525) (-2.394,5.475) (-2.436,5.425) (-2.475,5.376) (-2.514,5.325) (-2.550,5.276) (-2.586,5.225) (-2.619,5.175) (-2.650,5.126) (-2.681,5.075) (-2.710,5.025) (-2.737,4.975) (-2.763,4.925) (-2.788,4.875) (-2.812,4.825) (-2.834,4.775) (-2.856,4.725) (-2.875,4.676) (-2.894,4.625) (-2.912,4.575) (-2.929,4.525) (-2.944,4.475) (-2.959,4.425) (-2.972,4.375) (-2.984,4.325) (-2.995,4.275) (-3.006,4.225) (-3.015,4.175) (-3.023,4.125) (-3.030,4.075) (-3.036,4.025) (-3.041,3.975) (-3.046,3.925) (-3.049,3.875) (-3.051,3.825) (-3.052,3.775) (-3.052,3.725) (-3.052,3.675) (-3.050,3.625) (-3.047,3.575) (-3.044,3.525) (-3.039,3.475) (-3.033,3.425) (-3.026,3.375) (-3.019,3.325) (-3.010,3.275) (-3.000,3.225) (-2.990,3.175) (-2.978,3.125) (-2.965,3.075) (-2.951,3.025) (-2.936,2.975) (-2.920,2.925) (-2.903,2.875) (-2.885,2.825) (-2.866,2.775) (-2.845,2.725) (-2.825,2.677) (-2.801,2.625) (-2.778,2.575) (-2.753,2.525) (-2.726,2.475) (-2.700,2.427) (-2.675,2.383) (-2.641,2.325) (-2.610,2.275) (-2.578,2.225) (-2.544,2.175) (-2.509,2.125) (-2.475,2.078) (-2.436,2.025) (-2.400,1.979) (-2.357,1.925) (-2.315,1.875) (-2.275,1.828) (-2.228,1.775) (-2.183,1.725) (-2.136,1.675) (-2.100,1.638) (-2.063,1.600) (-2.025,1.562) (-1.987,1.525) (-1.950,1.489) (-1.909,1.450) (-1.854,1.400) (-1.800,1.351) (-1.750,1.307) (-1.700,1.264) (-1.655,1.225) (-1.600,1.179) (-1.550,1.138) (-1.504,1.100) (-1.450,1.057) (-1.410,1.025) (-1.350,0.978) (-1.300,0.940) (-1.250,0.901) (-1.200,0.863) (-1.150,0.826) (-1.100,0.788) (-1.050,0.751) (-1.000,0.714) (-0.950,0.678) (-0.900,0.641) (-0.850,0.605) (-0.809,0.575) (-0.750,0.533) (-0.705,0.500) (-0.650,0.461) (-0.600,0.425) (-0.550,0.390) (-0.500,0.354) (-0.450,0.319) (-0.400,0.284) (-0.354,0.250) (-0.300,0.213) (-0.250,0.177) (-0.220,0.155)}
\EQ{(-0.220,-0.155) (-0.275,-0.195) (-0.319,-0.225) (-0.375,-0.266) (-0.424,-0.300) (-0.475,-0.337) (-0.525,-0.372) (-0.575,-0.408) (-0.625,-0.444) (-0.670,-0.475) (-0.725,-0.515) (-0.774,-0.550) (-0.825,-0.587) (-0.875,-0.624) (-0.925,-0.660) (-0.975,-0.697) (-1.025,-0.734) (-1.075,-0.771) (-1.125,-0.809) (-1.175,-0.847) (-1.225,-0.885) (-1.275,-0.924) (-1.325,-0.963) (-1.372,-1.000) (-1.425,-1.043) (-1.465,-1.075) (-1.525,-1.124) (-1.575,-1.166) (-1.615,-1.200) (-1.672,-1.250) (-1.725,-1.297) (-1.775,-1.342) (-1.825,-1.389) (-1.863,-1.425) (-1.900,-1.461) (-1.940,-1.500) (-1.990,-1.550) (-2.025,-1.587) (-2.061,-1.625) (-2.100,-1.668) (-2.150,-1.724) (-2.194,-1.775) (-2.236,-1.825) (-2.275,-1.874) (-2.315,-1.925) (-2.350,-1.972) (-2.388,-2.025) (-2.423,-2.075) (-2.456,-2.125) (-2.488,-2.175) (-2.519,-2.225) (-2.548,-2.275) (-2.575,-2.324) (-2.600,-2.370) (-2.625,-2.419) (-2.650,-2.470) (-2.675,-2.524) (-2.697,-2.575) (-2.718,-2.625) (-2.737,-2.675) (-2.756,-2.725) (-2.773,-2.775) (-2.789,-2.825) (-2.804,-2.875) (-2.818,-2.925) (-2.830,-2.975) (-2.842,-3.025) (-2.852,-3.075) (-2.862,-3.125) (-2.870,-3.175) (-2.877,-3.225) (-2.884,-3.275) (-2.889,-3.325) (-2.893,-3.375) (-2.896,-3.425) (-2.898,-3.475) (-2.899,-3.525) (-2.899,-3.575) (-2.898,-3.625) (-2.896,-3.675) (-2.893,-3.725) (-2.889,-3.775) (-2.884,-3.825) (-2.878,-3.875) (-2.871,-3.925) (-2.862,-3.975) (-2.853,-4.025) (-2.842,-4.075) (-2.831,-4.125) (-2.818,-4.175) (-2.804,-4.225) (-2.789,-4.275) (-2.773,-4.325) (-2.756,-4.375) (-2.738,-4.425) (-2.718,-4.475) (-2.697,-4.525) (-2.675,-4.575) (-2.652,-4.625) (-2.627,-4.675) (-2.601,-4.725) (-2.575,-4.772) (-2.550,-4.815) (-2.514,-4.875) (-2.482,-4.925) (-2.450,-4.972) (-2.413,-5.025) (-2.376,-5.075) (-2.337,-5.125) (-2.300,-5.170) (-2.253,-5.225) (-2.208,-5.275) (-2.161,-5.325) (-2.125,-5.362) (-2.086,-5.400) (-2.032,-5.450) (-1.976,-5.500) (-1.925,-5.543) (-1.885,-5.575) (-1.825,-5.621) (-1.775,-5.658) (-1.725,-5.693) (-1.678,-5.725) (-1.625,-5.759) (-1.575,-5.790) (-1.525,-5.819) (-1.475,-5.847) (-1.425,-5.874) (-1.375,-5.900) (-1.325,-5.924) (-1.275,-5.947) (-1.225,-5.970) (-1.175,-5.991) (-1.125,-6.011) (-1.075,-6.030) (-1.025,-6.048) (-0.975,-6.065) (-0.925,-6.081) (-0.875,-6.096) (-0.825,-6.111) (-0.775,-6.124) (-0.725,-6.137) (-0.675,-6.148) (-0.625,-6.159) (-0.575,-6.169) (-0.525,-6.178) (-0.475,-6.186) (-0.425,-6.193) (-0.375,-6.200) (-0.325,-6.205) (-0.275,-6.210) (-0.225,-6.214) (-0.175,-6.217) (-0.125,-6.220) (-0.075,-6.221) (-0.025,-6.222) (0.050,-6.222) (0.100,-6.221) (0.150,-6.219) (0.200,-6.216) (0.250,-6.212) (0.300,-6.208) (0.350,-6.203) (0.400,-6.196) (0.450,-6.190) (0.500,-6.182) (0.550,-6.173) (0.600,-6.164) (0.650,-6.154) (0.700,-6.142) (0.750,-6.130) (0.800,-6.117) (0.850,-6.104) (0.900,-6.089) (0.950,-6.073) (1.000,-6.057) (1.050,-6.039) (1.100,-6.020) (1.150,-6.001) (1.200,-5.980) (1.250,-5.959) (1.300,-5.936) (1.350,-5.912) (1.400,-5.887) (1.450,-5.861) (1.500,-5.833) (1.550,-5.805) (1.599,-5.775) (1.650,-5.743) (1.700,-5.710) (1.750,-5.676) (1.800,-5.640) (1.850,-5.602) (1.900,-5.563) (1.946,-5.525) (2.000,-5.479) (2.050,-5.434) (2.100,-5.386) (2.137,-5.350) (2.175,-5.311) (2.225,-5.257) (2.275,-5.200) (2.317,-5.150) (2.350,-5.108) (2.394,-5.050) (2.425,-5.008) (2.465,-4.950) (2.498,-4.900) (2.525,-4.857) (2.559,-4.800) (2.587,-4.750) (2.614,-4.700) (2.639,-4.650) (2.663,-4.600) (2.686,-4.550) (2.708,-4.500) (2.728,-4.450) (2.747,-4.400) (2.765,-4.350) (2.781,-4.300) (2.797,-4.250) (2.811,-4.200) (2.825,-4.150) (2.837,-4.100) (2.848,-4.050) (2.858,-4.000) (2.867,-3.950) (2.874,-3.900) (2.881,-3.850) (2.887,-3.800) (2.891,-3.750) (2.895,-3.700) (2.898,-3.650) (2.899,-3.600) (2.900,-3.550) (2.899,-3.500) (2.897,-3.450) (2.895,-3.400) (2.891,-3.350) (2.887,-3.300) (2.881,-3.250) (2.874,-3.200) (2.866,-3.150) (2.857,-3.100) (2.847,-3.050) (2.836,-3.000) (2.824,-2.950) (2.811,-2.900) (2.796,-2.850) (2.781,-2.800) (2.764,-2.750) (2.747,-2.700) (2.728,-2.650) (2.708,-2.600) (2.687,-2.550) (2.664,-2.500) (2.640,-2.450) (2.616,-2.400) (2.589,-2.350) (2.562,-2.300) (2.533,-2.250) (2.503,-2.200) (2.475,-2.155) (2.440,-2.100) (2.406,-2.050) (2.375,-2.007) (2.334,-1.950) (2.300,-1.906) (2.256,-1.850) (2.215,-1.800) (2.175,-1.752) (2.129,-1.700) (2.084,-1.650) (2.050,-1.613) (2.014,-1.575) (1.975,-1.535) (1.925,-1.485) (1.889,-1.450) (1.850,-1.413) (1.811,-1.375) (1.756,-1.325) (1.700,-1.275) (1.650,-1.231) (1.600,-1.188) (1.556,-1.150) (1.500,-1.104) (1.450,-1.063) (1.404,-1.025) (1.350,-0.983) (1.309,-0.950) (1.250,-0.904) (1.200,-0.866) (1.150,-0.828) (1.100,-0.790) (1.050,-0.753) (1.000,-0.716) (0.950,-0.679) (0.900,-0.642) (0.850,-0.606) (0.809,-0.575) (0.750,-0.533) (0.705,-0.500) (0.650,-0.461) (0.600,-0.425) (0.550,-0.390) (0.500,-0.354) (0.450,-0.319) (0.400,-0.284) (0.354,-0.250) (0.300,-0.213) (0.250,-0.177) (0.206,-0.144)}
\EQ{(0.206,0.144) (0.248,0.175) (0.300,0.213) (0.350,0.248) (0.400,0.284) (0.450,0.319) (0.495,0.350) (0.550,0.390) (0.600,0.425) (0.650,0.461) (0.700,0.497) (0.750,0.533) (0.800,0.569) (0.843,0.600) (0.900,0.642) (0.946,0.675) (1.000,0.716) (1.047,0.750) (1.100,0.790) (1.147,0.825) (1.200,0.866) (1.245,0.900) (1.300,0.943) (1.341,0.975) (1.400,1.022) (1.450,1.063) (1.496,1.100) (1.550,1.145) (1.600,1.188) (1.644,1.225) (1.700,1.275) (1.750,1.320) (1.800,1.366) (1.837,1.400) (1.875,1.436) (1.915,1.475) (1.965,1.525) (2.000,1.561) (2.038,1.600) (2.075,1.640) (2.125,1.695) (2.173,1.750) (2.215,1.800) (2.250,1.843) (2.296,1.900) (2.334,1.950) (2.370,2.000) (2.400,2.042) (2.440,2.100) (2.472,2.150) (2.500,2.195) (2.533,2.250) (2.562,2.300) (2.589,2.350) (2.616,2.400) (2.640,2.450) (2.664,2.500) (2.687,2.550) (2.708,2.600) (2.728,2.650) (2.747,2.700) (2.764,2.750) (2.781,2.800) (2.796,2.850) (2.811,2.900) (2.824,2.950) (2.836,3.000) (2.847,3.050) (2.857,3.100) (2.866,3.150) (2.874,3.200) (2.881,3.250) (2.887,3.300) (2.891,3.350) (2.895,3.400) (2.897,3.450) (2.899,3.500) (2.900,3.550) (2.899,3.600) (2.898,3.650) (2.895,3.700) (2.891,3.750) (2.887,3.800) (2.881,3.850) (2.874,3.900) (2.867,3.950) (2.858,4.000) (2.848,4.050) (2.837,4.100) (2.825,4.150) (2.811,4.200) (2.797,4.250) (2.781,4.300) (2.765,4.350) (2.747,4.400) (2.728,4.450) (2.708,4.500) (2.686,4.550) (2.663,4.600) (2.639,4.650) (2.614,4.700) (2.587,4.750) (2.559,4.800) (2.529,4.850) (2.500,4.897) (2.465,4.950) (2.431,5.000) (2.400,5.043) (2.357,5.100) (2.325,5.140) (2.275,5.200) (2.231,5.250) (2.185,5.300) (2.150,5.337) (2.112,5.375) (2.075,5.410) (2.032,5.450) (1.976,5.500) (1.925,5.543) (1.885,5.575) (1.825,5.621) (1.775,5.658) (1.725,5.693) (1.678,5.725) (1.625,5.759) (1.575,5.790) (1.525,5.819) (1.475,5.847) (1.425,5.874) (1.375,5.900) (1.325,5.924) (1.275,5.947) (1.225,5.970) (1.175,5.991) (1.125,6.011) (1.075,6.030) (1.025,6.048) (0.975,6.065) (0.925,6.081) (0.875,6.096) (0.825,6.111) (0.775,6.124) (0.725,6.137) (0.675,6.148) (0.625,6.159) (0.575,6.169) (0.525,6.178) (0.475,6.186) (0.425,6.193) (0.375,6.200) (0.325,6.205) (0.275,6.210) (0.225,6.214) (0.175,6.217) (0.125,6.220) (0.075,6.221) (0.025,6.222) (-0.050,6.222) (-0.100,6.221) (-0.150,6.219) (-0.200,6.216) (-0.250,6.212) (-0.300,6.208) (-0.350,6.203) (-0.400,6.196) (-0.450,6.190) (-0.500,6.182) (-0.550,6.173) (-0.600,6.164) (-0.650,6.154) (-0.700,6.142) (-0.750,6.130) (-0.800,6.117) (-0.850,6.104) (-0.900,6.089) (-0.950,6.073) (-1.000,6.057) (-1.050,6.039) (-1.100,6.020) (-1.150,6.001) (-1.200,5.980) (-1.250,5.959) (-1.300,5.936) (-1.350,5.912) (-1.400,5.887) (-1.450,5.861) (-1.500,5.833) (-1.550,5.805) (-1.599,5.775) (-1.650,5.743) (-1.700,5.710) (-1.750,5.676) (-1.800,5.640) (-1.850,5.602) (-1.900,5.563) (-1.946,5.525) (-2.000,5.479) (-2.050,5.434) (-2.100,5.386) (-2.137,5.350) (-2.175,5.311) (-2.225,5.257) (-2.275,5.200) (-2.317,5.150) (-2.350,5.108) (-2.394,5.050) (-2.425,5.008) (-2.465,4.950) (-2.498,4.900) (-2.525,4.857) (-2.559,4.800) (-2.587,4.750) (-2.614,4.700) (-2.639,4.650) (-2.663,4.600) (-2.686,4.550) (-2.708,4.500) (-2.728,4.450) (-2.747,4.400) (-2.765,4.350) (-2.781,4.300) (-2.797,4.250) (-2.811,4.200) (-2.825,4.150) (-2.837,4.100) (-2.848,4.050) (-2.858,4.000) (-2.867,3.950) (-2.874,3.900) (-2.881,3.850) (-2.887,3.800) (-2.891,3.750) (-2.895,3.700) (-2.898,3.650) (-2.899,3.600) (-2.900,3.550) (-2.899,3.500) (-2.897,3.450) (-2.895,3.400) (-2.891,3.350) (-2.887,3.300) (-2.881,3.250) (-2.874,3.200) (-2.866,3.150) (-2.857,3.100) (-2.847,3.050) (-2.836,3.000) (-2.824,2.950) (-2.811,2.900) (-2.796,2.850) (-2.781,2.800) (-2.764,2.750) (-2.747,2.700) (-2.728,2.650) (-2.708,2.600) (-2.687,2.550) (-2.664,2.500) (-2.640,2.450) (-2.616,2.400) (-2.589,2.350) (-2.562,2.300) (-2.533,2.250) (-2.503,2.200) (-2.475,2.155) (-2.440,2.100) (-2.406,2.050) (-2.375,2.007) (-2.334,1.950) (-2.300,1.906) (-2.256,1.850) (-2.215,1.800) (-2.175,1.752) (-2.129,1.700) (-2.084,1.650) (-2.050,1.613) (-2.014,1.575) (-1.975,1.535) (-1.925,1.485) (-1.889,1.450) (-1.850,1.413) (-1.811,1.375) (-1.756,1.325) (-1.700,1.275) (-1.650,1.231) (-1.600,1.188) (-1.556,1.150) (-1.500,1.104) (-1.450,1.063) (-1.404,1.025) (-1.350,0.983) (-1.309,0.950) (-1.250,0.904) (-1.200,0.866) (-1.150,0.828) (-1.100,0.790) (-1.050,0.753) (-1.000,0.716) (-0.950,0.679) (-0.900,0.642) (-0.850,0.606) (-0.809,0.575) (-0.750,0.533) (-0.705,0.500) (-0.650,0.461) (-0.600,0.425) (-0.550,0.390) (-0.500,0.354) (-0.450,0.319) (-0.400,0.284) (-0.354,0.250) (-0.300,0.213) (-0.250,0.177) (-0.220,0.155)}
\EQ{(-0.220,-0.155) (-0.275,-0.195) (-0.319,-0.225) (-0.375,-0.266) (-0.424,-0.300) (-0.475,-0.337) (-0.525,-0.372) (-0.575,-0.408) (-0.625,-0.444) (-0.670,-0.475) (-0.725,-0.516) (-0.773,-0.550) (-0.825,-0.588) (-0.875,-0.624) (-0.925,-0.661) (-0.975,-0.698) (-1.025,-0.735) (-1.075,-0.773) (-1.125,-0.811) (-1.175,-0.849) (-1.225,-0.888) (-1.273,-0.925) (-1.325,-0.967) (-1.367,-1.000) (-1.425,-1.048) (-1.475,-1.089) (-1.518,-1.125) (-1.575,-1.174) (-1.625,-1.218) (-1.675,-1.263) (-1.715,-1.300) (-1.769,-1.350) (-1.821,-1.400) (-1.871,-1.450) (-1.920,-1.500) (-1.967,-1.550) (-2.013,-1.600) (-2.050,-1.641) (-2.100,-1.699) (-2.142,-1.750) (-2.175,-1.791) (-2.221,-1.850) (-2.258,-1.900) (-2.293,-1.950) (-2.325,-1.996) (-2.361,-2.050) (-2.392,-2.100) (-2.422,-2.150) (-2.450,-2.199) (-2.475,-2.245) (-2.500,-2.292) (-2.525,-2.343) (-2.550,-2.396) (-2.574,-2.450) (-2.595,-2.500) (-2.615,-2.550) (-2.633,-2.600) (-2.650,-2.648) (-2.667,-2.700) (-2.682,-2.750) (-2.696,-2.800) (-2.708,-2.850) (-2.720,-2.900) (-2.730,-2.950) (-2.739,-3.000) (-2.747,-3.050) (-2.754,-3.100) (-2.760,-3.150) (-2.765,-3.200) (-2.769,-3.250) (-2.771,-3.300) (-2.773,-3.350) (-2.773,-3.400) (-2.773,-3.450) (-2.771,-3.500) (-2.768,-3.550) (-2.764,-3.600) (-2.759,-3.650) (-2.753,-3.700) (-2.746,-3.750) (-2.738,-3.800) (-2.729,-3.850) (-2.718,-3.900) (-2.706,-3.950) (-2.694,-4.000) (-2.680,-4.050) (-2.664,-4.100) (-2.648,-4.150) (-2.630,-4.200) (-2.612,-4.250) (-2.592,-4.300) (-2.570,-4.350) (-2.547,-4.400) (-2.525,-4.446) (-2.500,-4.496) (-2.475,-4.542) (-2.450,-4.587) (-2.412,-4.650) (-2.381,-4.700) (-2.350,-4.746) (-2.312,-4.800) (-2.275,-4.850) (-2.237,-4.900) (-2.200,-4.945) (-2.153,-5.000) (-2.108,-5.050) (-2.061,-5.100) (-2.025,-5.136) (-1.984,-5.175) (-1.930,-5.225) (-1.875,-5.273) (-1.825,-5.315) (-1.781,-5.350) (-1.725,-5.392) (-1.679,-5.425) (-1.625,-5.462) (-1.575,-5.495) (-1.527,-5.525) (-1.475,-5.556) (-1.425,-5.585) (-1.375,-5.612) (-1.325,-5.638) (-1.275,-5.662) (-1.225,-5.686) (-1.175,-5.708) (-1.125,-5.729) (-1.075,-5.749) (-1.025,-5.769) (-0.975,-5.787) (-0.925,-5.804) (-0.875,-5.819) (-0.825,-5.834) (-0.775,-5.848) (-0.725,-5.862) (-0.675,-5.874) (-0.625,-5.885) (-0.575,-5.895) (-0.525,-5.905) (-0.475,-5.913) (-0.425,-5.921) (-0.375,-5.928) (-0.325,-5.934) (-0.275,-5.939) (-0.225,-5.943) (-0.175,-5.946) (-0.125,-5.949) (-0.075,-5.950) (-0.025,-5.951) (0.050,-5.951) (0.100,-5.950) (0.150,-5.948) (0.200,-5.945) (0.250,-5.941) (0.300,-5.936) (0.350,-5.931) (0.400,-5.924) (0.450,-5.917) (0.500,-5.909) (0.550,-5.900) (0.600,-5.890) (0.650,-5.879) (0.700,-5.868) (0.750,-5.855) (0.800,-5.842) (0.850,-5.827) (0.900,-5.812) (0.950,-5.795) (1.000,-5.778) (1.050,-5.759) (1.100,-5.740) (1.150,-5.719) (1.200,-5.697) (1.249,-5.675) (1.300,-5.650) (1.350,-5.625) (1.397,-5.600) (1.442,-5.575) (1.486,-5.550) (1.550,-5.511) (1.600,-5.479) (1.643,-5.450) (1.700,-5.410) (1.748,-5.375) (1.800,-5.335) (1.843,-5.300) (1.900,-5.252) (1.950,-5.207) (2.000,-5.160) (2.036,-5.125) (2.075,-5.085) (2.125,-5.031) (2.175,-4.975) (2.217,-4.925) (2.250,-4.883) (2.294,-4.825) (2.325,-4.782) (2.364,-4.725) (2.397,-4.675) (2.425,-4.629) (2.457,-4.575) (2.484,-4.525) (2.511,-4.475) (2.535,-4.425) (2.559,-4.375) (2.581,-4.325) (2.602,-4.275) (2.621,-4.225) (2.639,-4.175) (2.656,-4.125) (2.672,-4.075) (2.687,-4.025) (2.700,-3.975) (2.712,-3.925) (2.723,-3.875) (2.733,-3.825) (2.742,-3.775) (2.750,-3.725) (2.757,-3.675) (2.762,-3.625) (2.766,-3.575) (2.770,-3.525) (2.772,-3.475) (2.773,-3.425) (2.773,-3.375) (2.772,-3.325) (2.770,-3.275) (2.767,-3.225) (2.763,-3.175) (2.757,-3.125) (2.751,-3.075) (2.743,-3.025) (2.735,-2.975) (2.725,-2.926) (2.714,-2.875) (2.702,-2.825) (2.689,-2.775) (2.674,-2.725) (2.659,-2.675) (2.642,-2.625) (2.625,-2.577) (2.605,-2.525) (2.585,-2.475) (2.563,-2.425) (2.541,-2.375) (2.517,-2.325) (2.491,-2.275) (2.464,-2.225) (2.436,-2.175) (2.407,-2.125) (2.376,-2.075) (2.344,-2.025) (2.311,-1.975) (2.276,-1.925) (2.240,-1.875) (2.202,-1.825) (2.162,-1.775) (2.125,-1.729) (2.079,-1.675) (2.036,-1.625) (2.000,-1.586) (1.950,-1.532) (1.900,-1.480) (1.850,-1.429) (1.800,-1.380) (1.750,-1.333) (1.700,-1.286) (1.661,-1.250) (1.604,-1.200) (1.550,-1.153) (1.500,-1.110) (1.458,-1.075) (1.400,-1.027) (1.350,-0.987) (1.304,-0.950) (1.250,-0.907) (1.209,-0.875) (1.150,-0.830) (1.100,-0.792) (1.050,-0.754) (1.000,-0.717) (0.950,-0.680) (0.900,-0.643) (0.850,-0.606) (0.808,-0.575) (0.750,-0.534) (0.704,-0.500) (0.650,-0.462) (0.600,-0.425) (0.550,-0.390) (0.500,-0.354) (0.450,-0.319) (0.400,-0.284) (0.354,-0.250) (0.300,-0.213) (0.250,-0.177) (0.206,-0.144)}
\EQ{(0.206,0.144) (0.248,0.175) (0.300,0.213) (0.350,0.248) (0.400,0.284) (0.450,0.319) (0.495,0.350) (0.550,0.390) (0.599,0.425) (0.650,0.462) (0.700,0.497) (0.750,0.534) (0.800,0.570) (0.842,0.600) (0.900,0.643) (0.944,0.675) (1.000,0.717) (1.045,0.750) (1.100,0.792) (1.144,0.825) (1.200,0.868) (1.241,0.900) (1.300,0.947) (1.350,0.987) (1.398,1.025) (1.450,1.068) (1.500,1.110) (1.547,1.150) (1.600,1.196) (1.650,1.241) (1.688,1.275) (1.725,1.309) (1.769,1.350) (1.821,1.400) (1.871,1.450) (1.920,1.500) (1.967,1.550) (2.013,1.600) (2.050,1.641) (2.100,1.699) (2.142,1.750) (2.175,1.791) (2.221,1.850) (2.258,1.900) (2.293,1.950) (2.325,1.996) (2.361,2.050) (2.392,2.100) (2.422,2.150) (2.450,2.199) (2.475,2.245) (2.500,2.292) (2.525,2.343) (2.550,2.396) (2.574,2.450) (2.595,2.500) (2.615,2.550) (2.633,2.600) (2.650,2.648) (2.667,2.700) (2.682,2.750) (2.696,2.800) (2.708,2.850) (2.720,2.900) (2.730,2.950) (2.739,3.000) (2.747,3.050) (2.754,3.100) (2.760,3.150) (2.765,3.200) (2.769,3.250) (2.771,3.300) (2.773,3.350) (2.773,3.400) (2.773,3.450) (2.771,3.500) (2.768,3.550) (2.764,3.600) (2.759,3.650) (2.753,3.700) (2.746,3.750) (2.738,3.800) (2.729,3.850) (2.718,3.900) (2.706,3.950) (2.694,4.000) (2.680,4.050) (2.664,4.100) (2.648,4.150) (2.630,4.200) (2.612,4.250) (2.592,4.300) (2.570,4.350) (2.547,4.400) (2.525,4.446) (2.500,4.496) (2.475,4.542) (2.450,4.587) (2.412,4.650) (2.381,4.700) (2.350,4.746) (2.312,4.800) (2.275,4.850) (2.237,4.900) (2.200,4.945) (2.153,5.000) (2.108,5.050) (2.061,5.100) (2.025,5.136) (1.984,5.175) (1.930,5.225) (1.875,5.273) (1.825,5.315) (1.781,5.350) (1.725,5.392) (1.679,5.425) (1.625,5.462) (1.575,5.495) (1.527,5.525) (1.475,5.556) (1.425,5.585) (1.375,5.612) (1.325,5.638) (1.275,5.662) (1.225,5.686) (1.175,5.708) (1.125,5.729) (1.075,5.749) (1.025,5.769) (0.975,5.787) (0.925,5.804) (0.875,5.819) (0.825,5.834) (0.775,5.848) (0.725,5.862) (0.675,5.874) (0.625,5.885) (0.575,5.895) (0.525,5.905) (0.475,5.913) (0.425,5.921) (0.375,5.928) (0.325,5.934) (0.275,5.939) (0.225,5.943) (0.175,5.946) (0.125,5.949) (0.075,5.950) (0.025,5.951) (-0.050,5.951) (-0.100,5.950) (-0.150,5.948) (-0.200,5.945) (-0.250,5.941) (-0.300,5.936) (-0.350,5.931) (-0.400,5.924) (-0.450,5.917) (-0.500,5.909) (-0.550,5.900) (-0.600,5.890) (-0.650,5.879) (-0.700,5.868) (-0.750,5.855) (-0.800,5.842) (-0.850,5.827) (-0.900,5.812) (-0.950,5.795) (-1.000,5.778) (-1.050,5.759) (-1.100,5.740) (-1.150,5.719) (-1.200,5.697) (-1.249,5.675) (-1.300,5.650) (-1.350,5.625) (-1.397,5.600) (-1.442,5.575) (-1.486,5.550) (-1.550,5.511) (-1.600,5.479) (-1.643,5.450) (-1.700,5.410) (-1.748,5.375) (-1.800,5.335) (-1.843,5.300) (-1.900,5.252) (-1.950,5.207) (-2.000,5.160) (-2.036,5.125) (-2.075,5.085) (-2.125,5.031) (-2.175,4.975) (-2.217,4.925) (-2.250,4.883) (-2.294,4.825) (-2.325,4.782) (-2.364,4.725) (-2.397,4.675) (-2.425,4.629) (-2.457,4.575) (-2.484,4.525) (-2.511,4.475) (-2.535,4.425) (-2.559,4.375) (-2.581,4.325) (-2.602,4.275) (-2.621,4.225) (-2.639,4.175) (-2.656,4.125) (-2.672,4.075) (-2.687,4.025) (-2.700,3.975) (-2.712,3.925) (-2.723,3.875) (-2.733,3.825) (-2.742,3.775) (-2.750,3.725) (-2.757,3.675) (-2.762,3.625) (-2.766,3.575) (-2.770,3.525) (-2.772,3.475) (-2.773,3.425) (-2.773,3.375) (-2.772,3.325) (-2.770,3.275) (-2.767,3.225) (-2.763,3.175) (-2.757,3.125) (-2.751,3.075) (-2.743,3.025) (-2.735,2.975) (-2.725,2.926) (-2.714,2.875) (-2.702,2.825) (-2.689,2.775) (-2.674,2.725) (-2.659,2.675) (-2.642,2.625) (-2.625,2.577) (-2.605,2.525) (-2.585,2.475) (-2.563,2.425) (-2.541,2.375) (-2.517,2.325) (-2.491,2.275) (-2.464,2.225) (-2.436,2.175) (-2.407,2.125) (-2.376,2.075) (-2.344,2.025) (-2.311,1.975) (-2.276,1.925) (-2.240,1.875) (-2.202,1.825) (-2.162,1.775) (-2.125,1.729) (-2.079,1.675) (-2.036,1.625) (-2.000,1.586) (-1.950,1.532) (-1.900,1.480) (-1.850,1.429) (-1.800,1.380) (-1.750,1.333) (-1.700,1.286) (-1.661,1.250) (-1.604,1.200) (-1.550,1.153) (-1.500,1.110) (-1.458,1.075) (-1.400,1.027) (-1.350,0.987) (-1.304,0.950) (-1.250,0.907) (-1.209,0.875) (-1.150,0.830) (-1.100,0.792) (-1.050,0.754) (-1.000,0.717) (-0.950,0.680) (-0.900,0.643) (-0.850,0.606) (-0.808,0.575) (-0.750,0.534) (-0.704,0.500) (-0.650,0.462) (-0.600,0.425) (-0.550,0.390) (-0.500,0.354) (-0.450,0.319) (-0.400,0.284) (-0.354,0.250) (-0.300,0.213) (-0.250,0.177) (-0.220,0.155)}
\EQ{(-0.220,-0.155) (-0.275,-0.195) (-0.319,-0.225) (-0.375,-0.266) (-0.424,-0.300) (-0.475,-0.337) (-0.525,-0.372) (-0.575,-0.408) (-0.625,-0.444) (-0.669,-0.475) (-0.725,-0.516) (-0.773,-0.550) (-0.825,-0.589) (-0.875,-0.625) (-0.925,-0.662) (-0.975,-0.699) (-1.025,-0.737) (-1.075,-0.774) (-1.125,-0.813) (-1.173,-0.850) (-1.225,-0.891) (-1.269,-0.925) (-1.325,-0.970) (-1.375,-1.011) (-1.422,-1.050) (-1.475,-1.095) (-1.525,-1.139) (-1.567,-1.175) (-1.622,-1.225) (-1.675,-1.274) (-1.725,-1.322) (-1.775,-1.370) (-1.825,-1.421) (-1.875,-1.473) (-1.924,-1.525) (-1.969,-1.575) (-2.012,-1.625) (-2.050,-1.671) (-2.094,-1.725) (-2.125,-1.766) (-2.170,-1.825) (-2.200,-1.868) (-2.239,-1.925) (-2.272,-1.975) (-2.300,-2.020) (-2.333,-2.075) (-2.362,-2.125) (-2.389,-2.175) (-2.414,-2.225) (-2.439,-2.275) (-2.462,-2.325) (-2.483,-2.375) (-2.504,-2.425) (-2.523,-2.475) (-2.541,-2.525) (-2.558,-2.575) (-2.573,-2.625) (-2.587,-2.675) (-2.600,-2.725) (-2.612,-2.775) (-2.623,-2.825) (-2.632,-2.875) (-2.640,-2.925) (-2.647,-2.975) (-2.654,-3.025) (-2.658,-3.075) (-2.662,-3.125) (-2.665,-3.175) (-2.666,-3.225) (-2.667,-3.275) (-2.666,-3.325) (-2.664,-3.375) (-2.661,-3.425) (-2.657,-3.475) (-2.651,-3.525) (-2.645,-3.575) (-2.637,-3.625) (-2.629,-3.675) (-2.619,-3.725) (-2.608,-3.775) (-2.595,-3.825) (-2.582,-3.875) (-2.567,-3.925) (-2.551,-3.975) (-2.534,-4.025) (-2.515,-4.075) (-2.496,-4.125) (-2.475,-4.174) (-2.452,-4.225) (-2.428,-4.275) (-2.402,-4.325) (-2.375,-4.375) (-2.350,-4.420) (-2.317,-4.475) (-2.285,-4.525) (-2.251,-4.575) (-2.216,-4.625) (-2.179,-4.675) (-2.140,-4.725) (-2.100,-4.773) (-2.055,-4.825) (-2.009,-4.875) (-1.960,-4.925) (-1.909,-4.975) (-1.855,-5.025) (-1.800,-5.073) (-1.750,-5.114) (-1.704,-5.150) (-1.650,-5.191) (-1.602,-5.225) (-1.550,-5.260) (-1.500,-5.293) (-1.450,-5.323) (-1.405,-5.350) (-1.361,-5.375) (-1.314,-5.400) (-1.265,-5.425) (-1.214,-5.450) (-1.160,-5.475) (-1.103,-5.500) (-1.050,-5.522) (-1.000,-5.541) (-0.950,-5.559) (-0.900,-5.577) (-0.850,-5.593) (-0.800,-5.608) (-0.750,-5.622) (-0.700,-5.635) (-0.650,-5.647) (-0.600,-5.659) (-0.550,-5.669) (-0.500,-5.678) (-0.450,-5.687) (-0.400,-5.694) (-0.350,-5.701) (-0.300,-5.707) (-0.250,-5.711) (-0.200,-5.715) (-0.150,-5.718) (-0.100,-5.721) (-0.050,-5.722) (0.000,-5.722) (0.050,-5.722) (0.100,-5.721) (0.150,-5.718) (0.200,-5.715) (0.250,-5.711) (0.300,-5.707) (0.350,-5.701) (0.400,-5.694) (0.450,-5.687) (0.500,-5.678) (0.550,-5.669) (0.600,-5.659) (0.650,-5.647) (0.700,-5.635) (0.750,-5.622) (0.800,-5.608) (0.850,-5.593) (0.900,-5.577) (0.950,-5.559) (1.000,-5.541) (1.050,-5.522) (1.100,-5.501) (1.150,-5.480) (1.200,-5.457) (1.250,-5.433) (1.300,-5.407) (1.350,-5.381) (1.400,-5.353) (1.447,-5.325) (1.500,-5.293) (1.550,-5.260) (1.600,-5.226) (1.650,-5.191) (1.700,-5.153) (1.750,-5.114) (1.797,-5.075) (1.850,-5.029) (1.900,-4.983) (1.950,-4.935) (2.000,-4.884) (2.050,-4.830) (2.098,-4.775) (2.140,-4.725) (2.175,-4.680) (2.216,-4.625) (2.250,-4.577) (2.285,-4.525) (2.317,-4.475) (2.347,-4.425) (2.375,-4.376) (2.400,-4.330) (2.425,-4.281) (2.450,-4.229) (2.474,-4.175) (2.496,-4.125) (2.515,-4.075) (2.534,-4.025) (2.551,-3.975) (2.567,-3.925) (2.582,-3.875) (2.595,-3.825) (2.608,-3.775) (2.619,-3.725) (2.629,-3.675) (2.637,-3.625) (2.645,-3.575) (2.651,-3.525) (2.657,-3.475) (2.661,-3.425) (2.664,-3.375) (2.666,-3.325) (2.667,-3.275) (2.666,-3.225) (2.665,-3.175) (2.662,-3.125) (2.658,-3.075) (2.654,-3.025) (2.647,-2.975) (2.640,-2.925) (2.632,-2.875) (2.623,-2.825) (2.612,-2.775) (2.600,-2.725) (2.587,-2.675) (2.573,-2.625) (2.558,-2.575) (2.541,-2.525) (2.523,-2.475) (2.504,-2.425) (2.483,-2.375) (2.462,-2.325) (2.439,-2.275) (2.414,-2.225) (2.389,-2.175) (2.362,-2.125) (2.333,-2.075) (2.303,-2.025) (2.275,-1.980) (2.239,-1.925) (2.205,-1.875) (2.175,-1.833) (2.132,-1.775) (2.100,-1.733) (2.053,-1.675) (2.012,-1.625) (1.975,-1.583) (1.925,-1.527) (1.877,-1.475) (1.829,-1.425) (1.780,-1.375) (1.729,-1.325) (1.676,-1.275) (1.625,-1.228) (1.575,-1.183) (1.525,-1.139) (1.481,-1.100) (1.425,-1.053) (1.375,-1.011) (1.331,-0.975) (1.275,-0.930) (1.225,-0.891) (1.175,-0.851) (1.125,-0.813) (1.076,-0.775) (1.025,-0.737) (0.976,-0.700) (0.925,-0.662) (0.875,-0.625) (0.825,-0.589) (0.775,-0.552) (0.725,-0.516) (0.675,-0.480) (0.635,-0.450) (0.575,-0.408) (0.530,-0.375) (0.475,-0.337) (0.425,-0.301) (0.375,-0.266) (0.325,-0.230) (0.284,-0.200) (0.225,-0.160) (0.206,-0.144)}
\EQ{(0.206,0.144) (0.248,0.175) (0.300,0.213) (0.350,0.248) (0.400,0.284) (0.450,0.319) (0.495,0.350) (0.550,0.390) (0.599,0.425) (0.650,0.462) (0.700,0.498) (0.750,0.534) (0.800,0.570) (0.842,0.600) (0.900,0.644) (0.943,0.675) (1.000,0.718) (1.043,0.750) (1.100,0.794) (1.141,0.825) (1.200,0.871) (1.250,0.910) (1.300,0.950) (1.350,0.991) (1.392,1.025) (1.450,1.074) (1.500,1.117) (1.538,1.150) (1.595,1.200) (1.649,1.250) (1.700,1.298) (1.750,1.346) (1.800,1.395) (1.850,1.447) (1.900,1.499) (1.946,1.550) (1.990,1.600) (2.025,1.641) (2.074,1.700) (2.113,1.750) (2.150,1.799) (2.188,1.850) (2.222,1.900) (2.256,1.950) (2.288,2.000) (2.318,2.050) (2.348,2.100) (2.375,2.149) (2.400,2.197) (2.425,2.247) (2.450,2.299) (2.473,2.350) (2.494,2.400) (2.514,2.450) (2.532,2.500) (2.549,2.550) (2.565,2.600) (2.580,2.650) (2.594,2.700) (2.606,2.750) (2.617,2.800) (2.627,2.850) (2.636,2.900) (2.644,2.950) (2.651,3.000) (2.656,3.050) (2.660,3.100) (2.664,3.150) (2.666,3.200) (2.667,3.250) (2.666,3.300) (2.665,3.350) (2.663,3.400) (2.659,3.450) (2.654,3.500) (2.648,3.550) (2.641,3.600) (2.633,3.650) (2.624,3.700) (2.613,3.750) (2.602,3.800) (2.589,3.850) (2.575,3.899) (2.559,3.950) (2.543,4.000) (2.525,4.050) (2.506,4.100) (2.485,4.150) (2.463,4.200) (2.440,4.250) (2.415,4.300) (2.389,4.350) (2.361,4.400) (2.332,4.450) (2.301,4.500) (2.268,4.550) (2.234,4.600) (2.200,4.647) (2.160,4.700) (2.125,4.743) (2.077,4.800) (2.032,4.850) (1.985,4.900) (1.935,4.950) (1.882,5.000) (1.826,5.050) (1.775,5.094) (1.725,5.134) (1.675,5.172) (1.625,5.209) (1.575,5.244) (1.528,5.275) (1.475,5.308) (1.425,5.338) (1.375,5.367) (1.325,5.394) (1.275,5.420) (1.225,5.445) (1.175,5.468) (1.125,5.491) (1.075,5.512) (1.025,5.532) (0.976,5.550) (0.925,5.568) (0.875,5.585) (0.826,5.600) (0.775,5.615) (0.725,5.629) (0.675,5.641) (0.625,5.653) (0.575,5.664) (0.525,5.674) (0.475,5.683) (0.425,5.691) (0.375,5.698) (0.325,5.704) (0.275,5.709) (0.225,5.713) (0.175,5.717) (0.125,5.720) (0.075,5.721) (0.025,5.722) (-0.050,5.722) (-0.100,5.721) (-0.150,5.718) (-0.200,5.715) (-0.250,5.711) (-0.300,5.707) (-0.350,5.701) (-0.400,5.694) (-0.450,5.687) (-0.500,5.678) (-0.550,5.669) (-0.600,5.659) (-0.650,5.647) (-0.700,5.635) (-0.750,5.622) (-0.800,5.608) (-0.850,5.593) (-0.900,5.577) (-0.950,5.559) (-1.000,5.541) (-1.050,5.522) (-1.100,5.501) (-1.150,5.480) (-1.200,5.457) (-1.250,5.433) (-1.300,5.407) (-1.350,5.381) (-1.400,5.353) (-1.447,5.325) (-1.500,5.293) (-1.550,5.260) (-1.600,5.226) (-1.650,5.191) (-1.700,5.153) (-1.750,5.114) (-1.797,5.075) (-1.850,5.029) (-1.900,4.983) (-1.950,4.935) (-2.000,4.884) (-2.050,4.830) (-2.098,4.775) (-2.140,4.725) (-2.175,4.680) (-2.216,4.625) (-2.250,4.577) (-2.285,4.525) (-2.317,4.475) (-2.347,4.425) (-2.375,4.376) (-2.400,4.330) (-2.425,4.281) (-2.450,4.229) (-2.474,4.175) (-2.496,4.125) (-2.515,4.075) (-2.534,4.025) (-2.551,3.975) (-2.567,3.925) (-2.582,3.875) (-2.595,3.825) (-2.608,3.775) (-2.619,3.725) (-2.629,3.675) (-2.637,3.625) (-2.645,3.575) (-2.651,3.525) (-2.657,3.475) (-2.661,3.425) (-2.664,3.375) (-2.666,3.325) (-2.667,3.275) (-2.666,3.225) (-2.665,3.175) (-2.662,3.125) (-2.658,3.075) (-2.654,3.025) (-2.647,2.975) (-2.640,2.925) (-2.632,2.875) (-2.623,2.825) (-2.612,2.775) (-2.600,2.725) (-2.587,2.675) (-2.573,2.625) (-2.558,2.575) (-2.541,2.525) (-2.523,2.475) (-2.504,2.425) (-2.483,2.375) (-2.462,2.325) (-2.439,2.275) (-2.414,2.225) (-2.389,2.175) (-2.362,2.125) (-2.333,2.075) (-2.303,2.025) (-2.275,1.980) (-2.239,1.925) (-2.205,1.875) (-2.175,1.833) (-2.132,1.775) (-2.100,1.733) (-2.053,1.675) (-2.012,1.625) (-1.975,1.583) (-1.925,1.527) (-1.877,1.475) (-1.829,1.425) (-1.780,1.375) (-1.729,1.325) (-1.676,1.275) (-1.625,1.228) (-1.575,1.183) (-1.525,1.139) (-1.481,1.100) (-1.425,1.053) (-1.375,1.011) (-1.331,0.975) (-1.275,0.930) (-1.225,0.891) (-1.175,0.851) (-1.125,0.813) (-1.076,0.775) (-1.025,0.737) (-0.976,0.700) (-0.925,0.662) (-0.875,0.625) (-0.825,0.589) (-0.775,0.552) (-0.725,0.516) (-0.675,0.480) (-0.635,0.450) (-0.575,0.408) (-0.530,0.375) (-0.475,0.337) (-0.425,0.301) (-0.375,0.266) (-0.325,0.230) (-0.284,0.200) (-0.225,0.160) (-0.220,0.155)}
\EQ{(-0.220,-0.155) (-0.275,-0.195) (-0.319,-0.225) (-0.375,-0.266) (-0.424,-0.300) (-0.475,-0.337) (-0.525,-0.372) (-0.575,-0.408) (-0.625,-0.444) (-0.669,-0.475) (-0.725,-0.516) (-0.772,-0.550) (-0.825,-0.589) (-0.874,-0.625) (-0.925,-0.663) (-0.975,-0.700) (-1.025,-0.738) (-1.074,-0.775) (-1.125,-0.815) (-1.171,-0.850) (-1.225,-0.893) (-1.265,-0.925) (-1.325,-0.974) (-1.375,-1.016) (-1.416,-1.050) (-1.473,-1.100) (-1.525,-1.146) (-1.575,-1.191) (-1.612,-1.225) (-1.650,-1.261) (-1.691,-1.300) (-1.741,-1.350) (-1.790,-1.400) (-1.825,-1.438) (-1.860,-1.475) (-1.900,-1.520) (-1.947,-1.575) (-1.989,-1.625) (-2.025,-1.670) (-2.067,-1.725) (-2.100,-1.770) (-2.139,-1.825) (-2.173,-1.875) (-2.205,-1.925) (-2.236,-1.975) (-2.265,-2.025) (-2.293,-2.075) (-2.319,-2.125) (-2.344,-2.175) (-2.368,-2.225) (-2.390,-2.275) (-2.411,-2.325) (-2.431,-2.375) (-2.449,-2.425) (-2.466,-2.475) (-2.481,-2.525) (-2.496,-2.575) (-2.509,-2.625) (-2.521,-2.675) (-2.532,-2.725) (-2.541,-2.775) (-2.549,-2.825) (-2.557,-2.875) (-2.562,-2.925) (-2.567,-2.975) (-2.571,-3.025) (-2.573,-3.075) (-2.574,-3.125) (-2.574,-3.175) (-2.573,-3.225) (-2.571,-3.275) (-2.568,-3.325) (-2.563,-3.375) (-2.557,-3.425) (-2.550,-3.475) (-2.542,-3.525) (-2.533,-3.575) (-2.522,-3.625) (-2.510,-3.675) (-2.497,-3.725) (-2.483,-3.775) (-2.467,-3.825) (-2.450,-3.875) (-2.432,-3.925) (-2.413,-3.975) (-2.392,-4.025) (-2.369,-4.075) (-2.345,-4.125) (-2.320,-4.175) (-2.293,-4.225) (-2.264,-4.275) (-2.234,-4.325) (-2.202,-4.375) (-2.168,-4.425) (-2.133,-4.475) (-2.100,-4.518) (-2.055,-4.575) (-2.013,-4.625) (-1.975,-4.668) (-1.925,-4.722) (-1.875,-4.772) (-1.825,-4.820) (-1.775,-4.866) (-1.735,-4.900) (-1.675,-4.949) (-1.625,-4.988) (-1.575,-5.025) (-1.525,-5.060) (-1.475,-5.093) (-1.425,-5.124) (-1.375,-5.155) (-1.325,-5.183) (-1.275,-5.210) (-1.225,-5.236) (-1.175,-5.261) (-1.125,-5.284) (-1.075,-5.306) (-1.025,-5.327) (-0.975,-5.346) (-0.925,-5.365) (-0.875,-5.382) (-0.825,-5.398) (-0.775,-5.414) (-0.725,-5.428) (-0.675,-5.441) (-0.625,-5.453) (-0.575,-5.464) (-0.525,-5.474) (-0.475,-5.484) (-0.425,-5.492) (-0.375,-5.499) (-0.325,-5.506) (-0.275,-5.511) (-0.225,-5.516) (-0.175,-5.519) (-0.125,-5.522) (-0.075,-5.524) (-0.025,-5.525) (0.050,-5.524) (0.100,-5.523) (0.150,-5.521) (0.200,-5.518) (0.250,-5.513) (0.300,-5.508) (0.350,-5.503) (0.400,-5.496) (0.450,-5.488) (0.500,-5.479) (0.550,-5.469) (0.600,-5.459) (0.650,-5.447) (0.700,-5.434) (0.750,-5.421) (0.800,-5.406) (0.850,-5.390) (0.900,-5.374) (0.950,-5.356) (1.000,-5.337) (1.050,-5.316) (1.100,-5.295) (1.150,-5.272) (1.197,-5.250) (1.247,-5.225) (1.294,-5.200) (1.340,-5.175) (1.400,-5.140) (1.450,-5.109) (1.500,-5.076) (1.550,-5.042) (1.600,-5.006) (1.642,-4.975) (1.700,-4.929) (1.750,-4.887) (1.792,-4.850) (1.846,-4.800) (1.898,-4.750) (1.946,-4.700) (1.991,-4.650) (2.025,-4.611) (2.075,-4.550) (2.114,-4.500) (2.150,-4.451) (2.185,-4.400) (2.218,-4.350) (2.249,-4.300) (2.275,-4.257) (2.300,-4.212) (2.325,-4.165) (2.350,-4.115) (2.375,-4.062) (2.400,-4.005) (2.423,-3.950) (2.442,-3.900) (2.459,-3.850) (2.475,-3.801) (2.490,-3.750) (2.504,-3.700) (2.516,-3.650) (2.528,-3.600) (2.538,-3.550) (2.546,-3.500) (2.554,-3.450) (2.560,-3.400) (2.566,-3.350) (2.570,-3.300) (2.572,-3.250) (2.574,-3.200) (2.575,-3.150) (2.574,-3.100) (2.572,-3.050) (2.569,-3.000) (2.565,-2.950) (2.560,-2.900) (2.553,-2.850) (2.545,-2.800) (2.537,-2.750) (2.526,-2.700) (2.515,-2.650) (2.502,-2.600) (2.489,-2.550) (2.474,-2.500) (2.457,-2.450) (2.440,-2.400) (2.421,-2.350) (2.401,-2.300) (2.379,-2.250) (2.356,-2.200) (2.332,-2.150) (2.306,-2.100) (2.279,-2.050) (2.251,-2.000) (2.221,-1.950) (2.189,-1.900) (2.156,-1.850) (2.125,-1.805) (2.086,-1.750) (2.050,-1.702) (2.009,-1.650) (1.975,-1.608) (1.926,-1.550) (1.882,-1.500) (1.837,-1.450) (1.800,-1.411) (1.765,-1.375) (1.725,-1.334) (1.675,-1.285) (1.639,-1.250) (1.600,-1.214) (1.558,-1.175) (1.502,-1.125) (1.450,-1.080) (1.400,-1.037) (1.356,-1.000) (1.300,-0.954) (1.250,-0.913) (1.202,-0.875) (1.150,-0.834) (1.107,-0.800) (1.050,-0.757) (1.008,-0.725) (0.950,-0.682) (0.908,-0.650) (0.850,-0.607) (0.807,-0.575) (0.750,-0.534) (0.703,-0.500) (0.650,-0.462) (0.600,-0.426) (0.550,-0.390) (0.500,-0.355) (0.450,-0.319) (0.400,-0.284) (0.354,-0.250) (0.300,-0.213) (0.250,-0.177) (0.206,-0.144)}
\EQ{(0.206,0.144) (0.248,0.175) (0.300,0.213) (0.350,0.248) (0.400,0.284) (0.450,0.319) (0.495,0.350) (0.550,0.390) (0.599,0.425) (0.650,0.462) (0.700,0.498) (0.750,0.534) (0.800,0.571) (0.841,0.600) (0.900,0.644) (0.942,0.675) (1.000,0.719) (1.041,0.750) (1.100,0.795) (1.150,0.834) (1.200,0.873) (1.250,0.913) (1.296,0.950) (1.350,0.995) (1.400,1.037) (1.445,1.075) (1.500,1.123) (1.550,1.168) (1.600,1.214) (1.639,1.250) (1.675,1.285) (1.716,1.325) (1.765,1.375) (1.800,1.411) (1.837,1.450) (1.875,1.492) (1.925,1.549) (1.968,1.600) (2.000,1.639) (2.048,1.700) (2.086,1.750) (2.122,1.800) (2.156,1.850) (2.189,1.900) (2.221,1.950) (2.250,1.999) (2.275,2.043) (2.300,2.088) (2.325,2.136) (2.350,2.187) (2.375,2.241) (2.400,2.299) (2.421,2.350) (2.440,2.400) (2.457,2.450) (2.474,2.500) (2.489,2.550) (2.502,2.600) (2.515,2.650) (2.526,2.700) (2.537,2.750) (2.545,2.800) (2.553,2.850) (2.560,2.900) (2.565,2.950) (2.569,3.000) (2.572,3.050) (2.574,3.100) (2.575,3.150) (2.574,3.200) (2.572,3.250) (2.570,3.300) (2.566,3.350) (2.560,3.400) (2.554,3.450) (2.546,3.500) (2.538,3.550) (2.528,3.600) (2.516,3.650) (2.504,3.700) (2.490,3.750) (2.475,3.800) (2.459,3.850) (2.442,3.900) (2.423,3.950) (2.402,4.000) (2.381,4.050) (2.357,4.100) (2.333,4.150) (2.307,4.200) (2.279,4.250) (2.250,4.299) (2.218,4.350) (2.185,4.400) (2.151,4.450) (2.114,4.500) (2.075,4.550) (2.034,4.600) (2.000,4.640) (1.950,4.695) (1.900,4.747) (1.850,4.797) (1.800,4.843) (1.750,4.887) (1.705,4.925) (1.650,4.969) (1.609,5.000) (1.550,5.042) (1.502,5.075) (1.450,5.109) (1.400,5.140) (1.350,5.169) (1.300,5.197) (1.250,5.223) (1.200,5.249) (1.150,5.272) (1.100,5.295) (1.050,5.316) (1.000,5.337) (0.950,5.356) (0.900,5.374) (0.850,5.390) (0.800,5.406) (0.750,5.421) (0.700,5.434) (0.650,5.447) (0.600,5.459) (0.550,5.469) (0.500,5.479) (0.450,5.488) (0.400,5.496) (0.350,5.503) (0.300,5.508) (0.250,5.513) (0.200,5.518) (0.150,5.521) (0.100,5.523) (0.050,5.524) (0.000,5.525) (-0.050,5.524) (-0.100,5.523) (-0.150,5.521) (-0.200,5.518) (-0.250,5.513) (-0.300,5.508) (-0.350,5.503) (-0.400,5.496) (-0.450,5.488) (-0.500,5.479) (-0.550,5.469) (-0.600,5.459) (-0.650,5.447) (-0.700,5.434) (-0.750,5.421) (-0.800,5.406) (-0.850,5.390) (-0.900,5.374) (-0.950,5.356) (-1.000,5.337) (-1.050,5.316) (-1.100,5.295) (-1.150,5.272) (-1.197,5.250) (-1.247,5.225) (-1.294,5.200) (-1.340,5.175) (-1.400,5.140) (-1.450,5.109) (-1.500,5.076) (-1.550,5.042) (-1.600,5.006) (-1.642,4.975) (-1.700,4.929) (-1.750,4.887) (-1.792,4.850) (-1.846,4.800) (-1.898,4.750) (-1.946,4.700) (-1.991,4.650) (-2.025,4.611) (-2.075,4.550) (-2.114,4.500) (-2.150,4.451) (-2.185,4.400) (-2.218,4.350) (-2.249,4.300) (-2.275,4.257) (-2.300,4.212) (-2.325,4.165) (-2.350,4.115) (-2.375,4.062) (-2.400,4.005) (-2.423,3.950) (-2.442,3.900) (-2.459,3.850) (-2.475,3.801) (-2.490,3.750) (-2.504,3.700) (-2.516,3.650) (-2.528,3.600) (-2.538,3.550) (-2.546,3.500) (-2.554,3.450) (-2.560,3.400) (-2.566,3.350) (-2.570,3.300) (-2.572,3.250) (-2.574,3.200) (-2.575,3.150) (-2.574,3.100) (-2.572,3.050) (-2.569,3.000) (-2.565,2.950) (-2.560,2.900) (-2.553,2.850) (-2.545,2.800) (-2.537,2.750) (-2.526,2.700) (-2.515,2.650) (-2.502,2.600) (-2.489,2.550) (-2.474,2.500) (-2.457,2.450) (-2.440,2.400) (-2.421,2.350) (-2.401,2.300) (-2.379,2.250) (-2.356,2.200) (-2.332,2.150) (-2.306,2.100) (-2.279,2.050) (-2.251,2.000) (-2.221,1.950) (-2.189,1.900) (-2.156,1.850) (-2.125,1.805) (-2.086,1.750) (-2.050,1.702) (-2.009,1.650) (-1.975,1.608) (-1.926,1.550) (-1.882,1.500) (-1.837,1.450) (-1.800,1.411) (-1.765,1.375) (-1.725,1.334) (-1.675,1.285) (-1.639,1.250) (-1.600,1.214) (-1.558,1.175) (-1.502,1.125) (-1.450,1.080) (-1.400,1.037) (-1.356,1.000) (-1.300,0.954) (-1.250,0.913) (-1.202,0.875) (-1.150,0.834) (-1.107,0.800) (-1.050,0.757) (-1.008,0.725) (-0.950,0.682) (-0.908,0.650) (-0.850,0.607) (-0.807,0.575) (-0.750,0.534) (-0.703,0.500) (-0.650,0.462) (-0.600,0.426) (-0.550,0.390) (-0.500,0.355) (-0.450,0.319) (-0.400,0.284) (-0.354,0.250) (-0.300,0.213) (-0.250,0.177) (-0.220,0.155)}

    % Cuadrupolo puntual - líneas de campo
% Generado por scripts/generar_datos.py (no editar a mano).
\LC{(0.435,1.114) (0.489,1.167) (0.546,1.216) (0.605,1.261) (0.667,1.304) (0.732,1.342) (0.799,1.375) (0.868,1.405) (0.939,1.429) (1.011,1.448) (1.085,1.461) (1.160,1.468) (1.235,1.468) (1.310,1.461) (1.383,1.446) (1.455,1.424) (1.523,1.394) (1.588,1.355) (1.646,1.309) (1.698,1.255) (1.742,1.194) (1.777,1.128) (1.803,1.057) (1.817,0.984) (1.822,0.909) (1.817,0.834) (1.803,0.761) (1.780,0.689) (1.749,0.621) (1.712,0.556) (1.669,0.495) (1.621,0.437) (1.568,0.384) (1.511,0.335) (1.451,0.290) (1.388,0.249) (1.323,0.213) (1.255,0.179) (1.186,0.150)}
\LC{(-0.449,1.110) (-0.504,1.161) (-0.562,1.208) (-0.623,1.252) (-0.687,1.292) (-0.753,1.327) (-0.821,1.358) (-0.891,1.384) (-0.963,1.405) (-1.037,1.420) (-1.111,1.428) (-1.186,1.430) (-1.261,1.424) (-1.335,1.411) (-1.407,1.390) (-1.476,1.361) (-1.541,1.323) (-1.600,1.278) (-1.652,1.224) (-1.696,1.163) (-1.731,1.097) (-1.756,1.026) (-1.770,0.953) (-1.774,0.878) (-1.768,0.803) (-1.753,0.730) (-1.728,0.659) (-1.697,0.591) (-1.658,0.527) (-1.613,0.467) (-1.564,0.411) (-1.510,0.359) (-1.452,0.311) (-1.391,0.268) (-1.327,0.228) (-1.261,0.193) (-1.193,0.161) (-1.170,0.152)}
\LC{(-0.693,-0.961) (-0.768,-0.970) (-0.843,-0.968) (-0.916,-0.955) (-0.987,-0.930) (-1.052,-0.893) (-1.109,-0.844) (-1.154,-0.784) (-1.185,-0.716) (-1.201,-0.643) (-1.202,-0.568) (-1.189,-0.495) (-1.163,-0.424) (-1.126,-0.359)}
\LC{(0.153,-1.185) (0.173,-1.258) (0.193,-1.330) (0.215,-1.401) (0.239,-1.473) (0.263,-1.544) (0.289,-1.614) (0.316,-1.684) (0.343,-1.754) (0.373,-1.823) (0.403,-1.892) (0.434,-1.960) (0.467,-2.027) (0.501,-2.094) (0.536,-2.160) (0.572,-2.226) (0.609,-2.291) (0.648,-2.356) (0.687,-2.419) (0.728,-2.482) (0.770,-2.545) (0.813,-2.606) (0.857,-2.667) (0.902,-2.727) (0.948,-2.786) (0.995,-2.844) (1.044,-2.901) (1.094,-2.957) (1.144,-3.013) (1.196,-3.067) (1.249,-3.120) (1.303,-3.172) (1.358,-3.223) (1.413,-3.273) (1.470,-3.322) (1.528,-3.370) (1.587,-3.416) (1.647,-3.461) (1.708,-3.505) (1.770,-3.547) (1.833,-3.588) (1.897,-3.628) (1.961,-3.666) (2.027,-3.703) (2.093,-3.738) (2.160,-3.771) (2.228,-3.803) (2.297,-3.833) (2.366,-3.861) (2.436,-3.888) (2.507,-3.913) (2.578,-3.936) (2.651,-3.956) (2.723,-3.975) (2.796,-3.992) (2.870,-4.007) (2.944,-4.019) (3.018,-4.030) (3.093,-4.038) (3.167,-4.043) (3.242,-4.046) (3.317,-4.047) (3.392,-4.046) (3.467,-4.041) (3.542,-4.035) (3.616,-4.025) (3.690,-4.013) (3.764,-3.998) (3.836,-3.980) (3.909,-3.959) (3.980,-3.935) (4.050,-3.909) (4.119,-3.879) (4.186,-3.847) (4.253,-3.812) (4.317,-3.773) (4.380,-3.732) (4.440,-3.688) (4.499,-3.641) (4.555,-3.591) (4.608,-3.539) (4.659,-3.484) (4.707,-3.426) (4.752,-3.366) (4.794,-3.304) (4.833,-3.240) (4.868,-3.173) (4.900,-3.105) (4.928,-3.036) (4.953,-2.965) (4.973,-2.893) (4.991,-2.820) (5.004,-2.746) (5.014,-2.672) (5.020,-2.597) (5.022,-2.522) (5.021,-2.447) (5.016,-2.372) (5.008,-2.298) (4.996,-2.224) (4.981,-2.150) (4.963,-2.077) (4.942,-2.005) (4.918,-1.934) (4.891,-1.864) (4.862,-1.795) (4.830,-1.728) (4.795,-1.661) (4.758,-1.596) (4.719,-1.532) (4.677,-1.469) (4.634,-1.408) (4.588,-1.349) (4.541,-1.291) (4.492,-1.234) (4.441,-1.179) (4.389,-1.125) (4.335,-1.073) (4.280,-1.022) (4.223,-0.973) (4.165,-0.925) (4.106,-0.879) (4.046,-0.834) (3.985,-0.791) (3.923,-0.749) (3.859,-0.709) (3.795,-0.670) (3.730,-0.632) (3.665,-0.596) (3.598,-0.561) (3.531,-0.528) (3.463,-0.496) (3.395,-0.465) (3.326,-0.436) (3.256,-0.407) (3.186,-0.380) (3.116,-0.355) (3.045,-0.330) (2.974,-0.306) (2.902,-0.284) (2.830,-0.263) (2.758,-0.243) (2.686,-0.224) (2.613,-0.206) (2.540,-0.189) (2.467,-0.172) (2.393,-0.157) (2.319,-0.143) (2.246,-0.130) (2.172,-0.117) (2.098,-0.105) (2.023,-0.095) (1.949,-0.084) (1.875,-0.075) (1.800,-0.066) (1.751,-0.061)}
\LC{(0.283,-2.625) (0.299,-2.698) (0.316,-2.771) (0.333,-2.844) (0.350,-2.917) (0.368,-2.990) (0.387,-3.063) (0.406,-3.135) (0.425,-3.208) (0.445,-3.280) (0.466,-3.352) (0.486,-3.424) (0.507,-3.496) (0.529,-3.568) (0.551,-3.640) (0.574,-3.711) (0.597,-3.783) (0.620,-3.854) (0.644,-3.925) (0.669,-3.996) (0.693,-4.067) (0.719,-4.137) (0.744,-4.208) (0.770,-4.278) (0.797,-4.348) (0.824,-4.418) (0.852,-4.488) (0.879,-4.557) (0.908,-4.627) (0.937,-4.696) (0.966,-4.765) (0.996,-4.834) (1.026,-4.903) (1.056,-4.971) (1.087,-5.039) (1.119,-5.108) (1.151,-5.175) (1.183,-5.243) (1.216,-5.310) (1.249,-5.378) (1.283,-5.445) (1.317,-5.511) (1.352,-5.578) (1.387,-5.644) (1.422,-5.710) (1.458,-5.776) (1.494,-5.842) (1.531,-5.907) (1.568,-5.972) (1.606,-6.037) (1.644,-6.102) (1.682,-6.166) (1.721,-6.230) (1.761,-6.294) (1.801,-6.358) (1.841,-6.421) (1.882,-6.484) (1.923,-6.547) (1.964,-6.609) (2.006,-6.671) (2.048,-6.733) (2.091,-6.795) (2.135,-6.856) (2.178,-6.917) (2.222,-6.978) (2.267,-7.038) (2.312,-7.098) (2.357,-7.158) (2.403,-7.217) (2.449,-7.276) (2.496,-7.335) (2.543,-7.394) (2.590,-7.452) (2.638,-7.509) (2.686,-7.567) (2.735,-7.624) (2.784,-7.680) (2.834,-7.737) (2.884,-7.793) (2.934,-7.848) (2.985,-7.903) (3.036,-7.958) (3.088,-8.013) (3.140,-8.067) (3.192,-8.120) (3.245,-8.174) (3.298,-8.227) (3.352,-8.279) (3.406,-8.331) (3.460,-8.383) (3.515,-8.434) (3.570,-8.485) (3.626,-8.535) (3.682,-8.585) (3.738,-8.634) (3.795,-8.683) (3.852,-8.732) (3.909,-8.780) (3.967,-8.828) (4.026,-8.875) (4.084,-8.922) (4.143,-8.968) (4.203,-9.014) (4.262,-9.059) (4.323,-9.104) (4.383,-9.148) (4.444,-9.192) (4.505,-9.235) (4.567,-9.278) (4.629,-9.320) (4.691,-9.362) (4.754,-9.403) (4.817,-9.444) (4.880,-9.484) (4.944,-9.524) (5.008,-9.563) (5.072,-9.601) (5.137,-9.639) (5.202,-9.676) (5.267,-9.713) (5.333,-9.749) (5.399,-9.785) (5.466,-9.820) (5.532,-9.854) (5.599,-9.888) (5.666,-9.921) (5.734,-9.954) (5.802,-9.986) (5.847,-10.000)}
\LC{(-0.162,-1.698) (-0.176,-1.772) (-0.191,-1.845) (-0.207,-1.919) (-0.223,-1.992) (-0.240,-2.065) (-0.258,-2.138) (-0.276,-2.210) (-0.295,-2.283) (-0.314,-2.356) (-0.334,-2.428) (-0.355,-2.500) (-0.376,-2.572) (-0.398,-2.644) (-0.421,-2.715) (-0.444,-2.786) (-0.468,-2.858) (-0.492,-2.928) (-0.517,-2.999) (-0.543,-3.070) (-0.569,-3.140) (-0.596,-3.210) (-0.623,-3.280) (-0.651,-3.349) (-0.680,-3.419) (-0.709,-3.488) (-0.739,-3.557) (-0.769,-3.625) (-0.800,-3.693) (-0.832,-3.761) (-0.864,-3.829) (-0.897,-3.896) (-0.931,-3.963) (-0.965,-4.030) (-1.000,-4.097) (-1.035,-4.163) (-1.071,-4.229) (-1.107,-4.294) (-1.144,-4.359) (-1.182,-4.424) (-1.220,-4.489) (-1.259,-4.553) (-1.299,-4.617) (-1.339,-4.680) (-1.380,-4.743) (-1.421,-4.806) (-1.463,-4.868) (-1.505,-4.930) (-1.548,-4.991) (-1.592,-5.052) (-1.636,-5.113) (-1.681,-5.173) (-1.726,-5.233) (-1.772,-5.292) (-1.819,-5.351) (-1.866,-5.409) (-1.913,-5.467) (-1.961,-5.525) (-2.010,-5.582) (-2.060,-5.638) (-2.110,-5.694) (-2.160,-5.749) (-2.211,-5.804) (-2.263,-5.859) (-2.315,-5.913) (-2.368,-5.966) (-2.421,-6.019) (-2.475,-6.071) (-2.529,-6.123) (-2.584,-6.174) (-2.639,-6.224) (-2.695,-6.274) (-2.752,-6.324) (-2.809,-6.372) (-2.866,-6.420) (-2.924,-6.468) (-2.983,-6.515) (-3.042,-6.561) (-3.102,-6.606) (-3.162,-6.651) (-3.222,-6.695) (-3.284,-6.739) (-3.345,-6.782) (-3.407,-6.824) (-3.470,-6.865) (-3.533,-6.906) (-3.597,-6.945) (-3.661,-6.984) (-3.725,-7.023) (-3.790,-7.060) (-3.855,-7.097) (-3.921,-7.133) (-3.987,-7.168) (-4.054,-7.202) (-4.121,-7.236) (-4.189,-7.268) (-4.257,-7.300) (-4.325,-7.331) (-4.394,-7.361) (-4.463,-7.390) (-4.533,-7.418) (-4.602,-7.445) (-4.673,-7.472) (-4.743,-7.497) (-4.814,-7.521) (-4.886,-7.545) (-4.957,-7.567) (-5.029,-7.588) (-5.101,-7.609) (-5.174,-7.628) (-5.246,-7.646) (-5.319,-7.663) (-5.393,-7.679) (-5.466,-7.694) (-5.540,-7.708) (-5.614,-7.721) (-5.688,-7.732) (-5.762,-7.743) (-5.837,-7.752) (-5.911,-7.760) (-5.986,-7.767) (-6.061,-7.773) (-6.136,-7.777) (-6.210,-7.780) (-6.285,-7.782) (-6.360,-7.782) (-6.435,-7.782) (-6.510,-7.779) (-6.585,-7.776) (-6.660,-7.771) (-6.735,-7.765) (-6.810,-7.758) (-6.884,-7.749) (-6.958,-7.739) (-7.032,-7.727) (-7.106,-7.714) (-7.180,-7.699) (-7.253,-7.683) (-7.326,-7.665) (-7.399,-7.646) (-7.471,-7.626) (-7.542,-7.604) (-7.614,-7.581) (-7.684,-7.556) (-7.754,-7.529) (-7.824,-7.501) (-7.893,-7.471) (-7.961,-7.440) (-8.029,-7.408) (-8.095,-7.373) (-8.161,-7.338) (-8.227,-7.300) (-8.291,-7.262) (-8.354,-7.221) (-8.416,-7.180) (-8.477,-7.136) (-8.538,-7.091) (-8.597,-7.045) (-8.654,-6.997) (-8.711,-6.948) (-8.766,-6.898) (-8.820,-6.845) (-8.873,-6.792) (-8.924,-6.737) (-8.974,-6.681) (-9.022,-6.624) (-9.069,-6.565) (-9.114,-6.505) (-9.158,-6.444) (-9.200,-6.382) (-9.240,-6.319) (-9.278,-6.254) (-9.315,-6.189) (-9.350,-6.123) (-9.383,-6.055) (-9.415,-5.987) (-9.444,-5.918) (-9.472,-5.849) (-9.497,-5.778) (-9.521,-5.707) (-9.543,-5.635) (-9.563,-5.563) (-9.581,-5.490) (-9.597,-5.417) (-9.611,-5.343) (-9.623,-5.269) (-9.634,-5.195) (-9.642,-5.120) (-9.648,-5.046) (-9.653,-4.971) (-9.656,-4.896) (-9.656,-4.821) (-9.655,-4.746) (-9.652,-4.671) (-9.648,-4.596) (-9.641,-4.521) (-9.633,-4.447) (-9.622,-4.373) (-9.611,-4.298) (-9.597,-4.225) (-9.582,-4.151) (-9.565,-4.078) (-9.546,-4.006) (-9.526,-3.933) (-9.505,-3.861) (-9.482,-3.790) (-9.457,-3.719) (-9.431,-3.649) (-9.403,-3.579) (-9.374,-3.510) (-9.344,-3.442) (-9.312,-3.374) (-9.279,-3.306) (-9.245,-3.240) (-9.209,-3.174) (-9.172,-3.108) (-9.134,-3.044) (-9.095,-2.980) (-9.055,-2.916) (-9.013,-2.854) (-8.971,-2.792) (-8.927,-2.731) (-8.882,-2.671) (-8.837,-2.611) (-8.790,-2.553) (-8.743,-2.495) (-8.694,-2.437) (-8.645,-2.381) (-8.595,-2.325) (-8.544,-2.270) (-8.492,-2.216) (-8.439,-2.163) (-8.385,-2.110) (-8.331,-2.058) (-8.276,-2.007) (-8.221,-1.957) (-8.164,-1.908) (-8.107,-1.859) (-8.049,-1.811) (-7.991,-1.764) (-7.932,-1.718) (-7.872,-1.672) (-7.812,-1.627) (-7.751,-1.583) (-7.690,-1.540) (-7.628,-1.498) (-7.566,-1.456) (-7.503,-1.415) (-7.440,-1.375) (-7.376,-1.335) (-7.312,-1.297) (-7.247,-1.259) (-7.182,-1.221) (-7.117,-1.185) (-7.051,-1.149) (-6.985,-1.114) (-6.918,-1.079) (-6.851,-1.046) (-6.784,-1.012) (-6.716,-0.980) (-6.648,-0.948) (-6.580,-0.917) (-6.511,-0.887) (-6.442,-0.858) (-6.373,-0.829) (-6.304,-0.800) (-6.234,-0.773) (-6.164,-0.745) (-6.094,-0.719) (-6.023,-0.693) (-5.953,-0.668) (-5.882,-0.643) (-5.811,-0.620) (-5.740,-0.596) (-5.668,-0.573) (-5.596,-0.551) (-5.525,-0.530) (-5.453,-0.508) (-5.380,-0.488) (-5.308,-0.468) (-5.236,-0.449) (-5.163,-0.430) (-5.090,-0.412) (-5.018,-0.394) (-4.945,-0.376) (-4.871,-0.360) (-4.798,-0.343) (-4.725,-0.328) (-4.652,-0.312) (-4.578,-0.298) (-4.504,-0.283) (-4.431,-0.269) (-4.357,-0.256) (-4.283,-0.243) (-4.209,-0.231) (-4.135,-0.219) (-4.061,-0.207) (-3.987,-0.196) (-3.913,-0.185) (-3.838,-0.174) (-3.764,-0.164) (-3.690,-0.155) (-3.615,-0.146) (-3.541,-0.137) (-3.466,-0.128) (-3.392,-0.120) (-3.317,-0.112) (-3.242,-0.105) (-3.168,-0.098) (-3.093,-0.091) (-3.018,-0.085) (-2.944,-0.078) (-2.869,-0.073) (-2.794,-0.067) (-2.719,-0.062) (-2.644,-0.057) (-2.569,-0.052) (-2.495,-0.048) (-2.420,-0.044) (-2.345,-0.040) (-2.270,-0.036) (-2.195,-0.032) (-2.120,-0.029) (-2.045,-0.026) (-1.970,-0.023) (-1.895,-0.021) (-1.820,-0.019) (-1.770,-0.017)}
\LC{(-0.196,1.492) (-0.216,1.565) (-0.237,1.637) (-0.259,1.708) (-0.282,1.780) (-0.305,1.851) (-0.330,1.922) (-0.356,1.992) (-0.382,2.062) (-0.410,2.132) (-0.439,2.201) (-0.468,2.270) (-0.499,2.339) (-0.530,2.407) (-0.562,2.475) (-0.595,2.542) (-0.630,2.609) (-0.665,2.675) (-0.701,2.741) (-0.738,2.806) (-0.776,2.871) (-0.815,2.935) (-0.854,2.998) (-0.895,3.061) (-0.937,3.124) (-0.979,3.185) (-1.023,3.247) (-1.067,3.307) (-1.112,3.367) (-1.158,3.426) (-1.206,3.484) (-1.253,3.542) (-1.302,3.599) (-1.352,3.655) (-1.403,3.711) (-1.454,3.765) (-1.506,3.819) (-1.559,3.872) (-1.614,3.924) (-1.668,3.975) (-1.724,4.025) (-1.781,4.074) (-1.838,4.123) (-1.896,4.170) (-1.955,4.217) (-2.015,4.262) (-2.075,4.306) (-2.137,4.349) (-2.199,4.391) (-2.262,4.432) (-2.325,4.472) (-2.389,4.511) (-2.454,4.548) (-2.520,4.584) (-2.587,4.619) (-2.654,4.653) (-2.721,4.685) (-2.790,4.716) (-2.859,4.745) (-2.928,4.773) (-2.998,4.800) (-3.069,4.825) (-3.140,4.849) (-3.212,4.871) (-3.284,4.891) (-3.357,4.910) (-3.430,4.928) (-3.503,4.943) (-3.577,4.957) (-3.651,4.969) (-3.725,4.979) (-3.800,4.987) (-3.874,4.994) (-3.949,4.998) (-4.024,5.001) (-4.099,5.002) (-4.174,5.000) (-4.249,4.997) (-4.324,4.991) (-4.398,4.983) (-4.473,4.973) (-4.547,4.961) (-4.620,4.947) (-4.693,4.930) (-4.766,4.911) (-4.838,4.890) (-4.909,4.866) (-4.979,4.840) (-5.049,4.812) (-5.117,4.781) (-5.185,4.748) (-5.251,4.713) (-5.315,4.675) (-5.379,4.635) (-5.440,4.592) (-5.500,4.547) (-5.559,4.500) (-5.615,4.450) (-5.669,4.399) (-5.722,4.345) (-5.772,4.289) (-5.819,4.231) (-5.865,4.171) (-5.907,4.110) (-5.947,4.046) (-5.985,3.981) (-6.019,3.915) (-6.051,3.847) (-6.080,3.778) (-6.106,3.707) (-6.129,3.636) (-6.149,3.564) (-6.166,3.491) (-6.180,3.417) (-6.191,3.343) (-6.199,3.268) (-6.204,3.193) (-6.206,3.118) (-6.205,3.043) (-6.201,2.968) (-6.195,2.894) (-6.186,2.819) (-6.174,2.745) (-6.159,2.672) (-6.142,2.599) (-6.122,2.526) (-6.100,2.455) (-6.075,2.384) (-6.049,2.314) (-6.020,2.245) (-5.989,2.176) (-5.955,2.109) (-5.920,2.043) (-5.883,1.978) (-5.844,1.913) (-5.804,1.850) (-5.761,1.789) (-5.717,1.728) (-5.672,1.668) (-5.624,1.610) (-5.576,1.553) (-5.526,1.497) (-5.475,1.442) (-5.422,1.389) (-5.368,1.336) (-5.314,1.285) (-5.258,1.235) (-5.200,1.187) (-5.142,1.139) (-5.083,1.093) (-5.023,1.048) (-4.963,1.004) (-4.901,0.961) (-4.838,0.920) (-4.775,0.880) (-4.711,0.840) (-4.647,0.802) (-4.581,0.765) (-4.515,0.730) (-4.449,0.695) (-4.382,0.661) (-4.314,0.629) (-4.246,0.597) (-4.178,0.567) (-4.109,0.538) (-4.039,0.509) (-3.969,0.482) (-3.899,0.456) (-3.828,0.430) (-3.758,0.406) (-3.686,0.382) (-3.615,0.360) (-3.543,0.338) (-3.495,0.324)}
\LC{(0.436,1.693) (0.474,1.758) (0.513,1.822) (0.553,1.885) (0.595,1.947) (0.639,2.009) (0.684,2.069) (0.730,2.127) (0.778,2.185) (0.827,2.242) (0.878,2.297) (0.930,2.351) (0.984,2.403) (1.039,2.454) (1.095,2.504) (1.153,2.552) (1.212,2.598) (1.272,2.643) (1.333,2.686) (1.396,2.727) (1.460,2.766) (1.525,2.803) (1.591,2.839) (1.659,2.872) (1.727,2.903) (1.796,2.931) (1.866,2.957) (1.938,2.981) (2.009,3.003) (2.082,3.021) (2.155,3.037) (2.229,3.050) (2.304,3.060) (2.378,3.067) (2.453,3.071) (2.528,3.072) (2.603,3.070) (2.678,3.064) (2.752,3.054) (2.826,3.041) (2.899,3.024) (2.971,3.004) (3.042,2.980) (3.112,2.951) (3.180,2.919) (3.245,2.883) (3.309,2.844) (3.370,2.800) (3.428,2.753) (3.483,2.702) (3.534,2.647) (3.582,2.589) (3.626,2.528) (3.665,2.465) (3.701,2.398) (3.731,2.330) (3.757,2.259) (3.777,2.187) (3.793,2.114) (3.804,2.040) (3.811,1.965) (3.812,1.890) (3.809,1.815) (3.801,1.741) (3.788,1.667) (3.772,1.594) (3.751,1.521) (3.727,1.451) (3.699,1.381) (3.667,1.313) (3.632,1.247) (3.594,1.182) (3.553,1.119) (3.509,1.058) (3.463,0.999) (3.415,0.942) (3.364,0.887) (3.311,0.834) (3.256,0.782) (3.200,0.733) (3.141,0.686) (3.082,0.641) (3.020,0.597) (2.958,0.556) (2.894,0.516) (2.829,0.479) (2.763,0.443) (2.697,0.409) (2.629,0.376) (2.560,0.346) (2.491,0.317) (2.445,0.299)}
\LC{(0.289,2.296) (0.308,2.368) (0.328,2.441) (0.348,2.513) (0.369,2.585) (0.390,2.657) (0.412,2.729) (0.434,2.800) (0.458,2.872) (0.481,2.943) (0.506,3.014) (0.531,3.084) (0.556,3.155) (0.582,3.225) (0.609,3.295) (0.636,3.365) (0.664,3.435) (0.693,3.504) (0.722,3.573) (0.752,3.642) (0.782,3.711) (0.813,3.779) (0.844,3.847) (0.876,3.915) (0.909,3.983) (0.942,4.050) (0.976,4.117) (1.010,4.183) (1.045,4.250) (1.081,4.316) (1.117,4.381) (1.154,4.447) (1.191,4.512) (1.229,4.577) (1.267,4.641) (1.306,4.705) (1.346,4.769) (1.386,4.832) (1.427,4.895) (1.468,4.958) (1.510,5.020) (1.553,5.082) (1.596,5.143) (1.639,5.204) (1.684,5.265) (1.728,5.325) (1.774,5.385) (1.819,5.444) (1.866,5.503) (1.913,5.561) (1.960,5.619) (2.009,5.677) (2.057,5.734) (2.106,5.791) (2.156,5.847) (2.206,5.902) (2.257,5.957) (2.309,6.012) (2.361,6.066) (2.413,6.120) (2.466,6.173) (2.520,6.225) (2.574,6.277) (2.628,6.329) (2.683,6.380) (2.739,6.430) (2.795,6.480) (2.852,6.529) (2.909,6.577) (2.967,6.625) (3.025,6.673) (3.083,6.719) (3.143,6.765) (3.202,6.811) (3.262,6.856) (3.323,6.900) (3.384,6.943) (3.446,6.986) (3.508,7.028) (3.571,7.069) (3.634,7.110) (3.697,7.150) (3.761,7.189) (3.825,7.228) (3.890,7.265) (3.956,7.302) (4.021,7.338) (4.087,7.374) (4.154,7.408) (4.221,7.442) (4.288,7.475) (4.356,7.507) (4.424,7.538) (4.493,7.568) (4.562,7.598) (4.631,7.627) (4.701,7.654) (4.771,7.681) (4.841,7.707) (4.912,7.732) (4.983,7.756) (5.054,7.779) (5.126,7.801) (5.198,7.822) (5.270,7.842) (5.343,7.861) (5.416,7.879) (5.489,7.896) (5.562,7.912) (5.636,7.927) (5.709,7.941) (5.783,7.954) (5.857,7.965) (5.932,7.976) (6.006,7.985) (6.081,7.993) (6.155,8.000) (6.230,8.006) (6.305,8.011) (6.380,8.014) (6.455,8.016) (6.530,8.017) (6.605,8.017) (6.680,8.015) (6.755,8.012) (6.830,8.008) (6.904,8.002) (6.979,7.995) (7.054,7.987) (7.128,7.978) (7.202,7.967) (7.276,7.954) (7.350,7.941) (7.423,7.926) (7.496,7.909) (7.569,7.891) (7.642,7.872) (7.714,7.851) (7.785,7.828) (7.856,7.805) (7.927,7.779) (7.997,7.753) (8.067,7.724) (8.136,7.695) (8.204,7.664) (8.271,7.631) (8.338,7.597) (8.404,7.561) (8.469,7.524) (8.533,7.485) (8.597,7.445) (8.659,7.404) (8.721,7.360) (8.781,7.316) (8.840,7.270) (8.898,7.223) (8.955,7.174) (9.011,7.124) (9.066,7.072) (9.119,7.019) (9.171,6.965) (9.221,6.910) (9.270,6.853) (9.318,6.795) (9.364,6.736) (9.408,6.675) (9.451,6.614) (9.492,6.551) (9.532,6.487) (9.570,6.423) (9.606,6.357) (9.641,6.290) (9.673,6.223) (9.704,6.155) (9.733,6.086) (9.761,6.016) (9.786,5.945) (9.810,5.874) (9.832,5.802) (9.852,5.730) (9.870,5.657) (9.886,5.584) (9.900,5.510) (9.912,5.436) (9.923,5.362) (9.932,5.288) (9.938,5.213) (9.943,5.138) (9.946,5.063) (9.948,4.988) (9.947,4.913) (9.945,4.838) (9.941,4.763) (9.935,4.689) (9.927,4.614) (9.918,4.539) (9.907,4.465) (9.894,4.391) (9.880,4.318) (9.864,4.244) (9.846,4.172) (9.827,4.099) (9.806,4.027) (9.784,3.955) (9.761,3.884) (9.736,3.813) (9.709,3.743) (9.681,3.674) (9.652,3.605) (9.621,3.536) (9.589,3.468) (9.556,3.401) (9.522,3.334) (9.486,3.268) (9.449,3.203) (9.411,3.138) (9.372,3.074) (9.332,3.011) (9.291,2.948) (9.248,2.887) (9.205,2.825) (9.161,2.765) (9.115,2.705) (9.069,2.646) (9.022,2.588) (8.973,2.531) (8.924,2.474) (8.874,2.418) (8.824,2.363) (8.772,2.308) (8.720,2.254) (8.667,2.201) (8.613,2.149) (8.558,2.098) (8.503,2.047) (8.447,1.997) (8.390,1.948) (8.333,1.900) (8.275,1.852) (8.217,1.805) (8.157,1.759) (8.098,1.714) (8.037,1.669) (7.977,1.625) (7.915,1.582) (7.853,1.540) (7.791,1.498) (7.728,1.457) (7.665,1.417) (7.601,1.377) (7.537,1.338) (7.472,1.300) (7.407,1.263) (7.342,1.226) (7.276,1.190) (7.210,1.155) (7.143,1.120) (7.077,1.086) (7.009,1.053) (6.942,1.021) (6.874,0.989) (6.806,0.958) (6.737,0.927) (6.668,0.897) (6.599,0.868) (6.530,0.839) (6.461,0.811) (6.391,0.784) (6.321,0.757) (6.250,0.731) (6.180,0.705) (6.109,0.680) (6.038,0.656) (5.967,0.632) (5.896,0.609) (5.824,0.586) (5.753,0.564) (5.681,0.542) (5.609,0.521) (5.537,0.501) (5.464,0.481) (5.392,0.462) (5.319,0.443) (5.247,0.425) (5.174,0.407) (5.101,0.389) (5.028,0.373) (4.954,0.356) (4.881,0.340) (4.808,0.325) (4.734,0.310) (4.661,0.296) (4.587,0.282) (4.513,0.268) (4.439,0.255) (4.365,0.242) (4.291,0.230) (4.217,0.218) (4.143,0.207) (4.069,0.196) (3.995,0.185) (3.920,0.175) (3.846,0.165) (3.772,0.156) (3.697,0.147) (3.623,0.138) (3.548,0.130) (3.474,0.122) (3.399,0.114) (3.324,0.107) (3.250,0.100) (3.200,0.095)}
\LC{(1.013,3.088) (1.059,3.147) (1.106,3.206) (1.154,3.264) (1.202,3.321) (1.252,3.377) (1.303,3.432) (1.355,3.486) (1.407,3.540) (1.461,3.592) (1.515,3.644) (1.570,3.695) (1.627,3.744) (1.684,3.793) (1.742,3.840) (1.801,3.887) (1.861,3.932) (1.921,3.976) (1.983,4.019) (2.045,4.061) (2.108,4.102) (2.172,4.141) (2.236,4.179) (2.302,4.216) (2.368,4.251) (2.435,4.285) (2.503,4.317) (2.571,4.348) (2.640,4.378) (2.709,4.406) (2.780,4.432) (2.850,4.457) (2.922,4.480) (2.994,4.501) (3.066,4.521) (3.139,4.539) (3.212,4.555) (3.286,4.569) (3.360,4.581) (3.434,4.591) (3.509,4.600) (3.583,4.606) (3.658,4.610) (3.733,4.612) (3.808,4.612) (3.883,4.610) (3.958,4.605) (4.033,4.598) (4.107,4.589) (4.181,4.577) (4.255,4.563) (4.328,4.547) (4.401,4.528) (4.472,4.506) (4.543,4.482) (4.614,4.455) (4.683,4.426) (4.751,4.394) (4.817,4.360) (4.882,4.323) (4.946,4.283) (5.008,4.241) (5.068,4.196) (5.127,4.149) (5.183,4.100) (5.237,4.048) (5.289,3.993) (5.338,3.937) (5.385,3.878) (5.429,3.818) (5.470,3.755) (5.508,3.690) (5.544,3.624) (5.576,3.557) (5.605,3.488) (5.631,3.417) (5.654,3.346) (5.674,3.273) (5.690,3.200) (5.703,3.126) (5.713,3.052) (5.719,2.977) (5.723,2.902) (5.723,2.827) (5.720,2.752) (5.714,2.678) (5.705,2.603) (5.693,2.529) (5.678,2.456) (5.660,2.383) (5.640,2.311) (5.617,2.239) (5.592,2.169) (5.564,2.099) (5.533,2.030) (5.501,1.963) (5.466,1.896) (5.430,1.831) (5.391,1.767) (5.350,1.704) (5.308,1.642) (5.263,1.581) (5.217,1.522) (5.170,1.464) (5.121,1.408) (5.070,1.352) (5.018,1.298) (4.965,1.246) (4.910,1.194) (4.854,1.144) (4.797,1.095) (4.739,1.048) (4.680,1.002) (4.620,0.957) (4.559,0.913) (4.497,0.871) (4.434,0.830) (4.370,0.790) (4.306,0.752) (4.241,0.715) (4.175,0.679) (4.153,0.667)}
\LC{(-10.000,1.452) (-9.943,1.420) (-9.875,1.389) (-9.807,1.358) (-9.738,1.328) (-9.669,1.298) (-9.600,1.269) (-9.531,1.240) (-9.461,1.211) (-9.392,1.183) (-9.322,1.156) (-9.252,1.129) (-9.182,1.102) (-9.112,1.076) (-9.042,1.050) (-8.971,1.024) (-8.900,0.999) (-8.829,0.974) (-8.759,0.950) (-8.687,0.926) (-8.616,0.903) (-8.545,0.880) (-8.473,0.857) (-8.402,0.835) (-8.330,0.813) (-8.258,0.791) (-8.186,0.770) (-8.114,0.749) (-8.042,0.729) (-7.970,0.709) (-7.897,0.689) (-7.825,0.670) (-7.752,0.651) (-7.680,0.633) (-7.607,0.614) (-7.534,0.597) (-7.461,0.579) (-7.388,0.562) (-7.315,0.545) (-7.242,0.529) (-7.169,0.512) (-7.095,0.497) (-7.022,0.481) (-6.948,0.466) (-6.875,0.451) (-6.801,0.437) (-6.728,0.422) (-6.654,0.408) (-6.580,0.395) (-6.506,0.382) (-6.433,0.369) (-6.359,0.356) (-6.285,0.343) (-6.211,0.331) (-6.137,0.320) (-6.062,0.308) (-5.988,0.297) (-5.914,0.286) (-5.840,0.275) (-5.766,0.265) (-5.691,0.254) (-5.617,0.245) (-5.543,0.235) (-5.468,0.225) (-5.394,0.216) (-5.319,0.207) (-5.245,0.199) (-5.170,0.190) (-5.096,0.182) (-5.021,0.174) (-4.946,0.167) (-4.872,0.159) (-4.797,0.152) (-4.723,0.145) (-4.698,0.143)}
\LC{(1.673,-5.217) (1.718,-5.277) (1.763,-5.337) (1.809,-5.396) (1.856,-5.455) (1.903,-5.514) (1.950,-5.571) (1.999,-5.629) (2.047,-5.686) (2.097,-5.742) (2.147,-5.798) (2.197,-5.854) (2.248,-5.909) (2.300,-5.963) (2.352,-6.017) (2.404,-6.071) (2.458,-6.123) (2.511,-6.176) (2.566,-6.228) (2.620,-6.279) (2.676,-6.329) (2.732,-6.380) (2.788,-6.429) (2.845,-6.478) (2.902,-6.526) (2.960,-6.574) (3.019,-6.621) (3.077,-6.667) (3.137,-6.713) (3.197,-6.758) (3.257,-6.803) (3.318,-6.846) (3.380,-6.889) (3.441,-6.932) (3.504,-6.973) (3.567,-7.014) (3.630,-7.055) (3.694,-7.094) (3.758,-7.133) (3.823,-7.171) (3.888,-7.208) (3.953,-7.244) (4.019,-7.280) (4.086,-7.315) (4.153,-7.349) (4.220,-7.382) (4.287,-7.414) (4.356,-7.446) (4.424,-7.477) (4.493,-7.506) (4.562,-7.535) (4.632,-7.563) (4.702,-7.590) (4.772,-7.616) (4.843,-7.641) (4.914,-7.666) (4.985,-7.689) (5.056,-7.711) (5.128,-7.732) (5.201,-7.753) (5.273,-7.772) (5.346,-7.790) (5.419,-7.807) (5.492,-7.823) (5.566,-7.838) (5.639,-7.852) (5.713,-7.865) (5.787,-7.877) (5.862,-7.887) (5.936,-7.897) (6.011,-7.905) (6.085,-7.912) (6.160,-7.918) (6.235,-7.922) (6.310,-7.926) (6.385,-7.928) (6.460,-7.929) (6.535,-7.928) (6.610,-7.927) (6.685,-7.924) (6.760,-7.919) (6.834,-7.914) (6.909,-7.907) (6.983,-7.898) (7.058,-7.889) (7.132,-7.877) (7.206,-7.865) (7.280,-7.851) (7.353,-7.836) (7.426,-7.819) (7.499,-7.801) (7.571,-7.781) (7.643,-7.760) (7.715,-7.737) (7.786,-7.713) (7.856,-7.687) (7.926,-7.660) (7.996,-7.632) (8.064,-7.602) (8.132,-7.570) (8.200,-7.537) (8.266,-7.502) (8.332,-7.466) (8.397,-7.429) (8.461,-7.389) (8.524,-7.349) (8.586,-7.307) (8.647,-7.263) (8.707,-7.218) (8.766,-7.172) (8.823,-7.124) (8.880,-7.074) (8.935,-7.024) (8.989,-6.971) (9.041,-6.918) (9.093,-6.863) (9.142,-6.807) (9.191,-6.750) (9.238,-6.691) (9.283,-6.631) (9.327,-6.570) (9.369,-6.508) (9.409,-6.445) (9.448,-6.381) (9.485,-6.316) (9.520,-6.249) (9.553,-6.182) (9.585,-6.114) (9.615,-6.045) (9.643,-5.976) (9.669,-5.906) (9.694,-5.835) (9.716,-5.763) (9.736,-5.691) (9.755,-5.618) (9.772,-5.545) (9.787,-5.472) (9.800,-5.398) (9.811,-5.324) (9.820,-5.249) (9.827,-5.175) (9.833,-5.100) (9.836,-5.025) (9.838,-4.950) (9.838,-4.875) (9.836,-4.800) (9.832,-4.725) (9.826,-4.650) (9.819,-4.576) (9.810,-4.501) (9.799,-4.427) (9.787,-4.353) (9.773,-4.279) (9.757,-4.206) (9.739,-4.133) (9.720,-4.060) (9.700,-3.988) (9.678,-3.917) (9.654,-3.845) (9.629,-3.775) (9.603,-3.705) (9.575,-3.635) (9.546,-3.566) (9.515,-3.497) (9.483,-3.429) (9.450,-3.362) (9.416,-3.296) (9.380,-3.230) (9.343,-3.164) (9.305,-3.100) (9.266,-3.036) (9.226,-2.973) (9.184,-2.910) (9.142,-2.848) (9.098,-2.787) (9.054,-2.727) (9.008,-2.667) (8.962,-2.608) (8.914,-2.550) (8.866,-2.493) (8.817,-2.436) (8.767,-2.380) (8.716,-2.325) (8.664,-2.271) (8.612,-2.217) (8.558,-2.165) (8.504,-2.113) (8.450,-2.061) (8.394,-2.011) (8.338,-1.961) (8.281,-1.912) (8.224,-1.864) (8.165,-1.817) (8.107,-1.770) (8.047,-1.724) (7.987,-1.679) (7.927,-1.635) (7.866,-1.591) (7.804,-1.548) (7.742,-1.506) (7.680,-1.465) (7.617,-1.424) (7.553,-1.384) (7.489,-1.345) (7.425,-1.306) (7.360,-1.269) (7.295,-1.232) (7.229,-1.195) (7.163,-1.160) (7.097,-1.125) (7.030,-1.091) (6.963,-1.057) (6.896,-1.024) (6.828,-0.992) (6.760,-0.960) (6.692,-0.929) (6.623,-0.899) (6.554,-0.870) (6.485,-0.841) (6.415,-0.813) (6.346,-0.785) (6.276,-0.758) (6.206,-0.731) (6.135,-0.706) (6.064,-0.681) (5.994,-0.656) (5.923,-0.632) (5.851,-0.609) (5.780,-0.586) (5.708,-0.564) (5.636,-0.542) (5.564,-0.521) (5.492,-0.500) (5.420,-0.480) (5.348,-0.461) (5.275,-0.442) (5.202,-0.423) (5.130,-0.405) (5.057,-0.388) (4.983,-0.371) (4.910,-0.355) (4.837,-0.339) (4.764,-0.323) (4.690,-0.308) (4.617,-0.294) (4.543,-0.280) (4.469,-0.266) (4.395,-0.253) (4.321,-0.240) (4.272,-0.232)}
\LC{(-0.492,3.275) (-0.514,3.347) (-0.537,3.418) (-0.561,3.490) (-0.585,3.561) (-0.609,3.632) (-0.634,3.702) (-0.659,3.773) (-0.685,3.843) (-0.711,3.913) (-0.738,3.984) (-0.766,4.053) (-0.793,4.123) (-0.822,4.192) (-0.850,4.262) (-0.880,4.331) (-0.910,4.400) (-0.940,4.468) (-0.971,4.537) (-1.002,4.605) (-1.034,4.673) (-1.066,4.740) (-1.099,4.808) (-1.132,4.875) (-1.166,4.942) (-1.200,5.009) (-1.234,5.075) (-1.270,5.142) (-1.305,5.208) (-1.341,5.273) (-1.378,5.339) (-1.415,5.404) (-1.452,5.469) (-1.491,5.534) (-1.529,5.598) (-1.568,5.662) (-1.608,5.726) (-1.647,5.789) (-1.688,5.852) (-1.729,5.915) (-1.770,5.978) (-1.812,6.040) (-1.854,6.102) (-1.897,6.164) (-1.940,6.225) (-1.984,6.286) (-2.028,6.346) (-2.073,6.407) (-2.118,6.466) (-2.164,6.526) (-2.210,6.585) (-2.256,6.644) (-2.303,6.702) (-2.351,6.761) (-2.399,6.818) (-2.447,6.876) (-2.496,6.932) (-2.545,6.989) (-2.595,7.045) (-2.645,7.101) (-2.696,7.156) (-2.747,7.211) (-2.799,7.266) (-2.851,7.320) (-2.903,7.373) (-2.956,7.426) (-3.009,7.479) (-3.063,7.531) (-3.117,7.583) (-3.172,7.635) (-3.227,7.686) (-3.282,7.736) (-3.338,7.786) (-3.395,7.835) (-3.452,7.884) (-3.509,7.933) (-3.566,7.981) (-3.624,8.028) (-3.683,8.075) (-3.742,8.122) (-3.801,8.168) (-3.861,8.213) (-3.921,8.258) (-3.982,8.302) (-4.042,8.346) (-4.104,8.389) (-4.166,8.432) (-4.228,8.474) (-4.290,8.515) (-4.353,8.556) (-4.416,8.596) (-4.480,8.636) (-4.544,8.675) (-4.608,8.713) (-4.673,8.751) (-4.738,8.788) (-4.804,8.825) (-4.870,8.861) (-4.936,8.896) (-5.003,8.930) (-5.069,8.964) (-5.137,8.997) (-5.204,9.030) (-5.272,9.062) (-5.340,9.093) (-5.409,9.123) (-5.478,9.153) (-5.547,9.182) (-5.617,9.210) (-5.686,9.237) (-5.757,9.264) (-5.827,9.290) (-5.898,9.315) (-5.968,9.340) (-6.040,9.363) (-6.111,9.386) (-6.183,9.408) (-6.255,9.429) (-6.327,9.450) (-6.399,9.469) (-6.472,9.488) (-6.545,9.505) (-6.618,9.522) (-6.691,9.538) (-6.765,9.554) (-6.838,9.568) (-6.912,9.581) (-6.986,9.594) (-7.060,9.605) (-7.134,9.616) (-7.209,9.625) (-7.283,9.634) (-7.358,9.642) (-7.433,9.649) (-7.507,9.654) (-7.582,9.659) (-7.657,9.663) (-7.732,9.666) (-7.807,9.667) (-7.882,9.668) (-7.957,9.668) (-8.032,9.666) (-8.107,9.664) (-8.182,9.660) (-8.257,9.656) (-8.331,9.650) (-8.406,9.643) (-8.481,9.635) (-8.555,9.626) (-8.630,9.616) (-8.704,9.605) (-8.778,9.592) (-8.851,9.578) (-8.925,9.564) (-8.998,9.548) (-9.071,9.530) (-9.144,9.512) (-9.216,9.493) (-9.288,9.472) (-9.360,9.450) (-9.431,9.427) (-9.502,9.402) (-9.573,9.377) (-9.643,9.350) (-9.712,9.322) (-9.781,9.292) (-9.850,9.262) (-9.918,9.230) (-9.985,9.197) (-10.000,9.186)}
\LC{(-10.000,-1.331) (-9.939,-1.302) (-9.870,-1.273) (-9.801,-1.245) (-9.731,-1.217) (-9.661,-1.190) (-9.591,-1.163) (-9.521,-1.136) (-9.451,-1.110) (-9.380,-1.084) (-9.310,-1.059) (-9.239,-1.034) (-9.168,-1.009) (-9.097,-0.985) (-9.026,-0.961) (-8.955,-0.938) (-8.883,-0.915) (-8.812,-0.892) (-8.740,-0.870) (-8.669,-0.848) (-8.597,-0.826) (-8.525,-0.805) (-8.453,-0.784) (-8.381,-0.764) (-8.308,-0.744) (-8.236,-0.724) (-8.164,-0.704) (-8.091,-0.685) (-8.019,-0.667) (-7.946,-0.648) (-7.873,-0.630) (-7.800,-0.612) (-7.727,-0.595) (-7.654,-0.578) (-7.581,-0.561) (-7.508,-0.545) (-7.435,-0.529) (-7.361,-0.513) (-7.288,-0.498) (-7.214,-0.483) (-7.141,-0.468) (-7.067,-0.453) (-6.994,-0.439) (-6.920,-0.425) (-6.846,-0.411) (-6.772,-0.398) (-6.699,-0.385) (-6.625,-0.372) (-6.551,-0.360) (-6.477,-0.348) (-6.403,-0.336) (-6.328,-0.324)}
\LC{(1.709,4.492) (1.761,4.547) (1.813,4.600) (1.866,4.653) (1.920,4.706) (1.975,4.757) (2.030,4.808) (2.086,4.858) (2.142,4.907) (2.199,4.956) (2.257,5.003) (2.316,5.050) (2.375,5.096) (2.435,5.141) (2.496,5.186) (2.557,5.229) (2.619,5.271) (2.681,5.313) (2.744,5.354) (2.808,5.393) (2.872,5.432) (2.937,5.469) (3.003,5.506) (3.069,5.542) (3.135,5.576) (3.202,5.609) (3.270,5.642) (3.338,5.673) (3.407,5.703) (3.476,5.732) (3.546,5.760) (3.616,5.786) (3.687,5.811) (3.758,5.835) (3.829,5.858) (3.901,5.879) (3.973,5.899) (4.046,5.918) (4.119,5.935) (4.192,5.951) (4.266,5.966) (4.340,5.979) (4.414,5.990) (4.488,6.000) (4.563,6.009) (4.637,6.016) (4.712,6.021) (4.787,6.025) (4.862,6.027) (4.937,6.028) (5.012,6.027) (5.087,6.024) (5.162,6.019) (5.236,6.013) (5.311,6.004) (5.385,5.995) (5.459,5.983) (5.533,5.969) (5.607,5.953) (5.679,5.936) (5.752,5.917) (5.824,5.895) (5.895,5.872) (5.966,5.847) (6.036,5.820) (6.105,5.791) (6.173,5.760) (6.240,5.727) (6.307,5.692) (6.372,5.655) (6.436,5.616) (6.499,5.575) (6.560,5.532) (6.620,5.487) (6.679,5.440) (6.736,5.391) (6.792,5.341) (6.845,5.289) (6.897,5.235) (6.947,5.179) (6.996,5.121) (7.042,5.062) (7.086,5.002) (7.128,4.940) (7.168,4.876) (7.205,4.811) (7.241,4.745) (7.274,4.678) (7.304,4.609) (7.333,4.540) (7.358,4.469) (7.382,4.398) (7.402,4.326) (7.421,4.253) (7.437,4.180) (7.450,4.106) (7.461,4.032) (7.469,3.957) (7.475,3.883) (7.478,3.808) (7.479,3.733) (7.478,3.658) (7.474,3.583) (7.468,3.508) (7.460,3.433) (7.449,3.359) (7.436,3.285) (7.421,3.212) (7.404,3.139) (7.384,3.066) (7.363,2.995) (7.340,2.923) (7.315,2.853) (7.287,2.783) (7.258,2.714) (7.228,2.645) (7.195,2.577) (7.161,2.511) (7.125,2.445) (7.088,2.380) (7.049,2.316) (7.009,2.252) (6.967,2.190) (6.924,2.129) (6.880,2.068) (6.834,2.009) (6.787,1.950) (6.739,1.893) (6.689,1.837) (6.639,1.781) (6.587,1.727) (6.534,1.673) (6.481,1.621) (6.426,1.569) (6.371,1.519) (6.314,1.470) (6.257,1.421) (6.199,1.374) (6.140,1.328) (6.080,1.282) (6.020,1.238) (5.958,1.194) (5.897,1.152) (5.834,1.111) (5.771,1.070) (5.707,1.031) (5.686,1.018)}
\LC{(0.123,-3.792) (0.128,-3.866) (0.133,-3.941) (0.138,-4.016) (0.144,-4.091) (0.149,-4.166) (0.154,-4.241) (0.160,-4.315) (0.165,-4.390) (0.171,-4.465) (0.177,-4.540) (0.183,-4.614) (0.189,-4.689) (0.195,-4.764) (0.201,-4.839) (0.207,-4.913) (0.214,-4.988) (0.220,-5.063) (0.227,-5.138) (0.233,-5.212) (0.240,-5.287) (0.247,-5.362) (0.254,-5.436) (0.261,-5.511) (0.268,-5.586) (0.275,-5.660) (0.283,-5.735) (0.290,-5.810) (0.298,-5.884) (0.305,-5.959) (0.313,-6.033) (0.321,-6.108) (0.329,-6.183) (0.337,-6.257) (0.345,-6.332) (0.353,-6.406) (0.361,-6.481) (0.370,-6.555) (0.378,-6.630) (0.387,-6.704) (0.395,-6.779) (0.404,-6.853) (0.413,-6.928) (0.422,-7.002) (0.431,-7.077) (0.440,-7.151) (0.449,-7.226) (0.459,-7.300) (0.468,-7.374) (0.478,-7.449) (0.487,-7.523) (0.497,-7.598) (0.507,-7.672) (0.517,-7.746) (0.527,-7.821) (0.537,-7.895) (0.547,-7.969) (0.558,-8.044) (0.568,-8.118) (0.579,-8.192) (0.589,-8.266) (0.600,-8.341) (0.611,-8.415) (0.621,-8.489) (0.632,-8.563) (0.644,-8.637) (0.655,-8.712) (0.666,-8.786) (0.677,-8.860) (0.689,-8.934) (0.700,-9.008) (0.712,-9.082) (0.724,-9.156) (0.736,-9.230) (0.748,-9.304) (0.760,-9.378) (0.772,-9.452) (0.784,-9.526) (0.796,-9.600) (0.809,-9.674) (0.821,-9.748) (0.834,-9.822) (0.847,-9.896) (0.859,-9.970) (0.868,-10.000)}
\LC{(-0.270,-4.625) (-0.278,-4.699) (-0.287,-4.774) (-0.296,-4.848) (-0.306,-4.923) (-0.315,-4.997) (-0.324,-5.071) (-0.334,-5.146) (-0.344,-5.220) (-0.354,-5.295) (-0.364,-5.369) (-0.374,-5.443) (-0.384,-5.517) (-0.395,-5.592) (-0.405,-5.666) (-0.416,-5.740) (-0.427,-5.814) (-0.438,-5.889) (-0.449,-5.963) (-0.461,-6.037) (-0.472,-6.111) (-0.484,-6.185) (-0.495,-6.259) (-0.507,-6.333) (-0.519,-6.407) (-0.532,-6.481) (-0.544,-6.555) (-0.556,-6.629) (-0.569,-6.703) (-0.582,-6.777) (-0.594,-6.851) (-0.608,-6.925) (-0.621,-6.999) (-0.634,-7.072) (-0.647,-7.146) (-0.661,-7.220) (-0.675,-7.294) (-0.689,-7.367) (-0.703,-7.441) (-0.717,-7.515) (-0.731,-7.588) (-0.745,-7.662) (-0.760,-7.736) (-0.775,-7.809) (-0.789,-7.883) (-0.804,-7.956) (-0.820,-8.030) (-0.835,-8.103) (-0.850,-8.176) (-0.866,-8.250) (-0.881,-8.323) (-0.897,-8.396) (-0.913,-8.470) (-0.929,-8.543) (-0.946,-8.616) (-0.962,-8.689) (-0.979,-8.763) (-0.995,-8.836) (-1.012,-8.909) (-1.029,-8.982) (-1.046,-9.055) (-1.063,-9.128) (-1.081,-9.201) (-1.098,-9.274) (-1.116,-9.347) (-1.134,-9.419) (-1.151,-9.492) (-1.170,-9.565) (-1.188,-9.638) (-1.206,-9.711) (-1.225,-9.783) (-1.243,-9.856) (-1.262,-9.929) (-1.281,-10.000)}
\LC{(1.386,-3.828) (1.435,-3.884) (1.486,-3.940) (1.537,-3.994) (1.589,-4.048) (1.642,-4.101) (1.696,-4.153) (1.751,-4.205) (1.806,-4.256) (1.862,-4.305) (1.919,-4.354) (1.977,-4.402) (2.035,-4.449) (2.094,-4.495) (2.154,-4.540) (2.215,-4.584) (2.276,-4.627) (2.339,-4.669) (2.401,-4.710) (2.465,-4.750) (2.529,-4.789) (2.594,-4.827) (2.659,-4.863) (2.726,-4.899) (2.792,-4.933) (2.860,-4.965) (2.928,-4.997) (2.997,-5.027) (3.066,-5.056) (3.135,-5.084) (3.206,-5.110) (3.276,-5.135) (3.348,-5.158) (3.419,-5.180) (3.492,-5.201) (3.564,-5.220) (3.637,-5.237) (3.710,-5.253) (3.784,-5.267) (3.858,-5.279) (3.932,-5.290) (4.007,-5.299) (4.081,-5.306) (4.156,-5.312) (4.231,-5.316) (4.306,-5.318) (4.381,-5.317) (4.456,-5.316) (4.531,-5.312) (4.606,-5.306) (4.680,-5.298) (4.755,-5.288) (4.829,-5.276) (4.902,-5.262) (4.975,-5.245) (5.048,-5.227) (5.120,-5.206) (5.192,-5.184) (5.262,-5.158) (5.332,-5.131) (5.401,-5.102) (5.469,-5.070) (5.536,-5.036) (5.602,-5.000) (5.666,-4.961) (5.729,-4.921) (5.791,-4.878) (5.851,-4.833) (5.909,-4.785) (5.965,-4.736) (6.020,-4.685) (6.072,-4.631) (6.123,-4.576) (6.171,-4.518) (6.217,-4.459) (6.261,-4.398) (6.302,-4.336) (6.341,-4.272) (6.378,-4.206) (6.411,-4.139) (6.442,-4.071) (6.470,-4.001) (6.496,-3.931) (6.518,-3.859) (6.538,-3.787) (6.555,-3.714) (6.569,-3.640) (6.581,-3.566) (6.589,-3.491) (6.595,-3.417) (6.598,-3.342) (6.598,-3.267) (6.596,-3.192) (6.590,-3.117) (6.582,-3.042) (6.572,-2.968) (6.559,-2.894) (6.544,-2.821) (6.526,-2.748) (6.506,-2.676) (6.484,-2.604) (6.459,-2.533) (6.432,-2.463) (6.403,-2.394) (6.373,-2.326) (6.340,-2.258) (6.305,-2.192) (6.269,-2.126) (6.230,-2.062) (6.191,-1.998) (6.149,-1.936) (6.106,-1.874) (6.061,-1.814) (6.015,-1.755) (5.968,-1.697) (5.919,-1.640) (5.869,-1.584) (5.817,-1.529) (5.765,-1.476) (5.711,-1.424) (5.656,-1.372) (5.600,-1.322) (5.543,-1.274) (5.486,-1.226) (5.427,-1.179) (5.367,-1.134) (5.307,-1.089) (5.245,-1.046) (5.183,-1.004) (5.121,-0.963) (5.057,-0.923) (4.993,-0.884) (4.928,-0.846) (4.863,-0.809)}
\LC{(-2.350,6.023) (-2.402,6.076) (-2.456,6.129) (-2.509,6.182) (-2.563,6.234) (-2.618,6.285) (-2.673,6.336) (-2.729,6.386) (-2.785,6.435) (-2.842,6.484) (-2.900,6.533) (-2.957,6.580) (-3.016,6.628) (-3.075,6.674) (-3.134,6.720) (-3.194,6.765) (-3.254,6.810) (-3.315,6.853) (-3.376,6.897) (-3.438,6.939) (-3.500,6.981) (-3.563,7.022) (-3.626,7.062) (-3.690,7.102) (-3.754,7.141) (-3.819,7.179) (-3.884,7.216) (-3.949,7.253) (-4.015,7.289) (-4.081,7.324) (-4.148,7.358) (-4.215,7.391) (-4.283,7.424) (-4.351,7.455) (-4.419,7.486) (-4.488,7.516) (-4.557,7.545) (-4.627,7.573) (-4.697,7.600) (-4.767,7.627) (-4.838,7.652) (-4.908,7.677) (-4.980,7.700) (-5.051,7.723) (-5.123,7.744) (-5.195,7.765) (-5.268,7.784) (-5.340,7.802) (-5.413,7.820) (-5.486,7.836) (-5.560,7.851) (-5.634,7.866) (-5.707,7.879) (-5.781,7.891) (-5.856,7.901) (-5.930,7.911) (-6.005,7.920) (-6.079,7.927) (-6.154,7.933) (-6.229,7.938) (-6.304,7.942) (-6.379,7.944) (-6.454,7.945) (-6.529,7.945) (-6.604,7.944) (-6.679,7.941) (-6.753,7.937) (-6.828,7.932) (-6.903,7.925) (-6.978,7.917) (-7.052,7.908) (-7.126,7.897) (-7.200,7.885) (-7.274,7.872) (-7.348,7.857) (-7.421,7.840) (-7.494,7.823) (-7.566,7.803) (-7.638,7.783) (-7.710,7.760) (-7.781,7.737) (-7.852,7.712) (-7.922,7.685) (-7.991,7.657) (-8.060,7.627) (-8.128,7.596) (-8.196,7.563) (-8.263,7.529) (-8.329,7.494) (-8.394,7.456) (-8.458,7.418) (-8.521,7.378) (-8.584,7.336) (-8.645,7.293) (-8.705,7.248) (-8.765,7.202) (-8.823,7.155) (-8.880,7.106) (-8.935,7.056) (-8.990,7.004) (-9.043,6.951) (-9.094,6.897) (-9.145,6.841) (-9.193,6.784) (-9.241,6.726) (-9.287,6.667) (-9.331,6.606) (-9.374,6.544) (-9.415,6.481) (-9.454,6.418) (-9.491,6.353) (-9.527,6.287) (-9.561,6.220) (-9.594,6.152) (-9.624,6.084) (-9.653,6.015) (-9.680,5.945) (-9.705,5.874) (-9.728,5.803) (-9.749,5.731) (-9.769,5.658) (-9.786,5.585) (-9.802,5.512) (-9.816,5.438) (-9.827,5.364) (-9.837,5.290) (-9.845,5.215) (-9.851,5.140) (-9.856,5.066) (-9.858,4.991) (-9.859,4.916) (-9.857,4.841) (-9.854,4.766) (-9.850,4.691) (-9.843,4.616) (-9.835,4.542) (-9.825,4.467) (-9.813,4.393) (-9.799,4.319) (-9.784,4.246) (-9.768,4.173) (-9.749,4.100) (-9.729,4.028) (-9.708,3.956) (-9.685,3.884) (-9.661,3.814) (-9.635,3.743) (-9.608,3.673) (-9.579,3.604) (-9.549,3.535) (-9.518,3.467) (-9.485,3.399) (-9.451,3.333) (-9.416,3.266) (-9.380,3.201) (-9.342,3.136) (-9.304,3.072) (-9.264,3.008) (-9.223,2.945) (-9.181,2.883) (-9.138,2.822) (-9.094,2.761) (-9.049,2.701) (-9.003,2.642) (-8.956,2.583) (-8.908,2.525) (-8.859,2.468) (-8.810,2.412) (-8.759,2.357) (-8.708,2.302) (-8.656,2.248) (-8.603,2.195) (-8.549,2.143) (-8.495,2.091) (-8.439,2.040) (-8.384,1.990) (-8.327,1.941) (-8.270,1.892) (-8.212,1.844) (-8.154,1.797) (-8.095,1.751) (-8.035,1.706) (-7.975,1.661) (-7.914,1.617) (-7.853,1.574) (-7.791,1.531) (-7.729,1.489) (-7.666,1.448) (-7.603,1.408) (-7.539,1.368) (-7.475,1.330) (-7.410,1.291) (-7.345,1.254) (-7.280,1.217) (-7.214,1.181) (-7.148,1.146) (-7.082,1.111) (-7.015,1.077) (-6.947,1.044) (-6.880,1.012) (-6.812,0.980) (-6.744,0.949) (-6.675,0.918) (-6.607,0.888) (-6.537,0.859) (-6.491,0.840)}
\LC{(10.000,-1.736) (9.939,-1.698) (9.874,-1.662) (9.808,-1.625) (9.742,-1.590) (9.676,-1.554) (9.609,-1.519) (9.543,-1.485) (9.476,-1.451) (9.408,-1.418) (9.341,-1.385) (9.273,-1.353) (9.205,-1.321) (9.137,-1.290) (9.069,-1.259) (9.000,-1.229) (8.931,-1.199) (8.862,-1.170) (8.793,-1.141) (8.724,-1.113) (8.654,-1.085) (8.584,-1.057) (8.514,-1.030) (8.444,-1.004) (8.374,-0.978) (8.303,-0.952) (8.233,-0.927) (8.162,-0.902) (8.091,-0.878) (8.020,-0.854) (7.949,-0.831) (7.877,-0.808) (7.806,-0.785) (7.734,-0.763) (7.662,-0.741) (7.590,-0.720) (7.518,-0.699) (7.446,-0.679) (7.374,-0.658) (7.302,-0.639) (7.229,-0.619) (7.156,-0.600) (7.084,-0.582) (7.011,-0.564) (6.938,-0.546)}
\LC{(4.710,-8.828) (4.775,-8.865) (4.841,-8.902) (4.907,-8.938) (4.973,-8.973) (5.039,-9.008) (5.106,-9.042) (5.173,-9.075) (5.241,-9.107) (5.309,-9.139) (5.377,-9.171) (5.445,-9.201) (5.514,-9.231) (5.583,-9.260) (5.653,-9.288) (5.723,-9.316) (5.793,-9.343) (5.863,-9.369) (5.933,-9.394) (6.004,-9.419) (6.075,-9.443) (6.147,-9.466) (6.218,-9.488) (6.290,-9.509) (6.362,-9.530) (6.435,-9.550) (6.507,-9.569) (6.580,-9.587) (6.653,-9.604) (6.726,-9.620) (6.800,-9.636) (6.873,-9.650) (6.947,-9.664) (7.021,-9.677) (7.095,-9.689) (7.169,-9.700) (7.243,-9.710) (7.318,-9.719) (7.392,-9.727) (7.467,-9.734) (7.542,-9.741) (7.617,-9.746) (7.692,-9.750) (7.766,-9.753) (7.841,-9.755) (7.916,-9.757) (7.991,-9.757) (8.066,-9.756) (8.141,-9.754) (8.216,-9.751) (8.291,-9.747) (8.366,-9.742) (8.441,-9.736) (8.515,-9.728) (8.590,-9.720) (8.664,-9.710) (8.739,-9.700) (8.813,-9.688) (8.887,-9.675) (8.960,-9.661) (9.034,-9.646) (9.107,-9.629) (9.180,-9.612) (9.252,-9.593) (9.325,-9.573) (9.397,-9.552) (9.468,-9.529) (9.539,-9.506) (9.610,-9.481) (9.681,-9.455) (9.750,-9.428) (9.820,-9.400) (9.889,-9.370) (9.957,-9.339) (10.000,-9.318)}
\LC{(2.598,-5.621) (2.657,-5.667) (2.718,-5.712) (2.778,-5.756) (2.840,-5.799) (2.901,-5.842) (2.964,-5.883) (3.027,-5.924) (3.090,-5.964) (3.154,-6.003) (3.219,-6.041) (3.284,-6.078) (3.350,-6.114) (3.416,-6.149) (3.483,-6.184) (3.550,-6.217) (3.618,-6.249) (3.686,-6.281) (3.754,-6.311) (3.823,-6.340) (3.893,-6.369) (3.963,-6.396) (4.033,-6.422) (4.104,-6.447) (4.175,-6.470) (4.247,-6.493) (4.318,-6.514) (4.391,-6.534) (4.463,-6.553) (4.536,-6.571) (4.609,-6.588) (4.683,-6.603) (4.756,-6.617) (4.830,-6.629) (4.904,-6.641) (4.979,-6.651) (5.053,-6.659) (5.128,-6.666) (5.203,-6.672) (5.278,-6.676) (5.353,-6.679) (5.428,-6.680) (5.503,-6.680) (5.578,-6.678) (5.652,-6.675) (5.727,-6.670) (5.802,-6.664) (5.877,-6.656) (5.951,-6.646) (6.025,-6.635) (6.099,-6.622) (6.172,-6.607) (6.246,-6.591) (6.318,-6.572) (6.391,-6.553) (6.463,-6.531) (6.534,-6.508) (6.605,-6.483) (6.675,-6.456) (6.744,-6.427) (6.812,-6.397) (6.880,-6.364) (6.947,-6.330) (7.013,-6.295) (7.078,-6.257) (7.142,-6.218) (7.204,-6.177) (7.266,-6.134) (7.326,-6.089) (7.385,-6.043) (7.443,-5.995) (7.499,-5.945) (7.554,-5.894) (7.607,-5.841) (7.658,-5.786) (7.708,-5.730) (7.756,-5.672) (7.802,-5.613) (7.846,-5.553) (7.889,-5.491) (7.929,-5.428) (7.968,-5.363) (8.004,-5.298) (8.038,-5.231) (8.071,-5.163) (8.101,-5.095) (8.129,-5.025) (8.154,-4.955) (8.178,-4.883) (8.199,-4.812) (8.218,-4.739) (8.235,-4.666) (8.249,-4.592) (8.261,-4.518) (8.271,-4.444) (8.279,-4.369) (8.285,-4.295) (8.288,-4.220) (8.289,-4.145) (8.288,-4.070) (8.285,-3.995) (8.279,-3.920) (8.272,-3.845) (8.262,-3.771) (8.251,-3.697) (8.237,-3.623) (8.222,-3.550) (8.205,-3.477) (8.185,-3.404) (8.164,-3.332) (8.141,-3.261) (8.117,-3.190) (8.090,-3.120) (8.062,-3.050) (8.033,-2.981) (8.002,-2.913) (7.969,-2.846) (7.935,-2.779) (7.899,-2.713) (7.862,-2.648) (7.823,-2.583) (7.783,-2.520) (7.742,-2.457) (7.700,-2.395) (7.656,-2.335) (7.611,-2.274) (7.565,-2.215) (7.518,-2.157) (7.469,-2.100) (7.420,-2.043) (7.370,-1.987) (7.319,-1.933) (7.266,-1.879) (7.213,-1.826) (7.159,-1.774) (7.104,-1.723) (7.048,-1.673) (6.992,-1.624) (6.934,-1.575) (6.876,-1.528) (6.857,-1.512)}
\LC{(-0.778,5.179) (-0.800,5.251) (-0.823,5.322) (-0.846,5.394) (-0.869,5.465) (-0.893,5.536) (-0.916,5.608) (-0.941,5.679) (-0.965,5.749) (-0.990,5.820) (-1.015,5.891) (-1.041,5.961) (-1.066,6.032) (-1.092,6.102) (-1.119,6.172) (-1.145,6.242) (-1.172,6.312) (-1.200,6.382) (-1.227,6.452) (-1.255,6.522) (-1.284,6.591) (-1.312,6.660) (-1.341,6.730) (-1.370,6.799) (-1.400,6.868) (-1.429,6.937) (-1.459,7.005) (-1.490,7.074) (-1.521,7.142) (-1.552,7.211) (-1.583,7.279) (-1.615,7.347) (-1.647,7.415) (-1.679,7.482) (-1.711,7.550) (-1.744,7.617) (-1.777,7.684) (-1.811,7.752) (-1.845,7.818) (-1.879,7.885) (-1.913,7.952) (-1.948,8.018) (-1.983,8.085) (-2.018,8.151) (-2.054,8.217) (-2.090,8.283) (-2.126,8.348) (-2.163,8.414) (-2.200,8.479) (-2.237,8.544) (-2.274,8.609) (-2.312,8.674) (-2.350,8.739) (-2.388,8.803) (-2.427,8.867) (-2.466,8.931) (-2.505,8.995) (-2.545,9.059) (-2.585,9.123) (-2.625,9.186) (-2.666,9.249) (-2.706,9.312) (-2.747,9.375) (-2.789,9.437) (-2.830,9.500) (-2.872,9.562) (-2.915,9.624) (-2.957,9.685) (-3.000,9.747) (-3.043,9.808) (-3.087,9.869) (-3.131,9.930) (-3.175,9.991) (-3.190,10.000)}
\LC{(-3.035,-7.115) (-3.091,-7.164) (-3.147,-7.213) (-3.204,-7.262) (-3.262,-7.310) (-3.320,-7.358) (-3.379,-7.405) (-3.437,-7.451) (-3.497,-7.497) (-3.557,-7.542) (-3.617,-7.587) (-3.678,-7.631) (-3.739,-7.674) (-3.800,-7.717) (-3.862,-7.759) (-3.925,-7.801) (-3.988,-7.841) (-4.051,-7.882) (-4.115,-7.921) (-4.179,-7.960) (-4.243,-7.999) (-4.308,-8.036) (-4.374,-8.073) (-4.439,-8.109) (-4.505,-8.145) (-4.572,-8.179) (-4.639,-8.213) (-4.706,-8.247) (-4.773,-8.279) (-4.841,-8.311) (-4.910,-8.342) (-4.978,-8.372) (-5.047,-8.401) (-5.117,-8.430) (-5.186,-8.458) (-5.256,-8.485) (-5.327,-8.511) (-5.397,-8.536) (-5.468,-8.561) (-5.539,-8.584) (-5.611,-8.607) (-5.683,-8.629) (-5.755,-8.650) (-5.827,-8.670) (-5.899,-8.689) (-5.972,-8.707) (-6.045,-8.724) (-6.118,-8.741) (-6.192,-8.756) (-6.265,-8.770) (-6.339,-8.784) (-6.413,-8.796) (-6.487,-8.808) (-6.562,-8.818) (-6.636,-8.827) (-6.710,-8.835) (-6.785,-8.843) (-6.860,-8.849) (-6.935,-8.854) (-7.010,-8.858) (-7.085,-8.861) (-7.160,-8.862) (-7.235,-8.863) (-7.310,-8.862) (-7.384,-8.861) (-7.459,-8.858) (-7.534,-8.854) (-7.609,-8.849) (-7.684,-8.842) (-7.758,-8.835) (-7.833,-8.826) (-7.907,-8.816) (-7.981,-8.804) (-8.055,-8.792) (-8.129,-8.778) (-8.203,-8.763) (-8.276,-8.746) (-8.349,-8.729) (-8.421,-8.710) (-8.493,-8.689) (-8.565,-8.668) (-8.637,-8.645) (-8.708,-8.621) (-8.778,-8.595) (-8.848,-8.568) (-8.917,-8.540) (-8.986,-8.510) (-9.055,-8.479) (-9.122,-8.447) (-9.189,-8.413) (-9.256,-8.378) (-9.321,-8.342) (-9.386,-8.304) (-9.450,-8.265) (-9.513,-8.225) (-9.576,-8.183) (-9.637,-8.140) (-9.698,-8.095) (-9.757,-8.050) (-9.816,-8.003) (-9.873,-7.955) (-9.929,-7.905) (-9.985,-7.854) (-10.000,-7.837)}
\LC{(-0.793,-5.153) (-0.816,-5.224) (-0.839,-5.296) (-0.863,-5.367) (-0.886,-5.438) (-0.910,-5.509) (-0.935,-5.580) (-0.960,-5.651) (-0.985,-5.721) (-1.010,-5.792) (-1.036,-5.862) (-1.062,-5.933) (-1.088,-6.003) (-1.115,-6.073) (-1.142,-6.143) (-1.169,-6.213) (-1.197,-6.283) (-1.225,-6.352) (-1.253,-6.422) (-1.282,-6.491) (-1.311,-6.560) (-1.340,-6.629) (-1.370,-6.698) (-1.400,-6.767) (-1.430,-6.835) (-1.460,-6.904) (-1.491,-6.972) (-1.522,-7.041) (-1.554,-7.109) (-1.586,-7.177) (-1.618,-7.244) (-1.650,-7.312) (-1.683,-7.379) (-1.716,-7.447) (-1.749,-7.514) (-1.783,-7.581) (-1.817,-7.648) (-1.851,-7.714) (-1.886,-7.781) (-1.921,-7.847) (-1.956,-7.914) (-1.992,-7.980) (-2.028,-8.045) (-2.064,-8.111) (-2.101,-8.177) (-2.137,-8.242) (-2.175,-8.307) (-2.212,-8.372) (-2.250,-8.437) (-2.288,-8.501) (-2.326,-8.566) (-2.365,-8.630) (-2.404,-8.694) (-2.443,-8.758) (-2.483,-8.822) (-2.523,-8.885) (-2.563,-8.948) (-2.604,-9.011) (-2.645,-9.074) (-2.686,-9.137) (-2.727,-9.199) (-2.769,-9.262) (-2.811,-9.324) (-2.854,-9.386) (-2.897,-9.447) (-2.940,-9.509) (-2.983,-9.570) (-3.027,-9.631) (-3.071,-9.692) (-3.115,-9.752) (-3.159,-9.812) (-3.204,-9.873) (-3.249,-9.932) (-3.295,-9.992) (-3.310,-10.000)}
\LC{(-1.651,4.550) (-1.701,4.606) (-1.751,4.662) (-1.802,4.717) (-1.854,4.771) (-1.906,4.824) (-1.960,4.877) (-2.013,4.930) (-2.068,4.981) (-2.123,5.032) (-2.179,5.082) (-2.235,5.131) (-2.293,5.180) (-2.350,5.228) (-2.409,5.275) (-2.468,5.321) (-2.527,5.366) (-2.588,5.411) (-2.649,5.455) (-2.710,5.498) (-2.772,5.540) (-2.835,5.581) (-2.898,5.621) (-2.962,5.660) (-3.027,5.698) (-3.092,5.736) (-3.157,5.772) (-3.224,5.808) (-3.290,5.842) (-3.357,5.875) (-3.425,5.907) (-3.493,5.939) (-3.562,5.969) (-3.631,5.998) (-3.701,6.026) (-3.771,6.052) (-3.842,6.078) (-3.912,6.102) (-3.984,6.125) (-4.056,6.147) (-4.128,6.168) (-4.200,6.187) (-4.273,6.205) (-4.346,6.222) (-4.419,6.237) (-4.493,6.251) (-4.567,6.263) (-4.641,6.275) (-4.716,6.284) (-4.790,6.293) (-4.865,6.299) (-4.940,6.305) (-5.015,6.308) (-5.090,6.310) (-5.165,6.311) (-5.240,6.310) (-5.315,6.307) (-5.389,6.303) (-5.464,6.297) (-5.539,6.289) (-5.613,6.280) (-5.687,6.268) (-5.761,6.255) (-5.835,6.241) (-5.908,6.224) (-5.981,6.206) (-6.053,6.186) (-6.125,6.164) (-6.196,6.140) (-6.266,6.114) (-6.336,6.086) (-6.405,6.057) (-6.473,6.026) (-6.540,5.992) (-6.606,5.957) (-6.671,5.920) (-6.736,5.881) (-6.798,5.840) (-6.860,5.797) (-6.920,5.753) (-6.979,5.706) (-7.037,5.658) (-7.093,5.608) (-7.147,5.557) (-7.200,5.503) (-7.251,5.448) (-7.300,5.392) (-7.347,5.333) (-7.392,5.273) (-7.436,5.212) (-7.477,5.150) (-7.516,5.086) (-7.553,5.020) (-7.588,4.954) (-7.620,4.886) (-7.651,4.818) (-7.679,4.748) (-7.704,4.678) (-7.728,4.607) (-7.749,4.535) (-7.767,4.462) (-7.783,4.389) (-7.797,4.315) (-7.809,4.241) (-7.818,4.166) (-7.824,4.092) (-7.829,4.017) (-7.830,3.942) (-7.830,3.867) (-7.828,3.792) (-7.823,3.717) (-7.816,3.642) (-7.806,3.568) (-7.795,3.494) (-7.781,3.420) (-7.766,3.347) (-7.748,3.274) (-7.729,3.201) (-7.707,3.130) (-7.684,3.058) (-7.658,2.988) (-7.631,2.918) (-7.603,2.848) (-7.572,2.780) (-7.540,2.712) (-7.506,2.645) (-7.471,2.579) (-7.434,2.514) (-7.396,2.449) (-7.356,2.386) (-7.315,2.323) (-7.272,2.261) (-7.229,2.200) (-7.184,2.140) (-7.137,2.081) (-7.090,2.023) (-7.041,1.966) (-6.992,1.910) (-6.941,1.854) (-6.889,1.800) (-6.836,1.747) (-6.783,1.695) (-6.728,1.643) (-6.673,1.593) (-6.616,1.543) (-6.559,1.495) (-6.501,1.447) (-6.442,1.401) (-6.383,1.355) (-6.322,1.310) (-6.261,1.267) (-6.200,1.224) (-6.138,1.182) (-6.075,1.141) (-6.011,1.101) (-5.947,1.062) (-5.883,1.024) (-5.818,0.987) (-5.752,0.950) (-5.686,0.915) (-5.620,0.880) (-5.553,0.846) (-5.485,0.813) (-5.417,0.781)}
\LC{(10.000,-0.013) (9.943,-0.013) (9.868,-0.013) (9.793,-0.012) (9.718,-0.012) (9.643,-0.012) (9.568,-0.012) (9.493,-0.011) (9.418,-0.011) (9.343,-0.011) (9.268,-0.011) (9.193,-0.010) (9.118,-0.010) (9.043,-0.010) (8.968,-0.010) (8.893,-0.009) (8.818,-0.009) (8.743,-0.009) (8.668,-0.009) (8.593,-0.008) (8.518,-0.008) (8.443,-0.008) (8.368,-0.008) (8.293,-0.008) (8.218,-0.007) (8.143,-0.007) (8.068,-0.007) (7.993,-0.007) (7.918,-0.007) (7.843,-0.006) (7.768,-0.006) (7.693,-0.006) (7.618,-0.006) (7.543,-0.006) (7.468,-0.006) (7.393,-0.005) (7.318,-0.005) (7.243,-0.005) (7.168,-0.005) (7.093,-0.005) (7.018,-0.005) (6.943,-0.004) (6.868,-0.004) (6.793,-0.004) (6.718,-0.004) (6.643,-0.004) (6.568,-0.004) (6.493,-0.004) (6.418,-0.004) (6.343,-0.003) (6.268,-0.003) (6.193,-0.003) (6.118,-0.003) (6.043,-0.003) (5.968,-0.003) (5.893,-0.003) (5.818,-0.003) (5.743,-0.003) (5.668,-0.002) (5.593,-0.002) (5.518,-0.002) (5.443,-0.002) (5.418,-0.002)}
\LC{(-0.001,2.928) (-0.001,3.003) (-0.001,3.078) (-0.001,3.153) (-0.001,3.228) (-0.002,3.303) (-0.002,3.378) (-0.002,3.453) (-0.002,3.528) (-0.002,3.603) (-0.002,3.678) (-0.002,3.753) (-0.002,3.828) (-0.002,3.903) (-0.002,3.978) (-0.002,4.053) (-0.002,4.128) (-0.002,4.203) (-0.003,4.278) (-0.003,4.353) (-0.003,4.428) (-0.003,4.503) (-0.003,4.578) (-0.003,4.653) (-0.003,4.728) (-0.003,4.803) (-0.003,4.878) (-0.003,4.953) (-0.004,5.028) (-0.004,5.103) (-0.004,5.178) (-0.004,5.253) (-0.004,5.328) (-0.004,5.403) (-0.004,5.478) (-0.004,5.553) (-0.004,5.628) (-0.005,5.703) (-0.005,5.778) (-0.005,5.853) (-0.005,5.928) (-0.005,6.003) (-0.005,6.078) (-0.005,6.153) (-0.005,6.228) (-0.006,6.303) (-0.006,6.378) (-0.006,6.453) (-0.006,6.528) (-0.006,6.603) (-0.006,6.678) (-0.006,6.753) (-0.007,6.828) (-0.007,6.903) (-0.007,6.978) (-0.007,7.053) (-0.007,7.128) (-0.007,7.203) (-0.007,7.278) (-0.008,7.353) (-0.008,7.428) (-0.008,7.503) (-0.008,7.578) (-0.008,7.653) (-0.008,7.728) (-0.009,7.803) (-0.009,7.878) (-0.009,7.953) (-0.009,8.028) (-0.009,8.103) (-0.009,8.178) (-0.010,8.253) (-0.010,8.328) (-0.010,8.403) (-0.010,8.478) (-0.010,8.553) (-0.010,8.628) (-0.011,8.703) (-0.011,8.778) (-0.011,8.853) (-0.011,8.928) (-0.011,9.003) (-0.012,9.078) (-0.012,9.153) (-0.012,9.228) (-0.012,9.303) (-0.012,9.378) (-0.013,9.453) (-0.013,9.528) (-0.013,9.603) (-0.013,9.678) (-0.013,9.753) (-0.014,9.828) (-0.014,9.903) (-0.014,9.978) (-0.014,10.000)}
\LC{(4.007,7.839) (4.071,7.879) (4.135,7.919) (4.199,7.957) (4.264,7.995) (4.329,8.032) (4.394,8.069) (4.460,8.105) (4.527,8.140) (4.593,8.174) (4.660,8.208) (4.728,8.241) (4.795,8.273) (4.864,8.304) (4.932,8.335) (5.001,8.365) (5.070,8.394) (5.140,8.422) (5.209,8.449) (5.280,8.476) (5.350,8.501) (5.421,8.526) (5.492,8.550) (5.563,8.573) (5.635,8.596) (5.707,8.617) (5.779,8.637) (5.851,8.657) (5.924,8.675) (5.997,8.693) (6.070,8.710) (6.143,8.726) (6.217,8.740) (6.291,8.754) (6.364,8.767) (6.439,8.779) (6.513,8.789) (6.587,8.799) (6.662,8.808) (6.736,8.816) (6.811,8.822) (6.886,8.828) (6.961,8.832) (7.036,8.835) (7.110,8.837) (7.185,8.838) (7.260,8.838) (7.335,8.837) (7.410,8.835) (7.485,8.831) (7.560,8.826) (7.635,8.820) (7.710,8.813) (7.784,8.805) (7.859,8.795) (7.933,8.784) (8.007,8.772) (8.081,8.759) (8.154,8.744) (8.227,8.728) (8.300,8.711) (8.373,8.692) (8.445,8.672) (8.517,8.651) (8.589,8.629) (8.660,8.605) (8.731,8.580) (8.801,8.553) (8.870,8.525) (8.939,8.496) (9.008,8.466) (9.076,8.434) (9.143,8.401) (9.210,8.366) (9.275,8.330) (9.341,8.293) (9.405,8.254) (9.468,8.214) (9.531,8.173) (9.593,8.130) (9.653,8.086) (9.713,8.041) (9.772,7.994) (9.830,7.946) (9.886,7.897) (9.942,7.847) (9.996,7.795) (10.000,7.778)}
\LC{(-0.891,-2.699) (-0.937,-2.758) (-0.984,-2.816) (-1.033,-2.873) (-1.082,-2.930) (-1.133,-2.985) (-1.184,-3.039) (-1.237,-3.093) (-1.291,-3.145) (-1.346,-3.196) (-1.402,-3.246) (-1.459,-3.295) (-1.517,-3.342) (-1.576,-3.388) (-1.636,-3.434) (-1.697,-3.477) (-1.759,-3.520) (-1.822,-3.560) (-1.885,-3.600) (-1.950,-3.638) (-2.016,-3.674) (-2.082,-3.709) (-2.149,-3.743) (-2.217,-3.774) (-2.286,-3.804) (-2.356,-3.832) (-2.426,-3.858) (-2.497,-3.883) (-2.568,-3.905) (-2.640,-3.926) (-2.713,-3.944) (-2.786,-3.961) (-2.860,-3.975) (-2.934,-3.987) (-3.008,-3.997) (-3.083,-4.004) (-3.158,-4.009) (-3.233,-4.012) (-3.308,-4.012) (-3.383,-4.010) (-3.457,-4.005) (-3.532,-3.997) (-3.606,-3.987) (-3.680,-3.974) (-3.753,-3.958) (-3.826,-3.939) (-3.898,-3.918) (-3.969,-3.893) (-4.038,-3.866) (-4.107,-3.835) (-4.174,-3.802) (-4.240,-3.765) (-4.303,-3.726) (-4.365,-3.684) (-4.425,-3.638) (-4.483,-3.590) (-4.538,-3.539) (-4.590,-3.486) (-4.640,-3.430) (-4.687,-3.371) (-4.730,-3.310) (-4.771,-3.247) (-4.808,-3.182) (-4.842,-3.115) (-4.872,-3.046) (-4.898,-2.976) (-4.921,-2.904) (-4.940,-2.832) (-4.955,-2.758) (-4.966,-2.684) (-4.974,-2.610) (-4.978,-2.535) (-4.979,-2.460) (-4.975,-2.385) (-4.969,-2.310) (-4.959,-2.236) (-4.945,-2.162) (-4.928,-2.089) (-4.909,-2.017) (-4.886,-1.945) (-4.860,-1.875) (-4.831,-1.805) (-4.800,-1.737) (-4.767,-1.670) (-4.730,-1.605) (-4.692,-1.540) (-4.651,-1.477) (-4.608,-1.416) (-4.563,-1.356) (-4.517,-1.297) (-4.468,-1.240) (-4.418,-1.184) (-4.366,-1.130) (-4.313,-1.077) (-4.258,-1.026) (-4.201,-0.977) (-4.144,-0.929) (-4.085,-0.882) (-4.025,-0.837) (-3.964,-0.793) (-3.902,-0.751) (-3.839,-0.711) (-3.775,-0.671) (-3.710,-0.634) (-3.667,-0.609)}
\LC{(1.347,2.119) (1.418,2.141) (1.491,2.160) (1.564,2.175) (1.638,2.187) (1.713,2.194) (1.788,2.196) (1.863,2.194) (1.938,2.187) (2.012,2.176) (2.085,2.159) (2.156,2.137) (2.226,2.109) (2.294,2.076) (2.358,2.038) (2.419,1.994) (2.476,1.945) (2.528,1.891) (2.575,1.832) (2.616,1.770) (2.651,1.703) (2.679,1.634) (2.701,1.562) (2.716,1.489) (2.723,1.414) (2.725,1.339) (2.719,1.264) (2.708,1.190) (2.691,1.117) (2.667,1.046) (2.639,0.977) (2.606,0.909) (2.582,0.866)}
\LC{(-1.963,-4.791) (-2.018,-4.843) (-2.073,-4.893) (-2.129,-4.943) (-2.186,-4.992) (-2.243,-5.041) (-2.302,-5.088) (-2.360,-5.135) (-2.420,-5.181) (-2.479,-5.226) (-2.540,-5.270) (-2.601,-5.313) (-2.663,-5.356) (-2.726,-5.397) (-2.789,-5.438) (-2.853,-5.477) (-2.917,-5.516) (-2.982,-5.553) (-3.047,-5.590) (-3.113,-5.626) (-3.180,-5.660) (-3.247,-5.694) (-3.315,-5.726) (-3.383,-5.757) (-3.451,-5.787) (-3.521,-5.816) (-3.590,-5.844) (-3.660,-5.871) (-3.731,-5.896) (-3.802,-5.921) (-3.873,-5.944) (-3.945,-5.965) (-4.017,-5.986) (-4.090,-6.005) (-4.163,-6.022) (-4.236,-6.039) (-4.309,-6.054) (-4.383,-6.067) (-4.457,-6.079) (-4.531,-6.090) (-4.606,-6.099) (-4.680,-6.106) (-4.755,-6.112) (-4.830,-6.117) (-4.905,-6.119) (-4.980,-6.121) (-5.055,-6.120) (-5.130,-6.118) (-5.205,-6.114) (-5.280,-6.109) (-5.354,-6.102) (-5.429,-6.092) (-5.503,-6.082) (-5.577,-6.069) (-5.650,-6.054) (-5.724,-6.038) (-5.796,-6.020) (-5.869,-6.000) (-5.940,-5.978) (-6.011,-5.954) (-6.082,-5.928) (-6.152,-5.900) (-6.220,-5.871) (-6.288,-5.839) (-6.355,-5.805) (-6.421,-5.770) (-6.486,-5.732) (-6.550,-5.693) (-6.613,-5.651) (-6.674,-5.608) (-6.734,-5.563) (-6.792,-5.516) (-6.849,-5.467) (-6.904,-5.416) (-6.958,-5.363) (-7.009,-5.309) (-7.059,-5.253) (-7.107,-5.196) (-7.153,-5.136) (-7.197,-5.076) (-7.239,-5.013) (-7.279,-4.950) (-7.317,-4.885) (-7.352,-4.819) (-7.385,-4.752) (-7.416,-4.683) (-7.444,-4.614) (-7.470,-4.543) (-7.493,-4.472) (-7.515,-4.400) (-7.533,-4.327) (-7.549,-4.254) (-7.563,-4.180) (-7.574,-4.106) (-7.583,-4.032) (-7.589,-3.957) (-7.593,-3.882) (-7.595,-3.807) (-7.594,-3.732) (-7.591,-3.657) (-7.585,-3.582) (-7.577,-3.508) (-7.567,-3.434) (-7.555,-3.360) (-7.541,-3.286) (-7.524,-3.213) (-7.505,-3.140) (-7.485,-3.068) (-7.462,-2.996) (-7.438,-2.926) (-7.412,-2.855) (-7.383,-2.786) (-7.353,-2.717) (-7.322,-2.649) (-7.288,-2.582) (-7.254,-2.516) (-7.217,-2.450) (-7.179,-2.385) (-7.139,-2.322) (-7.098,-2.259) (-7.056,-2.197) (-7.012,-2.136) (-6.967,-2.076) (-6.921,-2.017) (-6.874,-1.959) (-6.825,-1.902) (-6.775,-1.846) (-6.724,-1.791) (-6.672,-1.737) (-6.619,-1.683) (-6.566,-1.631) (-6.511,-1.580) (-6.455,-1.530) (-6.398,-1.481) (-6.341,-1.433) (-6.282,-1.386) (-6.223,-1.339) (-6.163,-1.294) (-6.103,-1.250) (-6.042,-1.207) (-5.980,-1.164) (-5.959,-1.151)}
\LC{(-3.628,-6.337) (-3.695,-6.369) (-3.764,-6.400) (-3.832,-6.430) (-3.901,-6.460) (-3.971,-6.488) (-4.041,-6.515) (-4.111,-6.541) (-4.182,-6.565) (-4.253,-6.589) (-4.325,-6.612) (-4.397,-6.633) (-4.469,-6.653) (-4.541,-6.672) (-4.614,-6.690) (-4.687,-6.707) (-4.761,-6.722) (-4.834,-6.736) (-4.908,-6.749) (-4.983,-6.760) (-5.057,-6.770) (-5.131,-6.779) (-5.206,-6.787) (-5.281,-6.793) (-5.356,-6.797) (-5.430,-6.800) (-5.505,-6.802) (-5.580,-6.802) (-5.655,-6.801) (-5.730,-6.798) (-5.805,-6.794) (-5.880,-6.788) (-5.955,-6.781) (-6.029,-6.771) (-6.103,-6.761) (-6.177,-6.749) (-6.251,-6.735) (-6.324,-6.719) (-6.397,-6.702) (-6.470,-6.683) (-6.542,-6.662) (-6.614,-6.640) (-6.685,-6.616) (-6.755,-6.590) (-6.825,-6.562) (-6.894,-6.533) (-6.962,-6.502) (-7.029,-6.469) (-7.096,-6.434) (-7.162,-6.398) (-7.226,-6.360) (-7.290,-6.320) (-7.352,-6.278) (-7.413,-6.235) (-7.473,-6.190) (-7.532,-6.143) (-7.589,-6.095) (-7.645,-6.045) (-7.700,-5.993) (-7.752,-5.940) (-7.804,-5.885) (-7.853,-5.829) (-7.901,-5.771) (-7.947,-5.712) (-7.991,-5.652) (-8.034,-5.590) (-8.074,-5.527) (-8.113,-5.462) (-8.149,-5.397) (-8.184,-5.330) (-8.216,-5.262) (-8.246,-5.194) (-8.274,-5.124) (-8.300,-5.054) (-8.324,-4.983) (-8.346,-4.911) (-8.365,-4.838) (-8.382,-4.765) (-8.397,-4.692) (-8.410,-4.618) (-8.420,-4.544) (-8.428,-4.469) (-8.435,-4.394) (-8.439,-4.320) (-8.440,-4.245) (-8.440,-4.170) (-8.437,-4.095) (-8.433,-4.020) (-8.426,-3.945) (-8.418,-3.871) (-8.407,-3.796) (-8.394,-3.722) (-8.380,-3.649) (-8.363,-3.576) (-8.345,-3.503) (-8.325,-3.431) (-8.303,-3.359) (-8.279,-3.288) (-8.254,-3.217) (-8.227,-3.147) (-8.198,-3.078) (-8.168,-3.009) (-8.136,-2.941) (-8.103,-2.874) (-8.068,-2.808) (-8.032,-2.742) (-7.994,-2.677) (-7.955,-2.613) (-7.915,-2.550) (-7.873,-2.487) (-7.831,-2.426) (-7.787,-2.365) (-7.742,-2.305) (-7.695,-2.246) (-7.648,-2.188) (-7.600,-2.130) (-7.550,-2.074) (-7.500,-2.019) (-7.448,-1.964) (-7.431,-1.946)}
\LC{(3.506,6.475) (3.571,6.511) (3.638,6.546) (3.704,6.580) (3.772,6.614) (3.839,6.646) (3.907,6.678) (3.976,6.708) (4.045,6.738) (4.114,6.767) (4.184,6.794) (4.254,6.821) (4.324,6.846) (4.395,6.871) (4.467,6.894) (4.538,6.916) (4.610,6.937) (4.682,6.957) (4.755,6.976) (4.828,6.994) (4.901,7.010) (4.974,7.026) (5.048,7.040) (5.122,7.053) (5.196,7.064) (5.270,7.075) (5.345,7.084) (5.419,7.092) (5.494,7.098) (5.569,7.103) (5.644,7.107) (5.719,7.109) (5.794,7.110) (5.869,7.110) (5.944,7.108) (6.019,7.104) (6.094,7.099) (6.168,7.093) (6.243,7.085) (6.317,7.076) (6.391,7.065) (6.465,7.052) (6.539,7.038) (6.612,7.023) (6.685,7.005) (6.758,6.986) (6.830,6.966) (6.902,6.944) (6.973,6.920) (7.043,6.894) (7.113,6.867) (7.183,6.839) (7.251,6.808) (7.319,6.776) (7.386,6.742) (7.452,6.706) (7.517,6.669) (7.581,6.630) (7.644,6.590) (7.706,6.548) (7.767,6.504) (7.827,6.458) (7.885,6.411) (7.942,6.363) (7.998,6.312) (8.052,6.261) (8.105,6.207) (8.156,6.152) (8.206,6.096) (8.254,6.039) (8.300,5.980) (8.344,5.919) (8.387,5.858) (8.428,5.795) (8.467,5.731) (8.504,5.666) (8.540,5.599) (8.573,5.532) (8.604,5.464) (8.633,5.395) (8.660,5.325) (8.685,5.254) (8.708,5.183) (8.729,5.111) (8.748,5.038) (8.765,4.965) (8.779,4.892) (8.792,4.818) (8.802,4.743) (8.810,4.669) (8.816,4.594) (8.820,4.519) (8.822,4.444) (8.822,4.369) (8.820,4.294) (8.816,4.219) (8.810,4.145) (8.801,4.070) (8.791,3.996) (8.780,3.922) (8.766,3.848) (8.750,3.775) (8.733,3.702) (8.714,3.629) (8.693,3.557) (8.670,3.486) (8.646,3.415) (8.620,3.344) (8.593,3.274) (8.564,3.205) (8.533,3.137) (8.501,3.069) (8.468,3.002) (8.433,2.935) (8.397,2.870) (8.359,2.805) (8.321,2.740) (8.281,2.677) (8.239,2.614) (8.197,2.552) (8.153,2.491) (8.109,2.431) (8.063,2.372) (8.016,2.313) (7.968,2.256) (7.919,2.199) (7.869,2.143) (7.818,2.088) (7.767,2.033)}
\LC{(-1.864,3.540) (-1.929,3.577) (-1.995,3.614) (-2.061,3.648) (-2.129,3.681) (-2.197,3.712) (-2.266,3.741) (-2.336,3.768) (-2.406,3.794) (-2.477,3.818) (-2.549,3.839) (-2.622,3.859) (-2.695,3.877) (-2.768,3.892) (-2.842,3.905) (-2.916,3.916) (-2.991,3.925) (-3.065,3.931) (-3.140,3.935) (-3.215,3.936) (-3.290,3.934) (-3.365,3.930) (-3.440,3.924) (-3.514,3.914) (-3.588,3.902) (-3.662,3.887) (-3.734,3.869) (-3.806,3.848) (-3.877,3.824) (-3.947,3.797) (-4.016,3.767) (-4.084,3.734) (-4.149,3.698) (-4.213,3.659) (-4.276,3.617) (-4.336,3.572) (-4.393,3.525) (-4.449,3.474) (-4.501,3.420) (-4.551,3.364) (-4.598,3.306) (-4.641,3.245) (-4.682,3.181) (-4.718,3.116) (-4.752,3.049) (-4.781,2.980) (-4.807,2.910) (-4.830,2.838) (-4.848,2.765) (-4.863,2.692) (-4.873,2.618) (-4.880,2.543) (-4.883,2.468) (-4.883,2.393) (-4.879,2.318) (-4.871,2.244) (-4.860,2.169) (-4.845,2.096) (-4.827,2.023) (-4.806,1.951) (-4.782,1.880) (-4.755,1.810) (-4.726,1.741) (-4.693,1.673) (-4.658,1.607) (-4.621,1.542) (-4.581,1.478) (-4.540,1.416) (-4.496,1.355) (-4.450,1.296) (-4.402,1.238) (-4.352,1.182) (-4.301,1.127)}
\LC{(2.522,3.823) (2.594,3.843) (2.667,3.861) (2.740,3.877) (2.814,3.891) (2.888,3.902) (2.962,3.911) (3.037,3.918) (3.112,3.923) (3.187,3.924) (3.262,3.924) (3.337,3.920) (3.412,3.914) (3.486,3.906) (3.560,3.894) (3.634,3.880) (3.707,3.862) (3.779,3.842) (3.850,3.819) (3.921,3.793) (3.990,3.763) (4.057,3.731) (4.124,3.696) (4.188,3.657) (4.250,3.616) (4.311,3.572) (4.369,3.524) (4.425,3.474) (4.478,3.421) (4.528,3.366) (4.576,3.308) (4.620,3.247) (4.661,3.184) (4.698,3.119) (4.732,3.052) (4.763,2.984) (4.789,2.914) (4.812,2.842) (4.831,2.770) (4.846,2.696) (4.858,2.622) (4.865,2.548) (4.869,2.473) (4.869,2.398) (4.865,2.323) (4.858,2.248) (4.847,2.174) (4.833,2.100) (4.816,2.027) (4.795,1.955) (4.772,1.884) (4.745,1.814) (4.716,1.745) (4.684,1.677) (4.649,1.610) (4.612,1.545) (4.599,1.524)}
\LC{(-2.405,-4.586) (-2.470,-4.624) (-2.535,-4.661) (-2.601,-4.697) (-2.668,-4.731) (-2.735,-4.764) (-2.803,-4.796) (-2.871,-4.827) (-2.940,-4.856) (-3.010,-4.883) (-3.080,-4.910) (-3.151,-4.934) (-3.223,-4.958) (-3.294,-4.979) (-3.367,-5.000) (-3.439,-5.018) (-3.512,-5.035) (-3.586,-5.050) (-3.659,-5.064) (-3.733,-5.076) (-3.808,-5.086) (-3.882,-5.094) (-3.957,-5.101) (-4.032,-5.105) (-4.107,-5.108) (-4.182,-5.109) (-4.257,-5.108) (-4.332,-5.104) (-4.407,-5.099) (-4.481,-5.091) (-4.556,-5.082) (-4.630,-5.070) (-4.703,-5.056) (-4.777,-5.040) (-4.849,-5.022) (-4.921,-5.001) (-4.993,-4.978) (-5.063,-4.953) (-5.133,-4.925) (-5.202,-4.895) (-5.270,-4.863) (-5.336,-4.829) (-5.402,-4.792) (-5.465,-4.753) (-5.528,-4.711) (-5.589,-4.667) (-5.648,-4.621) (-5.705,-4.573) (-5.761,-4.522) (-5.814,-4.470) (-5.865,-4.415) (-5.915,-4.358) (-5.961,-4.300) (-6.006,-4.239) (-6.047,-4.177) (-6.087,-4.113) (-6.123,-4.048) (-6.157,-3.981) (-6.188,-3.912) (-6.216,-3.843) (-6.242,-3.772) (-6.264,-3.701) (-6.284,-3.628) (-6.300,-3.555) (-6.314,-3.481) (-6.325,-3.407) (-6.332,-3.333) (-6.337,-3.258) (-6.339,-3.183) (-6.338,-3.108) (-6.334,-3.033) (-6.328,-2.958) (-6.319,-2.884) (-6.307,-2.810) (-6.292,-2.736) (-6.276,-2.663) (-6.256,-2.591) (-6.234,-2.519) (-6.210,-2.448) (-6.184,-2.378) (-6.155,-2.308) (-6.125,-2.240) (-6.092,-2.172) (-6.058,-2.106) (-6.021,-2.040) (-5.983,-1.976) (-5.943,-1.912) (-5.901,-1.850) (-5.857,-1.789) (-5.812,-1.729) (-5.766,-1.670) (-5.718,-1.612) (-5.669,-1.556) (-5.618,-1.500) (-5.566,-1.446) (-5.513,-1.393) (-5.459,-1.342) (-5.403,-1.291) (-5.347,-1.242) (-5.290,-1.193) (-5.231,-1.146) (-5.172,-1.101) (-5.111,-1.056) (-5.050,-1.013) (-4.988,-0.970) (-4.926,-0.929) (-4.862,-0.889) (-4.798,-0.850) (-4.733,-0.812) (-4.668,-0.776)}
\LC{(-3.747,6.524) (-3.815,6.556) (-3.883,6.587) (-3.952,6.617) (-4.021,6.646) (-4.091,6.674) (-4.161,6.701) (-4.231,6.727) (-4.302,6.751) (-4.373,6.775) (-4.445,6.798) (-4.517,6.819) (-4.589,6.840) (-4.661,6.859) (-4.734,6.877) (-4.807,6.893) (-4.880,6.909) (-4.954,6.923) (-5.028,6.936) (-5.102,6.948) (-5.176,6.959) (-5.251,6.968) (-5.325,6.976) (-5.400,6.982) (-5.475,6.987) (-5.550,6.991) (-5.625,6.994) (-5.700,6.995) (-5.775,6.994) (-5.850,6.992) (-5.925,6.989) (-5.999,6.984) (-6.074,6.977) (-6.149,6.969) (-6.223,6.960) (-6.297,6.948) (-6.371,6.936) (-6.445,6.921) (-6.518,6.905) (-6.591,6.888) (-6.663,6.869) (-6.736,6.848) (-6.807,6.825) (-6.878,6.801) (-6.948,6.775) (-7.018,6.747) (-7.087,6.718) (-7.155,6.687) (-7.223,6.654) (-7.290,6.620) (-7.355,6.584) (-7.420,6.546) (-7.484,6.506) (-7.546,6.465) (-7.608,6.422) (-7.668,6.378) (-7.727,6.331) (-7.785,6.283) (-7.841,6.234) (-7.896,6.183) (-7.950,6.130) (-8.002,6.076) (-8.052,6.021) (-8.101,5.964) (-8.148,5.905) (-8.193,5.845) (-8.236,5.784) (-8.278,5.722) (-8.317,5.658) (-8.355,5.593) (-8.391,5.527) (-8.425,5.460) (-8.457,5.392) (-8.486,5.324) (-8.514,5.254) (-8.539,5.183) (-8.563,5.112) (-8.584,5.040) (-8.603,4.968) (-8.620,4.895) (-8.635,4.821) (-8.648,4.747) (-8.658,4.673) (-8.666,4.598) (-8.673,4.524) (-8.677,4.449) (-8.679,4.374) (-8.679,4.299) (-8.676,4.224) (-8.672,4.149) (-8.666,4.074) (-8.658,4.000) (-8.648,3.925) (-8.636,3.851) (-8.622,3.778) (-8.606,3.704) (-8.588,3.631) (-8.569,3.559) (-8.548,3.487) (-8.525,3.416) (-8.500,3.345) (-8.474,3.274) (-8.446,3.205) (-8.417,3.136) (-8.386,3.067) (-8.354,3.000) (-8.320,2.933) (-8.285,2.867) (-8.248,2.801) (-8.210,2.736) (-8.171,2.672) (-8.130,2.609) (-8.089,2.547) (-8.046,2.486) (-8.002,2.425) (-7.956,2.365) (-7.910,2.306) (-7.863,2.248) (-7.814,2.191) (-7.765,2.134) (-7.714,2.079) (-7.663,2.024) (-7.646,2.006)}
\LC{(-1.421,-5.598) (-1.458,-5.663) (-1.495,-5.729) (-1.532,-5.794) (-1.570,-5.858) (-1.608,-5.923) (-1.647,-5.987) (-1.686,-6.051) (-1.726,-6.115) (-1.766,-6.178) (-1.807,-6.241) (-1.848,-6.304) (-1.889,-6.366) (-1.931,-6.429) (-1.973,-6.490) (-2.016,-6.552) (-2.059,-6.613) (-2.103,-6.674) (-2.147,-6.735) (-2.192,-6.795) (-2.237,-6.855) (-2.282,-6.915) (-2.328,-6.974) (-2.374,-7.033) (-2.421,-7.092) (-2.468,-7.150) (-2.516,-7.208) (-2.564,-7.266) (-2.613,-7.323) (-2.661,-7.380) (-2.711,-7.436) (-2.761,-7.492) (-2.811,-7.548) (-2.861,-7.603) (-2.912,-7.658) (-2.964,-7.713) (-3.016,-7.767) (-3.068,-7.821) (-3.121,-7.874) (-3.174,-7.927) (-3.228,-7.980) (-3.281,-8.032) (-3.336,-8.083) (-3.391,-8.135) (-3.446,-8.185) (-3.501,-8.236) (-3.557,-8.286) (-3.614,-8.335) (-3.671,-8.384) (-3.728,-8.432) (-3.786,-8.480) (-3.844,-8.528) (-3.902,-8.575) (-3.961,-8.622) (-4.020,-8.668) (-4.080,-8.713) (-4.139,-8.758) (-4.200,-8.803) (-4.261,-8.847) (-4.322,-8.890) (-4.383,-8.933) (-4.445,-8.976) (-4.507,-9.018) (-4.570,-9.059) (-4.633,-9.100) (-4.696,-9.140) (-4.760,-9.180) (-4.824,-9.219) (-4.888,-9.257) (-4.953,-9.295) (-5.018,-9.333) (-5.083,-9.369) (-5.149,-9.405) (-5.215,-9.441) (-5.281,-9.476) (-5.348,-9.510) (-5.415,-9.544) (-5.482,-9.577) (-5.550,-9.609) (-5.618,-9.641) (-5.686,-9.672) (-5.755,-9.702) (-5.824,-9.732) (-5.893,-9.761) (-5.962,-9.789) (-6.032,-9.817) (-6.102,-9.843) (-6.172,-9.870) (-6.243,-9.895) (-6.314,-9.920) (-6.385,-9.944) (-6.456,-9.967) (-6.528,-9.989) (-6.576,-10.000)}
\LC{(1.060,5.164) (1.090,5.232) (1.121,5.301) (1.151,5.369) (1.183,5.437) (1.214,5.505) (1.247,5.573) (1.279,5.641) (1.312,5.708) (1.345,5.775) (1.379,5.842) (1.413,5.909) (1.448,5.975) (1.483,6.042) (1.518,6.108) (1.554,6.174) (1.590,6.239) (1.627,6.305) (1.664,6.370) (1.701,6.435) (1.739,6.500) (1.778,6.564) (1.816,6.629) (1.855,6.693) (1.895,6.756) (1.935,6.820) (1.975,6.883) (2.016,6.946) (2.057,7.009) (2.098,7.071) (2.140,7.134) (2.182,7.196) (2.225,7.257) (2.268,7.319) (2.312,7.380) (2.355,7.441) (2.400,7.501) (2.444,7.561) (2.489,7.621) (2.535,7.681) (2.581,7.740) (2.627,7.799) (2.674,7.858) (2.721,7.916) (2.768,7.975) (2.816,8.032) (2.864,8.090) (2.913,8.147) (2.962,8.204) (3.011,8.260) (3.061,8.316) (3.111,8.372) (3.162,8.428) (3.212,8.483) (3.264,8.537) (3.315,8.592) (3.367,8.646) (3.420,8.699) (3.473,8.753) (3.526,8.806) (3.579,8.858) (3.633,8.910) (3.687,8.962) (3.742,9.013) (3.797,9.064) (3.853,9.115) (3.908,9.165) (3.964,9.215) (4.021,9.264) (4.078,9.313) (4.135,9.362) (4.192,9.410) (4.250,9.457) (4.309,9.505) (4.367,9.551) (4.426,9.598) (4.486,9.644) (4.545,9.689) (4.605,9.734) (4.666,9.779) (4.726,9.823) (4.787,9.866) (4.849,9.909) (4.910,9.952) (4.972,9.994) (4.993,10.000)}
\LC{(-4.211,7.912) (-4.276,7.950) (-4.342,7.986) (-4.408,8.022) (-4.474,8.057) (-4.540,8.092) (-4.607,8.126) (-4.675,8.159) (-4.743,8.191) (-4.811,8.222) (-4.879,8.253) (-4.948,8.283) (-5.017,8.312) (-5.086,8.340) (-5.156,8.367) (-5.226,8.394) (-5.297,8.420) (-5.368,8.445) (-5.439,8.469) (-5.510,8.492) (-5.582,8.514) (-5.654,8.535) (-5.726,8.556) (-5.798,8.575) (-5.871,8.594) (-5.944,8.611) (-6.017,8.628) (-6.090,8.644) (-6.164,8.659) (-6.237,8.672) (-6.311,8.685) (-6.385,8.697) (-6.460,8.707) (-6.534,8.717) (-6.609,8.725) (-6.683,8.733) (-6.758,8.739) (-6.833,8.745) (-6.908,8.749) (-6.983,8.752) (-7.058,8.754) (-7.133,8.755) (-7.208,8.754) (-7.283,8.753) (-7.357,8.750) (-7.432,8.746) (-7.507,8.741) (-7.582,8.735) (-7.657,8.727) (-7.731,8.719) (-7.805,8.709) (-7.880,8.697) (-7.953,8.685) (-8.027,8.671) (-8.101,8.656) (-8.174,8.640) (-8.247,8.622) (-8.319,8.603) (-8.391,8.583) (-8.463,8.561) (-8.535,8.538) (-8.606,8.514) (-8.676,8.488) (-8.746,8.461) (-8.815,8.432) (-8.884,8.403) (-8.952,8.371) (-9.020,8.339) (-9.087,8.305) (-9.153,8.270) (-9.219,8.233) (-9.283,8.195) (-9.347,8.156) (-9.410,8.115) (-9.472,8.073) (-9.534,8.030) (-9.594,7.985) (-9.653,7.939) (-9.711,7.892) (-9.768,7.844) (-9.824,7.794) (-9.879,7.743) (-9.933,7.690) (-9.985,7.637) (-10.000,7.618)}
\LC{(4.811,-7.150) (4.883,-7.169) (4.956,-7.187) (5.029,-7.203) (5.103,-7.219) (5.176,-7.233) (5.250,-7.246) (5.324,-7.258) (5.398,-7.269) (5.473,-7.278) (5.547,-7.286) (5.622,-7.293) (5.697,-7.299) (5.772,-7.303) (5.847,-7.306) (5.922,-7.307) (5.997,-7.307) (6.072,-7.306) (6.146,-7.303) (6.221,-7.299) (6.296,-7.294) (6.371,-7.287) (6.445,-7.278) (6.520,-7.268) (6.594,-7.257) (6.668,-7.244) (6.741,-7.229) (6.815,-7.213) (6.887,-7.195) (6.960,-7.176) (7.032,-7.155) (7.104,-7.133) (7.175,-7.109) (7.245,-7.083) (7.315,-7.056) (7.384,-7.027) (7.453,-6.997) (7.520,-6.965) (7.587,-6.931) (7.653,-6.895) (7.719,-6.858) (7.783,-6.820) (7.846,-6.779) (7.908,-6.737) (7.969,-6.694) (8.029,-6.649) (8.088,-6.602) (8.145,-6.554) (8.202,-6.504) (8.256,-6.453) (8.310,-6.400) (8.361,-6.346) (8.412,-6.290) (8.460,-6.233) (8.507,-6.175) (8.553,-6.115) (8.597,-6.054) (8.639,-5.992) (8.679,-5.929) (8.717,-5.864) (8.753,-5.799) (8.788,-5.732) (8.821,-5.664) (8.851,-5.596) (8.880,-5.527) (8.907,-5.457) (8.931,-5.386) (8.954,-5.314) (8.974,-5.242) (8.993,-5.169) (9.009,-5.096) (9.024,-5.023) (9.036,-4.949) (9.046,-4.874) (9.054,-4.800) (9.060,-4.725) (9.065,-4.650) (9.067,-4.575) (9.067,-4.500) (9.065,-4.425) (9.061,-4.350) (9.055,-4.276) (9.048,-4.201) (9.038,-4.127) (9.027,-4.052) (9.014,-3.979) (8.999,-3.905) (8.982,-3.832) (8.963,-3.759) (8.943,-3.687) (8.921,-3.615) (8.898,-3.544) (8.873,-3.473) (8.846,-3.403) (8.818,-3.334) (8.788,-3.265) (8.757,-3.197) (8.725,-3.129) (8.691,-3.062) (8.656,-2.996) (8.619,-2.930) (8.581,-2.866) (8.555,-2.823)}
\LC{(-0.995,2.331) (-1.051,2.381) (-1.109,2.428) (-1.169,2.474) (-1.230,2.518) (-1.292,2.560) (-1.355,2.600) (-1.420,2.638) (-1.486,2.674) (-1.553,2.707) (-1.621,2.739) (-1.690,2.767) (-1.760,2.794) (-1.831,2.818) (-1.903,2.839) (-1.976,2.857) (-2.049,2.873) (-2.123,2.885) (-2.198,2.895) (-2.272,2.901) (-2.347,2.904) (-2.422,2.903) (-2.497,2.899) (-2.572,2.891) (-2.646,2.879) (-2.719,2.864) (-2.792,2.844) (-2.863,2.821) (-2.933,2.793) (-3.001,2.762) (-3.067,2.726) (-3.130,2.686) (-3.191,2.642) (-3.249,2.594) (-3.303,2.542) (-3.354,2.487) (-3.400,2.428) (-3.442,2.366) (-3.480,2.302) (-3.513,2.234) (-3.541,2.165) (-3.564,2.093) (-3.581,2.020) (-3.594,1.946) (-3.601,1.872) (-3.603,1.797) (-3.600,1.722) (-3.593,1.647) (-3.580,1.573) (-3.564,1.500) (-3.543,1.428) (-3.517,1.358) (-3.488,1.288) (-3.456,1.221) (-3.420,1.155) (-3.381,1.091) (-3.338,1.029) (-3.293,0.969) (-3.246,0.911) (-3.196,0.855) (-3.143,0.802) (-3.089,0.750) (-3.033,0.700) (-2.975,0.653) (-2.915,0.607) (-2.854,0.564) (-2.791,0.523) (-2.728,0.483) (-2.662,0.446) (-2.596,0.410) (-2.574,0.399)}
\LC{(-10.000,-0.013) (-9.942,-0.013) (-9.867,-0.013) (-9.792,-0.012) (-9.717,-0.012) (-9.642,-0.012) (-9.567,-0.012) (-9.492,-0.011) (-9.417,-0.011) (-9.342,-0.011) (-9.267,-0.010) (-9.192,-0.010) (-9.117,-0.010) (-9.042,-0.010) (-8.967,-0.010) (-8.892,-0.009) (-8.817,-0.009) (-8.742,-0.009) (-8.667,-0.009) (-8.592,-0.008) (-8.517,-0.008) (-8.442,-0.008) (-8.367,-0.008) (-8.292,-0.008) (-8.217,-0.007) (-8.142,-0.007) (-8.067,-0.007) (-7.992,-0.007) (-7.917,-0.007) (-7.842,-0.006) (-7.767,-0.006) (-7.692,-0.006) (-7.617,-0.006) (-7.542,-0.006) (-7.467,-0.005) (-7.392,-0.005) (-7.317,-0.005) (-7.242,-0.005) (-7.167,-0.005) (-7.142,-0.005)}
\LC{(1.556,-6.813) (1.589,-6.880) (1.622,-6.947) (1.656,-7.014) (1.690,-7.081) (1.725,-7.147) (1.760,-7.214) (1.795,-7.280) (1.830,-7.346) (1.866,-7.412) (1.902,-7.478) (1.939,-7.543) (1.976,-7.608) (2.013,-7.673) (2.051,-7.738) (2.089,-7.803) (2.127,-7.868) (2.165,-7.932) (2.204,-7.996) (2.244,-8.060) (2.283,-8.124) (2.323,-8.187) (2.363,-8.250) (2.404,-8.314) (2.445,-8.376) (2.486,-8.439) (2.528,-8.501) (2.570,-8.564) (2.612,-8.626) (2.654,-8.687) (2.697,-8.749) (2.741,-8.810) (2.784,-8.871) (2.828,-8.932) (2.872,-8.992) (2.917,-9.053) (2.962,-9.113) (3.007,-9.173) (3.053,-9.232) (3.099,-9.292) (3.145,-9.351) (3.191,-9.409) (3.238,-9.468) (3.286,-9.526) (3.333,-9.584) (3.381,-9.642) (3.429,-9.699) (3.478,-9.757) (3.526,-9.814) (3.576,-9.870) (3.625,-9.927) (3.675,-9.983) (3.691,-10.000)}
\LC{(0.769,5.621) (0.789,5.693) (0.810,5.765) (0.831,5.837) (0.852,5.909) (0.873,5.981) (0.895,6.053) (0.917,6.125) (0.939,6.196) (0.961,6.268) (0.984,6.339) (1.007,6.411) (1.031,6.482) (1.054,6.553) (1.078,6.624) (1.102,6.695) (1.126,6.766) (1.151,6.837) (1.176,6.908) (1.201,6.979) (1.226,7.049) (1.252,7.120) (1.278,7.190) (1.304,7.260) (1.330,7.331) (1.357,7.401) (1.384,7.471) (1.411,7.541) (1.439,7.610) (1.466,7.680) (1.494,7.750) (1.523,7.819) (1.551,7.888) (1.580,7.958) (1.609,8.027) (1.638,8.096) (1.668,8.165) (1.698,8.234) (1.728,8.302) (1.758,8.371) (1.789,8.439) (1.820,8.508) (1.851,8.576) (1.882,8.644) (1.914,8.712) (1.946,8.780) (1.978,8.848) (2.010,8.915) (2.043,8.983) (2.076,9.050) (2.109,9.117) (2.143,9.184) (2.176,9.251) (2.210,9.318) (2.245,9.385) (2.279,9.452) (2.314,9.518) (2.349,9.584) (2.384,9.651) (2.420,9.717) (2.456,9.783) (2.492,9.848) (2.528,9.914) (2.565,9.979) (2.577,10.000)}
\LC{(3.537,-7.525) (3.597,-7.570) (3.658,-7.614) (3.719,-7.658) (3.780,-7.701) (3.842,-7.743) (3.904,-7.785) (3.967,-7.826) (4.030,-7.867) (4.094,-7.907) (4.158,-7.946) (4.222,-7.984) (4.287,-8.022) (4.352,-8.059) (4.418,-8.095) (4.484,-8.131) (4.550,-8.166) (4.617,-8.200) (4.684,-8.234) (4.751,-8.266) (4.819,-8.298) (4.887,-8.330) (4.956,-8.360) (5.025,-8.390) (5.094,-8.419) (5.164,-8.447) (5.234,-8.474) (5.304,-8.500) (5.374,-8.526) (5.445,-8.551) (5.516,-8.574) (5.588,-8.597) (5.659,-8.619) (5.731,-8.641) (5.803,-8.661) (5.876,-8.680) (5.948,-8.699) (6.021,-8.716) (6.095,-8.733) (6.168,-8.749) (6.241,-8.763) (6.315,-8.777) (6.389,-8.790) (6.463,-8.801) (6.537,-8.812) (6.612,-8.821) (6.686,-8.830) (6.761,-8.838) (6.836,-8.844) (6.911,-8.849) (6.985,-8.854) (7.060,-8.857) (7.135,-8.859) (7.210,-8.860) (7.285,-8.860) (7.360,-8.858) (7.435,-8.856) (7.510,-8.852) (7.585,-8.847) (7.660,-8.841) (7.734,-8.834) (7.809,-8.825) (7.883,-8.816) (7.957,-8.805) (8.031,-8.792) (8.105,-8.779) (8.179,-8.764) (8.252,-8.748) (8.325,-8.731) (8.398,-8.712) (8.470,-8.692) (8.542,-8.671) (8.613,-8.648) (8.684,-8.624) (8.755,-8.599) (8.825,-8.573) (8.895,-8.545) (8.964,-8.515) (9.032,-8.485) (9.100,-8.453) (9.167,-8.420) (9.234,-8.385) (9.300,-8.349) (9.365,-8.312) (9.429,-8.273) (9.493,-8.233) (9.555,-8.192) (9.617,-8.149) (9.678,-8.105) (9.737,-8.060) (9.796,-8.013) (9.854,-7.965) (9.910,-7.916) (9.966,-7.866) (10.000,-7.831)}
\LC{(0.740,-5.756) (0.760,-5.829) (0.779,-5.901) (0.799,-5.974) (0.819,-6.046) (0.839,-6.118) (0.859,-6.190) (0.880,-6.262) (0.901,-6.334) (0.922,-6.406) (0.943,-6.478) (0.965,-6.550) (0.987,-6.622) (1.009,-6.694) (1.031,-6.765) (1.054,-6.837) (1.076,-6.908) (1.099,-6.979) (1.123,-7.051) (1.146,-7.122) (1.170,-7.193) (1.194,-7.264) (1.218,-7.335) (1.243,-7.406) (1.268,-7.477) (1.293,-7.547) (1.318,-7.618) (1.344,-7.689) (1.369,-7.759) (1.395,-7.829) (1.422,-7.900) (1.448,-7.970) (1.475,-8.040) (1.502,-8.110) (1.529,-8.180) (1.556,-8.250) (1.584,-8.319) (1.612,-8.389) (1.640,-8.458) (1.668,-8.528) (1.697,-8.597) (1.726,-8.666) (1.755,-8.736) (1.784,-8.805) (1.814,-8.873) (1.844,-8.942) (1.874,-9.011) (1.904,-9.080) (1.935,-9.148) (1.966,-9.216) (1.997,-9.285) (2.028,-9.353) (2.060,-9.421) (2.091,-9.489) (2.123,-9.557) (2.156,-9.624) (2.188,-9.692) (2.221,-9.759) (2.254,-9.827) (2.287,-9.894) (2.320,-9.961) (2.343,-10.000)}
\LC{(0.443,6.926) (0.453,7.000) (0.463,7.075) (0.472,7.149) (0.482,7.223) (0.492,7.298) (0.503,7.372) (0.513,7.446) (0.523,7.521) (0.534,7.595) (0.544,7.669) (0.555,7.743) (0.566,7.818) (0.576,7.892) (0.587,7.966) (0.598,8.040) (0.610,8.114) (0.621,8.189) (0.632,8.263) (0.644,8.337) (0.655,8.411) (0.667,8.485) (0.679,8.559) (0.691,8.633) (0.703,8.707) (0.715,8.781) (0.727,8.855) (0.739,8.929) (0.752,9.003) (0.764,9.077) (0.777,9.151) (0.790,9.225) (0.802,9.299) (0.815,9.373) (0.828,9.447) (0.841,9.520) (0.855,9.594) (0.868,9.668) (0.881,9.742) (0.895,9.816) (0.909,9.889) (0.922,9.963) (0.932,10.000)}
\LC{(-0.456,6.301) (-0.467,6.376) (-0.478,6.450) (-0.489,6.524) (-0.500,6.598) (-0.511,6.672) (-0.523,6.746) (-0.535,6.820) (-0.546,6.894) (-0.558,6.969) (-0.570,7.043) (-0.582,7.117) (-0.595,7.191) (-0.607,7.265) (-0.620,7.338) (-0.632,7.412) (-0.645,7.486) (-0.658,7.560) (-0.671,7.634) (-0.684,7.708) (-0.697,7.782) (-0.711,7.855) (-0.724,7.929) (-0.738,8.003) (-0.752,8.077) (-0.766,8.150) (-0.780,8.224) (-0.794,8.298) (-0.808,8.371) (-0.823,8.445) (-0.837,8.518) (-0.852,8.592) (-0.867,8.666) (-0.882,8.739) (-0.897,8.812) (-0.912,8.886) (-0.928,8.959) (-0.943,9.033) (-0.959,9.106) (-0.974,9.179) (-0.990,9.253) (-1.006,9.326) (-1.022,9.399) (-1.039,9.472) (-1.055,9.546) (-1.071,9.619) (-1.088,9.692) (-1.105,9.765) (-1.122,9.838) (-1.139,9.911) (-1.156,9.984) (-1.162,10.000)}
\LC{(-5.902,-9.274) (-5.973,-9.298) (-6.045,-9.322) (-6.116,-9.344) (-6.188,-9.365) (-6.260,-9.386) (-6.333,-9.406) (-6.405,-9.425) (-6.478,-9.443) (-6.551,-9.460) (-6.624,-9.477) (-6.698,-9.492) (-6.771,-9.507) (-6.845,-9.521) (-6.919,-9.533) (-6.993,-9.545) (-7.067,-9.556) (-7.141,-9.566) (-7.216,-9.575) (-7.290,-9.583) (-7.365,-9.590) (-7.440,-9.596) (-7.515,-9.601) (-7.590,-9.605) (-7.664,-9.608) (-7.739,-9.611) (-7.814,-9.612) (-7.889,-9.611) (-7.964,-9.610) (-8.039,-9.608) (-8.114,-9.605) (-8.189,-9.601) (-8.264,-9.595) (-8.339,-9.589) (-8.413,-9.581) (-8.488,-9.572) (-8.562,-9.562) (-8.636,-9.551) (-8.710,-9.539) (-8.784,-9.525) (-8.858,-9.511) (-8.931,-9.495) (-9.004,-9.478) (-9.077,-9.460) (-9.149,-9.441) (-9.221,-9.420) (-9.293,-9.398) (-9.365,-9.375) (-9.436,-9.351) (-9.506,-9.326) (-9.576,-9.299) (-9.646,-9.271) (-9.715,-9.242) (-9.784,-9.212) (-9.852,-9.180) (-9.919,-9.147) (-9.986,-9.113) (-10.000,-9.102)}
\LC{(0.919,-1.901) (0.981,-1.943) (1.045,-1.983) (1.110,-2.020) (1.177,-2.054) (1.245,-2.085) (1.315,-2.113) (1.386,-2.138) (1.458,-2.159) (1.530,-2.176) (1.604,-2.189) (1.679,-2.199) (1.754,-2.204) (1.829,-2.204) (1.903,-2.200) (1.978,-2.191) (2.051,-2.177) (2.124,-2.158) (2.195,-2.133) (2.264,-2.103) (2.330,-2.068) (2.393,-2.028) (2.452,-1.982) (2.507,-1.931) (2.557,-1.875) (2.602,-1.814) (2.640,-1.750) (2.673,-1.683) (2.698,-1.612) (2.717,-1.540) (2.730,-1.466) (2.735,-1.391) (2.734,-1.316) (2.726,-1.241) (2.712,-1.168) (2.693,-1.095) (2.668,-1.025) (2.638,-0.956) (2.603,-0.889) (2.564,-0.825) (2.521,-0.764) (2.475,-0.705) (2.425,-0.649) (2.408,-0.630)}
\LC{(2.081,-3.206) (2.153,-3.228) (2.226,-3.247) (2.299,-3.263) (2.373,-3.277) (2.447,-3.288) (2.521,-3.297) (2.596,-3.302) (2.671,-3.305) (2.746,-3.305) (2.821,-3.301) (2.896,-3.294) (2.970,-3.284) (3.044,-3.271) (3.117,-3.254) (3.189,-3.234) (3.260,-3.210) (3.330,-3.183) (3.398,-3.152) (3.465,-3.117) (3.529,-3.079) (3.592,-3.037) (3.651,-2.992) (3.708,-2.943) (3.762,-2.891) (3.813,-2.836) (3.860,-2.777) (3.903,-2.716) (3.943,-2.652) (3.978,-2.586) (4.009,-2.518) (4.036,-2.448) (4.058,-2.376) (4.075,-2.303) (4.088,-2.229) (4.097,-2.155) (4.101,-2.080) (4.100,-2.005) (4.095,-1.930) (4.086,-1.856) (4.073,-1.782) (4.056,-1.709) (4.036,-1.637) (4.011,-1.566) (3.983,-1.496) (3.952,-1.428) (3.918,-1.361) (3.881,-1.296) (3.855,-1.253)}
\LC{(-0.657,-1.870) (-0.706,-1.926) (-0.757,-1.982) (-0.809,-2.036) (-0.863,-2.088) (-0.918,-2.138) (-0.975,-2.187) (-1.034,-2.234) (-1.094,-2.279) (-1.155,-2.322) (-1.218,-2.363) (-1.283,-2.401) (-1.348,-2.438) (-1.415,-2.471) (-1.483,-2.503) (-1.553,-2.532) (-1.623,-2.558) (-1.694,-2.581) (-1.766,-2.601) (-1.839,-2.618) (-1.913,-2.632) (-1.987,-2.642) (-2.062,-2.649) (-2.137,-2.652) (-2.212,-2.652) (-2.287,-2.648) (-2.361,-2.639) (-2.435,-2.627) (-2.508,-2.610) (-2.580,-2.589) (-2.651,-2.563) (-2.720,-2.533) (-2.786,-2.499) (-2.850,-2.460) (-2.911,-2.416) (-2.969,-2.369) (-3.023,-2.317) (-3.073,-2.261) (-3.119,-2.201) (-3.159,-2.138) (-3.195,-2.072) (-3.225,-2.003) (-3.249,-1.932) (-3.268,-1.860) (-3.282,-1.786) (-3.289,-1.712) (-3.291,-1.637) (-3.288,-1.562) (-3.279,-1.487) (-3.266,-1.413) (-3.247,-1.341) (-3.224,-1.269) (-3.197,-1.200) (-3.165,-1.132) (-3.130,-1.065) (-3.091,-1.001) (-3.049,-0.939) (-3.004,-0.879) (-2.956,-0.822) (-2.905,-0.766) (-2.852,-0.713) (-2.797,-0.663) (-2.740,-0.614) (-2.681,-0.568) (-2.620,-0.524) (-2.599,-0.510)}
\LC{(3.005,4.986) (3.075,5.014) (3.145,5.040) (3.216,5.066) (3.287,5.089) (3.358,5.112) (3.431,5.133) (3.503,5.152) (3.576,5.170) (3.649,5.186) (3.723,5.200) (3.797,5.213) (3.871,5.224) (3.945,5.234) (4.020,5.241) (4.095,5.247) (4.169,5.251) (4.244,5.253) (4.319,5.253) (4.394,5.252) (4.469,5.248) (4.544,5.242) (4.619,5.234) (4.693,5.225) (4.767,5.213) (4.841,5.199) (4.914,5.182) (4.987,5.164) (5.059,5.143) (5.130,5.121) (5.201,5.095) (5.271,5.068) (5.340,5.039) (5.408,5.007) (5.474,4.973) (5.540,4.936) (5.604,4.897) (5.667,4.857) (5.728,4.813) (5.788,4.768) (5.846,4.720) (5.902,4.671) (5.956,4.619) (6.009,4.565) (6.059,4.509) (6.107,4.452) (6.152,4.392) (6.195,4.331) (6.236,4.268) (6.274,4.203) (6.310,4.137) (6.343,4.070) (6.373,4.001) (6.400,3.931) (6.425,3.860) (6.447,3.789) (6.465,3.716) (6.481,3.643) (6.495,3.569) (6.505,3.495) (6.512,3.420) (6.517,3.345) (6.519,3.270) (6.518,3.195) (6.514,3.120) (6.508,3.046) (6.499,2.971) (6.487,2.897) (6.473,2.823) (6.456,2.750) (6.437,2.678) (6.416,2.606) (6.392,2.535) (6.367,2.464) (6.339,2.395) (6.309,2.326) (6.277,2.258) (6.243,2.191) (6.207,2.125) (6.169,2.060) (6.130,1.996) (6.089,1.934) (6.047,1.872) (6.002,1.811) (5.972,1.771)}
\LC{(10.000,1.463) (9.942,1.431) (9.874,1.399) (9.805,1.368) (9.737,1.338) (9.668,1.308) (9.599,1.278) (9.530,1.249) (9.461,1.221) (9.391,1.192) (9.322,1.165) (9.252,1.137) (9.182,1.110) (9.112,1.084) (9.041,1.058) (8.971,1.032) (8.900,1.007) (8.829,0.982) (8.758,0.957) (8.687,0.933) (8.616,0.910) (8.545,0.886) (8.474,0.864) (8.402,0.841) (8.330,0.819) (8.258,0.797) (8.187,0.776) (8.115,0.755) (8.042,0.735) (7.970,0.714) (7.898,0.695) (7.825,0.675) (7.753,0.656) (7.680,0.638) (7.608,0.619) (7.535,0.601) (7.462,0.584) (7.389,0.566) (7.316,0.549) (7.243,0.533) (7.169,0.516) (7.096,0.500) (7.023,0.485) (6.949,0.470) (6.876,0.455) (6.851,0.450)}
\LC{(3.473,-5.665) (3.543,-5.692) (3.614,-5.717) (3.685,-5.741) (3.756,-5.764) (3.828,-5.786) (3.900,-5.806) (3.973,-5.825) (4.045,-5.843) (4.119,-5.859) (4.192,-5.874) (4.266,-5.887) (4.340,-5.898) (4.414,-5.909) (4.489,-5.917) (4.564,-5.924) (4.638,-5.930) (4.713,-5.934) (4.788,-5.936) (4.863,-5.936) (4.938,-5.935) (5.013,-5.932) (5.088,-5.927) (5.163,-5.921) (5.237,-5.913) (5.312,-5.902) (5.386,-5.890) (5.459,-5.876) (5.533,-5.861) (5.606,-5.843) (5.678,-5.823) (5.750,-5.801) (5.821,-5.778) (5.891,-5.752) (5.961,-5.725) (6.030,-5.695) (6.098,-5.663) (6.165,-5.630) (6.231,-5.594) (6.296,-5.556) (6.360,-5.517) (6.422,-5.475) (6.483,-5.432) (6.542,-5.386) (6.601,-5.338) (6.657,-5.289) (6.712,-5.238) (6.765,-5.185) (6.816,-5.130) (6.865,-5.073) (6.912,-5.015) (6.957,-4.955) (7.001,-4.894) (7.041,-4.831) (7.080,-4.767) (7.116,-4.701) (7.150,-4.634) (7.182,-4.566) (7.211,-4.497) (7.238,-4.427) (7.262,-4.356) (7.284,-4.284) (7.303,-4.212) (7.319,-4.138) (7.333,-4.065) (7.345,-3.991) (7.354,-3.916) (7.360,-3.842) (7.364,-3.767) (7.366,-3.692) (7.365,-3.617) (7.362,-3.542) (7.356,-3.467) (7.348,-3.392) (7.337,-3.318) (7.325,-3.244) (7.310,-3.171) (7.293,-3.098) (7.274,-3.025) (7.253,-2.953) (7.229,-2.882) (7.204,-2.811) (7.177,-2.741) (7.148,-2.672) (7.117,-2.604) (7.085,-2.536) (7.051,-2.469) (7.015,-2.403) (6.978,-2.338) (6.939,-2.274) (6.898,-2.211) (6.856,-2.149) (6.813,-2.088) (6.768,-2.028)}
\LC{(-3.826,5.589) (-3.899,5.607) (-3.972,5.622) (-4.046,5.636) (-4.119,5.649) (-4.194,5.660) (-4.268,5.670) (-4.343,5.677) (-4.417,5.684) (-4.492,5.688) (-4.567,5.691) (-4.642,5.692) (-4.717,5.691) (-4.792,5.688) (-4.867,5.684) (-4.942,5.678) (-5.016,5.670) (-5.091,5.660) (-5.165,5.647) (-5.238,5.634) (-5.312,5.618) (-5.384,5.600) (-5.457,5.580) (-5.528,5.558) (-5.599,5.533) (-5.670,5.507) (-5.739,5.479) (-5.808,5.449) (-5.875,5.416) (-5.942,5.382) (-6.007,5.345) (-6.071,5.306) (-6.134,5.265) (-6.196,5.222) (-6.256,5.177) (-6.314,5.130) (-6.371,5.081) (-6.426,5.031) (-6.480,4.978) (-6.531,4.923) (-6.580,4.867) (-6.628,4.808) (-6.673,4.749) (-6.716,4.687) (-6.756,4.624) (-6.794,4.559) (-6.830,4.494) (-6.864,4.426) (-6.895,4.358) (-6.923,4.289) (-6.949,4.218) (-6.972,4.147) (-6.992,4.075) (-7.010,4.002) (-7.026,3.929) (-7.038,3.855) (-7.048,3.780) (-7.055,3.706) (-7.060,3.631) (-7.062,3.556) (-7.062,3.481) (-7.059,3.406) (-7.053,3.331) (-7.046,3.257) (-7.035,3.182) (-7.023,3.108) (-7.008,3.035) (-6.990,2.962) (-6.971,2.889) (-6.949,2.818) (-6.926,2.746) (-6.900,2.676) (-6.872,2.606) (-6.843,2.537) (-6.811,2.469) (-6.778,2.402) (-6.743,2.336) (-6.707,2.270)}
\LC{(0.590,-1.027) (0.659,-1.056) (0.731,-1.079) (0.804,-1.094) (0.879,-1.100) (0.954,-1.098) (1.028,-1.086) (1.099,-1.064) (1.167,-1.031) (1.227,-0.987) (1.279,-0.933) (1.320,-0.870) (1.349,-0.801) (1.363,-0.728) (1.365,-0.653) (1.353,-0.579) (1.330,-0.507) (1.298,-0.440) (1.256,-0.377) (1.208,-0.320) (1.190,-0.302)}
\LC{(-2.212,7.292) (-2.255,7.354) (-2.298,7.415) (-2.341,7.476) (-2.385,7.537) (-2.429,7.598) (-2.474,7.658) (-2.519,7.718) (-2.564,7.778) (-2.610,7.837) (-2.656,7.896) (-2.703,7.955) (-2.750,8.013) (-2.797,8.071) (-2.845,8.129) (-2.893,8.187) (-2.942,8.244) (-2.991,8.301) (-3.040,8.357) (-3.090,8.414) (-3.140,8.469) (-3.190,8.525) (-3.241,8.580) (-3.292,8.635) (-3.344,8.689) (-3.395,8.744) (-3.448,8.797) (-3.500,8.851) (-3.553,8.904) (-3.607,8.956) (-3.661,9.009) (-3.715,9.061) (-3.769,9.112) (-3.824,9.163) (-3.879,9.214) (-3.935,9.264) (-3.991,9.314) (-4.047,9.364) (-4.104,9.413) (-4.161,9.462) (-4.218,9.510) (-4.276,9.558) (-4.334,9.605) (-4.392,9.652) (-4.451,9.699) (-4.510,9.745) (-4.570,9.791) (-4.630,9.836) (-4.690,9.881) (-4.750,9.925) (-4.811,9.969) (-4.872,10.000)}
\LC{(3.158,7.740) (3.213,7.791) (3.267,7.843) (3.322,7.894) (3.378,7.944) (3.434,7.994) (3.491,8.043) (3.547,8.092) (3.605,8.140) (3.662,8.188) (3.721,8.236) (3.779,8.283) (3.838,8.329) (3.897,8.375) (3.957,8.420) (4.017,8.465) (4.078,8.510) (4.138,8.553) (4.200,8.597) (4.261,8.639) (4.323,8.681) (4.386,8.723) (4.449,8.764) (4.512,8.804) (4.575,8.844) (4.639,8.883) (4.704,8.922) (4.768,8.960) (4.833,8.997) (4.899,9.034) (4.964,9.070) (5.030,9.106) (5.097,9.141) (5.164,9.175) (5.231,9.208) (5.298,9.241) (5.366,9.274) (5.434,9.305) (5.502,9.336) (5.571,9.366) (5.640,9.395) (5.709,9.424) (5.779,9.452) (5.849,9.479) (5.919,9.506) (5.989,9.532) (6.060,9.557) (6.131,9.581) (6.202,9.605) (6.274,9.627) (6.345,9.649) (6.417,9.670) (6.489,9.691) (6.562,9.710) (6.635,9.729) (6.707,9.747) (6.780,9.764) (6.854,9.780) (6.927,9.795) (7.001,9.809) (7.074,9.823) (7.148,9.836) (7.222,9.847) (7.297,9.858) (7.371,9.868) (7.445,9.877) (7.520,9.885) (7.595,9.892) (7.669,9.898) (7.744,9.903) (7.819,9.908) (7.894,9.911) (7.969,9.913) (8.044,9.914) (8.119,9.914) (8.194,9.914) (8.269,9.912) (8.344,9.909) (8.419,9.905) (8.494,9.900) (8.568,9.894) (8.643,9.886) (8.718,9.878) (8.792,9.869) (8.866,9.858) (8.940,9.847) (9.014,9.834) (9.088,9.820) (9.162,9.805) (9.235,9.789) (9.308,9.772) (9.380,9.753) (9.453,9.734) (9.525,9.713) (9.597,9.691) (9.668,9.668) (9.739,9.644) (9.809,9.618) (9.880,9.591) (9.949,9.564) (10.000,9.534)}
\LC{(-1.193,-1.771) (-1.266,-1.788) (-1.340,-1.800) (-1.415,-1.807) (-1.490,-1.809) (-1.565,-1.805) (-1.639,-1.796) (-1.713,-1.780) (-1.784,-1.758) (-1.853,-1.729) (-1.919,-1.693) (-1.982,-1.651) (-2.039,-1.603) (-2.090,-1.549) (-2.136,-1.489) (-2.174,-1.424) (-2.204,-1.356) (-2.226,-1.284) (-2.240,-1.210) (-2.245,-1.135) (-2.242,-1.061) (-2.232,-0.986) (-2.215,-0.913) (-2.190,-0.842) (-2.160,-0.774) (-2.124,-0.708) (-2.083,-0.645) (-2.037,-0.586) (-1.988,-0.529) (-1.935,-0.477)}
\LC{(-2.191,-3.252) (-2.264,-3.270) (-2.337,-3.286) (-2.411,-3.299) (-2.485,-3.309) (-2.560,-3.317) (-2.635,-3.321) (-2.710,-3.323) (-2.785,-3.321) (-2.860,-3.317) (-2.934,-3.309) (-3.009,-3.298) (-3.082,-3.284) (-3.155,-3.266) (-3.227,-3.245) (-3.298,-3.220) (-3.367,-3.191) (-3.435,-3.159) (-3.501,-3.124) (-3.565,-3.085) (-3.626,-3.042) (-3.686,-2.996) (-3.742,-2.946) (-3.795,-2.893) (-3.845,-2.837) (-3.891,-2.778) (-3.933,-2.716) (-3.972,-2.652) (-4.006,-2.585) (-4.036,-2.517) (-4.062,-2.446) (-4.083,-2.374) (-4.100,-2.301) (-4.112,-2.227) (-4.120,-2.153) (-4.123,-2.078) (-4.122,-2.003) (-4.116,-1.928) (-4.106,-1.853) (-4.093,-1.780) (-4.075,-1.707) (-4.054,-1.635) (-4.029,-1.564) (-4.001,-1.495) (-3.969,-1.427) (-3.935,-1.360) (-3.923,-1.338)}
\LC{(-10.000,-2.495) (-9.955,-2.443) (-9.899,-2.393) (-9.844,-2.342) (-9.787,-2.293) (-9.730,-2.244) (-9.673,-2.196) (-9.615,-2.148) (-9.556,-2.102) (-9.497,-2.055) (-9.438,-2.010) (-9.378,-1.965) (-9.317,-1.921) (-9.256,-1.877) (-9.195,-1.834) (-9.133,-1.791) (-9.070,-1.750) (-9.008,-1.709) (-8.945,-1.668) (-8.881,-1.628) (-8.817,-1.589) (-8.753,-1.550) (-8.688,-1.512) (-8.623,-1.475) (-8.558,-1.438) (-8.492,-1.402) (-8.426,-1.366) (-8.360,-1.331) (-8.293,-1.297) (-8.226,-1.263) (-8.159,-1.230) (-8.092,-1.197) (-8.024,-1.165) (-7.956,-1.133) (-7.888,-1.102) (-7.865,-1.092)}
\LC{(-1.547,2.185) (-1.621,2.198) (-1.695,2.206) (-1.770,2.211) (-1.845,2.210) (-1.920,2.206) (-1.994,2.196) (-2.068,2.181) (-2.140,2.161) (-2.211,2.136) (-2.279,2.105) (-2.345,2.069) (-2.408,2.028) (-2.466,1.981) (-2.521,1.930) (-2.570,1.873) (-2.614,1.813) (-2.652,1.748) (-2.684,1.680) (-2.709,1.609) (-2.727,1.537) (-2.739,1.463) (-2.743,1.388) (-2.741,1.313) (-2.733,1.238) (-2.719,1.165) (-2.699,1.092) (-2.673,1.022) (-2.654,0.976)}

    % Cuadrupolo puntual - puntas de flecha
% Generado por scripts/generar_datos.py (no editar a mano).
\FL{1.698}{1.255}{-50.0}
\FL{-1.668}{1.205}{-127.5}
\FL{-1.140}{-0.805}{125.1}
\FL{0.742}{-2.503}{-56.3}
\FL{3.093}{-4.038}{-5.2}
\FL{5.016}{-2.372}{95.0}
\FL{3.140}{-0.363}{160.0}
\FL{0.727}{-4.161}{-70.0}
\FL{2.222}{-6.978}{-53.8}
\FL{4.485}{-9.221}{-35.2}
\FL{-0.525}{-3.023}{-109.9}
\FL{-1.850}{-5.390}{-129.0}
\FL{-3.921}{-7.133}{-151.7}
\FL{-6.535}{-7.778}{177.6}
\FL{-8.941}{-6.719}{131.8}
\FL{-9.587}{-4.176}{78.1}
\FL{-8.164}{-1.908}{40.8}
\FL{-5.787}{-0.612}{18.2}
\FL{-3.118}{-0.093}{5.1}
\FL{-0.841}{2.977}{122.2}
\FL{-3.357}{4.910}{166.1}
\FL{-6.080}{3.778}{-111.5}
\FL{-4.942}{0.990}{-34.9}
\FL{1.636}{2.861}{25.9}
\FL{3.688}{1.358}{-114.5}
\FL{0.772}{3.688}{66.1}
\FL{2.378}{6.084}{45.7}
\FL{4.794}{7.690}{20.3}
\FL{7.642}{7.872}{-15.6}
\FL{9.761}{6.016}{-69.4}
\FL{9.462}{3.225}{-119.6}
\FL{7.364}{1.238}{-150.8}
\FL{4.636}{0.291}{-169.2}
\FL{2.129}{4.115}{31.8}
\FL{4.967}{4.269}{-33.8}
\FL{5.293}{1.622}{-125.9}
\FL{-8.711}{0.934}{-18.5}
\FL{-6.038}{0.304}{-8.7}
\FL{2.566}{-6.228}{-43.4}
\FL{4.795}{-7.625}{-19.7}
\FL{7.426}{-7.819}{13.5}
\FL{9.508}{-6.272}{62.2}
\FL{9.603}{-3.705}{111.2}
\FL{7.947}{-1.649}{143.8}
\FL{5.564}{-0.521}{163.9}
\FL{-1.077}{4.763}{115.8}
\FL{-2.868}{7.337}{134.3}
\FL{-5.432}{9.133}{156.6}
\FL{-8.506}{9.632}{-173.2}
\FL{-8.188}{-0.711}{15.0}
\FL{2.765}{5.367}{32.0}
\FL{5.361}{5.998}{-7.9}
\FL{7.382}{4.398}{-72.9}
\FL{6.722}{1.874}{-131.0}
\FL{0.247}{-5.362}{-84.7}
\FL{0.618}{-8.464}{-81.6}
\FL{-0.453}{-5.987}{-98.7}
\FL{-0.962}{-8.689}{-102.7}
\FL{2.681}{-4.875}{-28.2}
\FL{5.831}{-4.848}{37.3}
\FL{6.106}{-1.874}{125.8}
\FL{-3.542}{7.008}{146.9}
\FL{-6.429}{7.945}{179.2}
\FL{-9.161}{6.822}{-130.9}
\FL{-9.701}{3.932}{-72.4}
\FL{-7.832}{1.559}{-34.6}
\FL{8.491}{-1.022}{159.2}
\FL{5.981}{-9.411}{-19.0}
\FL{8.689}{-9.707}{8.0}
\FL{3.731}{-6.301}{-23.7}
\FL{6.270}{-6.585}{13.7}
\FL{8.162}{-4.931}{71.5}
\FL{7.742}{-2.457}{123.9}
\FL{-1.237}{6.475}{111.8}
\FL{-2.440}{8.889}{121.3}
\FL{-4.616}{-8.202}{-153.1}
\FL{-8.349}{-8.729}{165.8}
\FL{-1.263}{-6.445}{-112.4}
\FL{-2.523}{-8.885}{-122.3}
\FL{-2.814}{5.567}{146.9}
\FL{-5.712}{6.264}{-170.3}
\FL{-7.797}{4.315}{-99.7}
\FL{-6.691}{1.609}{-42.1}
\FL{7.693}{-0.006}{179.9}
\FL{-0.003}{4.703}{90.1}
\FL{-0.010}{8.253}{90.1}
\FL{5.468}{8.542}{18.5}
\FL{8.660}{8.605}{-19.1}
\FL{-2.426}{-3.858}{-160.2}
\FL{-4.909}{-2.017}{73.5}
\FL{2.476}{1.945}{-43.4}
\FL{-3.025}{-5.578}{-150.9}
\FL{-5.552}{-6.073}{170.1}
\FL{-7.493}{-4.472}{107.3}
\FL{-6.921}{-2.017}{51.3}
\FL{-5.057}{-6.770}{-172.8}
\FL{-7.770}{-5.922}{133.4}
\FL{-8.236}{-3.171}{68.6}
\FL{4.999}{7.031}{11.0}
\FL{7.998}{6.312}{-42.9}
\FL{8.620}{3.344}{-110.8}
\FL{-3.140}{3.935}{178.1}
\FL{-4.876}{2.293}{-84.5}
\FL{4.461}{3.439}{-45.4}
\FL{-3.684}{-5.068}{-170.7}
\FL{-6.074}{-4.134}{121.2}
\FL{-5.718}{-1.612}{49.6}
\FL{-5.201}{6.962}{172.8}
\FL{-7.984}{6.094}{-133.5}
\FL{-8.465}{3.251}{-68.5}
\FL{-2.390}{-7.053}{-128.5}
\FL{-4.996}{-9.320}{-150.2}
\FL{1.790}{6.586}{59.0}
\FL{3.760}{9.030}{42.9}
\FL{-5.654}{8.535}{163.8}
\FL{-8.723}{8.470}{-158.7}
\FL{6.717}{-7.234}{11.4}
\FL{9.066}{-4.625}{88.1}
\FL{-3.337}{2.506}{-131.8}
\FL{-8.567}{-0.008}{0.2}
\FL{2.514}{-8.481}{-56.3}
\FL{1.542}{7.865}{67.6}
\FL{5.071}{-8.409}{-22.5}
\FL{8.494}{-8.685}{16.3}
\FL{1.430}{-7.923}{-69.4}
\FL{0.667}{8.485}{81.0}
\FL{-0.771}{8.175}{100.8}
\FL{-7.989}{-9.610}{178.5}
\FL{2.413}{-2.013}{37.0}
\FL{3.762}{-2.891}{45.7}
\FL{-2.808}{-2.486}{149.3}
\FL{4.544}{5.242}{-5.2}
\FL{6.518}{3.195}{-91.8}
\FL{8.450}{0.856}{-162.5}
\FL{5.188}{-5.918}{6.1}
\FL{7.365}{-3.617}{91.6}
\FL{-5.263}{5.628}{-168.0}
\FL{-7.060}{3.631}{-92.6}
\FL{1.263}{-0.952}{47.8}
\FL{-3.430}{8.779}{134.3}
\FL{4.639}{8.883}{31.3}
\FL{8.194}{9.914}{-1.0}
\FL{-2.162}{-1.446}{119.3}
\FL{-3.778}{-2.911}{134.6}
\FL{-8.987}{-1.695}{32.8}
\FL{-2.521}{1.930}{-133.9}
}
  \icono
  \marcoMMA{10}{5}{$x$ (u.\,arb.)}{$z$ (u.\,arb.)}
\end{tikzpicture}\\[1pt]{\small\scshape (b)}
\end{minipage}\hfill
\begin{minipage}[t]{0.325\textwidth}\centering
{\footnotesize Campo eléctrico}\\[3pt]
\begin{tikzpicture}[x=0.198cm, y=0.198cm]
  \panelMMA{10}{mmafondo}{1.2}{% Cuadrupolo puntual - campo vectorial
% Generado por scripts/generar_datos.py (no editar a mano).
\VF{-10.000}{-10.000}{-0.370}{0.123}{0.000}
\VF{-9.000}{-10.000}{-0.386}{0.058}{0.000}
\VF{-8.000}{-10.000}{-0.424}{-0.013}{0.037}
\VF{-7.000}{-10.000}{-0.449}{-0.097}{0.076}
\VF{-6.000}{-10.000}{-0.455}{-0.192}{0.114}
\VF{-5.000}{-10.000}{-0.438}{-0.292}{0.150}
\VF{-4.000}{-10.000}{-0.395}{-0.391}{0.182}
\VF{-3.000}{-10.000}{-0.325}{-0.480}{0.209}
\VF{-2.000}{-10.000}{-0.232}{-0.551}{0.229}
\VF{-1.000}{-10.000}{-0.121}{-0.598}{0.242}
\VF{0.000}{-10.000}{0.000}{-0.614}{0.246}
\VF{1.000}{-10.000}{0.121}{-0.598}{0.242}
\VF{2.000}{-10.000}{0.232}{-0.551}{0.229}
\VF{3.000}{-10.000}{0.325}{-0.480}{0.209}
\VF{4.000}{-10.000}{0.395}{-0.391}{0.182}
\VF{5.000}{-10.000}{0.438}{-0.292}{0.150}
\VF{6.000}{-10.000}{0.455}{-0.192}{0.114}
\VF{7.000}{-10.000}{0.449}{-0.097}{0.076}
\VF{8.000}{-10.000}{0.424}{-0.013}{0.037}
\VF{9.000}{-10.000}{0.386}{0.058}{0.000}
\VF{10.000}{-10.000}{0.370}{0.123}{0.000}
\VF{-10.000}{-9.000}{-0.341}{0.189}{0.000}
\VF{-9.000}{-9.000}{-0.388}{0.129}{0.021}
\VF{-8.000}{-9.000}{-0.444}{0.058}{0.064}
\VF{-7.000}{-9.000}{-0.486}{-0.034}{0.107}
\VF{-6.000}{-9.000}{-0.507}{-0.143}{0.150}
\VF{-5.000}{-9.000}{-0.500}{-0.262}{0.192}
\VF{-4.000}{-9.000}{-0.460}{-0.383}{0.229}
\VF{-3.000}{-9.000}{-0.386}{-0.496}{0.262}
\VF{-2.000}{-9.000}{-0.279}{-0.588}{0.286}
\VF{-1.000}{-9.000}{-0.146}{-0.648}{0.302}
\VF{0.000}{-9.000}{0.000}{-0.670}{0.307}
\VF{1.000}{-9.000}{0.146}{-0.648}{0.302}
\VF{2.000}{-9.000}{0.279}{-0.588}{0.286}
\VF{3.000}{-9.000}{0.386}{-0.496}{0.262}
\VF{4.000}{-9.000}{0.460}{-0.383}{0.229}
\VF{5.000}{-9.000}{0.500}{-0.262}{0.192}
\VF{6.000}{-9.000}{0.507}{-0.143}{0.150}
\VF{7.000}{-9.000}{0.486}{-0.034}{0.107}
\VF{8.000}{-9.000}{0.444}{0.058}{0.064}
\VF{9.000}{-9.000}{0.388}{0.129}{0.021}
\VF{10.000}{-9.000}{0.341}{0.189}{0.000}
\VF{-10.000}{-8.000}{-0.292}{0.258}{0.000}
\VF{-9.000}{-8.000}{-0.371}{0.217}{0.044}
\VF{-8.000}{-8.000}{-0.447}{0.149}{0.090}
\VF{-7.000}{-8.000}{-0.513}{0.054}{0.138}
\VF{-6.000}{-8.000}{-0.556}{-0.067}{0.187}
\VF{-5.000}{-8.000}{-0.567}{-0.208}{0.235}
\VF{-4.000}{-8.000}{-0.536}{-0.357}{0.280}
\VF{-3.000}{-8.000}{-0.460}{-0.501}{0.319}
\VF{-2.000}{-8.000}{-0.338}{-0.622}{0.349}
\VF{-1.000}{-8.000}{-0.179}{-0.703}{0.369}
\VF{0.000}{-8.000}{0.000}{-0.732}{0.376}
\VF{1.000}{-8.000}{0.179}{-0.703}{0.369}
\VF{2.000}{-8.000}{0.338}{-0.622}{0.349}
\VF{3.000}{-8.000}{0.460}{-0.501}{0.319}
\VF{4.000}{-8.000}{0.536}{-0.357}{0.280}
\VF{5.000}{-8.000}{0.567}{-0.208}{0.235}
\VF{6.000}{-8.000}{0.556}{-0.067}{0.187}
\VF{7.000}{-8.000}{0.513}{0.054}{0.138}
\VF{8.000}{-8.000}{0.447}{0.149}{0.090}
\VF{9.000}{-8.000}{0.371}{0.217}{0.044}
\VF{10.000}{-8.000}{0.292}{0.258}{0.000}
\VF{-10.000}{-7.000}{-0.229}{0.338}{0.020}
\VF{-9.000}{-7.000}{-0.321}{0.315}{0.066}
\VF{-8.000}{-7.000}{-0.420}{0.261}{0.115}
\VF{-7.000}{-7.000}{-0.514}{0.171}{0.167}
\VF{-6.000}{-7.000}{-0.591}{0.043}{0.222}
\VF{-5.000}{-7.000}{-0.632}{-0.119}{0.278}
\VF{-4.000}{-7.000}{-0.623}{-0.303}{0.332}
\VF{-3.000}{-7.000}{-0.551}{-0.488}{0.380}
\VF{-2.000}{-7.000}{-0.415}{-0.650}{0.419}
\VF{-1.000}{-7.000}{-0.224}{-0.762}{0.445}
\VF{0.000}{-7.000}{0.000}{-0.803}{0.453}
\VF{1.000}{-7.000}{0.224}{-0.762}{0.445}
\VF{2.000}{-7.000}{0.415}{-0.650}{0.419}
\VF{3.000}{-7.000}{0.551}{-0.488}{0.380}
\VF{4.000}{-7.000}{0.623}{-0.303}{0.332}
\VF{5.000}{-7.000}{0.632}{-0.119}{0.278}
\VF{6.000}{-7.000}{0.591}{0.043}{0.222}
\VF{7.000}{-7.000}{0.514}{0.171}{0.167}
\VF{8.000}{-7.000}{0.420}{0.261}{0.115}
\VF{9.000}{-7.000}{0.321}{0.315}{0.066}
\VF{10.000}{-7.000}{0.229}{0.338}{0.020}
\VF{-10.000}{-6.000}{-0.131}{0.407}{0.041}
\VF{-9.000}{-6.000}{-0.228}{0.412}{0.088}
\VF{-8.000}{-6.000}{-0.344}{0.387}{0.140}
\VF{-7.000}{-6.000}{-0.471}{0.319}{0.196}
\VF{-6.000}{-6.000}{-0.592}{0.197}{0.257}
\VF{-5.000}{-6.000}{-0.682}{0.021}{0.321}
\VF{-4.000}{-6.000}{-0.714}{-0.201}{0.386}
\VF{-3.000}{-6.000}{-0.663}{-0.442}{0.447}
\VF{-2.000}{-6.000}{-0.517}{-0.665}{0.497}
\VF{-1.000}{-6.000}{-0.285}{-0.825}{0.531}
\VF{0.000}{-6.000}{0.000}{-0.884}{0.543}
\VF{1.000}{-6.000}{0.285}{-0.825}{0.531}
\VF{2.000}{-6.000}{0.517}{-0.665}{0.497}
\VF{3.000}{-6.000}{0.663}{-0.442}{0.447}
\VF{4.000}{-6.000}{0.714}{-0.201}{0.386}
\VF{5.000}{-6.000}{0.682}{0.021}{0.321}
\VF{6.000}{-6.000}{0.592}{0.197}{0.257}
\VF{7.000}{-6.000}{0.471}{0.319}{0.196}
\VF{8.000}{-6.000}{0.344}{0.387}{0.140}
\VF{9.000}{-6.000}{0.228}{0.412}{0.088}
\VF{10.000}{-6.000}{0.131}{0.407}{0.041}
\VF{-10.000}{-5.000}{0.000}{0.448}{0.064}
\VF{-9.000}{-5.000}{-0.086}{0.485}{0.113}
\VF{-8.000}{-5.000}{-0.203}{0.501}{0.166}
\VF{-7.000}{-5.000}{-0.352}{0.479}{0.225}
\VF{-6.000}{-5.000}{-0.522}{0.394}{0.290}
\VF{-5.000}{-5.000}{-0.683}{0.228}{0.363}
\VF{-4.000}{-5.000}{-0.790}{-0.024}{0.440}
\VF{-3.000}{-5.000}{-0.793}{-0.334}{0.517}
\VF{-2.000}{-5.000}{-0.655}{-0.649}{0.585}
\VF{-1.000}{-5.000}{-0.375}{-0.889}{0.632}
\VF{0.000}{-5.000}{0.000}{-0.981}{0.649}
\VF{1.000}{-5.000}{0.375}{-0.889}{0.632}
\VF{2.000}{-5.000}{0.655}{-0.649}{0.585}
\VF{3.000}{-5.000}{0.793}{-0.334}{0.517}
\VF{4.000}{-5.000}{0.790}{-0.024}{0.440}
\VF{5.000}{-5.000}{0.683}{0.228}{0.363}
\VF{6.000}{-5.000}{0.522}{0.394}{0.290}
\VF{7.000}{-5.000}{0.352}{0.479}{0.225}
\VF{8.000}{-5.000}{0.203}{0.501}{0.166}
\VF{9.000}{-5.000}{0.086}{0.485}{0.113}
\VF{10.000}{-5.000}{-0.000}{0.448}{0.064}
\VF{-10.000}{-4.000}{0.150}{0.445}{0.088}
\VF{-9.000}{-4.000}{0.092}{0.508}{0.138}
\VF{-8.000}{-4.000}{0.000}{0.566}{0.194}
\VF{-7.000}{-4.000}{-0.138}{0.607}{0.255}
\VF{-6.000}{-4.000}{-0.331}{0.600}{0.324}
\VF{-5.000}{-4.000}{-0.568}{0.501}{0.403}
\VF{-4.000}{-4.000}{-0.795}{0.265}{0.493}
\VF{-3.000}{-4.000}{-0.920}{-0.112}{0.590}
\VF{-2.000}{-4.000}{-0.841}{-0.561}{0.683}
\VF{-1.000}{-4.000}{-0.513}{-0.944}{0.752}
\VF{0.000}{-4.000}{0.000}{-1.099}{0.779}
\VF{1.000}{-4.000}{0.513}{-0.944}{0.752}
\VF{2.000}{-4.000}{0.841}{-0.561}{0.683}
\VF{3.000}{-4.000}{0.920}{-0.112}{0.590}
\VF{4.000}{-4.000}{0.795}{0.265}{0.493}
\VF{5.000}{-4.000}{0.568}{0.501}{0.403}
\VF{6.000}{-4.000}{0.331}{0.600}{0.324}
\VF{7.000}{-4.000}{0.138}{0.607}{0.255}
\VF{8.000}{-4.000}{-0.000}{0.566}{0.194}
\VF{9.000}{-4.000}{-0.092}{0.508}{0.138}
\VF{10.000}{-4.000}{-0.150}{0.445}{0.088}
\VF{-10.000}{-3.000}{0.296}{0.391}{0.110}
\VF{-9.000}{-3.000}{0.277}{0.462}{0.164}
\VF{-8.000}{-3.000}{0.234}{0.545}{0.223}
\VF{-7.000}{-3.000}{0.149}{0.635}{0.288}
\VF{-6.000}{-3.000}{0.000}{0.719}{0.361}
\VF{-5.000}{-3.000}{-0.243}{0.756}{0.445}
\VF{-4.000}{-3.000}{-0.587}{0.661}{0.543}
\VF{-3.000}{-3.000}{-0.940}{0.313}{0.660}
\VF{-2.000}{-3.000}{-1.067}{-0.300}{0.789}
\VF{-1.000}{-3.000}{-0.742}{-0.955}{0.900}
\VF{0.000}{-3.000}{0.000}{-1.251}{0.946}
\VF{1.000}{-3.000}{0.742}{-0.955}{0.900}
\VF{2.000}{-3.000}{1.067}{-0.300}{0.789}
\VF{3.000}{-3.000}{0.940}{0.313}{0.660}
\VF{4.000}{-3.000}{0.587}{0.661}{0.543}
\VF{5.000}{-3.000}{0.243}{0.756}{0.445}
\VF{6.000}{-3.000}{-0.000}{0.719}{0.361}
\VF{7.000}{-3.000}{-0.149}{0.635}{0.288}
\VF{8.000}{-3.000}{-0.234}{0.545}{0.223}
\VF{9.000}{-3.000}{-0.277}{0.462}{0.164}
\VF{10.000}{-3.000}{-0.296}{0.391}{0.110}
\VF{-10.000}{-2.000}{0.416}{0.289}{0.129}
\VF{-9.000}{-2.000}{0.436}{0.350}{0.186}
\VF{-8.000}{-2.000}{0.446}{0.427}{0.250}
\VF{-7.000}{-2.000}{0.436}{0.524}{0.320}
\VF{-6.000}{-2.000}{0.388}{0.646}{0.400}
\VF{-5.000}{-2.000}{0.266}{0.793}{0.491}
\VF{-4.000}{-2.000}{0.000}{0.933}{0.597}
\VF{-3.000}{-2.000}{-0.509}{0.921}{0.727}
\VF{-2.000}{-2.000}{-1.143}{0.381}{0.896}
\VF{-1.000}{-2.000}{-1.082}{-0.721}{1.000}
\VF{0.000}{-2.000}{0.000}{-1.300}{1.000}
\VF{1.000}{-2.000}{1.082}{-0.721}{1.000}
\VF{2.000}{-2.000}{1.143}{0.381}{0.896}
\VF{3.000}{-2.000}{0.509}{0.921}{0.727}
\VF{4.000}{-2.000}{-0.000}{0.933}{0.597}
\VF{5.000}{-2.000}{-0.266}{0.793}{0.491}
\VF{6.000}{-2.000}{-0.388}{0.646}{0.400}
\VF{7.000}{-2.000}{-0.436}{0.524}{0.320}
\VF{8.000}{-2.000}{-0.446}{0.427}{0.250}
\VF{9.000}{-2.000}{-0.436}{0.350}{0.186}
\VF{10.000}{-2.000}{-0.416}{0.289}{0.129}
\VF{-10.000}{-1.000}{0.495}{0.154}{0.141}
\VF{-9.000}{-1.000}{0.541}{0.188}{0.201}
\VF{-8.000}{-1.000}{0.590}{0.233}{0.268}
\VF{-7.000}{-1.000}{0.639}{0.294}{0.344}
\VF{-6.000}{-1.000}{0.684}{0.378}{0.431}
\VF{-5.000}{-1.000}{0.717}{0.499}{0.532}
\VF{-4.000}{-1.000}{0.710}{0.681}{0.653}
\VF{-3.000}{-1.000}{0.577}{0.961}{0.803}
\VF{-2.000}{-1.000}{0.000}{1.300}{1.000}
\VF{-1.000}{-1.000}{-1.233}{0.411}{1.000}
\VF{1.000}{-1.000}{1.233}{0.411}{1.000}
\VF{2.000}{-1.000}{-0.000}{1.300}{1.000}
\VF{3.000}{-1.000}{-0.577}{0.961}{0.803}
\VF{4.000}{-1.000}{-0.710}{0.681}{0.653}
\VF{5.000}{-1.000}{-0.717}{0.499}{0.532}
\VF{6.000}{-1.000}{-0.684}{0.378}{0.431}
\VF{7.000}{-1.000}{-0.639}{0.294}{0.344}
\VF{8.000}{-1.000}{-0.590}{0.233}{0.268}
\VF{9.000}{-1.000}{-0.541}{0.188}{0.201}
\VF{10.000}{-1.000}{-0.495}{0.154}{0.141}
\VF{-10.000}{0.000}{0.522}{-0.000}{0.145}
\VF{-9.000}{0.000}{0.578}{-0.000}{0.206}
\VF{-8.000}{0.000}{0.640}{-0.000}{0.275}
\VF{-7.000}{0.000}{0.711}{-0.000}{0.353}
\VF{-6.000}{0.000}{0.792}{-0.000}{0.442}
\VF{-5.000}{0.000}{0.889}{-0.000}{0.548}
\VF{-4.000}{0.000}{1.007}{-0.000}{0.678}
\VF{-3.000}{0.000}{1.159}{-0.000}{0.845}
\VF{-2.000}{0.000}{1.300}{-0.000}{1.000}
\VF{2.000}{0.000}{-1.300}{-0.000}{1.000}
\VF{3.000}{0.000}{-1.159}{-0.000}{0.845}
\VF{4.000}{0.000}{-1.007}{-0.000}{0.678}
\VF{5.000}{0.000}{-0.889}{-0.000}{0.548}
\VF{6.000}{0.000}{-0.792}{-0.000}{0.442}
\VF{7.000}{0.000}{-0.711}{-0.000}{0.353}
\VF{8.000}{0.000}{-0.640}{-0.000}{0.275}
\VF{9.000}{0.000}{-0.578}{-0.000}{0.206}
\VF{10.000}{0.000}{-0.522}{-0.000}{0.145}
\VF{-10.000}{1.000}{0.495}{-0.154}{0.141}
\VF{-9.000}{1.000}{0.541}{-0.188}{0.201}
\VF{-8.000}{1.000}{0.590}{-0.233}{0.268}
\VF{-7.000}{1.000}{0.639}{-0.294}{0.344}
\VF{-6.000}{1.000}{0.684}{-0.378}{0.431}
\VF{-5.000}{1.000}{0.717}{-0.499}{0.532}
\VF{-4.000}{1.000}{0.710}{-0.681}{0.653}
\VF{-3.000}{1.000}{0.577}{-0.961}{0.803}
\VF{-2.000}{1.000}{0.000}{-1.300}{1.000}
\VF{-1.000}{1.000}{-1.233}{-0.411}{1.000}
\VF{1.000}{1.000}{1.233}{-0.411}{1.000}
\VF{2.000}{1.000}{-0.000}{-1.300}{1.000}
\VF{3.000}{1.000}{-0.577}{-0.961}{0.803}
\VF{4.000}{1.000}{-0.710}{-0.681}{0.653}
\VF{5.000}{1.000}{-0.717}{-0.499}{0.532}
\VF{6.000}{1.000}{-0.684}{-0.378}{0.431}
\VF{7.000}{1.000}{-0.639}{-0.294}{0.344}
\VF{8.000}{1.000}{-0.590}{-0.233}{0.268}
\VF{9.000}{1.000}{-0.541}{-0.188}{0.201}
\VF{10.000}{1.000}{-0.495}{-0.154}{0.141}
\VF{-10.000}{2.000}{0.416}{-0.289}{0.129}
\VF{-9.000}{2.000}{0.436}{-0.350}{0.186}
\VF{-8.000}{2.000}{0.446}{-0.427}{0.250}
\VF{-7.000}{2.000}{0.436}{-0.524}{0.320}
\VF{-6.000}{2.000}{0.388}{-0.646}{0.400}
\VF{-5.000}{2.000}{0.266}{-0.793}{0.491}
\VF{-4.000}{2.000}{0.000}{-0.933}{0.597}
\VF{-3.000}{2.000}{-0.509}{-0.921}{0.727}
\VF{-2.000}{2.000}{-1.143}{-0.381}{0.896}
\VF{-1.000}{2.000}{-1.082}{0.721}{1.000}
\VF{0.000}{2.000}{0.000}{1.300}{1.000}
\VF{1.000}{2.000}{1.082}{0.721}{1.000}
\VF{2.000}{2.000}{1.143}{-0.381}{0.896}
\VF{3.000}{2.000}{0.509}{-0.921}{0.727}
\VF{4.000}{2.000}{-0.000}{-0.933}{0.597}
\VF{5.000}{2.000}{-0.266}{-0.793}{0.491}
\VF{6.000}{2.000}{-0.388}{-0.646}{0.400}
\VF{7.000}{2.000}{-0.436}{-0.524}{0.320}
\VF{8.000}{2.000}{-0.446}{-0.427}{0.250}
\VF{9.000}{2.000}{-0.436}{-0.350}{0.186}
\VF{10.000}{2.000}{-0.416}{-0.289}{0.129}
\VF{-10.000}{3.000}{0.296}{-0.391}{0.110}
\VF{-9.000}{3.000}{0.277}{-0.462}{0.164}
\VF{-8.000}{3.000}{0.234}{-0.545}{0.223}
\VF{-7.000}{3.000}{0.149}{-0.635}{0.288}
\VF{-6.000}{3.000}{0.000}{-0.719}{0.361}
\VF{-5.000}{3.000}{-0.243}{-0.756}{0.445}
\VF{-4.000}{3.000}{-0.587}{-0.661}{0.543}
\VF{-3.000}{3.000}{-0.940}{-0.313}{0.660}
\VF{-2.000}{3.000}{-1.067}{0.300}{0.789}
\VF{-1.000}{3.000}{-0.742}{0.955}{0.900}
\VF{0.000}{3.000}{0.000}{1.251}{0.946}
\VF{1.000}{3.000}{0.742}{0.955}{0.900}
\VF{2.000}{3.000}{1.067}{0.300}{0.789}
\VF{3.000}{3.000}{0.940}{-0.313}{0.660}
\VF{4.000}{3.000}{0.587}{-0.661}{0.543}
\VF{5.000}{3.000}{0.243}{-0.756}{0.445}
\VF{6.000}{3.000}{-0.000}{-0.719}{0.361}
\VF{7.000}{3.000}{-0.149}{-0.635}{0.288}
\VF{8.000}{3.000}{-0.234}{-0.545}{0.223}
\VF{9.000}{3.000}{-0.277}{-0.462}{0.164}
\VF{10.000}{3.000}{-0.296}{-0.391}{0.110}
\VF{-10.000}{4.000}{0.150}{-0.445}{0.088}
\VF{-9.000}{4.000}{0.092}{-0.508}{0.138}
\VF{-8.000}{4.000}{0.000}{-0.566}{0.194}
\VF{-7.000}{4.000}{-0.138}{-0.607}{0.255}
\VF{-6.000}{4.000}{-0.331}{-0.600}{0.324}
\VF{-5.000}{4.000}{-0.568}{-0.501}{0.403}
\VF{-4.000}{4.000}{-0.795}{-0.265}{0.493}
\VF{-3.000}{4.000}{-0.920}{0.112}{0.590}
\VF{-2.000}{4.000}{-0.841}{0.561}{0.683}
\VF{-1.000}{4.000}{-0.513}{0.944}{0.752}
\VF{0.000}{4.000}{0.000}{1.099}{0.779}
\VF{1.000}{4.000}{0.513}{0.944}{0.752}
\VF{2.000}{4.000}{0.841}{0.561}{0.683}
\VF{3.000}{4.000}{0.920}{0.112}{0.590}
\VF{4.000}{4.000}{0.795}{-0.265}{0.493}
\VF{5.000}{4.000}{0.568}{-0.501}{0.403}
\VF{6.000}{4.000}{0.331}{-0.600}{0.324}
\VF{7.000}{4.000}{0.138}{-0.607}{0.255}
\VF{8.000}{4.000}{-0.000}{-0.566}{0.194}
\VF{9.000}{4.000}{-0.092}{-0.508}{0.138}
\VF{10.000}{4.000}{-0.150}{-0.445}{0.088}
\VF{-10.000}{5.000}{0.000}{-0.448}{0.064}
\VF{-9.000}{5.000}{-0.086}{-0.485}{0.113}
\VF{-8.000}{5.000}{-0.203}{-0.501}{0.166}
\VF{-7.000}{5.000}{-0.352}{-0.479}{0.225}
\VF{-6.000}{5.000}{-0.522}{-0.394}{0.290}
\VF{-5.000}{5.000}{-0.683}{-0.228}{0.363}
\VF{-4.000}{5.000}{-0.790}{0.024}{0.440}
\VF{-3.000}{5.000}{-0.793}{0.334}{0.517}
\VF{-2.000}{5.000}{-0.655}{0.649}{0.585}
\VF{-1.000}{5.000}{-0.375}{0.889}{0.632}
\VF{0.000}{5.000}{0.000}{0.981}{0.649}
\VF{1.000}{5.000}{0.375}{0.889}{0.632}
\VF{2.000}{5.000}{0.655}{0.649}{0.585}
\VF{3.000}{5.000}{0.793}{0.334}{0.517}
\VF{4.000}{5.000}{0.790}{0.024}{0.440}
\VF{5.000}{5.000}{0.683}{-0.228}{0.363}
\VF{6.000}{5.000}{0.522}{-0.394}{0.290}
\VF{7.000}{5.000}{0.352}{-0.479}{0.225}
\VF{8.000}{5.000}{0.203}{-0.501}{0.166}
\VF{9.000}{5.000}{0.086}{-0.485}{0.113}
\VF{10.000}{5.000}{-0.000}{-0.448}{0.064}
\VF{-10.000}{6.000}{-0.131}{-0.407}{0.041}
\VF{-9.000}{6.000}{-0.228}{-0.412}{0.088}
\VF{-8.000}{6.000}{-0.344}{-0.387}{0.140}
\VF{-7.000}{6.000}{-0.471}{-0.319}{0.196}
\VF{-6.000}{6.000}{-0.592}{-0.197}{0.257}
\VF{-5.000}{6.000}{-0.682}{-0.021}{0.321}
\VF{-4.000}{6.000}{-0.714}{0.201}{0.386}
\VF{-3.000}{6.000}{-0.663}{0.442}{0.447}
\VF{-2.000}{6.000}{-0.517}{0.665}{0.497}
\VF{-1.000}{6.000}{-0.285}{0.825}{0.531}
\VF{0.000}{6.000}{0.000}{0.884}{0.543}
\VF{1.000}{6.000}{0.285}{0.825}{0.531}
\VF{2.000}{6.000}{0.517}{0.665}{0.497}
\VF{3.000}{6.000}{0.663}{0.442}{0.447}
\VF{4.000}{6.000}{0.714}{0.201}{0.386}
\VF{5.000}{6.000}{0.682}{-0.021}{0.321}
\VF{6.000}{6.000}{0.592}{-0.197}{0.257}
\VF{7.000}{6.000}{0.471}{-0.319}{0.196}
\VF{8.000}{6.000}{0.344}{-0.387}{0.140}
\VF{9.000}{6.000}{0.228}{-0.412}{0.088}
\VF{10.000}{6.000}{0.131}{-0.407}{0.041}
\VF{-10.000}{7.000}{-0.229}{-0.338}{0.020}
\VF{-9.000}{7.000}{-0.321}{-0.315}{0.066}
\VF{-8.000}{7.000}{-0.420}{-0.261}{0.115}
\VF{-7.000}{7.000}{-0.514}{-0.171}{0.167}
\VF{-6.000}{7.000}{-0.591}{-0.043}{0.222}
\VF{-5.000}{7.000}{-0.632}{0.119}{0.278}
\VF{-4.000}{7.000}{-0.623}{0.303}{0.332}
\VF{-3.000}{7.000}{-0.551}{0.488}{0.380}
\VF{-2.000}{7.000}{-0.415}{0.650}{0.419}
\VF{-1.000}{7.000}{-0.224}{0.762}{0.445}
\VF{0.000}{7.000}{0.000}{0.803}{0.453}
\VF{1.000}{7.000}{0.224}{0.762}{0.445}
\VF{2.000}{7.000}{0.415}{0.650}{0.419}
\VF{3.000}{7.000}{0.551}{0.488}{0.380}
\VF{4.000}{7.000}{0.623}{0.303}{0.332}
\VF{5.000}{7.000}{0.632}{0.119}{0.278}
\VF{6.000}{7.000}{0.591}{-0.043}{0.222}
\VF{7.000}{7.000}{0.514}{-0.171}{0.167}
\VF{8.000}{7.000}{0.420}{-0.261}{0.115}
\VF{9.000}{7.000}{0.321}{-0.315}{0.066}
\VF{10.000}{7.000}{0.229}{-0.338}{0.020}
\VF{-10.000}{8.000}{-0.292}{-0.258}{0.000}
\VF{-9.000}{8.000}{-0.371}{-0.217}{0.044}
\VF{-8.000}{8.000}{-0.447}{-0.149}{0.090}
\VF{-7.000}{8.000}{-0.513}{-0.054}{0.138}
\VF{-6.000}{8.000}{-0.556}{0.067}{0.187}
\VF{-5.000}{8.000}{-0.567}{0.208}{0.235}
\VF{-4.000}{8.000}{-0.536}{0.357}{0.280}
\VF{-3.000}{8.000}{-0.460}{0.501}{0.319}
\VF{-2.000}{8.000}{-0.338}{0.622}{0.349}
\VF{-1.000}{8.000}{-0.179}{0.703}{0.369}
\VF{0.000}{8.000}{0.000}{0.732}{0.376}
\VF{1.000}{8.000}{0.179}{0.703}{0.369}
\VF{2.000}{8.000}{0.338}{0.622}{0.349}
\VF{3.000}{8.000}{0.460}{0.501}{0.319}
\VF{4.000}{8.000}{0.536}{0.357}{0.280}
\VF{5.000}{8.000}{0.567}{0.208}{0.235}
\VF{6.000}{8.000}{0.556}{0.067}{0.187}
\VF{7.000}{8.000}{0.513}{-0.054}{0.138}
\VF{8.000}{8.000}{0.447}{-0.149}{0.090}
\VF{9.000}{8.000}{0.371}{-0.217}{0.044}
\VF{10.000}{8.000}{0.292}{-0.258}{0.000}
\VF{-10.000}{9.000}{-0.341}{-0.189}{0.000}
\VF{-9.000}{9.000}{-0.388}{-0.129}{0.021}
\VF{-8.000}{9.000}{-0.444}{-0.058}{0.064}
\VF{-7.000}{9.000}{-0.486}{0.034}{0.107}
\VF{-6.000}{9.000}{-0.507}{0.143}{0.150}
\VF{-5.000}{9.000}{-0.500}{0.262}{0.192}
\VF{-4.000}{9.000}{-0.460}{0.383}{0.229}
\VF{-3.000}{9.000}{-0.386}{0.496}{0.262}
\VF{-2.000}{9.000}{-0.279}{0.588}{0.286}
\VF{-1.000}{9.000}{-0.146}{0.648}{0.302}
\VF{0.000}{9.000}{0.000}{0.670}{0.307}
\VF{1.000}{9.000}{0.146}{0.648}{0.302}
\VF{2.000}{9.000}{0.279}{0.588}{0.286}
\VF{3.000}{9.000}{0.386}{0.496}{0.262}
\VF{4.000}{9.000}{0.460}{0.383}{0.229}
\VF{5.000}{9.000}{0.500}{0.262}{0.192}
\VF{6.000}{9.000}{0.507}{0.143}{0.150}
\VF{7.000}{9.000}{0.486}{0.034}{0.107}
\VF{8.000}{9.000}{0.444}{-0.058}{0.064}
\VF{9.000}{9.000}{0.388}{-0.129}{0.021}
\VF{10.000}{9.000}{0.341}{-0.189}{0.000}
\VF{-10.000}{10.000}{-0.370}{-0.123}{0.000}
\VF{-9.000}{10.000}{-0.386}{-0.058}{0.000}
\VF{-8.000}{10.000}{-0.424}{0.013}{0.037}
\VF{-7.000}{10.000}{-0.449}{0.097}{0.076}
\VF{-6.000}{10.000}{-0.455}{0.192}{0.114}
\VF{-5.000}{10.000}{-0.438}{0.292}{0.150}
\VF{-4.000}{10.000}{-0.395}{0.391}{0.182}
\VF{-3.000}{10.000}{-0.325}{0.480}{0.209}
\VF{-2.000}{10.000}{-0.232}{0.551}{0.229}
\VF{-1.000}{10.000}{-0.121}{0.598}{0.242}
\VF{0.000}{10.000}{0.000}{0.614}{0.246}
\VF{1.000}{10.000}{0.121}{0.598}{0.242}
\VF{2.000}{10.000}{0.232}{0.551}{0.229}
\VF{3.000}{10.000}{0.325}{0.480}{0.209}
\VF{4.000}{10.000}{0.395}{0.391}{0.182}
\VF{5.000}{10.000}{0.438}{0.292}{0.150}
\VF{6.000}{10.000}{0.455}{0.192}{0.114}
\VF{7.000}{10.000}{0.449}{0.097}{0.076}
\VF{8.000}{10.000}{0.424}{0.013}{0.037}
\VF{9.000}{10.000}{0.386}{-0.058}{0.000}
\VF{10.000}{10.000}{0.370}{-0.123}{0.000}
}
  \marcoMMA{10}{5}{$x$ (u.\,arb.)}{$z$ (u.\,arb.)}
\end{tikzpicture}\\[1pt]{\small\scshape (c)}
\end{minipage}
\caption{Término cuadrupolar con simetría de revolución, $\phimul{2}\propto Q_{zz}\Pl{2}(\cos\theta)/r^{3}$ con $Q_{zz}>0$ (\cref{teo:revolucion}). (a)~Equipotenciales: los
lóbulos polares (potencial positivo) son más extensos que los ecuatoriales (negativo) porque $\Pl{2}(\pm1)=1$ y $\Pl{2}(0)=-\frac12$. (b)~Líneas de campo: salen a lo largo
del eje $z$, entran por el plano ecuatorial y cruzan el cono nodal. (c)~Campo eléctrico, con la misma escala de color que la \cref{fig:dipolo-mma}: cae como $1/r^{4}$,
aún más deprisa que el dipolar. El icono rojo representa dos dipolos opuestos. Recreación de~\cite{FIGURAS-T2}.}
\label{fig:cuadrupolo-mma}
\end{figure}


\begin{ejemplo}[Cuadrupolo lineal]\label{ej:cuadrupolo-lineal}
Sean las cargas $+q$ en $z=\pm d$ y $-2q$ en el origen (dos dipolos opuestos, de momentos $\pm qd\,\uz$, unidos por su carga negativa). Entonces
$Q=q+q-2q=0$, $\vp=q\,d\,\uz+q(-d)\uz+(-2q)\vc{0}=\vc{0}$, y el sistema tiene simetría de revolución alrededor de $z$, con
\[
  Q_{zz}=\sum_iq_i\bigl(3z_i^{2}-r_i^{2}\bigr)=q\,(3d^{2}-d^{2})+q\,(3d^{2}-d^{2})+(-2q)\cdot0=4qd^{2}>0 .
\]
El primer término no nulo es el cuadrupolar, independiente del origen (\cref{prop:Q-origen}):
\[
  \phi(\vc{r})\simeq\phimul{2}(\vc{r})=\frac{4qd^{2}}{16\pi\eo r^{3}}\bigl(3\cos^{2}\theta-1\bigr)=\frac{qd^{2}}{4\pi\eo r^{3}}\bigl(3\cos^{2}\theta-1\bigr).
\]
La carga positiva está en los extremos del eje, «alargada» según $z$, y en efecto $Q_{zz}>0$: el potencial es positivo dentro del cono nodal que rodea al eje y
negativo cerca del plano ecuatorial (\cref{fig:cuadrupolo-fp}).
\end{ejemplo}

% =========================================================================
% Figura (clon de «Tema 2. Figuras»): cuadrupolo puntual frente a cuadrupolo
% físico lineal (+q en z = +-d, -2q en el origen, con 2qd^2 = 1 y d = 1/2).
% Datos: figuras/datos/cua_{puntual,fisico}_*.tex.
% =========================================================================
\begin{figure}[ht]
\centering
\begin{minipage}[t]{0.49\textwidth}\centering
\begin{tikzpicture}[x=0.29cm, y=0.29cm]
  \panelMMA{10}{mmafondo}{1.2}{%
    % Cuadrupolo puntual - equipotenciales
% Generado por scripts/generar_datos.py (no editar a mano).
\EQ{(-0.220,-0.155) (-0.275,-0.195) (-0.319,-0.225) (-0.375,-0.266) (-0.425,-0.300) (-0.475,-0.336) (-0.525,-0.371) (-0.568,-0.400) (-0.625,-0.441) (-0.675,-0.475) (-0.725,-0.510) (-0.775,-0.544) (-0.821,-0.575) (-0.875,-0.612) (-0.925,-0.646) (-0.969,-0.675) (-1.025,-0.713) (-1.075,-0.745) (-1.121,-0.775) (-1.175,-0.810) (-1.225,-0.841) (-1.275,-0.872) (-1.321,-0.900) (-1.375,-0.933) (-1.425,-0.963) (-1.475,-0.991) (-1.525,-1.020) (-1.575,-1.047) (-1.625,-1.074) (-1.674,-1.100) (-1.723,-1.125) (-1.774,-1.150) (-1.825,-1.175) (-1.875,-1.198) (-1.925,-1.220) (-1.975,-1.241) (-2.025,-1.262) (-2.075,-1.281) (-2.125,-1.300) (-2.175,-1.317) (-2.225,-1.334) (-2.275,-1.349) (-2.325,-1.364) (-2.375,-1.377) (-2.425,-1.389) (-2.475,-1.400) (-2.525,-1.410) (-2.575,-1.419) (-2.625,-1.426) (-2.675,-1.433) (-2.725,-1.438) (-2.775,-1.441) (-2.825,-1.444) (-2.875,-1.445) (-2.925,-1.445) (-2.975,-1.443) (-3.025,-1.440) (-3.075,-1.435) (-3.125,-1.429) (-3.175,-1.421) (-3.225,-1.412) (-3.275,-1.401) (-3.325,-1.388) (-3.375,-1.373) (-3.425,-1.357) (-3.475,-1.338) (-3.525,-1.318) (-3.575,-1.295) (-3.625,-1.270) (-3.675,-1.242) (-3.725,-1.212) (-3.775,-1.179) (-3.825,-1.142) (-3.875,-1.102) (-3.925,-1.058) (-3.975,-1.010) (-4.025,-0.956) (-4.072,-0.900) (-4.110,-0.850) (-4.145,-0.800) (-4.175,-0.753) (-4.206,-0.700) (-4.232,-0.650) (-4.256,-0.600) (-4.277,-0.550) (-4.297,-0.500) (-4.314,-0.450) (-4.329,-0.400) (-4.343,-0.350) (-4.354,-0.300) (-4.364,-0.250) (-4.371,-0.200) (-4.377,-0.150) (-4.382,-0.100) (-4.384,-0.050) (-4.385,0.000) (-4.384,0.050) (-4.382,0.100) (-4.377,0.150) (-4.371,0.200) (-4.364,0.250) (-4.354,0.300) (-4.343,0.350) (-4.329,0.400) (-4.314,0.450) (-4.297,0.500) (-4.277,0.550) (-4.256,0.600) (-4.232,0.650) (-4.206,0.700) (-4.177,0.750) (-4.150,0.792) (-4.110,0.850) (-4.075,0.896) (-4.030,0.950) (-3.985,1.000) (-3.934,1.050) (-3.878,1.100) (-3.825,1.142) (-3.780,1.175) (-3.725,1.212) (-3.675,1.242) (-3.625,1.270) (-3.575,1.295) (-3.525,1.318) (-3.475,1.338) (-3.425,1.357) (-3.375,1.373) (-3.325,1.388) (-3.275,1.401) (-3.225,1.412) (-3.175,1.421) (-3.125,1.429) (-3.075,1.435) (-3.025,1.440) (-2.975,1.443) (-2.925,1.445) (-2.875,1.445) (-2.825,1.444) (-2.775,1.441) (-2.725,1.438) (-2.675,1.433) (-2.625,1.426) (-2.575,1.419) (-2.525,1.410) (-2.475,1.400) (-2.425,1.389) (-2.375,1.377) (-2.325,1.364) (-2.275,1.349) (-2.225,1.334) (-2.175,1.317) (-2.126,1.300) (-2.075,1.281) (-2.025,1.262) (-1.975,1.241) (-1.925,1.220) (-1.875,1.198) (-1.826,1.175) (-1.775,1.151) (-1.725,1.126) (-1.675,1.100) (-1.627,1.075) (-1.580,1.050) (-1.535,1.025) (-1.491,1.000) (-1.447,0.975) (-1.400,0.948) (-1.350,0.918) (-1.300,0.888) (-1.250,0.857) (-1.200,0.825) (-1.150,0.794) (-1.100,0.762) (-1.050,0.729) (-1.007,0.700) (-0.950,0.663) (-0.900,0.629) (-0.858,0.600) (-0.800,0.561) (-0.750,0.527) (-0.700,0.493) (-0.650,0.458) (-0.604,0.425) (-0.550,0.388) (-0.500,0.353) (-0.450,0.318) (-0.400,0.283) (-0.354,0.250) (-0.300,0.213) (-0.250,0.177) (-0.220,0.155)}
\EQ{(0.206,0.144) (0.248,0.175) (0.300,0.213) (0.350,0.248) (0.400,0.283) (0.450,0.318) (0.496,0.350) (0.550,0.388) (0.600,0.423) (0.650,0.458) (0.700,0.493) (0.748,0.525) (0.800,0.561) (0.850,0.595) (0.895,0.625) (0.950,0.663) (1.000,0.696) (1.045,0.725) (1.100,0.762) (1.150,0.794) (1.200,0.825) (1.250,0.857) (1.300,0.888) (1.350,0.918) (1.400,0.948) (1.447,0.975) (1.491,1.000) (1.535,1.025) (1.580,1.050) (1.625,1.074) (1.674,1.100) (1.723,1.125) (1.774,1.150) (1.825,1.175) (1.875,1.198) (1.925,1.220) (1.975,1.241) (2.025,1.262) (2.075,1.281) (2.125,1.300) (2.175,1.317) (2.225,1.334) (2.275,1.349) (2.325,1.364) (2.375,1.377) (2.425,1.389) (2.475,1.400) (2.525,1.410) (2.575,1.419) (2.625,1.426) (2.675,1.433) (2.725,1.438) (2.775,1.441) (2.825,1.444) (2.875,1.445) (2.925,1.445) (2.975,1.443) (3.025,1.440) (3.075,1.435) (3.125,1.429) (3.175,1.421) (3.225,1.412) (3.275,1.401) (3.325,1.388) (3.375,1.373) (3.425,1.357) (3.475,1.338) (3.525,1.318) (3.575,1.295) (3.625,1.270) (3.675,1.242) (3.725,1.212) (3.775,1.179) (3.825,1.142) (3.875,1.102) (3.925,1.058) (3.975,1.010) (4.025,0.956) (4.072,0.900) (4.110,0.850) (4.145,0.800) (4.175,0.753) (4.206,0.700) (4.232,0.650) (4.256,0.600) (4.277,0.550) (4.297,0.500) (4.314,0.450) (4.329,0.400) (4.343,0.350) (4.354,0.300) (4.364,0.250) (4.371,0.200) (4.377,0.150) (4.382,0.100) (4.384,0.050) (4.385,0.000) (4.384,-0.050) (4.382,-0.100) (4.377,-0.150) (4.371,-0.200) (4.364,-0.250) (4.354,-0.300) (4.343,-0.350) (4.329,-0.400) (4.314,-0.450) (4.297,-0.500) (4.277,-0.550) (4.256,-0.600) (4.232,-0.650) (4.206,-0.700) (4.177,-0.750) (4.150,-0.792) (4.110,-0.850) (4.075,-0.896) (4.030,-0.950) (3.985,-1.000) (3.934,-1.050) (3.878,-1.100) (3.825,-1.142) (3.780,-1.175) (3.725,-1.212) (3.675,-1.242) (3.625,-1.270) (3.575,-1.295) (3.525,-1.318) (3.475,-1.338) (3.425,-1.357) (3.375,-1.373) (3.325,-1.388) (3.275,-1.401) (3.225,-1.412) (3.175,-1.421) (3.125,-1.429) (3.075,-1.435) (3.025,-1.440) (2.975,-1.443) (2.925,-1.445) (2.875,-1.445) (2.825,-1.444) (2.775,-1.441) (2.725,-1.438) (2.675,-1.433) (2.625,-1.426) (2.575,-1.419) (2.525,-1.410) (2.475,-1.400) (2.425,-1.389) (2.375,-1.377) (2.325,-1.364) (2.275,-1.349) (2.225,-1.334) (2.175,-1.317) (2.126,-1.300) (2.075,-1.281) (2.025,-1.262) (1.975,-1.241) (1.925,-1.220) (1.875,-1.198) (1.826,-1.175) (1.775,-1.151) (1.725,-1.126) (1.675,-1.100) (1.627,-1.075) (1.580,-1.050) (1.535,-1.025) (1.491,-1.000) (1.447,-0.975) (1.400,-0.948) (1.350,-0.918) (1.300,-0.888) (1.250,-0.857) (1.200,-0.825) (1.150,-0.794) (1.100,-0.762) (1.050,-0.729) (1.007,-0.700) (0.950,-0.663) (0.900,-0.629) (0.858,-0.600) (0.800,-0.561) (0.750,-0.527) (0.700,-0.493) (0.650,-0.458) (0.604,-0.425) (0.550,-0.388) (0.500,-0.353) (0.450,-0.318) (0.400,-0.283) (0.354,-0.250) (0.300,-0.213) (0.250,-0.177) (0.206,-0.144)}
\EQ{(-0.220,-0.155) (-0.275,-0.195) (-0.319,-0.225) (-0.375,-0.266) (-0.425,-0.300) (-0.475,-0.336) (-0.525,-0.371) (-0.568,-0.400) (-0.625,-0.441) (-0.675,-0.475) (-0.725,-0.510) (-0.775,-0.545) (-0.820,-0.575) (-0.875,-0.613) (-0.925,-0.647) (-0.968,-0.675) (-1.025,-0.714) (-1.075,-0.747) (-1.119,-0.775) (-1.175,-0.812) (-1.225,-0.844) (-1.275,-0.875) (-1.325,-0.906) (-1.375,-0.937) (-1.425,-0.967) (-1.475,-0.996) (-1.525,-1.025) (-1.569,-1.050) (-1.614,-1.075) (-1.660,-1.100) (-1.707,-1.125) (-1.755,-1.150) (-1.800,-1.173) (-1.850,-1.197) (-1.900,-1.221) (-1.950,-1.244) (-2.000,-1.266) (-2.050,-1.288) (-2.100,-1.308) (-2.150,-1.327) (-2.200,-1.346) (-2.250,-1.363) (-2.300,-1.380) (-2.350,-1.396) (-2.400,-1.410) (-2.450,-1.424) (-2.500,-1.436) (-2.550,-1.447) (-2.600,-1.458) (-2.650,-1.467) (-2.700,-1.475) (-2.750,-1.481) (-2.800,-1.487) (-2.850,-1.491) (-2.900,-1.494) (-2.950,-1.496) (-3.000,-1.497) (-3.050,-1.496) (-3.100,-1.493) (-3.150,-1.490) (-3.200,-1.485) (-3.250,-1.478) (-3.300,-1.470) (-3.350,-1.460) (-3.400,-1.449) (-3.450,-1.436) (-3.500,-1.421) (-3.550,-1.405) (-3.600,-1.386) (-3.650,-1.365) (-3.700,-1.343) (-3.750,-1.318) (-3.800,-1.290) (-3.850,-1.260) (-3.900,-1.228) (-3.950,-1.192) (-4.000,-1.153) (-4.050,-1.110) (-4.088,-1.075) (-4.125,-1.038) (-4.160,-1.000) (-4.200,-0.954) (-4.243,-0.900) (-4.275,-0.855) (-4.312,-0.800) (-4.342,-0.750) (-4.369,-0.700) (-4.395,-0.650) (-4.417,-0.600) (-4.438,-0.550) (-4.457,-0.500) (-4.473,-0.450) (-4.488,-0.400) (-4.501,-0.350) (-4.512,-0.300) (-4.521,-0.250) (-4.529,-0.200) (-4.534,-0.150) (-4.539,-0.100) (-4.541,-0.050) (-4.542,0.000) (-4.541,0.050) (-4.539,0.100) (-4.534,0.150) (-4.529,0.200) (-4.521,0.250) (-4.512,0.300) (-4.501,0.350) (-4.488,0.400) (-4.473,0.450) (-4.457,0.500) (-4.438,0.550) (-4.417,0.600) (-4.395,0.650) (-4.369,0.700) (-4.342,0.750) (-4.312,0.800) (-4.279,0.850) (-4.243,0.900) (-4.203,0.950) (-4.160,1.000) (-4.125,1.038) (-4.088,1.075) (-4.050,1.110) (-4.003,1.150) (-3.950,1.192) (-3.904,1.225) (-3.850,1.260) (-3.800,1.290) (-3.750,1.318) (-3.700,1.343) (-3.650,1.365) (-3.600,1.386) (-3.550,1.405) (-3.500,1.421) (-3.450,1.436) (-3.400,1.449) (-3.350,1.460) (-3.300,1.470) (-3.250,1.478) (-3.200,1.485) (-3.150,1.490) (-3.100,1.493) (-3.050,1.496) (-3.000,1.497) (-2.950,1.496) (-2.900,1.494) (-2.850,1.491) (-2.800,1.487) (-2.750,1.481) (-2.700,1.475) (-2.650,1.467) (-2.600,1.458) (-2.550,1.447) (-2.500,1.436) (-2.450,1.424) (-2.400,1.410) (-2.350,1.396) (-2.300,1.380) (-2.250,1.363) (-2.200,1.346) (-2.150,1.327) (-2.100,1.308) (-2.050,1.288) (-2.000,1.266) (-1.950,1.244) (-1.900,1.221) (-1.850,1.197) (-1.804,1.175) (-1.755,1.150) (-1.707,1.125) (-1.660,1.100) (-1.614,1.075) (-1.569,1.050) (-1.525,1.025) (-1.482,1.000) (-1.425,0.967) (-1.375,0.937) (-1.325,0.906) (-1.275,0.875) (-1.225,0.844) (-1.175,0.812) (-1.125,0.779) (-1.081,0.750) (-1.025,0.714) (-0.975,0.680) (-0.930,0.650) (-0.875,0.613) (-0.825,0.579) (-0.784,0.550) (-0.725,0.510) (-0.675,0.475) (-0.625,0.441) (-0.575,0.406) (-0.532,0.375) (-0.475,0.336) (-0.425,0.300) (-0.375,0.266) (-0.325,0.230) (-0.284,0.200) (-0.225,0.160) (-0.220,0.155)}
\EQ{(0.206,0.144) (0.248,0.175) (0.300,0.213) (0.350,0.248) (0.400,0.283) (0.450,0.318) (0.496,0.350) (0.550,0.388) (0.600,0.423) (0.650,0.458) (0.700,0.493) (0.747,0.525) (0.800,0.562) (0.850,0.596) (0.894,0.625) (0.950,0.664) (1.000,0.697) (1.043,0.725) (1.100,0.763) (1.150,0.796) (1.196,0.825) (1.250,0.859) (1.300,0.891) (1.350,0.921) (1.397,0.950) (1.450,0.982) (1.500,1.011) (1.550,1.040) (1.600,1.068) (1.650,1.095) (1.700,1.122) (1.750,1.148) (1.800,1.173) (1.850,1.197) (1.900,1.221) (1.950,1.244) (2.000,1.266) (2.050,1.288) (2.100,1.308) (2.150,1.327) (2.200,1.346) (2.250,1.363) (2.300,1.380) (2.350,1.396) (2.400,1.410) (2.450,1.424) (2.500,1.436) (2.550,1.447) (2.600,1.458) (2.650,1.467) (2.700,1.475) (2.750,1.481) (2.800,1.487) (2.850,1.491) (2.900,1.494) (2.950,1.496) (3.000,1.497) (3.050,1.496) (3.100,1.493) (3.150,1.490) (3.200,1.485) (3.250,1.478) (3.300,1.470) (3.350,1.460) (3.400,1.449) (3.450,1.436) (3.500,1.421) (3.550,1.405) (3.600,1.386) (3.650,1.365) (3.700,1.343) (3.750,1.318) (3.800,1.290) (3.850,1.260) (3.900,1.228) (3.950,1.192) (4.000,1.153) (4.050,1.110) (4.088,1.075) (4.125,1.038) (4.160,1.000) (4.200,0.954) (4.243,0.900) (4.275,0.855) (4.312,0.800) (4.342,0.750) (4.369,0.700) (4.395,0.650) (4.417,0.600) (4.438,0.550) (4.457,0.500) (4.473,0.450) (4.488,0.400) (4.501,0.350) (4.512,0.300) (4.521,0.250) (4.529,0.200) (4.534,0.150) (4.539,0.100) (4.541,0.050) (4.542,0.000) (4.541,-0.050) (4.539,-0.100) (4.534,-0.150) (4.529,-0.200) (4.521,-0.250) (4.512,-0.300) (4.501,-0.350) (4.488,-0.400) (4.473,-0.450) (4.457,-0.500) (4.438,-0.550) (4.417,-0.600) (4.395,-0.650) (4.369,-0.700) (4.342,-0.750) (4.312,-0.800) (4.279,-0.850) (4.243,-0.900) (4.203,-0.950) (4.160,-1.000) (4.125,-1.038) (4.088,-1.075) (4.050,-1.110) (4.003,-1.150) (3.950,-1.192) (3.904,-1.225) (3.850,-1.260) (3.800,-1.290) (3.750,-1.318) (3.700,-1.343) (3.650,-1.365) (3.600,-1.386) (3.550,-1.405) (3.500,-1.421) (3.450,-1.436) (3.400,-1.449) (3.350,-1.460) (3.300,-1.470) (3.250,-1.478) (3.200,-1.485) (3.150,-1.490) (3.100,-1.493) (3.050,-1.496) (3.000,-1.497) (2.950,-1.496) (2.900,-1.494) (2.850,-1.491) (2.800,-1.487) (2.750,-1.481) (2.700,-1.475) (2.650,-1.467) (2.600,-1.458) (2.550,-1.447) (2.500,-1.436) (2.450,-1.424) (2.400,-1.410) (2.350,-1.396) (2.300,-1.380) (2.250,-1.363) (2.200,-1.346) (2.150,-1.327) (2.100,-1.308) (2.050,-1.288) (2.000,-1.266) (1.950,-1.244) (1.900,-1.221) (1.850,-1.197) (1.804,-1.175) (1.755,-1.150) (1.707,-1.125) (1.660,-1.100) (1.614,-1.075) (1.569,-1.050) (1.525,-1.025) (1.482,-1.000) (1.425,-0.967) (1.375,-0.937) (1.325,-0.906) (1.275,-0.875) (1.225,-0.844) (1.175,-0.812) (1.125,-0.779) (1.081,-0.750) (1.025,-0.714) (0.975,-0.680) (0.930,-0.650) (0.875,-0.613) (0.825,-0.579) (0.784,-0.550) (0.725,-0.510) (0.675,-0.475) (0.625,-0.441) (0.575,-0.406) (0.532,-0.375) (0.475,-0.336) (0.425,-0.300) (0.375,-0.266) (0.325,-0.230) (0.284,-0.200) (0.225,-0.160) (0.206,-0.144)}
\EQ{(-0.220,-0.155) (-0.275,-0.195) (-0.319,-0.225) (-0.375,-0.266) (-0.425,-0.300) (-0.475,-0.336) (-0.525,-0.371) (-0.568,-0.400) (-0.625,-0.441) (-0.674,-0.475) (-0.725,-0.510) (-0.775,-0.545) (-0.819,-0.575) (-0.875,-0.614) (-0.925,-0.647) (-0.975,-0.681) (-1.025,-0.715) (-1.075,-0.748) (-1.125,-0.781) (-1.175,-0.814) (-1.225,-0.846) (-1.271,-0.875) (-1.325,-0.910) (-1.375,-0.941) (-1.425,-0.971) (-1.473,-1.000) (-1.525,-1.031) (-1.575,-1.060) (-1.625,-1.089) (-1.675,-1.116) (-1.725,-1.144) (-1.775,-1.170) (-1.825,-1.196) (-1.875,-1.221) (-1.925,-1.246) (-1.975,-1.270) (-2.025,-1.293) (-2.075,-1.315) (-2.125,-1.336) (-2.175,-1.357) (-2.222,-1.375) (-2.275,-1.395) (-2.325,-1.413) (-2.375,-1.430) (-2.425,-1.446) (-2.475,-1.461) (-2.525,-1.474) (-2.575,-1.487) (-2.625,-1.499) (-2.675,-1.510) (-2.725,-1.520) (-2.775,-1.529) (-2.825,-1.536) (-2.875,-1.542) (-2.925,-1.548) (-2.975,-1.552) (-3.025,-1.555) (-3.075,-1.556) (-3.125,-1.556) (-3.175,-1.556) (-3.225,-1.553) (-3.275,-1.550) (-3.325,-1.545) (-3.375,-1.538) (-3.425,-1.530) (-3.475,-1.521) (-3.525,-1.510) (-3.575,-1.497) (-3.625,-1.483) (-3.675,-1.467) (-3.725,-1.449) (-3.775,-1.429) (-3.825,-1.407) (-3.875,-1.383) (-3.925,-1.357) (-3.975,-1.329) (-4.021,-1.300) (-4.075,-1.264) (-4.125,-1.227) (-4.175,-1.187) (-4.217,-1.150) (-4.270,-1.100) (-4.318,-1.050) (-4.362,-1.000) (-4.400,-0.952) (-4.439,-0.900) (-4.473,-0.850) (-4.500,-0.807) (-4.533,-0.750) (-4.559,-0.700) (-4.583,-0.650) (-4.605,-0.600) (-4.624,-0.550) (-4.642,-0.500) (-4.658,-0.450) (-4.672,-0.400) (-4.684,-0.350) (-4.695,-0.300) (-4.704,-0.250) (-4.711,-0.200) (-4.717,-0.150) (-4.720,-0.100) (-4.723,-0.050) (-4.724,0.000) (-4.723,0.050) (-4.720,0.100) (-4.717,0.150) (-4.711,0.200) (-4.704,0.250) (-4.695,0.300) (-4.684,0.350) (-4.672,0.400) (-4.658,0.450) (-4.642,0.500) (-4.625,0.548) (-4.605,0.600) (-4.583,0.650) (-4.559,0.700) (-4.533,0.750) (-4.504,0.800) (-4.475,0.847) (-4.439,0.900) (-4.402,0.950) (-4.362,1.000) (-4.325,1.042) (-4.275,1.095) (-4.225,1.143) (-4.175,1.187) (-4.127,1.225) (-4.075,1.264) (-4.025,1.298) (-3.981,1.325) (-3.925,1.357) (-3.875,1.383) (-3.825,1.407) (-3.775,1.429) (-3.725,1.449) (-3.675,1.467) (-3.625,1.483) (-3.575,1.497) (-3.525,1.510) (-3.475,1.521) (-3.425,1.530) (-3.375,1.538) (-3.325,1.545) (-3.275,1.550) (-3.225,1.553) (-3.175,1.556) (-3.125,1.556) (-3.075,1.556) (-3.025,1.555) (-2.975,1.552) (-2.925,1.548) (-2.875,1.542) (-2.825,1.536) (-2.775,1.529) (-2.725,1.520) (-2.675,1.510) (-2.625,1.499) (-2.575,1.487) (-2.525,1.474) (-2.475,1.461) (-2.425,1.446) (-2.375,1.430) (-2.325,1.413) (-2.275,1.395) (-2.225,1.376) (-2.175,1.357) (-2.125,1.336) (-2.075,1.315) (-2.025,1.293) (-1.975,1.270) (-1.925,1.246) (-1.875,1.221) (-1.825,1.196) (-1.775,1.170) (-1.725,1.144) (-1.675,1.116) (-1.625,1.089) (-1.575,1.060) (-1.525,1.031) (-1.475,1.001) (-1.432,0.975) (-1.375,0.941) (-1.325,0.910) (-1.275,0.878) (-1.232,0.850) (-1.175,0.814) (-1.125,0.781) (-1.078,0.750) (-1.025,0.715) (-0.975,0.681) (-0.929,0.650) (-0.875,0.614) (-0.825,0.579) (-0.783,0.550) (-0.725,0.510) (-0.675,0.475) (-0.625,0.441) (-0.575,0.406) (-0.532,0.375) (-0.475,0.336) (-0.425,0.300) (-0.375,0.266) (-0.325,0.230) (-0.284,0.200) (-0.225,0.160) (-0.220,0.155)}
\EQ{(0.206,0.144) (0.248,0.175) (0.300,0.213) (0.350,0.248) (0.400,0.283) (0.450,0.318) (0.496,0.350) (0.550,0.389) (0.600,0.423) (0.650,0.458) (0.700,0.493) (0.747,0.525) (0.800,0.562) (0.850,0.596) (0.893,0.625) (0.950,0.665) (1.000,0.698) (1.050,0.732) (1.100,0.765) (1.150,0.797) (1.193,0.825) (1.250,0.862) (1.300,0.894) (1.350,0.925) (1.400,0.956) (1.450,0.986) (1.500,1.016) (1.550,1.046) (1.600,1.074) (1.646,1.100) (1.691,1.125) (1.737,1.150) (1.785,1.175) (1.833,1.200) (1.883,1.225) (1.934,1.250) (1.987,1.275) (2.042,1.300) (2.099,1.325) (2.150,1.346) (2.200,1.367) (2.250,1.386) (2.300,1.404) (2.350,1.421) (2.400,1.438) (2.450,1.453) (2.500,1.468) (2.550,1.481) (2.600,1.493) (2.650,1.505) (2.700,1.515) (2.750,1.524) (2.800,1.532) (2.850,1.539) (2.900,1.545) (2.950,1.550) (3.000,1.553) (3.050,1.556) (3.100,1.556) (3.150,1.556) (3.200,1.555) (3.250,1.552) (3.300,1.547) (3.350,1.542) (3.400,1.534) (3.450,1.526) (3.500,1.515) (3.550,1.504) (3.600,1.490) (3.650,1.475) (3.700,1.458) (3.750,1.439) (3.800,1.418) (3.850,1.396) (3.900,1.371) (3.950,1.343) (4.000,1.313) (4.050,1.281) (4.094,1.250) (4.150,1.207) (4.189,1.175) (4.244,1.125) (4.294,1.075) (4.340,1.025) (4.375,0.984) (4.421,0.925) (4.450,0.884) (4.489,0.825) (4.519,0.775) (4.546,0.725) (4.571,0.675) (4.594,0.625) (4.615,0.575) (4.633,0.525) (4.650,0.475) (4.665,0.425) (4.678,0.375) (4.690,0.325) (4.699,0.275) (4.708,0.225) (4.714,0.175) (4.719,0.125) (4.722,0.075) (4.723,0.025) (4.723,-0.050) (4.720,-0.100) (4.717,-0.150) (4.711,-0.200) (4.704,-0.250) (4.695,-0.300) (4.684,-0.350) (4.672,-0.400) (4.658,-0.450) (4.642,-0.500) (4.625,-0.548) (4.605,-0.600) (4.583,-0.650) (4.559,-0.700) (4.533,-0.750) (4.504,-0.800) (4.475,-0.847) (4.439,-0.900) (4.402,-0.950) (4.362,-1.000) (4.325,-1.042) (4.275,-1.095) (4.225,-1.143) (4.175,-1.187) (4.127,-1.225) (4.075,-1.264) (4.025,-1.298) (3.981,-1.325) (3.925,-1.357) (3.875,-1.383) (3.825,-1.407) (3.775,-1.429) (3.725,-1.449) (3.675,-1.467) (3.625,-1.483) (3.575,-1.497) (3.525,-1.510) (3.475,-1.521) (3.425,-1.530) (3.375,-1.538) (3.325,-1.545) (3.275,-1.550) (3.225,-1.553) (3.175,-1.556) (3.125,-1.556) (3.075,-1.556) (3.025,-1.555) (2.975,-1.552) (2.925,-1.548) (2.875,-1.542) (2.825,-1.536) (2.775,-1.529) (2.725,-1.520) (2.675,-1.510) (2.625,-1.499) (2.575,-1.487) (2.525,-1.474) (2.475,-1.461) (2.425,-1.446) (2.375,-1.430) (2.325,-1.413) (2.275,-1.395) (2.225,-1.376) (2.175,-1.357) (2.125,-1.336) (2.075,-1.315) (2.025,-1.293) (1.975,-1.270) (1.925,-1.246) (1.875,-1.221) (1.825,-1.196) (1.775,-1.170) (1.725,-1.144) (1.675,-1.116) (1.625,-1.089) (1.575,-1.060) (1.525,-1.031) (1.475,-1.001) (1.432,-0.975) (1.375,-0.941) (1.325,-0.910) (1.275,-0.878) (1.232,-0.850) (1.175,-0.814) (1.125,-0.781) (1.078,-0.750) (1.025,-0.715) (0.975,-0.681) (0.929,-0.650) (0.875,-0.614) (0.825,-0.579) (0.783,-0.550) (0.725,-0.510) (0.675,-0.475) (0.625,-0.441) (0.575,-0.406) (0.532,-0.375) (0.475,-0.336) (0.425,-0.300) (0.375,-0.266) (0.325,-0.230) (0.284,-0.200) (0.225,-0.160) (0.206,-0.144)}
\EQ{(-0.220,-0.155) (-0.275,-0.195) (-0.319,-0.225) (-0.375,-0.266) (-0.425,-0.300) (-0.475,-0.336) (-0.525,-0.371) (-0.568,-0.400) (-0.625,-0.441) (-0.674,-0.475) (-0.725,-0.511) (-0.775,-0.545) (-0.819,-0.575) (-0.875,-0.614) (-0.925,-0.648) (-0.975,-0.682) (-1.025,-0.716) (-1.075,-0.750) (-1.125,-0.783) (-1.175,-0.816) (-1.225,-0.848) (-1.275,-0.881) (-1.325,-0.913) (-1.375,-0.944) (-1.424,-0.975) (-1.475,-1.006) (-1.525,-1.037) (-1.575,-1.066) (-1.625,-1.096) (-1.675,-1.124) (-1.721,-1.150) (-1.766,-1.175) (-1.812,-1.200) (-1.859,-1.225) (-1.907,-1.250) (-1.957,-1.275) (-2.008,-1.300) (-2.060,-1.325) (-2.114,-1.350) (-2.171,-1.375) (-2.225,-1.398) (-2.275,-1.419) (-2.325,-1.438) (-2.375,-1.457) (-2.425,-1.475) (-2.475,-1.492) (-2.525,-1.508) (-2.575,-1.523) (-2.625,-1.538) (-2.675,-1.551) (-2.725,-1.563) (-2.775,-1.574) (-2.825,-1.585) (-2.875,-1.594) (-2.925,-1.602) (-2.975,-1.609) (-3.025,-1.615) (-3.075,-1.620) (-3.125,-1.623) (-3.175,-1.626) (-3.225,-1.627) (-3.275,-1.627) (-3.325,-1.626) (-3.375,-1.624) (-3.425,-1.620) (-3.475,-1.615) (-3.525,-1.609) (-3.575,-1.601) (-3.625,-1.592) (-3.675,-1.581) (-3.725,-1.569) (-3.775,-1.555) (-3.825,-1.539) (-3.875,-1.522) (-3.925,-1.503) (-3.975,-1.482) (-4.025,-1.459) (-4.075,-1.434) (-4.125,-1.407) (-4.175,-1.378) (-4.218,-1.350) (-4.275,-1.311) (-4.323,-1.275) (-4.375,-1.232) (-4.425,-1.188) (-4.464,-1.150) (-4.500,-1.113) (-4.535,-1.075) (-4.575,-1.028) (-4.616,-0.975) (-4.650,-0.928) (-4.685,-0.875) (-4.716,-0.825) (-4.744,-0.775) (-4.770,-0.725) (-4.793,-0.675) (-4.815,-0.625) (-4.835,-0.575) (-4.853,-0.525) (-4.869,-0.475) (-4.883,-0.425) (-4.895,-0.375) (-4.906,-0.325) (-4.916,-0.275) (-4.923,-0.225) (-4.929,-0.175) (-4.934,-0.125) (-4.937,-0.075) (-4.938,-0.025) (-4.938,0.050) (-4.936,0.100) (-4.932,0.150) (-4.926,0.200) (-4.920,0.250) (-4.911,0.300) (-4.901,0.350) (-4.889,0.400) (-4.876,0.450) (-4.861,0.500) (-4.844,0.550) (-4.825,0.600) (-4.804,0.650) (-4.782,0.700) (-4.757,0.750) (-4.730,0.800) (-4.701,0.850) (-4.669,0.900) (-4.635,0.950) (-4.600,0.996) (-4.556,1.050) (-4.512,1.100) (-4.475,1.139) (-4.438,1.175) (-4.400,1.210) (-4.354,1.250) (-4.300,1.292) (-4.255,1.325) (-4.200,1.362) (-4.150,1.393) (-4.100,1.421) (-4.050,1.447) (-4.000,1.471) (-3.950,1.493) (-3.900,1.513) (-3.850,1.531) (-3.800,1.547) (-3.750,1.562) (-3.700,1.575) (-3.650,1.586) (-3.600,1.596) (-3.550,1.605) (-3.500,1.612) (-3.450,1.618) (-3.400,1.622) (-3.350,1.625) (-3.300,1.627) (-3.250,1.627) (-3.200,1.627) (-3.150,1.625) (-3.100,1.622) (-3.050,1.617) (-3.000,1.612) (-2.950,1.605) (-2.900,1.598) (-2.850,1.589) (-2.800,1.580) (-2.750,1.569) (-2.700,1.557) (-2.650,1.544) (-2.600,1.531) (-2.550,1.516) (-2.500,1.500) (-2.450,1.484) (-2.400,1.466) (-2.350,1.448) (-2.300,1.429) (-2.250,1.409) (-2.200,1.388) (-2.150,1.366) (-2.100,1.344) (-2.050,1.320) (-2.000,1.296) (-1.950,1.272) (-1.900,1.246) (-1.850,1.220) (-1.800,1.194) (-1.750,1.167) (-1.700,1.139) (-1.650,1.110) (-1.600,1.081) (-1.550,1.051) (-1.506,1.025) (-1.450,0.991) (-1.400,0.960) (-1.350,0.929) (-1.305,0.900) (-1.250,0.865) (-1.200,0.832) (-1.151,0.800) (-1.100,0.766) (-1.050,0.733) (-1.001,0.700) (-0.950,0.665) (-0.900,0.631) (-0.855,0.600) (-0.800,0.563) (-0.750,0.528) (-0.700,0.493) (-0.650,0.459) (-0.603,0.425) (-0.550,0.389) (-0.500,0.353) (-0.450,0.318) (-0.400,0.283) (-0.354,0.250) (-0.300,0.213) (-0.250,0.177) (-0.220,0.155)}
\EQ{(0.206,0.144) (0.248,0.175) (0.300,0.213) (0.350,0.248) (0.400,0.283) (0.450,0.318) (0.496,0.350) (0.550,0.389) (0.600,0.423) (0.650,0.459) (0.700,0.493) (0.746,0.525) (0.800,0.563) (0.850,0.597) (0.891,0.625) (0.950,0.665) (1.000,0.699) (1.050,0.733) (1.100,0.766) (1.150,0.799) (1.200,0.832) (1.250,0.865) (1.300,0.897) (1.345,0.925) (1.400,0.960) (1.450,0.991) (1.500,1.021) (1.548,1.050) (1.600,1.081) (1.650,1.110) (1.700,1.139) (1.750,1.167) (1.800,1.194) (1.850,1.220) (1.900,1.246) (1.950,1.272) (2.000,1.296) (2.050,1.320) (2.100,1.344) (2.150,1.366) (2.200,1.388) (2.250,1.409) (2.300,1.429) (2.350,1.448) (2.400,1.466) (2.450,1.484) (2.500,1.500) (2.550,1.516) (2.600,1.531) (2.650,1.544) (2.700,1.557) (2.750,1.569) (2.800,1.580) (2.850,1.589) (2.900,1.598) (2.950,1.605) (3.000,1.612) (3.050,1.617) (3.100,1.622) (3.150,1.625) (3.200,1.627) (3.250,1.627) (3.300,1.627) (3.350,1.625) (3.400,1.622) (3.450,1.618) (3.500,1.612) (3.550,1.605) (3.600,1.596) (3.650,1.586) (3.700,1.575) (3.750,1.562) (3.800,1.547) (3.850,1.531) (3.900,1.513) (3.950,1.493) (4.000,1.471) (4.050,1.447) (4.100,1.421) (4.150,1.393) (4.200,1.362) (4.250,1.328) (4.300,1.292) (4.350,1.253) (4.400,1.210) (4.438,1.175) (4.475,1.139) (4.512,1.100) (4.550,1.058) (4.597,1.000) (4.635,0.950) (4.669,0.900) (4.700,0.851) (4.730,0.800) (4.757,0.750) (4.782,0.700) (4.804,0.650) (4.825,0.600) (4.844,0.550) (4.861,0.500) (4.876,0.450) (4.889,0.400) (4.901,0.350) (4.911,0.300) (4.920,0.250) (4.926,0.200) (4.932,0.150) (4.936,0.100) (4.938,0.050) (4.939,0.000) (4.938,-0.050) (4.936,-0.100) (4.932,-0.150) (4.926,-0.200) (4.920,-0.250) (4.911,-0.300) (4.901,-0.350) (4.889,-0.400) (4.876,-0.450) (4.861,-0.500) (4.844,-0.550) (4.825,-0.600) (4.804,-0.650) (4.782,-0.700) (4.757,-0.750) (4.730,-0.800) (4.701,-0.850) (4.669,-0.900) (4.635,-0.950) (4.600,-0.996) (4.556,-1.050) (4.512,-1.100) (4.475,-1.139) (4.438,-1.175) (4.400,-1.210) (4.354,-1.250) (4.300,-1.292) (4.255,-1.325) (4.200,-1.362) (4.150,-1.393) (4.100,-1.421) (4.050,-1.447) (4.000,-1.471) (3.950,-1.493) (3.900,-1.513) (3.850,-1.531) (3.800,-1.547) (3.750,-1.562) (3.700,-1.575) (3.650,-1.586) (3.600,-1.596) (3.550,-1.605) (3.500,-1.612) (3.450,-1.618) (3.400,-1.622) (3.350,-1.625) (3.300,-1.627) (3.250,-1.627) (3.200,-1.627) (3.150,-1.625) (3.100,-1.622) (3.050,-1.617) (3.000,-1.612) (2.950,-1.605) (2.900,-1.598) (2.850,-1.589) (2.800,-1.580) (2.750,-1.569) (2.700,-1.557) (2.650,-1.544) (2.600,-1.531) (2.550,-1.516) (2.500,-1.500) (2.450,-1.484) (2.400,-1.466) (2.350,-1.448) (2.300,-1.429) (2.250,-1.409) (2.200,-1.388) (2.150,-1.366) (2.100,-1.344) (2.050,-1.320) (2.000,-1.296) (1.950,-1.272) (1.900,-1.246) (1.850,-1.220) (1.800,-1.194) (1.750,-1.167) (1.700,-1.139) (1.650,-1.110) (1.600,-1.081) (1.550,-1.051) (1.506,-1.025) (1.450,-0.991) (1.400,-0.960) (1.350,-0.929) (1.305,-0.900) (1.250,-0.865) (1.200,-0.832) (1.151,-0.800) (1.100,-0.766) (1.050,-0.733) (1.001,-0.700) (0.950,-0.665) (0.900,-0.631) (0.855,-0.600) (0.800,-0.563) (0.750,-0.528) (0.700,-0.493) (0.650,-0.459) (0.603,-0.425) (0.550,-0.389) (0.500,-0.353) (0.450,-0.318) (0.400,-0.283) (0.354,-0.250) (0.300,-0.213) (0.250,-0.177) (0.206,-0.144)}
\EQ{(-0.220,-0.155) (-0.275,-0.195) (-0.319,-0.225) (-0.375,-0.266) (-0.425,-0.300) (-0.475,-0.336) (-0.525,-0.371) (-0.567,-0.400) (-0.625,-0.441) (-0.674,-0.475) (-0.725,-0.511) (-0.775,-0.546) (-0.818,-0.575) (-0.875,-0.615) (-0.925,-0.649) (-0.975,-0.683) (-1.025,-0.717) (-1.074,-0.750) (-1.125,-0.785) (-1.175,-0.818) (-1.224,-0.850) (-1.275,-0.884) (-1.325,-0.916) (-1.375,-0.948) (-1.417,-0.975) (-1.475,-1.011) (-1.525,-1.042) (-1.575,-1.073) (-1.620,-1.100) (-1.675,-1.133) (-1.725,-1.162) (-1.775,-1.190) (-1.825,-1.218) (-1.875,-1.246) (-1.925,-1.273) (-1.975,-1.299) (-2.025,-1.325) (-2.075,-1.350) (-2.125,-1.374) (-2.175,-1.398) (-2.225,-1.421) (-2.275,-1.443) (-2.325,-1.465) (-2.375,-1.486) (-2.425,-1.505) (-2.475,-1.525) (-2.525,-1.543) (-2.575,-1.560) (-2.625,-1.577) (-2.675,-1.593) (-2.725,-1.608) (-2.775,-1.622) (-2.825,-1.635) (-2.875,-1.647) (-2.925,-1.658) (-2.975,-1.668) (-3.025,-1.677) (-3.075,-1.686) (-3.125,-1.693) (-3.175,-1.699) (-3.225,-1.704) (-3.275,-1.708) (-3.325,-1.711) (-3.375,-1.713) (-3.425,-1.713) (-3.475,-1.713) (-3.525,-1.711) (-3.575,-1.708) (-3.625,-1.704) (-3.675,-1.698) (-3.725,-1.691) (-3.775,-1.683) (-3.825,-1.674) (-3.875,-1.663) (-3.925,-1.650) (-3.975,-1.636) (-4.025,-1.621) (-4.075,-1.604) (-4.125,-1.585) (-4.175,-1.564) (-4.225,-1.542) (-4.275,-1.518) (-4.325,-1.491) (-4.375,-1.462) (-4.425,-1.431) (-4.472,-1.400) (-4.525,-1.362) (-4.572,-1.325) (-4.625,-1.280) (-4.675,-1.234) (-4.725,-1.185) (-4.775,-1.130) (-4.822,-1.075) (-4.861,-1.025) (-4.897,-0.975) (-4.925,-0.932) (-4.961,-0.875) (-4.989,-0.825) (-5.015,-0.775) (-5.040,-0.725) (-5.062,-0.675) (-5.082,-0.625) (-5.100,-0.577) (-5.117,-0.525) (-5.133,-0.475) (-5.146,-0.425) (-5.158,-0.375) (-5.168,-0.325) (-5.177,-0.275) (-5.184,-0.225) (-5.190,-0.175) (-5.195,-0.125) (-5.197,-0.075) (-5.199,-0.025) (-5.198,0.050) (-5.196,0.100) (-5.193,0.150) (-5.187,0.200) (-5.181,0.250) (-5.173,0.300) (-5.163,0.350) (-5.152,0.400) (-5.140,0.450) (-5.125,0.500) (-5.109,0.550) (-5.092,0.600) (-5.072,0.650) (-5.051,0.700) (-5.028,0.750) (-5.002,0.800) (-4.975,0.850) (-4.946,0.900) (-4.914,0.950) (-4.879,1.000) (-4.842,1.050) (-4.801,1.100) (-4.758,1.150) (-4.710,1.200) (-4.659,1.250) (-4.602,1.300) (-4.550,1.343) (-4.507,1.375) (-4.450,1.415) (-4.400,1.447) (-4.353,1.475) (-4.309,1.500) (-4.260,1.525) (-4.208,1.550) (-4.150,1.575) (-4.100,1.595) (-4.050,1.613) (-4.000,1.629) (-3.950,1.644) (-3.900,1.657) (-3.850,1.668) (-3.800,1.679) (-3.750,1.687) (-3.700,1.695) (-3.650,1.701) (-3.600,1.706) (-3.550,1.710) (-3.500,1.712) (-3.450,1.713) (-3.400,1.713) (-3.350,1.712) (-3.300,1.710) (-3.250,1.706) (-3.200,1.702) (-3.150,1.696) (-3.100,1.689) (-3.050,1.682) (-3.000,1.673) (-2.950,1.663) (-2.900,1.653) (-2.850,1.641) (-2.800,1.628) (-2.750,1.615) (-2.700,1.600) (-2.650,1.585) (-2.600,1.569) (-2.550,1.552) (-2.500,1.534) (-2.450,1.515) (-2.400,1.496) (-2.350,1.475) (-2.300,1.454) (-2.250,1.432) (-2.200,1.410) (-2.150,1.386) (-2.100,1.362) (-2.050,1.338) (-2.000,1.312) (-1.950,1.286) (-1.900,1.259) (-1.850,1.232) (-1.800,1.204) (-1.750,1.176) (-1.705,1.150) (-1.650,1.118) (-1.600,1.088) (-1.550,1.058) (-1.500,1.027) (-1.457,1.000) (-1.400,0.964) (-1.350,0.932) (-1.300,0.900) (-1.250,0.867) (-1.200,0.835) (-1.150,0.801) (-1.100,0.768) (-1.050,0.734) (-1.000,0.700) (-0.950,0.666) (-0.900,0.632) (-0.854,0.600) (-0.800,0.563) (-0.750,0.528) (-0.700,0.494) (-0.650,0.459) (-0.603,0.425) (-0.550,0.389) (-0.500,0.353) (-0.450,0.318) (-0.400,0.283) (-0.354,0.250) (-0.300,0.213) (-0.250,0.177) (-0.220,0.155)}
\EQ{(0.206,0.144) (0.248,0.175) (0.300,0.213) (0.350,0.248) (0.400,0.283) (0.450,0.318) (0.496,0.350) (0.550,0.389) (0.600,0.423) (0.650,0.459) (0.700,0.494) (0.746,0.525) (0.800,0.563) (0.850,0.598) (0.900,0.632) (0.950,0.666) (1.000,0.700) (1.050,0.734) (1.100,0.768) (1.148,0.800) (1.200,0.835) (1.250,0.867) (1.300,0.900) (1.350,0.932) (1.400,0.964) (1.450,0.996) (1.497,1.025) (1.550,1.058) (1.600,1.088) (1.650,1.118) (1.700,1.147) (1.748,1.175) (1.792,1.200) (1.837,1.225) (1.883,1.250) (1.929,1.275) (1.975,1.299) (2.025,1.325) (2.075,1.350) (2.125,1.374) (2.175,1.398) (2.225,1.421) (2.275,1.443) (2.325,1.465) (2.375,1.486) (2.425,1.505) (2.475,1.525) (2.525,1.543) (2.575,1.560) (2.625,1.577) (2.675,1.593) (2.725,1.608) (2.775,1.622) (2.825,1.635) (2.875,1.647) (2.925,1.658) (2.975,1.668) (3.025,1.677) (3.075,1.686) (3.125,1.693) (3.175,1.699) (3.225,1.704) (3.275,1.708) (3.325,1.711) (3.375,1.713) (3.425,1.713) (3.475,1.713) (3.525,1.711) (3.575,1.708) (3.625,1.704) (3.675,1.698) (3.725,1.691) (3.775,1.683) (3.825,1.674) (3.875,1.663) (3.925,1.650) (3.975,1.636) (4.025,1.621) (4.075,1.604) (4.125,1.585) (4.175,1.564) (4.225,1.542) (4.275,1.518) (4.325,1.491) (4.375,1.462) (4.425,1.431) (4.472,1.400) (4.525,1.362) (4.572,1.325) (4.625,1.280) (4.675,1.234) (4.725,1.185) (4.775,1.130) (4.822,1.075) (4.861,1.025) (4.897,0.975) (4.925,0.932) (4.961,0.875) (4.989,0.825) (5.015,0.775) (5.040,0.725) (5.062,0.675) (5.082,0.625) (5.100,0.577) (5.117,0.525) (5.133,0.475) (5.146,0.425) (5.158,0.375) (5.168,0.325) (5.177,0.275) (5.184,0.225) (5.190,0.175) (5.195,0.125) (5.197,0.075) (5.199,0.025) (5.198,-0.050) (5.196,-0.100) (5.193,-0.150) (5.187,-0.200) (5.181,-0.250) (5.173,-0.300) (5.163,-0.350) (5.152,-0.400) (5.140,-0.450) (5.125,-0.500) (5.109,-0.550) (5.092,-0.600) (5.072,-0.650) (5.051,-0.700) (5.028,-0.750) (5.002,-0.800) (4.975,-0.850) (4.946,-0.900) (4.914,-0.950) (4.879,-1.000) (4.842,-1.050) (4.801,-1.100) (4.758,-1.150) (4.710,-1.200) (4.659,-1.250) (4.602,-1.300) (4.550,-1.343) (4.507,-1.375) (4.450,-1.415) (4.400,-1.447) (4.353,-1.475) (4.309,-1.500) (4.260,-1.525) (4.208,-1.550) (4.150,-1.575) (4.100,-1.595) (4.050,-1.613) (4.000,-1.629) (3.950,-1.644) (3.900,-1.657) (3.850,-1.668) (3.800,-1.679) (3.750,-1.687) (3.700,-1.695) (3.650,-1.701) (3.600,-1.706) (3.550,-1.710) (3.500,-1.712) (3.450,-1.713) (3.400,-1.713) (3.350,-1.712) (3.300,-1.710) (3.250,-1.706) (3.200,-1.702) (3.150,-1.696) (3.100,-1.689) (3.050,-1.682) (3.000,-1.673) (2.950,-1.663) (2.900,-1.653) (2.850,-1.641) (2.800,-1.628) (2.750,-1.615) (2.700,-1.600) (2.650,-1.585) (2.600,-1.569) (2.550,-1.552) (2.500,-1.534) (2.450,-1.515) (2.400,-1.496) (2.350,-1.475) (2.300,-1.454) (2.250,-1.432) (2.200,-1.410) (2.150,-1.386) (2.100,-1.362) (2.050,-1.338) (2.000,-1.312) (1.950,-1.286) (1.900,-1.259) (1.850,-1.232) (1.800,-1.204) (1.750,-1.176) (1.705,-1.150) (1.650,-1.118) (1.600,-1.088) (1.550,-1.058) (1.500,-1.027) (1.457,-1.000) (1.400,-0.964) (1.350,-0.932) (1.300,-0.900) (1.250,-0.867) (1.200,-0.835) (1.150,-0.801) (1.100,-0.768) (1.050,-0.734) (1.000,-0.700) (0.950,-0.666) (0.900,-0.632) (0.854,-0.600) (0.800,-0.563) (0.750,-0.528) (0.700,-0.494) (0.650,-0.459) (0.603,-0.425) (0.550,-0.389) (0.500,-0.353) (0.450,-0.318) (0.400,-0.283) (0.354,-0.250) (0.300,-0.213) (0.250,-0.177) (0.206,-0.144)}
\EQ{(-0.220,-0.155) (-0.275,-0.195) (-0.319,-0.225) (-0.375,-0.266) (-0.425,-0.300) (-0.475,-0.336) (-0.525,-0.371) (-0.567,-0.400) (-0.625,-0.441) (-0.674,-0.475) (-0.725,-0.511) (-0.775,-0.546) (-0.817,-0.575) (-0.875,-0.616) (-0.925,-0.650) (-0.975,-0.684) (-1.025,-0.719) (-1.072,-0.750) (-1.125,-0.787) (-1.175,-0.820) (-1.220,-0.850) (-1.275,-0.887) (-1.325,-0.920) (-1.372,-0.950) (-1.425,-0.985) (-1.475,-1.017) (-1.525,-1.048) (-1.568,-1.075) (-1.625,-1.110) (-1.675,-1.141) (-1.725,-1.171) (-1.774,-1.200) (-1.825,-1.230) (-1.875,-1.259) (-1.925,-1.287) (-1.975,-1.314) (-2.025,-1.342) (-2.075,-1.368) (-2.125,-1.394) (-2.175,-1.420) (-2.225,-1.444) (-2.275,-1.468) (-2.325,-1.492) (-2.375,-1.515) (-2.425,-1.537) (-2.475,-1.558) (-2.525,-1.579) (-2.575,-1.599) (-2.625,-1.618) (-2.675,-1.637) (-2.725,-1.654) (-2.775,-1.671) (-2.825,-1.687) (-2.875,-1.702) (-2.925,-1.717) (-2.975,-1.730) (-3.025,-1.743) (-3.075,-1.755) (-3.125,-1.765) (-3.175,-1.775) (-3.225,-1.784) (-3.275,-1.792) (-3.325,-1.799) (-3.375,-1.805) (-3.425,-1.811) (-3.475,-1.815) (-3.525,-1.818) (-3.575,-1.820) (-3.625,-1.820) (-3.675,-1.820) (-3.725,-1.819) (-3.775,-1.816) (-3.825,-1.813) (-3.875,-1.808) (-3.925,-1.802) (-3.975,-1.795) (-4.025,-1.786) (-4.075,-1.776) (-4.125,-1.765) (-4.175,-1.753) (-4.225,-1.739) (-4.275,-1.723) (-4.325,-1.706) (-4.375,-1.688) (-4.425,-1.667) (-4.475,-1.645) (-4.525,-1.622) (-4.575,-1.596) (-4.625,-1.568) (-4.675,-1.538) (-4.725,-1.506) (-4.770,-1.475) (-4.825,-1.435) (-4.868,-1.400) (-4.925,-1.352) (-4.975,-1.305) (-5.025,-1.255) (-5.075,-1.200) (-5.118,-1.150) (-5.150,-1.109) (-5.194,-1.050) (-5.225,-1.004) (-5.259,-0.950) (-5.289,-0.900) (-5.316,-0.850) (-5.341,-0.800) (-5.365,-0.750) (-5.386,-0.700) (-5.406,-0.650) (-5.424,-0.600) (-5.441,-0.550) (-5.456,-0.500) (-5.469,-0.450) (-5.481,-0.400) (-5.491,-0.350) (-5.500,-0.300) (-5.508,-0.250) (-5.514,-0.200) (-5.519,-0.150) (-5.522,-0.100) (-5.524,-0.050) (-5.525,0.000) (-5.524,0.050) (-5.522,0.100) (-5.519,0.150) (-5.514,0.200) (-5.508,0.250) (-5.500,0.300) (-5.491,0.350) (-5.481,0.400) (-5.469,0.450) (-5.456,0.500) (-5.441,0.550) (-5.424,0.600) (-5.406,0.650) (-5.386,0.700) (-5.365,0.750) (-5.341,0.800) (-5.316,0.850) (-5.289,0.900) (-5.259,0.950) (-5.228,1.000) (-5.194,1.050) (-5.157,1.100) (-5.125,1.141) (-5.075,1.200) (-5.030,1.250) (-4.980,1.300) (-4.927,1.350) (-4.875,1.395) (-4.825,1.435) (-4.775,1.472) (-4.725,1.506) (-4.675,1.538) (-4.625,1.568) (-4.575,1.596) (-4.525,1.622) (-4.475,1.645) (-4.425,1.667) (-4.375,1.688) (-4.325,1.706) (-4.275,1.723) (-4.225,1.739) (-4.175,1.753) (-4.125,1.765) (-4.075,1.776) (-4.025,1.786) (-3.975,1.795) (-3.925,1.802) (-3.875,1.808) (-3.825,1.813) (-3.775,1.816) (-3.725,1.819) (-3.675,1.820) (-3.625,1.820) (-3.575,1.820) (-3.525,1.818) (-3.475,1.815) (-3.425,1.811) (-3.375,1.805) (-3.325,1.799) (-3.275,1.792) (-3.225,1.784) (-3.175,1.775) (-3.125,1.765) (-3.075,1.755) (-3.025,1.743) (-2.975,1.730) (-2.925,1.717) (-2.875,1.702) (-2.825,1.687) (-2.775,1.671) (-2.725,1.654) (-2.675,1.637) (-2.625,1.618) (-2.578,1.600) (-2.525,1.579) (-2.475,1.558) (-2.425,1.537) (-2.375,1.515) (-2.325,1.492) (-2.275,1.468) (-2.225,1.444) (-2.175,1.420) (-2.125,1.394) (-2.075,1.368) (-2.025,1.342) (-1.975,1.314) (-1.925,1.287) (-1.875,1.259) (-1.825,1.230) (-1.775,1.201) (-1.725,1.171) (-1.675,1.141) (-1.625,1.110) (-1.575,1.079) (-1.528,1.050) (-1.475,1.017) (-1.425,0.985) (-1.375,0.952) (-1.325,0.920) (-1.275,0.887) (-1.225,0.854) (-1.183,0.825) (-1.125,0.787) (-1.075,0.753) (-1.025,0.719) (-0.975,0.684) (-0.925,0.650) (-0.875,0.616) (-0.825,0.581) (-0.781,0.550) (-0.725,0.511) (-0.675,0.476) (-0.625,0.441) (-0.575,0.406) (-0.532,0.375) (-0.475,0.336) (-0.425,0.300) (-0.375,0.266) (-0.325,0.230) (-0.284,0.200) (-0.225,0.160) (-0.220,0.155)}
\EQ{(0.206,0.144) (0.248,0.175) (0.300,0.213) (0.350,0.248) (0.400,0.283) (0.450,0.318) (0.496,0.350) (0.550,0.389) (0.600,0.424) (0.650,0.459) (0.700,0.494) (0.745,0.525) (0.800,0.564) (0.850,0.598) (0.900,0.633) (0.950,0.667) (0.998,0.700) (1.050,0.736) (1.100,0.770) (1.145,0.800) (1.200,0.837) (1.250,0.870) (1.296,0.900) (1.350,0.936) (1.400,0.968) (1.449,1.000) (1.500,1.032) (1.550,1.064) (1.600,1.095) (1.649,1.125) (1.700,1.156) (1.750,1.186) (1.800,1.215) (1.850,1.244) (1.900,1.273) (1.949,1.300) (1.995,1.325) (2.041,1.350) (2.088,1.375) (2.137,1.400) (2.186,1.425) (2.237,1.450) (2.289,1.475) (2.343,1.500) (2.398,1.525) (2.450,1.548) (2.500,1.569) (2.550,1.589) (2.600,1.609) (2.650,1.628) (2.700,1.646) (2.750,1.663) (2.800,1.679) (2.850,1.695) (2.900,1.710) (2.950,1.724) (3.000,1.737) (3.050,1.749) (3.100,1.760) (3.150,1.771) (3.200,1.780) (3.250,1.788) (3.300,1.796) (3.350,1.803) (3.400,1.808) (3.450,1.813) (3.500,1.816) (3.550,1.819) (3.600,1.820) (3.650,1.820) (3.700,1.820) (3.750,1.818) (3.800,1.815) (3.850,1.811) (3.900,1.805) (3.950,1.799) (4.000,1.791) (4.050,1.782) (4.100,1.771) (4.150,1.759) (4.200,1.746) (4.250,1.731) (4.300,1.715) (4.350,1.697) (4.400,1.678) (4.450,1.657) (4.500,1.634) (4.550,1.609) (4.600,1.582) (4.650,1.553) (4.696,1.525) (4.750,1.489) (4.800,1.453) (4.850,1.415) (4.898,1.375) (4.950,1.329) (5.000,1.280) (5.050,1.228) (5.097,1.175) (5.138,1.125) (5.175,1.076) (5.211,1.025) (5.244,0.975) (5.274,0.925) (5.300,0.880) (5.325,0.832) (5.350,0.782) (5.375,0.726) (5.396,0.675) (5.415,0.625) (5.433,0.575) (5.448,0.525) (5.462,0.475) (5.475,0.425) (5.486,0.375) (5.496,0.325) (5.504,0.275) (5.511,0.225) (5.516,0.175) (5.521,0.125) (5.523,0.075) (5.525,0.025) (5.524,-0.050) (5.522,-0.100) (5.519,-0.150) (5.514,-0.200) (5.508,-0.250) (5.500,-0.300) (5.491,-0.350) (5.481,-0.400) (5.469,-0.450) (5.456,-0.500) (5.441,-0.550) (5.424,-0.600) (5.406,-0.650) (5.386,-0.700) (5.365,-0.750) (5.341,-0.800) (5.316,-0.850) (5.289,-0.900) (5.259,-0.950) (5.228,-1.000) (5.194,-1.050) (5.157,-1.100) (5.125,-1.141) (5.075,-1.200) (5.030,-1.250) (4.980,-1.300) (4.927,-1.350) (4.875,-1.395) (4.825,-1.435) (4.775,-1.472) (4.725,-1.506) (4.675,-1.538) (4.625,-1.568) (4.575,-1.596) (4.525,-1.622) (4.475,-1.645) (4.425,-1.667) (4.375,-1.688) (4.325,-1.706) (4.275,-1.723) (4.225,-1.739) (4.175,-1.753) (4.125,-1.765) (4.075,-1.776) (4.025,-1.786) (3.975,-1.795) (3.925,-1.802) (3.875,-1.808) (3.825,-1.813) (3.775,-1.816) (3.725,-1.819) (3.675,-1.820) (3.625,-1.820) (3.575,-1.820) (3.525,-1.818) (3.475,-1.815) (3.425,-1.811) (3.375,-1.805) (3.325,-1.799) (3.275,-1.792) (3.225,-1.784) (3.175,-1.775) (3.125,-1.765) (3.075,-1.755) (3.025,-1.743) (2.975,-1.730) (2.925,-1.717) (2.875,-1.702) (2.825,-1.687) (2.775,-1.671) (2.725,-1.654) (2.675,-1.637) (2.625,-1.618) (2.578,-1.600) (2.525,-1.579) (2.475,-1.558) (2.425,-1.537) (2.375,-1.515) (2.325,-1.492) (2.275,-1.468) (2.225,-1.444) (2.175,-1.420) (2.125,-1.394) (2.075,-1.368) (2.025,-1.342) (1.975,-1.314) (1.925,-1.287) (1.875,-1.259) (1.825,-1.230) (1.775,-1.201) (1.725,-1.171) (1.675,-1.141) (1.625,-1.110) (1.575,-1.079) (1.528,-1.050) (1.475,-1.017) (1.425,-0.985) (1.375,-0.952) (1.325,-0.920) (1.275,-0.887) (1.225,-0.854) (1.183,-0.825) (1.125,-0.787) (1.075,-0.753) (1.025,-0.719) (0.975,-0.684) (0.925,-0.650) (0.875,-0.616) (0.825,-0.581) (0.781,-0.550) (0.725,-0.511) (0.675,-0.476) (0.625,-0.441) (0.575,-0.406) (0.532,-0.375) (0.475,-0.336) (0.425,-0.300) (0.375,-0.266) (0.325,-0.230) (0.284,-0.200) (0.225,-0.160) (0.206,-0.144)}
\EQ{(-0.220,-0.155) (-0.275,-0.195) (-0.319,-0.225) (-0.375,-0.266) (-0.425,-0.300) (-0.475,-0.336) (-0.525,-0.371) (-0.567,-0.400) (-0.625,-0.442) (-0.673,-0.475) (-0.725,-0.512) (-0.775,-0.547) (-0.817,-0.575) (-0.875,-0.616) (-0.924,-0.650) (-0.975,-0.686) (-1.025,-0.720) (-1.069,-0.750) (-1.125,-0.788) (-1.175,-0.822) (-1.225,-0.856) (-1.275,-0.890) (-1.325,-0.923) (-1.375,-0.956) (-1.425,-0.989) (-1.475,-1.022) (-1.519,-1.050) (-1.575,-1.086) (-1.625,-1.118) (-1.675,-1.149) (-1.725,-1.180) (-1.775,-1.211) (-1.825,-1.242) (-1.875,-1.271) (-1.924,-1.300) (-1.975,-1.330) (-2.025,-1.359) (-2.075,-1.387) (-2.125,-1.415) (-2.175,-1.442) (-2.225,-1.468) (-2.275,-1.495) (-2.325,-1.520) (-2.375,-1.545) (-2.425,-1.570) (-2.475,-1.593) (-2.525,-1.617) (-2.575,-1.639) (-2.625,-1.661) (-2.675,-1.682) (-2.725,-1.703) (-2.775,-1.723) (-2.825,-1.742) (-2.875,-1.761) (-2.925,-1.778) (-2.975,-1.796) (-3.025,-1.812) (-3.075,-1.827) (-3.125,-1.842) (-3.175,-1.856) (-3.225,-1.869) (-3.275,-1.881) (-3.325,-1.893) (-3.375,-1.904) (-3.425,-1.913) (-3.475,-1.922) (-3.525,-1.930) (-3.575,-1.937) (-3.625,-1.944) (-3.675,-1.949) (-3.725,-1.953) (-3.775,-1.957) (-3.825,-1.959) (-3.875,-1.961) (-3.925,-1.961) (-3.975,-1.960) (-4.025,-1.959) (-4.075,-1.956) (-4.125,-1.952) (-4.175,-1.948) (-4.225,-1.942) (-4.275,-1.934) (-4.325,-1.926) (-4.375,-1.917) (-4.425,-1.906) (-4.475,-1.894) (-4.525,-1.880) (-4.575,-1.866) (-4.624,-1.850) (-4.675,-1.832) (-4.725,-1.813) (-4.775,-1.793) (-4.825,-1.770) (-4.875,-1.746) (-4.925,-1.721) (-4.975,-1.693) (-5.025,-1.664) (-5.075,-1.632) (-5.123,-1.600) (-5.175,-1.562) (-5.223,-1.525) (-5.275,-1.482) (-5.325,-1.438) (-5.365,-1.400) (-5.415,-1.350) (-5.450,-1.312) (-5.500,-1.254) (-5.544,-1.200) (-5.575,-1.158) (-5.616,-1.100) (-5.649,-1.050) (-5.675,-1.007) (-5.708,-0.950) (-5.735,-0.900) (-5.759,-0.850) (-5.782,-0.800) (-5.804,-0.750) (-5.823,-0.700) (-5.842,-0.650) (-5.858,-0.600) (-5.874,-0.550) (-5.887,-0.500) (-5.900,-0.450) (-5.911,-0.400) (-5.920,-0.350) (-5.929,-0.300) (-5.936,-0.250) (-5.941,-0.200) (-5.946,-0.150) (-5.949,-0.100) (-5.951,-0.050) (-5.951,0.000) (-5.951,0.050) (-5.949,0.100) (-5.946,0.150) (-5.941,0.200) (-5.936,0.250) (-5.929,0.300) (-5.920,0.350) (-5.911,0.400) (-5.900,0.450) (-5.887,0.500) (-5.874,0.550) (-5.858,0.600) (-5.842,0.650) (-5.823,0.700) (-5.804,0.750) (-5.782,0.800) (-5.759,0.850) (-5.735,0.900) (-5.708,0.950) (-5.679,1.000) (-5.650,1.048) (-5.616,1.100) (-5.581,1.150) (-5.550,1.192) (-5.504,1.250) (-5.461,1.300) (-5.425,1.339) (-5.375,1.390) (-5.325,1.438) (-5.284,1.475) (-5.225,1.524) (-5.175,1.562) (-5.125,1.599) (-5.075,1.632) (-5.025,1.664) (-4.975,1.693) (-4.925,1.721) (-4.875,1.746) (-4.825,1.770) (-4.775,1.793) (-4.725,1.813) (-4.675,1.832) (-4.625,1.850) (-4.575,1.866) (-4.525,1.880) (-4.475,1.894) (-4.425,1.906) (-4.375,1.917) (-4.325,1.926) (-4.275,1.934) (-4.225,1.942) (-4.175,1.948) (-4.125,1.952) (-4.075,1.956) (-4.025,1.959) (-3.975,1.960) (-3.925,1.961) (-3.875,1.961) (-3.825,1.959) (-3.775,1.957) (-3.725,1.953) (-3.675,1.949) (-3.625,1.944) (-3.575,1.937) (-3.525,1.930) (-3.475,1.922) (-3.425,1.913) (-3.375,1.904) (-3.325,1.893) (-3.275,1.881) (-3.225,1.869) (-3.175,1.856) (-3.125,1.842) (-3.075,1.827) (-3.025,1.812) (-2.975,1.796) (-2.925,1.778) (-2.875,1.761) (-2.825,1.742) (-2.775,1.723) (-2.725,1.703) (-2.675,1.682) (-2.625,1.661) (-2.575,1.639) (-2.525,1.617) (-2.475,1.593) (-2.425,1.570) (-2.375,1.545) (-2.325,1.520) (-2.275,1.495) (-2.225,1.468) (-2.175,1.442) (-2.125,1.415) (-2.075,1.387) (-2.025,1.359) (-1.975,1.330) (-1.925,1.301) (-1.881,1.275) (-1.825,1.242) (-1.775,1.211) (-1.725,1.180) (-1.676,1.150) (-1.625,1.118) (-1.575,1.086) (-1.525,1.054) (-1.480,1.025) (-1.425,0.989) (-1.375,0.956) (-1.328,0.925) (-1.275,0.890) (-1.225,0.856) (-1.179,0.825) (-1.125,0.788) (-1.075,0.754) (-1.033,0.725) (-0.975,0.686) (-0.925,0.651) (-0.875,0.616) (-0.825,0.581) (-0.781,0.550) (-0.725,0.512) (-0.675,0.476) (-0.625,0.442) (-0.575,0.406) (-0.531,0.375) (-0.475,0.336) (-0.425,0.300) (-0.375,0.266) (-0.325,0.230) (-0.284,0.200) (-0.225,0.160) (-0.220,0.155)}
\EQ{(0.206,0.144) (0.248,0.175) (0.300,0.213) (0.350,0.248) (0.400,0.283) (0.450,0.318) (0.496,0.350) (0.550,0.389) (0.600,0.424) (0.650,0.459) (0.700,0.494) (0.745,0.525) (0.800,0.564) (0.850,0.599) (0.900,0.634) (0.950,0.668) (0.997,0.700) (1.050,0.737) (1.100,0.771) (1.143,0.800) (1.200,0.839) (1.250,0.873) (1.300,0.906) (1.350,0.940) (1.400,0.973) (1.442,1.000) (1.500,1.038) (1.550,1.070) (1.597,1.100) (1.650,1.134) (1.700,1.165) (1.750,1.196) (1.798,1.225) (1.850,1.257) (1.900,1.286) (1.950,1.316) (2.000,1.344) (2.050,1.373) (2.099,1.400) (2.144,1.425) (2.191,1.450) (2.238,1.475) (2.286,1.500) (2.335,1.525) (2.385,1.550) (2.437,1.575) (2.489,1.600) (2.544,1.625) (2.600,1.650) (2.650,1.672) (2.700,1.693) (2.750,1.713) (2.800,1.733) (2.850,1.751) (2.900,1.770) (2.950,1.787) (3.000,1.804) (3.050,1.820) (3.100,1.835) (3.150,1.849) (3.200,1.863) (3.249,1.875) (3.300,1.887) (3.350,1.898) (3.400,1.909) (3.450,1.918) (3.500,1.926) (3.550,1.934) (3.600,1.941) (3.650,1.946) (3.700,1.951) (3.750,1.955) (3.800,1.958) (3.850,1.960) (3.900,1.961) (3.950,1.961) (4.000,1.960) (4.050,1.958) (4.100,1.954) (4.150,1.950) (4.200,1.945) (4.250,1.938) (4.300,1.930) (4.350,1.922) (4.400,1.911) (4.450,1.900) (4.500,1.887) (4.550,1.873) (4.600,1.858) (4.650,1.841) (4.700,1.823) (4.750,1.803) (4.800,1.782) (4.850,1.759) (4.900,1.734) (4.950,1.707) (5.000,1.679) (5.047,1.650) (5.100,1.616) (5.150,1.581) (5.191,1.550) (5.250,1.503) (5.300,1.461) (5.339,1.425) (5.390,1.375) (5.438,1.325) (5.475,1.284) (5.524,1.225) (5.563,1.175) (5.599,1.125) (5.633,1.075) (5.664,1.025) (5.694,0.975) (5.721,0.925) (5.747,0.875) (5.771,0.825) (5.793,0.775) (5.814,0.725) (5.833,0.675) (5.850,0.625) (5.866,0.575) (5.881,0.525) (5.894,0.475) (5.905,0.425) (5.916,0.375) (5.925,0.325) (5.932,0.275) (5.939,0.225) (5.944,0.175) (5.947,0.125) (5.950,0.075) (5.951,0.025) (5.951,-0.050) (5.949,-0.100) (5.946,-0.150) (5.941,-0.200) (5.936,-0.250) (5.929,-0.300) (5.920,-0.350) (5.911,-0.400) (5.900,-0.450) (5.887,-0.500) (5.874,-0.550) (5.858,-0.600) (5.842,-0.650) (5.823,-0.700) (5.804,-0.750) (5.782,-0.800) (5.759,-0.850) (5.735,-0.900) (5.708,-0.950) (5.679,-1.000) (5.650,-1.048) (5.616,-1.100) (5.581,-1.150) (5.550,-1.192) (5.504,-1.250) (5.461,-1.300) (5.425,-1.339) (5.375,-1.390) (5.325,-1.438) (5.284,-1.475) (5.225,-1.524) (5.175,-1.562) (5.125,-1.599) (5.075,-1.632) (5.025,-1.664) (4.975,-1.693) (4.925,-1.721) (4.875,-1.746) (4.825,-1.770) (4.775,-1.793) (4.725,-1.813) (4.675,-1.832) (4.625,-1.850) (4.575,-1.866) (4.525,-1.880) (4.475,-1.894) (4.425,-1.906) (4.375,-1.917) (4.325,-1.926) (4.275,-1.934) (4.225,-1.942) (4.175,-1.948) (4.125,-1.952) (4.075,-1.956) (4.025,-1.959) (3.975,-1.960) (3.925,-1.961) (3.875,-1.961) (3.825,-1.959) (3.775,-1.957) (3.725,-1.953) (3.675,-1.949) (3.625,-1.944) (3.575,-1.937) (3.525,-1.930) (3.475,-1.922) (3.425,-1.913) (3.375,-1.904) (3.325,-1.893) (3.275,-1.881) (3.225,-1.869) (3.175,-1.856) (3.125,-1.842) (3.075,-1.827) (3.025,-1.812) (2.975,-1.796) (2.925,-1.778) (2.875,-1.761) (2.825,-1.742) (2.775,-1.723) (2.725,-1.703) (2.675,-1.682) (2.625,-1.661) (2.575,-1.639) (2.525,-1.617) (2.475,-1.593) (2.425,-1.570) (2.375,-1.545) (2.325,-1.520) (2.275,-1.495) (2.225,-1.468) (2.175,-1.442) (2.125,-1.415) (2.075,-1.387) (2.025,-1.359) (1.975,-1.330) (1.925,-1.301) (1.881,-1.275) (1.825,-1.242) (1.775,-1.211) (1.725,-1.180) (1.676,-1.150) (1.625,-1.118) (1.575,-1.086) (1.525,-1.054) (1.480,-1.025) (1.425,-0.989) (1.375,-0.956) (1.328,-0.925) (1.275,-0.890) (1.225,-0.856) (1.179,-0.825) (1.125,-0.788) (1.075,-0.754) (1.033,-0.725) (0.975,-0.686) (0.925,-0.651) (0.875,-0.616) (0.825,-0.581) (0.781,-0.550) (0.725,-0.512) (0.675,-0.476) (0.625,-0.442) (0.575,-0.406) (0.531,-0.375) (0.475,-0.336) (0.425,-0.300) (0.375,-0.266) (0.325,-0.230) (0.284,-0.200) (0.225,-0.160) (0.206,-0.144)}
\EQ{(-0.220,-0.155) (-0.275,-0.195) (-0.319,-0.225) (-0.375,-0.266) (-0.425,-0.300) (-0.475,-0.336) (-0.525,-0.371) (-0.567,-0.400) (-0.625,-0.442) (-0.673,-0.475) (-0.725,-0.512) (-0.775,-0.547) (-0.825,-0.582) (-0.875,-0.617) (-0.923,-0.650) (-0.975,-0.687) (-1.025,-0.721) (-1.067,-0.750) (-1.125,-0.790) (-1.175,-0.824) (-1.225,-0.859) (-1.275,-0.893) (-1.323,-0.925) (-1.375,-0.960) (-1.425,-0.994) (-1.472,-1.025) (-1.525,-1.060) (-1.575,-1.093) (-1.624,-1.125) (-1.675,-1.158) (-1.725,-1.190) (-1.775,-1.222) (-1.820,-1.250) (-1.875,-1.285) (-1.925,-1.316) (-1.975,-1.346) (-2.023,-1.375) (-2.075,-1.406) (-2.125,-1.436) (-2.175,-1.465) (-2.225,-1.493) (-2.275,-1.521) (-2.325,-1.549) (-2.372,-1.575) (-2.419,-1.600) (-2.466,-1.625) (-2.514,-1.650) (-2.563,-1.675) (-2.613,-1.700) (-2.664,-1.725) (-2.716,-1.750) (-2.770,-1.775) (-2.825,-1.800) (-2.875,-1.822) (-2.925,-1.844) (-2.975,-1.865) (-3.025,-1.885) (-3.075,-1.904) (-3.125,-1.923) (-3.175,-1.942) (-3.225,-1.959) (-3.275,-1.976) (-3.325,-1.993) (-3.375,-2.008) (-3.425,-2.023) (-3.475,-2.037) (-3.525,-2.051) (-3.575,-2.064) (-3.625,-2.076) (-3.675,-2.087) (-3.725,-2.097) (-3.775,-2.107) (-3.825,-2.116) (-3.875,-2.124) (-3.925,-2.131) (-3.975,-2.138) (-4.025,-2.143) (-4.075,-2.148) (-4.125,-2.152) (-4.175,-2.155) (-4.225,-2.157) (-4.275,-2.158) (-4.325,-2.158) (-4.375,-2.158) (-4.425,-2.156) (-4.475,-2.154) (-4.525,-2.150) (-4.575,-2.146) (-4.625,-2.140) (-4.675,-2.134) (-4.725,-2.126) (-4.775,-2.117) (-4.825,-2.108) (-4.875,-2.097) (-4.925,-2.085) (-4.975,-2.071) (-5.025,-2.057) (-5.075,-2.041) (-5.125,-2.024) (-5.175,-2.006) (-5.225,-1.986) (-5.275,-1.965) (-5.325,-1.942) (-5.375,-1.918) (-5.425,-1.892) (-5.475,-1.864) (-5.525,-1.835) (-5.575,-1.804) (-5.618,-1.775) (-5.675,-1.735) (-5.722,-1.700) (-5.775,-1.658) (-5.814,-1.625) (-5.869,-1.575) (-5.921,-1.525) (-5.970,-1.475) (-6.016,-1.425) (-6.050,-1.385) (-6.098,-1.325) (-6.136,-1.275) (-6.172,-1.225) (-6.200,-1.183) (-6.236,-1.125) (-6.266,-1.075) (-6.294,-1.025) (-6.320,-0.975) (-6.344,-0.925) (-6.367,-0.875) (-6.388,-0.825) (-6.408,-0.775) (-6.426,-0.725) (-6.443,-0.675) (-6.459,-0.625) (-6.473,-0.575) (-6.486,-0.525) (-6.498,-0.475) (-6.509,-0.425) (-6.518,-0.375) (-6.526,-0.325) (-6.533,-0.275) (-6.539,-0.225) (-6.543,-0.175) (-6.547,-0.125) (-6.549,-0.075) (-6.550,-0.025) (-6.550,0.037) (-6.548,0.100) (-6.545,0.150) (-6.541,0.200) (-6.536,0.250) (-6.530,0.300) (-6.522,0.350) (-6.513,0.400) (-6.503,0.450) (-6.492,0.500) (-6.480,0.550) (-6.466,0.600) (-6.451,0.650) (-6.435,0.700) (-6.417,0.750) (-6.398,0.800) (-6.378,0.850) (-6.356,0.900) (-6.332,0.950) (-6.307,1.000) (-6.280,1.050) (-6.251,1.100) (-6.225,1.143) (-6.189,1.200) (-6.154,1.250) (-6.118,1.300) (-6.079,1.350) (-6.037,1.400) (-6.000,1.442) (-5.950,1.496) (-5.900,1.546) (-5.850,1.593) (-5.800,1.637) (-5.754,1.675) (-5.700,1.717) (-5.654,1.750) (-5.600,1.787) (-5.550,1.819) (-5.500,1.850) (-5.456,1.875) (-5.409,1.900) (-5.360,1.925) (-5.307,1.950) (-5.251,1.975) (-5.200,1.996) (-5.150,2.015) (-5.100,2.033) (-5.050,2.049) (-5.000,2.064) (-4.950,2.078) (-4.900,2.091) (-4.850,2.102) (-4.800,2.113) (-4.750,2.122) (-4.700,2.130) (-4.650,2.137) (-4.600,2.143) (-4.550,2.148) (-4.500,2.152) (-4.450,2.155) (-4.400,2.157) (-4.350,2.158) (-4.300,2.158) (-4.250,2.158) (-4.200,2.156) (-4.150,2.153) (-4.100,2.150) (-4.050,2.146) (-4.000,2.140) (-3.950,2.134) (-3.900,2.128) (-3.850,2.120) (-3.800,2.111) (-3.750,2.102) (-3.700,2.092) (-3.650,2.081) (-3.600,2.070) (-3.550,2.057) (-3.500,2.044) (-3.450,2.030) (-3.400,2.016) (-3.350,2.001) (-3.300,1.985) (-3.250,1.968) (-3.200,1.951) (-3.150,1.933) (-3.100,1.914) (-3.050,1.895) (-3.001,1.875) (-2.950,1.854) (-2.900,1.833) (-2.850,1.811) (-2.800,1.789) (-2.750,1.766) (-2.700,1.742) (-2.650,1.718) (-2.600,1.694) (-2.550,1.669) (-2.500,1.643) (-2.450,1.617) (-2.400,1.590) (-2.350,1.563) (-2.300,1.535) (-2.250,1.507) (-2.200,1.479) (-2.150,1.450) (-2.100,1.421) (-2.050,1.391) (-2.000,1.361) (-1.950,1.331) (-1.900,1.300) (-1.850,1.269) (-1.800,1.238) (-1.750,1.206) (-1.702,1.175) (-1.650,1.142) (-1.600,1.109) (-1.550,1.076) (-1.500,1.044) (-1.450,1.010) (-1.400,0.977) (-1.350,0.943) (-1.300,0.910) (-1.250,0.876) (-1.200,0.842) (-1.150,0.807) (-1.103,0.775) (-1.050,0.739) (-1.000,0.704) (-0.950,0.669) (-0.900,0.634) (-0.851,0.600) (-0.800,0.565) (-0.750,0.529) (-0.709,0.500) (-0.650,0.459) (-0.602,0.425) (-0.550,0.389) (-0.500,0.354) (-0.450,0.319) (-0.400,0.283) (-0.354,0.250) (-0.300,0.213) (-0.250,0.177) (-0.220,0.155)}
\EQ{(0.206,0.144) (0.248,0.175) (0.300,0.213) (0.350,0.248) (0.400,0.283) (0.450,0.319) (0.496,0.350) (0.550,0.389) (0.600,0.424) (0.650,0.459) (0.700,0.494) (0.744,0.525) (0.800,0.565) (0.850,0.599) (0.900,0.634) (0.950,0.669) (0.995,0.700) (1.050,0.739) (1.100,0.773) (1.150,0.807) (1.200,0.842) (1.249,0.875) (1.300,0.910) (1.350,0.943) (1.397,0.975) (1.450,1.010) (1.500,1.044) (1.548,1.075) (1.600,1.109) (1.650,1.142) (1.700,1.174) (1.750,1.206) (1.800,1.238) (1.850,1.269) (1.900,1.300) (1.950,1.331) (2.000,1.361) (2.050,1.391) (2.100,1.421) (2.150,1.450) (2.193,1.475) (2.237,1.500) (2.282,1.525) (2.327,1.550) (2.372,1.575) (2.419,1.600) (2.466,1.625) (2.514,1.650) (2.563,1.675) (2.613,1.700) (2.664,1.725) (2.716,1.750) (2.770,1.775) (2.825,1.800) (2.875,1.822) (2.925,1.844) (2.975,1.865) (3.025,1.885) (3.075,1.904) (3.125,1.923) (3.175,1.942) (3.225,1.959) (3.275,1.976) (3.325,1.993) (3.375,2.008) (3.425,2.023) (3.475,2.037) (3.525,2.051) (3.575,2.064) (3.625,2.076) (3.675,2.087) (3.725,2.097) (3.775,2.107) (3.825,2.116) (3.875,2.124) (3.925,2.131) (3.975,2.138) (4.025,2.143) (4.075,2.148) (4.125,2.152) (4.175,2.155) (4.225,2.157) (4.275,2.158) (4.325,2.158) (4.375,2.158) (4.425,2.156) (4.475,2.154) (4.525,2.150) (4.575,2.146) (4.625,2.140) (4.675,2.134) (4.725,2.126) (4.775,2.117) (4.825,2.108) (4.875,2.097) (4.925,2.085) (4.975,2.071) (5.025,2.057) (5.075,2.041) (5.125,2.024) (5.175,2.006) (5.225,1.986) (5.275,1.965) (5.325,1.942) (5.375,1.918) (5.425,1.892) (5.475,1.864) (5.525,1.835) (5.575,1.804) (5.618,1.775) (5.675,1.735) (5.722,1.700) (5.775,1.658) (5.814,1.625) (5.869,1.575) (5.921,1.525) (5.970,1.475) (6.016,1.425) (6.050,1.385) (6.098,1.325) (6.136,1.275) (6.172,1.225) (6.200,1.183) (6.236,1.125) (6.266,1.075) (6.294,1.025) (6.320,0.975) (6.344,0.925) (6.367,0.875) (6.388,0.825) (6.408,0.775) (6.426,0.725) (6.443,0.675) (6.459,0.625) (6.473,0.575) (6.486,0.525) (6.498,0.475) (6.509,0.425) (6.518,0.375) (6.526,0.325) (6.533,0.275) (6.539,0.225) (6.543,0.175) (6.547,0.125) (6.549,0.075) (6.550,0.025) (6.550,-0.037) (6.548,-0.100) (6.545,-0.150) (6.541,-0.200) (6.536,-0.250) (6.530,-0.300) (6.522,-0.350) (6.513,-0.400) (6.503,-0.450) (6.492,-0.500) (6.480,-0.550) (6.466,-0.600) (6.451,-0.650) (6.435,-0.700) (6.417,-0.750) (6.398,-0.800) (6.378,-0.850) (6.356,-0.900) (6.332,-0.950) (6.307,-1.000) (6.280,-1.050) (6.251,-1.100) (6.225,-1.143) (6.189,-1.200) (6.154,-1.250) (6.118,-1.300) (6.079,-1.350) (6.037,-1.400) (6.000,-1.442) (5.950,-1.496) (5.900,-1.546) (5.850,-1.593) (5.800,-1.637) (5.754,-1.675) (5.700,-1.717) (5.654,-1.750) (5.600,-1.787) (5.550,-1.819) (5.500,-1.850) (5.456,-1.875) (5.409,-1.900) (5.360,-1.925) (5.307,-1.950) (5.251,-1.975) (5.200,-1.996) (5.150,-2.015) (5.100,-2.033) (5.050,-2.049) (5.000,-2.064) (4.950,-2.078) (4.900,-2.091) (4.850,-2.102) (4.800,-2.113) (4.750,-2.122) (4.700,-2.130) (4.650,-2.137) (4.600,-2.143) (4.550,-2.148) (4.500,-2.152) (4.450,-2.155) (4.400,-2.157) (4.350,-2.158) (4.300,-2.158) (4.250,-2.158) (4.200,-2.156) (4.150,-2.153) (4.100,-2.150) (4.050,-2.146) (4.000,-2.140) (3.950,-2.134) (3.900,-2.128) (3.850,-2.120) (3.800,-2.111) (3.750,-2.102) (3.700,-2.092) (3.650,-2.081) (3.600,-2.070) (3.550,-2.057) (3.500,-2.044) (3.450,-2.030) (3.400,-2.016) (3.350,-2.001) (3.300,-1.985) (3.250,-1.968) (3.200,-1.951) (3.150,-1.933) (3.100,-1.914) (3.050,-1.895) (3.001,-1.875) (2.950,-1.854) (2.900,-1.833) (2.850,-1.811) (2.800,-1.789) (2.750,-1.766) (2.700,-1.742) (2.650,-1.718) (2.600,-1.694) (2.550,-1.669) (2.500,-1.643) (2.450,-1.617) (2.400,-1.590) (2.350,-1.563) (2.300,-1.535) (2.250,-1.507) (2.200,-1.479) (2.150,-1.450) (2.100,-1.421) (2.050,-1.391) (2.000,-1.361) (1.950,-1.331) (1.900,-1.300) (1.850,-1.269) (1.800,-1.238) (1.750,-1.206) (1.702,-1.175) (1.650,-1.142) (1.600,-1.109) (1.550,-1.076) (1.500,-1.044) (1.450,-1.010) (1.400,-0.977) (1.350,-0.943) (1.300,-0.910) (1.250,-0.876) (1.200,-0.842) (1.150,-0.807) (1.103,-0.775) (1.050,-0.739) (1.000,-0.704) (0.950,-0.669) (0.900,-0.634) (0.851,-0.600) (0.800,-0.565) (0.750,-0.529) (0.709,-0.500) (0.650,-0.459) (0.602,-0.425) (0.550,-0.389) (0.500,-0.354) (0.450,-0.319) (0.400,-0.283) (0.354,-0.250) (0.300,-0.213) (0.250,-0.177) (0.206,-0.144)}
\EQ{(-0.220,-0.155) (-0.275,-0.195) (-0.319,-0.225) (-0.375,-0.266) (-0.424,-0.300) (-0.475,-0.336) (-0.525,-0.371) (-0.567,-0.400) (-0.625,-0.442) (-0.673,-0.475) (-0.725,-0.512) (-0.775,-0.547) (-0.825,-0.583) (-0.875,-0.618) (-0.922,-0.650) (-0.975,-0.688) (-1.025,-0.722) (-1.075,-0.757) (-1.125,-0.792) (-1.173,-0.825) (-1.225,-0.861) (-1.275,-0.896) (-1.318,-0.925) (-1.375,-0.964) (-1.425,-0.998) (-1.475,-1.032) (-1.525,-1.066) (-1.575,-1.100) (-1.625,-1.133) (-1.675,-1.167) (-1.725,-1.200) (-1.775,-1.233) (-1.825,-1.266) (-1.875,-1.298) (-1.925,-1.330) (-1.975,-1.362) (-2.025,-1.394) (-2.074,-1.425) (-2.125,-1.457) (-2.175,-1.488) (-2.225,-1.519) (-2.275,-1.549) (-2.325,-1.580) (-2.375,-1.609) (-2.425,-1.639) (-2.475,-1.668) (-2.525,-1.697) (-2.575,-1.725) (-2.619,-1.750) (-2.665,-1.775) (-2.710,-1.800) (-2.757,-1.825) (-2.803,-1.850) (-2.850,-1.874) (-2.900,-1.900) (-2.949,-1.925) (-2.999,-1.950) (-3.050,-1.975) (-3.100,-1.999) (-3.150,-2.022) (-3.200,-2.045) (-3.250,-2.068) (-3.300,-2.090) (-3.350,-2.111) (-3.400,-2.132) (-3.450,-2.152) (-3.500,-2.172) (-3.550,-2.191) (-3.600,-2.210) (-3.650,-2.228) (-3.700,-2.246) (-3.750,-2.263) (-3.800,-2.279) (-3.850,-2.295) (-3.900,-2.310) (-3.950,-2.324) (-4.000,-2.338) (-4.050,-2.352) (-4.100,-2.364) (-4.150,-2.376) (-4.200,-2.387) (-4.250,-2.398) (-4.300,-2.408) (-4.350,-2.417) (-4.400,-2.425) (-4.450,-2.433) (-4.500,-2.440) (-4.550,-2.447) (-4.600,-2.452) (-4.650,-2.457) (-4.700,-2.462) (-4.750,-2.465) (-4.800,-2.468) (-4.850,-2.469) (-4.900,-2.470) (-4.950,-2.471) (-5.000,-2.470) (-5.050,-2.469) (-5.100,-2.467) (-5.150,-2.464) (-5.200,-2.460) (-5.250,-2.455) (-5.300,-2.449) (-5.350,-2.443) (-5.400,-2.435) (-5.450,-2.427) (-5.500,-2.417) (-5.550,-2.407) (-5.600,-2.395) (-5.650,-2.383) (-5.700,-2.370) (-5.750,-2.355) (-5.800,-2.339) (-5.850,-2.323) (-5.900,-2.305) (-5.950,-2.286) (-6.000,-2.265) (-6.050,-2.244) (-6.100,-2.221) (-6.150,-2.196) (-6.200,-2.171) (-6.250,-2.144) (-6.300,-2.115) (-6.350,-2.085) (-6.400,-2.053) (-6.450,-2.019) (-6.500,-1.983) (-6.545,-1.950) (-6.600,-1.906) (-6.650,-1.864) (-6.695,-1.825) (-6.748,-1.775) (-6.799,-1.725) (-6.846,-1.675) (-6.892,-1.625) (-6.925,-1.586) (-6.974,-1.525) (-7.013,-1.475) (-7.049,-1.425) (-7.083,-1.375) (-7.116,-1.325) (-7.147,-1.275) (-7.175,-1.226) (-7.200,-1.181) (-7.225,-1.134) (-7.250,-1.084) (-7.275,-1.031) (-7.300,-0.975) (-7.320,-0.925) (-7.340,-0.875) (-7.358,-0.825) (-7.375,-0.775) (-7.391,-0.725) (-7.405,-0.675) (-7.419,-0.625) (-7.431,-0.575) (-7.443,-0.525) (-7.453,-0.475) (-7.462,-0.425) (-7.470,-0.375) (-7.477,-0.325) (-7.483,-0.275) (-7.488,-0.225) (-7.492,-0.175) (-7.495,-0.125) (-7.497,-0.075) (-7.498,-0.025) (-7.498,0.050) (-7.496,0.100) (-7.494,0.150) (-7.490,0.200) (-7.486,0.250) (-7.480,0.300) (-7.474,0.350) (-7.466,0.400) (-7.457,0.450) (-7.448,0.500) (-7.437,0.550) (-7.425,0.600) (-7.412,0.650) (-7.398,0.700) (-7.383,0.750) (-7.367,0.800) (-7.349,0.850) (-7.330,0.900) (-7.310,0.950) (-7.289,1.000) (-7.266,1.050) (-7.242,1.100) (-7.217,1.150) (-7.190,1.200) (-7.161,1.250) (-7.131,1.300) (-7.100,1.350) (-7.066,1.400) (-7.031,1.450) (-7.000,1.492) (-6.955,1.550) (-6.913,1.600) (-6.875,1.644) (-6.825,1.698) (-6.775,1.749) (-6.725,1.797) (-6.675,1.843) (-6.625,1.886) (-6.577,1.925) (-6.525,1.965) (-6.477,2.000) (-6.425,2.036) (-6.375,2.069) (-6.325,2.100) (-6.275,2.129) (-6.225,2.157) (-6.175,2.184) (-6.125,2.209) (-6.075,2.232) (-6.025,2.255) (-5.977,2.275) (-5.925,2.295) (-5.875,2.314) (-5.825,2.331) (-5.775,2.347) (-5.725,2.362) (-5.675,2.376) (-5.625,2.389) (-5.575,2.401) (-5.525,2.412) (-5.475,2.422) (-5.425,2.431) (-5.375,2.439) (-5.325,2.446) (-5.275,2.452) (-5.225,2.457) (-5.175,2.462) (-5.125,2.465) (-5.075,2.468) (-5.025,2.470) (-4.975,2.471) (-4.925,2.471) (-4.875,2.470) (-4.825,2.469) (-4.775,2.466) (-4.725,2.463) (-4.675,2.460) (-4.625,2.455) (-4.575,2.450) (-4.525,2.444) (-4.475,2.437) (-4.425,2.430) (-4.375,2.421) (-4.325,2.412) (-4.275,2.403) (-4.225,2.393) (-4.175,2.382) (-4.125,2.370) (-4.075,2.358) (-4.025,2.345) (-3.975,2.331) (-3.925,2.317) (-3.875,2.302) (-3.825,2.287) (-3.775,2.271) (-3.725,2.254) (-3.675,2.237) (-3.625,2.219) (-3.575,2.201) (-3.525,2.182) (-3.475,2.162) (-3.425,2.142) (-3.375,2.122) (-3.325,2.100) (-3.275,2.079) (-3.225,2.057) (-3.175,2.034) (-3.125,2.011) (-3.075,1.987) (-3.025,1.963) (-2.975,1.938) (-2.925,1.913) (-2.875,1.887) (-2.825,1.861) (-2.775,1.835) (-2.725,1.808) (-2.675,1.781) (-2.625,1.753) (-2.575,1.725) (-2.531,1.700) (-2.487,1.675) (-2.425,1.639) (-2.375,1.609) (-2.325,1.580) (-2.276,1.550) (-2.225,1.519) (-2.175,1.488) (-2.125,1.457) (-2.075,1.426) (-2.025,1.394) (-1.975,1.362) (-1.925,1.330) (-1.878,1.300) (-1.825,1.266) (-1.775,1.233) (-1.725,1.200) (-1.675,1.167) (-1.625,1.133) (-1.576,1.100) (-1.525,1.066) (-1.475,1.032) (-1.428,1.000) (-1.375,0.964) (-1.325,0.930) (-1.282,0.900) (-1.225,0.861) (-1.175,0.827) (-1.125,0.792) (-1.075,0.757) (-1.029,0.725) (-0.975,0.688) (-0.925,0.652) (-0.875,0.618) (-0.825,0.583) (-0.779,0.550) (-0.725,0.512) (-0.675,0.477) (-0.625,0.442) (-0.575,0.407) (-0.531,0.375) (-0.475,0.336) (-0.425,0.300) (-0.375,0.266) (-0.325,0.230) (-0.284,0.200) (-0.225,0.160) (-0.220,0.155)}
\EQ{(0.206,0.144) (0.248,0.175) (0.300,0.213) (0.350,0.248) (0.400,0.283) (0.450,0.319) (0.496,0.350) (0.550,0.389) (0.600,0.424) (0.650,0.460) (0.700,0.495) (0.744,0.525) (0.800,0.565) (0.850,0.600) (0.900,0.635) (0.950,0.670) (0.993,0.700) (1.050,0.740) (1.100,0.774) (1.150,0.809) (1.200,0.844) (1.245,0.875) (1.300,0.913) (1.350,0.947) (1.400,0.981) (1.450,1.015) (1.500,1.049) (1.550,1.083) (1.600,1.117) (1.650,1.150) (1.700,1.183) (1.750,1.216) (1.800,1.249) (1.850,1.282) (1.900,1.314) (1.950,1.346) (1.995,1.375) (2.050,1.410) (2.100,1.442) (2.150,1.473) (2.194,1.500) (2.250,1.534) (2.300,1.565) (2.350,1.594) (2.400,1.624) (2.444,1.650) (2.500,1.682) (2.550,1.711) (2.600,1.739) (2.650,1.767) (2.700,1.795) (2.750,1.822) (2.800,1.848) (2.850,1.874) (2.900,1.900) (2.949,1.925) (2.999,1.950) (3.050,1.975) (3.100,1.999) (3.150,2.022) (3.200,2.045) (3.250,2.068) (3.300,2.090) (3.350,2.111) (3.400,2.132) (3.450,2.152) (3.500,2.172) (3.550,2.191) (3.600,2.210) (3.650,2.228) (3.700,2.246) (3.750,2.263) (3.800,2.279) (3.850,2.295) (3.900,2.310) (3.950,2.324) (4.000,2.338) (4.050,2.352) (4.100,2.364) (4.150,2.376) (4.200,2.387) (4.250,2.398) (4.300,2.408) (4.350,2.417) (4.400,2.425) (4.450,2.433) (4.500,2.440) (4.550,2.447) (4.600,2.452) (4.650,2.457) (4.700,2.462) (4.750,2.465) (4.800,2.468) (4.850,2.469) (4.900,2.470) (4.950,2.471) (5.000,2.470) (5.050,2.469) (5.100,2.467) (5.150,2.464) (5.200,2.460) (5.250,2.455) (5.300,2.449) (5.350,2.443) (5.400,2.435) (5.450,2.427) (5.500,2.417) (5.550,2.407) (5.600,2.395) (5.650,2.383) (5.700,2.370) (5.750,2.355) (5.800,2.339) (5.850,2.323) (5.900,2.305) (5.950,2.286) (6.000,2.265) (6.050,2.244) (6.100,2.221) (6.150,2.196) (6.200,2.171) (6.250,2.144) (6.300,2.115) (6.350,2.085) (6.400,2.053) (6.450,2.019) (6.500,1.983) (6.545,1.950) (6.600,1.906) (6.650,1.864) (6.695,1.825) (6.748,1.775) (6.799,1.725) (6.846,1.675) (6.892,1.625) (6.925,1.586) (6.974,1.525) (7.013,1.475) (7.049,1.425) (7.083,1.375) (7.116,1.325) (7.147,1.275) (7.175,1.226) (7.200,1.181) (7.225,1.134) (7.250,1.084) (7.275,1.031) (7.300,0.975) (7.320,0.925) (7.340,0.875) (7.358,0.825) (7.375,0.775) (7.391,0.725) (7.405,0.675) (7.419,0.625) (7.431,0.575) (7.443,0.525) (7.453,0.475) (7.462,0.425) (7.470,0.375) (7.477,0.325) (7.483,0.275) (7.488,0.225) (7.492,0.175) (7.495,0.125) (7.497,0.075) (7.498,0.025) (7.498,-0.050) (7.496,-0.100) (7.494,-0.150) (7.490,-0.200) (7.486,-0.250) (7.480,-0.300) (7.474,-0.350) (7.466,-0.400) (7.457,-0.450) (7.448,-0.500) (7.437,-0.550) (7.425,-0.600) (7.412,-0.650) (7.398,-0.700) (7.383,-0.750) (7.367,-0.800) (7.349,-0.850) (7.330,-0.900) (7.310,-0.950) (7.289,-1.000) (7.266,-1.050) (7.242,-1.100) (7.217,-1.150) (7.190,-1.200) (7.161,-1.250) (7.131,-1.300) (7.100,-1.350) (7.066,-1.400) (7.031,-1.450) (7.000,-1.492) (6.955,-1.550) (6.913,-1.600) (6.875,-1.644) (6.825,-1.698) (6.775,-1.749) (6.725,-1.797) (6.675,-1.843) (6.625,-1.886) (6.577,-1.925) (6.525,-1.965) (6.477,-2.000) (6.425,-2.036) (6.375,-2.069) (6.325,-2.100) (6.275,-2.129) (6.225,-2.157) (6.175,-2.184) (6.125,-2.209) (6.075,-2.232) (6.025,-2.255) (5.977,-2.275) (5.925,-2.295) (5.875,-2.314) (5.825,-2.331) (5.775,-2.347) (5.725,-2.362) (5.675,-2.376) (5.625,-2.389) (5.575,-2.401) (5.525,-2.412) (5.475,-2.422) (5.425,-2.431) (5.375,-2.439) (5.325,-2.446) (5.275,-2.452) (5.225,-2.457) (5.175,-2.462) (5.125,-2.465) (5.075,-2.468) (5.025,-2.470) (4.975,-2.471) (4.925,-2.471) (4.875,-2.470) (4.825,-2.469) (4.775,-2.466) (4.725,-2.463) (4.675,-2.460) (4.625,-2.455) (4.575,-2.450) (4.525,-2.444) (4.475,-2.437) (4.425,-2.430) (4.375,-2.421) (4.325,-2.412) (4.275,-2.403) (4.225,-2.393) (4.175,-2.382) (4.125,-2.370) (4.075,-2.358) (4.025,-2.345) (3.975,-2.331) (3.925,-2.317) (3.875,-2.302) (3.825,-2.287) (3.775,-2.271) (3.725,-2.254) (3.675,-2.237) (3.625,-2.219) (3.575,-2.201) (3.525,-2.182) (3.475,-2.162) (3.425,-2.142) (3.375,-2.122) (3.325,-2.100) (3.275,-2.079) (3.225,-2.057) (3.175,-2.034) (3.125,-2.011) (3.075,-1.987) (3.025,-1.963) (2.975,-1.938) (2.925,-1.913) (2.875,-1.887) (2.825,-1.861) (2.775,-1.835) (2.725,-1.808) (2.675,-1.781) (2.625,-1.753) (2.575,-1.725) (2.531,-1.700) (2.487,-1.675) (2.425,-1.639) (2.375,-1.609) (2.325,-1.580) (2.276,-1.550) (2.225,-1.519) (2.175,-1.488) (2.125,-1.457) (2.075,-1.426) (2.025,-1.394) (1.975,-1.362) (1.925,-1.330) (1.878,-1.300) (1.825,-1.266) (1.775,-1.233) (1.725,-1.200) (1.675,-1.167) (1.625,-1.133) (1.576,-1.100) (1.525,-1.066) (1.475,-1.032) (1.428,-1.000) (1.375,-0.964) (1.325,-0.930) (1.282,-0.900) (1.225,-0.861) (1.175,-0.827) (1.125,-0.792) (1.075,-0.757) (1.029,-0.725) (0.975,-0.688) (0.925,-0.652) (0.875,-0.618) (0.825,-0.583) (0.779,-0.550) (0.725,-0.512) (0.675,-0.477) (0.625,-0.442) (0.575,-0.407) (0.531,-0.375) (0.475,-0.336) (0.425,-0.300) (0.375,-0.266) (0.325,-0.230) (0.284,-0.200) (0.225,-0.160) (0.206,-0.144)}
\EQ{(-0.220,-0.155) (-0.275,-0.195) (-0.319,-0.225) (-0.375,-0.266) (-0.424,-0.300) (-0.475,-0.336) (-0.525,-0.371) (-0.567,-0.400) (-0.625,-0.442) (-0.672,-0.475) (-0.725,-0.513) (-0.775,-0.548) (-0.825,-0.583) (-0.875,-0.618) (-0.921,-0.650) (-0.975,-0.689) (-1.025,-0.724) (-1.075,-0.759) (-1.125,-0.794) (-1.170,-0.825) (-1.225,-0.864) (-1.275,-0.899) (-1.325,-0.934) (-1.375,-0.968) (-1.421,-1.000) (-1.475,-1.038) (-1.525,-1.072) (-1.575,-1.107) (-1.625,-1.141) (-1.674,-1.175) (-1.725,-1.210) (-1.775,-1.244) (-1.821,-1.275) (-1.875,-1.312) (-1.925,-1.346) (-1.969,-1.375) (-2.025,-1.413) (-2.075,-1.446) (-2.118,-1.475) (-2.175,-1.513) (-2.225,-1.546) (-2.270,-1.575) (-2.325,-1.611) (-2.375,-1.644) (-2.424,-1.675) (-2.475,-1.708) (-2.525,-1.740) (-2.575,-1.772) (-2.620,-1.800) (-2.675,-1.834) (-2.725,-1.866) (-2.775,-1.896) (-2.822,-1.925) (-2.875,-1.957) (-2.925,-1.988) (-2.975,-2.017) (-3.025,-2.047) (-3.073,-2.075) (-3.125,-2.105) (-3.175,-2.134) (-3.225,-2.163) (-3.275,-2.191) (-3.325,-2.219) (-3.375,-2.246) (-3.425,-2.274) (-3.474,-2.300) (-3.521,-2.325) (-3.569,-2.350) (-3.617,-2.375) (-3.665,-2.400) (-3.715,-2.425) (-3.765,-2.450) (-3.816,-2.475) (-3.867,-2.500) (-3.920,-2.525) (-3.973,-2.550) (-4.025,-2.574) (-4.075,-2.596) (-4.125,-2.619) (-4.175,-2.640) (-4.225,-2.662) (-4.275,-2.683) (-4.325,-2.703) (-4.375,-2.723) (-4.425,-2.743) (-4.475,-2.762) (-4.525,-2.781) (-4.575,-2.799) (-4.625,-2.817) (-4.675,-2.834) (-4.725,-2.851) (-4.775,-2.867) (-4.825,-2.883) (-4.875,-2.899) (-4.925,-2.914) (-4.975,-2.928) (-5.025,-2.942) (-5.075,-2.956) (-5.125,-2.968) (-5.175,-2.981) (-5.225,-2.993) (-5.275,-3.004) (-5.325,-3.015) (-5.375,-3.025) (-5.425,-3.035) (-5.475,-3.044) (-5.525,-3.053) (-5.575,-3.061) (-5.625,-3.069) (-5.675,-3.075) (-5.725,-3.082) (-5.775,-3.088) (-5.825,-3.093) (-5.875,-3.098) (-5.925,-3.102) (-5.975,-3.105) (-6.025,-3.108) (-6.075,-3.110) (-6.125,-3.112) (-6.175,-3.113) (-6.225,-3.113) (-6.275,-3.113) (-6.325,-3.112) (-6.375,-3.110) (-6.425,-3.108) (-6.475,-3.105) (-6.525,-3.101) (-6.575,-3.097) (-6.625,-3.092) (-6.675,-3.086) (-6.725,-3.080) (-6.775,-3.072) (-6.825,-3.065) (-6.875,-3.056) (-6.925,-3.046) (-6.975,-3.036) (-7.025,-3.025) (-7.075,-3.013) (-7.125,-3.001) (-7.175,-2.987) (-7.225,-2.973) (-7.275,-2.958) (-7.325,-2.942) (-7.375,-2.925) (-7.425,-2.907) (-7.475,-2.888) (-7.525,-2.868) (-7.575,-2.848) (-7.625,-2.826) (-7.675,-2.803) (-7.725,-2.779) (-7.775,-2.754) (-7.825,-2.728) (-7.875,-2.701) (-7.920,-2.675) (-7.975,-2.642) (-8.025,-2.611) (-8.075,-2.579) (-8.117,-2.550) (-8.175,-2.509) (-8.222,-2.475) (-8.275,-2.434) (-8.318,-2.400) (-8.375,-2.352) (-8.425,-2.308) (-8.475,-2.262) (-8.514,-2.225) (-8.550,-2.189) (-8.589,-2.150) (-8.625,-2.111) (-8.675,-2.056) (-8.723,-2.000) (-8.765,-1.950) (-8.800,-1.905) (-8.842,-1.850) (-8.875,-1.804) (-8.912,-1.750) (-8.946,-1.700) (-8.975,-1.654) (-9.008,-1.600) (-9.037,-1.550) (-9.065,-1.500) (-9.092,-1.450) (-9.118,-1.400) (-9.142,-1.350) (-9.165,-1.300) (-9.188,-1.250) (-9.209,-1.200) (-9.229,-1.150) (-9.248,-1.100) (-9.267,-1.050) (-9.284,-1.000) (-9.300,-0.951) (-9.316,-0.900) (-9.330,-0.850) (-9.344,-0.800) (-9.357,-0.750) (-9.368,-0.700) (-9.379,-0.650) (-9.390,-0.600) (-9.399,-0.550) (-9.407,-0.500) (-9.415,-0.450) (-9.422,-0.400) (-9.428,-0.350) (-9.433,-0.300) (-9.437,-0.250) (-9.441,-0.200) (-9.444,-0.150) (-9.446,-0.100) (-9.447,-0.050) (-9.447,0.000) (-9.447,0.050) (-9.446,0.100) (-9.444,0.150) (-9.441,0.200) (-9.437,0.250) (-9.433,0.300) (-9.428,0.350) (-9.422,0.400) (-9.415,0.450) (-9.407,0.500) (-9.399,0.550) (-9.390,0.600) (-9.379,0.650) (-9.368,0.700) (-9.357,0.750) (-9.344,0.800) (-9.330,0.850) (-9.316,0.900) (-9.300,0.950) (-9.284,1.000) (-9.267,1.050) (-9.248,1.100) (-9.229,1.150) (-9.209,1.200) (-9.188,1.250) (-9.165,1.300) (-9.142,1.350) (-9.118,1.400) (-9.092,1.450) (-9.065,1.500) (-9.037,1.550) (-9.008,1.600) (-8.978,1.650) (-8.950,1.693) (-8.912,1.750) (-8.878,1.800) (-8.842,1.850) (-8.804,1.900) (-8.765,1.950) (-8.725,1.998) (-8.680,2.050) (-8.636,2.100) (-8.600,2.138) (-8.564,2.175) (-8.525,2.214) (-8.488,2.250) (-8.450,2.285) (-8.406,2.325) (-8.350,2.373) (-8.300,2.414) (-8.255,2.450) (-8.200,2.491) (-8.153,2.525) (-8.100,2.562) (-8.050,2.595) (-8.003,2.625) (-7.950,2.657) (-7.900,2.686) (-7.850,2.714) (-7.800,2.741) (-7.750,2.767) (-7.700,2.791) (-7.650,2.814) (-7.600,2.837) (-7.550,2.858) (-7.500,2.878) (-7.450,2.898) (-7.400,2.916) (-7.350,2.933) (-7.300,2.950) (-7.250,2.966) (-7.200,2.980) (-7.150,2.994) (-7.100,3.007) (-7.050,3.019) (-7.000,3.031) (-6.950,3.041) (-6.900,3.051) (-6.850,3.060) (-6.800,3.069) (-6.750,3.076) (-6.700,3.083) (-6.650,3.089) (-6.600,3.094) (-6.550,3.099) (-6.500,3.103) (-6.450,3.106) (-6.400,3.109) (-6.350,3.111) (-6.300,3.112) (-6.250,3.113) (-6.200,3.113) (-6.150,3.112) (-6.100,3.111) (-6.050,3.109) (-6.000,3.107) (-5.950,3.103) (-5.900,3.100) (-5.850,3.095) (-5.800,3.090) (-5.750,3.085) (-5.700,3.079) (-5.650,3.072) (-5.600,3.065) (-5.550,3.057) (-5.500,3.049) (-5.450,3.040) (-5.400,3.030) (-5.350,3.020) (-5.300,3.010) (-5.250,2.998) (-5.200,2.987) (-5.151,2.975) (-5.100,2.962) (-5.050,2.949) (-5.000,2.935) (-4.950,2.921) (-4.900,2.906) (-4.850,2.891) (-4.800,2.875) (-4.750,2.859) (-4.700,2.843) (-4.650,2.825) (-4.600,2.808) (-4.550,2.790) (-4.500,2.771) (-4.450,2.752) (-4.400,2.733) (-4.350,2.713) (-4.300,2.693) (-4.250,2.672) (-4.200,2.651) (-4.150,2.629) (-4.100,2.608) (-4.050,2.585) (-4.000,2.562) (-3.950,2.539) (-3.900,2.516) (-3.850,2.492) (-3.800,2.467) (-3.750,2.443) (-3.700,2.418) (-3.650,2.392) (-3.600,2.366) (-3.550,2.340) (-3.500,2.314) (-3.450,2.287) (-3.400,2.260) (-3.350,2.233) (-3.300,2.205) (-3.250,2.177) (-3.203,2.150) (-3.159,2.125) (-3.116,2.100) (-3.050,2.062) (-3.000,2.032) (-2.950,2.002) (-2.904,1.975) (-2.850,1.942) (-2.800,1.912) (-2.750,1.881) (-2.700,1.850) (-2.650,1.819) (-2.600,1.787) (-2.550,1.756) (-2.502,1.725) (-2.450,1.692) (-2.400,1.660) (-2.350,1.627) (-2.300,1.595) (-2.250,1.562) (-2.200,1.529) (-2.156,1.500) (-2.100,1.463) (-2.050,1.430) (-2.006,1.400) (-1.950,1.363) (-1.900,1.329) (-1.858,1.300) (-1.800,1.261) (-1.750,1.227) (-1.700,1.193) (-1.650,1.158) (-1.602,1.125) (-1.550,1.090) (-1.500,1.055) (-1.457,1.025) (-1.400,0.986) (-1.350,0.951) (-1.300,0.916) (-1.250,0.881) (-1.206,0.850) (-1.150,0.811) (-1.100,0.776) (-1.050,0.741) (-1.000,0.706) (-0.956,0.675) (-0.900,0.636) (-0.850,0.600) (-0.800,0.566) (-0.750,0.530) (-0.708,0.500) (-0.650,0.460) (-0.601,0.425) (-0.550,0.389) (-0.500,0.354) (-0.450,0.319) (-0.400,0.283) (-0.354,0.250) (-0.300,0.213) (-0.250,0.177) (-0.220,0.155)}
\EQ{(0.206,0.144) (0.248,0.175) (0.300,0.213) (0.350,0.248) (0.400,0.283) (0.450,0.319) (0.496,0.350) (0.550,0.389) (0.600,0.424) (0.650,0.460) (0.700,0.495) (0.743,0.525) (0.800,0.566) (0.849,0.600) (0.900,0.636) (0.950,0.671) (0.992,0.700) (1.050,0.741) (1.098,0.775) (1.150,0.811) (1.200,0.846) (1.242,0.875) (1.300,0.916) (1.349,0.950) (1.400,0.986) (1.450,1.020) (1.493,1.050) (1.550,1.090) (1.600,1.124) (1.650,1.158) (1.700,1.193) (1.748,1.225) (1.800,1.261) (1.850,1.295) (1.895,1.325) (1.950,1.363) (2.000,1.396) (2.043,1.425) (2.100,1.463) (2.150,1.496) (2.194,1.525) (2.250,1.562) (2.300,1.595) (2.347,1.625) (2.400,1.660) (2.450,1.692) (2.500,1.724) (2.550,1.756) (2.600,1.787) (2.650,1.819) (2.700,1.850) (2.750,1.881) (2.800,1.912) (2.850,1.942) (2.900,1.972) (2.946,2.000) (3.000,2.032) (3.050,2.062) (3.100,2.091) (3.150,2.120) (3.200,2.148) (3.247,2.175) (3.292,2.200) (3.337,2.225) (3.382,2.250) (3.428,2.275) (3.474,2.300) (3.521,2.325) (3.569,2.350) (3.617,2.375) (3.665,2.400) (3.715,2.425) (3.765,2.450) (3.816,2.475) (3.867,2.500) (3.920,2.525) (3.973,2.550) (4.025,2.574) (4.075,2.596) (4.125,2.619) (4.175,2.640) (4.225,2.662) (4.275,2.683) (4.325,2.703) (4.375,2.723) (4.425,2.743) (4.475,2.762) (4.525,2.781) (4.575,2.799) (4.625,2.817) (4.675,2.834) (4.725,2.851) (4.775,2.867) (4.825,2.883) (4.875,2.899) (4.925,2.914) (4.975,2.928) (5.025,2.942) (5.075,2.956) (5.125,2.968) (5.175,2.981) (5.225,2.993) (5.275,3.004) (5.325,3.015) (5.375,3.025) (5.425,3.035) (5.475,3.044) (5.525,3.053) (5.575,3.061) (5.625,3.069) (5.675,3.075) (5.725,3.082) (5.775,3.088) (5.825,3.093) (5.875,3.098) (5.925,3.102) (5.975,3.105) (6.025,3.108) (6.075,3.110) (6.125,3.112) (6.175,3.113) (6.225,3.113) (6.275,3.113) (6.325,3.112) (6.375,3.110) (6.425,3.108) (6.475,3.105) (6.525,3.101) (6.575,3.097) (6.625,3.092) (6.675,3.086) (6.725,3.080) (6.775,3.072) (6.825,3.065) (6.875,3.056) (6.925,3.046) (6.975,3.036) (7.025,3.025) (7.075,3.013) (7.125,3.001) (7.175,2.987) (7.225,2.973) (7.275,2.958) (7.325,2.942) (7.375,2.925) (7.425,2.907) (7.475,2.888) (7.525,2.868) (7.575,2.848) (7.625,2.826) (7.675,2.803) (7.725,2.779) (7.775,2.754) (7.825,2.728) (7.875,2.701) (7.920,2.675) (7.975,2.642) (8.025,2.611) (8.075,2.579) (8.117,2.550) (8.175,2.509) (8.222,2.475) (8.275,2.434) (8.318,2.400) (8.375,2.352) (8.425,2.308) (8.475,2.262) (8.514,2.225) (8.550,2.189) (8.589,2.150) (8.625,2.111) (8.675,2.056) (8.723,2.000) (8.765,1.950) (8.800,1.905) (8.842,1.850) (8.875,1.804) (8.912,1.750) (8.946,1.700) (8.975,1.654) (9.008,1.600) (9.037,1.550) (9.065,1.500) (9.092,1.450) (9.118,1.400) (9.142,1.350) (9.165,1.300) (9.188,1.250) (9.209,1.200) (9.229,1.150) (9.248,1.100) (9.267,1.050) (9.284,1.000) (9.300,0.951) (9.316,0.900) (9.330,0.850) (9.344,0.800) (9.357,0.750) (9.368,0.700) (9.379,0.650) (9.390,0.600) (9.399,0.550) (9.407,0.500) (9.415,0.450) (9.422,0.400) (9.428,0.350) (9.433,0.300) (9.437,0.250) (9.441,0.200) (9.444,0.150) (9.446,0.100) (9.447,0.050) (9.447,0.000) (9.447,-0.050) (9.446,-0.100) (9.444,-0.150) (9.441,-0.200) (9.437,-0.250) (9.433,-0.300) (9.428,-0.350) (9.422,-0.400) (9.415,-0.450) (9.407,-0.500) (9.399,-0.550) (9.390,-0.600) (9.379,-0.650) (9.368,-0.700) (9.357,-0.750) (9.344,-0.800) (9.330,-0.850) (9.316,-0.900) (9.300,-0.950) (9.284,-1.000) (9.267,-1.050) (9.248,-1.100) (9.229,-1.150) (9.209,-1.200) (9.188,-1.250) (9.165,-1.300) (9.142,-1.350) (9.118,-1.400) (9.092,-1.450) (9.065,-1.500) (9.037,-1.550) (9.008,-1.600) (8.978,-1.650) (8.950,-1.693) (8.912,-1.750) (8.878,-1.800) (8.842,-1.850) (8.804,-1.900) (8.765,-1.950) (8.725,-1.998) (8.680,-2.050) (8.636,-2.100) (8.600,-2.138) (8.564,-2.175) (8.525,-2.214) (8.488,-2.250) (8.450,-2.285) (8.406,-2.325) (8.350,-2.373) (8.300,-2.414) (8.255,-2.450) (8.200,-2.491) (8.153,-2.525) (8.100,-2.562) (8.050,-2.595) (8.003,-2.625) (7.950,-2.657) (7.900,-2.686) (7.850,-2.714) (7.800,-2.741) (7.750,-2.767) (7.700,-2.791) (7.650,-2.814) (7.600,-2.837) (7.550,-2.858) (7.500,-2.878) (7.450,-2.898) (7.400,-2.916) (7.350,-2.933) (7.300,-2.950) (7.250,-2.966) (7.200,-2.980) (7.150,-2.994) (7.100,-3.007) (7.050,-3.019) (7.000,-3.031) (6.950,-3.041) (6.900,-3.051) (6.850,-3.060) (6.800,-3.069) (6.750,-3.076) (6.700,-3.083) (6.650,-3.089) (6.600,-3.094) (6.550,-3.099) (6.500,-3.103) (6.450,-3.106) (6.400,-3.109) (6.350,-3.111) (6.300,-3.112) (6.250,-3.113) (6.200,-3.113) (6.150,-3.112) (6.100,-3.111) (6.050,-3.109) (6.000,-3.107) (5.950,-3.103) (5.900,-3.100) (5.850,-3.095) (5.800,-3.090) (5.750,-3.085) (5.700,-3.079) (5.650,-3.072) (5.600,-3.065) (5.550,-3.057) (5.500,-3.049) (5.450,-3.040) (5.400,-3.030) (5.350,-3.020) (5.300,-3.010) (5.250,-2.998) (5.200,-2.987) (5.151,-2.975) (5.100,-2.962) (5.050,-2.949) (5.000,-2.935) (4.950,-2.921) (4.900,-2.906) (4.850,-2.891) (4.800,-2.875) (4.750,-2.859) (4.700,-2.843) (4.650,-2.825) (4.600,-2.808) (4.550,-2.790) (4.500,-2.771) (4.450,-2.752) (4.400,-2.733) (4.350,-2.713) (4.300,-2.693) (4.250,-2.672) (4.200,-2.651) (4.150,-2.629) (4.100,-2.608) (4.050,-2.585) (4.000,-2.562) (3.950,-2.539) (3.900,-2.516) (3.850,-2.492) (3.800,-2.467) (3.750,-2.443) (3.700,-2.418) (3.650,-2.392) (3.600,-2.366) (3.550,-2.340) (3.500,-2.314) (3.450,-2.287) (3.400,-2.260) (3.350,-2.233) (3.300,-2.205) (3.250,-2.177) (3.203,-2.150) (3.159,-2.125) (3.116,-2.100) (3.050,-2.062) (3.000,-2.032) (2.950,-2.002) (2.904,-1.975) (2.850,-1.942) (2.800,-1.912) (2.750,-1.881) (2.700,-1.850) (2.650,-1.819) (2.600,-1.787) (2.550,-1.756) (2.502,-1.725) (2.450,-1.692) (2.400,-1.660) (2.350,-1.627) (2.300,-1.595) (2.250,-1.562) (2.200,-1.529) (2.156,-1.500) (2.100,-1.463) (2.050,-1.430) (2.006,-1.400) (1.950,-1.363) (1.900,-1.329) (1.858,-1.300) (1.800,-1.261) (1.750,-1.227) (1.700,-1.193) (1.650,-1.158) (1.602,-1.125) (1.550,-1.090) (1.500,-1.055) (1.457,-1.025) (1.400,-0.986) (1.350,-0.951) (1.300,-0.916) (1.250,-0.881) (1.206,-0.850) (1.150,-0.811) (1.100,-0.776) (1.050,-0.741) (1.000,-0.706) (0.956,-0.675) (0.900,-0.636) (0.850,-0.600) (0.800,-0.566) (0.750,-0.530) (0.708,-0.500) (0.650,-0.460) (0.601,-0.425) (0.550,-0.389) (0.500,-0.354) (0.450,-0.319) (0.400,-0.283) (0.354,-0.250) (0.300,-0.213) (0.250,-0.177) (0.206,-0.144)}
\EQ{(10.000,-7.071) (9.950,-7.036) (9.900,-7.000) (9.850,-6.965) (9.800,-6.930) (9.758,-6.900) (9.700,-6.859) (9.652,-6.825) (9.600,-6.788) (9.550,-6.753) (9.500,-6.718) (9.450,-6.682) (9.405,-6.650) (9.350,-6.611) (9.300,-6.576) (9.250,-6.541) (9.200,-6.505) (9.157,-6.475) (9.100,-6.435) (9.051,-6.400) (9.000,-6.364) (8.950,-6.329) (8.900,-6.293) (8.850,-6.258) (8.804,-6.225) (8.750,-6.187) (8.700,-6.152) (8.650,-6.117) (8.600,-6.081) (8.556,-6.050) (8.500,-6.010) (8.450,-5.975) (8.400,-5.940) (8.350,-5.904) (8.309,-5.875) (8.250,-5.834) (8.202,-5.800) (8.150,-5.763) (8.100,-5.728) (8.050,-5.692) (8.000,-5.657) (7.955,-5.625) (7.900,-5.586) (7.850,-5.551) (7.800,-5.515) (7.750,-5.480) (7.708,-5.450) (7.650,-5.409) (7.601,-5.375) (7.550,-5.339) (7.500,-5.303) (7.450,-5.268) (7.400,-5.233) (7.354,-5.200) (7.300,-5.162) (7.250,-5.127) (7.200,-5.091) (7.150,-5.056) (7.106,-5.025) (7.050,-4.985) (7.000,-4.950) (6.950,-4.914) (6.900,-4.879) (6.859,-4.850) (6.800,-4.808) (6.753,-4.775) (6.700,-4.738) (6.650,-4.702) (6.600,-4.667) (6.550,-4.632) (6.505,-4.600) (6.450,-4.561) (6.400,-4.525) (6.350,-4.490) (6.300,-4.455) (6.258,-4.425) (6.200,-4.384) (6.152,-4.350) (6.100,-4.313) (6.050,-4.278) (6.000,-4.243) (5.950,-4.207) (5.904,-4.175) (5.850,-4.137) (5.800,-4.101) (5.750,-4.066) (5.700,-4.031) (5.657,-4.000) (5.600,-3.960) (5.551,-3.925) (5.500,-3.889) (5.450,-3.854) (5.400,-3.818) (5.350,-3.783) (5.303,-3.750) (5.250,-3.712) (5.200,-3.677) (5.150,-3.642) (5.100,-3.606) (5.056,-3.575) (5.000,-3.536) (4.950,-3.500) (4.900,-3.465) (4.850,-3.429) (4.808,-3.400) (4.750,-3.359) (4.702,-3.325) (4.650,-3.288) (4.600,-3.253) (4.550,-3.217) (4.500,-3.182) (4.455,-3.150) (4.400,-3.111) (4.350,-3.076) (4.300,-3.041) (4.250,-3.005) (4.207,-2.975) (4.150,-2.935) (4.101,-2.900) (4.050,-2.864) (4.000,-2.828) (3.950,-2.793) (3.900,-2.758) (3.854,-2.725) (3.800,-2.687) (3.750,-2.652) (3.700,-2.616) (3.650,-2.581) (3.606,-2.550) (3.550,-2.510) (3.500,-2.475) (3.450,-2.440) (3.400,-2.404) (3.359,-2.375) (3.300,-2.334) (3.253,-2.300) (3.200,-2.263) (3.150,-2.227) (3.100,-2.192) (3.050,-2.157) (3.005,-2.125) (2.950,-2.086) (2.900,-2.051) (2.850,-2.015) (2.800,-1.980) (2.758,-1.950) (2.700,-1.909) (2.652,-1.875) (2.600,-1.839) (2.550,-1.803) (2.500,-1.768) (2.450,-1.732) (2.404,-1.700) (2.350,-1.662) (2.300,-1.626) (2.250,-1.591) (2.200,-1.556) (2.157,-1.525) (2.100,-1.485) (2.051,-1.450) (2.000,-1.414) (1.950,-1.379) (1.900,-1.344) (1.850,-1.308) (1.803,-1.275) (1.750,-1.238) (1.700,-1.202) (1.650,-1.167) (1.600,-1.131) (1.556,-1.100) (1.500,-1.061) (1.450,-1.025) (1.400,-0.990) (1.350,-0.955) (1.308,-0.925) (1.250,-0.884) (1.202,-0.850) (1.150,-0.813) (1.100,-0.778) (1.050,-0.743) (1.000,-0.707) (0.955,-0.675) (0.900,-0.637) (0.850,-0.601) (0.800,-0.566) (0.750,-0.531) (0.708,-0.500) (0.650,-0.460) (0.601,-0.425) (0.550,-0.389) (0.500,-0.354) (0.450,-0.319) (0.400,-0.283) (0.354,-0.250) (0.300,-0.213) (0.250,-0.177) (0.206,-0.144)}
\EQ{(-0.220,-0.155) (-0.275,-0.195) (-0.319,-0.225) (-0.375,-0.266) (-0.424,-0.300) (-0.475,-0.336) (-0.525,-0.371) (-0.566,-0.400) (-0.625,-0.442) (-0.672,-0.475) (-0.725,-0.513) (-0.775,-0.548) (-0.825,-0.584) (-0.875,-0.619) (-0.920,-0.650) (-0.975,-0.690) (-1.025,-0.725) (-1.075,-0.760) (-1.125,-0.796) (-1.167,-0.825) (-1.225,-0.866) (-1.273,-0.900) (-1.325,-0.937) (-1.375,-0.972) (-1.425,-1.008) (-1.475,-1.043) (-1.520,-1.075) (-1.575,-1.114) (-1.625,-1.149) (-1.675,-1.185) (-1.725,-1.220) (-1.768,-1.250) (-1.825,-1.291) (-1.874,-1.325) (-1.925,-1.361) (-1.975,-1.397) (-2.025,-1.432) (-2.075,-1.467) (-2.121,-1.500) (-2.175,-1.538) (-2.225,-1.573) (-2.275,-1.609) (-2.325,-1.644) (-2.369,-1.675) (-2.425,-1.715) (-2.475,-1.750) (-2.525,-1.786) (-2.575,-1.821) (-2.616,-1.850) (-2.675,-1.892) (-2.722,-1.925) (-2.775,-1.962) (-2.825,-1.998) (-2.875,-2.033) (-2.925,-2.068) (-2.970,-2.100) (-3.025,-2.139) (-3.075,-2.174) (-3.125,-2.210) (-3.175,-2.245) (-3.217,-2.275) (-3.275,-2.316) (-3.323,-2.350) (-3.375,-2.387) (-3.425,-2.422) (-3.475,-2.457) (-3.525,-2.493) (-3.571,-2.525) (-3.625,-2.563) (-3.675,-2.599) (-3.725,-2.634) (-3.775,-2.669) (-3.818,-2.700) (-3.875,-2.740) (-3.924,-2.775) (-3.975,-2.811) (-4.025,-2.846) (-4.066,-2.875) (-4.125,-2.917) (-4.172,-2.950) (-4.225,-2.988) (-4.275,-3.023) (-4.325,-3.058) (-4.375,-3.094) (-4.419,-3.125) (-4.475,-3.164) (-4.525,-3.200) (-4.575,-3.235) (-4.625,-3.270) (-4.667,-3.300) (-4.725,-3.341) (-4.773,-3.375) (-4.825,-3.412) (-4.875,-3.447) (-4.925,-3.483) (-4.975,-3.518) (-5.021,-3.550) (-5.075,-3.589) (-5.125,-3.624) (-5.175,-3.659) (-5.225,-3.695) (-5.268,-3.725) (-5.325,-3.765) (-5.374,-3.800) (-5.425,-3.836) (-5.475,-3.871) (-5.525,-3.907) (-5.575,-3.942) (-5.622,-3.975) (-5.675,-4.013) (-5.725,-4.048) (-5.775,-4.084) (-5.825,-4.119) (-5.869,-4.150) (-5.925,-4.190) (-5.975,-4.225) (-6.025,-4.260) (-6.075,-4.296) (-6.117,-4.325) (-6.175,-4.366) (-6.223,-4.400) (-6.275,-4.437) (-6.325,-4.472) (-6.375,-4.508) (-6.425,-4.543) (-6.470,-4.575) (-6.525,-4.614) (-6.575,-4.649) (-6.625,-4.685) (-6.675,-4.720) (-6.718,-4.750) (-6.775,-4.791) (-6.824,-4.825) (-6.875,-4.861) (-6.925,-4.897) (-6.975,-4.932) (-7.025,-4.967) (-7.071,-5.000) (-7.125,-5.038) (-7.175,-5.073) (-7.225,-5.109) (-7.275,-5.144) (-7.319,-5.175) (-7.375,-5.215) (-7.425,-5.250) (-7.475,-5.286) (-7.525,-5.321) (-7.566,-5.350) (-7.625,-5.392) (-7.672,-5.425) (-7.725,-5.462) (-7.775,-5.498) (-7.825,-5.533) (-7.875,-5.568) (-7.920,-5.600) (-7.975,-5.639) (-8.025,-5.675) (-8.075,-5.710) (-8.125,-5.745) (-8.167,-5.775) (-8.225,-5.816) (-8.273,-5.850) (-8.325,-5.887) (-8.375,-5.922) (-8.425,-5.957) (-8.475,-5.993) (-8.521,-6.025) (-8.575,-6.063) (-8.625,-6.099) (-8.675,-6.134) (-8.725,-6.170) (-8.768,-6.200) (-8.825,-6.240) (-8.874,-6.275) (-8.925,-6.311) (-8.975,-6.346) (-9.025,-6.382) (-9.075,-6.417) (-9.122,-6.450) (-9.175,-6.488) (-9.225,-6.523) (-9.275,-6.558) (-9.325,-6.594) (-9.369,-6.625) (-9.425,-6.665) (-9.475,-6.700) (-9.525,-6.735) (-9.575,-6.771) (-9.617,-6.800) (-9.675,-6.841) (-9.723,-6.875) (-9.775,-6.912) (-9.825,-6.947) (-9.875,-6.983) (-9.925,-7.018) (-9.970,-7.050) (-10.000,-7.071)}
\EQ{(0.206,0.144) (0.248,0.175) (0.300,0.213) (0.350,0.248) (0.400,0.283) (0.450,0.319) (0.496,0.350) (0.550,0.389) (0.600,0.424) (0.650,0.460) (0.700,0.495) (0.743,0.525) (0.800,0.566) (0.849,0.600) (0.900,0.637) (0.950,0.672) (1.000,0.707) (1.050,0.743) (1.096,0.775) (1.150,0.813) (1.200,0.849) (1.250,0.884) (1.300,0.919) (1.344,0.950) (1.400,0.990) (1.450,1.025) (1.500,1.061) (1.550,1.096) (1.591,1.125) (1.650,1.167) (1.697,1.200) (1.750,1.238) (1.800,1.273) (1.850,1.308) (1.900,1.344) (1.945,1.375) (2.000,1.414) (2.050,1.450) (2.100,1.485) (2.150,1.520) (2.192,1.550) (2.250,1.591) (2.298,1.625) (2.350,1.662) (2.400,1.697) (2.450,1.732) (2.500,1.768) (2.546,1.800) (2.600,1.839) (2.650,1.874) (2.700,1.909) (2.750,1.945) (2.793,1.975) (2.850,2.015) (2.899,2.050) (2.950,2.086) (3.000,2.121) (3.050,2.157) (3.100,2.192) (3.147,2.225) (3.200,2.263) (3.250,2.298) (3.300,2.334) (3.350,2.369) (3.394,2.400) (3.450,2.440) (3.500,2.475) (3.550,2.510) (3.600,2.546) (3.642,2.575) (3.700,2.616) (3.748,2.650) (3.800,2.687) (3.850,2.722) (3.900,2.758) (3.950,2.793) (3.995,2.825) (4.050,2.864) (4.100,2.899) (4.150,2.935) (4.200,2.970) (4.243,3.000) (4.300,3.041) (4.349,3.075) (4.400,3.111) (4.450,3.147) (4.500,3.182) (4.550,3.217) (4.596,3.250) (4.650,3.288) (4.700,3.323) (4.750,3.359) (4.800,3.394) (4.844,3.425) (4.900,3.465) (4.950,3.500) (5.000,3.536) (5.050,3.571) (5.091,3.600) (5.150,3.642) (5.197,3.675) (5.250,3.712) (5.300,3.748) (5.350,3.783) (5.400,3.818) (5.445,3.850) (5.500,3.889) (5.550,3.924) (5.600,3.960) (5.650,3.995) (5.692,4.025) (5.750,4.066) (5.798,4.100) (5.850,4.137) (5.900,4.172) (5.950,4.207) (6.000,4.243) (6.046,4.275) (6.100,4.313) (6.150,4.349) (6.200,4.384) (6.250,4.419) (6.293,4.450) (6.350,4.490) (6.399,4.525) (6.450,4.561) (6.500,4.596) (6.550,4.632) (6.600,4.667) (6.647,4.700) (6.700,4.738) (6.750,4.773) (6.800,4.808) (6.850,4.844) (6.894,4.875) (6.950,4.914) (7.000,4.950) (7.050,4.985) (7.100,5.020) (7.142,5.050) (7.200,5.091) (7.248,5.125) (7.300,5.162) (7.350,5.197) (7.400,5.233) (7.450,5.268) (7.495,5.300) (7.550,5.339) (7.600,5.374) (7.650,5.409) (7.700,5.445) (7.743,5.475) (7.800,5.515) (7.849,5.550) (7.900,5.586) (7.950,5.622) (8.000,5.657) (8.050,5.692) (8.096,5.725) (8.150,5.763) (8.200,5.798) (8.250,5.834) (8.300,5.869) (8.344,5.900) (8.400,5.940) (8.450,5.975) (8.500,6.010) (8.550,6.046) (8.591,6.075) (8.650,6.117) (8.697,6.150) (8.750,6.187) (8.800,6.223) (8.850,6.258) (8.900,6.293) (8.945,6.325) (9.000,6.364) (9.050,6.399) (9.100,6.435) (9.150,6.470) (9.192,6.500) (9.250,6.541) (9.298,6.575) (9.350,6.611) (9.400,6.647) (9.450,6.682) (9.500,6.718) (9.546,6.750) (9.600,6.788) (9.650,6.824) (9.700,6.859) (9.750,6.894) (9.793,6.925) (9.850,6.965) (9.899,7.000) (9.950,7.036) (10.000,7.071)}
\EQ{(-10.000,7.071) (-9.950,7.036) (-9.900,7.000) (-9.850,6.965) (-9.800,6.930) (-9.758,6.900) (-9.700,6.859) (-9.652,6.825) (-9.600,6.788) (-9.550,6.753) (-9.500,6.718) (-9.450,6.682) (-9.405,6.650) (-9.350,6.611) (-9.300,6.576) (-9.250,6.541) (-9.200,6.505) (-9.157,6.475) (-9.100,6.435) (-9.051,6.400) (-9.000,6.364) (-8.950,6.329) (-8.900,6.293) (-8.850,6.258) (-8.804,6.225) (-8.750,6.187) (-8.700,6.152) (-8.650,6.117) (-8.600,6.081) (-8.556,6.050) (-8.500,6.010) (-8.450,5.975) (-8.400,5.940) (-8.350,5.904) (-8.309,5.875) (-8.250,5.834) (-8.202,5.800) (-8.150,5.763) (-8.100,5.728) (-8.050,5.692) (-8.000,5.657) (-7.955,5.625) (-7.900,5.586) (-7.850,5.551) (-7.800,5.515) (-7.750,5.480) (-7.708,5.450) (-7.650,5.409) (-7.601,5.375) (-7.550,5.339) (-7.500,5.303) (-7.450,5.268) (-7.400,5.233) (-7.354,5.200) (-7.300,5.162) (-7.250,5.127) (-7.200,5.091) (-7.150,5.056) (-7.106,5.025) (-7.050,4.985) (-7.000,4.950) (-6.950,4.914) (-6.900,4.879) (-6.859,4.850) (-6.800,4.808) (-6.753,4.775) (-6.700,4.738) (-6.650,4.702) (-6.600,4.667) (-6.550,4.632) (-6.505,4.600) (-6.450,4.561) (-6.400,4.525) (-6.350,4.490) (-6.300,4.455) (-6.258,4.425) (-6.200,4.384) (-6.152,4.350) (-6.100,4.313) (-6.050,4.278) (-6.000,4.243) (-5.950,4.207) (-5.904,4.175) (-5.850,4.137) (-5.800,4.101) (-5.750,4.066) (-5.700,4.031) (-5.657,4.000) (-5.600,3.960) (-5.551,3.925) (-5.500,3.889) (-5.450,3.854) (-5.400,3.818) (-5.350,3.783) (-5.303,3.750) (-5.250,3.712) (-5.200,3.677) (-5.150,3.642) (-5.100,3.606) (-5.056,3.575) (-5.000,3.536) (-4.950,3.500) (-4.900,3.465) (-4.850,3.429) (-4.808,3.400) (-4.750,3.359) (-4.702,3.325) (-4.650,3.288) (-4.600,3.253) (-4.550,3.217) (-4.500,3.182) (-4.455,3.150) (-4.400,3.111) (-4.350,3.076) (-4.300,3.041) (-4.250,3.005) (-4.207,2.975) (-4.150,2.935) (-4.101,2.900) (-4.050,2.864) (-4.000,2.828) (-3.950,2.793) (-3.900,2.758) (-3.854,2.725) (-3.800,2.687) (-3.750,2.652) (-3.700,2.616) (-3.650,2.581) (-3.606,2.550) (-3.550,2.510) (-3.500,2.475) (-3.450,2.440) (-3.400,2.404) (-3.359,2.375) (-3.300,2.334) (-3.253,2.300) (-3.200,2.263) (-3.150,2.227) (-3.100,2.192) (-3.050,2.157) (-3.005,2.125) (-2.950,2.086) (-2.900,2.051) (-2.850,2.015) (-2.800,1.980) (-2.758,1.950) (-2.700,1.909) (-2.652,1.875) (-2.600,1.839) (-2.550,1.803) (-2.500,1.768) (-2.450,1.732) (-2.404,1.700) (-2.350,1.662) (-2.300,1.626) (-2.250,1.591) (-2.200,1.556) (-2.157,1.525) (-2.100,1.485) (-2.051,1.450) (-2.000,1.414) (-1.950,1.379) (-1.900,1.344) (-1.850,1.308) (-1.803,1.275) (-1.750,1.238) (-1.700,1.202) (-1.650,1.167) (-1.600,1.131) (-1.556,1.100) (-1.500,1.061) (-1.450,1.025) (-1.400,0.990) (-1.350,0.955) (-1.308,0.925) (-1.250,0.884) (-1.202,0.850) (-1.150,0.813) (-1.100,0.778) (-1.050,0.743) (-1.000,0.707) (-0.955,0.675) (-0.900,0.637) (-0.850,0.601) (-0.800,0.566) (-0.750,0.531) (-0.708,0.500) (-0.650,0.460) (-0.601,0.425) (-0.550,0.389) (-0.500,0.354) (-0.450,0.319) (-0.400,0.283) (-0.354,0.250) (-0.300,0.213) (-0.250,0.177) (-0.220,0.155)}
\EQ{(4.306,-10.000) (4.350,-9.950) (4.392,-9.900) (4.425,-9.860) (4.473,-9.800) (4.513,-9.750) (4.550,-9.701) (4.588,-9.650) (4.624,-9.600) (4.660,-9.550) (4.695,-9.500) (4.725,-9.455) (4.761,-9.400) (4.793,-9.350) (4.824,-9.300) (4.855,-9.250) (4.884,-9.200) (4.913,-9.150) (4.941,-9.100) (4.969,-9.050) (4.995,-9.000) (5.021,-8.950) (5.046,-8.900) (5.071,-8.850) (5.095,-8.800) (5.118,-8.750) (5.140,-8.700) (5.162,-8.650) (5.183,-8.600) (5.203,-8.550) (5.223,-8.500) (5.242,-8.450) (5.261,-8.400) (5.279,-8.350) (5.296,-8.300) (5.313,-8.250) (5.329,-8.200) (5.344,-8.150) (5.359,-8.100) (5.373,-8.050) (5.387,-8.000) (5.400,-7.951) (5.413,-7.900) (5.425,-7.850) (5.436,-7.800) (5.447,-7.750) (5.457,-7.700) (5.467,-7.650) (5.476,-7.600) (5.484,-7.550) (5.492,-7.500) (5.500,-7.450) (5.507,-7.400) (5.513,-7.350) (5.519,-7.300) (5.524,-7.250) (5.529,-7.200) (5.533,-7.150) (5.537,-7.100) (5.540,-7.050) (5.542,-7.000) (5.544,-6.950) (5.545,-6.900) (5.546,-6.850) (5.547,-6.800) (5.546,-6.750) (5.546,-6.700) (5.544,-6.650) (5.543,-6.600) (5.540,-6.550) (5.537,-6.500) (5.534,-6.450) (5.530,-6.400) (5.525,-6.350) (5.520,-6.300) (5.515,-6.250) (5.509,-6.200) (5.502,-6.150) (5.495,-6.100) (5.487,-6.050) (5.478,-6.000) (5.469,-5.950) (5.460,-5.900) (5.450,-5.850) (5.439,-5.800) (5.428,-5.750) (5.416,-5.700) (5.404,-5.650) (5.391,-5.600) (5.378,-5.550) (5.363,-5.500) (5.349,-5.450) (5.334,-5.400) (5.318,-5.350) (5.301,-5.300) (5.284,-5.250) (5.267,-5.200) (5.249,-5.150) (5.230,-5.100) (5.210,-5.050) (5.190,-5.000) (5.170,-4.950) (5.150,-4.903) (5.127,-4.850) (5.104,-4.800) (5.081,-4.750) (5.057,-4.700) (5.033,-4.650) (5.008,-4.600) (4.982,-4.550) (4.956,-4.500) (4.929,-4.450) (4.901,-4.400) (4.875,-4.354) (4.844,-4.300) (4.814,-4.250) (4.784,-4.200) (4.753,-4.150) (4.725,-4.107) (4.688,-4.050) (4.655,-4.000) (4.625,-3.955) (4.587,-3.900) (4.551,-3.850) (4.516,-3.800) (4.479,-3.750) (4.441,-3.700) (4.403,-3.650) (4.364,-3.600) (4.325,-3.551) (4.284,-3.500) (4.250,-3.458) (4.201,-3.400) (4.159,-3.350) (4.125,-3.311) (4.075,-3.254) (4.026,-3.200) (3.981,-3.150) (3.935,-3.100) (3.900,-3.063) (3.864,-3.025) (3.825,-2.985) (3.775,-2.934) (3.725,-2.883) (3.675,-2.834) (3.625,-2.785) (3.575,-2.736) (3.537,-2.700) (3.500,-2.665) (3.457,-2.625) (3.403,-2.575) (3.350,-2.526) (3.300,-2.481) (3.250,-2.437) (3.209,-2.400) (3.151,-2.350) (3.100,-2.306) (3.050,-2.263) (3.006,-2.225) (2.950,-2.178) (2.900,-2.137) (2.856,-2.100) (2.800,-2.054) (2.750,-2.014) (2.702,-1.975) (2.650,-1.933) (2.609,-1.900) (2.550,-1.854) (2.500,-1.814) (2.450,-1.775) (2.400,-1.737) (2.353,-1.700) (2.300,-1.660) (2.255,-1.625) (2.200,-1.583) (2.156,-1.550) (2.100,-1.508) (2.056,-1.475) (2.000,-1.433) (1.956,-1.400) (1.900,-1.359) (1.854,-1.325) (1.800,-1.285) (1.752,-1.250) (1.700,-1.212) (1.650,-1.175) (1.600,-1.139) (1.550,-1.103) (1.500,-1.067) (1.450,-1.031) (1.408,-1.000) (1.350,-0.959) (1.304,-0.925) (1.250,-0.887) (1.200,-0.851) (1.150,-0.815) (1.100,-0.780) (1.059,-0.750) (1.000,-0.708) (0.953,-0.675) (0.900,-0.637) (0.850,-0.602) (0.800,-0.566) (0.750,-0.531) (0.707,-0.500) (0.650,-0.460) (0.601,-0.425) (0.550,-0.389) (0.500,-0.354) (0.450,-0.319) (0.400,-0.283) (0.354,-0.250) (0.300,-0.213) (0.250,-0.177) (0.206,-0.144)}
\EQ{(-0.220,-0.155) (-0.275,-0.195) (-0.319,-0.225) (-0.375,-0.266) (-0.424,-0.300) (-0.475,-0.336) (-0.525,-0.372) (-0.566,-0.400) (-0.625,-0.442) (-0.672,-0.475) (-0.725,-0.513) (-0.775,-0.549) (-0.825,-0.584) (-0.875,-0.620) (-0.918,-0.650) (-0.975,-0.691) (-1.024,-0.725) (-1.075,-0.762) (-1.125,-0.797) (-1.175,-0.833) (-1.225,-0.869) (-1.269,-0.900) (-1.325,-0.941) (-1.373,-0.975) (-1.425,-1.013) (-1.475,-1.049) (-1.525,-1.085) (-1.575,-1.121) (-1.625,-1.157) (-1.675,-1.194) (-1.718,-1.225) (-1.775,-1.267) (-1.820,-1.300) (-1.875,-1.340) (-1.922,-1.375) (-1.975,-1.415) (-2.023,-1.450) (-2.075,-1.489) (-2.123,-1.525) (-2.175,-1.564) (-2.222,-1.600) (-2.275,-1.640) (-2.320,-1.675) (-2.375,-1.717) (-2.418,-1.750) (-2.475,-1.795) (-2.525,-1.834) (-2.575,-1.873) (-2.625,-1.913) (-2.671,-1.950) (-2.725,-1.993) (-2.764,-2.025) (-2.825,-2.075) (-2.875,-2.116) (-2.916,-2.150) (-2.975,-2.199) (-3.025,-2.241) (-3.065,-2.275) (-3.123,-2.325) (-3.175,-2.371) (-3.225,-2.415) (-3.265,-2.450) (-3.321,-2.500) (-3.375,-2.549) (-3.425,-2.595) (-3.475,-2.642) (-3.525,-2.689) (-3.563,-2.725) (-3.600,-2.760) (-3.641,-2.800) (-3.692,-2.850) (-3.742,-2.900) (-3.791,-2.950) (-3.840,-3.000) (-3.875,-3.037) (-3.911,-3.075) (-3.950,-3.117) (-4.000,-3.171) (-4.049,-3.225) (-4.093,-3.275) (-4.137,-3.325) (-4.175,-3.369) (-4.222,-3.425) (-4.264,-3.475) (-4.300,-3.519) (-4.345,-3.575) (-4.384,-3.625) (-4.422,-3.675) (-4.460,-3.725) (-4.497,-3.775) (-4.534,-3.825) (-4.569,-3.875) (-4.600,-3.919) (-4.638,-3.975) (-4.672,-4.025) (-4.700,-4.068) (-4.737,-4.125) (-4.768,-4.175) (-4.799,-4.225) (-4.825,-4.269) (-4.858,-4.325) (-4.887,-4.375) (-4.915,-4.425) (-4.942,-4.475) (-4.969,-4.525) (-4.995,-4.575) (-5.020,-4.625) (-5.045,-4.675) (-5.069,-4.725) (-5.093,-4.775) (-5.116,-4.825) (-5.138,-4.875) (-5.159,-4.925) (-5.180,-4.975) (-5.200,-5.024) (-5.220,-5.075) (-5.239,-5.125) (-5.258,-5.175) (-5.275,-5.223) (-5.293,-5.275) (-5.310,-5.325) (-5.325,-5.373) (-5.341,-5.425) (-5.356,-5.475) (-5.371,-5.525) (-5.384,-5.575) (-5.398,-5.625) (-5.410,-5.675) (-5.422,-5.725) (-5.434,-5.775) (-5.445,-5.825) (-5.455,-5.875) (-5.465,-5.925) (-5.474,-5.975) (-5.483,-6.025) (-5.491,-6.075) (-5.498,-6.125) (-5.505,-6.175) (-5.512,-6.225) (-5.518,-6.275) (-5.523,-6.325) (-5.528,-6.375) (-5.532,-6.425) (-5.536,-6.475) (-5.539,-6.525) (-5.542,-6.575) (-5.544,-6.625) (-5.545,-6.675) (-5.546,-6.725) (-5.547,-6.775) (-5.547,-6.825) (-5.546,-6.875) (-5.545,-6.925) (-5.543,-6.975) (-5.541,-7.025) (-5.538,-7.075) (-5.535,-7.125) (-5.531,-7.175) (-5.526,-7.225) (-5.522,-7.275) (-5.516,-7.325) (-5.510,-7.375) (-5.503,-7.425) (-5.496,-7.475) (-5.489,-7.525) (-5.480,-7.575) (-5.471,-7.625) (-5.462,-7.675) (-5.452,-7.725) (-5.442,-7.775) (-5.430,-7.825) (-5.419,-7.875) (-5.407,-7.925) (-5.394,-7.975) (-5.380,-8.025) (-5.366,-8.075) (-5.352,-8.125) (-5.337,-8.175) (-5.321,-8.225) (-5.305,-8.275) (-5.288,-8.325) (-5.270,-8.375) (-5.252,-8.425) (-5.233,-8.475) (-5.213,-8.525) (-5.193,-8.575) (-5.173,-8.625) (-5.151,-8.675) (-5.129,-8.725) (-5.106,-8.775) (-5.083,-8.825) (-5.059,-8.875) (-5.034,-8.925) (-5.008,-8.975) (-4.982,-9.025) (-4.955,-9.075) (-4.927,-9.125) (-4.900,-9.173) (-4.870,-9.225) (-4.840,-9.275) (-4.809,-9.325) (-4.777,-9.375) (-4.750,-9.417) (-4.711,-9.475) (-4.677,-9.525) (-4.642,-9.575) (-4.606,-9.625) (-4.575,-9.668) (-4.532,-9.725) (-4.500,-9.766) (-4.453,-9.825) (-4.413,-9.875) (-4.375,-9.920) (-4.328,-9.975) (-4.306,-10.000)}
\EQ{(0.206,0.144) (0.248,0.175) (0.300,0.213) (0.350,0.248) (0.400,0.283) (0.450,0.319) (0.495,0.350) (0.550,0.389) (0.600,0.424) (0.650,0.460) (0.700,0.495) (0.742,0.525) (0.800,0.566) (0.848,0.600) (0.900,0.637) (0.950,0.673) (1.000,0.708) (1.050,0.744) (1.094,0.775) (1.150,0.815) (1.199,0.850) (1.250,0.887) (1.300,0.923) (1.350,0.959) (1.400,0.995) (1.443,1.025) (1.500,1.067) (1.546,1.100) (1.600,1.139) (1.650,1.175) (1.700,1.212) (1.750,1.248) (1.800,1.285) (1.850,1.322) (1.900,1.359) (1.950,1.396) (2.000,1.433) (2.050,1.470) (2.100,1.508) (2.150,1.546) (2.200,1.583) (2.250,1.621) (2.300,1.660) (2.350,1.698) (2.400,1.737) (2.450,1.775) (2.500,1.814) (2.546,1.850) (2.600,1.893) (2.640,1.925) (2.700,1.973) (2.750,2.014) (2.795,2.050) (2.850,2.095) (2.900,2.137) (2.946,2.175) (3.000,2.220) (3.050,2.263) (3.094,2.300) (3.150,2.349) (3.200,2.393) (3.250,2.437) (3.293,2.475) (3.349,2.525) (3.400,2.572) (3.450,2.618) (3.500,2.665) (3.537,2.700) (3.575,2.736) (3.615,2.775) (3.666,2.825) (3.717,2.875) (3.767,2.925) (3.815,2.975) (3.864,3.025) (3.900,3.063) (3.935,3.100) (3.975,3.144) (4.025,3.198) (4.071,3.250) (4.115,3.300) (4.150,3.340) (4.200,3.399) (4.243,3.450) (4.275,3.489) (4.325,3.550) (4.364,3.600) (4.400,3.646) (4.441,3.700) (4.475,3.745) (4.516,3.800) (4.550,3.848) (4.587,3.900) (4.621,3.950) (4.650,3.992) (4.688,4.050) (4.721,4.100) (4.750,4.146) (4.784,4.200) (4.814,4.250) (4.844,4.300) (4.873,4.350) (4.900,4.398) (4.925,4.443) (4.950,4.489) (4.975,4.536) (5.000,4.585) (5.025,4.634) (5.050,4.685) (5.075,4.737) (5.100,4.791) (5.125,4.846) (5.149,4.900) (5.170,4.950) (5.190,5.000) (5.210,5.050) (5.230,5.100) (5.249,5.150) (5.267,5.200) (5.284,5.250) (5.301,5.300) (5.318,5.350) (5.334,5.400) (5.349,5.450) (5.363,5.500) (5.378,5.550) (5.391,5.600) (5.404,5.650) (5.416,5.700) (5.428,5.750) (5.439,5.800) (5.450,5.850) (5.460,5.900) (5.469,5.950) (5.478,6.000) (5.487,6.050) (5.495,6.100) (5.502,6.150) (5.509,6.200) (5.515,6.250) (5.520,6.300) (5.525,6.350) (5.530,6.400) (5.534,6.450) (5.537,6.500) (5.540,6.550) (5.543,6.600) (5.544,6.650) (5.546,6.700) (5.546,6.750) (5.547,6.800) (5.546,6.850) (5.545,6.900) (5.544,6.950) (5.542,7.000) (5.540,7.050) (5.537,7.100) (5.533,7.150) (5.529,7.200) (5.524,7.250) (5.519,7.300) (5.513,7.350) (5.507,7.400) (5.500,7.450) (5.492,7.500) (5.484,7.550) (5.476,7.600) (5.467,7.650) (5.457,7.700) (5.447,7.750) (5.436,7.800) (5.425,7.850) (5.413,7.900) (5.400,7.950) (5.387,8.000) (5.373,8.050) (5.359,8.100) (5.344,8.150) (5.329,8.200) (5.313,8.250) (5.296,8.300) (5.279,8.350) (5.261,8.400) (5.242,8.450) (5.223,8.500) (5.203,8.550) (5.183,8.600) (5.162,8.650) (5.140,8.700) (5.118,8.750) (5.095,8.800) (5.071,8.850) (5.046,8.900) (5.021,8.950) (4.995,9.000) (4.969,9.050) (4.941,9.100) (4.913,9.150) (4.884,9.200) (4.855,9.250) (4.825,9.299) (4.793,9.350) (4.761,9.400) (4.728,9.450) (4.700,9.492) (4.660,9.550) (4.625,9.599) (4.588,9.650) (4.551,9.700) (4.513,9.750) (4.475,9.798) (4.433,9.850) (4.400,9.890) (4.350,9.949) (4.306,10.000)}
\EQ{(-4.306,10.000) (-4.350,9.950) (-4.392,9.900) (-4.425,9.860) (-4.473,9.800) (-4.513,9.750) (-4.550,9.701) (-4.588,9.650) (-4.624,9.600) (-4.660,9.550) (-4.695,9.500) (-4.725,9.455) (-4.761,9.400) (-4.793,9.350) (-4.824,9.300) (-4.855,9.250) (-4.884,9.200) (-4.913,9.150) (-4.941,9.100) (-4.969,9.050) (-4.995,9.000) (-5.021,8.950) (-5.046,8.900) (-5.071,8.850) (-5.095,8.800) (-5.118,8.750) (-5.140,8.700) (-5.162,8.650) (-5.183,8.600) (-5.203,8.550) (-5.223,8.500) (-5.242,8.450) (-5.261,8.400) (-5.279,8.350) (-5.296,8.300) (-5.313,8.250) (-5.329,8.200) (-5.344,8.150) (-5.359,8.100) (-5.373,8.050) (-5.387,8.000) (-5.400,7.951) (-5.413,7.900) (-5.425,7.850) (-5.436,7.800) (-5.447,7.750) (-5.457,7.700) (-5.467,7.650) (-5.476,7.600) (-5.484,7.550) (-5.492,7.500) (-5.500,7.450) (-5.507,7.400) (-5.513,7.350) (-5.519,7.300) (-5.524,7.250) (-5.529,7.200) (-5.533,7.150) (-5.537,7.100) (-5.540,7.050) (-5.542,7.000) (-5.544,6.950) (-5.545,6.900) (-5.546,6.850) (-5.547,6.800) (-5.546,6.750) (-5.546,6.700) (-5.544,6.650) (-5.543,6.600) (-5.540,6.550) (-5.537,6.500) (-5.534,6.450) (-5.530,6.400) (-5.525,6.350) (-5.520,6.300) (-5.515,6.250) (-5.509,6.200) (-5.502,6.150) (-5.495,6.100) (-5.487,6.050) (-5.478,6.000) (-5.469,5.950) (-5.460,5.900) (-5.450,5.850) (-5.439,5.800) (-5.428,5.750) (-5.416,5.700) (-5.404,5.650) (-5.391,5.600) (-5.378,5.550) (-5.363,5.500) (-5.349,5.450) (-5.334,5.400) (-5.318,5.350) (-5.301,5.300) (-5.284,5.250) (-5.267,5.200) (-5.249,5.150) (-5.230,5.100) (-5.210,5.050) (-5.190,5.000) (-5.170,4.950) (-5.150,4.903) (-5.127,4.850) (-5.104,4.800) (-5.081,4.750) (-5.057,4.700) (-5.033,4.650) (-5.008,4.600) (-4.982,4.550) (-4.956,4.500) (-4.929,4.450) (-4.901,4.400) (-4.875,4.354) (-4.844,4.300) (-4.814,4.250) (-4.784,4.200) (-4.753,4.150) (-4.725,4.107) (-4.688,4.050) (-4.655,4.000) (-4.625,3.955) (-4.587,3.900) (-4.551,3.850) (-4.516,3.800) (-4.479,3.750) (-4.441,3.700) (-4.403,3.650) (-4.364,3.600) (-4.325,3.551) (-4.284,3.500) (-4.250,3.458) (-4.201,3.400) (-4.159,3.350) (-4.125,3.311) (-4.075,3.254) (-4.026,3.200) (-3.981,3.150) (-3.935,3.100) (-3.900,3.063) (-3.864,3.025) (-3.825,2.985) (-3.775,2.934) (-3.725,2.883) (-3.675,2.834) (-3.625,2.785) (-3.575,2.736) (-3.537,2.700) (-3.500,2.665) (-3.457,2.625) (-3.403,2.575) (-3.350,2.526) (-3.300,2.481) (-3.250,2.437) (-3.209,2.400) (-3.151,2.350) (-3.100,2.306) (-3.050,2.263) (-3.006,2.225) (-2.950,2.178) (-2.900,2.137) (-2.856,2.100) (-2.800,2.054) (-2.750,2.014) (-2.702,1.975) (-2.650,1.933) (-2.609,1.900) (-2.550,1.854) (-2.500,1.814) (-2.450,1.775) (-2.400,1.737) (-2.353,1.700) (-2.300,1.660) (-2.255,1.625) (-2.200,1.583) (-2.156,1.550) (-2.100,1.508) (-2.056,1.475) (-2.000,1.433) (-1.956,1.400) (-1.900,1.359) (-1.854,1.325) (-1.800,1.285) (-1.752,1.250) (-1.700,1.212) (-1.650,1.175) (-1.600,1.139) (-1.550,1.103) (-1.500,1.067) (-1.450,1.031) (-1.408,1.000) (-1.350,0.959) (-1.304,0.925) (-1.250,0.887) (-1.200,0.851) (-1.150,0.815) (-1.100,0.780) (-1.059,0.750) (-1.000,0.708) (-0.953,0.675) (-0.900,0.637) (-0.850,0.602) (-0.800,0.566) (-0.750,0.531) (-0.707,0.500) (-0.650,0.460) (-0.601,0.425) (-0.550,0.389) (-0.500,0.354) (-0.450,0.319) (-0.400,0.283) (-0.354,0.250) (-0.300,0.213) (-0.250,0.177) (-0.220,0.155)}
\EQ{(-0.220,-0.155) (-0.275,-0.195) (-0.319,-0.225) (-0.375,-0.266) (-0.424,-0.300) (-0.475,-0.337) (-0.525,-0.372) (-0.575,-0.407) (-0.625,-0.443) (-0.671,-0.475) (-0.725,-0.514) (-0.775,-0.549) (-0.825,-0.585) (-0.875,-0.620) (-0.917,-0.650) (-0.975,-0.692) (-1.022,-0.725) (-1.075,-0.763) (-1.125,-0.799) (-1.175,-0.835) (-1.225,-0.872) (-1.275,-0.908) (-1.325,-0.944) (-1.368,-0.975) (-1.425,-1.017) (-1.470,-1.050) (-1.525,-1.091) (-1.571,-1.125) (-1.625,-1.166) (-1.671,-1.200) (-1.725,-1.241) (-1.770,-1.275) (-1.825,-1.317) (-1.869,-1.350) (-1.925,-1.394) (-1.965,-1.425) (-2.025,-1.472) (-2.075,-1.512) (-2.123,-1.550) (-2.175,-1.592) (-2.216,-1.625) (-2.275,-1.674) (-2.325,-1.715) (-2.367,-1.750) (-2.425,-1.799) (-2.475,-1.843) (-2.525,-1.886) (-2.570,-1.925) (-2.625,-1.975) (-2.675,-2.020) (-2.725,-2.066) (-2.762,-2.100) (-2.800,-2.136) (-2.841,-2.175) (-2.893,-2.225) (-2.943,-2.275) (-2.993,-2.325) (-3.042,-2.375) (-3.090,-2.425) (-3.125,-2.463) (-3.160,-2.500) (-3.200,-2.545) (-3.249,-2.600) (-3.293,-2.650) (-3.336,-2.700) (-3.375,-2.747) (-3.418,-2.800) (-3.450,-2.840) (-3.497,-2.900) (-3.535,-2.950) (-3.572,-3.000) (-3.608,-3.050) (-3.644,-3.100) (-3.675,-3.146) (-3.712,-3.200) (-3.744,-3.250) (-3.775,-3.299) (-3.807,-3.350) (-3.837,-3.400) (-3.866,-3.450) (-3.894,-3.500) (-3.922,-3.550) (-3.948,-3.600) (-3.974,-3.650) (-3.999,-3.700) (-4.023,-3.750) (-4.047,-3.800) (-4.069,-3.850) (-4.091,-3.900) (-4.112,-3.950) (-4.132,-4.000) (-4.152,-4.050) (-4.170,-4.100) (-4.188,-4.150) (-4.206,-4.200) (-4.222,-4.250) (-4.238,-4.300) (-4.253,-4.350) (-4.267,-4.400) (-4.280,-4.450) (-4.293,-4.500) (-4.305,-4.550) (-4.316,-4.600) (-4.327,-4.650) (-4.337,-4.700) (-4.346,-4.750) (-4.354,-4.800) (-4.362,-4.850) (-4.369,-4.900) (-4.376,-4.950) (-4.381,-5.000) (-4.386,-5.050) (-4.391,-5.100) (-4.394,-5.150) (-4.397,-5.200) (-4.400,-5.250) (-4.401,-5.300) (-4.402,-5.350) (-4.402,-5.400) (-4.402,-5.450) (-4.401,-5.500) (-4.399,-5.550) (-4.396,-5.600) (-4.393,-5.650) (-4.389,-5.700) (-4.385,-5.750) (-4.380,-5.800) (-4.374,-5.850) (-4.367,-5.900) (-4.360,-5.950) (-4.352,-6.000) (-4.343,-6.050) (-4.334,-6.100) (-4.323,-6.150) (-4.313,-6.200) (-4.301,-6.250) (-4.289,-6.300) (-4.276,-6.350) (-4.262,-6.400) (-4.247,-6.450) (-4.232,-6.500) (-4.216,-6.550) (-4.200,-6.598) (-4.182,-6.650) (-4.163,-6.700) (-4.144,-6.750) (-4.125,-6.798) (-4.104,-6.850) (-4.082,-6.900) (-4.060,-6.950) (-4.036,-7.000) (-4.012,-7.050) (-3.987,-7.100) (-3.961,-7.150) (-3.934,-7.200) (-3.907,-7.250) (-3.878,-7.300) (-3.850,-7.347) (-3.818,-7.400) (-3.786,-7.450) (-3.753,-7.500) (-3.725,-7.542) (-3.684,-7.600) (-3.650,-7.648) (-3.611,-7.700) (-3.575,-7.747) (-3.533,-7.800) (-3.500,-7.841) (-3.450,-7.900) (-3.406,-7.950) (-3.361,-8.000) (-3.325,-8.039) (-3.275,-8.092) (-3.225,-8.142) (-3.175,-8.191) (-3.125,-8.238) (-3.085,-8.275) (-3.029,-8.325) (-2.975,-8.372) (-2.925,-8.414) (-2.880,-8.450) (-2.825,-8.494) (-2.784,-8.525) (-2.725,-8.569) (-2.682,-8.600) (-2.625,-8.640) (-2.575,-8.674) (-2.525,-8.707) (-2.475,-8.739) (-2.425,-8.770) (-2.375,-8.800) (-2.325,-8.829) (-2.275,-8.858) (-2.225,-8.886) (-2.175,-8.912) (-2.125,-8.938) (-2.075,-8.964) (-2.025,-8.988) (-1.975,-9.012) (-1.925,-9.035) (-1.875,-9.057) (-1.825,-9.078) (-1.775,-9.099) (-1.725,-9.119) (-1.675,-9.139) (-1.625,-9.158) (-1.577,-9.175) (-1.525,-9.193) (-1.475,-9.210) (-1.425,-9.226) (-1.375,-9.242) (-1.325,-9.257) (-1.275,-9.272) (-1.225,-9.285) (-1.175,-9.298) (-1.125,-9.311) (-1.075,-9.323) (-1.025,-9.335) (-0.975,-9.345) (-0.925,-9.356) (-0.875,-9.365) (-0.825,-9.375) (-0.775,-9.383) (-0.725,-9.391) (-0.675,-9.399) (-0.625,-9.406) (-0.575,-9.412) (-0.525,-9.418) (-0.475,-9.423) (-0.425,-9.428) (-0.375,-9.432) (-0.325,-9.436) (-0.275,-9.439) (-0.225,-9.442) (-0.175,-9.444) (-0.125,-9.446) (-0.075,-9.447) (-0.025,-9.447) (0.050,-9.447) (0.100,-9.446) (0.150,-9.445) (0.200,-9.443) (0.250,-9.441) (0.300,-9.438) (0.350,-9.434) (0.400,-9.430) (0.450,-9.426) (0.500,-9.421) (0.550,-9.415) (0.600,-9.409) (0.650,-9.402) (0.700,-9.395) (0.750,-9.387) (0.800,-9.379) (0.850,-9.370) (0.900,-9.361) (0.950,-9.351) (1.000,-9.340) (1.050,-9.329) (1.100,-9.317) (1.150,-9.305) (1.200,-9.292) (1.250,-9.278) (1.300,-9.264) (1.349,-9.250) (1.400,-9.234) (1.450,-9.218) (1.500,-9.202) (1.550,-9.185) (1.600,-9.167) (1.650,-9.148) (1.700,-9.129) (1.750,-9.109) (1.800,-9.089) (1.850,-9.068) (1.900,-9.046) (1.946,-9.025) (2.000,-9.000) (2.050,-8.976) (2.100,-8.951) (2.150,-8.925) (2.198,-8.900) (2.244,-8.875) (2.289,-8.850) (2.333,-8.825) (2.400,-8.785) (2.450,-8.755) (2.497,-8.725) (2.550,-8.691) (2.600,-8.657) (2.647,-8.625) (2.700,-8.587) (2.750,-8.551) (2.800,-8.513) (2.849,-8.475) (2.900,-8.434) (2.942,-8.400) (3.000,-8.351) (3.050,-8.307) (3.100,-8.262) (3.139,-8.225) (3.192,-8.175) (3.242,-8.125) (3.291,-8.075) (3.338,-8.025) (3.375,-7.985) (3.425,-7.929) (3.471,-7.875) (3.513,-7.825) (3.550,-7.779) (3.592,-7.725) (3.625,-7.682) (3.667,-7.625) (3.700,-7.578) (3.736,-7.525) (3.770,-7.475) (3.800,-7.428) (3.833,-7.375) (3.863,-7.325) (3.892,-7.275) (3.921,-7.225) (3.948,-7.175) (3.974,-7.125) (4.000,-7.075) (4.024,-7.025) (4.048,-6.975) (4.071,-6.925) (4.093,-6.875) (4.114,-6.825) (4.134,-6.775) (4.154,-6.725) (4.173,-6.675) (4.191,-6.625) (4.208,-6.575) (4.224,-6.525) (4.240,-6.475) (4.255,-6.425) (4.269,-6.375) (4.282,-6.325) (4.295,-6.275) (4.307,-6.225) (4.318,-6.175) (4.329,-6.125) (4.338,-6.075) (4.347,-6.025) (4.356,-5.975) (4.364,-5.925) (4.370,-5.875) (4.377,-5.825) (4.382,-5.775) (4.387,-5.725) (4.391,-5.675) (4.395,-5.625) (4.398,-5.575) (4.400,-5.525) (4.401,-5.475) (4.402,-5.425) (4.402,-5.375) (4.402,-5.325) (4.401,-5.275) (4.399,-5.225) (4.396,-5.175) (4.393,-5.125) (4.389,-5.075) (4.384,-5.025) (4.379,-4.975) (4.373,-4.925) (4.366,-4.875) (4.358,-4.825) (4.350,-4.775) (4.342,-4.725) (4.332,-4.675) (4.322,-4.625) (4.311,-4.575) (4.299,-4.525) (4.287,-4.475) (4.274,-4.425) (4.260,-4.375) (4.245,-4.325) (4.230,-4.275) (4.214,-4.225) (4.197,-4.175) (4.180,-4.125) (4.161,-4.075) (4.142,-4.025) (4.122,-3.975) (4.102,-3.925) (4.080,-3.875) (4.058,-3.825) (4.035,-3.775) (4.011,-3.725) (3.987,-3.675) (3.961,-3.625) (3.935,-3.575) (3.908,-3.525) (3.880,-3.475) (3.852,-3.425) (3.825,-3.380) (3.792,-3.325) (3.760,-3.275) (3.728,-3.225) (3.700,-3.183) (3.661,-3.125) (3.626,-3.075) (3.590,-3.025) (3.554,-2.975) (3.516,-2.925) (3.478,-2.875) (3.438,-2.825) (3.400,-2.778) (3.357,-2.725) (3.315,-2.675) (3.275,-2.629) (3.227,-2.575) (3.182,-2.525) (3.136,-2.475) (3.100,-2.436) (3.050,-2.384) (3.000,-2.332) (2.950,-2.282) (2.900,-2.232) (2.850,-2.184) (2.800,-2.136) (2.762,-2.100) (2.725,-2.066) (2.681,-2.025) (2.626,-1.975) (2.575,-1.930) (2.525,-1.886) (2.484,-1.850) (2.426,-1.800) (2.375,-1.757) (2.325,-1.715) (2.277,-1.675) (2.225,-1.633) (2.185,-1.600) (2.125,-1.552) (2.075,-1.512) (2.029,-1.475) (1.975,-1.433) (1.933,-1.400) (1.875,-1.355) (1.825,-1.317) (1.775,-1.279) (1.725,-1.241) (1.675,-1.203) (1.625,-1.166) (1.575,-1.128) (1.525,-1.091) (1.475,-1.054) (1.425,-1.017) (1.375,-0.981) (1.333,-0.950) (1.275,-0.908) (1.230,-0.875) (1.175,-0.835) (1.126,-0.800) (1.075,-0.763) (1.025,-0.727) (0.975,-0.692) (0.925,-0.656) (0.882,-0.625) (0.825,-0.585) (0.777,-0.550) (0.725,-0.514) (0.675,-0.478) (0.625,-0.443) (0.575,-0.407) (0.531,-0.375) (0.475,-0.337) (0.425,-0.301) (0.375,-0.266) (0.325,-0.230) (0.284,-0.200) (0.225,-0.160) (0.206,-0.144)}
\EQ{(0.206,0.144) (0.248,0.175) (0.300,0.213) (0.350,0.248) (0.400,0.283) (0.450,0.319) (0.495,0.350) (0.550,0.390) (0.600,0.425) (0.650,0.460) (0.700,0.496) (0.742,0.525) (0.800,0.567) (0.847,0.600) (0.900,0.638) (0.950,0.674) (1.000,0.710) (1.050,0.745) (1.092,0.775) (1.150,0.817) (1.196,0.850) (1.250,0.890) (1.299,0.925) (1.350,0.963) (1.400,0.999) (1.450,1.036) (1.500,1.073) (1.550,1.110) (1.600,1.147) (1.650,1.184) (1.700,1.222) (1.750,1.260) (1.800,1.298) (1.850,1.336) (1.900,1.374) (1.950,1.413) (1.997,1.450) (2.050,1.492) (2.092,1.525) (2.150,1.572) (2.200,1.612) (2.246,1.650) (2.300,1.694) (2.350,1.736) (2.396,1.775) (2.450,1.821) (2.500,1.864) (2.541,1.900) (2.598,1.950) (2.650,1.997) (2.700,2.043) (2.750,2.089) (2.789,2.125) (2.825,2.160) (2.867,2.200) (2.918,2.250) (2.968,2.300) (3.017,2.350) (3.066,2.400) (3.113,2.450) (3.150,2.490) (3.200,2.545) (3.249,2.600) (3.293,2.650) (3.336,2.700) (3.375,2.747) (3.418,2.800) (3.450,2.840) (3.497,2.900) (3.535,2.950) (3.572,3.000) (3.608,3.050) (3.644,3.100) (3.675,3.146) (3.712,3.200) (3.744,3.250) (3.775,3.299) (3.807,3.350) (3.837,3.400) (3.866,3.450) (3.894,3.500) (3.922,3.550) (3.948,3.600) (3.974,3.650) (3.999,3.700) (4.023,3.750) (4.047,3.800) (4.069,3.850) (4.091,3.900) (4.112,3.950) (4.132,4.000) (4.152,4.050) (4.170,4.100) (4.188,4.150) (4.206,4.200) (4.222,4.250) (4.238,4.300) (4.253,4.350) (4.267,4.400) (4.280,4.450) (4.293,4.500) (4.305,4.550) (4.316,4.600) (4.327,4.650) (4.337,4.700) (4.346,4.750) (4.354,4.800) (4.362,4.850) (4.369,4.900) (4.376,4.950) (4.381,5.000) (4.386,5.050) (4.391,5.100) (4.394,5.150) (4.397,5.200) (4.400,5.250) (4.401,5.300) (4.402,5.350) (4.402,5.400) (4.402,5.450) (4.401,5.500) (4.399,5.550) (4.396,5.600) (4.393,5.650) (4.389,5.700) (4.385,5.750) (4.380,5.800) (4.374,5.850) (4.367,5.900) (4.360,5.950) (4.352,6.000) (4.343,6.050) (4.334,6.100) (4.323,6.150) (4.313,6.200) (4.301,6.250) (4.289,6.300) (4.276,6.350) (4.262,6.400) (4.247,6.450) (4.232,6.500) (4.216,6.550) (4.200,6.598) (4.182,6.650) (4.163,6.700) (4.144,6.750) (4.125,6.798) (4.104,6.850) (4.082,6.900) (4.060,6.950) (4.036,7.000) (4.012,7.050) (3.987,7.100) (3.961,7.150) (3.934,7.200) (3.907,7.250) (3.878,7.300) (3.850,7.347) (3.818,7.400) (3.786,7.450) (3.753,7.500) (3.725,7.542) (3.684,7.600) (3.650,7.648) (3.611,7.700) (3.575,7.747) (3.533,7.800) (3.500,7.841) (3.450,7.900) (3.406,7.950) (3.361,8.000) (3.325,8.039) (3.275,8.092) (3.225,8.142) (3.175,8.191) (3.125,8.238) (3.085,8.275) (3.029,8.325) (2.975,8.372) (2.925,8.414) (2.880,8.450) (2.825,8.494) (2.784,8.525) (2.725,8.569) (2.682,8.600) (2.625,8.640) (2.575,8.674) (2.525,8.707) (2.475,8.739) (2.425,8.770) (2.375,8.800) (2.325,8.829) (2.275,8.858) (2.225,8.886) (2.175,8.912) (2.125,8.938) (2.075,8.964) (2.025,8.988) (1.975,9.012) (1.925,9.035) (1.875,9.057) (1.825,9.078) (1.775,9.099) (1.725,9.119) (1.675,9.139) (1.625,9.158) (1.577,9.175) (1.525,9.193) (1.475,9.210) (1.425,9.226) (1.375,9.242) (1.325,9.257) (1.275,9.272) (1.225,9.285) (1.175,9.298) (1.125,9.311) (1.075,9.323) (1.025,9.335) (0.975,9.345) (0.925,9.356) (0.875,9.365) (0.825,9.375) (0.775,9.383) (0.725,9.391) (0.675,9.399) (0.625,9.406) (0.575,9.412) (0.525,9.418) (0.475,9.423) (0.425,9.428) (0.375,9.432) (0.325,9.436) (0.275,9.439) (0.225,9.442) (0.175,9.444) (0.125,9.446) (0.075,9.447) (0.025,9.447) (-0.050,9.447) (-0.100,9.446) (-0.150,9.445) (-0.200,9.443) (-0.250,9.441) (-0.300,9.438) (-0.350,9.434) (-0.400,9.430) (-0.450,9.426) (-0.500,9.421) (-0.550,9.415) (-0.600,9.409) (-0.650,9.402) (-0.700,9.395) (-0.750,9.387) (-0.800,9.379) (-0.850,9.370) (-0.900,9.361) (-0.950,9.351) (-1.000,9.340) (-1.050,9.329) (-1.100,9.317) (-1.150,9.305) (-1.200,9.292) (-1.250,9.278) (-1.300,9.264) (-1.349,9.250) (-1.400,9.234) (-1.450,9.218) (-1.500,9.202) (-1.550,9.185) (-1.600,9.167) (-1.650,9.148) (-1.700,9.129) (-1.750,9.109) (-1.800,9.089) (-1.850,9.068) (-1.900,9.046) (-1.946,9.025) (-2.000,9.000) (-2.050,8.976) (-2.100,8.951) (-2.150,8.925) (-2.198,8.900) (-2.244,8.875) (-2.289,8.850) (-2.333,8.825) (-2.400,8.785) (-2.450,8.755) (-2.497,8.725) (-2.550,8.691) (-2.600,8.657) (-2.647,8.625) (-2.700,8.587) (-2.750,8.551) (-2.800,8.513) (-2.849,8.475) (-2.900,8.434) (-2.942,8.400) (-3.000,8.351) (-3.050,8.307) (-3.100,8.262) (-3.139,8.225) (-3.192,8.175) (-3.242,8.125) (-3.291,8.075) (-3.338,8.025) (-3.375,7.985) (-3.425,7.929) (-3.471,7.875) (-3.513,7.825) (-3.550,7.779) (-3.592,7.725) (-3.625,7.682) (-3.667,7.625) (-3.700,7.578) (-3.736,7.525) (-3.770,7.475) (-3.800,7.428) (-3.833,7.375) (-3.863,7.325) (-3.892,7.275) (-3.921,7.225) (-3.948,7.175) (-3.974,7.125) (-4.000,7.075) (-4.024,7.025) (-4.048,6.975) (-4.071,6.925) (-4.093,6.875) (-4.114,6.825) (-4.134,6.775) (-4.154,6.725) (-4.173,6.675) (-4.191,6.625) (-4.208,6.575) (-4.224,6.525) (-4.240,6.475) (-4.255,6.425) (-4.269,6.375) (-4.282,6.325) (-4.295,6.275) (-4.307,6.225) (-4.318,6.175) (-4.329,6.125) (-4.338,6.075) (-4.347,6.025) (-4.356,5.975) (-4.364,5.925) (-4.370,5.875) (-4.377,5.825) (-4.382,5.775) (-4.387,5.725) (-4.391,5.675) (-4.395,5.625) (-4.398,5.575) (-4.400,5.525) (-4.401,5.475) (-4.402,5.425) (-4.402,5.375) (-4.402,5.325) (-4.401,5.275) (-4.399,5.225) (-4.396,5.175) (-4.393,5.125) (-4.389,5.075) (-4.384,5.025) (-4.379,4.975) (-4.373,4.925) (-4.366,4.875) (-4.358,4.825) (-4.350,4.775) (-4.342,4.725) (-4.332,4.675) (-4.322,4.625) (-4.311,4.575) (-4.299,4.525) (-4.287,4.475) (-4.274,4.425) (-4.260,4.375) (-4.245,4.325) (-4.230,4.275) (-4.214,4.225) (-4.197,4.175) (-4.180,4.125) (-4.161,4.075) (-4.142,4.025) (-4.122,3.975) (-4.102,3.925) (-4.080,3.875) (-4.058,3.825) (-4.035,3.775) (-4.011,3.725) (-3.987,3.675) (-3.961,3.625) (-3.935,3.575) (-3.908,3.525) (-3.880,3.475) (-3.852,3.425) (-3.825,3.380) (-3.792,3.325) (-3.760,3.275) (-3.728,3.225) (-3.700,3.183) (-3.661,3.125) (-3.626,3.075) (-3.590,3.025) (-3.554,2.975) (-3.516,2.925) (-3.478,2.875) (-3.438,2.825) (-3.400,2.778) (-3.357,2.725) (-3.315,2.675) (-3.275,2.629) (-3.227,2.575) (-3.182,2.525) (-3.136,2.475) (-3.100,2.436) (-3.050,2.384) (-3.000,2.332) (-2.950,2.282) (-2.900,2.232) (-2.850,2.184) (-2.800,2.136) (-2.762,2.100) (-2.725,2.066) (-2.681,2.025) (-2.626,1.975) (-2.575,1.930) (-2.525,1.886) (-2.484,1.850) (-2.426,1.800) (-2.375,1.757) (-2.325,1.715) (-2.277,1.675) (-2.225,1.633) (-2.185,1.600) (-2.125,1.552) (-2.075,1.512) (-2.029,1.475) (-1.975,1.433) (-1.933,1.400) (-1.875,1.355) (-1.825,1.317) (-1.775,1.279) (-1.725,1.241) (-1.675,1.203) (-1.625,1.166) (-1.575,1.128) (-1.525,1.091) (-1.475,1.054) (-1.425,1.017) (-1.375,0.981) (-1.333,0.950) (-1.275,0.908) (-1.230,0.875) (-1.175,0.835) (-1.126,0.800) (-1.075,0.763) (-1.025,0.727) (-0.975,0.692) (-0.925,0.656) (-0.882,0.625) (-0.825,0.585) (-0.777,0.550) (-0.725,0.514) (-0.675,0.478) (-0.625,0.443) (-0.575,0.407) (-0.531,0.375) (-0.475,0.337) (-0.425,0.301) (-0.375,0.266) (-0.325,0.230) (-0.284,0.200) (-0.225,0.160) (-0.220,0.155)}
\EQ{(-0.220,-0.155) (-0.275,-0.195) (-0.319,-0.225) (-0.375,-0.266) (-0.424,-0.300) (-0.475,-0.337) (-0.525,-0.372) (-0.575,-0.407) (-0.625,-0.443) (-0.671,-0.475) (-0.725,-0.514) (-0.775,-0.549) (-0.825,-0.585) (-0.875,-0.621) (-0.916,-0.650) (-0.975,-0.693) (-1.020,-0.725) (-1.075,-0.765) (-1.124,-0.800) (-1.175,-0.838) (-1.225,-0.874) (-1.275,-0.911) (-1.325,-0.948) (-1.375,-0.985) (-1.425,-1.022) (-1.475,-1.060) (-1.525,-1.098) (-1.575,-1.136) (-1.625,-1.174) (-1.675,-1.213) (-1.723,-1.250) (-1.775,-1.291) (-1.818,-1.325) (-1.875,-1.370) (-1.925,-1.411) (-1.973,-1.450) (-2.025,-1.493) (-2.075,-1.535) (-2.122,-1.575) (-2.175,-1.621) (-2.225,-1.664) (-2.265,-1.700) (-2.321,-1.750) (-2.375,-1.800) (-2.425,-1.846) (-2.475,-1.894) (-2.525,-1.942) (-2.575,-1.991) (-2.625,-2.042) (-2.675,-2.093) (-2.725,-2.146) (-2.775,-2.200) (-2.820,-2.250) (-2.864,-2.300) (-2.900,-2.342) (-2.949,-2.400) (-2.989,-2.450) (-3.025,-2.495) (-3.068,-2.550) (-3.100,-2.593) (-3.142,-2.650) (-3.175,-2.697) (-3.212,-2.750) (-3.245,-2.800) (-3.275,-2.846) (-3.309,-2.900) (-3.340,-2.950) (-3.369,-3.000) (-3.398,-3.050) (-3.425,-3.099) (-3.450,-3.146) (-3.475,-3.195) (-3.500,-3.245) (-3.525,-3.298) (-3.549,-3.350) (-3.571,-3.400) (-3.592,-3.450) (-3.612,-3.500) (-3.631,-3.550) (-3.650,-3.600) (-3.667,-3.650) (-3.684,-3.700) (-3.700,-3.750) (-3.715,-3.800) (-3.729,-3.850) (-3.742,-3.900) (-3.755,-3.950) (-3.767,-4.000) (-3.777,-4.050) (-3.787,-4.100) (-3.797,-4.150) (-3.805,-4.200) (-3.813,-4.250) (-3.820,-4.300) (-3.826,-4.350) (-3.831,-4.400) (-3.835,-4.450) (-3.839,-4.500) (-3.842,-4.550) (-3.844,-4.600) (-3.845,-4.650) (-3.846,-4.700) (-3.846,-4.750) (-3.845,-4.800) (-3.843,-4.850) (-3.840,-4.900) (-3.837,-4.950) (-3.833,-5.000) (-3.828,-5.050) (-3.822,-5.100) (-3.816,-5.150) (-3.808,-5.200) (-3.800,-5.250) (-3.791,-5.300) (-3.782,-5.350) (-3.771,-5.400) (-3.760,-5.450) (-3.747,-5.500) (-3.734,-5.550) (-3.721,-5.600) (-3.706,-5.650) (-3.690,-5.700) (-3.674,-5.750) (-3.656,-5.800) (-3.638,-5.850) (-3.619,-5.900) (-3.600,-5.948) (-3.578,-6.000) (-3.556,-6.050) (-3.533,-6.100) (-3.509,-6.150) (-3.484,-6.200) (-3.458,-6.250) (-3.431,-6.300) (-3.403,-6.350) (-3.375,-6.399) (-3.344,-6.450) (-3.313,-6.500) (-3.280,-6.550) (-3.250,-6.594) (-3.211,-6.650) (-3.175,-6.699) (-3.137,-6.750) (-3.100,-6.797) (-3.057,-6.850) (-3.015,-6.900) (-2.975,-6.946) (-2.926,-7.000) (-2.879,-7.050) (-2.830,-7.100) (-2.779,-7.150) (-2.727,-7.200) (-2.675,-7.247) (-2.625,-7.291) (-2.585,-7.325) (-2.525,-7.375) (-2.475,-7.414) (-2.429,-7.450) (-2.375,-7.490) (-2.326,-7.525) (-2.275,-7.560) (-2.225,-7.594) (-2.176,-7.625) (-2.125,-7.657) (-2.075,-7.688) (-2.025,-7.717) (-1.975,-7.745) (-1.925,-7.772) (-1.875,-7.799) (-1.825,-7.824) (-1.775,-7.849) (-1.725,-7.873) (-1.675,-7.896) (-1.625,-7.918) (-1.575,-7.939) (-1.525,-7.959) (-1.475,-7.979) (-1.425,-7.998) (-1.375,-8.016) (-1.325,-8.034) (-1.276,-8.050) (-1.225,-8.066) (-1.175,-8.082) (-1.125,-8.096) (-1.075,-8.110) (-1.025,-8.123) (-0.975,-8.136) (-0.925,-8.148) (-0.875,-8.159) (-0.825,-8.170) (-0.775,-8.179) (-0.725,-8.189) (-0.675,-8.197) (-0.625,-8.205) (-0.575,-8.213) (-0.525,-8.219) (-0.475,-8.225) (-0.425,-8.231) (-0.375,-8.236) (-0.325,-8.240) (-0.275,-8.244) (-0.225,-8.247) (-0.175,-8.249) (-0.125,-8.251) (-0.075,-8.252) (-0.025,-8.253) (0.050,-8.253) (0.100,-8.252) (0.150,-8.250) (0.200,-8.248) (0.250,-8.245) (0.300,-8.242) (0.350,-8.238) (0.400,-8.234) (0.450,-8.228) (0.500,-8.223) (0.550,-8.216) (0.600,-8.209) (0.650,-8.201) (0.700,-8.193) (0.750,-8.184) (0.800,-8.175) (0.850,-8.164) (0.900,-8.153) (0.950,-8.142) (1.000,-8.130) (1.050,-8.117) (1.100,-8.103) (1.150,-8.089) (1.200,-8.074) (1.250,-8.058) (1.300,-8.042) (1.350,-8.025) (1.400,-8.007) (1.450,-7.989) (1.500,-7.969) (1.548,-7.950) (1.600,-7.928) (1.650,-7.907) (1.700,-7.884) (1.750,-7.861) (1.800,-7.837) (1.850,-7.812) (1.900,-7.786) (1.950,-7.759) (2.000,-7.731) (2.050,-7.702) (2.096,-7.675) (2.150,-7.642) (2.200,-7.610) (2.250,-7.577) (2.300,-7.543) (2.350,-7.508) (2.395,-7.475) (2.450,-7.434) (2.493,-7.400) (2.550,-7.354) (2.600,-7.313) (2.644,-7.275) (2.700,-7.225) (2.750,-7.178) (2.800,-7.130) (2.850,-7.080) (2.900,-7.028) (2.949,-6.975) (2.993,-6.925) (3.036,-6.875) (3.075,-6.828) (3.117,-6.775) (3.150,-6.733) (3.193,-6.675) (3.225,-6.630) (3.263,-6.575) (3.296,-6.525) (3.325,-6.480) (3.359,-6.425) (3.389,-6.375) (3.418,-6.325) (3.445,-6.275) (3.472,-6.225) (3.497,-6.175) (3.521,-6.125) (3.545,-6.075) (3.567,-6.025) (3.589,-5.975) (3.609,-5.925) (3.629,-5.875) (3.647,-5.825) (3.665,-5.775) (3.682,-5.725) (3.698,-5.675) (3.713,-5.625) (3.728,-5.575) (3.741,-5.525) (3.754,-5.475) (3.765,-5.425) (3.776,-5.375) (3.787,-5.325) (3.796,-5.275) (3.804,-5.225) (3.812,-5.175) (3.819,-5.125) (3.825,-5.075) (3.830,-5.025) (3.835,-4.975) (3.839,-4.925) (3.842,-4.875) (3.844,-4.825) (3.845,-4.775) (3.846,-4.725) (3.846,-4.675) (3.845,-4.625) (3.843,-4.575) (3.841,-4.525) (3.837,-4.475) (3.833,-4.425) (3.828,-4.375) (3.823,-4.325) (3.816,-4.275) (3.809,-4.225) (3.801,-4.175) (3.792,-4.125) (3.783,-4.075) (3.772,-4.025) (3.761,-3.975) (3.749,-3.925) (3.736,-3.875) (3.722,-3.825) (3.708,-3.775) (3.692,-3.725) (3.676,-3.675) (3.659,-3.625) (3.641,-3.575) (3.622,-3.525) (3.602,-3.475) (3.582,-3.425) (3.560,-3.375) (3.538,-3.325) (3.514,-3.275) (3.490,-3.225) (3.465,-3.175) (3.439,-3.125) (3.412,-3.075) (3.384,-3.025) (3.355,-2.975) (3.325,-2.926) (3.294,-2.875) (3.262,-2.825) (3.229,-2.775) (3.200,-2.733) (3.160,-2.675) (3.125,-2.627) (3.087,-2.575) (3.050,-2.527) (3.009,-2.475) (2.975,-2.432) (2.928,-2.375) (2.886,-2.325) (2.850,-2.284) (2.800,-2.228) (2.752,-2.175) (2.705,-2.125) (2.657,-2.075) (2.609,-2.025) (2.559,-1.975) (2.508,-1.925) (2.456,-1.875) (2.403,-1.825) (2.350,-1.777) (2.300,-1.731) (2.250,-1.687) (2.209,-1.650) (2.151,-1.600) (2.100,-1.556) (2.050,-1.514) (2.003,-1.475) (1.950,-1.431) (1.900,-1.391) (1.850,-1.350) (1.800,-1.311) (1.755,-1.275) (1.700,-1.232) (1.659,-1.200) (1.600,-1.155) (1.550,-1.117) (1.500,-1.079) (1.450,-1.041) (1.400,-1.004) (1.350,-0.966) (1.300,-0.929) (1.250,-0.893) (1.200,-0.856) (1.158,-0.825) (1.100,-0.783) (1.055,-0.750) (1.000,-0.711) (0.951,-0.675) (0.900,-0.639) (0.850,-0.603) (0.800,-0.567) (0.750,-0.532) (0.706,-0.500) (0.650,-0.461) (0.600,-0.425) (0.550,-0.390) (0.500,-0.354) (0.450,-0.319) (0.400,-0.284) (0.354,-0.250) (0.300,-0.213) (0.250,-0.177) (0.206,-0.144)}
\EQ{(0.206,0.144) (0.248,0.175) (0.300,0.213) (0.350,0.248) (0.400,0.284) (0.450,0.319) (0.495,0.350) (0.550,0.390) (0.600,0.425) (0.650,0.461) (0.700,0.496) (0.742,0.525) (0.800,0.567) (0.846,0.600) (0.900,0.639) (0.950,0.675) (1.000,0.711) (1.050,0.747) (1.100,0.783) (1.150,0.819) (1.192,0.850) (1.250,0.893) (1.294,0.925) (1.350,0.966) (1.396,1.000) (1.450,1.041) (1.495,1.075) (1.550,1.117) (1.594,1.150) (1.650,1.193) (1.691,1.225) (1.750,1.271) (1.800,1.311) (1.850,1.350) (1.900,1.391) (1.942,1.425) (2.000,1.472) (2.050,1.514) (2.093,1.550) (2.150,1.599) (2.200,1.643) (2.250,1.687) (2.293,1.725) (2.348,1.775) (2.400,1.823) (2.450,1.870) (2.500,1.918) (2.550,1.967) (2.600,2.017) (2.650,2.068) (2.700,2.120) (2.750,2.173) (2.797,2.225) (2.842,2.275) (2.875,2.313) (2.925,2.372) (2.969,2.425) (3.009,2.475) (3.049,2.525) (3.087,2.575) (3.124,2.625) (3.160,2.675) (3.195,2.725) (3.225,2.770) (3.262,2.825) (3.294,2.875) (3.325,2.925) (3.355,2.975) (3.384,3.025) (3.412,3.075) (3.439,3.125) (3.465,3.175) (3.490,3.225) (3.514,3.275) (3.538,3.325) (3.560,3.375) (3.582,3.425) (3.602,3.475) (3.622,3.525) (3.641,3.575) (3.659,3.625) (3.675,3.673) (3.692,3.725) (3.708,3.775) (3.722,3.825) (3.736,3.875) (3.749,3.925) (3.761,3.975) (3.772,4.025) (3.783,4.075) (3.792,4.125) (3.801,4.175) (3.809,4.225) (3.816,4.275) (3.823,4.325) (3.828,4.375) (3.833,4.425) (3.837,4.475) (3.841,4.525) (3.843,4.575) (3.845,4.625) (3.846,4.675) (3.846,4.725) (3.845,4.775) (3.844,4.825) (3.842,4.875) (3.839,4.925) (3.835,4.975) (3.830,5.025) (3.825,5.075) (3.819,5.125) (3.812,5.175) (3.804,5.225) (3.796,5.275) (3.787,5.325) (3.776,5.375) (3.765,5.425) (3.754,5.475) (3.741,5.525) (3.728,5.575) (3.713,5.625) (3.698,5.675) (3.682,5.725) (3.665,5.775) (3.647,5.825) (3.629,5.875) (3.609,5.925) (3.589,5.975) (3.567,6.025) (3.545,6.075) (3.521,6.125) (3.497,6.175) (3.472,6.225) (3.445,6.275) (3.418,6.325) (3.389,6.375) (3.359,6.425) (3.328,6.475) (3.300,6.519) (3.263,6.575) (3.229,6.625) (3.193,6.675) (3.156,6.725) (3.125,6.765) (3.077,6.825) (3.036,6.875) (3.000,6.917) (2.950,6.974) (2.903,7.025) (2.855,7.075) (2.805,7.125) (2.753,7.175) (2.700,7.225) (2.650,7.269) (2.600,7.313) (2.555,7.350) (2.500,7.395) (2.450,7.434) (2.400,7.471) (2.350,7.508) (2.300,7.543) (2.253,7.575) (2.200,7.610) (2.150,7.642) (2.100,7.673) (2.054,7.700) (2.000,7.731) (1.950,7.759) (1.900,7.786) (1.850,7.812) (1.800,7.837) (1.750,7.861) (1.700,7.884) (1.650,7.907) (1.600,7.928) (1.550,7.949) (1.500,7.969) (1.450,7.989) (1.400,8.007) (1.350,8.025) (1.300,8.042) (1.250,8.058) (1.200,8.074) (1.150,8.089) (1.100,8.103) (1.050,8.117) (1.000,8.130) (0.950,8.142) (0.900,8.153) (0.850,8.164) (0.800,8.175) (0.750,8.184) (0.700,8.193) (0.650,8.201) (0.600,8.209) (0.550,8.216) (0.500,8.223) (0.450,8.228) (0.400,8.234) (0.350,8.238) (0.300,8.242) (0.250,8.245) (0.200,8.248) (0.150,8.250) (0.100,8.252) (0.050,8.253) (0.000,8.253) (-0.050,8.253) (-0.100,8.252) (-0.150,8.250) (-0.200,8.248) (-0.250,8.245) (-0.300,8.242) (-0.350,8.238) (-0.400,8.234) (-0.450,8.228) (-0.500,8.223) (-0.550,8.216) (-0.600,8.209) (-0.650,8.201) (-0.700,8.193) (-0.750,8.184) (-0.800,8.175) (-0.850,8.164) (-0.900,8.153) (-0.950,8.142) (-1.000,8.130) (-1.050,8.117) (-1.100,8.103) (-1.150,8.089) (-1.200,8.074) (-1.250,8.058) (-1.300,8.042) (-1.350,8.025) (-1.400,8.007) (-1.450,7.989) (-1.500,7.969) (-1.548,7.950) (-1.600,7.928) (-1.650,7.907) (-1.700,7.884) (-1.750,7.861) (-1.800,7.837) (-1.850,7.812) (-1.900,7.786) (-1.950,7.759) (-2.000,7.731) (-2.050,7.702) (-2.096,7.675) (-2.150,7.642) (-2.200,7.610) (-2.250,7.577) (-2.300,7.543) (-2.350,7.508) (-2.395,7.475) (-2.450,7.434) (-2.493,7.400) (-2.550,7.354) (-2.600,7.313) (-2.644,7.275) (-2.700,7.225) (-2.750,7.178) (-2.800,7.130) (-2.850,7.080) (-2.900,7.028) (-2.949,6.975) (-2.993,6.925) (-3.036,6.875) (-3.075,6.828) (-3.117,6.775) (-3.150,6.733) (-3.193,6.675) (-3.225,6.630) (-3.263,6.575) (-3.296,6.525) (-3.325,6.480) (-3.359,6.425) (-3.389,6.375) (-3.418,6.325) (-3.445,6.275) (-3.472,6.225) (-3.497,6.175) (-3.521,6.125) (-3.545,6.075) (-3.567,6.025) (-3.589,5.975) (-3.609,5.925) (-3.629,5.875) (-3.647,5.825) (-3.665,5.775) (-3.682,5.725) (-3.698,5.675) (-3.713,5.625) (-3.728,5.575) (-3.741,5.525) (-3.754,5.475) (-3.765,5.425) (-3.776,5.375) (-3.787,5.325) (-3.796,5.275) (-3.804,5.225) (-3.812,5.175) (-3.819,5.125) (-3.825,5.075) (-3.830,5.025) (-3.835,4.975) (-3.839,4.925) (-3.842,4.875) (-3.844,4.825) (-3.845,4.775) (-3.846,4.725) (-3.846,4.675) (-3.845,4.625) (-3.843,4.575) (-3.841,4.525) (-3.837,4.475) (-3.833,4.425) (-3.828,4.375) (-3.823,4.325) (-3.816,4.275) (-3.809,4.225) (-3.801,4.175) (-3.792,4.125) (-3.783,4.075) (-3.772,4.025) (-3.761,3.975) (-3.749,3.925) (-3.736,3.875) (-3.722,3.825) (-3.708,3.775) (-3.692,3.725) (-3.676,3.675) (-3.659,3.625) (-3.641,3.575) (-3.622,3.525) (-3.602,3.475) (-3.582,3.425) (-3.560,3.375) (-3.538,3.325) (-3.514,3.275) (-3.490,3.225) (-3.465,3.175) (-3.439,3.125) (-3.412,3.075) (-3.384,3.025) (-3.355,2.975) (-3.325,2.926) (-3.294,2.875) (-3.262,2.825) (-3.229,2.775) (-3.200,2.733) (-3.160,2.675) (-3.125,2.627) (-3.087,2.575) (-3.050,2.527) (-3.009,2.475) (-2.975,2.432) (-2.928,2.375) (-2.886,2.325) (-2.850,2.284) (-2.800,2.228) (-2.752,2.175) (-2.705,2.125) (-2.657,2.075) (-2.609,2.025) (-2.559,1.975) (-2.508,1.925) (-2.456,1.875) (-2.403,1.825) (-2.350,1.777) (-2.300,1.731) (-2.250,1.687) (-2.209,1.650) (-2.151,1.600) (-2.100,1.556) (-2.050,1.514) (-2.003,1.475) (-1.950,1.431) (-1.900,1.391) (-1.850,1.350) (-1.800,1.311) (-1.755,1.275) (-1.700,1.232) (-1.659,1.200) (-1.600,1.155) (-1.550,1.117) (-1.500,1.079) (-1.450,1.041) (-1.400,1.004) (-1.350,0.966) (-1.300,0.929) (-1.250,0.893) (-1.200,0.856) (-1.158,0.825) (-1.100,0.783) (-1.055,0.750) (-1.000,0.711) (-0.951,0.675) (-0.900,0.639) (-0.850,0.603) (-0.800,0.567) (-0.750,0.532) (-0.706,0.500) (-0.650,0.461) (-0.600,0.425) (-0.550,0.390) (-0.500,0.354) (-0.450,0.319) (-0.400,0.284) (-0.354,0.250) (-0.300,0.213) (-0.250,0.177) (-0.220,0.155)}
\EQ{(-0.220,-0.155) (-0.275,-0.195) (-0.319,-0.225) (-0.375,-0.266) (-0.424,-0.300) (-0.475,-0.337) (-0.525,-0.372) (-0.575,-0.407) (-0.625,-0.443) (-0.671,-0.475) (-0.725,-0.514) (-0.775,-0.550) (-0.825,-0.586) (-0.875,-0.622) (-0.925,-0.658) (-0.975,-0.694) (-1.019,-0.725) (-1.075,-0.767) (-1.121,-0.800) (-1.175,-0.840) (-1.223,-0.875) (-1.275,-0.914) (-1.323,-0.950) (-1.375,-0.989) (-1.422,-1.025) (-1.475,-1.066) (-1.520,-1.100) (-1.575,-1.143) (-1.616,-1.175) (-1.675,-1.222) (-1.725,-1.263) (-1.771,-1.300) (-1.825,-1.344) (-1.875,-1.386) (-1.921,-1.425) (-1.975,-1.472) (-2.025,-1.515) (-2.064,-1.550) (-2.120,-1.600) (-2.174,-1.650) (-2.225,-1.698) (-2.275,-1.746) (-2.325,-1.795) (-2.375,-1.845) (-2.425,-1.896) (-2.475,-1.949) (-2.522,-2.000) (-2.567,-2.050) (-2.611,-2.100) (-2.650,-2.146) (-2.695,-2.200) (-2.735,-2.250) (-2.774,-2.300) (-2.812,-2.350) (-2.849,-2.400) (-2.885,-2.450) (-2.919,-2.500) (-2.950,-2.546) (-2.985,-2.600) (-3.016,-2.650) (-3.046,-2.700) (-3.075,-2.750) (-3.100,-2.794) (-3.125,-2.841) (-3.150,-2.889) (-3.175,-2.939) (-3.200,-2.991) (-3.225,-3.045) (-3.249,-3.100) (-3.270,-3.150) (-3.290,-3.200) (-3.309,-3.250) (-3.326,-3.300) (-3.343,-3.350) (-3.359,-3.400) (-3.374,-3.450) (-3.389,-3.500) (-3.402,-3.550) (-3.414,-3.600) (-3.426,-3.650) (-3.436,-3.700) (-3.446,-3.750) (-3.455,-3.800) (-3.462,-3.850) (-3.469,-3.900) (-3.476,-3.950) (-3.481,-4.000) (-3.485,-4.050) (-3.489,-4.100) (-3.491,-4.150) (-3.493,-4.200) (-3.494,-4.250) (-3.494,-4.300) (-3.493,-4.350) (-3.492,-4.400) (-3.489,-4.450) (-3.486,-4.500) (-3.482,-4.550) (-3.477,-4.600) (-3.471,-4.650) (-3.464,-4.700) (-3.456,-4.750) (-3.447,-4.800) (-3.438,-4.850) (-3.428,-4.900) (-3.416,-4.950) (-3.404,-5.000) (-3.391,-5.050) (-3.377,-5.100) (-3.362,-5.150) (-3.346,-5.200) (-3.329,-5.250) (-3.311,-5.300) (-3.292,-5.350) (-3.272,-5.400) (-3.251,-5.450) (-3.229,-5.500) (-3.206,-5.550) (-3.182,-5.600) (-3.157,-5.650) (-3.131,-5.700) (-3.103,-5.750) (-3.075,-5.799) (-3.044,-5.850) (-3.013,-5.900) (-2.981,-5.950) (-2.950,-5.995) (-2.912,-6.050) (-2.875,-6.100) (-2.837,-6.150) (-2.800,-6.196) (-2.755,-6.250) (-2.712,-6.300) (-2.675,-6.342) (-2.625,-6.396) (-2.575,-6.447) (-2.525,-6.496) (-2.475,-6.544) (-2.425,-6.589) (-2.384,-6.625) (-2.325,-6.675) (-2.275,-6.716) (-2.231,-6.750) (-2.175,-6.792) (-2.130,-6.825) (-2.075,-6.863) (-2.025,-6.897) (-1.982,-6.925) (-1.925,-6.961) (-1.875,-6.991) (-1.825,-7.020) (-1.775,-7.048) (-1.725,-7.074) (-1.676,-7.100) (-1.625,-7.125) (-1.575,-7.149) (-1.525,-7.172) (-1.475,-7.194) (-1.425,-7.215) (-1.375,-7.236) (-1.325,-7.255) (-1.275,-7.274) (-1.225,-7.292) (-1.175,-7.309) (-1.125,-7.325) (-1.075,-7.340) (-1.025,-7.355) (-0.975,-7.369) (-0.925,-7.382) (-0.875,-7.395) (-0.825,-7.406) (-0.775,-7.417) (-0.725,-7.427) (-0.675,-7.437) (-0.625,-7.446) (-0.575,-7.454) (-0.525,-7.461) (-0.475,-7.468) (-0.425,-7.474) (-0.375,-7.479) (-0.325,-7.484) (-0.275,-7.488) (-0.225,-7.492) (-0.175,-7.494) (-0.125,-7.496) (-0.075,-7.498) (-0.025,-7.498) (0.050,-7.498) (0.100,-7.497) (0.150,-7.495) (0.200,-7.493) (0.250,-7.490) (0.300,-7.486) (0.350,-7.482) (0.400,-7.477) (0.450,-7.471) (0.500,-7.465) (0.550,-7.458) (0.599,-7.450) (0.650,-7.441) (0.700,-7.432) (0.750,-7.422) (0.800,-7.412) (0.850,-7.400) (0.900,-7.388) (0.950,-7.376) (1.000,-7.362) (1.050,-7.348) (1.100,-7.333) (1.150,-7.317) (1.200,-7.300) (1.250,-7.283) (1.300,-7.265) (1.350,-7.246) (1.400,-7.226) (1.450,-7.205) (1.500,-7.183) (1.550,-7.161) (1.600,-7.137) (1.650,-7.113) (1.700,-7.088) (1.750,-7.061) (1.800,-7.034) (1.850,-7.005) (1.900,-6.976) (1.950,-6.945) (2.000,-6.913) (2.050,-6.880) (2.094,-6.850) (2.150,-6.811) (2.198,-6.775) (2.250,-6.735) (2.295,-6.700) (2.350,-6.654) (2.400,-6.611) (2.441,-6.575) (2.495,-6.525) (2.547,-6.475) (2.597,-6.425) (2.644,-6.375) (2.690,-6.325) (2.725,-6.286) (2.775,-6.227) (2.817,-6.175) (2.850,-6.133) (2.893,-6.075) (2.925,-6.031) (2.964,-5.975) (2.997,-5.925) (3.025,-5.881) (3.060,-5.825) (3.089,-5.775) (3.117,-5.725) (3.144,-5.675) (3.170,-5.625) (3.195,-5.575) (3.218,-5.525) (3.241,-5.475) (3.262,-5.425) (3.282,-5.375) (3.302,-5.325) (3.320,-5.275) (3.338,-5.225) (3.354,-5.175) (3.370,-5.125) (3.384,-5.075) (3.398,-5.025) (3.410,-4.975) (3.422,-4.925) (3.433,-4.875) (3.443,-4.825) (3.452,-4.775) (3.460,-4.725) (3.467,-4.675) (3.474,-4.625) (3.479,-4.575) (3.484,-4.525) (3.488,-4.475) (3.491,-4.425) (3.493,-4.375) (3.494,-4.325) (3.494,-4.275) (3.494,-4.225) (3.492,-4.175) (3.490,-4.125) (3.487,-4.075) (3.483,-4.025) (3.478,-3.975) (3.473,-3.925) (3.466,-3.875) (3.459,-3.825) (3.450,-3.775) (3.441,-3.725) (3.431,-3.675) (3.420,-3.625) (3.408,-3.575) (3.395,-3.525) (3.382,-3.475) (3.367,-3.425) (3.352,-3.375) (3.335,-3.325) (3.318,-3.275) (3.300,-3.227) (3.280,-3.175) (3.260,-3.125) (3.238,-3.075) (3.216,-3.025) (3.193,-2.975) (3.168,-2.925) (3.143,-2.875) (3.117,-2.825) (3.089,-2.775) (3.061,-2.725) (3.031,-2.675) (3.001,-2.625) (2.969,-2.575) (2.936,-2.525) (2.902,-2.475) (2.867,-2.425) (2.831,-2.375) (2.800,-2.334) (2.755,-2.275) (2.715,-2.225) (2.675,-2.176) (2.632,-2.125) (2.589,-2.075) (2.550,-2.031) (2.500,-1.976) (2.452,-1.925) (2.404,-1.875) (2.355,-1.825) (2.305,-1.775) (2.254,-1.725) (2.201,-1.675) (2.150,-1.628) (2.100,-1.582) (2.050,-1.537) (2.008,-1.500) (1.950,-1.450) (1.900,-1.407) (1.850,-1.365) (1.802,-1.325) (1.750,-1.283) (1.710,-1.250) (1.650,-1.202) (1.600,-1.163) (1.552,-1.125) (1.500,-1.085) (1.455,-1.050) (1.400,-1.008) (1.357,-0.975) (1.300,-0.933) (1.257,-0.900) (1.200,-0.858) (1.155,-0.825) (1.100,-0.785) (1.053,-0.750) (1.000,-0.712) (0.950,-0.676) (0.900,-0.640) (0.850,-0.604) (0.800,-0.568) (0.750,-0.532) (0.706,-0.500) (0.650,-0.461) (0.600,-0.425) (0.550,-0.390) (0.500,-0.354) (0.450,-0.319) (0.400,-0.284) (0.354,-0.250) (0.300,-0.213) (0.250,-0.177) (0.206,-0.144)}
\EQ{(0.206,0.144) (0.248,0.175) (0.300,0.213) (0.350,0.248) (0.400,0.284) (0.450,0.319) (0.495,0.350) (0.550,0.390) (0.600,0.425) (0.650,0.461) (0.700,0.496) (0.741,0.525) (0.800,0.568) (0.845,0.600) (0.900,0.640) (0.949,0.675) (1.000,0.712) (1.050,0.748) (1.100,0.785) (1.150,0.821) (1.200,0.858) (1.250,0.895) (1.300,0.933) (1.350,0.970) (1.400,1.008) (1.450,1.046) (1.500,1.085) (1.550,1.124) (1.600,1.163) (1.647,1.200) (1.700,1.242) (1.741,1.275) (1.800,1.324) (1.850,1.365) (1.892,1.400) (1.950,1.450) (2.000,1.493) (2.050,1.537) (2.092,1.575) (2.147,1.625) (2.200,1.674) (2.250,1.722) (2.300,1.770) (2.350,1.820) (2.400,1.871) (2.450,1.923) (2.499,1.975) (2.545,2.025) (2.589,2.075) (2.625,2.117) (2.674,2.175) (2.715,2.225) (2.750,2.269) (2.793,2.325) (2.825,2.367) (2.867,2.425) (2.900,2.472) (2.936,2.525) (2.969,2.575) (3.000,2.624) (3.031,2.675) (3.061,2.725) (3.089,2.775) (3.117,2.825) (3.143,2.875) (3.168,2.925) (3.193,2.975) (3.216,3.025) (3.238,3.075) (3.260,3.125) (3.280,3.175) (3.299,3.225) (3.318,3.275) (3.335,3.325) (3.352,3.375) (3.367,3.425) (3.382,3.475) (3.395,3.525) (3.408,3.575) (3.420,3.625) (3.431,3.675) (3.441,3.725) (3.450,3.775) (3.459,3.825) (3.466,3.875) (3.473,3.925) (3.478,3.975) (3.483,4.025) (3.487,4.075) (3.490,4.125) (3.492,4.175) (3.494,4.225) (3.494,4.275) (3.494,4.325) (3.493,4.375) (3.491,4.425) (3.488,4.475) (3.484,4.525) (3.479,4.575) (3.474,4.625) (3.467,4.675) (3.460,4.725) (3.452,4.775) (3.443,4.825) (3.433,4.875) (3.422,4.925) (3.410,4.975) (3.398,5.025) (3.384,5.075) (3.370,5.125) (3.354,5.175) (3.338,5.225) (3.320,5.275) (3.302,5.325) (3.282,5.375) (3.262,5.425) (3.241,5.475) (3.218,5.525) (3.195,5.575) (3.170,5.625) (3.144,5.675) (3.117,5.725) (3.089,5.775) (3.060,5.825) (3.029,5.875) (3.000,5.921) (2.964,5.975) (2.929,6.025) (2.900,6.066) (2.856,6.125) (2.825,6.165) (2.776,6.225) (2.734,6.275) (2.700,6.314) (2.650,6.369) (2.600,6.422) (2.550,6.472) (2.500,6.520) (2.450,6.567) (2.400,6.611) (2.355,6.650) (2.300,6.696) (2.250,6.735) (2.200,6.774) (2.150,6.811) (2.100,6.846) (2.058,6.875) (2.000,6.913) (1.950,6.945) (1.901,6.975) (1.850,7.005) (1.800,7.034) (1.750,7.061) (1.700,7.088) (1.650,7.113) (1.600,7.137) (1.550,7.161) (1.500,7.183) (1.450,7.205) (1.402,7.225) (1.350,7.246) (1.300,7.265) (1.250,7.283) (1.201,7.300) (1.150,7.317) (1.100,7.333) (1.050,7.348) (1.000,7.362) (0.950,7.376) (0.900,7.388) (0.850,7.400) (0.800,7.412) (0.750,7.422) (0.700,7.432) (0.650,7.441) (0.600,7.450) (0.550,7.458) (0.500,7.465) (0.450,7.471) (0.400,7.477) (0.350,7.482) (0.300,7.486) (0.250,7.490) (0.200,7.493) (0.150,7.495) (0.100,7.497) (0.050,7.498) (0.000,7.498) (-0.050,7.498) (-0.100,7.497) (-0.150,7.495) (-0.200,7.493) (-0.250,7.490) (-0.300,7.486) (-0.350,7.482) (-0.400,7.477) (-0.450,7.471) (-0.500,7.465) (-0.550,7.458) (-0.599,7.450) (-0.650,7.441) (-0.700,7.432) (-0.750,7.422) (-0.800,7.412) (-0.850,7.400) (-0.900,7.388) (-0.950,7.376) (-1.000,7.362) (-1.050,7.348) (-1.100,7.333) (-1.150,7.317) (-1.200,7.300) (-1.250,7.283) (-1.300,7.265) (-1.350,7.246) (-1.400,7.226) (-1.450,7.205) (-1.500,7.183) (-1.550,7.161) (-1.600,7.137) (-1.650,7.113) (-1.700,7.088) (-1.750,7.061) (-1.800,7.034) (-1.850,7.005) (-1.900,6.976) (-1.950,6.945) (-2.000,6.913) (-2.050,6.880) (-2.094,6.850) (-2.150,6.811) (-2.198,6.775) (-2.250,6.735) (-2.295,6.700) (-2.350,6.654) (-2.400,6.611) (-2.441,6.575) (-2.495,6.525) (-2.547,6.475) (-2.597,6.425) (-2.644,6.375) (-2.690,6.325) (-2.725,6.286) (-2.775,6.227) (-2.817,6.175) (-2.850,6.133) (-2.893,6.075) (-2.925,6.031) (-2.964,5.975) (-2.997,5.925) (-3.025,5.881) (-3.060,5.825) (-3.089,5.775) (-3.117,5.725) (-3.144,5.675) (-3.170,5.625) (-3.195,5.575) (-3.218,5.525) (-3.241,5.475) (-3.262,5.425) (-3.282,5.375) (-3.302,5.325) (-3.320,5.275) (-3.338,5.225) (-3.354,5.175) (-3.370,5.125) (-3.384,5.075) (-3.398,5.025) (-3.410,4.975) (-3.422,4.925) (-3.433,4.875) (-3.443,4.825) (-3.452,4.775) (-3.460,4.725) (-3.467,4.675) (-3.474,4.625) (-3.479,4.575) (-3.484,4.525) (-3.488,4.475) (-3.491,4.425) (-3.493,4.375) (-3.494,4.325) (-3.494,4.275) (-3.494,4.225) (-3.492,4.175) (-3.490,4.125) (-3.487,4.075) (-3.483,4.025) (-3.478,3.975) (-3.473,3.925) (-3.466,3.875) (-3.459,3.825) (-3.450,3.775) (-3.441,3.725) (-3.431,3.675) (-3.420,3.625) (-3.408,3.575) (-3.395,3.525) (-3.382,3.475) (-3.367,3.425) (-3.352,3.375) (-3.335,3.325) (-3.318,3.275) (-3.300,3.227) (-3.280,3.175) (-3.260,3.125) (-3.238,3.075) (-3.216,3.025) (-3.193,2.975) (-3.168,2.925) (-3.143,2.875) (-3.117,2.825) (-3.089,2.775) (-3.061,2.725) (-3.031,2.675) (-3.001,2.625) (-2.969,2.575) (-2.936,2.525) (-2.902,2.475) (-2.867,2.425) (-2.831,2.375) (-2.800,2.334) (-2.755,2.275) (-2.715,2.225) (-2.675,2.176) (-2.632,2.125) (-2.589,2.075) (-2.550,2.031) (-2.500,1.976) (-2.452,1.925) (-2.404,1.875) (-2.355,1.825) (-2.305,1.775) (-2.254,1.725) (-2.201,1.675) (-2.150,1.628) (-2.100,1.582) (-2.050,1.537) (-2.008,1.500) (-1.950,1.450) (-1.900,1.407) (-1.850,1.365) (-1.802,1.325) (-1.750,1.283) (-1.710,1.250) (-1.650,1.202) (-1.600,1.163) (-1.552,1.125) (-1.500,1.085) (-1.455,1.050) (-1.400,1.008) (-1.357,0.975) (-1.300,0.933) (-1.257,0.900) (-1.200,0.858) (-1.155,0.825) (-1.100,0.785) (-1.053,0.750) (-1.000,0.712) (-0.950,0.676) (-0.900,0.640) (-0.850,0.604) (-0.800,0.568) (-0.750,0.532) (-0.706,0.500) (-0.650,0.461) (-0.600,0.425) (-0.550,0.390) (-0.500,0.354) (-0.450,0.319) (-0.400,0.284) (-0.354,0.250) (-0.300,0.213) (-0.250,0.177) (-0.220,0.155)}
\EQ{(-0.220,-0.155) (-0.275,-0.195) (-0.319,-0.225) (-0.375,-0.266) (-0.424,-0.300) (-0.475,-0.337) (-0.525,-0.372) (-0.575,-0.408) (-0.625,-0.443) (-0.671,-0.475) (-0.725,-0.515) (-0.775,-0.550) (-0.825,-0.586) (-0.875,-0.622) (-0.925,-0.659) (-0.975,-0.695) (-1.017,-0.725) (-1.075,-0.768) (-1.119,-0.800) (-1.175,-0.842) (-1.219,-0.875) (-1.275,-0.917) (-1.319,-0.950) (-1.375,-0.994) (-1.416,-1.025) (-1.475,-1.071) (-1.525,-1.111) (-1.574,-1.150) (-1.625,-1.191) (-1.667,-1.225) (-1.725,-1.274) (-1.775,-1.316) (-1.815,-1.350) (-1.872,-1.400) (-1.925,-1.447) (-1.975,-1.492) (-2.025,-1.538) (-2.064,-1.575) (-2.117,-1.625) (-2.167,-1.675) (-2.217,-1.725) (-2.265,-1.775) (-2.300,-1.812) (-2.335,-1.850) (-2.375,-1.894) (-2.424,-1.950) (-2.467,-2.000) (-2.500,-2.040) (-2.548,-2.100) (-2.587,-2.150) (-2.624,-2.200) (-2.661,-2.250) (-2.696,-2.300) (-2.725,-2.343) (-2.762,-2.400) (-2.794,-2.450) (-2.824,-2.500) (-2.850,-2.545) (-2.875,-2.589) (-2.900,-2.635) (-2.925,-2.683) (-2.950,-2.733) (-2.975,-2.786) (-3.000,-2.841) (-3.025,-2.899) (-3.046,-2.950) (-3.065,-3.000) (-3.083,-3.050) (-3.100,-3.099) (-3.117,-3.150) (-3.132,-3.200) (-3.146,-3.250) (-3.159,-3.300) (-3.171,-3.350) (-3.183,-3.400) (-3.193,-3.450) (-3.202,-3.500) (-3.211,-3.550) (-3.218,-3.600) (-3.224,-3.650) (-3.230,-3.700) (-3.235,-3.750) (-3.238,-3.800) (-3.241,-3.850) (-3.243,-3.900) (-3.244,-3.950) (-3.244,-4.000) (-3.243,-4.050) (-3.241,-4.100) (-3.238,-4.150) (-3.234,-4.200) (-3.229,-4.250) (-3.224,-4.300) (-3.217,-4.350) (-3.210,-4.400) (-3.201,-4.450) (-3.192,-4.500) (-3.182,-4.550) (-3.170,-4.600) (-3.158,-4.650) (-3.145,-4.700) (-3.130,-4.750) (-3.115,-4.800) (-3.099,-4.850) (-3.081,-4.900) (-3.063,-4.950) (-3.043,-5.000) (-3.022,-5.050) (-3.001,-5.100) (-2.978,-5.150) (-2.954,-5.200) (-2.928,-5.250) (-2.902,-5.300) (-2.875,-5.348) (-2.845,-5.400) (-2.814,-5.450) (-2.782,-5.500) (-2.750,-5.549) (-2.715,-5.600) (-2.678,-5.650) (-2.640,-5.700) (-2.601,-5.750) (-2.560,-5.800) (-2.525,-5.840) (-2.475,-5.896) (-2.425,-5.949) (-2.375,-6.000) (-2.325,-6.049) (-2.275,-6.096) (-2.225,-6.140) (-2.184,-6.175) (-2.125,-6.224) (-2.075,-6.263) (-2.026,-6.300) (-1.975,-6.337) (-1.925,-6.372) (-1.875,-6.406) (-1.825,-6.438) (-1.775,-6.469) (-1.725,-6.499) (-1.679,-6.525) (-1.634,-6.550) (-1.587,-6.575) (-1.538,-6.600) (-1.487,-6.625) (-1.433,-6.650) (-1.377,-6.675) (-1.325,-6.697) (-1.275,-6.718) (-1.225,-6.737) (-1.175,-6.756) (-1.125,-6.773) (-1.075,-6.790) (-1.025,-6.806) (-0.975,-6.821) (-0.925,-6.835) (-0.875,-6.849) (-0.825,-6.861) (-0.775,-6.873) (-0.725,-6.884) (-0.675,-6.895) (-0.625,-6.904) (-0.575,-6.913) (-0.525,-6.921) (-0.475,-6.928) (-0.425,-6.935) (-0.375,-6.941) (-0.325,-6.946) (-0.275,-6.950) (-0.225,-6.954) (-0.175,-6.956) (-0.125,-6.959) (-0.075,-6.960) (-0.025,-6.961) (0.050,-6.960) (0.100,-6.959) (0.150,-6.958) (0.200,-6.955) (0.250,-6.952) (0.300,-6.948) (0.350,-6.943) (0.400,-6.938) (0.450,-6.932) (0.500,-6.925) (0.550,-6.917) (0.600,-6.909) (0.650,-6.899) (0.700,-6.890) (0.750,-6.879) (0.800,-6.867) (0.850,-6.855) (0.900,-6.842) (0.950,-6.828) (1.000,-6.814) (1.050,-6.798) (1.100,-6.782) (1.150,-6.765) (1.200,-6.747) (1.250,-6.728) (1.300,-6.708) (1.350,-6.687) (1.400,-6.665) (1.450,-6.642) (1.500,-6.619) (1.550,-6.594) (1.600,-6.568) (1.650,-6.541) (1.700,-6.513) (1.750,-6.484) (1.800,-6.454) (1.845,-6.425) (1.900,-6.389) (1.950,-6.355) (1.992,-6.325) (2.050,-6.282) (2.092,-6.250) (2.150,-6.204) (2.200,-6.162) (2.242,-6.125) (2.297,-6.075) (2.350,-6.025) (2.400,-5.975) (2.448,-5.925) (2.494,-5.875) (2.538,-5.825) (2.575,-5.782) (2.621,-5.725) (2.660,-5.675) (2.697,-5.625) (2.732,-5.575) (2.766,-5.525) (2.799,-5.475) (2.830,-5.425) (2.859,-5.375) (2.888,-5.325) (2.915,-5.275) (2.941,-5.225) (2.966,-5.175) (2.989,-5.125) (3.012,-5.075) (3.033,-5.025) (3.053,-4.975) (3.072,-4.925) (3.090,-4.875) (3.107,-4.825) (3.123,-4.775) (3.138,-4.725) (3.151,-4.675) (3.164,-4.625) (3.176,-4.575) (3.187,-4.525) (3.197,-4.475) (3.206,-4.425) (3.214,-4.375) (3.221,-4.325) (3.227,-4.275) (3.232,-4.225) (3.236,-4.175) (3.239,-4.125) (3.242,-4.075) (3.243,-4.025) (3.244,-3.975) (3.243,-3.925) (3.242,-3.875) (3.240,-3.825) (3.237,-3.775) (3.232,-3.725) (3.227,-3.675) (3.221,-3.625) (3.214,-3.575) (3.207,-3.525) (3.198,-3.475) (3.188,-3.425) (3.177,-3.375) (3.165,-3.325) (3.153,-3.275) (3.139,-3.225) (3.124,-3.175) (3.109,-3.125) (3.092,-3.075) (3.075,-3.027) (3.056,-2.975) (3.036,-2.925) (3.015,-2.875) (2.993,-2.825) (2.970,-2.775) (2.946,-2.725) (2.921,-2.675) (2.895,-2.625) (2.867,-2.575) (2.839,-2.525) (2.809,-2.475) (2.778,-2.425) (2.750,-2.381) (2.713,-2.325) (2.678,-2.275) (2.643,-2.225) (2.606,-2.175) (2.575,-2.135) (2.528,-2.075) (2.488,-2.025) (2.450,-1.980) (2.402,-1.925) (2.358,-1.875) (2.312,-1.825) (2.275,-1.785) (2.225,-1.733) (2.175,-1.683) (2.125,-1.633) (2.075,-1.585) (2.038,-1.550) (2.000,-1.515) (1.957,-1.475) (1.901,-1.425) (1.850,-1.381) (1.800,-1.337) (1.756,-1.300) (1.700,-1.253) (1.650,-1.212) (1.605,-1.175) (1.550,-1.131) (1.500,-1.091) (1.450,-1.052) (1.400,-1.013) (1.351,-0.975) (1.300,-0.936) (1.252,-0.900) (1.200,-0.861) (1.152,-0.825) (1.100,-0.787) (1.051,-0.750) (1.000,-0.713) (0.950,-0.677) (0.900,-0.640) (0.850,-0.604) (0.800,-0.568) (0.750,-0.532) (0.705,-0.500) (0.650,-0.461) (0.600,-0.425) (0.550,-0.390) (0.500,-0.354) (0.450,-0.319) (0.400,-0.284) (0.354,-0.250) (0.300,-0.213) (0.250,-0.177) (0.206,-0.144)}
\EQ{(0.206,0.144) (0.248,0.175) (0.300,0.213) (0.350,0.248) (0.400,0.284) (0.450,0.319) (0.495,0.350) (0.550,0.390) (0.600,0.425) (0.650,0.461) (0.700,0.497) (0.750,0.532) (0.800,0.568) (0.845,0.600) (0.900,0.640) (0.948,0.675) (1.000,0.713) (1.050,0.750) (1.100,0.787) (1.150,0.823) (1.200,0.861) (1.250,0.898) (1.300,0.936) (1.350,0.974) (1.400,1.013) (1.448,1.050) (1.500,1.091) (1.543,1.125) (1.600,1.171) (1.650,1.212) (1.697,1.250) (1.750,1.295) (1.800,1.337) (1.844,1.375) (1.900,1.424) (1.950,1.469) (2.000,1.515) (2.038,1.550) (2.075,1.585) (2.117,1.625) (2.167,1.675) (2.217,1.725) (2.265,1.775) (2.300,1.812) (2.335,1.850) (2.375,1.894) (2.424,1.950) (2.467,2.000) (2.500,2.040) (2.548,2.100) (2.587,2.150) (2.624,2.200) (2.661,2.250) (2.696,2.300) (2.725,2.343) (2.762,2.400) (2.794,2.450) (2.824,2.500) (2.850,2.545) (2.875,2.589) (2.900,2.635) (2.925,2.683) (2.950,2.733) (2.975,2.786) (3.000,2.841) (3.025,2.899) (3.046,2.950) (3.065,3.000) (3.083,3.050) (3.100,3.099) (3.117,3.150) (3.132,3.200) (3.146,3.250) (3.159,3.300) (3.171,3.350) (3.183,3.400) (3.193,3.450) (3.202,3.500) (3.211,3.550) (3.218,3.600) (3.224,3.650) (3.230,3.700) (3.235,3.750) (3.238,3.800) (3.241,3.850) (3.243,3.900) (3.244,3.950) (3.244,4.000) (3.243,4.050) (3.241,4.100) (3.238,4.150) (3.234,4.200) (3.229,4.250) (3.224,4.300) (3.217,4.350) (3.210,4.400) (3.201,4.450) (3.192,4.500) (3.182,4.550) (3.170,4.600) (3.158,4.650) (3.145,4.700) (3.130,4.750) (3.115,4.800) (3.099,4.850) (3.081,4.900) (3.063,4.950) (3.043,5.000) (3.022,5.050) (3.001,5.100) (2.978,5.150) (2.954,5.200) (2.928,5.250) (2.902,5.300) (2.875,5.348) (2.845,5.400) (2.814,5.450) (2.782,5.500) (2.750,5.549) (2.715,5.600) (2.678,5.650) (2.640,5.700) (2.601,5.750) (2.560,5.800) (2.525,5.840) (2.475,5.896) (2.425,5.949) (2.375,6.000) (2.325,6.049) (2.275,6.096) (2.225,6.140) (2.184,6.175) (2.125,6.224) (2.075,6.263) (2.026,6.300) (1.975,6.337) (1.925,6.372) (1.875,6.406) (1.825,6.438) (1.775,6.469) (1.725,6.499) (1.679,6.525) (1.634,6.550) (1.587,6.575) (1.538,6.600) (1.487,6.625) (1.433,6.650) (1.377,6.675) (1.325,6.697) (1.275,6.718) (1.225,6.737) (1.175,6.756) (1.125,6.773) (1.075,6.790) (1.025,6.806) (0.975,6.821) (0.925,6.835) (0.875,6.849) (0.825,6.861) (0.775,6.873) (0.725,6.884) (0.675,6.895) (0.625,6.904) (0.575,6.913) (0.525,6.921) (0.475,6.928) (0.425,6.935) (0.375,6.941) (0.325,6.946) (0.275,6.950) (0.225,6.954) (0.175,6.956) (0.125,6.959) (0.075,6.960) (0.025,6.961) (-0.050,6.960) (-0.100,6.959) (-0.150,6.958) (-0.200,6.955) (-0.250,6.952) (-0.300,6.948) (-0.350,6.943) (-0.400,6.938) (-0.450,6.932) (-0.500,6.925) (-0.550,6.917) (-0.600,6.909) (-0.650,6.899) (-0.700,6.890) (-0.750,6.879) (-0.800,6.867) (-0.850,6.855) (-0.900,6.842) (-0.950,6.828) (-1.000,6.814) (-1.050,6.798) (-1.100,6.782) (-1.150,6.765) (-1.200,6.747) (-1.250,6.728) (-1.300,6.708) (-1.350,6.687) (-1.400,6.665) (-1.450,6.642) (-1.500,6.619) (-1.550,6.594) (-1.600,6.568) (-1.650,6.541) (-1.700,6.513) (-1.750,6.484) (-1.800,6.454) (-1.845,6.425) (-1.900,6.389) (-1.950,6.355) (-1.992,6.325) (-2.050,6.282) (-2.092,6.250) (-2.150,6.204) (-2.200,6.162) (-2.242,6.125) (-2.297,6.075) (-2.350,6.025) (-2.400,5.975) (-2.448,5.925) (-2.494,5.875) (-2.538,5.825) (-2.575,5.782) (-2.621,5.725) (-2.660,5.675) (-2.697,5.625) (-2.732,5.575) (-2.766,5.525) (-2.799,5.475) (-2.830,5.425) (-2.859,5.375) (-2.888,5.325) (-2.915,5.275) (-2.941,5.225) (-2.966,5.175) (-2.989,5.125) (-3.012,5.075) (-3.033,5.025) (-3.053,4.975) (-3.072,4.925) (-3.090,4.875) (-3.107,4.825) (-3.123,4.775) (-3.138,4.725) (-3.151,4.675) (-3.164,4.625) (-3.176,4.575) (-3.187,4.525) (-3.197,4.475) (-3.206,4.425) (-3.214,4.375) (-3.221,4.325) (-3.227,4.275) (-3.232,4.225) (-3.236,4.175) (-3.239,4.125) (-3.242,4.075) (-3.243,4.025) (-3.244,3.975) (-3.243,3.925) (-3.242,3.875) (-3.240,3.825) (-3.237,3.775) (-3.232,3.725) (-3.227,3.675) (-3.221,3.625) (-3.214,3.575) (-3.207,3.525) (-3.198,3.475) (-3.188,3.425) (-3.177,3.375) (-3.165,3.325) (-3.153,3.275) (-3.139,3.225) (-3.124,3.175) (-3.109,3.125) (-3.092,3.075) (-3.075,3.027) (-3.056,2.975) (-3.036,2.925) (-3.015,2.875) (-2.993,2.825) (-2.970,2.775) (-2.946,2.725) (-2.921,2.675) (-2.895,2.625) (-2.867,2.575) (-2.839,2.525) (-2.809,2.475) (-2.778,2.425) (-2.750,2.381) (-2.713,2.325) (-2.678,2.275) (-2.643,2.225) (-2.606,2.175) (-2.575,2.135) (-2.528,2.075) (-2.488,2.025) (-2.450,1.980) (-2.402,1.925) (-2.358,1.875) (-2.312,1.825) (-2.275,1.785) (-2.225,1.733) (-2.175,1.683) (-2.125,1.633) (-2.075,1.585) (-2.038,1.550) (-2.000,1.515) (-1.957,1.475) (-1.901,1.425) (-1.850,1.381) (-1.800,1.337) (-1.756,1.300) (-1.700,1.253) (-1.650,1.212) (-1.605,1.175) (-1.550,1.131) (-1.500,1.091) (-1.450,1.052) (-1.400,1.013) (-1.351,0.975) (-1.300,0.936) (-1.252,0.900) (-1.200,0.861) (-1.152,0.825) (-1.100,0.787) (-1.051,0.750) (-1.000,0.713) (-0.950,0.677) (-0.900,0.640) (-0.850,0.604) (-0.800,0.568) (-0.750,0.532) (-0.705,0.500) (-0.650,0.461) (-0.600,0.425) (-0.550,0.390) (-0.500,0.354) (-0.450,0.319) (-0.400,0.284) (-0.354,0.250) (-0.300,0.213) (-0.250,0.177) (-0.220,0.155)}
\EQ{(-0.220,-0.155) (-0.275,-0.195) (-0.319,-0.225) (-0.375,-0.266) (-0.424,-0.300) (-0.475,-0.337) (-0.525,-0.372) (-0.575,-0.408) (-0.625,-0.443) (-0.670,-0.475) (-0.725,-0.515) (-0.774,-0.550) (-0.825,-0.587) (-0.875,-0.623) (-0.925,-0.659) (-0.975,-0.696) (-1.025,-0.733) (-1.075,-0.770) (-1.116,-0.800) (-1.175,-0.844) (-1.216,-0.875) (-1.275,-0.920) (-1.325,-0.959) (-1.375,-0.998) (-1.425,-1.037) (-1.472,-1.075) (-1.525,-1.118) (-1.565,-1.150) (-1.625,-1.200) (-1.675,-1.242) (-1.713,-1.275) (-1.771,-1.325) (-1.825,-1.374) (-1.875,-1.419) (-1.925,-1.466) (-1.962,-1.500) (-2.000,-1.537) (-2.038,-1.575) (-2.075,-1.612) (-2.112,-1.650) (-2.150,-1.690) (-2.200,-1.744) (-2.250,-1.799) (-2.294,-1.850) (-2.336,-1.900) (-2.375,-1.948) (-2.417,-2.000) (-2.450,-2.044) (-2.491,-2.100) (-2.525,-2.147) (-2.561,-2.200) (-2.594,-2.250) (-2.625,-2.299) (-2.656,-2.350) (-2.685,-2.400) (-2.713,-2.450) (-2.740,-2.500) (-2.765,-2.550) (-2.790,-2.600) (-2.813,-2.650) (-2.835,-2.700) (-2.856,-2.750) (-2.875,-2.799) (-2.894,-2.850) (-2.912,-2.900) (-2.928,-2.950) (-2.944,-3.000) (-2.958,-3.050) (-2.972,-3.100) (-2.984,-3.150) (-2.995,-3.200) (-3.005,-3.250) (-3.015,-3.300) (-3.023,-3.350) (-3.030,-3.400) (-3.036,-3.450) (-3.041,-3.500) (-3.045,-3.550) (-3.049,-3.600) (-3.051,-3.650) (-3.052,-3.700) (-3.052,-3.750) (-3.052,-3.800) (-3.050,-3.850) (-3.047,-3.900) (-3.044,-3.950) (-3.039,-4.000) (-3.033,-4.050) (-3.027,-4.100) (-3.019,-4.150) (-3.010,-4.200) (-3.001,-4.250) (-2.990,-4.300) (-2.978,-4.350) (-2.965,-4.400) (-2.951,-4.450) (-2.937,-4.500) (-2.921,-4.550) (-2.903,-4.600) (-2.885,-4.650) (-2.866,-4.700) (-2.845,-4.750) (-2.825,-4.796) (-2.800,-4.850) (-2.776,-4.900) (-2.751,-4.950) (-2.725,-4.998) (-2.700,-5.042) (-2.666,-5.100) (-2.635,-5.150) (-2.602,-5.200) (-2.568,-5.250) (-2.533,-5.300) (-2.500,-5.344) (-2.456,-5.400) (-2.415,-5.450) (-2.375,-5.497) (-2.328,-5.550) (-2.281,-5.600) (-2.231,-5.650) (-2.180,-5.700) (-2.125,-5.750) (-2.075,-5.794) (-2.025,-5.836) (-1.977,-5.875) (-1.925,-5.915) (-1.878,-5.950) (-1.825,-5.987) (-1.775,-6.021) (-1.731,-6.050) (-1.675,-6.085) (-1.625,-6.115) (-1.575,-6.143) (-1.525,-6.170) (-1.475,-6.197) (-1.425,-6.222) (-1.375,-6.246) (-1.325,-6.269) (-1.275,-6.291) (-1.225,-6.311) (-1.175,-6.331) (-1.126,-6.350) (-1.075,-6.368) (-1.025,-6.385) (-0.975,-6.401) (-0.925,-6.417) (-0.875,-6.431) (-0.825,-6.445) (-0.775,-6.457) (-0.725,-6.469) (-0.675,-6.480) (-0.625,-6.490) (-0.575,-6.499) (-0.525,-6.508) (-0.475,-6.516) (-0.425,-6.523) (-0.375,-6.529) (-0.325,-6.534) (-0.275,-6.539) (-0.225,-6.543) (-0.175,-6.546) (-0.125,-6.548) (-0.075,-6.549) (-0.025,-6.550) (0.047,-6.550) (0.100,-6.549) (0.150,-6.547) (0.200,-6.544) (0.250,-6.541) (0.300,-6.537) (0.350,-6.532) (0.400,-6.526) (0.450,-6.519) (0.500,-6.512) (0.550,-6.504) (0.600,-6.495) (0.650,-6.485) (0.700,-6.475) (0.750,-6.463) (0.800,-6.451) (0.850,-6.438) (0.900,-6.424) (0.950,-6.409) (1.000,-6.393) (1.050,-6.377) (1.100,-6.359) (1.150,-6.341) (1.200,-6.322) (1.250,-6.301) (1.300,-6.280) (1.350,-6.257) (1.400,-6.234) (1.450,-6.209) (1.500,-6.184) (1.550,-6.157) (1.600,-6.129) (1.650,-6.100) (1.700,-6.069) (1.750,-6.038) (1.800,-6.004) (1.842,-5.975) (1.900,-5.934) (1.945,-5.900) (2.000,-5.857) (2.050,-5.816) (2.097,-5.775) (2.150,-5.728) (2.200,-5.681) (2.250,-5.632) (2.300,-5.580) (2.350,-5.525) (2.394,-5.475) (2.436,-5.425) (2.475,-5.376) (2.514,-5.325) (2.550,-5.276) (2.586,-5.225) (2.619,-5.175) (2.650,-5.126) (2.681,-5.075) (2.710,-5.025) (2.737,-4.975) (2.763,-4.925) (2.788,-4.875) (2.812,-4.825) (2.834,-4.775) (2.856,-4.725) (2.875,-4.676) (2.894,-4.625) (2.912,-4.575) (2.929,-4.525) (2.944,-4.475) (2.959,-4.425) (2.972,-4.375) (2.984,-4.325) (2.995,-4.275) (3.006,-4.225) (3.015,-4.175) (3.023,-4.125) (3.030,-4.075) (3.036,-4.025) (3.041,-3.975) (3.046,-3.925) (3.049,-3.875) (3.051,-3.825) (3.052,-3.775) (3.052,-3.725) (3.052,-3.675) (3.050,-3.625) (3.047,-3.575) (3.044,-3.525) (3.039,-3.475) (3.033,-3.425) (3.026,-3.375) (3.019,-3.325) (3.010,-3.275) (3.000,-3.225) (2.990,-3.175) (2.978,-3.125) (2.965,-3.075) (2.951,-3.025) (2.936,-2.975) (2.920,-2.925) (2.903,-2.875) (2.885,-2.825) (2.866,-2.775) (2.845,-2.725) (2.825,-2.677) (2.801,-2.625) (2.778,-2.575) (2.753,-2.525) (2.726,-2.475) (2.700,-2.427) (2.675,-2.383) (2.641,-2.325) (2.610,-2.275) (2.578,-2.225) (2.544,-2.175) (2.509,-2.125) (2.475,-2.078) (2.436,-2.025) (2.400,-1.979) (2.357,-1.925) (2.315,-1.875) (2.275,-1.828) (2.228,-1.775) (2.183,-1.725) (2.136,-1.675) (2.100,-1.638) (2.063,-1.600) (2.025,-1.562) (1.987,-1.525) (1.950,-1.489) (1.909,-1.450) (1.854,-1.400) (1.800,-1.351) (1.750,-1.307) (1.700,-1.264) (1.655,-1.225) (1.600,-1.179) (1.550,-1.138) (1.504,-1.100) (1.450,-1.057) (1.410,-1.025) (1.350,-0.978) (1.300,-0.940) (1.250,-0.901) (1.200,-0.863) (1.150,-0.826) (1.100,-0.788) (1.050,-0.751) (1.000,-0.714) (0.950,-0.678) (0.900,-0.641) (0.850,-0.605) (0.809,-0.575) (0.750,-0.533) (0.705,-0.500) (0.650,-0.461) (0.600,-0.425) (0.550,-0.390) (0.500,-0.354) (0.450,-0.319) (0.400,-0.284) (0.354,-0.250) (0.300,-0.213) (0.250,-0.177) (0.206,-0.144)}
\EQ{(0.206,0.144) (0.248,0.175) (0.300,0.213) (0.350,0.248) (0.400,0.284) (0.450,0.319) (0.495,0.350) (0.550,0.390) (0.600,0.425) (0.650,0.461) (0.700,0.497) (0.750,0.533) (0.800,0.569) (0.844,0.600) (0.900,0.641) (0.947,0.675) (1.000,0.714) (1.049,0.750) (1.100,0.788) (1.149,0.825) (1.200,0.863) (1.249,0.900) (1.300,0.940) (1.346,0.975) (1.400,1.018) (1.441,1.050) (1.500,1.097) (1.550,1.138) (1.595,1.175) (1.650,1.221) (1.700,1.264) (1.742,1.300) (1.799,1.350) (1.850,1.396) (1.900,1.442) (1.950,1.489) (1.987,1.525) (2.025,1.562) (2.063,1.600) (2.100,1.638) (2.136,1.675) (2.175,1.717) (2.225,1.771) (2.273,1.825) (2.315,1.875) (2.350,1.917) (2.397,1.975) (2.436,2.025) (2.473,2.075) (2.509,2.125) (2.544,2.175) (2.575,2.221) (2.610,2.275) (2.641,2.325) (2.671,2.375) (2.699,2.425) (2.725,2.472) (2.750,2.520) (2.775,2.570) (2.800,2.622) (2.824,2.675) (2.845,2.725) (2.866,2.775) (2.885,2.825) (2.903,2.875) (2.920,2.925) (2.936,2.975) (2.951,3.025) (2.965,3.075) (2.978,3.125) (2.990,3.175) (3.000,3.225) (3.010,3.275) (3.019,3.325) (3.026,3.375) (3.033,3.425) (3.039,3.475) (3.044,3.525) (3.047,3.575) (3.050,3.625) (3.052,3.675) (3.052,3.725) (3.052,3.775) (3.051,3.825) (3.049,3.875) (3.046,3.925) (3.041,3.975) (3.036,4.025) (3.030,4.075) (3.023,4.125) (3.015,4.175) (3.006,4.225) (2.995,4.275) (2.984,4.325) (2.972,4.375) (2.959,4.425) (2.944,4.475) (2.929,4.525) (2.912,4.575) (2.894,4.625) (2.876,4.675) (2.856,4.725) (2.834,4.775) (2.812,4.825) (2.788,4.875) (2.763,4.925) (2.737,4.975) (2.710,5.025) (2.681,5.075) (2.651,5.125) (2.619,5.175) (2.586,5.225) (2.551,5.275) (2.514,5.325) (2.476,5.375) (2.436,5.425) (2.400,5.468) (2.350,5.525) (2.305,5.575) (2.256,5.625) (2.206,5.675) (2.153,5.725) (2.100,5.773) (2.050,5.816) (2.008,5.850) (1.950,5.896) (1.900,5.934) (1.850,5.970) (1.806,6.000) (1.750,6.038) (1.700,6.069) (1.650,6.100) (1.600,6.129) (1.550,6.157) (1.500,6.184) (1.450,6.209) (1.400,6.234) (1.350,6.257) (1.300,6.280) (1.253,6.300) (1.200,6.322) (1.150,6.341) (1.100,6.359) (1.050,6.377) (1.000,6.393) (0.950,6.409) (0.900,6.424) (0.850,6.438) (0.800,6.451) (0.750,6.463) (0.700,6.475) (0.650,6.485) (0.600,6.495) (0.550,6.504) (0.500,6.512) (0.450,6.519) (0.400,6.526) (0.350,6.532) (0.300,6.537) (0.250,6.541) (0.200,6.544) (0.150,6.547) (0.100,6.549) (0.050,6.550) (0.000,6.550) (-0.050,6.550) (-0.100,6.549) (-0.150,6.547) (-0.200,6.544) (-0.250,6.541) (-0.300,6.537) (-0.350,6.532) (-0.400,6.526) (-0.450,6.519) (-0.500,6.512) (-0.550,6.504) (-0.600,6.495) (-0.650,6.485) (-0.700,6.475) (-0.750,6.463) (-0.800,6.451) (-0.850,6.438) (-0.900,6.424) (-0.950,6.409) (-1.000,6.393) (-1.050,6.377) (-1.100,6.359) (-1.150,6.341) (-1.200,6.322) (-1.250,6.301) (-1.300,6.280) (-1.350,6.257) (-1.400,6.234) (-1.450,6.209) (-1.500,6.184) (-1.550,6.157) (-1.600,6.129) (-1.650,6.100) (-1.700,6.069) (-1.750,6.038) (-1.800,6.004) (-1.842,5.975) (-1.900,5.934) (-1.945,5.900) (-2.000,5.857) (-2.050,5.816) (-2.097,5.775) (-2.150,5.728) (-2.200,5.681) (-2.250,5.632) (-2.300,5.580) (-2.350,5.525) (-2.394,5.475) (-2.436,5.425) (-2.475,5.376) (-2.514,5.325) (-2.550,5.276) (-2.586,5.225) (-2.619,5.175) (-2.650,5.126) (-2.681,5.075) (-2.710,5.025) (-2.737,4.975) (-2.763,4.925) (-2.788,4.875) (-2.812,4.825) (-2.834,4.775) (-2.856,4.725) (-2.875,4.676) (-2.894,4.625) (-2.912,4.575) (-2.929,4.525) (-2.944,4.475) (-2.959,4.425) (-2.972,4.375) (-2.984,4.325) (-2.995,4.275) (-3.006,4.225) (-3.015,4.175) (-3.023,4.125) (-3.030,4.075) (-3.036,4.025) (-3.041,3.975) (-3.046,3.925) (-3.049,3.875) (-3.051,3.825) (-3.052,3.775) (-3.052,3.725) (-3.052,3.675) (-3.050,3.625) (-3.047,3.575) (-3.044,3.525) (-3.039,3.475) (-3.033,3.425) (-3.026,3.375) (-3.019,3.325) (-3.010,3.275) (-3.000,3.225) (-2.990,3.175) (-2.978,3.125) (-2.965,3.075) (-2.951,3.025) (-2.936,2.975) (-2.920,2.925) (-2.903,2.875) (-2.885,2.825) (-2.866,2.775) (-2.845,2.725) (-2.825,2.677) (-2.801,2.625) (-2.778,2.575) (-2.753,2.525) (-2.726,2.475) (-2.700,2.427) (-2.675,2.383) (-2.641,2.325) (-2.610,2.275) (-2.578,2.225) (-2.544,2.175) (-2.509,2.125) (-2.475,2.078) (-2.436,2.025) (-2.400,1.979) (-2.357,1.925) (-2.315,1.875) (-2.275,1.828) (-2.228,1.775) (-2.183,1.725) (-2.136,1.675) (-2.100,1.638) (-2.063,1.600) (-2.025,1.562) (-1.987,1.525) (-1.950,1.489) (-1.909,1.450) (-1.854,1.400) (-1.800,1.351) (-1.750,1.307) (-1.700,1.264) (-1.655,1.225) (-1.600,1.179) (-1.550,1.138) (-1.504,1.100) (-1.450,1.057) (-1.410,1.025) (-1.350,0.978) (-1.300,0.940) (-1.250,0.901) (-1.200,0.863) (-1.150,0.826) (-1.100,0.788) (-1.050,0.751) (-1.000,0.714) (-0.950,0.678) (-0.900,0.641) (-0.850,0.605) (-0.809,0.575) (-0.750,0.533) (-0.705,0.500) (-0.650,0.461) (-0.600,0.425) (-0.550,0.390) (-0.500,0.354) (-0.450,0.319) (-0.400,0.284) (-0.354,0.250) (-0.300,0.213) (-0.250,0.177) (-0.220,0.155)}
\EQ{(-0.220,-0.155) (-0.275,-0.195) (-0.319,-0.225) (-0.375,-0.266) (-0.424,-0.300) (-0.475,-0.337) (-0.525,-0.372) (-0.575,-0.408) (-0.625,-0.444) (-0.670,-0.475) (-0.725,-0.515) (-0.774,-0.550) (-0.825,-0.587) (-0.875,-0.624) (-0.925,-0.660) (-0.975,-0.697) (-1.025,-0.734) (-1.075,-0.771) (-1.125,-0.809) (-1.175,-0.847) (-1.225,-0.885) (-1.275,-0.924) (-1.325,-0.963) (-1.372,-1.000) (-1.425,-1.043) (-1.465,-1.075) (-1.525,-1.124) (-1.575,-1.166) (-1.615,-1.200) (-1.672,-1.250) (-1.725,-1.297) (-1.775,-1.342) (-1.825,-1.389) (-1.863,-1.425) (-1.900,-1.461) (-1.940,-1.500) (-1.990,-1.550) (-2.025,-1.587) (-2.061,-1.625) (-2.100,-1.668) (-2.150,-1.724) (-2.194,-1.775) (-2.236,-1.825) (-2.275,-1.874) (-2.315,-1.925) (-2.350,-1.972) (-2.388,-2.025) (-2.423,-2.075) (-2.456,-2.125) (-2.488,-2.175) (-2.519,-2.225) (-2.548,-2.275) (-2.575,-2.324) (-2.600,-2.370) (-2.625,-2.419) (-2.650,-2.470) (-2.675,-2.524) (-2.697,-2.575) (-2.718,-2.625) (-2.737,-2.675) (-2.756,-2.725) (-2.773,-2.775) (-2.789,-2.825) (-2.804,-2.875) (-2.818,-2.925) (-2.830,-2.975) (-2.842,-3.025) (-2.852,-3.075) (-2.862,-3.125) (-2.870,-3.175) (-2.877,-3.225) (-2.884,-3.275) (-2.889,-3.325) (-2.893,-3.375) (-2.896,-3.425) (-2.898,-3.475) (-2.899,-3.525) (-2.899,-3.575) (-2.898,-3.625) (-2.896,-3.675) (-2.893,-3.725) (-2.889,-3.775) (-2.884,-3.825) (-2.878,-3.875) (-2.871,-3.925) (-2.862,-3.975) (-2.853,-4.025) (-2.842,-4.075) (-2.831,-4.125) (-2.818,-4.175) (-2.804,-4.225) (-2.789,-4.275) (-2.773,-4.325) (-2.756,-4.375) (-2.738,-4.425) (-2.718,-4.475) (-2.697,-4.525) (-2.675,-4.575) (-2.652,-4.625) (-2.627,-4.675) (-2.601,-4.725) (-2.575,-4.772) (-2.550,-4.815) (-2.514,-4.875) (-2.482,-4.925) (-2.450,-4.972) (-2.413,-5.025) (-2.376,-5.075) (-2.337,-5.125) (-2.300,-5.170) (-2.253,-5.225) (-2.208,-5.275) (-2.161,-5.325) (-2.125,-5.362) (-2.086,-5.400) (-2.032,-5.450) (-1.976,-5.500) (-1.925,-5.543) (-1.885,-5.575) (-1.825,-5.621) (-1.775,-5.658) (-1.725,-5.693) (-1.678,-5.725) (-1.625,-5.759) (-1.575,-5.790) (-1.525,-5.819) (-1.475,-5.847) (-1.425,-5.874) (-1.375,-5.900) (-1.325,-5.924) (-1.275,-5.947) (-1.225,-5.970) (-1.175,-5.991) (-1.125,-6.011) (-1.075,-6.030) (-1.025,-6.048) (-0.975,-6.065) (-0.925,-6.081) (-0.875,-6.096) (-0.825,-6.111) (-0.775,-6.124) (-0.725,-6.137) (-0.675,-6.148) (-0.625,-6.159) (-0.575,-6.169) (-0.525,-6.178) (-0.475,-6.186) (-0.425,-6.193) (-0.375,-6.200) (-0.325,-6.205) (-0.275,-6.210) (-0.225,-6.214) (-0.175,-6.217) (-0.125,-6.220) (-0.075,-6.221) (-0.025,-6.222) (0.050,-6.222) (0.100,-6.221) (0.150,-6.219) (0.200,-6.216) (0.250,-6.212) (0.300,-6.208) (0.350,-6.203) (0.400,-6.196) (0.450,-6.190) (0.500,-6.182) (0.550,-6.173) (0.600,-6.164) (0.650,-6.154) (0.700,-6.142) (0.750,-6.130) (0.800,-6.117) (0.850,-6.104) (0.900,-6.089) (0.950,-6.073) (1.000,-6.057) (1.050,-6.039) (1.100,-6.020) (1.150,-6.001) (1.200,-5.980) (1.250,-5.959) (1.300,-5.936) (1.350,-5.912) (1.400,-5.887) (1.450,-5.861) (1.500,-5.833) (1.550,-5.805) (1.599,-5.775) (1.650,-5.743) (1.700,-5.710) (1.750,-5.676) (1.800,-5.640) (1.850,-5.602) (1.900,-5.563) (1.946,-5.525) (2.000,-5.479) (2.050,-5.434) (2.100,-5.386) (2.137,-5.350) (2.175,-5.311) (2.225,-5.257) (2.275,-5.200) (2.317,-5.150) (2.350,-5.108) (2.394,-5.050) (2.425,-5.008) (2.465,-4.950) (2.498,-4.900) (2.525,-4.857) (2.559,-4.800) (2.587,-4.750) (2.614,-4.700) (2.639,-4.650) (2.663,-4.600) (2.686,-4.550) (2.708,-4.500) (2.728,-4.450) (2.747,-4.400) (2.765,-4.350) (2.781,-4.300) (2.797,-4.250) (2.811,-4.200) (2.825,-4.150) (2.837,-4.100) (2.848,-4.050) (2.858,-4.000) (2.867,-3.950) (2.874,-3.900) (2.881,-3.850) (2.887,-3.800) (2.891,-3.750) (2.895,-3.700) (2.898,-3.650) (2.899,-3.600) (2.900,-3.550) (2.899,-3.500) (2.897,-3.450) (2.895,-3.400) (2.891,-3.350) (2.887,-3.300) (2.881,-3.250) (2.874,-3.200) (2.866,-3.150) (2.857,-3.100) (2.847,-3.050) (2.836,-3.000) (2.824,-2.950) (2.811,-2.900) (2.796,-2.850) (2.781,-2.800) (2.764,-2.750) (2.747,-2.700) (2.728,-2.650) (2.708,-2.600) (2.687,-2.550) (2.664,-2.500) (2.640,-2.450) (2.616,-2.400) (2.589,-2.350) (2.562,-2.300) (2.533,-2.250) (2.503,-2.200) (2.475,-2.155) (2.440,-2.100) (2.406,-2.050) (2.375,-2.007) (2.334,-1.950) (2.300,-1.906) (2.256,-1.850) (2.215,-1.800) (2.175,-1.752) (2.129,-1.700) (2.084,-1.650) (2.050,-1.613) (2.014,-1.575) (1.975,-1.535) (1.925,-1.485) (1.889,-1.450) (1.850,-1.413) (1.811,-1.375) (1.756,-1.325) (1.700,-1.275) (1.650,-1.231) (1.600,-1.188) (1.556,-1.150) (1.500,-1.104) (1.450,-1.063) (1.404,-1.025) (1.350,-0.983) (1.309,-0.950) (1.250,-0.904) (1.200,-0.866) (1.150,-0.828) (1.100,-0.790) (1.050,-0.753) (1.000,-0.716) (0.950,-0.679) (0.900,-0.642) (0.850,-0.606) (0.809,-0.575) (0.750,-0.533) (0.705,-0.500) (0.650,-0.461) (0.600,-0.425) (0.550,-0.390) (0.500,-0.354) (0.450,-0.319) (0.400,-0.284) (0.354,-0.250) (0.300,-0.213) (0.250,-0.177) (0.206,-0.144)}
\EQ{(0.206,0.144) (0.248,0.175) (0.300,0.213) (0.350,0.248) (0.400,0.284) (0.450,0.319) (0.495,0.350) (0.550,0.390) (0.600,0.425) (0.650,0.461) (0.700,0.497) (0.750,0.533) (0.800,0.569) (0.843,0.600) (0.900,0.642) (0.946,0.675) (1.000,0.716) (1.047,0.750) (1.100,0.790) (1.147,0.825) (1.200,0.866) (1.245,0.900) (1.300,0.943) (1.341,0.975) (1.400,1.022) (1.450,1.063) (1.496,1.100) (1.550,1.145) (1.600,1.188) (1.644,1.225) (1.700,1.275) (1.750,1.320) (1.800,1.366) (1.837,1.400) (1.875,1.436) (1.915,1.475) (1.965,1.525) (2.000,1.561) (2.038,1.600) (2.075,1.640) (2.125,1.695) (2.173,1.750) (2.215,1.800) (2.250,1.843) (2.296,1.900) (2.334,1.950) (2.370,2.000) (2.400,2.042) (2.440,2.100) (2.472,2.150) (2.500,2.195) (2.533,2.250) (2.562,2.300) (2.589,2.350) (2.616,2.400) (2.640,2.450) (2.664,2.500) (2.687,2.550) (2.708,2.600) (2.728,2.650) (2.747,2.700) (2.764,2.750) (2.781,2.800) (2.796,2.850) (2.811,2.900) (2.824,2.950) (2.836,3.000) (2.847,3.050) (2.857,3.100) (2.866,3.150) (2.874,3.200) (2.881,3.250) (2.887,3.300) (2.891,3.350) (2.895,3.400) (2.897,3.450) (2.899,3.500) (2.900,3.550) (2.899,3.600) (2.898,3.650) (2.895,3.700) (2.891,3.750) (2.887,3.800) (2.881,3.850) (2.874,3.900) (2.867,3.950) (2.858,4.000) (2.848,4.050) (2.837,4.100) (2.825,4.150) (2.811,4.200) (2.797,4.250) (2.781,4.300) (2.765,4.350) (2.747,4.400) (2.728,4.450) (2.708,4.500) (2.686,4.550) (2.663,4.600) (2.639,4.650) (2.614,4.700) (2.587,4.750) (2.559,4.800) (2.529,4.850) (2.500,4.897) (2.465,4.950) (2.431,5.000) (2.400,5.043) (2.357,5.100) (2.325,5.140) (2.275,5.200) (2.231,5.250) (2.185,5.300) (2.150,5.337) (2.112,5.375) (2.075,5.410) (2.032,5.450) (1.976,5.500) (1.925,5.543) (1.885,5.575) (1.825,5.621) (1.775,5.658) (1.725,5.693) (1.678,5.725) (1.625,5.759) (1.575,5.790) (1.525,5.819) (1.475,5.847) (1.425,5.874) (1.375,5.900) (1.325,5.924) (1.275,5.947) (1.225,5.970) (1.175,5.991) (1.125,6.011) (1.075,6.030) (1.025,6.048) (0.975,6.065) (0.925,6.081) (0.875,6.096) (0.825,6.111) (0.775,6.124) (0.725,6.137) (0.675,6.148) (0.625,6.159) (0.575,6.169) (0.525,6.178) (0.475,6.186) (0.425,6.193) (0.375,6.200) (0.325,6.205) (0.275,6.210) (0.225,6.214) (0.175,6.217) (0.125,6.220) (0.075,6.221) (0.025,6.222) (-0.050,6.222) (-0.100,6.221) (-0.150,6.219) (-0.200,6.216) (-0.250,6.212) (-0.300,6.208) (-0.350,6.203) (-0.400,6.196) (-0.450,6.190) (-0.500,6.182) (-0.550,6.173) (-0.600,6.164) (-0.650,6.154) (-0.700,6.142) (-0.750,6.130) (-0.800,6.117) (-0.850,6.104) (-0.900,6.089) (-0.950,6.073) (-1.000,6.057) (-1.050,6.039) (-1.100,6.020) (-1.150,6.001) (-1.200,5.980) (-1.250,5.959) (-1.300,5.936) (-1.350,5.912) (-1.400,5.887) (-1.450,5.861) (-1.500,5.833) (-1.550,5.805) (-1.599,5.775) (-1.650,5.743) (-1.700,5.710) (-1.750,5.676) (-1.800,5.640) (-1.850,5.602) (-1.900,5.563) (-1.946,5.525) (-2.000,5.479) (-2.050,5.434) (-2.100,5.386) (-2.137,5.350) (-2.175,5.311) (-2.225,5.257) (-2.275,5.200) (-2.317,5.150) (-2.350,5.108) (-2.394,5.050) (-2.425,5.008) (-2.465,4.950) (-2.498,4.900) (-2.525,4.857) (-2.559,4.800) (-2.587,4.750) (-2.614,4.700) (-2.639,4.650) (-2.663,4.600) (-2.686,4.550) (-2.708,4.500) (-2.728,4.450) (-2.747,4.400) (-2.765,4.350) (-2.781,4.300) (-2.797,4.250) (-2.811,4.200) (-2.825,4.150) (-2.837,4.100) (-2.848,4.050) (-2.858,4.000) (-2.867,3.950) (-2.874,3.900) (-2.881,3.850) (-2.887,3.800) (-2.891,3.750) (-2.895,3.700) (-2.898,3.650) (-2.899,3.600) (-2.900,3.550) (-2.899,3.500) (-2.897,3.450) (-2.895,3.400) (-2.891,3.350) (-2.887,3.300) (-2.881,3.250) (-2.874,3.200) (-2.866,3.150) (-2.857,3.100) (-2.847,3.050) (-2.836,3.000) (-2.824,2.950) (-2.811,2.900) (-2.796,2.850) (-2.781,2.800) (-2.764,2.750) (-2.747,2.700) (-2.728,2.650) (-2.708,2.600) (-2.687,2.550) (-2.664,2.500) (-2.640,2.450) (-2.616,2.400) (-2.589,2.350) (-2.562,2.300) (-2.533,2.250) (-2.503,2.200) (-2.475,2.155) (-2.440,2.100) (-2.406,2.050) (-2.375,2.007) (-2.334,1.950) (-2.300,1.906) (-2.256,1.850) (-2.215,1.800) (-2.175,1.752) (-2.129,1.700) (-2.084,1.650) (-2.050,1.613) (-2.014,1.575) (-1.975,1.535) (-1.925,1.485) (-1.889,1.450) (-1.850,1.413) (-1.811,1.375) (-1.756,1.325) (-1.700,1.275) (-1.650,1.231) (-1.600,1.188) (-1.556,1.150) (-1.500,1.104) (-1.450,1.063) (-1.404,1.025) (-1.350,0.983) (-1.309,0.950) (-1.250,0.904) (-1.200,0.866) (-1.150,0.828) (-1.100,0.790) (-1.050,0.753) (-1.000,0.716) (-0.950,0.679) (-0.900,0.642) (-0.850,0.606) (-0.809,0.575) (-0.750,0.533) (-0.705,0.500) (-0.650,0.461) (-0.600,0.425) (-0.550,0.390) (-0.500,0.354) (-0.450,0.319) (-0.400,0.284) (-0.354,0.250) (-0.300,0.213) (-0.250,0.177) (-0.220,0.155)}
\EQ{(-0.220,-0.155) (-0.275,-0.195) (-0.319,-0.225) (-0.375,-0.266) (-0.424,-0.300) (-0.475,-0.337) (-0.525,-0.372) (-0.575,-0.408) (-0.625,-0.444) (-0.670,-0.475) (-0.725,-0.516) (-0.773,-0.550) (-0.825,-0.588) (-0.875,-0.624) (-0.925,-0.661) (-0.975,-0.698) (-1.025,-0.735) (-1.075,-0.773) (-1.125,-0.811) (-1.175,-0.849) (-1.225,-0.888) (-1.273,-0.925) (-1.325,-0.967) (-1.367,-1.000) (-1.425,-1.048) (-1.475,-1.089) (-1.518,-1.125) (-1.575,-1.174) (-1.625,-1.218) (-1.675,-1.263) (-1.715,-1.300) (-1.769,-1.350) (-1.821,-1.400) (-1.871,-1.450) (-1.920,-1.500) (-1.967,-1.550) (-2.013,-1.600) (-2.050,-1.641) (-2.100,-1.699) (-2.142,-1.750) (-2.175,-1.791) (-2.221,-1.850) (-2.258,-1.900) (-2.293,-1.950) (-2.325,-1.996) (-2.361,-2.050) (-2.392,-2.100) (-2.422,-2.150) (-2.450,-2.199) (-2.475,-2.245) (-2.500,-2.292) (-2.525,-2.343) (-2.550,-2.396) (-2.574,-2.450) (-2.595,-2.500) (-2.615,-2.550) (-2.633,-2.600) (-2.650,-2.648) (-2.667,-2.700) (-2.682,-2.750) (-2.696,-2.800) (-2.708,-2.850) (-2.720,-2.900) (-2.730,-2.950) (-2.739,-3.000) (-2.747,-3.050) (-2.754,-3.100) (-2.760,-3.150) (-2.765,-3.200) (-2.769,-3.250) (-2.771,-3.300) (-2.773,-3.350) (-2.773,-3.400) (-2.773,-3.450) (-2.771,-3.500) (-2.768,-3.550) (-2.764,-3.600) (-2.759,-3.650) (-2.753,-3.700) (-2.746,-3.750) (-2.738,-3.800) (-2.729,-3.850) (-2.718,-3.900) (-2.706,-3.950) (-2.694,-4.000) (-2.680,-4.050) (-2.664,-4.100) (-2.648,-4.150) (-2.630,-4.200) (-2.612,-4.250) (-2.592,-4.300) (-2.570,-4.350) (-2.547,-4.400) (-2.525,-4.446) (-2.500,-4.496) (-2.475,-4.542) (-2.450,-4.587) (-2.412,-4.650) (-2.381,-4.700) (-2.350,-4.746) (-2.312,-4.800) (-2.275,-4.850) (-2.237,-4.900) (-2.200,-4.945) (-2.153,-5.000) (-2.108,-5.050) (-2.061,-5.100) (-2.025,-5.136) (-1.984,-5.175) (-1.930,-5.225) (-1.875,-5.273) (-1.825,-5.315) (-1.781,-5.350) (-1.725,-5.392) (-1.679,-5.425) (-1.625,-5.462) (-1.575,-5.495) (-1.527,-5.525) (-1.475,-5.556) (-1.425,-5.585) (-1.375,-5.612) (-1.325,-5.638) (-1.275,-5.662) (-1.225,-5.686) (-1.175,-5.708) (-1.125,-5.729) (-1.075,-5.749) (-1.025,-5.769) (-0.975,-5.787) (-0.925,-5.804) (-0.875,-5.819) (-0.825,-5.834) (-0.775,-5.848) (-0.725,-5.862) (-0.675,-5.874) (-0.625,-5.885) (-0.575,-5.895) (-0.525,-5.905) (-0.475,-5.913) (-0.425,-5.921) (-0.375,-5.928) (-0.325,-5.934) (-0.275,-5.939) (-0.225,-5.943) (-0.175,-5.946) (-0.125,-5.949) (-0.075,-5.950) (-0.025,-5.951) (0.050,-5.951) (0.100,-5.950) (0.150,-5.948) (0.200,-5.945) (0.250,-5.941) (0.300,-5.936) (0.350,-5.931) (0.400,-5.924) (0.450,-5.917) (0.500,-5.909) (0.550,-5.900) (0.600,-5.890) (0.650,-5.879) (0.700,-5.868) (0.750,-5.855) (0.800,-5.842) (0.850,-5.827) (0.900,-5.812) (0.950,-5.795) (1.000,-5.778) (1.050,-5.759) (1.100,-5.740) (1.150,-5.719) (1.200,-5.697) (1.249,-5.675) (1.300,-5.650) (1.350,-5.625) (1.397,-5.600) (1.442,-5.575) (1.486,-5.550) (1.550,-5.511) (1.600,-5.479) (1.643,-5.450) (1.700,-5.410) (1.748,-5.375) (1.800,-5.335) (1.843,-5.300) (1.900,-5.252) (1.950,-5.207) (2.000,-5.160) (2.036,-5.125) (2.075,-5.085) (2.125,-5.031) (2.175,-4.975) (2.217,-4.925) (2.250,-4.883) (2.294,-4.825) (2.325,-4.782) (2.364,-4.725) (2.397,-4.675) (2.425,-4.629) (2.457,-4.575) (2.484,-4.525) (2.511,-4.475) (2.535,-4.425) (2.559,-4.375) (2.581,-4.325) (2.602,-4.275) (2.621,-4.225) (2.639,-4.175) (2.656,-4.125) (2.672,-4.075) (2.687,-4.025) (2.700,-3.975) (2.712,-3.925) (2.723,-3.875) (2.733,-3.825) (2.742,-3.775) (2.750,-3.725) (2.757,-3.675) (2.762,-3.625) (2.766,-3.575) (2.770,-3.525) (2.772,-3.475) (2.773,-3.425) (2.773,-3.375) (2.772,-3.325) (2.770,-3.275) (2.767,-3.225) (2.763,-3.175) (2.757,-3.125) (2.751,-3.075) (2.743,-3.025) (2.735,-2.975) (2.725,-2.926) (2.714,-2.875) (2.702,-2.825) (2.689,-2.775) (2.674,-2.725) (2.659,-2.675) (2.642,-2.625) (2.625,-2.577) (2.605,-2.525) (2.585,-2.475) (2.563,-2.425) (2.541,-2.375) (2.517,-2.325) (2.491,-2.275) (2.464,-2.225) (2.436,-2.175) (2.407,-2.125) (2.376,-2.075) (2.344,-2.025) (2.311,-1.975) (2.276,-1.925) (2.240,-1.875) (2.202,-1.825) (2.162,-1.775) (2.125,-1.729) (2.079,-1.675) (2.036,-1.625) (2.000,-1.586) (1.950,-1.532) (1.900,-1.480) (1.850,-1.429) (1.800,-1.380) (1.750,-1.333) (1.700,-1.286) (1.661,-1.250) (1.604,-1.200) (1.550,-1.153) (1.500,-1.110) (1.458,-1.075) (1.400,-1.027) (1.350,-0.987) (1.304,-0.950) (1.250,-0.907) (1.209,-0.875) (1.150,-0.830) (1.100,-0.792) (1.050,-0.754) (1.000,-0.717) (0.950,-0.680) (0.900,-0.643) (0.850,-0.606) (0.808,-0.575) (0.750,-0.534) (0.704,-0.500) (0.650,-0.462) (0.600,-0.425) (0.550,-0.390) (0.500,-0.354) (0.450,-0.319) (0.400,-0.284) (0.354,-0.250) (0.300,-0.213) (0.250,-0.177) (0.206,-0.144)}
\EQ{(0.206,0.144) (0.248,0.175) (0.300,0.213) (0.350,0.248) (0.400,0.284) (0.450,0.319) (0.495,0.350) (0.550,0.390) (0.599,0.425) (0.650,0.462) (0.700,0.497) (0.750,0.534) (0.800,0.570) (0.842,0.600) (0.900,0.643) (0.944,0.675) (1.000,0.717) (1.045,0.750) (1.100,0.792) (1.144,0.825) (1.200,0.868) (1.241,0.900) (1.300,0.947) (1.350,0.987) (1.398,1.025) (1.450,1.068) (1.500,1.110) (1.547,1.150) (1.600,1.196) (1.650,1.241) (1.688,1.275) (1.725,1.309) (1.769,1.350) (1.821,1.400) (1.871,1.450) (1.920,1.500) (1.967,1.550) (2.013,1.600) (2.050,1.641) (2.100,1.699) (2.142,1.750) (2.175,1.791) (2.221,1.850) (2.258,1.900) (2.293,1.950) (2.325,1.996) (2.361,2.050) (2.392,2.100) (2.422,2.150) (2.450,2.199) (2.475,2.245) (2.500,2.292) (2.525,2.343) (2.550,2.396) (2.574,2.450) (2.595,2.500) (2.615,2.550) (2.633,2.600) (2.650,2.648) (2.667,2.700) (2.682,2.750) (2.696,2.800) (2.708,2.850) (2.720,2.900) (2.730,2.950) (2.739,3.000) (2.747,3.050) (2.754,3.100) (2.760,3.150) (2.765,3.200) (2.769,3.250) (2.771,3.300) (2.773,3.350) (2.773,3.400) (2.773,3.450) (2.771,3.500) (2.768,3.550) (2.764,3.600) (2.759,3.650) (2.753,3.700) (2.746,3.750) (2.738,3.800) (2.729,3.850) (2.718,3.900) (2.706,3.950) (2.694,4.000) (2.680,4.050) (2.664,4.100) (2.648,4.150) (2.630,4.200) (2.612,4.250) (2.592,4.300) (2.570,4.350) (2.547,4.400) (2.525,4.446) (2.500,4.496) (2.475,4.542) (2.450,4.587) (2.412,4.650) (2.381,4.700) (2.350,4.746) (2.312,4.800) (2.275,4.850) (2.237,4.900) (2.200,4.945) (2.153,5.000) (2.108,5.050) (2.061,5.100) (2.025,5.136) (1.984,5.175) (1.930,5.225) (1.875,5.273) (1.825,5.315) (1.781,5.350) (1.725,5.392) (1.679,5.425) (1.625,5.462) (1.575,5.495) (1.527,5.525) (1.475,5.556) (1.425,5.585) (1.375,5.612) (1.325,5.638) (1.275,5.662) (1.225,5.686) (1.175,5.708) (1.125,5.729) (1.075,5.749) (1.025,5.769) (0.975,5.787) (0.925,5.804) (0.875,5.819) (0.825,5.834) (0.775,5.848) (0.725,5.862) (0.675,5.874) (0.625,5.885) (0.575,5.895) (0.525,5.905) (0.475,5.913) (0.425,5.921) (0.375,5.928) (0.325,5.934) (0.275,5.939) (0.225,5.943) (0.175,5.946) (0.125,5.949) (0.075,5.950) (0.025,5.951) (-0.050,5.951) (-0.100,5.950) (-0.150,5.948) (-0.200,5.945) (-0.250,5.941) (-0.300,5.936) (-0.350,5.931) (-0.400,5.924) (-0.450,5.917) (-0.500,5.909) (-0.550,5.900) (-0.600,5.890) (-0.650,5.879) (-0.700,5.868) (-0.750,5.855) (-0.800,5.842) (-0.850,5.827) (-0.900,5.812) (-0.950,5.795) (-1.000,5.778) (-1.050,5.759) (-1.100,5.740) (-1.150,5.719) (-1.200,5.697) (-1.249,5.675) (-1.300,5.650) (-1.350,5.625) (-1.397,5.600) (-1.442,5.575) (-1.486,5.550) (-1.550,5.511) (-1.600,5.479) (-1.643,5.450) (-1.700,5.410) (-1.748,5.375) (-1.800,5.335) (-1.843,5.300) (-1.900,5.252) (-1.950,5.207) (-2.000,5.160) (-2.036,5.125) (-2.075,5.085) (-2.125,5.031) (-2.175,4.975) (-2.217,4.925) (-2.250,4.883) (-2.294,4.825) (-2.325,4.782) (-2.364,4.725) (-2.397,4.675) (-2.425,4.629) (-2.457,4.575) (-2.484,4.525) (-2.511,4.475) (-2.535,4.425) (-2.559,4.375) (-2.581,4.325) (-2.602,4.275) (-2.621,4.225) (-2.639,4.175) (-2.656,4.125) (-2.672,4.075) (-2.687,4.025) (-2.700,3.975) (-2.712,3.925) (-2.723,3.875) (-2.733,3.825) (-2.742,3.775) (-2.750,3.725) (-2.757,3.675) (-2.762,3.625) (-2.766,3.575) (-2.770,3.525) (-2.772,3.475) (-2.773,3.425) (-2.773,3.375) (-2.772,3.325) (-2.770,3.275) (-2.767,3.225) (-2.763,3.175) (-2.757,3.125) (-2.751,3.075) (-2.743,3.025) (-2.735,2.975) (-2.725,2.926) (-2.714,2.875) (-2.702,2.825) (-2.689,2.775) (-2.674,2.725) (-2.659,2.675) (-2.642,2.625) (-2.625,2.577) (-2.605,2.525) (-2.585,2.475) (-2.563,2.425) (-2.541,2.375) (-2.517,2.325) (-2.491,2.275) (-2.464,2.225) (-2.436,2.175) (-2.407,2.125) (-2.376,2.075) (-2.344,2.025) (-2.311,1.975) (-2.276,1.925) (-2.240,1.875) (-2.202,1.825) (-2.162,1.775) (-2.125,1.729) (-2.079,1.675) (-2.036,1.625) (-2.000,1.586) (-1.950,1.532) (-1.900,1.480) (-1.850,1.429) (-1.800,1.380) (-1.750,1.333) (-1.700,1.286) (-1.661,1.250) (-1.604,1.200) (-1.550,1.153) (-1.500,1.110) (-1.458,1.075) (-1.400,1.027) (-1.350,0.987) (-1.304,0.950) (-1.250,0.907) (-1.209,0.875) (-1.150,0.830) (-1.100,0.792) (-1.050,0.754) (-1.000,0.717) (-0.950,0.680) (-0.900,0.643) (-0.850,0.606) (-0.808,0.575) (-0.750,0.534) (-0.704,0.500) (-0.650,0.462) (-0.600,0.425) (-0.550,0.390) (-0.500,0.354) (-0.450,0.319) (-0.400,0.284) (-0.354,0.250) (-0.300,0.213) (-0.250,0.177) (-0.220,0.155)}
\EQ{(-0.220,-0.155) (-0.275,-0.195) (-0.319,-0.225) (-0.375,-0.266) (-0.424,-0.300) (-0.475,-0.337) (-0.525,-0.372) (-0.575,-0.408) (-0.625,-0.444) (-0.669,-0.475) (-0.725,-0.516) (-0.773,-0.550) (-0.825,-0.589) (-0.875,-0.625) (-0.925,-0.662) (-0.975,-0.699) (-1.025,-0.737) (-1.075,-0.774) (-1.125,-0.813) (-1.173,-0.850) (-1.225,-0.891) (-1.269,-0.925) (-1.325,-0.970) (-1.375,-1.011) (-1.422,-1.050) (-1.475,-1.095) (-1.525,-1.139) (-1.567,-1.175) (-1.622,-1.225) (-1.675,-1.274) (-1.725,-1.322) (-1.775,-1.370) (-1.825,-1.421) (-1.875,-1.473) (-1.924,-1.525) (-1.969,-1.575) (-2.012,-1.625) (-2.050,-1.671) (-2.094,-1.725) (-2.125,-1.766) (-2.170,-1.825) (-2.200,-1.868) (-2.239,-1.925) (-2.272,-1.975) (-2.300,-2.020) (-2.333,-2.075) (-2.362,-2.125) (-2.389,-2.175) (-2.414,-2.225) (-2.439,-2.275) (-2.462,-2.325) (-2.483,-2.375) (-2.504,-2.425) (-2.523,-2.475) (-2.541,-2.525) (-2.558,-2.575) (-2.573,-2.625) (-2.587,-2.675) (-2.600,-2.725) (-2.612,-2.775) (-2.623,-2.825) (-2.632,-2.875) (-2.640,-2.925) (-2.647,-2.975) (-2.654,-3.025) (-2.658,-3.075) (-2.662,-3.125) (-2.665,-3.175) (-2.666,-3.225) (-2.667,-3.275) (-2.666,-3.325) (-2.664,-3.375) (-2.661,-3.425) (-2.657,-3.475) (-2.651,-3.525) (-2.645,-3.575) (-2.637,-3.625) (-2.629,-3.675) (-2.619,-3.725) (-2.608,-3.775) (-2.595,-3.825) (-2.582,-3.875) (-2.567,-3.925) (-2.551,-3.975) (-2.534,-4.025) (-2.515,-4.075) (-2.496,-4.125) (-2.475,-4.174) (-2.452,-4.225) (-2.428,-4.275) (-2.402,-4.325) (-2.375,-4.375) (-2.350,-4.420) (-2.317,-4.475) (-2.285,-4.525) (-2.251,-4.575) (-2.216,-4.625) (-2.179,-4.675) (-2.140,-4.725) (-2.100,-4.773) (-2.055,-4.825) (-2.009,-4.875) (-1.960,-4.925) (-1.909,-4.975) (-1.855,-5.025) (-1.800,-5.073) (-1.750,-5.114) (-1.704,-5.150) (-1.650,-5.191) (-1.602,-5.225) (-1.550,-5.260) (-1.500,-5.293) (-1.450,-5.323) (-1.405,-5.350) (-1.361,-5.375) (-1.314,-5.400) (-1.265,-5.425) (-1.214,-5.450) (-1.160,-5.475) (-1.103,-5.500) (-1.050,-5.522) (-1.000,-5.541) (-0.950,-5.559) (-0.900,-5.577) (-0.850,-5.593) (-0.800,-5.608) (-0.750,-5.622) (-0.700,-5.635) (-0.650,-5.647) (-0.600,-5.659) (-0.550,-5.669) (-0.500,-5.678) (-0.450,-5.687) (-0.400,-5.694) (-0.350,-5.701) (-0.300,-5.707) (-0.250,-5.711) (-0.200,-5.715) (-0.150,-5.718) (-0.100,-5.721) (-0.050,-5.722) (0.000,-5.722) (0.050,-5.722) (0.100,-5.721) (0.150,-5.718) (0.200,-5.715) (0.250,-5.711) (0.300,-5.707) (0.350,-5.701) (0.400,-5.694) (0.450,-5.687) (0.500,-5.678) (0.550,-5.669) (0.600,-5.659) (0.650,-5.647) (0.700,-5.635) (0.750,-5.622) (0.800,-5.608) (0.850,-5.593) (0.900,-5.577) (0.950,-5.559) (1.000,-5.541) (1.050,-5.522) (1.100,-5.501) (1.150,-5.480) (1.200,-5.457) (1.250,-5.433) (1.300,-5.407) (1.350,-5.381) (1.400,-5.353) (1.447,-5.325) (1.500,-5.293) (1.550,-5.260) (1.600,-5.226) (1.650,-5.191) (1.700,-5.153) (1.750,-5.114) (1.797,-5.075) (1.850,-5.029) (1.900,-4.983) (1.950,-4.935) (2.000,-4.884) (2.050,-4.830) (2.098,-4.775) (2.140,-4.725) (2.175,-4.680) (2.216,-4.625) (2.250,-4.577) (2.285,-4.525) (2.317,-4.475) (2.347,-4.425) (2.375,-4.376) (2.400,-4.330) (2.425,-4.281) (2.450,-4.229) (2.474,-4.175) (2.496,-4.125) (2.515,-4.075) (2.534,-4.025) (2.551,-3.975) (2.567,-3.925) (2.582,-3.875) (2.595,-3.825) (2.608,-3.775) (2.619,-3.725) (2.629,-3.675) (2.637,-3.625) (2.645,-3.575) (2.651,-3.525) (2.657,-3.475) (2.661,-3.425) (2.664,-3.375) (2.666,-3.325) (2.667,-3.275) (2.666,-3.225) (2.665,-3.175) (2.662,-3.125) (2.658,-3.075) (2.654,-3.025) (2.647,-2.975) (2.640,-2.925) (2.632,-2.875) (2.623,-2.825) (2.612,-2.775) (2.600,-2.725) (2.587,-2.675) (2.573,-2.625) (2.558,-2.575) (2.541,-2.525) (2.523,-2.475) (2.504,-2.425) (2.483,-2.375) (2.462,-2.325) (2.439,-2.275) (2.414,-2.225) (2.389,-2.175) (2.362,-2.125) (2.333,-2.075) (2.303,-2.025) (2.275,-1.980) (2.239,-1.925) (2.205,-1.875) (2.175,-1.833) (2.132,-1.775) (2.100,-1.733) (2.053,-1.675) (2.012,-1.625) (1.975,-1.583) (1.925,-1.527) (1.877,-1.475) (1.829,-1.425) (1.780,-1.375) (1.729,-1.325) (1.676,-1.275) (1.625,-1.228) (1.575,-1.183) (1.525,-1.139) (1.481,-1.100) (1.425,-1.053) (1.375,-1.011) (1.331,-0.975) (1.275,-0.930) (1.225,-0.891) (1.175,-0.851) (1.125,-0.813) (1.076,-0.775) (1.025,-0.737) (0.976,-0.700) (0.925,-0.662) (0.875,-0.625) (0.825,-0.589) (0.775,-0.552) (0.725,-0.516) (0.675,-0.480) (0.635,-0.450) (0.575,-0.408) (0.530,-0.375) (0.475,-0.337) (0.425,-0.301) (0.375,-0.266) (0.325,-0.230) (0.284,-0.200) (0.225,-0.160) (0.206,-0.144)}
\EQ{(0.206,0.144) (0.248,0.175) (0.300,0.213) (0.350,0.248) (0.400,0.284) (0.450,0.319) (0.495,0.350) (0.550,0.390) (0.599,0.425) (0.650,0.462) (0.700,0.498) (0.750,0.534) (0.800,0.570) (0.842,0.600) (0.900,0.644) (0.943,0.675) (1.000,0.718) (1.043,0.750) (1.100,0.794) (1.141,0.825) (1.200,0.871) (1.250,0.910) (1.300,0.950) (1.350,0.991) (1.392,1.025) (1.450,1.074) (1.500,1.117) (1.538,1.150) (1.595,1.200) (1.649,1.250) (1.700,1.298) (1.750,1.346) (1.800,1.395) (1.850,1.447) (1.900,1.499) (1.946,1.550) (1.990,1.600) (2.025,1.641) (2.074,1.700) (2.113,1.750) (2.150,1.799) (2.188,1.850) (2.222,1.900) (2.256,1.950) (2.288,2.000) (2.318,2.050) (2.348,2.100) (2.375,2.149) (2.400,2.197) (2.425,2.247) (2.450,2.299) (2.473,2.350) (2.494,2.400) (2.514,2.450) (2.532,2.500) (2.549,2.550) (2.565,2.600) (2.580,2.650) (2.594,2.700) (2.606,2.750) (2.617,2.800) (2.627,2.850) (2.636,2.900) (2.644,2.950) (2.651,3.000) (2.656,3.050) (2.660,3.100) (2.664,3.150) (2.666,3.200) (2.667,3.250) (2.666,3.300) (2.665,3.350) (2.663,3.400) (2.659,3.450) (2.654,3.500) (2.648,3.550) (2.641,3.600) (2.633,3.650) (2.624,3.700) (2.613,3.750) (2.602,3.800) (2.589,3.850) (2.575,3.899) (2.559,3.950) (2.543,4.000) (2.525,4.050) (2.506,4.100) (2.485,4.150) (2.463,4.200) (2.440,4.250) (2.415,4.300) (2.389,4.350) (2.361,4.400) (2.332,4.450) (2.301,4.500) (2.268,4.550) (2.234,4.600) (2.200,4.647) (2.160,4.700) (2.125,4.743) (2.077,4.800) (2.032,4.850) (1.985,4.900) (1.935,4.950) (1.882,5.000) (1.826,5.050) (1.775,5.094) (1.725,5.134) (1.675,5.172) (1.625,5.209) (1.575,5.244) (1.528,5.275) (1.475,5.308) (1.425,5.338) (1.375,5.367) (1.325,5.394) (1.275,5.420) (1.225,5.445) (1.175,5.468) (1.125,5.491) (1.075,5.512) (1.025,5.532) (0.976,5.550) (0.925,5.568) (0.875,5.585) (0.826,5.600) (0.775,5.615) (0.725,5.629) (0.675,5.641) (0.625,5.653) (0.575,5.664) (0.525,5.674) (0.475,5.683) (0.425,5.691) (0.375,5.698) (0.325,5.704) (0.275,5.709) (0.225,5.713) (0.175,5.717) (0.125,5.720) (0.075,5.721) (0.025,5.722) (-0.050,5.722) (-0.100,5.721) (-0.150,5.718) (-0.200,5.715) (-0.250,5.711) (-0.300,5.707) (-0.350,5.701) (-0.400,5.694) (-0.450,5.687) (-0.500,5.678) (-0.550,5.669) (-0.600,5.659) (-0.650,5.647) (-0.700,5.635) (-0.750,5.622) (-0.800,5.608) (-0.850,5.593) (-0.900,5.577) (-0.950,5.559) (-1.000,5.541) (-1.050,5.522) (-1.100,5.501) (-1.150,5.480) (-1.200,5.457) (-1.250,5.433) (-1.300,5.407) (-1.350,5.381) (-1.400,5.353) (-1.447,5.325) (-1.500,5.293) (-1.550,5.260) (-1.600,5.226) (-1.650,5.191) (-1.700,5.153) (-1.750,5.114) (-1.797,5.075) (-1.850,5.029) (-1.900,4.983) (-1.950,4.935) (-2.000,4.884) (-2.050,4.830) (-2.098,4.775) (-2.140,4.725) (-2.175,4.680) (-2.216,4.625) (-2.250,4.577) (-2.285,4.525) (-2.317,4.475) (-2.347,4.425) (-2.375,4.376) (-2.400,4.330) (-2.425,4.281) (-2.450,4.229) (-2.474,4.175) (-2.496,4.125) (-2.515,4.075) (-2.534,4.025) (-2.551,3.975) (-2.567,3.925) (-2.582,3.875) (-2.595,3.825) (-2.608,3.775) (-2.619,3.725) (-2.629,3.675) (-2.637,3.625) (-2.645,3.575) (-2.651,3.525) (-2.657,3.475) (-2.661,3.425) (-2.664,3.375) (-2.666,3.325) (-2.667,3.275) (-2.666,3.225) (-2.665,3.175) (-2.662,3.125) (-2.658,3.075) (-2.654,3.025) (-2.647,2.975) (-2.640,2.925) (-2.632,2.875) (-2.623,2.825) (-2.612,2.775) (-2.600,2.725) (-2.587,2.675) (-2.573,2.625) (-2.558,2.575) (-2.541,2.525) (-2.523,2.475) (-2.504,2.425) (-2.483,2.375) (-2.462,2.325) (-2.439,2.275) (-2.414,2.225) (-2.389,2.175) (-2.362,2.125) (-2.333,2.075) (-2.303,2.025) (-2.275,1.980) (-2.239,1.925) (-2.205,1.875) (-2.175,1.833) (-2.132,1.775) (-2.100,1.733) (-2.053,1.675) (-2.012,1.625) (-1.975,1.583) (-1.925,1.527) (-1.877,1.475) (-1.829,1.425) (-1.780,1.375) (-1.729,1.325) (-1.676,1.275) (-1.625,1.228) (-1.575,1.183) (-1.525,1.139) (-1.481,1.100) (-1.425,1.053) (-1.375,1.011) (-1.331,0.975) (-1.275,0.930) (-1.225,0.891) (-1.175,0.851) (-1.125,0.813) (-1.076,0.775) (-1.025,0.737) (-0.976,0.700) (-0.925,0.662) (-0.875,0.625) (-0.825,0.589) (-0.775,0.552) (-0.725,0.516) (-0.675,0.480) (-0.635,0.450) (-0.575,0.408) (-0.530,0.375) (-0.475,0.337) (-0.425,0.301) (-0.375,0.266) (-0.325,0.230) (-0.284,0.200) (-0.225,0.160) (-0.220,0.155)}
\EQ{(-0.220,-0.155) (-0.275,-0.195) (-0.319,-0.225) (-0.375,-0.266) (-0.424,-0.300) (-0.475,-0.337) (-0.525,-0.372) (-0.575,-0.408) (-0.625,-0.444) (-0.669,-0.475) (-0.725,-0.516) (-0.772,-0.550) (-0.825,-0.589) (-0.874,-0.625) (-0.925,-0.663) (-0.975,-0.700) (-1.025,-0.738) (-1.074,-0.775) (-1.125,-0.815) (-1.171,-0.850) (-1.225,-0.893) (-1.265,-0.925) (-1.325,-0.974) (-1.375,-1.016) (-1.416,-1.050) (-1.473,-1.100) (-1.525,-1.146) (-1.575,-1.191) (-1.612,-1.225) (-1.650,-1.261) (-1.691,-1.300) (-1.741,-1.350) (-1.790,-1.400) (-1.825,-1.438) (-1.860,-1.475) (-1.900,-1.520) (-1.947,-1.575) (-1.989,-1.625) (-2.025,-1.670) (-2.067,-1.725) (-2.100,-1.770) (-2.139,-1.825) (-2.173,-1.875) (-2.205,-1.925) (-2.236,-1.975) (-2.265,-2.025) (-2.293,-2.075) (-2.319,-2.125) (-2.344,-2.175) (-2.368,-2.225) (-2.390,-2.275) (-2.411,-2.325) (-2.431,-2.375) (-2.449,-2.425) (-2.466,-2.475) (-2.481,-2.525) (-2.496,-2.575) (-2.509,-2.625) (-2.521,-2.675) (-2.532,-2.725) (-2.541,-2.775) (-2.549,-2.825) (-2.557,-2.875) (-2.562,-2.925) (-2.567,-2.975) (-2.571,-3.025) (-2.573,-3.075) (-2.574,-3.125) (-2.574,-3.175) (-2.573,-3.225) (-2.571,-3.275) (-2.568,-3.325) (-2.563,-3.375) (-2.557,-3.425) (-2.550,-3.475) (-2.542,-3.525) (-2.533,-3.575) (-2.522,-3.625) (-2.510,-3.675) (-2.497,-3.725) (-2.483,-3.775) (-2.467,-3.825) (-2.450,-3.875) (-2.432,-3.925) (-2.413,-3.975) (-2.392,-4.025) (-2.369,-4.075) (-2.345,-4.125) (-2.320,-4.175) (-2.293,-4.225) (-2.264,-4.275) (-2.234,-4.325) (-2.202,-4.375) (-2.168,-4.425) (-2.133,-4.475) (-2.100,-4.518) (-2.055,-4.575) (-2.013,-4.625) (-1.975,-4.668) (-1.925,-4.722) (-1.875,-4.772) (-1.825,-4.820) (-1.775,-4.866) (-1.735,-4.900) (-1.675,-4.949) (-1.625,-4.988) (-1.575,-5.025) (-1.525,-5.060) (-1.475,-5.093) (-1.425,-5.124) (-1.375,-5.155) (-1.325,-5.183) (-1.275,-5.210) (-1.225,-5.236) (-1.175,-5.261) (-1.125,-5.284) (-1.075,-5.306) (-1.025,-5.327) (-0.975,-5.346) (-0.925,-5.365) (-0.875,-5.382) (-0.825,-5.398) (-0.775,-5.414) (-0.725,-5.428) (-0.675,-5.441) (-0.625,-5.453) (-0.575,-5.464) (-0.525,-5.474) (-0.475,-5.484) (-0.425,-5.492) (-0.375,-5.499) (-0.325,-5.506) (-0.275,-5.511) (-0.225,-5.516) (-0.175,-5.519) (-0.125,-5.522) (-0.075,-5.524) (-0.025,-5.525) (0.050,-5.524) (0.100,-5.523) (0.150,-5.521) (0.200,-5.518) (0.250,-5.513) (0.300,-5.508) (0.350,-5.503) (0.400,-5.496) (0.450,-5.488) (0.500,-5.479) (0.550,-5.469) (0.600,-5.459) (0.650,-5.447) (0.700,-5.434) (0.750,-5.421) (0.800,-5.406) (0.850,-5.390) (0.900,-5.374) (0.950,-5.356) (1.000,-5.337) (1.050,-5.316) (1.100,-5.295) (1.150,-5.272) (1.197,-5.250) (1.247,-5.225) (1.294,-5.200) (1.340,-5.175) (1.400,-5.140) (1.450,-5.109) (1.500,-5.076) (1.550,-5.042) (1.600,-5.006) (1.642,-4.975) (1.700,-4.929) (1.750,-4.887) (1.792,-4.850) (1.846,-4.800) (1.898,-4.750) (1.946,-4.700) (1.991,-4.650) (2.025,-4.611) (2.075,-4.550) (2.114,-4.500) (2.150,-4.451) (2.185,-4.400) (2.218,-4.350) (2.249,-4.300) (2.275,-4.257) (2.300,-4.212) (2.325,-4.165) (2.350,-4.115) (2.375,-4.062) (2.400,-4.005) (2.423,-3.950) (2.442,-3.900) (2.459,-3.850) (2.475,-3.801) (2.490,-3.750) (2.504,-3.700) (2.516,-3.650) (2.528,-3.600) (2.538,-3.550) (2.546,-3.500) (2.554,-3.450) (2.560,-3.400) (2.566,-3.350) (2.570,-3.300) (2.572,-3.250) (2.574,-3.200) (2.575,-3.150) (2.574,-3.100) (2.572,-3.050) (2.569,-3.000) (2.565,-2.950) (2.560,-2.900) (2.553,-2.850) (2.545,-2.800) (2.537,-2.750) (2.526,-2.700) (2.515,-2.650) (2.502,-2.600) (2.489,-2.550) (2.474,-2.500) (2.457,-2.450) (2.440,-2.400) (2.421,-2.350) (2.401,-2.300) (2.379,-2.250) (2.356,-2.200) (2.332,-2.150) (2.306,-2.100) (2.279,-2.050) (2.251,-2.000) (2.221,-1.950) (2.189,-1.900) (2.156,-1.850) (2.125,-1.805) (2.086,-1.750) (2.050,-1.702) (2.009,-1.650) (1.975,-1.608) (1.926,-1.550) (1.882,-1.500) (1.837,-1.450) (1.800,-1.411) (1.765,-1.375) (1.725,-1.334) (1.675,-1.285) (1.639,-1.250) (1.600,-1.214) (1.558,-1.175) (1.502,-1.125) (1.450,-1.080) (1.400,-1.037) (1.356,-1.000) (1.300,-0.954) (1.250,-0.913) (1.202,-0.875) (1.150,-0.834) (1.107,-0.800) (1.050,-0.757) (1.008,-0.725) (0.950,-0.682) (0.908,-0.650) (0.850,-0.607) (0.807,-0.575) (0.750,-0.534) (0.703,-0.500) (0.650,-0.462) (0.600,-0.426) (0.550,-0.390) (0.500,-0.355) (0.450,-0.319) (0.400,-0.284) (0.354,-0.250) (0.300,-0.213) (0.250,-0.177) (0.206,-0.144)}
\EQ{(0.206,0.144) (0.248,0.175) (0.300,0.213) (0.350,0.248) (0.400,0.284) (0.450,0.319) (0.495,0.350) (0.550,0.390) (0.599,0.425) (0.650,0.462) (0.700,0.498) (0.750,0.534) (0.800,0.571) (0.841,0.600) (0.900,0.644) (0.942,0.675) (1.000,0.719) (1.041,0.750) (1.100,0.795) (1.150,0.834) (1.200,0.873) (1.250,0.913) (1.296,0.950) (1.350,0.995) (1.400,1.037) (1.445,1.075) (1.500,1.123) (1.550,1.168) (1.600,1.214) (1.639,1.250) (1.675,1.285) (1.716,1.325) (1.765,1.375) (1.800,1.411) (1.837,1.450) (1.875,1.492) (1.925,1.549) (1.968,1.600) (2.000,1.639) (2.048,1.700) (2.086,1.750) (2.122,1.800) (2.156,1.850) (2.189,1.900) (2.221,1.950) (2.250,1.999) (2.275,2.043) (2.300,2.088) (2.325,2.136) (2.350,2.187) (2.375,2.241) (2.400,2.299) (2.421,2.350) (2.440,2.400) (2.457,2.450) (2.474,2.500) (2.489,2.550) (2.502,2.600) (2.515,2.650) (2.526,2.700) (2.537,2.750) (2.545,2.800) (2.553,2.850) (2.560,2.900) (2.565,2.950) (2.569,3.000) (2.572,3.050) (2.574,3.100) (2.575,3.150) (2.574,3.200) (2.572,3.250) (2.570,3.300) (2.566,3.350) (2.560,3.400) (2.554,3.450) (2.546,3.500) (2.538,3.550) (2.528,3.600) (2.516,3.650) (2.504,3.700) (2.490,3.750) (2.475,3.800) (2.459,3.850) (2.442,3.900) (2.423,3.950) (2.402,4.000) (2.381,4.050) (2.357,4.100) (2.333,4.150) (2.307,4.200) (2.279,4.250) (2.250,4.299) (2.218,4.350) (2.185,4.400) (2.151,4.450) (2.114,4.500) (2.075,4.550) (2.034,4.600) (2.000,4.640) (1.950,4.695) (1.900,4.747) (1.850,4.797) (1.800,4.843) (1.750,4.887) (1.705,4.925) (1.650,4.969) (1.609,5.000) (1.550,5.042) (1.502,5.075) (1.450,5.109) (1.400,5.140) (1.350,5.169) (1.300,5.197) (1.250,5.223) (1.200,5.249) (1.150,5.272) (1.100,5.295) (1.050,5.316) (1.000,5.337) (0.950,5.356) (0.900,5.374) (0.850,5.390) (0.800,5.406) (0.750,5.421) (0.700,5.434) (0.650,5.447) (0.600,5.459) (0.550,5.469) (0.500,5.479) (0.450,5.488) (0.400,5.496) (0.350,5.503) (0.300,5.508) (0.250,5.513) (0.200,5.518) (0.150,5.521) (0.100,5.523) (0.050,5.524) (0.000,5.525) (-0.050,5.524) (-0.100,5.523) (-0.150,5.521) (-0.200,5.518) (-0.250,5.513) (-0.300,5.508) (-0.350,5.503) (-0.400,5.496) (-0.450,5.488) (-0.500,5.479) (-0.550,5.469) (-0.600,5.459) (-0.650,5.447) (-0.700,5.434) (-0.750,5.421) (-0.800,5.406) (-0.850,5.390) (-0.900,5.374) (-0.950,5.356) (-1.000,5.337) (-1.050,5.316) (-1.100,5.295) (-1.150,5.272) (-1.197,5.250) (-1.247,5.225) (-1.294,5.200) (-1.340,5.175) (-1.400,5.140) (-1.450,5.109) (-1.500,5.076) (-1.550,5.042) (-1.600,5.006) (-1.642,4.975) (-1.700,4.929) (-1.750,4.887) (-1.792,4.850) (-1.846,4.800) (-1.898,4.750) (-1.946,4.700) (-1.991,4.650) (-2.025,4.611) (-2.075,4.550) (-2.114,4.500) (-2.150,4.451) (-2.185,4.400) (-2.218,4.350) (-2.249,4.300) (-2.275,4.257) (-2.300,4.212) (-2.325,4.165) (-2.350,4.115) (-2.375,4.062) (-2.400,4.005) (-2.423,3.950) (-2.442,3.900) (-2.459,3.850) (-2.475,3.801) (-2.490,3.750) (-2.504,3.700) (-2.516,3.650) (-2.528,3.600) (-2.538,3.550) (-2.546,3.500) (-2.554,3.450) (-2.560,3.400) (-2.566,3.350) (-2.570,3.300) (-2.572,3.250) (-2.574,3.200) (-2.575,3.150) (-2.574,3.100) (-2.572,3.050) (-2.569,3.000) (-2.565,2.950) (-2.560,2.900) (-2.553,2.850) (-2.545,2.800) (-2.537,2.750) (-2.526,2.700) (-2.515,2.650) (-2.502,2.600) (-2.489,2.550) (-2.474,2.500) (-2.457,2.450) (-2.440,2.400) (-2.421,2.350) (-2.401,2.300) (-2.379,2.250) (-2.356,2.200) (-2.332,2.150) (-2.306,2.100) (-2.279,2.050) (-2.251,2.000) (-2.221,1.950) (-2.189,1.900) (-2.156,1.850) (-2.125,1.805) (-2.086,1.750) (-2.050,1.702) (-2.009,1.650) (-1.975,1.608) (-1.926,1.550) (-1.882,1.500) (-1.837,1.450) (-1.800,1.411) (-1.765,1.375) (-1.725,1.334) (-1.675,1.285) (-1.639,1.250) (-1.600,1.214) (-1.558,1.175) (-1.502,1.125) (-1.450,1.080) (-1.400,1.037) (-1.356,1.000) (-1.300,0.954) (-1.250,0.913) (-1.202,0.875) (-1.150,0.834) (-1.107,0.800) (-1.050,0.757) (-1.008,0.725) (-0.950,0.682) (-0.908,0.650) (-0.850,0.607) (-0.807,0.575) (-0.750,0.534) (-0.703,0.500) (-0.650,0.462) (-0.600,0.426) (-0.550,0.390) (-0.500,0.355) (-0.450,0.319) (-0.400,0.284) (-0.354,0.250) (-0.300,0.213) (-0.250,0.177) (-0.220,0.155)}

    % Cuadrupolo puntual - líneas de campo
% Generado por scripts/generar_datos.py (no editar a mano).
\LC{(0.435,1.114) (0.489,1.167) (0.546,1.216) (0.605,1.261) (0.667,1.304) (0.732,1.342) (0.799,1.375) (0.868,1.405) (0.939,1.429) (1.011,1.448) (1.085,1.461) (1.160,1.468) (1.235,1.468) (1.310,1.461) (1.383,1.446) (1.455,1.424) (1.523,1.394) (1.588,1.355) (1.646,1.309) (1.698,1.255) (1.742,1.194) (1.777,1.128) (1.803,1.057) (1.817,0.984) (1.822,0.909) (1.817,0.834) (1.803,0.761) (1.780,0.689) (1.749,0.621) (1.712,0.556) (1.669,0.495) (1.621,0.437) (1.568,0.384) (1.511,0.335) (1.451,0.290) (1.388,0.249) (1.323,0.213) (1.255,0.179) (1.186,0.150)}
\LC{(-0.449,1.110) (-0.504,1.161) (-0.562,1.208) (-0.623,1.252) (-0.687,1.292) (-0.753,1.327) (-0.821,1.358) (-0.891,1.384) (-0.963,1.405) (-1.037,1.420) (-1.111,1.428) (-1.186,1.430) (-1.261,1.424) (-1.335,1.411) (-1.407,1.390) (-1.476,1.361) (-1.541,1.323) (-1.600,1.278) (-1.652,1.224) (-1.696,1.163) (-1.731,1.097) (-1.756,1.026) (-1.770,0.953) (-1.774,0.878) (-1.768,0.803) (-1.753,0.730) (-1.728,0.659) (-1.697,0.591) (-1.658,0.527) (-1.613,0.467) (-1.564,0.411) (-1.510,0.359) (-1.452,0.311) (-1.391,0.268) (-1.327,0.228) (-1.261,0.193) (-1.193,0.161) (-1.170,0.152)}
\LC{(-0.693,-0.961) (-0.768,-0.970) (-0.843,-0.968) (-0.916,-0.955) (-0.987,-0.930) (-1.052,-0.893) (-1.109,-0.844) (-1.154,-0.784) (-1.185,-0.716) (-1.201,-0.643) (-1.202,-0.568) (-1.189,-0.495) (-1.163,-0.424) (-1.126,-0.359)}
\LC{(0.153,-1.185) (0.173,-1.258) (0.193,-1.330) (0.215,-1.401) (0.239,-1.473) (0.263,-1.544) (0.289,-1.614) (0.316,-1.684) (0.343,-1.754) (0.373,-1.823) (0.403,-1.892) (0.434,-1.960) (0.467,-2.027) (0.501,-2.094) (0.536,-2.160) (0.572,-2.226) (0.609,-2.291) (0.648,-2.356) (0.687,-2.419) (0.728,-2.482) (0.770,-2.545) (0.813,-2.606) (0.857,-2.667) (0.902,-2.727) (0.948,-2.786) (0.995,-2.844) (1.044,-2.901) (1.094,-2.957) (1.144,-3.013) (1.196,-3.067) (1.249,-3.120) (1.303,-3.172) (1.358,-3.223) (1.413,-3.273) (1.470,-3.322) (1.528,-3.370) (1.587,-3.416) (1.647,-3.461) (1.708,-3.505) (1.770,-3.547) (1.833,-3.588) (1.897,-3.628) (1.961,-3.666) (2.027,-3.703) (2.093,-3.738) (2.160,-3.771) (2.228,-3.803) (2.297,-3.833) (2.366,-3.861) (2.436,-3.888) (2.507,-3.913) (2.578,-3.936) (2.651,-3.956) (2.723,-3.975) (2.796,-3.992) (2.870,-4.007) (2.944,-4.019) (3.018,-4.030) (3.093,-4.038) (3.167,-4.043) (3.242,-4.046) (3.317,-4.047) (3.392,-4.046) (3.467,-4.041) (3.542,-4.035) (3.616,-4.025) (3.690,-4.013) (3.764,-3.998) (3.836,-3.980) (3.909,-3.959) (3.980,-3.935) (4.050,-3.909) (4.119,-3.879) (4.186,-3.847) (4.253,-3.812) (4.317,-3.773) (4.380,-3.732) (4.440,-3.688) (4.499,-3.641) (4.555,-3.591) (4.608,-3.539) (4.659,-3.484) (4.707,-3.426) (4.752,-3.366) (4.794,-3.304) (4.833,-3.240) (4.868,-3.173) (4.900,-3.105) (4.928,-3.036) (4.953,-2.965) (4.973,-2.893) (4.991,-2.820) (5.004,-2.746) (5.014,-2.672) (5.020,-2.597) (5.022,-2.522) (5.021,-2.447) (5.016,-2.372) (5.008,-2.298) (4.996,-2.224) (4.981,-2.150) (4.963,-2.077) (4.942,-2.005) (4.918,-1.934) (4.891,-1.864) (4.862,-1.795) (4.830,-1.728) (4.795,-1.661) (4.758,-1.596) (4.719,-1.532) (4.677,-1.469) (4.634,-1.408) (4.588,-1.349) (4.541,-1.291) (4.492,-1.234) (4.441,-1.179) (4.389,-1.125) (4.335,-1.073) (4.280,-1.022) (4.223,-0.973) (4.165,-0.925) (4.106,-0.879) (4.046,-0.834) (3.985,-0.791) (3.923,-0.749) (3.859,-0.709) (3.795,-0.670) (3.730,-0.632) (3.665,-0.596) (3.598,-0.561) (3.531,-0.528) (3.463,-0.496) (3.395,-0.465) (3.326,-0.436) (3.256,-0.407) (3.186,-0.380) (3.116,-0.355) (3.045,-0.330) (2.974,-0.306) (2.902,-0.284) (2.830,-0.263) (2.758,-0.243) (2.686,-0.224) (2.613,-0.206) (2.540,-0.189) (2.467,-0.172) (2.393,-0.157) (2.319,-0.143) (2.246,-0.130) (2.172,-0.117) (2.098,-0.105) (2.023,-0.095) (1.949,-0.084) (1.875,-0.075) (1.800,-0.066) (1.751,-0.061)}
\LC{(0.283,-2.625) (0.299,-2.698) (0.316,-2.771) (0.333,-2.844) (0.350,-2.917) (0.368,-2.990) (0.387,-3.063) (0.406,-3.135) (0.425,-3.208) (0.445,-3.280) (0.466,-3.352) (0.486,-3.424) (0.507,-3.496) (0.529,-3.568) (0.551,-3.640) (0.574,-3.711) (0.597,-3.783) (0.620,-3.854) (0.644,-3.925) (0.669,-3.996) (0.693,-4.067) (0.719,-4.137) (0.744,-4.208) (0.770,-4.278) (0.797,-4.348) (0.824,-4.418) (0.852,-4.488) (0.879,-4.557) (0.908,-4.627) (0.937,-4.696) (0.966,-4.765) (0.996,-4.834) (1.026,-4.903) (1.056,-4.971) (1.087,-5.039) (1.119,-5.108) (1.151,-5.175) (1.183,-5.243) (1.216,-5.310) (1.249,-5.378) (1.283,-5.445) (1.317,-5.511) (1.352,-5.578) (1.387,-5.644) (1.422,-5.710) (1.458,-5.776) (1.494,-5.842) (1.531,-5.907) (1.568,-5.972) (1.606,-6.037) (1.644,-6.102) (1.682,-6.166) (1.721,-6.230) (1.761,-6.294) (1.801,-6.358) (1.841,-6.421) (1.882,-6.484) (1.923,-6.547) (1.964,-6.609) (2.006,-6.671) (2.048,-6.733) (2.091,-6.795) (2.135,-6.856) (2.178,-6.917) (2.222,-6.978) (2.267,-7.038) (2.312,-7.098) (2.357,-7.158) (2.403,-7.217) (2.449,-7.276) (2.496,-7.335) (2.543,-7.394) (2.590,-7.452) (2.638,-7.509) (2.686,-7.567) (2.735,-7.624) (2.784,-7.680) (2.834,-7.737) (2.884,-7.793) (2.934,-7.848) (2.985,-7.903) (3.036,-7.958) (3.088,-8.013) (3.140,-8.067) (3.192,-8.120) (3.245,-8.174) (3.298,-8.227) (3.352,-8.279) (3.406,-8.331) (3.460,-8.383) (3.515,-8.434) (3.570,-8.485) (3.626,-8.535) (3.682,-8.585) (3.738,-8.634) (3.795,-8.683) (3.852,-8.732) (3.909,-8.780) (3.967,-8.828) (4.026,-8.875) (4.084,-8.922) (4.143,-8.968) (4.203,-9.014) (4.262,-9.059) (4.323,-9.104) (4.383,-9.148) (4.444,-9.192) (4.505,-9.235) (4.567,-9.278) (4.629,-9.320) (4.691,-9.362) (4.754,-9.403) (4.817,-9.444) (4.880,-9.484) (4.944,-9.524) (5.008,-9.563) (5.072,-9.601) (5.137,-9.639) (5.202,-9.676) (5.267,-9.713) (5.333,-9.749) (5.399,-9.785) (5.466,-9.820) (5.532,-9.854) (5.599,-9.888) (5.666,-9.921) (5.734,-9.954) (5.802,-9.986) (5.847,-10.000)}
\LC{(-0.162,-1.698) (-0.176,-1.772) (-0.191,-1.845) (-0.207,-1.919) (-0.223,-1.992) (-0.240,-2.065) (-0.258,-2.138) (-0.276,-2.210) (-0.295,-2.283) (-0.314,-2.356) (-0.334,-2.428) (-0.355,-2.500) (-0.376,-2.572) (-0.398,-2.644) (-0.421,-2.715) (-0.444,-2.786) (-0.468,-2.858) (-0.492,-2.928) (-0.517,-2.999) (-0.543,-3.070) (-0.569,-3.140) (-0.596,-3.210) (-0.623,-3.280) (-0.651,-3.349) (-0.680,-3.419) (-0.709,-3.488) (-0.739,-3.557) (-0.769,-3.625) (-0.800,-3.693) (-0.832,-3.761) (-0.864,-3.829) (-0.897,-3.896) (-0.931,-3.963) (-0.965,-4.030) (-1.000,-4.097) (-1.035,-4.163) (-1.071,-4.229) (-1.107,-4.294) (-1.144,-4.359) (-1.182,-4.424) (-1.220,-4.489) (-1.259,-4.553) (-1.299,-4.617) (-1.339,-4.680) (-1.380,-4.743) (-1.421,-4.806) (-1.463,-4.868) (-1.505,-4.930) (-1.548,-4.991) (-1.592,-5.052) (-1.636,-5.113) (-1.681,-5.173) (-1.726,-5.233) (-1.772,-5.292) (-1.819,-5.351) (-1.866,-5.409) (-1.913,-5.467) (-1.961,-5.525) (-2.010,-5.582) (-2.060,-5.638) (-2.110,-5.694) (-2.160,-5.749) (-2.211,-5.804) (-2.263,-5.859) (-2.315,-5.913) (-2.368,-5.966) (-2.421,-6.019) (-2.475,-6.071) (-2.529,-6.123) (-2.584,-6.174) (-2.639,-6.224) (-2.695,-6.274) (-2.752,-6.324) (-2.809,-6.372) (-2.866,-6.420) (-2.924,-6.468) (-2.983,-6.515) (-3.042,-6.561) (-3.102,-6.606) (-3.162,-6.651) (-3.222,-6.695) (-3.284,-6.739) (-3.345,-6.782) (-3.407,-6.824) (-3.470,-6.865) (-3.533,-6.906) (-3.597,-6.945) (-3.661,-6.984) (-3.725,-7.023) (-3.790,-7.060) (-3.855,-7.097) (-3.921,-7.133) (-3.987,-7.168) (-4.054,-7.202) (-4.121,-7.236) (-4.189,-7.268) (-4.257,-7.300) (-4.325,-7.331) (-4.394,-7.361) (-4.463,-7.390) (-4.533,-7.418) (-4.602,-7.445) (-4.673,-7.472) (-4.743,-7.497) (-4.814,-7.521) (-4.886,-7.545) (-4.957,-7.567) (-5.029,-7.588) (-5.101,-7.609) (-5.174,-7.628) (-5.246,-7.646) (-5.319,-7.663) (-5.393,-7.679) (-5.466,-7.694) (-5.540,-7.708) (-5.614,-7.721) (-5.688,-7.732) (-5.762,-7.743) (-5.837,-7.752) (-5.911,-7.760) (-5.986,-7.767) (-6.061,-7.773) (-6.136,-7.777) (-6.210,-7.780) (-6.285,-7.782) (-6.360,-7.782) (-6.435,-7.782) (-6.510,-7.779) (-6.585,-7.776) (-6.660,-7.771) (-6.735,-7.765) (-6.810,-7.758) (-6.884,-7.749) (-6.958,-7.739) (-7.032,-7.727) (-7.106,-7.714) (-7.180,-7.699) (-7.253,-7.683) (-7.326,-7.665) (-7.399,-7.646) (-7.471,-7.626) (-7.542,-7.604) (-7.614,-7.581) (-7.684,-7.556) (-7.754,-7.529) (-7.824,-7.501) (-7.893,-7.471) (-7.961,-7.440) (-8.029,-7.408) (-8.095,-7.373) (-8.161,-7.338) (-8.227,-7.300) (-8.291,-7.262) (-8.354,-7.221) (-8.416,-7.180) (-8.477,-7.136) (-8.538,-7.091) (-8.597,-7.045) (-8.654,-6.997) (-8.711,-6.948) (-8.766,-6.898) (-8.820,-6.845) (-8.873,-6.792) (-8.924,-6.737) (-8.974,-6.681) (-9.022,-6.624) (-9.069,-6.565) (-9.114,-6.505) (-9.158,-6.444) (-9.200,-6.382) (-9.240,-6.319) (-9.278,-6.254) (-9.315,-6.189) (-9.350,-6.123) (-9.383,-6.055) (-9.415,-5.987) (-9.444,-5.918) (-9.472,-5.849) (-9.497,-5.778) (-9.521,-5.707) (-9.543,-5.635) (-9.563,-5.563) (-9.581,-5.490) (-9.597,-5.417) (-9.611,-5.343) (-9.623,-5.269) (-9.634,-5.195) (-9.642,-5.120) (-9.648,-5.046) (-9.653,-4.971) (-9.656,-4.896) (-9.656,-4.821) (-9.655,-4.746) (-9.652,-4.671) (-9.648,-4.596) (-9.641,-4.521) (-9.633,-4.447) (-9.622,-4.373) (-9.611,-4.298) (-9.597,-4.225) (-9.582,-4.151) (-9.565,-4.078) (-9.546,-4.006) (-9.526,-3.933) (-9.505,-3.861) (-9.482,-3.790) (-9.457,-3.719) (-9.431,-3.649) (-9.403,-3.579) (-9.374,-3.510) (-9.344,-3.442) (-9.312,-3.374) (-9.279,-3.306) (-9.245,-3.240) (-9.209,-3.174) (-9.172,-3.108) (-9.134,-3.044) (-9.095,-2.980) (-9.055,-2.916) (-9.013,-2.854) (-8.971,-2.792) (-8.927,-2.731) (-8.882,-2.671) (-8.837,-2.611) (-8.790,-2.553) (-8.743,-2.495) (-8.694,-2.437) (-8.645,-2.381) (-8.595,-2.325) (-8.544,-2.270) (-8.492,-2.216) (-8.439,-2.163) (-8.385,-2.110) (-8.331,-2.058) (-8.276,-2.007) (-8.221,-1.957) (-8.164,-1.908) (-8.107,-1.859) (-8.049,-1.811) (-7.991,-1.764) (-7.932,-1.718) (-7.872,-1.672) (-7.812,-1.627) (-7.751,-1.583) (-7.690,-1.540) (-7.628,-1.498) (-7.566,-1.456) (-7.503,-1.415) (-7.440,-1.375) (-7.376,-1.335) (-7.312,-1.297) (-7.247,-1.259) (-7.182,-1.221) (-7.117,-1.185) (-7.051,-1.149) (-6.985,-1.114) (-6.918,-1.079) (-6.851,-1.046) (-6.784,-1.012) (-6.716,-0.980) (-6.648,-0.948) (-6.580,-0.917) (-6.511,-0.887) (-6.442,-0.858) (-6.373,-0.829) (-6.304,-0.800) (-6.234,-0.773) (-6.164,-0.745) (-6.094,-0.719) (-6.023,-0.693) (-5.953,-0.668) (-5.882,-0.643) (-5.811,-0.620) (-5.740,-0.596) (-5.668,-0.573) (-5.596,-0.551) (-5.525,-0.530) (-5.453,-0.508) (-5.380,-0.488) (-5.308,-0.468) (-5.236,-0.449) (-5.163,-0.430) (-5.090,-0.412) (-5.018,-0.394) (-4.945,-0.376) (-4.871,-0.360) (-4.798,-0.343) (-4.725,-0.328) (-4.652,-0.312) (-4.578,-0.298) (-4.504,-0.283) (-4.431,-0.269) (-4.357,-0.256) (-4.283,-0.243) (-4.209,-0.231) (-4.135,-0.219) (-4.061,-0.207) (-3.987,-0.196) (-3.913,-0.185) (-3.838,-0.174) (-3.764,-0.164) (-3.690,-0.155) (-3.615,-0.146) (-3.541,-0.137) (-3.466,-0.128) (-3.392,-0.120) (-3.317,-0.112) (-3.242,-0.105) (-3.168,-0.098) (-3.093,-0.091) (-3.018,-0.085) (-2.944,-0.078) (-2.869,-0.073) (-2.794,-0.067) (-2.719,-0.062) (-2.644,-0.057) (-2.569,-0.052) (-2.495,-0.048) (-2.420,-0.044) (-2.345,-0.040) (-2.270,-0.036) (-2.195,-0.032) (-2.120,-0.029) (-2.045,-0.026) (-1.970,-0.023) (-1.895,-0.021) (-1.820,-0.019) (-1.770,-0.017)}
\LC{(-0.196,1.492) (-0.216,1.565) (-0.237,1.637) (-0.259,1.708) (-0.282,1.780) (-0.305,1.851) (-0.330,1.922) (-0.356,1.992) (-0.382,2.062) (-0.410,2.132) (-0.439,2.201) (-0.468,2.270) (-0.499,2.339) (-0.530,2.407) (-0.562,2.475) (-0.595,2.542) (-0.630,2.609) (-0.665,2.675) (-0.701,2.741) (-0.738,2.806) (-0.776,2.871) (-0.815,2.935) (-0.854,2.998) (-0.895,3.061) (-0.937,3.124) (-0.979,3.185) (-1.023,3.247) (-1.067,3.307) (-1.112,3.367) (-1.158,3.426) (-1.206,3.484) (-1.253,3.542) (-1.302,3.599) (-1.352,3.655) (-1.403,3.711) (-1.454,3.765) (-1.506,3.819) (-1.559,3.872) (-1.614,3.924) (-1.668,3.975) (-1.724,4.025) (-1.781,4.074) (-1.838,4.123) (-1.896,4.170) (-1.955,4.217) (-2.015,4.262) (-2.075,4.306) (-2.137,4.349) (-2.199,4.391) (-2.262,4.432) (-2.325,4.472) (-2.389,4.511) (-2.454,4.548) (-2.520,4.584) (-2.587,4.619) (-2.654,4.653) (-2.721,4.685) (-2.790,4.716) (-2.859,4.745) (-2.928,4.773) (-2.998,4.800) (-3.069,4.825) (-3.140,4.849) (-3.212,4.871) (-3.284,4.891) (-3.357,4.910) (-3.430,4.928) (-3.503,4.943) (-3.577,4.957) (-3.651,4.969) (-3.725,4.979) (-3.800,4.987) (-3.874,4.994) (-3.949,4.998) (-4.024,5.001) (-4.099,5.002) (-4.174,5.000) (-4.249,4.997) (-4.324,4.991) (-4.398,4.983) (-4.473,4.973) (-4.547,4.961) (-4.620,4.947) (-4.693,4.930) (-4.766,4.911) (-4.838,4.890) (-4.909,4.866) (-4.979,4.840) (-5.049,4.812) (-5.117,4.781) (-5.185,4.748) (-5.251,4.713) (-5.315,4.675) (-5.379,4.635) (-5.440,4.592) (-5.500,4.547) (-5.559,4.500) (-5.615,4.450) (-5.669,4.399) (-5.722,4.345) (-5.772,4.289) (-5.819,4.231) (-5.865,4.171) (-5.907,4.110) (-5.947,4.046) (-5.985,3.981) (-6.019,3.915) (-6.051,3.847) (-6.080,3.778) (-6.106,3.707) (-6.129,3.636) (-6.149,3.564) (-6.166,3.491) (-6.180,3.417) (-6.191,3.343) (-6.199,3.268) (-6.204,3.193) (-6.206,3.118) (-6.205,3.043) (-6.201,2.968) (-6.195,2.894) (-6.186,2.819) (-6.174,2.745) (-6.159,2.672) (-6.142,2.599) (-6.122,2.526) (-6.100,2.455) (-6.075,2.384) (-6.049,2.314) (-6.020,2.245) (-5.989,2.176) (-5.955,2.109) (-5.920,2.043) (-5.883,1.978) (-5.844,1.913) (-5.804,1.850) (-5.761,1.789) (-5.717,1.728) (-5.672,1.668) (-5.624,1.610) (-5.576,1.553) (-5.526,1.497) (-5.475,1.442) (-5.422,1.389) (-5.368,1.336) (-5.314,1.285) (-5.258,1.235) (-5.200,1.187) (-5.142,1.139) (-5.083,1.093) (-5.023,1.048) (-4.963,1.004) (-4.901,0.961) (-4.838,0.920) (-4.775,0.880) (-4.711,0.840) (-4.647,0.802) (-4.581,0.765) (-4.515,0.730) (-4.449,0.695) (-4.382,0.661) (-4.314,0.629) (-4.246,0.597) (-4.178,0.567) (-4.109,0.538) (-4.039,0.509) (-3.969,0.482) (-3.899,0.456) (-3.828,0.430) (-3.758,0.406) (-3.686,0.382) (-3.615,0.360) (-3.543,0.338) (-3.495,0.324)}
\LC{(0.436,1.693) (0.474,1.758) (0.513,1.822) (0.553,1.885) (0.595,1.947) (0.639,2.009) (0.684,2.069) (0.730,2.127) (0.778,2.185) (0.827,2.242) (0.878,2.297) (0.930,2.351) (0.984,2.403) (1.039,2.454) (1.095,2.504) (1.153,2.552) (1.212,2.598) (1.272,2.643) (1.333,2.686) (1.396,2.727) (1.460,2.766) (1.525,2.803) (1.591,2.839) (1.659,2.872) (1.727,2.903) (1.796,2.931) (1.866,2.957) (1.938,2.981) (2.009,3.003) (2.082,3.021) (2.155,3.037) (2.229,3.050) (2.304,3.060) (2.378,3.067) (2.453,3.071) (2.528,3.072) (2.603,3.070) (2.678,3.064) (2.752,3.054) (2.826,3.041) (2.899,3.024) (2.971,3.004) (3.042,2.980) (3.112,2.951) (3.180,2.919) (3.245,2.883) (3.309,2.844) (3.370,2.800) (3.428,2.753) (3.483,2.702) (3.534,2.647) (3.582,2.589) (3.626,2.528) (3.665,2.465) (3.701,2.398) (3.731,2.330) (3.757,2.259) (3.777,2.187) (3.793,2.114) (3.804,2.040) (3.811,1.965) (3.812,1.890) (3.809,1.815) (3.801,1.741) (3.788,1.667) (3.772,1.594) (3.751,1.521) (3.727,1.451) (3.699,1.381) (3.667,1.313) (3.632,1.247) (3.594,1.182) (3.553,1.119) (3.509,1.058) (3.463,0.999) (3.415,0.942) (3.364,0.887) (3.311,0.834) (3.256,0.782) (3.200,0.733) (3.141,0.686) (3.082,0.641) (3.020,0.597) (2.958,0.556) (2.894,0.516) (2.829,0.479) (2.763,0.443) (2.697,0.409) (2.629,0.376) (2.560,0.346) (2.491,0.317) (2.445,0.299)}
\LC{(0.289,2.296) (0.308,2.368) (0.328,2.441) (0.348,2.513) (0.369,2.585) (0.390,2.657) (0.412,2.729) (0.434,2.800) (0.458,2.872) (0.481,2.943) (0.506,3.014) (0.531,3.084) (0.556,3.155) (0.582,3.225) (0.609,3.295) (0.636,3.365) (0.664,3.435) (0.693,3.504) (0.722,3.573) (0.752,3.642) (0.782,3.711) (0.813,3.779) (0.844,3.847) (0.876,3.915) (0.909,3.983) (0.942,4.050) (0.976,4.117) (1.010,4.183) (1.045,4.250) (1.081,4.316) (1.117,4.381) (1.154,4.447) (1.191,4.512) (1.229,4.577) (1.267,4.641) (1.306,4.705) (1.346,4.769) (1.386,4.832) (1.427,4.895) (1.468,4.958) (1.510,5.020) (1.553,5.082) (1.596,5.143) (1.639,5.204) (1.684,5.265) (1.728,5.325) (1.774,5.385) (1.819,5.444) (1.866,5.503) (1.913,5.561) (1.960,5.619) (2.009,5.677) (2.057,5.734) (2.106,5.791) (2.156,5.847) (2.206,5.902) (2.257,5.957) (2.309,6.012) (2.361,6.066) (2.413,6.120) (2.466,6.173) (2.520,6.225) (2.574,6.277) (2.628,6.329) (2.683,6.380) (2.739,6.430) (2.795,6.480) (2.852,6.529) (2.909,6.577) (2.967,6.625) (3.025,6.673) (3.083,6.719) (3.143,6.765) (3.202,6.811) (3.262,6.856) (3.323,6.900) (3.384,6.943) (3.446,6.986) (3.508,7.028) (3.571,7.069) (3.634,7.110) (3.697,7.150) (3.761,7.189) (3.825,7.228) (3.890,7.265) (3.956,7.302) (4.021,7.338) (4.087,7.374) (4.154,7.408) (4.221,7.442) (4.288,7.475) (4.356,7.507) (4.424,7.538) (4.493,7.568) (4.562,7.598) (4.631,7.627) (4.701,7.654) (4.771,7.681) (4.841,7.707) (4.912,7.732) (4.983,7.756) (5.054,7.779) (5.126,7.801) (5.198,7.822) (5.270,7.842) (5.343,7.861) (5.416,7.879) (5.489,7.896) (5.562,7.912) (5.636,7.927) (5.709,7.941) (5.783,7.954) (5.857,7.965) (5.932,7.976) (6.006,7.985) (6.081,7.993) (6.155,8.000) (6.230,8.006) (6.305,8.011) (6.380,8.014) (6.455,8.016) (6.530,8.017) (6.605,8.017) (6.680,8.015) (6.755,8.012) (6.830,8.008) (6.904,8.002) (6.979,7.995) (7.054,7.987) (7.128,7.978) (7.202,7.967) (7.276,7.954) (7.350,7.941) (7.423,7.926) (7.496,7.909) (7.569,7.891) (7.642,7.872) (7.714,7.851) (7.785,7.828) (7.856,7.805) (7.927,7.779) (7.997,7.753) (8.067,7.724) (8.136,7.695) (8.204,7.664) (8.271,7.631) (8.338,7.597) (8.404,7.561) (8.469,7.524) (8.533,7.485) (8.597,7.445) (8.659,7.404) (8.721,7.360) (8.781,7.316) (8.840,7.270) (8.898,7.223) (8.955,7.174) (9.011,7.124) (9.066,7.072) (9.119,7.019) (9.171,6.965) (9.221,6.910) (9.270,6.853) (9.318,6.795) (9.364,6.736) (9.408,6.675) (9.451,6.614) (9.492,6.551) (9.532,6.487) (9.570,6.423) (9.606,6.357) (9.641,6.290) (9.673,6.223) (9.704,6.155) (9.733,6.086) (9.761,6.016) (9.786,5.945) (9.810,5.874) (9.832,5.802) (9.852,5.730) (9.870,5.657) (9.886,5.584) (9.900,5.510) (9.912,5.436) (9.923,5.362) (9.932,5.288) (9.938,5.213) (9.943,5.138) (9.946,5.063) (9.948,4.988) (9.947,4.913) (9.945,4.838) (9.941,4.763) (9.935,4.689) (9.927,4.614) (9.918,4.539) (9.907,4.465) (9.894,4.391) (9.880,4.318) (9.864,4.244) (9.846,4.172) (9.827,4.099) (9.806,4.027) (9.784,3.955) (9.761,3.884) (9.736,3.813) (9.709,3.743) (9.681,3.674) (9.652,3.605) (9.621,3.536) (9.589,3.468) (9.556,3.401) (9.522,3.334) (9.486,3.268) (9.449,3.203) (9.411,3.138) (9.372,3.074) (9.332,3.011) (9.291,2.948) (9.248,2.887) (9.205,2.825) (9.161,2.765) (9.115,2.705) (9.069,2.646) (9.022,2.588) (8.973,2.531) (8.924,2.474) (8.874,2.418) (8.824,2.363) (8.772,2.308) (8.720,2.254) (8.667,2.201) (8.613,2.149) (8.558,2.098) (8.503,2.047) (8.447,1.997) (8.390,1.948) (8.333,1.900) (8.275,1.852) (8.217,1.805) (8.157,1.759) (8.098,1.714) (8.037,1.669) (7.977,1.625) (7.915,1.582) (7.853,1.540) (7.791,1.498) (7.728,1.457) (7.665,1.417) (7.601,1.377) (7.537,1.338) (7.472,1.300) (7.407,1.263) (7.342,1.226) (7.276,1.190) (7.210,1.155) (7.143,1.120) (7.077,1.086) (7.009,1.053) (6.942,1.021) (6.874,0.989) (6.806,0.958) (6.737,0.927) (6.668,0.897) (6.599,0.868) (6.530,0.839) (6.461,0.811) (6.391,0.784) (6.321,0.757) (6.250,0.731) (6.180,0.705) (6.109,0.680) (6.038,0.656) (5.967,0.632) (5.896,0.609) (5.824,0.586) (5.753,0.564) (5.681,0.542) (5.609,0.521) (5.537,0.501) (5.464,0.481) (5.392,0.462) (5.319,0.443) (5.247,0.425) (5.174,0.407) (5.101,0.389) (5.028,0.373) (4.954,0.356) (4.881,0.340) (4.808,0.325) (4.734,0.310) (4.661,0.296) (4.587,0.282) (4.513,0.268) (4.439,0.255) (4.365,0.242) (4.291,0.230) (4.217,0.218) (4.143,0.207) (4.069,0.196) (3.995,0.185) (3.920,0.175) (3.846,0.165) (3.772,0.156) (3.697,0.147) (3.623,0.138) (3.548,0.130) (3.474,0.122) (3.399,0.114) (3.324,0.107) (3.250,0.100) (3.200,0.095)}
\LC{(1.013,3.088) (1.059,3.147) (1.106,3.206) (1.154,3.264) (1.202,3.321) (1.252,3.377) (1.303,3.432) (1.355,3.486) (1.407,3.540) (1.461,3.592) (1.515,3.644) (1.570,3.695) (1.627,3.744) (1.684,3.793) (1.742,3.840) (1.801,3.887) (1.861,3.932) (1.921,3.976) (1.983,4.019) (2.045,4.061) (2.108,4.102) (2.172,4.141) (2.236,4.179) (2.302,4.216) (2.368,4.251) (2.435,4.285) (2.503,4.317) (2.571,4.348) (2.640,4.378) (2.709,4.406) (2.780,4.432) (2.850,4.457) (2.922,4.480) (2.994,4.501) (3.066,4.521) (3.139,4.539) (3.212,4.555) (3.286,4.569) (3.360,4.581) (3.434,4.591) (3.509,4.600) (3.583,4.606) (3.658,4.610) (3.733,4.612) (3.808,4.612) (3.883,4.610) (3.958,4.605) (4.033,4.598) (4.107,4.589) (4.181,4.577) (4.255,4.563) (4.328,4.547) (4.401,4.528) (4.472,4.506) (4.543,4.482) (4.614,4.455) (4.683,4.426) (4.751,4.394) (4.817,4.360) (4.882,4.323) (4.946,4.283) (5.008,4.241) (5.068,4.196) (5.127,4.149) (5.183,4.100) (5.237,4.048) (5.289,3.993) (5.338,3.937) (5.385,3.878) (5.429,3.818) (5.470,3.755) (5.508,3.690) (5.544,3.624) (5.576,3.557) (5.605,3.488) (5.631,3.417) (5.654,3.346) (5.674,3.273) (5.690,3.200) (5.703,3.126) (5.713,3.052) (5.719,2.977) (5.723,2.902) (5.723,2.827) (5.720,2.752) (5.714,2.678) (5.705,2.603) (5.693,2.529) (5.678,2.456) (5.660,2.383) (5.640,2.311) (5.617,2.239) (5.592,2.169) (5.564,2.099) (5.533,2.030) (5.501,1.963) (5.466,1.896) (5.430,1.831) (5.391,1.767) (5.350,1.704) (5.308,1.642) (5.263,1.581) (5.217,1.522) (5.170,1.464) (5.121,1.408) (5.070,1.352) (5.018,1.298) (4.965,1.246) (4.910,1.194) (4.854,1.144) (4.797,1.095) (4.739,1.048) (4.680,1.002) (4.620,0.957) (4.559,0.913) (4.497,0.871) (4.434,0.830) (4.370,0.790) (4.306,0.752) (4.241,0.715) (4.175,0.679) (4.153,0.667)}
\LC{(-10.000,1.452) (-9.943,1.420) (-9.875,1.389) (-9.807,1.358) (-9.738,1.328) (-9.669,1.298) (-9.600,1.269) (-9.531,1.240) (-9.461,1.211) (-9.392,1.183) (-9.322,1.156) (-9.252,1.129) (-9.182,1.102) (-9.112,1.076) (-9.042,1.050) (-8.971,1.024) (-8.900,0.999) (-8.829,0.974) (-8.759,0.950) (-8.687,0.926) (-8.616,0.903) (-8.545,0.880) (-8.473,0.857) (-8.402,0.835) (-8.330,0.813) (-8.258,0.791) (-8.186,0.770) (-8.114,0.749) (-8.042,0.729) (-7.970,0.709) (-7.897,0.689) (-7.825,0.670) (-7.752,0.651) (-7.680,0.633) (-7.607,0.614) (-7.534,0.597) (-7.461,0.579) (-7.388,0.562) (-7.315,0.545) (-7.242,0.529) (-7.169,0.512) (-7.095,0.497) (-7.022,0.481) (-6.948,0.466) (-6.875,0.451) (-6.801,0.437) (-6.728,0.422) (-6.654,0.408) (-6.580,0.395) (-6.506,0.382) (-6.433,0.369) (-6.359,0.356) (-6.285,0.343) (-6.211,0.331) (-6.137,0.320) (-6.062,0.308) (-5.988,0.297) (-5.914,0.286) (-5.840,0.275) (-5.766,0.265) (-5.691,0.254) (-5.617,0.245) (-5.543,0.235) (-5.468,0.225) (-5.394,0.216) (-5.319,0.207) (-5.245,0.199) (-5.170,0.190) (-5.096,0.182) (-5.021,0.174) (-4.946,0.167) (-4.872,0.159) (-4.797,0.152) (-4.723,0.145) (-4.698,0.143)}
\LC{(1.673,-5.217) (1.718,-5.277) (1.763,-5.337) (1.809,-5.396) (1.856,-5.455) (1.903,-5.514) (1.950,-5.571) (1.999,-5.629) (2.047,-5.686) (2.097,-5.742) (2.147,-5.798) (2.197,-5.854) (2.248,-5.909) (2.300,-5.963) (2.352,-6.017) (2.404,-6.071) (2.458,-6.123) (2.511,-6.176) (2.566,-6.228) (2.620,-6.279) (2.676,-6.329) (2.732,-6.380) (2.788,-6.429) (2.845,-6.478) (2.902,-6.526) (2.960,-6.574) (3.019,-6.621) (3.077,-6.667) (3.137,-6.713) (3.197,-6.758) (3.257,-6.803) (3.318,-6.846) (3.380,-6.889) (3.441,-6.932) (3.504,-6.973) (3.567,-7.014) (3.630,-7.055) (3.694,-7.094) (3.758,-7.133) (3.823,-7.171) (3.888,-7.208) (3.953,-7.244) (4.019,-7.280) (4.086,-7.315) (4.153,-7.349) (4.220,-7.382) (4.287,-7.414) (4.356,-7.446) (4.424,-7.477) (4.493,-7.506) (4.562,-7.535) (4.632,-7.563) (4.702,-7.590) (4.772,-7.616) (4.843,-7.641) (4.914,-7.666) (4.985,-7.689) (5.056,-7.711) (5.128,-7.732) (5.201,-7.753) (5.273,-7.772) (5.346,-7.790) (5.419,-7.807) (5.492,-7.823) (5.566,-7.838) (5.639,-7.852) (5.713,-7.865) (5.787,-7.877) (5.862,-7.887) (5.936,-7.897) (6.011,-7.905) (6.085,-7.912) (6.160,-7.918) (6.235,-7.922) (6.310,-7.926) (6.385,-7.928) (6.460,-7.929) (6.535,-7.928) (6.610,-7.927) (6.685,-7.924) (6.760,-7.919) (6.834,-7.914) (6.909,-7.907) (6.983,-7.898) (7.058,-7.889) (7.132,-7.877) (7.206,-7.865) (7.280,-7.851) (7.353,-7.836) (7.426,-7.819) (7.499,-7.801) (7.571,-7.781) (7.643,-7.760) (7.715,-7.737) (7.786,-7.713) (7.856,-7.687) (7.926,-7.660) (7.996,-7.632) (8.064,-7.602) (8.132,-7.570) (8.200,-7.537) (8.266,-7.502) (8.332,-7.466) (8.397,-7.429) (8.461,-7.389) (8.524,-7.349) (8.586,-7.307) (8.647,-7.263) (8.707,-7.218) (8.766,-7.172) (8.823,-7.124) (8.880,-7.074) (8.935,-7.024) (8.989,-6.971) (9.041,-6.918) (9.093,-6.863) (9.142,-6.807) (9.191,-6.750) (9.238,-6.691) (9.283,-6.631) (9.327,-6.570) (9.369,-6.508) (9.409,-6.445) (9.448,-6.381) (9.485,-6.316) (9.520,-6.249) (9.553,-6.182) (9.585,-6.114) (9.615,-6.045) (9.643,-5.976) (9.669,-5.906) (9.694,-5.835) (9.716,-5.763) (9.736,-5.691) (9.755,-5.618) (9.772,-5.545) (9.787,-5.472) (9.800,-5.398) (9.811,-5.324) (9.820,-5.249) (9.827,-5.175) (9.833,-5.100) (9.836,-5.025) (9.838,-4.950) (9.838,-4.875) (9.836,-4.800) (9.832,-4.725) (9.826,-4.650) (9.819,-4.576) (9.810,-4.501) (9.799,-4.427) (9.787,-4.353) (9.773,-4.279) (9.757,-4.206) (9.739,-4.133) (9.720,-4.060) (9.700,-3.988) (9.678,-3.917) (9.654,-3.845) (9.629,-3.775) (9.603,-3.705) (9.575,-3.635) (9.546,-3.566) (9.515,-3.497) (9.483,-3.429) (9.450,-3.362) (9.416,-3.296) (9.380,-3.230) (9.343,-3.164) (9.305,-3.100) (9.266,-3.036) (9.226,-2.973) (9.184,-2.910) (9.142,-2.848) (9.098,-2.787) (9.054,-2.727) (9.008,-2.667) (8.962,-2.608) (8.914,-2.550) (8.866,-2.493) (8.817,-2.436) (8.767,-2.380) (8.716,-2.325) (8.664,-2.271) (8.612,-2.217) (8.558,-2.165) (8.504,-2.113) (8.450,-2.061) (8.394,-2.011) (8.338,-1.961) (8.281,-1.912) (8.224,-1.864) (8.165,-1.817) (8.107,-1.770) (8.047,-1.724) (7.987,-1.679) (7.927,-1.635) (7.866,-1.591) (7.804,-1.548) (7.742,-1.506) (7.680,-1.465) (7.617,-1.424) (7.553,-1.384) (7.489,-1.345) (7.425,-1.306) (7.360,-1.269) (7.295,-1.232) (7.229,-1.195) (7.163,-1.160) (7.097,-1.125) (7.030,-1.091) (6.963,-1.057) (6.896,-1.024) (6.828,-0.992) (6.760,-0.960) (6.692,-0.929) (6.623,-0.899) (6.554,-0.870) (6.485,-0.841) (6.415,-0.813) (6.346,-0.785) (6.276,-0.758) (6.206,-0.731) (6.135,-0.706) (6.064,-0.681) (5.994,-0.656) (5.923,-0.632) (5.851,-0.609) (5.780,-0.586) (5.708,-0.564) (5.636,-0.542) (5.564,-0.521) (5.492,-0.500) (5.420,-0.480) (5.348,-0.461) (5.275,-0.442) (5.202,-0.423) (5.130,-0.405) (5.057,-0.388) (4.983,-0.371) (4.910,-0.355) (4.837,-0.339) (4.764,-0.323) (4.690,-0.308) (4.617,-0.294) (4.543,-0.280) (4.469,-0.266) (4.395,-0.253) (4.321,-0.240) (4.272,-0.232)}
\LC{(-0.492,3.275) (-0.514,3.347) (-0.537,3.418) (-0.561,3.490) (-0.585,3.561) (-0.609,3.632) (-0.634,3.702) (-0.659,3.773) (-0.685,3.843) (-0.711,3.913) (-0.738,3.984) (-0.766,4.053) (-0.793,4.123) (-0.822,4.192) (-0.850,4.262) (-0.880,4.331) (-0.910,4.400) (-0.940,4.468) (-0.971,4.537) (-1.002,4.605) (-1.034,4.673) (-1.066,4.740) (-1.099,4.808) (-1.132,4.875) (-1.166,4.942) (-1.200,5.009) (-1.234,5.075) (-1.270,5.142) (-1.305,5.208) (-1.341,5.273) (-1.378,5.339) (-1.415,5.404) (-1.452,5.469) (-1.491,5.534) (-1.529,5.598) (-1.568,5.662) (-1.608,5.726) (-1.647,5.789) (-1.688,5.852) (-1.729,5.915) (-1.770,5.978) (-1.812,6.040) (-1.854,6.102) (-1.897,6.164) (-1.940,6.225) (-1.984,6.286) (-2.028,6.346) (-2.073,6.407) (-2.118,6.466) (-2.164,6.526) (-2.210,6.585) (-2.256,6.644) (-2.303,6.702) (-2.351,6.761) (-2.399,6.818) (-2.447,6.876) (-2.496,6.932) (-2.545,6.989) (-2.595,7.045) (-2.645,7.101) (-2.696,7.156) (-2.747,7.211) (-2.799,7.266) (-2.851,7.320) (-2.903,7.373) (-2.956,7.426) (-3.009,7.479) (-3.063,7.531) (-3.117,7.583) (-3.172,7.635) (-3.227,7.686) (-3.282,7.736) (-3.338,7.786) (-3.395,7.835) (-3.452,7.884) (-3.509,7.933) (-3.566,7.981) (-3.624,8.028) (-3.683,8.075) (-3.742,8.122) (-3.801,8.168) (-3.861,8.213) (-3.921,8.258) (-3.982,8.302) (-4.042,8.346) (-4.104,8.389) (-4.166,8.432) (-4.228,8.474) (-4.290,8.515) (-4.353,8.556) (-4.416,8.596) (-4.480,8.636) (-4.544,8.675) (-4.608,8.713) (-4.673,8.751) (-4.738,8.788) (-4.804,8.825) (-4.870,8.861) (-4.936,8.896) (-5.003,8.930) (-5.069,8.964) (-5.137,8.997) (-5.204,9.030) (-5.272,9.062) (-5.340,9.093) (-5.409,9.123) (-5.478,9.153) (-5.547,9.182) (-5.617,9.210) (-5.686,9.237) (-5.757,9.264) (-5.827,9.290) (-5.898,9.315) (-5.968,9.340) (-6.040,9.363) (-6.111,9.386) (-6.183,9.408) (-6.255,9.429) (-6.327,9.450) (-6.399,9.469) (-6.472,9.488) (-6.545,9.505) (-6.618,9.522) (-6.691,9.538) (-6.765,9.554) (-6.838,9.568) (-6.912,9.581) (-6.986,9.594) (-7.060,9.605) (-7.134,9.616) (-7.209,9.625) (-7.283,9.634) (-7.358,9.642) (-7.433,9.649) (-7.507,9.654) (-7.582,9.659) (-7.657,9.663) (-7.732,9.666) (-7.807,9.667) (-7.882,9.668) (-7.957,9.668) (-8.032,9.666) (-8.107,9.664) (-8.182,9.660) (-8.257,9.656) (-8.331,9.650) (-8.406,9.643) (-8.481,9.635) (-8.555,9.626) (-8.630,9.616) (-8.704,9.605) (-8.778,9.592) (-8.851,9.578) (-8.925,9.564) (-8.998,9.548) (-9.071,9.530) (-9.144,9.512) (-9.216,9.493) (-9.288,9.472) (-9.360,9.450) (-9.431,9.427) (-9.502,9.402) (-9.573,9.377) (-9.643,9.350) (-9.712,9.322) (-9.781,9.292) (-9.850,9.262) (-9.918,9.230) (-9.985,9.197) (-10.000,9.186)}
\LC{(-10.000,-1.331) (-9.939,-1.302) (-9.870,-1.273) (-9.801,-1.245) (-9.731,-1.217) (-9.661,-1.190) (-9.591,-1.163) (-9.521,-1.136) (-9.451,-1.110) (-9.380,-1.084) (-9.310,-1.059) (-9.239,-1.034) (-9.168,-1.009) (-9.097,-0.985) (-9.026,-0.961) (-8.955,-0.938) (-8.883,-0.915) (-8.812,-0.892) (-8.740,-0.870) (-8.669,-0.848) (-8.597,-0.826) (-8.525,-0.805) (-8.453,-0.784) (-8.381,-0.764) (-8.308,-0.744) (-8.236,-0.724) (-8.164,-0.704) (-8.091,-0.685) (-8.019,-0.667) (-7.946,-0.648) (-7.873,-0.630) (-7.800,-0.612) (-7.727,-0.595) (-7.654,-0.578) (-7.581,-0.561) (-7.508,-0.545) (-7.435,-0.529) (-7.361,-0.513) (-7.288,-0.498) (-7.214,-0.483) (-7.141,-0.468) (-7.067,-0.453) (-6.994,-0.439) (-6.920,-0.425) (-6.846,-0.411) (-6.772,-0.398) (-6.699,-0.385) (-6.625,-0.372) (-6.551,-0.360) (-6.477,-0.348) (-6.403,-0.336) (-6.328,-0.324)}
\LC{(1.709,4.492) (1.761,4.547) (1.813,4.600) (1.866,4.653) (1.920,4.706) (1.975,4.757) (2.030,4.808) (2.086,4.858) (2.142,4.907) (2.199,4.956) (2.257,5.003) (2.316,5.050) (2.375,5.096) (2.435,5.141) (2.496,5.186) (2.557,5.229) (2.619,5.271) (2.681,5.313) (2.744,5.354) (2.808,5.393) (2.872,5.432) (2.937,5.469) (3.003,5.506) (3.069,5.542) (3.135,5.576) (3.202,5.609) (3.270,5.642) (3.338,5.673) (3.407,5.703) (3.476,5.732) (3.546,5.760) (3.616,5.786) (3.687,5.811) (3.758,5.835) (3.829,5.858) (3.901,5.879) (3.973,5.899) (4.046,5.918) (4.119,5.935) (4.192,5.951) (4.266,5.966) (4.340,5.979) (4.414,5.990) (4.488,6.000) (4.563,6.009) (4.637,6.016) (4.712,6.021) (4.787,6.025) (4.862,6.027) (4.937,6.028) (5.012,6.027) (5.087,6.024) (5.162,6.019) (5.236,6.013) (5.311,6.004) (5.385,5.995) (5.459,5.983) (5.533,5.969) (5.607,5.953) (5.679,5.936) (5.752,5.917) (5.824,5.895) (5.895,5.872) (5.966,5.847) (6.036,5.820) (6.105,5.791) (6.173,5.760) (6.240,5.727) (6.307,5.692) (6.372,5.655) (6.436,5.616) (6.499,5.575) (6.560,5.532) (6.620,5.487) (6.679,5.440) (6.736,5.391) (6.792,5.341) (6.845,5.289) (6.897,5.235) (6.947,5.179) (6.996,5.121) (7.042,5.062) (7.086,5.002) (7.128,4.940) (7.168,4.876) (7.205,4.811) (7.241,4.745) (7.274,4.678) (7.304,4.609) (7.333,4.540) (7.358,4.469) (7.382,4.398) (7.402,4.326) (7.421,4.253) (7.437,4.180) (7.450,4.106) (7.461,4.032) (7.469,3.957) (7.475,3.883) (7.478,3.808) (7.479,3.733) (7.478,3.658) (7.474,3.583) (7.468,3.508) (7.460,3.433) (7.449,3.359) (7.436,3.285) (7.421,3.212) (7.404,3.139) (7.384,3.066) (7.363,2.995) (7.340,2.923) (7.315,2.853) (7.287,2.783) (7.258,2.714) (7.228,2.645) (7.195,2.577) (7.161,2.511) (7.125,2.445) (7.088,2.380) (7.049,2.316) (7.009,2.252) (6.967,2.190) (6.924,2.129) (6.880,2.068) (6.834,2.009) (6.787,1.950) (6.739,1.893) (6.689,1.837) (6.639,1.781) (6.587,1.727) (6.534,1.673) (6.481,1.621) (6.426,1.569) (6.371,1.519) (6.314,1.470) (6.257,1.421) (6.199,1.374) (6.140,1.328) (6.080,1.282) (6.020,1.238) (5.958,1.194) (5.897,1.152) (5.834,1.111) (5.771,1.070) (5.707,1.031) (5.686,1.018)}
\LC{(0.123,-3.792) (0.128,-3.866) (0.133,-3.941) (0.138,-4.016) (0.144,-4.091) (0.149,-4.166) (0.154,-4.241) (0.160,-4.315) (0.165,-4.390) (0.171,-4.465) (0.177,-4.540) (0.183,-4.614) (0.189,-4.689) (0.195,-4.764) (0.201,-4.839) (0.207,-4.913) (0.214,-4.988) (0.220,-5.063) (0.227,-5.138) (0.233,-5.212) (0.240,-5.287) (0.247,-5.362) (0.254,-5.436) (0.261,-5.511) (0.268,-5.586) (0.275,-5.660) (0.283,-5.735) (0.290,-5.810) (0.298,-5.884) (0.305,-5.959) (0.313,-6.033) (0.321,-6.108) (0.329,-6.183) (0.337,-6.257) (0.345,-6.332) (0.353,-6.406) (0.361,-6.481) (0.370,-6.555) (0.378,-6.630) (0.387,-6.704) (0.395,-6.779) (0.404,-6.853) (0.413,-6.928) (0.422,-7.002) (0.431,-7.077) (0.440,-7.151) (0.449,-7.226) (0.459,-7.300) (0.468,-7.374) (0.478,-7.449) (0.487,-7.523) (0.497,-7.598) (0.507,-7.672) (0.517,-7.746) (0.527,-7.821) (0.537,-7.895) (0.547,-7.969) (0.558,-8.044) (0.568,-8.118) (0.579,-8.192) (0.589,-8.266) (0.600,-8.341) (0.611,-8.415) (0.621,-8.489) (0.632,-8.563) (0.644,-8.637) (0.655,-8.712) (0.666,-8.786) (0.677,-8.860) (0.689,-8.934) (0.700,-9.008) (0.712,-9.082) (0.724,-9.156) (0.736,-9.230) (0.748,-9.304) (0.760,-9.378) (0.772,-9.452) (0.784,-9.526) (0.796,-9.600) (0.809,-9.674) (0.821,-9.748) (0.834,-9.822) (0.847,-9.896) (0.859,-9.970) (0.868,-10.000)}
\LC{(-0.270,-4.625) (-0.278,-4.699) (-0.287,-4.774) (-0.296,-4.848) (-0.306,-4.923) (-0.315,-4.997) (-0.324,-5.071) (-0.334,-5.146) (-0.344,-5.220) (-0.354,-5.295) (-0.364,-5.369) (-0.374,-5.443) (-0.384,-5.517) (-0.395,-5.592) (-0.405,-5.666) (-0.416,-5.740) (-0.427,-5.814) (-0.438,-5.889) (-0.449,-5.963) (-0.461,-6.037) (-0.472,-6.111) (-0.484,-6.185) (-0.495,-6.259) (-0.507,-6.333) (-0.519,-6.407) (-0.532,-6.481) (-0.544,-6.555) (-0.556,-6.629) (-0.569,-6.703) (-0.582,-6.777) (-0.594,-6.851) (-0.608,-6.925) (-0.621,-6.999) (-0.634,-7.072) (-0.647,-7.146) (-0.661,-7.220) (-0.675,-7.294) (-0.689,-7.367) (-0.703,-7.441) (-0.717,-7.515) (-0.731,-7.588) (-0.745,-7.662) (-0.760,-7.736) (-0.775,-7.809) (-0.789,-7.883) (-0.804,-7.956) (-0.820,-8.030) (-0.835,-8.103) (-0.850,-8.176) (-0.866,-8.250) (-0.881,-8.323) (-0.897,-8.396) (-0.913,-8.470) (-0.929,-8.543) (-0.946,-8.616) (-0.962,-8.689) (-0.979,-8.763) (-0.995,-8.836) (-1.012,-8.909) (-1.029,-8.982) (-1.046,-9.055) (-1.063,-9.128) (-1.081,-9.201) (-1.098,-9.274) (-1.116,-9.347) (-1.134,-9.419) (-1.151,-9.492) (-1.170,-9.565) (-1.188,-9.638) (-1.206,-9.711) (-1.225,-9.783) (-1.243,-9.856) (-1.262,-9.929) (-1.281,-10.000)}
\LC{(1.386,-3.828) (1.435,-3.884) (1.486,-3.940) (1.537,-3.994) (1.589,-4.048) (1.642,-4.101) (1.696,-4.153) (1.751,-4.205) (1.806,-4.256) (1.862,-4.305) (1.919,-4.354) (1.977,-4.402) (2.035,-4.449) (2.094,-4.495) (2.154,-4.540) (2.215,-4.584) (2.276,-4.627) (2.339,-4.669) (2.401,-4.710) (2.465,-4.750) (2.529,-4.789) (2.594,-4.827) (2.659,-4.863) (2.726,-4.899) (2.792,-4.933) (2.860,-4.965) (2.928,-4.997) (2.997,-5.027) (3.066,-5.056) (3.135,-5.084) (3.206,-5.110) (3.276,-5.135) (3.348,-5.158) (3.419,-5.180) (3.492,-5.201) (3.564,-5.220) (3.637,-5.237) (3.710,-5.253) (3.784,-5.267) (3.858,-5.279) (3.932,-5.290) (4.007,-5.299) (4.081,-5.306) (4.156,-5.312) (4.231,-5.316) (4.306,-5.318) (4.381,-5.317) (4.456,-5.316) (4.531,-5.312) (4.606,-5.306) (4.680,-5.298) (4.755,-5.288) (4.829,-5.276) (4.902,-5.262) (4.975,-5.245) (5.048,-5.227) (5.120,-5.206) (5.192,-5.184) (5.262,-5.158) (5.332,-5.131) (5.401,-5.102) (5.469,-5.070) (5.536,-5.036) (5.602,-5.000) (5.666,-4.961) (5.729,-4.921) (5.791,-4.878) (5.851,-4.833) (5.909,-4.785) (5.965,-4.736) (6.020,-4.685) (6.072,-4.631) (6.123,-4.576) (6.171,-4.518) (6.217,-4.459) (6.261,-4.398) (6.302,-4.336) (6.341,-4.272) (6.378,-4.206) (6.411,-4.139) (6.442,-4.071) (6.470,-4.001) (6.496,-3.931) (6.518,-3.859) (6.538,-3.787) (6.555,-3.714) (6.569,-3.640) (6.581,-3.566) (6.589,-3.491) (6.595,-3.417) (6.598,-3.342) (6.598,-3.267) (6.596,-3.192) (6.590,-3.117) (6.582,-3.042) (6.572,-2.968) (6.559,-2.894) (6.544,-2.821) (6.526,-2.748) (6.506,-2.676) (6.484,-2.604) (6.459,-2.533) (6.432,-2.463) (6.403,-2.394) (6.373,-2.326) (6.340,-2.258) (6.305,-2.192) (6.269,-2.126) (6.230,-2.062) (6.191,-1.998) (6.149,-1.936) (6.106,-1.874) (6.061,-1.814) (6.015,-1.755) (5.968,-1.697) (5.919,-1.640) (5.869,-1.584) (5.817,-1.529) (5.765,-1.476) (5.711,-1.424) (5.656,-1.372) (5.600,-1.322) (5.543,-1.274) (5.486,-1.226) (5.427,-1.179) (5.367,-1.134) (5.307,-1.089) (5.245,-1.046) (5.183,-1.004) (5.121,-0.963) (5.057,-0.923) (4.993,-0.884) (4.928,-0.846) (4.863,-0.809)}
\LC{(-2.350,6.023) (-2.402,6.076) (-2.456,6.129) (-2.509,6.182) (-2.563,6.234) (-2.618,6.285) (-2.673,6.336) (-2.729,6.386) (-2.785,6.435) (-2.842,6.484) (-2.900,6.533) (-2.957,6.580) (-3.016,6.628) (-3.075,6.674) (-3.134,6.720) (-3.194,6.765) (-3.254,6.810) (-3.315,6.853) (-3.376,6.897) (-3.438,6.939) (-3.500,6.981) (-3.563,7.022) (-3.626,7.062) (-3.690,7.102) (-3.754,7.141) (-3.819,7.179) (-3.884,7.216) (-3.949,7.253) (-4.015,7.289) (-4.081,7.324) (-4.148,7.358) (-4.215,7.391) (-4.283,7.424) (-4.351,7.455) (-4.419,7.486) (-4.488,7.516) (-4.557,7.545) (-4.627,7.573) (-4.697,7.600) (-4.767,7.627) (-4.838,7.652) (-4.908,7.677) (-4.980,7.700) (-5.051,7.723) (-5.123,7.744) (-5.195,7.765) (-5.268,7.784) (-5.340,7.802) (-5.413,7.820) (-5.486,7.836) (-5.560,7.851) (-5.634,7.866) (-5.707,7.879) (-5.781,7.891) (-5.856,7.901) (-5.930,7.911) (-6.005,7.920) (-6.079,7.927) (-6.154,7.933) (-6.229,7.938) (-6.304,7.942) (-6.379,7.944) (-6.454,7.945) (-6.529,7.945) (-6.604,7.944) (-6.679,7.941) (-6.753,7.937) (-6.828,7.932) (-6.903,7.925) (-6.978,7.917) (-7.052,7.908) (-7.126,7.897) (-7.200,7.885) (-7.274,7.872) (-7.348,7.857) (-7.421,7.840) (-7.494,7.823) (-7.566,7.803) (-7.638,7.783) (-7.710,7.760) (-7.781,7.737) (-7.852,7.712) (-7.922,7.685) (-7.991,7.657) (-8.060,7.627) (-8.128,7.596) (-8.196,7.563) (-8.263,7.529) (-8.329,7.494) (-8.394,7.456) (-8.458,7.418) (-8.521,7.378) (-8.584,7.336) (-8.645,7.293) (-8.705,7.248) (-8.765,7.202) (-8.823,7.155) (-8.880,7.106) (-8.935,7.056) (-8.990,7.004) (-9.043,6.951) (-9.094,6.897) (-9.145,6.841) (-9.193,6.784) (-9.241,6.726) (-9.287,6.667) (-9.331,6.606) (-9.374,6.544) (-9.415,6.481) (-9.454,6.418) (-9.491,6.353) (-9.527,6.287) (-9.561,6.220) (-9.594,6.152) (-9.624,6.084) (-9.653,6.015) (-9.680,5.945) (-9.705,5.874) (-9.728,5.803) (-9.749,5.731) (-9.769,5.658) (-9.786,5.585) (-9.802,5.512) (-9.816,5.438) (-9.827,5.364) (-9.837,5.290) (-9.845,5.215) (-9.851,5.140) (-9.856,5.066) (-9.858,4.991) (-9.859,4.916) (-9.857,4.841) (-9.854,4.766) (-9.850,4.691) (-9.843,4.616) (-9.835,4.542) (-9.825,4.467) (-9.813,4.393) (-9.799,4.319) (-9.784,4.246) (-9.768,4.173) (-9.749,4.100) (-9.729,4.028) (-9.708,3.956) (-9.685,3.884) (-9.661,3.814) (-9.635,3.743) (-9.608,3.673) (-9.579,3.604) (-9.549,3.535) (-9.518,3.467) (-9.485,3.399) (-9.451,3.333) (-9.416,3.266) (-9.380,3.201) (-9.342,3.136) (-9.304,3.072) (-9.264,3.008) (-9.223,2.945) (-9.181,2.883) (-9.138,2.822) (-9.094,2.761) (-9.049,2.701) (-9.003,2.642) (-8.956,2.583) (-8.908,2.525) (-8.859,2.468) (-8.810,2.412) (-8.759,2.357) (-8.708,2.302) (-8.656,2.248) (-8.603,2.195) (-8.549,2.143) (-8.495,2.091) (-8.439,2.040) (-8.384,1.990) (-8.327,1.941) (-8.270,1.892) (-8.212,1.844) (-8.154,1.797) (-8.095,1.751) (-8.035,1.706) (-7.975,1.661) (-7.914,1.617) (-7.853,1.574) (-7.791,1.531) (-7.729,1.489) (-7.666,1.448) (-7.603,1.408) (-7.539,1.368) (-7.475,1.330) (-7.410,1.291) (-7.345,1.254) (-7.280,1.217) (-7.214,1.181) (-7.148,1.146) (-7.082,1.111) (-7.015,1.077) (-6.947,1.044) (-6.880,1.012) (-6.812,0.980) (-6.744,0.949) (-6.675,0.918) (-6.607,0.888) (-6.537,0.859) (-6.491,0.840)}
\LC{(10.000,-1.736) (9.939,-1.698) (9.874,-1.662) (9.808,-1.625) (9.742,-1.590) (9.676,-1.554) (9.609,-1.519) (9.543,-1.485) (9.476,-1.451) (9.408,-1.418) (9.341,-1.385) (9.273,-1.353) (9.205,-1.321) (9.137,-1.290) (9.069,-1.259) (9.000,-1.229) (8.931,-1.199) (8.862,-1.170) (8.793,-1.141) (8.724,-1.113) (8.654,-1.085) (8.584,-1.057) (8.514,-1.030) (8.444,-1.004) (8.374,-0.978) (8.303,-0.952) (8.233,-0.927) (8.162,-0.902) (8.091,-0.878) (8.020,-0.854) (7.949,-0.831) (7.877,-0.808) (7.806,-0.785) (7.734,-0.763) (7.662,-0.741) (7.590,-0.720) (7.518,-0.699) (7.446,-0.679) (7.374,-0.658) (7.302,-0.639) (7.229,-0.619) (7.156,-0.600) (7.084,-0.582) (7.011,-0.564) (6.938,-0.546)}
\LC{(4.710,-8.828) (4.775,-8.865) (4.841,-8.902) (4.907,-8.938) (4.973,-8.973) (5.039,-9.008) (5.106,-9.042) (5.173,-9.075) (5.241,-9.107) (5.309,-9.139) (5.377,-9.171) (5.445,-9.201) (5.514,-9.231) (5.583,-9.260) (5.653,-9.288) (5.723,-9.316) (5.793,-9.343) (5.863,-9.369) (5.933,-9.394) (6.004,-9.419) (6.075,-9.443) (6.147,-9.466) (6.218,-9.488) (6.290,-9.509) (6.362,-9.530) (6.435,-9.550) (6.507,-9.569) (6.580,-9.587) (6.653,-9.604) (6.726,-9.620) (6.800,-9.636) (6.873,-9.650) (6.947,-9.664) (7.021,-9.677) (7.095,-9.689) (7.169,-9.700) (7.243,-9.710) (7.318,-9.719) (7.392,-9.727) (7.467,-9.734) (7.542,-9.741) (7.617,-9.746) (7.692,-9.750) (7.766,-9.753) (7.841,-9.755) (7.916,-9.757) (7.991,-9.757) (8.066,-9.756) (8.141,-9.754) (8.216,-9.751) (8.291,-9.747) (8.366,-9.742) (8.441,-9.736) (8.515,-9.728) (8.590,-9.720) (8.664,-9.710) (8.739,-9.700) (8.813,-9.688) (8.887,-9.675) (8.960,-9.661) (9.034,-9.646) (9.107,-9.629) (9.180,-9.612) (9.252,-9.593) (9.325,-9.573) (9.397,-9.552) (9.468,-9.529) (9.539,-9.506) (9.610,-9.481) (9.681,-9.455) (9.750,-9.428) (9.820,-9.400) (9.889,-9.370) (9.957,-9.339) (10.000,-9.318)}
\LC{(2.598,-5.621) (2.657,-5.667) (2.718,-5.712) (2.778,-5.756) (2.840,-5.799) (2.901,-5.842) (2.964,-5.883) (3.027,-5.924) (3.090,-5.964) (3.154,-6.003) (3.219,-6.041) (3.284,-6.078) (3.350,-6.114) (3.416,-6.149) (3.483,-6.184) (3.550,-6.217) (3.618,-6.249) (3.686,-6.281) (3.754,-6.311) (3.823,-6.340) (3.893,-6.369) (3.963,-6.396) (4.033,-6.422) (4.104,-6.447) (4.175,-6.470) (4.247,-6.493) (4.318,-6.514) (4.391,-6.534) (4.463,-6.553) (4.536,-6.571) (4.609,-6.588) (4.683,-6.603) (4.756,-6.617) (4.830,-6.629) (4.904,-6.641) (4.979,-6.651) (5.053,-6.659) (5.128,-6.666) (5.203,-6.672) (5.278,-6.676) (5.353,-6.679) (5.428,-6.680) (5.503,-6.680) (5.578,-6.678) (5.652,-6.675) (5.727,-6.670) (5.802,-6.664) (5.877,-6.656) (5.951,-6.646) (6.025,-6.635) (6.099,-6.622) (6.172,-6.607) (6.246,-6.591) (6.318,-6.572) (6.391,-6.553) (6.463,-6.531) (6.534,-6.508) (6.605,-6.483) (6.675,-6.456) (6.744,-6.427) (6.812,-6.397) (6.880,-6.364) (6.947,-6.330) (7.013,-6.295) (7.078,-6.257) (7.142,-6.218) (7.204,-6.177) (7.266,-6.134) (7.326,-6.089) (7.385,-6.043) (7.443,-5.995) (7.499,-5.945) (7.554,-5.894) (7.607,-5.841) (7.658,-5.786) (7.708,-5.730) (7.756,-5.672) (7.802,-5.613) (7.846,-5.553) (7.889,-5.491) (7.929,-5.428) (7.968,-5.363) (8.004,-5.298) (8.038,-5.231) (8.071,-5.163) (8.101,-5.095) (8.129,-5.025) (8.154,-4.955) (8.178,-4.883) (8.199,-4.812) (8.218,-4.739) (8.235,-4.666) (8.249,-4.592) (8.261,-4.518) (8.271,-4.444) (8.279,-4.369) (8.285,-4.295) (8.288,-4.220) (8.289,-4.145) (8.288,-4.070) (8.285,-3.995) (8.279,-3.920) (8.272,-3.845) (8.262,-3.771) (8.251,-3.697) (8.237,-3.623) (8.222,-3.550) (8.205,-3.477) (8.185,-3.404) (8.164,-3.332) (8.141,-3.261) (8.117,-3.190) (8.090,-3.120) (8.062,-3.050) (8.033,-2.981) (8.002,-2.913) (7.969,-2.846) (7.935,-2.779) (7.899,-2.713) (7.862,-2.648) (7.823,-2.583) (7.783,-2.520) (7.742,-2.457) (7.700,-2.395) (7.656,-2.335) (7.611,-2.274) (7.565,-2.215) (7.518,-2.157) (7.469,-2.100) (7.420,-2.043) (7.370,-1.987) (7.319,-1.933) (7.266,-1.879) (7.213,-1.826) (7.159,-1.774) (7.104,-1.723) (7.048,-1.673) (6.992,-1.624) (6.934,-1.575) (6.876,-1.528) (6.857,-1.512)}
\LC{(-0.778,5.179) (-0.800,5.251) (-0.823,5.322) (-0.846,5.394) (-0.869,5.465) (-0.893,5.536) (-0.916,5.608) (-0.941,5.679) (-0.965,5.749) (-0.990,5.820) (-1.015,5.891) (-1.041,5.961) (-1.066,6.032) (-1.092,6.102) (-1.119,6.172) (-1.145,6.242) (-1.172,6.312) (-1.200,6.382) (-1.227,6.452) (-1.255,6.522) (-1.284,6.591) (-1.312,6.660) (-1.341,6.730) (-1.370,6.799) (-1.400,6.868) (-1.429,6.937) (-1.459,7.005) (-1.490,7.074) (-1.521,7.142) (-1.552,7.211) (-1.583,7.279) (-1.615,7.347) (-1.647,7.415) (-1.679,7.482) (-1.711,7.550) (-1.744,7.617) (-1.777,7.684) (-1.811,7.752) (-1.845,7.818) (-1.879,7.885) (-1.913,7.952) (-1.948,8.018) (-1.983,8.085) (-2.018,8.151) (-2.054,8.217) (-2.090,8.283) (-2.126,8.348) (-2.163,8.414) (-2.200,8.479) (-2.237,8.544) (-2.274,8.609) (-2.312,8.674) (-2.350,8.739) (-2.388,8.803) (-2.427,8.867) (-2.466,8.931) (-2.505,8.995) (-2.545,9.059) (-2.585,9.123) (-2.625,9.186) (-2.666,9.249) (-2.706,9.312) (-2.747,9.375) (-2.789,9.437) (-2.830,9.500) (-2.872,9.562) (-2.915,9.624) (-2.957,9.685) (-3.000,9.747) (-3.043,9.808) (-3.087,9.869) (-3.131,9.930) (-3.175,9.991) (-3.190,10.000)}
\LC{(-3.035,-7.115) (-3.091,-7.164) (-3.147,-7.213) (-3.204,-7.262) (-3.262,-7.310) (-3.320,-7.358) (-3.379,-7.405) (-3.437,-7.451) (-3.497,-7.497) (-3.557,-7.542) (-3.617,-7.587) (-3.678,-7.631) (-3.739,-7.674) (-3.800,-7.717) (-3.862,-7.759) (-3.925,-7.801) (-3.988,-7.841) (-4.051,-7.882) (-4.115,-7.921) (-4.179,-7.960) (-4.243,-7.999) (-4.308,-8.036) (-4.374,-8.073) (-4.439,-8.109) (-4.505,-8.145) (-4.572,-8.179) (-4.639,-8.213) (-4.706,-8.247) (-4.773,-8.279) (-4.841,-8.311) (-4.910,-8.342) (-4.978,-8.372) (-5.047,-8.401) (-5.117,-8.430) (-5.186,-8.458) (-5.256,-8.485) (-5.327,-8.511) (-5.397,-8.536) (-5.468,-8.561) (-5.539,-8.584) (-5.611,-8.607) (-5.683,-8.629) (-5.755,-8.650) (-5.827,-8.670) (-5.899,-8.689) (-5.972,-8.707) (-6.045,-8.724) (-6.118,-8.741) (-6.192,-8.756) (-6.265,-8.770) (-6.339,-8.784) (-6.413,-8.796) (-6.487,-8.808) (-6.562,-8.818) (-6.636,-8.827) (-6.710,-8.835) (-6.785,-8.843) (-6.860,-8.849) (-6.935,-8.854) (-7.010,-8.858) (-7.085,-8.861) (-7.160,-8.862) (-7.235,-8.863) (-7.310,-8.862) (-7.384,-8.861) (-7.459,-8.858) (-7.534,-8.854) (-7.609,-8.849) (-7.684,-8.842) (-7.758,-8.835) (-7.833,-8.826) (-7.907,-8.816) (-7.981,-8.804) (-8.055,-8.792) (-8.129,-8.778) (-8.203,-8.763) (-8.276,-8.746) (-8.349,-8.729) (-8.421,-8.710) (-8.493,-8.689) (-8.565,-8.668) (-8.637,-8.645) (-8.708,-8.621) (-8.778,-8.595) (-8.848,-8.568) (-8.917,-8.540) (-8.986,-8.510) (-9.055,-8.479) (-9.122,-8.447) (-9.189,-8.413) (-9.256,-8.378) (-9.321,-8.342) (-9.386,-8.304) (-9.450,-8.265) (-9.513,-8.225) (-9.576,-8.183) (-9.637,-8.140) (-9.698,-8.095) (-9.757,-8.050) (-9.816,-8.003) (-9.873,-7.955) (-9.929,-7.905) (-9.985,-7.854) (-10.000,-7.837)}
\LC{(-0.793,-5.153) (-0.816,-5.224) (-0.839,-5.296) (-0.863,-5.367) (-0.886,-5.438) (-0.910,-5.509) (-0.935,-5.580) (-0.960,-5.651) (-0.985,-5.721) (-1.010,-5.792) (-1.036,-5.862) (-1.062,-5.933) (-1.088,-6.003) (-1.115,-6.073) (-1.142,-6.143) (-1.169,-6.213) (-1.197,-6.283) (-1.225,-6.352) (-1.253,-6.422) (-1.282,-6.491) (-1.311,-6.560) (-1.340,-6.629) (-1.370,-6.698) (-1.400,-6.767) (-1.430,-6.835) (-1.460,-6.904) (-1.491,-6.972) (-1.522,-7.041) (-1.554,-7.109) (-1.586,-7.177) (-1.618,-7.244) (-1.650,-7.312) (-1.683,-7.379) (-1.716,-7.447) (-1.749,-7.514) (-1.783,-7.581) (-1.817,-7.648) (-1.851,-7.714) (-1.886,-7.781) (-1.921,-7.847) (-1.956,-7.914) (-1.992,-7.980) (-2.028,-8.045) (-2.064,-8.111) (-2.101,-8.177) (-2.137,-8.242) (-2.175,-8.307) (-2.212,-8.372) (-2.250,-8.437) (-2.288,-8.501) (-2.326,-8.566) (-2.365,-8.630) (-2.404,-8.694) (-2.443,-8.758) (-2.483,-8.822) (-2.523,-8.885) (-2.563,-8.948) (-2.604,-9.011) (-2.645,-9.074) (-2.686,-9.137) (-2.727,-9.199) (-2.769,-9.262) (-2.811,-9.324) (-2.854,-9.386) (-2.897,-9.447) (-2.940,-9.509) (-2.983,-9.570) (-3.027,-9.631) (-3.071,-9.692) (-3.115,-9.752) (-3.159,-9.812) (-3.204,-9.873) (-3.249,-9.932) (-3.295,-9.992) (-3.310,-10.000)}
\LC{(-1.651,4.550) (-1.701,4.606) (-1.751,4.662) (-1.802,4.717) (-1.854,4.771) (-1.906,4.824) (-1.960,4.877) (-2.013,4.930) (-2.068,4.981) (-2.123,5.032) (-2.179,5.082) (-2.235,5.131) (-2.293,5.180) (-2.350,5.228) (-2.409,5.275) (-2.468,5.321) (-2.527,5.366) (-2.588,5.411) (-2.649,5.455) (-2.710,5.498) (-2.772,5.540) (-2.835,5.581) (-2.898,5.621) (-2.962,5.660) (-3.027,5.698) (-3.092,5.736) (-3.157,5.772) (-3.224,5.808) (-3.290,5.842) (-3.357,5.875) (-3.425,5.907) (-3.493,5.939) (-3.562,5.969) (-3.631,5.998) (-3.701,6.026) (-3.771,6.052) (-3.842,6.078) (-3.912,6.102) (-3.984,6.125) (-4.056,6.147) (-4.128,6.168) (-4.200,6.187) (-4.273,6.205) (-4.346,6.222) (-4.419,6.237) (-4.493,6.251) (-4.567,6.263) (-4.641,6.275) (-4.716,6.284) (-4.790,6.293) (-4.865,6.299) (-4.940,6.305) (-5.015,6.308) (-5.090,6.310) (-5.165,6.311) (-5.240,6.310) (-5.315,6.307) (-5.389,6.303) (-5.464,6.297) (-5.539,6.289) (-5.613,6.280) (-5.687,6.268) (-5.761,6.255) (-5.835,6.241) (-5.908,6.224) (-5.981,6.206) (-6.053,6.186) (-6.125,6.164) (-6.196,6.140) (-6.266,6.114) (-6.336,6.086) (-6.405,6.057) (-6.473,6.026) (-6.540,5.992) (-6.606,5.957) (-6.671,5.920) (-6.736,5.881) (-6.798,5.840) (-6.860,5.797) (-6.920,5.753) (-6.979,5.706) (-7.037,5.658) (-7.093,5.608) (-7.147,5.557) (-7.200,5.503) (-7.251,5.448) (-7.300,5.392) (-7.347,5.333) (-7.392,5.273) (-7.436,5.212) (-7.477,5.150) (-7.516,5.086) (-7.553,5.020) (-7.588,4.954) (-7.620,4.886) (-7.651,4.818) (-7.679,4.748) (-7.704,4.678) (-7.728,4.607) (-7.749,4.535) (-7.767,4.462) (-7.783,4.389) (-7.797,4.315) (-7.809,4.241) (-7.818,4.166) (-7.824,4.092) (-7.829,4.017) (-7.830,3.942) (-7.830,3.867) (-7.828,3.792) (-7.823,3.717) (-7.816,3.642) (-7.806,3.568) (-7.795,3.494) (-7.781,3.420) (-7.766,3.347) (-7.748,3.274) (-7.729,3.201) (-7.707,3.130) (-7.684,3.058) (-7.658,2.988) (-7.631,2.918) (-7.603,2.848) (-7.572,2.780) (-7.540,2.712) (-7.506,2.645) (-7.471,2.579) (-7.434,2.514) (-7.396,2.449) (-7.356,2.386) (-7.315,2.323) (-7.272,2.261) (-7.229,2.200) (-7.184,2.140) (-7.137,2.081) (-7.090,2.023) (-7.041,1.966) (-6.992,1.910) (-6.941,1.854) (-6.889,1.800) (-6.836,1.747) (-6.783,1.695) (-6.728,1.643) (-6.673,1.593) (-6.616,1.543) (-6.559,1.495) (-6.501,1.447) (-6.442,1.401) (-6.383,1.355) (-6.322,1.310) (-6.261,1.267) (-6.200,1.224) (-6.138,1.182) (-6.075,1.141) (-6.011,1.101) (-5.947,1.062) (-5.883,1.024) (-5.818,0.987) (-5.752,0.950) (-5.686,0.915) (-5.620,0.880) (-5.553,0.846) (-5.485,0.813) (-5.417,0.781)}
\LC{(10.000,-0.013) (9.943,-0.013) (9.868,-0.013) (9.793,-0.012) (9.718,-0.012) (9.643,-0.012) (9.568,-0.012) (9.493,-0.011) (9.418,-0.011) (9.343,-0.011) (9.268,-0.011) (9.193,-0.010) (9.118,-0.010) (9.043,-0.010) (8.968,-0.010) (8.893,-0.009) (8.818,-0.009) (8.743,-0.009) (8.668,-0.009) (8.593,-0.008) (8.518,-0.008) (8.443,-0.008) (8.368,-0.008) (8.293,-0.008) (8.218,-0.007) (8.143,-0.007) (8.068,-0.007) (7.993,-0.007) (7.918,-0.007) (7.843,-0.006) (7.768,-0.006) (7.693,-0.006) (7.618,-0.006) (7.543,-0.006) (7.468,-0.006) (7.393,-0.005) (7.318,-0.005) (7.243,-0.005) (7.168,-0.005) (7.093,-0.005) (7.018,-0.005) (6.943,-0.004) (6.868,-0.004) (6.793,-0.004) (6.718,-0.004) (6.643,-0.004) (6.568,-0.004) (6.493,-0.004) (6.418,-0.004) (6.343,-0.003) (6.268,-0.003) (6.193,-0.003) (6.118,-0.003) (6.043,-0.003) (5.968,-0.003) (5.893,-0.003) (5.818,-0.003) (5.743,-0.003) (5.668,-0.002) (5.593,-0.002) (5.518,-0.002) (5.443,-0.002) (5.418,-0.002)}
\LC{(-0.001,2.928) (-0.001,3.003) (-0.001,3.078) (-0.001,3.153) (-0.001,3.228) (-0.002,3.303) (-0.002,3.378) (-0.002,3.453) (-0.002,3.528) (-0.002,3.603) (-0.002,3.678) (-0.002,3.753) (-0.002,3.828) (-0.002,3.903) (-0.002,3.978) (-0.002,4.053) (-0.002,4.128) (-0.002,4.203) (-0.003,4.278) (-0.003,4.353) (-0.003,4.428) (-0.003,4.503) (-0.003,4.578) (-0.003,4.653) (-0.003,4.728) (-0.003,4.803) (-0.003,4.878) (-0.003,4.953) (-0.004,5.028) (-0.004,5.103) (-0.004,5.178) (-0.004,5.253) (-0.004,5.328) (-0.004,5.403) (-0.004,5.478) (-0.004,5.553) (-0.004,5.628) (-0.005,5.703) (-0.005,5.778) (-0.005,5.853) (-0.005,5.928) (-0.005,6.003) (-0.005,6.078) (-0.005,6.153) (-0.005,6.228) (-0.006,6.303) (-0.006,6.378) (-0.006,6.453) (-0.006,6.528) (-0.006,6.603) (-0.006,6.678) (-0.006,6.753) (-0.007,6.828) (-0.007,6.903) (-0.007,6.978) (-0.007,7.053) (-0.007,7.128) (-0.007,7.203) (-0.007,7.278) (-0.008,7.353) (-0.008,7.428) (-0.008,7.503) (-0.008,7.578) (-0.008,7.653) (-0.008,7.728) (-0.009,7.803) (-0.009,7.878) (-0.009,7.953) (-0.009,8.028) (-0.009,8.103) (-0.009,8.178) (-0.010,8.253) (-0.010,8.328) (-0.010,8.403) (-0.010,8.478) (-0.010,8.553) (-0.010,8.628) (-0.011,8.703) (-0.011,8.778) (-0.011,8.853) (-0.011,8.928) (-0.011,9.003) (-0.012,9.078) (-0.012,9.153) (-0.012,9.228) (-0.012,9.303) (-0.012,9.378) (-0.013,9.453) (-0.013,9.528) (-0.013,9.603) (-0.013,9.678) (-0.013,9.753) (-0.014,9.828) (-0.014,9.903) (-0.014,9.978) (-0.014,10.000)}
\LC{(4.007,7.839) (4.071,7.879) (4.135,7.919) (4.199,7.957) (4.264,7.995) (4.329,8.032) (4.394,8.069) (4.460,8.105) (4.527,8.140) (4.593,8.174) (4.660,8.208) (4.728,8.241) (4.795,8.273) (4.864,8.304) (4.932,8.335) (5.001,8.365) (5.070,8.394) (5.140,8.422) (5.209,8.449) (5.280,8.476) (5.350,8.501) (5.421,8.526) (5.492,8.550) (5.563,8.573) (5.635,8.596) (5.707,8.617) (5.779,8.637) (5.851,8.657) (5.924,8.675) (5.997,8.693) (6.070,8.710) (6.143,8.726) (6.217,8.740) (6.291,8.754) (6.364,8.767) (6.439,8.779) (6.513,8.789) (6.587,8.799) (6.662,8.808) (6.736,8.816) (6.811,8.822) (6.886,8.828) (6.961,8.832) (7.036,8.835) (7.110,8.837) (7.185,8.838) (7.260,8.838) (7.335,8.837) (7.410,8.835) (7.485,8.831) (7.560,8.826) (7.635,8.820) (7.710,8.813) (7.784,8.805) (7.859,8.795) (7.933,8.784) (8.007,8.772) (8.081,8.759) (8.154,8.744) (8.227,8.728) (8.300,8.711) (8.373,8.692) (8.445,8.672) (8.517,8.651) (8.589,8.629) (8.660,8.605) (8.731,8.580) (8.801,8.553) (8.870,8.525) (8.939,8.496) (9.008,8.466) (9.076,8.434) (9.143,8.401) (9.210,8.366) (9.275,8.330) (9.341,8.293) (9.405,8.254) (9.468,8.214) (9.531,8.173) (9.593,8.130) (9.653,8.086) (9.713,8.041) (9.772,7.994) (9.830,7.946) (9.886,7.897) (9.942,7.847) (9.996,7.795) (10.000,7.778)}
\LC{(-0.891,-2.699) (-0.937,-2.758) (-0.984,-2.816) (-1.033,-2.873) (-1.082,-2.930) (-1.133,-2.985) (-1.184,-3.039) (-1.237,-3.093) (-1.291,-3.145) (-1.346,-3.196) (-1.402,-3.246) (-1.459,-3.295) (-1.517,-3.342) (-1.576,-3.388) (-1.636,-3.434) (-1.697,-3.477) (-1.759,-3.520) (-1.822,-3.560) (-1.885,-3.600) (-1.950,-3.638) (-2.016,-3.674) (-2.082,-3.709) (-2.149,-3.743) (-2.217,-3.774) (-2.286,-3.804) (-2.356,-3.832) (-2.426,-3.858) (-2.497,-3.883) (-2.568,-3.905) (-2.640,-3.926) (-2.713,-3.944) (-2.786,-3.961) (-2.860,-3.975) (-2.934,-3.987) (-3.008,-3.997) (-3.083,-4.004) (-3.158,-4.009) (-3.233,-4.012) (-3.308,-4.012) (-3.383,-4.010) (-3.457,-4.005) (-3.532,-3.997) (-3.606,-3.987) (-3.680,-3.974) (-3.753,-3.958) (-3.826,-3.939) (-3.898,-3.918) (-3.969,-3.893) (-4.038,-3.866) (-4.107,-3.835) (-4.174,-3.802) (-4.240,-3.765) (-4.303,-3.726) (-4.365,-3.684) (-4.425,-3.638) (-4.483,-3.590) (-4.538,-3.539) (-4.590,-3.486) (-4.640,-3.430) (-4.687,-3.371) (-4.730,-3.310) (-4.771,-3.247) (-4.808,-3.182) (-4.842,-3.115) (-4.872,-3.046) (-4.898,-2.976) (-4.921,-2.904) (-4.940,-2.832) (-4.955,-2.758) (-4.966,-2.684) (-4.974,-2.610) (-4.978,-2.535) (-4.979,-2.460) (-4.975,-2.385) (-4.969,-2.310) (-4.959,-2.236) (-4.945,-2.162) (-4.928,-2.089) (-4.909,-2.017) (-4.886,-1.945) (-4.860,-1.875) (-4.831,-1.805) (-4.800,-1.737) (-4.767,-1.670) (-4.730,-1.605) (-4.692,-1.540) (-4.651,-1.477) (-4.608,-1.416) (-4.563,-1.356) (-4.517,-1.297) (-4.468,-1.240) (-4.418,-1.184) (-4.366,-1.130) (-4.313,-1.077) (-4.258,-1.026) (-4.201,-0.977) (-4.144,-0.929) (-4.085,-0.882) (-4.025,-0.837) (-3.964,-0.793) (-3.902,-0.751) (-3.839,-0.711) (-3.775,-0.671) (-3.710,-0.634) (-3.667,-0.609)}
\LC{(1.347,2.119) (1.418,2.141) (1.491,2.160) (1.564,2.175) (1.638,2.187) (1.713,2.194) (1.788,2.196) (1.863,2.194) (1.938,2.187) (2.012,2.176) (2.085,2.159) (2.156,2.137) (2.226,2.109) (2.294,2.076) (2.358,2.038) (2.419,1.994) (2.476,1.945) (2.528,1.891) (2.575,1.832) (2.616,1.770) (2.651,1.703) (2.679,1.634) (2.701,1.562) (2.716,1.489) (2.723,1.414) (2.725,1.339) (2.719,1.264) (2.708,1.190) (2.691,1.117) (2.667,1.046) (2.639,0.977) (2.606,0.909) (2.582,0.866)}
\LC{(-1.963,-4.791) (-2.018,-4.843) (-2.073,-4.893) (-2.129,-4.943) (-2.186,-4.992) (-2.243,-5.041) (-2.302,-5.088) (-2.360,-5.135) (-2.420,-5.181) (-2.479,-5.226) (-2.540,-5.270) (-2.601,-5.313) (-2.663,-5.356) (-2.726,-5.397) (-2.789,-5.438) (-2.853,-5.477) (-2.917,-5.516) (-2.982,-5.553) (-3.047,-5.590) (-3.113,-5.626) (-3.180,-5.660) (-3.247,-5.694) (-3.315,-5.726) (-3.383,-5.757) (-3.451,-5.787) (-3.521,-5.816) (-3.590,-5.844) (-3.660,-5.871) (-3.731,-5.896) (-3.802,-5.921) (-3.873,-5.944) (-3.945,-5.965) (-4.017,-5.986) (-4.090,-6.005) (-4.163,-6.022) (-4.236,-6.039) (-4.309,-6.054) (-4.383,-6.067) (-4.457,-6.079) (-4.531,-6.090) (-4.606,-6.099) (-4.680,-6.106) (-4.755,-6.112) (-4.830,-6.117) (-4.905,-6.119) (-4.980,-6.121) (-5.055,-6.120) (-5.130,-6.118) (-5.205,-6.114) (-5.280,-6.109) (-5.354,-6.102) (-5.429,-6.092) (-5.503,-6.082) (-5.577,-6.069) (-5.650,-6.054) (-5.724,-6.038) (-5.796,-6.020) (-5.869,-6.000) (-5.940,-5.978) (-6.011,-5.954) (-6.082,-5.928) (-6.152,-5.900) (-6.220,-5.871) (-6.288,-5.839) (-6.355,-5.805) (-6.421,-5.770) (-6.486,-5.732) (-6.550,-5.693) (-6.613,-5.651) (-6.674,-5.608) (-6.734,-5.563) (-6.792,-5.516) (-6.849,-5.467) (-6.904,-5.416) (-6.958,-5.363) (-7.009,-5.309) (-7.059,-5.253) (-7.107,-5.196) (-7.153,-5.136) (-7.197,-5.076) (-7.239,-5.013) (-7.279,-4.950) (-7.317,-4.885) (-7.352,-4.819) (-7.385,-4.752) (-7.416,-4.683) (-7.444,-4.614) (-7.470,-4.543) (-7.493,-4.472) (-7.515,-4.400) (-7.533,-4.327) (-7.549,-4.254) (-7.563,-4.180) (-7.574,-4.106) (-7.583,-4.032) (-7.589,-3.957) (-7.593,-3.882) (-7.595,-3.807) (-7.594,-3.732) (-7.591,-3.657) (-7.585,-3.582) (-7.577,-3.508) (-7.567,-3.434) (-7.555,-3.360) (-7.541,-3.286) (-7.524,-3.213) (-7.505,-3.140) (-7.485,-3.068) (-7.462,-2.996) (-7.438,-2.926) (-7.412,-2.855) (-7.383,-2.786) (-7.353,-2.717) (-7.322,-2.649) (-7.288,-2.582) (-7.254,-2.516) (-7.217,-2.450) (-7.179,-2.385) (-7.139,-2.322) (-7.098,-2.259) (-7.056,-2.197) (-7.012,-2.136) (-6.967,-2.076) (-6.921,-2.017) (-6.874,-1.959) (-6.825,-1.902) (-6.775,-1.846) (-6.724,-1.791) (-6.672,-1.737) (-6.619,-1.683) (-6.566,-1.631) (-6.511,-1.580) (-6.455,-1.530) (-6.398,-1.481) (-6.341,-1.433) (-6.282,-1.386) (-6.223,-1.339) (-6.163,-1.294) (-6.103,-1.250) (-6.042,-1.207) (-5.980,-1.164) (-5.959,-1.151)}
\LC{(-3.628,-6.337) (-3.695,-6.369) (-3.764,-6.400) (-3.832,-6.430) (-3.901,-6.460) (-3.971,-6.488) (-4.041,-6.515) (-4.111,-6.541) (-4.182,-6.565) (-4.253,-6.589) (-4.325,-6.612) (-4.397,-6.633) (-4.469,-6.653) (-4.541,-6.672) (-4.614,-6.690) (-4.687,-6.707) (-4.761,-6.722) (-4.834,-6.736) (-4.908,-6.749) (-4.983,-6.760) (-5.057,-6.770) (-5.131,-6.779) (-5.206,-6.787) (-5.281,-6.793) (-5.356,-6.797) (-5.430,-6.800) (-5.505,-6.802) (-5.580,-6.802) (-5.655,-6.801) (-5.730,-6.798) (-5.805,-6.794) (-5.880,-6.788) (-5.955,-6.781) (-6.029,-6.771) (-6.103,-6.761) (-6.177,-6.749) (-6.251,-6.735) (-6.324,-6.719) (-6.397,-6.702) (-6.470,-6.683) (-6.542,-6.662) (-6.614,-6.640) (-6.685,-6.616) (-6.755,-6.590) (-6.825,-6.562) (-6.894,-6.533) (-6.962,-6.502) (-7.029,-6.469) (-7.096,-6.434) (-7.162,-6.398) (-7.226,-6.360) (-7.290,-6.320) (-7.352,-6.278) (-7.413,-6.235) (-7.473,-6.190) (-7.532,-6.143) (-7.589,-6.095) (-7.645,-6.045) (-7.700,-5.993) (-7.752,-5.940) (-7.804,-5.885) (-7.853,-5.829) (-7.901,-5.771) (-7.947,-5.712) (-7.991,-5.652) (-8.034,-5.590) (-8.074,-5.527) (-8.113,-5.462) (-8.149,-5.397) (-8.184,-5.330) (-8.216,-5.262) (-8.246,-5.194) (-8.274,-5.124) (-8.300,-5.054) (-8.324,-4.983) (-8.346,-4.911) (-8.365,-4.838) (-8.382,-4.765) (-8.397,-4.692) (-8.410,-4.618) (-8.420,-4.544) (-8.428,-4.469) (-8.435,-4.394) (-8.439,-4.320) (-8.440,-4.245) (-8.440,-4.170) (-8.437,-4.095) (-8.433,-4.020) (-8.426,-3.945) (-8.418,-3.871) (-8.407,-3.796) (-8.394,-3.722) (-8.380,-3.649) (-8.363,-3.576) (-8.345,-3.503) (-8.325,-3.431) (-8.303,-3.359) (-8.279,-3.288) (-8.254,-3.217) (-8.227,-3.147) (-8.198,-3.078) (-8.168,-3.009) (-8.136,-2.941) (-8.103,-2.874) (-8.068,-2.808) (-8.032,-2.742) (-7.994,-2.677) (-7.955,-2.613) (-7.915,-2.550) (-7.873,-2.487) (-7.831,-2.426) (-7.787,-2.365) (-7.742,-2.305) (-7.695,-2.246) (-7.648,-2.188) (-7.600,-2.130) (-7.550,-2.074) (-7.500,-2.019) (-7.448,-1.964) (-7.431,-1.946)}
\LC{(3.506,6.475) (3.571,6.511) (3.638,6.546) (3.704,6.580) (3.772,6.614) (3.839,6.646) (3.907,6.678) (3.976,6.708) (4.045,6.738) (4.114,6.767) (4.184,6.794) (4.254,6.821) (4.324,6.846) (4.395,6.871) (4.467,6.894) (4.538,6.916) (4.610,6.937) (4.682,6.957) (4.755,6.976) (4.828,6.994) (4.901,7.010) (4.974,7.026) (5.048,7.040) (5.122,7.053) (5.196,7.064) (5.270,7.075) (5.345,7.084) (5.419,7.092) (5.494,7.098) (5.569,7.103) (5.644,7.107) (5.719,7.109) (5.794,7.110) (5.869,7.110) (5.944,7.108) (6.019,7.104) (6.094,7.099) (6.168,7.093) (6.243,7.085) (6.317,7.076) (6.391,7.065) (6.465,7.052) (6.539,7.038) (6.612,7.023) (6.685,7.005) (6.758,6.986) (6.830,6.966) (6.902,6.944) (6.973,6.920) (7.043,6.894) (7.113,6.867) (7.183,6.839) (7.251,6.808) (7.319,6.776) (7.386,6.742) (7.452,6.706) (7.517,6.669) (7.581,6.630) (7.644,6.590) (7.706,6.548) (7.767,6.504) (7.827,6.458) (7.885,6.411) (7.942,6.363) (7.998,6.312) (8.052,6.261) (8.105,6.207) (8.156,6.152) (8.206,6.096) (8.254,6.039) (8.300,5.980) (8.344,5.919) (8.387,5.858) (8.428,5.795) (8.467,5.731) (8.504,5.666) (8.540,5.599) (8.573,5.532) (8.604,5.464) (8.633,5.395) (8.660,5.325) (8.685,5.254) (8.708,5.183) (8.729,5.111) (8.748,5.038) (8.765,4.965) (8.779,4.892) (8.792,4.818) (8.802,4.743) (8.810,4.669) (8.816,4.594) (8.820,4.519) (8.822,4.444) (8.822,4.369) (8.820,4.294) (8.816,4.219) (8.810,4.145) (8.801,4.070) (8.791,3.996) (8.780,3.922) (8.766,3.848) (8.750,3.775) (8.733,3.702) (8.714,3.629) (8.693,3.557) (8.670,3.486) (8.646,3.415) (8.620,3.344) (8.593,3.274) (8.564,3.205) (8.533,3.137) (8.501,3.069) (8.468,3.002) (8.433,2.935) (8.397,2.870) (8.359,2.805) (8.321,2.740) (8.281,2.677) (8.239,2.614) (8.197,2.552) (8.153,2.491) (8.109,2.431) (8.063,2.372) (8.016,2.313) (7.968,2.256) (7.919,2.199) (7.869,2.143) (7.818,2.088) (7.767,2.033)}
\LC{(-1.864,3.540) (-1.929,3.577) (-1.995,3.614) (-2.061,3.648) (-2.129,3.681) (-2.197,3.712) (-2.266,3.741) (-2.336,3.768) (-2.406,3.794) (-2.477,3.818) (-2.549,3.839) (-2.622,3.859) (-2.695,3.877) (-2.768,3.892) (-2.842,3.905) (-2.916,3.916) (-2.991,3.925) (-3.065,3.931) (-3.140,3.935) (-3.215,3.936) (-3.290,3.934) (-3.365,3.930) (-3.440,3.924) (-3.514,3.914) (-3.588,3.902) (-3.662,3.887) (-3.734,3.869) (-3.806,3.848) (-3.877,3.824) (-3.947,3.797) (-4.016,3.767) (-4.084,3.734) (-4.149,3.698) (-4.213,3.659) (-4.276,3.617) (-4.336,3.572) (-4.393,3.525) (-4.449,3.474) (-4.501,3.420) (-4.551,3.364) (-4.598,3.306) (-4.641,3.245) (-4.682,3.181) (-4.718,3.116) (-4.752,3.049) (-4.781,2.980) (-4.807,2.910) (-4.830,2.838) (-4.848,2.765) (-4.863,2.692) (-4.873,2.618) (-4.880,2.543) (-4.883,2.468) (-4.883,2.393) (-4.879,2.318) (-4.871,2.244) (-4.860,2.169) (-4.845,2.096) (-4.827,2.023) (-4.806,1.951) (-4.782,1.880) (-4.755,1.810) (-4.726,1.741) (-4.693,1.673) (-4.658,1.607) (-4.621,1.542) (-4.581,1.478) (-4.540,1.416) (-4.496,1.355) (-4.450,1.296) (-4.402,1.238) (-4.352,1.182) (-4.301,1.127)}
\LC{(2.522,3.823) (2.594,3.843) (2.667,3.861) (2.740,3.877) (2.814,3.891) (2.888,3.902) (2.962,3.911) (3.037,3.918) (3.112,3.923) (3.187,3.924) (3.262,3.924) (3.337,3.920) (3.412,3.914) (3.486,3.906) (3.560,3.894) (3.634,3.880) (3.707,3.862) (3.779,3.842) (3.850,3.819) (3.921,3.793) (3.990,3.763) (4.057,3.731) (4.124,3.696) (4.188,3.657) (4.250,3.616) (4.311,3.572) (4.369,3.524) (4.425,3.474) (4.478,3.421) (4.528,3.366) (4.576,3.308) (4.620,3.247) (4.661,3.184) (4.698,3.119) (4.732,3.052) (4.763,2.984) (4.789,2.914) (4.812,2.842) (4.831,2.770) (4.846,2.696) (4.858,2.622) (4.865,2.548) (4.869,2.473) (4.869,2.398) (4.865,2.323) (4.858,2.248) (4.847,2.174) (4.833,2.100) (4.816,2.027) (4.795,1.955) (4.772,1.884) (4.745,1.814) (4.716,1.745) (4.684,1.677) (4.649,1.610) (4.612,1.545) (4.599,1.524)}
\LC{(-2.405,-4.586) (-2.470,-4.624) (-2.535,-4.661) (-2.601,-4.697) (-2.668,-4.731) (-2.735,-4.764) (-2.803,-4.796) (-2.871,-4.827) (-2.940,-4.856) (-3.010,-4.883) (-3.080,-4.910) (-3.151,-4.934) (-3.223,-4.958) (-3.294,-4.979) (-3.367,-5.000) (-3.439,-5.018) (-3.512,-5.035) (-3.586,-5.050) (-3.659,-5.064) (-3.733,-5.076) (-3.808,-5.086) (-3.882,-5.094) (-3.957,-5.101) (-4.032,-5.105) (-4.107,-5.108) (-4.182,-5.109) (-4.257,-5.108) (-4.332,-5.104) (-4.407,-5.099) (-4.481,-5.091) (-4.556,-5.082) (-4.630,-5.070) (-4.703,-5.056) (-4.777,-5.040) (-4.849,-5.022) (-4.921,-5.001) (-4.993,-4.978) (-5.063,-4.953) (-5.133,-4.925) (-5.202,-4.895) (-5.270,-4.863) (-5.336,-4.829) (-5.402,-4.792) (-5.465,-4.753) (-5.528,-4.711) (-5.589,-4.667) (-5.648,-4.621) (-5.705,-4.573) (-5.761,-4.522) (-5.814,-4.470) (-5.865,-4.415) (-5.915,-4.358) (-5.961,-4.300) (-6.006,-4.239) (-6.047,-4.177) (-6.087,-4.113) (-6.123,-4.048) (-6.157,-3.981) (-6.188,-3.912) (-6.216,-3.843) (-6.242,-3.772) (-6.264,-3.701) (-6.284,-3.628) (-6.300,-3.555) (-6.314,-3.481) (-6.325,-3.407) (-6.332,-3.333) (-6.337,-3.258) (-6.339,-3.183) (-6.338,-3.108) (-6.334,-3.033) (-6.328,-2.958) (-6.319,-2.884) (-6.307,-2.810) (-6.292,-2.736) (-6.276,-2.663) (-6.256,-2.591) (-6.234,-2.519) (-6.210,-2.448) (-6.184,-2.378) (-6.155,-2.308) (-6.125,-2.240) (-6.092,-2.172) (-6.058,-2.106) (-6.021,-2.040) (-5.983,-1.976) (-5.943,-1.912) (-5.901,-1.850) (-5.857,-1.789) (-5.812,-1.729) (-5.766,-1.670) (-5.718,-1.612) (-5.669,-1.556) (-5.618,-1.500) (-5.566,-1.446) (-5.513,-1.393) (-5.459,-1.342) (-5.403,-1.291) (-5.347,-1.242) (-5.290,-1.193) (-5.231,-1.146) (-5.172,-1.101) (-5.111,-1.056) (-5.050,-1.013) (-4.988,-0.970) (-4.926,-0.929) (-4.862,-0.889) (-4.798,-0.850) (-4.733,-0.812) (-4.668,-0.776)}
\LC{(-3.747,6.524) (-3.815,6.556) (-3.883,6.587) (-3.952,6.617) (-4.021,6.646) (-4.091,6.674) (-4.161,6.701) (-4.231,6.727) (-4.302,6.751) (-4.373,6.775) (-4.445,6.798) (-4.517,6.819) (-4.589,6.840) (-4.661,6.859) (-4.734,6.877) (-4.807,6.893) (-4.880,6.909) (-4.954,6.923) (-5.028,6.936) (-5.102,6.948) (-5.176,6.959) (-5.251,6.968) (-5.325,6.976) (-5.400,6.982) (-5.475,6.987) (-5.550,6.991) (-5.625,6.994) (-5.700,6.995) (-5.775,6.994) (-5.850,6.992) (-5.925,6.989) (-5.999,6.984) (-6.074,6.977) (-6.149,6.969) (-6.223,6.960) (-6.297,6.948) (-6.371,6.936) (-6.445,6.921) (-6.518,6.905) (-6.591,6.888) (-6.663,6.869) (-6.736,6.848) (-6.807,6.825) (-6.878,6.801) (-6.948,6.775) (-7.018,6.747) (-7.087,6.718) (-7.155,6.687) (-7.223,6.654) (-7.290,6.620) (-7.355,6.584) (-7.420,6.546) (-7.484,6.506) (-7.546,6.465) (-7.608,6.422) (-7.668,6.378) (-7.727,6.331) (-7.785,6.283) (-7.841,6.234) (-7.896,6.183) (-7.950,6.130) (-8.002,6.076) (-8.052,6.021) (-8.101,5.964) (-8.148,5.905) (-8.193,5.845) (-8.236,5.784) (-8.278,5.722) (-8.317,5.658) (-8.355,5.593) (-8.391,5.527) (-8.425,5.460) (-8.457,5.392) (-8.486,5.324) (-8.514,5.254) (-8.539,5.183) (-8.563,5.112) (-8.584,5.040) (-8.603,4.968) (-8.620,4.895) (-8.635,4.821) (-8.648,4.747) (-8.658,4.673) (-8.666,4.598) (-8.673,4.524) (-8.677,4.449) (-8.679,4.374) (-8.679,4.299) (-8.676,4.224) (-8.672,4.149) (-8.666,4.074) (-8.658,4.000) (-8.648,3.925) (-8.636,3.851) (-8.622,3.778) (-8.606,3.704) (-8.588,3.631) (-8.569,3.559) (-8.548,3.487) (-8.525,3.416) (-8.500,3.345) (-8.474,3.274) (-8.446,3.205) (-8.417,3.136) (-8.386,3.067) (-8.354,3.000) (-8.320,2.933) (-8.285,2.867) (-8.248,2.801) (-8.210,2.736) (-8.171,2.672) (-8.130,2.609) (-8.089,2.547) (-8.046,2.486) (-8.002,2.425) (-7.956,2.365) (-7.910,2.306) (-7.863,2.248) (-7.814,2.191) (-7.765,2.134) (-7.714,2.079) (-7.663,2.024) (-7.646,2.006)}
\LC{(-1.421,-5.598) (-1.458,-5.663) (-1.495,-5.729) (-1.532,-5.794) (-1.570,-5.858) (-1.608,-5.923) (-1.647,-5.987) (-1.686,-6.051) (-1.726,-6.115) (-1.766,-6.178) (-1.807,-6.241) (-1.848,-6.304) (-1.889,-6.366) (-1.931,-6.429) (-1.973,-6.490) (-2.016,-6.552) (-2.059,-6.613) (-2.103,-6.674) (-2.147,-6.735) (-2.192,-6.795) (-2.237,-6.855) (-2.282,-6.915) (-2.328,-6.974) (-2.374,-7.033) (-2.421,-7.092) (-2.468,-7.150) (-2.516,-7.208) (-2.564,-7.266) (-2.613,-7.323) (-2.661,-7.380) (-2.711,-7.436) (-2.761,-7.492) (-2.811,-7.548) (-2.861,-7.603) (-2.912,-7.658) (-2.964,-7.713) (-3.016,-7.767) (-3.068,-7.821) (-3.121,-7.874) (-3.174,-7.927) (-3.228,-7.980) (-3.281,-8.032) (-3.336,-8.083) (-3.391,-8.135) (-3.446,-8.185) (-3.501,-8.236) (-3.557,-8.286) (-3.614,-8.335) (-3.671,-8.384) (-3.728,-8.432) (-3.786,-8.480) (-3.844,-8.528) (-3.902,-8.575) (-3.961,-8.622) (-4.020,-8.668) (-4.080,-8.713) (-4.139,-8.758) (-4.200,-8.803) (-4.261,-8.847) (-4.322,-8.890) (-4.383,-8.933) (-4.445,-8.976) (-4.507,-9.018) (-4.570,-9.059) (-4.633,-9.100) (-4.696,-9.140) (-4.760,-9.180) (-4.824,-9.219) (-4.888,-9.257) (-4.953,-9.295) (-5.018,-9.333) (-5.083,-9.369) (-5.149,-9.405) (-5.215,-9.441) (-5.281,-9.476) (-5.348,-9.510) (-5.415,-9.544) (-5.482,-9.577) (-5.550,-9.609) (-5.618,-9.641) (-5.686,-9.672) (-5.755,-9.702) (-5.824,-9.732) (-5.893,-9.761) (-5.962,-9.789) (-6.032,-9.817) (-6.102,-9.843) (-6.172,-9.870) (-6.243,-9.895) (-6.314,-9.920) (-6.385,-9.944) (-6.456,-9.967) (-6.528,-9.989) (-6.576,-10.000)}
\LC{(1.060,5.164) (1.090,5.232) (1.121,5.301) (1.151,5.369) (1.183,5.437) (1.214,5.505) (1.247,5.573) (1.279,5.641) (1.312,5.708) (1.345,5.775) (1.379,5.842) (1.413,5.909) (1.448,5.975) (1.483,6.042) (1.518,6.108) (1.554,6.174) (1.590,6.239) (1.627,6.305) (1.664,6.370) (1.701,6.435) (1.739,6.500) (1.778,6.564) (1.816,6.629) (1.855,6.693) (1.895,6.756) (1.935,6.820) (1.975,6.883) (2.016,6.946) (2.057,7.009) (2.098,7.071) (2.140,7.134) (2.182,7.196) (2.225,7.257) (2.268,7.319) (2.312,7.380) (2.355,7.441) (2.400,7.501) (2.444,7.561) (2.489,7.621) (2.535,7.681) (2.581,7.740) (2.627,7.799) (2.674,7.858) (2.721,7.916) (2.768,7.975) (2.816,8.032) (2.864,8.090) (2.913,8.147) (2.962,8.204) (3.011,8.260) (3.061,8.316) (3.111,8.372) (3.162,8.428) (3.212,8.483) (3.264,8.537) (3.315,8.592) (3.367,8.646) (3.420,8.699) (3.473,8.753) (3.526,8.806) (3.579,8.858) (3.633,8.910) (3.687,8.962) (3.742,9.013) (3.797,9.064) (3.853,9.115) (3.908,9.165) (3.964,9.215) (4.021,9.264) (4.078,9.313) (4.135,9.362) (4.192,9.410) (4.250,9.457) (4.309,9.505) (4.367,9.551) (4.426,9.598) (4.486,9.644) (4.545,9.689) (4.605,9.734) (4.666,9.779) (4.726,9.823) (4.787,9.866) (4.849,9.909) (4.910,9.952) (4.972,9.994) (4.993,10.000)}
\LC{(-4.211,7.912) (-4.276,7.950) (-4.342,7.986) (-4.408,8.022) (-4.474,8.057) (-4.540,8.092) (-4.607,8.126) (-4.675,8.159) (-4.743,8.191) (-4.811,8.222) (-4.879,8.253) (-4.948,8.283) (-5.017,8.312) (-5.086,8.340) (-5.156,8.367) (-5.226,8.394) (-5.297,8.420) (-5.368,8.445) (-5.439,8.469) (-5.510,8.492) (-5.582,8.514) (-5.654,8.535) (-5.726,8.556) (-5.798,8.575) (-5.871,8.594) (-5.944,8.611) (-6.017,8.628) (-6.090,8.644) (-6.164,8.659) (-6.237,8.672) (-6.311,8.685) (-6.385,8.697) (-6.460,8.707) (-6.534,8.717) (-6.609,8.725) (-6.683,8.733) (-6.758,8.739) (-6.833,8.745) (-6.908,8.749) (-6.983,8.752) (-7.058,8.754) (-7.133,8.755) (-7.208,8.754) (-7.283,8.753) (-7.357,8.750) (-7.432,8.746) (-7.507,8.741) (-7.582,8.735) (-7.657,8.727) (-7.731,8.719) (-7.805,8.709) (-7.880,8.697) (-7.953,8.685) (-8.027,8.671) (-8.101,8.656) (-8.174,8.640) (-8.247,8.622) (-8.319,8.603) (-8.391,8.583) (-8.463,8.561) (-8.535,8.538) (-8.606,8.514) (-8.676,8.488) (-8.746,8.461) (-8.815,8.432) (-8.884,8.403) (-8.952,8.371) (-9.020,8.339) (-9.087,8.305) (-9.153,8.270) (-9.219,8.233) (-9.283,8.195) (-9.347,8.156) (-9.410,8.115) (-9.472,8.073) (-9.534,8.030) (-9.594,7.985) (-9.653,7.939) (-9.711,7.892) (-9.768,7.844) (-9.824,7.794) (-9.879,7.743) (-9.933,7.690) (-9.985,7.637) (-10.000,7.618)}
\LC{(4.811,-7.150) (4.883,-7.169) (4.956,-7.187) (5.029,-7.203) (5.103,-7.219) (5.176,-7.233) (5.250,-7.246) (5.324,-7.258) (5.398,-7.269) (5.473,-7.278) (5.547,-7.286) (5.622,-7.293) (5.697,-7.299) (5.772,-7.303) (5.847,-7.306) (5.922,-7.307) (5.997,-7.307) (6.072,-7.306) (6.146,-7.303) (6.221,-7.299) (6.296,-7.294) (6.371,-7.287) (6.445,-7.278) (6.520,-7.268) (6.594,-7.257) (6.668,-7.244) (6.741,-7.229) (6.815,-7.213) (6.887,-7.195) (6.960,-7.176) (7.032,-7.155) (7.104,-7.133) (7.175,-7.109) (7.245,-7.083) (7.315,-7.056) (7.384,-7.027) (7.453,-6.997) (7.520,-6.965) (7.587,-6.931) (7.653,-6.895) (7.719,-6.858) (7.783,-6.820) (7.846,-6.779) (7.908,-6.737) (7.969,-6.694) (8.029,-6.649) (8.088,-6.602) (8.145,-6.554) (8.202,-6.504) (8.256,-6.453) (8.310,-6.400) (8.361,-6.346) (8.412,-6.290) (8.460,-6.233) (8.507,-6.175) (8.553,-6.115) (8.597,-6.054) (8.639,-5.992) (8.679,-5.929) (8.717,-5.864) (8.753,-5.799) (8.788,-5.732) (8.821,-5.664) (8.851,-5.596) (8.880,-5.527) (8.907,-5.457) (8.931,-5.386) (8.954,-5.314) (8.974,-5.242) (8.993,-5.169) (9.009,-5.096) (9.024,-5.023) (9.036,-4.949) (9.046,-4.874) (9.054,-4.800) (9.060,-4.725) (9.065,-4.650) (9.067,-4.575) (9.067,-4.500) (9.065,-4.425) (9.061,-4.350) (9.055,-4.276) (9.048,-4.201) (9.038,-4.127) (9.027,-4.052) (9.014,-3.979) (8.999,-3.905) (8.982,-3.832) (8.963,-3.759) (8.943,-3.687) (8.921,-3.615) (8.898,-3.544) (8.873,-3.473) (8.846,-3.403) (8.818,-3.334) (8.788,-3.265) (8.757,-3.197) (8.725,-3.129) (8.691,-3.062) (8.656,-2.996) (8.619,-2.930) (8.581,-2.866) (8.555,-2.823)}
\LC{(-0.995,2.331) (-1.051,2.381) (-1.109,2.428) (-1.169,2.474) (-1.230,2.518) (-1.292,2.560) (-1.355,2.600) (-1.420,2.638) (-1.486,2.674) (-1.553,2.707) (-1.621,2.739) (-1.690,2.767) (-1.760,2.794) (-1.831,2.818) (-1.903,2.839) (-1.976,2.857) (-2.049,2.873) (-2.123,2.885) (-2.198,2.895) (-2.272,2.901) (-2.347,2.904) (-2.422,2.903) (-2.497,2.899) (-2.572,2.891) (-2.646,2.879) (-2.719,2.864) (-2.792,2.844) (-2.863,2.821) (-2.933,2.793) (-3.001,2.762) (-3.067,2.726) (-3.130,2.686) (-3.191,2.642) (-3.249,2.594) (-3.303,2.542) (-3.354,2.487) (-3.400,2.428) (-3.442,2.366) (-3.480,2.302) (-3.513,2.234) (-3.541,2.165) (-3.564,2.093) (-3.581,2.020) (-3.594,1.946) (-3.601,1.872) (-3.603,1.797) (-3.600,1.722) (-3.593,1.647) (-3.580,1.573) (-3.564,1.500) (-3.543,1.428) (-3.517,1.358) (-3.488,1.288) (-3.456,1.221) (-3.420,1.155) (-3.381,1.091) (-3.338,1.029) (-3.293,0.969) (-3.246,0.911) (-3.196,0.855) (-3.143,0.802) (-3.089,0.750) (-3.033,0.700) (-2.975,0.653) (-2.915,0.607) (-2.854,0.564) (-2.791,0.523) (-2.728,0.483) (-2.662,0.446) (-2.596,0.410) (-2.574,0.399)}
\LC{(-10.000,-0.013) (-9.942,-0.013) (-9.867,-0.013) (-9.792,-0.012) (-9.717,-0.012) (-9.642,-0.012) (-9.567,-0.012) (-9.492,-0.011) (-9.417,-0.011) (-9.342,-0.011) (-9.267,-0.010) (-9.192,-0.010) (-9.117,-0.010) (-9.042,-0.010) (-8.967,-0.010) (-8.892,-0.009) (-8.817,-0.009) (-8.742,-0.009) (-8.667,-0.009) (-8.592,-0.008) (-8.517,-0.008) (-8.442,-0.008) (-8.367,-0.008) (-8.292,-0.008) (-8.217,-0.007) (-8.142,-0.007) (-8.067,-0.007) (-7.992,-0.007) (-7.917,-0.007) (-7.842,-0.006) (-7.767,-0.006) (-7.692,-0.006) (-7.617,-0.006) (-7.542,-0.006) (-7.467,-0.005) (-7.392,-0.005) (-7.317,-0.005) (-7.242,-0.005) (-7.167,-0.005) (-7.142,-0.005)}
\LC{(1.556,-6.813) (1.589,-6.880) (1.622,-6.947) (1.656,-7.014) (1.690,-7.081) (1.725,-7.147) (1.760,-7.214) (1.795,-7.280) (1.830,-7.346) (1.866,-7.412) (1.902,-7.478) (1.939,-7.543) (1.976,-7.608) (2.013,-7.673) (2.051,-7.738) (2.089,-7.803) (2.127,-7.868) (2.165,-7.932) (2.204,-7.996) (2.244,-8.060) (2.283,-8.124) (2.323,-8.187) (2.363,-8.250) (2.404,-8.314) (2.445,-8.376) (2.486,-8.439) (2.528,-8.501) (2.570,-8.564) (2.612,-8.626) (2.654,-8.687) (2.697,-8.749) (2.741,-8.810) (2.784,-8.871) (2.828,-8.932) (2.872,-8.992) (2.917,-9.053) (2.962,-9.113) (3.007,-9.173) (3.053,-9.232) (3.099,-9.292) (3.145,-9.351) (3.191,-9.409) (3.238,-9.468) (3.286,-9.526) (3.333,-9.584) (3.381,-9.642) (3.429,-9.699) (3.478,-9.757) (3.526,-9.814) (3.576,-9.870) (3.625,-9.927) (3.675,-9.983) (3.691,-10.000)}
\LC{(0.769,5.621) (0.789,5.693) (0.810,5.765) (0.831,5.837) (0.852,5.909) (0.873,5.981) (0.895,6.053) (0.917,6.125) (0.939,6.196) (0.961,6.268) (0.984,6.339) (1.007,6.411) (1.031,6.482) (1.054,6.553) (1.078,6.624) (1.102,6.695) (1.126,6.766) (1.151,6.837) (1.176,6.908) (1.201,6.979) (1.226,7.049) (1.252,7.120) (1.278,7.190) (1.304,7.260) (1.330,7.331) (1.357,7.401) (1.384,7.471) (1.411,7.541) (1.439,7.610) (1.466,7.680) (1.494,7.750) (1.523,7.819) (1.551,7.888) (1.580,7.958) (1.609,8.027) (1.638,8.096) (1.668,8.165) (1.698,8.234) (1.728,8.302) (1.758,8.371) (1.789,8.439) (1.820,8.508) (1.851,8.576) (1.882,8.644) (1.914,8.712) (1.946,8.780) (1.978,8.848) (2.010,8.915) (2.043,8.983) (2.076,9.050) (2.109,9.117) (2.143,9.184) (2.176,9.251) (2.210,9.318) (2.245,9.385) (2.279,9.452) (2.314,9.518) (2.349,9.584) (2.384,9.651) (2.420,9.717) (2.456,9.783) (2.492,9.848) (2.528,9.914) (2.565,9.979) (2.577,10.000)}
\LC{(3.537,-7.525) (3.597,-7.570) (3.658,-7.614) (3.719,-7.658) (3.780,-7.701) (3.842,-7.743) (3.904,-7.785) (3.967,-7.826) (4.030,-7.867) (4.094,-7.907) (4.158,-7.946) (4.222,-7.984) (4.287,-8.022) (4.352,-8.059) (4.418,-8.095) (4.484,-8.131) (4.550,-8.166) (4.617,-8.200) (4.684,-8.234) (4.751,-8.266) (4.819,-8.298) (4.887,-8.330) (4.956,-8.360) (5.025,-8.390) (5.094,-8.419) (5.164,-8.447) (5.234,-8.474) (5.304,-8.500) (5.374,-8.526) (5.445,-8.551) (5.516,-8.574) (5.588,-8.597) (5.659,-8.619) (5.731,-8.641) (5.803,-8.661) (5.876,-8.680) (5.948,-8.699) (6.021,-8.716) (6.095,-8.733) (6.168,-8.749) (6.241,-8.763) (6.315,-8.777) (6.389,-8.790) (6.463,-8.801) (6.537,-8.812) (6.612,-8.821) (6.686,-8.830) (6.761,-8.838) (6.836,-8.844) (6.911,-8.849) (6.985,-8.854) (7.060,-8.857) (7.135,-8.859) (7.210,-8.860) (7.285,-8.860) (7.360,-8.858) (7.435,-8.856) (7.510,-8.852) (7.585,-8.847) (7.660,-8.841) (7.734,-8.834) (7.809,-8.825) (7.883,-8.816) (7.957,-8.805) (8.031,-8.792) (8.105,-8.779) (8.179,-8.764) (8.252,-8.748) (8.325,-8.731) (8.398,-8.712) (8.470,-8.692) (8.542,-8.671) (8.613,-8.648) (8.684,-8.624) (8.755,-8.599) (8.825,-8.573) (8.895,-8.545) (8.964,-8.515) (9.032,-8.485) (9.100,-8.453) (9.167,-8.420) (9.234,-8.385) (9.300,-8.349) (9.365,-8.312) (9.429,-8.273) (9.493,-8.233) (9.555,-8.192) (9.617,-8.149) (9.678,-8.105) (9.737,-8.060) (9.796,-8.013) (9.854,-7.965) (9.910,-7.916) (9.966,-7.866) (10.000,-7.831)}
\LC{(0.740,-5.756) (0.760,-5.829) (0.779,-5.901) (0.799,-5.974) (0.819,-6.046) (0.839,-6.118) (0.859,-6.190) (0.880,-6.262) (0.901,-6.334) (0.922,-6.406) (0.943,-6.478) (0.965,-6.550) (0.987,-6.622) (1.009,-6.694) (1.031,-6.765) (1.054,-6.837) (1.076,-6.908) (1.099,-6.979) (1.123,-7.051) (1.146,-7.122) (1.170,-7.193) (1.194,-7.264) (1.218,-7.335) (1.243,-7.406) (1.268,-7.477) (1.293,-7.547) (1.318,-7.618) (1.344,-7.689) (1.369,-7.759) (1.395,-7.829) (1.422,-7.900) (1.448,-7.970) (1.475,-8.040) (1.502,-8.110) (1.529,-8.180) (1.556,-8.250) (1.584,-8.319) (1.612,-8.389) (1.640,-8.458) (1.668,-8.528) (1.697,-8.597) (1.726,-8.666) (1.755,-8.736) (1.784,-8.805) (1.814,-8.873) (1.844,-8.942) (1.874,-9.011) (1.904,-9.080) (1.935,-9.148) (1.966,-9.216) (1.997,-9.285) (2.028,-9.353) (2.060,-9.421) (2.091,-9.489) (2.123,-9.557) (2.156,-9.624) (2.188,-9.692) (2.221,-9.759) (2.254,-9.827) (2.287,-9.894) (2.320,-9.961) (2.343,-10.000)}
\LC{(0.443,6.926) (0.453,7.000) (0.463,7.075) (0.472,7.149) (0.482,7.223) (0.492,7.298) (0.503,7.372) (0.513,7.446) (0.523,7.521) (0.534,7.595) (0.544,7.669) (0.555,7.743) (0.566,7.818) (0.576,7.892) (0.587,7.966) (0.598,8.040) (0.610,8.114) (0.621,8.189) (0.632,8.263) (0.644,8.337) (0.655,8.411) (0.667,8.485) (0.679,8.559) (0.691,8.633) (0.703,8.707) (0.715,8.781) (0.727,8.855) (0.739,8.929) (0.752,9.003) (0.764,9.077) (0.777,9.151) (0.790,9.225) (0.802,9.299) (0.815,9.373) (0.828,9.447) (0.841,9.520) (0.855,9.594) (0.868,9.668) (0.881,9.742) (0.895,9.816) (0.909,9.889) (0.922,9.963) (0.932,10.000)}
\LC{(-0.456,6.301) (-0.467,6.376) (-0.478,6.450) (-0.489,6.524) (-0.500,6.598) (-0.511,6.672) (-0.523,6.746) (-0.535,6.820) (-0.546,6.894) (-0.558,6.969) (-0.570,7.043) (-0.582,7.117) (-0.595,7.191) (-0.607,7.265) (-0.620,7.338) (-0.632,7.412) (-0.645,7.486) (-0.658,7.560) (-0.671,7.634) (-0.684,7.708) (-0.697,7.782) (-0.711,7.855) (-0.724,7.929) (-0.738,8.003) (-0.752,8.077) (-0.766,8.150) (-0.780,8.224) (-0.794,8.298) (-0.808,8.371) (-0.823,8.445) (-0.837,8.518) (-0.852,8.592) (-0.867,8.666) (-0.882,8.739) (-0.897,8.812) (-0.912,8.886) (-0.928,8.959) (-0.943,9.033) (-0.959,9.106) (-0.974,9.179) (-0.990,9.253) (-1.006,9.326) (-1.022,9.399) (-1.039,9.472) (-1.055,9.546) (-1.071,9.619) (-1.088,9.692) (-1.105,9.765) (-1.122,9.838) (-1.139,9.911) (-1.156,9.984) (-1.162,10.000)}
\LC{(-5.902,-9.274) (-5.973,-9.298) (-6.045,-9.322) (-6.116,-9.344) (-6.188,-9.365) (-6.260,-9.386) (-6.333,-9.406) (-6.405,-9.425) (-6.478,-9.443) (-6.551,-9.460) (-6.624,-9.477) (-6.698,-9.492) (-6.771,-9.507) (-6.845,-9.521) (-6.919,-9.533) (-6.993,-9.545) (-7.067,-9.556) (-7.141,-9.566) (-7.216,-9.575) (-7.290,-9.583) (-7.365,-9.590) (-7.440,-9.596) (-7.515,-9.601) (-7.590,-9.605) (-7.664,-9.608) (-7.739,-9.611) (-7.814,-9.612) (-7.889,-9.611) (-7.964,-9.610) (-8.039,-9.608) (-8.114,-9.605) (-8.189,-9.601) (-8.264,-9.595) (-8.339,-9.589) (-8.413,-9.581) (-8.488,-9.572) (-8.562,-9.562) (-8.636,-9.551) (-8.710,-9.539) (-8.784,-9.525) (-8.858,-9.511) (-8.931,-9.495) (-9.004,-9.478) (-9.077,-9.460) (-9.149,-9.441) (-9.221,-9.420) (-9.293,-9.398) (-9.365,-9.375) (-9.436,-9.351) (-9.506,-9.326) (-9.576,-9.299) (-9.646,-9.271) (-9.715,-9.242) (-9.784,-9.212) (-9.852,-9.180) (-9.919,-9.147) (-9.986,-9.113) (-10.000,-9.102)}
\LC{(0.919,-1.901) (0.981,-1.943) (1.045,-1.983) (1.110,-2.020) (1.177,-2.054) (1.245,-2.085) (1.315,-2.113) (1.386,-2.138) (1.458,-2.159) (1.530,-2.176) (1.604,-2.189) (1.679,-2.199) (1.754,-2.204) (1.829,-2.204) (1.903,-2.200) (1.978,-2.191) (2.051,-2.177) (2.124,-2.158) (2.195,-2.133) (2.264,-2.103) (2.330,-2.068) (2.393,-2.028) (2.452,-1.982) (2.507,-1.931) (2.557,-1.875) (2.602,-1.814) (2.640,-1.750) (2.673,-1.683) (2.698,-1.612) (2.717,-1.540) (2.730,-1.466) (2.735,-1.391) (2.734,-1.316) (2.726,-1.241) (2.712,-1.168) (2.693,-1.095) (2.668,-1.025) (2.638,-0.956) (2.603,-0.889) (2.564,-0.825) (2.521,-0.764) (2.475,-0.705) (2.425,-0.649) (2.408,-0.630)}
\LC{(2.081,-3.206) (2.153,-3.228) (2.226,-3.247) (2.299,-3.263) (2.373,-3.277) (2.447,-3.288) (2.521,-3.297) (2.596,-3.302) (2.671,-3.305) (2.746,-3.305) (2.821,-3.301) (2.896,-3.294) (2.970,-3.284) (3.044,-3.271) (3.117,-3.254) (3.189,-3.234) (3.260,-3.210) (3.330,-3.183) (3.398,-3.152) (3.465,-3.117) (3.529,-3.079) (3.592,-3.037) (3.651,-2.992) (3.708,-2.943) (3.762,-2.891) (3.813,-2.836) (3.860,-2.777) (3.903,-2.716) (3.943,-2.652) (3.978,-2.586) (4.009,-2.518) (4.036,-2.448) (4.058,-2.376) (4.075,-2.303) (4.088,-2.229) (4.097,-2.155) (4.101,-2.080) (4.100,-2.005) (4.095,-1.930) (4.086,-1.856) (4.073,-1.782) (4.056,-1.709) (4.036,-1.637) (4.011,-1.566) (3.983,-1.496) (3.952,-1.428) (3.918,-1.361) (3.881,-1.296) (3.855,-1.253)}
\LC{(-0.657,-1.870) (-0.706,-1.926) (-0.757,-1.982) (-0.809,-2.036) (-0.863,-2.088) (-0.918,-2.138) (-0.975,-2.187) (-1.034,-2.234) (-1.094,-2.279) (-1.155,-2.322) (-1.218,-2.363) (-1.283,-2.401) (-1.348,-2.438) (-1.415,-2.471) (-1.483,-2.503) (-1.553,-2.532) (-1.623,-2.558) (-1.694,-2.581) (-1.766,-2.601) (-1.839,-2.618) (-1.913,-2.632) (-1.987,-2.642) (-2.062,-2.649) (-2.137,-2.652) (-2.212,-2.652) (-2.287,-2.648) (-2.361,-2.639) (-2.435,-2.627) (-2.508,-2.610) (-2.580,-2.589) (-2.651,-2.563) (-2.720,-2.533) (-2.786,-2.499) (-2.850,-2.460) (-2.911,-2.416) (-2.969,-2.369) (-3.023,-2.317) (-3.073,-2.261) (-3.119,-2.201) (-3.159,-2.138) (-3.195,-2.072) (-3.225,-2.003) (-3.249,-1.932) (-3.268,-1.860) (-3.282,-1.786) (-3.289,-1.712) (-3.291,-1.637) (-3.288,-1.562) (-3.279,-1.487) (-3.266,-1.413) (-3.247,-1.341) (-3.224,-1.269) (-3.197,-1.200) (-3.165,-1.132) (-3.130,-1.065) (-3.091,-1.001) (-3.049,-0.939) (-3.004,-0.879) (-2.956,-0.822) (-2.905,-0.766) (-2.852,-0.713) (-2.797,-0.663) (-2.740,-0.614) (-2.681,-0.568) (-2.620,-0.524) (-2.599,-0.510)}
\LC{(3.005,4.986) (3.075,5.014) (3.145,5.040) (3.216,5.066) (3.287,5.089) (3.358,5.112) (3.431,5.133) (3.503,5.152) (3.576,5.170) (3.649,5.186) (3.723,5.200) (3.797,5.213) (3.871,5.224) (3.945,5.234) (4.020,5.241) (4.095,5.247) (4.169,5.251) (4.244,5.253) (4.319,5.253) (4.394,5.252) (4.469,5.248) (4.544,5.242) (4.619,5.234) (4.693,5.225) (4.767,5.213) (4.841,5.199) (4.914,5.182) (4.987,5.164) (5.059,5.143) (5.130,5.121) (5.201,5.095) (5.271,5.068) (5.340,5.039) (5.408,5.007) (5.474,4.973) (5.540,4.936) (5.604,4.897) (5.667,4.857) (5.728,4.813) (5.788,4.768) (5.846,4.720) (5.902,4.671) (5.956,4.619) (6.009,4.565) (6.059,4.509) (6.107,4.452) (6.152,4.392) (6.195,4.331) (6.236,4.268) (6.274,4.203) (6.310,4.137) (6.343,4.070) (6.373,4.001) (6.400,3.931) (6.425,3.860) (6.447,3.789) (6.465,3.716) (6.481,3.643) (6.495,3.569) (6.505,3.495) (6.512,3.420) (6.517,3.345) (6.519,3.270) (6.518,3.195) (6.514,3.120) (6.508,3.046) (6.499,2.971) (6.487,2.897) (6.473,2.823) (6.456,2.750) (6.437,2.678) (6.416,2.606) (6.392,2.535) (6.367,2.464) (6.339,2.395) (6.309,2.326) (6.277,2.258) (6.243,2.191) (6.207,2.125) (6.169,2.060) (6.130,1.996) (6.089,1.934) (6.047,1.872) (6.002,1.811) (5.972,1.771)}
\LC{(10.000,1.463) (9.942,1.431) (9.874,1.399) (9.805,1.368) (9.737,1.338) (9.668,1.308) (9.599,1.278) (9.530,1.249) (9.461,1.221) (9.391,1.192) (9.322,1.165) (9.252,1.137) (9.182,1.110) (9.112,1.084) (9.041,1.058) (8.971,1.032) (8.900,1.007) (8.829,0.982) (8.758,0.957) (8.687,0.933) (8.616,0.910) (8.545,0.886) (8.474,0.864) (8.402,0.841) (8.330,0.819) (8.258,0.797) (8.187,0.776) (8.115,0.755) (8.042,0.735) (7.970,0.714) (7.898,0.695) (7.825,0.675) (7.753,0.656) (7.680,0.638) (7.608,0.619) (7.535,0.601) (7.462,0.584) (7.389,0.566) (7.316,0.549) (7.243,0.533) (7.169,0.516) (7.096,0.500) (7.023,0.485) (6.949,0.470) (6.876,0.455) (6.851,0.450)}
\LC{(3.473,-5.665) (3.543,-5.692) (3.614,-5.717) (3.685,-5.741) (3.756,-5.764) (3.828,-5.786) (3.900,-5.806) (3.973,-5.825) (4.045,-5.843) (4.119,-5.859) (4.192,-5.874) (4.266,-5.887) (4.340,-5.898) (4.414,-5.909) (4.489,-5.917) (4.564,-5.924) (4.638,-5.930) (4.713,-5.934) (4.788,-5.936) (4.863,-5.936) (4.938,-5.935) (5.013,-5.932) (5.088,-5.927) (5.163,-5.921) (5.237,-5.913) (5.312,-5.902) (5.386,-5.890) (5.459,-5.876) (5.533,-5.861) (5.606,-5.843) (5.678,-5.823) (5.750,-5.801) (5.821,-5.778) (5.891,-5.752) (5.961,-5.725) (6.030,-5.695) (6.098,-5.663) (6.165,-5.630) (6.231,-5.594) (6.296,-5.556) (6.360,-5.517) (6.422,-5.475) (6.483,-5.432) (6.542,-5.386) (6.601,-5.338) (6.657,-5.289) (6.712,-5.238) (6.765,-5.185) (6.816,-5.130) (6.865,-5.073) (6.912,-5.015) (6.957,-4.955) (7.001,-4.894) (7.041,-4.831) (7.080,-4.767) (7.116,-4.701) (7.150,-4.634) (7.182,-4.566) (7.211,-4.497) (7.238,-4.427) (7.262,-4.356) (7.284,-4.284) (7.303,-4.212) (7.319,-4.138) (7.333,-4.065) (7.345,-3.991) (7.354,-3.916) (7.360,-3.842) (7.364,-3.767) (7.366,-3.692) (7.365,-3.617) (7.362,-3.542) (7.356,-3.467) (7.348,-3.392) (7.337,-3.318) (7.325,-3.244) (7.310,-3.171) (7.293,-3.098) (7.274,-3.025) (7.253,-2.953) (7.229,-2.882) (7.204,-2.811) (7.177,-2.741) (7.148,-2.672) (7.117,-2.604) (7.085,-2.536) (7.051,-2.469) (7.015,-2.403) (6.978,-2.338) (6.939,-2.274) (6.898,-2.211) (6.856,-2.149) (6.813,-2.088) (6.768,-2.028)}
\LC{(-3.826,5.589) (-3.899,5.607) (-3.972,5.622) (-4.046,5.636) (-4.119,5.649) (-4.194,5.660) (-4.268,5.670) (-4.343,5.677) (-4.417,5.684) (-4.492,5.688) (-4.567,5.691) (-4.642,5.692) (-4.717,5.691) (-4.792,5.688) (-4.867,5.684) (-4.942,5.678) (-5.016,5.670) (-5.091,5.660) (-5.165,5.647) (-5.238,5.634) (-5.312,5.618) (-5.384,5.600) (-5.457,5.580) (-5.528,5.558) (-5.599,5.533) (-5.670,5.507) (-5.739,5.479) (-5.808,5.449) (-5.875,5.416) (-5.942,5.382) (-6.007,5.345) (-6.071,5.306) (-6.134,5.265) (-6.196,5.222) (-6.256,5.177) (-6.314,5.130) (-6.371,5.081) (-6.426,5.031) (-6.480,4.978) (-6.531,4.923) (-6.580,4.867) (-6.628,4.808) (-6.673,4.749) (-6.716,4.687) (-6.756,4.624) (-6.794,4.559) (-6.830,4.494) (-6.864,4.426) (-6.895,4.358) (-6.923,4.289) (-6.949,4.218) (-6.972,4.147) (-6.992,4.075) (-7.010,4.002) (-7.026,3.929) (-7.038,3.855) (-7.048,3.780) (-7.055,3.706) (-7.060,3.631) (-7.062,3.556) (-7.062,3.481) (-7.059,3.406) (-7.053,3.331) (-7.046,3.257) (-7.035,3.182) (-7.023,3.108) (-7.008,3.035) (-6.990,2.962) (-6.971,2.889) (-6.949,2.818) (-6.926,2.746) (-6.900,2.676) (-6.872,2.606) (-6.843,2.537) (-6.811,2.469) (-6.778,2.402) (-6.743,2.336) (-6.707,2.270)}
\LC{(0.590,-1.027) (0.659,-1.056) (0.731,-1.079) (0.804,-1.094) (0.879,-1.100) (0.954,-1.098) (1.028,-1.086) (1.099,-1.064) (1.167,-1.031) (1.227,-0.987) (1.279,-0.933) (1.320,-0.870) (1.349,-0.801) (1.363,-0.728) (1.365,-0.653) (1.353,-0.579) (1.330,-0.507) (1.298,-0.440) (1.256,-0.377) (1.208,-0.320) (1.190,-0.302)}
\LC{(-2.212,7.292) (-2.255,7.354) (-2.298,7.415) (-2.341,7.476) (-2.385,7.537) (-2.429,7.598) (-2.474,7.658) (-2.519,7.718) (-2.564,7.778) (-2.610,7.837) (-2.656,7.896) (-2.703,7.955) (-2.750,8.013) (-2.797,8.071) (-2.845,8.129) (-2.893,8.187) (-2.942,8.244) (-2.991,8.301) (-3.040,8.357) (-3.090,8.414) (-3.140,8.469) (-3.190,8.525) (-3.241,8.580) (-3.292,8.635) (-3.344,8.689) (-3.395,8.744) (-3.448,8.797) (-3.500,8.851) (-3.553,8.904) (-3.607,8.956) (-3.661,9.009) (-3.715,9.061) (-3.769,9.112) (-3.824,9.163) (-3.879,9.214) (-3.935,9.264) (-3.991,9.314) (-4.047,9.364) (-4.104,9.413) (-4.161,9.462) (-4.218,9.510) (-4.276,9.558) (-4.334,9.605) (-4.392,9.652) (-4.451,9.699) (-4.510,9.745) (-4.570,9.791) (-4.630,9.836) (-4.690,9.881) (-4.750,9.925) (-4.811,9.969) (-4.872,10.000)}
\LC{(3.158,7.740) (3.213,7.791) (3.267,7.843) (3.322,7.894) (3.378,7.944) (3.434,7.994) (3.491,8.043) (3.547,8.092) (3.605,8.140) (3.662,8.188) (3.721,8.236) (3.779,8.283) (3.838,8.329) (3.897,8.375) (3.957,8.420) (4.017,8.465) (4.078,8.510) (4.138,8.553) (4.200,8.597) (4.261,8.639) (4.323,8.681) (4.386,8.723) (4.449,8.764) (4.512,8.804) (4.575,8.844) (4.639,8.883) (4.704,8.922) (4.768,8.960) (4.833,8.997) (4.899,9.034) (4.964,9.070) (5.030,9.106) (5.097,9.141) (5.164,9.175) (5.231,9.208) (5.298,9.241) (5.366,9.274) (5.434,9.305) (5.502,9.336) (5.571,9.366) (5.640,9.395) (5.709,9.424) (5.779,9.452) (5.849,9.479) (5.919,9.506) (5.989,9.532) (6.060,9.557) (6.131,9.581) (6.202,9.605) (6.274,9.627) (6.345,9.649) (6.417,9.670) (6.489,9.691) (6.562,9.710) (6.635,9.729) (6.707,9.747) (6.780,9.764) (6.854,9.780) (6.927,9.795) (7.001,9.809) (7.074,9.823) (7.148,9.836) (7.222,9.847) (7.297,9.858) (7.371,9.868) (7.445,9.877) (7.520,9.885) (7.595,9.892) (7.669,9.898) (7.744,9.903) (7.819,9.908) (7.894,9.911) (7.969,9.913) (8.044,9.914) (8.119,9.914) (8.194,9.914) (8.269,9.912) (8.344,9.909) (8.419,9.905) (8.494,9.900) (8.568,9.894) (8.643,9.886) (8.718,9.878) (8.792,9.869) (8.866,9.858) (8.940,9.847) (9.014,9.834) (9.088,9.820) (9.162,9.805) (9.235,9.789) (9.308,9.772) (9.380,9.753) (9.453,9.734) (9.525,9.713) (9.597,9.691) (9.668,9.668) (9.739,9.644) (9.809,9.618) (9.880,9.591) (9.949,9.564) (10.000,9.534)}
\LC{(-1.193,-1.771) (-1.266,-1.788) (-1.340,-1.800) (-1.415,-1.807) (-1.490,-1.809) (-1.565,-1.805) (-1.639,-1.796) (-1.713,-1.780) (-1.784,-1.758) (-1.853,-1.729) (-1.919,-1.693) (-1.982,-1.651) (-2.039,-1.603) (-2.090,-1.549) (-2.136,-1.489) (-2.174,-1.424) (-2.204,-1.356) (-2.226,-1.284) (-2.240,-1.210) (-2.245,-1.135) (-2.242,-1.061) (-2.232,-0.986) (-2.215,-0.913) (-2.190,-0.842) (-2.160,-0.774) (-2.124,-0.708) (-2.083,-0.645) (-2.037,-0.586) (-1.988,-0.529) (-1.935,-0.477)}
\LC{(-2.191,-3.252) (-2.264,-3.270) (-2.337,-3.286) (-2.411,-3.299) (-2.485,-3.309) (-2.560,-3.317) (-2.635,-3.321) (-2.710,-3.323) (-2.785,-3.321) (-2.860,-3.317) (-2.934,-3.309) (-3.009,-3.298) (-3.082,-3.284) (-3.155,-3.266) (-3.227,-3.245) (-3.298,-3.220) (-3.367,-3.191) (-3.435,-3.159) (-3.501,-3.124) (-3.565,-3.085) (-3.626,-3.042) (-3.686,-2.996) (-3.742,-2.946) (-3.795,-2.893) (-3.845,-2.837) (-3.891,-2.778) (-3.933,-2.716) (-3.972,-2.652) (-4.006,-2.585) (-4.036,-2.517) (-4.062,-2.446) (-4.083,-2.374) (-4.100,-2.301) (-4.112,-2.227) (-4.120,-2.153) (-4.123,-2.078) (-4.122,-2.003) (-4.116,-1.928) (-4.106,-1.853) (-4.093,-1.780) (-4.075,-1.707) (-4.054,-1.635) (-4.029,-1.564) (-4.001,-1.495) (-3.969,-1.427) (-3.935,-1.360) (-3.923,-1.338)}
\LC{(-10.000,-2.495) (-9.955,-2.443) (-9.899,-2.393) (-9.844,-2.342) (-9.787,-2.293) (-9.730,-2.244) (-9.673,-2.196) (-9.615,-2.148) (-9.556,-2.102) (-9.497,-2.055) (-9.438,-2.010) (-9.378,-1.965) (-9.317,-1.921) (-9.256,-1.877) (-9.195,-1.834) (-9.133,-1.791) (-9.070,-1.750) (-9.008,-1.709) (-8.945,-1.668) (-8.881,-1.628) (-8.817,-1.589) (-8.753,-1.550) (-8.688,-1.512) (-8.623,-1.475) (-8.558,-1.438) (-8.492,-1.402) (-8.426,-1.366) (-8.360,-1.331) (-8.293,-1.297) (-8.226,-1.263) (-8.159,-1.230) (-8.092,-1.197) (-8.024,-1.165) (-7.956,-1.133) (-7.888,-1.102) (-7.865,-1.092)}
\LC{(-1.547,2.185) (-1.621,2.198) (-1.695,2.206) (-1.770,2.211) (-1.845,2.210) (-1.920,2.206) (-1.994,2.196) (-2.068,2.181) (-2.140,2.161) (-2.211,2.136) (-2.279,2.105) (-2.345,2.069) (-2.408,2.028) (-2.466,1.981) (-2.521,1.930) (-2.570,1.873) (-2.614,1.813) (-2.652,1.748) (-2.684,1.680) (-2.709,1.609) (-2.727,1.537) (-2.739,1.463) (-2.743,1.388) (-2.741,1.313) (-2.733,1.238) (-2.719,1.165) (-2.699,1.092) (-2.673,1.022) (-2.654,0.976)}

    % Cuadrupolo puntual - puntas de flecha
% Generado por scripts/generar_datos.py (no editar a mano).
\FL{1.698}{1.255}{-50.0}
\FL{-1.668}{1.205}{-127.5}
\FL{-1.140}{-0.805}{125.1}
\FL{0.742}{-2.503}{-56.3}
\FL{3.093}{-4.038}{-5.2}
\FL{5.016}{-2.372}{95.0}
\FL{3.140}{-0.363}{160.0}
\FL{0.727}{-4.161}{-70.0}
\FL{2.222}{-6.978}{-53.8}
\FL{4.485}{-9.221}{-35.2}
\FL{-0.525}{-3.023}{-109.9}
\FL{-1.850}{-5.390}{-129.0}
\FL{-3.921}{-7.133}{-151.7}
\FL{-6.535}{-7.778}{177.6}
\FL{-8.941}{-6.719}{131.8}
\FL{-9.587}{-4.176}{78.1}
\FL{-8.164}{-1.908}{40.8}
\FL{-5.787}{-0.612}{18.2}
\FL{-3.118}{-0.093}{5.1}
\FL{-0.841}{2.977}{122.2}
\FL{-3.357}{4.910}{166.1}
\FL{-6.080}{3.778}{-111.5}
\FL{-4.942}{0.990}{-34.9}
\FL{1.636}{2.861}{25.9}
\FL{3.688}{1.358}{-114.5}
\FL{0.772}{3.688}{66.1}
\FL{2.378}{6.084}{45.7}
\FL{4.794}{7.690}{20.3}
\FL{7.642}{7.872}{-15.6}
\FL{9.761}{6.016}{-69.4}
\FL{9.462}{3.225}{-119.6}
\FL{7.364}{1.238}{-150.8}
\FL{4.636}{0.291}{-169.2}
\FL{2.129}{4.115}{31.8}
\FL{4.967}{4.269}{-33.8}
\FL{5.293}{1.622}{-125.9}
\FL{-8.711}{0.934}{-18.5}
\FL{-6.038}{0.304}{-8.7}
\FL{2.566}{-6.228}{-43.4}
\FL{4.795}{-7.625}{-19.7}
\FL{7.426}{-7.819}{13.5}
\FL{9.508}{-6.272}{62.2}
\FL{9.603}{-3.705}{111.2}
\FL{7.947}{-1.649}{143.8}
\FL{5.564}{-0.521}{163.9}
\FL{-1.077}{4.763}{115.8}
\FL{-2.868}{7.337}{134.3}
\FL{-5.432}{9.133}{156.6}
\FL{-8.506}{9.632}{-173.2}
\FL{-8.188}{-0.711}{15.0}
\FL{2.765}{5.367}{32.0}
\FL{5.361}{5.998}{-7.9}
\FL{7.382}{4.398}{-72.9}
\FL{6.722}{1.874}{-131.0}
\FL{0.247}{-5.362}{-84.7}
\FL{0.618}{-8.464}{-81.6}
\FL{-0.453}{-5.987}{-98.7}
\FL{-0.962}{-8.689}{-102.7}
\FL{2.681}{-4.875}{-28.2}
\FL{5.831}{-4.848}{37.3}
\FL{6.106}{-1.874}{125.8}
\FL{-3.542}{7.008}{146.9}
\FL{-6.429}{7.945}{179.2}
\FL{-9.161}{6.822}{-130.9}
\FL{-9.701}{3.932}{-72.4}
\FL{-7.832}{1.559}{-34.6}
\FL{8.491}{-1.022}{159.2}
\FL{5.981}{-9.411}{-19.0}
\FL{8.689}{-9.707}{8.0}
\FL{3.731}{-6.301}{-23.7}
\FL{6.270}{-6.585}{13.7}
\FL{8.162}{-4.931}{71.5}
\FL{7.742}{-2.457}{123.9}
\FL{-1.237}{6.475}{111.8}
\FL{-2.440}{8.889}{121.3}
\FL{-4.616}{-8.202}{-153.1}
\FL{-8.349}{-8.729}{165.8}
\FL{-1.263}{-6.445}{-112.4}
\FL{-2.523}{-8.885}{-122.3}
\FL{-2.814}{5.567}{146.9}
\FL{-5.712}{6.264}{-170.3}
\FL{-7.797}{4.315}{-99.7}
\FL{-6.691}{1.609}{-42.1}
\FL{7.693}{-0.006}{179.9}
\FL{-0.003}{4.703}{90.1}
\FL{-0.010}{8.253}{90.1}
\FL{5.468}{8.542}{18.5}
\FL{8.660}{8.605}{-19.1}
\FL{-2.426}{-3.858}{-160.2}
\FL{-4.909}{-2.017}{73.5}
\FL{2.476}{1.945}{-43.4}
\FL{-3.025}{-5.578}{-150.9}
\FL{-5.552}{-6.073}{170.1}
\FL{-7.493}{-4.472}{107.3}
\FL{-6.921}{-2.017}{51.3}
\FL{-5.057}{-6.770}{-172.8}
\FL{-7.770}{-5.922}{133.4}
\FL{-8.236}{-3.171}{68.6}
\FL{4.999}{7.031}{11.0}
\FL{7.998}{6.312}{-42.9}
\FL{8.620}{3.344}{-110.8}
\FL{-3.140}{3.935}{178.1}
\FL{-4.876}{2.293}{-84.5}
\FL{4.461}{3.439}{-45.4}
\FL{-3.684}{-5.068}{-170.7}
\FL{-6.074}{-4.134}{121.2}
\FL{-5.718}{-1.612}{49.6}
\FL{-5.201}{6.962}{172.8}
\FL{-7.984}{6.094}{-133.5}
\FL{-8.465}{3.251}{-68.5}
\FL{-2.390}{-7.053}{-128.5}
\FL{-4.996}{-9.320}{-150.2}
\FL{1.790}{6.586}{59.0}
\FL{3.760}{9.030}{42.9}
\FL{-5.654}{8.535}{163.8}
\FL{-8.723}{8.470}{-158.7}
\FL{6.717}{-7.234}{11.4}
\FL{9.066}{-4.625}{88.1}
\FL{-3.337}{2.506}{-131.8}
\FL{-8.567}{-0.008}{0.2}
\FL{2.514}{-8.481}{-56.3}
\FL{1.542}{7.865}{67.6}
\FL{5.071}{-8.409}{-22.5}
\FL{8.494}{-8.685}{16.3}
\FL{1.430}{-7.923}{-69.4}
\FL{0.667}{8.485}{81.0}
\FL{-0.771}{8.175}{100.8}
\FL{-7.989}{-9.610}{178.5}
\FL{2.413}{-2.013}{37.0}
\FL{3.762}{-2.891}{45.7}
\FL{-2.808}{-2.486}{149.3}
\FL{4.544}{5.242}{-5.2}
\FL{6.518}{3.195}{-91.8}
\FL{8.450}{0.856}{-162.5}
\FL{5.188}{-5.918}{6.1}
\FL{7.365}{-3.617}{91.6}
\FL{-5.263}{5.628}{-168.0}
\FL{-7.060}{3.631}{-92.6}
\FL{1.263}{-0.952}{47.8}
\FL{-3.430}{8.779}{134.3}
\FL{4.639}{8.883}{31.3}
\FL{8.194}{9.914}{-1.0}
\FL{-2.162}{-1.446}{119.3}
\FL{-3.778}{-2.911}{134.6}
\FL{-8.987}{-1.695}{32.8}
\FL{-2.521}{1.930}{-133.9}
}
  \draw[-{Stealth[length=2.8pt,width=2.6pt]}, mmarojo, line width=0.9pt] (0,0) -- (0,0.95);
  \draw[-{Stealth[length=2.8pt,width=2.6pt]}, mmarojo, line width=0.9pt] (0,0) -- (0,-0.95);
  \fill[mmarojo] (0,0) circle (0.17);
  \marcoMMA{10}{5}{$x$ (u.\,arb.)}{$z$ (u.\,arb.)}
\end{tikzpicture}\\[1pt]{\small\scshape (a)\; cuadrupolo puntual}
\end{minipage}\hfill
\begin{minipage}[t]{0.49\textwidth}\centering
\begin{tikzpicture}[x=0.29cm, y=0.29cm]
  \panelMMA{10}{mmafondo}{1.2}{%
    % Cuadrupolo físico (lineal) - equipotenciales
% Generado por scripts/generar_datos.py (no editar a mano).
\EQ{(-1.975,-1.000) (-2.025,-0.998) (-2.075,-0.992) (-2.125,-0.985) (-2.175,-0.975) (-2.225,-0.963) (-2.275,-0.949) (-2.325,-0.931) (-2.375,-0.911) (-2.425,-0.887) (-2.475,-0.860) (-2.525,-0.829) (-2.567,-0.800) (-2.625,-0.753) (-2.675,-0.707) (-2.725,-0.653) (-2.768,-0.600) (-2.800,-0.555) (-2.834,-0.500) (-2.861,-0.450) (-2.884,-0.400) (-2.904,-0.350) (-2.921,-0.300) (-2.935,-0.250) (-2.947,-0.200) (-2.956,-0.150) (-2.962,-0.100) (-2.966,-0.050) (-2.967,0.000) (-2.966,0.050) (-2.962,0.100) (-2.956,0.150) (-2.947,0.200) (-2.935,0.250) (-2.921,0.300) (-2.904,0.350) (-2.884,0.400) (-2.861,0.450) (-2.834,0.500) (-2.803,0.550) (-2.768,0.600) (-2.728,0.650) (-2.682,0.700) (-2.629,0.750) (-2.575,0.794) (-2.531,0.825) (-2.475,0.860) (-2.425,0.887) (-2.375,0.911) (-2.325,0.931) (-2.275,0.949) (-2.225,0.963) (-2.175,0.975) (-2.125,0.985) (-2.075,0.992) (-2.025,0.998) (-1.975,1.000) (-1.925,1.001) (-1.875,1.000) (-1.825,0.998) (-1.775,0.993) (-1.725,0.986) (-1.675,0.978) (-1.625,0.968) (-1.575,0.957) (-1.525,0.944) (-1.475,0.930) (-1.425,0.914) (-1.375,0.897) (-1.325,0.878) (-1.275,0.859) (-1.225,0.838) (-1.175,0.816) (-1.125,0.793) (-1.075,0.770) (-1.025,0.745) (-0.975,0.720) (-0.925,0.694) (-0.875,0.668) (-0.825,0.641) (-0.775,0.614) (-0.725,0.587) (-0.675,0.560) (-0.625,0.533) (-0.575,0.506) (-0.525,0.480) (-0.475,0.454) (-0.425,0.430) (-0.375,0.407) (-0.325,0.385) (-0.275,0.365) (-0.225,0.348) (-0.175,0.333) (-0.125,0.321) (-0.075,0.313) (-0.025,0.309) (0.050,0.311) (0.100,0.317) (0.150,0.327) (0.200,0.340) (0.250,0.356) (0.300,0.375) (0.350,0.396) (0.400,0.418) (0.450,0.442) (0.500,0.467) (0.550,0.493) (0.600,0.519) (0.650,0.546) (0.700,0.573) (0.750,0.600) (0.796,0.625) (0.842,0.650) (0.889,0.675) (0.937,0.700) (0.986,0.725) (1.036,0.750) (1.087,0.775) (1.140,0.800) (1.196,0.825) (1.250,0.848) (1.300,0.869) (1.350,0.888) (1.400,0.905) (1.450,0.922) (1.500,0.937) (1.550,0.951) (1.600,0.963) (1.650,0.973) (1.700,0.982) (1.750,0.990) (1.800,0.995) (1.850,0.999) (1.900,1.001) (1.950,1.001) (2.000,0.999) (2.050,0.995) (2.100,0.989) (2.150,0.981) (2.200,0.970) (2.250,0.956) (2.300,0.940) (2.350,0.922) (2.399,0.900) (2.449,0.875) (2.492,0.850) (2.550,0.812) (2.599,0.775) (2.650,0.731) (2.700,0.681) (2.749,0.625) (2.786,0.575) (2.819,0.525) (2.848,0.475) (2.873,0.425) (2.895,0.375) (2.913,0.325) (2.929,0.275) (2.941,0.225) (2.952,0.175) (2.959,0.125) (2.964,0.075) (2.967,0.025) (2.966,-0.050) (2.962,-0.100) (2.956,-0.150) (2.947,-0.200) (2.935,-0.250) (2.921,-0.300) (2.904,-0.350) (2.884,-0.400) (2.861,-0.450) (2.834,-0.500) (2.803,-0.550) (2.768,-0.600) (2.728,-0.650) (2.682,-0.700) (2.629,-0.750) (2.575,-0.794) (2.531,-0.825) (2.475,-0.860) (2.425,-0.887) (2.375,-0.911) (2.325,-0.931) (2.275,-0.949) (2.225,-0.963) (2.175,-0.975) (2.125,-0.985) (2.075,-0.992) (2.025,-0.998) (1.975,-1.000) (1.925,-1.001) (1.875,-1.000) (1.825,-0.998) (1.775,-0.993) (1.725,-0.986) (1.675,-0.978) (1.625,-0.968) (1.575,-0.957) (1.525,-0.944) (1.475,-0.930) (1.425,-0.914) (1.375,-0.897) (1.325,-0.878) (1.275,-0.859) (1.225,-0.838) (1.175,-0.816) (1.125,-0.793) (1.075,-0.770) (1.025,-0.745) (0.975,-0.720) (0.925,-0.694) (0.875,-0.668) (0.825,-0.641) (0.775,-0.614) (0.725,-0.587) (0.675,-0.560) (0.625,-0.533) (0.575,-0.506) (0.525,-0.480) (0.475,-0.454) (0.425,-0.430) (0.375,-0.407) (0.325,-0.385) (0.275,-0.365) (0.225,-0.348) (0.175,-0.333) (0.125,-0.321) (0.075,-0.313) (0.025,-0.309) (-0.050,-0.311) (-0.100,-0.317) (-0.150,-0.327) (-0.200,-0.340) (-0.250,-0.356) (-0.300,-0.375) (-0.350,-0.396) (-0.400,-0.418) (-0.450,-0.442) (-0.500,-0.467) (-0.550,-0.493) (-0.600,-0.519) (-0.650,-0.546) (-0.700,-0.573) (-0.750,-0.600) (-0.796,-0.625) (-0.842,-0.650) (-0.889,-0.675) (-0.937,-0.700) (-0.986,-0.725) (-1.036,-0.750) (-1.087,-0.775) (-1.140,-0.800) (-1.196,-0.825) (-1.250,-0.848) (-1.300,-0.869) (-1.350,-0.888) (-1.400,-0.905) (-1.450,-0.922) (-1.500,-0.937) (-1.550,-0.951) (-1.600,-0.963) (-1.650,-0.973) (-1.700,-0.982) (-1.750,-0.990) (-1.800,-0.995) (-1.850,-0.999) (-1.900,-1.001) (-1.950,-1.001) (-1.975,-1.000)}
\EQ{(-2.200,-1.003) (-2.250,-0.992) (-2.300,-0.979) (-2.350,-0.963) (-2.400,-0.945) (-2.446,-0.925) (-2.497,-0.900) (-2.542,-0.875) (-2.600,-0.838) (-2.650,-0.801) (-2.700,-0.759) (-2.750,-0.711) (-2.800,-0.656) (-2.844,-0.600) (-2.875,-0.554) (-2.907,-0.500) (-2.933,-0.450) (-2.956,-0.400) (-2.975,-0.352) (-2.992,-0.300) (-3.006,-0.250) (-3.017,-0.200) (-3.026,-0.150) (-3.032,-0.100) (-3.036,-0.050) (-3.037,0.000) (-3.036,0.050) (-3.032,0.100) (-3.026,0.150) (-3.017,0.200) (-3.006,0.250) (-2.992,0.300) (-2.976,0.350) (-2.956,0.400) (-2.933,0.450) (-2.907,0.500) (-2.877,0.550) (-2.844,0.600) (-2.805,0.650) (-2.761,0.700) (-2.725,0.736) (-2.682,0.775) (-2.625,0.820) (-2.582,0.850) (-2.525,0.885) (-2.475,0.911) (-2.425,0.934) (-2.375,0.954) (-2.325,0.972) (-2.275,0.986) (-2.225,0.998) (-2.175,1.007) (-2.125,1.015) (-2.075,1.020) (-2.025,1.023) (-1.975,1.024) (-1.925,1.023) (-1.875,1.020) (-1.825,1.016) (-1.775,1.010) (-1.725,1.002) (-1.675,0.992) (-1.625,0.981) (-1.575,0.968) (-1.525,0.954) (-1.475,0.939) (-1.425,0.922) (-1.375,0.904) (-1.325,0.885) (-1.275,0.865) (-1.225,0.843) (-1.175,0.821) (-1.125,0.797) (-1.080,0.775) (-1.030,0.750) (-0.981,0.725) (-0.933,0.700) (-0.886,0.675) (-0.840,0.650) (-0.794,0.625) (-0.748,0.600) (-0.702,0.575) (-0.657,0.550) (-0.611,0.525) (-0.564,0.500) (-0.516,0.475) (-0.467,0.450) (-0.415,0.425) (-0.361,0.400) (-0.301,0.375) (-0.250,0.356) (-0.200,0.340) (-0.150,0.327) (-0.100,0.317) (-0.050,0.311) (0.000,0.309) (0.050,0.311) (0.100,0.317) (0.150,0.327) (0.200,0.340) (0.250,0.356) (0.300,0.375) (0.350,0.396) (0.400,0.418) (0.450,0.442) (0.500,0.467) (0.550,0.493) (0.600,0.520) (0.650,0.547) (0.700,0.574) (0.748,0.600) (0.794,0.625) (0.840,0.650) (0.886,0.675) (0.933,0.700) (0.981,0.725) (1.030,0.750) (1.075,0.773) (1.125,0.797) (1.175,0.821) (1.225,0.843) (1.275,0.865) (1.325,0.885) (1.375,0.904) (1.425,0.922) (1.475,0.939) (1.525,0.954) (1.575,0.968) (1.625,0.981) (1.675,0.992) (1.725,1.002) (1.775,1.010) (1.825,1.016) (1.875,1.020) (1.925,1.023) (1.975,1.024) (2.025,1.023) (2.075,1.020) (2.125,1.015) (2.175,1.007) (2.225,0.998) (2.275,0.986) (2.325,0.972) (2.375,0.954) (2.425,0.934) (2.475,0.911) (2.525,0.885) (2.575,0.854) (2.618,0.825) (2.675,0.781) (2.725,0.736) (2.761,0.700) (2.800,0.656) (2.844,0.600) (2.875,0.554) (2.907,0.500) (2.933,0.450) (2.956,0.400) (2.975,0.352) (2.992,0.300) (3.006,0.250) (3.017,0.200) (3.026,0.150) (3.032,0.100) (3.036,0.050) (3.037,0.000) (3.036,-0.050) (3.032,-0.100) (3.026,-0.150) (3.017,-0.200) (3.006,-0.250) (2.992,-0.300) (2.976,-0.350) (2.956,-0.400) (2.933,-0.450) (2.907,-0.500) (2.877,-0.550) (2.844,-0.600) (2.805,-0.650) (2.761,-0.700) (2.725,-0.736) (2.682,-0.775) (2.625,-0.820) (2.582,-0.850) (2.525,-0.885) (2.475,-0.911) (2.425,-0.934) (2.375,-0.954) (2.325,-0.972) (2.275,-0.986) (2.225,-0.998) (2.175,-1.007) (2.125,-1.015) (2.075,-1.020) (2.025,-1.023) (1.975,-1.024) (1.925,-1.023) (1.875,-1.020) (1.825,-1.016) (1.775,-1.010) (1.725,-1.002) (1.675,-0.992) (1.625,-0.981) (1.575,-0.968) (1.525,-0.954) (1.475,-0.939) (1.425,-0.922) (1.375,-0.904) (1.325,-0.885) (1.275,-0.865) (1.225,-0.843) (1.175,-0.821) (1.125,-0.797) (1.080,-0.775) (1.030,-0.750) (0.981,-0.725) (0.933,-0.700) (0.886,-0.675) (0.840,-0.650) (0.794,-0.625) (0.748,-0.600) (0.702,-0.575) (0.657,-0.550) (0.611,-0.525) (0.564,-0.500) (0.516,-0.475) (0.467,-0.450) (0.415,-0.425) (0.361,-0.400) (0.301,-0.375) (0.250,-0.356) (0.200,-0.340) (0.150,-0.327) (0.100,-0.317) (0.050,-0.311) (0.000,-0.309) (-0.050,-0.311) (-0.100,-0.317) (-0.150,-0.327) (-0.200,-0.340) (-0.250,-0.356) (-0.300,-0.375) (-0.350,-0.396) (-0.400,-0.418) (-0.450,-0.442) (-0.500,-0.467) (-0.550,-0.493) (-0.600,-0.520) (-0.650,-0.547) (-0.700,-0.574) (-0.748,-0.600) (-0.794,-0.625) (-0.840,-0.650) (-0.886,-0.675) (-0.933,-0.700) (-0.981,-0.725) (-1.030,-0.750) (-1.075,-0.773) (-1.125,-0.797) (-1.175,-0.821) (-1.225,-0.843) (-1.275,-0.865) (-1.325,-0.885) (-1.375,-0.904) (-1.425,-0.922) (-1.475,-0.939) (-1.525,-0.954) (-1.575,-0.968) (-1.625,-0.981) (-1.675,-0.992) (-1.725,-1.002) (-1.775,-1.010) (-1.825,-1.016) (-1.875,-1.020) (-1.925,-1.023) (-1.975,-1.024) (-2.025,-1.023) (-2.075,-1.020) (-2.125,-1.015) (-2.175,-1.007) (-2.200,-1.003)}
\EQ{(-2.250,-1.029) (-2.300,-1.018) (-2.350,-1.005) (-2.400,-0.990) (-2.450,-0.972) (-2.500,-0.951) (-2.550,-0.927) (-2.600,-0.900) (-2.650,-0.869) (-2.700,-0.834) (-2.743,-0.800) (-2.798,-0.750) (-2.847,-0.700) (-2.889,-0.650) (-2.925,-0.602) (-2.959,-0.550) (-2.988,-0.500) (-3.013,-0.450) (-3.035,-0.400) (-3.054,-0.350) (-3.070,-0.300) (-3.084,-0.250) (-3.095,-0.200) (-3.103,-0.150) (-3.109,-0.100) (-3.113,-0.050) (-3.114,0.000) (-3.113,0.050) (-3.109,0.100) (-3.103,0.150) (-3.095,0.200) (-3.084,0.250) (-3.070,0.300) (-3.054,0.350) (-3.035,0.400) (-3.013,0.450) (-2.988,0.500) (-2.959,0.550) (-2.926,0.600) (-2.889,0.650) (-2.850,0.696) (-2.800,0.748) (-2.750,0.794) (-2.700,0.834) (-2.650,0.869) (-2.600,0.900) (-2.555,0.925) (-2.503,0.950) (-2.450,0.972) (-2.400,0.990) (-2.350,1.005) (-2.300,1.018) (-2.250,1.029) (-2.200,1.037) (-2.150,1.043) (-2.100,1.047) (-2.050,1.048) (-2.000,1.049) (-1.950,1.047) (-1.900,1.043) (-1.850,1.038) (-1.800,1.031) (-1.750,1.022) (-1.700,1.012) (-1.650,1.000) (-1.600,0.987) (-1.550,0.973) (-1.500,0.957) (-1.450,0.940) (-1.400,0.921) (-1.350,0.902) (-1.300,0.881) (-1.250,0.859) (-1.200,0.837) (-1.150,0.813) (-1.100,0.789) (-1.050,0.764) (-1.000,0.738) (-0.950,0.711) (-0.900,0.685) (-0.850,0.657) (-0.800,0.630) (-0.750,0.602) (-0.701,0.575) (-0.656,0.550) (-0.610,0.525) (-0.563,0.500) (-0.516,0.475) (-0.466,0.450) (-0.415,0.425) (-0.361,0.400) (-0.300,0.375) (-0.250,0.356) (-0.200,0.340) (-0.150,0.327) (-0.100,0.317) (-0.050,0.311) (0.000,0.309) (0.050,0.311) (0.100,0.317) (0.150,0.327) (0.200,0.340) (0.250,0.356) (0.300,0.375) (0.350,0.396) (0.400,0.418) (0.450,0.442) (0.500,0.468) (0.550,0.494) (0.600,0.520) (0.650,0.547) (0.700,0.575) (0.746,0.600) (0.792,0.625) (0.837,0.650) (0.883,0.675) (0.929,0.700) (0.975,0.724) (1.024,0.750) (1.073,0.775) (1.123,0.800) (1.175,0.825) (1.225,0.848) (1.275,0.870) (1.325,0.892) (1.375,0.912) (1.425,0.931) (1.475,0.948) (1.525,0.965) (1.575,0.980) (1.625,0.994) (1.675,1.006) (1.725,1.017) (1.775,1.027) (1.825,1.035) (1.875,1.041) (1.925,1.045) (1.975,1.048) (2.025,1.049) (2.075,1.048) (2.125,1.045) (2.175,1.040) (2.225,1.033) (2.275,1.024) (2.325,1.012) (2.375,0.998) (2.425,0.982) (2.475,0.962) (2.525,0.940) (2.575,0.914) (2.625,0.885) (2.675,0.852) (2.725,0.815) (2.772,0.775) (2.823,0.725) (2.869,0.675) (2.900,0.636) (2.943,0.575) (2.974,0.525) (3.000,0.477) (3.024,0.425) (3.045,0.375) (3.063,0.325) (3.077,0.275) (3.090,0.225) (3.099,0.175) (3.106,0.125) (3.111,0.075) (3.113,0.025) (3.113,-0.050) (3.109,-0.100) (3.103,-0.150) (3.095,-0.200) (3.084,-0.250) (3.070,-0.300) (3.054,-0.350) (3.035,-0.400) (3.013,-0.450) (2.988,-0.500) (2.959,-0.550) (2.926,-0.600) (2.889,-0.650) (2.850,-0.696) (2.800,-0.748) (2.750,-0.794) (2.700,-0.834) (2.650,-0.869) (2.600,-0.900) (2.555,-0.925) (2.503,-0.950) (2.450,-0.972) (2.400,-0.990) (2.350,-1.005) (2.300,-1.018) (2.250,-1.029) (2.200,-1.037) (2.150,-1.043) (2.100,-1.047) (2.050,-1.048) (2.000,-1.049) (1.950,-1.047) (1.900,-1.043) (1.850,-1.038) (1.800,-1.031) (1.750,-1.022) (1.700,-1.012) (1.650,-1.000) (1.600,-0.987) (1.550,-0.973) (1.500,-0.957) (1.450,-0.940) (1.400,-0.921) (1.350,-0.902) (1.300,-0.881) (1.250,-0.859) (1.200,-0.837) (1.150,-0.813) (1.100,-0.789) (1.050,-0.764) (1.000,-0.738) (0.950,-0.711) (0.900,-0.685) (0.850,-0.657) (0.800,-0.630) (0.750,-0.602) (0.701,-0.575) (0.656,-0.550) (0.610,-0.525) (0.563,-0.500) (0.516,-0.475) (0.466,-0.450) (0.415,-0.425) (0.361,-0.400) (0.300,-0.375) (0.250,-0.356) (0.200,-0.340) (0.150,-0.327) (0.100,-0.317) (0.050,-0.311) (0.000,-0.309) (-0.050,-0.311) (-0.100,-0.317) (-0.150,-0.327) (-0.200,-0.340) (-0.250,-0.356) (-0.300,-0.375) (-0.350,-0.396) (-0.400,-0.418) (-0.450,-0.442) (-0.500,-0.468) (-0.550,-0.494) (-0.600,-0.520) (-0.650,-0.547) (-0.700,-0.575) (-0.746,-0.600) (-0.792,-0.625) (-0.837,-0.650) (-0.883,-0.675) (-0.929,-0.700) (-0.975,-0.724) (-1.024,-0.750) (-1.073,-0.775) (-1.123,-0.800) (-1.175,-0.825) (-1.225,-0.848) (-1.275,-0.870) (-1.325,-0.892) (-1.375,-0.912) (-1.425,-0.931) (-1.475,-0.948) (-1.525,-0.965) (-1.575,-0.980) (-1.625,-0.994) (-1.675,-1.006) (-1.725,-1.017) (-1.775,-1.027) (-1.825,-1.035) (-1.875,-1.041) (-1.925,-1.045) (-1.975,-1.048) (-2.025,-1.049) (-2.075,-1.048) (-2.125,-1.045) (-2.175,-1.040) (-2.225,-1.033) (-2.250,-1.029)}
\EQ{(-2.125,-1.076) (-2.175,-1.073) (-2.225,-1.068) (-2.275,-1.062) (-2.325,-1.053) (-2.375,-1.042) (-2.425,-1.029) (-2.475,-1.013) (-2.525,-0.995) (-2.572,-0.975) (-2.624,-0.950) (-2.670,-0.925) (-2.725,-0.891) (-2.775,-0.855) (-2.825,-0.815) (-2.870,-0.775) (-2.919,-0.725) (-2.962,-0.675) (-3.000,-0.625) (-3.034,-0.575) (-3.063,-0.525) (-3.089,-0.475) (-3.112,-0.425) (-3.132,-0.375) (-3.149,-0.325) (-3.164,-0.275) (-3.175,-0.225) (-3.185,-0.175) (-3.192,-0.125) (-3.196,-0.075) (-3.199,-0.025) (-3.198,0.050) (-3.194,0.100) (-3.189,0.150) (-3.180,0.200) (-3.170,0.250) (-3.157,0.300) (-3.141,0.350) (-3.123,0.400) (-3.101,0.450) (-3.077,0.500) (-3.050,0.548) (-3.017,0.600) (-2.982,0.650) (-2.950,0.689) (-2.900,0.745) (-2.850,0.793) (-2.800,0.836) (-2.750,0.874) (-2.700,0.907) (-2.650,0.936) (-2.600,0.962) (-2.550,0.985) (-2.500,1.004) (-2.450,1.021) (-2.400,1.036) (-2.350,1.048) (-2.300,1.058) (-2.250,1.065) (-2.200,1.071) (-2.150,1.074) (-2.100,1.076) (-2.050,1.076) (-2.000,1.074) (-1.950,1.070) (-1.900,1.065) (-1.850,1.058) (-1.800,1.049) (-1.750,1.039) (-1.700,1.027) (-1.650,1.014) (-1.601,1.000) (-1.550,0.984) (-1.500,0.967) (-1.450,0.949) (-1.400,0.929) (-1.350,0.909) (-1.300,0.887) (-1.250,0.865) (-1.200,0.842) (-1.150,0.817) (-1.100,0.792) (-1.050,0.767) (-1.000,0.740) (-0.950,0.714) (-0.900,0.686) (-0.850,0.659) (-0.800,0.631) (-0.750,0.603) (-0.700,0.575) (-0.654,0.550) (-0.609,0.525) (-0.563,0.500) (-0.515,0.475) (-0.466,0.450) (-0.415,0.425) (-0.360,0.400) (-0.300,0.375) (-0.250,0.356) (-0.200,0.340) (-0.150,0.327) (-0.100,0.317) (-0.050,0.311) (0.000,0.309) (0.050,0.311) (0.100,0.317) (0.150,0.327) (0.200,0.340) (0.250,0.356) (0.300,0.375) (0.350,0.396) (0.400,0.419) (0.450,0.443) (0.500,0.468) (0.550,0.494) (0.600,0.521) (0.650,0.548) (0.699,0.575) (0.745,0.600) (0.790,0.625) (0.835,0.650) (0.880,0.675) (0.925,0.700) (0.972,0.725) (1.019,0.750) (1.067,0.775) (1.116,0.800) (1.166,0.825) (1.218,0.850) (1.273,0.875) (1.325,0.898) (1.375,0.919) (1.425,0.939) (1.475,0.958) (1.523,0.975) (1.575,0.992) (1.625,1.007) (1.675,1.021) (1.725,1.033) (1.775,1.044) (1.825,1.054) (1.875,1.061) (1.925,1.068) (1.975,1.072) (2.025,1.075) (2.075,1.076) (2.125,1.076) (2.175,1.073) (2.225,1.068) (2.275,1.062) (2.325,1.053) (2.375,1.042) (2.425,1.029) (2.475,1.013) (2.525,0.995) (2.572,0.975) (2.624,0.950) (2.670,0.925) (2.725,0.891) (2.775,0.855) (2.825,0.815) (2.870,0.775) (2.919,0.725) (2.962,0.675) (3.000,0.625) (3.034,0.575) (3.063,0.525) (3.089,0.475) (3.112,0.425) (3.132,0.375) (3.149,0.325) (3.164,0.275) (3.175,0.225) (3.185,0.175) (3.192,0.125) (3.196,0.075) (3.199,0.025) (3.198,-0.050) (3.194,-0.100) (3.189,-0.150) (3.180,-0.200) (3.170,-0.250) (3.157,-0.300) (3.141,-0.350) (3.123,-0.400) (3.101,-0.450) (3.077,-0.500) (3.050,-0.548) (3.017,-0.600) (2.982,-0.650) (2.950,-0.689) (2.900,-0.745) (2.850,-0.793) (2.800,-0.836) (2.750,-0.874) (2.700,-0.907) (2.650,-0.936) (2.600,-0.962) (2.550,-0.985) (2.500,-1.004) (2.450,-1.021) (2.400,-1.036) (2.350,-1.048) (2.300,-1.058) (2.250,-1.065) (2.200,-1.071) (2.150,-1.074) (2.100,-1.076) (2.050,-1.076) (2.000,-1.074) (1.950,-1.070) (1.900,-1.065) (1.850,-1.058) (1.800,-1.049) (1.750,-1.039) (1.700,-1.027) (1.650,-1.014) (1.601,-1.000) (1.550,-0.984) (1.500,-0.967) (1.450,-0.949) (1.400,-0.929) (1.350,-0.909) (1.300,-0.887) (1.250,-0.865) (1.200,-0.842) (1.150,-0.817) (1.100,-0.792) (1.050,-0.767) (1.000,-0.740) (0.950,-0.714) (0.900,-0.686) (0.850,-0.659) (0.800,-0.631) (0.750,-0.603) (0.700,-0.575) (0.654,-0.550) (0.609,-0.525) (0.563,-0.500) (0.515,-0.475) (0.466,-0.450) (0.415,-0.425) (0.360,-0.400) (0.300,-0.375) (0.250,-0.356) (0.200,-0.340) (0.150,-0.327) (0.100,-0.317) (0.050,-0.311) (0.000,-0.309) (-0.050,-0.311) (-0.100,-0.317) (-0.150,-0.327) (-0.200,-0.340) (-0.250,-0.356) (-0.300,-0.375) (-0.350,-0.396) (-0.400,-0.419) (-0.450,-0.443) (-0.500,-0.468) (-0.550,-0.494) (-0.600,-0.521) (-0.650,-0.548) (-0.699,-0.575) (-0.745,-0.600) (-0.790,-0.625) (-0.835,-0.650) (-0.880,-0.675) (-0.925,-0.700) (-0.972,-0.725) (-1.019,-0.750) (-1.067,-0.775) (-1.116,-0.800) (-1.166,-0.825) (-1.218,-0.850) (-1.273,-0.875) (-1.325,-0.898) (-1.375,-0.919) (-1.425,-0.939) (-1.475,-0.958) (-1.523,-0.975) (-1.575,-0.992) (-1.625,-1.007) (-1.675,-1.021) (-1.725,-1.033) (-1.775,-1.044) (-1.825,-1.054) (-1.875,-1.061) (-1.925,-1.068) (-1.975,-1.072) (-2.025,-1.075) (-2.075,-1.076) (-2.125,-1.076)}
\EQ{(-2.275,-1.101) (-2.325,-1.095) (-2.375,-1.087) (-2.425,-1.077) (-2.475,-1.065) (-2.525,-1.050) (-2.575,-1.033) (-2.625,-1.014) (-2.675,-0.991) (-2.725,-0.965) (-2.775,-0.936) (-2.825,-0.904) (-2.875,-0.867) (-2.925,-0.825) (-2.975,-0.778) (-3.025,-0.725) (-3.066,-0.675) (-3.100,-0.628) (-3.134,-0.575) (-3.163,-0.525) (-3.188,-0.475) (-3.210,-0.425) (-3.229,-0.375) (-3.246,-0.325) (-3.260,-0.275) (-3.271,-0.225) (-3.280,-0.175) (-3.287,-0.125) (-3.292,-0.075) (-3.294,-0.025) (-3.293,0.050) (-3.290,0.100) (-3.284,0.150) (-3.276,0.200) (-3.266,0.250) (-3.253,0.300) (-3.238,0.350) (-3.220,0.400) (-3.200,0.449) (-3.176,0.500) (-3.150,0.548) (-3.119,0.600) (-3.084,0.650) (-3.050,0.695) (-3.002,0.750) (-2.953,0.800) (-2.900,0.847) (-2.850,0.886) (-2.800,0.921) (-2.752,0.950) (-2.707,0.975) (-2.656,1.000) (-2.600,1.024) (-2.550,1.042) (-2.500,1.058) (-2.450,1.071) (-2.400,1.082) (-2.350,1.091) (-2.300,1.098) (-2.250,1.103) (-2.200,1.106) (-2.150,1.107) (-2.100,1.106) (-2.050,1.104) (-2.000,1.100) (-1.950,1.094) (-1.900,1.087) (-1.850,1.078) (-1.800,1.068) (-1.750,1.056) (-1.700,1.043) (-1.650,1.028) (-1.600,1.013) (-1.550,0.996) (-1.500,0.977) (-1.450,0.958) (-1.400,0.938) (-1.350,0.916) (-1.300,0.894) (-1.250,0.871) (-1.200,0.846) (-1.150,0.821) (-1.100,0.796) (-1.050,0.770) (-1.000,0.743) (-0.950,0.716) (-0.900,0.688) (-0.850,0.660) (-0.800,0.632) (-0.750,0.604) (-0.700,0.576) (-0.653,0.550) (-0.608,0.525) (-0.562,0.500) (-0.514,0.475) (-0.466,0.450) (-0.415,0.425) (-0.360,0.400) (-0.300,0.375) (-0.250,0.356) (-0.200,0.340) (-0.150,0.327) (-0.100,0.317) (-0.050,0.311) (0.000,0.309) (0.050,0.311) (0.100,0.317) (0.150,0.327) (0.200,0.340) (0.250,0.356) (0.300,0.375) (0.350,0.396) (0.400,0.419) (0.450,0.443) (0.500,0.468) (0.550,0.494) (0.600,0.521) (0.650,0.548) (0.698,0.575) (0.743,0.600) (0.788,0.625) (0.832,0.650) (0.877,0.675) (0.922,0.700) (0.968,0.725) (1.014,0.750) (1.061,0.775) (1.109,0.800) (1.158,0.825) (1.208,0.850) (1.260,0.875) (1.314,0.900) (1.371,0.925) (1.425,0.948) (1.475,0.968) (1.525,0.987) (1.575,1.004) (1.625,1.021) (1.675,1.036) (1.725,1.049) (1.775,1.062) (1.825,1.073) (1.875,1.083) (1.925,1.091) (1.975,1.097) (2.025,1.102) (2.075,1.105) (2.125,1.107) (2.175,1.107) (2.225,1.105) (2.275,1.101) (2.325,1.095) (2.375,1.087) (2.425,1.077) (2.475,1.065) (2.525,1.050) (2.575,1.033) (2.625,1.014) (2.675,0.991) (2.725,0.965) (2.775,0.936) (2.825,0.904) (2.875,0.867) (2.925,0.825) (2.975,0.778) (3.025,0.725) (3.066,0.675) (3.100,0.628) (3.134,0.575) (3.163,0.525) (3.188,0.475) (3.210,0.425) (3.229,0.375) (3.246,0.325) (3.260,0.275) (3.271,0.225) (3.280,0.175) (3.287,0.125) (3.292,0.075) (3.294,0.025) (3.293,-0.050) (3.290,-0.100) (3.284,-0.150) (3.276,-0.200) (3.266,-0.250) (3.253,-0.300) (3.238,-0.350) (3.220,-0.400) (3.200,-0.449) (3.176,-0.500) (3.150,-0.548) (3.119,-0.600) (3.084,-0.650) (3.050,-0.695) (3.002,-0.750) (2.953,-0.800) (2.900,-0.847) (2.850,-0.886) (2.800,-0.921) (2.752,-0.950) (2.707,-0.975) (2.656,-1.000) (2.600,-1.024) (2.550,-1.042) (2.500,-1.058) (2.450,-1.071) (2.400,-1.082) (2.350,-1.091) (2.300,-1.098) (2.250,-1.103) (2.200,-1.106) (2.150,-1.107) (2.100,-1.106) (2.050,-1.104) (2.000,-1.100) (1.950,-1.094) (1.900,-1.087) (1.850,-1.078) (1.800,-1.068) (1.750,-1.056) (1.700,-1.043) (1.650,-1.028) (1.600,-1.013) (1.550,-0.996) (1.500,-0.977) (1.450,-0.958) (1.400,-0.938) (1.350,-0.916) (1.300,-0.894) (1.250,-0.871) (1.200,-0.846) (1.150,-0.821) (1.100,-0.796) (1.050,-0.770) (1.000,-0.743) (0.950,-0.716) (0.900,-0.688) (0.850,-0.660) (0.800,-0.632) (0.750,-0.604) (0.700,-0.576) (0.653,-0.550) (0.608,-0.525) (0.562,-0.500) (0.514,-0.475) (0.466,-0.450) (0.415,-0.425) (0.360,-0.400) (0.300,-0.375) (0.250,-0.356) (0.200,-0.340) (0.150,-0.327) (0.100,-0.317) (0.050,-0.311) (0.000,-0.309) (-0.050,-0.311) (-0.100,-0.317) (-0.150,-0.327) (-0.200,-0.340) (-0.250,-0.356) (-0.300,-0.375) (-0.350,-0.396) (-0.400,-0.419) (-0.450,-0.443) (-0.500,-0.468) (-0.550,-0.494) (-0.600,-0.521) (-0.650,-0.548) (-0.698,-0.575) (-0.743,-0.600) (-0.788,-0.625) (-0.832,-0.650) (-0.877,-0.675) (-0.922,-0.700) (-0.968,-0.725) (-1.014,-0.750) (-1.061,-0.775) (-1.109,-0.800) (-1.158,-0.825) (-1.208,-0.850) (-1.260,-0.875) (-1.314,-0.900) (-1.371,-0.925) (-1.425,-0.948) (-1.475,-0.968) (-1.525,-0.987) (-1.575,-1.004) (-1.625,-1.021) (-1.675,-1.036) (-1.725,-1.049) (-1.775,-1.062) (-1.825,-1.073) (-1.875,-1.083) (-1.925,-1.091) (-1.975,-1.097) (-2.025,-1.102) (-2.075,-1.105) (-2.125,-1.107) (-2.175,-1.107) (-2.225,-1.105) (-2.275,-1.101)}
\EQ{(-2.425,-1.126) (-2.475,-1.117) (-2.525,-1.106) (-2.575,-1.093) (-2.625,-1.078) (-2.675,-1.060) (-2.725,-1.040) (-2.775,-1.016) (-2.825,-0.990) (-2.875,-0.960) (-2.925,-0.927) (-2.975,-0.890) (-3.022,-0.850) (-3.075,-0.800) (-3.122,-0.750) (-3.163,-0.700) (-3.200,-0.650) (-3.233,-0.600) (-3.262,-0.550) (-3.287,-0.500) (-3.310,-0.450) (-3.330,-0.400) (-3.347,-0.350) (-3.362,-0.300) (-3.374,-0.250) (-3.384,-0.200) (-3.392,-0.150) (-3.397,-0.100) (-3.401,-0.050) (-3.402,0.000) (-3.401,0.050) (-3.397,0.100) (-3.392,0.150) (-3.384,0.200) (-3.374,0.250) (-3.362,0.300) (-3.347,0.350) (-3.330,0.400) (-3.310,0.450) (-3.287,0.500) (-3.262,0.550) (-3.233,0.600) (-3.200,0.650) (-3.163,0.700) (-3.125,0.746) (-3.075,0.800) (-3.025,0.848) (-2.975,0.890) (-2.928,0.925) (-2.875,0.960) (-2.825,0.990) (-2.775,1.016) (-2.725,1.040) (-2.675,1.060) (-2.625,1.078) (-2.575,1.093) (-2.525,1.106) (-2.475,1.117) (-2.425,1.126) (-2.375,1.133) (-2.325,1.138) (-2.275,1.141) (-2.225,1.142) (-2.175,1.141) (-2.125,1.139) (-2.075,1.135) (-2.025,1.130) (-1.975,1.123) (-1.925,1.114) (-1.875,1.104) (-1.825,1.093) (-1.775,1.080) (-1.725,1.066) (-1.675,1.051) (-1.625,1.034) (-1.575,1.017) (-1.525,0.998) (-1.475,0.978) (-1.425,0.957) (-1.375,0.935) (-1.325,0.912) (-1.275,0.888) (-1.225,0.864) (-1.175,0.839) (-1.125,0.813) (-1.075,0.786) (-1.025,0.759) (-0.975,0.732) (-0.925,0.704) (-0.875,0.676) (-0.830,0.650) (-0.786,0.625) (-0.742,0.600) (-0.697,0.575) (-0.652,0.550) (-0.607,0.525) (-0.561,0.500) (-0.514,0.475) (-0.465,0.450) (-0.414,0.425) (-0.360,0.400) (-0.300,0.375) (-0.250,0.356) (-0.200,0.340) (-0.150,0.327) (-0.100,0.317) (-0.050,0.311) (0.000,0.309) (0.050,0.311) (0.100,0.317) (0.150,0.327) (0.200,0.340) (0.250,0.356) (0.300,0.375) (0.350,0.396) (0.400,0.419) (0.450,0.443) (0.500,0.468) (0.550,0.495) (0.600,0.522) (0.650,0.549) (0.697,0.575) (0.742,0.600) (0.786,0.625) (0.830,0.650) (0.874,0.675) (0.919,0.700) (0.964,0.725) (1.009,0.750) (1.055,0.775) (1.100,0.799) (1.149,0.825) (1.198,0.850) (1.248,0.875) (1.300,0.900) (1.350,0.924) (1.400,0.946) (1.450,0.967) (1.500,0.988) (1.550,1.007) (1.599,1.025) (1.650,1.043) (1.700,1.059) (1.750,1.073) (1.800,1.087) (1.850,1.099) (1.900,1.109) (1.950,1.119) (2.000,1.126) (2.050,1.133) (2.100,1.137) (2.150,1.140) (2.200,1.142) (2.250,1.141) (2.300,1.139) (2.350,1.135) (2.400,1.130) (2.450,1.122) (2.500,1.112) (2.550,1.100) (2.600,1.086) (2.650,1.069) (2.700,1.050) (2.750,1.028) (2.800,1.004) (2.850,0.976) (2.900,0.944) (2.950,0.909) (2.993,0.875) (3.050,0.825) (3.099,0.775) (3.143,0.725) (3.175,0.684) (3.217,0.625) (3.247,0.575) (3.275,0.525) (3.299,0.475) (3.320,0.425) (3.339,0.375) (3.355,0.325) (3.369,0.275) (3.380,0.225) (3.388,0.175) (3.395,0.125) (3.399,0.075) (3.401,0.025) (3.401,-0.050) (3.397,-0.100) (3.392,-0.150) (3.384,-0.200) (3.374,-0.250) (3.362,-0.300) (3.347,-0.350) (3.330,-0.400) (3.310,-0.450) (3.287,-0.500) (3.262,-0.550) (3.233,-0.600) (3.200,-0.650) (3.163,-0.700) (3.125,-0.746) (3.075,-0.800) (3.025,-0.848) (2.975,-0.890) (2.928,-0.925) (2.875,-0.960) (2.825,-0.990) (2.775,-1.016) (2.725,-1.040) (2.675,-1.060) (2.625,-1.078) (2.575,-1.093) (2.525,-1.106) (2.475,-1.117) (2.425,-1.126) (2.375,-1.133) (2.325,-1.138) (2.275,-1.141) (2.225,-1.142) (2.175,-1.141) (2.125,-1.139) (2.075,-1.135) (2.025,-1.130) (1.975,-1.123) (1.925,-1.114) (1.875,-1.104) (1.825,-1.093) (1.775,-1.080) (1.725,-1.066) (1.675,-1.051) (1.625,-1.034) (1.575,-1.017) (1.525,-0.998) (1.475,-0.978) (1.425,-0.957) (1.375,-0.935) (1.325,-0.912) (1.275,-0.888) (1.225,-0.864) (1.175,-0.839) (1.125,-0.813) (1.075,-0.786) (1.025,-0.759) (0.975,-0.732) (0.925,-0.704) (0.875,-0.676) (0.830,-0.650) (0.786,-0.625) (0.742,-0.600) (0.697,-0.575) (0.652,-0.550) (0.607,-0.525) (0.561,-0.500) (0.514,-0.475) (0.465,-0.450) (0.414,-0.425) (0.360,-0.400) (0.300,-0.375) (0.250,-0.356) (0.200,-0.340) (0.150,-0.327) (0.100,-0.317) (0.050,-0.311) (0.000,-0.309) (-0.050,-0.311) (-0.100,-0.317) (-0.150,-0.327) (-0.200,-0.340) (-0.250,-0.356) (-0.300,-0.375) (-0.350,-0.396) (-0.400,-0.419) (-0.450,-0.443) (-0.500,-0.468) (-0.550,-0.495) (-0.600,-0.522) (-0.650,-0.549) (-0.697,-0.575) (-0.742,-0.600) (-0.786,-0.625) (-0.830,-0.650) (-0.874,-0.675) (-0.919,-0.700) (-0.964,-0.725) (-1.009,-0.750) (-1.055,-0.775) (-1.100,-0.799) (-1.149,-0.825) (-1.198,-0.850) (-1.248,-0.875) (-1.300,-0.900) (-1.350,-0.924) (-1.400,-0.946) (-1.450,-0.967) (-1.500,-0.988) (-1.550,-1.007) (-1.599,-1.025) (-1.650,-1.043) (-1.700,-1.059) (-1.750,-1.073) (-1.800,-1.087) (-1.850,-1.099) (-1.900,-1.109) (-1.950,-1.119) (-2.000,-1.126) (-2.050,-1.133) (-2.100,-1.137) (-2.150,-1.140) (-2.200,-1.142) (-2.250,-1.141) (-2.300,-1.139) (-2.350,-1.135) (-2.400,-1.130) (-2.425,-1.126)}
\EQ{(-2.425,-1.176) (-2.475,-1.171) (-2.525,-1.164) (-2.575,-1.154) (-2.625,-1.143) (-2.675,-1.130) (-2.725,-1.114) (-2.775,-1.096) (-2.825,-1.076) (-2.875,-1.052) (-2.925,-1.026) (-2.970,-1.000) (-3.025,-0.964) (-3.075,-0.927) (-3.125,-0.886) (-3.164,-0.850) (-3.200,-0.814) (-3.236,-0.775) (-3.275,-0.728) (-3.314,-0.675) (-3.347,-0.625) (-3.375,-0.577) (-3.400,-0.530) (-3.425,-0.477) (-3.446,-0.425) (-3.464,-0.375) (-3.480,-0.325) (-3.493,-0.275) (-3.503,-0.225) (-3.512,-0.175) (-3.518,-0.125) (-3.522,-0.075) (-3.524,-0.025) (-3.524,0.050) (-3.520,0.100) (-3.515,0.150) (-3.508,0.200) (-3.498,0.250) (-3.486,0.300) (-3.472,0.350) (-3.456,0.400) (-3.436,0.450) (-3.415,0.500) (-3.390,0.550) (-3.362,0.600) (-3.331,0.650) (-3.300,0.695) (-3.257,0.750) (-3.214,0.800) (-3.175,0.840) (-3.137,0.875) (-3.078,0.925) (-3.025,0.964) (-2.975,0.997) (-2.927,1.025) (-2.880,1.050) (-2.826,1.075) (-2.775,1.096) (-2.725,1.114) (-2.675,1.130) (-2.625,1.143) (-2.575,1.154) (-2.525,1.164) (-2.475,1.171) (-2.425,1.176) (-2.375,1.180) (-2.325,1.181) (-2.275,1.181) (-2.225,1.180) (-2.175,1.177) (-2.125,1.172) (-2.075,1.166) (-2.025,1.158) (-1.975,1.149) (-1.925,1.138) (-1.875,1.126) (-1.825,1.113) (-1.775,1.099) (-1.725,1.083) (-1.675,1.066) (-1.625,1.048) (-1.575,1.029) (-1.525,1.009) (-1.475,0.988) (-1.425,0.966) (-1.375,0.943) (-1.325,0.919) (-1.275,0.895) (-1.225,0.869) (-1.175,0.843) (-1.125,0.817) (-1.075,0.790) (-1.025,0.762) (-0.975,0.734) (-0.925,0.706) (-0.875,0.677) (-0.827,0.650) (-0.784,0.625) (-0.740,0.600) (-0.696,0.575) (-0.651,0.550) (-0.606,0.525) (-0.560,0.500) (-0.513,0.475) (-0.465,0.450) (-0.414,0.425) (-0.360,0.400) (-0.300,0.375) (-0.250,0.356) (-0.200,0.340) (-0.150,0.327) (-0.100,0.317) (-0.050,0.311) (0.000,0.309) (0.050,0.311) (0.100,0.317) (0.150,0.327) (0.200,0.340) (0.250,0.356) (0.300,0.375) (0.350,0.396) (0.400,0.419) (0.450,0.443) (0.500,0.469) (0.550,0.495) (0.600,0.522) (0.650,0.550) (0.696,0.575) (0.740,0.600) (0.784,0.625) (0.827,0.650) (0.871,0.675) (0.915,0.700) (0.959,0.725) (1.004,0.750) (1.049,0.775) (1.095,0.800) (1.141,0.825) (1.188,0.850) (1.237,0.875) (1.287,0.900) (1.338,0.925) (1.391,0.950) (1.446,0.975) (1.500,0.999) (1.550,1.019) (1.600,1.039) (1.650,1.057) (1.700,1.075) (1.750,1.091) (1.800,1.106) (1.850,1.120) (1.900,1.132) (1.950,1.144) (2.000,1.153) (2.050,1.162) (2.100,1.169) (2.150,1.174) (2.200,1.178) (2.250,1.181) (2.300,1.182) (2.350,1.181) (2.400,1.178) (2.450,1.174) (2.500,1.167) (2.550,1.159) (2.600,1.149) (2.650,1.137) (2.700,1.122) (2.750,1.105) (2.800,1.086) (2.850,1.064) (2.900,1.040) (2.950,1.012) (3.000,0.981) (3.045,0.950) (3.100,0.907) (3.150,0.863) (3.190,0.825) (3.225,0.787) (3.275,0.728) (3.314,0.675) (3.347,0.625) (3.375,0.577) (3.400,0.530) (3.425,0.477) (3.446,0.425) (3.464,0.375) (3.480,0.325) (3.493,0.275) (3.503,0.225) (3.512,0.175) (3.518,0.125) (3.522,0.075) (3.524,0.025) (3.524,-0.050) (3.520,-0.100) (3.515,-0.150) (3.508,-0.200) (3.498,-0.250) (3.486,-0.300) (3.472,-0.350) (3.456,-0.400) (3.436,-0.450) (3.415,-0.500) (3.390,-0.550) (3.362,-0.600) (3.331,-0.650) (3.300,-0.695) (3.257,-0.750) (3.214,-0.800) (3.175,-0.840) (3.137,-0.875) (3.078,-0.925) (3.025,-0.964) (2.975,-0.997) (2.927,-1.025) (2.880,-1.050) (2.826,-1.075) (2.775,-1.096) (2.725,-1.114) (2.675,-1.130) (2.625,-1.143) (2.575,-1.154) (2.525,-1.164) (2.475,-1.171) (2.425,-1.176) (2.375,-1.180) (2.325,-1.181) (2.275,-1.181) (2.225,-1.180) (2.175,-1.177) (2.125,-1.172) (2.075,-1.166) (2.025,-1.158) (1.975,-1.149) (1.925,-1.138) (1.875,-1.126) (1.825,-1.113) (1.775,-1.099) (1.725,-1.083) (1.675,-1.066) (1.625,-1.048) (1.575,-1.029) (1.525,-1.009) (1.475,-0.988) (1.425,-0.966) (1.375,-0.943) (1.325,-0.919) (1.275,-0.895) (1.225,-0.869) (1.175,-0.843) (1.125,-0.817) (1.075,-0.790) (1.025,-0.762) (0.975,-0.734) (0.925,-0.706) (0.875,-0.677) (0.827,-0.650) (0.784,-0.625) (0.740,-0.600) (0.696,-0.575) (0.651,-0.550) (0.606,-0.525) (0.560,-0.500) (0.513,-0.475) (0.465,-0.450) (0.414,-0.425) (0.360,-0.400) (0.300,-0.375) (0.250,-0.356) (0.200,-0.340) (0.150,-0.327) (0.100,-0.317) (0.050,-0.311) (0.000,-0.309) (-0.050,-0.311) (-0.100,-0.317) (-0.150,-0.327) (-0.200,-0.340) (-0.250,-0.356) (-0.300,-0.375) (-0.350,-0.396) (-0.400,-0.419) (-0.450,-0.443) (-0.500,-0.469) (-0.550,-0.495) (-0.600,-0.522) (-0.650,-0.550) (-0.696,-0.575) (-0.740,-0.600) (-0.784,-0.625) (-0.827,-0.650) (-0.871,-0.675) (-0.915,-0.700) (-0.959,-0.725) (-1.004,-0.750) (-1.049,-0.775) (-1.095,-0.800) (-1.141,-0.825) (-1.188,-0.850) (-1.237,-0.875) (-1.287,-0.900) (-1.338,-0.925) (-1.391,-0.950) (-1.446,-0.975) (-1.500,-0.999) (-1.550,-1.019) (-1.600,-1.039) (-1.650,-1.057) (-1.700,-1.075) (-1.750,-1.091) (-1.800,-1.106) (-1.850,-1.120) (-1.900,-1.132) (-1.950,-1.144) (-2.000,-1.153) (-2.050,-1.162) (-2.100,-1.169) (-2.150,-1.174) (-2.200,-1.178) (-2.250,-1.181) (-2.300,-1.182) (-2.350,-1.181) (-2.400,-1.178) (-2.425,-1.176)}
\EQ{(-2.475,-1.226) (-2.525,-1.222) (-2.575,-1.217) (-2.625,-1.210) (-2.675,-1.201) (-2.725,-1.190) (-2.775,-1.177) (-2.825,-1.161) (-2.875,-1.144) (-2.923,-1.125) (-2.975,-1.102) (-3.025,-1.077) (-3.073,-1.050) (-3.125,-1.017) (-3.175,-0.982) (-3.217,-0.950) (-3.275,-0.900) (-3.325,-0.851) (-3.372,-0.800) (-3.413,-0.750) (-3.449,-0.700) (-3.482,-0.650) (-3.512,-0.600) (-3.538,-0.550) (-3.562,-0.500) (-3.582,-0.450) (-3.600,-0.402) (-3.617,-0.350) (-3.630,-0.300) (-3.642,-0.250) (-3.651,-0.200) (-3.658,-0.150) (-3.663,-0.100) (-3.666,-0.050) (-3.667,0.000) (-3.666,0.050) (-3.663,0.100) (-3.658,0.150) (-3.651,0.200) (-3.642,0.250) (-3.630,0.300) (-3.617,0.350) (-3.601,0.400) (-3.582,0.450) (-3.562,0.500) (-3.538,0.550) (-3.512,0.600) (-3.482,0.650) (-3.450,0.699) (-3.413,0.750) (-3.375,0.796) (-3.326,0.850) (-3.275,0.900) (-3.225,0.943) (-3.185,0.975) (-3.125,1.017) (-3.075,1.049) (-3.028,1.075) (-2.979,1.100) (-2.925,1.124) (-2.875,1.144) (-2.825,1.161) (-2.775,1.177) (-2.725,1.190) (-2.675,1.201) (-2.625,1.210) (-2.575,1.217) (-2.525,1.222) (-2.475,1.226) (-2.425,1.227) (-2.375,1.228) (-2.325,1.226) (-2.275,1.223) (-2.225,1.219) (-2.175,1.213) (-2.125,1.206) (-2.075,1.197) (-2.025,1.187) (-1.975,1.176) (-1.925,1.163) (-1.875,1.149) (-1.825,1.134) (-1.775,1.118) (-1.725,1.100) (-1.675,1.082) (-1.625,1.063) (-1.575,1.042) (-1.525,1.021) (-1.479,1.000) (-1.425,0.975) (-1.375,0.951) (-1.325,0.926) (-1.275,0.901) (-1.226,0.875) (-1.179,0.850) (-1.133,0.825) (-1.088,0.800) (-1.043,0.775) (-0.999,0.750) (-0.955,0.725) (-0.912,0.700) (-0.868,0.675) (-0.825,0.650) (-0.775,0.622) (-0.725,0.593) (-0.675,0.565) (-0.625,0.536) (-0.575,0.509) (-0.525,0.482) (-0.475,0.456) (-0.425,0.431) (-0.375,0.408) (-0.325,0.386) (-0.275,0.366) (-0.225,0.348) (-0.175,0.333) (-0.125,0.322) (-0.075,0.313) (-0.025,0.309) (0.050,0.311) (0.100,0.317) (0.150,0.327) (0.200,0.340) (0.250,0.357) (0.300,0.375) (0.350,0.396) (0.400,0.419) (0.450,0.443) (0.500,0.469) (0.550,0.495) (0.600,0.522) (0.650,0.550) (0.694,0.575) (0.738,0.600) (0.782,0.625) (0.825,0.650) (0.875,0.679) (0.925,0.708) (0.975,0.737) (1.025,0.765) (1.075,0.793) (1.125,0.821) (1.175,0.848) (1.225,0.875) (1.274,0.900) (1.323,0.925) (1.373,0.950) (1.425,0.975) (1.475,0.998) (1.525,1.021) (1.575,1.042) (1.625,1.063) (1.675,1.082) (1.724,1.100) (1.775,1.118) (1.825,1.134) (1.875,1.149) (1.925,1.163) (1.975,1.176) (2.025,1.187) (2.075,1.197) (2.125,1.206) (2.175,1.213) (2.225,1.219) (2.275,1.223) (2.325,1.226) (2.375,1.228) (2.425,1.227) (2.475,1.226) (2.525,1.222) (2.575,1.217) (2.625,1.210) (2.675,1.201) (2.725,1.190) (2.775,1.177) (2.825,1.161) (2.875,1.144) (2.923,1.125) (2.975,1.102) (3.025,1.077) (3.073,1.050) (3.125,1.017) (3.175,0.982) (3.217,0.950) (3.275,0.900) (3.325,0.851) (3.372,0.800) (3.413,0.750) (3.449,0.700) (3.482,0.650) (3.512,0.600) (3.538,0.550) (3.562,0.500) (3.582,0.450) (3.600,0.402) (3.617,0.350) (3.630,0.300) (3.642,0.250) (3.651,0.200) (3.658,0.150) (3.663,0.100) (3.666,0.050) (3.667,0.000) (3.666,-0.050) (3.663,-0.100) (3.658,-0.150) (3.651,-0.200) (3.642,-0.250) (3.630,-0.300) (3.617,-0.350) (3.601,-0.400) (3.582,-0.450) (3.562,-0.500) (3.538,-0.550) (3.512,-0.600) (3.482,-0.650) (3.450,-0.699) (3.413,-0.750) (3.375,-0.796) (3.326,-0.850) (3.275,-0.900) (3.225,-0.943) (3.185,-0.975) (3.125,-1.017) (3.075,-1.049) (3.028,-1.075) (2.979,-1.100) (2.925,-1.124) (2.875,-1.144) (2.825,-1.161) (2.775,-1.177) (2.725,-1.190) (2.675,-1.201) (2.625,-1.210) (2.575,-1.217) (2.525,-1.222) (2.475,-1.226) (2.425,-1.227) (2.375,-1.228) (2.325,-1.226) (2.275,-1.223) (2.225,-1.219) (2.175,-1.213) (2.125,-1.206) (2.075,-1.197) (2.025,-1.187) (1.975,-1.176) (1.925,-1.163) (1.875,-1.149) (1.825,-1.134) (1.775,-1.118) (1.725,-1.100) (1.675,-1.082) (1.625,-1.063) (1.575,-1.042) (1.525,-1.021) (1.479,-1.000) (1.425,-0.975) (1.375,-0.951) (1.325,-0.926) (1.275,-0.901) (1.226,-0.875) (1.179,-0.850) (1.133,-0.825) (1.088,-0.800) (1.043,-0.775) (0.999,-0.750) (0.955,-0.725) (0.912,-0.700) (0.868,-0.675) (0.825,-0.650) (0.775,-0.622) (0.725,-0.593) (0.675,-0.565) (0.625,-0.536) (0.575,-0.509) (0.525,-0.482) (0.475,-0.456) (0.425,-0.431) (0.375,-0.408) (0.325,-0.386) (0.275,-0.366) (0.225,-0.348) (0.175,-0.333) (0.125,-0.322) (0.075,-0.313) (0.025,-0.309) (-0.050,-0.311) (-0.100,-0.317) (-0.150,-0.327) (-0.200,-0.340) (-0.250,-0.357) (-0.300,-0.375) (-0.350,-0.396) (-0.400,-0.419) (-0.450,-0.443) (-0.500,-0.469) (-0.550,-0.495) (-0.600,-0.522) (-0.650,-0.550) (-0.694,-0.575) (-0.738,-0.600) (-0.782,-0.625) (-0.825,-0.650) (-0.875,-0.679) (-0.925,-0.708) (-0.975,-0.737) (-1.025,-0.765) (-1.075,-0.793) (-1.125,-0.821) (-1.175,-0.848) (-1.225,-0.875) (-1.274,-0.900) (-1.323,-0.925) (-1.373,-0.950) (-1.425,-0.975) (-1.475,-0.998) (-1.525,-1.021) (-1.575,-1.042) (-1.625,-1.063) (-1.675,-1.082) (-1.724,-1.100) (-1.775,-1.118) (-1.825,-1.134) (-1.875,-1.149) (-1.925,-1.163) (-1.975,-1.176) (-2.025,-1.187) (-2.075,-1.197) (-2.125,-1.206) (-2.175,-1.213) (-2.225,-1.219) (-2.275,-1.223) (-2.325,-1.226) (-2.375,-1.228) (-2.425,-1.227) (-2.475,-1.226)}
\EQ{(-2.650,-1.276) (-2.700,-1.270) (-2.750,-1.263) (-2.800,-1.254) (-2.850,-1.243) (-2.900,-1.230) (-2.950,-1.215) (-3.000,-1.198) (-3.050,-1.179) (-3.100,-1.157) (-3.150,-1.132) (-3.200,-1.105) (-3.250,-1.075) (-3.300,-1.041) (-3.350,-1.004) (-3.400,-0.963) (-3.442,-0.925) (-3.492,-0.875) (-3.536,-0.825) (-3.575,-0.777) (-3.612,-0.725) (-3.645,-0.675) (-3.674,-0.625) (-3.700,-0.576) (-3.724,-0.525) (-3.745,-0.475) (-3.764,-0.425) (-3.780,-0.375) (-3.794,-0.325) (-3.806,-0.275) (-3.816,-0.225) (-3.824,-0.175) (-3.830,-0.125) (-3.833,-0.075) (-3.835,-0.025) (-3.835,0.050) (-3.832,0.100) (-3.827,0.150) (-3.820,0.200) (-3.811,0.250) (-3.800,0.300) (-3.787,0.350) (-3.772,0.400) (-3.755,0.450) (-3.735,0.500) (-3.713,0.550) (-3.688,0.600) (-3.660,0.650) (-3.629,0.700) (-3.600,0.743) (-3.557,0.800) (-3.515,0.850) (-3.475,0.892) (-3.425,0.941) (-3.386,0.975) (-3.325,1.023) (-3.275,1.059) (-3.225,1.090) (-3.175,1.119) (-3.125,1.145) (-3.075,1.168) (-3.025,1.189) (-2.975,1.207) (-2.925,1.223) (-2.875,1.237) (-2.825,1.249) (-2.775,1.259) (-2.725,1.267) (-2.675,1.273) (-2.625,1.278) (-2.575,1.281) (-2.525,1.282) (-2.475,1.282) (-2.425,1.280) (-2.375,1.277) (-2.325,1.273) (-2.275,1.267) (-2.225,1.259) (-2.175,1.251) (-2.125,1.241) (-2.075,1.229) (-2.025,1.217) (-1.975,1.203) (-1.925,1.188) (-1.875,1.172) (-1.825,1.155) (-1.775,1.137) (-1.725,1.118) (-1.675,1.098) (-1.625,1.077) (-1.575,1.055) (-1.525,1.032) (-1.475,1.009) (-1.425,0.984) (-1.375,0.959) (-1.325,0.934) (-1.275,0.907) (-1.225,0.880) (-1.175,0.853) (-1.126,0.825) (-1.082,0.800) (-1.038,0.775) (-0.994,0.750) (-0.951,0.725) (-0.900,0.695) (-0.850,0.666) (-0.800,0.637) (-0.750,0.608) (-0.700,0.579) (-0.650,0.551) (-0.604,0.525) (-0.559,0.500) (-0.512,0.475) (-0.464,0.450) (-0.413,0.425) (-0.359,0.400) (-0.300,0.375) (-0.250,0.357) (-0.200,0.340) (-0.150,0.327) (-0.100,0.317) (-0.050,0.311) (0.000,0.309) (0.050,0.311) (0.100,0.317) (0.150,0.327) (0.200,0.340) (0.250,0.357) (0.299,0.375) (0.350,0.396) (0.400,0.419) (0.450,0.444) (0.500,0.469) (0.550,0.496) (0.600,0.523) (0.649,0.550) (0.693,0.575) (0.737,0.600) (0.800,0.637) (0.850,0.666) (0.900,0.695) (0.950,0.724) (0.994,0.750) (1.038,0.775) (1.082,0.800) (1.126,0.825) (1.171,0.850) (1.216,0.875) (1.262,0.900) (1.309,0.925) (1.357,0.950) (1.407,0.975) (1.457,1.000) (1.510,1.025) (1.564,1.050) (1.621,1.075) (1.675,1.098) (1.725,1.118) (1.775,1.137) (1.825,1.155) (1.875,1.172) (1.925,1.188) (1.975,1.203) (2.025,1.217) (2.075,1.229) (2.125,1.241) (2.175,1.251) (2.225,1.259) (2.275,1.267) (2.325,1.273) (2.375,1.277) (2.425,1.280) (2.475,1.282) (2.525,1.282) (2.575,1.281) (2.625,1.278) (2.675,1.273) (2.725,1.267) (2.775,1.259) (2.825,1.249) (2.875,1.237) (2.925,1.223) (2.975,1.207) (3.025,1.189) (3.075,1.168) (3.125,1.145) (3.175,1.119) (3.225,1.090) (3.275,1.059) (3.323,1.025) (3.375,0.984) (3.415,0.950) (3.468,0.900) (3.515,0.850) (3.550,0.808) (3.595,0.750) (3.625,0.706) (3.660,0.650) (3.688,0.600) (3.713,0.550) (3.735,0.500) (3.755,0.450) (3.772,0.400) (3.787,0.350) (3.800,0.301) (3.811,0.250) (3.820,0.200) (3.827,0.150) (3.832,0.100) (3.835,0.050) (3.836,0.000) (3.835,-0.050) (3.832,-0.100) (3.827,-0.150) (3.820,-0.200) (3.811,-0.250) (3.800,-0.300) (3.787,-0.350) (3.772,-0.400) (3.755,-0.450) (3.735,-0.500) (3.713,-0.550) (3.688,-0.600) (3.660,-0.650) (3.629,-0.700) (3.600,-0.743) (3.557,-0.800) (3.515,-0.850) (3.475,-0.892) (3.425,-0.941) (3.386,-0.975) (3.325,-1.023) (3.275,-1.059) (3.225,-1.090) (3.175,-1.119) (3.125,-1.145) (3.075,-1.168) (3.025,-1.189) (2.975,-1.207) (2.925,-1.223) (2.875,-1.237) (2.825,-1.249) (2.775,-1.259) (2.725,-1.267) (2.675,-1.273) (2.625,-1.278) (2.575,-1.281) (2.525,-1.282) (2.475,-1.282) (2.425,-1.280) (2.375,-1.277) (2.325,-1.273) (2.275,-1.267) (2.225,-1.259) (2.175,-1.251) (2.125,-1.241) (2.075,-1.229) (2.025,-1.217) (1.975,-1.203) (1.925,-1.188) (1.875,-1.172) (1.825,-1.155) (1.775,-1.137) (1.725,-1.118) (1.675,-1.098) (1.625,-1.077) (1.575,-1.055) (1.525,-1.032) (1.475,-1.009) (1.425,-0.984) (1.375,-0.959) (1.325,-0.934) (1.275,-0.907) (1.225,-0.880) (1.175,-0.853) (1.126,-0.825) (1.082,-0.800) (1.038,-0.775) (0.994,-0.750) (0.951,-0.725) (0.900,-0.695) (0.850,-0.666) (0.800,-0.637) (0.750,-0.608) (0.700,-0.579) (0.650,-0.551) (0.604,-0.525) (0.559,-0.500) (0.512,-0.475) (0.464,-0.450) (0.413,-0.425) (0.359,-0.400) (0.300,-0.375) (0.250,-0.357) (0.200,-0.340) (0.150,-0.327) (0.100,-0.317) (0.050,-0.311) (0.000,-0.309) (-0.050,-0.311) (-0.100,-0.317) (-0.150,-0.327) (-0.200,-0.340) (-0.250,-0.357) (-0.299,-0.375) (-0.350,-0.396) (-0.400,-0.419) (-0.450,-0.444) (-0.500,-0.469) (-0.550,-0.496) (-0.600,-0.523) (-0.649,-0.550) (-0.693,-0.575) (-0.737,-0.600) (-0.800,-0.637) (-0.850,-0.666) (-0.900,-0.695) (-0.950,-0.724) (-0.994,-0.750) (-1.038,-0.775) (-1.082,-0.800) (-1.126,-0.825) (-1.171,-0.850) (-1.216,-0.875) (-1.262,-0.900) (-1.309,-0.925) (-1.357,-0.950) (-1.407,-0.975) (-1.457,-1.000) (-1.510,-1.025) (-1.564,-1.050) (-1.621,-1.075) (-1.675,-1.098) (-1.725,-1.118) (-1.775,-1.137) (-1.825,-1.155) (-1.875,-1.172) (-1.925,-1.188) (-1.975,-1.203) (-2.025,-1.217) (-2.075,-1.229) (-2.125,-1.241) (-2.175,-1.251) (-2.225,-1.259) (-2.275,-1.267) (-2.325,-1.273) (-2.375,-1.277) (-2.425,-1.280) (-2.475,-1.282) (-2.525,-1.282) (-2.575,-1.281) (-2.625,-1.278) (-2.650,-1.276)}
\EQ{(-2.900,-1.328) (-2.950,-1.319) (-3.000,-1.308) (-3.050,-1.295) (-3.100,-1.281) (-3.150,-1.264) (-3.200,-1.245) (-3.247,-1.225) (-3.300,-1.200) (-3.348,-1.175) (-3.392,-1.150) (-3.450,-1.113) (-3.500,-1.078) (-3.550,-1.039) (-3.596,-1.000) (-3.649,-0.950) (-3.696,-0.900) (-3.739,-0.850) (-3.775,-0.804) (-3.813,-0.750) (-3.845,-0.700) (-3.874,-0.650) (-3.900,-0.600) (-3.923,-0.550) (-3.944,-0.500) (-3.963,-0.450) (-3.979,-0.400) (-3.994,-0.350) (-4.006,-0.300) (-4.016,-0.250) (-4.025,-0.200) (-4.031,-0.150) (-4.036,-0.100) (-4.039,-0.050) (-4.039,0.000) (-4.039,0.050) (-4.036,0.100) (-4.031,0.150) (-4.025,0.200) (-4.016,0.250) (-4.006,0.300) (-3.994,0.350) (-3.979,0.400) (-3.963,0.450) (-3.944,0.500) (-3.923,0.550) (-3.900,0.599) (-3.875,0.648) (-3.850,0.691) (-3.813,0.750) (-3.778,0.800) (-3.739,0.850) (-3.700,0.896) (-3.650,0.949) (-3.600,0.996) (-3.550,1.039) (-3.504,1.075) (-3.450,1.113) (-3.400,1.145) (-3.350,1.174) (-3.300,1.200) (-3.250,1.224) (-3.200,1.245) (-3.150,1.264) (-3.100,1.281) (-3.050,1.295) (-3.000,1.308) (-2.950,1.319) (-2.900,1.328) (-2.850,1.335) (-2.800,1.341) (-2.750,1.345) (-2.700,1.348) (-2.650,1.349) (-2.600,1.348) (-2.550,1.346) (-2.500,1.343) (-2.450,1.338) (-2.400,1.332) (-2.350,1.325) (-2.300,1.316) (-2.250,1.307) (-2.200,1.296) (-2.150,1.283) (-2.100,1.270) (-2.050,1.256) (-2.000,1.240) (-1.950,1.223) (-1.900,1.206) (-1.850,1.187) (-1.800,1.168) (-1.750,1.147) (-1.700,1.126) (-1.650,1.103) (-1.600,1.080) (-1.550,1.057) (-1.500,1.032) (-1.450,1.007) (-1.400,0.981) (-1.350,0.954) (-1.300,0.927) (-1.251,0.900) (-1.206,0.875) (-1.162,0.850) (-1.119,0.825) (-1.075,0.800) (-1.025,0.771) (-0.975,0.742) (-0.925,0.712) (-0.875,0.683) (-0.825,0.653) (-0.778,0.625) (-0.725,0.595) (-0.675,0.566) (-0.625,0.538) (-0.575,0.510) (-0.525,0.483) (-0.475,0.456) (-0.425,0.431) (-0.375,0.408) (-0.325,0.386) (-0.275,0.366) (-0.225,0.348) (-0.175,0.333) (-0.125,0.322) (-0.075,0.313) (-0.025,0.309) (0.050,0.311) (0.100,0.317) (0.150,0.327) (0.200,0.340) (0.250,0.357) (0.299,0.375) (0.350,0.396) (0.400,0.419) (0.450,0.444) (0.500,0.469) (0.550,0.496) (0.600,0.523) (0.648,0.550) (0.692,0.575) (0.750,0.609) (0.800,0.638) (0.850,0.668) (0.900,0.697) (0.947,0.725) (1.000,0.756) (1.050,0.786) (1.100,0.815) (1.150,0.843) (1.200,0.872) (1.250,0.900) (1.296,0.925) (1.343,0.950) (1.390,0.975) (1.437,1.000) (1.487,1.025) (1.537,1.050) (1.589,1.075) (1.643,1.100) (1.699,1.125) (1.750,1.147) (1.800,1.168) (1.850,1.187) (1.900,1.206) (1.950,1.223) (2.000,1.240) (2.050,1.256) (2.100,1.270) (2.150,1.283) (2.200,1.296) (2.250,1.307) (2.300,1.316) (2.350,1.325) (2.400,1.332) (2.450,1.338) (2.500,1.343) (2.550,1.346) (2.600,1.348) (2.650,1.349) (2.700,1.348) (2.750,1.345) (2.800,1.341) (2.850,1.335) (2.900,1.328) (2.950,1.319) (3.000,1.308) (3.050,1.295) (3.100,1.281) (3.150,1.264) (3.200,1.245) (3.247,1.225) (3.300,1.200) (3.348,1.175) (3.392,1.150) (3.450,1.113) (3.500,1.078) (3.550,1.039) (3.596,1.000) (3.649,0.950) (3.696,0.900) (3.739,0.850) (3.775,0.804) (3.813,0.750) (3.845,0.700) (3.874,0.650) (3.900,0.600) (3.923,0.550) (3.944,0.500) (3.963,0.450) (3.979,0.400) (3.994,0.350) (4.006,0.300) (4.016,0.250) (4.025,0.200) (4.031,0.150) (4.036,0.100) (4.039,0.050) (4.039,0.000) (4.039,-0.050) (4.036,-0.100) (4.031,-0.150) (4.025,-0.200) (4.016,-0.250) (4.006,-0.300) (3.994,-0.350) (3.979,-0.400) (3.963,-0.450) (3.944,-0.500) (3.923,-0.550) (3.900,-0.599) (3.875,-0.648) (3.850,-0.691) (3.813,-0.750) (3.778,-0.800) (3.739,-0.850) (3.700,-0.896) (3.650,-0.949) (3.600,-0.996) (3.550,-1.039) (3.504,-1.075) (3.450,-1.113) (3.400,-1.145) (3.350,-1.174) (3.300,-1.200) (3.250,-1.224) (3.200,-1.245) (3.150,-1.264) (3.100,-1.281) (3.050,-1.295) (3.000,-1.308) (2.950,-1.319) (2.900,-1.328) (2.850,-1.335) (2.800,-1.341) (2.750,-1.345) (2.700,-1.348) (2.650,-1.349) (2.600,-1.348) (2.550,-1.346) (2.500,-1.343) (2.450,-1.338) (2.400,-1.332) (2.350,-1.325) (2.300,-1.316) (2.250,-1.307) (2.200,-1.296) (2.150,-1.283) (2.100,-1.270) (2.050,-1.256) (2.000,-1.240) (1.950,-1.223) (1.900,-1.206) (1.850,-1.187) (1.800,-1.168) (1.750,-1.147) (1.700,-1.126) (1.650,-1.103) (1.600,-1.080) (1.550,-1.057) (1.500,-1.032) (1.450,-1.007) (1.400,-0.981) (1.350,-0.954) (1.300,-0.927) (1.251,-0.900) (1.206,-0.875) (1.162,-0.850) (1.119,-0.825) (1.075,-0.800) (1.025,-0.771) (0.975,-0.742) (0.925,-0.712) (0.875,-0.683) (0.825,-0.653) (0.778,-0.625) (0.725,-0.595) (0.675,-0.566) (0.625,-0.538) (0.575,-0.510) (0.525,-0.483) (0.475,-0.456) (0.425,-0.431) (0.375,-0.408) (0.325,-0.386) (0.275,-0.366) (0.225,-0.348) (0.175,-0.333) (0.125,-0.322) (0.075,-0.313) (0.025,-0.309) (-0.050,-0.311) (-0.100,-0.317) (-0.150,-0.327) (-0.200,-0.340) (-0.250,-0.357) (-0.299,-0.375) (-0.350,-0.396) (-0.400,-0.419) (-0.450,-0.444) (-0.500,-0.469) (-0.550,-0.496) (-0.600,-0.523) (-0.648,-0.550) (-0.692,-0.575) (-0.750,-0.609) (-0.800,-0.638) (-0.850,-0.668) (-0.900,-0.697) (-0.947,-0.725) (-1.000,-0.756) (-1.050,-0.786) (-1.100,-0.815) (-1.150,-0.843) (-1.200,-0.872) (-1.250,-0.900) (-1.296,-0.925) (-1.343,-0.950) (-1.390,-0.975) (-1.437,-1.000) (-1.487,-1.025) (-1.537,-1.050) (-1.589,-1.075) (-1.643,-1.100) (-1.699,-1.125) (-1.750,-1.147) (-1.800,-1.168) (-1.850,-1.187) (-1.900,-1.206) (-1.950,-1.223) (-2.000,-1.240) (-2.050,-1.256) (-2.100,-1.270) (-2.150,-1.283) (-2.200,-1.296) (-2.250,-1.307) (-2.300,-1.316) (-2.350,-1.325) (-2.400,-1.332) (-2.450,-1.338) (-2.500,-1.343) (-2.550,-1.346) (-2.600,-1.348) (-2.650,-1.349) (-2.700,-1.348) (-2.750,-1.345) (-2.800,-1.341) (-2.850,-1.335) (-2.900,-1.328)}
\EQ{(-2.950,-1.426) (-3.000,-1.421) (-3.050,-1.415) (-3.100,-1.407) (-3.150,-1.398) (-3.200,-1.386) (-3.250,-1.373) (-3.300,-1.358) (-3.350,-1.342) (-3.400,-1.323) (-3.450,-1.302) (-3.500,-1.279) (-3.550,-1.253) (-3.600,-1.225) (-3.650,-1.194) (-3.700,-1.160) (-3.748,-1.125) (-3.800,-1.082) (-3.850,-1.037) (-3.888,-1.000) (-3.925,-0.961) (-3.975,-0.903) (-4.016,-0.850) (-4.050,-0.802) (-4.084,-0.750) (-4.113,-0.700) (-4.140,-0.650) (-4.164,-0.600) (-4.186,-0.550) (-4.205,-0.500) (-4.223,-0.450) (-4.238,-0.400) (-4.251,-0.350) (-4.263,-0.300) (-4.273,-0.250) (-4.281,-0.200) (-4.287,-0.150) (-4.291,-0.100) (-4.294,-0.050) (-4.294,0.000) (-4.294,0.050) (-4.291,0.100) (-4.287,0.150) (-4.281,0.200) (-4.273,0.250) (-4.263,0.300) (-4.251,0.350) (-4.238,0.400) (-4.223,0.450) (-4.205,0.500) (-4.186,0.550) (-4.164,0.600) (-4.140,0.650) (-4.113,0.700) (-4.084,0.750) (-4.052,0.800) (-4.016,0.850) (-3.978,0.900) (-3.935,0.950) (-3.900,0.988) (-3.863,1.025) (-3.825,1.060) (-3.779,1.100) (-3.725,1.142) (-3.679,1.175) (-3.625,1.210) (-3.575,1.240) (-3.525,1.266) (-3.475,1.291) (-3.425,1.313) (-3.375,1.333) (-3.326,1.350) (-3.275,1.366) (-3.225,1.380) (-3.175,1.392) (-3.125,1.403) (-3.075,1.411) (-3.025,1.418) (-2.975,1.424) (-2.925,1.428) (-2.875,1.430) (-2.825,1.432) (-2.775,1.431) (-2.725,1.430) (-2.675,1.427) (-2.625,1.422) (-2.575,1.417) (-2.525,1.410) (-2.475,1.402) (-2.425,1.393) (-2.375,1.383) (-2.325,1.371) (-2.275,1.359) (-2.225,1.345) (-2.175,1.330) (-2.125,1.315) (-2.075,1.298) (-2.025,1.280) (-1.975,1.262) (-1.925,1.242) (-1.875,1.222) (-1.825,1.200) (-1.775,1.178) (-1.725,1.155) (-1.675,1.132) (-1.625,1.107) (-1.575,1.082) (-1.525,1.057) (-1.475,1.030) (-1.425,1.003) (-1.375,0.976) (-1.328,0.950) (-1.284,0.925) (-1.241,0.900) (-1.197,0.875) (-1.150,0.848) (-1.100,0.818) (-1.050,0.789) (-1.000,0.759) (-0.950,0.729) (-0.902,0.700) (-0.850,0.669) (-0.800,0.640) (-0.750,0.610) (-0.700,0.581) (-0.650,0.552) (-0.603,0.525) (-0.558,0.500) (-0.511,0.475) (-0.463,0.450) (-0.413,0.425) (-0.359,0.400) (-0.300,0.375) (-0.250,0.357) (-0.200,0.340) (-0.150,0.327) (-0.100,0.317) (-0.050,0.311) (0.000,0.309) (0.050,0.311) (0.100,0.317) (0.150,0.327) (0.200,0.340) (0.250,0.357) (0.299,0.375) (0.350,0.397) (0.400,0.420) (0.450,0.444) (0.500,0.470) (0.550,0.496) (0.600,0.524) (0.647,0.550) (0.691,0.575) (0.750,0.610) (0.800,0.640) (0.850,0.669) (0.900,0.699) (0.944,0.725) (1.000,0.759) (1.050,0.789) (1.100,0.818) (1.150,0.848) (1.197,0.875) (1.241,0.900) (1.284,0.925) (1.328,0.950) (1.373,0.975) (1.419,1.000) (1.465,1.025) (1.513,1.050) (1.561,1.075) (1.611,1.100) (1.662,1.125) (1.714,1.150) (1.768,1.175) (1.825,1.200) (1.875,1.222) (1.925,1.242) (1.975,1.262) (2.025,1.280) (2.075,1.298) (2.125,1.315) (2.175,1.330) (2.225,1.345) (2.275,1.359) (2.325,1.371) (2.375,1.383) (2.425,1.393) (2.475,1.402) (2.525,1.410) (2.575,1.417) (2.625,1.422) (2.675,1.427) (2.725,1.430) (2.775,1.431) (2.825,1.432) (2.875,1.430) (2.925,1.428) (2.975,1.424) (3.025,1.418) (3.075,1.411) (3.125,1.403) (3.175,1.392) (3.225,1.380) (3.275,1.366) (3.325,1.350) (3.375,1.333) (3.425,1.313) (3.475,1.291) (3.525,1.266) (3.575,1.240) (3.625,1.210) (3.675,1.178) (3.725,1.142) (3.775,1.103) (3.825,1.060) (3.863,1.025) (3.900,0.988) (3.935,0.950) (3.975,0.903) (4.016,0.850) (4.050,0.802) (4.084,0.750) (4.113,0.700) (4.140,0.650) (4.164,0.600) (4.186,0.550) (4.205,0.500) (4.223,0.450) (4.238,0.400) (4.251,0.350) (4.263,0.300) (4.273,0.250) (4.281,0.200) (4.287,0.150) (4.291,0.100) (4.294,0.050) (4.294,0.000) (4.294,-0.050) (4.291,-0.100) (4.287,-0.150) (4.281,-0.200) (4.273,-0.250) (4.263,-0.300) (4.251,-0.350) (4.238,-0.400) (4.223,-0.450) (4.205,-0.500) (4.186,-0.550) (4.164,-0.600) (4.140,-0.650) (4.113,-0.700) (4.084,-0.750) (4.052,-0.800) (4.016,-0.850) (3.978,-0.900) (3.935,-0.950) (3.900,-0.988) (3.863,-1.025) (3.825,-1.060) (3.779,-1.100) (3.725,-1.142) (3.679,-1.175) (3.625,-1.210) (3.575,-1.240) (3.525,-1.266) (3.475,-1.291) (3.425,-1.313) (3.375,-1.333) (3.326,-1.350) (3.275,-1.366) (3.225,-1.380) (3.175,-1.392) (3.125,-1.403) (3.075,-1.411) (3.025,-1.418) (2.975,-1.424) (2.925,-1.428) (2.875,-1.430) (2.825,-1.432) (2.775,-1.431) (2.725,-1.430) (2.675,-1.427) (2.625,-1.422) (2.575,-1.417) (2.525,-1.410) (2.475,-1.402) (2.425,-1.393) (2.375,-1.383) (2.325,-1.371) (2.275,-1.359) (2.225,-1.345) (2.175,-1.330) (2.125,-1.315) (2.075,-1.298) (2.025,-1.280) (1.975,-1.262) (1.925,-1.242) (1.875,-1.222) (1.825,-1.200) (1.775,-1.178) (1.725,-1.155) (1.675,-1.132) (1.625,-1.107) (1.575,-1.082) (1.525,-1.057) (1.475,-1.030) (1.425,-1.003) (1.375,-0.976) (1.328,-0.950) (1.284,-0.925) (1.241,-0.900) (1.197,-0.875) (1.150,-0.848) (1.100,-0.818) (1.050,-0.789) (1.000,-0.759) (0.950,-0.729) (0.902,-0.700) (0.850,-0.669) (0.800,-0.640) (0.750,-0.610) (0.700,-0.581) (0.650,-0.552) (0.603,-0.525) (0.558,-0.500) (0.511,-0.475) (0.463,-0.450) (0.413,-0.425) (0.359,-0.400) (0.300,-0.375) (0.250,-0.357) (0.200,-0.340) (0.150,-0.327) (0.100,-0.317) (0.050,-0.311) (0.000,-0.309) (-0.050,-0.311) (-0.100,-0.317) (-0.150,-0.327) (-0.200,-0.340) (-0.250,-0.357) (-0.299,-0.375) (-0.350,-0.397) (-0.400,-0.420) (-0.450,-0.444) (-0.500,-0.470) (-0.550,-0.496) (-0.600,-0.524) (-0.647,-0.550) (-0.691,-0.575) (-0.750,-0.610) (-0.800,-0.640) (-0.850,-0.669) (-0.900,-0.699) (-0.944,-0.725) (-1.000,-0.759) (-1.050,-0.789) (-1.100,-0.818) (-1.150,-0.848) (-1.197,-0.875) (-1.241,-0.900) (-1.284,-0.925) (-1.328,-0.950) (-1.373,-0.975) (-1.419,-1.000) (-1.465,-1.025) (-1.513,-1.050) (-1.561,-1.075) (-1.611,-1.100) (-1.662,-1.125) (-1.714,-1.150) (-1.768,-1.175) (-1.825,-1.200) (-1.875,-1.222) (-1.925,-1.242) (-1.975,-1.262) (-2.025,-1.280) (-2.075,-1.298) (-2.125,-1.315) (-2.175,-1.330) (-2.225,-1.345) (-2.275,-1.359) (-2.325,-1.371) (-2.375,-1.383) (-2.425,-1.393) (-2.475,-1.402) (-2.525,-1.410) (-2.575,-1.417) (-2.625,-1.422) (-2.675,-1.427) (-2.725,-1.430) (-2.775,-1.431) (-2.825,-1.432) (-2.875,-1.430) (-2.925,-1.428) (-2.950,-1.426)}
\EQ{(-3.250,-1.528) (-3.300,-1.521) (-3.350,-1.513) (-3.400,-1.503) (-3.450,-1.492) (-3.500,-1.479) (-3.550,-1.464) (-3.600,-1.448) (-3.650,-1.430) (-3.700,-1.409) (-3.750,-1.387) (-3.800,-1.363) (-3.850,-1.336) (-3.900,-1.307) (-3.950,-1.275) (-4.000,-1.240) (-4.050,-1.203) (-4.100,-1.161) (-4.141,-1.125) (-4.192,-1.075) (-4.239,-1.025) (-4.275,-0.983) (-4.321,-0.925) (-4.357,-0.875) (-4.390,-0.825) (-4.421,-0.775) (-4.448,-0.725) (-4.474,-0.675) (-4.497,-0.625) (-4.518,-0.575) (-4.537,-0.525) (-4.554,-0.475) (-4.569,-0.425) (-4.582,-0.375) (-4.594,-0.325) (-4.604,-0.275) (-4.612,-0.225) (-4.618,-0.175) (-4.623,-0.125) (-4.626,-0.075) (-4.628,-0.025) (-4.627,0.050) (-4.625,0.100) (-4.621,0.150) (-4.615,0.200) (-4.608,0.250) (-4.599,0.300) (-4.588,0.350) (-4.576,0.400) (-4.562,0.450) (-4.546,0.500) (-4.528,0.550) (-4.508,0.600) (-4.486,0.650) (-4.461,0.700) (-4.435,0.750) (-4.406,0.800) (-4.375,0.849) (-4.340,0.900) (-4.302,0.950) (-4.261,1.000) (-4.225,1.040) (-4.175,1.092) (-4.125,1.139) (-4.084,1.175) (-4.025,1.222) (-3.975,1.258) (-3.925,1.291) (-3.875,1.322) (-3.825,1.350) (-3.776,1.375) (-3.725,1.399) (-3.675,1.420) (-3.625,1.439) (-3.575,1.456) (-3.525,1.472) (-3.475,1.486) (-3.425,1.498) (-3.375,1.508) (-3.325,1.517) (-3.275,1.524) (-3.225,1.530) (-3.175,1.535) (-3.125,1.538) (-3.075,1.540) (-3.025,1.540) (-2.975,1.540) (-2.925,1.538) (-2.875,1.534) (-2.825,1.530) (-2.775,1.524) (-2.725,1.517) (-2.675,1.509) (-2.625,1.500) (-2.575,1.490) (-2.525,1.479) (-2.475,1.467) (-2.425,1.453) (-2.375,1.439) (-2.325,1.424) (-2.275,1.408) (-2.225,1.391) (-2.175,1.373) (-2.125,1.354) (-2.075,1.334) (-2.025,1.314) (-1.975,1.292) (-1.925,1.270) (-1.875,1.247) (-1.828,1.225) (-1.776,1.200) (-1.726,1.175) (-1.677,1.150) (-1.629,1.125) (-1.582,1.100) (-1.536,1.075) (-1.491,1.050) (-1.446,1.025) (-1.402,1.000) (-1.359,0.975) (-1.300,0.941) (-1.250,0.912) (-1.200,0.882) (-1.150,0.852) (-1.105,0.825) (-1.050,0.792) (-1.000,0.762) (-0.950,0.731) (-0.900,0.701) (-0.850,0.671) (-0.800,0.641) (-0.750,0.611) (-0.700,0.582) (-0.650,0.553) (-0.602,0.525) (-0.557,0.500) (-0.511,0.475) (-0.463,0.450) (-0.412,0.425) (-0.359,0.400) (-0.300,0.375) (-0.250,0.357) (-0.200,0.340) (-0.150,0.327) (-0.100,0.317) (-0.050,0.311) (0.000,0.309) (0.050,0.311) (0.100,0.317) (0.150,0.327) (0.200,0.340) (0.250,0.357) (0.299,0.375) (0.350,0.397) (0.400,0.420) (0.450,0.444) (0.500,0.470) (0.550,0.497) (0.600,0.524) (0.646,0.550) (0.689,0.575) (0.750,0.611) (0.800,0.641) (0.850,0.671) (0.899,0.700) (0.950,0.731) (1.000,0.762) (1.050,0.792) (1.100,0.822) (1.147,0.850) (1.200,0.882) (1.250,0.912) (1.300,0.941) (1.350,0.970) (1.400,0.999) (1.446,1.025) (1.491,1.050) (1.536,1.075) (1.582,1.100) (1.629,1.125) (1.675,1.149) (1.725,1.175) (1.775,1.200) (1.825,1.224) (1.875,1.247) (1.925,1.270) (1.975,1.292) (2.025,1.314) (2.075,1.334) (2.125,1.354) (2.175,1.373) (2.225,1.391) (2.275,1.408) (2.325,1.424) (2.375,1.439) (2.425,1.453) (2.475,1.467) (2.525,1.479) (2.575,1.490) (2.625,1.500) (2.675,1.509) (2.725,1.517) (2.775,1.524) (2.825,1.530) (2.875,1.534) (2.925,1.538) (2.975,1.540) (3.025,1.540) (3.075,1.540) (3.125,1.538) (3.175,1.535) (3.225,1.530) (3.275,1.524) (3.325,1.517) (3.375,1.508) (3.425,1.498) (3.475,1.486) (3.525,1.472) (3.575,1.456) (3.625,1.439) (3.675,1.420) (3.722,1.400) (3.775,1.375) (3.824,1.350) (3.869,1.325) (3.925,1.291) (3.975,1.258) (4.021,1.225) (4.075,1.182) (4.113,1.150) (4.167,1.100) (4.216,1.050) (4.250,1.013) (4.300,0.952) (4.340,0.900) (4.374,0.850) (4.406,0.800) (4.435,0.750) (4.461,0.700) (4.486,0.650) (4.508,0.600) (4.528,0.550) (4.546,0.500) (4.562,0.450) (4.576,0.400) (4.588,0.350) (4.599,0.300) (4.608,0.250) (4.615,0.200) (4.621,0.150) (4.625,0.100) (4.627,0.050) (4.628,0.000) (4.627,-0.050) (4.625,-0.100) (4.621,-0.150) (4.615,-0.200) (4.608,-0.250) (4.599,-0.300) (4.588,-0.350) (4.576,-0.400) (4.562,-0.450) (4.546,-0.500) (4.528,-0.550) (4.508,-0.600) (4.486,-0.650) (4.461,-0.700) (4.435,-0.750) (4.406,-0.800) (4.375,-0.849) (4.340,-0.900) (4.302,-0.950) (4.261,-1.000) (4.225,-1.040) (4.175,-1.092) (4.125,-1.139) (4.084,-1.175) (4.025,-1.222) (3.975,-1.258) (3.925,-1.291) (3.875,-1.322) (3.825,-1.350) (3.776,-1.375) (3.725,-1.399) (3.675,-1.420) (3.625,-1.439) (3.575,-1.456) (3.525,-1.472) (3.475,-1.486) (3.425,-1.498) (3.375,-1.508) (3.325,-1.517) (3.275,-1.524) (3.225,-1.530) (3.175,-1.535) (3.125,-1.538) (3.075,-1.540) (3.025,-1.540) (2.975,-1.540) (2.925,-1.538) (2.875,-1.534) (2.825,-1.530) (2.775,-1.524) (2.725,-1.517) (2.675,-1.509) (2.625,-1.500) (2.575,-1.490) (2.525,-1.479) (2.475,-1.467) (2.425,-1.453) (2.375,-1.439) (2.325,-1.424) (2.275,-1.408) (2.225,-1.391) (2.175,-1.373) (2.125,-1.354) (2.075,-1.334) (2.025,-1.314) (1.975,-1.292) (1.925,-1.270) (1.875,-1.247) (1.828,-1.225) (1.776,-1.200) (1.726,-1.175) (1.677,-1.150) (1.629,-1.125) (1.582,-1.100) (1.536,-1.075) (1.491,-1.050) (1.446,-1.025) (1.402,-1.000) (1.359,-0.975) (1.300,-0.941) (1.250,-0.912) (1.200,-0.882) (1.150,-0.852) (1.105,-0.825) (1.050,-0.792) (1.000,-0.762) (0.950,-0.731) (0.900,-0.701) (0.850,-0.671) (0.800,-0.641) (0.750,-0.611) (0.700,-0.582) (0.650,-0.553) (0.602,-0.525) (0.557,-0.500) (0.511,-0.475) (0.463,-0.450) (0.412,-0.425) (0.359,-0.400) (0.300,-0.375) (0.250,-0.357) (0.200,-0.340) (0.150,-0.327) (0.100,-0.317) (0.050,-0.311) (0.000,-0.309) (-0.050,-0.311) (-0.100,-0.317) (-0.150,-0.327) (-0.200,-0.340) (-0.250,-0.357) (-0.299,-0.375) (-0.350,-0.397) (-0.400,-0.420) (-0.450,-0.444) (-0.500,-0.470) (-0.550,-0.497) (-0.600,-0.524) (-0.646,-0.550) (-0.689,-0.575) (-0.750,-0.611) (-0.800,-0.641) (-0.850,-0.671) (-0.899,-0.700) (-0.950,-0.731) (-1.000,-0.762) (-1.050,-0.792) (-1.100,-0.822) (-1.147,-0.850) (-1.200,-0.882) (-1.250,-0.912) (-1.300,-0.941) (-1.350,-0.970) (-1.400,-0.999) (-1.446,-1.025) (-1.491,-1.050) (-1.536,-1.075) (-1.582,-1.100) (-1.629,-1.125) (-1.675,-1.149) (-1.725,-1.175) (-1.775,-1.200) (-1.825,-1.224) (-1.875,-1.247) (-1.925,-1.270) (-1.975,-1.292) (-2.025,-1.314) (-2.075,-1.334) (-2.125,-1.354) (-2.175,-1.373) (-2.225,-1.391) (-2.275,-1.408) (-2.325,-1.424) (-2.375,-1.439) (-2.425,-1.453) (-2.475,-1.467) (-2.525,-1.479) (-2.575,-1.490) (-2.625,-1.500) (-2.675,-1.509) (-2.725,-1.517) (-2.775,-1.524) (-2.825,-1.530) (-2.875,-1.534) (-2.925,-1.538) (-2.975,-1.540) (-3.025,-1.540) (-3.075,-1.540) (-3.125,-1.538) (-3.175,-1.535) (-3.225,-1.530) (-3.250,-1.528)}
\EQ{(-3.600,-1.677) (-3.650,-1.670) (-3.700,-1.662) (-3.750,-1.652) (-3.800,-1.640) (-3.850,-1.628) (-3.900,-1.613) (-3.950,-1.597) (-4.000,-1.580) (-4.050,-1.560) (-4.100,-1.539) (-4.150,-1.516) (-4.200,-1.491) (-4.250,-1.464) (-4.300,-1.434) (-4.350,-1.402) (-4.400,-1.368) (-4.450,-1.331) (-4.500,-1.290) (-4.546,-1.250) (-4.599,-1.200) (-4.648,-1.150) (-4.692,-1.100) (-4.725,-1.061) (-4.772,-1.000) (-4.807,-0.950) (-4.839,-0.900) (-4.869,-0.850) (-4.897,-0.800) (-4.923,-0.750) (-4.946,-0.700) (-4.968,-0.650) (-4.988,-0.600) (-5.006,-0.550) (-5.022,-0.500) (-5.036,-0.450) (-5.049,-0.400) (-5.060,-0.350) (-5.070,-0.300) (-5.078,-0.250) (-5.085,-0.200) (-5.090,-0.150) (-5.094,-0.100) (-5.096,-0.050) (-5.097,0.000) (-5.096,0.050) (-5.094,0.100) (-5.090,0.150) (-5.085,0.200) (-5.078,0.250) (-5.070,0.300) (-5.060,0.350) (-5.049,0.400) (-5.036,0.450) (-5.022,0.500) (-5.006,0.550) (-4.988,0.600) (-4.968,0.650) (-4.946,0.700) (-4.925,0.746) (-4.900,0.795) (-4.875,0.840) (-4.839,0.900) (-4.807,0.950) (-4.775,0.995) (-4.734,1.050) (-4.700,1.091) (-4.650,1.147) (-4.600,1.199) (-4.550,1.247) (-4.500,1.290) (-4.457,1.325) (-4.400,1.368) (-4.354,1.400) (-4.300,1.434) (-4.250,1.464) (-4.200,1.491) (-4.150,1.516) (-4.100,1.539) (-4.050,1.560) (-4.000,1.580) (-3.950,1.597) (-3.900,1.613) (-3.850,1.628) (-3.800,1.640) (-3.750,1.652) (-3.700,1.662) (-3.650,1.670) (-3.600,1.677) (-3.550,1.683) (-3.500,1.687) (-3.450,1.691) (-3.400,1.693) (-3.350,1.693) (-3.300,1.693) (-3.250,1.691) (-3.200,1.689) (-3.150,1.685) (-3.100,1.680) (-3.050,1.674) (-3.000,1.667) (-2.950,1.659) (-2.900,1.650) (-2.850,1.640) (-2.800,1.629) (-2.750,1.617) (-2.700,1.604) (-2.650,1.591) (-2.600,1.576) (-2.550,1.560) (-2.500,1.544) (-2.450,1.527) (-2.400,1.509) (-2.350,1.490) (-2.300,1.470) (-2.251,1.450) (-2.200,1.428) (-2.150,1.406) (-2.100,1.384) (-2.050,1.361) (-2.000,1.337) (-1.950,1.312) (-1.900,1.287) (-1.850,1.261) (-1.800,1.235) (-1.750,1.208) (-1.700,1.181) (-1.650,1.153) (-1.600,1.125) (-1.557,1.100) (-1.500,1.068) (-1.450,1.038) (-1.400,1.009) (-1.350,0.979) (-1.303,0.950) (-1.250,0.918) (-1.200,0.888) (-1.150,0.857) (-1.100,0.826) (-1.050,0.795) (-1.000,0.764) (-0.950,0.734) (-0.900,0.703) (-0.855,0.675) (-0.800,0.642) (-0.750,0.612) (-0.700,0.583) (-0.650,0.553) (-0.601,0.525) (-0.556,0.500) (-0.510,0.475) (-0.462,0.450) (-0.412,0.425) (-0.358,0.400) (-0.300,0.376) (-0.250,0.357) (-0.200,0.340) (-0.150,0.327) (-0.100,0.317) (-0.050,0.311) (0.000,0.309) (0.050,0.311) (0.100,0.317) (0.150,0.327) (0.200,0.340) (0.250,0.357) (0.299,0.375) (0.350,0.397) (0.400,0.420) (0.450,0.444) (0.500,0.470) (0.550,0.497) (0.600,0.525) (0.645,0.550) (0.700,0.583) (0.750,0.612) (0.800,0.642) (0.850,0.672) (0.896,0.700) (0.950,0.734) (1.000,0.764) (1.050,0.795) (1.099,0.825) (1.150,0.857) (1.200,0.888) (1.250,0.918) (1.300,0.948) (1.344,0.975) (1.400,1.009) (1.450,1.038) (1.500,1.068) (1.550,1.096) (1.600,1.125) (1.645,1.150) (1.690,1.175) (1.735,1.200) (1.782,1.225) (1.829,1.250) (1.875,1.274) (1.925,1.300) (1.975,1.324) (2.025,1.349) (2.075,1.372) (2.125,1.395) (2.175,1.417) (2.225,1.439) (2.275,1.460) (2.325,1.480) (2.375,1.499) (2.425,1.518) (2.475,1.535) (2.525,1.552) (2.575,1.568) (2.625,1.583) (2.675,1.597) (2.725,1.611) (2.775,1.623) (2.825,1.635) (2.875,1.645) (2.925,1.655) (2.975,1.663) (3.025,1.671) (3.075,1.677) (3.125,1.683) (3.175,1.687) (3.225,1.690) (3.275,1.692) (3.325,1.693) (3.375,1.693) (3.425,1.692) (3.475,1.689) (3.525,1.685) (3.575,1.680) (3.625,1.674) (3.675,1.666) (3.725,1.657) (3.775,1.646) (3.825,1.634) (3.875,1.621) (3.925,1.605) (3.975,1.589) (4.025,1.570) (4.075,1.550) (4.125,1.528) (4.175,1.504) (4.225,1.478) (4.274,1.450) (4.325,1.419) (4.375,1.385) (4.425,1.350) (4.475,1.311) (4.518,1.275) (4.573,1.225) (4.624,1.175) (4.670,1.125) (4.713,1.075) (4.750,1.029) (4.790,0.975) (4.823,0.925) (4.855,0.875) (4.884,0.825) (4.910,0.775) (4.935,0.725) (4.957,0.675) (4.978,0.625) (4.997,0.575) (5.014,0.525) (5.029,0.475) (5.043,0.425) (5.055,0.375) (5.065,0.325) (5.074,0.275) (5.082,0.225) (5.088,0.175) (5.092,0.125) (5.095,0.075) (5.096,0.025) (5.096,-0.050) (5.094,-0.100) (5.090,-0.150) (5.085,-0.200) (5.078,-0.250) (5.070,-0.300) (5.060,-0.350) (5.049,-0.400) (5.036,-0.450) (5.022,-0.500) (5.006,-0.550) (4.988,-0.600) (4.968,-0.650) (4.946,-0.700) (4.925,-0.746) (4.900,-0.795) (4.875,-0.840) (4.839,-0.900) (4.807,-0.950) (4.775,-0.995) (4.734,-1.050) (4.700,-1.091) (4.650,-1.147) (4.600,-1.199) (4.550,-1.247) (4.500,-1.290) (4.457,-1.325) (4.400,-1.368) (4.354,-1.400) (4.300,-1.434) (4.250,-1.464) (4.200,-1.491) (4.150,-1.516) (4.100,-1.539) (4.050,-1.560) (4.000,-1.580) (3.950,-1.597) (3.900,-1.613) (3.850,-1.628) (3.800,-1.640) (3.750,-1.652) (3.700,-1.662) (3.650,-1.670) (3.600,-1.677) (3.550,-1.683) (3.500,-1.687) (3.450,-1.691) (3.400,-1.693) (3.350,-1.693) (3.300,-1.693) (3.250,-1.691) (3.200,-1.689) (3.150,-1.685) (3.100,-1.680) (3.050,-1.674) (3.000,-1.667) (2.950,-1.659) (2.900,-1.650) (2.850,-1.640) (2.800,-1.629) (2.750,-1.617) (2.700,-1.604) (2.650,-1.591) (2.600,-1.576) (2.550,-1.560) (2.500,-1.544) (2.450,-1.527) (2.400,-1.509) (2.350,-1.490) (2.300,-1.470) (2.251,-1.450) (2.200,-1.428) (2.150,-1.406) (2.100,-1.384) (2.050,-1.361) (2.000,-1.337) (1.950,-1.312) (1.900,-1.287) (1.850,-1.261) (1.800,-1.235) (1.750,-1.208) (1.700,-1.181) (1.650,-1.153) (1.600,-1.125) (1.557,-1.100) (1.500,-1.068) (1.450,-1.038) (1.400,-1.009) (1.350,-0.979) (1.303,-0.950) (1.250,-0.918) (1.200,-0.888) (1.150,-0.857) (1.100,-0.826) (1.050,-0.795) (1.000,-0.764) (0.950,-0.734) (0.900,-0.703) (0.855,-0.675) (0.800,-0.642) (0.750,-0.612) (0.700,-0.583) (0.650,-0.553) (0.601,-0.525) (0.556,-0.500) (0.510,-0.475) (0.462,-0.450) (0.412,-0.425) (0.358,-0.400) (0.300,-0.376) (0.250,-0.357) (0.200,-0.340) (0.150,-0.327) (0.100,-0.317) (0.050,-0.311) (0.000,-0.309) (-0.050,-0.311) (-0.100,-0.317) (-0.150,-0.327) (-0.200,-0.340) (-0.250,-0.357) (-0.299,-0.375) (-0.350,-0.397) (-0.400,-0.420) (-0.450,-0.444) (-0.500,-0.470) (-0.550,-0.497) (-0.600,-0.525) (-0.645,-0.550) (-0.700,-0.583) (-0.750,-0.612) (-0.800,-0.642) (-0.850,-0.672) (-0.896,-0.700) (-0.950,-0.734) (-1.000,-0.764) (-1.050,-0.795) (-1.099,-0.825) (-1.150,-0.857) (-1.200,-0.888) (-1.250,-0.918) (-1.300,-0.948) (-1.344,-0.975) (-1.400,-1.009) (-1.450,-1.038) (-1.500,-1.068) (-1.550,-1.096) (-1.600,-1.125) (-1.645,-1.150) (-1.690,-1.175) (-1.735,-1.200) (-1.782,-1.225) (-1.829,-1.250) (-1.875,-1.274) (-1.925,-1.300) (-1.975,-1.324) (-2.025,-1.349) (-2.075,-1.372) (-2.125,-1.395) (-2.175,-1.417) (-2.225,-1.439) (-2.275,-1.460) (-2.325,-1.480) (-2.375,-1.499) (-2.425,-1.518) (-2.475,-1.535) (-2.525,-1.552) (-2.575,-1.568) (-2.625,-1.583) (-2.675,-1.597) (-2.725,-1.611) (-2.775,-1.623) (-2.825,-1.635) (-2.875,-1.645) (-2.925,-1.655) (-2.975,-1.663) (-3.025,-1.671) (-3.075,-1.677) (-3.125,-1.683) (-3.175,-1.687) (-3.225,-1.690) (-3.275,-1.692) (-3.325,-1.693) (-3.375,-1.693) (-3.425,-1.692) (-3.475,-1.689) (-3.525,-1.685) (-3.575,-1.680) (-3.600,-1.677)}
\EQ{(-4.050,-1.926) (-4.100,-1.920) (-4.150,-1.914) (-4.200,-1.906) (-4.250,-1.898) (-4.300,-1.888) (-4.350,-1.876) (-4.400,-1.864) (-4.450,-1.850) (-4.500,-1.835) (-4.550,-1.818) (-4.599,-1.800) (-4.650,-1.780) (-4.700,-1.759) (-4.750,-1.736) (-4.800,-1.711) (-4.850,-1.684) (-4.900,-1.656) (-4.950,-1.625) (-5.000,-1.592) (-5.050,-1.557) (-5.093,-1.525) (-5.150,-1.480) (-5.200,-1.436) (-5.240,-1.400) (-5.291,-1.350) (-5.338,-1.300) (-5.375,-1.258) (-5.423,-1.200) (-5.461,-1.150) (-5.497,-1.100) (-5.525,-1.058) (-5.561,-1.000) (-5.590,-0.950) (-5.617,-0.900) (-5.642,-0.850) (-5.666,-0.800) (-5.688,-0.750) (-5.708,-0.700) (-5.725,-0.653) (-5.743,-0.600) (-5.758,-0.550) (-5.772,-0.500) (-5.785,-0.450) (-5.796,-0.400) (-5.806,-0.350) (-5.814,-0.300) (-5.821,-0.250) (-5.827,-0.200) (-5.832,-0.150) (-5.835,-0.100) (-5.837,-0.050) (-5.837,0.000) (-5.837,0.050) (-5.835,0.100) (-5.832,0.150) (-5.827,0.200) (-5.821,0.250) (-5.814,0.300) (-5.806,0.350) (-5.796,0.400) (-5.785,0.450) (-5.772,0.500) (-5.758,0.550) (-5.743,0.600) (-5.726,0.650) (-5.708,0.700) (-5.688,0.750) (-5.666,0.800) (-5.642,0.850) (-5.617,0.900) (-5.590,0.950) (-5.561,1.000) (-5.530,1.050) (-5.500,1.095) (-5.461,1.150) (-5.425,1.197) (-5.382,1.250) (-5.338,1.300) (-5.300,1.340) (-5.250,1.390) (-5.213,1.425) (-5.175,1.458) (-5.125,1.500) (-5.075,1.539) (-5.025,1.575) (-4.975,1.609) (-4.925,1.641) (-4.875,1.670) (-4.825,1.698) (-4.775,1.724) (-4.725,1.747) (-4.675,1.770) (-4.625,1.790) (-4.575,1.809) (-4.525,1.827) (-4.475,1.843) (-4.425,1.857) (-4.375,1.870) (-4.325,1.882) (-4.275,1.893) (-4.225,1.902) (-4.175,1.910) (-4.125,1.917) (-4.075,1.923) (-4.025,1.928) (-3.975,1.931) (-3.925,1.934) (-3.875,1.935) (-3.825,1.936) (-3.775,1.935) (-3.725,1.933) (-3.675,1.931) (-3.625,1.927) (-3.575,1.922) (-3.525,1.917) (-3.475,1.910) (-3.425,1.903) (-3.375,1.895) (-3.325,1.886) (-3.275,1.876) (-3.225,1.865) (-3.175,1.853) (-3.125,1.841) (-3.075,1.827) (-3.025,1.813) (-2.975,1.798) (-2.925,1.783) (-2.875,1.766) (-2.825,1.749) (-2.775,1.731) (-2.725,1.712) (-2.675,1.693) (-2.625,1.673) (-2.575,1.652) (-2.525,1.631) (-2.475,1.609) (-2.425,1.586) (-2.375,1.563) (-2.325,1.539) (-2.275,1.515) (-2.225,1.490) (-2.175,1.465) (-2.125,1.439) (-2.075,1.412) (-2.025,1.386) (-1.975,1.358) (-1.925,1.330) (-1.875,1.302) (-1.828,1.275) (-1.775,1.245) (-1.725,1.215) (-1.675,1.186) (-1.625,1.156) (-1.575,1.126) (-1.525,1.095) (-1.475,1.065) (-1.425,1.034) (-1.375,1.003) (-1.331,0.975) (-1.275,0.940) (-1.225,0.909) (-1.175,0.877) (-1.132,0.850) (-1.075,0.814) (-1.025,0.783) (-0.975,0.751) (-0.925,0.720) (-0.875,0.689) (-0.825,0.659) (-0.775,0.628) (-0.729,0.600) (-0.675,0.569) (-0.625,0.540) (-0.575,0.511) (-0.525,0.484) (-0.475,0.457) (-0.425,0.432) (-0.375,0.408) (-0.325,0.386) (-0.275,0.366) (-0.225,0.348) (-0.175,0.333) (-0.125,0.322) (-0.075,0.313) (-0.025,0.309) (0.050,0.311) (0.100,0.317) (0.150,0.327) (0.200,0.340) (0.250,0.357) (0.299,0.375) (0.350,0.397) (0.400,0.420) (0.450,0.445) (0.500,0.470) (0.550,0.497) (0.600,0.525) (0.644,0.550) (0.700,0.583) (0.750,0.613) (0.800,0.643) (0.850,0.674) (0.893,0.700) (0.950,0.736) (1.000,0.767) (1.050,0.798) (1.093,0.825) (1.150,0.862) (1.200,0.893) (1.250,0.924) (1.300,0.956) (1.350,0.987) (1.400,1.018) (1.450,1.049) (1.500,1.080) (1.550,1.110) (1.600,1.141) (1.650,1.171) (1.699,1.200) (1.750,1.230) (1.800,1.259) (1.850,1.288) (1.900,1.316) (1.950,1.344) (2.000,1.372) (2.050,1.399) (2.099,1.425) (2.147,1.450) (2.195,1.475) (2.245,1.500) (2.296,1.525) (2.347,1.550) (2.400,1.575) (2.450,1.598) (2.500,1.620) (2.550,1.642) (2.600,1.663) (2.650,1.683) (2.700,1.703) (2.750,1.722) (2.800,1.740) (2.850,1.758) (2.900,1.774) (2.950,1.790) (3.000,1.806) (3.050,1.820) (3.100,1.834) (3.150,1.847) (3.200,1.859) (3.250,1.870) (3.300,1.881) (3.350,1.890) (3.400,1.899) (3.450,1.907) (3.500,1.914) (3.550,1.920) (3.600,1.925) (3.650,1.929) (3.700,1.932) (3.750,1.934) (3.800,1.935) (3.850,1.936) (3.900,1.935) (3.950,1.933) (4.000,1.930) (4.050,1.926) (4.100,1.920) (4.150,1.914) (4.200,1.906) (4.250,1.898) (4.300,1.888) (4.350,1.876) (4.400,1.864) (4.450,1.850) (4.500,1.835) (4.550,1.818) (4.599,1.800) (4.650,1.780) (4.700,1.759) (4.750,1.736) (4.800,1.711) (4.850,1.684) (4.900,1.656) (4.950,1.625) (5.000,1.592) (5.050,1.557) (5.093,1.525) (5.150,1.480) (5.200,1.436) (5.240,1.400) (5.291,1.350) (5.338,1.300) (5.375,1.258) (5.423,1.200) (5.461,1.150) (5.497,1.100) (5.525,1.058) (5.561,1.000) (5.590,0.950) (5.617,0.900) (5.642,0.850) (5.666,0.800) (5.688,0.750) (5.708,0.700) (5.725,0.653) (5.743,0.600) (5.758,0.550) (5.772,0.500) (5.785,0.450) (5.796,0.400) (5.806,0.350) (5.814,0.300) (5.821,0.250) (5.827,0.200) (5.832,0.150) (5.835,0.100) (5.837,0.050) (5.837,0.000) (5.837,-0.050) (5.835,-0.100) (5.832,-0.150) (5.827,-0.200) (5.821,-0.250) (5.814,-0.300) (5.806,-0.350) (5.796,-0.400) (5.785,-0.450) (5.772,-0.500) (5.758,-0.550) (5.743,-0.600) (5.726,-0.650) (5.708,-0.700) (5.688,-0.750) (5.666,-0.800) (5.642,-0.850) (5.617,-0.900) (5.590,-0.950) (5.561,-1.000) (5.530,-1.050) (5.500,-1.095) (5.461,-1.150) (5.425,-1.197) (5.382,-1.250) (5.338,-1.300) (5.300,-1.340) (5.250,-1.390) (5.213,-1.425) (5.175,-1.458) (5.125,-1.500) (5.075,-1.539) (5.025,-1.575) (4.975,-1.609) (4.925,-1.641) (4.875,-1.670) (4.825,-1.698) (4.775,-1.724) (4.725,-1.747) (4.675,-1.770) (4.625,-1.790) (4.575,-1.809) (4.525,-1.827) (4.475,-1.843) (4.425,-1.857) (4.375,-1.870) (4.325,-1.882) (4.275,-1.893) (4.225,-1.902) (4.175,-1.910) (4.125,-1.917) (4.075,-1.923) (4.025,-1.928) (3.975,-1.931) (3.925,-1.934) (3.875,-1.935) (3.825,-1.936) (3.775,-1.935) (3.725,-1.933) (3.675,-1.931) (3.625,-1.927) (3.575,-1.922) (3.525,-1.917) (3.475,-1.910) (3.425,-1.903) (3.375,-1.895) (3.325,-1.886) (3.275,-1.876) (3.225,-1.865) (3.175,-1.853) (3.125,-1.841) (3.075,-1.827) (3.025,-1.813) (2.975,-1.798) (2.925,-1.783) (2.875,-1.766) (2.825,-1.749) (2.775,-1.731) (2.725,-1.712) (2.675,-1.693) (2.625,-1.673) (2.575,-1.652) (2.525,-1.631) (2.475,-1.609) (2.425,-1.586) (2.375,-1.563) (2.325,-1.539) (2.275,-1.515) (2.225,-1.490) (2.175,-1.465) (2.125,-1.439) (2.075,-1.412) (2.025,-1.386) (1.975,-1.358) (1.925,-1.330) (1.875,-1.302) (1.828,-1.275) (1.775,-1.245) (1.725,-1.215) (1.675,-1.186) (1.625,-1.156) (1.575,-1.126) (1.525,-1.095) (1.475,-1.065) (1.425,-1.034) (1.375,-1.003) (1.331,-0.975) (1.275,-0.940) (1.225,-0.909) (1.175,-0.877) (1.132,-0.850) (1.075,-0.814) (1.025,-0.783) (0.975,-0.751) (0.925,-0.720) (0.875,-0.689) (0.825,-0.659) (0.775,-0.628) (0.729,-0.600) (0.675,-0.569) (0.625,-0.540) (0.575,-0.511) (0.525,-0.484) (0.475,-0.457) (0.425,-0.432) (0.375,-0.408) (0.325,-0.386) (0.275,-0.366) (0.225,-0.348) (0.175,-0.333) (0.125,-0.322) (0.075,-0.313) (0.025,-0.309) (-0.050,-0.311) (-0.100,-0.317) (-0.150,-0.327) (-0.200,-0.340) (-0.250,-0.357) (-0.299,-0.375) (-0.350,-0.397) (-0.400,-0.420) (-0.450,-0.445) (-0.500,-0.470) (-0.550,-0.497) (-0.600,-0.525) (-0.644,-0.550) (-0.700,-0.583) (-0.750,-0.613) (-0.800,-0.643) (-0.850,-0.674) (-0.893,-0.700) (-0.950,-0.736) (-1.000,-0.767) (-1.050,-0.798) (-1.093,-0.825) (-1.150,-0.862) (-1.200,-0.893) (-1.250,-0.924) (-1.300,-0.956) (-1.350,-0.987) (-1.400,-1.018) (-1.450,-1.049) (-1.500,-1.080) (-1.550,-1.110) (-1.600,-1.141) (-1.650,-1.171) (-1.699,-1.200) (-1.750,-1.230) (-1.800,-1.259) (-1.850,-1.288) (-1.900,-1.316) (-1.950,-1.344) (-2.000,-1.372) (-2.050,-1.399) (-2.099,-1.425) (-2.147,-1.450) (-2.195,-1.475) (-2.245,-1.500) (-2.296,-1.525) (-2.347,-1.550) (-2.400,-1.575) (-2.450,-1.598) (-2.500,-1.620) (-2.550,-1.642) (-2.600,-1.663) (-2.650,-1.683) (-2.700,-1.703) (-2.750,-1.722) (-2.800,-1.740) (-2.850,-1.758) (-2.900,-1.774) (-2.950,-1.790) (-3.000,-1.806) (-3.050,-1.820) (-3.100,-1.834) (-3.150,-1.847) (-3.200,-1.859) (-3.250,-1.870) (-3.300,-1.881) (-3.350,-1.890) (-3.400,-1.899) (-3.450,-1.907) (-3.500,-1.914) (-3.550,-1.920) (-3.600,-1.925) (-3.650,-1.929) (-3.700,-1.932) (-3.750,-1.934) (-3.800,-1.935) (-3.850,-1.936) (-3.900,-1.935) (-3.950,-1.933) (-4.000,-1.930) (-4.050,-1.926)}
\EQ{(-5.075,-2.425) (-5.125,-2.421) (-5.175,-2.416) (-5.225,-2.409) (-5.275,-2.402) (-5.325,-2.394) (-5.375,-2.385) (-5.425,-2.375) (-5.475,-2.364) (-5.525,-2.351) (-5.575,-2.338) (-5.625,-2.324) (-5.675,-2.309) (-5.725,-2.292) (-5.774,-2.275) (-5.825,-2.256) (-5.875,-2.235) (-5.925,-2.214) (-5.975,-2.191) (-6.025,-2.167) (-6.075,-2.142) (-6.125,-2.115) (-6.175,-2.086) (-6.225,-2.056) (-6.274,-2.025) (-6.325,-1.990) (-6.375,-1.955) (-6.425,-1.917) (-6.475,-1.878) (-6.525,-1.836) (-6.566,-1.800) (-6.619,-1.750) (-6.670,-1.700) (-6.718,-1.650) (-6.763,-1.600) (-6.800,-1.556) (-6.845,-1.500) (-6.884,-1.450) (-6.920,-1.400) (-6.950,-1.356) (-6.986,-1.300) (-7.017,-1.250) (-7.046,-1.200) (-7.073,-1.150) (-7.099,-1.100) (-7.124,-1.050) (-7.147,-1.000) (-7.168,-0.950) (-7.189,-0.900) (-7.208,-0.850) (-7.225,-0.802) (-7.242,-0.750) (-7.258,-0.700) (-7.272,-0.650) (-7.285,-0.600) (-7.297,-0.550) (-7.308,-0.500) (-7.318,-0.450) (-7.327,-0.400) (-7.335,-0.350) (-7.341,-0.300) (-7.347,-0.250) (-7.351,-0.200) (-7.355,-0.150) (-7.358,-0.100) (-7.359,-0.050) (-7.360,0.000) (-7.359,0.050) (-7.358,0.100) (-7.355,0.150) (-7.351,0.200) (-7.347,0.250) (-7.341,0.300) (-7.335,0.350) (-7.327,0.400) (-7.318,0.450) (-7.308,0.500) (-7.297,0.550) (-7.285,0.600) (-7.272,0.650) (-7.258,0.700) (-7.242,0.750) (-7.226,0.800) (-7.208,0.850) (-7.189,0.900) (-7.168,0.950) (-7.147,1.000) (-7.125,1.047) (-7.100,1.098) (-7.075,1.147) (-7.050,1.193) (-7.025,1.236) (-6.986,1.300) (-6.954,1.350) (-6.925,1.392) (-6.884,1.450) (-6.850,1.494) (-6.805,1.550) (-6.763,1.600) (-6.725,1.642) (-6.675,1.695) (-6.625,1.745) (-6.575,1.792) (-6.538,1.825) (-6.478,1.875) (-6.425,1.917) (-6.382,1.950) (-6.325,1.990) (-6.275,2.024) (-6.225,2.056) (-6.175,2.086) (-6.125,2.115) (-6.075,2.142) (-6.025,2.167) (-5.975,2.191) (-5.925,2.214) (-5.875,2.235) (-5.825,2.256) (-5.775,2.274) (-5.725,2.292) (-5.675,2.309) (-5.625,2.324) (-5.575,2.338) (-5.525,2.351) (-5.475,2.364) (-5.425,2.375) (-5.375,2.385) (-5.325,2.394) (-5.275,2.402) (-5.225,2.409) (-5.175,2.416) (-5.125,2.421) (-5.075,2.425) (-5.025,2.429) (-4.975,2.432) (-4.925,2.434) (-4.875,2.435) (-4.825,2.435) (-4.775,2.434) (-4.725,2.433) (-4.675,2.430) (-4.625,2.427) (-4.575,2.424) (-4.525,2.419) (-4.475,2.414) (-4.425,2.408) (-4.375,2.401) (-4.325,2.394) (-4.275,2.385) (-4.225,2.376) (-4.175,2.367) (-4.125,2.357) (-4.075,2.346) (-4.025,2.334) (-3.975,2.322) (-3.925,2.309) (-3.875,2.295) (-3.825,2.281) (-3.775,2.266) (-3.725,2.250) (-3.675,2.234) (-3.625,2.218) (-3.575,2.200) (-3.525,2.182) (-3.475,2.164) (-3.425,2.145) (-3.375,2.125) (-3.325,2.105) (-3.275,2.085) (-3.225,2.063) (-3.175,2.042) (-3.125,2.019) (-3.075,1.997) (-3.028,1.975) (-2.976,1.950) (-2.925,1.926) (-2.875,1.901) (-2.825,1.876) (-2.775,1.850) (-2.726,1.825) (-2.679,1.800) (-2.632,1.775) (-2.586,1.750) (-2.540,1.725) (-2.495,1.700) (-2.450,1.675) (-2.406,1.650) (-2.362,1.625) (-2.300,1.589) (-2.250,1.560) (-2.200,1.530) (-2.150,1.500) (-2.100,1.470) (-2.050,1.440) (-2.000,1.409) (-1.950,1.378) (-1.905,1.350) (-1.850,1.316) (-1.800,1.284) (-1.750,1.253) (-1.707,1.225) (-1.650,1.189) (-1.600,1.157) (-1.551,1.125) (-1.500,1.093) (-1.450,1.060) (-1.400,1.028) (-1.357,1.000) (-1.300,0.963) (-1.250,0.931) (-1.203,0.900) (-1.150,0.866) (-1.100,0.834) (-1.050,0.802) (-1.000,0.770) (-0.950,0.738) (-0.900,0.707) (-0.850,0.675) (-0.800,0.645) (-0.750,0.614) (-0.700,0.584) (-0.650,0.555) (-0.600,0.526) (-0.555,0.500) (-0.509,0.475) (-0.462,0.450) (-0.411,0.425) (-0.358,0.400) (-0.300,0.376) (-0.250,0.357) (-0.200,0.340) (-0.150,0.327) (-0.100,0.317) (-0.050,0.311) (0.000,0.309) (0.050,0.311) (0.100,0.317) (0.150,0.327) (0.200,0.340) (0.250,0.357) (0.299,0.375) (0.350,0.397) (0.400,0.420) (0.450,0.445) (0.500,0.471) (0.550,0.498) (0.599,0.525) (0.643,0.550) (0.700,0.584) (0.750,0.614) (0.800,0.645) (0.850,0.675) (0.900,0.707) (0.950,0.738) (1.000,0.770) (1.047,0.800) (1.100,0.834) (1.150,0.866) (1.200,0.898) (1.250,0.931) (1.300,0.963) (1.350,0.996) (1.396,1.025) (1.450,1.060) (1.500,1.093) (1.550,1.125) (1.600,1.157) (1.650,1.189) (1.700,1.221) (1.746,1.250) (1.800,1.284) (1.850,1.316) (1.900,1.347) (1.945,1.375) (2.000,1.409) (2.050,1.440) (2.100,1.470) (2.150,1.500) (2.200,1.530) (2.250,1.560) (2.300,1.589) (2.350,1.618) (2.400,1.647) (2.450,1.675) (2.495,1.700) (2.540,1.725) (2.586,1.750) (2.632,1.775) (2.679,1.800) (2.725,1.825) (2.774,1.850) (2.823,1.875) (2.873,1.900) (2.924,1.925) (2.975,1.950) (3.025,1.973) (3.075,1.997) (3.125,2.019) (3.175,2.042) (3.225,2.063) (3.275,2.085) (3.325,2.105) (3.374,2.125) (3.425,2.145) (3.475,2.164) (3.525,2.182) (3.574,2.200) (3.625,2.218) (3.675,2.234) (3.724,2.250) (3.775,2.266) (3.825,2.281) (3.875,2.295) (3.925,2.309) (3.975,2.322) (4.025,2.334) (4.075,2.346) (4.125,2.357) (4.175,2.367) (4.225,2.376) (4.275,2.385) (4.325,2.394) (4.375,2.401) (4.425,2.408) (4.475,2.414) (4.525,2.419) (4.575,2.424) (4.625,2.427) (4.675,2.430) (4.725,2.433) (4.775,2.434) (4.825,2.435) (4.875,2.435) (4.925,2.434) (4.975,2.432) (5.025,2.429) (5.075,2.425) (5.125,2.421) (5.175,2.416) (5.225,2.409) (5.275,2.402) (5.325,2.394) (5.375,2.385) (5.425,2.375) (5.475,2.364) (5.525,2.351) (5.575,2.338) (5.625,2.324) (5.675,2.309) (5.725,2.292) (5.774,2.275) (5.825,2.256) (5.875,2.235) (5.925,2.214) (5.975,2.191) (6.025,2.167) (6.075,2.142) (6.125,2.115) (6.175,2.086) (6.225,2.056) (6.274,2.025) (6.325,1.990) (6.375,1.955) (6.425,1.917) (6.475,1.878) (6.525,1.836) (6.566,1.800) (6.619,1.750) (6.670,1.700) (6.718,1.650) (6.763,1.600) (6.800,1.556) (6.845,1.500) (6.884,1.450) (6.920,1.400) (6.950,1.356) (6.986,1.300) (7.017,1.250) (7.046,1.200) (7.073,1.150) (7.099,1.100) (7.124,1.050) (7.147,1.000) (7.168,0.950) (7.189,0.900) (7.208,0.850) (7.225,0.802) (7.242,0.750) (7.258,0.700) (7.272,0.650) (7.285,0.600) (7.297,0.550) (7.308,0.500) (7.318,0.450) (7.327,0.400) (7.335,0.350) (7.341,0.300) (7.347,0.250) (7.351,0.200) (7.355,0.150) (7.358,0.100) (7.359,0.050) (7.360,0.000) (7.359,-0.050) (7.358,-0.100) (7.355,-0.150) (7.351,-0.200) (7.347,-0.250) (7.341,-0.300) (7.335,-0.350) (7.327,-0.400) (7.318,-0.450) (7.308,-0.500) (7.297,-0.550) (7.285,-0.600) (7.272,-0.650) (7.258,-0.700) (7.242,-0.750) (7.226,-0.800) (7.208,-0.850) (7.189,-0.900) (7.168,-0.950) (7.147,-1.000) (7.125,-1.047) (7.100,-1.098) (7.075,-1.147) (7.050,-1.193) (7.025,-1.236) (6.986,-1.300) (6.954,-1.350) (6.925,-1.392) (6.884,-1.450) (6.850,-1.494) (6.805,-1.550) (6.763,-1.600) (6.725,-1.642) (6.675,-1.695) (6.625,-1.745) (6.575,-1.792) (6.538,-1.825) (6.478,-1.875) (6.425,-1.917) (6.382,-1.950) (6.325,-1.990) (6.275,-2.024) (6.225,-2.056) (6.175,-2.086) (6.125,-2.115) (6.075,-2.142) (6.025,-2.167) (5.975,-2.191) (5.925,-2.214) (5.875,-2.235) (5.825,-2.256) (5.775,-2.274) (5.725,-2.292) (5.675,-2.309) (5.625,-2.324) (5.575,-2.338) (5.525,-2.351) (5.475,-2.364) (5.425,-2.375) (5.375,-2.385) (5.325,-2.394) (5.275,-2.402) (5.225,-2.409) (5.175,-2.416) (5.125,-2.421) (5.075,-2.425) (5.025,-2.429) (4.975,-2.432) (4.925,-2.434) (4.875,-2.435) (4.825,-2.435) (4.775,-2.434) (4.725,-2.433) (4.675,-2.430) (4.625,-2.427) (4.575,-2.424) (4.525,-2.419) (4.475,-2.414) (4.425,-2.408) (4.375,-2.401) (4.325,-2.394) (4.275,-2.385) (4.225,-2.376) (4.175,-2.367) (4.125,-2.357) (4.075,-2.346) (4.025,-2.334) (3.975,-2.322) (3.925,-2.309) (3.875,-2.295) (3.825,-2.281) (3.775,-2.266) (3.725,-2.250) (3.675,-2.234) (3.625,-2.218) (3.575,-2.200) (3.525,-2.182) (3.475,-2.164) (3.425,-2.145) (3.375,-2.125) (3.325,-2.105) (3.275,-2.085) (3.225,-2.063) (3.175,-2.042) (3.125,-2.019) (3.075,-1.997) (3.028,-1.975) (2.976,-1.950) (2.925,-1.926) (2.875,-1.901) (2.825,-1.876) (2.775,-1.850) (2.726,-1.825) (2.679,-1.800) (2.632,-1.775) (2.586,-1.750) (2.540,-1.725) (2.495,-1.700) (2.450,-1.675) (2.406,-1.650) (2.362,-1.625) (2.300,-1.589) (2.250,-1.560) (2.200,-1.530) (2.150,-1.500) (2.100,-1.470) (2.050,-1.440) (2.000,-1.409) (1.950,-1.378) (1.905,-1.350) (1.850,-1.316) (1.800,-1.284) (1.750,-1.253) (1.707,-1.225) (1.650,-1.189) (1.600,-1.157) (1.551,-1.125) (1.500,-1.093) (1.450,-1.060) (1.400,-1.028) (1.357,-1.000) (1.300,-0.963) (1.250,-0.931) (1.203,-0.900) (1.150,-0.866) (1.100,-0.834) (1.050,-0.802) (1.000,-0.770) (0.950,-0.738) (0.900,-0.707) (0.850,-0.675) (0.800,-0.645) (0.750,-0.614) (0.700,-0.584) (0.650,-0.555) (0.600,-0.526) (0.555,-0.500) (0.509,-0.475) (0.462,-0.450) (0.411,-0.425) (0.358,-0.400) (0.300,-0.376) (0.250,-0.357) (0.200,-0.340) (0.150,-0.327) (0.100,-0.317) (0.050,-0.311) (0.000,-0.309) (-0.050,-0.311) (-0.100,-0.317) (-0.150,-0.327) (-0.200,-0.340) (-0.250,-0.357) (-0.299,-0.375) (-0.350,-0.397) (-0.400,-0.420) (-0.450,-0.445) (-0.500,-0.471) (-0.550,-0.498) (-0.599,-0.525) (-0.643,-0.550) (-0.700,-0.584) (-0.750,-0.614) (-0.800,-0.645) (-0.850,-0.675) (-0.900,-0.707) (-0.950,-0.738) (-1.000,-0.770) (-1.047,-0.800) (-1.100,-0.834) (-1.150,-0.866) (-1.200,-0.898) (-1.250,-0.931) (-1.300,-0.963) (-1.350,-0.996) (-1.396,-1.025) (-1.450,-1.060) (-1.500,-1.093) (-1.550,-1.125) (-1.600,-1.157) (-1.650,-1.189) (-1.700,-1.221) (-1.746,-1.250) (-1.800,-1.284) (-1.850,-1.316) (-1.900,-1.347) (-1.945,-1.375) (-2.000,-1.409) (-2.050,-1.440) (-2.100,-1.470) (-2.150,-1.500) (-2.200,-1.530) (-2.250,-1.560) (-2.300,-1.589) (-2.350,-1.618) (-2.400,-1.647) (-2.450,-1.675) (-2.495,-1.700) (-2.540,-1.725) (-2.586,-1.750) (-2.632,-1.775) (-2.679,-1.800) (-2.725,-1.825) (-2.774,-1.850) (-2.823,-1.875) (-2.873,-1.900) (-2.924,-1.925) (-2.975,-1.950) (-3.025,-1.973) (-3.075,-1.997) (-3.125,-2.019) (-3.175,-2.042) (-3.225,-2.063) (-3.275,-2.085) (-3.325,-2.105) (-3.374,-2.125) (-3.425,-2.145) (-3.475,-2.164) (-3.525,-2.182) (-3.574,-2.200) (-3.625,-2.218) (-3.675,-2.234) (-3.724,-2.250) (-3.775,-2.266) (-3.825,-2.281) (-3.875,-2.295) (-3.925,-2.309) (-3.975,-2.322) (-4.025,-2.334) (-4.075,-2.346) (-4.125,-2.357) (-4.175,-2.367) (-4.225,-2.376) (-4.275,-2.385) (-4.325,-2.394) (-4.375,-2.401) (-4.425,-2.408) (-4.475,-2.414) (-4.525,-2.419) (-4.575,-2.424) (-4.625,-2.427) (-4.675,-2.430) (-4.725,-2.433) (-4.775,-2.434) (-4.825,-2.435) (-4.875,-2.435) (-4.925,-2.434) (-4.975,-2.432) (-5.025,-2.429) (-5.075,-2.425)}
\EQ{(10.000,-7.078) (9.950,-7.043) (9.900,-7.007) (9.854,-6.975) (9.800,-6.937) (9.750,-6.901) (9.700,-6.866) (9.650,-6.831) (9.607,-6.800) (9.550,-6.760) (9.500,-6.725) (9.450,-6.689) (9.400,-6.654) (9.359,-6.625) (9.300,-6.584) (9.253,-6.550) (9.200,-6.513) (9.150,-6.478) (9.100,-6.442) (9.050,-6.407) (9.005,-6.375) (8.950,-6.336) (8.900,-6.301) (8.850,-6.266) (8.800,-6.230) (8.757,-6.200) (8.700,-6.160) (8.651,-6.125) (8.600,-6.089) (8.550,-6.054) (8.500,-6.019) (8.450,-5.983) (8.403,-5.950) (8.350,-5.913) (8.300,-5.877) (8.250,-5.842) (8.200,-5.807) (8.155,-5.775) (8.100,-5.736) (8.050,-5.701) (8.000,-5.665) (7.950,-5.630) (7.907,-5.600) (7.850,-5.560) (7.801,-5.525) (7.750,-5.489) (7.700,-5.454) (7.650,-5.418) (7.600,-5.383) (7.553,-5.350) (7.500,-5.312) (7.450,-5.277) (7.400,-5.242) (7.350,-5.207) (7.305,-5.175) (7.250,-5.136) (7.200,-5.101) (7.150,-5.065) (7.100,-5.030) (7.057,-5.000) (7.000,-4.960) (6.951,-4.925) (6.900,-4.889) (6.850,-4.854) (6.800,-4.818) (6.750,-4.783) (6.703,-4.750) (6.650,-4.713) (6.600,-4.677) (6.550,-4.642) (6.500,-4.607) (6.455,-4.575) (6.400,-4.536) (6.350,-4.501) (6.300,-4.466) (6.250,-4.430) (6.207,-4.400) (6.150,-4.360) (6.101,-4.325) (6.050,-4.289) (6.000,-4.254) (5.959,-4.225) (5.900,-4.184) (5.852,-4.150) (5.800,-4.113) (5.750,-4.078) (5.700,-4.043) (5.650,-4.007) (5.604,-3.975) (5.550,-3.937) (5.500,-3.902) (5.450,-3.866) (5.400,-3.831) (5.356,-3.800) (5.300,-3.761) (5.250,-3.725) (5.200,-3.690) (5.150,-3.655) (5.108,-3.625) (5.050,-3.585) (5.001,-3.550) (4.950,-3.514) (4.900,-3.479) (4.859,-3.450) (4.800,-3.408) (4.753,-3.375) (4.700,-3.338) (4.650,-3.303) (4.600,-3.268) (4.550,-3.232) (4.504,-3.200) (4.450,-3.162) (4.400,-3.127) (4.350,-3.092) (4.300,-3.057) (4.255,-3.025) (4.200,-2.986) (4.150,-2.951) (4.100,-2.916) (4.050,-2.881) (4.006,-2.850) (3.950,-2.810) (3.900,-2.775) (3.850,-2.740) (3.800,-2.705) (3.757,-2.675) (3.700,-2.635) (3.650,-2.600) (3.600,-2.565) (3.550,-2.530) (3.508,-2.500) (3.450,-2.459) (3.401,-2.425) (3.350,-2.389) (3.300,-2.354) (3.258,-2.325) (3.200,-2.284) (3.151,-2.250) (3.100,-2.214) (3.050,-2.179) (3.009,-2.150) (2.950,-2.109) (2.901,-2.075) (2.850,-2.039) (2.800,-2.004) (2.758,-1.975) (2.700,-1.935) (2.651,-1.900) (2.600,-1.865) (2.550,-1.830) (2.507,-1.800) (2.450,-1.760) (2.400,-1.725) (2.350,-1.691) (2.300,-1.656) (2.256,-1.625) (2.200,-1.587) (2.150,-1.552) (2.100,-1.517) (2.050,-1.483) (2.003,-1.450) (1.950,-1.414) (1.900,-1.379) (1.858,-1.350) (1.800,-1.311) (1.750,-1.276) (1.700,-1.242) (1.650,-1.208) (1.602,-1.175) (1.550,-1.140) (1.500,-1.106) (1.455,-1.075) (1.400,-1.038) (1.350,-1.004) (1.307,-0.975) (1.250,-0.937) (1.200,-0.904) (1.157,-0.875) (1.100,-0.838) (1.050,-0.805) (1.004,-0.775) (0.950,-0.741) (0.900,-0.709) (0.850,-0.677) (0.807,-0.650) (0.750,-0.615) (0.700,-0.585) (0.650,-0.555) (0.600,-0.526) (0.554,-0.500) (0.509,-0.475) (0.461,-0.450) (0.411,-0.425) (0.358,-0.400) (0.300,-0.376) (0.250,-0.357) (0.200,-0.341) (0.150,-0.327) (0.100,-0.317) (0.050,-0.311) (0.000,-0.309) (-0.050,-0.311) (-0.100,-0.317) (-0.150,-0.327) (-0.200,-0.341) (-0.250,-0.357) (-0.298,-0.375) (-0.350,-0.397) (-0.400,-0.420) (-0.450,-0.445) (-0.500,-0.471) (-0.550,-0.498) (-0.598,-0.525) (-0.642,-0.550) (-0.700,-0.585) (-0.750,-0.615) (-0.800,-0.646) (-0.847,-0.675) (-0.900,-0.709) (-0.950,-0.741) (-1.000,-0.773) (-1.042,-0.800) (-1.100,-0.838) (-1.150,-0.871) (-1.194,-0.900) (-1.250,-0.937) (-1.300,-0.971) (-1.344,-1.000) (-1.400,-1.038) (-1.450,-1.072) (-1.492,-1.100) (-1.550,-1.140) (-1.600,-1.174) (-1.650,-1.208) (-1.700,-1.242) (-1.748,-1.275) (-1.800,-1.311) (-1.850,-1.345) (-1.894,-1.375) (-1.950,-1.414) (-2.000,-1.448) (-2.050,-1.483) (-2.100,-1.517) (-2.147,-1.550) (-2.200,-1.587) (-2.250,-1.621) (-2.292,-1.650) (-2.350,-1.691) (-2.399,-1.725) (-2.450,-1.760) (-2.500,-1.795) (-2.543,-1.825) (-2.600,-1.865) (-2.650,-1.900) (-2.700,-1.935) (-2.750,-1.969) (-2.794,-2.000) (-2.850,-2.039) (-2.900,-2.074) (-2.950,-2.109) (-3.000,-2.144) (-3.044,-2.175) (-3.100,-2.214) (-3.150,-2.249) (-3.200,-2.284) (-3.250,-2.319) (-3.294,-2.350) (-3.350,-2.389) (-3.400,-2.424) (-3.450,-2.459) (-3.500,-2.494) (-3.544,-2.525) (-3.600,-2.565) (-3.650,-2.600) (-3.700,-2.635) (-3.750,-2.670) (-3.793,-2.700) (-3.850,-2.740) (-3.900,-2.775) (-3.950,-2.810) (-4.000,-2.846) (-4.042,-2.875) (-4.100,-2.916) (-4.149,-2.950) (-4.200,-2.986) (-4.250,-3.021) (-4.300,-3.057) (-4.350,-3.092) (-4.397,-3.125) (-4.450,-3.162) (-4.500,-3.197) (-4.550,-3.232) (-4.600,-3.268) (-4.646,-3.300) (-4.700,-3.338) (-4.750,-3.373) (-4.800,-3.408) (-4.850,-3.444) (-4.895,-3.475) (-4.950,-3.514) (-5.000,-3.549) (-5.050,-3.585) (-5.100,-3.620) (-5.143,-3.650) (-5.200,-3.690) (-5.249,-3.725) (-5.300,-3.761) (-5.350,-3.796) (-5.391,-3.825) (-5.450,-3.866) (-5.498,-3.900) (-5.550,-3.937) (-5.600,-3.972) (-5.650,-4.007) (-5.700,-4.043) (-5.746,-4.075) (-5.800,-4.113) (-5.850,-4.148) (-5.900,-4.184) (-5.950,-4.219) (-5.994,-4.250) (-6.050,-4.289) (-6.100,-4.325) (-6.150,-4.360) (-6.200,-4.395) (-6.242,-4.425) (-6.300,-4.466) (-6.349,-4.500) (-6.400,-4.536) (-6.450,-4.572) (-6.500,-4.607) (-6.550,-4.642) (-6.597,-4.675) (-6.650,-4.713) (-6.700,-4.748) (-6.750,-4.783) (-6.800,-4.818) (-6.845,-4.850) (-6.900,-4.889) (-6.950,-4.924) (-7.000,-4.960) (-7.050,-4.995) (-7.093,-5.025) (-7.150,-5.065) (-7.199,-5.100) (-7.250,-5.136) (-7.300,-5.171) (-7.350,-5.207) (-7.400,-5.242) (-7.447,-5.275) (-7.500,-5.312) (-7.550,-5.348) (-7.600,-5.383) (-7.650,-5.418) (-7.695,-5.450) (-7.750,-5.489) (-7.800,-5.524) (-7.850,-5.560) (-7.900,-5.595) (-7.943,-5.625) (-8.000,-5.665) (-8.049,-5.700) (-8.100,-5.736) (-8.150,-5.771) (-8.200,-5.807) (-8.250,-5.842) (-8.297,-5.875) (-8.350,-5.913) (-8.400,-5.948) (-8.450,-5.983) (-8.500,-6.019) (-8.545,-6.050) (-8.600,-6.089) (-8.650,-6.124) (-8.700,-6.160) (-8.750,-6.195) (-8.792,-6.225) (-8.850,-6.266) (-8.899,-6.300) (-8.950,-6.336) (-9.000,-6.372) (-9.050,-6.407) (-9.100,-6.442) (-9.146,-6.475) (-9.200,-6.513) (-9.250,-6.548) (-9.300,-6.584) (-9.350,-6.619) (-9.394,-6.650) (-9.450,-6.689) (-9.500,-6.725) (-9.550,-6.760) (-9.600,-6.795) (-9.642,-6.825) (-9.700,-6.866) (-9.748,-6.900) (-9.800,-6.937) (-9.850,-6.972) (-9.900,-7.007) (-9.950,-7.043) (-9.996,-7.075) (-10.000,-7.078)}
\EQ{(-10.000,7.078) (-9.950,7.043) (-9.900,7.007) (-9.854,6.975) (-9.800,6.937) (-9.750,6.901) (-9.700,6.866) (-9.650,6.831) (-9.607,6.800) (-9.550,6.760) (-9.500,6.725) (-9.450,6.689) (-9.400,6.654) (-9.359,6.625) (-9.300,6.584) (-9.253,6.550) (-9.200,6.513) (-9.150,6.478) (-9.100,6.442) (-9.050,6.407) (-9.005,6.375) (-8.950,6.336) (-8.900,6.301) (-8.850,6.266) (-8.800,6.230) (-8.757,6.200) (-8.700,6.160) (-8.651,6.125) (-8.600,6.089) (-8.550,6.054) (-8.500,6.019) (-8.450,5.983) (-8.403,5.950) (-8.350,5.913) (-8.300,5.877) (-8.250,5.842) (-8.200,5.807) (-8.155,5.775) (-8.100,5.736) (-8.050,5.701) (-8.000,5.665) (-7.950,5.630) (-7.907,5.600) (-7.850,5.560) (-7.801,5.525) (-7.750,5.489) (-7.700,5.454) (-7.650,5.418) (-7.600,5.383) (-7.553,5.350) (-7.500,5.312) (-7.450,5.277) (-7.400,5.242) (-7.350,5.207) (-7.305,5.175) (-7.250,5.136) (-7.200,5.101) (-7.150,5.065) (-7.100,5.030) (-7.057,5.000) (-7.000,4.960) (-6.951,4.925) (-6.900,4.889) (-6.850,4.854) (-6.800,4.818) (-6.750,4.783) (-6.703,4.750) (-6.650,4.713) (-6.600,4.677) (-6.550,4.642) (-6.500,4.607) (-6.455,4.575) (-6.400,4.536) (-6.350,4.501) (-6.300,4.466) (-6.250,4.430) (-6.207,4.400) (-6.150,4.360) (-6.101,4.325) (-6.050,4.289) (-6.000,4.254) (-5.959,4.225) (-5.900,4.184) (-5.852,4.150) (-5.800,4.113) (-5.750,4.078) (-5.700,4.043) (-5.650,4.007) (-5.604,3.975) (-5.550,3.937) (-5.500,3.902) (-5.450,3.866) (-5.400,3.831) (-5.356,3.800) (-5.300,3.761) (-5.250,3.725) (-5.200,3.690) (-5.150,3.655) (-5.108,3.625) (-5.050,3.585) (-5.001,3.550) (-4.950,3.514) (-4.900,3.479) (-4.859,3.450) (-4.800,3.408) (-4.753,3.375) (-4.700,3.338) (-4.650,3.303) (-4.600,3.268) (-4.550,3.232) (-4.504,3.200) (-4.450,3.162) (-4.400,3.127) (-4.350,3.092) (-4.300,3.057) (-4.255,3.025) (-4.200,2.986) (-4.150,2.951) (-4.100,2.916) (-4.050,2.881) (-4.006,2.850) (-3.950,2.810) (-3.900,2.775) (-3.850,2.740) (-3.800,2.705) (-3.757,2.675) (-3.700,2.635) (-3.650,2.600) (-3.600,2.565) (-3.550,2.530) (-3.508,2.500) (-3.450,2.459) (-3.401,2.425) (-3.350,2.389) (-3.300,2.354) (-3.258,2.325) (-3.200,2.284) (-3.151,2.250) (-3.100,2.214) (-3.050,2.179) (-3.009,2.150) (-2.950,2.109) (-2.901,2.075) (-2.850,2.039) (-2.800,2.004) (-2.758,1.975) (-2.700,1.935) (-2.651,1.900) (-2.600,1.865) (-2.550,1.830) (-2.507,1.800) (-2.450,1.760) (-2.400,1.725) (-2.350,1.691) (-2.300,1.656) (-2.256,1.625) (-2.200,1.587) (-2.150,1.552) (-2.100,1.517) (-2.050,1.483) (-2.003,1.450) (-1.950,1.414) (-1.900,1.379) (-1.858,1.350) (-1.800,1.311) (-1.750,1.276) (-1.700,1.242) (-1.650,1.208) (-1.602,1.175) (-1.550,1.140) (-1.500,1.106) (-1.455,1.075) (-1.400,1.038) (-1.350,1.004) (-1.307,0.975) (-1.250,0.937) (-1.200,0.904) (-1.157,0.875) (-1.100,0.838) (-1.050,0.805) (-1.004,0.775) (-0.950,0.741) (-0.900,0.709) (-0.850,0.677) (-0.807,0.650) (-0.750,0.615) (-0.700,0.585) (-0.650,0.555) (-0.600,0.526) (-0.554,0.500) (-0.509,0.475) (-0.461,0.450) (-0.411,0.425) (-0.358,0.400) (-0.300,0.376) (-0.250,0.357) (-0.200,0.341) (-0.150,0.327) (-0.100,0.317) (-0.050,0.311) (0.000,0.309) (0.050,0.311) (0.100,0.317) (0.150,0.327) (0.200,0.341) (0.250,0.357) (0.298,0.375) (0.350,0.397) (0.400,0.420) (0.450,0.445) (0.500,0.471) (0.550,0.498) (0.598,0.525) (0.642,0.550) (0.700,0.585) (0.750,0.615) (0.800,0.646) (0.847,0.675) (0.900,0.709) (0.950,0.741) (1.000,0.773) (1.042,0.800) (1.100,0.838) (1.150,0.871) (1.194,0.900) (1.250,0.937) (1.300,0.971) (1.344,1.000) (1.400,1.038) (1.450,1.072) (1.492,1.100) (1.550,1.140) (1.600,1.174) (1.650,1.208) (1.700,1.242) (1.748,1.275) (1.800,1.311) (1.850,1.345) (1.894,1.375) (1.950,1.414) (2.000,1.448) (2.050,1.483) (2.100,1.517) (2.147,1.550) (2.200,1.587) (2.250,1.621) (2.292,1.650) (2.350,1.691) (2.399,1.725) (2.450,1.760) (2.500,1.795) (2.543,1.825) (2.600,1.865) (2.650,1.900) (2.700,1.935) (2.750,1.969) (2.794,2.000) (2.850,2.039) (2.900,2.074) (2.950,2.109) (3.000,2.144) (3.044,2.175) (3.100,2.214) (3.150,2.249) (3.200,2.284) (3.250,2.319) (3.294,2.350) (3.350,2.389) (3.400,2.424) (3.450,2.459) (3.500,2.494) (3.544,2.525) (3.600,2.565) (3.650,2.600) (3.700,2.635) (3.750,2.670) (3.793,2.700) (3.850,2.740) (3.900,2.775) (3.950,2.810) (4.000,2.846) (4.042,2.875) (4.100,2.916) (4.149,2.950) (4.200,2.986) (4.250,3.021) (4.300,3.057) (4.350,3.092) (4.397,3.125) (4.450,3.162) (4.500,3.197) (4.550,3.232) (4.600,3.268) (4.646,3.300) (4.700,3.338) (4.750,3.373) (4.800,3.408) (4.850,3.444) (4.895,3.475) (4.950,3.514) (5.000,3.549) (5.050,3.585) (5.100,3.620) (5.143,3.650) (5.200,3.690) (5.249,3.725) (5.300,3.761) (5.350,3.796) (5.391,3.825) (5.450,3.866) (5.498,3.900) (5.550,3.937) (5.600,3.972) (5.650,4.007) (5.700,4.043) (5.746,4.075) (5.800,4.113) (5.850,4.148) (5.900,4.184) (5.950,4.219) (5.994,4.250) (6.050,4.289) (6.100,4.325) (6.150,4.360) (6.200,4.395) (6.242,4.425) (6.300,4.466) (6.349,4.500) (6.400,4.536) (6.450,4.572) (6.500,4.607) (6.550,4.642) (6.597,4.675) (6.650,4.713) (6.700,4.748) (6.750,4.783) (6.800,4.818) (6.845,4.850) (6.900,4.889) (6.950,4.924) (7.000,4.960) (7.050,4.995) (7.093,5.025) (7.150,5.065) (7.199,5.100) (7.250,5.136) (7.300,5.171) (7.350,5.207) (7.400,5.242) (7.447,5.275) (7.500,5.312) (7.550,5.348) (7.600,5.383) (7.650,5.418) (7.695,5.450) (7.750,5.489) (7.800,5.524) (7.850,5.560) (7.900,5.595) (7.943,5.625) (8.000,5.665) (8.049,5.700) (8.100,5.736) (8.150,5.771) (8.200,5.807) (8.250,5.842) (8.297,5.875) (8.350,5.913) (8.400,5.948) (8.450,5.983) (8.500,6.019) (8.545,6.050) (8.600,6.089) (8.650,6.124) (8.700,6.160) (8.750,6.195) (8.792,6.225) (8.850,6.266) (8.899,6.300) (8.950,6.336) (9.000,6.372) (9.050,6.407) (9.100,6.442) (9.146,6.475) (9.200,6.513) (9.250,6.548) (9.300,6.584) (9.350,6.619) (9.394,6.650) (9.450,6.689) (9.500,6.725) (9.550,6.760) (9.600,6.795) (9.642,6.825) (9.700,6.866) (9.748,6.900) (9.800,6.937) (9.850,6.972) (9.900,7.007) (9.950,7.043) (9.996,7.075) (10.000,7.078)}
\EQ{(-0.375,-9.277) (-0.325,-9.281) (-0.275,-9.284) (-0.225,-9.287) (-0.175,-9.289) (-0.125,-9.291) (-0.075,-9.292) (-0.025,-9.292) (0.050,-9.292) (0.100,-9.291) (0.150,-9.290) (0.200,-9.288) (0.250,-9.285) (0.300,-9.282) (0.350,-9.279) (0.400,-9.275) (0.450,-9.270) (0.500,-9.265) (0.550,-9.259) (0.600,-9.253) (0.650,-9.246) (0.700,-9.239) (0.750,-9.231) (0.800,-9.223) (0.850,-9.213) (0.900,-9.204) (0.950,-9.194) (1.000,-9.183) (1.050,-9.171) (1.100,-9.160) (1.150,-9.147) (1.200,-9.134) (1.250,-9.120) (1.300,-9.106) (1.350,-9.091) (1.400,-9.075) (1.450,-9.059) (1.500,-9.042) (1.548,-9.025) (1.600,-9.006) (1.650,-8.987) (1.700,-8.968) (1.750,-8.947) (1.800,-8.926) (1.850,-8.905) (1.900,-8.882) (1.950,-8.859) (2.000,-8.835) (2.050,-8.811) (2.100,-8.785) (2.150,-8.759) (2.200,-8.732) (2.250,-8.704) (2.300,-8.675) (2.350,-8.646) (2.400,-8.615) (2.450,-8.583) (2.500,-8.551) (2.550,-8.518) (2.600,-8.483) (2.647,-8.450) (2.700,-8.411) (2.748,-8.375) (2.800,-8.335) (2.843,-8.300) (2.900,-8.253) (2.950,-8.211) (2.991,-8.175) (3.046,-8.125) (3.099,-8.075) (3.150,-8.026) (3.200,-7.976) (3.249,-7.925) (3.296,-7.875) (3.341,-7.825) (3.375,-7.786) (3.425,-7.727) (3.468,-7.675) (3.500,-7.634) (3.545,-7.575) (3.583,-7.525) (3.619,-7.475) (3.650,-7.430) (3.687,-7.375) (3.720,-7.325) (3.750,-7.277) (3.782,-7.225) (3.811,-7.175) (3.840,-7.125) (3.867,-7.075) (3.894,-7.025) (3.920,-6.975) (3.944,-6.925) (3.968,-6.875) (3.991,-6.825) (4.013,-6.775) (4.035,-6.725) (4.055,-6.675) (4.075,-6.625) (4.093,-6.575) (4.111,-6.525) (4.129,-6.475) (4.145,-6.425) (4.161,-6.375) (4.176,-6.325) (4.190,-6.275) (4.203,-6.225) (4.216,-6.175) (4.228,-6.125) (4.239,-6.075) (4.249,-6.025) (4.259,-5.975) (4.268,-5.925) (4.276,-5.875) (4.284,-5.825) (4.290,-5.775) (4.296,-5.725) (4.302,-5.675) (4.307,-5.625) (4.311,-5.575) (4.314,-5.525) (4.317,-5.475) (4.318,-5.425) (4.320,-5.375) (4.320,-5.325) (4.320,-5.275) (4.319,-5.225) (4.318,-5.175) (4.315,-5.125) (4.312,-5.075) (4.309,-5.025) (4.304,-4.975) (4.299,-4.925) (4.294,-4.875) (4.287,-4.825) (4.280,-4.775) (4.272,-4.725) (4.263,-4.675) (4.254,-4.625) (4.244,-4.575) (4.233,-4.525) (4.222,-4.475) (4.210,-4.425) (4.197,-4.375) (4.183,-4.325) (4.169,-4.275) (4.153,-4.225) (4.138,-4.175) (4.121,-4.125) (4.103,-4.075) (4.085,-4.025) (4.066,-3.975) (4.046,-3.925) (4.026,-3.875) (4.004,-3.825) (3.982,-3.775) (3.959,-3.725) (3.935,-3.675) (3.911,-3.625) (3.885,-3.575) (3.859,-3.525) (3.832,-3.475) (3.803,-3.425) (3.775,-3.376) (3.745,-3.325) (3.714,-3.275) (3.683,-3.225) (3.650,-3.175) (3.617,-3.125) (3.582,-3.075) (3.550,-3.029) (3.511,-2.975) (3.475,-2.926) (3.436,-2.875) (3.400,-2.829) (3.357,-2.775) (3.325,-2.736) (3.275,-2.676) (3.232,-2.625) (3.188,-2.575) (3.150,-2.533) (3.100,-2.478) (3.050,-2.425) (3.003,-2.375) (2.954,-2.325) (2.904,-2.275) (2.854,-2.225) (2.802,-2.175) (2.750,-2.126) (2.700,-2.079) (2.650,-2.033) (2.600,-1.988) (2.558,-1.950) (2.501,-1.900) (2.450,-1.856) (2.400,-1.814) (2.354,-1.775) (2.300,-1.730) (2.250,-1.689) (2.202,-1.650) (2.150,-1.608) (2.109,-1.575) (2.050,-1.529) (2.000,-1.490) (1.950,-1.452) (1.900,-1.413) (1.850,-1.375) (1.800,-1.338) (1.750,-1.301) (1.700,-1.264) (1.650,-1.227) (1.600,-1.191) (1.550,-1.155) (1.509,-1.125) (1.450,-1.083) (1.403,-1.050) (1.350,-1.013) (1.300,-0.978) (1.250,-0.944) (1.200,-0.910) (1.150,-0.876) (1.100,-0.842) (1.050,-0.809) (1.000,-0.776) (0.950,-0.743) (0.900,-0.711) (0.850,-0.679) (0.805,-0.650) (0.750,-0.616) (0.700,-0.586) (0.650,-0.556) (0.600,-0.527) (0.553,-0.500) (0.508,-0.475) (0.461,-0.450) (0.411,-0.425) (0.357,-0.400) (0.300,-0.376) (0.250,-0.357) (0.200,-0.341) (0.150,-0.327) (0.100,-0.317) (0.050,-0.311) (0.000,-0.309) (-0.050,-0.311) (-0.100,-0.317) (-0.150,-0.327) (-0.200,-0.341) (-0.250,-0.357) (-0.298,-0.375) (-0.350,-0.397) (-0.400,-0.420) (-0.450,-0.445) (-0.500,-0.471) (-0.550,-0.498) (-0.598,-0.525) (-0.641,-0.550) (-0.700,-0.586) (-0.750,-0.616) (-0.800,-0.647) (-0.845,-0.675) (-0.900,-0.711) (-0.950,-0.743) (-0.999,-0.775) (-1.050,-0.809) (-1.100,-0.842) (-1.149,-0.875) (-1.200,-0.910) (-1.250,-0.944) (-1.295,-0.975) (-1.350,-1.013) (-1.400,-1.048) (-1.450,-1.083) (-1.500,-1.119) (-1.544,-1.150) (-1.600,-1.191) (-1.647,-1.225) (-1.700,-1.264) (-1.749,-1.300) (-1.800,-1.338) (-1.849,-1.375) (-1.900,-1.413) (-1.948,-1.450) (-2.000,-1.490) (-2.045,-1.525) (-2.100,-1.568) (-2.140,-1.600) (-2.200,-1.648) (-2.250,-1.689) (-2.294,-1.725) (-2.350,-1.772) (-2.400,-1.814) (-2.443,-1.850) (-2.500,-1.900) (-2.550,-1.943) (-2.600,-1.988) (-2.641,-2.025) (-2.696,-2.075) (-2.749,-2.125) (-2.800,-2.173) (-2.850,-2.222) (-2.900,-2.271) (-2.950,-2.321) (-3.000,-2.372) (-3.050,-2.425) (-3.097,-2.475) (-3.143,-2.525) (-3.188,-2.575) (-3.225,-2.618) (-3.274,-2.675) (-3.316,-2.725) (-3.350,-2.766) (-3.397,-2.825) (-3.436,-2.875) (-3.474,-2.925) (-3.511,-2.975) (-3.547,-3.025) (-3.582,-3.075) (-3.617,-3.125) (-3.650,-3.175) (-3.683,-3.225) (-3.714,-3.275) (-3.745,-3.325) (-3.775,-3.375) (-3.800,-3.419) (-3.825,-3.463) (-3.850,-3.509) (-3.875,-3.556) (-3.900,-3.604) (-3.923,-3.650) (-3.947,-3.700) (-3.971,-3.750) (-3.993,-3.800) (-4.015,-3.850) (-4.036,-3.900) (-4.056,-3.950) (-4.075,-3.998) (-4.094,-4.050) (-4.112,-4.100) (-4.129,-4.150) (-4.146,-4.200) (-4.161,-4.250) (-4.176,-4.300) (-4.190,-4.350) (-4.203,-4.400) (-4.216,-4.450) (-4.228,-4.500) (-4.239,-4.550) (-4.249,-4.600) (-4.259,-4.650) (-4.268,-4.700) (-4.276,-4.750) (-4.284,-4.800) (-4.290,-4.850) (-4.297,-4.900) (-4.302,-4.950) (-4.307,-5.000) (-4.311,-5.050) (-4.314,-5.100) (-4.317,-5.150) (-4.318,-5.200) (-4.320,-5.250) (-4.320,-5.300) (-4.320,-5.350) (-4.319,-5.400) (-4.318,-5.450) (-4.315,-5.500) (-4.312,-5.550) (-4.309,-5.600) (-4.304,-5.650) (-4.299,-5.700) (-4.294,-5.750) (-4.287,-5.800) (-4.280,-5.850) (-4.272,-5.900) (-4.263,-5.950) (-4.254,-6.000) (-4.244,-6.050) (-4.233,-6.100) (-4.222,-6.150) (-4.210,-6.200) (-4.197,-6.250) (-4.183,-6.300) (-4.168,-6.350) (-4.153,-6.400) (-4.137,-6.450) (-4.120,-6.500) (-4.103,-6.550) (-4.084,-6.600) (-4.065,-6.650) (-4.045,-6.700) (-4.025,-6.748) (-4.002,-6.800) (-3.980,-6.850) (-3.956,-6.900) (-3.932,-6.950) (-3.907,-7.000) (-3.881,-7.050) (-3.854,-7.100) (-3.826,-7.150) (-3.800,-7.194) (-3.767,-7.250) (-3.736,-7.300) (-3.704,-7.350) (-3.675,-7.393) (-3.636,-7.450) (-3.601,-7.500) (-3.564,-7.550) (-3.526,-7.600) (-3.487,-7.650) (-3.450,-7.697) (-3.406,-7.750) (-3.363,-7.800) (-3.325,-7.843) (-3.275,-7.897) (-3.225,-7.950) (-3.176,-8.000) (-3.125,-8.050) (-3.075,-8.098) (-3.025,-8.144) (-2.975,-8.189) (-2.933,-8.225) (-2.875,-8.274) (-2.825,-8.315) (-2.780,-8.350) (-2.725,-8.392) (-2.681,-8.425) (-2.625,-8.466) (-2.576,-8.500) (-2.525,-8.534) (-2.475,-8.567) (-2.425,-8.599) (-2.375,-8.630) (-2.325,-8.660) (-2.275,-8.690) (-2.225,-8.718) (-2.175,-8.745) (-2.125,-8.772) (-2.075,-8.798) (-2.025,-8.823) (-1.975,-8.847) (-1.925,-8.871) (-1.875,-8.894) (-1.825,-8.916) (-1.775,-8.937) (-1.725,-8.958) (-1.675,-8.977) (-1.625,-8.997) (-1.575,-9.015) (-1.525,-9.033) (-1.476,-9.050) (-1.425,-9.067) (-1.375,-9.083) (-1.325,-9.098) (-1.275,-9.113) (-1.225,-9.127) (-1.175,-9.140) (-1.125,-9.153) (-1.075,-9.166) (-1.025,-9.177) (-0.975,-9.188) (-0.925,-9.199) (-0.875,-9.209) (-0.825,-9.218) (-0.775,-9.227) (-0.725,-9.235) (-0.675,-9.243) (-0.625,-9.250) (-0.575,-9.256) (-0.525,-9.262) (-0.475,-9.268) (-0.425,-9.273) (-0.375,-9.277)}
\EQ{(-0.125,0.322) (-0.075,0.313) (-0.025,0.309) (0.050,0.311) (0.100,0.317) (0.150,0.327) (0.200,0.341) (0.250,0.357) (0.298,0.375) (0.350,0.397) (0.400,0.420) (0.450,0.445) (0.500,0.471) (0.550,0.498) (0.598,0.525) (0.641,0.550) (0.700,0.586) (0.750,0.616) (0.800,0.647) (0.845,0.675) (0.900,0.711) (0.950,0.743) (0.999,0.775) (1.050,0.809) (1.100,0.842) (1.149,0.875) (1.200,0.910) (1.250,0.944) (1.295,0.975) (1.350,1.013) (1.400,1.048) (1.450,1.083) (1.500,1.119) (1.544,1.150) (1.600,1.191) (1.647,1.225) (1.700,1.264) (1.749,1.300) (1.800,1.338) (1.849,1.375) (1.900,1.413) (1.948,1.450) (2.000,1.490) (2.045,1.525) (2.100,1.568) (2.140,1.600) (2.200,1.648) (2.250,1.689) (2.294,1.725) (2.350,1.772) (2.400,1.814) (2.443,1.850) (2.500,1.900) (2.550,1.943) (2.600,1.988) (2.641,2.025) (2.696,2.075) (2.749,2.125) (2.800,2.173) (2.850,2.222) (2.900,2.271) (2.950,2.321) (3.000,2.372) (3.050,2.425) (3.097,2.475) (3.143,2.525) (3.188,2.575) (3.225,2.618) (3.274,2.675) (3.316,2.725) (3.350,2.766) (3.397,2.825) (3.436,2.875) (3.474,2.925) (3.511,2.975) (3.547,3.025) (3.582,3.075) (3.617,3.125) (3.650,3.175) (3.683,3.225) (3.714,3.275) (3.745,3.325) (3.775,3.375) (3.800,3.419) (3.825,3.463) (3.850,3.509) (3.875,3.556) (3.900,3.604) (3.923,3.650) (3.947,3.700) (3.971,3.750) (3.993,3.800) (4.015,3.850) (4.036,3.900) (4.056,3.950) (4.075,3.998) (4.094,4.050) (4.112,4.100) (4.129,4.150) (4.146,4.200) (4.161,4.250) (4.176,4.300) (4.190,4.350) (4.203,4.400) (4.216,4.450) (4.228,4.500) (4.239,4.550) (4.249,4.600) (4.259,4.650) (4.268,4.700) (4.276,4.750) (4.284,4.800) (4.290,4.850) (4.297,4.900) (4.302,4.950) (4.307,5.000) (4.311,5.050) (4.314,5.100) (4.317,5.150) (4.318,5.200) (4.320,5.250) (4.320,5.300) (4.320,5.350) (4.319,5.400) (4.318,5.450) (4.315,5.500) (4.312,5.550) (4.309,5.600) (4.304,5.650) (4.299,5.700) (4.294,5.750) (4.287,5.800) (4.280,5.850) (4.272,5.900) (4.263,5.950) (4.254,6.000) (4.244,6.050) (4.233,6.100) (4.222,6.150) (4.210,6.200) (4.197,6.250) (4.183,6.300) (4.168,6.350) (4.153,6.400) (4.137,6.450) (4.120,6.500) (4.103,6.550) (4.084,6.600) (4.065,6.650) (4.045,6.700) (4.025,6.748) (4.002,6.800) (3.980,6.850) (3.956,6.900) (3.932,6.950) (3.907,7.000) (3.881,7.050) (3.854,7.100) (3.826,7.150) (3.800,7.194) (3.767,7.250) (3.736,7.300) (3.704,7.350) (3.675,7.393) (3.636,7.450) (3.601,7.500) (3.564,7.550) (3.526,7.600) (3.487,7.650) (3.450,7.697) (3.406,7.750) (3.363,7.800) (3.325,7.843) (3.275,7.897) (3.225,7.950) (3.176,8.000) (3.125,8.050) (3.075,8.098) (3.025,8.144) (2.975,8.189) (2.933,8.225) (2.875,8.274) (2.825,8.315) (2.780,8.350) (2.725,8.392) (2.681,8.425) (2.625,8.466) (2.576,8.500) (2.525,8.534) (2.475,8.567) (2.425,8.599) (2.375,8.630) (2.325,8.660) (2.275,8.690) (2.225,8.718) (2.175,8.745) (2.125,8.772) (2.075,8.798) (2.025,8.823) (1.975,8.847) (1.925,8.871) (1.875,8.894) (1.825,8.916) (1.775,8.937) (1.725,8.958) (1.675,8.977) (1.625,8.997) (1.575,9.015) (1.525,9.033) (1.476,9.050) (1.425,9.067) (1.375,9.083) (1.325,9.098) (1.275,9.113) (1.225,9.127) (1.175,9.140) (1.125,9.153) (1.075,9.166) (1.025,9.177) (0.975,9.188) (0.925,9.199) (0.875,9.209) (0.825,9.218) (0.775,9.227) (0.725,9.235) (0.675,9.243) (0.625,9.250) (0.575,9.256) (0.525,9.262) (0.475,9.268) (0.425,9.273) (0.375,9.277) (0.325,9.281) (0.275,9.284) (0.225,9.287) (0.175,9.289) (0.125,9.291) (0.075,9.292) (0.025,9.292) (-0.050,9.292) (-0.100,9.291) (-0.150,9.290) (-0.200,9.288) (-0.250,9.285) (-0.300,9.282) (-0.350,9.279) (-0.400,9.275) (-0.450,9.270) (-0.500,9.265) (-0.550,9.259) (-0.600,9.253) (-0.650,9.246) (-0.700,9.239) (-0.750,9.231) (-0.800,9.223) (-0.850,9.213) (-0.900,9.204) (-0.950,9.194) (-1.000,9.183) (-1.050,9.171) (-1.100,9.160) (-1.150,9.147) (-1.200,9.134) (-1.250,9.120) (-1.300,9.106) (-1.350,9.091) (-1.400,9.075) (-1.450,9.059) (-1.500,9.042) (-1.548,9.025) (-1.600,9.006) (-1.650,8.987) (-1.700,8.968) (-1.750,8.947) (-1.800,8.926) (-1.850,8.905) (-1.900,8.882) (-1.950,8.859) (-2.000,8.835) (-2.050,8.811) (-2.100,8.785) (-2.150,8.759) (-2.200,8.732) (-2.250,8.704) (-2.300,8.675) (-2.350,8.646) (-2.400,8.615) (-2.450,8.583) (-2.500,8.551) (-2.550,8.518) (-2.600,8.483) (-2.647,8.450) (-2.700,8.411) (-2.748,8.375) (-2.800,8.335) (-2.843,8.300) (-2.900,8.253) (-2.950,8.211) (-2.991,8.175) (-3.046,8.125) (-3.099,8.075) (-3.150,8.026) (-3.200,7.976) (-3.249,7.925) (-3.296,7.875) (-3.341,7.825) (-3.375,7.786) (-3.425,7.727) (-3.468,7.675) (-3.500,7.634) (-3.545,7.575) (-3.583,7.525) (-3.619,7.475) (-3.650,7.430) (-3.687,7.375) (-3.720,7.325) (-3.750,7.277) (-3.782,7.225) (-3.811,7.175) (-3.840,7.125) (-3.867,7.075) (-3.894,7.025) (-3.920,6.975) (-3.944,6.925) (-3.968,6.875) (-3.991,6.825) (-4.013,6.775) (-4.035,6.725) (-4.055,6.675) (-4.075,6.625) (-4.093,6.575) (-4.111,6.525) (-4.129,6.475) (-4.145,6.425) (-4.161,6.375) (-4.176,6.325) (-4.190,6.275) (-4.203,6.225) (-4.216,6.175) (-4.228,6.125) (-4.239,6.075) (-4.249,6.025) (-4.259,5.975) (-4.268,5.925) (-4.276,5.875) (-4.284,5.825) (-4.290,5.775) (-4.296,5.725) (-4.302,5.675) (-4.307,5.625) (-4.311,5.575) (-4.314,5.525) (-4.317,5.475) (-4.318,5.425) (-4.320,5.375) (-4.320,5.325) (-4.320,5.275) (-4.319,5.225) (-4.318,5.175) (-4.315,5.125) (-4.312,5.075) (-4.309,5.025) (-4.304,4.975) (-4.299,4.925) (-4.294,4.875) (-4.287,4.825) (-4.280,4.775) (-4.272,4.725) (-4.263,4.675) (-4.254,4.625) (-4.244,4.575) (-4.233,4.525) (-4.222,4.475) (-4.210,4.425) (-4.197,4.375) (-4.183,4.325) (-4.169,4.275) (-4.153,4.225) (-4.138,4.175) (-4.121,4.125) (-4.103,4.075) (-4.085,4.025) (-4.066,3.975) (-4.046,3.925) (-4.026,3.875) (-4.004,3.825) (-3.982,3.775) (-3.959,3.725) (-3.935,3.675) (-3.911,3.625) (-3.885,3.575) (-3.859,3.525) (-3.832,3.475) (-3.803,3.425) (-3.775,3.376) (-3.745,3.325) (-3.714,3.275) (-3.683,3.225) (-3.650,3.175) (-3.617,3.125) (-3.582,3.075) (-3.550,3.029) (-3.511,2.975) (-3.475,2.926) (-3.436,2.875) (-3.400,2.829) (-3.357,2.775) (-3.325,2.736) (-3.275,2.676) (-3.232,2.625) (-3.188,2.575) (-3.150,2.533) (-3.100,2.478) (-3.050,2.425) (-3.003,2.375) (-2.954,2.325) (-2.904,2.275) (-2.854,2.225) (-2.802,2.175) (-2.750,2.126) (-2.700,2.079) (-2.650,2.033) (-2.600,1.988) (-2.558,1.950) (-2.501,1.900) (-2.450,1.856) (-2.400,1.814) (-2.354,1.775) (-2.300,1.730) (-2.250,1.689) (-2.202,1.650) (-2.150,1.608) (-2.109,1.575) (-2.050,1.529) (-2.000,1.490) (-1.950,1.452) (-1.900,1.413) (-1.850,1.375) (-1.800,1.338) (-1.750,1.301) (-1.700,1.264) (-1.650,1.227) (-1.600,1.191) (-1.550,1.155) (-1.509,1.125) (-1.450,1.083) (-1.403,1.050) (-1.350,1.013) (-1.300,0.978) (-1.250,0.944) (-1.200,0.910) (-1.150,0.876) (-1.100,0.842) (-1.050,0.809) (-1.000,0.776) (-0.950,0.743) (-0.900,0.711) (-0.850,0.679) (-0.805,0.650) (-0.750,0.616) (-0.700,0.586) (-0.650,0.556) (-0.600,0.527) (-0.553,0.500) (-0.508,0.475) (-0.461,0.450) (-0.411,0.425) (-0.357,0.400) (-0.300,0.376) (-0.250,0.357) (-0.200,0.341) (-0.150,0.327) (-0.125,0.322)}
\EQ{(-0.175,-7.375) (-0.125,-7.377) (-0.075,-7.379) (-0.025,-7.379) (0.050,-7.379) (0.100,-7.378) (0.150,-7.376) (0.200,-7.374) (0.250,-7.371) (0.300,-7.367) (0.350,-7.363) (0.400,-7.358) (0.450,-7.352) (0.500,-7.345) (0.550,-7.338) (0.600,-7.330) (0.650,-7.321) (0.700,-7.312) (0.750,-7.302) (0.800,-7.291) (0.850,-7.280) (0.900,-7.267) (0.950,-7.254) (1.000,-7.240) (1.050,-7.226) (1.100,-7.210) (1.150,-7.194) (1.200,-7.177) (1.250,-7.159) (1.300,-7.141) (1.350,-7.121) (1.400,-7.101) (1.450,-7.080) (1.500,-7.057) (1.550,-7.034) (1.600,-7.010) (1.650,-6.985) (1.700,-6.959) (1.750,-6.932) (1.800,-6.904) (1.850,-6.875) (1.900,-6.844) (1.950,-6.813) (2.000,-6.780) (2.044,-6.750) (2.100,-6.711) (2.149,-6.675) (2.200,-6.636) (2.246,-6.600) (2.300,-6.555) (2.350,-6.513) (2.392,-6.475) (2.446,-6.425) (2.498,-6.375) (2.548,-6.325) (2.596,-6.275) (2.641,-6.225) (2.675,-6.187) (2.725,-6.127) (2.767,-6.075) (2.800,-6.033) (2.843,-5.975) (2.875,-5.931) (2.913,-5.875) (2.946,-5.825) (2.975,-5.780) (3.008,-5.725) (3.037,-5.675) (3.065,-5.625) (3.091,-5.575) (3.117,-5.525) (3.141,-5.475) (3.164,-5.425) (3.186,-5.375) (3.207,-5.325) (3.227,-5.275) (3.246,-5.225) (3.264,-5.175) (3.281,-5.125) (3.297,-5.075) (3.312,-5.025) (3.326,-4.975) (3.339,-4.925) (3.351,-4.875) (3.362,-4.825) (3.372,-4.775) (3.382,-4.725) (3.390,-4.675) (3.398,-4.625) (3.404,-4.575) (3.410,-4.525) (3.415,-4.475) (3.419,-4.425) (3.422,-4.375) (3.424,-4.325) (3.426,-4.275) (3.426,-4.225) (3.426,-4.175) (3.425,-4.125) (3.422,-4.075) (3.419,-4.025) (3.415,-3.975) (3.411,-3.925) (3.405,-3.875) (3.398,-3.825) (3.391,-3.775) (3.383,-3.725) (3.373,-3.675) (3.363,-3.625) (3.352,-3.575) (3.340,-3.525) (3.327,-3.475) (3.313,-3.425) (3.298,-3.375) (3.283,-3.325) (3.266,-3.275) (3.248,-3.225) (3.230,-3.175) (3.210,-3.125) (3.189,-3.075) (3.168,-3.025) (3.145,-2.975) (3.121,-2.925) (3.096,-2.875) (3.070,-2.825) (3.044,-2.775) (3.016,-2.725) (2.986,-2.675) (2.956,-2.625) (2.925,-2.575) (2.892,-2.525) (2.859,-2.475) (2.825,-2.427) (2.788,-2.375) (2.751,-2.325) (2.713,-2.275) (2.675,-2.228) (2.632,-2.175) (2.600,-2.137) (2.550,-2.079) (2.502,-2.025) (2.457,-1.975) (2.410,-1.925) (2.375,-1.889) (2.337,-1.850) (2.300,-1.813) (2.261,-1.775) (2.225,-1.740) (2.183,-1.700) (2.129,-1.650) (2.075,-1.601) (2.025,-1.557) (1.975,-1.514) (1.930,-1.475) (1.875,-1.429) (1.825,-1.387) (1.779,-1.350) (1.725,-1.307) (1.686,-1.275) (1.625,-1.228) (1.575,-1.190) (1.525,-1.152) (1.475,-1.114) (1.425,-1.077) (1.375,-1.040) (1.325,-1.004) (1.275,-0.968) (1.225,-0.933) (1.178,-0.900) (1.125,-0.863) (1.075,-0.829) (1.033,-0.800) (0.975,-0.762) (0.925,-0.729) (0.881,-0.700) (0.825,-0.664) (0.775,-0.633) (0.725,-0.602) (0.682,-0.575) (0.625,-0.542) (0.575,-0.513) (0.525,-0.485) (0.475,-0.458) (0.425,-0.433) (0.375,-0.409) (0.325,-0.386) (0.275,-0.366) (0.225,-0.348) (0.175,-0.333) (0.125,-0.322) (0.075,-0.314) (0.025,-0.309) (-0.050,-0.311) (-0.100,-0.317) (-0.150,-0.327) (-0.200,-0.341) (-0.250,-0.357) (-0.298,-0.375) (-0.350,-0.397) (-0.400,-0.420) (-0.450,-0.445) (-0.500,-0.471) (-0.550,-0.499) (-0.597,-0.525) (-0.650,-0.556) (-0.700,-0.587) (-0.750,-0.617) (-0.800,-0.648) (-0.843,-0.675) (-0.900,-0.713) (-0.950,-0.745) (-0.995,-0.775) (-1.050,-0.812) (-1.100,-0.846) (-1.142,-0.875) (-1.200,-0.915) (-1.249,-0.950) (-1.300,-0.986) (-1.350,-1.022) (-1.400,-1.059) (-1.450,-1.095) (-1.500,-1.133) (-1.550,-1.170) (-1.600,-1.209) (-1.650,-1.247) (-1.700,-1.287) (-1.748,-1.325) (-1.800,-1.367) (-1.841,-1.400) (-1.900,-1.450) (-1.950,-1.492) (-1.989,-1.525) (-2.045,-1.575) (-2.100,-1.624) (-2.150,-1.670) (-2.200,-1.717) (-2.250,-1.764) (-2.287,-1.800) (-2.325,-1.838) (-2.362,-1.875) (-2.400,-1.915) (-2.450,-1.968) (-2.500,-2.022) (-2.547,-2.075) (-2.590,-2.125) (-2.625,-2.166) (-2.673,-2.225) (-2.713,-2.275) (-2.750,-2.324) (-2.788,-2.375) (-2.824,-2.425) (-2.859,-2.475) (-2.892,-2.525) (-2.925,-2.575) (-2.956,-2.625) (-2.986,-2.675) (-3.016,-2.725) (-3.044,-2.775) (-3.070,-2.825) (-3.096,-2.875) (-3.121,-2.925) (-3.145,-2.975) (-3.168,-3.025) (-3.189,-3.075) (-3.210,-3.125) (-3.230,-3.175) (-3.248,-3.225) (-3.266,-3.275) (-3.283,-3.325) (-3.298,-3.375) (-3.313,-3.425) (-3.327,-3.475) (-3.340,-3.525) (-3.352,-3.575) (-3.363,-3.625) (-3.373,-3.675) (-3.383,-3.725) (-3.391,-3.775) (-3.398,-3.825) (-3.405,-3.875) (-3.411,-3.925) (-3.415,-3.975) (-3.419,-4.025) (-3.422,-4.075) (-3.425,-4.125) (-3.426,-4.175) (-3.426,-4.225) (-3.426,-4.275) (-3.424,-4.325) (-3.422,-4.375) (-3.419,-4.425) (-3.415,-4.475) (-3.410,-4.525) (-3.404,-4.575) (-3.398,-4.625) (-3.390,-4.675) (-3.382,-4.725) (-3.372,-4.775) (-3.362,-4.825) (-3.351,-4.875) (-3.339,-4.925) (-3.326,-4.975) (-3.312,-5.025) (-3.297,-5.075) (-3.281,-5.125) (-3.264,-5.175) (-3.246,-5.225) (-3.227,-5.275) (-3.207,-5.325) (-3.186,-5.375) (-3.164,-5.425) (-3.141,-5.475) (-3.117,-5.525) (-3.091,-5.575) (-3.065,-5.625) (-3.037,-5.675) (-3.008,-5.725) (-2.978,-5.775) (-2.950,-5.819) (-2.913,-5.875) (-2.879,-5.925) (-2.850,-5.966) (-2.806,-6.025) (-2.775,-6.065) (-2.727,-6.125) (-2.685,-6.175) (-2.650,-6.215) (-2.600,-6.270) (-2.550,-6.323) (-2.500,-6.373) (-2.450,-6.422) (-2.400,-6.468) (-2.350,-6.513) (-2.306,-6.550) (-2.250,-6.596) (-2.200,-6.636) (-2.150,-6.674) (-2.100,-6.711) (-2.050,-6.746) (-2.008,-6.775) (-1.950,-6.813) (-1.900,-6.844) (-1.850,-6.875) (-1.807,-6.900) (-1.763,-6.925) (-1.717,-6.950) (-1.670,-6.975) (-1.625,-6.998) (-1.575,-7.022) (-1.525,-7.046) (-1.475,-7.069) (-1.425,-7.090) (-1.375,-7.111) (-1.325,-7.131) (-1.275,-7.150) (-1.225,-7.168) (-1.175,-7.186) (-1.125,-7.202) (-1.075,-7.218) (-1.025,-7.233) (-0.975,-7.247) (-0.925,-7.261) (-0.875,-7.274) (-0.825,-7.285) (-0.775,-7.297) (-0.725,-7.307) (-0.675,-7.317) (-0.625,-7.326) (-0.575,-7.334) (-0.525,-7.342) (-0.475,-7.349) (-0.425,-7.355) (-0.375,-7.360) (-0.325,-7.365) (-0.275,-7.369) (-0.225,-7.372) (-0.175,-7.375)}
\EQ{(-0.125,0.322) (-0.075,0.314) (-0.025,0.309) (0.050,0.311) (0.100,0.317) (0.150,0.327) (0.200,0.341) (0.250,0.357) (0.298,0.375) (0.350,0.397) (0.400,0.420) (0.450,0.445) (0.500,0.471) (0.550,0.499) (0.597,0.525) (0.650,0.556) (0.700,0.587) (0.750,0.617) (0.800,0.648) (0.843,0.675) (0.900,0.713) (0.950,0.745) (0.995,0.775) (1.050,0.812) (1.100,0.846) (1.142,0.875) (1.200,0.915) (1.249,0.950) (1.300,0.986) (1.350,1.022) (1.400,1.059) (1.450,1.095) (1.500,1.133) (1.550,1.170) (1.600,1.209) (1.650,1.247) (1.700,1.287) (1.748,1.325) (1.800,1.367) (1.841,1.400) (1.900,1.450) (1.950,1.492) (1.989,1.525) (2.045,1.575) (2.100,1.624) (2.150,1.670) (2.200,1.717) (2.250,1.764) (2.287,1.800) (2.325,1.838) (2.362,1.875) (2.400,1.915) (2.450,1.968) (2.500,2.022) (2.547,2.075) (2.590,2.125) (2.625,2.166) (2.673,2.225) (2.713,2.275) (2.750,2.324) (2.788,2.375) (2.824,2.425) (2.859,2.475) (2.892,2.525) (2.925,2.575) (2.956,2.625) (2.986,2.675) (3.016,2.725) (3.044,2.775) (3.070,2.825) (3.096,2.875) (3.121,2.925) (3.145,2.975) (3.168,3.025) (3.189,3.075) (3.210,3.125) (3.230,3.175) (3.248,3.225) (3.266,3.275) (3.283,3.325) (3.298,3.375) (3.313,3.425) (3.327,3.475) (3.340,3.525) (3.352,3.575) (3.363,3.625) (3.373,3.675) (3.383,3.725) (3.391,3.775) (3.398,3.825) (3.405,3.875) (3.411,3.925) (3.415,3.975) (3.419,4.025) (3.422,4.075) (3.425,4.125) (3.426,4.175) (3.426,4.225) (3.426,4.275) (3.424,4.325) (3.422,4.375) (3.419,4.425) (3.415,4.475) (3.410,4.525) (3.404,4.575) (3.398,4.625) (3.390,4.675) (3.382,4.725) (3.372,4.775) (3.362,4.825) (3.351,4.875) (3.339,4.925) (3.326,4.975) (3.312,5.025) (3.297,5.075) (3.281,5.125) (3.264,5.175) (3.246,5.225) (3.227,5.275) (3.207,5.325) (3.186,5.375) (3.164,5.425) (3.141,5.475) (3.117,5.525) (3.091,5.575) (3.065,5.625) (3.037,5.675) (3.008,5.725) (2.978,5.775) (2.950,5.819) (2.913,5.875) (2.879,5.925) (2.850,5.966) (2.806,6.025) (2.775,6.065) (2.727,6.125) (2.685,6.175) (2.650,6.215) (2.600,6.270) (2.550,6.323) (2.500,6.373) (2.450,6.422) (2.400,6.468) (2.350,6.513) (2.306,6.550) (2.250,6.596) (2.200,6.636) (2.150,6.674) (2.100,6.711) (2.050,6.746) (2.008,6.775) (1.950,6.813) (1.900,6.844) (1.850,6.875) (1.807,6.900) (1.763,6.925) (1.717,6.950) (1.670,6.975) (1.625,6.998) (1.575,7.022) (1.525,7.046) (1.475,7.069) (1.425,7.090) (1.375,7.111) (1.325,7.131) (1.275,7.150) (1.225,7.168) (1.175,7.186) (1.125,7.202) (1.075,7.218) (1.025,7.233) (0.975,7.247) (0.925,7.261) (0.875,7.274) (0.825,7.285) (0.775,7.297) (0.725,7.307) (0.675,7.317) (0.625,7.326) (0.575,7.334) (0.525,7.342) (0.475,7.349) (0.425,7.355) (0.375,7.360) (0.325,7.365) (0.275,7.369) (0.225,7.372) (0.175,7.375) (0.125,7.377) (0.075,7.379) (0.025,7.379) (-0.050,7.379) (-0.100,7.378) (-0.150,7.376) (-0.200,7.374) (-0.250,7.371) (-0.300,7.367) (-0.350,7.363) (-0.400,7.358) (-0.450,7.352) (-0.500,7.345) (-0.550,7.338) (-0.600,7.330) (-0.650,7.321) (-0.700,7.312) (-0.750,7.302) (-0.800,7.291) (-0.850,7.280) (-0.900,7.267) (-0.950,7.254) (-1.000,7.240) (-1.050,7.226) (-1.100,7.210) (-1.150,7.194) (-1.200,7.177) (-1.250,7.159) (-1.300,7.141) (-1.350,7.121) (-1.400,7.101) (-1.450,7.080) (-1.500,7.057) (-1.550,7.034) (-1.600,7.010) (-1.650,6.985) (-1.700,6.959) (-1.750,6.932) (-1.800,6.904) (-1.850,6.875) (-1.900,6.844) (-1.950,6.813) (-2.000,6.780) (-2.044,6.750) (-2.100,6.711) (-2.149,6.675) (-2.200,6.636) (-2.246,6.600) (-2.300,6.555) (-2.350,6.513) (-2.392,6.475) (-2.446,6.425) (-2.498,6.375) (-2.548,6.325) (-2.596,6.275) (-2.641,6.225) (-2.675,6.187) (-2.725,6.127) (-2.767,6.075) (-2.800,6.033) (-2.843,5.975) (-2.875,5.931) (-2.913,5.875) (-2.946,5.825) (-2.975,5.780) (-3.008,5.725) (-3.037,5.675) (-3.065,5.625) (-3.091,5.575) (-3.117,5.525) (-3.141,5.475) (-3.164,5.425) (-3.186,5.375) (-3.207,5.325) (-3.227,5.275) (-3.246,5.225) (-3.264,5.175) (-3.281,5.125) (-3.297,5.075) (-3.312,5.025) (-3.326,4.975) (-3.339,4.925) (-3.351,4.875) (-3.362,4.825) (-3.372,4.775) (-3.382,4.725) (-3.390,4.675) (-3.398,4.625) (-3.404,4.575) (-3.410,4.525) (-3.415,4.475) (-3.419,4.425) (-3.422,4.375) (-3.424,4.325) (-3.426,4.275) (-3.426,4.225) (-3.426,4.175) (-3.425,4.125) (-3.422,4.075) (-3.419,4.025) (-3.415,3.975) (-3.411,3.925) (-3.405,3.875) (-3.398,3.825) (-3.391,3.775) (-3.383,3.725) (-3.373,3.675) (-3.363,3.625) (-3.352,3.575) (-3.340,3.525) (-3.327,3.475) (-3.313,3.425) (-3.298,3.375) (-3.283,3.325) (-3.266,3.275) (-3.248,3.225) (-3.230,3.175) (-3.210,3.125) (-3.189,3.075) (-3.168,3.025) (-3.145,2.975) (-3.121,2.925) (-3.096,2.875) (-3.070,2.825) (-3.044,2.775) (-3.016,2.725) (-2.986,2.675) (-2.956,2.625) (-2.925,2.575) (-2.892,2.525) (-2.859,2.475) (-2.825,2.427) (-2.788,2.375) (-2.751,2.325) (-2.713,2.275) (-2.675,2.228) (-2.632,2.175) (-2.600,2.137) (-2.550,2.079) (-2.502,2.025) (-2.457,1.975) (-2.410,1.925) (-2.375,1.889) (-2.337,1.850) (-2.300,1.813) (-2.261,1.775) (-2.225,1.740) (-2.183,1.700) (-2.129,1.650) (-2.075,1.601) (-2.025,1.557) (-1.975,1.514) (-1.930,1.475) (-1.875,1.429) (-1.825,1.387) (-1.779,1.350) (-1.725,1.307) (-1.686,1.275) (-1.625,1.228) (-1.575,1.190) (-1.525,1.152) (-1.475,1.114) (-1.425,1.077) (-1.375,1.040) (-1.325,1.004) (-1.275,0.968) (-1.225,0.933) (-1.178,0.900) (-1.125,0.863) (-1.075,0.829) (-1.033,0.800) (-0.975,0.762) (-0.925,0.729) (-0.881,0.700) (-0.825,0.664) (-0.775,0.633) (-0.725,0.602) (-0.682,0.575) (-0.625,0.542) (-0.575,0.513) (-0.525,0.485) (-0.475,0.458) (-0.425,0.433) (-0.375,0.409) (-0.325,0.386) (-0.275,0.366) (-0.225,0.348) (-0.175,0.333) (-0.125,0.322)}
\EQ{(-0.375,-6.428) (-0.325,-6.433) (-0.275,-6.438) (-0.225,-6.442) (-0.175,-6.445) (-0.125,-6.447) (-0.075,-6.449) (-0.025,-6.449) (0.050,-6.449) (0.100,-6.448) (0.150,-6.446) (0.200,-6.443) (0.250,-6.440) (0.300,-6.436) (0.350,-6.430) (0.400,-6.425) (0.450,-6.418) (0.500,-6.410) (0.550,-6.402) (0.600,-6.393) (0.650,-6.383) (0.700,-6.372) (0.750,-6.361) (0.800,-6.348) (0.850,-6.335) (0.900,-6.320) (0.950,-6.305) (1.000,-6.289) (1.050,-6.272) (1.100,-6.254) (1.150,-6.235) (1.200,-6.216) (1.250,-6.195) (1.300,-6.173) (1.349,-6.150) (1.400,-6.126) (1.450,-6.101) (1.498,-6.075) (1.544,-6.050) (1.588,-6.025) (1.650,-5.988) (1.700,-5.956) (1.748,-5.925) (1.800,-5.889) (1.850,-5.854) (1.900,-5.816) (1.950,-5.777) (2.000,-5.737) (2.043,-5.700) (2.099,-5.650) (2.150,-5.603) (2.200,-5.554) (2.250,-5.502) (2.298,-5.450) (2.342,-5.400) (2.375,-5.361) (2.424,-5.300) (2.463,-5.250) (2.499,-5.200) (2.534,-5.150) (2.567,-5.100) (2.599,-5.050) (2.625,-5.007) (2.658,-4.950) (2.685,-4.900) (2.711,-4.850) (2.736,-4.800) (2.759,-4.750) (2.782,-4.700) (2.802,-4.650) (2.822,-4.600) (2.841,-4.550) (2.858,-4.500) (2.874,-4.450) (2.890,-4.400) (2.904,-4.350) (2.917,-4.300) (2.929,-4.250) (2.939,-4.200) (2.949,-4.150) (2.958,-4.100) (2.966,-4.050) (2.972,-4.000) (2.978,-3.950) (2.983,-3.900) (2.986,-3.850) (2.989,-3.800) (2.991,-3.750) (2.991,-3.700) (2.991,-3.650) (2.989,-3.600) (2.987,-3.550) (2.984,-3.500) (2.979,-3.450) (2.974,-3.400) (2.967,-3.350) (2.960,-3.300) (2.951,-3.250) (2.942,-3.200) (2.931,-3.150) (2.920,-3.100) (2.907,-3.050) (2.893,-3.000) (2.878,-2.950) (2.862,-2.900) (2.845,-2.850) (2.827,-2.800) (2.808,-2.750) (2.788,-2.700) (2.766,-2.650) (2.743,-2.600) (2.719,-2.550) (2.694,-2.500) (2.668,-2.450) (2.640,-2.400) (2.611,-2.350) (2.581,-2.300) (2.550,-2.251) (2.517,-2.200) (2.483,-2.150) (2.450,-2.104) (2.411,-2.050) (2.375,-2.003) (2.333,-1.950) (2.300,-1.910) (2.250,-1.850) (2.206,-1.800) (2.161,-1.750) (2.125,-1.712) (2.090,-1.675) (2.050,-1.634) (2.000,-1.584) (1.950,-1.536) (1.913,-1.500) (1.875,-1.465) (1.831,-1.425) (1.775,-1.375) (1.725,-1.332) (1.675,-1.290) (1.628,-1.250) (1.575,-1.207) (1.536,-1.175) (1.475,-1.127) (1.425,-1.089) (1.375,-1.050) (1.325,-1.013) (1.275,-0.976) (1.225,-0.939) (1.175,-0.903) (1.125,-0.868) (1.075,-0.833) (1.028,-0.800) (0.975,-0.765) (0.925,-0.731) (0.878,-0.700) (0.825,-0.666) (0.775,-0.634) (0.725,-0.602) (0.681,-0.575) (0.625,-0.542) (0.575,-0.513) (0.525,-0.485) (0.475,-0.459) (0.425,-0.433) (0.375,-0.409) (0.325,-0.387) (0.275,-0.366) (0.225,-0.348) (0.175,-0.333) (0.125,-0.322) (0.075,-0.314) (0.025,-0.309) (-0.050,-0.311) (-0.100,-0.317) (-0.150,-0.327) (-0.200,-0.341) (-0.250,-0.357) (-0.298,-0.375) (-0.350,-0.397) (-0.400,-0.421) (-0.450,-0.445) (-0.500,-0.472) (-0.550,-0.499) (-0.596,-0.525) (-0.650,-0.557) (-0.700,-0.587) (-0.750,-0.618) (-0.800,-0.649) (-0.850,-0.682) (-0.900,-0.715) (-0.950,-0.748) (-1.000,-0.781) (-1.050,-0.816) (-1.100,-0.850) (-1.150,-0.886) (-1.200,-0.921) (-1.250,-0.958) (-1.300,-0.994) (-1.342,-1.025) (-1.400,-1.069) (-1.440,-1.100) (-1.500,-1.147) (-1.550,-1.187) (-1.597,-1.225) (-1.650,-1.269) (-1.700,-1.311) (-1.746,-1.350) (-1.800,-1.397) (-1.850,-1.442) (-1.900,-1.489) (-1.939,-1.525) (-1.991,-1.575) (-2.041,-1.625) (-2.090,-1.675) (-2.125,-1.712) (-2.161,-1.750) (-2.200,-1.794) (-2.250,-1.850) (-2.292,-1.900) (-2.325,-1.940) (-2.373,-2.000) (-2.411,-2.050) (-2.448,-2.100) (-2.483,-2.150) (-2.517,-2.200) (-2.550,-2.250) (-2.581,-2.300) (-2.611,-2.350) (-2.640,-2.400) (-2.668,-2.450) (-2.694,-2.500) (-2.719,-2.550) (-2.743,-2.600) (-2.766,-2.650) (-2.788,-2.700) (-2.808,-2.750) (-2.827,-2.800) (-2.845,-2.850) (-2.862,-2.900) (-2.878,-2.950) (-2.893,-3.000) (-2.907,-3.050) (-2.920,-3.100) (-2.931,-3.150) (-2.942,-3.200) (-2.951,-3.250) (-2.960,-3.300) (-2.967,-3.350) (-2.974,-3.400) (-2.979,-3.450) (-2.984,-3.500) (-2.987,-3.550) (-2.989,-3.600) (-2.991,-3.650) (-2.991,-3.700) (-2.991,-3.750) (-2.989,-3.800) (-2.986,-3.850) (-2.983,-3.900) (-2.978,-3.950) (-2.972,-4.000) (-2.966,-4.050) (-2.958,-4.100) (-2.949,-4.150) (-2.939,-4.200) (-2.929,-4.250) (-2.917,-4.300) (-2.904,-4.350) (-2.890,-4.400) (-2.875,-4.448) (-2.858,-4.500) (-2.841,-4.550) (-2.822,-4.600) (-2.802,-4.650) (-2.782,-4.700) (-2.759,-4.750) (-2.736,-4.800) (-2.711,-4.850) (-2.685,-4.900) (-2.658,-4.950) (-2.629,-5.000) (-2.600,-5.048) (-2.567,-5.100) (-2.534,-5.150) (-2.500,-5.199) (-2.463,-5.250) (-2.425,-5.299) (-2.384,-5.350) (-2.350,-5.391) (-2.300,-5.448) (-2.252,-5.500) (-2.204,-5.550) (-2.153,-5.600) (-2.100,-5.649) (-2.050,-5.694) (-2.000,-5.737) (-1.953,-5.775) (-1.900,-5.816) (-1.855,-5.850) (-1.800,-5.889) (-1.750,-5.924) (-1.700,-5.956) (-1.650,-5.988) (-1.600,-6.018) (-1.550,-6.047) (-1.500,-6.074) (-1.451,-6.100) (-1.401,-6.125) (-1.350,-6.150) (-1.300,-6.173) (-1.250,-6.195) (-1.200,-6.216) (-1.150,-6.235) (-1.100,-6.254) (-1.050,-6.272) (-1.000,-6.289) (-0.950,-6.305) (-0.900,-6.320) (-0.850,-6.335) (-0.800,-6.348) (-0.750,-6.361) (-0.700,-6.372) (-0.650,-6.383) (-0.600,-6.393) (-0.550,-6.402) (-0.500,-6.410) (-0.450,-6.418) (-0.400,-6.425) (-0.375,-6.428)}
\EQ{(-0.125,0.322) (-0.075,0.314) (-0.025,0.309) (0.050,0.311) (0.100,0.317) (0.150,0.327) (0.200,0.341) (0.250,0.357) (0.298,0.375) (0.350,0.397) (0.400,0.421) (0.450,0.445) (0.500,0.472) (0.550,0.499) (0.596,0.525) (0.650,0.557) (0.700,0.587) (0.750,0.618) (0.800,0.649) (0.850,0.682) (0.900,0.715) (0.950,0.748) (1.000,0.781) (1.050,0.816) (1.100,0.850) (1.150,0.886) (1.200,0.921) (1.250,0.958) (1.300,0.994) (1.342,1.025) (1.400,1.069) (1.440,1.100) (1.500,1.147) (1.550,1.187) (1.597,1.225) (1.650,1.269) (1.700,1.311) (1.746,1.350) (1.800,1.397) (1.850,1.442) (1.900,1.489) (1.939,1.525) (1.991,1.575) (2.041,1.625) (2.090,1.675) (2.125,1.712) (2.161,1.750) (2.200,1.794) (2.250,1.850) (2.292,1.900) (2.325,1.940) (2.373,2.000) (2.411,2.050) (2.448,2.100) (2.483,2.150) (2.517,2.200) (2.550,2.250) (2.581,2.300) (2.611,2.350) (2.640,2.400) (2.668,2.450) (2.694,2.500) (2.719,2.550) (2.743,2.600) (2.766,2.650) (2.788,2.700) (2.808,2.750) (2.827,2.800) (2.845,2.850) (2.862,2.900) (2.878,2.950) (2.893,3.000) (2.907,3.050) (2.920,3.100) (2.931,3.150) (2.942,3.200) (2.951,3.250) (2.960,3.300) (2.967,3.350) (2.974,3.400) (2.979,3.450) (2.984,3.500) (2.987,3.550) (2.989,3.600) (2.991,3.650) (2.991,3.700) (2.991,3.750) (2.989,3.800) (2.986,3.850) (2.983,3.900) (2.978,3.950) (2.972,4.000) (2.966,4.050) (2.958,4.100) (2.949,4.150) (2.939,4.200) (2.929,4.250) (2.917,4.300) (2.904,4.350) (2.890,4.400) (2.875,4.448) (2.858,4.500) (2.841,4.550) (2.822,4.600) (2.802,4.650) (2.782,4.700) (2.759,4.750) (2.736,4.800) (2.711,4.850) (2.685,4.900) (2.658,4.950) (2.629,5.000) (2.600,5.048) (2.567,5.100) (2.534,5.150) (2.500,5.199) (2.463,5.250) (2.425,5.299) (2.384,5.350) (2.350,5.391) (2.300,5.448) (2.252,5.500) (2.204,5.550) (2.153,5.600) (2.100,5.649) (2.050,5.694) (2.000,5.737) (1.953,5.775) (1.900,5.816) (1.855,5.850) (1.800,5.889) (1.750,5.924) (1.700,5.956) (1.650,5.988) (1.600,6.018) (1.550,6.047) (1.500,6.074) (1.451,6.100) (1.401,6.125) (1.350,6.150) (1.300,6.173) (1.250,6.195) (1.200,6.216) (1.150,6.235) (1.100,6.254) (1.050,6.272) (1.000,6.289) (0.950,6.305) (0.900,6.320) (0.850,6.335) (0.800,6.348) (0.750,6.361) (0.700,6.372) (0.650,6.383) (0.600,6.393) (0.550,6.402) (0.500,6.410) (0.450,6.418) (0.400,6.425) (0.350,6.430) (0.300,6.436) (0.250,6.440) (0.200,6.443) (0.150,6.446) (0.100,6.448) (0.050,6.449) (0.000,6.450) (-0.050,6.449) (-0.100,6.448) (-0.150,6.446) (-0.200,6.443) (-0.250,6.440) (-0.300,6.436) (-0.350,6.430) (-0.400,6.425) (-0.450,6.418) (-0.500,6.410) (-0.550,6.402) (-0.600,6.393) (-0.650,6.383) (-0.700,6.372) (-0.750,6.361) (-0.800,6.348) (-0.850,6.335) (-0.900,6.320) (-0.950,6.305) (-1.000,6.289) (-1.050,6.272) (-1.100,6.254) (-1.150,6.235) (-1.200,6.216) (-1.250,6.195) (-1.300,6.173) (-1.349,6.150) (-1.400,6.126) (-1.450,6.101) (-1.498,6.075) (-1.544,6.050) (-1.588,6.025) (-1.650,5.988) (-1.700,5.956) (-1.748,5.925) (-1.800,5.889) (-1.850,5.854) (-1.900,5.816) (-1.950,5.777) (-2.000,5.737) (-2.043,5.700) (-2.099,5.650) (-2.150,5.603) (-2.200,5.554) (-2.250,5.502) (-2.298,5.450) (-2.342,5.400) (-2.375,5.361) (-2.424,5.300) (-2.463,5.250) (-2.499,5.200) (-2.534,5.150) (-2.567,5.100) (-2.599,5.050) (-2.625,5.007) (-2.658,4.950) (-2.685,4.900) (-2.711,4.850) (-2.736,4.800) (-2.759,4.750) (-2.782,4.700) (-2.802,4.650) (-2.822,4.600) (-2.841,4.550) (-2.858,4.500) (-2.874,4.450) (-2.890,4.400) (-2.904,4.350) (-2.917,4.300) (-2.929,4.250) (-2.939,4.200) (-2.949,4.150) (-2.958,4.100) (-2.966,4.050) (-2.972,4.000) (-2.978,3.950) (-2.983,3.900) (-2.986,3.850) (-2.989,3.800) (-2.991,3.750) (-2.991,3.700) (-2.991,3.650) (-2.989,3.600) (-2.987,3.550) (-2.984,3.500) (-2.979,3.450) (-2.974,3.400) (-2.967,3.350) (-2.960,3.300) (-2.951,3.250) (-2.942,3.200) (-2.931,3.150) (-2.920,3.100) (-2.907,3.050) (-2.893,3.000) (-2.878,2.950) (-2.862,2.900) (-2.845,2.850) (-2.827,2.800) (-2.808,2.750) (-2.788,2.700) (-2.766,2.650) (-2.743,2.600) (-2.719,2.550) (-2.694,2.500) (-2.668,2.450) (-2.640,2.400) (-2.611,2.350) (-2.581,2.300) (-2.550,2.251) (-2.517,2.200) (-2.483,2.150) (-2.450,2.104) (-2.411,2.050) (-2.375,2.003) (-2.333,1.950) (-2.300,1.910) (-2.250,1.850) (-2.206,1.800) (-2.161,1.750) (-2.125,1.712) (-2.090,1.675) (-2.050,1.634) (-2.000,1.584) (-1.950,1.536) (-1.913,1.500) (-1.875,1.465) (-1.831,1.425) (-1.775,1.375) (-1.725,1.332) (-1.675,1.290) (-1.628,1.250) (-1.575,1.207) (-1.536,1.175) (-1.475,1.127) (-1.425,1.089) (-1.375,1.050) (-1.325,1.013) (-1.275,0.976) (-1.225,0.939) (-1.175,0.903) (-1.125,0.868) (-1.075,0.833) (-1.028,0.800) (-0.975,0.765) (-0.925,0.731) (-0.878,0.700) (-0.825,0.666) (-0.775,0.634) (-0.725,0.602) (-0.681,0.575) (-0.625,0.542) (-0.575,0.513) (-0.525,0.485) (-0.475,0.459) (-0.425,0.433) (-0.375,0.409) (-0.325,0.387) (-0.275,0.366) (-0.225,0.348) (-0.175,0.333) (-0.125,0.322)}
\EQ{(-0.250,-5.852) (-0.200,-5.855) (-0.150,-5.858) (-0.100,-5.861) (-0.050,-5.862) (0.000,-5.862) (0.050,-5.862) (0.100,-5.861) (0.150,-5.858) (0.200,-5.855) (0.250,-5.852) (0.300,-5.847) (0.350,-5.841) (0.400,-5.835) (0.450,-5.827) (0.500,-5.819) (0.550,-5.810) (0.600,-5.800) (0.650,-5.789) (0.700,-5.777) (0.750,-5.764) (0.800,-5.750) (0.850,-5.735) (0.900,-5.719) (0.950,-5.702) (1.000,-5.685) (1.050,-5.666) (1.100,-5.646) (1.149,-5.625) (1.200,-5.602) (1.250,-5.579) (1.300,-5.554) (1.350,-5.528) (1.400,-5.501) (1.445,-5.475) (1.500,-5.442) (1.550,-5.410) (1.600,-5.377) (1.650,-5.343) (1.700,-5.306) (1.741,-5.275) (1.800,-5.228) (1.850,-5.186) (1.891,-5.150) (1.945,-5.100) (1.997,-5.050) (2.045,-5.000) (2.091,-4.950) (2.125,-4.911) (2.175,-4.852) (2.216,-4.800) (2.250,-4.754) (2.289,-4.700) (2.323,-4.650) (2.350,-4.607) (2.385,-4.550) (2.414,-4.500) (2.441,-4.450) (2.467,-4.400) (2.491,-4.350) (2.514,-4.300) (2.536,-4.250) (2.556,-4.200) (2.575,-4.150) (2.593,-4.100) (2.609,-4.050) (2.624,-4.000) (2.638,-3.950) (2.651,-3.900) (2.663,-3.850) (2.673,-3.800) (2.682,-3.750) (2.690,-3.700) (2.697,-3.650) (2.703,-3.600) (2.708,-3.550) (2.712,-3.500) (2.714,-3.450) (2.716,-3.400) (2.716,-3.350) (2.715,-3.300) (2.713,-3.250) (2.710,-3.200) (2.706,-3.150) (2.701,-3.100) (2.695,-3.050) (2.687,-3.000) (2.679,-2.950) (2.669,-2.900) (2.658,-2.850) (2.646,-2.800) (2.633,-2.750) (2.618,-2.700) (2.603,-2.650) (2.586,-2.600) (2.568,-2.550) (2.548,-2.500) (2.528,-2.450) (2.506,-2.400) (2.483,-2.350) (2.459,-2.300) (2.433,-2.250) (2.406,-2.200) (2.377,-2.150) (2.350,-2.105) (2.316,-2.050) (2.283,-2.000) (2.250,-1.952) (2.213,-1.900) (2.175,-1.850) (2.137,-1.800) (2.100,-1.755) (2.054,-1.700) (2.011,-1.650) (1.975,-1.610) (1.925,-1.557) (1.875,-1.505) (1.825,-1.455) (1.775,-1.406) (1.725,-1.359) (1.675,-1.313) (1.633,-1.275) (1.575,-1.225) (1.525,-1.183) (1.486,-1.150) (1.425,-1.100) (1.375,-1.061) (1.330,-1.025) (1.275,-0.983) (1.231,-0.950) (1.175,-0.909) (1.129,-0.875) (1.075,-0.837) (1.025,-0.802) (0.975,-0.767) (0.925,-0.733) (0.876,-0.700) (0.825,-0.667) (0.775,-0.635) (0.725,-0.603) (0.679,-0.575) (0.625,-0.543) (0.575,-0.514) (0.525,-0.486) (0.475,-0.459) (0.425,-0.433) (0.375,-0.409) (0.325,-0.387) (0.275,-0.366) (0.225,-0.348) (0.175,-0.334) (0.125,-0.322) (0.075,-0.314) (0.025,-0.309) (-0.050,-0.311) (-0.100,-0.317) (-0.150,-0.327) (-0.200,-0.341) (-0.250,-0.357) (-0.298,-0.375) (-0.350,-0.397) (-0.400,-0.421) (-0.450,-0.446) (-0.500,-0.472) (-0.550,-0.499) (-0.595,-0.525) (-0.650,-0.558) (-0.700,-0.588) (-0.750,-0.619) (-0.799,-0.650) (-0.850,-0.683) (-0.900,-0.716) (-0.950,-0.750) (-1.000,-0.784) (-1.050,-0.819) (-1.094,-0.850) (-1.150,-0.891) (-1.197,-0.925) (-1.250,-0.964) (-1.297,-1.000) (-1.350,-1.041) (-1.394,-1.075) (-1.450,-1.121) (-1.500,-1.162) (-1.546,-1.200) (-1.600,-1.247) (-1.650,-1.291) (-1.689,-1.325) (-1.743,-1.375) (-1.795,-1.425) (-1.846,-1.475) (-1.895,-1.525) (-1.943,-1.575) (-1.989,-1.625) (-2.025,-1.666) (-2.075,-1.724) (-2.117,-1.775) (-2.150,-1.817) (-2.194,-1.875) (-2.225,-1.917) (-2.266,-1.975) (-2.299,-2.025) (-2.332,-2.075) (-2.362,-2.125) (-2.391,-2.175) (-2.419,-2.225) (-2.446,-2.275) (-2.471,-2.325) (-2.495,-2.375) (-2.517,-2.425) (-2.538,-2.475) (-2.558,-2.525) (-2.577,-2.575) (-2.594,-2.625) (-2.611,-2.675) (-2.625,-2.723) (-2.640,-2.775) (-2.652,-2.825) (-2.664,-2.875) (-2.674,-2.925) (-2.683,-2.975) (-2.691,-3.025) (-2.698,-3.075) (-2.704,-3.125) (-2.708,-3.175) (-2.712,-3.225) (-2.714,-3.275) (-2.716,-3.325) (-2.716,-3.375) (-2.715,-3.425) (-2.713,-3.475) (-2.710,-3.525) (-2.706,-3.575) (-2.700,-3.625) (-2.694,-3.675) (-2.686,-3.725) (-2.678,-3.775) (-2.668,-3.825) (-2.657,-3.875) (-2.645,-3.925) (-2.631,-3.975) (-2.617,-4.025) (-2.601,-4.075) (-2.584,-4.125) (-2.566,-4.175) (-2.546,-4.225) (-2.525,-4.275) (-2.503,-4.325) (-2.479,-4.375) (-2.454,-4.425) (-2.428,-4.475) (-2.400,-4.524) (-2.370,-4.575) (-2.339,-4.625) (-2.306,-4.675) (-2.275,-4.720) (-2.235,-4.775) (-2.200,-4.820) (-2.156,-4.875) (-2.113,-4.925) (-2.075,-4.968) (-2.025,-5.021) (-1.975,-5.071) (-1.925,-5.119) (-1.875,-5.164) (-1.834,-5.200) (-1.775,-5.249) (-1.725,-5.288) (-1.675,-5.325) (-1.625,-5.360) (-1.575,-5.394) (-1.527,-5.425) (-1.475,-5.457) (-1.425,-5.486) (-1.375,-5.514) (-1.325,-5.541) (-1.275,-5.566) (-1.225,-5.590) (-1.175,-5.613) (-1.125,-5.635) (-1.075,-5.656) (-1.026,-5.675) (-0.975,-5.694) (-0.925,-5.711) (-0.875,-5.727) (-0.825,-5.743) (-0.775,-5.757) (-0.725,-5.770) (-0.675,-5.783) (-0.625,-5.794) (-0.575,-5.805) (-0.525,-5.815) (-0.475,-5.823) (-0.425,-5.831) (-0.375,-5.838) (-0.325,-5.844) (-0.275,-5.849) (-0.250,-5.852)}
\EQ{(-0.125,0.322) (-0.075,0.314) (-0.025,0.309) (0.050,0.311) (0.100,0.317) (0.150,0.327) (0.200,0.341) (0.250,0.357) (0.298,0.375) (0.350,0.397) (0.400,0.421) (0.450,0.446) (0.500,0.472) (0.550,0.499) (0.595,0.525) (0.650,0.558) (0.700,0.588) (0.750,0.619) (0.799,0.650) (0.850,0.683) (0.900,0.716) (0.950,0.750) (1.000,0.784) (1.050,0.819) (1.094,0.850) (1.150,0.891) (1.197,0.925) (1.250,0.964) (1.297,1.000) (1.350,1.041) (1.394,1.075) (1.450,1.121) (1.500,1.162) (1.546,1.200) (1.600,1.247) (1.650,1.291) (1.689,1.325) (1.743,1.375) (1.795,1.425) (1.846,1.475) (1.895,1.525) (1.943,1.575) (1.989,1.625) (2.025,1.666) (2.075,1.724) (2.117,1.775) (2.150,1.817) (2.194,1.875) (2.225,1.917) (2.266,1.975) (2.299,2.025) (2.332,2.075) (2.362,2.125) (2.391,2.175) (2.419,2.225) (2.446,2.275) (2.471,2.325) (2.495,2.375) (2.517,2.425) (2.538,2.475) (2.558,2.525) (2.577,2.575) (2.594,2.625) (2.611,2.675) (2.625,2.723) (2.640,2.775) (2.652,2.825) (2.664,2.875) (2.674,2.925) (2.683,2.975) (2.691,3.025) (2.698,3.075) (2.704,3.125) (2.708,3.175) (2.712,3.225) (2.714,3.275) (2.716,3.325) (2.716,3.375) (2.715,3.425) (2.713,3.475) (2.710,3.525) (2.706,3.575) (2.700,3.625) (2.694,3.675) (2.686,3.725) (2.678,3.775) (2.668,3.825) (2.657,3.875) (2.645,3.925) (2.631,3.975) (2.617,4.025) (2.601,4.075) (2.584,4.125) (2.566,4.175) (2.546,4.225) (2.525,4.275) (2.503,4.325) (2.479,4.375) (2.454,4.425) (2.428,4.475) (2.400,4.524) (2.370,4.575) (2.339,4.625) (2.306,4.675) (2.275,4.720) (2.235,4.775) (2.200,4.820) (2.156,4.875) (2.113,4.925) (2.075,4.968) (2.025,5.021) (1.975,5.071) (1.925,5.119) (1.875,5.164) (1.834,5.200) (1.775,5.249) (1.725,5.288) (1.675,5.325) (1.625,5.360) (1.575,5.394) (1.527,5.425) (1.475,5.457) (1.425,5.486) (1.375,5.514) (1.325,5.541) (1.275,5.566) (1.225,5.590) (1.175,5.613) (1.125,5.635) (1.075,5.656) (1.026,5.675) (0.975,5.694) (0.925,5.711) (0.875,5.727) (0.825,5.743) (0.775,5.757) (0.725,5.770) (0.675,5.783) (0.625,5.794) (0.575,5.805) (0.525,5.815) (0.475,5.823) (0.425,5.831) (0.375,5.838) (0.325,5.844) (0.275,5.849) (0.225,5.854) (0.175,5.857) (0.125,5.860) (0.075,5.861) (0.025,5.862) (-0.050,5.862) (-0.100,5.861) (-0.150,5.858) (-0.200,5.855) (-0.250,5.852) (-0.300,5.847) (-0.350,5.841) (-0.400,5.835) (-0.450,5.827) (-0.500,5.819) (-0.550,5.810) (-0.600,5.800) (-0.650,5.789) (-0.700,5.777) (-0.750,5.764) (-0.800,5.750) (-0.850,5.735) (-0.900,5.719) (-0.950,5.702) (-1.000,5.685) (-1.050,5.666) (-1.100,5.646) (-1.149,5.625) (-1.200,5.602) (-1.250,5.579) (-1.300,5.554) (-1.350,5.528) (-1.400,5.501) (-1.445,5.475) (-1.500,5.442) (-1.550,5.410) (-1.600,5.377) (-1.650,5.343) (-1.700,5.306) (-1.741,5.275) (-1.800,5.228) (-1.850,5.186) (-1.891,5.150) (-1.945,5.100) (-1.997,5.050) (-2.045,5.000) (-2.091,4.950) (-2.125,4.911) (-2.175,4.852) (-2.216,4.800) (-2.250,4.754) (-2.289,4.700) (-2.323,4.650) (-2.350,4.607) (-2.385,4.550) (-2.414,4.500) (-2.441,4.450) (-2.467,4.400) (-2.491,4.350) (-2.514,4.300) (-2.536,4.250) (-2.556,4.200) (-2.575,4.150) (-2.593,4.100) (-2.609,4.050) (-2.624,4.000) (-2.638,3.950) (-2.651,3.900) (-2.663,3.850) (-2.673,3.800) (-2.682,3.750) (-2.690,3.700) (-2.697,3.650) (-2.703,3.600) (-2.708,3.550) (-2.712,3.500) (-2.714,3.450) (-2.716,3.400) (-2.716,3.350) (-2.715,3.300) (-2.713,3.250) (-2.710,3.200) (-2.706,3.150) (-2.701,3.100) (-2.695,3.050) (-2.687,3.000) (-2.679,2.950) (-2.669,2.900) (-2.658,2.850) (-2.646,2.800) (-2.633,2.750) (-2.618,2.700) (-2.603,2.650) (-2.586,2.600) (-2.568,2.550) (-2.548,2.500) (-2.528,2.450) (-2.506,2.400) (-2.483,2.350) (-2.459,2.300) (-2.433,2.250) (-2.406,2.200) (-2.377,2.150) (-2.350,2.105) (-2.316,2.050) (-2.283,2.000) (-2.250,1.952) (-2.213,1.900) (-2.175,1.850) (-2.137,1.800) (-2.100,1.755) (-2.054,1.700) (-2.011,1.650) (-1.975,1.610) (-1.925,1.557) (-1.875,1.505) (-1.825,1.455) (-1.775,1.406) (-1.725,1.359) (-1.675,1.313) (-1.633,1.275) (-1.575,1.225) (-1.525,1.183) (-1.486,1.150) (-1.425,1.100) (-1.375,1.061) (-1.330,1.025) (-1.275,0.983) (-1.231,0.950) (-1.175,0.909) (-1.129,0.875) (-1.075,0.837) (-1.025,0.802) (-0.975,0.767) (-0.925,0.733) (-0.876,0.700) (-0.825,0.667) (-0.775,0.635) (-0.725,0.603) (-0.679,0.575) (-0.625,0.543) (-0.575,0.514) (-0.525,0.486) (-0.475,0.459) (-0.425,0.433) (-0.375,0.409) (-0.325,0.387) (-0.275,0.366) (-0.225,0.348) (-0.175,0.334) (-0.125,0.322)}
\EQ{(-0.300,-5.427) (-0.250,-5.433) (-0.200,-5.437) (-0.150,-5.440) (-0.100,-5.442) (-0.050,-5.444) (0.000,-5.444) (0.050,-5.444) (0.100,-5.442) (0.150,-5.440) (0.200,-5.437) (0.250,-5.433) (0.300,-5.427) (0.350,-5.421) (0.400,-5.414) (0.450,-5.406) (0.500,-5.398) (0.550,-5.388) (0.600,-5.377) (0.650,-5.365) (0.700,-5.352) (0.750,-5.338) (0.800,-5.323) (0.850,-5.307) (0.900,-5.289) (0.950,-5.271) (1.000,-5.251) (1.050,-5.231) (1.100,-5.209) (1.150,-5.186) (1.200,-5.161) (1.250,-5.135) (1.300,-5.108) (1.350,-5.079) (1.398,-5.050) (1.450,-5.017) (1.500,-4.984) (1.548,-4.950) (1.600,-4.911) (1.646,-4.875) (1.700,-4.831) (1.750,-4.787) (1.791,-4.750) (1.843,-4.700) (1.892,-4.650) (1.938,-4.600) (1.975,-4.558) (2.023,-4.500) (2.062,-4.450) (2.099,-4.400) (2.134,-4.350) (2.167,-4.300) (2.199,-4.250) (2.225,-4.206) (2.250,-4.161) (2.275,-4.114) (2.300,-4.065) (2.325,-4.012) (2.350,-3.955) (2.372,-3.900) (2.391,-3.850) (2.409,-3.800) (2.425,-3.750) (2.439,-3.700) (2.453,-3.650) (2.465,-3.600) (2.476,-3.550) (2.486,-3.500) (2.494,-3.450) (2.502,-3.400) (2.508,-3.350) (2.513,-3.300) (2.516,-3.250) (2.519,-3.200) (2.520,-3.150) (2.520,-3.100) (2.519,-3.050) (2.516,-3.000) (2.513,-2.950) (2.508,-2.900) (2.502,-2.850) (2.495,-2.800) (2.487,-2.750) (2.477,-2.700) (2.466,-2.650) (2.454,-2.600) (2.441,-2.550) (2.426,-2.500) (2.410,-2.450) (2.393,-2.400) (2.375,-2.352) (2.355,-2.300) (2.333,-2.250) (2.311,-2.200) (2.286,-2.150) (2.261,-2.100) (2.234,-2.050) (2.205,-2.000) (2.175,-1.950) (2.144,-1.900) (2.111,-1.850) (2.076,-1.800) (2.040,-1.750) (2.002,-1.700) (1.962,-1.650) (1.925,-1.605) (1.877,-1.550) (1.833,-1.500) (1.786,-1.450) (1.750,-1.413) (1.713,-1.375) (1.675,-1.338) (1.636,-1.300) (1.582,-1.250) (1.527,-1.200) (1.475,-1.155) (1.425,-1.113) (1.380,-1.075) (1.325,-1.031) (1.275,-0.991) (1.225,-0.952) (1.175,-0.914) (1.125,-0.877) (1.075,-0.841) (1.025,-0.805) (0.983,-0.775) (0.925,-0.735) (0.875,-0.701) (0.825,-0.669) (0.775,-0.636) (0.725,-0.604) (0.678,-0.575) (0.625,-0.543) (0.575,-0.514) (0.525,-0.486) (0.475,-0.459) (0.425,-0.433) (0.375,-0.409) (0.325,-0.387) (0.275,-0.366) (0.225,-0.348) (0.175,-0.334) (0.125,-0.322) (0.075,-0.314) (0.025,-0.309) (-0.050,-0.311) (-0.100,-0.317) (-0.150,-0.327) (-0.200,-0.341) (-0.250,-0.357) (-0.298,-0.375) (-0.350,-0.397) (-0.400,-0.421) (-0.450,-0.446) (-0.500,-0.472) (-0.550,-0.500) (-0.595,-0.525) (-0.650,-0.558) (-0.700,-0.589) (-0.750,-0.620) (-0.797,-0.650) (-0.850,-0.685) (-0.900,-0.718) (-0.947,-0.750) (-1.000,-0.787) (-1.050,-0.823) (-1.100,-0.859) (-1.150,-0.896) (-1.200,-0.933) (-1.250,-0.971) (-1.300,-1.011) (-1.349,-1.050) (-1.400,-1.092) (-1.440,-1.125) (-1.498,-1.175) (-1.550,-1.221) (-1.600,-1.267) (-1.650,-1.314) (-1.688,-1.350) (-1.725,-1.387) (-1.762,-1.425) (-1.800,-1.465) (-1.850,-1.519) (-1.899,-1.575) (-1.942,-1.625) (-1.975,-1.667) (-2.021,-1.725) (-2.058,-1.775) (-2.093,-1.825) (-2.125,-1.872) (-2.160,-1.925) (-2.191,-1.975) (-2.220,-2.025) (-2.248,-2.075) (-2.274,-2.125) (-2.299,-2.175) (-2.322,-2.225) (-2.344,-2.275) (-2.365,-2.325) (-2.384,-2.375) (-2.402,-2.425) (-2.418,-2.475) (-2.434,-2.525) (-2.448,-2.575) (-2.460,-2.625) (-2.472,-2.675) (-2.482,-2.725) (-2.491,-2.775) (-2.499,-2.825) (-2.505,-2.875) (-2.511,-2.925) (-2.515,-2.975) (-2.518,-3.025) (-2.519,-3.075) (-2.520,-3.125) (-2.519,-3.175) (-2.518,-3.225) (-2.514,-3.275) (-2.510,-3.325) (-2.505,-3.375) (-2.498,-3.425) (-2.490,-3.475) (-2.481,-3.525) (-2.471,-3.575) (-2.459,-3.625) (-2.446,-3.675) (-2.432,-3.725) (-2.417,-3.775) (-2.400,-3.825) (-2.382,-3.875) (-2.362,-3.925) (-2.341,-3.975) (-2.319,-4.025) (-2.295,-4.075) (-2.270,-4.125) (-2.242,-4.175) (-2.214,-4.225) (-2.183,-4.275) (-2.151,-4.325) (-2.117,-4.375) (-2.081,-4.425) (-2.050,-4.466) (-2.003,-4.525) (-1.960,-4.575) (-1.925,-4.614) (-1.875,-4.668) (-1.825,-4.718) (-1.775,-4.765) (-1.736,-4.800) (-1.677,-4.850) (-1.625,-4.892) (-1.582,-4.925) (-1.525,-4.966) (-1.476,-5.000) (-1.425,-5.033) (-1.375,-5.064) (-1.325,-5.094) (-1.275,-5.122) (-1.225,-5.148) (-1.175,-5.174) (-1.125,-5.197) (-1.075,-5.220) (-1.025,-5.241) (-0.975,-5.261) (-0.925,-5.280) (-0.875,-5.298) (-0.825,-5.315) (-0.775,-5.330) (-0.725,-5.345) (-0.675,-5.358) (-0.625,-5.371) (-0.575,-5.382) (-0.525,-5.393) (-0.475,-5.402) (-0.425,-5.411) (-0.375,-5.418) (-0.325,-5.425) (-0.300,-5.427)}
\EQ{(-0.125,0.322) (-0.075,0.314) (-0.025,0.309) (0.050,0.311) (0.100,0.317) (0.150,0.327) (0.200,0.341) (0.250,0.357) (0.298,0.375) (0.350,0.397) (0.400,0.421) (0.450,0.446) (0.500,0.472) (0.550,0.500) (0.595,0.525) (0.650,0.558) (0.700,0.589) (0.750,0.620) (0.797,0.650) (0.850,0.685) (0.900,0.718) (0.947,0.750) (1.000,0.787) (1.050,0.823) (1.100,0.859) (1.150,0.896) (1.200,0.933) (1.250,0.971) (1.300,1.011) (1.349,1.050) (1.400,1.092) (1.440,1.125) (1.498,1.175) (1.550,1.221) (1.600,1.267) (1.650,1.314) (1.688,1.350) (1.725,1.387) (1.762,1.425) (1.800,1.465) (1.850,1.519) (1.899,1.575) (1.942,1.625) (1.975,1.667) (2.021,1.725) (2.058,1.775) (2.093,1.825) (2.125,1.872) (2.160,1.925) (2.191,1.975) (2.220,2.025) (2.248,2.075) (2.274,2.125) (2.299,2.175) (2.322,2.225) (2.344,2.275) (2.365,2.325) (2.384,2.375) (2.402,2.425) (2.418,2.475) (2.434,2.525) (2.448,2.575) (2.460,2.625) (2.472,2.675) (2.482,2.725) (2.491,2.775) (2.499,2.825) (2.505,2.875) (2.511,2.925) (2.515,2.975) (2.518,3.025) (2.519,3.075) (2.520,3.125) (2.519,3.175) (2.518,3.225) (2.514,3.275) (2.510,3.325) (2.505,3.375) (2.498,3.425) (2.490,3.475) (2.481,3.525) (2.471,3.575) (2.459,3.625) (2.446,3.675) (2.432,3.725) (2.417,3.775) (2.400,3.825) (2.382,3.875) (2.362,3.925) (2.341,3.975) (2.319,4.025) (2.295,4.075) (2.270,4.125) (2.242,4.175) (2.214,4.225) (2.183,4.275) (2.151,4.325) (2.117,4.375) (2.081,4.425) (2.050,4.466) (2.003,4.525) (1.960,4.575) (1.925,4.614) (1.875,4.668) (1.825,4.718) (1.775,4.765) (1.736,4.800) (1.677,4.850) (1.625,4.892) (1.582,4.925) (1.525,4.966) (1.476,5.000) (1.425,5.033) (1.375,5.064) (1.325,5.094) (1.275,5.122) (1.225,5.148) (1.175,5.174) (1.125,5.197) (1.075,5.220) (1.025,5.241) (0.975,5.261) (0.925,5.280) (0.875,5.298) (0.825,5.315) (0.775,5.330) (0.725,5.345) (0.675,5.358) (0.625,5.371) (0.575,5.382) (0.525,5.393) (0.475,5.402) (0.425,5.411) (0.375,5.418) (0.325,5.425) (0.275,5.430) (0.225,5.435) (0.175,5.439) (0.125,5.441) (0.075,5.443) (0.025,5.444) (-0.050,5.444) (-0.100,5.442) (-0.150,5.440) (-0.200,5.437) (-0.250,5.433) (-0.300,5.427) (-0.350,5.421) (-0.400,5.414) (-0.450,5.406) (-0.500,5.398) (-0.550,5.388) (-0.600,5.377) (-0.650,5.365) (-0.700,5.352) (-0.750,5.338) (-0.800,5.323) (-0.850,5.307) (-0.900,5.289) (-0.950,5.271) (-1.000,5.251) (-1.050,5.231) (-1.100,5.209) (-1.150,5.186) (-1.200,5.161) (-1.250,5.135) (-1.300,5.108) (-1.350,5.079) (-1.398,5.050) (-1.450,5.017) (-1.500,4.984) (-1.548,4.950) (-1.600,4.911) (-1.646,4.875) (-1.700,4.831) (-1.750,4.787) (-1.791,4.750) (-1.843,4.700) (-1.892,4.650) (-1.938,4.600) (-1.975,4.558) (-2.023,4.500) (-2.062,4.450) (-2.099,4.400) (-2.134,4.350) (-2.167,4.300) (-2.199,4.250) (-2.225,4.206) (-2.250,4.161) (-2.275,4.114) (-2.300,4.065) (-2.325,4.012) (-2.350,3.955) (-2.372,3.900) (-2.391,3.850) (-2.409,3.800) (-2.425,3.750) (-2.439,3.700) (-2.453,3.650) (-2.465,3.600) (-2.476,3.550) (-2.486,3.500) (-2.494,3.450) (-2.502,3.400) (-2.508,3.350) (-2.513,3.300) (-2.516,3.250) (-2.519,3.200) (-2.520,3.150) (-2.520,3.100) (-2.519,3.050) (-2.516,3.000) (-2.513,2.950) (-2.508,2.900) (-2.502,2.850) (-2.495,2.800) (-2.487,2.750) (-2.477,2.700) (-2.466,2.650) (-2.454,2.600) (-2.441,2.550) (-2.426,2.500) (-2.410,2.450) (-2.393,2.400) (-2.375,2.352) (-2.355,2.300) (-2.333,2.250) (-2.311,2.200) (-2.286,2.150) (-2.261,2.100) (-2.234,2.050) (-2.205,2.000) (-2.175,1.950) (-2.144,1.900) (-2.111,1.850) (-2.076,1.800) (-2.040,1.750) (-2.002,1.700) (-1.962,1.650) (-1.925,1.605) (-1.877,1.550) (-1.833,1.500) (-1.786,1.450) (-1.750,1.413) (-1.713,1.375) (-1.675,1.338) (-1.636,1.300) (-1.582,1.250) (-1.527,1.200) (-1.475,1.155) (-1.425,1.113) (-1.380,1.075) (-1.325,1.031) (-1.275,0.991) (-1.225,0.952) (-1.175,0.914) (-1.125,0.877) (-1.075,0.841) (-1.025,0.805) (-0.983,0.775) (-0.925,0.735) (-0.875,0.701) (-0.825,0.669) (-0.775,0.636) (-0.725,0.604) (-0.678,0.575) (-0.625,0.543) (-0.575,0.514) (-0.525,0.486) (-0.475,0.459) (-0.425,0.433) (-0.375,0.409) (-0.325,0.387) (-0.275,0.366) (-0.225,0.348) (-0.175,0.334) (-0.125,0.322)}
\EQ{(0.000,-5.125) (0.050,-5.125) (0.100,-5.123) (0.150,-5.121) (0.200,-5.117) (0.250,-5.113) (0.300,-5.107) (0.350,-5.101) (0.400,-5.093) (0.450,-5.085) (0.500,-5.075) (0.550,-5.065) (0.600,-5.053) (0.650,-5.040) (0.700,-5.027) (0.750,-5.012) (0.800,-4.995) (0.850,-4.978) (0.900,-4.960) (0.950,-4.940) (1.000,-4.919) (1.050,-4.897) (1.096,-4.875) (1.146,-4.850) (1.193,-4.825) (1.238,-4.800) (1.300,-4.764) (1.350,-4.732) (1.399,-4.700) (1.450,-4.664) (1.500,-4.628) (1.550,-4.589) (1.597,-4.550) (1.650,-4.504) (1.700,-4.458) (1.750,-4.408) (1.800,-4.356) (1.850,-4.300) (1.892,-4.250) (1.925,-4.209) (1.969,-4.150) (2.000,-4.107) (2.038,-4.050) (2.069,-4.000) (2.099,-3.950) (2.125,-3.902) (2.150,-3.854) (2.175,-3.803) (2.199,-3.750) (2.220,-3.700) (2.240,-3.650) (2.258,-3.600) (2.275,-3.550) (2.290,-3.500) (2.304,-3.450) (2.316,-3.400) (2.328,-3.350) (2.337,-3.300) (2.346,-3.250) (2.353,-3.200) (2.359,-3.150) (2.364,-3.100) (2.367,-3.050) (2.369,-3.000) (2.370,-2.950) (2.370,-2.900) (2.368,-2.850) (2.365,-2.800) (2.361,-2.750) (2.355,-2.700) (2.349,-2.650) (2.341,-2.600) (2.331,-2.550) (2.320,-2.500) (2.308,-2.450) (2.295,-2.400) (2.280,-2.350) (2.264,-2.300) (2.247,-2.250) (2.228,-2.200) (2.207,-2.150) (2.185,-2.100) (2.162,-2.050) (2.137,-2.000) (2.111,-1.950) (2.083,-1.900) (2.053,-1.850) (2.025,-1.806) (1.989,-1.750) (1.954,-1.700) (1.918,-1.650) (1.880,-1.600) (1.840,-1.550) (1.800,-1.503) (1.754,-1.450) (1.709,-1.400) (1.662,-1.350) (1.625,-1.313) (1.587,-1.275) (1.550,-1.239) (1.508,-1.200) (1.453,-1.150) (1.400,-1.103) (1.350,-1.061) (1.308,-1.025) (1.250,-0.979) (1.200,-0.939) (1.150,-0.901) (1.100,-0.863) (1.050,-0.826) (1.000,-0.790) (0.950,-0.755) (0.907,-0.725) (0.850,-0.687) (0.800,-0.653) (0.757,-0.625) (0.700,-0.590) (0.650,-0.559) (0.600,-0.529) (0.550,-0.500) (0.505,-0.475) (0.459,-0.450) (0.409,-0.425) (0.356,-0.400) (0.300,-0.376) (0.250,-0.357) (0.200,-0.341) (0.150,-0.327) (0.100,-0.317) (0.050,-0.311) (0.000,-0.309) (-0.050,-0.311) (-0.100,-0.317) (-0.150,-0.327) (-0.200,-0.341) (-0.250,-0.357) (-0.298,-0.375) (-0.350,-0.398) (-0.400,-0.421) (-0.450,-0.446) (-0.500,-0.472) (-0.550,-0.500) (-0.594,-0.525) (-0.650,-0.559) (-0.700,-0.590) (-0.750,-0.621) (-0.795,-0.650) (-0.850,-0.687) (-0.900,-0.720) (-0.943,-0.750) (-1.000,-0.790) (-1.049,-0.825) (-1.100,-0.863) (-1.149,-0.900) (-1.200,-0.939) (-1.246,-0.975) (-1.300,-1.019) (-1.350,-1.061) (-1.396,-1.100) (-1.450,-1.147) (-1.500,-1.193) (-1.550,-1.239) (-1.587,-1.275) (-1.625,-1.313) (-1.662,-1.350) (-1.700,-1.391) (-1.750,-1.445) (-1.798,-1.500) (-1.840,-1.550) (-1.875,-1.594) (-1.918,-1.650) (-1.950,-1.695) (-1.989,-1.750) (-2.022,-1.800) (-2.050,-1.846) (-2.083,-1.900) (-2.111,-1.950) (-2.137,-2.000) (-2.162,-2.050) (-2.185,-2.100) (-2.207,-2.150) (-2.228,-2.200) (-2.247,-2.250) (-2.264,-2.300) (-2.280,-2.350) (-2.295,-2.400) (-2.308,-2.450) (-2.320,-2.500) (-2.331,-2.550) (-2.341,-2.600) (-2.349,-2.650) (-2.355,-2.700) (-2.361,-2.750) (-2.365,-2.800) (-2.368,-2.850) (-2.370,-2.900) (-2.370,-2.950) (-2.369,-3.000) (-2.367,-3.050) (-2.364,-3.100) (-2.359,-3.150) (-2.353,-3.200) (-2.346,-3.250) (-2.337,-3.300) (-2.328,-3.350) (-2.316,-3.400) (-2.304,-3.450) (-2.290,-3.500) (-2.275,-3.549) (-2.258,-3.600) (-2.240,-3.650) (-2.220,-3.700) (-2.200,-3.748) (-2.176,-3.800) (-2.152,-3.850) (-2.126,-3.900) (-2.100,-3.947) (-2.069,-4.000) (-2.038,-4.050) (-2.005,-4.100) (-1.975,-4.142) (-1.932,-4.200) (-1.900,-4.241) (-1.850,-4.300) (-1.806,-4.350) (-1.758,-4.400) (-1.708,-4.450) (-1.654,-4.500) (-1.600,-4.547) (-1.550,-4.589) (-1.503,-4.625) (-1.450,-4.664) (-1.400,-4.699) (-1.350,-4.732) (-1.300,-4.764) (-1.250,-4.793) (-1.200,-4.821) (-1.150,-4.848) (-1.100,-4.873) (-1.050,-4.897) (-1.000,-4.919) (-0.950,-4.940) (-0.900,-4.960) (-0.850,-4.978) (-0.800,-4.995) (-0.750,-5.012) (-0.700,-5.027) (-0.650,-5.040) (-0.600,-5.053) (-0.550,-5.065) (-0.500,-5.075) (-0.450,-5.085) (-0.400,-5.093) (-0.350,-5.101) (-0.300,-5.107) (-0.250,-5.113) (-0.200,-5.117) (-0.150,-5.121) (-0.100,-5.123) (-0.050,-5.125) (0.000,-5.125)}
\EQ{(-0.125,0.322) (-0.075,0.314) (-0.025,0.309) (0.050,0.311) (0.100,0.317) (0.150,0.327) (0.200,0.341) (0.250,0.357) (0.298,0.375) (0.350,0.398) (0.400,0.421) (0.450,0.446) (0.500,0.472) (0.550,0.500) (0.594,0.525) (0.650,0.559) (0.700,0.590) (0.750,0.621) (0.795,0.650) (0.850,0.687) (0.900,0.720) (0.943,0.750) (1.000,0.790) (1.049,0.825) (1.100,0.863) (1.149,0.900) (1.200,0.939) (1.246,0.975) (1.300,1.019) (1.350,1.061) (1.396,1.100) (1.450,1.147) (1.500,1.193) (1.550,1.239) (1.587,1.275) (1.625,1.313) (1.662,1.350) (1.700,1.391) (1.750,1.445) (1.798,1.500) (1.840,1.550) (1.875,1.594) (1.918,1.650) (1.950,1.695) (1.989,1.750) (2.022,1.800) (2.050,1.846) (2.083,1.900) (2.111,1.950) (2.137,2.000) (2.162,2.050) (2.185,2.100) (2.207,2.150) (2.228,2.200) (2.247,2.250) (2.264,2.300) (2.280,2.350) (2.295,2.400) (2.308,2.450) (2.320,2.500) (2.331,2.550) (2.341,2.600) (2.349,2.650) (2.355,2.700) (2.361,2.750) (2.365,2.800) (2.368,2.850) (2.370,2.900) (2.370,2.950) (2.369,3.000) (2.367,3.050) (2.364,3.100) (2.359,3.150) (2.353,3.200) (2.346,3.250) (2.337,3.300) (2.328,3.350) (2.316,3.400) (2.304,3.450) (2.290,3.500) (2.275,3.549) (2.258,3.600) (2.240,3.650) (2.220,3.700) (2.200,3.748) (2.176,3.800) (2.152,3.850) (2.126,3.900) (2.100,3.947) (2.069,4.000) (2.038,4.050) (2.005,4.100) (1.975,4.142) (1.932,4.200) (1.900,4.241) (1.850,4.300) (1.806,4.350) (1.758,4.400) (1.708,4.450) (1.654,4.500) (1.600,4.547) (1.550,4.589) (1.503,4.625) (1.450,4.664) (1.400,4.699) (1.350,4.732) (1.300,4.764) (1.250,4.793) (1.200,4.821) (1.150,4.848) (1.100,4.873) (1.050,4.897) (1.000,4.919) (0.950,4.940) (0.900,4.960) (0.850,4.978) (0.800,4.995) (0.750,5.012) (0.700,5.027) (0.650,5.040) (0.600,5.053) (0.550,5.065) (0.500,5.075) (0.450,5.085) (0.400,5.093) (0.350,5.101) (0.300,5.107) (0.250,5.113) (0.200,5.117) (0.150,5.121) (0.100,5.123) (0.050,5.125) (0.000,5.125) (-0.050,5.125) (-0.100,5.123) (-0.150,5.121) (-0.200,5.117) (-0.250,5.113) (-0.300,5.107) (-0.350,5.101) (-0.400,5.093) (-0.450,5.085) (-0.500,5.075) (-0.550,5.065) (-0.600,5.053) (-0.650,5.040) (-0.700,5.027) (-0.750,5.012) (-0.800,4.995) (-0.850,4.978) (-0.900,4.960) (-0.950,4.940) (-1.000,4.919) (-1.050,4.897) (-1.096,4.875) (-1.146,4.850) (-1.193,4.825) (-1.238,4.800) (-1.300,4.764) (-1.350,4.732) (-1.399,4.700) (-1.450,4.664) (-1.500,4.628) (-1.550,4.589) (-1.597,4.550) (-1.650,4.504) (-1.700,4.458) (-1.750,4.408) (-1.800,4.356) (-1.850,4.300) (-1.892,4.250) (-1.925,4.209) (-1.969,4.150) (-2.000,4.107) (-2.038,4.050) (-2.069,4.000) (-2.099,3.950) (-2.125,3.902) (-2.150,3.854) (-2.175,3.803) (-2.199,3.750) (-2.220,3.700) (-2.240,3.650) (-2.258,3.600) (-2.275,3.550) (-2.290,3.500) (-2.304,3.450) (-2.316,3.400) (-2.328,3.350) (-2.337,3.300) (-2.346,3.250) (-2.353,3.200) (-2.359,3.150) (-2.364,3.100) (-2.367,3.050) (-2.369,3.000) (-2.370,2.950) (-2.370,2.900) (-2.368,2.850) (-2.365,2.800) (-2.361,2.750) (-2.355,2.700) (-2.349,2.650) (-2.341,2.600) (-2.331,2.550) (-2.320,2.500) (-2.308,2.450) (-2.295,2.400) (-2.280,2.350) (-2.264,2.300) (-2.247,2.250) (-2.228,2.200) (-2.207,2.150) (-2.185,2.100) (-2.162,2.050) (-2.137,2.000) (-2.111,1.950) (-2.083,1.900) (-2.053,1.850) (-2.025,1.806) (-1.989,1.750) (-1.954,1.700) (-1.918,1.650) (-1.880,1.600) (-1.840,1.550) (-1.800,1.503) (-1.754,1.450) (-1.709,1.400) (-1.662,1.350) (-1.625,1.313) (-1.587,1.275) (-1.550,1.239) (-1.508,1.200) (-1.453,1.150) (-1.400,1.103) (-1.350,1.061) (-1.308,1.025) (-1.250,0.979) (-1.200,0.939) (-1.150,0.901) (-1.100,0.863) (-1.050,0.826) (-1.000,0.790) (-0.950,0.755) (-0.907,0.725) (-0.850,0.687) (-0.800,0.653) (-0.757,0.625) (-0.700,0.590) (-0.650,0.559) (-0.600,0.529) (-0.550,0.500) (-0.505,0.475) (-0.459,0.450) (-0.409,0.425) (-0.356,0.400) (-0.300,0.376) (-0.250,0.357) (-0.200,0.341) (-0.150,0.327) (-0.125,0.322)}
\EQ{(-0.300,-4.851) (-0.250,-4.857) (-0.200,-4.862) (-0.150,-4.865) (-0.100,-4.868) (-0.050,-4.870) (0.000,-4.870) (0.050,-4.870) (0.100,-4.868) (0.150,-4.865) (0.200,-4.862) (0.250,-4.857) (0.300,-4.851) (0.350,-4.845) (0.400,-4.837) (0.450,-4.828) (0.500,-4.818) (0.550,-4.806) (0.600,-4.794) (0.650,-4.781) (0.700,-4.766) (0.750,-4.750) (0.800,-4.733) (0.850,-4.715) (0.900,-4.695) (0.948,-4.675) (1.000,-4.652) (1.050,-4.628) (1.100,-4.603) (1.150,-4.576) (1.196,-4.550) (1.250,-4.517) (1.300,-4.485) (1.350,-4.452) (1.400,-4.416) (1.450,-4.378) (1.500,-4.338) (1.544,-4.300) (1.599,-4.250) (1.650,-4.201) (1.700,-4.150) (1.745,-4.100) (1.788,-4.050) (1.825,-4.004) (1.866,-3.950) (1.900,-3.902) (1.934,-3.850) (1.966,-3.800) (1.995,-3.750) (2.022,-3.700) (2.048,-3.650) (2.072,-3.600) (2.094,-3.550) (2.115,-3.500) (2.134,-3.450) (2.151,-3.400) (2.167,-3.350) (2.182,-3.300) (2.195,-3.250) (2.207,-3.200) (2.217,-3.150) (2.226,-3.100) (2.233,-3.050) (2.239,-3.000) (2.244,-2.950) (2.248,-2.900) (2.250,-2.850) (2.250,-2.800) (2.250,-2.750) (2.248,-2.700) (2.245,-2.650) (2.240,-2.600) (2.234,-2.550) (2.227,-2.500) (2.218,-2.450) (2.208,-2.400) (2.197,-2.350) (2.184,-2.300) (2.169,-2.250) (2.154,-2.200) (2.137,-2.150) (2.118,-2.100) (2.097,-2.050) (2.076,-2.000) (2.052,-1.950) (2.027,-1.900) (2.001,-1.850) (1.975,-1.804) (1.943,-1.750) (1.911,-1.700) (1.877,-1.650) (1.842,-1.600) (1.805,-1.550) (1.766,-1.500) (1.725,-1.450) (1.682,-1.400) (1.638,-1.350) (1.600,-1.310) (1.550,-1.258) (1.500,-1.209) (1.450,-1.161) (1.411,-1.125) (1.355,-1.075) (1.300,-1.028) (1.250,-0.986) (1.206,-0.950) (1.150,-0.906) (1.110,-0.875) (1.050,-0.830) (1.010,-0.800) (0.950,-0.757) (0.904,-0.725) (0.850,-0.688) (0.800,-0.655) (0.755,-0.625) (0.700,-0.590) (0.650,-0.560) (0.600,-0.530) (0.550,-0.500) (0.505,-0.475) (0.458,-0.450) (0.409,-0.425) (0.356,-0.400) (0.300,-0.376) (0.250,-0.357) (0.200,-0.341) (0.150,-0.327) (0.100,-0.317) (0.050,-0.311) (0.000,-0.309) (-0.050,-0.311) (-0.100,-0.317) (-0.150,-0.327) (-0.200,-0.341) (-0.250,-0.357) (-0.297,-0.375) (-0.350,-0.398) (-0.400,-0.421) (-0.450,-0.446) (-0.500,-0.473) (-0.549,-0.500) (-0.593,-0.525) (-0.650,-0.560) (-0.700,-0.590) (-0.750,-0.622) (-0.793,-0.650) (-0.850,-0.688) (-0.900,-0.722) (-0.950,-0.757) (-1.000,-0.793) (-1.044,-0.825) (-1.100,-0.868) (-1.143,-0.900) (-1.200,-0.945) (-1.250,-0.986) (-1.297,-1.025) (-1.350,-1.071) (-1.400,-1.115) (-1.438,-1.150) (-1.475,-1.185) (-1.517,-1.225) (-1.567,-1.275) (-1.615,-1.325) (-1.650,-1.364) (-1.700,-1.421) (-1.746,-1.475) (-1.786,-1.525) (-1.824,-1.575) (-1.860,-1.625) (-1.894,-1.675) (-1.925,-1.722) (-1.958,-1.775) (-1.987,-1.825) (-2.014,-1.875) (-2.040,-1.925) (-2.064,-1.975) (-2.087,-2.025) (-2.108,-2.075) (-2.127,-2.125) (-2.145,-2.175) (-2.162,-2.225) (-2.177,-2.275) (-2.190,-2.325) (-2.203,-2.375) (-2.213,-2.425) (-2.223,-2.475) (-2.231,-2.525) (-2.237,-2.575) (-2.243,-2.625) (-2.246,-2.675) (-2.249,-2.725) (-2.250,-2.775) (-2.250,-2.825) (-2.249,-2.875) (-2.246,-2.925) (-2.242,-2.975) (-2.237,-3.025) (-2.230,-3.075) (-2.221,-3.125) (-2.212,-3.175) (-2.201,-3.225) (-2.189,-3.275) (-2.175,-3.324) (-2.160,-3.375) (-2.143,-3.425) (-2.125,-3.474) (-2.105,-3.525) (-2.083,-3.575) (-2.060,-3.625) (-2.035,-3.675) (-2.009,-3.725) (-1.981,-3.775) (-1.950,-3.825) (-1.918,-3.875) (-1.884,-3.925) (-1.850,-3.971) (-1.808,-4.025) (-1.775,-4.065) (-1.725,-4.122) (-1.676,-4.175) (-1.626,-4.225) (-1.575,-4.273) (-1.525,-4.317) (-1.485,-4.350) (-1.425,-4.397) (-1.375,-4.434) (-1.325,-4.469) (-1.277,-4.500) (-1.225,-4.533) (-1.175,-4.562) (-1.125,-4.590) (-1.075,-4.616) (-1.025,-4.640) (-0.975,-4.663) (-0.925,-4.685) (-0.875,-4.705) (-0.825,-4.724) (-0.775,-4.742) (-0.725,-4.758) (-0.675,-4.773) (-0.625,-4.788) (-0.575,-4.800) (-0.525,-4.812) (-0.475,-4.823) (-0.425,-4.832) (-0.375,-4.841) (-0.325,-4.848) (-0.300,-4.851)}
\EQ{(-0.125,0.322) (-0.075,0.314) (-0.025,0.309) (0.050,0.311) (0.100,0.317) (0.150,0.327) (0.200,0.341) (0.250,0.357) (0.297,0.375) (0.350,0.398) (0.400,0.421) (0.450,0.446) (0.500,0.473) (0.549,0.500) (0.593,0.525) (0.650,0.560) (0.700,0.590) (0.750,0.622) (0.793,0.650) (0.850,0.688) (0.900,0.722) (0.950,0.757) (1.000,0.793) (1.044,0.825) (1.100,0.868) (1.143,0.900) (1.200,0.945) (1.250,0.986) (1.297,1.025) (1.350,1.071) (1.400,1.115) (1.438,1.150) (1.475,1.185) (1.517,1.225) (1.567,1.275) (1.615,1.325) (1.650,1.364) (1.700,1.421) (1.746,1.475) (1.786,1.525) (1.824,1.575) (1.860,1.625) (1.894,1.675) (1.925,1.722) (1.958,1.775) (1.987,1.825) (2.014,1.875) (2.040,1.925) (2.064,1.975) (2.087,2.025) (2.108,2.075) (2.127,2.125) (2.145,2.175) (2.162,2.225) (2.177,2.275) (2.190,2.325) (2.203,2.375) (2.213,2.425) (2.223,2.475) (2.231,2.525) (2.237,2.575) (2.243,2.625) (2.246,2.675) (2.249,2.725) (2.250,2.775) (2.250,2.825) (2.249,2.875) (2.246,2.925) (2.242,2.975) (2.237,3.025) (2.230,3.075) (2.221,3.125) (2.212,3.175) (2.201,3.225) (2.189,3.275) (2.175,3.324) (2.160,3.375) (2.143,3.425) (2.125,3.474) (2.105,3.525) (2.083,3.575) (2.060,3.625) (2.035,3.675) (2.009,3.725) (1.981,3.775) (1.950,3.825) (1.918,3.875) (1.884,3.925) (1.850,3.971) (1.808,4.025) (1.775,4.065) (1.725,4.122) (1.676,4.175) (1.626,4.225) (1.575,4.273) (1.525,4.317) (1.485,4.350) (1.425,4.397) (1.375,4.434) (1.325,4.469) (1.277,4.500) (1.225,4.533) (1.175,4.562) (1.125,4.590) (1.075,4.616) (1.025,4.640) (0.975,4.663) (0.925,4.685) (0.875,4.705) (0.825,4.724) (0.775,4.742) (0.725,4.758) (0.675,4.773) (0.625,4.788) (0.575,4.800) (0.525,4.812) (0.475,4.823) (0.425,4.832) (0.375,4.841) (0.325,4.848) (0.275,4.854) (0.225,4.860) (0.175,4.864) (0.125,4.867) (0.075,4.869) (0.025,4.870) (-0.050,4.870) (-0.100,4.868) (-0.150,4.865) (-0.200,4.862) (-0.250,4.857) (-0.300,4.851) (-0.350,4.845) (-0.400,4.837) (-0.450,4.828) (-0.500,4.818) (-0.550,4.806) (-0.600,4.794) (-0.650,4.781) (-0.700,4.766) (-0.750,4.750) (-0.800,4.733) (-0.850,4.715) (-0.900,4.695) (-0.948,4.675) (-1.000,4.652) (-1.050,4.628) (-1.100,4.603) (-1.150,4.576) (-1.196,4.550) (-1.250,4.517) (-1.300,4.485) (-1.350,4.452) (-1.400,4.416) (-1.450,4.378) (-1.500,4.338) (-1.544,4.300) (-1.599,4.250) (-1.650,4.201) (-1.700,4.150) (-1.745,4.100) (-1.788,4.050) (-1.825,4.004) (-1.866,3.950) (-1.900,3.902) (-1.934,3.850) (-1.966,3.800) (-1.995,3.750) (-2.022,3.700) (-2.048,3.650) (-2.072,3.600) (-2.094,3.550) (-2.115,3.500) (-2.134,3.450) (-2.151,3.400) (-2.167,3.350) (-2.182,3.300) (-2.195,3.250) (-2.207,3.200) (-2.217,3.150) (-2.226,3.100) (-2.233,3.050) (-2.239,3.000) (-2.244,2.950) (-2.248,2.900) (-2.250,2.850) (-2.250,2.800) (-2.250,2.750) (-2.248,2.700) (-2.245,2.650) (-2.240,2.600) (-2.234,2.550) (-2.227,2.500) (-2.218,2.450) (-2.208,2.400) (-2.197,2.350) (-2.184,2.300) (-2.169,2.250) (-2.154,2.200) (-2.137,2.150) (-2.118,2.100) (-2.097,2.050) (-2.076,2.000) (-2.052,1.950) (-2.027,1.900) (-2.001,1.850) (-1.975,1.804) (-1.943,1.750) (-1.911,1.700) (-1.877,1.650) (-1.842,1.600) (-1.805,1.550) (-1.766,1.500) (-1.725,1.450) (-1.682,1.400) (-1.638,1.350) (-1.600,1.310) (-1.550,1.258) (-1.500,1.209) (-1.450,1.161) (-1.411,1.125) (-1.355,1.075) (-1.300,1.028) (-1.250,0.986) (-1.206,0.950) (-1.150,0.906) (-1.110,0.875) (-1.050,0.830) (-1.010,0.800) (-0.950,0.757) (-0.904,0.725) (-0.850,0.688) (-0.800,0.655) (-0.755,0.625) (-0.700,0.590) (-0.650,0.560) (-0.600,0.530) (-0.550,0.500) (-0.505,0.475) (-0.458,0.450) (-0.409,0.425) (-0.356,0.400) (-0.300,0.376) (-0.250,0.357) (-0.200,0.341) (-0.150,0.327) (-0.125,0.322)}
\EQ{(-0.200,-4.651) (-0.150,-4.655) (-0.100,-4.657) (-0.050,-4.659) (0.000,-4.660) (0.050,-4.659) (0.100,-4.657) (0.150,-4.655) (0.200,-4.651) (0.250,-4.646) (0.300,-4.640) (0.350,-4.633) (0.400,-4.625) (0.450,-4.615) (0.500,-4.605) (0.550,-4.593) (0.600,-4.580) (0.650,-4.566) (0.700,-4.550) (0.750,-4.534) (0.800,-4.516) (0.850,-4.497) (0.900,-4.476) (0.950,-4.454) (1.000,-4.430) (1.050,-4.405) (1.100,-4.378) (1.149,-4.350) (1.200,-4.319) (1.250,-4.287) (1.300,-4.253) (1.350,-4.217) (1.400,-4.178) (1.450,-4.137) (1.493,-4.100) (1.547,-4.050) (1.597,-4.000) (1.644,-3.950) (1.689,-3.900) (1.725,-3.856) (1.769,-3.800) (1.800,-3.757) (1.839,-3.700) (1.871,-3.650) (1.900,-3.601) (1.925,-3.557) (1.950,-3.509) (1.975,-3.458) (2.000,-3.403) (2.022,-3.350) (2.041,-3.300) (2.058,-3.250) (2.074,-3.200) (2.088,-3.150) (2.101,-3.100) (2.112,-3.050) (2.122,-3.000) (2.131,-2.950) (2.138,-2.900) (2.143,-2.850) (2.147,-2.800) (2.150,-2.750) (2.151,-2.700) (2.151,-2.650) (2.150,-2.600) (2.147,-2.550) (2.143,-2.500) (2.137,-2.450) (2.130,-2.400) (2.121,-2.350) (2.111,-2.300) (2.100,-2.250) (2.087,-2.200) (2.073,-2.150) (2.057,-2.100) (2.039,-2.050) (2.020,-2.000) (2.000,-1.951) (1.977,-1.900) (1.953,-1.850) (1.928,-1.800) (1.900,-1.750) (1.875,-1.707) (1.840,-1.650) (1.808,-1.600) (1.775,-1.553) (1.736,-1.500) (1.700,-1.453) (1.657,-1.400) (1.615,-1.350) (1.575,-1.305) (1.525,-1.252) (1.475,-1.200) (1.425,-1.151) (1.375,-1.104) (1.325,-1.059) (1.287,-1.025) (1.228,-0.975) (1.175,-0.931) (1.125,-0.891) (1.075,-0.853) (1.025,-0.815) (0.975,-0.778) (0.925,-0.742) (0.875,-0.707) (0.829,-0.675) (0.775,-0.640) (0.725,-0.607) (0.675,-0.575) (0.625,-0.545) (0.575,-0.515) (0.525,-0.487) (0.475,-0.460) (0.425,-0.434) (0.375,-0.409) (0.325,-0.387) (0.275,-0.367) (0.225,-0.349) (0.175,-0.334) (0.125,-0.322) (0.075,-0.314) (0.025,-0.309) (-0.050,-0.311) (-0.100,-0.317) (-0.150,-0.327) (-0.200,-0.341) (-0.250,-0.357) (-0.297,-0.375) (-0.350,-0.398) (-0.400,-0.421) (-0.450,-0.446) (-0.500,-0.473) (-0.549,-0.500) (-0.592,-0.525) (-0.650,-0.560) (-0.700,-0.591) (-0.750,-0.623) (-0.800,-0.656) (-0.850,-0.690) (-0.900,-0.724) (-0.950,-0.760) (-1.000,-0.796) (-1.050,-0.834) (-1.100,-0.872) (-1.150,-0.911) (-1.198,-0.950) (-1.250,-0.994) (-1.300,-1.037) (-1.343,-1.075) (-1.398,-1.125) (-1.449,-1.175) (-1.499,-1.225) (-1.547,-1.275) (-1.593,-1.325) (-1.636,-1.375) (-1.675,-1.422) (-1.717,-1.475) (-1.750,-1.519) (-1.791,-1.575) (-1.824,-1.625) (-1.856,-1.675) (-1.886,-1.725) (-1.914,-1.775) (-1.941,-1.825) (-1.966,-1.875) (-1.989,-1.925) (-2.010,-1.975) (-2.030,-2.025) (-2.048,-2.075) (-2.065,-2.125) (-2.080,-2.175) (-2.094,-2.225) (-2.106,-2.275) (-2.117,-2.325) (-2.126,-2.375) (-2.134,-2.425) (-2.140,-2.475) (-2.145,-2.525) (-2.149,-2.575) (-2.151,-2.625) (-2.151,-2.675) (-2.151,-2.725) (-2.149,-2.775) (-2.145,-2.825) (-2.141,-2.875) (-2.134,-2.925) (-2.127,-2.975) (-2.118,-3.025) (-2.107,-3.075) (-2.095,-3.125) (-2.081,-3.175) (-2.066,-3.225) (-2.050,-3.274) (-2.031,-3.325) (-2.012,-3.375) (-1.990,-3.425) (-1.967,-3.475) (-1.942,-3.525) (-1.915,-3.575) (-1.886,-3.625) (-1.855,-3.675) (-1.825,-3.721) (-1.787,-3.775) (-1.750,-3.825) (-1.710,-3.875) (-1.675,-3.916) (-1.625,-3.971) (-1.575,-4.023) (-1.525,-4.071) (-1.475,-4.116) (-1.435,-4.150) (-1.375,-4.198) (-1.325,-4.235) (-1.275,-4.270) (-1.230,-4.300) (-1.175,-4.335) (-1.125,-4.364) (-1.075,-4.392) (-1.025,-4.418) (-0.975,-4.442) (-0.925,-4.465) (-0.875,-4.486) (-0.825,-4.506) (-0.775,-4.525) (-0.725,-4.542) (-0.675,-4.558) (-0.625,-4.573) (-0.575,-4.587) (-0.525,-4.599) (-0.475,-4.610) (-0.425,-4.620) (-0.375,-4.629) (-0.325,-4.637) (-0.275,-4.643) (-0.225,-4.649) (-0.200,-4.651)}
\EQ{(-0.125,0.322) (-0.075,0.314) (-0.025,0.309) (0.050,0.311) (0.100,0.317) (0.150,0.327) (0.200,0.341) (0.250,0.357) (0.297,0.375) (0.350,0.398) (0.400,0.421) (0.450,0.446) (0.500,0.473) (0.549,0.500) (0.592,0.525) (0.650,0.560) (0.700,0.591) (0.750,0.623) (0.800,0.656) (0.850,0.690) (0.900,0.724) (0.950,0.760) (1.000,0.796) (1.050,0.834) (1.100,0.872) (1.150,0.911) (1.198,0.950) (1.250,0.994) (1.300,1.037) (1.343,1.075) (1.398,1.125) (1.449,1.175) (1.499,1.225) (1.547,1.275) (1.593,1.325) (1.636,1.375) (1.675,1.422) (1.717,1.475) (1.750,1.519) (1.791,1.575) (1.824,1.625) (1.856,1.675) (1.886,1.725) (1.914,1.775) (1.941,1.825) (1.966,1.875) (1.989,1.925) (2.010,1.975) (2.030,2.025) (2.048,2.075) (2.065,2.125) (2.080,2.175) (2.094,2.225) (2.106,2.275) (2.117,2.325) (2.126,2.375) (2.134,2.425) (2.140,2.475) (2.145,2.525) (2.149,2.575) (2.151,2.625) (2.151,2.675) (2.151,2.725) (2.149,2.775) (2.145,2.825) (2.141,2.875) (2.134,2.925) (2.127,2.975) (2.118,3.025) (2.107,3.075) (2.095,3.125) (2.081,3.175) (2.066,3.225) (2.050,3.274) (2.031,3.325) (2.012,3.375) (1.990,3.425) (1.967,3.475) (1.942,3.525) (1.915,3.575) (1.886,3.625) (1.855,3.675) (1.825,3.721) (1.787,3.775) (1.750,3.825) (1.710,3.875) (1.675,3.916) (1.625,3.971) (1.575,4.023) (1.525,4.071) (1.475,4.116) (1.435,4.150) (1.375,4.198) (1.325,4.235) (1.275,4.270) (1.230,4.300) (1.175,4.335) (1.125,4.364) (1.075,4.392) (1.025,4.418) (0.975,4.442) (0.925,4.465) (0.875,4.486) (0.825,4.506) (0.775,4.525) (0.725,4.542) (0.675,4.558) (0.625,4.573) (0.575,4.587) (0.525,4.599) (0.475,4.610) (0.425,4.620) (0.375,4.629) (0.325,4.637) (0.275,4.643) (0.225,4.649) (0.175,4.653) (0.125,4.656) (0.075,4.658) (0.025,4.659) (-0.050,4.659) (-0.100,4.657) (-0.150,4.655) (-0.200,4.651) (-0.250,4.646) (-0.300,4.640) (-0.350,4.633) (-0.400,4.625) (-0.450,4.615) (-0.500,4.605) (-0.550,4.593) (-0.600,4.580) (-0.650,4.566) (-0.700,4.550) (-0.750,4.534) (-0.800,4.516) (-0.850,4.497) (-0.900,4.476) (-0.950,4.454) (-1.000,4.430) (-1.050,4.405) (-1.100,4.378) (-1.149,4.350) (-1.200,4.319) (-1.250,4.287) (-1.300,4.253) (-1.350,4.217) (-1.400,4.178) (-1.450,4.137) (-1.493,4.100) (-1.547,4.050) (-1.597,4.000) (-1.644,3.950) (-1.689,3.900) (-1.725,3.856) (-1.769,3.800) (-1.800,3.757) (-1.839,3.700) (-1.871,3.650) (-1.900,3.601) (-1.925,3.557) (-1.950,3.509) (-1.975,3.458) (-2.000,3.403) (-2.022,3.350) (-2.041,3.300) (-2.058,3.250) (-2.074,3.200) (-2.088,3.150) (-2.101,3.100) (-2.112,3.050) (-2.122,3.000) (-2.131,2.950) (-2.138,2.900) (-2.143,2.850) (-2.147,2.800) (-2.150,2.750) (-2.151,2.700) (-2.151,2.650) (-2.150,2.600) (-2.147,2.550) (-2.143,2.500) (-2.137,2.450) (-2.130,2.400) (-2.121,2.350) (-2.111,2.300) (-2.100,2.250) (-2.087,2.200) (-2.073,2.150) (-2.057,2.100) (-2.039,2.050) (-2.020,2.000) (-2.000,1.951) (-1.977,1.900) (-1.953,1.850) (-1.928,1.800) (-1.900,1.750) (-1.875,1.707) (-1.840,1.650) (-1.808,1.600) (-1.775,1.553) (-1.736,1.500) (-1.700,1.453) (-1.657,1.400) (-1.615,1.350) (-1.575,1.305) (-1.525,1.252) (-1.475,1.200) (-1.425,1.151) (-1.375,1.104) (-1.325,1.059) (-1.287,1.025) (-1.228,0.975) (-1.175,0.931) (-1.125,0.891) (-1.075,0.853) (-1.025,0.815) (-0.975,0.778) (-0.925,0.742) (-0.875,0.707) (-0.829,0.675) (-0.775,0.640) (-0.725,0.607) (-0.675,0.575) (-0.625,0.545) (-0.575,0.515) (-0.525,0.487) (-0.475,0.460) (-0.425,0.434) (-0.375,0.409) (-0.325,0.387) (-0.275,0.367) (-0.225,0.349) (-0.175,0.334) (-0.125,0.322)}
\EQ{(-0.150,-4.476) (-0.100,-4.479) (-0.050,-4.481) (0.000,-4.482) (0.050,-4.481) (0.100,-4.479) (0.150,-4.476) (0.200,-4.473) (0.250,-4.467) (0.300,-4.461) (0.350,-4.454) (0.400,-4.445) (0.450,-4.435) (0.500,-4.424) (0.550,-4.412) (0.600,-4.399) (0.650,-4.384) (0.700,-4.368) (0.750,-4.350) (0.800,-4.332) (0.850,-4.311) (0.900,-4.290) (0.950,-4.266) (1.000,-4.241) (1.050,-4.215) (1.100,-4.187) (1.150,-4.157) (1.199,-4.125) (1.250,-4.090) (1.300,-4.054) (1.350,-4.015) (1.399,-3.975) (1.450,-3.930) (1.500,-3.883) (1.550,-3.832) (1.600,-3.778) (1.645,-3.725) (1.685,-3.675) (1.722,-3.625) (1.757,-3.575) (1.790,-3.525) (1.820,-3.475) (1.848,-3.425) (1.875,-3.375) (1.899,-3.325) (1.921,-3.275) (1.942,-3.225) (1.961,-3.175) (1.979,-3.125) (1.994,-3.075) (2.009,-3.025) (2.021,-2.975) (2.032,-2.925) (2.042,-2.875) (2.050,-2.825) (2.056,-2.775) (2.061,-2.725) (2.065,-2.675) (2.067,-2.625) (2.068,-2.575) (2.067,-2.525) (2.065,-2.475) (2.061,-2.425) (2.056,-2.375) (2.049,-2.325) (2.041,-2.275) (2.032,-2.225) (2.020,-2.175) (2.008,-2.125) (1.994,-2.075) (1.978,-2.025) (1.960,-1.975) (1.941,-1.925) (1.921,-1.875) (1.900,-1.829) (1.875,-1.777) (1.850,-1.729) (1.825,-1.683) (1.791,-1.625) (1.760,-1.575) (1.726,-1.525) (1.691,-1.475) (1.653,-1.425) (1.614,-1.375) (1.575,-1.328) (1.529,-1.275) (1.483,-1.225) (1.435,-1.175) (1.385,-1.125) (1.332,-1.075) (1.277,-1.025) (1.225,-0.980) (1.175,-0.937) (1.130,-0.900) (1.075,-0.857) (1.035,-0.825) (0.975,-0.781) (0.934,-0.750) (0.875,-0.709) (0.826,-0.675) (0.775,-0.641) (0.725,-0.608) (0.675,-0.576) (0.625,-0.545) (0.575,-0.516) (0.525,-0.487) (0.475,-0.460) (0.425,-0.434) (0.375,-0.410) (0.325,-0.387) (0.275,-0.367) (0.225,-0.349) (0.175,-0.334) (0.125,-0.322) (0.075,-0.314) (0.025,-0.309) (-0.050,-0.311) (-0.100,-0.317) (-0.150,-0.327) (-0.200,-0.341) (-0.250,-0.357) (-0.297,-0.375) (-0.350,-0.398) (-0.400,-0.421) (-0.450,-0.447) (-0.500,-0.473) (-0.548,-0.500) (-0.592,-0.525) (-0.650,-0.561) (-0.700,-0.592) (-0.750,-0.624) (-0.800,-0.657) (-0.850,-0.692) (-0.898,-0.725) (-0.950,-0.762) (-1.000,-0.799) (-1.050,-0.837) (-1.099,-0.875) (-1.150,-0.917) (-1.191,-0.950) (-1.249,-1.000) (-1.300,-1.046) (-1.350,-1.092) (-1.400,-1.140) (-1.450,-1.191) (-1.500,-1.244) (-1.550,-1.299) (-1.594,-1.350) (-1.625,-1.389) (-1.672,-1.450) (-1.709,-1.500) (-1.743,-1.550) (-1.775,-1.599) (-1.806,-1.650) (-1.835,-1.700) (-1.861,-1.750) (-1.886,-1.800) (-1.910,-1.850) (-1.931,-1.900) (-1.950,-1.948) (-1.969,-2.000) (-1.986,-2.050) (-2.001,-2.100) (-2.014,-2.150) (-2.026,-2.200) (-2.037,-2.250) (-2.045,-2.300) (-2.053,-2.350) (-2.059,-2.400) (-2.063,-2.450) (-2.066,-2.500) (-2.068,-2.550) (-2.068,-2.600) (-2.066,-2.650) (-2.063,-2.700) (-2.059,-2.750) (-2.053,-2.800) (-2.046,-2.850) (-2.037,-2.900) (-2.027,-2.950) (-2.015,-3.000) (-2.002,-3.050) (-1.987,-3.100) (-1.970,-3.150) (-1.952,-3.200) (-1.932,-3.250) (-1.910,-3.300) (-1.887,-3.350) (-1.862,-3.400) (-1.834,-3.450) (-1.805,-3.500) (-1.775,-3.548) (-1.740,-3.600) (-1.704,-3.650) (-1.665,-3.700) (-1.625,-3.749) (-1.580,-3.800) (-1.533,-3.850) (-1.482,-3.900) (-1.428,-3.950) (-1.375,-3.995) (-1.325,-4.035) (-1.275,-4.072) (-1.225,-4.108) (-1.175,-4.141) (-1.125,-4.172) (-1.077,-4.200) (-1.031,-4.225) (-0.983,-4.250) (-0.932,-4.275) (-0.876,-4.300) (-0.825,-4.322) (-0.775,-4.341) (-0.725,-4.359) (-0.675,-4.376) (-0.625,-4.391) (-0.575,-4.406) (-0.525,-4.418) (-0.475,-4.430) (-0.425,-4.440) (-0.375,-4.450) (-0.325,-4.458) (-0.275,-4.465) (-0.225,-4.470) (-0.175,-4.475) (-0.150,-4.476)}
\EQ{(-0.125,0.322) (-0.075,0.314) (-0.025,0.309) (0.050,0.311) (0.100,0.317) (0.150,0.327) (0.200,0.341) (0.250,0.357) (0.297,0.375) (0.350,0.398) (0.400,0.421) (0.450,0.447) (0.500,0.473) (0.548,0.500) (0.592,0.525) (0.650,0.561) (0.700,0.592) (0.750,0.624) (0.800,0.657) (0.850,0.692) (0.898,0.725) (0.950,0.762) (1.000,0.799) (1.050,0.837) (1.099,0.875) (1.150,0.917) (1.191,0.950) (1.249,1.000) (1.300,1.046) (1.350,1.092) (1.400,1.140) (1.450,1.191) (1.500,1.244) (1.550,1.299) (1.594,1.350) (1.625,1.389) (1.672,1.450) (1.709,1.500) (1.743,1.550) (1.775,1.599) (1.806,1.650) (1.835,1.700) (1.861,1.750) (1.886,1.800) (1.910,1.850) (1.931,1.900) (1.950,1.948) (1.969,2.000) (1.986,2.050) (2.001,2.100) (2.014,2.150) (2.026,2.200) (2.037,2.250) (2.045,2.300) (2.053,2.350) (2.059,2.400) (2.063,2.450) (2.066,2.500) (2.068,2.550) (2.068,2.600) (2.066,2.650) (2.063,2.700) (2.059,2.750) (2.053,2.800) (2.046,2.850) (2.037,2.900) (2.027,2.950) (2.015,3.000) (2.002,3.050) (1.987,3.100) (1.970,3.150) (1.952,3.200) (1.932,3.250) (1.910,3.300) (1.887,3.350) (1.862,3.400) (1.834,3.450) (1.805,3.500) (1.775,3.548) (1.740,3.600) (1.704,3.650) (1.665,3.700) (1.625,3.749) (1.580,3.800) (1.533,3.850) (1.482,3.900) (1.428,3.950) (1.375,3.995) (1.325,4.035) (1.275,4.072) (1.225,4.108) (1.175,4.141) (1.125,4.172) (1.077,4.200) (1.031,4.225) (0.983,4.250) (0.932,4.275) (0.876,4.300) (0.825,4.322) (0.775,4.341) (0.725,4.359) (0.675,4.376) (0.625,4.391) (0.575,4.406) (0.525,4.418) (0.475,4.430) (0.425,4.440) (0.375,4.450) (0.325,4.458) (0.275,4.465) (0.225,4.470) (0.175,4.475) (0.125,4.478) (0.075,4.480) (0.025,4.481) (-0.050,4.481) (-0.100,4.479) (-0.150,4.476) (-0.200,4.473) (-0.250,4.467) (-0.300,4.461) (-0.350,4.454) (-0.400,4.445) (-0.450,4.435) (-0.500,4.424) (-0.550,4.412) (-0.600,4.399) (-0.650,4.384) (-0.700,4.368) (-0.750,4.350) (-0.800,4.332) (-0.850,4.311) (-0.900,4.290) (-0.950,4.266) (-1.000,4.241) (-1.050,4.215) (-1.100,4.187) (-1.150,4.157) (-1.199,4.125) (-1.250,4.090) (-1.300,4.054) (-1.350,4.015) (-1.399,3.975) (-1.450,3.930) (-1.500,3.883) (-1.550,3.832) (-1.600,3.778) (-1.645,3.725) (-1.685,3.675) (-1.722,3.625) (-1.757,3.575) (-1.790,3.525) (-1.820,3.475) (-1.848,3.425) (-1.875,3.375) (-1.899,3.325) (-1.921,3.275) (-1.942,3.225) (-1.961,3.175) (-1.979,3.125) (-1.994,3.075) (-2.009,3.025) (-2.021,2.975) (-2.032,2.925) (-2.042,2.875) (-2.050,2.825) (-2.056,2.775) (-2.061,2.725) (-2.065,2.675) (-2.067,2.625) (-2.068,2.575) (-2.067,2.525) (-2.065,2.475) (-2.061,2.425) (-2.056,2.375) (-2.049,2.325) (-2.041,2.275) (-2.032,2.225) (-2.020,2.175) (-2.008,2.125) (-1.994,2.075) (-1.978,2.025) (-1.960,1.975) (-1.941,1.925) (-1.921,1.875) (-1.900,1.829) (-1.875,1.777) (-1.850,1.729) (-1.825,1.683) (-1.791,1.625) (-1.760,1.575) (-1.726,1.525) (-1.691,1.475) (-1.653,1.425) (-1.614,1.375) (-1.575,1.328) (-1.529,1.275) (-1.483,1.225) (-1.435,1.175) (-1.385,1.125) (-1.332,1.075) (-1.277,1.025) (-1.225,0.980) (-1.175,0.937) (-1.130,0.900) (-1.075,0.857) (-1.035,0.825) (-0.975,0.781) (-0.934,0.750) (-0.875,0.709) (-0.826,0.675) (-0.775,0.641) (-0.725,0.608) (-0.675,0.576) (-0.625,0.545) (-0.575,0.516) (-0.525,0.487) (-0.475,0.460) (-0.425,0.434) (-0.375,0.410) (-0.325,0.387) (-0.275,0.367) (-0.225,0.349) (-0.175,0.334) (-0.125,0.322)}
\EQ{(-0.100,-4.326) (-0.050,-4.328) (0.000,-4.328) (0.050,-4.328) (0.100,-4.326) (0.150,-4.323) (0.200,-4.319) (0.250,-4.314) (0.300,-4.307) (0.350,-4.299) (0.400,-4.291) (0.450,-4.280) (0.500,-4.269) (0.550,-4.256) (0.600,-4.242) (0.650,-4.227) (0.700,-4.210) (0.750,-4.192) (0.800,-4.172) (0.850,-4.151) (0.900,-4.128) (0.950,-4.104) (1.000,-4.078) (1.050,-4.050) (1.100,-4.021) (1.150,-3.989) (1.200,-3.955) (1.242,-3.925) (1.300,-3.880) (1.350,-3.839) (1.395,-3.800) (1.448,-3.750) (1.497,-3.700) (1.543,-3.650) (1.586,-3.600) (1.625,-3.551) (1.663,-3.500) (1.698,-3.450) (1.730,-3.400) (1.760,-3.350) (1.788,-3.300) (1.814,-3.250) (1.838,-3.200) (1.860,-3.150) (1.881,-3.100) (1.899,-3.050) (1.916,-3.000) (1.931,-2.950) (1.944,-2.900) (1.956,-2.850) (1.966,-2.800) (1.975,-2.750) (1.982,-2.700) (1.988,-2.650) (1.992,-2.600) (1.994,-2.550) (1.996,-2.500) (1.995,-2.450) (1.993,-2.400) (1.990,-2.350) (1.985,-2.300) (1.978,-2.250) (1.970,-2.200) (1.961,-2.150) (1.949,-2.100) (1.937,-2.050) (1.922,-2.000) (1.906,-1.950) (1.889,-1.900) (1.869,-1.850) (1.848,-1.800) (1.825,-1.750) (1.800,-1.700) (1.775,-1.652) (1.750,-1.608) (1.715,-1.550) (1.683,-1.500) (1.650,-1.453) (1.612,-1.400) (1.575,-1.352) (1.533,-1.300) (1.490,-1.250) (1.450,-1.206) (1.400,-1.153) (1.350,-1.103) (1.300,-1.055) (1.250,-1.009) (1.212,-0.975) (1.154,-0.925) (1.100,-0.881) (1.050,-0.841) (1.000,-0.802) (0.950,-0.765) (0.900,-0.728) (0.850,-0.693) (0.800,-0.659) (0.750,-0.625) (0.700,-0.593) (0.650,-0.561) (0.600,-0.531) (0.550,-0.502) (0.503,-0.475) (0.457,-0.450) (0.408,-0.425) (0.355,-0.400) (0.300,-0.376) (0.250,-0.357) (0.200,-0.341) (0.150,-0.327) (0.100,-0.317) (0.050,-0.311) (0.000,-0.309) (-0.050,-0.311) (-0.100,-0.317) (-0.150,-0.327) (-0.200,-0.341) (-0.250,-0.357) (-0.297,-0.375) (-0.350,-0.398) (-0.400,-0.422) (-0.450,-0.447) (-0.500,-0.473) (-0.548,-0.500) (-0.600,-0.531) (-0.650,-0.561) (-0.700,-0.593) (-0.750,-0.625) (-0.800,-0.659) (-0.850,-0.693) (-0.896,-0.725) (-0.950,-0.765) (-0.997,-0.800) (-1.050,-0.841) (-1.093,-0.875) (-1.150,-0.922) (-1.200,-0.965) (-1.240,-1.000) (-1.295,-1.050) (-1.347,-1.100) (-1.397,-1.150) (-1.444,-1.200) (-1.490,-1.250) (-1.525,-1.291) (-1.573,-1.350) (-1.612,-1.400) (-1.648,-1.450) (-1.683,-1.500) (-1.715,-1.550) (-1.745,-1.600) (-1.774,-1.650) (-1.800,-1.699) (-1.825,-1.750) (-1.848,-1.800) (-1.869,-1.850) (-1.889,-1.900) (-1.906,-1.950) (-1.922,-2.000) (-1.937,-2.050) (-1.949,-2.100) (-1.961,-2.150) (-1.970,-2.200) (-1.978,-2.250) (-1.985,-2.300) (-1.990,-2.350) (-1.993,-2.400) (-1.995,-2.450) (-1.996,-2.500) (-1.994,-2.550) (-1.992,-2.600) (-1.988,-2.650) (-1.982,-2.700) (-1.975,-2.750) (-1.966,-2.800) (-1.956,-2.850) (-1.944,-2.900) (-1.931,-2.950) (-1.916,-3.000) (-1.899,-3.050) (-1.881,-3.100) (-1.860,-3.150) (-1.838,-3.200) (-1.814,-3.250) (-1.788,-3.300) (-1.760,-3.350) (-1.730,-3.400) (-1.700,-3.447) (-1.663,-3.500) (-1.626,-3.550) (-1.586,-3.600) (-1.550,-3.642) (-1.500,-3.697) (-1.450,-3.748) (-1.400,-3.795) (-1.350,-3.839) (-1.307,-3.875) (-1.250,-3.919) (-1.207,-3.950) (-1.150,-3.989) (-1.100,-4.021) (-1.050,-4.050) (-1.006,-4.075) (-0.958,-4.100) (-0.907,-4.125) (-0.853,-4.150) (-0.800,-4.172) (-0.750,-4.192) (-0.700,-4.210) (-0.650,-4.227) (-0.600,-4.242) (-0.550,-4.256) (-0.500,-4.269) (-0.450,-4.280) (-0.400,-4.291) (-0.350,-4.299) (-0.300,-4.307) (-0.250,-4.314) (-0.200,-4.319) (-0.150,-4.323) (-0.100,-4.326)}
\EQ{(-0.125,0.322) (-0.075,0.314) (-0.025,0.309) (0.050,0.311) (0.100,0.317) (0.150,0.327) (0.200,0.341) (0.250,0.357) (0.297,0.375) (0.350,0.398) (0.400,0.422) (0.450,0.447) (0.500,0.473) (0.548,0.500) (0.600,0.531) (0.650,0.561) (0.700,0.593) (0.750,0.625) (0.800,0.659) (0.850,0.693) (0.896,0.725) (0.950,0.765) (0.997,0.800) (1.050,0.841) (1.093,0.875) (1.150,0.922) (1.200,0.965) (1.240,1.000) (1.295,1.050) (1.347,1.100) (1.397,1.150) (1.444,1.200) (1.490,1.250) (1.525,1.291) (1.573,1.350) (1.612,1.400) (1.648,1.450) (1.683,1.500) (1.715,1.550) (1.745,1.600) (1.774,1.650) (1.800,1.699) (1.825,1.750) (1.848,1.800) (1.869,1.850) (1.889,1.900) (1.906,1.950) (1.922,2.000) (1.937,2.050) (1.949,2.100) (1.961,2.150) (1.970,2.200) (1.978,2.250) (1.985,2.300) (1.990,2.350) (1.993,2.400) (1.995,2.450) (1.996,2.500) (1.994,2.550) (1.992,2.600) (1.988,2.650) (1.982,2.700) (1.975,2.750) (1.966,2.800) (1.956,2.850) (1.944,2.900) (1.931,2.950) (1.916,3.000) (1.899,3.050) (1.881,3.100) (1.860,3.150) (1.838,3.200) (1.814,3.250) (1.788,3.300) (1.760,3.350) (1.730,3.400) (1.700,3.447) (1.663,3.500) (1.626,3.550) (1.586,3.600) (1.550,3.642) (1.500,3.697) (1.450,3.748) (1.400,3.795) (1.350,3.839) (1.307,3.875) (1.250,3.919) (1.207,3.950) (1.150,3.989) (1.100,4.021) (1.050,4.050) (1.006,4.075) (0.958,4.100) (0.907,4.125) (0.853,4.150) (0.800,4.172) (0.750,4.192) (0.700,4.210) (0.650,4.227) (0.600,4.242) (0.550,4.256) (0.500,4.269) (0.450,4.280) (0.400,4.291) (0.350,4.299) (0.300,4.307) (0.250,4.314) (0.200,4.319) (0.150,4.323) (0.100,4.326) (0.050,4.328) (0.000,4.328) (-0.050,4.328) (-0.100,4.326) (-0.150,4.323) (-0.200,4.319) (-0.250,4.314) (-0.300,4.307) (-0.350,4.299) (-0.400,4.291) (-0.450,4.280) (-0.500,4.269) (-0.550,4.256) (-0.600,4.242) (-0.650,4.227) (-0.700,4.210) (-0.750,4.192) (-0.800,4.172) (-0.850,4.151) (-0.900,4.128) (-0.950,4.104) (-1.000,4.078) (-1.050,4.050) (-1.100,4.021) (-1.150,3.989) (-1.200,3.955) (-1.242,3.925) (-1.300,3.880) (-1.350,3.839) (-1.395,3.800) (-1.448,3.750) (-1.497,3.700) (-1.543,3.650) (-1.586,3.600) (-1.625,3.551) (-1.663,3.500) (-1.698,3.450) (-1.730,3.400) (-1.760,3.350) (-1.788,3.300) (-1.814,3.250) (-1.838,3.200) (-1.860,3.150) (-1.881,3.100) (-1.899,3.050) (-1.916,3.000) (-1.931,2.950) (-1.944,2.900) (-1.956,2.850) (-1.966,2.800) (-1.975,2.750) (-1.982,2.700) (-1.988,2.650) (-1.992,2.600) (-1.994,2.550) (-1.996,2.500) (-1.995,2.450) (-1.993,2.400) (-1.990,2.350) (-1.985,2.300) (-1.978,2.250) (-1.970,2.200) (-1.961,2.150) (-1.949,2.100) (-1.937,2.050) (-1.922,2.000) (-1.906,1.950) (-1.889,1.900) (-1.869,1.850) (-1.848,1.800) (-1.825,1.750) (-1.800,1.700) (-1.775,1.652) (-1.750,1.608) (-1.715,1.550) (-1.683,1.500) (-1.650,1.453) (-1.612,1.400) (-1.575,1.352) (-1.533,1.300) (-1.490,1.250) (-1.450,1.206) (-1.400,1.153) (-1.350,1.103) (-1.300,1.055) (-1.250,1.009) (-1.212,0.975) (-1.154,0.925) (-1.100,0.881) (-1.050,0.841) (-1.000,0.802) (-0.950,0.765) (-0.900,0.728) (-0.850,0.693) (-0.800,0.659) (-0.750,0.625) (-0.700,0.593) (-0.650,0.561) (-0.600,0.531) (-0.550,0.502) (-0.503,0.475) (-0.457,0.450) (-0.408,0.425) (-0.355,0.400) (-0.300,0.376) (-0.250,0.357) (-0.200,0.341) (-0.150,0.327) (-0.125,0.322)}

    % Cuadrupolo físico (lineal) - líneas de campo
% Generado por scripts/generar_datos.py (no editar a mano).
\LC{(0.337,1.139) (0.386,1.196) (0.437,1.251) (0.491,1.302) (0.548,1.351) (0.607,1.398) (0.669,1.440) (0.733,1.480) (0.799,1.515) (0.867,1.547) (0.936,1.574) (1.008,1.597) (1.081,1.615) (1.155,1.627) (1.229,1.634) (1.304,1.635) (1.379,1.630) (1.453,1.617) (1.526,1.598) (1.596,1.572) (1.663,1.538) (1.726,1.498) (1.783,1.450) (1.834,1.395) (1.878,1.334) (1.914,1.268) (1.940,1.198) (1.958,1.125) (1.966,1.051) (1.965,0.976) (1.956,0.901) (1.938,0.828) (1.913,0.758) (1.881,0.690) (1.843,0.625) (1.799,0.564) (1.751,0.507) (1.699,0.453) (1.643,0.403) (1.585,0.356) (1.523,0.314) (1.459,0.274) (1.393,0.239) (1.325,0.206) (1.256,0.177) (1.186,0.151)}
\LC{(-0.638,1.015) (-0.712,1.023) (-0.787,1.022) (-0.861,1.010) (-0.933,0.987) (-0.999,0.952) (-1.058,0.906) (-1.106,0.848) (-1.141,0.782) (-1.162,0.711) (-1.168,0.636) (-1.159,0.562) (-1.137,0.490) (-1.104,0.423)}
\LC{(-0.638,-1.015) (-0.712,-1.023) (-0.787,-1.022) (-0.861,-1.010) (-0.933,-0.987) (-0.999,-0.952) (-1.058,-0.906) (-1.106,-0.848) (-1.141,-0.782) (-1.162,-0.711) (-1.168,-0.636) (-1.159,-0.562) (-1.137,-0.490) (-1.104,-0.423)}
\LC{(0.447,-1.112) (0.508,-1.156) (0.572,-1.195) (0.639,-1.229) (0.708,-1.258) (0.779,-1.282) (0.852,-1.300) (0.926,-1.311) (1.001,-1.315) (1.075,-1.311) (1.149,-1.298) (1.221,-1.278) (1.290,-1.248) (1.354,-1.208) (1.412,-1.160) (1.461,-1.104) (1.501,-1.041) (1.530,-0.971) (1.547,-0.899) (1.553,-0.824) (1.547,-0.749) (1.531,-0.676) (1.506,-0.605) (1.473,-0.538) (1.432,-0.475) (1.385,-0.417) (1.332,-0.363) (1.276,-0.314) (1.256,-0.299)}
\LC{(0.068,-1.372) (0.077,-1.447) (0.086,-1.521) (0.095,-1.596) (0.105,-1.670) (0.115,-1.744) (0.126,-1.818) (0.138,-1.893) (0.149,-1.967) (0.162,-2.041) (0.174,-2.115) (0.187,-2.188) (0.201,-2.262) (0.215,-2.336) (0.229,-2.409) (0.244,-2.483) (0.260,-2.556) (0.276,-2.630) (0.292,-2.703) (0.309,-2.776) (0.326,-2.849) (0.344,-2.922) (0.362,-2.995) (0.380,-3.067) (0.399,-3.140) (0.419,-3.212) (0.439,-3.285) (0.459,-3.357) (0.480,-3.429) (0.501,-3.501) (0.523,-3.572) (0.545,-3.644) (0.568,-3.716) (0.591,-3.787) (0.614,-3.858) (0.638,-3.929) (0.662,-4.000) (0.687,-4.071) (0.713,-4.142) (0.738,-4.212) (0.765,-4.282) (0.791,-4.352) (0.818,-4.422) (0.846,-4.492) (0.874,-4.562) (0.902,-4.631) (0.931,-4.700) (0.961,-4.769) (0.990,-4.838) (1.021,-4.907) (1.051,-4.975) (1.082,-5.043) (1.114,-5.112) (1.146,-5.179) (1.178,-5.247) (1.211,-5.314) (1.245,-5.382) (1.278,-5.449) (1.312,-5.515) (1.347,-5.582) (1.382,-5.648) (1.418,-5.714) (1.454,-5.780) (1.490,-5.846) (1.527,-5.911) (1.564,-5.976) (1.602,-6.041) (1.640,-6.105) (1.679,-6.170) (1.718,-6.234) (1.757,-6.298) (1.797,-6.361) (1.837,-6.424) (1.878,-6.487) (1.919,-6.550) (1.961,-6.612) (2.003,-6.674) (2.045,-6.736) (2.088,-6.798) (2.132,-6.859) (2.175,-6.920) (2.219,-6.981) (2.264,-7.041) (2.309,-7.101) (2.355,-7.160) (2.400,-7.220) (2.447,-7.279) (2.493,-7.337) (2.541,-7.396) (2.588,-7.454) (2.636,-7.511) (2.684,-7.569) (2.733,-7.626) (2.782,-7.682) (2.832,-7.739) (2.882,-7.794) (2.933,-7.850) (2.983,-7.905) (3.035,-7.960) (3.086,-8.014) (3.138,-8.068) (3.191,-8.122) (3.244,-8.175) (3.297,-8.228) (3.351,-8.280) (3.405,-8.332) (3.459,-8.384) (3.514,-8.435) (3.570,-8.485) (3.625,-8.536) (3.681,-8.585) (3.738,-8.635) (3.795,-8.684) (3.852,-8.732) (3.909,-8.780) (3.967,-8.828) (4.026,-8.875) (4.084,-8.922) (4.144,-8.968) (4.203,-9.013) (4.263,-9.059) (4.323,-9.103) (4.384,-9.147) (4.445,-9.191) (4.506,-9.234) (4.568,-9.277) (4.630,-9.319) (4.692,-9.361) (4.755,-9.402) (4.818,-9.442) (4.882,-9.482) (4.945,-9.522) (5.009,-9.561) (5.074,-9.599) (5.139,-9.637) (5.204,-9.674) (5.269,-9.710) (5.335,-9.746) (5.401,-9.782) (5.468,-9.817) (5.534,-9.851) (5.601,-9.884) (5.669,-9.917) (5.736,-9.950) (5.804,-9.982) (5.850,-10.000)}
\LC{(-0.267,-2.215) (-0.286,-2.288) (-0.305,-2.360) (-0.326,-2.432) (-0.346,-2.504) (-0.368,-2.576) (-0.390,-2.648) (-0.413,-2.720) (-0.436,-2.791) (-0.460,-2.862) (-0.484,-2.933) (-0.510,-3.003) (-0.535,-3.074) (-0.562,-3.144) (-0.589,-3.214) (-0.616,-3.284) (-0.644,-3.353) (-0.673,-3.422) (-0.703,-3.491) (-0.733,-3.560) (-0.763,-3.629) (-0.795,-3.697) (-0.827,-3.765) (-0.859,-3.832) (-0.892,-3.900) (-0.926,-3.967) (-0.960,-4.033) (-0.995,-4.100) (-1.031,-4.166) (-1.067,-4.231) (-1.104,-4.297) (-1.141,-4.362) (-1.179,-4.427) (-1.217,-4.491) (-1.256,-4.555) (-1.296,-4.619) (-1.336,-4.682) (-1.377,-4.745) (-1.419,-4.807) (-1.461,-4.869) (-1.504,-4.931) (-1.547,-4.992) (-1.591,-5.053) (-1.635,-5.113) (-1.680,-5.173) (-1.726,-5.233) (-1.772,-5.292) (-1.819,-5.351) (-1.866,-5.409) (-1.914,-5.467) (-1.962,-5.524) (-2.012,-5.580) (-2.061,-5.637) (-2.111,-5.692) (-2.162,-5.748) (-2.213,-5.802) (-2.265,-5.857) (-2.318,-5.910) (-2.371,-5.963) (-2.424,-6.016) (-2.478,-6.068) (-2.533,-6.119) (-2.588,-6.170) (-2.644,-6.220) (-2.700,-6.270) (-2.757,-6.319) (-2.814,-6.367) (-2.872,-6.415) (-2.930,-6.462) (-2.989,-6.508) (-3.049,-6.554) (-3.109,-6.599) (-3.169,-6.644) (-3.230,-6.688) (-3.291,-6.731) (-3.353,-6.773) (-3.416,-6.815) (-3.478,-6.855) (-3.542,-6.896) (-3.606,-6.935) (-3.670,-6.974) (-3.735,-7.011) (-3.800,-7.048) (-3.866,-7.085) (-3.932,-7.120) (-3.998,-7.155) (-4.065,-7.188) (-4.133,-7.221) (-4.200,-7.253) (-4.269,-7.285) (-4.337,-7.315) (-4.406,-7.344) (-4.476,-7.373) (-4.545,-7.400) (-4.616,-7.427) (-4.686,-7.452) (-4.757,-7.477) (-4.828,-7.501) (-4.900,-7.523) (-4.971,-7.545) (-5.043,-7.566) (-5.116,-7.585) (-5.188,-7.604) (-5.261,-7.621) (-5.335,-7.638) (-5.408,-7.653) (-5.482,-7.667) (-5.556,-7.680) (-5.630,-7.692) (-5.704,-7.703) (-5.778,-7.712) (-5.853,-7.720) (-5.927,-7.727) (-6.002,-7.733) (-6.077,-7.738) (-6.152,-7.741) (-6.227,-7.743) (-6.302,-7.744) (-6.377,-7.744) (-6.452,-7.742) (-6.527,-7.739) (-6.602,-7.734) (-6.677,-7.728) (-6.751,-7.721) (-6.826,-7.713) (-6.900,-7.703) (-6.974,-7.691) (-7.048,-7.678) (-7.122,-7.664) (-7.195,-7.648) (-7.268,-7.631) (-7.340,-7.612) (-7.413,-7.592) (-7.484,-7.570) (-7.556,-7.546) (-7.627,-7.522) (-7.697,-7.495) (-7.766,-7.467) (-7.835,-7.438) (-7.904,-7.407) (-7.971,-7.375) (-8.038,-7.340) (-8.104,-7.305) (-8.169,-7.268) (-8.233,-7.229) (-8.297,-7.189) (-8.359,-7.147) (-8.420,-7.104) (-8.481,-7.059) (-8.540,-7.013) (-8.597,-6.965) (-8.654,-6.916) (-8.709,-6.865) (-8.764,-6.813) (-8.816,-6.760) (-8.867,-6.705) (-8.917,-6.649) (-8.966,-6.592) (-9.012,-6.533) (-9.057,-6.473) (-9.101,-6.412) (-9.143,-6.350) (-9.183,-6.287) (-9.222,-6.222) (-9.258,-6.157) (-9.293,-6.091) (-9.326,-6.023) (-9.357,-5.955) (-9.387,-5.886) (-9.414,-5.816) (-9.440,-5.746) (-9.463,-5.675) (-9.485,-5.603) (-9.505,-5.530) (-9.523,-5.458) (-9.539,-5.384) (-9.553,-5.311) (-9.565,-5.237) (-9.575,-5.162) (-9.583,-5.088) (-9.589,-5.013) (-9.593,-4.938) (-9.596,-4.863) (-9.596,-4.788) (-9.595,-4.713) (-9.592,-4.638) (-9.587,-4.563) (-9.580,-4.489) (-9.571,-4.414) (-9.561,-4.340) (-9.548,-4.266) (-9.535,-4.192) (-9.519,-4.119) (-9.502,-4.046) (-9.483,-3.973) (-9.463,-3.901) (-9.441,-3.829) (-9.417,-3.758) (-9.392,-3.687) (-9.366,-3.617) (-9.338,-3.548) (-9.308,-3.479) (-9.278,-3.410) (-9.246,-3.342) (-9.212,-3.275) (-9.178,-3.209) (-9.142,-3.143) (-9.105,-3.078) (-9.066,-3.013) (-9.027,-2.950) (-8.986,-2.887) (-8.944,-2.824) (-8.901,-2.763) (-8.858,-2.702) (-8.813,-2.642) (-8.767,-2.583) (-8.720,-2.524) (-8.672,-2.466) (-8.623,-2.409) (-8.574,-2.353) (-8.523,-2.298) (-8.472,-2.243) (-8.419,-2.189) (-8.366,-2.136) (-8.313,-2.084) (-8.258,-2.032) (-8.203,-1.982) (-8.147,-1.932) (-8.090,-1.882) (-8.033,-1.834) (-7.975,-1.787) (-7.916,-1.740) (-7.857,-1.694) (-7.797,-1.649) (-7.737,-1.604) (-7.676,-1.560) (-7.614,-1.517) (-7.552,-1.475) (-7.490,-1.434) (-7.427,-1.393) (-7.363,-1.353) (-7.299,-1.314) (-7.235,-1.276) (-7.170,-1.238) (-7.105,-1.201) (-7.039,-1.165) (-6.973,-1.129) (-6.907,-1.095) (-6.840,-1.060) (-6.773,-1.027) (-6.705,-0.994) (-6.637,-0.962) (-6.569,-0.931) (-6.501,-0.900) (-6.432,-0.870) (-6.363,-0.841) (-6.294,-0.812) (-6.224,-0.784) (-6.154,-0.757) (-6.084,-0.730) (-6.014,-0.704) (-5.943,-0.679) (-5.873,-0.654) (-5.802,-0.630) (-5.731,-0.606) (-5.659,-0.583) (-5.588,-0.561) (-5.516,-0.539) (-5.444,-0.517) (-5.372,-0.497) (-5.300,-0.476) (-5.227,-0.457) (-5.155,-0.438) (-5.082,-0.419) (-5.009,-0.401) (-4.936,-0.384) (-4.863,-0.367) (-4.790,-0.350) (-4.717,-0.334) (-4.643,-0.319) (-4.570,-0.304) (-4.496,-0.289) (-4.423,-0.275) (-4.349,-0.262) (-4.275,-0.249) (-4.201,-0.236) (-4.127,-0.224) (-4.053,-0.212) (-3.979,-0.201) (-3.905,-0.190) (-3.830,-0.179) (-3.756,-0.169) (-3.682,-0.159) (-3.607,-0.150) (-3.533,-0.141) (-3.458,-0.132) (-3.384,-0.124) (-3.309,-0.116) (-3.235,-0.108) (-3.160,-0.101) (-3.085,-0.094) (-3.011,-0.088) (-2.936,-0.082) (-2.861,-0.076) (-2.786,-0.070) (-2.712,-0.065) (-2.637,-0.060) (-2.562,-0.055) (-2.487,-0.050) (-2.412,-0.046) (-2.337,-0.042) (-2.262,-0.038) (-2.187,-0.035) (-2.112,-0.031) (-2.038,-0.028) (-1.963,-0.025) (-1.888,-0.023) (-1.813,-0.020) (-1.738,-0.018) (-1.663,-0.016) (-1.588,-0.014) (-1.513,-0.012) (-1.438,-0.011) (-1.363,-0.009) (-1.288,-0.008)}
\LC{(0.756,-2.701) (0.797,-2.764) (0.839,-2.826) (0.883,-2.887) (0.927,-2.948) (0.972,-3.008) (1.018,-3.067) (1.065,-3.125) (1.113,-3.183) (1.162,-3.240) (1.213,-3.295) (1.264,-3.350) (1.316,-3.404) (1.369,-3.457) (1.423,-3.509) (1.478,-3.560) (1.533,-3.611) (1.590,-3.660) (1.648,-3.708) (1.706,-3.754) (1.766,-3.800) (1.826,-3.845) (1.887,-3.888) (1.950,-3.930) (2.012,-3.971) (2.076,-4.010) (2.141,-4.048) (2.206,-4.085) (2.272,-4.121) (2.339,-4.155) (2.407,-4.187) (2.475,-4.218) (2.544,-4.248) (2.614,-4.275) (2.684,-4.302) (2.755,-4.326) (2.826,-4.349) (2.898,-4.370) (2.971,-4.389) (3.044,-4.407) (3.117,-4.422) (3.191,-4.436) (3.265,-4.447) (3.339,-4.457) (3.414,-4.464) (3.489,-4.470) (3.564,-4.473) (3.639,-4.474) (3.714,-4.472) (3.788,-4.469) (3.863,-4.463) (3.938,-4.454) (4.012,-4.443) (4.086,-4.430) (4.159,-4.414) (4.232,-4.396) (4.304,-4.375) (4.375,-4.351) (4.445,-4.325) (4.514,-4.296) (4.582,-4.264) (4.649,-4.230) (4.714,-4.193) (4.778,-4.153) (4.840,-4.111) (4.900,-4.066) (4.958,-4.019) (5.014,-3.969) (5.068,-3.916) (5.119,-3.862) (5.168,-3.805) (5.214,-3.745) (5.257,-3.684) (5.297,-3.621) (5.334,-3.556) (5.368,-3.489) (5.399,-3.421) (5.427,-3.351) (5.452,-3.280) (5.473,-3.208) (5.490,-3.135) (5.505,-3.062) (5.516,-2.987) (5.523,-2.913) (5.527,-2.838) (5.528,-2.763) (5.526,-2.688) (5.521,-2.613) (5.512,-2.539) (5.501,-2.465) (5.486,-2.391) (5.469,-2.318) (5.448,-2.246) (5.426,-2.174) (5.400,-2.104) (5.372,-2.034) (5.342,-1.966) (5.309,-1.898) (5.274,-1.832) (5.237,-1.767) (5.197,-1.703) (5.156,-1.640) (5.113,-1.579) (5.068,-1.519) (5.022,-1.460) (4.974,-1.402) (4.924,-1.346) (4.873,-1.292) (4.820,-1.238) (4.766,-1.186) (4.711,-1.135) (4.654,-1.086) (4.596,-1.038) (4.538,-0.992) (4.478,-0.946) (4.417,-0.902) (4.355,-0.860) (4.293,-0.819) (4.229,-0.779) (4.165,-0.740) (4.100,-0.703) (4.034,-0.666) (3.968,-0.631) (3.901,-0.598) (3.833,-0.565) (3.765,-0.534) (3.697,-0.504) (3.627,-0.475) (3.558,-0.447) (3.488,-0.420) (3.417,-0.394) (3.346,-0.370) (3.275,-0.346) (3.204,-0.324) (3.132,-0.302) (3.060,-0.282) (2.987,-0.262) (2.915,-0.243) (2.842,-0.226) (2.769,-0.209) (2.695,-0.193) (2.622,-0.178) (2.548,-0.163) (2.475,-0.150) (2.401,-0.137) (2.327,-0.125) (2.252,-0.114) (2.178,-0.103) (2.104,-0.093) (2.029,-0.084) (1.955,-0.076) (1.880,-0.068) (1.806,-0.060) (1.781,-0.058)}
\LC{(-0.106,1.195) (-0.122,1.268) (-0.139,1.341) (-0.157,1.414) (-0.175,1.487) (-0.195,1.559) (-0.216,1.631) (-0.238,1.703) (-0.260,1.774) (-0.284,1.846) (-0.308,1.916) (-0.334,1.987) (-0.360,2.057) (-0.388,2.127) (-0.416,2.196) (-0.445,2.265) (-0.476,2.334) (-0.507,2.402) (-0.539,2.470) (-0.572,2.537) (-0.606,2.604) (-0.641,2.671) (-0.677,2.737) (-0.713,2.802) (-0.751,2.867) (-0.790,2.931) (-0.829,2.995) (-0.869,3.058) (-0.911,3.121) (-0.953,3.183) (-0.996,3.244) (-1.040,3.305) (-1.085,3.365) (-1.131,3.424) (-1.178,3.483) (-1.226,3.541) (-1.274,3.598) (-1.323,3.654) (-1.374,3.710) (-1.425,3.765) (-1.477,3.819) (-1.530,3.872) (-1.583,3.924) (-1.638,3.976) (-1.693,4.026) (-1.750,4.076) (-1.807,4.125) (-1.864,4.172) (-1.923,4.219) (-1.983,4.265) (-2.043,4.310) (-2.104,4.353) (-2.166,4.396) (-2.228,4.437) (-2.291,4.477) (-2.355,4.517) (-2.420,4.555) (-2.485,4.591) (-2.552,4.627) (-2.618,4.661) (-2.686,4.694) (-2.754,4.725) (-2.822,4.755) (-2.892,4.784) (-2.962,4.812) (-3.032,4.837) (-3.103,4.862) (-3.174,4.885) (-3.246,4.906) (-3.319,4.926) (-3.391,4.944) (-3.465,4.960) (-3.538,4.975) (-3.612,4.987) (-3.686,4.999) (-3.761,5.008) (-3.835,5.015) (-3.910,5.021) (-3.985,5.025) (-4.060,5.026) (-4.135,5.026) (-4.210,5.024) (-4.285,5.019) (-4.359,5.013) (-4.434,5.004) (-4.508,4.993) (-4.582,4.980) (-4.655,4.965) (-4.728,4.947) (-4.801,4.927) (-4.872,4.905) (-4.943,4.881) (-5.013,4.854) (-5.082,4.824) (-5.150,4.793) (-5.217,4.759) (-5.283,4.722) (-5.347,4.684) (-5.409,4.642) (-5.471,4.599) (-5.530,4.553) (-5.588,4.505) (-5.643,4.455) (-5.697,4.402) (-5.748,4.348) (-5.798,4.291) (-5.844,4.233) (-5.889,4.172) (-5.931,4.110) (-5.970,4.046) (-6.006,3.981) (-6.040,3.914) (-6.071,3.845) (-6.099,3.776) (-6.124,3.705) (-6.146,3.633) (-6.165,3.561) (-6.181,3.488) (-6.194,3.414) (-6.204,3.339) (-6.212,3.265) (-6.216,3.190) (-6.217,3.115) (-6.215,3.040) (-6.211,2.965) (-6.203,2.890) (-6.193,2.816) (-6.181,2.742) (-6.165,2.669) (-6.148,2.596) (-6.127,2.524) (-6.105,2.452) (-6.080,2.382) (-6.052,2.312) (-6.023,2.243) (-5.991,2.175) (-5.958,2.108) (-5.922,2.042) (-5.884,1.977) (-5.845,1.913) (-5.804,1.850) (-5.761,1.788) (-5.717,1.728) (-5.671,1.669) (-5.624,1.611) (-5.575,1.554) (-5.525,1.498) (-5.473,1.443) (-5.421,1.390) (-5.367,1.338) (-5.311,1.287) (-5.255,1.237) (-5.198,1.189) (-5.140,1.142) (-5.081,1.096) (-5.021,1.051) (-4.960,1.007) (-4.898,0.965) (-4.835,0.923) (-4.772,0.883) (-4.708,0.844) (-4.643,0.806) (-4.578,0.769) (-4.512,0.734) (-4.445,0.699) (-4.378,0.666) (-4.310,0.633) (-4.242,0.602) (-4.174,0.572) (-4.105,0.542) (-4.035,0.514) (-3.965,0.487) (-3.895,0.460) (-3.825,0.435) (-3.754,0.411) (-3.682,0.387) (-3.611,0.365) (-3.539,0.343) (-3.467,0.322)}
\LC{(0.526,1.891) (0.568,1.953) (0.612,2.014) (0.657,2.074) (0.704,2.132) (0.752,2.190) (0.801,2.247) (0.852,2.302) (0.904,2.355) (0.958,2.408) (1.013,2.459) (1.069,2.508) (1.127,2.556) (1.186,2.603) (1.246,2.647) (1.308,2.690) (1.371,2.731) (1.435,2.770) (1.500,2.807) (1.566,2.842) (1.633,2.875) (1.702,2.906) (1.771,2.935) (1.841,2.961) (1.912,2.985) (1.984,3.006) (2.057,3.024) (2.130,3.040) (2.204,3.053) (2.279,3.063) (2.353,3.070) (2.428,3.074) (2.503,3.074) (2.578,3.072) (2.653,3.066) (2.727,3.056) (2.801,3.043) (2.874,3.026) (2.946,3.005) (3.017,2.981) (3.087,2.952) (3.154,2.920) (3.220,2.884) (3.283,2.844) (3.344,2.800) (3.402,2.752) (3.457,2.701) (3.508,2.646) (3.556,2.588) (3.599,2.527) (3.638,2.463) (3.673,2.397) (3.703,2.328) (3.729,2.258) (3.749,2.186) (3.765,2.112) (3.775,2.038) (3.781,1.963) (3.782,1.888) (3.778,1.813) (3.770,1.739) (3.757,1.665) (3.740,1.592) (3.719,1.520) (3.694,1.449) (3.666,1.380) (3.634,1.312) (3.598,1.246) (3.560,1.181) (3.519,1.119) (3.475,1.058) (3.428,0.999) (3.380,0.942) (3.329,0.887) (3.275,0.834) (3.220,0.783) (3.164,0.734) (3.105,0.687) (3.045,0.642) (2.984,0.599) (2.921,0.558) (2.857,0.519) (2.792,0.481) (2.726,0.446) (2.659,0.412) (2.592,0.380) (2.523,0.349) (2.454,0.321)}
\LC{(0.293,2.351) (0.312,2.424) (0.332,2.496) (0.353,2.568) (0.374,2.640) (0.396,2.712) (0.418,2.783) (0.442,2.855) (0.465,2.926) (0.489,2.997) (0.514,3.067) (0.540,3.138) (0.566,3.208) (0.592,3.278) (0.620,3.348) (0.647,3.418) (0.676,3.487) (0.705,3.557) (0.734,3.626) (0.765,3.694) (0.795,3.763) (0.827,3.831) (0.859,3.899) (0.891,3.966) (0.924,4.033) (0.958,4.100) (0.992,4.167) (1.027,4.233) (1.063,4.300) (1.099,4.365) (1.136,4.431) (1.173,4.496) (1.211,4.561) (1.249,4.625) (1.288,4.689) (1.327,4.753) (1.368,4.816) (1.408,4.879) (1.450,4.942) (1.491,5.004) (1.534,5.066) (1.577,5.127) (1.620,5.188) (1.664,5.249) (1.709,5.309) (1.754,5.369) (1.800,5.429) (1.846,5.488) (1.893,5.546) (1.941,5.604) (1.989,5.662) (2.038,5.719) (2.087,5.776) (2.136,5.832) (2.187,5.887) (2.237,5.943) (2.289,5.997) (2.341,6.051) (2.393,6.105) (2.446,6.158) (2.499,6.211) (2.553,6.263) (2.608,6.314) (2.663,6.365) (2.719,6.416) (2.775,6.465) (2.831,6.515) (2.888,6.563) (2.946,6.611) (3.004,6.658) (3.063,6.705) (3.122,6.751) (3.182,6.797) (3.242,6.842) (3.302,6.886) (3.364,6.929) (3.425,6.972) (3.487,7.014) (3.550,7.056) (3.613,7.096) (3.676,7.136) (3.740,7.175) (3.805,7.214) (3.869,7.252) (3.935,7.289) (4.000,7.325) (4.066,7.360) (4.133,7.395) (4.200,7.428) (4.267,7.461) (4.335,7.494) (4.403,7.525) (4.472,7.555) (4.541,7.585) (4.610,7.613) (4.680,7.641) (4.750,7.668) (4.820,7.694) (4.891,7.719) (4.962,7.743) (5.033,7.766) (5.105,7.788) (5.177,7.809) (5.249,7.829) (5.322,7.848) (5.394,7.866) (5.468,7.883) (5.541,7.899) (5.614,7.914) (5.688,7.928) (5.762,7.941) (5.836,7.952) (5.910,7.963) (5.985,7.972) (6.059,7.980) (6.134,7.987) (6.209,7.993) (6.284,7.998) (6.359,8.001) (6.433,8.003) (6.508,8.004) (6.583,8.004) (6.658,8.002) (6.733,7.999) (6.808,7.995) (6.883,7.989) (6.958,7.982) (7.032,7.974) (7.107,7.965) (7.181,7.954) (7.255,7.941) (7.329,7.927) (7.402,7.912) (7.475,7.896) (7.548,7.878) (7.620,7.858) (7.692,7.837) (7.764,7.815) (7.835,7.791) (7.906,7.766) (7.976,7.739) (8.045,7.711) (8.114,7.681) (8.182,7.650) (8.250,7.617) (8.316,7.583) (8.382,7.547) (8.447,7.510) (8.512,7.471) (8.575,7.431) (8.637,7.389) (8.699,7.346) (8.759,7.301) (8.818,7.255) (8.876,7.208) (8.933,7.159) (8.989,7.109) (9.043,7.057) (9.096,7.004) (9.148,6.950) (9.198,6.894) (9.247,6.837) (9.295,6.779) (9.340,6.720) (9.385,6.659) (9.428,6.598) (9.469,6.535) (9.508,6.471) (9.546,6.406) (9.582,6.341) (9.616,6.274) (9.649,6.206) (9.680,6.138) (9.708,6.069) (9.736,5.999) (9.761,5.928) (9.784,5.857) (9.806,5.785) (9.826,5.713) (9.843,5.640) (9.859,5.567) (9.873,5.493) (9.885,5.419) (9.896,5.345) (9.904,5.270) (9.911,5.195) (9.915,5.121) (9.918,5.046) (9.919,4.971) (9.918,4.896) (9.916,4.821) (9.911,4.746) (9.905,4.671) (9.897,4.596) (9.888,4.522) (9.876,4.448) (9.863,4.374) (9.849,4.300) (9.832,4.227) (9.815,4.154) (9.795,4.082) (9.774,4.010) (9.752,3.938) (9.728,3.867) (9.703,3.797) (9.676,3.727) (9.648,3.657) (9.618,3.588) (9.588,3.520) (9.555,3.452) (9.522,3.385) (9.487,3.318) (9.452,3.252) (9.415,3.187) (9.376,3.123) (9.337,3.059) (9.296,2.996) (9.255,2.933) (9.212,2.872) (9.169,2.810) (9.124,2.750) (9.079,2.691) (9.032,2.632) (8.985,2.574) (8.936,2.516) (8.887,2.460) (8.837,2.404) (8.786,2.349) (8.734,2.295) (8.682,2.241) (8.628,2.188) (8.574,2.136) (8.520,2.085) (8.464,2.035) (8.408,1.985) (8.351,1.936) (8.294,1.888) (8.236,1.840) (8.177,1.793) (8.118,1.748) (8.058,1.702) (7.998,1.658) (7.937,1.614) (7.875,1.571) (7.813,1.529) (7.751,1.487) (7.688,1.447) (7.624,1.407) (7.560,1.367) (7.496,1.329) (7.431,1.291) (7.366,1.254) (7.301,1.217) (7.235,1.181) (7.169,1.146) (7.102,1.112) (7.035,1.078) (6.968,1.045) (6.900,1.013) (6.832,0.981) (6.764,0.950) (6.695,0.920) (6.626,0.890) (6.557,0.861) (6.488,0.832) (6.418,0.804) (6.348,0.777) (6.278,0.750) (6.208,0.724) (6.137,0.699) (6.066,0.674) (5.996,0.650) (5.924,0.626) (5.853,0.603) (5.781,0.581) (5.710,0.559) (5.638,0.538) (5.566,0.517) (5.494,0.496) (5.421,0.477) (5.349,0.458) (5.276,0.439) (5.203,0.421) (5.130,0.403) (5.057,0.386) (4.984,0.369) (4.911,0.353) (4.838,0.337) (4.764,0.322) (4.691,0.307) (4.617,0.293) (4.543,0.279) (4.470,0.266) (4.396,0.253) (4.322,0.240) (4.248,0.228) (4.174,0.216) (4.100,0.205) (4.025,0.194) (3.951,0.184) (3.877,0.174) (3.802,0.164) (3.728,0.155) (3.654,0.146) (3.579,0.137) (3.504,0.129) (3.430,0.121) (3.355,0.113) (3.330,0.111)}
\LC{(1.019,3.129) (1.066,3.188) (1.113,3.246) (1.161,3.304) (1.211,3.360) (1.261,3.416) (1.312,3.471) (1.364,3.525) (1.417,3.578) (1.471,3.630) (1.526,3.681) (1.582,3.732) (1.638,3.781) (1.696,3.829) (1.754,3.876) (1.813,3.922) (1.874,3.967) (1.935,4.010) (1.996,4.053) (2.059,4.094) (2.122,4.134) (2.187,4.173) (2.252,4.210) (2.317,4.247) (2.384,4.281) (2.451,4.315) (2.519,4.347) (2.587,4.377) (2.656,4.406) (2.726,4.433) (2.797,4.459) (2.868,4.483) (2.939,4.506) (3.011,4.527) (3.084,4.546) (3.157,4.563) (3.230,4.578) (3.304,4.592) (3.378,4.603) (3.452,4.613) (3.527,4.621) (3.602,4.626) (3.677,4.630) (3.752,4.631) (3.827,4.630) (3.902,4.627) (3.977,4.622) (4.051,4.615) (4.125,4.605) (4.199,4.592) (4.273,4.577) (4.346,4.560) (4.418,4.540) (4.490,4.518) (4.561,4.493) (4.631,4.466) (4.699,4.436) (4.767,4.404) (4.833,4.369) (4.898,4.331) (4.961,4.291) (5.023,4.248) (5.083,4.203) (5.141,4.155) (5.196,4.105) (5.250,4.052) (5.301,3.997) (5.350,3.940) (5.396,3.881) (5.439,3.820) (5.480,3.757) (5.518,3.692) (5.552,3.626) (5.584,3.558) (5.613,3.488) (5.638,3.418) (5.660,3.346) (5.679,3.274) (5.695,3.200) (5.707,3.126) (5.716,3.052) (5.722,2.977) (5.725,2.902) (5.724,2.827) (5.721,2.752) (5.714,2.678) (5.704,2.603) (5.692,2.529) (5.677,2.456) (5.659,2.383) (5.638,2.311) (5.615,2.240) (5.589,2.169) (5.560,2.100) (5.530,2.031) (5.497,1.964) (5.462,1.898) (5.425,1.832) (5.386,1.768) (5.345,1.706) (5.302,1.644) (5.258,1.584) (5.211,1.525) (5.164,1.467) (5.114,1.410) (5.063,1.355) (5.011,1.301) (4.958,1.249) (4.903,1.198) (4.847,1.148) (4.790,1.099) (4.732,1.052) (4.672,1.006) (4.612,0.961) (4.551,0.918) (4.489,0.876) (4.426,0.835) (4.362,0.795) (4.298,0.757) (4.233,0.720) (4.167,0.684) (4.145,0.672)}
\LC{(-10.000,1.451) (-9.943,1.420) (-9.875,1.389) (-9.807,1.358) (-9.738,1.328) (-9.669,1.298) (-9.600,1.269) (-9.531,1.240) (-9.461,1.212) (-9.392,1.184) (-9.322,1.156) (-9.252,1.129) (-9.182,1.102) (-9.112,1.076) (-9.041,1.050) (-8.971,1.025) (-8.900,1.000) (-8.829,0.975) (-8.758,0.951) (-8.687,0.927) (-8.616,0.904) (-8.545,0.881) (-8.473,0.858) (-8.401,0.836) (-8.330,0.814) (-8.258,0.792) (-8.186,0.771) (-8.114,0.750) (-8.042,0.730) (-7.969,0.710) (-7.897,0.691) (-7.824,0.671) (-7.752,0.652) (-7.679,0.634) (-7.606,0.616) (-7.534,0.598) (-7.461,0.580) (-7.388,0.563) (-7.314,0.546) (-7.241,0.530) (-7.168,0.514) (-7.095,0.498) (-7.021,0.482) (-6.948,0.467) (-6.874,0.453) (-6.801,0.438) (-6.727,0.424) (-6.653,0.410) (-6.580,0.396) (-6.506,0.383) (-6.432,0.370) (-6.358,0.357) (-6.284,0.345) (-6.210,0.333) (-6.136,0.321) (-6.062,0.310) (-5.988,0.298) (-5.914,0.287) (-5.839,0.277) (-5.765,0.266) (-5.691,0.256) (-5.616,0.246) (-5.542,0.236) (-5.468,0.227) (-5.393,0.218) (-5.319,0.209) (-5.244,0.200) (-5.170,0.192) (-5.095,0.184) (-5.021,0.176) (-4.946,0.168) (-4.871,0.161) (-4.797,0.154) (-4.722,0.147) (-4.647,0.140)}
\LC{(1.701,-5.279) (1.746,-5.339) (1.792,-5.398) (1.838,-5.457) (1.885,-5.516) (1.933,-5.574) (1.981,-5.631) (2.030,-5.689) (2.079,-5.745) (2.129,-5.801) (2.179,-5.857) (2.230,-5.912) (2.281,-5.967) (2.333,-6.021) (2.386,-6.074) (2.439,-6.127) (2.492,-6.180) (2.546,-6.232) (2.601,-6.283) (2.656,-6.334) (2.712,-6.384) (2.768,-6.434) (2.825,-6.483) (2.882,-6.531) (2.940,-6.579) (2.998,-6.627) (3.057,-6.673) (3.116,-6.719) (3.176,-6.764) (3.236,-6.809) (3.297,-6.853) (3.358,-6.896) (3.420,-6.939) (3.482,-6.981) (3.544,-7.022) (3.608,-7.063) (3.671,-7.102) (3.735,-7.141) (3.800,-7.180) (3.865,-7.217) (3.930,-7.254) (3.996,-7.290) (4.062,-7.325) (4.129,-7.359) (4.196,-7.393) (4.264,-7.425) (4.331,-7.457) (4.400,-7.488) (4.468,-7.518) (4.538,-7.547) (4.607,-7.576) (4.677,-7.603) (4.747,-7.630) (4.817,-7.655) (4.888,-7.680) (4.959,-7.703) (5.031,-7.726) (5.103,-7.748) (5.175,-7.768) (5.247,-7.788) (5.320,-7.807) (5.393,-7.824) (5.466,-7.841) (5.539,-7.856) (5.613,-7.871) (5.687,-7.884) (5.761,-7.896) (5.835,-7.907) (5.909,-7.917) (5.984,-7.926) (6.058,-7.933) (6.133,-7.940) (6.208,-7.945) (6.283,-7.949) (6.358,-7.952) (6.433,-7.953) (6.508,-7.953) (6.583,-7.952) (6.658,-7.950) (6.733,-7.946) (6.807,-7.941) (6.882,-7.935) (6.957,-7.927) (7.031,-7.918) (7.106,-7.908) (7.180,-7.896) (7.253,-7.883) (7.327,-7.868) (7.400,-7.852) (7.473,-7.834) (7.546,-7.816) (7.618,-7.795) (7.690,-7.773) (7.761,-7.750) (7.832,-7.725) (7.902,-7.699) (7.972,-7.671) (8.041,-7.642) (8.109,-7.611) (8.177,-7.579) (8.244,-7.545) (8.310,-7.510) (8.375,-7.473) (8.440,-7.434) (8.503,-7.395) (8.566,-7.353) (8.627,-7.311) (8.688,-7.266) (8.747,-7.221) (8.806,-7.174) (8.863,-7.125) (8.919,-7.075) (8.974,-7.024) (9.027,-6.971) (9.079,-6.917) (9.130,-6.862) (9.179,-6.805) (9.227,-6.747) (9.273,-6.688) (9.317,-6.628) (9.360,-6.566) (9.402,-6.504) (9.441,-6.440) (9.480,-6.376) (9.516,-6.310) (9.550,-6.243) (9.583,-6.176) (9.614,-6.108) (9.643,-6.039) (9.671,-5.969) (9.696,-5.898) (9.720,-5.827) (9.741,-5.755) (9.761,-5.683) (9.779,-5.610) (9.795,-5.537) (9.810,-5.463) (9.822,-5.389) (9.832,-5.315) (9.841,-5.240) (9.847,-5.166) (9.852,-5.091) (9.855,-5.016) (9.856,-4.941) (9.855,-4.866) (9.853,-4.791) (9.848,-4.716) (9.842,-4.641) (9.834,-4.567) (9.825,-4.492) (9.813,-4.418) (9.800,-4.344) (9.786,-4.271) (9.769,-4.198) (9.751,-4.125) (9.732,-4.052) (9.711,-3.980) (9.688,-3.909) (9.664,-3.838) (9.639,-3.767) (9.612,-3.697) (9.584,-3.628) (9.554,-3.559) (9.523,-3.490) (9.491,-3.423) (9.458,-3.356) (9.423,-3.289) (9.387,-3.223) (9.350,-3.158) (9.311,-3.094) (9.272,-3.030) (9.231,-2.967) (9.189,-2.905) (9.147,-2.843) (9.103,-2.782) (9.058,-2.722) (9.012,-2.663) (8.966,-2.604) (8.918,-2.546) (8.869,-2.489) (8.820,-2.432) (8.770,-2.377) (8.719,-2.322) (8.667,-2.268) (8.614,-2.214) (8.561,-2.162) (8.506,-2.110) (8.451,-2.059) (8.396,-2.009) (8.339,-1.959) (8.282,-1.910) (8.225,-1.862) (8.167,-1.815) (8.108,-1.769) (8.048,-1.723) (7.988,-1.678) (7.928,-1.634) (7.867,-1.590) (7.805,-1.547) (7.743,-1.505) (7.680,-1.464) (7.617,-1.424) (7.553,-1.384) (7.489,-1.345) (7.425,-1.307) (7.360,-1.269) (7.295,-1.232) (7.229,-1.196) (7.163,-1.160) (7.097,-1.125) (7.030,-1.091) (6.963,-1.058) (6.895,-1.025) (6.828,-0.993) (6.759,-0.962) (6.691,-0.931) (6.622,-0.901) (6.553,-0.871) (6.484,-0.842) (6.415,-0.814) (6.345,-0.787) (6.275,-0.760) (6.205,-0.733) (6.134,-0.708) (6.064,-0.682) (5.993,-0.658) (5.922,-0.634) (5.850,-0.611) (5.779,-0.588) (5.707,-0.566) (5.635,-0.544) (5.563,-0.523) (5.491,-0.503) (5.419,-0.483) (5.347,-0.463) (5.274,-0.444) (5.201,-0.426) (5.128,-0.408) (5.055,-0.391) (4.982,-0.374) (4.909,-0.357) (4.836,-0.341) (4.763,-0.326) (4.689,-0.311) (4.615,-0.297) (4.542,-0.283) (4.468,-0.269) (4.394,-0.256) (4.320,-0.243) (4.246,-0.231)}
\LC{(-0.493,3.308) (-0.515,3.380) (-0.539,3.451) (-0.562,3.522) (-0.586,3.593) (-0.611,3.664) (-0.636,3.735) (-0.662,3.805) (-0.688,3.876) (-0.714,3.946) (-0.741,4.016) (-0.769,4.085) (-0.797,4.155) (-0.825,4.224) (-0.855,4.293) (-0.884,4.362) (-0.914,4.431) (-0.945,4.500) (-0.976,4.568) (-1.007,4.636) (-1.039,4.704) (-1.071,4.771) (-1.104,4.839) (-1.138,4.906) (-1.172,4.973) (-1.206,5.039) (-1.241,5.106) (-1.276,5.172) (-1.312,5.238) (-1.349,5.303) (-1.385,5.369) (-1.423,5.434) (-1.461,5.499) (-1.499,5.563) (-1.538,5.627) (-1.577,5.691) (-1.616,5.755) (-1.657,5.818) (-1.697,5.881) (-1.738,5.944) (-1.780,6.006) (-1.822,6.069) (-1.865,6.130) (-1.908,6.192) (-1.951,6.253) (-1.995,6.314) (-2.039,6.374) (-2.084,6.434) (-2.130,6.494) (-2.176,6.553) (-2.222,6.612) (-2.269,6.671) (-2.316,6.729) (-2.363,6.787) (-2.412,6.845) (-2.460,6.902) (-2.509,6.958) (-2.559,7.015) (-2.609,7.071) (-2.659,7.126) (-2.710,7.181) (-2.761,7.236) (-2.813,7.290) (-2.865,7.344) (-2.918,7.398) (-2.971,7.451) (-3.024,7.503) (-3.078,7.555) (-3.133,7.607) (-3.188,7.658) (-3.243,7.709) (-3.299,7.759) (-3.355,7.809) (-3.411,7.858) (-3.468,7.907) (-3.526,7.955) (-3.584,8.003) (-3.642,8.050) (-3.701,8.097) (-3.760,8.143) (-3.819,8.188) (-3.879,8.234) (-3.939,8.278) (-4.000,8.322) (-4.061,8.366) (-4.123,8.409) (-4.185,8.451) (-4.247,8.493) (-4.310,8.534) (-4.373,8.574) (-4.436,8.614) (-4.500,8.654) (-4.564,8.693) (-4.629,8.731) (-4.694,8.768) (-4.759,8.805) (-4.825,8.841) (-4.891,8.877) (-4.957,8.912) (-5.024,8.946) (-5.091,8.980) (-5.158,9.012) (-5.226,9.045) (-5.294,9.076) (-5.363,9.107) (-5.431,9.137) (-5.500,9.166) (-5.570,9.195) (-5.639,9.223) (-5.709,9.250) (-5.779,9.276) (-5.850,9.302) (-5.921,9.327) (-5.992,9.351) (-6.063,9.374) (-6.135,9.396) (-6.206,9.418) (-6.278,9.439) (-6.351,9.459) (-6.423,9.478) (-6.496,9.496) (-6.569,9.513) (-6.642,9.530) (-6.715,9.546) (-6.789,9.560) (-6.863,9.574) (-6.937,9.587) (-7.011,9.599) (-7.085,9.610) (-7.159,9.620) (-7.234,9.630) (-7.308,9.638) (-7.383,9.645) (-7.457,9.652) (-7.532,9.657) (-7.607,9.661) (-7.682,9.664) (-7.757,9.667) (-7.832,9.668) (-7.907,9.668) (-7.982,9.667) (-8.057,9.665) (-8.132,9.662) (-8.207,9.658) (-8.282,9.653) (-8.356,9.647) (-8.431,9.640) (-8.506,9.631) (-8.580,9.622) (-8.654,9.611) (-8.728,9.599) (-8.802,9.586) (-8.876,9.572) (-8.949,9.556) (-9.022,9.540) (-9.095,9.522) (-9.168,9.503) (-9.240,9.483) (-9.312,9.462) (-9.383,9.439) (-9.455,9.416) (-9.525,9.391) (-9.596,9.365) (-9.665,9.337) (-9.735,9.308) (-9.803,9.279) (-9.872,9.248) (-9.939,9.215) (-10.000,9.182)}
\LC{(-4.639,-7.914) (-4.708,-7.944) (-4.777,-7.973) (-4.846,-8.002) (-4.916,-8.030) (-4.986,-8.057) (-5.056,-8.083) (-5.127,-8.108) (-5.198,-8.132) (-5.269,-8.156) (-5.341,-8.178) (-5.413,-8.199) (-5.485,-8.220) (-5.557,-8.240) (-5.630,-8.258) (-5.703,-8.276) (-5.776,-8.292) (-5.849,-8.308) (-5.923,-8.323) (-5.997,-8.336) (-6.071,-8.349) (-6.145,-8.360) (-6.219,-8.370) (-6.293,-8.380) (-6.368,-8.388) (-6.443,-8.395) (-6.517,-8.401) (-6.592,-8.405) (-6.667,-8.409) (-6.742,-8.411) (-6.817,-8.413) (-6.892,-8.413) (-6.967,-8.412) (-7.042,-8.409) (-7.117,-8.406) (-7.192,-8.401) (-7.267,-8.395) (-7.341,-8.387) (-7.416,-8.379) (-7.490,-8.369) (-7.564,-8.357) (-7.638,-8.345) (-7.712,-8.331) (-7.785,-8.315) (-7.858,-8.299) (-7.931,-8.281) (-8.004,-8.261) (-8.076,-8.241) (-8.147,-8.219) (-8.219,-8.195) (-8.289,-8.170) (-8.360,-8.144) (-8.429,-8.116) (-8.498,-8.087) (-8.567,-8.056) (-8.635,-8.024) (-8.702,-7.991) (-8.768,-7.956) (-8.834,-7.920) (-8.899,-7.882) (-8.963,-7.843) (-9.026,-7.802) (-9.088,-7.761) (-9.149,-7.717) (-9.209,-7.673) (-9.269,-7.627) (-9.327,-7.579) (-9.384,-7.530) (-9.440,-7.480) (-9.494,-7.429) (-9.548,-7.376) (-9.600,-7.322) (-9.650,-7.267) (-9.700,-7.211) (-9.748,-7.153) (-9.795,-7.094) (-9.840,-7.035) (-9.883,-6.974) (-9.926,-6.911) (-9.966,-6.848) (-10.000,-6.784)}
\LC{(-10.000,-1.330) (-9.940,-1.301) (-9.870,-1.272) (-9.801,-1.244) (-9.731,-1.216) (-9.661,-1.189) (-9.591,-1.162) (-9.521,-1.136) (-9.451,-1.109) (-9.380,-1.084) (-9.310,-1.058) (-9.239,-1.033) (-9.168,-1.009) (-9.097,-0.985) (-9.026,-0.961) (-8.955,-0.938) (-8.884,-0.915) (-8.812,-0.892) (-8.740,-0.870) (-8.669,-0.848) (-8.597,-0.826) (-8.525,-0.805) (-8.453,-0.784) (-8.381,-0.764) (-8.308,-0.744) (-8.236,-0.724) (-8.164,-0.705) (-8.091,-0.686) (-8.018,-0.667) (-7.946,-0.649) (-7.873,-0.631) (-7.800,-0.613) (-7.727,-0.596) (-7.654,-0.579) (-7.581,-0.562) (-7.508,-0.546) (-7.434,-0.530) (-7.361,-0.514) (-7.288,-0.499) (-7.214,-0.483) (-7.141,-0.469) (-7.067,-0.454) (-6.993,-0.440) (-6.920,-0.426) (-6.846,-0.413) (-6.772,-0.399) (-6.698,-0.386) (-6.624,-0.374) (-6.550,-0.361) (-6.476,-0.349) (-6.402,-0.337) (-6.328,-0.325) (-6.279,-0.318)}
\LC{(1.716,4.522) (1.768,4.576) (1.820,4.629) (1.874,4.682) (1.928,4.734) (1.982,4.785) (2.038,4.836) (2.094,4.886) (2.151,4.935) (2.208,4.983) (2.266,5.030) (2.325,5.077) (2.384,5.123) (2.445,5.167) (2.505,5.211) (2.567,5.254) (2.629,5.297) (2.691,5.338) (2.755,5.378) (2.819,5.417) (2.883,5.456) (2.948,5.493) (3.014,5.529) (3.080,5.565) (3.147,5.599) (3.214,5.632) (3.282,5.664) (3.350,5.695) (3.419,5.725) (3.488,5.753) (3.558,5.780) (3.628,5.807) (3.699,5.831) (3.770,5.855) (3.842,5.877) (3.914,5.898) (3.986,5.918) (4.059,5.937) (4.132,5.954) (4.205,5.969) (4.279,5.983) (4.353,5.996) (4.427,6.007) (4.502,6.017) (4.576,6.025) (4.651,6.032) (4.726,6.037) (4.801,6.040) (4.876,6.042) (4.951,6.042) (5.026,6.040) (5.100,6.037) (5.175,6.032) (5.250,6.025) (5.324,6.017) (5.399,6.006) (5.473,5.994) (5.546,5.980) (5.620,5.964) (5.693,5.946) (5.765,5.926) (5.837,5.905) (5.908,5.881) (5.978,5.856) (6.048,5.828) (6.117,5.799) (6.185,5.767) (6.252,5.734) (6.318,5.698) (6.383,5.661) (6.447,5.621) (6.510,5.580) (6.571,5.537) (6.631,5.492) (6.689,5.445) (6.746,5.396) (6.801,5.345) (6.855,5.292) (6.906,5.238) (6.956,5.182) (7.004,5.124) (7.050,5.065) (7.094,5.004) (7.135,4.941) (7.175,4.878) (7.212,4.813) (7.247,4.746) (7.280,4.679) (7.310,4.610) (7.338,4.540) (7.363,4.470) (7.386,4.398) (7.406,4.326) (7.424,4.253) (7.440,4.180) (7.453,4.106) (7.463,4.032) (7.471,3.957) (7.477,3.883) (7.480,3.808) (7.480,3.733) (7.479,3.658) (7.474,3.583) (7.468,3.508) (7.459,3.434) (7.448,3.359) (7.435,3.286) (7.420,3.212) (7.402,3.139) (7.383,3.067) (7.361,2.995) (7.338,2.924) (7.312,2.853) (7.285,2.783) (7.256,2.714) (7.225,2.646) (7.192,2.578) (7.158,2.512) (7.122,2.446) (7.084,2.381) (7.045,2.317) (7.005,2.254) (6.963,2.192) (6.920,2.130) (6.875,2.070) (6.829,2.011) (6.782,1.952) (6.734,1.895) (6.684,1.839) (6.634,1.783) (6.582,1.729) (6.529,1.676) (6.475,1.623) (6.421,1.572) (6.365,1.522) (6.309,1.472) (6.251,1.424) (6.193,1.377) (6.134,1.331) (6.074,1.285) (6.014,1.241) (5.952,1.198) (5.891,1.155) (5.828,1.114) (5.765,1.073) (5.701,1.034) (5.680,1.021)}
\LC{(0.124,-3.817) (0.129,-3.892) (0.134,-3.966) (0.139,-4.041) (0.144,-4.116) (0.150,-4.191) (0.155,-4.266) (0.161,-4.340) (0.166,-4.415) (0.172,-4.490) (0.178,-4.565) (0.184,-4.640) (0.190,-4.714) (0.196,-4.789) (0.202,-4.864) (0.208,-4.938) (0.215,-5.013) (0.221,-5.088) (0.228,-5.163) (0.235,-5.237) (0.242,-5.312) (0.248,-5.387) (0.256,-5.461) (0.263,-5.536) (0.270,-5.611) (0.277,-5.685) (0.284,-5.760) (0.292,-5.835) (0.300,-5.909) (0.307,-5.984) (0.315,-6.058) (0.323,-6.133) (0.331,-6.208) (0.339,-6.282) (0.347,-6.357) (0.355,-6.431) (0.364,-6.506) (0.372,-6.580) (0.381,-6.655) (0.389,-6.729) (0.398,-6.804) (0.407,-6.878) (0.416,-6.953) (0.425,-7.027) (0.434,-7.102) (0.443,-7.176) (0.453,-7.250) (0.462,-7.325) (0.472,-7.399) (0.481,-7.474) (0.491,-7.548) (0.501,-7.622) (0.511,-7.697) (0.521,-7.771) (0.531,-7.845) (0.541,-7.920) (0.551,-7.994) (0.562,-8.068) (0.572,-8.143) (0.583,-8.217) (0.593,-8.291) (0.604,-8.365) (0.615,-8.439) (0.626,-8.514) (0.637,-8.588) (0.648,-8.662) (0.659,-8.736) (0.671,-8.810) (0.682,-8.884) (0.694,-8.958) (0.705,-9.033) (0.717,-9.107) (0.729,-9.181) (0.741,-9.255) (0.753,-9.329) (0.765,-9.403) (0.777,-9.477) (0.790,-9.551) (0.802,-9.625) (0.815,-9.699) (0.827,-9.773) (0.840,-9.847) (0.853,-9.920) (0.866,-9.994) (0.870,-10.000)}
\LC{(-0.271,-4.650) (-0.280,-4.724) (-0.289,-4.799) (-0.298,-4.873) (-0.307,-4.948) (-0.316,-5.022) (-0.326,-5.097) (-0.336,-5.171) (-0.346,-5.245) (-0.355,-5.320) (-0.366,-5.394) (-0.376,-5.468) (-0.386,-5.542) (-0.397,-5.617) (-0.408,-5.691) (-0.418,-5.765) (-0.429,-5.839) (-0.441,-5.914) (-0.452,-5.988) (-0.463,-6.062) (-0.475,-6.136) (-0.486,-6.210) (-0.498,-6.284) (-0.510,-6.358) (-0.522,-6.432) (-0.535,-6.506) (-0.547,-6.580) (-0.560,-6.654) (-0.572,-6.728) (-0.585,-6.802) (-0.598,-6.876) (-0.611,-6.950) (-0.624,-7.023) (-0.638,-7.097) (-0.651,-7.171) (-0.665,-7.245) (-0.679,-7.318) (-0.693,-7.392) (-0.707,-7.466) (-0.721,-7.539) (-0.735,-7.613) (-0.750,-7.687) (-0.764,-7.760) (-0.779,-7.834) (-0.794,-7.907) (-0.809,-7.981) (-0.824,-8.054) (-0.840,-8.128) (-0.855,-8.201) (-0.871,-8.274) (-0.887,-8.348) (-0.903,-8.421) (-0.919,-8.494) (-0.935,-8.567) (-0.951,-8.641) (-0.968,-8.714) (-0.984,-8.787) (-1.001,-8.860) (-1.018,-8.933) (-1.035,-9.006) (-1.052,-9.079) (-1.070,-9.152) (-1.087,-9.225) (-1.105,-9.298) (-1.122,-9.371) (-1.140,-9.444) (-1.158,-9.516) (-1.176,-9.589) (-1.195,-9.662) (-1.213,-9.735) (-1.232,-9.807) (-1.250,-9.880) (-1.269,-9.952) (-1.282,-10.000)}
\LC{(2.185,-4.584) (2.247,-4.628) (2.308,-4.670) (2.371,-4.712) (2.434,-4.752) (2.498,-4.791) (2.563,-4.829) (2.628,-4.866) (2.694,-4.901) (2.761,-4.936) (2.828,-4.969) (2.896,-5.001) (2.964,-5.032) (3.033,-5.061) (3.103,-5.090) (3.173,-5.116) (3.243,-5.142) (3.315,-5.166) (3.386,-5.188) (3.458,-5.209) (3.530,-5.228) (3.603,-5.246) (3.677,-5.263) (3.750,-5.278) (3.824,-5.291) (3.898,-5.302) (3.972,-5.312) (4.047,-5.320) (4.122,-5.326) (4.197,-5.331) (4.271,-5.333) (4.346,-5.334) (4.421,-5.333) (4.496,-5.330) (4.571,-5.325) (4.646,-5.317) (4.720,-5.308) (4.795,-5.297) (4.868,-5.284) (4.942,-5.269) (5.015,-5.251) (5.087,-5.231) (5.159,-5.210) (5.230,-5.185) (5.300,-5.159) (5.369,-5.131) (5.438,-5.100) (5.505,-5.067) (5.571,-5.032) (5.636,-4.994) (5.700,-4.955) (5.762,-4.913) (5.823,-4.869) (5.882,-4.822) (5.939,-4.774) (5.995,-4.724) (6.048,-4.671) (6.100,-4.616) (6.149,-4.560) (6.196,-4.502) (6.241,-4.442) (6.284,-4.380) (6.324,-4.316) (6.361,-4.251) (6.396,-4.185) (6.428,-4.117) (6.458,-4.048) (6.485,-3.978) (6.509,-3.907) (6.530,-3.835) (6.548,-3.763) (6.564,-3.689) (6.577,-3.615) (6.586,-3.541) (6.594,-3.466) (6.598,-3.392) (6.599,-3.317) (6.598,-3.242) (6.594,-3.167) (6.588,-3.092) (6.579,-3.018) (6.567,-2.943) (6.553,-2.870) (6.537,-2.797) (6.518,-2.724) (6.497,-2.652) (6.473,-2.581) (6.448,-2.510) (6.420,-2.441) (6.390,-2.372) (6.358,-2.304) (6.325,-2.237) (6.289,-2.171) (6.252,-2.106) (6.213,-2.042) (6.172,-1.979) (6.130,-1.917) (6.086,-1.856) (6.041,-1.796) (5.994,-1.738) (5.946,-1.680) (5.896,-1.624) (5.846,-1.568) (5.794,-1.514) (5.741,-1.461) (5.687,-1.409) (5.631,-1.359) (5.575,-1.309) (5.518,-1.261) (5.459,-1.213) (5.400,-1.167) (5.340,-1.122) (5.320,-1.107)}
\LC{(-2.347,6.032) (-2.399,6.085) (-2.453,6.138) (-2.506,6.191) (-2.561,6.242) (-2.615,6.294) (-2.671,6.344) (-2.727,6.394) (-2.783,6.444) (-2.840,6.493) (-2.897,6.541) (-2.955,6.589) (-3.013,6.636) (-3.072,6.682) (-3.132,6.728) (-3.192,6.773) (-3.252,6.817) (-3.313,6.861) (-3.374,6.904) (-3.436,6.947) (-3.499,6.988) (-3.561,7.029) (-3.625,7.069) (-3.689,7.109) (-3.753,7.148) (-3.817,7.186) (-3.882,7.223) (-3.948,7.259) (-4.014,7.295) (-4.080,7.330) (-4.147,7.364) (-4.214,7.397) (-4.282,7.430) (-4.350,7.461) (-4.419,7.492) (-4.487,7.522) (-4.557,7.551) (-4.626,7.579) (-4.696,7.606) (-4.766,7.632) (-4.837,7.657) (-4.908,7.681) (-4.979,7.705) (-5.051,7.727) (-5.123,7.748) (-5.195,7.769) (-5.267,7.788) (-5.340,7.806) (-5.413,7.824) (-5.486,7.840) (-5.560,7.855) (-5.633,7.869) (-5.707,7.882) (-5.781,7.894) (-5.856,7.904) (-5.930,7.914) (-6.005,7.922) (-6.079,7.929) (-6.154,7.935) (-6.229,7.940) (-6.304,7.943) (-6.379,7.946) (-6.454,7.947) (-6.529,7.946) (-6.604,7.945) (-6.679,7.942) (-6.754,7.938) (-6.828,7.933) (-6.903,7.926) (-6.978,7.918) (-7.052,7.908) (-7.126,7.897) (-7.200,7.885) (-7.274,7.871) (-7.347,7.856) (-7.421,7.839) (-7.493,7.821) (-7.566,7.802) (-7.638,7.781) (-7.709,7.759) (-7.780,7.735) (-7.851,7.709) (-7.921,7.682) (-7.990,7.654) (-8.059,7.624) (-8.127,7.593) (-8.195,7.560) (-8.262,7.526) (-8.327,7.490) (-8.392,7.452) (-8.457,7.414) (-8.520,7.373) (-8.582,7.331) (-8.643,7.288) (-8.703,7.243) (-8.762,7.197) (-8.820,7.149) (-8.877,7.100) (-8.933,7.050) (-8.987,6.998) (-9.040,6.945) (-9.091,6.890) (-9.141,6.834) (-9.190,6.777) (-9.237,6.719) (-9.282,6.659) (-9.326,6.599) (-9.369,6.537) (-9.410,6.474) (-9.449,6.410) (-9.486,6.345) (-9.522,6.279) (-9.556,6.212) (-9.588,6.144) (-9.618,6.076) (-9.646,6.006) (-9.673,5.936) (-9.698,5.865) (-9.721,5.794) (-9.742,5.722) (-9.761,5.649) (-9.778,5.576) (-9.793,5.503) (-9.807,5.429) (-9.818,5.355) (-9.828,5.281) (-9.835,5.206) (-9.841,5.131) (-9.845,5.056) (-9.847,4.981) (-9.848,4.906) (-9.846,4.831) (-9.843,4.756) (-9.838,4.682) (-9.831,4.607) (-9.822,4.532) (-9.812,4.458) (-9.800,4.384) (-9.786,4.310) (-9.771,4.237) (-9.754,4.164) (-9.735,4.091) (-9.715,4.019) (-9.694,3.947) (-9.671,3.876) (-9.646,3.805) (-9.620,3.735) (-9.593,3.665) (-9.564,3.596) (-9.534,3.527) (-9.502,3.459) (-9.469,3.391) (-9.435,3.325) (-9.400,3.258) (-9.363,3.193) (-9.326,3.128) (-9.287,3.064) (-9.247,3.001) (-9.206,2.938) (-9.163,2.876) (-9.120,2.814) (-9.076,2.754) (-9.031,2.694) (-8.985,2.635) (-8.938,2.577) (-8.890,2.519) (-8.841,2.462) (-8.791,2.406) (-8.740,2.351) (-8.689,2.296) (-8.636,2.242) (-8.583,2.189) (-8.530,2.137) (-8.475,2.086) (-8.420,2.035) (-8.364,1.985) (-8.307,1.936) (-8.250,1.887) (-8.192,1.840) (-8.134,1.793) (-8.074,1.747) (-8.015,1.701) (-7.954,1.657) (-7.894,1.613) (-7.832,1.570) (-7.770,1.527) (-7.708,1.485) (-7.645,1.445) (-7.582,1.404) (-7.518,1.365) (-7.454,1.326) (-7.389,1.288) (-7.324,1.251) (-7.259,1.214) (-7.193,1.178) (-7.127,1.143) (-7.060,1.109) (-6.993,1.075) (-6.926,1.042) (-6.858,1.009) (-6.790,0.977) (-6.722,0.946) (-6.654,0.916) (-6.585,0.886) (-6.516,0.857) (-6.469,0.838)}
\LC{(10.000,-1.735) (9.939,-1.698) (9.874,-1.661) (9.808,-1.625) (9.742,-1.589) (9.676,-1.554) (9.610,-1.519) (9.543,-1.485) (9.476,-1.451) (9.409,-1.418) (9.341,-1.385) (9.273,-1.353) (9.205,-1.321) (9.137,-1.290) (9.069,-1.260) (9.000,-1.229) (8.931,-1.200) (8.862,-1.170) (8.793,-1.142) (8.723,-1.113) (8.654,-1.085) (8.584,-1.058) (8.514,-1.031) (8.444,-1.005) (8.373,-0.979) (8.303,-0.953) (8.232,-0.928) (8.161,-0.903) (8.091,-0.879) (8.019,-0.855) (7.948,-0.832) (7.877,-0.809) (7.805,-0.786) (7.734,-0.764) (7.662,-0.742) (7.590,-0.721) (7.518,-0.700) (7.446,-0.680) (7.373,-0.660) (7.301,-0.640) (7.229,-0.621) (7.156,-0.602) (7.083,-0.583) (7.010,-0.565) (6.938,-0.548)}
\LC{(4.774,-8.871) (4.840,-8.907) (4.906,-8.943) (4.972,-8.978) (5.038,-9.013) (5.105,-9.047) (5.173,-9.080) (5.240,-9.112) (5.308,-9.144) (5.376,-9.175) (5.445,-9.206) (5.514,-9.236) (5.583,-9.265) (5.652,-9.293) (5.722,-9.320) (5.792,-9.347) (5.862,-9.373) (5.933,-9.399) (6.004,-9.423) (6.075,-9.447) (6.147,-9.470) (6.218,-9.492) (6.290,-9.513) (6.362,-9.534) (6.435,-9.553) (6.507,-9.572) (6.580,-9.590) (6.653,-9.607) (6.726,-9.624) (6.800,-9.639) (6.873,-9.653) (6.947,-9.667) (7.021,-9.680) (7.095,-9.692) (7.169,-9.702) (7.244,-9.712) (7.318,-9.721) (7.393,-9.729) (7.467,-9.736) (7.542,-9.743) (7.617,-9.748) (7.692,-9.752) (7.767,-9.755) (7.842,-9.757) (7.917,-9.758) (7.992,-9.758) (8.067,-9.757) (8.142,-9.755) (8.216,-9.752) (8.291,-9.748) (8.366,-9.743) (8.441,-9.736) (8.516,-9.729) (8.590,-9.720) (8.664,-9.711) (8.739,-9.700) (8.813,-9.688) (8.887,-9.675) (8.960,-9.661) (9.034,-9.645) (9.107,-9.629) (9.180,-9.611) (9.252,-9.592) (9.324,-9.572) (9.396,-9.551) (9.468,-9.528) (9.539,-9.505) (9.610,-9.480) (9.680,-9.454) (9.750,-9.426) (9.819,-9.398) (9.888,-9.368) (9.956,-9.337) (10.000,-9.316)}
\LC{(2.556,-5.601) (2.615,-5.647) (2.675,-5.692) (2.736,-5.736) (2.797,-5.780) (2.858,-5.823) (2.920,-5.865) (2.983,-5.906) (3.046,-5.946) (3.110,-5.986) (3.174,-6.024) (3.239,-6.062) (3.305,-6.098) (3.371,-6.134) (3.437,-6.169) (3.504,-6.203) (3.571,-6.236) (3.639,-6.268) (3.708,-6.298) (3.777,-6.328) (3.846,-6.357) (3.915,-6.385) (3.986,-6.411) (4.056,-6.437) (4.127,-6.461) (4.198,-6.484) (4.270,-6.506) (4.342,-6.527) (4.415,-6.546) (4.487,-6.565) (4.560,-6.582) (4.634,-6.598) (4.707,-6.612) (4.781,-6.626) (4.855,-6.638) (4.929,-6.648) (5.004,-6.657) (5.078,-6.665) (5.153,-6.672) (5.228,-6.677) (5.303,-6.680) (5.378,-6.682) (5.453,-6.683) (5.528,-6.682) (5.603,-6.679) (5.678,-6.675) (5.752,-6.670) (5.827,-6.662) (5.902,-6.653) (5.976,-6.643) (6.050,-6.631) (6.124,-6.617) (6.197,-6.601) (6.270,-6.584) (6.342,-6.565) (6.415,-6.545) (6.486,-6.522) (6.557,-6.498) (6.627,-6.472) (6.697,-6.444) (6.766,-6.415) (6.834,-6.383) (6.902,-6.350) (6.968,-6.315) (7.033,-6.279) (7.098,-6.240) (7.161,-6.200) (7.223,-6.158) (7.284,-6.114) (7.344,-6.069) (7.402,-6.022) (7.459,-5.973) (7.514,-5.922) (7.568,-5.870) (7.620,-5.816) (7.671,-5.761) (7.720,-5.704) (7.767,-5.646) (7.812,-5.586) (7.856,-5.525) (7.897,-5.462) (7.937,-5.399) (7.974,-5.334) (8.010,-5.267) (8.043,-5.200) (8.074,-5.132) (8.103,-5.063) (8.130,-4.993) (8.154,-4.922) (8.177,-4.850) (8.197,-4.778) (8.215,-4.705) (8.230,-4.632) (8.244,-4.558) (8.255,-4.484) (8.264,-4.410) (8.270,-4.335) (8.275,-4.260) (8.277,-4.185) (8.277,-4.110) (8.275,-4.035) (8.270,-3.960) (8.264,-3.885) (8.255,-3.811) (8.245,-3.737) (8.232,-3.663) (8.218,-3.589) (8.201,-3.516) (8.183,-3.443) (8.163,-3.371) (8.141,-3.299) (8.117,-3.228) (8.092,-3.158) (8.064,-3.088) (8.036,-3.018) (8.005,-2.950) (7.973,-2.882) (7.940,-2.815) (7.905,-2.749) (7.868,-2.683) (7.830,-2.618) (7.791,-2.555) (7.750,-2.491) (7.709,-2.429) (7.666,-2.368) (7.621,-2.307) (7.576,-2.248) (7.529,-2.189) (7.481,-2.131) (7.433,-2.074) (7.383,-2.018) (7.332,-1.963) (7.280,-1.909) (7.227,-1.855) (7.174,-1.803) (7.119,-1.751) (7.064,-1.701) (7.008,-1.651) (6.951,-1.602) (6.893,-1.555) (6.835,-1.508) (6.815,-1.492)}
\LC{(-0.782,5.205) (-0.804,5.277) (-0.827,5.348) (-0.850,5.419) (-0.873,5.491) (-0.897,5.562) (-0.921,5.633) (-0.945,5.704) (-0.970,5.775) (-0.995,5.845) (-1.020,5.916) (-1.046,5.987) (-1.072,6.057) (-1.098,6.127) (-1.125,6.197) (-1.151,6.267) (-1.179,6.337) (-1.206,6.407) (-1.234,6.477) (-1.262,6.546) (-1.290,6.616) (-1.319,6.685) (-1.348,6.754) (-1.377,6.823) (-1.407,6.892) (-1.437,6.961) (-1.467,7.029) (-1.498,7.098) (-1.529,7.166) (-1.560,7.234) (-1.591,7.303) (-1.623,7.370) (-1.655,7.438) (-1.688,7.506) (-1.720,7.573) (-1.754,7.641) (-1.787,7.708) (-1.821,7.775) (-1.854,7.842) (-1.889,7.908) (-1.923,7.975) (-1.958,8.041) (-1.993,8.108) (-2.029,8.174) (-2.065,8.240) (-2.101,8.305) (-2.137,8.371) (-2.174,8.436) (-2.211,8.501) (-2.248,8.566) (-2.286,8.631) (-2.324,8.696) (-2.362,8.761) (-2.401,8.825) (-2.440,8.889) (-2.479,8.953) (-2.518,9.017) (-2.558,9.080) (-2.598,9.144) (-2.638,9.207) (-2.679,9.270) (-2.720,9.333) (-2.761,9.395) (-2.803,9.458) (-2.845,9.520) (-2.887,9.582) (-2.929,9.644) (-2.972,9.706) (-3.015,9.767) (-3.059,9.828) (-3.102,9.889) (-3.146,9.950) (-3.191,10.000)}
\LC{(-0.789,-5.154) (-0.812,-5.226) (-0.836,-5.297) (-0.859,-5.368) (-0.883,-5.440) (-0.907,-5.511) (-0.932,-5.581) (-0.957,-5.652) (-0.982,-5.723) (-1.007,-5.793) (-1.033,-5.864) (-1.059,-5.934) (-1.086,-6.004) (-1.112,-6.074) (-1.139,-6.144) (-1.167,-6.214) (-1.195,-6.284) (-1.223,-6.353) (-1.251,-6.423) (-1.280,-6.492) (-1.309,-6.561) (-1.338,-6.630) (-1.368,-6.699) (-1.398,-6.768) (-1.428,-6.836) (-1.458,-6.905) (-1.489,-6.973) (-1.521,-7.041) (-1.552,-7.110) (-1.584,-7.177) (-1.616,-7.245) (-1.649,-7.313) (-1.682,-7.380) (-1.715,-7.447) (-1.748,-7.515) (-1.782,-7.582) (-1.816,-7.648) (-1.851,-7.715) (-1.885,-7.781) (-1.920,-7.848) (-1.956,-7.914) (-1.991,-7.980) (-2.027,-8.046) (-2.064,-8.111) (-2.100,-8.177) (-2.137,-8.242) (-2.174,-8.307) (-2.212,-8.372) (-2.250,-8.437) (-2.288,-8.501) (-2.327,-8.566) (-2.365,-8.630) (-2.404,-8.694) (-2.444,-8.758) (-2.484,-8.821) (-2.524,-8.885) (-2.564,-8.948) (-2.605,-9.011) (-2.646,-9.074) (-2.687,-9.136) (-2.729,-9.199) (-2.771,-9.261) (-2.813,-9.323) (-2.855,-9.385) (-2.898,-9.446) (-2.941,-9.508) (-2.985,-9.569) (-3.029,-9.630) (-3.073,-9.690) (-3.117,-9.751) (-3.162,-9.811) (-3.207,-9.871) (-3.252,-9.931) (-3.297,-9.990) (-3.313,-10.000)}
\LC{(-1.680,4.601) (-1.731,4.657) (-1.782,4.712) (-1.833,4.766) (-1.886,4.820) (-1.939,4.873) (-1.992,4.925) (-2.047,4.977) (-2.102,5.028) (-2.158,5.078) (-2.214,5.127) (-2.271,5.176) (-2.329,5.224) (-2.387,5.271) (-2.446,5.317) (-2.506,5.363) (-2.566,5.408) (-2.627,5.451) (-2.688,5.494) (-2.750,5.537) (-2.813,5.578) (-2.876,5.618) (-2.940,5.658) (-3.004,5.696) (-3.069,5.734) (-3.135,5.770) (-3.201,5.806) (-3.267,5.840) (-3.334,5.874) (-3.402,5.906) (-3.470,5.937) (-3.539,5.968) (-3.608,5.997) (-3.677,6.025) (-3.747,6.052) (-3.818,6.077) (-3.889,6.102) (-3.960,6.125) (-4.032,6.147) (-4.104,6.168) (-4.176,6.188) (-4.249,6.206) (-4.322,6.223) (-4.395,6.238) (-4.469,6.252) (-4.543,6.265) (-4.617,6.277) (-4.691,6.287) (-4.766,6.295) (-4.841,6.302) (-4.915,6.308) (-4.990,6.312) (-5.065,6.314) (-5.140,6.315) (-5.215,6.314) (-5.290,6.312) (-5.365,6.308) (-5.440,6.302) (-5.514,6.294) (-5.589,6.285) (-5.663,6.274) (-5.737,6.262) (-5.811,6.247) (-5.884,6.231) (-5.957,6.213) (-6.029,6.193) (-6.101,6.172) (-6.172,6.148) (-6.243,6.123) (-6.312,6.095) (-6.382,6.066) (-6.450,6.035) (-6.517,6.002) (-6.584,5.967) (-6.649,5.930) (-6.713,5.892) (-6.776,5.851) (-6.838,5.809) (-6.899,5.765) (-6.958,5.719) (-7.016,5.671) (-7.072,5.621) (-7.127,5.570) (-7.179,5.517) (-7.231,5.462) (-7.280,5.405) (-7.328,5.348) (-7.373,5.288) (-7.417,5.227) (-7.459,5.165) (-7.498,5.101) (-7.536,5.036) (-7.571,4.970) (-7.604,4.902) (-7.634,4.834) (-7.663,4.764) (-7.689,4.694) (-7.713,4.623) (-7.734,4.551) (-7.753,4.479) (-7.770,4.405) (-7.784,4.332) (-7.796,4.258) (-7.805,4.183) (-7.812,4.109) (-7.817,4.034) (-7.819,3.959) (-7.819,3.884) (-7.817,3.809) (-7.813,3.734) (-7.806,3.659) (-7.797,3.585) (-7.786,3.511) (-7.773,3.437) (-7.758,3.363) (-7.740,3.290) (-7.721,3.218) (-7.700,3.146) (-7.677,3.075) (-7.652,3.004) (-7.625,2.934) (-7.597,2.864) (-7.566,2.796) (-7.535,2.728) (-7.501,2.661) (-7.466,2.595) (-7.429,2.529) (-7.391,2.464) (-7.352,2.401) (-7.311,2.338) (-7.269,2.276) (-7.225,2.215) (-7.180,2.155) (-7.134,2.095) (-7.087,2.037) (-7.039,1.980) (-6.989,1.924) (-6.939,1.868) (-6.887,1.814) (-6.834,1.760) (-6.781,1.708) (-6.726,1.656) (-6.671,1.606) (-6.615,1.556) (-6.558,1.507) (-6.500,1.460) (-6.441,1.413) (-6.382,1.367) (-6.322,1.322) (-6.261,1.278) (-6.199,1.235) (-6.137,1.193) (-6.075,1.152) (-6.011,1.112) (-5.947,1.073) (-5.883,1.035) (-5.818,0.997) (-5.752,0.961) (-5.686,0.925) (-5.620,0.890) (-5.553,0.856) (-5.486,0.823) (-5.441,0.801)}
\LC{(10.000,-0.013) (9.943,-0.013) (9.868,-0.013) (9.793,-0.012) (9.718,-0.012) (9.643,-0.012) (9.568,-0.012) (9.493,-0.011) (9.418,-0.011) (9.343,-0.011) (9.268,-0.011) (9.193,-0.010) (9.118,-0.010) (9.043,-0.010) (8.968,-0.010) (8.893,-0.009) (8.818,-0.009) (8.743,-0.009) (8.668,-0.009) (8.593,-0.008) (8.518,-0.008) (8.443,-0.008) (8.368,-0.008) (8.293,-0.008) (8.218,-0.007) (8.143,-0.007) (8.068,-0.007) (7.993,-0.007) (7.918,-0.007) (7.843,-0.006) (7.768,-0.006) (7.693,-0.006) (7.618,-0.006) (7.543,-0.006) (7.468,-0.006) (7.393,-0.005) (7.318,-0.005) (7.243,-0.005) (7.168,-0.005) (7.093,-0.005) (7.018,-0.005) (6.943,-0.004) (6.868,-0.004) (6.793,-0.004) (6.718,-0.004) (6.643,-0.004) (6.568,-0.004) (6.493,-0.004) (6.418,-0.004) (6.343,-0.003) (6.268,-0.003) (6.193,-0.003) (6.118,-0.003) (6.043,-0.003) (5.968,-0.003) (5.893,-0.003) (5.818,-0.003) (5.743,-0.003) (5.668,-0.002) (5.593,-0.002) (5.518,-0.002) (5.443,-0.002) (5.368,-0.002)}
\LC{(-0.001,2.953) (-0.001,3.028) (-0.001,3.103) (-0.001,3.178) (-0.001,3.253) (-0.002,3.328) (-0.002,3.403) (-0.002,3.478) (-0.002,3.553) (-0.002,3.628) (-0.002,3.703) (-0.002,3.778) (-0.002,3.853) (-0.002,3.928) (-0.002,4.003) (-0.002,4.078) (-0.002,4.153) (-0.002,4.228) (-0.003,4.303) (-0.003,4.378) (-0.003,4.453) (-0.003,4.528) (-0.003,4.603) (-0.003,4.678) (-0.003,4.753) (-0.003,4.828) (-0.003,4.903) (-0.003,4.978) (-0.004,5.053) (-0.004,5.128) (-0.004,5.203) (-0.004,5.278) (-0.004,5.353) (-0.004,5.428) (-0.004,5.503) (-0.004,5.578) (-0.004,5.653) (-0.005,5.728) (-0.005,5.803) (-0.005,5.878) (-0.005,5.953) (-0.005,6.028) (-0.005,6.103) (-0.005,6.178) (-0.005,6.253) (-0.006,6.328) (-0.006,6.403) (-0.006,6.478) (-0.006,6.553) (-0.006,6.628) (-0.006,6.703) (-0.006,6.778) (-0.007,6.853) (-0.007,6.928) (-0.007,7.003) (-0.007,7.078) (-0.007,7.153) (-0.007,7.228) (-0.008,7.303) (-0.008,7.378) (-0.008,7.453) (-0.008,7.528) (-0.008,7.603) (-0.008,7.678) (-0.008,7.753) (-0.009,7.828) (-0.009,7.903) (-0.009,7.978) (-0.009,8.053) (-0.009,8.128) (-0.009,8.203) (-0.010,8.278) (-0.010,8.353) (-0.010,8.428) (-0.010,8.503) (-0.010,8.578) (-0.011,8.653) (-0.011,8.728) (-0.011,8.803) (-0.011,8.878) (-0.011,8.953) (-0.011,9.028) (-0.012,9.103) (-0.012,9.178) (-0.012,9.253) (-0.012,9.328) (-0.012,9.403) (-0.013,9.478) (-0.013,9.553) (-0.013,9.628) (-0.013,9.703) (-0.013,9.778) (-0.014,9.853) (-0.014,9.928) (-0.014,10.000)}
\LC{(3.943,7.805) (4.006,7.846) (4.070,7.886) (4.134,7.925) (4.198,7.963) (4.263,8.001) (4.328,8.038) (4.393,8.075) (4.459,8.111) (4.526,8.146) (4.592,8.180) (4.659,8.213) (4.727,8.246) (4.795,8.278) (4.863,8.309) (4.931,8.340) (5.000,8.370) (5.070,8.399) (5.139,8.427) (5.209,8.454) (5.279,8.480) (5.350,8.506) (5.420,8.531) (5.492,8.554) (5.563,8.577) (5.635,8.600) (5.706,8.621) (5.779,8.641) (5.851,8.661) (5.924,8.679) (5.997,8.697) (6.070,8.713) (6.143,8.729) (6.217,8.743) (6.291,8.757) (6.364,8.770) (6.439,8.781) (6.513,8.792) (6.587,8.802) (6.662,8.810) (6.736,8.818) (6.811,8.824) (6.886,8.830) (6.961,8.834) (7.036,8.837) (7.111,8.839) (7.186,8.840) (7.261,8.840) (7.336,8.838) (7.411,8.836) (7.485,8.832) (7.560,8.827) (7.635,8.821) (7.710,8.814) (7.784,8.805) (7.859,8.795) (7.933,8.784) (8.007,8.772) (8.080,8.758) (8.154,8.743) (8.227,8.727) (8.300,8.710) (8.373,8.691) (8.445,8.671) (8.517,8.650) (8.588,8.627) (8.660,8.603) (8.730,8.578) (8.800,8.551) (8.870,8.523) (8.939,8.494) (9.007,8.463) (9.075,8.431) (9.142,8.398) (9.209,8.363) (9.274,8.327) (9.339,8.289) (9.404,8.251) (9.467,8.210) (9.529,8.169) (9.591,8.126) (9.652,8.082) (9.711,8.037) (9.770,7.990) (9.827,7.942) (9.884,7.892) (9.939,7.842) (9.994,7.790) (10.000,7.773)}
\LC{(-0.905,-2.753) (-0.953,-2.811) (-1.001,-2.869) (-1.050,-2.925) (-1.100,-2.981) (-1.152,-3.036) (-1.204,-3.089) (-1.258,-3.142) (-1.312,-3.193) (-1.368,-3.243) (-1.425,-3.293) (-1.482,-3.341) (-1.541,-3.387) (-1.601,-3.433) (-1.661,-3.477) (-1.723,-3.520) (-1.785,-3.561) (-1.849,-3.601) (-1.913,-3.639) (-1.978,-3.676) (-2.045,-3.712) (-2.111,-3.746) (-2.179,-3.778) (-2.248,-3.808) (-2.317,-3.837) (-2.387,-3.864) (-2.458,-3.889) (-2.529,-3.912) (-2.601,-3.934) (-2.673,-3.953) (-2.746,-3.970) (-2.820,-3.985) (-2.894,-3.998) (-2.968,-4.009) (-3.042,-4.017) (-3.117,-4.023) (-3.192,-4.027) (-3.267,-4.028) (-3.342,-4.027) (-3.417,-4.023) (-3.492,-4.016) (-3.566,-4.007) (-3.640,-3.995) (-3.714,-3.981) (-3.787,-3.963) (-3.859,-3.943) (-3.930,-3.919) (-4.000,-3.893) (-4.069,-3.864) (-4.137,-3.832) (-4.203,-3.797) (-4.268,-3.759) (-4.331,-3.718) (-4.392,-3.674) (-4.450,-3.627) (-4.506,-3.577) (-4.560,-3.525) (-4.611,-3.470) (-4.659,-3.412) (-4.704,-3.353) (-4.746,-3.290) (-4.785,-3.226) (-4.820,-3.160) (-4.852,-3.092) (-4.880,-3.022) (-4.905,-2.952) (-4.925,-2.880) (-4.942,-2.807) (-4.956,-2.733) (-4.965,-2.658) (-4.971,-2.584) (-4.973,-2.509) (-4.972,-2.434) (-4.967,-2.359) (-4.958,-2.284) (-4.946,-2.210) (-4.931,-2.137) (-4.913,-2.064) (-4.892,-1.992) (-4.867,-1.921) (-4.840,-1.851) (-4.810,-1.782) (-4.778,-1.715) (-4.743,-1.649) (-4.706,-1.583) (-4.666,-1.520) (-4.624,-1.457) (-4.581,-1.397) (-4.535,-1.337) (-4.487,-1.279) (-4.438,-1.223) (-4.387,-1.168) (-4.334,-1.114) (-4.280,-1.063) (-4.224,-1.012) (-4.168,-0.963) (-4.109,-0.916) (-4.050,-0.870) (-3.990,-0.826) (-3.928,-0.783) (-3.866,-0.741) (-3.802,-0.701) (-3.738,-0.663) (-3.673,-0.625) (-3.651,-0.613)}
\LC{(-2.011,-4.857) (-2.066,-4.908) (-2.122,-4.958) (-2.179,-5.007) (-2.236,-5.055) (-2.294,-5.102) (-2.353,-5.149) (-2.412,-5.195) (-2.472,-5.240) (-2.533,-5.284) (-2.594,-5.327) (-2.656,-5.370) (-2.718,-5.411) (-2.782,-5.452) (-2.845,-5.491) (-2.910,-5.530) (-2.974,-5.568) (-3.040,-5.604) (-3.106,-5.640) (-3.172,-5.674) (-3.240,-5.708) (-3.307,-5.740) (-3.375,-5.772) (-3.444,-5.802) (-3.513,-5.831) (-3.583,-5.859) (-3.653,-5.885) (-3.723,-5.911) (-3.794,-5.935) (-3.866,-5.958) (-3.938,-5.980) (-4.010,-6.000) (-4.082,-6.019) (-4.155,-6.037) (-4.228,-6.053) (-4.302,-6.068) (-4.376,-6.082) (-4.450,-6.094) (-4.524,-6.105) (-4.598,-6.114) (-4.673,-6.122) (-4.748,-6.128) (-4.822,-6.132) (-4.897,-6.135) (-4.972,-6.136) (-5.047,-6.136) (-5.122,-6.134) (-5.197,-6.130) (-5.272,-6.125) (-5.347,-6.118) (-5.421,-6.109) (-5.495,-6.098) (-5.569,-6.086) (-5.643,-6.072) (-5.716,-6.055) (-5.789,-6.037) (-5.861,-6.017) (-5.933,-5.996) (-6.004,-5.972) (-6.075,-5.946) (-6.144,-5.919) (-6.213,-5.889) (-6.281,-5.858) (-6.349,-5.824) (-6.415,-5.789) (-6.480,-5.752) (-6.544,-5.712) (-6.606,-5.671) (-6.668,-5.628) (-6.728,-5.583) (-6.786,-5.536) (-6.843,-5.488) (-6.899,-5.437) (-6.953,-5.385) (-7.005,-5.331) (-7.055,-5.275) (-7.103,-5.218) (-7.149,-5.159) (-7.194,-5.098) (-7.236,-5.036) (-7.276,-4.973) (-7.314,-4.908) (-7.350,-4.842) (-7.383,-4.775) (-7.414,-4.707) (-7.443,-4.637) (-7.469,-4.567) (-7.493,-4.496) (-7.514,-4.424) (-7.533,-4.352) (-7.550,-4.278) (-7.564,-4.205) (-7.576,-4.131) (-7.585,-4.056) (-7.591,-3.982) (-7.596,-3.907) (-7.598,-3.832) (-7.597,-3.757) (-7.594,-3.682) (-7.589,-3.607) (-7.582,-3.532) (-7.572,-3.458) (-7.560,-3.384) (-7.546,-3.310) (-7.530,-3.237) (-7.512,-3.164) (-7.492,-3.092) (-7.470,-3.020) (-7.446,-2.949) (-7.420,-2.879) (-7.392,-2.809) (-7.362,-2.740) (-7.331,-2.672) (-7.298,-2.605) (-7.263,-2.538) (-7.227,-2.473) (-7.189,-2.408) (-7.150,-2.344) (-7.110,-2.281) (-7.068,-2.219) (-7.024,-2.158) (-6.979,-2.098) (-6.933,-2.038) (-6.886,-1.980) (-6.838,-1.923) (-6.788,-1.866) (-6.738,-1.811) (-6.686,-1.757) (-6.633,-1.703) (-6.580,-1.651) (-6.525,-1.600) (-6.469,-1.549) (-6.413,-1.500) (-6.356,-1.451) (-6.297,-1.404) (-6.239,-1.358) (-6.179,-1.312) (-6.119,-1.268) (-6.057,-1.224) (-5.996,-1.182)}
\LC{(-3.675,-6.371) (-3.743,-6.403) (-3.812,-6.433) (-3.881,-6.462) (-3.951,-6.490) (-4.020,-6.518) (-4.091,-6.544) (-4.161,-6.569) (-4.233,-6.593) (-4.304,-6.615) (-4.376,-6.637) (-4.448,-6.657) (-4.520,-6.677) (-4.593,-6.695) (-4.666,-6.711) (-4.740,-6.727) (-4.813,-6.741) (-4.887,-6.754) (-4.961,-6.766) (-5.036,-6.776) (-5.110,-6.785) (-5.185,-6.793) (-5.259,-6.799) (-5.334,-6.804) (-5.409,-6.808) (-5.484,-6.810) (-5.559,-6.810) (-5.634,-6.809) (-5.709,-6.807) (-5.784,-6.803) (-5.859,-6.797) (-5.933,-6.790) (-6.008,-6.781) (-6.082,-6.771) (-6.156,-6.759) (-6.230,-6.745) (-6.303,-6.730) (-6.376,-6.713) (-6.449,-6.694) (-6.521,-6.674) (-6.593,-6.652) (-6.664,-6.628) (-6.735,-6.603) (-6.805,-6.576) (-6.874,-6.547) (-6.942,-6.516) (-7.010,-6.483) (-7.076,-6.449) (-7.142,-6.413) (-7.207,-6.375) (-7.271,-6.336) (-7.333,-6.295) (-7.395,-6.252) (-7.455,-6.207) (-7.514,-6.161) (-7.572,-6.113) (-7.628,-6.063) (-7.683,-6.012) (-7.736,-5.959) (-7.787,-5.904) (-7.837,-5.848) (-7.886,-5.791) (-7.932,-5.732) (-7.977,-5.672) (-8.019,-5.610) (-8.060,-5.547) (-8.099,-5.483) (-8.136,-5.418) (-8.171,-5.351) (-8.204,-5.284) (-8.235,-5.216) (-8.263,-5.146) (-8.289,-5.076) (-8.314,-5.005) (-8.336,-4.933) (-8.356,-4.861) (-8.373,-4.788) (-8.389,-4.715) (-8.402,-4.641) (-8.413,-4.567) (-8.422,-4.492) (-8.428,-4.418) (-8.433,-4.343) (-8.435,-4.268) (-8.435,-4.193) (-8.433,-4.118) (-8.429,-4.043) (-8.423,-3.968) (-8.415,-3.894) (-8.404,-3.819) (-8.392,-3.745) (-8.378,-3.672) (-8.362,-3.598) (-8.344,-3.526) (-8.324,-3.453) (-8.303,-3.381) (-8.279,-3.310) (-8.254,-3.239) (-8.228,-3.169) (-8.199,-3.100) (-8.170,-3.031) (-8.138,-2.963) (-8.105,-2.896) (-8.071,-2.829) (-8.035,-2.763) (-7.998,-2.698) (-7.959,-2.634) (-7.919,-2.570) (-7.878,-2.508) (-7.835,-2.446) (-7.792,-2.385) (-7.747,-2.325) (-7.701,-2.265) (-7.654,-2.207) (-7.606,-2.149) (-7.557,-2.093) (-7.506,-2.037) (-7.455,-1.982)}
\LC{(3.548,6.507) (3.614,6.542) (3.681,6.577) (3.748,6.610) (3.816,6.643) (3.884,6.675) (3.952,6.705) (4.021,6.735) (4.090,6.764) (4.160,6.792) (4.230,6.818) (4.300,6.844) (4.371,6.869) (4.442,6.892) (4.514,6.915) (4.586,6.936) (4.658,6.956) (4.731,6.975) (4.803,6.993) (4.877,7.010) (4.950,7.026) (5.024,7.040) (5.097,7.053) (5.171,7.065) (5.246,7.075) (5.320,7.085) (5.395,7.093) (5.469,7.099) (5.544,7.105) (5.619,7.109) (5.694,7.111) (5.769,7.112) (5.844,7.112) (5.919,7.111) (5.994,7.107) (6.069,7.103) (6.144,7.097) (6.218,7.089) (6.293,7.080) (6.367,7.069) (6.441,7.057) (6.515,7.043) (6.588,7.028) (6.661,7.011) (6.734,6.992) (6.806,6.972) (6.878,6.950) (6.949,6.927) (7.020,6.902) (7.090,6.875) (7.159,6.846) (7.228,6.816) (7.296,6.784) (7.363,6.751) (7.429,6.716) (7.494,6.679) (7.558,6.640) (7.622,6.600) (7.684,6.558) (7.745,6.514) (7.805,6.469) (7.863,6.422) (7.921,6.374) (7.977,6.324) (8.031,6.273) (8.084,6.220) (8.136,6.165) (8.186,6.109) (8.234,6.052) (8.281,5.993) (8.325,5.933) (8.369,5.871) (8.410,5.809) (8.449,5.745) (8.487,5.680) (8.522,5.614) (8.556,5.547) (8.587,5.479) (8.617,5.410) (8.645,5.340) (8.670,5.270) (8.693,5.199) (8.715,5.127) (8.734,5.054) (8.751,4.981) (8.766,4.908) (8.778,4.834) (8.789,4.759) (8.798,4.685) (8.804,4.610) (8.808,4.535) (8.811,4.460) (8.811,4.385) (8.809,4.310) (8.805,4.235) (8.800,4.161) (8.792,4.086) (8.782,4.012) (8.771,3.938) (8.757,3.864) (8.742,3.790) (8.725,3.717) (8.706,3.645) (8.685,3.573) (8.663,3.501) (8.639,3.430) (8.614,3.359) (8.586,3.290) (8.558,3.220) (8.527,3.152) (8.496,3.084) (8.463,3.016) (8.428,2.950) (8.392,2.884) (8.355,2.819) (8.316,2.755) (8.276,2.691) (8.235,2.628) (8.193,2.566) (8.150,2.505) (8.105,2.445) (8.060,2.385) (8.013,2.327) (7.965,2.269) (7.916,2.212) (7.867,2.156) (7.816,2.100) (7.764,2.046)}
\LC{(-1.844,3.549) (-1.909,3.586) (-1.974,3.623) (-2.041,3.657) (-2.108,3.690) (-2.177,3.721) (-2.246,3.750) (-2.315,3.778) (-2.386,3.804) (-2.457,3.827) (-2.529,3.849) (-2.601,3.869) (-2.674,3.887) (-2.747,3.902) (-2.821,3.916) (-2.895,3.927) (-2.970,3.935) (-3.045,3.942) (-3.119,3.946) (-3.194,3.947) (-3.269,3.946) (-3.344,3.942) (-3.419,3.936) (-3.494,3.927) (-3.568,3.915) (-3.641,3.900) (-3.714,3.882) (-3.786,3.862) (-3.857,3.838) (-3.927,3.812) (-3.996,3.782) (-4.064,3.749) (-4.130,3.714) (-4.194,3.675) (-4.256,3.633) (-4.317,3.589) (-4.375,3.541) (-4.430,3.491) (-4.483,3.438) (-4.533,3.382) (-4.580,3.324) (-4.624,3.263) (-4.665,3.200) (-4.702,3.135) (-4.736,3.068) (-4.766,2.999) (-4.793,2.929) (-4.816,2.858) (-4.834,2.785) (-4.850,2.712) (-4.861,2.637) (-4.868,2.563) (-4.872,2.488) (-4.872,2.413) (-4.868,2.338) (-4.861,2.263) (-4.850,2.189) (-4.836,2.116) (-4.819,2.043) (-4.798,1.970) (-4.775,1.899) (-4.748,1.829) (-4.719,1.760) (-4.687,1.692) (-4.652,1.626) (-4.615,1.560) (-4.576,1.497) (-4.534,1.434) (-4.491,1.373) (-4.445,1.314) (-4.398,1.256) (-4.348,1.199) (-4.297,1.144) (-4.280,1.126)}
\LC{(2.501,3.832) (2.574,3.853) (2.646,3.871) (2.720,3.887) (2.793,3.901) (2.867,3.913) (2.942,3.922) (3.016,3.929) (3.091,3.934) (3.166,3.936) (3.241,3.935) (3.316,3.932) (3.391,3.926) (3.465,3.918) (3.540,3.907) (3.613,3.893) (3.686,3.876) (3.759,3.856) (3.830,3.833) (3.900,3.807) (3.970,3.778) (4.037,3.746) (4.104,3.711) (4.168,3.673) (4.231,3.632) (4.292,3.588) (4.350,3.541) (4.406,3.491) (4.460,3.438) (4.511,3.383) (4.558,3.325) (4.603,3.265) (4.644,3.203) (4.682,3.138) (4.717,3.071) (4.747,3.003) (4.775,2.933) (4.798,2.862) (4.817,2.789) (4.833,2.716) (4.845,2.642) (4.853,2.567) (4.857,2.492) (4.858,2.417) (4.855,2.343) (4.848,2.268) (4.838,2.194) (4.824,2.120) (4.807,2.047) (4.787,1.975) (4.764,1.903) (4.738,1.833) (4.709,1.764) (4.677,1.696) (4.643,1.629) (4.606,1.564) (4.580,1.521)}
\LC{(-2.404,-4.600) (-2.469,-4.638) (-2.535,-4.675) (-2.601,-4.710) (-2.668,-4.744) (-2.735,-4.777) (-2.803,-4.809) (-2.872,-4.839) (-2.941,-4.868) (-3.011,-4.895) (-3.081,-4.921) (-3.152,-4.945) (-3.223,-4.968) (-3.295,-4.990) (-3.368,-5.010) (-3.440,-5.028) (-3.513,-5.045) (-3.587,-5.060) (-3.661,-5.073) (-3.735,-5.084) (-3.809,-5.094) (-3.884,-5.102) (-3.959,-5.108) (-4.033,-5.112) (-4.108,-5.115) (-4.183,-5.115) (-4.258,-5.113) (-4.333,-5.110) (-4.408,-5.104) (-4.483,-5.096) (-4.557,-5.086) (-4.631,-5.074) (-4.705,-5.060) (-4.778,-5.043) (-4.850,-5.024) (-4.922,-5.003) (-4.994,-4.980) (-5.064,-4.954) (-5.134,-4.926) (-5.202,-4.896) (-5.270,-4.863) (-5.336,-4.828) (-5.401,-4.791) (-5.465,-4.751) (-5.527,-4.710) (-5.588,-4.665) (-5.646,-4.619) (-5.703,-4.570) (-5.759,-4.519) (-5.812,-4.466) (-5.862,-4.411) (-5.911,-4.354) (-5.957,-4.295) (-6.001,-4.234) (-6.043,-4.172) (-6.082,-4.108) (-6.118,-4.042) (-6.151,-3.975) (-6.182,-3.906) (-6.209,-3.837) (-6.234,-3.766) (-6.256,-3.694) (-6.275,-3.622) (-6.291,-3.548) (-6.304,-3.474) (-6.314,-3.400) (-6.322,-3.325) (-6.326,-3.251) (-6.327,-3.176) (-6.326,-3.101) (-6.322,-3.026) (-6.315,-2.951) (-6.305,-2.877) (-6.293,-2.803) (-6.278,-2.729) (-6.261,-2.656) (-6.241,-2.584) (-6.219,-2.512) (-6.194,-2.441) (-6.167,-2.371) (-6.138,-2.302) (-6.108,-2.234) (-6.075,-2.166) (-6.040,-2.100) (-6.003,-2.035) (-5.964,-1.970) (-5.924,-1.907) (-5.882,-1.845) (-5.838,-1.784) (-5.793,-1.724) (-5.746,-1.666) (-5.698,-1.608) (-5.648,-1.552) (-5.597,-1.497) (-5.545,-1.443) (-5.492,-1.390) (-5.438,-1.339) (-5.382,-1.288) (-5.325,-1.239) (-5.268,-1.191) (-5.209,-1.144) (-5.150,-1.099) (-5.089,-1.054) (-5.028,-1.011) (-4.966,-0.969) (-4.903,-0.928) (-4.840,-0.888) (-4.775,-0.849) (-4.710,-0.811) (-4.667,-0.787)}
\LC{(-3.726,6.525) (-3.794,6.557) (-3.862,6.588) (-3.931,6.618) (-4.000,6.647) (-4.069,6.676) (-4.139,6.703) (-4.210,6.729) (-4.280,6.754) (-4.352,6.778) (-4.423,6.800) (-4.495,6.822) (-4.567,6.842) (-4.639,6.862) (-4.712,6.880) (-4.785,6.897) (-4.858,6.913) (-4.932,6.927) (-5.006,6.941) (-5.080,6.953) (-5.154,6.963) (-5.229,6.973) (-5.303,6.981) (-5.378,6.988) (-5.453,6.993) (-5.527,6.997) (-5.602,7.000) (-5.677,7.001) (-5.752,7.001) (-5.827,6.999) (-5.902,6.996) (-5.977,6.991) (-6.052,6.985) (-6.127,6.978) (-6.201,6.968) (-6.275,6.958) (-6.349,6.945) (-6.423,6.931) (-6.496,6.915) (-6.569,6.898) (-6.642,6.879) (-6.714,6.859) (-6.785,6.837) (-6.857,6.813) (-6.927,6.787) (-6.997,6.760) (-7.066,6.731) (-7.135,6.700) (-7.202,6.668) (-7.269,6.634) (-7.335,6.598) (-7.400,6.560) (-7.464,6.521) (-7.527,6.480) (-7.588,6.438) (-7.649,6.394) (-7.708,6.348) (-7.766,6.300) (-7.823,6.251) (-7.878,6.200) (-7.932,6.148) (-7.984,6.094) (-8.035,6.039) (-8.084,5.982) (-8.131,5.924) (-8.177,5.864) (-8.221,5.804) (-8.263,5.741) (-8.303,5.678) (-8.341,5.614) (-8.377,5.548) (-8.411,5.481) (-8.444,5.413) (-8.474,5.345) (-8.502,5.275) (-8.528,5.205) (-8.552,5.134) (-8.573,5.062) (-8.593,4.990) (-8.610,4.917) (-8.626,4.843) (-8.639,4.769) (-8.650,4.695) (-8.659,4.621) (-8.665,4.546) (-8.670,4.471) (-8.672,4.396) (-8.673,4.321) (-8.671,4.246) (-8.667,4.171) (-8.662,4.097) (-8.654,4.022) (-8.644,3.948) (-8.633,3.873) (-8.619,3.800) (-8.604,3.726) (-8.586,3.653) (-8.567,3.581) (-8.547,3.509) (-8.524,3.437) (-8.500,3.366) (-8.474,3.296) (-8.447,3.226) (-8.418,3.157) (-8.387,3.088) (-8.355,3.020) (-8.322,2.953) (-8.287,2.887) (-8.250,2.821) (-8.213,2.756) (-8.174,2.692) (-8.134,2.629) (-8.092,2.567) (-8.050,2.505) (-8.006,2.444) (-7.961,2.384) (-7.915,2.325) (-7.868,2.266) (-7.819,2.209) (-7.770,2.152) (-7.720,2.097) (-7.669,2.042) (-7.617,1.988)}
\LC{(-1.428,-5.625) (-1.465,-5.691) (-1.502,-5.756) (-1.540,-5.821) (-1.578,-5.885) (-1.617,-5.950) (-1.656,-6.014) (-1.695,-6.077) (-1.735,-6.141) (-1.775,-6.204) (-1.816,-6.267) (-1.857,-6.330) (-1.899,-6.392) (-1.941,-6.454) (-1.984,-6.516) (-2.027,-6.577) (-2.070,-6.639) (-2.114,-6.699) (-2.158,-6.760) (-2.203,-6.820) (-2.248,-6.880) (-2.294,-6.939) (-2.340,-6.999) (-2.387,-7.057) (-2.434,-7.116) (-2.481,-7.174) (-2.529,-7.232) (-2.577,-7.289) (-2.626,-7.346) (-2.675,-7.403) (-2.724,-7.459) (-2.774,-7.515) (-2.825,-7.571) (-2.876,-7.626) (-2.927,-7.681) (-2.978,-7.735) (-3.031,-7.789) (-3.083,-7.842) (-3.136,-7.896) (-3.189,-7.948) (-3.243,-8.001) (-3.297,-8.052) (-3.352,-8.104) (-3.407,-8.155) (-3.462,-8.205) (-3.518,-8.256) (-3.574,-8.305) (-3.631,-8.354) (-3.688,-8.403) (-3.745,-8.451) (-3.803,-8.499) (-3.861,-8.546) (-3.920,-8.593) (-3.979,-8.640) (-4.038,-8.685) (-4.098,-8.731) (-4.158,-8.776) (-4.219,-8.820) (-4.280,-8.864) (-4.341,-8.907) (-4.402,-8.950) (-4.465,-8.992) (-4.527,-9.033) (-4.590,-9.074) (-4.653,-9.115) (-4.716,-9.155) (-4.780,-9.194) (-4.844,-9.233) (-4.909,-9.271) (-4.974,-9.309) (-5.039,-9.346) (-5.104,-9.382) (-5.170,-9.418) (-5.237,-9.453) (-5.303,-9.488) (-5.370,-9.522) (-5.437,-9.555) (-5.505,-9.588) (-5.573,-9.620) (-5.641,-9.651) (-5.709,-9.682) (-5.778,-9.712) (-5.847,-9.741) (-5.916,-9.770) (-5.986,-9.798) (-6.056,-9.825) (-6.126,-9.852) (-6.196,-9.877) (-6.267,-9.902) (-6.338,-9.927) (-6.409,-9.950) (-6.481,-9.973) (-6.552,-9.995) (-6.576,-10.000)}
\LC{(1.055,5.167) (1.085,5.236) (1.116,5.304) (1.147,5.372) (1.178,5.440) (1.210,5.508) (1.242,5.576) (1.275,5.644) (1.308,5.711) (1.341,5.778) (1.375,5.845) (1.409,5.912) (1.444,5.978) (1.479,6.045) (1.515,6.111) (1.550,6.177) (1.587,6.242) (1.623,6.308) (1.661,6.373) (1.698,6.438) (1.736,6.502) (1.774,6.567) (1.813,6.631) (1.852,6.695) (1.892,6.759) (1.932,6.822) (1.972,6.886) (2.013,6.948) (2.054,7.011) (2.096,7.074) (2.138,7.136) (2.180,7.198) (2.223,7.259) (2.266,7.321) (2.309,7.382) (2.353,7.442) (2.398,7.503) (2.443,7.563) (2.488,7.623) (2.533,7.683) (2.579,7.742) (2.625,7.801) (2.672,7.859) (2.719,7.918) (2.767,7.976) (2.815,8.034) (2.863,8.091) (2.912,8.148) (2.961,8.205) (3.010,8.261) (3.060,8.317) (3.110,8.373) (3.161,8.428) (3.212,8.483) (3.263,8.538) (3.315,8.592) (3.367,8.646) (3.420,8.700) (3.473,8.753) (3.526,8.806) (3.579,8.858) (3.633,8.910) (3.688,8.962) (3.743,9.013) (3.798,9.064) (3.853,9.114) (3.909,9.164) (3.965,9.214) (4.022,9.263) (4.079,9.312) (4.136,9.360) (4.194,9.408) (4.252,9.456) (4.310,9.503) (4.369,9.550) (4.428,9.596) (4.487,9.642) (4.547,9.687) (4.607,9.732) (4.667,9.776) (4.728,9.820) (4.789,9.864) (4.851,9.907) (4.913,9.949) (4.975,9.991) (4.996,10.000)}
\LC{(-4.233,7.934) (-4.299,7.971) (-4.364,8.007) (-4.430,8.043) (-4.497,8.078) (-4.563,8.112) (-4.631,8.145) (-4.698,8.178) (-4.766,8.210) (-4.834,8.241) (-4.903,8.271) (-4.972,8.301) (-5.041,8.329) (-5.111,8.357) (-5.181,8.384) (-5.251,8.411) (-5.322,8.436) (-5.393,8.460) (-5.464,8.484) (-5.535,8.507) (-5.607,8.529) (-5.679,8.549) (-5.751,8.569) (-5.824,8.589) (-5.897,8.607) (-5.970,8.624) (-6.043,8.640) (-6.116,8.655) (-6.190,8.670) (-6.264,8.683) (-6.338,8.695) (-6.412,8.706) (-6.486,8.717) (-6.561,8.726) (-6.635,8.734) (-6.710,8.741) (-6.785,8.747) (-6.859,8.751) (-6.934,8.755) (-7.009,8.758) (-7.084,8.759) (-7.159,8.760) (-7.234,8.759) (-7.309,8.757) (-7.384,8.753) (-7.459,8.749) (-7.534,8.743) (-7.608,8.737) (-7.683,8.729) (-7.757,8.719) (-7.832,8.709) (-7.906,8.697) (-7.980,8.684) (-8.053,8.669) (-8.127,8.654) (-8.200,8.637) (-8.272,8.619) (-8.345,8.599) (-8.417,8.578) (-8.488,8.556) (-8.560,8.532) (-8.630,8.507) (-8.701,8.481) (-8.770,8.453) (-8.840,8.424) (-8.908,8.394) (-8.976,8.362) (-9.043,8.329) (-9.110,8.295) (-9.176,8.259) (-9.241,8.222) (-9.305,8.183) (-9.369,8.143) (-9.431,8.102) (-9.493,8.060) (-9.554,8.016) (-9.614,7.970) (-9.673,7.924) (-9.730,7.876) (-9.787,7.827) (-9.843,7.776) (-9.897,7.725) (-9.950,7.672) (-10.000,7.618)}
\LC{(4.719,-7.134) (4.791,-7.155) (4.863,-7.174) (4.936,-7.192) (5.009,-7.209) (5.083,-7.224) (5.156,-7.239) (5.230,-7.252) (5.304,-7.264) (5.378,-7.275) (5.453,-7.285) (5.527,-7.293) (5.602,-7.300) (5.677,-7.306) (5.751,-7.310) (5.826,-7.313) (5.901,-7.315) (5.976,-7.316) (6.051,-7.315) (6.126,-7.312) (6.201,-7.308) (6.276,-7.303) (6.351,-7.296) (6.425,-7.288) (6.500,-7.279) (6.574,-7.267) (6.648,-7.255) (6.721,-7.240) (6.795,-7.225) (6.868,-7.207) (6.940,-7.188) (7.012,-7.168) (7.084,-7.146) (7.155,-7.122) (7.226,-7.097) (7.296,-7.070) (7.365,-7.041) (7.434,-7.011) (7.502,-6.979) (7.569,-6.946) (7.635,-6.911) (7.701,-6.874) (7.765,-6.836) (7.828,-6.796) (7.891,-6.754) (7.952,-6.711) (8.012,-6.666) (8.071,-6.620) (8.129,-6.572) (8.185,-6.523) (8.240,-6.472) (8.294,-6.419) (8.346,-6.365) (8.397,-6.310) (8.446,-6.253) (8.493,-6.195) (8.539,-6.136) (8.583,-6.075) (8.626,-6.013) (8.666,-5.950) (8.705,-5.886) (8.742,-5.821) (8.777,-5.754) (8.810,-5.687) (8.841,-5.619) (8.870,-5.550) (8.897,-5.480) (8.922,-5.409) (8.945,-5.338) (8.966,-5.266) (8.985,-5.193) (9.002,-5.120) (9.017,-5.047) (9.030,-4.973) (9.040,-4.898) (9.049,-4.824) (9.056,-4.749) (9.060,-4.674) (9.063,-4.599) (9.064,-4.524) (9.062,-4.449) (9.059,-4.374) (9.053,-4.300) (9.046,-4.225) (9.037,-4.151) (9.026,-4.076) (9.013,-4.002) (8.999,-3.929) (8.983,-3.856) (8.964,-3.783) (8.945,-3.711) (8.923,-3.639) (8.900,-3.567) (8.876,-3.497) (8.849,-3.426) (8.821,-3.357) (8.792,-3.288) (8.762,-3.219) (8.729,-3.151) (8.696,-3.084) (8.661,-3.018) (8.625,-2.952) (8.587,-2.887) (8.548,-2.823) (8.508,-2.760)}
\LC{(-1.000,2.375) (-1.058,2.423) (-1.117,2.470) (-1.177,2.515) (-1.238,2.558) (-1.301,2.599) (-1.365,2.638) (-1.430,2.675) (-1.497,2.710) (-1.564,2.742) (-1.633,2.773) (-1.702,2.801) (-1.773,2.826) (-1.844,2.849) (-1.917,2.869) (-1.990,2.886) (-2.063,2.900) (-2.137,2.912) (-2.212,2.920) (-2.287,2.925) (-2.362,2.926) (-2.437,2.924) (-2.512,2.919) (-2.586,2.909) (-2.660,2.896) (-2.733,2.880) (-2.805,2.859) (-2.876,2.834) (-2.945,2.805) (-3.012,2.772) (-3.077,2.735) (-3.140,2.694) (-3.200,2.649) (-3.257,2.600) (-3.310,2.547) (-3.360,2.491) (-3.405,2.431) (-3.446,2.368) (-3.482,2.303) (-3.514,2.235) (-3.540,2.164) (-3.562,2.093) (-3.578,2.019) (-3.589,1.945) (-3.595,1.870) (-3.596,1.795) (-3.592,1.721) (-3.583,1.646) (-3.570,1.572) (-3.552,1.500) (-3.530,1.428) (-3.504,1.357) (-3.474,1.289) (-3.441,1.222) (-3.404,1.156) (-3.364,1.093) (-3.321,1.031) (-3.276,0.971) (-3.228,0.914) (-3.177,0.858) (-3.125,0.805) (-3.070,0.754) (-3.014,0.704) (-2.955,0.657) (-2.895,0.612) (-2.834,0.569) (-2.771,0.528) (-2.707,0.489) (-2.642,0.452) (-2.576,0.417) (-2.553,0.405)}
\LC{(-10.000,-0.013) (-9.942,-0.013) (-9.867,-0.013) (-9.792,-0.012) (-9.717,-0.012) (-9.642,-0.012) (-9.567,-0.012) (-9.492,-0.011) (-9.417,-0.011) (-9.342,-0.011) (-9.267,-0.010) (-9.192,-0.010) (-9.117,-0.010) (-9.042,-0.010) (-8.967,-0.010) (-8.892,-0.009) (-8.817,-0.009) (-8.742,-0.009) (-8.667,-0.009) (-8.592,-0.008) (-8.517,-0.008) (-8.442,-0.008) (-8.367,-0.008) (-8.292,-0.008) (-8.217,-0.007) (-8.142,-0.007) (-8.067,-0.007) (-7.992,-0.007) (-7.917,-0.007) (-7.842,-0.006) (-7.767,-0.006) (-7.692,-0.006) (-7.617,-0.006) (-7.542,-0.006) (-7.467,-0.006) (-7.392,-0.005) (-7.317,-0.005) (-7.242,-0.005) (-7.167,-0.005) (-7.117,-0.005)}
\LC{(1.553,-6.814) (1.586,-6.882) (1.620,-6.949) (1.654,-7.016) (1.688,-7.082) (1.722,-7.149) (1.757,-7.215) (1.793,-7.282) (1.828,-7.348) (1.864,-7.413) (1.900,-7.479) (1.937,-7.544) (1.974,-7.610) (2.011,-7.675) (2.049,-7.740) (2.087,-7.804) (2.125,-7.869) (2.164,-7.933) (2.203,-7.997) (2.242,-8.061) (2.282,-8.125) (2.322,-8.188) (2.362,-8.251) (2.403,-8.314) (2.444,-8.377) (2.485,-8.440) (2.527,-8.502) (2.569,-8.564) (2.611,-8.626) (2.654,-8.688) (2.697,-8.749) (2.740,-8.811) (2.784,-8.871) (2.828,-8.932) (2.872,-8.993) (2.917,-9.053) (2.962,-9.113) (3.007,-9.173) (3.053,-9.232) (3.099,-9.291) (3.145,-9.350) (3.192,-9.409) (3.239,-9.468) (3.286,-9.526) (3.334,-9.584) (3.381,-9.642) (3.430,-9.699) (3.478,-9.756) (3.527,-9.813) (3.576,-9.869) (3.626,-9.926) (3.676,-9.982) (3.693,-10.000)}
\LC{(0.766,5.622) (0.786,5.694) (0.807,5.766) (0.828,5.838) (0.849,5.910) (0.871,5.982) (0.893,6.054) (0.915,6.126) (0.937,6.197) (0.959,6.269) (0.982,6.340) (1.005,6.411) (1.029,6.483) (1.052,6.554) (1.076,6.625) (1.100,6.696) (1.125,6.767) (1.149,6.838) (1.174,6.909) (1.199,6.979) (1.225,7.050) (1.250,7.120) (1.276,7.191) (1.303,7.261) (1.329,7.331) (1.356,7.401) (1.383,7.471) (1.410,7.541) (1.438,7.611) (1.466,7.680) (1.494,7.750) (1.522,7.819) (1.551,7.889) (1.579,7.958) (1.609,8.027) (1.638,8.096) (1.668,8.165) (1.698,8.234) (1.728,8.302) (1.758,8.371) (1.789,8.439) (1.820,8.508) (1.851,8.576) (1.883,8.644) (1.914,8.712) (1.946,8.780) (1.979,8.847) (2.011,8.915) (2.044,8.982) (2.077,9.050) (2.110,9.117) (2.144,9.184) (2.178,9.251) (2.212,9.318) (2.246,9.384) (2.281,9.451) (2.315,9.517) (2.351,9.584) (2.386,9.650) (2.422,9.716) (2.457,9.782) (2.494,9.847) (2.530,9.913) (2.567,9.978) (2.579,10.000)}
\LC{(3.636,-7.607) (3.697,-7.650) (3.758,-7.694) (3.820,-7.736) (3.882,-7.778) (3.945,-7.819) (4.008,-7.860) (4.071,-7.900) (4.135,-7.939) (4.199,-7.978) (4.264,-8.015) (4.329,-8.053) (4.395,-8.089) (4.461,-8.125) (4.527,-8.160) (4.594,-8.195) (4.661,-8.228) (4.728,-8.261) (4.796,-8.293) (4.864,-8.325) (4.932,-8.355) (5.001,-8.385) (5.070,-8.414) (5.140,-8.442) (5.210,-8.470) (5.280,-8.496) (5.350,-8.522) (5.421,-8.547) (5.492,-8.571) (5.563,-8.594) (5.635,-8.616) (5.707,-8.638) (5.779,-8.658) (5.851,-8.678) (5.924,-8.696) (5.997,-8.714) (6.070,-8.731) (6.143,-8.747) (6.217,-8.762) (6.291,-8.776) (6.364,-8.788) (6.439,-8.800) (6.513,-8.811) (6.587,-8.821) (6.662,-8.830) (6.736,-8.837) (6.811,-8.844) (6.886,-8.850) (6.961,-8.854) (7.035,-8.858) (7.110,-8.860) (7.185,-8.861) (7.260,-8.861) (7.335,-8.860) (7.410,-8.858) (7.485,-8.854) (7.560,-8.850) (7.635,-8.844) (7.710,-8.837) (7.784,-8.828) (7.858,-8.819) (7.933,-8.808) (8.007,-8.796) (8.081,-8.783) (8.154,-8.769) (8.227,-8.753) (8.301,-8.736) (8.373,-8.717) (8.446,-8.698) (8.518,-8.677) (8.589,-8.654) (8.660,-8.631) (8.731,-8.606) (8.801,-8.580) (8.871,-8.552) (8.940,-8.523) (9.009,-8.493) (9.077,-8.461) (9.144,-8.428) (9.211,-8.394) (9.277,-8.358) (9.342,-8.321) (9.406,-8.283) (9.470,-8.243) (9.533,-8.202) (9.595,-8.159) (9.656,-8.116) (9.716,-8.071) (9.775,-8.024) (9.833,-7.977) (9.889,-7.928) (9.945,-7.878) (10.000,-7.826) (10.000,-7.809)}
\LC{(1.310,-2.704) (1.373,-2.746) (1.436,-2.785) (1.501,-2.822) (1.567,-2.858) (1.635,-2.891) (1.703,-2.922) (1.772,-2.951) (1.842,-2.978) (1.913,-3.002) (1.985,-3.024) (2.058,-3.043) (2.131,-3.059) (2.204,-3.073) (2.279,-3.083) (2.353,-3.091) (2.428,-3.096) (2.503,-3.097) (2.578,-3.095) (2.653,-3.090) (2.727,-3.081) (2.801,-3.069) (2.875,-3.053) (2.947,-3.033) (3.018,-3.010) (3.088,-2.983) (3.156,-2.952) (3.223,-2.917) (3.287,-2.878) (3.349,-2.836) (3.408,-2.789) (3.464,-2.739) (3.516,-2.686) (3.565,-2.629) (3.610,-2.569) (3.651,-2.506) (3.687,-2.441) (3.719,-2.373) (3.747,-2.303) (3.769,-2.231) (3.787,-2.158) (3.799,-2.085) (3.807,-2.010) (3.810,-1.935) (3.808,-1.860) (3.802,-1.785) (3.791,-1.711) (3.776,-1.638) (3.757,-1.565) (3.734,-1.494) (3.707,-1.424) (3.677,-1.355) (3.643,-1.288) (3.607,-1.223) (3.567,-1.159) (3.524,-1.098) (3.479,-1.038) (3.431,-0.980) (3.382,-0.924) (3.330,-0.870) (3.294,-0.835)}
\LC{(0.737,-5.757) (0.757,-5.830) (0.776,-5.902) (0.796,-5.975) (0.816,-6.047) (0.836,-6.119) (0.856,-6.191) (0.877,-6.263) (0.898,-6.335) (0.919,-6.407) (0.941,-6.479) (0.962,-6.551) (0.984,-6.623) (1.006,-6.694) (1.029,-6.766) (1.051,-6.838) (1.074,-6.909) (1.097,-6.980) (1.121,-7.052) (1.144,-7.123) (1.168,-7.194) (1.192,-7.265) (1.217,-7.336) (1.241,-7.407) (1.266,-7.478) (1.291,-7.548) (1.316,-7.619) (1.342,-7.689) (1.368,-7.760) (1.394,-7.830) (1.420,-7.900) (1.447,-7.970) (1.473,-8.041) (1.500,-8.111) (1.528,-8.180) (1.555,-8.250) (1.583,-8.320) (1.611,-8.389) (1.639,-8.459) (1.667,-8.528) (1.696,-8.598) (1.725,-8.667) (1.754,-8.736) (1.784,-8.805) (1.813,-8.874) (1.843,-8.943) (1.873,-9.011) (1.904,-9.080) (1.934,-9.148) (1.965,-9.217) (1.996,-9.285) (2.028,-9.353) (2.059,-9.421) (2.091,-9.489) (2.123,-9.557) (2.155,-9.624) (2.188,-9.692) (2.221,-9.759) (2.254,-9.827) (2.287,-9.894) (2.321,-9.961) (2.343,-10.000)}
\LC{(0.446,6.951) (0.455,7.025) (0.465,7.100) (0.475,7.174) (0.485,7.248) (0.495,7.323) (0.505,7.397) (0.515,7.471) (0.526,7.546) (0.536,7.620) (0.547,7.694) (0.558,7.768) (0.569,7.842) (0.579,7.917) (0.590,7.991) (0.602,8.065) (0.613,8.139) (0.624,8.213) (0.636,8.287) (0.647,8.362) (0.659,8.436) (0.670,8.510) (0.682,8.584) (0.694,8.658) (0.706,8.732) (0.719,8.806) (0.731,8.880) (0.743,8.954) (0.756,9.028) (0.768,9.102) (0.781,9.176) (0.794,9.250) (0.807,9.323) (0.820,9.397) (0.833,9.471) (0.846,9.545) (0.859,9.619) (0.872,9.693) (0.886,9.766) (0.900,9.840) (0.913,9.914) (0.927,9.988) (0.932,10.000)}
\LC{(-0.458,6.326) (-0.469,6.400) (-0.480,6.475) (-0.491,6.549) (-0.503,6.623) (-0.514,6.697) (-0.526,6.771) (-0.537,6.845) (-0.549,6.919) (-0.561,6.993) (-0.573,7.067) (-0.586,7.141) (-0.598,7.215) (-0.610,7.289) (-0.623,7.363) (-0.636,7.437) (-0.649,7.511) (-0.662,7.585) (-0.675,7.659) (-0.688,7.733) (-0.701,7.806) (-0.715,7.880) (-0.729,7.954) (-0.742,8.028) (-0.756,8.101) (-0.770,8.175) (-0.784,8.249) (-0.799,8.322) (-0.813,8.396) (-0.828,8.469) (-0.842,8.543) (-0.857,8.617) (-0.872,8.690) (-0.887,8.764) (-0.902,8.837) (-0.917,8.910) (-0.933,8.984) (-0.948,9.057) (-0.964,9.130) (-0.980,9.204) (-0.996,9.277) (-1.012,9.350) (-1.028,9.424) (-1.044,9.497) (-1.061,9.570) (-1.077,9.643) (-1.094,9.716) (-1.111,9.789) (-1.128,9.862) (-1.145,9.935) (-1.162,10.000)}
\LC{(-5.950,-9.296) (-6.021,-9.319) (-6.093,-9.342) (-6.165,-9.364) (-6.237,-9.384) (-6.309,-9.404) (-6.381,-9.424) (-6.454,-9.442) (-6.527,-9.459) (-6.600,-9.476) (-6.674,-9.492) (-6.747,-9.507) (-6.821,-9.520) (-6.895,-9.533) (-6.969,-9.545) (-7.043,-9.557) (-7.117,-9.567) (-7.192,-9.576) (-7.266,-9.584) (-7.341,-9.592) (-7.415,-9.598) (-7.490,-9.603) (-7.565,-9.608) (-7.640,-9.611) (-7.715,-9.613) (-7.790,-9.614) (-7.865,-9.615) (-7.940,-9.614) (-8.015,-9.612) (-8.090,-9.609) (-8.165,-9.605) (-8.240,-9.599) (-8.314,-9.593) (-8.389,-9.586) (-8.464,-9.577) (-8.538,-9.567) (-8.612,-9.557) (-8.686,-9.545) (-8.760,-9.532) (-8.834,-9.517) (-8.907,-9.502) (-8.980,-9.485) (-9.053,-9.467) (-9.126,-9.448) (-9.198,-9.428) (-9.270,-9.407) (-9.341,-9.384) (-9.412,-9.360) (-9.483,-9.335) (-9.553,-9.308) (-9.623,-9.281) (-9.692,-9.252) (-9.761,-9.222) (-9.829,-9.191) (-9.896,-9.158) (-9.963,-9.124) (-10.000,-9.101)}
\LC{(0.638,-1.711) (0.692,-1.763) (0.748,-1.813) (0.806,-1.861) (0.865,-1.907) (0.927,-1.950) (0.990,-1.991) (1.054,-2.029) (1.120,-2.064) (1.188,-2.097) (1.257,-2.126) (1.327,-2.152) (1.399,-2.175) (1.471,-2.194) (1.545,-2.209) (1.619,-2.221) (1.693,-2.228) (1.768,-2.231) (1.843,-2.229) (1.918,-2.223) (1.992,-2.212) (2.065,-2.196) (2.137,-2.174) (2.207,-2.148) (2.275,-2.116) (2.340,-2.078) (2.402,-2.035) (2.459,-1.987) (2.512,-1.935) (2.560,-1.877) (2.603,-1.815) (2.639,-1.750) (2.670,-1.681) (2.693,-1.610) (2.710,-1.537) (2.720,-1.463) (2.723,-1.388) (2.720,-1.313) (2.710,-1.238) (2.694,-1.165) (2.673,-1.093) (2.647,-1.023) (2.615,-0.955) (2.579,-0.889) (2.539,-0.826) (2.495,-0.765) (2.448,-0.707) (2.398,-0.651) (2.344,-0.598) (2.326,-0.581)}
\LC{(-0.370,-1.511) (-0.409,-1.575) (-0.449,-1.638) (-0.492,-1.700) (-0.536,-1.760) (-0.582,-1.820) (-0.629,-1.878) (-0.678,-1.934) (-0.729,-1.990) (-0.782,-2.043) (-0.836,-2.095) (-0.891,-2.146) (-0.948,-2.194) (-1.007,-2.241) (-1.067,-2.285) (-1.129,-2.328) (-1.192,-2.369) (-1.256,-2.407) (-1.322,-2.443) (-1.389,-2.477) (-1.458,-2.508) (-1.527,-2.536) (-1.597,-2.562) (-1.669,-2.585) (-1.741,-2.604) (-1.814,-2.621) (-1.888,-2.635) (-1.962,-2.645) (-2.037,-2.651) (-2.112,-2.654) (-2.187,-2.653) (-2.262,-2.648) (-2.336,-2.640) (-2.410,-2.626) (-2.483,-2.609) (-2.555,-2.588) (-2.625,-2.561) (-2.694,-2.531) (-2.760,-2.496) (-2.824,-2.456) (-2.884,-2.412) (-2.942,-2.364) (-2.995,-2.311) (-3.044,-2.255) (-3.089,-2.195) (-3.129,-2.131) (-3.164,-2.065) (-3.193,-1.996) (-3.217,-1.924) (-3.235,-1.852) (-3.247,-1.778) (-3.254,-1.703) (-3.255,-1.628) (-3.251,-1.553) (-3.241,-1.479) (-3.227,-1.405) (-3.207,-1.333) (-3.183,-1.262) (-3.155,-1.192) (-3.123,-1.125) (-3.086,-1.059) (-3.047,-0.995) (-3.004,-0.934) (-2.958,-0.874) (-2.910,-0.817) (-2.859,-0.762) (-2.805,-0.710) (-2.749,-0.660) (-2.692,-0.612) (-2.632,-0.566) (-2.571,-0.522) (-2.550,-0.508)}
\LC{(2.984,4.989) (3.053,5.017) (3.123,5.044) (3.194,5.070) (3.265,5.094) (3.336,5.116) (3.408,5.137) (3.481,5.157) (3.554,5.175) (3.627,5.191) (3.700,5.206) (3.774,5.219) (3.848,5.230) (3.923,5.240) (3.997,5.248) (4.072,5.254) (4.147,5.258) (4.222,5.261) (4.297,5.261) (4.372,5.260) (4.447,5.256) (4.521,5.251) (4.596,5.243) (4.670,5.234) (4.745,5.222) (4.818,5.208) (4.892,5.193) (4.964,5.175) (5.037,5.154) (5.108,5.132) (5.179,5.107) (5.249,5.080) (5.318,5.051) (5.386,5.019) (5.453,4.986) (5.519,4.950) (5.583,4.911) (5.646,4.871) (5.708,4.828) (5.768,4.783) (5.826,4.736) (5.883,4.686) (5.937,4.635) (5.990,4.582) (6.040,4.526) (6.089,4.469) (6.135,4.409) (6.178,4.348) (6.219,4.286) (6.258,4.221) (6.294,4.156) (6.327,4.089) (6.358,4.020) (6.386,3.950) (6.411,3.880) (6.433,3.808) (6.453,3.736) (6.469,3.663) (6.483,3.589) (6.493,3.515) (6.501,3.440) (6.506,3.365) (6.509,3.290) (6.508,3.215) (6.505,3.140) (6.499,3.066) (6.491,2.991) (6.479,2.917) (6.466,2.843) (6.450,2.770) (6.431,2.697) (6.410,2.625) (6.387,2.554) (6.362,2.483) (6.334,2.413) (6.305,2.345) (6.273,2.277) (6.239,2.210) (6.204,2.143) (6.167,2.078) (6.128,2.014) (6.087,1.951) (6.045,1.889) (6.001,1.829) (5.955,1.769)}
\LC{(10.000,1.462) (9.942,1.430) (9.874,1.399) (9.805,1.368) (9.737,1.338) (9.668,1.308) (9.599,1.279) (9.530,1.249) (9.461,1.221) (9.391,1.193) (9.321,1.165) (9.252,1.138) (9.182,1.111) (9.111,1.084) (9.041,1.058) (8.971,1.033) (8.900,1.007) (8.829,0.982) (8.758,0.958) (8.687,0.934) (8.616,0.910) (8.545,0.887) (8.473,0.864) (8.402,0.842) (8.330,0.820) (8.258,0.798) (8.186,0.777) (8.114,0.756) (8.042,0.736) (7.970,0.716) (7.897,0.696) (7.825,0.676) (7.752,0.657) (7.680,0.639) (7.607,0.620) (7.534,0.602) (7.461,0.585) (7.388,0.568) (7.315,0.551) (7.242,0.534) (7.169,0.518) (7.096,0.502) (7.022,0.486) (6.949,0.471) (6.875,0.456) (6.826,0.446)}
\LC{(3.477,-5.679) (3.547,-5.706) (3.618,-5.731) (3.689,-5.755) (3.760,-5.778) (3.832,-5.799) (3.904,-5.819) (3.977,-5.838) (4.050,-5.855) (4.123,-5.871) (4.197,-5.885) (4.271,-5.898) (4.345,-5.910) (4.419,-5.920) (4.494,-5.928) (4.569,-5.935) (4.643,-5.940) (4.718,-5.944) (4.793,-5.945) (4.868,-5.946) (4.943,-5.944) (5.018,-5.941) (5.093,-5.936) (5.168,-5.929) (5.242,-5.920) (5.316,-5.910) (5.390,-5.898) (5.464,-5.883) (5.537,-5.867) (5.610,-5.849) (5.682,-5.829) (5.754,-5.807) (5.825,-5.783) (5.896,-5.757) (5.965,-5.730) (6.034,-5.700) (6.102,-5.668) (6.169,-5.634) (6.234,-5.598) (6.299,-5.560) (6.363,-5.520) (6.425,-5.478) (6.486,-5.434) (6.545,-5.388) (6.603,-5.340) (6.659,-5.291) (6.714,-5.239) (6.766,-5.186) (6.817,-5.131) (6.866,-5.074) (6.913,-5.016) (6.958,-4.956) (7.001,-4.894) (7.041,-4.831) (7.080,-4.766) (7.116,-4.701) (7.149,-4.634) (7.181,-4.565) (7.209,-4.496) (7.236,-4.426) (7.260,-4.355) (7.281,-4.283) (7.300,-4.210) (7.316,-4.137) (7.330,-4.063) (7.341,-3.989) (7.350,-3.915) (7.356,-3.840) (7.360,-3.765) (7.361,-3.690) (7.359,-3.615) (7.356,-3.540) (7.350,-3.465) (7.341,-3.391) (7.331,-3.317) (7.318,-3.243) (7.303,-3.169) (7.286,-3.096) (7.266,-3.024) (7.245,-2.952) (7.221,-2.881) (7.196,-2.810) (7.169,-2.740) (7.140,-2.671) (7.109,-2.603) (7.076,-2.535) (7.042,-2.469) (7.006,-2.403) (6.968,-2.338) (6.929,-2.274) (6.888,-2.211) (6.846,-2.149) (6.803,-2.088) (6.773,-2.048)}
\LC{(-3.854,5.607) (-3.927,5.624) (-4.000,5.639) (-4.074,5.652) (-4.148,5.664) (-4.223,5.674) (-4.297,5.683) (-4.372,5.690) (-4.447,5.695) (-4.521,5.699) (-4.596,5.701) (-4.671,5.701) (-4.746,5.699) (-4.821,5.696) (-4.896,5.690) (-4.971,5.683) (-5.045,5.674) (-5.119,5.663) (-5.193,5.650) (-5.267,5.635) (-5.340,5.618) (-5.412,5.599) (-5.484,5.578) (-5.556,5.555) (-5.626,5.530) (-5.696,5.503) (-5.765,5.473) (-5.833,5.442) (-5.901,5.408) (-5.967,5.373) (-6.031,5.335) (-6.095,5.295) (-6.157,5.253) (-6.218,5.210) (-6.277,5.164) (-6.335,5.116) (-6.391,5.066) (-6.445,5.014) (-6.498,4.960) (-6.548,4.905) (-6.596,4.847) (-6.643,4.788) (-6.687,4.728) (-6.729,4.666) (-6.768,4.602) (-6.805,4.537) (-6.840,4.470) (-6.872,4.402) (-6.902,4.334) (-6.929,4.264) (-6.954,4.193) (-6.975,4.121) (-6.995,4.049) (-7.011,3.976) (-7.025,3.902) (-7.037,3.828) (-7.046,3.753) (-7.052,3.678) (-7.055,3.604) (-7.056,3.529) (-7.054,3.454) (-7.050,3.379) (-7.044,3.304) (-7.035,3.230) (-7.023,3.155) (-7.010,3.082) (-6.994,3.008) (-6.975,2.936) (-6.955,2.864) (-6.932,2.792) (-6.908,2.721) (-6.881,2.651) (-6.853,2.582) (-6.822,2.513) (-6.790,2.445) (-6.756,2.379) (-6.720,2.313) (-6.708,2.291)}
\LC{(-2.223,7.316) (-2.266,7.377) (-2.309,7.438) (-2.353,7.499) (-2.397,7.560) (-2.441,7.620) (-2.486,7.681) (-2.531,7.740) (-2.577,7.800) (-2.623,7.859) (-2.670,7.918) (-2.716,7.977) (-2.764,8.035) (-2.811,8.093) (-2.859,8.151) (-2.907,8.208) (-2.956,8.265) (-3.005,8.322) (-3.055,8.378) (-3.104,8.434) (-3.155,8.490) (-3.205,8.545) (-3.256,8.600) (-3.308,8.655) (-3.359,8.709) (-3.412,8.763) (-3.464,8.817) (-3.517,8.870) (-3.570,8.923) (-3.624,8.975) (-3.678,9.027) (-3.732,9.079) (-3.787,9.130) (-3.842,9.181) (-3.897,9.232) (-3.953,9.282) (-4.009,9.332) (-4.066,9.381) (-4.122,9.430) (-4.180,9.478) (-4.237,9.526) (-4.295,9.574) (-4.353,9.621) (-4.412,9.668) (-4.471,9.714) (-4.530,9.760) (-4.590,9.806) (-4.650,9.851) (-4.710,9.895) (-4.771,9.939) (-4.832,9.983) (-4.873,10.000)}
\LC{(3.192,7.782) (3.247,7.834) (3.302,7.885) (3.358,7.935) (3.414,7.985) (3.470,8.034) (3.527,8.083) (3.584,8.132) (3.642,8.180) (3.700,8.228) (3.758,8.275) (3.817,8.321) (3.876,8.367) (3.936,8.413) (3.996,8.458) (4.056,8.502) (4.117,8.546) (4.178,8.589) (4.240,8.632) (4.302,8.674) (4.364,8.716) (4.427,8.757) (4.490,8.798) (4.553,8.838) (4.617,8.877) (4.681,8.916) (4.746,8.954) (4.811,8.991) (4.876,9.028) (4.942,9.065) (5.008,9.100) (5.074,9.135) (5.141,9.170) (5.208,9.203) (5.275,9.236) (5.343,9.269) (5.411,9.300) (5.479,9.331) (5.548,9.362) (5.617,9.391) (5.686,9.420) (5.755,9.448) (5.825,9.476) (5.895,9.502) (5.966,9.528) (6.036,9.554) (6.107,9.578) (6.178,9.602) (6.250,9.625) (6.321,9.647) (6.393,9.668) (6.465,9.688) (6.538,9.708) (6.610,9.727) (6.683,9.745) (6.756,9.762) (6.829,9.779) (6.903,9.794) (6.976,9.809) (7.050,9.822) (7.124,9.835) (7.198,9.847) (7.272,9.858) (7.347,9.868) (7.421,9.877) (7.496,9.886) (7.570,9.893) (7.645,9.899) (7.720,9.905) (7.795,9.909) (7.870,9.913) (7.944,9.915) (8.019,9.916) (8.094,9.917) (8.169,9.916) (8.244,9.915) (8.319,9.912) (8.394,9.908) (8.469,9.903) (8.544,9.898) (8.619,9.891) (8.693,9.883) (8.768,9.873) (8.842,9.863) (8.916,9.852) (8.990,9.840) (9.064,9.826) (9.137,9.811) (9.211,9.795) (9.284,9.778) (9.356,9.760) (9.429,9.741) (9.501,9.720) (9.573,9.699) (9.644,9.676) (9.715,9.652) (9.786,9.627) (9.856,9.600) (9.926,9.573) (9.995,9.544) (10.000,9.534)}
\LC{(2.312,-3.626) (2.383,-3.648) (2.456,-3.668) (2.528,-3.686) (2.602,-3.702) (2.676,-3.715) (2.750,-3.726) (2.824,-3.735) (2.899,-3.741) (2.974,-3.744) (3.049,-3.745) (3.124,-3.743) (3.199,-3.739) (3.273,-3.731) (3.348,-3.721) (3.421,-3.707) (3.495,-3.691) (3.567,-3.671) (3.638,-3.648) (3.709,-3.623) (3.778,-3.594) (3.846,-3.561) (3.912,-3.526) (3.976,-3.488) (4.039,-3.446) (4.099,-3.401) (4.156,-3.353) (4.211,-3.302) (4.264,-3.248) (4.313,-3.192) (4.359,-3.133) (4.402,-3.071) (4.441,-3.007) (4.477,-2.941) (4.508,-2.873) (4.536,-2.804) (4.560,-2.733) (4.580,-2.660) (4.596,-2.587) (4.608,-2.513) (4.616,-2.438) (4.620,-2.364) (4.620,-2.289) (4.616,-2.214) (4.608,-2.139) (4.597,-2.065) (4.582,-1.991) (4.564,-1.919) (4.543,-1.847) (4.518,-1.776) (4.490,-1.706) (4.460,-1.638) (4.426,-1.571) (4.390,-1.505) (4.352,-1.441) (4.325,-1.398)}
\LC{(-2.176,-3.270) (-2.249,-3.288) (-2.322,-3.303) (-2.396,-3.316) (-2.470,-3.327) (-2.545,-3.334) (-2.620,-3.339) (-2.695,-3.341) (-2.770,-3.339) (-2.845,-3.335) (-2.919,-3.327) (-2.993,-3.316) (-3.067,-3.302) (-3.140,-3.284) (-3.212,-3.263) (-3.283,-3.239) (-3.352,-3.210) (-3.420,-3.179) (-3.486,-3.143) (-3.550,-3.104) (-3.612,-3.062) (-3.672,-3.016) (-3.728,-2.967) (-3.781,-2.914) (-3.832,-2.858) (-3.878,-2.800) (-3.921,-2.738) (-3.960,-2.674) (-3.995,-2.607) (-4.025,-2.539) (-4.051,-2.469) (-4.073,-2.397) (-4.090,-2.324) (-4.103,-2.250) (-4.111,-2.175) (-4.115,-2.101) (-4.114,-2.026) (-4.109,-1.951) (-4.100,-1.876) (-4.086,-1.802) (-4.069,-1.729) (-4.049,-1.657) (-4.024,-1.586) (-3.996,-1.517) (-3.965,-1.449) (-3.931,-1.382) (-3.907,-1.338)}
\LC{(-0.672,1.555) (-0.733,1.599) (-0.795,1.641) (-0.860,1.679) (-0.926,1.713) (-0.995,1.745) (-1.064,1.772) (-1.136,1.795) (-1.208,1.814) (-1.282,1.828) (-1.356,1.838) (-1.431,1.842) (-1.506,1.841) (-1.581,1.834) (-1.655,1.821) (-1.727,1.802) (-1.798,1.777) (-1.866,1.745) (-1.930,1.707) (-1.990,1.662) (-2.045,1.611) (-2.094,1.554) (-2.136,1.492) (-2.171,1.426) (-2.198,1.356) (-2.217,1.284) (-2.228,1.209) (-2.231,1.134) (-2.225,1.060) (-2.212,0.986) (-2.192,0.914) (-2.166,0.843) (-2.134,0.776) (-2.096,0.711) (-2.054,0.649) (-2.007,0.590) (-1.957,0.535) (-1.903,0.483) (-1.846,0.434)}
\LC{(-10.000,-2.494) (-9.955,-2.443) (-9.900,-2.392) (-9.844,-2.342) (-9.787,-2.293) (-9.730,-2.244) (-9.673,-2.196) (-9.615,-2.149) (-9.556,-2.102) (-9.497,-2.056) (-9.437,-2.010) (-9.377,-1.965) (-9.317,-1.921) (-9.256,-1.877) (-9.194,-1.835) (-9.132,-1.792) (-9.070,-1.751) (-9.007,-1.709) (-8.944,-1.669) (-8.880,-1.629) (-8.816,-1.590) (-8.752,-1.551) (-8.687,-1.513) (-8.622,-1.476) (-8.557,-1.439) (-8.491,-1.403) (-8.425,-1.368) (-8.359,-1.333) (-8.292,-1.298) (-8.225,-1.264) (-8.158,-1.231) (-8.091,-1.199) (-8.023,-1.166) (-7.955,-1.135) (-7.886,-1.104) (-7.864,-1.094)}

    % Cuadrupolo físico (lineal) - puntas de flecha
% Generado por scripts/generar_datos.py (no editar a mano).
\FL{1.818}{1.414}{-48.1}
\FL{-1.091}{0.869}{-128.2}
\FL{-1.091}{-0.869}{128.2}
\FL{1.393}{-1.177}{41.3}
\FL{0.297}{-2.727}{-77.2}
\FL{1.189}{-5.269}{-64.1}
\FL{2.668}{-7.550}{-49.8}
\FL{4.671}{-9.347}{-33.6}
\FL{-0.703}{-3.491}{-113.4}
\FL{-2.162}{-5.748}{-132.9}
\FL{-4.337}{-7.315}{-156.6}
\FL{-6.974}{-7.691}{170.6}
\FL{-9.196}{-6.265}{121.1}
\FL{-9.383}{-3.664}{69.5}
\FL{-7.717}{-1.589}{35.8}
\FL{-5.275}{-0.470}{15.2}
\FL{-2.612}{-0.058}{3.7}
\FL{1.726}{-3.770}{-37.7}
\FL{4.398}{-4.343}{20.2}
\FL{5.362}{-2.011}{113.6}
\FL{3.180}{-0.316}{163.2}
\FL{-0.539}{2.470}{115.7}
\FL{-2.270}{4.464}{147.7}
\FL{-4.801}{4.927}{-163.7}
\FL{-6.206}{2.915}{-84.1}
\FL{-4.686}{0.831}{-30.5}
\FL{1.771}{2.935}{21.4}
\FL{3.655}{1.357}{-114.8}
\FL{0.775}{3.717}{65.9}
\FL{2.375}{6.087}{45.6}
\FL{4.773}{7.677}{20.3}
\FL{7.596}{7.865}{-15.2}
\FL{9.708}{6.069}{-68.1}
\FL{9.476}{3.296}{-118.4}
\FL{7.431}{1.291}{-150.0}
\FL{4.740}{0.317}{-168.6}
\FL{2.144}{4.147}{31.3}
\FL{4.982}{4.277}{-34.4}
\FL{5.287}{1.624}{-126.2}
\FL{-8.711}{0.935}{-18.5}
\FL{-6.013}{0.302}{-8.6}
\FL{2.601}{-6.283}{-43.0}
\FL{4.865}{-7.672}{-19.0}
\FL{7.473}{-7.834}{14.1}
\FL{9.539}{-6.266}{62.8}
\FL{9.603}{-3.674}{111.9}
\FL{7.948}{-1.648}{144.0}
\FL{5.539}{-0.516}{164.0}
\FL{-1.071}{4.771}{115.9}
\FL{-2.865}{7.344}{134.3}
\FL{-5.431}{9.137}{156.7}
\FL{-8.506}{9.631}{-173.1}
\FL{-6.071}{-8.349}{-170.8}
\FL{-8.963}{-7.843}{147.9}
\FL{-8.164}{-0.705}{14.8}
\FL{2.776}{5.391}{31.8}
\FL{5.374}{6.010}{-8.2}
\FL{7.386}{4.398}{-73.2}
\FL{6.717}{1.876}{-131.1}
\FL{0.248}{-5.387}{-84.7}
\FL{0.619}{-8.464}{-81.6}
\FL{-0.456}{-6.012}{-98.7}
\FL{-0.962}{-8.689}{-102.7}
\FL{3.434}{-5.202}{-16.1}
\FL{6.031}{-4.689}{44.8}
\FL{6.252}{-2.106}{120.6}
\FL{-3.540}{7.016}{147.0}
\FL{-6.429}{7.947}{179.4}
\FL{-9.158}{6.816}{-130.7}
\FL{-9.686}{3.923}{-72.2}
\FL{-7.812}{1.555}{-34.5}
\FL{8.491}{-1.022}{159.3}
\FL{6.028}{-9.431}{-18.5}
\FL{8.714}{-9.704}{8.4}
\FL{3.685}{-6.288}{-24.1}
\FL{6.246}{-6.590}{13.5}
\FL{8.154}{-4.922}{71.8}
\FL{7.709}{-2.429}{124.5}
\FL{-1.234}{6.477}{111.9}
\FL{-2.440}{8.889}{121.3}
\FL{-1.260}{-6.446}{-112.4}
\FL{-2.524}{-8.885}{-122.4}
\FL{-2.834}{5.591}{147.3}
\FL{-5.712}{6.266}{-170.1}
\FL{-7.788}{4.307}{-99.4}
\FL{-6.690}{1.622}{-42.2}
\FL{7.668}{-0.006}{179.9}
\FL{-0.003}{4.728}{90.1}
\FL{-0.010}{8.253}{90.1}
\FL{5.420}{8.531}{18.9}
\FL{8.636}{8.611}{-18.8}
\FL{-2.458}{-3.889}{-161.2}
\FL{-4.899}{-2.016}{73.1}
\FL{-3.084}{-5.628}{-151.8}
\FL{-5.594}{-6.081}{169.2}
\FL{-7.493}{-4.496}{107.6}
\FL{-6.949}{-2.058}{52.0}
\FL{-5.085}{-6.783}{-173.3}
\FL{-7.770}{-5.923}{133.2}
\FL{-8.237}{-3.193}{68.9}
\FL{5.048}{7.044}{10.2}
\FL{8.013}{6.290}{-43.6}
\FL{8.614}{3.359}{-110.6}
\FL{-3.144}{3.947}{178.6}
\FL{-4.866}{2.313}{-84.9}
\FL{4.460}{3.438}{-46.0}
\FL{-3.685}{-5.077}{-170.9}
\FL{-6.056}{-4.151}{121.5}
\FL{-5.714}{-1.627}{49.8}
\FL{-5.179}{6.967}{172.6}
\FL{-7.984}{6.094}{-133.3}
\FL{-8.447}{3.226}{-67.9}
\FL{-2.402}{-7.077}{-128.7}
\FL{-4.995}{-9.321}{-150.3}
\FL{1.787}{6.588}{59.0}
\FL{3.761}{9.030}{42.8}
\FL{-5.679}{8.549}{164.2}
\FL{-8.724}{8.472}{-158.6}
\FL{6.648}{-7.255}{10.3}
\FL{9.063}{-4.574}{89.3}
\FL{-3.343}{2.510}{-130.6}
\FL{-8.542}{-0.008}{0.2}
\FL{2.513}{-8.481}{-56.2}
\FL{1.541}{7.866}{67.5}
\FL{5.163}{-8.451}{-21.5}
\FL{8.542}{-8.669}{17.1}
\FL{3.408}{-2.789}{40.0}
\FL{1.429}{-7.924}{-69.3}
\FL{0.667}{8.485}{81.0}
\FL{-0.770}{8.175}{100.8}
\FL{-8.015}{-9.612}{178.1}
\FL{2.297}{-2.104}{29.1}
\FL{-1.300}{-2.431}{-151.6}
\FL{-3.251}{-1.553}{84.7}
\FL{4.521}{5.251}{-4.9}
\FL{6.508}{3.190}{-92.1}
\FL{8.426}{0.849}{-162.7}
\FL{5.193}{-5.926}{6.4}
\FL{7.360}{-3.640}{91.2}
\FL{-5.291}{5.629}{-167.2}
\FL{-7.053}{3.654}{-93.0}
\FL{-3.429}{8.781}{134.3}
\FL{4.681}{8.916}{30.8}
\FL{8.219}{9.915}{-1.4}
\FL{4.229}{-3.285}{45.3}
\FL{-3.781}{-2.914}{133.7}
\FL{-1.990}{1.662}{-140.3}
\FL{-8.986}{-1.696}{32.7}
}
  \fill[mmarojo] (0,0.5) circle (0.17) (0,-0.5) circle (0.17) (0,0) circle (0.17);
  \node[font=\tiny, text=mmarojo, anchor=west, inner sep=0.5pt] at (0.22,0.55) {$+q$};
  \node[font=\tiny, text=mmarojo, anchor=west, inner sep=0.5pt] at (0.22,-0.55) {$+q$};
  \node[font=\tiny, text=mmarojo, anchor=east, inner sep=0.5pt] at (-0.22,0) {$-2q$};
  \marcoMMA{10}{5}{$x$ (u.\,arb.)}{$z$ (u.\,arb.)}
\end{tikzpicture}\\[1pt]{\small\scshape (b)\; cuadrupolo físico}
\end{minipage}
\caption{Equipotenciales (gris) y líneas de campo (azul) de (a)~el término cuadrupolar puntual y (b)~el cuadrupolo físico lineal del \cref{ej:cuadrupolo-lineal}, con
$+q$ en $z=\pm d$, $-2q$ en el origen y el mismo $Q_{zz}=4qd^{2}$. Como en el caso del dipolo (\cref{fig:fisico-puntual}), lejos de las cargas ambos campos coinciden; el cuadrupolo
físico difiere del puntual en términos de orden $\ell=4$ (hexadecapolares), cuyo peso relativo es $(d/r)^{2}$. Recreación de~\cite{FIGURAS-T2}.}
\label{fig:cuadrupolo-fp}
\end{figure}


% -------------------------------------------------------------------------
\subsection{Un ejemplo sin simetría de revolución: el cuadrupolo plano}\label{sec:cuadrupolo-plano}
% -------------------------------------------------------------------------
\begin{ejemplo}[Cuadrupolo plano de cuatro cargas]\label{ej:cuadrupolo-plano}
Hallar todas las componentes del tensor cuadrupolar del sistema de cuatro cargas situadas en el plano $xy$ a distancia $a$ del origen:
$q_1=+q$ en $(a,0,0)$, $q_2=+q$ en $(-a,0,0)$, $q_3=-q$ en $(0,a,0)$ y $q_4=-q$ en $(0,-a,0)$.

\emph{Momentos de orden inferior.} $Q=q+q-q-q=0$ y $\vp=q(a\ux)+q(-a\ux)-q(a\uy)-q(-a\uy)=\vc{0}$, de modo que el tensor cuadrupolar no depende del origen y el término
cuadrupolar es el dominante.

\emph{Diagonal principal.} Todas las cargas están en el plano $z=0$ y a distancia $r_i=a$ del origen. Con \cref{eq:Qjk-discreto},
\begin{align*}
  Q_{xx}&=q\,(3a^{2}-a^{2})+q\bigl[3(-a)^{2}-a^{2}\bigr]-q\,(0-a^{2})-q\,(0-a^{2})=6qa^{2},\\
  Q_{yy}&=q\,(0-a^{2})+q\,(0-a^{2})-q\,(3a^{2}-a^{2})-q\,(3a^{2}-a^{2})=-6qa^{2},\\
  Q_{zz}&=q\,(0-a^{2})+q\,(0-a^{2})-q\,(0-a^{2})-q\,(0-a^{2})=0 .
\end{align*}
\emph{Fuera de la diagonal.} Cada carga está sobre un eje y en el plano $z=0$, así que los productos cruzados $x_iy_i$, $x_iz_i$, $y_iz_i$ se anulan para todas:
$Q_{xy}=Q_{xz}=Q_{yz}=0$. En consecuencia,
\[
  \Qt=\begin{pmatrix}6qa^{2}&0&0\\0&-6qa^{2}&0\\0&0&0\end{pmatrix},\qquad \Tr\Qt=6qa^{2}-6qa^{2}+0=0\ \checkmark
\]
y el potencial a grandes distancias es $\phi\simeq\frac{1}{8\pi\eo r^{3}}\bigl(6qa^{2}\cd{x}^{2}-6qa^{2}\cd{y}^{2}\bigr)=\frac{3qa^{2}}{4\pi\eo r^{3}}\,\sen^{2}\theta\cos2\varphi$.
\end{ejemplo}

\begin{observacion}[Lectura de los signos]\label{obs:cuadrupolo-plano}
Los signos opuestos de la diagonal dibujan la geometría del campo (\cref{fig:cuadrupolo-plano}): $Q_{xx}=6qa^{2}>0$ indica que el eje $x$, donde están las cargas
positivas, es el que «manda» la repulsión (allí el potencial es positivo y las líneas salen), mientras que $Q_{yy}=-6qa^{2}<0$ hace del eje $y$ el «sumidero»
negativo. Es un cuadrupolo puramente bidimensional, que estira el campo según $x$ y lo comprime según $y$; el sistema no tiene simetría de revolución, y por eso
$Q_{xx}\neq Q_{yy}$ aunque los ejes $x$, $y$, $z$ son ya sus ejes principales.
\end{observacion}

% =========================================================================
% Figura (apuntes manuscritos, pág. 11): líneas de campo del cuadrupolo plano
% de cuatro cargas (+q en (+-a,0), -q en (0,+-a)), en el plano z = 0, con la
% paleta del documento. Datos: figuras/datos/cua_plano_*.tex (L = 4a).
% =========================================================================
\begin{figure}[ht]
\centering
\begin{tikzpicture}[x=0.78cm, y=0.78cm]
  \renewcommand{\LC}[1]{\draw[magacero, line width=0.5pt] plot coordinates {#1};}
  \renewcommand{\FL}[3]{\draw[-{Stealth[length=3.6pt,width=3pt]}, magacero, line width=0.5pt] (#1,#2) -- ++(#3:0.01);}
  \fill[magfondo] (-4,-4) rectangle (4,4);
  \begin{scope}\clip (-4,-4) rectangle (4,4);
    \draw[pizarra!60, line width=0.45pt, densely dashed] (-4,-4) -- (4,4) (-4,4) -- (4,-4);
    % Cuadrupolo plano (cuatro cargas) - líneas de campo
% Generado por scripts/generar_datos.py (no editar a mano).
\LC{(1.069,0.012) (1.128,0.023) (1.187,0.033) (1.246,0.044) (1.305,0.055) (1.364,0.066) (1.423,0.078) (1.482,0.090) (1.540,0.102) (1.599,0.114) (1.658,0.127) (1.716,0.141) (1.775,0.155) (1.833,0.169) (1.891,0.184) (1.949,0.199) (2.007,0.215) (2.064,0.232) (2.122,0.249) (2.179,0.266) (2.237,0.284) (2.294,0.303) (2.350,0.323) (2.407,0.342) (2.463,0.363) (2.520,0.384) (2.575,0.406) (2.631,0.428) (2.687,0.451) (2.742,0.475) (2.797,0.499) (2.851,0.524) (2.905,0.549) (2.959,0.576) (3.013,0.603) (3.066,0.630) (3.119,0.658) (3.172,0.687) (3.224,0.717) (3.276,0.747) (3.327,0.778) (3.378,0.810) (3.429,0.842) (3.479,0.875) (3.528,0.909) (3.578,0.943) (3.626,0.978) (3.674,1.014) (3.722,1.051) (3.769,1.088) (3.815,1.126) (3.861,1.165) (3.907,1.204) (3.951,1.244) (3.995,1.285) (4.000,1.292)}
\LC{(1.061,0.035) (1.113,0.065) (1.164,0.096) (1.216,0.126) (1.267,0.158) (1.318,0.190) (1.368,0.223) (1.417,0.256) (1.466,0.291) (1.514,0.327) (1.562,0.364) (1.608,0.402) (1.654,0.441) (1.698,0.482) (1.741,0.523) (1.783,0.566) (1.823,0.611) (1.862,0.657) (1.899,0.704) (1.934,0.753) (1.967,0.802) (1.998,0.854) (2.027,0.906) (2.054,0.960) (2.078,1.015) (2.099,1.071) (2.118,1.128) (2.133,1.186) (2.145,1.245) (2.153,1.304) (2.158,1.364) (2.159,1.424) (2.156,1.484) (2.149,1.544) (2.137,1.602) (2.121,1.660) (2.101,1.717) (2.076,1.771) (2.046,1.823) (2.012,1.873) (1.974,1.919) (1.932,1.962) (1.887,2.001) (1.838,2.036) (1.787,2.067) (1.733,2.094) (1.677,2.116) (1.620,2.133) (1.561,2.146) (1.502,2.154) (1.442,2.159) (1.382,2.159) (1.322,2.155) (1.262,2.148) (1.203,2.137) (1.145,2.122) (1.088,2.105) (1.031,2.085) (0.976,2.061) (0.922,2.035) (0.869,2.007) (0.817,1.977) (0.767,1.944) (0.718,1.909) (0.671,1.873) (0.624,1.835) (0.579,1.795) (0.536,1.753) (0.494,1.711) (0.453,1.667) (0.413,1.622) (0.375,1.576) (0.338,1.529) (0.302,1.481) (0.267,1.432) (0.233,1.382) (0.199,1.332) (0.167,1.282) (0.136,1.231) (0.105,1.179) (0.074,1.128) (0.044,1.076) (0.039,1.067)}
\LC{(1.045,0.054) (1.083,0.100) (1.121,0.146) (1.159,0.193) (1.195,0.241) (1.231,0.289) (1.265,0.338) (1.298,0.389) (1.328,0.440) (1.357,0.493) (1.383,0.547) (1.407,0.602) (1.428,0.658) (1.445,0.716) (1.458,0.774) (1.467,0.834) (1.472,0.893) (1.471,0.953) (1.466,1.013) (1.454,1.072) (1.436,1.129) (1.413,1.184) (1.383,1.236) (1.347,1.285) (1.306,1.328) (1.260,1.367) (1.210,1.399) (1.156,1.426) (1.100,1.446) (1.042,1.461) (0.982,1.469) (0.922,1.472) (0.862,1.470) (0.803,1.463) (0.744,1.452) (0.686,1.436) (0.629,1.417) (0.573,1.395) (0.519,1.370) (0.465,1.342) (0.413,1.313) (0.362,1.281) (0.313,1.247) (0.264,1.213) (0.216,1.176) (0.169,1.139) (0.122,1.102) (0.076,1.064) (0.060,1.051)}
\LC{(1.024,0.066) (1.044,0.122) (1.064,0.179) (1.083,0.236) (1.100,0.293) (1.116,0.351) (1.129,0.410) (1.140,0.469) (1.148,0.528) (1.152,0.588) (1.152,0.648) (1.147,0.708) (1.136,0.767) (1.119,0.824) (1.097,0.880) (1.067,0.932) (1.031,0.980) (0.990,1.023) (0.942,1.060) (0.891,1.091) (0.836,1.115) (0.779,1.133) (0.720,1.145) (0.660,1.151) (0.600,1.152) (0.541,1.149) (0.481,1.142) (0.422,1.132) (0.363,1.119) (0.305,1.103) (0.248,1.086) (0.191,1.068) (0.134,1.048) (0.077,1.028) (0.068,1.025)}
\LC{(1.000,0.070) (1.000,0.130) (0.999,0.190) (0.997,0.250) (0.993,0.310) (0.988,0.370) (0.980,0.429) (0.969,0.488) (0.955,0.546) (0.936,0.603) (0.912,0.658) (0.883,0.711) (0.849,0.760) (0.809,0.805) (0.765,0.845) (0.716,0.880) (0.664,0.910) (0.609,0.934) (0.552,0.953) (0.494,0.968) (0.435,0.979) (0.376,0.987) (0.316,0.993) (0.256,0.997) (0.196,0.999) (0.136,1.000) (0.076,1.000)}
\LC{(0.976,0.066) (0.955,0.122) (0.934,0.178) (0.912,0.234) (0.889,0.289) (0.864,0.344) (0.837,0.398) (0.809,0.451) (0.778,0.502) (0.744,0.551) (0.707,0.599) (0.667,0.643) (0.624,0.685) (0.578,0.724) (0.530,0.759) (0.479,0.792) (0.427,0.822) (0.374,0.849) (0.320,0.875) (0.265,0.899) (0.209,0.922) (0.153,0.944) (0.097,0.965) (0.069,0.975)}
\LC{(0.955,0.054) (0.916,0.099) (0.877,0.145) (0.838,0.190) (0.798,0.235) (0.757,0.279) (0.716,0.323) (0.675,0.366) (0.633,0.409) (0.591,0.452) (0.549,0.495) (0.506,0.537) (0.464,0.580) (0.421,0.622) (0.378,0.664) (0.335,0.705) (0.291,0.746) (0.247,0.787) (0.202,0.827) (0.157,0.867) (0.112,0.906) (0.066,0.945) (0.058,0.951)}
\LC{(0.939,0.035) (0.887,0.065) (0.835,0.095) (0.783,0.124) (0.731,0.153) (0.678,0.183) (0.626,0.213) (0.575,0.243) (0.524,0.276) (0.475,0.310) (0.428,0.348) (0.384,0.389) (0.344,0.433) (0.307,0.480) (0.272,0.529) (0.240,0.580) (0.210,0.631) (0.180,0.683) (0.151,0.736) (0.121,0.788) (0.092,0.840) (0.062,0.892) (0.037,0.936)}
\LC{(0.931,0.012) (0.872,0.023) (0.813,0.033) (0.754,0.043) (0.695,0.053) (0.636,0.064) (0.577,0.075) (0.518,0.086) (0.459,0.099) (0.401,0.115) (0.344,0.134) (0.290,0.160) (0.240,0.193) (0.197,0.235) (0.163,0.284) (0.137,0.338) (0.117,0.394) (0.101,0.452) (0.088,0.511) (0.076,0.570) (0.065,0.629) (0.055,0.688) (0.044,0.747) (0.034,0.806) (0.024,0.865) (0.013,0.924)}
\LC{(0.931,-0.012) (0.872,-0.023) (0.813,-0.033) (0.754,-0.043) (0.695,-0.053) (0.636,-0.064) (0.577,-0.075) (0.518,-0.086) (0.459,-0.099) (0.401,-0.115) (0.344,-0.134) (0.290,-0.160) (0.240,-0.193) (0.197,-0.235) (0.163,-0.284) (0.137,-0.338) (0.117,-0.394) (0.101,-0.452) (0.088,-0.511) (0.076,-0.570) (0.065,-0.629) (0.055,-0.688) (0.044,-0.747) (0.034,-0.806) (0.024,-0.865) (0.013,-0.924)}
\LC{(0.939,-0.035) (0.887,-0.065) (0.835,-0.095) (0.783,-0.124) (0.731,-0.153) (0.678,-0.183) (0.626,-0.213) (0.575,-0.243) (0.524,-0.276) (0.475,-0.310) (0.428,-0.348) (0.384,-0.389) (0.344,-0.433) (0.307,-0.480) (0.272,-0.529) (0.240,-0.580) (0.210,-0.631) (0.180,-0.683) (0.151,-0.736) (0.121,-0.788) (0.092,-0.840) (0.062,-0.892) (0.037,-0.936)}
\LC{(0.955,-0.054) (0.916,-0.099) (0.877,-0.145) (0.838,-0.190) (0.798,-0.235) (0.757,-0.279) (0.716,-0.323) (0.675,-0.366) (0.633,-0.409) (0.591,-0.452) (0.549,-0.495) (0.506,-0.537) (0.464,-0.580) (0.421,-0.622) (0.378,-0.664) (0.335,-0.705) (0.291,-0.746) (0.247,-0.787) (0.202,-0.827) (0.157,-0.867) (0.112,-0.906) (0.066,-0.945) (0.058,-0.951)}
\LC{(0.976,-0.066) (0.955,-0.122) (0.934,-0.178) (0.912,-0.234) (0.889,-0.289) (0.864,-0.344) (0.837,-0.398) (0.809,-0.451) (0.778,-0.502) (0.744,-0.551) (0.707,-0.599) (0.667,-0.643) (0.624,-0.685) (0.578,-0.724) (0.530,-0.759) (0.479,-0.792) (0.427,-0.822) (0.374,-0.849) (0.320,-0.875) (0.265,-0.899) (0.209,-0.922) (0.153,-0.944) (0.097,-0.965) (0.069,-0.975)}
\LC{(1.000,-0.070) (1.000,-0.130) (0.999,-0.190) (0.997,-0.250) (0.993,-0.310) (0.988,-0.370) (0.980,-0.429) (0.969,-0.488) (0.955,-0.546) (0.936,-0.603) (0.912,-0.658) (0.883,-0.711) (0.849,-0.760) (0.809,-0.805) (0.765,-0.845) (0.716,-0.880) (0.664,-0.910) (0.609,-0.934) (0.552,-0.953) (0.494,-0.968) (0.435,-0.979) (0.376,-0.987) (0.316,-0.993) (0.256,-0.997) (0.196,-0.999) (0.136,-1.000) (0.076,-1.000)}
\LC{(1.024,-0.066) (1.044,-0.122) (1.064,-0.179) (1.083,-0.236) (1.100,-0.293) (1.116,-0.351) (1.129,-0.410) (1.140,-0.469) (1.148,-0.528) (1.152,-0.588) (1.152,-0.648) (1.147,-0.708) (1.136,-0.767) (1.119,-0.824) (1.097,-0.880) (1.067,-0.932) (1.031,-0.980) (0.990,-1.023) (0.942,-1.060) (0.891,-1.091) (0.836,-1.115) (0.779,-1.133) (0.720,-1.145) (0.660,-1.151) (0.600,-1.152) (0.541,-1.149) (0.481,-1.142) (0.422,-1.132) (0.363,-1.119) (0.305,-1.103) (0.248,-1.086) (0.191,-1.068) (0.134,-1.048) (0.077,-1.028) (0.068,-1.025)}
\LC{(1.045,-0.054) (1.083,-0.100) (1.121,-0.146) (1.159,-0.193) (1.195,-0.241) (1.231,-0.289) (1.265,-0.338) (1.298,-0.389) (1.328,-0.440) (1.357,-0.493) (1.383,-0.547) (1.407,-0.602) (1.428,-0.658) (1.445,-0.716) (1.458,-0.774) (1.467,-0.834) (1.472,-0.893) (1.471,-0.953) (1.466,-1.013) (1.454,-1.072) (1.436,-1.129) (1.413,-1.184) (1.383,-1.236) (1.347,-1.285) (1.306,-1.328) (1.260,-1.367) (1.210,-1.399) (1.156,-1.426) (1.100,-1.446) (1.042,-1.461) (0.982,-1.469) (0.922,-1.472) (0.862,-1.470) (0.803,-1.463) (0.744,-1.452) (0.686,-1.436) (0.629,-1.417) (0.573,-1.395) (0.519,-1.370) (0.465,-1.342) (0.413,-1.313) (0.362,-1.281) (0.313,-1.247) (0.264,-1.213) (0.216,-1.176) (0.169,-1.139) (0.122,-1.102) (0.076,-1.064) (0.060,-1.051)}
\LC{(1.061,-0.035) (1.113,-0.065) (1.164,-0.096) (1.216,-0.126) (1.267,-0.158) (1.318,-0.190) (1.368,-0.223) (1.417,-0.256) (1.466,-0.291) (1.514,-0.327) (1.562,-0.364) (1.608,-0.402) (1.654,-0.441) (1.698,-0.482) (1.741,-0.523) (1.783,-0.566) (1.823,-0.611) (1.862,-0.657) (1.899,-0.704) (1.934,-0.753) (1.967,-0.802) (1.998,-0.854) (2.027,-0.906) (2.054,-0.960) (2.078,-1.015) (2.099,-1.071) (2.118,-1.128) (2.133,-1.186) (2.145,-1.245) (2.153,-1.304) (2.158,-1.364) (2.159,-1.424) (2.156,-1.484) (2.149,-1.544) (2.137,-1.602) (2.121,-1.660) (2.101,-1.717) (2.076,-1.771) (2.046,-1.823) (2.012,-1.873) (1.974,-1.919) (1.932,-1.962) (1.887,-2.001) (1.838,-2.036) (1.787,-2.067) (1.733,-2.094) (1.677,-2.116) (1.620,-2.133) (1.561,-2.146) (1.502,-2.154) (1.442,-2.159) (1.382,-2.159) (1.322,-2.155) (1.262,-2.148) (1.203,-2.137) (1.145,-2.122) (1.088,-2.105) (1.031,-2.085) (0.976,-2.061) (0.922,-2.035) (0.869,-2.007) (0.817,-1.977) (0.767,-1.944) (0.718,-1.909) (0.671,-1.873) (0.624,-1.835) (0.579,-1.795) (0.536,-1.753) (0.494,-1.711) (0.453,-1.667) (0.413,-1.622) (0.375,-1.576) (0.338,-1.529) (0.302,-1.481) (0.267,-1.432) (0.233,-1.382) (0.199,-1.332) (0.167,-1.282) (0.136,-1.231) (0.105,-1.179) (0.074,-1.128) (0.044,-1.076) (0.039,-1.067)}
\LC{(1.069,-0.012) (1.128,-0.023) (1.187,-0.033) (1.246,-0.044) (1.305,-0.055) (1.364,-0.066) (1.423,-0.078) (1.482,-0.090) (1.540,-0.102) (1.599,-0.114) (1.658,-0.127) (1.716,-0.141) (1.775,-0.155) (1.833,-0.169) (1.891,-0.184) (1.949,-0.199) (2.007,-0.215) (2.064,-0.232) (2.122,-0.249) (2.179,-0.266) (2.237,-0.284) (2.294,-0.303) (2.350,-0.323) (2.407,-0.342) (2.463,-0.363) (2.520,-0.384) (2.575,-0.406) (2.631,-0.428) (2.687,-0.451) (2.742,-0.475) (2.797,-0.499) (2.851,-0.524) (2.905,-0.549) (2.959,-0.576) (3.013,-0.603) (3.066,-0.630) (3.119,-0.658) (3.172,-0.687) (3.224,-0.717) (3.276,-0.747) (3.327,-0.778) (3.378,-0.810) (3.429,-0.842) (3.479,-0.875) (3.528,-0.909) (3.578,-0.943) (3.626,-0.978) (3.674,-1.014) (3.722,-1.051) (3.769,-1.088) (3.815,-1.126) (3.861,-1.165) (3.907,-1.204) (3.951,-1.244) (3.995,-1.285) (4.000,-1.292)}
\LC{(-0.931,0.012) (-0.872,0.023) (-0.813,0.033) (-0.754,0.043) (-0.695,0.053) (-0.636,0.064) (-0.577,0.075) (-0.518,0.086) (-0.459,0.099) (-0.401,0.115) (-0.344,0.134) (-0.290,0.160) (-0.240,0.193) (-0.197,0.235) (-0.163,0.284) (-0.137,0.338) (-0.117,0.394) (-0.101,0.452) (-0.088,0.511) (-0.076,0.570) (-0.065,0.629) (-0.055,0.688) (-0.044,0.747) (-0.034,0.806) (-0.024,0.865) (-0.013,0.924)}
\LC{(-0.939,0.035) (-0.887,0.065) (-0.835,0.095) (-0.783,0.124) (-0.731,0.153) (-0.678,0.183) (-0.626,0.213) (-0.575,0.243) (-0.524,0.276) (-0.475,0.310) (-0.428,0.348) (-0.384,0.389) (-0.344,0.433) (-0.307,0.480) (-0.272,0.529) (-0.240,0.580) (-0.210,0.631) (-0.180,0.683) (-0.151,0.736) (-0.121,0.788) (-0.092,0.840) (-0.062,0.892) (-0.037,0.936)}
\LC{(-0.955,0.054) (-0.916,0.099) (-0.877,0.145) (-0.838,0.190) (-0.798,0.235) (-0.757,0.279) (-0.716,0.323) (-0.675,0.366) (-0.633,0.409) (-0.591,0.452) (-0.549,0.495) (-0.506,0.537) (-0.464,0.580) (-0.421,0.622) (-0.378,0.664) (-0.335,0.705) (-0.291,0.746) (-0.247,0.787) (-0.202,0.827) (-0.157,0.867) (-0.112,0.906) (-0.066,0.945) (-0.058,0.951)}
\LC{(-0.976,0.066) (-0.955,0.122) (-0.934,0.178) (-0.912,0.234) (-0.889,0.289) (-0.864,0.344) (-0.837,0.398) (-0.809,0.451) (-0.778,0.502) (-0.744,0.551) (-0.707,0.599) (-0.667,0.643) (-0.624,0.685) (-0.578,0.724) (-0.530,0.759) (-0.479,0.792) (-0.427,0.822) (-0.374,0.849) (-0.320,0.875) (-0.265,0.899) (-0.209,0.922) (-0.153,0.944) (-0.097,0.965) (-0.069,0.975)}
\LC{(-1.000,0.070) (-1.000,0.130) (-0.999,0.190) (-0.997,0.250) (-0.993,0.310) (-0.988,0.370) (-0.980,0.429) (-0.969,0.488) (-0.955,0.546) (-0.936,0.603) (-0.912,0.658) (-0.883,0.711) (-0.849,0.760) (-0.809,0.805) (-0.765,0.845) (-0.716,0.880) (-0.664,0.910) (-0.609,0.934) (-0.552,0.953) (-0.494,0.968) (-0.435,0.979) (-0.376,0.987) (-0.316,0.993) (-0.256,0.997) (-0.196,0.999) (-0.136,1.000) (-0.076,1.000)}
\LC{(-1.024,0.066) (-1.044,0.122) (-1.064,0.179) (-1.083,0.236) (-1.100,0.293) (-1.116,0.351) (-1.129,0.410) (-1.140,0.469) (-1.148,0.528) (-1.152,0.588) (-1.152,0.648) (-1.147,0.708) (-1.136,0.767) (-1.119,0.824) (-1.097,0.880) (-1.067,0.932) (-1.031,0.980) (-0.990,1.023) (-0.942,1.060) (-0.891,1.091) (-0.836,1.115) (-0.779,1.133) (-0.720,1.145) (-0.660,1.151) (-0.600,1.152) (-0.541,1.149) (-0.481,1.142) (-0.422,1.132) (-0.363,1.119) (-0.305,1.103) (-0.248,1.086) (-0.191,1.068) (-0.134,1.048) (-0.077,1.028) (-0.068,1.025)}
\LC{(-1.045,0.054) (-1.083,0.100) (-1.121,0.146) (-1.159,0.193) (-1.195,0.241) (-1.231,0.289) (-1.265,0.338) (-1.298,0.389) (-1.328,0.440) (-1.357,0.493) (-1.383,0.547) (-1.407,0.602) (-1.428,0.658) (-1.445,0.716) (-1.458,0.774) (-1.467,0.834) (-1.472,0.893) (-1.471,0.953) (-1.466,1.013) (-1.454,1.072) (-1.436,1.129) (-1.413,1.184) (-1.383,1.236) (-1.347,1.285) (-1.306,1.328) (-1.260,1.367) (-1.210,1.399) (-1.156,1.426) (-1.100,1.446) (-1.042,1.461) (-0.982,1.469) (-0.922,1.472) (-0.862,1.470) (-0.803,1.463) (-0.744,1.452) (-0.686,1.436) (-0.629,1.417) (-0.573,1.395) (-0.519,1.370) (-0.465,1.342) (-0.413,1.313) (-0.362,1.281) (-0.313,1.247) (-0.264,1.213) (-0.216,1.176) (-0.169,1.139) (-0.122,1.102) (-0.076,1.064) (-0.060,1.051)}
\LC{(-1.061,0.035) (-1.113,0.065) (-1.164,0.096) (-1.216,0.126) (-1.267,0.158) (-1.318,0.190) (-1.368,0.223) (-1.417,0.256) (-1.466,0.291) (-1.514,0.327) (-1.562,0.364) (-1.608,0.402) (-1.654,0.441) (-1.698,0.482) (-1.741,0.523) (-1.783,0.566) (-1.823,0.611) (-1.862,0.657) (-1.899,0.704) (-1.934,0.753) (-1.967,0.802) (-1.998,0.854) (-2.027,0.906) (-2.054,0.960) (-2.078,1.015) (-2.099,1.071) (-2.118,1.128) (-2.133,1.186) (-2.145,1.245) (-2.153,1.304) (-2.158,1.364) (-2.159,1.424) (-2.156,1.484) (-2.149,1.544) (-2.137,1.602) (-2.121,1.660) (-2.101,1.717) (-2.076,1.771) (-2.046,1.823) (-2.012,1.873) (-1.974,1.919) (-1.932,1.962) (-1.887,2.001) (-1.838,2.036) (-1.787,2.067) (-1.733,2.094) (-1.677,2.116) (-1.620,2.133) (-1.561,2.146) (-1.502,2.154) (-1.442,2.159) (-1.382,2.159) (-1.322,2.155) (-1.262,2.148) (-1.203,2.137) (-1.145,2.122) (-1.088,2.105) (-1.031,2.085) (-0.976,2.061) (-0.922,2.035) (-0.869,2.007) (-0.817,1.977) (-0.767,1.944) (-0.718,1.909) (-0.671,1.873) (-0.624,1.835) (-0.579,1.795) (-0.536,1.753) (-0.494,1.711) (-0.453,1.667) (-0.413,1.622) (-0.375,1.576) (-0.338,1.529) (-0.302,1.481) (-0.267,1.432) (-0.233,1.382) (-0.199,1.332) (-0.167,1.282) (-0.136,1.231) (-0.105,1.179) (-0.074,1.128) (-0.044,1.076) (-0.039,1.067)}
\LC{(-1.069,0.012) (-1.128,0.023) (-1.187,0.033) (-1.246,0.044) (-1.305,0.055) (-1.364,0.066) (-1.423,0.078) (-1.482,0.090) (-1.540,0.102) (-1.599,0.114) (-1.658,0.127) (-1.716,0.141) (-1.775,0.155) (-1.833,0.169) (-1.891,0.184) (-1.949,0.199) (-2.007,0.215) (-2.064,0.232) (-2.122,0.249) (-2.179,0.266) (-2.237,0.284) (-2.294,0.303) (-2.350,0.323) (-2.407,0.342) (-2.463,0.363) (-2.520,0.384) (-2.575,0.406) (-2.631,0.428) (-2.687,0.451) (-2.742,0.475) (-2.797,0.499) (-2.851,0.524) (-2.905,0.549) (-2.959,0.576) (-3.013,0.603) (-3.066,0.630) (-3.119,0.658) (-3.172,0.687) (-3.224,0.717) (-3.276,0.747) (-3.327,0.778) (-3.378,0.810) (-3.429,0.842) (-3.479,0.875) (-3.528,0.909) (-3.578,0.943) (-3.626,0.978) (-3.674,1.014) (-3.722,1.051) (-3.769,1.088) (-3.815,1.126) (-3.861,1.165) (-3.907,1.204) (-3.951,1.244) (-3.995,1.285) (-4.000,1.292)}
\LC{(-1.069,-0.012) (-1.128,-0.023) (-1.187,-0.033) (-1.246,-0.044) (-1.305,-0.055) (-1.364,-0.066) (-1.423,-0.078) (-1.482,-0.090) (-1.540,-0.102) (-1.599,-0.114) (-1.658,-0.127) (-1.716,-0.141) (-1.775,-0.155) (-1.833,-0.169) (-1.891,-0.184) (-1.949,-0.199) (-2.007,-0.215) (-2.064,-0.232) (-2.122,-0.249) (-2.179,-0.266) (-2.237,-0.284) (-2.294,-0.303) (-2.350,-0.323) (-2.407,-0.342) (-2.463,-0.363) (-2.520,-0.384) (-2.575,-0.406) (-2.631,-0.428) (-2.687,-0.451) (-2.742,-0.475) (-2.797,-0.499) (-2.851,-0.524) (-2.905,-0.549) (-2.959,-0.576) (-3.013,-0.603) (-3.066,-0.630) (-3.119,-0.658) (-3.172,-0.687) (-3.224,-0.717) (-3.276,-0.747) (-3.327,-0.778) (-3.378,-0.810) (-3.429,-0.842) (-3.479,-0.875) (-3.528,-0.909) (-3.578,-0.943) (-3.626,-0.978) (-3.674,-1.014) (-3.722,-1.051) (-3.769,-1.088) (-3.815,-1.126) (-3.861,-1.165) (-3.907,-1.204) (-3.951,-1.244) (-3.995,-1.285) (-4.000,-1.292)}
\LC{(-1.061,-0.035) (-1.113,-0.065) (-1.164,-0.096) (-1.216,-0.126) (-1.267,-0.158) (-1.318,-0.190) (-1.368,-0.223) (-1.417,-0.256) (-1.466,-0.291) (-1.514,-0.327) (-1.562,-0.364) (-1.608,-0.402) (-1.654,-0.441) (-1.698,-0.482) (-1.741,-0.523) (-1.783,-0.566) (-1.823,-0.611) (-1.862,-0.657) (-1.899,-0.704) (-1.934,-0.753) (-1.967,-0.802) (-1.998,-0.854) (-2.027,-0.906) (-2.054,-0.960) (-2.078,-1.015) (-2.099,-1.071) (-2.118,-1.128) (-2.133,-1.186) (-2.145,-1.245) (-2.153,-1.304) (-2.158,-1.364) (-2.159,-1.424) (-2.156,-1.484) (-2.149,-1.544) (-2.137,-1.602) (-2.121,-1.660) (-2.101,-1.717) (-2.076,-1.771) (-2.046,-1.823) (-2.012,-1.873) (-1.974,-1.919) (-1.932,-1.962) (-1.887,-2.001) (-1.838,-2.036) (-1.787,-2.067) (-1.733,-2.094) (-1.677,-2.116) (-1.620,-2.133) (-1.561,-2.146) (-1.502,-2.154) (-1.442,-2.159) (-1.382,-2.159) (-1.322,-2.155) (-1.262,-2.148) (-1.203,-2.137) (-1.145,-2.122) (-1.088,-2.105) (-1.031,-2.085) (-0.976,-2.061) (-0.922,-2.035) (-0.869,-2.007) (-0.817,-1.977) (-0.767,-1.944) (-0.718,-1.909) (-0.671,-1.873) (-0.624,-1.835) (-0.579,-1.795) (-0.536,-1.753) (-0.494,-1.711) (-0.453,-1.667) (-0.413,-1.622) (-0.375,-1.576) (-0.338,-1.529) (-0.302,-1.481) (-0.267,-1.432) (-0.233,-1.382) (-0.199,-1.332) (-0.167,-1.282) (-0.136,-1.231) (-0.105,-1.179) (-0.074,-1.128) (-0.044,-1.076) (-0.039,-1.067)}
\LC{(-1.045,-0.054) (-1.083,-0.100) (-1.121,-0.146) (-1.159,-0.193) (-1.195,-0.241) (-1.231,-0.289) (-1.265,-0.338) (-1.298,-0.389) (-1.328,-0.440) (-1.357,-0.493) (-1.383,-0.547) (-1.407,-0.602) (-1.428,-0.658) (-1.445,-0.716) (-1.458,-0.774) (-1.467,-0.834) (-1.472,-0.893) (-1.471,-0.953) (-1.466,-1.013) (-1.454,-1.072) (-1.436,-1.129) (-1.413,-1.184) (-1.383,-1.236) (-1.347,-1.285) (-1.306,-1.328) (-1.260,-1.367) (-1.210,-1.399) (-1.156,-1.426) (-1.100,-1.446) (-1.042,-1.461) (-0.982,-1.469) (-0.922,-1.472) (-0.862,-1.470) (-0.803,-1.463) (-0.744,-1.452) (-0.686,-1.436) (-0.629,-1.417) (-0.573,-1.395) (-0.519,-1.370) (-0.465,-1.342) (-0.413,-1.313) (-0.362,-1.281) (-0.313,-1.247) (-0.264,-1.213) (-0.216,-1.176) (-0.169,-1.139) (-0.122,-1.102) (-0.076,-1.064) (-0.060,-1.051)}
\LC{(-1.024,-0.066) (-1.044,-0.122) (-1.064,-0.179) (-1.083,-0.236) (-1.100,-0.293) (-1.116,-0.351) (-1.129,-0.410) (-1.140,-0.469) (-1.148,-0.528) (-1.152,-0.588) (-1.152,-0.648) (-1.147,-0.708) (-1.136,-0.767) (-1.119,-0.824) (-1.097,-0.880) (-1.067,-0.932) (-1.031,-0.980) (-0.990,-1.023) (-0.942,-1.060) (-0.891,-1.091) (-0.836,-1.115) (-0.779,-1.133) (-0.720,-1.145) (-0.660,-1.151) (-0.600,-1.152) (-0.541,-1.149) (-0.481,-1.142) (-0.422,-1.132) (-0.363,-1.119) (-0.305,-1.103) (-0.248,-1.086) (-0.191,-1.068) (-0.134,-1.048) (-0.077,-1.028) (-0.068,-1.025)}
\LC{(-1.000,-0.070) (-1.000,-0.130) (-0.999,-0.190) (-0.997,-0.250) (-0.993,-0.310) (-0.988,-0.370) (-0.980,-0.429) (-0.969,-0.488) (-0.955,-0.546) (-0.936,-0.603) (-0.912,-0.658) (-0.883,-0.711) (-0.849,-0.760) (-0.809,-0.805) (-0.765,-0.845) (-0.716,-0.880) (-0.664,-0.910) (-0.609,-0.934) (-0.552,-0.953) (-0.494,-0.968) (-0.435,-0.979) (-0.376,-0.987) (-0.316,-0.993) (-0.256,-0.997) (-0.196,-0.999) (-0.136,-1.000) (-0.076,-1.000)}
\LC{(-0.976,-0.066) (-0.955,-0.122) (-0.934,-0.178) (-0.912,-0.234) (-0.889,-0.289) (-0.864,-0.344) (-0.837,-0.398) (-0.809,-0.451) (-0.778,-0.502) (-0.744,-0.551) (-0.707,-0.599) (-0.667,-0.643) (-0.624,-0.685) (-0.578,-0.724) (-0.530,-0.759) (-0.479,-0.792) (-0.427,-0.822) (-0.374,-0.849) (-0.320,-0.875) (-0.265,-0.899) (-0.209,-0.922) (-0.153,-0.944) (-0.097,-0.965) (-0.069,-0.975)}
\LC{(-0.955,-0.054) (-0.916,-0.099) (-0.877,-0.145) (-0.838,-0.190) (-0.798,-0.235) (-0.757,-0.279) (-0.716,-0.323) (-0.675,-0.366) (-0.633,-0.409) (-0.591,-0.452) (-0.549,-0.495) (-0.506,-0.537) (-0.464,-0.580) (-0.421,-0.622) (-0.378,-0.664) (-0.335,-0.705) (-0.291,-0.746) (-0.247,-0.787) (-0.202,-0.827) (-0.157,-0.867) (-0.112,-0.906) (-0.066,-0.945) (-0.058,-0.951)}
\LC{(-0.939,-0.035) (-0.887,-0.065) (-0.835,-0.095) (-0.783,-0.124) (-0.731,-0.153) (-0.678,-0.183) (-0.626,-0.213) (-0.575,-0.243) (-0.524,-0.276) (-0.475,-0.310) (-0.428,-0.348) (-0.384,-0.389) (-0.344,-0.433) (-0.307,-0.480) (-0.272,-0.529) (-0.240,-0.580) (-0.210,-0.631) (-0.180,-0.683) (-0.151,-0.736) (-0.121,-0.788) (-0.092,-0.840) (-0.062,-0.892) (-0.037,-0.936)}
\LC{(-0.931,-0.012) (-0.872,-0.023) (-0.813,-0.033) (-0.754,-0.043) (-0.695,-0.053) (-0.636,-0.064) (-0.577,-0.075) (-0.518,-0.086) (-0.459,-0.099) (-0.401,-0.115) (-0.344,-0.134) (-0.290,-0.160) (-0.240,-0.193) (-0.197,-0.235) (-0.163,-0.284) (-0.137,-0.338) (-0.117,-0.394) (-0.101,-0.452) (-0.088,-0.511) (-0.076,-0.570) (-0.065,-0.629) (-0.055,-0.688) (-0.044,-0.747) (-0.034,-0.806) (-0.024,-0.865) (-0.013,-0.924)}
\LC{(-2.445,-0.677) (-2.492,-0.715) (-2.538,-0.753) (-2.583,-0.793) (-2.627,-0.834) (-2.670,-0.875) (-2.712,-0.918) (-2.754,-0.962) (-2.794,-1.006) (-2.832,-1.052) (-2.870,-1.099) (-2.906,-1.146) (-2.941,-1.195) (-2.975,-1.245) (-3.007,-1.296) (-3.038,-1.347) (-3.067,-1.400) (-3.094,-1.453) (-3.119,-1.507) (-3.143,-1.563) (-3.165,-1.618) (-3.185,-1.675) (-3.203,-1.732) (-3.218,-1.790) (-3.232,-1.849) (-3.243,-1.908) (-3.252,-1.967) (-3.258,-2.027) (-3.262,-2.087) (-3.263,-2.147) (-3.261,-2.207) (-3.257,-2.266) (-3.249,-2.326) (-3.239,-2.385) (-3.226,-2.444) (-3.210,-2.501) (-3.191,-2.558) (-3.169,-2.614) (-3.144,-2.669) (-3.116,-2.722) (-3.085,-2.773) (-3.051,-2.822) (-3.014,-2.870) (-2.975,-2.915) (-2.933,-2.958) (-2.888,-2.998) (-2.842,-3.036) (-2.793,-3.071) (-2.742,-3.103) (-2.690,-3.133) (-2.636,-3.159) (-2.581,-3.182) (-2.524,-3.203) (-2.467,-3.220) (-2.409,-3.234) (-2.350,-3.246) (-2.290,-3.254) (-2.231,-3.260) (-2.171,-3.262) (-2.111,-3.262) (-2.051,-3.260) (-1.991,-3.255) (-1.932,-3.247) (-1.872,-3.237) (-1.814,-3.224) (-1.756,-3.209) (-1.698,-3.192) (-1.641,-3.173) (-1.585,-3.152) (-1.530,-3.129) (-1.475,-3.104) (-1.421,-3.078) (-1.368,-3.050) (-1.316,-3.020) (-1.265,-2.988) (-1.215,-2.955) (-1.166,-2.921) (-1.118,-2.885) (-1.071,-2.848) (-1.025,-2.809) (-0.979,-2.770) (-0.935,-2.729) (-0.892,-2.687) (-0.850,-2.645) (-0.809,-2.601) (-0.769,-2.556) (-0.730,-2.510) (-0.692,-2.464) (-0.656,-2.416) (-0.620,-2.368) (-0.585,-2.320) (-0.551,-2.270) (-0.518,-2.220) (-0.486,-2.169) (-0.455,-2.118) (-0.426,-2.066) (-0.421,-2.057)}
\LC{(2.872,-0.869) (2.916,-0.908) (2.960,-0.949) (3.004,-0.991) (3.046,-1.033) (3.087,-1.077) (3.127,-1.121) (3.167,-1.167) (3.205,-1.213) (3.242,-1.260) (3.278,-1.308) (3.312,-1.357) (3.346,-1.407) (3.378,-1.458) (3.408,-1.510) (3.438,-1.562) (3.465,-1.615) (3.492,-1.669) (3.516,-1.724) (3.539,-1.779) (3.560,-1.835) (3.580,-1.892) (3.597,-1.950) (3.613,-2.007) (3.627,-2.066) (3.638,-2.125) (3.648,-2.184) (3.655,-2.244) (3.660,-2.303) (3.663,-2.363) (3.663,-2.423) (3.661,-2.483) (3.657,-2.543) (3.650,-2.603) (3.640,-2.662) (3.628,-2.721) (3.613,-2.779) (3.596,-2.836) (3.575,-2.893) (3.553,-2.948) (3.527,-3.002) (3.499,-3.055) (3.468,-3.107) (3.434,-3.156) (3.399,-3.204) (3.360,-3.251) (3.320,-3.295) (3.277,-3.337) (3.232,-3.376) (3.185,-3.414) (3.136,-3.449) (3.085,-3.481) (3.033,-3.511) (2.980,-3.538) (2.925,-3.562) (2.869,-3.584) (2.812,-3.603) (2.755,-3.620) (2.696,-3.634) (2.637,-3.645) (2.578,-3.653) (2.518,-3.659) (2.458,-3.663) (2.398,-3.664) (2.338,-3.662) (2.278,-3.658) (2.219,-3.652) (2.159,-3.644) (2.100,-3.634) (2.042,-3.621) (1.983,-3.607) (1.926,-3.590) (1.869,-3.572) (1.812,-3.552) (1.756,-3.530) (1.701,-3.506) (1.647,-3.481) (1.593,-3.454) (1.540,-3.426) (1.488,-3.396) (1.437,-3.365) (1.386,-3.332) (1.337,-3.298) (1.288,-3.263) (1.241,-3.227) (1.194,-3.189) (1.148,-3.151) (1.103,-3.111) (1.059,-3.070) (1.016,-3.028) (0.973,-2.986) (0.932,-2.942) (0.892,-2.898) (0.852,-2.853) (0.814,-2.807) (0.776,-2.760) (0.740,-2.712) (0.704,-2.664) (0.669,-2.615) (0.635,-2.566) (0.602,-2.516) (0.570,-2.465) (0.539,-2.414) (0.509,-2.362) (0.480,-2.309) (0.451,-2.257) (0.423,-2.203) (0.397,-2.150) (0.371,-2.096) (0.345,-2.041) (0.321,-1.986) (0.305,-1.950)}
\LC{(0.641,-4.000) (0.620,-3.952) (0.598,-3.896) (0.577,-3.839) (0.557,-3.783) (0.537,-3.726) (0.517,-3.670) (0.498,-3.613) (0.479,-3.556) (0.461,-3.499) (0.443,-3.442) (0.425,-3.384) (0.408,-3.327) (0.391,-3.269) (0.375,-3.211) (0.359,-3.154) (0.343,-3.096) (0.328,-3.038) (0.313,-2.980) (0.299,-2.921) (0.284,-2.863) (0.271,-2.805) (0.257,-2.746) (0.244,-2.688) (0.232,-2.629) (0.219,-2.570) (0.207,-2.511) (0.196,-2.453) (0.186,-2.403)}
\LC{(-3.072,0.978) (-3.116,1.019) (-3.158,1.062) (-3.200,1.105) (-3.241,1.149) (-3.281,1.194) (-3.319,1.240) (-3.357,1.286) (-3.393,1.334) (-3.429,1.382) (-3.463,1.432) (-3.496,1.482) (-3.528,1.533) (-3.558,1.584) (-3.587,1.637) (-3.615,1.690) (-3.641,1.744) (-3.665,1.799) (-3.688,1.855) (-3.709,1.911) (-3.729,1.967) (-3.746,2.025) (-3.762,2.083) (-3.776,2.141) (-3.789,2.200) (-3.799,2.259) (-3.807,2.318) (-3.812,2.378) (-3.816,2.438) (-3.818,2.498) (-3.817,2.558) (-3.813,2.618) (-3.808,2.677) (-3.800,2.737) (-3.789,2.796) (-3.776,2.855) (-3.760,2.912) (-3.742,2.970) (-3.721,3.026) (-3.698,3.081) (-3.672,3.135) (-3.643,3.188) (-3.612,3.239) (-3.579,3.289) (-3.543,3.337) (-3.505,3.383) (-3.464,3.428) (-3.422,3.470) (-3.377,3.510) (-3.330,3.548) (-3.282,3.584) (-3.232,3.617) (-3.181,3.647) (-3.128,3.675) (-3.073,3.701) (-3.018,3.724) (-2.961,3.745) (-2.904,3.763) (-2.846,3.778) (-2.788,3.791) (-2.728,3.801) (-2.669,3.809) (-2.609,3.814) (-2.549,3.817) (-2.489,3.817) (-2.429,3.816) (-2.369,3.812) (-2.310,3.806) (-2.250,3.797) (-2.191,3.787) (-2.133,3.775) (-2.074,3.760) (-2.017,3.744) (-1.959,3.726) (-1.903,3.706) (-1.847,3.685) (-1.791,3.662) (-1.737,3.637) (-1.683,3.611) (-1.629,3.583) (-1.577,3.554) (-1.525,3.523) (-1.475,3.491) (-1.425,3.458) (-1.375,3.424) (-1.327,3.388) (-1.280,3.352) (-1.233,3.314) (-1.187,3.275) (-1.142,3.235) (-1.099,3.194) (-1.055,3.152) (-1.013,3.110) (-0.972,3.066) (-0.932,3.022) (-0.892,2.977) (-0.854,2.931) (-0.816,2.884) (-0.779,2.836) (-0.743,2.788) (-0.709,2.740) (-0.674,2.690) (-0.641,2.640) (-0.609,2.590) (-0.577,2.539) (-0.547,2.487) (-0.517,2.435) (-0.488,2.382) (-0.460,2.329) (-0.433,2.276) (-0.407,2.222) (-0.381,2.168) (-0.356,2.113) (-0.332,2.058) (-0.309,2.003) (-0.286,1.947) (-0.275,1.919)}
\LC{(2.400,0.661) (2.447,0.699) (2.493,0.737) (2.538,0.776) (2.582,0.817) (2.626,0.859) (2.668,0.901) (2.709,0.945) (2.749,0.989) (2.788,1.035) (2.825,1.082) (2.862,1.130) (2.897,1.179) (2.930,1.228) (2.962,1.279) (2.992,1.331) (3.021,1.383) (3.048,1.437) (3.074,1.491) (3.097,1.547) (3.119,1.603) (3.138,1.659) (3.155,1.717) (3.171,1.775) (3.184,1.833) (3.194,1.893) (3.202,1.952) (3.208,2.012) (3.211,2.072) (3.211,2.132) (3.208,2.192) (3.203,2.251) (3.195,2.311) (3.183,2.370) (3.169,2.428) (3.152,2.485) (3.131,2.542) (3.108,2.597) (3.081,2.651) (3.052,2.703) (3.019,2.754) (2.984,2.802) (2.946,2.848) (2.905,2.892) (2.862,2.934) (2.816,2.973) (2.768,3.009) (2.719,3.042) (2.667,3.073) (2.613,3.100) (2.559,3.125) (2.503,3.146) (2.445,3.164) (2.387,3.179) (2.329,3.192) (2.269,3.201) (2.210,3.207) (2.150,3.210) (2.090,3.211) (2.030,3.209) (1.970,3.204) (1.911,3.197) (1.851,3.187) (1.793,3.175) (1.734,3.160) (1.677,3.144) (1.620,3.125) (1.564,3.104) (1.508,3.081) (1.453,3.056) (1.400,3.030) (1.347,3.001) (1.295,2.971) (1.244,2.940) (1.194,2.907) (1.144,2.872) (1.096,2.837) (1.049,2.799) (1.003,2.761) (0.958,2.721) (0.914,2.680) (0.871,2.639) (0.829,2.596) (0.789,2.552) (0.749,2.507) (0.710,2.461) (0.672,2.414) (0.636,2.367) (0.600,2.318) (0.565,2.269) (0.532,2.220) (0.499,2.169) (0.468,2.118) (0.442,2.075)}
\LC{(-3.089,1.434) (-3.115,1.488) (-3.139,1.543) (-3.162,1.598) (-3.183,1.654) (-3.201,1.711) (-3.218,1.769) (-3.232,1.827) (-3.244,1.886) (-3.254,1.945) (-3.261,2.005) (-3.266,2.065) (-3.269,2.125) (-3.268,2.185) (-3.265,2.245) (-3.259,2.304) (-3.250,2.364) (-3.238,2.422) (-3.223,2.481) (-3.206,2.538) (-3.185,2.594) (-3.161,2.649) (-3.134,2.703) (-3.105,2.755) (-3.072,2.805) (-3.036,2.854) (-2.998,2.900) (-2.958,2.944) (-2.914,2.986) (-2.869,3.025) (-2.821,3.061) (-2.771,3.094) (-2.720,3.125) (-2.666,3.153) (-2.612,3.178) (-2.556,3.199) (-2.499,3.218) (-2.441,3.234) (-2.382,3.247) (-2.323,3.256) (-2.264,3.263) (-2.204,3.267) (-2.144,3.269) (-2.084,3.267) (-2.024,3.263) (-1.964,3.257) (-1.905,3.248) (-1.846,3.236) (-1.788,3.223) (-1.730,3.207) (-1.673,3.189) (-1.616,3.169) (-1.560,3.147) (-1.505,3.123) (-1.451,3.097) (-1.433,3.088)}
\LC{(-1.162,-4.000) (-1.124,-3.959) (-1.088,-3.912) (-1.052,-3.864) (-1.016,-3.816) (-0.981,-3.767) (-0.947,-3.717) (-0.914,-3.667) (-0.881,-3.617) (-0.849,-3.566) (-0.818,-3.515) (-0.787,-3.464) (-0.757,-3.412) (-0.727,-3.360) (-0.698,-3.307) (-0.670,-3.254) (-0.643,-3.201) (-0.616,-3.147) (-0.589,-3.093) (-0.564,-3.039) (-0.539,-2.985) (-0.514,-2.930) (-0.490,-2.875) (-0.467,-2.819) (-0.445,-2.764) (-0.423,-2.708) (-0.401,-2.652) (-0.380,-2.596) (-0.360,-2.539) (-0.350,-2.511)}
\LC{(-0.859,4.000) (-0.831,3.952) (-0.803,3.899) (-0.775,3.846) (-0.748,3.792) (-0.722,3.738) (-0.696,3.684) (-0.670,3.630) (-0.646,3.575) (-0.621,3.520) (-0.597,3.465) (-0.574,3.410) (-0.551,3.354) (-0.529,3.299) (-0.507,3.243) (-0.486,3.187) (-0.465,3.130) (-0.445,3.074) (-0.425,3.017) (-0.406,2.961) (-0.387,2.904) (-0.369,2.846) (-0.351,2.789) (-0.333,2.732) (-0.325,2.703)}
\LC{(3.348,1.368) (3.382,1.418) (3.414,1.468) (3.445,1.520) (3.474,1.572) (3.502,1.625) (3.529,1.679) (3.554,1.734) (3.577,1.789) (3.598,1.845) (3.618,1.901) (3.636,1.959) (3.652,2.016) (3.667,2.075) (3.679,2.133) (3.689,2.193) (3.697,2.252) (3.703,2.312) (3.706,2.372) (3.708,2.432) (3.706,2.492) (3.703,2.551) (3.697,2.611) (3.688,2.671) (3.677,2.729) (3.663,2.788) (3.647,2.846) (3.628,2.902) (3.606,2.958) (3.582,3.013) (3.555,3.067) (3.525,3.119) (3.493,3.170) (3.458,3.219) (3.421,3.266) (3.382,3.311) (3.340,3.354) (3.296,3.395) (3.250,3.434) (3.202,3.470) (3.153,3.504) (3.102,3.535) (3.049,3.564) (2.995,3.590) (2.940,3.614) (2.884,3.634) (2.826,3.653) (2.768,3.668) (2.710,3.681) (2.651,3.691) (2.591,3.699) (2.532,3.704) (2.472,3.707) (2.412,3.707) (2.352,3.705) (2.292,3.701) (2.232,3.695) (2.173,3.686) (2.114,3.675) (2.055,3.662) (1.997,3.647) (1.939,3.631) (1.882,3.612) (1.826,3.591) (1.770,3.569) (1.715,3.545) (1.661,3.520) (1.607,3.493) (1.555,3.465) (1.503,3.435) (1.451,3.403) (1.401,3.371) (1.368,3.348)}
\LC{(-2.404,-0.022) (-2.464,-0.023) (-2.524,-0.024) (-2.584,-0.026) (-2.644,-0.027) (-2.704,-0.029) (-2.764,-0.030) (-2.824,-0.032) (-2.883,-0.033) (-2.943,-0.035) (-3.003,-0.037) (-3.063,-0.038) (-3.123,-0.040) (-3.183,-0.042) (-3.243,-0.044) (-3.303,-0.046) (-3.363,-0.048) (-3.423,-0.050) (-3.483,-0.052) (-3.543,-0.054) (-3.603,-0.056) (-3.663,-0.058) (-3.723,-0.061) (-3.783,-0.063) (-3.843,-0.065) (-3.903,-0.068) (-3.963,-0.070) (-4.000,-0.072)}
\LC{(-3.229,-1.082) (-3.270,-1.125) (-3.311,-1.168) (-3.351,-1.213) (-3.390,-1.259) (-3.428,-1.305) (-3.465,-1.353) (-3.501,-1.401) (-3.535,-1.450) (-3.569,-1.500) (-3.601,-1.550) (-3.632,-1.602) (-3.661,-1.654) (-3.689,-1.707) (-3.716,-1.761) (-3.741,-1.815) (-3.765,-1.870) (-3.787,-1.926) (-3.808,-1.983) (-3.826,-2.040) (-3.843,-2.097) (-3.859,-2.155) (-3.872,-2.214) (-3.883,-2.273) (-3.893,-2.332) (-3.900,-2.391) (-3.905,-2.451) (-3.908,-2.511) (-3.909,-2.571) (-3.908,-2.631) (-3.904,-2.691) (-3.898,-2.751) (-3.890,-2.810) (-3.879,-2.869) (-3.866,-2.927) (-3.850,-2.985) (-3.831,-3.042) (-3.811,-3.099) (-3.787,-3.154) (-3.761,-3.208) (-3.733,-3.261) (-3.702,-3.312) (-3.669,-3.362) (-3.633,-3.411) (-3.595,-3.457) (-3.555,-3.502) (-3.513,-3.545) (-3.469,-3.585) (-3.423,-3.624) (-3.375,-3.660) (-3.326,-3.694) (-3.275,-3.725) (-3.222,-3.754) (-3.168,-3.781) (-3.113,-3.805) (-3.057,-3.826) (-3.000,-3.845) (-2.943,-3.862) (-2.884,-3.876) (-2.825,-3.887) (-2.766,-3.896) (-2.707,-3.903) (-2.647,-3.907) (-2.587,-3.909) (-2.527,-3.909) (-2.467,-3.906) (-2.407,-3.902) (-2.347,-3.895) (-2.288,-3.886) (-2.229,-3.875) (-2.170,-3.862) (-2.112,-3.848) (-2.055,-3.831) (-1.997,-3.813) (-1.941,-3.793) (-1.885,-3.771) (-1.830,-3.748) (-1.775,-3.723) (-1.721,-3.697) (-1.668,-3.669) (-1.615,-3.640) (-1.564,-3.609) (-1.513,-3.577) (-1.463,-3.544) (-1.414,-3.510) (-1.365,-3.474) (-1.318,-3.438) (-1.271,-3.400) (-1.225,-3.362) (-1.180,-3.322) (-1.136,-3.281) (-1.093,-3.240) (-1.078,-3.226)}
\LC{(1.209,4.000) (1.171,3.959) (1.133,3.913) (1.095,3.866) (1.059,3.818) (1.023,3.770) (0.987,3.722) (0.953,3.673) (0.919,3.623) (0.886,3.573) (0.853,3.523) (0.821,3.472) (0.790,3.421) (0.759,3.369) (0.729,3.317) (0.700,3.265) (0.671,3.212) (0.643,3.159) (0.616,3.106) (0.589,3.052) (0.563,2.998) (0.538,2.944) (0.513,2.889) (0.489,2.834) (0.466,2.779) (0.443,2.723) (0.421,2.668) (0.399,2.612) (0.378,2.556) (0.357,2.499)}
\LC{(2.398,-0.019) (2.458,-0.020) (2.518,-0.022) (2.578,-0.023) (2.638,-0.024) (2.698,-0.025) (2.758,-0.027) (2.818,-0.028) (2.878,-0.030) (2.937,-0.031) (2.997,-0.033) (3.057,-0.034) (3.117,-0.036) (3.177,-0.038) (3.237,-0.039) (3.297,-0.041) (3.357,-0.043) (3.417,-0.044) (3.477,-0.046) (3.537,-0.048) (3.597,-0.050) (3.657,-0.052) (3.717,-0.054) (3.777,-0.056) (3.837,-0.058) (3.897,-0.060) (3.957,-0.063) (4.000,-0.064)}
\LC{(-2.508,-1.142) (-2.535,-1.195) (-2.560,-1.250) (-2.583,-1.305) (-2.603,-1.362) (-2.621,-1.419) (-2.636,-1.477) (-2.648,-1.536) (-2.658,-1.595) (-2.664,-1.655) (-2.668,-1.715) (-2.668,-1.775) (-2.665,-1.835) (-2.658,-1.894) (-2.648,-1.953) (-2.634,-2.012) (-2.617,-2.069) (-2.596,-2.125) (-2.571,-2.180) (-2.543,-2.233) (-2.511,-2.284) (-2.476,-2.332) (-2.437,-2.378) (-2.395,-2.421) (-2.350,-2.461) (-2.303,-2.498) (-2.253,-2.531) (-2.201,-2.561) (-2.147,-2.587) (-2.091,-2.609) (-2.034,-2.628) (-1.976,-2.643) (-1.917,-2.655) (-1.858,-2.663) (-1.798,-2.667) (-1.738,-2.668) (-1.678,-2.666) (-1.618,-2.661) (-1.559,-2.652) (-1.500,-2.641) (-1.442,-2.627) (-1.384,-2.610) (-1.327,-2.591) (-1.271,-2.569) (-1.216,-2.545) (-1.162,-2.519) (-1.136,-2.505)}
\LC{(2.686,-1.108) (2.718,-1.158) (2.749,-1.209) (2.778,-1.262) (2.806,-1.315) (2.831,-1.370) (2.855,-1.425) (2.876,-1.481) (2.895,-1.538) (2.912,-1.595) (2.926,-1.654) (2.938,-1.712) (2.948,-1.772) (2.954,-1.831) (2.958,-1.891) (2.959,-1.951) (2.957,-2.011) (2.952,-2.071) (2.944,-2.130) (2.933,-2.189) (2.918,-2.247) (2.900,-2.304) (2.879,-2.361) (2.854,-2.415) (2.826,-2.468) (2.795,-2.520) (2.760,-2.569) (2.723,-2.616) (2.683,-2.660) (2.640,-2.702) (2.594,-2.741) (2.546,-2.777) (2.496,-2.810) (2.444,-2.839) (2.390,-2.866) (2.335,-2.889) (2.278,-2.909) (2.220,-2.925) (2.162,-2.938) (2.102,-2.948) (2.043,-2.955) (1.983,-2.959) (1.923,-2.959) (1.863,-2.957) (1.803,-2.952) (1.744,-2.944) (1.685,-2.933) (1.626,-2.920) (1.568,-2.904) (1.511,-2.886) (1.455,-2.866) (1.399,-2.844) (1.344,-2.819) (1.290,-2.793) (1.237,-2.765) (1.185,-2.735) (1.134,-2.703) (1.101,-2.681)}
\LC{(2.560,1.233) (2.583,1.288) (2.604,1.344) (2.623,1.401) (2.639,1.459) (2.653,1.517) (2.664,1.576) (2.671,1.636) (2.676,1.695) (2.678,1.755) (2.676,1.815) (2.671,1.875) (2.662,1.934) (2.650,1.993) (2.635,2.051) (2.615,2.108) (2.592,2.163) (2.566,2.217) (2.535,2.269) (2.502,2.318) (2.464,2.366) (2.424,2.410) (2.381,2.451) (2.334,2.490) (2.286,2.524) (2.235,2.556) (2.181,2.584) (2.127,2.608) (2.070,2.629) (2.013,2.646) (1.954,2.659) (1.895,2.669) (1.835,2.675) (1.775,2.678) (1.715,2.677) (1.655,2.673) (1.596,2.666) (1.537,2.657) (1.478,2.644) (1.420,2.629) (1.363,2.611) (1.306,2.591) (1.251,2.568) (1.233,2.560)}
\LC{(0.253,4.000) (0.244,3.947) (0.236,3.888) (0.227,3.828) (0.219,3.769) (0.211,3.710) (0.203,3.650) (0.195,3.591) (0.188,3.531) (0.181,3.471) (0.173,3.412) (0.166,3.352) (0.159,3.293) (0.153,3.233) (0.146,3.173) (0.140,3.114) (0.134,3.054) (0.127,2.994) (0.122,2.935) (0.116,2.875) (0.110,2.815) (0.105,2.755) (0.099,2.696) (0.097,2.666)}
\LC{(-2.303,0.784) (-2.344,0.829) (-2.383,0.875) (-2.420,0.921) (-2.456,0.969) (-2.491,1.019) (-2.523,1.069) (-2.554,1.120) (-2.583,1.173) (-2.611,1.226) (-2.636,1.281) (-2.659,1.336) (-2.679,1.392) (-2.698,1.449) (-2.714,1.507) (-2.727,1.566) (-2.737,1.625) (-2.745,1.684) (-2.749,1.744) (-2.751,1.804) (-2.749,1.864) (-2.744,1.924) (-2.736,1.983) (-2.724,2.042) (-2.709,2.100) (-2.690,2.157) (-2.668,2.213) (-2.642,2.267) (-2.612,2.319) (-2.579,2.369) (-2.543,2.417) (-2.504,2.462) (-2.461,2.505) (-2.416,2.544) (-2.368,2.580) (-2.318,2.613) (-2.266,2.642) (-2.211,2.668) (-2.156,2.691) (-2.099,2.709) (-2.041,2.725) (-1.982,2.736) (-1.922,2.745) (-1.863,2.749) (-1.803,2.751) (-1.743,2.749) (-1.683,2.745) (-1.623,2.737) (-1.564,2.726) (-1.506,2.713) (-1.448,2.697) (-1.391,2.679) (-1.335,2.658) (-1.279,2.635) (-1.225,2.610) (-1.171,2.583) (-1.119,2.554) (-1.068,2.522) (-1.017,2.490) (-0.968,2.455) (-0.920,2.419) (-0.873,2.382) (-0.828,2.343) (-0.798,2.316)}

    % Cuadrupolo plano (cuatro cargas) - puntas de flecha
% Generado por scripts/generar_datos.py (no editar a mano).
\FL{2.640}{0.432}{22.3}
\FL{1.947}{1.948}{135.1}
\FL{1.321}{1.314}{134.5}
\FL{1.004}{1.009}{135.5}
\FL{0.809}{0.805}{134.7}
\FL{0.653}{0.658}{135.2}
\FL{0.520}{0.523}{135.0}
\FL{0.384}{0.389}{134.8}
\FL{0.218}{0.213}{135.7}
\FL{0.218}{-0.213}{-135.7}
\FL{0.384}{-0.389}{-134.8}
\FL{0.520}{-0.523}{-135.0}
\FL{0.653}{-0.658}{-135.2}
\FL{0.809}{-0.805}{-134.7}
\FL{1.004}{-1.009}{-135.5}
\FL{1.321}{-1.314}{-134.5}
\FL{1.947}{-1.948}{-135.1}
\FL{2.640}{-0.432}{-22.3}
\FL{-0.218}{0.213}{44.3}
\FL{-0.384}{0.389}{45.2}
\FL{-0.520}{0.523}{45.0}
\FL{-0.653}{0.658}{44.8}
\FL{-0.809}{0.805}{45.3}
\FL{-1.004}{1.009}{44.5}
\FL{-1.321}{1.314}{45.5}
\FL{-1.947}{1.948}{44.9}
\FL{-2.640}{0.432}{157.7}
\FL{-2.640}{-0.432}{-157.7}
\FL{-1.947}{-1.948}{-44.9}
\FL{-1.321}{-1.314}{-45.5}
\FL{-1.004}{-1.009}{-44.5}
\FL{-0.809}{-0.805}{-45.3}
\FL{-0.653}{-0.658}{-44.8}
\FL{-0.520}{-0.523}{-45.0}
\FL{-0.384}{-0.389}{-45.2}
\FL{-0.218}{-0.213}{-44.3}
\FL{-3.230}{-1.839}{-102.2}
\FL{-1.395}{-3.064}{28.2}
\FL{3.659}{-2.283}{-85.5}
\FL{1.280}{-3.257}{143.2}
\FL{0.375}{-3.211}{105.7}
\FL{-3.818}{2.498}{90.3}
\FL{-1.248}{3.327}{-39.3}
\FL{3.175}{1.794}{77.2}
\FL{1.417}{3.039}{-153.4}
\FL{-2.943}{2.958}{44.4}
\FL{-0.689}{-3.289}{61.8}
\FL{-0.555}{3.364}{-67.9}
\FL{3.340}{3.354}{135.5}
\FL{-3.203}{-0.043}{-178.2}
\FL{-3.882}{-2.263}{-100.3}
\FL{-2.249}{-3.879}{10.8}
\FL{0.710}{3.283}{-119.1}
\FL{3.207}{-0.038}{-1.6}
\FL{-2.402}{-2.414}{-44.4}
\FL{2.662}{-2.681}{-135.9}
\FL{2.417}{2.417}{135.0}
\FL{0.164}{3.332}{-96.6}
\FL{-2.490}{2.477}{45.6}

  \end{scope}
  \node[cargapos] at (1,0) {$+$};
  \node[cargapos] at (-1,0) {$+$};
  \node[carganeg] at (0,1) {$-$};
  \node[carganeg] at (0,-1) {$-$};
  \colorlet{marcoc}{grispizarra!70}
  \marcoMMA{4}{2}{$x/a$}{$y/a$}
\end{tikzpicture}
\caption{Cuadrupolo plano del \cref{ej:cuadrupolo-plano}: $+q$ en $(\pm a,0,0)$ y $-q$ en $(0,\pm a,0)$, en el plano $z=0$. Líneas de campo (magenta, con su sentido), trazadas desde las cargas positivas a
intervalos angulares iguales y completadas con las líneas que llegan desde lejos; las diagonales $y=\pm x$ (a trazos) forman la equipotencial nula. El campo se «estira» según $x$ (donde están las cargas positivas, $Q_{xx}>0$) y se «comprime» según
$y$ ($Q_{yy}<0$).}
\label{fig:cuadrupolo-plano}
\end{figure}


% -------------------------------------------------------------------------
\subsection{Resumen: los tres primeros momentos}\label{sec:resumen}
% -------------------------------------------------------------------------
A distancias suficientemente grandes de una distribución dada podemos escribir el potencial como
\begin{equation}\label{eq:resumen}
  \boxed{\;\phi(\vc{r})=\phimul{0}+\phimul{1}+\phimul{2}+\cdots=\frac{1}{4\pi\eo}\Biggl[\frac{Q}{r}+\frac{\vp\cdot\ur}{r^{2}}+\frac{1}{2r^{3}}\sum_{j,k}\Qjk\,\cd{j}\cd{k}+\cdots\Biggr]\;}
\end{equation}
donde los momentos multipolares se calculan según la \cref{tab:momentos}, para distribuciones continuas o discretas.

\begin{table}[ht]
\centering
\caption{Momentos multipolares de una distribución de carga respecto del origen $O$. Para distribuciones superficiales o lineales se sustituye $\rho\,\dVp$ por
$\sigma\,\dSp$ o $\lambda\,\dlp$.}
\label{tab:momentos}
\renewcommand{\arraystretch}{1.7}
\small\setlength{\tabcolsep}{4pt}
\begin{tabular}{@{}lccc@{}}
  \toprule
  \textsc{Momento} & \textsc{Distribución continua} & \textsc{Distribución discreta} & \textsc{Término}\\
  \midrule
  Monopolar   & $Q=\displaystyle\int_{V'}\rho(\rp)\,\dVp$ & $Q=\displaystyle\sum_{i=1}^{N}q_i$ & $\dfrac{Q}{4\pi\eo r}$\\
  Dipolar     & $\vp=\displaystyle\int_{V'}\rp\rho(\rp)\,\dVp$ & $\vp=\displaystyle\sum_{i=1}^{N}q_i\,\rp_i$ & $\dfrac{\vp\cdot\ur}{4\pi\eo r^{2}}$\\
  Cuadrupolar & $\Qjk=\displaystyle\int_{V'}\rho\,\bigl(3x'_jx'_k-\delt{j}{k}r'^{2}\bigr)\dVp$ & $\Qjk=\displaystyle\sum_{i=1}^{N}q_i\bigl(3x'_{i,j}x'_{i,k}-\delt{j}{k}r_i'^{\,2}\bigr)$
                         & $\dfrac{\sum_{j,k}\Qjk\cd{j}\cd{k}}{8\pi\eo r^{3}}$\\
  \bottomrule
\end{tabular}
\end{table}

% =========================================================================
\section{Estimación del error al truncar el desarrollo}\label{sec:error}
% =========================================================================
La pregunta obligada ahora es: ¿cuántos términos del desarrollo debemos considerar? Todo depende del error con el que queramos determinar el
potencial. Pero conocer el error exige conocer el potencial exacto, que es precisamente lo que no sabemos calcular y la razón por la que hacemos el
desarrollo. Empezamos, por eso, con una situación en la que sí lo conocemos, y después veremos que, bajo una hipótesis muy natural, el error puede
acotarse sin conocer el potencial exacto.

\begin{definicion}[Error relativo de truncamiento]\label{def:error}
Sea $\phi^{(N)}=\sum_{\ell=0}^{N}\phimul{\ell}$ el desarrollo multipolar truncado en el orden $N$. Su \emph{error relativo} en el punto $\vc{r}$ es
\begin{equation}\label{eq:error}
  e_{N}(\vc{r})=\left\lvert\frac{\phi_{\mathrm{exacto}}(\vc{r})-\phi^{(N)}(\vc{r})}{\phi_{\mathrm{exacto}}(\vc{r})}\right\rvert\times100\ (\%).
\end{equation}
\end{definicion}

% -------------------------------------------------------------------------
\subsection{Un caso con solución exacta: la carga puntual}\label{sec:error-carga}
% -------------------------------------------------------------------------
\begin{ejemplo}[Carga puntual en el origen]\label{ej:carga-origen}
Si la carga $q$ está en el origen, $\phi_{\mathrm{exacto}}=q/(4\pi\eo r)$ coincide con el término monopolar: todos los demás momentos son nulos ($\vp=q\,\vc{0}=\vc{0}$,
$\Qjk=q\,(0-0)=0$, \dots) y, tomando un solo término, el error es del $0\%$.
\end{ejemplo}

\begin{ejemplo}[Carga puntual fuera del origen]\label{ej:carga-fuera}
Sea ahora la carga $q$ en $\rp=a\,\uz$ (\cref{fig:error-carga}). El potencial exacto es
\[
  \phi_{\mathrm{exacto}}(\vc{r})=\frac{1}{4\pi\eo}\,\frac{q}{R}=\frac{q}{4\pi\eo\sqrt{r^{2}+a^{2}-2ra\cos\theta}},
\]
y, por el \cref{teo:1R-legendre} con $r'=a$ y $\theta'=\theta$, su desarrollo multipolar es $\phi=\frac{q}{4\pi\eo}\sum_{\ell}\frac{a^{\ell}}{r^{\ell+1}}\Pl{\ell}(\cos\theta)$. Los
dos primeros términos son
\[
  \phimul{0}=\frac{q}{4\pi\eo r},\qquad \phimul{1}=\frac{\vp\cdot\ur}{4\pi\eo r^{2}}=\frac{q\,\rp\cdot\ur}{4\pi\eo r^{2}}=\frac{q\,a\cos\theta}{4\pi\eo r^{2}},
  \qquad \phi^{(1)}=\frac{q}{4\pi\eo}\Bigl(\frac1r+\frac{a\cos\theta}{r^{2}}\Bigr).
\]
La \cref{tab:error-carga} recoge el error cometido, con $q=1\,\mathrm{nC}$ y $a=1\,\mathrm{m}$, para $\theta=\pi/2$ y $\theta=0$.
\end{ejemplo}

\begin{table}[ht]
\centering
\caption{Potencial exacto y aproximado (términos monopolar y dipolar) de una carga $q=1\,\mathrm{nC}$ situada en $a\,\uz$ con $a=1\,\mathrm{m}$, y error relativo
$e_{1}$ (\cref{def:error}). Se ha usado $1/(4\pi\eo)=8{,}98755\cdot10^{9}\,\mathrm{N\,m^{2}/C^{2}}$.}
\label{tab:error-carga}
\renewcommand{\arraystretch}{1.25}
\small\setlength{\tabcolsep}{3.5pt}
\begin{tabular}{@{}c@{\quad}ccc@{\quad}ccc@{}}
  \toprule
  & \multicolumn{3}{c}{$\theta=\pi/2$} & \multicolumn{3}{c}{$\theta=0$}\\
  \cmidrule(r{0.8em}){2-4}\cmidrule{5-7}
  $r$ & $\phi_{\mathrm{exacto}}$ (V) & $\phi^{(1)}$ (V) & $e_{1}$ (\%) & $\phi_{\mathrm{exacto}}$ (V) & $\phi^{(1)}$ (V) & $e_{1}$ (\%)\\
  \midrule
  $10\,a$   & $0{,}894295$  & $0{,}898755$  & $0{,}499$              & $0{,}998617$  & $0{,}988631$  & $1{,}00$\\
  $100\,a$  & $0{,}0898710$ & $0{,}0898755$ & $5{,}00\cdot10^{-3}$   & $0{,}0907834$ & $0{,}0907743$ & $1{,}00\cdot10^{-2}$\\
  $1000\,a$ & $8{,}987547\cdot10^{-3}$ & $8{,}987552\cdot10^{-3}$ & $5{,}00\cdot10^{-5}$ & $8{,}996548\cdot10^{-3}$ & $8{,}996539\cdot10^{-3}$ & $1{,}00\cdot10^{-4}$\\
  \bottomrule
\end{tabular}
\end{table}

% =========================================================================
% Figura: error relativo e_1 (términos monopolar + dipolar) del potencial de
% una carga q en a u_z, en función de r/a, para theta = 0, pi/4, pi/2 (escala
% logarítmica). Datos exactos: figuras/datos/error_carga.dat.
% =========================================================================
\begin{figure}[ht]
\centering
\begin{minipage}[c]{0.3\textwidth}\centering
\begin{tikzpicture}[x=0.88cm, y=0.88cm, >=Stealth]
  \draw[eje] (0,-0.3) -- (0,2.6) node[above, etiq] {$z$};
  \draw[eje] (-0.3,0) -- (2.9,0) node[right, etiq] {$y$};
  \draw[eje] (0,0) -- (-0.8,-0.7) node[below left=-2pt, etiq] {$x$};
  \coordinate (O) at (0,0);
  \coordinate (Q) at (0,1.0);
  \coordinate (P) at (2.5,1.7);
  \draw[vect, rojoteo] (O) -- (Q) node[pos=0.5, left=1pt, etiq, text=rojoteo] {$\rp=a\,\uz$};
  \draw[vect] (O) -- (P) node[pos=0.55, below right=-2pt, etiq] {$\vc{r}$};
  \draw[->, line width=0.6pt, grispizarra, densely dashed] (Q) -- (P) node[pos=0.5, above=1pt, etiq] {$\vc{R}$};
  \pic[draw=magacero, line width=0.5pt, angle radius=0.55cm, "$\theta$", angle eccentricity=1.45, font=\small] {angle=P--O--Q};
  \node[cargapos] at (Q) {$+$};
  \fill[magprof] (P) circle (1.5pt) node[right=2pt, etiq] {$P$};
\end{tikzpicture}
\end{minipage}\hfill
\begin{minipage}[c]{0.68\textwidth}\centering
\begin{tikzpicture}
\begin{loglogaxis}[marcoMMA, width=8.2cm, height=5.6cm, xmin=1.5, xmax=1000, ymin=1e-5, ymax=60,
                   xlabel={$r/a$}, ylabel={$e_{1}$ (\%)}, every axis plot/.style={line width=1.1pt},
                   log basis y=10, ytick={1e-4,1e-2,1,10}, xtick={2,10,100,1000},
                   xticklabels={$2$,$10$,$100$,$1000$},
                   legend pos=north east, legend cell align=left]
  \addplot[mmaverde!85!black] table[x=x, y=e0]  {figuras/datos/error_carga.dat};
  \addplot[mmanaranja]        table[x=x, y=e45] {figuras/datos/error_carga.dat};
  \addplot[mmaazul]           table[x=x, y=e90] {figuras/datos/error_carga.dat};
  \legend{$\theta=0$, $\theta=\pi/4$, $\theta=\pi/2$}
\end{loglogaxis}
\end{tikzpicture}
\end{minipage}
\caption{Error relativo $e_{1}$ al aproximar el potencial de una carga puntual situada en $a\,\uz$ por sus términos monopolar y dipolar (\cref{ej:carga-fuera}). En escala
doble logarítmica las curvas son rectas de pendiente $-2$: el error decrece como $(a/r)^{2}$, con un coeficiente $\abs{\Pl{2}(\cos\theta)}$ que vale $1$ sobre el eje, $1/4$
para $\theta=\pi/4$ y $1/2$ en el plano ecuatorial. La curva $\theta=0$ coincide con la cota $100\,(a/r)^{2}$ del primer término despreciado, y la caída brusca de la curva $\theta=\pi/4$ cerca de $r=2a$ corresponde a un cambio de signo del error. Los valores para $r=10a,100a,1000a$ son los de la
\cref{tab:error-carga}.}
\label{fig:error-carga}
\end{figure}


\begin{observacion}[Lectura de los resultados]\label{obs:error-carga}
Al multiplicar la distancia por diez, el error se divide por cien: decrece como $(a/r)^{2}$, que es el peso relativo del primer término despreciado, el
cuadrupolar ($\phimul{2}/\phimul{0}=(a/r)^{2}\Pl{2}(\cos\theta)$). El coeficiente es $\abs{\Pl{2}(\cos\theta)}$: vale $1$ sobre el eje y $\frac12$ en el plano ecuatorial, de
ahí que el error sobre el eje sea el doble (\cref{fig:error-carga}). En la dirección del cono nodal, $\cos^{2}\theta=1/3$, el término cuadrupolar se anula y el error cae
un orden más, al $(a/r)^{3}$.
\end{observacion}

% -------------------------------------------------------------------------
\subsection{Cota del error para distribuciones de un solo signo}\label{sec:cota}
% -------------------------------------------------------------------------
Existe una situación en la que, aunque no conozcamos el potencial exacto, podemos estimar el error cometido: cuando toda la carga de la distribución es del
mismo signo. En ese caso el término monopolar no se anula y es el dominante, y medimos el peso de cada término respecto de él.

\begin{teorema}[Estimación del error]\label{teo:cota-error}
Sea $\rho$ una distribución de carga \emph{de un solo signo} (por ejemplo, $\rho\ge0$) con carga neta $Q\neq0$, contenida en la bola de radio $a$ centrada en el origen
($a$ es el \emph{tamaño característico} de la distribución). Entonces, para todo $r>a$ y todo $n\in\N$,
\begin{equation}\label{eq:cota-error}
  \boxed{\;\left\lvert\frac{\phimul{n}(\vc{r})}{\phimul{0}(\vc{r})}\right\rvert\le\Bigl(\frac{a}{r}\Bigr)^{n}\;}
\end{equation}
\end{teorema}
\begin{proof}
Por \cref{eq:desarrollo-multipolar} y \cref{eq:monopolar},
\[
  \frac{\phimul{n}(\vc{r})}{\phimul{0}(\vc{r})}=\frac{\dfrac{1}{4\pi\eo}\dfrac{1}{r^{n+1}}\displaystyle\int_{V'}r'^{\,n}\Pl{n}(\cosp)\,\rho(\rp)\,\dVp}{\dfrac{1}{4\pi\eo}\dfrac{Q}{r}}
  =\frac{1}{Q\,r^{n}}\int_{V'}r'^{\,n}\Pl{n}(\cosp)\,\rho(\rp)\,\dVp .
\]
Como $\rho\ge0$, $\abs{\rho}=\rho$; además $\abs{\Pl{n}(\cosp)}\le1$ (\cref{lem:legendre-cota}) y $r'^{\,n}\le a^{n}$ en el soporte de $\rho$. Por tanto,
\[
  \left\lvert\int_{V'}r'^{\,n}\Pl{n}(\cosp)\,\rho\,\dVp\right\rvert\le\int_{V'}r'^{\,n}\,\abs{\Pl{n}(\cosp)}\,\rho\,\dVp\le a^{n}\int_{V'}\rho\,\dVp=Q\,a^{n},
\]
y dividiendo por $Q\,r^{n}>0$ se obtiene \cref{eq:cota-error}. Si $\rho\le0$ se aplica lo anterior a $-\rho$.
\end{proof}

\begin{corolario}[Error al truncar en el orden $N$]\label{cor:resto}
En las condiciones del \cref{teo:cota-error}, el error relativo del desarrollo truncado (respecto del término dominante) satisface
\begin{equation}\label{eq:resto}
  \frac{\abs{\phi-\phi^{(N)}}}{\abs{\phimul{0}}}\le\sum_{n=N+1}^{\infty}\Bigl(\frac{a}{r}\Bigr)^{n}=\frac{(a/r)^{N+1}}{1-a/r}.
\end{equation}
\end{corolario}
\begin{proof}
$\abs{\phi-\phi^{(N)}}=\bigl\lvert\sum_{n>N}\phimul{n}\bigr\rvert\le\sum_{n>N}\abs{\phimul{n}}\le\abs{\phimul{0}}\sum_{n>N}(a/r)^{n}$, y la serie geométrica de razón $a/r<1$ se suma
directamente.
\end{proof}

\begin{observacion}[Uso práctico: cuántos términos tomar]\label{obs:cuantos}
El criterio del formulario~\cite{FORMULARIO} identifica el error con el peso del término $n$-ésimo, $\phimul{n}/\phimul{0}<(a/r)^{n}$. Si, por ejemplo, $r\sim10a$ y
queremos un error inferior al $1\%$,
\[
  \Bigl(\frac{a}{r}\Bigr)^{n}=\Bigl(\frac{a}{10a}\Bigr)^{n}=10^{-n}\le0{,}01\quad\Longrightarrow\quad n=2,
\]
y habrá que tomar las contribuciones $n=0,1,2$: monopolar, dipolar y cuadrupolar. El \cref{cor:resto} lo confirma con todo rigor: truncando en $N=2$, el error es menor
que $(0{,}1)^{3}/0{,}9\simeq0{,}11\%$, mientras que truncando en $N=1$ solo se garantiza un $1{,}1\%$. En la práctica se usa también la estimación con el tamaño
\emph{medio} de la distribución, $\langle r'\rangle$, en lugar del máximo $a$:
\begin{equation}\label{eq:estimacion-media}
  \frac{\phimul{n}}{\phimul{0}}\lesssim\frac{\langle r'\rangle^{n}}{r^{n}}\,\frac{1}{Q}\int_{V'}\rho\,\dVp=\Bigl(\frac{\langle r'\rangle}{r}\Bigr)^{n},
\end{equation}
que es menos rigurosa pero más ajustada: es la que usaremos en los seminarios.
\end{observacion}

\begin{observacion}[Alcance de la cota]\label{obs:alcance}
\begin{enumerate}[label=(\alph*)]
  \item \emph{La cota es óptima.} Para la carga puntual del \cref{ej:carga-fuera}, sobre el eje ($\theta=0$) se tiene $\phimul{n}/\phimul{0}=(a/r)^{n}\Pl{n}(1)=(a/r)^{n}$: la
        desigualdad \cref{eq:cota-error} se alcanza.
  \item \emph{La hipótesis de un solo signo es esencial.} Si hay cargas de ambos signos, $Q$ puede ser pequeña o nula (un dipolo, un cuadrupolo\dots) y entonces el término
        monopolar deja de ser el dominante: no tiene sentido comparar con él. En ese caso el error debe estimarse respecto del primer término no nulo, como hicimos al comparar el
        dipolo físico con el puntual (\cref{prop:discrepancia}).
\end{enumerate}
\end{observacion}

% =========================================================================
\section{Seminarios prácticos}\label{sec:seminarios}
% =========================================================================
Esta sección resuelve, con todo detalle, los cinco problemas de la relación del tema~\cite{PROBLEMAS-T2}. Cada uno sigue la misma estructura que en
los temas anteriores: el \emph{enunciado}, la \emph{idea} que guía la resolución, la \emph{resolución} paso a paso y un breve comentario sobre la
\emph{interpretación} del resultado. Todos se reducen a calcular unos pocos momentos multipolares, de modo que la estrategia es siempre la misma:
\begin{enumerate}[label=(\arabic*)]
  \item elegir el origen (si el enunciado no lo fija, el centro de simetría de la distribución) y los ejes adaptados a la geometría;
  \item calcular $Q$, $\vp$ y $\Qjk$ con la \cref{tab:momentos}, aprovechando las simetrías para anular integrales antes de hacerlas
        (\cref{teo:revolucion}: con simetría de revolución basta $Q_{zz}$ y se puede partir directamente de \cref{eq:cuadrupolar-revolucion});
  \item decidir con la \cref{sec:error} cuántos términos hacen falta y construir el potencial; el campo se obtiene, si se pide, como $\vE=-\nab\phi$.
\end{enumerate}
La \cref{tab:resultados}, al final de la sección, recoge todos los resultados.

% =========================================================================
\subsection{Problema 2.1. Barra cargada vista desde lejos}\label{sem:barra}
% =========================================================================
\begin{ejercicio}[Potencial y campo de una barra con error menor del $1\%$]\label{ejer:barra}
Calcular, con un error menor del $1\%$, el campo eléctrico y el potencial creados por una distribución lineal de carga de densidad uniforme $\lambda$ sobre una línea
recta finita de longitud $\ell$, para puntos de observación situados a una distancia mayor que $10$ veces la longitud de la barra.
\end{ejercicio}

\begin{resolucion}
\textbf{Idea.} Tenemos una barra de longitud $\ell$ y densidad uniforme $\lambda$, y la observamos desde $r>10\ell$. No hace falta resolver la integral exacta: basta
desarrollar el potencial y quedarse con los términos que garantizan un error inferior al $1\%$. Tomamos el origen en el centro de la barra y el eje $z$ a lo largo de
ella (\figpanel{fig:sem-barra}{a}); así el problema tiene simetría de revolución alrededor de $z$ y simetría de inversión $z\to-z$.

\textbf{Paso 1: cuántos términos.} Toda la carga es del mismo signo, de modo que podemos usar la estimación del \cref{teo:cota-error}, en su versión con el tamaño
característico (\cref{eq:estimacion-media}): $\phimul{n}/\phimul{0}\lesssim(\langle r'\rangle/r)^{n}$. Para esta distribución, la longitud característica es del orden de la
propia barra, $\langle r'\rangle\sim\ell$, y como $r>10\ell$,
\[
  \frac{\langle r'\rangle}{r}\lesssim\frac{\ell}{r}<\frac{1}{10}\quad\Longrightarrow\quad\frac{\phimul{n}}{\phimul{0}}\lesssim\Bigl(\frac{1}{10}\Bigr)^{n}=10^{-n}.
\]
El término dipolar pesaría a lo sumo un $10\%$ y el cuadrupolar un $1\%$; para garantizar un error menor del $1\%$ basta calcular hasta el término cuadrupolar,
$\phi\simeq\phimul{0}+\phimul{1}+\phimul{2}$ (y veremos que $\phimul{1}=0$ por simetría).

\textbf{Paso 2: el desarrollo para la barra.} El elemento de carga es $\dif q=\lambda\,\dif z'$, en $\rp=z'\,\uz$ con $z'\in[-\ell/2,\ell/2]$. Si $z'>0$, el ángulo entre $\rp$ y
$\vc{r}$ es $\theta'=\theta$ y $r'=z'$; si $z'<0$, $\theta'=\pi-\theta$ y $r'=-z'$. Por la paridad de los polinomios de Legendre, $\Pl{n}(-x)=(-1)^{n}\Pl{n}(x)$, en ambos casos
$r'^{\,n}\Pl{n}(\cosp)=z'^{\,n}\Pl{n}(\cos\theta)$, y el \cref{teo:multipolar} da
\[
  \phi(\vc{r})=\frac{1}{4\pi\eo}\int_{-\ell/2}^{\ell/2}\frac{\lambda\,\dif z'}{\norm{\vc{r}-\rp}}
  =\frac{\lambda}{4\pi\eo}\sum_{n=0}^{\infty}\frac{\Pl{n}(\cos\theta)}{r^{n+1}}\int_{-\ell/2}^{\ell/2}z'^{\,n}\,\dif z' .
\]
\emph{Término monopolar} ($n=0$): $\int_{-\ell/2}^{\ell/2}\dif z'=\ell$ y $\Pl{0}=1$, luego
\[
  \phimul{0}(\vc{r})=\frac{\lambda\ell}{4\pi\eo r}=\frac{Q}{4\pi\eo r},\qquad Q=\lambda\ell .
\]
\emph{Término dipolar} ($n=1$): $\int_{-\ell/2}^{\ell/2}z'\,\dif z'=\bigl[\tfrac{z'^{2}}{2}\bigr]_{-\ell/2}^{\ell/2}=\tfrac{(\ell/2)^{2}}{2}-\tfrac{(\ell/2)^{2}}{2}=0$. Es consecuencia de haber
colocado el origen en el centro de la barra: el momento dipolar,
\[
  \vp=\int\rp\,\dif q=\lambda\int_{-\ell/2}^{\ell/2}z'\,\uz\,\dif z'=\lambda\,\uz\Bigl[\frac{z'^{2}}{2}\Bigr]_{-\ell/2}^{\ell/2}=\vc{0},
\]
se anula porque por cada elemento en $+z'$ hay otro en $-z'$ con la misma carga, $(+z')\lambda\,\dif z'+(-z')\lambda\,\dif z'=0$. Dicho con simetrías: la de revolución obliga a
que $\vp\parallel\uz$ (\cref{teo:revolucion}), y la de inversión $z\to-z$ a que $p_z=-p_z$, es decir, $\vp=\vc{0}$.

\emph{Término cuadrupolar} ($n=2$): $\int_{-\ell/2}^{\ell/2}z'^{2}\,\dif z'=\bigl[\tfrac{z'^{3}}{3}\bigr]_{-\ell/2}^{\ell/2}=\tfrac{\ell^{3}}{24}+\tfrac{\ell^{3}}{24}=\tfrac{\ell^{3}}{12}$ y
$\Pl{2}(\cos\theta)=\frac12(3\cos^{2}\theta-1)$, luego
\[
  \phimul{2}(\vc{r})=\frac{\lambda}{4\pi\eo r^{3}}\,\frac12\bigl(3\cos^{2}\theta-1\bigr)\frac{\ell^{3}}{12}=\frac{\lambda\ell^{3}}{96\pi\eo r^{3}}\bigl(3\cos^{2}\theta-1\bigr).
\]
\emph{Comprobación con el tensor cuadrupolar.} Como hay simetría de revolución, basta $Q_{zz}=\int(3z'^{2}-r'^{2})\,\dif q=\int2z'^{2}\,\lambda\,\dif z'=2\lambda\frac{\ell^{3}}{12}=\frac{\lambda\ell^{3}}{6}$,
y \cref{eq:cuadrupolar-revolucion} da $\phimul{2}=\frac{\lambda\ell^{3}/6}{16\pi\eo r^{3}}(3\cos^{2}\theta-1)$, el mismo resultado.

\textbf{Paso 3: el potencial.} Con un error menor del $1\%$,
\[
  \boxed{\;\phi(\vc{r})\simeq\frac{\lambda\ell}{4\pi\eo r}+\frac{\lambda\ell^{3}\bigl(3\cos^{2}\theta-1\bigr)}{96\pi\eo r^{3}}\qquad(r>10\ell)\;}
\]

\textbf{Paso 4: el campo.} El potencial solo depende de $r$ y $\theta$, de modo que, con el gradiente en esféricas,
$\vE=-\nab\phi=-\pd{\phi}{r}\,\ur-\frac1r\pd{\phi}{\theta}\,\uth$ ($E_\varphi=0$). Como hemos aproximado el potencial, el campo resultante es también aproximado.
\emph{Componente radial:}
\[
  E_r=-\pd{\phi}{r}=-\Bigl[-\frac{\lambda\ell}{4\pi\eo r^{2}}+\frac{\lambda\ell^{3}}{96\pi\eo}\bigl(3\cos^{2}\theta-1\bigr)(-3)r^{-4}\Bigr]
  =\frac{\lambda\ell}{4\pi\eo r^{2}}+\frac{\lambda\ell^{3}\bigl(3\cos^{2}\theta-1\bigr)}{32\pi\eo r^{4}} .
\]
\emph{Componente angular:} como $\pd{}{\theta}\bigl(3\cos^{2}\theta-1\bigr)=3\cdot2\cos\theta(-\sen\theta)=-6\sen\theta\cos\theta$,
\[
  \pd{\phi}{\theta}=\frac{\lambda\ell^{3}}{96\pi\eo r^{3}}\bigl(-6\sen\theta\cos\theta\bigr)=-\frac{\lambda\ell^{3}\sen\theta\cos\theta}{16\pi\eo r^{3}},\qquad
  E_\theta=-\frac1r\pd{\phi}{\theta}=\frac{\lambda\ell^{3}\sen\theta\cos\theta}{16\pi\eo r^{4}} .
\]
En definitiva,
\[
  \boxed{\;\vE(\vc{r})\simeq\Bigl[\frac{\lambda\ell}{4\pi\eo r^{2}}+\frac{\lambda\ell^{3}\bigl(3\cos^{2}\theta-1\bigr)}{32\pi\eo r^{4}}\Bigr]\ur+\frac{\lambda\ell^{3}\sen\theta\cos\theta}{16\pi\eo r^{4}}\,\uth\;}
\]
\end{resolucion}

% =========================================================================
% Figura del problema 2.1: (a) geometría de la barra (apuntes manuscritos,
% pág. 23) y (b) error relativo de la aproximación multipolar frente al
% potencial exacto. Datos: figuras/datos/error_barra.dat.
% =========================================================================
\begin{figure}[ht]
\centering
\begin{minipage}[c]{0.36\textwidth}\centering
\begin{tikzpicture}[x=1cm, y=1cm, >=Stealth]
  \draw[eje] (0,-1.9) -- (0,2.6) node[above, etiq] {$z$};
  \draw[eje] (0,0) -- (2.9,-0.55) node[right, etiq] {$y$};
  \draw[eje] (0,0) -- (-1.1,-0.95) node[below left=-2pt, etiq] {$x$};
  % barra
  \fill[magacero!85] (-0.06,-1.3) rectangle (0.06,1.3);
  \draw[fino] (0.15,1.3) -- (0.75,1.3) (0.15,-1.3) -- (0.75,-1.3);
  \draw[<->, fino] (0.6,-1.3) -- node[pos=0.22, right=1pt, etiq, font=\footnotesize] {$\ell$} (0.6,1.3);
  \node[etiq, font=\footnotesize, anchor=east] at (-0.12,1.45) {$\ell/2$};
  \node[etiq, font=\footnotesize, anchor=east] at (-0.12,-1.45) {$-\ell/2$};
  % elemento de carga
  \coordinate (Q) at (0,0.95);
  \fill[white, draw=magprof, line width=0.5pt] (-0.1,0.85) rectangle (0.1,1.05);
  \node[etiq, font=\footnotesize, anchor=east] at (-0.12,0.95) {$\dif q$};
  % punto campo
  \coordinate (O) at (0,0);
  \coordinate (Zp) at (0,1);
  \coordinate (P) at (2.4,2.45);
  \draw[vect] (O) -- (P) node[pos=0.55, below right=-1pt, etiq] {$\vc{r}$};
  \draw[->, line width=0.6pt, grispizarra, densely dashed] (Q) -- (P) node[pos=0.55, above=1pt, etiq] {$\vc{R}$};
  \pic[draw=magacero, line width=0.5pt, angle radius=0.42cm, "$\theta$", angle eccentricity=1.45, font=\small] {angle=P--O--Zp};
  \fill[magprof] (P) circle (1.5pt) node[right=2pt, etiq] {$P(r,\theta,\varphi)$};
\end{tikzpicture}\\[2pt]{\small\scshape (a)}
\end{minipage}\hfill
\begin{minipage}[c]{0.62\textwidth}\centering
\begin{tikzpicture}
\begin{loglogaxis}[marcoMMA, width=8.6cm, height=5.6cm, xmin=1, xmax=100, ymin=1e-9, ymax=40,
                   xlabel={$r/\ell$}, ylabel={error relativo (\%)},
                   every axis plot/.style={line width=1.05pt},
                   ytick={1e-8,1e-6,1e-4,1e-2,1}, xtick={1,10,100}, xticklabels={$1$,$10$,$100$},
                   legend pos=south west, legend cell align=left, legend style={font=\tiny}]
  \addplot[mmanaranja]                table[x=x, y=e0eje] {figuras/datos/error_barra.dat};
  \addplot[mmanaranja, densely dashed] table[x=x, y=e0ecu] {figuras/datos/error_barra.dat};
  \addplot[mmaazul]                   table[x=x, y=e2eje] {figuras/datos/error_barra.dat};
  \addplot[mmaazul, densely dashed]    table[x=x, y=e2ecu] {figuras/datos/error_barra.dat};
  \draw[rojoteo, line width=0.6pt] (axis cs:1,1) -- (axis cs:100,1);
  \draw[grispizarra, line width=0.6pt, densely dotted] (axis cs:10,1e-9) -- (axis cs:10,40);
  \node[font=\tiny, text=rojoteo, anchor=south east] at (axis cs:95,1) {$1\%$};
  \node[font=\tiny, text=grispizarra, anchor=south west, rotate=90, xshift=2pt] at (axis cs:10,3e-9) {$r=10\ell$};
  \legend{$\phimul{0}$; eje, $\phimul{0}$; ecuador, $\phimul{0}+\phimul{2}$; eje, $\phimul{0}+\phimul{2}$; ecuador}
\end{loglogaxis}
\end{tikzpicture}\\[2pt]{\small\scshape (b)}
\end{minipage}
\caption{Problema~2.1. (a)~Barra de longitud $\ell$ y densidad $\lambda$ sobre el eje $z$, centrada en el origen; el elemento de carga $\dif q=\lambda\,\dif z'$, en $z'$, y el punto campo $P$.
(b)~Error relativo, frente al potencial exacto, de la aproximación monopolar (naranja) y monopolar más cuadrupolar (azul), sobre el eje ($\theta=0$, continua) y en el plano
ecuatorial ($\theta=\pi/2$, a trazos). Para $r>10\ell$ (punteada) ambas aproximaciones quedan por debajo del $1\%$; la que incluye el término cuadrupolar mejora el resultado en
unos tres órdenes de magnitud.}
\label{fig:sem-barra}
\end{figure}


\begin{observacion}[Interpretación y un análisis más fino]\label{obs:sem-barra}
Lejos de la barra el término dominante es el de una carga puntual $Q=\lambda\ell$; la corrección cuadrupolar es positiva cerca del eje (la carga está «alargada» según $z$,
$Q_{zz}>0$) y negativa cerca del plano ecuatorial. La estimación $\langle r'\rangle\sim\ell$ es conservadora: la carga más alejada del centro está a $a=\ell/2$, de modo que
el \cref{teo:cota-error} garantiza $\abs{\phimul{2}/\phimul{0}}\le(1/20)^{2}=0{,}25\%$ y, de hecho, $\abs{\phimul{2}/\phimul{0}}=\frac{\ell^{2}}{24r^{2}}\abs{3\cos^{2}\theta-1}\le\frac{\ell^{2}}{12r^{2}}<0{,}084\%$.
Así pues, para $r>10\ell$ ya el término monopolar \emph{solo} cumple el requisito del $1\%$, y al añadir el cuadrupolar el error cae hasta el orden del término
siguiente no nulo, $\phimul{4}/\phimul{0}=\frac{\ell^{4}}{80r^{4}}\Pl{4}(\cos\theta)\sim10^{-6}$. La \figpanel{fig:sem-barra}{b} lo confirma comparando con el potencial exacto de la
barra, $\phi=\frac{\lambda}{4\pi\eo}\bigl[\operatorname{arcsenh}\frac{z+\ell/2}{s}-\operatorname{arcsenh}\frac{z-\ell/2}{s}\bigr]$ (con $s$ la distancia al eje).
\end{observacion}

% =========================================================================
\subsection{Problema 2.2. Momentos de una carga puntual}\label{sem:carga}
% =========================================================================
\begin{ejercicio}[Momentos de una carga puntual fuera del origen]\label{ejer:carga}
Sea una carga puntual $q$ situada en el punto $(a,b,c)$. Calcular el momento monopolar, dipolar y cuadrupolar para ese sistema. ¿Cuál de esas cantidades varía si se añade una
carga $-q$ en el origen de coordenadas?
\end{ejercicio}

\begin{resolucion}
\textbf{Idea.} Es un sistema discreto de una sola carga: los momentos se obtienen directamente de las expresiones discretas de la \cref{tab:momentos}, con
$\rp_1=a\,\ux+b\,\uy+c\,\uz$ y, por tanto, $r_1'^{\,2}=a^{2}+b^{2}+c^{2}$ (\figpanel{fig:sem-carga}{a}).

\textbf{1) Momentos del sistema inicial.} \emph{Monopolar} ($\ell=0$): la carga neta del sistema, $Q=\sum_iq_i=q$.

\emph{Dipolar} ($\ell=1$): el primer momento espacial de la distribución,
\[
  \vp=\sum_iq_i\,\rp_i=q\bigl(a\,\ux+b\,\uy+c\,\uz\bigr).
\]
\emph{Cuadrupolar} ($\ell=2$): $\Qjk=\sum_iq_i\bigl(3x'_{i,j}x'_{i,k}-\delt{j}{k}r_i'^{\,2}\bigr)$. En la diagonal principal ($j=k$, $\delt{j}{k}=1$),
\[
  Q_{xx}=q\bigl[3a^{2}-(a^{2}+b^{2}+c^{2})\bigr]=q\bigl(2a^{2}-b^{2}-c^{2}\bigr),
\]
\[
  Q_{yy}=q\bigl[3b^{2}-(a^{2}+b^{2}+c^{2})\bigr]=q\bigl(2b^{2}-a^{2}-c^{2}\bigr),\qquad
  Q_{zz}=q\bigl(2c^{2}-a^{2}-b^{2}\bigr).
\]
Fuera de la diagonal ($j\neq k$, $\delt{j}{k}=0$), y usando que el tensor es simétrico,
\[
  Q_{xy}=Q_{yx}=3qab,\qquad Q_{xz}=Q_{zx}=3qac,\qquad Q_{yz}=Q_{zy}=3qbc .
\]
En forma matricial,
\[
  \Qt=q\begin{pmatrix}2a^{2}-b^{2}-c^{2}&3ab&3ac\\3ab&2b^{2}-a^{2}-c^{2}&3bc\\3ac&3bc&2c^{2}-a^{2}-b^{2}\end{pmatrix},
\]
con $\Tr\Qt=q\,(2-1-1)(a^{2}+b^{2}+c^{2})=0$, como debe ser.

\textbf{2) Variaciones al añadir $-q$ en el origen} (\figpanel{fig:sem-carga}{b}). La nueva carga $q_2=-q$ tiene vector de posición nulo, $\rp_2=0\,\ux+0\,\uy+0\,\uz$.

\emph{Momento monopolar:} $Q_{\mathrm{nuevo}}=q_1+q_2=q+(-q)=0$. \emph{Es la única cantidad que varía} (y se anula).

\emph{Momentos dipolar y cuadrupolar:} sus contribuciones son proporcionales a las coordenadas de cada carga ($x'_j$) o a productos de ellas ($x'_jx'_k$, $r'^{2}$), y todas
se anulan para la carga situada en el origen:
\[
  \vp_{q_2}=(-q)\,\vc{0}=\vc{0},\qquad \bigl(\Qjk\bigr)_{q_2}=-q\bigl(3\cdot0\cdot0-\delt{j}{k}\cdot0\bigr)=0 .
\]
Por tanto, la carga añadida no perturba en absoluto los momentos de orden superior calculados antes.
\end{resolucion}

% =========================================================================
% Figura del problema 2.2: (a) carga q en (a,b,c); (b) tras añadir -q en el
% origen, el sistema es un dipolo de momento p = q r'_1. Perspectiva tdplot.
% =========================================================================
\begin{figure}[ht]
\centering
\newcommand{\ejesIII}{%
  \draw[eje] (0,0,0) -- (2.6,0,0) node[anchor=north east, etiq] {$x$};
  \draw[eje] (0,0,0) -- (0,2.7,0) node[anchor=west, etiq] {$y$};
  \draw[eje] (0,0,0) -- (0,0,2.3) node[anchor=south, etiq] {$z$};}
\newcommand{\proyecciones}{%
  \draw[punteado] (1.6,1.9,0) -- (1.6,1.9,1.5);
  \draw[punteado] (1.6,0,0) -- (1.6,1.9,0) -- (0,1.9,0);
  \node[etiq, font=\footnotesize, anchor=north] at (1.6,0,0) {$a$};
  \node[etiq, font=\footnotesize, anchor=north west] at (0,1.9,0) {$b$};
  \draw[punteado] (0,0,1.5) -- (1.6,1.9,1.5);
  \node[etiq, font=\footnotesize, anchor=east] at (0,0,1.5) {$c$};}
\tdplotsetmaincoords{70}{120}
\begin{minipage}[b]{0.47\textwidth}\centering
\begin{tikzpicture}[tdplot_main_coords, scale=1.25, >=Stealth]
  \ejesIII \proyecciones
  \draw[vect, rojoteo] (0,0,0) -- (1.6,1.9,1.5) node[pos=0.55, anchor=south east, etiq, text=rojoteo] {$\rp_1$};
  \node[cargapos] at (1.6,1.9,1.5) {$+$};
  \node[etiq, anchor=south west] at (1.6,1.9,1.62) {$q\ (a,b,c)$};
\end{tikzpicture}\\[2pt]{\small\scshape (a)}
\end{minipage}\hfill
\begin{minipage}[b]{0.47\textwidth}\centering
\begin{tikzpicture}[tdplot_main_coords, scale=1.25, >=Stealth]
  \ejesIII \proyecciones
  \draw[->, line width=1.5pt, magacero] (0,0,0) -- (1.6,1.9,1.5) node[pos=0.45, anchor=south east, xshift=-3pt, yshift=4pt, etiq, text=magacero] {$\vp=q\,\rp_1$};
  \node[cargapos] at (1.6,1.9,1.5) {$+$};
  \node[carganeg] at (0,0,0) {$-$};
  \node[etiq, anchor=south west] at (1.6,1.9,1.62) {$q$};
  \node[etiq, anchor=north west] at (0,0.1,-0.15) {$-q$};
\end{tikzpicture}\\[2pt]{\small\scshape (b)}
\end{minipage}
\caption{Problema~2.2. (a)~Carga puntual $q$ en $\rp_1=(a,b,c)$: $Q=q$, $\vp=q\,\rp_1$ y $\Qjk=q\,(3x'_{1,j}x'_{1,k}-\delt{j}{k}r_1'^{\,2})$. (b)~Al añadir $-q$ en el origen, el sistema se
convierte en un dipolo físico: solo cambia la carga neta ($Q=0$), mientras que $\vp$ y $\Qt$ no varían porque la nueva carga está en $\rp_2=\vc{0}$.}
\label{fig:sem-carga}
\end{figure}


\begin{observacion}[Interpretación]\label{obs:sem-carga}
Tras añadir la carga, el sistema es un dipolo físico de momento $\vp=q\,\rp_1$ y, al ser nula la carga neta, entra en juego el resultado de independencia del origen
(\cref{cor:centro-carga}): $\vp$ ya no cambia al trasladar los ejes. Antes de añadirla, en cambio, $\vp$ dependía del origen, y tomando el origen en la propia carga, que es su
centro de carga $\vc{r}_c=\vp/Q=(a,b,c)$, se anulaban $\vp$ y $\Qt$: una carga puntual vista desde su posición es un monopolo puro (\cref{ej:carga-origen}).
\end{observacion}

% =========================================================================
\subsection{Problema 2.3. Paralelepípedo uniformemente cargado}\label{sem:paralelepipedo}
% =========================================================================
\begin{ejercicio}[Potencial de un paralelepípedo con error menor del $1\%$]\label{ejer:paralelepipedo}
Sea una distribución de carga $Q$, uniforme, en un paralelepípedo de lados $2a$, $2b$, $2c$. Siendo todos los lados del mismo orden de magnitud $d$, calcular el potencial con un
error $<1\%$ cuando se observa ese recinto cargado a distancias mayores que $10d$.
\end{ejercicio}

\begin{resolucion}
\textbf{Idea.} El enunciado está incompleto: no indica respecto de qué origen, y como la carga neta no es nula el momento dipolar depende de esa elección
(\cref{prop:p-origen}). Lo natural es tomar el origen en el centro del paralelepípedo y los ejes paralelos a sus aristas (\cref{fig:sem-paralelepipedo}), de modo que ocupa la región
$[-a,a]\times[-b,b]\times[-c,c]$ con densidad $\rho=Q/\bigl[(2a)(2b)(2c)\bigr]=Q/(8abc)$.

\textbf{Paso 1: cuántos términos.} La carga es de un solo signo y su tamaño característico es $\langle r'\rangle\sim d$; con $r>10d$, $\phimul{n}/\phimul{0}\lesssim10^{-n}$, y para un
error menor del $1\%$ basta llegar al término cuadrupolar:
\[
  \phi\simeq\phimul{0}+\phimul{1}+\phimul{2}=\frac{Q}{4\pi\eo r}+\frac{\vp\cdot\ur}{4\pi\eo r^{2}}+\frac{1}{8\pi\eo r^{3}}\sum_{j,k}\Qjk\,\cd{j}\cd{k}.
\]

\textbf{Paso 2: momento dipolar.} Con $\rp=x'\,\ux+y'\,\uy+z'\,\uz$,
\[
  \vp=\rho\int_{-c}^{c}\!\!\int_{-b}^{b}\!\!\int_{-a}^{a}\bigl(x'\,\ux+y'\,\uy+z'\,\uz\bigr)\,\dif x'\,\dif y'\,\dif z'=\vc{0},
\]
porque cada componente es la integral de una función impar en un intervalo simétrico (por ejemplo, $\int_{-a}^{a}x'\,\dif x'=0$).

\textbf{Paso 3: tensor cuadrupolar.} En la diagonal, $3x'^{2}-r'^{2}=2x'^{2}-y'^{2}-z'^{2}$, y como las integrales se separan,
\begin{align*}
  Q_{xx}&=\rho\int_{-c}^{c}\!\!\int_{-b}^{b}\!\!\int_{-a}^{a}\bigl(2x'^{2}-y'^{2}-z'^{2}\bigr)\dif x'\,\dif y'\,\dif z'\\
         &=\rho\Bigl[2\cdot2b\cdot2c\Bigl[\frac{x'^{3}}{3}\Bigr]_{-a}^{a}-2a\cdot2c\Bigl[\frac{y'^{3}}{3}\Bigr]_{-b}^{b}-2a\cdot2b\Bigl[\frac{z'^{3}}{3}\Bigr]_{-c}^{c}\Bigr]\\
        &=\rho\Bigl[2\cdot4bc\,\frac{2a^{3}}{3}-4ac\,\frac{2b^{3}}{3}-4ab\,\frac{2c^{3}}{3}\Bigr]\\
        &=\underbrace{8\rho\,abc}_{Q}\,\frac{2a^{2}-b^{2}-c^{2}}{3}=\frac{Q}{3}\bigl(2a^{2}-b^{2}-c^{2}\bigr),
\end{align*}
y, de forma análoga (permutando los papeles de $a$, $b$, $c$),
\[
  Q_{yy}=\frac{Q}{3}\bigl(2b^{2}-a^{2}-c^{2}\bigr),\qquad Q_{zz}=\frac{Q}{3}\bigl(2c^{2}-a^{2}-b^{2}\bigr).
\]
Fuera de la diagonal, $Q_{xy}=3\rho\iiint x'y'\,\dif x'\,\dif y'\,\dif z'=0$ (integrando impar en $x'$), e igual $Q_{xz}=Q_{yz}=0$: los ejes de las aristas son los ejes principales.
Comprobación: $\Tr\Qt=\frac{Q}{3}\bigl[(2-1-1)a^{2}+(2-1-1)b^{2}+(2-1-1)c^{2}\bigr]=0$.

\textbf{Paso 4: el potencial.} Sustituyendo,
\[
  \phi(\vc{r})\simeq\frac{Q}{4\pi\eo r}+\frac{Q}{24\pi\eo r^{3}}\Bigl[\bigl(2a^{2}-b^{2}-c^{2}\bigr)\cd{x}^{2}+\bigl(2b^{2}-a^{2}-c^{2}\bigr)\cd{y}^{2}+\bigl(2c^{2}-a^{2}-b^{2}\bigr)\cd{z}^{2}\Bigr],
\]
con $\cd{x}^{2}=\sen^{2}\theta\cos^{2}\varphi=x^{2}/r^{2}$, $\cd{y}^{2}=\sen^{2}\theta\sen^{2}\varphi=y^{2}/r^{2}$ y $\cd{z}^{2}=\cos^{2}\theta=z^{2}/r^{2}$. En cartesianas se aprecia mejor la simetría:
\[
  \boxed{\;\phi(\vc{r})\simeq\frac{Q}{4\pi\eo r}+\frac{Q}{24\pi\eo r^{5}}\Bigl[x^{2}\bigl(2a^{2}-b^{2}-c^{2}\bigr)+y^{2}\bigl(2b^{2}-a^{2}-c^{2}\bigr)+z^{2}\bigl(2c^{2}-a^{2}-b^{2}\bigr)\Bigr]\;}
\]
\end{resolucion}

% =========================================================================
% Figura del problema 2.3 (apuntes manuscritos, pág. 19): paralelepípedo
% [-a,a]x[-b,b]x[-c,c] uniformemente cargado, con el origen en su centro.
% =========================================================================
\begin{figure}[ht]
\centering
\tdplotsetmaincoords{68}{122}
\begin{tikzpicture}[tdplot_main_coords, scale=1.05, >=Stealth]
  \def\A{1.5}\def\B{2.1}\def\C{1.0}
  % caras ocultas (aristas a trazos)
  \draw[magacero!70, densely dashed, line width=0.5pt] (-\A,-\B,-\C) -- (\A,-\B,-\C) (-\A,-\B,-\C) -- (-\A,\B,-\C) (-\A,-\B,-\C) -- (-\A,-\B,\C);
  % caras visibles
  \fill[rosaclaro, opacity=0.55] (\A,-\B,-\C) -- (\A,\B,-\C) -- (\A,\B,\C) -- (\A,-\B,\C) -- cycle;
  \fill[rosaclaro!70, opacity=0.55] (-\A,\B,-\C) -- (\A,\B,-\C) -- (\A,\B,\C) -- (-\A,\B,\C) -- cycle;
  \fill[rosaclaro!45, opacity=0.6] (-\A,-\B,\C) -- (\A,-\B,\C) -- (\A,\B,\C) -- (-\A,\B,\C) -- cycle;
  \draw[magacero, line width=0.7pt] (\A,-\B,-\C) -- (\A,\B,-\C) -- (-\A,\B,-\C) -- (-\A,\B,\C) -- (-\A,-\B,\C) -- (\A,-\B,\C) -- cycle;
  \draw[magacero, line width=0.7pt] (\A,\B,-\C) -- (\A,\B,\C) -- (-\A,\B,\C) (\A,\B,\C) -- (\A,-\B,\C);
  % ejes por el centro
  \draw[eje] (0,0,0) -- (2.9,0,0) node[anchor=north east, etiq] {$x$};
  \draw[eje] (0,0,0) -- (0,3.3,0) node[anchor=west, etiq] {$y$};
  \draw[eje] (0,0,0) -- (0,0,2.2) node[anchor=south, etiq] {$z$};
  % elemento de volumen
  \coordinate (Q) at (1.0,0.6,0.3);
  \draw[vect] (0,0,0) -- (Q) node[pos=0.55, anchor=north west, etiq] {$\rp$};
  \fill[white, draw=magprof, line width=0.5pt] (Q) circle (2pt);
  \node[etiq, font=\footnotesize, anchor=south west] at ($(Q)+(0,0.08,0.05)$) {$\dVp$};
  % cotas
  \draw[<->, fino] (\A+0.25,-\B,-\C) -- node[anchor=north, etiq, font=\footnotesize] {$2b$} (\A+0.25,\B,-\C);
  \draw[<->, fino] (\A,\B+0.25,-\C) -- node[anchor=north west, etiq, font=\footnotesize] {$2a$} (-\A,\B+0.25,-\C);
  \draw[<->, fino] (\A,\B+0.25,-\C) -- node[anchor=west, etiq, font=\footnotesize] {$2c$} (\A,\B+0.25,\C);
  \node[etiq, text=magprof, anchor=south east] at (-\A,-\B,\C) {$\rho=\dfrac{Q}{8abc}$};
\end{tikzpicture}
\caption{Problema~2.3. Paralelepípedo de lados $2a$, $2b$, $2c$ con carga $Q$ uniforme, centrado en el origen y con las aristas paralelas a los ejes. La simetría de inversión anula
el momento dipolar y todas las componentes no diagonales del tensor cuadrupolar.}
\label{fig:sem-paralelepipedo}
\end{figure}


\begin{observacion}[Interpretación: el cubo y la paridad]\label{obs:sem-paralelepipedo}
\begin{enumerate}[label=(\alph*)]
  \item La corrección cuadrupolar es positiva en la dirección de la arista más larga (si $a>b,c$, $Q_{xx}>0$): la caja se comporta como una carga «alargada» según ese eje.
  \item Para un \emph{cubo}, $a=b=c$, todas las componentes se anulan, $\Qt=0$: la simetría cúbica es suficiente para que el potencial no tenga corrección cuadrupolar, aunque la
        distribución no sea esférica.
  \item La distribución es simétrica frente a la inversión $\rp\to-\rp$, de modo que \emph{todos} los momentos de orden impar se anulan (sus integrandos, $r'^{\,\ell}\Pl{\ell}(\cosp)$
        con $\ell$ impar, cambian de signo). El primer término despreciado no es, pues, el octupolar sino el hexadecapolar ($\ell=4$), y el error real es del orden de $(d/r)^{4}\sim10^{-4}$,
        muy por debajo del $1\%$ exigido.
\end{enumerate}
\end{observacion}

% =========================================================================
\subsection{Problema 2.4. Esfera con densidad superficial \texorpdfstring{$\sigma_0\cos\theta$}{σ0 cos θ}}\label{sem:esfera-sigma}
% =========================================================================
\begin{ejercicio}[Momento dipolar de una esfera polarizada]\label{ejer:esfera-sigma}
Calcular el momento dipolar de una distribución superficial de carga sobre una esfera de radio $a$ cuando la densidad varía con el ángulo $\theta$ de la forma
$\sigma(\theta)=\sigma_0\cos\theta$.
\end{ejercicio}

\begin{resolucion}
\textbf{Idea.} La carga es positiva en el hemisferio norte y negativa en el sur (\figpanel{fig:sem-esfera-sigma}{a}): esperamos un momento dipolar dirigido según $+\uz$. Antes de
calcularlo comprobamos que la carga neta es nula, porque entonces el resultado no depende del origen y podemos tomar como referencia el centro de la esfera:
con $\dSp=a^{2}\sen\theta'\,\dif\theta'\,\dif\varphi'$,
\[
  Q=\oiint_{S'}\sigma\,\dSp=\sigma_0a^{2}\int_{0}^{2\pi}\dif\varphi'\int_{0}^{\pi}\cos\theta'\sen\theta'\,\dif\theta'=\sigma_0a^{2}\cdot2\pi\cdot\Bigl[\frac{\sen^{2}\theta'}{2}\Bigr]_{0}^{\pi}=0 .
\]

\textbf{Paso 1: el momento dipolar de una distribución superficial.} El momento dipolar de una distribución volúmica es $\vp=\int_{V'}\rp\rho(\rp)\,\dVp$. Al pasar a una distribución
estrictamente bidimensional sobre la superficie $S'$, sustituimos la densidad volúmica por la superficial y el diferencial de volumen por el de área:
\[
  \vp=\oiint_{S'}\rp\,\sigma(\theta')\,\dSp .
\]
\emph{Justificación mediante la delta de Dirac.} Una distribución superficial no es más que el caso límite de una volúmica en la que toda la carga está confinada en una capa de espesor
nulo: en esféricas, $\rho(r',\theta',\varphi')=\sigma(\theta',\varphi')\,\delta(r'-a)$, dimensionalmente correcta porque $[\mathrm{C/m^{3}}]=[\mathrm{C/m^{2}}][\mathrm{m^{-1}}]$.
Sustituyendo, con $\rp=r'\,\ur$,
\[
  \vp=\int_{0}^{2\pi}\!\!\int_{0}^{\pi}\!\!\int_{0}^{\infty}\bigl(r'\,\ur\bigr)\,\sigma(\theta',\varphi')\,\delta(r'-a)\,r'^{2}\sen\theta'\,\dif r'\,\dif\theta'\,\dif\varphi' .
\]
Reagrupando los factores radiales, la integral en $r'$ es $\int_{0}^{\infty}r'^{3}\delta(r'-a)\,\dif r'=a^{3}$, de modo que
$\vp=\int_{0}^{2\pi}\!\int_{0}^{\pi}(a\,\ur)\,\sigma\,\bigl(a^{2}\sen\theta'\,\dif\theta'\,\dif\varphi'\bigr)=\oiint_{S'}\rp\big|_{r'=a}\,\sigma\,\dSp$, como queríamos.

\textbf{Paso 2: parametrización.} Esta integral es \emph{vectorial} (no es un flujo $\iint\vc{F}\cdot\un\,\dsc$): cada componente cartesiana es una integral escalar. Usamos
\[
  \rp=a\,\ur=a\bigl(\sen\theta'\cos\varphi'\,\ux+\sen\theta'\sen\varphi'\,\uy+\cos\theta'\,\uz\bigr),\qquad \dSp=a^{2}\sen\theta'\,\dif\theta'\,\dif\varphi' .
\]

\textbf{Paso 3: componentes transversales.} Por el teorema de Fubini (producto de integrales simples iteradas),
\begin{align*}
  p_x&=\int_{0}^{2\pi}\!\!\int_{0}^{\pi}\bigl(a\sen\theta'\cos\varphi'\bigr)\bigl(\sigma_0\cos\theta'\bigr)\bigl(a^{2}\sen\theta'\bigr)\dif\theta'\,\dif\varphi'\\
  &=a^{3}\sigma_0\int_{0}^{\pi}\sen^{2}\theta'\cos\theta'\,\dif\theta'\underbrace{\int_{0}^{2\pi}\cos\varphi'\,\dif\varphi'}_{[\sen\varphi']_{0}^{2\pi}=0}=0,
\end{align*}
\[
  p_y=a^{3}\sigma_0\int_{0}^{\pi}\sen^{2}\theta'\cos\theta'\,\dif\theta'\underbrace{\int_{0}^{2\pi}\sen\varphi'\,\dif\varphi'}_{[-\cos\varphi']_{0}^{2\pi}=0}=0 .
\]
Las componentes transversales al eje de revolución se anulan idénticamente porque $\sigma=\sigma_0\cos\theta'$ tiene simetría azimutal perfecta (no depende de $\varphi'$),
como asegura el \cref{teo:revolucion}.

\textbf{Paso 4: componente axial.} Separando las variables angulares,
\[
  \vp=p_z\,\uz=a^{3}\sigma_0\,\uz\int_{0}^{2\pi}\dif\varphi'\int_{0}^{\pi}\cos^{2}\theta'\sen\theta'\,\dif\theta' .
\]
Con el cambio $u=\cos\theta'$, $\dif u=-\sen\theta'\,\dif\theta'$, que lleva $\theta'=0\mapsto u=1$ y $\theta'=\pi\mapsto u=-1$,
\[
  \int_{0}^{\pi}\cos^{2}\theta'\sen\theta'\,\dif\theta'=-\int_{1}^{-1}u^{2}\,\dif u=\int_{-1}^{1}u^{2}\,\dif u=\Bigl[\frac{u^{3}}{3}\Bigr]_{-1}^{1}=\frac13-\Bigl(-\frac13\Bigr)=\frac23 .
\]
Por tanto,
\[
  \boxed{\;\vp=a^{3}\sigma_0\,(2\pi)\,\frac23\,\uz=\frac43\pi a^{3}\sigma_0\,\uz\;}
\]
\end{resolucion}

% =========================================================================
% Figura del problema 2.4 (apuntes manuscritos, pág. 19): (a) esfera con
% sigma = sigma_0 cos(theta) (signos con densidad proporcional a |cos t|) y
% su momento dipolar; (b) líneas de campo exactas (analíticas): dipolares,
% r = C sen^2(t), fuera, y uniformes (verticales, hacia -z) dentro, con las
% equipotenciales a trazos (r^2 = a^2 cos t / cos t1 fuera, z = cte dentro).
% =========================================================================
\begin{figure}[ht]
\centering
\begin{minipage}[b]{0.45\textwidth}\centering
\begin{tikzpicture}[x=1cm, y=1cm, >=Stealth]
  \def\ra{1.55}
  \draw[eje] (0,-2.1) -- (0,2.4) node[above, etiq] {$z$};
  \draw[eje] (-2.1,0) -- (2.4,0) node[right, etiq] {$x$};
  \shade[inner color=white, outer color=rosaclaro] (0,0) circle (\ra);
  \draw[magacero, line width=0.8pt] (0,0) circle (\ra);
  % signos: más densos cerca de los polos (|sigma| ~ |cos t|)
  \foreach \t in {8,22,36,50,64}{%
    \foreach \s in {1,-1}{%
      \pgfmathsetmacro{\tam}{0.35+0.65*cos(\t)}%
      \node[font=\scriptsize\bfseries, text=rojoteo, scale=\tam] at ({\s*(\ra+0.22)*sin(\t)},{(\ra+0.22)*cos(\t)}) {$+$};
      \node[font=\scriptsize\bfseries, text=pizarra, scale=\tam] at ({\s*(\ra+0.22)*sin(\t)},{-(\ra+0.22)*cos(\t)}) {$-$};}}
  % momento dipolar
  \draw[->, line width=1.6pt, magprof] (0,-0.55) -- (0,0.75) node[left=1pt, etiq, text=magprof] {$\vp$};
  % punto fuente y ángulo
  \coordinate (O) at (0,0);
  \coordinate (Zp) at (0,1);
  \coordinate (S) at ({\ra*sin(52)},{\ra*cos(52)});
  \draw[vect, rojoteo] (O) -- (S);
  \node[etiq, text=rojoteo, font=\footnotesize, anchor=west] at (1.74,0.55) {$\rp=a\,\ur$};
  \draw[rojoteo!60, line width=0.3pt] (1.76,0.55) -- (0.72,0.52);
  \pic[draw=magacero, line width=0.5pt, angle radius=0.62cm, "$\theta'$", angle eccentricity=1.35, font=\scriptsize] {angle=S--O--Zp};
  \fill[magprof] (S) circle (1.3pt) node[above right=0pt, etiq, font=\footnotesize] {$\dSp$};
\end{tikzpicture}\\[2pt]{\small\scshape (a)}
\end{minipage}\hfill
\begin{minipage}[b]{0.5\textwidth}\centering
\begin{tikzpicture}[x=0.95cm, y=0.95cm, >=Stealth]
  \def\ra{1.2}
  \begin{scope}
    \clip (-3,-2.6) rectangle (3,2.6);
    % equipotenciales (a trazos)
    \foreach \tu in {40,70}{%
      \foreach \s in {1,-1}{%
        \draw[lequi] ({-\ra*sin(\tu)},{\s*\ra*cos(\tu)}) -- ({\ra*sin(\tu)},{\s*\ra*cos(\tu)});
        \draw[lequi, domain=0:89, samples=60, variable=\t]
          plot ({\ra*sqrt(cos(\t)/cos(\tu))*sin(\t)},{\s*\ra*sqrt(cos(\t)/cos(\tu))*cos(\t)});
        \draw[lequi, domain=0:89, samples=60, variable=\t]
          plot ({-\ra*sqrt(cos(\t)/cos(\tu))*sin(\t)},{\s*\ra*sqrt(cos(\t)/cos(\tu))*cos(\t)});}}
    \draw[lequi] (-3,0) -- (3,0);
    % líneas de campo exteriores: r = a sen^2(t)/sen^2(t0), de t0 a 180-t0
    \foreach \tc in {22,38,56}{%
      \pgfmathsetmacro{\C}{\ra/(sin(\tc)^2)}%
      \foreach \s in {1,-1}{%
        \draw[lcampo, domain=\tc:180-\tc, samples=90, variable=\t]
          plot ({\s*\C*sin(\t)^3},{\C*sin(\t)^2*cos(\t)});}}
    % líneas de campo interiores: verticales, hacia -z
    \foreach \tc in {22,38,56}{%
      \foreach \s in {1,-1}{%
        \draw[lcampo] ({\s*\ra*sin(\tc)},{\ra*cos(\tc)}) -- ({\s*\ra*sin(\tc)},{-\ra*cos(\tc)});}}
    \draw[lcampo] (0,\ra) -- (0,-\ra);
    \draw[lcampoinv] (0,2.6) -- (0,\ra);
    \draw[lcampo] (0,-\ra) -- (0,-2.6);
  \end{scope}
  \draw[magprof, line width=0.8pt] (0,0) circle (\ra);
  \draw[grispizarra!60, line width=0.5pt] (-3,-2.6) rectangle (3,2.6);
  \node[etiq, font=\scriptsize, fill=white, inner sep=0.8pt] at (0.25,-0.32) {$\vE_{\mathrm{int}}$};
\end{tikzpicture}\\[2pt]{\small\scshape (b)}
\end{minipage}
\caption{Problema~2.4. (a)~Esfera de radio $a$ con densidad superficial $\sigma=\sigma_0\cos\theta$: positiva en el hemisferio norte y negativa en el sur (el tamaño de cada signo
crece con $\abs{\sigma}$), y su momento dipolar $\vp=\frac43\pi a^{3}\sigma_0\,\uz$. (b)~Líneas de campo exactas (magenta) y equipotenciales (a trazos) en un plano meridiano: fuera de
la esfera el campo es el de un dipolo puntual en el centro, $r=C\sen^{2}\theta$; dentro, uniforme y dirigido según $-\uz$. Las líneas nacen en el hemisferio positivo y mueren en
el negativo, tanto por fuera como por dentro.}
\label{fig:sem-esfera-sigma}
\end{figure}


\begin{observacion}[Interpretación: un dipolo exacto]\label{obs:sem-esfera-sigma}
El momento es el de una esfera de volumen $\frac43\pi a^{3}$ con polarización uniforme $\vP=\sigma_0\,\uz$: en efecto, por el \cref{teo:cargas-pol}, una esfera uniformemente polarizada
tiene $\rhoP=-\nabp\cdot\vP=0$ y $\sigP=\vP\cdot\ur=\sigma_0\cos\theta$, que es justamente nuestra distribución. Más aún, como $\sigma\propto\Pl{1}(\cos\theta')$, la ortogonalidad de los
polinomios de Legendre hace que todos los momentos con $\ell\neq1$ se anulen, y el potencial es \emph{exactamente} dipolar fuera de la esfera y el de un campo uniforme
dentro~\cite{GRIFFITHS}:
\[
  \phi(\vc{r})=\begin{cases}\dfrac{\sigma_0}{3\eo}\,r\cos\theta=\dfrac{\sigma_0}{3\eo}\,z,& r\le a,\\[10pt]
  \dfrac{\sigma_0a^{3}}{3\eo}\,\dfrac{\cos\theta}{r^{2}}=\dfrac{\vp\cdot\ur}{4\pi\eo r^{2}},& r\ge a,\end{cases}
  \qquad \vE_{\mathrm{int}}=-\frac{\sigma_0}{3\eo}\,\uz .
\]
Ambas expresiones coinciden en $r=a$, como exige la continuidad del potencial~\cite{TEMA1}, y la \figpanel{fig:sem-esfera-sigma}{b} muestra las líneas de campo
correspondientes: dipolares fuera y paralelas dentro.
\end{observacion}

% =========================================================================
\subsection{Problema 2.5. Distribución con simetría esférica}\label{sem:esfera-rho}
% =========================================================================
\begin{ejercicio}[Momentos de una distribución esféricamente simétrica]\label{ejer:esfera-rho}
Demostrar que una distribución de carga esférica y esféricamente simétrica da lugar a momentos dipolar y cuadrupolar nulos cuando se elige un sistema de coordenadas centrado en la
esfera.
\end{ejercicio}

\begin{resolucion}
\textbf{Idea.} En una distribución esféricamente simétrica centrada en el origen la densidad de carga solo depende de la coordenada radial, $\rho=\rho(r')$. Sea la carga confinada en una
esfera de radio $R$ (\cref{fig:sem-esfera-rho}). Expresamos el vector de posición del punto fuente y el elemento de volumen en esféricas,
\[
  \rp=r'\bigl(\sen\theta'\cos\varphi'\,\ux+\sen\theta'\sen\varphi'\,\uy+\cos\theta'\,\uz\bigr),\qquad \dVp=r'^{2}\sen\theta'\,\dif r'\,\dif\theta'\,\dif\varphi',
\]
y aprovechamos que todas las integrales se separan en un factor radial y dos angulares.

\textbf{Momento dipolar.} Separando $\vp=\int_{V'}\rp\rho(r')\,\dVp$ en sus componentes cartesianas,
\[
  p_x=\Bigl[\int_{0}^{R}r'^{3}\rho(r')\,\dif r'\Bigr]\Bigl[\int_{0}^{\pi}\sen^{2}\theta'\,\dif\theta'\Bigr]\Bigl[\int_{0}^{2\pi}\cos\varphi'\,\dif\varphi'\Bigr]=0,
\]
porque la integral azimutal es $[\sen\varphi']_{0}^{2\pi}=0$ (es nula al ser una función sinusoidal integrada en un ciclo completo). Por la misma razón,
$p_y\propto\int_{0}^{2\pi}\sen\varphi'\,\dif\varphi'=0$. En la componente $z$ es la integral \emph{polar} la que se anula:
\[
  p_z=\Bigl[\int_{0}^{R}r'^{3}\rho\,\dif r'\Bigr]\Bigl[\int_{0}^{\pi}\cos\theta'\sen\theta'\,\dif\theta'\Bigr]\Bigl[\int_{0}^{2\pi}\dif\varphi'\Bigr],
\]
\[
  \int_{0}^{\pi}\cos\theta'\sen\theta'\,\dif\theta'=\Bigl[\frac{\sen^{2}\theta'}{2}\Bigr]_{0}^{\pi}=\frac12\bigl(\sen^{2}\pi-\sen^{2}0\bigr)=0,
\]
donde hemos usado $\frac{\dif}{\dif\theta'}\bigl(\sen^{2}\theta'\bigr)=2\sen\theta'\cos\theta'$ y el teorema fundamental del cálculo. Concluimos que $\vp=\vc{0}$.

\textbf{Momento cuadrupolar.} Una distribución esféricamente simétrica tiene, en particular, simetría de revolución alrededor de $z$, así que por el \cref{teo:revolucion} (dos
ligaduras: traza nula y $Q_{xx}=Q_{yy}$)
\[
  \Qjk=0\ (j\neq k),\qquad Q_{xx}=Q_{yy}=-\tfrac12Q_{zz}.
\]
(Se comprueba también directamente: por ejemplo, $Q_{xy}=3\int r'^{2}\sen^{2}\theta'\cos\varphi'\sen\varphi'\,\rho\,\dVp\propto\int_0^{2\pi}\cos\varphi'\sen\varphi'\,\dif\varphi'=0$.) Basta,
pues, demostrar que $Q_{zz}=0$. Con $z'=r'\cos\theta'$ y $\delt{z}{z}=1$,
\[
  Q_{zz}=\int_{V'}\bigl(3z'^{2}-r'^{2}\bigr)\rho\,\dVp=\Bigl[\int_{0}^{R}r'^{4}\rho(r')\,\dif r'\Bigr]\Bigl[\int_{0}^{\pi}\bigl(3\cos^{2}\theta'-1\bigr)\sen\theta'\,\dif\theta'\Bigr]\Bigl[\int_{0}^{2\pi}\dif\varphi'\Bigr].
\]
El factor radial es finito y el azimutal vale $2\pi$; todo se decide en el factor polar. Con $u=\cos\theta'$,
\[
  \int_{0}^{\pi}\bigl(3\cos^{2}\theta'-1\bigr)\sen\theta'\,\dif\theta'=\int_{-1}^{1}\bigl(3u^{2}-1\bigr)\dif u=\bigl[u^{3}-u\bigr]_{-1}^{1}=0-0=0 .
\]
Concluimos que $Q_{zz}=0$ y, con ello, $Q_{xx}=Q_{yy}=-\frac12Q_{zz}=0$ y $\Qjk=0$ para todo $j,k$.
\end{resolucion}

% =========================================================================
% Figura del problema 2.5 (apuntes manuscritos, pág. 27): distribución con
% simetría esférica rho(r) confinada en una esfera de radio R (sombreado
% radial) y geometría del cálculo de los momentos.
% =========================================================================
\begin{figure}[ht]
\centering
\begin{tikzpicture}[x=1cm, y=1cm, >=Stealth]
  \def\Rr{1.7}
  \shade[inner color=magacero!75, outer color=rosaclaro!60] (0,0) circle (\Rr);
  \draw[magacero, line width=0.8pt] (0,0) circle (\Rr);
  \draw[magacero!60, line width=0.4pt, densely dashed] (0,0) ellipse ({\Rr} and {0.42*\Rr});
  \draw[eje] (0,0) -- (0,2.5) node[above, etiq] {$z$};
  \draw[eje] (0,0) -- (2.6,-0.35) node[right, etiq] {$y$};
  \draw[eje] (0,0) -- (-1.3,-1.25) node[below left=-2pt, etiq] {$x$};
  \coordinate (O) at (0,0);
  \coordinate (Zp) at (0,1);
  \coordinate (Q) at (-0.75,1.05);
  \coordinate (P) at (4.6,1.8);
  \draw[vect] (O) -- ($(Q)+(0.06,-0.1)$) node[pos=0.55, left=5pt, etiq, text=white] {$\rp$};
  \cubito{Q}
  \node[etiq, font=\footnotesize, anchor=south east, text=magprof] at ($(Q)+(-0.12,0.12)$) {$\dif q$};
  \draw[vect] (O) -- (P) node[pos=0.62, below right=-2pt, etiq] {$\vc{r}$};
  \draw[vectrojo] ($(Q)+(0.11,0.02)$) -- (P) node[pos=0.62, above=2pt, etiq, text=rojoteo] {$\vc{R}=\vc{r}-\rp$};
  \pic[draw=white, line width=0.5pt, angle radius=0.45cm, "$\theta'$", angle eccentricity=1.75, font=\scriptsize, text=white] {angle=Zp--O--Q};
  \fill[magprof] (P) circle (1.5pt) node[right=2pt, etiq] {$P$};
  \draw[<-, fino] (1.2,-1.2) -- (2.0,-1.75) node[right, etiq, font=\footnotesize] {$\rho=\rho(r')$, \ $r'\le R$};
\end{tikzpicture}
\caption{Problema~2.5. Distribución esféricamente simétrica, $\rho=\rho(r')$, confinada en una esfera de radio $R$ centrada en el origen (el sombreado representa una densidad
decreciente con $r'$). Al usar coordenadas esféricas para el punto fuente, todas las integrales de los momentos se separan en un factor radial y dos angulares, y son estos
últimos los que se anulan.}
\label{fig:sem-esfera-rho}
\end{figure}


\begin{observacion}[Interpretación: el teorema de Gauss, de nuevo]\label{obs:sem-esfera-rho}
El factor polar del término cuadrupolar es $\int_{-1}^{1}\Pl{2}(u)\,\dif u=0$, y lo mismo ocurre para todo $\ell\ge1$: $\int_{-1}^{1}\Pl{\ell}(u)\,\dif u=\int_{-1}^{1}\Pl{\ell}\Pl{0}\,\dif u=0$ por la
ortogonalidad de los polinomios de Legendre. Así, \emph{todos} los momentos multipolares con $\ell\ge1$ de una distribución esféricamente simétrica se anulan, y fuera de ella el
potencial es exactamente $\phi=Q/(4\pi\eo r)$: el de una carga puntual en el centro. Es el mismo resultado que obtuvimos con la ley de Gauss en el Tema~1~\cite{TEMA1}, ahora
deducido término a término.
\end{observacion}

% =========================================================================
\subsection{Resumen de resultados}\label{sem:resumen}
% =========================================================================
\begin{table}[ht]
\centering
\caption{Resultados de los seminarios del Tema~2.}
\label{tab:resultados}
\renewcommand{\arraystretch}{1.75}
\small
\begin{tabular}{@{}>{\raggedright\arraybackslash}p{0.2\textwidth}>{\raggedright\arraybackslash}p{0.76\textwidth}@{}}
  \toprule
  \textsc{Problema} & \textsc{Resultado}\\
  \midrule
  2.1 Barra ($r>10\ell$) & $\phi\simeq\dfrac{\lambda\ell}{4\pi\eo r}+\dfrac{\lambda\ell^{3}(3\cos^{2}\theta-1)}{96\pi\eo r^{3}}$\\
                         & $\vE\simeq\Bigl[\dfrac{\lambda\ell}{4\pi\eo r^{2}}+\dfrac{\lambda\ell^{3}(3\cos^{2}\theta-1)}{32\pi\eo r^{4}}\Bigr]\ur+\dfrac{\lambda\ell^{3}\sen\theta\cos\theta}{16\pi\eo r^{4}}\,\uth$\\
  2.2 Carga en $(a,b,c)$ & $Q=q$, \ $\vp=q(a\ux+b\uy+c\uz)$, \ $Q_{xx}=q(2a^{2}-b^{2}-c^{2})$ (y cíclicos), \ $Q_{xy}=3qab$, $Q_{xz}=3qac$, $Q_{yz}=3qbc$; con $-q$ en el origen solo cambia $Q\to0$\\
  2.3 Paralelepípedo     & $\phi\simeq\dfrac{Q}{4\pi\eo r}+\dfrac{Q}{24\pi\eo r^{5}}\bigl[x^{2}(2a^{2}-b^{2}-c^{2})+y^{2}(2b^{2}-a^{2}-c^{2})+z^{2}(2c^{2}-a^{2}-b^{2})\bigr]$\\
  2.4 Esfera $\sigma_0\cos\theta$ & $\vp=\frac43\pi a^{3}\sigma_0\,\uz$ (independiente del origen, porque $Q=0$)\\
  2.5 Simetría esférica  & $\vp=\vc{0}$, \ $\Qjk=0$; en general, todos los momentos con $\ell\ge1$ son nulos\\
  \bottomrule
\end{tabular}
\end{table}

% =========================================================================
\clearpage
\appendix
\section{Formulario del tema y herramientas auxiliares}\label{app:formulario}
% =========================================================================
Recogemos aquí, en el orden en que aparecen en el formulario del curso~\cite{FORMULARIO}, las expresiones del Tema~2, junto con las propiedades de los polinomios de
Legendre y las integrales angulares que se han usado repetidamente.

% -------------------------------------------------------------------------
\subsection{Formulario del Tema 2: multipolos eléctricos}\label{app:tabla}
% -------------------------------------------------------------------------
\begin{table}[ht]
\centering
\caption{Formulario del Tema~2. Entre paréntesis, el resultado del texto donde se obtiene cada expresión.}
\label{tab:formulario}
\renewcommand{\arraystretch}{2.05}
\small
\begin{tabular}{@{}>{\centering\arraybackslash}p{0.5\textwidth}>{\raggedright\arraybackslash}p{0.43\textwidth}@{}}
  \toprule
  \textsc{Expresión} & \textsc{Significado}\\
  \midrule
  $\displaystyle\vp=\lim_{q\to\infty,\ \ell\to0}q\,\vl$ & Dipolo eléctrico puntual (\cref{def:dipolo-puntual})\\
  $\displaystyle\phi(\vc{r})=\frac{\vp\cdot\ur}{4\pi\eo r^{2}}$,\quad $\displaystyle\vE(\vc{r})=\frac{p}{4\pi\eo r^{3}}\bigl(2\cos\theta\,\ur+\sen\theta\,\uth\bigr)$\newline
  $\displaystyle\vE(\vc{r})=\frac{1}{4\pi\eo r^{3}}\bigl[3(\vp\cdot\ur)\,\ur-\vp\bigr]$
     & Potencial y campo creados por un dipolo (\cref{prop:pot-dipolo,teo:campo-dipolo})\\
  $\displaystyle\vP(\rp)=\frac{\dif\vp(\rp)}{\dVp}$ & Polarización, densidad de momento dipolar (\cref{def:polarizacion})\\
  $\rhoP(\rp)=-\nabp\cdot\vP(\rp)$,\quad $\sigP(\rp)=\vP(\rp)\cdot\un$ & Densidades de carga de polarización (\cref{teo:cargas-pol})\\
  $\displaystyle\phi(\vc{r})=\frac{1}{4\pi\eo}\sum_{\ell=0}^{\infty}\frac{1}{r^{\ell+1}}\int_{V'}r'^{\,\ell}\rho(\rp)\,\Pl{\ell}(\cosp)\,\dVp$ & Desarrollo multipolar del potencial (\cref{teo:multipolar})\\
  $\displaystyle\phimul{0}(\vc{r})=\frac{Q}{4\pi\eo r}$,\quad $\displaystyle Q=\int_{V'}\rho(\rp)\,\dVp$ & Término monopolar, $\ell=0$ (\cref{prop:monopolar})\\
  $\displaystyle\phimul{1}(\vc{r})=\frac{\vp\cdot\ur}{4\pi\eo r^{2}}$,\quad $\displaystyle\vp=\int_{V'}\rp\rho(\rp)\,\dVp$ & Término dipolar, $\ell=1$ (\cref{prop:dipolar})\\
  $\vp_{O'}=\vp_{O}-Q\,\OOp$ & Efecto de la elección del origen de coordenadas en el momento dipolar (\cref{prop:p-origen})\\
  $\displaystyle\phimul{2}(\vc{r})=\frac{1}{8\pi\eo r^{3}}\sum_{j,k}\Qjk\,\cd{j}\cd{k}$,\quad $\displaystyle\Qjk=\int_{V'}\rho(\rp)\bigl(3x'_jx'_k-\delt{j}{k}r'^{2}\bigr)\dVp$
     & Término cuadrupolar, $\ell=2$ (\cref{prop:cuadrupolar})\\
  $\displaystyle\phimul{2}(\vc{r})=\frac{Q_{zz}}{16\pi\eo r^{3}}\bigl(3\cos^{2}\theta-1\bigr)$ & Término cuadrupolar para sistemas con simetría de revolución (\cref{teo:revolucion})\\
  $\displaystyle\frac{\phimul{n}(\vc{r})}{\phimul{0}(\vc{r})}<\Bigl(\frac{a}{r}\Bigr)^{n}$ & Estimación del error al truncar el desarrollo a orden $n$, en sistemas con carga neta no nula y toda del mismo signo (\cref{teo:cota-error})\\
  \bottomrule
\end{tabular}
\end{table}

% -------------------------------------------------------------------------
\subsection{Polinomios de Legendre}\label{app:legendre}
% -------------------------------------------------------------------------
Los polinomios de Legendre son las soluciones polinómicas de la ecuación de Legendre, $(1-x^{2})y''-2xy'+\ell(\ell+1)y=0$, que aparece al separar variables en la ecuación de
Laplace en coordenadas esféricas~\cite{ARFKEN,JACKSON}; de ahí su papel central en la electrostática. Además de la función generatriz y la fórmula de Rodrigues
(\cref{def:legendre}), usamos estas propiedades:
\begin{enumerate}[label=(\alph*)]
  \item \emph{Paridad}: $\Pl{\ell}(-x)=(-1)^{\ell}\Pl{\ell}(x)$ (se sigue de $g(-x,-t)=g(x,t)$).
  \item \emph{Valores extremos y acotación}: $\Pl{\ell}(1)=1$, $\Pl{\ell}(-1)=(-1)^{\ell}$ y $\abs{\Pl{\ell}(x)}\le1$ en $[-1,1]$ (\cref{lem:legendre-cota}).
  \item \emph{Ortogonalidad}: $\displaystyle\int_{-1}^{1}\Pl{\ell}(x)\,\Pl{m}(x)\,\dif x=\frac{2}{2\ell+1}\,\delt{\ell}{m}$. En particular, $\int_{-1}^{1}\Pl{\ell}(x)\,\dif x=0$ si $\ell\ge1$.
  \item \emph{Recurrencia de Bonnet}: $(\ell+1)\Pl{\ell+1}(x)=(2\ell+1)\,x\,\Pl{\ell}(x)-\ell\,\Pl{\ell-1}(x)$, que permite generar todos a partir de $\Pl{0}=1$ y $\Pl{1}=x$.
\end{enumerate}

% =========================================================================
% Figura del apéndice: polinomios de Legendre P_0 ... P_4 en [-1,1]
% (analíticos, pgfplots), con la banda |P_l| <= 1 sombreada.
% =========================================================================
\begin{figure}[ht]
\centering
\begin{tikzpicture}
\begin{axis}[tema, width=10.5cm, height=6cm, xmin=-1, xmax=1, ymin=-1.08, ymax=1.08,
             xlabel={$x$}, ylabel={$\Pl{\ell}(x)$}, domain=-1:1, samples=161,
             every axis plot/.style={line width=1.1pt},
             xtick={-1,-0.5,0,0.5,1}, ytick={-1,-0.5,0,0.5,1},
             legend pos=outer north east, legend cell align=left,
             legend style={font=\scriptsize, draw=none}]
  \addplot[magacero]                 {1};
  \addplot[rojoclaro]                {x};
  \addplot[pizarra]                  {0.5*(3*x^2-1)};
  \addplot[verdeteo]                 {0.5*(5*x^3-3*x)};
  \addplot[magprof, densely dashed]  {(35*x^4-30*x^2+3)/8};
  \legend{$\Pl{0}$, $\Pl{1}$, $\Pl{2}$, $\Pl{3}$, $\Pl{4}$}
\end{axis}
\end{tikzpicture}
\caption{Los cinco primeros polinomios de Legendre en $[-1,1]$. Todos valen $1$ en $x=1$, tienen la paridad de $\ell$, $\ell$ ceros en $(-1,1)$ y están acotados por $1$ en valor
absoluto (\cref{lem:legendre-cota}). El cero de $\Pl{2}$ en $x=1/\sqrt3$ define el cono nodal del término cuadrupolar.}
\label{fig:legendre-pol}
\end{figure}


% -------------------------------------------------------------------------
\subsection{Integrales angulares de uso frecuente}\label{app:integrales}
% -------------------------------------------------------------------------
En coordenadas esféricas, $\dVp=r'^{2}\sen\theta'\,\dif r'\,\dif\theta'\,\dif\varphi'$ y, sobre una esfera de radio $a$, $\dSp=a^{2}\sen\theta'\,\dif\theta'\,\dif\varphi'$. Las integrales azimutales
que anulan los momentos de las distribuciones con simetría de revolución son
\[
  \int_{0}^{2\pi}\cos\varphi'\,\dif\varphi'=\int_{0}^{2\pi}\sen\varphi'\,\dif\varphi'=\int_{0}^{2\pi}\sen\varphi'\cos\varphi'\,\dif\varphi'=0,
\]
\[
  \int_{0}^{2\pi}\cos^{2}\varphi'\,\dif\varphi'=\int_{0}^{2\pi}\sen^{2}\varphi'\,\dif\varphi'=\pi ,
\]
y las polares, con el cambio $u=\cos\theta'$, $\dif u=-\sen\theta'\,\dif\theta'$,
\[
  \int_{0}^{\pi}\sen\theta'\,\dif\theta'=2,\quad
  \int_{0}^{\pi}\cos\theta'\sen\theta'\,\dif\theta'=0,\quad
  \int_{0}^{\pi}\cos^{2}\theta'\sen\theta'\,\dif\theta'=\frac23,
\]
\[
  \int_{0}^{\pi}\bigl(3\cos^{2}\theta'-1\bigr)\sen\theta'\,\dif\theta'=0 .
\]
Por último, para el cálculo de momentos de distribuciones superficiales y lineales mediante la delta de Dirac~\cite{TEMA0,TEMA1},
$\int_{0}^{\infty}f(r')\,\delta(r'-a)\,\dif r'=f(a)$ para $a>0$.


% =========================================================================
%                              BIBLIOGRAFÍA
% =========================================================================
\printbibliography[title={Referencias}]

\end{document}
